\documentclass[10pt,oneside]{book}

\usepackage{fontspec}
\setmainfont{Malgun Gothic}[
  ItalicFont={Malgun Gothic},
  ItalicFeatures={FakeSlant=0.2},
  BoldItalicFont={Malgun Gothic Bold},
  BoldItalicFeatures={FakeSlant=0.2}
]
\setsansfont{Malgun Gothic}
\XeTeXlinebreaklocale "ko"
\XeTeXlinebreakskip=0pt plus 1pt
\usepackage{amsmath,amssymb,mathrsfs}
\usepackage{longtable}
\usepackage[all]{xy}
\usepackage{tikz-cd}
\usepackage[a4paper,margin=22mm]{geometry}
\usepackage[
  colorlinks=true,
  linkcolor=blue,
  hypertexnames=false,
  unicode=true,
  pdfencoding=auto
]{hyperref}
\hypersetup{
  pdftitle={대수기하학 원론 -- 한국어 누적판},
  pdfauthor={알렉산더 그로텐디크와 장 디외도네},
  pdfsubject={NUMDAM 프랑스어 권위본에 따른 EGA 한국어 누적 번역}
}

\renewcommand{\thesection}{\arabic{section}}
\renewcommand{\thesubsection}{\thesection.\arabic{subsection}}
\renewcommand{\bibname}{참고문헌}
\renewcommand{\contentsname}{목차}
\makeatletter
\renewcommand{\@pnumwidth}{2em}
\renewcommand{\@tocrmarg}{3em}
\makeatother
% Preserve the historical printed-page rhythm in the cumulative reader.  The
% first marker anchors the opening page; every later marker forces the line
% carrying it onto a fresh output page without splitting the source paragraph.
\newif\ifagoldpageseen
\agoldpageseenfalse
\newcommand{\oldpage}[2][]{%
  \ifagoldpageseen
    \ifhmode
      \vadjust pre{\penalty -10000}%
    \else
      \ifvmode\penalty -10000\fi
    \fi
  \else
    \global\agoldpageseentrue
  \fi
  \setcounter{footnote}{0}%
  \marginpar{\scriptsize\textbf{#1}~|~#2}\ignorespaces}
\newcommand{\agref}[2]{%
  \ifcsname r@#1\endcsname\hyperref[#1]{#2}\else#2\fi}
% Compatibility forms retained by the source-synchronized EGA II indexes and
% addenda.  They are defined centrally so the editable source stays legible.
\providecommand{\sh}[1]{\mathcal{#1}}
\providecommand{\Spec}{\operatorname{Spec}}
\providecommand{\Proj}{\operatorname{Proj}}
\providecommand{\Hom}{\operatorname{Hom}}
\providecommand{\Hyp}{\operatorname{Hyp}}
\providecommand{\vphi}{\varphi}
\providecommand{\kres}{\kappa}
\providecommand{\isoto}{\xrightarrow{\sim}}
\providecommand{\supertilde}{^{\sim}}
\newenvironment{env}[1][]{\par\medskip\noindent\textbf{(#1)}\ }{\par}
\newenvironment{definition}[1][]{%
  \par\medskip\noindent\itshape 정의 (#1). ---\ }{\par}
\newenvironment{proposition}[1][]{%
  \par\medskip\noindent\itshape 명제 (#1). ---\ }{\par}
\newenvironment{theorem}[1][]{%
  \par\medskip\noindent\itshape 정리 (#1). ---\ }{\par}
\newenvironment{corollary}[1][]{%
  \par\medskip\noindent\itshape 따름정리 (#1). ---\ }{\par}
\newenvironment{remark}[1][]{%
  \par\medskip\noindent\itshape 주 (#1). ---\ }{\par}
\newenvironment{remarks}[1][]{%
  \par\medskip\noindent\itshape 주들 (#1). ---\ }{\par}
\newenvironment{lemma}[1][]{%
  \par\medskip\noindent\itshape 보조정리 (#1). ---\ }{\par}
\newenvironment{example}[1][]{%
  \par\medskip\noindent\textit{예 (#1). ---\ }}{\par}
\newenvironment{examples}[1][]{%
  \par\medskip\noindent\textit{예들 (#1). ---\ }}{\par}
\newenvironment{notation}[1][]{%
  \par\medskip\noindent\textit{표기 (#1). ---\ }}{\par}
\newenvironment{proof}{%
  \par\medskip\noindent\textit{증명. ---\ }}{\par}
\newenvironment{namedtheorem}[2][]{%
  \par\medskip\noindent\itshape #2 (#1). ---\ }{\par}
\newenvironment{namedlemma}[2][]{%
  \par\medskip\noindent\itshape #2 (#1). ---\ }{\par}

\begin{document}

% ===== BEGIN EXACT INLINE INPUT: front.tex =====
% Authority: EGA I frontmatter-fr.tex, SHA-256 D506B87684E2136E3F87495190EECDF40B79DB8887ED71093C4CD56648E282A9.
% Give the unnumbered title leaf a distinct logical PDF page label; the
% Arabic sequence begins with the following edition-information page.
\pagenumbering{roman}
\begin{titlepage}
\centering

{\Large 프랑스 고등과학연구소\par}
\vspace{0.4em}
{\large INSTITUT DES HAUTES ÉTUDES SCIENTIFIQUES\par}

% 원 인쇄본에서는 이 자리에 연구소 문장이 놓여 있다.
\vfill

{\Large\bfseries 대수기하학 원론\par}
\vspace{1.5em}
{\large\itshape A. 그로텐디크 지음\par}
\vspace{0.7em}
{\itshape J. 디외도네의 협력으로 편집\par}

\vspace{2em}
I

\vspace{1.5em}
{\Large\bfseries 스킴의 언어\par}

\vspace{1.6em}
{\large 한국어 누적판\par}
\vspace{0.5em}
{\small 현재 작업 범위: EGA $0_{\mathrm I}$, EGA I, EGA II 및 EGA III-1 완역\par}

\vfill
{\Large 1960\par}
\vspace{1.2em}
{\Large PUBLICATIONS MATHÉMATIQUES, 제4호\par}
\vspace{0.5em}
{\large 5, ROND-POINT BUGEAUD --- PARIS (XVI\textsuperscript{e})\par}
\end{titlepage}
\pagenumbering{arabic}

\thispagestyle{empty}
\begin{center}
{\Large\bfseries 한국어판 정보\par}
\end{center}

\vspace{1.5em}
이 책은 알렉산더 그로텐디크와 장 디외도네의
\emph{Éléments de géométrie algébrique}를 한국어로 옮기는 누적판이다.
이 판의 맨 앞 독자용 산출물인
\texttt{00\_EGA\_ko\_CUMULATIVE\_READER.pdf}는 현재까지 완성된
한국어 번역을 빠짐없이 담는다. 수록 범위는 EGA $0_{\mathrm I}$,
EGA I와 EGA II의 프로그램, 본문, 서지 및 부록 전체와, EGA III의
첫째 분책인 EGA III-1 전체이다. EGA III-1에는 제0장의 제8--13절,
제III장의 제1--5절, 참고문헌과 원 목차가 각 정본 파일의 끝까지
들어 있다. 이 판은 이 네 범위의 완역을 주장하지만, 뒤이은 EGA III
분책이나 EGA 전집 전체의 완성을 주장하지 않는다.
EGA 전체의 한국어 번역은 같은 DOI 계열과 공공 보존소들에서 끊김 없이 계속된다.
이 교정판은 수록 범위를 조밀한 연속 조판으로 압축하지 않고, 확인된
581개의 원 인쇄면 경계를 출력 쪽 리듬으로 보존한다. 판 정보와 목차를 포함한 독자용 PDF의
총 쪽수와 최종 바이트 식별자는 각 판의 결정적 빌드 영수증에 기록한다.

\medskip
\noindent\textbf{한국어판 안정 DOI:}
\href{https://doi.org/10.5281/zenodo.21921513}{10.5281/zenodo.21921513}

\smallskip
\noindent\textbf{이 판의 정확 DOI:}
\href{https://doi.org/10.5281/zenodo.22864049}{10.5281/zenodo.22864049}

\smallskip
\noindent\textbf{GitHub 저장소:}
\href{https://github.com/KokunoYumeto/ega-ko}{github.com/KokunoYumeto/ega-ko}

\smallskip
\noindent\textbf{전역 EGA 자료실:}
\href{https://doi.org/10.5281/zenodo.20414353}{10.5281/zenodo.20414353}

\medskip
한국어 번역과 조판의 프로젝트 기여자는
\emph{AI typesetting \& translation}이다.
편집 가능한 TeX, 프랑스어 원문의 정확한 식별자와 해시, 모든 번역\,·\,용어\,·\,조판 결정,
난점과 미해결 항목, 빌드 기록, 추출 검사, 렌더 검사 및 SHA-256 목록은 이 판의
출처 및 증거 묶음에 들어 있다.

\medskip
CC BY 4.0은 권리를 보유한 범위의 한국어 번역, 한국어 조판,
프로젝트 작성 메타데이터\,·\,색인\,·\,결정\,·\,QA 증거에만 적용된다.
바탕 수학 저작, 역사적 프랑스어 판, 제3자 서지 자료는 각자의 권리와 귀속을 유지한다.
역사적 저자, NUMDAM과 원출판사는 이 독립 유지 번역을 보증하지 않는다.

\medskip
이 판은 원문 동기화와 모델 기반 검사를 거친
누적 작업판이며, 비평판 또는 사람의 외부 인증을 주장하지 않는다.

\vfill
\noindent\textbf{언어:} 한국어 (ko; Zenodo: kor)\par
\noindent\textbf{작업 식별자:} 2026-09-20-ega3-iii1-complete\par
\clearpage

\thispagestyle{empty}
\vspace*{\fill}
\begin{center}
\fbox{%
  \begin{minipage}{0.50\textwidth}
  \centering
  납본

  \smallskip
  초판 \dotfill 1960년 제3사분기

  \smallskip
  모든 권리

  모든 국가에서 유보됨

  \smallskip
  \copyright\ 1960, \emph{Institut des Hautes Études Scientifiques}
  \end{minipage}%
}
\end{center}
\vspace*{\fill}
\clearpage
% ===== END EXACT INLINE INPUT: front.tex =====

\tableofcontents
\clearpage
\markboth{}{}

% ===== BEGIN EXACT INLINE INPUT: intro.tex =====
% Authority: EGA I intro-fr.tex, SHA-256 CE78400CD9DBC36A4D11CC933B7BE18BEAC0C69AE6A04FBA3ECFA86053572980.
\section*{서론}

\begin{center}
\emph{오스카 자리스키와 앙드레 베유에게.}
\end{center}

\oldpage[I]{5}

이 논문과 뒤이어 나올 많은 논문은 대수기하학의 기초에 관한 하나의 논고를 이루기 위한 것이다.
원칙적으로 이 분야에 관한 어떤 특별한 지식도 전제하지 않는다. 오히려 그러한 지식은 분명한 장점이 있음에도 불구하고, 그것이 수반하는 쌍유리적 관점에 지나치게 익숙해진 탓에 여기서 제시하는 관점과 기법에 익숙해지려는 이에게 때로 방해가 될 수 있음이 드러났다.
반면 독자가 다음 주제들을 잘 알고 있다고 가정한다.

\begin{enumerate}
  \item[a)] \emph{가환대수학}. 예를 들어 현재 준비 중인 N.~부르바키의 \emph{원론} 여러 권에서 전개되는 내용이다(이 책들이 출간되기 전까지는 사뮈엘-자리스키~\cite{I-13}와 사뮈엘~\cite{I-11},~\cite{I-12}을 참조할 수 있다).
  \item[b)] \emph{호몰로지 대수학}. 카르탕-아일렌베르크~\cite{I-2}(이하 (M)), 고드망~\cite{I-4}(이하 (G)), 그리고 A.~그로텐디크의 최근 논문~\cite{I-6}(이하 (T))을 참조한다.
  \item[c)] \emph{층 이론}. 주된 참고문헌은 (G)와 (T)이다. 이 이론은 가환대수학의 핵심 개념들을 ``기하학적'' 용어로 해석하고 ``전역화''하는 데 필수적인 언어를 제공한다.
  \item[d)] 끝으로 독자가 \emph{함자 언어}에 어느 정도 익숙하면 유용할 것이다. 이 논고에서는 이 언어를 계속 사용한다. (M), (G), 특히 (T)를 참조할 수 있으며, 그 원리와 일반 함자론의 주요 결과들은 이 논고의 저자들이 준비 중인 별도의 저서에서 더 자세히 설명할 예정이다.
\end{enumerate}

\bigskip
\begin{center}*\quad*\quad*\end{center}
\bigskip

이 서론은 대수기하학에서 ``스킴''이라는 관점이 무엇인지 개략적으로 설명하거나, 이 관점을 채택해야 했던 수많은 이유를 열거할 자리가 아니다. 특히 우리가 다루는 ``다양체''의 국소환에서 멱영원을 체계적으로 받아들여야 하는 이유도 여기서는 설명하지 않는다. 이 수용은 필연적으로 유리 사상의 개념을 뒤로 물리고, 정칙 사상 또는 ``사상''의 개념을 앞세운다.
이 논고의 목적은 바로 ``스킴''의 언어를 체계적으로 전개하고, 바라건대 그 필요성을 입증하는 데 있다.
그렇게 하기는 어렵지 않겠지만, 제1장에서 전개하는 개념들에 대한 ``직관적'' 입문도 여기서는 시도하지 않겠다.

\oldpage[I]{6}
이 논고의 내용을 미리 훑어보고 싶은 독자는 A.~그로텐디크가 1958년 에든버러 국제수학자대회에서 한 강연~\cite{I-7}과 같은 저자의 발표~\cite{I-8}을 참조할 수 있다.
J.-P.~세르의 논문~\cite{I-14}(이하 (FAC))도 대수기하학의 고전적 관점과 스킴의 관점 사이를 잇는 설명으로 볼 수 있다. 따라서 이 논문을 읽는 것은 우리의 \emph{원론}을 읽기 위한 훌륭한 준비가 될 수 있다.

\bigskip
\begin{center}*\quad*\quad*\end{center}
\bigskip

참고로 이 논고에서 예정한 전체 구성을 아래에 제시한다. 특히 뒤쪽 장들의 구성은 나중에 바뀔 수 있다.

\medskip
\begin{tabular}{rl}
  제1장. & --- 스킴의 언어.\\
  --- \textbf{II}. & --- 몇몇 사상 부류에 대한 초보적 전역 연구.\\
  --- \textbf{III}. & --- 연접 대수적 층의 코호몰로지. 응용.\\
  --- \textbf{IV}. & --- 사상의 국소적 연구.\\
  --- \textbf{V}. & --- 스킴의 초보적 구성법.\\
  --- \textbf{VI}. & --- 강하 기법. 스킴 구성의 일반적 방법.\\
  --- \textbf{VII}. & --- 군 스킴, 주다발 공간.\\
  --- \textbf{VIII}. & --- 다발 공간의 미분론적 연구.\\
  --- \textbf{IX}. & --- 기본군.\\
  --- \textbf{X}. & --- 유수와 쌍대성.\\
  --- \textbf{XI}. & --- 교차이론, 천 특성류, 리만-로흐 정리.\\
  --- \textbf{XII}. & --- 아벨 스킴과 피카르 스킴.\\
  --- \textbf{XIII}. & --- 베유 코호몰로지.\\
\end{tabular}
\medskip

원칙적으로 모든 장은 열려 있는 것으로 간주하며, 언제든지 절을 덧붙일 수 있다. 채택한 출판 방식의 불편을 줄이기 위해 그러한 절은 별도의 분책으로 출간한다.
한 장이 출간될 때 어떤 절이 예정되어 있거나 준비 중이면, 긴급성에 따른 순서 때문에 실제 출간이 훨씬 뒤로 미뤄지더라도 그 장의 목차에 이를 적어 둔다.
독자의 편의를 위해 이 논고의 여러 장에서 사용하는 가환대수학, 호몰로지 대수학, 층 이론의 여러 보충 결과를 ``제0장''에 모았다. 이 결과들은 어느 정도 알려져 있지만 편리한 참고문헌을 제시할 수 없었던 것들이다.
독자는 논고 본문을 읽다가 우리가 인용하는 결과가 충분히 익숙하지 않을 때에만 제0장을 참조하기를 권한다.

\oldpage[I]{7}
이렇게 읽으면 이 논고는 초심자가 가환대수학과 호몰로지 대수학에 익숙해지는 좋은 방법이 될 수 있다고 생각한다. 구체적인 응용을 동반하지 않는 이 두 분야의 공부는 적지 않은 이들에게 따분하고 심지어 침울하기까지 한 것으로 여겨지기 때문이다.

\bigskip
\begin{center}*\quad*\quad*\end{center}
\bigskip

이 서론에서 여기 제시하는 개념과 결과의 역사를, 아무리 간략하게라도 개관하는 일은 우리의 능력 밖이다.
본문에는 이해에 특히 유용하다고 판단한 참고문헌만 싣고, 가장 중요한 결과들에 대해서만 그 기원을 밝힐 것이다.
적어도 형식적으로는 이 저서가 다루는 주제들이 상당히 새롭다. 그래서 우리가 그 연구를 전해 들었을 뿐인 19세기와 20세기 초 대수기하학의 선구자들을 거의 인용하지 않는 것이다.
그러나 저자들에게 가장 직접적인 영향을 주고 스킴 관점의 발전에 기여한 저작들에 관해서는 여기서 몇 마디 할 필요가 있다.
무엇보다 먼저 J.-P.~세르의 기초적인 논문 (FAC)을 들 수 있다. 이 논문은 A.~베유의 고전적인 \emph{Foundations}~\cite{I-18}가 너무 건조하다고 느낀 여러 젊은 입문자들(그 가운데에는 이 논고의 저자 한 명도 있다)에게 대수기하학의 입문서 역할을 했다.
여기서 처음으로 ``추상'' 대수다양체의 ``자리스키 위상''이 대수적 위상수학의 특정 기법들을 적용하기에 완전히 적합하며, 특히 코호몰로지 이론을 낳는다는 사실이 증명되었다.
더 나아가 그곳에서 주어진 대수다양체의 정의는 우리가 여기서 전개하는 개념의 확장에 가장 자연스럽게 맞는다\footnote{J.-P.~세르가 지적했듯이, 환의 층을 주어 다양체의 구조를 정의한다는 발상은 H.~카르탕에게서 비롯되었으며, 카르탕은 이를 해석공간 이론의 출발점으로 삼았다.
물론 대수기하학에서와 마찬가지로 ``해석기하학''에서도 해석공간의 국소환 안에 있는 멱영원에 정당한 지위를 부여해야 한다.
H.~카르탕과 J.-P.~세르의 정의를 이렇게 확장하는 문제는 최근 H.~그라우어트~\cite{I-5}가 다루기 시작했으며, 이 일반적 틀 안에서 해석기하학을 체계적으로 설명하는 저서가 머지않아 나오기를 기대할 수 있다.
이 논고에서 전개하는 개념과 기법이 해석기하학에서도 의미를 유지한다는 점은 분명하다. 다만 후자의 이론에서는 기술적 어려움이 훨씬 더 클 것으로 예상해야 한다.
방법이 단순한 대수기하학은 앞으로 해석공간 이론이 발전하는 데 일종의 형식적 모형이 될 수 있을 것이다.}.
또한 세르는 체 위의 아핀 대수 대신 임의의 가환환을 쓰면 아핀 대수다양체의 코호몰로지 이론을 어려움 없이 옮겨 적을 수 있음을 이미 알아차렸다.
따라서 이 논고의 제I장과 제II장, 그리고 제III장의 첫 두 절은 본질적으로 (FAC)와 같은 저자의 뒤이은 논문~\cite{I-15}의 주요 결과를 이 확장된 틀로 곧바로 옮긴 것으로 볼 수 있다.
우리는 C.~슈발레의 \emph{대수기하학 세미나}~\cite{I-1}에서도 큰 도움을 얻었다. 특히 그가 도입한 ``구성가능 집합''을 체계적으로 사용하는 것이 스킴 이론에서 매우 유용하다는 사실이 드러났다(제IV장 참조).
또한 차원의 관점에서 사상을 연구하는 방법(제IV장)도 그에게서 가져왔다. 이 방법은 스킴의 틀로 거의 아무 변화 없이 옮겨진다.

\oldpage[I]{8}
한편 슈발레가 도입한 ``국소환의 스킴''이라는 개념은 대수기하학을 자연스럽게 확장할 수 있게 한다. 다만 우리가 여기서 부여하려는 유연성과 일반성을 모두 갖추지는 못한다. 이 개념과 우리 이론의 관계는 제1장 \S~8을 보라.
M.~나가타는 데데킨트 환 위의 대수기하학에 관한 많은 특수한 결과를 담은 일련의 논문~\cite{I-9}에서 이러한 확장을 전개하였다\footnote{대수기하학에 관한 우리의 관점에 가까운 연구로는 E.~켈러의 중요한 저작~\cite{I-22}과 저우와 이구사의 최근 노트~\cite{I-3}을 들 수 있다. 이들은 나가타-슈발레 이론의 틀 안에서 (FAC)의 몇몇 결과를 다시 다루고 퀸네트 공식도 제시한다.}.

\bigskip
\begin{center}*\quad*\quad*\end{center}
\bigskip

끝으로 대수기하학에 관한 책, 특히 그 기초에 관한 책이 O.~자리스키와 A.~베유 같은 수학자들에게서, 설령 다른 사람들을 거쳐서라도, 필연적으로 영향을 받는다는 것은 말할 필요도 없다.
특히 자리스키의 \emph{정칙함수론}~\cite{I-20}은 코호몰로지 방법으로 적절히 유연하게 만들고 존재정리(제III장 \S\S~4, 5)로 보완하면, 제VI장에서 설명하는 강하 기법과 더불어 이 논고에서 사용하는 주요 도구 가운데 하나가 된다. 우리가 보기에 이는 대수기하학에서 쓸 수 있는 가장 강력한 도구 가운데 하나이다.

이 이론이 속하는 일반적 기법은 다음과 같이 개략적으로 설명할 수 있다. 전형적인 예는 제IX장에서 기본군을 연구할 때 나타난다.
한 대수다양체에서 다른 대수다양체로 가는 고유 사상(제II장) $f:X\to Y$가 있고(더 일반적으로는 한 스킴에서 다른 스킴으로 가는 사상), 점 $y\in Y$의 근방에 관한 문제 P를 풀기 위해 $y$의 근방에서 이 사상을 연구하려 한다.
다음 단계를 차례로 밟는다.
\begin{enumerate}
  \item[$1^\circ$] $Y$가 아핀이라고 가정할 수 있다. 그러면 $X$는 $Y$의 아핀환 $A$ 위에서 정의된 스킴이 되며, 나아가 $A$를 $y$의 국소환으로 바꿀 수 있다.
  이 축약은 실제로 언제나 쉽고(제V장), $A$가 \emph{국소환}인 경우로 문제를 돌린다.
  \item[$2^\circ$] $A$가 \emph{아르틴 국소환}인 경우에 문제를 연구한다.
  $A$가 정역이라고 가정하지 않을 때에도 문제가 실제로 의미를 유지하게 하려면 문제 P를 다시 정식화해야 할 때가 있다. 이렇게 하면 흔히 이 단계에서 ``무한소적'' 성격을 띠는 문제를 더 잘 이해하게 된다.
  \item[$3^\circ$] 형식 스킴 이론(제III장 \S\S~3, 4, 5)을 이용하면 아르틴 환의 경우에서 \emph{완비 국소환}의 경우로 넘어갈 수 있다.
  \item[$4^\circ$] 끝으로 $A$가 임의의 국소환이면, $X$ 위의 적당한 스킴에서 주어진 ``형식적'' 절단에 가까운 ``다중 절단''을 고려함으로써(제IV장), $A$의 완비화로 스칼라를 확대하여 $X$에서 얻은 스킴에 대해 알려진 결과를 $A$의 충분히 단순한 유한 확대(예를 들면 비분기 확대)에 대한 유사한 결과로 흔히 옮길 수 있다.
\end{enumerate}

이 개요는 아르틴 환 $A$ 위에서 정의된 스킴을 체계적으로 연구하는 일이 중요함을 보여 준다.
세르가 국소 유체론을 정식화할 때 취한 관점과 그린버그의 최근 연구는, 그러한 스킴 $X$에 $A$의 잉여체 $k$(완전체라고 가정한다) 위의 스킴 $X'$를 함자적으로 대응시키는 방식으로 이 연구를 수행할 수 있음을 시사하는 듯하다. 여기서 $X'$의 차원은 좋은 경우 $n\dim X$이고, $n$은 $A$의 길이이다.

A.~베유의 영향에 관해서는 다음과 같이 말하는 것으로 충분하다. ``베유 코호몰로지''의 정의를 필요한 일반성 전체를 갖추어 정식화하고, 디오판토스 기하학의 유명한 추측들~\cite{I-19}을 확립하는 데 필요한 모든 형식적 성질의 증명\footnote{오해를 피하기 위해 덧붙이면, 이 서론을 쓰는 시점에는 이 과제가 이제 막 시작되었을 뿐이며, 따라서 아직 베유 추측의 증명에 이르지 못했다.}에 착수하려면 필요한 도구를 개발해야 한다. 이것이 이 논고를 집필하게 된 주된 동기 가운데 하나였다. 대수기하학에서 통상 사용하는 개념과 방법의 자연스러운 틀을 찾고, 저자들 자신이 그 개념과 기법을 이해할 기회를 얻으려는 바람도 똑같이 중요한 동기였다.

\bigskip
\begin{center}*\quad*\quad*\end{center}
\bigskip

끝으로 독자들에게 미리 알려 두는 것이 유익하다고 생각한다. 저자들 자신이 그랬듯이, 독자들도 스킴의 언어에 익숙해지고 기하학적 직관이 제안하는 통상적 구성들을 본질적으로 단 하나의 합리적인 방식으로 이 언어에 옮길 수 있음을 확신하기까지는 어느 정도 어려움을 겪을 것이다.
현대수학의 많은 분야에서 그렇듯이, 최초의 직관은 그것을 필요한 정밀성과 일반성을 모두 갖추어 표현하는 고유한 언어에서 겉보기에는 점점 더 멀어지고 있다.
여기서 심리적 어려움은 집합에 익숙한 개념들, 곧 데카르트 곱, 군·환·가군의 연산, 다발, 주동차다발 등을 집합의 범주와 이미 상당히 다른 범주의 대상들(즉 준스킴의 범주 또는 주어진 준스킴 위의 준스킴 범주)로 옮겨야 한다는 데 있다.
미래의 수학자가 이러한 새로운 추상화의 노력을 피하기는 어려울 것이다. 그러나 집합론에 익숙해지면서 우리의 선배들이 기울였던 노력과 비교하면, 결국 이 노력은 꽤 작은 것일지도 모른다.

\bigskip
\begin{center}*\quad*\quad*\end{center}
\bigskip

참조는 십진법 체계에 따라 표시한다. 예를 들어 III, 4.9.3에서 III은 장, 4는 절, 9는 그 절의 항을 나타낸다.
같은 장 안에서는 장 표시를 생략한다.
% ===== END EXACT INLINE INPUT: intro.tex =====

\clearpage

% ===== BEGIN EXACT INLINE INPUT: c0.tex =====
% Authority: EGA I ega0-1-fr.tex, SHA-256 F86663178E620A678BF0857B18BB4AC8FDB34E44BF5D241AE853EA7171E06C16.
% Sequential translation begins at authority lines 1-63.
\section{분수환}
\label{section:0.1-ko}

\setcounter{subsection}{-1}
\subsection{환과 대수}
\label{subsection:0.1.0-ko}

\oldpage[0\textsubscript{I}]{11}
\begin{env}[1.0.1]
\label{0.1.0.1-ko}
이 논고에서 다루는 모든 환은 \emph{단위원}을 갖는다.
그러한 환 위의 모든 가군은 \emph{단위원 가군}이라고 가정한다. 환 준동형은 언제나 단위원을 단위원으로 보낸다고 가정하며, 명시적으로 달리 말하지 않는 한 환 $A$의 부분환은 $A$의 \emph{단위원을 포함}한다고 가정한다.
주로 \emph{가환환}을 다루며, 아무 수식 없이 환이라고 말할 때에는 가환환을 뜻하는 것으로 이해한다.
$A$가 반드시 가환일 필요가 없는 환이면, 명시적으로 달리 말하지 않는 한 $A$-가군은 언제나 \emph{왼쪽} 가군을 뜻한다.
\end{env}

\begin{env}[1.0.2]
\label{0.1.0.2-ko}
$A$, $B$를 반드시 가환일 필요가 없는 두 환이라 하고, $\varphi:A\to B$를 준동형이라 하자.
모든 왼쪽(각각 오른쪽) $B$-가군 $M$에는 $a.m=\varphi(a).m$(각각 $m.a=m.\varphi(a)$)로 놓아 왼쪽(각각 오른쪽) $A$-가군 구조를 줄 수 있다. $M$ 위의 $A$-가군 구조와 $B$-가군 구조를 구별해야 할 때에는 이렇게 정의한 왼쪽(각각 오른쪽) $A$-가군을 $M_{[\varphi]}$로 나타낸다.
$L$이 $A$-가군이면 준동형 $u:L\to M_{[\varphi]}$는 곧 $a\in A$, $x\in L$에 대하여 $u(a.x)=\varphi(a).u(x)$를 만족하는 가환군 준동형이다. 이를 $L\to M$인 \emph{$\varphi$-준동형}이라고도 하며, 쌍 $(\varphi,u)$(또는 언어를 남용하여 $u$)를 $(A,L)$에서 $(B,M)$로 가는 \emph{디-준동형}이라고 한다.
따라서 환 $A$와 $A$-가군 $L$로 이루어진 쌍 $(A,L)$들은 디-준동형을 사상으로 하는 하나의 \emph{범주}를 이룬다.
\end{env}

\begin{env}[1.0.3]
\label{0.1.0.3-ko}
(1.0.2)의 가정 아래에서 $\mathfrak{J}$가 $A$의 왼쪽(각각 오른쪽) 아이디얼이면, $\varphi(\mathfrak{J})$가 생성하는 $B$의 왼쪽(각각 오른쪽) 아이디얼 $B\varphi(\mathfrak{J})$(각각 $\varphi(\mathfrak{J})B$)를 $B\mathfrak{J}$(각각 $\mathfrak{J}B$)로 나타낸다. 이는 왼쪽(각각 오른쪽) $B$-가군의 표준 준동형 $B\otimes_A\mathfrak{J}\to B$(각각 $\mathfrak{J}\otimes_A B\to B$)의 상이기도 하다.
\end{env}

\begin{env}[1.0.4]
\label{0.1.0.4-ko}
$A$가 (가환)환이고 $B$가 반드시 가환일 필요가 없는 환이면, $B$에 \emph{$A$-대수} 구조를 주는 것은 $\varphi(A)$가 $B$의 중심에 포함되도록 하는 환 준동형 $\varphi:A\to B$를 주는 것과 동치이다.
$A$의 모든 아이디얼 $\mathfrak{J}$에 대하여 $\mathfrak{J}B=B\mathfrak{J}$는 $B$의 양쪽 아이디얼이고, 모든 $B$-가군 $M$에 대하여 $\mathfrak{J}M$은 $(B\mathfrak{J})M$과 같은 $B$-부분가군이다.
\end{env}

\begin{env}[1.0.5]
\label{0.1.0.5-ko}
\emph{유한 생성 가군}과 \emph{유한 생성} (가환) \emph{대수}의 개념은 다시 설명하지 않겠다. $A$-가군 $M$이 유한 생성이라는 것은 완전열

\oldpage[0\textsubscript{I}]{12}
$A^p\to M\to 0$이 존재한다는 뜻이다.
$A$-가군 $M$이 어떤 준동형 $A^p\to A^q$의 여핵과 동형일 때, 다시 말해 완전열 $A^p\to A^q\to M\to 0$이 존재할 때 $M$이 \emph{유한 표시}를 갖는다고 한다.
\emph{뇌터 환} $A$ 위의 모든 유한 생성 $A$-가군은 유한 표시를 갖는다는 점에 유의하자.

$A$-대수 $B$의 모든 원소가 $A$에 계수를 갖는 일계수 다항식의 $B$ 안에서의 근일 때 $B$가 \emph{$A$ 위에서 정수적}이라고 함을 상기하자. 이는 $B$의 모든 원소가 유한 생성 \emph{$A$-가군}인 $B$의 어떤 부분대수에 들어간다는 것과 동치이다.
이 조건이 성립하고 $B$가 가환이면, $B$의 유한 부분집합이 생성하는 $B$의 부분대수는 유한 생성 $A$-가군이다. 따라서 가환대수 $B$가 $A$ 위에서 정수적이면서 유한 생성일 필요충분조건은 $B$가 유한 생성 $A$-가군인 것이다. 이때 $B$를 \emph{유한 정수적} $A$-대수라고도 한다(혼동의 우려가 없으면 단순히 \emph{유한} $A$-대수라고 한다).
이 정의들에서는 $A$-대수 구조를 정하는 준동형 $A\to B$가 단사라고 가정하지 않는다는 점에 유의하자.
\end{env}

\begin{env}[1.0.6]
\label{0.1.0.6-ko}
\emph{정역}은 $0$이 아닌 원소들로 이루어진 유한족의 곱이 $0$이 아닌 환이다. 이는 그러한 환에서 $0\neq 1$이고, $0$이 아닌 두 원소의 곱이 $0$이 아니라는 것과 동치이다.
환 $A$의 \emph{소 아이디얼}은 $A/\mathfrak{p}$가 정역이 되게 하는 아이디얼 $\mathfrak{p}$이다. 따라서 $\mathfrak{p}\neq A$이다.
환 $A$가 적어도 하나의 소 아이디얼을 가질 필요충분조건은 $A\neq\{0\}$이다.
\end{env}

\begin{env}[1.0.7]
\label{0.1.0.7-ko}
\emph{국소환}은 유일한 극대 아이디얼을 갖는 환 $A$이다. 이 극대 아이디얼은 가역원들의 여집합이며 $A$가 아닌 모든 아이디얼을 포함한다.
$A$, $B$가 각각 극대 아이디얼 $\mathfrak{m}$, $\mathfrak{n}$을 갖는 두 국소환일 때, 준동형 $\varphi:A\to B$가 $\varphi(\mathfrak{m})\subset\mathfrak{n}$을 만족하면(또는 이와 동치로 $\varphi^{-1}(\mathfrak{n})=\mathfrak{m}$이면) $\varphi$를 \emph{국소 준동형}이라고 한다.
몫으로 내려가면 이러한 준동형은 잉여체 $A/\mathfrak{m}$에서 잉여체 $B/\mathfrak{n}$으로 가는 단사 준동형을 정한다.
두 국소 준동형의 합성은 국소 준동형이다.
\end{env}

\subsection{아이디얼의 근기. 환의 멱영근기와 근기}
\label{subsection:0.1.1-ko}

\begin{env}[1.1.1]
\label{0.1.1.1-ko}
$\mathfrak{a}$를 환 $A$의 아이디얼이라 하자. $\mathfrak{a}$의 \emph{근기} $\mathfrak{r}(\mathfrak{a})$는 어떤 정수 $n>0$에 대하여 $x^n\in\mathfrak{a}$가 되는 $x\in A$ 전체의 집합이다. 이는 $\mathfrak{a}$를 포함하는 아이디얼이다.
$\mathfrak{r}(\mathfrak{r}(\mathfrak{a}))=\mathfrak{r}(\mathfrak{a})$이고, 관계 $\mathfrak{a}\subset\mathfrak{b}$는 $\mathfrak{r}(\mathfrak{a})\subset\mathfrak{r}(\mathfrak{b})$를 함의하며, 유한개의 아이디얼의 교집합의 근기는 각 근기의 교집합이다.
$\varphi$가 환 $A'$에서 $A$로 가는 준동형이면, 모든 아이디얼 $\mathfrak{a}\subset A$에 대하여 $\mathfrak{r}(\varphi^{-1}(\mathfrak{a}))=\varphi^{-1}(\mathfrak{r}(\mathfrak{a}))$이다.
한 아이디얼이 어떤 아이디얼의 근기일 필요충분조건은 그것이 소 아이디얼들의 교집합인 것이다.
아이디얼 $\mathfrak{a}$의 근기는 $\mathfrak{a}$를 포함하는 소 아이디얼들 가운데 \emph{극소}인 것들의 교집합이다. $A$가 뇌터 환이면 이러한 극소 소 아이디얼은 유한개이다.

영 아이디얼 $(0)$의 근기를 $A$의 \emph{멱영근기}라고도 한다. 이는 $A$의 멱영원 전체의 집합 $\mathfrak{N}$이다.
$\mathfrak{N}=(0)$일 때 환 $A$를 \emph{축소환}이라고 한다. 모든 환 $A$에 대하여 멱영근기로 나눈 몫환 $A/\mathfrak{N}$은 축소환이다.
\end{env}

\begin{env}[1.1.2]
\label{0.1.1.2-ko}
환 $A$(반드시 가환일 필요는 없다)의 \emph{근기} $\mathfrak{R}(A)$는 $A$의 극대 왼쪽 아이디얼 전체의 교집합(또한 극대 오른쪽 아이디얼 전체의 교집합)임을 상기하자.
$A/\mathfrak{R}(A)$의 근기는 $(0)$이다.
\end{env}

\subsection{분수가군과 분수환}
\label{subsection:0.1.2-ko}

\oldpage[0\textsubscript{I}]{13}
\begin{env}[1.2.1]
\label{0.1.2.1-ko}
환 $A$의 부분집합 $S$가 $1\in S$를 만족하고 $S$의 두 원소의 곱이 다시 $S$에 속하면 $S$를 \emph{곱셈적}이라고 한다.
이후 가장 중요한 예는 다음과 같다. $1^\circ$ 원소 $f\in A$의 거듭제곱 $f^n$($n\geq 0$) 전체의 집합 $S_f$. $2^\circ$ \emph{소} 아이디얼 $\mathfrak{p}$의 여집합 $A-\mathfrak{p}$.
\end{env}

\begin{env}[1.2.2]
\label{0.1.2.2-ko}
$S$를 환 $A$의 곱셈적 부분집합이라 하고, $M$을 $A$-가군이라 하자. 집합 $M\times S$에서 두 쌍 $(m_1,s_1)$, $(m_2,s_2)$ 사이의 관계
\[
  \text{``$s(s_1m_2-s_2m_1)=0$이 되는 $s\in S$가 존재한다''}
\]
는 동치관계이다.
이 관계에 의한 $M\times S$의 몫집합을 $S^{-1}M$으로 나타내고, 쌍 $(m,s)$의 $S^{-1}M$에서의 표준 상을 $m/s$로 나타낸다. 사상 $i_M^S:m\mapsto m/1$($i^S$로도 쓴다)을 $M$에서 $S^{-1}M$으로 가는 \emph{표준 사상}이라고 한다.
이 사상은 일반적으로 단사도 전사도 아니다. 그 핵은 어떤 $s\in S$에 대하여 $sm=0$이 되는 $m\in M$ 전체의 집합이다.

$S^{-1}M$ 위에
\[
  (m_1/s_1)+(m_2/s_2)=(s_2m_1+s_1m_2)/(s_1s_2)
\]
로 놓아 덧셈군 연산을 정의한다(이는 선택한 $S^{-1}M$의 원소의 표현에 무관함을 확인할 수 있다).
또 $S^{-1}A$ 위에 $(a_1/s_1)(a_2/s_2)=(a_1a_2)/(s_1s_2)$로 곱셈을 정의하고, 끝으로 $(a/s)(m/s')=(am)/(ss')$로 놓아 $S^{-1}A$를 작용소 집합으로 하는 $S^{-1}M$ 위의 외부 연산을 정의한다.
이로써 $S^{-1}A$에는 환 구조가 주어지고($S$에 분모를 갖는 $A$의 \emph{분수환}이라 한다), $S^{-1}M$에는 $S^{-1}A$-가군 구조가 주어진다($S$에 분모를 갖는 $M$의 \emph{분수가군}이라 한다). 모든 $s\in S$에 대하여 $s/1$은 $S^{-1}A$에서 가역이고 그 역원은 $1/s$이다.
표준 사상 $i_A^S$(각각 $i_M^S$)은 환 준동형(각각 $A$-가군 준동형)이다. 여기서 $S^{-1}M$은 준동형 $i_A^S:A\to S^{-1}A$를 통해 $A$-가군으로 본다.
\end{env}

\begin{env}[1.2.3]
\label{0.1.2.3-ko}
$f\in A$에 대하여 $S_f=\{f^n\}_{n\geq 0}$이면, $S_f^{-1}A$와 $S_f^{-1}M$ 대신 $A_f$와 $M_f$를 쓴다. $A_f$를 $A$ 위의 대수로 볼 때에는 $A_f=A[1/f]$라고 쓸 수 있다.
$A_f$는 몫대수 $A[T]/(fT-1)A[T]$와 동형이다.
$f=1$이면 $A_f$와 $M_f$는 각각 $A$와 $M$에 표준적으로 동일시된다. $f$가 멱영원이면 $A_f$와 $M_f$는 모두 영 대상이다.

$S=A-\mathfrak{p}$이고 $\mathfrak{p}$가 $A$의 소 아이디얼이면, $S^{-1}A$와 $S^{-1}M$ 대신 $A_{\mathfrak{p}}$와 $M_{\mathfrak{p}}$를 쓴다. $A_{\mathfrak{p}}$는 $i_A^S(\mathfrak{p})$가 생성하는 극대 아이디얼 $\mathfrak{q}$를 갖는 \emph{국소환}이고 $(i_A^S)^{-1}(\mathfrak{q})=\mathfrak{p}$이다. 몫으로 내려가면 $i_A^S$는 정역 $A/\mathfrak{p}$에서 체 $A_{\mathfrak{p}}/\mathfrak{q}$로 가는 단사 준동형을 주며, 이 체는 $A/\mathfrak{p}$의 분수체와 동일시된다.
\end{env}

\begin{env}[1.2.4]
\label{0.1.2.4-ko}
분수환 $S^{-1}A$와 표준 준동형 $i_A^S$는 하나의 \emph{보편 사상 문제}의 해이다. $u(S)$가 $B$의 \emph{가역원}들로 이루어지도록 하는 $A$에서 환 $B$로 가는 모든 준동형 $u$는 유일한 방식으로
\[
  u:A\xrightarrow{i_A^S}S^{-1}A\xrightarrow{u^*}B
\]

\oldpage[0\textsubscript{I}]{14}
와 같이 분해된다. 여기서 $u^*$는 환 준동형이다.
같은 가정 아래에서 $M$을 $A$-가군, $N$을 $B$-가군이라 하고, $v:M\to N$을 $A$-가군 준동형이라 하자($N$의 $A$-가군 구조는 $u:A\to B$로 정한다). 그러면 $v$는 유일한 방식으로
\[
  v:M\xrightarrow{i_M^S}S^{-1}M\xrightarrow{v^*}N
\]
와 같이 분해된다. 여기서 $v^*$는 $S^{-1}A$-가군 준동형이다($N$의 $S^{-1}A$-가군 구조는 $u^*$로 정한다).
\end{env}

\begin{env}[1.2.5]
\label{0.1.2.5-ko}
원소 $(a/s)\otimes m$에 원소 $(am)/s$를 대응시켜 $S^{-1}A$-가군의 표준 동형 $S^{-1}A\otimes_A M\xrightarrow{\sim}S^{-1}M$을 정의한다. 그 역동형은 $m/s$를 $(1/s)\otimes m$으로 보낸다.
\end{env}

\begin{env}[1.2.6]
\label{0.1.2.6-ko}
$S^{-1}A$의 모든 아이디얼 $\mathfrak{a}'$에 대하여 $\mathfrak{a}=(i_A^S)^{-1}(\mathfrak{a}')$은 $A$의 아이디얼이고, $\mathfrak{a}'$은 $i_A^S(\mathfrak{a})$가 생성하는 $S^{-1}A$의 아이디얼이며 $S^{-1}\mathfrak{a}$와 동일시된다(1.3.2).
사상 $\mathfrak{p}'\mapsto(i_A^S)^{-1}(\mathfrak{p}')$은 $S^{-1}A$의 \emph{소} 아이디얼 전체의 순서집합에서 $\mathfrak{p}\cap S=\varnothing$을 만족하는 $A$의 소 아이디얼 $\mathfrak{p}$ 전체의 순서집합으로 가는 순서동형이다.
또한 이때 국소환 $A_{\mathfrak{p}}$와 $(S^{-1}A)_{S^{-1}\mathfrak{p}}$는 표준적으로 동형이다(1.5.1).
\end{env}

\begin{env}[1.2.7]
\label{0.1.2.7-ko}
$A$가 \emph{정역}이고 그 분수체를 $K$로 나타내면, $0$을 포함하지 않는 모든 곱셈적 부분집합 $S$에 대하여 표준 사상 $i_A^S:A\to S^{-1}A$는 단사이고, $S^{-1}A$는 $A$를 포함하는 $K$의 부분환과 표준적으로 동일시된다.
특히 $A$의 모든 소 아이디얼 $\mathfrak{p}$에 대하여 $A_{\mathfrak{p}}$는 $A$를 포함하고 극대 아이디얼 $\mathfrak{p}A_{\mathfrak{p}}$를 갖는 국소환이며 $\mathfrak{p}A_{\mathfrak{p}}\cap A=\mathfrak{p}$이다.
\end{env}

\begin{env}[1.2.8]
\label{0.1.2.8-ko}
$A$가 \emph{축소환}이면(1.1.1), $S^{-1}A$도 축소환이다. 실제로 $x\in A$, $s\in S$에 대하여 $(x/s)^n=0$이면 $s'x^n=0$이 되는 $s'\in S$가 존재한다. 따라서 $(s'x)^n=0$이고, 가정에 따라 $s'x=0$이므로 $x/s=0$이다.
\end{env}

\subsection{함자적 성질}
\label{subsection:0.1.3-ko}

\begin{env}[1.3.1]
\label{0.1.3.1-ko}
$M$, $N$을 두 $A$-가군이라 하고, $u$를 $A$-준동형 $M\to N$이라 하자.
$S$가 $A$의 곱셈적 부분집합이면, $(S^{-1}u)(m/s)=u(m)/s$로 놓아 $S^{-1}u$로 나타내는 $S^{-1}A$-준동형 $S^{-1}M\to S^{-1}N$을 정의한다. $S^{-1}M$과 $S^{-1}N$을 각각 $S^{-1}A\otimes_A M$과 $S^{-1}A\otimes_A N$에 표준적으로 동일시하면(1.2.5), $S^{-1}u$는 $1\otimes u$와 동일시된다.
$P$가 세 번째 $A$-가군이고 $v$가 $A$-준동형 $N\to P$이면 $S^{-1}(v\circ u)=(S^{-1}v)\circ(S^{-1}u)$이다. 다시 말해 $A$와 $S$를 고정하면 $S^{-1}M$은 $A$-가군의 범주에서 $S^{-1}A$-가군의 범주로 가는, \emph{$M$에 관한 공변 함자}이다.
\end{env}

\begin{env}[1.3.2]
\label{0.1.3.2-ko}
함자 $S^{-1}M$은 \emph{완전}하다. 다시 말해 열
\[
  M\xrightarrow{u}N\xrightarrow{v}P
\]
이 완전하면 열
\[
  S^{-1}M\xrightarrow{S^{-1}u}S^{-1}N\xrightarrow{S^{-1}v}S^{-1}P
\]
도 완전하다.
특히 $u:M\to N$이 단사(각각 전사)이면 $S^{-1}u$도 단사(각각 전사)이다.

\oldpage[0\textsubscript{I}]{15}
$N$과 $P$가 $M$의 두 부분가군이면, $S^{-1}N$과 $S^{-1}P$는 $S^{-1}M$의 부분가군에 표준적으로 동일시되며
\[
  S^{-1}(N+P)=S^{-1}N+S^{-1}P
  \quad\text{이고}\quad
  S^{-1}(N\cap P)=(S^{-1}N)\cap(S^{-1}P)
\]
이다.
\end{env}

\begin{env}[1.3.3]
\label{0.1.3.3-ko}
$(M_\alpha,\varphi_{\beta\alpha})$를 $A$-가군의 귀납계라 하자. 그러면 $(S^{-1}M_\alpha,S^{-1}\varphi_{\beta\alpha})$는 $S^{-1}A$-가군의 귀납계이다.
$S^{-1}M_\alpha$와 $S^{-1}\varphi_{\beta\alpha}$를 텐서곱으로 나타내고(1.2.5, 1.3.1), 텐서곱과 귀납극한을 취하는 연산들이 서로 교환한다는 사실을 쓰면 표준 동형
\[
  S^{-1}\varinjlim M_\alpha\xrightarrow{\sim}\varinjlim S^{-1}M_\alpha
\]
를 얻는다. 이를 함자 $S^{-1}M$($M$에 관한 함자)이 \emph{귀납극한과 가환한다}고도 표현한다.
\end{env}

\begin{env}[1.3.4]
\label{0.1.3.4-ko}
$M$, $N$을 두 $A$-가군이라 하자. $(M,N)$에 관해 \emph{함자적인} 표준 동형
\[
  (S^{-1}M)\otimes_{S^{-1}A}(S^{-1}N)
  \xrightarrow{\sim}
  S^{-1}(M\otimes_A N)
\]
이 존재하며, 이는 $(m/s)\otimes(n/t)$를 $(m\otimes n)/st$로 보낸다.
\end{env}

\begin{env}[1.3.5]
\label{0.1.3.5-ko}
마찬가지로 $(M,N)$에 관해 \emph{함자적인} 준동형
\[
  S^{-1}\operatorname{Hom}_A(M,N)
  \longrightarrow
  \operatorname{Hom}_{S^{-1}A}(S^{-1}M,S^{-1}N)
\]
이 존재하며, 이는 $u/s$에 준동형 $m/t\mapsto u(m)/st$를 대응시킨다.
$M$이 \emph{유한 표시}를 가지면 앞의 준동형은 \emph{동형}이다. $M$이 $A^r$ 꼴일 때에는 바로 알 수 있고, 일반적인 경우에는 완전열 $A^p\to A^q\to M\to 0$에서 출발하여 함자 $S^{-1}M$의 완전성과 $M$에 관한 함자 $\operatorname{Hom}_A(M,N)$의 왼쪽 완전성을 사용하면 된다.
$A$가 \emph{뇌터 환}이고 $A$-가군 $M$이 \emph{유한 생성}이면 언제나 이 경우에 해당한다는 점에 유의하자.
\end{env}

\subsection{곱셈적 부분집합의 변경}
\label{subsection:0.1.4-ko}

\begin{env}[1.4.1]
\label{0.1.4.1-ko}
$S$, $T$를 $S\subset T$인 환 $A$의 두 곱셈적 부분집합이라 하자. $S^{-1}A$의 $a/s$로 나타낸 원소를 $T^{-1}A$에서 역시 $a/s$로 나타낸 원소에 대응시키는 표준 준동형 $\rho_A^{T,S}$($\rho^{T,S}$로도 쓴다)이 $S^{-1}A$에서 $T^{-1}A$로 존재하며
\[
  i_A^T=\rho_A^{T,S}\circ i_A^S
\]
이다.
모든 $A$-가군 $M$에 대하여 $S^{-1}M$에서 $T^{-1}M$으로 가는 $S^{-1}A$-선형 사상도 존재한다. 여기서 $T^{-1}M$은 준동형 $\rho_A^{T,S}$를 통해 $S^{-1}A$-가군으로 보며, 이 사상은 $S^{-1}M$의 원소 $m/s$를 $T^{-1}M$의 원소 $m/s$로 보낸다. 이 사상을 $\rho_M^{T,S}$ 또는 단순히 $\rho^{T,S}$로 나타내며
\[
  i_M^T=\rho_M^{T,S}\circ i_M^S
\]
이다. 표준 동일시 (1.2.5) 아래에서 $\rho_M^{T,S}$는 $\rho_A^{T,S}\otimes 1$과 동일시된다.
준동형 $\rho_M^{T,S}$는 함자 $S^{-1}M$에서 함자 $T^{-1}M$으로 가는 \emph{함자적 사상}(또는 자연변환)이다. 다시 말해 그림
\[
\begin{tikzcd}[column sep=large,row sep=large]
S^{-1}M \arrow[r,"S^{-1}u"] \arrow[d,"\rho_M^{T,S}"'] & S^{-1}N \arrow[d,"\rho_N^{T,S}"] \\
T^{-1}M \arrow[r,"T^{-1}u"] & T^{-1}N
\end{tikzcd}
\]

\oldpage[0\textsubscript{I}]{16}
은 모든 준동형 $u:M\to N$에 대하여 가환한다. 또한 $T^{-1}u$는 $S^{-1}u$에 의해 완전히 결정된다는 점에 유의하자. 실제로 $m\in M$, $t\in T$에 대하여
\[
  (T^{-1}u)(m/t)=(t/1)^{-1}\rho^{T,S}\bigl((S^{-1}u)(m/1)\bigr)
\]
이다.
\end{env}

\begin{env}[1.4.2]
\label{0.1.4.2-ko}
같은 표기 아래에서 두 $A$-가군 $M$, $N$에 관한 다음 그림들은 가환한다(1.3.4, 1.3.5 참조).
\[
\begin{tikzcd}[column sep=large,row sep=large]
(S^{-1}M)\otimes_{S^{-1}A}(S^{-1}N) \arrow[r,"\sim"] \arrow[d]
  & S^{-1}(M\otimes_A N) \arrow[d]
  \\
(T^{-1}M)\otimes_{T^{-1}A}(T^{-1}N) \arrow[r,"\sim"]
  & T^{-1}(M\otimes_A N)
\end{tikzcd}
\]
\[
\begin{tikzcd}[column sep=large,row sep=large]
S^{-1}\operatorname{Hom}_A(M,N) \arrow[r] \arrow[d]
  & \operatorname{Hom}_{S^{-1}A}(S^{-1}M,S^{-1}N) \arrow[d] \\
T^{-1}\operatorname{Hom}_A(M,N) \arrow[r]
  & \operatorname{Hom}_{T^{-1}A}(T^{-1}M,T^{-1}N)
\end{tikzcd}
\]
\end{env}

\begin{env}[1.4.3]
\label{0.1.4.3-ko}
준동형 $\rho^{T,S}$가 \emph{전단사}인 중요한 경우가 있다. 곧 $T$의 모든 원소가 $S$의 어떤 원소의 약수인 경우이다. 이때 $\rho^{T,S}$를 통해 가군 $S^{-1}M$과 $T^{-1}M$을 동일시한다.
$S$의 원소의 $A$ 안에서의 모든 약수가 다시 $S$에 속할 때 $S$를 \emph{포화}되었다고 한다. $S$를 $S$의 원소들의 약수 전체의 집합 $T$로 바꾸면($T$는 곱셈적이고 포화되어 있다), 원한다면 언제나 $S$가 포화된 분수가군 $S^{-1}M$만을 고려해도 됨을 알 수 있다.
\end{env}

\begin{env}[1.4.4]
\label{0.1.4.4-ko}
$S$, $T$, $U$가 $S\subset T\subset U$인 $A$의 세 곱셈적 부분집합이면
\[
  \rho^{U,S}=\rho^{U,T}\circ\rho^{T,S}
\]
이다.
\end{env}

\begin{env}[1.4.5]
\label{0.1.4.5-ko}
$A$의 곱셈적 부분집합들의 \emph{증가 여과족} $(S_\alpha)$를 생각하자($S_\alpha\subset S_\beta$일 때 $\alpha\leq\beta$라고 쓴다). 곱셈적 부분집합 $S=\bigcup_\alpha S_\alpha$로 놓고
\[
  \rho_{\beta\alpha}=\rho_A^{S_\beta,S_\alpha}\quad\text{($\alpha\leq\beta$일 때)}
\]
로 놓자.
(1.4.4)에 따라 준동형 $\rho_{\beta\alpha}$들은 환의 귀납계 $(S_\alpha^{-1}A,\rho_{\beta\alpha})$의 \emph{귀납극한}인 환 $A'$을 정의한다.
$\rho_\alpha$를 표준 사상 $S_\alpha^{-1}A\to A'$이라 하고 $\varphi_\alpha=\rho_A^{S,S_\alpha}$로 놓자. (1.4.4)에 따라 $\alpha\leq\beta$이면 $\varphi_\alpha=\varphi_\beta\circ\rho_{\beta\alpha}$이므로, 그림
\[
\begin{tikzcd}[column sep=huge,row sep=large]
  & S_\alpha^{-1}A
      \arrow[ddl,"\rho_\alpha"']
      \arrow[d,"\rho_{\beta\alpha}"]
      \arrow[ddr,"\varphi_\alpha"] & \\
  & S_\beta^{-1}A
      \arrow[dl,"\rho_\beta"]
      \arrow[dr,"\varphi_\beta"'] & \\
A' \arrow[rr,"\varphi"'] & & S^{-1}A
\end{tikzcd}
\qquad(\alpha\leq\beta)
\]
이 가환하도록 하는 준동형 $\varphi:A'\to S^{-1}A$를 유일하게 정의할 수 있다.
실제로 $\varphi$는 \emph{동형}이다. 구성으로부터 $\varphi$가 전사라는 것은 바로 알 수 있다.
한편 $\varphi(\rho_\alpha(a/s_\alpha))=0$인 $\rho_\alpha(a/s_\alpha)\in A'$이 있다고 하자. 이는 $S^{-1}A$에서 $a/s_\alpha=0$, 즉 $sa=0$인 $s\in S$가 존재한다는 뜻이다. 그런데 $s\in S_\beta$인 $\beta\geq\alpha$가 존재하므로 $\rho_\alpha(a/s_\alpha)=\rho_\beta(sa/ss_\alpha)=0$이다. 따라서 $\varphi$는 단사이다.
$A$-가군 $M$의 경우도 똑같이 다루어 표준 동형
\[
  \varinjlim S_\alpha^{-1}A\xrightarrow{\sim}
  (\varinjlim S_\alpha)^{-1}A,
  \qquad
  \varinjlim S_\alpha^{-1}M\xrightarrow{\sim}
  (\varinjlim S_\alpha)^{-1}M
\]
을 얻는다. 둘째 동형은 \emph{$M$에 관해 함자적}이다.
\end{env}

\oldpage[0\textsubscript{I}]{17}
\begin{env}[1.4.6]
\label{0.1.4.6-ko}
$S_1$, $S_2$를 $A$의 두 곱셈적 부분집합이라 하자. 그러면 $S_1S_2$도 $A$의 곱셈적 부분집합이다.
$S_2'$을 환 $S_1^{-1}A$ 안에서의 $S_2$의 표준 상이라 하자. 이는 이 환의 곱셈적 부분집합이다.
모든 $A$-가군 $M$에 대하여 함자적 동형
\[
  S_2'^{-1}(S_1^{-1}M)\xrightarrow{\sim}(S_1S_2)^{-1}M
\]
이 존재하며, 이는 $(m/s_1)/(s_2/1)$을 $m/(s_1s_2)$에 대응시킨다.
\end{env}

\subsection{환의 변경}
\label{subsection:0.1.5-ko}

\begin{env}[1.5.1]
\label{0.1.5.1-ko}
$A$, $A'$을 두 환, $\varphi$를 준동형 $A'\to A$, $S$(각각 $S'$)를 $\varphi(S')\subset S$를 만족하는 $A$(각각 $A'$)의 곱셈적 부분집합이라 하자. 합성 준동형
\[
  A'\xrightarrow{\varphi}A\longrightarrow S^{-1}A
\]
은 (1.2.4)에 따라
\[
  A'\longrightarrow S'^{-1}A'\xrightarrow{\varphi^{S'}}S^{-1}A
\]
와 같이 분해되며, $\varphi^{S'}(a'/s')=\varphi(a')/\varphi(s')$이다.
$A=\varphi(A')$이고 $S=\varphi(S')$이면 $\varphi^{S'}$는 \emph{전사}이다.
$A'=A$이고 $\varphi$가 항등사상이면 $\varphi^{S'}$는 (1.4.1)에서 정의한 준동형 $\rho_A^{S,S'}$와 다르지 않다.
\end{env}

\begin{env}[1.5.2]
\label{0.1.5.2-ko}
(1.5.1)의 가정 아래에서 $M$을 $A$-가군이라 하자.
$S'^{-1}A'$-가군의 표준적인 함자적 준동형
\[
  \sigma:S'^{-1}(M_{[\varphi]})\longrightarrow
  (S^{-1}M)_{[\varphi^{S'}]}
\]
이 존재하며, 이는 $S'^{-1}(M_{[\varphi]})$의 모든 원소 $m/s'$을 $(S^{-1}M)_{[\varphi^{S'}]}$의 원소 $m/\varphi(s')$에 대응시킨다. 이 정의가 해당 원소의 표현 $m/s'$에 무관하다는 것은 바로 확인할 수 있다.
$S=\varphi(S')$이면 준동형 $\sigma$는 \emph{전단사}이다.
$A'=A$이고 $\varphi$가 항등사상이면 $\sigma$는 (1.4.1)에서 정의한 준동형 $\rho_M^{S,S'}$와 다르지 않다.

특히 $M=A$로 놓으면 준동형 $\varphi$는 $A$ 위에 $A'$-대수 구조를 정의한다. 이때 $S'^{-1}(A_{[\varphi]})$에는 $(\varphi(S'))^{-1}A$와 동일시되는 환 구조가 주어지고, 준동형
\[
  \sigma:S'^{-1}(A_{[\varphi]})\longrightarrow S^{-1}A
\]
는 $S'^{-1}A'$-대수의 준동형이다.
\end{env}

\begin{env}[1.5.3]
\label{0.1.5.3-ko}
$M$, $N$을 두 $A$-가군이라 하자. (1.3.4)와 (1.5.2)에서 정의한 준동형들을 합성하면 준동형
\[
  (S^{-1}M\otimes_{S^{-1}A}S^{-1}N)_{[\varphi^{S'}]}
  \longleftarrow
  S'^{-1}((M\otimes_A N)_{[\varphi]})
\]
을 얻으며, $\varphi(S')=S$이면 이는 동형이다.
마찬가지로 준동형 (1.3.5)와 (1.5.2)를 합성하면 준동형
\[
  S'^{-1}((\operatorname{Hom}_A(M,N))_{[\varphi]})
  \longrightarrow
  (\operatorname{Hom}_{S^{-1}A}(S^{-1}M,S^{-1}N))_{[\varphi^{S'}]}
\]
을 얻으며, $\varphi(S')=S$이고 $M$이 유한 표시를 가지면 이는 동형이다.
\end{env}

\begin{env}[1.5.4]
\label{0.1.5.4-ko}
이제 $A'$-가군 $N'$을 생각하고 텐서곱 $N'\otimes_{A'}A_{[\varphi]}$를 만들자. 이를
\[
  a.(n'\otimes b)=n'\otimes(ab)
\]
로 놓아 $A$-가군으로 볼 수 있다.
$S^{-1}A$-가군의 함자적 동형
\[
  \tau:(S'^{-1}N')\otimes_{S'^{-1}A'}(S^{-1}A)_{[\varphi^{S'}]}
  \xrightarrow{\sim}
  S^{-1}(N'\otimes_{A'}A_{[\varphi]})
\]
이 존재한다.

\oldpage[0\textsubscript{I}]{18}
이는 원소 $(n'/s')\otimes(a/s)$를 원소 $(n'\otimes a)/(\varphi(s')s)$에 대응시킨다. $n'/s'$(각각 $a/s$)를 같은 원소의 다른 표현으로 바꾸어도 $(n'\otimes a)/(\varphi(s')s)$가 변하지 않는다는 것을 각각 확인할 수 있다. 한편 $(n'\otimes a)/s$에 $(n'/1)\otimes(a/s)$를 대응시켜 $\tau$의 역준동형을 정의할 수 있다. 여기서는 $S^{-1}(N'\otimes_{A'}A_{[\varphi]})$가 $(N'\otimes_{A'}A_{[\varphi]})\otimes_A S^{-1}A$와 표준적으로 동형이고(1.2.5), 따라서 $A'$에서 $S^{-1}A$로 가는 합성 준동형 $a'\mapsto\varphi(a')/1$을 $\psi$로 나타낼 때 $N'\otimes_{A'}(S^{-1}A)_{[\psi]}$와도 표준적으로 동형이라는 사실을 사용한다.
\end{env}

\begin{env}[1.5.5]
\label{0.1.5.5-ko}
$M'$, $N'$이 두 $A'$-가군이면, (1.3.4)와 (1.5.4)의 동형들을 합성하여 동형
\[
  S'^{-1}M'\otimes_{S'^{-1}A'}S'^{-1}N'
  \otimes_{S'^{-1}A'}S^{-1}A
  \xrightarrow{\sim}
  S^{-1}(M'\otimes_{A'}N'\otimes_{A'}A)
\]
을 얻는다.
마찬가지로 $M'$이 유한 표시를 가지면 (1.3.5)와 (1.5.4)에 따라 동형
\[
  \operatorname{Hom}_{S'^{-1}A'}(S'^{-1}M',S'^{-1}N')
  \otimes_{S'^{-1}A'}S^{-1}A
  \xrightarrow{\sim}
  S^{-1}(\operatorname{Hom}_{A'}(M',N')\otimes_{A'}A)
\]
을 얻는다.
\end{env}

\begin{env}[1.5.6]
\label{0.1.5.6-ko}
(1.5.1)의 가정 아래에서 $T$(각각 $T'$)를 $S\subset T$(각각 $S'\subset T'$)이고 $\varphi(T')\subset T$인 $A$(각각 $A'$)의 두 번째 곱셈적 부분집합이라 하자. 그러면 그림
\[
\begin{tikzcd}[column sep=huge,row sep=large]
S'^{-1}A' \arrow[r,"\varphi^{S'}"] \arrow[d,"\rho^{T',S'}"']
  & S^{-1}A \arrow[d,"\rho^{T,S}"] \\
T'^{-1}A' \arrow[r,"\varphi^{T'}"']
  & T^{-1}A
\end{tikzcd}
\]
은 가환한다. $M$이 $A$-가군이면 그림
\[
\begin{tikzcd}[column sep=huge,row sep=large]
S'^{-1}(M_{[\varphi]}) \arrow[r,"\sigma"] \arrow[d,"\rho^{T',S'}"']
  & (S^{-1}M)_{[\varphi^{S'}]} \arrow[d,"\rho^{T,S}"] \\
T'^{-1}(M_{[\varphi]}) \arrow[r,"\sigma"']
  & (T^{-1}M)_{[\varphi^{T'}]}
\end{tikzcd}
\]
도 가환한다. 끝으로 $N'$이 $A'$-가군이면 그림
\[
\begin{tikzcd}[column sep=huge,row sep=large]
(S'^{-1}N')\otimes_{S'^{-1}A'}(S^{-1}A)_{[\varphi^{S'}]}
  \arrow[r,"\tau","{\sim}"'] \arrow[d]
  & S^{-1}(N'\otimes_{A'}A_{[\varphi]}) \arrow[d,"\rho^{T,S}"] \\
(T'^{-1}N')\otimes_{T'^{-1}A'}(T^{-1}A)_{[\varphi^{T'}]}
  \arrow[r,"\tau"',"{\sim}"]
  & T^{-1}(N'\otimes_{A'}A_{[\varphi]})
\end{tikzcd}
\]
도 가환한다. 왼쪽의 수직 화살표는 $S'^{-1}N'$에 $\rho_{N'}^{T',S'}$를, $S^{-1}A$에 $\rho_A^{T,S}$를 적용하여 얻는다.
\end{env}

\oldpage[0\textsubscript{I}]{19}

\begin{env}[1.5.7]
\label{0.1.5.7-ko}
$A''$을 세 번째 환, $\varphi':A''\to A'$을 환 준동형, $S''$을 $\varphi'(S'')\subset S'$인 $A''$의 곱셈적 부분집합이라 하자. $\varphi''=\varphi\circ\varphi'$로 놓으면
\[
  \varphi''^{S''}=\varphi^{S'}\circ\varphi'^{S''}
\]
이다.
$M$을 $A$-가군이라 하자. 분명히 $M_{[\varphi'']}=(M_{[\varphi]})_{[\varphi']}$이다. $\sigma$가 (1.5.2)에서 $\varphi$로부터 정의된 것과 같은 방식으로 $\varphi'$과 $\varphi''$에서 정의된 준동형을 각각 $\sigma'$과 $\sigma''$이라 하면 추이 공식
\[
  \sigma''=\sigma\circ\sigma'
\]
이 성립한다.
끝으로 $N''$을 $A''$-가군이라 하자. $A$-가군 $N''\otimes_{A''}A_{[\varphi'']}$는
\[
  (N''\otimes_{A''}A'_{[\varphi']})\otimes_{A'}A_{[\varphi]}
\]
와 표준적으로 동일시된다.
마찬가지로 $S^{-1}A$-가군
\[
  (S''^{-1}N'')\otimes_{S''^{-1}A''}
  (S^{-1}A)_{[\varphi''^{S''}]}
\]
는
\[
  \bigl((S''^{-1}N'')\otimes_{S''^{-1}A''}
  (S'^{-1}A')_{[\varphi'^{S''}]}\bigr)
  \otimes_{S'^{-1}A'}(S^{-1}A)_{[\varphi^{S'}]}
\]
와 표준적으로 동일시된다.
이 동일시들 아래에서 $\tau$가 (1.5.4)에서 $\varphi$로부터 정의된 것과 같은 방식으로 $\varphi'$과 $\varphi''$에서 정의된 동형을 각각 $\tau'$과 $\tau''$이라 하면 추이 공식
\[
  \tau''=\tau\circ(\tau'\otimes 1)
\]
이 성립한다.
\end{env}

\begin{env}[1.5.8]
\label{0.1.5.8-ko}
$A$를 환 $B$의 부분환이라 하자. $A$의 모든 \emph{극소} 소 아이디얼 $\mathfrak p$에 대하여 $\mathfrak p=A\cap\mathfrak q$가 되는 $B$의 극소 소 아이디얼 $\mathfrak q$가 존재한다. 실제로 $A_{\mathfrak p}$는 $B_{\mathfrak p}$의 부분환이고(1.3.2), 유일한 소 아이디얼 $\mathfrak p'$을 갖는다(1.2.6). $B_{\mathfrak p}$는 영환이 아니므로 적어도 하나의 소 아이디얼 $\mathfrak q'$을 가지며, 반드시 $\mathfrak q'\cap A_{\mathfrak p}=\mathfrak p'$이다. 따라서 $\mathfrak q'$의 역상인 $B$의 소 아이디얼 $\mathfrak q_1$은 $\mathfrak q_1\cap A=\mathfrak p$를 만족한다. 그러므로 특히 $\mathfrak q_1$에 포함된 $B$의 모든 극소 소 아이디얼 $\mathfrak q$에 대하여 $\mathfrak q\cap A=\mathfrak p$이다.
\end{env}

\subsection{가군 $M_f$와 귀납극한의 동일시}
\label{subsection:0.1.6-ko}

\begin{env}[1.6.1]
\label{0.1.6.1-ko}
$M$을 $A$-가군, $f$를 $A$의 원소라 하자. 모두 $M$과 같은 $A$-가군들의 열 $(M_n)$을 생각하고, 정수의 모든 쌍 $m\leq n$에 대하여 $\varphi_{nm}$을 $M_m$에서 $M_n$으로 가는 준동형 $z\mapsto f^{n-m}z$라 하자. $((M_n),(\varphi_{nm}))$이 $A$-가군의 \emph{귀납계}임은 자명하다. 이 계의 귀납극한을 $N=\varinjlim M_n$이라 하자. 이제 $N$에서 $M_f$ 위로 가는 함자적인 표준 $A$-동형을 정의하겠다. 이를 위해 모든 $n$에 대하여
\[
  \theta_n:z\longmapsto z/f^n
\]
은 $M=M_n$에서 $M_f$로 가는 $A$-준동형이고, 정의로부터 $m\leq n$일 때 $\theta_n\circ\varphi_{nm}=\theta_m$임에 유의하자. 따라서 표준 준동형 $M_n\to N$을 $\varphi_n$이라 할 때, 모든 $n$에 대하여 $\theta_n=\theta\circ\varphi_n$을 만족하는 $A$-준동형 $\theta:N\to M_f$가 존재한다. 가정에 따라 $M_f$의 모든 원소는 어떤 $n$에 대하여 $z/f^n$의 꼴이므로 $\theta$가 전사임은 분명하다. 한편 $\theta(\varphi_n(z))=0$, 즉 $z/f^n=0$이면 $f^kz=0$인 정수 $k>0$이 존재한다. 따라서 $\varphi_{n+k,n}(z)=0$이고, 이로부터 $\varphi_n(z)=0$이 따른다. 그러므로 $\theta$를 통해 $M_f$와 $\varinjlim M_n$을 동일시할 수 있다.
\end{env}

\begin{env}[1.6.2]
\label{0.1.6.2-ko}
이제 $M_n$, $\varphi_{nm}$, $\varphi_n$ 대신 각각 $M_{f,n}$, $\varphi_{nm}^f$, $\varphi_n^f$라 쓰자. $g$를 $A$의 두 번째 원소라 하자. $f^n$이 $f^ng^n$의 약수이므로 함자적 준동형
\[
  \rho_{fg,f}:M_f\longrightarrow M_{fg}\qquad\text{(1.4.1과 1.4.3)}
\]
이 존재한다.

\oldpage[0\textsubscript{I}]{20}

$M_f$와 $M_{fg}$를 각각 $\varinjlim M_{f,n}$과 $\varinjlim M_{fg,n}$으로 동일시하면, $\rho_{fg,f}$는
\[
  \rho_{fg,f}^n:M_{f,n}\longrightarrow M_{fg,n},
  \qquad \rho_{fg,f}^n(z)=g^nz
\]
로 정의되는 사상들의 \emph{귀납극한}과 동일시된다. 실제로 이는 그림
\[
\begin{tikzcd}[column sep=huge,row sep=large]
M_{f,n} \arrow[r,"\rho_{fg,f}^n"] \arrow[d,"\varphi_n^f"']
  & M_{fg,n} \arrow[d,"\varphi_n^{fg}"] \\
M_f \arrow[r,"\rho_{fg,f}"]
  & M_{fg}
\end{tikzcd}
\]
의 가환성에서 바로 따른다.
\end{env}

\subsection{가군의 지지집합}
\label{subsection:0.1.7-ko}

\begin{env}[1.7.1]
\label{0.1.7.1-ko}
$A$-가군 $M$이 주어졌을 때, $M_{\mathfrak p}\neq0$인 $A$의 소 아이디얼 $\mathfrak p$들의 집합을 $M$의 \emph{지지집합}이라 하고 $\operatorname{Supp}(M)$으로 나타낸다. $M=0$이기 위한 필요충분조건은 $\operatorname{Supp}(M)=\varnothing$인 것이다. 실제로 모든 $\mathfrak p$에 대하여 $M_{\mathfrak p}=0$이면, 임의의 원소 $x\in M$의 소멸자는 $A$의 어떤 소 아이디얼에도 포함될 수 없으므로 $A$ 전체와 같다.
\end{env}

\begin{env}[1.7.2]
\label{0.1.7.2-ko}
$A$-가군의 완전열 $0\to N\to M\to P\to0$이 있으면
\[
  \operatorname{Supp}(M)=\operatorname{Supp}(N)\cup\operatorname{Supp}(P)
\]
이다. 실제로 $A$의 모든 소 아이디얼 $\mathfrak p$에 대하여 열
\[
  0\longrightarrow N_{\mathfrak p}\longrightarrow M_{\mathfrak p}
  \longrightarrow P_{\mathfrak p}\longrightarrow0
\]
은 완전하고(1.3.2), $M_{\mathfrak p}=0$이기 위한 필요충분조건은 $N_{\mathfrak p}=P_{\mathfrak p}=0$인 것이다.
\end{env}

\begin{env}[1.7.3]
\label{0.1.7.3-ko}
$M$이 부분가군들의 족 $(M_\lambda)$의 합이면, $A$의 모든 소 아이디얼 $\mathfrak p$에 대하여 $M_{\mathfrak p}$는 $(M_\lambda)_{\mathfrak p}$들의 합이므로(1.3.3과 1.3.2)
\[
  \operatorname{Supp}(M)=\bigcup_\lambda\operatorname{Supp}(M_\lambda)
\]
이다.
\end{env}

\begin{env}[1.7.4]
\label{0.1.7.4-ko}
$M$이 \emph{유한 생성} $A$-가군이면, $\operatorname{Supp}(M)$은 \emph{$M$의 소멸자를 포함하는} 소 아이디얼들의 집합이다. 실제로 $M$이 한 원소로 생성되고 $x$가 그 생성원이면, $M_{\mathfrak p}=0$이라는 것은 $s.x=0$인 $s\notin\mathfrak p$가 존재한다는 뜻이고, 따라서 $\mathfrak p$가 $x$의 소멸자를 포함하지 않는다는 뜻이다. 이제 $M$이 유한 생성계 $(x_i)_{1\leq i\leq n}$을 가지고 $\mathfrak a_i$가 $x_i$의 소멸자라 하자. (1.7.3)에 따라 $\operatorname{Supp}(M)$은 $\mathfrak a_i$들 가운데 하나를 포함하는 $\mathfrak p$들의 집합이며, 이는 $M$의 소멸자인 $\mathfrak a=\bigcap_i\mathfrak a_i$를 포함하는 $\mathfrak p$들의 집합과 같다.
\end{env}

\begin{env}[1.7.5]
\label{0.1.7.5-ko}
$M$, $N$이 두 \emph{유한 생성} $A$-가군이면
\[
  \operatorname{Supp}(M\otimes_A N)
  =\operatorname{Supp}(M)\cap\operatorname{Supp}(N)
\]
이다. 이를 위해서는 $A$의 소 아이디얼 $\mathfrak p$에 대하여 조건
\[
  M_{\mathfrak p}\otimes_{A_{\mathfrak p}}N_{\mathfrak p}\neq0
\]
이 ``$M_{\mathfrak p}\neq0$이고 $N_{\mathfrak p}\neq0$이다''와 동치임을 보이면 충분하다((1.3.4) 참조). 다시 말해 국소환 $B$ 위의 두 유한 생성 가군 $P$, $Q$가 영가군이 아니면 $P\otimes_BQ\neq0$임을 보이면 된다. $B$의 극대 아이디얼을 $\mathfrak m$이라 하자. 나카야마 보조정리에 따라 벡터공간 $P/\mathfrak mP$와 $Q/\mathfrak mQ$는 영공간이 아니므로, 그 텐서곱
\[
  (P/\mathfrak mP)\otimes_{B/\mathfrak m}(Q/\mathfrak mQ)
  =(P\otimes_B Q)\otimes_B(B/\mathfrak m)
\]
도 영공간이 아니다. 이로부터 결론이 따른다.

특히 $M$이 유한 생성 $A$-가군이고 $\mathfrak a$가 $A$의 아이디얼이면, $\operatorname{Supp}(M/\mathfrak aM)$은 $\mathfrak a$와 $M$의 소멸자 $\mathfrak n$을 모두 포함하는 소 아이디얼들의 집합이며(1.7.4), 다시 말해 $\mathfrak a+\mathfrak n$을 포함하는 소 아이디얼들의 집합이다.
\end{env}

\oldpage[0\textsubscript{I}]{21}

\section{기약공간과 뇌터 공간}
\label{section:0.2-ko}

\subsection{기약공간}
\label{subsection:0.2.1-ko}

\begin{env}[2.1.1]
\label{0.2.1.1-ko}
위상공간 $X$가 공집합이 아니고 $X$ 자체와 다른 두 닫힌 부분공간의 합집합으로 나타나지 않으면 $X$를 \emph{기약}이라 한다. 이는 $X\neq\varnothing$이고 $X$의 공집합이 아닌 두 열린집합(따라서 유한 개의 열린집합)의 교집합이 공집합이 아니라고 하는 것, 또는 공집합이 아닌 모든 열린집합이 모든 곳에서 조밀하다고 하는 것, 또는 $X$와 다른 모든 닫힌 부분집합이 \emph{어디에서도 조밀하지 않다}고 하는 것, 또는 끝으로 $X$의 모든 열린집합이 \emph{연결}되어 있다고 하는 것과 동치이다.
\end{env}

\begin{env}[2.1.2]
\label{0.2.1.2-ko}
위상공간 $X$의 부분공간 $Y$가 기약이기 위한 필요충분조건은 그 폐포 $\overline Y$가 기약인 것이다. 특히 한 점만으로 이루어진 부분공간의 폐포 $\overline{\{x\}}$인 모든 부분공간은 기약이다. 관계
\[
  y\in\overline{\{x\}}
  \qquad\bigl(\text{동치로 }\overline{\{y\}}\subset\overline{\{x\}}\bigr)
\]
를 $y$가 $x$의 \emph{특수화}라고 하거나 $x$가 $y$의 \emph{일반화}라고 표현하겠다. 기약공간 $X$에 $X=\overline{\{x\}}$인 점 $x$가 존재하면 $x$를 $X$의 \emph{일반점}이라 한다. 이때 $X$의 공집합이 아닌 모든 열린집합은 $x$를 포함하고, $x$를 포함하는 모든 부분공간은 $x$를 일반점으로 갖는다.
\end{env}

\begin{env}[2.1.3]
\label{0.2.1.3-ko}
다음 분리 공리를 만족하는 위상공간 $X$를 \emph{콜모고로프 공간}이라 함을 상기하자.

\smallskip
\noindent
\textup{(T$_0$)} $X$의 서로 다른 두 점 $x\neq y$가 주어지면, $x$, $y$ 중 한 점을 포함하고 다른 점은 포함하지 않는 열린집합이 존재한다.

기약 콜모고로프 공간에 일반점이 존재하면 그것은 \emph{유일}하다. 실제로 공집합이 아닌 열린집합은 모든 일반점을 포함한다.

위상공간 $X$의 모든 열린 덮개에서 $X$의 유한 열린 덮개를 뽑아낼 수 있으면 $X$를 \emph{준콤팩트}라 함을 상기하자. 이는 공집합이 아닌 닫힌집합들의 모든 감소 여과족이 공집합이 아닌 교집합을 갖는다는 것과 동치이다. $X$가 준콤팩트 공간이면 $X$의 공집합이 아닌 모든 닫힌 부분집합 $A$는 공집합이 아닌 \emph{극소} 닫힌집합 $M$을 포함한다. 실제로 $A$의 공집합이 아닌 닫힌 부분집합들의 집합은 포함관계 $\supset$에 관하여 귀납적이다. 더 나아가 $X$가 콜모고로프 공간이면 $M$은 반드시 한 점만으로 이루어진다. 이러한 $M$을 관용적으로 \emph{닫힌점}이라고도 한다.
\end{env}

\begin{env}[2.1.4]
\label{0.2.1.4-ko}
기약공간 $X$에서 공집합이 아닌 모든 열린 부분공간 $U$는 기약이고, $X$가 일반점 $x$를 가지면 $x$는 $U$의 일반점이기도 하다.

공집합이 아닌 열린집합들로 이루어진 위상공간 $X$의 덮개 $(U_\alpha)$를 생각하자. 이때 지표집합도 공집합이 아니라고 한다. $X$가 기약이기 위한 필요충분조건은 모든 $\alpha$에 대하여 $U_\alpha$가 기약이고 임의의 $\alpha$, $\beta$에 대하여 $U_\alpha\cap U_\beta\neq\varnothing$인 것이다. 이 조건이 필요함은 자명하다. 충분함을 보이려면 $X$의 공집합이 아닌 열린집합 $V$에 대하여 모든 $\alpha$에서 $V\cap U_\alpha$가 공집합이 아님을 보이면 충분하다. 그러면 모든 $\alpha$에서 $V\cap U_\alpha$가 $U_\alpha$에 조밀하고, 따라서 $V$가 $X$에 조밀하기 때문이다. 그런데 $V\cap U_\gamma\neq\varnothing$인 지표 $\gamma$가 적어도 하나 존재하므로 $V\cap U_\gamma$는 $U_\gamma$에 조밀하다. 또한 모든 $\alpha$에 대하여 $U_\alpha\cap U_\gamma\neq\varnothing$이므로 $V\cap U_\alpha\cap U_\gamma\neq\varnothing$이다.
\end{env}

\oldpage[0\textsubscript{I}]{22}

\begin{env}[2.1.5]
\label{0.2.1.5-ko}
$X$를 기약공간, $f$를 $X$에서 위상공간 $Y$로 가는 연속사상이라 하자. 그러면 $f(X)$는 기약이고, $x$가 $X$의 일반점이면 $f(x)$는 $f(X)$의 일반점이며 따라서 $\overline{f(X)}$의 일반점이기도 하다. 특히 $Y$도 기약이고 유일한 일반점 $y$를 가지면, $f(X)$가 모든 곳에서 조밀하기 위한 필요충분조건은 $f(x)=y$인 것이다.
\end{env}

\begin{env}[2.1.6]
\label{0.2.1.6-ko}
위상공간 $X$의 모든 기약 부분공간은 \emph{극대} 기약 부분공간에 포함되고, 이러한 극대 기약 부분공간은 반드시 닫혀 있다. $X$의 극대 기약 부분공간들을 $X$의 \emph{기약 성분}이라 한다. $Z_1$, $Z_2$가 $X$의 서로 다른 두 기약 성분이면 $Z_1\cap Z_2$는 부분공간 $Z_1$, $Z_2$ 각각에서 닫혀 있고 \emph{어디에서도 조밀하지 않은} 집합이다. 특히 $X$의 한 기약 성분이 일반점을 가지면(2.1.2), 그 일반점은 다른 어떤 기약 성분에도 속할 수 없다. $X$가 유한 개의 기약 성분 $Z_i$ $(1\leq i\leq n)$만을 가지고 각 $i$에 대하여
\[
  U_i=Z_i\cap\complement\Bigl(\bigcup_{j\neq i}Z_j\Bigr)
\]
로 놓으면, $U_i$들은 열려 있고 기약이며 서로소이고 그 합집합은 $X$에 조밀하다.

$U$를 위상공간 $X$의 열린 부분집합이라 하자. $Z$가 $U$와 만나는 $X$의 기약 부분집합이면 $Z\cap U$는 $Z$에서 열려 있고 조밀하며 따라서 기약이다. 역으로 $Y$가 $U$의 닫힌 기약 부분집합이면, $X$에서의 $Y$의 폐포 $\overline Y$는 기약이고 $\overline Y\cap U=Y$이다. 따라서 $U$의 기약 성분들과 $U$와 만나는 $X$의 기약 성분들 사이에는 \emph{전단사 대응}이 있다.
\end{env}

\begin{env}[2.1.7]
\label{0.2.1.7-ko}
위상공간 $X$가 유한 개의 닫힌 기약 부분공간 $Y_i$의 합집합이면, $X$의 기약 성분들은 $Y_i$들의 집합에서 극대인 원소들이다. 실제로 $Z$가 $X$의 닫힌 기약 부분집합이면 $Z$는 $Z\cap Y_i$들의 합집합이고, 이로부터 $Z$가 어떤 $Y_i$에 포함되어야 함이 곧바로 따른다. $Y$를 위상공간 $X$의 부분공간이라 하고 $Y$가 유한 개의 기약 성분 $Y_i$ $(1\leq i\leq n)$만을 가진다고 하자. 그러면 $X$에서의 폐포 $\overline{Y_i}$들은 $\overline Y$의 기약 성분들이다.
\end{env}

\begin{env}[2.1.8]
\label{0.2.1.8-ko}
$Y$를 유일한 일반점 $y$를 갖는 기약공간이라 하고, $X$를 위상공간, $f$를 $X$에서 $Y$로 가는 연속사상이라 하자. 그러면 $f^{-1}(y)$와 만나는 $X$의 모든 기약 성분 $Z$에 대하여 $f(Z)$는 $Y$에 조밀하다. 그 역은 반드시 참인 것은 아니다. 그러나 $Z$가 일반점 $z$를 가지고 $f(Z)$가 $Y$에 조밀하면 반드시 $f(z)=y$이다(2.1.5). 또한 이때 $Z\cap f^{-1}(y)$는 $f^{-1}(y)$에서 $\{z\}$의 폐포이므로 기약이고, $z$를 포함하는 $f^{-1}(y)$의 모든 기약 부분집합은 반드시 $Z$에 포함되므로(2.1.6) $z$는 $Z\cap f^{-1}(y)$의 일반점이다. $f^{-1}(y)$의 모든 기약 성분은 $X$의 한 기약 성분에 포함되므로, $f^{-1}(y)$와 만나는 $X$의 모든 기약 성분 $Z$가 일반점을 가진다면 이 기약 성분들의 집합과 $f^{-1}(y)$의 기약 성분들 $Z\cap f^{-1}(y)$의 집합 사이에는 \emph{전단사 대응}이 있다. 이때 $Z$의 일반점과 $Z\cap f^{-1}(y)$의 일반점은 같다.
\end{env}

\oldpage[0\textsubscript{I}]{23}

\subsection{뇌터 공간}
\label{subsection:0.2.2-ko}

\begin{env}[2.2.1]
\label{0.2.2.1-ko}
위상공간 $X$의 열린집합들의 집합이 \emph{극대 조건}을 만족하면, 또는 이와 동치로 $X$의 닫힌집합들의 집합이 \emph{극소 조건}을 만족하면 $X$를 \emph{뇌터}라 한다. 모든 $x\in X$가 뇌터 부분공간인 근방을 가지면 $X$를 \emph{국소 뇌터}라 한다.
\end{env}

\begin{env}[2.2.2]
\label{0.2.2.2-ko}
$E$를 \emph{극소 조건}을 만족하는 순서집합이라 하고, $P$를 $E$의 원소들에 관한 다음 조건을 만족하는 성질이라 하자. $a\in E$에 대하여 모든 $x<a$에서 $P(x)$가 참이면 $P(a)$도 참이다. 그러면 \emph{모든 $x\in E$에 대하여 $P(x)$가 참이다}. 이를 ``뇌터 귀납법의 원리''라 한다. 실제로 $P(x)$가 거짓인 $x\in E$들의 집합을 $F$라 하자. $F$가 공집합이 아니면 극소 원소 $a$를 갖는다. 그러면 모든 $x<a$에 대하여 $P(x)$가 참이므로 $P(a)$도 참이어야 하며, 이는 모순이다.

특히 $E$가 \emph{뇌터 공간의 닫힌 부분집합들의 집합}인 경우에 이 원리를 적용하겠다.
\end{env}

\begin{env}[2.2.3]
\label{0.2.2.3-ko}
뇌터 공간의 모든 부분공간은 뇌터이다. 역으로 유한 개의 뇌터 부분공간의 합집합인 모든 위상공간은 뇌터이다.
\end{env}

\begin{env}[2.2.4]
\label{0.2.2.4-ko}
모든 뇌터 공간은 준콤팩트이다. 역으로 모든 열린집합이 준콤팩트인 위상공간은 뇌터이다.
\end{env}

\begin{env}[2.2.5]
\label{0.2.2.5-ko}
뇌터 공간은 유한 개의 기약 성분만을 가진다. 이는 뇌터 귀납법으로 알 수 있다.
\end{env}

\section{층에 관한 보충}
\label{section:0.3-ko}

\subsection{범주에 값을 갖는 층}
\label{subsection:0.3.1-ko}

\begin{env}[3.1.1]
\label{0.3.1.1-ko}
$\mathbf K$를 범주라 하고, $(A_\alpha)_{\alpha\in I}$와 $(A_{\alpha\beta})_{(\alpha,\beta)\in I\times I}$를 $A_{\beta\alpha}=A_{\alpha\beta}$를 만족하는 $\mathbf K$의 두 대상족이라 하자. 또한 $(\rho_{\alpha\beta})_{(\alpha,\beta)\in I\times I}$를 사상
\[
  \rho_{\alpha\beta}:A_\alpha\longrightarrow A_{\alpha\beta}
\]
들의 족이라 하자. $\mathbf K$의 대상 $A$와 사상족 $\rho_\alpha:A\to A_\alpha$로 이루어진 쌍이 다음 조건을 만족하면, 이 쌍을 족 $(A_\alpha)$, $(A_{\alpha\beta})$, $(\rho_{\alpha\beta})$가 정의하는 \emph{보편 문제의 해}라 한다. $\mathbf K$의 모든 대상 $B$에 대하여 각 $f\in\operatorname{Hom}(B,A)$에 족
\[
  (\rho_\alpha\circ f)\in
  \prod_\alpha\operatorname{Hom}(B,A_\alpha)
\]
를 대응시키는 사상은 $\operatorname{Hom}(B,A)$에서 모든 지표쌍 $(\alpha,\beta)$에 대하여
\[
  \rho_{\alpha\beta}\circ f_\alpha
  =\rho_{\beta\alpha}\circ f_\beta
\]
를 만족하는 족 $(f_\alpha)$들의 집합 위로 가는 \emph{전단사}이다. 이러한 해가 존재하면 동형을 제외하고 유일함은 곧바로 알 수 있다.
\end{env}

\begin{env}[3.1.2]
\label{0.3.1.2-ko}
위상공간 $X$ 위에서 범주 $\mathbf K$에 값을 갖는 \emph{준층} $U\mapsto\mathcal F(U)$의 정의는 상기하지 않겠다(G, I, 1.9). 이러한 준층이 다음 공리를 만족하면 \emph{$\mathbf K$에 값을 갖는 층}이라 한다.

\smallskip
\noindent
\textup{(F)} \emph{$X$의 열린집합 $U$를 $U$에 포함된 열린집합 $U_\alpha$들로 덮는 임의의 덮개 $(U_\alpha)$에 대하여, $\rho_\alpha$(각각 $\rho_{\alpha\beta}$)를 제한사상}
\[
  \mathcal F(U)\longrightarrow\mathcal F(U_\alpha)
  \qquad
  \bigl(\text{각각 }\mathcal F(U_\alpha)
  \longrightarrow\mathcal F(U_\alpha\cap U_\beta)\bigr)
\]

\oldpage[0\textsubscript{I}]{24}

\emph{이라 하면, $\mathcal F(U)$와 족 $(\rho_\alpha)$로 이루어진 쌍은 $(\mathcal F(U_\alpha))$, $(\mathcal F(U_\alpha\cap U_\beta))$, $(\rho_{\alpha\beta})$에 관한 보편 문제의 해이다(3.1.1).}%
\footnote{이는 여과되지 않아도 되는 일반적인 \emph{사영극한} 개념의 특수한 경우이다((T, I, 1.8) 및 서론에서 예고한 집필 중인 책 참조).}

이는 $\mathbf K$의 모든 대상 $T$에 대하여 족
\[
  U\longmapsto\operatorname{Hom}(T,\mathcal F(U))
\]
이 \emph{집합의 층}이라고 하는 것과 동치이다.
\end{env}

\begin{env}[3.1.3]
\label{0.3.1.3-ko}
$\mathbf K$가 ``사상을 갖는 구조의 한 종류'' $\Sigma$로 정의되는 범주라고 하자. 따라서 $\mathbf K$의 대상은 $\Sigma$형 구조가 주어진 집합들이고 사상은 $\Sigma$의 사상들이다. 더 나아가 범주 $\mathbf K$가 다음 조건을 만족한다고 하자.

\smallskip
\noindent
\textup{(E)} $(A,(\rho_\alpha))$가 족 $(A_\alpha)$, $(A_{\alpha\beta})$, $(\rho_{\alpha\beta})$에 관하여 \emph{범주 $\mathbf K$ 안에서} 보편 사상 문제의 해이면, 같은 족에 관하여 \emph{집합의 범주 안에서도} 보편 사상 문제의 해이다. 여기서는 $A$, $A_\alpha$, $A_{\alpha\beta}$를 집합으로, $\rho_\alpha$, $\rho_{\alpha\beta}$를 사상으로 본다.%
\footnote{이는 표준 함자 $\mathbf K\to(\mathrm{Ens})$가 여과될 필요가 없는 \emph{사영극한과 교환한다}는 뜻이기도 함을 증명할 수 있다.}

이 조건 아래에서 (F)는 $U\mapsto\mathcal F(U)$를 \emph{집합의 준층}으로 보았을 때 \emph{층}이 되게 한다. 또한 (F)에 따라 사상 $u:T\to\mathcal F(U)$가 $\mathbf K$의 사상이기 위한 필요충분조건은 각 합성 $\rho_\alpha\circ u$가 $T\to\mathcal F(U_\alpha)$인 $\mathbf K$의 사상인 것이다. 이는 $\mathcal F(U)$ 위의 $\Sigma$형 구조가 사상 $\rho_\alpha$들에 관한 \emph{시초 구조}라는 뜻이다. 역으로 $X$ 위에서 $\mathbf K$에 값을 갖는 준층 $U\mapsto\mathcal F(U)$가 \emph{집합의 층}이고 앞의 조건을 만족한다고 하자. 그러면 이 준층은 (F)를 만족하므로 \emph{$\mathbf K$에 값을 갖는 층}이다.
\end{env}

\begin{env}[3.1.4]
\label{0.3.1.4-ko}
$\Sigma$가 군 또는 환 구조의 종류이면, $\mathbf K$에 값을 갖는 준층 $U\mapsto\mathcal F(U)$가 \emph{집합의 층}이라는 사실만으로 곧바로 \emph{$\mathbf K$에 값을 갖는 층}, 즉 (G)의 의미에서 군의 층 또는 환의 층이 된다.%
\footnote{이는 범주 $\mathbf K$에서 집합의 사상으로서 \emph{전단사}인 모든 사상이 \emph{동형}이기 때문이다. 예를 들어 $\mathbf K$가 위상공간의 범주이면 더 이상 그렇지 않다.}
그러나 예를 들어 $\mathbf K$가 연속 준동형을 사상으로 하는 \emph{위상환}의 범주이면 더 이상 그렇지 않다. $\mathbf K$에 값을 갖는 층은 다음 조건을 만족하는 환의 층 $U\mapsto\mathcal F(U)$이다. 모든 열린집합 $U$와 $U$에 포함된 열린집합 $U_\alpha$들에 의한 모든 덮개에 대하여, 환 $\mathcal F(U)$의 위상은 준동형
\[
  \mathcal F(U)\longrightarrow\mathcal F(U_\alpha)
\]
들을 연속이 되게 하는 \emph{가장 약한} 위상이다. 이 경우 위상을 잊고 환의 층으로 본 $U\mapsto\mathcal F(U)$를 위상환의 층 $U\mapsto\mathcal F(U)$의 \emph{바탕} 층이라 한다. 따라서 위상환의 층 사이의 사상 $u_V:\mathcal F(V)\to\mathcal G(V)$($V$는 $X$의 임의의 열린집합)은 바탕 환의 층 사이의 준동형으로서, 모든 열린집합 $V\subset X$에 대하여 $u_V$가 \emph{연속}인 것들이다. 이를 바탕 환의 층 사이의 임의의 준동형과 구별하기 위해 위상환의 층 사이의 \emph{연속} 준동형이라 부르겠다. 위상공간의 층이나 위상군의 층에 대해서도 비슷한 정의와 규약을 사용한다.
\end{env}

\oldpage[0\textsubscript{I}]{25}

\begin{env}[3.1.5]
\label{0.3.1.5-ko}
임의의 범주 $\mathbf K$에 대하여 $\mathcal F$가 $X$ 위에서 $\mathbf K$에 값을 갖는 준층(각각 층)이고 $U$가 $X$의 열린집합이면, 열린집합 $V\subset U$에 대한 $\mathcal F(V)$들은 $\mathbf K$에 값을 갖는 준층(각각 층)을 이룬다. 이를 $U$ 위에서 $\mathcal F$가 \emph{유도하는} 준층(각각 층)이라 하고 $\mathcal F|U$로 나타낸다.

$X$ 위에서 $\mathbf K$에 값을 갖는 준층 사이의 모든 사상 $u:\mathcal F\to\mathcal G$에 대하여, $V\subset U$인 $u_V$들로 이루어진 사상 $\mathcal F|U\to\mathcal G|U$를 $u|U$로 나타낸다.
\end{env}

\begin{env}[3.1.6]
\label{0.3.1.6-ko}
이제 범주 $\mathbf K$가 \emph{귀납극한}을 갖는다고 하자(T, 1.8). 그러면 $X$ 위에서 $\mathbf K$에 값을 갖는 모든 준층(특히 모든 층) $\mathcal F$와 모든 $x\in X$에 대하여, $\mathbf K$의 대상인 \emph{줄기} $\mathcal F_x$를 다음 귀납극한으로 정의할 수 있다. $x$의 열린 근방 $U$들의 집합을 $\supset$에 관한 여과집합으로 보고, 사상 $\rho_U^V:\mathcal F(V)\to\mathcal F(U)$에 따른 $\mathcal F(U)$들의 귀납극한을 취한다. $u:\mathcal F\to\mathcal G$가 $\mathbf K$에 값을 갖는 준층 사이의 사상이면, 모든 $x\in X$에 대하여 사상
\[
  u_x:\mathcal F_x\longrightarrow\mathcal G_x
\]
를 $x$의 열린 근방들의 집합을 따라 $u_U:\mathcal F(U)\to\mathcal G(U)$들의 귀납극한으로 정의한다. 이렇게 하여 모든 $x\in X$에 대하여 $\mathcal F_x$는 $\mathcal F$에 관한 공변함자로 정의되며 $\mathbf K$에 값을 갖는다.

더 나아가 $\mathbf K$가 사상을 갖는 구조의 종류 $\Sigma$로 정의되는 경우, $\mathbf K$에 값을 갖는 \emph{층 $\mathcal F$}의 $U$ \emph{위의 단면}은 여전히 $\mathcal F(U)$의 원소를 뜻하며, 이때 $\mathcal F(U)$ 대신 $\Gamma(U,\mathcal F)$라 쓴다. $s\in\Gamma(U,\mathcal F)$이고 $V$가 $U$에 포함된 열린집합이면 $\rho_V^U(s)$ 대신 $s|V$라 쓴다. 모든 $x\in U$에 대하여 $\mathcal F_x$에서 $s$의 표준 상을 점 $x$에서 $s$의 \emph{싹}이라 하고 $s_x$로 나타낸다. 이 의미로는 \emph{표기 $s(x)$를 결코 사용하지 않겠다}. 이 표기는 이 논고에서 뒤에 다룰 특정 층에 관한 다른 개념을 위해 남겨 둔다(5.5.1).

$u:\mathcal F\to\mathcal G$가 $\mathbf K$에 값을 갖는 층 사이의 사상이면, 모든 $s\in\Gamma(U,\mathcal F)$에 대하여 $u_V(s)$ 대신 $u(s)$라 쓰겠다.

$\mathcal F$가 가환군, 환 또는 가군의 층이면 $\mathcal F_x\neq\{0\}$인 $x\in X$들의 집합을 $\mathcal F$의 \emph{지지집합}이라 하고 $\operatorname{Supp}(\mathcal F)$로 나타낸다. 이 집합은 $X$에서 반드시 닫혀 있지는 않다.

$\mathbf K$가 사상을 갖는 구조의 한 종류로 정의되는 경우, $\mathbf K$에 값을 갖는 층에 관해서는 \emph{``에탈레 공간''의 관점을 체계적으로 사용하지 않겠다}. 다시 말해 층을 위상공간으로, 심지어 그 줄기들의 합집합인 집합으로도 보지 않을 것이며, $X$ 위의 이러한 층 사이의 사상 $u:\mathcal F\to\mathcal G$를 위상공간 사이의 연속사상으로 보지도 않을 것이다.
\end{env}

\subsection{열린집합의 기저 위의 준층}
\label{subsection:0.3.2-ko}

\begin{env}[3.2.1]
\label{0.3.2.1-ko}
앞으로는 (일반화된, 즉 반드시 여과될 필요가 없는 준순서집합에 대응하는) \emph{사영극한}을 갖는 범주 $\mathbf K$만을 다루겠다((T, 1.8) 참조). $X$를 위상공간, $\mathfrak B$를 $X$의 위상의 열린집합 기저라 하자. 각 $U\in\mathfrak B$에 대응하는 대상 $\mathcal F(U)\in\mathbf K$의 족과, $U\subset V$인 $\mathfrak B$의 모든 원소쌍 $(U,V)$에 대하여 정의된 사상
\[
  \rho_U^V:\mathcal F(V)\longrightarrow\mathcal F(U)
\]
들의 족을 \emph{$\mathbf K$에 값을 갖는 $\mathfrak B$ 위의 준층}이라 한다.

\oldpage[0\textsubscript{I}]{26}

여기에는 $\rho_U^U=\text{항등사상}$이라는 조건과, $U\subset V\subset W$인 $\mathfrak B$의 원소 $U,V,W$에 대하여
\[
  \rho_U^W=\rho_U^V\circ\rho_V^W
\]
라는 조건을 부과한다. 이 준층에는 보통 의미에서 $\mathbf K$에 값을 갖는 준층 $U\mapsto\mathcal F'(U)$를 다음과 같이 대응시킬 수 있다. 모든 열린집합 $U$에 대하여
\[
  \mathcal F'(U)=\varprojlim_{V\subset U}\mathcal F(V)
\]
로 놓는다. 여기서 $V$는 $V\subset U$인 $V\in\mathfrak B$들의 $\subset$에 관한 순서집합을 움직이며, 이 순서집합은 일반적으로 \emph{여과되지 않는다}. 실제로 $\mathcal F(V)$들은 $V\subset W\subset U$, $V,W\in\mathfrak B$일 때의 $\rho_V^W$에 관하여 사영계를 이룬다. $U\subset U'$인 $X$의 두 열린집합 $U,U'$가 주어지면, $\rho'_U{}^{U'}$를 $V\subset U$에 관한 표준 사상 $\mathcal F'(U')\to\mathcal F(V)$들의 사영극한으로 정의한다. 다시 말해 이는 표준 사상 $\mathcal F'(U)\to\mathcal F(V)$들과 합성했을 때 표준 사상 $\mathcal F'(U')\to\mathcal F(V)$가 되는 유일한 사상
\[
  \mathcal F'(U')\longrightarrow\mathcal F'(U)
\]
이다. 그러면 $\rho'_U{}^{U'}$의 추이성은 곧바로 확인된다. 더 나아가 $U\in\mathfrak B$이면 표준 사상 $\mathcal F'(U)\to\mathcal F(U)$는 동형이므로 두 대상을 동일시할 수 있다.%
\footnote{$X$가 \emph{뇌터 공간}이면, $\mathbf K$가 \emph{유한} 사영계의 사영극한만을 갖는다고 가정해도 $\mathcal F'(U)$를 정의하고 이것이 보통 의미의 준층임을 보일 수 있다. 실제로 $U$가 $X$의 임의의 열린집합이면 $\mathfrak B$의 원소들로 이루어진 $U$의 \emph{유한} 덮개 $(V_i)$가 존재한다. 각 지표쌍 $(i,j)$에 대하여 $(V_{ijk})$를 $\mathfrak B$의 원소들로 이루어진 $V_i\cap V_j$의 유한 덮개라 하자. $I$를 지표 $i$와 삼중항 $(i,j,k)$들로 이루어진 집합으로 하고, 오직 $i>(i,j,k)$와 $j>(i,j,k)$라는 관계로 순서를 주자. 이때 $\mathcal F'(U)$를 $\mathcal F(V_i)$와 $\mathcal F(V_{ijk})$로 이루어진 계의 사영극한으로 취한다. 이것이 덮개 $(V_i)$와 $(V_{ijk})$에 의존하지 않으며 $U\mapsto\mathcal F'(U)$가 준층임은 쉽게 확인된다.}
\end{env}

\begin{env}[3.2.2]
\label{0.3.2.2-ko}
이렇게 정의한 준층 $\mathcal F'$가 \emph{층}이기 위한 필요충분조건은 $\mathfrak B$ 위의 준층 $\mathcal F$가 다음 조건을 만족하는 것이다.

\smallskip
\noindent
\textup{(F\textsubscript{0})} \emph{$U\in\mathfrak B$를 $U$에 포함되는 $U_\alpha\in\mathfrak B$들로 덮는 모든 덮개 $(U_\alpha)$와 모든 대상 $T\in\mathbf K$에 대하여, 각 $f\in\operatorname{Hom}(T,\mathcal F(U))$에 족}
\[
  (\rho_{U_\alpha}^{U}\circ f)
  \in\prod_\alpha\operatorname{Hom}(T,\mathcal F(U_\alpha))
\]
\emph{을 대응시키는 사상은 $\operatorname{Hom}(T,\mathcal F(U))$에서 다음 조건을 만족하는 족 $(f_\alpha)$들의 집합 위로 가는 전단사이다. 모든 지표쌍 $(\alpha,\beta)$와 $V\subset U_\alpha\cap U_\beta$인 모든 $V\in\mathfrak B$에 대하여}
\[
  \rho_V^{U_\alpha}\circ f_\alpha
  =\rho_V^{U_\beta}\circ f_\beta.
\]%
\footnote{이는 $\mathcal F(U)$와 $\rho_\alpha=\rho_{U_\alpha}^{U}$들로 이루어진 쌍이 (3.1.1)에서 다음 자료로 정의한 보편 문제의 해라는 뜻이기도 하다. $A_\alpha=\mathcal F(U_\alpha)$이고,
$A_{\alpha\beta}=\prod\mathcal F(V)$이다. 여기서 곱은 $V\subset U_\alpha\cap U_\beta$인 $V\in\mathfrak B$ 전체에 걸친다. 또한
$\rho_{\alpha\beta}=(\rho_V''):\mathcal F(U_\alpha)\to\prod\mathcal F(V)$이며, 이는 $V,V',W\in\mathfrak B$, $V\cup V'\subset U_\alpha\cap U_\beta$, $W\subset V\cap V'$일 때
$\rho_W^V\circ\rho_V''=\rho_W^{V'}\circ\rho_{V'}''$가 되도록 정의한다.}

이 조건이 필요함은 자명하다. 충분함을 보이기 위해 먼저 $\mathfrak B$에 \emph{포함되는} $X$의 위상의 두 번째 기저 $\mathfrak B'$를 생각하자. 부분족 $(\mathcal F(V))_{V\in\mathfrak B'}$에서 얻은 준층을 $\mathcal F''$라 하면 $\mathcal F''$는 $\mathcal F'$와 \emph{표준적으로 동형}임을 보이겠다. 우선 $V\in\mathfrak B'$, $V\subset U$에 관한 표준 사상 $\mathcal F'(U)\to\mathcal F(V)$들의 사영극한은 모든 열린집합 $U$에 대하여 사상 $\mathcal F'(U)\to\mathcal F''(U)$를 준다. $U\in\mathfrak B$이면 이 사상은 동형이다. 실제로 가정에 따라 $V\in\mathfrak B'$, $V\subset U$에 대한 표준 사상 $\mathcal F''(U)\to\mathcal F(V)$들은
\[
  \mathcal F''(U)\longrightarrow\mathcal F(U)
  \longrightarrow\mathcal F(V)
\]
로 분해되며, 이렇게 정의된 사상 $\mathcal F(U)\to\mathcal F''(U)$와 $\mathcal F''(U)\to\mathcal F(U)$의 두 합성이 항등사상임은 곧바로 알 수 있다. 이제 모든 열린집합 $U$에 대하여, $W\in\mathfrak B$, $W\subset U$일 때의 사상 $\mathcal F''(U)\to\mathcal F''(W)=\mathcal F(W)$들은 $\mathcal F(W)$들의 사영극한을 특징짓는 조건을 만족한다. 따라서 사영극한이 동형을 제외하고 유일하다는 사실에 의해 주장이 증명된다.

이제 $U$를 $X$의 임의의 열린집합, $(U_\alpha)$를 $U$에 포함되는 열린집합들에 의한 $U$의 덮개라 하고, $\mathfrak B'$를 다음 원소들로 이루어진 부분족이라 하자.

\oldpage[0\textsubscript{I}]{27}

그 원소들은 $\mathfrak B$에 속하며 적어도 하나의 $U_\alpha$에 포함된다. $\mathfrak B'$도 여전히 $X$의 위상의 기저임은 명백하다. 따라서 $\mathcal F'(U)$는 $V\in\mathfrak B'$, $V\subset U$인 $\mathcal F(V)$들의 사영극한이고, $\mathcal F''(U_\alpha)$는 $V\in\mathfrak B'$, $V\subset U_\alpha$인 $\mathcal F(V)$들의 사영극한이다. 그러므로 사영극한의 정의에 의해 공리 (F)가 곧바로 성립한다.

(F\textsubscript{0})가 성립할 때에는 용어를 남용하여 $\mathfrak B$ 위의 준층 $\mathcal F$를 \emph{층}이라 하겠다.
\end{env}

\begin{env}[3.2.3]
\label{0.3.2.3-ko}
$\mathcal F$, $\mathcal G$를 $\mathbf K$에 값을 갖는 $\mathfrak B$ 위의 두 준층이라 하자. \emph{사상} $u:\mathcal F\to\mathcal G$를 제한사상 $\rho_V^W$와의 통상적인 호환조건을 만족하는 사상
\[
  u_V:\mathcal F(V)\longrightarrow\mathcal G(V)
  \qquad(V\in\mathfrak B)
\]
들의 족 $(u_V)_{V\in\mathfrak B}$로 정의한다. (3.2.1)의 표기에 따라, $V\in\mathfrak B$, $V\subset U$인 $u_V$들의 사영극한을 $u_U'$로 취하면 대응하는 보통 준층 사이의 사상 $u':\mathcal F'\to\mathcal G'$를 얻는다. $\rho'_U{}^{U'}$와의 호환조건은 사영극한의 함자성에서 따른다.
\end{env}

\begin{env}[3.2.4]
\label{0.3.2.4-ko}
범주 $\mathbf K$가 귀납극한을 갖고 $\mathcal F$가 $\mathbf K$에 값을 갖는 $\mathfrak B$ 위의 준층이라고 하자. 모든 $x\in X$에 대하여 $\mathfrak B$에 속하는 $x$의 근방들은 $x$의 모든 근방들의 집합에서 $\supset$에 관한 공종집합을 이룬다. 따라서 $\mathcal F'$가 $\mathcal F$에 대응하는 보통 준층이면 줄기 $\mathcal F_x'$는
\[
  \varinjlim_{\mathfrak B}\mathcal F(V)
\]
와 같다. 여기서 극한은 $x$를 포함하는 $V\in\mathfrak B$들의 집합을 따라 취한다. $u:\mathcal F\to\mathcal G$가 $\mathbf K$에 값을 갖는 $\mathfrak B$ 위의 준층 사이의 사상이고 $u':\mathcal F'\to\mathcal G'$가 대응하는 보통 준층 사이의 사상이면, $u_x'$도 마찬가지로 $V\in\mathfrak B$, $x\in V$인 사상 $u_V:\mathcal F(V)\to\mathcal G(V)$들의 귀납극한이다.
\end{env}

\begin{env}[3.2.5]
\label{0.3.2.5-ko}
(3.2.1)의 일반적인 조건으로 돌아가자. $\mathcal F$가 $\mathbf K$에 값을 갖는 보통 \emph{층}이고 $\mathcal F_1$이 $\mathcal F$를 $\mathfrak B$에 제한하여 얻은 \emph{$\mathfrak B$ 위의 층}이면, (3.2.1)의 절차로 $\mathcal F_1$에서 얻는 보통 층 $\mathcal F_1'$은 조건 (F)와 사영극한의 유일성에 따라 $\mathcal F$와 표준적으로 동형이다. 보통 $\mathcal F$와 $\mathcal F_1'$을 동일시하겠다.

$\mathcal G$가 $X$ 위에서 $\mathbf K$에 값을 갖는 두 번째 보통 층이고 $u:\mathcal F\to\mathcal G$가 사상이라 하자. 앞의 설명에 따르면 \emph{$V\in\mathfrak B$인 $u_V:\mathcal F(V)\to\mathcal G(V)$들만} 주어도 $u$가 완전히 결정된다. 역으로 $V\in\mathfrak B$인 $u_V$들이 주어졌을 때, $V,W\in\mathfrak B$, $V\subset W$인 제한사상 $\rho_V^W$들과의 가환도만 확인하면, 모든 $V\in\mathfrak B$에 대하여 $u_V'=u_V$인 $\mathcal F$에서 $\mathcal G$로 가는 사상 $u'$가 유일하게 존재한다(3.2.3).
\end{env}

\begin{env}[3.2.6]
\label{0.3.2.6-ko}
계속해서 $\mathbf K$가 사영극한을 갖는다고 하자. 그러면 \emph{$\mathbf K$에 값을 갖는 $X$ 위의 층}들의 범주도 \emph{사영극한}을 갖는다. 실제로 $(\mathcal F_\lambda)$가 $X$ 위에서 $\mathbf K$에 값을 갖는 층들의 사영계이면
\[
  \mathcal F(U)=\varprojlim_\lambda\mathcal F_\lambda(U)
\]
들은 $\mathbf K$에 값을 갖는 준층을 정의한다. 공리 (F)는 사영극한의 추이성으로부터 성립하며, 이때 $\mathcal F$가 $\mathcal F_\lambda$들의 사영극한이라는 사실은 자명하다.

$\mathbf K$가 집합의 범주일 때, 다음 조건을 만족하는 모든 사영계 $(\mathcal H_\lambda)$에 대하여

\oldpage[0\textsubscript{I}]{28}

모든 $\lambda$에 대하여 $\mathcal H_\lambda$는 $\mathcal F_\lambda$의 \emph{부분층}이고, $\varprojlim_\lambda\mathcal H_\lambda$는 $\varprojlim_\lambda\mathcal F_\lambda$의 \emph{부분층}과 표준적으로 동일시된다. $\mathbf K$가 가환군의 범주이면 공변함자 $\varprojlim_\lambda\mathcal F_\lambda$는 \emph{가법적이고 왼쪽 완전}이다.
\end{env}

\subsection{층의 붙이기}
\label{subsection:0.3.3-ko}

\begin{env}[3.3.1]
\label{0.3.3.1-ko}
계속해서 범주 $\mathbf K$가 일반화된 사영극한을 갖는다고 하자. $X$를 위상공간, $\mathfrak U=(U_\lambda)_{\lambda\in L}$를 $X$의 열린 덮개라 하고, 각 $\lambda\in L$에 대하여 $\mathcal F_\lambda$를 $U_\lambda$ 위에서 $\mathbf K$에 값을 갖는 층이라 하자. 모든 지표쌍 $(\lambda,\mu)$에 대하여 \emph{동형}
\[
  \theta_{\lambda\mu}:
  \mathcal F_\mu|(U_\lambda\cap U_\mu)
  \xrightarrow{\sim}
  \mathcal F_\lambda|(U_\lambda\cap U_\mu)
\]
이 주어졌다고 하자. 더 나아가 모든 삼중항 $(\lambda,\mu,\nu)$에 대하여 $\theta'_{\lambda\mu}$, $\theta'_{\mu\nu}$, $\theta'_{\lambda\nu}$를 각각 $\theta_{\lambda\mu}$, $\theta_{\mu\nu}$, $\theta_{\lambda\nu}$의 $U_\lambda\cap U_\mu\cap U_\nu$로의 제한이라 할 때
\[
  \theta'_{\lambda\nu}
  =\theta'_{\lambda\mu}\circ\theta'_{\mu\nu}
\]
가 성립한다고 하자. 이를 $\theta_{\lambda\mu}$들에 관한 \emph{붙이기 조건}이라 한다. 그러면 $X$ 위에서 $\mathbf K$에 값을 갖는 층 $\mathcal F$와 각 $\lambda$에 대한 동형
\[
  \eta_\lambda:\mathcal F|U_\lambda
  \xrightarrow{\sim}\mathcal F_\lambda
\]
가 존재하여 다음 조건을 만족한다. 모든 지표쌍 $(\lambda,\mu)$에 대하여 $\eta'_\lambda$, $\eta'_\mu$를 각각 $\eta_\lambda$, $\eta_\mu$의 $U_\lambda\cap U_\mu$로의 제한이라 하면
\[
  \theta_{\lambda\mu}
  =\eta'_\lambda\circ{\eta'_\mu}^{-1}.
\]
더 나아가 이 조건들에 의해 $\mathcal F$와 $\eta_\lambda$들은 유일한 동형을 제외하고 결정된다. 유일성은 (3.2.5)에서 곧바로 따른다. $\mathcal F$의 존재를 보이기 위해, 적어도 하나의 $U_\lambda$에 포함되는 열린집합들로 이루어진 열린집합 기저를 $\mathfrak B$라 하자. 모든 $U\in\mathfrak B$에 대하여 $U\subset U_\lambda$인 $\lambda$ 하나를 골라, 그 $\mathcal F_\lambda(U)$ 가운데 하나를 힐베르트의 $\tau$ 함수로 선택한다. 이 대상을 $\mathcal F(U)$라 하면 $U\subset V$, $U,V\in\mathfrak B$인 $\rho_U^V$들은 $\theta_{\lambda\mu}$를 사용하여 명백한 방식으로 정의된다. 추이성 조건은 붙이기 조건에서 따르고, (F\textsubscript{0})도 곧바로 확인된다. 따라서 이렇게 정의한 $\mathfrak B$ 위의 준층은 층이며, 일반적인 절차 (3.2.1)에 의해 요구한 성질을 갖는 보통 층을 얻는다. 이 층도 $\mathcal F$라 쓰겠다. $\mathcal F$는 \emph{$\theta_{\lambda\mu}$들을 사용하여 $\mathcal F_\lambda$들을 붙여 얻었다}고 말하며, 보통 $\eta_\lambda$를 통해 $\mathcal F_\lambda$와 $\mathcal F|U_\lambda$를 동일시하겠다.

$X$ 위에서 $\mathbf K$에 값을 갖는 모든 층 $\mathcal F$는 다음과 같이 붙여 얻은 것으로 볼 수 있음이 명백하다. $X$의 임의의 열린 덮개 $(U_\lambda)$에 대하여 $\mathcal F_\lambda=\mathcal F|U_\lambda$로 놓고, $\theta_{\lambda\mu}$를 항등사상으로 취한다.
\end{env}

\begin{env}[3.3.2]
\label{0.3.3.2-ko}
같은 표기 아래에서, 모든 $\lambda\in L$에 대하여 $\mathcal G_\lambda$를 $U_\lambda$ 위에서 $\mathbf K$에 값을 갖는 두 번째 층이라 하자. 모든 지표쌍 $(\lambda,\mu)$에 대하여 붙이기 조건을 만족하는 동형
\[
  \omega_{\lambda\mu}:
  \mathcal G_\mu|(U_\lambda\cap U_\mu)
  \xrightarrow{\sim}
  \mathcal G_\lambda|(U_\lambda\cap U_\mu)
\]
이 주어졌다고 하자. 끝으로 모든 $\lambda$에 대하여 사상 $u_\lambda:\mathcal F_\lambda\to\mathcal G_\lambda$가 주어지고, 가환도
\begin{equation}
\begin{tikzcd}[column sep=huge,row sep=large]
\mathcal F_\mu|(U_\lambda\cap U_\mu)
  \arrow[r,"u_\mu"] \arrow[d]
  & \mathcal G_\mu|(U_\lambda\cap U_\mu) \arrow[d] \\
\mathcal F_\lambda|(U_\lambda\cap U_\mu)
  \arrow[r,"u_\lambda"']
  & \mathcal G_\lambda|(U_\lambda\cap U_\mu)
\end{tikzcd}
\tag{3.3.2.1}\label{0.3.3.2.1-ko}
\end{equation}
들이 가환이라고 하자. $\mathcal G$가 $\omega_{\lambda\mu}$들을 사용하여 $\mathcal G_\lambda$들을 붙여 얻은 층이면, 사상 $u:\mathcal F\to\mathcal G$가 유일하게 존재한다. 그 조건은 다음 가환도들이

\oldpage[0\textsubscript{I}]{29}

\[
\begin{tikzcd}[column sep=huge,row sep=large]
\mathcal F|U_\lambda
  \arrow[r,"u|U_\lambda"] \arrow[d]
  & \mathcal G|U_\lambda \arrow[d] \\
\mathcal F_\lambda
  \arrow[r,"u_\lambda"']
  & \mathcal G_\lambda
\end{tikzcd}
\]
가환인 것이다. 이는 (3.2.3)에서 곧바로 따른다. 족 $(u_\lambda)$와 $u$ 사이의 대응은 조건 (3.3.2.1)을 만족하는
\[
  \prod_\lambda
  \operatorname{Hom}(\mathcal F_\lambda,\mathcal G_\lambda)
\]
의 부분에서 $\operatorname{Hom}(\mathcal F,\mathcal G)$ 위로 가는 함자적 전단사이다.
\end{env}

\begin{env}[3.3.3]
\label{0.3.3.3-ko}
(3.3.1)의 표기에서 $V$를 $X$의 열린집합이라 하자. $\theta_{\lambda\mu}$들의 $V\cap U_\lambda\cap U_\mu$로의 제한들은 유도된 층 $\mathcal F_\lambda|(V\cap U_\lambda)$들에 관한 붙이기 조건을 만족한다. 이 층들을 붙여 얻은 $V$ 위의 층은 $\mathcal F|V$와 표준적으로 동일시됨이 곧바로 따른다.
\end{env}

\subsection{준층의 직상}
\label{subsection:0.3.4-ko}

\begin{env}[3.4.1]
\label{0.3.4.1-ko}
$X$, $Y$를 두 위상공간, $\psi:X\to Y$를 연속사상이라 하자. $\mathcal F$를 $X$ 위에서 범주 $\mathbf K$에 값을 갖는 준층이라 하자. 모든 열린집합 $U\subset Y$에 대하여
\[
  \mathcal G(U)=\mathcal F(\psi^{-1}(U))
\]
로 놓고, $U\subset V$인 $Y$의 두 열린 부분집합 $U,V$에 대하여 $\rho_U^V$를 사상
\[
  \mathcal F(\psi^{-1}(V))
  \longrightarrow\mathcal F(\psi^{-1}(U))
\]
이라 하자. $\mathcal G(U)$들과 $\rho_U^V$들은 $Y$ 위에서 $\mathbf K$에 값을 갖는 \emph{준층}을 정의함이 곧바로 따른다. 이를 \emph{$\psi$에 의한 $\mathcal F$의 직상}이라 하고 $\psi_*(\mathcal F)$로 나타낸다. $\mathcal F$가 층이면 준층 $\mathcal G$에 관한 공리 (F)가 곧바로 확인되므로 $\psi_*(\mathcal F)$도 층이다.
\end{env}

\begin{env}[3.4.2]
\label{0.3.4.2-ko}
$\mathcal F_1$, $\mathcal F_2$를 $X$ 위에서 $\mathbf K$에 값을 갖는 두 준층이라 하고, $u:\mathcal F_1\to\mathcal F_2$를 사상이라 하자. $U$가 $Y$의 열린 부분집합 전체를 움직일 때 사상족
\[
  u_{\psi^{-1}(U)}:
  \mathcal F_1(\psi^{-1}(U))
  \longrightarrow\mathcal F_2(\psi^{-1}(U))
\]
은 제한사상과의 호환조건을 만족하므로 사상
\[
  \psi_*(u):\psi_*(\mathcal F_1)
  \longrightarrow\psi_*(\mathcal F_2)
\]
를 정의한다. $\mathcal F_3$가 $X$ 위에서 $\mathbf K$에 값을 갖는 세 번째 준층이고 $v:\mathcal F_2\to\mathcal F_3$가 사상이면
\[
  \psi_*(v\circ u)=\psi_*(v)\circ\psi_*(u).
\]
다시 말해 $\psi_*(\mathcal F)$는 $\mathcal F$에 관한 \emph{공변함자}이며, $X$ 위에서 $\mathbf K$에 값을 갖는 준층(각각 층)의 범주에서 $Y$ 위에서 $\mathbf K$에 값을 갖는 준층(각각 층)의 범주로 간다.
\end{env}

\begin{env}[3.4.3]
\label{0.3.4.3-ko}
$Z$를 세 번째 위상공간, $\psi':Y\to Z$를 연속사상이라 하고 $\psi''=\psi'\circ\psi$로 놓자. $X$ 위에서 $\mathbf K$에 값을 갖는 모든 준층 $\mathcal F$에 대하여
\[
  \psi_*''(\mathcal F)
  =\psi_*'(\psi_*(\mathcal F))
\]
임은 명백하다. 또한 이러한 준층 사이의 모든 사상 $u:\mathcal F\to\mathcal G$에 대하여
\[
  \psi_*''(u)=\psi_*'(\psi_*(u)).
\]
다시 말해 $\psi_*''$은 함자 $\psi_*'$과 $\psi_*$의 \emph{합성}이며, 이를
\[
  (\psi'\circ\psi)_*=\psi_*'\circ\psi_*
\]
로 쓸 수 있다.

더 나아가 $Y$의 모든 열린집합 $U$에 대하여, 유도된 준층 $\mathcal F|\psi^{-1}(U)$의 제한사상 $\psi|\psi^{-1}(U)$에 의한 직상은 바로 유도된 준층 $\psi_*(\mathcal F)|U$이다.
\end{env}

\begin{env}[3.4.4]
\label{0.3.4.4-ko}
범주 $\mathbf K$가 귀납극한을 갖는다고 하고, $\mathcal F$를 $X$ 위에서 $\mathbf K$에 값을 갖는 준층이라 하자. 모든 $x\in X$에 대하여 사상
\[
  \Gamma(\psi^{-1}(U),\mathcal F)
  \longrightarrow\mathcal F_x
\]
들은 귀납계를 이룬다. 여기서 $U$는 $Y$에서 $\psi(x)$의 열린 근방을 움직인다. 이 귀납계에서

\oldpage[0\textsubscript{I}]{30}

귀납극한으로 넘어가면 줄기 사이의 사상
\[
  \psi_x:(\psi_*(\mathcal F))_{\psi(x)}
  \longrightarrow\mathcal F_x
\]
을 얻는다. 일반적으로 이 사상은 \emph{단사도 전사도 아니다}. 이 사상은 함자적이다. 실제로 $u:\mathcal F_1\to\mathcal F_2$가 $X$ 위에서 $\mathbf K$에 값을 갖는 준층 사이의 사상이면 가환도
\[
\begin{tikzcd}[column sep=huge,row sep=large]
(\psi_*(\mathcal F_1))_{\psi(x)}
  \arrow[r,"\psi_x"] \arrow[d,"{(\psi_*(u))_{\psi(x)}}"']
  & (\mathcal F_1)_x \arrow[d,"u_x"] \\
(\psi_*(\mathcal F_2))_{\psi(x)}
  \arrow[r,"\psi_x"']
  & (\mathcal F_2)_x
\end{tikzcd}
\]
는 가환이다. $Z$가 세 번째 위상공간이고 $\psi':Y\to Z$가 연속사상이며 $\psi''=\psi'\circ\psi$이면, 모든 $x\in X$에 대하여
\[
  \psi_x''=\psi_x\circ\psi'_{\psi(x)}.
\]
\end{env}

\begin{env}[3.4.5]
\label{0.3.4.5-ko}
(3.4.4)의 가정 아래에서, 더 나아가 $\psi$가 $X$에서 $Y$의 부분공간 $\psi(X)$ 위로 가는 \emph{위상동형}이라고 하자. 그러면 모든 $x\in X$에 대하여 $\psi_x$는 \emph{동형}이다. 이는 특히 $Y$의 부분집합 $X$에서 $Y$로 가는 표준 단사사상 $j$에 적용된다.
\end{env}

\begin{env}[3.4.6]
\label{0.3.4.6-ko}
$\mathbf K$가 군, 환 등의 범주라고 하자. $\mathcal F$가 $X$ 위에서 $\mathbf K$에 값을 갖고 지지집합이 $S$인 층이며 $y\notin\overline{\psi(S)}$이면, $\psi_*(\mathcal F)$의 정의로부터
\[
  (\psi_*(\mathcal F))_y=\{0\}
\]
을 얻는다. 다시 말해 $\psi_*(\mathcal F)$의 지지집합은 $\overline{\psi(S)}$에 포함되지만 반드시 $\psi(S)$에 포함되는 것은 아니다. 같은 가정 아래에서 $j$가 $Y$의 부분집합 $X$에서 $Y$로 가는 표준 단사사상이면, 층 $j_*(\mathcal F)$가 $X$ 위에 유도하는 층은 $\mathcal F$이다. 더 나아가 $X$가 $Y$에서 \emph{닫혀} 있으면 $j_*(\mathcal F)$는 $X$ 위에 $\mathcal F$를, $Y-X$ 위에 $0$을 유도하는 $Y$ 위의 층이다(G, II, 2.9.2). 그러나 $X$가 국소 닫혔지만 닫히지는 않았다고 가정하면 일반적으로 $j_*(\mathcal F)$는 이 마지막 층과 다르다.
\end{env}

\subsection{준층의 역상}
\label{subsection:0.3.5-ko}

\begin{env}[3.5.1]
\label{0.3.5.1-ko}
(3.4.1)의 가정 아래에서 $\mathcal F$가 $X$ 위의 준층이고 $\mathcal G$가 $Y$ 위의 준층이며 둘 다 $\mathbf K$에 값을 갖는다고 하자. $Y$ 위의 준층 사이의 모든 사상
\[
  u:\mathcal G\longrightarrow\psi_*(\mathcal F)
\]
을 역시 $\mathcal G$에서 $\mathcal F$로 가는 $\psi$-\emph{사상}이라 하고 $\mathcal G\to\mathcal F$로도 쓴다. 또한 $\mathcal G$에서 $\mathcal F$로 가는 $\psi$-사상들의 집합
\[
  \operatorname{Hom}_Y(\mathcal G,\psi_*(\mathcal F))
\]
을 $\operatorname{Hom}_\psi(\mathcal G,\mathcal F)$로 나타낸다.

$U$가 $X$의 열린집합이고 $V$가 $\psi(U)\subset V$를 만족하는 $Y$의 열린집합인 모든 쌍 $(U,V)$에 대하여, 제한사상
\[
  \mathcal F(\psi^{-1}(V))\longrightarrow\mathcal F(U)
\]
과 사상
\[
  u_V:\mathcal G(V)\longrightarrow
  \psi_*(\mathcal F)(V)=\mathcal F(\psi^{-1}(V))
\]
을 합성하여 사상 $u_{U,V}:\mathcal G(V)\to\mathcal F(U)$를 얻는다. 이 사상들은 $U'\subset U$, $V'\subset V$, $\psi(U')\subset V'$일 때 가환도
\begin{equation*}
\begin{tikzcd}[column sep=huge,row sep=large]
\mathcal G(V) \arrow[r,"u_{U,V}"] \arrow[d]
  & \mathcal F(U) \arrow[d] \\
\mathcal G(V') \arrow[r,"u_{U',V'}"']
  & \mathcal F(U')
\end{tikzcd}
\tag{3.5.1.1}\label{0.3.5.1.1-ko}
\end{equation*}
들을 가환이 되게 함이 곧바로 따른다. 역으로 가환도 (3.5.1.1)을 가환이 되게 하는 사상족 $(u_{U,V})$가 주어지면
\[
  u_V=u_{\psi^{-1}(V),V}
\]
로 취함으로써 $\psi$-사상 $u$가 정의된다.

\oldpage[0\textsubscript{I}]{31}

범주 $\mathbf K$가 일반화된 사영극한을 갖고 $\mathfrak B$, $\mathfrak B'$이 각각 $X$, $Y$의 위상의 기저라고 하자. 층 사이의 $\psi$-사상 $u$를 정의하려면 $U\in\mathfrak B$, $V\in\mathfrak B'$, $\psi(U)\subset V$인 $u_{U,V}$들만 주고, $U,U'\in\mathfrak B$, $V,V'\in\mathfrak B'$에 대하여 호환조건 (3.5.1.1)을 확인하면 된다. 실제로 모든 열린집합 $W\subset Y$에 대하여 $u_W$를 $V\in\mathfrak B'$, $V\subset W$, $U\in\mathfrak B$, $\psi(U)\subset V$인 $u_{U,V}$들의 사영극한으로 정의하면 된다.

범주 $\mathbf K$가 귀납극한을 가질 때에는 모든 $x\in X$와 $Y$에서 $\psi(x)$의 모든 열린 근방 $V$에 대하여 사상
\[
  \mathcal G(V)\longrightarrow
  \mathcal F(\psi^{-1}(V))\longrightarrow\mathcal F_x
\]
을 얻는다. 이 사상들은 귀납계를 이루며, 귀납극한으로 넘어가면 사상
\[
  \mathcal G_{\psi(x)}\longrightarrow\mathcal F_x
\]
을 얻는다.
\end{env}

\begin{env}[3.5.2]
\label{0.3.5.2-ko}
(3.4.3)의 가정 아래에서 $\mathcal F$, $\mathcal G$, $\mathcal H$를 각각 $X$, $Y$, $Z$ 위에서 $\mathbf K$에 값을 갖는 준층이라 하고,
\[
  u:\mathcal G\longrightarrow\psi_*(\mathcal F),
  \qquad
  v:\mathcal H\longrightarrow\psi'_*(\mathcal G)
\]
를 각각 $\psi$-사상과 $\psi'$-사상이라 하자. 이들로부터 $\psi''$-사상
\[
  w:\mathcal H\xrightarrow{v}\psi'_*(\mathcal G)
  \xrightarrow{\psi'_*(u)}
  \psi'_*(\psi_*(\mathcal F))=\psi''_*(\mathcal F)
\]
을 얻으며, 이를 정의에 따라 $u$와 $v$의 \emph{합성}이라 한다. 따라서 위상공간 $X$와 그 위에서 $\mathbf K$에 값을 갖는 준층 $\mathcal F$의 쌍 $(X,\mathcal F)$들을 하나의 \emph{범주}를 이루는 것으로 볼 수 있다. 여기서 사상
\[
  (\psi,\theta):(X,\mathcal F)\longrightarrow(Y,\mathcal G)
\]
은 연속사상 $\psi:X\to Y$와 $\psi$-사상 $\theta:\mathcal G\to\mathcal F$의 쌍이다.
\end{env}

\begin{env}[3.5.3]
\label{0.3.5.3-ko}
$\psi:X\to Y$를 연속사상, $\mathcal G$를 $Y$ 위에서 $\mathbf K$에 값을 갖는 \emph{준층}이라 하자. $\psi$에 의한 $\mathcal G$의 \emph{역상}을 쌍 $(\mathcal G',\rho)$로 정의한다. 여기서 $\mathcal G'$은 $X$ 위에서 $\mathbf K$에 값을 갖는 \emph{층}이고, $\rho:\mathcal G\to\mathcal G'$은 $\psi$-사상, 다시 말해 준층의 사상 $\mathcal G\to\psi_*(\mathcal G')$이며, 다음 조건을 만족한다. $X$ 위에서 $\mathbf K$에 값을 갖는 모든 \emph{층} $\mathcal F$에 대하여 $v$를 $\psi_*(v)\circ\rho$로 보내는 사상
\begin{equation*}
  \operatorname{Hom}_X(\mathcal G',\mathcal F)
  \longrightarrow \operatorname{Hom}_\psi(\mathcal G,\mathcal F)
  =\operatorname{Hom}_Y(\mathcal G,\psi_*(\mathcal F))
\tag{3.5.3.1}\label{0.3.5.3.1-ko}
\end{equation*}
이 \emph{전단사}이다. 이 사상은 $\mathcal F$에 관하여 함자적이므로 $\mathcal F$에 관한 함자들의 동형을 정의한다. 쌍 $(\mathcal G',\rho)$은 보편문제의 해이므로, 존재할 때 \emph{유일한 동형을 제외하고 결정된다}. 이때
\[
  \mathcal G'=\psi^*(\mathcal G),
  \qquad
  \rho=\rho_{\mathcal G}
\]
로 쓰고, 표현을 약간 남용하여 $\psi^*(\mathcal G)$를 $\psi$에 의한 $\mathcal G$의 \emph{역상층}이라 하겠다. 여기서는 $\psi^*(\mathcal G)$에 표준 $\psi$-사상
\[
  \rho_{\mathcal G}:\mathcal G\longrightarrow\psi^*(\mathcal G),
\]
즉 $Y$ 위의 준층 사이의 \emph{표준 사상}
\begin{equation*}
  \rho_{\mathcal G}:\mathcal G
  \longrightarrow\psi_*(\psi^*(\mathcal G))
\tag{3.5.3.2}\label{0.3.5.3.2-ko}
\end{equation*}
이 갖추어져 있다고 이해한다.

$v:\psi^*(\mathcal G)\to\mathcal F$가 $X$ 위에서 $\mathbf K$에 값을 갖는 층 사이의 사상이면
\[
  v^b=\psi_*(v)\circ\rho_{\mathcal G}:
  \mathcal G\longrightarrow\psi_*(\mathcal F)
\]
로 놓는다. 정의에 따라 준층 사이의 \emph{모든} 사상 $u:\mathcal G\to\psi_*(\mathcal F)$는 유일한 하나의 $v$에 대하여 $v^b$의 꼴이고, 이 $v$를 $u^\sharp$이라 쓴다. 다시 말해 준층 사이의 모든 사상 $u:\mathcal G\to\psi_*(\mathcal F)$는 유일하게
\begin{equation*}
  u:\mathcal G\xrightarrow{\rho_{\mathcal G}}
  \psi_*(\psi^*(\mathcal G))
  \xrightarrow{\psi_*(u^\sharp)}\psi_*(\mathcal F)
\tag{3.5.3.3}\label{0.3.5.3.3-ko}
\end{equation*}
로 분해된다.
\end{env}

\oldpage[0\textsubscript{I}]{32}

\begin{env}[3.5.4]
\label{0.3.5.4-ko}
이제 범주 $\mathbf K$가 다음 성질을 갖는다고 하자.\footnote{서론에서 인용한 책에서, $\mathbf K$에 값을 갖는 준층의 역상이 존재함을 보장하는 범주 $\mathbf K$에 대한 매우 일반적인 조건을 제시할 것이다.} $Y$ 위에서 $\mathbf K$에 값을 갖는 \emph{모든} 준층 $\mathcal G$는 $\psi$에 의한 역상을 가지며, 이를 $\psi^*(\mathcal G)$로 쓴다.

$\psi^*(\mathcal G)$를 $\mathcal G$에 관한 \emph{공변함자}, 즉 $Y$ 위에서 $\mathbf K$에 값을 갖는 준층의 범주에서 $X$ 위에서 $\mathbf K$에 값을 갖는 층의 범주로 가는 함자로 정의할 수 있음을 보이겠다. 이때 동형 $v\mapsto v^b$는 $\mathcal G$와 $\mathcal F$에 관한 \emph{쌍함자의 동형}
\begin{equation*}
  \operatorname{Hom}_X(\psi^*(\mathcal G),\mathcal F)
  \xrightarrow{\sim}
  \operatorname{Hom}_Y(\mathcal G,\psi_*(\mathcal F))
\tag{3.5.4.1}\label{0.3.5.4.1-ko}
\end{equation*}
이 된다.

실제로 $w:\mathcal G_1\to\mathcal G_2$가 $Y$ 위에서 $\mathbf K$에 값을 갖는 준층 사이의 사상이면 합성사상
\[
  \mathcal G_1\xrightarrow{w}\mathcal G_2
  \xrightarrow{\rho_{\mathcal G_2}}
  \psi_*(\psi^*(\mathcal G_2))
\]
을 생각하자. 이에 대응하는 사상
\[
  (\rho_{\mathcal G_2}\circ w)^\sharp:
  \psi^*(\mathcal G_1)\longrightarrow\psi^*(\mathcal G_2)
\]
을 $\psi^*(w)$라 쓰겠다. (3.5.3.3)에 따라
\begin{equation*}
  \psi_*(\psi^*(w))\circ\rho_{\mathcal G_1}
  =\rho_{\mathcal G_2}\circ w
\tag{3.5.4.2}\label{0.3.5.4.2-ko}
\end{equation*}
이다. 여기서 $\mathcal F$가 $X$ 위에서 $\mathbf K$에 값을 갖는 층이고 $u:\mathcal G_2\to\psi_*(\mathcal F)$가 사상이면, (3.5.3.3), (3.5.4.2), 그리고 $v^b$의 정의로부터
\[
  (u^\sharp\circ\psi^*(w))^b
  =\psi_*(u^\sharp)\circ\psi_*(\psi^*(w))\circ\rho_{\mathcal G_1}
  =\psi_*(u^\sharp)\circ\rho_{\mathcal G_2}\circ w
  =u\circ w,
\]
즉
\begin{equation*}
  (u\circ w)^\sharp=u^\sharp\circ\psi^*(w)
\tag{3.5.4.3}\label{0.3.5.4.3-ko}
\end{equation*}
을 얻는다. 특히 $u$를 사상
\[
  \mathcal G_2\xrightarrow{w'}\mathcal G_3
  \xrightarrow{\rho_{\mathcal G_3}}
  \psi_*(\psi^*(\mathcal G_3))
\]
으로 취하면
\[
  \psi^*(w'\circ w)
  =(\rho_{\mathcal G_3}\circ w'\circ w)^\sharp
  =(\rho_{\mathcal G_3}\circ w')^\sharp\circ\psi^*(w)
  =\psi^*(w')\circ\psi^*(w)
\]
을 얻는다. 이로써 주장이 증명된다.

끝으로 $X$ 위에서 $\mathbf K$에 값을 갖는 모든 층 $\mathcal F$에 대하여 $i_{\mathcal F}$를 $\psi_*(\mathcal F)$의 항등사상이라 하고,
\[
  \sigma_{\mathcal F}:
  \psi^*(\psi_*(\mathcal F))\longrightarrow\mathcal F
\]
를 사상 $(i_{\mathcal F})^\sharp$이라 쓰자. 식 (3.5.4.3)은 모든 사상 $u:\mathcal G\to\psi_*(\mathcal F)$에 대하여 특히 분해
\begin{equation*}
  u^\sharp:\psi^*(\mathcal G)
  \xrightarrow{\psi^*(u)}
  \psi^*(\psi_*(\mathcal F))
  \xrightarrow{\sigma_{\mathcal F}}\mathcal F
\tag{3.5.4.4}\label{0.3.5.4.4-ko}
\end{equation*}
를 준다. 사상 $\sigma_{\mathcal F}$를 \emph{표준적}이라고 하겠다.
\end{env}

\begin{env}[3.5.5]
\label{0.3.5.5-ko}
$\psi':Y\to Z$를 연속사상이라 하고, $Z$ 위에서 $\mathbf K$에 값을 갖는 모든 준층 $\mathcal H$가 $\psi'$에 의한 역상 $\psi'^*(\mathcal H)$을 갖는다고 하자. 그러면 (3.5.4)의 가정 아래에서 $Z$ 위에서 $\mathbf K$에 값을 갖는 모든 준층 $\mathcal H$는 $\psi''=\psi'\circ\psi$에 의한 역상을 가지며, 함자적인 표준 동형
\begin{equation*}
  \psi''^*(\mathcal H)\xrightarrow{\sim}
  \psi^*(\psi'^*(\mathcal H))
\tag{3.5.5.1}\label{0.3.5.5.1-ko}
\end{equation*}
이 존재한다.

\oldpage[0\textsubscript{I}]{33}

이는 $\psi''_*=\psi'_*\circ\psi_*$임을 고려하면 정의에서 곧바로 따른다. 또한 $u:\mathcal G\to\psi_*(\mathcal F)$가 $\psi$-사상이고, $v:\mathcal H\to\psi'_*(\mathcal G)$가 $\psi'$-사상이며, $w=\psi'_*(u)\circ v$가 이들의 합성이라 하자(3.5.2). 그러면 $w^\sharp$은 합성사상
\[
  w^\sharp:\psi^*(\psi'^*(\mathcal H))
  \xrightarrow{\psi^*(v^\sharp)}
  \psi^*(\mathcal G)
  \xrightarrow{u^\sharp}\mathcal F
\]
임이 곧바로 확인된다.
\end{env}

\begin{env}[3.5.6]
\label{0.3.5.6-ko}
특히 $\psi$를 항등사상 $1_X:X\to X$로 취하자. $X$ 위에서 $\mathbf K$에 값을 갖는 준층 $\mathcal F$의 $\psi$에 의한 역상이 존재하면, 이 역상을 \emph{준층 $\mathcal F$의 층화}라 한다. $\mathcal F$에서 $\mathbf K$에 값을 갖는 \emph{층} $\mathcal F'$으로 가는 모든 사상 $u:\mathcal F\to\mathcal F'$은 유일하게
\[
  \mathcal F\xrightarrow{\rho_{\mathcal F}}
  1_X^*(\mathcal F)\xrightarrow{u^\sharp}\mathcal F'
\]
으로 분해된다.
\end{env}

\subsection{상수층과 국소 상수층}
\label{subsection:0.3.6-ko}

\begin{env}[3.6.1]
\label{0.3.6.1-ko}
$X$ 위에서 $\mathbf K$에 값을 갖는 \emph{준층} $\mathcal F$를 생각하자. 모든 공집합이 아닌 열린집합 $U\subset X$에 대하여 표준 사상
\[
  \mathcal F(X)\longrightarrow\mathcal F(U)
\]
이 \emph{동형}이면 $\mathcal F$를 \emph{상수 준층}이라 하겠다. $\mathcal F$가 반드시 층인 것은 아님에 유의하자. 상수 준층의 층화인 \emph{층}을 \emph{상수층}이라 한다(원문의 «단순층»). 층 $\mathcal F$가 \emph{국소 상수}라는 것은 모든 $x\in X$가 $\mathcal F|U$가 상수층이 되는 열린 근방 $U$를 갖는다는 뜻이다.
\end{env}

\begin{env}[3.6.2]
\label{0.3.6.2-ko}
$X$가 \emph{기약}이라고 하자(2.1.1). 그러면 다음 성질들은 서로 동치이다.
\begin{enumerate}
  \item[a)] \emph{$\mathcal F$는 $X$ 위의 상수 준층이다.}
  \item[b)] \emph{$\mathcal F$는 $X$ 위의 상수층이다.}
  \item[c)] \emph{$\mathcal F$는 $X$ 위의 국소 상수층이다.}
\end{enumerate}

실제로 $\mathcal F$를 $X$ 위의 상수 준층이라 하자. $U$, $V$가 $X$의 공집합이 아닌 두 열린집합이면 $U\cap V$도 공집합이 아니다. 따라서
\[
  \mathcal F(X)\longrightarrow\mathcal F(U)
  \longrightarrow\mathcal F(U\cap V)
\]
와 $\mathcal F(X)\to\mathcal F(U)$가 동형이므로 $\mathcal F(U)\to\mathcal F(U\cap V)$도 동형이고, 마찬가지로 $\mathcal F(V)\to\mathcal F(U\cap V)$도 동형이다. 이로부터 (3.1.2)의 공리 (F)가 성립함이 곧바로 따른다. 따라서 $\mathcal F$는 그 층화와 동형이고, a)는 b)를 함의한다.

이제 $(U_\alpha)$를 공집합이 아닌 열린집합들로 이루어진 $X$의 열린 덮개라 하고, $\mathcal F$를 모든 $\alpha$에 대하여 $\mathcal F|U_\alpha$가 상수층인 $X$ 위의 층이라 하자. $U_\alpha$는 기약이므로 앞의 결과에 따라 $\mathcal F|U_\alpha$는 상수 준층이다. $U_\alpha\cap U_\beta$가 공집합이 아니므로
\[
  \mathcal F(U_\alpha)\longrightarrow
  \mathcal F(U_\alpha\cap U_\beta),
  \qquad
  \mathcal F(U_\beta)\longrightarrow
  \mathcal F(U_\alpha\cap U_\beta)
\]
는 동형이다. 따라서 모든 지표쌍에 대하여 표준 동형
\[
  \theta_{\alpha\beta}:
  \mathcal F(U_\alpha)\xrightarrow{\sim}\mathcal F(U_\beta)
\]
을 얻는다. 그런데 $U=X$에 조건 (F)를 적용하면, 모든 지표 $\alpha_0$에 대하여 $\mathcal F(U_{\alpha_0})$와 $\theta_{\alpha_0\alpha}$들이 보편문제의 해임을 알 수 있다. 유일성에 따라
\[
  \mathcal F(X)\longrightarrow\mathcal F(U_{\alpha_0})
\]
는 동형이다. 그러므로 c)는 a)를 함의한다.
\end{env}

\subsection{군 또는 환 준층의 역상}
\label{subsection:0.3.7-ko}

\oldpage[0\textsubscript{I}]{34}

\begin{env}[3.7.1]
\label{0.3.7.1-ko}
$\mathbf K$를 집합의 범주로 취할 때, $\mathbf K$에 값을 갖는 모든 준층 $\mathcal G$는 $\psi$에 의한 역상을 \emph{항상 가짐}을 보이겠다. $X$, $Y$, $\psi$에 관한 표기와 가정은 (3.5.3)의 것을 사용한다. 모든 열린집합 $U\subset X$에 대하여 $\mathcal G'(U)$를 다음과 같이 정의하자. $\mathcal G'(U)$의 원소 $s'$은 족 $(s'_x)_{x\in U}$이고, 모든 $x\in U$에 대하여
\[
  s'_x\in\mathcal G_{\psi(x)}
\]
이며 다음 조건을 만족한다. 모든 $x\in U$에 대하여 $Y$에서 $\psi(x)$의 열린 근방 $V$, $X$에서 $x$의 근방 $W\subset\psi^{-1}(V)\cap U$, 그리고 원소 $s\in\mathcal G(V)$가 존재하여 모든 $z\in W$에 대하여
\[
  s'_z=s_{\psi(z)}
\]
이다. $U\mapsto\mathcal G'(U)$가 \emph{층}의 공리들을 만족함은 곧바로 확인된다.

이제 $\mathcal F$를 $X$ 위의 집합의 층이라 하고
\[
  u:\mathcal G\longrightarrow\psi_*(\mathcal F),
  \qquad
  v:\mathcal G'\longrightarrow\mathcal F
\]
를 사상이라 하자. $u^\sharp$과 $v^b$를 다음과 같이 정의한다. $s'$이 $x\in X$의 근방 $U$ 위의 $\mathcal G'$의 단면이고, $V$가 $\psi(x)$의 열린 근방이며, $s\in\mathcal G(V)$가 $x$의 어떤 근방 $W\subset\psi^{-1}(V)\cap U$에서 모든 $z\in W$에 대하여 $s'_z=s_{\psi(z)}$를 만족한다고 하자. 이때
\[
  u_x^\sharp(s'_x)=u_{\psi(x)}(s_{\psi(x)})
\]
로 취한다. 마찬가지로 $V$가 $Y$의 열린집합이고 $s\in\mathcal G(V)$이면, $v^b(s)$는 $\psi^{-1}(V)$ 위의 $\mathcal G'$의 단면
\[
  s'=(s_{\psi(x)})_{x\in\psi^{-1}(V)}
\]
의 $v$에 의한 상인 $\psi^{-1}(V)$ 위의 $\mathcal F$의 단면이다.

더 나아가 표준 사상 (3.5.3)
\[
  \rho:\mathcal G\longrightarrow\psi_*(\psi^*(\mathcal G))
\]
은 다음과 같이 정의된다. 모든 열린집합 $V\subset Y$와 모든 단면 $s\in\Gamma(V,\mathcal G)$에 대하여 $\rho(s)$는 $\psi^{-1}(V)$ 위의 $\psi^*(\mathcal G)$의 단면
\[
  (s_{\psi(x)})_{x\in\psi^{-1}(V)}
\]
이다. 관계식
\[
  (u^\sharp)^b=u,\qquad
  (v^b)^\sharp=v,\qquad
  v^b=\psi_*(v)\circ\rho
\]
은 곧바로 확인되며, 이로써 주장이 증명된다.

$w:\mathcal G_1\to\mathcal G_2$가 $Y$ 위의 집합의 준층 사이의 사상이면 $\psi^*(w)$는 다음과 같이 구체적으로 주어진다. $s'=(s'_x)_{x\in U}$가 $X$의 열린집합 $U$ 위의 $\psi^*(\mathcal G_1)$의 단면이면
\[
  (\psi^*(w))(s')
  =(w_{\psi(x)}(s'_x))_{x\in U}.
\]
끝으로 $Y$의 모든 열린집합 $V$에 대하여, $\psi$의 $\psi^{-1}(V)$로의 제한에 의한 $\mathcal G|V$의 역상은 유도된 층
\[
  \psi^*(\mathcal G)|\psi^{-1}(V)
\]
과 동일하다.

$\psi$가 항등사상 $1_X$일 때에는 준층에 대응하는 집합의 층의 정의를 다시 얻는다(G, II, 1.2). 앞의 논의는 $\mathbf K$가 군 또는 반드시 가환일 필요는 없는 환의 범주일 때에도 아무런 변경 없이 적용된다.

$X$가 위상공간 $Y$의 임의의 부분집합이고 $j:X\to Y$가 표준 단사사상이라고 하자. 범주 $\mathbf K$에 값을 갖는 $Y$ 위의 모든 층 $\mathcal G$에 대하여, 역상 $j^*(\mathcal G)$이 존재할 때 이를 $\mathcal G$가 $X$ 위에 \emph{유도하는 층}이라 한다. 집합, 군 또는 환의 층에 대해서는 통상적인 정의를 다시 얻는다(G, II, 1.5).
\end{env}

\begin{env}[3.7.2]
\label{0.3.7.2-ko}
(3.5.3)의 표기와 가정을 유지하고, $\mathcal G$가 $Y$ 위의 군의 층(각각 환의 층)이라고 하자. $\psi^*(\mathcal G)$의 단면에 관한 정의 (3.7.1)와 (3.4.4)에 따르면 줄기 사상
\[
  \psi_x\circ\rho_{\psi(x)}:
  \mathcal G_{\psi(x)}
  \longrightarrow(\psi^*(\mathcal G))_x
\]
은 $\mathcal G$에 관하여 함자적인 \emph{동형}이다. 이 동형을 통해 두 줄기를 동일시할 수 있다. 이 동일시 아래에서 $u_x^\sharp$은 (3.5.1)에서 정의한 준동형과 같고, 특히
\[
  \operatorname{Supp}(\psi^*(\mathcal G))
  =\psi^{-1}(\operatorname{Supp}(\mathcal G))
\]
이다.

이 결과로부터 가환군의 층들의 아벨 범주에서 \emph{함자 $\psi^*(\mathcal G)$는 $\mathcal G$에 관하여 완전하다}는 것이 곧바로 따른다.
\end{env}

\subsection{의사이산 공간의 층}
\label{subsection:0.3.8-ko}

\oldpage[0\textsubscript{I}]{35}

\begin{env}[3.8.1]
\label{0.3.8.1-ko}
$X$를 위상이 \emph{준콤팩트} 열린집합들로 이루어진 기저 $\mathfrak B$를 갖는 위상공간이라 하자. $\mathcal F$를 $X$ 위의 \emph{집합의 층}이라 하자. 각 $\mathcal F(U)$에 \emph{이산위상}을 주면
\[
  U\longmapsto\mathcal F(U)
\]
는 \emph{위상공간의 준층}이 된다. 모든 \emph{준콤팩트} 열린집합 $U$에 대하여 $\Gamma(U,\mathcal F')$가 이산공간 $\mathcal F(U)$인, $\mathcal F$에 대응하는 \emph{위상공간의 층 $\mathcal F'$}이 존재함을 보이겠다(3.5.6).

이를 위해서는 이산 위상공간의 준층 $U\mapsto\mathcal F(U)$가 $\mathfrak B$ \emph{위에서} (3.2.2)의 조건 $(\mathrm F_0)$를 만족함을 보이면 충분하다. 더 일반적으로 $U$가 준콤팩트 열린집합이고 $(U_\alpha)$가 $\mathfrak B$의 원소들로 이루어진 $U$의 덮개이면, 사상
\[
  \Gamma(U,\mathcal F)\longrightarrow
  \Gamma(U_\alpha,\mathcal F)
\]
들을 연속이 되게 하는 $\Gamma(U,\mathcal F)$ 위의 가장 엉성한 위상 $\mathcal T$가 \emph{이산위상}임을 보이면 충분하다. 실제로
\[
  U=\bigcup_iU_{\alpha_i}
\]
가 되게 하는 유한개의 지표 $\alpha_i$가 존재한다. $s\in\Gamma(U,\mathcal F)$라 하고 $s_i$를 $\Gamma(U_{\alpha_i},\mathcal F)$에서의 그 상이라 하자. 집합 $\{s_i\}$들의 역상들의 교집합은 정의에 따라 $\mathcal T$에서 $s$의 근방이다. 그런데 $\mathcal F$는 집합의 층이고 $U_{\alpha_i}$들이 $U$를 덮으므로 이 교집합은 $\{s\}$와 같다. 이로써 주장이 증명된다.

$U$가 $X$의 준콤팩트하지 않은 열린집합이면 위상공간 $\Gamma(U,\mathcal F')$의 바탕집합은 여전히 $\Gamma(U,\mathcal F)$이지만, 그 위상은 일반적으로 이산위상이 아니다. 그 위상은 $V\in\mathfrak B$, $V\subset U$인 사상
\[
  \Gamma(U,\mathcal F)\longrightarrow\Gamma(V,\mathcal F)
\]
들을 연속이 되게 하는 가장 엉성한 위상이다. 여기서 각 $\Gamma(V,\mathcal F)$는 이산공간이다.

앞의 논의는 반드시 가환일 필요는 없는 군 또는 환의 층에도 아무런 변경 없이 적용되어, 각각 \emph{위상군} 또는 \emph{위상환}의 층을 대응시킨다. 줄여서 $\mathcal F'$을 집합(각각 군, 환)의 층 $\mathcal F$에 대응하는 \emph{의사이산 공간}(각각 \emph{의사이산 군}, \emph{의사이산 환})의 층이라 하겠다.
\end{env}

\begin{env}[3.8.2]
\label{0.3.8.2-ko}
$\mathcal F$, $\mathcal G$를 $X$ 위의 집합(각각 군, 환)의 두 층이라 하고 $u:\mathcal F\to\mathcal G$를 사상이라 하자. $\mathcal F'$, $\mathcal G'$을 각각 $\mathcal F$, $\mathcal G$에 대응하는 의사이산 층이라 하면, $u$는 \emph{연속인} 사상
\[
  \mathcal F'\longrightarrow\mathcal G'
\]
이기도 하다. 이는 (3.2.5)에서 따른다.
\end{env}

\begin{env}[3.8.3]
\label{0.3.8.3-ko}
$\mathcal F$를 집합의 층, $\mathcal H$를 $\mathcal F$의 부분층이라 하고, $\mathcal F'$, $\mathcal H'$을 각각 $\mathcal F$, $\mathcal H$에 대응하는 의사이산 층이라 하자. 그러면 모든 열린집합 $U\subset X$에 대하여
\[
  \Gamma(U,\mathcal H')
\]
은 $\Gamma(U,\mathcal F')$에서 \emph{닫혀} 있다. 실제로 이는 $V\in\mathfrak B$, $V\subset U$인 연속사상
\[
  \Gamma(U,\mathcal F)\longrightarrow\Gamma(V,\mathcal F)
\]
에 의한 $\Gamma(V,\mathcal H)$들의 역상들의 교집합이며, $\Gamma(V,\mathcal H)$는 이산공간 $\Gamma(V,\mathcal F)$에서 닫혀 있다.
\end{env}

\section{환 달린 공간}
\label{section:0.4-ko}

\subsection{환 달린 공간, $\mathcal A$-가군, $\mathcal A$-대수}
\label{subsection:0.4.1-ko}

\begin{env}[4.1.1]
\label{0.4.1.1-ko}
\emph{환 달린 공간}(각각 \emph{위상환 달린 공간})은 위상공간 $X$와 $X$ 위의 반드시 가환일 필요는 없는 환의 층(각각 위상환의 층) $\mathcal A$로 이루어진 쌍 $(X,\mathcal A)$이다. $X$를 환 달린 공간 $(X,\mathcal A)$의 바탕
\oldpage[0\textsubscript{I}]{36}
위상공간이라 하고, $\mathcal A$를 \emph{구조층}이라 한다. 구조층은 $\mathcal O_X$로 나타내며, 한 점 $x\in X$에서의 그 줄기는 $\mathcal O_{X,x}$ 또는 혼동할 염려가 없을 때에는 간단히 $\mathcal O_x$로 나타낸다.

$X$ 위의 $\mathcal O_X$의 \emph{단위원 단면}, 곧 $\Gamma(X,\mathcal O_X)$의 단위원은 $1$ 또는 $e$로 나타낸다.

이 논고에서는 주로 \emph{가환}환의 층을 다룰 것이므로, 별도의 명시 없이 환 달린 공간 $(X,\mathcal A)$을 말할 때에는 $\mathcal A$가 가환환의 층임을 전제로 한다.

구조층이 반드시 가환일 필요가 없는 환 달린 공간들(각각 위상환 달린 공간들)은 다음과 같이 \emph{범주}를 이룬다. \emph{사상}
\[
  (X,\mathcal A)\longrightarrow(Y,\mathcal B)
\]
을 연속사상 $\psi:X\to Y$와 환의 층(각각 위상환의 층)의 $\psi$-\emph{사상} $\theta:\mathcal B\to\mathcal A$ (3.5.1)로 이루어진 쌍 $(\psi,\theta)=\Psi$로 정의한다. 두 번째 사상
\[
  \Psi'=(\psi',\theta'):(Y,\mathcal B)\longrightarrow(Z,\mathcal C)
\]
과 $\Psi$의 \emph{합성}은 $\Psi''=\Psi'\circ\Psi$로 나타내며, 이는 $\psi''=\psi'\circ\psi$이고 $\theta''$가 $\theta$와 $\theta'$의 합성인 사상 $(\psi'',\theta'')$이다. 이때 $\theta''=\psi'_*(\theta)\circ\theta'$이다(3.5.2 참조). 환 달린 공간에 대해서는
\[
  \theta''^\sharp
  =\theta^\sharp\circ\psi^*(\theta'^\sharp)
\]
임을 상기하자(3.5.5). 따라서 $\theta'^\sharp$과 $\theta^\sharp$이 \emph{단사}(각각 \emph{전사}) 준동형이면 $\theta''^\sharp$도 그러하다. 여기서는 모든 $x\in X$에 대하여 $\psi_x\circ\rho_{\psi(x)}$가 동형이라는 사실을 함께 사용한다(3.7.2). 앞의 사실들로부터 $\psi$가 \emph{단사} 연속사상이고 $\theta^\sharp$이 환의 층의 \emph{전사} 준동형이면 사상 $(\psi,\theta)$가 환 달린 공간의 범주에서 \emph{단사사상}(T, 1.1)임을 곧바로 확인할 수 있다.

혼동할 염려가 없을 때에는 표기에서 흔히 $\psi$를 $\Psi$로 바꾸어 쓰겠다. 예를 들어 $Y$의 부분집합 $U$에 대하여 $\psi^{-1}(U)$ 대신 $\Psi^{-1}(U)$라고 쓴다.
\end{env}

\begin{env}[4.1.2]
\label{0.4.1.2-ko}
$X$의 임의의 부분집합 $M$에 대하여 쌍 $(M,\mathcal A|M)$은 자명하게 환 달린 공간이다. 이를 환 달린 공간 $(X,\mathcal A)$이 $M$ 위에 \emph{유도하는} 환 달린 공간이라 하며, $(X,\mathcal A)$의 $M$으로의 \emph{제한}이라고도 한다. $j:M\to X$가 표준 단사사상이고 $\omega$가 $\mathcal A|M$의 항등사상이면,
\[
  (j,\omega^b):(M,\mathcal A|M)\longrightarrow(X,\mathcal A)
\]
은 환 달린 공간의 단사사상이며 이를 \emph{표준 단사사상}이라 한다. 사상 $\Psi:(X,\mathcal A)\to(Y,\mathcal B)$와 이 표준 단사사상의 합성을 $\Psi$의 $M$으로의 \emph{제한}이라 한다.
\end{env}

\begin{env}[4.1.3]
\label{0.4.1.3-ko}
환 달린 공간 $(X,\mathcal A)$ 위의 $\mathcal A$-가군 또는 대수적 층의 정의(G, II, 2.2)는 다시 설명하지 않겠다. $\mathcal A$가 반드시 가환일 필요가 없는 환의 층이면, 반대의 것을 명시하지 않는 한 $\mathcal A$-가군은 항상 ``왼쪽 $\mathcal A$-가군''을 뜻한다. $\mathcal A$의 부분-$\mathcal A$-가군은 $\mathcal A$ 안의 (왼쪽, 오른쪽 또는 양쪽) \emph{아이디얼층} 또는 $\mathcal A$-아이디얼이라 한다.

$\mathcal A$가 가환환의 층일 때 $\mathcal A$-가군의 정의에서 가군 구조를 모두 대수 구조로 바꾸면 $X$ 위의 $\mathcal A$-대수의 정의를 얻는다. 이는 반드시 가환일 필요는 없는 $\mathcal A$-대수가 $\mathcal A$-가군 $\mathcal C$에 $\mathcal A$-가군 준동형
\[
  \varphi:
  \mathcal C\otimes_{\mathcal A}\mathcal C
  \longrightarrow\mathcal C
\]
과 $X$ 위의 단면 $e$를 주어 다음 조건들을 만족시키는 것과 같다. 1\textsuperscript{o} 도표
\oldpage[0\textsubscript{I}]{37}
\[
\begin{tikzcd}[column sep=large,row sep=large]
\mathcal{C}\otimes_{\mathcal{A}}\mathcal{C}\otimes_{\mathcal{A}}\mathcal{C}
  \arrow[r,"\varphi\otimes 1"]
  \arrow[d,"1\otimes\varphi"']
&
\mathcal{C}\otimes_{\mathcal{A}}\mathcal{C}
  \arrow[d,"\varphi"]
\\
\mathcal{C}\otimes_{\mathcal{A}}\mathcal{C}
  \arrow[r,"\varphi"']
&
\mathcal{C}
\end{tikzcd}
\]
가 가환이다. 2\textsuperscript{o} 모든 열린집합 $U\subset X$와 모든 단면 $s\in\Gamma(U,\mathcal C)$에 대하여
\[
  \varphi((e|U)\otimes s)
  =\varphi(s\otimes(e|U))
  =s
\]
이다. $\mathcal C$가 가환 $\mathcal A$-대수라는 것은 이에 더하여 도표
\[
\begin{tikzcd}[column sep=large,row sep=large]
\mathcal{C}\otimes_{\mathcal{A}}\mathcal{C}
  \arrow[rr,"\sigma"]
  \arrow[dr,"\varphi"']
&&
\mathcal{C}\otimes_{\mathcal{A}}\mathcal{C}
  \arrow[dl,"\varphi"]
\\
& \mathcal{C} &
\end{tikzcd}
\]
가 가환이라는 뜻이다. 여기서 $\sigma$는 텐서곱 $\mathcal C\otimes_{\mathcal A}\mathcal C$의 표준 대칭이다.

$\mathcal A$-대수의 준동형도 (G, II, 2.2)의 $\mathcal A$-가군 준동형과 같은 방식으로 정의하지만, 이제 그 준동형들은 당연히 아벨군을 이루지 않는다.

$\mathcal M$이 $\mathcal A$-대수 $\mathcal C$의 부분-$\mathcal A$-가군이면, $\mathcal M$이 생성하는 $\mathcal C$의 부분-$\mathcal A$-대수는 $n\geqslant0$에 대한 준동형
\[
  \bigotimes^n\mathcal M\longrightarrow\mathcal C
\]
들의 상들의 합이다. 이는 대수의 준층 $U\mapsto\mathcal B(U)$의 층화이기도 하다. 여기서 $\mathcal B(U)$는 부분가군 $\Gamma(U,\mathcal M)$이 생성하는 $\Gamma(U,\mathcal C)$의 부분대수이다.
\end{env}

\begin{env}[4.1.4]
\label{0.4.1.4-ko}
위상공간 $X$ 위의 환의 층 $\mathcal A$에 대하여 줄기 $\mathcal A_x$가 축소환(1.1.1)이면 $\mathcal A$가 $X$의 점 $x$에서 \emph{축소}라고 한다. $\mathcal A$가 $X$의 모든 점에서 축소이면 $\mathcal A$를 \emph{축소}라고 한다. 환 $A$의 모든 소 아이디얼 $\mathfrak p$에 대하여 국소화환 $A_{\mathfrak p}$가 정칙 국소환이면 $A$를 \emph{정칙환}이라 한다는 것을 상기하자. $\mathcal A_x$가 정칙환이면 $\mathcal A$가 점 $x$에서 \emph{정칙}이라고 하고, 모든 점에서 정칙이면 $\mathcal A$를 \emph{정칙}이라고 한다. 끝으로 $\mathcal A_x$가 정수적으로 닫힌 정역이면 $\mathcal A$가 점 $x$에서 \emph{정규}라고 하고, 모든 점에서 정규이면 $\mathcal A$를 \emph{정규}라고 한다. 환 달린 공간 $(X,\mathcal A)$이 앞의 성질 가운데 하나를 갖는다는 것은 환의 층 $\mathcal A$가 그 성질을 갖는다는 뜻이다.

\emph{등급환의 층} $\mathcal A$는 정의에 따라 아벨군의 층들의 족 $(\mathcal A_n)_{n\in\mathbf Z}$의 직합(G, II, 2.7)인 환의 층으로서
\[
  \mathcal A_m\mathcal A_n
  \subset\mathcal A_{m+n}
\]
을 만족하는 것이다. \emph{등급 $\mathcal A$-가군}은 아벨군의 층들의 족 $(\mathcal F_n)_{n\in\mathbf Z}$의 직합인 $\mathcal A$-가군 $\mathcal F$로서
\[
  \mathcal A_m\mathcal F_n
  \subset\mathcal F_{m+n}
\]
을 만족하는 것이다. 이는 모든 점 $x$에서
\[
  (\mathcal A_m)_x(\mathcal A_n)_x
  \subset(\mathcal A_{m+n})_x
\]
(각각
\[
  (\mathcal A_m)_x(\mathcal F_n)_x
  \subset(\mathcal F_{m+n})_x
\]
)가 성립한다는 것과 같다.
\end{env}

\begin{env}[4.1.5]
\label{0.4.1.5-ko}
반드시 가환일 필요가 없는 환 달린 공간 $(X,\mathcal A)$을 생각하자. 경우에 따라 왼쪽 또는 오른쪽 $\mathcal A$-가군의 범주에서 정의되고 아벨군의 층의 범주(또는 더
\oldpage[0\textsubscript{I}]{38}
일반적으로 $\mathcal C$가 $\mathcal A$의 중심이면 $\mathcal C$-가군의 범주)에 값을 갖는 쌍함자
\[
  \mathcal F\otimes_{\mathcal A}\mathcal G,
  \qquad
  \mathcal H\!om_{\mathcal A}(\mathcal F,\mathcal G),
  \qquad
  \operatorname{Hom}_{\mathcal A}(\mathcal F,\mathcal G)
\]
의 정의(G, II, 2.8 및 2.2)는 여기서 되풀이하지 않겠다. 모든 점 $x\in X$에서 줄기
\[
  (\mathcal F\otimes_{\mathcal A}\mathcal G)_x
\]
는 표준적으로 $\mathcal F_x\otimes_{\mathcal A_x}\mathcal G_x$와 동일시되며, 함자적인 표준 준동형
\[
  (\mathcal H\!om_{\mathcal A}(\mathcal F,\mathcal G))_x
  \longrightarrow
  \operatorname{Hom}_{\mathcal A_x}(\mathcal F_x,\mathcal G_x)
\]
이 정의된다. 이 준동형은 일반적으로 단사도 전사도 아니다.

앞의 쌍함자들은 가법적이며, 특히 유한 직합과 가환한다. $\mathcal F\otimes_{\mathcal A}\mathcal G$는 $\mathcal F$와 $\mathcal G$ 각각에 관하여 오른쪽 완전하고 귀납극한과 가환하며,
\[
  \mathcal A\otimes_{\mathcal A}\mathcal G
  \simeq\mathcal G,
  \qquad
  \mathcal F\otimes_{\mathcal A}\mathcal A
  \simeq\mathcal F
\]
인 표준 동형이 있다. 함자 $\mathcal H\!om_{\mathcal A}(\mathcal F,\mathcal G)$와 $\operatorname{Hom}_{\mathcal A}(\mathcal F,\mathcal G)$은 $\mathcal F$와 $\mathcal G$에 관하여 \emph{왼쪽 완전}하다. 더 정확히 말하면 완전열
\[
  0\longrightarrow\mathcal G'
  \longrightarrow\mathcal G
  \longrightarrow\mathcal G''
\]
이 주어질 때 열
\[
  0\longrightarrow
  \mathcal H\!om_{\mathcal A}(\mathcal F,\mathcal G')
  \longrightarrow
  \mathcal H\!om_{\mathcal A}(\mathcal F,\mathcal G)
  \longrightarrow
  \mathcal H\!om_{\mathcal A}(\mathcal F,\mathcal G'')
\]
은 완전하고, 완전열
\[
  \mathcal F'\longrightarrow\mathcal F
  \longrightarrow\mathcal F''\longrightarrow0
\]
이 주어질 때 열
\[
  0\longrightarrow
  \mathcal H\!om_{\mathcal A}(\mathcal F'',\mathcal G)
  \longrightarrow
  \mathcal H\!om_{\mathcal A}(\mathcal F,\mathcal G)
  \longrightarrow
  \mathcal H\!om_{\mathcal A}(\mathcal F',\mathcal G)
\]
은 완전하다. $\operatorname{Hom}$ 함자에도 같은 성질들이 성립한다. 또한 $\mathcal H\!om_{\mathcal A}(\mathcal A,\mathcal G)$는 표준적으로 $\mathcal G$와 동일시되며, 모든 열린집합 $U\subset X$에 대하여
\[
  \Gamma\bigl(U,\mathcal H\!om_{\mathcal A}(\mathcal F,\mathcal G)\bigr)
  =
  \operatorname{Hom}_{\mathcal A|U}(\mathcal F|U,\mathcal G|U)
\]
이다.

모든 왼쪽(각각 오른쪽) $\mathcal A$-가군 $\mathcal F$에 대하여 오른쪽(각각 왼쪽) $\mathcal A$-가군
\[
  \check{\mathcal F}
  =\mathcal H\!om_{\mathcal A}(\mathcal F,\mathcal A)
\]
를 $\mathcal F$의 \emph{쌍대}라 한다.

끝으로 $\mathcal A$가 가환환의 층이고 $\mathcal F$가 $\mathcal A$-가군이면
\[
  U\longmapsto\bigwedge^p\Gamma(U,\mathcal F)
\]
는 준층을 이룬다. 이 준층의 층화는 $\bigwedge^p\mathcal F$로 나타내는 $\mathcal A$-가군이며, 이를 $\mathcal F$의 \emph{$p$차 외승}이라 한다. 준층 $U\mapsto\bigwedge^p\Gamma(U,\mathcal F)$에서 그 층화 $\bigwedge^p\mathcal F$로 가는 표준 사상은 \emph{단사}이고, 모든 $x\in X$에 대하여
\[
  (\bigwedge^p\mathcal F)_x
  =\bigwedge^p(\mathcal F_x)
\]
이다. $\bigwedge^p\mathcal F$가 $\mathcal F$에 관한 공변함자임은 자명하다.
\end{env}

\begin{env}[4.1.6]
\label{0.4.1.6-ko}
$\mathcal A$를 반드시 가환일 필요는 없는 환의 층, $\mathcal J$를 $\mathcal A$의 왼쪽 아이디얼층, $\mathcal F$를 왼쪽 $\mathcal A$-가군이라 하자. 이때 $\mathcal J\mathcal F$는 표준 사상
\[
  \mathcal J\otimes_{\mathbf Z}\mathcal F
  \longrightarrow\mathcal F
\]
의 상인 $\mathcal F$의 부분-$\mathcal A$-가군을 나타낸다. 여기서 $\mathbf Z$는 상수 준층 $U\mapsto\mathbf Z$의 층화이다. 모든 $x\in X$에 대하여
\[
  (\mathcal J\mathcal F)_x
  =\mathcal J_x\mathcal F_x
\]
임은 자명하다. $\mathcal A$가 가환이면 $\mathcal J\mathcal F$는 표준 사상
\[
  \mathcal J\otimes_{\mathcal A}\mathcal F
  \longrightarrow\mathcal F
\]
의 상이기도 하다. 또한 $\mathcal J\mathcal F$는 준층
\[
  U\longmapsto
  \Gamma(U,\mathcal J)\Gamma(U,\mathcal F)
\]
의 층화이다. $\mathcal J_1$, $\mathcal J_2$가 $\mathcal A$의 두 왼쪽 아이디얼층이면
\[
  \mathcal J_1(\mathcal J_2\mathcal F)
  =(\mathcal J_1\mathcal J_2)\mathcal F
\]
이다.
\end{env}

\begin{env}[4.1.7]
\label{0.4.1.7-ko}
$(X_\lambda,\mathcal A_\lambda)_{\lambda\in L}$를 환 달린 공간들의 족이라 하자. 각 $(\lambda,\mu)$에 대하여 $X_\lambda$의 열린부분집합 $V_{\lambda\mu}$와 환 달린 공간의 동형사상
\[
  \varphi_{\lambda\mu}:
  (V_{\mu\lambda},\mathcal A_\mu|V_{\mu\lambda})
  \xrightarrow{\sim}
  (V_{\lambda\mu},\mathcal A_\lambda|V_{\lambda\mu})
\]
이 주어졌다고 하자. 여기서 $V_{\lambda\lambda}=X_\lambda$이고 $\varphi_{\lambda\lambda}$는 항등사상이다. 더 나아가 모든 세쌍 $(\lambda,\mu,\nu)$에 대하여, $\varphi'_{\mu\lambda}$가 $\varphi_{\mu\lambda}$의 $V_{\lambda\mu}\cap V_{\lambda\nu}$에 대한 제한을 뜻할 때, $\varphi'_{\mu\lambda}$가
\[
  (V_{\lambda\mu}\cap V_{\lambda\nu},
   \mathcal A_\lambda|(V_{\lambda\mu}\cap V_{\lambda\nu}))
\]
에서
\[
  (V_{\mu\nu}\cap V_{\mu\lambda},
   \mathcal A_\mu|(V_{\mu\nu}\cap V_{\mu\lambda}))
\]
로 가는 동형사상이고
$\varphi'_{\lambda\nu}=\varphi'_{\lambda\mu}\circ\varphi'_{\mu\nu}$
라고 하자. 이를 $\varphi_{\lambda\mu}$들에 관한 \emph{붙이기 조건}이라 한다. 그러면 먼저 $X_\lambda$들을

\oldpage[0\textsubscript{I}]{39}
$V_{\lambda\mu}$들을 따라 $\varphi_{\lambda\mu}$들로 붙여서 얻는 위상공간 $X$를 생각할 수 있다. $X_\lambda$를 $X$ 안의 대응하는 열린부분집합 $X'_\lambda$와 동일시하면, 가정에 따라 세 집합
$V_{\lambda\mu}\cap V_{\lambda\nu}$,
$V_{\mu\nu}\cap V_{\mu\lambda}$,
$V_{\nu\lambda}\cap V_{\nu\mu}$는 모두
$X'_\lambda\cap X'_\mu\cap X'_\nu$와 동일시된다. 이제 $X_\lambda$의 환 달린 공간 구조를 $X'_\lambda$로 옮길 수 있다. $\mathcal A'_\lambda$가 $\mathcal A_\lambda$로부터 옮겨진 환의 층이면, $\mathcal A'_\lambda$들은 붙이기 조건 (3.3.1)을 만족하므로 $X$ 위의 환의 층 $\mathcal A$를 정의한다. 이때 $(X,\mathcal A)$를 $\varphi_{\lambda\mu}$들을 사용하여 $V_{\lambda\mu}$들을 따라 $(X_\lambda,\mathcal A_\lambda)$들을 \emph{붙여서 얻은 환 달린 공간}이라 한다.
\end{env}

\subsection{$\mathcal A$-가군의 직상}
\label{subsection:0.4.2-ko}

\begin{env}[4.2.1]
\label{0.4.2.1-ko}
$(X,\mathcal A)$와 $(Y,\mathcal B)$를 두 환 달린 공간이라 하고,
$\Psi=(\psi,\theta)$를
$(X,\mathcal A)$에서 $(Y,\mathcal B)$로 가는 사상이라 하자.
그러면 $\psi_*(\mathcal A)$는 $Y$ 위의 환의 층이고,
$\theta$는 환의 층의 준동형
$\mathcal B\to\psi_*(\mathcal A)$이다.
이제 $\mathcal F$를 $\mathcal A$-가군이라 하자. 그 직상
$\psi_*(\mathcal F)$는 $Y$ 위의 아벨군의 층이다. 더 나아가 모든
열린집합 $U\subset Y$에 대하여
\[
  \Gamma(U,\psi_*(\mathcal F))
  =\Gamma(\psi^{-1}(U),\mathcal F)
\]
에는 환
$\Gamma(U,\psi_*(\mathcal A))=\Gamma(\psi^{-1}(U),\mathcal A)$에
관한 가군 구조가 주어진다. 이 구조들을 정의하는 쌍선형 사상들은
제한연산과 양립하므로 $\psi_*(\mathcal F)$ 위에
$\psi_*(\mathcal A)$-가군 구조를 정의한다. 따라서 준동형
$\theta:\mathcal B\to\psi_*(\mathcal A)$를 통해
$\psi_*(\mathcal F)$ 위에 $\mathcal B$-가군 구조도 정의할 수 있다.
이 $\mathcal B$-가군을 
\emph{사상 $\Psi$에 의한 $\mathcal F$의 직상}이라 하고
$\Psi_*(\mathcal F)$로 나타낸다.

$\mathcal F_1$, $\mathcal F_2$가 $X$ 위의 두 $\mathcal A$-가군이고
$u:\mathcal F_1\to\mathcal F_2$가 $\mathcal A$-준동형이면,
$Y$의 열린집합 위의 절단들을 생각함으로써 $\psi_*(u)$가
$\psi_*(\mathcal A)$-준동형
\[
  \psi_*(\mathcal F_1)\longrightarrow\psi_*(\mathcal F_2)
\]
임을 곧바로 알 수 있다. 따라서 이는 특히 $\mathcal B$-준동형
\[
  \Psi_*(\mathcal F_1)\longrightarrow\Psi_*(\mathcal F_2)
\]
이며, $\mathcal B$-준동형으로 볼 때에도 이를 $\Psi_*(u)$로 나타낸다.
따라서 $\Psi_*$는 $\mathcal A$-가군의 범주에서
$\mathcal B$-가군의 범주로 가는 
\emph{공변함자}이다. 또한 이 함자는 
\emph{왼쪽 완전}함이 자명하다(G, II, 2.12).

$\psi_*(\mathcal A)$ 위의 $\mathcal B$-가군 구조와 환의 층 구조는
$\mathcal B$-대수 구조를 정의한다. 이 $\mathcal B$-대수를
$\Psi_*(\mathcal A)$로 나타낸다.
\end{env}

\begin{env}[4.2.2]
\label{0.4.2.2-ko}
$\mathcal M$, $\mathcal N$을 두 $\mathcal A$-가군이라 하자.
$Y$의 모든 열린집합 $U$에 대하여 표준 사상
\[
  \Gamma(\psi^{-1}(U),\mathcal M)
  \times\Gamma(\psi^{-1}(U),\mathcal N)
  \longrightarrow
  \Gamma(\psi^{-1}(U),\mathcal M\otimes_{\mathcal A}\mathcal N)
\]
이 있다. 이 사상은 환
$\Gamma(\psi^{-1}(U),\mathcal A)=\Gamma(U,\psi_*(\mathcal A))$
위에서 쌍선형이고, 따라서 특히 $\Gamma(U,\mathcal B)$ 위에서
쌍선형이다. 그러므로 이는 준동형
\[
  \Gamma(U,\Psi_*(\mathcal M))
  \otimes_{\Gamma(U,\mathcal B)}
  \Gamma(U,\Psi_*(\mathcal N))
  \longrightarrow
  \Gamma(U,\Psi_*(\mathcal M\otimes_{\mathcal A}\mathcal N))
\]
을 정의한다. 이 준동형들이 제한연산과 양립함을 곧바로 확인할 수
있으므로, 이들은 $\mathcal B$-가군의 함자적인 표준 준동형
\[
  \Psi_*(\mathcal M)\otimes_{\mathcal B}\Psi_*(\mathcal N)
  \longrightarrow
  \Psi_*(\mathcal M\otimes_{\mathcal A}\mathcal N)
  \tag{4.2.2.1}\label{0.4.2.2.1-ko}
\]
을 준다.

\oldpage[0\textsubscript{I}]{40}
이 준동형은 일반적으로 단사도 전사도 아니다. $\mathcal P$가 세 번째
$\mathcal A$-가군이면 다음 도표가 가환함을 곧바로 확인할 수 있다.
\[
\begin{tikzcd}[column sep=large,row sep=large]
\Psi_*(\mathcal M)\otimes_{\mathcal B}\Psi_*(\mathcal N)
  \otimes_{\mathcal B}\Psi_*(\mathcal P)
  \arrow[r]
  \arrow[d]
&
\Psi_*(\mathcal M\otimes_{\mathcal A}\mathcal N)
  \otimes_{\mathcal B}\Psi_*(\mathcal P)
  \arrow[d]
\\
\Psi_*(\mathcal M)\otimes_{\mathcal B}
  \Psi_*(\mathcal N\otimes_{\mathcal A}\mathcal P)
  \arrow[r]
&
\Psi_*(\mathcal M\otimes_{\mathcal A}\mathcal N
  \otimes_{\mathcal A}\mathcal P)
\end{tikzcd}
\tag{4.2.2.2}\label{0.4.2.2.2-ko}
\]
\end{env}

\begin{env}[4.2.3]
\label{0.4.2.3-ko}
$\mathcal M$, $\mathcal N$을 두 $\mathcal A$-가군이라 하자.
모든 열린집합 $U\subset Y$에 대하여 정의상
\[
  \Gamma\bigl(\psi^{-1}(U),
    \mathcal H\!om_{\mathcal A}(\mathcal M,\mathcal N)\bigr)
  =
  \operatorname{Hom}_{\mathcal A|V}
    (\mathcal M|V,\mathcal N|V)
\]
이다. 여기서 표기를 간단히 하려고 $V=\psi^{-1}(U)$로 놓았다.
사상 $u\mapsto\Psi_*(u)$는 $\Gamma(U,\mathcal B)$-가군 구조에 관한
준동형
\[
  \operatorname{Hom}_{\mathcal A|V}(\mathcal M|V,\mathcal N|V)
  \longrightarrow
  \operatorname{Hom}_{\mathcal B|U}
    (\Psi_*(\mathcal M)|U,\Psi_*(\mathcal N)|U)
\]
이다. 이 준동형들은 제한연산과 양립하므로 $\mathcal B$-가군의
함자적인 표준 준동형
\[
  \Psi_*\bigl(\mathcal H\!om_{\mathcal A}(\mathcal M,\mathcal N)\bigr)
  \longrightarrow
  \mathcal H\!om_{\mathcal B}
    (\Psi_*(\mathcal M),\Psi_*(\mathcal N))
  \tag{4.2.3.1}\label{0.4.2.3.1-ko}
\]
을 정의한다.
\end{env}

\begin{env}[4.2.4]
\label{0.4.2.4-ko}
$\mathcal C$가 $\mathcal A$-대수이면 합성 준동형
\[
  \Psi_*(\mathcal C)\otimes_{\mathcal B}\Psi_*(\mathcal C)
  \longrightarrow
  \Psi_*(\mathcal C\otimes_{\mathcal A}\mathcal C)
  \longrightarrow
  \Psi_*(\mathcal C)
\]
은 $\Psi_*(\mathcal C)$ 위에 $\mathcal B$-대수 구조를 정의한다.
이는 (4.2.2.2)에서 곧바로 따른다. 마찬가지로 $\mathcal M$이
$\mathcal C$-가군이면 $\Psi_*(\mathcal M)$에는 표준적으로
$\Psi_*(\mathcal C)$-가군 구조가 주어짐을 알 수 있다.
\end{env}

\begin{env}[4.2.5]
\label{0.4.2.5-ko}
특히 $X$가 $Y$의 \emph{닫힌} 부분공간이고 $\psi$가 표준 단사사상
$j:X\to Y$인 경우를 생각하자. 환의 층 $\mathcal B$의 $X$로의 제한을
\[
  \mathcal B'=\mathcal B|X=j^*(\mathcal B)
\]
로 놓으면, $\mathcal A$-가군 $\mathcal M$은 준동형
$\theta^\sharp:\mathcal B'\to\mathcal A$를 통해
$\mathcal B'$-가군으로 볼 수 있다. 이때 $\Psi_*(\mathcal M)$은
$X$ 위에는 $\mathcal M$을, 그 밖에는 $0$을 유도하는
$\mathcal B$-가군이다. $\mathcal N$이 두 번째 $\mathcal A$-가군이면
\[
  \Psi_*(\mathcal M)\otimes_{\mathcal B}\Psi_*(\mathcal N)
\]
은 $\Psi_*(\mathcal M\otimes_{\mathcal B'}\mathcal N)$과 동일시되고,
\[
  \mathcal H\!om_{\mathcal B}
    (\Psi_*(\mathcal M),\Psi_*(\mathcal N))
\]
은
\[
  \Psi_*\bigl(
    \mathcal H\!om_{\mathcal B'}(\mathcal M,\mathcal N)\bigr)
\]
과 동일시된다.
\end{env}

\begin{env}[4.2.6]
\label{0.4.2.6-ko}
$(Z,\mathcal C)$를 세 번째 환 달린 공간이라 하고,
$\Psi'=(\psi',\theta')$를 $(Y,\mathcal B)$에서
$(Z,\mathcal C)$로 가는 사상이라 하자. $\Psi''$이 합성사상
$\Psi'\circ\Psi$이면
\[
  \Psi''_*=\Psi'_*\circ\Psi_*
\]
임은 자명하다.
\end{env}

\oldpage[0\textsubscript{I}]{41}
\subsection{$\mathcal B$-가군의 역상}
\label{subsection:0.4.3-ko}

\begin{env}[4.3.1]
\label{0.4.3.1-ko}
(4.2.1)의 가정과 표기를 사용하자. $\mathcal G$를 $\mathcal B$-가군이라
하고, 그 역상 $\psi^*(\mathcal G)$(3.7.1)를 생각하자. 이는 $X$ 위의
아벨군의 층이다. $\psi^*(\mathcal G)$와
$\psi^*(\mathcal B)$의 절단들의 정의(3.7.1)로부터
$\psi^*(\mathcal G)$에는 표준적으로 $\psi^*(\mathcal B)$-가군 구조가
주어짐을 알 수 있다. 한편 준동형
\[
  \theta^\sharp:\psi^*(\mathcal B)\longrightarrow\mathcal A
\]
는 $\mathcal A$에 $\psi^*(\mathcal B)$-가군 구조를 준다. 혼동을
피할 필요가 있을 때에는 이 가군을 $\mathcal A_{[\theta]}$로 나타낸다.
그러면 텐서곱
\[
  \psi^*(\mathcal G)
  \otimes_{\psi^*(\mathcal B)}\mathcal A_{[\theta]}
\]
에는 $\mathcal A$-가군 구조가 주어진다. 이 $\mathcal A$-가군을
사상 $\Psi$에 의한 $\mathcal G$의 \emph{역상}이라 하고
$\Psi^*(\mathcal G)$로 나타낸다.

$\mathcal G_1$, $\mathcal G_2$가 $Y$ 위의 두 $\mathcal B$-가군이고
$v:\mathcal G_1\to\mathcal G_2$가 $\mathcal B$-준동형이면,
$\psi^*(v)$는 $\psi^*(\mathcal G_1)$에서
$\psi^*(\mathcal G_2)$로 가는 $\psi^*(\mathcal B)$-준동형이다.
따라서 $\psi^*(v)\otimes1$은 $\mathcal A$-준동형
\[
  \Psi^*(\mathcal G_1)\longrightarrow\Psi^*(\mathcal G_2)
\]
이며 이를 $\Psi^*(v)$로 나타낸다. 이로써 $\Psi^*$를
$\mathcal B$-가군의 범주에서 $\mathcal A$-가군의 범주로 가는
\emph{공변함자}로 정의하였다. 여기서 이 함자는 $\psi^*$와 달리
일반적으로 완전하지 않고 \emph{오른쪽 완전}할 뿐이다. 실제로
$\mathcal A$와의 텐서곱은 $\psi^*(\mathcal B)$-가군의 범주에서
오른쪽 완전함자이다.

모든 $x\in X$에 대하여 (3.7.2)에 의해
\[
  (\Psi^*(\mathcal G))_x
  =\mathcal G_{\psi(x)}
   \otimes_{\mathcal B_{\psi(x)}}\mathcal A_x
\]
이다. 따라서 $\Psi^*(\mathcal G)$의 지지집합은
$\psi^{-1}(\operatorname{Supp}\mathcal G)$에 포함된다.
\end{env}

\begin{env}[4.3.2]
\label{0.4.3.2-ko}
$(\mathcal G_\lambda)$를 $\mathcal B$-가군의 귀납계라 하고
$\mathcal G=\varinjlim\mathcal G_\lambda$를 그 귀납극한이라 하자.
표준 준동형 $\mathcal G_\lambda\to\mathcal G$는
$\psi^*(\mathcal B)$-준동형
$\psi^*(\mathcal G_\lambda)\to\psi^*(\mathcal G)$를 정의하며,
따라서 표준 준동형
\[
  \varinjlim\psi^*(\mathcal G_\lambda)
  \longrightarrow
  \psi^*(\mathcal G)
\]
을 준다. 층의 귀납극한의 한 점에서의 줄기는 같은 점에서의 줄기들의
귀납극한이므로(G, II, 1.11), 앞의 표준 준동형은
\emph{전단사}이다(3.7.2). 또한 텐서곱은 층의 귀납극한과 가환하므로
$\mathcal A$-가군의 \emph{함자적인 표준 동형}
\[
  \varinjlim\Psi^*(\mathcal G_\lambda)
  \xrightarrow{\sim}
  \Psi^*\bigl(\varinjlim\mathcal G_\lambda\bigr)
\]
을 얻는다.

한편 $\mathcal B$-가군들의 유한 직합
$\bigoplus_i\mathcal G_i$에 대하여
\[
  \psi^*\bigl(\bigoplus_i\mathcal G_i\bigr)
  =\bigoplus_i\psi^*(\mathcal G_i)
\]
임은 자명하다. 따라서 $\mathcal A_{[\theta]}$와 텐서곱하면
\[
  \Psi^*\bigl(\bigoplus_i\mathcal G_i\bigr)
  =\bigoplus_i\Psi^*(\mathcal G_i)
  \tag{4.3.2.1}\label{0.4.3.2.1-ko}
\]
을 얻는다. 앞의 결과에 따라 귀납극한으로 넘어가면 이 등식은
\emph{임의의} 직합에 대해서도 성립한다.
\end{env}

\begin{env}[4.3.3]
\label{0.4.3.3-ko}
$\mathcal G_1$, $\mathcal G_2$를 두 $\mathcal B$-가군이라 하자.
아벨군의 층의 역상에 관한 정의(3.7.1)로부터
$\psi^*(\mathcal B)$-가군의 표준 준동형
\[
  \psi^*(\mathcal G_1)
  \otimes_{\psi^*(\mathcal B)}\psi^*(\mathcal G_2)
  \longrightarrow
  \psi^*(\mathcal G_1\otimes_{\mathcal B}\mathcal G_2)
\]
을 곧바로 얻는다. 층들의 텐서곱의 한 점에서의 줄기는 그 점에서의
줄기들의 텐서곱이므로(G, II, 2.8), (3.7.2)에 따라 이 준동형은
실제로 \emph{동형}이다. $\mathcal A$와 텐서곱하여
\emph{함자적인 표준 동형}
\[
  \Psi^*(\mathcal G_1)\otimes_{\mathcal A}\Psi^*(\mathcal G_2)
  \xrightarrow{\sim}
  \Psi^*(\mathcal G_1\otimes_{\mathcal B}\mathcal G_2)
  \tag{4.3.3.1}\label{0.4.3.3.1-ko}
\]
을 얻는다.
\end{env}

\begin{env}[4.3.4]
\label{0.4.3.4-ko}
$\mathcal C$를 $\mathcal B$-대수라 하자. $\mathcal C$ 위의 대수 구조를
주는 것은 각 줄기에서 확인되는 결합법칙과 가환법칙을 만족하는
$\mathcal B$-준동형
\[
  \mathcal C\otimes_{\mathcal B}\mathcal C
  \longrightarrow\mathcal C
\]
을 주는 것과 같다. 앞의 동형을 사용하면 이 준동형을 같은 조건들을
만족하는 $\mathcal A$-가군 준동형
\[
  \Psi^*(\mathcal C)\otimes_{\mathcal A}\Psi^*(\mathcal C)
  \longrightarrow\Psi^*(\mathcal C)
\]
으로 옮길 수 있다. 따라서 $\Psi^*(\mathcal C)$에는
$\mathcal A$-대수 구조가 주어진다. 특히 정의에서 곧바로
$\mathcal A$-대수 $\Psi^*(\mathcal B)$가 표준 동형을 제외하면
$\mathcal A$와 같음을 알 수 있다.

마찬가지로 $\mathcal M$이 $\mathcal C$-가군이면, 이 가군 구조를
주는 것은
\oldpage[0\textsubscript{I}]{42}
결합법칙을 만족하는 $\mathcal B$-준동형
\[
  \mathcal C\otimes_{\mathcal B}\mathcal M
  \longrightarrow\mathcal M
\]
을 주는 것과 같다. 이를 옮겨서 $\Psi^*(\mathcal M)$ 위에
$\Psi^*(\mathcal C)$-가군 구조를 얻는다.
\end{env}

\begin{env}[4.3.5]
\label{0.4.3.5-ko}
$\mathcal J$를 $\mathcal B$의 아이디얼층이라 하자. 함자 $\psi^*$가
완전하므로 $\psi^*(\mathcal B)$-가군 $\psi^*(\mathcal J)$는
$\psi^*(\mathcal B)$의 아이디얼층과 표준적으로 동일시된다.
따라서 표준 단사사상
$\psi^*(\mathcal J)\to\psi^*(\mathcal B)$는
$\mathcal A$-가군 준동형
\[
  \Psi^*(\mathcal J)
  =\psi^*(\mathcal J)
   \otimes_{\psi^*(\mathcal B)}\mathcal A_{[\theta]}
  \longrightarrow\mathcal A
\]
을 준다. 이 준동형에 의한 $\Psi^*(\mathcal J)$의 상을
$\Psi^*(\mathcal J)\mathcal A$로 나타내며, 혼동의 우려가 없으면
$\mathcal J\mathcal A$로도 나타낸다. 따라서 정의상
\[
  \mathcal J\mathcal A
  =\theta^\sharp\bigl(\psi^*(\mathcal J)\bigr)\mathcal A
\]
이고, 특히 모든 $x\in X$에 대하여
\[
  (\mathcal J\mathcal A)_x
  =\theta_x^\sharp(\mathcal J_{\psi(x)})\mathcal A_x
\]
이다. 여기서는 $\psi^*(\mathcal J)$의 줄기와
$\mathcal J$의 대응하는 줄기의 표준적인 동일시(3.7.2)를 사용하였다.
$\mathcal J_1$, $\mathcal J_2$가 $\mathcal B$의 두 아이디얼이면
\[
  (\mathcal J_1\mathcal J_2)\mathcal A
  =\mathcal J_1(\mathcal J_2\mathcal A)
  =(\mathcal J_1\mathcal A)(\mathcal J_2\mathcal A)
\]
이다. $\mathcal F$가 $\mathcal A$-가군이면
$\mathcal J\mathcal F=(\mathcal J\mathcal A)\mathcal F$로 놓는다.
\end{env}

\begin{env}[4.3.6]
\label{0.4.3.6-ko}
$(Z,\mathcal C)$를 세 번째 환 달린 공간이라 하고,
$\Psi'=(\psi',\theta')$를 $(Y,\mathcal B)$에서
$(Z,\mathcal C)$로 가는 사상이라 하자. $\Psi''$이 합성사상
$\Psi'\circ\Psi$이면, 정의 (4.3.1)와 (4.3.3.1)에 따라
\[
  {\Psi''}^*=\Psi^*\circ{\Psi'}^*
\]
이다.
\end{env}

\subsection{직상과 역상의 관계}
\label{subsection:0.4.4-ko}

\begin{env}[4.4.1]
\label{0.4.4.1-ko}
(4.2.1)의 가정과 표기를 사용하고, $\mathcal G$를 $\mathcal B$-가군이라
하자. 정의상 $\mathcal B$-가군 준동형
$u:\mathcal G\to\Psi_*(\mathcal F)$를 여전히 $\mathcal G$에서
$\mathcal F$로 가는 \emph{$\Psi$-사상}이라 한다. 혼동의 우려가 없으면
단순히 $\mathcal G$에서 $\mathcal F$로 가는 \emph{준동형}이라 하고
$u:\mathcal G\to\mathcal F$로 쓴다.
이와 같은 준동형을 주는 것은 $U$가 $X$의 열린집합이고 $V$가 $Y$의
열린집합이며 $\psi(U)\subset V$인 모든 쌍 $(U,V)$에 대하여
\emph{준동형}
\[
  u_{U,V}:\Gamma(V,\mathcal G)\longrightarrow\Gamma(U,\mathcal F)
\]
을 주는 것과 같다. 여기서 이는 \emph{$\Gamma(V,\mathcal B)$-가군의
준동형}이고, $\Gamma(U,\mathcal F)$의 $\Gamma(V,\mathcal B)$-가군 구조는
환 준동형
$\theta_{U,V}:\Gamma(V,\mathcal B)\to\Gamma(U,\mathcal A)$를 통하여
주어진다. 또한 $u_{U,V}$들은 도식 (3.5.1.1)을 가환하게 해야 한다.
사실 $u$를 정의하려면 $U$와 $V$가 각각 $X$와 $Y$의 위상기저
$\mathfrak B$와 $\mathfrak B'$를 달릴 때의 $u_{U,V}$만 주고,
이러한 제한들에 대하여 (3.5.1.1)의 가환성을 확인하면 충분하다.
\end{env}

\begin{env}[4.4.2]
\label{0.4.4.2-ko}
(4.2.1)과 (4.2.6)의 가정 아래에서 $\mathcal H$를 $\mathcal C$-가군이라
하고 $v:\mathcal H\to\Psi'_*(\mathcal G)$를 $\Psi'$-사상이라 하자.
그러면
\[
  w:\mathcal H\xrightarrow{v}\Psi'_*(\mathcal G)
  \xrightarrow{\Psi'_*(u)}\Psi'_*(\Psi_*(\mathcal F))
\]
는 $\Psi''$-사상이며, 이를 $u$와 $v$의 \emph{합성}이라 한다.
\end{env}

\begin{env}[4.4.3]
\label{0.4.4.3-ko}
이제 $\mathcal F$와 $\mathcal G$에 관한 \emph{쌍함자}들의 표준
\emph{동형}
\[
  \operatorname{Hom}_{\mathcal A}(\Psi^*(\mathcal G),\mathcal F)
  \xrightarrow{\sim}
  \operatorname{Hom}_{\mathcal B}(\mathcal G,\Psi_*(\mathcal F))
  \tag{4.4.3.1}\label{0.4.4.3.1-ko}
\]
을 정의할 수 있음을 보이자. 이 동형을
$v\mapsto v_\theta^\flat$으로 나타내며, 혼동의 우려가 없으면
$v\mapsto v^\flat$으로도 나타낸다. 그 역동형은
$u\mapsto u_\theta^\sharp$, 또는 $u\mapsto u^\sharp$으로 나타낸다.
정의는 다음과 같다. 준동형
$v:\Psi^*(\mathcal G)\to\mathcal F$를 표준 사상
$\psi^*(\mathcal G)\to\Psi^*(\mathcal G)$와 합성하면
$\psi^*(\mathcal B)$-가군의 준동형이기도 한 아벨군의 층의 준동형
$v':\psi^*(\mathcal G)\to\mathcal F$를 얻는다. 이에 따라 (3.7.1)에서
준동형
${v'}^\flat:\mathcal G\to\psi_*(\mathcal F)=\Psi_*(\mathcal F)$를
얻으며, 쉽게 확인할 수 있듯이 이는 $\mathcal B$-가군의 준동형이다.
\oldpage[0\textsubscript{I}]{43}
$v_\theta^\flat={v'}^\flat$으로 놓는다. 마찬가지로
$\mathcal B$-가군 준동형 $u:\mathcal G\to\Psi_*(\mathcal F)$로부터
(3.7.1)에 의해 $\psi^*(\mathcal B)$-가군 준동형
$u^\sharp:\psi^*(\mathcal G)\to\mathcal F$를 얻는다. 이를
$\mathcal A$와 텐서곱하여 $\mathcal A$-가군 준동형
$\Psi^*(\mathcal G)\to\mathcal F$를 얻고, 이를
$u_\theta^\sharp$으로 나타낸다. 다음 등식들은 곧바로 확인된다.
\[
  (u_\theta^\sharp)_\theta^\flat=u,
  \qquad
  (v_\theta^\flat)_\theta^\sharp=v.
\]
또한 $v\mapsto v_\theta^\flat$은 $\mathcal F$에 관하여 함자적이다.
$u\mapsto u_\theta^\sharp$이 $\mathcal G$에 관하여 함자적이라는 것은
(3.5.4)에서와 같은 형식적인 논증으로부터 따른다. 이 논증은
(4.3.1)에서 직접 증명한 $\Psi^*$의 함자성도 증명한다.

$v$로 $\Psi^*(\mathcal G)$의 항등준동형을 취하면
$v_\theta^\flat$은 준동형
\[
  \rho_{\mathcal G}:\mathcal G\longrightarrow
  \Psi_*(\Psi^*(\mathcal G))
  \tag{4.4.3.2}\label{0.4.4.3.2-ko}
\]
이 된다. $u$로 $\Psi_*(\mathcal F)$의 항등준동형을 취하면
$u_\theta^\sharp$은 준동형
\[
  \sigma_{\mathcal F}:\Psi^*(\Psi_*(\mathcal F))
  \longrightarrow\mathcal F
  \tag{4.4.3.3}\label{0.4.4.3.3-ko}
\]
이 된다. 이들을 \emph{표준} 준동형이라 한다. 이들은 일반적으로
단사도 전사도 아니다. (3.5.3.3)과 (3.5.4.4)에 대응하는 표준 분해도
성립한다.

$s$가 $Y$의 열린집합 $V$ 위에서 $\mathcal G$의 절단이면,
$\rho_{\mathcal G}(s)$는 $\psi^{-1}(V)$ 위의 $\Psi^*(\mathcal G)$의
절단 $s'\otimes1$이다. 여기서 모든 $x\in\psi^{-1}(V)$에 대하여
$s'_x=s_{\psi(x)}$이다. 또한 준동형
$u:\mathcal G\to\psi_*(\mathcal F)$는 모든 $x\in X$에 대하여 줄기
위의 준동형 $u_x:\mathcal G_{\psi(x)}\to\mathcal F_x$를 정의한다.
이는 $(u^\sharp)_x:(\Psi^*(\mathcal G))_x\to\mathcal F_x$와 표준
준동형
\[
  \mathcal G_{\psi(x)}\longrightarrow(\Psi^*(\mathcal G))_x,
  \qquad s_x\longmapsto s_x\otimes1
\]
을 합성하여 얻는다. 여기서
\[
  (\Psi^*(\mathcal G))_x
  =\mathcal G_{\psi(x)}
   \otimes_{\mathcal B_{\psi(x)}}\mathcal A_x.
\]
또한 $u_x$는 $V$가 $\psi(x)$의 근방들을 달릴 때 준동형
\[
  \Gamma(V,\mathcal G)\xrightarrow{u}
  \Gamma(\psi^{-1}(V),\mathcal F)\longrightarrow\mathcal F_x
\]
의 귀납극한을 취하여 얻을 수도 있다.
\end{env}

\begin{env}[4.4.4]
\label{0.4.4.4-ko}
$\mathcal F_1$, $\mathcal F_2$를 $\mathcal A$-가군,
$\mathcal G_1$, $\mathcal G_2$를 $\mathcal B$-가군이라 하고,
$u_i$ $(i=1,2)$를 $\mathcal G_i$에서 $\mathcal F_i$로 가는
준동형이라 하자. $u_1\otimes u_2$로 준동형
\[
  u:\mathcal G_1\otimes_{\mathcal B}\mathcal G_2
  \longrightarrow
  \mathcal F_1\otimes_{\mathcal A}\mathcal F_2
\]
중에서 (4.3.3.1)의 동일시 아래
$u^\sharp=(u_1)^\sharp\otimes(u_2)^\sharp$을 만족하는 것을 나타낸다.
또한 $u$는 합성
\[
  \mathcal G_1\otimes_{\mathcal B}\mathcal G_2
  \longrightarrow
  \Psi_*(\mathcal F_1)\otimes_{\mathcal B}\Psi_*(\mathcal F_2)
  \longrightarrow
  \Psi_*(\mathcal F_1\otimes_{\mathcal A}\mathcal F_2)
\]
과 같음을 확인할 수 있다. 여기서 첫째 화살표는 보통의 텐서곱
$u_1\otimes_{\mathcal B}u_2$이고, 둘째 화살표는 (4.2.2.1)의 표준
준동형이다.
\end{env}

\begin{env}[4.4.5]
\label{0.4.4.5-ko}
$(\mathcal G_\lambda)_{\lambda\in L}$를 $\mathcal B$-가군의 귀납계라
하고, 모든 $\lambda\in L$에 대하여
$u_\lambda:\mathcal G_\lambda\to\Psi_*(\mathcal F)$가 귀납계를 이루는
준동형이라 하자. $\mathcal G=\varinjlim\mathcal G_\lambda$ 및
$u=\varinjlim u_\lambda$로 놓는다. 그러면 $(u_\lambda)^\sharp$들은
준동형 $\Psi^*(\mathcal G_\lambda)\to\mathcal F$의 귀납계를 이루고,
이 계의 귀납극한은 바로 $u^\sharp$이다.
\end{env}

\begin{env}[4.4.6]
\label{0.4.4.6-ko}
$\mathcal M$, $\mathcal N$을 두 $\mathcal B$-가군이라 하고, $V$를
$Y$의 열린집합, $U=\psi^{-1}(V)$라 하자. 사상
$v\mapsto\Psi^*(v)$는 $\Gamma(V,\mathcal B)$-가군의 준동형
\[
  \operatorname{Hom}_{\mathcal B|V}(\mathcal M|V,\mathcal N|V)
  \longrightarrow
  \operatorname{Hom}_{\mathcal A|U}
  (\Psi^*(\mathcal M)|U,\Psi^*(\mathcal N)|U)
\]
이다.
\begin{sloppypar}
실제로 $\operatorname{Hom}_{\mathcal A|U}
(\Psi^*(\mathcal M)|U,\allowbreak\Psi^*(\mathcal N)|U)$에는 본래
$\Gamma(U,\psi^*(\mathcal B))$-가군 구조가 주어지며, 표준 준동형
\oldpage[0\textsubscript{I}]{44}
$\Gamma(V,\mathcal B)\to\Gamma(U,\psi^*(\mathcal B))$(3.7.2)을 통하여
$\Gamma(V,\mathcal B)$-가군이기도 하다.
\end{sloppypar}
이 준동형들은 제한사상과 양립함을 곧바로 알 수 있다. 따라서
함자적인 표준 준동형
\[
  \gamma:\mathcal H\!om_{\mathcal B}(\mathcal M,\mathcal N)
  \longrightarrow
  \Psi_*\bigl(\mathcal H\!om_{\mathcal A}
  (\Psi^*(\mathcal M),\Psi^*(\mathcal N))\bigr)
\]
을 정의한다. 이에 대응하여 준동형
\[
  \gamma^\sharp:
  \Psi^*\bigl(\mathcal H\!om_{\mathcal B}(\mathcal M,\mathcal N)\bigr)
  \longrightarrow
  \mathcal H\!om_{\mathcal A}
  (\Psi^*(\mathcal M),\Psi^*(\mathcal N))
\]
도 얻는다. 이 표준 준동형들은 $\mathcal M$과 $\mathcal N$에 관하여
함자적이다.
\end{env}

\begin{env}[4.4.7]
\label{0.4.4.7-ko}
$\mathcal F$와 $\mathcal G$가 각각 $\mathcal A$-대수와
$\mathcal B$-대수라고 하자. $u:\mathcal G\to\Psi_*(\mathcal F)$가
$\mathcal B$-대수의 준동형이면
$u^\sharp:\Psi^*(\mathcal G)\to\mathcal F$는 $\mathcal A$-대수의
준동형이다. 이는 도식
\[
\begin{tikzcd}[column sep=huge,row sep=large]
\mathcal G\otimes_{\mathcal B}\mathcal G
  \arrow[r] \arrow[d]
  & \mathcal G \arrow[d,"u"] \\
\Psi_*(\mathcal F\otimes_{\mathcal A}\mathcal F)
  \arrow[r]
  & \Psi_*(\mathcal F)
\end{tikzcd}
\]
의 가환성과 (4.4.4)에서 따른다. 마찬가지로
$v:\Psi^*(\mathcal G)\to\mathcal F$가 $\mathcal A$-대수의 준동형이면
$v^\flat:\mathcal G\to\Psi_*(\mathcal F)$는 $\mathcal B$-대수의
준동형이다.
\end{env}

\begin{env}[4.4.8]
\label{0.4.4.8-ko}
$(Z,\mathcal C)$를 세 번째 환 달린 공간이라 하고,
$\Psi'=(\psi',\theta')$를 $(Y,\mathcal B)$에서 $(Z,\mathcal C)$로 가는
사상, $\Psi'':(X,\mathcal A)\to(Z,\mathcal C)$를 합성사상
$\Psi'\circ\Psi$라 하자. $\mathcal H$를 $\mathcal C$-가군이라 하고,
$v'$를 $\mathcal H$에서 $\mathcal G$로 가는 준동형이라 하자. 정의상
합성 $v''=v\circ v'$는 $\mathcal H$에서 $\mathcal F$로 가는 준동형
\[
  \mathcal H\xrightarrow{v'}\Psi'_*(\mathcal G)
  \xrightarrow{\Psi'_*(v)}\Psi'_*(\Psi_*(\mathcal F))
\]
으로 주어진다. 그러면 ${v''}^\sharp$은 준동형
\[
  \Psi^*({\Psi'}^*(\mathcal H))
  \xrightarrow{\Psi^*({v'}^\sharp)}
  \Psi^*(\mathcal G)
  \xrightarrow{v^\sharp}\mathcal F
\]
과 같음을 확인할 수 있다.
\end{env}

\section{준연접층과 연접층}
\label{section:0.5-ko}

\subsection{준연접층}
\label{subsection:0.5.1-ko}

\begin{env}[5.1.1]
\label{0.5.1.1-ko}
$(X,\mathcal O_X)$를 환 달린 공간이라 하고, $\mathcal F$를
$\mathcal O_X$-가군이라 하자. $\mathcal O_X$-가군 준동형
$u:\mathcal O_X\to\mathcal F$를 주는 것은 절단
$s=u(1)\in\Gamma(X,\mathcal F)$를 주는 것과 같다. 실제로 $s$가
주어졌을 때 모든 절단 $t\in\Gamma(U,\mathcal O_X)$에 대하여 반드시
$u(t)=t\cdot(s|U)$이다. 이때 $u$가 \emph{절단 $s$에 의해 정의된다}고
한다.

이제 $I$를 임의의 지표집합이라 하고, 직합층
$\mathcal O_X^{(I)}$를 생각하자. 모든 $i\in I$에 대하여 $h_i$를
$i$번째 인자에서 $\mathcal O_X^{(I)}$로 가는 표준 단사사상이라 하자.
그러면 $u\mapsto(u\circ h_i)$는
$\operatorname{Hom}_{\mathcal O_X}(\mathcal O_X^{(I)},\mathcal F)$에서
곱
\[
  \bigl(\operatorname{Hom}_{\mathcal O_X}
  (\mathcal O_X,\mathcal F)\bigr)^I
\]
로 가는 동형이다. 따라서 준동형
$u:\mathcal O_X^{(I)}\to\mathcal F$와 $X$ 위에서 $\mathcal F$의
\emph{절단족 $(s_i)_{i\in I}$} 사이에는 표준적인 일대일 대응이 있다.
$(s_i)$에 대응하는 준동형 $u$는
$(a_i)\in(\Gamma(U,\mathcal O_X))^{(I)}$를
\[
  \sum_{i\in I}a_i\cdot(s_i|U)
\]
로 보낸다.

$(s_i)$가 정의하는 준동형
\oldpage[0\textsubscript{I}]{45}
$\mathcal O_X^{(I)}\to\mathcal F$가 \emph{전사}이면
$\mathcal F$가 \emph{족 $(s_i)$에 의해 생성된다}고 한다. 다시 말해,
모든 $x\in X$에 대하여 $\mathcal F_x$가 $(s_i)_x$들로 생성되는
$\mathcal O_x$-가군이라는 뜻이다. $\mathcal F$가 $X$ 위의 모든 절단의
족, 또는 그 한 부분족에 의해 생성되면 $\mathcal F$가
\emph{$X$ 위의 절단들에 의해 생성된다}고 한다. 즉, 알맞은 $I$에
대하여 전사 준동형
\[
  \mathcal O_X^{(I)}\longrightarrow\mathcal F
\]
가 존재한다는 뜻이다.

어떤 $\mathcal O_X$-가군 $\mathcal F$에는 점 $x_0\in X$가 있어서,
$x_0$의 열린 근방 $U$를 \emph{어떻게} 택하더라도
$\mathcal F|U$가 $U$ 위의 절단들로 생성되지 않을 수 있다.
실제로 $X=\mathbb R$, $\mathcal O_X$를 상수층 $\mathbb Z$로 놓고,
$\mathcal F_0=\{0\}$ 및 $x\ne0$일 때
$\mathcal F_x=\mathbb Z$인 $\mathcal O_X$의 대수적 부분층
$\mathcal F$를 취하며 $x_0=0$으로 놓으면 된다. $0$의 근방 $U$ 위에서
$\mathcal F|U$의 절단은 $0$뿐이다.
\end{env}

\begin{env}[5.1.2]
\label{0.5.1.2-ko}
$f:X\to Y$를 환 달린 공간의 사상이라 하자. $\mathcal F$가 $X$ 위의
절단들로 생성되는 $\mathcal O_X$-가군이면, 표준 준동형
\[
  f^*(f_*(\mathcal F))\longrightarrow\mathcal F
\]
(4.4.3.3)은 \emph{전사}이다. 실제로 (5.1.1)의 표기에서
$s_i\otimes1$은 $X$ 위의 $f^*(f_*(\mathcal F))$의 절단이고,
$\mathcal F$에서의 그 상은 $s_i$이다. (5.1.1)의 예에서 $f$를
항등사상으로 취하면 이 명제의 역은 일반적으로 거짓임을 알 수 있다.
\end{env}

\begin{env}[5.1.3]
\label{0.5.1.3-ko}
$\mathcal O_X$-가군 $\mathcal F$가 다음 조건을 만족하면
\emph{준연접}이라고 한다. 모든 $x\in X$에 대하여 $x$의 열린 근방
$U$가 존재하여, $\mathcal F|U$가 임의의 지표집합 $I$, $J$에 대한
꼴
\[
  \mathcal O_X^{(I)}|U\longrightarrow\mathcal O_X^{(J)}|U
\]
의 준동형의 \emph{여핵}과 동형이다. $\mathcal O_X$ 자체가 준연접
$\mathcal O_X$-가군임은 자명하며, 준연접 $\mathcal O_X$-가군들의
모든 직합도 준연접 $\mathcal O_X$-가군이다. $\mathcal O_X$-대수
$\mathcal A$가 $\mathcal O_X$-가군으로서 준연접이면
$\mathcal A$를 \emph{준연접}이라고 한다.
\end{env}

\begin{env}[5.1.4]
\label{0.5.1.4-ko}
$f:X\to Y$를 환 달린 공간의 사상이라 하자. $\mathcal G$가 준연접
$\mathcal O_Y$-가군이면 $f^*(\mathcal G)$는 준연접
$\mathcal O_X$-가군이다. 실제로 모든 $x\in X$에 대하여 $Y$에서
$f(x)$의 열린 근방 $V$가 존재하여 $\mathcal G|V$가 준동형
\[
  \mathcal O_Y^{(I)}|V\longrightarrow\mathcal O_Y^{(J)}|V
\]
의 여핵이다. $U=f^{-1}(V)$라 하고 $f_U$를 $f$의 $U$로의 제한이라
하면 $f^*(\mathcal G)|U=f_U^*(\mathcal G|V)$이다. $f_U^*$는 오른쪽
완전하고 직합과 가환하므로 $f_U^*(\mathcal G|V)$는 준동형
\[
  \mathcal O_X^{(I)}|U\longrightarrow\mathcal O_X^{(J)}|U
\]
의 여핵이다.
\end{env}

\subsection{유한 생성층}
\label{subsection:0.5.2-ko}

\begin{env}[5.2.1]
\label{0.5.2.1-ko}
$\mathcal O_X$-가군 $\mathcal F$가 다음 조건을 만족하면
\emph{유한 생성}이라고 한다. 모든 $x\in X$에 대하여 $x$의 열린
근방 $U$가 존재하여 $\mathcal F|U$가 $U$ 위의 \emph{유한} 절단족에
의해 생성된다. 동치로, 유한한 $p$에 대하여 $\mathcal F|U$가
$(\mathcal O_X|U)^p$ 꼴의 층의 몫층과 동형이다. 유한 생성층의 모든
몫층은 유한 생성이고, 유한 생성층들의 모든 유한 직합과 유한
텐서곱도 유한 생성이다. 유한 생성 $\mathcal O_X$-가군이 반드시
준연접인 것은 아니다. (5.1.1)의 예의 $\mathcal F$에 대하여
$\mathcal O_X/\mathcal F$를 생각하면 이를 알 수 있다.
$\mathcal F$가 유한 생성이면 모든 $x\in X$에 대하여
$\mathcal F_x$는 유한 생성 $\mathcal O_x$-가군이지만, (5.1.1)의 예는
이 필요조건이 일반적으로 충분하지 않음을 보여 준다.
\end{env}

\begin{env}[5.2.2]
\label{0.5.2.2-ko}
$\mathcal F$를 \emph{유한 생성} $\mathcal O_X$-가군이라 하자.
$s_i$ $(1\leq i\leq n)$가 점 $x\in X$의 열린 근방 $U$ 위에서
$\mathcal F$의 절단들이고 $(s_i)_x$들이 $\mathcal F_x$를 생성하면,
$x$의 열린 근방 $V\subset U$가 존재하여 모든 $y\in V$에 대하여
$(s_i)_y$들이 $\mathcal F_y$를 생성한다(FAC, I, 2, 12, 명제 1).
특히 $\mathcal F$의 지지집합은 \emph{닫혀 있다}.

\oldpage[0\textsubscript{I}]{46}

마찬가지로 $u:\mathcal F\to\mathcal G$가 $u_x=0$인 준동형이면,
$x$의 근방 $U$가 존재하여 모든 $y\in U$에 대하여 $u_y=0$이다.
\end{env}

\begin{env}[5.2.3]
\label{0.5.2.3-ko}
$X$가 \emph{준콤팩트}라고 하자. $\mathcal F$, $\mathcal G$를 두
$\mathcal O_X$-가군, $\mathcal G$를 \emph{유한 생성} 가군이라 하고,
$u:\mathcal F\to\mathcal G$를 \emph{전사} 준동형이라 하자. 또한
$\mathcal F$가 $\mathcal O_X$-가군의 귀납계
$(\mathcal F_\lambda)$의 귀납극한이라고 하자. 그러면 어떤 지표
$\mu$가 존재하여 준동형 $\mathcal F_\mu\to\mathcal G$가
\emph{전사}이다. 실제로 모든 $x\in X$에 대하여 $x$의 열린 근방
$U(x)$ 위에 $\mathcal G$의 유한한 절단계 $s_i$가 존재하여, 모든
$y\in U(x)$에 대하여 $(s_i)_y$들이 $\mathcal G_y$를 생성한다.
따라서 $x$의 열린 근방 $V(x)\subset U(x)$와 $V(x)$ 위의
$\mathcal F$의 절단 $t_i$들이 존재하여 모든 $i$에 대하여
$s_i|V(x)=u(t_i)$이다. 또한 $t_i$들이 $V(x)$ 위의 같은 층
$\mathcal F_{\lambda(x)}$의 절단들의 표준적인 상이라고 가정할 수
있다. 이제 유한개의 근방 $V(x_k)$로 $X$를 덮고, 모든
$\lambda(x_k)$보다 큰 지표 $\mu$를 택하면 이 $\mu$가 조건을
만족함은 자명하다.

여전히 $X$가 준콤팩트라고 하고, $\mathcal F$를 $X$ 위의 절단들에
의해 생성되는 유한 생성 $\mathcal O_X$-가군이라 하자(5.1.1).
그러면 $\mathcal F$는 이 절단들의 어떤 \emph{유한} 부분족에 의해
생성된다. 실제로 유한개의 열린 근방 $U_k$로 $X$를 덮되, 각 $k$에
대하여 $X$ 위의 $\mathcal F$의 유한개의 절단 $s_{ik}$의 $U_k$로의
제한들이 $\mathcal F|U_k$를 생성하도록 하면 된다. 그러면
$s_{ik}$들이 $\mathcal F$를 생성함은 자명하다.
\end{env}

\begin{env}[5.2.4]
\label{0.5.2.4-ko}
$f:X\to Y$를 환 달린 공간의 사상이라 하자. $\mathcal G$가 유한 생성
$\mathcal O_Y$-가군이면 $f^*(\mathcal G)$는 유한 생성
$\mathcal O_X$-가군이다. 실제로 모든 $x\in X$에 대하여 $Y$에서
$f(x)$의 열린 근방 $V$와 전사 준동형
$v:\mathcal O_Y^p|V\to\mathcal G|V$가 존재한다. $U=f^{-1}(V)$라 하고
$f_U$를 $f$의 $U$로의 제한이라 하면
$f^*(\mathcal G)|U=f_U^*(\mathcal G|V)$이다. $f_U^*$는 오른쪽
완전하고(4.3.1) 직합과 가환하므로(4.3.2), $f_U^*(v)$는 전사 준동형
\[
  \mathcal O_X^p|U\longrightarrow f^*(\mathcal G)|U
\]
이다.
\end{env}

\begin{env}[5.2.5]
\label{0.5.2.5-ko}
$\mathcal O_X$-가군 $\mathcal F$가 다음 조건을 만족하면
\emph{유한 표시를 갖는다}고 한다. 모든 $x\in X$에 대하여 $x$의
열린 근방 $U$가 존재하여 $\mathcal F|U$가 양의 정수 $p$, $q$에 대한
\emph{$(\mathcal O_X|U)$-준동형}
\[
  \mathcal O_X^p|U\longrightarrow\mathcal O_X^q|U
\]
의 \emph{여핵}과 동형이다. 이러한 $\mathcal O_X$-가군은 유한
생성이고 준연접이다. $f:X\to Y$가 환 달린 공간의 사상이고
$\mathcal G$가 유한 표시를 갖는 $\mathcal O_Y$-가군이면,
(5.1.4)의 논증이 보여 주듯이 $f^*(\mathcal G)$도 유한 표시를 갖는다.
\end{env}

\begin{env}[5.2.6]
\label{0.5.2.6-ko}
$\mathcal F$를 유한 표시를 갖는 $\mathcal O_X$-가군이라 하자(5.2.5).
그러면 모든 $\mathcal O_X$-가군 $\mathcal H$에 대하여 함자적인 표준
준동형
\[
  \bigl(\mathcal H\!om_{\mathcal O_X}(\mathcal F,\mathcal H)\bigr)_x
  \longrightarrow
  \operatorname{Hom}_{\mathcal O_x}(\mathcal F_x,\mathcal H_x)
\]
은 \emph{전단사}이다(T, 4.1.1).
\end{env}

\begin{env}[5.2.7]
\label{0.5.2.7-ko}
$\mathcal F$, $\mathcal G$를 유한 표시를 갖는 두
$\mathcal O_X$-가군이라 하자. 어떤 $x\in X$에 대하여
$\mathcal F_x$와 $\mathcal G_x$가 동형인 두 $\mathcal O_x$-가군이면,
$x$의 열린 근방 $U$가 존재하여 $\mathcal F|U$와 $\mathcal G|U$는
\emph{동형}이다. 실제로
$\varphi:\mathcal F_x\to\mathcal G_x$와
$\psi:\mathcal G_x\to\mathcal F_x$가 서로 역인 두 동형이라 하자.
(5.2.6)에 따라 $x$의 열린 근방 $V$ 위에서
$\mathcal H\!om_{\mathcal O_X}(\mathcal F,\mathcal G)$의 절단 $u$와
$\mathcal H\!om_{\mathcal O_X}(\mathcal G,\mathcal F)$의 절단 $v$가
존재하여

\oldpage[0\textsubscript{I}]{47}

$u_x=\varphi$ 및 $v_x=\psi$이다. $(u\circ v)_x$와
$(v\circ u)_x$는 항등자동사상이므로, $x$의 열린 근방
$U\subset V$가 존재하여 $(u\circ v)|U$와 $(v\circ u)|U$가
항등자동사상이다. 이로부터 명제가 따른다.
\end{env}

\subsection{연접층}
\label{subsection:0.5.3-ko}

\begin{env}[5.3.1]
\label{0.5.3.1-ko}
$\mathcal O_X$-가군 $\mathcal F$가 다음 두 조건을 만족하면
\emph{연접}이라고 한다.
\begin{enumerate}
  \item[$a)$] $\mathcal F$는 유한 생성이다.
  \item[$b)$] 모든 열린집합 $U\subset X$, 모든 정수 $n>0$, 그리고
    모든 준동형
    $u:\mathcal O_X^n|U\to\mathcal F|U$에 대하여 $u$의 핵은 유한
    생성이다.
\end{enumerate}
이 두 조건은 \emph{국소적} 성질임에 유의하자.

이하에서 상기하는 연접층의 성질 대부분의 증명에 대해서는
(FAC), I, 2를 보라.
\end{env}

\begin{env}[5.3.2]
\label{0.5.3.2-ko}
모든 연접 $\mathcal O_X$-가군은 유한 표시를 갖는다(5.2.5). 그 역은
반드시 참인 것은 아니다. 실제로 $\mathcal O_X$ 자체가 반드시 연접
$\mathcal O_X$-가군인 것은 아니다.

연접 $\mathcal O_X$-가군의 \emph{유한 생성} 부분
$\mathcal O_X$-가군은 모두 연접이고, 연접 $\mathcal O_X$-가군들의
모든 \emph{유한} 직합도 연접 $\mathcal O_X$-가군이다.
\end{env}

\begin{env}[5.3.3]
\label{0.5.3.3-ko}
$0\to\mathcal F\to\mathcal G\to\mathcal H\to0$이
$\mathcal O_X$-가군들의 완전열이고 이 세 $\mathcal O_X$-가군 가운데
둘이 연접이면 나머지 하나도 연접이다.
\end{env}

\begin{env}[5.3.4]
\label{0.5.3.4-ko}
$\mathcal F$, $\mathcal G$가 두 연접 $\mathcal O_X$-가군이고
$u:\mathcal F\to\mathcal G$가 준동형이면,
$\operatorname{Im}(u)$, $\operatorname{Ker}(u)$,
$\operatorname{Coker}(u)$는 연접 $\mathcal O_X$-가군이다. 특히
$\mathcal F$, $\mathcal G$가 어떤 연접 $\mathcal O_X$-가군의 연접
부분 $\mathcal O_X$-가군이면 $\mathcal F+\mathcal G$와
$\mathcal F\cap\mathcal G$도 연접이다.
\end{env}

\begin{env}[5.3.5]
\label{0.5.3.5-ko}
$\mathcal F$, $\mathcal G$가 두 연접 $\mathcal O_X$-가군이면
$\mathcal F\otimes_{\mathcal O_X}\mathcal G$와
$\mathcal H\!om_{\mathcal O_X}(\mathcal F,\mathcal G)$도 연접이다.
\end{env}

\begin{env}[5.3.6]
\label{0.5.3.6-ko}
$\mathcal F$를 연접 $\mathcal O_X$-가군, $\mathcal J$를
$\mathcal O_X$의 연접 아이디얼층이라 하자. 그러면
$\mathcal O_X$-가군 $\mathcal J\mathcal F$는 표준 준동형
\[
  \mathcal J\otimes_{\mathcal O_X}\mathcal F\longrightarrow\mathcal F
\]
에 의한 $\mathcal J\otimes_{\mathcal O_X}\mathcal F$의 상이므로
연접이다(5.3.4, 5.3.5).
\end{env}

\begin{env}[5.3.7]
\label{0.5.3.7-ko}
$\mathcal O_X$-대수 $\mathcal A$가 $\mathcal O_X$-가군으로서
연접이면 $\mathcal A$를 \emph{연접}이라고 한다. 특히
$\mathcal O_X$가 \emph{연접인 환의 층}일 필요충분조건은 모든
열린집합 $U\subset X$와 모든 꼴의 준동형
$u:\mathcal O_X^p|U\to\mathcal O_X|U$에 대하여 $u$의 핵이
유한 생성 $(\mathcal O_X|U)$-가군인 것이다.

$\mathcal O_X$가 연접인 환의 층이면, 유한 표시를 갖는 모든
$\mathcal O_X$-가군 $\mathcal F$는 (5.3.4)에 따라 연접이다.

$\mathcal O_X$-가군 $\mathcal F$의 \emph{소멸자}는 표준 준동형
\[
  \mathcal O_X\longrightarrow
  \mathcal H\!om_{\mathcal O_X}(\mathcal F,\mathcal F)
\]
의 핵 $\mathcal J$이다. 이 준동형은 각 절단
$s\in\Gamma(U,\mathcal O_X)$에
$\operatorname{Hom}(\mathcal F|U,\mathcal F|U)$에서 $s$에 의한
곱셈을 대응시킨다. $\mathcal O_X$가 연접이고 $\mathcal F$가 연접
$\mathcal O_X$-가군이면 $\mathcal J$는 연접이며(5.3.4, 5.3.5),
모든 $x\in X$에 대하여 $\mathcal J_x$는 $\mathcal F_x$의
소멸자이다(5.2.6).
\end{env}

\oldpage[0\textsubscript{I}]{48}

\begin{env}[5.3.8]
\label{0.5.3.8-ko}
$\mathcal O_X$가 연접이라고 하자. $\mathcal F$를 연접
$\mathcal O_X$-가군, $x$를 $X$의 한 점, $M$을 $\mathcal F_x$의
유한 생성 부분가군이라 하자. 그러면 $x$의 열린 근방 $U$와
$\mathcal F|U$의 연접 부분 $(\mathcal O_X|U)$-가군 $\mathcal G$가
존재하여 $\mathcal G_x=M$이다(T, 4.1, 보조정리 1).

이 결과와 연접 $\mathcal O_X$-가군의 부분 $\mathcal O_X$-가군들에
관한 성질을 함께 적용하면, $\mathcal O_X$가 연접이기 위하여 환
$\mathcal O_x$들이 만족해야 할 필요조건을 얻는다. 예를 들어
(5.3.4)에 따라 $\mathcal O_x$의 두 유한 생성 아이디얼의 교집합은
다시 유한 생성 아이디얼이어야 한다.
\end{env}

\begin{env}[5.3.9]
\label{0.5.3.9-ko}
$\mathcal O_X$가 연접이고 $M$이 유한 표시를 갖는
$\mathcal O_x$-가군이라고 하자. 따라서 $M$은 어떤 준동형
$\varphi:\mathcal O_x^p\to\mathcal O_x^q$의 여핵과 동형이다. 그러면
$x$의 열린 근방 $U$와 연접 $(\mathcal O_X|U)$-가군 $\mathcal F$가
존재하여 $\mathcal F_x$는 $M$과 동형이다. 실제로 (5.2.6)에 따라
$x$의 열린 근방 $U$ 위의
$\mathcal H\!om_{\mathcal O_X}(\mathcal O_X^p,\mathcal O_X^q)$의
절단 $u$가 존재하여 $u_x=\varphi$이다. 준동형
\[
  u:\mathcal O_X^p|U\longrightarrow\mathcal O_X^q|U
\]
의 여핵 $\mathcal F$가 조건을 만족한다(5.3.4).
\end{env}

\begin{env}[5.3.10]
\label{0.5.3.10-ko}
$\mathcal O_X$가 연접이고 $\mathcal J$가 $\mathcal O_X$의 연접
아이디얼층이라고 하자. $(\mathcal O_X/\mathcal J)$-가군
$\mathcal F$가 연접일 필요충분조건은 그것이 $\mathcal O_X$-가군으로서
연접인 것이다. 특히 $\mathcal O_X/\mathcal J$는 연접인 환의 층이다.
\end{env}

\begin{env}[5.3.11]
\label{0.5.3.11-ko}
$f:X\to Y$를 환 달린 공간의 사상이라 하고 $\mathcal O_X$가
연접이라고 하자. 그러면 모든 연접 $\mathcal O_Y$-가군
$\mathcal G$에 대하여 $f^*(\mathcal G)$는 연접
$\mathcal O_X$-가군이다. 실제로 (5.2.4)의 표기에서
$\mathcal G|V$가 준동형
$v:\mathcal O_Y^q|V\to\mathcal O_Y^p|V$의 여핵이라고 가정할 수 있다.
$f_U^*$가 오른쪽 완전하므로
$f^*(\mathcal G)|U=f_U^*(\mathcal G|V)$는 준동형
\[
  f_U^*(v):\mathcal O_X^q|U\longrightarrow\mathcal O_X^p|U
\]
의 여핵이고, 이로부터 주장이 따른다.
\end{env}

\begin{env}[5.3.12]
\label{0.5.3.12-ko}
$Y$를 $X$의 닫힌 부분집합, $j:Y\to X$를 표준 포함사상,
$\mathcal O_Y$를 $Y$ 위의 환의 층이라 하고
$\mathcal O_X=j_*(\mathcal O_Y)$로 놓자. $\mathcal O_Y$-가군
$\mathcal G$가 유한 생성일(각각 준연접일, 연접일) 필요충분조건은
$j_*(\mathcal G)$가 유한 생성(각각 준연접, 연접)
$\mathcal O_X$-가군인 것이다.
\end{env}

\subsection{국소 자유층}
\label{subsection:0.5.4-ko}

\begin{env}[5.4.1]
\label{0.5.4.1-ko}
$X$를 환 달린 공간이라 하자. $\mathcal O_X$-가군 $\mathcal F$가
\emph{국소 자유}라는 것은 모든 $x\in X$에 대하여 $x$의 열린 근방
$U$가 존재하여 $\mathcal F|U$가
$\mathcal O_X^{(I)}|U$ 꼴의 $(\mathcal O_X|U)$-가군과 동형이라는
뜻이다. 여기서 $I$는 $U$에 따라 달라도 된다. 모든 $U$에서 $I$가
유한이면 $\mathcal F$가 \emph{유한 계수}라고 하고, 모든 $U$에서
$I$의 원소 수가 같은 유한수 $n$이면 $\mathcal F$가
\emph{계수 $n$}이라고 한다. 계수 1인 국소 자유
$\mathcal O_X$-가군을 \emph{가역}이라고도 한다((5.4.3) 참조).
$\mathcal F$가 유한 계수의 국소 자유 $\mathcal O_X$-가군이면 모든
$x\in X$에 대하여 $\mathcal F_x$는 유한 계수 $n(x)$의 자유
$\mathcal O_x$-가군이고, $x$의 어떤 근방 $U$에서 $\mathcal F|U$는
계수 $n(x)$이다. 따라서 $X$가 연결이면 $n(x)$는 \emph{상수}이다.

모든 국소 자유층이 준연접임은 자명하다. 또한 $\mathcal O_X$가
연접인 환의 층이면 유한 계수의 모든 국소 자유 $\mathcal O_X$-가군은
연접이다.

$\mathcal L$이 국소 자유이면 $\mathcal O_X$-가군의 범주에서
$\mathcal F$에 관한 함자
$\mathcal L\otimes_{\mathcal O_X}\mathcal F$는 \emph{완전}하다.

우리는 주로 유한 계수의 국소 자유 $\mathcal O_X$-가군을 다룰 것이다.

\oldpage[0\textsubscript{I}]{49}

따라서 별도의 말이 없이 국소 자유층이라고 할 때에는
\emph{유한 계수}임을 전제한다.
\end{env}

\begin{env}[5.4.2]
\label{0.5.4.2-ko}
$\mathcal L$, $\mathcal F$가 두 $\mathcal O_X$-가군이면 함자적인 표준
준동형
\[
  \check{\mathcal L}\otimes_{\mathcal O_X}\mathcal F
  =\mathcal H\!om_{\mathcal O_X}(\mathcal L,\mathcal O_X)
   \otimes_{\mathcal O_X}\mathcal F
  \longrightarrow
  \mathcal H\!om_{\mathcal O_X}(\mathcal L,\mathcal F)
  \tag{5.4.2.1}\label{0.5.4.2.1-ko}
\]
이 존재한다. 이는 다음과 같이 정의된다. 모든 열린집합 $U$와 모든 쌍
$(u,t)$, 즉
\[
  u\in\Gamma\bigl(U,
    \mathcal H\!om_{\mathcal O_X}(\mathcal L,\mathcal O_X)\bigr)
  =\operatorname{Hom}(\mathcal L|U,\mathcal O_X|U),
  \qquad
  t\in\Gamma(U,\mathcal F)
\]
에 대하여, 모든 $x\in U$에서 $s_x\in\mathcal L_x$를
$u_x(s_x)t_x\in\mathcal F_x$로 보내는
$\operatorname{Hom}(\mathcal L|U,\mathcal F|U)$의 원소를 대응시킨다.
$\mathcal L$이 \emph{유한 계수의 국소 자유}이면 이 준동형은
\emph{전단사}이다. 이 성질은 국소적이므로
$\mathcal L=\mathcal O_X^n$인 경우로 한정할 수 있다. 모든
$\mathcal O_X$-가군 $\mathcal G$에 대하여
$\mathcal H\!om_{\mathcal O_X}(\mathcal O_X^n,\mathcal G)$가
$\mathcal G^n$과 표준적으로 동형이므로, 다시 자명한 경우
$\mathcal L=\mathcal O_X$로 환원된다.
\end{env}

\begin{env}[5.4.3]
\label{0.5.4.3-ko}
$\mathcal L$이 가역이면 그 쌍대
$\check{\mathcal L}=\mathcal H\!om_{\mathcal O_X}
(\mathcal L,\mathcal O_X)$도 가역이다. 실제로 문제는 국소적이므로 곧
$\mathcal L=\mathcal O_X$인 경우로 환원된다. 또한 표준 동형
\[
  \mathcal H\!om_{\mathcal O_X}(\mathcal L,\mathcal O_X)
  \otimes_{\mathcal O_X}\mathcal L
  \xrightarrow{\sim}
  \mathcal O_X
  \tag{5.4.3.1}\label{0.5.4.3.1-ko}
\]
이 존재한다. 실제로 (5.4.2)에 따라 표준 동형
$\mathcal H\!om_{\mathcal O_X}(\mathcal L,\mathcal L)
\xrightarrow{\sim}\mathcal O_X$를 정의하면 충분하다. 그런데 모든
$\mathcal O_X$-가군 $\mathcal F$에 대하여 표준 준동형
$\mathcal O_X\longrightarrow
\mathcal H\!om_{\mathcal O_X}(\mathcal F,\mathcal F)$가 존재한다
(5.3.7). $\mathcal F=\mathcal L$이 가역일 때 이 준동형이 전단사임을
보이면 된다. 이 문제도 국소적이므로 자명한 경우
$\mathcal L=\mathcal O_X$로 환원된다.

이에 따라
$\mathcal L^{-1}=\mathcal H\!om_{\mathcal O_X}
(\mathcal L,\mathcal O_X)$로 놓고, $\mathcal L^{-1}$을
$\mathcal L$의 \emph{역}이라고 한다. $X$가 한 점으로만 이루어지고
$\mathcal O_X$가 극대 아이디얼 $\mathfrak m$을 갖는 \emph{국소환}
$A$인 경우에는 ``가역층''이라는 용어를 다음과 같이 정당화할 수 있다.
$M$, $M'$을 두 $A$-가군이라 하고($M$은 유한 생성),
$M\otimes_A M'$이 $A$와 동형이라고 하자. 그러면
$(A/\mathfrak m)\otimes_A(M\otimes_A M')$은
$(M/\mathfrak mM)\otimes_{A/\mathfrak m}(M'/\mathfrak mM')$과
동일시된다. 체 $A/\mathfrak m$ 위의 이 벡터공간 텐서곱이
$A/\mathfrak m$과 동형이므로 $M/\mathfrak mM$과
$M'/\mathfrak mM'$은 모두 1차원이어야 한다. 따라서
$\mathfrak mM$에 속하지 않는 모든 $z\in M$에 대하여
$M=Az+\mathfrak mM$이고, $M$이 유한 생성이므로 나카야마 보조정리에
따라 $M=Az$이다. 더욱이 $z$의 소멸자는 $A$와 동형인
$M\otimes_A M'$을 소멸시키므로 그 소멸자는 $\{0\}$이다. 따라서
$M$은 \emph{$A$와 동형}이다. 일반적인 경우, 이는 $\mathcal L$이
유한 생성 $\mathcal O_X$-가군이고 어떤 $\mathcal O_X$-가군
$\mathcal F$에 대하여
$\mathcal L\otimes_{\mathcal O_X}\mathcal F$가 $\mathcal O_X$와
동형이며, 각 환 $\mathcal O_x$가 국소환이면 모든 $x\in X$에 대하여
$\mathcal L_x$가 $\mathcal O_x$와 동형임을 보여 준다. 더 나아가
$\mathcal O_X$와 $\mathcal L$이 \emph{연접}이면 (5.2.7)에 따라
$\mathcal L$은 가역이다.
\end{env}

\begin{env}[5.4.4]
\label{0.5.4.4-ko}
$\mathcal L$, $\mathcal L'$이 두 가역 $\mathcal O_X$-가군이면
$\mathcal L\otimes_{\mathcal O_X}\mathcal L'$도 가역이다. 실제로
문제는 국소적이므로 $\mathcal L=\mathcal O_X$라고 가정할 수 있고,
그 경우 결과는 자명하다. 모든 정수 $n\geq1$에 대하여
$\mathcal L^{\otimes n}$은 $\mathcal L$과 같은 $n$개의 층의
텐서곱을 뜻한다.

\oldpage[0\textsubscript{I}]{50}

관례상 $\mathcal L^{\otimes0}=\mathcal O_X$로 놓고, $n\geq1$일 때
$\mathcal L^{\otimes(-n)}=(\mathcal L^{-1})^{\otimes n}$으로 놓는다.
이 표기 아래 임의의 정수 $m,n$에 대하여 \emph{함자적인 표준 동형}
\[
  \mathcal L^{\otimes m}\otimes_{\mathcal O_X}\mathcal L^{\otimes n}
  \xrightarrow{\sim}
  \mathcal L^{\otimes(n+m)}
  \tag{5.4.4.1}\label{0.5.4.4.1-ko}
\]
이 존재한다. 실제로 정의에 따라 곧 $m=-1$, $n=1$인 경우로
환원되며, 이 경우의 동형은 (5.4.3)에서 정의하였다.
\end{env}

\begin{env}[5.4.5]
\label{0.5.4.5-ko}
$f:Y\to X$를 환 달린 공간의 사상이라 하자. $\mathcal L$이 국소 자유
(각각 가역) $\mathcal O_X$-가군이면 $f^*(\mathcal L)$은 국소 자유
(각각 가역) $\mathcal O_Y$-가군이다. 이는 국소적으로 동형인 두
$\mathcal O_X$-가군의 역상들이 국소적으로 동형이라는 사실,
$f^*$가 유한 직합과 가환한다는 사실, 그리고
$f^*(\mathcal O_X)=\mathcal O_Y$라는 사실(4.3.4)에서 곧 따른다.
또한 함자적인 표준 준동형
\[
  f^*(\check{\mathcal L})
  \longrightarrow
  \check{\bigl(f^*(\mathcal L)\bigr)}
\]
이 존재한다(4.4.6). $\mathcal L$이 국소 자유이면 이 준동형은
\emph{전단사}이다. 실제로 다시 자명한 경우
$\mathcal L=\mathcal O_X$로 환원된다. 따라서 $\mathcal L$이
가역이면 모든 정수 $n$에 대하여 $f^*(\mathcal L^{\otimes n})$은
$(f^*(\mathcal L))^{\otimes n}$과 표준적으로 동일시된다.
\end{env}

\begin{env}[5.4.6]
\label{0.5.4.6-ko}
$\mathcal L$을 가역 $\mathcal O_X$-가군이라 하자. 직합 아벨군
\[
  \Gamma_*(X,\mathcal L)=\Gamma_*(\mathcal L)
  =\bigoplus_{n\in\mathbb Z}\Gamma(X,\mathcal L^{\otimes n})
\]
을 생각하자. $s_n\in\Gamma(X,\mathcal L^{\otimes n})$과
$s_m\in\Gamma(X,\mathcal L^{\otimes m})$의 쌍 $(s_n,s_m)$에
(5.4.4.1)에 따라 $s_n\otimes s_m$에 표준적으로 대응하는
$X$ 위의 $\mathcal L^{\otimes(n+m)}$의 절단을 대응시켜 이 군에
\emph{등급환}의 구조를 준다. 이 곱셈의 결합법칙은 곧 확인된다.
$\Gamma_*(X,\mathcal L)$은 $\mathcal L$에 관한 공변함자로서
등급환의 범주에 값을 갖는다.

이제 $\mathcal F$가 임의의 $\mathcal O_X$-가군이면
\[
  \Gamma_*(\mathcal L,\mathcal F)
  =\bigoplus_{n\in\mathbb Z}
   \Gamma\bigl(X,
     \mathcal F\otimes_{\mathcal O_X}\mathcal L^{\otimes n}\bigr)
\]
로 놓는다. 이 아벨군에 등급환 $\Gamma_*(\mathcal L)$ 위의
\emph{등급 가군} 구조를 다음과 같이 준다. 두 절단
$s_n\in\Gamma(X,\mathcal L^{\otimes n})$과
$u_m\in\Gamma(X,
\mathcal F\otimes_{\mathcal O_X}\mathcal L^{\otimes m})$의 쌍에
(5.4.4.1)에 따라 $s_n\otimes u_m$에 표준적으로 대응하는
$\mathcal F\otimes_{\mathcal O_X}\mathcal L^{\otimes(m+n)}$의 절단을
대응시킨다. 가군의 공리들은 곧 확인된다. $X$, $\mathcal L$을
고정하면 $\Gamma_*(\mathcal L,\mathcal F)$는 $\mathcal F$에 관한
공변함자로서 $\Gamma_*(\mathcal L)$-등급 가군의 범주에 값을 갖는다.
$X$, $\mathcal F$를 고정하면 이는 $\mathcal L$에 관한 공변함자로서
아벨군의 범주에 값을 갖는다.

$f:Y\to X$를 환 달린 공간의 사상이라 하자. 표준 준동형 (4.4.3.2)
\[
  \rho:\mathcal L^{\otimes n}
  \longrightarrow
  f_*\bigl(f^*(\mathcal L^{\otimes n})\bigr)
\]
은 아벨군의 준동형
\[
  \Gamma(X,\mathcal L^{\otimes n})
  \longrightarrow
  \Gamma\bigl(Y,f^*(\mathcal L^{\otimes n})\bigr)
\]
을 정의한다. $f^*(\mathcal L^{\otimes n})
=(f^*(\mathcal L))^{\otimes n}$이므로, 표준 준동형 (4.4.3.2)와
(5.4.4.1)의 정의로부터 앞의 준동형들이 \emph{함자적인 등급환
준동형}
\[
  \Gamma_*(\mathcal L)\longrightarrow
  \Gamma_*(f^*(\mathcal L))
\]
을 정의함을 알 수 있다. 같은 표준 준동형 (4.4.3)은 아벨군의
준동형
\[
  \Gamma\bigl(X,
    \mathcal F\otimes_{\mathcal O_X}\mathcal L^{\otimes n}\bigr)
  \longrightarrow
  \Gamma\bigl(Y,
    f^*(\mathcal F\otimes_{\mathcal O_X}\mathcal L^{\otimes n})\bigr)
\]
도 정의한다. 또한
\[
  f^*(\mathcal F\otimes_{\mathcal O_X}\mathcal L^{\otimes n})
  =f^*(\mathcal F)\otimes_{\mathcal O_Y}
   (f^*(\mathcal L))^{\otimes n}
  \tag{4.3.3.1}
\]
이다.

\oldpage[0\textsubscript{I}]{51}

따라서 이 준동형들은 $n$이 변할 때 \emph{등급 가군의 쌍준동형}
\[
  \Gamma_*(\mathcal L,\mathcal F)
  \longrightarrow
  \Gamma_*(f^*(\mathcal L),f^*(\mathcal F))
\]
을 정의한다.
\end{env}

\begin{env}[5.4.7]
\label{0.5.4.7-ko}
모든 가역 $\mathcal O_X$-가군이 $\mathfrak M$의 정확히 한 원소와
동형이 되도록 하는 가역 $\mathcal O_X$-가군들의 \emph{집합}
$\mathfrak M$($\mathfrak M(X)$라고도 쓴다)이 존재함을 보일 수
있다\footnote{서론에서 인용한 준비 중인 책을 보라.}. $\mathfrak M$의
두 원소 $\mathcal L$, $\mathcal L'$에
$\mathcal L\otimes_{\mathcal O_X}\mathcal L'$과 동형인
$\mathfrak M$의 유일한 원소를 대응시켜 $\mathfrak M$ 위에 합성법칙을
정의한다. 이 합성법칙에 관하여 $\mathfrak M$은 \emph{코호몰로지군}
$H^1(X,\mathcal O_X^*)$과 동형인 군이다. 여기서 $\mathcal O_X^*$는
모든 열린집합 $U\subset X$에 대하여
$\Gamma(U,\mathcal O_X^*)$가 환 $\Gamma(U,\mathcal O_X)$의 가역원군이
되는 $\mathcal O_X$의 부분층이다. 따라서 $\mathcal O_X^*$는
\emph{곱셈적} 아벨군의 층이다.

이를 위해 먼저 모든 열린집합 $U\subset X$에서 절단군
$\Gamma(U,\mathcal O_X^*)$가 $(\mathcal O_X|U)$-가군
$\mathcal O_X|U$의 \emph{자동사상군}과 표준적으로 동일시됨에
유의하자. 이 동일시는 $U$ 위의 $\mathcal O_X^*$의 절단
$\varepsilon$에, 모든 $x\in U$와 모든 $s_x\in\mathcal O_x$에 대하여
$u_x(s_x)=\varepsilon_xs_x$인 $\mathcal O_X|U$의 자동사상 $u$를
대응시킨다. 이제 $\mathfrak U=(U_\lambda)$를 $X$의 열린 덮개라 하자.
모든 지표쌍 $(\lambda,\mu)$에 대하여
$\mathcal O_X|(U_\lambda\cap U_\mu)$의 자동사상
$\theta_{\lambda\mu}$를 주는 것은 $\mathcal O_X^*$에 값을 갖는 덮개
$\mathfrak U$의 $1$-\emph{공사슬}을 주는 것과 같다. 이
$\theta_{\lambda\mu}$들이 붙이기 조건(3.3.1)을 만족한다는 것은
대응하는 공사슬이 \emph{공순환}이라는 뜻이다. 마찬가지로 모든
$\lambda$에 대하여 $\mathcal O_X|U_\lambda$의 자동사상
$\omega_\lambda$를 주는 것은 $\mathcal O_X^*$에 값을 갖는
$\mathfrak U$의 $0$-공사슬을 주는 것과 같고, 그 \emph{공경계}는
자동사상족
\[
  (\omega_\lambda|U_\lambda\cap U_\mu)\circ
  (\omega_\mu|U_\lambda\cap U_\mu)^{-1}
\]
에 대응한다.

$\mathcal O_X^*$에 값을 갖는 $\mathfrak U$의 각 1-공순환에, 그
공순환에 대응하는 자동사상족 $(\theta_{\lambda\mu})$로부터 붙여서
얻는 가역 $\mathcal O_X$-가군과 동형인 $\mathfrak M$의 원소를
대응시킬 수 있다. 서로 코호몰로그인 두 공순환에는 같은
$\mathfrak M$의 원소가 대응한다(3.3.2). 즉, 사상
$\varphi_{\mathfrak U}:H^1(\mathfrak U,\mathcal O_X^*)\to\mathfrak M$을
정의하였다. $\mathfrak B$가 $\mathfrak U$보다 세밀한 $X$의 두 번째
열린 덮개이면 도식
\[
\begin{tikzcd}[column sep=huge,row sep=large]
H^1(\mathfrak U,\mathcal O_X^*)
  \arrow[dd]
  \arrow[rrd,"\varphi_{\mathfrak U}"] & & \\
& & \mathfrak M \\
H^1(\mathfrak B,\mathcal O_X^*)
  \arrow[rru,swap,"\varphi_{\mathfrak B}"] & &
\end{tikzcd}
\]
은 가환한다. 여기서 세로 화살표는 표준 준동형(G, II, 5.7)이며,
가환성은 (3.3.3)에서 따른다. 귀납극한을 취하면 사상
$H^1(X,\mathcal O_X^*)\to\mathfrak M$을 얻는다. 실제로 체흐
코호몰로지군 $\check H^1(X,\mathcal O_X^*)$는 잘 알려진 대로 첫째
코호몰로지군 $H^1(X,\mathcal O_X^*)$와 동일시된다(G, II, 5.9,
정리 5.9.1의 따름정리). 이 사상은 \emph{전사}이다. 실제로 정의에
따라 모든 가역 $\mathcal O_X$-가군 $\mathcal L$에 대하여 $X$의 열린
덮개 $(U_\lambda)$가 존재하여 $\mathcal L$은 층들
$\mathcal O_X|U_\lambda$를 붙여서 얻어진다(3.3.1). 또한 이 사상은
\emph{단사}이다. 이를 보이기 위해서는 각 사상
$H^1(\mathfrak U,\mathcal O_X^*)\to\mathfrak M$에 대하여 증명하면
충분하고, 이는 (3.3.2)에서 따른다. 이제 이렇게 정의한 전단사가
군 준동형임을 보이는 일만 남았다.

\oldpage[0\textsubscript{I}]{52}

두 가역 $\mathcal O_X$-가군 $\mathcal L$, $\mathcal L'$이 주어졌다고
하자. 모든 $\lambda$에 대하여 $\mathcal L|U_\lambda$와
$\mathcal L'|U_\lambda$가 $\mathcal O_X|U_\lambda$와 동형이 되게 하는
열린 덮개 $(U_\lambda)$가 존재한다. 따라서 각 지표 $\lambda$마다
$\Gamma(U_\lambda,\mathcal L)$의 원소 $a_\lambda$(각각
$\Gamma(U_\lambda,\mathcal L')$의 원소 $a'_\lambda$)가 존재하여,
$s_\lambda$가 $\Gamma(U_\lambda,\mathcal O_X)$를 달릴 때
$\Gamma(U_\lambda,\mathcal L)$의 원소들은
$s_\lambda a_\lambda$ 꼴이고, $\Gamma(U_\lambda,\mathcal L')$의
원소들은 $s_\lambda a'_\lambda$ 꼴이다. 대응하는 두 공순환을
$(\varepsilon_{\lambda\mu})$, $(\varepsilon'_{\lambda\mu})$라 하자.
$U_\lambda\cap U_\mu$ 위에서
$s_\lambda a_\lambda=s_\mu a_\mu$(각각
$s_\lambda a'_\lambda=s_\mu a'_\mu$)인 것은
$s_\lambda=\varepsilon_{\lambda\mu}s_\mu$(각각
$s_\lambda=\varepsilon'_{\lambda\mu}s_\mu$)인 것과 동치이다.
$U_\lambda$ 위의
$\mathcal L\otimes_{\mathcal O_X}\mathcal L'$의 절단들은
$s_\lambda$, $s'_\lambda$가 $\Gamma(U_\lambda,\mathcal O_X)$를 달릴
때 $s_\lambda s'_\lambda(a_\lambda\otimes a'_\lambda)$들의 유한합이다.
따라서 공순환
$(\varepsilon_{\lambda\mu}\varepsilon'_{\lambda\mu})$가
$\mathcal L\otimes_{\mathcal O_X}\mathcal L'$에 대응함은 자명하고,
이로써 증명이 끝난다\footnote{이 결과의 일반형에 대해서는 51쪽의
주에서 인용한 책을 보라.}.
\end{env}

\begin{env}[5.4.8]
\label{0.5.4.8-ko}
$f=(\psi,\omega)$를 환 달린 공간의 사상 $Y\to X$라 하자. 역상함자
$f^*$는 집합 $\mathfrak M(X)$에서 집합 $\mathfrak M(Y)$로 가는
사상을 정의한다. 이 사상도 언어의 남용으로 $f^*$라고 쓴다. 한편
표준 준동형(T, 3.2.2)
\[
  H^1(X,\mathcal O_X^*)\longrightarrow H^1(Y,\mathcal O_Y^*)
  \tag{5.4.8.1}\label{0.5.4.8.1-ko}
\]
이 존재한다. $\mathfrak M(X)$와 $H^1(X,\mathcal O_X^*)$(각각
$\mathfrak M(Y)$와 $H^1(Y,\mathcal O_Y^*)$)를 (5.4.7)에 따라
표준적으로 동일시하면, 준동형 (5.4.8.1)은 \emph{사상 $f^*$와
동일시된다}. 실제로 $\mathcal L$이 $X$의 열린 덮개
$(U_\lambda)$에 대응하는 공순환 $(\varepsilon_{\lambda\mu})$에서
나왔다고 하자. $f^*(\mathcal L)$이
$(\varepsilon_{\lambda\mu})$의 코호몰로지류의 (5.4.8.1)에 의한 상을
코호몰로지류로 갖는 공순환에서 나옴을 보이면 충분하다. 이제
$\theta_{\lambda\mu}$를 $\varepsilon_{\lambda\mu}$에 대응하는
$\mathcal O_X|(U_\lambda\cap U_\mu)$의 자동사상이라 하자.
$f^*(\mathcal L)$은 자동사상 $f^*(\theta_{\lambda\mu})$를 이용하여
$\mathcal O_Y|\psi^{-1}(U_\lambda)$들을 붙여서 얻어진다. 따라서 이
자동사상들이 공순환
$(\omega^\#(\varepsilon_{\lambda\mu}))$에 대응함을 확인하면 충분하고,
이는 정의에서 곧 따른다. 여기서 $\varepsilon_{\lambda\mu}$는
$\rho$(3.7.2)에 의한 그 표준적인 상, 즉
$\psi^{-1}(U_\lambda\cap U_\mu)$ 위의
$\psi^*(\mathcal O_X^*)$의 절단과 동일시하였다.
\end{env}

\begin{env}[5.4.9]
\label{0.5.4.9-ko}
$\mathcal E$, $\mathcal F$를 두 $\mathcal O_X$-가군이라 하고
$\mathcal F$가 \emph{국소 자유}라고 하자. $\mathcal G$가
\emph{$\mathcal E$에 의한 $\mathcal F$의 확장}, 즉 완전열
\[
  0\longrightarrow\mathcal E\xrightarrow{i}\mathcal G
  \xrightarrow{p}\mathcal F\longrightarrow0
\]
이 존재하는 $\mathcal O_X$-가군이라고 하자. 그러면 모든 $x\in X$에
대하여 $x$의 열린 근방 $U$가 존재하여 $\mathcal G|U$는
\emph{직합} $\mathcal E|U\oplus\mathcal F|U$와 동형이다. 실제로
$\mathcal F=\mathcal O_X^n$인 경우로 한정할 수 있다.
$e_i$ $(1\leq i\leq n)$를 $\mathcal O_X^n$의 표준 절단들(5.5.5)이라
하자. 그러면 $x$의 열린 근방 $U$와 $U$ 위의 $\mathcal G$의 절단
$s_i$ $(1\leq i\leq n)$가 존재하여 $p(s_i|U)=e_i|U$이다. 이제
$f$를 절단 $s_i|U$들이 정의하는 준동형
$\mathcal F|U\to\mathcal G|U$라 하자(5.1.1). 모든 열린집합
$V\subset U$와 모든 절단 $s\in\Gamma(V,\mathcal G)$에 대하여
$s-f(p(s))\in\Gamma(V,\mathcal E)$임은 곧 확인되고, 이로부터 주장이
따른다.
\end{env}

\begin{env}[5.4.10]
\label{0.5.4.10-ko}
$f:X\to Y$를 환 달린 공간의 사상, $\mathcal F$를
$\mathcal O_X$-가군, $\mathcal L$을 유한 계수의 국소 자유
$\mathcal O_Y$-가군이라 하자. 그러면 표준 동형
\[
  f_*(\mathcal F)\otimes_{\mathcal O_Y}\mathcal L
  \xrightarrow{\sim}
  f_*\bigl(\mathcal F\otimes_{\mathcal O_X}f^*(\mathcal L)\bigr)
  \tag{5.4.10.1}\label{0.5.4.10.1-ko}
\]
이 존재한다.

\oldpage[0\textsubscript{I}]{53}

실제로 임의의 $\mathcal O_Y$-가군 $\mathcal L$에 대하여 표준 준동형
\[
  f_*(\mathcal F)\otimes_{\mathcal O_Y}\mathcal L
  \xrightarrow{1\otimes\rho}
  f_*(\mathcal F)\otimes_{\mathcal O_Y}f_*(f^*(\mathcal L))
  \xrightarrow{\alpha}
  f_*\bigl(\mathcal F\otimes_{\mathcal O_X}f^*(\mathcal L)\bigr)
\]
이 있다. 여기서 $\rho$는 (4.4.3.2)의 준동형이고 $\alpha$는
(4.2.2.1)의 준동형이다. $\mathcal L$이 국소 자유일 때 이 준동형이
전단사임을 보이기 위해서는 문제가 $Y$ 위에서 국소적이므로
$\mathcal L=\mathcal O_Y^n$인 경우만 생각하면 충분하다. 더 나아가
$f_*$와 $f^*$는 유한 직합과 가환하므로 $n=1$이라고 가정할 수 있다.
이 경우 명제는 정의와 관계식 $f^*(\mathcal O_Y)=\mathcal O_X$에서
곧 따른다.
\end{env}

\subsection{국소환 달린 공간 위의 층}
\label{subsection:0.5.5-ko}

\begin{env}[5.5.1]
\label{0.5.5.1-ko}
환 달린 공간 $(X,\mathcal O_X)$가 \emph{국소환 달린 공간}이라는 것은
모든 $x\in X$에 대하여 $\mathcal O_x$가 국소환이라는 뜻이다. 이
논고에서 다룰 환 달린 공간은 거의 모두 이런 공간이다. 이때
$\mathcal O_x$의 \emph{극대 아이디얼}을 $\mathfrak m_x$로,
\emph{잉여체} $\mathcal O_x/\mathfrak m_x$를 $k(x)$로 쓴다. 모든
$\mathcal O_X$-가군 $\mathcal F$, $X$의 모든 열린집합 $U$, 모든 점
$x\in U$, 모든 절단 $f\in\Gamma(U,\mathcal F)$에 대하여 $f(x)$는
싹 $f_x$의
\[
  \mathcal F_x/\mathfrak m_x\mathcal F_x
\]
에서의 류를 뜻하며, 이를 점 $x$에서 $f$의 \emph{값}이라고 한다.
따라서 $f(x)=0$은
$f_x\in\mathfrak m_x\mathcal F_x$라는 뜻이다. 이 관계가 성립하면
언어의 남용으로 $f$가 \emph{$x$에서 소멸한다}고 한다. 이 관계와
$f_x=0$을 혼동하지 않도록 주의해야 한다.
\end{env}

\begin{env}[5.5.2]
\label{0.5.5.2-ko}
$X$를 국소환 달린 공간, $\mathcal L$을 가역
$\mathcal O_X$-가군, $f$를 $X$ 위의 $\mathcal L$의 절단이라 하자.
점 $x\in X$에서 다음 세 조건은 \emph{서로 동치}이다.
\begin{enumerate}
  \item[$a)$] \emph{$f_x$는 $\mathcal L_x$의 생성원이다.}
  \item[$b)$] \emph{$f_x\notin\mathfrak m_x\mathcal L_x$이다.} 다시
    말해 $f(x)\neq0$이다.
  \item[$c)$] $x$의 열린 근방 $V$ 위의 $\mathcal L^{-1}$의 절단
    $g$가 존재하여, $f\otimes g$의
    $\Gamma(V,\mathcal O_X)$에서의 표준적인 상(5.4.3)이
    \emph{단위 절단}이다.
\end{enumerate}

실제로 문제는 국소적이므로 $\mathcal L=\mathcal O_X$인 경우로
한정할 수 있다. 그러면 $a)$와 $b)$의 동치는 자명하고, $c)$가
$b)$를 함의함도 자명하다. 한편
$f_x\notin\mathfrak m_x$이면 $f_x$는 $\mathcal O_x$에서 가역이므로
$f_xg_x=1_x$인 $g_x$가 존재한다. 절단의 싹의 정의에 따라 이는
$x$의 근방 $V$와 $V$ 위의 $\mathcal O_X$의 절단 $g$가 존재하여
$V$에서 $fg=1$이라는 뜻이고, 따라서 $c)$가 성립한다.

조건 $c)$에서 곧, 동치인 조건 $a)$, $b)$, $c)$를 만족하는
$x$들의 집합 $X_f$가 $X$에서 \emph{열린집합}임을 알 수 있다.
(5.5.1)의 용어로 이는 $f$가 \emph{소멸하지 않는} 점들의 집합이다.
\end{env}

\begin{env}[5.5.3]
\label{0.5.5.3-ko}
(5.5.2)의 가정 아래 $\mathcal L'$을 두 번째 가역
$\mathcal O_X$-가군이라 하자. 그러면
$f\in\Gamma(X,\mathcal L)$, $g\in\Gamma(X,\mathcal L')$에 대하여
\[
  X_f\cap X_g=X_{f\otimes g}
\]
이다. 실제로 문제는 국소적이므로 곧
$\mathcal L=\mathcal L'=\mathcal O_X$인 경우로 환원된다.
이 경우 $f\otimes g$는 곱 $fg$와 표준적으로 동일시되므로 명제는
자명하다.
\end{env}

\oldpage[0\textsubscript{I}]{54}

\begin{env}[5.5.4]
\label{0.5.5.4-ko}
$\mathcal F$를 계수 $n$의 국소 자유 $\mathcal O_X$-가군이라 하자.
$p\leq n$이면 $\bigwedge^p\mathcal F$는 계수
$\binom{n}{p}$의 국소 자유 $\mathcal O_X$-가군이고, $p>n$이면
영가군임은 곧 따른다. 실제로 문제는 국소적이므로
$\mathcal F=\mathcal O_X^n$인 경우로 환원된다. 또한 모든 $x\in X$에
대하여
\[
  (\bigwedge^p\mathcal F)_x/
  \mathfrak m_x(\bigwedge^p\mathcal F)_x
\]
는 $k(x)$ 위의 차원 $\binom{n}{p}$인 벡터공간이며,
\[
  \bigwedge^p(\mathcal F_x/\mathfrak m_x\mathcal F_x)
\]
와 표준적으로 동일시된다. $s_1,\ldots,s_p$를 $X$의 열린집합 $U$
위의 $\mathcal F$의 절단들이라 하고,
$s=s_1\wedge\cdots\wedge s_p$라 하자. 이는 $U$ 위의
$\bigwedge^p\mathcal F$의 절단이다(4.1.5). 이때
$s(x)=s_1(x)\wedge\cdots\wedge s_p(x)$이므로,
$s_1(x),\ldots,s_p(x)$가 \emph{선형 종속}이라는 것은 $s(x)=0$이라는
뜻이다. 따라서 \emph{$s_1(x),\ldots,s_p(x)$가 선형 독립인
$x\in U$들의 집합은 $X$에서 열린집합이다.} 실제로
$\mathcal F=\mathcal O_X^n$인 경우로 환원한 뒤,
\[
  \bigwedge^p\mathcal F
  =\mathcal O_X^{\binom{n}{p}}
\]
에서 $\binom{n}{p}$개의 성분으로 가는 사영 중 하나에 의한 $s$의
상에 (5.5.2)를 적용하면 충분하다.

특히 $s_1,\ldots,s_n$이 $U$ 위의 $\mathcal F$의 $n$개 절단이고 모든
$x\in U$에서 $s_1(x),\ldots,s_n(x)$가 선형 독립이면, 이 절단들이
정의하는 준동형(5.1.1)
\[
  u:\mathcal O_X^n|U\longrightarrow\mathcal F|U
\]
는 \emph{동형}이다. 실제로
$\mathcal F=\mathcal O_X^n$이고
$\bigwedge^n\mathcal F$를 $\mathcal O_X$와 표준적으로 동일시한
경우로 한정할 수 있다. 그러면
$s=s_1\wedge\cdots\wedge s_n$은 $U$ 위의 $\mathcal O_X$의
\emph{가역 절단}이고, 크라메르 공식으로 $u$의 역준동형을 정의할 수
있다.
\end{env}

\begin{env}[5.5.5]
\label{0.5.5.5-ko}
$\mathcal E$, $\mathcal F$를 유한 계수의 두 국소 자유
$\mathcal O_X$-가군, $u:\mathcal E\to\mathcal F$를 준동형이라 하자.
점 $x\in X$의 근방 $U$가 존재하여 $u|U$가 \emph{단사}이고
$\mathcal F|U$가 $u(\mathcal E)|U$와 어떤 국소 자유
$(\mathcal O_X|U)$-부분가군 $\mathcal G$의 \emph{직합}일
필요충분조건은, $u_x:\mathcal E_x\to\mathcal F_x$가 몫을 취하여
벡터공간의 \emph{단사} 준동형
\[
  \mathcal E_x/\mathfrak m_x\mathcal E_x
  \longrightarrow
  \mathcal F_x/\mathfrak m_x\mathcal F_x
\]
을 유도하는 것이다.

이 조건은 \emph{필요}하다. 실제로 이때 $\mathcal F_x$는 자유
$\mathcal O_x$-가군 $u_x(\mathcal E_x)$와 $\mathcal G_x$의 직합이고,
따라서 $\mathcal F_x/\mathfrak m_x\mathcal F_x$는
\[
  u_x(\mathcal E_x)/\mathfrak m_xu_x(\mathcal E_x)
  \quad\text{및}\quad
  \mathcal G_x/\mathfrak m_x\mathcal G_x
\]
의 직합이다. 이 조건은 \emph{충분}하기도 하다. 실제로
$\mathcal E=\mathcal O_X^m$인 경우로 한정할 수 있다.
$e_i$를 모든 $y\in X$에서 $(e_i)_y$가 $\mathcal O_y^m$의 표준기저의
$i$번째 원소가 되는 $\mathcal O_X^m$의 절단, 즉
$\mathcal O_X^m$의 \emph{표준 절단}이라 하고,
$s_1,\ldots,s_m$을 $u$에 의한 $e_i$들의 상이라 하자. 가정에 따라
$s_1(x),\ldots,s_m(x)$는 선형 독립이다. $\mathcal F$의 계수가
$n$이면 $x$의 어떤 근방 $V$ 위의 $\mathcal F$의 절단
$s_{m+1},\ldots,s_n$이 존재하여 $s_i(x)$ $(1\leq i\leq n)$가
$\mathcal F_x/\mathfrak m_x\mathcal F_x$의 기저를 이룬다. 그러면
(5.5.4)에 따라 $x$의 근방 $U\subset V$가 존재하여 모든 $y\in U$에서
$s_i(y)$ $(1\leq i\leq n)$가
$\mathcal F_y/\mathfrak m_y\mathcal F_y$의 기저를 이룬다. 다시
(5.5.4)에 따라 $s_i|U$ $(1\leq i\leq m)$를 $e_i|U$로 보내는 동형
\[
  \mathcal F|U\xrightarrow{\sim}\mathcal O_X^n|U
\]
이 존재하고, 이로써 증명이 끝난다.
\end{env}

\section{평탄성}
\label{section:0.6-ko}

\begin{env}[6.0]
\label{0.6.0.0-ko}
평탄성의 개념은 J.-P.~세르[16]에게서 비롯되었다. 이하에서는
N.~부르바키의 \emph{가환대수}에 제시된 결과들의 증명을 생략하고
독자를 그 책으로 보낸다. 모든 환은 가환이라고 가정한다.\footnote{그
결과들의 대부분을 비가환인 경우로 일반화한 내용은 앞서 인용한
N.~부르바키의 강연을 보라.}

\oldpage[0\textsubscript{I}]{55}

$M$, $N$을 두 $A$-가군, $M'$을 $M$의 부분가군, $N'$을 $N$의
부분가군이라 하자. 표준 사상
\[
  M'\otimes_A N'\longrightarrow M\otimes_A N
\]
의 상인 $M\otimes_A N$의 부분가군을
$\operatorname{Im}(M'\otimes_A N')$로 쓴다.
\end{env}

\subsection{평탄 가군}
\label{subsection:0.6.1-ko}

\begin{env}[6.1.1]
\label{0.6.1.1-ko}
$M$을 $A$-가군이라 하자. 다음 조건들은 서로 동치이다.
\begin{enumerate}
  \item[$a)$] $N$에 관한 함자 $M\otimes_A N$은 $A$-가군의 범주에서
    완전하다.
  \item[$b)$] 모든 $i>0$과 모든 $A$-가군 $N$에 대하여
    $\operatorname{Tor}_i^A(M,N)=0$이다.
  \item[$c)$] 모든 $A$-가군 $N$에 대하여
    $\operatorname{Tor}_1^A(M,N)=0$이다.
\end{enumerate}

$M$이 이 조건들을 만족하면 $M$을 \emph{평탄 $A$-가군}이라 한다.
모든 자유 $A$-가군은 평탄함이 자명하다.

모든 유한 생성 아이디얼 $\mathfrak J$에 대하여 표준 사상
\[
  M\otimes_A\mathfrak J\longrightarrow M\otimes_A A=M
\]
이 \emph{단사}이면 $M$은 평탄 $A$-가군이다.
\end{env}

\begin{env}[6.1.2]
\label{0.6.1.2-ko}
평탄 $A$-가군들의 모든 귀납극한은 평탄 $A$-가군이다.
$A$-가군들의 직합 $\bigoplus_{\lambda\in L}M_\lambda$가 평탄
$A$-가군일 필요충분조건은 각 $A$-가군 $M_\lambda$가 평탄한
것이다. 특히 모든 사영 $A$-가군은 평탄하다.

$A$-가군들의 완전열
\[
  0\longrightarrow M'\longrightarrow M\longrightarrow M''
  \longrightarrow 0
\]
에서 $M''$이 \emph{평탄}하다고 하자. 그러면 모든 $A$-가군 $N$에
대하여 완전열
\[
  0\longrightarrow M'\otimes N\longrightarrow M\otimes N
  \longrightarrow M''\otimes N\longrightarrow 0
\]
을 얻는다. 또한 이때 $M$이 평탄할 필요충분조건은 $M'$이 평탄한
것이다. 다만 $M$과 $M'$이 평탄하지만 $M''=M/M'$은 평탄하지 않을
수 있다.
\end{env}

\begin{env}[6.1.3]
\label{0.6.1.3-ko}
$M$을 평탄 $A$-가군, $N$을 임의의 $A$-가군이라 하자. $N$의 두
부분가군 $N'$, $N''$에 대하여
\[
  \operatorname{Im}(M\otimes(N'+N''))
  =\operatorname{Im}(M\otimes N')+\operatorname{Im}(M\otimes N''),
\]
\[
  \operatorname{Im}(M\otimes(N'\cap N''))
  =\operatorname{Im}(M\otimes N')\cap\operatorname{Im}(M\otimes N'')
\]
이다. 여기서 상들은 $M\otimes N$ 안에서 취한다.
\end{env}

\begin{env}[6.1.4]
\label{0.6.1.4-ko}
$M$, $N$을 두 $A$-가군, $M'$을 $M$의 부분가군, $N'$을 $N$의
부분가군이라 하고, $M/M'$과 $N/N'$ 가운데 하나가 평탄하다고
가정하자. 그러면
\[
  \operatorname{Im}(M'\otimes N')
  =\operatorname{Im}(M'\otimes N)\cap
   \operatorname{Im}(M\otimes N')
\]
이다. 여기서 상들은 $M\otimes N$ 안에서 취한다. 특히
$\mathfrak J$가 $A$의 아이디얼이고 $M/M'$이 평탄하면
$\mathfrak JM'=M'\cap\mathfrak JM$이다.
\end{env}

\subsection{환의 변경}
\label{subsection:0.6.2-ko}

가법군 $M$에 환 $A$, $B$, \ldots{}에 관한 여러 가군 구조가 주어진
경우, $M$이 $A$-가군, $B$-가군, \ldots{}으로서 평탄하다고 하는 대신
때로는 $M$이 각각 \emph{$A$-평탄}, \emph{$B$-평탄}, \ldots{}이라고도
한다.

\begin{env}[6.2.1]
\label{0.6.2.1-ko}
$A$, $B$를 두 환, $M$을 $A$-가군, $N$을 $(A,B)$-쌍가군이라 하자.
$M$이 평탄하고 $N$이 $B$-평탄하면 $M\otimes_A N$은
$B$-평탄하다. 특히 $M$, $N$이 두 평탄 $A$-가군이면
$M\otimes_A N$은 평탄 $A$-가군이다. $B$가 $A$-대수이고 $M$이

\oldpage[0\textsubscript{I}]{56}

평탄 $A$-가군이면 $B$-가군 $M_{(B)}=M\otimes_A B$는 평탄하다.
끝으로 $B$가 $A$-가군으로서 평탄한 $A$-대수이고 $N$이 평탄
$B$-가군이면 $N$은 $A$-평탄이기도 하다.
\end{env}

\begin{env}[6.2.2]
\label{0.6.2.2-ko}
$A$를 환, $B$를 $A$-가군으로서 평탄한 $A$-대수라 하자.
$M$, $N$을 두 $A$-가군이라 하고 $M$이 유한 표시를 갖는다고 하자.
그러면 표준 준동형
\[
  \operatorname{Hom}_A(M,N)\otimes_A B
  \longrightarrow
  \operatorname{Hom}_B(M\otimes_A B,N\otimes_A B)
  \tag{6.2.2.1}
\]
은 동형이다. 이 준동형은 $u\otimes b$를
\[
  m\otimes b'\longmapsto u(m)\otimes b'b
\]
인 준동형으로 보낸다.
\end{env}

\begin{env}[6.2.3]
\label{0.6.2.3-ko}
$(A_\lambda,\varphi_{\mu\lambda})$를 환들의 여과 귀납계라 하고
$A=\varinjlim A_\lambda$라 하자. 또한 각 $\lambda$에 대하여
$M_\lambda$를 $A_\lambda$-가군이라 하고, $\lambda\leq\mu$일 때
$\theta_{\mu\lambda}:M_\lambda\to M_\mu$를
$\varphi_{\mu\lambda}$-준동형이라 하되
$(M_\lambda,\theta_{\mu\lambda})$가 귀납계를 이룬다고 하자.
그러면 $M=\varinjlim M_\lambda$는 $A$-가군이다.

이제 모든 $\lambda$에 대하여 $M_\lambda$가 평탄
$A_\lambda$-가군이면 $M$은 평탄 $A$-가군이다. 실제로
$\mathfrak J$를 $A$의 \emph{유한 생성} 아이디얼이라 하자.
귀납극한의 정의에 따라 어떤 지표 $\lambda$와
$A_\lambda$의 아이디얼 $\mathfrak J_\lambda$가 존재하여
$\mathfrak J=\mathfrak J_\lambda A$이다.
$\mu\geq\lambda$일 때
$\mathfrak J'_\mu=\mathfrak J_\lambda A_\mu$로 놓으면
$\mathfrak J=\varinjlim\mathfrak J'_\mu$이기도 하다. 여기서
$\mu$는 $\lambda$ 이상인 지표들을 달린다. 귀납극한 함자가
완전하고 텐서곱과 가환하므로
\[
  M\otimes_A\mathfrak J
  =\varinjlim(M_\mu\otimes_{A_\mu}\mathfrak J'_\mu)
  =\varinjlim\mathfrak J'_\mu M_\mu
  =\mathfrak JM.
\]
\end{env}

\subsection{평탄성의 국소화}
\label{subsection:0.6.3-ko}

\begin{env}[6.3.1]
\label{0.6.3.1-ko}
$A$를 환, $S$를 $A$의 곱셈적 부분집합이라 하면 $S^{-1}A$는
\emph{평탄} $A$-가군이다. 실제로 모든 $A$-가군 $N$에 대하여
$N\otimes_A S^{-1}A$는 $S^{-1}N$과 동일시되고(1.2.5),
$N$에 관한 함자 $S^{-1}N$은 완전함을 안다(1.3.2).

이제 $M$이 평탄 $A$-가군이면
$S^{-1}M=M\otimes_A S^{-1}A$는 평탄 $S^{-1}A$-가군이고(6.2.1),
앞의 결과와 (6.2.1)에 따라 $A$-평탄이기도 하다. 특히 $P$가
$S^{-1}A$-가군이면 $P$를 $S^{-1}P$와 동형인 $A$-가군으로 볼 수
있으며, $P$가 $A$-평탄일 필요충분조건은 $S^{-1}A$-평탄인 것이다.
\end{env}

\begin{env}[6.3.2]
\label{0.6.3.2-ko}
$A$를 환, $B$를 $A$-대수, $T$를 $B$의 곱셈적 부분집합이라 하자.
$P$가 $A$-평탄인 $B$-가군이면 $T^{-1}P$는 $A$-평탄이다. 실제로
모든 $A$-가군 $N$에 대하여
\[
  (T^{-1}P)\otimes_A N
  =(T^{-1}B\otimes_B P)\otimes_A N
  =T^{-1}B\otimes_B(P\otimes_A N)
  =T^{-1}(P\otimes_A N).
\]
그런데 $T^{-1}(P\otimes_A N)$은 두 완전 함자, 즉 $N$에 관한
$P\otimes_A N$과 $Q$에 관한 $T^{-1}Q$의 합성이므로 $N$에 관한
완전 함자이다. $S$가 $A$의 곱셈적 부분집합이고 $B$에서의 그 상이
\emph{$T$에 포함}되면
$T^{-1}P=S^{-1}(T^{-1}P)$이므로 (6.3.1)에 따라
$T^{-1}P$는 $S^{-1}A$-평탄이기도 하다.
\end{env}

\begin{env}[6.3.3]
\label{0.6.3.3-ko}
$\varphi:A\to B$를 환의 준동형, $M$을 $B$-가군이라 하자.
다음 조건들은 서로 동치이다.
\begin{enumerate}
  \item[$a)$] $M$은 평탄 $A$-가군이다.
  \item[$b)$] $B$의 모든 극대 아이디얼 $\mathfrak n$에 대하여
    $M_\mathfrak n$은 평탄 $A$-가군이다.
  \item[$c)$] $B$의 모든 극대 아이디얼 $\mathfrak n$에 대하여
    $\mathfrak m=\varphi^{-1}(\mathfrak n)$으로 놓으면
    $M_\mathfrak n$은 평탄 $A_\mathfrak m$-가군이다.
\end{enumerate}

$M_\mathfrak n=(M_\mathfrak n)_\mathfrak m$이므로 $b)$와 $c)$의
동치는 (6.3.1)에서 따르고, $a)$가 $b)$를 함의하는 것은 (6.3.2)의
특수한 경우이다. 이제 $b)$가 $a)$를 함의함을 보이면 된다. 즉

\oldpage[0\textsubscript{I}]{57}

$A$-가군의 모든 단사 준동형 $u:N'\to N$에 대하여 준동형
\[
  v=1\otimes u:M\otimes_A N'\longrightarrow M\otimes_A N
\]
이 단사임을 보이면 된다. $v$는 $B$-가군의 준동형이기도 하므로,
모든 극대 아이디얼 $\mathfrak n$에 대하여
\[
  v_\mathfrak n:(M\otimes_A N')_\mathfrak n
  \longrightarrow(M\otimes_A N)_\mathfrak n
\]
이 단사이면 $v$가 단사임을 안다. 그런데
\[
  (M\otimes_A N)_\mathfrak n
  =B_\mathfrak n\otimes_B(M\otimes_A N)
  =M_\mathfrak n\otimes_A N
\]
이므로 $v_\mathfrak n$은 바로
\[
  1\otimes u:M_\mathfrak n\otimes_A N'
  \longrightarrow M_\mathfrak n\otimes_A N
\]
이고, $M_\mathfrak n$이 $A$-평탄이므로 이는 단사이다.

특히 $B=A$로 놓으면, $A$-가군 $M$이 평탄일 필요충분조건은
$A$의 모든 극대 아이디얼 $\mathfrak m$에 대하여
$M_\mathfrak m$이 $A_\mathfrak m$-평탄인 것이다.
\end{env}

\begin{env}[6.3.4]
\label{0.6.3.4-ko}
$M$을 $A$-가군이라 하자. $M$이 평탄하고 $f\in A$가 $A$에서
영인자가 아니면 $f$는 $M$의 영이 아닌 어떤 원소도 소멸시키지
않는다. 실제로 준동형 $m\mapsto f\cdot m$은 $1\otimes u$로
쓸 수 있다. 여기서 $u$는 $A$에서의 곱셈
$a\mapsto f\cdot a$이고 $M$은 $M\otimes_A A$와 동일시한다.
$u$가 단사이면 $M$의 평탄성에 따라 $1\otimes u$도 단사이다.
특히 $A$가 \emph{정역}이면 $M$은 \emph{꼬임이 없다}.

역으로 $A$가 정역이고 $M$에 꼬임이 없으며, $A$의 모든 극대
아이디얼 $\mathfrak m$에 대하여 $A_\mathfrak m$이
\emph{이산 값매김환}이라고 하자. 그러면 $M$은 $A$-평탄이다.
실제로 (6.3.3)에 따라 $M_\mathfrak m$이
$A_\mathfrak m$-평탄임을 보이면 충분하므로, $A$ 자체가 이미
이산 값매김환이라고 가정할 수 있다. 또한 $M$은 그 유한 생성
부분가군들의 귀납극한이고 이 부분가군들에도 꼬임이 없으므로,
(6.1.2)에 따라 $M$이 유한 생성인 경우로 한정할 수 있다.
이 경우 $M$이 자유 $A$-가군이므로 명제가 따른다.

특히 $A$가 \emph{정역}이고 환의 준동형
$\varphi:A\to B$가 $B$를 $0$이 아닌 평탄 $A$-가군으로 만들면,
$\varphi$는 반드시 \emph{단사}이다. 역으로 $B$가 정역이고
$A$가 $B$의 부분환이며, $A$의 모든 극대 아이디얼
$\mathfrak m$에 대하여 $A_\mathfrak m$이 이산 값매김환이면
$B$는 $A$-평탄이다.
\end{env}

\subsection{충실 평탄 가군}
\label{subsection:0.6.4-ko}

\begin{env}[6.4.1]
\label{0.6.4.1-ko}
$A$-가군 $M$에 관한 다음 네 조건은 서로 동치이다.
\begin{enumerate}
  \item[$a)$] $A$-가군의 열 $N'\to N\to N''$이 완전할 필요충분조건은
    열
    $M\otimes_A N'\to M\otimes_A N\to M\otimes_A N''$이 완전한 것이다.
  \item[$b)$] $M$은 평탄하고, 모든 $A$-가군 $N$에 대하여
    $M\otimes_A N=0$이면 $N=0$이다.
  \item[$c)$] $M$은 평탄하고, $A$-가군의 모든 준동형 $v:N\to N'$에
    대하여 $1_M\otimes v=0$이면 $v=0$이다. 여기서 $1_M$은 $M$의
    항등 자기동형이다.
  \item[$d)$] $M$은 평탄하고, $A$의 모든 극대 아이디얼
    $\mathfrak m$에 대하여 $\mathfrak mM\neq M$이다.
\end{enumerate}

$M$이 이 조건들을 만족할 때 $M$을 \emph{충실 평탄} $A$-가군이라
한다. 이때 $M$은 반드시 \emph{충실} 가군이다. 또한
$u:N\to N'$이 $A$-가군의 준동형이면, $u$가 단사(각각 전사,
전단사)일 필요충분조건은
$1\otimes u:M\otimes_A N\to M\otimes_A N'$이 그러한 것이다.
\end{env}

\oldpage[0\textsubscript{I}]{58}

\begin{env}[6.4.2]
\label{0.6.4.2-ko}
$0$이 아닌 자유 가군은 충실 평탄하다. 평탄 가군과 충실 평탄
가군의 직합도 충실 평탄하다. $S$가 $A$의 곱셈적 부분집합이면,
$S^{-1}A$가 충실 평탄 $A$-가군일 수 있는 것은 $S$가 가역원들로
이루어진 경우뿐이다. 따라서 이 경우 $S^{-1}A=A$이다.
\end{env}

\begin{env}[6.4.3]
\label{0.6.4.3-ko}
$A$-가군의 완전열
\[
  0\longrightarrow M'\longrightarrow M\longrightarrow M''\longrightarrow 0
\]
이 주어졌다고 하자. $M'$과 $M''$이 평탄하고 그 가운데 하나가
충실 평탄이면 $M$은 충실 평탄이다.
\end{env}

\begin{env}[6.4.4]
\label{0.6.4.4-ko}
$A,B$를 두 환, $M$을 $A$-가군, $N$을 $(A,B)$-쌍가군이라 하자.
$M$이 충실 평탄이고 $N$이 충실 평탄 $B$-가군이면
$M\otimes_A N$은 충실 평탄 $B$-가군이다. 특히 $M,N$이 두 충실
평탄 $A$-가군이면 $M\otimes_A N$도 충실 평탄하다. $B$가
$A$-대수이고 $M$이 충실 평탄 $A$-가군이면 $B$-가군
$M_{(B)}$는 충실 평탄하다.
\end{env}

\begin{env}[6.4.5]
\label{0.6.4.5-ko}
$M$이 충실 평탄 $A$-가군이고 $S$가 $A$의 곱셈적 부분집합이면,
$S^{-1}M$은 충실 평탄 $S^{-1}A$-가군이다. 실제로
$S^{-1}M=M\otimes_A(S^{-1}A)$이므로 이는 (6.4.4)에서 따른다.
역으로 $A$의 모든 극대 아이디얼 $\mathfrak m$에 대하여
$M_\mathfrak m$이 충실 평탄 $A_\mathfrak m$-가군이면 $M$은 충실
평탄 $A$-가군이다. 실제로 $M$은 $A$-평탄이고(6.3.3),
\[
  M_\mathfrak m/\mathfrak mM_\mathfrak m
  =(M\otimes_A A_\mathfrak m)
  \otimes_{A_\mathfrak m}
  (A_\mathfrak m/\mathfrak mA_\mathfrak m)
  =M\otimes_A(A/\mathfrak m)
  =M/\mathfrak mM.
\]
따라서 가정에 의해 $A$의 모든 극대 아이디얼 $\mathfrak m$에
대하여 $M/\mathfrak mM\neq0$이고, 이로써 명제가 따른다(6.4.1).
\end{env}

\subsection{스칼라 제한}
\label{subsection:0.6.5-ko}

\begin{env}[6.5.1]
\label{0.6.5.1-ko}
$A$를 환, $\varphi:A\to B$를 $B$에 $A$-대수 구조를 주는 환의
준동형이라 하자. 어떤 $B$-가군 $N$이 충실 평탄 $A$-가군이라고
가정하자. 그러면 모든 $A$-가군 $M$에 대하여
$x\mapsto1\otimes x$로 주어지는 $M$에서
$B\otimes_A M=M_{(B)}$로의 준동형은
\emph{단사}이다. 특히 $\varphi$는 단사이고, $A$의 모든 아이디얼
$\mathfrak a$에 대하여
$\varphi^{-1}(\mathfrak aB)=\mathfrak a$이다. 또한 $A$의 모든 극대
아이디얼(각각 소 아이디얼) $\mathfrak m$에 대하여
$\varphi^{-1}(\mathfrak n)=\mathfrak m$을 만족하는 $B$의 극대
아이디얼(각각 소 아이디얼) $\mathfrak n$이 존재한다.
\end{env}

\begin{env}[6.5.2]
\label{0.6.5.2-ko}
(6.5.1)의 조건이 만족되면 $\varphi$를 통하여 $A$를 $B$의 부분환과
동일시하고, 더 일반적으로 모든 $A$-가군 $M$을 $M_{(B)}$의
$A$-부분가군과 동일시한다. 이때 $B$가 \emph{뇌터} 환이면 $A$도
뇌터 환이다. 실제로 사상
$\mathfrak a\mapsto\mathfrak aB$는 $A$의 아이디얼 전체에서
$B$의 아이디얼 전체로 가는 순서 보존
단사이다. 따라서 $A$의 아이디얼들로 이루어진 무한한 엄밀 증가열이
존재하면 $B$의 아이디얼들로 이루어진 같은 종류의 증가열도
존재하게 된다.
\end{env}

\subsection{충실 평탄환}
\label{subsection:0.6.6-ko}

\begin{env}[6.6.1]
\label{0.6.6.1-ko}
$\varphi:A\to B$를 $B$에 $A$-대수 구조를 주는 환의 준동형이라
하자. 다음 다섯 조건은 서로 동치이다.
\begin{enumerate}
  \item[$a)$] $B$는 충실 평탄 $A$-가군이다. 다시 말해 $M$에 관한
    함자 $M_{(B)}$는 완전하고 충실하다.
  \item[$b)$] 준동형 $\varphi$는 단사이고 $A$-가군
    $B/\varphi(A)$는 평탄하다.

\oldpage[0\textsubscript{I}]{59}

  \item[$c)$] $A$-가군 $B$는 평탄하다. 다시 말해 함자 $M_{(B)}$는
    완전하다. 또한 모든 $A$-가군 $M$에 대하여 $M$에서
    $M_{(B)}$로 가는 준동형 $x\mapsto1\otimes x$는 단사이다.
  \item[$d)$] $A$-가군 $B$는 평탄하고, $A$의 모든 아이디얼
    $\mathfrak a$에 대하여
    $\varphi^{-1}(\mathfrak aB)=\mathfrak a$이다.
  \item[$e)$] $A$-가군 $B$는 평탄하고, $A$의 모든 극대 아이디얼
    $\mathfrak m$에 대하여
    $\varphi^{-1}(\mathfrak n)=\mathfrak m$을 만족하는 $B$의 극대
    아이디얼 $\mathfrak n$이 존재한다.
\end{enumerate}

이 조건들이 만족되면 $A$를 $B$의 부분환과 동일시한다.
\end{env}

\begin{env}[6.6.2]
\label{0.6.6.2-ko}
$A$를 국소환, $\mathfrak m$을 그 극대 아이디얼, $B$를
$\mathfrak mB\neq B$인 $A$-대수라 하자. 예를 들어 $B$가
국소환이고 구조 준동형 $\varphi:A\to B$가 국소 준동형이면 이
조건이 성립한다. $B$가 평탄 $A$-가군이면 $B$는 충실 평탄
$A$-가군이다. 실제로 $\mathfrak mB\neq B$이므로
$\mathfrak mB$를 포함하는 $B$의 극대 아이디얼 $\mathfrak n$이
존재한다. $\varphi^{-1}(\mathfrak n)\cap A$는 $\mathfrak m$을
포함하고 $1$을 포함하지 않으므로
$\varphi^{-1}(\mathfrak n)=\mathfrak m$이다. 따라서 (6.6.1)의
조건 $e)$를 적용할 수 있다. 이 조건들 아래에서 $B$가 뇌터 환이면
$A$도 뇌터 환임을 알 수 있다(6.5.2).
\end{env}

\begin{env}[6.6.3]
\label{0.6.6.3-ko}
$B$를 충실 평탄 $A$-가군인 $A$-대수라 하자. 모든 $A$-가군 $M$과
그 $A$-부분가군 $M'$에 대하여, $M$을 $M_{(B)}$의
$A$-부분가군과 동일시하면 $M'=M\cap M'_{(B)}$이다.
$M$이 평탄(각각 충실 평탄) $A$-가군일 필요충분조건은
$M_{(B)}$가 평탄(각각 충실 평탄) $B$-가군인 것이다.
\end{env}

\begin{env}[6.6.4]
\label{0.6.6.4-ko}
$B$를 $A$-대수, $N$을 충실 평탄 $B$-가군이라 하자. $B$가
평탄(각각 충실 평탄) $A$-가군일 필요충분조건은 $N$이
평탄(각각 충실 평탄) $A$-가군인 것이다.

특히 $C$를 $B$-대수라 하자. $C$가 $B$ 위에서 충실 평탄이고
$B$가 $A$ 위에서 충실 평탄이면 $C$는 $A$ 위에서 충실 평탄이다.
또한 $C$가 $B$ 위와 $A$ 위에서 모두 충실 평탄이면 $B$는 $A$
위에서 충실 평탄이다.
\end{env}

\subsection{환 달린 공간의 평탄 사상}
\label{subsection:0.6.7-ko}

\begin{env}[6.7.1]
\label{0.6.7.1-ko}
$f:X\to Y$를 환 달린 공간의 사상, $\mathcal F$를
$\mathcal O_X$-가군이라 하자. 점 $x\in X$에서 $\mathcal F_x$가
평탄 $\mathcal O_{f(x)}$-가군이면 $\mathcal F$가 $x$에서
\emph{$f$-평탄}이라고 한다. $f$에 관해 혼동할 염려가 없을 때에는
$Y$-평탄이라고도 한다. $y\in Y$ 위에서 $\mathcal F$가
$f$-평탄하다는 것은 모든 $x\in f^{-1}(y)$에서 $f$-평탄하다는
뜻이고, $\mathcal F$가 $f$-평탄하다는 것은 $X$의 모든 점에서
$f$-평탄하다는 뜻이다. $\mathcal O_X$가 각각 $x\in X$에서
$f$-평탄, $y\in Y$ 위에서 $f$-평탄, 또는 $f$-평탄일 때 사상
$f$가 각각 $x$에서 평탄, $y$ 위에서 평탄, 또는 평탄하다고 한다.
\end{env}

\begin{env}[6.7.2]
\label{0.6.7.2-ko}
(6.7.1)의 표기를 사용하자. $\mathcal F$가 $x$에서 $f$-평탄이면,
$y=f(x)$의 모든 열린 근방 $U$에 대하여 $\mathcal G$에 관한 함자
$\bigl(f^*(\mathcal G)\otimes_{\mathcal O_X}\mathcal F\bigr)_x$는
$(\mathcal O_Y|U)$-가군의 범주에서 완전하다. 실제로 이 줄기는
$\mathcal G_y\otimes_{\mathcal O_y}\mathcal F_x$와 표준적으로
동일시되므로 이는 정의에서 바로 따른다. 특히
$f$가 평탄 사상이면 역상함자 $f^*$는 $\mathcal O_Y$-가군의
범주에서 완전하다.
\end{env}

\begin{env}[6.7.3]
\label{0.6.7.3-ko}
역으로 환층 $\mathcal O_Y$가 연접이라고 가정하고, $y$의 모든 열린
근방 $U$에 대하여 $\mathcal G$에 관한 함자
$\bigl(f^*(\mathcal G)\otimes_{\mathcal O_X}\mathcal F\bigr)_x$가
연접 $(\mathcal O_Y|U)$-가군의 범주에서 완전하다고 가정하자.
그러면 $\mathcal F$는 $x$에서 $f$-평탄이다. 실제로
$\mathcal O_y$의 모든 유한 생성 아이디얼 $\mathfrak J$에 대하여
표준 준동형
$\mathfrak J\otimes_{\mathcal O_y}\mathcal F_x\to\mathcal F_x$가
단사임을 보이면 충분하다(6.1.1). 그런데 (5.3.8)에 따라

\oldpage[0\textsubscript{I}]{60}

$y$의 어떤 열린 근방 $U$와 $\mathcal O_Y|U$의 연접 아이디얼층
$\mathcal J$가 존재하여 $\mathcal J_y=\mathfrak J$이므로 결론이
따른다.
\end{env}

\begin{env}[6.7.4]
\label{0.6.7.4-ko}
평탄 가군에 관한 (6.1)의 결과들은 한 점에서 $f$-평탄한 층들에
관한 명제로 곧바로 옮겨진다.

$\mathcal O_X$-가군의 완전열
$0\to\mathcal F'\to\mathcal F\to\mathcal F''\to0$이 주어지고
$\mathcal F''$가 점 $x\in X$에서 $f$-평탄이라고 하자.
그러면 $y=f(x)$의 모든 열린 근방 $U$와 모든
$(\mathcal O_Y|U)$-가군 $\mathcal G$에 대하여 열
\[
  0\longrightarrow
  \bigl(f^*(\mathcal G)\otimes_{\mathcal O_X}\mathcal F'\bigr)_x
  \longrightarrow
  \bigl(f^*(\mathcal G)\otimes_{\mathcal O_X}\mathcal F\bigr)_x
  \longrightarrow
  \bigl(f^*(\mathcal G)\otimes_{\mathcal O_X}\mathcal F''\bigr)_x
  \longrightarrow0
\]
은 완전하다. 이때 $\mathcal F$가 $x$에서 $f$-평탄일
필요충분조건은 $\mathcal F'$가 그러한 것이다. $y\in Y$ 위에서
$f$-평탄한 $\mathcal O_X$-가군과 전역적으로 $f$-평탄한
$\mathcal O_X$-가군에 대해서도 각각 같은 명제가 성립한다.
\end{env}

\begin{env}[6.7.5]
\label{0.6.7.5-ko}
$f:X\to Y$와 $g:Y\to Z$를 환 달린 공간의 두 사상이라 하고,
$x\in X$, $y=f(x)$라 하자. $\mathcal F$를 $\mathcal O_X$-가군이라
하자. $\mathcal F$가 점 $x$에서 $f$-평탄이고 사상 $g$가 점
$y$에서 평탄이면 $\mathcal F$는 $x$에서 $(g\circ f)$-평탄이다
(6.2.1). 특히 $f$와 $g$가 평탄 사상이면 $g\circ f$도 평탄하다.
\end{env}

\begin{env}[6.7.6]
\label{0.6.7.6-ko}
$X,Y$를 환 달린 두 공간, $f:X\to Y$를 평탄 사상이라 하자.
쌍함자들의 표준 준동형(4.4.6)
\[
  f^*(\mathcal{H}om_{\mathcal O_Y}(\mathcal F,\mathcal G))
  \longrightarrow
  \mathcal{H}om_{\mathcal O_X}(f^*(\mathcal F),f^*(\mathcal G))
  \tag{6.7.6.1}
  \label{0.6.7.6.1-ko}
\]
은 $\mathcal F$가 유한 표시를 가지면 동형이다(5.2.5).

실제로 문제는 국소적이므로 완전열
$\mathcal O_Y^m\to\mathcal O_Y^n\to\mathcal F\to0$이 존재한다고
가정할 수 있다. $f$에 관한 가정에 따라 (6.7.6.1)의 양변은
$\mathcal F$에 관한 왼쪽 완전 함자이다. 따라서
$\mathcal F=\mathcal O_Y$인 경우로 귀착되며, 이 경우 명제는
자명하다.
\end{env}

\begin{env}[6.7.8]
\label{0.6.7.8-ko}
환 달린 공간의 사상 $f:X\to Y$가 전사이고, 모든 $x\in X$에
대하여 $\mathcal O_x$가 충실 평탄 $\mathcal O_{f(x)}$-가군이면
$f$를 \emph{충실 평탄} 사상이라고 한다. $X,Y$가 국소환 달린
공간이면 이는 $f$가 전사이고 평탄하다는 것과 동치이다(6.6.2).
$f$가 충실 평탄이면 $f^*$는 $\mathcal O_Y$-가군의 범주에서
완전하고 충실한 함자이고(6.6.1, $a)$), $\mathcal O_Y$-가군
$\mathcal G$가 $Y$-평탄일 필요충분조건은 $f^*(\mathcal G)$가
$X$-평탄인 것이다(6.6.3).
\end{env}

\section{아딕환}
\label{section:0.7-ko}

\subsection{허용환}
\label{subsection:0.7.1-ko}

\begin{env}[7.1.1]
\label{0.7.1.1-ko}
위상환 $A$(반드시 분리일 필요는 없다)에서 수열
$(x^n)_{n\geq0}$이 $0$을 극한으로 가지면 원소 $x$가
\emph{위상적으로 멱영}이라고 한다. $A$에서 $0$의 근방 기본계
가운데 아이디얼들(따라서 열린 아이디얼들)로 이루어진 것이 있으면
위상환 $A$를 \emph{선형 위상환}이라고 한다.
\end{env}

\begin{definition}[7.1.2]
\label{0.7.1.2-ko}
선형 위상환 $A$에서 아이디얼 $\mathfrak J$가 열려 있고, $0$의
모든 근방 $V$에 대하여 어떤 정수 $n>0$가 존재하여

\oldpage[0\textsubscript{I}]{61}

$\mathfrak J^n\subset V$이면 $\mathfrak J$를 \emph{정의 아이디얼}이라고
한다. 이를 흔히 수열 $(\mathfrak J^n)$이 $0$으로 간다고 표현한다.
$A$에 정의 아이디얼이 존재하면 선형 위상환 $A$를
\emph{준허용환}이라고 한다. $A$가 준허용환이고 더 나아가 분리이며
완비이면 $A$를 \emph{허용환}이라고 한다.
\end{definition}

$\mathfrak J$가 정의 아이디얼이고 $\mathfrak L$이 $A$의 열린
아이디얼이면 $\mathfrak J\cap\mathfrak L$도 정의 아이디얼임은
명백하다. 따라서 준허용환 $A$의 정의 아이디얼들은 $0$의 근방
기본계를 이룬다.

\begin{lemma}[7.1.3]
\label{0.7.1.3-ko}
$A$를 선형 위상환이라 하자.
\begin{enumerate}
  \item[{\upshape(i)}] $x\in A$가 위상적으로 멱영일 필요충분조건은
    $A$의 모든 열린 아이디얼 $\mathfrak J$에 대하여 $x$의
    $A/\mathfrak J$에서의 표준상이 멱영인 것이다. $A$의 위상적으로
    멱영인 원소 전체의 집합 $\mathfrak T$는 아이디얼이다.
  \item[{\upshape(ii)}] 더 나아가 $A$가 준허용환이고 $\mathfrak J$가
    $A$의 정의 아이디얼이라고 하자. $x\in A$가 위상적으로 멱영일
    필요충분조건은 그 표준상이 $A/\mathfrak J$에서 멱영인 것이다.
    아이디얼 $\mathfrak T$는 $A/\mathfrak J$의 멱영근기의 역상이고,
    따라서 열린 아이디얼이다.
\end{enumerate}
\end{lemma}

(i)은 정의에서 바로 따른다. (ii)를 보이기 위해서는 $A$에서 $0$의
모든 근방 $V$에 대하여 $\mathfrak J^n\subset V$인 $n>0$가 존재함을
주목하면 충분하다. $x^m\in\mathfrak J$이면 $q\geq n$에 대하여
$x^{mq}\in V$이므로 $x$는 위상적으로 멱영이다.

\begin{proposition}[7.1.4]
\label{0.7.1.4-ko}
$A$를 준허용환, $\mathfrak J$를 $A$의 정의 아이디얼이라 하자.
\begin{enumerate}
  \item[{\upshape(i)}] $A$의 아이디얼 $\mathfrak J'$이 어떤 정의
    아이디얼에 포함될 필요충분조건은 어떤 정수 $n>0$에 대하여
    ${\mathfrak J'}^n\subset\mathfrak J$인 것이다.
  \item[{\upshape(ii)}] $x\in A$가 어떤 정의 아이디얼에 포함될
    필요충분조건은 $x$가 위상적으로 멱영인 것이다.
\end{enumerate}
\end{proposition}

(i) ${\mathfrak J'}^n\subset\mathfrak J$이라고 하자. $A$에서 $0$의
모든 열린 근방 $V$에 대하여 $\mathfrak J^m\subset V$인 $m$이
존재하므로 ${\mathfrak J'}^{mn}\subset V$이다.

(ii) 이 조건이 필요함은 명백하다. 역으로 이 조건이 성립하면
$x^n\in\mathfrak J$인 $n$이 존재한다. 따라서
$\mathfrak J'=\mathfrak J+Ax$는 열린 아이디얼이고
${\mathfrak J'}^n\subset\mathfrak J$이므로 정의 아이디얼이다.

\begin{corollary}[7.1.5]
\label{0.7.1.5-ko}
준허용환 $A$에서 열린 소 아이디얼은 모든 정의 아이디얼을 포함한다.
\end{corollary}

\begin{corollary}[7.1.6]
\label{0.7.1.6-ko}
(7.1.4)의 표기와 가정 아래에서 $A$의 아이디얼 $\mathfrak J_0$에
관한 다음 성질들은 동치이다.
\begin{enumerate}
  \item[$a)$] $\mathfrak J_0$는 $A$의 가장 큰 정의 아이디얼이다.
  \item[$b)$] $\mathfrak J_0$는 극대 정의 아이디얼이다.
  \item[$c)$] $\mathfrak J_0$는 정의 아이디얼이고
    $A/\mathfrak J_0$는 축소환이다.
\end{enumerate}
이 성질들을 갖는 아이디얼 $\mathfrak J_0$가 존재할 필요충분조건은
$A/\mathfrak J$의 멱영근기가 멱영인 것이다. 이때
$\mathfrak J_0$는 $A$의 위상적으로 멱영인 원소들의 아이디얼
$\mathfrak T$와 같다.
\end{corollary}

$a)$가 $b)$를 함의함은 명백하고, (7.1.4, (ii))와 (7.1.3, (ii))에
따라 $b)$는 $c)$를 함의한다. 같은 이유로 $c)$는 $a)$를 함의한다.
마지막 명제는 (7.1.4, (i))와 (7.1.3, (ii))에서 따른다.

$A/\mathfrak J$의 멱영근기 $\mathfrak T/\mathfrak J$가 멱영이면
축소 몫환 $A/\mathfrak T$를 $A_{\mathrm{red}}$로 쓴다.

\oldpage[0\textsubscript{I}]{62}

\begin{corollary}[7.1.7]
\label{0.7.1.7-ko}
뇌터 준허용환에는 가장 큰 정의 아이디얼이 존재한다.
\end{corollary}

\begin{corollary}[7.1.8]
\label{0.7.1.8-ko}
준허용환 $A$에서 어떤 정의 아이디얼 $\mathfrak J$의 거듭제곱
$\mathfrak J^n$ ($n>0$)들이 $0$의 근방 기본계를 이룬다면, 모든
정의 아이디얼 $\mathfrak J'$의 거듭제곱 ${\mathfrak J'}^n$들도
$0$의 근방 기본계를 이룬다.
\end{corollary}

\begin{definition}[7.1.9]
\label{0.7.1.9-ko}
준허용환 $A$에서 어떤 정의 아이디얼 $\mathfrak J$의 거듭제곱
$\mathfrak J^n$들이 $A$에서 $0$의 근방 기본계를 이루면(동치로,
$\mathfrak J^n$들이 열린 아이디얼이면) $A$를 \emph{준아딕환}이라고
한다. 분리이고 완비인 준아딕환을 \emph{아딕환}이라고 한다.
\end{definition}

$\mathfrak J$가 준아딕환(각각 아딕환) $A$의 정의 아이디얼이면
$A$를 \emph{$\mathfrak J$-준아딕환}(각각
\emph{$\mathfrak J$-아딕환})이라고도 하고, 그 위상을
\emph{$\mathfrak J$-준아딕 위상}(각각
\emph{$\mathfrak J$-아딕 위상})이라고 한다. 더 일반적으로 $M$이
$A$-가군이면 부분가군 $\mathfrak J^nM$들을 $0$의 근방 기본계로
갖는 $M$의 위상을 \emph{$\mathfrak J$-준아딕 위상}(각각
\emph{$\mathfrak J$-아딕 위상})이라고 한다. (7.1.8)에 따라 이
위상들은 정의 아이디얼 $\mathfrak J$의 선택과 무관하다.

\begin{proposition}[7.1.10]
\label{0.7.1.10-ko}
$A$를 허용환, $\mathfrak J$를 $A$의 정의 아이디얼이라 하자. 그러면
$\mathfrak J$는 환 $A$의 근기에 포함된다.
\end{proposition}

이 명제는 다음 따름정리들 가운데 어느 것과도 동치이다.

\begin{corollary}[7.1.11]
\label{0.7.1.11-ko}
모든 $x\in\mathfrak J$에 대하여 $1+x$는 $A$에서 가역이다.
\end{corollary}

\begin{corollary}[7.1.12]
\label{0.7.1.12-ko}
$f\in A$가 $A$에서 가역일 필요충분조건은 그 표준상이
$A/\mathfrak J$에서 가역인 것이다.
\end{corollary}

\begin{corollary}[7.1.13]
\label{0.7.1.13-ko}
모든 유한형 $A$-가군 $M$에 대하여 관계 $M=\mathfrak JM$(동치로
$M\otimes_A(A/\mathfrak J)=0$)는 $M=0$을 함의한다.
\end{corollary}

\begin{corollary}[7.1.14]
\label{0.7.1.14-ko}
$u:M\to N$을 $A$-가군의 준동형이라 하고 $N$은 유한형이라고 하자.
$u$가 전사일 필요충분조건은
$u\otimes1:M\otimes_A(A/\mathfrak J)\to
N\otimes_A(A/\mathfrak J)$가 전사인 것이다.
\end{corollary}

실제로 (7.1.10)과 (7.1.11)의 동치는 Bourbaki, \emph{Alg.},
chap.~VIII, \textsection6, n\textsuperscript{o}~3, th.~1에서 따르고,
(7.1.10)과 (7.1.13)의 동치는 \emph{loc. cit.}, th.~2에서 따른다.
(7.1.10)이 (7.1.14)를 함의한다는 것은 \emph{loc. cit.}, prop.~6의
cor.~4에서 따른다. 한편 영 준동형에 (7.1.14)를 적용하면
(7.1.14)가 (7.1.13)을 함의함을 알 수 있다. 마지막으로 (7.1.10)에
따르면 $f$가 $A/\mathfrak J$에서 가역일 때 $f$는 $A$의 어떤 극대
아이디얼에도 속하지 않으므로 $A$에서 가역이다. 즉 (7.1.10)은
(7.1.12)를 함의한다. 역으로 (7.1.12)는 (7.1.11)을 함의한다.

따라서 (7.1.11)을 증명하면 충분하다. $A$는 분리이고 완비이며
수열 $(\mathfrak J^n)$은 $0$으로 가므로 급수
$\sum_{n=0}^{\infty}(-1)^n x^n$은 $A$에서 수렴한다. 그 합을 $y$라
하면 $y(1+x)=1$이다.

\subsection{아딕환과 사영극한}
\label{subsection:0.7.2-ko}

\begin{env}[7.2.1]
\label{0.7.2.1-ko}
이산환들의 모든 사영극한은 자명하게 선형 위상환이며 분리이고
완비이다. 역으로 $A$를 선형 위상환이라 하고, 아이디얼들로 이루어진
$A$의 $0$의 열린 근방 기본계를 $(\mathfrak J_\lambda)$라 하자.
\oldpage[0\textsubscript{I}]{63}
표준 사상
$\varphi_\lambda:A\to A/\mathfrak J_\lambda$들은 연속 준동형의
사영계를 이루므로 연속 준동형
$\varphi:A\to\varprojlim A/\mathfrak J_\lambda$를 정의한다. $A$가
분리이면 $\varphi$는 $A$에서
$\varprojlim A/\mathfrak J_\lambda$의 조밀한 부분환으로 가는
위상동형이다. 더 나아가 $A$가 완비이면 $\varphi$는 $A$에서
$\varprojlim A/\mathfrak J_\lambda$로 가는 위상동형이다.
\end{env}

\begin{lemma}[7.2.2]
\label{0.7.2.2-ko}
선형 위상환이 허용환일 필요충분조건은 이산환들의 사영계
$(A_\lambda,u_{\lambda\mu})$의 사영극한
$A=\varprojlim A_\lambda$와 동형인 것이다. 여기서 지표집합은
$\leq$에 관한 여과 순서집합 $L$이고, $L$은 $0$으로 쓰는 최소원을
가지며 다음 조건들을 만족한다.
1\textsuperscript{o} 사상 $u_\lambda:A\to A_\lambda$는 전사이다.
2\textsuperscript{o} 준동형
$u_{0\lambda}:A_\lambda\to A_0$의 핵 $\mathfrak J_\lambda$는
멱영이다. 이때 준동형 $u_0:A\to A_0$의 핵 $\mathfrak J$는
$\varprojlim\mathfrak J_\lambda$와 같다.
\end{lemma}

이 조건의 필요성은 (7.2.1)에서 어떤 정의 아이디얼
$\mathfrak J_0$에 포함되는 정의 아이디얼들로 이루어진 $0$의 근방
기본계 $(\mathfrak J_\lambda)$를 취하고 (7.1.4, (i))를 적용하면
따른다. 충분성은 사영극한의 정의와 (7.1.2)에서 따르고, 마지막
명제는 자명하다.

\begin{env}[7.2.3]
\label{0.7.2.3-ko}
$A$를 허용 위상환이라 하고, $\mathfrak J$를 $A$의 어떤 정의
아이디얼에 포함되는 아이디얼, 다시 말해 (7.1.4)에 따라 수열
$(\mathfrak J^n)$이 $0$으로 가는 아이디얼이라 하자. $0$의 근방
기본계가 거듭제곱 $\mathfrak J^n$ ($n>0$)들인 환 위상을 $A$에
줄 수 있으며, 이를 다시 \emph{$\mathfrak J$-준아딕 위상}이라고
하겠다. $A$가 허용환이라는 가정에 따라
$\bigcap_n\mathfrak J^n=0$이므로 $A$의 $\mathfrak J$-준아딕
위상은 분리이다. 이 위상에 관한 $A$의 완비화를
$\widehat A=\varprojlim A/\mathfrak J^n$이라 하자. 여기서
$A/\mathfrak J^n$에는 이산 위상을 준다. 또한 준동형들의 열
$u_n:A\to A/\mathfrak J^n$의 사영극한인 환 준동형
$u:A\to\widehat A$를 생각하자. 이 준동형은 반드시 연속일 필요는
없다. 한편 $A$의 $\mathfrak J$-준아딕 위상은 주어진 위상
$\mathcal T$보다 세밀하다. $A$는 $\mathcal T$에 관하여 분리이고
완비이므로, $\mathfrak J$-준아딕 위상을 준 $A$에서
$\mathcal T$를 준 $A$로 가는 항등사상을 연속적으로 확장할 수 있다.
따라서 연속 준동형 $v:\widehat A\to A$를 얻는다.
\end{env}

\begin{proposition}[7.2.4]
\label{0.7.2.4-ko}
$A$가 허용환이고 $\mathfrak J$가 $A$의 어떤 정의 아이디얼에
포함되면, $A$는 $\mathfrak J$-준아딕 위상에 관하여 분리이고
완비이다.
\end{proposition}

실제로 (7.2.3)의 표기를 사용하면 $v\circ u$가 $A$의 항등사상임은
바로 알 수 있다. 한편
$u_n\circ v:\widehat A\to A/\mathfrak J^n$은
$A$의 $\mathfrak J$-준아딕 위상과 $A/\mathfrak J^n$의 이산
위상에 관하여 표준 사상 $u_n$을 연속적으로 확장한 것이다. 다시
말해 이는
$\widehat A=\varprojlim_k A/\mathfrak J^k$에서
$A/\mathfrak J^n$으로 가는 표준 사상이다. 따라서 $u\circ v$는 이
사상열의 사영극한, 즉 정의에 따라 $\widehat A$의 항등사상이다.
이로써 명제가 증명된다.

\begin{corollary}[7.2.5]
\label{0.7.2.5-ko}
(7.2.3)의 가정 아래에서 다음 조건들은 동치이다.
\begin{enumerate}
  \item[$a)$] 준동형 $u$는 연속이다.

\oldpage[0\textsubscript{I}]{64}

  \item[$b)$] 준동형 $v$는 쌍연속이다.
  \item[$c)$] $A$는 $\mathfrak J$-아딕환이다.
\end{enumerate}
\end{corollary}

\begin{corollary}[7.2.6]
\label{0.7.2.6-ko}
$A$를 허용환, $\mathfrak J$를 $A$의 정의 아이디얼이라 하자.
$A$가 뇌터 환일 필요충분조건은 $A/\mathfrak J$가 뇌터 환이고
$\mathfrak J/\mathfrak J^2$가 유한형
$(A/\mathfrak J)$-가군인 것이다.
\end{corollary}

이 조건들이 필요함은 자명하다. 역으로 이 조건들이 성립한다고 하자.
(7.2.4)에 따라 $A$는 $\mathfrak J$-준아딕 위상에 관하여
완비이므로, $A$가 뇌터 환일 필요충분조건은 $\mathfrak J^n$의
여과에 관한 연관 등급환 $\operatorname{grad}(A)$가 뇌터 환인
것이다([I], p.~18-07, th.~4). $\mathfrak J$의 원소
$a_1,\ldots,a_n$의 $\mathfrak J^2$에 대한 유들이
$\mathfrak J/\mathfrak J^2$를 $(A/\mathfrak J)$-가군으로서
생성한다고 하자. 귀납법으로, $a_i$들($1\leq i\leq n$)에 관한
전체 차수 $m$인 단항식들의 $\mathfrak J^{m+1}$에 대한 유들이
$(A/\mathfrak J)$-가군
$\mathfrak J^m/\mathfrak J^{m+1}$의 생성계를 이룸을 바로 알 수
있다. 따라서 $\operatorname{grad}(A)$는 다항식환
$(A/\mathfrak J)[T_1,\ldots,T_n]$의 몫환과 동형이다. 여기서
$T_i$들은 부정원이다. 이로써 증명이 끝난다.

\begin{proposition}[7.2.7]
\label{0.7.2.7-ko}
$(A_i,u_{ij})$를 이산환들의 사영계($i\in\mathbf N$)라 하고, 모든
정수 $i$에 대하여 $\mathfrak J_i$를 준동형
$u_{0i}:A_i\to A_0$의 $A_i$에서의 핵이라 하자. 다음을 가정한다.
\begin{enumerate}
  \item[$a)$] $i\leq j$이면 $u_{ij}$는 전사이고 그 핵은
    $\mathfrak J_j^{i+1}$이다. 따라서 $A_i$는
    $A_j/\mathfrak J_j^{i+1}$과 동형이다.
  \item[$b)$] $\mathfrak J_1/\mathfrak J_1^2$
    $(=\mathfrak J_1)$은 $A_0=A_1/\mathfrak J_1$ 위의 유한형
    가군이다.
\end{enumerate}
$A=\varprojlim_i A_i$라 하고, 모든 정수 $n\geq0$에 대하여
$u_n:A\to A_n$을 표준 준동형, $\mathfrak J^{(n+1)}\subset A$를
그 핵이라 하자. 그러면 다음이 성립한다.
\begin{enumerate}
  \item[(i)] $A$는 $\mathfrak J=\mathfrak J^{(1)}$을 정의
    아이디얼로 갖는 아딕환이다.
  \item[(ii)] 모든 $n\geq1$에 대하여
    $\mathfrak J^{(n)}=\mathfrak J^n$이다.
  \item[(iii)] $\mathfrak J/\mathfrak J^2$는
    $\mathfrak J_1=\mathfrak J_1/\mathfrak J_1^2$과 동형이고,
    따라서 $A_0=A/\mathfrak J$ 위의 유한형 가군이다.
\end{enumerate}
\end{proposition}

$u_{ij}$들이 전사라는 가정에 따라 $u_n$은 전사이다. 또한 가정
$a)$에 따라 $\mathfrak J_j^{j+1}=0$이므로 $A$는 허용환이다
(7.2.2). 정의상 $\mathfrak J^{(n)}$들은 $A$의 $0$의 근방
기본계를 이루므로 (ii)는 (i)을 함의한다. 더 나아가
$\mathfrak J=\varprojlim_i\mathfrak J_i$이고 사상
$\mathfrak J\to\mathfrak J_i$들은 전사이므로 (ii)는 (iii)을
함의한다. 따라서 (ii)를 증명하면 충분하다.

정의에 따라 $\mathfrak J^{(n)}$은 $k<n$일 때 $x_k=0$인 $A$의
원소 $(x_k)_{k\geq0}$들로 이루어진다. 그러므로
$\mathfrak J^{(n)}\mathfrak J^{(m)}\subset
\mathfrak J^{(n+m)}$이고, 다시 말해 $\mathfrak J^{(n)}$들은
$A$의 \emph{여과}를 이룬다. 한편
$\mathfrak J^{(n)}/\mathfrak J^{(n+1)}$은
$\mathfrak J^{(n)}$의 $A_n$ 위로의 사영과 동형이다.
$\mathfrak J^{(n)}=\varprojlim_{i>n}\mathfrak J_i^n$이므로 이
사영은 $\mathfrak J_n^n$이고, 이는
$A_0=A_n/\mathfrak J_n$ 위의 가군이다.

$a_j=(a_{jk})_{k\geq0}$ ($1\leq j\leq r$)를
$\mathfrak J=\mathfrak J^{(1)}$의 원소들이라 하고,
$a_{11},\ldots,a_{r1}$이 $A_0$ 위에서 $\mathfrak J_1$의 생성계를
이룬다고 하자. $a_j$들에 관한 전체 차수 $n$인 단항식들의 집합
$S_n$이 $A$의 아이디얼 $\mathfrak J^{(n)}$을 생성함을 보이겠다.
$\mathfrak J_i^{i+1}=0$이므로 우선
$S_n\subset\mathfrak J^{(n)}$임은 명백하다. $A$는 여과
$(\mathfrak J^{(m)})$에 관하여 완비이므로, $S_n$의 원소들의
$\mathfrak J^{(n+1)}$에 대한 유들의 집합 $\overline S_n$이 앞의
여과에 관한 등급환 $\operatorname{grad}(A)$ 위에서 등급가군
$\operatorname{grad}(\mathfrak J^{(n)})$을 생성함을 보이면
충분하다([I], p.~18-06, lemme). $\operatorname{grad}(A)$의 곱셈
정의에 따라

\oldpage[0\textsubscript{I}]{65}

모든 $m$에 대하여 $\overline S_m$이 $A_0$-가군
$\mathfrak J^{(m)}/\mathfrak J^{(m+1)}$의 생성계임을 보이면
충분하다. 동치로 $\mathfrak J_m^m$이 $a_{jm}$들
($1\leq j\leq r$)에 관한 차수 $m$인 단항식들로 생성됨을 보이면
된다. 이를 위해서는 $\mathfrak J_m$이 $A_m$-가군으로서
$a_{jm}$들에 관한 차수 $\leq m$인 단항식들로 생성됨을 보이면
된다. $m=1$에서는 정의에 따라 자명하므로 $m$에 관하여 귀납하자.
$\mathfrak J'_m$을 이 단항식들이 생성하는
$\mathfrak J_m$의 $A_m$-부분가군이라 하자. 관계
$\mathfrak J_{m-1}=\mathfrak J_m/\mathfrak J_m^m$과 귀납가정에
따라 $\mathfrak J_m=\mathfrak J'_m+\mathfrak J_m^m$이다.
$\mathfrak J_m^{m+1}=0$이므로 여기서
$\mathfrak J_m^m={\mathfrak J'_m}^{\,m}$을 얻고, 마침내
$\mathfrak J_m=\mathfrak J'_m$을 얻는다.

\begin{corollary}[7.2.8]
\label{0.7.2.8-ko}
(7.2.7)의 조건 아래에서 $A$가 뇌터 환일 필요충분조건은 $A_0$가
뇌터 환인 것이다.
\end{corollary}

이는 (7.2.6)에서 바로 따른다.

\begin{proposition}[7.2.9]
\label{0.7.2.9-ko}
(7.2.7)의 가정들이 성립한다고 하자. 모든 정수 $i$에 대하여
$M_i$를 $A_i$-가군이라 하고, $i\leq j$일 때
$v_{ij}:M_j\to M_i$를 $u_{ij}$-준동형이라 하여
$(M_i,v_{ij})$가 사영계를 이룬다고 하자. 더 나아가 $M_0$가 유한형
$A_0$-가군이고 $v_{ij}$가 전사이며 그 핵이
$\mathfrak J_j^{i+1}M_j$라고 하자. 그러면
$M=\varprojlim M_i$는 유한형 $A$-가군이고, 전사
$u_n$-준동형 $v_n:M\to M_n$의 핵은
$\mathfrak J^{n+1}M$이다. 따라서 $M_n$은
$M/\mathfrak J^{n+1}M=M\otimes_A(A/\mathfrak J^{n+1})$과
동일시된다.
\end{proposition}

$z_h=(z_{hk})_{k\geq0}$ ($1\leq h\leq s$)를 $M$의 원소들이라
하고, $z_{h0}$들이 $M_0$의 생성계를 이룬다고 하자. $z_h$들이
$A$-가군 $M$을 생성함을 보이겠다. $M^{(n)}$을 $k<n$일 때
$y_k=0$인 $M$의 원소 $y=(y_k)_{k\geq0}$ 전체의 집합이라 하자.
$A$-가군 $M$은 $M^{(n)}$들이 이루는 여과에 관하여 분리이고
완비이다. 또한
$\mathfrak J^{(n)}M\subset M^{(n)}$이고
$M^{(n)}/M^{(n+1)}=\mathfrak J_n^nM_n$임은 명백하다.

따라서 $z_h$들의 $M^{(1)}$에 대한 유들이 앞의 여과에 관한
등급환 $\operatorname{grad}(A)$ 위에서 등급가군
$\operatorname{grad}(M)$을 생성함을 보이면 충분하다
([I], p.~18-06, lemme). 이를 위해서는 $z_{hn}$들
($1\leq h\leq s$)이 $A_n$-가군 $M_n$을 생성함을 보이면 충분함을
쉽게 알 수 있다. $n=0$에서는 정의에 따라 자명하므로 다시 $n$에
관하여 귀납하자. 관계
$M_{n-1}=M_n/\mathfrak J_n^nM_n$과 귀납가정에 따라, $M'_n$을
$z_{hn}$들이 생성하는 $M_n$의 부분가군이라 하면
$M_n=M'_n+\mathfrak J_n^nM_n$이다. $\mathfrak J_n$은 멱영이므로
여기서 $M_n=M'_n$을 얻는다. 연관 등급가군으로 넘어가는 같은
논법에 따라 $\mathfrak J^{(n)}M$에서 $M^{(n)}$으로 가는 표준
사상은 전사이고, 따라서 전단사이다. 다시 말해
$\mathfrak J^{(n)}M=\mathfrak J^nM$은 $M\to M_n$의 핵이다.

\begin{corollary}[7.2.10]
\label{0.7.2.10-ko}
$(N_i,w_{ij})$를 (7.2.9)의 조건을 만족하는 두 번째
$A_i$-가군의 사영계라 하고 $N=\varprojlim N_i$라 하자.
$A_i$-준동형 $h_i:M_i\to N_i$의 사영계 $(h_i)$와
$A$-가군 준동형 $h:M\to N$ 사이에는 전단사 대응이 있다. 이
$h$들은 반드시 $\mathfrak J$-아딕 위상에 관하여 연속이다.
\end{corollary}

$h:M\to N$이 $A$-준동형이면
$h(\mathfrak J^nM)\subset\mathfrak J^nN$이므로 $h$는 연속이다.
몫으로 내려가면 $h$에는 $A_i$-준동형 $h_i:M_i\to N_i$의 사영계가
대응하고 $h$는 그 사영극한이다. 따라서 따름정리가 성립한다.

\begin{remark}[7.2.11]
\label{0.7.2.11-ko}
$A$를 아딕환, $\mathfrak J$를
$\mathfrak J/\mathfrak J^2$가 유한형
$(A/\mathfrak J)$-가군인 정의 아이디얼이라 하자.
$A_i=A/\mathfrak J^{i+1}$들이 (7.2.7)의 조건을 만족함은 명백하다.

\oldpage[0\textsubscript{I}]{66}

$A$는 $A_i$들의 사영극한이므로 명제 (7.2.7)은 이러한 종류의
\emph{모든} 아딕환, 특히 모든 \emph{뇌터 아딕환}을 기술한다.
\end{remark}

\begin{example}[7.2.12]
\label{0.7.2.12-ko}
$B$를 환, $\mathfrak J$를 $B$의 아이디얼이라 하고,
$\mathfrak J/\mathfrak J^2$가 $B/\mathfrak J$ 위의 유한형
가군이라고 하자. 이는 $B$ 위의 유한형 가군이라는 것과 동치이다.
$A=\varprojlim_n B/\mathfrak J^{n+1}$로 놓자. $A$는
$\mathfrak J$-준아딕 위상을 준 $B$의 분리 완비화이다.
$A_n=B/\mathfrak J^{n+1}$로 놓으면 $A_n$들이 (7.2.7)의 조건을
만족함은 바로 알 수 있다. 따라서 $A$는 아딕환이다.
$\overline{\mathfrak J}$를 $\mathfrak J$의 표준상의 $A$에서의
폐포라 하면 $\overline{\mathfrak J}$는 $A$의 정의 아이디얼이고,
$\overline{\mathfrak J^n}$은 $\mathfrak J^n$의 표준상의 폐포이며,
$A/\overline{\mathfrak J^n}$은 $B/\mathfrak J^n$과 동일시된다.
또한 $\overline{\mathfrak J}/\overline{\mathfrak J^2}$은
$(A/\overline{\mathfrak J})$-가군으로서
$\mathfrak J/\mathfrak J^2$과 동형이다.

마찬가지로 $N/\mathfrak JN$이 유한형 $B$-가군이고
$M_i=N/\mathfrak J^{i+1}N$, $M=\varprojlim M_i$로 놓으면 $M$은
유한형 $A$-가군이고 $\mathfrak J$-준아딕 위상에 관한 $N$의 분리
완비화와 동형이다. $\overline{\mathfrak J^n}M$은
$\mathfrak J^nN$의 표준상의 폐포와 동일시되고,
$M/\overline{\mathfrak J^n}M$은 $N/\mathfrak J^nN$과
동일시된다.
\end{example}

\subsection{뇌터 준아딕환}
\label{subsection:0.7.3-ko}

\begin{env}[7.3.1]
\label{0.7.3.1-ko}
$A$를 환, $\mathfrak J$를 $A$의 아이디얼, $M$을 $A$-가군이라
하자. $\mathfrak J$-준아딕 위상에 관한 $A$(각각 $M$)의 분리
완비화를
$\widehat A=\varprojlim_n A/\mathfrak J^n$(각각
$\widehat M=\varprojlim_n M/\mathfrak J^nM$)으로 쓰겠다.
$A$-가군의 완전열
\[
  M'\xrightarrow{u}M\xrightarrow{v}M''\to0
\]
이 주어졌다고 하자. 등식
$M/\mathfrak J^nM=M\otimes_A(A/\mathfrak J^n)$에 따라 모든
$n$에 대하여 열
\[
  M'/\mathfrak J^nM'\xrightarrow{u_n}
  M/\mathfrak J^nM\xrightarrow{v_n}
  M''/\mathfrak J^nM''\to0
\]
은 완전하다. 또한
$v(\mathfrak J^nM)=\mathfrak J^nv(M)=\mathfrak J^nM''$이므로
$\widehat v=\varprojlim v_n$은 전사이다(Bourbaki,
\emph{Top. gén.}, chap.~IX, 2\textsuperscript{e} éd., p.~60,
cor.~2). 한편 $z=(z_k)$가 $\widehat v$의 핵의 원소이면 모든 정수
$k$에 대하여 $u_k(z'_k)=z_k$인
$z'_k\in M'/\mathfrak J^kM'$가 존재한다. 따라서
$z'=(z'_n)\in\widehat{M'}$가 존재하여 $\widehat u(z')$의 처음
$k$개 성분이 $z$의 처음 $k$개 성분과 일치한다. 다시 말해
$\widehat u$에 의한 $\widehat{M'}$의 상은 $\widehat v$의 핵에
\emph{조밀}하다.

$A$가 \emph{뇌터 환}이라고 가정하면
$\mathfrak J/\mathfrak J^2$가 유한형 $A$-가군이므로 (7.2.12)에
따라 $\widehat A$도 뇌터 환이다. 더 나아가 다음이 성립한다.
\end{env}

\begin{namedtheorem}[7.3.2]{크룰 정리}
\label{0.7.3.2-ko}
$A$를 뇌터 환, $\mathfrak J$를 $A$의 아이디얼, $M$을 유한형
$A$-가군, $M'$을 $M$의 부분가군이라 하자. $M$의
$\mathfrak J$-준아딕 위상이 $M'$에 유도하는 위상은 $M'$의
$\mathfrak J$-준아딕 위상과 같다.
\end{namedtheorem}

이는 다음 보조정리에서 바로 따른다.

\begin{namedlemma}[7.3.2.1]{아르틴--리스 보조정리}
\label{0.7.3.2.1-ko}
(7.3.2)의 가정 아래에서 어떤 정수 $p$가 존재하여 $n\geq p$이면
\[
  M'\cap\mathfrak J^nM=\mathfrak J^{n-p}
  (M'\cap\mathfrak J^pM)
\]
이다.
\end{namedlemma}

증명은 ([I], p.~2-04)를 보라.

\oldpage[0\textsubscript{I}]{67}

\begin{corollary}[7.3.3]
\label{0.7.3.3-ko}
(7.3.2)의 가정 아래에서 표준 사상
$M\otimes_A\widehat A\to\widehat M$은 전단사이고, 유한형
$A$-가군의 범주에서 $M$에 관한 함자
$M\mapsto M\otimes_A\widehat A$는 완전하다. 따라서
$\mathfrak J$-아딕 분리 완비화 $\widehat A$는 평탄 $A$-가군이다
(6.1.1).
\end{corollary}

먼저 유한형 $A$-가군의 범주에서 $\widehat M$이 $M$에 관한
\emph{완전} 함자임을 주목하자. 실제로 완전열
\[
  0\to M'\xrightarrow{u}M\xrightarrow{v}M''\to0
\]
이 주어졌다고 하자. $\widehat v:\widehat M\to\widehat{M''}$가
전사임은 이미 알고 있다(7.3.1). 한편 $i$를 표준 준동형
$M\to\widehat M$이라 하면 크룰 정리에 따라 $\widehat M$에서
$i(u(M'))$의 폐포는 $M'$의 $\mathfrak J$-준아딕 분리 완비화와
동일시된다. 따라서 $\widehat u$는 단사이고, (7.3.1)에 따라
$\widehat u$의 상은 $\widehat v$의 핵과 같다.

이제 표준 사상 $M\otimes_A\widehat A\to\widehat M$은 사상들
\[
  M\otimes_A\widehat A\to
  M\otimes_A(A/\mathfrak J^n)=M/\mathfrak J^nM
\]
의 사영극한을 취하여 얻는다. $M=A^p$ 꼴일 때 이 사상이
전단사임은 명백하다. $M$이 유한형 $A$-가군이면 완전열
$A^p\to A^q\to M\to0$이 존재한다. 유한형 $A$-가군의 범주에서
$M$에 관한 두 함자 $M\otimes_A\widehat A$와 $\widehat M$은
\emph{오른쪽 완전}이므로 다음 가환 그림을 얻는다.
\[
\begin{tikzcd}[column sep=large,row sep=large]
  A^p\otimes_A\widehat A \arrow[r] \arrow[d] &
  A^q\otimes_A\widehat A \arrow[r] \arrow[d] &
  M\otimes_A\widehat A \arrow[r] \arrow[d] &
  0 \\
  \widehat A^{\,p} \arrow[r] &
  \widehat A^{\,q} \arrow[r] &
  \widehat M \arrow[r] &
  0
\end{tikzcd}
\]
두 행은 완전하고 처음 두 수직 화살표는 동형이므로 결론이 바로
따른다.

\begin{corollary}[7.3.4]
\label{0.7.3.4-ko}
$A$를 뇌터 환, $\mathfrak J$를 $A$의 아이디얼, $M,N$을 두 유한형
$A$-가군이라 하자. 다음과 같은 함자적 표준 동형들이 있다.
\[
  (M\otimes_A N)^\wedge
  \xrightarrow{\sim}
  \widehat M\otimes_{\widehat A}\widehat N,
  \qquad
  \bigl(\operatorname{Hom}_A(M,N)\bigr)^\wedge
  \xrightarrow{\sim}
  \operatorname{Hom}_{\widehat A}(\widehat M,\widehat N).
\]
\end{corollary}

이는 (7.3.3), (6.2.1), (6.2.2)에서 따른다.

\begin{corollary}[7.3.5]
\label{0.7.3.5-ko}
$A$를 뇌터 환, $\mathfrak J$를 $A$의 아이디얼이라 하자. 다음
조건들은 동치이다.
\begin{enumerate}
  \item[$a)$] $\mathfrak J$는 $A$의 근기에 포함된다.
  \item[$b)$] $\widehat A$는 충실 평탄 $A$-가군이다(6.4.1).
  \item[$c)$] 모든 유한형 $A$-가군은 $\mathfrak J$-준아딕 위상에
    관하여 분리이다.
  \item[$d)$] 유한형 $A$-가군의 모든 부분가군은
    $\mathfrak J$-준아딕 위상에 관하여 닫혀 있다.
\end{enumerate}
\end{corollary}

$\widehat A$는 평탄 $A$-가군이므로 조건 $b)$와 $c)$는 동치이다.
실제로 $b)$는 모든 유한형 $A$-가군 $M$에 대하여 표준 사상
$M\to\widehat M=M\otimes_A\widehat A$가 단사라는 것과 동치이다
(6.6.1, $c)$). $c)$가 $d)$를 함의함은 바로 알 수 있다. 유한형
$A$-가군 $M$의 부분가군이 $N$이면 $M/N$은
$\mathfrak J$-준아딕 위상에 관하여 분리이므로 $N$은 $M$에서
닫혀 있기 때문이다.

$d)$가 $a)$를 함의함을 보이자. $\mathfrak m$이 $A$의 극대
아이디얼이면 $\mathfrak m$은 $A$의 $\mathfrak J$-준아딕 위상에
관하여 닫혀 있으므로
\[
  \mathfrak m=\bigcap_{p\geq0}(\mathfrak m+\mathfrak J^p)
\]
이다. $\mathfrak m+\mathfrak J^p$는 반드시 $A$ 또는
$\mathfrak m$과 같으므로 충분히 큰 $p$에 대하여
$\mathfrak m+\mathfrak J^p=\mathfrak m$이다.

\oldpage[0\textsubscript{I}]{68}

따라서 $\mathfrak J^p\subset\mathfrak m$이고, $\mathfrak m$은
소 아이디얼이므로 $\mathfrak J\subset\mathfrak m$이다.

마지막으로 $a)$가 $b)$를 함의함을 보이자. 유한형 $A$-가군
$M$에서 $\mathfrak J$-준아딕 위상에 관한 $\{0\}$의 폐포를 $P$라
하자. 크룰 정리 (7.3.2)에 따라 $M$의
$\mathfrak J$-준아딕 위상이 $P$에 유도하는 위상은 $P$의
$\mathfrak J$-준아딕 위상이므로 $\mathfrak JP=P$이다. $P$는
유한형이고 $\mathfrak J$는 $A$의 근기에 포함되므로 나카야마
보조정리에 따라 $P=0$이다.

$A$가 \emph{뇌터 국소환}이고 $\mathfrak J\neq A$가 $A$의 임의의
아이디얼이면 (7.3.5)의 조건들이 만족된다는 점에 주의하자.

\begin{corollary}[7.3.6]
\label{0.7.3.6-ko}
$A$가 뇌터 $\mathfrak J$-아딕환이면 모든 유한형 $A$-가군은
$\mathfrak J$-준아딕 위상에 관하여 분리이고 완비이다.
\end{corollary}

이때 $\widehat A=A$이므로 이는 (7.3.3)에서 바로 따른다.

따라서 명제 (7.2.9)는 뇌터 아딕환 위의 \emph{모든} 유한형
가군을 기술한다.

\begin{corollary}[7.3.7]
\label{0.7.3.7-ko}
(7.3.2)의 가정 아래에서 표준 사상
$M\to\widehat M=M\otimes_A\widehat A$의 핵은
$1+\mathfrak J$의 어떤 원소에 의하여 소거되는 $x\in M$ 전체의
집합이다.
\end{corollary}

실제로 $x\in M$이 이 핵에 속할 필요충분조건은 (7.3.2)에 따라
부분가군 $Ax$의 분리 완비화가 $0$인 것이다. 이는
$x\in\mathfrak J(Ax)$와 동치이다.

\subsection{국소환 위의 준유한 가군}
\label{subsection:0.7.4-ko}

\begin{definition}[7.4.1]
\label{0.7.4.1-ko}
국소환 $A$의 극대 아이디얼을 $\mathfrak m$이라 하자. $A$-가군
$M$에 대하여 $M/\mathfrak mM$이 잉여체 $k=A/\mathfrak m$ 위의
유한 차원 벡터공간이면 $M$을 ($A$ 위의) \emph{준유한} 가군이라고
한다.
\end{definition}

$A$가 뇌터 환이면 $\mathfrak m$-준아딕 위상에 관한 $M$의 분리
완비화 $\widehat M$은 유한형 $\widehat A$-가군이다. 실제로
$\mathfrak m/\mathfrak m^2$가 유한형 $A$-가군이므로 이는
(7.2.12)와 $M/\mathfrak mM$에 관한 가정에서 따른다.

특히 $A$가 $\mathfrak m$-준아딕 위상에 관하여 \emph{완비}이고
$M$이 \emph{분리}라고 더 가정하면, 다시 말해
$\bigcap_n\mathfrak m^nM=0$이면 $M$ 자체가 유한형 $A$-가군이다.
실제로 $\widehat M$은 유한형 $A$-가군이고 $M$은 $\widehat M$의
부분가군과 동일시되므로 $M$도 유한형이다(더욱이 (7.3.6)에 따라
$M$은 자신의 완비화와 같다).

\begin{proposition}[7.4.2]
\label{0.7.4.2-ko}
$A,B$를 두 국소환, $\mathfrak m,\mathfrak n$을 각각의 극대
아이디얼이라 하고 $B$가 뇌터 환이라고 가정하자. 국소 준동형
$\varphi:A\to B$와 유한형 $B$-가군 $M$이 주어졌다고 하자. $M$이
준유한 $A$-가군이면 $M$ 위의 $\mathfrak m$-준아딕 위상과
$\mathfrak n$-준아딕 위상은 같고, 따라서 분리이다.
\end{proposition}

가정에 따라 $M/\mathfrak mM$은 $A$-가군으로서 \emph{유한 길이}이고,
따라서 더더욱 $B$-가군으로서도 유한 길이이다. 그러므로
$M/\mathfrak mM$의 소멸자를 포함하는 $B$의 소 아이디얼은
$\mathfrak n$ 하나뿐이다. 실제로 (I.7.4)와 (I.7.2)에 따라 곧바로
$M/\mathfrak mM$이 \emph{단순}인 경우로 환원되며, 이 경우에는
반드시 $B/\mathfrak n$과 동형이므로 주장은 명백하다.

한편 $M$은 유한형 $B$-가군이므로 $M/\mathfrak mM$의 소멸자를
포함하는 소 아이디얼들은 $\mathfrak mB+\mathfrak b$를 포함하는
것들이다. 여기서 $\mathfrak b$는 $B$-가군 $M$의 소멸자이다
(I.7.5). $B$가 뇌터 환이므로 ([III], p.~127, cor.~4)에 따라
$\mathfrak mB+\mathfrak b$는 $B$의 정의 아이디얼이다.

\oldpage[0\textsubscript{I}]{69}

다시 말해 어떤 $k>0$에 대하여
$\mathfrak n^k\subset\mathfrak mB+\mathfrak b\subset\mathfrak n$이다.
따라서 모든 $h>0$에 대하여
\[
  \mathfrak n^{hk}M
  \subset(\mathfrak mB+\mathfrak b)^hM
  =\mathfrak m^hM
  \subset\mathfrak n^hM
\]
이고, 이는 $M$ 위의 $\mathfrak m$-준아딕 위상과
$\mathfrak n$-준아딕 위상이 같음을 보인다. 후자는 (7.3.5)에 따라
분리이기도 하다.

\begin{corollary}[7.4.3]
\label{0.7.4.3-ko}
(7.4.2)의 가정 아래에서 $A$가 뇌터 환이고 $\mathfrak m$-준아딕
위상에 관하여 완비이면 $M$은 유한형 $A$-가군이다.
\end{corollary}

실제로 이때 $M$은 $\mathfrak m$-준아딕 위상에 관하여 분리이고,
주장은 (7.4.1)의 주의에서 따른다.

\begin{env}[7.4.4]
\label{0.7.4.4-ko}
(7.4.2)의 가장 중요한 적용은 $B$ 자체가 준유한 $A$-가군인
경우이다. 이는 $B/\mathfrak mB$가 $k=A/\mathfrak m$ 위의
\emph{유한 차원 대수}라는 것과 동치이다. 앞에서 본 바에 따라 이
조건은 다음 두 조건의 논리곱으로 분해된다.
\begin{enumerate}
  \item[(i)] $\mathfrak mB$는 $B$의 \emph{정의 아이디얼이다}.
  \item[(ii)] $B/\mathfrak n$은 체 $A/\mathfrak m$의
    \emph{유한 차원 확대체이다}.
\end{enumerate}
이 조건들이 성립하면 모든 유한형 $B$-가군은 명백히 준유한
$A$-가군이다.
\end{env}

\begin{corollary}[7.4.5]
\label{0.7.4.5-ko}
(7.4.2)의 가정 아래에서 $\mathfrak b$가 $B$-가군 $M$의
소멸자이면 $B/\mathfrak b$는 준유한 $A$-가군이다.
\end{corollary}

$M\neq0$이라고 가정하자(그렇지 않으면 따름정리는 명백하다).
$M$을 뇌터 국소환 $B/\mathfrak b$ 위의 가군으로 볼 수 있다. 이때
그 소멸자는 $0$이므로 (7.4.2)의 증명에 따라
$\mathfrak m(B/\mathfrak b)$는 $B/\mathfrak b$의 정의 아이디얼이다.

한편 $M/\mathfrak nM$은 $M/\mathfrak mM$의 몫이므로
$A/\mathfrak m$ 위의 유한 차원 벡터공간이다. $M\neq0$이므로
나카야마 보조정리에 따라 $M\neq\mathfrak nM$이다. 따라서
$M/\mathfrak nM$은 $B/\mathfrak n$ 위의 영이 아닌 벡터공간이고,
그것이 $A/\mathfrak m$ 위에서 유한 차원이라는 사실로부터
$B/\mathfrak n$도 $A/\mathfrak m$ 위에서 유한 차원임이 따른다.
그러므로 $B/\mathfrak b$에 (7.4.4)를 적용하면 결론을 얻는다.

\subsection{제한 형식적 멱급수환}
\label{subsection:0.7.5-ko}

\begin{env}[7.5.1]
\label{0.7.5.1-ko}
$A$를 선형 위상을 갖는 분리 완비 위상환이라 하고,
$(\mathfrak J_\lambda)$를 (열린) 아이디얼들로 이루어진 $A$의 $0$의
근방 기본계라 하자. 그러면 $A$는
$\varprojlim A/\mathfrak J_\lambda$와 표준적으로 동일시된다
(7.2.1). 각 $\lambda$에 대하여
\[
  B_\lambda=(A/\mathfrak J_\lambda)[T_1,\ldots,T_r]
\]
라 하자. 여기서 $T_i$들은 부정원이다. $B_\lambda$들이 이산환의
사영계를 이룸은 명백하다. 이제
\[
  \varprojlim B_\lambda=A\{T_1,\ldots,T_r\}
\]
로 놓겠다. 이 위상환이 선택한 아이디얼 기본계
$(\mathfrak J_\lambda)$와 무관함을 보이자.

정확히 말해 $A'$를 형식적 멱급수환
$A[\![T_1,\ldots,T_r]\!]$의 부분환 가운데, 형식적 멱급수
$\sum_\alpha c_\alpha T^\alpha$로서
$\alpha=(\alpha_1,\ldots,\alpha_r)\in\mathbf N^r$이고
\[
  \lim_\alpha c_\alpha=0
\]
인 것들로 이루어진 부분환이라 하자. 여기서 극한은 $\mathbf N^r$의
유한 부분집합들의 여집합이 이루는 필터에 관한 것이다. 이러한
급수들을 $A$에 계수를 갖는 $T_i$들에 관한 \emph{제한 형식적
멱급수}라고 하겠다.

\oldpage[0\textsubscript{I}]{70}

$A$의 $0$의 각 근방 $V$에 대하여 $V'$을, 모든 $\alpha$에 대하여
$c_\alpha\in V$인
$x=\sum_\alpha c_\alpha T^\alpha\in A'$ 전체의 집합이라 하자.
$V'$들이 $A'$ 위에 \emph{분리} 환위상을 정의하는 $0$의 근방
기본계를 이룸은 곧 확인된다. 이제 환 $A\{T_1,\ldots,T_r\}$에서
$A'$로 가는 표준적인 \emph{위상동형}을 정의하겠다.

각 $\alpha\in\mathbf N^r$와 각 $\lambda$에 대하여
$\varphi_{\lambda,\alpha}$를
$(A/\mathfrak J_\lambda)[T_1,\ldots,T_r]$에서
$A/\mathfrak J_\lambda$로 가는 사상 가운데, 첫째 환의 각
다항식에 그 다항식에서 $T^\alpha$의 계수를 대응시키는 사상이라
하자. $\varphi_{\lambda,\alpha}$들은
$(A/\mathfrak J_\lambda)$-가군 준동형들의 사영계를 이루며, 그
사영극한은 연속 준동형
\[
  \varphi_\alpha:A\{T_1,\ldots,T_r\}\longrightarrow A
\]
이다.

모든 $y\in A\{T_1,\ldots,T_r\}$에 대하여 형식적 멱급수
$\sum_\alpha\varphi_\alpha(y)T^\alpha$가 제한됨을 보이자.
$y_\lambda$를 $B_\lambda$에서 $y$의 성분이라 하고,
$H_\lambda$를 다항식 $y_\lambda$의 계수가 영이 아닌
$\alpha\in\mathbf N^r$ 전체의 유한집합이라 하자.
$\mathfrak J_\mu\subset\mathfrak J_\lambda$이고
$\alpha\notin H_\lambda$이면
$\varphi_{\lambda,\alpha}(y_\mu)\in\mathfrak J_\lambda$이고, 극한으로
넘기면 $\alpha\notin H_\lambda$에 대하여
$\varphi_\alpha(y)\in\mathfrak J_\lambda$이다.

따라서
\[
  \varphi(y)=\sum_\alpha\varphi_\alpha(y)T^\alpha
\]
로 놓아 환 준동형
$\varphi:A\{T_1,\ldots,T_r\}\to A'$를 정의할 수 있고,
$\varphi$가 연속임은 바로 알 수 있다. 역으로
$\theta_\lambda:A\to A/\mathfrak J_\lambda$를 표준 준동형이라 하자.
각 $z=\sum_\alpha c_\alpha T^\alpha\in A'$와 각 $\lambda$에 대하여
$\theta_\lambda(c_\alpha)\neq0$인 지표 $\alpha$는 유한 개뿐이므로
\[
  \psi_\lambda(z)=
  \sum_\alpha\theta_\lambda(c_\alpha)T^\alpha
  \in B_\lambda.
\]
$\psi_\lambda$들은 연속이고 준동형들의 사영계를 이루므로 그
사영극한은 연속 준동형
$\psi:A'\to A\{T_1,\ldots,T_r\}$이다. 끝으로
$\varphi\circ\psi$와 $\psi\circ\varphi$가 항등 자기동형임을 확인하면
되는데, 이는 명백하다.
\end{env}

\begin{env}[7.5.2]
\label{0.7.5.2-ko}
(7.5.1)에서 정의한 동형을 통하여 $A\{T_1,\ldots,T_r\}$를 제한
형식적 멱급수환 $A'$와 동일시하겠다. 표준 동형
\[
  ((A/\mathfrak J_\lambda)[T_1,\ldots,T_r])
  [T_{r+1},\ldots,T_s]
  \xrightarrow{\sim}
  (A/\mathfrak J_\lambda)[T_1,\ldots,T_s]
\]
들은 사영극한으로 넘겨 표준 동형
\[
  (A\{T_1,\ldots,T_r\})\{T_{r+1},\ldots,T_s\}
  \xrightarrow{\sim}
  A\{T_1,\ldots,T_s\}
\]
을 정의한다.
\end{env}

\begin{env}[7.5.3]
\label{0.7.5.3-ko}
$A$에서 선형 위상을 갖는 분리 완비환 $B$로 가는 임의의 연속
준동형 $u:A\to B$와 $B$의 $r$개 원소로 이루어진 임의의 계
$(b_1,\ldots,b_r)$에 대하여, 모든 $a\in A$와 $1\leq j\leq r$에
대해
\[
  \overline u(a)=u(a),\qquad \overline u(T_j)=b_j
\]
를 만족하는 연속 준동형
\[
  \overline u:A\{T_1,\ldots,T_r\}\longrightarrow B
\]
가 \emph{유일하게} 존재한다. 실제로
\[
  \overline u\Bigl(\sum_\alpha c_\alpha T^\alpha\Bigr)
  =\sum_\alpha u(c_\alpha)b_1^{\alpha_1}\cdots b_r^{\alpha_r}
\]
로 잡으면 충분하다. 족
$(u(c_\alpha)b_1^{\alpha_1}\cdots b_r^{\alpha_r})$가 $B$에서
가합이고 $\overline u$가 연속임을 확인하는 일은 즉시 되므로
독자에게 맡긴다. $B$와 $b_j$들이 임의라는 이 성질은 위상환
$A\{T_1,\ldots,T_r\}$를 유일한 동형을 제외하고
\emph{특징짓는다}는 점에 주의하자.
\end{env}

\begin{proposition}[7.5.4]
\label{0.7.5.4-ko}
\begin{enumerate}
  \item[(i)] $A$가 허용환이면
    $A'=A\{T_1,\ldots,T_r\}$도 허용환이다.
  \item[(ii)] $A$를 아딕환이라 하고, $\mathfrak J$를 $A$의 정의
    아이디얼로서 $\mathfrak J/\mathfrak J^2$가
    $A/\mathfrak J$ 위에서 유한형인 것이라 하자.

\oldpage[0\textsubscript{I}]{71}

    $\mathfrak J'=\mathfrak JA'$로 놓으면 $A'$는
    $\mathfrak J'$-아딕환이고
    $\mathfrak J'/\mathfrak J'^{\,2}$는
    $A'/\mathfrak J'$ 위에서 유한형이다. 더 나아가 $A$가 뇌터
    환이면 $A'$도 뇌터 환이다.
\end{enumerate}
\end{proposition}

\begin{enumerate}
  \item[(i)] $\mathfrak J$가 $A$의 아이디얼일 때, 모든
    $\alpha$에 대하여 $c_\alpha\in\mathfrak J$인
    $\sum_\alpha c_\alpha T^\alpha$들로 이루어진 $A'$의
    아이디얼을 $\mathfrak J'$라 하자. 그러면
    $(\mathfrak J')^n\subset(\mathfrak J^n)'$이다. 따라서
    $\mathfrak J$가 $A$의 정의 아이디얼이면 $\mathfrak J'$도
    $A'$의 정의 아이디얼이다.

  \item[(ii)] $A_i=A/\mathfrak J^{i+1}$로 놓고, $i\leq j$일 때
    $u_{ij}$를 표준 준동형
    $A/\mathfrak J^{j+1}\to A/\mathfrak J^{i+1}$이라 하자. 또한
    $A'_i=A_i[T_1,\ldots,T_r]$로 놓고, $u'_{ij}:A'_j\to A'_i$를
    $A'_j$의 다항식 계수들에 $u_{ij}$를 적용하여 얻는 준동형이라
    하자.

    사영계 $(A'_i,u'_{ij})$가 (7.2.7)의 조건들을 만족함을 보이자.
    $\mathfrak J'$는 $A'\to A'_0$의 핵이므로, 이 사실은 (ii)의 첫째
    주장을 증명할 것이다. $u'_{ij}$들이 전사임은 명백하다.
    $u'_{0i}$의 핵 $\mathfrak J'_i$는
    $A_i[T_1,\ldots,T_r]$의 다항식 가운데 그 계수들이
    $\mathfrak J/\mathfrak J^{i+1}$에 속하는 것들 전체이다. 특히
    $\mathfrak J'_1$은 $A_1[T_1,\ldots,T_r]$의 다항식 가운데
    계수들이 $\mathfrak J/\mathfrak J^2$에 속하는 것들 전체이다.
    $\mathfrak J/\mathfrak J^2$가
    $A_1=A/\mathfrak J^2$ 위에서 유한형이므로
    $\mathfrak J'_1/\mathfrak J_1'^{\,2}$는 $A'_1$ 위에서 유한형
    가군이다(동치로 $A'_0=A'_1/\mathfrak J'_1$ 위에서 유한형이다).

    $u'_{ij}$의 핵이 $\mathfrak J_j'^{\,i+1}$임을 보이자.
    $\mathfrak J_j'^{\,i+1}$가 이 핵에 포함됨은 명백하다. 한편
    $a_1,\ldots,a_m$을 $\mathfrak J$의 원소들 가운데
    $\mathfrak J^2$을 법으로 한 유들이
    $\mathfrak J/\mathfrak J^2$를 생성하는 것들이라 하자.
    $a_k$들($1\leq k\leq m$)로 이루어진 차수 $\leq j$인 단항식들의
    $\mathfrak J^{j+1}$을 법으로 한 유들이
    $\mathfrak J/\mathfrak J^{j+1}$을 생성함은 바로 확인된다.
    따라서 차수가 $>i$이고 $\leq j$인 단항식들의 유들은
    $\mathfrak J^{i+1}/\mathfrak J^{j+1}$을 생성한다. 그러한 원소를
    계수로 갖는 $T_k$들에 관한 단항식은
    $\mathfrak J'_j$의 $i+1$개 원소의 곱이고, 이것으로 주장이
    증명된다.

    끝으로 $A$가 뇌터 환이면
    $A'/\mathfrak J'=(A/\mathfrak J)[T_1,\ldots,T_r]$도 뇌터
    환이므로 $A'$가 뇌터 환이다(7.2.8).
\end{enumerate}

\begin{proposition}[7.5.5]
\label{0.7.5.5-ko}
$A$를 뇌터 $\mathfrak J$-아딕환, $B$를 허용 위상환이라 하고,
$\varphi:A\to B$를 $B$에 $A$-대수 구조를 주는 연속 준동형이라
하자. 다음 조건들은 동치이다.
\begin{enumerate}
  \item[a)] $B$는 뇌터 $\mathfrak JB$-아딕환이고,
    $B/\mathfrak JB$는 $A/\mathfrak J$ 위의 유한형 대수이다.

  \item[b)] $B$는 $\varprojlim B_n$과 위상적으로 $A$-동형이다.
    여기서 $m\geq n$에 대하여
    $B_n=B_m/\mathfrak J^{n+1}B_m$이고, $B_1$은
    $A_1=A/\mathfrak J^2$ 위의 유한형 대수이다.

  \item[c)] $B$는 $A\{T_1,\ldots,T_r\}$ 꼴의 대수를 어떤
    아이디얼로 나눈 몫과 위상적으로 $A$-동형이다. 이 아이디얼은
    (7.3.6)과 (7.5.4 (ii))에 따라 반드시 닫혀 있다.
\end{enumerate}
\end{proposition}

$A$가 뇌터 환이므로 $A'=A\{T_1,\ldots,T_r\}$도 뇌터 환이다
(7.5.4). 따라서 c)는 $B$가 뇌터 환임을 함의한다. 또한
$\mathfrak J'=\mathfrak JA'$는 $A'$의 $0$의 열린 근방이고
$\mathfrak J'^{n}$들이 $0$의 근방 기본계를 이룬다. 그 상들인
$\mathfrak J^nB$들이 $B$에서 $0$의 근방 기본계를 이루며, $B$가
분리 완비임을 알고 있으므로 $B$는 $\mathfrak JB$-아딕환이다.
끝으로 $B/\mathfrak JB$는
\[
  A'/\mathfrak JA'=(A/\mathfrak J)[T_1,\ldots,T_r]
\]
의 몫인 $A/\mathfrak J$-대수이므로 유한형이다. 이로써 c)가 a)를
함의함을 모두 증명하였다.

$B$가 뇌터 $\mathfrak JB$-아딕환이면
\[
  B\simeq\varprojlim B_n,
  \qquad B_n=B/\mathfrak J^{n+1}B
\]
이다(7.2.11). 또한
$\mathfrak JB/\mathfrak J^2B$는 $B/\mathfrak JB$ 위의 유한형
가군이다. $(a_j)_{1\leq j\leq s}$를
$B/\mathfrak JB$-가군 $\mathfrak JB/\mathfrak J^2B$의 생성계라
하고, $(c_i)_{1\leq i\leq r}$를 $B/\mathfrak J^2B$의 원소들
가운데 그 유들이 $\mathfrak JB/\mathfrak J^2B$를 법으로 하여
$A/\mathfrak J$-대수 $B/\mathfrak JB$를 생성하는 것들이라 하자.

\oldpage[0\textsubscript{I}]{72}

곱 $c_i a_j$들은 $A/\mathfrak J^2$-대수
$B/\mathfrak J^2B$의 생성계를 이룸을 곧 알 수 있다. 따라서 a)는
b)를 함의한다.

b)가 c)를 함의함을 보이는 일만 남았다. 가정에 따라 $B_1$은 뇌터
환이다. $B_1=B_2/\mathfrak J^2B_2$이므로
$\mathfrak J^2B_1=0$이고, 따라서
\[
  \mathfrak JB_1
  =\mathfrak JB_1/\mathfrak J^2B_1
\]
은 유한형 $B_0$-가군이다. 그러므로 사영계 $(B_n)$은 (7.2.7)의
조건들을 만족하고 $B$는 $\mathfrak JB$-아딕환이다.

$(c_i)_{1\leq i\leq r}$를 $B$의 원소들의 유한계로서, 그
$\mathfrak JB$를 법으로 한 유들이 $A/\mathfrak J$-대수
$B/\mathfrak JB$를 생성하고, 그 원소들의 $\mathfrak J$-계수
선형결합들의 $\mathfrak J^2B$를 법으로 한 유들이
$B_0$-가군 $\mathfrak JB/\mathfrak J^2B$를 생성하도록 잡자.
(7.5.3)에 따라 $A$에서는 $\varphi$로 제한되고
$1\leq i\leq r$에 대하여 $u(T_i)=c_i$인 연속 $A$-준동형
\[
  u:A'=A\{T_1,\ldots,T_r\}\longrightarrow B
\]
가 존재한다.

$u$가 \emph{전사}임을 보이면 c)가 증명된다. 실제로
$u(A')=B$에서 $u(\mathfrak J^nA')=\mathfrak J^nB$가 따르므로 $u$는
위상환의 엄밀 사상이고, 따라서 $B$는 $A'$를 어떤 닫힌 아이디얼로
나눈 몫과 동형이다. $B$는 $\mathfrak JB$-아딕 위상에 관하여
완비이므로, $\mathfrak J$-아딕 여과를 준 $A'$와 $B$에 대하여
$u$에서 표준적으로 유도되는 준동형
\[
  \operatorname{grad}(A')\longrightarrow\operatorname{grad}(B)
\]
가 전사임을 보이면 충분하다([I], p.~18-07).

정의에 따라 $u$에서 유도되는 준동형
\[
  A'/\mathfrak JA'\longrightarrow B/\mathfrak JB,
  \qquad
  \mathfrak JA'/\mathfrak J^2A'\longrightarrow
  \mathfrak JB/\mathfrak J^2B
\]
들은 전사이다. $n$에 관한 귀납법으로
\[
  \mathfrak JA'/\mathfrak J^nA'\longrightarrow
  \mathfrak JB/\mathfrak J^nB
\]
도 전사임을 곧 얻고, 따라서 더더욱
\[
  \mathfrak J^nA'/\mathfrak J^{n+1}A'\longrightarrow
  \mathfrak J^nB/\mathfrak J^{n+1}B
\]
도 전사이다. 이것으로 증명이 끝난다.

\subsection{완비 분수환}
\label{subsection:0.7.6-ko}

\begin{env}[7.6.1]
\label{0.7.6.1-ko}
$A$를 선형 위상을 갖는 환, $(\mathfrak J_\lambda)$를 아이디얼들로
이루어진 $A$의 $0$의 근방 기본계, $S$를 $A$의 곱셈적 부분집합이라
하자. $u_\lambda$를 표준 준동형
$A\to A_\lambda=A/\mathfrak J_\lambda$라 하고,
$\mathfrak J_\mu\subset\mathfrak J_\lambda$일 때 $u_{\lambda\mu}$를
표준 준동형 $A_\mu\to A_\lambda$라 하자. 또한
$S_\lambda=u_\lambda(S)$로 놓으면
$u_{\lambda\mu}(S_\mu)=S_\lambda$이다. $u_{\lambda\mu}$들은 표준적으로
전사 준동형
\[
  S_\mu^{-1}A_\mu\longrightarrow S_\lambda^{-1}A_\lambda
\]
를 주며, 이 환들은 그 준동형들에 관하여 사영계를 이룬다. 이 계의
사영극한을 $A\{S^{-1}\}$로 나타내겠다. 이 정의는 선택한 근방
기본계 $(\mathfrak J_\lambda)$에 의존하지 않는다. 실제로 다음이
성립한다.
\end{env}

\begin{proposition}[7.6.2]
\label{0.7.6.2-ko}
환 $A\{S^{-1}\}$는 $S^{-1}\mathfrak J_\lambda$들이 $0$의 근방
기본계를 이루는 위상을 준 환 $S^{-1}A$의 분리 완비화와 위상적으로
동형이다.
\end{proposition}

실제로 $v_\lambda$를 $u_\lambda$에서 유도되는 표준 준동형
$S^{-1}A\to S_\lambda^{-1}A_\lambda$라 하면, $v_\lambda$의 핵은
$S^{-1}\mathfrak J_\lambda$이고 $v_\lambda$는 전사이다. 따라서
명제는 (7.2.1)에서 따른다.

\begin{corollary}[7.6.3]
\label{0.7.6.3-ko}
$S'$을 $A$의 분리 완비화 $\widehat A$에서 $S$의 표준상이라 하면
$A\{S^{-1}\}$는 $\widehat A\{{S'}^{-1}\}$와 표준적으로
동일시된다.
\end{corollary}

$A$가 분리 완비환이어도 $S^{-1}A$는
$S^{-1}\mathfrak J_\lambda$들이 정의하는 위상에 관하여 분리
완비환일 필요가 없다. 예를 들어 $f$가 위상적 멱영이지만 멱영은 아닌
원소이고, $S$가 $f^n$들($n\geq0$)의 집합인 경우를 생각하자.
$S^{-1}A$는 영환이 아니다. 한편 각 $\lambda$에 대하여
$f^n\in\mathfrak J_\lambda$인 $n$이 존재하므로
\[
  1=\frac{f^n}{f^n}\in S^{-1}\mathfrak J_\lambda,
  \qquad S^{-1}\mathfrak J_\lambda=S^{-1}A.
\]

\oldpage[0\textsubscript{I}]{73}

\begin{corollary}[7.6.4]
\label{0.7.6.4-ko}
$A$에서 $0$이 $S$의 폐포에 속하지 않으면 환 $A\{S^{-1}\}$는
영환이 아니다.
\end{corollary}

실제로 환 $S^{-1}A$에서 $0$은 $\{1\}$의 폐포에 속하지 않는다.
그렇지 않으면 $A$의 모든 열린 아이디얼 $\mathfrak J_\lambda$에
대하여 $1\in S^{-1}\mathfrak J_\lambda$일 것이고, 따라서 모든
$\lambda$에 대하여
$\mathfrak J_\lambda\cap S\neq\varnothing$가 되어 가정에 모순이다.

\begin{env}[7.6.5]
\label{0.7.6.5-ko}
$A\{S^{-1}\}$를 $S$에 분모를 갖는 $A$의 \emph{완비 분수환}이라
하겠다. 앞의 표기에서 $A$에서 $S^{-1}\mathfrak J_\lambda$의
역상은 $\mathfrak J_\lambda$를 포함함이 명백하므로 표준 사상
$A\to S^{-1}A$는 연속이다. 이를 표준 사상
$S^{-1}A\to A\{S^{-1}\}$와 합성하면 표준 연속 준동형
\[
  A\longrightarrow A\{S^{-1}\}
\]
를 얻는데, 이는 준동형 $A\to S_\lambda^{-1}A_\lambda$들의
사영극한이다.
\end{env}

\begin{env}[7.6.6]
\label{0.7.6.6-ko}
$A\{S^{-1}\}$와 표준 사상 $A\to A\{S^{-1}\}$의 쌍은 다음
\emph{보편 성질}로 특징지어진다. $A$에서 선형 위상을 갖는 분리
완비환 $B$로 가는 연속 준동형 $u$가 $u(S)$의 원소들을 $B$의
가역원으로 보내면, $u$는 유일한 방식으로
\[
  A\longrightarrow A\{S^{-1}\}\xrightarrow{u'}B
\]
로 분해된다. 여기서 $u'$은 연속이다.

실제로 $u$는 유일한 방식으로
$A\to S^{-1}A\xrightarrow{v'}B$로 분해된다. $B$의 각 열린
아이디얼 $\mathfrak K$에 대하여 $u^{-1}(\mathfrak K)$는 어떤
$\mathfrak J_\lambda$를 포함하므로
${v'}^{-1}(\mathfrak K)$는 반드시
$S^{-1}\mathfrak J_\lambda$를 포함한다. 따라서 $v'$은 연속이다.
$B$가 분리 완비환이므로 $v'$은 유일한 방식으로
\[
  S^{-1}A\longrightarrow A\{S^{-1}\}\xrightarrow{u'}B
\]
로 분해되며 $u'$은 연속이다. 이로써 주장이 증명된다.
\end{env}

\begin{env}[7.6.7]
\label{0.7.6.7-ko}
$B$를 또 하나의 선형 위상환, $T$를 $B$의 곱셈적 부분집합,
$\varphi:A\to B$를 $\varphi(S)\subset T$를 만족하는 연속
준동형이라 하자. 앞의 결과에 따라 연속 준동형
\[
  A\xrightarrow{\varphi}B\longrightarrow B\{T^{-1}\}
\]
은 유일한 방식으로
\[
  A\longrightarrow A\{S^{-1}\}
  \xrightarrow{\varphi'}B\{T^{-1}\}
\]
로 분해되며 $\varphi'$은 연속이다. 특히 $B=A$이고 $\varphi$가
항등사상일 때, $S\subset T$이면
$S^{-1}A\to T^{-1}A$를 분리 완비화하여 얻는 연속 준동형
\[
  \rho^{T,S}:A\{S^{-1}\}\longrightarrow A\{T^{-1}\}
\]
가 존재한다. $S\subset T\subset U$인 세 번째 곱셈적 부분집합
$U$에 대하여
\[
  \rho^{U,S}=\rho^{U,T}\circ\rho^{T,S}
\]
이다.
\end{env}

\begin{env}[7.6.8]
\label{0.7.6.8-ko}
$S_1,S_2$를 $A$의 두 곱셈적 부분집합, $S_2'$을
$A\{S_1^{-1}\}$에서 $S_2$의 표준상이라 하자. 그러면 표준
위상동형
\[
  A\{(S_1S_2)^{-1}\}
  \xrightarrow{\sim}
  A\{S_1^{-1}\}\{{S_2'}^{-1}\}
\]
이 존재한다. 실제로 $S_2''$을 $S_1^{-1}A$에서 $S_2$의 표준상이라
하면 표준 동형
\[
  (S_1S_2)^{-1}A\xrightarrow{\sim}{S_2''}^{-1}(S_1^{-1}A)
\]
은 양방향 연속이고, 이 동형에서 결론이 따른다.
\end{env}

\begin{env}[7.6.9]
\label{0.7.6.9-ko}
$\mathfrak a$를 $A$의 \emph{열린} 아이디얼이라 하자. 모든
$\lambda$에 대하여
$\mathfrak J_\lambda\subset\mathfrak a$라고 가정해도 된다. 따라서
$S^{-1}A$에서
$S^{-1}\mathfrak J_\lambda\subset S^{-1}\mathfrak a$이고,
$S^{-1}\mathfrak a$는 $S^{-1}A$의 \emph{열린} 아이디얼이다. 그
분리 완비화를
\[
  \mathfrak a\{S^{-1}\}
  =\varprojlim
  (S^{-1}\mathfrak a/S^{-1}\mathfrak J_\lambda)
\]
로 나타내겠다. 이는 $A\{S^{-1}\}$의 열린 아이디얼이며
$S^{-1}\mathfrak a$의 표준상의 폐포와 동형이다. 더 나아가
\emph{이산환 $A\{S^{-1}\}/\mathfrak a\{S^{-1}\}$은
$S^{-1}A/S^{-1}\mathfrak a=S^{-1}(A/\mathfrak a)$와 표준적으로
동형이다.}

역으로 $\mathfrak a'$이 $A\{S^{-1}\}$의 열린 아이디얼이면
$\mathfrak a'$은 $\mathfrak J_\lambda\{S^{-1}\}$ 꼴의 아이디얼을
포함한다. 따라서 $\mathfrak a'$은
$S^{-1}A/S^{-1}\mathfrak J_\lambda$의 어떤 아이디얼의 역상이고,
그 아이디얼은 반드시 (1.2.6)에 따라
$\mathfrak a\supset\mathfrak J_\lambda$인 어떤 $\mathfrak a$에 대한
$S^{-1}\mathfrak a$ 꼴이다. 그러므로
$\mathfrak a'=\mathfrak a\{S^{-1}\}$이다. 특히 (1.2.6)에 따라
다음이 성립한다.
\end{env}

\begin{proposition}[7.6.10]
\label{0.7.6.10-ko}
대응
\[
  \mathfrak p\longmapsto\mathfrak p\{S^{-1}\}
\]
은 $\mathfrak p\cap S=\varnothing$인 $A$의 열린 소 아이디얼
$\mathfrak p$들의 집합에서 $A\{S^{-1}\}$의 열린 소 아이디얼들의
집합으로 가는 포함관계를 보존하는 전단사이다. 더 나아가

\oldpage[0\textsubscript{I}]{74}

$A\{S^{-1}\}/\mathfrak p\{S^{-1}\}$의 분수체는
$A/\mathfrak p$의 분수체와 표준적으로 동형이다.
\end{proposition}

\begin{proposition}[7.6.11]
\label{0.7.6.11-ko}
\begin{enumerate}
  \item[(i)] $A$가 허용환이면 $A'=A\{S^{-1}\}$도 허용환이고,
    $A$의 모든 정의 아이디얼 $\mathfrak J$에 대하여
    $\mathfrak J'=\mathfrak J\{S^{-1}\}$은 $A'$의 정의
    아이디얼이다.
  \item[(ii)] $A$를 아딕환, $\mathfrak J$를 $A$의 정의
    아이디얼로서 $\mathfrak J/\mathfrak J^2$가
    $A/\mathfrak J$ 위에서 유한형인 것이라 하자. 그러면 $A'$는
    $\mathfrak J'$-아딕환이고
    $\mathfrak J'/\mathfrak J'^2$는 $A'/\mathfrak J'$ 위에서
    유한형이다. 더 나아가 $A$가 뇌터 환이면 $A'$도 뇌터 환이다.
\end{enumerate}
\end{proposition}

\begin{enumerate}
  \item[(i)] $\mathfrak J$가 $A$의 정의 아이디얼이면
    \[
      (S^{-1}\mathfrak J)^n=S^{-1}\mathfrak J^n
    \]
    이므로 $S^{-1}\mathfrak J$는 위상환 $S^{-1}A$의 정의
    아이디얼임이 명백하다. $A''$을 $S^{-1}A$에 연관된 분리환,
    $\mathfrak J''$을 $A''$에서 $S^{-1}\mathfrak J$의 상이라 하자.
    $S^{-1}\mathfrak J^n$의 상은 $\mathfrak J''^n$이므로
    $\mathfrak J''^n$은 $A''$에서 $0$으로 수렴한다.
    $\mathfrak J'$은 $A'$에서 $\mathfrak J''$의 폐포이므로
    $\mathfrak J'^n$은 $\mathfrak J''^n$의 폐포에 포함되고, 따라서
    $A'$에서 $0$으로 수렴한다.

  \item[(ii)] $A_i=A/\mathfrak J^{i+1}$로 놓고, $i\leq j$일 때
    $u_{ij}$를 표준 준동형
    $A/\mathfrak J^{j+1}\to A/\mathfrak J^{i+1}$이라 하자. $S_i$를
    $A_i$에서 $S$의 표준상이라 하고 $A_i'=S_i^{-1}A_i$로 놓자.
    끝으로 $u_{ij}':A_j'\to A_i'$을 $u_{ij}$에서 표준적으로
    유도되는 준동형이라 하자.

    사영계 $(A_i',u_{ij}')$가 명제 (7.2.7)의 조건들을 만족함을
    보이자. $u_{ij}'$들이 전사임은 명백하다. 한편 $u_{ij}'$의 핵은
    (1.3.2)에 따라
    \[
      S_j^{-1}(\mathfrak J^{i+1}/\mathfrak J^{j+1})
      =\mathfrak J_j'^{\,i+1}
    \]
    이다. 여기서
    \[
      \mathfrak J_j'=S_j^{-1}(\mathfrak J/\mathfrak J^{j+1}).
    \]
    또한
    \[
      \mathfrak J_1'/\mathfrak J_1'^2
      =S_1^{-1}(\mathfrak J/\mathfrak J^2).
    \]
    $\mathfrak J/\mathfrak J^2$가 $A/\mathfrak J^2$ 위에서
    유한형이므로 $\mathfrak J_1'/\mathfrak J_1'^2$는 $A_1'$ 위에서
    유한형이다. 끝으로 $A$가 뇌터 환이면
    $A_0'=S_0^{-1}(A/\mathfrak J)$도 뇌터 환이다. 이로써 명제가
    증명된다(7.2.8).
\end{enumerate}

\begin{corollary}[7.6.12]
\label{0.7.6.12-ko}
(7.6.11 (ii))의 가정 아래에서
\[
  (\mathfrak J\{S^{-1}\})^n=\mathfrak J^n\{S^{-1}\}
\]
이다.
\end{corollary}

실제로 이는 (7.2.7)과 (7.6.11)의 증명에서 따른다.

\begin{proposition}[7.6.13]
\label{0.7.6.13-ko}
$A$를 뇌터 아딕환, $S$를 $A$의 곱셈적 부분집합이라 하자. 그러면
$A\{S^{-1}\}$는 평탄 $A$-가군이다.
\end{proposition}

실제로 $\mathfrak J$가 $A$의 정의 아이디얼이면
$A\{S^{-1}\}$는 $S^{-1}\mathfrak J$-준아딕 위상을 준 뇌터 환
$S^{-1}A$의 분리 완비화이다. 따라서 (7.3.3)에 따라
$A\{S^{-1}\}$는 평탄 $S^{-1}A$-가군이다. $S^{-1}A$는 평탄
$A$-가군이므로(6.3.1), 명제는 평탄성의 추이성에서 따른다(6.2.1).

\begin{corollary}[7.6.14]
\label{0.7.6.14-ko}
(7.6.13)의 가정 아래에서 $S'\subset S$인 $A$의 또 하나의 곱셈적
부분집합 $S'$을 잡자. 그러면 $A\{S^{-1}\}$는 평탄
$A\{{S'}^{-1}\}$-가군이다.
\end{corollary}

실제로 (7.6.8)에 따라 $A\{S^{-1}\}$는
$A\{{S'}^{-1}\}\{S_0^{-1}\}$와 표준적으로 동일시된다. 여기서
$S_0$은 $A\{{S'}^{-1}\}$에서 $S$의 표준상이고,
$A\{{S'}^{-1}\}$는 뇌터 환이다(7.6.11).

\begin{env}[7.6.15]
\label{0.7.6.15-ko}
선형 위상환 $A$의 각 원소 $f$에 대하여 $S_f$를
$f^n$들($n\geq0$)의 곱셈적 집합이라 하고, 완비 분수환
$A\{S_f^{-1}\}$를 $A_{\{f\}}$로 나타내겠다. $A$의 각 열린
아이디얼 $\mathfrak a$에 대하여
$\mathfrak a\{S_f^{-1}\}$ 대신 $\mathfrak a_{\{f\}}$라고 쓰겠다.
$g$가 $A$의 또 하나의 원소이면 표준 연속 준동형
\[
  A_{\{f\}}\longrightarrow A_{\{fg\}}
\]
이 존재한다(7.6.7). $f$가 $A$의 곱셈적 부분집합 $S$를 움직일 때
$A_{\{f\}}$들은 이 준동형들에 관한 환들의 여과 귀납계를 이룬다.
다음과 같이 놓자.
\[
  A_{\{S\}}=\varinjlim_{f\in S}A_{\{f\}}.
\]
각 $f\in S$에 대하여 준동형
$A_{\{f\}}\to A\{S^{-1}\}$가 있고(7.6.7), 이 준동형들은

\oldpage[0\textsubscript{I}]{75}

귀납계를 이룬다. 따라서 귀납극한으로 넘기면 표준 준동형
\[
  A_{\{S\}}\longrightarrow A\{S^{-1}\}
\]
을 얻는다.
\end{env}

\begin{proposition}[7.6.16]
\label{0.7.6.16-ko}
$A$가 뇌터 환이면 $A\{S^{-1}\}$는 평탄 $A_{\{S\}}$-가군이다.
\end{proposition}

실제로 (7.6.14)에 따라 $A\{S^{-1}\}$는 각 $f\in S$에 대하여
평탄 $A_{\{f\}}$-가군이고, 결론은 (6.2.3)에서 따른다.

\begin{proposition}[7.6.17]
\label{0.7.6.17-ko}
$\mathfrak p$를 허용환 $A$의 열린 소 아이디얼이라 하고
$S=A-\mathfrak p$로 놓자. 그러면 환 $A\{S^{-1}\}$와
$A_{\{S\}}$는 국소환이고, 표준 준동형
\[
  A_{\{S\}}\longrightarrow A\{S^{-1}\}
\]
은 국소 준동형이며, $A_{\{S\}}$와 $A\{S^{-1}\}$의 잉여체들은
$A/\mathfrak p$의 분수체와 표준적으로 동형이다.
\end{proposition}

실제로 $\mathfrak J\subset\mathfrak p$인 $A$의 정의 아이디얼
$\mathfrak J$를 잡자. 그러면
\[
  S^{-1}\mathfrak J
  \subset S^{-1}\mathfrak p
  =\mathfrak pA_\mathfrak p
\]
이므로 $A_\mathfrak p/S^{-1}\mathfrak J$는 국소환이다.
(7.1.12), (7.6.9), (7.6.11 (i))에 따라 $A\{S^{-1}\}$가
국소환임이 따른다.

\[
  \mathfrak m=\varinjlim_{f\in S}\mathfrak p_{\{f\}}
\]
로 놓자. 이는 $A_{\{S\}}$의 아이디얼이다. $\mathfrak m$에 속하지
않는 $A_{\{S\}}$의 모든 원소가 가역임을 보이겠다. 그러한 원소는
어떤 $f\in S$에 대하여
$\mathfrak p_{\{f\}}$에 속하지 않는 $z\in A_{\{f\}}$의
$A_{\{S\}}$에서의 상이다. $z$의
\[
  A_{\{f\}}/\mathfrak J_{\{f\}}
  =S_f^{-1}(A/\mathfrak J)
\]
에서의 표준상 $z_0$는
$S_f^{-1}(\mathfrak p/\mathfrak J)$에 속하지 않는다(7.6.9).
따라서
\[
  z_0=\frac{\overline x}{\overline f^{\,k}},
  \qquad x\notin\mathfrak p,
\]
이다. 여기서 $\overline x,\overline f$는 $\mathfrak J$를 법으로 한
$x,f$의 유이다. $x\in S$이므로 $g=xf\in S$이다. $S_g^{-1}A$에서
\[
  y_0=\frac{x^{k+1}}{g^k},
\]
즉 $x/f^k\in S_f^{-1}A$의 표준상은
\[
  \frac{x^{k-1}f^{2k}}{g^k}
\]
를 역원으로 갖는다. 따라서 더더욱 $S_g^{-1}A/S_g^{-1}\mathfrak J$에서
$y_0$의 상은 가역이고, (7.6.9)와 (7.1.12)에 따라
$A_{\{g\}}$에서 $z$의 표준상 $y$도 가역이다. 그러므로
$A_{\{S\}}$에서 $z$의 상, 즉 $y$의 상은 가역이다.

따라서 $A_{\{S\}}$는 극대 아이디얼 $\mathfrak m$을 갖는
국소환이다. 더 나아가 $A\{S^{-1}\}$에서
$\mathfrak p_{\{f\}}$의 상은 이 환의 극대 아이디얼
$\mathfrak p\{S^{-1}\}$에 포함된다. 따라서
$A\{S^{-1}\}$에서 $\mathfrak m$의 상도
$\mathfrak p\{S^{-1}\}$에 포함되고, 표준 준동형
$A_{\{S\}}\to A\{S^{-1}\}$는 국소 준동형이다.

끝으로 $A\{S^{-1}\}/\mathfrak p\{S^{-1}\}$의 모든 원소는 적당한
$f\in S$에 대한 환 $S_f^{-1}A$의 어떤 원소의 상이다. 따라서
준동형
\[
  A_{\{S\}}\longrightarrow
  A\{S^{-1}\}/\mathfrak p\{S^{-1}\}
\]
은 전사이고, 몫으로 내려가면 잉여체들의 동형을 준다.

\begin{corollary}[7.6.18]
\label{0.7.6.18-ko}
(7.6.17)의 가정에 더하여 $A$가 뇌터 아딕환이라고 하자. 그러면
국소환 $A\{S^{-1}\}$와 $A_{\{S\}}$는 뇌터 환이고,
$A\{S^{-1}\}$는 충실 평탄 $A_{\{S\}}$-가군이다.
\end{corollary}

$A\{S^{-1}\}$가 뇌터 환임은 이미 (7.6.11 (ii))에서 알고 있고,
이는 $A_{\{S\}}$-평탄이다(7.6.16). 준동형
$A_{\{S\}}\to A\{S^{-1}\}$가 국소 준동형이므로
$A\{S^{-1}\}$는 충실 평탄 $A_{\{S\}}$-가군이고(6.6.2), 따라서
$A_{\{S\}}$는 뇌터 환이다(6.5.2).

\subsection{완비 텐서곱}
\label{subsection:0.7.7-ko}

\begin{env}[7.7.1]
\label{0.7.7.1-ko}
$A$를 선형 위상을 갖는 환, $M,N$을 선형 위상을 갖는 두
$A$-가군이라 하자. $\mathfrak J,V,W$를 각각 $A,M,N$에서 $0$의
열린 근방으로서 $A$-부분가군이고
$\mathfrak JM\subset V$, $\mathfrak JN\subset W$를 만족하는 것이라
하자. 그러면 $M/V$와 $N/W$를 $(A/\mathfrak J)$-가군으로 볼 수 있다.
$\mathfrak J,V,W$가 앞의 조건들을 만족하는 열린 근방들을
움직일 때, 가군
\[
  (M/V)\otimes_{A/\mathfrak J}(N/W)
\]

\oldpage[0\textsubscript{I}]{76}

들은 환들의 사영계 $A/\mathfrak J$ 위의 가군들의 사영계를 이룸이
명백하다. 따라서 사영극한을 취하면 $A$의 분리 완비화
$\widehat A$ 위의 가군을 얻는다. 이를 $M$과 $N$의
\emph{완비 텐서곱}이라 하고 $(M\otimes_A N)^\wedge$로 쓴다.
$M/V$는 $\widehat M/\overline V$와 표준적으로 동형이다. 여기서
$\widehat M$은 $M$의 분리 완비화이고 $\overline V$는 $\widehat M$에서
$V$의 상의 폐포이다. 따라서 완비 텐서곱
$(M\otimes_A N)^\wedge$는
$(\widehat M\otimes_{\widehat A}\widehat N)^\wedge$와 표준적으로
동일시되며, 이를
$\widehat M\widehat\otimes_{\widehat A}\widehat N$으로도 쓴다.
\end{env}

\begin{env}[7.7.2]
\label{0.7.7.2-ko}
앞의 표기 아래에서 텐서곱
$(M/V)\otimes_A(N/W)$와
$(M/V)\otimes_{A/\mathfrak J}(N/W)$는 표준적으로 동일시된다. 따라서
이들은 또한
\[
  (M\otimes_A N)/
  \bigl(\Im(V\otimes_A N)+\Im(M\otimes_A W)\bigr)
\]
와 동일시된다. 그러므로 $(M\otimes_A N)^\wedge$는
\emph{$A$-가군 $M\otimes_A N$에 다음 위상을 주었을 때의 분리
완비화이다. 즉 부분가군}
\[
  \Im(V\otimes_A N)+\Im(M\otimes_A W)
\]
\emph{들이 $0$의 근방 기본계를 이룬다}. 여기서 $V$와 $W$는 각각
$M$과 $N$의 열린 부분가군 전체를 움직인다. 줄여서 이 위상을
$M$과 $N$에 주어진 위상들의 \emph{텐서곱}이라 하겠다.
\end{env}

\begin{env}[7.7.3]
\label{0.7.7.3-ko}
$M',N'$을 선형 위상을 갖는 두 $A$-가군,
$u:M\to M'$, $v:N\to N'$을 두 연속 준동형이라 하자. $M\otimes_A N$과
$M'\otimes_A N'$에 각각 텐서곱 위상을 주면 $u\otimes v$가 연속임은
명백하다. 분리 완비화로 옮겨 가면 연속 준동형
\[
  (M\otimes_A N)^\wedge\longrightarrow
  (M'\otimes_A N')^\wedge
\]
을 얻으며, 이를 $u\widehat\otimes v$로 쓰겠다. 따라서
$(M\otimes_A N)^\wedge$는 선형 위상을 갖는 $A$-가군들의 범주에서
$M$과 $N$에 관한 \emph{쌍함자}이다.
\end{env}

\begin{env}[7.7.4]
\label{0.7.7.4-ko}
선형 위상을 갖는 유한 개의 $A$-가군들의 완비 텐서곱도 같은 방식으로
정의한다. 이 텐서곱이 보통의 결합법칙과 교환법칙을 만족함은
명백하다.
\end{env}

\begin{env}[7.7.5]
\label{0.7.7.5-ko}
$B,C$가 선형 위상을 갖는 두 위상 $A$-대수이면, $B\otimes_A C$의
텐서곱 위상에서 $0$의 근방 기본계는 $B\otimes_A C$의 아이디얼
\[
  \Im(\mathfrak K\otimes_A C)+
  \Im(B\otimes_A\mathfrak L)
\]
들로 이루어진다. 여기서 $\mathfrak K$와 $\mathfrak L$은 각각
$B$와 $C$의 열린 아이디얼 전체를 움직인다. 따라서
$(B\otimes_A C)^\wedge$는 $(A/\mathfrak J)$-대수들의 사영계
\[
  (B/\mathfrak K)\otimes_{A/\mathfrak J}(C/\mathfrak L)
\]
의 사영극한으로서 \emph{$\widehat A$ 위의 위상대수} 구조를 갖는다.
여기서 $\mathfrak J$는
$\mathfrak J B\subset\mathfrak K$,
$\mathfrak J C\subset\mathfrak L$을 만족하는 $A$의 열린 아이디얼이며,
그러한 아이디얼은 언제나 존재한다. 이 대수를 $B$와 $C$의
\emph{완비 텐서곱}이라 한다.
\end{env}

\begin{env}[7.7.6]
\label{0.7.7.6-ko}
$B$와 $C$에서 $B\otimes_A C$로 가는 $A$-대수 준동형
\[
  b\longmapsto b\otimes1,
  \qquad
  c\longmapsto1\otimes c
\]
는 뒤의 대수에 텐서곱 위상을 주면 연속이다. 따라서 이를
$B\otimes_A C$에서 그 분리 완비화로 가는 표준 준동형과 합성하면
\emph{표준 준동형}
\[
  \rho:B\longrightarrow(B\otimes_A C)^\wedge,
  \qquad
  \sigma:C\longrightarrow(B\otimes_A C)^\wedge
\]
를 얻는다. 더 나아가 대수 $(B\otimes_A C)^\wedge$와 준동형
$\rho,\sigma$는 다음 \emph{보편 성질}을 갖는다. 선형 위상을 갖고
분리이며 완비인 임의의 $A$-대수 $D$와 두 연속 $A$-준동형
$u:B\to D$, $v:C\to D$에 대하여
\[
  w:(B\otimes_A C)^\wedge\longrightarrow D
\]
인 유일한 연속 $A$-준동형 $w$가 존재하여
$u=w\circ\rho$, $v=w\circ\sigma$를 만족한다.

실제로 이미 $u(b)=w_0(b\otimes1)$,
$v(c)=w_0(1\otimes c)$를 만족하는 유일한 $A$-준동형
$w_0:B\otimes_A C\to D$가 존재한다. 따라서 $w_0$가 연속임을
보이면 충분하다. 그러면 분리 완비화로 옮겨 가서 연속 준동형 $w$를
얻기 때문이다.

\oldpage[0\textsubscript{I}]{77}

이제 $\mathfrak M$이 $D$의 열린 아이디얼이면, 가정에
따라
\[
  u(\mathfrak K)\subset\mathfrak M,
  \qquad
  v(\mathfrak L)\subset\mathfrak M
\]
을 만족하는 열린 아이디얼
$\mathfrak K\subset B$, $\mathfrak L\subset C$가 존재한다. 그러면
\[
  \Im(\mathfrak K\otimes_A C)+
  \Im(B\otimes_A\mathfrak L)
\]
의 $w_0$에 의한 상도 $\mathfrak M$에 포함되므로 주장이 따른다.
\end{env}

\begin{proposition}[7.7.7]
\label{0.7.7.7-ko}
$B$와 $C$가 두 준허용 $A$-대수이면
$(B\otimes_A C)^\wedge$는 허용환이다. 또한 $\mathfrak K$와
$\mathfrak L$이 각각 $B$와 $C$의 정의 아이디얼이면,
$(B\otimes_A C)^\wedge$에서
\[
  \mathfrak H=
  \Im(\mathfrak K\otimes_A C)+
  \Im(B\otimes_A\mathfrak L)
\]
의 표준상의 폐포는 정의 아이디얼이다.
\end{proposition}

텐서곱 위상에 관하여 $\mathfrak H^n$이 $0$으로 수렴함을 보이면
충분하다. 이는 포함관계
\[
  \mathfrak H^{2n}\subset
  \Im(\mathfrak K^n\otimes_A C)+
  \Im(B\otimes_A\mathfrak L^n)
\]
에서 곧바로 따른다.

\begin{proposition}[7.7.8]
\label{0.7.7.8-ko}
$A$를 준아딕환, $\mathfrak J$를 $A$의 정의 아이디얼,
$M$을 $\mathfrak J$-준아딕 위상을 준 \emph{유한형} $A$-가군이라
하자. 임의의 아딕이고 뇌터인 위상 $A$-대수 $B$에 대하여
$B\otimes_A M$은 완비 텐서곱
$(B\otimes_A M)^\wedge$와 동일시된다.
\end{proposition}

$\mathfrak K$가 $B$의 정의 아이디얼이면, 가정에 따라
$\mathfrak J^mB\subset\mathfrak K$를 만족하는 정수 $m$이 존재한다.
따라서
\[
  \Im(B\otimes_A\mathfrak J^{nm}M)
  =\Im(\mathfrak J^{nm}B\otimes_A M)
  \subset\Im(\mathfrak K^nB\otimes_A M)
  =\mathfrak K^n(B\otimes_A M).
\]
그러므로 $B\otimes_A M$ 위에서 $B$와 $M$의 위상들의 텐서곱은
$\mathfrak K$-준아딕 위상이다. $B\otimes_A M$은 유한형
$B$-가군이므로 명제는 (7.3.6)에서 곧바로 따른다.

\subsection{준동형 가군의 위상}
\label{subsection:0.7.8-ko}

\begin{env}[7.8.1]
\label{0.7.8.1-ko}
$A$를 뇌터 $\mathfrak J$-아딕환, $M,N$을 $\mathfrak J$-준아딕 위상을
준 두 유한형 $A$-가군이라 하자. 이들은 분리이며 완비임을 알고
있다(7.3.6). 더 나아가 모든 $A$-준동형 $M\to N$은 자동으로
연속이고, $A$-가군 $\operatorname{Hom}_A(M,N)$은 유한형이다.
각 정수 $i\geq0$에 대하여
\[
  A_i=A/\mathfrak J^{i+1},
  \qquad
  M_i=M/\mathfrak J^{i+1}M,
  \qquad
  N_i=N/\mathfrak J^{i+1}N
\]
으로 놓자. $i\leq j$일 때 모든 준동형 $u_j:M_j\to N_j$는
$\mathfrak J^{i+1}M_j$를 $\mathfrak J^{i+1}N_j$ 안으로 보내므로,
몫으로 내려가면 준동형 $u_i:M_i\to N_i$를 준다. 이로부터 표준
준동형
\[
  \operatorname{Hom}_{A_j}(M_j,N_j)
  \longrightarrow
  \operatorname{Hom}_{A_i}(M_i,N_i)
\]
이 정의된다. 이 준동형들에 관하여
$\operatorname{Hom}_{A_i}(M_i,N_i)$들은 \emph{사영계}를 이루며,
(7.2.10)에 따라 표준 동형
\[
  \varphi:\operatorname{Hom}_A(M,N)
  \longrightarrow
  \varprojlim_i\operatorname{Hom}_{A_i}(M_i,N_i)
\]
이 존재한다. 더 나아가 다음이 성립한다.
\end{env}

\begin{proposition}[7.8.2]
\label{0.7.8.2-ko}
$M,N$이 뇌터 $\mathfrak J$-아딕환 $A$ 위의 유한형 가군이면,
부분가군
\[
  \operatorname{Hom}_A(M,\mathfrak J^{i+1}N)
\]
들은 $\operatorname{Hom}_A(M,N)$의 $\mathfrak J$-아딕 위상에서
$0$의 근방 기본계를 이루며, 표준 준동형
\[
  \varphi:\operatorname{Hom}_A(M,N)
  \longrightarrow
  \varprojlim_i\operatorname{Hom}_{A_i}(M_i,N_i)
\]
은 위상동형이다.
\end{proposition}

실제로 $M$을 유한형 자유 $A$-가군 $L$의 몫으로 볼 수 있고, 따라서
$\operatorname{Hom}_A(M,N)$을 $\operatorname{Hom}_A(L,N)$의
부분가군과 동일시할 수 있다. 이 동일시에서
$\operatorname{Hom}_A(M,\mathfrak J^{i+1}N)$은
$\operatorname{Hom}_A(M,N)$과
$\operatorname{Hom}_A(L,\mathfrak J^{i+1}N)$의 교집합이다.
$\operatorname{Hom}_A(L,N)$의 $\mathfrak J$-아딕 위상이
$\operatorname{Hom}_A(M,N)$에 유도하는 위상은
$\mathfrak J$-아딕 위상이므로(7.3.2), 첫째 주장은 $M=L=A^m$인 경우로
귀착된다.

\oldpage[0\textsubscript{I}]{78}

그런데 이 경우
\[
  \operatorname{Hom}_A(L,N)=N^m,
  \qquad
  \operatorname{Hom}_A(L,\mathfrak J^{i+1}N)
  =(\mathfrak J^{i+1}N)^m
  =\mathfrak J^{i+1}N^m,
\]
이므로 결과는 명백하다. 둘째 주장을 보이기 위하여
$\operatorname{Hom}_A(M,\mathfrak J^{i+1}N)$의
$\operatorname{Hom}_{A_j}(M_j,N_j)$에서의 상은 $j\leq i$일 때
$0$임을 주목하자. 따라서 $\varphi$는 연속이다. 역으로
$\operatorname{Hom}_{A_i}(M_i,N_i)$의 $0$의
$\operatorname{Hom}_A(M,N)$에서의 역상은
$\operatorname{Hom}_A(M,\mathfrak J^{i+1}N)$이므로
$\varphi$는 쌍연속이다.

$A$가 뇌터 $\mathfrak J$-준아딕환이고 $M,N$이
$\mathfrak J$-준아딕 위상에 관하여 분리인 두 유한형 $A$-가군이라고만
가정해도, 앞의 증명은 (7.8.2)의 첫째 주장이 그대로 성립하며
$\varphi$가 $\operatorname{Hom}_A(M,N)$에서
\[
  \varprojlim_i\operatorname{Hom}_{A_i}(M_i,N_i)
\]
의 한 부분가군 위로 가는 위상동형임을 보인다.

\begin{proposition}[7.8.3]
\label{0.7.8.3-ko}
(7.8.2)의 가정 아래에서 $M$에서 $N$으로 가는 단사 준동형들의 집합은
$\operatorname{Hom}_A(M,N)$의 열린 부분집합이다. 전사 준동형들과
전단사 준동형들의 집합도 마찬가지이다.
\end{proposition}

실제로 (7.3.5)와 (7.1.14)에 따라 $u$가 전사일 필요충분조건은 대응하는
준동형
\[
  u_0:M/\mathfrak JM\longrightarrow N/\mathfrak JN
\]
이 전사인 것이다. 따라서 $M$에서 $N$으로 가는 전사 준동형들의
집합은 연속 사상
\[
  \operatorname{Hom}_A(M,N)
  \longrightarrow
  \operatorname{Hom}_{A_0}(M_0,N_0)
\]
에 의한 이산공간의 한 부분집합의 역상이다.

이제 단사 준동형들의 집합이 열림을 보이자. 그러한 준동형 $v$를 잡고
$M'=v(M)$으로 놓자. 아르틴--리스 보조정리(7.3.2.1)에 따라 모든
$m>0$에 대하여
\[
  M'\cap\mathfrak J^{m+k}N\subset\mathfrak J^mM'
\]
를 만족하는 정수 $k\geq0$이 존재한다. 이제
$w\in\mathfrak J^{k+1}\operatorname{Hom}_A(M,N)$이면
$u=v+w$가 단사임을 보이겠다. 그러면 증명이 끝난다.

실제로 $u(x)=0$인 $x\in M$을 잡자. 모든 $i\geq0$에 대하여
$x\in\mathfrak J^iM$이면 $x\in\mathfrak J^{i+1}M$임을 보이겠다.
그러면
\[
  x\in\bigcap_{i\geq0}\mathfrak J^iM=(0)
\]
이 되어 충분하다. 그런데
$w(x)\in\mathfrak J^{i+k+1}N$이고 다른 한편
$w(x)=-v(x)\in M'$이다. 따라서
\[
  v(x)\in
  M'\cap\mathfrak J^{i+k+1}N
  \subset\mathfrak J^{i+1}M'.
\]
$v$는 $M$에서 $M'$ 위로 가는 동형이므로
$x\in\mathfrak J^{i+1}M$이다. 증명이 끝났다.

\begin{flushright}
\emph{(계속.)}
\end{flushright}
% ===== END EXACT INLINE INPUT: c0.tex =====


% ===== BEGIN EXACT INLINE INPUT: c1front.tex =====
% Authority: EGA I chapter1-frontmatter-fr.tex, SHA-256 7B2D0F8F812EBA3121202F0AE6415FFC6C281B8428DA8F0F72D89DF1CEC01708.
\clearpage
\oldpage[I]{79}
\phantomsection
\label{chapter:I-ko}

\begin{center}
{\large 제1장\par}
\vspace{1.5em}
{\Large\bfseries 스킴의 언어\par}
\vspace{1.8em}
{\bfseries 목차\par}
\end{center}

\medskip
\begin{tabular}{@{}r@{\ }l@{}}
\S~1.  & 아핀 스킴.\\
\S~2.  & 준스킴과 준스킴의 사상.\\
\S~3.  & 준스킴의 곱.\\
\S~4.  & 부분준스킴과 몰입 사상.\\
\S~5.  & 축소 준스킴; 분리 조건.\\
\S~6.  & 유한성 조건.\\
\S~7.  & 유리 사상.\\
\S~8.  & 슈발레 스킴.\\
\S~9.  & 준연접층에 관한 보충.\\
\S~10. & 형식 스킴.
\end{tabular}

\medskip
\S\S~1에서~8까지는 이후 전체에서 사용할 언어를 전개하는 것에서
크게 벗어나지 않는다. 다만 이 논고의 전반적인 정신에 따라
\S\S~7과~8은 다른 절들보다 덜 쓰이고, 쓰이더라도 덜 본질적이라는
점을 밝혀 둔다. 슈발레 스킴을 다룬 것도 슈발레~\cite{I-1}와
나가타~\cite{I-9}의 언어와 연결하기 위해서일 뿐이다. \S~9에서는
준연접층에 관한 정의와 결과들을 제시한다. 그 가운데에는 알려진
가환대수의 개념들을 ``기하학적'' 언어로 옮기는 데 그치지 않고 이미
전역적인 성격을 띠는 것도 있으며, 이러한 결과들은 바로 다음 장들부터
사상을 전역적으로 연구하는 데 반드시 필요하다. 끝으로 \S~10에서는
스킴 개념의 한 일반화를 도입한다. 이는 제III장에서 고유 사상의 코호몰로지
연구에 관한 기본 결과들을 편리하게 서술하고 증명하는 중간 도구로 쓰일
것이다. 한편 형식 스킴의 개념은 ``모듈라이 이론''의 몇몇 사실, 곧
대수다양체의 분류 문제를 표현하는 데 없어서는 안 될 것으로 보인다.
\S~10의 결과들은 제III장 \S~3 이전에는 사용하지 않으므로, 그때까지는
이 절을 건너뛰어 읽기를 권한다.

\clearpage
% ===== END EXACT INLINE INPUT: c1front.tex =====


% ===== BEGIN EXACT INLINE INPUT: c1s1.tex =====
% Authority: EGA I ega1-1-fr.tex, whole-file SHA-256 25840D1D77C5284A158CDC78D622A458221DA3FE9ECAECC5EA9382A5C4EC595B.
% Sequential coverage in this file: authority lines 1-1624, complete subsections 1.1-1.7.
\setcounter{section}{0}
\section{아핀 스킴}
\label{section:I.1-ko}

\setcounter{subsection}{0}
\subsection{환의 소 스펙트럼}
\label{subsection:I.1.1-ko}

\oldpage[I]{80}
\begin{env}[1.1.1]
\label{I.1.1.1-ko}
\emph{표기}: $A$를 (가환)환, $M$을 $A$-가군이라 하자. 이 장과
이어지는 장들에서는 다음 표기들을 계속 사용한다.

\smallskip
\noindent
$\operatorname{Spec}(A)={}$ $A$의 \emph{소 아이디얼들의 집합}이며,
$A$의 \emph{소 스펙트럼}이라고도 한다. 흔히
$x\in X=\operatorname{Spec}(A)$ 대신 $\mathfrak{j}_x$라고 쓰는 것이
편리하다. $\operatorname{Spec}(A)$가 \emph{공집합}일 필요충분조건은
$A$가 영환인 것이다.

\smallskip
\noindent
$A_x=A_{\mathfrak{j}_x}={}$ \emph{$S^{-1}A$라는 (국소) 분수환},
여기서 $S=A-\mathfrak{j}_x$이다.

\smallskip
\noindent
$\mathfrak{m}_x=\mathfrak{j}_xA_{\mathfrak{j}_x}={}$ $A_x$의
\emph{극대 아이디얼}.

\smallskip
\noindent
$k(x)=A_x/\mathfrak{m}_x={}$ $A_x$의 \emph{잉여체}. 이는 정역
$A/\mathfrak{j}_x$의 분수체와 표준적으로 동형이며, 그 분수체와
동일시한다.

\smallskip
\noindent
$f(x)={}$ $A/\mathfrak{j}_x\subset k(x)$ 안에서 취한 \emph{$f$의
$\mathfrak{j}_x$ 법 동치류} ($f\in A$, $x\in X$). $f(x)$를
$x\in\operatorname{Spec}(A)$에서의 $f$의 \emph{값}이라고도 한다.
관계 $f(x)=0$과 $f\in\mathfrak{j}_x$는 \emph{동치}이다.

\smallskip
\noindent
$M_x=M\otimes_AA_x={}$ $A-\mathfrak{j}_x$에 분모를 갖는
\emph{분수가군}.

\smallskip
\noindent
$\mathfrak{r}(E)={}$ $A$의 부분집합 $E$가 생성하는 아이디얼의
\emph{근기}.

\smallskip
\noindent
$V(E)={}$ \emph{$E\subset\mathfrak{j}_x$를 만족하는 $x\in X$들의 집합}
(여기서 $E\subset A$)
(또는 같은 말로, 모든 $f\in E$에 대하여 $f(x)=0$인 $x\in X$들의
집합). 따라서
\begin{equation}
  \mathfrak{r}(E)=\bigcap_{x\in V(E)}\mathfrak{j}_x.
  \tag{1.1.1.1}\label{I.1.1.1.1-ko}
\end{equation}

\smallskip
\noindent
$V(f)=V(\{f\})$ ($f\in A$).

\smallskip
\noindent
$D(f)=X-V(f)={}$ \emph{$f(x)\neq0$인 $x\in X$들의 집합}.
\end{env}

\begin{proposition}[1.1.2]
\label{I.1.1.2-ko}
다음 성질들이 성립한다.
\begin{enumerate}
  \item[\rm(i)] $V(0)=X$, $V(1)=\varnothing$.
  \item[\rm(ii)] $E\subset E'$이면 $V(E)\supset V(E')$이다.
  \item[\rm(iii)] $A$의 부분집합들의 임의의 족 $(E_\lambda)$에 대하여
  \[
    V\left(\bigcup_\lambda E_\lambda\right)
      =V\left(\sum_\lambda E_\lambda\right)
      =\bigcap_\lambda V(E_\lambda).
  \]
  \item[\rm(iv)] $V(EE')=V(E)\cup V(E')$.
  \item[\rm(v)] $V(E)=V(\mathfrak{r}(E))$.
\end{enumerate}
\end{proposition}

성질 (i), (ii), (iii)은 자명하고, (v)는 (ii)와 공식
(\hyperref[I.1.1.1.1-ko]{1.1.1.1})에서 따른다.
$V(EE')\supset V(E)\cup V(E')$임은 자명하다. 역으로
$x\notin V(E)$이고 $x\notin V(E')$이면 $k(x)$ 안에서
$f(x)\neq0$, $f'(x)\neq0$을 만족하는 $f\in E$, $f'\in E'$가
존재한다. 따라서 $f(x)f'(x)\neq0$, 즉 $x\notin V(EE')$이고,
이로써 (iv)가 증명된다.

명제 (\hyperref[I.1.1.2-ko]{1.1.2})에서 특히 $E$가 $A$의 모든
부분집합을 달릴 때 $V(E)$ 꼴의 집합들이 $X$ 위의 한 위상의
\emph{닫힌집합}들이 됨을 알 수 있다. 이 위상을
\emph{스펙트럼 위상}\footnote{대수기하학에 이 위상을 도입한 사람은
자리스키이다. 따라서 흔히 $X$의 ``자리스키 위상''이라고도 한다.}이라
하겠다. 명시적으로 달리 말하지 않는 한
$X=\operatorname{Spec}(A)$에는 언제나 스펙트럼 위상을 준다.

\oldpage[I]{81}
\begin{env}[1.1.3]
\label{I.1.1.3-ko}
$X$의 임의의 부분집합 $Y$에 대하여 $\mathfrak{j}(Y)$는 모든
$y\in Y$에서 $f(y)=0$인 $f\in A$들의 집합을 뜻한다. 같은 말로
$\mathfrak{j}(Y)$는 $y\in Y$에 대한 소 아이디얼
$\mathfrak{j}_y$들의 교집합이다. $Y\subset Y'$이면
$\mathfrak{j}(Y)\supset\mathfrak{j}(Y')$임은 자명하며, $X$의
부분집합들의 임의의 족 $(Y_\lambda)$에 대하여
\begin{equation}
  \mathfrak{j}\left(\bigcup_\lambda Y_\lambda\right)
    =\bigcap_\lambda\mathfrak{j}(Y_\lambda)
  \tag{1.1.3.1}\label{I.1.1.3.1-ko}
\end{equation}
이다. 끝으로
\begin{equation}
  \mathfrak{j}(\{x\})=\mathfrak{j}_x
  \tag{1.1.3.2}\label{I.1.1.3.2-ko}
\end{equation}
이다.
\end{env}

\begin{proposition}[1.1.4]
\label{I.1.1.4-ko}
\begin{enumerate}
  \item[\rm(i)] $A$의 모든 부분집합 $E$에 대하여
  $\mathfrak{j}(V(E))=\mathfrak{r}(E)$이다.
  \item[\rm(ii)] $X$의 모든 부분집합 $Y$에 대하여
  $V(\mathfrak{j}(Y))=\overline{Y}$이며, 여기서
  $\overline{Y}$는 $X$에서의 $Y$의 폐포이다.
\end{enumerate}
\end{proposition}

(i)는 정의와 (\hyperref[I.1.1.1.1-ko]{1.1.1.1})에서 바로 따른다.
한편 $V(\mathfrak{j}(Y))$는 닫혀 있고 $Y$를 포함한다. 역으로
$Y\subset V(E)$이면 모든 $f\in E$, $y\in Y$에 대하여 $f(y)=0$이므로
$E\subset\mathfrak{j}(Y)$이고
$V(E)\supset V(\mathfrak{j}(Y))$이다. 이로써 (ii)가 증명된다.

\begin{corollary}[1.1.5]
\label{I.1.1.5-ko}
$X=\operatorname{Spec}(A)$의 닫힌 부분집합들과 자기 근기와 같은
$A$의 아이디얼들(다시 말해 소 아이디얼들의 교집합들)은 포함관계를
뒤집는 사상
$Y\mapsto\mathfrak{j}(Y)$, $\mathfrak{a}\mapsto V(\mathfrak{a})$에
의하여 서로 전단사 대응한다. 두 닫힌 부분집합의 합집합
$Y_1\cup Y_2$에는 $\mathfrak{j}(Y_1)\cap\mathfrak{j}(Y_2)$가
대응하고, 닫힌 부분집합들의 임의의 족 $(Y_\lambda)$의 교집합에는
$\mathfrak{j}(Y_\lambda)$들의 합의 근기가 대응한다.
\end{corollary}

\begin{corollary}[1.1.6]
\label{I.1.1.6-ko}
$A$가 뇌터 환이면 $X=\operatorname{Spec}(A)$는 뇌터 공간이다.
\end{corollary}

이 따름정리의 역은 거짓이다. 예를 들어 영 아이디얼이 아닌 소
아이디얼을 단 하나만 갖는 비뇌터 정역, 이를테면 랭크 1인 비이산
값매김환이 반례가 된다.

스펙트럼이 뇌터 공간이 아닌 환 $A$의 예로는 무한 콤팩트 공간
$Y$ 위의 실수값 연속함수들의 환 $\mathcal{C}(Y)$를 들 수 있다.
집합으로서 $Y$는 $A$의 극대 아이디얼들의 집합과 동일시된다는 것이
알려져 있고, $X=\operatorname{Spec}(A)$의 위상이 $Y$ 위에 유도하는
위상은 $Y$의 본래 위상임을 쉽게 알 수 있다. $Y$가 뇌터 공간이
아니므로 $X$도 뇌터 공간이 아니다.

\begin{corollary}[1.1.7]
\label{I.1.1.7-ko}
모든 $x\in X$에 대하여 $\{x\}$의 폐포는
$\mathfrak{j}_x\subset\mathfrak{j}_y$를 만족하는 $y\in X$들의
집합이다. $\{x\}$가 닫혀 있을 필요충분조건은
$\mathfrak{j}_x$가 극대인 것이다.
\end{corollary}

\begin{corollary}[1.1.8]
\label{I.1.1.8-ko}
공간 $X=\operatorname{Spec}(A)$는 콜모고로프 공간이다.
\end{corollary}

실제로 $x$, $y$가 $X$의 서로 다른 두 점이면
$\mathfrak{j}_x\not\subset\mathfrak{j}_y$이거나
$\mathfrak{j}_y\not\subset\mathfrak{j}_x$이다. 따라서 두 점
$x$, $y$ 가운데 하나는 다른 하나의 폐포에 속하지 않는다.

\begin{env}[1.1.9]
\label{I.1.1.9-ko}
명제 (\hyperref[I.1.1.2-ko]{1.1.2}, (iv))에 따라 $A$의 두 원소
$f$, $g$에 대하여
\begin{equation}
  D(fg)=D(f)\cap D(g)
  \tag{1.1.9.1}\label{I.1.1.9.1-ko}
\end{equation}
이다. 또한 명제 (\hyperref[I.1.1.4-ko]{1.1.4}, (i))와
명제 (\hyperref[I.1.1.2-ko]{1.1.2}, (v))에 따르면
$D(f)=D(g)$라는 관계는
$\mathfrak{r}(f)=\mathfrak{r}(g)$라는 뜻이며, 같은 말로
$(f)$와 $(g)$를 포함하는 극소 소 아이디얼들이 같다는 뜻이다.
특히 $u$가 가역원이고 $f=ug$이면 그러하다.
\end{env}

\oldpage[I]{82}
\begin{proposition}[1.1.10]
\label{I.1.1.10-ko}
\begin{enumerate}
  \item[\rm(i)] $f$가 $A$를 달릴 때 집합 $D(f)$들은 $X$의 위상의
  한 기저를 이룬다.
  \item[\rm(ii)] 모든 $f\in A$에 대하여 $D(f)$는 준콤팩트이다.
  특히 $X=D(1)$은 준콤팩트이다.
\end{enumerate}
\end{proposition}

(i) $U$를 $X$의 열린집합이라 하자. 정의에 따라 $A$의 한 부분집합
$E$에 대하여 $U=X-V(E)$이고,
$V(E)=\bigcap_{f\in E}V(f)$이다. 따라서
$U=\bigcup_{f\in E}D(f)$이다.

(ii) (i)에 따라 다음을 보이면 충분하다. $A$의 원소들의 족
$(f_\lambda)_{\lambda\in L}$이
$D(f)\subset\bigcup_{\lambda\in L}D(f_\lambda)$를 만족하면,
$D(f)\subset\bigcup_{\lambda\in J}D(f_\lambda)$가 되는 $L$의
유한 부분집합 $J$가 존재한다. $\mathfrak{a}$를 $f_\lambda$들이
생성하는 $A$의 아이디얼이라 하자. 가정에 따라
$V(f)\supset V(\mathfrak{a})$이므로
$\mathfrak{r}(f)\subset\mathfrak{r}(\mathfrak{a})$이다.
$f\in\mathfrak{r}(f)$이므로 어떤 정수 $n\geq0$에 대하여
$f^n\in\mathfrak{a}$이다. 그러면 $f^n$은 한 유한 부분족
$(f_\lambda)_{\lambda\in J}$이 생성하는 아이디얼
$\mathfrak{b}$에 속하며,
$V(f)=V(f^n)\supset V(\mathfrak{b})
=\bigcap_{\lambda\in J}V(f_\lambda)$이다. 다시 말해
$D(f)\subset\bigcup_{\lambda\in J}D(f_\lambda)$이다.

\begin{proposition}[1.1.11]
\label{I.1.1.11-ko}
$A$의 모든 아이디얼 $\mathfrak{a}$에 대하여
$\operatorname{Spec}(A/\mathfrak{a})$는
$\operatorname{Spec}(A)$의 닫힌 부분공간 $V(\mathfrak{a})$와
표준적으로 동일시된다.
\end{proposition}

실제로 $A/\mathfrak{a}$의 아이디얼들(각각 소 아이디얼들)과
$\mathfrak{a}$를 포함하는 $A$의 아이디얼들(각각 소 아이디얼들)
사이에는 포함관계의 순서구조를 보존하는 표준 전단사 대응이 있다.

$A$의 멱영원들의 집합 $\mathfrak{N}$, 곧 $A$의
\emph{멱영근기}는 $\mathfrak{r}(0)$과 같은 아이디얼이고
$A$의 모든 소 아이디얼의 교집합임을 상기하자
(\hyperref[0.1.1.1-ko]{0, 1.1.1}).

\begin{corollary}[1.1.12]
\label{I.1.1.12-ko}
위상공간 $\operatorname{Spec}(A)$와
$\operatorname{Spec}(A/\mathfrak{N})$은 표준적으로 동형이다.
\end{corollary}

\begin{proposition}[1.1.13]
\label{I.1.1.13-ko}
$X=\operatorname{Spec}(A)$가 기약일
(\hyperref[0.2.1.1-ko]{0, 2.1.1}) 필요충분조건은
$A/\mathfrak{N}$이 정역인 것이다. 이는 $\mathfrak{N}$이
소 아이디얼인 것과도 동치이다.
\end{proposition}

따름정리 (\hyperref[I.1.1.12-ko]{1.1.12})에 따라
$\mathfrak{N}=0$인 경우만 생각하면 된다. $X$가 기약이 아니면
$X$와 다른 두 닫힌 부분집합 $Y_1$, $Y_2$가 존재하여
$X=Y_1\cup Y_2$이고, 따라서
$\mathfrak{j}(X)=\mathfrak{j}(Y_1)\cap\mathfrak{j}(Y_2)=0$이다.
아이디얼 $\mathfrak{j}(Y_1)$과 $\mathfrak{j}(Y_2)$는 모두
$(0)$과 다르므로(\hyperref[I.1.1.5-ko]{1.1.5}) $A$는 정역이
아니다. 역으로 $A$에 $f\neq0$, $g\neq0$, $fg=0$인 두 원소가
있으면, $A$의 모든 소 아이디얼의 교집합이 $(0)$이므로
$V(f)\neq X$, $V(g)\neq X$이고
$X=V(fg)=V(f)\cup V(g)$이다.

\begin{corollary}[1.1.14]
\label{I.1.1.14-ko}
\begin{enumerate}
  \item[\rm(i)] $X=\operatorname{Spec}(A)$의 닫힌 부분집합들과 자기
  근기와 같은 $A$의 아이디얼들 사이의 전단사 대응에서, $X$의 닫힌
  기약 부분집합들은 $A$의 소 아이디얼들과 대응한다. 특히 $X$의
  기약 성분들은 $A$의 극소 소 아이디얼들과 대응한다.
  \item[\rm(ii)] 사상 $x\mapsto\overline{\{x\}}$는 $X$와 $X$의
  닫힌 기약 부분집합들의 집합 사이에 전단사 대응을 정한다.
  다시 말해 $X$의 모든 닫힌 기약 부분집합은 일반점을 하나만 갖는다.
\end{enumerate}
\end{corollary}

(i)는 (\hyperref[I.1.1.13-ko]{1.1.13})과
(\hyperref[I.1.1.11-ko]{1.1.11})에서 바로 따른다. (ii)를
증명할 때에도 (\hyperref[I.1.1.11-ko]{1.1.11})에 따라 $X$가
기약인 경우만 생각하면 된다. 그러면 명제
(\hyperref[I.1.1.13-ko]{1.1.13})에 따라 $A$에는 가장 작은 소
아이디얼 $\mathfrak{N}$이 존재하고, 이는 $X$의 한 일반점에
대응한다.
\oldpage[I]{83}
더욱이 $X$는 콜모고로프 공간이므로
((\hyperref[I.1.1.8-ko]{1.1.8})과
(\hyperref[0.2.1.3-ko]{0, 2.1.3})) 일반점을 하나만 갖는다.

\begin{proposition}[1.1.15]
\label{I.1.1.15-ko}
$\mathfrak{J}$가 $A$의 근기 $\mathfrak{R}(A)$에 포함되는
아이디얼이면, $X=\operatorname{Spec}(A)$에서
$V(\mathfrak{J})$의 유일한 근방은 $X$ 전체이다.
\end{proposition}

실제로 $A$의 모든 극대 아이디얼은 정의에 따라
$V(\mathfrak{J})$에 속한다. $A$의 모든 아이디얼
$\mathfrak{a}$는 어떤 극대 아이디얼에 포함되므로
$V(\mathfrak{a})\cap V(\mathfrak{J})\neq\varnothing$이고,
이로부터 명제가 따른다.

\subsection{환의 소 스펙트럼의 함자적 성질}
\label{subsection:I.1.2-ko}

\begin{env}[1.2.1]
\label{I.1.2.1-ko}
$A$, $A'$을 두 환이라 하고
\[
  \varphi:A'\longrightarrow A
\]
를 환 준동형이라 하자. 모든 소 아이디얼
$x=\mathfrak{j}_x\in\operatorname{Spec}(A)=X$에 대하여 환
$A'/\varphi^{-1}(\mathfrak{j}_x)$는
$A/\mathfrak{j}_x$의 한 부분환과 표준적으로 동형이므로 정역이다.
다시 말해 $\varphi^{-1}(\mathfrak{j}_x)$는 $A'$의 소 아이디얼이다.
이를 ${}^a\varphi(x)$로 나타내면 사상
\[
  {}^a\varphi:X=\operatorname{Spec}(A)\longrightarrow
  X'=\operatorname{Spec}(A')
\]
가 정의된다. 이 사상은 $\operatorname{Spec}(\varphi)$라고도 쓰며,
준동형 $\varphi$에 대응하는 사상이라고 한다. 몫으로 내려가서
$\varphi$로부터 얻는
$A'/\varphi^{-1}(\mathfrak{j}_x)$에서
$A/\mathfrak{j}_x$로 가는 단사 준동형과, 이를 체의 단사 준동형으로
표준 연장한 것을 모두 $\varphi^x$로 나타낸다. 따라서
\[
  \varphi^x:k({}^a\varphi(x))\longrightarrow k(x).
\]
정의에 따라 모든 $f'\in A'$에 대하여
\begin{equation}
  \varphi^x\bigl(f'({}^a\varphi(x))\bigr)=\bigl(\varphi(f')\bigr)(x)
  \qquad(x\in X).
  \tag{1.2.1.1}\label{I.1.2.1.1-ko}
\end{equation}
\end{env}

\begin{proposition}[1.2.2]
\label{I.1.2.2-ko}
\begin{enumerate}
  \item[\rm(i)] $A'$의 임의의 부분집합 $E'$에 대하여
  \begin{equation}
    {}^a\varphi^{-1}(V(E'))=V(\varphi(E'))
    \tag{1.2.2.1}\label{I.1.2.2.1-ko}
  \end{equation}
  이고, 특히 모든 $f'\in A'$에 대하여
  \begin{equation}
    {}^a\varphi^{-1}(D(f'))=D(\varphi(f')).
    \tag{1.2.2.2}\label{I.1.2.2.2-ko}
  \end{equation}
  \item[\rm(ii)] $A$의 모든 아이디얼 $\mathfrak{a}$에 대하여
  \begin{equation}
    \overline{{}^a\varphi(V(\mathfrak{a}))}
    =V(\varphi^{-1}(\mathfrak{a})).
    \tag{1.2.2.3}\label{I.1.2.2.3-ko}
  \end{equation}
\end{enumerate}
\end{proposition}

실제로 관계 ${}^a\varphi(x)\in V(E')$는 정의에 따라
$E'\subset\varphi^{-1}(\mathfrak{j}_x)$와 동치이고, 이는 다시
$\varphi(E')\subset\mathfrak{j}_x$, 끝으로
$x\in V(\varphi(E'))$와 동치이다. 이로써 (i)가 증명된다. (ii)를
증명할 때에는 $V(\mathfrak{r}(\mathfrak{a}))=V(\mathfrak{a})$
((\hyperref[I.1.1.2-ko]{1.1.2 (v)}))이고
\[
  \varphi^{-1}(\mathfrak{r}(\mathfrak{a}))
  =\mathfrak{r}(\varphi^{-1}(\mathfrak{a}))
\]
이므로 $\mathfrak{a}$가 자기 근기와 같다고 가정해도 된다.
$Y=V(\mathfrak{a})$와
$\mathfrak{a}'=\mathfrak{j}({}^a\varphi(Y))$로 놓으면 명제
(\hyperref[I.1.1.4-ko]{1.1.4 (ii)})에 따라
$\overline{{}^a\varphi(Y)}=V(\mathfrak{a}')$이다. 관계
$f'\in\mathfrak{a}'$은 정의에 따라 모든
$x'\in{}^a\varphi(Y)$에 대하여 $f'(x')=0$이라는 것과 동치이다.
공식 (\hyperref[I.1.2.1.1-ko]{1.2.1.1})에 따라 이는 모든
$x\in Y$에 대하여 $\varphi(f')(x)=0$이라는 것과도 동치이고,
$\mathfrak{a}$가 자기 근기와 같으므로 다시
$\varphi(f')\in\mathfrak{j}(Y)=\mathfrak{a}$와 동치이다. 이로써
(ii)가 증명된다.

\begin{corollary}[1.2.3]
\label{I.1.2.3-ko}
사상 ${}^a\varphi$는 연속이다.
\end{corollary}

$A''$이 세 번째 환이고 $\varphi':A''\to A'$이 환 준동형이면
\[
  {}^a(\varphi\circ\varphi')={}^a\varphi'\circ{}^a\varphi.
\]
이 결과와 따름정리 (\hyperref[I.1.2.3-ko]{1.2.3})은
$\operatorname{Spec}(A)$가 $A$에 관하여 환의 범주에서 위상공간의
범주로 가는 반변 함자임을 뜻한다.

\begin{corollary}[1.2.4]
\label{I.1.2.4-ko}
\oldpage[I]{84}
$\varphi$가 다음 성질을 갖는다고 하자. 모든 $f\in A$를
\[
  f=h\varphi(f')
\]
꼴로 쓸 수 있고, 여기서 $h$는 $A$의 가역원이다. 특히
$\varphi$가 전사이면 이 성질이 성립한다. 그러면 ${}^a\varphi$는
$X$에서 ${}^a\varphi(X)$로 가는 위상동형이다.
\end{corollary}

임의의 부분집합 $E\subset A$에 대하여
$V(E)=V(\varphi(E'))$가 되는 $A'$의 부분집합 $E'$이 존재함을
보이자. 공리 $(T_0)$
((\hyperref[I.1.1.8-ko]{1.1.8}))과 공식
(\hyperref[I.1.2.2.1-ko]{1.2.2.1})에 따라 이 사실은 먼저
${}^a\varphi$가 단사임을 보이고, 다시 같은 공식에 따라
${}^a\varphi$가 위상동형임을 보인다. 실제로 각 $f\in E$마다
$h\varphi(f')=f$이고 $h$가 $A$의 가역원이 되도록 $f'\in A'$을
하나 택하면 충분하다. 이렇게 얻은 원소 $f'$들의 집합 $E'$이 원하는
조건을 만족한다.

\begin{env}[1.2.5]
\label{I.1.2.5-ko}
특히 $\varphi$가 $A$에서 몫환 $A/\mathfrak{a}$로 가는 표준
준동형이면 (\hyperref[I.1.1.12-ko]{1.1.12})를 다시 얻는다. 이때
${}^a\varphi$는 $\operatorname{Spec}(A/\mathfrak{a})$와 동일시한
$V(\mathfrak{a})$에서 $X=\operatorname{Spec}(A)$로 가는 표준
단사 사상이다.
\end{env}

(\hyperref[I.1.2.4-ko]{1.2.4})의 또 다른 특수한 경우로 다음을
얻는다.

\begin{corollary}[1.2.6]
\label{I.1.2.6-ko}
$S$가 $A$의 곱셈적 부분집합이면 스펙트럼
$\operatorname{Spec}(S^{-1}A)$는, 위상까지 포함하여,
$X=\operatorname{Spec}(A)$에서
$\mathfrak{j}_x\cap S=\varnothing$을 만족하는 점 $x$들로 이루어진
부분공간과 표준적으로 동일시된다.
\end{corollary}

실제로 $S^{-1}A$의 소 아이디얼들이
$\mathfrak{j}_x\cap S=\varnothing$을 만족하는
$S^{-1}\mathfrak{j}_x$들이고
\[
  \mathfrak{j}_x=({}^S i_A)^{-1}(S^{-1}\mathfrak{j}_x)
\]
임은 (\hyperref[0.1.2.6-ko]{0, 1.2.6})에서 이미 알고 있다. 따라서
${}^S i_A$에 따름정리 (\hyperref[I.1.2.4-ko]{1.2.4})를 적용하면
된다.

\begin{corollary}[1.2.7]
\label{I.1.2.7-ko}
${}^a\varphi(X)$가 $X'$에서 모든 곳에서 조밀하기 위한 필요충분조건은
$\operatorname{Ker}\varphi$의 모든 원소가 멱영원인 것이다.
\end{corollary}

실제로 공식 (\hyperref[I.1.2.2.3-ko]{1.2.2.3})을 아이디얼
$\mathfrak{a}=(0)$에 적용하면
\[
  \overline{{}^a\varphi(X)}=V(\operatorname{Ker}\varphi)
\]
를 얻는다. $V(\operatorname{Ker}\varphi)=X'$이기 위한
필요충분조건은 $\operatorname{Ker}\varphi$가 $A'$의 모든 소
아이디얼, 곧 $A'$의 멱영근기 $\mathfrak{N}$에 포함되는 것이다.

\subsection{가군에 대응하는 층}
\label{subsection:I.1.3-ko}
\label{I.1.3-ko}

\begin{env}[1.3.1]
\label{I.1.3.1-ko}
$A$를 가환환, $M$을 $A$-가군, $f$를 $A$의 원소라 하고,
$S_f$를 $n\geq0$일 때의 $f^n$들로 이루어진 곱셈적 집합이라 하자.
$A_f=S_f^{-1}A$, $M_f=S_f^{-1}M$으로 놓는다는 것을 상기하자.
$S'_f$가 $S_f$의 한 원소를 나누는 $g\in A$들로 이루어진 $A$의
포화 곱셈적 부분집합이면, $A_f$와 $M_f$는 각각
${S'_f}^{-1}A$와 ${S'_f}^{-1}M$에 표준적으로 동일시된다
(\hyperref[0.1.4.3-ko]{0, 1.4.3}).
\end{env}

\begin{lemma}[1.3.2]
\label{I.1.3.2-ko}
다음 조건들은 서로 동치이다.
\begin{enumerate}
  \item[\rm(a)] $g\in S'_f$;
  \item[\rm(b)] $S'_g\subset S'_f$;
  \item[\rm(c)] $f\in\mathfrak{r}(g)$;
  \item[\rm(d)] $\mathfrak{r}(f)\subset\mathfrak{r}(g)$;
  \item[\rm(e)] $V(g)\subset V(f)$;
  \item[\rm(f)] $D(f)\subset D(g)$.
\end{enumerate}
\end{lemma}

이는 정의와 (\hyperref[I.1.1.5-ko]{1.1.5})에서 바로 따른다.

\begin{env}[1.3.3]
\label{I.1.3.3-ko}
$D(f)=D(g)$이면 보조정리 (\hyperref[I.1.3.2-ko]{1.3.2, b)})에
따라 $M_f=M_g$이다. 더 일반적으로 $D(f)\supset D(g)$, 따라서
$S'_f\subset S'_g$이면 함자적인 표준 준동형
\[
  \rho_{g,f}:M_f\longrightarrow M_g
\]
가 존재함을 알고 있다(\hyperref[0.1.4.1-ko]{0, 1.4.1}). 또한
$D(f)\supset D(g)\supset D(h)$이면
(\hyperref[0.1.4.4-ko]{0, 1.4.4})
\begin{equation}
  \rho_{h,g}\circ\rho_{g,f}=\rho_{h,f}.
  \tag{1.3.3.1}\label{I.1.3.3.1-ko}
\end{equation}
\end{env}

\oldpage[I]{85}
주어진 $x\in X=\operatorname{Spec}(A)$에 대하여 $f$가
$A-\mathfrak{j}_x$를 움직일 때, $S'_f$들은
$A-\mathfrak{j}_x$의 부분집합들로 이루어진 증가 여과족을 이룬다.
실제로 $A-\mathfrak{j}_x$의 두 원소 $f$, $g$에 대하여 $S'_f$와
$S'_g$는 $S'_{fg}$에 포함되고, $f\in A-\mathfrak{j}_x$일 때의
$S'_f$들의 합집합은 $A-\mathfrak{j}_x$이다. 따라서
(\hyperref[0.1.4.5-ko]{0, 1.4.5}) $A_x$-가군 $M_x$는 준동형족
$(\rho_{g,f})$에 관한 귀납극한 $\varinjlim M_f$와 표준적으로
동일시된다. $f\in A-\mathfrak{j}_x$, 또는 같은 말로
$x\in D(f)$일 때의 표준 준동형을
\[
  \rho_x^f:M_f\longrightarrow M_x
\]
로 나타내겠다.

\begin{definition}[1.3.4]
\label{I.1.3.4-ko}
$D(f)$($f\in A$)들로 이루어진 $X$의 기저 $\mathfrak{B}$ 위의
준층 $D(f)\mapsto A_f$(각각 $D(f)\mapsto M_f$)에 대응하는 환의
층(각각 $\widetilde{A}$-가군)을 소 스펙트럼
$X=\operatorname{Spec}(A)$의 \emph{구조층}(각각 \emph{$A$-가군
$M$에 대응하는 층})이라 하고, $\widetilde{A}$ 또는
$\mathscr{O}_X$(각각 $\widetilde{M}$)로 나타낸다
((\hyperref[I.1.1.10-ko]{1.1.10}),
(\hyperref[0.3.2.1-ko]{0, 3.2.1} 및
\hyperref[0.3.5.6-ko]{3.5.6})).
\end{definition}

줄기 $\widetilde{A}_x$(각각 $\widetilde{M}_x$)가 환 $A_x$(각각
$A_x$-가군 $M_x$)와 동일시됨은 이미 보았다
(\hyperref[0.3.2.4-ko]{0, 3.2.4}). 표준 사상을
\[
  \theta_f:A_f\longrightarrow\Gamma(D(f),\widetilde{A})
\]
\[
  \text{(각각 }\theta_f:M_f\longrightarrow
  \Gamma(D(f),\widetilde{M})\text{)}
\]
로 나타내겠다. 그러면 모든 $x\in D(f)$와 모든 $\xi\in M_f$에
대하여
\begin{equation}
  (\theta_f(\xi))_x=\rho_x^f(\xi).
  \tag{1.3.4.1}\label{I.1.3.4.1-ko}
\end{equation}

\begin{proposition}[1.3.5]
\label{I.1.3.5-ko}
$\widetilde{M}$은 $M$에 관한 완전 공변 함자로서, $A$-가군의
범주에서 $\widetilde{A}$-가군의 범주로 간다.
\end{proposition}

실제로 $M$, $N$을 두 $A$-가군, $u:M\to N$을 준동형이라 하자.
모든 $f\in A$에 대하여 $u$에는 $A_f$-가군 $M_f$에서
$A_f$-가군 $N_f$로 가는 준동형 $u_f$가 표준적으로 대응하고,
$D(g)\subset D(f)$일 때 그림
\[
  \xymatrix{
    M_f\ar[r]^{u_f}\ar[d]_{\rho_{g,f}} & N_f\ar[d]^{\rho_{g,f}}\\
    M_g\ar[r]_{u_g} & N_g
  }
\]
은 가환한다(\hyperref[0.1.4.1-ko]{0, 1.4.1}). 따라서 이 준동형들은
$\widetilde{A}$-가군 준동형
$\widetilde{u}:\widetilde{M}\to\widetilde{N}$을 정의한다
(\hyperref[0.3.2.3-ko]{0, 3.2.3}). 더욱이 모든 $x\in X$에 대하여
$\widetilde{u}_x$는 $x\in D(f)$($f\in A$)일 때의 $u_f$들의
귀납극한이다. 따라서 $\widetilde{M}_x$와 $\widetilde{N}_x$를 각각
$M_x$와 $N_x$에 표준적으로 동일시하면
(\hyperref[0.1.4.5-ko]{0, 1.4.5}), $\widetilde{u}_x$는 $u$에서
표준적으로 얻는 준동형 $u_x$와 동일시된다. 이제 $P$를 세 번째
$A$-가군, $v:N\to P$를 준동형이라 하고 $w=v\circ u$로 놓으면
$w_x=v_x\circ u_x$이고, 따라서
$\widetilde{w}=\widetilde{v}\circ\widetilde{u}$임은 자명하다. 이로써
$M\mapsto\widetilde{M}$을 $A$-가군의 범주에서
$\widetilde{A}$-가군의 범주로 가는 \emph{공변 함자}로 정의하였다.
이 함자는 \emph{완전}하다. 실제로 모든 $x\in X$에 대하여 $M_x$는
$M$에 관한 완전 함자이기 때문이다
(\hyperref[0.1.3.2-ko]{0, 1.3.2}). 또한 두 지지집합의 정의에 따라
(\hyperref[0.1.7.1-ko]{0, 1.7.1} 및
\hyperref[0.3.1.6-ko]{3.1.6})
$\operatorname{Supp}(M)=\operatorname{Supp}(\widetilde{M})$이다.

\oldpage[I]{86}
\begin{proposition}[1.3.6]
\label{I.1.3.6-ko}
모든 $f\in A$에 대하여 열린집합 $D(f)\subset X$는 소 스펙트럼
$\operatorname{Spec}(A_f)$와 표준적으로 동일시되고, $A_f$-가군
$M_f$에 대응하는 층 $\widetilde{M_f}$는 제한
$\widetilde{M}|D(f)$와 표준적으로 동일시된다.
\end{proposition}

첫째 명제는 (\hyperref[I.1.2.6-ko]{1.2.6})의 특수한 경우이다.
더욱이 $D(g)\subset D(f)$인 $g\in A$에 대하여 $M_g$는 분모가
$A_f$에서 $g$의 표준 상의 거듭제곱인 $M_f$의 분수가군과
표준적으로 동일시된다(\hyperref[0.1.4.6-ko]{0, 1.4.6}). 이제
$\widetilde{M_f}$와 $\widetilde{M}|D(f)$의 표준적인 동일시는
정의에서 따른다.

\begin{theorem}[1.3.7]
\label{I.1.3.7-ko}
모든 $A$-가군 $M$과 모든 $f\in A$에 대하여 준동형
\[
  \theta_f:M_f\longrightarrow\Gamma(D(f),\widetilde{M})
\]
는 전단사이다. 다시 말해 준층 $D(f)\mapsto M_f$는 \emph{층}이다.
특히 $M$은 $\theta_1$에 의해 $\Gamma(X,\widetilde{M})$와
동일시된다.
\end{theorem}

$M=A$이면 $\theta_f$가 환 구조의 준동형임에 유의하자. 따라서
정리 (\hyperref[I.1.3.7-ko]{1.3.7})에 따라 $\theta_f$를 통하여
환 $A_f$와 $\Gamma(D(f),\widetilde{A})$를 동일시하면, 준동형
$\theta_f:M_f\to\Gamma(D(f),\widetilde{M})$는 \emph{가군의 동형}이
된다.

먼저 $\theta_f$가 \emph{단사}임을 보이자. 실제로
$\theta_f(\xi)=0$인 $\xi\in M_f$에 대하여, 이는 $A_f$의 모든 소
아이디얼 $\mathfrak{p}$마다 $h\notin\mathfrak{p}$이고 $h\xi=0$인
$h$가 존재한다는 뜻이다. 따라서 $\xi$의 소멸자는 $A_f$의 어느 소
아이디얼에도 포함되지 않으므로 $A_f$ 전체이고, 결국 $\xi=0$이다.

이제 $\theta_f$가 \emph{전사}임을 보이면 된다. 일반적인 경우는
(\hyperref[I.1.3.6-ko]{1.3.6})을 써서 ``국소화''함으로써 얻으므로
$f=1$인 경우로 귀착할 수 있다. $s$를 $X$ 위의
$\widetilde{M}$의 절단이라 하자. (\hyperref[I.1.3.4-ko]{1.3.4})와
(\hyperref[I.1.1.10-ko]{1.1.10, (ii)})에 따라 $X$의 \emph{유한}
덮개 $(D(f_i))_{i\in I}$($f_i\in A$)가 존재하여, 모든 $i\in I$에
대하여 제한 $s_i=s|D(f_i)$는 어떤 $\xi_i\in M_{f_i}$에 대한
$\theta_{f_i}(\xi_i)$ 꼴이다. $i$, $j$가 $I$의 두 지표이면 $s_i$와
$s_j$의 $D(f_i)\cap D(f_j)=D(f_i f_j)$로의 제한들이 같다는 조건은
$\widetilde{M}$의 정의에 따라
\begin{equation}
  \rho_{f_i f_j,f_i}(\xi_i)=\rho_{f_i f_j,f_j}(\xi_j).
  \tag{1.3.7.1}\label{I.1.3.7.1-ko}
\end{equation}
이 된다. 정의에 따라 모든 $i\in I$에 대하여
$\xi_i=z_i/f_i^{n_i}$($z_i\in M$)로 쓸 수 있다. $I$가 유한이므로
각 $z_i$에 $f_i$의 거듭제곱을 곱하여 모든 $n_i$가 같은 $n$이라고
가정할 수 있다. 그러면 (\hyperref[I.1.3.7.1-ko]{1.3.7.1})은
정의에 따라 어떤 정수 $m_{ij}\geq0$에 대하여
$(f_i f_j)^{m_{ij}}(f_j^n z_i-f_i^n z_j)=0$이라는 뜻이다. 모든
$m_{ij}$도 같은 정수 $m$이라고 가정할 수 있다. 이제 $z_i$를
$f_i^m z_i$로 바꾸면 $m=0$인 경우, 즉 모든 $i$, $j$에 대하여
\begin{equation}
  f_j^n z_i=f_i^n z_j
  \tag{1.3.7.2}\label{I.1.3.7.2-ko}
\end{equation}
인 경우로 귀착된다. 한편 $D(f_i^n)=D(f_i)$이고 $D(f_i)$들이 $X$를
덮으므로, $f_i^n$들이 생성하는 아이디얼은 $A$이다. 다시 말해
$\sum_i g_i f_i^n=1$인 원소 $g_i\in A$들이 존재한다. 이제
$M$의 원소 $z=\sum_i g_i z_i$를 생각하자. 공식
(\hyperref[I.1.3.7.2-ko]{1.3.7.2})에 따라
$f_i^n z=\sum_j g_j f_i^n z_j=(\sum_j g_j f_j^n)z_i=z_i$이고,
따라서 정의에 따라 $M_{f_i}$에서 $\xi_i=z/1$이다.
\oldpage[I]{87}
그러므로 $s_i$는 $\theta_1(z)$의 $D(f_i)$로의 제한이다. 따라서
$s=\theta_1(z)$이고 증명이 끝난다.

\begin{corollary}[1.3.8]
\label{I.1.3.8-ko}
$M$, $N$을 두 $A$-가군이라 하자. 표준 준동형
\[
  u\longmapsto\widetilde{u}:
  \operatorname{Hom}_A(M,N)\longrightarrow
  \operatorname{Hom}_{\widetilde{A}}(\widetilde{M},\widetilde{N})
\]
는 전단사이다. 특히 $M=0$과 $\widetilde{M}=0$은 서로 동치이다.
\end{corollary}

실제로 표준 준동형
\[
  v\longmapsto\Gamma(v):
  \operatorname{Hom}_{\widetilde{A}}(\widetilde{M},\widetilde{N})
  \longrightarrow
  \operatorname{Hom}_{\Gamma(\widetilde{A})}
  (\Gamma(\widetilde{M}),\Gamma(\widetilde{N}))
\]
을 생각하자. 정리 (\hyperref[I.1.3.7-ko]{1.3.7})에 따라 끝의
가군은 $\operatorname{Hom}_A(M,N)$과 표준적으로 동일시된다.
$u\mapsto\widetilde{u}$와 $v\mapsto\Gamma(v)$가 서로 역임을 확인하면
된다. $\widetilde{u}$의 정의에 따라
$\Gamma(\widetilde{u})=u$임은 자명하다. 반대로
$v\in\operatorname{Hom}_{\widetilde{A}}(\widetilde{M},\widetilde{N})$에
대하여 $u=\Gamma(v)$로 놓으면, $v$에서 표준적으로 얻는 사상
\[
  w:\Gamma(D(f),\widetilde{M})\longrightarrow
  \Gamma(D(f),\widetilde{N})
\]
에 대하여 그림
\[
  \xymatrix{
    M\ar[r]^u\ar[d]_{\rho_{f,1}} & N\ar[d]^{\rho_{f,1}}\\
    M_f\ar[r]_w & N_f
  }
\]
은 가환한다. 따라서 모든 $f\in A$에 대하여 반드시 $w=u_f$이고
(\hyperref[0.1.2.4-ko]{0, 1.2.4}), 이로부터
$\widetilde{\Gamma(v)}=v$가 따른다.

\begin{corollary}[1.3.9]
\label{I.1.3.9-ko}
\begin{enumerate}
  \item[\rm(i)] $u:M\to N$을 $A$-가군의 준동형이라 하자. 그러면
  $\operatorname{Ker}u$, $\operatorname{Im}u$,
  $\operatorname{Coker}u$에 대응하는 층들은 각각
  $\operatorname{Ker}\widetilde{u}$,
  $\operatorname{Im}\widetilde{u}$,
  $\operatorname{Coker}\widetilde{u}$이다. 특히 $\widetilde{u}$가
  단사(각각 전사, 전단사)이기 위한 필요충분조건은 $u$가
  단사(각각 전사, 전단사)인 것이다.
  \item[\rm(ii)] $M$이 $A$-가군족 $(M_\lambda)$의
  귀납극한(각각 직합)이면, $\widetilde{M}$은 표준 동형을 제외하고
  $(\widetilde{M_\lambda})$의 귀납극한(각각 직합)이다.
\end{enumerate}
\end{corollary}

(i) $\widetilde{M}$이 $M$에 관한 완전 함자라는 사실
(\hyperref[I.1.3.5-ko]{1.3.5})을 두 $A$-가군의 완전열
\[
  0\to\operatorname{Ker}u\to M\to\operatorname{Im}u\to0
\]
과
\[
  0\to\operatorname{Im}u\to N\to\operatorname{Coker}u\to0
\]
에 적용하면 된다. 두 번째 명제는 정리
(\hyperref[I.1.3.7-ko]{1.3.7})에서 따른다.

(ii) $(M_\lambda,g_{\mu\lambda})$를 귀납극한이 $M$인 $A$-가군의
귀납계라 하고, $g_\lambda:M_\lambda\to M$을 표준 준동형이라 하자.
$\lambda\leq\mu\leq\nu$에 대하여
\[
  \widetilde{g_{\nu\mu}}\circ\widetilde{g_{\mu\lambda}}
  =\widetilde{g_{\nu\lambda}},
  \qquad
  \widetilde{g_\lambda}=\widetilde{g_\mu}\circ
  \widetilde{g_{\mu\lambda}}
\]
이므로 $(\widetilde{M_\lambda},\widetilde{g_{\mu\lambda}})$는
$X$ 위 층들의 귀납계이다. 표준 준동형
\[
  h_\lambda:\widetilde{M_\lambda}\longrightarrow
  \varinjlim\widetilde{M_\lambda}
\]
를 생각하면
$v\circ h_\lambda=\widetilde{g_\lambda}$를 만족하는 유일한 준동형
\[
  v:\varinjlim\widetilde{M_\lambda}\longrightarrow\widetilde{M}
\]
가 존재한다. $v$가 전단사임을 보이려면 모든 $x\in X$에 대하여
$v_x$가 $(\varinjlim\widetilde{M_\lambda})_x$에서
$\widetilde{M}_x$ 위로 가는 전단사임을 확인하면 된다. 그런데
$\widetilde{M}_x=M_x$이고
\[
  (\varinjlim\widetilde{M_\lambda})_x
  =\varinjlim(\widetilde{M_\lambda})_x
  =\varinjlim(M_\lambda)_x=M_x
  \quad(\hyperref[0.1.3.3-ko]{0, 1.3.3}).
\]
또한 정의에 따라 $(\widetilde{g_\lambda})_x$와 $(h_\lambda)_x$는
모두 $(M_\lambda)_x$에서 $M_x$로 가는 표준 사상이다. 따라서
$(\widetilde{g_\lambda})_x=v_x\circ(h_\lambda)_x$이므로 $v_x$는
항등사상이다.
\oldpage[I]{88}
끝으로 $M$이 두 $A$-가군 $N$, $P$의 직합이면
$\widetilde{M}=\widetilde{N}\oplus\widetilde{P}$임은 자명하다. 모든
직합은 유한 직합들의 귀납극한이므로 (ii)가 증명되었다.

(\hyperref[I.1.3.8-ko]{1.3.8})에 따라 $A$-가군에 대응하는 층들과
동형인 층들은 \emph{아벨 범주}를 이룸에 유의하자(T, I, 1.4).

또한 (\hyperref[I.1.3.9-ko]{1.3.9})에서 다음이 따른다. $M$이
\emph{유한 생성} $A$-가군, 즉 전사 준동형 $A^n\to M$이 존재하는
가군이면 전사 준동형
$\widetilde{A^n}\to\widetilde{M}$이 존재한다. 다시 말해
$\widetilde{A}$-가군 $\widetilde{M}$은 \emph{$X$ 위의 유한한
절단족에 의해 생성된다}(\hyperref[0.5.1.1-ko]{0, 5.1.1}). 그 역도
성립한다.

\begin{env}[1.3.10]
\label{I.1.3.10-ko}
$N$이 $A$-가군 $M$의 부분가군이면 표준 단사 준동형 $j:N\to M$은
(\hyperref[I.1.3.9-ko]{1.3.9})에 따라 단사 준동형
$\widetilde{N}\to\widetilde{M}$을 준다. 이를 통하여
$\widetilde{N}$을 $\widetilde{M}$의
\emph{부분-$\widetilde{A}$-가군}과 표준적으로 동일시하며, 앞으로
항상 이 동일시를 사용하겠다. $N$, $P$가 $M$의 두 부분가군이면
\begin{equation}
  \widetilde{(N+P)}=\widetilde{N}+\widetilde{P}
  \tag{1.3.10.1}\label{I.1.3.10.1-ko}
\end{equation}
이고
\begin{equation}
  \widetilde{(N\cap P)}=\widetilde{N}\cap\widetilde{P}.
  \tag{1.3.10.2}\label{I.1.3.10.2-ko}
\end{equation}
실제로 $N+P$와 $N\cap P$는 각각 표준 준동형
$N\oplus P\to M$의 상과 표준 준동형
$M\to(M/N)\oplus(M/P)$의 핵이므로
(\hyperref[I.1.3.9-ko]{1.3.9})를 적용하면 된다.

(\hyperref[I.1.3.10.1-ko]{1.3.10.1})과
(\hyperref[I.1.3.10.2-ko]{1.3.10.2})에 따라
$\widetilde{N}=\widetilde{P}$이면 $N=P$이다.
\end{env}

\begin{corollary}[1.3.11]
\label{I.1.3.11-ko}
$A$-가군에 대응하는 층들과 동형인 층들의 범주에서 함자 $\Gamma$는
완전하다.
\end{corollary}

실제로 $A$-가군 준동형 $u:M\to N$, $v:N\to P$에 대응하는 완전열
\[
  \widetilde{M}\xrightarrow{\widetilde{u}}\widetilde{N}
  \xrightarrow{\widetilde{v}}\widetilde{P}
\]
을 생각하자. $Q=\operatorname{Im}u$, $R=\operatorname{Ker}v$로
놓으면 따름정리 (\hyperref[I.1.3.9-ko]{1.3.9})에 따라
\[
  \widetilde{Q}=\operatorname{Im}\widetilde{u}
  =\operatorname{Ker}\widetilde{v}=\widetilde{R},
\]
따라서 $Q=R$이다.

\begin{corollary}[1.3.12]
\label{I.1.3.12-ko}
$M$, $N$을 두 $A$-가군이라 하자.
\begin{enumerate}
  \item[\rm(i)] $M\otimes_A N$에 대응하는 층은
  $\widetilde{M}\otimes_{\widetilde{A}}\widetilde{N}$과 표준적으로
  동일시된다.
  \item[\rm(ii)] 더욱이 $M$이 유한 표시를 가지면
  $\operatorname{Hom}_A(M,N)$에 대응하는 층은
  $\mathscr{H}\!om_{\widetilde{A}}(\widetilde{M},\widetilde{N})$과
  표준적으로 동일시된다.
\end{enumerate}
\end{corollary}

(i) 층
$\mathscr{F}=\widetilde{M}\otimes_{\widetilde{A}}\widetilde{N}$은
$X$에서 $D(f)$($f\in A$)들로 이루어진 기저
(\hyperref[I.1.1.10-ko]{1.1.10, (i)}) 위의 준층
\[
  U\longmapsto\mathscr{F}(U)=\Gamma(U,\widetilde{M})
  \otimes_{\Gamma(U,\widetilde{A})}\Gamma(U,\widetilde{N})
\]
에 대응한다. 그런데 (\hyperref[I.1.3.7-ko]{1.3.7})과
(\hyperref[I.1.3.6-ko]{1.3.6})에 따라 $\mathscr{F}(D(f))$는
$M_f\otimes_{A_f}N_f$와 표준적으로 동일시된다. 한편 이
$A_f$-가군은 $(M\otimes_A N)_f$와 표준적으로 동형이고
(\hyperref[0.1.3.4-ko]{0, 1.3.4}), 다시 이는
$\Gamma(D(f),\widetilde{(M\otimes_A N)})$과 표준적으로 동형이다
(\hyperref[I.1.3.7-ko]{1.3.7} 및
\hyperref[I.1.3.6-ko]{1.3.6}). 또한 이렇게 얻은 표준 동형
\[
  \mathscr{F}(D(f))\simeq
  \Gamma(D(f),\widetilde{(M\otimes_A N)})
\]
\oldpage[I]{89}
들은 제한 준동형과의 양립 조건을 만족한다
(\hyperref[0.1.4.2-ko]{0, 1.4.2}). 따라서 이들은 함자적인 표준 동형
\[
  \widetilde{M}\otimes_{\widetilde{A}}\widetilde{N}
  \simeq\widetilde{(M\otimes_A N)}
\]
을 정의한다.

(ii) 층
$\mathscr{G}=\mathscr{H}\!om_{\widetilde{A}}(\widetilde{M},
\widetilde{N})$은 $D(f)$들로 이루어진 $X$의 기저 위의 준층
\[
  U\longmapsto\mathscr{G}(U)
  =\operatorname{Hom}_{\widetilde{A}|U}
  (\widetilde{M}|U,\widetilde{N}|U)
\]
에 대응한다. 그런데 $\mathscr{G}(D(f))$는
$\operatorname{Hom}_{A_f}(M_f,N_f)$와 표준적으로 동일시된다
(\hyperref[I.1.3.6-ko]{1.3.6} 및
\hyperref[I.1.3.8-ko]{1.3.8}). 이는 $M$에 관한 가정에 따라 다시
$(\operatorname{Hom}_A(M,N))_f$와 표준적으로 동일시되고
(\hyperref[0.1.3.5-ko]{0, 1.3.5}), 끝으로
$(\operatorname{Hom}_A(M,N))_f$는
$\Gamma(D(f),\widetilde{(\operatorname{Hom}_A(M,N))})$과
표준적으로 동일시된다(\hyperref[I.1.3.6-ko]{1.3.6} 및
\hyperref[I.1.3.7-ko]{1.3.7}). 이렇게 얻은 표준 동형
\[
  \mathscr{G}(D(f))\simeq
  \Gamma(D(f),\widetilde{(\operatorname{Hom}_A(M,N))})
\]
들은 제한 준동형과 양립하므로
(\hyperref[0.1.4.2-ko]{0, 1.4.2}) 표준 동형
\[
  \mathscr{H}\!om_{\widetilde{A}}(\widetilde{M},\widetilde{N})
  \simeq\widetilde{(\operatorname{Hom}_A(M,N))}
\]
을 정의한다.

\begin{env}[1.3.13]
\label{I.1.3.13-ko}
이제 $B$를 (가환) $A$-대수라 하자. 이는 $B$가 $A$-가군이고,
한 원소 $e\in B$와 $A$-준동형
$\varphi:B\otimes_A B\to B$가 주어져 다음 조건들을 만족한다고
해석할 수 있다. $1^\circ$ 그림
\[
  \xymatrix{
    B\otimes_A B\otimes_A B\ar[r]^{\varphi\otimes 1}
      \ar[d]_{1\otimes\varphi} &
    B\otimes_A B\ar[d]^\varphi & &
    B\otimes_A B\ar[rr]^\sigma\ar[rd]_\varphi & &
    B\otimes_A B\ar[dl]^\varphi\\
    B\otimes_A B\ar[r]_\varphi &
    B & & & B
  }
\]
들이 가환한다($\sigma$는 표준 대칭사상). $2^\circ$
$\varphi(e\otimes x)=\varphi(x\otimes e)=x$이다. 따름정리
(\hyperref[I.1.3.12-ko]{1.3.12})에 따라 $\widetilde{A}$-가군의
준동형
\[
  \widetilde{\varphi}:
  \widetilde{B}\otimes_{\widetilde{A}}\widetilde{B}
  \longrightarrow\widetilde{B}
\]
도 같은 조건들을 만족하므로 $\widetilde{B}$에
$\widetilde{A}$-대수 구조를 정의한다. 마찬가지로 $B$-가군 $N$을
주는 것은 $A$-가군 $N$과, 그림
\[
  \xymatrix{
    B\otimes_A B\otimes_A N\ar[r]^{\varphi\otimes 1}
      \ar[d]_{1\otimes\psi} &
    B\otimes_A N\ar[d]^\psi\\
    B\otimes_A N\ar[r]_\psi & N
  }
\]
이 가환하고 $\psi(e\otimes n)=n$을 만족하는 $A$-준동형
$\psi:B\otimes_A N\to N$을 주는 것과 같다. 준동형
\[
  \widetilde{\psi}:
  \widetilde{B}\otimes_{\widetilde{A}}\widetilde{N}
  \longrightarrow\widetilde{N}
\]
도 같은 조건을 만족하므로 $\widetilde{N}$에
$\widetilde{B}$-가군 구조를 정의한다.

같은 방식으로 다음을 알 수 있다. $u:B\to B'$가 $A$-대수의
준동형(각각 $v:N\to N'$이 $B$-가군의 준동형)이면
$\widetilde{u}$는 $\widetilde{A}$-대수의 준동형(각각
$\widetilde{v}$는 $\widetilde{B}$-가군의 준동형)이고,
$\operatorname{Ker}\widetilde{u}$는 $\widetilde{B}$-아이디얼이며
(각각 $\operatorname{Ker}\widetilde{v}$,
$\operatorname{Coker}\widetilde{v}$,
$\operatorname{Im}\widetilde{v}$는 $\widetilde{B}$-가군이다).
$N$이 $B$-가군이면 $\widetilde{N}$이 유한 생성
$\widetilde{B}$-가군이기 위한 필요충분조건은 $N$이 유한 생성
$B$-가군인 것이다(\hyperref[0.5.2.3-ko]{0, 5.2.3}).
\oldpage[I]{90}
$M$, $N$이 두 $B$-가군이면 $\widetilde{B}$-가군
$\widetilde{M}\otimes_{\widetilde{B}}\widetilde{N}$은
$\widetilde{(M\otimes_B N)}$과 표준적으로 동일시된다. 또한 $M$이
유한 표시를 가지면
$\mathscr{H}\!om_{\widetilde{B}}(\widetilde{M},\widetilde{N})$은
$\widetilde{(\operatorname{Hom}_B(M,N))}$과 표준적으로 동일시된다.
증명은 (\hyperref[I.1.3.12-ko]{1.3.12})와 같다.

$\mathfrak{J}$가 $B$의 아이디얼이고 $N$이 $B$-가군이면
\[
  \widetilde{(\mathfrak{J}N)}
  =\widetilde{\mathfrak{J}}\cdot\widetilde{N}.
\]

끝으로 $B$가 부분-$A$-가군 $B_n$($n\in\mathbb{Z}$)들로 등급이
매겨진 $A$-대수이면, $\widetilde{A}$-가군
$\widetilde{B_n}$들의 직합인 $\widetilde{A}$-대수
$\widetilde{B}$(\hyperref[I.1.3.9-ko]{1.3.9})도 이
부분-$\widetilde{A}$-가군들로 등급이 매겨진다. 실제로 등급 공리는
준동형 $B_m\otimes_A B_n\to B$의 상이 $B_{m+n}$에 포함된다는
조건으로 표현된다. 마찬가지로 $M$이 부분가군 $M_n$들로 등급이
매겨진 $B$-가군이면, $\widetilde{M}$은
$\widetilde{M_n}$들로 등급이 매겨진 $\widetilde{B}$-가군이다.
\end{env}

\begin{env}[1.3.14]
\label{I.1.3.14-ko}
$B$가 $A$-대수이고 $M$이 $B$의 부분가군이면,
$\widetilde{M}$이 생성하는 $\widetilde{B}$의
부분-$\widetilde{A}$-대수(\hyperref[0.4.1.3-ko]{0, 4.1.3})는
$\widetilde{C}$이다. 여기서 $C$는 $M$이 생성하는 $B$의 부분대수이다.
실제로 $C$는 준동형 $\bigotimes_A^n M\to B$($n\geq0$)의 상인
$B$의 부분가군들의 합이고, (\hyperref[I.1.3.9-ko]{1.3.9})와
(\hyperref[I.1.3.12-ko]{1.3.12})를 적용하면 된다.
\end{env}

\subsection{소 스펙트럼 위의 준연접층}
\label{subsection:I.1.4-ko}

\begin{theorem}[1.4.1]
\label{I.1.4.1-ko}
$X$를 환 $A$의 소 스펙트럼, $V$를 $X$의 준콤팩트 열린 부분집합,
$\mathscr{F}$를 $(\mathscr{O}_X|V)$-가군이라 하자. 다음 네 조건은
서로 동치이다.
\begin{enumerate}
  \item[\textit{a})] $\mathscr{F}$가 $\widetilde{M}|V$와 동형이 되는
  $A$-가군 $M$이 존재한다.
  \item[\textit{b})] $V$에 포함되고 $D(f_i)$($f_i\in A$) 꼴인
  집합들로 이루어진 $V$의 유한 열린 덮개 $(V_i)$가 존재하여, 모든
  $i$에 대하여 $\mathscr{F}|V_i$는 어떤 $A_{f_i}$-가군 $M_i$에
  대응하는 층 $\widetilde{M_i}$와 동형이다.
  \item[\textit{c})] 층 $\mathscr{F}$는 준연접이다
  (\hyperref[0.5.1.3-ko]{0, 5.1.3}).
  \item[\textit{d})] 다음 두 성질이 성립한다.
  \begin{enumerate}
    \item[\textit{d}\,1)] $D(f)\subset V$인 모든 $f\in A$와 모든 절단
    $s\in\Gamma(D(f),\mathscr{F})$에 대하여, $f^n s$가 $V$ 위의
    $\mathscr{F}$의 절단으로 연장되는 정수 $n\geq0$이 존재한다.
    \item[\textit{d}\,2)] $D(f)\subset V$인 모든 $f\in A$와,
    $D(f)$로의 제한이 $0$인 모든 절단
    $t\in\Gamma(V,\mathscr{F})$에 대하여, $f^n t=0$인 정수
    $n\geq0$이 존재한다.
  \end{enumerate}
\end{enumerate}
\end{theorem}

(조건 \textit{d}\,1)과 \textit{d}\,2)를 쓸 때에는
(\hyperref[I.1.3.7-ko]{1.3.7})에 따라 $A$와
$\Gamma(\widetilde{A})$를 암묵적으로 동일시하였다.)

\textit{a})가 \textit{b})를 함의함은
(\hyperref[I.1.3.6-ko]{1.3.6})과 $D(f_i)$들이 $X$의 위상의 기저를
이룬다는 사실(\hyperref[I.1.1.10-ko]{1.1.10})에서 바로 따른다.
모든 $A$-가군은 어떤 준동형 $A^{(I)}\to A^{(J)}$의 여핵과
동형이므로, (\hyperref[I.1.3.9-ko]{1.3.9})에 따라 모든 $A$-가군에
대응하는 층은 준연접이다. 따라서 \textit{b})는 \textit{c})를
함의한다. 역으로 $\mathscr{F}$가 준연접이면 모든 $x\in V$는 다음
성질을 갖는 $D(f)\subset V$ 꼴의 근방을 가진다. 즉,
$\mathscr{F}|D(f)$는 어떤 준동형
$\widetilde{A_f}^{(I)}\to\widetilde{A_f}^{(J)}$의 여핵과 동형이고,
따라서 이에 대응하는 준동형 $A_f^{(I)}\to A_f^{(J)}$의 여핵인
가군 $N$에 대응하는 층 $\widetilde{N}$과 동형이다
(\hyperref[I.1.3.8-ko]{1.3.8} 및
\hyperref[I.1.3.9-ko]{1.3.9}). $V$가 준콤팩트이므로
\textit{c})가 \textit{b})를 함의함은 자명하다.
\oldpage[I]{91}
\textit{b})가 \textit{d}\,1)과 \textit{d}\,2)를 함의함을
증명하기 위해 먼저 어떤 $g\in A$에 대하여 $V=D(g)$이고,
$\mathscr{F}$가 $A_g$-가군 $N$에 대응하는 층
$\widetilde{N}$과 동형이라고 하자. $X$를 $V$로, $A$를 $A_g$로
바꾸면(\hyperref[I.1.3.6-ko]{1.3.6}) $g=1$인 경우로 귀착할 수
있다. 그러면 $\Gamma(D(f),\widetilde{N})$과 $N_f$는 표준적으로
동일시된다(\hyperref[I.1.3.6-ko]{1.3.6} 및
\hyperref[I.1.3.7-ko]{1.3.7}). 따라서 절단
$s\in\Gamma(D(f),\widetilde{N})$는 $z\in N$인 $z/f^n$ 꼴의
원소와 동일시된다. 절단 $f^n s$는 $N_f$의 원소 $z/1$과
동일시되고, 따라서 $z\in N$과 동일시되는 $X$ 위의
$\widetilde{N}$의 절단을 $D(f)$로 제한한 것이다. 이로써 이 경우의
\textit{d}\,1)을 얻는다. 마찬가지로
$t\in\Gamma(X,\widetilde{N})$는 어떤 원소 $z'\in N$과 동일시된다.
$t$의 $D(f)$로의 제한은 $N_f$에서 $z'$의 상 $z'/1$과 동일시되고,
이 상이 $0$이라는 것은 $N$에서 $f^n z'=0$인 어떤 $n\geq0$이
존재한다는 뜻이다. 이는 $f^n t=0$과 같다.

\textit{b})가 \textit{d}\,1)과 \textit{d}\,2)를 함의함을
마치려면 다음 보조정리를 증명하면 충분하다.

\begin{lemma}[1.4.1.1]
\label{I.1.4.1.1-ko}
$V$가 $D(g_i)$ 꼴의 집합들의 유한 합집합이고, 각 층
$\mathscr{F}|D(g_i)$와
$\mathscr{F}|(D(g_i)\cap D(g_j))=\mathscr{F}|D(g_i g_j)$가
\textit{d}\,1)과 \textit{d}\,2)를 만족한다고 하자. 그러면
$\mathscr{F}$는 다음 두 성질을 갖는다.
\begin{enumerate}
  \item[\textit{d}'\,1)] 모든 $f\in A$와 모든 절단
  $s\in\Gamma(D(f)\cap V,\mathscr{F})$에 대하여, $f^n s$가
  $V$ 위의 $\mathscr{F}$의 절단으로 연장되는 정수 $n\geq0$이
  존재한다.
  \item[\textit{d}'\,2)] 모든 $f\in A$와, $D(f)\cap V$로의 제한이
  $0$인 모든 절단 $t\in\Gamma(V,\mathscr{F})$에 대하여,
  $f^n t=0$인 정수 $n\geq0$이 존재한다.
\end{enumerate}
\end{lemma}

먼저 \textit{d}'\,2)를 증명하자. $D(f)\cap D(g_i)=D(fg_i)$이므로
모든 $i$에 대하여 $(fg_i)^{n_i}t$의 $D(g_i)$로의 제한이 $0$인
정수 $n_i$가 존재한다. $A_{g_i}$에서 $g_i$의 상은 가역원이므로
$f^{n_i}t$의 $D(g_i)$로의 제한도 $0$이다. $n$을 $n_i$들의
최댓값으로 취하면 \textit{d}'\,2)를 얻는다.

\textit{d}'\,1)을 증명하기 위해 층 $\mathscr{F}|D(g_i)$에
\textit{d}\,1)을 적용하자. 그러면 정수 $n_i\geq0$과
$D(g_i)$ 위의 $\mathscr{F}$의 절단 $s_i'$이 존재하여,
$s_i'$은 $(fg_i)^{n_i}s$의 $D(fg_i)$로의 제한을 연장한다.
$A_{g_i}$에서 $g_i$의 상이 가역원이므로, $s_i'=g_i^{n_i}s_i$이고
$f^{n_i}s$의 $D(fg_i)$로의 제한을 연장하는 $D(g_i)$ 위의 절단
$s_i$가 존재한다. 모든 $n_i$가 같은 정수 $n$이라고 가정할 수도
있다. 구성에 따라 $s_i-s_j$의
$D(f)\cap D(g_i)\cap D(g_j)=D(fg_i g_j)$로의 제한은 $0$이다.
층 $\mathscr{F}|D(g_i g_j)$에 \textit{d}\,2)를 적용하면,
$(fg_i g_j)^{m_{ij}}(s_i-s_j)$의 $D(g_i g_j)$로의 제한이 $0$인
정수 $m_{ij}\geq0$이 존재한다. $A_{g_i g_j}$에서 $g_i g_j$의 상이
가역원이므로 $f^{m_{ij}}(s_i-s_j)$의 $D(g_i g_j)$로의 제한도
$0$이다. 모든 $m_{ij}$가 같은 정수 $m$이라고 가정할 수 있다.
그러면 $f^m s_i$들을 연장하는 절단
$s'\in\Gamma(V,\mathscr{F})$가 존재하고, 이 절단은
$f^{n+m}s$를 연장한다. 이로써 \textit{d}'\,1)을 얻는다.

이제 \textit{d}\,1)과 \textit{d}\,2)가 \textit{a})를
함의함을 증명하면 된다. 먼저 이 두 조건이, $D(g)\subset V$인 모든
$g\in A$에 대하여 층 $\mathscr{F}|D(g)$에도 성립함을 보이자.
\textit{d}\,1)에 대해서는 자명하다. 한편
$t\in\Gamma(D(g),\mathscr{F})$이고 그 $D(f)\subset D(g)$로의
제한이 $0$이라고 하자. \textit{d}\,1)에 따라 $g^m t$가 $V$ 위의
$\mathscr{F}$의 절단 $s$로 연장되는 정수 $m\geq0$이 존재한다.
\oldpage[I]{92}
\textit{d}\,2)를 적용하면 $f^n g^m t=0$인 정수 $n\geq0$이
존재한다. $A_g$에서 $g$의 상이 가역원이므로 $f^n t=0$이다.

이제 $V$가 준콤팩트이므로 보조정리
(\hyperref[I.1.4.1.1-ko]{1.4.1.1})에 따라 조건
\textit{d}'\,1)과 \textit{d}'\,2)가 성립한다. $A$-가군
$M=\Gamma(V,\mathscr{F})$을 생각하자. $j:V\to X$를 표준
포함사상이라 하고 $\widetilde{A}$-가군 준동형
\[
  u:\widetilde{M}\longrightarrow j_*(\mathscr{F})
\]
을 정의하겠다. $D(f)$들이 $X$의 위상의 기저를 이루므로, 모든
$f\in A$에 대하여 통상적인 양립 조건
(\hyperref[0.3.2.5-ko]{0, 3.2.5})을 만족하는 준동형
\[
  u_f:M_f\longrightarrow
  \Gamma(D(f),j_*(\mathscr{F}))
  =\Gamma(D(f)\cap V,\mathscr{F})
\]
을 정의하면 충분하다. $A_f$에서 $f$의 표준 상이 가역원이므로 제한
준동형
\[
  M=\Gamma(V,\mathscr{F})\longrightarrow
  \Gamma(D(f)\cap V,\mathscr{F})
\]
은 다음과 같이 분해된다
(\hyperref[0.1.2.4-ko]{0, 1.2.4}).
\[
  M\longrightarrow M_f\xrightarrow{u_f}
  \Gamma(D(f)\cap V,\mathscr{F}).
\]
$D(g)\subset D(f)$에 대한 양립 조건은 바로 확인된다.
\textit{d}'\,1)(각각 \textit{d}'\,2))이 모든 $u_f$가
전사(각각 단사)임을 함의함을 보이면 $u$가 \emph{전단사}이고,
따라서 $\mathscr{F}$는 $\widetilde{M}|V$와 동형이다. 그런데
$s\in\Gamma(D(f)\cap V,\mathscr{F})$이면 \textit{d}'\,1)에 따라
$f^n s$가 어떤 절단 $z\in M$으로 연장되는 정수 $n\geq0$이
존재한다. 그러면 $u_f(z/f^n)=s$이므로 $u_f$는 전사이다. 마찬가지로
$u_f(z/1)=0$인 $z\in M$에 대하여, 이는 절단 $z$의
$D(f)\cap V$로의 제한이 $0$이라는 뜻이다. \textit{d}'\,2)에 따라
$f^n z=0$인 정수 $n\geq0$이 존재하므로 $M_f$에서 $z/1=0$이고,
따라서 $u_f$는 단사이다.

\par\hfill C.Q.F.D.\par

\begin{corollary}[1.4.2]
\label{I.1.4.2-ko}
$X$의 준콤팩트 열린집합 위의 모든 준연접층은 $X$ 위의 어떤
준연접층으로부터 유도된다.
\end{corollary}

\begin{corollary}[1.4.3]
\label{I.1.4.3-ko}
$X=\operatorname{Spec}(A)$ 위의 모든 준연접
$\mathscr{O}_X$-대수는 어떤 $A$-대수 $B$에 대한
$\widetilde{B}$ 꼴의 $\mathscr{O}_X$-대수와 동형이다. 모든 준연접
$\widetilde{B}$-가군은 어떤 $B$-가군 $N$에 대한
$\widetilde{N}$ 꼴의 $\widetilde{B}$-가군과 동형이다.
\end{corollary}

실제로 준연접 $\mathscr{O}_X$-대수는 준연접
$\mathscr{O}_X$-가군이므로, 어떤 $A$-가군 $B$에 대한
$\widetilde{B}$ 꼴이다. $B$가 $A$-대수라는 것은
$\widetilde{A}$-가군 준동형
\[
  \widetilde{B}\otimes_{\widetilde{A}}\widetilde{B}
  \longrightarrow\widetilde{B}
\]
을 써서 $\mathscr{O}_X$-대수 구조를 특징짓는 것과
(\hyperref[I.1.3.12-ko]{1.3.12})에서 따른다. $\mathscr{G}$가 준연접
$\widetilde{B}$-가군이면 $\mathscr{G}$가 준연접
$\widetilde{A}$-가군임을 보이기만 하면 같은 방식으로 결론을 얻는다.
문제는 국소적이므로 $X$의 $D(f)$ 꼴 열린집합으로 제한하여
$\mathscr{G}$가 $\widetilde{B}$-가군(따라서 특히
$\widetilde{A}$-가군)의 준동형
\[
  \widetilde{B}^{(I)}\longrightarrow\widetilde{B}^{(J)}
\]
의 여핵이라고 가정할 수 있다. 이제 명제는
(\hyperref[I.1.3.8-ko]{1.3.8})과
(\hyperref[I.1.3.9-ko]{1.3.9})에서 따른다.

\subsection{소 스펙트럼 위의 연접층}
\label{subsection:I.1.5-ko}

\begin{theorem}[1.5.1]
\label{I.1.5.1-ko}
$A$를 \emph{뇌터 환}, $X=\operatorname{Spec}(A)$를 그 소 스펙트럼,
$V$를 $X$의 열린 부분집합, $\mathscr{F}$를
$(\mathscr{O}_X|V)$-가군이라 하자. 다음 조건들은 서로 동치이다.
\begin{enumerate}
  \item[\textit{a})] $\mathscr{F}$는 연접이다.
  \item[\textit{b})] $\mathscr{F}$는 유한 생성이고 준연접이다.
  \item[\textit{c})] $\mathscr{F}$가 층 $\widetilde{M}|V$와 동형이
  되는 유한 생성 $A$-가군 $M$이 존재한다.
\end{enumerate}
\end{theorem}

\oldpage[I]{93}
\textit{a})가 \textit{b})를 함의함은 자명하다. \textit{b})가
\textit{c})를 함의함을 보이자. $V$는 준콤팩트이므로
(\hyperref[0.2.2.3-ko]{0, 2.2.3}), 먼저 $\mathscr{F}$는 어떤
$A$-가군 $N$에 대한 층 $\widetilde{N}|V$와 동형이다
(\hyperref[I.1.4.1-ko]{1.4.1}). 그런데 $M_\lambda$가 $N$의 유한 생성
부분-$A$-가군 전체를 달릴 때
$N=\varinjlim M_\lambda$이므로, (\hyperref[I.1.3.9-ko]{1.3.9})에 따라
\[
  \mathscr{F}=\widetilde{N}|V
  =\varinjlim\widetilde{M_\lambda}|V.
\]
한편 $\mathscr{F}$는 유한 생성이고 $V$는 준콤팩트이므로 어떤 지표
$\lambda$에 대하여
$\mathscr{F}=\widetilde{M_\lambda}|V$이다
(\hyperref[0.5.2.3-ko]{0, 5.2.3}).

끝으로 \textit{c})가 \textit{a})를 함의함을 보이자.
$\mathscr{F}$가 유한 생성임은 분명하다
(\hyperref[I.1.3.6-ko]{1.3.6} 및
\hyperref[I.1.3.9-ko]{1.3.9}). 또한 문제는 국소적이므로
$V=D(f)$($f\in A$)인 경우만 생각하면 된다. $A_f$는 뇌터 환이므로
$A$를 $A_f$로 바꾸어, 결국 $M$이 $A$-가군일 때 준동형
$\widetilde{A^n}\to\widetilde{M}$의 핵이 유한 생성임을 증명하면
충분하다. 이러한 준동형은 어떤 준동형 $u:A^n\to M$에 대한
$\widetilde{u}$ 꼴이다(\hyperref[I.1.3.8-ko]{1.3.8}).
$P=\operatorname{Ker}u$라 놓으면
\[
  \widetilde{P}=\operatorname{Ker}\widetilde{u}
\]
이다(\hyperref[I.1.3.9-ko]{1.3.9}). $A$가 뇌터 환이므로 $P$는
유한 생성이고, 이로써 증명이 끝난다.

\begin{corollary}[1.5.2]
\label{I.1.5.2-ko}
(\hyperref[I.1.5.1-ko]{1.5.1})의 가정 아래에서 층
$\mathscr{O}_X$는 연접인 환의 층이다.
\end{corollary}

\begin{corollary}[1.5.3]
\label{I.1.5.3-ko}
(\hyperref[I.1.5.1-ko]{1.5.1})의 가정 아래에서 $X$의 열린집합 위의
모든 연접층은 $X$ 위의 어떤 연접층으로부터 유도된다.
\end{corollary}

\begin{corollary}[1.5.4]
\label{I.1.5.4-ko}
(\hyperref[I.1.5.1-ko]{1.5.1})의 가정 아래에서 모든 준연접
$\mathscr{O}_X$-가군 $\mathscr{F}$는 $\mathscr{F}$의 연접
부분-$\mathscr{O}_X$-가군들의 귀납극한이다.
\end{corollary}

실제로 $\mathscr{F}=\widetilde{M}$인 $A$-가군 $M$이 존재하고,
$M$은 그 유한 생성 부분가군들의 귀납극한이다. 이제
(\hyperref[I.1.3.9-ko]{1.3.9})와
(\hyperref[I.1.5.1-ko]{1.5.1})을 적용하면 된다.

\subsection{소 스펙트럼 위의 준연접층의 함자적 성질}
\label{subsection:I.1.6-ko}

\begin{env}[1.6.1]
\label{I.1.6.1-ko}
$A$, $A'$을 두 환,
\[
  \varphi:A'\to A
\]
를 준동형,
\[
  {}^a\varphi:X=\operatorname{Spec}(A)\to
  X'=\operatorname{Spec}(A')
\]
를 $\varphi$에 대응하는 연속사상이라 하자
(\hyperref[I.1.2.1-ko]{1.2.1}). 환의 층의 \emph{표준 준동형}
\[
  \widetilde{\varphi}:\mathscr{O}_{X'}\to
  {}^a\varphi_*(\mathscr{O}_X)
\]
을 정의하겠다. 모든 $f'\in A'$에 대하여 $f=\varphi(f')$로 놓자.
그러면
${}^a\varphi^{-1}(D(f'))=D(f)$이다
(\hyperref[I.1.2.2.2-ko]{1.2.2.2}). 환
$\Gamma(D(f'),\widetilde{A'})$와
$\Gamma(D(f),\widetilde{A})$는 각각 $A'_{f'}$와 $A_f$와 동일시된다
(\hyperref[I.1.3.6-ko]{1.3.6} 및
\hyperref[I.1.3.7-ko]{1.3.7}). 한편 준동형 $\varphi$는 표준적으로
준동형 $\varphi_{f'}:A'_{f'}\to A_f$를 정의한다
(\hyperref[0.1.5.1-ko]{0, 1.5.1}). 다시 말해 환 준동형
\[
  \Gamma(D(f'),\widetilde{A'})\to
  \Gamma({}^a\varphi^{-1}(D(f')),\widetilde{A})
  =\Gamma(D(f'),{}^a\varphi_*(\widetilde{A}))
\]
이 존재한다.

\oldpage[I]{94}
또한 이 준동형들은 통상적인 양립 조건을 만족한다. 즉,
$D(f')\supset D(g')$일 때 그림
\[
  \xymatrix{
    \Gamma(D(f'),\widetilde{A'})\ar[r]\ar[d] &
    \Gamma(D(f'),{}^a\varphi_*(\widetilde{A}))\ar[d]\\
    \Gamma(D(g'),\widetilde{A'})\ar[r] &
    \Gamma(D(g'),{}^a\varphi_*(\widetilde{A}))
  }
\]
은 가환이다(\hyperref[0.1.5.1-ko]{0, 1.5.1}). $D(f')$들이
$X'$의 위상의 기저를 이루므로
(\hyperref[0.3.2.3-ko]{0, 3.2.3}), 이로써
$\mathscr{O}_{X'}$-대수의 준동형을 정의하였다. 따라서 쌍
$\Phi=({}^a\varphi,\widetilde{\varphi})$는 환 달린 공간의 사상
\[
  \Phi:(X,\mathscr{O}_X)\to(X',\mathscr{O}_{X'})
\]
이다(\hyperref[0.4.1.1-ko]{0, 4.1.1}).

더욱이 $x'={}^a\varphi(x)$로 놓으면 준동형
$\widetilde{\varphi}_x^\#$
(\hyperref[0.3.7.1-ko]{0, 3.7.1})은 바로 $\varphi:A'\to A$에서
표준적으로 유도되는 준동형
\[
  \varphi_x:A'_{x'}\to A_x
\]
이다(\hyperref[0.1.5.1-ko]{0, 1.5.1}). 실제로 모든
$z'\in A'_{x'}$은 $f',g'\in A'$이고
$f'\notin\mathfrak{j}_{x'}$인 $g'/f'$ 꼴로 쓸 수 있다. 따라서
$D(f')$는 $X'$에서 $x'$의 근방이고, $\widetilde{\varphi}$에서
유도되는 준동형
$\Gamma(D(f'),\widetilde{A'})\to
\Gamma({}^a\varphi^{-1}(D(f')),\widetilde{A})$은
$\varphi_{f'}$와 같다. $g'/f'\in A'_{f'}$에 대응하는 절단
$s'\in\Gamma(D(f'),\widetilde{A'})$을 생각하면 $A_x$에서
$\widetilde{\varphi}_x^\#(z')=\varphi(g')/\varphi(f')$을 얻는다.
\end{env}

\begin{example}[1.6.2]
\label{I.1.6.2-ko}
$S$를 $A$의 곱셈적 부분집합, $\varphi$를 표준 준동형
$A\to S^{-1}A$라 하자. 앞서
(\hyperref[I.1.2.6-ko]{1.2.6})에서 ${}^a\varphi$가
$Y=\operatorname{Spec}(S^{-1}A)$에서
$X=\operatorname{Spec}(A)$의 부분공간
\[
  \{x\in X\mid\mathfrak{j}_x\cap S=\varnothing\}
\]
으로 가는 \emph{위상동형}임을 보았다. 또한 이 부분공간의 모든
$x={}^a\varphi(y)$($y\in Y$)에 대하여 준동형
$\widetilde{\varphi}_y^\#:\mathscr{O}_x\to\mathscr{O}_y$는
\emph{전단사}이다(\hyperref[0.1.2.6-ko]{0, 1.2.6}). 다시 말해
$\mathscr{O}_Y$는 $\mathscr{O}_X$가 $Y$ 위에 유도하는 층과
동일시된다.
\end{example}

\begin{proposition}[1.6.3]
\label{I.1.6.3-ko}
모든 $A$-가군 $M$에 대하여 $\mathscr{O}_{X'}$-가군
$(M_{[\varphi]})^{\sim}$에서 직상 $\Phi_*(\widetilde{M})$으로 가는
함자적인 표준 동형이 존재한다.
\end{proposition}

간단히 $M'=M_{[\varphi]}$로 놓고, 모든 $f'\in A'$에 대하여 다시
$f=\varphi(f')$로 놓자. 절단가군
$\Gamma(D(f'),\widetilde{M'})$와
$\Gamma(D(f),\widetilde{M})$는 각각 $A'_{f'}$-가군 $M'_{f'}$와
$A_f$-가군 $M_f$와 동일시된다. 또한 $A'_{f'}$-가군
$(M_f)_{[\varphi_{f'}]}$는 $M'_{f'}$와 표준적으로 동형이다
(\hyperref[0.1.5.2-ko]{0, 1.5.2}). 따라서
$\Gamma(D(f'),\widetilde{A'})$-가군의 함자적 동형
\[
  \Gamma(D(f'),\widetilde{M'})\xrightarrow{\sim}
  \Gamma({}^a\varphi^{-1}(D(f')),\widetilde{M})_{[\varphi_{f'}]}
\]
이 존재한다. 이 동형들은 제한사상과의 통상적인 양립 조건을
만족하므로(\hyperref[0.1.5.6-ko]{0, 1.5.6}) 앞에서 말한 함자적 동형을
정의한다. 더 정확히, $u:M_1\to M_2$가 $A$-가군 준동형이면 이를
$A'$-가군 준동형
$(M_1)_{[\varphi]}\to(M_2)_{[\varphi]}$로 볼 수 있다. 이 준동형을
$u_{[\varphi]}$라 하면 $\Phi_*(\widetilde{u})$는
$(u_{[\varphi]})^{\sim}$와 동일시된다.

이 증명은 모든 $A$-대수 $B$에 대하여 함자적인 표준 동형
\oldpage[I]{95}
\[
  (B_{[\varphi]})^{\sim}\xrightarrow{\sim}\Phi_*(\widetilde{B})
\]
이 $\mathscr{O}_{X'}$-대수의 동형임도 보여 준다. $M$이
$B$-가군이면 함자적인 표준 동형
\[
  (M_{[\varphi]})^{\sim}\xrightarrow{\sim}\Phi_*(\widetilde{M})
\]
은 $\Phi_*(\widetilde{B})$-가군의 동형이다.

\begin{corollary}[1.6.4]
\label{I.1.6.4-ko}
직상함자 $\Phi_*$는 준연접 $\mathscr{O}_X$-가군의 범주에서
완전하다.
\end{corollary}

실제로 $M_{[\varphi]}$는 $M$에 관한 완전 함자이고,
$\widetilde{M'}$은 $M'$에 관한 완전 함자이다
(\hyperref[I.1.3.5-ko]{1.3.5}).

\begin{proposition}[1.6.5]
\label{I.1.6.5-ko}
$N'$을 $A'$-가군이라 하고 $N=N'\otimes_{A'}A_{[\varphi]}$를
$A$-가군이라 하자. $\mathscr{O}_X$-가군
$\Phi^*(\widetilde{N'})$에서 $\widetilde{N}$으로 가는 함자적인
표준 동형이 존재한다.
\end{proposition}

먼저 $j:z'\mapsto z'\otimes1$이 $N'$에서 $N_{[\varphi]}$으로 가는
$A'$-준동형임을 보자. 실제로 $f'\in A'$에 대하여 정의상
$(f'z')\otimes1=z'\otimes\varphi(f')=\varphi(f')(z'\otimes1)$이다.
따라서 (\hyperref[I.1.3.8-ko]{1.3.8})에 의해
$\mathscr{O}_{X'}$-가군 준동형
$\widetilde{j}:\widetilde{N'}\to(N_{[\varphi]})^{\sim}$을 얻는다.
(\hyperref[I.1.6.3-ko]{1.6.3})에 따라 $\widetilde{j}$가
$\widetilde{N'}$을 $\Phi_*(\widetilde{N})$으로 보낸다고 보아도 된다.
이 준동형 $\widetilde{j}$에는 표준적으로 준동형
$h=\widetilde{j}^{\#}:\Phi^*(\widetilde{N'})\to\widetilde{N}$이
대응한다(\hyperref[0.4.4.3-ko]{0, 4.4.3}). 모든 줄기에서 $h_x$가
\emph{전단사}임을 보이자. $x'={}^a\varphi(x)$로 놓고
$x'\in D(f')$인 $f'\in A'$를 잡은 뒤 $f=\varphi(f')$로 놓자. 환
$\Gamma(D(f),\widetilde{A})$는 $A_f$와, 가군
$\Gamma(D(f),\widetilde{N})$와
$\Gamma(D(f'),\widetilde{N'})$은 각각 $N_f$와 $N'_{f'}$와
동일시된다. 절단 $s'\in\Gamma(D(f'),\widetilde{N'})$이
$n'/{f'}^p$($n'\in N'$)와 동일시된다고 하고, $s$를
$\widetilde{j}$에 의한 그 상이라 하자. 그러면 $s$는
$(n'\otimes1)/f^p$와 동일시된다. 한편
$t\in\Gamma(D(f),\widetilde{A})$가 $g/f^q$($g\in A$)와
동일시된다고 하자. 정의에 따라
\[
  h_x(s'_{x'}\otimes t_x)=t_x\cdot s_x
\]
이다(\hyperref[0.4.4.3-ko]{0, 4.4.3}). 그런데
$N_f$는 표준적으로
\[
  N'_{f'}\otimes_{A'_{f'}}(A_f)_{[\varphi_{f'}]}
\]
와 동일시된다(\hyperref[0.1.5.4-ko]{0, 1.5.4}). 이때 $s$는
$(n'/{f'}^p)\otimes1$에 대응하고, 절단 $y\mapsto t_y\cdot s_y$는
$(n'/{f'}^p)\otimes(g/f^q)$에 대응한다.
(\hyperref[0.1.5.6-ko]{0, 1.5.6})의 양립 그림들은 $h_x$가 바로
표준 동형
\[
  \label{I.1.6.5.1-ko}
  N'_{x'}\otimes_{A'_{x'}}(A_x)_{[\varphi_x]}
  \xrightarrow{\sim}N_x=(N'\otimes_{A'}A_{[\varphi]})_x
  \tag{1.6.5.1}
\]
임을 보여 준다.

또한 $v:N'_1\to N'_2$가 $A'$-가군 준동형이면 모든 $x'\in X'$에
대하여 $\widetilde{v}_{x'}=v_{x'}$이다. 그러므로 앞의 결과에서 곧바로
$\Phi^*(\widetilde{v})$가 $(v\otimes1)^{\sim}$와 표준적으로
동일시됨을 얻으며, 이로써
(\hyperref[I.1.6.5-ko]{1.6.5})의 증명이 끝난다.

$B'$가 $A'$-대수이면 $\Phi^*(\widetilde{B'})$에서
$(B'\otimes_{A'}A_{[\varphi]})^{\sim}$로 가는 표준 동형은
$\mathscr{O}_X$-대수의 동형이다. 더욱이 $N'$이 $B'$-가군이면
$\Phi^*(\widetilde{N'})$에서
$(N'\otimes_{A'}A_{[\varphi]})^{\sim}$로 가는 표준 동형은
$\Phi^*(\widetilde{B'})$-가군의 동형이다.

\begin{corollary}[1.6.6]
\label{I.1.6.6-ko}
$s'$이 $A'$-가군 $\Gamma(\widetilde{N'})$을 달릴 때, 절단 $s'$의
$\Phi^*(\widetilde{N'})$ 안의 표준 상들은 $A$-가군
$\Gamma(X,\Phi^*(\widetilde{N'}))$을 생성한다.
\end{corollary}

실제로 $N'$과 $N$을 각각 $\Gamma(\widetilde{N'})$과
$\Gamma(\widetilde{N})$와 동일시하면
(\hyperref[I.1.3.7-ko]{1.3.7}), 이 상들은 $z'$이 $N'$을 달릴 때의
$N$의 원소 $z'\otimes1$과 동일시된다.

\begin{env}[1.6.7]
\label{I.1.6.7-ko}
(\hyperref[I.1.6.5-ko]{1.6.5})의 증명에서 표준 사상
(\hyperref[0.4.4.3.2-ko]{0, 4.4.3.2})
$\rho:\widetilde{N'}\to\Phi_*(\Phi^*(\widetilde{N'}))$이 바로
준동형 $\widetilde{j}$임을 함께 증명하였다.
\oldpage[I]{96}
여기서 $j:N'\to N'\otimes_{A'}A_{[\varphi]}$는
$z'\mapsto z'\otimes1$인 준동형이다. 마찬가지로 표준 사상
(\hyperref[0.4.4.3.3-ko]{0, 4.4.3.3})
$\sigma:\Phi^*(\Phi_*(\widetilde{M}))\to\widetilde{M}$은
$\widetilde{p}$와 같다. 여기서
$p:M_{[\varphi]}\otimes_{A'}A_{[\varphi]}\to M$은 모든 텐서곱
$z\otimes a$($z\in M$, $a\in A$)에 $a\cdot z$를 대응시키는 표준
준동형이다. 이는 정의
(\hyperref[0.3.7.1-ko]{0, 3.7.1},
\hyperref[0.4.4.3-ko]{0, 4.4.3} 및
\hyperref[I.1.3.7-ko]{I, 1.3.7})에서 곧바로 따른다.

따라서 (\hyperref[0.4.4.3-ko]{0, 4.4.3})과
(\hyperref[0.3.5.4.4-ko]{3.5.4.4})에 의해
$v:N'\to M_{[\varphi]}$가 $A'$-준동형이면
$\widetilde{v}^{\#}=(v\otimes1)^{\sim}$이다.
\end{env}

\begin{env}[1.6.8]
\label{I.1.6.8-ko}
$N'_1$, $N'_2$를 두 $A'$-가군이라 하고 $N'_1$이 \emph{유한 표시}를
갖는다고 하자. 그러면
(\hyperref[I.1.6.7-ko]{1.6.7})과
(\hyperref[I.1.3.12-ko]{1.3.12, (ii)})에 따라 표준 준동형
(\hyperref[0.4.4.6-ko]{0, 4.4.6})
\[
  \Phi^*\bigl(\mathscr{H}\!om_{\widetilde{A'}}
    (\widetilde{N'_1},\widetilde{N'_2})\bigr)
  \longrightarrow
  \mathscr{H}\!om_{\widetilde A}
    \bigl(\Phi^*(\widetilde{N'_1}),\Phi^*(\widetilde{N'_2})\bigr)
\]
은 $\widetilde{\gamma}$와 같다. 여기서 $\gamma$는 $A$-가군의
표준 준동형
\[
  \operatorname{Hom}_{A'}(N'_1,N'_2)\otimes_{A'}A
  \longrightarrow
  \operatorname{Hom}_A(N'_1\otimes_{A'}A,N'_2\otimes_{A'}A)
\]
이다.
\end{env}

\begin{env}[1.6.9]
\label{I.1.6.9-ko}
$\mathfrak{J}'$을 $A'$의 아이디얼, $M$을 $A$-가군이라 하자. 정의에
따라 $\widetilde{\mathfrak{J}'}\widetilde M$은 표준 준동형
$\Phi^*(\widetilde{\mathfrak{J}'})
\otimes_{\widetilde A}\widetilde M\to\widetilde M$의 상이다. 따라서
(\hyperref[I.1.6.5-ko]{1.6.5})와
(\hyperref[I.1.3.12-ko]{1.3.12, (i)})에 의해
$\widetilde{\mathfrak{J}'}\widetilde M$은
$(\mathfrak{J}'M)^{\sim}$와 표준적으로 동일시된다. 특히
$\Phi^*(\widetilde{\mathfrak{J}'})\widetilde A$는
$(\mathfrak{J}'A)^{\sim}$와 동일시된다. 또한 함자 $\Phi^*$의 오른쪽
완전성을 고려하면 $\widetilde A$-대수
$\Phi^*((A'/\mathfrak{J}')^{\sim})$는
$(A/\mathfrak{J}'A)^{\sim}$와 동일시된다.
\end{env}

\begin{env}[1.6.10]
\label{I.1.6.10-ko}
$A''$을 세 번째 환, $\varphi'$을 준동형 $A''\to A'$이라 하고
$\varphi''=\varphi\circ\varphi'$로 놓자. 정의에서 곧바로
\[
  {}^a\varphi''=({}^a\varphi')\circ({}^a\varphi),
  \qquad
  \widetilde{\varphi''}=\widetilde\varphi\circ\widetilde{\varphi'}
\]
를 얻는다(\hyperref[0.1.5.7-ko]{0, 1.5.7}). 따라서
$\Phi''=\Phi'\circ\Phi$이다. 다시 말해
$(\operatorname{Spec}A,\widetilde A)$는 $A$에 관하여 환의 범주에서
환 달린 공간의 범주로 가는 \emph{반변 함자}이다.
\end{env}

\subsection{아핀 스킴의 사상의 특징짓기}
\label{subsection:I.1.7-ko}

\begin{definition}[1.7.1]
\label{I.1.7.1-ko}
환 달린 공간 $(X,\mathscr{O}_X)$가 어떤 환 $A$에 대한
$(\operatorname{Spec}(A),\widetilde A)$ 꼴의 환 달린 공간과 동형이면
이를 \emph{아핀 스킴}이라 한다. 이때 $\Gamma(X,\mathscr{O}_X)$는
환 $A$와 표준적으로 동일시되며
(\hyperref[I.1.3.7-ko]{1.3.7}), 이를 아핀 스킴
$(X,\mathscr{O}_X)$의 환이라 한다. 혼동의 염려가 없을 때에는
$A(X)$로 쓴다.
\end{definition}

앞으로 \emph{아핀 스킴} $\operatorname{Spec}(A)$라고 하면 항상 환 달린
공간 $(\operatorname{Spec}(A),\widetilde A)$를 뜻한다.

\begin{env}[1.7.2]
\label{I.1.7.2-ko}
$A$, $B$를 두 환, $(X,\mathscr{O}_X)$와 $(Y,\mathscr{O}_Y)$를 각각
소 스펙트럼 $X=\operatorname{Spec}(A)$와 $Y=\operatorname{Spec}(B)$에
대응하는 아핀 스킴이라 하자. 앞서
(\hyperref[I.1.6.1-ko]{1.6.1})에서 모든 환 준동형
$\varphi:B\to A$에 사상
\[
  \Phi=({}^a\varphi,\widetilde\varphi)
  =\operatorname{Spec}(\varphi):
  (X,\mathscr{O}_X)\to(Y,\mathscr{O}_Y)
\]
이 대응함을 보았다. $\varphi$는 $\Phi$에 의해 완전히 결정된다. 실제로
정의상
\[
  \varphi=\Gamma(\widetilde\varphi):
  \Gamma(\widetilde B)\longrightarrow
  \Gamma({}^a\varphi_*(\widetilde A))=\Gamma(\widetilde A)
\]
이다.
\end{env}

\begin{theorem}[1.7.3]
\label{I.1.7.3-ko}
$(X,\mathscr{O}_X)$와 $(Y,\mathscr{O}_Y)$를 두 아핀 스킴이라 하자.
환 달린 공간의 사상
$(\psi,\theta):(X,\mathscr{O}_X)\to(Y,\mathscr{O}_Y)$가 어떤 환 준동형
$\varphi:A(Y)\to A(X)$에 대하여
$({}^a\varphi,\widetilde\varphi)$ 꼴이기 위한 필요충분조건은 모든
$x\in X$에 대하여
\[
  \theta_x^{\#}:\mathscr{O}_{\psi(x)}\to\mathscr{O}_x
\]
가 국소 준동형인 것이다.
\end{theorem}

\oldpage[I]{97}
$A=A(X)$, $B=A(Y)$로 놓자. 이 조건은 필요하다. 실제로
(\hyperref[I.1.6.1-ko]{1.6.1})에서
$\widetilde\varphi_x^{\#}$이 $\varphi$에서 표준적으로 유도되는 준동형
$B_{{}^a\varphi(x)}\to A_x$임을 보았다. 정의상
${}^a\varphi(x)=\varphi^{-1}(\mathfrak{j}_x)$이므로 이 준동형은 국소
준동형이다.

이제 조건이 충분함을 증명하자. 정의상 $\theta$는 준동형
$\mathscr{O}_Y\to\psi_*(\mathscr{O}_X)$이고, 여기서 표준적으로 환 준동형
\[
  \varphi=\Gamma(\theta):
  B=\Gamma(Y,\mathscr{O}_Y)
  \longrightarrow
  \Gamma(Y,\psi_*(\mathscr{O}_X))
  =\Gamma(X,\mathscr{O}_X)=A
\]
을 얻는다.

$\theta_x^{\#}$에 대한 가정으로부터 몫을 취하여 잉여체
$k(\psi(x))$에서 잉여체 $k(x)$로 가는 단사 준동형 $\theta^x$를 얻는다.
모든 절단 $f\in\Gamma(Y,\mathscr{O}_Y)=B$에 대하여
$\theta^x(f(\psi(x)))=\varphi(f)(x)$이다. 따라서 관계
$f(\psi(x))=0$과 $\varphi(f)(x)=0$은 동치이다. 이는
$\mathfrak{j}_{\psi(x)}=\mathfrak{j}_{{}^a\varphi(x)}$라는 뜻이므로 모든
$x\in X$에 대하여 $\psi(x)={}^a\varphi(x)$, 곧
$\psi={}^a\varphi$이다. 또한 그림
\phantomsection\label{I.1.7.3.diagram-ko}
\[
  \xymatrix{
    B=\Gamma(Y,\mathscr{O}_Y)\ar[r]^\varphi\ar[d] &
    \Gamma(X,\mathscr{O}_X)=A\ar[d]\\
    B_{\psi(x)}\ar[r]_{\theta_x^{\#}} & A_x
  }
\]
은 가환이다(\hyperref[0.3.7.2-ko]{0, 3.7.2}). 이는
$\theta_x^{\#}$이 $\varphi$에서 표준적으로 유도되는 준동형
$\varphi_x:B_{\psi(x)}\to A_x$와 같다는 뜻이다
(\hyperref[0.1.5.1-ko]{0, 1.5.1}). $\theta_x^{\#}$들의 자료가
$\theta^{\#}$, 따라서 $\theta$도 완전히 결정하므로
(\hyperref[0.3.7.1-ko]{0, 3.7.1}),
$\widetilde\varphi$의 정의
(\hyperref[I.1.6.1-ko]{1.6.1})에 의해
$\theta=\widetilde\varphi$이다.

(\hyperref[I.1.7.3-ko]{1.7.3})의 조건을 만족하는 환 달린 공간의 사상
$(\psi,\theta)$를 \emph{아핀 스킴의 사상}이라 하겠다.

\begin{corollary}[1.7.4]
\label{I.1.7.4-ko}
$(X,\mathscr{O}_X)$와 $(Y,\mathscr{O}_Y)$가 두 아핀 스킴이면 아핀
스킴의 사상들의 집합과
$\operatorname{Hom}_{\mathrm{Ring}}(B,A)$ 사이에 표준 전단사가
존재한다. 여기서 $A=\Gamma(\mathscr{O}_X)$이고
$B=\Gamma(\mathscr{O}_Y)$이다.
\end{corollary}

달리 말하면, $A$에 관한 $(\operatorname{Spec}(A),\widetilde A)$와
$(X,\mathscr{O}_X)$에 관한 $\Gamma(X,\mathscr{O}_X)$는 가환환의 범주와
아핀 스킴의 범주의 반대 범주 사이의 \emph{동치}를 정의한다
(T, I, 1.2).

\begin{corollary}[1.7.5]
\label{I.1.7.5-ko}
$\varphi:B\to A$가 전사이면 이에 대응하는 사상
$({}^a\varphi,\widetilde\varphi)$는 환 달린 공간의 단사사상이다
\emph{(4.1.7 참조)}.
\end{corollary}

실제로 ${}^a\varphi$는 단사이다
(\hyperref[I.1.2.5-ko]{1.2.5}). 또한 $\varphi$가 전사이므로 모든
$x\in X$에 대하여 $\varphi$에서 분수환을 취하여 얻는
$\varphi_x^{\#}:B_{{}^a\varphi(x)}\to A_x$도 전사이다
(\hyperref[0.1.5.1-ko]{0, 1.5.1}). 따라서 결론은
(\hyperref[0.4.1.1-ko]{0, 4.1.1})에서 따른다.
% ===== END EXACT INLINE INPUT: c1s1.tex =====


% ===== BEGIN EXACT INLINE INPUT: c1s2.tex =====
% Authority: EGA I ega1-2-fr.tex, whole-file SHA-256 AE6B128092ACBB8C1AFB4899EEA003FB966B6FF6669A264B59FD5F095AF4F029.
% Sequential coverage in this file: authority lines 1-296, complete subsections 2.1-2.2.
\section{준스킴과 준스킴의 사상}
\label{section:I.2-ko}

\subsection{준스킴의 정의}
\label{subsection:I.2.1-ko}

\begin{env}[2.1.1]
\label{I.2.1.1-ko}
환 달린 공간 $(X,\mathscr{O}_X)$가 주어졌을 때, $X$의 열린 부분집합
$V$에 유도된 환 달린 공간 $(V,\mathscr{O}_X|V)$가 아핀 스킴이면
$V$를 \emph{아핀 열린집합}이라 한다
(\hyperref[I.1.7.1-ko]{1.7.1}).
\end{env}

\begin{definition}[2.1.2]
\label{I.2.1.2-ko}
모든 점이 아핀 열린 근방을 갖는 환 달린 공간 $(X,\mathscr{O}_X)$를
\emph{준스킴}이라 한다.
\end{definition}

\oldpage[I]{98}

\begin{proposition}[2.1.3]
\label{I.2.1.3-ko}
$(X,\mathscr{O}_X)$가 준스킴이면 아핀 열린집합들은 $X$의 위상의
기저를 이룬다.
\end{proposition}

실제로 $V$가 $x\in X$의 임의의 열린 근방이면, 가정에 따라 $x$의
어떤 열린 근방 $W$에 대하여 $(W,\mathscr{O}_X|W)$가 아핀 스킴이다.
그 환을 $A$라 하자. 공간 $W$ 안에서 $V\cap W$는 $x$의 열린
근방이므로, $x$의 열린 근방이고 $V\cap W$에 포함되는 $D(f)$가
존재하도록 하는 $f\in A$가 있다
(\hyperref[I.1.1.10-ko]{1.1.10 (i)}). 이때 환 달린 공간
$(D(f),\mathscr{O}_X|D(f))$는 $A_f$와 동형인 환을 갖는 아핀
스킴이므로(\hyperref[I.1.3.6-ko]{1.3.6}), 명제가 따른다.

\begin{proposition}[2.1.4]
\label{I.2.1.4-ko}
준스킴의 바탕공간은 콜모고로프 공간이다.
\end{proposition}

실제로 준스킴 $X$의 서로 다른 두 점 $x,y$가 하나의 아핀
열린집합에 함께 들어 있지 않으면, 두 점 가운데 하나를 포함하고
다른 하나를 포함하지 않는 열린 근방이 존재함은 자명하다. 두 점이
하나의 아핀 열린집합에 함께 들어 있으면
(\hyperref[I.1.1.8-ko]{1.1.8})에서 결론이 따른다.

\begin{proposition}[2.1.5]
\label{I.2.1.5-ko}
$(X,\mathscr{O}_X)$가 준스킴이면 $X$의 모든 기약 닫힌 부분집합은
유일한 일반점을 갖는다. 따라서 사상
\[
  x\longmapsto\overline{\{x\}}
\]
은 $X$에서 그 기약 닫힌 부분집합들의 집합으로 가는 전단사이다.
\end{proposition}

실제로 $Y$를 $X$의 기약 닫힌 부분집합, $y\in Y$, $U$를 $X$에서
$y$의 아핀 열린 근방이라 하자. $U\cap Y$는 $Y$에서 조밀하고
기약이다
(\hyperref[0.2.1.1-ko]{0, 2.1.1} 및
\hyperref[0.2.1.4-ko]{2.1.4}). 따라서
(\hyperref[I.1.1.14-ko]{1.1.14})에 의해 $U\cap Y$는 어떤 점 $x$의
$U$ 안에서의 폐포이다. 그러므로 $Y\subset\overline U$는 $X$에서
$x$의 폐포이다. $Y$의 일반점의 유일성은
(\hyperref[I.2.1.4-ko]{2.1.4})와
(\hyperref[0.2.1.3-ko]{0, 2.1.3})에서 따른다.

\begin{env}[2.1.6]
\label{I.2.1.6-ko}
$Y$가 $X$의 기약 닫힌 부분집합이고 $y$가 그 일반점이면, 국소환
$\mathscr{O}_y$를 $\mathscr{O}_{X/Y}$라고도 쓰며 이를
\emph{$Y$를 따라 취한 $X$의 국소환}, 또는 \emph{$X$ 안의 $Y$의
국소환}이라 한다.

$X$ 자체가 기약이고 $x$가 그 일반점이면, $\mathscr{O}_x$를
\emph{$X$ 위의 유리함수환}이라고도 한다(§~7 참조).
\end{env}

\begin{proposition}[2.1.7]
\label{I.2.1.7-ko}
$(X,\mathscr{O}_X)$가 준스킴이면 $X$의 모든 열린 부분집합 $U$에
대하여 환 달린 공간 $(U,\mathscr{O}_X|U)$도 준스킴이다.
\end{proposition}

이는 정의 (\hyperref[I.2.1.2-ko]{2.1.2})와 명제
(\hyperref[I.2.1.3-ko]{2.1.3})에서 곧바로 따른다.

$(U,\mathscr{O}_X|U)$를 $(X,\mathscr{O}_X)$가 $U$ 위에
\emph{유도하는} 준스킴, 또는 $(X,\mathscr{O}_X)$의 $U$에 대한
\emph{제한}이라 한다.

\begin{env}[2.1.8]
\label{I.2.1.8-ko}
준스킴 $(X,\mathscr{O}_X)$의 바탕공간 $X$가 기약(각각 연결)이면
이 준스킴을 \emph{기약}(각각 \emph{연결})이라 한다. 준스킴이
\emph{기약이고 축소}이면 이를 \emph{정역 준스킴}이라 한다
\emph{(5.1.4 참조)}. 준스킴 $(X,\mathscr{O}_X)$의 모든 점
$x\in X$가, 그 위에 유도된 준스킴이 정역 준스킴인 어떤 열린 근방
$U$를 가지면 $(X,\mathscr{O}_X)$를 \emph{국소 정역 준스킴}이라
한다.
\end{env}

\subsection{준스킴의 사상}
\label{subsection:I.2.2-ko}
\label{I.2.2-ko}

\begin{definition}[2.2.1]
\label{I.2.2.1-ko}
두 준스킴 $(X,\mathscr{O}_X)$, $(Y,\mathscr{O}_Y)$가 주어졌다고 하자.
$(X,\mathscr{O}_X)$에서 $(Y,\mathscr{O}_Y)$로 가는 환 달린 공간의
사상 $(\psi,\theta)$ 가운데, 모든 $x\in X$에 대하여
\[
  \theta_x^\sharp:\mathscr{O}_{\psi(x)}\longrightarrow\mathscr{O}_x
\]
가 국소 준동형인 것을 \emph{준스킴의 사상}이라 한다.
\end{definition}

\oldpage[I]{99}

따라서 몫을 취하면
$\theta_x^\sharp:\mathscr{O}_{\psi(x)}\to\mathscr{O}_x$는 단사 준동형
$\theta^x:k(\psi(x))\to k(x)$를 유도하며, 이에 따라 $k(x)$를 체
$k(\psi(x))$의 \emph{확대체}로 볼 수 있다.

\begin{env}[2.2.2]
\label{I.2.2.2-ko}
준스킴의 두 사상 $(\psi,\theta)$와 $(\psi',\theta')$의 합성
$(\psi'',\theta'')$도 준스킴의 사상이다. 이는 공식
${\theta''}^\sharp=\theta^\sharp\circ\psi^*({\theta'}^\sharp)$
(\hyperref[0.3.5.5-ko]{0, 3.5.5})에서 따른다. 따라서 준스킴들은 하나의
\emph{범주}를 이룬다. 일반 표기에 따라 준스킴 $X$에서 준스킴 $Y$로
가는 사상들의 집합을 $\operatorname{Hom}(X,Y)$로 쓴다.
\end{env}

\begin{example}[2.2.3]
\label{I.2.2.3-ko}
$U$가 $X$의 열린 부분집합이면, 유도된 준스킴
$(U,\mathscr{O}_X|U)$에서 $(X,\mathscr{O}_X)$로 가는 표준 단사사상
(\hyperref[0.4.1.2-ko]{0, 4.1.2})은 준스킴의 사상이다. 더구나 이는
환 달린 공간의 단사사상이며(따라서 특히 준스킴의 단사사상이며),
이는 (\hyperref[0.4.1.1-ko]{0, 4.1.1})에서 곧바로 따른다.
\end{example}

\begin{proposition}[2.2.4]
\label{I.2.2.4-ko}
$(X,\mathscr{O}_X)$를 준스킴, $(S,\mathscr{O}_S)$를 환 $A$의 아핀
스킴이라 하자. 준스킴 $(X,\mathscr{O}_X)$에서 준스킴
$(S,\mathscr{O}_S)$로 가는 사상들과 $A$에서 환
$\Gamma(X,\mathscr{O}_X)$로 가는 준동형들 사이에는 표준 전단사 대응이
존재한다.
\end{proposition}

먼저 $(X,\mathscr{O}_X)$와 $(Y,\mathscr{O}_Y)$가 임의의 두 환 달린
공간이면, $(X,\mathscr{O}_X)$에서 $(Y,\mathscr{O}_Y)$로 가는 사상
$(\psi,\theta)$는 표준적으로 환 준동형
\[
  \Gamma(\theta):\Gamma(Y,\mathscr{O}_Y)\longrightarrow
  \Gamma(Y,\psi_*(\mathscr{O}_X))=\Gamma(X,\mathscr{O}_X)
\]
을 정의함을 주목하자. 지금의 경우에는 임의의 준동형
$\varphi:A\to\Gamma(X,\mathscr{O}_X)$가 유일한 $\theta$에 대하여
$\Gamma(\theta)$ 꼴임을 보이면 충분하다. 가정에 따라 $X$는 아핀
열린집합들 $(V_\alpha)$로 덮인다. $\varphi$를 제한 준동형
\[
  \Gamma(X,\mathscr{O}_X)\longrightarrow
  \Gamma(V_\alpha,\mathscr{O}_X|V_\alpha)
\]
과 합성하면 준동형
$\varphi_\alpha:A\to\Gamma(V_\alpha,\mathscr{O}_X|V_\alpha)$를 얻는다.
(\hyperref[I.1.7.3-ko]{1.7.3})에 따라 이는 준스킴
$(V_\alpha,\mathscr{O}_X|V_\alpha)$에서 $(S,\mathscr{O}_S)$로 가는
유일한 사상 $(\psi_\alpha,\theta_\alpha)$에 대응한다. 또한 모든 지표쌍
$(\alpha,\beta)$와 모든 점 $x\in V_\alpha\cap V_\beta$에 대하여,
$x$를 포함하고 $V_\alpha\cap V_\beta$에 포함되는 아핀 열린집합 $W$가
존재한다(\hyperref[I.2.1.3-ko]{2.1.3}). $\varphi_\alpha$와
$\varphi_\beta$를 $W$로 가는 제한 준동형과 합성하면 같은 준동형
\[
  \Gamma(S,\mathscr{O}_S)\longrightarrow
  \Gamma(W,\mathscr{O}_X|W)
\]
을 얻는다. 따라서 모든 $x\in V_\alpha$와 모든 $\alpha$에 대한 관계
$(\theta_\alpha^\sharp)_x=(\varphi_\alpha)_x$
(\hyperref[I.1.6.1-ko]{1.6.1})와 함께 보면,
$(\psi_\alpha,\theta_\alpha)$와 $(\psi_\beta,\theta_\beta)$의 $W$로의
제한은 일치한다. 그러므로 각 $V_\alpha$로의 제한이
$(\psi_\alpha,\theta_\alpha)$인 환 달린 공간의 사상
\[
  (\psi,\theta):(X,\mathscr{O}_X)\longrightarrow(S,\mathscr{O}_S)
\]
이 유일하게 존재한다. 이 사상이 준스킴의 사상이며
$\Gamma(\theta)=\varphi$임은 명백하다.

$u:A\to\Gamma(X,\mathscr{O}_X)$를 환 준동형이라 하고,
$v=(\psi,\theta)$를 이에 대응하는 사상
$(X,\mathscr{O}_X)\to(S,\mathscr{O}_S)$라 하자. 모든 $f\in A$에 대하여
\begin{equation}
  \psi^{-1}(D(f))=X_{u(f)}
  \tag{2.2.4.1}\label{I.2.2.4.1-ko}
\end{equation}
이다. 여기서 표기는 국소 자유층 $\mathscr{O}_X$에 관한
(\hyperref[0.5.5.2-ko]{0, 5.5.2})의 것을 사용한다. 실제로 이 공식을
$X$ 자체가 아핀인 경우에 확인하면 충분하며, 그때에는 바로
(\hyperref[I.1.2.2.2-ko]{1.2.2.2})이다.

\begin{proposition}[2.2.5]
\label{I.2.2.5-ko}
(\hyperref[I.2.2.4-ko]{2.2.4})의 가정 아래에서, 환 준동형
$\varphi:A\to\Gamma(X,\mathscr{O}_X)$와 이에 대응하는 준스킴의 사상
$f:(X,\mathscr{O}_X)\to(S,\mathscr{O}_S)$를 생각하자.
$\mathscr{G}$를 준연접 $\mathscr{O}_X$-가군,
$\mathscr{F}$를 준연접 $\mathscr{O}_S$-가군이라 하고
$M=\Gamma(S,\mathscr{F})$로 놓자.
\oldpage[I]{100}
그러면 $f$-사상
$\mathscr{F}\to\mathscr{G}$
(\hyperref[0.4.4.1-ko]{0, 4.4.1})들과 $A$-가군 준동형
\[
  M\longrightarrow(\Gamma(X,\mathscr{G}))_{[\varphi]}
\]
들 사이에는 표준 전단사 대응이 존재한다.
\end{proposition}

실제로 (\hyperref[I.2.2.4-ko]{2.2.4})에서와 같이 논하면 곧 $X$가
아핀인 경우로 환원되며, 그 경우 명제는
(\hyperref[I.1.6.3-ko]{1.6.3})과
(\hyperref[I.1.3.8-ko]{1.3.8})에서 따른다.

\begin{env}[2.2.6]
\label{I.2.2.6-ko}
준스킴의 사상
$(\psi,\theta):(X,\mathscr{O}_X)\to(Y,\mathscr{O}_Y)$가 $X$의 모든
열린 부분집합 $U$에 대하여 $\psi(U)$가 $Y$에서 열려 있으면 이 사상을
\emph{열린 사상}이라 하고, $X$의 모든 닫힌 부분집합 $F$에 대하여
$\psi(F)$가 $Y$에서 닫혀 있으면 \emph{닫힌 사상}이라 한다.
$\psi(X)$가 $Y$에서 조밀하면 \emph{우세 사상}, $\psi$가 전사이면
\emph{전사 사상}이라 한다. 이 조건들은 연속사상 $\psi$에만 관련됨을
주목하자.
\end{env}

\begin{proposition}[2.2.7]
\label{I.2.2.7-ko}
준스킴의 두 사상
\[
  f=(\psi,\theta):(X,\mathscr{O}_X)\to(Y,\mathscr{O}_Y),
  \qquad
  g=(\psi',\theta'):(Y,\mathscr{O}_Y)\to(Z,\mathscr{O}_Z)
\]
를 생각하자.
\begin{enumerate}
  \item[\rm(i)] $f$와 $g$가 모두 열린(각각 닫힌, 우세, 전사) 사상이면
    $g\circ f$도 그러하다.
  \item[\rm(ii)] $f$가 전사이고 $g\circ f$가 닫힌 사상이면 $g$도 닫힌
    사상이다.
  \item[\rm(iii)] $g\circ f$가 전사이면 $g$도 전사이다.
\end{enumerate}
\end{proposition}

(i)와 (iii)은 명백하다. $g\circ f=(\psi'',\theta'')$로 놓자. $F$가
$Y$에서 닫혀 있으면 $\psi^{-1}(F)$는 $X$에서 닫혀 있으므로
$\psi''(\psi^{-1}(F))$는 $Z$에서 닫혀 있다. 그런데 $\psi$가 전사이므로
$\psi(\psi^{-1}(F))=F$이고, 따라서
$\psi''(\psi^{-1}(F))=\psi'(F)$이다. 이로써 (ii)가 증명된다.

\begin{proposition}[2.2.8]
\label{I.2.2.8-ko}
$f=(\psi,\theta)$를
$(X,\mathscr{O}_X)\to(Y,\mathscr{O}_Y)$인 사상이라 하고,
$(U_\alpha)$를 $Y$의 열린 덮개라 하자. $f$가 열린(각각 닫힌, 전사,
우세) 사상이기 위한 필요충분조건은, 각각의 유도된 준스킴
$(\psi^{-1}(U_\alpha),\mathscr{O}_X|\psi^{-1}(U_\alpha))$에 대한 $f$의
제한을 이 준스킴에서 유도된 준스킴
$(U_\alpha,\mathscr{O}_Y|U_\alpha)$로 가는 사상으로 보았을 때, 그
제한이 열린(각각 닫힌, 전사, 우세) 사상인 것이다.
\end{proposition}

이 명제는 정의에서 곧바로 따른다. 여기서는 $Y$의 부분집합 $F$가
$Y$에서 닫힌(각각 열린, 조밀한) 부분집합이기 위한 필요충분조건이
모든 $\alpha$에 대하여 $F\cap U_\alpha$가 $U_\alpha$에서 닫힌(각각
열린, 조밀한) 부분집합인 것임을 사용한다.

\begin{env}[2.2.9]
\label{I.2.2.9-ko}
$(X,\mathscr{O}_X)$와 $(Y,\mathscr{O}_Y)$를 두 준스킴이라 하고, $X$와
$Y$가 같은 유한 개수의 기약 성분 $X_i$와 $Y_i$
$(1\leq i\leq n)$를 갖는다고 하자. $X_i$와 $Y_i$의 일반점을 각각
$\xi_i$와 $\eta_i$라 하자
(\hyperref[I.2.1.5-ko]{2.1.5}). 사상
\[
  f=(\psi,\theta):(X,\mathscr{O}_X)\longrightarrow(Y,\mathscr{O}_Y)
\]
이 모든 $i$에 대하여 $\psi^{-1}(\eta_i)=\{\xi_i\}$를 만족하고
\[
  \theta_{\xi_i}^{\sharp}:\mathscr{O}_{\eta_i}\longrightarrow
  \mathscr{O}_{\xi_i}
\]
가 동형이면 $f$를 \emph{쌍유리 사상}이라 한다. 쌍유리 사상이 우세
사상임은 명백하다(\hyperref[0.2.1.8-ko]{0, 2.1.8}). 따라서 닫힌
사상이기도 하면 전사이다.
\end{env}

\begin{notation}[2.2.10]
\label{I.2.2.10-ko}
이 책의 앞으로의 모든 부분에서는, 혼동의 염려가 없을 때 준스킴의
표기에서는 구조층을, 사상의 표기에서는 구조층의 준동형을
\emph{생략}한다. 준스킴 $X$의 바탕공간의 열린 부분집합 $U$를
준스킴이라고 할 때에는 언제나 $U$ 위에 유도된 준스킴을 뜻한다.
\end{notation}

\subsection{준스킴의 붙이기}
\label{subsection:I.2.3-ko}

\begin{env}[2.3.1]
\label{I.2.3.1-ko}
\oldpage[I]{101}
정의 (\hyperref[I.2.1.2-ko]{2.1.2})에 따라, 준스킴들을
\emph{붙여서} 얻은 모든 환 달린 공간
(\hyperref[0.4.1.6-ko]{0, 4.1.6})은 다시 준스킴이다. 특히 모든
준스킴은 정의상 아핀 열린집합들로 이루어진 덮개를 가지므로, 모든
준스킴은 \emph{아핀 스킴들을 붙여서} 얻을 수 있다.
\end{env}

\begin{example}[2.3.2]
\label{I.2.3.2-ko}
$K$를 체라 하고, $B=K[s]$, $C=K[t]$를 각각 $K$ 위의 하나의 부정원에
대한 다항식환이라 하자. 또한 $X_1=\operatorname{Spec}(B)$,
$X_2=\operatorname{Spec}(C)$라 하자. 이들은 서로 동형인 두 아핀
스킴이다. $X_1$에서(각각 $X_2$에서) $U_{12}$를(각각 $U_{21}$을)
아핀 열린집합 $D(s)$로(각각 $D(t)$로) 잡자. 그 환 $B_s$는(각각
$C_t$는) $f\in B$인(각각 $g\in C$인) 유리분수 $f(s)/s^m$의
꼴로(각각 $g(t)/t^n$의 꼴로) 이루어진다. $B_s$에서 $C_t$로 가며
$f(s)/s^m$을 유리분수 $f(1/t)/(1/t^m)$에 대응시키는 동형에
(\hyperref[I.2.2.4-ko]{2.2.4})에 따라 대응하는 준스킴의 동형
\[
  u_{12}:U_{21}\longrightarrow U_{12}
\]
를 생각하자. 붙이기 조건이 전혀 없음은 명백하므로, $u_{12}$를
사용하여 $U_{12}$와 $U_{21}$을 따라 $X_1$과 $X_2$를 붙일 수 있다.
이렇게 얻은 준스킴 $X$는 뒤에서 일반적인 구성법의 특수한 경우로
다시 만나게 된다(II, 2.4.3). 여기서는 $X$가 \emph{아핀 스킴이
아님}만을 보이자. 실제로 환 $\Gamma(X,\mathscr{O}_X)$는 $K$와
\emph{동형}이므로 그 스펙트럼은 한 점으로 이루어진다. $X$ 위의
$\mathscr{O}_X$의 한 단면을 $X$의 아핀 열린집합과 동일시되는
$X_1$에(각각 $X_2$에) 제한하면 다항식 $f(s)$를(각각 $g(t)$를)
얻는다. 정의에 따라 $g(t)=f(1/t)$이어야 하며, 이는
$f=g\in K$일 때에만 가능하다.
\end{example}

\subsection{국소 스킴}
\label{subsection:I.2.4-ko}

\begin{env}[2.4.1]
\label{I.2.4.1-ko}
환 $A$가 국소환인 아핀 스킴을 \emph{국소 스킴}이라 한다. 그러면
$X=\operatorname{Spec}(A)$에는 유일한 \emph{닫힌점} $a$가 존재하며,
다른 모든 점 $b\in X$에 대하여
$a\in\overline{\{b\}}$이다
(\hyperref[I.1.1.7-ko]{1.1.7}).
\end{env}

준스킴 $Y$의 점 $y$마다 국소 스킴
$\operatorname{Spec}(\mathscr{O}_y)$를 \emph{$Y$의 점 $y$에서의 국소
스킴}이라 한다. $V$를 $y$를 포함하는 $Y$의 아핀 열린집합, $B$를 아핀
스킴 $V$의 환이라 하자. $\mathscr{O}_y$는 $B_y$와 표준적으로
동일시되고(\hyperref[I.1.3.4-ko]{1.3.4}), 따라서 표준 준동형
$B\to B_y$는 (\hyperref[I.1.6.1-ko]{1.6.1})에 따라 준스킴의 사상
$\operatorname{Spec}(\mathscr{O}_y)\to V$에 대응한다. 이 사상을 표준
단사사상 $V\to Y$와 합성하면 사상
$\operatorname{Spec}(\mathscr{O}_y)\to Y$를 얻으며, 이는 선택한
$y$의 아핀 열린 근방 $V$에 \emph{무관하다}. 실제로 $V'$이 $y$를
포함하는 또 다른 아핀 열린집합이면, $y$를 포함하고
$W\subset V\cap V'$인 세 번째 아핀 열린집합 $W$가 존재한다
(\hyperref[I.2.1.3-ko]{2.1.3}). 따라서 $V\subset V'$인 경우만
생각하면 충분하다. $B'$을 $V'$의 환이라 하면, 이는 도식
\phantomsection\label{I.2.4.1.diagram-ko}
\[
  \xymatrix{
    B'\ar[rr]\ar[dr] & & B\ar[dl]\\
    & \mathscr{O}_y
  }
\]
이 가환임을 확인하는 것으로 귀착된다
(\hyperref[0.1.5.1-ko]{0, 1.5.1}). 이렇게 정의한 사상
\[
  \operatorname{Spec}(\mathscr{O}_y)\longrightarrow Y
\]
을 \emph{표준 사상}이라 한다.

\begin{proposition}[2.4.2]
\label{I.2.4.2-ko}
\oldpage[I]{102}
$(Y,\mathscr{O}_Y)$를 준스킴이라 하자. 각 $y\in Y$에 대하여
$(\psi,\theta)$를 표준 사상
\[
  (\operatorname{Spec}(\mathscr{O}_y),\widetilde{\mathscr{O}}_y)
  \longrightarrow(Y,\mathscr{O}_Y)
\]
이라 하자. 그러면 $\psi$는 $\operatorname{Spec}(\mathscr{O}_y)$에서
$y\in\overline{\{z\}}$를 만족하는 점 $z$들로 이루어진 $Y$의 부분공간
$S_y$ 위로 가는 위상동형이다. 다시 말해 이 점들은 $y$의
\emph{일반화들}이다(\hyperref[0.2.1.2-ko]{0, 2.1.2}). 또한
$z=\psi(\mathfrak{p})$이면
\[
  \theta_z^{\sharp}:\mathscr{O}_z\longrightarrow
  (\mathscr{O}_y)_{\mathfrak{p}}
\]
는 동형이다. 따라서 $(\psi,\theta)$는 환 달린 공간의 단사사상이다.
\end{proposition}

$\operatorname{Spec}(\mathscr{O}_y)$의 유일한 닫힌점 $a$는 이 공간의
모든 점의 폐포에 속하고 $\psi(a)=y$이므로, 연속사상 $\psi$에 의한
$\operatorname{Spec}(\mathscr{O}_y)$의 상은 $S_y$에 포함된다. 한편
$S_y$는 $y$를 포함하는 모든 아핀 열린집합에 포함되므로, $Y$가 아핀
스킴인 경우로 환원할 수 있다. 이 경우 명제는
(\hyperref[I.1.6.2-ko]{1.6.2})에서 따른다.

\emph{따라서 (\hyperref[I.2.1.5-ko]{2.1.5})에 의해
$\operatorname{Spec}(\mathscr{O}_y)$와 $y$를 포함하는 $Y$의 기약
닫힌 부분집합들의 집합 사이에는 전단사 대응이 있다.}

\begin{corollary}[2.4.3]
\label{I.2.4.3-ko}
$y\in Y$가 $Y$의 한 기약 성분의 일반점이기 위한 필요충분조건은,
국소환 $\mathscr{O}_y$의 유일한 소 아이디얼이 그 극대 아이디얼인
것이다. 다시 말해 $\mathscr{O}_y$의 차원이 $0$인 것이다.
\end{corollary}

\begin{proposition}[2.4.4]
\label{I.2.4.4-ko}
$(X,\mathscr{O}_X)$를 환이 $A$인 국소 스킴, $a$를 그 유일한 닫힌점,
$(Y,\mathscr{O}_Y)$를 준스킴이라 하자. 모든 사상
$u=(\psi,\theta):(X,\mathscr{O}_X)\to(Y,\mathscr{O}_Y)$는 유일하게
\[
  X\longrightarrow\operatorname{Spec}(\mathscr{O}_{\psi(a)})
  \longrightarrow Y
\]
로 분해된다. 여기서 두 번째 화살표는 표준 사상이고, 첫 번째 화살표는
국소 준동형 $\mathscr{O}_{\psi(a)}\to A$에 대응한다. 이로써 사상들의
집합
$\operatorname{Hom}((X,\mathscr{O}_X),(Y,\mathscr{O}_Y))$와 국소
준동형 $\mathscr{O}_y\to A$들의 집합($y\in Y$) 사이에 표준 전단사
대응이 주어진다.
\end{proposition}

실제로 모든 $x\in X$에 대하여 $a\in\overline{\{x\}}$이므로
$\psi(a)\in\overline{\{\psi(x)\}}$이다. 따라서 $\psi(X)$는
$\psi(a)$를 포함하는 모든 아핀 열린집합에 포함된다. 그러므로
$(Y,\mathscr{O}_Y)$가 환 $B$의 아핀 스킴인 경우로 환원할 수 있다.
이때 $\varphi\in\operatorname{Hom}(B,A)$에 대하여
$u=({}^a\varphi,\widetilde{\varphi})$이다
(\hyperref[I.1.7.3-ko]{1.7.3}). 또한
$\varphi^{-1}(\mathfrak{j}_a)=\mathfrak{j}_{\psi(a)}$이므로,
$B-\mathfrak{j}_{\psi(a)}$의 모든 원소의 $\varphi$에 의한 상은
국소환 $A$에서 가역이다. 따라서 명제의 분해는 분수환의 보편 성질
(\hyperref[0.1.2.4-ko]{0, 1.2.4})에서 따른다. 역으로, 모든 국소
준동형 $\mathscr{O}_y\to A$에는 $\psi(a)=y$인 유일한 사상
$(\psi,\theta):X\to\operatorname{Spec}(\mathscr{O}_y)$가 대응한다
(\hyperref[I.1.7.3-ko]{1.7.3}). 이를 표준 사상
$\operatorname{Spec}(\mathscr{O}_y)\to Y$와 합성하면 사상 $X\to Y$를
얻으며, 이로써 명제가 증명된다.

\begin{env}[2.4.5]
\label{I.2.4.5-ko}
환이 \emph{체} $K$인 아핀 스킴의 바탕공간은 한 점으로 이루어진다. $A$가
극대 아이디얼 $\mathfrak{m}$을 갖는 국소환이면, 모든 국소 준동형
$A\to K$의 핵은 $\mathfrak{m}$과 같으므로
\[
  A\longrightarrow A/\mathfrak{m}\longrightarrow K
\]
로 분해된다. 여기서 두 번째 화살표는 단사 준동형이다. 따라서 사상
$\operatorname{Spec}(K)\to\operatorname{Spec}(A)$들은 체의 단사
준동형 $A/\mathfrak{m}\to K$들과 전단사 대응한다.
\end{env}

$(Y,\mathscr{O}_Y)$를 준스킴이라 하자. 모든 $y\in Y$와
$\mathscr{O}_y$의 모든 아이디얼 $\mathfrak{a}_y$에 대하여 표준 준동형
$\mathscr{O}_y\to\mathscr{O}_y/\mathfrak{a}_y$는 사상
$\operatorname{Spec}(\mathscr{O}_y/\mathfrak{a}_y)
\to\operatorname{Spec}(\mathscr{O}_y)$를 정의한다. 이를 표준 사상
$\operatorname{Spec}(\mathscr{O}_y)\to Y$와 합성하면 사상
$\operatorname{Spec}(\mathscr{O}_y/\mathfrak{a}_y)\to Y$를 얻으며,
이 사상도 \emph{표준 사상}이라 한다. $\mathfrak{a}_y=\mathfrak{m}_y$,
즉 $\mathscr{O}_y/\mathfrak{a}_y=k(y)$인 경우에는 명제
(\hyperref[I.2.4.4-ko]{2.4.4})로부터 다음을 얻는다.

\begin{corollary}[2.4.6]
\label{I.2.4.6-ko}
\oldpage[I]{103}
$(X,\mathscr{O}_X)$를 그 환이 체 $K$인 국소 스킴, $\xi$를 $X$의
유일한 점, $(Y,\mathscr{O}_Y)$를 준스킴이라 하자. 모든 사상
$u=(\psi,\theta):(X,\mathscr{O}_X)\to(Y,\mathscr{O}_Y)$는 유일하게
\[
  X\longrightarrow\operatorname{Spec}(k(\psi(\xi)))
  \longrightarrow Y
\]
로 분해된다. 여기서 두 번째 화살표는 표준 사상이고, 첫 번째 화살표는
단사 준동형 $k(\psi(\xi))\to K$에 대응한다. 이로써 사상들의 집합
$\operatorname{Hom}((X,\mathscr{O}_X),(Y,\mathscr{O}_Y))$와 체의 단사
준동형 $k(y)\to K$들의 집합($y\in Y$) 사이에 표준 전단사 대응이
주어진다.
\end{corollary}

\begin{corollary}[2.4.7]
\label{I.2.4.7-ko}
모든 $y\in Y$에 대하여, 모든 표준 사상
$\operatorname{Spec}(\mathscr{O}_y/\mathfrak{a}_y)\to Y$는 환 달린
공간의 단사사상이다.
\end{corollary}

$\mathfrak{a}_y=0$인 경우에는 이미 이를 보았으며
(\hyperref[I.2.4.2-ko]{2.4.2}), 일반적인 경우에는
(\hyperref[I.1.7.5-ko]{1.7.5})를 적용하면 충분하다.

\begin{remark}[2.4.8]
\label{I.2.4.8-ko}
$X$를 국소 스킴, $a$를 그 유일한 닫힌점이라 하자. $a$를 포함하는
모든 아핀 열린집합은 반드시 $X$ 전체이므로, 모든 \emph{가역}
$\mathscr{O}_X$-가군(\hyperref[0.5.4.1-ko]{0, 5.4.1})은 반드시
\emph{$\mathscr{O}_X$와 동형}이다. 다시
말해 이 가군은 \emph{자명하다}. 이 성질은 임의의 아핀 스킴
$\operatorname{Spec}(A)$에 대해서는 일반적으로 성립하지 않는다.
$A$가 정규환일 때에는, $A$가 \emph{유일인수분해정역}이면 이 성질이
성립함을 제V장에서 보게 된다.
\end{remark}

\subsection{준스킴 위의 준스킴}
\label{subsection:I.2.5-ko}

\begin{definition}[2.5.1]
\label{I.2.5.1-ko}
준스킴 $S$가 주어졌을 때, 준스킴 $X$와 준스킴의 사상
$\varphi:X\to S$를 주는 것을
\emph{준스킴 $S$ 위의} 준스킴 $X$, 또는 \emph{$S$-준스킴}을
정의한다고 한다. 이때 $S$를 $S$-준스킴 $X$의 \emph{밑준스킴}이라
한다. 사상 $\varphi$를 $S$-준스킴 $X$의 \emph{구조 사상}이라 한다.
$S$가 환 $A$의 아핀 스킴이면, $\varphi$를 갖춘 $X$를 환 $A$ 위의
준스킴(또는 $A$-준스킴)이라고도 한다.
\end{definition}

(\hyperref[I.2.2.4-ko]{2.2.4})에 따르면, 환 $A$ 위의 준스킴을 주는
것은 구조층 $\mathscr{O}_X$가 $A$-\emph{대수들의 층}인 준스킴
$(X,\mathscr{O}_X)$를 주는 것과 동치이다. \emph{따라서 임의의
준스킴은 언제나 유일한 방식으로 $\mathbf{Z}$-준스킴으로 볼 수 있다.}

$\varphi:X\to S$가 $S$-준스킴 $X$의 구조 사상이면, 점 $x\in X$가
점 $s\in S$ \emph{위에 있다}는 것은 $\varphi(x)=s$라는 뜻이다.
$\varphi$가 우세 사상이면(\hyperref[I.2.2.6-ko]{2.2.6}), $X$가
$S$에 대해 \emph{우세하다}고 한다.

\begin{env}[2.5.2]
\label{I.2.5.2-ko}
$X$, $Y$를 두 $S$-준스킴이라 하자. 준스킴의 사상 $u:X\to Y$가
\emph{$S$ 위의 준스킴의 사상}(또는 \emph{$S$-사상})이라는 것은 도식
\phantomsection\label{I.2.5.2.diagram-ko}
\[
  \xymatrix{
    X\ar[rr]^u\ar[dr] & & Y\ar[dl]\\
    & S &
  }
\]
이 가환이라는 뜻이다. 여기서 빗금 화살표들은 구조 사상이다. 이
조건으로부터 모든 $s\in S$와 $s$ 위에 있는 모든 $x\in X$에 대하여
$u(x)$도 $s$ 위에 있어야 함이 따른다.
\end{env}

이 정의에서 두 $S$-사상의 합성이 $S$-사상임이 곧바로 보인다.
따라서 $S$-준스킴들은 하나의 \emph{범주}를 이룬다.

$S$-준스킴 $X$에서 $S$-준스킴 $Y$로 가는 $S$-사상들의 집합을
$\operatorname{Hom}_S(X,Y)$로 표기한다. $S$-준스킴 $X$의 항등사상은
$1_X$로 나타낸다.

$S$가 환 $A$의 아핀 스킴이면 $S$-사상을 $A$-\emph{사상}이라고도
한다.

\begin{env}[2.5.3]
\label{I.2.5.3-ko}
\oldpage[I]{104}
$X$가 $S$-준스킴이고 $v:X'\to X$가 준스킴의 사상이면, 합성 사상
$X'\xrightarrow{v}X\to S$는 $X'$를 $S$-준스킴으로 정의한다. 특히
$X$의 열린집합 $U$ 위에 유도된 모든 준스킴은 표준 단사사상을
통하여 $S$-준스킴으로 볼 수 있다.

$u:X\to Y$가 $S$-준스킴들의 $S$-사상이면, $X$의 열린집합 $U$ 위에
유도된 모든 준스킴에 대한 $u$의 제한은 $S$-사상 $U\to Y$이다.
역으로 $(U_\alpha)$가 $X$의 열린 덮개이고, 각 $\alpha$에 대하여
$u_\alpha:U_\alpha\to Y$가 $S$-사상이라고 하자. 모든 지표쌍
$(\alpha,\beta)$에 대하여 $U_\alpha\cap U_\beta$에 대한
$u_\alpha$와 $u_\beta$의 제한이 일치하면, 각 $U_\alpha$에 대한
제한이 $u_\alpha$인 $S$-사상 $X\to Y$가 유일하게 존재한다.

$u:X\to Y$가 $u(X)\subset V$를 만족하는 $S$-사상이고 $V$가 $Y$의
열린 부분집합이면, $u$를 $X$에서 $V$로 가는 사상으로 보아도 여전히
$S$-사상이다.
\end{env}

\begin{env}[2.5.4]
\label{I.2.5.4-ko}
$S'\to S$를 준스킴의 사상이라 하자. 모든 $S'$-준스킴 $X$에 대하여
합성 사상 $X\to S'\to S$는 $X$를 $S$-준스킴으로 정의한다.

역으로 $S'$이 $S$의 한 열린집합 위에 유도된 준스킴이라고 하자.
$X$가 $S$-준스킴이고 그 구조 사상 $f:X\to S$가
$f(X)\subset S'$을 만족하면, $X$를 $S'$-준스킴으로 볼 수 있다.
이 후자의 경우, $Y$가 바탕공간을 $S'$ 안으로 보내는 구조 사상을
가진 $S$-준스킴이면, $X$에서 $Y$로 가는 모든 $S$-사상은
$S'$-사상이기도 하다.
\end{env}

\begin{env}[2.5.5]
\label{I.2.5.5-ko}
$X$가 $S$-준스킴이고 $\varphi:X\to S$가 그 구조 사상이면,
$S$에서 $X$로 가는 $S$-사상, 즉 $\varphi\circ\psi$가 $S$의
항등사상인 준스킴의 사상 $\psi:S\to X$를 $X$의
\emph{$S$-단면}이라 한다. $X$의 $S$-단면들의 집합을
$\Gamma(X/S)$로 나타낸다.
\end{env}
% ===== END EXACT INLINE INPUT: c1s2.tex =====


% ===== BEGIN EXACT INLINE INPUT: c1s3.tex =====
\section{준스킴의 곱}
\label{section:I.3-ko}

\subsection{준스킴의 합}
\label{subsection:I.3.1-ko}

\phantomsection\label{I.3.1.1-ko}
$(X_\alpha)$를 임의의 준스킴족이라 하고, $X$를 그 바탕공간
$X_\alpha$들의 위상적 \emph{합}인 위상공간이라 하자. 그러면 $X$는
서로소인 열린 부분공간들 $X'_\alpha$의 합집합이고, 각 $\alpha$에
대하여 $X_\alpha$에서 $X'_\alpha$ 위로 가는 위상동형
$\varphi_\alpha$가 존재한다. 각 $X'_\alpha$에 층
$(\varphi_\alpha)_*(\mathscr{O}_{X_\alpha})$를 주면, $X$는 분명히
준스킴이 된다. 이를 준스킴족 $(X_\alpha)$의 \emph{합}이라 하며
$\bigsqcup_\alpha X_\alpha$로 표기한다. $Y$가 준스킴이면 사상
$f\mapsto(f\circ\varphi_\alpha)$는 집합 $\operatorname{Hom}(X,Y)$에서
곱집합 $\prod_\alpha\operatorname{Hom}(X_\alpha,Y)$ 위로 가는
\emph{전단사}이다. 특히 $X_\alpha$들이 구조 사상
$\psi_\alpha$를 가진 $S$-준스킴들이면, 모든 $\alpha$에 대하여
$\psi\circ\varphi_\alpha=\psi_\alpha$를 만족하는 유일한 사상
$\psi:X\to S$에 의해 $X$는 $S$-준스킴이 된다. 두 준스킴 $X$, $Y$의
합은 $X\sqcup Y$로 표기한다. $X=\operatorname{Spec}(A)$,
$Y=\operatorname{Spec}(B)$이면 $X\sqcup Y$가
$\operatorname{Spec}(A\times B)$와 표준적으로 동일시됨은 자명하다.

\subsection{준스킴의 곱}
\label{subsection:I.3.2-ko}

\begin{definition}[3.2.1]
\label{I.3.2.1-ko}
두 $S$-준스킴 $X$, $Y$가 주어졌다고 하자. $S$-준스킴 $Z$와 두
$S$-사상 $p_1:Z\to X$, $p_2:Z\to Y$로 이루어진 삼중항
$(Z,p_1,p_2)$가 다음 조건을 만족하면 이를 $S$-준스킴 $X$와 $Y$의
곱이라 한다. 모든 $S$-준스킴 $T$에 대하여 사상
\oldpage[I]{105}
$f\mapsto(p_1\circ f,p_2\circ f)$는 $T$에서 $Z$로 가는 $S$-사상들의
집합에서, $S$-사상 $T\to X$ 하나와 $S$-사상 $T\to Y$ 하나로
이루어진 쌍들의 집합 위로 가는 전단사이다. 다시 말해 다음은
전단사이다.
\[
  \operatorname{Hom}_S(T,Z)\xrightarrow{\sim}
  \operatorname{Hom}_S(T,X)\times\operatorname{Hom}_S(T,Y).
\]
\end{definition}

이는 범주의 두 대상의 \emph{곱}이라는 일반적인 개념을
$S$-준스킴들의 범주에 적용한 것이다(T, I, 1.1). 특히 두
$S$-준스킴의 곱은 유일한 $S$-동형까지 \emph{유일}하다. 이 유일성
때문에 두 $S$-준스킴 $X$, $Y$의 곱은 보통 $X\times_S Y$로
표기하고(혼동의 염려가 없으면 단순히 $X\times Y$로 표기한다),
$X\times_S Y$에서 각각 $X$와 $Y$로 가는 사상 $p_1$, $p_2$는
\emph{표준 사영}이라 부르되 표기에서는 생략한다. $g:T\to X$,
$h:T\to Y$가 두 $S$-사상이면 $p_1\circ f=g$, $p_2\circ f=h$를
만족하는 $S$-사상 $f:T\to X\times_S Y$를 $(g,h)_S$ 또는 단순히
$(g,h)$로 나타내겠다. $X'$, $Y'$이 두 $S$-준스킴이고 $p'_1$,
$p'_2$가 존재한다고 가정한 $X'\times_S Y'$의 표준 사영이며,
$u:X'\to X$, $v:Y'\to Y$가 두 $S$-사상이면, $X'\times_S Y'$에서
$X\times_S Y$로 가는 $S$-사상 $(u\circ p'_1,v\circ p'_2)_S$를
$u\times_S v$ 또는 단순히 $u\times v$로 쓰겠다.

$S$가 환 $A$의 아핀 스킴이면 앞의 표기에서 $S$를 $A$로 바꾸어
쓰는 경우가 많다.

\begin{proposition}[3.2.2]
\label{I.3.2.2-ko}
$X$, $Y$, $S$를 세 아핀 스킴이라 하고 각각의 환을 $B$, $C$, $A$라
하자. $Z=\operatorname{Spec}(B\otimes_A C)$라 하고, $p_1$, $p_2$를
$B$와 $C$에서 $B\otimes_A C$로 가는 표준 $A$-준동형
$u:b\mapsto b\otimes1$, $v:c\mapsto1\otimes c$에 대응하는
(\hyperref[I.2.2.4-ko]{2.2.4}) $S$-사상이라 하자. 그러면
$(Z,p_1,p_2)$는 $X$와 $Y$의 곱이다.
\end{proposition}

(\hyperref[I.2.2.4-ko]{2.2.4})에 따라, 모든 $A$-준동형
$f:B\otimes_A C\to L$($L$은 $A$-대수)에 쌍
$(f\circ u,f\circ v)$를 대응시키는 사상이 전단사
\[
  \operatorname{Hom}_A(B\otimes_A C,L)\xrightarrow{\sim}
  \operatorname{Hom}_A(B,L)\times\operatorname{Hom}_A(C,L)\footnotemark
\]
\footnotetext{여기서 표기 $\operatorname{Hom}_A$는 $A$-대수의
준동형들의 집합을 뜻한다.}를 정의함을 확인하면 충분하다. 이는 정의와
관계 $b\otimes c=(b\otimes1)(1\otimes c)$에서 곧바로 따른다.

\begin{corollary}[3.2.3]
\label{I.3.2.3-ko}
$T$를 환 $D$의 아핀 스킴이라 하자.
$\alpha=({}^a\rho,\widetilde{\rho})$
(각각 $\beta=({}^a\sigma,\widetilde{\sigma})$)를 $S$-사상 $T\to X$
(각각 $T\to Y$)라 하자. 여기서 $\rho$(각각 $\sigma$)는 $B$(각각
$C$)에서 $D$로 가는 $A$-준동형이다. 그러면
$(\alpha,\beta)_S=({}^a\tau,\widetilde{\tau})$이다. 여기서 $\tau$는
$\tau(b\otimes c)=\rho(b)\sigma(c)$를 만족하는 준동형
$B\otimes_A C\to D$이다.
\end{corollary}

\begin{proposition}[3.2.4]
\label{I.3.2.4-ko}
$f:S'\to S$를 준스킴의 단사사상이라 하고(T, I, 1.1), $X$, $Y$를
$f$에 의해 $S$-준스킴으로도 간주하는 두 $S'$-준스킴이라 하자.
$S$-준스킴 $X$, $Y$의 모든 곱은 $S'$-준스킴 $X$, $Y$의 곱이고,
그 역도 성립한다.
\end{proposition}

$\varphi:X\to S'$, $\psi:Y\to S'$를 구조 사상이라 하자. $T$가
$S$-준스킴이고 $u:T\to X$, $v:T\to Y$가 두 $S$-사상이면, 정의에
따라 $f\circ\varphi\circ u=f\circ\psi\circ v=\theta$이다. 여기서
$\theta$는 $T$의 구조 사상이다. $f$에 대한 가정으로부터
$\varphi\circ u=\psi\circ v=\theta'$을 얻으며, $T$를 구조 사상
$\theta'$을 가진 $S'$-준스킴으로, $u$와 $v$를 $S'$-사상으로 간주할
수 있다. 그러므로 (\hyperref[I.3.2.1-ko]{3.2.1})을 고려하면 명제의
결론이 곧바로 따른다.

\begin{corollary}[3.2.5]
\label{I.3.2.5-ko}
$X$, $Y$를 구조 사상 $\varphi:X\to S$, $\psi:Y\to S$를 가진 두
$S$-준스킴이라 하고, $S'$을 $\varphi(X)\subset S'$,
$\psi(Y)\subset S'$을 만족하는 $S$의 열린 부분집합이라 하자.
$S$-준스킴 $X$, $Y$의 모든 곱은 $S'$-준스킴 $X$, $Y$의 곱이고,
그 역도 성립한다.
\end{corollary}

\oldpage[I]{106}
표준 단사사상 $S'\to S$에 (\hyperref[I.3.2.4-ko]{3.2.4})를
적용하면 충분하다.

\begin{theorem}[3.2.6]
\label{I.3.2.6-ko}
두 $S$-준스킴 $X$, $Y$가 주어지면 곱 $X\times_S Y$가 존재한다.
\end{theorem}

여러 단계로 나누어 증명하겠다.

\begin{lemma}[3.2.6.1]
\label{I.3.2.6.1-ko}
$(Z,p,q)$를 $X$와 $Y$의 곱이라 하고, $U$, $V$를 각각 $X$, $Y$의
열린 부분집합이라 하자. $W=p^{-1}(U)\cap q^{-1}(V)$라 놓으면,
$W$와 $p$, $q$의 $W$로의 제한(각각 사상 $W\to U$, $W\to V$로
간주한다)로 이루어진 삼중항은 $U$와 $V$의 곱이다.
\end{lemma}

실제로 $T$가 $S$-준스킴이면, $S$-사상 $T\to W$를 $T$를 $W$ 안으로
보내는 $S$-사상 $T\to Z$와 동일시할 수 있다. 이제 $g:T\to U$,
$h:T\to V$가 임의의 두 $S$-사상이면, 이를 각각 $T$에서 $X$, $Y$로
가는 $S$-사상으로 간주할 수 있다. 가정에 따라 $g=p\circ f$,
$h=q\circ f$를 만족하는 유일한 $S$-사상 $f:T\to Z$가 존재한다.
$p(f(T))\subset U$, $q(f(T))\subset V$이므로
\[
  f(T)\subset p^{-1}(U)\cap q^{-1}(V)=W,
\]
이고, 주장이 성립한다.

\begin{lemma}[3.2.6.2]
\label{I.3.2.6.2-ko}
$Z$를 $S$-준스킴, $p:Z\to X$, $q:Z\to Y$를 두 $S$-사상,
$(U_\alpha)$를 $X$의 열린 덮개, $(V_\lambda)$를 $Y$의 열린 덮개라
하자. 모든 쌍 $(\alpha,\lambda)$에 대하여 $S$-준스킴
$W_{\alpha\lambda}=p^{-1}(U_\alpha)\cap q^{-1}(V_\lambda)$와 $p$, $q$의
$W_{\alpha\lambda}$로의 제한이 $U_\alpha$와 $V_\lambda$의 곱을
이룬다고 가정하자. 그러면 $(Z,p,q)$는 $X$와 $Y$의 곱이다.
\end{lemma}

먼저 $f_1$, $f_2$가 두 $S$-사상 $T\to Z$일 때, 관계
$p\circ f_1=p\circ f_2$, $q\circ f_1=q\circ f_2$로부터
$f_1=f_2$가 따름을 보이자. 실제로 $Z$는 $W_{\alpha\lambda}$들의
합집합이므로 $f_1^{-1}(W_{\alpha\lambda})$들은 $T$의 열린 덮개를
이루며, $f_2^{-1}(W_{\alpha\lambda})$들도 마찬가지이다. 더 나아가
가정에 따라
\[
\begin{aligned}
  f_1^{-1}(W_{\alpha\lambda})
  &=f_1^{-1}(p^{-1}(U_\alpha))\cap
    f_1^{-1}(q^{-1}(V_\lambda)) \\
  &=f_2^{-1}(p^{-1}(U_\alpha))\cap
    f_2^{-1}(q^{-1}(V_\lambda))
   =f_2^{-1}(W_{\alpha\lambda})
\end{aligned}
\]
이다. 따라서 모든 지표쌍에 대하여 $f_1$과 $f_2$의
$f_1^{-1}(W_{\alpha\lambda})=f_2^{-1}(W_{\alpha\lambda})$로의 제한이
같음을 보이면 충분하다. 그런데 이 제한들은
$f_1^{-1}(W_{\alpha\lambda})$에서 $W_{\alpha\lambda}$로 가는
$S$-사상으로 간주할 수 있으므로, 가정과 정의
(\hyperref[I.3.2.1-ko]{3.2.1})에서 주장이 따른다.

이제 두 $S$-사상 $g:T\to X$, $h:T\to Y$가 주어졌다고 하자.
$T_{\alpha\lambda}=g^{-1}(U_\alpha)\cap h^{-1}(V_\lambda)$라 놓으면
$T_{\alpha\lambda}$들은 $T$의 열린 덮개를 이룬다. 가정에 따라
$p\circ f_{\alpha\lambda}$와 $q\circ f_{\alpha\lambda}$가 각각
$g$와 $h$의 $T_{\alpha\lambda}$로의 제한이 되는 $S$-사상
$f_{\alpha\lambda}$가 존재한다. 또한 $f_{\alpha\lambda}$와
$f_{\beta\mu}$의 $T_{\alpha\lambda}\cap T_{\beta\mu}$로의 제한이
일치함을 보이자. 이것으로 (\hyperref[I.3.2.6.2-ko]{3.2.6.2})의
증명이 끝난다. 정의에 따라 $f_{\alpha\lambda}$와
$f_{\beta\mu}$에 의한 $T_{\alpha\lambda}\cap T_{\beta\mu}$의 상은
$W_{\alpha\lambda}\cap W_{\beta\mu}$에 포함된다. 한편
\[
  W_{\alpha\lambda}\cap W_{\beta\mu}
  =p^{-1}(U_\alpha\cap U_\beta)\cap
   q^{-1}(V_\lambda\cap V_\mu)
\]
이므로, (\hyperref[I.3.2.6.1-ko]{3.2.6.1})에 따라
$W_{\alpha\lambda}\cap W_{\beta\mu}$와 이 준스킴으로 제한한 $p$,
$q$는 $U_\alpha\cap U_\beta$와 $V_\lambda\cap V_\mu$의 곱을
이룬다. $p\circ f_{\alpha\lambda}$와 $p\circ f_{\beta\mu}$가
$T_{\alpha\lambda}\cap T_{\beta\mu}$에서 일치하고,
$q\circ f_{\alpha\lambda}$와 $q\circ f_{\beta\mu}$도 마찬가지이므로
$f_{\alpha\lambda}$와 $f_{\beta\mu}$는
$T_{\alpha\lambda}\cap T_{\beta\mu}$에서 일치한다. 증명이 끝났다.

\oldpage[I]{107}
\begin{lemma}[3.2.6.3]
\label{I.3.2.6.3-ko}
$(U_\alpha)$를 $X$의 열린 덮개, $(V_\lambda)$를 $Y$의 열린 덮개라
하고, 모든 쌍 $(\alpha,\lambda)$에 대하여 $U_\alpha$와
$V_\lambda$의 곱이 존재한다고 가정하자. 그러면 $X$와 $Y$의 곱이
존재한다.
\end{lemma}

(\hyperref[I.3.2.6.1-ko]{3.2.6.1})을 열린집합
$U_\alpha\cap U_\beta$와 $V_\lambda\cap V_\mu$에 적용하면 $X$와
$Y$가 각각 이 열린집합들 위에 유도하는 $S$-준스킴들의 곱이
존재함을 알 수 있다. 또한 곱의 유일성으로부터, $i=(\alpha,\lambda)$,
$j=(\beta,\mu)$라 놓을 때 이 곱에서
$U_\alpha\times_S V_\lambda$(각각 $U_\beta\times_S V_\mu$)가 한
열린집합 위에 유도하는 $S$-준스킴 $W_{ij}$(각각 $W_{ji}$) 위로 가는
표준 동형 $h_{ij}$(각각 $h_{ji}$)가 존재함을 알 수 있다. 따라서
$f_{ij}=h_{ij}\circ h_{ji}^{-1}$은 $W_{ji}$에서 $W_{ij}$ 위로 가는
동형이다. 더 나아가 세 번째 쌍 $k=(\gamma,\nu)$에 대하여
$W_{ki}\cap W_{kj}$에서 $f_{ik}=f_{ij}\circ f_{jk}$이다. 이는
(\hyperref[I.3.2.6.1-ko]{3.2.6.1})을 각각 $U_\beta$, $V_\mu$ 안의
열린집합 $U_\alpha\cap U_\beta\cap U_\gamma$와
$V_\lambda\cap V_\mu\cap V_\nu$에 적용하여 얻는다. 따라서 준스킴
$Z$, 그 바탕공간의 열린 덮개 $(Z_i)$, 그리고 각 $i$에 대하여 유도
준스킴 $Z_i$에서 준스킴 $U_\alpha\times_S V_\lambda$ 위로 가는
동형 $g_i$가 존재하여 모든 쌍 $(i,j)$에 대하여
$f_{ij}=g_i\circ g_j^{-1}$이 성립한다
(\hyperref[I.2.3.1-ko]{2.3.1}). 또한
$g_i(Z_i\cap Z_j)=W_{ij}$이다. $p_i$, $q_i$, $\theta_i$를
$S$-준스킴 $U_\alpha\times_S V_\lambda$의 사영들과 구조 사상이라
하면, $Z_i\cap Z_j$에서 $p_i\circ g_i=p_j\circ g_j$이고 나머지 두
사상에 대해서도 마찬가지임을 곧바로 확인할 수 있다. 따라서 각
$Z_i$에서 $p$(각각 $q$, $\theta$)가 $p_i\circ g_i$(각각
$q_i\circ g_i$, $\theta_i\circ g_i$)와 일치한다는 조건으로 준스킴의
사상 $p:Z\to X$(각각 $q:Z\to Y$, $\theta:Z\to S$)를 정의할 수
있다. $\theta$를 갖춘 $Z$는 $S$-준스킴이다. 이제
$Z'_i=p^{-1}(U_\alpha)\cap q^{-1}(V_\lambda)$가 $Z_i$와 같음을
보이자. 모든 지표 $j=(\beta,\mu)$에 대하여
$Z_j\cap Z'_i=g_j^{-1}(p_j^{-1}(U_\alpha)\cap q_j^{-1}(V_\lambda))$이다.
그런데
\[
  p_j^{-1}(U_\alpha)\cap q_j^{-1}(V_\lambda)
  =p_j^{-1}(U_\alpha\cap U_\beta)\cap
   q_j^{-1}(V_\lambda\cap V_\mu).
\]
(\hyperref[I.3.2.6.1-ko]{3.2.6.1})에 따라 $p_j$와 $q_j$의
$p_j^{-1}(U_\alpha)\cap q_j^{-1}(V_\lambda)$로의 제한은 이
$S$-준스킴에 $U_\alpha\cap U_\beta$와 $V_\lambda\cap V_\mu$의 곱
구조를 준다. 그런데 곱의 유일성으로부터
$p_j^{-1}(U_\alpha)\cap q_j^{-1}(V_\lambda)=W_{ji}$가 따른다.
그러므로 모든 $j$에 대하여 $Z_j\cap Z'_i=Z_j\cap Z_i$이고,
$Z'_i=Z_i$이다. 이제 (\hyperref[I.3.2.6.2-ko]{3.2.6.2})에 따라
$(Z,p,q)$는 $X$와 $Y$의 곱이다.

\begin{lemma}[3.2.6.4]
\label{I.3.2.6.4-ko}
$\varphi:X\to S$, $\psi:Y\to S$를 $X$, $Y$의 구조 사상이라 하고,
$(S_i)$를 $S$의 열린 덮개라 하자. $X_i=\varphi^{-1}(S_i)$,
$Y_i=\psi^{-1}(S_i)$라 놓자. 각 곱 $X_i\times_S Y_i$가 존재하면
$X\times_S Y$가 존재한다.
\end{lemma}

(\hyperref[I.3.2.6.3-ko]{3.2.6.3})에 따라 모든 $i$, $j$에 대하여
곱 $X_i\times_S Y_j$가 존재함을 보이면 충분하다.
$X_{ij}=X_i\cap X_j=\varphi^{-1}(S_i\cap S_j)$,
$Y_{ij}=Y_i\cap Y_j=\psi^{-1}(S_i\cap S_j)$라 놓자.
(\hyperref[I.3.2.6.1-ko]{3.2.6.1})에 따라 곱
$Z_{ij}=X_{ij}\times_S Y_{ij}$가 존재한다. 이제 $T$가 $S$-준스킴이고
$g:T\to X_i$, $h:T\to Y_j$가 $S$-사상이면, $S$-사상의 정의에 따라
반드시 $\varphi(g(T))=\psi(h(T))\subset S_i\cap S_j$이다. 따라서
$g(T)\subset X_{ij}$, $h(T)\subset Y_{ij}$이고, $Z_{ij}$가 $X_i$와
$Y_j$의 곱임은 자명하다.

\begin{env}[3.2.6.5]
\label{I.3.2.6.5-ko}
이제 정리 (\hyperref[I.3.2.6-ko]{3.2.6})의 증명을 끝낼 수 있다.
$S$가 \emph{아핀 스킴}이면, $X$와 $Y$에는 각각 아핀 열린집합들로
이루어진 덮개 $(U_\alpha)$, $(V_\lambda)$가 존재한다.
(\hyperref[I.3.2.2-ko]{3.2.2})에 따라
$U_\alpha\times_S V_\lambda$가 존재하므로
(\hyperref[I.3.2.6.3-ko]{3.2.6.3})에 따라 $X\times_S Y$도 존재한다.
$S$가 임의의 준스킴이면 아핀 열린집합들로 이루어진 $S$의 덮개
$(S_i)$가 존재한다. $\varphi:X\to S$, $\psi:Y\to S$가 구조 사상이고
$X_i=\varphi^{-1}(S_i)$, $Y_i=\psi^{-1}(S_i)$라 놓으면 앞에서 보인
바에 따라 곱 $X_i\times_{S_i}Y_i$가 존재한다.
\oldpage[I]{108}
그러면 (\hyperref[I.3.2.5-ko]{3.2.5})에 따라 곱
$X_i\times_S Y_i$도 존재하고, (\hyperref[I.3.2.6.4-ko]{3.2.6.4})에
따라 $X\times_S Y$도 존재한다.
\end{env}

\begin{corollary}[3.2.7]
\label{I.3.2.7-ko}
$Z=X\times_S Y$를 두 $S$-준스킴의 곱, $p$, $q$를 $Z$에서 $X$, $Y$로
가는 사영, $\varphi$(각각 $\psi$)를 $X$(각각 $Y$)의 구조 사상이라
하자. $S'$을 $S$의 열린 부분집합, $U$(각각 $V$)를
$\varphi^{-1}(S')$(각각 $\psi^{-1}(S')$)에 포함되는 $X$(각각
$Y$)의 열린 부분집합이라 하자. 그러면 곱 $U\times_{S'}V$는
$p^{-1}(U)\cap q^{-1}(V)$ 위에 $Z$가 유도하는 준스킴($S'$-준스킴으로
간주한다)과 표준적으로 동일시된다. 또한 $f:T\to X$, $g:T\to Y$가
$f(T)\subset U$, $g(T)\subset V$를 만족하는 $S$-사상이면,
$S'$-사상 $(f,g)_{S'}$는 $(f,g)_S:T\to Z$의 공역을
$p^{-1}(U)\cap q^{-1}(V)$로 제한하여 얻은 사상과 동일시된다.
\end{corollary}

이는 (\hyperref[I.3.2.5-ko]{3.2.5})와
(\hyperref[I.3.2.6.1-ko]{3.2.6.1})에서 따른다.

\begin{env}[3.2.8]
\label{I.3.2.8-ko}
$(X_\alpha)$, $(Y_\lambda)$를 두 $S$-준스킴족이라 하고, $X$(각각
$Y$)를 족 $(X_\alpha)$(각각 $(Y_\lambda)$)의 합이라 하자
(\hyperref[subsection:I.3.1-ko]{3.1}). 그러면 $X\times_S Y$는 족
$(X_\alpha\times_S Y_\lambda)$의 \emph{합}과 동일시된다. 이는
(\hyperref[I.3.2.6.3-ko]{3.2.6.3})에서 곧바로 따른다.
\end{env}

\subsection{곱의 형식적 성질~; 밑준스킴의 변경}
\label{subsection:I.3.3-ko}

\begin{env}[3.3.1]
\label{I.3.3.1-ko}
독자는 이 절에서 진술하는 모든 성질이
(\hyperref[I.3.3.13-ko]{3.3.13})과
(\hyperref[I.3.3.15-ko]{3.3.15})을 제외하면, 그 진술에 나오는
곱들이 \emph{존재}할 때 \emph{임의의 범주}에서 아무 변경 없이
성립함을 알 수 있을 것이다. 실제로 범주의 임의의 대상 $S$에
대해서도 $S$-대상과 $S$-사상의 개념을
(\hyperref[subsection:I.2.5-ko]{2.5})에서와 정확히 같은 방식으로
정의할 수 있음은 분명하다.
\end{env}

\begin{env}[3.3.2]
\label{I.3.3.2-ko}
먼저 $X\times_S Y$는 $S$-준스킴의 범주에서 $X$와 $Y$에 관한
\emph{공변 쌍함자}이다. 실제로 다음 그림식이 가환임을 확인하면
충분하다.
\[
  \xymatrix{
    X\times Y\ar[r]^{f\times 1}\ar[d] &
    X'\times Y\ar[r]^{f'\times 1}\ar[d] &
    X''\times Y\ar[d]\\
    X\ar[r]_f & X'\ar[r]_{f'} & X''
  }
\]
\end{env}

\begin{proposition}[3.3.3]
\label{I.3.3.3-ko}
모든 $S$-준스킴 $X$에 대하여 $X\times_S S$(각각 $S\times_S X$)의
첫째 사영(각각 둘째 사영)은 $X\times_S S$(각각 $S\times_S X$)에서
$X$ 위로 가는 함자적 동형이고, 그 역동형은
$(1_X,\varphi)_S$(각각 $(\varphi,1_X)_S$)이다. 여기서 $\varphi$는
구조 사상 $X\to S$이다. 따라서 표준 동형까지
\[
  X\times_S S=S\times_S X=X
\]
라고 쓸 수 있다.
\end{proposition}

삼중항 $(X,1_X,\varphi)$가 $X$와 $S$의 곱임을 보이면 충분하다.
그런데 $T$가 $S$-준스킴이면, $T$에서 $S$로 가는 유일한 $S$-사상은
반드시 구조 사상 $\psi:T\to S$이다. $f$가 $T$에서 $X$로 가는
$S$-사상이면 반드시 $\psi=\varphi\circ f$이므로 주장이 따른다.

\begin{corollary}[3.3.4]
\label{I.3.3.4-ko}
$X$, $Y$를 두 $S$-준스킴, $\varphi:X\to S$, $\psi:Y\to S$를 그
구조 사상이라 하자. $X$를 $X\times_S S$와, $Y$를 $S\times_S Y$와
표준적으로 동일시하면, 사영 $X\times_S Y\to X$와
$X\times_S Y\to Y$는 각각 $1_X\times\psi$와
$\varphi\times1_Y$와 동일시된다.
\end{corollary}

확인은 즉시 되므로 독자에게 맡긴다.

\begin{env}[3.3.5]
\label{I.3.3.5-ko}
임의의 유한개 $n$개의 $S$-준스킴의 곱도
(\hyperref[subsection:I.3.2-ko]{3.2})에서와 같은 방식으로 정의할
수 있다.
\oldpage[I]{109}
이러한 곱들의 존재는 $n$에 관한 귀납법으로
(\hyperref[I.3.2.6-ko]{3.2.6})에서 따른다. 실제로
$(X_1\times_S X_2\times_S\cdots\times_S X_{n-1})\times_S X_n$은
곱의 정의를 만족한다. 곱의 유일성에서, 임의의 범주에서와 마찬가지로
그 \emph{교환법칙}과 \emph{결합법칙}이 따른다. 예를 들어
$p_1$, $p_2$, $p_3$가 $X_1\times_S X_2\times_S X_3$의 사영들이고
이 준스킴을 $(X_1\times_S X_2)\times_S X_3$와 동일시하면,
$X_1\times_S X_2$로 가는 사영은 $(p_1,p_2)_S$와 동일시된다.
\end{env}

\begin{env}[3.3.6]
\label{I.3.3.6-ko}
$S$, $S'$을 두 준스킴, $\varphi:S'\to S$를 $S'$에 $S$-준스킴
구조를 주는 사상이라 하자. 모든 $S$-준스킴 $X$에 대하여 곱
$X\times_S S'$을 생각하고, 각각 $X$와 $S'$로 가는 그 사영을
$p$와 $\pi'$이라 하자. $\pi'$을 갖춘 이 곱은 $S'$-준스킴이다.
이를 $S'$-준스킴으로 간주할 때 $X_{(S')}$ 또는
$X_{(\varphi)}$로 나타내고, $\varphi$를 사용하여 밑준스킴을
$S$에서 $S'$으로 \emph{확장하여 얻은 준스킴}, 또는 $\varphi$에
의한 $X$의 \emph{역상}이라 한다. $\pi$가 $X$의 구조 사상이고,
$\theta$가 $X\times_S S'$을 $S$-준스킴으로 간주했을 때의 구조
사상이면, 다음 그림식이 가환임에 유의하자.
\phantomsection\label{I.3.3.6.base-change-diagram-ko}
\[
  \xymatrix{
    X\ar[d]_{\pi} &
    X_{(S')}\ar[l]_{p}\ar[ld]_{\theta}\ar[d]^{\pi'}\\
    S & S'\ar[l]^{\varphi}
  }
\]
\end{env}

\begin{env}[3.3.7]
\label{I.3.3.7-ko}
(\hyperref[I.3.3.6-ko]{3.3.6})의 표기를 사용하자. 모든 $S$-사상
$f:X\to Y$에 대하여 $S'$-사상
$f\times_S1:X_{(S')}\to Y_{(S')}$도 $f_{(S')}$로 나타내고,
$\varphi$에 의한 사상 $f$의 \emph{역상}이라 한다. 따라서
$X_{(S')}$은 $X$에 관한 \emph{공변 함자}이며, $S$-준스킴의
범주에서 $S'$-준스킴의 범주로 간다.
\end{env}

\begin{env}[3.3.8]
\label{I.3.3.8-ko}
준스킴 $X_{(S')}$은 하나의 \emph{보편 사상 문제}의 해로도 볼 수
있다. 모든 $S'$-준스킴 $T$는 $\varphi$를 통하여 $S$-준스킴이기도
하다. 모든 $S$-사상 $g:T\to X$는 $g=p\circ f$라는 꼴로 유일하게
쓰인다. 여기서 $f$는 $S'$-사상 $T\to X_{(S')}$이다. 이는 곱의
정의를 $S$-사상 $g$와 $\psi:T\to S'$($T$의 구조 사상)에 적용하면
따른다.
\end{env}

\begin{proposition}[3.3.9]
\label{I.3.3.9-ko}
\textup{(「밑준스킴 확장의 추이성」).} $S''$을 준스킴,
$\varphi':S''\to S'$을 사상이라 하자. 모든 $S$-준스킴 $X$에
대하여 $S''$-준스킴 $(X_{(\varphi)})_{(\varphi')}$에서
$S''$-준스킴 $X_{(\varphi\circ\varphi')}$ 위로 가는 표준적이고
함자적인 동형이 존재한다.
\end{proposition}

실제로 $T$를 $S''$-준스킴, $\psi$를 그 구조 사상, $g$를 $T$에서
$X$로 가는 $S$-사상이라 하자. 여기서 $T$는 구조 사상이
$\varphi\circ\varphi'\circ\psi$인 $S$-준스킴으로 간주한다.
$T$는 구조 사상이 $\varphi'\circ\psi$인 $S'$-준스킴이기도 하므로
$g=p\circ g'$으로 쓸 수 있다. 여기서 $g'$은
$S'$-사상 $T\to X_{(\varphi)}$이다. 이어서
$g'=p'\circ g''$으로 쓸 수 있다. 여기서 $g''$은
$S''$-사상 $T\to(X_{(\varphi)})_{(\varphi')}$이다.
\phantomsection\label{I.3.3.9.transitivity-diagram-ko}
\[
  \xymatrix{
    X\ar[d]_{\pi} &
    X_{(\varphi)}\ar[l]_{p}\ar[d]_{\pi'} &
    (X_{(\varphi)})_{(\varphi')}\ar[l]_{p'}\ar[d]_{\pi''}\\
    S & S'\ar[l]^{\varphi} & S''\ar[l]^{\varphi'}
  }
\]
\oldpage[I]{110}
그러므로 보편 사상 문제의 해의 유일성에 따라 명제가 성립한다.

이 결과는 혼동의 염려가 없을 때 표준 동형까지
$(X_{(S')})_{(S'')}=X_{(S'')}$라고 써도 되며, 다음과 같이 쓸
수도 있다.
\begin{equation}
  (X\times_S S')\times_{S'}S''=X\times_S S''~;
  \tag{3.3.9.1}\label{I.3.3.9.1-ko}
\end{equation}
(\hyperref[I.3.3.9-ko]{3.3.9})에서 정의한 동형의 함자성도 모든
$S$-사상 $f:X\to Y$에 대하여 성립하는
\emph{사상의 역상의 추이성} 공식
\begin{equation}
  (f_{(S')})_{(S'')}=f_{(S'')}
  \tag{3.3.9.2}\label{I.3.3.9.2-ko}
\end{equation}
으로 나타난다.

\begin{corollary}[3.3.10]
\label{I.3.3.10-ko}
$X$와 $Y$가 두 $S$-준스킴이면, $S'$-준스킴
$X_{(S')}\times_{S'}Y_{(S')}$에서 $S'$-준스킴
$(X\times_S Y)_{(S')}$ 위로 가는 표준적이고 함자적인 동형이
존재한다.
\end{corollary}

실제로 표준 동형까지
\[
  (X\times_S S')\times_{S'}(Y\times_S S')
  =X\times_S(Y\times_S S')=(X\times_S Y)\times_S S'
\]
이다. 이는 (\hyperref[I.3.3.9.1-ko]{3.3.9.1})과 $S$-준스킴의
곱의 결합법칙에서 따른다.

(\hyperref[I.3.3.10-ko]{3.3.10})에서 정의한 동형의 함자성은 모든
$S$-사상쌍 $u:T\to X$, $v:T\to Y$에 대하여 성립하는 공식
\begin{equation}
  (u_{(S')},v_{(S')})_{S'}=((u,v)_S)_{(S')}
  \tag{3.3.10.1}\label{I.3.3.10.1-ko}
\end{equation}
으로 나타난다.

다시 말해 역상함자 $X_{(S')}$은 \emph{곱의 형성과 가환}한다.
또한 합의 형성과도 가환한다(\hyperref[I.3.2.8-ko]{3.2.8}).

\begin{corollary}[3.3.11]
\label{I.3.3.11-ko}
$Y$를 $S$-준스킴, $f:X\to Y$를 $X$에 $Y$-준스킴 구조를 주는
사상이라 하자. 따라서 $X$는 $S$-준스킴이기도 하다. 그러면 준스킴
$X_{(S')}$은 곱 $X\times_Y Y_{(S')}$와 동일시되고, 사영
$X\times_Y Y_{(S')}\to Y_{(S')}$은 $f_{(S')}$와 동일시된다.
\end{corollary}

$\psi:Y\to S$를 $Y$의 구조 사상이라 하자. 다음 그림식은
가환이다.
\phantomsection\label{I.3.3.11.diagram-ko}
\[
  \xymatrix{
    S'\ar[d] &
    Y_{(S')}\ar[l]_{\psi_{(S')}}\ar[d] &
    X_{(S')}\ar[l]_{f_{(S')}}\ar[d]\\
    S & Y\ar[l]^{\psi} & X\ar[l]^{f}
  }
\]
그런데 $Y_{(S')}$은 $S'_{(\psi)}$와, $X_{(S')}$은
$S'_{(\psi\circ f)}$와 동일시된다.
(\hyperref[I.3.3.9-ko]{3.3.9})와
(\hyperref[I.3.3.4-ko]{3.3.4})를 고려하면 따름정리가 따른다.

\begin{env}[3.3.12]
\label{I.3.3.12-ko}
$f:X\to X'$, $g:Y\to Y'$을 준스킴의 \emph{단사사상}인 두
$S$-사상이라 하자(T, I, 1.1). 그러면 $f\times_S g$도
\emph{단사사상}이다. 실제로 $p$, $q$를 $X\times_S Y$의 사영,
$p'$, $q'$을 $X'\times_S Y'$의 사영, $u$, $v$를 두
$S$-사상 $T\to X\times_S Y$라 하자. 관계
$(f\times_S g)\circ u=(f\times_S g)\circ v$에서
$p'\circ(f\times_S g)\circ u=p'\circ(f\times_S g)\circ v$, 즉
$f\circ p\circ u=f\circ p\circ v$가 따른다. $f$가
단사사상이므로 $p\circ u=p\circ v$이다. 마찬가지로 $g$가
단사사상임을 사용하면 $q\circ u=q\circ v$를 얻으므로 $u=v$이다.
\oldpage[I]{111}
따라서 밑준스킴의 모든 확장 $S'\to S$에 대하여
\phantomsection\label{I.3.3.12.base-change-monomorphism-ko}
\[
  f_{(S')}:X_{(S')}\to Y_{(S')}
\]
은 단사사상이다.
\end{env}

\begin{env}[3.3.13]
\label{I.3.3.13-ko}
$S$, $S'$을 각각 환 $A$, $A'$의 두 아핀 스킴이라 하자. 따라서
사상 $S'\to S$은 환 준동형 $A\to A'$에 대응한다. $X$가
$S$-준스킴이면 $S'$-준스킴 $X_{(S')}$을
$X_{(A')}$ 또는 $X\otimes_A A'$로도 나타내겠다. $X$ 자체가 환
$B$의 아핀 스킴이면, $X_{(A')}$은 환
$B_{(A')}=B\otimes_A A'$의 아핀 스킴이다. 이 환은
$A$-대수 $B$의 스칼라환을 $A'$으로 확장하여 얻는다.
\end{env}

\begin{env}[3.3.14]
\label{I.3.3.14-ko}
(\hyperref[I.3.3.6-ko]{3.3.6})의 표기를 사용하자. 모든
\emph{$S$-사상} $f:S'\to X$에 대하여
$f'=(f,1_{S'})_S$는 $p\circ f'=f$, $\pi'\circ f'=1_{S'}$을
만족하는 $S'$-사상 $S'\to X'=X_{(S')}$이다. 다시 말해 $f'$은
\emph{$X'$의 $S'$-단면}이다. 역으로 $f'$이 이러한 $S'$-단면이면
$f=p\circ f'$은 $S$-사상 $S'\to X$이다. 이로써 표준적인
\emph{전단사 대응}
\phantomsection\label{I.3.3.14.correspondence-ko}
\[
  \operatorname{Hom}_S(S',X)\overset{\sim}{\longrightarrow}
  \operatorname{Hom}_{S'}(S',X')
\]
을 정의한다. $f'$을 $f$의 \emph{그래프 사상}이라 하고
$\Gamma_f$로 나타낸다.
\end{env}

\begin{env}[3.3.15]
\label{I.3.3.15-ko}
준스킴 $X$가 주어지면 언제나 이를 $\mathbf Z$-준스킴으로 간주할
수 있다. 특히 (\hyperref[I.3.3.14-ko]{3.3.14})에 따라
$X\otimes_{\mathbf Z}\mathbf Z[T]$의 $X$-단면들($T$는
부정원)은 $\mathbf Z[T]$에서 $X$로 가는 \emph{사상}들과
전단사로 대응한다. 이 $X$-단면들이
\emph{$X$ 위의 구조층 $\mathcal O_X$의 단면들}과도 전단사로
대응함을 보이자. $(U_\alpha)$를 아핀 열린집합들로 이루어진 $X$의
덮개라 하자. $u:X\to X\otimes_{\mathbf Z}\mathbf Z[T]$를
$X$-사상, $u_\alpha$를 그 $U_\alpha$로의 제한이라 하자.
$A_\alpha$가 아핀 스킴 $U_\alpha$의 환이면
$U_\alpha\otimes_{\mathbf Z}\mathbf Z[T]$는 환
$A_\alpha[T]$의 아핀 스킴이고
(\hyperref[I.3.2.2-ko]{3.2.2}), $u_\alpha$는 표준적으로
$A_\alpha$-준동형 $A_\alpha[T]\to A_\alpha$에 대응한다(1.7.3).
그런 준동형은 $T$의 $A_\alpha$에서의 상, 즉
$s_\alpha\in A_\alpha=\Gamma(U_\alpha,\mathcal O_X)$를 주면
완전히 결정된다. $u_\alpha$와 $u_\beta$의 제한이
$V\subset U_\alpha\cap U_\beta$인 아핀 열린집합 $V$에서
일치한다는 조건을 쓰면, $s_\alpha$와 $s_\beta$도 $V$에서
일치함을 곧바로 알 수 있다. 따라서 족 $(s_\alpha)$은
$X$ 위의 $\mathcal O_X$의 한 단면 $s$를 각 $U_\alpha$로 제한하여
얻은 족이다. 역으로 이러한 단면은, 어떤 $X$-사상
$X\to X\otimes_{\mathbf Z}\mathbf Z[T]$을 각 $U_\alpha$로
제한하여 얻은 사상족 $(u_\alpha)$를 정의함이 분명하다. 이 결과는
II, 1.7.12에서 일반화될 것이다.
\end{env}

\subsection{준스킴에 값을 갖는 준스킴의 점; 기하적 점}
\label{subsection:I.3.4-ko}

\begin{env}[3.4.1]
\label{I.3.4.1-ko}
$X$를 준스킴이라 하자. 모든 준스킴 $T$에 대하여 $X(T)$로
사상 $T\to X$들의 집합 $\operatorname{Hom}(T,X)$도 나타내며, 이 집합의
원소를 \emph{$T$에 값을 갖는 $X$의 점}이라고도 한다. 각 사상
$f:T\to T'$에 $X(T')$에서 $X(T)$으로 가는 사상
$u'\mapsto u'\circ f$를 대응시키면, $X$를 고정했을 때 $X(T)$가
준스킴의 범주에서 집합의 범주로 가는 \emph{$T$에 관한 반변함자}임을
알 수 있다. 또한 준스킴의 모든 사상 $g:X\to Y$는 $v\in X(T)$에
$g\circ v$를 대응시키는 함자적 사상 $X(T)\to Y(T)$를 정의한다.
\end{env}

\begin{env}[3.4.2]
\label{I.3.4.2-ko}
세 집합 $P$, $Q$, $R$와 두 사상 $\varphi:P\to R$, $\psi:Q\to R$가
주어졌다고 하자. $\varphi(p)=\psi(q)$를 만족하는 순서쌍 $(p,q)$들로
이루어진 곱집합 $P\times Q$의 부분집합을 ($\varphi$와 $\psi$에 관한)
\emph{$R$ 위에서의 $P$와 $Q$의 섬유곱}이라 하고
\oldpage[I]{112}
$P\times_R Q$로 나타낸다. $S$-준스킴의 곱의 정의
(\hyperref[I.3.2.1-ko]{3.2.1})는
(\hyperref[I.3.4.1-ko]{3.4.1})의 표기를 사용하면 다음 식으로도
해석할 수 있다.
\phantomsection\label{I.3.4.2.1-ko}
\[
  (X\times_S Y)(T)=X(T)\times_{S(T)}Y(T)
  \tag{3.4.2.1}
\]
여기서 사상 $X(T)\to S(T)$와 $Y(T)\to S(T)$는 각각 구조 사상
$X\to S$와 $Y\to S$에 대응한다.
\end{env}

\begin{env}[3.4.3]
\label{I.3.4.3-ko}
준스킴 $S$가 주어져 있고 $S$-준스킴과 $S$-사상만을 생각할 때,
$S$-사상 $T\to X$들의 집합 $\operatorname{Hom}_S(T,X)$을
$X(T)_S$로도 나타내며 혼동의 염려가 없으면 첨자 $S$를 생략한다.
$X(T)_S$의 원소를 \emph{점}(혼동할 염려가 있으면 \emph{$S$-점})이라고도
하며, 이는 \emph{$T$에 값을 갖는 $S$-준스킴 $X$의 점}이다. 특히 $X$의
\emph{$S$-단면}은 다름 아닌 \emph{$S$에 값을 갖는 $X$의 점}이다.
(\hyperref[I.3.4.2.1-ko]{3.4.2.1})의 식은 다음과 같이도 쓸 수 있다.
\phantomsection\label{I.3.4.3.1-ko}
\[
  (X\times_S Y)(T)_S=X(T)_S\times Y(T)_S~;
  \tag{3.4.3.1}
\]
더 일반적으로 $Z$가 $S$-준스킴이고 $X$, $Y$, $T$가
$Z$-준스킴(따라서 \emph{그 자체로} $S$-준스킴)이면
\phantomsection\label{I.3.4.3.2-ko}
\[
  (X\times_Z Y)(T)_S=X(T)_S\times_{Z(T)_S}Y(T)_S.
  \tag{3.4.3.2}
\]

$S$-준스킴 $W$와 두 $S$-사상 $r:W\to X$, $s:W\to Y$로 이루어진
삼중항 $(W,r,s)$가 ($Z$ 위에서) $X$와 $Y$의 곱임을 보이려면,
정의에 따라 \emph{모든 $S$-준스킴 $T$}에 대하여 다음 그림에서
\phantomsection\label{I.3.4.3.product-diagram-ko}
\[
  \xymatrix{
    W(T)_S\ar[r]^{r'}\ar[d]_{s'} & X(T)_S\ar[d]^{\varphi'}\\
    Y(T)_S\ar[r]_{\psi'} & Z(T)_S
  }
\]
$W(T)_S$가 $Z(T)_S$ 위에서 $X(T)_S$와 $Y(T)_S$의 섬유곱임을
확인하면 충분하다. 여기서 $r'$과 $s'$은 $r$과 $s$에,
$\varphi'$과 $\psi'$은 구조 사상 $\varphi:X\to Z$와
$\psi:Y\to Z$에 각각 대응한다.
\end{env}

\begin{env}[3.4.4]
\label{I.3.4.4-ko}
위에서 $T$(각각 $S$)가 환 $B$(각각 $A$)의 아핀 스킴이면 앞의
표기에서 $T$(각각 $S$)를 $B$(각각 $A$)로 바꾼다. 이때 $X(B)$의
원소를 \emph{환 $B$에 값을 갖는 $X$의 점}, $X(B)_A$의 원소를
\emph{$A$-대수 $B$에 값을 갖는 $A$-준스킴 $X$의 점}이라고 한다.
$X(B)$와 $X(B)_A$는 이제 \emph{$B$에 관한 공변함자}임에 유의하자.
마찬가지로 $A$-준스킴 $T$에 값을 갖는 $A$-준스킴 $X$의 점들의
집합을 $X(T)_A$로 쓴다.
\end{env}

\begin{env}[3.4.5]
\label{I.3.4.5-ko}
특히 $T=\operatorname{Spec}(A)$이고 $A$가 \emph{국소}환인 경우를
생각하자. 그러면 $X(A)$의 원소들은 $x\in X$에 대한 \emph{국소}
준동형 $\mathcal O_x\to A$들과 전단사로 대응한다
(\hyperref[I.2.4.4-ko]{2.4.4}). 이때 $X$의 바탕공간의 점 $x$를
그에 대응하는 $A$에 값을 갖는 $X$의 점의 \emph{소재점}이라 한다.

더욱 특별히 준스킴 $X$의 \emph{기하적 점}이란
\emph{어떤 체 $K$에 값을 갖는 $X$의 점}을 말한다. 따라서 이러한
점 하나를 주는 것은 그 점의 $X$의 바탕공간에서의
\oldpage[I]{113}
소재점 $x$와 $k(x)$의 \emph{확대체} $K$를 주는 것과 같다. $K$를
그 기하적 점의 \emph{값체}라고 하며, 이 기하적 점이
\emph{$x$에 위치한다}고도 한다. 이로써 $K$에 값을 갖는 기하적 점을
그 소재점에 대응시키는 사상 $X(K)\to X$를 정의한다.

$S'=\operatorname{Spec}(K)$가 $S$-준스킴이고(다시 말해
$s\in S$인 어떤 잉여체 $k(s)$의 확대체로 $K$를 보고 있고),
$X$가 $S$-준스킴이라고 하자. $X(K)_S$의 원소, 즉
\emph{$K$에 값을 갖고 $s$ 위에 있는 $X$의 기하적 점}은
$s$ \emph{위에 있는} $X$의 점 $x$에 대하여 잉여체 $k(x)$에서
$K$로 가는 $k(s)$-단사 준동형을 주는 것과 같다. 이때 $k(x)$는
$k(s)$의 확대체이다.

더욱 특별히 $S=\operatorname{Spec}(K)=\{\xi\}$이면
\emph{$K$에 값을 갖는 $X$의 기하적 점들은 $k(x)=K$인 점
$x\in X$와 동일시된다}. 이러한 점을 \emph{$K$ 위에서 유리적인
$K$-준스킴 $X$의 점}, 곧 $K$-유리점이라고도 한다. $K'$이 $K$의
확대체이면 $K'$에 값을 갖는 $X$의 기하적 점들은
$K'$ 위에서 유리적인 $X'=X_{(K')}$의 점들과 전단사로 대응한다
(\hyperref[I.3.3.14-ko]{3.3.14}).
\end{env}

\begin{lemma}[3.4.6]
\label{I.3.4.6-ko}
$X_i$ ($1\leq i\leq n$)를 $S$-준스킴들, $s$를 $S$의 한 점,
$x_i$ ($1\leq i\leq n$)를 $s$ 위에 있는 $X_i$의 한 점이라 하자.
그러면 $k(s)$의 확대체 $K$와, 곱
$Y=X_1\times_S X_2\times_S\cdots\times_S X_n$의 $K$에 값을 갖는
기하적 점이 존재하여 그 각 사영은 $x_i$에 위치한다.
\end{lemma}

실제로 어떤 하나의 $k(s)$의 확대체 $K$ 안에 $k(s)$-단사 준동형
$k(x_i)\to K$들을 잡을 수 있다(Bourbaki, \emph{Alg.}, chap. V,
\S~4, prop. 2). 합성 $k(s)\to k(x_i)\to K$들은 모두 같으므로
$k(x_i)\to K$에 대응하는 사상 $\operatorname{Spec}(K)\to X_i$들은
$S$-사상이다. 따라서 이들은 유일한 사상
$\operatorname{Spec}(K)\to Y$를 정의한다. $y$를 이에 대응하는
$Y$의 점이라 하면, $y$의 각 $X_i$로의 사영은 분명히 $x_i$이다.

\begin{proposition}[3.4.7]
\label{I.3.4.7-ko}
$X_i$ ($1\leq i\leq n$)를 $S$-준스킴들이라 하고, 각 첨자 $i$에
대하여 $x_i$를 $X_i$의 한 점이라 하자. $x_i$가
$Y=X_1\times_S X_2\times_S\cdots\times_S X_n$의 어떤 점 $y$의
$i$번째 사영이 될 필요충분조건은 모든 $x_i$가 $S$의 같은 한 점
$s$ 위에 있는 것이다.
\end{proposition}

이 조건이 필요함은 자명하고, 보조정리
(\hyperref[I.3.4.6-ko]{3.4.6})이 충분함을 보인다.

\phantomsection\label{I.3.4.7.underlying-map-ko}
다시 말해 $X$의 바탕집합을 $(X)$로 나타내면 표준적인 \emph{전사}
사상 $(X\times_S Y)\to (X)\times_{(S)}(Y)$이 있음을 알 수 있다.
이 사상은 일반적으로 \emph{단사이지 않음}에 유의해야 한다. 즉
\emph{$X\times_S Y$에서 서로 다른 여러 점 $z$가 같은 사영
$x\in X$, $y\in Y$를 가질 수 있다}. 실제로 $S$, $X$, $Y$가 각각
체 $k$, $K$, $K'$의 소 스펙트럼인 경우에도 이를 볼 수 있는데,
텐서곱 $K\otimes_k K'$은 일반적으로 서로 다른 여러 소 아이디얼을
갖기 때문이다(\hyperref[I.3.4.9-ko]{3.4.9} 참조).

\begin{corollary}[3.4.8]
\label{I.3.4.8-ko}
$f:X\to Y$를 $S$-사상, $f_{(S')}:X_{(S')}\to Y_{(S')}$를
밑준스킴의 확장 $S'\to S$로 $f$에서 얻는 $S'$-사상이라 하자.
$p$(각각 $q$)를 사영 $X_{(S')}\to X$(각각 $Y_{(S')}\to Y$)라
하자. $X$의 모든 부분집합 $M$에 대하여
\phantomsection\label{I.3.4.8.base-change-image-ko}
\[
  q^{-1}(f(M))=f_{(S')}(p^{-1}(M)).
\]
\end{corollary}

\oldpage[I]{114}
실제로 (\hyperref[I.3.3.11-ko]{3.3.11})에 따라 다음 가환 그림으로
$X_{(S')}$을 곱 $X\times_Y Y_{(S')}$과 동일시한다.
\phantomsection\label{I.3.4.8.proof-diagram-ko}
\[
  \xymatrix{
    X\ar[d]_f & X_{(S')}\ar[l]_p\ar[d]^{f_{(S')}}\\
    Y & Y_{(S')}\ar[l]^q
  }
\]

(\hyperref[I.3.4.7-ko]{3.4.7})에 따라 $x\in M$,
$y'\in Y_{(S')}$에 대한 관계 $q(y')=f(x)$은
$p(x')=x$이고 $f_{(S')}(x')=y'$인 어떤 $x'\in X_{(S')}$가
존재함과 동치이다. 이로부터 따름정리가 얻어진다.

보조정리 (\hyperref[I.3.4.6-ko]{3.4.6})을 다음과 같이 정밀화할 수
있다.

\begin{proposition}[3.4.9]
\label{I.3.4.9-ko}
$X$, $Y$를 두 $S$-준스킴, $x$를 $X$의 점, $y$를 $Y$의 점이라
하고, $x$와 $y$가 같은 점 $s\in S$ 위에 있다고 하자.
$X\times_S Y$에서 $x$와 $y$를 사영으로 갖는 점들의 집합은
$k(s)$의 확대체로 본 $k(x)$와 $k(y)$의 합성 확대체의 유형들의
집합과 표준적으로 전단사 대응한다(Bourbaki, \emph{Alg.}, chap. VIII,
\S~8, prop. 2).
\end{proposition}

$p$(각각 $q$)를 $X\times_S Y$에서 $X$(각각 $Y$)로 가는 사영이라
하고, $E$를 $X\times_S Y$의 바탕공간의 부분공간
$p^{-1}(x)\cap q^{-1}(y)$라 하자. 먼저 사상
$\operatorname{Spec}(k(x))\to S$와 $\operatorname{Spec}(k(y))\to S$가
$\operatorname{Spec}(k(x))\to\operatorname{Spec}(k(s))\to S$와
$\operatorname{Spec}(k(y))\to\operatorname{Spec}(k(s))\to S$로
각각 인수분해됨에 유의하자.
$\operatorname{Spec}(k(s))\to S$는 단사사상이므로
(\hyperref[I.2.4.7-ko]{2.4.7}),
(\hyperref[I.3.2.4-ko]{3.2.4})에 따라
\phantomsection\label{I.3.4.9.tensor-product-fibre-ko}
\[
  P=\operatorname{Spec}(k(x))\times_S\operatorname{Spec}(k(y))
  =\operatorname{Spec}(k(x))\times_{\operatorname{Spec}(k(s))}
   \operatorname{Spec}(k(y))
  =\operatorname{Spec}\bigl(k(x)\otimes_{k(s)}k(y)\bigr).
\]

서로 역인 두 사상 $\alpha:P_0\to E$, $\beta:E\to P_0$를 정의하자.
여기서 $P_0$는 준스킴 $P$의 바탕집합을 뜻한다.
$i:\operatorname{Spec}(k(x))\to X$와
$j:\operatorname{Spec}(k(y))\to Y$가 표준 사상들이면
(\hyperref[I.2.4.5-ko]{2.4.5}), $\alpha$를 사상 $i\times_S j$에
대응하는 바탕공간의 사상으로 잡는다. 한편 모든 $z\in E$는 가정에
따라 두 $k(s)$-단사 준동형 $k(x)\to k(z)$와 $k(y)\to k(z)$를
정의하고, 따라서 유도된 $k(s)$-대수 준동형
$k(x)\otimes_{k(s)}k(y)\to k(z)$와 이어서 사상
$\operatorname{Spec}(k(z))\to P$를 정의한다. $\beta(z)$는 이 사상에
의해 $P_0$에 정해지는 점으로 둔다. $\alpha\circ\beta$와
$\beta\circ\alpha$가 항등사상임은
(\hyperref[I.2.4.5-ko]{2.4.5})와 곱의 정의
(\hyperref[I.3.2.1-ko]{3.2.1})에서 따른다. 끝으로 $P_0$가
$k(x)$와 $k(y)$의 합성 확대체의 유형들의 집합과 전단사 대응함은
알려져 있다(Bourbaki, \emph{Alg.}, chap. VIII, \S~8, prop. 1).

\subsection{전사와 단사}
\label{subsection:I.3.5-ko}

\begin{env}[3.5.1]
\label{I.3.5.1-ko}
일반적으로 준스킴의 사상들이 갖는 한 성질 $\mathbf{P}$를 생각하고,
다음 두 명제를 생각하자.
\begin{enumerate}
  \item[(i)] $f:X\to X'$, $g:Y\to Y'$가 성질 $\mathbf{P}$를 갖는 두
  $S$-사상이면 $f\times_S g$도 성질 $\mathbf{P}$를 갖는다.
  \item[(ii)] $f:X\to Y$가 성질 $\mathbf{P}$를 갖는 $S$-사상이면,
  밑준스킴의 확장으로 $f$에서 얻는 모든 $S'$-사상
  $f_{(S')}:X_{(S')}\to Y_{(S')}$도 성질 $\mathbf{P}$를 갖는다.
\end{enumerate}

\oldpage[I]{115}
$f_{(S')}=f\times_S 1_{S'}$이므로 모든 준스킴 $X$에 대하여
\emph{항등사상} $1_X$가 성질 $\mathbf{P}$를 가지면 (i)이 (ii)를
함의함을 알 수 있다. 한편 $f\times_S g$는 합성사상
\phantomsection\label{I.3.5.1.product-composition-ko}
\[
  X\times_S Y\xrightarrow{f\times 1_Y}X'\times_S Y
  \xrightarrow{1_{X'}\times g}X'\times_S Y'
\]
이므로 성질 $\mathbf{P}$를 갖는 두 사상의 \emph{합성}도 이 성질을
가지면 (ii)가 (i)을 함의함을 알 수 있다.
\end{env}

이 관찰의 첫 응용은 다음 명제이다.

\begin{proposition}[3.5.2]
\label{I.3.5.2-ko}
\begin{enumerate}
  \item[(i)] $f:X\to X'$, $g:Y\to Y'$가 전사인 $S$-사상들이면
  $f\times_S g$도 전사이다.
  \item[(ii)] $f:X\to Y$가 전사인 $S$-사상이면 밑준스킴의 모든 확장
  $S'$에 대하여 $f_{(S')}$도 전사이다.
\end{enumerate}
\end{proposition}

두 전사의 합성은 전사이므로 (ii)만 증명하면 충분하다. 그런데 이
명제는 $M=X$에 (\hyperref[I.3.4.8-ko]{3.4.8})을 적용하면 곧바로
따른다.

\begin{proposition}[3.5.3]
\label{I.3.5.3-ko}
사상 $f:X\to Y$가 전사일 필요충분조건은 모든 체 $K$와 모든 사상
$\operatorname{Spec}(K)\to Y$에 대하여 $K$의 확대체 $K'$과 다음
그림을 가환이게 하는 사상 $\operatorname{Spec}(K')\to X$가 존재하는
것이다.
\phantomsection\label{I.3.5.3.surjectivity-diagram-ko}
\[
  \xymatrix{
    X\ar[d]_f & \operatorname{Spec}(K')\ar[l]\ar[d]\\
    Y & \operatorname{Spec}(K)\ar[l]
  }
\]
\end{proposition}

이 조건은 충분하다. 실제로 모든 $y\in Y$에 대하여 $K$가 $k(y)$의
확대체일 때 단사 준동형 $k(y)\to K$에 대응하는 사상
$\operatorname{Spec}(K)\to Y$에 이 조건을 적용하면 된다
(\hyperref[I.2.4.6-ko]{2.4.6}). 역으로 $f$가 전사라고 하고,
$y\in Y$를 $\operatorname{Spec}(K)$의 유일한 점의 상이라 하자.
$f(x)=y$인 $x\in X$가 존재한다. 이에 대응하는 단사 준동형
$k(y)\to k(x)$를 생각하자(\hyperref[I.2.2.1-ko]{2.2.1}).
$k(x)$와 $K$에서 $K'$으로 가는 $k(y)$-단사 준동형들이 존재하도록
$K'$을 $k(y)$의 확대체로 잡으면 충분하다(Bourbaki, \emph{Alg.},
chap. V, \S~4, prop. 2). 이때 $k(x)\to K'$에 대응하는 사상
$\operatorname{Spec}(K')\to X$가 요구한 조건을 만족한다.

(\hyperref[I.3.4.5-ko]{3.4.5})에서 도입한 말로 표현하면,
\emph{$K$에 값을 갖는 $Y$의 모든 기하적 점은 $K$의 어떤 확대체에
값을 갖는 $X$의 기하적 점에서 온다}.

\begin{definition}[3.5.4]
\label{I.3.5.4-ko}
준스킴의 사상 $f:X\to Y$가 모든 체 $K$에 대하여 대응하는 사상
$X(K)\to Y(K)$을 단사로 만들면, $f$를 \emph{보편적으로 단사}라고
하거나 \emph{라디시엘 사상}이라고 한다.
\end{definition}

정의에서 곧바로 준스킴의 모든 \emph{단사사상}(T, 1.1)은 라디시엘
사상임을 알 수 있다.

\begin{env}[3.5.5]
\label{I.3.5.5-ko}
사상 $f:X\to Y$가 라디시엘 사상임을 보이려면 정의
(\hyperref[I.3.5.4-ko]{3.5.4})의 조건을 모든 \emph{대수적으로
닫힌} 체에 대하여 확인하는 것으로 충분하다. 실제로 $K$가 임의의
체이고 $K'$이 $K$의 대수적으로 닫힌 확대체이면 그림
\phantomsection\label{I.3.5.5.algebraic-closure-diagram-ko}
\[
  \xymatrix{
    X(K)\ar[r]^\alpha\ar[d]_\varphi & Y(K)\ar[d]^{\varphi'}\\
    X(K')\ar[r]_{\alpha'} & Y(K')
  }
\]

\oldpage[I]{116}
은 가환이다. 여기서 $\varphi$와 $\varphi'$은 사상
$\operatorname{Spec}(K')\to\operatorname{Spec}(K)$에서 오고,
$\alpha$와 $\alpha'$은 $f$에 대응한다. 그런데 $\varphi$는 단사이고
가정에 따라 $\alpha'$도 단사이므로 $\alpha$도 반드시 단사이다.
\end{env}

\begin{proposition}[3.5.6]
\label{I.3.5.6-ko}
$f:X\to Y$, $g:Y\to Z$를 준스킴의 두 사상이라 하자.
\begin{enumerate}
  \item[(i)] $f$와 $g$가 라디시엘 사상이면 $g\circ f$도
  라디시엘 사상이다.
  \item[(ii)] 역으로 $g\circ f$가 라디시엘 사상이면 $f$도
  라디시엘 사상이다.
\end{enumerate}
\end{proposition}

정의 (\hyperref[I.3.5.4-ko]{3.5.4})를 고려하면 이 명제는 사상들
$X(K)\to Y(K)\to Z(K)$에 대한 대응하는 자명한 명제들로 귀착된다.

\begin{proposition}[3.5.7]
\label{I.3.5.7-ko}
\begin{enumerate}
  \item[(i)] $S$-사상 $f:X\to X'$, $g:Y\to Y'$가 라디시엘
  사상들이면 $f\times_S g$도 라디시엘 사상이다.
  \item[(ii)] $S$-사상 $f:X\to Y$가 라디시엘 사상이면 밑준스킴의
  모든 확장 $S'\to S$에 대하여
  $f_{(S')}:X_{(S')}\to Y_{(S')}$도 라디시엘 사상이다.
\end{enumerate}
\end{proposition}

(\hyperref[I.3.5.1-ko]{3.5.1})에 따라 (i)만 증명하면 충분하다.
(\hyperref[I.3.4.2.1-ko]{3.4.2.1})에서
\phantomsection\label{I.3.5.7.field-valued-product-ko}
\[
  (X\times_S Y)(K)=X(K)\times_{S(K)}Y(K),\qquad
  (X'\times_S Y')(K)=X'(K)\times_{S(K)}Y'(K)
\]
임을 보았다. 이때 $f\times_S g$에 대응하는 사상
$(X\times_S Y)(K)\to(X'\times_S Y')(K)$은
$(u,v)\mapsto(f\circ u,g\circ v)$와 동일시되므로 명제가 곧바로
따른다.

\begin{proposition}[3.5.8]
\label{I.3.5.8-ko}
사상 $f=(\psi,\theta):X\to Y$가 라디시엘 사상일 필요충분조건은
$\psi$가 단사이고, 모든 $x\in X$에 대하여 단사 준동형
$\theta^x:k(\psi(x))\to k(x)$가 $k(x)$을 $k(\psi(x))$의 순수 비분리
확대체로 만드는 것이다.
\end{proposition}

$f$가 라디시엘 사상이라고 하고, 먼저 관계
$\psi(x_1)=\psi(x_2)=y$가 반드시 $x_1=x_2$를 함의함을 보이자.
실제로 $k(y)$의 확대체인 어떤 체 $K$와 $k(y)$-단사 준동형
$k(x_1)\to K$, $k(x_2)\to K$가 존재한다(Bourbaki, \emph{Alg.},
chap. V, \S~4, prop. 2). 따라서 이에 대응하는 두 사상
$u_1:\operatorname{Spec}(K)\to X$,
$u_2:\operatorname{Spec}(K)\to X$는 $f\circ u_1=f\circ u_2$를
만족한다. 가정에 따라 $u_1=u_2$이고, 특히 $x_1=x_2$이다. 이제
$\theta^x$를 통하여 $k(x)$을 $k(\psi(x))$의 확대체로 보자. $k(x)$이
$k(\psi(x))$의 순수 비분리 확대체가 아니면 $k(\psi(x))$의 어떤
대수적으로 닫힌 확대체 $K$ 안으로 가는 서로 다른 두
$k(\psi(x))$-단사 준동형 $k(x)\to K$가 존재하고, 이에 대응하는 두
사상 $\operatorname{Spec}(K)\to X$는 가정에 모순된다. 역으로
(\hyperref[I.2.4.6-ko]{2.4.6})을 고려하면 명제의 조건들이 $f$가
라디시엘 사상일 충분조건임은 곧바로 알 수 있다.

\begin{corollary}[3.5.9]
\label{I.3.5.9-ko}
$A$를 환, $S$를 $A$의 곱셈적 부분집합이라 하면 표준 사상
$\operatorname{Spec}(S^{-1}A)\to\operatorname{Spec}(A)$는 라디시엘
사상이다.
\end{corollary}

실제로 이 사상은 단사사상이다(\hyperref[I.1.6.2-ko]{1.6.2}).

\begin{corollary}[3.5.10]
\label{I.3.5.10-ko}
$f:X\to Y$를 라디시엘 사상, $g:Y'\to Y$를 사상이라 하고
$X'=X_{(Y')}=X\times_Y Y'$라 하자. 그러면 라디시엘 사상
$f_{(Y')}$ (\hyperref[I.3.5.7-ko]{3.5.7 (ii)})은 $X'$의 바탕공간을
$g^{-1}(f(X))$에 전단사로 대응시킨다. 또한 모든 체 $K$에 대하여
$X'(K)$는 $g$에 대응하는 사상 $Y'(K)\to Y(K)$에 의한, $Y(K)$의
부분집합 $X(K)$의 역상인 $Y'(K)$의 부분집합과 동일시된다.
\end{corollary}

첫째 명제는 (\hyperref[I.3.5.8-ko]{3.5.8})과
(\hyperref[I.3.4.8-ko]{3.4.8})에서 곧바로 따르고, 둘째 명제는
다음 그림의 가환성에서 따른다.
\oldpage[I]{117}
\phantomsection\label{I.3.5.10.base-change-points-diagram-ko}
\[
  \xymatrix{
    X'(K)\ar[r]\ar[d] & Y'(K)\ar[d]\\
    X(K)\ar[r] & Y(K)
  }
\]

\begin{remark}[3.5.11]
\label{I.3.5.11-ko}
준스킴의 사상 $f=(\psi,\theta)$에서 바탕공간의 사상 $\psi$가
단사이면 $f$를
\emph{단사}라고 하겠다. 사상 $f=(\psi,\theta):X\to Y$가 라디시엘
사상일 필요충분조건은 모든 사상 $Y'\to Y$에 대하여 사상
$f_{(Y')}:X_{(Y')}\to Y'$가 단사인 것이다. 이것이 \emph{보편적으로
단사인 사상}이라는 용어를 정당화한다. 실제로 조건의 필요성은
(\hyperref[I.3.5.7-ko]{3.5.7, (ii)})와
(\hyperref[I.3.5.8-ko]{3.5.8})에서 따른다. 역으로 이 조건은 먼저
$\psi$가 단사임을 함의한다. 어떤 $x\in X$에 대하여 단사 준동형
$\theta^x:k(\psi(x))\to k(x)$가 $k(x)$을 $k(\psi(x))$의 순수
비분리 확대체로 만들지 않는다면
$k(\psi(x))$의 어떤 확대체 $K$와 같은 사상
$\operatorname{Spec}(K)\to Y$에 대응하는 서로 다른 두 사상
$\operatorname{Spec}(K)\to X$가 존재할 것이다
(\hyperref[I.3.5.8-ko]{3.5.8}). 그러나 $Y'=\operatorname{Spec}(K)$로
두면 $X_{(Y')}$의 서로 다른 두 $Y'$-단면이 생긴다
(\hyperref[I.3.3.14-ko]{3.3.14}). 이는 $f_{(Y')}$가 단사라는
가정에 모순된다.
\end{remark}

\subsection{섬유}
\label{subsection:I.3.6-ko}

\begin{proposition}[3.6.1]
\label{I.3.6.1-ko}
$f:X\to Y$를 사상, $y$를 $Y$의 점, $\mathfrak{a}_y$를
$\mathscr{O}_y$의 $\mathfrak{m}_y$-준아딕 위상에 대한 정의
아이디얼이라 하자. 사영
$p:X\times_Y\operatorname{Spec}(\mathscr{O}_y/\mathfrak{a}_y)\to X$는
준스킴 $X\times_Y\operatorname{Spec}(\mathscr{O}_y/\mathfrak{a}_y)$의
바탕공간을, $X$의 바탕공간 위상에서 유도된 위상을 준 섬유
$f^{-1}(y)$ 위로 위상동형시킨다.
\end{proposition}

$\operatorname{Spec}(\mathscr{O}_y/\mathfrak{a}_y)\to Y$는 라디시엘
사상이고(\hyperref[I.3.5.4-ko]{3.5.4} 및
\hyperref[I.2.4.7-ko]{2.4.7}), 아이디얼
$\mathfrak{m}_y/\mathfrak{a}_y$가 가정에 따라 멱영이므로
(\hyperref[I.1.1.12-ko]{1.1.12})
$\operatorname{Spec}(\mathscr{O}_y/\mathfrak{a}_y)$에는 점이 하나뿐이다.
따라서 이미 (\hyperref[I.3.5.10-ko]{3.5.10} 및
\hyperref[I.3.3.4-ko]{3.3.4})에서 $p$가
$X\times_Y\operatorname{Spec}(\mathscr{O}_y/\mathfrak{a}_y)$의
바탕공간을 $f^{-1}(y)$와 \emph{집합으로서} 동일시함을 알고 있다.
그러므로 $p$가 위상동형임을 보이면 충분하다.
(\hyperref[I.3.2.7-ko]{3.2.7})에 따라 이 문제는 $X$와 $Y$ 위에서
국소적이므로 $X=\operatorname{Spec}(B)$,
$Y=\operatorname{Spec}(A)$이고 $B$가 $A$-대수라고 가정할 수 있다.
이때 사상 $p$는 준동형
$1\otimes\varphi:B\to B\otimes_A A'$에 대응한다. 여기서
$A'=A_y/\mathfrak{a}_y$이고 $\varphi$는 $A$에서 $A'$으로 가는 표준
사상이다. 그런데 $B\otimes_A A'$의 모든 원소는
\phantomsection\label{I.3.6.1.localization-fraction-ko}
\[
  \sum_i b_i\otimes\varphi(a_i)/\varphi(s)
  =\left(\sum_i(a_i b_i\otimes 1)\right)(1\otimes\varphi(s))^{-1}
\]
꼴로 쓸 수 있다. 여기서 $s\notin\mathfrak{j}_y$이며 명제
(\hyperref[I.1.2.4-ko]{1.2.4})를 적용할 수 있다.

\begin{env}[3.6.2]
\label{I.3.6.2-ko}
이 책의 앞으로의 모든 부분에서 사상의 섬유 $f^{-1}(y)$에
$k(y)$-준스킴의 구조가 주어졌다고 할 때에는,
\emph{언제나 $X\times_Y\operatorname{Spec}(k(y))$의 구조를 $X$로의
사영을 통하여 옮겨 얻은 준스킴을 뜻한다}. 이 마지막 곱을
$X\otimes_Y k(y)$ 또는 $X\otimes_{\mathscr{O}_Y}k(y)$라고도 쓰겠다.
더 일반적으로 $B$가 $\mathscr{O}_y$-대수이면 곱
$X\times_Y\operatorname{Spec}(B)$를 $X\otimes_Y B$ 또는
$X\otimes_{\mathscr{O}_Y}B$로 나타내겠다.
\end{env}

앞의 규약 아래에서 (\hyperref[I.3.5.10-ko]{3.5.10})에 따라
$k(y)$의 확대체 $K$에 값을 갖는 $X$의 점들은
\emph{$K$에 값을 갖는 $f^{-1}(y)$의 점들}과 동일시된다.

\begin{env}[3.6.3]
\label{I.3.6.3-ko}
$f:X\to Y$, $g:Y\to Z$를 두 사상, $h=g\circ f$를 그 합성이라 하자.
모든 $z\in Z$에 대하여 섬유 $h^{-1}(z)$는 다음 준스킴과 동형이다.
\phantomsection\label{I.3.6.3.fibre-composition-ko}
\[
  X\times_Z\operatorname{Spec}(k(z))
  =(X\times_Y Y)\times_Z\operatorname{Spec}(k(z))
  =X\times_Y g^{-1}(z)
\]
\oldpage[I]{118}
특히 $U$가 $X$의 열린부분이면, 섬유 준스킴 $f^{-1}(y)$가
$U\cap f^{-1}(y)$ 위에 유도하는 준스킴은 $f_U^{-1}(y)$와 동형이다.
여기서 $f_U$는 $f$의 $U$로의 제한이다.
\end{env}

\begin{proposition}[3.6.4) (``섬유의 추이성'']
\label{I.3.6.4-ko}
$f:X\to Y$, $g:Y'\to Y$를 두 사상이라 하자. $X'=X\times_Y Y'
=X_{(Y')}$, $f'=f_{(Y')}:X'\to Y'$로 놓자. 모든 $y'\in Y'$에
대하여 $y=g(y')$로 놓으면 준스킴 $f'^{-1}(y')$는
$f^{-1}(y)\otimes_{k(y)}k(y')$와 동형이다.
\end{proposition}

실제로 두 준스킴
$(X\otimes_Y k(y))\otimes_{k(y)}k(y')$와
$(X\times_Y Y')\otimes_{Y'}k(y')$가 모두
(\hyperref[I.3.3.9.1-ko]{3.3.9.1})에 따라
$X\times_Y\operatorname{Spec}(k(y'))$와 표준적으로 동형임을
관찰하면 된다.

특히 $V$가 $Y$에서 $y$의 열린 근방이고, $f_V$가 $f$를
$f^{-1}(V)$ 위에 유도된 준스킴에 제한한 사상이라면 준스킴 $f^{-1}(y)$와
$f_V^{-1}(y)$는 표준적으로 동일시된다.

\begin{proposition}[3.6.5]
\label{I.3.6.5-ko}
$f:X\to Y$를 사상, $y$를 $Y$의 점,
$Z=\operatorname{Spec}(\mathscr{O}_y)$를 국소 스킴,
$p=(\psi,\theta):X\times_Y Z\to X$를 사영이라 하자. $Z$의
바탕공간을 $Y$의 부분공간과 동일시하면
(\hyperref[I.2.4.2-ko]{2.4.2}), $p$는 $X\times_Y Z$의 바탕공간을
$X$의 부분공간 $f^{-1}(Z)$ 위로 위상동형시킨다. 또한 모든
$t\in X\times_Y Z$에 대하여 $z=\psi(t)$로 놓으면
$\theta_t^\sharp$은 $\mathscr{O}_z$에서 $\mathscr{O}_t$로 가는
동형이다.
\end{proposition}

$Z$는 $Y$의 부분공간으로 동일시했을 때 $y$를 포함하는 모든 아핀
열린집합에 포함되므로(\hyperref[I.2.4.2-ko]{2.4.2}),
(\hyperref[I.3.6.1-ko]{3.6.1})에서와 같이
$X=\operatorname{Spec}(A)$와 $Y=\operatorname{Spec}(B)$가 아핀
스킴이고 $A$가 $B$-대수인 경우로 귀착시킬 수 있다. 그러면
$X\times_Y Z$는 $A\otimes_B B_y$의 소 스펙트럼이고, 이 환은
$S^{-1}A$와 표준적으로 동일시된다. 여기서 $S$는
$B-\mathfrak{j}_y$의 $A$ 안에서의 상이다(0, 1.5.2). 이때 $p$는
표준 준동형 $A\to S^{-1}A$에 대응하므로 명제는
(\hyperref[I.1.6.2-ko]{1.6.2})에서 따른다.

\subsection{응용: 준스킴의 mod.
\texorpdfstring{$\mathfrak{J}$ 환원\protect\footnotemark}{J 환원}}
\label{subsection:I.3.7-ko}
\label{I.3.7-ko}
\footnotetext{이 항은 제I장의 뒷부분과 제II장의 개념과 결과를 사용하지만
뒤에서는 쓰지 않을 것이며, 고전 대수기하학에 익숙한 독자만을 위한 것이다.}

\begin{env}[3.7.1]
\label{I.3.7.1-ko}
$A$를 환, $X$를 $A$-준스킴, $\mathfrak{J}$를 $A$의 아이디얼이라 하자.
그러면 $X_0=X\otimes_A(A/\mathfrak{J})$는
$(A/\mathfrak{J})$-준스킴이며, 때로는 $X$에서 mod. $\mathfrak{J}$
\emph{환원}으로 얻어진다고 한다.
\end{env}

\begin{env}[3.7.2]
\label{I.3.7.2-ko}
이 용어는 주로 $A$가 \emph{국소환}이고 $\mathfrak{J}$가 그 극대
아이디얼인 경우에 쓰인다. 이때 $X_0$는 $A$의 잉여체
$k=A/\mathfrak{J}$ 위의 준스킴이다.

더 나아가 $A$가 분수체 $K$를 갖는 정역이면 $K$-준스킴
$X'=X\otimes_A K$도 생각할 수 있다. 우리가 사용하지 않을 용어의
남용으로서, 종래에는 흔히 $X_0$가 mod. $\mathfrak{J}$ 환원으로
\emph{$X'$에서 얻어진다}고 하였다. 이 말을 쓰던 경우에는 $A$가
차원 1인 국소환이었고(대개 이산 값매김환이었다), 주어진
$K$-준스킴 $X'$가 어떤 $K$-준스킴 $P'$의 닫힌 부분준스킴이라는
사실을 다소 명시적으로 전제하였다. 실제로 $P'$은
$\mathbf{P}_K^r$형의 사영공간이었고(II, 4.1.1 참조),
$P'=P\otimes_A K$ 꼴이었다. 여기서 $P$는 주어진 $A$-준스킴,
실제로는 II, 4.1.1의 표기에서 $A$-스킴 $\mathbf{P}_A^r$이었다.
우리의 언어로 $X'$에서 $X_0$를 정의하는 방법은 다음과 같다.

아핀 스킴 $Y=\operatorname{Spec}(A)$를 생각하자. 이는 유일한
\oldpage[I]{119}
닫힌점 $y=\mathfrak{J}$와 일반점 $(0)$의 두 점으로 이루어진다. 따라서 일반점
하나로 이루어진 집합 $U$는 $Y$ 안의 열린집합
$U=\operatorname{Spec}(K)$이다.
$X$가 $A$-준스킴, 다시 말해 $Y$-준스킴이고 $\psi:X\to Y$가 구조
사상이면 $X\otimes_A K=X'$은 $X$가 $\psi^{-1}(U)$ 위에 유도하는
준스킴과 다름없다. 특히 $\varphi:P\to Y$가 구조 사상이면
$P'=\varphi^{-1}(U)$의 닫힌 부분준스킴 $X'$은 $P$의 (국소 닫힌)
부분준스킴이다. $P$가 뇌터이면(예를 들어 $A$가 뇌터이고 $P$가
$A$ 위에서 유한형이면), $X'$을 포함하는 $P$의 가장 작은 닫힌
부분준스킴 $X=\overline{X'}$가 존재한다(9.5.10). 또한 $X'$은
$X$가 열린집합 $\varphi^{-1}(U)\cap X$ 위에 유도하는 준스킴이므로
$X\otimes_A K$와 동형이다(9.5.10).
\emph{따라서 $X'$의 $P'=P\otimes_A K$ 안으로의 몰입은 표준적인
방식으로 $X'$를 $X'=X\otimes_A K$ 꼴로 볼 수 있게 한다. 여기서
$X$는 $A$-준스킴이다.} 이제 mod. $\mathfrak{J}$로 환원한 준스킴
$X_0=X\otimes_A k$를 생각할 수 있으며, 이는 닫힌점 $y$의 섬유
$\psi^{-1}(y)$와 다름없다. 종래에는 적절한 용어가 없어서
$A$-준스킴 $X$를 명시적으로 도입하기를 피하였다. 그러나 통상
``mod. $\mathfrak{J}$로 환원한'' 준스킴 $X_0$에 관하여 말하는 모든
명제는 $X$ 자체에 관한 더 완전한 명제들의 귀결로 보아야 하며,
그렇게 해석해야만 만족스럽게 정식화하고 이해할 수 있다. 더구나
가정들은 언제나 $X'$의 $\mathbf{P}_K^r$ 안으로의 몰입을 미리 주는
일과는 무관하게 $X$ 자체에 관한 가정들로 귀착되는 듯하다. 이로써
더 내재적인 명제들을 얻을 수 있다.
\end{env}

\begin{env}[3.7.3]
\label{I.3.7.3-ko}
끝으로 여기서 생각한 상황의 개념적 정리가 늦어진 까닭이 되었을 법한
매우 특수한 사실을 지적하자. $A$가 이산 값매김환이고 $X$가
\emph{$A$ 위에서 고유}이면($X$가 $\mathbf{P}_A^r$의 닫힌
부분준스킴일 때에는 실제로 그러하다. II, 5.5.4 참조), $A$에 값을 갖는
$X$의 점들과 $K$에 값을 갖는 $X'$의 점들은 전단사 대응한다
(II, 7.3.8). 이 때문에 실제로는 $X$에 관한 명제들을 증명하면서도
흔히 $X'$에 관한 결과를 증명한다고 여겨 왔다. 그 명제들은 이 형태로
바탕 국소환의 차원이 1이라는 가정을 없애도 여전히 성립한다.
\end{env}
% ===== END EXACT INLINE INPUT: c1s3.tex =====


% ===== BEGIN EXACT INLINE INPUT: c1s4.tex =====
% Authority: EGA I ega1-4-fr.tex, whole-file SHA-256 A5F5CDAC81E654E9ABA75BBFEAD24F3010B122452C689B9EBAB5A7711B454EC1.
% Sequential coverage in this file: authority lines 1-395, complete subsections 4.1-4.2.
\section{부분준스킴과 몰입 사상}
\label{section:I.4-ko}

\subsection{부분준스킴}
\label{subsection:I.4.1-ko}

\begin{env}[4.1.1]
\label{I.4.1.1-ko}
준연접층이라는 개념은 국소적이므로
(\hyperref[0.5.1.3-ko]{0, 5.1.3}), 준스킴 $X$ 위의 준연접
$\mathscr{O}_X$-가군 $\mathscr{F}$는 다음 조건으로 정의할 수 있다.
$X$의 모든 아핀 열린집합 $V$에 대하여 $\mathscr{F}|V$는 어떤
$\Gamma(V,\mathscr{O}_X)$-가군에 대응하는 층과 동형이다
(\hyperref[I.1.4.1-ko]{1.4.1}). 준스킴 $X$ 위에서 구조층
$\mathscr{O}_X$가 준연접이고, 준연접 $\mathscr{O}_X$-가군 사이의
준동형의 핵, 여핵, 상과 준연접 $\mathscr{O}_X$-가군들의 귀납극한과
직합이 모두 준연접임은 자명하다
(\hyperref[I.1.3.7-ko]{1.3.7} 및
\hyperref[I.1.3.9-ko]{1.3.9}).
\end{env}

\begin{proposition}[4.1.2]
\label{I.4.1.2-ko}
$X$를 준스킴, $\mathscr{I}$를 $\mathscr{O}_X$의 준연접 아이디얼층이라
하자. 그러면 층 $\mathscr{O}_X/\mathscr{I}$의 지지집합 $Y$는 닫혀
있다. $\mathscr{O}_Y$를 $\mathscr{O}_X/\mathscr{I}$의 $Y$로의
제한이라 하면 $(Y,\mathscr{O}_Y)$는 준스킴이다.
\end{proposition}

\oldpage[I]{120}
(\hyperref[I.2.1.3-ko]{2.1.3})에 따라 $X$가 아핀 스킴인 경우만
생각하고, 이 경우 $Y$가 $X$에서 닫혀 있으며 \emph{아핀 스킴}임을
보이면 충분하다. 실제로 $X=\operatorname{Spec}(A)$이면
$\mathscr{O}_X=\widetilde{A}$이고
$\mathscr{I}=\widetilde{\mathfrak{I}}$이다. 여기서 $\mathfrak{I}$는
$A$의 아이디얼이다(\hyperref[I.1.4.1-ko]{1.4.1}). 이때 $Y$는
$X$의 닫힌 부분집합 $V(\mathfrak{I})$와 같고, 환
$B=A/\mathfrak{I}$의 소 스펙트럼과 동일시된다
(\hyperref[I.1.1.11-ko]{1.1.11}). 또한 $\varphi$가 표준 준동형
$A\to B=A/\mathfrak{I}$이면, 직상
${}^a\varphi_*(\widetilde{B})$는 층
$\widetilde{A}/\widetilde{\mathfrak{I}}
=\mathscr{O}_X/\mathscr{I}$와 표준적으로 동일시된다
(\hyperref[I.1.6.3-ko]{1.6.3} 및
\hyperref[I.1.3.9-ko]{1.3.9}). 이로써 증명이 끝난다.

$(Y,\mathscr{O}_Y)$를 \emph{아이디얼층 $\mathscr{I}$에 의해 정의되는}
$(X,\mathscr{O}_X)$의 \emph{부분준스킴}이라 한다. 이는 다음과 같은
일반적인 \emph{부분준스킴} 개념의 특수한 경우이다.

\begin{definition}[4.1.3]
\label{I.4.1.3-ko}
환 달린 공간 $(Y,\mathscr{O}_Y)$가 다음 조건들을 만족하면 이를
준스킴 $(X,\mathscr{O}_X)$의 \emph{부분준스킴}이라 한다.
\begin{enumerate}
  \item[$1^\circ$] $Y$는 $X$의 국소 닫힌 부분공간이다.
  \item[$2^\circ$] $U$를 $Y$를 포함하며 그 안에서 $Y$가 닫혀 있는
  $X$의 가장 큰 열린집합이라 하자. 다시 말해 $U$는 $X$에서
  $\overline{Y}$ 안에서의 $Y$의 경계를 뺀 집합이다. 그러면
  $(Y,\mathscr{O}_Y)$는 $\mathscr{O}_X|U$의 어떤 준연접 아이디얼층에
  의해 정의되는 $(U,\mathscr{O}_X|U)$의 닫힌 부분준스킴이다.
\end{enumerate}
$X$의 부분준스킴 $(Y,\mathscr{O}_Y)$에서 $Y$가 $X$ 안에서 닫혀
있으면 $(Y,\mathscr{O}_Y)$를 $X$의 \emph{닫힌} 부분준스킴이라 한다.
이 경우 $U=X$이다.
\end{definition}

이 정의와 (\hyperref[I.4.1.2-ko]{4.1.2})에서 $X$의 닫힌
부분준스킴들과 $\mathscr{O}_X$의 준연접 아이디얼층
$\mathscr{J}$들 사이에 \emph{표준적인 전단사 대응}이 있음을 곧바로
알 수 있다. 실제로 그러한 두 층 $\mathscr{J}$, $\mathscr{J}'$의
지지집합이 같은 닫힌집합 $Y$이고, $\mathscr{O}_X/\mathscr{J}$와
$\mathscr{O}_X/\mathscr{J}'$의 $Y$로의 제한이 서로 같으면 곧
$\mathscr{J}'=\mathscr{J}$이다.

\begin{env}[4.1.4]
\label{I.4.1.4-ko}
$(Y,\mathscr{O}_Y)$를 $X$의 부분준스킴, $U$를 $Y$를 포함하며 그
안에서 $Y$가 닫혀 있는 $X$의 가장 큰 열린집합, $V$를 $U$에
포함되는 $X$의 열린집합이라 하자. 그러면 $V\cap Y$는 $V$ 안에서
닫혀 있다. 또한 $Y$가 $\mathscr{O}_X|U$의 준연접 아이디얼층
$\mathscr{J}$에 의해 정의되면, $\mathscr{J}|V$는
$\mathscr{O}_X|V$의 준연접 아이디얼층이고, $Y$가 $Y\cap V$ 위에
유도하는 준스킴은 아이디얼층 $\mathscr{J}|V$에 의해 정의되는
$V$의 닫힌 부분준스킴임이 자명하다. 역으로 다음이 성립한다.
\end{env}

\begin{proposition}[4.1.5]
\label{I.4.1.5-ko}
$(Y,\mathscr{O}_Y)$를 환 달린 공간이라 하자. $Y$가 $X$의 부분공간이고,
$Y$를 덮는 $X$의 열린집합족 $(V_\alpha)$가 존재하여 모든
$\alpha$에 대하여 $Y\cap V_\alpha$가 $V_\alpha$ 안에서 닫혀 있으며
환 달린 공간
$(Y\cap V_\alpha,\mathscr{O}_Y|(Y\cap V_\alpha))$가 $X$가
$V_\alpha$ 위에 유도하는 준스킴의 닫힌 부분준스킴이라고 하자.
그러면 $(Y,\mathscr{O}_Y)$는 $X$의 부분준스킴이다.
\end{proposition}

가정에 따라 $Y$는 $X$ 안에서 국소 닫혀 있다. 또한 $Y$를 포함하며
그 안에서 $Y$가 닫혀 있는 가장 큰 열린집합 $U$는 모든
$V_\alpha$를 포함한다. 따라서 $U=X$이고 $Y$가 $X$ 안에서 닫힌
경우로 귀착할 수 있다. 이제 $\mathscr{O}_X$의 준연접 아이디얼층
$\mathscr{J}$를 다음과 같이 정의한다. $\mathscr{J}|V_\alpha$는 닫힌
부분준스킴
$(Y\cap V_\alpha,\mathscr{O}_Y|(Y\cap V_\alpha))$를 정의하는
$\mathscr{O}_X|V_\alpha$의 아이디얼층으로 잡고, $Y$와 만나지 않는
$X$의 모든 열린집합 $W$에 대해서는
$\mathscr{J}|W=\mathscr{O}_X|W$로 잡는다.
(\hyperref[I.4.1.3-ko]{4.1.3})과
(\hyperref[I.4.1.4-ko]{4.1.4})에 따라 이 조건들을 만족하는
아이디얼층 $\mathscr{J}$가 하나뿐임을 곧바로 확인할 수 있으며, 이
층은 닫힌 부분준스킴 $(Y,\mathscr{O}_Y)$를 정의한다.

특히 $X$가 $X$의 어떤 \emph{열린집합} 위에 \emph{유도하는}
준스킴은 $X$의 \emph{부분준스킴}이다.

\begin{proposition}[4.1.6]
\label{I.4.1.6-ko}
$X$의 부분\oldpage[I]{121}준스킴(각각 닫힌 부분준스킴)의
부분준스킴(각각 닫힌 부분준스킴)은 $X$의 부분준스킴(각각 닫힌
부분준스킴)과 표준적으로 동일시된다.
\end{proposition}

$X$의 국소 닫힌 부분공간의 국소 닫힌 부분집합은 $X$의 국소 닫힌
부분공간이므로, 문제가 국소적임은 자명하다
(\hyperref[I.4.1.5-ko]{4.1.5}). 따라서 $X$가 아핀이라고 가정해도
된다. 이제 명제는 환 $A$의 두 아이디얼 $\mathfrak{J}$,
$\mathfrak{J}'$가 $\mathfrak{J}\subset\mathfrak{J}'$을 만족할 때
$A/\mathfrak{J}'$와
$(A/\mathfrak{J})/(\mathfrak{J}'/\mathfrak{J})$를 표준적으로
동일시할 수 있다는 사실에서 따른다.

이후에는 언제나 앞의 동일시를 사용한다.

\begin{env}[4.1.7]
\label{I.4.1.7-ko}
$Y$를 준스킴 $X$의 부분준스킴이라 하고, $Y$의 \emph{바탕공간}에서
$X$의 바탕공간으로 가는 표준 단사사상을 $\psi:Y\to X$라 하자.
역상 $\psi^*(\mathscr{O}_X)$는 제한 $\mathscr{O}_X|Y$임이 알려져
있다(\hyperref[0.3.7.1-ko]{0, 3.7.1}). 모든 $y\in Y$에 대하여
$\omega_y$를 표준 준동형
$(\mathscr{O}_X)_y\to(\mathscr{O}_Y)_y$라 하자. 이 준동형들은 환의
층 사이의 \emph{전사} 준동형
$\omega:\mathscr{O}_X|Y\to\mathscr{O}_Y$가 줄기들에 유도하는
준동형들이다. 실제로 이를 $Y$ 위에서 국소적으로 확인하면 충분하므로,
$X$가 아핀이고 부분준스킴 $Y$가 닫힌 경우만 생각해도 된다. 이 경우
$\mathscr{I}$가 $Y$를 정의하는 $\mathscr{O}_X$의 아이디얼층이면,
$\omega_y$들은 준동형
$\mathscr{O}_X|Y\to(\mathscr{O}_X/\mathscr{I})|Y$가 줄기들에
유도하는 준동형들이다.

따라서 \emph{환 달린 공간의 단사사상}
$j=(\psi,\omega^b)$를 정의하였다
(\hyperref[0.4.1.1-ko]{0, 4.1.1}). 이는 준스킴의 사상
$Y\to X$임이 자명하며(\hyperref[I.2.2.1-ko]{2.2.1}), 이를
\emph{표준 포함 사상}이라 한다.

$f:X\to Z$가 사상이면 합성사상
$Y\xrightarrow{j}X\xrightarrow{f}Z$도 부분준스킴 $Y$에 대한
$f$의 \emph{제한}이라 한다.
\end{env}

\begin{env}[4.1.8]
\label{I.4.1.8-ko}
일반적인 정의(T, I, 1.1)에 따라, 준스킴의 사상 $f:Z\to X$가
$X$의 부분준스킴 $Y$의 포함 사상 $j:Y\to X$를 \emph{통하여
인수분해된다}고 함은
$f$가 $Z\xrightarrow{g}Y\xrightarrow{j}X$로 인수분해된다는 뜻이다.
여기서 $g$는 준스킴의 사상이며, $j$가 단사사상이므로 $g$는 반드시
\emph{유일}하다.
\end{env}

\begin{proposition}[4.1.9]
\label{I.4.1.9-ko}
사상 $f:Z\to X$가 \emph{포함 사상 $j:Y\to X$를 통하여
인수분해되기} 위한 필요충분조건은 다음과 같다. $f(Z)\subset Y$이고,
모든 $z\in Z$에 대하여 $y=f(z)$로 놓으면 $f$에 대응하는 준동형
$(\mathscr{O}_X)_y\to\mathscr{O}_z$가
$(\mathscr{O}_X)_y\to(\mathscr{O}_Y)_y\to\mathscr{O}_z$로
인수분해된다. 동치로, $(\mathscr{O}_X)_y\to\mathscr{O}_z$의 핵은
$(\mathscr{O}_X)_y\to(\mathscr{O}_Y)_y$의 핵을 포함한다.
\end{proposition}

이 조건들이 필요함은 자명하다. 충분함을 보이기 위해, 필요하다면
$X$를 $Y$가 그 안에서 닫히는 열린집합 $U$로 바꾸어
(\hyperref[I.4.1.3-ko]{4.1.3}) $Y$가 $X$의 \emph{닫힌}
부분준스킴인 경우로 귀착할 수 있다. 이때 $Y$는
$\mathscr{O}_X$의 준연접 아이디얼층 $\mathscr{I}$에 의해 정의된다.
$f=(\psi,\theta)$로 놓고, $\mathscr{J}$를
$\theta^\sharp:\psi^*(\mathscr{O}_X)\to\mathscr{O}_Z$의 핵인
$\psi^*(\mathscr{O}_X)$의 아이디얼층이라 하자. 함자 $\psi^*$의
성질(\hyperref[0.3.7.2-ko]{0, 3.7.2})을 고려하면, 가정에 따라 모든
$z\in Z$에 대하여
$(\psi^*(\mathscr{I}))_z\subset\mathscr{J}_z$이고, 따라서
$\psi^*(\mathscr{I})\subset\mathscr{J}$이다. 그러므로
$\theta^\sharp$은 다음과 같이 인수분해된다.
\[
  \psi^*(\mathscr{O}_X)\longrightarrow
  \psi^*(\mathscr{O}_X)/\psi^*(\mathscr{I})
  =\psi^*(\mathscr{O}_X/\mathscr{I})
  \xrightarrow{\omega}\mathscr{O}_Z.
\]
여기서 첫째 화살표는 표준 준동형이다. $\psi'$을 $Y$의 $X$ 안으로의
포함과 합성하면 $\psi$가 되는 연속사상 $Z\to Y$라 하자. 그러면
${\psi'}^*(\mathscr{O}_Y)=\psi^*(\mathscr{O}_X/\mathscr{I})$임은
자명하다. 한편 $\omega$는 국소 준동형임이 자명하므로
$g=(\psi',\omega^b)$는 준스킴의 사상 $Z\to Y$이다
\oldpage[I]{122}
(\hyperref[I.2.2.1-ko]{2.2.1}). 앞에서 본 바에 따라 $f=j\circ g$이므로
명제가 증명되었다.

\begin{corollary}[4.1.10]
\label{I.4.1.10-ko}
포함 사상 $Z\to X$가 \emph{포함 사상 $Y\to X$를 통하여
인수분해되기} 위한 필요충분조건은 $Z$가 $Y$의 부분준스킴인 것이다.
\end{corollary}

이때 $Z\leq Y$라고 쓰며, 이 관계는 $X$의 부분준스킴들의 집합 위의
\emph{순서관계}임이 자명하다.

\subsection{몰입 사상}
\label{subsection:I.4.2-ko}

\begin{definition}[4.2.1]
\label{I.4.2.1-ko}
사상 $f:Y\to X$가
$Y\xrightarrow{g}Z\xrightarrow{j}X$로 인수분해되고, 여기서 $g$가
동형사상, $Z$가 $X$의 부분준스킴(각각 닫힌 부분준스킴, 열린집합 위에
유도된 부분준스킴), $j$가 포함 사상이면 $f$를 \emph{몰입 사상}(각각
\emph{닫힌 몰입 사상}, \emph{열린 몰입 사상})이라 한다.
\end{definition}

이때 부분준스킴 $Z$와 동형사상 $g$는 \emph{유일하게} 결정된다. 실제로
$Z'$이 $X$의 또 다른 부분준스킴이고, $j'$이 포함 사상 $Z'\to X$이며,
$g'$이 $j\circ g=j'\circ g'$을 만족하는 동형사상 $Y\to Z'$이라고 하자.
그러면 $j'=j\circ g\circ {g'}^{-1}$이므로
$Z'\leq Z$이고(\hyperref[I.4.1.10-ko]{4.1.10}), 같은 방법으로
$Z\leq Z'$를 보일 수 있다. 따라서 $Z'=Z$이고, $j$가 준스킴의
단사사상이므로 $g'=g$이다.

$f=j\circ g$를 몰입 사상 $f$의 \emph{표준 인수분해}라 하고,
부분준스킴 $Z$와 동형사상 $g$를 \emph{$f$에 대응하는 것들}이라 한다.

몰입 사상은 준스킴의 \emph{단사사상}이고
(\hyperref[I.4.1.7-ko]{4.1.7}), \emph{따라서 특히} \emph{라디시엘 사상}임은
자명하다(\hyperref[I.3.5.4-ko]{3.5.4}).

\begin{proposition}[4.2.2]
\label{I.4.2.2-ko}
\begin{enumerate}
  \item[a)] 사상 $f=(\psi,\theta):Y\to X$가 \emph{열린 몰입 사상}이기
  위한 필요충분조건은, $\psi$가 $Y$에서 $X$의 열린 부분집합으로 가는
  위상동형이고 모든 $y\in Y$에 대하여 준동형
  $\theta_y^\sharp:\mathscr{O}_{\psi(y)}\to\mathscr{O}_y$가
  전단사인 것이다.
  \item[b)] 사상 $f=(\psi,\theta):Y\to X$가 \emph{몰입 사상}(각각
  \emph{닫힌 몰입 사상})이기 위한 필요충분조건은, $\psi$가 $Y$에서
  $X$의 국소 닫힌(각각 닫힌) 부분집합으로 가는 위상동형이고 모든
  $y\in Y$에 대하여 준동형
  $\theta_y^\sharp:\mathscr{O}_{\psi(y)}\to\mathscr{O}_y$가 전사인
  것이다.
\end{enumerate}
\end{proposition}

\begin{enumerate}
  \item[a)] 조건들이 필요함은 자명하다. 역으로 이 조건들이 성립하면
  $\theta^\sharp:\psi^*(\mathscr{O}_X)\to\mathscr{O}_Y$가 동형임이
  분명하다. 또한 $\psi^*(\mathscr{O}_X)$는
  $\mathscr{O}_X|\psi(Y)$에서 $\psi^{-1}$을 써서 구조를 옮겨 얻은
  층이므로 결론이 따른다.
  \item[b)] 조건들이 필요함은 자명하므로, 이제 충분함을 증명하자.
  먼저 $X$가 아핀 스킴이고 $Z=\psi(Y)$가 $X$ 안에서 \emph{닫혀}
  있다고 가정하는 특수한 경우를 생각하자.
  (\hyperref[0.3.4.6-ko]{0, 3.4.6})에 따라
  $\psi_*(\mathscr{O}_Y)$의 지지집합은 $Z$이고, 이 층의 $Z$로의
  제한을 $\mathscr{O}'_Z$라 하면 환 달린 공간
  $(Z,\mathscr{O}'_Z)$는 $Y$에서 $Z$로 가는 위상동형으로 본 $\psi$를
  써서 $(Y,\mathscr{O}_Y)$의 구조를 옮겨 얻는다.

  이제 $f_*(\mathscr{O}_Y)=\psi_*(\mathscr{O}_Y)$가 \emph{준연접}
  $\mathscr{O}_X$-가군임을 보이자. 실제로 $x\notin Z$이면
  $\psi_*(\mathscr{O}_Y)$는 $x$의 적당한 근방에 제한했을 때 영이다.
  반대로 $x\in Z$이면 유일하게 정해지는 $y\in Y$에 대하여
  $x=\psi(y)$이다. $V$를 $Y$에서 $y$의 아핀 열린 근방이라 하자.
  그러면 $\psi(V)$는 $Z$에서 열려 있으므로, $X$의 어떤 열린집합
  $U$와 $Z$의 교집합이다. $\psi_*(\mathscr{O}_Y)$의 $U$로의 제한은
  직\oldpage[I]{123}상 $(\psi_V)_*(\mathscr{O}_Y|V)$의 $U$로의 제한과 같다. 여기서
  $\psi_V$는 $\psi$의 $V$로의 제한이다.
  한편 사상 $(\psi,\theta)$의 $(V,\mathscr{O}_Y|V)$로의 제한은 이
  준스킴에서 $(X,\mathscr{O}_X)$로 가는 사상이므로
  $({}^a\varphi,\widetilde{\varphi})$ 꼴이다. 여기서 $\varphi$는 환
  $A=\Gamma(X,\mathscr{O}_X)$에서 환
  $\Gamma(V,\mathscr{O}_Y)$로 가는 준동형이다
  (\hyperref[I.1.7.3-ko]{1.7.3}). 따라서
  $(\psi_V)_*(\mathscr{O}_Y|V)$는 준연접 $\mathscr{O}_X$-가군이고
  (\hyperref[I.1.6.3-ko]{1.6.3}), 준연접층이라는 조건이 국소적이므로
  앞의 주장이 증명된다.

  더 나아가 $\psi$가 위상동형이라는 가정에서 모든 $y\in Y$에 대하여
  $\psi_y:(\psi_*(\mathscr{O}_Y))_{\psi(y)}\to\mathscr{O}_y$가
  동형임이 따른다(\hyperref[0.3.4.5-ko]{0, 3.4.5}). 도표
  \phantomsection\label{I.4.2.2.theta-diagram-ko}
  \[
    \xymatrix{
      \mathscr{O}_{\psi(y)}\ar[r]^{\theta_{\psi(y)}}
        \ar[d]_{\psi_y\circ\rho_{\psi(y)}} &
      (\psi_*(\mathscr{O}_Y))_{\psi(y)}\ar[d]^{\psi_y}\\
      (\psi^*(\mathscr{O}_X))_y\ar[r]^{\theta_y^\sharp} &
      \mathscr{O}_y
    }
  \]
  는 가환하고 두 세로 화살표는 동형이므로
  (\hyperref[0.3.7.2-ko]{0, 3.7.2}), $\theta_y^\sharp$이 전사라는
  가정에서 $\theta_{\psi(y)}$도 전사임이 따른다. 한편
  $\psi_*(\mathscr{O}_Y)$의 지지집합이 $Z=\psi(Y)$이므로,
  $\theta$는 $\mathscr{O}_X=\widetilde{A}$에서 준연접
  $\mathscr{O}_X$-가군 $f_*(\mathscr{O}_Y)$로 가는 \emph{전사}
  준동형이다. 따라서 어떤 $A$의 아이디얼 $\mathfrak{J}$와 유일한
  동형 $\omega:\widetilde{A}/\widetilde{\mathfrak{J}}
  \to f_*(\mathscr{O}_Y)$이 존재하여, $\omega$를 표준 준동형
  $\widetilde{A}\to\widetilde{A}/\widetilde{\mathfrak{J}}$와 합성하면
  $\theta$가 된다(\hyperref[I.1.3.8-ko]{1.3.8}).
  $\mathscr{O}_Z$를
  $\widetilde{A}/\widetilde{\mathfrak{J}}$의 $Z$로의 제한이라 하면
  $(Z,\mathscr{O}_Z)$는 $(X,\mathscr{O}_X)$의 부분준스킴이다.
  또한 $f$는 이 부분준스킴의 $X$ 안으로의 표준 포함 사상과 동형사상
  $(\psi_0,\omega_0)$을 거쳐 인수분해된다. 여기서 $\psi_0$는
  $Y$에서 $Z$로 가는 사상으로 본 $\psi$이고, $\omega_0$은
  $\omega$의 $\mathscr{O}_Z$로의 제한이다.

  이제 일반적인 경우를 생각하자. $U$를 $X$의 아핀 열린집합으로서
  $U\cap\psi(Y)$가 $U$ 안에서 닫혀 있고 공집합이 아닌 것으로 잡자.
  $f$를 열린집합 $\psi^{-1}(U)$ 위에서 $Y$가 유도하는 준스킴으로
  제한하고, 이를 $U$ 위에서 $X$가 유도하는 준스킴으로 가는 사상으로
  보면 첫째 경우로 귀착된다. 따라서 $f$의 $\psi^{-1}(U)$로의 제한은
  닫힌 몰입 사상 $\psi^{-1}(U)\to U$이고, 표준적으로
  $j_U\circ g_U$로 인수분해된다. 여기서 $g_U$는
  $\psi^{-1}(U)$에서 $U$의 부분준스킴 $Z_U$로 가는 동형사상이고,
  $j_U$는 표준 포함 사상 $Z_U\to U$이다.

  $V\subset U$인 $X$의 또 다른 아핀 열린집합 $V$를 잡자. $Z_U$의
  $V$로의 제한 $Z'_V$는 준스킴 $V$의 부분준스킴이므로, $f$의
  $\psi^{-1}(V)$로의 제한은 $j'_V\circ g'_V$로 인수분해된다. 여기서
  $j'_V$는 표준 포함 사상 $Z'_V\to V$이고, $g'_V$는
  $\psi^{-1}(V)$에서 $Z'_V$로 가는 동형사상이다. 몰입 사상의 표준
  인수분해가 유일하므로(\hyperref[I.4.2.1-ko]{4.2.1}) 반드시
  $Z'_V=Z_V$이고 $g'_V=g_V$이다. 따라서
  (\hyperref[I.4.1.5-ko]{4.1.5})에 따라 바탕공간이 $\psi(Y)$이고
  모든 $U\cap\psi(Y)$로의 제한이 $Z_U$인 $X$의 부분준스킴 $Z$가
  존재한다. 이때 $g_U$들은 각 $\psi^{-1}(U)$ 위에서 어떤 동형사상
  $g:Y\to Z$의 제한이고, $f=j\circ g$이다. 여기서 $j$는 표준 포함
  사상 $Z\to X$이다.
\end{enumerate}

\begin{corollary}[4.2.3]
\label{I.4.2.3-ko}
$X$를 아핀 스킴이라 하자. 사상 $f=(\psi,\theta):Y\to X$가 닫힌
몰입 사상이기 위한 필요충분조건은 $Y$가 아핀 스킴이고 준동형
$\Gamma(\theta):\Gamma(\mathscr{O}_X)\to\Gamma(\mathscr{O}_Y)$가
전사인 것이다.
\end{corollary}

\begin{corollary}[4.2.4]
\label{I.4.2.4-ko}
\begin{enumerate}
  \item[a)] $f:Y\to X$를 사상, $(V_\lambda)$를 $X$의 열린집합들로
  이루어진 $f(Y)$의 덮개라 하자. $f$가 몰입 사상(각각 열린 몰입
  사상)이기 위한 필요충분조건은
  \oldpage[I]{124}
  각 유도 준스킴
  $f^{-1}(V_\lambda)$로의 제한이 $V_\lambda$ 안으로의 몰입 사상
  (각각 열린 몰입 사상)인 것이다.
  \item[b)] $f:Y\to X$를 사상, $(V_\lambda)$를 $X$의 열린덮개라
  하자. $f$가 닫힌 몰입 사상이기 위한 필요충분조건은 각 유도
  준스킴 $f^{-1}(V_\lambda)$로의 제한이 $V_\lambda$ 안으로의 닫힌
  몰입 사상인 것이다.
\end{enumerate}
\end{corollary}

$f=(\psi,\theta)$라 하자. a)의 경우 모든 $y\in Y$에 대하여
$\theta_y^\sharp$은 전사(각각 전단사)이고, b)의 경우 모든 $y\in Y$에
대하여 $\theta_y^\sharp$은 전사이다. 따라서 a)의 경우 $\psi$가
$Y$에서 $X$의 국소 닫힌(각각 열린) 부분집합으로 가는 위상동형이고,
b)의 경우 $Y$에서 $X$의 닫힌 부분집합으로 가는 위상동형임만 확인하면
충분하다. 그런데 가정에 따라 $\psi$는 자명하게 단사이고, 모든
$y\in Y$에 대하여 $Y$에서 $y$의 모든 근방을 $\psi(Y)$에서
$\psi(y)$의 근방으로 보낸다. a)의 경우
$\psi(Y)\cap V_\lambda$는 $V_\lambda$에서 국소 닫혀(각각 열려)
있으므로 $\psi(Y)$는 $V_\lambda$들의 합집합에서 국소 닫혀(각각 열려)
있고, \emph{따라서 특히} $X$에서도 그러하다. b)의 경우
$\psi(Y)\cap V_\lambda$는 $V_\lambda$에서 닫혀 있으므로
$X=\bigcup_\lambda V_\lambda$에서 $\psi(Y)$는 닫혀 있다.

\begin{proposition}[4.2.5]
\label{I.4.2.5-ko}
두 몰입 사상(각각 두 열린 몰입 사상, 두 닫힌 몰입 사상)의 합성은
몰입 사상(각각 열린 몰입 사상, 닫힌 몰입 사상)이다.
\end{proposition}

이는 (\hyperref[I.4.1.6-ko]{4.1.6})에서 곧바로 따른다.

\subsection{몰입 사상의 곱}
\label{subsection:I.4.3-ko}

\begin{proposition}[4.3.1]
\label{I.4.3.1-ko}
$\alpha:X'\to X$, $\beta:Y'\to Y$를 두 $S$-사상이라 하자. $\alpha$와
$\beta$가 몰입 사상(각각 열린 몰입 사상, 닫힌 몰입 사상)이면
$\alpha\times_S\beta$도 몰입 사상(각각 열린 몰입 사상, 닫힌 몰입
사상)이다. 더욱이 $\alpha$(각각 $\beta$)가 $X'$(각각 $Y'$)를
$X$(각각 $Y$)의 부분준스킴 $X''$(각각 $Y''$)와 동일시한다면,
$\alpha\times_S\beta$는 $X'\times_S Y'$의 바탕공간을
$X\times_S Y$의 바탕공간의 부분공간
$p^{-1}(X'')\cap q^{-1}(Y'')$와 동일시한다. 여기서 $p$와 $q$는 각각
$X\times_S Y$에서 $X$와 $Y$로 가는 사영이다.
\end{proposition}

정의 (\hyperref[I.4.2.1-ko]{4.2.1})을 고려하면 $X'$와 $Y'$가
부분준스킴이고 $\alpha$와 $\beta$가 포함 사상인 경우로 한정할 수
있다. 열린집합 위에 유도된 부분준스킴에 대해서는 이 명제가 이미
(\hyperref[I.3.2.7-ko]{3.2.7})에서 증명되었다. 모든 부분준스킴은 어떤
열린집합 위에 유도된 준스킴의 닫힌 부분준스킴이므로
(\hyperref[I.4.1.3-ko]{4.1.3}), $X'$와 $Y'$가 \emph{닫힌}
부분준스킴인 경우로 귀착된다.

먼저 $S$가 \emph{아핀}이라고 가정할 수 있음을 보이자. 실제로
$(S_\lambda)$를 아핀 열린집합들로 이루어진 $S$의 덮개라 하고,
$\varphi$와 $\psi$를 $X$와 $Y$의 구조 사상이라 하자. 또
$X_\lambda=\varphi^{-1}(S_\lambda)$,
$Y_\lambda=\psi^{-1}(S_\lambda)$라 하자. $X'$의
$X_\lambda\cap X'$로의 제한 $X'_\lambda$(각각 $Y'$의
$Y_\lambda\cap Y'$로의 제한 $Y'_\lambda$)은 $X_\lambda$(각각
$Y_\lambda$)의 닫힌 부분준스킴이다. 준스킴 $X_\lambda$,
$Y_\lambda$, $X'_\lambda$, $Y'_\lambda$는 $S_\lambda$-준스킴으로
볼 수 있고, 곱 $X_\lambda\times_S Y_\lambda$와
$X_\lambda\times_{S_\lambda}Y_\lambda$(각각
$X'_\lambda\times_S Y'_\lambda$와
$X'_\lambda\times_{S_\lambda}Y'_\lambda$)는 동일하다
(\hyperref[I.3.2.5-ko]{3.2.5}). $S$가 아핀일 때 명제가 참이라면,
$\alpha\times_S\beta$를 각 $X'_\lambda\times_S Y'_\lambda$에
제한한 사상은 몰입 사상이다(\hyperref[I.3.2.7-ko]{3.2.7}). 또한 곱
$X'_\lambda\times_S Y'_\mu$(각각 $X_\lambda\times_S Y_\mu$)는
$(X'_\lambda\cap X'_\mu)\times_S
(Y'_\lambda\cap Y'_\mu)$(각각
$(X_\lambda\cap X_\mu)\times_S(Y_\lambda\cap Y_\mu)$)와
동일시되므로(\hyperref[I.3.2.6.4-ko]{3.2.6.4}),
$\alpha\times_S\beta$
\oldpage[I]{125}
를 각 $X'_\lambda\times_S Y'_\mu$에 제한한 사상도 몰입 사상이다.
따라서 (\hyperref[I.4.2.4-ko]{4.2.4})에 의해
$\alpha\times_S\beta$ 자체도 몰입 사상이다.

다음으로 $X$와 $Y$도 \emph{아핀}이라고 가정할 수 있음을 증명하자.
실제로 $(U_i)$(각각 $(V_j)$)를 아핀 열린집합들로 이루어진 $X$(각각
$Y$)의 덮개라 하고, $X'_i$(각각 $Y'_j$)를 $X'$의
$X'\cap U_i$로의 제한(각각 $Y'$의 $Y'\cap V_j$로의 제한)이라 하자.
이는 $U_i$(각각 $V_j$)의 닫힌 부분준스킴이다.
$U_i\times_S V_j$는 $X\times_S Y$를
$p^{-1}(U_i)\cap q^{-1}(V_j)$에 제한하여 얻은 준스킴과 동일시된다
(\hyperref[I.3.2.7-ko]{3.2.7}). 마찬가지로 $p'$, $q'$가
$X'\times_S Y'$의 사영이면 $X'_i\times_S Y'_j$는
$X'\times_S Y'$를
${p'}^{-1}(X'_i)\cap{q'}^{-1}(Y'_j)$에 제한하여 얻은 준스킴과
동일시된다. $\gamma=\alpha\times_S\beta$라 놓자. 정의에 따라
$p\circ\gamma=\alpha\circ p'$, $q\circ\gamma=\beta\circ q'$이고,
$X'_i=\alpha^{-1}(U_i)$, $Y'_j=\beta^{-1}(V_j)$이므로
${p'}^{-1}(X'_i)=\gamma^{-1}(p^{-1}(U_i))$,
${q'}^{-1}(Y'_j)=\gamma^{-1}(q^{-1}(V_j))$도 성립한다. 따라서
\[
  {p'}^{-1}(X'_i)\cap{q'}^{-1}(Y'_j)
  =\gamma^{-1}\bigl(p^{-1}(U_i)\cap q^{-1}(V_j)\bigr)
  =\gamma^{-1}(U_i\times_S V_j).
\]
이제 첫째 부분의 논증과 같이 결론을 얻는다.

그러므로 $X$, $Y$, $S$가 아핀이라고 가정하고, 그 환을 각각 $B$,
$C$, $A$라 하자. 이때 $B$와 $C$는 $A$-대수이고, $X'$와 $Y'$는
각각 $B$와 $C$의 몫대수 $B'$, $C'$를 환으로 갖는 아핀 스킴이다.
또한 $\alpha=({}^a\rho,\widetilde{\rho})$,
$\beta=({}^a\sigma,\widetilde{\sigma})$이고, 여기서 $\rho$와
$\sigma$는 각각 표준 준동형 $B\to B'$, $C\to C'$이다
(\hyperref[I.1.7.3-ko]{1.7.3}). 한편 $X\times_S Y$(각각
$X'\times_S Y'$)는 $B\otimes_A C$(각각 $B'\otimes_A C'$)를 환으로
갖는 아핀 스킴이고,
$\alpha\times_S\beta=({}^a\tau,\widetilde{\tau})$이다. 여기서 $\tau$는
$B\otimes_A C$에서 $B'\otimes_A C'$로 가는 준동형
$\rho\otimes\sigma$이다(\hyperref[I.3.2.2-ko]{3.2.2}와
\hyperref[I.3.2.3-ko]{3.2.3}). 이 준동형은 전사이므로
$\alpha\times_S\beta$는 실제로 몰입 사상이다. 더욱이
$\mathfrak{b}$(각각 $\mathfrak{c}$)가 $\rho$(각각 $\sigma$)의
핵이면 $\tau$의 핵은
$\operatorname{Im}(\mathfrak{b}\otimes_A C\to B\otimes_A C)
+\operatorname{Im}(B\otimes_A\mathfrak{c}\to B\otimes_A C)$이며, 이는
$u(\mathfrak{b})$와 $v(\mathfrak{c})$가 생성하는 아이디얼들의 합과
같다. 여기서 $u$(각각 $v$)는 준동형 $b\mapsto b\otimes1$(각각
$c\mapsto1\otimes c$)이다. $p=({}^a u,\widetilde{u})$이고
$q=({}^a v,\widetilde{v})$이므로, 이 핵은 $B\otimes_A C$의 소
스펙트럼에서 닫힌집합 $p^{-1}(X')\cap q^{-1}(Y')$에 대응한다
(\hyperref[I.1.2.2.1-ko]{1.2.2.1}과
\hyperref[I.1.1.2-ko]{1.1.2 (iii)}). 이로써 증명이 끝난다.

\begin{corollary}[4.3.2]
\label{I.4.3.2-ko}
$f:X\to Y$가 몰입 사상(각각 열린 몰입 사상, 닫힌 몰입 사상)이면서
$S$-사상이면, 밑준스킴의 임의의 확장 $S'\to S$에 대하여
$f_{(S')}$는 몰입 사상(각각 열린 몰입 사상, 닫힌 몰입 사상)이다.
\end{corollary}

\subsection{부분준스킴의 역상}
\label{subsection:I.4.4-ko}

\begin{proposition}[4.4.1]
\label{I.4.4.1-ko}
$f:X\to Y$를 사상, $Y'$를 $Y$의 부분준스킴(각각 닫힌 부분준스킴,
열린집합 위에 유도된 준스킴), $j:Y'\to Y$를 포함 사상이라 하자.
그러면 사영 $p:X\times_Y Y'\to X$는 몰입 사상(각각 닫힌 몰입 사상,
열린 몰입 사상)이다. $p$에 대응하는 $X$의 부분준스킴은
$f^{-1}(Y')$를 바탕공간으로 갖는다. 더욱이 $j'$를 이 부분준스킴에서
$X$로 가는 포함 사상이라 하면, 사상 $h:Z\to X$에 대하여
$f\circ h:Z\to Y$가 $j$를 통하여 인수분해되기 위한 필요충분조건은
$h$가 $j'$를 통하여 인수분해되는 것이다.
\end{proposition}

$p=1_X\times_Y j$이므로(\hyperref[I.3.3.4-ko]{3.3.4}), 첫째 주장은
(\hyperref[I.4.3.1-ko]{4.3.1})에서 따른다. 둘째 주장은
(\hyperref[I.3.5.10-ko]{3.5.10})의 특수한 경우이다(여기서는 $X$와
$Y'$의 역할을 서로 바꾸어야 한다). 끝으로 $h':Z\to Y'$인 사상
$h'$에 대하여 $f\circ h=j\circ h'$이라면, 곱의 정의에서 어떤 사상
$u:Z\to X\times_Y Y'$에 대하여 $h=p\circ u$임이 따른다. 따라서
마지막 주장도 얻는다.

이렇게 정의한 $X$의 부분준스킴을 사상 $f$에 의한 $Y'$의
\emph{역상}이라 부른다. 이 용어는 앞서
\oldpage[I]{126}
(\hyperref[I.3.3.6-ko]{3.3.6})에서 더 일반적으로 도입한 용어와
일치한다. $f^{-1}(Y')$를 $X$의 부분준스킴으로 말할 때에는 언제나 이
부분준스킴을 뜻한다.

준스킴 $f^{-1}(Y')$와 $X$가 동일하면 $j'$는 항등사상이고, 따라서
모든 사상 $h:Z\to X$가 $j'$를 통하여 인수분해된다. 그러므로
\emph{사상 $f:X\to Y$는 이때 다음과 같이 인수분해된다}:
$X\xrightarrow{g}Y'\xrightarrow{j}Y$.

$y$가 $Y$의 \emph{닫힌} 점이고
$Y'=\operatorname{Spec}(k(y))$가 $\{y\}$를 바탕공간으로 갖는 $Y$의
가장 작은 닫힌 부분준스킴이면(\hyperref[I.4.1.9-ko]{4.1.9}), 닫힌
부분준스킴 $f^{-1}(Y')$는
(\hyperref[I.3.6.2-ko]{3.6.2})에서 정의한 \emph{섬유} $f^{-1}(y)$와
표준적으로 동형이며, 앞으로 둘을 동일시한다.

\begin{corollary}[4.4.2]
\label{I.4.4.2-ko}
$f:X\to Y$, $g:Y\to Z$를 두 사상, $h=g\circ f$를 그 합성이라 하자.
$Z$의 모든 부분준스킴 $Z'$에 대하여 $X$의 부분준스킴
$f^{-1}(g^{-1}(Z'))$와 $h^{-1}(Z')$는 동일하다.
\end{corollary}

이는 표준 동형사상
$X\times_Y(Y\times_Z Z')\xrightarrow{\sim}X\times_Z Z'$의 존재에서
따른다(\hyperref[I.3.3.9.1-ko]{3.3.9.1}).

\begin{corollary}[4.4.3]
\label{I.4.4.3-ko}
$X'$, $X''$를 $X$의 두 부분준스킴,
$j':X'\to X$, $j'':X''\to X$를 포함 사상이라 하자. 그러면
${j'}^{-1}(X'')$와 ${j''}^{-1}(X')$는 둘 다 부분준스킴 사이의
순서관계에 관한 $X'$와 $X''$의 하한 $\inf(X',X'')$과 같고, 이
하한은 $X'\times_X X''$와 표준적으로 동형이다.
\end{corollary}

이는 (\hyperref[I.4.4.1-ko]{4.4.1})과
(\hyperref[I.4.1.10-ko]{4.1.10})에서 곧바로 따른다.

\begin{corollary}[4.4.4]
\label{I.4.4.4-ko}
$f:X\to Y$를 사상, $Y'$, $Y''$를 $Y$의 두 부분준스킴이라 하자.
그러면
\[
  f^{-1}(\inf(Y',Y''))=\inf(f^{-1}(Y'),f^{-1}(Y'')).
\]
\end{corollary}

이는 $(X\times_Y Y')\times_X(X\times_Y Y'')$와
$X\times_Y(Y'\times_Y Y'')$ 사이의 표준 동형사상이 존재함에서
따른다(\hyperref[I.3.3.9.1-ko]{3.3.9.1}).

\begin{proposition}[4.4.5]
\label{I.4.4.5-ko}
$f:X\to Y$를 사상, $Y'$를 $\mathscr{O}_Y$의 준연접 아이디얼층
$\mathscr{K}$로 정의된 $Y$의 닫힌 부분준스킴이라 하자
(\hyperref[I.4.1.3-ko]{4.1.3}). 그러면 $X$의 닫힌 부분준스킴
$f^{-1}(Y')$는 $\mathscr{O}_X$의 준연접 아이디얼층
$f^*(\mathscr{K})\mathscr{O}_X$로 정의된다.
\end{proposition}

{\raggedright
문제는 $X$와 $Y$ 위에서 자명하게 국소적이다. 따라서 $A$가
$B$-대수이고 $\mathfrak{K}$가 $B$의 아이디얼이면
$A\otimes_B(B/\mathfrak{K})=A/\mathfrak{K}A$임을 관찰하고
(\hyperref[I.1.6.9-ko]{1.6.9})를 적용하면 충분하다.\par}

\begin{corollary}[4.4.6]
\label{I.4.4.6-ko}
$X'$를 $\mathscr{O}_X$의 준연접 아이디얼층 $\mathscr{J}$로 정의된
$X$의 닫힌 부분준스킴, $i:X'\to X$를 포함 사상이라 하자. $f$를
$X'$에 제한한 사상 $f\circ i$가 포함 사상 $j:Y'\to Y$를 통하여
인수분해되기 위한(다시 말해 어떤 사상 $g:X'\to Y'$에 대하여
$j\circ g$로 인수분해되기 위한) 필요충분조건은
$f^*(\mathscr{K})\mathscr{O}_X\subset\mathscr{J}$인 것이다.
\end{corollary}

이는 명제 (\hyperref[I.4.4.1-ko]{4.4.1})을 $i$에 적용하되
(\hyperref[I.4.4.5-ko]{4.4.5})를 고려하면 된다.

\subsection{국소 몰입 사상과 국소 동형사상}
\label{subsection:I.4.5-ko}

\begin{definition}[4.5.1]
\label{I.4.5.1-ko}
준스킴의 사상 $f:X\to Y$가 \emph{점 $x\in X$에서 국소 몰입
사상}이라는 것은 $X$에서 $x$의 열린 근방 $U$와 $Y$에서 $f(x)$의
열린 근방 $V$가 존재하여, $U$ 위에 유도된 준스킴으로 $f$를 제한한
사상이 $U$에서 $V$ 위에 유도된 준스킴으로 가는 닫힌 몰입 사상인
것을 뜻한다. $f$가 $X$의 모든 점에서 국소 몰입 사상이면 $f$를 국소
몰입 사상이라 한다.
\end{definition}

\begin{definition}[4.5.2]
\label{I.4.5.2-ko}
사상 $f:X\to Y$가
\oldpage[I]{127}
점 $x\in X$에서 국소 동형사상이라는 것은 $X$에서 $x$의 열린 근방
$U$가 존재하여, $U$ 위에 유도된 준스킴으로 $f$를 제한한 사상이
$U$에서 $Y$로 가는 열린 몰입 사상인 것을 뜻한다. $f$가 $X$의 모든
점에서 국소 동형사상이면 $f$를 국소 동형사상이라 한다.
\end{definition}

\begin{env}[4.5.3]
\label{I.4.5.3-ko}
따라서 몰입 사상(각각 닫힌 몰입 사상) $f:X\to Y$는 국소 몰입
사상이면서 $X$의 바탕공간에서 $Y$의 한 부분집합(각각 닫힌
부분집합) 위로 가는 위상동형사상인 사상으로 특징지을 수 있다. 열린
몰입 사상 $f$는 바탕사상이 \emph{단사인} 국소 동형사상으로 특징지을
수 있다.
\end{env}

\begin{proposition}[4.5.4]
\label{I.4.5.4-ko}
$X$를 기약 준스킴, $f:X\to Y$를 바탕사상이 단사인 우세 사상이라
하자. $f$가 국소 몰입 사상이면 $f$는 몰입 사상이고 $f(X)$는 $Y$에서
열려 있다.
\end{proposition}

실제로 $x\in X$라 하고, $U$를 $x$의 열린 근방, $V$를 $Y$에서
$f(x)$의 열린 근방이라 하여 $f$의 $U$로의 제한이 $V$ 안으로의 닫힌
몰입 사상이 되게 하자. $U$는 $X$에서 조밀하고 가정에 따라 $f(U)$는
$Y$에서 조밀하므로 $f(U)=V$이며, $f$는 $U$에서 $V$ 위로 가는
위상동형사상이다. $f$의 바탕사상이 단사라는 가정에서
$f^{-1}(V)=U$가 따르므로 명제를 곧바로 얻는다.

\begin{proposition}[4.5.5]
\label{I.4.5.5-ko}
\begin{enumerate}
  \item[\rm(i)] 두 국소 몰입 사상(각각 두 국소 동형사상)의 합성은
    국소 몰입 사상(각각 국소 동형사상)이다.
  \item[\rm(ii)] $f:X\to X'$, $g:Y\to Y'$를 두 $S$-사상이라 하자.
    $f$와 $g$가 국소 몰입 사상(각각 국소 동형사상)이면
    $f\times_S g$도 그러하다.
  \item[\rm(iii)] $S$-사상 $f$가 국소 몰입 사상(각각 국소
    동형사상)이면, 밑준스킴의 모든 확장 $S'\to S$에 대하여
    $f_{(S')}$도 그러하다.
\end{enumerate}
\end{proposition}

(\hyperref[I.3.5.1-ko]{3.5.1})에 따라 (i)와 (ii)만 증명하면 충분하다.

(i)는 닫힌 몰입 사상(각각 열린 몰입 사상)의 추이성
(\hyperref[I.4.2.5-ko]{4.2.5})과 다음 사실에서 곧바로 따른다. $f$가
$X$에서 $Y$의 닫힌 부분집합 위로 가는 위상동형사상이면, 모든
열린집합 $U\subset X$에 대하여 $f(U)$는 $f(X)$에서 열려 있으므로
$f(U)=V\cap f(X)$가 되는 $Y$의 열린 부분집합 $V$가 존재하고, 따라서
$f(U)$는 $V$에서 닫혀 있다.

(ii)를 증명하자. $z\in X\times_S Y$를 임의의 점으로 택하고
$z'=(f\times_S g)(z)\in X'\times_S Y'$라 하자. $p$, $q$를
$X\times_S Y$의 사영, $p'$, $q'$를 $X'\times_S Y'$의 사영이라 하자.
가정에 따라 $x=p(z)$, $x'=p'(z')$, $y=q(z)$, $y'=q'(z')$의 열린
근방 $U$, $U'$, $V$, $V'$가 각각 존재하여, $f$와 $g$를 각각 $U$와
$V$로 제한한 사상이 각각 $U'$와 $V'$ 안으로의 닫힌 몰입 사상(각각
열린 몰입 사상)이 된다. $U\times_S V$와 $U'\times_S V'$의
바탕공간은 각각 $z$와 $z'$의 열린 근방
$p^{-1}(U)\cap q^{-1}(V)$와
${p'}^{-1}(U')\cap{q'}^{-1}(V')$로 동일시되므로
(\hyperref[I.3.2.7-ko]{3.2.7}), 명제는
(\hyperref[I.4.3.1-ko]{4.3.1})에서 따른다.
% ===== END EXACT INLINE INPUT: c1s4.tex =====


% ===== BEGIN EXACT INLINE INPUT: c1s5.tex =====
\section{축소 준스킴; 분리 조건}
\label{section:I.5-ko}

\subsection{축소 준스킴}
\label{subsection:I.5.1-ko}

\begin{proposition}[5.1.1]
\label{I.5.1.1-ko}
$(X,\mathscr{O}_X)$를 준스킴, $\mathscr{B}$를 준연접
$\mathscr{O}_X$-대수라 하자. 모든 $x\in X$에서 줄기 $\mathscr{N}_x$가
환 $\mathscr{B}_x$의 멱영근기가 되는 유일한 준연접
$\mathscr{O}_X$-가군 $\mathscr{N}$이 존재한다. $X$가 아핀이고 따라서
$\mathscr{B}=\widetilde{B}$이며 $B$가 $A(X)$ 위의 대수이면,
$\mathscr{N}=\widetilde{\mathfrak{N}}$이다. 여기서 $\mathfrak{N}$은
$B$의 멱영근기이다.
\end{proposition}

\oldpage[I]{128}
문제는 국소적이므로 마지막 주장만 증명하면 된다.
$\widetilde{\mathfrak{N}}$은 준연접 $\mathscr{O}_X$-가군이고
(\hyperref[I.1.4.1-ko]{1.4.1}), 점 $x\in X$에서 그 줄기는 분수환
$B_x$의 아이디얼 $\mathfrak{N}_x$이다. 따라서 반대쪽 포함관계는
자명하므로, $B_x$의 멱영근기가 $\mathfrak{N}_x$에 포함됨을 보이면
충분하다. $z\in B$, $s\notin\mathfrak{j}_x$이고 $z/s$가 $B_x$의
멱영원이라고 하자. 가정에 따라 $(z/s)^k=0$인 정수 $k$가 존재하고,
이는 $tz^k=0$인 $t\notin\mathfrak{j}_x$가 존재한다는 뜻이다. 그러면
$(tz)^k=0$이고, 따라서 $z/s=(tz)/(ts)$는 $\mathfrak{N}_x$에 속한다.

이렇게 정의한 준연접 $\mathscr{O}_X$-가군 $\mathscr{N}$을
$\mathscr{O}_X$-대수 $\mathscr{B}$의 \emph{멱영근기}라 부른다. 특히
$\mathscr{O}_X$의 멱영근기를 $\mathscr{N}_X$로 쓴다.

\begin{corollary}[5.1.2]
\label{I.5.1.2-ko}
$X$를 준스킴이라 하자. 아이디얼층 $\mathscr{N}_X$로 정의된 $X$의
닫힌 부분준스킴은 $X$를 바탕공간으로 갖는 유일한 축소 부분준스킴
(\hyperref[0.4.1.4-ko]{0, 4.1.4})이다. 또한 이는 $X$를 바탕공간으로
갖는 $X$의 가장 작은 부분준스킴이다.
\end{corollary}

$\mathscr{N}_X$로 정의된 닫힌 부분준스킴 $Y$의 구조층은
$\mathscr{O}_X/\mathscr{N}_X$이다. 모든 $x\in X$에 대하여
$\mathscr{N}_x\ne\mathscr{O}_x$이므로 $Y$가 축소이고 $X$를
바탕공간으로 갖는다는 것은 자명하다. 나머지 주장을 증명하기 위하여,
$X$를 바탕공간으로 갖는 $X$의 부분준스킴 $Z$가 모든 $x\in X$에
대하여 $\mathscr{J}_x\ne\mathscr{O}_x$인 아이디얼층 $\mathscr{J}$로
정의됨을 주목하자(\hyperref[I.4.1.3-ko]{4.1.3}). $X$가 아핀인 경우로
환원할 수 있다. $X=\operatorname{Spec}(A)$,
$\mathscr{J}=\widetilde{\mathfrak{J}}$라 하고 $\mathfrak{J}$를 $A$의
아이디얼이라 하자. 그러면 모든 $x\in X$에 대하여
$\mathfrak{J}\subset\mathfrak{j}_x$이다. 따라서 $\mathfrak{J}$는
$A$의 모든 소 아이디얼, 곧 그 교집합인 $A$의 멱영근기
$\mathfrak{N}$에 포함된다. 이로부터 $Y$가 $X$를 바탕공간으로 갖는
$X$의 가장 작은 부분준스킴임이 따른다
(\hyperref[I.4.1.9-ko]{4.1.9}). 더욱이 $Z$가 $Y$와 다르면 적어도 한
$x\in X$에 대하여 $\mathscr{J}_x\ne\mathscr{N}_x$이고, 따라서
(\hyperref[I.5.1.1-ko]{5.1.1})에 의하여 $Z$는 축소가 아니다.

\begin{definition}[5.1.3]
\label{I.5.1.3-ko}
준스킴 $X$에 대응하며 $X$를 바탕공간으로 갖는 유일한 축소
부분준스킴을 $X$에 대응하는 축소 준스킴이라 하고
$X_{\mathrm{red}}$로 쓴다.
\end{definition}

그러므로 준스킴 $X$가 축소라는 것은 $X=X_{\mathrm{red}}$임을 뜻한다.

\begin{proposition}[5.1.4]
\label{I.5.1.4-ko}
환 $A$의 소 스펙트럼이 축소 준스킴(각각 정역 준스킴)
(\hyperref[I.2.1.8-ko]{2.1.8})이기 위한 필요충분조건은 $A$가
축소환(각각 정역)인 것이다.
\end{proposition}

실제로 (\hyperref[I.5.1.1-ko]{5.1.1})에서
$X=\operatorname{Spec}(A)$가 축소이기 위한 필요충분조건이
$\mathscr{N}=(0)$임을 곧바로 얻는다. 정역에 관한 주장은
(\hyperref[I.1.1.13-ko]{1.1.13})에서 따른다.

정역의 영환이 아닌 모든 분수환은 정역이므로,
(\hyperref[I.5.1.4-ko]{5.1.4})에서 모든 \emph{국소 정역 준스킴}
$X$와 모든 $x\in X$에 대하여 $\mathscr{O}_x$가 \emph{정역}임이
따른다. $X$의 바탕공간이 \emph{국소 뇌터}이면 그 역도 성립한다.
실제로 이때 $X$는 축소이다. $X$의 아핀 열린집합 $U$가 뇌터 공간이면
$U$의 기약 성분은 유한개뿐이고, 따라서 그 환 $A$의 극소 소
아이디얼도 유한개뿐이다(\hyperref[I.1.1.14-ko]{1.1.14}). 이 성분들
$U_i$ 가운데 둘이 공통점 $x$를 가지면 $\mathscr{O}_x$는 서로 다른
극소 소 아이디얼을 적어도 둘 가지므로 정역이 아니다. 따라서
$U_i$들은 서로소인 열린집합이고, 각각 정역 준스킴이다.

\begin{env}[5.1.5]
\label{I.5.1.5-ko}
$f=(\psi,\theta):X\to Y$를 준스킴의 사상이라 하자. 준동형
\oldpage[I]{129}
$\theta_x^\sharp:\mathscr{O}_{\psi(x)}\to\mathscr{O}_x$는
$\mathscr{O}_{\psi(x)}$의 모든 멱영원을 $\mathscr{O}_x$의 멱영원으로
보낸다. 그러므로 몫을 취하면 $\theta^\sharp$에서 준동형
\[
  \omega:\psi^*(\mathscr{O}_Y/\mathscr{N}_Y)
  \longrightarrow \mathscr{O}_X/\mathscr{N}_X
\]
를 얻는다. 모든 $x\in X$에 대하여
$\omega_x:\mathscr{O}_{\psi(x)}/\mathscr{N}_{\psi(x)}
\to\mathscr{O}_x/\mathscr{N}_x$가 국소 준동형임은 명백하다. 따라서
$(\psi,\omega^b)$는 준스킴의 사상
$X_{\mathrm{red}}\to Y_{\mathrm{red}}$이다. 이를 $f_{\mathrm{red}}$로
쓰고 $f$에 대응하는 \emph{축소} 사상이라 부른다. 두 사상
$f:X\to Y$, $g:Y\to Z$에 대하여
$(g\circ f)_{\mathrm{red}}=g_{\mathrm{red}}\circ f_{\mathrm{red}}$임은
자명하다. 따라서 $X_{\mathrm{red}}$는 $X$에 관한 \emph{공변 함자}로
정의되었다.

앞의 정의에서 도식
\[
  \xymatrix{
    X_{\mathrm{red}}\ar[r]^{f_{\mathrm{red}}}\ar[d] &
    Y_{\mathrm{red}}\ar[d]\\
    X\ar[r]_f & Y
  }
\]
이 가환임이 따른다. 여기서 세로 화살표는 포함 사상이다. 다시 말해
$X_{\mathrm{red}}\to X$는 \emph{함자적}인 사상이다. 특히 $X$가
\emph{축소}이면 모든 사상 $f:X\to Y$는
$X\xrightarrow{f_{\mathrm{red}}}Y_{\mathrm{red}}\to Y$로
인수분해된다. 다시 말해 $f$는 포함 사상
$Y_{\mathrm{red}}\to Y$를 통하여 \emph{인수분해된다}.
\end{env}

\begin{proposition}[5.1.6]
\label{I.5.1.6-ko}
$f:X\to Y$를 사상이라 하자. $f$가 전사(각각 라디시엘 사상, 몰입
사상, 닫힌 몰입 사상, 열린 몰입 사상, 국소 몰입 사상, 국소
동형사상)이면 $f_{\mathrm{red}}$도 그러하다. 역으로
$f_{\mathrm{red}}$가 전사(각각 라디시엘 사상)이면 $f$도 그러하다.
\end{proposition}

$f$가 전사인 경우 명제는 자명하다. $f$가 라디시엘 사상이면, 모든
$x\in X$에 대하여 체 $k(x)$가 준스킴 $X$와 $X_{\mathrm{red}}$에서
동일하다는 사실에서 명제가 따른다
(\hyperref[I.3.5.8-ko]{3.5.8}). 끝으로 $f=(\psi,\theta)$가 몰입 사상,
닫힌 몰입 사상 또는 국소 몰입 사상(각각 열린 몰입 사상 또는 국소
동형사상)이면, $\theta_x^\sharp$가 전사(각각 전단사)일 때
$\mathscr{O}_{\psi(x)}$와 $\mathscr{O}_x$의 멱영근기로 몫을 취하여
얻는 준동형도 전사(각각 전단사)라는 사실에서 명제가 따른다
(\hyperref[I.5.1.2-ko]{5.1.2}와
\hyperref[I.4.2.2-ko]{4.2.2}; 또한
(\agref{I.5.5.12-ko}{5.5.12}) 참조).

\begin{proposition}[5.1.7]
\label{I.5.1.7-ko}
$X$, $Y$가 두 $S$-준스킴이면, 준스킴
$X_{\mathrm{red}}\times_{S_{\mathrm{red}}}Y_{\mathrm{red}}$와
$X_{\mathrm{red}}\times_S Y_{\mathrm{red}}$는 동일하며,
$X\times_S Y$와 같은 바탕공간을 갖는 $X\times_S Y$의 부분준스킴과 표준적으로
동일시된다.
\end{proposition}

$X_{\mathrm{red}}\times_S Y_{\mathrm{red}}$가 $X\times_S Y$와 같은
바탕공간을 갖는 $X\times_S Y$의 부분준스킴과 표준적으로 동일시되는 것은
(\hyperref[I.4.3.1-ko]{4.3.1})에서 따른다. 한편 $\varphi$와 $\psi$가
구조 사상 $X_{\mathrm{red}}\to S$, $Y_{\mathrm{red}}\to S$이면, 이들은
$S_{\mathrm{red}}$를 통하여 인수분해된다
(\hyperref[I.5.1.5-ko]{5.1.5}). $S_{\mathrm{red}}\to S$는
단사사상이므로 첫째 주장은 (\hyperref[I.3.2.4-ko]{3.2.4})에서
따른다.

\begin{corollary}[5.1.8]
\label{I.5.1.8-ko}
준스킴 $(X\times_S Y)_{\mathrm{red}}$와
$(X_{\mathrm{red}}\times_{S_{\mathrm{red}}}Y_{\mathrm{red}})_{\mathrm{red}}$는
표준적으로 동일시된다.
\end{corollary}

이는 (\hyperref[I.5.1.2-ko]{5.1.2}와
\hyperref[I.5.1.7-ko]{5.1.7})에서 따른다.

$X$와 $Y$가 축소 준스킴이어도 $X\times_S Y$가 반드시 축소인 것은
아니다. 두 축소 대수의 텐서곱에 멱영원이 있을 수 있기 때문이다.

\begin{proposition}[5.1.9]
\label{I.5.1.9-ko}
\oldpage[I]{130}
$X$를 준스킴, $\mathscr{J}$를 어떤 정수 $n>0$에 대하여
$\mathscr{J}^n=0$인 $\mathscr{O}_X$의 준연접 아이디얼층이라 하자.
$X_0$를 $X$의 닫힌 부분준스킴
$(X,\mathscr{O}_X/\mathscr{J})$라 하자. $X$가 아핀 스킴이기 위한
필요충분조건은 $X_0$가 아핀 스킴인 것이다.
\end{proposition}

조건이 필요하다는 것은 자명하므로 충분함을 증명하자.
$X_k=(X,\mathscr{O}_X/\mathscr{J}^{k+1})$라 두면, $k$에 관한 귀납법으로
각 $X_k$가 아핀임을 증명하면 충분하다. 따라서 $\mathscr{J}^2=0$인
경우로 환원된다. 다음과 같이 놓자.
\[
  A=\Gamma(X,\mathscr{O}_X),\qquad
  A_0=\Gamma(X_0,\mathscr{O}_{X_0})
      =\Gamma(X,\mathscr{O}_X/\mathscr{J}).
\]
표준 준동형 $\mathscr{O}_X\to\mathscr{O}_X/\mathscr{J}$에서 환의
준동형 $\varphi:A\to A_0$를 얻는다. 아래에서 $\varphi$가
\emph{전사}임을 보일 것이다. 그러면 열
\[
  0\to\Gamma(X,\mathscr{J})\to\Gamma(X,\mathscr{O}_X)
  \to\Gamma(X,\mathscr{O}_X/\mathscr{J})\to 0
  \tag{5.1.9.1}\label{I.5.1.9.1-ko}
\]
은 \emph{완전}하다. 이 점을 안다고 하고, 그것이 명제를 함의함을
보이자. $\mathfrak{K}=\Gamma(X,\mathscr{J})$는 $A$의 제곱영
아이디얼이고, 따라서 $A_0=A/\mathfrak{K}$ 위의 가군이다. 가정에
따라 $X_0=\operatorname{Spec}(A_0)$이고, 바탕 위상공간 $X_0$와 $X$는
동일하므로 $\mathfrak{K}=\Gamma(X_0,\mathscr{J})$이다. 또한
$\mathscr{J}^2=0$이므로 $\mathscr{J}$는 준연접
$(\mathscr{O}_X/\mathscr{J})$-가군이다. 따라서
$\mathscr{J}\simeq\widetilde{\mathfrak{K}}$이고, 모든 $x\in X_0$에
대하여 $\mathfrak{K}_x=\mathscr{J}_x$이다
(\hyperref[I.1.4.1-ko]{1.4.1}).

이제 $X'=\operatorname{Spec}(A)$라 하고, 항등 사상
$A\to\Gamma(X,\mathscr{O}_X)$에 대응하는 준스킴의 사상
$f=(\psi,\theta):X\to X'$를 생각하자
(\hyperref[I.2.2.4-ko]{2.2.4}). $X$의 모든 아핀 열린집합 $V$에 대하여
도식
\[
  \xymatrix{
    A\ar[r]\ar[d] & \Gamma(V,\mathscr{O}_X|V)\ar[d]\\
    A_0=A/\mathfrak{K}\ar[r] & \Gamma(V,\mathscr{O}_{X_0}|V)
  }
\]
은 가환이다. 따라서 도식
\[
  \xymatrix{
    X' & X\ar[l]_f\\
    X'_0\ar[u]^{j'} & X_0\ar[l]^{f_0}\ar[u]_j
  }
\]
도 가환이다. 여기서 $X'_0$는 준연접 아이디얼층
$\widetilde{\mathfrak{K}}$로 정의된 $X'$의 닫힌 부분준스킴이고,
$j$, $j'$는 표준 포함 사상이다. $X_0$는 아핀이므로 $f_0$는
동형사상이다. $j$와 $j'$의 바탕 연속사상은 항등사상이므로 먼저
$\psi:X\to X'$가 위상동형사상임을 알 수 있다. 또한
$\mathfrak{K}_x=\mathscr{J}_x$라는 관계에서
$\theta^\sharp:\psi^*(\mathscr{O}_{X'})\to\mathscr{O}_X$의 제한이
$\psi^*(\widetilde{\mathfrak{K}})$에서 $\mathscr{J}$ 위로 가는
\emph{동형사상}임이 따른다. 한편 몫을 취하면 $f_0$가 동형사상이므로
$\theta^\sharp$에서 \emph{동형사상}
$\psi^*(\mathscr{O}_{X'}/\widetilde{\mathfrak{K}})
\to\mathscr{O}_X/\mathscr{J}$를 얻는다. 따라서 5-보조정리
$(\mathrm{M},\mathrm{I},1.1)$에서 $\theta^\sharp$ 자체가
동형사상임이 곧바로 따른다. 그러므로 $f$는 \emph{동형사상}이고
$X$는 아핀이다.

이제 (\hyperref[I.5.1.9.1-ko]{5.1.9.1})의 완전성을 증명하면
충분하다. 이는

\oldpage[I]{131}
$H^1(X,\mathscr{J})=0$에서 따른다. 그런데
$H^1(X,\mathscr{J})=H^1(X_0,\mathscr{J})$이고, 앞에서
$\mathscr{J}$가 준연접 $\mathscr{O}_{X_0}$-가군임을 보았다. 따라서
우리의 주장은 다음 보조정리에서 따른다.

\begin{lemma}[5.1.9.2]
\label{I.5.1.9.2-ko}
$Y$가 아핀 스킴이고 $\mathscr{F}$가 준연접
$\mathscr{O}_Y$-가군이면 $H^1(Y,\mathscr{F})=0$이다.
\end{lemma}

이 보조정리는 제III장 §1에서, 모든 $i>0$에 대하여
$H^i(Y,\mathscr{F})=0$이라는 더 일반적인 정리의 결과로 증명할
것이다. 독립적인 증명을 주기 위하여, $H^1(Y,\mathscr{F})$가
$\mathscr{F}$에 의한 $\mathscr{O}_Y$-가군 $\mathscr{O}_Y$의 확장류로
이루어진 가군
$\operatorname{Ext}^1_{\mathscr{O}_Y}
(Y;\mathscr{O}_Y,\mathscr{F})$와 동일시됨을 주목하자
$(\mathrm{T},4.2.3)$. 따라서 그러한 확장 $\mathscr{G}$가 자명함을
보이면 충분하다. 모든 $y\in Y$에 대하여 $Y$에서 $y$의 근방 $V$가
존재하여 $\mathscr{G}|V$가
$\mathscr{F}|V\oplus\mathscr{O}_Y|V$와 동형이다
(\hyperref[0.5.4.9-ko]{0, 5.4.9}). 따라서 $\mathscr{G}$는 준연접
$\mathscr{O}_Y$-가군이다. $A$를 $Y$의 환이라 하면
$\mathscr{F}=\widetilde{M}$, $\mathscr{G}=\widetilde{N}$이고, 여기서
$M$, $N$은 $A$-가군이다. 가정에 따라 $N$은 $A$-가군 $M$에 의한
$A$-가군 $A$의 확장이다(\hyperref[I.1.3.11-ko]{1.3.11}). 이 확장은
반드시 자명하므로 보조정리가 증명되고, 따라서
(\hyperref[I.5.1.9-ko]{5.1.9})도 증명된다.

\begin{corollary}[5.1.10]
\label{I.5.1.10-ko}
$X$를 $\mathscr{N}_X$가 멱영인 준스킴이라 하자. $X$가 아핀 스킴이기
위한 필요충분조건은 $X_{\mathrm{red}}$가 아핀 스킴인 것이다.
\end{corollary}

\subsection{주어진 바탕공간을 갖는 부분준스킴의 존재}
\label{subsection:I.5.2-ko}

\begin{proposition}[5.2.1]
\label{I.5.2.1-ko}
준스킴 $X$의 바탕공간의 모든 국소 닫힌 부분공간 $Y$에 대하여,
$Y$를 바탕공간으로 갖는 $X$의 축소 부분준스킴이 유일하게 존재한다.
\end{proposition}

유일성은 (\hyperref[I.5.1.2-ko]{5.1.2})에서 따르므로, 문제의
부분준스킴의 존재만 증명하면 된다.

$X$가 환 $A$의 아핀 준스킴이고 $Y$가 $X$에서 닫혀 있으면 명제는
곧바로 따른다. 실제로 $\mathfrak{j}(Y)$는
$V(\mathfrak{a})=Y$를 만족하는 아이디얼 $\mathfrak{a}\subset A$ 가운데
가장 큰 것이며 자기 근기와 같다
(\hyperref[I.1.1.4-ko]{1.1.4 (i)}). 따라서
$A/\mathfrak{j}(Y)$는 축소환이다.

일반적인 경우, $U\cap Y$가 $U$에서 닫히도록 하는 모든 아핀
열린집합 $U\subset X$에 대하여, $A(U)$의 아이디얼
$\mathfrak{j}(U\cap Y)$에 대응하는 아이디얼층으로 정의된 $U$의 닫힌
부분준스킴 $Y_U$를 생각하자. 이 부분준스킴은 축소이다. 이제 $V$가
$U$에 포함되는 $X$의 아핀 열린집합이면, $Y_V$는 $V\cap Y$ 위에서
$Y_U$에 의해 \emph{유도된} 부분준스킴임을 보이자. 그런데 이렇게
유도된 준스킴은 $V$의 축소 닫힌 부분준스킴이고 $V\cap Y$를
바탕공간으로 가지므로, $Y_V$의 유일성에 의해 주장이 따른다.

\begin{proposition}[5.2.2]
\label{I.5.2.2-ko}
$X$를 축소 준스킴이라 하고, $f:X\to Y$를 사상, $Z$를
$f(X)\subset Z$를 만족하는 $Y$의 닫힌 부분준스킴이라 하자. 그러면 $f$는
$X\xrightarrow{g}Z\xrightarrow{j}Y$로 인수분해되며, 여기서 $j$는 포함
사상이다.
\end{proposition}

가정에 따라 $X$의 닫힌 부분준스킴 $f^{-1}(Z)$의 바탕공간은 $X$
전체이다 (\hyperref[I.4.4.1-ko]{4.4.1}). $X$가 축소이므로 이 닫힌
부분준스킴은 $X$와 동일하다 (\hyperref[I.5.1.2-ko]{5.1.2}). 따라서
명제는 (\hyperref[I.4.4.1-ko]{4.4.1})에서 따른다.

\begin{corollary}[5.2.3]
\label{I.5.2.3-ko}
$X$를 준스킴 $Y$의 축소 부분준스킴이라 하자. $Z$가 $\overline{X}$를
바탕공간으로 갖는 $Y$의 축소 닫힌 부분준스킴이면, $X$는 $Z$의 어떤
열린집합 위에 유도된 부분준스킴이다.
\end{corollary}

\oldpage[I]{132}
실제로 $Y$의 어떤 열린집합 $U$에 대하여, 바탕공간으로서
$X=U\cap\overline{X}$이다. (\hyperref[I.5.2.2-ko]{5.2.2})에 따라 $X$는
$Z$의 축소 부분준스킴이고, 유일성
(\hyperref[I.5.2.1-ko]{5.2.1})에 따라 부분준스킴 $X$는 열린
부분공간 $X$ 위에서 $Z$에 의해 유도된다.

\begin{corollary}[5.2.4]
\label{I.5.2.4-ko}
$f:X\to Y$를 사상이라 하고, $X'$ (각각 $Y'$)를 $\mathscr{O}_X$의
(각각 $\mathscr{O}_Y$의) 준연접 아이디얼층 $\mathscr{J}$
(각각 $\mathscr{K}$)로 정의된 $X$의 (각각 $Y$의) 닫힌
부분준스킴이라 하자. $X'$가 축소이고 $f(X')\subset Y'$라고 하자.
그러면 $f^*(\mathscr{K})\mathscr{O}_X\subset\mathscr{J}$이다.
\end{corollary}

(\hyperref[I.5.2.2-ko]{5.2.2})에 따라 $f$를 $X'$에 제한한 사상이
$X'\to Y'\to Y$로 인수분해되므로,
(\hyperref[I.4.4.6-ko]{4.4.6})을 적용하면 충분하다.

\subsection{대각선; 사상의 그래프}
\label{subsection:I.5.3-ko}

\begin{env}[5.3.1]
\label{I.5.3.1-ko}
$X$를 $S$-준스킴이라 하자. $X$에서 $X\times_S X$로 가는
$S$-사상 $(1_X,1_X)_S$, 다시 말해
\[
  p_1\circ\Delta_X=p_2\circ\Delta_X=1_X
  \tag{5.3.1.1}\label{I.5.3.1.1-ko}
\]
을 만족하는 유일한 $S$-사상 $\Delta_X$를 $X$의 \emph{대각 사상}이라
하고 $\Delta_{X|S}$ 또는 $\Delta_X$, 혼동의 우려가 없으면 간단히
$\Delta$로 쓴다. 여기서 $p_1$, $p_2$는 $X\times_S X$의 사영이다
(정의 \hyperref[I.3.2.1-ko]{3.2.1}). $f:T\to X$, $g:T\to Y$가 두
$S$-사상이면 곧바로
\[
  (f,g)_S=(f\times_S g)\circ\Delta_{T|S}
  \tag{5.3.1.2}\label{I.5.3.1.2-ko}
\]
임을 확인할 수 있다.

독자는 앞의 정의와 (\hyperref[I.5.3.1-ko]{5.3.1})부터 (5.3.8)까지
서술한 결과가 어느 범주에서도
\emph{그 안에 나타나는 곱들이 그 범주에 존재하기만 하면} 성립함을
주목할 것이다.
\end{env}

\begin{proposition}[5.3.2]
\label{I.5.3.2-ko}
$X$, $Y$를 두 $S$-준스킴이라 하자. 곱
$(X\times Y)\times(X\times Y)$를
$(X\times X)\times(Y\times Y)$와 표준적으로 동일시하면, 사상
$\Delta_{X\times Y}$는 $\Delta_X\times\Delta_Y$와 동일시된다.
\end{proposition}

실제로 $p_1$, $q_1$이 각각 $X\times X\to X$,
$Y\times Y\to Y$인 첫째 사영이면,
$(X\times Y)\times(X\times Y)\to X\times Y$인 첫째 사영은
$p_1\times q_1$과 동일시되고
\[
  (p_1\times q_1)\circ(\Delta_X\times\Delta_Y)
  =(p_1\circ\Delta_X)\times(q_1\circ\Delta_Y)=1_{X\times Y}
\]
이다. 둘째 사영들에 대해서도 마찬가지이다.

\begin{corollary}[5.3.4]
\label{I.5.3.4-ko}
밑준스킴의 모든 확장 $S'\to S$에 대하여,
$\Delta_{X_{(S')}}$는 $(\Delta_X)_{(S')}$와 표준적으로 동일시된다.
\end{corollary}

실제로 $(X\times_S X)_{(S')}$가
$X_{(S')}\times_{S'}X_{(S')}$와 표준적으로 동일시됨을 주목하면
충분하다 (\hyperref[I.3.3.10-ko]{3.3.10}).

\begin{proposition}[5.3.5]
\label{I.5.3.5-ko}
$X$, $Y$를 두 $S$-준스킴이라 하고, $\varphi:S\to T$를 모든
$S$-준스킴에 $T$-준스킴의 구조를 주는 사상이라 하자. $f:X\to S$,
$g:Y\to S$를 구조 사상, $p$, $q$를 $X\times_S Y$의 사영,
$\pi=f\circ p=g\circ q$를 구조 사상 $X\times_S Y\to S$라 하자.
그러면 그림
\[
  \xymatrix{
    X\times_S Y\ar[r]^{(p,q)_T}\ar[d]_\pi &
    X\times_T Y\ar[d]^{f\times_T g}\\
    S\ar[r]_{\Delta_{S|T}} &
    S\times_T S
  }
  \tag{5.3.5.1}\label{I.5.3.5.1-ko}
\]
\oldpage[I]{133}
은 가환이고, $X\times_S Y$를 $(S\times_T S)$-준스킴 $S$와
$X\times_T Y$의 곱과 동일시한다. 이때 사영들은 각각 $\pi$와
$(p,q)_T$로 동일시된다.
\end{proposition}

(\hyperref[I.3.4.3-ko]{3.4.3})에 따라, $Z$를 임의의 $T$-준스킴이라
하고 $X$, $Y$, $S$를 각각 $X(Z)_T$, $Y(Z)_T$, $S(Z)_T$로 바꾸어
얻는 \emph{집합의 범주 안에서}의 대응하는 명제를 증명하면 된다.
집합의 범주에서는 확인이 즉각적이므로 독자에게 맡긴다.

\begin{corollary}[5.3.6]
\label{I.5.3.6-ko}
$P=S\times_T S$라 놓으면 사상 $(p,q)_T$는
$1_{X\times_T Y}\times_P\Delta_S$와 동일시된다.
\end{corollary}

이는 (\hyperref[I.5.3.5-ko]{5.3.5})와
(\hyperref[I.3.3.4-ko]{3.3.4})에서 따른다.

\begin{corollary}[5.3.7]
\label{I.5.3.7-ko}
$f:X\to Y$가 $S$-사상이면 그림
\[
  \xymatrix{
    X\ar[r]^{(1_X,f)_S}\ar[d]_f &
    X\times_S Y\ar[d]^{f\times_S 1_Y}\\
    Y\ar[r]_{\Delta_Y} &
    Y\times_S Y
  }
\]
은 가환이고, $X$를 $(Y\times_S Y)$-준스킴 $Y$와
$X\times_S Y$의 곱과 동일시한다.
\end{corollary}

(\hyperref[I.5.3.5-ko]{5.3.5})에서 $S$를 $Y$로, $T$를 $S$로
바꾸고 $X\times_Y Y=X$임을 주목하면 충분하다
(\hyperref[I.3.3.3-ko]{3.3.3}).

\begin{proposition}[5.3.8]
\label{I.5.3.8-ko}
$f:X\to Y$가 준스킴의 단사사상이기 위한 필요충분조건은
$\Delta_{X|Y}$가 $X$에서 $X\times_Y X$ 위로 가는 동형사상인 것이다.
\end{proposition}

실제로 $f$가 단사사상이라는 것은 모든 $Y$-준스킴 $Z$에 대하여
대응하는 사상 $f':X(Z)_Y\to Y(Z)_Y$가 단사라는 뜻이다.
$Y(Z)_Y$에는 원소가 하나뿐이므로, 이는 $X(Z)_Y$의 원소가 많아야
하나라는 뜻이다. 이는 다시 $X(Z)_Y\times X(Z)_Y$가
$X(Z)_Y$와 표준적으로 동형이라는 말과 같다. 첫째 집합은
(\hyperref[I.3.4.3.1-ko]{3.4.3.1})에 따라
$(X\times_Y X)(Z)_Y$이므로, 이는 $\Delta_{X|Y}$가 동형사상이라는
뜻이다.

\begin{proposition}[5.3.9]
\label{I.5.3.9-ko}
대각 사상 $\Delta_X$는 $X$에서 $X\times_S X$ 안으로 가는 몰입
사상이다.
\end{proposition}

(\hyperref[I.4.2.4-ko]{4.2.4 (a)})를 사용하면 충분하다. 실제로
$X$의 모든 $x$와 $x$의 모든 아핀 근방 $U$에 대하여,
$U\times_S U$는 $\Delta_X(x)$의 아핀 근방이다.
(\hyperref[I.5.3.16-ko]{5.3.16})의 증명은 정의
(\hyperref[I.5.3.1.1-ko]{5.3.1.1})만을 사용하므로, $S=\operatorname{Spec}(B)$,
$X=\operatorname{Spec}(A)$가 아핀 스킴인 경우만 증명하면 된다. 이때
(\hyperref[I.5.3.1.1-ko]{5.3.1.1})에서 $\Delta_X$가
\[
A\otimes_B A\longrightarrow A,\qquad x\otimes y\longmapsto xy
\]
인 표준 준동형에 대응함이 곧바로 따른다. 이 준동형은 전사이므로,
$\Delta_X$는 이 경우 닫힌 몰입 사상이다
(\hyperref[I.4.2.3-ko]{4.2.3}). 이것으로 증명이 끝난다.

몰입 사상 $\Delta_X$에 대응하는 $X\times_S X$의 부분준스킴
(\hyperref[I.4.2.1-ko]{4.2.1})을 $X\times_S X$의
\emph{대각선}이라 한다.

\begin{corollary}[5.3.10]
\label{I.5.3.10-ko}
(\hyperref[I.5.3.5-ko]{5.3.5})의 가정 아래에서 $(p,q)_T$는 몰입
사상이다.
\end{corollary}

이는 (\hyperref[I.5.3.6-ko]{5.3.6})과
(\hyperref[I.4.3.1-ko]{4.3.1})에서 따른다.

(\hyperref[I.5.3.5-ko]{5.3.5})의 가정 아래에서 $(p,q)_T$를
$X\times_S Y$에서 $X\times_T Y$ 안으로 가는 \emph{표준 몰입}이라
한다.

\begin{corollary}[5.3.11]
\label{I.5.3.11-ko}
$X$, $Y$를 두 $S$-준스킴, $f:X\to Y$를 $S$-사상이라 하자. 그러면
$f$의 그래프 사상 $\Gamma_f=(1_X,f)_S$
(\hyperref[I.3.3.14-ko]{3.3.14})는 $X$에서 $X\times_S Y$ 안으로 가는
몰입 사상이다.
\end{corollary}

이는 (\hyperref[I.5.3.10-ko]{5.3.10})에서 $S$를 $Y$로, $T$를 $S$로
바꾼 특별한 경우이다 (\hyperref[I.5.3.7-ko]{5.3.7} 참조).

\oldpage[I]{134}
몰입 사상 $\Gamma_f$에 대응하는 $X\times_S Y$의 부분준스킴
(\hyperref[I.4.2.1-ko]{4.2.1})을 사상 $f$의 \emph{그래프}라 한다.
사상 $X\to Y$의 그래프인 $X\times_S Y$의 부분준스킴들은 다음
성질로 특징지어진다. 첫째 사영 $p_1:X\times_S Y\to X$를 그러한
부분준스킴 $G$에 제한한 사상은 $G$에서 $X$ 위로 가는
\emph{동형사상} $g$이고, 이때 $G$는 사상 $p_2\circ g^{-1}$의
그래프이다. 여기서 $p_2$는 사영 $X\times_S Y\to Y$이다.

특히 $X=S$로 취하면, $S$-사상 $S\to Y$, 곧 $Y$의 $S$-단면
(\hyperref[I.2.5.5-ko]{2.5.5})은 자신의 그래프 사상과 같다.
$S$-단면의 그래프인 $Y$의 부분준스킴들, 다시 말해 구조 사상
$Y\to S$의 제한에 의해 $S$와 동형인 부분준스킴들을 이 단면들의
\emph{상}, 또는 관용적으로 $Y$의 \emph{$S$-단면}이라고도 한다.

\begin{corollary}[5.3.12]
\label{I.5.3.12-ko}
(\hyperref[I.5.3.11-ko]{5.3.11})의 가정과 표기를 사용하자. 모든 사상
$g:S'\to S$에 대하여, $f'$을 $g$에 의한 $f$의 역상이라 하면
(\hyperref[I.3.3.7-ko]{3.3.7}), $\Gamma_{f'}$은 $g$에 의한
$\Gamma_f$의 역상이다.
\end{corollary}

이는 공식 (\hyperref[I.3.3.10.1-ko]{3.3.10.1})의 특별한 경우이다.

\begin{corollary}[5.3.13]
\label{I.5.3.13-ko}
$f:X\to Y$, $g:Y\to Z$를 두 사상이라 하자. $g\circ f$가 몰입
사상(각각 국소 몰입 사상)이면 $f$도 그러하다.
\end{corollary}

실제로 $f$는
\[
  X\xrightarrow{\Gamma_f}X\times_Z Y\xrightarrow{p_2}Y
\]
로 인수분해된다. 한편 $p_2$는 $(g\circ f)\times_Z 1_Y$와
동일시된다 (\hyperref[I.3.3.4-ko]{3.3.4}). $g\circ f$가 몰입
사상(각각 국소 몰입 사상)이면 $p_2$도 그러하다
(\hyperref[I.4.3.1-ko]{4.3.1}과 \hyperref[I.4.5.5-ko]{4.5.5}). 또한
$\Gamma_f$는 몰입 사상이므로 (\hyperref[I.5.3.11-ko]{5.3.11}),
(\hyperref[I.4.2.5-ko]{4.2.5}) (각각
\hyperref[I.4.5.5-ko]{4.5.5})에서 결론이 따른다.

\begin{corollary}[5.3.14]
\label{I.5.3.14-ko}
$j:X\to Y$, $g:X\to Z$를 두 $S$-사상이라 하자. $j$가 몰입
사상(각각 국소 몰입 사상)이면 $(j,g)_S$도 그러하다.
\end{corollary}

실제로 $p:Y\times_S Z\to Y$가 첫째 사영이면
$j=p\circ(j,g)_S$이고, (\hyperref[I.5.3.13-ko]{5.3.13})을 적용하면
충분하다.

\begin{proposition}[5.3.15]
\label{I.5.3.15-ko}
$f:X\to Y$가 $S$-사상이면 그림
\[
  \xymatrix{
    X\ar[r]^{\Delta_X}\ar[d]_f &
    X\times_S X\ar[d]^{f\times_S f}\\
    Y\ar[r]_{\Delta_Y} &
    Y\times_S Y
  }
  \tag{5.3.15.1}\label{I.5.3.15.1-ko}
\]
은 가환이다. 다시 말해 $\Delta_X$는 준스킴의 범주에서 함자적
사상이다.
\end{proposition}

확인은 즉각적이므로 독자에게 맡긴다.

\begin{corollary}[5.3.16]
\label{I.5.3.16-ko}
$X$가 $Y$의 부분준스킴이면 대각선 $\Delta_X(X)$는
$\Delta_Y(Y)$의 부분준스킴과 동일시되며, 그 바탕공간은
\[
  \Delta_Y(Y)\cap p_1^{-1}(X)=\Delta_Y(Y)\cap p_2^{-1}(X)
\]
와 동일시된다. 여기서 $p_1$, $p_2$는 $Y\times_S Y$의 사영이다.
\end{corollary}

포함 사상 $f:X\to Y$에 (\hyperref[I.5.3.15-ko]{5.3.15})를 적용하자.
그러면 $f\times_S f$는 몰입 사상이고, $X\times_S X$의 바탕공간을
$Y\times_S Y$의 부분공간
$p_1^{-1}(X)\cap p_2^{-1}(X)$와 동일시한다
(\hyperref[I.4.3.1-ko]{4.3.1}). 또한
$z\in\Delta_Y(Y)\cap p_1^{-1}(X)$이면 $z=\Delta_Y(y)$이고

\oldpage[I]{135}
$y=p_1(z)\in X$이다. 따라서 포함 사상 아래에서 $y=f(y)$로
동일시되며, 그림 (\hyperref[I.5.3.15.1-ko]{5.3.15.1})의 가환성에
따라 $z=\Delta_Y(f(y))$는 $\Delta_X(X)$에 속한다.

\begin{corollary}[5.3.17]
\label{I.5.3.17-ko}
$f_1:Y\to X$, $f_2:Y\to X$를 두 $S$-사상이라 하고, $y$를
$f_1(y)=f_2(y)=x$이면서 $f_1$, $f_2$에 대응하는 준동형
$k(x)\to k(y)$가 서로 같은 $Y$의 점이라 하자. 그러면
$f=(f_1,f_2)_S$라 할 때, 점 $f(y)$는 대각선
$\Delta_{X|S}(X)$에 속한다.
\end{corollary}

$f_i$ ($i=1,2$)에 대응하는 두 준동형 $k(x)\to k(y)$는 두
$S$-사상
$g_i:\operatorname{Spec}(k(y))\to\operatorname{Spec}(k(x))$를 정의하며,
그림
\[
  \xymatrix{
    \operatorname{Spec}(k(y))\ar[r]^{g_i}\ar[d] &
    \operatorname{Spec}(k(x))\ar[d]\\
    Y\ar[r]_{f_i} & X
  }
\]
은 가환이다. 따라서 그림
\[
  \xymatrix{
    \operatorname{Spec}(k(y))\ar[rr]^{(g_1,g_2)_S}\ar[d] & &
    \operatorname{Spec}(k(x))\times_S\operatorname{Spec}(k(x))\ar[d]\\
    Y\ar[rr]_{(f_1,f_2)_S} & & X\times_S X
  }
\]
도 가환이다. 그런데 $g_1=g_2$이므로, $(g_1,g_2)_S$에 의한
$\operatorname{Spec}(k(y))$의 유일한 점의 상은
$\operatorname{Spec}(k(x))\times_S\operatorname{Spec}(k(x))$의 대각선에
속한다. 따라서 결론은 (\hyperref[I.5.3.15-ko]{5.3.15})에서 따른다.

\subsection{분리 사상과 분리 준스킴}
\label{subsection:I.5.4-ko}

\begin{definition}[5.4.1]
\label{I.5.4.1-ko}
준스킴의 사상 $f:X\to Y$에서 대각 사상
$X\to X\times_Y X$가 \emph{닫힌 몰입 사상}이면 $f$를
\emph{분리} 사상이라고 한다. 이때 $X$를 \emph{$Y$ 위에서 분리된
준스킴} 또는 \emph{$Y$-스킴}이라고도 한다. 준스킴 $X$가
$\operatorname{Spec}(\mathbf{Z})$ 위에서 분리되어 있으면 $X$를 분리
준스킴이라 하고, 이때 $X$를 \emph{스킴}이라고도 한다
(\agref{I.5.5.7-ko}{5.5.7} 참조).
\end{definition}

(\hyperref[I.5.3.9-ko]{5.3.9})에 따라 $X$가 $Y$ 위에서 분리되어 있기
위한 필요충분조건은 $\Delta_X(X)$가 $X\times_Y X$의 바탕공간의
\emph{닫힌 부분공간}인 것이다.

\begin{proposition}[5.4.2]
\label{I.5.4.2-ko}
$S\to T$를 분리 사상이라 하자. $X$, $Y$가 두 $S$-준스킴이면 표준
몰입 $X\times_S Y\to X\times_T Y$
(\hyperref[I.5.3.10-ko]{5.3.10})는 닫힌 몰입 사상이다.
\end{proposition}

실제로 그림 (\hyperref[I.5.3.5.1-ko]{5.3.5.1})을 보면 $(p,q)_T$는
$\Delta_{S|T}$에서, 사상
$f\times_T g:X\times_T Y\to S\times_T S$에 의한 밑준스킴
$S\times_T S$의 확장을 통해 얻어진 것으로 볼 수 있다. 따라서 명제는
(\hyperref[I.4.3.2-ko]{4.3.2})에서 따른다.

\begin{corollary}[5.4.3]
\label{I.5.4.3-ko}
$Y$를 $S$-스킴, $f:X\to Y$를 $S$-사상이라 하자. 그러면 $f$의
그래프 사상 $\Gamma_f:X\to X\times_S Y$
(\hyperref[I.5.3.11-ko]{5.3.11})는 닫힌 몰입 사상이다.
\end{corollary}

이는 (\hyperref[I.5.4.2-ko]{5.4.2})에서 $S$를 $Y$로, $T$를 $S$로
바꾼 특별한 경우이다.

\begin{corollary}[5.4.4]
\label{I.5.4.4-ko}
$f:X\to Y$, $g:Y\to Z$를 두 사상이라 하고 $g$가 분리 사상이라고
하자. $g\circ f$가 닫힌 몰입 사상이면 $f$도 닫힌 몰입 사상이다.
\end{corollary}

(\hyperref[I.5.4.3-ko]{5.4.3})에서 출발하는 증명은
(\hyperref[I.5.3.11-ko]{5.3.11})에서 출발하는
(\hyperref[I.5.3.13-ko]{5.3.13})의 증명과 같다.

\oldpage[I]{136}
\begin{corollary}[5.4.5]
\label{I.5.4.5-ko}
$Z$를 $S$-스킴, $j:X\to Y$, $g:X\to Z$를 두 $S$-사상이라 하자.
$j$가 닫힌 몰입 사상이면
$(j,g)_S:X\to Y\times_S Z$도 닫힌 몰입 사상이다.
\end{corollary}

(\hyperref[I.5.4.4-ko]{5.4.4})에서 출발하는 증명은
(\hyperref[I.5.3.13-ko]{5.3.13})에서 출발하는
(\hyperref[I.5.3.14-ko]{5.3.14})의 증명과 같다.

\begin{corollary}[5.4.6]
\label{I.5.4.6-ko}
$X$가 $S$-스킴이면 $X$의 모든 $S$-단면
(\hyperref[I.2.5.5-ko]{2.5.5})은 닫힌 몰입 사상이다.
\end{corollary}

$\varphi:X\to S$가 구조 사상이고 $\psi:S\to X$가 $X$의
$S$-단면이면, $\varphi\circ\psi=1_S$에
(\hyperref[I.5.4.5-ko]{5.4.5})를 적용하면 충분하다.

\begin{corollary}[5.4.7]
\label{I.5.4.7-ko}
$S$를 정역 준스킴, $s$를 그 일반점, $X$를 $S$-스킴이라 하자.
$X$의 두 $S$-단면 $f$, $g$가 $f(s)=g(s)$를 만족하면 $f=g$이다.
\end{corollary}

실제로 $x=f(s)=g(s)$라 하면, $f$와 $g$에 대응하는 준동형
$k(x)\to k(s)$는 반드시 서로 같다. $h=(f,g)_S$라 하면
(\hyperref[I.5.3.17-ko]{5.3.17})에 따라 $h(s)$는 대각선
$Z=\Delta_X(X)$에 속한다. 그런데 $S=\overline{\{s\}}$이고 가정에
따라 $Z$가 닫혀 있으므로 $h(S)\subset Z$이다. 따라서
(\hyperref[I.5.2.2-ko]{5.2.2})에 따라 $h$는
$S\to Z\to X\times_S X$로 인수분해되고, 대각선의 정의에서
$f=g$가 따른다.

\begin{remark}[5.4.8]
\label{I.5.4.8-ko}
역으로 (\hyperref[I.5.4.3-ko]{5.4.3})의 결론이 $f=1_Y$일 때
성립한다고 가정하면 $Y$가 $S$ 위에서 분리되어 있음을 얻는다.
마찬가지로 (\hyperref[I.5.4.5-ko]{5.4.5})의 결론이 두 사상
\[
  Y\xrightarrow{\Delta_Y}Y\times_Z Y\xrightarrow{p_1}Y
\]
에 적용된다고 가정하면 $\Delta_Y$가 닫힌 몰입 사상이고, 따라서
$Y$가 $Z$ 위에서 분리되어 있음을 얻는다. 마지막으로
(\hyperref[I.5.4.6-ko]{5.4.6})의 결론이 $Y$-준스킴
$Y\times_S Y\to Y$의 $Y$-단면 $\Delta_Y$에 대하여 성립하면,
$Y$는 $S$ 위에서 분리되어 있다.
\end{remark}

\subsection{분리 판정법}
\label{subsection:I.5.5-ko}

\begin{proposition}[5.5.1]
\label{I.5.5.1-ko}
\begin{enumerate}
  \item[\rm(i)] 준스킴의 모든 단사사상(특히 모든 몰입 사상)은 분리
    사상이다.
  \item[\rm(ii)] 두 분리 사상의 합성은 분리 사상이다.
  \item[\rm(iii)] $f:X\to X'$, $g:Y\to Y'$가 두 분리 $S$-사상이면
    $f\times_S g$도 분리 사상이다.
  \item[\rm(iv)] $f:X\to Y$가 분리 $S$-사상이면, 밑준스킴의 모든 확장
    $S'\to S$에 대하여 $S'$-사상 $f_{(S')}$도 분리 사상이다.
  \item[\rm(v)] 두 사상의 합성 $g\circ f$가 분리 사상이면 $f$도 분리
    사상이다.
  \item[\rm(vi)] 사상 $f$가 분리 사상이기 위한 필요충분조건은
    $f_{\mathrm{red}}$ (\hyperref[I.5.1.5-ko]{5.1.5})가 분리 사상인
    것이다.
\end{enumerate}
\end{proposition}

(i)는 (\hyperref[I.5.3.8-ko]{5.3.8})에서 곧바로 따른다. 두 사상
$f:X\to Y$, $g:Y\to Z$에 대하여 그림
\[
  \xymatrix{
    X\ar[rr]^{\Delta_{X|Z}}\ar[dr]_{\Delta_{X|Y}} & &
    X\times_Z X\\
    & X\times_Y X\ar[ur]_j
  }
  \tag{5.5.1.1}\label{I.5.5.1.1-ko}
\]
은 가환이다. 여기서 $j$는 표준 몰입 사상
(\hyperref[I.5.3.10-ko]{5.3.10})이며, 가환성은 곧바로 확인된다.
$f$와 $g$가 분리 사상이면 정의에 따라 $\Delta_{X|Y}$는 닫힌 몰입
사상이고, (\hyperref[I.5.4.2-ko]{5.4.2})에 따라 $j$도 닫힌 몰입
사상이다. 따라서 (\hyperref[I.4.2.5-ko]{4.2.5})에 따라
$\Delta_{X|Z}$도 닫힌 몰입 사상이며, 이로써

\oldpage[I]{137}
(ii)가 증명된다. (i)와 (ii)를 고려하면 (iii)와 (iv)는 서로
동치이고 (\hyperref[I.3.5.1-ko]{3.5.1}), 따라서 (iv)만 증명하면
충분하다. 그런데 (\hyperref[I.3.3.11-ko]{3.3.11})과
(\hyperref[I.3.3.9.1-ko]{3.3.9.1})에 따라
$X_{(S')}\times_{Y_{(S')}}X_{(S')}$은
$(X\times_Y X)\times_Y Y_{(S')}$과 표준적으로 동일시된다. 이때 대각
사상 $\Delta_{X_{(S')}}$은
$\Delta_X\times_Y 1_{Y_{(S')}}$과 동일시됨을 곧바로 확인할 수 있다.
따라서 명제는 (\hyperref[I.4.3.1-ko]{4.3.1})에서 따른다.

(v)를 증명하기 위하여 (\hyperref[I.5.3.13-ko]{5.3.13})에서와 같이
$f$의 인수분해
$X\xrightarrow{\Gamma_f}X\times_Z Y\xrightarrow{p_2}Y$를 생각하자.
여기서 $p_2=(g\circ f)\times_Z 1_Y$이다. $g\circ f$가 분리
사상이라는 가정과 (iii), (i)에 따라 $p_2$는 분리 사상이다.
$\Gamma_f$는 몰입 사상이므로 (i)에 따라 분리 사상이고, 따라서
(ii)에 따라 $f$도 분리 사상이다. 마지막으로 (vi)를 증명하기 위해,
준스킴 $X_{\mathrm{red}}\times_{Y_{\mathrm{red}}}X_{\mathrm{red}}$와
$X_{\mathrm{red}}\times_Y X_{\mathrm{red}}$가 표준적으로 동일시됨을
(\hyperref[I.5.1.7-ko]{5.1.7}) 상기하자. $j$를 포함 사상
$X_{\mathrm{red}}\to X$라 하면 그림
\[
  \xymatrix{
    X_{\mathrm{red}}\ar[r]^{\Delta_{X_{\mathrm{red}}}}\ar[d]_j &
    X_{\mathrm{red}}\times_Y X_{\mathrm{red}}\ar[d]^{j\times_Y j}\\
    X\ar[r]_{\Delta_X} & X\times_Y X
  }
\]
은 가환이다 (\hyperref[I.5.3.15-ko]{5.3.15}). 수직 화살표들이
바탕공간의 위상동형이라는 사실
(\hyperref[I.4.3.1-ko]{4.3.1})에서 명제가 따른다.

\begin{corollary}[5.5.2]
\label{I.5.5.2-ko}
$f:X\to Y$가 분리 사상이면, $X$의 모든 부분준스킴에 대한 $f$의
제한도 분리 사상이다.
\end{corollary}

이는 (\hyperref[I.5.5.1-ko]{5.5.1, (i)와 (ii)})에서 따른다.

\begin{corollary}[5.5.3]
\label{I.5.5.3-ko}
$X$, $Y$가 두 $S$-준스킴이고 $Y$가 $S$ 위에서 분리되어 있으면,
$X\times_S Y$는 $X$ 위에서 분리되어 있다.
\end{corollary}

이는 (\hyperref[I.5.5.1-ko]{5.5.1, (iv)})의 특별한 경우이다.

\begin{proposition}[5.5.4]
\label{I.5.5.4-ko}
$X$를 준스킴이라 하고, 그 바탕공간이 유한개의 닫힌 부분집합
$X_k$ $(1\leq k\leq n)$의 합집합이라고 하자. 각 $k$에 대하여
$X_k$를 바탕공간으로 갖는 $X$의 축소 부분준스킴
(\hyperref[I.5.2.1-ko]{5.2.1})을 생각하고, 이 부분준스킴도 $X_k$로
표시한다. $f:X\to Y$를 사상이라 하고, 각 $k$에 대하여
$f(X_k)\subset Y_k$를 만족하는 $Y$의 닫힌 부분집합 $Y_k$를
택하자. $Y_k$를 바탕공간으로 갖는 $Y$의 축소 부분준스킴도
$Y_k$로 표시하면, $X_k$에 대한 $f$의 제한 $X_k\to Y$는
$X_k\xrightarrow{f_k}Y_k\to Y$로 인수분해된다
(\hyperref[I.5.2.2-ko]{5.2.2}). 이때 $f$가 분리 사상이기 위한
필요충분조건은 모든 $f_k$가 분리 사상인 것이다.
\end{proposition}

필요성은 (\hyperref[I.5.5.1-ko]{5.5.1, (i), (ii), (v)})에서 따른다.
역으로 명제의 조건이 성립한다고 하자. 그러면 $X_k$에 대한 $f$의
각 제한 $X_k\to Y$는 분리 사상이다
(\hyperref[I.5.5.1-ko]{5.5.1, (i)와 (ii)}). $p_1$, $p_2$를
$X\times_Y X$의 사영이라 하면, 부분공간 $\Delta_{X_k}(X_k)$는
$X\times_Y X$의 바탕공간의 부분공간
$\Delta_X(X)\cap p_1^{-1}(X_k)$와 동일시된다
(\hyperref[I.5.3.16-ko]{5.3.16}). 이 부분공간들은
$X\times_Y X$에서 닫혀 있으므로, 그 합집합 $\Delta_X(X)$도 닫혀
있다.

특히 $X_k$들이 $X$의 \emph{기약 성분}이라고 하자. 그러면
$Y_k$들도 $Y$의 기약 성분이라고 가정할 수 있다
(\hyperref[0.2.1.5-ko]{0, 2.1.5}). 따라서 이 경우 명제
(\hyperref[I.5.5.4-ko]{5.5.4})는 분리의 개념을 \emph{정역 준스킴}
(\hyperref[I.2.1.8-ko]{2.1.8})의 경우로 환원한다.

\oldpage[I]{138}

\begin{proposition}[5.5.5]
\label{I.5.5.5-ko}
$(Y_\lambda)$를 준스킴 $Y$의 열린 덮개라 하자. 사상 $f:X\to Y$가
분리 사상이기 위한 필요충분조건은 각 제한
$f^{-1}(Y_\lambda)\to Y_\lambda$가 분리 사상인 것이다.
\end{proposition}

$X_\lambda=f^{-1}(Y_\lambda)$라 놓자. 곱
$X_\lambda\times_Y X_\lambda$와
$X_\lambda\times_{Y_\lambda}X_\lambda$가 동일하다는 사실
(\hyperref[I.3.2.5-ko]{3.2.5})과
(\hyperref[I.4.2.4-ko]{4.2.4, b)})를 고려하면, 모든 것은
$X_\lambda\times_Y X_\lambda$들이 $X\times_Y X$의 덮개를 이룸을
증명하는 것으로 귀착된다. 그런데
$Y_{\lambda\mu}=Y_\lambda\cap Y_\mu$,
$X_{\lambda\mu}=X_\lambda\cap X_\mu=f^{-1}(Y_{\lambda\mu})$라 놓으면,
$X_\lambda\times_Y X_\mu$는 곱
$X_{\lambda\mu}\times_{Y_{\lambda\mu}}X_{\lambda\mu}$와 동일시되고
(\hyperref[I.3.2.6.4-ko]{3.2.6.4}), 따라서
$X_{\lambda\mu}\times_Y X_{\lambda\mu}$와도 동일시된다
(\hyperref[I.3.2.5-ko]{3.2.5}). 결국 이것은
$X_\lambda\times_Y X_\lambda$의 열린집합과 동일시되므로, 우리의
주장이 증명된다 (\hyperref[I.3.2.7-ko]{3.2.7}).

명제 (\hyperref[I.5.5.5-ko]{5.5.5})에 따라 $Y$를 아핀 열린집합들로
덮으면, 분리 사상의 연구를 아핀 스킴에 값을 갖는 분리 사상의
연구로 환원할 수 있다.

\begin{proposition}[5.5.6]
\label{I.5.5.6-ko}
$Y$를 아핀 스킴, $X$를 준스킴, $(U_\alpha)$를 $X$의 아핀 열린
덮개라 하자. 사상 $f:X\to Y$가 분리 사상이기 위한 필요충분조건은
모든 지표쌍 $(\alpha,\beta)$에 대하여 $U_\alpha\cap U_\beta$가 아핀
열린집합이고, 환 $\Gamma(U_\alpha\cap U_\beta,\mathscr{O}_X)$가 두 환
$\Gamma(U_\alpha,\mathscr{O}_X)$와
$\Gamma(U_\beta,\mathscr{O}_X)$의 표준 상들의 합집합으로 생성되는
것이다.
\end{proposition}

$U_\alpha\times_Y U_\beta$들은 $X\times_Y X$의 열린 덮개를 이룬다
(\hyperref[I.3.2.7-ko]{3.2.7}). $p$, $q$를 $X\times_Y X$의 사영이라
하면
\[
  \begin{aligned}
    \Delta_X^{-1}(U_\alpha\times_Y U_\beta)
      &=\Delta_X^{-1}\bigl(p^{-1}(U_\alpha)\cap q^{-1}(U_\beta)\bigr)\\
      &=\Delta_X^{-1}\bigl(p^{-1}(U_\alpha)\bigr)
        \cap\Delta_X^{-1}\bigl(q^{-1}(U_\beta)\bigr)
        =U_\alpha\cap U_\beta.
  \end{aligned}
\]
따라서 모든 것은 $U_\alpha\cap U_\beta$에 대한 $\Delta_X$의 제한이
$U_\alpha\times_Y U_\beta$ 안의 닫힌 몰입 사상임을 표현하는 것으로
귀착된다. 그런데 이 제한은 정의에 따라
$(j_\alpha,j_\beta)_Y$와 다르지 않다. 여기서 $j_\alpha$ (각각
$j_\beta$)는 $U_\alpha\cap U_\beta$에서 $U_\alpha$ (각각
$U_\beta$)로 가는 포함 사상이다.
$U_\alpha\times_Y U_\beta$는 그 환이
$\Gamma(U_\alpha,\mathscr{O}_X)
 \otimes_{\Gamma(Y,\mathscr{O}_Y)}
 \Gamma(U_\beta,\mathscr{O}_X)$와 표준적으로 동형인 아핀 스킴이므로
(\hyperref[I.3.2.2-ko]{3.2.2}), $U_\alpha\cap U_\beta$는 아핀
스킴이어야 하고, 환 $A(U_\alpha\times_Y U_\beta)$에서
$\Gamma(U_\alpha\cap U_\beta,\mathscr{O}_X)$로 가는 사상
$h_\alpha\otimes h_\beta\to h_\alpha h_\beta$는 전사이어야 한다
(\hyperref[I.4.2.3-ko]{4.2.3}). 이로써 증명이 끝난다.

\begin{corollary}[5.5.7]
\label{I.5.5.7-ko}
아핀 스킴은 분리되어 있다(따라서 스킴이며, 이것이
(\hyperref[I.5.4.1-ko]{5.4.1})의 용어법을 정당화한다).
\end{corollary}

\begin{corollary}[5.5.8]
\label{I.5.5.8-ko}
$Y$를 아핀 스킴이라 하자. 사상 $f:X\to Y$가 분리 사상이기 위한
필요충분조건은 $X$가 분리되어 있는 것(다시 말해 $X$가 스킴인
것)이다.
\end{corollary}

실제로 (\hyperref[I.5.5.6-ko]{5.5.6})의 판정법은 $f$에 의존하지
않는다.

\begin{corollary}[5.5.9]
\label{I.5.5.9-ko}
사상 $f:X\to Y$가 분리 사상이기 위한 필요조건은 $Y$가 분리
준스킴을 유도하는 모든 열린집합 $U$에 대하여 유도 준스킴
$f^{-1}(U)$가 분리되어 있는 것이다. 이것이 모든 아핀 열린집합
$U\subset Y$에 대하여 성립하면 충분하다.
\end{corollary}

조건의 필요성은 (\hyperref[I.5.5.5-ko]{5.5.5})와
(\hyperref[I.5.5.1-ko]{5.5.1, (ii)})에서 따른다. 충분성은 $Y$에 아핀
열린 덮개가 존재한다는 사실을 고려하면
(\hyperref[I.5.5.5-ko]{5.5.5})와
(\hyperref[I.5.5.8-ko]{5.5.8})에서 따른다.

특히 $X$와 $Y$가 아핀 스킴이면, \emph{모든} 사상 $X\to Y$는 분리
사상이다.

\begin{proposition}[5.5.10]
\label{I.5.5.10-ko}
$Y$를 스킴, $f:X\to Y$를 사상이라 하자. $X$의 모든 아핀 열린집합
$U$와 $Y$의 모든 아핀 열린집합 $V$에 대하여
$U\cap f^{-1}(V)$는 아핀이다.
\end{proposition}

$p_1$, $p_2$를 $X\times_{\mathbf{Z}}Y$의 사영이라 하자. 부분공간
$U\cap f^{-1}(V)$는
$\Gamma_f(X)\cap p_1^{-1}(U)\cap p_2^{-1}(V)$의 $p_1$에 의한
상이다. 그런데 $p_1^{-1}(U)\cap p_2^{-1}(V)$는

\oldpage[I]{139}
준스킴 $U\times_{\mathbf{Z}}V$의 바탕공간과 동일시되므로
(\hyperref[I.3.2.7-ko]{3.2.7}) 아핀 스킴이다
(\hyperref[I.3.2.2-ko]{3.2.2}). 또한 $\Gamma_f(X)$는
$X\times_{\mathbf{Z}}Y$에서 닫혀 있으므로
(\hyperref[I.5.4.3-ko]{5.4.3}),
$\Gamma_f(X)\cap p_1^{-1}(U)\cap p_2^{-1}(V)$는
$U\times_{\mathbf{Z}}V$에서 닫혀 있다. 따라서 $\Gamma_f$에 대응하는
$X\times_{\mathbf{Z}}Y$의 부분준스킴
(\hyperref[I.4.2.1-ko]{4.2.1})이 그 바탕공간의 열린 부분집합
$\Gamma_f(X)\cap p_1^{-1}(U)\cap p_2^{-1}(V)$ 위에 유도하는 준스킴은
아핀 스킴의 닫힌 부분준스킴이고, 따라서 아핀 스킴이다
(\hyperref[I.4.2.3-ko]{4.2.3}). 이제 $\Gamma_f$가 몰입 사상이라는
사실에서 명제가 따른다.

\begin{examples}[5.5.11]
\label{I.5.5.11-ko}
예 (\hyperref[I.2.3.2-ko]{2.3.2})의 준스킴(<<체 $K$ 위의 사영
직선>>)은 \emph{분리되어 있다}. 실제로 $X$의 아핀 열린 덮개
$(X_1,X_2)$에 대하여 $X_1\cap X_2=U_{12}$는 아핀이고,
$\Gamma(U_{12},\mathscr{O}_X)$는 $f\in K[s]$인
$f(s)/s^m$ 꼴의 유리함수들로 이루어진 환으로서 $K[s]$와 $1/s$로
생성된다. 따라서 (\hyperref[I.5.5.6-ko]{5.5.6})의 조건들이
성립한다.

예 (\hyperref[I.2.3.2-ko]{2.3.2})에서와 같은 $X_1$, $X_2$,
$U_{12}$, $U_{21}$을 택하되, 이번에는 $u_{12}$로서 $f(s)$를
$f(t)$에 대응시키는 동형을 택하자. 그러면 붙이기로 \emph{분리되지
않은 정역 준스킴} $X$를 얻는다. 실제로
(\hyperref[I.5.5.6-ko]{5.5.6})의 첫째 조건은 성립하지만 둘째 조건은
성립하지 않는다. 여기서
$\Gamma(X,\mathscr{O}_X)\to\Gamma(X_1,\mathscr{O}_X)=K[s]$가
동형임은 곧바로 알 수 있다. 그 역동형은 전사 사상
$f:X\to\operatorname{Spec}(K[s])$를 정의한다. 또한
$\mathfrak{j}_y\neq(s)$인 모든
$y\in\operatorname{Spec}(K[s])$에 대하여 $f^{-1}(y)$는 한 점으로만
이루어지지만, $\mathfrak{j}_y=(s)$이면 $f^{-1}(y)$는 서로 다른
\emph{두} 점으로 이루어진다($X$를 <<점 $0$을 이중화한 $K$ 위의
아핀 직선>>이라고 한다).

(\hyperref[I.5.5.6-ko]{5.5.6})의 둘째 조건은 성립하지만
\emph{첫째 조건은 성립하지 않는} 예도 줄 수 있다. 먼저 체 $K$ 위의 두 부정원 다항식환
$A=K[s,t]$의 소 스펙트럼 $Y$에서 $D(s)$와 $D(t)$의 합집합인
열린집합 $U$는 \emph{아핀 열린집합이 아니다}. 실제로 $z$가 $U$ 위의
$\mathscr{O}_Y$의 단면이면, $s^m z$와 $t^n z$가 각각 $s$, $t$의
다항식의 $U$에 대한 제한이 되는 두 정수 $m\geq0$, $n\geq0$이
존재한다 (\hyperref[I.1.4.1-ko]{1.4.1}). 이는 단면 $z$가 $Y$ 전체
위의 단면으로 연장되어 $s$, $t$의 다항식과 동일시되는 경우에만
가능함이 분명하다. 만일 $U$가 아핀 열린집합이면 포함 사상
$U\to Y$는 동형이어야 하는데 (\hyperref[I.1.7.3-ko]{1.7.3}), 이는
$U\neq Y$이므로 모순이다.

이제 두 아핀 스킴 $Y_1$, $Y_2$를 각각 환
$A_1=K[s_1,t_1]$, $A_2=K[s_2,t_2]$의 소 스펙트럼이라 하자.
$U_{12}=D(s_1)\cup D(t_1)$,
$U_{21}=D(s_2)\cup D(t_2)$로 놓고, $u_{12}$로서
$f(s_1,t_1)$을 $f(s_2,t_2)$에 대응시키는 환의 동형에 대응하는
동형 $Y_2\to Y_1$의 $U_{21}$에 대한 제한을 택하자. 이렇게 하여
(\hyperref[I.5.5.6-ko]{5.5.6})의 첫째 조건은 성립하지 않고 둘째
조건은 성립하는 예를 얻는다(이렇게 얻은 정역 준스킴을 <<점 $0$을
이중화한 $K$ 위의 아핀 평면>>이라고 한다).
\end{examples}

\begin{remark}[5.5.12]
\label{I.5.5.12-ko}
\normalfont 준스킴의 사상에 대한 성질 \textbf{P}가 주어졌다고 하고,
다음 명제들을 생각하자.
\begin{enumerate}
  \item[\rm(i)] \emph{모든 닫힌 몰입 사상은 성질 \textbf{P}를 갖는다.}
  \item[\rm(ii)] \emph{성질 \textbf{P}를 갖는 두 사상의 합성은 성질
    \textbf{P}를 갖는다.}
  \item[\rm(iii)] \emph{$f:X\to X'$, $g:Y\to Y'$가 성질 \textbf{P}를
    갖는 두 $S$-사상이면, $f\times_S g$는 성질 \textbf{P}를 갖는다.}

\oldpage[I]{140}
  \item[\rm(iv)] \emph{$f:X\to Y$가 성질 \textbf{P}를 갖는
    $S$-사상이면, 밑준스킴의 확장 $S'\to S$로부터 얻는 모든
    $S'$-사상 $f_{(S')}$는 성질 \textbf{P}를 갖는다.}
  \item[\rm(v)] \emph{두 사상 $f:X\to Y$, $g:Y\to Z$의 합성
    $g\circ f$가 성질 \textbf{P}를 갖고 $g$가 분리 사상이면,
    $f$는 성질 \textbf{P}를 갖는다.}
  \item[\rm(vi)] \emph{사상 $f:X\to Y$가 성질 \textbf{P}를 가지면
    $f_{\mathrm{red}}$ (\hyperref[I.5.1.5-ko]{5.1.5})도 성질
    \textbf{P}를 갖는다.}
\end{enumerate}

이때 (i)와 (ii)가 성립한다고 가정하면 (iii)와 (iv)는
\emph{동치}이고, (v)와 (vi)는 (i), (ii), (iii)의
\emph{귀결}이다.

첫째 주장은 이미 증명하였다 (\hyperref[I.3.5.1-ko]{3.5.1}).
(\hyperref[I.5.3.13-ko]{5.3.13})의 $f$의 인수분해
$X\xrightarrow{\Gamma_f}X\times_Z Y\xrightarrow{p_2}Y$를 생각하자.
관계식 $p_2=(g\circ f)\times_Z 1_Y$에 따라 $g\circ f$가 성질
\textbf{P}를 가지면 (iii)에 의해 $p_2$도 성질 \textbf{P}를 갖는다.
$g$가 분리 사상이면 $\Gamma_f$는 닫힌 몰입 사상이고
(\hyperref[I.5.4.3-ko]{5.4.3}), 따라서 (i)에 의해 성질 \textbf{P}를
갖는다. 마지막으로 (ii)에 의해 $f$도 성질 \textbf{P}를 갖는다.

끝으로 가환 그림
\[
  \xymatrix{
    X_{\mathrm{red}}\ar[r]^{f_{\mathrm{red}}}\ar[d] &
    Y_{\mathrm{red}}\ar[d]\\
    X\ar[r]_f &
    Y
  }
\]
을 생각하자. 수직 화살표들은 닫힌 몰입 사상이고
(\hyperref[I.5.1.5-ko]{5.1.5}), 따라서 (i)에 의해 성질
\textbf{P}를 갖는다. $f$가 성질 \textbf{P}를 갖는다는 가정과
(ii)에 따라 합성
$X_{\mathrm{red}}\xrightarrow{f_{\mathrm{red}}}Y_{\mathrm{red}}\to Y$가
성질 \textbf{P}를 갖는다. 마지막으로 닫힌 몰입 사상은 분리
사상이므로 (\hyperref[I.5.5.1-ko]{5.5.1, (i)}), (v)에 의해
$f_{\mathrm{red}}$도 성질 \textbf{P}를 갖는다.

다음 명제들을 생각하면
\begin{enumerate}
  \item[\rm(i')] \emph{모든 몰입 사상은 성질 \textbf{P}를 갖는다.}
  \item[\rm(v')] \emph{$g\circ f$가 성질 \textbf{P}를 가지면 $f$도
    성질 \textbf{P}를 갖는다.}
\end{enumerate}
위의 논증에서 (v')은 (i'), (ii), (iii)의 귀결임을 알 수 있다.
\end{remark}

\begin{env}[5.5.13]
\label{I.5.5.13-ko}
(v)와 (vi)는 (i), (iii), 그리고 다음 명제의 귀결이기도 하다.
\begin{enumerate}
  \item[\rm(ii')] \emph{$j:X\to Y$가 닫힌 몰입 사상이고 $g:Y\to Z$가
    성질 \textbf{P}를 갖는 사상이면, $g\circ j$는 성질 \textbf{P}를
    갖는다.}
\end{enumerate}
마찬가지로 (v')은 (i'), (iii), 그리고 다음 명제의 귀결이다.
\begin{enumerate}
  \item[\rm(ii'')] \emph{$j:X\to Y$가 몰입 사상이고 $g:Y\to Z$가
    성질 \textbf{P}를 갖는 사상이면, $g\circ j$는 성질 \textbf{P}를
    갖는다.}
\end{enumerate}
실제로 이는 (\hyperref[I.5.5.12-ko]{5.5.12})의 논증에서 곧바로
따른다.
\end{env}
% ===== END EXACT INLINE INPUT: c1s5.tex =====


% ===== BEGIN EXACT INLINE INPUT: c1s6.tex =====
\section{유한성 조건}
\label{section:I.6-ko}

\subsection{뇌터 준스킴과 국소 뇌터 준스킴}
\label{subsection:I.6.1-ko}

\begin{definition}[6.1.1]
\label{I.6.1.1-ko}
준스킴 $X$가 환이 뇌터인 아핀 열린집합 $V_\alpha$들의 유한 합집합
(각각 합집합)이면 $X$를 \emph{뇌터}(각각 \emph{국소 뇌터})라 한다.
여기서 각 환은 $V_\alpha$ 위에 유도된 스킴의 환을 뜻한다.
\end{definition}

$X$가 국소 뇌터이면 구조층 $\mathscr{O}_X$가
\emph{연접인 환의 층}이라는 것은 문제가 국소적임과
(\hyperref[I.1.5.2-ko]{1.5.2})에서 곧바로 따른다.
$\mathscr{O}_X$-가군 $\mathscr{F}$가 \emph{연접}이면,
$\mathscr{F}$의 \emph{모든 준연접 $\mathscr{O}_X$-부분가군}
\oldpage[I]{141}
\emph{(각각 모든 준연접 $\mathscr{O}_X$-몫가군)}은
\emph{연접이다}.
실제로 이 문제도 국소적이며, 뇌터 환 위의 유한 생성 가군의 부분가군
(각각 몫가군)이 유한 생성이라는 사실과
(\hyperref[I.1.5.1-ko]{1.5.1}),
(\hyperref[I.1.4.1-ko]{1.4.1}),
(\hyperref[I.1.3.10-ko]{1.3.10})을 적용하면 충분하다. 특히
$\mathscr{O}_X$의 \emph{모든 준연접 아이디얼층은} \emph{연접이다}.

준스킴 $X$가 열린집합 $W_\lambda$들의 유한 합집합(각각 합집합)이고,
각 $W_\lambda$ 위에 유도된 준스킴이 뇌터(각각 국소 뇌터)이면 $X$가
뇌터(각각 국소 뇌터)임은 자명하다.

\begin{proposition}[6.1.2]
\label{I.6.1.2-ko}
준스킴 $X$가 뇌터이기 위한 필요충분조건은 $X$가 국소 뇌터이고 그
바탕공간이 준콤팩트인 것이다. 이때 $X$의 바탕공간은 뇌터이다.
\end{proposition}

\begin{proof}
첫째 주장은 정의와 (\hyperref[I.1.1.10-ko]{1.1.10 (ii)})에서 곧바로
따른다. 둘째 주장은 (\hyperref[I.1.1.6-ko]{1.1.6})과 뇌터
부분공간들의 유한 합집합인 모든 공간이 뇌터라는 사실
(\hyperref[0.2.2.3-ko]{0, 2.2.3})에서 따른다.
\end{proof}

\begin{proposition}[6.1.3]
\label{I.6.1.3-ko}
$X$를 환 $A$의 아핀 스킴이라 하자. 다음 조건들은 동치이다.
a) $X$는 뇌터이다. b) $X$는 국소 뇌터이다. c) $A$는 뇌터이다.
\end{proposition}

\begin{proof}
a)와 b)의 동치는 (\hyperref[I.6.1.2-ko]{6.1.2})와 모든 아핀 스킴의
바탕공간이 준콤팩트라는 사실
(\hyperref[I.1.1.10-ko]{1.1.10})에서 따른다. 또한 c)가 a)를
함의함은 자명하다. a)가 c)를 함의함을 보이자. 환이 뇌터인 아핀
열린집합 $V_i$들의 유한 덮개를 $X$에서 잡고, $V_i$ 위에 유도된
준스킴의 환을 $A_i$라 하자. $A$의 아이디얼들의 증가열
$(\mathfrak{a}_n)$을 잡으면 이에 표준적으로 전단사 대응하는
(\hyperref[I.1.3.7-ko]{1.3.7})
$\widetilde{A}\equiv\mathscr{O}_X$의 아이디얼층들의 증가열
$(\widetilde{\mathfrak{a}}_n)$이 존재한다. $(\mathfrak{a}_n)$이
어느 단계 이후 일정함을 보이려면 $(\widetilde{\mathfrak{a}}_n)$이
어느 단계 이후 일정함을
보이면 충분하다. 그런데 $\widetilde{\mathfrak{a}}_n|V_i$는 표준
포함 사상 $V_i\to X$에 의한 $\widetilde{\mathfrak{a}}_n$의
역상이므로 (\hyperref[0.5.1.4-ko]{0, 5.1.4})
$\mathscr{O}_X|V_i$의 준연접 아이디얼층이다. 따라서
$\widetilde{\mathfrak{a}}_n|V_i=\widetilde{\mathfrak{a}}_{ni}$이고,
여기서 $\mathfrak{a}_{ni}$는 $A_i$의 아이디얼이다
(\hyperref[I.1.3.7-ko]{1.3.7}). $A_i$가 뇌터이므로 모든 $i$에
대하여 $(\mathfrak{a}_{ni})$는 어느 단계 이후 일정하며, 이로부터
명제가 따른다.
\end{proof}

앞의 논증은 \emph{$X$가 뇌터 준스킴이면 $\mathscr{O}_X$의 연접
아이디얼층들의 모든 증가열은 어느 단계 이후 일정하다}라는 것도
증명한다.

\begin{proposition}[6.1.4]
\label{I.6.1.4-ko}
뇌터(각각 국소 뇌터) 준스킴의 모든 부분준스킴은 뇌터(각각 국소
뇌터)이다.
\end{proposition}

\begin{proof}
뇌터 준스킴 $X$에 대해서만 증명하면 충분하다. 또한 정의
(\hyperref[I.6.1.1-ko]{6.1.1})에 따라 $X$가 아핀 스킴인 경우로
곧바로 환원된다. $X$의 모든 부분준스킴은 어떤 열린집합 위에 유도된
준스킴의 닫힌 부분준스킴이므로
(\hyperref[I.4.1.3-ko]{4.1.3}), $Y$가 $X$의 닫힌 부분준스킴이거나
열린집합 위에 유도된 부분준스킴인 경우만 생각하면 된다. $Y$가
닫힌 경우는 자명하다. 실제로 $A$가 $X$의 환이면 $Y$는 $A$의 어떤
아이디얼 $\mathfrak{J}$에 대하여 환 $A/\mathfrak{J}$의 아핀
스킴이다(\hyperref[I.4.2.3-ko]{4.2.3}). $A$가 뇌터이므로
(\hyperref[I.6.1.3-ko]{6.1.3}) $A/\mathfrak{J}$도 뇌터이다.

이제 $Y$가 $X$의 열린집합이라고 하자. 바탕공간 $Y$는 뇌터이므로
(\hyperref[I.6.1.2-ko]{6.1.2}) 준콤팩트이고, 따라서
$f_i\in A$인 열린집합 $D(f_i)$들의 유한 합집합이다.
\oldpage[I]{142}
그러므로 $f\in A$에 대하여 $Y=D(f)$인 경우만 증명하면 된다. 이때
$Y$는 환이 $A_f$와 동형인 아핀 스킴이다
(\hyperref[I.1.3.6-ko]{1.3.6}). $A$가 뇌터이므로
(\hyperref[I.6.1.3-ko]{6.1.3}) $A_f$도 뇌터이다.
\end{proof}

\begin{env}[6.1.5]
\label{I.6.1.5-ko}
두 뇌터 $S$-준스킴의 \emph{곱}은 반드시 뇌터인 것은 아니다. 이는
두 뇌터 대수의 텐서곱이 반드시 뇌터 환인 것은 아니므로, 두
준스킴이 아핀인 경우에도 마찬가지이다
(\agref{I.6.3.8-ko}{6.3.8} 참조).
\end{env}

\begin{proposition}[6.1.6]
\label{I.6.1.6-ko}
$X$가 뇌터 준스킴이면 $\mathscr{O}_X$의 멱영근기
$\mathscr{N}_X$는 멱영이다.
\end{proposition}

실제로 $X$를 유한개의 아핀 열린집합 $U_i$로 덮을 수 있고,
$(\mathscr{N}_X|U_i)^{n_i}=0$인 정수 $n_i$가 존재함을 보이면
충분하다. $n$을 $n_i$들의 최댓값이라 하면
$\mathscr{N}_X^n=0$이다. 따라서 $A$가 뇌터 환이고
$X=\operatorname{Spec}(A)$가 아핀인 경우로 환원된다. 이때
(\hyperref[I.5.1.1-ko]{5.1.1})과
(\hyperref[I.1.3.13-ko]{1.3.13})에 따라 $A$의 멱영근기가 멱영이라는
사실(\cite{I-11}, p.~127, cor.~4)을 확인하면 충분하다.

\begin{corollary}[6.1.7]
\label{I.6.1.7-ko}
$X$를 뇌터 준스킴이라 하자. $X$가 아핀 스킴이기 위한
필요충분조건은 $X_{\mathrm{red}}$가 아핀 스킴인 것이다.
\end{corollary}

이는 (\hyperref[I.6.1.6-ko]{6.1.6})과
(\hyperref[I.5.1.10-ko]{5.1.10})에서 따른다.

\begin{lemma}[6.1.8]
\label{I.6.1.8-ko}
$X$를 위상공간, $x$를 $X$의 한 점, $U$를 기약 성분이 유한개뿐인
$x$의 열린 근방이라 하자. 그러면 $x$의 근방 $V$가 존재하여,
$V$에 포함되는 $x$의 모든 열린 근방이 연결이다.
\end{lemma}

$x$를 포함하지 않는 $U$의 기약 성분들을 $U_i$($1\leq i\leq m$)라
하자. 그 합집합의 $U$에서의 여집합은 $U$ 안에서, 따라서 $X$ 안에서
$x$의 열린 근방 $V$이다. 또한 이는 $U$와, $x$를 포함하지 않는
$X$의 기약 성분들의 합집합을 $X$에서 뺀 여집합의 교집합이다
(\hyperref[0.2.1.6-ko]{0, 2.1.6}). 이제 $V$에 포함되는 $x$의 열린
근방 $W$를 잡자. $W$의 기약 성분들은 $W$와 만나는 $U$의 기약
성분들과 $W$의 교집합들이므로
(\hyperref[0.2.1.6-ko]{0, 2.1.6}), 모두 $x$를 포함한다. 이 성분들은
연결이므로 $W$도 연결이다.

\begin{corollary}[6.1.9]
\label{I.6.1.9-ko}
국소 뇌터 위상공간은 국소 연결이다. 따라서 특히 그 연결 성분들은
열린집합이다.
\end{corollary}

\begin{proposition}[6.1.10]
\label{I.6.1.10-ko}
$X$를 국소 뇌터 위상공간이라 하자. 다음 조건들은 동치이다.
\begin{enumerate}
  \item[a)] $X$의 기약 성분들은 열린집합이다.
  \item[b)] $X$의 기약 성분들은 연결 성분들과 일치한다.
  \item[c)] $X$의 연결 성분들은 기약이다.
  \item[d)] $X$의 서로 다른 두 기약 성분은 만나지 않는다.
\end{enumerate}

또한 $X$가 준스킴이면 이 조건들은 다음 조건과도 동치이다.
\begin{enumerate}
  \item[e)] 모든 $x\in X$에 대하여
    $\operatorname{Spec}(\mathscr{O}_x)$는 기약이다. 다시 말해
    $\mathscr{O}_x$의 멱영근기는 소 아이디얼이다.
\end{enumerate}
\end{proposition}

기약공간은 연결이므로 a)는 b)를 함의한다. 실제로 a)에 따라 $X$의
기약 성분들은 동시에 열린집합이고 닫힌집합이다. b)가 c)를 함의함은
자명하다. 역으로 닫힌집합 $F$가
\oldpage[I]{143}
$X$의 연결 성분 $C$를 포함하면서 $C$와 다르면 $F$는 기약일 수
없다. 실제로 이 집합은 연결이 아니므로
$F$ 안에서 동시에 열리고 닫힌 서로소인 두 공집합 아닌 집합의
합집합이며, 이 두 집합은 $X$에서 닫혀 있다. 따라서 c)는 b)를
함의한다. 이로부터 서로 다른 두 연결 성분은 만나지 않으므로 c)가
d)를 함의함이 곧바로 따른다.

여기까지는 $X$가 국소 뇌터라는 사실을 사용하지 않았다. 이제 이
가정을 하고 d)가 a)를 함의함을 보이자. (\hyperref[0.2.1.6-ko]{0,
2.1.6})에 따라 $X$가 뇌터인 경우로 환원할 수 있고, 이때 $X$의 기약
성분은 유한개뿐이다. 이 성분들이 닫혀 있고 서로소이므로 열린집합이다.

끝으로 d)와 e)의 동치는 준스킴 $X$의 바탕공간이 국소 뇌터라는
가정 없이도 성립한다. (\hyperref[0.2.1.6-ko]{0, 2.1.6})에 따라
$X=\operatorname{Spec}(A)$가 아핀인 경우로 환원할 수 있다. $x$가
$X$의 기약 성분 하나에만 포함된다는 것은
$\mathfrak{j}_x$가 $A$의 극소 소 아이디얼 하나만을 포함한다는 뜻이고
(\hyperref[I.1.1.14-ko]{1.1.14}), 이는
$\mathfrak{j}_x\mathscr{O}_x$가 $\mathscr{O}_x$의 극소 소 아이디얼
하나만을 포함한다는 것과 동치이다. 이로부터 결론이 따른다.

\begin{corollary}[6.1.11]
\label{I.6.1.11-ko}
$X$를 국소 뇌터 공간이라 하자. $X$가 기약이기 위한 필요충분조건은
$X$가 연결이고 공집합이 아니며, $X$의 서로 다른 두 기약 성분이
만나지 않는 것이다. $X$가 준스킴이면 마지막 조건은 모든 $x\in X$에
대하여 $\operatorname{Spec}(\mathscr{O}_x)$가 기약이라는 것과
동치이다.
\end{corollary}

마지막 주장은 (\hyperref[I.6.1.10-ko]{6.1.10})에서 보였다. 따라서
첫째 주장의 조건들이 충분함만 증명하면 된다. 그런데
(\hyperref[I.6.1.10-ko]{6.1.10})에 따라 이 조건들은 $X$의 기약
성분들이 연결 성분들과 일치함을 함의하고, $X$가 연결이고 공집합이
아니므로 $X$는 기약이다.

\begin{corollary}[6.1.12]
\label{I.6.1.12-ko}
$X$를 국소 뇌터 준스킴이라 하자. $X$가 정역 준스킴이기 위한
필요충분조건은 $X$가 공집합이 아니고 연결이며 모든 $x\in X$에 대하여
$\mathscr{O}_x$가 정역인 것이다.
\end{corollary}

\begin{proposition}[6.1.13]
\label{I.6.1.13-ko}
$X$를 국소 뇌터 준스킴, $x\in X$를 한 점이라 하자. 국소환
$\mathscr{O}_x$의 멱영근기 $\mathscr{N}_x$가 소 아이디얼이면(각각
$\mathscr{O}_x$가 축소환이면, 정역이면), $x$의 기약인(각각 축소인,
정역인) 열린 근방 $U$가 존재한다.
\end{proposition}

$\mathscr{N}_x$가 소 아이디얼인 경우와 $\mathscr{N}_x=0$인 경우만
생각하면 충분하다. 셋째 가정은 이 두 가정을 함께 말한 것이기
때문이다. $\mathscr{N}_x$가 소 아이디얼이면 $x$는 $X$의 기약 성분
$Y$ 하나에만 속한다(\hyperref[I.6.1.10-ko]{6.1.10}). $x$를 포함하지
않는 $X$의 기약 성분들의 집합은 국소 유한이므로 그 합집합은 닫혀
있다. 따라서 이 합집합의 여집합 $U$는 열린집합이고 $Y$에 포함되므로
기약이다(\hyperref[0.2.1.6-ko]{0, 2.1.6}). $\mathscr{N}_x=0$이면
$x$의 어떤 근방에 속하는 모든 $y$에 대하여 $\mathscr{N}_y=0$이다.
실제로 $\mathscr{N}$은 준연접이고
(\hyperref[I.5.1.1-ko]{5.1.1}), $X$가 국소 뇌터이므로 연접이다.
따라서 결론은 (\hyperref[0.5.2.2-ko]{0, 5.2.2})에서 따른다.

\subsection{아르틴 준스킴}
\label{subsection:I.6.2-ko}

\begin{definition}[6.2.1]
\label{I.6.2.1-ko}
준스킴이 아핀이고 그 환이 아르틴 환이면 그 준스킴을
\emph{아르틴}이라고 한다.
\end{definition}

\oldpage[I]{144}
\begin{proposition}[6.2.2]
\label{I.6.2.2-ko}
준스킴 $X$에 대하여 다음 조건들은 동치이다.
\begin{enumerate}
  \item[a)] $X$는 아르틴 스킴이다.
  \item[b)] $X$는 뇌터이고 그 바탕공간은 이산이다.
  \item[c)] $X$는 뇌터이고 그 바탕공간의 모든 점은 닫힌점이다
    ($\mathrm{T}_1$ 조건).
\end{enumerate}

이 조건들이 성립하면 $X$의 바탕공간은 유한하고, $X$의 환 $A$는
$X$의 점들의 (아르틴) 국소환들의 직접곱이다.
\end{proposition}

a)가 마지막 주장을 함의함은 알려져 있다(\cite{I-13}, p.~205,
th.~3). 이때 $A$의 모든 소 아이디얼은 극대 아이디얼이고, $A$의 한
국소환 성분의 극대 아이디얼의 역상이다. 따라서 공간 $X$는 유한하고
이산이다. 그러므로 a)는 b)를 함의하고, b)가 c)를 함의함은 자명하다.
c)가 a)를 함의함을 보이기 위해 먼저 $X$가 유한함을 보이자. 실제로
$X$가 아핀인 경우로 환원할 수 있고, 모든 소 아이디얼이 극대
아이디얼인 뇌터 환은 아르틴 환임이 알려져 있다(\cite{I-13}, p.~203).
따라서 주장이 따른다. 이제 $X$의 바탕공간은 이산이고, 유한개의 점
$x_i$의 위상적 합이다. 국소환 $\mathscr{O}_{x_i}=A_i$들은 아르틴
환이며, $A\simeq\prod_i A_i$이고 $X$는 이 환 $A$의 소 스펙트럼인
아핀 스킴과 동형임이 자명하다(\hyperref[I.1.7.3-ko]{1.7.3}).

\subsection{유한형 사상}
\label{subsection:I.6.3-ko}

\begin{definition}[6.3.1]
\label{I.6.3.1-ko}
$Y$가 다음 성질을 갖는 아핀 열린집합들의 족 $(V_\alpha)$의 합집합이면
사상 $f:X\to Y$를 \emph{유한형}이라고 한다.
\begin{enumerate}
  \item[(P)] $f^{-1}(V_\alpha)$는 유한개의 아핀 열린집합
    $U_{\alpha i}$의 합집합이고, 각 환 $A(U_{\alpha i})$는
    $A(V_\alpha)$ 위의 유한 생성 대수이다.
\end{enumerate}
이때 $X$를 $Y$ 위의 유한형 준스킴 또는 유한형 $Y$-준스킴이라고도
한다.
\end{definition}

\begin{proposition}[6.3.2]
\label{I.6.3.2-ko}
$f:X\to Y$가 유한형 사상이면 $Y$의 모든 아핀 열린집합 $W$는
정의 (\hyperref[I.6.3.1-ko]{6.3.1})의 성질 (P)를 갖는다.
\end{proposition}

먼저 다음 보조정리를 증명하자.

\begin{lemma}[6.3.2.1]
\label{I.6.3.2.1-ko}
$T\subset Y$가 성질 (P)를 갖는 아핀 열린집합이면, 모든
$g\in A(T)$에 대하여 $D(g)$도 성질 (P)를 갖는다.
\end{lemma}

실제로 가정에 따라 $f^{-1}(T)$는 유한개의 아핀 열린집합 $Z_j$의
합집합이고, 각 $A(Z_j)$는 $A(T)$ 위의 유한 생성 대수이다. $f$를
$Z_j$에 제한한 사상에 대응하는 환 준동형을
$\varphi_j:A(T)\to A(Z_j)$라 하고
(\hyperref[I.2.2.4-ko]{2.2.4}), $g_j=\varphi_j(g)$라 놓자. 그러면
$f^{-1}(D(g))\cap Z_j=D(g_j)$이다
(\hyperref[I.1.2.2.2-ko]{1.2.2.2}). 한편
$A(D(g_j))=A(Z_j)_{g_j}=A(Z_j)[1/g_j]$는 $A(Z_j)$ 위의 유한 생성
대수이고, 따라서 가정에 의해 \emph{특히} $A(T)$ 위의 유한 생성
대수이며, 그러므로 $A(D(g))=A(T)[1/g]$ 위의 유한 생성 대수이기도 하다.
이로써 보조정리가 증명된다.

이제 $W$가 준콤팩트이므로(\hyperref[I.1.1.10-ko]{1.1.10}), 각
$g_i$가 어떤 환 $A(V_{\alpha(i)})$에 속하는
$D(g_i)\subset W$들로 이루어진 $W$의 유한 덮개가 존재한다. 각
$D(g_i)$도 준콤팩트이므로
$h_{ik}\in A(W)$인 집합 $D(h_{ik})$들의 유한 합집합이다. 표준 사상을
$\varphi_i:A(W)\to A(D(g_i))$라 하면
(\hyperref[I.1.2.2.2-ko]{1.2.2.2})에 따라
$D(h_{ik})=D(\varphi_i(h_{ik}))$이다.
(\hyperref[I.6.3.2.1-ko]{6.3.2.1})에 따라 각
$f^{-1}(D(h_{ik}))$에는 유한 아핀 열린 덮개 $(U_{ijk})$가 존재하고,
각 $A(U_{ijk})$는
$A(D(h_{ik}))=A(W)[1/h_{ik}]$ 위의 유한 생성 대수이다. 이로부터 명제가
따른다.

\oldpage[I]{145}
따라서 $Y$ 위의 유한형 준스킴이라는 개념은 \emph{$Y$ 위에서
국소적}이라고 할 수 있다.

\begin{proposition}[6.3.3]
\label{I.6.3.3-ko}
$X$, $Y$를 두 아핀 스킴이라 하자. $X$가 $Y$ 위에서 유한형이기 위한
필요충분조건은 $A(X)$가 $A(Y)$ 위의 유한 생성 대수인 것이다.
\end{proposition}

이 조건이 충분함은 자명하므로 필요함을 증명하자. $A=A(Y)$,
$B=A(X)$라 놓자. (\hyperref[I.6.3.2-ko]{6.3.2})에 따라 $X$에는 유한
아핀 열린 덮개 $(V_i)$가 존재하여, 각 환 $A(V_i)$는 $A$ 위의 유한 생성
대수이다. 또한 $V_i$들이 준콤팩트이므로 각각은 유한개의 열린집합
$D(g_{ij})\subset V_i$로 덮이며, 여기서 $g_{ij}\in B$이다. 표준 포함
사상 $V_i\to X$에 대응하는 준동형을
$\varphi_i:B\to A(V_i)$라 하면
\[
  B_{g_{ij}}=(A(V_i))_{\varphi_i(g_{ij})}
  =A(V_i)[1/\varphi_i(g_{ij})].
\]
따라서 $B_{g_{ij}}$는 $A$ 위의 유한 생성 대수이다. 그러므로
$V_i=D(g_i)$이고 $g_i\in B$인 경우로 환원할 수 있다. 가정에 따라
$B_{g_i}$가 $A$ 위에서, $b_i$가 유한집합 $F_i\subset B$를 움직일 때의
원소 $b_i/g_i^{n_i}$들로 생성되도록 하는 정수 $n_i\geq0$가 존재한다.
$g_i$들이 유한개뿐이므로 모든 $n_i$가 같은 정수 $n$이라고 해도 된다.
또한 $D(g_i)$들이 $X$를 덮으므로 $g_i$들이 $B$에서 생성하는
아이디얼은 $B$이다. 다시 말해
$\sum_i h_i g_i=1$인 $h_i\in B$들이 존재한다. 이제 $F_i$들, $g_i$들,
$h_i$들의 합집합인 $B$의 유한 부분집합을 $F$라 하자. $B$의 부분환
$B'=A[F]$가 $B$와 같음을 보이겠다. 가정에 따라 모든 $b\in B$와 모든
$i$에 대하여 $B_{g_i}$에서 $b$의 표준상은 $b'_i/g_i^{m_i}$ 꼴이고,
여기서 $b'_i\in B'$이다. $b'_i$에 $g_i$의 적당한 거듭제곱을 곱하여
모든 $m_i$가 같은 정수 $m$이라고 해도 된다. 따라서 분수환의 정의에
의해 $N\geq m$이고 모든 $i$에 대하여 $g_i^N b\in B'$인 정수
$N$이 존재한다($N$은 $b$에 의존한다). 한편 \emph{환 $B'$ 안에서}
$g_i^N$들은 아이디얼 $B'$를 생성한다. 실제로 $g_i$들이 그러하고
$h_i\in B'$이기 때문이다. 따라서
$\sum_i c_i g_i^N=1$인 $c_i\in B'$들이 존재하고,
$b=\sum_i c_i g_i^N b\in B'$이다. 이로써 증명이 끝난다.

\begin{proposition}[6.3.4]
\label{I.6.3.4-ko}
\begin{enumerate}
  \item[(i)] 모든 닫힌 몰입 사상은 유한형이다.
  \item[(ii)] 두 유한형 사상의 합성은 유한형이다.
  \item[(iii)] $f:X\to X'$, $g:Y\to Y'$가 두 유한형 $S$-사상이면
    $f\times_S g:X\times_S Y\to X'\times_S Y'$도 유한형이다.
  \item[(iv)] $f:X\to Y$가 유한형 $S$-사상이면, 밑준스킴의 모든 확장
    $g:S'\to S$에 대하여 $f_{(S')}:X_{(S')}\to Y_{(S')}$도 유한형이다.
  \item[(v)] 두 사상의 합성 $g\circ f$가 유한형이고 $g$가 분리
    사상이면, $f$는 유한형이다.
  \item[(vi)] 사상 $f$가 유한형이면 $f_{\mathrm{red}}$도 유한형이다.
\end{enumerate}
\end{proposition}

(\hyperref[I.5.5.12-ko]{5.5.12})에 따라 (i), (ii), (iv)만 증명하면
충분하다.

(i)을 보이기 위해서는 $X$가 $Y$의 닫힌 부분준스킴이고
$X\to Y$가 표준 포함 사상인 경우만 다루면 된다. 또한
(\hyperref[I.6.3.2-ko]{6.3.2})에 따라 $Y$가 아핀이라고 해도 된다.
이때 $X$도 아핀이고(\hyperref[I.4.2.3-ko]{4.2.3}), 그 환은
$A/\mathfrak{J}$ 꼴의 몫환과 동형이다. 여기서 $A$는 $Y$의 환이고
$\mathfrak{J}$는 $A$의 아이디얼이다. $A/\mathfrak{J}$는 $A$ 위의
유한 생성 대수이므로 결론이 따른다.

이제 (ii)를 증명하자. $f:X\to Y$, $g:Y\to Z$를 두 유한형 사상이라
하고, $U$를 $Z$의 아핀 열린집합이라 하자. $g^{-1}(U)$에는 유한 아핀
열린 덮개 $(V_i)$가 존재하여 각 $A(V_i)$는 $A(U)$ 위의 유한 생성
대수이다(\hyperref[I.6.3.2-ko]{6.3.2}). 마찬가지로

\oldpage[I]{146}
각 $f^{-1}(V_i)$에는 유한 아핀 열린 덮개 $(W_{ij})$가 존재하여
$A(W_{ij})$는 $A(V_i)$ 위의 유한 생성 대수이고, 따라서 $A(U)$ 위의
유한 생성 대수이기도 하다. 이로부터 결론이 따른다.

마지막으로 (iv)를 증명하기 위해 $S=Y$인 경우만 다루면 된다. 실제로
$f$를 $Y$-사상으로 보고 밑준스킴의 확장을 $Y_{(S')}\to Y$로 보면
$f_{(S')}$는 $f_{(Y_{(S')})}$와도 같다
(\hyperref[I.3.3.9-ko]{3.3.9}). 이제 $p$, $q$를 각각 사영
$X_{(S')}\to X$, $X_{(S')}\to S'$라 하자. $V$를 $S$의 아핀
열린집합이라 하자. $f^{-1}(V)$는 유한개의 아핀 열린집합 $W_i$의
합집합이고, 각 $A(W_i)$는 $A(V)$ 위의 유한 생성 대수이다
(\hyperref[I.6.3.2-ko]{6.3.2}). $V'$를 $g^{-1}(V)$에 포함되는 $S'$의
아핀 열린집합이라 하자. $f\circ p=g\circ q$이므로 $q^{-1}(V')$는
$p^{-1}(W_i)$들의 합집합에 포함된다. 한편 교집합
$p^{-1}(W_i)\cap q^{-1}(V')$는 곱 $W_i\times_V V'$와 동일시된다
(\hyperref[I.3.2.7-ko]{3.2.7}). 이것은 환
$A(W_i)\otimes_{A(V)}A(V')$와 동형인 환을 갖는 아핀 스킴이다
(\hyperref[I.3.2.2-ko]{3.2.2}). 이 환은 가정에 따라 $A(V')$ 위의
유한 생성 대수이므로 명제가 증명된다.

\begin{corollary}[6.3.5]
\label{I.6.3.5-ko}
$f:X\to Y$를 몰입 사상이라 하자. $Y$의 바탕공간이 국소 뇌터이거나
$X$의 바탕공간이 뇌터이면 $f$는 유한형이다.
\end{corollary}

(\hyperref[I.6.3.2-ko]{6.3.2})에 따라 언제나 $Y$가 아핀이라고 해도
된다. $Y$의 바탕공간이 국소 뇌터이면 더 나아가 뇌터라고 해도 되고,
그 부분공간인 $X$의 바탕공간도 뇌터이다. 따라서 $Y$가 아핀이고
$X$의 바탕공간이 뇌터인 경우만 다루면 된다. 이때 $X$에는
$g_i\in A(Y)$인 유한개의 아핀 열린집합 $D(g_i)\subset Y$로 된 덮개가
존재하여, 각 $X\cap D(g_i)$는 $D(g_i)$에서 닫혀 있고 따라서 아핀
스킴이다(\hyperref[I.4.2.3-ko]{4.2.3}). 그런 덮개가 존재하는 것은 $X$가
$Y$에서 국소 닫혀 있기 때문이다(\hyperref[I.4.1.3-ko]{4.1.3}).
(\hyperref[I.6.3.4-ko]{6.3.4, (i)})와
(\hyperref[I.6.3.3-ko]{6.3.3})에 따라 $A(X\cap D(g_i))$는
$A(D(g_i))$ 위의 유한 생성 대수이다. 또한
$A(D(g_i))=A(Y)_{g_i}=A(Y)[1/g_i]$는 $A(Y)$ 위의 유한 생성 대수이므로
증명이 끝난다.

\begin{corollary}[6.3.6]
\label{I.6.3.6-ko}
$f:X\to Y$, $g:Y\to Z$를 두 사상이라 하자. $g\circ f$가 유한형이고,
$X$가 뇌터이거나 $X\times_Z Y$가 국소 뇌터이면 $f$는 유한형이다.
\end{corollary}

이는 (\hyperref[I.5.5.12-ko]{5.5.12})의 증명과 몰입 사상
$\Gamma_f$에 적용한 (\hyperref[I.6.3.5-ko]{6.3.5})에서 곧바로 따른다.

\begin{proposition}[6.3.7]
\label{I.6.3.7-ko}
$f:X\to Y$를 유한형 사상이라 하자. $Y$가 뇌터이면 $X$도 뇌터이고,
$Y$가 국소 뇌터이면 $X$도 국소 뇌터이다.
\end{proposition}

$Y$가 뇌터인 경우만 증명하면 된다. 이때 $Y$는 환 $A(V_i)$가 뇌터
환인 유한개의 아핀 열린집합 $V_i$의 합집합이다.
(\hyperref[I.6.3.2-ko]{6.3.2})에 따라 각 $f^{-1}(V_i)$는 유한개의
아핀 열린집합 $W_{ij}$의 합집합이고, 각 $A(W_{ij})$는 $A(V_i)$ 위의
유한 생성 대수이므로 뇌터 환이다. 따라서 $X$는 뇌터이다.

\begin{corollary}[6.3.8]
\label{I.6.3.8-ko}
$X$를 $S$ 위의 유한형 준스킴이라 하자. $S'$가 뇌터인(각각 국소
뇌터인) 밑준스킴의 모든 확장 $S'\to S$에 대하여 $X_{(S')}$는
뇌터이다(각각 국소 뇌터이다).
\end{corollary}

이는 (\hyperref[I.6.3.7-ko]{6.3.7})에서 따른다. 실제로
(\hyperref[I.6.3.4-ko]{6.3.4, (iv)})에 따라 $X_{(S')}$는 $S'$ 위에서
유한형이다.

또한 $S$-준스킴의 곱 $X\times_S Y$에서, 인자 $X$, $Y$ 가운데
\emph{하나가}
\oldpage[I]{147}
$S$ 위에서 \emph{유한형}이고
\emph{다른 하나가 뇌터이면}(각각 \emph{국소 뇌터이면}),
$X\times_S Y$는 \emph{뇌터이다}(각각 \emph{국소 뇌터이다}).

\begin{corollary}[6.3.9]
\label{I.6.3.9-ko}
$X$를 국소 뇌터 준스킴 $S$ 위의 유한형 준스킴이라 하자. 그러면 모든
$S$-사상 $f:X\to Y$는 유한형이다.
\end{corollary}

실제로 $S$가 뇌터라고 해도 된다. $\varphi:X\to S$,
$\psi:Y\to S$가 구조 사상이면 $\varphi=\psi\circ f$이고,
(\hyperref[I.6.3.7-ko]{6.3.7})에 따라 $X$는 뇌터이다. 따라서
(\hyperref[I.6.3.6-ko]{6.3.6})에 따라 $f$는 유한형이다.

\begin{proposition}[6.3.10]
\label{I.6.3.10-ko}
$f:X\to Y$를 유한형 사상이라 하자. $f$가 전사이기 위한 필요충분조건은
모든 대수적으로 닫힌 체 $\Omega$에 대하여 $f$에 대응하는 사상
$X(\Omega)\to Y(\Omega)$가 전사인 것이다
(\hyperref[I.3.4.1-ko]{3.4.1}).
\end{proposition}

이 조건이 충분함은 각 $y\in Y$에 대하여 $k(y)$의 대수적으로 닫힌
확대체 $\Omega$와 다음 가환 그림을 고려하면 알 수 있다.
\[
  \xymatrix{
    X\ar[dd]_f\\
    & \operatorname{Spec}(\Omega)\ar[ul]\ar[dl]\\
    Y
  }
\]
(\emph{참조} (\hyperref[I.3.5.3-ko]{3.5.3})). 역으로 $f$가 전사라고
가정하고, $\Omega$가 대수적으로 닫힌 체인 사상
$g:\{\xi\}=\operatorname{Spec}(\Omega)\to Y$를 택하자. 다음 그림을
고려하면
\[
  \xymatrix{
    X\ar[d]_f &
    X_{(\Omega)}\ar[l]\ar[d]^{f_{(\Omega)}}\\
    Y &
    \operatorname{Spec}(\Omega)\ar[l]
  }
\]
$X_{(\Omega)}$에 \emph{$\Omega$-유리점}이 존재함을 보이면 충분하다
(\hyperref[I.3.3.14-ko]{3.3.14}, \hyperref[I.3.4.3-ko]{3.4.3},
\hyperref[I.3.4.4-ko]{3.4.4}). $f$가 전사이므로 $X_{(\Omega)}$는
공집합이 아니고(\hyperref[I.3.5.10-ko]{3.5.10}), $f$가 유한형이므로
(\hyperref[I.6.3.4-ko]{6.3.4, (iv)})에 따라 $f_{(\Omega)}$도 유한형이다.
따라서 $X_{(\Omega)}$에는 $A(Z)$가 $\Omega$ 위의 영환이 아닌 유한 생성
대수인 공집합이 아닌 아핀 열린집합 $Z$가 존재한다. 힐베르트
영점정리(\cite{I-21})에 따라 $\Omega$-대수 준동형
$A(Z)\to\Omega$가 존재하고, 따라서 $\operatorname{Spec}(\Omega)$ 위의
$X_{(\Omega)}$의 단면이 존재한다. 이로써 명제가 증명된다.

\subsection{대수적 준스킴}
\label{subsection:I.6.4-ko}

\begin{definition}[6.4.1]
\label{I.6.4.1-ko}
체 $K$가 주어졌을 때 $K$ 위의 유한형 준스킴 $X$를
\emph{$K$-대수적 준스킴}이라고 한다. $K$를 $X$의 \emph{기초체}라고 한다.
더 나아가 $X$가 스킴이면(또는 이와 동치로
(\hyperref[I.5.5.8-ko]{5.5.8}), $X$가 $K$-스킴이면) $X$를
\emph{$K$-대수적 스킴}이라고도 한다.
\end{definition}

모든 $K$-대수적 준스킴은 \emph{뇌터이다}
(\hyperref[I.6.3.7-ko]{6.3.7}).

\begin{proposition}[6.4.2]
\label{I.6.4.2-ko}
$X$를 $K$-대수적 준스킴이라 하자. 점 $x\in X$가 닫힌점이기 위한
필요충분조건은 $k(x)$가 $K$의 유한 차수 대수적 확대체인 것이다.
\end{proposition}

$X$가 아핀이고 그 환 $A$가 $K$ 위의 유한 생성 대수인 경우만 다루면
된다. 실제로 $A(U)$가 $K$ 위의 유한 생성 대수인 $X$의 아핀 열린집합
$U$들은 $X$를 덮는다(\hyperref[I.6.3.1-ko]{6.3.1}). 이때 $X$의 닫힌점들은
$\mathfrak{j}_x$가
\oldpage[I]{148}
$A$의 극대 아이디얼인 점들, 다시 말해 $A/\mathfrak{j}_x$가 체인 점들이다.
이 체는 반드시 $k(x)$와 같다. $A/\mathfrak{j}_x$가 $K$ 위의 유한 생성
대수이므로, $x$가 닫힌점이면 $k(x)$는 $K$ 위의 유한 생성 대수인 체이고,
따라서 반드시 $K$ 위에서 \emph{유한 차원}이다 \cite{I-21}. 역으로
$k(x)$가 $K$ 위에서 유한 차원이면 그 부분환
$A/\mathfrak{j}_x\subset k(x)$도 그러하다. $K$ 위의 유한 차원 대수인
정역은 모두 체이므로 $A/\mathfrak{j}_x=k(x)$이고, 따라서 $x$는 닫힌점이다.

\begin{corollary}[6.4.3]
\label{I.6.4.3-ko}
$K$를 대수적으로 닫힌 체, $X$를 $K$-대수적 준스킴이라 하자. 그러면
$X$의 닫힌점들은 $K$-유리점들이고(\hyperref[I.3.4.4-ko]{3.4.4}),
$K$에 값을 갖는 $X$의 점들과 표준적으로 동일시된다.
\end{corollary}

\begin{proposition}[6.4.4]
\label{I.6.4.4-ko}
$X$를 체 $K$ 위의 대수적 준스킴이라 하자. 다음 성질들은 동치이다.
\begin{enumerate}
  \item[\rm(a)] $X$는 아르틴이다.
  \item[\rm(b)] $X$의 바탕공간은 이산공간이다.
  \item[\rm(c)] $X$의 바탕공간에는 닫힌점이 유한개뿐이다.
  \item[\rm(c')] $X$의 바탕공간은 유한집합이다.
  \item[\rm(d)] $X$의 모든 점은 닫힌점이다.
  \item[\rm(e)] $X$는 $K$ 위의 유한 차원 대수 $A$에 대하여
    $\operatorname{Spec}(A)$와 동형이다.
\end{enumerate}
\end{proposition}

$X$가 뇌터이므로 (\hyperref[I.6.2.2-ko]{6.2.2})에 따라 조건 (a), (b),
(d)는 동치이고 (c)와 (c')를 함의한다. 또한 (e)가 (a)를 함의함은
명백하다. (c)가 (d)와 (e)를 함의함을 보이는 일만 남는다. $X$가 아핀인
경우만 다루면 된다. 이때 $A(X)$는 $K$ 위의 유한 생성 대수이고
(\hyperref[I.6.3.3-ko]{6.3.3}), 따라서 야콥슨 환이다
(\cite[p.~3-11 및 3-12]{I-1}). 가정에 따라 이 환에는 극대 아이디얼이
유한개뿐이다. 소 아이디얼들의 유한 교집합이 소 아이디얼이면 그 가운데
하나와 같아야 하므로, $A(X)$의 모든 소 아이디얼은 극대 아이디얼이다.
따라서 (d)가 따른다. 또한 (\hyperref[I.6.2.2-ko]{6.2.2})에 따라
$A(X)$는 유한 생성 아르틴 $K$-대수이고, 따라서 반드시
\emph{유한 차원}이다 \cite{I-21}.

\begin{env}[6.4.5]
\label{I.6.4.5-ko}
(\hyperref[I.6.4.4-ko]{6.4.4})의 조건들이 성립할 때 $X$를
\emph{$K$ 위의 유한 스킴}(\emph{참조}
(\agref{II.6.1.1-ko}{II, 6.1.1})) 또는 \emph{유한 $K$-스킴}이라고 한다.
그 \emph{계수(階數)}는 $[A:K]$이며 $\operatorname{rg}_K(X)$로도 쓴다.
$X$, $Y$가 $K$ 위의 두 유한 스킴이면
\[
  \operatorname{rg}_K(X\amalg Y)
  =\operatorname{rg}_K(X)+\operatorname{rg}_K(Y)
  \tag{6.4.5.1}\label{I.6.4.5.1-ko}
\]
\[
  \operatorname{rg}_K(X\times_K Y)
  =\operatorname{rg}_K(X)\operatorname{rg}_K(Y)
  \tag{6.4.5.2}\label{I.6.4.5.2-ko}
\]
이다. 이는 (\hyperref[I.3.2.2-ko]{3.2.2})에서 따른다.
\end{env}

\begin{corollary}[6.4.6]
\label{I.6.4.6-ko}
$X$를 체 $K$ 위의 유한 스킴이라 하자. $K$의 모든 확대체 $K'$에 대하여
$X\otimes_K K'$는 $K'$ 위의 유한 스킴이고, $K'$ 위의 계수는 $X$의
$K$ 위의 계수와 같다.
\end{corollary}

실제로 $A=A(X)$라 하면 $[A\otimes_K K':K']=[A:K]$이다.

\begin{corollary}[6.4.7]
\label{I.6.4.7-ko}
$X$를 체 $K$ 위의 유한 스킴이라 하고
$n=\sum_{x\in X}[k(x):K]_s$라 놓자.
\emph{(확대체 $K'/K$에 대하여 $[K':K]_s$는 $K'$에 포함되는 $K$의
가장 큰 대수적 분리확대체의 $K$ 위의 \emph{분리차수}임을 상기하자.)}
\oldpage[I]{149}
그러면 $K$의 모든 대수적으로 닫힌 확대체 $\Omega$에 대하여
$X\otimes_K\Omega$의 바탕공간에는 점이 정확히 $n$개 있고, 이 점들은
$X(\Omega)_K$의 원소들과 동일시된다.
\end{corollary}

(\hyperref[I.6.2.2-ko]{6.2.2})에 따라 환 $A=A(X)$가 국소환인 경우만
다루면 된다. $\mathfrak m$을 그 극대 아이디얼, $L=A/\mathfrak m$을
$K$의 대수적 확대체인 잉여체라 하자. 그러면 $X(\Omega)_K$의 원소들은
$X\otimes_K\Omega$의 $\Omega$-단면들과 전단사로 대응하고
(\hyperref[I.3.4.1-ko]{3.4.1}, \hyperref[I.3.3.14-ko]{3.3.14}), 또한
$L$에서 $\Omega$로 가는 $K$-준동형들과 전단사로 대응한다
(\hyperref[I.1.7.3-ko]{1.7.3}). 따라서
(Bourbaki, \emph{Alg.}, 제V장, \textsection7, n\textsuperscript{o}~5,
명제~8)과 (\hyperref[I.6.4.3-ko]{6.4.3})을 고려하면 결론이 따른다.

\begin{env}[6.4.8]
\label{I.6.4.8-ko}
(\hyperref[I.6.4.7-ko]{6.4.7})에서 정의한 수 $n$을 $A$(또는 $X$)의
$K$ 위의 \emph{분리차수}, 또는 \emph{$X$의 점의 기하적 개수}라고 한다.
따라서 이는 $X(\Omega)_K$의 원소 수와 같다. 이 정의에서 곧바로,
$K$의 모든 확대체 $K'$에 대하여 $X\otimes_K K'$의 점의 기하적 개수가
$X$의 점의 기하적 개수와 같음이 따른다. 이 수를 $n(X)$라 하자. $X$, $Y$가
$K$ 위의 두 유한 스킴이면 명백히
\[
  n(X\amalg Y)=n(X)+n(Y).
  \tag{6.4.8.1}\label{I.6.4.8.1-ko}
\]
이다.

같은 가정 아래에서
\[
  n(X\times_K Y)=n(X)n(Y)
  \tag{6.4.8.2}\label{I.6.4.8.2-ko}
\]
도 성립한다. 이는 $n(X)$를 $X(\Omega)_K$의 원소 수로 해석하고 공식
(\hyperref[I.3.4.3.1-ko]{3.4.3.1})을 적용하면 곧바로 따른다.
\end{env}

\begin{proposition}[6.4.9]
\label{I.6.4.9-ko}
$K$를 체, $X$, $Y$를 두 $K$-대수적 준스킴,
$f:X\to Y$를 $K$-사상이라 하자. 또한 $\Omega$를 $K$ 위의 초월차수가
무한인 대수적으로 닫힌 확대체라 하자. $f$가 전사이기 위한 필요충분조건은
$f$에 대응하는 사상 $X(\Omega)_K\to Y(\Omega)_K$가 전사인 것이다
(\hyperref[I.3.4.1-ko]{3.4.1}).
\end{proposition}

필요성은 (\hyperref[I.6.3.10-ko]{6.3.10})에서 따른다. 여기서 $f$가
반드시 유한형인 것은 (\hyperref[I.6.3.9-ko]{6.3.9})에 따른다.
충분성을 보이기 위해 (\hyperref[I.6.3.10-ko]{6.3.10})에서와 같이
논증한다. 실제로 모든 $y\in Y$에 대하여 $k(y)$는 $K$의 유한 생성
확대체이고, 따라서 $\Omega$의 부분체와 $K$-동형이다.

\begin{remark}[6.4.10]
\label{I.6.4.10-ko}
제IV장에서 (\hyperref[I.6.4.9-ko]{6.4.9})의 결론이 $\Omega$의 $K$ 위의
초월차수에 관한 가정 없이도 성립함을 보일 것이다.
\end{remark}

\begin{proposition}[6.4.11]
\label{I.6.4.11-ko}
$f:X\to Y$가 유한형 사상이면, 모든 $y\in Y$에 대하여 섬유
$f^{-1}(y)$는 잉여체 $k(y)$ 위의 대수적 준스킴이고, 모든
$x\in f^{-1}(y)$에 대하여 $k(x)$는 $k(y)$의 유한 생성 확대체이다.
\end{proposition}

$f^{-1}(y)=X\otimes_Y k(y)$이므로(\hyperref[I.3.6.3-ko]{3.6.3}), 이 명제는
(\hyperref[I.6.3.4-ko]{6.3.4, (iv)})와
(\hyperref[I.6.3.3-ko]{6.3.3})에서 따른다.

\begin{proposition}[6.4.12]
\label{I.6.4.12-ko}
$f:X\to Y$, $g:Y'\to Y$를 두 사상이라 하자.
$X'=X\times_Y Y'$라 놓고 $f'=f_{(Y')}:X'\to Y'$라 하자.
$y'\in Y'$, $y=g(y')$라 하자. 섬유 $f^{-1}(y)$가 $k(y)$ 위의 유한
대수적 스킴이면, 섬유 $f'^{-1}(y')$도 $k(y')$ 위의 유한 대수적
스킴이고, 그 계수와 점의 기하적 개수는 $f^{-1}(y)$의 것과 같다.
\end{proposition}

섬유의 추이성(\hyperref[I.3.6.5-ko]{3.6.5})을 고려하면 이는
(\hyperref[I.6.4.6-ko]{6.4.6})과 (\hyperref[I.6.4.8-ko]{6.4.8})에서
곧바로 따른다.

\begin{env}[6.4.13]
\label{I.6.4.13-ko}
명제 (\hyperref[I.6.4.11-ko]{6.4.11})은 유한형 사상이 직관적으로
``대수다양체들의 대수적 족''에 해당함을 보여 준다. 여기서 $Y$의 점들은
\oldpage[I]{150}
``매개변수''의 역할을 하며, 이것이 이 사상들에 ``기하적'' 의미를 준다.
유한형이 아닌 사상들은 뒤에서 주로, 예를 들어 국소화나 완비화에 의한
``밑준스킴의 변경'' 문제에 나타날 것이다.
\end{env}

\subsection{사상의 국소적 결정}
\label{subsection:I.6.5-ko}

\begin{proposition}[6.5.1]
\label{I.6.5.1-ko}
$X$, $Y$를 두 $S$-준스킴이라 하고, $Y$가 $S$ 위의 유한형이라고 하자.
또한 $x\in X$, $y\in Y$가 같은 점 $s\in S$ 위에 있다고 하자.
\begin{enumerate}
  \item[\rm(i)] $X$에서 $Y$로 가는 두 $S$-사상
    $f=(\psi,\theta)$, $f'=(\psi',\theta')$가
    $\psi(x)=\psi'(x)=y$를 만족하고, $\mathscr{O}_y$에서
    $\mathscr{O}_x$로 가는 (국소) $\mathscr{O}_s$-준동형
    $\theta_x^\sharp$와 ${\theta'}_x^\sharp$가 같으면, $f$와 $f'$는
    $x$의 어떤 열린 근방에서 일치한다.
  \item[\rm(ii)] 더 나아가 $S$가 국소 뇌터라고 하자. 모든 국소
    $\mathscr{O}_s$-준동형
    $\varphi:\mathscr{O}_y\to\mathscr{O}_x$에 대하여 $X$에서 $x$의
    열린 근방 $U$와 $U$에서 $Y$로 가는 $S$-사상
    $f=(\psi,\theta)$가 존재하여
    $\psi(x)=y$이고 $\theta_x^\sharp=\varphi$이다.
\end{enumerate}
\end{proposition}

\begin{enumerate}
  \item[\rm(i)] 문제는 $S$, $X$, $Y$ 위에서 국소적이므로 $S$, $X$, $Y$가
    각각 환 $A$, $B$, $C$의 아핀 준스킴이라고 해도 된다. 이때 $f$, $f'$는
    각각 $({}^a\varphi,\widetilde{\varphi})$,
    $({}^a\varphi',\widetilde{\varphi}')$의 꼴이며, $\varphi$, $\varphi'$는
    $C$에서 $B$로 가는 두 $A$-준동형으로서
    $\varphi^{-1}(\mathfrak{j}_x)
    ={\varphi'}^{-1}(\mathfrak{j}_x)=\mathfrak{j}_y$를 만족한다. 또한
    $\varphi$, $\varphi'$에서 얻어지는 $C_y$에서 $B_x$로
    가는 준동형 $\varphi_x$, $\varphi'_x$는 같다. 더 나아가 $C$가 $A$ 위의
    유한 생성 대수라고 해도 된다. $c_i$($1\leq i\leq n$)를 $A$-대수 $C$의
    생성원이라 하고 $b_i=\varphi(c_i)$, $b'_i=\varphi'(c_i)$라 놓자. 가정에
    따라 분수환 $B_x$에서 $b_i/1=b'_i/1$이다($1\leq i\leq n$). 이는
    $1\leq i\leq n$에 대하여 $s_i(b_i-b'_i)=0$이 되는 원소
    $s_i\in B-\mathfrak{j}_x$가 존재한다는 뜻이며, 모든 $s_i$가 같은 원소
    $g\in B-\mathfrak{j}_x$라고 해도 된다. 따라서 분수환 $B_g$에서
    $1\leq i\leq n$에 대하여 $b_i/1=b'_i/1$이다. $i_g:B\to B_g$를 표준
    준동형이라 하면 $i_g\circ\varphi=i_g\circ\varphi'$이고, 그러므로
    $f$와 $f'$의 $D(g)$로의 제한은 같다.
  \item[\rm(ii)] (i)와 같은 상황으로 환원할 수 있고, 더 나아가 환 $A$가
    뇌터라고 해도 된다. $c_i$($1\leq i\leq n$)를 $A$-대수 $C$의 생성원이라
    하고, $\alpha:A[X_1,\ldots,X_n]\to C$를 $X_i$를 $c_i$로 보내는
    다항식 대수 $A[X_1,\ldots,X_n]$에서 $C$로 가는 전사 준동형이라 하자
    ($1\leq i\leq n$). 한편 $i_y:C\to C_y$를 표준 준동형이라 하고, 합성
    준동형
    \[
      \beta:A[X_1,\ldots,X_n]
      \xrightarrow{\alpha} C
      \xrightarrow{i_y} C_y
      \xrightarrow{\varphi} B_x
    \]
    을 생각하자. $\mathfrak{a}$를 $\beta$의 핵이라 하자. $A$가 뇌터이므로
    $A[X_1,\ldots,X_n]$도 뇌터이고, 따라서 $\mathfrak{a}$는 유한 생성계
    $Q_j(X_1,\ldots,X_n)$($1\leq j\leq m$)를 갖는다. 한편 각 원소
    $\varphi(i_y(c_i))$는 $b_i\in B$, $s_i\notin\mathfrak{j}_x$에 대하여
    $b_i/s_i$로 쓸 수 있고, 모든 $s_i$가 같은 원소
    $g\in B-\mathfrak{j}_x$라고 해도 된다. 이제 가정에 따라 $B_x$에서
    $Q_j(b_1/g,\ldots,b_n/g)=0$이다. 다음과 같이 놓자.
    \[
      Q_j(X_1/T,\ldots,X_n/T)
      =P_j(X_1,\ldots,X_n,T)/T^{k_j},
    \]
    여기서 $P_j$는 $k_j$차 동차다항식이다. 또한
    $d_j=P_j(b_1,\ldots,b_n,g)\in B$라 하자. 가정에 따라
    $t_j\in B-\mathfrak{j}_x$에 대하여 $t_jd_j=0$이다
    ($1\leq j\leq m$). 모든 $t_j$가
\oldpage[I]{151}
    같은 원소 $h\in B-\mathfrak{j}_x$라고 해도 된다. 따라서
    $1\leq j\leq m$에 대하여
    $P_j(hb_1,\ldots,hb_n,hg)=0$이다. 이제 $X_i$를 $hb_i/hg$로 보내는
    $A[X_1,\ldots,X_n]$에서 분수환 $B_{hg}$로 가는 준동형 $\rho$를
    생각하자($1\leq i\leq n$). 이 준동형에 의한 $\mathfrak a$의 상은
    0이고, 따라서 $\rho$에 의한 $\alpha^{-1}(0)$의 상도
    \emph{특히} 0이다. 그러므로 $\rho$는
    $A[X_1,\ldots,X_n]\xrightarrow{\alpha}C
    \xrightarrow{\gamma}B_{hg}$로 인수분해되며,
    $\gamma(c_i)=hb_i/hg$이다. $i_x:B_{hg}\to B_x$를 표준
    준동형이라 하면 그림
    \[
      \xymatrix{
        C\ar[r]^\gamma\ar[d]_{i_y} & B_{hg}\ar[d]^{i_x}\\
        C_y\ar[r]_\varphi & B_x
      }
      \tag{6.5.1.1}\label{I.6.5.1.1-ko}
    \]
    은 가환임이 명백하다. 따라서 $\varphi=\gamma_x$이고, $\varphi$가 국소
    준동형이므로 ${}^a\gamma(x)=y$이다. 그러므로
    $f=({}^a\gamma,\widetilde{\gamma})$는 $x$의 근방 $D(hg)$에서 $Y$로
    가는, 문제의 조건을 만족하는 $S$-사상이다.
\end{enumerate}

\begin{corollary}[6.5.2]
\label{I.6.5.2-ko}
(\hyperref[I.6.5.1-ko]{6.5.1, (ii)})의 가정 아래에서, 더 나아가 $X$가
$S$ 위의 유한형이면 사상 $f$가 유한형이라고 할 수 있다.
\end{corollary}

이는 (\hyperref[I.6.3.6-ko]{6.3.6})에서 따른다.

\begin{corollary}[6.5.3]
\label{I.6.5.3-ko}
(\hyperref[I.6.5.1-ko]{6.5.1, (ii)})의 가정들이 성립한다고 하자. 더 나아가
$Y$가 정역 준스킴이고 $\varphi$가 단사 준동형이면,
$f=({}^a\gamma,\widetilde{\gamma})$이고 $\gamma$가 단사 준동형이라고 할 수
있다.
\end{corollary}

실제로 $C$가 정역이라고 해도 되고(\hyperref[I.5.1.4-ko]{5.1.4}), 따라서
$i_y$는 단사 준동형이다. 그러므로 그림
(\hyperref[I.6.5.1.1-ko]{6.5.1.1})에서 $\gamma$가 단사 준동형임이 따른다.

\begin{proposition}[6.5.4]
\label{I.6.5.4-ko}
$f=(\psi,\theta):X\to Y$를 유한형 사상, $x$를 $X$의 점,
$y=\psi(x)$라 하자.
\begin{enumerate}
  \item[\rm(i)] $f$가 점 $x$에서 국소 몰입 사상이기 위한
    (\hyperref[I.4.5.1-ko]{4.5.1}) 필요충분조건은
    $\theta_x^\sharp:\mathscr{O}_y\to\mathscr{O}_x$가 전사 준동형인 것이다.
  \item[\rm(ii)] 더 나아가 $Y$가 국소 뇌터라고 하자. $f$가 점 $x$에서
    국소 동형사상이기 위한(\hyperref[I.4.5.2-ko]{4.5.2}) 필요충분조건은
    $\theta_x^\sharp$가 동형인 것이다.
\end{enumerate}
\end{proposition}

\begin{enumerate}
  \item[\rm(ii)] (\hyperref[I.6.5.1-ko]{6.5.1})에 따라 $y$의 열린 근방
    $V$와 사상 $g:V\to X$가 존재하여 $g\circ f$(각각 $f\circ g$)가
    정의되고 $x$(각각 $y$)의 어떤 근방에서 항등사상과 일치한다. 이로부터
    $f$가 국소 동형사상임이 곧바로 따른다.
  \item[\rm(i)] 문제는 $X$, $Y$ 위에서 국소적이므로 $X$, $Y$가 각각 환
    $A$, $B$의 아핀 준스킴이라고 해도 된다. 이때
    $f=({}^a\varphi,\widetilde{\varphi})$이고, $\varphi:B\to A$는 $A$를
    유한 생성 $B$-대수로 만드는 환 준동형이다. 또한
    $\varphi^{-1}(\mathfrak{j}_x)=\mathfrak{j}_y$이고, $\varphi$에서
    얻어지는 준동형 $\varphi_x:B_y\to A_x$는 전사이다. $(t_i)$
    ($1\leq i\leq n$)를 $B$-대수 $A$의 생성계라 하자. $\varphi_x$에 관한
    가정에 따라 $b_i\in B$와 $c\in B-\mathfrak{j}_y$가 존재하여 분수환
    $A_x$에서 $1\leq i\leq n$에 대하여
    $t_i/1=\varphi(b_i)/\varphi(c)$이다. 따라서
    (\hyperref[I.1.3.3-ko]{1.3.3})에 따라 $a\in A-\mathfrak{j}_x$가
    존재하여, $g=a\varphi(c)$라 놓으면
    $t_i/1=a\varphi(b_i)/g$가 \emph{분수환 $A_g$ 안에서}도 성립한다.
    한편 가정에 따라 $\varphi(B)$를 계수환으로 하는 다항식
    $Q(X_1,\ldots,X_n)$가
\oldpage[I]{152}
    존재하여 $a=Q(t_1,\ldots,t_n)$이다. 다음과 같이 놓자.
    \[
      Q(X_1/T,\ldots,X_n/T)
      =P(X_1,\ldots,X_n,T)/T^m,
    \]
    여기서 $P$는 $m$차 동차다항식이다. 환 $A_g$에서
    \[
      a/1
      =a^mP(\varphi(b_1),\ldots,\varphi(b_n),\varphi(c))/g^m
      =a^m\varphi(d)/g^m
    \]
    이며, 여기서 $d\in B$이다. $A_g$에서
    $g/1=(a/1)(\varphi(c)/1)$은 정의에 따라 가역이므로 $a/1$과
    $\varphi(c)/1$도 가역이고, 따라서
    \[
      a/1=(\varphi(d)/1)(\varphi(c)/1)^{-m}
    \]
    으로 쓸 수 있다. 그러므로 $\varphi(d)/1$도 $A_g$에서 가역이다.
    이제 $h=cd$라 놓자. $\varphi(h)/1$이 $A_g$에서 가역이므로 합성
    준동형
    \[
      B\xrightarrow{\varphi}A\longrightarrow A_g
    \]
    은 다음과 같이 인수분해된다.
    \[
      B\longrightarrow B_h\xrightarrow{\gamma}A_g
    \]
    (\hyperref[0.1.2.4-ko]{0, 1.2.4}). $\gamma$가
    \emph{전사}임을 보이자. $A_g$에서 $B_h$의 상이 $t_i/1$과
    $(g/1)^{-1}$을 포함함을 보이면 충분하다. 그런데
    \[
      (g/1)^{-1}
      =(\varphi(c)/1)^{m-1}(\varphi(d)/1)^{-1}
      =\gamma(c^m/h),
    \]
    이고
    \[
      a/1=\gamma(d^{m+1}/h^m),
    \]
    이므로
    \[
      (a\varphi(b_i))/1=\gamma(b_i d^{m+1}/h^m)
    \]
    이다. 또한
    $t_i/1=(a\varphi(b_i)/1)(g/1)^{-1}$이므로 주장이 증명된다. $h$의
    선택에 따라 $\psi(D(g))\subset D(h)$이고, $f$를 $D(g)$에 제한한
    사상은 $({}^a\gamma,\widetilde{\gamma})$와 같다. $\gamma$가 전사이므로
    이 제한은 $D(g)$에서 $D(h)$로 가는 닫힌 몰입 사상이다
    (\hyperref[I.4.2.3-ko]{4.2.3}).
\end{enumerate}

\begin{corollary}[6.5.5]
\label{I.6.5.5-ko}
$f=(\psi,\theta):X\to Y$를 유한형 사상이라 하자. $X$가 기약이라고 하고,
그 일반점을 $x$라 하며 $y=\psi(x)$라 놓자.
\begin{enumerate}
  \item[\rm(i)] $f$가 $X$의 어떤 점에서 국소 몰입 사상이기 위한
    필요충분조건은
    $\theta_x^\sharp:\mathscr{O}_y\to\mathscr{O}_x$가 전사 준동형인 것이다.
  \item[\rm(ii)] 더 나아가 $Y$가 기약이고 국소 뇌터라고 하자. $f$가
    $X$의 어떤 점에서 국소 동형사상이기 위한 필요충분조건은 $y$가 $Y$의
    일반점인 것, 다시 말해(\hyperref[0.2.1.4-ko]{0, 2.1.4}) $f$가
    \emph{우세 사상}인 것과 $\theta_x^\sharp$가 동형인 것, 다시 말해
    $f$가 \emph{쌍유리 사상}인 것이다
    (\hyperref[I.2.2.9-ko]{2.2.9}).
\end{enumerate}
\end{corollary}

(i)는 $X$의 공집합이 아닌 모든 열린집합이 $x$를 포함한다는 사실을
고려하면 (\hyperref[I.6.5.4-ko]{6.5.4, (i)})에서 명백히 따른다. 마찬가지로
(ii)는 (\hyperref[I.6.5.4-ko]{6.5.4, (ii)})에서 따른다.

\subsection{준콤팩트 사상과 국소 유한형 사상}
\label{subsection:I.6.6-ko}

\begin{definition}[6.6.1]
\label{I.6.6.1-ko}
사상 $f:X\to Y$에 의한 $Y$의 모든 준콤팩트 열린집합의 역상이
준콤팩트이면, $f$를 \emph{준콤팩트} 사상이라 한다.
\end{definition}

$\mathfrak{B}$를 준콤팩트 열린집합들(예를 들어 아핀 열린집합들)로
이루어진 $Y$의 위상의 기저라 하자. $f$가 준콤팩트이기 위한 필요충분조건은
$\mathfrak{B}$의 모든 원소의 $f$에 의한 역상이 준콤팩트인 것, 또는 같은
말로 아핀 열린집합들의 \emph{유한} 합집합인 것이다. 실제로 $Y$의 모든
준콤팩트 열린집합은 $\mathfrak{B}$의 유한개의 원소의 합집합이다. 예를 들어
$X$가 \emph{준콤팩트}이고 $Y$가 \emph{아핀}이면 \emph{모든} 사상
$f:X\to Y$는 준콤팩트이다. 실제로 $X$는 유한개의 아핀 열린집합 $U_i$의
합집합이고, $Y$의 모든 아핀 열린집합 $V$에 대하여
$U_i\cap f^{-1}(V)$는 아핀이므로
(\hyperref[I.5.5.10-ko]{5.5.10}) 준콤팩트이다.

$f:X\to Y$가 준콤팩트 사상이면, $Y$의 모든 열린집합 $V$에 대하여
$f$를 $f^{-1}(V)$에 제한한 사상 $f^{-1}(V)\to V$가 준콤팩트임은
명백하다. 역으로 $(U_\alpha)$가 $Y$의 열린 덮개이고, 사상 $f:X\to Y$의
제한들 $f^{-1}(U_\alpha)\to U_\alpha$가 준콤팩트이면, $f$도
준콤팩트이다.

\begin{definition}[6.6.2]
\label{I.6.6.2-ko}
사상 $f:X\to Y$가 다음 조건을 만족하면 \emph{국소 유한형} 사상이라 한다.
모든 $x\in X$에 대하여 $x$의 열린 근방 $U$와
$f(U)$를 포함하는 $y=f(x)$의 열린 근방 $V$가 존재하고, $f$의
\oldpage[I]{153}
$U$로의 제한이 $U$에서 $V$로 가는 유한형 사상이다. 이때 $X$를 $Y$
위의 국소 유한형 준스킴 또는 국소 유한형 $Y$-준스킴이라고도 한다.
\end{definition}

(\hyperref[I.6.3.2-ko]{6.3.2})에 따라 $f$가 국소 유한형이면, $Y$의 모든
열린집합 $W$에 대하여 $f$를 $f^{-1}(W)$에 제한한 사상
$f^{-1}(W)\to W$도 국소 유한형임이 곧바로 따른다.

$Y$가 국소 뇌터이고 $X$가 $Y$ 위의 국소 유한형이면,
(\hyperref[I.6.3.7-ko]{6.3.7})에 따라 $X$는 국소 뇌터이다.

\begin{proposition}[6.6.3]
\label{I.6.6.3-ko}
사상 $f:X\to Y$가 유한형이기 위한 필요충분조건은 $f$가 준콤팩트이고
국소 유한형인 것이다.
\end{proposition}

조건의 필요성은 (\hyperref[I.6.3.1-ko]{6.3.1})과
(\hyperref[I.6.6.1-ko]{6.6.1}) 뒤의 주에서 곧바로 따른다. 역으로 이
조건들이 성립한다고 하고, $U$를 $Y$의 아핀 열린집합, $A$를 그 환이라
하자. 모든
$x\in f^{-1}(U)$에 대하여 가정에 따라 $x$의 근방
$V(x)\subset f^{-1}(U)$와 $y=f(x)$의 근방 $W(x)\subset U$가 존재하여,
$W(x)$는 $f(V(x))$를 포함하고 $f$를 $V(x)$에 제한한 사상
$V(x)\to W(x)$는 유한형이다. $W(x)$를 $y=f(x)$의
$D(g)$ 꼴인 근방 $W_1(x)\subset W(x)$로($g\in A$), $V(x)$를
$V(x)\cap f^{-1}(W_1(x))$로 바꾸면, $W(x)$가 $D(g)$ 꼴이라고 해도 된다.
따라서 그 환이 $A[1/g]$로 쓰이므로 $W(x)$는 $U$ 위의 유한형이고, 그러므로
$V(x)$도 $U$ 위의 유한형이다. 더 나아가 $f^{-1}(U)$는 가정에 따라
준콤팩트이므로 유한개의 열린집합 $V(x_i)$의 합집합이다. 이로써 증명이
끝난다.

\begin{proposition}[6.6.4]
\label{I.6.6.4-ko}
\begin{enumerate}
  \item[\rm(i)] 몰입 사상 $X\to Y$가 닫힌 몰입 사상이거나, $Y$의
    바탕공간이 국소 뇌터이거나, $X$의 바탕공간이 뇌터이면 이 몰입 사상은
    준콤팩트이다.
  \item[\rm(ii)] 두 준콤팩트 사상의 합성은 준콤팩트이다.
  \item[\rm(iii)] $f:X\to Y$가 준콤팩트 $S$-사상이면, 밑준스킴의 모든
    확장 $g:S'\to S$에 대하여
    $f_{(S')}:X_{(S')}\to Y_{(S')}$도 준콤팩트이다.
  \item[\rm(iv)] $f:X\to X'$와 $g:Y\to Y'$가 두 준콤팩트 $S$-사상이면
    $f\times_S g$도 준콤팩트이다.
  \item[\rm(v)] 두 사상 $f:X\to Y$, $g:Y\to Z$의 합성이 준콤팩트이고,
    또한 $g$가 분리 사상이거나 $X$의 바탕공간이 국소 뇌터이면 $f$는
    준콤팩트이다.
  \item[\rm(vi)] 사상 $f$가 준콤팩트이기 위한 필요충분조건은
    $f_{\mathrm{red}}$가 준콤팩트인 것이다.
\end{enumerate}
\end{proposition}

(vi)는 사상의 준콤팩트성이 바탕공간 사이의 대응하는 연속사상에만
의존하므로 명백하다. 마찬가지로 $X$의 바탕공간이 국소 뇌터라고 가정한
경우에 해당하는 (v)를 증명하자. $h=g\circ f$라 놓고, $U$를 $Y$의
준콤팩트 열린집합이라 하자. $g(U)$는 $Z$에서 준콤팩트이지만 반드시
열린집합인 것은 아니므로, 유한개의 준콤팩트 열린집합 $V_j$의 합집합에
포함된다(\hyperref[I.2.1.3-ko]{2.1.3}). 따라서 $f^{-1}(U)$는
$h^{-1}(V_j)$들의 합집합에 포함된다. 이들은 $X$의 준콤팩트
부분공간들이므로 뇌터 부분공간들이다. 그러므로
(\hyperref[0.2.2.3-ko]{0, 2.2.3})에 따라 $f^{-1}(U)$는 뇌터 공간이고,
\emph{따라서 특히} 준콤팩트이다.

나머지 주장들을 증명하기 위해서는 (i), (ii), (iii)를 증명하면 충분하다
(\hyperref[I.5.5.12-ko]{5.5.12}). 그런데 (ii)는 명백하고, (i)는 $Y$의
바탕공간이 국소 뇌터이거나 $X$의 바탕공간이 뇌터인 경우에는
(\hyperref[I.6.3.5-ko]{6.3.5})에서 따르며 닫힌 몰입 사상의 경우에는
명백하다. (iii)를 증명하기 위하여,
\oldpage[I]{154}
$S=Y$인 경우로 환원할 수 있다(\hyperref[I.3.3.11-ko]{3.3.11}).
$f'=f_{(S')}$라 놓고, $U'$를 $S'$의 준콤팩트 열린집합이라 하자. 모든
$s'\in U'$에 대하여 $T$를 $S$에서 $g(s')$의 아핀 열린 근방이라 하고,
$W$를 $U'\cap g^{-1}(T)$에 포함되는 $s'$의 아핀 열린 근방이라 하자.
$f'^{-1}(W)$가 준콤팩트임을 보이면 충분하다. 다시 말해 $S$와 $S'$가
아핀인 경우에 $X\times_S S'$의 바탕공간이 준콤팩트임을 보이는 문제로
환원된다. 그런데 이때 $X$는 가정에 따라 유한개의 아핀 열린집합 $V_j$의
합집합이고, $X\times_S S'$의 바탕공간은 아핀 스킴
$V_j\times_S S'$들의 바탕공간의 합집합이므로
(\hyperref[I.3.2.2-ko]{3.2.2}와 \hyperref[I.3.2.7-ko]{3.2.7})에 따라 명제의
증명이 끝난다.

또한 $X=X'\amalg X''$가 두 준스킴의 합이면, 사상 $f:X\to Y$가
준콤팩트이기 위한 필요충분조건은 $f$를 $X'$와 $X''$에 제한한 사상들이
준콤팩트인 것임을 주목하자.

\begin{proposition}[6.6.5]
\label{I.6.6.5-ko}
$f:X\to Y$를 준콤팩트 사상이라 하자. $f$가 우세 사상이기 위한
필요충분조건은 $Y$의 모든 기약 성분의 일반점 $y$에 대하여
$f^{-1}(y)$가 $X$의 어떤 기약 성분의 일반점을 포함하는 것이다.
\end{proposition}

이 조건이 충분함은($f$가 준콤팩트 사상이라고 가정하지 않아도) 명백하다.
필요성을 보이기 위하여 $y$의 아핀 열린 근방 $U$를 생각하자.
$f^{-1}(U)$는 준콤팩트이므로 아핀 열린집합 $V_i$들의 \emph{유한}
합집합이고, $f$가 우세 사상이라는 가정에 따라 $y$는 $U$에서 어떤
$f(V_i)$의 폐포에 속한다. $X$와 $Y$가 축소라고 해도 됨은 명백하다.
$V_i$의 기약 성분의 $X$에서의 폐포는 $X$의 기약 성분이므로
(\hyperref[0.2.1.6-ko]{0, 2.1.6}), $X$를 $V_i$로, $Y$를
$\overline{f(V_i)}\cap U$를 바탕공간으로 갖는 $U$의 축소 닫힌
부분준스킴으로 바꿀 수 있다
(\hyperref[I.5.2.1-ko]{5.2.1}). 따라서 $X=\operatorname{Spec}(A)$와
$Y=\operatorname{Spec}(B)$가 아핀이고 축소인 경우의 명제를 증명하면 된다.
$f$가 우세 사상이므로 $B$는 $A$의 부분환이고
(\hyperref[I.1.2.7-ko]{1.2.7}), 이제 명제는 $B$의 모든 극소 소 아이디얼이
$B$와 $A$의 어떤 극소 소 아이디얼의 교집합이라는 사실에서 따른다
(\hyperref[0.1.5.8-ko]{0, 1.5.8}).

\begin{proposition}[6.6.6]
\label{I.6.6.6-ko}
\begin{enumerate}
  \item[\rm(i)] 모든 국소 몰입 사상은 국소 유한형이다.
  \item[\rm(ii)] 두 사상 $f:X\to Y$, $g:Y\to Z$가 국소 유한형이면
    $g\circ f$도 국소 유한형이다.
  \item[\rm(iii)] $f:X\to Y$가 국소 유한형 $S$-사상이면, 밑준스킴의
    모든 확장 $S'\to S$에 대하여
    $f_{(S')}:X_{(S')}\to Y_{(S')}$도 국소 유한형이다.
  \item[\rm(iv)] $f:X\to X'$와 $g:Y\to Y'$가 두 국소 유한형
    $S$-사상이면 $f\times_S g$도 국소 유한형이다.
  \item[\rm(v)] 두 사상의 합성 $g\circ f$가 국소 유한형이면 $f$도
    국소 유한형이다.
  \item[\rm(vi)] 사상 $f$가 국소 유한형이면 $f_{\mathrm{red}}$도
    국소 유한형이다.
\end{enumerate}
\end{proposition}

(\hyperref[I.5.5.12-ko]{5.5.12})에 따라 (i), (ii), (iii)를 증명하면
충분하다. $j:X\to Y$가 국소 몰입 사상이면, 모든 $x\in X$에 대하여
$Y$에서 $j(x)$의 열린 근방 $V$와 $X$에서 $x$의 열린 근방 $U$가
존재하여 $j$를 $U$에 제한한 사상은 닫힌 몰입 사상 $U\to V$이다
(\hyperref[I.4.5.1-ko]{4.5.1}). 따라서 이 제한은 유한형이다. (ii)를
증명하기 위하여 점 $x\in X$를 생각하자. 가정에 따라 $g(f(x))$의 열린
근방 $W$와 $f(x)$의 열린 근방 $V$가 존재하여 $g(V)\subset W$이고
$V$는 $W$ 위의 유한형이다. 더 나아가 $f^{-1}(V)$는 $V$ 위의 국소
유한형이므로(\hyperref[I.6.6.2-ko]{6.6.2}), $x$의 열린 근방 $U$가
\oldpage[I]{155}
$f^{-1}(V)$에 포함되고 $V$ 위의 유한형이 되게 할 수 있다. 따라서
$g(f(U))\subset W$이고 $U$는 $W$ 위의 유한형이다
(\hyperref[I.6.3.4-ko]{6.3.4, (ii)}). 끝으로 (iii)를 증명하기 위해서는
$Y=S$인 경우로 환원할 수 있다(\hyperref[I.3.3.11-ko]{3.3.11}). 모든
$x'\in X'=X_{(S')}$에 대하여 $x$를 $X$에서 $x'$의 상, $s$를 $S$에서
$x$의 상, $T$를 $s$의 열린 근방, $T'$를 $S'$에서 $T$의 역상, $U$를
$x$의 열린 근방으로서 그 상이 $T$에 포함되고 $T$ 위의 유한형인 것으로
잡자. 그러면 $U\times_S T'=U\times_T T'$는 $x'$의 열린 근방이고
(\hyperref[I.3.2.7-ko]{3.2.7}), $T'$ 위의 유한형이다
(\hyperref[I.6.3.4-ko]{6.3.4, (iv)}).

\begin{corollary}[6.6.7]
\label{I.6.6.7-ko}
$X$, $Y$를 $S$ 위의 국소 유한형인 두 $S$-준스킴이라 하자. $S$가 국소
뇌터이면 $X\times_S Y$는 국소 뇌터이다.
\end{corollary}

실제로 $X$는 $S$ 위의 국소 유한형이므로 국소 뇌터이고,
$X\times_S Y$는 $X$ 위의 국소 유한형이므로 역시 국소 뇌터이다.

\begin{remark}[6.6.8]
\label{I.6.6.8-ko}
명제 (\hyperref[I.6.3.10-ko]{6.3.10})과 그 증명은 사상 $f$가 국소
유한형이라고만 가정하는 경우로 곧바로 확장된다. 마찬가지로 명제
(\hyperref[I.6.4.2-ko]{6.4.2})와 (\hyperref[I.6.4.9-ko]{6.4.9})는 그
명제들에 나타나는 준스킴 $X$, $Y$가 체 $K$ 위의 국소 유한형이라고만
가정할 때에도 성립한다.
\end{remark}
% ===== END EXACT INLINE INPUT: c1s6.tex =====


% ===== BEGIN EXACT INLINE INPUT: c1s7.tex =====
\section{유리 사상}
\label{section:I.7-ko}

\subsection{유리 사상과 유리함수}
\label{subsection:I.7.1-ko}

\begin{env}[7.1.1]
\label{I.7.1.1-ko}
$X$, $Y$를 두 준스킴, $U$, $V$를 $X$의 두 조밀한 열린집합,
$f$(각각 $g$)를 $U$(각각 $V$)에서 $Y$로 가는 사상이라 하자.
$f$와 $g$가 $U\cap V$ 안의 어떤 조밀한 열린집합에서 일치하면
이들을 \emph{동치}라고 하겠다. $X$의 조밀한 열린집합들의 유한
교집합은 $X$의 조밀한 열린집합이므로 이 관계가
\emph{동치관계}임은 자명하다.
\end{env}

\begin{definition}[7.1.2]
\label{I.7.1.2-ko}
두 준스킴 $X$, $Y$가 주어졌다고 하자.
(\hyperref[I.7.1.1-ko]{7.1.1})에서 정의한 관계에 따른 $X$의 조밀한
열린 부분집합에서 $Y$로 가는 사상들의 동치류를
\emph{$X$에서 $Y$로 가는 유리 사상}이라 한다. $X$와 $Y$가
$S$-준스킴일 때, 이 동치류의 어떤 대표가 $S$-사상이면 이를
\emph{$S$-유리 사상}이라 한다. $S$에서 $X$로 가는 모든
$S$-유리 사상을 \emph{$X$의 $S$-유리 단면}이라 한다. $X$-준스킴
$X\otimes_{\mathbf{Z}}\mathbf{Z}[T]$($T$는 부정원)의 모든
$X$-유리 단면을 \emph{준스킴 $X$ 위의 유리함수}라 한다.

$S$-준스킴만을 다룰 때에는 혼동의 염려가 없으면 용어를 남용하여
``$S$-유리 사상'' 대신 ``유리 사상''이라고 하겠다.

$f$를 $X$에서 $Y$로 가는 유리 사상, $U$를 $X$의 열린집합이라 하자.
$f$의 동치류에 속하는 사상 $f_1$, $f_2$가 각각 $X$의 조밀한
열린집합 $V$, $W$ 위에서 정의되어 있으면, 그 제한
$f_1|(U\cap V)$와 $f_2|(U\cap W)$는 $U\cap V\cap W$ 안의,
$U$에서 조밀한 어떤 열린집합에서 일치하므로 동치이다. 따라서
동치류 $f$는
$U$에서 $Y$로 가는 유리 사상을 정의한다. 이를
\emph{$f$의 $U$로의 제한}이라 하고 $f|U$로 쓴다.

각 $S$-사상 $f:X\to Y$에 $S$-유리 사상 하나를 대응시킬 때 그 사상을
\oldpage[I]{156}
$f$가 속하는 동치류로 정하면,
$\operatorname{Hom}_S(X,Y)$에서 $X$에서 $Y$로 가는 $S$-유리 사상들의
집합으로 가는 표준 사상이 정의된다. $X$의 $Y$-유리 단면들의 집합을
$\Gamma_{\mathrm{rat}}(X/Y)$로 나타내면 표준 사상
$\Gamma(X/Y)\to\Gamma_{\mathrm{rat}}(X/Y)$를 얻는다. 또한 $X$와 $Y$가
두 $S$-준스킴이면, $X$에서 $Y$로 가는 $S$-유리 사상들의 집합은
$\Gamma_{\mathrm{rat}}((X\times_S Y)/X)$와 표준적으로 동일시됨이
자명하다
(\hyperref[I.3.3.14-ko]{3.3.14}).
\end{definition}

\begin{env}[7.1.3]
\label{I.7.1.3-ko}
(\hyperref[I.7.1.2-ko]{7.1.2})와
(\hyperref[I.3.3.14-ko]{3.3.14})에서 곧바로 알 수 있듯이,
$X$ 위의 유리함수들은 $X$의 모든 곳에서 조밀한 열린집합 위의
\emph{구조층 $\mathscr{O}_X$의 단면들의 동치류}와 표준적으로
동일시된다. 여기서 두 단면은 그 정의역들의 교집합에 포함된 모든
곳에서 조밀한 열린집합 위에서 일치할 때 동치이다. 특히 $X$ 위의
유리함수들은 \emph{환} $R(X)$을 이룬다.
\end{env}

\begin{env}[7.1.4]
\label{I.7.1.4-ko}
$X$가 \emph{기약} 준스킴이면 공집합이 아닌 모든 열린집합은 $X$에서
조밀하다. 달리 말하면 $X$의 공집합이 아닌 열린집합들은 $X$의
\emph{일반점 $x$의 열린 근방}들이다. 그러므로 이 경우 $X$의 공집합이
아닌 열린 부분집합들에서 $Y$로 가는 두 사상이 동치라는 것은 두
사상이 $x$에서 \emph{같은 싹}을 갖는다는 뜻이다. 다시 말해,
$X\to Y$인 유리 사상(각각 $S$-유리 사상)들은 $X$의 일반점
$x$에서, $X$의 공집합이 아닌 열린 부분집합들에서 $Y$로 가는
\emph{사상의 싹}(각각 $S$-사상의 싹)들과 동일시된다. 특히 다음을
얻는다.
\end{env}

\begin{proposition}[7.1.5]
\label{I.7.1.5-ko}
$X$가 기약 준스킴이면 $X$ 위의 유리함수환 $R(X)$은 $X$의 일반점
$x$의 국소환 $\mathscr{O}_x$와 표준적으로 동일시된다. 이는 차원
0의 국소환이며, 따라서 $X$가 뇌터 준스킴이면 아르틴 국소환이다.
$X$가 정역 준스킴이면 이는 체이고, 나아가 $X$가 아핀 스킴이면 $A(X)$의
분수체와 동일시된다.
\end{proposition}

앞의 내용과 유리함수를 모든 곳에서 조밀한 열린집합 위의
$\mathscr{O}_X$의 단면과 동일시한 것을 생각하면, 첫째 주장은 한
점에서 층의 줄기를 정의한 것에 지나지 않는다. 나머지 주장들은
$X$가 환 $A$의 아핀 스킴인 경우만 생각하면 충분하다. 이때
$\mathfrak{j}_x$는 $A$의 멱영근기이므로 $\mathscr{O}_x$는 차원
0이다. $A$가 정역이면 $\mathfrak{j}_x=(0)$이므로 $\mathscr{O}_x$는
$A$의 분수체이다. 끝으로 $A$가 뇌터이면 $\mathfrak{j}_x$가
멱영이고(\cite{I-11}, p.~127, cor.~4),
$\mathscr{O}_x=A_{\mathfrak{j}_x}$가 아르틴임이 알려져 있다.

$X$가 \emph{정역 준스킴}이면 모든 $z\in X$에 대하여 환 $\mathscr{O}_z$는
정역이다. $z$를 포함하는 모든 아핀 열린집합 $U$는 $x$도 포함하며,
$A(U)$의 분수체인 $R(U)$는 따라서 $R(X)$와 동일시된다.
그러므로 $R(X)$은 \emph{$\mathscr{O}_z$의 분수체}와도 동일시된다.
$\mathscr{O}_z$를 $R(X)$의 부분환과 표준적으로 동일시하는 사상은
단면의 모든 싹 $s\in\mathscr{O}_z$에, $z$에서 싹 $s$를 갖는
$\mathscr{O}_X$의 단면(반드시 모든 곳에서 조밀한 열린집합 위에서
정의된다)의 동치류인 $X$ 위의 유일한 유리함수를 대응시킨다.

\begin{env}[7.1.6]
\label{I.7.1.6-ko}
이제 $X$가 \emph{유한} 개의 기약 성분 $X_i$($1\leq i\leq n$)를 갖는다고
하자($X$의 바탕공간이 \emph{뇌터}이면 그러하다). $X'_i$를
$X_i$에서 $j\neq i$인 $X_j\cap X_i$들의 합집합을 뺀 $X$의
열린집합이라 하자. $X'_i$는 기약이고, 그 일반점 $x_i$는 $X_i$의
일반점이며, $X'_i$들은 서로소이고 그 합집합은 $X$에서 조밀하다
(\hyperref[0.2.1.6-ko]{0, 2.1.6}).
\oldpage[I]{157}
$X$의 모든 곳에서 조밀한 임의의 열린집합 $U$에 대하여
$U_i=U\cap X'_i$는 $X'_i$의 공집합이
아닌 조밀한 열린집합이고, $U_i$들은 서로소이므로
$U'=\bigcup_i U_i$는 $X$에서 조밀하다. $U'$에서 $Y$로 가는 사상을
주는 것은 각 $U_i$에서 $Y$로 가는 사상을 임의로 하나씩 주는 것과
같다. 따라서 다음을 얻는다.
\end{env}

\begin{proposition}[7.1.7]
\label{I.7.1.7-ko}
$X$, $Y$를 두 준스킴(각각 $S$-준스킴)이라 하고, $X$가 일반점
$x_i$를 갖는 유한 개의 기약 성분 $X_i$($1\leq i\leq n$)를 갖는다고
하자. $R_i$를 $x_i$에서, $X$의 열린 부분집합에서 $Y$로 가는 사상
(각각 $S$-사상)의 싹들의 집합이라 하면, $X$에서 $Y$로 가는
유리 사상(각각 $S$-유리 사상)들의 집합은 $R_i$들의 곱
($1\leq i\leq n$)과 동일시된다.
\end{proposition}

\begin{corollary}[7.1.8]
\label{I.7.1.8-ko}
$X$를 뇌터 준스킴이라 하자. $X$ 위의 유리함수환은 아르틴 환이고,
그 국소환 성분들은 $X$의 기약 성분들의 일반점 $x_i$의 국소환
$\mathscr{O}_{x_i}$이다.
\end{corollary}

\begin{corollary}[7.1.9]
\label{I.7.1.9-ko}
$A$를 뇌터 환, $X=\operatorname{Spec}(A)$라 하자. $Q$가 $A$의
극소 소 아이디얼들의 합집합의 여집합이면 $X$ 위의 유리함수환은
분수환 $Q^{-1}A$와 표준적으로 동일시된다.
\end{corollary}

이는 다음 보조정리에서 따른다.

\begin{lemma}[7.1.9.1]
\label{I.7.1.9.1-ko}
$f\in A$에 대하여 $D(f)$가 $X$에서 모든 곳에서 조밀하기 위한
필요충분조건은 $f\in Q$인 것이다. $X$의 모든 조밀한 열린집합은
$f\in Q$인 $D(f)$ 꼴의 열린집합을 하나 포함한다.
\end{lemma}

실제로 이 보조정리가 증명되었다고 하자. 단면환
$\Gamma(D(f),\mathscr{O}_X)$는 $A_f$와 동일시된다
(\hyperref[I.1.3.6-ko]{1.3.6}과 \hyperref[I.1.3.7-ko]{1.3.7}).
$f\in Q$인 $D(f)$들은 $\supset$에 관해 순서가 주어진 $X$의 조밀한
열린집합들의 집합에서 공종집합을 이루므로, 정의
(\hyperref[I.7.1.1-ko]{7.1.1})에 따라 $X$ 위의 유리함수환은
$f\in Q$인 $A_f$들의 귀납극한과 동일시된다. 여기서 준순서관계는
``$g$가 $f$의 배수이다''이고, 이 귀납극한은 곧 $Q^{-1}A$이다
(\hyperref[0.1.4.5-ko]{0, 1.4.5}).

(\hyperref[I.7.1.9.1-ko]{7.1.9.1})을 증명하기 위하여 $X$의 기약
성분들도 $X_i$($1\leq i\leq n$)로 나타내자. $D(f)$가 $X$에서
조밀하면 모든 $1\leq i\leq n$에 대하여
$D(f)\cap X_i\neq\varnothing$이고 그 역도 성립한다. 그런데
$\mathfrak{p}_i=\mathfrak{j}(X_i)$라 놓으면 이는 모든
$1\leq i\leq n$에 대하여 $f\notin\mathfrak{p}_i$라는 뜻이다.
$\mathfrak{p}_i$들은 $A$의 극소 소 아이디얼들이므로
(\hyperref[I.1.1.14-ko]{1.1.14}), 관계들
$f\notin\mathfrak{p}_i$($1\leq i\leq n$)는 $f\in Q$와 동치이다.
이로써 보조정리의 첫째 주장이 증명되었다. 한편 $U$가 $X$의 조밀한
열린집합이면 $U$의 여집합은 $V(\mathfrak{a})$ 꼴이고, 여기서
$\mathfrak{a}$는 어느 $\mathfrak{p}_i$에도 포함되지 않는
아이디얼이다. 따라서 이 아이디얼은 그 합집합에도 포함되지 않고
(\cite{I-10}, p.~13), 어떤 $f\in\mathfrak{a}$는 $Q$에도 속한다.
그러므로 $D(f)\subset U$이고 증명이 끝난다.

\begin{env}[7.1.10]
\label{I.7.1.10-ko}
다시 $X$가 일반점 $x$를 갖는 기약 준스킴이라고 하자. $X$의
공집합이 아닌 모든 열린집합 $U$는 $x$를 포함하므로
$x\in\overline{\{z\}}$인 모든 $z\in X$도 포함한다. 따라서 모든
사상 $U\to Y$는 표준 사상 $\operatorname{Spec}(\mathscr{O}_x)\to X$
(\hyperref[I.2.4.1-ko]{2.4.1})와 합성할 수 있다. 또한 $X$의 공집합이
아닌 두 열린 부분집합에서 $Y$로 가는 두 사상이 $X$의 공집합이 아닌
어떤 열린 부분집합에서 일치하면, 이 두 사상을 각각 그 표준 사상과
합성하여 얻는 두 사상
$\operatorname{Spec}(\mathscr{O}_x)\to Y$는 같다. 다시 말해,
$X$에서 $Y$로 가는 모든 유리 사상에는 잘 정해진 사상
$\operatorname{Spec}(\mathscr{O}_x)\to Y$가 이렇게 대응한다.
\end{env}

\oldpage[I]{158}
\begin{proposition}[7.1.11]
\label{I.7.1.11-ko}
$X$, $Y$를 두 $S$-준스킴이라 하자. $X$는 일반점 $x$를 갖는
기약 준스킴이고 $Y$는 $S$ 위에서 유한형이라고 하자. 같은
$S$-사상 $\operatorname{Spec}(\mathscr{O}_x)\to Y$가 대응하는,
$X$에서 $Y$로 가는 두 $S$-유리 사상은 같다. 나아가 $S$가 국소
뇌터이면 $\operatorname{Spec}(\mathscr{O}_x)$에서 $Y$로 가는 모든
$S$-사상에는 $X$에서 $Y$로 가는 유일한 $S$-유리 사상이 대응한다.
\end{proposition}

$X$의 공집합이 아닌 모든 열린집합이 모든 곳에서 조밀함을 고려하면
이는 (\hyperref[I.6.5.1-ko]{6.5.1})에서 곧바로 따른다.

\begin{corollary}[7.1.12]
\label{I.7.1.12-ko}
$S$가 국소 뇌터이고 (\hyperref[I.7.1.11-ko]{7.1.11})의 나머지
가정들이 성립한다고 하자. 그러면 $X$에서 $Y$로 가는 $S$-유리
사상들은 $S$-준스킴 $\operatorname{Spec}(\mathscr{O}_x)$에 값을 갖는
$S$-준스킴 $Y$의 점들과 동일시된다.
\end{corollary}

이는 (\hyperref[I.7.1.11-ko]{7.1.11})을
(\hyperref[I.3.4.1-ko]{3.4.1})에서 도입한 용어로 다시 쓴 것에
지나지 않는다.

\begin{corollary}[7.1.13]
\label{I.7.1.13-ko}
(\hyperref[I.7.1.12-ko]{7.1.12})의 조건들이 성립한다고 하자.
$s$를 $S$에서 $x$의 상이라 하자. $X$에서 $Y$로 가는 $S$-유리
사상을 주는 것은 $s$ 위에 있는 $Y$의 점 $y$와 국소
$\mathscr{O}_s$-준동형
$\mathscr{O}_y\to\mathscr{O}_x=R(X)$을 주는 것과 같다.
\end{corollary}

이는 (\hyperref[I.7.1.11-ko]{7.1.11})과
(\hyperref[I.2.4.4-ko]{2.4.4})에서 따른다.

특히 다음을 얻는다.

\begin{corollary}[7.1.14]
\label{I.7.1.14-ko}
(\hyperref[I.7.1.12-ko]{7.1.12})의 조건 아래에서, $X$에서 $Y$로 가는
$S$-유리 사상들은 ($Y$가 주어졌을 때) $S$-준스킴
$\operatorname{Spec}(\mathscr{O}_x)$에만 의존한다. 특히 모든
$z\in X$에 대하여 $X$를 $\operatorname{Spec}(\mathscr{O}_z)$로
바꾸어도 이 유리 사상들은 변하지 않는다.
\end{corollary}

실제로 $z\in\overline{\{x\}}$이므로 $x$는
$Z=\operatorname{Spec}(\mathscr{O}_z)$의 일반점이고
$\mathscr{O}_{X,x}=\mathscr{O}_{Z,x}$이다.

$X$가 정역 준스킴이면 $R(X)=\mathscr{O}_x=k(x)$는 체이다
(\hyperref[I.7.1.5-ko]{7.1.5}). 따라서 앞의 따름정리들은 다음과 같이
특수화된다.

\begin{corollary}[7.1.15]
\label{I.7.1.15-ko}
(\hyperref[I.7.1.12-ko]{7.1.12})의 조건들이 성립하고, 나아가 $X$가
정역 준스킴이라고 하자. $s$를 $S$에서 $x$의 상이라 하자. 그러면 $X$에서
$Y$로 가는 $S$-유리 사상들은 $k(s)$의 확대체 $R(X)$에 값을 갖는
$Y\otimes_S k(s)$의 기하적 점들과 동일시된다. 다시 말해, 이러한
유리 사상 하나를 주는 것은 $s$ 위에 있는 점 $y\in Y$와
$k(y)$에서 $k(x)=R(X)$로 가는 $k(s)$-단사 준동형을 주는 것과 같다.
\end{corollary}

실제로 $s$ 위에 있는 $Y$의 점들은 $Y\otimes_S k(s)$의 점들과
동일시되고(\hyperref[I.3.6.3-ko]{3.6.3}), 국소
$\mathscr{O}_s$-준동형 $\mathscr{O}_y\to R(X)$들은
$k(s)$-단사 준동형 $k(y)\to R(X)$들과 동일시된다.

더욱 특별히 다음을 얻는다.

\begin{corollary}[7.1.16]
\label{I.7.1.16-ko}
$k$를 체, $X$, $Y$를 $k$ 위의 두 대수적 준스킴
(\hyperref[I.6.4.1-ko]{6.4.1})이라 하고, 나아가 $X$가 정역
준스킴이라고 하자. 그러면 $X$에서 $Y$로 가는 $k$-유리 사상들은 $k$의 확대체
$R(X)$에 값을 갖는 $Y$의 기하적 점들과 동일시된다
(\hyperref[I.3.4.4-ko]{3.4.4}).
\end{corollary}

\subsection{유리 사상의 정의역}
\label{subsection:I.7.2-ko}

\begin{env}[7.2.1]
\label{I.7.2.1-ko}
$X$, $Y$를 두 준스킴, $f$를 $X$에서 $Y$로 가는 유리 사상이라 하자.
$x$를 포함하는 모든 곳에서 조밀한 열린집합 $U$와 동치류 $f$에 속하는
사상 $U\to Y$가 존재하면 $f$가 \emph{점 $x\in X$에서 정의된다}고 한다.
$f$가 정의되는 점 $x\in X$들의 집합을 \emph{$f$의 정의역}이라 한다.
이는 $X$에서 모든 곳에서 조밀한 열린집합임이 자명하다.
\end{env}

\oldpage[I]{159}
\begin{proposition}[7.2.2]
\label{I.7.2.2-ko}
$X$, $Y$를 두 $S$-준스킴이라 하고, $X$는 축소 준스킴이며 $Y$는 $S$
위에서 분리되어 있다고 하자. $f$를 $X$에서 $Y$로 가는 $S$-유리
사상, $U_0$를 그 정의역이라 하자. 그러면 동치류 $f$에 속하는
$S$-사상 $U_0\to Y$가 유일하게 존재한다.
\end{proposition}

동치류 $f$에 속하는 모든 사상 $U\to Y$에 대하여 반드시
$U\subset U_0$이므로, 이 명제는 다음 보조정리에서 바로 따른다.

\begin{lemma}[7.2.2.1]
\label{I.7.2.2.1-ko}
(\hyperref[I.7.2.2-ko]{7.2.2})의 가정 아래에서, $U_1$, $U_2$를 $X$의
모든 곳에서 조밀한 두 열린집합, $f_i:U_i\to Y$($i=1,2$)를 두
$S$-사상이라 하자. $X$에서 조밀한 열린집합
$V\subset U_1\cap U_2$가 존재하여 $f_1$과 $f_2$가 그 위에서
일치하면, $f_1$과 $f_2$는 $U_1\cap U_2$ 전체에서 일치한다.
\end{lemma}

$X=U_1=U_2$인 경우만 생각하면 충분함이 자명하다. $X$(따라서 $V$)가
축소 준스킴이므로, $X$는 $V$를 포함하는 $X$의 가장 작은 닫힌 부분준스킴이다
(\hyperref[I.5.2.2-ko]{5.2.2}).
$g=(f_1,f_2)_S:X\to Y\times_S Y$라 하자. 가정에 따라 대각선
$T=\Delta_Y(Y)$는 $Y\times_S Y$의 닫힌 부분준스킴이므로,
$Z=g^{-1}(T)$는 $X$의 닫힌 부분준스킴이다
(\hyperref[I.4.4.1-ko]{4.4.1}). $h:V\to Y$를 $f_1$과 $f_2$의 $V$로의
공통 제한이라 하면, $g$의 $V$로의 제한은 $g'=(h,h)_S$이고, 이는
$g'=\Delta_Y\circ h$로 인수분해된다. $\Delta_Y^{-1}(T)=Y$이므로
$g'^{-1}(T)=V$이다. 따라서 $Z$는 $V$를 유도하는 $X$의 닫힌
부분준스킴이고 $V$를 포함하므로 $Z=X$이다. 관계 $g^{-1}(T)=X$에서
(\hyperref[I.4.4.1-ko]{4.4.1})에 따라
$g$는 $\Delta_Y\circ f$로 인수분해되며, 여기서 $f$는 $X\to Y$인
사상이다. 대각
사상의 정의에 따라 $f_1=f_2=f$이다.

(\hyperref[I.7.2.2-ko]{7.2.2})에서 정의한 사상 $U_0\to Y$는 $f$의
동치류에 속하면서 $U_0$을 진정하게 포함하는 $X$의 열린 부분집합으로
\emph{연장할 수 없는} 유일한 사상임이 자명하다.
\emph{(\hyperref[I.7.2.2-ko]{7.2.2})의 가정 아래에서}, $X$에서 $Y$로
가는 유리 사상들을 $X$의 모든 곳에서 조밀한 열린집합에서 $Y$로 가는,
더 큰 열린집합으로 \emph{연장할 수 없는} 사상들과 \emph{동일시}할 수
있다. 이 동일시 아래에서 명제
(\hyperref[I.7.2.2-ko]{7.2.2})로부터 다음이 따른다.

\begin{corollary}[7.2.3]
\label{I.7.2.3-ko}
$X$, $Y$에 대한 가정이 (\hyperref[I.7.2.2-ko]{7.2.2})와 같고,
$U$를 $X$의 모든 곳에서 조밀한 열린집합이라 하자. $U$에서 $Y$로
가는 $S$-사상들과 $U$의 모든 점에서 정의되는, $X$에서 $Y$로 가는
$S$-유리 사상들 사이에는 표준 전단사 대응이 존재한다.
\end{corollary}

실제로 (\hyperref[I.7.2.2-ko]{7.2.2})에 따라, $U$에서 $Y$로 가는 모든
$S$-사상 $f$에 대하여, $f$를 연장하는 $X$에서 $Y$로 가는 유일한 $S$-유리 사상
$\overline f$가 존재한다.

\begin{corollary}[7.2.4]
\label{I.7.2.4-ko}
$S$를 스킴, $X$를 축소 $S$-준스킴, $Y$를 $S$-스킴이라 하고,
$f:U\to Y$를 $X$의 조밀한 열린집합 $U$에서 $Y$로 가는 $S$-사상이라
하자. $\overline f$가 $f$를 연장하는, $X$에서 $Y$로 가는
$\mathbf{Z}$-유리 사상이면, $\overline f$는 그 정의역 위에서
$S$-사상이다. 따라서 이는 $f$를 연장하는, $X$에서 $Y$로 가는
$S$-유리 사상이다.
\end{corollary}

$\varphi:X\to S$, $\psi:Y\to S$를 구조 사상, $U_0$을 $\overline f$의
정의역, $j$를 포함 사상 $U_0\to X$라 하자. 그러면
$\psi\circ\overline f=\varphi\circ j$임을 보이면 충분하다. 이는 $f$가
$S$-사상이므로 (\hyperref[I.7.2.2.1-ko]{7.2.2.1})에서 바로 따른다.

\begin{corollary}[7.2.5]
\label{I.7.2.5-ko}
$X$, $Y$를 두 $S$-준스킴이라 하고, $X$는 축소 준스킴이며 $X$와 $Y$는
$S$ 위에서 분리되어 있다고 하자. $p:Y\to X$를 $S$-사상($Y$를
$X$-준스킴으로 만드는 사상), $U$를 $X$의 모든 곳에서 조밀한
열린집합, $f$를 $Y$의 $U$-단면이라 하자. 그러면 $f$를 연장하는,
$X$에서 $Y$로 가는 유리 사상 $\overline f$는 $Y$의 $X$-유리
단면이다.
\end{corollary}

\oldpage[I]{160}
$\overline f$의 정의역에서 $p\circ\overline f$가 항등사상임을 보이면
된다. $X$가 $S$ 위에서 분리되어 있으므로, 이 역시
(\hyperref[I.7.2.2.1-ko]{7.2.2.1})에서 따른다.

\begin{corollary}[7.2.6]
\label{I.7.2.6-ko}
$X$를 축소 준스킴, $U$를 $X$의 모든 곳에서 조밀한 열린집합이라 하자.
$U$ 위의 구조층 $\mathscr{O}_X$의 단면들과 $U$의 모든 점에서 정의되는
$X$ 위의 유리함수 $f$들 사이에는 표준 전단사 대응이 존재한다.
\end{corollary}

(\hyperref[I.7.2.3-ko]{7.2.3}), (\hyperref[I.7.1.2-ko]{7.1.2}),
(\hyperref[I.7.1.3-ko]{7.1.3})을 고려하면,
$X$-준스킴 $X\otimes_{\mathbf Z}\mathbf Z[T]$가 $X$ 위에서 분리되어
있음을 지적하면 충분하다(\hyperref[I.5.5.1-ko]{5.5.1, (iv)}).

\begin{corollary}[7.2.7]
\label{I.7.2.7-ko}
$Y$를 축소 준스킴, $f:X\to Y$를 분리 사상, $U$를 $Y$의 모든 곳에서
조밀한 열린집합, $g:U\to f^{-1}(U)$를 $f^{-1}(U)$의 $U$-단면이라
하자. $Z$를 $\overline{g(U)}$를 바탕공간으로 갖는 $X$의 축소
부분준스킴이라 하자(\hyperref[I.5.2.1-ko]{5.2.1}). $g$가 $X$의
$Y$-단면의 제한이기 위한 필요충분조건, 다시 말해
(\hyperref[I.7.2.5-ko]{7.2.5})에서 $g$를 연장하는 $Y$에서 $X$로 가는
유리 사상이 모든 곳에서 정의되기 위한 필요충분조건은 $f$의 $Z$로의
제한이 $Z$에서 $Y$로 가는 동형사상인 것이다.
\end{corollary}

$f$의 $f^{-1}(U)$로의 제한은 분리 사상이다
(\hyperref[I.5.5.1-ko]{5.5.1, (i)}). 따라서 $g$는 닫힌 몰입 사상이고
(\hyperref[I.5.4.6-ko]{5.4.6}),
$g(U)=Z\cap f^{-1}(U)$이다. 또한 $Z$의 열린집합 $g(U)$ 위에 $Z$가
유도하는 부분준스킴은 $g$에 대응하는 $f^{-1}(U)$의 닫힌 부분준스킴과
동일하다(\hyperref[I.5.2.1-ko]{5.2.1}). 이제 명제의 조건이 충분함은
자명하다. 실제로 그 조건이 성립하고 $f_Z:Z\to Y$가 $f$의 $Z$로의
제한이며 $\overline g:Y\to Z$가 그 역동형사상이면, $\overline g$는
$g$를 연장한다. 반대로 $g$가 $X$의 $Y$-단면 $h$의 $U$로의 제한이면,
$h$는 닫힌 몰입 사상이다
(\hyperref[I.5.4.6-ko]{5.4.6}). 따라서 $h(Y)$는 닫혀 있고 $Z$에
포함되므로 $Z$와 같다. 그러므로 (\hyperref[I.5.2.1-ko]{5.2.1})에 따라
$h$는 반드시 $Y$에서 $X$의 닫힌 부분준스킴 $Z$로 가는 동형사상이다.

\begin{env}[7.2.8]
\label{I.7.2.8-ko}
$X$, $Y$를 두 $S$-준스킴이라 하고, $X$는 축소 준스킴이며 $Y$는 $S$
위에서 분리되어 있다고 하자. $f$를 $X$에서 $Y$로 가는 $S$-유리
사상, $x$를 $X$의 점이라 하자. $f$를 표준 $S$-사상
$\operatorname{Spec}(\mathscr{O}_x)\to X$ 뒤에 합성할 수 있다
(\hyperref[I.2.4.1-ko]{2.4.1}). 다만 $\operatorname{Spec}(\mathscr{O}_x)$과
$f$의 정의역의 교집합이 이 국소 스펙트럼에서 조밀해야 하며, 여기서
$\operatorname{Spec}(\mathscr{O}_x)$는 $x\in\overline{\{z\}}$인 $z\in X$들의 집합과 동일시된다
(\hyperref[I.2.4.2-ko]{2.4.2}). 이 조밀성은 다음 두 경우에 성립한다.
\begin{enumerate}
  \item[1\textsuperscript{o}] $X$가 기약이다. 이때 앞의 축소 가정과
    함께 정역 준스킴이다. 실제로 $X$의 일반점 $\xi$는
    $\operatorname{Spec}(\mathscr{O}_x)$의 일반점이기도 하다. $f$의
    정의역 $U$가 $\xi$를 포함하므로
    $U\cap\operatorname{Spec}(\mathscr{O}_x)$도 $\xi$를 포함하고,
    따라서 $\operatorname{Spec}(\mathscr{O}_x)$에서 조밀하다.
  \item[2\textsuperscript{o}] $X$가 국소 뇌터 준스킴이다. 이 경우의
    주장은 다음 보조정리에서 따른다.
\end{enumerate}
\end{env}

\begin{lemma}[7.2.8.1]
\label{I.7.2.8.1-ko}
$X$를 그 바탕공간이 국소 뇌터인 준스킴, $x$를 $X$의 점이라 하자.
$\operatorname{Spec}(\mathscr{O}_x)$의 기약 성분들은 $x$를 포함하는
$X$의 기약 성분들과 $\operatorname{Spec}(\mathscr{O}_x)$의
교집합들이다. 열린집합 $U\subset X$에 대하여
$U\cap\operatorname{Spec}(\mathscr{O}_x)$가
$\operatorname{Spec}(\mathscr{O}_x)$에서 조밀하기 위한 필요충분조건은
이 열린집합이 $X$의 기약 성분들 가운데 $x$를 포함하는 모든 성분과
만나는 것이다. 특히
$U$가 $X$에서 조밀하면 이 조건이 성립한다.
\end{lemma}

둘째 주장은 첫째 주장에서 자명하게 따르므로 첫째 주장만 증명하면
된다. $\operatorname{Spec}(\mathscr{O}_x)$는 $x$를 포함하는 모든 아핀
열린집합 $U$에 포함되고, $x$를 포함하는 $U$의 기약 성분들은 $x$를
포함하는 $X$의 기약 성분들과 $U$의 교집합들이므로
(\hyperref[0.2.1.6-ko]{0, 2.1.6}), $X$가 환 $A$의 아핀 스킴인 경우만
생각하면 충분하다. $A_x$의 소 아이디얼들은
\oldpage[I]{161}
$A$의 소 아이디얼들 중 $\mathfrak{j}_x$에 포함되는 것들과 전단사로 대응하므로
(\hyperref[0.1.2.6-ko]{0, 1.2.6}), $A_x$의 극소 소 아이디얼들은
$\mathfrak{j}_x$에 포함되는 $A$의 극소 소 아이디얼들과 대응한다.
이로써 보조정리가 따른다(\hyperref[I.1.1.14-ko]{1.1.14}).

이제 앞의 두 경우 가운데 하나가 성립한다고 하자. $U$가 $S$-유리
사상 $f$의 정의역이면, $U\cap\operatorname{Spec}(\mathscr{O}_x)$에서
$f$와 일치하는(\hyperref[I.2.4.2-ko]{2.4.2})
$\operatorname{Spec}(\mathscr{O}_x)$에서 $Y$로 가는 유리 사상을
$f'$이라 하겠다. 이 유리 사상을 \emph{$f$가 유도한} 유리 사상이라
한다.

\begin{proposition}[7.2.9]
\label{I.7.2.9-ko}
$S$를 국소 뇌터 준스킴, $X$를 축소 $S$-준스킴, $Y$를 $S$ 위의
유한형 스킴이라 하자. 나아가 $X$가 기약이거나 국소 뇌터라고 하자.
$f$를 $X$에서 $Y$로 가는 $S$-유리 사상, $x$를 $X$의 점이라 하자.
$f$가 점 $x$에서 정의되기 위한 필요충분조건은 $f$가 유도한
(\hyperref[I.7.2.8-ko]{7.2.8})
$\operatorname{Spec}(\mathscr{O}_x)$에서 $Y$로 가는 유리 사상 $f'$이
사상인 것이다.
\end{proposition}

$\operatorname{Spec}(\mathscr{O}_x)$는 $x$를 포함하는 모든
열린집합에 포함되므로 이 조건이 필요함은 자명하다. 충분함을 증명하자.
(\hyperref[I.6.5.1-ko]{6.5.1})에 따라, $X$에서 $x$의 열린 근방 $U$와
$S$-사상 $g$가 존재하며, 이 사상은 $U$에서 $Y$로 가고
$\operatorname{Spec}(\mathscr{O}_x)$ 위에서 $f'$을 유도한다. $X$가 기약이면 $U$는 $X$에서
조밀하므로, (\hyperref[I.7.2.3-ko]{7.2.3})에 따라 $g$가 정하는
$S$-유리 사상으로 이를 보겠다. 또한 $X$의 일반점은
$\operatorname{Spec}(\mathscr{O}_x)$와 $f$의 정의역에 모두 속한다.
따라서 $f$와 $g$는 이 점에서 일치하고, 그러므로 $X$의 공집합이 아닌
어떤 열린집합에서 일치한다(\hyperref[I.6.5.1-ko]{6.5.1}). 그런데
$f$와 $g$는 $S$-유리 사상이므로 서로 같다
(\hyperref[I.7.2.3-ko]{7.2.3}). 따라서 $f$는 $x$에서 정의된다.

이제 $X$가 국소 뇌터라고 하자. $U$가 뇌터라고 가정할 수 있다.
그러면 $x$를 포함하는 $X$의 기약 성분 $X_i$는 유한 개뿐이고
(\hyperref[I.7.2.8.1-ko]{7.2.8.1}), 필요하면 $U$를 더 작은
열린집합으로 바꾸어 이 성분들만 $U$와 만나게 할 수 있다. 실제로
$U$가 뇌터이므로 $U$와 만나는 $X$의 기약 성분은 유한 개뿐이다.
앞에서와 같이 $f$와 $g$는 각 $X_i$의 공집합이 아닌 열린집합에서
일치한다. 각 $X_i$가 $\overline U$에 포함됨을 고려하여,
$U\cup(X-\overline U)$의 조밀한 열린 부분집합 위에서 정의되고,
$U$에서는 $g$와 같으며, $X-\overline U$와 $f$의 정의역의 교집합에서는
$f$와 같은 사상 $f_1$을 생각하자. $U\cup(X-\overline U)$가 $X$에서
조밀하므로 $f_1$과 $f$는 $X$의 조밀한 열린집합에서 일치한다.
$f$가 유리 사상이므로, $f$는 $f_1$의 연장이다
(\hyperref[I.7.2.3-ko]{7.2.3}). 따라서 원래 유리 사상은 점 $x$에서 정의된다.

\subsection{유리함수층}
\label{subsection:I.7.3-ko}

\begin{env}[7.3.1]
\label{I.7.3.1-ko}
$X$를 준스킴이라 하자. 모든 열린집합 $U\subset X$에 대하여 $R(U)$를
$U$ 위의 유리함수환이라 하자(\hyperref[I.7.1.3-ko]{7.1.3}). 이는
$\Gamma(U,\mathscr{O}_X)$-대수이다. 또한 $V\subset U$가 $X$의 또 다른
열린집합이면, $U$의 모든 곳에서 조밀한 열린 부분집합 위의
$\mathscr{O}_X$의 단면은 $V$로 제한하여 $V$의 모든 곳에서 조밀한
열린 부분집합 위의 단면을 주고, 두 단면이 $U$의 모든 곳에서 조밀한
열린 부분집합 위에서 일치하면 그들의 $V$로의 제한은 $V$의 모든 곳에서 조밀한
열린 부분집합 위에서 일치한다. 따라서 대수의 디-준동형
$\rho_{V,U}:R(U)\to R(V)$가 정의된다. $U\supset V\supset W$가 $X$의
세 열린집합이면
$\rho_{W,U}=\rho_{W,V}\circ\rho_{V,U}$임이 자명하다. 그러므로
$R(U)$들은 $X$ 위의 대수의 \emph{준층}을 정의한다.
\end{env}

\oldpage[I]{162}
\begin{definition}[7.3.2]
\label{I.7.3.2-ko}
준스킴 $X$ 위의 \emph{유리함수층}이라 함은 $R(U)$들로 이루어진
준층의 층화인 $\mathscr{O}_X$-대수를 말하며, 이를
$\mathscr{R}(X)$로 나타낸다.
\end{definition}

모든 준스킴 $X$와 모든 열린집합 $U\subset X$에 대하여 유도된 층
$\mathscr{R}(X)|U$는 바로 $\mathscr{R}(U)$이다.

\begin{proposition}[7.3.3]
\label{I.7.3.3-ko}
$X$를 그 기약 성분들의 족 $(X_\lambda)$가 국소 유한인 준스킴이라
하자. 특히 $X$의 바탕공간이 국소 뇌터이면 이 조건이 성립한다.
그러면 $\mathscr{O}_X$-가군 $\mathscr{R}(X)$는 준연접이고,
유한 개의 성분 $X_\lambda$만 만나는 $X$의 모든 열린집합 $U$에
대하여 $R(U)$는 $\Gamma(U,\mathscr{R}(X))$와 같으며, 이 환은
$U\cap X_\lambda\neq\varnothing$인 $X_\lambda$들의 일반점의
국소환들의 직접곱과 표준적으로 동일시된다.
\end{proposition}

$X$가 유한 개의 기약 성분 $X_i$만 갖고 그 일반점을 $x_i$
($1\leq i\leq n$)라 하는 경우만 생각하면 충분하다. 이때 $R(U)$가
$U\cap X_i\neq\varnothing$인 $\mathscr{O}_{x_i}=R(X_i)$들의
직접곱과 표준적으로 동일시된다는 것은
(\hyperref[I.7.1.7-ko]{7.1.7})에서 따른다. 또한 준층
$U\to R(U)$가 층의 공리를 만족함을 보이자. 그러면
$R(U)=\Gamma(U,\mathscr{R}(X))$가 증명된다. 실제로 앞의 결과에 따라
(F 1)이 성립한다. (F 2)를 보이기 위하여 $X$의 열린집합 $U$의
열린집합 $V_\alpha\subset U$들로 이루어진 열린 덮개를 생각하자. $s_\alpha\in R(V_\alpha)$들이
주어지고 모든 지표쌍에 대하여 $s_\alpha$와 $s_\beta$의
$V_\alpha\cap V_\beta$로의 제한이 일치한다고 하자. 그러면
$U\cap X_i\neq\varnothing$인 모든 지표 $i$에 대하여
$V_\alpha\cap X_i\neq\varnothing$인 모든 $s_\alpha$의
$R(X_i)$-성분은 같다. 이 성분을 $t_i$라 하면, $t_i$들을 성분으로
갖는 $R(U)$의 원소는 각 $V_\alpha$로 제한하면 $s_\alpha$가 된다.
마지막으로 $\mathscr{R}(X)$가 준연접임을 보이기 위해
$X=\operatorname{Spec}(A)$가 아핀인 경우만 생각하면 된다.
$f\in A$인 $D(f)$ 꼴의 아핀 열린집합을 $U$로 취하면 앞의 결과와
정의 (\hyperref[I.1.3.4-ko]{1.3.4})에 따라
$\mathscr{R}(X)=\widetilde M$이다. 여기서 $M$은 $A$-가군
$A_{x_i}$들의 직합이다.

\begin{corollary}[7.3.4]
\label{I.7.3.4-ko}
$X$를 유한 개의 기약 성분만 갖는 축소 준스킴이라 하고,
$X_i$ ($1\leq i\leq n$)를 $X$의 기약 성분을 바탕공간으로 갖는
$X$의 축소 닫힌 부분준스킴이라 하자
(\hyperref[I.5.2.1-ko]{5.2.1}). $h_i$가 $X_i\to X$인 표준 포함
사상이면 $\mathscr{R}(X)$는 $\mathscr{O}_X$-대수
$(h_i)_*(\mathscr{R}(X_i))$들의 직접곱이다.
\end{corollary}

\begin{corollary}[7.3.5]
\label{I.7.3.5-ko}
$X$가 기약이면 모든 준연접 $\mathscr{R}(X)$-가군
$\mathscr{F}$는 상수층이다.
\end{corollary}

모든 $x\in X$가 $\mathscr{F}|U$가 상수층이 되는 근방 $U$를
가짐을 보이면 충분하다(\hyperref[0.3.6.2-ko]{0, 3.6.2}).
다시 말해 $X$가 아핀인 경우로 환원된다. 더 나아가
$\mathscr{F}$가 어떤 준동형
$(\mathscr{R}(X))^{(I)}\to(\mathscr{R}(X))^{(J)}$의 여핵이라고
가정할 수 있다(\hyperref[0.5.1.3-ko]{0, 5.1.3}). 따라서
$\mathscr{R}(X)$가 상수층임만 보이면 된다. 그러나 $X$의 공집합이
아닌 임의의 열린집합을 $U$라 하면 그 $U$는 일반점을 포함하므로
$\Gamma(U,\mathscr{R}(X))=R(X)$이고, 이는 자명하다.

\begin{corollary}[7.3.6]
\label{I.7.3.6-ko}
$X$가 기약이면 모든 준연접 $\mathscr{O}_X$-가군
$\mathscr{F}$에 대하여
$\mathscr{F}\otimes_{\mathscr{O}_X}\mathscr{R}(X)$는 상수층이다.
더 나아가 $X$가 축소 준스킴이면(따라서 정역 준스킴이다),
$\mathscr{F}\otimes_{\mathscr{O}_X}\mathscr{R}(X)$는
$(\mathscr{R}(X))^{(I)}$ 꼴의 층과 동형이다.
\end{corollary}

두 번째 주장은 이 경우 $R(X)$가 체라는 사실에서 따른다.

\begin{proposition}[7.3.7]
\label{I.7.3.7-ko}
준스킴 $X$가 국소 정역 준스킴이거나 국소
\oldpage[I]{163}
뇌터 준스킴이라고 하자. 그러면 $\mathscr{R}(X)$는 준연접
$\mathscr{O}_X$-대수이다. 더 나아가 $X$가 축소 준스킴이면
(특히 $X$가 국소 정역 준스킴이면 그러하다), 표준 준동형
$\mathscr{O}_X\to\mathscr{R}(X)$는 단사이다.
\end{proposition}

이 문제는 국소적이므로 첫 번째 주장은
(\hyperref[I.7.3.3-ko]{7.3.3})에서 따르고, 두 번째 주장은
(\hyperref[I.7.2.3-ko]{7.2.3})에서 바로 따른다.

\begin{env}[7.3.8]
\label{I.7.3.8-ko}%
\footnote{\emph{[번역 주] 인쇄된 이 문단은 EGA II의 정오표에서
오류로 판정되어 전면 대체되었다. 여기서는 그 공식 대체문을 옮긴다.}}
$X$와 $Y$를 두 \emph{정역} 준스킴이라 하자. 그러면
$\mathscr{R}(X)$는 준연접 $\mathscr{O}_X$-가군이고,
$\mathscr{R}(Y)$는 준연접 $\mathscr{O}_Y$-가군이다
(\hyperref[I.7.3.3-ko]{7.3.3}). $f:X\to Y$를 \emph{우세} 사상이라
하자. 그러면 표준 $\mathscr{O}_X$-가군 준동형
\[
  \tau:f^*(\mathscr{R}(Y))\longrightarrow\mathscr{R}(X)
  \tag{7.3.8.1}\label{I.7.3.8.1-ko}
\]
이 존재한다.
\end{env}

먼저 $X=\operatorname{Spec}(A)$와 $Y=\operatorname{Spec}(B)$가
정역 $A$, $B$로 주어진 아핀 준스킴이라고 하자. 그러면 $f$는
단사 준동형 $B\to A$에 대응하고, 이 준동형은 $B$의 분수체
$L$에서 $A$의 분수체 $K$로 가는 단사 준동형 $L\to K$로
연장된다. 이때 (7.3.8.1)의 준동형은 표준 준동형
$L\otimes_B A\to K$에 대응한다
(\hyperref[I.1.6.5-ko]{1.6.5}).

일반적인 경우에는 $f(U)\subset V$인 공집합이 아닌 아핀
열린집합 $U\subset X$, $V\subset Y$의 각 쌍에 대하여 앞과 같이
준동형 $\tau_{U,V}$를 정의한다. $U'\subset U$, $V'\subset V$이고
$f(U')\subset V'$이면 $\tau_{U,V}$는 $\tau_{U',V'}$의 연장이다.
따라서 이 준동형들을 붙이면 앞의 표준
준동형을 얻는다. $x$, $y$가 각각 $X$, $Y$의 일반점이면
$f(x)=y$이고,
\[
  (f^*(\mathscr{R}(Y)))_x
  =\mathscr{O}_y\otimes_{\mathscr{O}_y}\mathscr{O}_x
  =\mathscr{O}_x
\]
이다(\hyperref[0.4.3.1-ko]{0, 4.3.1}). 따라서 $\tau_x$는
\emph{동형}이다.

\subsection{꼬임층과 꼬임 없는 층}
\label{subsection:I.7.4-ko}

\begin{env}[7.4.1]
\label{I.7.4.1-ko}
$X$를 \emph{정역} 준스킴이라 하자. 임의의 $\mathscr{O}_X$-가군
$\mathscr{F}$에 대하여 표준 준동형
$\mathscr{O}_X\to\mathscr{R}(X)$은 텐서곱을 취함으로써 준동형
(역시 \emph{표준}이라 한다)
\[
  \mathscr{F}\to
  \mathscr{F}\otimes_{\mathscr{O}_X}\mathscr{R}(X)
\]
을 정의한다. 각 줄기에서 이는 $\mathscr{F}_x$에서
$\mathscr{F}_x\otimes_{\mathscr{O}_x}R(X)$로 가는 준동형
$z\to z\otimes 1$이다. 이 준동형의 \emph{핵 $\mathscr{T}$}은
$\mathscr{F}$의 $\mathscr{O}_X$-부분가군이며, 이를
\emph{$\mathscr{F}$의 꼬임 부분층}이라 한다. $\mathscr{F}$가
준연접이면 이 부분층도 준연접이다
(\hyperref[I.4.1.1-ko]{4.1.1}과
\hyperref[I.7.3.6-ko]{7.3.6}). $\mathscr{T}=0$이면
$\mathscr{F}$에 \emph{꼬임이 없다}고 하고,
$\mathscr{T}=\mathscr{F}$이면 $\mathscr{F}$를 \emph{꼬임층}이라 한다.
임의의 $\mathscr{O}_X$-가군 $\mathscr{F}$에 대하여
$\mathscr{F}/\mathscr{T}$에는 꼬임이 없다.
(\hyperref[I.7.3.5-ko]{7.3.5})에서 다음을 얻는다.
\end{env}

\begin{proposition}[7.4.2]
\label{I.7.4.2-ko}
$X$가 정역 준스킴이면, 꼬임이 없는 모든 준연접
$\mathscr{O}_X$-가군 $\mathscr{F}$는 $(\mathscr{R}(X))^{(I)}$ 꼴의
상수층의 부분층 $\mathscr{G}$와 동형이다. 이 상수층은
$\mathscr{G}$에 의해 $\mathscr{R}(X)$-가군으로 생성된다.
\end{proposition}

$I$의 기수를 \emph{$\mathscr{F}$의 계수}라 한다. $X$의
공집합이 아닌 임의의 아핀 열린집합 $U$에 대하여 $\mathscr{F}$의
계수는 $\Gamma(U,\mathscr{F})$의 $\Gamma(U,\mathscr{O}_X)$-가군으로서의
계수와 같다. 이는 $U$에 들어 있는 $X$의 일반점을 고려하면 곧 알 수
있다. 특히 다음을 얻는다.

\begin{corollary}[7.4.3]
\label{I.7.4.3-ko}
정역 준스킴 $X$ 위에서 꼬임이 없고 계수가 $1$인 모든 준연접
$\mathscr{O}_X$-가군(특히 모든 가역 $\mathscr{O}_X$-가군)은
$\mathscr{R}(X)$의 $\mathscr{O}_X$-부분가군과 동형이며, 그 역도
성립한다.
\end{corollary}

\begin{corollary}[7.4.4]
\label{I.7.4.4-ko}
$X$를 정역 준스킴, $\mathscr{L}$과 $\mathscr{L}'$을 꼬임이 없는 두
$\mathscr{O}_X$-가군, $f$(각각 $f'$)를 $X$ 위의
$\mathscr{L}$(각각 $\mathscr{L}'$)의 단면이라 하자.
$f\otimes f'=0$일 필요충분조건은 두 단면 $f$, $f'$ 가운데 하나가
영인 것이다.
\end{corollary}

$x$를 $X$의 일반점이라 하자. 가정에 따라
$(f\otimes f')_x=f_x\otimes f'_x=0$이다. $\mathscr{L}_x$와
$\mathscr{L}'_x$는 체 $\mathscr{O}_x$의 $\mathscr{O}_x$-부분가군으로
동일시되므로, 앞의 관계에서 $f_x=0$ 또는 $f'_x=0$을 얻는다.
$\mathscr{L}$과 $\mathscr{L}'$에는 꼬임이 없으므로
(\hyperref[I.7.3.5-ko]{7.3.5}) 결국 $f=0$ 또는 $f'=0$이다.

\begin{proposition}[7.4.5]
\label{I.7.4.5-ko}
$X$, $Y$를 두 정역 준스킴, $f:X\to Y$를 우세 사상이라 하자.
꼬임이 없는 모든 준연접 $\mathscr{O}_X$-가군 $\mathscr{F}$에 대하여
$f_*(\mathscr{F})$는 꼬임이 없는 $\mathscr{O}_Y$-가군이다.
\end{proposition}

\oldpage[I]{164}
\begin{proof}
$f_*$는 왼쪽 완전하므로
(\hyperref[0.4.2.1-ko]{0, 4.2.1}),
(\hyperref[I.7.4.2-ko]{7.4.2})에 따라
$\mathscr{F}=(\mathscr{R}(X))^{(I)}$인 경우에 명제를 증명하면 충분하다.
그런데 $Y$의 공집합이 아닌 모든 열린집합 $U$는 $Y$의 일반점을
포함하므로, $f^{-1}(U)$는 $X$의 일반점을 포함한다
(\hyperref[0.2.1.5-ko]{0, 2.1.5}). 따라서
\[
  \Gamma(U,f_*(\mathscr{F}))
  =\Gamma(f^{-1}(U),\mathscr{F})=(R(X))^{(I)}.
\]
다시 말해 $f_*(\mathscr{F})$는 $(R(X))^{(I)}$를 값으로 하는
상수층이다. 이 층을 $\mathscr{R}(Y)$-가군으로 보면 명백히 꼬임이
없다.
\end{proof}

\begin{proposition}[7.4.6]
\label{I.7.4.6-ko}
$X$를 정역 준스킴, $x$를 그 일반점이라 하자. 유한 생성인 임의의
준연접 $\mathscr{O}_X$-가군 $\mathscr{F}$에 대하여 다음 조건들은
동치이다. a) $\mathscr{F}$는 꼬임층이다. b) $\mathscr{F}_x=0$이다.
c) $\operatorname{Supp}(\mathscr{F})\neq X$이다.
\end{proposition}

(\hyperref[I.7.3.5-ko]{7.3.5})와
(\hyperref[I.7.4.1-ko]{7.4.1})에 따라 관계 $\mathscr{F}_x=0$과
$\mathscr{F}\otimes_{\mathscr{O}_X}\mathscr{R}(X)=0$은 동치이므로
a)와 b)는 동치이다. 한편 $\operatorname{Supp}(\mathscr{F})$는
$X$에서 닫혀 있고(\hyperref[0.5.2.2-ko]{0, 5.2.2}), $X$의 공집합이
아닌 모든 열린집합은 $x$를 포함하므로 b)와 c)는 동치이다.

\begin{env}[7.4.7]
\label{I.7.4.7-ko}
$X$가 \emph{축소} 준스킴이고 기약성분을 \emph{유한}개만 갖는 경우에도
용어를 남용하여 (\hyperref[I.7.4.1-ko]{7.4.1})의 정의들을 사용한다.
그러면 (\hyperref[I.7.3.4-ko]{7.3.4})에 따라 이러한 준스킴에 대해서도
(\hyperref[I.7.4.6-ko]{7.4.6})의 a)와 c)는 동치이다.
\end{env}
% ===== END EXACT INLINE INPUT: c1s7.tex =====


% ===== BEGIN EXACT INLINE INPUT: c1s8.tex =====
\section{슈발레 스킴}
\label{section:I.8-ko}

\subsection{서로 연관된 국소환}
\label{subsection:I.8.1-ko}

임의의 국소환 $A$에 대하여 $A$의 극대 아이디얼을
$\mathfrak m(A)$로 나타낸다.

\begin{lemma}[8.1.1]
\label{I.8.1.1-ko}
$A$, $B$를 $A\subset B$인 두 국소환이라 하자. 다음 조건들은
동치이다.
\begin{enumerate}
  \item[(i)] $\mathfrak m(B)\cap A=\mathfrak m(A)$이다.
  \item[(ii)] $\mathfrak m(A)\subset\mathfrak m(B)$이다.
  \item[(iii)] $1$은 $\mathfrak m(A)$에 의해 생성되는 $B$의
    아이디얼에 속하지 않는다.
\end{enumerate}
\end{lemma}

(i)가 (ii)를 함의하고 (ii)가 (iii)을 함의함은 명백하다. 끝으로
(iii)이 성립하면 $\mathfrak m(B)\cap A$는 $\mathfrak m(A)$를 포함하고
$1$을 포함하지 않으므로 $\mathfrak m(A)$와 같다.

(\hyperref[I.8.1.1-ko]{8.1.1})의 동치인 조건들이 성립할 때 $B$가
$A$를 \emph{지배한다}고 한다. 이는 포함 준동형 $A\to B$가
\emph{국소} 준동형이라고 하는 것과 같다. 환 $R$의 국소 부분환들의
집합에서 지배관계가 순서관계임은 명백하다.

\begin{env}[8.1.2]
\label{I.8.1.2-ko}
이제 \emph{체} $R$을 생각하자. $R$의 임의의 부분환 $A$에 대하여
$\mathfrak p$가 $A$의 소 스펙트럼을 달릴 때의 국소환
$A_\mathfrak p$들의 집합을 $L(A)$로 나타낸다. 이 환들은 $A$를
포함하는 $R$의 부분환으로 동일시된다.
$\mathfrak p=(\mathfrak p A_\mathfrak p)\cap A$이므로,
$\operatorname{Spec}(A)$에서 $L(A)$로 가는 사상
$\mathfrak p\mapsto A_\mathfrak p$는 전단사이다.
\end{env}

\begin{lemma}[8.1.3]
\label{I.8.1.3-ko}
$R$을 체, $A$를 $R$의 부분환이라 하자. $R$의 국소 부분환 $M$이
환 $A_\mathfrak p\in L(A)$를 지배할 필요충분조건은
$A\subset M$인 것이다. 이때 $M$이 지배하는 국소환
$A_\mathfrak p$는 유일하며,
$\mathfrak p=\mathfrak m(M)\cap A$에 대응한다.
\end{lemma}

실제로 $M$이 $A_\mathfrak p$를 지배하면
(\hyperref[I.8.1.1-ko]{8.1.1})에 따라
$\mathfrak m(M)\cap A_\mathfrak p=\mathfrak p A_\mathfrak p$이므로
$\mathfrak p$는 유일하다. 다른 한편 $A\subset M$이면
$\mathfrak m(M)\cap A=\mathfrak p$는 $A$의 소 아이디얼이고,
$A-\mathfrak p\subset M$이므로 $A_\mathfrak p\subset M$이고
$\mathfrak p A_\mathfrak p\subset\mathfrak m(M)$이다. 따라서 $M$은
$A_\mathfrak p$를 지배한다.

\oldpage[I]{165}
\begin{lemma}[8.1.4]
\label{I.8.1.4-ko}
$R$을 체, $M$, $N$을 $R$의 두 국소 부분환, $P$를 $M\cup N$에 의해
생성되는 $R$의 부분환이라 하자. 다음 조건들은 동치이다.
\begin{enumerate}
  \item[(i)] $P$의 소 아이디얼 $\mathfrak p$가 존재하여
    $\mathfrak m(M)=\mathfrak p\cap M$,
    $\mathfrak m(N)=\mathfrak p\cap N$이다.
  \item[(ii)] $\mathfrak m(M)\cup\mathfrak m(N)$에 의해 $P$에서
    생성되는 아이디얼 $\mathfrak a$는 $P$와 다르다.
  \item[(iii)] $M$과 $N$을 모두 지배하는 $R$의 국소 부분환 $Q$가
    존재한다.
\end{enumerate}
\end{lemma}

(i)가 (ii)를 함의함은 명백하다. 역으로 $\mathfrak a\ne P$이면
$\mathfrak a$는 $P$의 어떤 극대 아이디얼 $\mathfrak n$에 포함된다.
$1\notin\mathfrak n$이므로 $\mathfrak n\cap M$은 $\mathfrak m(M)$을
포함하면서 $M$과 다르며, 따라서
$\mathfrak n\cap M=\mathfrak m(M)$이다. 마찬가지로
$\mathfrak n\cap N=\mathfrak m(N)$이다. $Q$가 $M$과 $N$을 지배하면
$P\subset Q$이고
\[
  \mathfrak m(M)=\mathfrak m(Q)\cap M
  =(\mathfrak m(Q)\cap P)\cap M,\qquad
  \mathfrak m(N)=(\mathfrak m(Q)\cap P)\cap N,
\]
이므로 (iii)은 (i)를 함의한다. 역은 $Q=P_\mathfrak p$로 두면
명백하다.

(\hyperref[I.8.1.4-ko]{8.1.4})의 조건들이 성립할 때 C.~슈발레를 따라
국소환 $M$과 $N$이 \emph{서로 연관되어 있다}고 한다.

\begin{proposition}[8.1.5]
\label{I.8.1.5-ko}
$A$, $B$를 체 $R$의 두 부분환, $C$를 $A\cup B$에 의해 생성되는
$R$의 부분환이라 하자. 다음 조건들은 동치이다.
\begin{enumerate}
  \item[(i)] $A$와 $B$를 포함하는 모든 국소환 $Q$에 대하여
    $\mathfrak p=\mathfrak m(Q)\cap A$,
    $\mathfrak q=\mathfrak m(Q)\cap B$로 놓으면
    $A_\mathfrak p=B_\mathfrak q$이다.
  \item[(ii)] $C$의 모든 소 아이디얼 $\mathfrak r$에 대하여
    $\mathfrak p=\mathfrak r\cap A$, $\mathfrak q=\mathfrak r\cap B$로
    놓으면 $A_\mathfrak p=B_\mathfrak q$이다.
  \item[(iii)] $M\in L(A)$와 $N\in L(B)$가 서로 연관되어 있으면
    두 환은 같다.
  \item[(iv)] $L(A)\cap L(B)=L(C)$이다.
\end{enumerate}
\end{proposition}

보조정리 (\hyperref[I.8.1.3-ko]{8.1.3})과
(\hyperref[I.8.1.4-ko]{8.1.4})는 (i)와 (iii)이 동치임을 보인다.
(i)를 $Q=C_\mathfrak r$에 적용하면 (i)가 (ii)를 함의함은 명백하다.
역으로 $Q$가 $A\cup B$를 포함하면 $C$도 포함한다. 이때
$\mathfrak r=\mathfrak m(Q)\cap C$로 놓으면
(\hyperref[I.8.1.3-ko]{8.1.3})에 따라
$\mathfrak p=\mathfrak r\cap A$이고 $\mathfrak q=\mathfrak r\cap B$이므로
(ii)는 (i)를 함의한다. (iv)가 (i)를 함의함도 곧 알 수 있다. 실제로
$Q$가 $A\cup B$를 포함하면 (\hyperref[I.8.1.3-ko]{8.1.3})에 따라
어떤 국소환 $C_\mathfrak r\in L(C)$를 지배한다. 가정에 따라
$C_\mathfrak r\in L(A)\cap L(B)$이고,
(\hyperref[I.8.1.1-ko]{8.1.1})과
(\hyperref[I.8.1.3-ko]{8.1.3})에 따라
$C_\mathfrak r=A_\mathfrak p=B_\mathfrak q$이다.

끝으로 (iii)이 (iv)를 함의함을 보이자. $Q\in L(C)$라 하자.
(\hyperref[I.8.1.3-ko]{8.1.3})에 따라 $Q$는 어떤 $M\in L(A)$와
$N\in L(B)$를 지배한다. 따라서 서로 연관된 $M$과 $N$은 가정에 따라
같다. 이제 $C\subset M$이므로 (\hyperref[I.8.1.3-ko]{8.1.3})에 따라
$M$은 어떤 $Q'\in L(C)$를 지배한다. 따라서 $Q$는 $Q'$을 지배한다.
그러므로 (\hyperref[I.8.1.3-ko]{8.1.3})에 따라 반드시
$Q=Q'=M$이고, $Q\in L(A)\cap L(B)$이다. 역으로
$Q\in L(A)\cap L(B)$이면 $C\subset Q$이므로
(\hyperref[I.8.1.3-ko]{8.1.3})에 따라 $Q$는 어떤
$Q''\in L(C)\subset L(A)\cap L(B)$를 지배한다. 서로 연관된 $Q$와
$Q''$은 같으므로 $Q''=Q\in L(C)$이다. 이로써 증명이 끝난다.

\subsection{정역 스킴의 국소환}
\label{subsection:I.8.2-ko}

\begin{env}[8.2.1]
\label{I.8.2.1-ko}
$X$를 \emph{정역} 준스킴이라 하자. 그 유리함수체 $R$은 $X$의 일반점
$a$의 국소환과 같다. 모든 $x\in X$에 대하여
$\mathscr{O}_x$는 $R$의 부분환과 표준적으로 동일시된다
(\hyperref[I.7.1.5-ko]{7.1.5}). 또한 임의의 유리함수 $f\in R$에
대하여 $f$의 정의역 $\delta(f)$는 $f\in\mathscr{O}_x$인 점
$x\in X$들의 열린집합이다. 따라서
(\hyperref[I.7.2.6-ko]{7.2.6}) 모든 열린집합 $U\subset X$에 대하여
\begin{equation}
  \Gamma(U,\mathscr{O}_X)=\bigcap_{x\in U}\mathscr{O}_x.
  \tag{8.2.1.1}\label{I.8.2.1.1-ko}
\end{equation}
\end{env}

\oldpage[I]{166}
\begin{proposition}[8.2.2]
\label{I.8.2.2-ko}
$X$를 정역 준스킴, $R$을 그 유리함수체라 하자. $X$가 스킴이기 위한
필요충분조건은 $X$의 두 점 $x$, $y$ 사이의 관계
``$\mathscr{O}_x$와 $\mathscr{O}_y$가 서로 연관되어 있다''
(\hyperref[I.8.1.4-ko]{8.1.4})가 $x=y$를 함의하는 것이다.
\end{proposition}

이 조건이 성립한다고 가정하고 $X$가 분리임을 보이자. $X$의 서로 다른
두 아핀 열린집합을 $U$, $V$라 하고, 그 환을 각각 $A$, $B$라 하여
이들을 $R$의 부분환과 동일시하자. 그러면 $U$와 $V$는 각각
(\hyperref[I.8.1.2-ko]{8.1.2})에 따라 $L(A)$와 $L(B)$로 동일시된다.
가정과 (\hyperref[I.8.1.5-ko]{8.1.5})에 따라 $C$가 $A\cup B$에 의해
$R$에서 생성되는 부분환이면 $W=U\cap V$는
$L(A)\cap L(B)=L(C)$와 동일시된다. 한편 $R$의 모든 부분환 $E$는
$L(E)$에 속하는 국소환들의 교집합과 같음이 알려져 있다
(\cite{I-1}, p.~5-03, 명제~4 \emph{bis}). 따라서 $C$는 $z\in W$에
대한 국소환 $\mathscr{O}_z$들의 교집합, 즉
(\hyperref[I.8.2.1.1-ko]{8.2.1.1})에 따라
$\Gamma(W,\mathscr{O}_X)$와 동일시된다. 이제 $W$ 위에 $X$가 유도하는
부분준스킴을 생각하자. 항등 준동형
$\varphi:C\to\Gamma(W,\mathscr{O}_X)$에 대응하여
(\hyperref[I.2.2.4-ko]{2.2.4}) 사상
$\Phi=(\psi,\theta):W\to\operatorname{Spec}(C)$가 존재한다. $\Phi$가
준스킴의 \emph{동형}임을 보이면 $W$가 \emph{아핀} 열린집합임이
따른다. $W$와 $L(C)=\operatorname{Spec}(C)$의 동일시에 따라
$\psi$는 \emph{전단사}이다. 또한 모든 $x\in W$에 대하여
$\mathfrak r=\mathfrak m_x\cap C$로 놓으면 $\theta_x^\sharp$는 포함
준동형 $C_\mathfrak r\to\mathscr{O}_x$이고, 정의에 따라
$C_\mathfrak r$는 $\mathscr{O}_x$와 동일시되므로
$\theta_x^\sharp$는 전단사이다. 따라서 $\psi$가 \emph{위상동형}, 즉
모든 닫힌 부분집합 $F\subset W$에 대하여 $\psi(F)$가
$\operatorname{Spec}(C)$에서 닫혔음을 보이는 일만 남는다. $F$는
$A$의 어떤 아이디얼 $\mathfrak a$에 대하여 $V(\mathfrak a)$ 꼴인
$U$의 닫힌 부분집합이 $W$ 위에 남기는 자취이다. 이제
$\psi(F)=V(\mathfrak a C)$임을 보이자. 그러면 주장이 따른다. 실제로
$\mathfrak a C$를 포함하는 $C$의 소 아이디얼들은 $\mathfrak a$를
포함하는 $C$의 소 아이디얼들, 즉 $x\in W$이고
$\mathfrak a\subset\mathfrak m_x$인
$\psi(x)=\mathfrak m_x\cap C$ 꼴의 아이디얼들이다. 이러한 점에 대하여
$\mathfrak a\subset\mathfrak m_x$인 것은
$x\in W\cap V(\mathfrak a)=F$인 것과 동치이다.\footnote{\emph{[번역 주]
인쇄 원문은 여기서 불가능한 등식
$x\in V(\mathfrak a)=W\cap F$를 적는다. 앞 문장의 정의에 맞는
등식 $F=W\cap V(\mathfrak a)$을 사용하였다.}} 따라서
$\psi(F)=V(\mathfrak a C)$이다.

$U\cap V$가 아핀이고 그 환 $C$가 $U$와 $V$의 두 환의 합집합
$A\cup B$에 의해 생성되므로 $X$는 분리이다
(\hyperref[I.5.5.6-ko]{5.5.6}).

역으로 $X$가 분리라고 하고, $\mathscr{O}_x$와 $\mathscr{O}_y$가 서로
연관되어 있는 $X$의 두 점을 $x$, $y$라 하자. $x$를 포함하는 아핀
열린집합 $U$의 환을 $A$, $y$를 포함하는 아핀 열린집합 $V$의 환을
$B$라 하자. 그러면 $U\cap V$는 아핀이고 그 환 $C$는 $A\cup B$에
의해 생성된다(\hyperref[I.5.5.6-ko]{5.5.6}).
$\mathfrak p=\mathfrak m_x\cap A$,
$\mathfrak q=\mathfrak m_y\cap B$라 놓으면
$A_\mathfrak p=\mathscr{O}_x$, $B_\mathfrak q=\mathscr{O}_y$이다.
$A_\mathfrak p$와 $B_\mathfrak q$가 서로 연관되어 있으므로
(\hyperref[I.8.1.4-ko]{8.1.4})
$\mathfrak p=\mathfrak r\cap A$,
$\mathfrak q=\mathfrak r\cap B$인 $C$의 소 아이디얼
$\mathfrak r$가 존재한다. $U\cap V$가 아핀이므로
$\mathfrak r=\mathfrak m_z\cap C$인 점 $z\in U\cap V$가 존재한다.
그러면 자명하게 $x=z$이고 $y=z$이므로 $x=y$이다.

\begin{corollary}[8.2.3]
\label{I.8.2.3-ko}
$X$를 정역 스킴, $x$, $y$를 $X$의 두 점이라 하자.
$x\in\overline{\{y\}}$이기 위한 필요충분조건은
$\mathscr{O}_x\subset\mathscr{O}_y$인 것, 즉 $x$에서 정의되는 모든
유리함수가 $y$에서도 정의되는 것이다.
\end{corollary}

유리함수 $f\in R$의 정의역 $\delta(f)$가 열려 있으므로 이 조건이
필요함은 자명하다. 충분함을 보이자. $\mathscr{O}_x\subset
\mathscr{O}_y$이면 (\hyperref[I.8.1.3-ko]{8.1.3})
$\mathscr{O}_y$가 $(\mathscr{O}_x)_\mathfrak p$를 지배하는
$\mathscr{O}_x$의 소 아이디얼 $\mathfrak p$가 존재한다. 그런데
(\hyperref[I.2.4.2-ko]{2.4.2})
$x\in\overline{\{z\}}$이고
$\mathscr{O}_z=(\mathscr{O}_x)_\mathfrak p$인 점 $z\in X$가 존재한다.
$\mathscr{O}_z$와 $\mathscr{O}_y$가 서로 연관되어 있으므로
(\hyperref[I.8.2.2-ko]{8.2.2})에 따라 $z=y$이고, 따름정리가 증명된다.

\begin{corollary}[8.2.4]
\label{I.8.2.4-ko}
$X$가 정역 스킴이면 사상 $x\mapsto\mathscr{O}_x$는 단사이다. 즉
$x$, $y$가 $X$의 서로 다른 두 점이면 그 가운데 한 점에서는 정의되고
다른 점에서는 정의되지 않는 유리함수가 존재한다.
\end{corollary}

\oldpage[I]{167}
이는 (\hyperref[I.8.2.3-ko]{8.2.3})과 공리
$(T_0)$ (\hyperref[I.2.1.4-ko]{2.1.4})에서 따른다.

\begin{corollary}[8.2.5]
\label{I.8.2.5-ko}
$X$를 바탕공간이 뇌터인 정역 스킴이라 하자. $f$가 $X$ 위의
유리함수체 $R$을 달릴 때 집합 $\delta(f)$들은 $X$의 위상을 생성한다.
\end{corollary}

실제로 $X$의 모든 닫힌 부분집합은 유한개의 기약 닫힌 부분집합, 즉
$\overline{\{y\}}$ 꼴인 집합들의 합집합이다
(\hyperref[I.2.1.5-ko]{2.1.5}). 그런데
$x\notin\overline{\{y\}}$이면 $x$에서는 정의되고 $y$에서는 정의되지
않는 유리함수 $f$가 존재한다(\hyperref[I.8.2.3-ko]{8.2.3}). 다시
말해 $x\in\delta(f)$이고 $\delta(f)$는 $\overline{\{y\}}$와 만나지
않는다. 따라서 $\overline{\{y\}}$의 여집합은 $\delta(f)$ 꼴인
집합들의 합집합이다. 첫 관찰에 따라 $X$의 모든 열린집합은
$\delta(f)$ 꼴인 열린집합들의 유한 교집합들의 합집합이다.

\begin{env}[8.2.6]
\label{I.8.2.6-ko}
따름정리 (\hyperref[I.8.2.5-ko]{8.2.5})에 따라 $X$의 위상은 $R$을
분수체로 갖는 국소환들의 족 $(\mathscr{O}_x)_{x\in X}$으로 완전히
특징지어진다. 동치로 $X$의 닫힌 부분집합들은 다음과 같이 정의된다.
$X$의 유한 부분집합 $\{x_1,\ldots,x_n\}$을 주고, 어떤 첨자 $i$에
대하여 $\mathscr{O}_y\subset\mathscr{O}_{x_i}$인 점 $y\in X$들의
집합을 생각하자. $\{x_1,\ldots,x_n\}$의 모든 선택에 대하여 얻는
이 집합들이 바로 $X$의 닫힌 부분집합들이다. 더욱이 $X$의 위상을
알고 나면 구조층 $\mathscr{O}_X$도 $\mathscr{O}_x$들의 족으로
결정된다. 실제로
$\Gamma(U,\mathscr{O}_X)=\bigcap_{x\in U}\mathscr{O}_x$이다
(\hyperref[I.8.2.1.1-ko]{8.2.1.1}). 따라서 $X$가 바탕공간이 뇌터인
정역 스킴이면 족 $(\mathscr{O}_x)_{x\in X}$가 준스킴 $X$를 완전히
결정한다.
\end{env}

\begin{proposition}[8.2.7]
\label{I.8.2.7-ko}
$X$, $Y$를 두 정역 스킴, $f:X\to Y$를 우세 사상
(\hyperref[I.2.2.6-ko]{2.2.6})이라 하자. $X$와 $Y$ 위의 유리함수체를
각각 $K$, $L$이라 하자. 그러면 $L$은 $K$의 부분체와 동일시되고,
모든 $x\in X$에 대하여 $\mathscr{O}_{f(x)}$는 $\mathscr{O}_x$가
지배하는 $Y$의 유일한 국소환이다.
\end{proposition}

실제로 $f=(\psi,\theta)$이고 $a$가 $X$의 일반점이면 $\psi(a)$는
$Y$의 일반점이다(\hyperref[0.2.1.5-ko]{0, 2.1.5}). 따라서
$\theta_a^\sharp$는 체 $L=\mathscr{O}_{\psi(a)}$에서 체
$K=\mathscr{O}_a$로 가는 단사 준동형이다. $Y$의 공집합이 아닌 모든
아핀 열린집합 $U$가 $\psi(a)$를 포함하므로
(\hyperref[I.2.2.4-ko]{2.2.4})에 따라 $f$에 대응하는 준동형
$\Gamma(U,\mathscr{O}_Y)\to
\Gamma(\psi^{-1}(U),\mathscr{O}_X)$는
$\theta_a^\sharp$를 $\Gamma(U,\mathscr{O}_Y)$에 제한한 준동형이다.
따라서 모든 $x\in X$에 대하여 $\theta_x^\sharp$는
$\theta_a^\sharp$를 $\mathscr{O}_{\psi(x)}$에 제한한 것이므로
단사이다. 또한 $\theta_x^\sharp$는 국소 준동형이다. 그러므로
$\theta_a^\sharp$를 통해 $L$을 $K$의 부분체와 동일시하면
$\mathscr{O}_{\psi(x)}$는 $\mathscr{O}_x$가 지배한다
(\hyperref[I.8.1.1-ko]{8.1.1}). 이것은 $\mathscr{O}_x$가 지배하는
$Y$의 유일한 국소환이다. 실제로 $Y$의 서로 연관된 두 국소환은
같기 때문이다(\hyperref[I.8.2.2-ko]{8.2.2}).

\begin{proposition}[8.2.8]
\label{I.8.2.8-ko}
$X$를 기약 준스킴, $f:X\to Y$를 국소 몰입 사상(각각 국소
동형사상)이라 하고, 더 나아가 사상 $f$가 분리라고 하자. 그러면
$f$는 몰입 사상(각각 열린 몰입 사상)이다.
\end{proposition}

$f=(\psi,\theta)$라 하자. 두 경우 모두 $\psi$가 $X$에서
$\psi(X)$로 가는 \emph{위상동형}임을 보이면 충분하다
(\hyperref[I.4.5.3-ko]{4.5.3}). $f$를 $f_{\mathrm{red}}$로 바꾸면
(\hyperref[I.5.1.6-ko]{5.1.6} 및
\hyperref[I.5.5.1-ko]{5.5.1, (vi)})에 따라 $X$와 $Y$가 축소라고 가정할 수
있다. $Y'$을 바탕공간이 $\overline{\psi(X)}$인 $Y$의 축소 닫힌
부분준스킴이라 하자. 그러면 $f$는
$X\xrightarrow{f'}Y'\xrightarrow{j}Y$로 인수분해되며 $j$는 표준 포함
사상이다(\hyperref[I.5.2.2-ko]{5.2.2}).
(\hyperref[I.5.5.1-ko]{5.5.1, (v)})에 따라 $f'$도 여전히 분리
사상이다. 또한 $f'$은 여전히
\oldpage[I]{168}
국소 몰입 사상(각각 국소 동형사상)이다. 실제로 이 문제는 $X$와
$Y$ 위에서 국소적이므로, 이를 보일 때 $f$가 닫힌 몰입 사상(각각
열린 몰입 사상)인 경우만 다루면 되고, 이때 주장은
(\hyperref[I.4.2.2-ko]{4.2.2})에서 곧바로 따른다.

따라서 $f$가 \emph{우세} 사상이라고 가정할 수 있다. 그러면 $Y$도
기약이다(\hyperref[0.2.1.5-ko]{0, 2.1.5}). 따라서 $X$와 $Y$는 모두
\emph{정역}이다. 또한 이 문제는 $Y$ 위에서 국소적이므로 $Y$가 아핀
스킴이라고 가정할 수 있다. $f$가 분리이므로 $X$는 스킴이다
(\hyperref[I.5.5.1-ko]{5.5.1, (ii)}). 이제
(\hyperref[I.8.2.7-ko]{8.2.7})의 가정들이 모두 성립한다. 그러면 모든
$x\in X$에 대하여 $\theta_x^\sharp$는 단사이다. 한편 $f$가 국소
몰입 사상이라는 가정에 따라 $\theta_x^\sharp$는 전사이다
(\hyperref[I.4.2.2-ko]{4.2.2}). 따라서 $\theta_x^\sharp$는
전단사이고, (\hyperref[I.8.2.7-ko]{8.2.7})의 동일시 아래
$\mathscr{O}_{\psi(x)}=\mathscr{O}_x$이다. 그러므로
(\hyperref[I.8.2.4-ko]{8.2.4})에 따라 $\psi$는 \emph{단사}이다. $f$가 국소
동형사상인 경우에는 이것으로 명제가 증명된다
(\hyperref[I.4.5.3-ko]{4.5.3}). $f$가 국소 몰입 사상이라고만
가정하는 경우, 모든 $x\in X$에 대하여 $X$에서 $x$의 열린 근방
$U$와 $Y$에서 $\psi(x)$의 열린 근방 $V$가 존재하여 $\psi$를 $U$에
제한한 사상은 $U$에서 $V$의 어떤 \emph{닫힌} 부분집합으로 가는
위상동형이다. 그런데 $U$는 $X$에서 조밀하므로 $\psi(U)$는 $Y$에서
조밀하고, \emph{따라서 특히} $V$에서 조밀하다. 그러므로
$\psi(U)=V$이다. $\psi$가 단사이므로 $\psi^{-1}(V)=U$이고, 이로써
$\psi$가 $X$에서 $\psi(X)$로 가는 위상동형임이 증명된다.

\subsection{슈발레 스킴}
\label{subsection:I.8.3-ko}

\begin{env}[8.3.1]
\label{I.8.3.1-ko}
$X$를 \emph{뇌터} 정역 스킴, $R$을 그 유리함수체라 하자. $x$가 $X$를
달릴 때 국소 부분환 $\mathscr{O}_x\subset R$들의 집합을
$X'$이라 쓰자. 집합 $X'$은 다음 세 조건을 만족한다.
\begin{enumerate}
  \item[(Sch. 1)] \emph{모든 $M\in X'$에 대하여 $R$은 $M$의
    분수체이다.}
  \item[(Sch. 2)] \emph{$R$의 뇌터 부분환 $A_i$가 유한개 존재하여
    $X'=\bigcup_i L(A_i)$이고, 모든 첨자쌍 $i$, $j$에 대하여
    $A_i\cup A_j$에 의해 $R$에서 생성되는 부분환 $A_{ij}$는 $A_i$
    위의 유한 생성 대수이다.}
  \item[(Sch. 3)] \emph{$X'$의 서로 연관된 두 원소 $M$, $N$은
    같다.}
\end{enumerate}
\end{env}

실제로 (Sch. 1)은 (\hyperref[I.8.2.1-ko]{8.2.1})에서 보았고,
(Sch. 3)은 (\hyperref[I.8.2.2-ko]{8.2.2})에서 따른다. (Sch. 2)를
보이기 위해서는 $X$를 환이 뇌터인 유한개의 아핀 열린집합 $U_i$로
덮고 $A_i=\Gamma(U_i,\mathscr{O}_X)$로 두면 충분하다. $X$가 스킴이라는
가정에 따라 $U_i\cap U_j$는 아핀이고
$\Gamma(U_i\cap U_j,\mathscr{O}_X)=A_{ij}$이다
(\hyperref[I.5.5.6-ko]{5.5.6}). 더욱이 $U_i$의 바탕공간이 뇌터이므로
몰입 사상 $U_i\cap U_j\to U_i$는 유한형이다
(\hyperref[I.6.3.5-ko]{6.3.5}). 따라서 $A_{ij}$는 $A_i$ 위의 유한
생성 대수이다(\hyperref[I.6.3.3-ko]{6.3.3}).

\begin{env}[8.3.2]
\label{I.8.3.2-ko}
공리 (Sch. 1), (Sch. 2), (Sch. 3)을 갖는 구조는 C.~슈발레가 말한
의미의 「스킴」을 일반화한다. 슈발레는 이에 더하여 $R$이 어떤 체
$K$의 유한 생성 확대체이고 $A_i$들이 $K$ 위의 유한 생성
대수라고 가정한다. 이 가정 아래에서는 (Sch. 2)의 일부가 불필요하다
\cite{I-1}. 역으로 집합 $X'$ 위에 이러한 구조가 주어지면
(\hyperref[I.8.2.6-ko]{8.2.6})의 관찰을 이용하여 이에 정역 스킴 $X$를
대응시킬 수 있다. $X$의 바탕공간은
(\hyperref[I.8.2.6-ko]{8.2.6})에서 정의한 위상을 갖춘 $X'$이고,
구조층 $\mathscr{O}_X$는

\oldpage[I]{169}
모든 열린집합 $U\subset X$에 대하여
$\Gamma(U,\mathscr{O}_X)=\bigcap_{x\in U}\mathscr{O}_x$를 만족하며
제한 준동형은 자명하게 정의된다. 이로써 국소환들이 정확히 $X'$의
원소들인 정역 스킴을 실제로 얻는다는 확인은 독자에게 맡긴다. 우리는
뒤에서 이 결과를 사용하지 않는다.
\end{env}
% ===== END EXACT INLINE INPUT: c1s8.tex =====


% ===== BEGIN EXACT INLINE INPUT: c1s9.tex =====
\section{준연접층에 관한 보충}
\label{section:I.9-ko}

\subsection{준연접층의 텐서곱}
\label{subsection:I.9.1-ko}

\begin{proposition}[9.1.1]
\label{I.9.1.1-ko}
$X$를 준스킴(각각 국소 뇌터 준스킴)이라 하자. $\mathscr{F}$와
$\mathscr{G}$를 두 준연접(각각 연접) $\mathscr{O}_X$-가군이라 하자.
그러면 $\mathscr{F}\otimes_{\mathscr{O}_X}\mathscr{G}$는 준연접(각각
연접)이고, $\mathscr{F}$와 $\mathscr{G}$가 유한 생성이면 이것도 유한
생성이다. $\mathscr{F}$가 유한 표시를 갖고 $\mathscr{G}$가 준연접(각각
연접)이면 $\mathscr{H}\!om_{\mathscr{O}_X}(\mathscr{F},\mathscr{G})$는
준연접(각각 연접)이다.
\end{proposition}

문제가 국소적이므로 $X$가 아핀(각각 뇌터 아핀)이라고 가정할 수 있다.
더욱이 $\mathscr{F}$가 연접이면 어떤 준동형
$\mathscr{O}_X^m\to\mathscr{O}_X^n$의 여핵이라고 가정할 수 있다.
그러면 준연접층에 관한 주장들은
(\hyperref[I.1.3.12-ko]{1.3.12})와
(\hyperref[I.1.3.9-ko]{1.3.9})에서 따르고, 연접층에 관한 주장들은
(\hyperref[I.1.5.1-ko]{1.5.1})과 다음 사실에서 따른다. 뇌터 환 $A$
위의 유한 생성 가군 $M$, $N$에 대하여 $M\otimes_A N$과
$\operatorname{Hom}_A(M,N)$은 유한 생성 $A$-가군이다.

\begin{definition}[9.1.2]
\label{I.9.1.2-ko}
$X$, $Y$를 두 $S$-준스킴, $p$, $q$를 $X\times_S Y$의 사영,
$\mathscr{F}$(각각 $\mathscr{G}$)를 준연접 $\mathscr{O}_X$-가군(각각
$\mathscr{O}_Y$-가군)이라 하자. $X\times_S Y$ 위의 텐서곱
$p^*(\mathscr{F})\otimes_{\mathscr{O}_{X\times_S Y}}q^*(\mathscr{G})$을
$\mathscr{F}$와 $\mathscr{G}$의 $\mathscr{O}_S$ 위의(또는 $S$
위의) 텐서곱이라 하고
$\mathscr{F}\otimes_{\mathscr{O}_S}\mathscr{G}$(또는
$\mathscr{F}\otimes_S\mathscr{G}$)로 쓴다.
\end{definition}

$X_i$($1\leq i\leq n$)가 $S$-준스킴이고 $\mathscr{F}_i$가 준연접
$\mathscr{O}_{X_i}$-가군($1\leq i\leq n$)이면, 같은 방식으로 준스킴
$Z=X_1\times_S X_2\times_S\cdots\times_S X_n$ 위의 텐서곱
$\mathscr{F}_1\otimes_S\mathscr{F}_2\otimes_S\cdots\otimes_S\mathscr{F}_n$
을 정의한다. (\hyperref[I.9.1.1-ko]{9.1.1})과
(\hyperref[0.5.1.4-ko]{0, 5.1.4})에 따라 이는
\emph{준연접} $\mathscr{O}_Z$-가군이다. $\mathscr{F}_i$들이 연접이고
$Z$가 \emph{국소 뇌터}이면 (\hyperref[I.9.1.1-ko]{9.1.1}),
(\hyperref[0.5.3.11-ko]{0, 5.3.11}),
(\hyperref[I.6.1.1-ko]{6.1.1})에 따라 이 가군은 \emph{연접}이다.

$X=Y=S$로 두면 정의 (\hyperref[I.9.1.2-ko]{9.1.2})에서
$\mathscr{O}_S$-가군의 텐서곱을 다시 얻는다. 또한
$q^*(\mathscr{O}_Y)=\mathscr{O}_{X\times_S Y}$이므로
(\hyperref[0.4.3.4-ko]{0, 4.3.4}),
$\mathscr{F}\otimes_S\mathscr{O}_Y$는 $p^*(\mathscr{F})$와 표준적으로
동형이고, 마찬가지로 $\mathscr{O}_X\otimes_S\mathscr{G}$는
$q^*(\mathscr{G})$와 표준적으로 동형이다. 특히 $Y=S$로 두고
$f:X\to Y$를 구조 사상이라 하면
$\mathscr{O}_X\otimes_Y\mathscr{G}=f^*(\mathscr{G})$이다. 따라서 보통의
텐서곱과 역상층은 일반적인 텐서곱의 특수한 경우로 나타난다.

정의 (\hyperref[I.9.1.2-ko]{9.1.2})에서 곧바로, $X$와 $Y$를 고정하면
$\mathscr{F}\otimes_S\mathscr{G}$가 $\mathscr{F}$와 $\mathscr{G}$ 각각에
관하여 \emph{공변이고 가법적이며 오른쪽 완전한 쌍함자}임을 얻는다.

\begin{proposition}[9.1.3]
\label{I.9.1.3-ko}
$S$, $X$, $Y$를 각각 환 $A$, $B$, $C$의 아핀 스킴이라 하고, $B$와
$C$를 $A$-대수라 하자. $M$(각각 $N$)을 $B$-가군(각각 $C$-가군),
$\mathscr{F}=\widetilde{M}$(각각 $\mathscr{G}=\widetilde{N}$)을 이에
대응하는 준연접층이라 하자. 그러면 $\mathscr{F}\otimes_S\mathscr{G}$는
$(B\otimes_A C)$-가군 $M\otimes_A N$에 대응하는 층과 표준적으로
동형이다.
\end{proposition}

\begin{proof}
\oldpage[I]{170}
(\hyperref[I.1.6.5-ko]{1.6.5})에 따라
$\mathscr{F}\otimes_S\mathscr{G}$는 $(B\otimes_A C)$-가군
\[
  \bigl(M\otimes_B(B\otimes_A C)\bigr)
  \otimes_{B\otimes_A C}
  \bigl((B\otimes_A C)\otimes_C N\bigr)
\]
에 대응하는 층과 표준적으로 동형이다. 텐서곱 사이의 표준 동형들에
따라 이 가군은
\[
  M\otimes_B(B\otimes_A C)\otimes_C N
  =(M\otimes_B B)\otimes_A(C\otimes_C N)=M\otimes_A N
\]
과 동형이다.
\end{proof}

\begin{proposition}[9.1.4]
\label{I.9.1.4-ko}
$f:T\to X$, $g:T\to Y$를 두 $S$-사상, $\mathscr{F}$(각각
$\mathscr{G}$)를 준연접 $\mathscr{O}_X$-가군(각각
$\mathscr{O}_Y$-가군)이라 하자. 그러면
\[
  (f,g)_S^*(\mathscr{F}\otimes_S\mathscr{G})
  =f^*(\mathscr{F})\otimes_{\mathscr{O}_T}g^*(\mathscr{G}).
\]
\end{proposition}

\begin{proof}
$p$, $q$가 $X\times_S Y$의 사영이면, 이 공식은
$(f,g)_S^*\circ p^*=f^*$, $(f,g)_S^*\circ q^*=g^*$
(\hyperref[0.3.5.5-ko]{0, 3.5.5})와 대수적 층들의 텐서곱의 역상층이
그 역상층들의 텐서곱이라는 사실
(\hyperref[0.4.3.3-ko]{0, 4.3.3})에서 따른다.
\end{proof}

\begin{corollary}[9.1.5]
\label{I.9.1.5-ko}
$f:X\to X'$, $g:Y\to Y'$을 두 $S$-사상, $\mathscr{F}'$(각각
$\mathscr{G}'$)을 준연접 $\mathscr{O}_{X'}$-가군(각각
$\mathscr{O}_{Y'}$-가군)이라 하자. 그러면
\[
  (f\times_S g)^*(\mathscr{F}'\otimes_S\mathscr{G}')
  =f^*(\mathscr{F}')\otimes_S g^*(\mathscr{G}').
\]
\end{corollary}

\begin{proof}
$p$, $q$를 $X\times_S Y$의 사영이라 하면
$f\times_S g=(f\circ p,g\circ q)_S$이므로, 따름정리는
(\hyperref[I.9.1.4-ko]{9.1.4})에서 따른다.
\end{proof}

\begin{corollary}[9.1.6]
\label{I.9.1.6-ko}
$X$, $Y$, $Z$를 세 $S$-준스킴, $\mathscr{F}$(각각 $\mathscr{G}$,
$\mathscr{H}$)를 준연접 $\mathscr{O}_X$-가군(각각
$\mathscr{O}_Y$-가군, $\mathscr{O}_Z$-가군)이라 하자. 층
$\mathscr{F}\otimes_S\mathscr{G}\otimes_S\mathscr{H}$는
$X\times_S Y\times_S Z$에서 $(X\times_S Y)\times_S Z$로 가는 표준
동형사상에 의한
$(\mathscr{F}\otimes_S\mathscr{G})\otimes_S\mathscr{H}$의 역상층이다.
\end{corollary}

\begin{proof}
$p_1$, $p_2$, $p_3$을 $X\times_S Y\times_S Z$의 사영이라 하면 이
동형사상은 $(p_1,p_2)_S\times_S p_3$으로 쓸 수 있다.

마찬가지로 $X\times_S Y$에서 $Y\times_S X$로 가는 표준 동형사상에
의한 $\mathscr{G}\otimes_S\mathscr{F}$의 역상층은
$\mathscr{F}\otimes_S\mathscr{G}$이다.
\end{proof}

\begin{corollary}[9.1.7]
\label{I.9.1.7-ko}
$X$가 $S$-준스킴이면, 모든 준연접 $\mathscr{O}_X$-가군
$\mathscr{F}$는 $X$에서 $X\times_S S$로 가는 표준 동형사상
(\hyperref[I.3.3.3-ko]{3.3.3})에 의한
$\mathscr{F}\otimes_S\mathscr{O}_S$의 역상층이다.
\end{corollary}

\begin{proof}
$\varphi:X\to S$가 구조 사상이면 이 동형사상은 $(1_X,\varphi)_S$이다.
따라서 따름정리는 (\hyperref[I.9.1.4-ko]{9.1.4})와
$\varphi^*(\mathscr{O}_S)=\mathscr{O}_X$에서 따른다.
\end{proof}

\begin{env}[9.1.8]
\label{I.9.1.8-ko}
$X$를 $S$-준스킴, $\mathscr{F}$를 준연접 $\mathscr{O}_X$-가군,
$\varphi:S'\to S$를 사상이라 하자. $X\times_S S'
=X_{(\varphi)}=X_{(S')}$ 위의 준연접층
$\mathscr{F}\otimes_S\mathscr{O}_{S'}$를
$\mathscr{F}_{(\varphi)}$ 또는 $\mathscr{F}_{(S')}$로 쓴다. 따라서
$p:X_{(S')}\to X$가 사영이면
$\mathscr{F}_{(S')}=p^*(\mathscr{F})$이다.
\end{env}

\begin{proposition}[9.1.9]
\label{I.9.1.9-ko}
$\varphi':S''\to S'$을 사상이라 하자. $S$-준스킴 $X$ 위의 모든 준연접
$\mathscr{O}_X$-가군 $\mathscr{F}$에 대하여
$(\mathscr{F}_{(\varphi)})_{(\varphi')}$은 표준 동형사상
$(X_{(\varphi)})_{(\varphi')}\xrightarrow{\sim}
X_{(\varphi\circ\varphi')}$
(\hyperref[I.3.3.9-ko]{3.3.9})에 의한
$\mathscr{F}_{(\varphi\circ\varphi')}$의 역상층이다.
\end{proposition}

\begin{proof}
이는 정의와 (\hyperref[I.3.3.9-ko]{3.3.9})에서 곧바로 따르며, 다음과
같이 쓸 수도 있다.
\[
  (\mathscr{F}\otimes_S\mathscr{O}_{S'})
  \otimes_{S'}\mathscr{O}_{S''}
  =\mathscr{F}\otimes_S\mathscr{O}_{S''}.
  \tag{9.1.9.1}\label{I.9.1.9.1-ko}
\]
\end{proof}

\begin{proposition}[9.1.10]
\label{I.9.1.10-ko}
$Y$를 $S$-준스킴, $f:X\to Y$를 $S$-사상이라 하자. 모든 준연접
$\mathscr{O}_Y$-가군 $\mathscr{G}$와 모든 사상 $S'\to S$에 대하여
\[
  (f_{(S')})^*(\mathscr{G}_{(S')})
  =(f^*(\mathscr{G}))_{(S')}
\]
이다.
\end{proposition}

\begin{proof}
이는 다음 그림의 가환성에서 곧바로 따른다.
\oldpage[I]{171}
\[
  \xymatrix{
    X_{(S')}\ar[r]^{f_{(S')}}\ar[d] &
    Y_{(S')}\ar[d]\\
    X\ar[r]_f &
    Y
  }
\]
\end{proof}

\begin{corollary}[9.1.11]
\label{I.9.1.11-ko}
$X$, $Y$를 두 $S$-준스킴, $\mathscr{F}$(각각 $\mathscr{G}$)를 준연접
$\mathscr{O}_X$-가군(각각 $\mathscr{O}_Y$-가군)이라 하자. 표준
동형사상 $(X\times_S Y)_{(S')}\xrightarrow{\sim}
(X_{(S')})\times_{S'}(Y_{(S')})$
(\hyperref[I.3.3.10-ko]{3.3.10})에 의한
$\mathscr{F}_{(S')}\otimes_{S'}\mathscr{G}_{(S')}$의 역상층은
$(\mathscr{F}\otimes_S\mathscr{G})_{(S')}$과 같다.
\end{corollary}

\begin{proof}
$p$, $q$가 $X\times_S Y$의 사영이면 위 동형사상은
$(p_{(S')},q_{(S')})_{S'}$이다. 따라서 따름정리는
(\hyperref[I.9.1.4-ko]{9.1.4})와
(\hyperref[I.9.1.10-ko]{9.1.10})에서 따른다.
\end{proof}

\begin{proposition}[9.1.12]
\label{I.9.1.12-ko}
(\hyperref[I.9.1.2-ko]{9.1.2})의 표기 아래에서
$z\in X\times_S Y$, $x=p(z)$, $y=q(z)$라 하자. 줄기
$(\mathscr{F}\otimes_S\mathscr{G})_z$는
$(\mathscr{F}_x\otimes_{\mathscr{O}_x}\mathscr{O}_z)
\otimes_{\mathscr{O}_z}
(\mathscr{G}_y\otimes_{\mathscr{O}_y}\mathscr{O}_z)
=\mathscr{F}_x\otimes_{\mathscr{O}_x}\mathscr{O}_z
\otimes_{\mathscr{O}_y}\mathscr{G}_y$와 동형이다.
\end{proposition}

\begin{proof}
아핀인 경우로 귀착할 수 있으므로 명제는 공식
(\hyperref[I.1.6.5.1-ko]{1.6.5.1})에서 따른다.
\end{proof}

\begin{corollary}[9.1.13]
\label{I.9.1.13-ko}
$\mathscr{F}$와 $\mathscr{G}$가 유한 생성이면
\[
  \operatorname{Supp}(\mathscr{F}\otimes_S\mathscr{G})
  =p^{-1}(\operatorname{Supp}(\mathscr{F}))
  \cap q^{-1}(\operatorname{Supp}(\mathscr{G})).
\]
\end{corollary}

\begin{proof}
$p^*(\mathscr{F})$와 $q^*(\mathscr{G})$는 유한 생성
$\mathscr{O}_{X\times_S Y}$-가군이므로,
(\hyperref[I.9.1.12-ko]{9.1.12})와
(\hyperref[0.1.7.5-ko]{0, 1.7.5})에 따라
$\mathscr{G}=\mathscr{O}_Y$인 경우, 다시 말해 다음 공식을 보이는
경우로 귀착한다.
\[
  \operatorname{Supp}(p^*(\mathscr{F}))
  =p^{-1}(\operatorname{Supp}(\mathscr{F})).
  \tag{9.1.13.1}\label{I.9.1.13.1-ko}
\]\footnote{\textbf{번역 주.} 인쇄 원문은 왼쪽에
$\operatorname{Supp}(p^{-1}(\mathscr{F}))$를 쓴다. 그러나 바로 앞의
귀착은 $\mathscr{F}\otimes_S\mathscr{O}_Y=p^*(\mathscr{F})$를 사용하고,
이어지는 줄기 계산도 역상층 $p^*(\mathscr{F})$에 관한 것이다. 따라서
형에 맞는 $p^*$를 사용하였다.}

(\hyperref[0.1.7.5-ko]{0, 1.7.5})에서와 같은 논증으로, 모든
$z\in X\times_S Y$에 대하여 $x=p(z)$라 할 때
$\mathscr{O}_z/\mathfrak{m}_x\mathscr{O}_z\neq0$임을 확인하면 충분하다.
이는 준동형 $\mathscr{O}_x\to\mathscr{O}_z$가 가정에 따라
\emph{국소}라는 사실에서 따른다.
\end{proof}

이 항에서 두 인자의 곱에 대하여 증명한 결과들을 임의의 유한 개수의
인자의 곱으로 확장하는 일은 독자에게 맡긴다.

\subsection{준연접층의 직상}
\label{subsection:I.9.2-ko}

\begin{proposition}[9.2.1]
\label{I.9.2.1-ko}
$f:X\to Y$를 준스킴의 사상이라 하자. $Y$가 다음 성질을 갖는 아핀
열린집합들의 덮개 $(Y_\alpha)$를 갖는다고 가정하자. 각
$f^{-1}(Y_\alpha)$는 $f^{-1}(Y_\alpha)$ 안에 포함된 아핀 열린집합들의 유한 덮개
$(X_{\alpha i})$를 가지며, 각 교집합
$X_{\alpha i}\cap X_{\alpha j}$ 자체가 아핀 열린집합들의 유한
합집합이다. 이 조건 아래에서 모든 준연접 $\mathscr{O}_X$-가군
$\mathscr{F}$에 대하여 $f_*(\mathscr{F})$는 준연접
$\mathscr{O}_Y$-가군이다.
\end{proposition}

\begin{proof}
문제는 $Y$ 위에서 국소적이므로 $Y$가 $Y_\alpha$들 가운데 하나와
같다고 가정하고 첨자 $\alpha$를 생략할 수 있다.

\begin{enumerate}
\item[a)] 먼저 $X_i\cap X_j$ 자체가 \emph{아핀} 열린집합이라고
가정하자. $\mathscr{F}_i=\mathscr{F}|X_i$,
$\mathscr{F}_{ij}=\mathscr{F}|(X_i\cap X_j)$라 두고,
$\mathscr{F}'_i$와 $\mathscr{F}'_{ij}$를 각각 $f$를 $X_i$와
$X_i\cap X_j$에 제한한 사상에 의한 $\mathscr{F}_i$와
$\mathscr{F}_{ij}$의 직상이라 하자. $\mathscr{F}'_i$와
$\mathscr{F}'_{ij}$는 준연접임이 알려져 있다
(\hyperref[I.1.6.3-ko]{1.6.3}). 이제
$\mathscr{G}=\bigoplus_i\mathscr{F}'_i$,
$\mathscr{H}=\bigoplus_{i,j}\mathscr{F}'_{ij}$라 두자.
$\mathscr{G}$와 $\mathscr{H}$는 준연접 $\mathscr{O}_Y$-가군이다.
준동형 $u:\mathscr{G}\to\mathscr{H}$를 정의하여
$f_*(\mathscr{F})$가 $u$의 \emph{핵}이 되게 할 것이다. 그러면
(\hyperref[I.1.3.9-ko]{1.3.9})에 따라 $f_*(\mathscr{F})$가 준연접임이
따른다. $u$를
\oldpage[I]{172}
준층의 준동형으로 정의하면 충분하다. 따라서 $\mathscr{G}$와
$\mathscr{H}$의 정의를 고려하면, 모든 열린집합 $W\subset Y$에 대하여
$W$가 변할 때 통상적 양립 조건들을 만족하는 준동형
\[
  u_W:\bigoplus_i\Gamma(f^{-1}(W)\cap X_i,\mathscr{F})
  \longrightarrow
  \bigoplus_{i,j}\Gamma(f^{-1}(W)\cap X_i\cap X_j,\mathscr{F})
\]
을 정의하면 된다. 각 단면
$s_i\in\Gamma(f^{-1}(W)\cap X_i,\mathscr{F})$에 대하여
$f^{-1}(W)\cap X_i\cap X_j$로의 제한을 $s_{i|j}$로 나타내고
\[
  u_W((s_i))=(s_{i|j}-s_{j|i})
\]
라 두면 양립 조건들이 자명하게 성립한다. $u$의 핵 $\mathscr{R}$이
$f_*(\mathscr{F})$임을 보이기 위하여, 모든 단면
$s\in\Gamma(f^{-1}(W),\mathscr{F})$에 그 제한들로 이루어진 족
$(s_i)$를 대응시켜 $f_*(\mathscr{F})$에서 $\mathscr{R}$로 가는
준동형을 정의하자. 여기서 $s_i$는 $s$를
$f^{-1}(W)\cap X_i$에 제한한 것이다. 층의 공리 (F 1)과 (F 2)
(G, II, 1.1)에 따라 이 준동형은 \emph{전단사}이고, 이 경우의 증명이
끝난다.

\item[b)] 일반적인 경우에는 $\mathscr{F}'_{ij}$가 준연접임을 보인
뒤에 같은 논증을 적용한다. 가정에 따라 $X_i\cap X_j$는 아핀
열린집합 $X_{ijk}$들의 유한 합집합이다. 그런데 $X_{ijk}$들은
\emph{어떤 스킴 안의} 아핀 열린집합들이므로 그 가운데 임의의 두
열린집합의 교집합도 아핀 열린집합이다
(\hyperref[I.5.5.6-ko]{5.5.6}). 따라서 첫째 경우로 귀착되고,
(9.2.1)이 증명된다.
\end{enumerate}
\end{proof}

\begin{corollary}[9.2.2]
\label{I.9.2.2-ko}
(9.2.1)의 결론은 다음 각 경우에 성립한다.
\begin{enumerate}
\item[a)] $f$가 분리 사상이고 준콤팩트이다.
\item[b)] $f$가 분리 사상이고 유한형이다.
\item[c)] $f$가 준콤팩트이고 $X$의 바탕공간이 국소 뇌터이다.
\end{enumerate}
\end{corollary}

\begin{proof}
a)의 경우에는 $X_{\alpha i}\cap X_{\alpha j}$가 아핀이다
(\hyperref[I.5.5.6-ko]{5.5.6}). b)는 a)의 특수한 경우이다
(\hyperref[I.6.6.3-ko]{6.6.3}). 마지막으로 c)의 경우에는 $Y$가
아핀이고 $X$의 바탕공간이 뇌터인 경우로 귀착할 수 있다. 이때 $X$는
아핀 열린집합들의 유한 덮개 $(X_i)$를 갖고, $X_i\cap X_j$는
준콤팩트이므로 아핀 열린집합들의 유한 합집합이다
(\hyperref[I.2.1.3-ko]{2.1.3}).
\end{proof}

\subsection{준연접층의 단면 연장}
\label{subsection:I.9.3-ko}

\begin{theorem}[9.3.1]
\label{I.9.3.1-ko}
$X$를 바탕공간이 뇌터인 준스킴이거나 바탕공간이 준콤팩트인 스킴이라
하자. $\mathscr{L}$을 가역 $\mathscr{O}_X$-가군
(\hyperref[0.5.4.1-ko]{0, 5.4.1}), $f$를 $X$ 위의 $\mathscr{L}$의
단면, $X_f$를 $f(x)\neq0$인 $x\in X$들의 열린집합
(\hyperref[0.5.5.1-ko]{0, 5.5.1}), $\mathscr{F}$를 준연접
$\mathscr{O}_X$-가군이라 하자.
\begin{enumerate}
\item[(i)] $s\in\Gamma(X,\mathscr{F})$이고 $s|X_f=0$이면,
  $s\otimes f^{\otimes n}=0$인 정수 $n>0$이 존재한다.
\item[(ii)] 모든 단면 $s\in\Gamma(X_f,\mathscr{F})$에 대하여,
  $s\otimes f^{\otimes n}$이 $X$ 위의
  $\mathscr{F}\otimes\mathscr{L}^{\otimes n}$의 단면으로 연장되는 정수
  $n>0$이 존재한다.
\end{enumerate}
\end{theorem}

\begin{proof}
\begin{enumerate}
\item[(i)] $X$의 바탕공간은 준콤팩트이므로,
  $\mathscr{L}|U_i$가 $\mathscr{O}_X|U_i$와 동형인 아핀 열린집합
  $U_i$들의 유한 합집합이다. 따라서 $X$가 아핀이고
  $\mathscr{L}=\mathscr{O}_X$인 경우로 귀착된다. 이 경우 $f$는
  $A(X)$의 한 원소와 동일시되고 $X_f=D(f)$이다. 또한 $s$는 어떤
  $A(X)$-가군 $M$의 한 원소와 동일시되며, $s|X_f$는 $M_f$의 대응하는
  원소와 동일시된다. 그러므로 결과는 분수가군의 정의에서 곧바로
  따른다.
\oldpage[I]{173}

\item[(ii)] 다시 $X$는 $\mathscr{L}|U_i\simeq\mathscr{O}_X|U_i$인 아핀
  열린집합 $U_i$들($1\leq i\leq r$)의 유한 합집합이다. 각 $i$에
  대하여, 앞의 동형에 의해
  $(s\otimes f^{\otimes n})|(U_i\cap X_f)$는
  $(f|(U_i\cap X_f))^n(s|(U_i\cap X_f))$와 동일시된다.
  그러면 (\hyperref[I.1.4.1-ko]{1.4.1})에 따라 모든 $i$에 대하여
  $(s\otimes f^{\otimes n})|(U_i\cap X_f)$가 $U_i$ 위의
  $\mathscr{F}\otimes\mathscr{L}^{\otimes n}$의 단면 $s_i$로 연장되는
  정수 $n>0$이 존재한다. $s_{i|j}$를 $s_i$의 $U_i\cap U_j$로의
  제한이라 하자. 정의에 따라 $X_f\cap U_i\cap U_j$에서
  $s_{i|j}-s_{j|i}=0$이다. 그런데 $X$의 바탕공간이 뇌터이면
  $U_i\cap U_j$는 준콤팩트이고, $X$가 스킴이면 $U_i\cap U_j$는 아핀
  열린집합이므로 (\hyperref[I.5.5.6-ko]{5.5.6}) 역시 준콤팩트이다.
  따라서 (i)에 의해 $i$와 $j$에 무관한 정수 $m$이 존재하여
  $(s_{i|j}-s_{j|i})\otimes f^{\otimes m}=0$이다. 이에 따라 $X$ 위의
  $\mathscr{F}\otimes\mathscr{L}^{\otimes(n+m)}$의 단면 $s'$가 존재하여,
  각 $U_i$ 위에서의 제한은 $s_i\otimes f^{\otimes m}$이고, 따라서
  $X_f$ 위에서의 제한은 $s\otimes f^{\otimes(n+m)}$이다.
\end{enumerate}
\end{proof}

다음 따름정리들은 정리 (9.3.1)을 더 대수적인 언어로 해석한다.

\begin{corollary}[9.3.2]
\label{I.9.3.2-ko}
(9.3.1)의 가정 아래에서 등급환 $A_* = \Gamma_*(\mathscr{L})$과
등급 $A_*$-가군 $M_* = \Gamma_*(\mathscr{L},\mathscr{F})$를 생각하자
(\hyperref[0.5.4.6-ko]{0, 5.4.6}). $f\in A_n$이고
$n\in\mathbb{Z}$이면 표준 동형
\[
  \Gamma(X_f,\mathscr{F})\xrightarrow{\sim}((M_*)_f)_0
\]
이 존재한다. 여기서 오른쪽은 분수가군 $(M_*)_f$의 차수 $0$인
원소들로 이루어진 부분군이다.
\end{corollary}

\begin{corollary}[9.3.3]
\label{I.9.3.3-ko}
(9.3.1)의 가정에 더하여 $\mathscr{L}=\mathscr{O}_X$라고 하자. 이때
$A=\Gamma(X,\mathscr{O}_X)$, $M=\Gamma(X,\mathscr{F})$로 놓으면,
$A_f$-가군 $\Gamma(X_f,\mathscr{F})$는 $M_f$와 표준적으로 동형이다.
\end{corollary}

\begin{proposition}[9.3.4]
\label{I.9.3.4-ko}
$X$를 뇌터 준스킴, $\mathscr{F}$를 연접 $\mathscr{O}_X$-가군,
$\mathscr{J}$를 $\mathscr{O}_X$의 연접 아이디얼층이라 하자.
$\mathscr{F}$의 지지집합이 $\mathscr{O}_X/\mathscr{J}$의 지지집합에
포함되면, $\mathscr{J}^n\mathscr{F}=0$인 정수 $n>0$이 존재한다.
\end{proposition}

\begin{proof}
$X$는 환이 뇌터인 아핀 열린집합들의 유한 합집합이므로, $X$가 뇌터
환 $A$의 아핀 스킴인 경우로 귀착할 수 있다. 그러면
$\mathscr{F}=\widetilde M$이고, 여기서
$M=\Gamma(X,\mathscr{F})$는 유한 생성 $A$-가군이다. 또한
$\mathscr{J}=\widetilde{\mathfrak{J}}$이고, 여기서
$\mathfrak{J}=\Gamma(X,\mathscr{J})$는 $A$의 아이디얼이다
(\hyperref[I.1.4.1-ko]{1.4.1}과 \hyperref[I.1.5.1-ko]{1.5.1}).
$A$가 뇌터이므로 $\mathfrak{J}$는 유한 생성계 $f_i$들
($1\leq i\leq m$)을 갖는다. 가정에 따라 $X$ 위의 $\mathscr{F}$의 모든
단면은 각 $D(f_i)$ 위에서 영단면이다. $s_j$들($1\leq j\leq q$)이
$M$을 생성하는 $\mathscr{F}$의 단면들이면,
(\hyperref[I.1.4.1-ko]{1.4.1})에 따라 $i$와 $j$에 무관한 정수 $h$가
존재하여 $f_i^h s_j=0$이다. 따라서 모든 $s\in M$에 대하여
$f_i^h s=0$이다. 그러므로 $n=mh$로 놓으면
$\mathfrak{J}^nM=0$이고, 이에 대응하는 $\mathscr{O}_X$-가군
$\mathscr{J}^n\mathscr{F}=\widetilde{\mathfrak{J}^nM}$
(\hyperref[I.1.3.13-ko]{1.3.13})은 영가군이다.
\end{proof}

\begin{corollary}[9.3.5]
\label{I.9.3.5-ko}
(9.3.4)의 가정 아래에서, 바탕공간이
$\mathscr{O}_X/\mathscr{J}$의 지지집합인 $X$의 닫힌 부분준스킴 $Y$가
존재한다. 또한 $j:Y\to X$가 표준 포함 사상이면
$\mathscr{F}=j_*(j^*(\mathscr{F}))$이다.
\end{corollary}

\begin{proof}
먼저 $\mathscr{O}_X/\mathscr{J}$와
$\mathscr{O}_X/\mathscr{J}^n$의 지지집합은 같다. 실제로
$\mathscr{J}_x=\mathscr{O}_x$이면
$\mathscr{J}_x^n=\mathscr{O}_x$이고, 한편 모든 $x\in X$에 대하여
$\mathscr{J}_x^n\subset\mathscr{J}_x$이다. 따라서 (9.3.4)에 의해
$\mathscr{J}\mathscr{F}=0$이라고 가정할 수 있다. 이제 $Y$를
$\mathscr{J}$로 정의되는 $X$의 닫힌 부분준스킴으로 잡으면,
$\mathscr{F}$가 $(\mathscr{O}_X/\mathscr{J})$-가군이므로 결론은
곧바로 따른다.
\end{proof}

\oldpage[I]{174}

\subsection{준연접층의 연장}
\label{subsection:I.9.4-ko}

\begin{env}[9.4.1]
\label{I.9.4.1-ko}
$X$를 위상공간, $\mathscr{F}$를 $X$ 위의 집합의 층(각각 군의 층,
환의 층), $U$를 $X$의 열린부분집합, $\psi:U\to X$를 표준 포함 사상,
$\mathscr{G}$를 $\mathscr{F}|U=\psi^*(\mathscr{F})$의 부분층이라 하자.
$\psi_*$는 왼쪽 완전하므로 $\psi_*(\mathscr{G})$는
$\psi_*(\psi^*(\mathscr{F}))$의 부분층이다. 표준 준동형
$\rho:\mathscr{F}\to\psi_*(\psi^*(\mathscr{F}))$
(\hyperref[0.3.5.3-ko]{0, 3.5.3})을 생각할 때,
$\mathscr{F}$의 부분층
$\rho^{-1}(\psi_*(\mathscr{G}))$를
$\overline{\mathscr{G}}$로 나타내겠다. 정의로부터 곧바로, $X$의 모든
열린집합 $V$에 대하여 $\Gamma(V,\overline{\mathscr{G}})$는
$V\cap U$로의 제한이 $V\cap U$ 위의 $\mathscr{G}$의 단면인
$s\in\Gamma(V,\mathscr{F})$들로 이루어짐을 알 수 있다. 따라서
$\overline{\mathscr{G}}|U=\psi^*(\overline{\mathscr{G}})=\mathscr{G}$이고,
$\overline{\mathscr{G}}$는 $U$ 위에 $\mathscr{G}$를 유도하는
$\mathscr{F}$의 \emph{가장 큰} 부분층이다. 이
$\overline{\mathscr{G}}$를 $\mathscr{F}|U$의 부분층
$\mathscr{G}$를 $\mathscr{F}$의 부분층으로 만드는
\emph{표준 연장}이라 하겠다.
\end{env}

\begin{proposition}[9.4.2]
\label{I.9.4.2-ko}
$X$를 준스킴, $U$를 표준 포함 사상 $j:U\to X$가 준콤팩트 사상인
$X$의 열린부분집합이라 하자
\emph{(이는 $X$의 바탕공간이 \emph{국소 뇌터}이면
\emph{모든} $U$에 대하여 성립한다
(\hyperref[I.6.6.4-ko]{6.6.4, (i)}))}. 그러면 다음이 성립한다.
\begin{enumerate}
\item[(i)] 모든 준연접 $(\mathscr{O}_X|U)$-가군 $\mathscr{G}$에 대하여
  $j_*(\mathscr{G})$는 준연접 $\mathscr{O}_X$-가군이고
  $j_*(\mathscr{G})|U=j^*(j_*(\mathscr{G}))=\mathscr{G}$이다.
\item[(ii)] 모든 준연접 $\mathscr{O}_X$-가군 $\mathscr{F}$와
  그 제한의 모든 준연접 부분-$(\mathscr{O}_X|U)$-가군
  $\mathscr{G}$에 대하여, $\mathscr{G}$의 표준 연장
  $\overline{\mathscr{G}}$(9.4.1)는 $\mathscr{F}$의 준연접
  부분-$\mathscr{O}_X$-가군이다.
\end{enumerate}
\end{proposition}

\begin{proof}
$j=(\psi,\theta)$라 하자. 여기서 $\psi$는 바탕공간의 포함 사상
$U\to X$이다. 모든 $(\mathscr{O}_X|U)$-가군 $\mathscr{G}$에 대하여
정의상 $j_*(\mathscr{G})=\psi_*(\mathscr{G})$이고, 열린집합 위에 유도된
준스킴의 정의에 따라 모든 $\mathscr{O}_X$-가군 $\mathscr{H}$에 대하여
$j^*(\mathscr{H})=\psi^*(\mathscr{H})=\mathscr{H}|U$이다. 따라서 (i)는
(\hyperref[I.9.2.2-ko]{9.2.2, a)})의 특수한 경우이다. 같은 이유로
$j_*(j^*(\mathscr{F}))$는 준연접이고,
$\overline{\mathscr{G}}$는 준동형
$\rho:\mathscr{F}\to j_*(j^*(\mathscr{F}))$에 의한
$j_*(\mathscr{G})$의 역상이므로 (ii)는
(\hyperref[I.4.1.1-ko]{4.1.1})에서 따른다.
\end{proof}

사상 $j:U\to X$가 준콤팩트라는 가정은 열린집합 $U$가
\emph{준콤팩트}이고 $X$가 \emph{스킴}일 때에도 성립한다. 실제로
$U$는 유한개의 아핀 열린집합 $U_i$들의 합집합이고, $X$의 모든 아핀
열린집합 $V$에 대하여 $V\cap U_i$는 아핀 열린집합
(\hyperref[I.5.5.6-ko]{5.5.6})이므로 준콤팩트이다.

\begin{corollary}[9.4.3]
\label{I.9.4.3-ko}
$X$를 준스킴, $U$를 포함 사상 $j:U\to X$가 준콤팩트인 $X$의
준콤팩트 열린집합이라 하자. 또한 모든 준연접 $\mathscr{O}_X$-가군이
그 유한 생성 준연접 부분-$\mathscr{O}_X$-가군들의 귀납극한이라고
가정하자 \emph{(이는 $X$가 \emph{아핀 스킴}이면 성립한다)}.
$\mathscr{F}$를 준연접 $\mathscr{O}_X$-가군,
$\mathscr{G}$를 $\mathscr{F}|U$의 유한 생성 준연접
부분-$(\mathscr{O}_X|U)$-가군이라 하자. 그러면 $\mathscr{F}$의
유한 생성 준연접 부분-$\mathscr{O}_X$-가군 $\mathscr{G}'$이 존재하여
$\mathscr{G}'|U=\mathscr{G}$이다.
\end{corollary}

\begin{proof}
실제로 $\mathscr{G}=\overline{\mathscr{G}}|U$이고, (9.4.2)에 따라
$\overline{\mathscr{G}}$는 준연접이므로 그 유한 생성 준연접
부분-$\mathscr{O}_X$-가군 $\mathscr{H}_\lambda$들의 귀납극한이다.
따라서 $\mathscr{G}$는 $\mathscr{H}_\lambda|U$들의 귀납극한이고,
유한 생성이므로 (\hyperref[0.5.2.3-ko]{0, 5.2.3})에 따라 어떤
$\mathscr{H}_\lambda|U$와 같다.
\end{proof}

\begin{remark}[9.4.4]
\label{I.9.4.4-ko}
\emph{모든} 아핀 열린집합 $U\subset X$에 대하여 포함 사상
$U\to X$가 준콤팩트라고 가정하자. 모든 아핀 열린집합 $U$와
$\mathscr{F}|U$의 모든 유한 생성 준연접
부분-$(\mathscr{O}_X|U)$-가군 $\mathscr{G}$에 대하여 (9.4.3)의 결론이
성립한다면,
\oldpage[I]{175}
$\mathscr{F}$는 그 유한 생성 준연접 부분-$\mathscr{O}_X$-가군들의
귀납극한이다. 실제로 모든 아핀 열린집합 $U\subset X$에 대하여
$\mathscr{F}|U=\widetilde{M}$이고, 여기서 $M$은 $A(U)$-가군이다.
이 가군은 그 유한 생성 부분가군들의 귀납극한이므로,
$\mathscr{F}|U$는 그 유한 생성 준연접 부분-$(\mathscr{O}_X|U)$-가군들의
귀납극한이다(\hyperref[I.1.3.9-ko]{1.3.9}). 그런데 가정에 따라 이러한
부분가군 각각은 $\mathscr{F}$의 유한 생성 준연접
부분-$\mathscr{O}_X$-가군 $\mathscr{G}_{\lambda,U}$가 $U$ 위에 유도하는
가군이다. $\mathscr{G}_{\lambda,U}$들의 유한합도 다시 유한 생성
준연접 $\mathscr{O}_X$-가군이다. 실제로 이 문제는 국소적이고,
$X$가 아핀인 경우는 (\hyperref[I.1.3.10-ko]{1.3.10})에서 다루었다.
이제 $\mathscr{F}$가 이러한 유한합들의 귀납극한임은 분명하며,
주장이 증명되었다.
\end{remark}

\begin{corollary}[9.4.5]
\label{I.9.4.5-ko}
(9.4.3)의 가정 아래에서 모든 유한 생성 준연접
$(\mathscr{O}_X|U)$-가군 $\mathscr{G}$에 대하여 유한 생성 준연접
$\mathscr{O}_X$-가군 $\mathscr{G}'$이 존재하여
$\mathscr{G}'|U=\mathscr{G}$이다.
\end{corollary}

\begin{proof}
$\mathscr{F}=j_*(\mathscr{G})$는 (9.4.2)에 따라 준연접이고
$\mathscr{F}|U=\mathscr{G}$이므로, $\mathscr{F}$에 (9.4.3)을 적용하면
된다.
\end{proof}

\begin{lemma}[9.4.6]
\label{I.9.4.6-ko}
$X$를 준스킴, $L$을 정렬집합,
$(V_\lambda)_{\lambda\in L}$를 $X$의 아핀 열린집합들로 이루어진 덮개,
$U$를 $X$의 열린집합이라 하자. 모든 $\lambda\in L$에 대하여
$W_\lambda=\bigcup_{\mu<\lambda}V_\mu$로 놓고, 다음을 가정하자.
\begin{enumerate}
  \item[1\textsuperscript{o}] 모든 $\lambda\in L$에 대하여
    $V_\lambda\cap W_\lambda$는 준콤팩트이다.
  \item[2\textsuperscript{o}] 몰입 사상 $U\to X$는 준콤팩트이다.
\end{enumerate}
그러면 모든 준연접 $\mathscr{O}_X$-가군 $\mathscr{F}$와
$\mathscr{F}|U$의 모든 유한 생성 준연접
부분-$(\mathscr{O}_X|U)$-가군 $\mathscr{G}$에 대하여 $\mathscr{F}$의
유한 생성 준연접 부분-$\mathscr{O}_X$-가군 $\mathscr{G}'$이 존재하여
$\mathscr{G}'|U=\mathscr{G}$이다.
\end{lemma}

\begin{proof}
$L$에 그 모든 원소보다 큰 새 원소 $\infty$를 덧붙인 정렬집합을
$\widehat L=L\cup\{\infty\}$라 하고, $\lambda\in\widehat L$에 대하여
$U_\lambda=U\cup\bigcup_{\substack{\mu\in L\\ \mu<\lambda}}V_\mu$로 놓자.
초한 귀납법으로 족 $(\mathscr{G}'_\lambda)_{\lambda\in\widehat L}$를
정의하겠다. 여기서
$\mathscr{G}'_\lambda$는 $\mathscr{F}|U_\lambda$의 유한 생성 준연접
부분-$(\mathscr{O}_X|U_\lambda)$-가군이고,
$\mu<\lambda$이면
$\mathscr{G}'_\lambda|U_\mu=\mathscr{G}'_\mu$이며
$\mathscr{G}'_\lambda|U=\mathscr{G}$이다. 이 족은
(\hyperref[0.3.3.1-ko]{0, 3.3.1})에 따라 서로 붙으며,
$U_\infty=X$이므로 $\mathscr{G}'=\mathscr{G}'_\infty$가 결론을 준다.
이제 $\widehat L$의 최소 원소 $\lambda_0$에 대해서는
$U_{\lambda_0}=U$이고 $\mathscr{G}'_{\lambda_0}=\mathscr{G}$로 둔다.
그 뒤
$\mu<\lambda$인 $\mathscr{G}'_\mu$들이 정의되어 앞의 성질들을
만족한다고 가정하자. $\lambda$가 최소 원소가 아닌 극한 원소이면,
모든 $\mu<\lambda$에 대하여
$\mathscr{G}'_\lambda|U_\mu=\mathscr{G}'_\mu$인
$\mathscr{F}|U_\lambda$의 유일한 부분-$(\mathscr{O}_X|U_\lambda)$-가군을
$\mathscr{G}'_\lambda$로 잡는다. $\mu<\lambda$인 $U_\mu$들이
$U_\lambda$를 덮으므로 이렇게 할 수 있다. 반대로
$\lambda=\mu+1$이면 $U_\lambda=U_\mu\cup V_\mu$이고,
$\mathscr{F}|V_\mu$의 유한 생성 준연접
부분-$(\mathscr{O}_X|V_\mu)$-가군 $\mathscr{G}''_\mu$를
\[
  \mathscr{G}''_\mu|(U_\mu\cap V_\mu)
  =\mathscr{G}'_\mu|(U_\mu\cap V_\mu)~;
\]
가 되도록 정의하면 충분하다. 그런 다음
$\mathscr{G}'_\lambda|U_\mu=\mathscr{G}'_\mu$이고
$\mathscr{G}'_\lambda|V_\mu=\mathscr{G}''_\mu$인
$\mathscr{F}|U_\lambda$의 부분-$(\mathscr{O}_X|U_\lambda)$-가군을
$\mathscr{G}'_\lambda$로 잡는다
(\hyperref[0.3.3.1-ko]{0, 3.3.1}). $V_\mu$가 아핀이므로,
$U_\mu\cap V_\mu$가 준콤팩트임을 보이기만 하면
$\mathscr{G}''_\mu$의 존재는 (9.4.3)에서 따른다. 그런데
$U_\mu\cap V_\mu$는 $U\cap V_\mu$와 $W_\mu\cap V_\mu$의 합집합이고,
둘 다 가정에 따라 준콤팩트이다.\footnote{\emph{[번역 주] 인쇄 원문의
귀납 지표는 $L$뿐이어서
$L$이 최대 원소를 가지면 $U_\lambda$들이 마지막 $V_\lambda$를 덮지
못하고, 최소 원소 단계에서는 앞선 $U_\mu$들이 $U$를 덮는다는 주장도
성립하지 않는다. 위에서는 새 종단 원소 $\infty$를 붙이고 최소 원소
단계를 따로 정의하여 두 지표 결함을 바로잡았다.}}
\end{proof}

\begin{theorem}[9.4.7]
\label{I.9.4.7-ko}
$X$를 준스킴, $U$를 $X$의 열린집합이라 하자. 다음 조건 중 하나가
성립한다고 가정하자.
\begin{enumerate}
  \item[(a)] $X$의 바탕공간은 국소 뇌터이다.
  \item[(b)] $X$는 준콤팩트 스킴이고 $U$는 준콤팩트 열린집합이다.
\end{enumerate}
그러면 모든 준연접 $\mathscr{O}_X$-가군 $\mathscr{F}$와
$\mathscr{F}|U$의 모든 유한 생성 준연접
부분-$(\mathscr{O}_X|U)$-가군 $\mathscr{G}$에 대하여 $\mathscr{F}$의
유한 생성 준연접 부분-$\mathscr{O}_X$-가군 $\mathscr{G}'$이 존재하여
$\mathscr{G}'|U=\mathscr{G}$이다.
\end{theorem}

\oldpage[I]{176}

\begin{proof}
$(V_\lambda)_{\lambda\in L}$를 $X$의 아핀 열린집합들로 이루어진
덮개라 하자. 경우 b)에서는 $L$이 유한하다고 가정한다. $L$에
정렬 순서를 주면, 보조정리 (9.4.6)의 조건들이 성립함을 확인하면
충분하다. 경우 a)에서는 공간 $V_\lambda$들이 뇌터이므로 이는
명백하다. 경우 b)에서는 $V_\lambda\cap V_\mu$가 아핀
(\hyperref[I.5.5.6-ko]{5.5.6})이므로 준콤팩트이고, $L$이 유한하므로
$V_\lambda\cap W_\lambda$는 준콤팩트이다. 이로써 정리가 증명되었다.
\end{proof}

\begin{corollary}[9.4.8]
\label{I.9.4.8-ko}
(9.4.7)의 가정 아래에서 모든 유한 생성 준연접
$(\mathscr{O}_X|U)$-가군 $\mathscr{G}$에 대하여 유한 생성 준연접
$\mathscr{O}_X$-가군 $\mathscr{G}'$이 존재하여
$\mathscr{G}'|U=\mathscr{G}$이다.
\end{corollary}

\begin{proof}
$\mathscr{F}=j_*(\mathscr{G})$에 (9.4.7)을 적용하면 된다. 이 가군은
(9.4.2)에 따라 준연접이고 $\mathscr{F}|U=\mathscr{G}$이다.
\end{proof}

\begin{corollary}[9.4.9]
\label{I.9.4.9-ko}
$X$를 바탕공간이 국소 뇌터인 준스킴 또는 준콤팩트 스킴이라 하자.
그러면 모든 준연접 $\mathscr{O}_X$-가군은 그 유한 생성 준연접
부분-$\mathscr{O}_X$-가군들의 귀납극한이다.
\end{corollary}

\begin{proof}
이는 (9.4.7)과 주 (9.4.4)에서 따른다.
\end{proof}

\begin{corollary}[9.4.10]
\label{I.9.4.10-ko}
(9.4.9)의 가정 아래에서 준연접 $\mathscr{O}_X$-가군
$\mathscr{F}$를 생각하자. $\mathscr{F}$의 모든 유한 생성 준연접
부분-$\mathscr{O}_X$-가군이 $X$ 위의 단면들로 생성된다고 가정하면,
$\mathscr{F}$도 $X$ 위의 단면들로 생성된다.
\end{corollary}

\begin{proof}
$U$를 점 $x\in X$의 아핀 열린 근방, $s$를 $U$ 위의
$\mathscr{F}$의 단면이라 하자. $s$로 생성되는 $\mathscr{F}|U$의
부분-$(\mathscr{O}_X|U)$-가군 $\mathscr{G}$는 준연접이고 유한 생성이므로,
$\mathscr{F}$의 유한 생성 준연접 부분-$\mathscr{O}_X$-가군
$\mathscr{G}'$이 존재하여 $\mathscr{G}'|U=\mathscr{G}$이다(9.4.7).
가정에 따라 $X$ 위의 $\mathscr{G}'$의 유한개의 단면 $t_i$와
$x$의 근방 $V\subset U$ 위의 $\mathscr{O}_X$의 단면 $a_i$들이 존재하여
$s|V=\sum_i a_i\mathbin{\cdot}(t_i|V)$이다. 이로써 따름정리가
증명되었다.
\end{proof}

\subsection{준스킴의 닫힌 상. 부분준스킴의 폐포.}
\label{subsection:I.9.5-ko}

\begin{proposition}[9.5.1]
\label{I.9.5.1-ko}
$f:X\to Y$를 $f_*(\mathscr{O}_X)$가 준연접
$\mathscr{O}_Y$-가군인 준스킴의 사상이라 하자(이는 $f$가 준콤팩트이고,
또한 $f$가 분리 사상이거나 $X$가 국소 뇌터이면 성립한다
(\hyperref[I.9.2.2-ko]{9.2.2})). 그러면 $Y$의 가장 작은
닫힌 부분준스킴\footnote{\textbf{번역 주.} 인쇄 원문은 ‘닫힌’ 없이
단순히 “가장 작은 부분준스킴”이라고 쓴다.
그러나 이어지는 (9.5.2)의 구성과 (9.5.3)의 정의는 닫힌
부분준스킴 가운데의 최소성을 증명하고 정의한다. 실제로 조밀한
준콤팩트 열린 몰입 $D(t)\to\operatorname{Spec}(k[t])$에서는 핵으로
정의되는 닫힌 부분준스킴이 $\operatorname{Spec}(k[t])$ 전체인 반면,
$f$는 진부분 열린 부분준스킴 $D(t)$를 통하여 인수분해된다. 따라서
핵으로 정의된 대상이 모든 부분준스킴 가운데 최소라는 인쇄문과
(9.5.2)의 결합된 주장은 성립하지 않으며, 여기서는 논증과 정의에
필수적인 ‘닫힌’을 보충했다. 이는 공식 정오표가 아니라 뒤의
구성·정의와 위 반례가 강제하는 편집 교정이다.}
$Y'$이 존재하여, 표준 포함 사상 $j:Y'\to Y$를 통하여 $f$가
인수분해된다(또는 동치로
(\hyperref[I.4.4.1-ko]{4.4.1}), $X$의 부분준스킴
$f^{-1}(Y')$이 $X$와 \emph{동일하다}).
\end{proposition}

좀 더 정확히 말하면 다음과 같다.

\begin{corollary}[9.5.2]
\label{I.9.5.2-ko}
(9.5.1)의 가정 아래에서 $f=(\psi,\theta)$라 하고, 준동형
$\theta:\mathscr{O}_Y\to f_*(\mathscr{O}_X)$의 준연접 핵을
$\mathscr{J}$라 하자. 그러면 $\mathscr{J}$로 정의된 $Y$의 닫힌
부분준스킴 $Y'$은 (9.5.1)의 조건을 만족한다.
\end{corollary}

\begin{proof}
함자 $\psi^*$가 완전하므로, 표준 인수분해
$\theta:\mathscr{O}_Y\to\mathscr{O}_Y/\mathscr{J}
\xrightarrow{\theta'}\psi_*(\mathscr{O}_X)$는
(\hyperref[0.3.5.4.3-ko]{0, 3.5.4.3})에 따라 다음 인수분해를 준다.
\[
  \theta^\sharp:\psi^*(\mathscr{O}_Y)
  \to\psi^*(\mathscr{O}_Y)/\psi^*(\mathscr{J})
  \xrightarrow{{\theta'}^\sharp}\mathscr{O}_X~;
\]
모든 $x\in X$에 대하여 $\theta_x^\sharp$가 국소 준동형이므로
${\theta'_x}^\sharp$도 그러하다. 연속사상 $\psi$를 $X$에서
$Y'$로 가는 사상으로 보아 $\psi_0$라 하고,\footnote{\textbf{번역 주.}
인쇄 원문은 이 증명의 해당 구간에서
$X'$, $X''$를 쓰지만, 출판된 EGA II 정오표는 이들을 전부
$Y'$, $Y''$로 고치도록 지시한다. 여기서는 그 공식 교정을
적용하였다.}
$\theta_0$를 제한 사상
$\theta'|Y':(\mathscr{O}_Y/\mathscr{J})|Y'
\to\psi_*(\mathscr{O}_X)|Y'=(\psi_0)_*(\mathscr{O}_X)$라 하자.
그러면 $f_0=(\psi_0,\theta_0)$는 준스킴의 사상 $X\to Y'$이고
(\hyperref[I.2.2.1-ko]{2.2.1}), $f=j\circ f_0$이다. 이제
$\mathscr{O}_Y$의 준연접 아이디얼층 $\mathscr{J}'$로 정의된
$Y$의 두 번째 닫힌 부분준스킴을 $Y''$
\oldpage[I]{177}
라 하고, 표준 포함 사상 $j':Y''\to Y$를 통하여 $f$가
인수분해된다고 하자. 우선 $Y''\supset\psi(X)$이어야 하므로,
$Y''$가 닫혀 있다는 데서 $Y'\subset Y''$가 따른다. 또 모든
$y\in Y''$에 대하여 $\theta$는
$\mathscr{O}_y\to\mathscr{O}_y/\mathscr{J}'_y
\to(\psi_*(\mathscr{O}_X))_y$로 인수분해되어야 한다.
따라서 정의에 의해 $\mathscr{J}'_y\subset\mathscr{J}_y$이고,
그러므로 $Y'$은 $Y''$의 닫힌 부분준스킴이다
(\hyperref[I.4.1.10-ko]{4.1.10}).
\end{proof}

\begin{definition}[9.5.3]
\label{I.9.5.3-ko}
표준 포함 사상 $j:Y'\to Y$를 통하여 $f$가 인수분해되는 $Y$의
가장 작은 닫힌 부분준스킴 $Y'$이 존재할 때, $Y'$을 사상 $f$에
의한 $X$의 닫힌 상 준스킴이라고 한다.
\end{definition}

\begin{proposition}[9.5.4]
\label{I.9.5.4-ko}
$f_*(\mathscr{O}_X)$가 준연접 $\mathscr{O}_Y$-가군이면, $f$에 의한
$X$의 닫힌 상의 바탕공간은 $Y$에서 집합론적 상의 폐포
$\overline{f(X)}$이다.
\end{proposition}

\begin{proof}
$f_*(\mathscr{O}_X)$의 지지집합은 $\overline{f(X)}$에 포함되므로,
(\hyperref[I.9.5.2-ko]{9.5.2})의 표기에서
$y\notin\overline{f(X)}$이면 $\mathscr{J}_y=\mathscr{O}_y$이다.
따라서 $\mathscr{O}_Y/\mathscr{J}$의 지지집합은
$\overline{f(X)}$에 포함된다. 한편 이 지지집합은 닫혀 있고
$f(X)$를 포함한다. 실제로 $y\in f(X)$이면 환
$(\psi_*(\mathscr{O}_X))_y$의 단위원은 영이 아닌데, 이는 단면
\[
  1\in\Gamma(X,\mathscr{O}_X)=\Gamma(Y,\psi_*(\mathscr{O}_X))~;
\]
의 $y$에서의 싹이기 때문이다. 이 단위원은 $\mathscr{O}_y$의
단위원의 $\theta$에 의한 상이므로, 후자는 $\mathscr{J}_y$에 속하지
않는다. 따라서 $\mathscr{O}_y/\mathscr{J}_y\neq0$이고, 이로써 증명이
끝난다.
\end{proof}

\begin{proposition}[9.5.5]
\label{I.9.5.5-ko}
\emph{(닫힌 상의 추이성).} $f:X\to Y$와 $g:Y\to Z$를 준스킴의
두 사상이라 하자. $f$에 의한 $X$의 닫힌 상 $Y'$이 존재하고,
$g'$이 $g$를 $Y'$에 제한한 사상일 때 $g'$에 의한 $Y'$의 닫힌 상
$Z'$이 존재한다고 가정하자. 그러면 $g\circ f$에 의한 $X$의 닫힌
상이 존재하며 $Z'$과 같다.
\end{proposition}

\begin{proof}
(\hyperref[I.9.5.1-ko]{9.5.1})에 따라, $Z'$이 다음 조건을 만족하는
$Z$의 가장 작은 닫힌 부분준스킴 $Z_1$임을 보이면 충분하다. 그
조건은 $X$의 닫힌 부분준스킴 $(g\circ f)^{-1}(Z_1)$이 $X$와 같다는
것이며, 이 부분준스킴은 (\hyperref[I.4.4.2-ko]{4.4.2})에 따라
$f^{-1}(g^{-1}(Z_1))$과 같다. 이는 $Z'$이 $f$가 포함 사상
$g^{-1}(Z_1)\to Y$를 통하여 인수분해되게 하는 $Z$의 가장 작은
닫힌 부분준스킴이라는 것과 동치이다
(\hyperref[I.4.4.1-ko]{4.4.1}). 그런데 닫힌 상 $Y'$의 존재에 따라,
이 조건을 만족하는 모든 $Z_1$에 대하여 $g^{-1}(Z_1)$은 $Y'$을
부분준스킴으로 포함한다. $j$를 포함 사상 $Y'\to Y$라 하면, 이는
$j^{-1}(g^{-1}(Z_1))=g'^{-1}(Z_1)=Y'$이라는 것과 동치이다. 이제
$Z'$의 정의로부터 $Z'$이 앞의 조건을 만족하는 $Z$의 가장 작은
닫힌 부분준스킴임이 따른다.
\end{proof}

\begin{corollary}[9.5.6]
\label{I.9.5.6-ko}
$f:X\to Y$를 $f$에 의한 $X$의 닫힌 상이 $Y$인 $S$-사상이라
하자. $Z$를 $S$-스킴이라 하자. $Y$에서 $Z$로 가는 두 $S$-사상
$g_1$, $g_2$가 $g_1\circ f=g_2\circ f$를 만족하면
$g_1=g_2$이다.
\end{corollary}

\begin{proof}
$h=(g_1,g_2)_S:Y\to Z\times_S Z$라 하자. 대각선
$T=\Delta_Z(Z)$은 $Z\times_S Z$의 닫힌 부분준스킴이므로,
$Y'=h^{-1}(T)$는 $Y$의 닫힌 부분준스킴이다
(\hyperref[I.4.4.1-ko]{4.4.1}). $u=g_1\circ f=g_2\circ f$라 두면,
곱의 정의에 따라 $h'=h\circ f=(u,u)_S$이고, 따라서
$h\circ f=\Delta_Z\circ u$이다. $\Delta_Z^{-1}(T)=Z$이므로
$h'^{-1}(T)=u^{-1}(Z)=X$이고, 따라서 $f^{-1}(Y')=X$이다. 이로부터
(\hyperref[I.4.4.1-ko]{4.4.1})에 따라 표준 포함 사상 $Y'\to Y$를 통하여
$f$가 인수분해되므로, 가정에 따라 $Y'=Y$이다. 따라서
(\hyperref[I.4.4.1-ko]{4.4.1})에 따라 $h$는 $\Delta_Z\circ v$로
인수분해된다. 여기서 $v$는 $Y\to Z$인 사상이며, 이로부터
$g_1=g_2=v$가 따른다.
\end{proof}

\begin{remark}[9.5.7]
\label{I.9.5.7-ko}
$X$와 $Y$가 $S$-스킴이면, 명제
(\hyperref[I.9.5.6-ko]{9.5.6})이 뜻하는 바는
\oldpage[I]{178}
$f$에 의한 $X$의 닫힌 상이 $Y$일 때 $f$가 $S$-스킴의 범주에서
\emph{에피사상}이라는 것이다(T, 1.1). 제V장에서 역으로, $f$에
의한 $X$의 닫힌 상 $Y'$이 존재하고 $f$가 $S$-스킴의 에피사상이면
반드시 $Y'=Y$임을 증명할 것이다.
\end{remark}

\begin{proposition}[9.5.8]
\label{I.9.5.8-ko}
(\hyperref[I.9.5.1-ko]{9.5.1})의 가정들이 성립하고, $Y'$을 $f$에
의한 $X$의 닫힌 상이라 하자. $Y$의 임의의 열린집합 $V$에 대하여,
$f_V:f^{-1}(V)\to V$를 $f$의 제한이라 하자. 그러면 $f_V$에 의한
$f^{-1}(V)$의 $V$ 안에서의 닫힌 상은 존재하며, $Y'$의 열린집합
$V\cap Y'$ 위에 $Y'$로부터 유도된 준스킴과 같다. 다시 말해 이는
$Y$의 부분준스킴 $\inf(V,Y')$와 같다
(\hyperref[I.4.4.3-ko]{4.4.3}).
\end{proposition}

\begin{proof}
$X'=f^{-1}(V)$라 하자. $f_V$에 의한 $\mathscr{O}_{X'}$의 직상은
$f_*(\mathscr{O}_X)$를 $V$에 제한한 것과 같으므로, 준동형
$\mathscr{O}_V\to(f_V)_*(\mathscr{O}_{X'})$의 핵
$\mathscr{J}'$은 $\mathscr{J}$를 $V$에 제한한 것이다. 이로부터
명제가 곧바로 따른다.
\end{proof}

이 결과는 닫힌 상을 만드는 것이 밑준스킴의 확장 $Y_1\to Y$가
\emph{열린 몰입 사상}인 경우 그 확장과 가환한다는 뜻으로 해석할
수 있다. 제IV장에서 $f$가 분리 사상이고 준콤팩트일 때에는
$Y_1\to Y$가 \emph{평탄} 사상인 경우에도 마찬가지임을 볼 것이다.

\begin{proposition}[9.5.9]
\label{I.9.5.9-ko}
$f:X\to Y$를 $f$에 의한 $X$의 닫힌 상 $Y'$이 존재하는 사상이라
하자.
\begin{enumerate}
  \item[(i)] $X$가 축소 준스킴이면 $Y'$도 축소 준스킴이다.
  \item[(ii)] (\hyperref[I.9.5.1-ko]{9.5.1})의 가정들이 성립하고
    $X$가 기약이면(각각 정역 준스킴이면), $Y'$도 기약이다
    (각각 정역 준스킴이다).
\end{enumerate}
\end{proposition}

\begin{proof}
가정에 따라 $f$는 $X\xrightarrow{g}Y'\xrightarrow{j}Y$로
인수분해되며, 여기서 $j$는 표준 포함 사상이다. $X$가 축소이므로
$g$는 $X\xrightarrow{h}Y'_{\mathrm{red}}\xrightarrow{j'}Y'$로
인수분해된다. 여기서 $j'$은 표준 포함 사상이다
(\hyperref[I.5.2.2-ko]{5.2.2}). 따라서 $Y'$의 정의로부터
$Y'_{\mathrm{red}}=Y'$이 따른다. 또한
(\hyperref[I.9.5.1-ko]{9.5.1})의 조건들이 성립하면,
(\hyperref[I.9.5.4-ko]{9.5.4})에 따라 $f(X)$는 $Y'$에서 조밀하다.
그러므로 $X$가 기약이면 $Y'$도 기약이다
(\hyperref[0.2.1.5-ko]{0, 2.1.5}). 정역 준스킴에 관한 주장은
앞의 두 주장을 함께 적용하여 얻는다.
\end{proof}

\begin{proposition}[9.5.10]
\label{I.9.5.10-ko}
$Y$를 준스킴 $X$의 부분준스킴이라 하고, 표준 포함 사상
$i:Y\to X$가 준콤팩트 사상이라고 하자. 그러면 $Y$를 부분준스킴으로
포함하는 $X$의 가장 작은 닫힌 부분준스킴 $\overline{Y}$이 존재한다.
그 바탕공간은 $Y$의 바탕공간의 폐포이고, 후자는 이
폐포에서 열려 있으며, 이 열린집합 위에서 $\overline{Y}$로부터
유도된 준스킴은 $Y$이다.
\end{proposition}

\begin{proof}
포함 사상 $i$에 (\hyperref[I.9.5.1-ko]{9.5.1})을 적용하면
된다.\footnote{\textbf{번역 주.} 인쇄 원문은 이 증명에서 앞서 정의되지
않은 $j$에 (9.5.1)을 적용한다고 쓰지만, 명제에서 정의된 표준 포함
사상은 $i:Y\to X$이다. 여기서는 명제에서 정의된 표기에 맞게 $i$로
바로잡았다.} 이 사상은
분리 사상이고 (\hyperref[I.5.5.1-ko]{5.5.1}), 가정에 따라
준콤팩트이다. 따라서 (\hyperref[I.9.5.1-ko]{9.5.1})은
$\overline{Y}$의 존재를 보이고, (\hyperref[I.9.5.4-ko]{9.5.4})는
그 바탕공간이 $X$에서 $Y$의 폐포임을 보인다. $Y$는 $X$에서 국소
닫혀 있으므로 $\overline{Y}$에서 열려 있다. 마지막 주장은 $Y$가
$X$의 어떤 열린집합 $V$ 안에서 닫혀 있도록 $V$를 택한 뒤
(\hyperref[I.9.5.8-ko]{9.5.8})을 적용하면 얻는다.
\end{proof}

이 표기 아래에서, 포함 사상 $V\to X$가 준콤팩트이고
$\mathscr{J}$가 $V$의 닫힌 부분준스킴 $Y$를 정의하는
$\mathscr{O}_X|V$의 준연접 아이디얼층이면,
(\hyperref[I.9.5.1-ko]{9.5.1})에 따라 $\overline{Y}$을 정의하는
$\mathscr{O}_X$의 준연접 아이디얼층은 $\mathscr{J}$의 표준 연장
(\hyperref[I.9.4.1-ko]{9.4.1}) $\overline{\mathscr{J}}$이다. 실제로
이는 $V$ 위에 $\mathscr{J}$를 유도하는 $\mathscr{O}_X$의 준연접
아이디얼 부분층 가운데 가장 큰 것이다.
\oldpage[I]{179}
\begin{corollary}[9.5.11]
\label{I.9.5.11-ko}
(\hyperref[I.9.5.10-ko]{9.5.10})의 가정 아래에서,
$\overline{Y}$의 열린집합 $V$ 위의 $\mathscr{O}_{\overline{Y}}$의
단면이 $V\cap Y$에서 영이면 그 단면은 영이다.
\end{corollary}

\begin{proof}
(\hyperref[I.9.5.8-ko]{9.5.8})에 따라 $V=\overline{Y}$인 경우로
환원할 수 있다. $\overline{Y}$ 위의 $\mathscr{O}_{\overline{Y}}$의
단면들이 준스킴
$\overline{Y}\otimes_{\mathbb{Z}}\mathbb{Z}[T]$의
$\overline{Y}$-단면들과 표준적으로 대응하고
(\hyperref[I.3.3.15-ko]{3.3.15}), 이 준스킴은 $\overline{Y}$ 위에서
분리되어 있으므로, 이 따름정리는
(\hyperref[I.9.5.6-ko]{9.5.6})의 특별한 경우이다.
\end{proof}

$X$의 부분준스킴 $Y$를 부분준스킴으로 포함하는 $X$의 가장 작은 닫힌
부분준스킴 $\overline{Y}$이 존재할 때, 혼동할 우려가 없으면
$\overline{Y}$을 $X$에서 $Y$의 \emph{폐포}라고 한다.

\subsection{준연접 대수; 구조층의 변경.}
\label{subsection:I.9.6-ko}

\begin{proposition}[9.6.1]
\label{I.9.6.1-ko}
$X$를 준스킴, $\mathscr{B}$를 준연접 $\mathscr{O}_X$-대수라 하자
(\hyperref[0.5.1.3-ko]{0, 5.1.3}). $\mathscr{B}$-가군 $\mathscr{F}$가
(환 달린 공간 $(X,\mathscr{B})$ 위에서) 준연접일 필요충분조건은
$\mathscr{F}$가 준연접 $\mathscr{O}_X$-가군인 것이다.
\end{proposition}

\begin{proof}
문제가 국소적이므로 $X$가 환 $A$를 갖는 아핀 스킴이라고 해도 된다.
그러면 $\mathscr{B}=\widetilde{B}$이고, 여기서 $B$는 $A$-대수이다
(\hyperref[I.1.4.3-ko]{1.4.3}). $\mathscr{F}$가 환 달린 공간
$(X,\mathscr{B})$ 위에서 준연접이면, $\mathscr{F}$가 어떤
$\mathscr{B}$-준동형 $\mathscr{B}^{(I)}\to\mathscr{B}^{(J)}$의
여핵이라고 해도 된다. 이 준동형은 $\mathscr{O}_X$-가군의
$\mathscr{O}_X$-준동형이기도 하고, $\mathscr{B}^{(I)}$와
$\mathscr{B}^{(J)}$는 준연접 $\mathscr{O}_X$-가군이므로
(\hyperref[I.1.3.9-ko]{1.3.9, (ii)}), $\mathscr{F}$도 준연접
$\mathscr{O}_X$-가군이다
(\hyperref[I.1.3.9-ko]{1.3.9, (i)}).

거꾸로 $\mathscr{F}$가 준연접 $\mathscr{O}_X$-가군이면
$\mathscr{F}=\widetilde{M}$이고, 여기서 $M$은 $B$-가군이다
(\hyperref[I.1.4.3-ko]{1.4.3}). $M$은 어떤 $B$-준동형
$B^{(I)}\to B^{(J)}$의 여핵과 동형이므로, $\mathscr{F}$는 이에
대응하는 준동형 $\mathscr{B}^{(I)}\to\mathscr{B}^{(J)}$의 여핵과
동형인 $\mathscr{B}$-가군이다
(\hyperref[I.1.3.13-ko]{1.3.13}). 이로써 증명이 끝난다.
\end{proof}

특히 $\mathscr{F}$와 $\mathscr{G}$가 두 준연접
$\mathscr{B}$-가군이면 $\mathscr{F}\otimes_{\mathscr{B}}\mathscr{G}$는
준연접 $\mathscr{B}$-가군이다. 더욱이 $\mathscr{F}$가 유한 표시를
갖는다고 가정하면 $\mathscr{H}\!om_{\mathscr{B}}(\mathscr{F},\mathscr{G})$도
그러하다(\hyperref[I.1.3.13-ko]{1.3.13}).

\begin{env}[9.6.2]
\label{I.9.6.2-ko}
준스킴 $X$가 주어졌을 때, 준연접 $\mathscr{O}_X$-대수
$\mathscr{B}$가 \emph{유한 생성}이라는 것은 모든 $x\in X$에 대하여
$x$의 \emph{아핀} 열린 근방 $U$가 존재하여
$\Gamma(U,\mathscr{B})=B$가 $\Gamma(U,\mathscr{O}_X)=A$ 위의 유한 생성
대수라는 뜻이다. 이때 $\mathscr{B}|U=\widetilde{B}$이고, 모든 $f\in A$에
대하여 그 원소에 대응하는 주 열린집합 위에 유도된
$(\mathscr{O}_X|D(f))$-대수
$\mathscr{B}|D(f)$도 유한 생성이다. 실제로 이는
$\widetilde{B_f}$와 동형이고, $B_f=B\otimes_A A_f$는 분명히 $A_f$ 위의
유한 생성 대수이다. $D(f)$들이 $U$의 위상의 기저를 이루므로, 준연접
유한 생성 $\mathscr{O}_X$-대수 $\mathscr{B}$와 $X$의 모든 열린집합
$V$에 대하여 $\mathscr{B}|V$는 준연접 유한 생성
$(\mathscr{O}_X|V)$-대수이다.
\end{env}

\begin{proposition}[9.6.3]
\label{I.9.6.3-ko}
$X$를 국소 뇌터 준스킴이라 하자. 모든 준연접 유한 생성
$\mathscr{O}_X$-대수 $\mathscr{B}$는 연접인 환의 층이다
(\hyperref[0.5.3.7-ko]{0, 5.3.7}).
\end{proposition}

\begin{proof}
다시 $X$가 뇌터 환 $A$를 갖는 아핀 스킴이고
$\mathscr{B}=\widetilde{B}$이며 $B$가 $A$ 위의 유한 생성 대수인 경우로
한정해도 된다. 그러면 $B$는 뇌터 환이다. 따라서
$\mathscr{B}$-준동형 $\mathscr{B}^m\to\mathscr{B}$의 핵
$\mathscr{N}$이 유한 생성 $\mathscr{B}$-
\oldpage[I]{180}
가군임을 보이면 된다. 그런데 이 핵은 이에 대응하는 $B$-가군의
준동형 $B^m\to B$의 핵 $N$에 대응하는 층 $\widetilde{N}$과
($\mathscr{B}$-가군으로서) 동형이다
(\hyperref[I.1.3.13-ko]{1.3.13}). $B$가 뇌터이므로 $B^m$의 부분가군
$N$은 유한 생성 $B$-가군이다. 따라서 상이 $N$인 준동형
$B^p\to B^m$이 존재한다. 열 $B^p\to B^m\to B$이 완전이므로 이에
대응하는 열 $\mathscr{B}^p\to\mathscr{B}^m\to\mathscr{B}$도 완전이고
(\hyperref[I.1.3.5-ko]{1.3.5}), $\mathscr{N}$은
$\mathscr{B}^p\to\mathscr{B}^m$의 상이므로
(\hyperref[I.1.3.9-ko]{1.3.9, (i)}) 명제가 증명된다.
\end{proof}

\begin{corollary}[9.6.4]
\label{I.9.6.4-ko}
(\hyperref[I.9.6.3-ko]{9.6.3})의 가정 아래에서, $\mathscr{B}$-가군
$\mathscr{F}$가 연접일 필요충분조건은 그 가군이 준연접
$\mathscr{O}_X$-가군이면서 유한 생성 $\mathscr{B}$-가군인 것이다.
이 조건들이 성립하고 $\mathscr{G}$가 $\mathscr{F}$의
부분-$\mathscr{B}$-가군 또는 $\mathscr{B}$-몫가군이면,
$\mathscr{G}$가 연접 $\mathscr{B}$-가군일 필요충분조건은
$\mathscr{G}$가 준연접 $\mathscr{O}_X$-가군인 것이다.
\end{corollary}

\begin{proof}
(\hyperref[I.9.6.1-ko]{9.6.1})을 고려하면 $\mathscr{F}$에 관한 조건들이
필요함은 분명하다. 충분함을 보이자. $X$가 뇌터 환 $A$를 갖는
아핀 스킴이고, $\mathscr{B}=\widetilde{B}$이며, $B$가 $A$ 위의 유한 생성
대수이고, $\mathscr{F}=\widetilde{M}$이며, 여기서 $M$은 $B$-가군이고,
전사 $\mathscr{B}$-준동형 $\mathscr{B}^m\to\mathscr{F}\to 0$이 존재하는
경우로 한정해도 된다. 그러면 이에 대응하는 완전열
$B^m\to M\to 0$이 있으므로 $M$은 유한 생성 $B$-가군이다. 또한
$B^m\to M$의 핵 $P$도 $B$가 뇌터이므로 유한 생성 $B$-가군이다.
따라서 (\hyperref[I.1.3.13-ko]{1.3.13})에 의해 $\mathscr{F}$는 어떤
$\mathscr{B}$-준동형 $\mathscr{B}^n\to\mathscr{B}^m$의 여핵이고,
$\mathscr{B}$가 연접인 환의 층이므로 연접이다
(\hyperref[0.5.3.4-ko]{0, 5.3.4}). 같은 논증으로 $\mathscr{F}$의 준연접
부분-$\mathscr{B}$-가군(각각 $\mathscr{B}$-몫가군)은 유한 생성임을
알 수 있으며, 이로부터 따름정리의 두 번째 주장이 따른다.
\end{proof}

\begin{proposition}[9.6.5]
\label{I.9.6.5-ko}
$X$를 준콤팩트 스킴 또는 바탕공간이 뇌터인 준스킴이라 하자.
모든 준연접 유한 생성 $\mathscr{O}_X$-대수 $\mathscr{B}$에 대하여,
$\mathscr{B}$의 유한 생성 준연접 부분-$\mathscr{O}_X$-가군
$\mathscr{E}$가 존재하여 $\mathscr{E}$가 $\mathscr{O}_X$-대수
$\mathscr{B}$를 생성한다(\hyperref[0.4.1.4-ko]{0, 4.1.4}).
\end{proposition}

\begin{proof}
실제로 가정에 따라 $X$는 아핀 열린집합 $(U_i)$의 유한 덮개를 가지며,
$\Gamma(U_i,\mathscr{B})=B_i$는 $\Gamma(U_i,\mathscr{O}_X)=A_i$ 위의
유한 생성 대수이다. $B_i$의 유한 생성 부분-$A_i$-가군 $E_i$를
$A_i$-대수 $B_i$를 생성하도록 잡자. (\hyperref[I.9.4.7-ko]{9.4.7})에
따르면, $\mathscr{B}$의 유한 생성 준연접 부분-$\mathscr{O}_X$-가군
$\mathscr{E}_i$가 존재하여 $\mathscr{E}_i|U_i=\widetilde{E_i}$이다.
$\mathscr{E}_i$들의 합 $\mathscr{E}$가 요구된 조건을 만족함은 분명하다.
\end{proof}

\begin{proposition}[9.6.6]
\label{I.9.6.6-ko}
$X$를 바탕공간이 국소 뇌터인 준스킴 또는 준콤팩트 스킴이라 하자.
그러면 모든 준연접 $\mathscr{O}_X$-대수 $\mathscr{B}$는 그 준연접 유한
생성 부분-$\mathscr{O}_X$-대수들의 귀납극한이다.
\end{proposition}

\begin{proof}
(\hyperref[I.9.4.9-ko]{9.4.9})에 따라 $\mathscr{B}$는
$\mathscr{O}_X$-가군으로서 그 유한 생성 준연접
부분-$\mathscr{O}_X$-가군들의 귀납극한이다. 이 부분가군들은
$\mathscr{B}$의 준연접 유한 생성 부분-$\mathscr{O}_X$-대수들을 생성하고
(\hyperref[I.1.3.14-ko]{1.3.14}), 따라서 $\mathscr{B}$는
\emph{특히} 이 부분대수들의 귀납극한이다.
\end{proof}
% ===== END EXACT INLINE INPUT: c1s9.tex =====


% ===== BEGIN EXACT INLINE INPUT: c1s10.tex =====
\section{형식적 스킴}
\label{section:I.10-ko}

\subsection{아핀 형식적 스킴.}
\label{subsection:I.10.1-ko}

\begin{env}[10.1.1]
\label{I.10.1.1-ko}
$A$를 위상환인 \emph{허용환}이라 하자
(\hyperref[0.7.1.2-ko]{0, 7.1.2}). $A$의 모든 정의 아이디얼
$\mathfrak{J}$에 대하여 $\operatorname{Spec}(A/\mathfrak{J})$는
$\operatorname{Spec}(A)$의 닫힌 부분공간 $V(\mathfrak{J})$와
동일시된다(\hyperref[I.1.1.11-ko]{1.1.11}). 이는 $A$의
\emph{열린} 소 아이디얼들의 집합이다. 이 위상공간은 고려한
\oldpage[I]{181}
정의 아이디얼 $\mathfrak{J}$의 선택과 무관하며, 이를
$\mathfrak{X}$로 나타내자. $(\mathfrak{J}_\lambda)$를 정의
아이디얼들로 이루어진 $A$에서 0의 근방들의 기본계라 하고, 모든
$\lambda$에 대하여 $\mathscr{O}_\lambda$를
$\operatorname{Spec}(A/\mathfrak{J}_\lambda)$의 구조층이라 하자. 이 층은
$\widetilde{A}/\widetilde{\mathfrak{J}}_\lambda$에 의해
$\mathfrak{X}$ 위에 유도된다(이 몫층은 $\mathfrak{X}$ 밖에서 영이다).

$\mathfrak{J}_\mu\subset\mathfrak{J}_\lambda$이면 표준 준동형
$A/\mathfrak{J}_\mu\to A/\mathfrak{J}_\lambda$는 환의 층의 준동형
$u_{\lambda\mu}:\mathscr{O}_\mu\to\mathscr{O}_\lambda$를 정의하므로
(\hyperref[I.1.6.1-ko]{1.6.1}), $(\mathscr{O}_\lambda)$는 이 준동형들에
관하여 \emph{환의 층들의 사영계}를 이룬다. $\mathfrak{X}$의 위상에는
준콤팩트 열린집합들로 이루어진 기저가 있으므로, 각
$\mathscr{O}_\lambda$에 그 바탕 환의 층(위상을 잊은 층)이
$\mathscr{O}_\lambda$인 \emph{의사이산 위상환의 층}을 대응시킬 수
있다(\hyperref[0.3.8.1-ko]{0, 3.8.1}). 이 층도 여전히
$\mathscr{O}_\lambda$로 나타내겠다. 또한 $\mathscr{O}_\lambda$들은
여전히 \emph{위상환의 층들의 사영계}를 이룬다
(\hyperref[0.3.8.2-ko]{0, 3.8.2}). 이 계
$(\mathscr{O}_\lambda)$의 사영극한인 $\mathfrak{X}$ 위의
\emph{위상환의 층}을 $\mathscr{O}_{\mathfrak{X}}$로 나타내겠다. 따라서
$\mathfrak{X}$의 모든 \emph{준콤팩트} 열린집합 $U$에 대하여
$\Gamma(U,\mathscr{O}_{\mathfrak{X}})$는 \emph{이산}환들의 계
$\Gamma(U,\mathscr{O}_\lambda)$의 사영극한인 위상환이다
(\hyperref[0.3.2.6-ko]{0, 3.2.6}).
\end{env}

\begin{definition}[10.1.2]
\label{I.10.1.2-ko}
허용환 $A$가 주어졌을 때, $A$의 열린 소 아이디얼들로 이루어진
$\operatorname{Spec}(A)$의 닫힌 부분공간 $\mathfrak{X}$를 $A$의
형식적 스펙트럼이라 하고 $\operatorname{Spf}(A)$로 나타낸다. 위상환 달린
공간이 형식적 스펙트럼 $\operatorname{Spf}(A)=\mathfrak{X}$에 다음의
구조층을 갖춘 것과 동형이면 이를 아핀 형식적 스킴이라고 한다. 여기서
$\mathfrak{J}_\lambda$가 $A$의 정의 아이디얼들의 여과집합을 달릴 때
$\mathscr{O}_{\mathfrak{X}}$는
의사이산 위상환의 층
$(\widetilde{A}/\widetilde{\mathfrak{J}}_\lambda)|\mathfrak{X}$들의
사영극한이다.
\end{definition}

\emph{형식적 스펙트럼} $\mathfrak{X}=\operatorname{Spf}(A)$를 아핀 형식적
스킴으로 말할 때에는 언제나 위에서 정의한
$\mathscr{O}_{\mathfrak{X}}$를 갖는 위상환 달린 공간
$(\mathfrak{X},\mathscr{O}_{\mathfrak{X}})$을 뜻한다.

모든 \emph{아핀 스킴} $X=\operatorname{Spec}(A)$는 $A$를 이산
위상환으로 보아 유일한 방식으로 아핀 형식적 스킴으로 볼 수 있다. 이때
$U$가 준콤팩트이면 위상환 $\Gamma(U,\mathscr{O}_X)$는 이산이지만,
$U$가 $X$의 임의의 열린집합일 때에는 일반적으로 그렇지 않다.

\begin{proposition}[10.1.3]
\label{I.10.1.3-ko}
$A$가 허용환이고 $\mathfrak{X}=\operatorname{Spf}(A)$이면,
$\Gamma(\mathfrak{X},\mathscr{O}_{\mathfrak{X}})$는 $A$와 위상적으로
동형이다.
\end{proposition}

\begin{proof}
실제로 $\mathfrak{X}$는 $\operatorname{Spec}(A)$에서 닫혀 있으므로
준콤팩트이다. 따라서
$\Gamma(\mathfrak{X},\mathscr{O}_{\mathfrak{X}})$는 이산환
$\Gamma(\mathfrak{X},\mathscr{O}_\lambda)$들의 사영극한과 위상적으로
동형이다. 그런데 $\Gamma(\mathfrak{X},\mathscr{O}_\lambda)$는
$A/\mathfrak{J}_\lambda$와 동형이다
(\hyperref[I.1.3.7-ko]{1.3.7}). $A$는 분리이고 완비이므로
$\varprojlim A/\mathfrak{J}_\lambda$와 위상적으로 동형이다
(\hyperref[0.7.2.1-ko]{0, 7.2.1}). 이로부터 명제가 따른다.
\end{proof}

\begin{proposition}[10.1.4]
\label{I.10.1.4-ko}
$A$를 허용환, $\mathfrak{X}=\operatorname{Spf}(A)$라 하고, 모든
$f\in A$에 대하여 $\mathfrak{D}(f)=D(f)\cap\mathfrak{X}$라 하자. 그러면
위상환 달린 공간
$(\mathfrak{D}(f),\mathscr{O}_{\mathfrak{X}}|\mathfrak{D}(f))$는 아핀
형식적 스펙트럼 $\operatorname{Spf}(A_{\{f\}})$와 동형이다
(\hyperref[0.7.6.15-ko]{0, 7.6.15}).
\end{proposition}

\begin{proof}
$A$의 모든 정의 아이디얼 $\mathfrak{J}$에 대하여 이산환
$S_f^{-1}A/S_f^{-1}\mathfrak{J}$는
$A_{\{f\}}/\mathfrak{J}_{\{f\}}$와 표준적으로 동일시된다
(\hyperref[0.7.6.9-ko]{0, 7.6.9}). 따라서
(\hyperref[I.1.2.5-ko]{1.2.5}와 \hyperref[I.1.2.6-ko]{1.2.6})에 따라 위상공간
$\operatorname{Spf}(A_{\{f\}})$는 $\mathfrak{D}(f)$와 표준적으로
동일시된다. 더 나아가 $\mathfrak{D}(f)$에 포함되는
$\mathfrak{X}$의 모든 준콤팩트 열린집합 $U$에 대하여
$\Gamma(U,\mathscr{O}_\lambda)$는
$\operatorname{Spec}(S_f^{-1}A/S_f^{-1}\mathfrak{J}_\lambda)$의 구조층의
$U$ 위의 단면가군과 동일시된다
(\hyperref[I.1.3.6-ko]{1.3.6}). 따라서
$\mathfrak{Y}=\operatorname{Spf}(A_{\{f\}})$로 놓으면
$\Gamma(U,\mathscr{O}_{\mathfrak{X}})$는 단면가군
$\Gamma(U,\mathscr{O}_{\mathfrak{Y}})$와 동일시된다. 이로써 명제가
증명된다.
\end{proof}

\oldpage[I]{182}
\begin{env}[10.1.5]
\label{I.10.1.5-ko}
\emph{위상을 잊은} 환의 층으로서 $\operatorname{Spf}(A)$의 구조층
$\mathscr{O}_{\mathfrak{X}}$는 모든 $x\in\mathfrak{X}$에서
(\hyperref[I.10.1.4-ko]{10.1.4})에 따라 귀납극한
$\varinjlim_{f\notin\mathfrak{j}_x}A_{\{f\}}$와 동일시되는 줄기를 갖는다.
따라서 (\hyperref[0.7.6.17-ko]{0, 7.6.17}과
\hyperref[0.7.6.18-ko]{7.6.18}):
\end{env}

\begin{proposition}[10.1.6]
\label{I.10.1.6-ko}
모든 $x\in\mathfrak{X}=\operatorname{Spf}(A)$에 대하여 줄기
$\mathscr{O}_x$는 국소환이고, 그 잉여체는
$k(x)=A_x/\mathfrak{j}_xA_x$와 동형이다. 더 나아가 $A$가 아딕이고
뇌터이면 $\mathscr{O}_x$는 뇌터 환이다.
\end{proposition}

$k(x)$가 영환이 아니므로, 이 결과로부터 환의 층
$\mathscr{O}_{\mathfrak{X}}$의 \emph{지지집합}은
\emph{$\mathfrak{X}$와 같다}는 결론을 얻는다.

\subsection{아핀 형식적 스킴의 사상.}
\label{subsection:I.10.2-ko}

\begin{env}[10.2.1]
\label{I.10.2.1-ko}
$A$, $B$를 두 허용환이라 하고,
$\varphi:B\to A$를 \emph{연속} 준동형이라 하자. 그러면 연속사상
$ {}^a\varphi:\operatorname{Spec}(A)\to\operatorname{Spec}(B)$
(\hyperref[I.1.2.1-ko]{1.2.1})은
$\mathfrak{X}=\operatorname{Spf}(A)$를
$\mathfrak{Y}=\operatorname{Spf}(B)$ 안으로 보낸다. 실제로 $A$의 열린
소 아이디얼의 $\varphi$에 의한 역상은 $B$의 열린 소 아이디얼이다.
한편 모든 $g\in B$에 대하여 $\varphi$는 연속 준동형
\[
  \Gamma(\mathfrak{D}(g),\mathscr{O}_{\mathfrak{Y}})
  \longrightarrow
  \Gamma(\mathfrak{D}(\varphi(g)),\mathscr{O}_{\mathfrak{X}})
\]
을 정의한다. 이는 (\hyperref[I.10.1.4-ko]{10.1.4}),
(\hyperref[I.10.1.3-ko]{10.1.3}),
(\hyperref[0.7.6.7-ko]{0, 7.6.7})에 따른다. 이 준동형들은 $g$를 $g$의
배수로 바꿀 때의 제한사상들과 양립하고,
$\mathfrak{D}(\varphi(g))={}^a\varphi^{-1}(\mathfrak{D}(g))$이므로,
위상환의 층 사이의 \emph{연속} 준동형
$\mathscr{O}_{\mathfrak{Y}}\to
{}^a\varphi_*(\mathscr{O}_{\mathfrak{X}})$
(\hyperref[0.3.2.5-ko]{0, 3.2.5})을 정의한다. 이를 다시
$\widetilde{\varphi}$로 나타내겠다. 이렇게 하여 위상환 달린 공간의 사상
$\Phi=({}^a\varphi,\widetilde{\varphi})$를
$\mathfrak{X}\to\mathfrak{Y}$로 얻는다. 위상을 잊은 환의 층의
준동형으로서 $\widetilde{\varphi}$는 모든 $x\in\mathfrak{X}$에서 줄기
사이의 준동형
$\widetilde{\varphi}_x^\sharp:
\mathscr{O}_{{}^a\varphi(x)}\to\mathscr{O}_x$를 정의한다.
\end{env}

\begin{proposition}[10.2.2]
\label{I.10.2.2-ko}
$A$, $B$를 두 허용환이라 하고,
$\mathfrak{X}=\operatorname{Spf}(A)$,
$\mathfrak{Y}=\operatorname{Spf}(B)$라 하자. 위상환 달린 공간의 사상
$u=(\psi,\theta):\mathfrak{X}\to\mathfrak{Y}$가 $B\to A$인 어떤 연속 환 준동형
$\varphi$로부터 얻어지는
$({}^a\varphi,\widetilde{\varphi})$ 꼴이기 위한 필요충분조건은 모든
$x\in\mathfrak{X}$에 대하여 $\theta_x^\sharp$가 국소 준동형
$\mathscr{O}_{\psi(x)}\to\mathscr{O}_x$인 것이다.
\end{proposition}

\begin{proof}
이 조건은 필요하다. 실제로
$\mathfrak{p}=\mathfrak{j}_x\in\operatorname{Spf}(A)$,
$\mathfrak{q}=\varphi^{-1}(\mathfrak{j}_x)$라 하자.
$g\notin\mathfrak{q}$이면 $\varphi(g)\notin\mathfrak{p}$이고,
$\varphi$에서 얻어지는 준동형
$B_{\{g\}}\to A_{\{\varphi(g)\}}$
(\hyperref[0.7.6.7-ko]{0, 7.6.7})은
$\mathfrak{q}_{\{g\}}$를
$\mathfrak{p}_{\{\varphi(g)\}}$의 부분집합으로 보낸다. 따라서 귀납극한을
취하고 (\hyperref[I.10.1.5-ko]{10.1.5})와
(\hyperref[0.7.6.17-ko]{0, 7.6.17})을 고려하면
$\widetilde{\varphi}_x^\sharp$가 국소 준동형임을 알 수 있다.

역으로 $(\psi,\theta)$가 명제의 조건을 만족하는 사상이라고 하자.
(\hyperref[I.10.1.3-ko]{10.1.3})에 따라 $\theta$는 다음 연속 환 준동형을
정의한다:
\[
  \varphi=\Gamma(\theta):
  B=\Gamma(\mathfrak{Y},\mathscr{O}_{\mathfrak{Y}})
  \longrightarrow
  \Gamma(\mathfrak{X},\mathscr{O}_{\mathfrak{X}})=A.
\]
$\theta$에 대한 가정에 의하여 $\mathfrak{X}$ 위의
$\mathscr{O}_{\mathfrak{X}}$의 단면 $\varphi(g)$의 $x$에서의 싹이
가역일 필요충분조건은 $g$의 $\psi(x)$에서의 싹이 가역인 것이다. 그런데
(\hyperref[0.7.6.17-ko]{0, 7.6.17})에 따라
$\mathfrak{X}$ 위의 $\mathscr{O}_{\mathfrak{X}}$의 단면(각각
$\mathfrak{Y}$ 위의 $\mathscr{O}_{\mathfrak{Y}}$의 단면) 가운데
$x$에서의 싹(각각 $\psi(x)$에서의 싹)이 가역이 아닌 것들은 정확히
$\mathfrak{j}_x$
\oldpage[I]{183}
(각각 $\mathfrak{j}_{\psi(x)}$)의 원소들이다. 따라서 앞의 관찰로부터
${}^a\varphi=\psi$를 얻는다. 마지막으로 모든 $g\in B$에 대하여 도식
\[
  \xymatrix{
    B=\Gamma(\mathfrak{Y},\mathscr{O}_{\mathfrak{Y}})
      \ar[r]^\varphi\ar[d] &
    \Gamma(\mathfrak{X},\mathscr{O}_{\mathfrak{X}})=A\ar[d]\\
    B_{\{g\}}=\Gamma(\mathfrak{D}(g),\mathscr{O}_{\mathfrak{Y}})
      \ar[r]_{\Gamma(\theta_{\mathfrak{D}(g)})} &
    \Gamma(\mathfrak{D}(\varphi(g)),\mathscr{O}_{\mathfrak{X}})
      =A_{\{\varphi(g)\}}
  }
\]
은 가환이다. 완비 분수환의 보편 성질
(\hyperref[0.7.6.6-ko]{0, 7.6.6})에 따라 모든 $g\in B$에 대하여
$\theta_{\mathfrak{D}(g)}$는
$\widetilde{\varphi}_{\mathfrak{D}(g)}$와 같고, 따라서
(\hyperref[0.3.2.5-ko]{0, 3.2.5})에 따라
$\theta=\widetilde{\varphi}$이다.
\end{proof}

(\hyperref[I.10.2.2-ko]{10.2.2})의 조건을 만족하는 위상환 달린 공간의
사상 $(\psi,\theta)$를 \emph{아핀 형식적 스킴의 사상}이라 하겠다.
$A$에 $\operatorname{Spf}(A)$를 대응시키는 함자와 $\mathfrak{X}$에
$\Gamma(\mathfrak{X},\mathscr{O}_{\mathfrak{X}})$를 대응시키는 함자는
허용환의 범주와 아핀 형식적 스킴의 범주의 반대 범주 사이의
\emph{동치}를 정의한다고 할 수 있다(T, I, 1.2).

\begin{env}[10.2.3]
\label{I.10.2.3-ko}
(\hyperref[I.10.2.2-ko]{10.2.2})의 특수한 경우로서, $f\in A$에 대하여
$\mathfrak{X}$가 $\mathfrak{D}(f)$ 위에 유도하는 아핀 형식적 스킴의
표준 포함 사상은 표준 연속 준동형 $A\to A_{\{f\}}$에 대응한다. 이제
(\hyperref[I.10.2.2-ko]{10.2.2})의 가정 아래에서 $h$를 $B$의 원소,
$g$를 $\varphi(h)$의 배수인 $A$의 원소라 하자. 그러면
$\psi(\mathfrak{D}(g))\subset\mathfrak{D}(h)$이다. $u$를
$\mathfrak{D}(g)$에 제한한 사상을 $\mathfrak{D}(g)$에서
$\mathfrak{D}(h)$로 가는 사상으로 보면, 이는 도식
\[
  \xymatrix{
    \mathfrak{D}(g)\ar[r]^v\ar[d] &
    \mathfrak{D}(h)\ar[d]\\
    \mathfrak{X}\ar[r]_u &
    \mathfrak{Y}
  }
\]
을 가환이게 하는 유일한 사상 $v$이다. 이 사상은 도식
\[
  \xymatrix{
    A\ar[d] &
    B\ar[l]_\varphi\ar[d]\\
    A_{\{g\}} &
    B_{\{h\}}\ar[l]_{\varphi'}
  }
\]
을 가환이게 하는 유일한 연속 준동형
$\varphi':B_{\{h\}}\to A_{\{g\}}$
(\hyperref[0.7.6.7-ko]{0, 7.6.7})에 대응한다.
\end{env}

\subsection{아핀 형식적 스킴의 정의 아이디얼.}
\label{subsection:I.10.3-ko}

\begin{env}[10.3.1]
\label{I.10.3.1-ko}
$A$를 허용환, $\mathfrak{J}$를 $A$의 열린 아이디얼,
$\mathfrak{X}$를 아핀 형식적 스킴 $\operatorname{Spf}(A)$라 하자.
$(\mathfrak{J}_\lambda)$를 $\mathfrak{J}$에 포함되는 $A$의 정의
아이디얼들의 집합이라 하자. 그러면
$\widetilde{\mathfrak{J}}/\widetilde{\mathfrak{J}}_\lambda$는
$\widetilde{A}/\widetilde{\mathfrak{J}}_\lambda$의 아이디얼층이다.
$\mathfrak{X}$ 위에
$\widetilde{\mathfrak{J}}/\widetilde{\mathfrak{J}}_\lambda$가 유도하는
층들의 사영극한을 $\mathfrak{J}^\Delta$로 나타내자. 이는
$\mathscr{O}_{\mathfrak{X}}$의 \emph{아이디얼층}과 동일시된다
(\hyperref[0.3.2.6-ko]{0, 3.2.6}). 모든 $f\in A$에 대하여
$\Gamma(\mathfrak{D}(f),\mathfrak{J}^\Delta)$는
$S_f^{-1}\mathfrak{J}/S_f^{-1}\mathfrak{J}_\lambda$들의 사영극한, 즉
환 $A_{\{f\}}$의 열린 아이디얼 $\mathfrak{J}_{\{f\}}$와 동일시된다
(\hyperref[0.7.6.9-ko]{0, 7.6.9}). 특히
$\Gamma(\mathfrak{X},\mathfrak{J}^\Delta)=\mathfrak{J}$이다.
$\mathfrak{D}(f)$들이 $\mathfrak{X}$의 위상의 기저를 이루므로 이로부터
다음을 얻는다:
\[
  \mathfrak{J}^\Delta|\mathfrak{D}(f)
  =(\mathfrak{J}_{\{f\}})^\Delta.
  \tag{10.3.1.1}\label{I.10.3.1.1-ko}
\]
\end{env}

\oldpage[I]{184}
\begin{env}[10.3.2]
\label{I.10.3.2-ko}
(\hyperref[I.10.3.1-ko]{10.3.1})의 표기 아래에서, 모든 $f\in A$에
대하여
$A_{\{f\}}=\Gamma(\mathfrak{D}(f),\mathscr{O}_{\mathfrak{X}})$에서
$\Gamma\bigl(\mathfrak{D}(f),
(\widetilde{A}/\widetilde{\mathfrak{J}})|\mathfrak{X}\bigr)
=S_f^{-1}A/S_f^{-1}\mathfrak{J}$로 가는 표준 사상은 \emph{전사}이고,
그 핵은
$\Gamma(\mathfrak{D}(f),\mathfrak{J}^\Delta)=\mathfrak{J}_{\{f\}}$이다
(\hyperref[0.7.6.9-ko]{0, 7.6.9}). 따라서 이 사상들은 위상환의 층
$\mathscr{O}_{\mathfrak{X}}$에서 이산환의 층
$(\widetilde{A}/\widetilde{\mathfrak{J}})|\mathfrak{X}$로 가는 연속
\emph{전사} 준동형을 정의하며, 이를 \emph{표준}이라 한다. 그 핵은
$\mathfrak{J}^\Delta$이다. 이 준동형은
$\varphi$가 연속 준동형 $A\to A/\mathfrak{J}$일 때의
$\widetilde{\varphi}$
(\hyperref[I.10.2.1-ko]{10.2.1})와 다르지 않다. 아핀 형식적 스킴의 사상
$({}^a\varphi,\widetilde{\varphi}):
\operatorname{Spec}(A/\mathfrak{J})\to\mathfrak{X}$
(여기서 ${}^a\varphi$는 닫힌 부분공간
$V(\mathfrak{J})\subset\mathfrak{X}$로의 표준 위상동형사상에 그
포함사상을 합성한 것이다)도 \emph{표준}이라 한다.\footnote{\textbf{번역 주.} 인쇄 원문은
${}^a\varphi$가 $\mathfrak{X}$에서 자기 자신으로 가는 항등
위상동형사상이라고 쓰지만, 이는 $\mathfrak{J}$가 임의의 열린 아이디얼일
때 일반적으로 성립하지 않는다. 예를 들어 체 $k$에 대하여 이산 허용환
$A=k\times k$와 열린 아이디얼 $\mathfrak{J}=k\times\{0\}$에 대하여
$\operatorname{Spf}(A)$는 두 점이지만
$\operatorname{Spec}(A/\mathfrak{J})$는 한 점이다. 올바른 상은
$V(\mathfrak{J})\subset\mathfrak{X}$이며, $\mathfrak{J}$가 정의
아이디얼일 때에는 이 상이 $\mathfrak{X}$ 전체이다. 이는 공식 정오표에
실린 항목은 아니며, 위와 같이 형에 맞게 바로잡았다.} 따라서 앞의
결과에 의하여 \emph{표준 동형사상}
\[
  \mathscr{O}_{\mathfrak{X}}/\mathfrak{J}^\Delta
  \xrightarrow{\sim}
  (\widetilde{A}/\widetilde{\mathfrak{J}})|\mathfrak{X}.
  \tag{10.3.2.1}\label{I.10.3.2.1-ko}
\]
을 얻는다.

$\Gamma(\mathfrak{X},\mathfrak{J}^\Delta)=\mathfrak{J}$이므로
$\mathfrak{J}\to\mathfrak{J}^\Delta$가
\emph{포함관계를 엄격히 보존한다}는 것은 명백하다. 앞의 결과에 따라
$\mathfrak{J}\subset\mathfrak{J}'$이면 층
${\mathfrak{J}'}^\Delta/\mathfrak{J}^\Delta$는
$\widetilde{\mathfrak{J}'}/\widetilde{\mathfrak{J}}
=\widetilde{\mathfrak{J}'/\mathfrak{J}}$와 표준적으로 동형이다.
\end{env}

\begin{env}[10.3.3]
\label{I.10.3.3-ko}
가정과 표기는 계속
(\hyperref[I.10.3.1-ko]{10.3.1})의 것이라 하자.
$\mathscr{O}_{\mathfrak{X}}$의 아이디얼층 $\mathscr{J}$가 다음 조건을
만족할 때 이를 \emph{$\mathfrak{X}$의 정의 아이디얼층}(또는
\emph{$\mathfrak{X}$의 정의 아이디얼})이라 하겠다. 모든
$x\in\mathfrak{X}$에 대하여 $x$의 열린 근방 $\mathfrak{D}(f)$가
존재하고, 여기서 $f\in A$이며,
$\mathscr{J}|\mathfrak{D}(f)$는 $\mathfrak{H}^\Delta$ 꼴이다. 이때
$\mathfrak{H}$는 $A_{\{f\}}$의 정의 아이디얼이다.
\end{env}

\begin{proposition}[10.3.4]
\label{I.10.3.4-ko}
모든 $f\in A$에 대하여 $\mathfrak{X}$의 모든 정의 아이디얼은
$\mathfrak{D}(f)$의 정의 아이디얼을 유도한다.
\end{proposition}

\begin{proof}
이는 (\hyperref[I.10.3.1.1-ko]{10.3.1.1})에서 따른다.
\end{proof}

\begin{proposition}[10.3.5]
\label{I.10.3.5-ko}
$A$가 허용환이면 $\mathfrak{X}=\operatorname{Spf}(A)$의 모든 정의
아이디얼은 $A$의 유일하게 결정되는 정의 아이디얼 $\mathfrak{J}$에 대한
$\mathfrak{J}^\Delta$ 꼴이다.
\end{proposition}

\begin{proof}
실제로 $\mathscr{J}$를 $\mathfrak{X}$의 정의 아이디얼이라 하자.
가정과 $\mathfrak{X}$의 준콤팩트성에 따라 유한 개의 원소
$f_i\in A$가 존재하여 $\mathfrak{D}(f_i)$들이 $\mathfrak{X}$를 덮고,
$\mathscr{J}|\mathfrak{D}(f_i)=\mathfrak{H}_i^\Delta$이다. 여기서
$\mathfrak{H}_i$는 $A_{\{f_i\}}$의 정의 아이디얼이다. 모든 $i$에
대하여 $(\mathfrak{K}_i)_{\{f_i\}}=\mathfrak{H}_i$인 $A$의 열린
아이디얼 $\mathfrak{K}_i$가 존재한다
(\hyperref[0.7.6.9-ko]{0, 7.6.9}). 모든 $\mathfrak{K}_i$에 포함되는
$A$의 정의 아이디얼을 $\mathfrak{K}$라 하자.
$\mathscr{J}/\mathfrak{K}^\Delta$의
$\operatorname{Spec}(A/\mathfrak{K})$의 구조층
$\widetilde{A/\mathfrak{K}}$ 안에서의 표준 상
(\hyperref[I.10.3.2-ko]{10.3.2})은 $\mathfrak{D}(f_i)$ 위로 제한하면
$\widetilde{\mathfrak{K}_i/\mathfrak{K}}$의 제한과 같다. 따라서 이
표준 상은 $\operatorname{Spec}(A/\mathfrak{K})$ 위의 준연접층이고,
$\mathfrak{K}$를 포함하는 $A$의 아이디얼 $\mathfrak{J}$에 대하여
$\widetilde{\mathfrak{J}/\mathfrak{K}}$ 꼴이다
(\hyperref[I.1.4.1-ko]{1.4.1}). 그러므로
$\mathscr{J}=\mathfrak{J}^\Delta$이다
(\hyperref[I.10.3.2-ko]{10.3.2}). 또한 모든 $i$에 대하여
$\mathfrak{H}_i^{n_i}\subset\mathfrak{K}_{\{f_i\}}$인 정수 $n_i$가
존재한다. $n$을 $n_i$들의 최댓값이라 하면
$(\mathscr{J}/\mathfrak{K}^\Delta)^n=0$이고, 따라서
(\hyperref[I.10.3.2-ko]{10.3.2})에 의하여
$(\widetilde{\mathfrak{J}/\mathfrak{K}})^n=0$이다. 그러므로 마침내
$(\mathfrak{J}/\mathfrak{K})^n=0$이고
(\hyperref[I.1.3.13-ko]{1.3.13}), 이는 $\mathfrak{J}$가 $A$의 정의
아이디얼임을 보인다 (\hyperref[0.7.1.4-ko]{0, 7.1.4}).
\end{proof}

\begin{proposition}[10.3.6]
\label{I.10.3.6-ko}
$A$를 아딕환, $\mathfrak{J}$를 $A$의 정의 아이디얼이라 하고,
$\mathfrak{J}/\mathfrak{J}^2$가 유한 생성
$A/\mathfrak{J}$-가군이라고 하자. 그러면 모든 정수 $n>0$에 대하여
$(\mathfrak{J}^\Delta)^n=(\mathfrak{J}^n)^\Delta$이다.
\end{proposition}

\begin{proof}
실제로 모든 $f\in A$에 대하여 $\mathfrak{J}^n$은 열린 아이디얼이므로
\[
  \bigl(\Gamma(\mathfrak{D}(f),\mathfrak{J}^\Delta)\bigr)^n
  =(\mathfrak{J}_{\{f\}})^n
  =(\mathfrak{J}^n)_{\{f\}}
  =\Gamma\bigl(\mathfrak{D}(f^n),(\mathfrak{J}^n)^\Delta\bigr)
\]
\oldpage[I]{185}
이다. 이는 (\hyperref[I.10.3.1.1-ko]{10.3.1.1})과
(\hyperref[0.7.6.12-ko]{0, 7.6.12})에서 따른다.
$(\mathfrak{J}^\Delta)^n$은 준층
$U\mapsto(\Gamma(U,\mathfrak{J}^\Delta))^n$의 층화이므로
(\hyperref[0.4.1.6-ko]{0, 4.1.6}), $\mathfrak{D}(f)$들이
$\mathfrak{X}$의 위상의 기저를 이룬다는 사실로부터 결론이
따른다.\footnote{\textbf{번역 주.} 인쇄 원문은 이 명제의 증명에서
“따름정리가 따른다”라고 쓰지만, 이곳에는 따름정리가 없으므로 “결론이
따른다”로 바로잡았다.}
\end{proof}

\begin{env}[10.3.7]
\label{I.10.3.7-ko}
$\mathfrak{X}$의 정의 아이디얼들의 족 $(\mathscr{J}_\lambda)$가 다음
조건을 만족할 때 이를 \emph{정의 아이디얼들의 기본계}라고 한다.
$\mathfrak{X}$의 모든 정의 아이디얼은 어떤 $\mathscr{J}_\lambda$를
포함한다. $\mathscr{J}_\lambda=\mathfrak{J}_\lambda^\Delta$이므로 이는
$\mathfrak{J}_\lambda$들이 $A$에서 \emph{0의 근방 기본계}를 이룬다는
것과 같다. $\mathfrak{D}(f_\alpha)$들이 $\mathfrak{X}$를 덮도록
$A$의 원소들의 족 $(f_\alpha)$를 잡자.
$(\mathscr{J}_\lambda)$가 $\mathscr{O}_{\mathfrak{X}}$의 아이디얼들의
감소 여과족이고, 모든 $\alpha$에 대하여
$(\mathscr{J}_\lambda|\mathfrak{D}(f_\alpha))$가
$\mathfrak{D}(f_\alpha)$의 정의 아이디얼들의 기본계를 이룬다면,
$(\mathscr{J}_\lambda)$는 $\mathfrak{X}$의 정의 아이디얼들의 기본계이다.
실제로 $\mathfrak{X}$의 모든 정의 아이디얼 $\mathscr{J}$에 대하여
$\mathfrak{D}(f_i)$들로 이루어진 $\mathfrak{X}$의 유한 덮개가 존재하여,
모든 $i$에 대하여
$\mathscr{J}_{\lambda_i}|\mathfrak{D}(f_i)$는
$\mathfrak{D}(f_i)$의 정의 아이디얼이고
$\mathscr{J}|\mathfrak{D}(f_i)$에 포함된다. 모든 $i$에 대하여
$\mathscr{J}_\mu\subset\mathscr{J}_{\lambda_i}$인 지표 $\mu$를 잡으면,
(\hyperref[I.10.3.3-ko]{10.3.3})에 따라 $\mathscr{J}_\mu$는
$\mathfrak{X}$의 정의 아이디얼이고, 명백히 $\mathscr{J}$에 포함된다.
이로부터 주장이 따른다.
\end{env}

\subsection{형식적 준스킴과 형식적 준스킴의 사상.}
\label{subsection:I.10.4-ko}

\begin{env}[10.4.1]
\label{I.10.4.1-ko}
위상환 달린 공간 $\mathfrak{X}$가 주어졌다고 하자. 열린집합
$U\subset\mathfrak{X}$ 위에 $\mathfrak{X}$가 유도하는 위상환 달린 공간이
아핀 형식적 스킴(각각 그 환이 아딕인 아핀 형식적 스킴, 아딕이고 뇌터인
아핀 형식적 스킴)일 때, $U$를 \emph{아핀 형식적 열린집합}(각각
\emph{아딕 아핀 형식적 열린집합}, \emph{뇌터 아핀 형식적 열린집합})이라
한다.
\end{env}

\begin{definition}[10.4.2]
\label{I.10.4.2-ko}
모든 점이 아핀 형식적 열린 근방을 갖는 위상환 달린 공간
$\mathfrak{X}$를 형식적 준스킴이라 한다. 형식적 준스킴
$\mathfrak{X}$가 아딕(각각 국소 뇌터)이라는 것은 $\mathfrak{X}$의 모든
점이 아딕 아핀 형식적 열린 근방(각각 뇌터 아핀 형식적 열린 근방)을
갖는다는 뜻이다.
국소 뇌터인 $\mathfrak{X}$의 바탕공간이 준콤팩트일 때(따라서 그
바탕공간은 뇌터이다) 이를 뇌터라 한다.
\end{definition}

\begin{proposition}[10.4.3]
\label{I.10.4.3-ko}
$\mathfrak{X}$가 형식적 준스킴(각각 국소 뇌터 형식적 준스킴)이면,
아핀 형식적 열린집합들(각각 뇌터 아핀 형식적 열린집합들)은
$\mathfrak{X}$의 위상의 기저를 이룬다.
\end{proposition}

\begin{proof}
이는 (\hyperref[I.10.4.2-ko]{10.4.2})와
(\hyperref[I.10.1.4-ko]{10.1.4})에서 따른다. 여기서 $A$가 뇌터
아딕환이면 $A_{\{f\}}$도 모든 $f\in A$에 대하여 그러하다는 사실
(\hyperref[0.7.6.11-ko]{0, 7.6.11})을 사용한다.
\end{proof}

\begin{corollary}[10.4.4]
\label{I.10.4.4-ko}
$\mathfrak{X}$가 형식적 준스킴(각각 국소 뇌터 형식적 준스킴, 뇌터
형식적 준스킴)이면, $\mathfrak{X}$의 모든 열린집합 위에 유도된
위상환 달린 공간도 형식적 준스킴(각각 국소 뇌터 형식적 준스킴,
뇌터 형식적 준스킴)이다.
\end{corollary}

\begin{definition}[10.4.5]
\label{I.10.4.5-ko}
두 형식적 준스킴 $\mathfrak{X}$, $\mathfrak{Y}$가 주어졌다고 하자.
$\mathfrak{X}$에서 $\mathfrak{Y}$로 가는 위상환 달린 공간의 사상
$(\psi,\theta)$ 가운데 모든 $x\in\mathfrak{X}$에 대하여
$\theta_x^\sharp$가 국소 준동형
$\mathscr{O}_{\psi(x)}\to\mathscr{O}_x$인 것을 (형식적 준스킴의)
사상이라 한다.
\end{definition}

두 형식적 준스킴의 사상의 합성도 형식적 준스킴의 사상임은 자명하다.
따라서 형식적 준스킴들은 하나의 \emph{범주}를 이룬다.
$\operatorname{Hom}(\mathfrak{X},\mathfrak{Y})$로 형식적 준스킴
$\mathfrak{X}$에서 형식적 준스킴 $\mathfrak{Y}$로 가는 사상들의 집합을
나타낸다.

\oldpage[I]{186}
$U$가 $\mathfrak{X}$의 열린부분이면, $U$ 위에 $\mathfrak{X}$가 유도하는
형식적 준스킴에서 $\mathfrak{X}$로 가는 표준 포함 사상은 형식적
준스킴의 사상이다(더 나아가 위상환 달린 공간의 \emph{단사사상}이다
(\hyperref[0.4.1.1-ko]{0, 4.1.1})).

\begin{proposition}[10.4.6]
\label{I.10.4.6-ko}
$\mathfrak{X}$를 형식적 준스킴이라 하고,
$\mathfrak{S}=\operatorname{Spf}(A)$를 아핀 형식적 스킴이라 하자.
형식적 준스킴 $\mathfrak{X}$에서 형식적 준스킴 $\mathfrak{S}$로 가는
사상들과 환 $A$에서 위상환
$\Gamma(\mathfrak{X},\mathscr{O}_{\mathfrak{X}})$로 가는 연속
준동형들 사이에는 표준 전단사 대응이 존재한다.
\end{proposition}

증명은 (\hyperref[I.2.2.4-ko]{2.2.4})의 증명과 같다. 다만
“준동형”을 “연속 준동형”으로, “아핀 열린집합”을 “아핀 형식적
열린집합”으로 바꾸고, (\hyperref[I.10.2.2-ko]{10.2.2})를
(\hyperref[I.1.7.3-ko]{1.7.3}) 대신 사용한다. 자세한 내용은 독자에게
맡긴다.

\begin{env}[10.4.7]
\label{I.10.4.7-ko}
형식적 준스킴 $\mathfrak{S}$가 주어졌다고 하자. 형식적 준스킴
$\mathfrak{X}$와 사상 $\varphi:\mathfrak{X}\to\mathfrak{S}$의 자료가
주어졌을 때, 형식적 준스킴 $\mathfrak{X}$를
\emph{$\mathfrak{S}$ 위의 형식적 준스킴} 또는
\emph{$\mathfrak{S}$-형식적 준스킴}이라 하고,
$\varphi$를 $\mathfrak{S}$-형식적 준스킴 $\mathfrak{X}$의
\emph{구조 사상}이라 한다. $\mathfrak{S}=\operatorname{Spf}(A)$이고
$A$가 허용환이면, $\mathfrak{S}$-형식적 준스킴 $\mathfrak{X}$를
\emph{$A$-형식적 준스킴} 또는 \emph{$A$ 위의 형식적 준스킴}이라고도
한다. 모든 형식적 준스킴은 (이산 위상을 갖춘) $\mathbf{Z}$ 위의
형식적 준스킴으로 볼 수 있다.

$\mathfrak{X}$, $\mathfrak{Y}$를 두 $\mathfrak{S}$-형식적 준스킴이라
하자. 사상 $u:\mathfrak{X}\to\mathfrak{Y}$에 대하여 도식
\[
  \xymatrix{
    \mathfrak{X}\ar[rr]^u\ar[rd] & &
    \mathfrak{Y}\ar[ld]\\
    & \mathfrak{S}
  }
\]
(빗금 화살표들은 구조 사상이다)이 가환일 때, 이 사상을
\emph{$\mathfrak{S}$-사상}이라 한다. 이 정의에 따라 (고정된
$\mathfrak{S}$에 대한) $\mathfrak{S}$-형식적 준스킴들은 하나의
\emph{범주}를 이룬다.
$\operatorname{Hom}_{\mathfrak{S}}(\mathfrak{X},\mathfrak{Y})$로
$\mathfrak{S}$-형식적 준스킴 $\mathfrak{X}$에서
$\mathfrak{S}$-형식적 준스킴 $\mathfrak{Y}$로 가는
$\mathfrak{S}$-사상들의 집합을 나타낸다.
$\mathfrak{S}=\operatorname{Spf}(A)$이면 \emph{$A$-사상}이라는 말도
\emph{$\mathfrak{S}$-사상} 대신 사용한다.
\end{env}

\begin{env}[10.4.8]
\label{I.10.4.8-ko}
모든 아핀 스킴은 아핀 형식적 스킴으로 볼 수 있으므로
(\hyperref[I.10.1.2-ko]{10.1.2}), 모든 (일반적인) 준스킴은 형식적
준스킴으로 볼 수 있다. 또한 (\hyperref[I.10.4.5-ko]{10.4.5})에 따라
\emph{일반적인} 준스킴에 대해서는 \emph{형식적} 준스킴의 사상(각각
$S$-사상)이 §~\hyperref[section:I.2-ko]{2}에서 정의한 사상(각각
$S$-사상)과 일치한다.
\end{env}

\subsection{형식적 준스킴의 정의 아이디얼.}
\label{subsection:I.10.5-ko}

\begin{env}[10.5.1]
\label{I.10.5.1-ko}
형식적 준스킴 $\mathfrak{X}$가 주어졌다고 하자.
$\mathscr{O}_{\mathfrak{X}}$-아이디얼 $\mathscr{J}$가 다음 조건을 만족할
때 이를 $\mathfrak{X}$의 \emph{정의 아이디얼층}(또는
\emph{정의 아이디얼})이라 한다. 모든 $x\in\mathfrak{X}$에 대하여 아핀
형식적 열린 근방 $U$가 존재하여, $\mathscr{J}|U$는 $U$ 위에
$\mathfrak{X}$가 유도하는 아핀 형식적 스킴의 정의 아이디얼이다
(\hyperref[I.10.3.3-ko]{10.3.3}). 이때
(\hyperref[I.10.3.1.1-ko]{10.3.1.1})과
(\hyperref[I.10.4.3-ko]{10.4.3})에 따라 모든 열린집합
$V\subset\mathfrak{X}$에 대하여 $\mathscr{J}|V$는 $V$ 위에
$\mathfrak{X}$가 유도하는 형식적 준스킴의 정의 아이디얼이다.

$\mathfrak{X}$의 정의 아이디얼들로 이루어진 족
$(\mathscr{J}_\lambda)$에 대하여, $\mathfrak{X}$의 아핀 형식적
열린집합들로 이루어진 덮개 $(U_\alpha)$가 존재하여 모든 $\alpha$에
대해 $\mathscr{J}_\lambda|U_\alpha$들로 이루어진 족이 $U_\alpha$ 위에
$\mathfrak{X}$가 유도하는 아핀 형식적 스킴의 정의 아이디얼들의
기본계를 이룰 때(\hyperref[I.10.3.6-ko]{10.3.6}), 이 족을
\emph{정의 아이디얼들의}
\oldpage[I]{187}
\emph{기본계}라고 한다. (\hyperref[I.10.3.7-ko]{10.3.7})의 마지막
주에 따라, $\mathfrak{X}$가 아핀 형식적 스킴이면 이 정의는
(\hyperref[I.10.3.7-ko]{10.3.7})에서 주어진 정의와 일치한다. $V$를
$\mathfrak{X}$의 임의의 열린집합이라 하면, $\mathscr{J}_\lambda|V$들은
$V$ 위에 유도된 형식적 준스킴의 정의 아이디얼들의 기본계를 이룬다. 이는
(\hyperref[I.10.3.1.1-ko]{10.3.1.1})에서 따른다. $\mathfrak{X}$가
\emph{국소 뇌터} 형식적 준스킴이고 $\mathscr{J}$가 $\mathfrak{X}$의
정의 아이디얼이면, (\hyperref[I.10.3.6-ko]{10.3.6})에 따라 거듭제곱들
$\mathscr{J}^n$은 $\mathfrak{X}$의 정의 아이디얼들의 기본계를 이룬다.
\end{env}

\begin{env}[10.5.2]
\label{I.10.5.2-ko}
$\mathfrak{X}$를 형식적 준스킴이라 하고, $\mathscr{J}$를
$\mathfrak{X}$의 정의 아이디얼이라 하자. 그러면 환 달린 공간
$(\mathfrak{X},\mathscr{O}_{\mathfrak{X}}/\mathscr{J})$는
\emph{준스킴}(일반적인 의미에서)이다. $\mathfrak{X}$가 아핀 형식적
스킴(각각 국소 뇌터 형식적 준스킴, 뇌터 형식적 준스킴)이면 이 준스킴은
아핀(각각 국소 뇌터, 뇌터)이다. 실제로 곧바로 아핀인 경우로 환원되며,
그 경우에는 (\hyperref[I.10.3.2-ko]{10.3.2})에서 이미 이 명제를
증명했다. 또한
$\theta:\mathscr{O}_{\mathfrak{X}}\to
\mathscr{O}_{\mathfrak{X}}/\mathscr{J}$가 표준 준동형이면,
$u=(1_{\mathfrak{X}},\theta)$는 형식적 준스킴의 \emph{사상}
(\emph{표준 사상}이라 한다)
$(\mathfrak{X},\mathscr{O}_{\mathfrak{X}}/\mathscr{J})\to
(\mathfrak{X},\mathscr{O}_{\mathfrak{X}})$이다. 이 사실도 곧바로 아핀인
경우로 환원되며, 그 경우에는
(\hyperref[I.10.3.2-ko]{10.3.2})에서 이미 확인했다.
\end{env}

\begin{proposition}[10.5.3]
\label{I.10.5.3-ko}
$\mathfrak{X}$를 형식적 준스킴이라 하고, $(\mathscr{J}_\lambda)$를
$\mathfrak{X}$의 정의 아이디얼들의 기본계라 하자. 그러면 위상환의 층
$\mathscr{O}_{\mathfrak{X}}$는 의사이산 위상환의 층들
(\hyperref[0.3.8.1-ko]{0, 3.8.1})
$\mathscr{O}_{\mathfrak{X}}/\mathscr{J}_\lambda$의 사영극한이다.
\end{proposition}

\begin{proof}
$\mathfrak{X}$의 위상에는 준콤팩트 아핀 형식적 열린집합들로 이루어진
기저가 있으므로(\hyperref[I.10.4.3-ko]{10.4.3}), 아핀인 경우로
환원된다. 이 경우 명제는 (\hyperref[I.10.3.5-ko]{10.3.5}),
(\hyperref[I.10.3.2-ko]{10.3.2})와
(\hyperref[I.10.1.1-ko]{10.1.1})의 정의에서 따른다.
\end{proof}

모든 형식적 준스킴이 정의 아이디얼을 갖는지는 확실하지 않다. 그러나
다음은 성립한다.

\begin{proposition}[10.5.4]
\label{I.10.5.4-ko}
$\mathfrak{X}$를 국소 뇌터 형식적 준스킴이라 하자. 그러면
정의 아이디얼 $\mathscr{T}$가 $\mathfrak{X}$의 가장 큰 정의
아이디얼로서 존재한다. 이는 정의 아이디얼 $\mathscr{J}$ 가운데 준스킴
$(\mathfrak{X},\mathscr{O}_{\mathfrak{X}}/\mathscr{J})$가 축소인
유일한 것이다. $\mathscr{J}$가
$\mathfrak{X}$의 정의 아이디얼이면, $\mathscr{T}$는 표준 준동형
$\mathscr{O}_{\mathfrak{X}}\to
\mathscr{O}_{\mathfrak{X}}/\mathscr{J}$에 의한
$\mathscr{O}_{\mathfrak{X}}/\mathscr{J}$의 멱영근기의 역상이다.
\end{proposition}

\begin{proof}
먼저 $\mathfrak{X}=\operatorname{Spf}(A)$이고 $A$가 뇌터 아딕환이라고
하자. $\mathscr{T}$의 존재성과 성질들은
(\hyperref[I.10.3.5-ko]{10.3.5})\footnote{\textbf{번역 주.} 인쇄 원문은
이곳에서 (10.3.4)를 참조하지만, 출판된 EGA II 정오표는 이를
(10.3.5)로 고치도록 지시한다. 그 공식 교정을 적용했다.}와
(\hyperref[I.5.1.1-ko]{5.1.1})에서 바로 따른다. 이때 $A$의 가장 큰
정의 아이디얼의 존재성과 성질들
(\hyperref[0.7.1.6-ko]{0, 7.1.6}과
\hyperref[0.7.1.7-ko]{7.1.7})을 이용한다.

일반적인 경우에 $\mathscr{T}$의 존재성과 성질들을 증명하려면,
$U\supset V$를 $\mathfrak{X}$의 두 뇌터 아핀 형식적 열린집합이라 할 때
$U$의 가장 큰 정의 아이디얼 $\mathscr{T}_U$가 $V$의 가장 큰 정의
아이디얼 $\mathscr{T}_V$를 유도함을 보이면 충분하다. 그런데
$(V,(\mathscr{O}_{\mathfrak{X}}|V)/(\mathscr{T}_U|V))$는 축소이므로,
이는 앞에서 보인 결과에서 따른다.
\end{proof}

$\mathfrak{X}_{\mathrm{red}}$는 축소 준스킴(일반적인 의미에서)
$(\mathfrak{X},\mathscr{O}_{\mathfrak{X}}/\mathscr{T})$를 나타낸다.

\begin{corollary}[10.5.5]
\label{I.10.5.5-ko}
$\mathfrak{X}$를 국소 뇌터 형식적 준스킴이라 하고, $\mathscr{T}$를
$\mathfrak{X}$의 가장 큰 정의 아이디얼이라 하자. $V$가
$\mathfrak{X}$의 임의의 열린집합이면, $\mathscr{T}|V$는
$\mathfrak{X}$가 $V$ 위에 유도하는 형식적 준스킴의 가장 큰 정의
아이디얼이다.
\end{corollary}

\oldpage[I]{188}

\begin{proposition}[10.5.6]
\label{I.10.5.6-ko}
$\mathfrak{X}$, $\mathfrak{Y}$를 두 형식적 준스킴이라 하고,
$\mathscr{J}$와 $\mathscr{K}$를 각각 $\mathfrak{X}$와
$\mathfrak{Y}$의 정의 아이디얼이라 하며,
$f:\mathfrak{X}\to\mathfrak{Y}$를 형식적 준스킴의 사상이라 하자.
\begin{enumerate}
  \item[{\upshape(i)}]
    $f^*(\mathscr{K})\mathscr{O}_{\mathfrak{X}}\subset\mathscr{J}$이면,
    일반적인 의미에서 준스킴의 사상인
    $f':(\mathfrak{X},\mathscr{O}_{\mathfrak{X}}/\mathscr{J})\to
    (\mathfrak{Y},\mathscr{O}_{\mathfrak{Y}}/\mathscr{K})$가 유일하게
    존재하고 다음 도식을 가환이게 한다.
    \[
      \label{I.10.5.6.1-ko}
      \xymatrix{
        (\mathfrak{X},\mathscr{O}_{\mathfrak{X}})\ar[r]^f &
        (\mathfrak{Y},\mathscr{O}_{\mathfrak{Y}})\\
        (\mathfrak{X},\mathscr{O}_{\mathfrak{X}}/\mathscr{J})
          \ar[r]_{f'}\ar[u] &
        (\mathfrak{Y},\mathscr{O}_{\mathfrak{Y}}/\mathscr{K})\ar[u]
      }
      \tag{10.5.6.1}
    \]
    여기서 세로 화살표들은 표준 사상이다.
  \item[{\upshape(ii)}]
    $\mathfrak{X}=\operatorname{Spf}(A)$,
    $\mathfrak{Y}=\operatorname{Spf}(B)$가 아핀 형식적 스킴이고,
    $\mathscr{J}=\mathfrak{J}^{\Delta}$,
    $\mathscr{K}=\mathfrak{K}^{\Delta}$라 하자. 여기서
    $\mathfrak{J}$와 $\mathfrak{K}$는 각각 $A$와 $B$의 정의
    아이디얼이다. 또한
    $f=({}^a\varphi,\widetilde{\varphi})$이고
    $\varphi:B\to A$가 연속 준동형이라 하자. 이때
    $f^*(\mathscr{K})\mathscr{O}_{\mathfrak{X}}\subset\mathscr{J}$일
    필요충분조건은
    $\varphi(\mathfrak{K})\subset\mathfrak{J}$인 것이다. 이 경우
    $f'$은 사상 $({}^a\varphi',\widetilde{\varphi'})$이며, 여기서
    $\varphi':B/\mathfrak{K}\to A/\mathfrak{J}$은 $\varphi$에서
    몫으로 내려가 얻은 준동형이다.
\end{enumerate}
\end{proposition}

\begin{proof}
\medskip\noindent
\begin{enumerate}
  \item[{\upshape(i)}] $f=(\psi,\theta)$라 하면, 가정에 따라 준동형
    $\theta^\sharp:\psi^*(\mathscr{O}_{\mathfrak{Y}})\to
    \mathscr{O}_{\mathfrak{X}}$에 의한
    $\psi^*(\mathscr{O}_{\mathfrak{Y}})$의 아이디얼층
    $\psi^*(\mathscr{K})$의 상은 $\mathscr{J}$에 포함된다
    (\hyperref[0.4.3.5-ko]{0, 4.3.5}). 따라서 몫으로 내려가면
    $\theta^\sharp$에서 환의 층의 준동형
    \[
      \omega:\psi^*(\mathscr{O}_{\mathfrak{Y}}/\mathscr{K})
      =\psi^*(\mathscr{O}_{\mathfrak{Y}})/\psi^*(\mathscr{K})
      \to\mathscr{O}_{\mathfrak{X}}/\mathscr{J}~;
    \]
    을 얻는다. 또한 모든 $x\in\mathfrak{X}$에 대하여
    $\theta_x^\sharp$가 \emph{국소} 준동형이므로 $\omega_x$도 그러하다.
    따라서 환 달린 공간의 사상 $(\psi,\omega^b)$는
    (\hyperref[I.2.2.1-ko]{2.2.1}) 문제의 요구를 만족하는 환 달린 공간의
    유일한 사상 $f'$이다.
  \item[{\upshape(ii)}] 아핀 형식적 스킴의 사상들과 환의 연속 준동형들
    사이의 표준적이고 함자적인 대응
    (\hyperref[I.10.2.2-ko]{10.2.2})에 따르면, 지금의 경우 관계
    $f^*(\mathscr{K})\mathscr{O}_{\mathfrak{X}}\subset\mathscr{J}$로부터
    $f'=({}^a\varphi',\widetilde{\varphi'})$를 얻는다. 여기서
    $\varphi':B/\mathfrak{K}\to A/\mathfrak{J}$은 다음 도식을 가환이게
    하는 유일한 준동형이다.
    \[
      \label{I.10.5.6.2-ko}
      \xymatrix{
        B\ar[r]^\varphi\ar[d] & A\ar[d]\\
        B/\mathfrak{K}\ar[r]_{\varphi'} & A/\mathfrak{J}
      }
      \tag{10.5.6.2}
    \]
    따라서 $\varphi'$이 존재하면
    $\varphi(\mathfrak{K})\subset\mathfrak{J}$이다. 역으로 이 조건이
    성립한다고 하자. $\varphi'$을 도식
    (\hyperref[I.10.5.6.2-ko]{10.5.6.2})을 가환이게 하는 유일한
    준동형이라 하고 $f'=({}^a\varphi',\widetilde{\varphi'})$로 놓으면,
    도식 (\hyperref[I.10.5.6.1-ko]{10.5.6.1})이 가환임은 명백하다. 이제
    각각 $f$와 $f'$에 대응하는 준동형
    ${}^a\varphi^*(\mathscr{O}_{\mathfrak{Y}})\to
    \mathscr{O}_{\mathfrak{X}}$와
    ${}^a{\varphi'}^*(\mathscr{O}_{\mathfrak{Y}}/\mathscr{K})\to
    \mathscr{O}_{\mathfrak{X}}/\mathscr{J}$을 고려하면, 이 가환성으로부터
    $f^*(\mathscr{K})\mathscr{O}_{\mathfrak{X}}\subset\mathscr{J}$가
    따름을 알 수 있다.
\end{enumerate}
\end{proof}

위에서 정의한 대응 $f\mapsto f'$이 \emph{함자적}임은 명백하다.

\subsection{준스킴들의 귀납극한으로서의 형식적 준스킴.}
\label{subsection:I.10.6-ko}

\begin{env}[10.6.1]
\label{I.10.6.1-ko}
$\mathfrak{X}$를 형식적 준스킴이라 하고,
$(\mathscr{J}_\lambda)$를 $\mathfrak{X}$의 정의 아이디얼들의 기본계라
하자. 각 $\lambda$에 대하여 $f_\lambda$를 표준 사상
$(\mathfrak{X},\mathscr{O}_{\mathfrak{X}}/\mathscr{J}_\lambda)
\to\mathfrak{X}$라 하자(\hyperref[I.10.5.2-ko]{10.5.2}).
$\mathscr{J}_\mu\subset\mathscr{J}_\lambda$이면 표준 준동형
$\mathscr{O}_{\mathfrak{X}}/\mathscr{J}_\mu\to
\mathscr{O}_{\mathfrak{X}}/\mathscr{J}_\lambda$는 보통 의미의 준스킴
사이의 표준 사상

\oldpage[I]{189}

$f_{\mu\lambda}:(\mathfrak{X},
\mathscr{O}_{\mathfrak{X}}/\mathscr{J}_\lambda)\to
(\mathfrak{X},\mathscr{O}_{\mathfrak{X}}/\mathscr{J}_\mu)$를 정의하며,
$f_\lambda=f_\mu\circ f_{\mu\lambda}$가 성립한다. 따라서 준스킴들
$X_\lambda=(\mathfrak{X},
\mathscr{O}_{\mathfrak{X}}/\mathscr{J}_\lambda)$와 사상들
$f_{\mu\lambda}$는 (\hyperref[I.10.4.8-ko]{10.4.8})에 따라 형식적
준스킴의 범주에서 하나의 \emph{귀납계}를 이룬다.
\end{env}

\begin{proposition}[10.6.2]
\label{I.10.6.2-ko}
(\hyperref[I.10.6.1-ko]{10.6.1})의 표기 아래에서 형식적 준스킴
$\mathfrak{X}$와 사상들 $f_\lambda$는 형식적 준스킴의 범주에서 계
$(X_\lambda,f_{\mu\lambda})$의 귀납극한(T, I, 1.8)을 이룬다.
\end{proposition}

\begin{proof}
$\mathfrak{Y}$를 형식적 준스킴이라 하고, 각 지표 $\lambda$에 대하여
사상
\[
  g_\lambda=(\psi_\lambda,\theta_\lambda):X_\lambda\to\mathfrak{Y}
\]
가 주어졌으며, $\mathscr{J}_\mu\subset\mathscr{J}_\lambda$이면
$g_\lambda=g_\mu\circ f_{\mu\lambda}$가 성립한다고 하자. 이 조건과
$X_\lambda$의 정의로부터 우선 모든 $\psi_\lambda$가 바탕공간 사이의
동일한 연속사상 $\psi:\mathfrak{X}\to\mathfrak{Y}$와 일치함을 알 수 있다.
또한 준동형
$\theta_\lambda^\sharp:\psi^*(\mathscr{O}_{\mathfrak{Y}})\to
\mathscr{O}_{X_\lambda}=
\mathscr{O}_{\mathfrak{X}}/\mathscr{J}_\lambda$들은 환의 층의 준동형들의
\emph{사영계}를 이룬다. 따라서 사영극한을 취하면 준동형
\[
  \omega:\psi^*(\mathscr{O}_{\mathfrak{Y}})\to
  \varprojlim\mathscr{O}_{\mathfrak{X}}/\mathscr{J}_\lambda
  =\mathscr{O}_{\mathfrak{X}}
\]
을 얻는다. \emph{환 달린 공간}의 사상 $g=(\psi,\omega^b)$가 도식들
\[
  \label{I.10.6.2.1-ko}
  \xymatrix{
    X_\lambda\ar[rr]^{g_\lambda}\ar[rd]_{f_\lambda} &&
    \mathfrak{Y}\\
    & \mathfrak{X}\ar[ru]_g
  }
  \tag{10.6.2.1}
\]
을 가환이게 하는 \emph{유일한} 사상임은 명백하다. 그러므로 이제 $g$가
\emph{형식적 준스킴}의 사상임을 보이면 된다. 이 문제는
$\mathfrak{X}$와 $\mathfrak{Y}$ 위에서 국소적이므로
$\mathfrak{X}=\operatorname{Spf}(A)$,
$\mathfrak{Y}=\operatorname{Spf}(B)$라 해도 된다. 여기서 $A$와 $B$는
허용환이고 $\mathscr{J}_\lambda=\mathfrak{J}_\lambda^{\Delta}$이며,
$(\mathfrak{J}_\lambda)$는 $A$의 정의 아이디얼들의 기본계이다
(\hyperref[I.10.3.5-ko]{10.3.5}). 이제
$A=\varprojlim A/\mathfrak{J}_\lambda$이므로, 도식
(\hyperref[I.10.6.2.1-ko]{10.6.2.1})을 가환이게 하는 아핀 형식적
스킴의 사상 $g$의 존재는 아핀 형식적 스킴의 사상과 환의 연속
준동형 사이의 전단사 대응 (\hyperref[I.10.2.2-ko]{10.2.2}) 및
사영극한의 정의에서 따른다. 한편 환 달린 공간의 사상으로서 $g$가
유일하므로, 이 사상은 증명 앞부분에서 같은 기호로 나타낸 사상과
일치한다.
\end{proof}

다음 명제는 몇 가지 추가 조건 아래, 보통 의미의 준스킴들로 주어진
귀납계의 귀납극한이 형식적 준스킴의 범주에서 존재함을 보인다.

\begin{proposition}[10.6.3]
\label{I.10.6.3-ko}
$\mathfrak{X}$를 위상공간이라 하고,
$(\mathscr{O}_i,u_{ji})$를 $\mathfrak{X}$ 위의 환의 층들의 사영계라
하자. 그 지표집합은 $\mathbf{N}$이다. $\mathscr{J}_i$를
$u_{0i}:\mathscr{O}_i\to\mathscr{O}_0$의 핵이라 하자. 다음을 가정한다.
\begin{enumerate}
  \item[a)] 환 달린 공간 $(\mathfrak{X},\mathscr{O}_i)$는 준스킴
    $X_i$이다.
  \item[b)] 모든 $x\in\mathfrak{X}$와 모든 $i$에 대하여,
    $\mathfrak{X}$ 안에서 $x$의 열린 근방 $U_i$가 존재하여 제한
    $\mathscr{J}_i|U_i$는 멱영이다.\footnote{\textbf{번역 주.} 인쇄
    원문은 여기서 로만체 $X$를 쓰지만, 출판된 EGA II 정오표는 이를
    프락투어체 $\mathfrak{X}$로 고치도록 지시한다. 그 공식 교정을
    적용했다.}
  \item[c)] 준동형들 $u_{ji}$는 전사이다.
\end{enumerate}

\oldpage[I]{190}

$\mathscr{O}_{\mathfrak{X}}$를 의사이산 환의 층들 $\mathscr{O}_i$의
사영극한인 위상환의 층이라 하고,
$u_i:\mathscr{O}_{\mathfrak{X}}\to\mathscr{O}_i$를 표준 준동형이라
하자. 그러면 위상환 달린 공간
$(\mathfrak{X},\mathscr{O}_{\mathfrak{X}})$는 형식적 준스킴이다.
준동형들 $u_i$는 전사이고, 그 핵들 $\mathscr{J}^{(i)}$는
$\mathfrak{X}$의 정의 아이디얼들의 기본계를 이루며,
$\mathscr{J}^{(0)}$은 아이디얼층들 $\mathscr{J}_i$의 사영극한이다.
\end{proposition}

\begin{proof}
먼저 각 줄기에서 $u_{ji}$는 전사 준동형이며
\emph{따라서 특히} 국소 준동형임에 유의하자. 그러므로
$v_{ij}=(1_{\mathfrak{X}},u_{ji})$는 준스킴의 사상
$X_j\to X_i$ ($i\geq j$)이다
(\hyperref[I.2.2.1-ko]{2.2.1}). 우선 각 $X_i$가
환 $A_i$의 아핀 스킴이라고 가정하자.
\emph{환 준동형}
$\varphi_{ji}:A_i\to A_j$가 존재하여
$u_{ji}=\widetilde{\varphi_{ji}}$를 만족한다
(\hyperref[I.1.7.3-ko]{1.7.3}). 따라서
(\hyperref[I.1.6.3-ko]{1.6.3}), 층 $\mathscr{O}_j$는
$\mathscr{O}_i$-가군으로서 $X_i$ 위에서 준연접이다(그 스칼라
작용은 $u_{ji}$로 정의된다). 이 층은 $A_j$를 $A_i$-가군으로 보아
얻는 가군에 대응하며, 이때 그 가군 구조는 $\varphi_{ji}$로 정한다.
모든 $f\in A_i$에 대하여
$f'=\varphi_{ji}(f)$라 하자. 가정에 따라 열린집합 $D(f)$와
$D(f')$는 $\mathfrak{X}$에서 동일하며,
$\Gamma(D(f),\mathscr{O}_i)=(A_i)_f$에서
$\Gamma(D(f),\mathscr{O}_j)=(A_j)_{f'}$로 가는 준동형 가운데
$u_{ji}$에 대응하는 것은 바로
$(\varphi_{ji})_f$이다
(\hyperref[I.1.6.1-ko]{1.6.1}). 한편 $A_j$를 $A_i$-가군으로 보면,
$(A_j)_{f'}$는 $(A_i)_f$-가군 $(A_j)_f$이다. 따라서
$u_{ji}=\widetilde{\varphi_{ji}}$가 여전히 성립하는데, 이번에는
$\varphi_{ji}$를 $A_i$-가군 준동형으로 본다. 이제 $u_{ji}$가
전사이므로 $\varphi_{ji}$도 전사임을 알 수 있다
(\hyperref[I.1.3.9-ko]{1.3.9}). $\mathfrak{J}_{ji}$가
$\varphi_{ji}$의 핵이면, $u_{ji}$의 핵은 준연접
$\mathscr{O}_i$-가군이고 $\widetilde{\mathfrak{J}_{ji}}$와 같다. 특히
$\mathscr{J}_i=\widetilde{\mathfrak{J}_i}$이고, 여기서
$\mathfrak{J}_i$는 $\varphi_{0i}:A_i\to A_0$의 핵이다. 가정 b)로부터
$\mathscr{J}_i$가 \emph{멱영}임이 따른다. 실제로
$\mathfrak{X}$가 준콤팩트이므로, $\mathfrak{X}$를 유한 개의
열린집합 $U_k$로 덮을 수 있고, 각각
$(\mathscr{J}_i|U_k)^{n_k}=0$을 만족한다.
$n$을 $n_k$들 가운데 가장 큰 수로 택하면
$\mathscr{J}_i^n=0$이다. 따라서 $\mathfrak{J}_i$는 멱영이다
(\hyperref[I.1.3.13-ko]{1.3.13}). 이제 환
$A=\varprojlim A_i$는 허용환이다
(\hyperref[0.7.2.2-ko]{0, 7.2.2}). 또한 표준 준동형
$\varphi_i:A\to A_i$는 전사이다. 그 핵 $\mathfrak{J}^{(i)}$는
$\mathfrak{J}_{ik}$들($k\geq i$)의 사영극한이며,
$\mathfrak{J}^{(i)}$들은 $A$에서 0의 근방들의 기본계를 이룬다.
이 경우 (\hyperref[I.10.6.3-ko]{10.6.3})의 주장들은
(\hyperref[I.10.1.1-ko]{10.1.1})과
(\hyperref[I.10.3.2-ko]{10.3.2})에서 따른다. 실제로
$(\mathfrak{X},\mathscr{O}_{\mathfrak{X}})$는 다름 아닌
$\operatorname{Spf}(A)$이다.

여전히 이 특수한 경우에, $f=(f_i)$가 사영극한
$A=\varprojlim A_i$의 원소라고 하자. 모든 열린집합 $D(f_i)$는
$X_i$에서의 아핀 열린집합이며, $\mathfrak{D}(f)$라는
$\mathfrak{X}$의 열린집합과 동일시된다. 따라서 $X_i$가
$\mathfrak{D}(f)$에 유도하는 준스킴은 아핀 스킴
$\operatorname{Spec}((A_i)_{f_i})$와 동일시된다.

일반적인 경우에는 먼저, 임의의 준콤팩트 열린집합 $U$를
$\mathfrak{X}$에서 택하면 각 $\mathscr{J}_i|U$가 멱영임에
유의하자. 이는 위의
논법으로 알 수 있다. 이제 모든 $x\in\mathfrak{X}$에 대하여,
어떤 열린집합 $U$가 존재하여 $x$를 포함하고 $\mathfrak{X}$에서 열린 근방이며,
\emph{모든 $X_i$에 대하여 아핀 열린집합}임을 보이자. 실제로
$U$를 $X_0$의 아핀 열린집합으로 택하고
$\mathscr{O}_{X_0}=\mathscr{O}_{X_i}/\mathscr{J}_i$임에 유의하자.
앞에서 본 바와 같이 $\mathscr{J}_i|U$가 멱영이므로,
(\hyperref[I.5.1.9-ko]{5.1.9})에 따라 $U$는 각 $X_i$에 대하여도
아핀 열린집합이다. 이제 앞의 조건을 만족하는 모든 $U$에 대하여,
위에서 다룬 아핀 경우에 의해 $(U,\mathscr{O}_{\mathfrak{X}}|U)$는
형식적 준스킴이고, $\mathscr{J}^{(i)}|U$들은 정의 아이디얼들의
기본계를 이루며, $\mathscr{J}^{(0)}|U$는
$\mathscr{J}_i|U$들의 사영극한이다. 이로써 결론을 얻는다.
\end{proof}

\begin{corollary}[10.6.4]
\label{I.10.6.4-ko}
$i\geq j$일 때 $u_{ji}$의 핵이
$\mathscr{J}_i^{j+1}$이고 $\mathscr{J}_1/\mathscr{J}_1^2$가

\oldpage[I]{191}

$\mathscr{O}_0=\mathscr{O}_1/\mathscr{J}_1$-가군으로서 유한
생성이라고 가정하자. 그러면 $\mathfrak{X}$는 아딕 형식적 준스킴이다.
또한 $\mathscr{J}^{(n)}$을
$\mathscr{O}_{\mathfrak{X}}\to\mathscr{O}_n$의 핵이라 하면
$\mathscr{J}^{(n)}=\mathscr{J}^{n+1}$이고
$\mathscr{J}/\mathscr{J}^2$는 $\mathscr{J}_1$과 동형이다. 더욱이
$X_0$가 국소 뇌터(각각 뇌터)이면 $\mathfrak{X}$도 국소 뇌터(각각
뇌터)이다.
\end{corollary}

\begin{proof}
$\mathfrak{X}$와 $X_0$의 바탕공간은 같으므로 문제는 국소적이고, 모든
$X_i$가 아핀이라 해도 된다. 관계
$\mathscr{J}_{ij}=\widetilde{\mathfrak{J}_{ji}}$
((\hyperref[I.10.6.3-ko]{10.6.3})의 표기)를 고려하면
(\hyperref[0.7.2.7-ko]{0, 7.2.7}과
\hyperref[0.7.2.8-ko]{7.2.8})의 대응하는 주장들로 곧바로 환원된다.
여기서 $\mathfrak{J}_1/\mathfrak{J}_1^2$는 유한 생성
$A_0$-가군임에 유의한다
(\hyperref[I.1.3.9-ko]{1.3.9}).
\end{proof}

특히 \emph{모든 국소 뇌터 형식적 준스킴
$\mathfrak{X}$}는
(\hyperref[I.10.6.3-ko]{10.6.3})과
(\hyperref[I.10.6.4-ko]{10.6.4})의 조건을 만족하는 (보통 의미의)
국소 뇌터 준스킴들의 수열 $(X_n)$의 귀납극한이다. 실제로
$\mathscr{J}$를 $\mathfrak{X}$의 정의 아이디얼로 택하고
(\hyperref[I.10.5.4-ko]{10.5.4})
$X_n=(\mathfrak{X},
\mathscr{O}_{\mathfrak{X}}/\mathscr{J}^{n+1})$
로 두면 충분하다
((\hyperref[I.10.5.1-ko]{10.5.1})과
(\hyperref[I.10.6.2-ko]{10.6.2})).

\begin{corollary}[10.6.5]
\label{I.10.6.5-ko}
$A$를 허용환이라 하자. 아핀 형식적 스킴
$\mathfrak{X}=\operatorname{Spf}(A)$가 뇌터이기 위한 필요충분조건은
$A$가 아딕이고 뇌터인 것이다.
\end{corollary}

\begin{proof}
그 조건이 충분함은 명백하다. 역으로 $\mathfrak{X}$가 뇌터라고
가정하고, $\mathfrak{J}$를 $A$의 정의 아이디얼로,
$\mathscr{J}=\mathfrak{J}^\Delta$를 이에 대응하는
$\mathfrak{X}$의 정의 아이디얼로 택하자. 그러면 (보통 의미의) 준스킴
$X_n=(\mathfrak{X},
\mathscr{O}_{\mathfrak{X}}/\mathscr{J}^{n+1})$은 아핀이고 뇌터이다.
따라서 환 $A_n=A/\mathfrak{J}^{n+1}$도 뇌터이고
(\hyperref[I.6.1.3-ko]{6.1.3}), 이로부터
$\mathfrak{J}/\mathfrak{J}^2$가 유한 생성
$A/\mathfrak{J}$-가군임을 알 수 있다. $\mathscr{J}^n$들이
$\mathfrak{X}$의 정의 아이디얼들의 기본계를 이루므로
(\hyperref[I.10.5.1-ko]{10.5.1}),
$\mathscr{O}_{\mathfrak{X}}
=\varprojlim(\mathscr{O}_{\mathfrak{X}}/\mathscr{J}^n)$이다
(\hyperref[I.10.5.3-ko]{10.5.3}). 따라서
(\hyperref[I.10.1.3-ko]{10.1.3})에 의하여 $A$는
$\varprojlim A/\mathfrak{J}^n$과 위상적으로 동형이고, 그러므로
아딕이고 뇌터이다
(\hyperref[0.7.2.8-ko]{0, 7.2.8}).
\end{proof}

\begin{remark}[10.6.6]
\label{I.10.6.6-ko}
(\hyperref[I.10.6.3-ko]{10.6.3})의 표기에서
$\mathscr{F}_i$를 $\mathscr{O}_i$-가군이라 하자. 또한
$i\geq j$에 대하여 $v_{ij}$-사상
$\theta_{ji}:\mathscr{F}_i\to\mathscr{F}_j$가 주어졌고,
$\theta_{kj}\circ\theta_{ji}=\theta_{ki}$가
$k\leq j\leq i$에 대하여 성립한다고 하자. $v_{ij}$의 바탕 연속사상은
항등사상이므로, $\theta_{ji}$는 공간 $\mathfrak{X}$ 위의 아벨군의
층의 준동형이다. 또한 $\mathscr{F}$가 아벨군의 층들의 사영계
$(\mathscr{F}_i)$의 사영극한이면, $\theta_{ji}$들이
$v_{ij}$-사상이라는 사실로부터 사영극한을 취하여 $\mathscr{F}$에
$\mathscr{O}_{\mathfrak{X}}$-가군 구조를 정의할 수 있다. 이 구조를
갖춘 $\mathscr{F}$를, $\theta_{ji}$들에 관한
$\mathscr{O}_i$-가군들의 계 $(\mathscr{F}_i)$의
\emph{사영극한}이라 하겠다. 특별히
$v_{ij}^*(\mathscr{F}_i)=\mathscr{F}_j$이고 $\theta_{ji}$가
\emph{항등사상}인 경우에는, 줄여서 $\mathscr{F}$를 계
$(\mathscr{F}_i)$의 사영극한이라 하겠다. 이때
$v_{ij}^*(\mathscr{F}_i)=\mathscr{F}_j$는 $j\leq i$에 대하여
성립하며, $\theta_{ji}$들은 언급하지 않는다.
\end{remark}

\begin{env}[10.6.7]
\label{I.10.6.7-ko}
$\mathfrak{X}$와 $\mathfrak{Y}$를 두 형식적 준스킴이라 하고,
$\mathscr{J}$와 $\mathscr{K}$를 각각 $\mathfrak{X}$와
$\mathfrak{Y}$의 정의 아이디얼이라 하자. 또한
$f:\mathfrak{X}\to\mathfrak{Y}$를
$f^*(\mathscr{K})\mathscr{O}_{\mathfrak{X}}\subset\mathscr{J}$를
만족하는 형식적 준스킴의 사상이라 하자. 그러면 모든 정수 $n>0$에 대하여
$f^*(\mathscr{K}^n)\mathscr{O}_{\mathfrak{X}}
=(f^*(\mathscr{K})\mathscr{O}_{\mathfrak{X}})^n
\subset\mathscr{J}^n$이다. 따라서
(\hyperref[I.10.5.6-ko]{10.5.6})에 의하여 $f$로부터 (보통 의미의)
준스킴의 사상 $f_n:X_n\to Y_n$을 유도할 수 있다. 여기서
$X_n=(\mathfrak{X},
\mathscr{O}_{\mathfrak{X}}/\mathscr{J}^{n+1})$,
$Y_n=(\mathfrak{Y},
\mathscr{O}_{\mathfrak{Y}}/\mathscr{K}^{n+1})$로 둔다. 정의로부터
곧바로 다음 도식들이
\[
  \label{I.10.6.7.1-ko}
  \xymatrix{
    X_m\ar[r]^{f_m}\ar[d] &
    Y_m\ar[d]\\
    X_n\ar[r]_{f_n} &
    Y_n
  }
  \tag{10.6.7.1}
\]

\oldpage[I]{192}

$m\leq n$일 때 가환함을 알 수 있다. 다시 말해, $(f_n)$은 사상들의
\emph{귀납계}이다.
\end{env}

\begin{env}[10.6.8]
\label{I.10.6.8-ko}
역으로, $(X_n)$과 $(Y_n)$을 각각
(\hyperref[I.10.6.3-ko]{10.6.3})의 조건 b)와 c)를 만족하는 (보통
의미의) 준스킴들의 귀납계라 하고, $\mathfrak{X}$와
$\mathfrak{Y}$를 각각 그 귀납극한이라 하자. 귀납극한의 정의에 따라,
임의의 수열 $(f_n)$이 $X_n\to Y_n$인 사상들로 이루어진 귀납계이면 그
\emph{귀납극한} $f:\mathfrak{X}\to\mathfrak{Y}$가 존재한다. 이 사상은
다음 도식들을 가환이게 하는 유일한 형식적 준스킴의 사상이다.
\[
  \xymatrix{
    X_n\ar[r]^{f_n}\ar[d] &
    Y_n\ar[d]\\
    \mathfrak{X}\ar[r]_f &
    \mathfrak{Y}
  }
\]
\end{env}

\begin{proposition}[10.6.9]
\label{I.10.6.9-ko}
$\mathfrak{X}$와 $\mathfrak{Y}$를 두 국소 뇌터 형식적 준스킴이라
하고, $\mathscr{J}$와 $\mathscr{K}$를 각각 $\mathfrak{X}$와
$\mathfrak{Y}$의 정의 아이디얼이라 하자.
(\hyperref[I.10.6.7-ko]{10.6.7})에서 정의한 대응
$f\mapsto(f_n)$은
사상 $f:\mathfrak{X}\to\mathfrak{Y}$ 가운데
$f^*(\mathscr{K})\mathscr{O}_{\mathfrak{X}}\subset\mathscr{J}$를
만족하는 것들의 집합과, 도식들
(\hyperref[I.10.6.7.1-ko]{10.6.7.1})을 가환이게 하는 사상들의 수열
$(f_n)$의 집합 사이의 전단사 대응이다.
\end{proposition}

\begin{proof}
$f$가 그러한 사상 수열의 귀납극한이면
$f^*(\mathscr{K})\mathscr{O}_{\mathfrak{X}}\subset\mathscr{J}$임을
보여야 한다. 문제는 $\mathfrak{X}$와 $\mathfrak{Y}$ 위에서
국소적이므로, $\mathfrak{X}=\operatorname{Spf}(A)$,
$\mathfrak{Y}=\operatorname{Spf}(B)$인 아핀 경우로 한정해도 된다.
이때 $A$와 $B$는 아딕이고 뇌터이다. 또한
$\mathscr{J}=\mathfrak{J}^\Delta$,
$\mathscr{K}=\mathfrak{K}^\Delta$이다. 여기서 $\mathfrak{J}$와
$\mathfrak{K}$는 각각 $A$와 $B$의 정의 아이디얼이다. 그러면
$X_n=\operatorname{Spec}(A_n)$,
$Y_n=\operatorname{Spec}(B_n)$이고, 여기서
$A_n=A/\mathfrak{J}^{n+1}$,
$B_n=B/\mathfrak{K}^{n+1}$이며, 이는
(\hyperref[I.10.3.6-ko]{10.3.6})과
(\hyperref[I.10.3.2-ko]{10.3.2})에 따른다. 또한
$f_n=({}^a\varphi_n,\widetilde{\varphi_n})$이며, 환 준동형
$\varphi_n:B_n\to A_n$들은 사영계를 이룬다. 따라서
$f=({}^a\varphi,\widetilde{\varphi})$이고,
$\varphi=\varprojlim\varphi_n$이다. 이제 $m=0$일 때 도식
(\hyperref[I.10.6.7.1-ko]{10.6.7.1})의 가환성으로부터 조건
$\varphi_n(\mathfrak{K}/\mathfrak{K}^{n+1})
\subset\mathfrak{J}/\mathfrak{J}^{n+1}$이 모든 $n$에 대하여
성립한다. 따라서 사영극한을 취하면
$\varphi(\mathfrak{K})\subset
\mathfrak{J}$이고, 이로부터
$f^*(\mathscr{K})\mathscr{O}_{\mathfrak{X}}\subset\mathscr{J}$가
따른다
(\hyperref[I.10.5.6-ko]{10.5.6, (ii)}).
\end{proof}

\begin{corollary}[10.6.10]
\label{I.10.6.10-ko}
$\mathfrak{X}$와 $\mathfrak{Y}$를 두 국소 뇌터 형식적 준스킴이라 하고,
$\mathscr{T}$를 $\mathfrak{X}$의 가장 큰 정의 아이디얼이라 하자
(\hyperref[I.10.5.4-ko]{10.5.4}).
\begin{enumerate}
  \item[{\upshape(i)}] $\mathfrak{Y}$의 임의의 정의 아이디얼
    $\mathscr{K}$와 임의의 사상
    $f:\mathfrak{X}\to\mathfrak{Y}$에 대하여,
    $f^*(\mathscr{K})\mathscr{O}_{\mathfrak{X}}\subset\mathscr{T}$이다.
  \item[{\upshape(ii)}]
    $\operatorname{Hom}(\mathfrak{X},\mathfrak{Y})$과 사상 수열
    $(f_n)$ 가운데 도식들
    (\hyperref[I.10.6.7.1-ko]{10.6.7.1})을 가환이게 하는 것들의 집합
    사이에는 표준 전단사 대응이 있다. 여기서
    $X_n=(\mathfrak{X},
    \mathscr{O}_{\mathfrak{X}}/\mathscr{T}^{n+1})$,
    $Y_n=(\mathfrak{Y},
    \mathscr{O}_{\mathfrak{Y}}/\mathscr{K}^{n+1})$이다.
\end{enumerate}
\end{corollary}

\begin{proof}
(ii)는 (i)와 (\hyperref[I.10.6.9-ko]{10.6.9})에서 곧바로 따른다.
(i)를 보이기 위해서는 $\mathfrak{X}=\operatorname{Spf}(A)$,
$\mathfrak{Y}=\operatorname{Spf}(B)$이고 $A$와 $B$가 뇌터이며,
$\mathscr{T}=\mathfrak{T}^\Delta$,
$\mathscr{K}=\mathfrak{K}^\Delta$인 경우로 한정해도 된다. 여기서
$\mathfrak{T}$는 $A$의 가장 큰 정의 아이디얼이고
$\mathfrak{K}$는 $B$의 정의 아이디얼이다.
$f=({}^a\varphi,\widetilde{\varphi})$라 하자. 여기서
$\varphi:B\to A$는 연속 준동형이다. $\mathfrak{K}$의 원소들은
위상적으로 멱영이고
(\hyperref[0.7.1.4-ko]{0, 7.1.4, (ii)}), 따라서
$\varphi(\mathfrak{K})$의 원소들도 위상적으로 멱영이다. 그러므로
$\mathfrak{T}$가 $A$의 위상적으로 멱영인 원소들의 집합이므로
(\hyperref[0.7.1.6-ko]{0, 7.1.6}),
$\varphi(\mathfrak{K})\subset\mathfrak{T}$이다. 따라서
(\hyperref[I.10.5.6-ko]{10.5.6, (ii)})에 의해 결론이 따른다.
\end{proof}

\begin{corollary}[10.6.11]
\label{I.10.6.11-ko}
$\mathfrak{S}$, $\mathfrak{X}$, $\mathfrak{Y}$를 세 국소 뇌터
형식적 준스킴이라 하고,
$f:\mathfrak{X}\to\mathfrak{S}$,
$g:\mathfrak{Y}\to\mathfrak{S}$를 $\mathfrak{X}$와
$\mathfrak{Y}$를 $\mathfrak{S}$-형식적 준스킴으로 만드는 사상이라 하자.
$\mathscr{J}$ (각각 $\mathscr{K}$, $\mathscr{L}$)를
$\mathfrak{S}$ (각각 $\mathfrak{X}$, $\mathfrak{Y}$)의 정의 아이디얼이라 하고,
$f^*(\mathscr{J})\mathscr{O}_{\mathfrak{X}}\subset\mathscr{K}$,
$g^*(\mathscr{J})\mathscr{O}_{\mathfrak{Y}}=\mathscr{L}$이라 가정하자.
다음과 같이 둔다.
$S_n=(\mathfrak{S},
\mathscr{O}_{\mathfrak{S}}/\mathscr{J}^{n+1})$,
$X_n=(\mathfrak{X},
\mathscr{O}_{\mathfrak{X}}/\mathscr{K}^{n+1})$,
$Y_n=(\mathfrak{Y},
\mathscr{O}_{\mathfrak{Y}}/\mathscr{L}^{n+1})$.
그러면

\oldpage[I]{193}

$\operatorname{Hom}_{\mathfrak{S}}(\mathfrak{X},\mathfrak{Y})$과 수열
$(u_n)$ 가운데 각 항이 $S_n$-사상
$u_n:X_n\to Y_n$이고 도식들
(\hyperref[I.10.6.7.1-ko]{10.6.7.1})을 가환이게 하는 것들의 집합
사이에는 표준 전단사 대응이 있다.
\end{corollary}

\begin{proof}
임의의 $\mathfrak{S}$-사상
$u:\mathfrak{X}\to\mathfrak{Y}$에 대하여, 정의에 의해
$f=g\circ u$이므로
\[
  u^*(\mathscr{L})\mathscr{O}_{\mathfrak{X}}
  =u^*(g^*(\mathscr{J})\mathscr{O}_{\mathfrak{Y}})
    \mathscr{O}_{\mathfrak{X}}
  =f^*(\mathscr{J})\mathscr{O}_{\mathfrak{X}}
  \subset\mathscr{K}
\]
이다. 따라서 이 따름정리는
(\hyperref[I.10.6.9-ko]{10.6.9})에서 따른다.
\end{proof}

$m\leq n$일 때, 사상 $f_n:X_n\to Y_n$ 하나가 주어지면
도식 (\hyperref[I.10.6.7.1-ko]{10.6.7.1})을 가환이게 하는 사상
$f_m:X_m\to Y_m$가 유일하게 정해진다. 이는 아핀 경우로 환원하면
곧바로 알 수 있다. 이로써 사상
$\varphi_{mn}:\operatorname{Hom}_{S_n}(X_n,Y_n)
\to\operatorname{Hom}_{S_m}(X_m,Y_m)$가 정의되고,
$\operatorname{Hom}_{S_n}(X_n,Y_n)$들은 $\varphi_{mn}$들에 관하여
\emph{집합들의 사영계}를 이룬다.
(\hyperref[I.10.6.11-ko]{10.6.11})은 또한 다음과 같은 표준 전단사 사상이
존재한다고 진술할 수도 있다.
\[
  \operatorname{Hom}_{\mathfrak{S}}(\mathfrak{X},\mathfrak{Y})
  \xrightarrow{\sim}
  \varprojlim_n\operatorname{Hom}_{S_n}(X_n,Y_n)
\]

\subsection{형식적 준스킴의 곱.}
\label{subsection:I.10.7-ko}

\begin{env}[10.7.1]
\label{I.10.7.1-ko}
$\mathfrak{S}$를 형식적 준스킴이라 하자.
$\mathfrak{S}$-형식적 준스킴들은 하나의 범주를 이루므로,
$\mathfrak{S}$-형식적 준스킴의 \emph{곱}이라는 개념을 정의할 수 있다.
\end{env}

\begin{proposition}[10.7.2]
\label{I.10.7.2-ko}
$\mathfrak{X}=\operatorname{Spf}(B)$,
$\mathfrak{Y}=\operatorname{Spf}(C)$를 아핀 형식적 스킴
$\mathfrak{S}=\operatorname{Spf}(A)$ 위의 두 아핀 형식적 스킴이라 하자.
$\mathfrak{Z}=\operatorname{Spf}(B\widehat{\otimes}_A C)$라 하고,
$p_1$, $p_2$를 $B$와 $C$에서
$B\widehat{\otimes}_A C$로 가는 표준 연속 $A$-준동형
$\rho$, $\sigma$에 각각 대응하는
(\hyperref[I.10.2.2-ko]{10.2.2}) $\mathfrak{S}$-사상이라 하자.
그러면 $(\mathfrak{Z},p_1,p_2)$는 $\mathfrak{S}$ 위의 아핀 형식적 스킴
$\mathfrak{X}$와 $\mathfrak{Y}$의 곱이다.
\end{proposition}

\begin{proof}
(\hyperref[I.10.4.6-ko]{10.4.6})에 따라, 위상 $A$-대수인 허용환 $D$와
임의의 연속 $A$-준동형
$\varphi:B\widehat{\otimes}_A C\to D$에 대하여, 쌍
$(\varphi\circ\rho,\varphi\circ\sigma)$를 대응시키는 사상
\[
  \operatorname{Hom}_A(B\widehat{\otimes}_A C,D)
  \xrightarrow{\sim}
  \operatorname{Hom}_A(B,D)\times\operatorname{Hom}_A(C,D)
\]
이 전단사임을 확인하면 충분하다. 이는 다름 아닌 완비 텐서곱의 보편 성질
(\hyperref[0.7.7.6-ko]{0, 7.7.6})이다.
\end{proof}

\begin{proposition}[10.7.3]
\label{I.10.7.3-ko}
두 $\mathfrak{S}$-형식적 준스킴 $\mathfrak{X}$,
$\mathfrak{Y}$가 주어지면, 곱
$\mathfrak{X}\times_{\mathfrak{S}}\mathfrak{Y}$가 존재한다.
\end{proposition}

\begin{proof}
증명은 (\hyperref[I.3.2.6-ko]{3.2.6})의 증명에서 아핀 스킴(각각
아핀 열린집합)을 아핀 형식적 스킴(각각 아핀 형식적 열린집합)으로
바꾸고, 명제 (\hyperref[I.3.2.2-ko]{3.2.2})를
(\hyperref[I.10.7.2-ko]{10.7.2})로 바꾸면 그대로 얻어진다.
\end{proof}

준스킴의 곱에 관한 모든 형식적 성질
(\hyperref[I.3.2.7-ko]{3.2.7}과
\hyperref[I.3.2.8-ko]{3.2.8},
\hyperref[I.3.3.1-ko]{3.3.1}부터
\hyperref[I.3.3.12-ko]{3.3.12}까지)은
형식적 준스킴의 곱에도 아무 변경 없이 성립한다.

\begin{env}[10.7.4]
\label{I.10.7.4-ko}
$\mathfrak{S}$, $\mathfrak{X}$, $\mathfrak{Y}$를 세 형식적 준스킴이라
하고, $f:\mathfrak{X}\to\mathfrak{S}$,
$g:\mathfrak{Y}\to\mathfrak{S}$를 두 사상이라 하자.
$\mathfrak{S}$, $\mathfrak{X}$, $\mathfrak{Y}$ 안에 각각 정의
아이디얼들의 기본계 $(\mathscr{J}_\lambda)$,
$(\mathscr{K}_\lambda)$, $(\mathscr{L}_\lambda)$가 존재한다고 가정하자.
이 세 기본계의 지표집합은 모두 $I$이며,
$f^*(\mathscr{J}_\lambda)\mathscr{O}_{\mathfrak{X}}
\subset\mathscr{K}_\lambda$ 및
$g^*(\mathscr{J}_\lambda)\mathscr{O}_{\mathfrak{Y}}
\subset\mathscr{L}_\lambda$가 모든 $\lambda$에 대하여 성립한다고 하자.
다음과 같이 놓자.
$S_\lambda=(\mathfrak{S},
\mathscr{O}_{\mathfrak{S}}/\mathscr{J}_\lambda)$,
$X_\lambda=(\mathfrak{X},
\mathscr{O}_{\mathfrak{X}}/\mathscr{K}_\lambda)$,
$Y_\lambda=(\mathfrak{Y},
\mathscr{O}_{\mathfrak{Y}}/\mathscr{L}_\lambda)$.
$\mathscr{J}_\mu\subset\mathscr{J}_\lambda$,
$\mathscr{K}_\mu\subset\mathscr{K}_\lambda$,
$\mathscr{L}_\mu\subset\mathscr{L}_\lambda$인 경우,
$S_\lambda$, $X_\lambda$, $Y_\lambda$는 각각
$S_\mu$, $X_\mu$, $Y_\mu$의 닫힌 부분준스킴이고, 각각은 대응하는
준스킴과 같은

\oldpage[I]{194}

바탕공간을 갖는다(\hyperref[I.10.6.1-ko]{10.6.1}).
$S_\lambda\to S_\mu$가 준스킴의 단사사상이므로, 먼저 두 곱
$X_\lambda\times_{S_\lambda}Y_\lambda$와
$X_\lambda\times_{S_\mu}Y_\lambda$가 동일함을 알 수 있고
(\hyperref[I.3.2.4-ko]{3.2.4}), 이어서
$X_\lambda\times_{S_\mu}Y_\lambda$는
$X_\mu\times_{S_\mu}Y_\mu$의 닫힌 부분준스킴으로 동일시되며, 이
부분준스킴은 후자와 같은 바탕공간을 갖는다
(\hyperref[I.4.3.1-ko]{4.3.1}). 이제 곱
$\mathfrak{X}\times_{\mathfrak{S}}\mathfrak{Y}$는 (보통 의미의) 준스킴
$X_\lambda\times_{S_\lambda}Y_\lambda$들의 \emph{귀납극한}이다. 실제로
(\hyperref[I.10.6.2-ko]{10.6.2})에서와 같이
$\mathfrak{S}$, $\mathfrak{X}$, $\mathfrak{Y}$가 아핀 형식적 스킴인
경우로 환원할 수 있다. (\hyperref[I.10.5.6-ko]{10.5.6, (ii)})와
$\mathfrak{S}$, $\mathfrak{X}$, $\mathfrak{Y}$의 정의 아이디얼들의
기본계에 관한 가정을 고려하면, 이 주장은 두 대수의 완비 텐서곱의
정의(\hyperref[0.7.7.1-ko]{0, 7.7.1})에서 곧바로 따른다.

또한 $\mathfrak{Z}$를 $\mathfrak{S}$-형식적 준스킴이라 하고,
$(\mathscr{M}_\lambda)$를 $\mathfrak{Z}$의 정의 아이디얼들의 기본계라
하되 그 지표집합은 $I$라 하자. 또
$u:\mathfrak{Z}\to\mathfrak{X}$,
$v:\mathfrak{Z}\to\mathfrak{Y}$를 두 $\mathfrak{S}$-사상이라 하고,
$u^*(\mathscr{K}_\lambda)\mathscr{O}_{\mathfrak{Z}}
\subset\mathscr{M}_\lambda$ 및
$v^*(\mathscr{L}_\lambda)\mathscr{O}_{\mathfrak{Z}}
\subset\mathscr{M}_\lambda$라고 하자.
$Z_\lambda=(\mathfrak{Z},
\mathscr{O}_{\mathfrak{Z}}/\mathscr{M}_\lambda)$라 놓고,
$u_\lambda:Z_\lambda\to X_\lambda$와
$v_\lambda:Z_\lambda\to Y_\lambda$를 $S_\lambda$-사상 가운데 각각
$u$와 $v$에 대응하는 것이라 하면
(\hyperref[I.10.5.6-ko]{10.5.6}), $(u,v)_{\mathfrak{S}}$가
$S_\lambda$-사상 $(u_\lambda,v_\lambda)_{S_\lambda}$들의
귀납극한임을 곧바로 확인할 수 있다.

이 번호의 고찰은 특히 $\mathfrak{S}$, $\mathfrak{X}$,
$\mathfrak{Y}$가 국소 뇌터인 경우, 각각 하나의 정의 아이디얼의
거듭제곱들로 이루어진 계를 정의 아이디얼들의 기본계로 택하면 적용된다
(\hyperref[I.10.5.1-ko]{10.5.1}). 그러나
$\mathfrak{X}\times_{\mathfrak{S}}\mathfrak{Y}$가 반드시 국소 뇌터인
것은 아니다(다만 (\agref{I.10.13.5-ko}{10.13.5})를 보라).
\end{env}

\subsection{닫힌 부분집합을 따른 준스킴의 형식적 완비화.}
\label{subsection:I.10.8-ko}

\begin{env}[10.8.1]
\label{I.10.8.1-ko}
$X$를 (보통 의미의) 국소 뇌터 준스킴이라 하고, $X'$을 $X$의
바탕공간의 닫힌 부분집합이라 하자. $\Phi$로 다음 조건을 만족하는
아이디얼층들의 집합을 나타내자. $\mathscr{J}$는
$\mathscr{O}_X$ 안의 연접 아이디얼층이고, $\mathscr{O}_X/\mathscr{J}$의
지지집합은 $X'$이다. 집합 $\Phi$는 비어 있지
않으며(\hyperref[I.5.2.1-ko]{5.2.1},
\hyperref[I.4.1.4-ko]{4.1.4}와
\hyperref[I.6.1.1-ko]{6.1.1}), 관계 $\supset$로 순서를 주기로 한다.
\end{env}

\begin{lemma}[10.8.2]
\label{I.10.8.2-ko}
순서집합 $\Phi$는 여과집합이다. $X$가 뇌터이면, 모든
$\mathscr{J}_0\in\Phi$에 대하여 거듭제곱 $\mathscr{J}_0^n$
($n>0$)으로 이루어진 집합은 $\Phi$에서 공종이다.
\end{lemma}

\begin{proof}
실제로 $\mathscr{J}_1$과 $\mathscr{J}_2$가 $\Phi$에 속하고
$\mathscr{J}=\mathscr{J}_1\cap\mathscr{J}_2$라 놓으면,
$\mathscr{O}_X$가 연접이므로 $\mathscr{J}$도 연접이고
(\hyperref[I.6.1.1-ko]{6.1.1}과
\hyperref[0.5.3.4-ko]{0, 5.3.4}), 모든 $x\in X$에 대하여
$\mathscr{J}_x=(\mathscr{J}_1)_x\cap(\mathscr{J}_2)_x$이다. 따라서
$x\notin X'$이면 $\mathscr{J}_x=\mathscr{O}_x$이고
$x\in X'$이면 $\mathscr{J}_x\neq\mathscr{O}_x$이므로
$\mathscr{J}\in\Phi$이다. 한편 $X$가 뇌터이고
$\mathscr{J}_0$와 $\mathscr{J}$가 $\Phi$에 속하면,
$\mathscr{J}_0^n(\mathscr{O}_X/\mathscr{J})=0$이 되는 정수
$n>0$이 존재한다(\hyperref[I.9.3.4-ko]{9.3.4}). 이는
$\mathscr{J}_0^n\subset\mathscr{J}$를 뜻한다.
\end{proof}

\begin{env}[10.8.3]
\label{I.10.8.3-ko}
이제 $\mathscr{F}$를 연접 $\mathscr{O}_X$-가군이라 하자. 모든
$\mathscr{J}\in\Phi$에 대하여
$\mathscr{F}\otimes_{\mathscr{O}_X}(\mathscr{O}_X/\mathscr{J})$는
연접 $\mathscr{O}_X$-가군이고
(\hyperref[I.9.1.1-ko]{9.1.1}), 그 지지집합은 $X'$에 포함된다. 우리는
이 층을 대개 $X'$으로 제한한 층과 동일시하겠다. $\mathscr{J}$가
$\Phi$를 달릴 때 이 층들은 아벨군의 층들의
\emph{사영계}를 이룬다.
\end{env}

\begin{definition}[10.8.4]
\label{I.10.8.4-ko}
국소 뇌터 준스킴 $X$의 닫힌 부분집합 $X'$과 연접
$\mathscr{O}_X$-가군 $\mathscr{F}$가 주어졌을 때, 층

\oldpage[I]{195}

$\varprojlim_{\Phi}
(\mathscr{F}\otimes_{\mathscr{O}_X}
(\mathscr{O}_X/\mathscr{J}))$의 $X'$으로의 제한을
\emph{$X'$을 따른 $\mathscr{F}$의 완비화}라 하고
$\mathscr{F}_{/X'}$로 나타내며, 혼동의 우려가 없으면
$\widehat{\mathscr{F}}$로도 나타낸다. 이 완비화의 $X'$ 위 단면들을
\emph{$X'$을 따른 $\mathscr{F}$의 형식적 단면}이라 한다.
\end{definition}

모든 열린집합 $U\subset X$에 대하여 다음이 곧바로 성립한다.
$(\mathscr{F}|U)_{/(U\cap X')}=(\mathscr{F}_{/X'})|(U\cap X')$.

사영극한을 취하면 $(\mathscr{O}_X)_{/X'}$가 환의 층이고
$\mathscr{F}_{/X'}$를 $(\mathscr{O}_X)_{/X'}$-가군으로 볼 수 있음은
명백하다. 또한 $X'$의 위상에는 준콤팩트 열린집합들로 이루어진 기저가
있으므로, $(\mathscr{O}_X)_{/X'}$ (각각 $\mathscr{F}_{/X'}$)를
\emph{의사이산} 환의 층들 (각각 의사이산 군의 층들)
$\mathscr{O}_X/\mathscr{J}$ (각각
$\mathscr{F}\otimes_{\mathscr{O}_X}
(\mathscr{O}_X/\mathscr{J})=\mathscr{F}/\mathscr{J}\mathscr{F}$)의
사영극한인 \emph{위상환의 층} (각각 \emph{위상군의 층})으로 볼 수
있다. 이때 사영극한을 취하면 $\mathscr{F}_{/X'}$는
\emph{위상 $(\mathscr{O}_X)_{/X'}$-가군}이 된다
(\hyperref[0.3.8.1-ko]{0, 3.8.1}과
\hyperref[0.3.8.2-ko]{3.8.2}). 모든 \emph{준콤팩트} 열린집합
$U\subset X$에 대하여
$\Gamma(U\cap X',(\mathscr{O}_X)_{/X'})$ (각각
$\Gamma(U\cap X',\mathscr{F}_{/X'})$)는 이산환들 (각각 이산군들)
$\Gamma(U,\mathscr{O}_X/\mathscr{J})$ (각각
$\Gamma(U,\mathscr{F}/\mathscr{J}\mathscr{F})$)의 사영극한임을
상기하자.

이제 $u:\mathscr{F}\to\mathscr{G}$가 $\mathscr{O}_X$-가군 사이의
준동형이면, 모든 $\mathscr{J}\in\Phi$에 대하여 표준적으로 준동형
$u_{\mathscr{J}}:
\mathscr{F}\otimes_{\mathscr{O}_X}(\mathscr{O}_X/\mathscr{J})
\to
\mathscr{G}\otimes_{\mathscr{O}_X}(\mathscr{O}_X/\mathscr{J})$를 얻으며,
이 준동형들은 사영계를 이룬다. 사영극한을 취하고 $X'$에 제한하면
연속 $(\mathscr{O}_X)_{/X'}$-준동형
$\mathscr{F}_{/X'}\to\mathscr{G}_{/X'}$를 얻는다. 이를
$u_{/X'}$ 또는 $\widehat u$로 나타내며,
\emph{$X'$을 따른 준동형 $u$의 완비화}라 한다. 또한
$v:\mathscr{G}\to\mathscr{H}$가 $\mathscr{O}_X$-가군 사이의
또 하나의 준동형이면
$(v\circ u)_{/X'}=(v_{/X'})\circ(u_{/X'})$임은 명백하다. 따라서
$\mathscr{F}_{/X'}$는 $\mathscr{F}$에 관한
\emph{가법적 공변 함자}이며, 연접 $\mathscr{O}_X$-가군의 범주에서
위상 $(\mathscr{O}_X)_{/X'}$-가군의 범주에 값을 갖는다.

\begin{proposition}[10.8.5]
\label{I.10.8.5-ko}
$(\mathscr{O}_X)_{/X'}$의 지지집합은 $X'$이다. 위상환 달린 공간
$(X',(\mathscr{O}_X)_{/X'})$는 국소 뇌터 형식적 준스킴이며, 모든
$\mathscr{J}\in\Phi$에 대하여 $\mathscr{J}_{/X'}$는 이 형식적
준스킴의 정의 아이디얼이다. $X=\operatorname{Spec}(A)$가 뇌터 환을
갖는 아핀 스킴이고, $\mathscr{J}=\widetilde{\mathfrak{J}}$이며,
$\mathfrak{J}$가 $A$의 아이디얼이고 $X'=V(\mathfrak{J})$이면,
$(X',(\mathscr{O}_X)_{/X'})$는 표준적으로
$\operatorname{Spf}(\widehat A)$와 동일시된다. 여기서 $\widehat A$는
$A$에 $\mathfrak{J}$-준아딕 위상을 주어 얻는 분리 완비화이다.
\end{proposition}

\begin{proof}
마지막 주장만 증명하면 충분함은 명백하다.
(\hyperref[0.7.3.3-ko]{0, 7.3.3})에 의해, $\mathfrak{J}$의
$\mathfrak{J}$-준아딕 위상에 관한 분리 완비화
$\widehat{\mathfrak{J}}$는 $\widehat A$의 아이디얼
$\mathfrak{J}\widehat A$와 동일시된다. 또한 $\widehat A$는
뇌터 $\widehat{\mathfrak{J}}$-아딕환이며
$\widehat A/\widehat{\mathfrak{J}}^n=A/\mathfrak{J}^n$임을 안다
(\hyperref[0.7.2.6-ko]{0, 7.2.6}). 이 마지막 관계에서
$\widehat A$의 열린 소 아이디얼들은
$\widehat{\mathfrak p}=\mathfrak p\widehat A$ 꼴의 아이디얼들임을 알 수
있다. 여기서 $\mathfrak p$는 $A$의 소 아이디얼이고
$\mathfrak{J}$를 포함하며, $\widehat{\mathfrak p}\cap A=\mathfrak p$이다. 따라서
바탕공간으로서 $\operatorname{Spf}(\widehat A)=X'$이다. 또한
$\mathscr{O}_X/\mathscr{J}^n=(A/\mathfrak{J}^n)^{\sim}$이므로, 명제는
정의들에서 곧바로 따른다.
\end{proof}

이렇게 정의한 형식적 준스킴을 \emph{$X'$을 따른 $X$의 완비화}라 하고
$X_{/X'}$로 나타내며, 혼동의 우려가 없으면 $\widehat X$로도 나타낸다.
$X'=X$로 잡을 때에는 $\mathscr{J}=0$을 택할 수 있으므로
$X_{/X}=X$이다.

$U$가 $X$의 열린집합 위에 유도된 부분준스킴이면,
$U_{/(U\cap X')}$는 $X_{/X'}$가 $X'$의 열린집합 $U\cap X'$ 위에
유도하는 형식적 부분준스킴과 표준적으로 동일시됨은 명백하다.

\begin{corollary}[10.8.6]
\label{I.10.8.6-ko}
(보통 의미의) 준스킴 $\widehat X_{\mathrm{red}}$는 $X$의 유일한 축소
부분준스킴으로서, 그 바탕공간은 $X'$이다
(\hyperref[I.5.2.1-ko]{5.2.1}). $\widehat X$가 뇌터이기 위한
필요충분조건은 $\widehat X_{\mathrm{red}}$가 뇌터인 것이며, 더 나아가
$X$가 뇌터이면 전자도 뇌터이다.
\end{corollary}

\oldpage[I]{196}

$\widehat X_{\mathrm{red}}$의 결정은 국소적이므로
(\hyperref[I.10.5.4-ko]{10.5.4}), 다시 $X$가 뇌터 환을 갖는 아핀 스킴인
경우로 한정할 수 있다. (\hyperref[I.10.8.5-ko]{10.8.5})의 표기를
쓰면, 아이디얼 $\mathfrak{T}$는 $\widehat A$의 위상적으로 멱영인
원소들로 이루어지며, 표준 사상
$\widehat A\to\widehat A/\widehat{\mathfrak{J}}=A/\mathfrak{J}$에 의한
$A/\mathfrak{J}$의 멱영근기의 역상이다
(\hyperref[0.7.1.3-ko]{0, 7.1.3}). 따라서
$\widehat A/\mathfrak{T}$는 $A/\mathfrak{J}$를 그 멱영근기로 나눈
몫과 동형이다. 그러므로 첫 번째 주장은
(\hyperref[I.10.5.4-ko]{10.5.4})와
(\hyperref[I.5.1.1-ko]{5.1.1})에서 따른다.
$\widehat X_{\mathrm{red}}$가 뇌터이면 그 바탕공간 $X'$도 뇌터이므로,
$X'_n=\operatorname{Spec}(\mathscr{O}_X/\mathscr{J}^n)$은 뇌터이고
(\hyperref[I.6.1.2-ko]{6.1.2}), $\widehat X$도 뇌터이다
(\hyperref[I.10.6.4-ko]{10.6.4}). 역은
(\hyperref[I.6.1.2-ko]{6.1.2})에 의해 명백하다.

\begin{env}[10.8.7]
\label{I.10.8.7-ko}
표준 준동형 $\mathscr{O}_X\to\mathscr{O}_X/\mathscr{J}$
($\mathscr{J}\in\Phi$)들은 사영계를 이루므로, 사영극한을 취하면 환의
층의 준동형
$\theta:\mathscr{O}_X\to\psi_*((\mathscr{O}_X)_{/X'})
=\varprojlim_\Phi(\mathscr{O}_X/\mathscr{J})$를 얻는다. 여기서 $\psi$는
바탕공간 사이의 표준 포함 사상 $X'\to X$이다. 환 달린 공간의 사상
\[
  (\psi,\theta):X_{/X'}\to X
\]
을 $i$ (또는 $i_X$)로 나타내고 표준 사상이라 하겠다.

텐서곱을 취하면, 모든 연접 $\mathscr{O}_X$-가군 $\mathscr{F}$에 대하여
표준 준동형 $\mathscr{O}_X\to\mathscr{O}_X/\mathscr{J}$들로부터
$\mathscr{O}_X$-가군 준동형
$\mathscr{F}\to
\mathscr{F}\otimes_{\mathscr{O}_X}(\mathscr{O}_X/\mathscr{J})$들을 얻으며,
이들도 사영계를 이룬다. 따라서 사영극한을 취하면 함자적인 표준
$\mathscr{O}_X$-가군 준동형
$\gamma:\mathscr{F}\to\psi_*(\mathscr{F}_{/X'})$를 얻는다.
\end{env}

\begin{proposition}[10.8.8]
\label{I.10.8.8-ko}
\medskip\noindent
\begin{enumerate}
  \item[\rm{(i)}] 함자 $\mathscr{F}_{/X'}$는
    ($\mathscr{F}$에 관하여) 완전하다.
  \item[\rm{(ii)}] 함자적인 준동형
    $\gamma^\sharp:i^*(\mathscr{F})\to\mathscr{F}_{/X'}$는
    $(\mathscr{O}_X)_{/X'}$-가군의 동형이다.
\end{enumerate}
\end{proposition}

\medskip\noindent
\begin{enumerate}
  \item[\rm{(i)}] 다음을 증명하면 충분하다.
    $0\to\mathscr{F}'\to\mathscr{F}\to\mathscr{F}''\to0$이 연접
    $\mathscr{O}_X$-가군의 완전열이고, $U$가 $X$의 아핀 열린집합이며
    그 환이 뇌터 환 $A$이면, 다음 열은 완전하다.
    \[
      0\to\Gamma(U\cap X',\mathscr{F}'_{/X'})
      \to\Gamma(U\cap X',\mathscr{F}_{/X'})
      \to\Gamma(U\cap X',\mathscr{F}''_{/X'})\to0
    \]
    이때 $\mathscr{F}|U=\widetilde M$,
    $\mathscr{F}'|U=\widetilde{M'}$, $\mathscr{F}''|U=\widetilde{M''}$이고,
    여기서 $M$, $M'$, $M''$는 세 유한 생성 $A$-가군이며 열
    $0\to M'\to M\to M''\to0$은 완전하다
    (\hyperref[I.1.5.1-ko]{1.5.1}과
    \hyperref[I.1.3.11-ko]{1.3.11})~; $\mathscr{J}\in\Phi$라 하고,
    $\mathfrak{J}$를 $A$의 아이디얼로서
    $\mathscr{J}|U=\widetilde{\mathfrak{J}}$인 것으로 택하자. 그러면
    \[
      \Gamma(U\cap X',
      \mathscr{F}\otimes_{\mathscr{O}_X}(\mathscr{O}_X/\mathscr{J}^n))
      =M\otimes_A(A/\mathfrak{J}^n)
    \]
    이다(\hyperref[I.1.3.12-ko]{1.3.12})~; 따라서 사영극한의 정의에 의해
    \[
      \Gamma(U\cap X',\mathscr{F}_{/X'})
      =\varprojlim_n(M\otimes_A(A/\mathfrak{J}^n))=\widehat M
    \]
    이고, 이는 $M$의 $\mathfrak{J}$-준아딕 위상에 관한 분리
    완비화이다. 마찬가지로
    \[
      \Gamma(U\cap X',\mathscr{F}'_{/X'})=\widehat{M'},\qquad
      \Gamma(U\cap X',\mathscr{F}''_{/X'})=\widehat{M''}~;
    \]
    이제 $A$가 뇌터이면 함자 $\widehat M$이 $M$에 관하여 유한 생성
    $A$-가군의 범주에서 완전하다는 사실에서 우리의 주장이 따른다
    (\hyperref[0.7.3.3-ko]{0, 7.3.3}).

\oldpage[I]{197}
  \item[\rm{(ii)}] 문제는 국소적이므로 완전열
    $\mathscr{O}_X^m\to\mathscr{O}_X^n\to\mathscr{F}\to0$이
    있다고 가정할 수 있다
    (\hyperref[0.5.3.2-ko]{0, 5.3.2})~; $\gamma^\sharp$가 함자적이고,
    함자 $i^*(\mathscr{F})$와 $\mathscr{F}_{/X'}$가 각각
    (\hyperref[0.4.3.1-ko]{0, 4.3.1})과 (i)에 의해 오른쪽 완전하므로
    다음 가환도를 얻는다.
    \[
      \xymatrix{
        i^*(\mathscr{O}_X^m)\ar[r]\ar[d]_{\gamma^\sharp} &
        i^*(\mathscr{O}_X^n)\ar[r]\ar[d]_{\gamma^\sharp} &
        i^*(\mathscr{F})\ar[r]\ar[d]_{\gamma^\sharp} &
        0\\
        (\mathscr{O}_X^m)_{/X'}\ar[r] &
        (\mathscr{O}_X^n)_{/X'}\ar[r] &
        \mathscr{F}_{/X'}\ar[r] &
        0
      }
      \tag{10.8.8.1}\label{I.10.8.8.1-ko}
    \]
    이 가환도의 두 행은 완전하다. 또한 두 함자
    $i^*(\mathscr{F})$와 $\mathscr{F}_{/X'}$는 유한 직합과 가환하므로
    (\hyperref[0.3.2.6-ko]{0, 3.2.6}과
    \hyperref[0.4.3.2-ko]{4.3.2}), 우리의 주장을
    $\mathscr{F}=\mathscr{O}_X$인 경우에 증명하는 것으로 환원된다.
    이때
    $i^*(\mathscr{O}_X)=(\mathscr{O}_X)_{/X'}=
    \mathscr{O}_{\widehat X}$
    이고(\hyperref[0.4.3.4-ko]{0, 4.3.4}), $\gamma^\sharp$는
    $\mathscr{O}_{\widehat X}$-가군 준동형이다. 따라서
    $\gamma^\sharp$가 $\mathscr{O}_{\widehat X}$의 단위원 단면을
    $X'$의 한 열린집합 위에서 자기 자신으로 보내는지 확인하면 충분하다.
    이는 명백하며, 따라서 이 경우 $\gamma^\sharp$는 항등준동형이다.
\end{enumerate}

\begin{corollary}[10.8.9]
\label{I.10.8.9-ko}
환 달린 공간의 사상 $i:X_{/X'}\to X$는 평탄하다.
\end{corollary}

이는 실제로 (\hyperref[0.6.7.3-ko]{0, 6.7.3})과
(\hyperref[I.10.8.8-ko]{10.8.8, (i)})에서 따른다.

\begin{corollary}[10.8.10]
\label{I.10.8.10-ko}
$\mathscr{F}$와 $\mathscr{G}$가 연접 $\mathscr{O}_X$-가군이면, 다음
표준 동형사상들이 존재하며 $\mathscr{F}$와 $\mathscr{G}$에 관하여 함자적이다.
\[
  (\mathscr{F}_{/X'})\otimes_{(\mathscr{O}_X)_{/X'}}
  (\mathscr{G}_{/X'})
  \xrightarrow{\sim}
  (\mathscr{F}\otimes_{\mathscr{O}_X}\mathscr{G})_{/X'}
  \tag{10.8.10.1}\label{I.10.8.10.1-ko}
\]
\[
  (\mathscr{H}\!om_{\mathscr{O}_X}(\mathscr{F},\mathscr{G}))_{/X'}
  \xrightarrow{\sim}
  \mathscr{H}\!om_{(\mathscr{O}_X)_{/X'}}
  (\mathscr{F}_{/X'},\mathscr{G}_{/X'})
  \tag{10.8.10.2}\label{I.10.8.10.2-ko}
\]
\end{corollary}

이는 $i^*(\mathscr{F})$와 $\mathscr{F}_{/X'}$의 표준 동일시에서 따른다~;
이때 첫 번째 동형의 존재는 모든 환 달린 공간의 사상에 대하여 성립하는
결과이고 (\hyperref[0.4.3.3.1-ko]{0, 4.3.3.1}), 두 번째 동형의 존재는
모든 평탄 사상에 대하여 성립하는 결과이므로
(\hyperref[0.6.7.6-ko]{0, 6.7.6}),
(\hyperref[I.10.8.9-ko]{10.8.9})에서 따른다.

\begin{proposition}[10.8.11]
\label{I.10.8.11-ko}
모든 연접 $\mathscr{O}_X$-가군 $\mathscr{F}$에 대하여, 표준 준동형
$\Gamma(X,\mathscr{F})\to\Gamma(X',\mathscr{F}_{/X'})$는
$\mathscr{F}\to\mathscr{F}_{/X'}$에서 유도되며, 그 핵은 $X'$의 어떤
근방에서 영인 단면들로 이루어진다.
\end{proposition}

그러한 단면의 표준 상이 영임은 $\mathscr{F}_{/X'}$의 정의에서 따른다.
역으로 $s\in\Gamma(X,\mathscr{F})$의 상이
$\Gamma(X',\mathscr{F}_{/X'})$에서 영이라고 하자. 모든 $x\in X'$가
$X$ 안에 하나의 근방을 가져 그곳에서 $s$가 영임을 보이면 충분하고, 따라서
$X=\operatorname{Spec}(A)$가 아핀 스킴이고 $A$가 뇌터 환이며,
$X'=V(\mathfrak{J})$이고 $\mathfrak{J}$가 $A$의 아이디얼이며,
$\mathscr{F}=\widetilde M$이고 $M$이 유한 생성 $A$-가군인 경우로
환원할 수 있다. 이때 $\Gamma(X',\mathscr{F}_{/X'})$는
분리 완비화 $\widehat M$이고, 이는 $M$을 $\mathfrak{J}$-준아딕 위상에
관하여 분리 완비화한 것이다. 또한
준동형 $\Gamma(X,\mathscr{F})\to\Gamma(X',\mathscr{F}_{/X'})$는
표준 준동형 $M\to\widehat M$이다. 이 준동형의 핵은
$z\in M$ 가운데 $1+\mathfrak{J}$의 한 원소에 의해 소멸되는 것들의 집합임이
알려져 있다 (\hyperref[0.7.3.7-ko]{0, 7.3.7}). 따라서
$(1+f)s=0$인 어떤 $f\in\mathfrak{J}$가 존재한다~; 모든 $x\in X'$에
대하여 이로부터 $(1_x+f_x)s_x=0$을 얻고,
$1_x+f_x$는 $\mathscr{O}_x$에서 가역인데(실제로
$\mathfrak{J}_x\mathscr{O}_x$는 $\mathscr{O}_x$의 극대 아이디얼에
포함된다). 따라서
$s_x=0$이고, 이는 명제를 증명한다.

\begin{corollary}[10.8.12]
\label{I.10.8.12-ko}
$\mathscr{F}_{/X'}$의 지지집합은
$\operatorname{Supp}(\mathscr{F})\cap X'$과 같다.
\end{corollary}

$\mathscr{F}_{/X'}$가 유한 생성 $(\mathscr{O}_X)_{/X'}$-가군임은
(\hyperref[I.10.8.8-ko]{10.8.8, (ii)})와
(\hyperref[0.5.2.4-ko]{0, 5.2.4})에서 명백하므로,

\oldpage[I]{198}

그 지지집합은 닫혀 있고
(\hyperref[0.5.2.2-ko]{0, 5.2.2}), 명백히
$\operatorname{Supp}(\mathscr{F})\cap X'$에 포함된다. 이 마지막
집합과 같음을 보이기 위해서는 등식
$\Gamma(X',\mathscr{F}_{/X'})=0$이
$\operatorname{Supp}(\mathscr{F})\cap X'=\varnothing$을 함의함을
증명하는 것으로 즉시 환원된다~; 그런데 이는
(\hyperref[I.10.8.11-ko]{10.8.11})과
(\hyperref[I.1.4.1-ko]{1.4.1})에서 따른다.

\begin{corollary}[10.8.13]
\label{I.10.8.13-ko}
$u:\mathscr{F}\to\mathscr{G}$를 연접 $\mathscr{O}_X$-가군 사이의 준동형이라
하자. $u_{/X'}:\mathscr{F}_{/X'}\to\mathscr{G}_{/X'}$가 영이기 위한
필요충분조건은 $u$가 $X'$의 어떤 근방에서 영 준동형인 것이다.
\end{corollary}

실제로 (\hyperref[I.10.8.8-ko]{10.8.8, (ii)})에 따라 $u_{/X'}$는
$i^*(u)$와 동일시된다. 따라서 $u$를 $X$ 위에서 층
$\mathscr{H}=\mathscr{H}\!om_{\mathscr{O}_X}(\mathscr{F},\mathscr{G})$의
단면으로 보면, $u_{/X'}$는
$i^*(\mathscr{H})=\mathscr{H}_{/X'}$의 $X'$ 위의 단면이며 앞서 본
단면에 표준적으로 대응한다
((\hyperref[I.10.8.10.2-ko]{10.8.10.2})와
(\hyperref[0.4.4.6-ko]{0, 4.4.6})). 그러므로 연접
$\mathscr{O}_X$-가군 $\mathscr{H}$에
(\hyperref[I.10.8.11-ko]{10.8.11})을 적용하면 충분하다.

\begin{corollary}[10.8.14]
\label{I.10.8.14-ko}
$u:\mathscr{F}\to\mathscr{G}$를 연접 $\mathscr{O}_X$-가군 사이의
준동형이라 하자. $u_{/X'}$가 단사(각각 전사)이기 위한 필요충분조건은
$u$가 $X'$의 어떤 근방에서 단사(각각 전사)인 것이다.
\end{corollary}

$\mathscr{P}$와 $\mathscr{N}$을 각각 $u$의 여핵과 핵이라 하자.
그러면 다음 완전열이 성립한다.
\[
  0\to\mathscr{N}\xrightarrow{v}\mathscr{F}\xrightarrow{u}
  \mathscr{G}\xrightarrow{w}\mathscr{P}\to0,
\]
따라서 (\hyperref[I.10.8.8-ko]{10.8.8, (i)})에 의해 다음 완전열이
성립한다.
\[
  0\to\mathscr{N}_{/X'}\xrightarrow{v_{/X'}}
  \mathscr{F}_{/X'}\xrightarrow{u_{/X'}}
  \mathscr{G}_{/X'}\xrightarrow{w_{/X'}}
  \mathscr{P}_{/X'}\to0.
\]
$u_{/X'}$가 단사(각각 전사)이면 $v_{/X'}=0$ (각각
$w_{/X'}=0$)이다. 따라서 (\hyperref[I.10.8.13-ko]{10.8.13})에 의해
$X'$의 어떤 근방에서는 $v=0$ (각각 $w=0$)이다.

\subsection{완비화로의 사상의 연장.}
\label{subsection:I.10.9-ko}

\begin{env}[10.9.1]
\label{I.10.9.1-ko}
$X$, $Y$를 (보통 의미의) 국소 뇌터 준스킴이라 하고, $f:X\to Y$를
사상이라 하며, $X'$ (각각 $Y'$)를 $X$ (각각 $Y$)의 바탕공간의 닫힌
부분집합이라 하자. $f(X')\subset Y'$라고 가정한다. $\mathscr{J}$
(각각 $\mathscr{K}$)를 $\mathscr{O}_X$ (각각 $\mathscr{O}_Y$)의
아이디얼층이라 하고, $\mathscr{O}_X/\mathscr{J}$ (각각
$\mathscr{O}_Y/\mathscr{K}$)의 지지집합은 $X'$ (각각 $Y'$)이며
$f^*(\mathscr{K})\mathscr{O}_X\subset\mathscr{J}$라고 하자. 이러한
아이디얼층은 항상 존재한다는 점을 지적해 두자. 예를 들어 $\mathscr{J}$로
$\mathscr{O}_X$의 아이디얼층 가운데 $X$의 부분준스킴을 정의하고
$X'$을 바탕공간으로 갖는 가장 큰 것을 택할 수 있고
(\hyperref[I.5.2.1-ko]{5.2.1}),
$f(X')\subset Y'$라는 가정에서
$f^*(\mathscr{K})\mathscr{O}_X\subset\mathscr{J}$가 따른다
(\hyperref[I.5.2.4-ko]{5.2.4}). 따라서 모든 정수 $n>0$에 대하여
$f^*(\mathscr{K}^n)\mathscr{O}_X\subset\mathscr{J}^n$
(\hyperref[0.4.3.5-ko]{0, 4.3.5})이고, 따라서
(\hyperref[I.4.4.6-ko]{4.4.6})에 의해 다음과 같이 놓으면
\[
  X'_n=(X',\mathscr{O}_X/\mathscr{J}^{n+1}),\qquad
  Y'_n=(Y',\mathscr{O}_Y/\mathscr{K}^{n+1}),
\]
$f$로부터 사상 $f_n:X'_n\to Y'_n$을 얻는다. 또한 $f_n$들은
귀납계를 이룬다. 그 귀납극한
(\hyperref[I.10.6.8-ko]{10.6.8})을
$\widehat f:X_{/X'}\to Y_{/Y'}$로 나타내고, (표현을 다소 남용하여)
$\widehat f$를 $f$의, $X$와 $Y$를 각각 $X'$와 $Y'$를 따라 완비화한
것들 사이로의 연장이라 부르겠다. 이 사상이 위 조건을
만족하는 아이디얼층 $\mathscr{J}$, $\mathscr{K}$의 선택에 의존하지
않음은 즉시 확인할 수 있다. 실제로 이를 확인하기 위해서는 $X$, $Y$가
각각 환 $A$, $B$의 아핀 뇌터 준스킴인 경우만 보면 충분하다. 이때
$\mathscr{J}=\widetilde{\mathfrak{J}}$,
$\mathscr{K}=\widetilde{\mathfrak{K}}$이고, $\mathfrak{J}$ (각각
$\mathfrak{K}$)는 $A$ (각각 $B$)의 아이디얼이며, $f$는
$\varphi:B\to A$라는 환 준동형에 대응하고
$\varphi(\mathfrak{K})\subset\mathfrak{J}$이다
(\hyperref[I.4.4.6-ko]{4.4.6} 및
\hyperref[I.1.7.4-ko]{1.7.4}). 이때
\emph{$\widehat f$는 (\hyperref[I.10.2.2-ko]{10.2.2})에 따라 연속
준동형 $\widehat\varphi:\widehat B\to\widehat A$에 대응하는
사상이다}. 여기서 $\widehat A$ (각각 $\widehat B$)는 $A$ (각각
$B$)에 $\mathfrak{J}$-준아딕 (각각 $\mathfrak{K}$-준아딕) 위상을
주어 얻는 분리 완비화이다
(\hyperref[I.10.6.8-ko]{10.6.8}). 또한
$\mathscr{J}$를 다른
\oldpage[I]{199}
아이디얼층 $\mathscr{J}'=\widetilde{\mathfrak{J}'}$로 바꾸어도,
$\mathscr{O}_X/\mathscr{J}'$의 지지집합이 여전히 $X'$이면,
$\mathfrak{J}$-준아딕 위상과 $\mathfrak{J}'$-준아딕 위상은 $A$ 위에서
같음을 알고 있다
(\hyperref[I.10.8.2-ko]{10.8.2}).
이 정의에 따르면, $X_{/X'}$와 $Y_{/Y'}$의 바탕공간 사이의
연속사상 $X'\to Y'$는 $\widehat f$에 대응하며, 다름 아닌 $X'$에 대한
$f$의 제한임을 지적하자.
\end{env}

\begin{env}[10.9.2]
\label{I.10.9.2-ko}
앞의 정의로부터 환 달린 공간의 사상들로 이루어진 도식
\[
  \label{I.10.9.2.1-ko}
  \xymatrix{
    \widehat X\ar[r]^{\widehat f}\ar[d]_{i_X} &
    \widehat Y\ar[d]^{i_Y}\\
    X\ar[r]_f &
    Y
  }
\]
이 가환이고, 세로 화살표들은 표준 사상
(\hyperref[I.10.8.7-ko]{10.8.7})임이 곧바로 따른다.
\end{env}

\begin{env}[10.9.3]
\label{I.10.9.3-ko}
$Z$를 세 번째 준스킴, $g:Y\to Z$를 사상, $Z'$을
$g(Y')\subset Z'$인 $Z$의 닫힌 부분집합이라 하자. $Y'$과 $Z'$을
따른 사상 $g$의 완비화를 $\widehat g$로 나타내면,
(\hyperref[I.10.9.1-ko]{10.9.1})로부터
$\widehat{(g\circ f)}=\widehat g\circ\widehat f$임이 곧바로 따른다.
\end{env}

\begin{proposition}[10.9.4]
\label{I.10.9.4-ko}
$X$, $Y$를 두 국소 뇌터 $S$-준스킴이라 하고, $Y$가 $S$ 위의
유한형이라고 하자. $f$, $g$를 $X$에서 $Y$로 가는 두 $S$-사상이라
하고, $f(X')\subset Y'$, $g(X')\subset Y'$라고 하자.
$\widehat f=\widehat g$일 필요충분조건은 $f$와 $g$가 $X'$의 어떤
근방에서 일치하는 것이다.
\end{proposition}

$Y$에 대한 유한성 가정 없이도 이 조건은 자명하게 충분하다.
필요성을 보기 위해, 먼저 가정 $\widehat f=\widehat g$에서 모든
$x\in X'$에 대하여 $f(x)=g(x)$가 따름을 지적하자. 한편 문제는
국소적이므로, $X$와 $Y$가 각각 $x$와 $y=f(x)=g(x)$의 아핀 열린
근방이고 그 좌표환들이 뇌터이며, $S$가 아핀이고
$\Gamma(Y,\mathscr{O}_Y)$가 $\Gamma(S,\mathscr{O}_S)$ 위의 유한 생성
대수라고 가정할 수 있다
(\hyperref[I.6.3.3-ko]{6.3.3}). 이때 $f$와 $g$는
$\Gamma(Y,\mathscr{O}_Y)$에서 $\Gamma(X,\mathscr{O}_X)$로 가는 두
$\Gamma(S,\mathscr{O}_S)$-대수 준동형 $\rho$, $\sigma$에 대응한다
(\hyperref[I.1.7.3-ko]{1.7.3}). 또한 가정에 따라 이 준동형들을
$\Gamma(Y,\mathscr{O}_Y)$의 분리 완비화에 연속적으로 연장한
준동형들은 서로 같다. 따라서 (\hyperref[I.10.8.11-ko]{10.8.11})에서
모든 단면 $s\in\Gamma(Y,\mathscr{O}_Y)$에 대하여 단면 $\rho(s)$와
$\sigma(s)$가 $X'$의 어떤 근방에서 일치한다는 결론을 얻는다(이
근방은 $s$에 의존한다)~; $\Gamma(Y,\mathscr{O}_Y)$가
$\Gamma(S,\mathscr{O}_S)$ 위의 유한 생성 대수이므로, $X'$의 한 근방
$V$가 존재하여 \emph{모든} 단면 $s\in\Gamma(Y,\mathscr{O}_Y)$에
대하여 $\rho(s)$와 $\sigma(s)$가 $V$에서 일치함이 곧바로 따른다.
$D(h)$가 $V$에 포함되는 $x$의 근방이 되도록
$h\in\Gamma(X,\mathscr{O}_X)$를 택하면, 앞의 결론과
(\hyperref[I.1.4.1-ko]{1.4.1, d)})에서 $f$와 $g$가 $D(h)$에서
일치함을 얻는다.\footnote{\textbf{번역 주.} 인쇄 원문은 여기서 형에
맞지 않는 $h\in\Gamma(Y,\mathscr{O}_X)$라고 쓴다. 그러나
$\mathscr{O}_X$는 $X$ 위의 층이고 이어지는 $D(h)$는 $X$ 안에서
$x$의 근방이므로, $h\in\Gamma(X,\mathscr{O}_X)$로 바로잡았다.}

\begin{proposition}[10.9.5]
\label{I.10.9.5-ko}
(\hyperref[I.10.9.1-ko]{10.9.1})의 가정 아래에서, 모든 연접
$\mathscr{O}_Y$-가군 $\mathscr{G}$에 대하여 다음과 같은
$(\mathscr{O}_X)_{/X'}$-가군의 함자적인 표준 동형이 존재한다.
\[
  (f^*(\mathscr{G}))_{/X'}\xrightarrow{\sim}
  \widehat f^*(\mathscr{G}_{/Y'})
\]
\end{proposition}

$(f^*(\mathscr{G}))_{/X'}$를 $i_X^*(f^*(\mathscr{G}))$와,
$\widehat f^*(\mathscr{G}_{/Y'})$를
$\widehat f^*(i_Y^*(\mathscr{G}))$와 표준적으로 동일시하면
(\hyperref[I.10.8.8-ko]{10.8.8}), 명제는
(\hyperref[I.10.9.2-ko]{10.9.2})의 도식이 가환이라는 사실에서 곧바로
따른다.

\begin{env}[10.9.6]
\label{I.10.9.6-ko}
이제 $\mathscr{F}$를 연접 $\mathscr{O}_X$-가군,
$\mathscr{G}$를 연접 $\mathscr{O}_Y$-가군이라 하자.
$u:\mathscr{G}\to\mathscr{F}$가 $\mathscr{G}$에서 $\mathscr{F}$로
가는 $f$-사상이면, 이에 대응하는 $\mathscr{O}_X$-준동형
$u^\sharp:f^*(\mathscr{G})\to\mathscr{F}$가 있다. 따라서 이 준동형을
완비화하면 연속 $(\mathscr{O}_X)_{/X'}$-준동형
$(u^\sharp)_{/X'}:(f^*(\mathscr{G}))_{/X'}\to\mathscr{F}_{/X'}$를
얻고, (\hyperref[I.10.9.5-ko]{10.9.5})에 의해
$\widehat f$-사상 $v:\mathscr{G}_{/Y'}\to\mathscr{F}_{/X'}$가 존재하고,
\oldpage[I]{200}

유일하며, $v^\sharp=(u^\sharp)_{/X'}$를 만족한다. 삼중항
$(\mathscr{F},X,X')$들($\mathscr{F}$는 연접 $\mathscr{O}_X$-가군이고
$X'$은 $X$의 닫힌 부분집합이다)을 하나의 \emph{범주}의 대상으로
생각하고, 사상
$(\mathscr{F},X,X')\to(\mathscr{G},Y,Y')$들이
$f(X')\subset Y'$인 준스킴의 사상 $f:X\to Y$와
$f$-사상 $u:\mathscr{G}\to\mathscr{F}$로 이루어진다고 하자. 그러면
$(X_{/X'},\mathscr{F}_{/X'})$는 $(\mathscr{F},X,X')$에 관한
\emph{함자}라고 말할 수 있다. 이 함자는 국소 뇌터 형식적 준스킴
$\mathfrak{Z}$와 $\mathscr{O}_{\mathfrak{Z}}$-가군 $\mathscr{H}$로
이루어진 쌍 $(\mathfrak{Z},\mathscr{H})$들의 범주에 값을 갖는다.
후자의 범주에서 사상은 형식적 준스킴의 사상 $g$와 $g$-사상으로
이루어진 쌍이다.
\end{env}

\begin{proposition}[10.9.7]
\label{I.10.9.7-ko}
$S$, $X$, $Y$를 세 국소 뇌터 준스킴,
$g:X\to S$, $h:Y\to S$를 두 사상, $S'$을 $S$의 닫힌 부분집합,
$X'$ (각각 $Y'$)을 $g(X')\subset S'$ (각각
$h(Y')\subset S'$)을 만족하는 $X$ (각각 $Y$)의 닫힌 부분집합이라
하자~; $Z=X\times_S Y$라 하고, $Z$가 국소 뇌터라고 가정하며,
$Z'=p^{-1}(X')\cap q^{-1}(Y')$라 하자. 여기서 $p$와 $q$는
$X\times_S Y$의 사영이다. 이 조건에서 완비화 $Z_{/Z'}$는
$S_{/S'}$-형식적 준스킴들의 곱
$(X_{/X'})\times_{S_{/S'}}(Y_{/Y'})$와 동일시된다. 구조 사상들은
각각 $\widehat g$, $\widehat h$와 동일시되고, 사영들은 각각
$\widehat p$, $\widehat q$와 동일시된다.
\end{proposition}

문제가 $S$, $X$, $Y$에 관하여 국소적임은 명백하므로,
$S=\operatorname{Spec}(A)$,
$X=\operatorname{Spec}(B)$, $Y=\operatorname{Spec}(C)$,
$S'=V(\mathfrak{J})$, $X'=V(\mathfrak{K})$,
$Y'=V(\mathfrak{L})$인 경우로 환원된다. 여기서
$\mathfrak{J}$, $\mathfrak{K}$, $\mathfrak{L}$은
$\varphi(\mathfrak{J})\subset\mathfrak{K}$ 및
$\psi(\mathfrak{J})\subset\mathfrak{L}$을 만족하는 세 아이디얼이고,
$\varphi:A\to B$와 $\psi:A\to C$는 각각 $g$와 $h$에 대응하는
준동형이다. 이때 $Z=\operatorname{Spec}(B\otimes_A C)$이고
$Z'=V(\mathfrak{M})$임을 알고 있다. 여기서 $\mathfrak{M}$은
아이디얼
$\operatorname{Im}(\mathfrak{K}\otimes_A C)
+\operatorname{Im}(B\otimes_A\mathfrak{L})$이다.
(\hyperref[I.10.7.2-ko]{10.7.2})에 의해, 결론은 완비 텐서곱
$(\widehat B\otimes_{\widehat A}\widehat C)^\wedge$
(여기서 $\widehat A$, $\widehat B$, $\widehat C$는 각각
$A$, $B$, $C$에 $\mathfrak{J}$-, $\mathfrak{K}$-,
$\mathfrak{L}$-준아딕 위상을 주어 얻는 분리 완비화이다)이
텐서곱 $B\otimes_A C$의 $\mathfrak{M}$-준아딕 위상에 관한
분리 완비화라는 사실에서 따른다
(\hyperref[0.7.7.2-ko]{0, 7.7.2}).

또한 $T$가 국소 뇌터 $S$-준스킴이고,
$u:T\to X$, $v:T\to Y$가 두 $S$-사상이며, $T'$이
$u(T')\subset X'$, $v(T')\subset Y'$을 만족하는 $T$의 닫힌
부분집합이면, 완비화한 것들 사이로의 연장
$((u,v)_S)^\wedge$는
$(\widehat u,\widehat v)_{S_{/S'}}$와 동일시됨을 유의하자.

\begin{corollary}[10.9.8]
\label{I.10.9.8-ko}
$X$, $Y$를 두 국소 뇌터 $S$-준스킴이라 하고,
$X\times_S Y$도 국소 뇌터라고 하자~; $S'$을 $S$의 닫힌
부분집합이라 하고, $X'$ (각각 $Y'$)을 $X$ (각각 $Y$)의 닫힌
부분집합이라 하되 그 구조 사상에 의한 상이 $S'$에 포함된다고 하자.
$f(X')\subset Y'$을 만족하는 모든 $S$-사상 $f:X\to Y$에 대하여,
그래프 사상 $\Gamma_{\widehat f}$는 $f$의 그래프 사상의 연장
$(\Gamma_f)^\wedge$와 동일시된다.
\end{corollary}

\begin{corollary}[10.9.9]
\label{I.10.9.9-ko}
$X$, $Y$를 두 국소 뇌터 준스킴, $f:X\to Y$를 사상,
$Y'$을 $Y$의 닫힌 부분집합이라 하고, $X'=f^{-1}(Y')$라 하자.
그러면 형식적 준스킴\footnote{\textbf{번역 주.} 인쇄 원문은 여기서
\emph{préschéma}라고 쓰지만, $X_{/X'}$는 정의상 형식적 준스킴이고
우변도 형식적 준스킴들의 곱이므로 \emph{préschéma formel}에서
\emph{formel}이 빠진 것으로 보아 바로잡았다.} $X_{/X'}$는 다음 가환도식
\[
  \label{I.10.9.9.1-ko}
  \xymatrix{
    X\ar[d]_f &
    X_{/X'}\ar[l]\ar[d]^{\widehat f}\\
    Y &
    Y_{/Y'}\ar[l]
  }
\]
에 의해 형식적 준스킴들의 곱
$X\times_Y(Y_{/Y'})$와 동일시된다.
\end{corollary}

$S$와 $S'$을 모두 $Y$로, $X$와 $X'$을 모두 $X$로 바꾸어
(\hyperref[I.10.9.7-ko]{10.9.7})을 적용하면 충분하다.
\oldpage[I]{201}

\begin{remark}[10.9.10]
\label{I.10.9.10-ko}
$X$가 합 $X_1\amalg X_2$
(\hyperref[subsection:I.3.1-ko]{3.1})이고, $X'$이 합집합
$X'_1\cup X'_2$이며, 여기서 $X'_i$가 $X_i$의 닫힌 부분집합이면
($i=1,2$), 곧바로
$X_{/X'}={X_1}_{/X'_1}\amalg {X_2}_{/X'_2}$임을 알 수 있다.
\end{remark}

\subsection{아핀 형식적 스킴 위의 연접층에 대한 응용.}
\label{subsection:I.10.10-ko}

\begin{env}[10.10.1]
\label{I.10.10.1-ko}
이 절 전체에서 $A$는 \emph{뇌터 아딕환}, $\mathfrak{J}$는 $A$의
정의 아이디얼을 나타낸다. $X=\operatorname{Spec}(A)$,
$\mathfrak{X}=\operatorname{Spf}(A)$라 하자. 후자의 바탕공간은
$V(\mathfrak{J})$와 동일시되며, 이는 $X$의 닫힌 부분집합이다
(\hyperref[I.10.1.2-ko]{10.1.2}). 또한 정의
(\hyperref[I.10.1.2-ko]{10.1.2})와 정의
(\hyperref[I.10.8.4-ko]{10.8.4})에서 \emph{아핀 형식적 스킴
$\mathfrak{X}$}는 \emph{완비화 $X_{/\mathfrak{X}}$}, 곧 아핀 스킴
$X$를 그 바탕공간의 닫힌 부분집합 $\mathfrak{X}$을 따라 완비화한
것과 동일함을 알 수 있다. 따라서 각 연접 $\mathscr{O}_X$-가군 $\mathscr{F}$에는
유한 생성 $\mathscr{O}_{\mathfrak{X}}$-가군
$\mathscr{F}_{/\mathfrak{X}}$가 대응하며, 이는 또한 위상환의 층
$\mathscr{O}_{\mathfrak{X}}$ 위의 위상 가군의 층이다. 한편 모든 연접
$\mathscr{O}_X$-가군 $\mathscr{F}$는 $\widetilde M$ 꼴이다. 여기서
$M$은 유한 생성 $A$-가군이다
(\hyperref[I.1.5.1-ko]{1.5.1})~; 우리는
$(\widetilde M)_{/\mathfrak{X}}=M^\Delta$로 두겠다. 또한
$u:M\to N$이 유한 생성 $A$-가군 사이의 $A$-준동형이면, 이에
준동형 $\widetilde u:\widetilde M\to\widetilde N$이 대응하고, 따라서
연속 준동형
$\widetilde u_{/\mathfrak{X}}:(\widetilde M)_{/\mathfrak{X}}\to
(\widetilde N)_{/\mathfrak{X}}$도 대응한다. 이를 $u^\Delta$로
나타내겠다. 곧바로
$(v\circ u)^\Delta=v^\Delta\circ u^\Delta$임을 알 수 있다~; 이로써
\emph{$M^\Delta$를 주는 가법적 공변 함자}를 정의하였다. 이 함자는
유한 생성 $A$-가군의 범주에서 유한 생성
$\mathscr{O}_{\mathfrak{X}}$-가군의 범주로 간다. $A$가 \emph{이산}환이면
$M^\Delta=\widetilde M$이다.
\end{env}

\begin{proposition}[10.10.2]
\label{I.10.10.2-ko}
\medskip\noindent
\begin{enumerate}
  \item[\rm{(i)}] $M^\Delta$는 $M$에 관한 완전 함자이고, $A$-가군의
    함자적인 표준 동형
    $\Gamma(\mathfrak{X},M^\Delta)\xrightarrow{\sim}M$이 존재한다.
  \item[\rm{(ii)}] $M$과 $N$이 두 유한 생성 $A$-가군이면, 다음
    함자적인 표준 동형들이 존재한다.
    \[
      (M\otimes_A N)^\Delta
      \xrightarrow{\sim}
      M^\Delta\otimes_{\mathscr{O}_{\mathfrak{X}}}N^\Delta
      \tag{10.10.2.1}\label{I.10.10.2.1-ko}
    \]
    \[
      (\operatorname{Hom}_A(M,N))^\Delta
      \xrightarrow{\sim}
      \mathscr{H}\!om_{\mathscr{O}_{\mathfrak{X}}}
      (M^\Delta,N^\Delta)
      \tag{10.10.2.2}\label{I.10.10.2.2-ko}
    \]
  \item[\rm{(iii)}] 사상 $u\to u^\Delta$는 다음 함자적인 동형이다.
    \[
      \operatorname{Hom}_A(M,N)
      \xrightarrow{\sim}
      \operatorname{Hom}_{\mathscr{O}_{\mathfrak{X}}}
      (M^\Delta,N^\Delta)
      \tag{10.10.2.3}\label{I.10.10.2.3-ko}
    \]
\end{enumerate}
\end{proposition}

$M^\Delta$를 주는 함자의 완전성은 $\widetilde M$을 주는 함자
(\hyperref[I.1.3.5-ko]{1.3.5})와 $\mathscr{F}_{/\mathfrak{X}}$
(\hyperref[I.10.8.8-ko]{10.8.8})의 완전성에서 따른다. 정의에 따라
$\Gamma(\mathfrak{X},M^\Delta)$는 $A$-가군
$\Gamma(X,\widetilde M)=M$의 $\mathfrak{J}$-준아딕 위상에 관한 분리
완비화이다~; 그러나 $A$가 완비이고 $M$이 유한 생성인 까닭에,
$M$이 분리이고 완비임을 알고 있다
(\hyperref[0.7.3.6-ko]{0, 7.3.6}). 이로써 (i)의 증명이 끝난다. 동형
(\hyperref[I.10.10.2.1-ko]{10.10.2.1}) (각각
\hyperref[I.10.10.2.2-ko]{10.10.2.2})는 동형
(\hyperref[I.1.3.12-ko]{1.3.12, (i)})와
(\hyperref[I.10.8.10.1-ko]{10.8.10.1}) (각각
(\hyperref[I.1.3.12-ko]{1.3.12, (ii)})와
(\hyperref[I.10.8.10.2-ko]{10.8.10.2}))의 합성에서 얻어진다. 끝으로
$\operatorname{Hom}_A(M,N)$은 유한 생성 $A$-가군이므로 여기에 (i)를
적용할 수 있다. 그러면
$\Gamma(\mathfrak{X},(\operatorname{Hom}_A(M,N))^\Delta)$가
$\operatorname{Hom}_A(M,N)$과 동일시되고,
(\hyperref[I.10.10.2.2-ko]{10.10.2.2})를 이용하면 준동형
(\hyperref[I.10.10.2.3-ko]{10.10.2.3})이 동형임을 알 수 있다.

(\hyperref[I.10.10.2-ko]{10.10.2})에서
(\hyperref[I.1.3.7-ko]{1.3.7})과
(\hyperref[I.1.3.12-ko]{1.3.12})에서 얻은 것들과 유사한 일련의
귀결들을 얻으며, 그 정식화는 독자에게 맡긴다.
\oldpage[I]{202}

$M^\Delta$의 완전성을 완전열
$0\to\mathfrak{J}\to A\to A/\mathfrak{J}\to0$에 적용하면, 여기서
$\mathscr{O}_{\mathfrak{X}}$의 아이디얼층 중
$\mathfrak{J}^\Delta$로 나타낸 것이
(\hyperref[I.10.3.1-ko]{10.3.1})에서 같은 기호로 나타낸
아이디얼층과 일치함을 알 수 있다는 점에 유의하자. 이는
(\hyperref[I.10.3.2-ko]{10.3.2})에 따른다.

\begin{proposition}[10.10.3]
\label{I.10.10.3-ko}
(\hyperref[I.10.10.1-ko]{10.10.1})의 가정 아래에서
$\mathscr{O}_{\mathfrak{X}}$는 연접인 환의 층이다.
\end{proposition}

$f\in A$이면 $A_{\{f\}}$가 뇌터 아딕환임을 알고 있다
(\hyperref[0.7.6.11-ko]{0, 7.6.11}). 문제는 국소적이므로
(\hyperref[I.10.1.4-ko]{10.1.4})에 따라 준동형
$v:\mathscr{O}_{\mathfrak{X}}^n\to\mathscr{O}_{\mathfrak{X}}$의 핵이
유한 생성 $\mathscr{O}_{\mathfrak{X}}$-가군임을 보이면 충분하다.
이때 $v=u^\Delta$이다. 여기서 $u$는 $A$-준동형 $A^n\to A$이다
(\hyperref[I.10.10.2-ko]{10.10.2}). $A$가 뇌터이므로 $u$의 핵은
유한 생성이다. 다시 말해, 어떤 준동형
$A^m\xrightarrow{w}A^n$이 있어서 열
$A^m\xrightarrow{w}A^n\xrightarrow{u}A$가 완전하다. 따라서
(\hyperref[I.10.10.2-ko]{10.10.2})에 의해 열
$\mathscr{O}_{\mathfrak{X}}^m\xrightarrow{w^\Delta}
\mathscr{O}_{\mathfrak{X}}^n\xrightarrow{v}
\mathscr{O}_{\mathfrak{X}}$도 완전하다. 이로써 $v$의 핵이 유한
생성임을 보였다.

\begin{env}[10.10.4]
\label{I.10.10.4-ko}
앞의 표기 아래에서 $A_n=A/\mathfrak{J}^{n+1}$로 두고, $X_n$을 아핀
스킴
$\operatorname{Spec}(A_n)=(\mathfrak{X},
\mathscr{O}_{\mathfrak{X}}/\mathscr{J}^{n+1})$이라 하자. 여기서
$\mathscr{J}=\mathfrak{J}^\Delta$는 아이디얼 $\mathfrak{J}$에 대응하는
$\mathscr{O}_{\mathfrak{X}}$의 정의 아이디얼층이다. $m\leq n$일 때
표준 준동형 $A_n\to A_m$에 대응하는 준스킴의 사상 $X_m\to X_n$을
$u_{mn}$이라 하자~; 형식적 스킴 $\mathfrak{X}$는 $u_{mn}$들에 관한
$X_n$들의 귀납극한이다
(\hyperref[I.10.6.3-ko]{10.6.3}).
\end{env}

\begin{proposition}[10.10.5]
\label{I.10.10.5-ko}
(\hyperref[I.10.10.1-ko]{10.10.1})의 가정 아래에서
$\mathscr{F}$를 $\mathscr{O}_{\mathfrak{X}}$-가군이라 하자. 다음 조건들은
서로 동치이다~:
\begin{enumerate}
  \item[a)] $\mathscr{F}$는 연접 $\mathscr{O}_{\mathfrak{X}}$-가군이다.
  \item[b)] $\mathscr{F}$는 열 $(\mathscr{F}_n)$의 사영극한
    (\hyperref[I.10.6.6-ko]{10.6.6})과 동형이다. 이 열은 연접
    $\mathscr{O}_{X_n}$-가군들로 이루어지고
    $u_{mn}^*(\mathscr{F}_n)=\mathscr{F}_m$을 만족한다.
  \item[c)] 유한 생성 $A$-가군 $M$이 존재한다. 그 가군은
    (\hyperref[I.10.10.2-ko]{10.10.2, (i)})에 의해 표준 동형을
    제외하고 결정되며, $\mathscr{F}$는 $M^\Delta$와 동형이다.
\end{enumerate}
\end{proposition}

먼저 b)가 c)를 함의함을 보이자. $\mathscr{F}_n=\widetilde{M_n}$이고,
여기서 $M_n$은 유한 생성 $A_n$-가군이다. 가정에 의해
$M_m=M_n\otimes_{A_n}A_m$이 $m\leq n$일 때 성립한다
(\hyperref[I.1.6.5-ko]{1.6.5})~; 따라서 $M_n$들은 표준 디-준동형
$M_n\to M_m$ ($m\leq n$)에 관하여 사영계를 이룬다. $A_n$들의
정의로부터 이 사영계가
(\hyperref[0.7.2.9-ko]{0, 7.2.9})의 조건들을 만족함이 곧바로 따른다~;
그러므로 그 사영극한 $M$은 유한 생성 $A$-가군이고,
$M_n=M\otimes_A A_n$이 모든 $n$에 대하여 성립한다. 따라서
$\mathscr{F}_n$은 $X_n$ 위에서
$\widetilde M\otimes_{\mathscr{O}_X}
(\mathscr{O}_X/\widetilde{\mathfrak{J}}^{n+1})$에 의해 유도되는
가군이고, 정의 (\hyperref[I.10.8.4-ko]{10.8.4})에 의해
$\mathscr{F}=M^\Delta$이다.

역으로 c)는 b)를 함의한다. 실제로 $u_n$이 몰입 사상 $X_n\to X$이면,
$u_n^*(\widetilde M)=(M\otimes_A A_n)^{\sim}$은 $X_n$ 위에서
$\widetilde M\otimes_{\mathscr{O}_X}
(\mathscr{O}_X/\widetilde{\mathfrak{J}}^{n+1})$에 의해 유도되는
가군이고, 정의 (\hyperref[I.10.8.4-ko]{10.8.4})에 의해
$M^\Delta=\varprojlim u_n^*(\widetilde M)$이다~;
$u_m=u_n\circ u_{mn}$이 $m\leq n$일 때 성립하므로,
$\mathscr{F}_n=u_n^*(\widetilde M)$들은 b)의 조건들을 만족한다. 이로써
주장이 따른다.

이제 c)가 a)를 함의함을 보이자. 정의에 의해
$\mathscr{O}_{\mathfrak{X}}=A^\Delta$이다~; $M$은 어떤 준동형
$A^m\to A^n$의 여핵이므로, (\hyperref[I.10.10.2-ko]{10.10.2})에 의해
$M^\Delta$는 준동형
$\mathscr{O}_{\mathfrak{X}}^m\to\mathscr{O}_{\mathfrak{X}}^n$의
여핵이다. $\mathscr{O}_{\mathfrak{X}}$가 연접인 환의 층이므로
(\hyperref[I.10.10.3-ko]{10.10.3}), $M^\Delta$도 연접이다
(\hyperref[0.5.3.4-ko]{0, 5.3.4}).
\oldpage[I]{203}

끝으로 a)는 b)를 함의한다. $\mathscr{O}_{\mathfrak{X}}$-가군으로
보면,
$\mathscr{O}_{X_n}=\mathscr{O}_{\mathfrak{X}}/
\mathscr{J}^{n+1}=A_n^\Delta$이다~; $\mathscr{F}_n=\mathscr{F}
\otimes_{\mathscr{O}_{\mathfrak{X}}}\mathscr{O}_{X_n}$은 연접
$\mathscr{O}_{\mathfrak{X}}$-가군이다
(\hyperref[0.5.3.5-ko]{0, 5.3.5}). 또한 이것은
$\mathscr{O}_{X_n}$-가군이고 $\mathscr{J}^{n+1}$이 연접이므로,
$\mathscr{F}_n$은 연접 $\mathscr{O}_{X_n}$-가군이다
(\hyperref[0.5.3.10-ko]{0, 5.3.10}).
$u_{mn}^*(\mathscr{F}_n)=\mathscr{F}_m$이 $m\leq n$일 때 성립함도
곧바로 알 수 있다
($X_m\to X_n$에서 얻는 바탕공간 사이의 연속사상은
$\mathfrak{X}$의 항등사상임을 기억하면 된다). 따라서 층
$\mathscr{G}=\varprojlim\mathscr{F}_n$은 연접
$\mathscr{O}_{\mathfrak{X}}$-가군이다. 이는 b)가 a)를 함의함을 이미
보였기 때문이다. 표준 준동형 $\mathscr{F}\to\mathscr{F}_n$들은
사영계를 이루며, 사영극한을 취하면 표준 준동형
$w:\mathscr{F}\to\mathscr{G}$를 얻는다. 이제 남은 것은 $w$가 전단사임을
보이는 일뿐이다. 이 문제는 국소적이므로 $\mathscr{F}$가 어떤 준동형
$\mathscr{O}_{\mathfrak{X}}^p\to\mathscr{O}_{\mathfrak{X}}^q$의
여핵인 경우로 국한할 수 있다~; 이 준동형은 $v^\Delta$ 꼴이고, 여기서
$v$는 준동형 $A^p\to A^q$이다\footnote{\textbf{번역 주.} 인쇄 원문은
앞의 층 준동형을
$\mathscr{O}_{\mathfrak{X}}^p\to\mathscr{O}_{\mathfrak{X}}^q$로
둔 뒤, 이를 유도하는 $v$를 $A^m\to A^n$으로 쓴다. 그러나
$\Delta$ 대응에서는 $(A^r)^\Delta=\mathscr{O}_{\mathfrak{X}}^r$이므로
유한 자유 가군의 계수(階數)가 보존된다. 따라서 $v:A^p\to A^q$로
바로잡았다.}
(\hyperref[I.10.10.2-ko]{10.10.2}). 그러므로 $\mathscr{F}$는
$M^\Delta$와 동형이다. 여기서 $M=\operatorname{Coker}v$이다
(\hyperref[I.10.10.2-ko]{10.10.2}). 이제
(\hyperref[I.10.10.2-ko]{10.10.2})에 의해
$\mathscr{F}_n=M^\Delta\otimes_{\mathscr{O}_{\mathfrak{X}}}A_n^\Delta
=(M\otimes_A A_n)^\Delta$이다. 또한 $\mathfrak{J}$-아딕 위상은
$M\otimes_A A_n$ 위에서 이산이므로,
$(M\otimes_A A_n)^\Delta=(M\otimes_A A_n)^{\sim}$이다
($\mathscr{O}_{X_n}$-가군으로서). 위에서
$M^\Delta=\varprojlim\mathscr{F}_n$임을 보았으므로, 이 경우 $w$는
실제로 항등사상이다.
\par\hfill C.Q.F.D.\par

\begin{corollary}[10.10.6]
\label{I.10.10.6-ko}
$\mathscr{F}$가 (\hyperref[I.10.10.5-ko]{10.10.5})의 조건 b)를
만족하면, 사영계 $(\mathscr{F}_n)$은
$\mathscr{F}\otimes_{\mathscr{O}_{\mathfrak{X}}}\mathscr{O}_{X_n}$들로
이루어진 계와 동형이다.
\end{corollary}

\begin{env}[10.10.7]
\label{I.10.10.7-ko}
이제 $A$, $B$를 두 뇌터 아딕환이라 하고,
$\varphi:B\to A$를 연속 준동형이라 하자. $\mathfrak{J}$ (각각
$\mathfrak{K}$)를 $A$ (각각 $B$)의 정의 아이디얼이라 하되
$\varphi(\mathfrak{K})\subset\mathfrak{J}$라고 하고,
$X=\operatorname{Spec}(A)$, $Y=\operatorname{Spec}(B)$,
$\mathfrak{X}=\operatorname{Spf}(A)$,
$\mathfrak{Y}=\operatorname{Spf}(B)$라 하자. $f:X\to Y$를
$\varphi$에 대응하는 준스킴의 사상이라 하고
(\hyperref[I.1.6.1-ko]{1.6.1}), 완비화한 것들 사이로의 연장을
$\widehat f:\mathfrak{X}\to\mathfrak{Y}$라 하자
(\hyperref[I.10.9.1-ko]{10.9.1}). 이는 또한 $\varphi$에 대응하는
형식적 준스킴의 사상이다 (\hyperref[I.10.2.2-ko]{10.2.2}).
\end{env}

\begin{proposition}[10.10.8]
\label{I.10.10.8-ko}
모든 유한 생성 $B$-가군 $N$에 대하여 다음과 같은
$\mathscr{O}_{\mathfrak{X}}$-가군의 함자적인 표준 동형이 존재한다~:
\[
  \widehat f^*(N^\Delta)
  \xrightarrow{\sim}
  (N\otimes_B A)^\Delta.
\]
\end{proposition}

실제로 표준 사상들을 $i_X:\mathfrak{X}\to X$와
$i_Y:\mathfrak{Y}\to Y$로 나타내면, 함자적인 표준 동형을 통해
동일시할 때 (\hyperref[I.10.8.8-ko]{10.8.8})
$N^\Delta=i_Y^*(\widetilde N)$이고
\[
  (N\otimes_B A)^\Delta
  =i_X^*((N\otimes_B A)^{\sim})
  =i_X^*(f^*(\widetilde N))
\]
이다 (\hyperref[I.1.6.5-ko]{1.6.5})~; 따라서 명제는 도표
(\hyperref[I.10.9.2-ko]{10.9.2})의 가환성에서 따른다.

\begin{corollary}[10.10.9]
\label{I.10.10.9-ko}
$\mathfrak b$를 $B$의 임의의 아이디얼이라 하면,
$\widehat f^*(\mathfrak b^\Delta)\mathscr{O}_{\mathfrak{X}}
=(\mathfrak bA)^\Delta$이다.
\end{corollary}

실제로 $j$를 표준 단사 준동형 $\mathfrak b\to B$라 하자. 이에 대응하는
$j^\Delta:\mathfrak b^\Delta\to\mathscr{O}_{\mathfrak{Y}}$는
$\mathscr{O}_{\mathfrak{Y}}$-가군의 표준 단사 준동형이다~;
정의에 따라
$\widehat f^*(\mathfrak b^\Delta)\mathscr{O}_{\mathfrak{X}}$는 준동형
$\widehat f^*(j^\Delta):\widehat f^*(\mathfrak b^\Delta)\to
\mathscr{O}_{\mathfrak{X}}=\widehat f^*(\mathscr{O}_{\mathfrak{Y}})$의
상이다~; 그런데 (\hyperref[I.10.10.8-ko]{10.10.8})에 의해 이 준동형은
$(j\otimes1)^\Delta:(\mathfrak b\otimes_B A)^\Delta\to
\mathscr{O}_{\mathfrak{X}}=(B\otimes_B A)^\Delta$와 동일시된다.
$j\otimes1$의 상은 아이디얼 $\mathfrak bA$이고 이는 $A$의 아이디얼이므로,
(\hyperref[I.10.10.2-ko]{10.10.2})에 따라 $(j\otimes1)^\Delta$의 상은
$(\mathfrak bA)^\Delta$이다. 이로써 결론이 따른다.
\oldpage[I]{204}

\subsection{형식적 준스킴 위의 연접층.}
\label{subsection:I.10.11-ko}

\begin{proposition}[10.11.1]
\label{I.10.11.1-ko}
$\mathfrak{X}$가 국소 뇌터 형식적 준스킴이면,
$\mathscr{O}_{\mathfrak{X}}$는 연접인 환의 층이고 $\mathfrak{X}$의
모든 정의 아이디얼층은 연접이다.
\end{proposition}

문제는 국소적이므로 뇌터 아핀 형식적 스킴인 경우로 환원되며, 따라서
명제는 (\hyperref[I.10.10.3-ko]{10.10.3})과
(\hyperref[I.10.10.5-ko]{10.10.5})에서 따른다.

\begin{env}[10.11.2]
\label{I.10.11.2-ko}
$\mathfrak{X}$를 국소 뇌터 형식적 준스킴,
$\mathscr{J}$를 $\mathfrak{X}$의 정의 아이디얼층이라 하자. $X_n$을
(보통 의미의) 국소 뇌터 준스킴
$(\mathfrak{X},\mathscr{O}_{\mathfrak{X}}/\mathscr{J}^{n+1})$이라 하자.
그러면 $\mathfrak{X}$는 열 $(X_n)$의 귀납극한이고, 그 전이사상들은
표준 사상 $u_{mn}:X_m\to X_n$이다
(\hyperref[I.10.6.3-ko]{10.6.3}). 이 표기 아래에서 다음이 성립한다~:
\end{env}

\begin{theorem}[10.11.3]
\label{I.10.11.3-ko}
$\mathscr{O}_{\mathfrak{X}}$-가군 $\mathscr{F}$가 연접이기 위한
필요충분조건은 그것이 열 $(\mathscr{F}_n)$의 사영극한과 동형인 것이다.
여기서 각 $\mathscr{F}_n$은 연접
$\mathscr{O}_{X_n}$-가군이며,
$u_{mn}^*(\mathscr{F}_n)=\mathscr{F}_m$이 $m\leq n$일 때 성립한다
(\hyperref[I.10.6.6-ko]{10.6.6}). 이때 사영계 $(\mathscr{F}_n)$은
$u_n^*(\mathscr{F})=\mathscr{F}
\otimes_{\mathscr{O}_{\mathfrak{X}}}\mathscr{O}_{X_n}$들로 이루어진
계와 동형이다. 여기서 $u_n$은 표준 사상 $X_n\to\mathfrak{X}$이다.
\end{theorem}

문제는 국소적이므로 $\mathfrak{X}$가 뇌터 아핀 형식적 스킴인 경우로
환원되며, 이 경우 정리는 (\hyperref[I.10.10.5-ko]{10.10.5})와
(\hyperref[I.10.10.6-ko]{10.10.6})의 귀결이다.

따라서 \emph{연접 $\mathscr{O}_{\mathfrak{X}}$-가군을 주는 것은
사영계 $(\mathscr{F}_n)$를 주는 것과 동치이다. 여기서 각 항은 연접
$\mathscr{O}_{X_n}$-가군이고,
$u_{mn}^*(\mathscr{F}_n)=\mathscr{F}_m$이 $m\leq n$일 때 성립한다}.

\begin{corollary}[10.11.4]
\label{I.10.11.4-ko}
$\mathscr{F}$와 $\mathscr{G}$가 두 연접
$\mathscr{O}_{\mathfrak{X}}$-가군이면,
(\hyperref[I.10.11.3-ko]{10.11.3})의 표기를 쓰면 다음 함자적인 표준
동형을 정의할 수 있다~:
\[
  \label{I.10.11.4.1-ko}
  \operatorname{Hom}_{\mathscr{O}_{\mathfrak{X}}}
  (\mathscr{F},\mathscr{G})
  \xrightarrow{\sim}
  \varprojlim_n
  \operatorname{Hom}_{\mathscr{O}_{X_n}}
  (\mathscr{F}_n,\mathscr{G}_n).
  \tag{10.11.4.1}
\]
\end{corollary}

우변의 사영극한은 다음 사상들에 관하여 취한 것으로 이해해야 한다~:
$\theta_n\mapsto u_{mn}^*(\theta_n)$ ($m\leq n$). 이들은
$\operatorname{Hom}_{\mathscr{O}_{X_n}}
(\mathscr{F}_n,\mathscr{G}_n)$에서
$\operatorname{Hom}_{\mathscr{O}_{X_m}}
(\mathscr{F}_m,\mathscr{G}_m)$으로 가는 사상이다. 준동형
(\hyperref[I.10.11.4.1-ko]{10.11.4.1})은 원소
$\theta\in\operatorname{Hom}_{\mathscr{O}_{\mathfrak{X}}}
(\mathscr{F},\mathscr{G})$에 열 $(u_n^*(\theta))$을 대응시킨다~;
(\hyperref[I.10.11.3-ko]{10.11.3})을 고려하면, 사영계
$(\theta_n)\in\varprojlim_n
\operatorname{Hom}_{\mathscr{O}_{X_n}}
(\mathscr{F}_n,\mathscr{G}_n)$에
$\operatorname{Hom}_{\mathscr{O}_{\mathfrak{X}}}
(\mathscr{F},\mathscr{G})$ 안에서의 그 사영극한을 대응시킴으로써
앞의 준동형의 역준동형을 정의할 수 있음을 곧 알 수 있다.

\begin{corollary}[10.11.5]
\label{I.10.11.5-ko}
준동형 $\theta:\mathscr{F}\to\mathscr{G}$가 전사일 필요충분조건은
이에 대응하는 준동형
$\theta_0=u_0^*(\theta):\mathscr{F}_0\to\mathscr{G}_0$가 전사인 것이다.
\end{corollary}

문제는 국소적이므로
$\mathfrak{X}=\operatorname{Spf}(A)$이고 $A$가 뇌터 아딕환이며,
$\mathscr{F}=M^\Delta$, $\mathscr{G}=N^\Delta$, $\theta=u^\Delta$인
경우로 환원된다. 여기서 $M$과 $N$은 유한 생성 $A$-가군이고
$u$는 준동형 $M\to N$이다~; 또한 이때
$\theta_0=\widetilde{u_0}$이고, 여기서 $u_0$는 준동형
$u\otimes1:M\otimes_A A/\mathfrak{J}\to
N\otimes_A A/\mathfrak{J}$이다~;
결론은 $\theta$와 $u$가 동시에 전사이고, 또 $\theta_0$와
$u_0$가 동시에 전사라는 사실
(\hyperref[I.1.3.9-ko]{1.3.9} 및
\hyperref[I.10.10.2-ko]{10.10.2})과 $u$와 $u_0$가 동시에 전사라는
사실 (\hyperref[0.7.1.14-ko]{0, 7.1.14})에서 따른다.

\begin{env}[10.11.6]
\label{I.10.11.6-ko}
정리 (\hyperref[I.10.11.3-ko]{10.11.3})은 모든 연접
$\mathscr{O}_{\mathfrak{X}}$-가군 $\mathscr{F}$를
\emph{위상} $\mathscr{O}_{\mathfrak{X}}$-가군으로 볼 수 있음을
보여 준다. 즉 이를 \emph{의사이산} 군의 층 $\mathscr{F}_n$들의
사영극한으로 본다
(\hyperref[0.3.8.1-ko]{0, 3.8.1}). 따라서
(\hyperref[I.10.11.4-ko]{10.11.4})에 의해 준동형
$u:\mathscr{F}\to\mathscr{G}$가 연접
$\mathscr{O}_{\mathfrak{X}}$-가군 사이의 준동형이면 자동으로
\emph{연속}이다

\oldpage[I]{205}

(\hyperref[0.3.8.2-ko]{0, 3.8.2}). 또한 $\mathscr{H}$가 연접
$\mathscr{O}_{\mathfrak{X}}$-부분가군으로서 연접
$\mathscr{O}_{\mathfrak{X}}$-가군 $\mathscr{F}$에 들어 있으면, 모든 열린집합
$U\subset\mathfrak{X}$에 대하여 $\Gamma(U,\mathscr{H})$는 위상군
$\Gamma(U,\mathscr{F})$의 \emph{닫힌} 부분군이다. 실제로 함자
$\Gamma$가 왼쪽 완전이므로 $\Gamma(U,\mathscr{H})$는 준동형
$\Gamma(U,\mathscr{F})\to\Gamma(U,\mathscr{F}/\mathscr{H})$의
핵이다. 이 준동형은 앞서 본 바에 따라 \emph{연속}인데,
$\mathscr{F}/\mathscr{H}$가 연접이기 때문이다
(\hyperref[0.5.3.4-ko]{0, 5.3.4})~; 따라서 이 주장은
$\Gamma(U,\mathscr{F}/\mathscr{H})$가 분리 위상군이라는 사실에서
따른다.
\end{env}

\begin{proposition}[10.11.7]
\label{I.10.11.7-ko}
$\mathscr{F}$와 $\mathscr{G}$를 두 연접
$\mathscr{O}_{\mathfrak{X}}$-가군이라 하자.
(\hyperref[I.10.11.3-ko]{10.11.3})의 표기를 쓰면 다음 위상
$\mathscr{O}_{\mathfrak{X}}$-가군의 함자적인 표준 동형들을 정의할 수 있다
(\hyperref[I.10.11.6-ko]{10.11.6})~:
\[
  \label{I.10.11.7.1-ko}
  \mathscr{F}\otimes_{\mathscr{O}_{\mathfrak{X}}}\mathscr{G}
  \xrightarrow{\sim}
  \varprojlim_n
  (\mathscr{F}_n\otimes_{\mathscr{O}_{X_n}}\mathscr{G}_n)
  \tag{10.11.7.1}
\]
\[
  \label{I.10.11.7.2-ko}
  \mathscr{H}\!om_{\mathscr{O}_{\mathfrak{X}}}
  (\mathscr{F},\mathscr{G})
  \xrightarrow{\sim}
  \varprojlim_n
  \mathscr{H}\!om_{\mathscr{O}_{X_n}}(\mathscr{F}_n,\mathscr{G}_n)
  \tag{10.11.7.2}
\]
\end{proposition}

동형 (\hyperref[I.10.11.7.1-ko]{10.11.7.1})의 존재는 공식
\[
  \mathscr{F}_n\otimes_{\mathscr{O}_{X_n}}\mathscr{G}_n
  =
  (\mathscr{F}\otimes_{\mathscr{O}_{\mathfrak{X}}}
  \mathscr{O}_{X_n})
  \otimes_{\mathscr{O}_{X_n}}
  (\mathscr{G}\otimes_{\mathscr{O}_{\mathfrak{X}}}
  \mathscr{O}_{X_n})
  =
  (\mathscr{F}\otimes_{\mathscr{O}_{\mathfrak{X}}}\mathscr{G})
  \otimes_{\mathscr{O}_{\mathfrak{X}}}\mathscr{O}_{X_n}
\]
과 (\hyperref[I.10.11.3-ko]{10.11.3})에서 따른다. 두 변을 위상을
고려하지 않은 가군의 층으로 보았을 때, 동형
(\hyperref[I.10.11.7.2-ko]{10.11.7.2})은
$\mathscr{H}\!om_{\mathscr{O}_{\mathfrak{X}}}
(\mathscr{F},\mathscr{G})$와
$\mathscr{H}\!om_{\mathscr{O}_{X_n}}(\mathscr{F}_n,\mathscr{G}_n)$의
단면의 정의와 $\mathfrak{X}$의 임의의 뇌터 아핀 형식적 열린집합 위에
유도된 준스킴에 (\hyperref[I.10.11.4.1-ko]{10.11.4.1})을 적용하여 얻는
동형의 존재에서 따른다. 이제 준콤팩트 집합 위에서 동형
(\hyperref[I.10.11.7.2-ko]{10.11.7.2})이 쌍연속임을 보이는 것만 남는다.
따라서 $\mathfrak{X}=\operatorname{Spf}(A)$이고 $A$가 뇌터 아딕환인
경우로 환원되며, 이때 (\hyperref[I.10.10.5-ko]{10.10.5})에 따라
$\mathscr{F}=M^\Delta$, $\mathscr{G}=N^\Delta$이고 $M$, $N$은 유한 생성
$A$-가군이다~; (\hyperref[I.10.10.2.1-ko]{10.10.2.1}),
(\hyperref[I.10.10.2.3-ko]{10.10.2.3}) 및
(\hyperref[I.1.3.12-ko]{1.3.12, (ii)})을 고려하면, 표준 동형
$\operatorname{Hom}_A(M,N)\xrightarrow{\sim}
\varprojlim_n\operatorname{Hom}_{A_n}(M_n,N_n)$
($M_n=M\otimes_A A_n$, $N_n=N\otimes_A A_n$)이 연속임을 보이는 문제로
환원되는데, 이는 (\hyperref[0.7.8.2-ko]{0, 7.8.2})에서 증명되었다.

\begin{env}[10.11.8]
\label{I.10.11.8-ko}
$\operatorname{Hom}_{\mathscr{O}_{\mathfrak{X}}}
(\mathscr{F},\mathscr{G})$는 위상군의 층
$\mathscr{H}\!om_{\mathscr{O}_{\mathfrak{X}}}
(\mathscr{F},\mathscr{G})$의 단면군이므로 군 위상을 갖는다.
$\mathfrak{X}$가 \emph{뇌터}이면,
(\hyperref[I.10.11.7.2-ko]{10.11.7.2})에 따라 이 군에서 $0$의 근방들의
기본계는 부분군
$\operatorname{Hom}_{\mathscr{O}_{\mathfrak{X}}}
(\mathscr{F},\mathscr{J}^n\mathscr{G})$ ($n$은 임의)을 취하여 얻는다.
\end{env}

\begin{proposition}[10.11.9]
\label{I.10.11.9-ko}
$\mathfrak{X}$를 뇌터 형식적 준스킴,
$\mathscr{F}$와 $\mathscr{G}$를 두 연접
$\mathscr{O}_{\mathfrak{X}}$-가군이라 하자. 위상군
$\operatorname{Hom}_{\mathscr{O}_{\mathfrak{X}}}
(\mathscr{F},\mathscr{G})$에서 전사인 준동형들(각각 단사인 준동형들,
전단사인 준동형들)은 열린 부분집합을 이룬다.
\end{proposition}

(\hyperref[I.10.11.5-ko]{10.11.5})에 따라
$\operatorname{Hom}_{\mathscr{O}_{\mathfrak{X}}}
(\mathscr{F},\mathscr{G})$ 안의 전사 준동형들의 집합은 다음 연속사상
$\operatorname{Hom}_{\mathscr{O}_{\mathfrak{X}}}
(\mathscr{F},\mathscr{G})\to\allowbreak
\operatorname{Hom}_{\mathscr{O}_{X_0}}(\mathscr{F}_0,\mathscr{G}_0)$에
의한 이산군
$\operatorname{Hom}_{\mathscr{O}_{X_0}}(\mathscr{F}_0,\mathscr{G}_0)$의
어떤 부분집합의 역상이므로 첫 번째 주장이 따른다. 두 번째 주장을
증명하기 위해 $\mathfrak{X}$를 유한 개의 뇌터 아핀 형식적 열린집합
$U_i$로 덮자. $\theta\in
\operatorname{Hom}_{\mathscr{O}_{\mathfrak{X}}}
(\mathscr{F},\mathscr{G})$가 단사일 필요충분조건은 각 연속 제한사상
$\operatorname{Hom}_{\mathscr{O}_{\mathfrak{X}}}
(\mathscr{F},\mathscr{G})\to
\operatorname{Hom}_{\mathscr{O}_{\mathfrak{X}}|U_i}
(\mathscr{F}|U_i,\mathscr{G}|U_i)$ 아래에서 얻는 상이 모두 단사인 것이다~;
따라서 아핀인 경우로 환원되며, 이 경우에는 이미
(\hyperref[0.7.8.3-ko]{0, 7.8.3})에서 증명되었다.

\oldpage[I]{206}

\subsection{형식적 준스킴의 아딕 사상.}
\label{subsection:I.10.12-ko}

\begin{env}[10.12.1]
\label{I.10.12.1-ko}
$\mathfrak{X}$, $\mathfrak{S}$를 두 \emph{국소 뇌터} 형식적
준스킴이라 하자~; 사상 $f:\mathfrak{X}\to\mathfrak{S}$가
\emph{아딕}이라 함은 아이디얼 $\mathscr{J}$가 존재하여 그것이
$\mathfrak{S}$의 정의 아이디얼이고
$\mathscr{K}=f^*(\mathscr{J})\mathscr{O}_{\mathfrak{X}}$가
$\mathfrak{X}$의 정의 아이디얼인 것을 뜻한다~; 이때
$\mathfrak{X}$를
\emph{$\mathfrak{S}$-아딕 형식적 준스킴}이라고도 한다
($f$에 관하여). 이 경우 $\mathscr{J}_1$을 $\mathfrak{S}$의
\emph{임의의} 정의 아이디얼이라 하면,
$\mathscr{K}_1=f^*(\mathscr{J}_1)\mathscr{O}_{\mathfrak{X}}$는
$\mathfrak{X}$의 정의 아이디얼이다. 실제로 문제는 국소적이므로
$\mathfrak{X}$와 $\mathfrak{S}$가 뇌터 아핀 형식적 준스킴이라고
가정해도 된다~; 따라서 정수 $n$이 존재하여
$\mathscr{J}^n\subset\mathscr{J}_1$이고
$\mathscr{J}_1^n\subset\mathscr{J}$이다
(\hyperref[I.10.3.6-ko]{10.3.6} 및
\hyperref[0.7.1.4-ko]{0, 7.1.4}). 그러므로
$\mathscr{K}^n\subset\mathscr{K}_1$이고
$\mathscr{K}_1^n\subset\mathscr{K}$이다. 첫 번째 관계는
$\mathscr{K}_1=\mathfrak{K}_1^\Delta$임을 보인다. 여기서
$\mathfrak{K}_1$은
$A=\Gamma(\mathfrak{X},\mathscr{O}_{\mathfrak{X}})$의 열린
아이디얼이다. 두 번째 관계는 $\mathfrak{K}_1$이 $A$의 정의
아이디얼임을 보이므로
(\hyperref[0.7.1.4-ko]{0, 7.1.4}), 이로써 주장이 따른다.
\end{env}

앞의 내용에서 $\mathfrak{X}$와 $\mathfrak{Y}$가 두
$\mathfrak{S}$-아딕 형식적 준스킴이면,
\emph{모든 $\mathfrak{S}$-사상
$u:\mathfrak{X}\to\mathfrak{Y}$은 아딕이다}라는 것이 곧바로 따른다~:
실제로 $f:\mathfrak{X}\to\mathfrak{S}$,
$g:\mathfrak{Y}\to\mathfrak{S}$를 각각 구조 사상이라 하고,
$\mathscr{J}$를 $\mathfrak{S}$의 정의 아이디얼이라 하자. 그러면
$f=g\circ u$이므로
$u^*(g^*(\mathscr{J})\mathscr{O}_{\mathfrak{Y}})
\mathscr{O}_{\mathfrak{X}}=f^*(\mathscr{J})
\mathscr{O}_{\mathfrak{X}}$는 $\mathfrak{X}$의 정의 아이디얼이고,
가정에 따라 $g^*(\mathscr{J})\mathscr{O}_{\mathfrak{Y}}$는
$\mathfrak{Y}$의 정의 아이디얼이다.

\begin{env}[10.12.2]
\label{I.10.12.2-ko}
이하에서는 국소 뇌터 형식적 준스킴 $\mathfrak{S}$와 아이디얼
$\mathscr{J}$를 고정하되, 후자는 $\mathfrak{S}$의 정의 아이디얼이라고
가정하자~; 또한
$S_n=(\mathfrak{S},\mathscr{O}_{\mathfrak{S}}/\mathscr{J}^{n+1})$로
두겠다. $\mathfrak{S}$-아딕 형식적 준스킴들(국소 뇌터인 것들)은
자명하게 하나의 \emph{범주}를 이룬다. 귀납계 $(X_n)$이 (보통 의미의)
국소 뇌터 $S_n$-준스킴들로 이루어져 있을 때, 다음 조건을 만족하면
이를 \emph{아딕 $(S_n)$-귀납계}라 하겠다. 구조 사상
$f_n:X_n\to S_n$들에 대하여 $m\leq n$이면 도식
\[
  \label{I.10.12.2.1-ko}
  \xymatrix{
    X_n \ar[d]_{f_n}
    & X_m \ar[l] \ar[d]^{f_m}\\
    S_n
    & S_m \ar[l]
  }
  \tag{10.12.2.1}
\]
은 가환이고, \emph{$X_m$을 곱
$X_n\times_{S_n}S_m=(X_n)_{(S_m)}$과 동일시한다}. 아딕 귀납계들도
하나의 \emph{범주}를 이룬다~: 실제로 이러한 계들의 사상
$(X_n)\to(Y_n)$을 \emph{$S_n$-사상 $u_n:X_n\to Y_n$들로 이루어진
귀납계}로 정의하되, $u_m$이 $(u_n)_{(S_m)}$과 동일시되는 조건을
모든 $m\leq n$에 대하여 만족하도록 하면 충분하다. 그러면~:
\end{env}

\begin{theorem}[10.12.3]
\label{I.10.12.3-ko}
$\mathfrak{S}$-아딕 형식적 준스킴들의 범주와 아딕
$(S_n)$-귀납계들의 범주 사이에는 표준 동치가 있다.
\end{theorem}

이 동치는 다음과 같이 얻는다~: $\mathfrak{X}$가
$\mathfrak{S}$-아딕 형식적 준스킴이고,
$f:\mathfrak{X}\to\mathfrak{S}$가 구조 사상이라 하면,
$\mathscr{K}=f^*(\mathscr{J})\mathscr{O}_{\mathfrak{X}}$는
$\mathfrak{X}$의 정의 아이디얼이며, $\mathfrak{X}$에 귀납계
$X_n=(\mathfrak{X},\mathscr{O}_{\mathfrak{X}}/\mathscr{K}^{n+1})$을
대응시킨다. 여기서 $f_n:X_n\to S_n$은 $f$에 대응하는 구조 사상이다
(\hyperref[I.10.5.6-ko]{10.5.6}). 먼저 $(X_n)$이
\emph{아딕 귀납계}임을 보이자~: $f=(\psi,\theta)$라 쓰면
$\psi^*(\mathscr{J})\mathscr{O}_{\mathfrak{X}}=\mathscr{K}$이고, 따라서
$\psi^*(\mathscr{J}^n)\mathscr{O}_{\mathfrak{X}}=\mathscr{K}^n$이며, 이는
모든 $n$에 대하여 성립한다.
또한 함자 $\psi^*$의 완전성에 의해
$\mathscr{K}^{m+1}/\mathscr{K}^{n+1}
=\psi^*(\mathscr{J}^{m+1}/\mathscr{J}^{n+1})
(\mathscr{O}_{\mathfrak{X}}/\mathscr{K}^{n+1})$이 $m\leq n$일 때
성립한다. 그러므로 원하는 결론은
(\hyperref[I.4.4.5-ko]{4.4.5})에서 따른다. 더욱이
$\mathfrak{S}$-사상 $u:\mathfrak{X}\to\mathfrak{Y}$이 두
$\mathfrak{S}$-아딕 형식적 준스킴 사이의 사상일 때에는, 명백한 표기를 쓰면,

\oldpage[I]{207}

$S_n$-사상 $u_n:X_n\to Y_n$들의 귀납계로서 $u_m$이
$(u_n)_{(S_m)}$와 동일시되는 조건이 $m\leq n$일 때 성립하는 것이
대응함도 곧바로 확인된다.

이렇게 동치를 올바르게 정의했음은 다음의 더 정밀한 명제에서 따른다~:

\begin{proposition}[10.12.3.1]
\label{I.10.12.3.1-ko}
$(X_n)$을 $S_n$-준스킴들의 귀납계라 하자. 구조 사상
$f_n:X_n\to S_n$들이 도식
(\hyperref[I.10.12.2.1-ko]{10.12.2.1})을 가환이게 하고, $X_m$을
$X_n\times_{S_n}S_m$와 동일시하는 조건이 $m\leq n$일 때 성립한다고
가정하자.
그러면 귀납계 $(X_n)$은 (\hyperref[I.10.6.3-ko]{10.6.3})의 조건
b)와 c)를 만족한다. 그 귀납극한을 $\mathfrak{X}$라 하고,
$f:\mathfrak{X}\to\mathfrak{S}$를 귀납계 $(f_n)$의 귀납극한 사상이라
하자. 이때 $X_0$이 국소 뇌터이면
$\mathfrak{X}$도 국소 뇌터이고 $f$는 아딕 사상이다.
\end{proposition}

$\mathscr{O}_{S_n}$에서 부분준스킴 $S_m$을 $S_n$ 안에 정의하는
아이디얼층은 멱영이다. 따라서 (\hyperref[I.4.4.5-ko]{4.4.5})에
의해 $\mathscr{O}_{X_n}$에서 부분준스킴 $X_m$을 $X_n$ 안에 정의하는
아이디얼층도 멱영이므로, (\hyperref[I.10.6.3-ko]{10.6.3})의 조건들은
실제로 만족된다. 그러므로 문제는 $\mathfrak{X}$와 $\mathfrak{S}$에서
국소적이며, $\mathfrak{S}=\operatorname{Spf}(A)$,
$\mathscr{J}=\mathfrak{J}^\Delta$라 가정해도 된다. 여기서 $A$는
뇌터 $\mathfrak{J}$-아딕환이고, $X_n=\operatorname{Spec}(B_n)$이다.
$A_n=A/\mathfrak{J}^{n+1}$이라 하면,
가정에 의해 $B_0$은 뇌터이고,
$\mathfrak{J}_n=\mathfrak{J}/\mathfrak{J}^{n+1}$이라 놓을 때
$B_m=B_n/\mathfrak{J}_n^{m+1}B_n$이다. 따라서 $B_n\to B_0$의 핵은
$\mathfrak{K}_n=\mathfrak{J}_nB_n$이고, $B_n\to B_m$의 핵은
$\mathfrak{K}_n^{m+1}$이며, 이는 $m\leq n$일 때 성립한다. 또한 $A_1$이
뇌터이므로 $\mathfrak{J}_1$은 유한 생성 $A_1$-가군이다. 따라서
$\overline{\mathfrak{K}}_1=\mathfrak{K}_1/\mathfrak{K}_1^2$은
유한 생성 $B_1$-가군이고, \emph{따라서 특히}
$B_0=B_1/\mathfrak{K}_1$-가군으로도 유한 생성이다. 이제
$\mathfrak{X}$가 뇌터임은 (\hyperref[I.10.6.4-ko]{10.6.4})에서
따른다. $B=\varprojlim B_n$이라 하면
$\mathfrak{X}=\operatorname{Spf}(B)$이고, $\mathfrak{K}$를
$B\to B_0$의 핵이라 하면 $B_n=B/\mathfrak{K}^{n+1}$이다.
$\rho_n:A/\mathfrak{J}^{n+1}\to B/\mathfrak{K}^{n+1}$을 $f_n$에
대응하는 준동형이라 하면
\[
  \mathfrak{K}/\mathfrak{K}^{n+1}
  =(B/\mathfrak{K}^{n+1})\rho_n(\mathfrak{J}/\mathfrak{J}^{n+1})
\]
이다. 한편 준동형 $\rho:A\to B$가 $f$에 대응하며
$\varprojlim\rho_n$과 같으므로, 아이디얼 $\mathfrak{J}B$는 $B$의
아이디얼로서 $\mathfrak{K}$ 안에서 조밀하다. 그런데 $B$의 모든
아이디얼은 닫혀 있으므로 (\hyperref[0.7.3.5-ko]{0, 7.3.5}),
$\mathfrak{K}=\mathfrak{J}B$이다.
$\mathscr{K}=\mathfrak{K}^\Delta$라 하면 관계
$f^*(\mathscr{J})\mathscr{O}_{\mathfrak{X}}=\mathscr{K}$는
(\hyperref[I.10.10.9-ko]{10.10.9})에서 따르며, 이것으로 증명이 끝난다.

\begin{env}[10.12.3.2]
\label{I.10.12.3.2-ko}
앞의 동치는 두 $\mathfrak{S}$-아딕 형식적 준스킴
$\mathfrak{X}$, $\mathfrak{Y}$에 대하여 다음 \emph{표준 전단사}를 준다~:
\[
  \operatorname{Hom}_{\mathfrak{S}}(\mathfrak{X},\mathfrak{Y})
  \xrightarrow{\sim}
  \varprojlim_n\operatorname{Hom}_{S_n}(X_n,Y_n)
\]
여기서 사영극한은 사상 $u_n\to(u_n)_{(S_m)}$들에 관하여 취하며,
$m\leq n$이다.
\end{env}

\subsection{유한형 사상.}
\label{subsection:I.10.13-ko}

\begin{proposition}[10.13.1]
\label{I.10.13.1-ko}
$\mathfrak{Y}$를 국소 뇌터 형식적 준스킴,
$\mathscr{K}$를 $\mathfrak{Y}$의 정의 아이디얼,
$f:\mathfrak{X}\to\mathfrak{Y}$를 형식적 준스킴의 사상이라 하자. 다음
조건들은 동치이다~:
\begin{enumerate}
  \item[a)] $\mathfrak{X}$는 국소 뇌터이고, $f$는 아딕 사상이며
    (\hyperref[I.10.12.1-ko]{10.12.1}),
    $\mathscr{J}=f^*(\mathscr{K})\mathscr{O}_{\mathfrak{X}}$라 놓을 때
    사상 $f_0:(\mathfrak{X},\mathscr{O}_{\mathfrak{X}}/\mathscr{J})
    \to(\mathfrak{Y},\mathscr{O}_{\mathfrak{Y}}/\mathscr{K})$는 $f$로부터
    유도되며 유한형이다.
  \item[b)] $\mathfrak{X}$는 국소 뇌터이고, 아딕
    $(Y_n)$-귀납계 $(X_n)$의 귀납극한이며, 사상 $X_0\to Y_0$는
    유한형이다.

\oldpage[I]{208}

  \item[c)] $\mathfrak{Y}$의 모든 점은 다음 성질을 갖는 뇌터 아핀
    형식적 열린집합인 근방 $V$를 가진다~:
    \begin{enumerate}
      \item[(\textbf{Q})] $f^{-1}(V)$는 뇌터 아핀 형식적 열린집합
        $U_i$들의 유한 합집합이고, 각각에 대하여 뇌터 아딕환
        $\Gamma(U_i,\mathscr{O}_{\mathfrak{X}})$는
        $\Gamma(V,\mathscr{O}_{\mathfrak{Y}})$ 위의 제한 형식적 멱급수
        대수를 어떤 (필연적으로 닫힌) 아이디얼로 나눈 몫과 위상적으로
        동형이다 (\hyperref[0.7.5.1-ko]{0, 7.5.1}).
    \end{enumerate}
\end{enumerate}
\end{proposition}

a)가 b)를 함의하는 것은 (\hyperref[I.10.12.3-ko]{10.12.3})에 따라
명백하다. b)가 c)를 함의함을 보이기 위해, 문제는 $\mathfrak{Y}$ 위에서
국소적이므로 $\mathfrak{Y}=\operatorname{Spf}(B)$이고 $B$가 뇌터
아딕환이라고 가정할 수 있다. $\mathscr{K}=\mathfrak{K}^\Delta$라
하되, $\mathfrak{K}$는 $B$의 정의 아이디얼이라 하자. 가정에 의해 $X_0$는
$Y_0$ 위에서 유한형이므로, $X_0$는 아핀 열린집합 $U_i$들의 유한
합집합이고, 환 $A_{i0}$, 즉 $X_0$가 $U_i$ 위에 유도하는 아핀
스킴의 환은 $B/\mathfrak{K}$, 즉 $Y_0$의 환 위의 유한 생성 대수이다
(\hyperref[I.6.3.2-ko]{6.3.2}).
(\hyperref[I.5.1.9-ko]{5.1.9})에 따라 $U_i$는 각 뇌터 준스킴
$X_n$에서도 아핀 열린집합이며, 환 $A_{in}$을 $X_n$이 $U_i$ 위에
유도하는 아핀 스킴의 환이라 하면 가정 b)에 의해 $m\leq n$일 때
$A_{im}$은 $A_{in}/\mathfrak{K}^{m+1}A_{in}$과 동형이다. 따라서
$U_i$ 위에서 $\mathfrak{X}$가 유도하는 형식적 준스킴은
$\operatorname{Spf}(A_i)$와 동형이며, 여기서
$A_i=\varprojlim_n A_{in}$이다 (\hyperref[I.10.6.4-ko]{10.6.4}).
$A_i$는 $\mathfrak{K}A_i$-아딕환이고, $A_i/\mathfrak{K}A_i$는
$A_{i0}$과 동형이며 $B/\mathfrak{K}$ 위의 유한 생성 대수이다.
그러므로 (\hyperref[0.7.5.5-ko]{0, 7.5.5})에 따라 $A_i$는 $B$ 위의
제한 형식적 멱급수 대수를 어떤 아이디얼로 나눈 몫과 위상적으로
동형이다. 이 아이디얼은 반드시 닫혀 있는데, 그러한 대수는 뇌터이기 때문이다
(\hyperref[0.7.5.4-ko]{0, 7.5.4}).

c)가 a)를 함의함을 증명하기 위해, $\mathfrak{X}=\operatorname{Spf}(A)$
도 아핀인 경우로 한정할 수 있다. 여기서 $A$는 뇌터 아딕환이고, $B$
위의 제한 형식적 멱급수 대수를 어떤 닫힌 아이디얼로 나눈 몫과
동형이다. 그러면 (\hyperref[0.7.5.5-ko]{0, 7.5.5})에 따라
$A/\mathfrak{K}A$는 $B/\mathfrak{K}$ 위의 유한 생성 대수이고,
$\mathfrak{K}A=\mathfrak{J}$는 $A$의 정의 아이디얼이다. 따라서
(\hyperref[I.10.10.9-ko]{10.10.9})에 의해 a)의 조건들이 성립한다.

명제 (\hyperref[I.10.13.1-ko]{10.13.1})의 조건들이 성립하면, a)는
\emph{모든 정의 아이디얼 $\mathscr{K}$}에 대하여 성립한다. 여기서
정의 아이디얼은 $\mathfrak{Y}$의 것이며, 이는 c)에 따른다. 따라서
b)에서는 \emph{모든 $f_n$이 유한형 사상}임에 유의하자.

\begin{corollary}[10.13.2]
\label{I.10.13.2-ko}
(\hyperref[I.10.13.1-ko]{10.13.1})의 조건들이 성립하면, 모든 뇌터
아핀 형식적 열린집합 $V$는 $\mathfrak{Y}$ 안에서 성질
(\textbf{Q})를 가지며, $\mathfrak{Y}$가 뇌터이면 $\mathfrak{X}$도
뇌터이다.
\end{corollary}

이는 (\hyperref[I.10.13.1-ko]{10.13.1})과
(\hyperref[I.6.3.2-ko]{6.3.2})에서 곧바로 따른다.

\begin{definition}[10.13.3]
\label{I.10.13.3-ko}
(\hyperref[I.10.13.1-ko]{10.13.1})의 동치인 성질 a), b), c)가
성립하면, 사상 $f$를 \emph{유한형}이라고 하거나 $\mathfrak{X}$를
\emph{유한형 $\mathfrak{Y}$-형식적 준스킴} 또는
\emph{$\mathfrak{Y}$ 위의 유한형 형식적 준스킴}이라고 한다.
\end{definition}

\begin{corollary}[10.13.4]
\label{I.10.13.4-ko}
$\mathfrak{X}=\operatorname{Spf}(A)$,
$\mathfrak{Y}=\operatorname{Spf}(B)$를 두 뇌터 아핀 형식적 스킴이라
하자. $\mathfrak{X}$가 $\mathfrak{Y}$ 위에서 유한형이기 위한
필요충분조건은 뇌터 아딕환 $A$가 $B$ 위의 제한 형식적 멱급수 대수를
어떤 닫힌 아이디얼로 나눈 몫과 동형인 것이다.
\end{corollary}

실제로 (\hyperref[I.10.13.1-ko]{10.13.1})의 표기를 사용하자.
$\mathfrak{X}$가 $\mathfrak{Y}$ 위에서 유한형이면
(\hyperref[I.6.3.3-ko]{6.3.3})에 따라 $A/\mathfrak{K}A$는
$(B/\mathfrak{K})$ 위의 유한 생성 대수이고,
$\mathfrak{K}A$는 $A$의 정의 아이디얼이다
(\hyperref[I.10.10.9-ko]{10.10.9}). 그러므로
(\hyperref[0.7.5.5-ko]{0, 7.5.5})에서 결론이 따른다.

\begin{proposition}[10.13.5]
\label{I.10.13.5-ko}
\medskip\noindent
\begin{enumerate}
  \item[\rm{(i)}] 형식적 준스킴의 두 유한형 사상의 합성은
    유한형이다.

\oldpage[I]{209}

  \item[\rm{(ii)}] $\mathfrak{X}$, $\mathfrak{S}$,
    $\mathfrak{S}'$을 세 국소 뇌터(각각 뇌터) 형식적 준스킴이라 하고,
    $f:\mathfrak{X}\to\mathfrak{S}$,
    $g:\mathfrak{S}'\to\mathfrak{S}$를 두 사상이라 하자. $f$가
    유한형이면 $\mathfrak{X}\times_{\mathfrak{S}}\mathfrak{S}'$은
    국소 뇌터(각각 뇌터)이고 $\mathfrak{S}'$ 위에서 유한형이다.
  \item[\rm{(iii)}] $\mathfrak{S}$를 국소 뇌터 형식적 준스킴,
    $\mathfrak{X}'$, $\mathfrak{Y}'$을 두 국소 뇌터
    $\mathfrak{S}$-형식적 준스킴이라 하되,
    $\mathfrak{X}'\times_{\mathfrak{S}}\mathfrak{Y}'$은 국소 뇌터라
    하자.
    $\mathfrak{X}$, $\mathfrak{Y}$가 국소 뇌터
    $\mathfrak{S}$-형식적 준스킴이고,
    $f:\mathfrak{X}\to\mathfrak{X}'$,
    $g:\mathfrak{Y}\to\mathfrak{Y}'$이 두 유한형
    $\mathfrak{S}$-사상이면,
    $\mathfrak{X}\times_{\mathfrak{S}}\mathfrak{Y}$은 국소 뇌터이고
    $f\times_{\mathfrak{S}}g$는 유한형 $\mathfrak{S}$-사상이다.
\end{enumerate}
\end{proposition}

(iii)은 (\hyperref[I.3.5.1-ko]{3.5.1})의 형식적 논증에 의해 (i)과
(ii)에서 따르므로, (i)과 (ii)만 증명하면 충분하다.

$\mathfrak{X}$, $\mathfrak{Y}$, $\mathfrak{Z}$를 세 국소 뇌터 형식적
준스킴이라 하고, $f:\mathfrak{X}\to\mathfrak{Y}$,
$g:\mathfrak{Y}\to\mathfrak{Z}$를 두 유한형 사상이라 하자.
$\mathscr{L}$이 $\mathfrak{Z}$의 정의 아이디얼이면,
$\mathscr{K}=g^*(\mathscr{L})\mathscr{O}_{\mathfrak{Y}}$는
$\mathfrak{Y}$의 정의 아이디얼이고,
$\mathscr{J}=f^*(g^*(\mathscr{L}))\mathscr{O}_{\mathfrak{X}}$는
$\mathfrak{X}$의 정의 아이디얼이다. 다음과 같이 놓자~:
$X_0=(\mathfrak{X},\mathscr{O}_{\mathfrak{X}}/\mathscr{J})$,
$Y_0=(\mathfrak{Y},\mathscr{O}_{\mathfrak{Y}}/\mathscr{K})$,
$Z_0=(\mathfrak{Z},\mathscr{O}_{\mathfrak{Z}}/\mathscr{L})$.
또 $f_0:X_0\to Y_0$, $g_0:Y_0\to Z_0$를 각각 $f$, $g$에 대응하는
사상이라 하자. 가정에 의해 $f_0$와 $g_0$가 유한형이므로,
$g_0\circ f_0$도 유한형이다 (\hyperref[I.6.3.4-ko]{6.3.4}). 이
사상은 $g\circ f$에 대응하므로, (\hyperref[I.10.13.1-ko]{10.13.1})에
따라 $g\circ f$는 유한형이다.

(ii)의 조건 아래에서 $\mathfrak{S}$ (각각 $\mathfrak{X}$,
$\mathfrak{S}'$)는 국소 뇌터 준스킴들의 열 $(S_n)$ (각각 $(X_n)$,
$(S'_n)$)의 귀납극한이고, (\hyperref[I.10.13.1-ko]{10.13.1})에 따라
$X_m=X_n\times_{S_n}S_m$라고 가정할 수 있으며, 이는 모든
$m\leq n$에 대해 성립한다. 그러면
형식적 준스킴
$\mathfrak{X}\times_{\mathfrak{S}}\mathfrak{S}'$은 준스킴
$X_n\times_{S_n}S'_n$들의 귀납극한이고
(\hyperref[I.10.7.4-ko]{10.7.4}), 다음이 성립한다~:
\[
  X_m\times_{S_m}S'_m
  =(X_n\times_{S_n}S_m)\times_{S_m}S'_m
  =(X_n\times_{S_n}S'_n)\times_{S'_n}S'_m.
\]
또한 $X_0\times_{S_0}S'_0$은 $X_0$가 $S_0$ 위에서 유한형이므로
국소 뇌터이다
(\hyperref[I.6.3.8-ko]{6.3.8}). 따라서 먼저
(\hyperref[I.10.12.3.1-ko]{10.12.3.1})에 의해
$\mathfrak{X}\times_{\mathfrak{S}}\mathfrak{S}'$은 국소 뇌터이다.
더 나아가 $X_0\times_{S_0}S'_0$은 $S'_0$ 위에서 유한형이므로
(\hyperref[I.6.3.8-ko]{6.3.8}),
(\hyperref[I.10.12.3.1-ko]{10.12.3.1})과
(\hyperref[I.10.13.1-ko]{10.13.1})에 따라
$\mathfrak{X}\times_{\mathfrak{S}}\mathfrak{S}'$은
$\mathfrak{S}'$ 위에서 유한형이다. 이것으로 (ii)의 증명이 끝난다.
뇌터 준스킴에 관한 주장은 (\hyperref[I.6.3.8-ko]{6.3.8})의 즉각적인
귀결이다.

\begin{corollary}[10.13.6]
\label{I.10.13.6-ko}
(\hyperref[I.10.9.9-ko]{10.9.9})의 가정 아래에서 $f$가 유한형
사상이면, 완비화한 것들 사이로의 연장 $\widehat f$도 유한형이다.
\end{corollary}

\subsection{형식적 준스킴의 닫힌 부분준스킴.}
\label{subsection:I.10.14-ko}

\begin{proposition}[10.14.1]
\label{I.10.14.1-ko}
$\mathfrak{X}$를 국소 뇌터 형식적 준스킴,
$\mathscr{A}$를 $\mathscr{O}_{\mathfrak{X}}$의 연접 아이디얼층이라
하자. $\mathfrak{Y}$가 $\mathscr{O}_{\mathfrak{X}}/\mathscr{A}$의
(닫힌) 지지집합이면, 위상환 달린 공간
$(\mathfrak{Y},(\mathscr{O}_{\mathfrak{X}}/\mathscr{A})|\mathfrak{Y})$은
국소 뇌터 형식적 준스킴이며, $\mathfrak{X}$가 뇌터이면 이 형식적
준스킴도 뇌터이다.
\end{proposition}

(\hyperref[I.10.10.3-ko]{10.10.3})과
(\hyperref[0.5.3.4-ko]{0, 5.3.4})에 따라
$\mathscr{O}_{\mathfrak{X}}/\mathscr{A}$는 구조층 위의 연접 가군이고, 따라서 그
지지집합 $\mathfrak{Y}$는 닫혀 있다
(\hyperref[0.5.2.2-ko]{0, 5.2.2}). $\mathscr{J}$를
$\mathfrak{X}$의 정의 아이디얼이라 하고,
$X_n=(\mathfrak{X},\mathscr{O}_{\mathfrak{X}}/\mathscr{J}^{n+1})$이라
하자~; 환의 층 $\mathscr{O}_{\mathfrak{X}}/\mathscr{A}$는 층들
$\mathscr{O}_{\mathfrak{X}}/(\mathscr{A}+\mathscr{J}^{n+1})
=(\mathscr{O}_{\mathfrak{X}}/\mathscr{A})
\otimes_{\mathscr{O}_{\mathfrak{X}}}
(\mathscr{O}_{\mathfrak{X}}/\mathscr{J}^{n+1})$의 사영극한이다
(\hyperref[I.10.11.3-ko]{10.11.3}). 이 층들은 모두 지지집합
$\mathfrak{Y}$를 갖는다. 층
$(\mathscr{A}+\mathscr{J}^{n+1})/\mathscr{J}^{n+1}$은
$\mathscr{O}_{\mathfrak{X}}$-가군으로서 연접이다. 실제로
$\mathscr{J}^{n+1}$이 연접이므로,
$(\mathscr{A}+\mathscr{J}^{n+1})/\mathscr{J}^{n+1}$은 또한 연접
$(\mathscr{O}_{\mathfrak{X}}/\mathscr{J}^{n+1})$-가군이다
(\hyperref[0.5.3.10-ko]{0, 5.3.10})~; $Y_n$을 이 아이디얼층으로
정의되는 $X_n$의 닫힌 부분준스킴이라 하면, 위상환 달린 공간
$(\mathfrak{Y},(\mathscr{O}_{\mathfrak{X}}/\mathscr{A})|\mathfrak{Y})$

\oldpage[I]{210}
은 $Y_n$들의 귀납극한인 형식적 준스킴임이 곧바로 확인된다.
(\hyperref[I.10.6.4-ko]{10.6.4})의 조건들이 만족되므로, 이 형식적
준스킴은 국소 뇌터이고, $\mathfrak{X}$가 뇌터이면 뇌터이다. 실제로
이 경우에는 (\hyperref[I.6.1.4-ko]{6.1.4})에 따라 $Y_0$가 뇌터이다.

\begin{definition}[10.14.2]
\label{I.10.14.2-ko}
형식적 준스킴 $\mathfrak{X}$의 \emph{닫힌 부분준스킴}이라 함은
$(\mathfrak{Y},(\mathscr{O}_{\mathfrak{X}}/\mathscr{A})|\mathfrak{Y})$
꼴의 모든 형식적 준스킴을 말한다. 여기서 $\mathscr{A}$는
$\mathscr{O}_{\mathfrak{X}}$의 연접 아이디얼층이다~; 이 형식적
준스킴을 $\mathscr{A}$가 정의하는 닫힌 부분준스킴이라고 한다.
\footnote{\textbf{번역 주.} 인쇄 원문은 여기와 바로 다음 문장에서
단지 연접 가군이라고 쓰지만, 몫환의 구성과 뒤따르는 대응에는 연접
아이디얼층이 필요하므로 두 곳 모두 이를 교정했다.}
\end{definition}

이와 같이 정의된 $\mathscr{O}_{\mathfrak{X}}$의 연접 아이디얼층들과
$\mathfrak{X}$의 닫힌 부분준스킴들 사이의 대응은 전단사임이 명백하다.

위상환 달린 공간의 사상
$j=(\psi,\theta):\mathfrak{Y}\to\mathfrak{X}$에서 $\psi$는 단사사상
$\mathfrak{Y}\to\mathfrak{X}$이고 $\theta^\sharp$는 표준 준동형
$\mathscr{O}_{\mathfrak{X}}\to\mathscr{O}_{\mathfrak{X}}/\mathscr{A}$이다.
이는 명백히 (\hyperref[I.10.4.5-ko]{10.4.5}) 형식적 준스킴의
사상이며, 이를 $\mathfrak{Y}$에서 $\mathfrak{X}$로 가는
\emph{표준 포함 사상}이라 한다. $\mathfrak{X}=\operatorname{Spf}(A)$이고
$A$가 뇌터 아딕환이면, $\mathscr{A}=\mathfrak{a}^\Delta$이고 여기서
$\mathfrak{a}$는 $A$의 아이디얼이다
(\hyperref[I.10.10.5-ko]{10.10.5}). 앞의 내용에서 곧바로
$\mathfrak{Y}=\operatorname{Spf}(A/\mathfrak{a})$임이 동형을 제외하면
따르고, $j$는 (\hyperref[I.10.2.2-ko]{10.2.2}) 표준 준동형
$A\to A/\mathfrak{a}$에 대응한다.

국소 뇌터 형식적 준스킴의 사상
$f:\mathfrak{Z}\to\mathfrak{X}$가 \emph{닫힌 몰입 사상}이라 함은
이 사상이 $\mathfrak{Z}\xrightarrow{g}\mathfrak{Y}\xrightarrow{j}\mathfrak{X}$로
인수분해되고, 여기서 $g$는 $\mathfrak{Z}$에서 닫힌 부분준스킴
$\mathfrak{Y}$로 가는 동형사상이며, 이 부분준스킴은 $\mathfrak{X}$의
닫힌 부분준스킴이고 $j$는 표준 포함
사상인 것을 뜻한다. $j$는 환 달린 공간의 단사사상이므로 $g$와
$\mathfrak{Y}$는 반드시 \emph{유일}하다.

\begin{proposition}[10.14.3]
\label{I.10.14.3-ko}
닫힌 몰입 사상은 유한형 사상이다.
\end{proposition}

$\mathfrak{X}$가 아핀 형식적 스킴 $\operatorname{Spf}(A)$이고
$\mathfrak{Y}=\operatorname{Spf}(A/\mathfrak{a})$인 경우로 곧바로
환원된다~; 명제는 (\hyperref[I.10.13.1-ko]{10.13.1, c)})에서 따른다.

\begin{lemma}[10.14.4]
\label{I.10.14.4-ko}
$f:\mathfrak{Y}\to\mathfrak{X}$를 국소 뇌터 형식적 준스킴의
사상이라 하고, $(U_\alpha)$를 $f(\mathfrak{Y})$의 덮개라 하자. 이
덮개의 원소들은 $\mathfrak{X}$의 뇌터 아핀 형식적 열린집합이고, 각
$f^{-1}(U_\alpha)$는
$\mathfrak{Y}$의 뇌터 아핀 형식적 열린집합이라고 하자. $f$가 닫힌
몰입 사상이기 위한 필요충분조건은 $f(\mathfrak{Y})$가
$\mathfrak{X}$의 닫힌 부분집합이고, 모든 $\alpha$에 대하여 $f$를
$f^{-1}(U_\alpha)$에 제한한 사상이
(\hyperref[I.10.4.6-ko]{10.4.6})에 따라 전사 준동형
$\Gamma(U_\alpha,\mathscr{O}_{\mathfrak{X}})\to
\Gamma(f^{-1}(U_\alpha),\mathscr{O}_{\mathfrak{Y}})$에 대응하는 것이다.
\end{lemma}

이 조건들이 필요함은 명백하다. 역으로 이 조건들이 성립한다고 하고,
$\mathfrak{a}_\alpha$를 준동형
$\Gamma(U_\alpha,\mathscr{O}_{\mathfrak{X}})\to
\Gamma(f^{-1}(U_\alpha),\mathscr{O}_{\mathfrak{Y}})$의 핵이라 하자.
연접 아이디얼층 $\mathscr{A}$를 $\mathscr{O}_{\mathfrak{X}}$ 위에
다음과 같이 정의한다. $\mathscr{A}|U_\alpha=\mathfrak{a}_\alpha^\Delta$로
두고, $\mathscr{A}$는 $U_\alpha$들의 합집합의 여집합에서 단위
아이디얼로 둔다. 실제로 $f(\mathfrak{Y})$가 닫혀 있고
$(\mathscr{O}_{\mathfrak{X}}|U_\alpha)/\mathfrak{a}_\alpha^\Delta$의
지지집합이 $U_\alpha\cap f(\mathfrak{Y})$이므로,
$\mathfrak{a}_\alpha^\Delta$와 $\mathfrak{a}_\beta^\Delta$가 뇌터
아핀 형식적 열린집합 $V\subset U_\alpha\cap U_\beta$ 위에서 같은
층을 유도함을 확인하면 충분하다.\footnote{\textbf{번역 주.} 인쇄
원문은 여집합에서 영 아이디얼로 둔다고 쓰지만, 닫힌 부분집합인
사상의 상 밖에서 몫층이 영층이 되려면 단위 아이디얼이어야 한다.
또한 그 닫힌 부분집합은 핵 아이디얼층 자체의 지지집합이 아니라 그
몫층의 지지집합이므로 두 곳을 함께 교정했다.}
$f$를 $f^{-1}(U_\alpha)$에 제한한 사상은 이 형식적 준스킴에서
$U_\alpha$로 가는 닫힌 몰입 사상이므로, $f^{-1}(V)$는
$f^{-1}(U_\alpha)$의 뇌터 아핀 형식적 열린집합이고 $f$를
$f^{-1}(V)$에 제한한 사상도 닫힌 몰입 사상이다~; 이 제한에 대응하는
전사 준동형의 핵을 $\mathfrak{b}$라 하자. 그 준동형은
$\Gamma(V,\mathscr{O}_{\mathfrak{X}})\to
\Gamma(f^{-1}(V),\mathscr{O}_{\mathfrak{Y}})$이다. 그러면
(\hyperref[I.10.10.2-ko]{10.10.2})에 따라
$\mathfrak{a}_\alpha^\Delta$가 $\mathfrak{b}^\Delta$를 $V$ 위에
유도함은 명백하다. 이렇게 아이디얼층 $\mathscr{A}$를 정의하면
$f=j\circ g$임이 명백하다. 여기서
$j:\mathfrak{Z}\to\mathfrak{X}$는 닫힌 부분준스킴 $\mathfrak{Z}$에서
$\mathfrak{X}$로 가는 표준 포함 사상이고, 이 부분준스킴은
$\mathscr{A}$가 정의하며,
$g$는 $\mathfrak{Y}$에서 $\mathfrak{Z}$로 가는 동형사상이다.
\footnote{\textbf{번역 주.} 인쇄 원문은 합성의 순서를 거꾸로 쓰지만,
선언된 정의역과 공역에 맞추어 바로잡았다.}

\begin{proposition}[10.14.5]
\label{I.10.14.5-ko}
\begin{enumerate}
  \item[\rm{(i)}] $f:\mathfrak{Z}\to\mathfrak{Y}$,
  $g:\mathfrak{Y}\to\mathfrak{X}$가 국소 뇌터 형식적 준스킴의
  닫힌 몰입 사상이면, $g\circ f$는 닫힌 몰입 사상이다.

\oldpage[I]{211}
  \item[\rm{(ii)}] $\mathfrak{X}$, $\mathfrak{Y}$,
  $\mathfrak{S}$를 세 국소 뇌터 형식적 준스킴이라 하고,
  $f:\mathfrak{X}\to\mathfrak{S}$를 닫힌 몰입 사상,
  $g:\mathfrak{Y}\to\mathfrak{S}$를 사상이라 하자. 그러면 사상
  $\mathfrak{X}\times_{\mathfrak{S}}\mathfrak{Y}\to\mathfrak{Y}$는
  닫힌 몰입 사상이다.
  \item[\rm{(iii)}] $\mathfrak{S}$를 국소 뇌터 형식적 준스킴,
  $\mathfrak{X}'$, $\mathfrak{Y}'$을 두 국소 뇌터
  $\mathfrak{S}$-형식적 준스킴이라 하되,
  $\mathfrak{X}'\times_{\mathfrak{S}}\mathfrak{Y}'$은 국소 뇌터라고
  하자. $\mathfrak{X}$, $\mathfrak{Y}$가 두 국소 뇌터
  $\mathfrak{S}$-형식적 준스킴이고,
  $f:\mathfrak{X}\to\mathfrak{X}'$, $g:\mathfrak{Y}\to\mathfrak{Y}'$이
  닫힌 몰입 사상인 두 $\mathfrak{S}$-사상이면,
  $f\times_{\mathfrak{S}}g$는 닫힌 몰입 사상이다.
\end{enumerate}
\end{proposition}

(\hyperref[I.3.5.1-ko]{3.5.1})에 따라 여기서도 (i)과 (ii)만 증명하면
충분하다.

(i)을 증명하기 위해 $\mathfrak{Y}$ (각각 $\mathfrak{Z}$)가
$\mathfrak{X}$ (각각 $\mathfrak{Y}$)의 닫힌 부분준스킴이고, 이들이
각각 연접 아이디얼층 $\mathscr{J}$ (각각 $\mathscr{K}$)로 정의된다고
가정할 수 있다. 이 두 층은 각각 $\mathscr{O}_{\mathfrak{X}}$ (각각
$\mathscr{O}_{\mathfrak{Y}}$)의 아이디얼층이다~; $\psi$가 바탕공간의 단사사상
$\mathfrak{Y}\to\mathfrak{X}$이면, $\psi_*(\mathscr{K})$는
$\psi_*(\mathscr{O}_{\mathfrak{Y}})
=\mathscr{O}_{\mathfrak{X}}/\mathscr{J}$의 연접 아이디얼층이다
(\hyperref[0.5.3.12-ko]{0, 5.3.12}). 따라서 이는 또한 연접
$\mathscr{O}_{\mathfrak{X}}$-가군이다
(\hyperref[0.5.3.10-ko]{0, 5.3.10})~; 핵 $\mathscr{K}_1$, 즉 준동형
$\mathscr{O}_{\mathfrak{X}}\to
(\mathscr{O}_{\mathfrak{X}}/\mathscr{J})/\psi_*(\mathscr{K})$의 핵은
$\mathscr{O}_{\mathfrak{X}}$의 연접 아이디얼층이고
(\hyperref[0.5.3.4-ko]{0, 5.3.4}),
$\mathscr{O}_{\mathfrak{X}}/\mathscr{K}_1$은
$\psi_*(\mathscr{O}_{\mathfrak{Y}}/\mathscr{K})$와 동형이다. 따라서
$\mathfrak{Z}$는 $\mathfrak{X}$의 닫힌 부분준스킴과 동형이다.

(ii)를 증명하기 위해 $\mathfrak{S}=\operatorname{Spf}(A)$,
$\mathfrak{X}=\operatorname{Spf}(B)$,
$\mathfrak{Y}=\operatorname{Spf}(C)$인 경우로 한정할 수 있음은
명백하다. 여기서 $A$는 뇌터 $\mathfrak{J}$-아딕환이고,
$B=A/\mathfrak{a}$이며 $\mathfrak{a}$는 $A$의 아이디얼이고,
$C$는 아딕이고 뇌터인 위상 $A$-대수이다. 이제 준동형
$C\to C\widehat{\otimes}_A(A/\mathfrak{a})$이 \emph{전사}임을
보이면 충분하다~: 실제로 $A/\mathfrak{a}$는 유한 생성 $A$-가군이고
그 위상은 $\mathfrak{J}$-아딕 위상이다~; 따라서
(\hyperref[0.7.7.8-ko]{0, 7.7.8})에 의해
$C\widehat{\otimes}_A(A/\mathfrak{a})$는
$C\otimes_A(A/\mathfrak{a})=C/\mathfrak{a}C$와 동일시되므로 주장이
따른다.

\begin{corollary}[10.14.6]
\label{I.10.14.6-ko}
(\hyperref[I.10.14.5-ko]{10.14.5, (ii)})의 가정 아래에서
$p:\mathfrak{X}\times_{\mathfrak{S}}\mathfrak{Y}\to\mathfrak{X}$,
$q:\mathfrak{X}\times_{\mathfrak{S}}\mathfrak{Y}\to\mathfrak{Y}$를
표준 사영들이라 하자. 그러면 도식
\[
  \xymatrix{
    \mathfrak{X} \ar[d]_{f}
    & \mathfrak{X}\times_{\mathfrak{S}}\mathfrak{Y}
      \ar[l]_{p} \ar[d]^{q}\\
    \mathfrak{S}
    & \mathfrak{Y} \ar[l]^{g}
  }
\]
은 가환이다. 모든 연접 $\mathscr{O}_{\mathfrak{X}}$-가군
$\mathscr{F}$에 대하여 다음 표준 $\mathscr{O}_{\mathfrak{Y}}$-가군
동형사상이 존재한다~:
\[
  \label{I.10.14.6.1-ko}
  u:g^*(f_*(\mathscr{F}))
  \xrightarrow{\sim}
  q_*(p^*(\mathscr{F})).
  \tag{10.14.6.1}
\]
\end{corollary}

준동형 $g^*(f_*(\mathscr{F}))\to q_*(p^*(\mathscr{F}))$을 정의하는
것은 준동형
$f_*(\mathscr{F})\to g_*(q_*(p^*(\mathscr{F})))
=f_*(p_*(p^*(\mathscr{F})))$을 정의하는 것과 같다
(\hyperref[0.4.4.3-ko]{0, 4.4.3})~: 앞의 수반 대응 아래에서 원문의
표기를 따라 $u=f_*(\rho)$로 쓰겠다. 여기서 $\rho$는 표준 준동형
$\mathscr{F}\to p_*(p^*(\mathscr{F}))$이다
(\hyperref[0.4.4.3-ko]{0, 4.4.3}).\footnote{\textbf{번역 주.} 이
등식은 서로 다른 준동형 집합의 원소들을 수반 전단사로 대응시켜 쓴
표기이며, 두 사상이 문자 그대로 같은 형을 갖는다는 뜻이 아니다.}
$u$가 동형사상임을 보이기 위해 $\mathfrak{S}$, $\mathfrak{X}$,
$\mathfrak{Y}$가 각각 뇌터 아딕환 $A$, $B$, $C$의 형식적
스펙트럼이고, (\hyperref[I.10.14.5-ko]{10.14.5, (ii)})에서 본
조건들이 성립하는 경우로 곧바로 환원된다~; 이때
$\mathscr{F}=M^\Delta$이고, $M$은 유한 생성
$(A/\mathfrak{a})$-가군이다
(\hyperref[I.10.10.5-ko]{10.10.5}).
(\hyperref[I.10.14.6.1-ko]{10.14.6.1})의 양변은
(\hyperref[I.10.10.8-ko]{10.10.8})에 따라 각각
$(C\otimes_A M)^\Delta$와
$((C/\mathfrak{a}C)\otimes_{A/\mathfrak{a}}M)^\Delta$로 동일시된다.
그런데
$(C/\mathfrak{a}C)\otimes_{A/\mathfrak{a}}M
=(C\otimes_A(A/\mathfrak{a}))\otimes_{A/\mathfrak{a}}M$은
$C\otimes_A M$과 표준적으로 동일시되므로 따름정리가 증명된다.

\oldpage[I]{212}
\begin{corollary}[10.14.7]
\label{I.10.14.7-ko}
$X$를 (보통 의미의) 국소 뇌터 준스킴, $Y$를 $X$의 닫힌 부분준스킴,
$j$를 표준 포함 사상 $Y\to X$, $X'$을 $X$의 닫힌 부분집합이라 하고
$Y'=Y\cap X'$이라 하자~; 그러면
$\widehat{j}:Y_{/Y'}\to X_{/X'}$은 닫힌 몰입 사상이고, 모든 연접
$\mathscr{O}_Y$-가군 $\mathscr{F}$에 대하여
\[
  \widehat{j}_*(\mathscr{F}_{/Y'})=(j_*(\mathscr{F}))_{/X'}.
\]
\end{corollary}

$Y'=j^{-1}(X')$이므로 (\hyperref[I.10.9.9-ko]{10.9.9})를 사용하고
(\hyperref[I.10.14.5-ko]{10.14.5})와
(\hyperref[I.10.14.6-ko]{10.14.6})을 적용하면 충분하다.
\subsection{분리된 형식적 준스킴.}
\label{subsection:I.10.15-ko}

\begin{definition}[10.15.1]
\label{I.10.15.1-ko}
형식적 준스킴 $\mathfrak{S}$, $\mathfrak{S}$-형식적 준스킴
$\mathfrak{X}$, 구조 사상 $f:\mathfrak{X}\to\mathfrak{S}$가 주어졌다고
하자. 대각 사상
$\Delta_{\mathfrak{X}|\mathfrak{S}}:\mathfrak{X}\to
\mathfrak{X}\times_{\mathfrak{S}}\mathfrak{X}$ (또한
$\Delta_{\mathfrak{X}}$라고도 쓴다)을 사상
$(1_{\mathfrak{X}},1_{\mathfrak{X}})_{\mathfrak{S}}$라 한다. 대각 사상
$\Delta_{\mathfrak{X}}$에 의한 바탕공간 $\mathfrak{X}$의 상이
$\mathfrak{X}\times_{\mathfrak{S}}\mathfrak{X}$의 바탕공간의 닫힌 부분집합이면,
$\mathfrak{X}$가 $\mathfrak{S}$ 위에서 분리되어 있다고, 또는
$\mathfrak{S}$-형식적 스킴이라고, 또는 $f$가 분리 사상이라고 한다. 형식적
준스킴 $\mathfrak{X}$가 분리되어 있다고, 또는 형식적 스킴이라고 하는 것은
$\mathbf{Z}$ 위에서 분리되어 있다는 뜻이다.
\end{definition}

\begin{proposition}[10.15.2]
\label{I.10.15.2-ko}
형식적 준스킴 $\mathfrak{S}$, $\mathfrak{X}$가 각각 보통 준스킴들의 열
$(S_n)$, $(X_n)$의 귀납극한이고, 사상
$f:\mathfrak{X}\to\mathfrak{S}$가 사상들의 열
$f_n:X_n\to S_n$의 귀납극한이라고 하자. $f$가 분리 사상이기 위한
필요충분조건은 $f_0:X_0\to S_0$가 분리 사상인 것이다.
\end{proposition}

실제로, $\Delta_{\mathfrak{X}|\mathfrak{S}}$는 사상들의 열
$\Delta_{X_n|S_n}$의 귀납극한이다
(\hyperref[I.10.7.4-ko]{10.7.4}). 또한 바탕공간 $\mathfrak{X}$의
$\Delta_{\mathfrak{X}|\mathfrak{S}}$에 의한 상(각각 바탕공간
$\mathfrak{X}\times_{\mathfrak{S}}\mathfrak{X}$)은 바탕공간 $X_0$의
$\Delta_{X_0|S_0}$에 의한 상(각각 바탕공간
$X_0\times_{S_0}X_0$)과 동일하다~; 따라서 결론이 따른다.

\begin{proposition}[10.15.3]
\label{I.10.15.3-ko}
이하에서는 고려하는 모든 형식적 준스킴(각각 형식적 준스킴의 모든 사상)이
보통 준스킴들의 열(각각 보통 준스킴의 사상들의 열)의 귀납극한이라고
가정한다.
\begin{enumerate}
  \item[\rm{(i)}] 두 분리 사상의 합성은 분리 사상이다.
  \item[\rm{(ii)}] $f:\mathfrak{X}\to\mathfrak{X}'$,
  $g:\mathfrak{Y}\to\mathfrak{Y}'$가 두 분리
  $\mathfrak{S}$-사상이면 $f\times_{\mathfrak{S}}g$도 분리 사상이다.
  \item[\rm{(iii)}] $f:\mathfrak{X}\to\mathfrak{Y}$가 분리
  $\mathfrak{S}$-사상이면, 밑 형식적 준스킴의 임의의 확장
  $\mathfrak{S}'\to\mathfrak{S}$에 대하여 $\mathfrak{S}'$-사상
  $f_{(\mathfrak{S}')}$도 분리 사상이다.
  \item[\rm{(iv)}] 두 사상의 합성 $g\circ f$가 분리 사상이면 $f$도 분리
  사상이다.
\end{enumerate}
\emph{(이 명제에서 동일한 형식적 준스킴 $\mathfrak{Z}$가 하나의 명제 안에
여러 번 나타나는 경우에는, 나타나는 모든 곳에서 그것을 보통 준스킴들의
동일한 열 $(Z_n)$의 귀납극한으로 간주한다는 것을 함의한다. 또한
형식적 준스킴 $\mathfrak{Z}$에서 형식적 준스킴으로 가는 사상(각각 형식적
준스킴에서 $\mathfrak{Z}$로 가는 사상)은 보통 준스킴의 사상들
 $Z_n$에서 보통 준스킴으로 가는 사상(각각 보통 준스킴에서 $Z_n$으로
가는 사상)의 귀납극한으로 간주한다.)}
\end{proposition}

실제로 (\hyperref[I.10.15.2-ko]{10.15.2})의 표기를 사용하면
$(g\circ f)_0=g_0\circ f_0$이고,
$(f\times_{\mathfrak{S}}g)_0=f_0\times_{S_0}g_0$이다~; 따라서
(\hyperref[I.10.15.3-ko]{10.15.3})의 주장들은 보통 준스킴에 대한
(\hyperref[I.10.15.2-ko]{10.15.2})와
(\hyperref[I.5.5.1-ko]{5.5.1})의 대응하는 주장들의 즉각적인
귀결이다.

독자에게는 (\hyperref[I.10.15.3-ko]{10.15.3})과 같은 종류의 형식적
준스킴 및 사상에 대하여
(\hyperref[I.5.5.5-ko]{5.5.5}),

\oldpage[I]{213}
(\hyperref[I.5.5.9-ko]{5.5.9}) 및
(\hyperref[I.5.5.10-ko]{5.5.10})에 대응하는 명제들을 같은 방식으로
서술하는 일을 맡긴다(여기서 «~아핀 열린집합~»을
«~(\hyperref[I.10.6.3-ko]{10.6.3})의 조건 b)를 만족하는 아핀 형식적
열린집합~»으로 바꾼다).

유사한 논증으로 모든 뇌터 아핀 형식적 스킴이 분리되어 있음도 보인다.
따라서 이 용어가 정당화된다.

\begin{proposition}[10.15.4]
\label{I.10.15.4-ko}
형식적 준스킴 $\mathfrak{S}$를 국소 뇌터라고 하고,
$\mathfrak{X}$, $\mathfrak{Y}$를 두 국소 뇌터
$\mathfrak{S}$-형식적 준스킴이라 하자. $\mathfrak{X}$ 또는
$\mathfrak{Y}$가 $\mathfrak{S}$ 위에서 유한형이라고 가정하자(따라서
$\mathfrak{X}\times_{\mathfrak{S}}\mathfrak{Y}$는 국소 뇌터이다).
또 $\mathfrak{Y}$가 $\mathfrak{S}$ 위에서 분리되어 있다고 하자. 그러면
$\mathfrak{S}$-사상 $f:\mathfrak{X}\to\mathfrak{Y}$의 그래프 사상
$\Gamma_f=(1_{\mathfrak{X}},f)_{\mathfrak{S}}:\mathfrak{X}\to
\mathfrak{X}\times_{\mathfrak{S}}\mathfrak{Y}$는 닫힌 몰입 사상이다.
\end{proposition}

형식적 준스킴 $\mathfrak{S}$가 국소 뇌터 준스킴들의 열
$(S_n)$의 귀납극한이고, $\mathfrak{X}$(각각 $\mathfrak{Y}$)가 열
$(X_n)$(각각 $(Y_n)$)의 $S_n$-준스킴들의 귀납극한이며, $f$가 $S_n$-사상들의 열
$f_n:X_n\to Y_n$의 귀납극한이라고 가정할 수 있다. 그러면
$\mathfrak{X}\times_{\mathfrak{S}}\mathfrak{Y}$는 열
$(X_n\times_{S_n}Y_n)$의 귀납극한이고, $\Gamma_f$는 열
$(\Gamma_{f_n})$의 귀납극한이다
(\hyperref[I.10.7.4-ko]{10.7.4}). 가정에 따라 $Y_0$는 $S_0$ 위에서
분리되어 있다(\hyperref[I.10.15.2-ko]{10.15.2}). 따라서 공간
$\Gamma_{f_0}(X_0)$는 $X_0\times_{S_0}Y_0$의 닫힌 부분공간이다.
$\mathfrak{X}\times_{\mathfrak{S}}\mathfrak{Y}$의 바탕공간(각각
$\Gamma_f(\mathfrak{X})$)과 $X_0\times_{S_0}Y_0$(각각
$\Gamma_{f_0}(X_0)$)의 바탕공간이 서로 같으므로, 이미
$\Gamma_f(\mathfrak{X})$가
$\mathfrak{X}\times_{\mathfrak{S}}\mathfrak{Y}$의 닫힌
부분공간임을 알 수 있다. 이제 $(U,V)$가 $\mathfrak{X}$의 뇌터 아핀
형식적 열린집합 $U$(각각 $\mathfrak{Y}$의 뇌터 아핀 형식적 열린집합
$V$)로 이루어진 순서쌍들 중 $f(U)\subset V$를 만족하는 것들을
취하면, 열린집합 $U\times_{\mathfrak{S}}V$들은
$\mathfrak{X}\times_{\mathfrak{S}}\mathfrak{Y}$ 안에서
$\Gamma_f(\mathfrak{X})$의 덮개를 이룬다. $f_U:U\to V$를 $f$의
$U$로의 제한이라 하면, $\Gamma_{f_U}:U\to
U\times_{\mathfrak{S}}V$는 $U$ 위에서 $\Gamma_f$를 제한한 것이다.
따라서 $\Gamma_{f_U}$가 닫힌 몰입 사상임을 보이면
(\hyperref[I.10.14.4-ko]{10.14.4})에 따라 $\Gamma_f$도 닫힌 몰입
사상이 되며, 문제는 다음 경우로 환원된다:
$\mathfrak{S}=\operatorname{Spf}(A)$,
$\mathfrak{X}=\operatorname{Spf}(B)$,
$\mathfrak{Y}=\operatorname{Spf}(C)$가 아핀이고
($A$, $B$, $C$는 뇌터 아딕환), $f$가 연속 $A$-준동형
$\varphi:C\to B$에 대응하는 경우이다. 이때 $\Gamma_f$는 유일한 연속
준동형
$\omega:B\widehat{\otimes}_A C\to B$에 대응한다. 표준 준동형
$B\to B\widehat{\otimes}_A C$ 및
$C\to B\widehat{\otimes}_A C$를 이 유일한 준동형과 각각 합성하면 항등사상과
$\varphi$를 각각 얻는다. 그런데 $\omega$가
\emph{전사}임은 명백하므로 주장이 따른다.

\begin{corollary}[10.15.5]
\label{I.10.15.5-ko}
형식적 준스킴 $\mathfrak{S}$를 국소 뇌터라고 하고,
$\mathfrak{X}$를 $\mathfrak{S}$ 위의 유한형 형식적 준스킴이라 하자.
$\mathfrak{X}$가 $\mathfrak{S}$ 위에서 분리되어 있기 위한 필요충분조건은
대각 사상
$\mathfrak{X}\to\mathfrak{X}\times_{\mathfrak{S}}\mathfrak{X}$가
닫힌 몰입 사상인 것이다.
\end{corollary}

\begin{proposition}[10.15.6]
\label{I.10.15.6-ko}
국소 뇌터 형식적 준스킴들의 닫힌 몰입 사상
$j:\mathfrak{Y}\to\mathfrak{X}$는 분리 사상이다.
\end{proposition}

(\hyperref[I.10.14.2-ko]{10.14.2})의 표기를 사용하면
$j_0:Y_0\to X_0$는 닫힌 몰입 사상이고, 따라서 분리 사상이다. 이제
(\hyperref[I.10.15.2-ko]{10.15.2})를 적용하면 충분하다.

\begin{proposition}[10.15.7]
\label{I.10.15.7-ko}
$X$를 (보통 의미의) 국소 뇌터 준스킴, $X'$을 $X$의 닫힌 부분집합,
$\widehat{X}=X_{/X'}$라 하자. $\widehat{X}$가 분리되기 위한
필요충분조건은 $\widehat{X}_{\mathrm{red}}$가 분리되는 것이며,
특히 $X$가 분리되어 있으면 이 조건이 충족된다.
\end{proposition}

실제로, (\hyperref[I.10.8.5-ko]{10.8.5})의 표기를 사용하면
$\widehat{X}$가 분리되기 위한 필요충분조건은 $X'_0$가 분리되는
것이다(\hyperref[I.10.15.2-ko]{10.15.2}). 그런데
$\widehat{X}_{\mathrm{red}}=(X'_0)_{\mathrm{red}}$이므로,
$\widehat{X}_{\mathrm{red}}$가 분리된다고 말하는 것과 동치이다
(\hyperref[I.5.5.1-ko]{5.5.1, (vi)}).
% ===== END EXACT INLINE INPUT: c1s10.tex =====


% ===== BEGIN EXACT INLINE INPUT: bib.tex =====
% Bibliographic data inherited verbatim from the current French authority.
\begin{thebibliography}{22}

\bibitem{I-1}
\oldpage[I]{214}
H.~\textsc{Cartan} et C.~\textsc{Chevalley},
\emph{Séminaire de l'École Normale Supérieure},
8\textsuperscript{e} année (1955-56), \emph{Géométrie algébrique}.

\bibitem{I-2}
H.~\textsc{Cartan} and S.~\textsc{Eilenberg},
\emph{Homological Algebra}, Princeton Math. Series (Princeton University
Press), 1956.

\bibitem{I-3}
W.~L.~\textsc{Chow} and J.~\textsc{Igusa},
\emph{Cohomology theory of varieties over rings},
\emph{Proc. Nat. Acad. Sci. U.S.A.}, t.~XLIV (1958), p.~1244-1248.

\bibitem{I-4}
R.~\textsc{Godement}, \emph{Théorie des faisceaux}, Actual. Scient. et Ind.,
n\textsuperscript{o}~1252, Paris (Hermann), 1958.

\bibitem{I-5}
H.~\textsc{Grauert},
\emph{Ein Theorem der analytischen Garbentheorie und die Modulräume komplexer
Strukturen}, \emph{Publ. Math. Inst. Hautes Études Scient.},
n\textsuperscript{o}~5, 1960.

\bibitem{I-6}
A.~\textsc{Grothendieck}, \emph{Sur quelques points d'algèbre homologique},
\emph{Tôhoku Math. Journ.}, t.~IX (1957), p.~119-221.

\bibitem{I-7}
A.~\textsc{Grothendieck},
\emph{Cohomology theory of abstract algebraic varieties},
\emph{Proc. Intern. Congress of Math.}, p.~103-118, Edinburgh (1958).

\bibitem{I-8}
A.~\textsc{Grothendieck},
\emph{Géométrie formelle et géométrie algébrique},
\emph{Séminaire Bourbaki}, 11\textsuperscript{e} année (1958-59), exposé 182.

\bibitem{I-9}
M.~\textsc{Nagata},
\emph{A general theory of algebraic geometry over Dedekind domains},
\emph{Amer. Math. Journ.}~: I, t.~LXXVIII, p.~78-116 (1956)~; II, t.~LXXX,
p.~382-420 (1958).

\bibitem{I-10}
D.~G.~\textsc{Northcott}, \emph{Ideal theory}, Cambridge Univ. Press, 1953.

\bibitem{I-11}
P.~\textsc{Samuel}, \emph{Commutative algebra} (Notes by D.~Herzig), Cornell
Univ., 1953.

\bibitem{I-12}
P.~\textsc{Samuel}, \emph{Algèbre locale}, Mém. Sci. Math.,
n\textsuperscript{o}~123, Paris, 1953.

\bibitem{I-13}
P.~\textsc{Samuel} and O.~\textsc{Zariski},
\emph{Commutative algebra}, 2 vol., New York (Van Nostrand), 1958-60.

\bibitem{I-14}
J.-P.~\textsc{Serre}, \emph{Faisceaux algébriques cohérents},
\emph{Ann. of Math.}, t.~LXI (1955), p.~197-278.

\bibitem{I-15}
J.-P.~\textsc{Serre},
\emph{Sur la cohomologie des variétés algébriques},
\emph{Journ. de Math.} (9), t.~XXXVI (1957), p.~1-16.

\bibitem{I-16}
J.-P.~\textsc{Serre},
\emph{Géométrie algébrique et géométrie analytique},
\emph{Ann. Inst. Fourier}, t.~VI (1955-56), p.~1-42.

\bibitem{I-17}
J.-P.~\textsc{Serre},
\emph{Sur la dimension homologique des anneaux et des modules noethériens},
\emph{Proc. Intern. Symp. on Alg. Number theory}, p.~176-189, Tokyo-Nikko,
1955.

\bibitem{I-18}
A.~\textsc{Weil}, \emph{Foundations of algebraic geometry}, Amer. Math. Soc.
Coll. Publ., n\textsuperscript{o}~29, 1946.

\bibitem{I-19}
A.~\textsc{Weil}, \emph{Numbers of solutions of equations in finite fields},
\emph{Bull. Amer. Math. Soc.}, t.~LV (1949), p.~497-508.

\bibitem{I-20}
O.~\textsc{Zariski},
\emph{Theory and applications of holomorphic functions on algebraic varieties
over arbitrary ground fields}, Mem. Amer. Math. Soc., n\textsuperscript{o}~5
(1951).

\bibitem{I-21}
O.~\textsc{Zariski}, \emph{A new proof of Hilbert's Nullstellensatz},
\emph{Bull. Amer. Math. Soc.}, t.~LIII (1947), p.~362-368.

\bibitem{I-22}
E.~\textsc{Kähler}, \emph{Geometria Arithmetica},
\emph{Ann. di Mat.} (4), t.~XLV (1958), p.~1-368.

\end{thebibliography}
% ===== END EXACT INLINE INPUT: bib.tex =====


% ===== BEGIN EXACT INLINE INPUT: c2front.tex =====
% Authority: EGA II front-fr.tex, SHA-256 F199DE0C2C8C49E2BDAEF48D764086952806E77D8066795B97B3595F7EA03198.
\clearpage
\thispagestyle{plain}
\markboth{제II장}{제II장}
\oldpage[II]{5}
\phantomsection
\label{chapter:II-ko}

\begin{center}
{\large 제II장\par}
\vspace{1.5em}
{\Large\bfseries 몇몇 사상 부류에 대한\par}
{\Large\bfseries 초보적 전역 연구\par}
\vspace{1.8em}
{\bfseries 목차\par}
\end{center}

\medskip
\begin{tabular}{@{}r@{\ }p{0.82\textwidth}@{}}
\S~1. & 아핀 사상.\\
\S~2. & 동차 소 스펙트럼.\\
\S~3. & 등급 대수의 층의 동차 스펙트럼.\\
\S~4. & 사영 다발. 풍부한 층.\\
\S~5. & 준아핀 사상; 준사영 사상; 고유 사상; 사영 사상.\\
\S~6. & 정수적 사상과 유한 사상.\\
\S~7. & 값매김 판정법.\\
\S~8. & 블로업 스킴; 사영 원뿔; 사영 폐포.
\end{tabular}

\medskip
이 장에서 다루는 여러 사상 부류는 코호몰로지적 방법을 크게 사용하지
않고 연구한다. 이 방법을 사용하여 그 사상 부류들을 더 깊이 연구하는
일은 제III장에서 이루어질 것이며, 거기서는 제II장의 \S\S~2, 4, 5를
주로 사용할 것이다. \S~8은 처음 읽을 때 생략해도 된다. 이 절은
\S\S~1에서~3까지 전개한 형식 체계에 몇 가지 보충을 더하는데, 그
내용은 이 형식 체계의 쉬운 응용으로 귀착된다. 우리는 이 절의 결과들을
이 장의 다른 결과들만큼 자주 사용하지는 않을 것이다.
% ===== END EXACT INLINE INPUT: c2front.tex =====


% ===== BEGIN EXACT INLINE INPUT: c2s1.tex =====
% Authority slice: EGA II ega2-1-fr.tex lines 1-18087, 820503 LF UTF-8 bytes, SHA-256 AEF273683A786E968FA47C92CD7C14E5EE2D85056ABC451CD4398BA05ECD0B74.
\setcounter{section}{0}
\section{아핀 사상}
\label{section:II.1-ko}

이 절의 결과 대부분은 제I장 \S~1의 ``전역적'' 대응물이다. 따라서
본질적으로 새로운 것은 아니며, 이후를 위한 편리한 언어를 제공할 뿐이다.

\setcounter{subsection}{0}
\subsection{$S$-준스킴과 $\mathcal{O}_S$-대수}
\label{subsection:II.1.1-ko}

\begin{env}[1.1.1]
\label{II.1.1.1-ko}
$S$를 준스킴, $X$를 $S$-준스킴, $f:X\to S$를 그 구조 사상이라 하자.
(\hyperref[0.4.2.4-ko]{0, 4.2.4})에 따라 직상
$f_*(\mathcal{O}_X)$는 $\mathcal{O}_S$-대수이다. 혼동의 우려가 없으면
이를
\oldpage[II]{6}
$\mathcal{A}(X)$로 나타내기로 한다. $U$가 $S$의 열린집합이면
\[
  \mathcal{A}(f^{-1}(U))=\mathcal{A}(X)|U.
\]
마찬가지로 임의의 $\mathcal{O}_X$-가군 $\mathcal{F}$(각각
$\mathcal{O}_X$-대수 $\mathcal{B}$)에 대하여, 직상
$f_*(\mathcal{F})$(각각 $f_*(\mathcal{B})$)를
$\mathcal{A}(\mathcal{F})$(각각 $\mathcal{A}(\mathcal{B})$)로
나타내기로 한다. 이는 단순히 $\mathcal{O}_S$-가군(각각
$\mathcal{O}_S$-대수)일 뿐만 아니라 $\mathcal{A}(X)$-가군(각각
$\mathcal{A}(X)$-대수)이다.
\end{env}

\begin{env}[1.1.2]
\label{II.1.1.2-ko}
$Y$를 또 하나의 $S$-준스킴, $g:Y\to S$를 그 구조 사상,
$h:X\to Y$를 $S$-사상이라 하자. 그러면 다음 도식은 가환한다.
\[
  \xymatrix{
    X\ar[rr]^h\ar[rd]_f && Y\ar[ld]^g\\
    & S
  }
  \tag{1.1.2.1}\label{II.1.1.2.1-ko}
\]

정의상 $h=(\psi,\theta)$이고, 여기서
$\theta:\mathcal{O}_Y\to h_*(\mathcal{O}_X)=\psi_*(\mathcal{O}_X)$는
환의 층의 준동형이다. 이로부터
(\hyperref[0.4.2.2-ko]{0, 4.2.2})에 따라 다음
$\mathcal{O}_S$-대수 준동형을 얻는다.
\[
  g_*(\theta):g_*(\mathcal{O}_Y)\to
  g_*(h_*(\mathcal{O}_X))=f_*(\mathcal{O}_X),
\]
다시 말해 $\mathcal{O}_S$-대수 준동형
$\mathcal{A}(Y)\to\mathcal{A}(X)$를 얻으며, 이를
$\mathcal{A}(h)$로도 나타낸다. $h':Y\to Z$가 또 하나의
$S$-사상이면
$\mathcal{A}(h'\circ h)=\mathcal{A}(h)\circ\mathcal{A}(h')$임은
곧바로 알 수 있다. 따라서 $S$-준스킴의 범주 위에 정의되어
$\mathcal{O}_S$-대수의 범주에 값을 갖는 \emph{반변 함자}
$\mathcal{A}(X)$를 정의하였다.

이제 $\mathcal{F}$를 $\mathcal{O}_X$-가군, $\mathcal{G}$를
$\mathcal{O}_Y$-가군이라 하고, $u:\mathcal{G}\to\mathcal{F}$를
$h$-사상이라 하자. 다시 말해
(\hyperref[0.4.4.1-ko]{0, 4.4.1})에 따라 $u$는
$\mathcal{O}_Y$-가군 준동형
$\mathcal{G}\to h_*(\mathcal{F})$이다. 그러면
\[
  g_*(u):g_*(\mathcal{G})\to g_*(h_*(\mathcal{F}))=f_*(\mathcal{F})
\]
는 $\mathcal{O}_S$-가군 준동형
$\mathcal{A}(\mathcal{G})\to\mathcal{A}(\mathcal{F})$이며, 이를
$\mathcal{A}(u)$로 나타낸다. 더욱이 쌍
$(\mathcal{A}(h),\mathcal{A}(u))$는
$\mathcal{A}(Y)$-가군 $\mathcal{A}(\mathcal{G})$에서
$\mathcal{A}(X)$-가군 $\mathcal{A}(\mathcal{F})$로 가는
\emph{디-준동형}을 이룬다.
\end{env}

\begin{env}[1.1.3]
\label{II.1.1.3-ko}
$S$를 고정하면, $X$가 $S$-준스킴이고 $\mathcal{F}$가
$\mathcal{O}_X$-가군인 쌍 $(X,\mathcal{F})$들이 하나의
\emph{범주}를 이룬다고 볼 수 있다. 여기서
$(X,\mathcal{F})$에서 $(Y,\mathcal{G})$로 가는 \emph{사상}은 쌍
$(h,u)$로 정의한다. 이때 $h:X\to Y$는 $S$-사상이고
$u:\mathcal{G}\to\mathcal{F}$는 $h$-사상이다. 그러면
$(\mathcal{A}(X),\mathcal{A}(\mathcal{F}))$는
$\mathcal{O}_S$-대수와 그 대수 위의 가군으로 이루어진 쌍들을
대상으로 하고 디-준동형들을 사상으로 하는 범주에 값을 갖는
\emph{반변 함자}라고 할 수 있다.
\end{env}

\subsection{준스킴 위에서 아핀인 준스킴}
\label{subsection:II.1.2-ko}

\begin{env}[1.2.1]
\label{II.1.2.1-ko}
$X$를 $S$-준스킴, $f:X\to S$를 그 구조 사상이라 하자.
$S$의 아핀 열린집합들로 이루어진 덮개 $(S_\alpha)$가 존재하여,
모든 $\alpha$에 대하여 열린집합 $f^{-1}(S_\alpha)$ 위에 $X$가
유도하는 준스킴이 아핀이면, $X$가 \emph{$S$ 위에서 아핀}이라고
한다.
\end{env}

\begin{env}[1.2.2]
\label{II.1.2.2-ko}
$S$의 모든 닫힌 부분준스킴은 $S$ 위에서 아핀인
$S$-준스킴이다
(\hyperref[I.4.2.3-ko]{I, 4.2.3}과
\hyperref[I.4.2.4-ko]{4.2.4}).
\end{env}

\begin{env}[1.2.3]
\label{II.1.2.3-ko}
$S$ 위에서 아핀인 준스킴 $X$가 반드시 아핀 스킴인 것은 아니다. 예
$X=S$가 이를 보여 준다(\hyperref[II.1.2.2-ko]{1.2.2}). 한편 아핀 스킴
$X$가 $S$-준스킴이라 해도, $X$가 반드시 $S$ 위에서 아핀인
\oldpage[II]{7}
것은 아니다(\hyperref[II.1.3.3-ko]{1.3.3} 참조). 그러나 $S$가
\emph{스킴}이면, 아핀 스킴인 모든 $S$-준스킴은 $S$ 위에서 아핀임을
상기하자(\hyperref[I.5.5.10-ko]{I, 5.5.10}).
\end{env}

\begin{env}[1.2.4]
\label{II.1.2.4-ko}
$S$ 위에서 아핀인 모든 $S$-준스킴은 $S$ 위에서 분리되어 있다
(다시 말해 $S$-스킴이다).
\end{env}

이는 (\hyperref[I.5.5.5-ko]{I, 5.5.5})와
(\hyperref[I.5.5.8-ko]{I, 5.5.8})에서 곧바로 따른다.

\begin{env}[1.2.5]
\label{II.1.2.5-ko}
$X$를 $S$ 위에서 아핀인 $S$-스킴, $f:X\to S$를 그 구조 사상이라
하자. 모든 열린집합 $U\subset S$에 대하여, $f^{-1}(U)$는 $U$ 위에서
아핀이다.
\end{env}

정의 (\hyperref[II.1.2.1-ko]{1.2.1})에 따라,
$S=\operatorname{Spec}(A)$와 $X=\operatorname{Spec}(B)$가 아핀인
경우로 환원된다. 이때 $f=({}^a\varphi,\widetilde{\varphi})$이고,
여기서 $\varphi:A\to B$는 환 준동형이다. $g\in A$인 $D(g)$들이
$S$의 기저를 이루므로, $U=D(g)$인 경우로 환원된다. 그런데 이때
$f^{-1}(U)=D(\varphi(g))$임을 알고 있으므로
(\hyperref[I.1.2.2.2-ko]{I, 1.2.2.2}), 명제가 따른다.

\begin{env}[1.2.6]
\label{II.1.2.6-ko}
$X$를 $S$ 위에서 아핀인 $S$-스킴, $f:X\to S$를 그 구조 사상이라
하자. 임의의 준연접 $\mathcal{O}_X$-가군 $\mathcal{F}$에 대하여,
$f_*(\mathcal{F})$는 준연접 $\mathcal{O}_S$-가군이다.
\end{env}

(\hyperref[II.1.2.4-ko]{1.2.4})를 고려하면, 이는
(\hyperref[I.9.2.2-ko]{I, 9.2.2, a)})에서 따른다.

특히 $\mathcal{O}_S$-대수
$\mathcal{A}(X)=f_*(\mathcal{O}_X)$는 \emph{준연접}이다.

\begin{env}[1.2.7]
\label{II.1.2.7-ko}
$X$를 $S$ 위에서 아핀인 $S$-스킴이라 하자. 임의의 $S$-준스킴
$Y$에 대하여, 집합 $\operatorname{Hom}_S(Y,X)$에서 집합
$\operatorname{Hom}(\mathcal{A}(X),\mathcal{A}(Y))$로 가는 사상
$h\mapsto\mathcal{A}(h)$
(\hyperref[II.1.1.2-ko]{1.1.2})은 전단사이다.
\end{env}

$f:X\to S$와 $g:Y\to S$를 구조 사상이라 하자. 먼저
$S=\operatorname{Spec}(A)$와 $X=\operatorname{Spec}(B)$가
아핀이라고 하자. 임의의 $\mathcal{O}_S$-대수 준동형
$\omega:f_*(\mathcal{O}_X)\to g_*(\mathcal{O}_Y)$에 대하여,
$\mathcal{A}(h)=\omega$인 $S$-사상 $h:Y\to X$가 유일하게 존재함을
보여야 한다. 정의에 따라, 모든 열린집합 $U\subset S$에 대하여
$\omega$는 다음 $\Gamma(U,\mathcal{O}_S)$-대수 준동형을 정한다.
\[
  \omega_U=\Gamma(U,\omega):
  \Gamma(f^{-1}(U),\mathcal{O}_X)\to
  \Gamma(g^{-1}(U),\mathcal{O}_Y)
\]
특히 $U=S$일 때, 이는 $\Gamma(S,\mathcal{O}_S)$-대수 준동형
$\varphi:\Gamma(X,\mathcal{O}_X)\to\Gamma(Y,\mathcal{O}_Y)$를
준다. $X$가 아핀이므로, 이 준동형에는 유일하게 정해지는
$S$-사상 $h:Y\to X$가 대응한다
(\hyperref[I.2.2.4-ko]{I, 2.2.4}). 이제
$\mathcal{A}(h)=\omega$임을 보여야 한다. 다시 말해, $S$의 한 기저에
속하는 모든 열린집합 $U$에 대하여 $\omega_U$가 $h$의 제한인
$S$-사상 $g^{-1}(U)\to f^{-1}(U)$에 대응하는 대수 준동형
$\varphi_U$와 일치함을 보이면 된다. $\lambda\in A$인
$U=D(\lambda)$의 경우만 살피면 충분하다.\footnote{\textbf{번역 주.}
인쇄 원문은 $\lambda\in S$라고 쓰지만, $S=\operatorname{Spec}(A)$이고
$\rho:A\to B$이며 뒤에서 $\rho(\lambda)$를 사용하므로
EGA2-CANON-DISP-0001에 따라 $\lambda\in A$로 교정했다.} 이때
$f=({}^a\rho,\widetilde{\rho})$이고 $\rho:A\to B$가 환 준동형이면,
$f^{-1}(U)=D(\mu)$이며 여기서 $\mu=\rho(\lambda)$이다. 또한
$\Gamma(f^{-1}(U),\mathcal{O}_X)$는 분수환 $B_\mu$이다. 그런데 다음
도식은 가환한다.
\[
  \xymatrix{
    B\ar[r]^\varphi\ar[d] & \Gamma(Y,\mathcal{O}_Y)\ar[d]\\
    B_\mu\ar[r]_{\varphi_U} & \Gamma(g^{-1}(U),\mathcal{O}_Y)
  }
\]
$\varphi_U$를 $\omega_U$로 바꾼 유사한 도식도 가환한다. 따라서
분수환의 보편 성질
(\hyperref[0.1.2.4-ko]{0, 1.2.4})에 의해
$\varphi_U=\omega_U$이다.

이제 일반적인 경우를 보자. $(S_\alpha)$를 $S$의 아핀 열린집합들로
이루어진 덮개라 하자.
\oldpage[II]{8}
여기서 각 $f^{-1}(S_\alpha)$는 아핀이라고 하자. 그러면 임의의
$\mathcal{O}_S$-대수 준동형
$\omega:\mathcal{A}(X)\to\mathcal{A}(Y)$는 제한에 의해 다음
$\mathcal{O}_{S_\alpha}$-대수 준동형들의 족을 정한다.
\[
  \omega_\alpha:
  \mathcal{A}(f^{-1}(S_\alpha))\to
  \mathcal{A}(g^{-1}(S_\alpha))
\]
따라서 앞에서 보인 바에 의해 $S_\alpha$-사상들의 족
$h_\alpha:g^{-1}(S_\alpha)\to f^{-1}(S_\alpha)$가 정해진다. 이제
$S_\alpha\cap S_\beta$의 한 기저에 속하는 모든 아핀 열린집합 $U$에
대하여, $h_\alpha$와 $h_\beta$를 $g^{-1}(U)$에 제한한 것들이
일치함을 보이면 된다. 이는 명백하다. 실제로 앞에서 보인 바에 의해
이 두 제한은 모두 $\omega$의 제한인 준동형
$\mathcal{A}(X)|U\to\mathcal{A}(Y)|U$에 대응한다.

\begin{env}[1.2.8]
\label{II.1.2.8-ko}
$X$와 $Y$를 $S$ 위에서 아핀인 두 $S$-스킴이라 하자. $S$-사상
$h:Y\to X$가 동형사상이기 위한 필요충분조건은
$\mathcal{A}(h):\mathcal{A}(X)\to\mathcal{A}(Y)$가 동형사상인 것이다.
\end{env}

이는 (\hyperref[II.1.2.7-ko]{1.2.7})과 $\mathcal{A}(X)$의 함자성에서
곧바로 따른다.

\subsection{$\mathcal{O}_S$-대수에 대응하는 $S$ 위에서 아핀인 준스킴}
\label{subsection:II.1.3-ko}

\begin{env}[1.3.1]
\label{II.1.3.1-ko}
$S$를 준스킴이라 하자. 임의의 준연접 $\mathcal{O}_S$-대수
$\mathcal{B}$에 대하여, $\mathcal{A}(X)=\mathcal{B}$를 만족하는
$S$ 위에서 아핀인 준스킴 $X$가 존재하며, $X$는 유일한
$S$-동형까지 유일하게 결정된다.
\end{env}

유일성은 (\hyperref[II.1.2.8-ko]{1.2.8})에서 곧바로 따른다.
$X$의 존재를 증명하자. 임의의 아핀 열린집합 $U\subset S$에 대하여
$X_U$를 준스킴
$\operatorname{Spec}(\Gamma(U,\mathcal{B}))$라 하자.
$\Gamma(U,\mathcal{B})$는 $\Gamma(U,\mathcal{O}_S)$-대수이므로
$X_U$는 $S$-준스킴이다
(\hyperref[I.1.6.1-ko]{I, 1.6.1}). 또한 $\mathcal{B}$가
준연접이므로 $\mathcal{O}_S$-대수 $\mathcal{A}(X_U)$는
$\mathcal{B}|U$와 표준적으로 동일시된다
(\hyperref[I.1.3.7-ko]{I, 1.3.7},
\hyperref[I.1.3.13-ko]{1.3.13}과
\hyperref[I.1.6.3-ko]{1.6.3}).

$V$를 $S$의 두 번째 아핀 열린집합이라 하고, $f_U$가 구조 사상
$X_U\to S$를 나타낼 때 $X_{U,V}$를 $X_U$가
$f_U^{-1}(U\cap V)$ 위에 유도하는 준스킴이라 하자.
$X_{U,V}$와 $X_{V,U}$는 $U\cap V$ 위에서 아핀이고
(\hyperref[II.1.2.5-ko]{1.2.5}), 정의에 따라
$\mathcal{A}(X_{U,V})$와 $\mathcal{A}(X_{V,U})$는
$\mathcal{B}|(U\cap V)$와 표준적으로 동일시된다. 그러므로
(\hyperref[II.1.2.8-ko]{1.2.8})에 따라 표준 $S$-동형사상
$\theta_{U,V}:X_{V,U}\to X_{U,V}$가 존재한다. 더욱이 $W$를
$S$의 세 번째 아핀 열린집합이라 하고, 구조 사상들에 의해
$X_V$, $X_W$, $X_W$ 안에서 각각 얻는 $U\cap V\cap W$의
역상들로 $\theta_{U,V}$, $\theta_{V,W}$, $\theta_{U,W}$를
제한한 사상을 차례로 $\theta'_{U,V}$, $\theta'_{V,W}$,
$\theta'_{U,W}$라 하면
$\theta'_{U,V}\circ\theta'_{V,W}=\theta'_{U,W}$이다.

따라서 준스킴 $X$, 아핀 열린집합들로 이루어진 $X$의 덮개
$(T_U)$, 그리고 각 $U$에 대한 동형사상
$\varphi_U:X_U\to T_U$가 존재하여, $\varphi_U$는
$f_U^{-1}(U\cap V)$를 $T_U\cap T_V$ 위로 보내고
$\theta_{U,V}=\varphi_U^{-1}\circ\varphi_V$가 성립한다
(\hyperref[I.2.3.1-ko]{I, 2.3.1}). 사상
$g_U=f_U\circ\varphi_U^{-1}$는 $T_U$에 $S$-준스킴 구조를 주며,
$g_U$와 $g_V$는 $T_U\cap T_V$에서 일치한다. 따라서 $X$는
$S$-준스킴이다. 또한 정의상 $X$가 $S$ 위에서 아핀이고
$\mathcal{A}(T_U)=\mathcal{B}|U$임은 명백하므로
$\mathcal{A}(X)=\mathcal{B}$이다.

이렇게 정의한 $S$-스킴 $X$가
\emph{$\mathcal{O}_S$-대수 $\mathcal{B}$에 대응한다}고 하거나
$X$를 \emph{$\mathcal{B}$의 스펙트럼}이라 하며, 이를
$\operatorname{Spec}(\mathcal{B})$로 나타낸다.

\begin{env}[1.3.2]
\label{II.1.3.2-ko}
$X$를 $S$ 위에서 아핀인 준스킴, $f:X\to S$를 그 구조 사상이라
하자. 임의의 아핀 열린집합 $U\subset S$에 대하여
$f^{-1}(U)$ 위에 유도된 준스킴은 환
$\Gamma(U,\mathcal{A}(X))$의 아핀 스킴이다.
\end{env}

\oldpage[II]{9}
(\hyperref[II.1.2.6-ko]{1.2.6})과
(\hyperref[II.1.3.1-ko]{1.3.1})에 따라 $X$가 어떤
$\mathcal{O}_S$-대수에 대응한다고 가정해도 되므로, 이 따름정리는
(\hyperref[II.1.3.1-ko]{1.3.1})에서 서술한 $X$의 구성에서 따른다.

\begin{env}[1.3.3]
\label{II.1.3.3-ko}
$S$를 점 $0$을 이중화한 체 $K$ 위의 아핀 평면이라 하자
(\hyperref[I.5.5.11-ko]{I, 5.5.11}). 그 표기에서 $S$는 두 아핀
열린집합 $Y_1$, $Y_2$의 합집합이다
(\hyperref[I.5.5.11-ko]{I, 5.5.11}). $f$가 열린 몰입 사상
$Y_1\to S$이면
$f^{-1}(Y_2)=Y_1\cap Y_2$는 $Y_1$ 안의 아핀 열린집합이 아니다
(\emph{loc. cit.}). 따라서 이로써 아핀 스킴이지만 $S$ 위에서
아핀이 아닌 예를 얻는다.
\end{env}

\begin{env}[1.3.4]
\label{II.1.3.4-ko}
$S$를 아핀 스킴이라 하자. $S$-준스킴 $X$가 $S$ 위에서
아핀이기 위한 필요충분조건은 $X$가 아핀 스킴인 것이다.
\end{env}

\begin{env}[1.3.5]
\label{II.1.3.5-ko}
$S$를 준스킴, $X$를 $S$ 위에서 아핀인 준스킴, $Y$를
$X$-준스킴이라 하자. $Y$가 $S$ 위에서 아핀이기 위한
필요충분조건은 $Y$가 $X$ 위에서 아핀인 것이다.
\end{env}

곧바로 $S$가 아핀 스킴인 경우로 환원된다. 그러면 $X$도 아핀
스킴이다(\hyperref[II.1.3.4-ko]{1.3.4}). 이때 명제의 두 조건은
모두 $Y$가 아핀 스킴이라는 것을 나타낸다
(\hyperref[II.1.3.4-ko]{1.3.4}).

\begin{env}[1.3.6]
\label{II.1.3.6-ko}
$X$를 $S$ 위에서 아핀인 준스킴이라 하자.
(\hyperref[II.1.3.5-ko]{1.3.5})에 따라, $X$ 위에서 아핀인
준스킴 $Y$를 정하는 것은 $S$ 위에서 아핀인 준스킴 $Y$와
$S$-사상 $g:Y\to X$를 정하는 것과 같다. 다시 말해
(\hyperref[II.1.3.1-ko]{1.3.1}과
\hyperref[II.1.2.7-ko]{1.2.7}), 이는 준연접
$\mathcal{O}_S$-대수 $\mathcal{B}$와 $\mathcal{O}_S$-대수 준동형
$\mathcal{A}(X)\to\mathcal{B}$를 정하는 것과 같다(이를
$\mathcal{B}$에 $\mathcal{A}(X)$-대수 구조를 정하는 것으로 볼 수
있다). $f:X\to S$가 구조 사상이면, 이때
$\mathcal{B}=f_*(g_*(\mathcal{O}_Y))$이다.
\end{env}

\begin{env}[1.3.7]
\label{II.1.3.7-ko}
$X$를 $S$ 위에서 아핀인 준스킴이라 하자. $X$가 $S$ 위에서
유한형이기 위한 필요충분조건은 준연접 $\mathcal{O}_S$-대수
$\mathcal{A}(X)$가 유한 생성인 것이다
(\hyperref[I.9.6.2-ko]{I, 9.6.2}).
\end{env}

정의 (\hyperref[I.9.6.2-ko]{I, 9.6.2})에 따라 $S$가 아핀인
경우로 환원된다. 그러면 $X$는 아핀 스킴이고
(\hyperref[II.1.3.4-ko]{1.3.4}), $S=\operatorname{Spec}(A)$,
$X=\operatorname{Spec}(B)$이면 $\mathcal{A}(X)$는
$\mathcal{O}_S$-대수 $\widetilde{B}$이다. 그런데
$\Gamma(S,\widetilde{B})=B$이므로,\footnote{\textbf{번역 주.} 인쇄
원문은 이 증명에서 정의되지 않은 $U$를 쓰지만,
EGA2-CANON-DISP-0030에 따라 $S$로 교정했다.} 이 따름정리는
(\hyperref[I.9.6.2-ko]{I, 9.6.2})와
(\hyperref[I.6.3.3-ko]{I, 6.3.3})에서 따른다.

\begin{env}[1.3.8]
\label{II.1.3.8-ko}
$X$를 $S$ 위에서 아핀인 준스킴이라 하자. $X$가 축소이기 위한
필요충분조건은 준연접 $\mathcal{O}_S$-대수 $\mathcal{A}(X)$가
축소인 것이다
(\hyperref[0.4.1.4-ko]{0, 4.1.4}).
\end{env}

실제로 이 문제는 명백히 $S$ 위에서 국소적이다.
(\hyperref[II.1.3.2-ko]{1.3.2})에 따라, 이 따름정리는
(\hyperref[I.5.1.1-ko]{I, 5.1.1})과
(\hyperref[I.5.1.4-ko]{I, 5.1.4})에서 따른다.

\subsection{$S$ 위에서 아핀인 준스킴 위의 준연접층}
\label{subsection:II.1.4-ko}

\begin{env}[1.4.1]
\label{II.1.4.1-ko}
$X$를 $S$ 위에서 아핀인 준스킴, $Y$를 $S$-준스킴,
$\mathcal{F}$(각각 $\mathcal{G}$)를 준연접
$\mathcal{O}_X$-가군(각각 $\mathcal{O}_Y$-가군)이라 하자. 그러면
$(h,u)\mapsto(\mathcal{A}(h),\mathcal{A}(u))$로 주어지는, 쌍
$(Y,\mathcal{G})$에서 $(X,\mathcal{F})$로 가는 사상들의 집합에서 쌍
$(\mathcal{A}(X),\mathcal{A}(\mathcal{F}))$에서
$(\mathcal{A}(Y),\mathcal{A}(\mathcal{G}))$로 가는 디-준동형들의
집합으로 가는 사상
(\hyperref[II.1.1.2-ko]{1.1.2}와
\hyperref[II.1.1.3-ko]{1.1.3})은 전단사이다.
\end{env}

증명은 (\hyperref[I.2.2.5-ko]{I, 2.2.5})와
(\hyperref[I.2.2.4-ko]{I, 2.2.4})를 사용하여
(\hyperref[II.1.2.7-ko]{1.2.7})의 증명과 정확히 같은 절차를 따르며,
세부사항은 독자에게 맡긴다.

\begin{env}[1.4.2]
\label{II.1.4.2-ko}
(\hyperref[II.1.4.1-ko]{1.4.1})의 가정들에 더하여 $Y$가 $S$ 위에서
아핀이라고 가정하면, $(h,u)$가 동형이기 위한 필요충분조건은
$(\mathcal{A}(h),\mathcal{A}(u))$가 디-동형인 것이다.
\end{env}

\oldpage[II]{10}
\begin{env}[1.4.3]
\label{II.1.4.3-ko}
준연접 $\mathcal{O}_S$-대수 $\mathcal{B}$와 준연접
$\mathcal{B}$-가군 $\mathcal{M}$으로 이루어진 모든 쌍
$(\mathcal{B},\mathcal{M})$에 대하여
($\mathcal{M}$을 $\mathcal{O}_S$-가군으로 보든
$\mathcal{B}$-가군으로 보든 준연접성은 같은 조건이다
(\hyperref[I.9.6.1-ko]{I, 9.6.1})), $S$ 위에서 아핀인 준스킴 $X$와
준연접 $\mathcal{O}_X$-가군 $\mathcal{F}$로 이루어진 쌍
$(X,\mathcal{F})$가 존재하여
$\mathcal{A}(X)=\mathcal{B}$이고
$\mathcal{A}(\mathcal{F})=\mathcal{M}$이다. 더욱이 이 쌍은 유일한
동형까지 유일하게 결정된다.
\end{env}

유일성은 (\hyperref[II.1.4.1-ko]{1.4.1})과
(\hyperref[II.1.4.2-ko]{1.4.2})에서 따른다. 존재성은
(\hyperref[II.1.3.1-ko]{1.3.1})에서와 같이 증명되며, 여기서도
세부사항은 독자에게 맡긴다. $\mathcal{O}_X$-가군 $\mathcal{F}$를
$\widetilde{\mathcal{M}}$로 나타내고, 이 가군을 준연접
$\mathcal{B}$-가군 $\mathcal{M}$에 \emph{대응하는} 가군이라고 한다.
모든 아핀 열린집합 $U\subset S$에 대하여
$\widetilde{\mathcal{M}}|p^{-1}(U)$는
($p$가 구조 사상 $X\to S$를 나타낼 때)
$(\Gamma(U,\mathcal{M}))^{\sim}$와 표준적으로 동형이다.

\begin{env}[1.4.4]
\label{II.1.4.4-ko}
준연접 $\mathcal{B}$-가군들의 범주에서
$\widetilde{\mathcal{M}}$는 $\mathcal{M}$에 관한 공변 가법 완전
함자이고, 귀납극한과 직합과 가환한다.
\end{env}

실제로 곧 $S$가 아핀인 경우로 환원되며, 그 경우 따름정리는
(\hyperref[I.1.3.5-ko]{I, 1.3.5}),
(\hyperref[I.1.3.9-ko]{I, 1.3.9})와
(\hyperref[I.1.3.11-ko]{I, 1.3.11})에서 따른다.

\begin{env}[1.4.5]
\label{II.1.4.5-ko}
(\hyperref[II.1.4.3-ko]{1.4.3})의 가정 아래에서
$\widetilde{\mathcal{M}}$가 유한 생성 $\mathcal{O}_X$-가군이기 위한
필요충분조건은 $\mathcal{M}$이 유한 생성 $\mathcal{B}$-가군인 것이다.
\end{env}

곧 $S=\operatorname{Spec}(A)$가 아핀 스킴인 경우로 환원된다. 이때
$\mathcal{B}=\widetilde{B}$이고, 여기서 $B$는 $A$-대수이다.\footnote{
\textbf{번역 주.} 인쇄 원문은 여기서 $B$가 유한 생성 $A$-대수라고
하고 (I, 9.6.2)를 인용하지만, 이는 명제의 가정에 없으므로
EGA2-CANON-DISP-0031에 따라 그 구절과 인용을 삭제했다.} 또한
$\mathcal{M}=\widetilde{M}$이고, 여기서 $M$은 $B$-가군이다
(\hyperref[I.1.3.13-ko]{I, 1.3.13}). \emph{준스킴 $X$ 위에서}
$\mathcal{O}_X$는 환 $B$에 대응하고
$\widetilde{\mathcal{M}}$는 $B$-가군 $M$에 대응한다. 따라서
$\widetilde{\mathcal{M}}$가 유한 생성이기 위한 필요충분조건은 $M$이
유한 생성인 것이다(\hyperref[I.1.3.13-ko]{I, 1.3.13}). 이로써
주장이 증명된다.

\begin{env}[1.4.6]
\label{II.1.4.6-ko}
$Y$를 $S$ 위에서 아핀인 준스킴, $X$와 $X'$을 $Y$ 위에서 아핀인
두 준스킴이라 하자(따라서 이들은 $S$ 위에서도 아핀이다
(\hyperref[II.1.3.5-ko]{1.3.5})). 다음과 같이 놓자.
$\mathcal{B}=\mathcal{A}(Y)$, $\mathcal{A}=\mathcal{A}(X)$,
$\mathcal{A}'=\mathcal{A}(X')$. 그러면 $X\times_Y X'$은 $Y$ 위에서
아핀이고(따라서 $S$ 위에서도 아핀이다),
$\mathcal{A}(X\times_Y X')$는
$\mathcal{A}\otimes_{\mathcal{B}}\mathcal{A}'$과 표준적으로
동일시된다.
\end{env}

실제로 (\hyperref[I.9.6.1-ko]{I, 9.6.1})에 따라
$\mathcal{A}\otimes_{\mathcal{B}}\mathcal{A}'$은 준연접
$\mathcal{B}$-대수이고, 따라서 준연접 $\mathcal{O}_S$-대수이기도
하다(\hyperref[I.9.6.1-ko]{I, 9.6.1}).
$Z$를 $\mathcal{A}\otimes_{\mathcal{B}}\mathcal{A}'$의 스펙트럼이라
하자(\hyperref[II.1.3.1-ko]{1.3.1}). 표준
$\mathcal{B}$-대수 준동형
$\mathcal{A}\to\mathcal{A}\otimes_{\mathcal{B}}\mathcal{A}'$과
$\mathcal{A}'\to\mathcal{A}\otimes_{\mathcal{B}}\mathcal{A}'$은
(\hyperref[II.1.2.7-ko]{1.2.7})에 따라 $Y$-사상 $p:Z\to X$와
$p':Z\to X'$에 대응한다. 삼중항 $(Z,p,p')$가 곱
$X\times_Y X'$임을 보이기 위해, $S$가 환 $C$의 아핀 스킴인
경우로 환원할 수 있다(\hyperref[I.3.2.6.4-ko]{I, 3.2.6.4}).
그러면 $Y$, $X$, $X'$은 아핀 스킴이고
(\hyperref[II.1.3.4-ko]{1.3.4}), 그 환 $B$, $A$, $A'$은
$C$-대수이며
$\mathcal{B}=\widetilde{B}$, $\mathcal{A}=\widetilde{A}$,
$\mathcal{A}'=\widetilde{A'}$이다.
(\hyperref[I.1.3.13-ko]{I, 1.3.13})에 따라
$\mathcal{A}\otimes_{\mathcal{B}}\mathcal{A}'$은
$\mathcal{O}_S$-대수 $\widetilde{A\otimes_B A'}$와 표준적으로
동일시된다. 따라서 환 $A(Z)$는 $A\otimes_B A'$과 동일시되고,
사상 $p$, $p'$은 표준 준동형
$A\to A\otimes_B A'$, $A'\to A\otimes_B A'$에 대응한다.
이제 명제는 (\hyperref[I.3.2.2-ko]{I, 3.2.2})에서 따른다.

\begin{env}[1.4.7]
\label{II.1.4.7-ko}
$\mathcal{F}$(각각 $\mathcal{F}'$)를 준연접
$\mathcal{O}_X$-가군(각각 $\mathcal{O}_{X'}$-가군)이라 하자. 그러면
$\mathcal{A}(\mathcal{F}\otimes_Y\mathcal{F}')$는
$\mathcal{A}(\mathcal{F})\otimes_{\mathcal{A}(Y)}
\mathcal{A}(\mathcal{F}')$과 표준적으로 동일시된다.
\end{env}

$\mathcal{F}\otimes_Y\mathcal{F}'$은 $X\times_Y X'$ 위에서
준연접이다(\hyperref[I.9.1.2-ko]{I, 9.1.2}).
$g:Y\to S$, $f:X\to Y$, $f':X'\to Y$를 구조 사상이라 하자.
그러면 구조 사상
\oldpage[II]{11}
$h:Z\to S$는 $g\circ f\circ p$ 및 $g\circ f'\circ p'$과 같다.
표준 준동형
\[
  \mathcal{A}(\mathcal{F})\otimes_{\mathcal{A}(Y)}\mathcal{A}(\mathcal{F}')
  \longrightarrow \mathcal{A}(\mathcal{F}\otimes_Y\mathcal{F}')
\]
을 다음과 같이 정의한다. 임의의 열린집합 $U\subset S$에 대하여
표준 준동형
\[
  \begin{aligned}
  \Gamma(f^{-1}(g^{-1}(U)),\mathcal{F})
  &\longrightarrow \Gamma(h^{-1}(U),p^*(\mathcal{F}))\quad\text{및}\quad
  \Gamma(f^{\prime-1}(g^{-1}(U)),\mathcal{F}')
  \longrightarrow \Gamma(h^{-1}(U),p^{\prime*}(\mathcal{F}'))
  \end{aligned}
\]
들이 있다(\hyperref[0.4.4.3-ko]{0, 4.4.3}). 따라서 표준 준동형
\[
  \begin{aligned}
  &\Gamma(f^{-1}(g^{-1}(U)),\mathcal{F})
    \otimes_{\Gamma(g^{-1}(U),\mathcal{O}_Y)}
    \Gamma(f^{\prime-1}(g^{-1}(U)),\mathcal{F}')\longrightarrow\\
  &\hspace{17em}
    \Gamma(h^{-1}(U),p^*(\mathcal{F}))
    \otimes_{\Gamma(h^{-1}(U),\mathcal{O}_Z)}
    \Gamma(h^{-1}(U),p^{\prime*}(\mathcal{F}'))
  \end{aligned}
\]
을 얻는다.

이것이 실제로 $\mathcal{A}(Z)$-가군의 동형임을 보이기 위해,
$S$(따라서 $X$, $X'$, $Y$, $X\times_Y X'$도)가 아핀 스킴인
경우로 환원할 수 있다. 이때
(\hyperref[II.1.4.6-ko]{1.4.6})의 표기를 사용하여
$\mathcal{F}=\widetilde{M}$, $\mathcal{F}'=\widetilde{M'}$라 쓸 수
있으며, $M$(각각 $M'$)은 $A$-가군(각각 $A'$-가군)이다. 그러면
$\mathcal{F}\otimes_Y\mathcal{F}'$은 $X\times_Y X'$ 위에서
$(A\otimes_B A')$-가군 $M\otimes_B M'$에 대응하는 층과
동일시된다(\hyperref[I.9.1.3-ko]{I, 9.1.3}). 따라서 따름정리는
$\mathcal{O}_S$-가군 $\widetilde{M\otimes_B M'}$와
$\widetilde{M}\otimes_{\widetilde{B}}\widetilde{M'}$의 표준적
동일시에서 따른다. 여기서 $M$, $M'$, $B$는 $C$-가군으로 본다
(\hyperref[I.1.3.13-ko]{I, 1.3.13}과
\hyperref[I.1.6.3-ko]{I, 1.6.3}).

특히 (\hyperref[II.1.4.7-ko]{1.4.7})을 $X=Y$이고
$\mathcal{F}'=\mathcal{O}_{X'}$인 경우에 적용하면,
$\mathcal{A}'$-가군 $\mathcal{A}(f^{\prime*}(\mathcal{F}))$가
$\mathcal{A}(\mathcal{F})\otimes_{\mathcal{B}}\mathcal{A}'$과
동일시됨을 알 수 있다.

\begin{env}[1.4.8]
\label{II.1.4.8-ko}
특히 $X=X'=Y$일 때($X$는 $S$ 위에서 아핀이다),
$\mathcal{F}$와 $\mathcal{G}$가 두 준연접
$\mathcal{O}_X$-가군이면
\[
  \mathcal{A}(\mathcal{F}\otimes_{\mathcal{O}_X}\mathcal{G})
  =\mathcal{A}(\mathcal{F})\otimes_{\mathcal{A}(X)}\mathcal{A}(\mathcal{G})
  \tag{1.4.8.1}\label{II.1.4.8.1-ko}
\]
이라는 표준적이고 함자적인 동형이 있다. 더욱이 $\mathcal{F}$가
유한 표시를 가지면
(\hyperref[I.1.6.3-ko]{I, 1.6.3}과
\hyperref[I.1.3.12-ko]{I, 1.3.12})에 따라
\[
  \mathcal{A}(\mathcal{H}om_X(\mathcal{F},\mathcal{G}))
  =\mathcal{H}om_{\mathcal{A}(X)}
    (\mathcal{A}(\mathcal{F}),\mathcal{A}(\mathcal{G}))
  \tag{1.4.8.2}\label{II.1.4.8.2-ko}
\]
이라는 표준적 동형이 있다.
\end{env}

\begin{env}[1.4.9]
\label{II.1.4.9-ko}
$X$, $X'$가 $S$ 위에서 아핀인 두 준스킴이면, 합
$X\sqcup X'$도 $S$ 위에서 아핀이다. 두 아핀 스킴의 합은 아핀
스킴이기 때문이다.
\end{env}

\begin{env}[1.4.10]
\label{II.1.4.10-ko}
$S$를 준스킴, $\mathcal{B}$를 준연접 $\mathcal{O}_S$-대수라 하고,
$X=\operatorname{Spec}(\mathcal{B})$라 하자. $\mathcal{J}$를
$\mathcal{B}$의 임의의 준연접 아이디얼층이라 하자. 그러면
$\widetilde{\mathcal{J}}$는 $\mathcal{O}_X$의 준연접 아이디얼층이다.
$Y$를 $X$에서 $\widetilde{\mathcal{J}}$로 정의되는 닫힌 부분준스킴이라
하면, 이는 $\operatorname{Spec}(\mathcal{B}/\mathcal{J})$와 표준적으로
동형이다.
\end{env}

실제로 (\hyperref[I.4.2.3-ko]{I, 4.2.3})에서 $Y$가 $S$ 위에서
아핀임이 곧바로 따른다. 따라서 (\hyperref[II.1.3.1-ko]{1.3.1})에 의해
$S$가 아핀인 경우로 환원되고, 이 경우 명제는
(\hyperref[I.4.1.2-ko]{I, 4.1.2})에서 곧바로 따른다.

(\hyperref[II.1.4.10-ko]{1.4.10})의 결과는 다음과 같이도 표현할 수
있다. $h:\mathcal{B}\to\mathcal{B}'$가 준연접 $\mathcal{O}_S$-대수들의
\emph{전사} 준동형이면,
${}^a h:\operatorname{Spec}(\mathcal{B}')\to
\operatorname{Spec}(\mathcal{B})$는 \emph{닫힌 몰입 사상}이다.

\oldpage[II]{12}
\begin{env}[1.4.11]
\label{II.1.4.11-ko}
$S$를 준스킴, $\mathcal{B}$를 준연접 $\mathcal{O}_S$-대수라 하고,
$X=\operatorname{Spec}(\mathcal{B})$라 하자. $\mathcal{K}$를
$\mathcal{O}_S$의 임의의 준연접 아이디얼층이라 하자. 그러면
($f$가 구조 사상 $X\to S$를 나타낼 때)
$f^*(\mathcal{K})\mathcal{O}_X=\widetilde{\mathcal{K}\mathcal{B}}$가
표준 동형까지 성립한다.
\end{env}

이 문제는 $S$ 위에서 국소적이므로 $S=\operatorname{Spec}(A)$가 아핀인
경우로 환원된다. 이 경우 명제는
(\hyperref[I.1.6.9-ko]{I, 1.6.9})에 다름 아니다.

\subsection{밑준스킴의 변경}
\label{II.1.5-ko}

\begin{env}[1.5.1]
\label{II.1.5.1-ko}
$X$를 $S$ 위에서 아핀인 준스킴이라 하자. 밑준스킴의 임의의 확장
$g:S'\to S$에 대하여 $X'=X_{(S')}=X\times_S S'$은 $S'$ 위에서
아핀이다.
\end{env}

$f'$을 사영 $X'\to S'$이라 하자. $g(U')$가 $S$의 어떤 아핀
열린집합 $U$에 포함되는 $S'$의 모든 아핀 열린집합 $U'$에 대하여
$f^{\prime-1}(U')$가 아핀 열린집합임을 보이면 충분하다
(\hyperref[II.1.2.1-ko]{1.2.1}). 따라서 $S$와 $S'$가 아핀인 경우로
한정할 수 있고, 이 경우 $X'$이 아핀 스킴임을 보이면 충분하다
(\hyperref[II.1.3.4-ko]{1.3.4}). 그런데 이때
(\hyperref[II.1.3.4-ko]{1.3.4}) $X$는 아핀 스킴이다. $S$, $S'$, $X$의
환을 각각 $A$, $A'$, $B$라 하면, $X'$은 환 $A'\otimes_A B$의 아핀
스킴임을 안다(\hyperref[I.3.2.2-ko]{I, 3.2.2}).

\begin{env}[1.5.2]
\label{II.1.5.2-ko}
(\hyperref[II.1.5.1-ko]{1.5.1})의 가정 아래에서 $f:X\to S$를 구조
사상, $f':X'\to S'$과 $g':X'\to X$를 사영이라 하자. 그러면 다음
도식은 가환한다.
\[
  \xymatrix{
    X\ar[d] & X'\ar[l]_{g'}\ar[d]^{f'}\\
    S & S'\ar[l]_g
  }
\]
모든 준연접 $\mathcal{O}_X$-가군 $\mathcal{F}$에 대하여 다음 표준
$\mathcal{O}_{S'}$-가군 동형이 존재한다.
\[
  u:g^*(f_*(\mathcal{F}))\xrightarrow{\sim}
  f'_*(g^{\prime*}(\mathcal{F})).
  \tag{1.5.2.1}\label{II.1.5.2.1-ko}
\]
특히 $\mathcal{A}(X')$에서 $g^*(\mathcal{A}(X))$로 가는 표준 동형이
존재한다.
\end{env}

$u$를 정의하려면 준동형
\[
  v:f_*(\mathcal{F})\longrightarrow
  g_*(f'_*(g^{\prime*}(\mathcal{F})))
  =f_*(g'_*(g^{\prime*}(\mathcal{F})))
\]
를 정의하고 $u=v^\sharp$로 두면 충분하다
(\hyperref[0.4.4.3-ko]{0, 4.4.3}). $v=f_*(\rho)$로 두겠다. 여기서
$\rho$는 표준 준동형
$\mathcal{F}\to g'_*(g^{\prime*}(\mathcal{F}))$이다
(\hyperref[0.4.4.3-ko]{0, 4.4.3}). $u$가 동형임을 보이기 위해서는
$S$와 $S'$, 따라서 $X$와 $X'$도 아핀인 경우로 한정할 수 있다.
(\hyperref[II.1.5.1-ko]{1.5.1})의 표기를 사용하면 이때
$\mathcal{F}=\widetilde{M}$이고, 여기서 $M$은 $B$-가군이다. 그러면
$g^*(f_*(\mathcal{F}))$와 $f'_*(g^{\prime*}(\mathcal{F}))$는 둘 다
$A'$-가군 $A'\otimes_A M$에 대응하는 $\mathcal{O}_{S'}$-가군과 같다
(여기서 $M$은 $A$-가군으로 본다). 또한 $u$는 항등사상에 대응하는
준동형임을 곧바로 알 수 있다
(\hyperref[I.1.6.3-ko]{I, 1.6.3},
\hyperref[I.1.6.5-ko]{1.6.5}와 \hyperref[I.1.6.7-ko]{1.6.7}).

\begin{env}[1.5.3]
\label{II.1.5.3-ko}
$X$가 더 이상 $S$ 위에서 아핀이라고 가정되지 않을 때에는
(\hyperref[II.1.5.2-ko]{1.5.2})가 여전히 성립한다고 생각해서는 안
된다. 이는 심지어 $S'=\operatorname{Spec}(k(s))$ ($s\in S$)이고
$S'\to S$가 표준 사상인 경우에도 마찬가지이다
(\hyperref[I.2.4.5-ko]{I, 2.4.5}) --- 이 경우 $X'$은 바로
\emph{섬유} $f^{-1}(s)$이다(\hyperref[I.3.6.2-ko]{I, 3.6.2}). 다시
말해 $X$가 $S$ 위에서 아핀이 아닐 때에는
\oldpage[II]{13}
``준연접층의 직상'' 연산은 ``섬유 취하기'' 연산과 교환하지 않는다.
그러나 제III장(\agref{III.4.2.4-ko}{III, 4.2.4})에서 $f$가 고유
사상이고(\agref{II.5.4-ko}{5.4}) $S$가 뇌터일 때 $X$ 위의
\emph{연접층}에 대하여 성립하는, 이 방향의 ``점근적'' 성격의 결과를
보게 될 것이다.
\end{env}

\begin{env}[1.5.4]
\label{II.1.5.4-ko}
$S$ 위에서 아핀인 모든 준스킴 $X$와 모든 $s\in S$에 대하여, 섬유
$f^{-1}(s)$는 아핀 스킴이다(여기서 $f$는 구조 사상 $X\to S$를
나타낸다).
\end{env}

(\hyperref[II.1.5.1-ko]{1.5.1})을 $S'=\operatorname{Spec}(k(s))$에
적용하고 (\hyperref[II.1.3.4-ko]{1.3.4})을 사용하면 충분하다.

\begin{corollary}[1.5.5]
\label{II.1.5.5-ko}
$X$를 $S$-준스킴, $S'$을 $S$ 위에서 아핀인 준스킴이라 하자. 그러면
$X'=X_{(S')}$은 $X$ 위에서 아핀인 준스킴이다. 또한 $f:X\to S$가
구조 사상이면 다음 표준 $\mathcal{O}_X$-대수 동형이 존재한다.
$\mathcal{A}(X')\xrightarrow{\sim}f^*(\mathcal{A}(S'))$.
그리고 모든 준연접 $\mathcal{A}(S')$-가군 $\mathcal{M}$에 대하여 다음
표준 디-동형이 존재한다.
$f^*(\mathcal{M})\xrightarrow{\sim}
\mathcal{A}(f^{\prime*}(\widetilde{\mathcal{M}}))$.
여기서 $f'=f_{(S')}$은 구조 사상 $X'\to S'$을 나타낸다.
\end{corollary}

(\hyperref[II.1.5.1-ko]{1.5.1})과
(\hyperref[II.1.5.2-ko]{1.5.2})에서 $X$와 $S'$의 역할을 서로 바꾸면
충분하다.

\begin{env}[1.5.6]
\label{II.1.5.6-ko}
이제 $S$, $S'$을 두 준스킴, $q:S'\to S$를 사상,
$\mathcal{B}$(각각 $\mathcal{B}'$)를 준연접
$\mathcal{O}_S$-대수(각각 준연접 $\mathcal{O}_{S'}$-대수),
$u:\mathcal{B}\to\mathcal{B}'$를 $q$-사상, 즉 $\mathcal{O}_S$-대수
준동형 $\mathcal{B}\to q_*(\mathcal{B}')$이라 하자.
$X=\operatorname{Spec}(\mathcal{B})$,
$X'=\operatorname{Spec}(\mathcal{B}')$라 하면, 다음 사상을 표준적으로
얻는다.
\[
  v=\operatorname{Spec}(u):X'\to X
\]
이때 도식
\[
  \xymatrix{
    X'\ar[r]^v\ar[d] & X\ar[d]\\
    S'\ar[r]_q & S
  }
  \tag{1.5.6.1}\label{II.1.5.6.1-ko}
\]
은 가환한다(수직 화살표들은 구조 사상이다). 실제로 $u$를 주는 것은
준연접 $\mathcal{O}_{S'}$-대수들의 준동형
$u^\sharp:q^*(\mathcal{B})\to\mathcal{B}'$를 주는 것과 동치이다
(\hyperref[0.4.4.3-ko]{0, 4.4.3}). 따라서 이는
$\mathcal{A}(w)=u^\sharp$인 표준 $S'$-사상
\[
  w:\operatorname{Spec}(\mathcal{B}')\to
  \operatorname{Spec}(q^*(\mathcal{B}))
\]
을 정의한다(\hyperref[II.1.2.7-ko]{1.2.7}). 한편
(\hyperref[II.1.5.2-ko]{1.5.2})에 따라
$\operatorname{Spec}(q^*(\mathcal{B}))$는 $X\times_S S'$과 표준적으로
동일시된다. 사상 $v$는 $w$에 이어 첫째 사영을 취한 합성
$X'\xrightarrow{w}X\times_S S'\xrightarrow{p_1}X$이고,
(\hyperref[II.1.5.6.1-ko]{1.5.6.1})의 가환성은 정의에서 따른다.
$U$(각각 $U'$)를 $q(U')\subset U$인 $S$(각각 $S'$)의 아핀
열린집합이라 하고,
$A=\Gamma(U,\mathcal{O}_S)$,
$A'=\Gamma(U',\mathcal{O}_{S'})$를 그 환들,
$B=\Gamma(U,\mathcal{B})$,
$B'=\Gamma(U',\mathcal{B}')$라 하자. $u$의 제한인
$(q|U')$-사상 $\mathcal{B}|U\to\mathcal{B}'|U'$은 대수의
디-준동형 $B\to B'$에 대응한다. $V$와 $V'$을 각각 구조 사상
$X\to S$와 $X'\to S'$에 의한 $U$와 $U'$의 역상이라 하면, $v$의 제한인 사상
$V'\to V$는 (\hyperref[I.1.7.3-ko]{I, 1.7.3})에 따라 앞의
디-준동형에 대응한다.
\end{env}

\begin{env}[1.5.7]
\label{II.1.5.7-ko}
(\hyperref[II.1.5.6-ko]{1.5.6})과 같은 가정 아래에서 $\mathcal{M}$을
준연접 $\mathcal{B}$-가군이라 하자. 그러면 다음 표준
$\mathcal{O}_{X'}$-가군 동형이 존재한다.
\[
  v^*(\widetilde{\mathcal{M}})\xrightarrow{\sim}
  \bigl(q^*(\mathcal{M})\otimes_{q^*(\mathcal{B})}\mathcal{B}'\bigr)^{\sim}
  \tag{1.5.7.1}\label{II.1.5.7.1-ko}
\]

\oldpage[II]{14}
실제로 (\hyperref[II.1.5.2.1-ko]{1.5.2.1})의 표준 동형은
$p_1^*(\widetilde{\mathcal{M}})$에서
$\operatorname{Spec}(q^*(\mathcal{B}))$ 위의
$q^*(\mathcal{B})$-가군 $q^*(\mathcal{M})$에 대응하는 층으로 가는
표준 동형을 준다. 이제 (\hyperref[II.1.4.7-ko]{1.4.7})을 적용하면
충분하다.
\end{env}

\subsection{아핀 사상}
\label{II.1.6-ko}

\begin{env}[1.6.1]
\label{II.1.6.1-ko}
준스킴의 사상 $f:X\to Y$를 구조 사상으로 보았을 때 $X$가 $Y$ 위에서
아핀인 준스킴이면, $f$를 \emph{아핀} 사상이라 한다. 한 준스킴 위에서
아핀인 준스킴들의 성질은 이 용어로 다음과 같이 표현된다.
\end{env}

\begin{proposition}[1.6.2]
\label{II.1.6.2-ko}
\begin{enumerate}
  \item[\rm(i)] 닫힌 몰입 사상은 아핀 사상이다.
  \item[\rm(ii)] 두 아핀 사상의 합성은 아핀 사상이다.
  \item[\rm(iii)] $f:X\to Y$가 아핀 $S$-사상이면, 밑준스킴의 모든 확장
    $S'\to S$에 대하여 $f_{(S')}:X_{(S')}\to Y_{(S')}$도 아핀 사상이다.
  \item[\rm(iv)] $f:X\to Y$와 $f':X'\to Y'$이 두 아핀 $S$-사상이면,
    \[
      f\times_S f':X\times_S X'\to Y\times_S Y'
    \]
    도 아핀 사상이다.
  \item[\rm(v)] $f:X\to Y$와 $g:Y\to Z$가 두 사상이고 $g\circ f$가
    아핀 사상이며 $g$가 분리 사상이면, $f$는 아핀 사상이다.
  \item[\rm(vi)] $f$가 아핀 사상이면 $f_{\mathrm{red}}$도 아핀 사상이다.
\end{enumerate}
\end{proposition}

(\hyperref[I.5.5.12-ko]{I, 5.5.12})에 따라 (i), (ii), (iii)만 증명하면
충분하다. 그런데 (i)는 (\hyperref[II.1.2.2-ko]{1.2.2})의 예이고,
(ii)는 (\hyperref[II.1.3.5-ko]{1.3.5})이며, 마지막으로 (iii)은
$X_{(S')}$이 곱 $X\times_Y Y_{(S')}$과 동일시되므로
(\hyperref[II.1.5.1-ko]{1.5.1})에서 따른다
(\hyperref[I.3.3.11-ko]{I, 3.3.11}).

\begin{corollary}[1.6.3]
\label{II.1.6.3-ko}
$X$가 아핀 스킴이고 $Y$가 스킴이면, 모든 사상 $X\to Y$는 아핀
사상이다.
\end{corollary}

\begin{proposition}[1.6.4]
\label{II.1.6.4-ko}
$Y$를 국소 뇌터 준스킴, $f:X\to Y$를 유한형 사상이라 하자. $f$가
아핀 사상이기 위한 필요충분조건은 $f_{\mathrm{red}}$가 아핀 사상인
것이다.
\end{proposition}

(\hyperref[II.1.6.2-ko]{1.6.2, (vi)})에 따라 조건의 충분성만 증명하면
된다. $Y$가 아핀이고 뇌터일 때 $X$가 아핀임을 보이면 충분하다. 이때
$Y_{\mathrm{red}}$는 아핀이므로, 가정에 따라 $X_{\mathrm{red}}$도
아핀이다. 한편 $X$는 뇌터이므로, 결론은
(\hyperref[I.6.1.7-ko]{I, 6.1.7})에서 따른다.

\subsection{가군의 층에 대응하는 벡터 다발}
\label{II.1.7-ko}

\begin{env}[1.7.1]
\label{II.1.7.1-ko}
$A$를 환, $E$를 $A$-가군이라 하자. $E$의 \emph{대칭 대수}라 하고
$\mathbf{S}(E)$(또는 $\mathbf{S}_A(E)$)로 나타내는 것은 텐서 대수
$\mathbf{T}(E)$를 $x,y\in E$에 대한 원소
$x\otimes y-y\otimes x$들이 생성하는 양쪽 아이디얼로 나눈 몫대수임을
상기하자. 대수 $\mathbf{S}(E)$는 다음 보편 성질로 특징지어진다.
$\sigma$를 표준 사상 $E\to\mathbf{S}(E)$, 즉
$E\to\mathbf{T}(E)$와 표준 사상
$\mathbf{T}(E)\to\mathbf{S}(E)$의 합성이라 하자. $B$가 가환
$A$-대수이면, 모든 $A$-선형 사상 $E\to B$는 유일하게
$E\xrightarrow{\sigma}\mathbf{S}(E)\xrightarrow{g}B$로 인수분해된다. 여기서
$g$는 $A$-대수 준동형이다. 이 특징지음에서 두 $A$-가군 $E$, $F$에
대하여
\[
  \mathbf{S}(E\oplus F)=\mathbf{S}(E)\otimes\mathbf{S}(F)
\]

\oldpage[II]{15}
가 표준 동형을 제외하면 성립함이 곧바로 따른다. 또한
$E\mapsto\mathbf{S}(E)$는 $A$-가군의 범주에서 가환 $A$-대수의 범주로
가는 공변 함자이다. 마지막으로 앞의 특징지음에 따라
$E=\varinjlim E_\lambda$이면
$\mathbf{S}(E)=\varinjlim\mathbf{S}(E_\lambda)$가 표준 동형을 제외하면
성립한다. 말의 남용으로 $x_i\in E$일 때의 곱
$\sigma(x_1)\sigma(x_2)\cdots\sigma(x_n)$을 혼동의 우려가 없으면 흔히
$x_1x_2\cdots x_n$으로 나타낸다. 대수 $\mathbf{S}(E)$는
\emph{등급} 대수이며, $\mathbf{S}_n(E)$는 $E$의 $n$개 원소의 곱들의
$A$-선형 결합 전체이다($n\geq 0$). 대수 $\mathbf{S}(A)$는 한 부정원의
다항식 대수 $A[T]$와 표준적으로 동형이고, 대수 $\mathbf{S}(A^n)$은
$n$개 부정원의 다항식 대수 $A[T_1,\ldots,T_n]$과 표준적으로 동형이다.
\end{env}

\begin{env}[1.7.2]
\label{II.1.7.2-ko}
$\varphi:A\to B$를 환 준동형이라 하자. $F$가 $B$-가군이면, 표준 사상
$F\to\mathbf{S}(F)$는 표준 사상
$F_{[\varphi]}\to\mathbf{S}(F)_{[\varphi]}$를 주고, 따라서 이 사상은
$F_{[\varphi]}\to\mathbf{S}(F_{[\varphi]})\to
\mathbf{S}(F)_{[\varphi]}$로 인수분해된다. 표준 $A$-대수 준동형
$\mathbf{S}(F_{[\varphi]})\to\mathbf{S}(F)_{[\varphi]}$는 전사이지만
반드시 전단사인 것은 아니다. $E$가 $A$-가군이면, 모든 디-준동형
$E\to F$, 즉 모든 $A$-가군 준동형 $E\to F_{[\varphi]}$는 따라서
표준적으로 $A$-대수 준동형
$\mathbf{S}(E)\to\mathbf{S}(F_{[\varphi]})\to
\mathbf{S}(F)_{[\varphi]}$를 준다.
이는 곧 대수의 디-준동형 $\mathbf{S}(E)\to\mathbf{S}(F)$이다.

같은 표기 아래에서 모든 $A$-가군 $E$에 대하여
$\mathbf{S}(E\otimes_A B)$는 대수
$\mathbf{S}(E)\otimes_A B$와 표준적으로 동일시된다. 이는
$\mathbf{S}(E)$의 보편 성질에서 곧바로 따른다
(\hyperref[II.1.7.1-ko]{1.7.1}).
\end{env}

\begin{env}[1.7.3]
\label{II.1.7.3-ko}
$R$을 환 $A$의 곱셈적 부분집합이라 하자.
(\hyperref[II.1.7.2-ko]{1.7.2})를 환 $B=R^{-1}A$에 적용하고
$R^{-1}E=E\otimes_A R^{-1}A$임을 상기하면,
$\mathbf{S}(R^{-1}E)=R^{-1}\mathbf{S}(E)$가 표준 동형을 제외하면
성립함을 알 수 있다. 또한 $R'\supset R$이 $A$의 또 하나의 곱셈적
부분집합이면 도식
\[
  \xymatrix{
    R^{-1}E\ar[r]\ar[d] & {R'}^{-1}E\ar[d]\\
    \mathbf{S}(R^{-1}E)\ar[r] & \mathbf{S}({R'}^{-1}E)
  }
\]
은 가환한다.
\end{env}

\begin{env}[1.7.4]
\label{II.1.7.4-ko}
이제 $(S,\mathcal{A})$를 환 달린 공간, $\mathcal{E}$를 $S$ 위의
$\mathcal{A}$-가군이라 하자. 모든 열린집합 $U\subset S$에
$\Gamma(U,\mathcal{A})$-가군
$\mathbf{S}(\Gamma(U,\mathcal{E}))$을 대응시키면,
$\mathbf{S}(E)$의 함자성에 따라 대수의 준층이 정의된다
(\hyperref[II.1.7.2-ko]{1.7.2}). 이 준층의 층화를
$\mathbf{S}(\mathcal{E})$ 또는
$\mathbf{S}_{\mathcal{A}}(\mathcal{E})$로 나타내고,
$\mathcal{A}$-가군 $\mathcal{E}$의 \emph{대칭 대수}라 한다.
(\hyperref[II.1.7.1-ko]{1.7.1})에서
$\mathbf{S}(\mathcal{E})$가 보편 문제의 해임이 곧바로 따른다.
$\mathcal{B}$가 $\mathcal{A}$-대수일 때 모든 $\mathcal{A}$-가군
준동형 $\mathcal{E}\to\mathcal{B}$는 유일하게
$\mathcal{E}\to\mathbf{S}(\mathcal{E})\to\mathcal{B}$로 인수분해되며,
둘째 화살표는 $\mathcal{A}$-대수 준동형이다. 따라서
$\mathcal{A}$-가군 준동형 $\mathcal{E}\to\mathcal{B}$들과
$\mathcal{A}$-대수 준동형
$\mathbf{S}(\mathcal{E})\to\mathcal{B}$들 사이에는 전단사 대응이
있다. 특히 모든 $\mathcal{A}$-가군 준동형
$u:\mathcal{E}\to\mathcal{F}$는 $\mathcal{A}$-대수 준동형
$\mathbf{S}(u):\mathbf{S}(\mathcal{E})\to
\mathbf{S}(\mathcal{F})$를 정의하며, 따라서
$\mathcal{E}\mapsto\mathbf{S}(\mathcal{E})$는 공변 함자이다.

\oldpage[II]{16}
(\hyperref[II.1.7.2-ko]{1.7.2})와 $\mathbf{S}$가 귀납극한과 가환한다는
사실에 따라, 모든 점 $s\in S$에 대하여
$(\mathbf{S}(\mathcal{E}))_s=\mathbf{S}(\mathcal{E}_s)$이다.
$\mathcal{E}$, $\mathcal{F}$가 두 $\mathcal{A}$-가군이면
$\mathbf{S}(\mathcal{E}\oplus\mathcal{F})$는
$\mathbf{S}(\mathcal{E})\otimes_{\mathcal{A}}
\mathbf{S}(\mathcal{F})$와 표준적으로 동일시된다. 이는 대응하는
준층들에서 확인된다.

또한 $\mathbf{S}(\mathcal{E})$는 $\mathbf{S}_n(\mathcal{E})$들의 무한
직합인 등급 $\mathcal{A}$-대수이다. 여기서 $\mathcal{A}$-가군
$\mathbf{S}_n(\mathcal{E})$는 준층
$U\mapsto\mathbf{S}_n(\Gamma(U,\mathcal{E}))$의 층화이다. 특히
$\mathcal{E}=\mathcal{A}$로 두면
$\mathbf{S}_{\mathcal{A}}(\mathcal{A})$는
$\mathcal{A}[T]=\mathcal{A}\otimes_{\mathbf{Z}}\mathbf{Z}[T]$와
동일시된다($T$는 부정원이고, $\mathbf{Z}$는 상수층으로 본다).
\end{env}

\begin{env}[1.7.5]
\label{II.1.7.5-ko}
$(T,\mathcal{B})$를 또 하나의 환 달린 공간,
$f:(S,\mathcal{A})\to(T,\mathcal{B})$를 환 달린 공간의 사상이라 하자.
$\mathcal{F}$가 $\mathcal{B}$-가군이면
$\mathbf{S}(f^*(\mathcal{F}))$는
$f^*(\mathbf{S}(\mathcal{F}))$와 표준적으로 동일시된다. 실제로
$f=(\psi,\theta)$이면 정의상
(\hyperref[0.4.3.1-ko]{0, 4.3.1})
\[
  \mathbf{S}(f^*(\mathcal{F}))
  =\mathbf{S}(\psi^*(\mathcal{F})
    \otimes_{\psi^*(\mathcal{B})}\mathcal{A})
  =\mathbf{S}(\psi^*(\mathcal{F}))
    \otimes_{\psi^*(\mathcal{B})}\mathcal{A}
\]
이다(\hyperref[II.1.7.2-ko]{1.7.2}). $S$의 임의의 열린집합 $U$와
$U$ 위의 $\mathbf{S}(\psi^*(\mathcal{F}))$의 임의의 단면 $h$에
대하여, 각 점 $s\in U$의 어떤 근방 $V$에서 $h$는
$\mathbf{S}(\Gamma(V,\psi^*(\mathcal{F})))$의 한 원소와 일치한다.
$\psi^*(\mathcal{F})$의 정의
(\hyperref[0.3.7.1-ko]{0, 3.7.1})를 참조하고, 가군 $E$에 대한
$\mathbf{S}(E)$의 모든 원소가 $E$의 원소들의 곱을 유한 개만 사용한
선형 결합임을 고려하면, $T$에서 $\psi(s)$의 어떤 근방 $W$,
$W$ 위의 $\mathbf{S}(\mathcal{F})$의 어떤 단면 $h'$, 그리고 $s$의
어떤 근방 $V'\subset V\cap\psi^{-1}(W)$가 존재하여 $V'$ 위에서
$h$는 $t\mapsto h'(\psi(t))$와 일치함을 알 수 있다. 따라서 주장이
따른다.
\end{env}

\begin{proposition}[1.7.6]
\label{II.1.7.6-ko}
$A$를 환, $S=\operatorname{Spec}(A)$를 그 소 스펙트럼,
$\mathcal{E}=\widetilde{M}$을 $A$-가군 $M$에 대응하는
$\mathcal{O}_S$-가군이라 하자. 그러면 $\mathcal{O}_S$-대수
$\mathbf{S}(\mathcal{E})$는 $A$-대수 $\mathbf{S}(M)$에 대응하는
$\mathcal{O}_S$-대수이다.
\end{proposition}

\begin{proof}
실제로 모든 $f\in A$에 대하여
$\mathbf{S}(M_f)=(\mathbf{S}(M))_f$이다
(\hyperref[II.1.7.3-ko]{1.7.3}). 따라서 명제는 정의들
(\hyperref[I.1.3.4-ko]{I, 1.3.4})에서 따른다.
\end{proof}

\begin{corollary}[1.7.7]
\label{II.1.7.7-ko}
$S$가 준스킴이고 $\mathcal{E}$가 준연접 $\mathcal{O}_S$-가군이면,
$\mathcal{O}_S$-대수 $\mathbf{S}(\mathcal{E})$는 준연접이다. 또한
$\mathcal{E}$가 유한 생성 $\mathcal{O}_S$-가군이면, 각
$\mathcal{O}_S$-가군 $\mathbf{S}_n(\mathcal{E})$도 유한 생성이다.
\end{corollary}

\begin{proof}
첫째 주장은 (\hyperref[II.1.7.6-ko]{1.7.6})과
(\hyperref[I.1.4.1-ko]{I, 1.4.1})에서 곧바로 따른다. 둘째 주장은
$E$가 유한 생성 $A$-가군이면 $\mathbf{S}_n(E)$도 유한 생성
$A$-가군이라는 사실에서 따른다. 이제
(\hyperref[I.1.3.13-ko]{I, 1.3.13})을 적용하면 된다.
\end{proof}

\begin{definition}[1.7.8]
\label{II.1.7.8-ko}
$\mathcal{E}$를 준연접 $\mathcal{O}_S$-가군이라 하자. 준연접
$\mathcal{O}_S$-대수 $\mathbf{S}(\mathcal{E})$의 스펙트럼
(\hyperref[II.1.3.1-ko]{1.3.1})을 $\mathbf{V}(\mathcal{E})$로 나타내고,
이를 $\mathcal{E}$로 정의되는 $S$ 위의 \emph{벡터 다발}이라 한다.
\end{definition}

(\hyperref[II.1.2.7-ko]{1.2.7})에 따라 임의의 $S$-준스킴 $X$에 대하여,
$S$-사상 $X\to\mathbf{V}(\mathcal{E})$들과 $\mathcal{O}_S$-대수 준동형
$\mathbf{S}(\mathcal{E})\to\mathcal{A}(X)$들 사이에는 표준적인 전단사
대응이 있고, 따라서 이 $S$-사상들과 $\mathcal{O}_S$-가군 준동형
$\mathcal{E}\to\mathcal{A}(X)=f_*(\mathcal{O}_X)$들 사이에도 그러한
대응이 있다(여기서 $f$는 구조 사상 $X\to S$이다). 특히 다음을 얻는다.

\begin{env}[1.7.9]
\label{II.1.7.9-ko}
$X$로 $S$가 열린집합 $U\subset S$ 위에 유도하는 부분준스킴을 잡자.
그러면 $S$-사상 $U\to\mathbf{V}(\mathcal{E})$들은 열린집합
$p^{-1}(U)$ 위에 $\mathbf{V}(\mathcal{E})$가 유도하는 $U$-준스킴의
$U$-단면들(\hyperref[I.2.5.5-ko]{I, 2.5.5})에 지나지 않는다(여기서
$p$는 구조 사상 $\mathbf{V}(\mathcal{E})\to S$이다). 방금 본 바에
따르면 이 $U$-단면들은 $\mathcal{O}_S$-가군 준동형
$\mathcal{E}\to j_*(\mathcal{O}_S|U)$들과 전단사로 대응한다(여기서
$j$는 표준 포함 사상 $U\to S$이다). 또는, 이와 동치로

\oldpage[II]{17}
(\hyperref[0.4.4.3-ko]{0, 4.4.3}), 이들은 $(\mathcal{O}_S|U)$-가군
준동형
$j^*(\mathcal{E})=\mathcal{E}|U\to\mathcal{O}_S|U$들과 대응한다. 또한
$S$-사상 $U\to\mathbf{V}(\mathcal{E})$을 열린집합 $U'\subset U$로
제한한 것에는 대응하는 준동형
$\mathcal{E}|U\to\mathcal{O}_S|U$를 $U'$로 제한한 것이 대응함은
자명하다. 따라서 $\mathbf{V}(\mathcal{E})$의 \emph{$S$-단면의 싹으로
이루어진 층}은 $\mathcal{E}$의 \emph{쌍대 $\check{\mathcal{E}}$}와
표준적으로 동일시된다.

특히 $X=U=S$로 두면 영 준동형
$\mathcal{E}\to\mathcal{O}_S$에는 $\mathbf{V}(\mathcal{E})$의 표준
$S$-단면 하나가 대응하며, 이를 \emph{영 $S$-단면}이라 한다
(8.3.3 참조).
\end{env}

\begin{env}[1.7.10]
\label{II.1.7.10-ko}
이제 $X$로 체 $K$의 스펙트럼 $\{\xi\}$를 잡자. 그러면 구조 사상
$f:X\to S$는 $s=f(\xi)$일 때 체의 단사 준동형
$k(s)\to K$에 대응한다(\hyperref[I.2.4.6-ko]{I, 2.4.6}). 따라서
$S$-사상 $\{\xi\}\to\mathbf{V}(\mathcal{E})$들은 $k(s)$의 확대체
$K$를 값체로 하는 $\mathbf{V}(\mathcal{E})$의 \emph{기하적 점들}
(\hyperref[I.3.4.5-ko]{I, 3.4.5})에 지나지 않으며, 이 기하적 점들은
$p^{-1}(s)$의 점들에 놓인다. 이 점들의 집합은 점 $s$ 위의
$\mathbf{V}(\mathcal{E})$의 \emph{$K$-유리 기하적 섬유}라고도 할 수
있으며, (\hyperref[II.1.7.8-ko]{1.7.8})에 따라 $\mathcal{O}_S$-가군
준동형 $\mathcal{E}\to f_*(\mathcal{O}_X)$들의 집합과 동일시된다.
또는, 이와 동치로 (\hyperref[0.4.4.3-ko]{0, 4.4.3}),
$\mathcal{O}_X$-가군 준동형
$f^*(\mathcal{E})\to\mathcal{O}_X=K$들의 집합과 동일시된다. 한편 정의상
(\hyperref[0.4.3.1-ko]{0, 4.3.1})
$f^*(\mathcal{E})=\mathcal{E}_s\otimes_{\mathcal{O}_s}K
=\mathcal{E}^s\otimes_{k(s)}K$이며, 여기서
$\mathcal{E}^s=\mathcal{E}_s/\mathfrak{m}_s\mathcal{E}_s$로 둔다.
따라서 점 $s$ 위의 $\mathbf{V}(\mathcal{E})$의 $K$-유리 기하적
섬유는 \emph{$K$-벡터 공간
$\mathcal{E}^s\otimes_{k(s)}K$}의 \emph{쌍대}와 동일시된다.
$\mathcal{E}^s$ 또는 $K$가 $k(s)$ 위에서 유한 차원이면 이 쌍대는
$(\mathcal{E}^s)^\vee\otimes_{k(s)}K$와도 동일시된다. 여기서
$(\mathcal{E}^s)^\vee$는 $k(s)$-벡터 공간 $\mathcal{E}^s$의 쌍대를
나타낸다.
\end{env}

\begin{proposition}[1.7.11]
\label{II.1.7.11-ko}
\begin{enumerate}
  \item[\rm(i)] $\mathbf{V}(\mathcal{E})$는 준연접
    $\mathcal{O}_S$-가군의 범주에서 아핀 $S$-스킴의 범주로 가는,
    $\mathcal{E}$에 관한 반변 함자이다.
  \item[\rm(ii)] $\mathcal{E}$가 유한 생성 $\mathcal{O}_S$-가군이면,
    $\mathbf{V}(\mathcal{E})$는 $S$ 위에서 유한형이다.
  \item[\rm(iii)] $\mathcal{E}$와 $\mathcal{F}$가 두 준연접
    $\mathcal{O}_S$-가군이면,
    $\mathbf{V}(\mathcal{E}\oplus\mathcal{F})$는
    $\mathbf{V}(\mathcal{E})\times_S\mathbf{V}(\mathcal{F})$와
    표준적으로 동일시된다.
  \item[\rm(iv)] $g:S'\to S$를 사상이라 하자. 임의의 준연접
    $\mathcal{O}_S$-가군 $\mathcal{E}$에 대하여,
    $\mathbf{V}(g^*(\mathcal{E}))$는
    $\mathbf{V}(\mathcal{E})_{(S')}
    =\mathbf{V}(\mathcal{E})\times_S S'$와 표준적으로 동일시된다.
  \item[\rm(v)] 준연접 $\mathcal{O}_S$-가군의 전사 준동형
    $\mathcal{E}\to\mathcal{F}$에는 닫힌 몰입 사상
    $\mathbf{V}(\mathcal{F})\to\mathbf{V}(\mathcal{E})$가 대응한다.
\end{enumerate}
\end{proposition}

\begin{proof}
임의의 $\mathcal{O}_S$-가군 준동형 $\mathcal{E}\to\mathcal{F}$가
$\mathcal{O}_S$-대수 준동형
$\mathbf{S}(\mathcal{E})\to\mathbf{S}(\mathcal{F})$를 표준적으로
정한다는 점을 고려하면, (i)는 (\hyperref[II.1.2.7-ko]{1.2.7})의
직접적인 귀결이다. (ii)는 정의들
(\hyperref[I.6.3.1-ko]{I, 6.3.1})과, $E$가 유한 생성 $A$-가군이면
$\mathbf{S}(E)$가 유한 생성 $A$-대수라는 사실에서 곧바로 따른다.
(iii)를 증명하려면 표준 동형
$\mathbf{S}(\mathcal{E}\oplus\mathcal{F})
\simeq\mathbf{S}(\mathcal{E})
\otimes_{\mathcal{O}_S}\mathbf{S}(\mathcal{F})$
(\hyperref[II.1.7.4-ko]{1.7.4})에서 출발하여
(\hyperref[II.1.4.6-ko]{1.4.6})을 적용하면 충분하다. 마찬가지로
(iv)를 증명하려면 표준 동형
$\mathbf{S}(g^*(\mathcal{E}))
\simeq g^*(\mathbf{S}(\mathcal{E}))$
(\hyperref[II.1.7.5-ko]{1.7.5})에서 출발하여
(\hyperref[II.1.5.2-ko]{1.5.2})를 적용하면 충분하다. 끝으로 (v)를
보이려면 준동형 $\mathcal{E}\to\mathcal{F}$가 전사일 때 이에 대응하는
$\mathcal{O}_S$-대수 준동형
$\mathbf{S}(\mathcal{E})\to\mathbf{S}(\mathcal{F})$도 전사임을
관찰하면 충분하며, 결론은 (\hyperref[II.1.4.10-ko]{1.4.10})에서
따른다.
\end{proof}

\begin{env}[1.7.12]
\label{II.1.7.12-ko}
특히 $\mathcal{E}=\mathcal{O}_S$로 잡자. 준스킴
$\mathbf{V}(\mathcal{O}_S)$는 $\mathcal{O}_S$-대수
$\mathbf{S}(\mathcal{O}_S)$의 스펙트럼인 아핀 $S$-스킴이며, 이
대수는 $\mathcal{O}_S$-대수
$\mathcal{O}_S[T]=\mathcal{O}_S\otimes_{\mathbf{Z}}\mathbf{Z}[T]$와
\oldpage[II]{18}
동일시된다($T$는 부정원이다). $S=\operatorname{Spec}(\mathbf{Z})$일
때에는 (\hyperref[II.1.7.6-ko]{1.7.6})에 따라 자명하며, 구조 사상
$S\to\operatorname{Spec}(\mathbf{Z})$를 고려하고
(\hyperref[II.1.7.11-ko]{1.7.11, (iv)})를 사용하면 여기서 일반적인
경우로 넘어갈 수 있다. 이 결과에 따라
$\mathbf{V}(\mathcal{O}_S)=S[T]$로도 쓰며, 따라서 다음 공식을 얻는다.
\[
  S[T]=S\otimes_{\mathbf{Z}}\mathbf{Z}[T].
  \tag{1.7.12.1}\label{II.1.7.12.1-ko}
\]

(\hyperref[I.3.3.15-ko]{I, 3.3.15})에서 이미 본, $S[T]$의
$S$-단면의 싹으로 이루어진 층과 $\mathcal{O}_S$의 동일시는 여기에서
(\hyperref[II.1.7.9-ko]{1.7.9})의 특수한 경우로서 더 일반적인 맥락
안에서 다시 얻어진다.
\end{env}

\begin{env}[1.7.13]
\label{II.1.7.13-ko}
임의의 $S$-준스킴 $X$에 대하여,
(\hyperref[II.1.7.8-ko]{1.7.8})에서
$\operatorname{Hom}_S(X,S[T])$가
$\operatorname{Hom}_{\mathcal{O}_S}(\mathcal{O}_S,\mathcal{A}(X))$와
표준적으로 동일시됨을 보았다. 이 집합은
$\Gamma(S,\mathcal{A}(X))$와 표준적으로 동형이므로 환 구조가 주어진다.
또한 임의의 $S$-사상 $h:X\to Y$에는 이 환 구조들에 관한 사상
$\Gamma(\mathcal{A}(h)):\Gamma(S,\mathcal{A}(Y))
\to\Gamma(S,\mathcal{A}(X))$가 대응한다
(\hyperref[II.1.1.2-ko]{1.1.2}).
$\operatorname{Hom}_S(X,S[T])$에 이와 같이 정의된 환 구조를 주면,
따라서 $\operatorname{Hom}_S(X,S[T])$를 $S$-준스킴의 범주에서 환의
범주로 가는 \emph{$X$에 관한 반변 함자}로 볼 수 있다.
한편 $\operatorname{Hom}_S(X,\mathbf{V}(\mathcal{E}))$도 마찬가지로
$\operatorname{Hom}_{\mathcal{O}_S}(\mathcal{E},\mathcal{A}(X))$와
동일시된다(여기서는 $\mathcal{A}(X)$를
\emph{$\mathcal{O}_S$-가군}으로 본다). 따라서 여기에 환
$\operatorname{Hom}_S(X,S[T])$ 위의 \emph{가군} 구조를 표준적으로
줄 수 있으며, 위와 마찬가지로 쌍
\[
  \bigl(\operatorname{Hom}_S(X,S[T]),
  \operatorname{Hom}_S(X,\mathbf{V}(\mathcal{E}))\bigr)
\]
은 $X$에 관한 반변 함자임을 알 수 있다. 이 함자는 환 $A$와
$A$-가군 $M$으로 이루어진 쌍 $(A,M)$을 대상으로 하고 디-준동형들을
사상으로 하는 범주에 값을 갖는다.

이 사실들을 $S[T]$가 \emph{$S$ 위의 환 스킴}이고
$\mathbf{V}(\mathcal{E})$가 $S$ 위의 환 스킴 $S[T]$ 위의
\emph{$S$ 위의 가군 스킴}이라는 말로 해석하겠다
(0장, \textsection8 참조).
\end{env}

\begin{env}[1.7.14]
\label{II.1.7.14-ko}
$S$-스킴 $\mathbf{V}(\mathcal{E})$ 위에 정의된 $S$ 위의 가군 스킴
구조를 사용하여 $\mathcal{O}_S$-가군 $\mathcal{E}$를 유일한 동형을
제외하고 복원할 수 있음을 보이겠다. 이를 위해 $\mathcal{E}$가 이
구조를 사용하여 정의되는
$\mathbf{S}(\mathcal{E})=\mathcal{A}(\mathbf{V}(\mathcal{E}))$의 한
부분-$\mathcal{O}_S$-가군과 표준적으로 동형임을 보이겠다. 실제로
(\hyperref[II.1.7.4-ko]{1.7.4})에 따라
\emph{$\mathcal{O}_S$-대수} 준동형들의 집합
$\operatorname{Hom}_{\mathcal{O}_S}(\mathbf{S}(\mathcal{E}),
\mathcal{A}(X))$는 \emph{$\mathcal{O}_S$-가군} 준동형들의 집합
$\operatorname{Hom}_{\mathcal{O}_S}(\mathcal{E},\mathcal{A}(X))$와
표준적으로 동일시된다. $h$, $h'$을 이 뒤의 집합의 두 원소,
$s_i$ ($1\leq i\leq n$)를 열린집합 $U\subset S$ 위의
$\mathcal{E}$의 단면들, $t$를 $U$ 위의 $\mathcal{A}(X)$의 단면이라
하면, 정의에 따라
\[
  (h+h')(s_1s_2\cdots s_n)
  =\prod_{i=1}^n\bigl(h(s_i)+h'(s_i)\bigr)
\]
이고
\[
  (t.h)(s_1s_2\cdots s_n)=t^n\prod_{i=1}^n h(s_i).
\]

이제 $z$가 $U$ 위의 $\mathbf{S}(\mathcal{E})$의 단면이면,
$h\mapsto h(z)$는
$\operatorname{Hom}_S(X,\mathbf{V}(\mathcal{E}))
=\operatorname{Hom}_{\mathcal{O}_S}(\mathbf{S}(\mathcal{E}),
\mathcal{A}(X))$에서 $\Gamma(U,\mathcal{A}(X))$로 가는 사상이다.
이제 다음을
\oldpage[II]{19}
보이겠다. $\mathcal{E}$는 다음 성질을 갖는
$\mathbf{S}(\mathcal{E})$의 부분가군과 동일시된다.
\emph{모든 열린집합 $U\subset S$, 이 부분-$\mathcal{O}_S$-가군의
$U$ 위의 모든 단면 $z$, 모든 $S$-준스킴 $X$에 대하여,
$\operatorname{Hom}_{\mathcal{O}_S|U}(\mathbf{S}(\mathcal{E})|U,
\mathcal{A}(X)|U)$에서 $\Gamma(U,\mathcal{A}(X))$로 가는 사상
$h\mapsto h(z)$가 $\Gamma(U,\mathcal{A}(X))$-가군 준동형이다.}

실제로 $\mathcal{E}$가 이 성질을 갖는다는 것은 곧바로 알 수 있다.
그 역을 증명하려면, $S=\operatorname{Spec}(A)$,
$\mathcal{E}=\widetilde{M}$일 때 $S$ 위의
$\mathbf{S}(\mathcal{E})$의 단면 $z$가 ($U=S$에 대하여) 위의
성질을 가지면 반드시 $\mathcal{E}$의 단면임을 증명하는 것으로
충분하다. 이때 $z=\sum_{n=0}^{\infty}z_n$이고
$z_n\in\mathbf{S}_n(M)$이며, $n\ne1$이면 $z_n=0$임을 보이면 된다.
$B=\mathbf{S}(M)$으로 놓고 $X$로 준스킴
$\operatorname{Spec}(B[T])$를 택하자. 여기서 $T$는 부정원이다.
집합
$\operatorname{Hom}_{\mathcal{O}_S}(\mathbf{S}(\mathcal{E}),
\mathcal{A}(X))$는 환 준동형 $h:B\to B[T]$들의 집합과 동일시된다
(\hyperref[I.1.3.13-ko]{I, 1.3.13}). 위에서 본 바에 따라
$(T.h)(z)=\sum_{n=0}^{\infty}T^nh(z_n)$이며, $z$에 관한 가정은 모든
준동형 $h$에 대하여
$\sum_{n=0}^{\infty}T^nh(z_n)=T.\sum_{n=0}^{\infty}h(z_n)$임을
뜻한다. 특히 $h$로 표준 포함 사상을 택하면
$\sum_{n=0}^{\infty}T^nz_n=T.\sum_{n=0}^{\infty}z_n$을 얻으며,
여기서 실제로 $n\ne1$이면 $z_n=0$이라는 결론이 따른다.
\end{env}

\begin{proposition}[1.7.15]
\label{II.1.7.15-ko}
$Y$를 바탕공간이 뇌터인 준스킴 또는 준콤팩트 스킴이라 하자.
$Y$ 위에서 유한형인 임의의 아핀 $Y$-스킴 $X$는
$\mathbf{V}(\mathcal{E})$ 꼴의 $Y$-스킴의 닫힌 부분-$Y$-스킴과
$Y$-동형이다. 여기서 $\mathcal{E}$는 유한 생성 준연접
$\mathcal{O}_Y$-가군이다.
\end{proposition}

\begin{proof}
실제로 준연접 $\mathcal{O}_Y$-대수 $\mathcal{A}(X)$는 유한 생성이다
(\hyperref[II.1.3.7-ko]{1.3.7}). 가정에 따라 $\mathcal{A}(X)$는
유한 생성 준연접 부분-$\mathcal{O}_Y$-가군 $\mathcal{E}$에 의해
대수로서 생성된다(\hyperref[I.9.6.5-ko]{I, 9.6.5}). 정의에 따라
이는 포함 사상 $\mathcal{E}\to\mathcal{A}(X)$를 표준적으로 연장하는
표준 준동형 $\mathbf{S}(\mathcal{E})\to\mathcal{A}(X)$가
\emph{전사}임을 뜻한다. 그러면 결론은
(\hyperref[II.1.4.10-ko]{1.4.10})에서 따른다.
\end{proof}

\section{동차 소 스펙트럼}
\label{section:II.2-ko}

\subsection{등급환과 등급가군에 관한 일반 사항.}
\label{subsection:II.2.1-ko}

\begin{notation}[2.1.1]
\label{II.2.1.1-ko}
음이 아닌 차수로 \emph{등급}이 주어진 환 $S$에 대하여,
차수가 $n$인 동차 원소들로 이루어진 $S$의 부분을 $S_n$으로
($n\geq0$), $n>0$인 $S_n$들의 합(직합)을 $S_+$로 나타내겠다.
$1\in S_0$이고, $S_0$는 $S$의 부분환이며, $S_+$는 $S$의 동차
아이디얼이고, $S$는 $S_0$와 $S_+$의 직합이다. $M$이 $S$ 위의
\emph{등급가군}이면(음의 차수도 허용한다), 마찬가지로
$M$의 차수가 $n$인 동차 원소들로 이루어진 $S_0$-가군을
$M_n$으로 나타내겠다(여기서 $n\in\mathbf{Z}$).

임의의 정수 $d>0$에 대하여, $S_{nd}$들의 직합을 $S^{(d)}$로
나타낸다. $S_{nd}$의 원소들을 차수가 $n$인 동차 원소로 보면,
$S_{nd}$들은 $S^{(d)}$ 위에 등급환 구조를 정의한다.

$0\leq k\leq d-1$인 임의의 정수 $k$에 대하여, $M^{(d,k)}$로
다음 직합을 나타내겠다.
\oldpage[II]{20}
$M_{nd+k}$ ($n\in\mathbf{Z}$)들의 직합이다.
$M_{nd+k}$의 원소들을 차수가 $n$인 동차 원소로 보면, 이는 등급
$S^{(d)}$-가군이다. $M^{(d,0)}$ 대신 $M^{(d)}$로 쓰겠다.

앞의 표기법 아래에서, 임의의 정수 $n$(음수도 허용한다)에 대하여,
모든 $k\in\mathbf{Z}$에 대해 $(M(n))_k=M_{n+k}$로 정의되는 등급
$S$-가군을 $M(n)$으로 나타내겠다. 특히 $S(n)$은
$(S(n))_k=S_{n+k}$를 만족하는 등급 $S$-가군이며,
$n<0$이면 $S_n=0$으로 놓기로 한다. 등급 $S$-가군 $M$이
\emph{등급가군으로서} $S(n)$ 꼴의 가군들의 직합과 동형이면,
\emph{자유}라고 한다. $S(n)$은 $S$의 원소 $1$을 차수가
$-n$인 원소로 보아 이를 생성원으로 갖는
$S$-가군이므로, 이는 $M$이 \emph{동차} 원소들로 이루어진
$S$ 위의 \emph{기저}를 갖는다고 말하는 것과 같다.

등급 $S$-가군 $M$이 \emph{유한 표시를 갖는다}는 것은, $P$와 $Q$가
$S(n)$ 꼴의 가군들의 유한 직합이고 준동형들의 차수가 $0$인 완전열
$P\to Q\to M\to0$이 존재한다는 뜻이다
(\agref{II.2.1.2-ko}{2.1.2} 참조).
\end{notation}

\begin{env}[2.1.2]
\phantomsection
\label{II.2.1.2-ko}
$M$, $N$을 두 등급 $S$-가군이라 하자. $M\otimes_S N$ 위에 다음과
같이 \emph{등급} $S$-가군 구조를 정의한다. 텐서곱
$M\otimes_{\mathbf{Z}}N$ 위에는
$\mathbf{Z}$-가군의 등급 구조($\mathbf{Z}$에는
$\mathbf{Z}_0=\mathbf{Z}$, $n\ne0$일 때 $\mathbf{Z}_n=0$인 등급을
준다)를 다음과 같이 정의할 수 있다.
\[
  (M\otimes_{\mathbf{Z}}N)_q
  =\bigoplus_{m+n=q}M_m\otimes_{\mathbf{Z}}N_n
\]
($M$과 $N$은 각각 $M_m$들과 $N_n$들의 직합이므로,
$M\otimes_{\mathbf{Z}}N$을 모든 $M_m\otimes_{\mathbf{Z}}N_n$의
직합과 표준적으로 동일시할 수 있음을 알고 있다.) 이제
$M\otimes_S N=(M\otimes_{\mathbf{Z}}N)/P$이다. 여기서 $P$는
$x\in M$, $y\in N$, $s\in S$에 대하여
$(xs)\otimes y-x\otimes(sy)$ 꼴의 원소들로 생성되는
$M\otimes_{\mathbf{Z}}N$의 부분-$\mathbf{Z}$-가군이다. $P$가
$M\otimes_{\mathbf{Z}}N$의 \emph{등급} 부분-$\mathbf{Z}$-가군임은
명백하며, 몫을 취하면 실제로 $M\otimes_S N$ 위에 등급 $S$-가군
구조가 얻어짐을 곧바로 알 수 있다.

두 등급 $S$-가군 $M$, $N$에 대하여, $S$-가군 준동형
$u:M\to N$이 모든 $j\in\mathbf{Z}$에 대하여
$u(M_j)\subset N_{j+k}$를 만족하면 \emph{차수 $k$}라고 함을
상기하자. $M$에서 $N$으로 가는 차수 $n$인 모든 준동형의 집합을
$H_n$이라 하고, $M$에서 $N$으로 가는 모든 ($S$-가군) 준동형들의
$S$-가군 $H$ 안에서 $H_n$들의 \emph{합}(직합)($n\in\mathbf{Z}$)을
$\operatorname{Hom}_S(M,N)$으로 나타낸다. 일반적으로
$\operatorname{Hom}_S(M,N)$은 앞서 말한 가군 전체와 같지 않다. 다만 $M$이
\emph{유한 생성}이면 $H=\operatorname{Hom}_S(M,N)$이다. 실제로 이때
$M$이 유한 개의 동차 원소 $x_i$ ($1\leq i\leq n$)로 생성된다고
할 수 있고, 모든 준동형 $u\in H$는 유일하게
$\sum_{k\in\mathbf{Z}}u_k$로 쓰인다. 여기서 각 $k$에 대하여
$u_k(x_i)$는 $u(x_i)$의 차수가 $k+\deg(x_i)$인 동차 성분과 같고
($1\leq i\leq n$), 따라서 유한 개의 지표를 제외하면 $u_k=0$이다.
정의에 따라 $u_k\in H_k$이므로 결론을 얻는다.

$\operatorname{Hom}_S(M,N)$의 차수 $0$인 원소들을 \emph{등급
$S$-가군의 준동형}이라 한다. $S_mH_n\subset H_{m+n}$임은 명백하므로,
$H_n$들은 $\operatorname{Hom}_S(M,N)$ 위에 등급 $S$-가군 구조를
정의한다.

이 정의들로부터 두 등급 $S$-가군 $M$, $N$에 대하여 다음을 곧바로
얻는다.
\[
  \label{II.2.1.2.1-ko}
  M(m)\otimes_S N(n)=(M\otimes_S N)(m+n),
  \tag{2.1.2.1}
\]
\[
  \label{II.2.1.2.2-ko}
  \operatorname{Hom}_S(M(m),N(n))
  =(\operatorname{Hom}_S(M,N))(n-m)
  \tag{2.1.2.2}
\]
\oldpage[II]{21}
$S$, $S'$을 두 등급환이라 하자. \emph{등급환}의 준동형
$\varphi:S\to S'$은 모든 $n\in\mathbf{Z}$에 대하여
$\varphi(S_n)\subset S'_n$을 만족하는 환 준동형이다(다시 말해,
$\varphi$는 등급 $\mathbf{Z}$-가군의 \emph{차수 $0$인} 준동형이어야
한다). 이러한 준동형 하나를 주면 $S'$ 위에 \emph{등급} $S$-가군 구조가
정의된다. 이 구조와 원래의 등급환 구조를 갖춘 $S'$을
\emph{등급 $S$-대수}라고 한다.

이제 $M$이 등급 $S$-가군이면, 등급 $S$-가군의 텐서곱
$M\otimes_S S'$ 위에 \emph{등급} $S'$-가군 구조가 자연스럽게 주어지며,
그 등급은 위에서 정의한 것이다.
\end{env}

\begin{lemma}[2.1.3]
\phantomsection
\label{II.2.1.3-ko}
$S$를 음이 아닌 차수로 등급이 주어진 환이라 하자. 동차 원소들로
이루어진 $S_+$의 부분집합 $E$가 $S_+$를 $S$-가군으로서 생성할
필요충분조건은 $E$가 $S$를 $S_0$-대수로서 생성하는 것이다.
\end{lemma}

\begin{proof}
조건이 충분함은 명백하다. 필요함을 보이자. $E$의 차수가 $n$인
원소들의 집합(각각 차수가 $\leq n$인 원소들의 집합)을 $E_n$
(각각 $E^n$)이라 하자. $n>0$에 대한 귀납법으로, $S_n$이
$E^n$의 원소들의 곱인 차수 $n$의 원소들로 생성되는 $S_0$-가군임을
보이면 충분하다. 그런데 $n=1$일 때 이는 가정에 의해 명백하다.
또한 가정으로부터
$S_n=\sum_{p=0}^{n-1}S_pE_{n-p}$이고, 이제 귀납 논법은 즉시
따른다.
\end{proof}

\begin{corollary}[2.1.4]
\phantomsection
\label{II.2.1.4-ko}
$S_+$가 유한 생성 아이디얼이기 위한 필요충분조건은 $S$가
$S_0$ 위의 유한 생성 대수인 것이다.
\end{corollary}

\begin{proof}
실제로 $S_0$-대수 $S$ (각각 $S$-아이디얼 $S_+$)의 유한 생성계가
동차 원소들로 이루어져 있다고 언제나 가정할 수 있다. 주어진 각 생성원을
그 동차 성분들로 바꾸면 되기 때문이다.
\end{proof}

\begin{corollary}[2.1.5]
\phantomsection
\label{II.2.1.5-ko}
$S$가 뇌터이기 위한 필요충분조건은 $S_0$가 뇌터이고 $S$가
$S_0$ 위의 유한 생성 대수인 것이다.
\end{corollary}

\begin{proof}
조건이 충분함은 명백하다. 필요함은 $S_0$가 $S/S_+$와 동형이고
$S_+$가 유한 생성 아이디얼이어야 한다는 데서 따른다
(\hyperref[II.2.1.4-ko]{2.1.4}).
\end{proof}

\begin{lemma}[2.1.6]
\phantomsection
\label{II.2.1.6-ko}
$S$를 음이 아닌 차수로 등급이 주어진 환이며 $S_0$ 위의 유한 생성
대수라고 하자. $M$을 유한 생성 등급 $S$-가군이라 하자. 그러면 다음이
성립한다.
\begin{enumerate}
  \item[\rm(i)] $M_n$들은 유한 생성 $S_0$-가군이고, 모든
    $n\leq n_0$에 대하여 $M_n=0$이 되는 정수 $n_0$가 존재한다.
  \item[\rm(ii)] 모든 정수 $n\geq n_1$에 대하여
    $M_{n+h}=S_hM_n$이 되는 정수 $n_1$과 정수 $h>0$이 존재한다.
  \item[\rm(iii)] $d>0$, $0\leq k\leq d-1$을 만족하는 모든 정수쌍
    $(d,k)$에 대하여 $M^{(d,k)}$는 유한 생성 $S^{(d)}$-가군이다.
  \item[\rm(iv)] 모든 정수 $d>0$에 대하여 $S^{(d)}$는 $S_0$ 위의
    유한 생성 대수이다.
  \item[\rm(v)] 모든 $m>0$에 대하여 $S_{mh}=(S_h)^m$이 되는 정수
    $h>0$이 존재한다.
  \item[\rm(vi)] 모든 정수 $n>0$에 대하여, 모든 $m\geq m_0$에 대해
    $S_m\subset S_+^n$이 되는 정수 $m_0$가 존재한다.
\end{enumerate}
\end{lemma}

\begin{proof}
$S$가 $S_0$-대수로서 차수가 $h_i$인 동차 원소 $f_i$
($1\leq i\leq r$)들로 생성되고, $M$이 $S$-가군으로서 차수가 $k_j$인
동차 원소 $x_j$ ($1\leq j\leq s$)들로 생성된다고 가정할 수 있다.
$M_n$은 $\alpha_i$들이 $\geq0$인 정수이고
$k_j+\sum_i\alpha_i h_i=n$을 만족할 때의 원소
\oldpage[II]{22}
$f_1^{\alpha_1}\cdots f_r^{\alpha_r}x_j$들의, 계수가 $S_0$에 속하는
선형결합들로
이루어짐이 명백하다. 각 $j$에 대하여 이 방정식을 만족하는
$(\alpha_i)$는 유한 개뿐이다. 이는 $h_i$들이 $>0$이기 때문이며, 따라서
(i)의 첫 번째 주장이 따른다. 두 번째 주장은 명백하다. 한편 $h$를
$h_i$들의 최소공배수라 하고, 모든 $g_i$의 차수가 $h$가 되도록
$g_i=f_i^{h/h_i}$ ($1\leq i\leq r$)로 놓자. 또 $z_\mu$들을
$0\leq\alpha_i<h/h_i$ ($1\leq i\leq r$)인
$f_1^{\alpha_1}\cdots f_r^{\alpha_r}x_j$ 꼴의 $M$의 원소들이라 하자.
이 원소들은 유한 개이고, 그 차수들 가운데 가장 큰 것을 $n_1$이라 하자.
$n\geq n_1$이면 $M_{n+h}$의 모든 원소가 $z_\mu$들의 선형결합이며 그
계수들은 $g_i$들에 관한 차수가 $>0$인 단항식임이 명백하다. 따라서
$M_{n+h}=S_hM_n$이고, 이로써 (ii)가 증명된다. 같은 방식으로 (모든
$d>0$에 대하여) $M^{(d,k)}$의 원소는 $g$가 $S$의 동차 원소이고
$0\leq\alpha_i<d$일 때의
$g^df_1^{\alpha_1}\cdots f_r^{\alpha_r}x_j$ 꼴의 원소들의, 계수가
$S_0$에 속하는 선형결합임을 알 수 있다. 따라서 (iii)가 성립한다. 이제
$M=S_+$로 놓고 $(S_+)^{(d)}=(S^{(d)})_+$임을 이용하면, (iv)는
(iii)와 보조정리 (\hyperref[II.2.1.3-ko]{2.1.3})에서 따른다. (v)는
(ii)에서 $M=S$로 놓으면 얻어진다. 마지막으로 주어진 $n$에 대하여
$\alpha_i\geq0$이고 $\sum_i\alpha_i<n$인 계 $(\alpha_i)$는 유한 개뿐이다.
따라서 이 계들에 대한 합 $\sum_i\alpha_i h_i$의 최댓값을 $m_0$라 하면
$m>m_0$에 대하여 $S_m\subset S_+^n$이고, 이로써 (vi)가 증명된다.
\footnote{\textbf{번역 주.} 인쇄 원문의 (vi)는 $m\geq m_0$라고 서술하지만,
여기서 $m_0$로 정의한 바로 그 최댓값을 쓰면 증명이 직접 주는 범위는
$m>m_0$이다. 존재명제 자체는 이 최댓값에 $1$을 더한 수를 임계값으로
취하면 원문의 $m\geq m_0$ 형태로 성립한다.}
\end{proof}

\begin{corollary}[2.1.7]
\phantomsection
\label{II.2.1.7-ko}
$S$가 뇌터이면 모든 정수 $d>0$에 대하여 $S^{(d)}$도 뇌터이다.
\end{corollary}

\begin{proof}
이는 (\hyperref[II.2.1.5-ko]{2.1.5})와
(\hyperref[II.2.1.6-ko]{2.1.6, (iv)})에서 따른다.
\end{proof}

\begin{env}[2.1.8]
\phantomsection
\label{II.2.1.8-ko}
$\mathfrak{p}$를 등급환 $S$의 동차 소아이디얼이라 하자. 그러면
$\mathfrak{p}$는 부분군 $\mathfrak{p}_n=\mathfrak{p}\cap S_n$들의
직합이다. $\mathfrak{p}$가 $S_+$를 포함하지 않는다고 가정하자. 그러면
$\mathfrak{p}$에 속하지 않는 $f\in S_+$에 대하여, 관계
$f^nx\in\mathfrak{p}$는 $x\in\mathfrak{p}$와 동치이다. 특히
$f\in S_d$ ($d>0$)이면, 모든 $x\in S_{m-nd}$에 대하여 관계
$f^nx\in\mathfrak{p}_m$은 $x\in\mathfrak{p}_{m-nd}$와 동치이다.
\end{env}

\begin{proposition}[2.1.9]
\phantomsection
\label{II.2.1.9-ko}
$n_0$를 $>0$인 정수라 하고, 각 $n\geq n_0$에 대하여
$\mathfrak{p}_n$을 $S_n$의 부분군이라 하자. $S_+$를 포함하지 않으며
모든 $n\geq n_0$에 대하여
$\mathfrak{p}\cap S_n=\mathfrak{p}_n$을 만족하는 $S$의 동차
소아이디얼 $\mathfrak{p}$가 존재하기 위한 필요충분조건은 다음 조건들이
성립하는 것이다.
\begin{enumerate}
  \item[$1^\circ$] 모든 $m\geq 0$과 모든 $n\geq n_0$에 대하여
    $S_m\mathfrak{p}_n\subset\mathfrak{p}_{m+n}$이다.
  \item[$2^\circ$] $m\geq n_0$, $n\geq n_0$, $f\in S_m$, $g\in S_n$일 때,
    관계 $fg\in\mathfrak{p}_{m+n}$은 $f\in\mathfrak{p}_m$ 또는
    $g\in\mathfrak{p}_n$을 함의한다.
  \item[$3^\circ$] 적어도 하나의 $n\geq n_0$에 대하여
    $\mathfrak{p}_n\neq S_n$이다.
\end{enumerate}
또한 이때 동차 소아이디얼 $\mathfrak{p}$는 유일하다.
\end{proposition}

\begin{proof}
조건 $1^\circ$와 $2^\circ$가 필요함은 명백하다. 또한
$\mathfrak{p}\not\supset S_+$이면
$\mathfrak{p}\cap S_k\neq S_k$인 $k>0$이 적어도 하나 존재한다.
$f\in S_k$가 $\mathfrak{p}$에 속하지 않으면, 관계
$\mathfrak{p}\cap S_n=S_n$은
(\hyperref[II.2.1.8-ko]{2.1.8})에 따라
$\mathfrak{p}\cap S_{n-mk}=S_{n-mk}$를 함의한다. 따라서 어떤 값 이상의
모든 $n$에 대하여 $\mathfrak{p}\cap S_n=S_n$이라면
$\mathfrak{p}\supset S_+$가 되어 가정에 모순이다. 이로써 $3^\circ$가
필요함이 증명된다. 역으로 조건 $1^\circ$, $2^\circ$, $3^\circ$가
성립한다고 가정하자. 어떤 정수 $d\geq n_0$에 대하여
$f\in S_d$가 $\mathfrak{p}_d$에 속하지 않으면, $\mathfrak{p}$가
존재할 때 $m<n_0$에 대한 $\mathfrak{p}_m$은 유한 개의 값을 제외한 모든
$r$에 대하여 $f^rx\in\mathfrak{p}_{m+rd}$를 만족하는
$x\in S_m$들의 집합과 반드시 같음에 유의하자. 이는 $\mathfrak{p}$가 존재한다면
유일함을 이미 증명한다. 이제 앞의 조건으로 $m<n_0$에 대한
$\mathfrak{p}_m$들을 정의하면
$\mathfrak{p}=\sum_{n=0}^{\infty}\mathfrak{p}_n$이 소아이디얼임을 보이는
일만 남는다. 먼저 $2^\circ$에 따라, $m\geq n_0$에 대해서도
$\mathfrak{p}_m$은 유한 개의 값을 제외한 모든 $r$에 대하여
$f^rx\in\mathfrak{p}_{m+rd}$를 만족하는 $x\in S_m$들의 집합으로
정의됨에 유의하자. 이제
\oldpage[II]{23}
$g\in S_m$, $x\in\mathfrak{p}_n$이면, 유한 개의 값을 제외한 모든
$r$에 대하여 $f^rgx\in\mathfrak{p}_{m+n+rd}$이다. 따라서
$gx\in\mathfrak{p}_{m+n}$이고, 이는 $\mathfrak{p}$가 $S$의
아이디얼임을 증명한다. 이 아이디얼이 소아이디얼임을, 다시 말해 부분군
$S_n/\mathfrak{p}_n$들로 등급이 주어진 환 $S/\mathfrak{p}$가 정역임을
보이기 위해서는 ($S/\mathfrak{p}$의 두 원소의 최고차 동차성분들을
고려하여) $x\in S_m$, $y\in S_n$이
$x\notin\mathfrak{p}_m$, $y\notin\mathfrak{p}_n$을 만족하면
$xy\notin\mathfrak{p}_{m+n}$임을 보이면 충분하다. 그렇지 않다면 충분히
큰 $r$에 대하여 $f^{2r}xy\in\mathfrak{p}_{m+n+2rd}$일 것이다. 그러나
모든 $r>0$에 대하여 $f^ry\notin\mathfrak{p}_{n+rd}$이다. 그러면
$2^\circ$에 따라 유한 개의 값을 제외한 모든 $r$에 대하여
$f^rx\in\mathfrak{p}_{m+rd}$이고, 이로부터
$x\in\mathfrak{p}_m$을 얻어 가정에 모순이다.
\end{proof}

\begin{env}[2.1.10]
\phantomsection
\label{II.2.1.10-ko}
$S_+$의 부분집합 $\mathfrak{J}$가 $S$의 아이디얼이면 이를
\emph{$S_+$의 아이디얼}이라 한다. 또한 $\mathfrak{J}$가 $S_+$와,
\emph{$S_+$를 포함하지 않는} $S$의 동차 소아이디얼의 교집합이면 이를
\emph{$S_+$의 동차 소아이디얼}이라 한다(더욱이 이 소아이디얼은
(\hyperref[II.2.1.9-ko]{2.1.9})에 따라 유일하다). $\mathfrak{J}$가
$S_+$의 아이디얼이면, \emph{$S_+$에서의 $\mathfrak{J}$의 근기}는 어떤
거듭제곱이 $\mathfrak{J}$에 속하는 $S_+$의 원소들로 이루어진 집합이다.
다시 말해 이는 $r_+(\mathfrak{J})=r(\mathfrak{J})\cap S_+$이다. 특히
$S_+$에서 $0$의 근기는 다시 $S_+$의 \emph{멱영근기}라 하고
$\mathfrak{N}_+$로 나타낸다. 따라서 이는 $S_+$의 멱영원 전체의
집합이다. $\mathfrak{J}$가 $S_+$의 \emph{동차} 아이디얼이면 그 근기
$r_+(\mathfrak{J})$도 동차 아이디얼이다. 실제로 몫환
$S/\mathfrak{J}$로 넘어가면 $\mathfrak{J}=0$인 경우로 환원할 수 있고,
$x=x_h+x_{h+1}+\cdots+x_k$가 멱영원이면
$x_i\in S_i$ ($1\leq h\leq i\leq k$)들도 모두 멱영원임을 보이면
충분하다. $x_k\neq 0$이라고 가정할 수 있고, 이때 $x^n$의 최고차
성분은 $x_k^n$이므로 $x_k$는 멱영원이다. 이제 $k$에 대한 귀납법을
적용한다. 등급환 $S$가 \emph{본질적으로 축소}라는 것은
$\mathfrak{N}_+=0$이라는 뜻이다. 다시 말해 $S_+$가 $\neq 0$인 멱영원을
포함하지 않는다는 뜻이다.
\end{env}

\begin{env}[2.1.11]
\phantomsection
\label{II.2.1.11-ko}
등급환 $S$에서 원소 $x$가 어떤 영이 아닌 원소와 곱하여 $0$이 되는
영인자이면 그 최고차 동차성분도 영인자임에 유의하자. 환 $S$가
\emph{본질적으로 정역}이라는 것은 $S_+$가
(\emph{단위원을 갖지 않는 환으로서}) 영인자를 포함하지 않아 영이 아닌
두 원소의 곱이 $0$일 수 없고, 또한 $\neq 0$인 것을 뜻한다. 이를 위해서는
$S_+$의 동차 원소 가운데 $\neq 0$인 것은 모두 이 환에서 영인자가 아니어서
어떤 영이 아닌 원소와 곱해도 $0$이 되지 않으면 충분하다.
$\mathfrak{p}$가 $S_+$의
동차 소아이디얼이면 $S/\mathfrak{p}$가 본질적으로 정역임은 명백하다.

$S$를 본질적으로 정역인 등급환이라 하고 $x_0\in S_0$라 하자. 이때
$S_+$의 동차 원소 $f\neq 0$가 \emph{하나라도} 존재하여
$x_0f=0$이면 $x_0S_+=0$이다. 실제로 모든 $g\in S_+$에
대하여 $(x_0g)f=(x_0f)g=0$이고, 가정에 따라 $x_0g=0$이다. 따라서
$S$가 정역이기 위한 필요충분조건은 $S_0$가 정역이고 $S_0$에서
$S_+$의 소멸자가 $0$뿐인 것이다.
\end{env}

\subsection{등급환의 분수환.}
\label{subsection:II.2.2-ko}

\begin{env}[2.2.1]
\phantomsection
\label{II.2.2.1-ko}
$S$를 음이 아닌 차수로 등급이 주어진 환이라 하고, $f$를 차수가
$d>0$인 $S$의 \emph{동차} 원소라 하자. 그러면 분수환 $S'=S_f$는
$S'_n$을 $x\in S_{n+kd}$이고 $k\geq 0$일 때의 $x/f^k$ 전체의
집합으로 놓음으로써 등급환이 된다(여기서 $n$은 임의의 음의 값도 취할
수 있음에 유의하자). $S'$의 \emph{차수가 $0$인} 원소들로 이루어진
부분환 $S'_0=(S_f)_0$을 $S_{(f)}$로 나타내겠다.

$f\in S_d$이면 $S_f$의 단항식 $(f/1)^h$들($h$는 영을 포함한 임의의
정수)은 환 $S_{(f)}$ 위의 \emph{자유계}를 이루며, 그 선형결합 전체의
집합은 바로
\oldpage[II]{24}
환 $(S^{(d)})_f$이다. 따라서 이 환은
\emph{$S_{(f)}[T,T^{-1}]=S_{(f)}\otimes_{\mathbf{Z}}\mathbf{Z}[T,T^{-1}]$와
동형이다}($T$는 부정원이다). 실제로
$\sum_{h=-a}^{b}z_h(f/1)^h=0$인 관계가 있고 $z_h=x_h/f^m$이며 모든
$x_h$가 $S_{md}$에 속한다고 하자. 정의에 따라 이 관계는
어떤 $k\geq a$가 존재하여
$\sum_{h=-a}^{b}f^{h+k}x_h=0$이 되는 것과 동치이다. 이 합의 항들은
차수가 서로 다르므로 모든 $h$에 대하여 $f^{h+k}x_h=0$이고, 따라서
모든 $h$에 대하여 $z_h=0$이다.

$M$이 등급 $S$-가군이면 $M'=M_f$는 등급 $S_f$-가군이다. 여기서
$M'_n$은 $z/f^k$ 꼴의 원소들로 이루어진 집합이며, 여기서
$z\in M_{n+kd}$이고 $k\geq 0$이다.
$M'$의 차수가 $0$인 동차 원소들의 집합을 $M_{(f)}$로 나타낸다.
$M_{(f)}$가 $S_{(f)}$-가군이고
$(M^{(d)})_f=M_{(f)}\otimes_{S_{(f)}}(S^{(d)})_f$임은 바로 알 수 있다.
\end{env}

\begin{lemma}[2.2.2]
\phantomsection
\label{II.2.2.2-ko}
$d$, $e$를 $>0$인 두 정수, $f\in S_d$, $g\in S_e$라 하자. 다음과 같은
환의 표준 동형이 존재한다.
\[
  S_{(fg)}\xrightarrow{\sim}(S_{(f)})_{g^d/f^e}~;
\]
이 두 환을 표준적으로 동일시하면, 다음과 같은 가군의 표준 동형이
존재한다.
\[
  M_{(fg)}\xrightarrow{\sim}(M_{(f)})_{g^d/f^e}.
\]
\end{lemma}

\begin{proof}
실제로 $fg$는 $f^e g^d$의 약수이고, 후자는 $(fg)^{de}$의 약수이므로
등급환 $S_{fg}$와 $S_{f^e g^d}$는 표준적으로 동일시된다. 다른 한편
$S_{f^e g^d}$는 $(S_{f^e})_{g^d/1}$와도 동일시되고
(\hyperref[0.1.4.6-ko]{0, I.4.6}), $f^e/1$은 $S_{f^e}$에서 가역이므로
$S_{f^e g^d}$는 $(S_{f^e})_{g^d/f^e}$와도 동일시된다. 그런데
$g^d/f^e$는 차수가 $0$인 $S_{f^e}$의 원소이다. 따라서
$(S_{f^e})_{g^d/f^e}$에서 차수가 $0$인 원소들로 이루어진 부분환은
$(S_{(f^e)})_{g^d/f^e}$임을 곧바로 얻는다. 또 명백히
$S_{(f^e)}=S_{(f)}$이므로, 이는 명제의 첫 번째 부분을 증명한다.
두 번째 부분도 같은 방식으로 확립된다.
\end{proof}

\begin{env}[2.2.3]
\phantomsection
\label{II.2.2.3-ko}
(\hyperref[II.2.2.2-ko]{2.2.2})의 가정 아래에서, 표준 준동형
$S_f\to S_{fg}$ (\hyperref[0.1.4.1-ko]{0, I.4.1})은 $x/f^k$를
$g^k x/(fg)^k$로 보내며 차수가 $0$임은 명백하다. 따라서
제한함으로써
\emph{표준 준동형 $S_{(f)}\to S_{(fg)}$}을 얻으며, 다음 도식은
가환한다.
\[
  \xymatrix{
    & S_{(f)}\ar[dl]\ar[dr]\\
    S_{(fg)}\ar[rr]^-{\sim} && (S_{(f)})_{g^d/f^e}
  }
\]
마찬가지로 표준 준동형 $M_{(f)}\to M_{(fg)}$을 정의한다.
\end{env}

\begin{lemma}[2.2.4]
\phantomsection
\label{II.2.2.4-ko}
$f$, $g$가 $S_+$의 두 동차 원소이면, 환 $S_{(fg)}$는 $S_{(f)}$와
$S_{(g)}$의 표준 상들의 합집합에 의해 생성된다.
\end{lemma}

\begin{proof}
(\hyperref[II.2.2.2-ko]{2.2.2})에 따라
$1/(g^d/f^e)=f^{d+e}/(fg)^d$가 $S_{(g)}$의 $S_{(fg)}$ 안에서의 표준
상에 속함을 보이면 충분하며, 이는 정의상 명백하다.
\end{proof}

\begin{proposition}[2.2.5]
\phantomsection
\label{II.2.2.5-ko}
$d$를 $>0$인 정수, $f\in S_d$라 하자. 그러면 다음과 같은 환의
표준 동형이 존재한다.
$S_{(f)}\xrightarrow{\sim}S^{(d)}/(f-1)S^{(d)}$~; 이 동형으로 두 환을
표준적으로 동일시하면, 다음과 같은 가군의 표준 동형이 존재한다.
$M_{(f)}\xrightarrow{\sim}M^{(d)}/(f-1)M^{(d)}$.
\end{proposition}

\begin{proof}
이 동형들 중 첫 번째 것은 $x/f^n$($x\in S_{nd}$)에 원소
$\overline{x}$, 곧 $x$의 $(f-1)S^{(d)}$에 대한 잉여류를 대응시켜 정의한다.
이 사상은 잘 정의된다. 실제로 모든 양의 정수 지수에 대하여 합동식
$f^h x\equiv x\pmod{(f-1)S^{(d)}}$이 모든 $x\in S^{(d)}$에 대해
성립한다. 따라서 $f^h x=0$인 $h>0$가 있으면
\oldpage[II]{25}
$\overline{x}=0$이다. 한편 $x\in S_{nd}$이고 $x=(f-1)y$라 하자. 여기서
$y=y_{hd}+y_{(h+1)d}+\cdots+y_{kd}$이고 $y_{jd}\in S_{jd}$이며
$y_{hd}\neq 0$이다. 그러면 반드시 $h=n$이고 $x=-y_{hd}$이며, 또한
관계 $y_{(j+1)d}=fy_{jd}$가 $h\leq j\leq k-1$에 대해 성립하고
$fy_{kd}=0$이다.
이로부터 마침내 $f^{k-n+1}x=0$을 얻는다. 따라서 $x\in S_{nd}$인 원소가
나타내는 $(f-1)S^{(d)}$에 대한 각 잉여류 $\overline{x}$에 $S_{(f)}$의 원소
$x/f^n$을 대응시킬 수 있다.
앞의 고찰에 따라 이 사상은 잘 정의된다. 이렇게 정의된 두 사상이
환 준동형이며 서로 역임을 곧바로 알 수 있다. $M$에 대해서도 완전히
같은 방식으로 한다.
\end{proof}

\begin{corollary}[2.2.6]
\phantomsection
\label{II.2.2.6-ko}
$S$가 뇌터이면, 차수가 $>0$인 임의의 동차 원소 $f$에 대하여
$S_{(f)}$도 뇌터이다.
\end{corollary}

\begin{proof}
이는 (\hyperref[II.2.1.7-ko]{2.1.7})과
(\hyperref[II.2.2.5-ko]{2.2.5})에서 곧바로 따른다.
\end{proof}

\begin{env}[2.2.7]
\phantomsection
\label{II.2.2.7-ko}
$T$를 \emph{동차} 원소들로 이루어진 $S_+$의 곱셈적 부분집합이라 하자.
그러면 $T_0=T\cup\{1\}$은 $S$의 곱셈적 부분집합이다. $T_0$의
원소들은 동차이므로, 환 $T_0^{-1}S$에는 명백한 방식으로 여전히 등급이
주어진다. $S_{(T)}$로 $T_0^{-1}S$에서 차수가 $0$인 원소들, 즉
$x/h$ 꼴의 원소들로 이루어진 부분환을 나타내겠다. 여기서 $h\in T$이고
$x$는 $h$와 차수가 같은 동차 원소이다.
(\hyperref[0.1.4.5-ko]{0, 1.4.5})에 따라 $T_0^{-1}S$는 환
$S_f$들($f$가 $T$를 달리며 전이사상은 표준 준동형
$S_f\to S_{fg}$이다)의 귀납극한과 표준적으로 동일시됨이 알려져 있다.
이 동일시는 차수를 보존하므로 $S_{(T)}$를
$S_{(f)}$들의($f\in T$) \emph{귀납극한}과 동일시한다. 모든 등급
$S$-가군 $M$에 대하여도 마찬가지로 가군 $M_{(T)}$(환 $S_{(T)}$ 위의
가군)를 차수가 $0$인 $T_0^{-1}M$의 원소들로 이루어진 것으로 정의한다.
이 가군은 $M_{(f)}$들의($f\in T$) 귀납극한임을 알 수 있다.

$\mathfrak{p}$가 $S_+$의 동차 소아이디얼이면, 각각
$S_{(\mathfrak{p})}$와 $M_{(\mathfrak{p})}$로 환 $S_{(T)}$와 가군
$M_{(T)}$를 나타내겠다. 여기서 $T$는 $S_+$에서 $\mathfrak{p}$에
속하지 않는 \emph{동차} 원소들의 집합이다.
\end{env}

\subsection{등급환의 동차 소 스펙트럼.}
\label{subsection:II.2.3-ko}

\begin{env}[2.3.1]
\phantomsection
\label{II.2.3.1-ko}
음이 아닌 차수로 등급이 주어진 환 $S$에 대하여, $S$의
\emph{동차 소 스펙트럼}이라 하고 $\operatorname{Proj}(S)$로 나타내는 것은
$S_+$의 동차 소아이디얼들의 집합
(\hyperref[II.2.1.10-ko]{2.1.10}), 또는 같은 말로 $S$의 동차
소아이디얼 가운데 $S_+$를 \emph{포함하지 않는} 것들의 집합이다. 이제
$\operatorname{Proj}(S)$를 바탕집합으로 갖는 \emph{스킴} 구조를 정의하겠다.
\end{env}

\begin{env}[2.3.2]
\phantomsection
\label{II.2.3.2-ko}
$E$를 $S$의 임의의 부분집합이라 하고, $V_+(E)$를 $S$의 동차
소아이디얼 가운데 $E$를 포함하고 $S_+$를 포함하지 않는 것들의 집합이라
하자. 그러므로 이는 $V(E)\cap\operatorname{Proj}(S)$라는
$\operatorname{Spec}(S)$의 부분집합이다. 따라서
(\hyperref[I.1.1.2-ko]{I, 1.1.2})에서 다음을 얻는다.
\[
  \label{II.2.3.2.1-ko}
  V_+(0)=\operatorname{Proj}(S),\quad
  V_+(S)=V_+(S_+)=\varnothing
  \tag{2.3.2.1}
\]
\[
  \label{II.2.3.2.2-ko}
  V_+\left(\bigcup_\lambda E_\lambda\right)
  =\bigcap_\lambda V_+(E_\lambda)
  \tag{2.3.2.2}
\]
\[
  \label{II.2.3.2.3-ko}
  V_+(EE')=V_+(E)\cup V_+(E')
  \tag{2.3.2.3}
\]
$V_+(E)$는 집합 $E$를 $E$가 생성하는 동차 아이디얼로 바꾸어도 변하지
않는다. 또한 $\mathfrak{J}$가 $S$의 동차 아이디얼이면 다음이 성립한다.
\[
  \label{II.2.3.2.4-ko}
  V_+(\mathfrak{J})
  =V_+\left(\bigcup_{q\geq n}(\mathfrak{J}\cap S_q)\right)
  \tag{2.3.2.4}
\]
\oldpage[II]{26}
이는 모든 $n>0$에 대하여 성립한다. 실제로
$\mathfrak{p}\in\operatorname{Proj}(S)$가 $\mathfrak{J}$의 차수가
$\geq n$인 동차 원소들을 포함한다고 하자. 가정에 따라 동차 원소
$f\in S_d$ 가운데 $\mathfrak{p}$에 속하지 않는 것이 존재하므로, 모든
$m\geq 0$과 모든 $x\in S_m\cap\mathfrak{J}$에 대하여 관계
$f^r x\in\mathfrak{J}\cap S_{m+rd}$는 모든 음이 아닌 정수 $r$에 대해
성립한다. 한편 유한 개의 값을 제외한 모든 지수에 대해서는
$m+rd\geq n$이므로 위 교집합은 해당 소아이디얼에 포함되고, 따라서
$f^r x\in\mathfrak{p}\cap S_{m+rd}$이다. 앞서 고른 동차 원소가 해당
소아이디얼에 속하지 않으므로 소아이디얼의 성질로 그 거듭제곱을 소거하면
$x\in\mathfrak{p}\cap S_m$을 얻는다
(\hyperref[II.2.1.8-ko]{2.1.8}).

마지막으로 $\mathfrak{J}$가 $S$의 임의의 동차 아이디얼이면
\[
  \label{II.2.3.2.5-ko}
  V_+(\mathfrak{J})=V_+(r_+(\mathfrak{J})).
  \tag{2.3.2.5}
\]
\end{env}

\begin{env}[2.3.3]
\phantomsection
\label{II.2.3.3-ko}
정의에 따라 $V_+(E)$ 꼴의 집합들은 $X=\operatorname{Proj}(S)$ 위에서
$\operatorname{Spec}(S)$의 스펙트럼 위상으로부터 유도된 위상의
닫힌집합들이다. 이 위상도 $X$ 위의 \emph{스펙트럼 위상}이라 하겠다.
모든 $f\in S$에 대하여 다음과 같이 놓는다.
\[
  \label{II.2.3.3.1-ko}
  D_+(f)=D(f)\cap\operatorname{Proj}(S)
  =\operatorname{Proj}(S)\setminus V_+(f)
  \tag{2.3.3.1}
\]
따라서 $f$, $g$가 $S$의 두 원소이면
(\hyperref[I.1.1.9.1-ko]{I, 1.1.9.1})
\[
  \label{II.2.3.3.2-ko}
  D_+(fg)=D_+(f)\cap D_+(g).
  \tag{2.3.3.2}
\]
\end{env}

\begin{proposition}[2.3.4]
\phantomsection
\label{II.2.3.4-ko}
$f$가 $S_+$의 동차 원소 전체를 달릴 때, $D_+(f)$들은
$X=\operatorname{Proj}(S)$의 위상의 기저를 이룬다.
\end{proposition}

\begin{proof}
실제로 (\hyperref[II.2.3.2.2-ko]{2.3.2.2})와
(\hyperref[II.2.3.2.4-ko]{2.3.2.4})에 따라 $X$의 모든 닫힌 부분집합은
$V_+(f)$ 꼴 집합들의 교집합이다. 여기서 $f$는 차수가 $>0$인 동차
원소이다.
\end{proof}

\begin{env}[2.3.5]
\phantomsection
\label{II.2.3.5-ko}
$f$를 $S_+$의 차수가 $d>0$인 \emph{동차} 원소라 하자.
$\mathfrak{p}$가 $S$의 동차 소아이디얼이고 $f$를 포함하지 않을 때,
$x/f^n$ 꼴 원소
가운데 $x\in\mathfrak{p}$이고 $n\geq 0$인 것들의 집합은 분수환
$S_f$의 소아이디얼임이 알려져 있다
(\hyperref[0.1.2.6-ko]{0, 1.2.6}). 그 아이디얼과 $S_{(f)}$의 교집합도
이 환의 소아이디얼이므로 이를 $\psi_f(\mathfrak{p})$로 나타내겠다.
이는 $x/f^n$ 가운데 $n\geq 0$이고
$x\in\mathfrak{p}\cap S_{nd}$인 것들의 집합이다. 이로써 사상
\[
  \psi_f:D_+(f)\longrightarrow\operatorname{Spec}(S_{(f)})~;
\]
을 정의하였다. 또한 $g\in S_e$가 $S_+$의 또 하나의 동차 원소이면 다음
도식은 가환한다.
\[
  \label{II.2.3.5.1-ko}
  \xymatrix{
    D_+(f)\ar[r]^{\psi_f}
    & \operatorname{Spec}(S_{(f)})\\
    D_+(fg)\ar[u]\ar[r]_{\psi_{fg}}
    & \operatorname{Spec}(S_{(fg)})\ar[u]
  }
  \tag{2.3.5.1}
\]
여기서 왼쪽 수직 화살표는 포함 사상이고, 오른쪽 수직 화살표는 사상
${}^a\omega_{fg,f}$이다. 이는 표준 준동형
$\omega=\omega_{fg,f}:S_{(f)}\to S_{(fg)}$에서 유도된다
(\hyperref[I.1.2.1-ko]{I, 1.2.1}). 실제로
$x/f^n\in\omega^{-1}(\psi_{fg}(\mathfrak{p}))$이고
$fg\notin\mathfrak{p}$라 하자. 정의에 따라
$g^n x/(fg)^n\in\psi_{fg}(\mathfrak{p})$이므로
$g^n x\in\mathfrak{p}$이고, 따라서 $x\in\mathfrak{p}$이다. 역은
자명하다.
\end{env}

\begin{proposition}[2.3.6]
\phantomsection
\label{II.2.3.6-ko}
사상 $\psi_f$는 $D_+(f)$에서 $\operatorname{Spec}(S_{(f)})$로의
위상동형이다.
\end{proposition}

\begin{proof}
먼저 $\psi_f$는 연속이다. 실제로 $h\in S_{nd}$이고
$h/f^n\in\psi_f(\mathfrak{p})$이면 정의에 따라
$h\in\mathfrak{p}$이며, 그 역도 성립한다. 따라서
\[
  \psi_f^{-1}(D(h/f^n))=D_+(hf)
\]
이고, 주장은 공식
(\hyperref[II.2.3.3.2-ko]{2.3.3.2})에서 따른다. 또한 $D_+(hf)$들, 즉
$h$가 여러 집합 $S_{nd}$의 원소 전체를 달릴 때 얻는 이 집합들은,
(\hyperref[II.2.3.4-ko]{2.3.4})와 공식
(\hyperref[II.2.3.3.2-ko]{2.3.3.2})에 따라 $D_+(f)$의 위상의 기저를
이룬다.
\oldpage[II]{27}
그러므로 공리 $(T_0)$가 $D_+(f)$와 $\operatorname{Spec}(S_{(f)})$에서
성립함을 위의 논의와 함께 고려하면, $\psi_f$는 단사이고 그 역인 사상
$\psi_f(D_+(f))\longrightarrow D_+(f)$는 연속이다. 마지막으로
$\psi_f$가 전사임을 보이자. $\mathfrak{q}_0$를 $S_{(f)}$의
소아이디얼이라 하고, 모든 $n>0$에 대하여 $\mathfrak{p}_n$을
$x\in S_n$ 가운데 $x^d/f^n\in\mathfrak{q}_0$인 것들의 집합이라 하자.
그러면 $\mathfrak{p}_n$들은 (\hyperref[II.2.1.9-ko]{2.1.9})의 조건들을
만족한다. 실제로 $x\in S_n$, $y\in S_n$이고
$x^d/f^n\in\mathfrak{q}_0$, $y^d/f^n\in\mathfrak{q}_0$라 하자. 그러면
$(x+y)^{2d}/f^{2n}\in\mathfrak{q}_0$이고, 따라서
$(x+y)^d/f^n\in\mathfrak{q}_0$이다. 이는 $\mathfrak{q}_0$가
소아이디얼이기 때문이다. 이로써
$\mathfrak{p}_n$들이 $S_n$의 부분군임이 증명된다.
(\hyperref[II.2.1.9-ko]{2.1.9})의 다른 조건들도 $\mathfrak{q}_0$가
소아이디얼임을 고려하면 자명하게 확인된다. 이렇게 정의한 $\mathfrak{p}$는 $S$의
동차 소아이디얼이고 실제로
$\psi_f(\mathfrak{p})=\mathfrak{q}_0$이다. 왜냐하면 $x\in S_{nd}$일 때
두 관계 $x/f^n\in\mathfrak{q}_0$와
$x^d/f^{nd}\in\mathfrak{q}_0$는 $\mathfrak{q}_0$가 소아이디얼이므로
서로 동치이기 때문이다.
\end{proof}

\begin{corollary}[2.3.7]
\phantomsection
\label{II.2.3.7-ko}
$D_+(f)=\varnothing$이기 위한 필요충분조건은 $f$가 멱영원인 것이다.
\end{corollary}

\begin{proof}
실제로 $\operatorname{Spec}(S_{(f)})=\varnothing$이기 위한 필요충분조건은
$S_{(f)}=0$인 것이며, 이는 다시 $1=0$이 $S_f$에서 성립하는 것과 동치이다.
정의에 따라 이는 $f$가 멱영원이라는 뜻이다.
\end{proof}

\begin{corollary}[2.3.8]
\phantomsection
\label{II.2.3.8-ko}
$E$를 $S_+$의 부분집합이라 하자. 다음 조건들은 서로 동치이다.
\begin{enumerate}
  \item[\rm a)] $V_+(E)=X=\operatorname{Proj}(S)$.
  \item[\rm b)] $E$의 모든 원소는 멱영원이다.
  \item[\rm c)] $E$의 각 원소의 모든 동차 성분은 멱영원이다.
\end{enumerate}
\end{corollary}

\begin{proof}
c)가 b)를 함의하고 b)가 a)를 함의함은 명백하다. $\mathfrak{J}$가
$S$의 동차 아이디얼이고 $E$로 생성된다면, 조건 a)는
$V_+(\mathfrak{J})=X$와 동치이다. \emph{특히}, a)는 모든 동차 원소
$f\in\mathfrak{J}$가 $V_+(f)=X$를 만족함을
함의하므로, (\hyperref[II.2.3.7-ko]{2.3.7})에 따라 $f$는 멱영원이다.
\end{proof}

\begin{corollary}[2.3.9]
\phantomsection
\label{II.2.3.9-ko}
$\mathfrak{J}$가 $S_+$의 동차 아이디얼이면,
$r_+(\mathfrak{J})$는 $S_+$의 동차 소아이디얼들 가운데
$\mathfrak{J}$를 포함하는 것들의 교집합이다.
\end{corollary}

\begin{proof}
등급환 $S/\mathfrak{J}$를 고려하여 $\mathfrak{J}=0$인 경우로 환원할 수
있다. $f\in S_+$가 멱영원이 아니면 $S$의 동차 소아이디얼 가운데
$f$를 포함하지 않는 것이 존재함을 증명해야 한다. 그런데 $f$의 동차
성분 중 적어도 하나는 멱영원이 아니므로 $f$가 동차라고 가정할 수 있고,
이때 주장은 (\hyperref[II.2.3.7-ko]{2.3.7})에서 따른다.
\end{proof}

\begin{env}[2.3.10]
\phantomsection
\label{II.2.3.10-ko}
$Y$를 $X=\operatorname{Proj}(S)$의 임의의 부분집합이라 하자.
$\mathfrak{j}_+(Y)$를 모든 $f\in S_+$ 가운데 $Y\subset V_+(f)$를
만족하는 것들의 집합으로 둔다. 다시 말해,
$\mathfrak{j}_+(Y)=\mathfrak{j}(Y)\cap S_+$이다. 따라서
$\mathfrak{j}_+(Y)$는 $S_+$의 아이디얼이며, $S_+$에서 취한 그 근기와
같다.
\end{env}

\begin{proposition}[2.3.11]
\phantomsection
\label{II.2.3.11-ko}
\begin{enumerate}
  \item[\rm(i)] $E$를 $S_+$의 임의의 부분집합이라 하면,
    $\mathfrak{j}_+(V_+(E))$는 $S_+$에서 취한, $S_+$에서 $E$가 생성하는
    동차 아이디얼의 근기이다.
  \item[\rm(ii)] $Y$를 $X$의 임의의 부분집합이라 하면,
    $V_+(\mathfrak{j}_+(Y))=\overline{Y}$이다. 즉 우변은 $Y$의 폐포이며,
    이 폐포는 $X$ 안에서 취한다.
\end{enumerate}
\end{proposition}

\begin{proof}
\begin{enumerate}
  \item[\rm(i)] $\mathfrak{J}$가 $S_+$의 동차 아이디얼이고 $E$로
    생성된다면 $V_+(E)=V_+(\mathfrak{J})$이고, 이때 주장은
    (\hyperref[II.2.3.9-ko]{2.3.9})에서 따른다.
  \item[\rm(ii)] $V_+(\mathfrak{J})=\bigcap_{f\in\mathfrak{J}}V_+(f)$이므로,
    관계 $Y\subset V_+(\mathfrak{J})$는 $Y\subset V_+(f)$가 모든
    $f\in\mathfrak{J}$에 대하여 성립함을 함의한다. 따라서
    $\mathfrak{j}_+(Y)\supset\mathfrak{J}$이고, 이로부터
    $V_+(\mathfrak{j}_+(Y))\subset V_+(\mathfrak{J})$를 얻는다. 따라서
    닫힌 부분집합의 정의에 의해 (ii)가 증명된다.
\end{enumerate}
\end{proof}

\begin{corollary}[2.3.12]
\phantomsection
\label{II.2.3.12-ko}
$Y$가 $X=\operatorname{Proj}(S)$의 닫힌 부분집합일 때, 그러한 부분집합
전체와 $S_+$의 동차 아이디얼 가운데 $S_+$에서 취한 자신의 근기와 같은
것 전체는 포함관계를 역전하는 두 사상
$Y\longmapsto\mathfrak{j}_+(Y)$,
$\mathfrak{J}\longmapsto V_+(\mathfrak{J})$에 의해 서로 전단사로 대응한다.
또 $X$의 두 닫힌 부분집합의 합집합 $Y_1\cup Y_2$에는
\oldpage[II]{28}
$\mathfrak{j}_+(Y_1)\cap\mathfrak{j}_+(Y_2)$가 대응하고, 임의의 닫힌
부분집합족 $(Y_\lambda)$의 교집합에는 $S_+$에서
$\mathfrak{j}_+(Y_\lambda)$들의 합의 근기가 대응한다.
\end{corollary}

\begin{corollary}[2.3.13]
\phantomsection
\label{II.2.3.13-ko}
$\mathfrak{J}$를 $S_+$의 동차 아이디얼이라 하자.
$V_+(\mathfrak{J})=\varnothing$이기 위한 필요충분조건은 $S_+$의 각
원소마다 그 원소의 양의 정수 거듭제곱 중 하나가 $\mathfrak{J}$에 속하는
것이다.
\end{corollary}

마지막 따름정리는 다음과 같은 서로 동치인 형태들로도 표현된다.

\begin{corollary}[2.3.14]
\phantomsection
\label{II.2.3.14-ko}
$(f_\alpha)$를 $S_+$의 동차 원소족이라 하자. $D_+(f_\alpha)$들이
$X=\operatorname{Proj}(S)$의 덮개를 이루기 위한 필요충분조건은 $S_+$의
각 원소마다 그 원소의 양의 정수 거듭제곱 중 하나가 $f_\alpha$들이
생성하는 아이디얼에 속하는 것이다.
\end{corollary}

\begin{corollary}[2.3.15]
\phantomsection
\label{II.2.3.15-ko}
$(f_\alpha)$를 $S_+$의 동차 원소족, $f$를 $S_+$의 한 원소라 하자.
다음 관계들은 서로 동치이다.
\begin{enumerate}
  \item[\rm a)] $D_+(f)\subset\bigcup_\alpha D_+(f_\alpha)$.
  \item[\rm b)] $V_+(f)\supset\bigcap_\alpha V_+(f_\alpha)$.
  \item[\rm c)] $f$의 어떤 양의 정수 거듭제곱이 $f_\alpha$들이 생성하는
    아이디얼에 속한다.
\end{enumerate}
\end{corollary}

\begin{corollary}[2.3.16]
\phantomsection
\label{II.2.3.16-ko}
$X=\operatorname{Proj}(S)$가 공집합이기 위한 필요충분조건은 $S_+$의
모든 원소가 멱영원인 것이다.
\end{corollary}

\begin{corollary}[2.3.17]
\phantomsection
\label{II.2.3.17-ko}
(\hyperref[II.2.3.12-ko]{2.3.12})에서 서술한 전단사 대응에서 $X$의
기약 닫힌 부분집합들에는 $S_+$의 동차 소아이디얼들이 대응한다.
\end{corollary}

\begin{proof}
실제로 $Y=Y_1\cup Y_2$이고 $Y_1$, $Y_2$가 $Y$의 진부분집합인
닫힌집합들이면
\[
  \mathfrak{j}_+(Y)
  =\mathfrak{j}_+(Y_1)\cap\mathfrak{j}_+(Y_2)
\]
이며, 아이디얼 $\mathfrak{j}_+(Y_1)$과 $\mathfrak{j}_+(Y_2)$는 둘 다
$\mathfrak{j}_+(Y)$와 다르므로 $\mathfrak{j}_+(Y)$는 소아이디얼이
아니다. 역으로 $\mathfrak{J}$가 $S_+$의 소아이디얼이 아닌 동차
아이디얼이면 두 원소 $f$, $g$가 $S_+$에 존재하여
$f\notin\mathfrak{J}$, $g\notin\mathfrak{J}$, $fg\in\mathfrak{J}$를
만족한다. 그러면
$V_+(f)\not\supset V_+(\mathfrak{J})$,
$V_+(g)\not\supset V_+(\mathfrak{J})$이지만
(\hyperref[II.2.3.2.3-ko]{2.3.2.3})에 따라
$V_+(\mathfrak{J})\subset V_+(f)\cup V_+(g)$이다. 따라서
$V_+(\mathfrak{J})$는 닫힌집합
$V_+(f)\cap V_+(\mathfrak{J})$와
$V_+(g)\cap V_+(\mathfrak{J})$의 합집합이고, 이 두 닫힌집합은 모두
$V_+(\mathfrak{J})$와 다르다.
\end{proof}

\subsection{$\operatorname{Proj}(S)$의 스킴 구조}
\label{subsection:II.2.4-ko}

\begin{env}[2.4.1]
\phantomsection
\label{II.2.4.1-ko}
$f$, $g$를 $S_+$의 두 동차 원소라 하자. 아핀 스킴
$Y_f=\operatorname{Spec}(S_{(f)})$,
$Y_g=\operatorname{Spec}(S_{(g)})$,
$Y_{fg}=\operatorname{Spec}(S_{(fg)})$를 생각하자.
(\hyperref[II.2.2.2-ko]{2.2.2})에 따라, 사상
$w_{fg,f}=({}^a\omega_{fg,f},\widetilde{\omega}_{fg,f})$는 $Y_{fg}$에서
$Y_f$로 가며, 표준 준동형
$\omega_{fg,f}:S_{(f)}\to S_{(fg)}$에 대응하고,
\emph{열린 몰입 사상}이다
(\hyperref[I.1.3.6-ko]{I, 1.3.6}).
위상동형 $\psi_f:D_+(f)\to Y_f$의 역
(\hyperref[II.2.3.6-ko]{2.3.6})을 이용하여 $D_+(f)$에 $Y_f$의 아핀
스킴 구조를 옮길 수 있다. 도식
(\hyperref[II.2.3.5.1-ko]{2.3.5.1})의 가환성에 따라, 이렇게 얻은 아핀
스킴 $D_+(fg)$는 아핀 스킴 $D_+(f)$의 바탕공간의 열린 부분집합
$D_+(fg)$에 유도된 스킴과 동일시된다. 그러면
(\hyperref[II.2.3.4-ko]{2.3.4})를 고려할 때
$X=\operatorname{Proj}(S)$에는 유일한 \emph{준스킴} 구조가 주어지며,
그 구조를 각 $D_+(f)$에 제한하면 방금 정의한 아핀 스킴이 된다. 또한
다음이 성립한다.
\end{env}

\begin{proposition}[2.4.2]
\phantomsection
\label{II.2.4.2-ko}
준스킴 $\operatorname{Proj}(S)$는 스킴이다.
\end{proposition}

\begin{proof}
(\hyperref[I.5.5.6-ko]{I, 5.5.6})에 따라, $f$, $g$가 $S_+$의 임의의
동차 원소일 때 $D_+(f)\cap D_+(g)=D_+(fg)$가 아핀이고 그 환이
$D_+(f)$와 $D_+(g)$의 환들의 표준 상들에 의해 생성됨을 보이면
충분하다. 첫째 주장은 정의상 명백하고, 둘째 주장은
(\hyperref[II.2.2.4-ko]{2.2.4})에서 따른다.
\end{proof}
\oldpage[II]{29}
$\operatorname{Proj}(S)$라는 동차 소 스펙트럼을 \emph{스킴}이라고 할
때에는 언제나 방금 정의한 구조를 뜻한다.

\begin{example}[2.4.3]
\phantomsection
\label{II.2.4.3-ko}
$S=K[T_1,T_2]$라 하자. 여기서 $K$는 체이고, $T_1$, $T_2$는 두
부정원이며, $S$에는 전체 차수에 따른 등급이 주어져 있다.
(\hyperref[II.2.3.14-ko]{2.3.14})에 따라
$\operatorname{Proj}(S)$는 $D_+(T_1)$과 $D_+(T_2)$의 합집합이다.
이 두 아핀 스킴은 각각 $\operatorname{Spec}(K[T])$와 표준적으로
동형이며, $\operatorname{Proj}(S)$는
(\hyperref[I.2.3.2-ko]{I, 2.3.2})에서 서술한 대로 이 두 아핀 스킴을
붙여서 얻어진다((7.4.14) 참조).
\end{example}

\begin{proposition}[2.4.4]
\phantomsection
\label{II.2.4.4-ko}
$S$를 음이 아닌 차수로 등급이 주어진 환이라 하고, $X$를 스킴
$\operatorname{Proj}(S)$라 하자.
\begin{enumerate}
  \item[\rm(i)] $\mathfrak{N}_+$가 $S_+$의 멱영근기이면
    (\hyperref[II.2.1.10-ko]{2.1.10}), 스킴 $X_{\mathrm{red}}$는
    $\operatorname{Proj}(S/\mathfrak{N}_+)$와 표준적으로 동형이다. 특히
    $S$가 본질적으로 축소이면 $\operatorname{Proj}(S)$는 축소이다.
  \item[\rm(ii)] $S$가 본질적으로 축소라고 가정하자. 이때 $X$가 정역
    스킴이기 위한 필요충분조건은 $S$가 본질적으로 정역인 것이다.
\end{enumerate}
\end{proposition}

\begin{proof}
\begin{enumerate}
  \item[\rm(i)] $\overline{S}$를 등급환 $S/\mathfrak{N}_+$라 하고, 차수가
    $0$인 표준 준동형 $S\to\overline{S}$를
    $x\mapsto\overline{x}$로 나타내자. 모든 $f\in S_d$ ($d>0$)에 대하여
    표준 준동형
    $S_f\to\overline{S}_{\overline{f}}$
    (\hyperref[0.1.5.1-ko]{0, 1.5.1})는 전사이고 차수가 $0$이므로, 이를
    제한하면 전사 준동형
    $S_{(f)}\to\overline{S}_{(\overline{f})}$를 얻는다.
    $f\notin\mathfrak{N}_+$라고 가정하면
    $\overline{S}_{(\overline{f})}$가 축소환이고 앞의 준동형의 핵이
    $S_{(f)}$의 멱영근기임을 바로 확인할 수 있다. 다시 말해
    $\overline{S}_{(\overline{f})}=(S_{(f)})_{\mathrm{red}}$이다. 따라서 이
    준동형에는 $D_+(\overline{f})$를 $(D_+(f))_{\mathrm{red}}$와 동일시하는
    닫힌 몰입 사상
    $D_+(\overline{f})\to D_+(f)$가 대응한다
    (\hyperref[I.5.1.2-ko]{I, 5.1.2}). 특히 이 사상은 두 아핀 스킴의
    바탕공간 사이의 위상동형이다. 더욱이
    $g\notin\mathfrak{N}_+$가 $S_+$의 또 다른 동차 원소이면 도식
    \[
      \xymatrix{
        S_{(f)}\ar[r]\ar[d]
        & \overline{S}_{(\overline{f})}\ar[d]\\
        S_{(fg)}\ar[r]
        & \overline{S}_{(\overline{fg})}
      }
    \]
    은 가환한다. 한편 $D_+(f)$들은 $f$가 차수가 $>0$인 동차 원소이고
    $f\notin\mathfrak{N}_+$일 때
    $X=\operatorname{Proj}(S)$의 덮개를 이룬다
    (\hyperref[II.2.3.7-ko]{2.3.7}). 따라서 사상
    $D_+(\overline{f})\to D_+(f)$들은 닫힌 몰입 사상
    $\operatorname{Proj}(\overline{S})\to\operatorname{Proj}(S)$의
    제한들이고, 이 닫힌 몰입 사상은 바탕공간 사이의 위상동형이다. 이제
    (\hyperref[I.5.1.2-ko]{I, 5.1.2})에 의해 결론이 따른다.
  \item[\rm(ii)] $S$가 본질적으로 정역이라고 가정하자. 다시 말해
    $(0)$은 $S_+$와 다른 $S_+$의 등급 소아이디얼이다. 그러면 (i)에 따라
    $X$는 축소이고, (\hyperref[II.2.3.17-ko]{2.3.17})에 따라 기약이다.
    반대로 $S$가 본질적으로 축소이고 $X$가 정역 스킴이라고 가정하자.
    그러면 $S_+$의 동차 원소 $f\neq0$에 대하여
    $D_+(f)\neq\varnothing$이다
    (\hyperref[II.2.3.7-ko]{2.3.7}). $X$가 기약이라는 가정에 따라
    $S_+$의 동차 원소 $f$, $g$가 모두 $\neq0$이면
    $D_+(f)\cap D_+(g)\neq\varnothing$이다. 따라서
    (\hyperref[II.2.3.3.2-ko]{2.3.3.2})에 의해 $fg\neq0$이고, 이로부터
    $S_+$에는 영인자가 없음을 얻으므로 처음의 주장이 따른다.
\end{enumerate}
\end{proof}

\begin{env}[2.4.5]
\phantomsection
\label{II.2.4.5-ko}
$A$를 가환환이라 하자. 등급환 $S$에 각 부분군 $S_n$이 $A$-가군이
되도록 하는 $A$-대수 구조가 주어져 있으면 이를
\emph{등급 $A$-대수}라고 한다는 것을 상기하자. 이를 위해서는 $S_0$가
\oldpage[II]{30}
$A$-대수이면 충분하다. 다시 말해 $S_0$에 $A$-대수 구조를 정의하고
$\alpha\in A$, $x\in S_n$에 대하여
$\alpha\cdot x=(\alpha\cdot1)x$로 놓으면 $S$에 등급 $A$-대수 구조가
정의된다.
\end{env}

\begin{proposition}[2.4.6]
\phantomsection
\label{II.2.4.6-ko}
$S$가 등급 $A$-대수라고 가정하자. 그러면
$X=\operatorname{Proj}(S)$ 위에서 구조층 $\mathcal{O}_X$는
$A$-대수들의 층이다($A$는 $X$ 위의 상수층으로 본다). 다시 말해 $X$는
$\operatorname{Spec}(A)$ 위의 스킴이다.
\end{proposition}

\begin{proof}
$S_+$의 모든 동차 원소 $f$에 대하여 $S_{(f)}$가 $A$ 위의 대수이고,
$S_+$의 동차 원소 $f$, $g$에 대하여 도식
\[
  \xymatrix{
    S_{(f)}\ar[rr] && S_{(fg)}\\
    & A\ar[ul]\ar[ur] &
  }
\]
이 가환함에 유의하면 충분하다.
\end{proof}

\begin{proposition}[2.4.7]
\phantomsection
\label{II.2.4.7-ko}
$S$를 음이 아닌 차수로 등급이 주어진 환이라 하자.
\begin{enumerate}
  \item[\rm(i)] 모든 정수 $d>0$에 대하여 스킴
    $\operatorname{Proj}(S)$에서 스킴
    $\operatorname{Proj}(S^{(d)})$로 가는 표준 동형이 존재한다.
  \item[\rm(ii)] $S'$를 다음 조건을 만족하는 등급환이라 하자:
    $S'_0=\mathbf{Z}$이고, $S'_n=S_n$은 ($\mathbf{Z}$-가군으로서)
    $n>0$일 때 성립한다. 그러면 스킴 $\operatorname{Proj}(S)$에서 스킴
    $\operatorname{Proj}(S')$로 가는 표준 동형이 존재한다.
\end{enumerate}
\end{proposition}

\begin{proof}
\begin{enumerate}
  \item[\rm(i)] 먼저 사상
    $\mathfrak{p}\mapsto\mathfrak{p}\cap S^{(d)}$가 집합
    $\operatorname{Proj}(S)$에서
    $\operatorname{Proj}(S^{(d)})$로 가는 전단사임을 보이자. 실제로 동차
    소아이디얼 $\mathfrak{p}'\in\operatorname{Proj}(S^{(d)})$가 주어졌다고
    하고 $\mathfrak{p}_{nd}=\mathfrak{p}'\cap S_{nd}$ ($n\geq0$)로 놓자.
    $n>0$이고 $d$의 배수가 아닌 모든 경우에 $\mathfrak{p}_n$을
    $x\in S_n$이고 $x^d\in\mathfrak{p}_{nd}$인 원소들의 집합으로 정의하자.
    $x\in\mathfrak{p}_n$, $y\in\mathfrak{p}_n$이면
    $(x+y)^{2d}\in\mathfrak{p}_{2nd}$이고, 따라서
    $(x+y)^d\in\mathfrak{p}_{nd}$이다. 이는 $\mathfrak{p}'$이
    소아이디얼이기 때문이다. 이와 같이 정의된 $\mathfrak{p}_n$들이 모든
    $n\geq0$에 대하여
    (\hyperref[II.2.1.9-ko]{2.1.9})의 조건을 만족함은 자명하다. 따라서
    유일한 소아이디얼 $\mathfrak{p}\in\operatorname{Proj}(S)$가 존재하며,
    이 소아이디얼은 $\mathfrak{p}\cap S^{(d)}=\mathfrak{p}'$을 만족한다. 한편 $f$가
    $S_+$의 동차 원소일 때마다
    $V_+(f)=V_+(f^d)$이므로
    (\hyperref[II.2.3.2.3-ko]{2.3.2.3}), 앞의 전단사는 위상공간들의
    위상동형이다. 마지막으로 같은 표기 아래에서 $S_{(f)}$와
    $S^{(d)}_{(f^d)}$는 표준적으로 동일시되므로
    (\hyperref[II.2.2.2-ko]{2.2.2}), $\operatorname{Proj}(S)$와
    $\operatorname{Proj}(S^{(d)})$도 스킴으로서 표준적으로 동일시된다.
  \item[\rm(ii)] 각 $\mathfrak{p}\in\operatorname{Proj}(S)$에 유일한
    소아이디얼 $\mathfrak{p}'\in\operatorname{Proj}(S')$을 대응시키되,
    $\mathfrak{p}'\cap S_n=\mathfrak{p}\cap S_n$이 모든 $n>0$에 대하여
    성립하게 하면
    (\hyperref[II.2.1.9-ko]{2.1.9}), 바탕공간들의 표준 위상동형
    $\operatorname{Proj}(S)\xrightarrow{\sim}\operatorname{Proj}(S')$을
    정의함이 분명하다. 실제로 $V_+(f)$는 $S$에 대해서나 $S'$에 대해서나
    같은 집합이다($f$가 $S_+$의 동차 원소일 때). 더욱이
    $S_{(f)}=S'_{(f)}$이므로 $\operatorname{Proj}(S)$와
    $\operatorname{Proj}(S')$은 스킴으로서 동일시된다.
\end{enumerate}
\end{proof}

\begin{corollary}[2.4.8]
\phantomsection
\label{II.2.4.8-ko}
$S$가 등급 $A$-대수이고, $S'_A$가 다음 조건을 만족하는 등급
$A$-대수라고 하자: $(S'_A)_0=A$, $(S'_A)_n=S_n$ ($n>0$).
$\operatorname{Proj}(S)$에서 $\operatorname{Proj}(S'_A)$로 가는 표준
동형이 존재한다.
\end{corollary}

\begin{proof}
실제로 (\hyperref[II.2.4.7-ko]{2.4.7, (ii)})의 표기 아래에서 이 두
스킴은 모두 $\operatorname{Proj}(S')$와 표준적으로 동형이다.
\end{proof}

\subsection{등급가군에 대응하는 층.}
\label{subsection:II.2.5-ko}

\begin{env}[2.5.1]
\phantomsection
\label{II.2.5.1-ko}
$M$을 등급 $S$-가군이라 하자. $f$가 $S_+$의 동차 원소이면
$M_{(f)}$는 $S_{(f)}$-가군이고, 따라서 이 가군에 대응하는 준연접층
$\widetilde{M_{(f)}}$가 아핀 스킴
$\operatorname{Spec}(S_{(f)})$ 위에 있다. 이 스킴은 $D_+(f)$와
동일시된다(\hyperref[I.1.3.4-ko]{I, 1.3.4}).
\end{env}
\oldpage[II]{31}

\begin{proposition}[2.5.2]
\phantomsection
\label{II.2.5.2-ko}
$X=\operatorname{Proj}(S)$ 위에 다음 조건을 만족하는 준연접
$\mathcal{O}_X$-가군 $\widetilde{M}$이 유일하게 존재한다. 즉 $f$가
$S_+$의 동차 원소일 때마다
$\Gamma(D_+(f),\widetilde{M})=M_{(f)}$이고, 제한 준동형
$\Gamma(D_+(f),\widetilde{M})\to
\Gamma(D_+(fg),\widetilde{M})$은 $f$, $g$가 $S_+$의 동차 원소일 때
표준 준동형 $M_{(f)}\to M_{(fg)}$이다
(\hyperref[II.2.2.3-ko]{2.2.3}).
\end{proposition}

\begin{proof}
$f\in S_d$, $g\in S_e$라 하자.
(\hyperref[II.2.2.2-ko]{2.2.2})에 따라 $D_+(fg)$는
$(S_{(f)})_{g^d/f^e}$의 소 스펙트럼과 동일시된다. 따라서
$D_+(fg)$로 제한한 층 $\widetilde{M_{(f)}}$---여기서 제한 전의 층은
$D_+(f)$ 위에 있다---은 가군 $(M_{(f)})_{g^d/f^e}$에 대응하는 층과 표준적으로
동일시되고(\hyperref[I.1.3.6-ko]{I, 1.3.6}), 따라서
$\widetilde{M_{(fg)}}$와도 동일시된다
(\hyperref[II.2.2.2-ko]{2.2.2}). 그러므로 표준 동형
\[
  \theta_{g,f}:\widetilde{M_{(f)}}|D_+(fg)
  \xrightarrow{\sim}\widetilde{M_{(g)}}|D_+(fg)
\]
이 존재하며, $h$가 $S_+$의 세 번째 동차 원소이면
$\theta_{f,h}=\theta_{f,g}\circ\theta_{g,h}$가 $D_+(fgh)$에서 성립한다.
따라서
(\hyperref[0.3.3.1-ko]{0, 3.3.1}) $X$ 위에 준연접
$\mathcal{O}_X$-가군 $\mathcal{F}$가 존재하고, 각 $f$가 $S_+$의 동차
원소일 때 동형 $\eta_f$가 $\mathcal{F}|D_+(f)$에서
$\widetilde{M_{(f)}}$로 가며 $\theta_{g,f}=\eta_g\circ\eta_f^{-1}$을
만족한다.
이제 층 $\mathcal{G}$를 생각하자. 이는 $X$의 위상의 기저인
$D_+(f)$들 위에서 $D_+(f)\to M_{(f)}$로 정의되고, 표준 준동형
$M_{(f)}\to M_{(fg)}$를 제한 준동형으로 갖는 준층을 층화하여 얻는
층이다. 앞의 논의로부터 $\mathcal{F}$와 $\mathcal{G}$가 동형임을 안다
((\hyperref[I.1.3.7-ko]{I, 1.3.7})도 고려한다). 층 $\mathcal{G}$를
$\widetilde{M}$으로 나타내면 이 층은 명제의 조건들을 만족한다. 특히
$\widetilde{S}=\mathcal{O}_X$이다.
\end{proof}

\begin{definition}[2.5.3]
\phantomsection
\label{II.2.5.3-ko}
(\hyperref[II.2.5.2-ko]{2.5.2})에서 정의한 준연접
$\mathcal{O}_X$-가군 $\widetilde{M}$을 등급 $S$-가군 $M$에
\emph{대응하는} 층이라 한다.
\end{definition}

등급 $S$-가군들은 등급가군의 준동형을 \emph{차수 $0$의} 준동형으로
제한하면 하나의 범주를 이룸을 상기하자. 이 약속 아래에서 다음이
성립한다.

\begin{proposition}[2.5.4]
\phantomsection
\label{II.2.5.4-ko}
함자 $M\mapsto\widetilde{M}$은 등급 $S$-가군의 범주에서 준연접
$\mathcal{O}_X$-가군의 범주로 가는 공변 가법 완전 함자이고,
귀납극한 및 직합과 가환한다.
\end{proposition}

\begin{proof}
실제로 이 성질들은 국소적이므로 층
$\widetilde{M}|D_+(f)=\widetilde{M_{(f)}}$에 대하여 확인하면 충분하다.
그런데 함자 $M\mapsto M_f$, 함자 $N\mapsto N_0$ (등급
$S_f$-가군의 범주에서) 및 함자 $P\mapsto\widetilde{P}$
($S_{(f)}$-가군의 범주에서)는 모두 완전성과 귀납극한 및 직합과의
가환성을 갖는다(\hyperref[I.1.3.5-ko]{I, 1.3.5} 및
\hyperref[I.1.3.9-ko]{I, 1.3.9}). 이로부터 명제가 따른다.
\end{proof}

$\widetilde{u}$로 준동형 $\widetilde{M}\to\widetilde{N}$을 나타내기로
한다. 이는 차수 $0$의 준동형 $u:M\to N$에 대응하는 준동형이다.
(\hyperref[II.2.5.4-ko]{2.5.4})에서 곧바로, 등급 $S$-가군과 차수
$0$의 준동형에 대해서도 (여기서 정한 $\widetilde{M}$의 뜻으로)
(\hyperref[I.1.3.9-ko]{I, 1.3.9} 및
\hyperref[I.1.3.10-ko]{I, 1.3.10})의 결과들이 여전히 성립함을 알 수
있다. 실제로 그 증명들은 순전히 형식적이다.

\begin{proposition}[2.5.5]
\phantomsection
\label{II.2.5.5-ko}
모든 $\mathfrak{p}\in X=\operatorname{Proj}(S)$에 대하여
$\widetilde{M}_{\mathfrak{p}}=M_{(\mathfrak{p})}$이다.
\end{proposition}

\begin{proof}
실제로 정의에 따라
$\widetilde{M}_{\mathfrak{p}}=
\varinjlim\Gamma(D_+(f),\widetilde{M})$이다. 여기서 $f$의 범위는
$S_+$의 동차 원소 가운데 $f\notin\mathfrak{p}$인 원소 전체이다. 한편
$\Gamma(D_+(f),\widetilde{M})=M_{(f)}$이므로, 명제는
$M_{(\mathfrak{p})}$의 정의
(\hyperref[II.2.2.7-ko]{2.2.7})에서 따른다.
\end{proof}
\oldpage[II]{32}

특히 \emph{국소환 $(\mathcal{O}_X)_\mathfrak{p}$}은 바로 환
$S_{(\mathfrak{p})}$이다. 이는 분수 $x/f$들의 집합으로, 여기서 $f$는
$S_+$의 동차 원소이고 $\mathfrak{p}$에 속하지 않으며, $x$는 $f$와
같은 차수의 동차 원소이다.

더욱 구체적으로, $S$가 \emph{본질적으로 정역}이어서
$\operatorname{Proj}(S)=X$가 \emph{정역 스킴}이고
(\hyperref[II.2.4.4-ko]{2.4.4}), $\xi=(0)$가 $X$의 일반점이면,
\emph{유리함수체} $R(X)=\mathcal{O}_\xi$는 $fg^{-1}$ 꼴의 원소들로 이루어진
체이다. 여기서 $f$와 $g$는 $S_+$에 속하는 같은 차수의 동차 원소이고
$g\neq0$이다.

\begin{proposition}[2.5.6]
\phantomsection
\label{II.2.5.6-ko}
모든 $z\in M$에 대하여, $f$가 $S_+$의 동차 원소일 때마다 $f$의 어떤
거듭제곱이 $z$를 소멸시키면, $\widetilde{M}=0$이다. 이 충분조건은
$S_0$-대수 $S$가 집합 $S_1$에 의해 생성되고, 이 집합이 차수 $1$의
동차 원소 전체로 이루어질 때에는 필요조건이기도 하다.
\end{proposition}

\begin{proof}
실제로 조건 $\widetilde{M}=0$은 $M_{(f)}=0$이 $f$가 $S_+$의 동차
원소일 때마다 성립하는 것과 동치이다. 한편 $f\in S_d$일 때
$M_{(f)}=0$이라는 것은 모든 동차 $z\in M$ 가운데 차수가 $d$의 배수인
각 원소에 대하여 어떤 거듭제곱 $f^n$이 존재하여 $f^nz=0$임을 뜻한다.
$M_{(f)}=0$이 모든 $f\in S_1$에 대해 성립하면 명제의 조건은 모든
$f\in S_1$에 대하여 성립한다. 또한 앞의 생성 가정 아래에서는 $f$가
$S_+$의 동차 원소일 때마다 이 조건이 \emph{더 나아가}
성립한다. 이 생성 가정은 $S_1$이 $S$를 생성한다는 것이며, 실제로
$S_+$의 모든 동차 원소는 $S_1$의 원소들의 곱들의 선형결합이기 때문이다.
\end{proof}

\begin{proposition}[2.5.7]
\phantomsection
\label{II.2.5.7-ko}
$d$를 $>0$인 정수라 하고 $f\in S_d$라 하자. 그러면 모든
$n\in\mathbf{Z}$에 대하여 $(\mathcal{O}_X|D_+(f))$-가군
$\widetilde{S(nd)}|D_+(f)$는 $\mathcal{O}_X|D_+(f)$와 표준적으로
동형이다.
\end{proposition}

\begin{proof}
실제로 $(f/1)^n$은 $S_f$의 가역원이며, 이를 곱하는 사상은
$S_{(f)}=(S_f)_0$에서
$(S_f)_{nd}=(S_f(nd))_0=(S(nd)_f)_0=S(nd)_{(f)}$로 가는 전단사이다.
다시 말해 $S_{(f)}$-가군 $S_{(f)}$와 $S(nd)_{(f)}$는 표준적으로
동형이고, 따라서 명제가 성립한다.
\end{proof}

\begin{corollary}[2.5.8]
\phantomsection
\label{II.2.5.8-ko}
열린부분집합
\[
  U=\bigcup_{f\in S_d}D_+(f),
\]
위에서 $\mathcal{O}_X$-가군 $\widetilde{S(nd)}$의 제한은 가역인
$(\mathcal{O}_X|U)$-가군이다
(\hyperref[0.5.4.1-ko]{0, 5.4.1}).
\end{corollary}

\begin{corollary}[2.5.9]
\phantomsection
\label{II.2.5.9-ko}
아이디얼 $S_+$가 $S$에서 집합 $S_1$에 의해 생성되고, 이 집합이 차수
$1$의 동차 원소 전체로 이루어지면, $\mathcal{O}_X$-가군
$\widetilde{S(n)}$은 모든 $n\in\mathbf{Z}$에 대하여 가역이다.
\end{corollary}

\begin{proof}
$X=\bigcup_{f\in S_1}D_+(f)$임은 가정과
(\hyperref[II.2.3.14-ko]{2.3.14})에서 따른다. 이제 $U=X$로 두고
(\hyperref[II.2.5.8-ko]{2.5.8})을 적용하면 충분하다.
\end{proof}

\begin{env}[2.5.10]
\phantomsection
\label{II.2.5.10-ko}
이 절의 나머지 부분 전체에서
\[
\label{II.2.5.10.1-ko}
  \mathcal{O}_X(n)=\widetilde{S(n)}
\tag{2.5.10.1}
\]
로 두기로 한다. 이는 모든 $n\in\mathbf{Z}$에 대한 약속이다. 또한
$U$를 $X$의 임의의 열린부분집합이라 하고, 임의의
$(\mathcal{O}_X|U)$-가군 $\mathcal{F}$에 대하여
\[
\label{II.2.5.10.2-ko}
  \mathcal{F}(n)=
  \mathcal{F}\otimes_{\mathcal{O}_X|U}(\mathcal{O}_X(n)|U)
\tag{2.5.10.2}
\]
로 두기로 한다. 이는 모든 $n\in\mathbf{Z}$에 대한 약속이다. 아이디얼
$S_+$가 $S_1$에 의해 생성되면, $\mathcal{F}(n)$을 대응시키는 함자는
$\mathcal{F}$에 관해 모든 $n\in\mathbf{Z}$에 대하여 \emph{완전}하다. 이 경우
$\mathcal{O}_X(n)$이 가역인 $\mathcal{O}_X$-가군이기 때문이다.
\end{env}

\begin{env}[2.5.11]
\phantomsection
\label{II.2.5.11-ko}
$M$과 $N$을 두 등급 $S$-가군이라 하자. 모든 $f\in S_d$ ($d>0$)에
대하여 다음과 같은 $S_{(f)}$-가군의 함자적인 표준 준동형을 정의한다.
\[
\label{II.2.5.11.1-ko}
  \lambda_f:M_{(f)}\otimes_{S_{(f)}}N_{(f)}
  \longrightarrow(M\otimes_S N)_{(f)}
\tag{2.5.11.1}
\]
\oldpage[II]{33}

정의는 준동형
$M_{(f)}\otimes_{S_{(f)}}N_{(f)}\to M_f\otimes_{S_f}N_f$
---이는 단사 준동형들 $M_{(f)}\to M_f$, $N_{(f)}\to N_f$ 및
$S_{(f)}\to S_f$에서 유도된다---과 표준 동형
$M_f\otimes_{S_f}N_f\xrightarrow{\sim}(M\otimes_S N)_f$
(\hyperref[0.1.3.4-ko]{0, 1.3.4})을 합성하여 얻는다. 또한 두 등급가군의
텐서곱의 정의에 따라 뒤의 동형이 차수를 보존한다는 점을 사용한다. 따라서
$x\in M_{md}$, $y\in N_{nd}$ ($m\geq0$, $n\geq0$)에 대하여
\[
  \lambda_f((x/f^m)\otimes(y/f^n))=(x\otimes y)/f^{m+n}.
\]

이 정의에서 곧바로, $g\in S_e$ ($e>0$)이면 다음 도식이 가환함을 알 수
있다.
\[
  \xymatrix{
    M_{(f)}\otimes_{S_{(f)}}N_{(f)}\ar[r]^{\lambda_f}\ar[d]
    & (M\otimes_S N)_{(f)}\ar[d]\\
    M_{(fg)}\otimes_{S_{(fg)}}N_{(fg)}\ar[r]_{\lambda_{fg}}
    & (M\otimes_S N)_{(fg)}
  }
\]
여기서 오른쪽 세로 화살표는 표준 준동형이고 왼쪽 세로 화살표는 표준
준동형들에서 유도된다. 따라서 이들 $\lambda$로부터
$\mathcal{O}_X$-가군의 함자적인 표준 준동형
\[
\label{II.2.5.11.2-ko}
  \lambda:\widetilde{M}\otimes_{\mathcal{O}_X}\widetilde{N}
  \longrightarrow\widetilde{M\otimes_S N}.
\tag{2.5.11.2}
\]
을 얻는다.

특히 두 동차 아이디얼 $\mathfrak{J}$, $\mathfrak{K}$를 생각하자. 이들은
$S$의 아이디얼이다. $\widetilde{\mathfrak{J}}$와
$\widetilde{\mathfrak{K}}$는 $\mathcal{O}_X$의 아이디얼층이므로 표준 준동형
$\widetilde{\mathfrak{J}}\otimes_{\mathcal{O}_X}
\widetilde{\mathfrak{K}}\to\mathcal{O}_X$가 존재하며, 다음 도식은 가환한다.
\[
\label{II.2.5.11.3-ko}
  \xymatrix{
    \widetilde{\mathfrak{J}}\otimes_{\mathcal{O}_X}
    \widetilde{\mathfrak{K}}\ar[rr]^\lambda\ar[dr]
    && \widetilde{\mathfrak{J}\otimes_S\mathfrak{K}}\ar[dl]\\
    & \mathcal{O}_X &
  }
\tag{2.5.11.3}
\]
실제로 각 열린집합 $D_+(f)$ ($f$는 $S_+$의 동차 원소)에서 이를 확인하는
경우로 귀착할 수 있고, 이는 $\lambda_f$의 정의
(\hyperref[II.2.5.11.1-ko]{2.5.11.1})와
(\hyperref[I.1.3.13-ko]{I, 1.3.13})에서 곧바로 따른다.

끝으로 $M$, $N$, $P$를 세 등급 $S$-가군이라 하면 다음 도식을 얻는다.
\[
\label{II.2.5.11.4-ko}
  \xymatrix{
    \widetilde{M}\otimes_{\mathcal{O}_X}\widetilde{N}
    \otimes_{\mathcal{O}_X}\widetilde{P}
    \ar[r]^{\lambda\otimes1}\ar[d]_{1\otimes\lambda}
    & \widetilde{M\otimes_S N}\otimes_{\mathcal{O}_X}\widetilde{P}
    \ar[d]^\lambda\\
    \widetilde{M}\otimes_{\mathcal{O}_X}\widetilde{N\otimes_S P}
    \ar[r]_\lambda
    & \widetilde{M\otimes_S N\otimes_S P}
  }
\tag{2.5.11.4}
\]
\oldpage[II]{34}

이 도식은 가환한다. 다시 각 $D_+(f)$에서 확인하면 충분하며, 이는 정의와
(\hyperref[I.1.3.13-ko]{I, 1.3.13})에서 곧바로 따른다.
\end{env}

\begin{env}[2.5.12]
\phantomsection
\label{II.2.5.12-ko}
(\hyperref[II.2.5.11-ko]{2.5.11})의 가정 아래에서 다음과 같은
$S_{(f)}$-가군의 함자적인 표준 준동형을 정의한다.
\[
\label{II.2.5.12.1-ko}
  \mu_f:\bigl(\operatorname{Hom}_S(M,N)\bigr)_{(f)}
  \longrightarrow
  \operatorname{Hom}_{S_{(f)}}(M_{(f)},N_{(f)})
\tag{2.5.12.1}
\]
이는 $u/f^n$---여기서 $u$는 차수 $nd$인 준동형이다---에 준동형
$M_{(f)}\to N_{(f)}$을 대응시켜 정의한다. 이 준동형은 $x/f^m$
($x\in M_{md}$)을 $u(x)/f^{m+n}$으로 보낸다. $g\in S_e$ ($e>0$)에
대해서도 다음 도식은 가환한다.
\[
  \xymatrix{
    \bigl(\operatorname{Hom}_S(M,N)\bigr)_{(f)}
    \ar[r]^{\mu_f}\ar[d]
    & \operatorname{Hom}_{S_{(f)}}(M_{(f)},N_{(f)})\ar[d]\\
    \bigl(\operatorname{Hom}_S(M,N)\bigr)_{(fg)}
    \ar[r]_{\mu_{fg}}
    & \operatorname{Hom}_{S_{(fg)}}(M_{(fg)},N_{(fg)})
  }
\]
여기서 왼쪽 화살표는 표준 준동형이고 오른쪽 화살표는 표준 준동형들에서
유도된다. 따라서 (\hyperref[I.1.3.8-ko]{I, 1.3.8})을 고려하면
$\mu_f$들이 다음과 같은 $\mathcal{O}_X$-가군의 함자적인 표준 준동형을
정의함을 다시 얻는다.
\[
\label{II.2.5.12.2-ko}
  \mu:\widetilde{\operatorname{Hom}_S(M,N)}
  \longrightarrow
  \mathcal{H}om_{\mathcal{O}_X}(\widetilde{M},\widetilde{N}).
\tag{2.5.12.2}
\]
\end{env}

\begin{proposition}[2.5.13]
\phantomsection
\label{II.2.5.13-ko}
아이디얼 $S_+$가 $S_1$에 의해 생성된다고 가정하자. 그러면 준동형
$\lambda$(\hyperref[II.2.5.11.2-ko]{2.5.11.2})는 동형사상이고, 준동형
$\mu$(\hyperref[II.2.5.12.2-ko]{2.5.12.2})도 등급 $S$-가군 $M$이
유한 표시를 가지는 경우(\hyperref[II.2.1.1-ko]{2.1.1})에는 동형사상이다.
\end{proposition}

\begin{proof}
$X$가 $D_+(f)$ ($f\in S_1$)들의 합집합이므로
(\hyperref[II.2.3.14-ko]{2.3.14}), 고려하는 각각의 가정 아래에서
$\lambda_f$와 $\mu_f$가 $f$가 동차이고 \emph{차수가 $1$}일 때
동형사상임을 보이는 것으로 귀착된다. 이때 $\mathbf{Z}$-쌍선형 사상
\[
  M_m\times N_n\longrightarrow
  M_{(f)}\otimes_{S_{(f)}}N_{(f)}
\]
을 $(x,y)$에 원소 $(x/f^m)\otimes(y/f^n)$을 대응시켜 정의한다
($m<0$이면 $x/f^m$은 $f^{-m}x/1$을 뜻한다). 이 사상들로
$\mathbf{Z}$-선형 사상
\[
  M\otimes_{\mathbf{Z}}N\longrightarrow
  M_{(f)}\otimes_{S_{(f)}}N_{(f)},
\]
을 정의할 수 있다. $s\in S_q$이면 이 사상은 $(sx)\otimes y$를
\[
  (s/f^q)\bigl((x/f^m)\otimes(y/f^n)\bigr)
\]
으로 보낸다($x\in M_m$, $y\in N_n$). 따라서 가군의 쌍준동형
\[
  \gamma_f:M\otimes_S N\longrightarrow
  M_{(f)}\otimes_{S_{(f)}}N_{(f)},
\]
을 얻는다. 이는 표준 준동형 $S\to S_{(f)}$---$s\in S_q$를
$s/f^q$로 보내는 준동형---에 관한 것이다. 더 나아가 원소
$\sum_i(x_i\otimes y_i)$가 $M\otimes_S N$에 속하고, 여기서
$x_i$, $y_i$는 각각 차수가
$m_i$, $n_i$인 동차 원소이다---에 대하여
\[
  f^r\sum_i(x_i\otimes y_i)=0,
\]
즉 $\sum_i(f^rx_i\otimes y_i)=0$이라고 가정하자.
(\hyperref[0.1.3.4-ko]{0, 1.3.4})에 따라
\[
  \sum_i(f^rx_i/f^{m_i+r})\otimes(y_i/f^{n_i})=0,
\]
을 얻는다. 다시 말해 $\gamma_f(\sum_i(x_i\otimes y_i))=0$이다. 따라서
$\gamma_f$는 다음과 같이 인수분해된다.
\[
  M\otimes_S N\longrightarrow(M\otimes_S N)_f
  \xrightarrow{\gamma'_f}
  M_{(f)}\otimes_{S_{(f)}}N_{(f)}~;
\]
$\lambda'_f$를 $\gamma'_f$의
\oldpage[II]{35}
$(M\otimes_S N)_{(f)}$로의 제한이라 하면,
$\lambda_f$와 $\lambda'_f$가 서로 역인 $S_{(f)}$-준동형임을 곧바로
확인할 수 있다. 이로써 명제의 첫 번째 부분이 증명된다.

두 번째 부분을 증명하기 위해, $M$이 준동형 $P\to Q$의 여핵이라고
가정하자. 이는 등급 $S$-가군의 준동형이며, $P$와 $Q$는 $S(n)$ 꼴인
가군들의 유한 직합이다. $\operatorname{Hom}_S(L,N)$의 $L$에 관한
왼쪽 완전성과 $M_{(f)}$의 $M$에 관한 완전성을 사용하면, $\mu_f$가
$M=S(n)$일 때 동형사상임을 보이는 것으로 곧바로 귀착된다. 이제 임의의
동차 원소 $z$를 $N$에서 택하고, $u_z$를 $S(n)$에서 $N$으로 가며
$u_z(1)=z$를 만족하는 준동형이라 하자. 그러면
$\eta:z\mapsto u_z$는 차수 $0$이고 $N(-n)$에서
$\operatorname{Hom}_S(S(n),N)$으로 가는 동형사상임을 곧바로
알 수 있다. 이에 대응하는 동형사상
\[
  \eta_f:\bigl(N(-n)\bigr)_{(f)}
  \longrightarrow
  \bigl(\operatorname{Hom}_S(S(n),N)\bigr)_{(f)}.
\]
이 있다. 한편 $\eta'_f$를 다음 동형사상이라 하자.
\[
  N_{(f)}\longrightarrow
  \operatorname{Hom}_{S_{(f)}}\bigl(S(n)_{(f)},N_{(f)}\bigr)
\]
이는 각 $z'\in N_{(f)}$에 준동형 $v_{z'}$를 대응시키며, 이 준동형은
$v_{z'}(s/f^k)=sz'/f^{n+k}$를 만족한다
($s\in S_{n+k}=(S(n))_k$). 합성 사상
\[
  \bigl(N(-n)\bigr)_{(f)}
  \xrightarrow{\eta_f}
  \bigl(\operatorname{Hom}_S(S(n),N)\bigr)_{(f)}
  \xrightarrow{\mu_f}
  \operatorname{Hom}_{S_{(f)}}\bigl(S(n)_{(f)},N_{(f)}\bigr)
  \xrightarrow{{\eta'_f}^{-1}}N_{(f)}
\]
은 $z/f^h\mapsto z/f^{h-n}$으로 주어지는
$\bigl(N(-n)\bigr)_{(f)}$에서 $N_{(f)}$로 가는 동형사상임을 쉽게 확인할 수
있다. 따라서 $\mu_f$는 동형사상이다.
\end{proof}

아이디얼 $S_+$가 $S_1$에 의해 생성되면,
(\hyperref[II.2.5.13-ko]{2.5.13})에서 $S$의 모든 등급 아이디얼
$\mathfrak{J}$와 모든 등급 $S$-가군 $M$에 대하여
\[
\label{II.2.5.13.1-ko}
  \widetilde{\mathfrak{J}}\cdot\widetilde{M}
  =\widetilde{\mathfrak{J}\cdot M}
\tag{2.5.13.1}
\]
가 표준 동형까지 성립함을 얻는다. 실제로 이는 다음 도식의 가환성에서
따른다.
\[
  \xymatrix{
    \widetilde{\mathfrak{J}}\otimes_{\mathcal{O}_X}\widetilde{M}
    \ar[rr]^\lambda\ar[dr]
    && \widetilde{\mathfrak{J}\otimes_S M}\ar[dl]\\
    & \widetilde{M} &
  }
\]
이 가환성은 (\hyperref[II.2.5.11.3-ko]{2.5.11.3})의 경우와 마찬가지로
확인한다.

\begin{corollary}[2.5.14]
\phantomsection
\label{II.2.5.14-ko}
$S$가 $S_1$에 의해 생성된다고 하자. 어떠한 $m,n\in\mathbf{Z}$에
대해서도 다음 공식들이 성립한다.
\[
\label{II.2.5.14.1-ko}
  \mathcal{O}_X(m)\otimes_{\mathcal{O}_X}\mathcal{O}_X(n)
  =\mathcal{O}_X(m+n)
\tag{2.5.14.1}
\]
\[
\label{II.2.5.14.2-ko}
  \mathcal{O}_X(n)=\bigl(\mathcal{O}_X(1)\bigr)^{\otimes n}
\tag{2.5.14.2}
\]
여기서 등식들은 표준 동형까지 성립한다.
\end{corollary}

\begin{proof}
첫째 공식은 (\hyperref[II.2.5.13-ko]{2.5.13})과 차수 $0$인 표준 동형
$S(m)\otimes_S S(n)\xrightarrow{\sim}S(m+n)$의 존재에서 따른다. 이 동형은
$1\otimes1$---여기서 첫째 인자 $1$은 $(S(m))_{-m}$에, 둘째 인자 $1$은
$(S(n))_{-n}$에 속한다---을 원소 $1\in(S(m+n))_{-m-n}$에 대응시킨다.
이제 둘째 공식은 $n=-1$인 경우만 증명하면 충분하다. 그리고
(\hyperref[II.2.5.13-ko]{2.5.13})에 의하면 이는
$\operatorname{Hom}_S(S(1),S)$가 $S(-1)$과 표준적으로 동형임을 보이는
것으로 귀착한다. 정의들 (\hyperref[II.2.1.2-ko]{2.1.2})을 되짚어 보고
$S(1)$이 한 원소로 생성되는 $S$-가군임을 상기하면 곧바로 확인된다.
\end{proof}

\oldpage[II]{36}
\begin{corollary}[2.5.15]
\phantomsection
\label{II.2.5.15-ko}
$S$가 $S_1$에 의해 생성된다고 하자. 모든 등급 $S$-가군 $M$과 모든
$n\in\mathbf{Z}$에 대하여
\[
\label{II.2.5.15.1-ko}
  \widetilde{M(n)}=\widetilde{M}(n)
\tag{2.5.15.1}
\]
가 표준 동형까지 성립한다.
\end{corollary}

\begin{proof}
이는 정의들 (\hyperref[II.2.5.10.2-ko]{2.5.10.2})와
(\hyperref[II.2.5.10.1-ko]{2.5.10.1}), 명제
(\hyperref[II.2.5.13-ko]{2.5.13}), 그리고 차수 $0$인 표준 동형
$M(n)\xrightarrow{\sim}M\otimes_S S(n)$의 존재에서 따른다. 이 동형은 모든
$z\in(M(n))_h=M_{n+h}$를
$z\otimes1\in M_{n+h}\otimes(S(n))_{-n}\subset(M\otimes_S S(n))_h$에
대응시킨다.
\end{proof}

\begin{env}[2.5.16]
\phantomsection
\label{II.2.5.16-ko}
$S'_0=\mathbf{Z}$이고 $n>0$에 대하여 $S'_n=S_n$인 등급환을 $S'$라 하자.
그러면 $f\in S_d$ ($d>0$)일 때 모든 $n\in\mathbf{Z}$에 대하여
$(S(n))_{(f)}=(S'(n))_{(f)}$이다. 실제로 $(S'(n))_{(f)}$의 원소는
$x\in S'_{n+kd}$ ($k>0$)인 $x/f^k$ 꼴이며, 항상 $n+kd\neq0$이 되도록
$k$를 택할 수 있다. $X=\operatorname{Proj}(S)$와
$X'=\operatorname{Proj}(S')$가 표준적으로 동일시되므로
(\hyperref[II.2.4.7-ko]{2.4.7, (ii)}), 모든 $n\in\mathbf{Z}$에 대하여
$\mathcal{O}_X(n)$과 $\mathcal{O}_{X'}(n)$은 앞의 동일시에서 서로
대응한다.

한편 모든 $d>0$과 모든 $n\in\mathbf{Z}$에 대하여
\[
  (S^{(d)}(n))_h=S_{(n+h)d}=(S(nd))_{hd}
\]
이며, $f\in S_d$이면 따라서
$(S^{(d)}(n))_{(f)}=(S(nd))_{(f)}$이다. 스킴
$X=\operatorname{Proj}(S)$와 $X^{(d)}=\operatorname{Proj}(S^{(d)})$가
표준적으로 동일시됨을 알고 있다
(\hyperref[II.2.4.7-ko]{2.4.7, (i)}). 앞의 논의에 따르면 $S_0$-대수
$S^{(d)}$가 $S_d$에 의해 생성될 때, 모든 $n\in\mathbf{Z}$에 대하여
$\mathcal{O}_X(nd)$와 $\mathcal{O}_{X^{(d)}}(n)$은 이 동일시에서 서로
대응한다.
\end{env}

\begin{proposition}[2.5.17]
\phantomsection
\label{II.2.5.17-ko}
$d$를 $>0$인 정수라 하고
$U=\bigcup_{f\in S_d}D_+(f)$라 하자. 표준 준동형
$\mathcal{O}_X(nd)\otimes_{\mathcal{O}_X}\mathcal{O}_X(-nd)
\to\mathcal{O}_X$의 $U$에 대한 제한은 모든 정수 $n$에 대하여
동형이다.
\end{proposition}

\begin{proof}
(\hyperref[II.2.5.16-ko]{2.5.16})에 의해 $d=1$인 경우로 환원할 수 있고,
그 경우 결론은 (\hyperref[II.2.5.13-ko]{2.5.13})의 증명에서 따른다.
\end{proof}

\subsection{$\operatorname{Proj}(S)$ 위의 층에 대응하는 등급
$S$-가군.}
\label{subsection:II.2.6-ko}

\emph{이 절 전체에서 $S$의 아이디얼 $S_+$가 차수 $1$의 동차 원소들의
집합 $S_1$에 의해 생성된다고 가정한다.}

\begin{env}[2.6.1]
\phantomsection
\label{II.2.6.1-ko}
그러면 $\mathcal{O}_X$-가군 $\mathcal{O}_X(1)$은
\emph{가역}이다(\hyperref[II.2.5.9-ko]{2.5.9}). 이때 모든
$\mathcal{O}_X$-가군 $\mathcal{F}$
(\hyperref[0.5.4.6-ko]{0, 5.4.6})에 대하여
\[
\label{II.2.6.1.1-ko}
  \Gamma_*(\mathcal{F})
  =\Gamma_*(\mathcal{O}_X(1),\mathcal{F})
  =\bigoplus_{n\in\mathbf{Z}}\Gamma(X,\mathcal{F}(n))
\tag{2.6.1.1}
\]
로 둔다. 여기서는 (\hyperref[II.2.5.14.2-ko]{2.5.14.2})를 고려하였다.
(\hyperref[0.5.4.6-ko]{0, 5.4.6})에서 상기한 바와 같이
$\Gamma_*(\mathcal{O}_X)$는 \emph{등급환}의 구조를 가지며,
$\Gamma_*(\mathcal{F})$는 $\Gamma_*(\mathcal{O}_X)$ 위의
\emph{등급가군} 구조를 가진다.

$\mathcal{O}_X(n)$이 국소 자유이므로
$\mathcal{F}\mapsto\Gamma_*(\mathcal{F})$는 $\mathcal{F}$에 관한
공변 가법 \emph{왼쪽 완전} 함자이다. 특히 $\mathcal{J}$가
$\mathcal{O}_X$의 아이디얼층이면, $\Gamma_*(\mathcal{J})$는
$\Gamma_*(\mathcal{O}_X)$의 \emph{등급 아이디얼}과 표준적으로
동일시된다.
\end{env}

\begin{env}[2.6.2]
\phantomsection
\label{II.2.6.2-ko}
$M$을 등급 $S$-가군이라 하자. 모든 $f\in S_d$ ($d>0$)에 대하여
$x\mapsto x/1$은 아벨군의 준동형 $M_0\to M_{(f)}$이고,
$M_{(f)}$는 표준적으로

\oldpage[II]{37}

$\Gamma(D_+(f),\widetilde{M})$와 동일시되므로 아벨군의 준동형
\[
  \alpha_0^f:M_0\longrightarrow\Gamma(D_+(f),\widetilde{M}).
\]
을 얻는다. 모든 $g\in S_e$ ($e>0$)에 대하여 다음 도식이 가환함은
명백하다.
\[
  \xymatrix{
    & \Gamma(D_+(f),\widetilde{M})\ar[dd]\\
    M_0\ar[ur]^{\alpha_0^f}\ar[dr]_{\alpha_0^{fg}}\\
    & \Gamma(D_+(fg),\widetilde{M})
  }
\]
이는 모든 $x\in M_0$에 대하여 $\widetilde{M}$의 단면
$\alpha_0^f(x)$와 $\alpha_0^g(x)$가 $D_+(f)\cap D_+(g)$에서
일치함을 뜻한다. 따라서 각 $D_+(f)$로의 제한이 $\alpha_0^f(x)$인
유일한 단면 $\alpha_0(x)\in\Gamma(X,\widetilde{M})$이 존재한다.
이로써 ($S$가 $S_1$에 의해 생성된다는 가정을 사용하지 않고) 아벨군의
준동형
\[
\label{II.2.6.2.1-ko}
  \alpha_0:M_0\longrightarrow\Gamma(X,\widetilde{M}).
\tag{2.6.2.1}
\]
을 정의하였다.

이 결과를 모든 $n\in\mathbf{Z}$에 대하여 등급 $S$-가군 $M(n)$에
적용하면, 각 $n\in\mathbf{Z}$에 대하여 아벨군의 준동형
\[
\label{II.2.6.2.2-ko}
  \alpha_n:M_n=(M(n))_0\longrightarrow\Gamma(X,\widetilde{M}(n))
\tag{2.6.2.2}
\]
을 얻는다((\hyperref[II.2.5.15-ko]{2.5.15})를 고려한다). 따라서 등급
아벨군의 함자적인 준동형(차수 $0$)
\[
\label{II.2.6.2.3-ko}
  \alpha:M\longrightarrow\Gamma_*(\widetilde{M})
\tag{2.6.2.3}
\]
을 얻으며, 이를 $\alpha_M$으로도 나타낸다. 이 준동형은 각 $M_n$에서
$\alpha_n$과 일치한다.

특히 $M=S$로 두면, $\Gamma_*(\mathcal{O}_X)$의 곱셈의 정의
(\hyperref[0.5.4.6-ko]{0, 5.4.6})를 고려하여
$\alpha:S\to\Gamma_*(\mathcal{O}_X)$가 등급환의 준동형임을 곧바로
확인할 수 있다. 또한 모든 등급 $S$-가군 $M$에 대하여
(\hyperref[II.2.6.2.3-ko]{2.6.2.3})은 등급 가군의 쌍준동형이다.
\end{env}

\begin{proposition}[2.6.3]
\phantomsection
\label{II.2.6.3-ko}
모든 $f\in S_d$ ($d>0$)에 대하여, $D_+(f)$는 $\mathcal{O}_X(d)$의
단면 $\alpha_d(f)$가 소멸하지 않는 $\mathfrak{p}\in X$들의 집합과
일치한다(\hyperref[0.5.5.2-ko]{0, 5.5.2}).
\end{proposition}

\begin{proof}
가정에 따라 $X=\bigcup_{g\in S_1}D_+(g)$이므로, 모든 $g\in S_1$에
대하여 $\alpha_d(f)$가 소멸하지 않는 $\mathfrak{p}\in D_+(g)$들의
집합이 $D_+(fg)$와 일치함을 보이면 충분하다. 그런데 $\alpha_d(f)$의
$D_+(g)$에 대한 제한은 정의상 $(S(d))_{(g)}$의 원소 $f/1$에 대응하는
단면이다. 표준 동형
$(S(d))_{(g)}\xrightarrow{\sim}S_{(g)}$
(\hyperref[II.2.5.7-ko]{2.5.7}) 아래에서, $D_+(g)$ 위의
$\mathcal{O}_X(d)$의 이 단면은 $S_{(g)}$의 원소 $f/g^d$에 대응하는
$D_+(g)$ 위의 $\mathcal{O}_X$의 단면과 동일시된다. 이 단면이
$\mathfrak{p}\in D_+(g)$에서 소멸한다는 것은 $f/g^d\in\mathfrak{q}$임을
뜻한다. 여기서 $\mathfrak{q}$는 $\mathfrak{p}$에 대응하는 $S_{(g)}$의
소아이디얼이다(\hyperref[II.2.3.6-ko]{2.3.6}). 정의상 이는
$f\in\mathfrak{p}$라는 뜻이므로 명제가 성립한다.
\end{proof}

\begin{env}[2.6.4]
\phantomsection
\label{II.2.6.4-ko}
이제 $\mathcal{F}$를 $\mathcal{O}_X$-가군이라 하고
$M=\Gamma_*(\mathcal{F})$로 두자. 등급환 준동형
$\alpha:S\to\Gamma_*(\mathcal{O}_X)$가 존재하므로, $M$을 등급
$S$-가군으로 볼 수 있다. 모든 $f\in S_d$ ($d>0$)에 대하여
(\hyperref[II.2.6.3-ko]{2.6.3})에서 $\mathcal{O}_X(d)$의 단면
$\alpha_d(f)$의 $D_+(f)$에 대한 제한이 가역임을 알 수 있다. 따라서
모든 $n>0$에 대하여 $\mathcal{O}_X(nd)$의 단면
$\alpha_{nd}(f^n)$\footnote{\textbf{번역 주.} 인쇄 원문은 이 문단의 세
곳에서 각각 $\alpha_d(f^n)$, $\alpha_d(f^k)$, $\alpha_d(f^n)$을
쓰지만, $f^n\in S_{nd}$와 $f^k\in S_{kd}$이므로
EGA2-CANON-DISP-0002에 따라 각각 $\alpha_{nd}(f^n)$,
$\alpha_{kd}(f^k)$, $\alpha_{nd}(f^n)$으로 교정했다.}의 $D_+(f)$에
대한 제한도 가역이다. 이제
$z\in M_{nd}=\Gamma(X,\mathcal{F}(nd))$ ($n>0$)라 하자. 어떤 정수
$k\geq0$에 대하여 $f^kz$의 $D_+(f)$에 대한 제한, 즉
\oldpage[II]{38}
단면
$(z|D_+(f))(\alpha_{kd}(f^k)|D_+(f))$가
$\mathcal{F}((n+k)d)$의 단면으로서 영이면, 앞의 설명에 따라
$z|D_+(f)=0$이기도 하다. 이로써 원소 $z/f^n$에 $D_+(f)$ 위의
$\mathcal{F}$의 단면
$(z|D_+(f))(\alpha_{nd}(f^n)|D_+(f))^{-1}$을 대응시켜
$S_{(f)}$-가군 준동형
$\beta_f:M_{(f)}\to\Gamma(D_+(f),\mathcal{F})$를 정의할 수 있다.
또한 $g\in S_e$ ($e>0$)일 때 다음 도식이 가환함을 곧바로 확인할 수
있다.
\[
\label{II.2.6.4.1-ko}
  \xymatrix{
    M_{(f)}\ar[r]^{\beta_f}\ar[d]
    & \Gamma(D_+(f),\mathcal{F})\ar[d]\\
    M_{(fg)}\ar[r]_{\beta_{fg}}
    & \Gamma(D_+(fg),\mathcal{F})
  }
\tag{2.6.4.1}
\]
$M_{(f)}$가 $\Gamma(D_+(f),\widetilde{M})$과 표준적으로 동일시되고
$D_+(f)$들이 $X$의 위상의 기저를 이룬다는 사실
(\hyperref[II.2.3.4-ko]{2.3.4})을 상기하면, $\beta_f$들이 유일한 표준
$\mathcal{O}_X$-가군 준동형
\[
\label{II.2.6.4.2-ko}
  \beta:\widetilde{\Gamma_*(\mathcal{F})}
  \longrightarrow\mathcal{F}
\tag{2.6.4.2}
\]
에서 유도됨을 알 수 있다. 이 준동형은 $\beta_{\mathcal{F}}$로도
표기하며 명백히 함자적이다.
\end{env}

\begin{proposition}[2.6.5]
\phantomsection
\label{II.2.6.5-ko}
$M$을 등급 $S$-가군, $\mathcal{F}$를 $\mathcal{O}_X$-가군이라 하자.
그러면 다음 합성 준동형들은 각각 항등 준동형이다.
\[
\label{II.2.6.5.1-ko}
  \widetilde{M}\xrightarrow{\widetilde{\alpha}}
  \widetilde{\Gamma_*(\widetilde{M})}
  \xrightarrow{\beta}\widetilde{M}
\tag{2.6.5.1}
\]
\[
\label{II.2.6.5.2-ko}
  \Gamma_*(\mathcal{F})\xrightarrow{\alpha}
  \Gamma_*(\widetilde{\Gamma_*(\mathcal{F})})
  \xrightarrow{\Gamma_*(\beta)}\Gamma_*(\mathcal{F})
\tag{2.6.5.2}
\]
\end{proposition}

\begin{proof}
(\hyperref[II.2.6.5.1-ko]{2.6.5.1})의 확인은 국소적이다. 열린집합
$D_+(f)$에서는 정의들과, 준연접 층에 적용한 $\beta$가 $D_+(f)$ 위의
단면들에 대한 작용으로 정해진다는 사실
(\hyperref[I.1.3.8-ko]{I, 1.3.8})에서 곧바로 따른다.
(\hyperref[II.2.6.5.2-ko]{2.6.5.2})의 확인은 각 차수에서 따로 이루어진다.
$M=\Gamma_*(\mathcal{F})$로 두면
$M_n=\Gamma(X,\mathcal{F}(n))$이고
$(\Gamma_*(\widetilde{M}))_n=\Gamma(X,\widetilde{M}(n))
=\Gamma(X,\widetilde{M(n)})$이다. 그런데 $f\in S_1$이고 $z\in M_n$이면,
$\alpha_n^f(z)$는 $(M(n))_{(f)}$의 원소 $z/1$이며
$(f/1)^n(z/f^n)$과 같다. 이에 $\beta_f$로 대응하는 것은 $D_+(f)$ 위의
단면
\[
  ((\alpha_1(f))^n|D_+(f))
  ((z|D_+(f))((\alpha_1(f))^n|D_+(f))^{-1})
\]
이다. 이는 곧 $z$의 $D_+(f)$에 대한 제한이므로
(\hyperref[II.2.6.5.2-ko]{2.6.5.2})가 확인된다.
\end{proof}

\subsection{유한성 조건.}
\label{subsection:II.2.7-ko}

\begin{proposition}[2.7.1]
\phantomsection
\label{II.2.7.1-ko}
\begin{enumerate}
  \item[\rm(i)] $S$가 뇌터인 등급환이면
    $X=\operatorname{Proj}(S)$는 뇌터 스킴이다.
  \item[\rm(ii)] $S$가 유한형 등급 $A$-대수이면
    $X=\operatorname{Proj}(S)$는 $Y=\operatorname{Spec}(A)$ 위의
    유한형 스킴이다.
\end{enumerate}
\end{proposition}

\oldpage[II]{39}
\begin{proof}
\begin{enumerate}
  \item[\rm(i)] $S$가 뇌터이면 아이디얼 $S_+$는 동차 원소들
    $f_i$ ($1\leq i\leq p$)로 이루어진 유한 생성계를 갖는다. 따라서
    (\hyperref[II.2.3.14-ko]{2.3.14}) 바탕공간 $X$는
    $D_+(f_i)=\operatorname{Spec}(S_{(f_i)})$들의 합집합이고, 각
    $S_{(f_i)}$가 뇌터임을 보이면 충분하다. 이는
    (\hyperref[II.2.2.6-ko]{2.2.6})에서 따른다.
  \item[\rm(ii)] 가정으로부터 $S_0$가 유한형 $A$-대수이고 $S$가
    유한형 $S_0$-대수임이 따른다. 따라서 $S_+$는 유한 생성
    아이디얼이다(\hyperref[II.2.1.4-ko]{2.1.4}). 그러면 (i)에서와 같이
    $f\in S_d$일 때 $S_{(f)}$가 유한형 $A$-대수임을 증명하면
    충분하다. (\hyperref[II.2.2.5-ko]{2.2.5})에 따라 $S^{(d)}$가 유한형
    $A$-대수임을 보이면 충분하고, 이는
    (\hyperref[II.2.1.6-ko]{2.1.6})에서 따른다.
\end{enumerate}
\end{proof}

\begin{env}[2.7.2]
\phantomsection
\label{II.2.7.2-ko}
이하에서는 등급 $S$-가군 $M$에 관한 다음 유한성 조건들을 고려하겠다.
\begin{enumerate}
  \item[\rm(TF)] 어떤 정수 $n$이 존재하여 부분가군
    $\bigoplus_{k\geq n}M_k$가 유한 생성 $S$-가군이다.
  \item[\rm(TN)] 어떤 정수 $n$이 존재하여 $k\geq n$이면 $M_k=0$이다.
\end{enumerate}

$M$이 (TN)을 만족하면 $S_+$의 모든 동차 원소 $f$에 대하여
$M_{(f)}=0$이고, 따라서 $\widetilde{M}=0$이다.

$M$, $N$을 두 등급 $S$-가군이라 하자. 차수 $0$의 준동형
$u:M\to N$에 대하여, 어떤 정수 $n$이 존재하여 $k\geq n$일 때
$u_k:M_k\to N_k$가 단사(각각 전사, 전단사)이면 $u$를
(TN)-\emph{단사}(각각 (TN)-\emph{전사}, (TN)-\emph{전단사})라고 한다.
따라서 $u$가 (TN)-단사(각각 (TN)-전사)라는 것은
$\operatorname{Ker}u$(각각 $\operatorname{Coker}u$)가 (TN)을
만족한다는 것과 같다. (\hyperref[II.2.5.4-ko]{2.5.4})에 따라 $u$가
(TN)-단사(각각 (TN)-전사, (TN)-전단사)이면 $\widetilde{u}$는
단사(각각 전사, 전단사)이다. $u$가 (TN)-전단사일 때 $u$를
(TN)-\emph{동형사상}이라고도 한다.
\end{env}

\begin{proposition}[2.7.3]
\phantomsection
\label{II.2.7.3-ko}
$S$를 아이디얼 $S_+$가 유한 생성인 등급환이라 하고, $M$을 등급
$S$-가군이라 하자.
\begin{enumerate}
  \item[\rm(i)] $M$이 조건 (TF)를 만족하면 $\mathcal{O}_X$-가군
    $\widetilde{M}$은 유한 생성이다.
  \item[\rm(ii)] $M$이 (TF)를 만족한다고 가정하자. $\widetilde{M}=0$일
    필요충분조건은 $M$이 (TN)을 만족하는 것이다.
\end{enumerate}
\end{proposition}

\begin{proof}
방금 본 바와 같이 조건 (TN)은 $\widetilde{M}=0$을 함의한다. $M$이
(TF)를 만족하면, 가정에 따라 유한 생성인 등급 부분가군
$M'=\bigoplus_{k\geq n}M_k$에 대하여 $M/M'$은 (TN)을 만족한다. 따라서
$\widetilde{M/M'}=0$이고, 함자 $M\mapsto\widetilde{M}$의 완전성
(\hyperref[II.2.5.4-ko]{2.5.4})으로부터
$\widetilde{M}=\widetilde{M'}$을 얻는다. 그러므로 $\widetilde{M}$이 유한
생성임을 증명하기 위해 $M$이 유한 생성인 경우로 귀착할 수 있다. 문제는
국소적이므로 $M_{(f)}$가 유한 생성 $S_{(f)}$-가군임을 보이면 충분하다
(\hyperref[I.1.3.9-ko]{I, 1.3.9}). 그런데 $M^{(d)}$는 유한 생성
$S^{(d)}$-가군이고(\hyperref[II.2.1.6-ko]{2.1.6, (iii)}), 이제 주장은
(\hyperref[II.2.2.5-ko]{2.2.5})에서 따른다.

이제 $M$이 (TF)를 만족하고 $\widetilde{M}=0$이라고 가정하자. 위와 같은
표기에서 $\widetilde{M'}=0$이고, $M'$에 대한 조건 (TN)은 $M$에 대한
조건 (TN)과 동치이다. 따라서 $\widetilde{M}=0$이 $M$의 조건 (TN)을
함의함을 증명하기 위해, 다시 $M$이 유한 개의 동차 원소 $x_i$
($1\leq i\leq p$)에 의해 생성되는 경우만 고려하면 된다. 또
$(f_j)_{1\leq j\leq q}$를 아이디얼 $S_+$의 동차 생성계라 하자. 가정에
따라 모든 $j$에 대하여 $M_{(f_j)}=0$이므로, 어떤 정수 $n$이 존재하여
모든 $i$, $j$에 대하여 $f_j^nx_i=0$이다. $n_j$를 $f_j$의 차수라 하고,
$\sum_jr_j\leq nq$를 만족하는 정수족 $(r_j)$에 대하여
$\sum_jr_jn_j$가 취하는 최댓값을 $m$이라 하자. 그러면

\oldpage[II]{40}
$k>m$이면 모든 $i$에 대하여 $S_kx_i=0$임이 명백하다. $x_i$들의 차수 중
최댓값을 $h$라 하면 $k>h+m$일 때 $M_k=0$이고, 이것으로 증명이 끝난다.
\end{proof}

\begin{corollary}[2.7.4]
\phantomsection
\label{II.2.7.4-ko}
$S$를 아이디얼 $S_+$가 유한 생성인 등급환이라 하자. 그러면
$X=\operatorname{Proj}(S)=\varnothing$이기 위한 필요충분조건은 어떤 $n$이
존재하여 $k\geq n$일 때 $S_k=0$인 것이다.
\end{corollary}

\begin{proof}
실제로 조건 $X=\varnothing$은
$\mathcal{O}_X=\widetilde{S}=0$과 동치이고, $S$는 한 원소로 생성되는
$S$-가군이다.
\end{proof}

\begin{theorem}[2.7.5]
\phantomsection
\label{II.2.7.5-ko}
아이디얼 $S_+$가 차수 $1$인 유한 개의 동차 원소로 생성된다고 가정하고,
$X=\operatorname{Proj}(S)$라 하자. 그러면 모든 준연접
$\mathcal{O}_X$-가군 $\mathcal{F}$에 대하여 표준 준동형
$\beta:\widetilde{\Gamma_*(\mathcal{F})}
\to\mathcal{F}$ (\hyperref[II.2.6.4-ko]{2.6.4})는 동형이다.
\end{theorem}

\begin{proof}
실제로 $S_+$가 유한 개의 원소 $f_i\in S_1$로 생성되면, $X$는 준콤팩트인
부분공간 $\operatorname{Spec}(S_{(f_i)})$
(\hyperref[II.2.3.6-ko]{2.3.6})들의 합집합이므로 준콤팩트이다. 또한 $X$는
스킴이다(\hyperref[II.2.4.2-ko]{2.4.2}).
(\hyperref[I.9.3.2-ko]{I, 9.3.2}),
(\hyperref[II.2.5.14.2-ko]{2.5.14.2}) 및
(\hyperref[II.2.6.3-ko]{2.6.3})에 따라 모든 $f\in S_d$ ($d>0$)에 대하여
표준 동형
$(\Gamma_*(\mathcal{F}))_{(\alpha_d(f))}\xrightarrow{\sim}
\Gamma(D_+(f),\mathcal{F})$을 얻는다. 더욱이
$(\Gamma_*(\mathcal{F}))_{(\alpha_d(f))}$는
$\Gamma_*(\mathcal{F})$를 $\Gamma_*(\mathcal{O}_X)$-가군으로 볼 때의
가군인데, 이는 $\Gamma_*(\mathcal{F})$를 $S$-가군으로 볼 때의
$(\Gamma_*(\mathcal{F}))_{(f)}$와 다름없다. 앞의 표준 동형의 정의
(\hyperref[I.9.3.1-ko]{I, 9.3.1})를 참조하면 이 동형이 준동형
$\beta_f$와 일치함을 알 수 있고, 정리가 따른다.
\end{proof}

\begin{remark}[2.7.6]
\phantomsection
\label{II.2.7.6-ko}
등급환 $S$가 \emph{뇌터}라고 가정하면, 아이디얼 $S_+$가 차수 $1$의 동차
원소 전체의 집합 $S_1$에 의해 생성된다고 가정하는 것만으로
(\hyperref[II.2.7.5-ko]{2.7.5})의 조건이 \emph{자동으로} 성립한다.
\end{remark}

\begin{corollary}[2.7.7]
\phantomsection
\label{II.2.7.7-ko}
(\hyperref[II.2.7.5-ko]{2.7.5})의 가정 아래에서 모든 준연접
$\mathcal{O}_X$-가군 $\mathcal{F}$는 $\widetilde{M}$ 꼴의
$\mathcal{O}_X$-가군과 동형이다. 여기서 $M$은 등급 $S$-가군이다.
\end{corollary}

\begin{corollary}[2.7.8]
\phantomsection
\label{II.2.7.8-ko}
(\hyperref[II.2.7.5-ko]{2.7.5})의 가정 아래에서 모든 유한 생성 준연접
$\mathcal{O}_X$-가군 $\mathcal{F}$는 $\widetilde{N}$ 꼴의
$\mathcal{O}_X$-가군과 동형이다. 여기서 $N$은 유한 생성 등급
$S$-가군이다.
\end{corollary}

\begin{proof}
(\hyperref[II.2.7.7-ko]{2.7.7})에 따라 $\mathcal{F}=\widetilde{M}$이라고
가정할 수 있다. 여기서 $M$은 등급 $S$-가군이다.
$(f_\lambda)_{\lambda\in L}$를 $M$의 동차 생성계라 하자. $L$의 모든 유한
부분집합 $H$에 대하여, $\lambda\in H$인 $f_\lambda$들로 생성되는 $M$의
등급 부분가군을 $M_H$라 하자. $M$이 부분가군 $M_H$들의 귀납극한임은
명백하므로, $\mathcal{F}$는 그 $\mathcal{O}_X$-부분가군
$\widetilde{M_H}$들의 귀납극한이다
(\hyperref[II.2.5.4-ko]{2.5.4}). 그런데 $\mathcal{F}$가 유한 생성이고
바탕공간 $X$가 준콤팩트이므로, (\hyperref[0.5.2.3-ko]{0, 5.2.3})에서
$L$의 어떤 유한 부분집합 $H$에 대하여
$\mathcal{F}=\widetilde{M_H}$임이 따른다.
\end{proof}

\begin{corollary}[2.7.9]
\phantomsection
\label{II.2.7.9-ko}
(\hyperref[II.2.7.5-ko]{2.7.5})의 가정 아래에서 $\mathcal{F}$를 유한 생성
준연접 $\mathcal{O}_X$-가군이라 하자. 그러면 어떤 정수 $n_0$가 존재하여,
모든 $n\geq n_0$에 대하여 $\mathcal{F}(n)$은 $\mathcal{O}_X^k$의 어떤
몫가군과 동형이다. 여기서 $k>0$은 $n$에 의존한다. 따라서
$\mathcal{F}(n)$은 $X$ 위의 유한 개 단면으로 생성된다
(\hyperref[0.5.1.1-ko]{0, 5.1.1}).
\end{corollary}

\begin{proof}
(\hyperref[II.2.7.8-ko]{2.7.8})에 따라 $\mathcal{F}=\widetilde{M}$이라고
가정할 수 있다. 여기서 $M$은 $S(m_i)$ 꼴의 $S$-가군들의 유한 직합의
몫가군이다. 따라서 (\hyperref[II.2.5.4-ko]{2.5.4})에 따라 $M=S(m)$인
경우로 귀착할 수 있고, 이때
$\mathcal{F}(n)=\widetilde{S(m+n)}
=\mathcal{O}_X(m+n)$이다
(\hyperref[II.2.5.15-ko]{2.5.15}). 그러므로 다음 보조정리를 증명하면
충분하다.

\oldpage[II]{41}
\begin{lemma}[2.7.9.1]
\phantomsection
\label{II.2.7.9.1-ko}
(\hyperref[II.2.7.5-ko]{2.7.5})의 가정 아래에서 모든 $n\geq0$에 대하여,
어떤 정수 $k$가 존재하여(이는 $n$에 의존한다) 전사 준동형
$\mathcal{O}_X^k\to\mathcal{O}_X(n)$이 존재한다.
\end{lemma}

실제로 (\hyperref[II.2.7.2-ko]{2.7.2})에 따라, 적당한 $k$에 대하여 등급
$S$-가군 곱 $S^k$에서 $S(n)$으로 가는 \emph{차수 $0$}의
(TN)-전사 준동형 $u$가 존재함을 보이면 충분하다. 그런데
$(S(n))_0=S_n$이고 가정에 따라 모든 $h>0$에 대하여 $S_h=S_1^h$이므로
$SS_n=S_n+S_{n+1}+\cdots$이다. $S_n$은 유한 생성 $S_0$-가군이므로
(\hyperref[II.2.1.5-ko]{2.1.5} 및
\hyperref[II.2.1.6-ko]{2.1.6, (i)}), 이 가군의 생성계
$(a_i)_{1\leq i\leq k}$를 택하자. 준동형 $u$는 $S^k$의 $i$번째 표준 기저
원소를 $a_i$에 대응시킨다($1\leq i\leq k$). 이때
$\operatorname{Coker}u$는
$(S(n))_{-n}+\cdots+(S(n))_{-1}$과 동일시되므로, $u$는 요구한 조건을
충족한다.
\end{proof}

\begin{corollary}[2.7.10]
\phantomsection
\label{II.2.7.10-ko}
(\hyperref[II.2.7.5-ko]{2.7.5})의 가정 아래에서, $\mathcal{F}$를 유한 생성
준연접 $\mathcal{O}_X$-가군이라 하자. 그러면 어떤 정수 $n_0$가 존재하여,
모든 $n\geq n_0$에 대하여 $\mathcal{F}$는
$(\mathcal{O}_X(-n))^k$ 꼴의 $\mathcal{O}_X$-가군의 몫과 동형이다
($k$는 $n$에 의존한다).
\end{corollary}

\begin{proposition}[2.7.11]
\phantomsection
\label{II.2.7.11-ko}
(\hyperref[II.2.7.5-ko]{2.7.5})의 가정이 성립한다고 하고, $M$을 등급
$S$-가군이라 하자. 그러면 다음이 성립한다.
\begin{enumerate}
  \item[\rm(i)] 표준 준동형
    $\widetilde{\alpha}:\widetilde{M}\to
    \widetilde{\Gamma_*(\widetilde{M})}$은 동형이다.
  \item[\rm(ii)] $\mathcal{G}$를 $\widetilde{M}$의 준연접
    $\mathcal{O}_X$-부분가군이라 하고, $N$을 $\alpha$에 의한
    $\Gamma_*(\mathcal{G})$의 역상인 $M$의 등급 $S$-부분가군이라 하자.
    그러면 $\widetilde{N}=\mathcal{G}$이다. 여기서
    (\hyperref[II.2.5.4-ko]{2.5.4})에 따라 $\widetilde{N}$을
    $\widetilde{M}$의 $\mathcal{O}_X$-부분가군으로 동일시한다.
\end{enumerate}
\end{proposition}

\begin{proof}
(\hyperref[II.2.7.5-ko]{2.7.5})에 따라
$\beta:\widetilde{\Gamma_*(\widetilde{M})}\to\widetilde{M}$가 동형이므로,
(\hyperref[II.2.6.5.1-ko]{2.6.5.1})에 따라 $\widetilde{\alpha}$는 그
역동형이다. 이로써 (i)가 따른다. $P$를 $\Gamma_*(\widetilde{M})$의 등급
부분가군 $\alpha(M)$이라 하자. $M\mapsto\widetilde{M}$은 완전함자이므로
(\hyperref[II.2.5.4-ko]{2.5.4}), $\widetilde{\alpha}$에 의한
$\widetilde{M}$의 상은 $\widetilde{P}$와 같다. 따라서 (i)에 따라
$\widetilde{P}=\widetilde{\Gamma_*(\widetilde{M})}$이다. 이제
$Q=\Gamma_*(\mathcal{G})\cap P$로 놓자. 앞의 결과와
(\hyperref[II.2.5.4-ko]{2.5.4})에 따라
$\widetilde{Q}=\widetilde{\Gamma_*(\mathcal{G})}$이므로,
(\hyperref[II.2.7.5-ko]{2.7.5})에 따라 $\beta$를 $\widetilde{Q}$에 제한한
사상은 이 $\mathcal{O}_X$-가군에서 $\mathcal{G}$ 위로 가는
\emph{동형}이다. 한편 $N$의 정의와
(\hyperref[II.2.5.4-ko]{2.5.4})에 따라, 동형 $\widetilde{\alpha}$를
$\widetilde{N}$에 제한한 사상은 $\widetilde{N}$에서
$\widetilde{Q}$ 위로 가는 동형이다. 따라서
(\hyperref[II.2.6.5.1-ko]{2.6.5.1})에 의해 결론이 따른다.
\end{proof}

\subsection{함자적 성질.}
\label{subsection:II.2.8-ko}

\begin{env}[2.8.1]
\phantomsection
\label{II.2.8.1-ko}
$S$, $S'$을 음이 아닌 차수로 등급이 주어진 두 환이라 하고,
$\varphi:S'\to S$를 등급환의 준동형이라 하자. $X=\operatorname{Proj}(S)$에서
$V_+(\varphi(S'_+))$의 여집합인 열린부분집합, 또는 같은 말로 $S'_+$의
동차 원소 $f'$ 전체에 대한 $D_+(\varphi(f'))$들의 합집합을
$G(\varphi)$로 나타내겠다. $\operatorname{Spec}(S)$에서
$\operatorname{Spec}(S')$로 가는 연속사상 ${}^a\varphi$
(\hyperref[I.1.2.1-ko]{I, 1.2.1})를 $G(\varphi)$에 제한하면
$G(\varphi)$에서 $\operatorname{Proj}(S')$로 가는 연속사상이 된다. 표기의
남용으로 이 사상도 ${}^a\varphi$로 나타내겠다. $f'\in S'_+$가 동차
원소이면
\[
\label{II.2.8.1.1-ko}
  {}^a\varphi^{-1}(D_+(f'))=D_+(\varphi(f'))
\tag{2.8.1.1}
\]
이다. 이는 ${}^a\varphi$가 $G(\varphi)$를 $\operatorname{Proj}(S')$로
보낸다는 사실과 (\hyperref[I.1.2.2.2-ko]{I, 1.2.2.2})에서 따른다. 한편
같은 표기 아래에서 준동형 $\varphi$는 등급환의 준동형
$S'_{f'}\to S_f$를 표준적으로 정의하며, 이를 차수 $0$인 원소들에
제한하면 다음 준동형을 얻는다.

\oldpage[II]{42}
$S'_{(f')}\to S_{(f)}$; 이를 $\varphi_{(f)}$로 나타내겠다. 이에 대응하여
(\hyperref[I.1.6.1-ko]{I, 1.6.1}) 아핀 스킴의 사상
$({}^a\varphi_{(f)},\widetilde{\varphi}_{(f)}):
\operatorname{Spec}(S_{(f)})\to\operatorname{Spec}(S'_{(f')})$을 얻는다.
$\operatorname{Spec}(S_{(f)})$를 $D_+(f)$ 위에서
$\operatorname{Proj}(S)$로부터 유도된 스킴과 표준적으로 동일시하면
(\hyperref[II.2.3.6-ko]{2.3.6}), 이로써 사상
$\Phi_f:D_+(f)\to D_+(f')$를 정의하였고, ${}^a\varphi_{(f)}$는
${}^a\varphi$를 $D_+(f)$에 제한한 사상과 동일시된다. 한편 $g'$이
$S'_+$의 두 번째 동차 원소이고 $g=\varphi(g')$이면 다음 도식이
가환임을 곧 알 수 있다.
\[
  \xymatrix{
    D_+(f)\ar[r]^{\Phi_f} & D_+(f')\\
    D_+(fg)\ar[u]\ar[r]_{\Phi_{fg}} & D_+(f'g')\ar[u]
  }
\]
이는 다음 도식이 가환하기 때문이다.
\[
  \xymatrix{
    S'_{(f')}\ar[r]^{\varphi_{(f)}}\ar[d]_{\omega_{f'g',f'}}
    & S_{(f)}\ar[d]^{\omega_{fg,f}}\\
    S'_{(f'g')}\ar[r]_{\varphi_{(fg)}} & S_{(fg)}
  }
\]
$G(\varphi)$의 정의와 (\hyperref[II.2.3.3.2-ko]{2.3.3.2})을 고려하면
따라서 다음을 알 수 있다.
\end{env}

\begin{proposition}[2.8.2]
\phantomsection
\label{II.2.8.2-ko}
등급환 준동형 $\varphi:S'\to S$가 주어지면, 유도 준스킴
$G(\varphi)$에서 $\operatorname{Proj}(S')$로 가는 유일한 사상
$({}^a\varphi,\widetilde{\varphi})$가 존재한다. 이 사상을
\emph{$\varphi$에 대응하는 사상}이라 하고 $\operatorname{Proj}(\varphi)$로
나타낸다. 모든 동차 원소 $f'\in S'_+$에 대하여, 이 사상을
$D_+(\varphi(f'))$에 제한한 사상은 $\varphi$가 정하는 준동형
$S'_{(f')}\to S_{(\varphi(f'))}$에 대응하는 사상과 일치한다.
\end{proposition}

\begin{proof}
앞의 표기에서 $f'\in S'_d$이면 다음 도식이 가환함에 유의하자.
\[
\label{II.2.8.2.1-ko}
  \xymatrix{
    S'_{(f')}\ar[r]^{\varphi_{(f)}}\ar[d]_{\sim}
    & S_{(f)}\ar[d]^{\sim}\\
    {S'}^{(d)}/(f'-1){S'}^{(d)}\ar[r]
    & S^{(d)}/(f-1)S^{(d)}
  }
\tag{2.8.2.1}
\]
여기서 세로 화살표들은
(\hyperref[II.2.2.5-ko]{2.2.5})의 동형들이다.
\end{proof}

\begin{corollary}[2.8.3]
\phantomsection
\label{II.2.8.3-ko}
\begin{enumerate}
  \item[\rm(i)] 사상 $\operatorname{Proj}(\varphi)$는 아핀이다.
  \item[\rm(ii)] $\operatorname{Ker}(\varphi)$가 멱영이면(특히
    $\varphi$가 단사이면), 사상 $\operatorname{Proj}(\varphi)$는
    우세 사상이다.
\end{enumerate}
\end{corollary}

\begin{proof}
(i)은 (\hyperref[II.2.8.2-ko]{2.8.2})와
(\hyperref[II.2.8.1.1-ko]{2.8.1.1})에서 곧바로 따른다. 한편
$\operatorname{Ker}(\varphi)$가 멱영이면, 모든 동차 원소
$f'\in S'_+$에 대하여 $f=\varphi(f')$로 둘 때
$\operatorname{Ker}(\varphi_f)$가 멱영임을 곧바로 확인할 수 있고,
따라서 $\operatorname{Ker}(\varphi_{(f)})$도 멱영이다. 그러므로 결론은
(\hyperref[II.2.8.2-ko]{2.8.2})와
(\hyperref[I.1.2.7-ko]{I, 1.2.7})에서 따른다.
\end{proof}

\oldpage[II]{43}
일반적으로 $\operatorname{Proj}(S)$에서 $\operatorname{Proj}(S')$로 가는
사상들 가운데 아핀이 아닌 것이 있으며, 따라서 이러한 사상은 등급환 준동형
$S'\to S$에서 나오지 않는다. 예를 들어 $A$가 체일 때의 구조 사상
$\operatorname{Proj}(S)\to\operatorname{Spec}(A)$이 그러하다. 여기서는
$\operatorname{Spec}(A)$를 $\operatorname{Proj}(A[T])$와 동일시한다
(\agref{II.3.1.7-ko}{3.1.7} 참조). 실제로 이는
(\hyperref[I.2.3.2-ko]{I, 2.3.2})에서 따른다.

\begin{env}[2.8.4]
\phantomsection
\label{II.2.8.4-ko}
$S''$을 음이 아닌 차수로 등급이 주어진 세 번째 환이라 하고,
$\varphi':S''\to S'$을 등급환 준동형이라 하며,
$\varphi''=\varphi\circ\varphi'$으로 놓자.
(\hyperref[II.2.8.1.1-ko]{2.8.1.1})과 공식
${}^a\varphi''=({}^a\varphi')\circ({}^a\varphi)$에 따라
$G(\varphi'')\subset G(\varphi)$임을 곧바로 확인할 수 있다. 또한
$\Phi$, $\Phi'$, $\Phi''$이 각각 $\varphi$, $\varphi'$, $\varphi''$에
대응하는 사상이면
$\Phi''=\Phi'\circ(\Phi|G(\varphi''))$이다.
\end{env}

\begin{env}[2.8.5]
\phantomsection
\label{II.2.8.5-ko}
$S$가 등급 $A$-대수이고 $S'$이 등급 $A'$-대수라고 하자. 또
$\psi:A'\to A$를 다음 도식이 가환하도록 하는 환 준동형이라 하자.
\[
\label{II.2.8.5.1-ko}
  \xymatrix{
    A'\ar[r]^{\psi}\ar[d] & A\ar[d]\\
    S'\ar[r]_{\varphi} & S
  }
\tag{2.8.5.1}
\]
그러면 $G(\varphi)$와 $\operatorname{Proj}(S')$을 각각
$\operatorname{Spec}(A)$와 $\operatorname{Spec}(A')$ 위의 스킴으로 볼 수
있다. $\Phi$와 $\Psi$가 각각 $\varphi$와 $\psi$에 대응하는 사상이면
다음 도식은 가환한다.
\[
  \xymatrix{
    G(\varphi)\ar[r]^{\Phi}\ar[d] & \operatorname{Proj}(S')\ar[d]\\
    \operatorname{Spec}(A)\ar[r]_{\Psi} & \operatorname{Spec}(A')
  }
\]
실제로 $f=\varphi(f')$이고 $f'$이 $S'_+$의 동차 원소일 때, $\Phi$를
$D_+(f)$에 제한한 사상에 대해 이를 보이면 충분하다. 이 경우에는 다음
도식의 가환성에서 결론이 따른다.
\[
  \xymatrix{
    A'\ar[r]^{\psi}\ar[d] & A\ar[d]\\
    S'_{(f')}\ar[r]_{\varphi_{(f)}} & S_{(f)}
  }
\]
\end{env}

\begin{env}[2.8.6]
\phantomsection
\label{II.2.8.6-ko}
이제 $M$을 등급 $S$-가군이라 하고, 자명하게 등급이 매겨지는
$S'$-가군 $M_{[\varphi]}$를 생각하자. $f'$을 $S'_+$의 동차 원소라 하고
$f=\varphi(f')$이라 하자. 표준 동형
$(M_{[\varphi]})_{f'}\xrightarrow{\sim}(M_f)_{[\varphi_f]}$이 존재함을
알고 있으며(\hyperref[0.1.5.2-ko]{0, 1.5.2}), 이 동형이 차수를
보존함은 자명하다. 따라서 표준 동형
$(M_{[\varphi]})_{(f')}\xrightarrow{\sim}
(M_{(f)})_{[\varphi_{(f)}]}$이 존재한다. 이 동형에 표준적으로 대응하는
층의 동형
$\widetilde{M_{[\varphi]}}|D_+(f')
\xrightarrow{\sim}(\Phi_f)_*(\widetilde{M}|D_+(f))$
이 존재한다
(\hyperref[II.2.5.2-ko]{2.5.2} 및
\hyperref[I.1.6.3-ko]{I, 1.6.3}). 더욱이,

\oldpage[II]{44}
$g'$가 $S'_+$의 두 번째 동차 원소이고 $g=\varphi(g')$이면 다음 도식은
가환한다.
\[
  \xymatrix{
    (M_{[\varphi]})_{(f')}\ar[r]^{\sim}\ar[d]
    & (M_{(f)})_{[\varphi_{(f)}]}\ar[d]\\
    (M_{[\varphi]})_{(f'g')}\ar[r]^{\sim}
    & (M_{(fg)})_{[\varphi_{(fg)}]}
  }
\]
따라서 동형
\[
  \widetilde{M_{[\varphi]}}|D_+(f'g')
  \xrightarrow{\sim}(\Phi_{fg})_*(\widetilde{M}|D_+(fg))
\]
은 동형
$\widetilde{M_{[\varphi]}}|D_+(f')
\xrightarrow{\sim}(\Phi_f)_*(\widetilde{M}|D_+(f))$을
$D_+(f'g')$로 제한한 것임을 곧바로 알 수 있다. $\Phi_f$는 사상
$\Phi$를 $D_+(f)$로 제한한 것이므로,
(\hyperref[II.2.8.1.1-ko]{2.8.1.1})을 고려하고
$X'=\operatorname{Proj}(S')$으로 두면 다음을 얻는다.
\end{env}

\begin{proposition}[2.8.7]
\phantomsection
\label{II.2.8.7-ko}
$\mathcal{O}_{X'}$-가군 $\widetilde{M_{[\varphi]}}$에서
$\mathcal{O}_{X'}$-가군 $\Phi_*(\widetilde{M}|G(\varphi))$로 가는
함자적인 표준 동형이 존재한다.
\end{proposition}

이로부터 등급 $S'$-가군 $M'$에서 등급 $S$-가군 $M$으로 가는
$\varphi$-사상들의 집합에서 $\Phi$-사상
$\widetilde{M'}\to\widetilde{M}|G(\varphi)$들의 집합으로 가는 함자적인
표준 사상을 곧바로 얻는다. (\hyperref[II.2.8.4-ko]{2.8.4})의 표기
아래에서 $M''$가 등급 $S''$-가군이면, $\varphi$-사상 $M'\to M$과
$\varphi'$-사상 $M''\to M'$의 합성에
$\widetilde{M'}|G(\varphi')\to\widetilde{M}|G(\varphi'')$와
$\widetilde{M''}\to\widetilde{M'}|G(\varphi')$의 합성이 표준적으로
대응한다.

\begin{proposition}[2.8.8]
\phantomsection
\label{II.2.8.8-ko}
(\hyperref[II.2.8.1-ko]{2.8.1})의 가정 아래에서 $M'$을 등급
$S'$-가군이라 하자. $(\mathcal{O}_X|G(\varphi))$-가군
$\Phi^*(\widetilde{M'})$에서 $(\mathcal{O}_X|G(\varphi))$-가군
$\widetilde{M'\otimes_{S'}S}|G(\varphi)$로 가는 함자적인 표준 준동형
$\nu$가 존재한다. 아이디얼 $S'_+$가 $S'_1$에 의해 생성되면
$\nu$는 동형이다.
\end{proposition}

\begin{proof}
실제로 $f'\in S'_d$ ($d>0$)에 대하여 다음과 같은
$S_{(f)}$-가군의 함자적인 표준 준동형을 정의한다. 여기서
$f=\varphi(f')$이다.
\[
\label{II.2.8.8.1-ko}
  \nu_f:M'_{(f')}\otimes_{S'_{(f')}}S_{(f)}
  \longrightarrow(M'\otimes_{S'}S)_{(f)}
\tag{2.8.8.1}
\]
이는 준동형
$M'_{(f')}\otimes_{S'_{(f')}}S_{(f)}\to
M'_{f'}\otimes_{S'_{f'}}S_f$와 표준 동형
$M'_{f'}\otimes_{S'_{f'}}S_f\xrightarrow{\sim}(M'\otimes_{S'}S)_f$
(\hyperref[0.1.5.4-ko]{0, 1.5.4})을 합성하여 정의하며, 뒤의 동형이
차수를 보존한다는 점을 사용한다. 두 번째 원소 $g'\in S'_+$와
$g=\varphi(g')$에 대하여 $D_+(f)$에서 $D_+(fg)$로 가는 제한
준동형과 $\nu_f$들이 양립함을 곧바로 확인할 수 있다. 따라서
(\hyperref[I.1.6.5-ko]{I, 1.6.5})을 고려하면 준동형
\[
  \nu:\Phi^*(\widetilde{M'})
  \longrightarrow\widetilde{M'\otimes_{S'}S}|G(\varphi)
\]
가 정의된다. 두 번째 주장을 증명하려면 모든 $f'\in S'_1$에 대하여
$\nu_f$가 동형임을 보이면 충분하다. 실제로 이 경우 $G(\varphi)$는
$D_+(\varphi(f'))$들의 합집합이다. 먼저 $\mathbf{Z}$-쌍선형 사상
$M'_m\times S_n\to M'_{(f')}\otimes_{S'_{(f')}}S_{(f)}$을 정의하자.
이 사상은 $(x',s)$에 원소
$(x'/{f'}^m)\otimes(s/f^n)$을 대응시킨다. 여기서 $m<0$일 때에는
$x'/{f'}^m$이 ${f'}^{-m}x'/1$이라는 약속을 쓴다.

\oldpage[II]{45}
(\hyperref[II.2.5.13-ko]{2.5.13})의 증명에서와 같이 이 사상으로부터
가군의 쌍준동형
\[
  \eta_f:M'\otimes_{S'}S
  \longrightarrow M'_{(f')}\otimes_{S'_{(f')}}S_{(f)}.
\]
을 얻는다. 또한 어떤 $r>0$에 대하여
$f^r\sum_i(x'_i\otimes s_i)=0$이면, 이를
$\sum_i({f'}^rx'_i\otimes s_i)=0$으로도 쓸 수 있다. 따라서
(\hyperref[0.1.5.4-ko]{0, 1.5.4})에 의해
$\sum_i({f'}^rx'_i/{f'}^{m_i+r})\otimes(s_i/f^{n_i})=0$이다. 즉
$\eta_f(\sum_i x'_i\otimes s_i)=0$이므로 $\eta_f$는 다음과 같이
인수분해된다.
\[
  M'\otimes_{S'}S\longrightarrow(M'\otimes_{S'}S)_f
  \xrightarrow{\eta'_f}M'_{(f')}\otimes_{S'_{(f')}}S_{(f)}~;
\]
끝으로 $\eta'_f$와 $\nu_f$가 서로 역인 동형임을 확인할 수 있다.

특히 (\hyperref[II.2.1.2.1-ko]{2.1.2.1})에서 모든
$n\in\mathbf{Z}$에 대하여 표준 준동형
\[
\label{II.2.8.8.2-ko}
  \Phi^*(\mathcal{O}_{X'}(n))\xrightarrow{\sim}
  \mathcal{O}_X(n)|G(\varphi)
\tag{2.8.8.2}
\]
을 얻는다.
\end{proof}

\begin{env}[2.8.9]
\phantomsection
\label{II.2.8.9-ko}
$A$, $A'$을 두 환이라 하고, $\psi:A'\to A$를 환 준동형이라 하자. 이
준동형은 사상
$\Psi:\operatorname{Spec}(A)\to\operatorname{Spec}(A')$를 정의한다.
$S'$를 음이 아닌 차수로 등급이 주어진 $A'$-대수라 하고
$S=S'\otimes_{A'}A$로 두자. 이는 자명하게 $S'_n\otimes_{A'}A$들을
동차 성분으로 갖는 등급 $A$-대수이다. 그러면 사상
$\varphi:s'\mapsto s'\otimes1$은 도식
(\hyperref[II.2.8.5.1-ko]{2.8.5.1})을 가환하게 하는 등급환
준동형이다. 여기서는 $S_+$가 $\varphi(S'_+)$에 의해 생성되는
$A$-가군이므로 $G(\varphi)=\operatorname{Proj}(S)=X$이다. 따라서
$X'=\operatorname{Proj}(S')$로 두면 다음 가환 도식을 얻는다.
\[
\label{II.2.8.9.1-ko}
  \xymatrix{
    X\ar[r]^{\Phi}\ar[d]_p & X'\ar[d]\\
    Y\ar[r]_{\Psi} & Y'
  }
\tag{2.8.9.1}
\]

이제 $M'$을 등급 $S'$-가군이라 하고
$M=M'\otimes_{A'}A=M'\otimes_{S'}S$로 두자. 이 조건 아래에서는
다음이 성립한다.
\end{env}

\begin{proposition}[2.8.10]
\phantomsection
\label{II.2.8.10-ko}
도식 (\hyperref[II.2.8.9.1-ko]{2.8.9.1})은 스킴 $X$를 곱
$X'\times_{Y'}Y$와 동일시한다. 또한 표준 준동형
$\nu:\Phi^*(\widetilde{M'})\to\widetilde{M}$
(\hyperref[II.2.8.8-ko]{2.8.8})은 동형이다.
\end{proposition}

\begin{proof}
첫째 주장은 $S'_+$의 모든 동차 원소 $f'$에 대하여
$f=\varphi(f')$로 둘 때, $\Phi$와 $p$를 $D_+(f)$에 제한한 사상들이
이 스킴을 곱 $D_+(f')\times_{Y'}Y$와 동일시함을 보이면 증명된다
(\hyperref[I.3.2.6.2-ko]{I, 3.2.6.2}). 다시 말해
(\hyperref[I.3.2.2-ko]{I, 3.2.2})에 따라 $S_{(f)}$가
$S'_{(f')}\otimes_{A'}A$와 표준적으로 동일시됨을 보이면 충분하다.
이는 차수를 보존하는 표준 동형
$S_f\xrightarrow{\sim}S'_{f'}\otimes_{A'}A$가 존재하므로 곧바로
따른다(\hyperref[0.1.5.4-ko]{0, 1.5.4}). 둘째 주장은 앞에서 본 바에
따라 $M'_{(f')}\otimes_{S'_{(f')}}S_{(f)}$가
$M'_{(f')}\otimes_{A'}A$와 동일시되고, 후자가 $M_{(f)}$와 동일시되는
사실에서 따른다. 실제로 $M_f$는 차수를 보존하는 동형에 의해
$M'_{f'}\otimes_{A'}A$와 표준적으로 동일시된다.
\end{proof}

\begin{corollary}[2.8.11]
\phantomsection
\label{II.2.8.11-ko}
모든 정수 $n\in\mathbf{Z}$에 대하여 $\widetilde{M(n)}$은
$\Phi^*(\widetilde{M'(n)})
=\widetilde{M'(n)}\otimes_{Y'}\mathcal{O}_Y$와 동일시된다. 특히
$\mathcal{O}_X(n)=\Phi^*(\mathcal{O}_{X'}(n))
=\mathcal{O}_{X'}(n)\otimes_{Y'}\mathcal{O}_Y$이다.
\end{corollary}

\begin{proof}
이는 (\hyperref[II.2.8.10-ko]{2.8.10})과
(\hyperref[II.2.5.15-ko]{2.5.15})에서 따른다.
\end{proof}

\oldpage[II]{46}

\begin{env}[2.8.12]
\phantomsection
\label{II.2.8.12-ko}
(\hyperref[II.2.8.9-ko]{2.8.9})의 가정 아래에서
$f'\in S'_d$ ($d>0$) 및 $f=\varphi(f')$에 대하여 다음 도식은
가환한다.
\[
\label{II.2.8.12.1-ko}
  \xymatrix{
    M'_{(f')}\ar[r]^{\sim}\ar[d]
    & M'^{(d)}/(f'-1)M'^{(d)}\ar[d]\\
    M_{(f)}\ar[r]^{\sim}
    & M^{(d)}/(f-1)M^{(d)}
  }
\tag{2.8.12.1}
\]
(\hyperref[II.2.2.5-ko]{2.2.5} 참조).
\end{env}

\begin{env}[2.8.13]
\phantomsection
\label{II.2.8.13-ko}
(\hyperref[II.2.8.9-ko]{2.8.9})의 표기와 가정을 그대로 두고,
$\mathcal{F}'$를 $\mathcal{O}_{X'}$-가군이라 하자.
$\mathcal{F}=\Phi^*(\mathcal{F}')$로 두면,
(\hyperref[II.2.8.11-ko]{2.8.11})과
(\hyperref[0.4.3.3-ko]{0, 4.3.3})에 의해 모든 $n\in\mathbf{Z}$에
대하여 $\mathcal{F}(n)=\Phi^*(\mathcal{F}'(n))$이다. 따라서
(\hyperref[0.3.7.1-ko]{0, 3.7.1})에 의해 표준 준동형
\[
  \Gamma(\rho):\Gamma(X',\mathcal{F}'(n))
  \longrightarrow\Gamma(X,\mathcal{F}(n))
\]
이 존재하며, 이로부터 등급가군의 표준 쌍준동형
\[
  \Gamma_\bullet(\mathcal{F}')\longrightarrow
  \Gamma_\bullet(\mathcal{F}).
\]
을 얻는다.

아이디얼 $S_+$가 $S_1$에 의해 생성된다고 가정하고,
$\mathcal{F}'=\widetilde{M'}$라 하자. 그러면
$M=M'\otimes_{A'}A$로 둘 때 $\mathcal{F}=\widetilde{M}$이다.
$f'$가 $S'_+$의 동차 원소이고 $f=\varphi(f')$이면,
$M_{(f)}=M'_{(f')}\otimes_{A'}A$임을 이미 보았으며 다음 도식은
가환한다.
\[
  \xymatrix{
    M'_0\ar[r]\ar[d]
    & M'_{(f')}\ar[d]
    & =\Gamma(D_+(f'),\widetilde{M'})\\
    M_0\ar[r]
    & M_{(f)}
    & =\Gamma(D_+(f),\widetilde{M})
  }
\]
이 관찰과 준동형 $\alpha$의 정의
(\hyperref[II.2.6.2-ko]{2.6.2})에서 다음 도식이 가환함을 곧바로
얻는다.
\[
\label{II.2.8.13.1-ko}
  \xymatrix{
    M'\ar[r]^{\alpha_{M'}}\ar[d]
    & \Gamma_\bullet(\widetilde{M'})\ar[d]\\
    M\ar[r]_{\alpha_M}
    & \Gamma_\bullet(\widetilde{M})
  }
\tag{2.8.13.1}
\]
마찬가지로 다음 도식도 가환한다.
\[
\label{II.2.8.13.2-ko}
  \xymatrix{
    \widetilde{\Gamma_\bullet(\mathcal{F}')}\ar[r]^{\beta_{\mathcal{F}'}}
      \ar[d]
    & \mathcal{F}'\ar[d]\\
    \widetilde{\Gamma_\bullet(\mathcal{F})}\ar[r]_{\beta_{\mathcal{F}}}
    & \mathcal{F}
  }
\tag{2.8.13.2}
\]
(오른쪽 세로 화살표는 표준 $\Phi$-사상
$\mathcal{F}'\to\Phi^*(\mathcal{F}')=\mathcal{F}$이다.)
\end{env}

\oldpage[II]{47}

\begin{env}[2.8.14]
\phantomsection
\label{II.2.8.14-ko}
계속해서 (\hyperref[II.2.8.9-ko]{2.8.9})의 가정과 표기법을 유지하고,
$N'$을 또 하나의 등급 $S'$-가군이라 하며
$N=N'\otimes_{A'}A$라 하자. 표준 쌍준동형 $M'\to M$, $N'\to N$은
(표준 환 준동형 $S'\to S$에 상대적인) 쌍준동형
$M'\otimes_{S'}N'\to M\otimes_S N$을 주며, 따라서 차수~$0$의
$S$-준동형
\[
  (M'\otimes_{S'}N')\otimes_{A'}A\longrightarrow M\otimes_S N,
\]
도 준다는 것은 곧바로 알 수 있다. 이에
(\hyperref[II.2.8.10-ko]{2.8.10})을 고려하여 다음과 같은
$\mathcal{O}_X$-가군의 준동형이 대응한다.
\[
\label{II.2.8.14.1-ko}
  \Phi^*(\widetilde{M'\otimes_{S'}N'})
  \longrightarrow\widetilde{M\otimes_S N}.
\tag{2.8.14.1}
\]

더 나아가 다음 도식이 가환함을 곧바로 확인할 수 있다.
\[
\label{II.2.8.14.2-ko}
  \xymatrix{
    \Phi^*(\widetilde{M'}\otimes_{\mathcal{O}_{X'}}\widetilde{N'})
      \ar[r]^{\sim}\ar[d]_{\Phi^*(\lambda)}
    & \widetilde{M}\otimes_{\mathcal{O}_X}\widetilde{N}
      \ar[d]^{\lambda}
    & =\Phi^*(\widetilde{M'})\otimes_{\mathcal{O}_X}
      \Phi^*(\widetilde{N'})\\
    \Phi^*(\widetilde{M'\otimes_{S'}N'})\ar[r]
    & \widetilde{M\otimes_S N}
  }
\tag{2.8.14.2}
\]
여기서 첫째 가로줄은 표준 동형
(\hyperref[0.4.3.3-ko]{0, 4.3.3})이다. 아이디얼 $S'_+$가
$S'_1$에 의해 생성되면, $S_+$가 $S_1$에 의해 생성됨은 명백하다.
이 경우 (\hyperref[II.2.8.14.2-ko]{2.8.14.2})의 두 세로 화살표는
동형이며(\hyperref[II.2.5.13-ko]{2.5.13}), 따라서
(\hyperref[II.2.8.14.1-ko]{2.8.14.1})도 동형이다.

마찬가지로, 차수 $k$의 준동형 $u'$에 역시 차수 $k$인 준동형
$u'\otimes1$을 대응시키면 표준 쌍준동형
$\operatorname{Hom}_{S'}(M',N')\to\operatorname{Hom}_S(M,N)$을 얻는다.
이로부터 다시 차수~$0$의 $S$-준동형
\[
  (\operatorname{Hom}_{S'}(M',N'))\otimes_{A'}A
  \longrightarrow\operatorname{Hom}_S(M,N)
\]
을 얻으며, 따라서 $\mathcal{O}_X$-가군의 준동형
\[
\label{II.2.8.14.3-ko}
  \Phi^*(\widetilde{\operatorname{Hom}_{S'}(M',N')})
  \longrightarrow\widetilde{\operatorname{Hom}_S(M,N)}.
\tag{2.8.14.3}
\]
을 얻는다.

더 나아가 다음 도식은 가환한다.
\[
  \xymatrix{
    \Phi^*(\widetilde{\operatorname{Hom}_{S'}(M',N')})\ar[r]
      \ar[d]_{\Phi^*(\mu)}
    & \widetilde{\operatorname{Hom}_S(M,N)}\ar[d]^{\mu}\\
    \Phi^*(\mathcal{H}om_{\mathcal{O}_{X'}}
      (\widetilde{M'},\widetilde{N'}))\ar[r]
    & \mathcal{H}om_{\mathcal{O}_X}(\widetilde{M},\widetilde{N})
  }
\]
여기서 둘째 가로줄은 표준 준동형
(\hyperref[0.4.4.6-ko]{0, 4.4.6})이다.
\end{env}

\oldpage[II]{48}

\begin{env}[2.8.15]
\phantomsection
\label{II.2.8.15-ko}
표기법과 가정이 (\hyperref[II.2.8.1-ko]{2.8.1})의 것이라 하자.
(\hyperref[II.2.4.7-ko]{2.4.7})에 의하면, $S_0$과 $S'_0$을
$\mathbf{Z}$로 바꾸고 $\varphi_0$을 항등사상으로 바꾼 뒤, 각각 $S$와
$S'$을 $S^{(d)}$와 ${S'}^{(d)}$로 바꾸고($d>0$) $\varphi$를
$S^{(d)}$에 대한 제한 $\varphi^{(d)}$로 바꾸어도 사상 $\Phi$는 동형까지
변하지 않는다.
\end{env}

\subsection{$\operatorname{Proj}(S)$ 스킴의 닫힌 부분스킴.}
\label{subsection:II.2.9-ko}

\begin{env}[2.9.1]
\phantomsection
\label{II.2.9.1-ko}
$\varphi:S\to S'$이 등급환의 준동형이라고 하자. 어떤 정수 $n$이 존재하여
$k\geq n$일 때 $\varphi_k:S_k\to S'_k$가 전사이면(각각 단사이면,
전단사이면), $\varphi$를 (TN)-\emph{전사}(각각 (TN)-\emph{단사},
(TN)-\emph{전단사})라고 한다. 그러면
(\hyperref[II.2.8.15-ko]{2.8.15})에 따라 $\Phi$의 연구는 $\varphi$가
\emph{전사}(각각 \emph{단사}, \emph{전단사})인 경우로 환원된다.
$\varphi$가 (TN)-전단사라고 하는 대신, 이때 이를
(TN)-\emph{동형사상}이라고도 한다.
\end{env}

\begin{proposition}[2.9.2]
\phantomsection
\label{II.2.9.2-ko}
$S$를 음이 아닌 차수로 등급이 주어진 환이라 하고
$X=\operatorname{Proj}(S)$라 하자.
\begin{enumerate}
  \item[\rm(i)] $\varphi:S\to S'$이 등급환의 (TN)-전사 준동형이면, 이에
    대응하는 사상 $\Phi$ (\hyperref[II.2.8.1-ko]{2.8.1})는
    $\operatorname{Proj}(S')$ 전체에서 정의되며
    $\operatorname{Proj}(S')$에서 $X$로 가는 닫힌 몰입 사상이다.
    $\mathfrak{J}$가 $\varphi$의 핵이면, $\Phi$에 결부된 $X$의 닫힌
    부분스킴은 $\mathcal{O}_X$의 준연접 아이디얼층
    $\widetilde{\mathfrak{J}}$로 정의된다.
  \item[\rm(ii)] 더 나아가 아이디얼 $S_+$가 차수~$1$인 유한 개의 동차
    원소로 생성된다고 가정하자. $X'$을 $\mathcal{O}_X$의 준연접
    아이디얼층 $\mathcal{J}$로 정의되는 $X$의 닫힌 부분스킴이라 하자.
    $\mathfrak{J}$를 표준 준동형
    $\alpha:S\to\Gamma_\bullet(\mathcal{O}_X)$
    (\hyperref[II.2.6.2-ko]{2.6.2})에 의한
    $\Gamma_\bullet(\mathcal{J})$의 역상인 $S$의 등급 아이디얼이라 하고,
    $S'=S/\mathfrak{J}$라 놓자. 그러면 $X'$은 등급환의 표준 준동형
    $S\to S'$에 대응하는 닫힌 몰입 사상
    $\operatorname{Proj}(S')\to X$에 결부된 부분스킴이다.
\end{enumerate}
\end{proposition}

\begin{proof}
\begin{enumerate}
  \item[\rm(i)] $\varphi$가 전사라고 가정해도 된다
    (\hyperref[II.2.9.1-ko]{2.9.1}). 가정에 따라 $\varphi(S_+)$가
    $S'_+$를 생성하므로 $G(\varphi)=\operatorname{Proj}(S')$이다.
    한편 둘째 주장은 $X$ 위에서 국소적으로 확인할 수 있다. 따라서
    $f$를 $S_+$의 동차 원소라 하고 $f'=\varphi(f)$라 놓자.
    $\varphi$는 등급환의 전사 준동형이므로
    $\varphi_{(f')}:S_{(f)}\to S'_{(f')}$가 전사이고 그 핵이
    $\mathfrak{J}_{(f)}$임을 곧바로 확인할 수 있다. 이로써 (i)가
    증명된다 (\hyperref[I.4.2.3-ko]{I, 4.2.3}).
  \item[\rm(ii)] (i)에 따라, 표준 단사 준동형
    $j:\mathfrak{J}\to S$에서 유도되는 준동형
    $\widetilde{j}:\widetilde{\mathfrak{J}}\to\mathcal{O}_X$가
    $\widetilde{\mathfrak{J}}$에서 $\mathcal{J}$로 가는 동형임을
    확인하면 충분하다. 이는 (\hyperref[II.2.7.11-ko]{2.7.11})에서
    따른다.
\end{enumerate}
\end{proof}

$\mathfrak{J}$는
$\widetilde{j}(\widetilde{\mathfrak{J}'})=\mathcal{J}$를 만족하는
$S$의 등급 아이디얼 $\mathfrak{J}'$들 가운데 \emph{가장 큰} 것이다.
실제로 정의로 거슬러 올라가 곧바로 확인하면
(\hyperref[II.2.6.2-ko]{2.6.2}), 이 관계에서
$\alpha(\mathfrak{J}')\subset\Gamma_\bullet(\mathcal{J})$가 따른다.

\begin{corollary}[2.9.3]
\phantomsection
\label{II.2.9.3-ko}
(\hyperref[II.2.9.2-ko]{2.9.2}, (i))의 가정이 성립하고, 더 나아가
아이디얼 $S_+$가 $S_1$로 생성된다고 가정하자. 그러면 모든
$n\in\mathbf{Z}$에 대하여 $\Phi^*(\widetilde{S(n)})$은
$\widetilde{S'(n)}$과 표준적으로 동형이고, 따라서 모든
$\mathcal{O}_X$-가군 $\mathcal{F}$에 대하여
$\Phi^*(\mathcal{F}(n))$은 $\Phi^*(\mathcal{F})(n)$과 표준적으로
동형이다.
\end{corollary}

\begin{proof}
이는 (\hyperref[II.2.8.8-ko]{2.8.8})의 특수한 경우이며,
(\hyperref[II.2.5.10.2-ko]{2.5.10.2})를 고려한다.
\end{proof}

\begin{corollary}[2.9.4]
\phantomsection
\label{II.2.9.4-ko}
(\hyperref[II.2.9.2-ko]{2.9.2}, (ii))의 가정이 성립한다고 하자.
$X$의 닫힌 부분준스킴 $X'$이 정역 스킴이기 위한 필요충분조건은 등급
아이디얼 $\mathfrak{J}$가 $S$의 소아이디얼인 것이다.
\end{corollary}

\oldpage[II]{49}

\begin{proof}
$X'$은 $\operatorname{Proj}(S/\mathfrak{J})$와 동형이므로,
(\hyperref[II.2.4.4-ko]{2.4.4})에 따라 조건은 충분하다. 필요성을
보이기 위하여 완전열
$0\to\mathcal{J}\to\mathcal{O}_X\to\mathcal{O}_X/\mathcal{J}$를
생각하자. 이로부터 완전열
\[
  0\longrightarrow\Gamma_\bullet(\mathcal{J})
  \longrightarrow\Gamma_\bullet(\mathcal{O}_X)
  \longrightarrow\Gamma_\bullet(\mathcal{O}_X/\mathcal{J}).
\]
을 얻는다. $f\in S_m$, $g\in S_n$이고
$\Gamma_\bullet(\mathcal{O}_X/\mathcal{J})$에서
$\alpha_{n+m}(fg)$의 상이 영일 때, $\alpha_m(f)$와
$\alpha_n(g)$의 상 가운데 하나가 영임을 보이면 충분하다. 그런데
정의에 따라 이 상들은 정역 스킴 $X'$ 위의 가역
$(\mathcal{O}_X/\mathcal{J})$-가군
$\mathcal{L}=(\mathcal{O}_X/\mathcal{J})(m)$과
$\mathcal{L}'=(\mathcal{O}_X/\mathcal{J})(n)$의 단면들이다. 가정에
의하면 이 두 단면의 곱은 $\mathcal{L}\otimes\mathcal{L}'$에서 영이다
((\hyperref[II.2.9.3-ko]{2.9.3}) 및
(\hyperref[II.2.5.14.1-ko]{2.5.14.1})). 따라서
(\hyperref[I.7.4.4-ko]{I, 7.4.4})에 따라 둘 중 하나가 영이다.
\end{proof}

\begin{corollary}[2.9.5]
\phantomsection
\label{II.2.9.5-ko}
$A$를 환, $M$을 $A$-가군, $S$를 차수~$1$인 동차 원소들의 집합
$S_1$로 생성되는 등급 $A$-대수라 하자. $u:M\to S_1$을 $A$-가군의
전사 준동형이라 하고, $\overline{u}:\mathbf{S}(M)\to S$를 $u$를
확장하는 $M$의 대칭대수 $\mathbf{S}(M)$에서 $S$로 가는
($A$-대수의) 준동형이라 하자. 그러면 $\overline{u}$에 대응하는 사상은
$\operatorname{Proj}(S)$에서 $\operatorname{Proj}(\mathbf{S}(M))$로
가는 닫힌 몰입 사상이다.
\end{corollary}

\begin{proof}
가정에 따라 $\overline{u}$는 전사이므로
(\hyperref[II.2.9.2-ko]{2.9.2})를 적용하면 된다.
\end{proof}

\section{등급 대수의 층의 동차 스펙트럼}
\label{section:II.3-ko}

\subsection{준연접 등급 $\mathcal{O}_Y$-대수의 동차 스펙트럼.}
\label{subsection:II.3.1-ko}
\label{II.3.1-ko}

\begin{env}[3.1.1]
\phantomsection
\label{II.3.1.1-ko}
$Y$를 준스킴, $\mathcal{S}$를 등급 $\mathcal{O}_Y$-대수,
$\mathcal{M}$을 등급 $\mathcal{S}$-가군이라 하자. $\mathcal{S}$가
\emph{준연접}이면, 각각의 동차 성분 $\mathcal{S}_n$은 $\mathcal{S}$에서
자기 자신으로 가는 준동형에 의한 $\mathcal{S}$의 상이므로 \emph{준연접}
$\mathcal{O}_Y$-가군이다
((\hyperref[I.1.3.8-ko]{I, 1.3.8}) 및
(\hyperref[I.1.3.9-ko]{I, 1.3.9})). 마찬가지로 $\mathcal{M}$이
$\mathcal{O}_Y$-가군으로서 준연접이면 그 동차 성분 $\mathcal{M}_n$도
준연접이며, 역도 성립한다. $d$가 $>0$인 정수이면,
$\mathcal{O}_Y$-가군들 $\mathcal{S}_{nd}$ ($n\in\mathbf{Z}$)의 직합을
$\mathcal{S}^{(d)}$로 나타내겠다. 이 직합은 $\mathcal{S}$가 준연접이면
준연접이다 (\hyperref[I.1.3.9-ko]{I, 1.3.9}).
$0\leq k\leq d-1$인 모든 정수 $k$에 대하여, $\mathcal{M}_{nd+k}$
($n\in\mathbf{Z}$)의 직합을 $\mathcal{M}^{(d,k)}$로(또는 $k=0$일
때에는 $\mathcal{M}^{(d)}$로) 나타내겠다. 이는 등급
$\mathcal{S}^{(d)}$-가군이며, $\mathcal{S}$와 $\mathcal{M}$이 준연접이면
준연접이다 (\hyperref[I.9.6.1-ko]{I, 9.6.1}). 모든
$k\in\mathbf{Z}$에 대하여
$(\mathcal{M}(n))_k=\mathcal{M}_{n+k}$를 만족하는 등급
$\mathcal{S}$-가군을 $\mathcal{M}(n)$으로 나타내겠다. $\mathcal{S}$와
$\mathcal{M}$이 준연접이면 $\mathcal{M}(n)$은 준연접 등급
$\mathcal{S}$-가군이다 (\hyperref[I.9.6.1-ko]{I, 9.6.1}).

$\mathcal{M}$이 등급 $\mathcal{S}$-가군으로서 \emph{유한 생성}이라는
것은(각각 \emph{유한 표시를 갖는다}는 것은), 모든 $y\in Y$에 대하여
$y$의 열린 근방 $U$와 정수 $n_i$(각각 정수 $m_i$와 $n_j$)가 존재하여
차수~$0$인 전사 준동형
\[
  \bigoplus_{i=1}^r(\mathcal{S}(n_i)|U)\longrightarrow\mathcal{M}|U
\]
이 존재한다는 뜻이다(각각 $\mathcal{M}|U$가 차수~$0$인 다음 준동형의
여핵과 동형이어야 한다:
\[
  \bigoplus_{i=1}^r(\mathcal{S}(m_i)|U)
  \longrightarrow\bigoplus_{j=1}^s(\mathcal{S}(n_j)|U)).
\]

$U$를 $Y$의 아핀 열린집합, $A=\Gamma(U,\mathcal{O}_Y)$를 그 환이라
하자. 가정에 따라 등급 $(\mathcal{O}_Y|U)$-대수 $\mathcal{S}|U$는
$\widetilde{S}$와 동형이다. 여기서 $S=\Gamma(U,\mathcal{S})$는
$A$-대수이며,

\oldpage[II]{50}

등급이 주어져 있다 (\hyperref[I.1.4.3-ko]{I, 1.4.3}).
$X_U=\operatorname{Proj}(\Gamma(U,\mathcal{S}))$라 놓자.
$U'\subset U$를 $Y$의 두 번째 아핀 열린집합, $A'$를 그 환,
$j:U'\to U$를 제한 준동형 $A\to A'$에 대응하는 표준 포함 사상이라
하자. 이때 $\mathcal{S}|U'=j^*(\mathcal{S}|U)$이고, 따라서
$S'=\Gamma(U',\mathcal{S})$는 $S\otimes_A A'$과 표준적으로 동일시된다
(\hyperref[I.1.6.4-ko]{I, 1.6.4}). 이로부터
(\hyperref[II.2.8.10-ko]{2.8.10})에 따라 $X_{U'}$은
$X_U\times_U U'$과 표준적으로 동일시되고, 따라서 $f_U$가 구조 사상
$X_U\to U$를 나타낼 때 $f_U^{-1}(U')$와도 표준적으로 동일시된다
(\hyperref[I.4.4.1-ko]{I, 4.4.1}). 이렇게 정의된 표준 동형사상
$f_U^{-1}(U')\xrightarrow{\sim}X_{U'}$을 $\sigma_{U',U}$로 나타내고,
$\sigma_{U',U}^{-1}$과 표준 포함 사상 $f_U^{-1}(U')\to X_U$를 차례로
합성하여 얻는 열린 몰입 사상 $X_{U'}\to X_U$를 $\rho_{U',U}$로
나타내겠다. $U''\subset U'$이 $Y$의 세 번째 아핀 열린집합이면
$\rho_{U'',U}=\rho_{U',U}\circ\rho_{U'',U'}$임은 곧바로 알 수 있다.
\end{env}

\begin{proposition}[3.1.2]
\phantomsection
\label{II.3.1.2-ko}
$Y$를 준스킴이라 하자. 음이 아닌 차수로 등급이 주어진 임의의 준연접
$\mathcal{O}_Y$-대수 $\mathcal{S}$에 대하여, $Y$ 위의 준스킴 $X$가
존재하며 $Y$-동형까지 유일하게 결정되고, 다음 성질을 갖는다. $f$가
구조 사상 $X\to Y$이면, $Y$의 모든 아핀 열린집합 $U$에 대하여
$f^{-1}(U)$ 위에 유도된 준스킴에서
$X_U=\operatorname{Proj}(\Gamma(U,\mathcal{S}))$로 가는
\emph{동형사상} $\eta_U$가 존재하며, $U$에 포함되는 $Y$의 두 번째
아핀 열린집합 $V$에 대하여 다음 도식이 가환한다.
\[
\label{II.3.1.2.1-ko}
  \xymatrix{
    f^{-1}(V)\ar[r]^{\eta_V}\ar[d]
    & X_V\ar[d]^{\rho_{V,U}}\\
    f^{-1}(U)\ar[r]_{\eta_U}
    & X_U
  }
\tag{3.1.2.1}
\]
\end{proposition}

\begin{proof}
$Y$의 두 아핀 열린집합 $U$, $V$에 대하여, $X_{U,V}$를 $X_U$가
$f_U^{-1}(U\cap V)$ 위에 유도하는 준스킴이라 하자. $Y$-동형사상
$\theta_{U,V}:X_{V,U}\xrightarrow{\sim}X_{U,V}$를 정의하고자 한다.
이를 위해 아핀 열린집합 $W\subset U\cap V$를 생각하자. 다음
동형사상들을 합성하여
\[
  f_U^{-1}(W)\xrightarrow{\sigma_{W,U}}X_W
  \xrightarrow{\sigma_{W,V}^{-1}}f_V^{-1}(W),
\]
동형사상 $\tau_W$를 얻는다. $W'\subset W$가 아핀 열린집합이면,
$\tau_{W'}$는 $\tau_W$를 $f_U^{-1}(W')$로 제한한 것임을 곧바로
확인할 수 있다. 따라서 $\tau_W$들은 실제로 하나의 $Y$-동형사상
$\theta_{V,U}$의 제한들이다. 또한 $U$, $V$, $W$가 $Y$의 세 아핀
열린집합이고, $\theta'_{U,V}$, $\theta'_{V,W}$,
$\theta'_{U,W}$가 각각 $\theta_{U,V}$, $\theta_{V,W}$,
$\theta_{U,W}$를 $X_V$, $X_W$, $X_W$ 안에서의
$U\cap V\cap W$의 역상들로 제한한 것들이라 하면, 앞의 정의들로부터
$\theta'_{U,V}\circ\theta'_{V,W}=\theta'_{U,W}$임이 따른다. 따라서
명시된 성질들을 만족하는 $X$의 존재는
(\hyperref[I.2.3.1-ko]{I, 2.3.1})에서 따르며, $Y$-동형까지 유일함은
(\hyperref[II.3.1.2.1-ko]{3.1.2.1})을 고려하면 자명하다.
\end{proof}

\begin{env}[3.1.3]
\phantomsection
\label{II.3.1.3-ko}
(\hyperref[II.3.1.2-ko]{3.1.2})에서 정의한 준스킴 $X$를 준연접인
등급 $\mathcal{O}_Y$-대수 $\mathcal{S}$의 \emph{동차 스펙트럼}이라
하고, 이를 $\operatorname{Proj}(\mathcal{S})$로 나타낸다.
$\operatorname{Proj}(\mathcal{S})$가 \emph{$Y$ 위에서 분리되어 있음}은
명백하다 ((\hyperref[II.2.4.2-ko]{2.4.2}) 및
(\hyperref[I.5.5.5-ko]{I, 5.5.5})). $\mathcal{S}$가
$\mathcal{O}_Y$-대수로서 \emph{유한 생성}이면
(\hyperref[I.9.6.2-ko]{I, 9.6.2}),
$\operatorname{Proj}(\mathcal{S})$는 $Y$ 위에서 \emph{유한형}이다
((\hyperref[II.2.7.1-ko]{2.7.1, (ii)}) 및
(\hyperref[I.6.3.1-ko]{I, 6.3.1})).

$f$가 구조 사상 $X\to Y$이면, $Y$의 임의의 열린집합 $U$를 $Y$가
그 위에 유도하는 준스킴으로 볼 때 $f^{-1}(U)$는 동차 스펙트럼
$\operatorname{Proj}(\mathcal{S}|U)$와 동일시됨이 명백하다.
\end{env}

\begin{proposition}[3.1.4]
\phantomsection
\label{II.3.1.4-ko}
$d>0$이고 $f\in\Gamma(Y,\mathcal{S}_d)$라 하자.
$X=\operatorname{Proj}(\mathcal{S})$의 바탕공간에는 다음 성질을 갖는
열린 부분집합 $X_f$가 존재한다. 즉, $Y$의 모든 아핀 열린집합 $U$에 대하여,

\oldpage[II]{51}

$X_f\cap\varphi^{-1}(U)=D_+(f|U)$이다. 이 등식에서는
$\varphi^{-1}(U)$를
$X_U=\operatorname{Proj}(\Gamma(U,\mathcal{S}))$와 동일시하며
($\varphi$는 구조 사상 $X\to Y$를 나타낸다). 더 나아가 $X_f$ 위에
유도된 준스킴은 $Y$ 위에서 아핀이고
$\operatorname{Spec}(\mathcal{S}^{(d)}/(f-1)\mathcal{S}^{(d)})$와
표준적으로 동형이다
(\hyperref[II.1.3.1-ko]{1.3.1}).
\end{proposition}

\begin{proof}
$f|U\in\Gamma(U,\mathcal{S}_d)=(\Gamma(U,\mathcal{S}))_d$이다.
$U$, $U'$이 $U'\subset U$인 $Y$의 두 아핀 열린집합이면, $f|U'$은
제한 준동형
\[
  \Gamma(U,\mathcal{S})\longrightarrow\Gamma(U',\mathcal{S}),
\]
에 의한 $f|U$의 상이다. 따라서
(\hyperref[II.3.1.1-ko]{3.1.1})의 표기에서 $D_+(f|U')$은
$X_{U'}$가 역상 $\rho_{U',U}^{-1}(D_+(f|U))$ 위에 유도하는
준스킴과 같다
(\hyperref[II.2.8.1-ko]{2.8.1}). 이로부터 첫째 주장이 따른다. 또한
$X_U$가 $D_+(f|U)$ 위에 유도하는 준스킴은
$\operatorname{Spec}((\Gamma(U,\mathcal{S}))_{(f|U)})$와 표준적으로
동일시되며, 이러한 동일시들은 제한 준동형들과 양립한다
(\hyperref[II.2.8.1-ko]{2.8.1}). 그러면 둘째 주장은
(\hyperref[II.2.2.5-ko]{2.2.5})와 도식
(\hyperref[II.2.8.2.1-ko]{2.8.2.1})의 가환성에서 따른다.
\end{proof}

또한 $X_f$를 ($X$의 바탕공간 안의 열린집합으로 보아) $f$가
\emph{소멸하지 않는} $x\in X$들의 집합이라고 한다.

\begin{corollary}[3.1.5]
\phantomsection
\label{II.3.1.5-ko}
$f\in\Gamma(Y,\mathcal{S}_d)$이고
$g\in\Gamma(Y,\mathcal{S}_e)$이면 다음이 성립한다.
\[
\label{II.3.1.5.1-ko}
  X_{fg}=X_f\cap X_g.
\tag{3.1.5.1}
\]
\end{corollary}

\begin{proof}
실제로 $U$가 $Y$의 아핀 열린집합일 때, 두 변과
$\varphi^{-1}(U)$의 교집합을 각각 고찰하고 공식
(\hyperref[II.2.3.3.2-ko]{2.3.3.2})을 적용하면 충분하다.
\end{proof}

\begin{corollary}[3.1.6]
\phantomsection
\label{II.3.1.6-ko}
$(f_\alpha)$를 $Y$ 위의 $\mathcal{S}$의 단면들로 이루어진 족이라 하고,
각 $\alpha$에 대하여
$f_\alpha\in\Gamma(Y,\mathcal{S}_{d_\alpha})$라 하자. 이 족이 생성하는
$\mathcal{S}$의 아이디얼층
(\hyperref[0.5.1.1-ko]{0, 5.1.1})이 어떤 값 이상의 모든 $n$에 대하여
$\mathcal{S}_n$을 포함하면, 바탕공간 $X$는 $X_{f_\alpha}$들의
합집합이다.
\end{corollary}

\begin{proof}
실제로 $Y$의 모든 아핀 열린집합 $U$에 대하여 $\varphi^{-1}(U)$는
$X_{f_\alpha}\cap\varphi^{-1}(U)$들의 합집합이다
(\hyperref[II.2.3.14-ko]{2.3.14}).
\end{proof}

\begin{corollary}[3.1.7]
\label{II.3.1.7-ko}
$\mathcal{A}$를 준연접 $\mathcal{O}_Y$-대수라 하고
\[
  \mathcal{S}=\mathcal{A}[T]
  =\mathcal{A}\otimes_{\mathbf{Z}}\mathbf{Z}[T]
\]
로 놓자. 여기서 $T$는 부정원이다($\mathbf{Z}$와 $\mathbf{Z}[T]$은
$Y$ 위의 상수층으로 본다). 그러면
$X=\operatorname{Proj}(\mathcal{S})$는
$\operatorname{Spec}(\mathcal{A})$와 표준적으로 동일시된다. 특히
$\operatorname{Proj}(\mathcal{O}_Y[T])$는 $Y$와 동일시된다.
\end{corollary}

\begin{proof}
$Y$의 각 점에서 $T$와 같은 유일한 단면
$f\in\Gamma(Y,\mathcal{S})$에 (\hyperref[II.3.1.6-ko]{3.1.6})을 적용하면
$X_f=X$임을 알 수 있다. 또한 여기서는 $d=1$이고
$\mathcal{S}^{(1)}/(f-1)\mathcal{S}^{(1)}
=\mathcal{S}/(f-1)\mathcal{S}$는 $\mathcal{A}$와 표준적으로 동형이다.
이로부터 따름정리가 따른다(\hyperref[II.1.2.2-ko]{1.2.2}).
\end{proof}

$g\in\Gamma(Y,\mathcal{O}_Y)$라 하자. $\mathcal{S}=\mathcal{O}_Y[T]$로
택하면 $g\in\Gamma(Y,\mathcal{S}_0)$이고, 다음과 같이 놓자.
\[
  h=gT\in\Gamma(Y,\mathcal{S}_1).
\]
$X=\operatorname{Proj}(\mathcal{S})$라 하면
(\hyperref[II.3.1.7-ko]{3.1.7})에서 정의한 표준 동일시는 $X_h$를 $Y$의
열린집합 $Y_g$와 동일시한다((\hyperref[0.5.5.2-ko]{0, 5.5.2})의
의미에서). 실제로 $Y=\operatorname{Spec}(A)$가 아핀인 경우로 한정해도
되고, 그러면 모든 것은 (\hyperref[II.2.2.5-ko]{2.2.5})를 고려할 때
분수환 $A_g$가 $A[T]/(gT-1)A[T]$와 표준적으로 동일시된다는 사실로
환원된다(\hyperref[0.1.2.3-ko]{0, 1.2.3}).

\begin{proposition}[3.1.8]
\label{II.3.1.8-ko}
$\mathcal{S}$를 음이 아닌 차수로 등급이 주어진 준연접
$\mathcal{O}_Y$-대수라 하자.
\begin{enumerate}
  \item[\rm(i)] 모든 $d>0$에 대하여 $\operatorname{Proj}(\mathcal{S})$에서
    $\operatorname{Proj}(\mathcal{S}^{(d)})$로 가는 표준 $Y$-동형이
    존재한다.

\oldpage[II]{52}

  \item[\rm(ii)] $\mathcal{S}'$을 $\mathcal{O}_Y$와 $\mathcal{S}_n$들
    ($n\geq0$)의 직합인 등급 $\mathcal{O}_Y$-대수라 하자. 그러면
    $\operatorname{Proj}(\mathcal{S}')$과
    $\operatorname{Proj}(\mathcal{S})$는 표준적으로 $Y$-동형이다.
  \item[\rm(iii)] $\mathcal{L}$을 가역 $\mathcal{O}_Y$-가군
    (\hyperref[0.5.4.1-ko]{0, 5.4.1})이라 하고,
    $\mathcal{S}_{(\mathcal{L})}$을
    $\mathcal{S}_d\otimes\mathcal{L}^{\otimes d}$들($d\geq0$)의 직합인
    등급 $\mathcal{O}_Y$-대수라 하자. 그러면
    $\operatorname{Proj}(\mathcal{S})$와
    $\operatorname{Proj}(\mathcal{S}_{(\mathcal{L})})$는 표준적으로
    $Y$-동형이다.
\end{enumerate}
\end{proposition}

\begin{proof}
세 경우 각각에서 동형을 $Y$ 위에서 국소적으로 정의하면 충분하다.
열린집합에서 그보다 작은 열린집합으로의 제한 연산과 양립함을 확인하는
것은 자명하기 때문이다. 따라서 $Y$가 아핀이라고 가정할 수 있으며,
이때 (i)은 (\hyperref[II.2.4.7-ko]{2.4.7, (i)})에서, (ii)는
(\hyperref[II.2.4.8-ko]{2.4.8})에서 따른다. (iii)의 경우에는 더 나아가
$\mathcal{L}$이 $\mathcal{O}_Y$와 동형이라고 가정하면(문제가 $Y$ 위에서
국소적이므로 이렇게 해도 된다), $\operatorname{Proj}(\mathcal{S})$와
$\operatorname{Proj}(\mathcal{S}_{(\mathcal{L})})$가 서로 동형임은
자명하다. \emph{표준} 동형을 정의하기 위하여
$Y=\operatorname{Spec}(A)$, $\mathcal{S}=\widetilde{S}$라 하자. 여기서
$S$는 등급 $A$-대수이다. 또한 $\mathcal{L}=\widetilde{L}$인 자유
$A$-가군 $L$의 생성원을 $c$라 하자. 그러면 모든 $n>0$에 대하여
$x_n\mapsto x_n\otimes c^{\otimes n}$은 $S_n$에서
$S_n\otimes L^{\otimes n}$으로 가는 $A$-동형이고, 이 $A$-동형들은
다음과 같은 등급 $A$-대수의 동형을 정의한다.
\[
  \varphi_c:S\longrightarrow S_{(L)}
  =\bigoplus_{n\geq0}S_n\otimes L^{\otimes n}.
\]
이제 $f\in S_+$가 차수 $d$인 동차 원소라 하자. 모든 $x\in S_{nd}$에
대하여
\[
  \frac{x\otimes c^{\otimes nd}}{(f\otimes c^{\otimes d})^n}
  =\frac{x\otimes(\varepsilon c)^{\otimes nd}}
  {(f\otimes(\varepsilon c)^{\otimes d})^n}
\]
가 모든 가역원 $\varepsilon\in A$에 대하여 성립한다. 따라서
$\varphi_c$에서 유도되는 동형
$S_{(f)}\to(S_{(L)})_{(f\otimes c^d)}$는 택한 $L$의 생성원 $c$에
\emph{무관}하며, 이로써 증명이 끝난다.
\end{proof}

\begin{env}[3.1.9]
\label{II.3.1.9-ko}
((\hyperref[0.4.1.3-ko]{0, 4.1.3})과
(\hyperref[I.1.3.14-ko]{I, 1.3.14}))에 따라, 등급이 주어진 준연접
$\mathcal{O}_Y$-대수 $\mathcal{S}$가
\emph{$\mathcal{O}_Y$-가군 $\mathcal{S}_1$에 의해 생성되기} 위한
필요충분조건은 다음과 같음을 상기하자. $Y$의 아핀 열린집합들로 이루어진
덮개 $(U_\alpha)$가 존재하여, 모든 $\alpha$에 대하여
$\Gamma(U_\alpha,\mathcal{S}_0)$ 위의 등급 대수
$\Gamma(U_\alpha,\mathcal{S})$가 차수~$1$인 동차 원소들의 집합
$\Gamma(U_\alpha,\mathcal{S}_1)$에 의해 생성되어야 한다. 그러면 $Y$의
모든 열린집합 $V$에 대하여 $\mathcal{S}|V$는
$(\mathcal{O}_Y|V)$-가군 $\mathcal{S}_1|V$에 의해 생성된다.
\end{env}

\begin{proposition}[3.1.10]
\label{II.3.1.10-ko}
$Y$에 유한 개의 아핀 열린집합들로 이루어진 덮개 $(U_i)$가 존재하여
각 등급 대수 $\Gamma(U_i,\mathcal{S})$가
$\Gamma(U_i,\mathcal{O}_Y)$ 위의 유한 생성 대수라고 가정하자. 그러면
어떤 $d>0$가 존재하여 $\mathcal{S}^{(d)}$는 $\mathcal{S}_d$로 생성되고,
$\mathcal{S}_d$는 유한 생성 $\mathcal{O}_Y$-가군이다.
\end{proposition}

\begin{proof}
실제로 각 $i$에 대하여 어떤 정수 $m_i$가 존재하여 모든 $n>0$에 대해
$\Gamma(U_i,\mathcal{S}_{nm_i})=(\Gamma(U_i,\mathcal{S}_{m_i}))^n$임이
(\hyperref[II.2.1.6-ko]{2.1.6, (v)})에서 따른다. 따라서
(\hyperref[II.2.1.6-ko]{2.1.6, (i)})을 고려하면, $d$로 $m_i$들의
공배수를 택하면 충분하다.
\end{proof}

\begin{corollary}[3.1.11]
\label{II.3.1.11-ko}
(\hyperref[II.3.1.10-ko]{3.1.10})의 가정 아래에서
$\operatorname{Proj}(\mathcal{S})$는 동차 스펙트럼
$\operatorname{Proj}(\mathcal{S}')$과 $Y$-동형이다. 여기서
$\mathcal{S}'$은 $\mathcal{S}'_1$로 생성되는 등급
$\mathcal{O}_Y$-대수이고, $\mathcal{S}'_1$은 유한 생성
$\mathcal{O}_Y$-가군이다.
\end{corollary}

\begin{proof}
실제로 (\hyperref[II.3.1.10-ko]{3.1.10})의 성질로 정해지는 $d$에 대하여
$\mathcal{S}'=\mathcal{S}^{(d)}$로 놓고
(\hyperref[II.3.1.8-ko]{3.1.8, (i)})을 적용하면 충분하다.
\end{proof}

\begin{env}[3.1.12]
\label{II.3.1.12-ko}
$\mathcal{S}$가 음이 아닌 차수로 등급이 주어진 준연접
$\mathcal{O}_Y$-대수이면, 그 \emph{멱영근기} $\mathcal{N}$이 준연접
$\mathcal{O}_Y$-가군임을 알고 있다
(\hyperref[I.5.1.1-ko]{I, 5.1.1}).
$\mathcal{N}_+=\mathcal{N}\cap\mathcal{S}_+$를
\emph{$\mathcal{S}_+$의 멱영근기}라 한다. 이는 준연접 등급
$\mathcal{S}_0$-가군이다. 실제로 $Y$가 아핀인 경우로 곧바로 환원되며,
그 경우 이 명제는 (\hyperref[II.2.1.10-ko]{2.1.10})에서 따른다. 모든
$y\in Y$에 대하여 $(\mathcal{N}_+)_y$는
$(\mathcal{S}_+)_y=(\mathcal{S}_y)_+$의 멱영근기이다
(\hyperref[I.5.1.1-ko]{I, 5.1.1}). 등급 $\mathcal{O}_Y$-대수
$\mathcal{S}$가 \emph{본질적으로 축소}라는 것은
$\mathcal{N}_+=0$이라는 뜻이며, 이는 모든 $y\in Y$에 대하여

\oldpage[II]{53}

$\mathcal{S}_y$가 본질적으로 축소인 등급 $\mathcal{O}_y$-대수라는 것과
동치이다. 임의의 등급 $\mathcal{O}_Y$-대수 $\mathcal{S}$에 대하여
$\mathcal{S}/\mathcal{N}_+$는 본질적으로 축소이다.

$\mathcal{S}$가 \emph{정역}이라는 것은 모든 $y\in Y$에 대하여
$\mathcal{S}_y$가 정역이고, 더욱이 모든 $y\in Y$에 대하여
$(\mathcal{S}_y)_+=(\mathcal{S}_+)_y\neq0$이라는 뜻이다.
\end{env}

\begin{proposition}[3.1.13]
\phantomsection
\label{II.3.1.13-ko}
$\mathcal{S}$를 음이 아닌 차수로 등급이 주어진 $\mathcal{O}_Y$-대수라
하자. $X=\operatorname{Proj}(\mathcal{S})$라 하면, $Y$-스킴
$X_{\mathrm{red}}$는 $\operatorname{Proj}(\mathcal{S}/\mathcal{N}_+)$와
표준적으로 동형이다. 특히 $\mathcal{S}$가 본질적으로 축소이면 $X$는
축소이다.
\end{proposition}

\begin{proof}
$X'=\operatorname{Proj}(\mathcal{S}/\mathcal{N}_+)$가 축소라는 사실은 이
성질이 국소적이므로 (\hyperref[II.2.4.4-ko]{2.4.4, (i)})에서 곧바로
따른다. 더욱이 모든 아핀 열린집합 $U\subset Y$에 대하여
${\varphi'}^{-1}(U)$는 $(\varphi^{-1}(U))_{\mathrm{red}}$와 같다(여기서
$\varphi$와 $\varphi'$는 각각 구조 사상 $X\to Y$와 $X'\to Y$를
나타낸다). 표준 $U$-사상
${\varphi'}^{-1}(U)\to\varphi^{-1}(U)$들이 제한 연산과 양립함을 곧바로
확인할 수 있으므로, 이들은 바탕공간 사이의 위상동형인 닫힌 몰입 사상
$X'\to X$를 정의한다. 따라서 (\hyperref[I.5.1.2-ko]{I, 5.1.2})에 의해
결론이 따른다.
\end{proof}

\begin{proposition}[3.1.14]
\phantomsection
\label{II.3.1.14-ko}
$Y$를 정역 준스킴, $\mathcal{S}$를 등급이 주어진 준연접
$\mathcal{O}_Y$-대수라 하고 $\mathcal{S}_0=\mathcal{O}_Y$라 하자.
\begin{enumerate}
  \item[\rm(i)] $\mathcal{S}$가 정역이면
    (\hyperref[II.3.1.12-ko]{3.1.12}),
    $X=\operatorname{Proj}(\mathcal{S})$는 정역 스킴이고 구조 사상
    $\varphi:X\to Y$는 우세 사상이다.
  \item[\rm(ii)] 더 나아가 $\mathcal{S}$가 본질적으로 축소라고
    가정하자. 역으로 $X$가 정역 스킴이고 $\varphi$가 우세 사상이면
    $\mathcal{S}$는 정역이다.
\end{enumerate}
\end{proposition}

\begin{proof}
\begin{enumerate}
  \item[\rm(i)] $(U_\alpha)$가 공집합이 아닌 아핀 열린집합들로 이루어진
    $Y$의 기저라 하자. $Y$를 어느 한 $U_\alpha$로, $\mathcal{S}$를
    $\mathcal{S}|U_\alpha$로 바꾼 경우에 명제를 증명하면 충분하다.
    실제로 이로부터 한편으로 바탕공간 $\varphi^{-1}(U_\alpha)$들은
    $X$의 기약 열린집합들(따라서 공집합이 아닌 열린집합들)이고, 모든
    지표쌍에 대하여
    $\varphi^{-1}(U_\alpha)\cap\varphi^{-1}(U_\beta)\neq\varnothing$이다
    ($U_\alpha\cap U_\beta$가 어떤 $U_\gamma$를 포함하기 때문이다).
    따라서 $X$는 기약이다 (\hyperref[0.2.1.4-ko]{0, 2.1.4}). 다른 한편
    축소성은 국소 성질이므로 $X$는 축소이다. 그러므로 $X$는 정역
    스킴이고 $\varphi(X)$는 $Y$에서 조밀하다.

    이제 $Y=\operatorname{Spec}(A)$이고 $A$가 정역이라고 가정하자
    (\hyperref[I.5.1.4-ko]{I, 5.1.4}). 또한 $\mathcal{S}=\widetilde{S}$라
    하자. 여기서 $S$는 등급 $A$-대수이다. 가정에 따르면 모든 $y\in Y$에
    대하여 $\widetilde{S}_y=S_y$는 $(S_y)_+\neq0$을 만족하는 정역인
    등급환이다. $S$가 \emph{정역}임을 증명하면 충분하다. 그러면
    $S_+\neq0$이고 (\hyperref[II.2.4.4-ko]{2.4.4, (ii)})를 적용할 수 있기
    때문이다. 이제 $f$, $g$를 $\neq0$인 $S$의 두 원소라 하고
    $fg=0$이라고 가정하자. 모든 $y\in Y$에 대하여
    $(f/1)(g/1)=0$이 $S_y$에서 성립하므로, 가정에 따라 $f/1=0$이거나
    $g/1=0$이다. 예를 들어 $f/1=0$이 $S_y$에서 성립한다고 하자. 이는
    어떤 $a\in A$에 대하여 $a\notin\mathfrak{j}_y$이고 $af=0$이라는
    뜻이다. 그러면 모든 $z\in Y$에 대하여 $(a/1)(f/1)=0$이 \emph{정역}
    $S_z$에서 성립하고, $a/1\neq0$이다($A$가 정역이기 때문이다).
    따라서 $f/1=0$이고, 이는 $f=0$을 함의한다.
  \item[\rm(ii)] 문제는 $Y$ 위에서 국소적이므로, 다시
    $Y=\operatorname{Spec}(A)$이고 $A$가 정역이며
    $\mathcal{S}=\widetilde{S}$라고 가정할 수 있다. 가정에 따라 모든
    $y\in Y$에 대하여 $(S_y)_+$는 $\neq0$인 멱영원을 포함하지 않으며,
    $(S_0)_y=A_y$도 마찬가지이다. 따라서 $S_y$는 모든 $y\in Y$에
    대하여 축소환이고, 이로부터 먼저 $S$ 자체가 축소환임을 얻는다
    (\hyperref[I.5.1.1-ko]{I, 5.1.1}). 다른 한편 $X$가 정역 스킴이라는
    가정으로부터 $S$가 본질적으로 정역임을 얻는다
    (\hyperref[II.2.4.4-ko]{2.4.4, (ii)}). 따라서 모든 것은 소멸자
    $\mathfrak{J}$, 즉 $S_+$의 $A=S_0$에서의 소멸자가

\oldpage[II]{54}

    $0$임을 보이는 것으로 귀착된다
    (\hyperref[II.2.1.11-ko]{2.1.11}). 그렇지 않으면
    $(S_h)_+=0$인 어떤 $h\neq0$가 $\mathfrak{J}$에 존재하고, 따라서
    (\hyperref[II.3.1.1-ko]{3.1.1})에 의해
    $\varphi^{-1}(D(h))=\varnothing$이다. 그러면 가정과 달리
    $\varphi(X)$는 $Y$에서 조밀하지 않다($D(h)\neq\varnothing$인데,
    이는 $h$가 멱영원이 아니기 때문이다).
\end{enumerate}
\end{proof}

\subsection{등급 $\mathcal{S}$-가군에 대응하는
$\operatorname{Proj}(\mathcal{S})$ 위의 층.}
\label{subsection:II.3.2-ko}

\begin{env}[3.2.1]
\label{II.3.2.1-ko}
$Y$를 준스킴, $\mathcal{S}$를 음이 아닌 차수로 등급이 주어진 준연접
$\mathcal{O}_Y$-대수, $\mathcal{M}$을 준연접인 등급 $\mathcal{S}$-가군이라
하자(준연접성을 $(Y,\mathcal{O}_Y)$ 위에서 생각하든 환 달린 공간
$(Y,\mathcal{S})$ 위에서 생각하든 이는 같은 조건이다
(\hyperref[I.9.6.1-ko]{I, 9.6.1})).
(\hyperref[II.3.1.1-ko]{3.1.1})의 표기법으로, 준연접
$\mathcal{O}_{X_U}$-가군 $\widetilde{\Gamma(U,\mathcal{M})}$을
$\widetilde{\mathcal{M}}_U$로 나타내자. $U'\subset U$이면
$\Gamma(U',\mathcal{M})$은
$\Gamma(U,\mathcal{M})\otimes_A A'$과 표준적으로 동일시된다
(\hyperref[I.1.6.4-ko]{I, 1.6.4}). 따라서
$\widetilde{\mathcal{M}}_{U'}
=\rho_{U',U}^*(\widetilde{\mathcal{M}}_U)$이다
(\hyperref[II.2.8.11-ko]{2.8.11}).
\end{env}

\begin{proposition}[3.2.2]
\label{II.3.2.2-ko}
$\operatorname{Proj}(\mathcal{S})=X$ 위에 준연접
$\mathcal{O}_X$-가군 $\widetilde{\mathcal{M}}$이 유일하게 존재하여,
$Y$의 모든 아핀 열린집합 $U$에 대하여
\[
  \eta_U^*(\widetilde{\Gamma(U,\mathcal{M})})
  =\widetilde{\mathcal{M}}|f^{-1}(U)
\]
가 성립한다(여기서 $\eta_U$로 $Y$-준스킴들의 동형사상
$f^{-1}(U)\xrightarrow{\sim}\operatorname{Proj}(\Gamma(U,\mathcal{S}))$을
나타내며, $f$는 구조 사상 $X\to Y$이다).
\end{proposition}

\begin{proof}
$\rho_{U',U}$가 포함 사상 $f^{-1}(U')\to f^{-1}(U)$과 동일시되므로
(\hyperref[II.3.1.2.1-ko]{3.1.2.1}), 명제는 관계식
$\widetilde{\mathcal{M}}_{U'}=\rho_{U',U}^*(\widetilde{\mathcal{M}}_U)$과
층의 붙이기 원리 (\hyperref[0.3.3.1-ko]{0, 3.3.1})에서 곧바로 따른다.
\end{proof}

$\widetilde{\mathcal{M}}$을 준연접인 등급 $\mathcal{S}$-가군
$\mathcal{M}$에 대응하는 $\mathcal{O}_X$-가군이라 한다.

\begin{proposition}[3.2.3]
\label{II.3.2.3-ko}
$\mathcal{M}$을 준연접인 등급 $\mathcal{S}$-가군이라 하고,
$f\in\Gamma(Y,\mathcal{S}_d)$ ($d>0$)라 하자. $\xi_f$가 $X_f$에서
$Y$-준스킴
$Z_f=\operatorname{Spec}(\mathcal{S}^{(d)}/(f-1)\mathcal{S}^{(d)})$로
가는 표준 동형사상이면 (\hyperref[II.3.1.4-ko]{3.1.4}),
$(\xi_f)_*(\widetilde{\mathcal{M}}|X_f)$는 $\mathcal{O}_{Z_f}$-가군
$\widetilde{\mathcal{M}^{(d)}/(f-1)\mathcal{M}^{(d)}}$이다
(\hyperref[II.1.4.3-ko]{1.4.3}).
\end{proposition}

\begin{proof}
이 문제는 $Y$ 위에서 국소적이므로, 도식
(\hyperref[II.2.8.12.1-ko]{2.8.12.1})의 가환성을 고려하면 곧바로
(\hyperref[II.2.2.5-ko]{2.2.5})로 환원된다.
\end{proof}

\begin{proposition}[3.2.4]
\phantomsection
\label{II.3.2.4-ko}
함자 $\mathcal{M}\mapsto\widetilde{\mathcal{M}}$은 준연접 등급
$\mathcal{S}$-가군의 범주에서 준연접 $\mathcal{O}_X$-가군의 범주로
가는 공변 가법 완전 함자이고, 귀납극한 및 직합과 가환한다.
\end{proposition}

\begin{proof}
이 문제는 $Y$ 위에서 국소적이므로
(\hyperref[I.1.3.11-ko]{I, 1.3.11} 및
\hyperref[I.1.3.9-ko]{I, 1.3.9})와
(\hyperref[II.2.5.4-ko]{2.5.4})로 환원된다.
\end{proof}

특히 $\mathcal{N}$이 $\mathcal{M}$의 준연접 등급
부분-$\mathcal{S}$-가군이면, $\widetilde{\mathcal{N}}$은
$\widetilde{\mathcal{M}}$의 준연접 부분-$\mathcal{O}_X$-가군과
표준적으로 동일시된다. 그중에서도 $\mathcal{S}$의 임의의 준연접 등급
아이디얼층 $\mathcal{J}$에 대하여, $\widetilde{\mathcal{J}}$는
$\mathcal{O}_X$의 준연접 아이디얼층이다.

$\mathcal{M}$이 준연접 등급 $\mathcal{S}$-가군이고
$\mathcal{I}$가 $\mathcal{O}_Y$의 준연접 아이디얼층이면,
$\mathcal{I}\mathcal{M}$은 $\mathcal{M}$의 준연접 등급
부분-$\mathcal{S}$-가군이며 다음이 성립한다.
\[
\label{II.3.2.4.1-ko}
  \widetilde{\mathcal{I}\mathcal{M}}
  =\mathcal{I}\cdot\widetilde{\mathcal{M}}
\tag{3.2.4.1}
\]
(여기서 우변은 (\hyperref[0.4.3.5-ko]{0, 4.3.5})에서 정의한 뜻을
갖는다.) 실제로 $Y=\operatorname{Spec}(A)$가 아핀이고,
$\mathcal{S}=\widetilde{S}$이며 $S$는 등급 $A$-대수이고,
$\mathcal{M}=\widetilde{M}$이고,

\oldpage[II]{55}

$M$은 등급 $S$-가군이며, $\mathcal{I}=\widetilde{\mathfrak{I}}$이고
$\mathfrak{I}$는 $A$의 아이디얼인 경우에 이 공식을 확인하면 충분하다.
$S_+$의 임의의 동차 원소 $f$에 대하여,
(\hyperref[II.3.2.4.1-ko]{3.2.4.1})의 좌변을
$D_+(f)=\operatorname{Spec}(S_{(f)})$에 제한한 것은
$(\mathfrak{I}M)_{(f)}=\mathfrak{I}\cdot M_{(f)}$에 대응하는 층이다.
(\hyperref[I.1.3.13-ko]{I, 1.3.13} 및
\hyperref[I.1.6.9-ko]{I, 1.6.9})에 따라 우변을 제한한 것도 마찬가지이다.

\begin{proposition}[3.2.5]
\phantomsection
\label{II.3.2.5-ko}
$f\in\Gamma(Y,\mathcal{S}_d)$ ($d>0$)라 하자. 열린집합 $X_f$ 위에서
$(\mathcal{O}_X|X_f)$-가군 $\widetilde{\mathcal{S}(nd)}|X_f$는 모든
$n\in\mathbf{Z}$에 대하여 $\mathcal{O}_X|X_f$와 표준적으로 동형이다.
특히 $\mathcal{O}_Y$-대수 $\mathcal{S}$가 $\mathcal{S}_1$에 의해
생성되면 (\hyperref[II.3.1.9-ko]{3.1.9}), $\mathcal{O}_X$-가군
$\widetilde{\mathcal{S}(n)}$은 모든 $n\in\mathbf{Z}$에 대하여
가역이다.
\end{proposition}

\begin{proof}
실제로 $Y$의 모든 아핀 열린집합 $U$에 대하여,
(\hyperref[II.2.5.7-ko]{2.5.7})에서
$\widetilde{\mathcal{S}(nd)}|(X_f\cap\varphi^{-1}(U))$에서
$\mathcal{O}_X|(X_f\cap\varphi^{-1}(U))$로 가는 표준 동형을
(\hyperref[II.3.1.4-ko]{3.1.4})를 고려하여 정의하였다(여기서
$\varphi$는 구조 사상 $X\to Y$이다). $U$에서 아핀 열린집합
$U'\subset U$로 제한할 때 이 동형들이 서로 양립함은 곧바로 확인할 수
있으므로 첫째 주장이 따른다. 둘째 주장을 보이기 위해서는,
$\mathcal{S}$가 $\mathcal{S}_1$에 의해 생성되면 $Y$에 아핀 열린집합들로
이루어진 덮개 $(U_\alpha)$가 존재하여
$\Gamma(U_\alpha,\mathcal{S})$가
$\Gamma(U_\alpha,\mathcal{S})_1=\Gamma(U_\alpha,\mathcal{S}_1)$에
의해 생성됨을 관찰하면 충분하다. 그러면
(\hyperref[II.2.5.9-ko]{2.5.9})의 결과를 적용하는 경우로 환원되는데,
가역이라는 성질은 국소적이기 때문이다.
\end{proof}

또 모든 $n\in\mathbf{Z}$에 대하여
\[
\label{II.3.2.5.1-ko}
  \mathcal{O}_X(n)=\widetilde{\mathcal{S}(n)}
\tag{3.2.5.1}
\]
로 두고, 모든 $\mathcal{O}_X$-가군 $\mathcal{F}$에 대하여
\[
\label{II.3.2.5.2-ko}
  \mathcal{F}(n)=\mathcal{F}\otimes_{\mathcal{O}_X}\mathcal{O}_X(n).
\tag{3.2.5.2}
\]
로 둔다. 이 정의들에서 $Y$의 모든 열린집합 $U$에 대하여 다음이 곧바로
따른다.
\[
  \widetilde{(\mathcal{S}|U)(n)}
  =\mathcal{O}_X(n)|f^{-1}(U)
\]
여기서 $f$는 구조 사상 $X\to Y$이다.

\begin{proposition}[3.2.6]
\phantomsection
\label{II.3.2.6-ko}
$\mathcal{M}$과 $\mathcal{N}$을 두 준연접 등급 $\mathcal{S}$-가군이라
하자. 그러면 ($\mathcal{M}$과 $\mathcal{N}$에 관해) 함자적인 표준 준동형
\[
\label{II.3.2.6.1-ko}
  \lambda:\widetilde{\mathcal{M}}
  \otimes_{\mathcal{O}_X}\widetilde{\mathcal{N}}
  \longrightarrow\widetilde{\mathcal{M}\otimes_{\mathcal{S}}\mathcal{N}}
\tag{3.2.6.1}
\]
이 존재한다. 또한 ($\mathcal{M}$과 $\mathcal{N}$에 관해) 함자적인 표준
준동형
\[
\label{II.3.2.6.2-ko}
  \mu:\widetilde{\mathcal{H}om_{\mathcal{S}}(\mathcal{M},\mathcal{N})}
  \longrightarrow
  \mathcal{H}om_{\mathcal{O}_X}
  (\widetilde{\mathcal{M}},\widetilde{\mathcal{N}}).
\tag{3.2.6.2}
\]
이 존재한다. 더 나아가 $\mathcal{S}$가 $\mathcal{S}_1$에 의해 생성되면
(\hyperref[II.3.1.9-ko]{3.1.9}), $\lambda$는 동형사상이다~; 또한
$\mathcal{M}$이 유한 표시를 가지면
(\hyperref[II.3.1.1-ko]{3.1.1}), $\mu$는 동형사상이다.
\end{proposition}

\begin{proof}
$Y$가 아핀일 때 준동형 $\lambda$와 $\mu$는
(\hyperref[II.2.5.11.2-ko]{2.5.11.2})와
(\hyperref[II.2.5.12.2-ko]{2.5.12.2})에서 정의되었다~; 이 정의들은
국소적이므로 (\hyperref[II.2.8.14-ko]{2.8.14})을 고려하면 여기서 다루는
일반적인 경우로 곧바로 옮겨진다.
\end{proof}

\begin{corollary}[3.2.7]
\phantomsection
\label{II.3.2.7-ko}
$\mathcal{S}$가 $\mathcal{S}_1$에 의해 생성되면, $\mathbf{Z}$에 속하는
임의의 $m$, $n$에 대하여
\[
\label{II.3.2.7.1-ko}
  \mathcal{O}_X(m)\otimes_{\mathcal{O}_X}\mathcal{O}_X(n)
  =\mathcal{O}_X(m+n)
\tag{3.2.7.1}
\]
\[
\label{II.3.2.7.2-ko}
  \mathcal{O}_X(n)=(\mathcal{O}_X(1))^{\otimes n}
\tag{3.2.7.2}
\]
가 성립한다. 여기서 등식들은 표준 동형까지 성립한다.
\end{corollary}

\oldpage[II]{56}

\begin{corollary}[3.2.8]
\phantomsection
\label{II.3.2.8-ko}
$\mathcal{S}$가 $\mathcal{S}_1$에 의해 생성된다고 가정하자. 모든 등급
$\mathcal{S}$-가군 $\mathcal{M}$과 모든 $n\in\mathbf{Z}$에 대하여
\[
\label{II.3.2.8.1-ko}
  (\mathcal{M}(n))^{\sim}=\widetilde{\mathcal{M}}(n)
\tag{3.2.8.1}
\]
가 표준 동형까지 성립한다.
\end{corollary}

\begin{proof}
이는 $Y$가 아핀일 때의 대응하는 성질들
(\hyperref[II.2.5.14-ko]{2.5.14}와
\hyperref[II.2.5.15-ko]{2.5.15}) 및
(\hyperref[II.2.8.11-ko]{2.8.11})에서 따른다.
\end{proof}

\par\medskip\noindent\textit{비고 (3.2.9). ---\ }
\label{II.3.2.9-ko}
\begin{enumerate}
  \item[\rm(i)] $\mathcal{S}=\mathcal{A}[T]$이고 $\mathcal{A}$가 준연접
    $\mathcal{O}_Y$-대수이면
    (\hyperref[II.3.1.7-ko]{3.1.7}), 모든 가역
    $\mathcal{O}_X$-가군 $\mathcal{O}_X(n)$은 $\mathcal{O}_X$와
    표준적으로 동형임을 곧바로 확인할 수 있다.

    더 나아가 $\mathcal{N}$을 준연접 $\mathcal{A}$-가군이라 하고
    $\mathcal{M}=\mathcal{N}\otimes_{\mathcal{A}}\mathcal{A}[T]$라 놓자.
    그러면 (\hyperref[II.3.2.3-ko]{3.2.3})과
    (\hyperref[II.3.1.7-ko]{3.1.7})에 따라
    $X=\operatorname{Proj}(\mathcal{A}[T])$와
    $X'=\operatorname{Spec}(\mathcal{A})$의 표준적 동일시에서
    $\mathcal{O}_X$-가군 $\widetilde{\mathcal{M}}$은 $\mathcal{N}$에
    대응하는 $\mathcal{O}_{X'}$-가군 $\widetilde{\mathcal{N}}$과
    동일시된다((\hyperref[II.1.4.3-ko]{1.4.3})의 뜻에서).
  \item[\rm(ii)] $\mathcal{S}$를 임의의 등급 $\mathcal{O}_Y$-대수라 하고,
    $\mathcal{S}'_0=\mathcal{O}_Y$이며 모든 $n>0$에 대하여
    $\mathcal{S}'_n=\mathcal{S}_n$인 등급 $\mathcal{O}_Y$-대수를
    $\mathcal{S}'$이라 하자~; $X=\operatorname{Proj}(\mathcal{S})$에서
    $X'=\operatorname{Proj}(\mathcal{S}')$로 가는 표준 동형사상
    (\hyperref[II.3.1.8-ko]{3.1.8, (ii)})은 모든 $n\in\mathbf{Z}$에 대하여
    $\mathcal{O}_X(n)$과 $\mathcal{O}_{X'}(n)$을 동일시한다~: 이는 아핀
    경우에 대한 같은 명제 (\hyperref[II.2.5.16-ko]{2.5.16})와 $Y$의 아핀
    열린집합들에 대한 이 동일시들이 제한 연산과 가환한다는 사실에서 따른다.
    마찬가지로 $X^{(d)}=\operatorname{Proj}(\mathcal{S}^{(d)})$라 하자~;
    $X$에서 $X^{(d)}$로 가는 표준 동형사상
    (\hyperref[II.3.1.8-ko]{3.1.8, (i)})은 모든 $n\in\mathbf{Z}$에 대하여
    $\mathcal{O}_X(nd)$와 $\mathcal{O}_{X^{(d)}}(n)$을 동일시한다.
\end{enumerate}

\begin{proposition}[3.2.10]
\phantomsection
\label{II.3.2.10-ko}
$\mathcal{L}$을 가역 $\mathcal{O}_Y$-가군이라 하고, $g$를 표준 동형사상
$X_{(\mathcal{L})}=\operatorname{Proj}(\mathcal{S}_{(\mathcal{L})})
\to X=\operatorname{Proj}(\mathcal{S})$
(\hyperref[II.3.1.8-ko]{3.1.8, (iii)})라 하자. 모든 정수
$n\in\mathbf{Z}$에 대하여
$g_*(\mathcal{O}_{X_{(\mathcal{L})}}(n))$은
$\mathcal{O}_X(n)\otimes_Y\mathcal{L}^{\otimes n}$과 표준적으로 동형이다.
\end{proposition}

\begin{proof}
먼저 $Y$가 환 $A$의 아핀 스킴이고 $\mathcal{L}=\widetilde{L}$이라고
가정하자. 여기서 $L$은 생성원 하나를 갖는 자유 $A$-가군이다.
(\hyperref[II.3.1.8-ko]{3.1.8, (iii)})의 증명에서 사용한 표기법 아래에서,
$f\in S_d$에 대하여
$(S(n))_{(f)}\otimes_A L^{\otimes n}$에서
$(S_{(L)}(n))_{(f\otimes c^{\otimes d})}$로 가는 동형을 다음과 같이
정의한다. 즉 $x\in S_{kd+n}$일 때
$(x/f^k)\otimes c^{\otimes n}$을
$(x\otimes c^{\otimes(n+kd)})/(f\otimes c^{\otimes d})^k$에 대응시킨다.
이 동형이 $L$의 생성원으로 택한 $c$와 무관함은 명백하다. 또한 각
$f\in S_+$에 대해 이와 같이 정의된 동형들은 제한 연산
$D_+(f)\to D_+(fg)$과 양립한다. 마지막으로 일반적인 경우에는 정의
(\hyperref[II.3.1.1-ko]{3.1.1})로 돌아가 보면, $Y$의 각 아핀 열린집합
$U$에 대해 이와 같이 정의된 동형들이 아핀 열린집합 $U'\subset U$로의
제한과 양립함을 쉽게 알 수 있다.
\end{proof}

\subsection{$\operatorname{Proj}(\mathcal{S})$ 위의 층에 대응하는 등급
$\mathcal{S}$-가군.}
\label{subsection:II.3.3-ko}

\emph{이 절 전체에서 등급 $\mathcal{O}_Y$-대수 $\mathcal{S}$가
$\mathcal{S}_1$에 의해 생성된다고 가정한다
(\hyperref[II.3.1.9-ko]{3.1.9}).}
(\hyperref[II.3.1.8-ko]{3.1.8, (i)})에 따라, 유한성 조건
(\hyperref[II.3.1.10-ko]{3.1.10})을 전제로 하면 이 제약은 본질적이지
않음을 상기하자.

\begin{env}[3.3.1]
\phantomsection
\label{II.3.3.1-ko}
$p$를 구조 사상 $X=\operatorname{Proj}(\mathcal{S})\to Y$라 하자.
모든 $\mathcal{O}_X$-가군 $\mathcal{F}$에 대하여 다음과 같이 둔다.
\[
\label{II.3.3.1.1-ko}
  \Gamma_*(\mathcal{F})
  =\bigoplus_{n\in\mathbf{Z}}p_*(\mathcal{F}(n))
\tag{3.3.1.1}
\]

\oldpage[II]{57}

그리고 특히
\[
\label{II.3.3.1.2-ko}
  \Gamma_*(\mathcal{O}_X)
  =\bigoplus_{n\in\mathbf{Z}}p_*(\mathcal{O}_X(n)).
\tag{3.3.1.2}
\]
두 $\mathcal{O}_X$-가군 $\mathcal{F}$와 $\mathcal{G}$에 대하여 표준
준동형
\[
  p_*(\mathcal{F})\otimes_{\mathcal{O}_Y}p_*(\mathcal{G})
  \longrightarrow
  p_*(\mathcal{F}\otimes_{\mathcal{O}_X}\mathcal{G})
\]
이 존재함이 알려져 있다
(\hyperref[0.4.2.2-ko]{0, 4.2.2}). 따라서
(\hyperref[II.3.2.7.1-ko]{3.2.7.1})로부터
$\Gamma_*(\mathcal{O}_X)$에 \emph{등급 $\mathcal{O}_Y$-대수} 구조가 주어지며,
(\hyperref[II.3.2.5.2-ko]{3.2.5.2})는 마찬가지로
$\Gamma_*(\mathcal{F})$에 $\Gamma_*(\mathcal{O}_X)$ 위의 \emph{등급가군}
구조를 정의한다.

(\hyperref[II.3.2.5-ko]{3.2.5})와 함자 $p_*$의 왼쪽 완전성
(\hyperref[0.4.2.1-ko]{0, 4.2.1})에 따라,
$\mathcal{F}\mapsto\Gamma_*(\mathcal{F})$는 $\mathcal{O}_X$-가군의
범주에서 등급 $\mathcal{O}_Y$-가군의 범주로 가는 공변 가법
\emph{왼쪽 완전} 함자이다. 여기서 사상들은 차수 $0$의 준동형이다. 특히
$\mathcal{J}$가 $\mathcal{O}_X$의 아이디얼층이면,
$\Gamma_*(\mathcal{J})$는 $\Gamma_*(\mathcal{O}_X)$의
\emph{등급 아이디얼층}과 동일시된다.
\end{env}

\begin{env}[3.3.2]
\label{II.3.3.2-ko}
$\mathcal{M}$을 준연접 등급 $\mathcal{S}$-가군이라 하자. $Y$의 모든
아핀 열린집합 $U$에 대하여, (\hyperref[II.2.6.2-ko]{2.6.2})에서
아벨군의 준동형
\[
  \alpha_{0,U}:\Gamma(U,\mathcal{M}_0)
  \longrightarrow\Gamma(p^{-1}(U),\widetilde{\mathcal{M}}).
\]
을 정의하였다. 이 준동형들이 제한 연산과 가환함은 명백하다
(\hyperref[II.2.8.13.1-ko]{2.8.13.1}). 따라서 ($\mathcal{S}$가
$\mathcal{S}_1$에 의해 생성된다는 가정을 사용하지 않고) 아벨군의 층의
준동형
\[
\label{II.3.3.2.1-ko}
  \alpha_0:\mathcal{M}_0\longrightarrow p_*(\widetilde{\mathcal{M}}).
\tag{3.3.2.1}
\]
을 정의한다. 이 결과를 각각의
$\mathcal{M}_n=(\mathcal{M}(n))_0$에 적용하고
(\hyperref[II.3.2.8.1-ko]{3.2.8.1})을 고려하면, 아벨군의 층의 준동형
\[
\label{II.3.3.2.2-ko}
  \alpha_n:\mathcal{M}_n\longrightarrow
  p_*(\widetilde{\mathcal{M}}(n))
  \qquad\text{모든 }n\in\mathbf{Z}\text{에 대하여},
\tag{3.3.2.2}
\]
을 정의하며, 이로부터 등급 아벨군의 층의 함자적인 준동형(차수~$0$)
\[
\label{II.3.3.2.3-ko}
  \alpha:\mathcal{M}\longrightarrow\Gamma_*(\widetilde{\mathcal{M}})
\tag{3.3.2.3}
\]
을 얻는다(이를 $\alpha_{\mathcal{M}}$으로도 나타낸다).

특히 $\mathcal{M}=\mathcal{S}$로 두면,
$\alpha:\mathcal{S}\to\Gamma_*(\mathcal{O}_X)$는 등급
$\mathcal{O}_Y$-대수의 준동형이고,
(\hyperref[II.3.3.2.3-ko]{3.3.2.3})은 이 등급 대수 준동형에 상대적인
등급 가군의 쌍준동형임을 확인할 수 있다.

또한 각각의 $\alpha_n$에는 (\hyperref[0.4.4.3-ko]{0, 4.4.3})에 따라
표준 $\mathcal{O}_X$-가군 준동형
\[
\label{II.3.3.2.4-ko}
  \alpha_n^\sharp:p^*(\mathcal{M}_n)
  \longrightarrow\widetilde{\mathcal{M}}(n).
\tag{3.3.2.4}
\]
이 대응함을 지적하자. 이 준동형은 다름 아니라
(\hyperref[II.3.2.4-ko]{3.2.4})에 따라 다음 등급
$\mathcal{O}_Y$-가군의 표준 준동형(차수~$0$)에 함자적으로 대응하는
준동형임을 쉽게 확인할 수 있다.
\[
\label{II.3.3.2.5-ko}
  \mathcal{M}_n\otimes_{\mathcal{O}_Y}\mathcal{S}
  \longrightarrow\mathcal{M}(n)
\tag{3.3.2.5}
\]

\oldpage[II]{58}

여기서 둘째 항의 등급은 $\mathcal{S}$의 등급에서 자연스럽게 유도된다.
실제로 $Y=\operatorname{Spec}(A)$가 아핀이고,
$\mathcal{M}=\widetilde{M}$이며 $\mathcal{S}=\widetilde{S}$인 경우로
한정할 수 있다. 이때 등급 $A$-대수 $S$는 $S_1$에 의해 생성되므로,
$f$가 $S_1$을 달릴 때 $D_+(f)$들은 $X$의 덮개를 이룬다. 이제 정의
(\hyperref[II.2.6.2-ko]{2.6.2})로 돌아가서
(\hyperref[I.1.6.7-ko]{I, 1.6.7})을 고려하면, 준동형
(\hyperref[II.3.3.2.4-ko]{3.3.2.4})의 $D_+(f)$로의 제한은
(\hyperref[I.1.3.8-ko]{I, 1.3.8})에 따라 $S_{(f)}$-가군 준동형
$M_n\otimes_A S_{(f)}\to(M(n))_{(f)}$에 대응함을 알 수 있다. 이 준동형은
$x\otimes1$ ($x\in M_n$)에 $x/1$을 대응시킨다. 이로써 주장이 증명된다.
\end{env}

\begin{proposition}[3.3.3]
\phantomsection
\label{II.3.3.3-ko}
임의의 단면 $f\in\Gamma(Y,\mathcal{S}_d)$ ($d>0$)에 대하여, $X_f$는
$\mathcal{O}_X(d)$의 단면으로 본 $\alpha_d(f)$가 소멸하지 않는 $X$의
점들의 집합과 일치한다
(\hyperref[0.5.5.2-ko]{0, 5.5.2}).
\end{proposition}

\begin{proof}
($\alpha_d(f)$는 $Y$ 위의 $p_*(\mathcal{O}_X(d))$의 단면이지만, 정의상
그러한 단면은 $X$ 위의 $\mathcal{O}_X(d)$의 단면이기도 하다
(\hyperref[0.4.2.1-ko]{0, 4.2.1})). $X_f$의 정의
(\hyperref[II.3.1.4-ko]{3.1.4})에 따라 $Y$가 아핀인 경우로 환원되며,
이는 (\hyperref[II.2.6.3-ko]{2.6.3})에서 다루었다.
\end{proof}

\begin{env}[3.3.4]
\phantomsection
\label{II.3.3.4-ko}
이제부터는 이 번호의 처음에 둔 가정 외에도, 모든 준연접
$\mathcal{O}_X$-가군 $\mathcal{F}$에 대하여 $p_*(\mathcal{F}(n))$들이
$Y$ 위에서 \emph{준연접}이라고 가정하자. 그러면
$\Gamma_*(\mathcal{F})=\bigoplus_{n\in\mathbf{Z}}p_*(\mathcal{F}(n))$도
준연접 $\mathcal{O}_Y$-가군이다
(\hyperref[I.1.4.1-ko]{I, 1.4.1} 및
\hyperref[I.1.3.9-ko]{I, 1.3.9}). 특히 $X$가 $Y$ 위에서
\emph{유한형}이면 이 상황은 언제나 성립한다
(\hyperref[I.9.2.2-ko]{I, 9.2.2}). 따라서
$(\Gamma_*(\mathcal{F}))^{\sim}$가 정의되며 이는 준연접
$\mathcal{O}_X$-가군이다. $Y$의 임의의 아핀 열린집합 $U$에 대하여
다음 등식들이 성립한다(\hyperref[I.1.3.9-ko]{I, 1.3.9} 및
\hyperref[II.2.5.4-ko]{2.5.4}).
\begin{align*}
  \left(\Gamma\left(U,\bigoplus_{n\in\mathbf{Z}}
  p_*(\mathcal{F}(n))\right)\right)^{\sim}
  &=\bigoplus_{n\in\mathbf{Z}}
  \bigl(\Gamma(U,p_*(\mathcal{F}(n)))\bigr)^{\sim}\\
  &=\bigoplus_{n\in\mathbf{Z}}
  \bigl(\Gamma(p^{-1}(U),\mathcal{F}(n))\bigr)^{\sim}\\
  &=\left(\bigoplus_{n\in\mathbf{Z}}
  \Gamma(p^{-1}(U),\mathcal{F}(n))\right)^{\sim}\\
  &=\bigl(\Gamma_*(\mathcal{F}|p^{-1}(U))\bigr)^{\sim}
\end{align*}
따라서 (\hyperref[II.2.6.4-ko]{2.6.4})에 의해 표준 준동형
\[
  \beta_U:\left(\Gamma\left(U,\bigoplus_{n\in\mathbf{Z}}
  p_*(\mathcal{F}(n))\right)\right)^{\sim}
  \longrightarrow\mathcal{F}|p^{-1}(U).
\]
을 얻는다.

더 나아가 (\hyperref[II.2.8.13.2-ko]{2.8.13.2})의 가환성으로부터 이
준동형들이 $Y$ 위의 제한 연산들과 교환함을 알 수 있다. 따라서 준연접
$\mathcal{O}_X$-가군에 대한 표준적이고 함자적인 준동형
\[
\label{II.3.3.4.1-ko}
  \beta:(\Gamma_*(\mathcal{F}))^{\sim}\longrightarrow\mathcal{F}
\tag{3.3.4.1}
\]
을 얻으며, 이를 $\beta_{\mathcal{F}}$로도 표기한다.
\end{env}

\begin{proposition}[3.3.5]
\phantomsection
\label{II.3.3.5-ko}
$\mathcal{M}$을 준연접 등급 $\mathcal{S}$-가군,
$\mathcal{F}$를 준연접 $\mathcal{O}_X$-가군이라 하자. 그러면 다음 합성
준동형들은 각각 항등 준동형이다.
\[
\label{II.3.3.5.1-ko}
  \widetilde{\mathcal{M}}
  \xrightarrow{\widetilde{\alpha}}
  (\Gamma_*(\widetilde{\mathcal{M}}))^{\sim}
  \xrightarrow{\beta}\widetilde{\mathcal{M}}
\tag{3.3.5.1}
\]
\[
\label{II.3.3.5.2-ko}
  \Gamma_*(\mathcal{F})\xrightarrow{\alpha}
  \Gamma_*((\Gamma_*(\mathcal{F}))^{\sim})
  \xrightarrow{\Gamma_*(\beta)}\Gamma_*(\mathcal{F})
\tag{3.3.5.2}
\]
\end{proposition}

\begin{proof}
실제로 이 문제는 $Y$ 위에서 국소적이므로
(\hyperref[II.2.6.5-ko]{2.6.5})로 환원된다.
\end{proof}

\subsection{유한성 조건.}
\label{subsection:II.3.4-ko}

\oldpage[II]{59}

\begin{proposition}[3.4.1]
\label{II.3.4.1-ko}
$Y$를 준스킴, $\mathcal{S}$를 $\mathcal{S}_1$에 의해 생성되는 준연접
$\mathcal{O}_Y$-대수라 하자
(\hyperref[II.3.1.9-ko]{3.1.9}). 또한 $\mathcal{S}_1$은 유한 생성이라고
가정한다. 그러면 $X=\operatorname{Proj}(\mathcal{S})$는 $Y$ 위에서
유한형이다.
\end{proposition}

\begin{proof}
실제로 문제는 $Y$ 위에서 국소적이므로, $Y$가 환 $A$의 아핀 스킴이라고
가정할 수 있다. 이때 $\mathcal{S}=\widetilde{S}$이고
$S=\Gamma(Y,\mathcal{S})$이다. 가정에 따라 $S$는
$S_1=\Gamma(Y,\mathcal{S}_1)$에 의해 생성되는 $A$-대수이며, 또한
$S_1$이 유한 생성 $A$-가군이라고 가정할 수 있다
(\hyperref[I.1.3.9-ko]{I, 1.3.9} 및
\hyperref[I.1.3.12-ko]{I, 1.3.12}). 따라서 $S$는 유한 생성 등급
$A$-대수이고, 문제는
(\hyperref[II.2.7.1-ko]{2.7.1, (ii)})로 환원된다.
\end{proof}

\begin{env}[3.4.2]
\label{II.3.4.2-ko}
$\mathcal{S}$를 준연접 등급 $\mathcal{O}_Y$-대수라 하자. 준연접 등급
$\mathcal{S}$-가군 $\mathcal{M}$에 대하여 다음 유한성 조건들을
고려하겠다:
\begin{enumerate}
  \item[(\textbf{TF})] \emph{어떤 정수 $n$이 존재하여
    $\mathcal{S}$-가군 $\bigoplus_{k\geq n}\mathcal{M}_k$가 유한 생성이다.}
  \item[(\textbf{TN})] \emph{어떤 정수 $n$이 존재하여
    $k\geq n$이면 $\mathcal{M}_k=0$이다.}
\end{enumerate}
$\mathcal{M}$이 (\textbf{TN})을 만족하면
$\widetilde{\mathcal{M}}=0$이다. 이는 문제가 $Y$ 위에서 국소적이기
때문이다(\hyperref[II.2.7.2-ko]{2.7.2}).

$\mathcal{M}$, $\mathcal{N}$을 두 준연접 등급 $\mathcal{S}$-가군이라
하자. 차수 $0$의 준동형 $u:\mathcal{M}\to\mathcal{N}$에 대하여, 어떤
정수 $n$이 존재하여 $k\geq n$일 때
$u_k:\mathcal{M}_k\to\mathcal{N}_k$가 단사(각각 전사, 전단사)이면
$u$를 (\textbf{TN})-\emph{단사}(각각 (\textbf{TN})-\emph{전사},
(\textbf{TN})-\emph{전단사})라고 한다. 그러면 문제가 $Y$ 위에서
국소적이라는 점과 (\hyperref[I.1.3.9-ko]{I, 1.3.9})를 고려하면,
(\hyperref[II.2.7.2-ko]{2.7.2})에 따라
$\widetilde{u}:\widetilde{\mathcal{M}}\to\widetilde{\mathcal{N}}$는
단사(각각 전사, 전단사)이다. $u$가 (\textbf{TN})-전단사일 때 $u$를
(\textbf{TN})-\emph{동형사상}이라고도 한다.
\end{env}

\begin{proposition}[3.4.3]
\label{II.3.4.3-ko}
$Y$를 준스킴, $\mathcal{S}$를 $\mathcal{S}_1$에 의해 생성되는 준연접
등급 $\mathcal{O}_Y$-대수라 하고, $\mathcal{S}_1$은 유한 생성이라고
가정하자. $\mathcal{M}$을 준연접 등급 $\mathcal{S}$-가군이라 하자.
\begin{enumerate}
  \item[\rm(i)] $\mathcal{M}$이 조건 (\textbf{TF})를 만족하면
    $\widetilde{\mathcal{M}}$은 유한 생성이다.
  \item[\rm(ii)] $\mathcal{M}$이 (\textbf{TF})를 만족한다고 가정하자.
    $\widetilde{\mathcal{M}}=0$일 필요충분조건은 $\mathcal{M}$이
    (\textbf{TN})을 만족하는 것이다.
\end{enumerate}
\end{proposition}

\begin{proof}
두 주장은 $Y$ 위에서 국소적이므로, $Y$가 환 $A$의 아핀 스킴이고
$\mathcal{S}=\widetilde{S}$이며, 여기서 $S$는 아이디얼 $S_+$가 유한
생성인 등급 $A$-대수이고, $\mathcal{M}=\widetilde{M}$이며, 여기서
$M$은 등급 $S$-가군인 경우로 환원된다. 이때 명제는
(\hyperref[II.2.7.3-ko]{2.7.3})에서 따른다.
\end{proof}

\begin{theorem}[3.4.4]
\label{II.3.4.4-ko}
$Y$를 준스킴이라 하고, $\mathcal{S}$를 $\mathcal{S}_1$로 생성되는
준연접 등급 $\mathcal{O}_Y$-대수라 하자. 또한 $\mathcal{S}_1$은 유한
생성 $\mathcal{O}_Y$-가군이라고 가정하고,
$X=\operatorname{Proj}(\mathcal{S})$라 하자. 모든
준연접 $\mathcal{O}_X$-가군 $\mathcal{F}$에 대하여 표준 준동형
$\beta$ (\hyperref[II.3.3.4-ko]{3.3.4})는 동형이다.
\end{theorem}

\begin{proof}
먼저 $\beta$는 (\hyperref[II.3.4.1-ko]{3.4.1})에 의해 정의됨에 유의하자.
$\beta$가 동형임을 보이기 위해서는 $Y$가 환 $A$의 아핀 스킴이고,
$\mathcal{S}=\widetilde{S}$이며, $S$가 $S_1$로 생성되는 등급 $A$-대수이고
$S_1$이 유한 생성 $A$-가군인 경우로 환원하면 된다. 이제
(\hyperref[II.2.7.5-ko]{2.7.5})를 적용하면 충분하다.
\end{proof}

\begin{corollary}[3.4.5]
\label{II.3.4.5-ko}
(\hyperref[II.3.4.4-ko]{3.4.4})의 가정 아래에서 모든 준연접
$\mathcal{O}_X$-가군 $\mathcal{F}$는 $\widetilde{\mathcal{M}}$ 꼴의
$\mathcal{O}_X$-가군과 동형이다. 여기서 $\mathcal{M}$은 준연접 등급
$\mathcal{S}$-가군이다. 더 나아가 $\mathcal{F}$가 유한 생성이고,
$Y$가 준콤팩트 스킴이거나 $Y$의 바탕공간이 뇌터라고 가정하면,
$\mathcal{M}$을 유한 생성 등급 $\mathcal{S}$-가군으로 택할 수 있다.
\end{corollary}

\oldpage[II]{60}

\begin{proof}
첫째 주장은 $\mathcal{M}=\Gamma_*(\mathcal{F})$로 놓고
(\hyperref[II.3.4.4-ko]{3.4.4})를 적용하면 곧바로 따른다. 둘째 주장을
증명하기 위해서는 $\mathcal{M}$이 그 유한 생성 등급
부분-$\mathcal{S}$-가군 $\mathcal{N}_\lambda$들의 귀납극한임을 보이면
충분하다. 실제로 이로부터 $\widetilde{\mathcal{M}}$이
$\widetilde{\mathcal{N}}_\lambda$들의 귀납극한임이 따르고
(\hyperref[II.3.2.4-ko]{3.2.4}), 따라서 $\mathcal{F}$는
$\beta(\widetilde{\mathcal{N}}_\lambda)$들의 귀납극한이다. 그런데 $X$가
준콤팩트이고 ((\hyperref[II.3.4.1-ko]{3.4.1}) 및
\hyperref[I.6.3.1-ko]{I, 6.3.1}) $\mathcal{F}$가 유한 생성이므로,
$\mathcal{F}$는 반드시 $\beta(\widetilde{\mathcal{N}}_\lambda)$ 중 하나와
같다 (\hyperref[0.5.2.3-ko]{0, 5.2.3}).

$\mathcal{M}$을 귀납극한으로 갖는 $\mathcal{N}_\lambda$들을 정의하기
위해서는 각 $n\in\mathbf{Z}$에 대하여 준연접 $\mathcal{O}_Y$-가군
$\mathcal{M}_n$을 생각하면 충분하다. $Y$에 관한 가정에 의해 이 가군은
그 유한 생성 부분-$\mathcal{O}_Y$-가군
$\mathcal{M}_n^{(\mu_n)}$들의 귀납극한이다
(\hyperref[I.9.4.9-ko]{I, 9.4.9}).
$\mathcal{P}_{\mu_n}=\mathcal{S}\cdot\mathcal{M}_n^{(\mu_n)}$이 유한 생성
등급 $\mathcal{S}$-가군임은 명백하다. 이제 $\mathcal{N}_\lambda$로
$\mathcal{P}_{\mu_n}$ 꼴의 $\mathcal{S}$-가군들의 유한합들을 취하면
원하는 조건이 충족됨을 곧바로 확인할 수 있다.
\end{proof}

\begin{corollary}[3.4.6]
\label{II.3.4.6-ko}
(\hyperref[II.3.4.4-ko]{3.4.4})의 가정이 성립하고, 더 나아가 $Y$의
바탕공간이 준콤팩트라고 가정하자. $\mathcal{F}$를 유한 생성 준연접
$\mathcal{O}_X$-가군이라 하자. 그러면 어떤 정수 $n_0$가 존재하여 모든
$n\geq n_0$에 대하여 표준 준동형
$\sigma:p^*(p_*(\mathcal{F}(n)))\to\mathcal{F}(n)$
(\hyperref[0.4.4.3-ko]{0, 4.4.3})가 전사이다.
\end{corollary}

\begin{proof}
실제로 모든 $y\in Y$에 대하여 $U$를 $Y$에서 $y$의 아핀 열린 근방이라
하자. 어떤 정수 $n_0(U)$가 존재하여 $n\geq n_0(U)$일 때
$\mathcal{F}(n)|p^{-1}(U)$는 $p^{-1}(U)$ 위의 그 단면들 가운데 유한 개로
생성된다 (\hyperref[II.2.7.9-ko]{2.7.9}). 그런데 이 단면들은
$p^*(p_*(\mathcal{F}(n)))$의 $p^{-1}(U)$ 위의 단면들의 표준 상들이므로
(\hyperref[0.3.7.1-ko]{0, 3.7.1} 및
\hyperref[0.4.4.3-ko]{0, 4.4.3}),
$\mathcal{F}(n)|p^{-1}(U)$는
$p^*(p_*(\mathcal{F}(n)))|p^{-1}(U)$의 표준 상과 같다. 마지막으로 $Y$가
준콤팩트이므로 아핀 열린집합 $U_i$들로 이루어진 $Y$의 유한 덮개가
존재한다. $n_0$로 $n_0(U_i)$들 가운데 가장 큰 것을 취하면 증명이
끝난다.
\end{proof}

\par\medskip\noindent\textit{비고 (3.4.7). ---\ }
\label{II.3.4.7-ko}
환 달린 공간의 사상 $p=(\psi,\theta):X\to Y$와
$\mathcal{O}_X$-가군 $\mathcal{F}$가 주어졌다고 하자. 표준 준동형
$\sigma:p^*(p_*(\mathcal{F}))\to\mathcal{F}$가 전사라는 사실은 다음과
같이 풀어 쓸 수 있다(\hyperref[0.4.4.1-ko]{0, 4.4.1}). 임의의 $x\in X$와
$x$의 열린 근방 $V$ 위의 $\mathcal{F}$의 임의의 단면 $s$에 대하여,
$Y$에서 $p(x)$의 열린 근방 $U$, $p^{-1}(U)$ 위의 $\mathcal{F}$의 유한 개
단면 $t_i$ ($1\leq i\leq m$), $x$의 열린 근방
$W\subset V\cap p^{-1}(U)$, 그리고 $W$ 위의 $\mathcal{O}_X$의 단면 $a_i$
($1\leq i\leq m$)가 존재하여
\[
  s|W=\sum_i a_i\cdot(t_i|W).
\]
$Y$가 \emph{아핀 스킴}이고 $p_*(\mathcal{F})$가 \emph{준연접}이면, 이 조건은
$\mathcal{F}$가 \emph{$X$ 위의 단면들로 생성된다}는 것과 동치이다
(\hyperref[0.5.5.1-ko]{0, 5.5.1}). 실제로
$Y=\operatorname{Spec}(A)$이면 $f\in A$인 $U=D(f)$라고 가정할 수 있다.
어떤 정수 $n>0$과 $X$ 위의 $\mathcal{F}$의 단면 $s_i$들이 존재하여,
$g=\theta(f)$라 놓으면 $t_i$는 $s_ig^n$을 $p^{-1}(U)$에 제한한 단면이다
($p_*(\mathcal{F})$에 (\hyperref[I.1.4.1-ko]{I, 1.4.1})을 적용한다).
$g$가 $p^{-1}(U)$ 위에서 가역이므로
\[
  s|W=\sum_i b_i\cdot(s_i|W)
\]
이고, 여기서 $b_i=a_i(g|W)^{-n}$이다. 이로써 주장이 따른다. 따라서
$Y$가 아핀이면 따름정리 (\hyperref[II.3.4.6-ko]{3.4.6})는
(\hyperref[II.2.7.9-ko]{2.7.9})를 다시 준다.

\oldpage[II]{61}

따라서 $Y$가 임의의 준스킴일 때, $p_*(\mathcal{F})$가
\emph{준연접}인 준연접 $\mathcal{O}_X$-가군 $\mathcal{F}$에 대하여 다음 세
조건은 동치이다.
\begin{enumerate}
  \item[a)] \emph{표준 준동형
    $\sigma:p^*(p_*(\mathcal{F}))\to\mathcal{F}$는 전사이다.}
  \item[b)] \emph{준연접 $\mathcal{O}_Y$-가군 $\mathcal{G}$와 전사 준동형
    $p^*(\mathcal{G})\to\mathcal{F}$가 존재한다.}
  \item[c)] \emph{$Y$의 임의의 아핀 열린집합 $U$에 대하여,
    $\mathcal{F}|p^{-1}(U)$는 $p^{-1}(U)$ 위의 단면들로 생성된다.}
\end{enumerate}
실제로 방금 \emph{a)}와 \emph{c)}가 동치임을 증명하였다. 한편 가정에 따라
$p_*(\mathcal{F})$가 준연접이므로 \emph{a)}가 \emph{b)}를 함의함은 명백하다.
반대로 모든 준동형 $u:p^*(\mathcal{G})\to\mathcal{F}$는
$p^*(\mathcal{G})\to p^*(p_*(\mathcal{F}))\xrightarrow{\sigma}\mathcal{F}$로
인수분해된다(\hyperref[0.3.5.4.4-ko]{0, 3.5.4.4}). 따라서 $u$가 전사이면
$\sigma$도 전사이고, 이로써 \emph{b)}가 \emph{a)}를 함의함이 증명된다.

\begin{corollary}[3.4.8]
\label{II.3.4.8-ko}
(\hyperref[II.3.4.4-ko]{3.4.4})의 가정이 성립한다고 하고, 더 나아가
$Y$가 준콤팩트 스킴이거나 $Y$의 바탕공간이 뇌터라고 가정하자.
$\mathcal{F}$를 유한 생성 준연접 $\mathcal{O}_X$-가군이라 하자. 그러면
어떤 정수 $n_0$가 존재하여, 모든 $n\geq n_0$에 대하여 $\mathcal{F}$는
$(p^*(\mathcal{G}))(-n)$ 꼴의 $\mathcal{O}_X$-가군의 몫과 동형이다. 여기서
$\mathcal{G}$는 유한 생성 준연접 $\mathcal{O}_Y$-가군이다
($n$에 의존한다).
\end{corollary}

\begin{proof}
구조 사상 $X\to Y$가 분리 사상이고 유한형이므로,
$p_*(\mathcal{F}(n))$은 준연접이다
(\hyperref[I.9.2.2-ko]{I, 9.2.2, b)}). 따라서 $Y$에 관한 가정에 의해 이는
그 유한 생성 준연접 $\mathcal{O}_Y$-부분가군들의 귀납극한이다
(\hyperref[I.9.4.9-ko]{I, 9.4.9}). 이로부터
(\hyperref[II.3.4.6-ko]{3.4.6}),
(\hyperref[0.4.3.2-ko]{0, 4.3.2}) 및
(\hyperref[0.5.2.3-ko]{0, 5.2.3})에 의해 $\mathcal{F}(n)$은
$p^*(\mathcal{G})$ 꼴의 $\mathcal{O}_X$-가군의 표준적인 상임을 얻는다.
여기서 $\mathcal{G}$는 $p_*(\mathcal{F}(n))$의 유한 생성 준연접
$\mathcal{O}_Y$-부분가군이다. 그러면 따름정리는
(\hyperref[II.3.2.5.2-ko]{3.2.5.2})와
(\hyperref[II.3.2.7.1-ko]{3.2.7.1})에서 따른다.
\end{proof}

\subsection{함자적 성질.}
\label{subsection:II.3.5-ko}

\begin{env}[3.5.1]
\label{II.3.5.1-ko}
$Y$를 준스킴이라 하고, $\mathcal{S}$, $\mathcal{S}'$을 음이 아닌 차수로
등급이 주어진 두 준연접 $\mathcal{O}_Y$-대수라 하자. 또한
$X=\operatorname{Proj}(\mathcal{S})$, $X'=\operatorname{Proj}(\mathcal{S}')$로
놓고, $p$, $p'$을 각각 $X$, $X'$에서 $Y$로 가는 구조 사상이라 하자.
$\varphi:\mathcal{S}'\to\mathcal{S}$를 등급 $\mathcal{O}_Y$-대수의
준동형이라 하자. $Y$의 모든 아핀 열린집합 $U$에 대하여
$S_U=\Gamma(U,\mathcal{S})$, $S'_U=\Gamma(U,\mathcal{S}')$로 놓자.
준동형 $\varphi$는 등급 $A_U$-대수의 준동형
$\varphi_U:S'_U\to S_U$를 정의한다. 여기서
$A_U=\Gamma(U,\mathcal{O}_Y)$로 놓는다. 이에 대응하여 $p^{-1}(U)$ 안에
열린집합 $G(\varphi_U)$와 사상
$\Phi_U:G(\varphi_U)\to p'^{-1}(U)$가 존재한다
(\hyperref[II.2.8.1-ko]{2.8.1}). 더욱이 $V\subset U$가 아핀
열린집합이면 다음 도식은 가환한다.
\[
\label{II.3.5.1.1-ko}
  \xymatrix{
    S'_U\ar[r]^{\varphi_U}\ar[d]
    & S_U\ar[d]\\
    S'_V\ar[r]_{\varphi_V}
    & S_V
  }
\tag{3.5.1.1}
\]
또한 정의들(\hyperref[II.2.8.1-ko]{2.8.1})로부터
\[
  G(\varphi_V)=G(\varphi_U)\cap p^{-1}(V)
\]
임과 $\Phi_V$가 $\Phi_U$를 $G(\varphi_V)$로 제한한 것임을 곧바로
확인할 수 있다. 이로써 모든 아핀 열린집합 $U\subset Y$에 대하여
$G(\varphi)\cap p^{-1}(U)=G(\varphi_U)$를 만족하는 $X$의 열린부분집합
$G(\varphi)$와 $Y$-사상

\oldpage[II]{62}

$\Phi:G(\varphi)\to X'$을 정의하였다. 이 사상은 아핀이며,
\emph{$\varphi$에 대응하는 사상}이라 하고
$\operatorname{Proj}(\varphi)$로 나타내겠다. 모든 $y\in Y$에 대하여
$\Gamma(U,\mathcal{O}_Y)$-가군 $\Gamma(U,\mathcal{S}_+)$이
$\varphi(\Gamma(U,\mathcal{S}'_+))$로 생성되도록 하는 $y$의 아핀 근방
$U$가 존재하면 $G(\varphi_U)=p^{-1}(U)$이고, 따라서
$G(\varphi)=X$이다.
\end{env}

\begin{proposition}[3.5.2]
\label{II.3.5.2-ko}
\begin{enumerate}
  \item[\rm(i)] $\mathcal{M}$이 준연접 등급 $\mathcal{S}$-가군이면,
    $\mathcal{O}_{X'}$-가군 $(\mathcal{M}_{[\varphi]})^{\sim}$에서
    $\mathcal{O}_{X'}$-가군
    $\Phi_*(\widetilde{\mathcal{M}}|G(\varphi))$로 가는 함자적인 표준
    동형이 존재한다.
  \item[\rm(ii)] $\mathcal{M}'$이 준연접 등급 $\mathcal{S}'$-가군이면,
    $(\mathcal{O}_X|G(\varphi))$-가군
    $\Phi^*(\widetilde{\mathcal{M}'})$에서
    $(\mathcal{O}_X|G(\varphi))$-가군
    $(\mathcal{M}'\otimes_{\mathcal{S}'}\mathcal{S})^{\sim}|G(\varphi)$로
    가는 함자적인 표준 준동형 $\nu$가 존재한다. $\mathcal{S}'$이
    $\mathcal{S}'_1$로 생성되면 $\nu$는 동형이다.
\end{enumerate}
\end{proposition}

\begin{proof}
실제로 $Y$가 아핀일 때에는 여기서 고려하는 준동형들이 이미
(\hyperref[II.2.8.7-ko]{2.8.7} 및
\hyperref[II.2.8.8-ko]{2.8.8})에서 정의되었다. 일반적인 경우에는 이들이
$Y$의 아핀 열린집합에서 더 작은 열린집합으로 가는 제한 연산과 양립함을
확인하면 충분하며, 이는
(\hyperref[II.3.5.1.1-ko]{3.5.1.1})의 가환성에서 곧바로 따른다.
\end{proof}

특히 모든 $n\in\mathbf{Z}$에 대하여 다음 표준 준동형이 존재한다.
\[
\label{II.3.5.2.1-ko}
  \Phi^*(\mathcal{O}_{X'}(n))\longrightarrow
  \mathcal{O}_X(n)|G(\varphi).
\tag{3.5.2.1}
\]

\begin{proposition}[3.5.3]
\label{II.3.5.3-ko}
$Y$, $Y'$을 두 준스킴, $\psi:Y'\to Y$를 사상,
$\mathcal{S}$를 준연접 등급 $\mathcal{O}_Y$-대수라 하고,
$\mathcal{S}'=\psi^*(\mathcal{S})$라 놓자. 그러면 $Y'$-스킴
$X'=\operatorname{Proj}(\mathcal{S}')$은
$\operatorname{Proj}(\mathcal{S})\times_Y Y'$과 표준적으로 동일시된다.
또한 $\mathcal{M}$이 준연접 등급 $\mathcal{S}$-가군이면,
$\mathcal{O}_{X'}$-가군 $(\psi^*(\mathcal{M}))^{\sim}$은
$\widetilde{\mathcal{M}}\otimes_Y\mathcal{O}_{Y'}$과 표준적으로 동일시된다.
\end{proposition}

\begin{proof}
먼저 $\psi^*(\mathcal{S})$와 $\psi^*(\mathcal{M})$은 그 동차 성분들과
마찬가지로 준연접 $\mathcal{O}_{Y'}$-가군임을 지적해 두자
(\hyperref[0.5.1.4-ko]{0, 5.1.4}). $U$를 $Y$의 아핀 열린집합,
$U'\subset\psi^{-1}(U)$를 $Y'$의 아핀 열린집합, $A$, $A'$을 각각
$U$, $U'$의 환이라 하자. 그러면 $\mathcal{S}|U=\widetilde{S}$이고,
여기서 $S$는 등급 $A$-대수이며, $\mathcal{S}'|U'$은
$(S\otimes_A A')^{\sim}$와 동일시된다
(\hyperref[I.1.6.5-ko]{I, 1.6.5}). 따라서 첫째 주장은
(\hyperref[II.2.8.10-ko]{2.8.10})과
(\hyperref[I.3.2.6.2-ko]{I, 3.2.6.2})에서 따른다. 실제로 앞의
동일시로 정의되는 사영
$\operatorname{Proj}(\mathcal{S}'|U')\to
\operatorname{Proj}(\mathcal{S}|U)$은 $U$, $U'$에 대한 제한 연산과
호환됨을 곧바로 확인할 수 있으므로, 사상
$\operatorname{Proj}(\mathcal{S}')\to\operatorname{Proj}(\mathcal{S})$를
잘 정의한다. 이제
\[
  q:\operatorname{Proj}(\mathcal{S})\to Y,
  \qquad q':\operatorname{Proj}(\mathcal{S}')\to Y'
\]
을 구조 사상이라 하자. 그러면 $q'^{-1}(U')$은
$q^{-1}(U)\times_U U'$과 동일시되며, 두 층
$(\psi^*(\mathcal{M}))^{\sim}|q'^{-1}(U')$과
$(\widetilde{\mathcal{M}}\otimes_Y\mathcal{O}_{Y'})|q'^{-1}(U')$은
(\hyperref[II.2.8.10-ko]{2.8.10})과
(\hyperref[I.1.6.5-ko]{I, 1.6.5})에 따라 모두
$(M\otimes_A A')^{\sim}$와 표준적으로 동일시된다. 여기서
$M=\Gamma(U,\mathcal{M})$라 놓았다. 앞의 동일시들이 제한 연산과
호환됨도 다시 곧바로 확인되므로 둘째 주장이 따른다.
\end{proof}

\begin{corollary}[3.5.4]
\label{II.3.5.4-ko}
(\hyperref[II.3.5.3-ko]{3.5.3})의 표기 아래, 모든
$n\in\mathbf{Z}$에 대하여 $\mathcal{O}_{X'}(n)$은
$\mathcal{O}_X(n)\otimes_Y\mathcal{O}_{Y'}$과 표준적으로 동일시된다
($X=\operatorname{Proj}(\mathcal{S})$).
\end{corollary}

\begin{proof}
실제로 (\hyperref[II.3.5.3-ko]{3.5.3})의 표기 아래, 모든
$n\in\mathbf{Z}$에 대하여
$\psi^*(\mathcal{S}(n))=\mathcal{S}'(n)$임은 명백하다.
\end{proof}

\begin{env}[3.5.5]
\label{II.3.5.5-ko}
앞의 표기를 유지하고, 표준 사영 $X'\to X$를 $\Psi$로 나타내며
$\mathcal{M}'=\psi^*(\mathcal{M})$라 놓자. 또한 $\mathcal{S}$가
$\mathcal{S}_1$로 생성되고

\oldpage[II]{63}

$X$가 $Y$ 위에서 유한형이라고 가정하자. 그러면 $\mathcal{S}'$은
$\mathcal{S}'_1$로 생성되고(이는 $Y$, $Y'$이 아핀인 경우로 환원하여
알 수 있다), $X'$은 $Y'$ 위에서 유한형이다
(\hyperref[I.6.3.4-ko]{I, 6.3.4}). $\mathcal{F}$를
$\mathcal{O}_X$-가군이라 하고 $\mathcal{F}'=\Psi^*(\mathcal{F})$라
놓자. 그러면 (\hyperref[II.3.5.4-ko]{3.5.4})와
(\hyperref[0.4.3.3-ko]{0, 4.3.3})에 따라 모든 $n\in\mathbf{Z}$에
대하여 $\mathcal{F}'(n)=\Psi^*(\mathcal{F}(n))$이다. 또한 표준
$\psi$-준동형
$\theta_n:q_*(\mathcal{F}(n))\to q'_*(\mathcal{F}'(n))$을 다음과 같이
정의한다. 도식
\[
  \xymatrix{
    X\ar[d]_{q}
    & X'\ar[l]_{\Psi}\ar[d]^{q'}\\
    Y
    & Y'\ar[l]^{\psi}
  }
\]
이 가환하므로, 정의해야 하는 것은 준동형
$q_*(\mathcal{F}(n))\to
\psi_*(q'_*(\Psi^*(\mathcal{F}(n))))
=q_*(\Psi_*(\Psi^*(\mathcal{F}(n))))$이다. 따라서
$\theta_n=q_*(\rho_n)$으로 잡으면 충분하다. 여기서 $\rho_n$은 표준
준동형
$\mathcal{F}(n)\to\Psi_*(\Psi^*(\mathcal{F}(n)))$이다
(\hyperref[0.4.4.3-ko]{0, 4.4.3}). $Y$의 임의의 아핀 열린집합
$U$와 $U'\subset\psi^{-1}(U)$인 $Y'$의 임의의 아핀 열린집합 $U'$에
대하여, $\theta_n$이 단면 위에서 표준 준동형
(\hyperref[0.3.7.2-ko]{0, 3.7.2})
$\Gamma(q^{-1}(U),\mathcal{F}(n))\to
\Gamma(q'^{-1}(U'),\mathcal{F}'(n))$을 준다는 것은 곧바로 확인된다.

(\hyperref[II.2.8.13.2-ko]{2.8.13.2})의 가환성에 따라,
$\mathcal{F}$가 준연접이면 도식
\[
  \xymatrix{
    \mathcal{F}\ar[r]^{\rho}
    & \mathcal{F}'\\
    (\Gamma_*(\mathcal{F}))^{\sim}
      \ar[u]^{\beta_{\mathcal{F}}}\ar[r]_{\widetilde{\theta}}
    & (\Gamma_*(\mathcal{F}'))^{\sim}
      \ar[u]_{\beta_{\mathcal{F}'}}
  }
\]
은 가환한다(위쪽 수평 화살표는 표준 $\Psi$-사상
$\mathcal{F}\to\Psi^*(\mathcal{F})$이다).

마찬가지로 (\hyperref[II.2.8.13.1-ko]{2.8.13.1})의 가환성에 따라 도식
\[
  \xymatrix{
    \Gamma_*(\widetilde{\mathcal{M}})\ar[r]^{\theta}
    & \Gamma_*(\widetilde{\mathcal{M}'})\\
    \mathcal{M}\ar[r]_{\rho}\ar[u]^{\alpha_{\mathcal{M}}}
    & \mathcal{M}'\ar[u]_{\alpha_{\mathcal{M}'}}
  }
\]
도 가환한다(아래쪽 수평 화살표는 표준 $\psi$-사상
$\mathcal{M}\to\psi^*(\mathcal{M})$이다).
\end{env}

\begin{env}[3.5.6]
\label{II.3.5.6-ko}
이제 두 준스킴 $Y$, $Y'$, 사상
$g:Y'\to Y$, 준연접 등급 $\mathcal{O}_Y$-대수(각각 준연접 등급
$\mathcal{O}_{Y'}$-대수) $\mathcal{S}$(각각 $\mathcal{S}'$), 그리고
등급 대수의 $g$-사상 $u:\mathcal{S}\to\mathcal{S}'$, 즉 등급
$\mathcal{O}_Y$-대수의 준동형
$\mathcal{S}\to g_*(\mathcal{S}')$을 생각하자. 또한 $u$를 주는 것은 등급
$\mathcal{O}_{Y'}$-대수의 준동형
$u^\sharp:g^*(\mathcal{S})\to\mathcal{S}'$을 주는 것과 동치임이 알려져 있다.
그러면 $u^\sharp$으로부터 표준적으로 $Y'$-사상
$W=\operatorname{Proj}(u^\sharp):G(u^\sharp)\to
\operatorname{Proj}(g^*(\mathcal{S}))$을 얻는다. 여기서 $G(u^\sharp)$은
$X'=\operatorname{Proj}(\mathcal{S}')$의 열린집합이다
(\hyperref[II.3.5.1-ko]{3.5.1}). 한편
$X''=\operatorname{Proj}(g^*(\mathcal{S}))$는

\oldpage[II]{64}

$X=\operatorname{Proj}(\mathcal{S})$라 놓을 때 $X\times_Y Y'$과 표준적으로
동일시된다(\hyperref[II.3.5.3-ko]{3.5.3}). 따라서
$\operatorname{Proj}(u^\sharp)$에 이어 첫째 사영
$p:X\times_Y Y'\to X$를 취하면 사상
$v:G(u^\sharp)\to X$를 얻는다. 이를 $\operatorname{Proj}(u)$로 나타내며,
이때 다음 도식은 가환한다.
\[
  \xymatrix{
    G(u^\sharp)\ar[r]^{v}\ar[d]
    & X\ar[d]\\
    Y'\ar[r]_{g}
    & Y
  }
\]

더 나아가 임의의 준연접 등급 $\mathcal{S}$-가군 $\mathcal{M}$에 대하여
다음 표준 $v$-사상이 존재한다.
\[
\label{II.3.5.6.1-ko}
  \nu:\widetilde{\mathcal{M}}\longrightarrow
  (g^*(\mathcal{M})\otimes_{g^*(\mathcal{S})}\mathcal{S}')^{\sim}
  |G(u^\sharp).
\tag{3.5.6.1}
\]

실제로 $\nu^\sharp$은 다음 준동형들을 차례로 합성하여 얻는다.
\[
  v^*(\widetilde{\mathcal{M}})
  =W^*(p^*(\widetilde{\mathcal{M}}))
  \longrightarrow W^*((g^*(\mathcal{M}))^{\sim})
  \longrightarrow
  (g^*(\mathcal{M})\otimes_{g^*(\mathcal{S})}\mathcal{S}')^{\sim}
  |G(u^\sharp)
\]
여기서 첫째 화살표는 동형
(\hyperref[II.3.5.3-ko]{3.5.3})에서 유도되고, 둘째 화살표는 준동형
(\hyperref[II.3.5.2-ko]{3.5.2, (ii)})이다. $\mathcal{S}$가
$\mathcal{S}_1$에 의해 생성되면 (\hyperref[II.3.5.2-ko]{3.5.2})에 따라
$\nu^\sharp$은 동형이다.

(\hyperref[II.3.5.6.1-ko]{3.5.6.1})의 특수한 경우로서, 모든
$n\in\mathbf{Z}$에 대하여 다음 표준 $v$-사상을 얻는다.
\[
\label{II.3.5.6.2-ko}
  \nu:\mathcal{O}_X(n)\longrightarrow
  \mathcal{O}_{X'}(n)|G(u^\sharp).
\tag{3.5.6.2}
\]
\end{env}

\subsection{준스킴 $\operatorname{Proj}(\mathcal{S})$의 닫힌 부분준스킴.}
\label{subsection:II.3.6-ko}

\begin{env}[3.6.1]
\label{II.3.6.1-ko}
$Y$를 준스킴,
$\varphi:\mathcal{S}\to\mathcal{S}'$을 차수 $0$인 준연접 등급
$\mathcal{O}_Y$-대수의 준동형이라 하자. 어떤 $n$이 존재하여
$k\geq n$일 때 $\varphi_k:\mathcal{S}_k\to\mathcal{S}'_k$가 전사이면
(각각 단사이면, 전단사이면), $\varphi$를
(\textbf{TN})-\emph{전사}(각각 (\textbf{TN})-\emph{단사},
(\textbf{TN})-\emph{전단사})라고 한다. 그러면 이에 대응하는 사상
$\Phi:\operatorname{Proj}(\mathcal{S}')\to\operatorname{Proj}(\mathcal{S})$
의 연구는 $\varphi$가 전사인(각각 단사인, 전단사인) 경우로 환원된다.
이는 (\hyperref[II.3.1.8-ko]{3.1.8})을 사용하여
(\hyperref[II.2.9.1-ko]{2.9.1})에서와 같이 증명된다(후자는 $Y$가
아핀인 특수한 경우이다). $\varphi$가 (\textbf{TN})-전단사라고 하는
대신, 이때 이를 (\textbf{TN})-동형사상이라고도 한다.
\end{env}

\begin{proposition}[3.6.2]
\label{II.3.6.2-ko}
$Y$를 준스킴, $\mathcal{S}$를 준연접 등급 $\mathcal{O}_Y$-대수라 하고,
$X=\operatorname{Proj}(\mathcal{S})$라 하자.
\begin{enumerate}
  \item[\rm(i)] $\varphi:\mathcal{S}\to\mathcal{S}'$이 등급
    $\mathcal{O}_Y$-대수의 (\textbf{TN})-전사 준동형이면, 이에 대응하는
    사상 $\Phi=\operatorname{Proj}(\varphi)$
    (\hyperref[II.3.5.1-ko]{3.5.1})는
    $\operatorname{Proj}(\mathcal{S}')$ 전체에서 정의되며
    $\operatorname{Proj}(\mathcal{S}')$에서 $X$로 가는 닫힌 몰입 사상이다.
    $\mathcal{J}$가 $\varphi$의 핵이면, $\Phi$에 결부된 $X$의 닫힌
    부분준스킴은 $\mathcal{O}_X$의 준연접 아이디얼층
    $\widetilde{\mathcal{J}}$로 정의된다.
  \item[\rm(ii)] 더 나아가 $\mathcal{S}_0=\mathcal{O}_Y$이고,
    $\mathcal{S}$가 $\mathcal{S}_1$로 생성되며 $\mathcal{S}_1$이
    유한형이라고 가정하자. $X=\operatorname{Proj}(\mathcal{S})$의
    닫힌 부분준스킴 $X'$이 $\mathcal{O}_X$의 준연접 아이디얼층
    $\mathcal{I}$로 정의된다고 하자. $\mathcal{J}$를

\oldpage[II]{65}

    표준 준동형 $\alpha:\mathcal{S}\to\Gamma_*(\mathcal{O}_X)$
    (\hyperref[II.3.3.2-ko]{3.3.2})에 의한
    $\Gamma_*(\mathcal{I})$의 역상인 $\mathcal{S}$의 준연접 등급
    아이디얼층이라 하고, $\mathcal{S}'=\mathcal{S}/\mathcal{J}$라 놓자.
    그러면 $X'$은 등급 $\mathcal{O}_Y$-대수의 표준 준동형
    $\mathcal{S}\to\mathcal{S}'$에 대응하는 닫힌 몰입 사상
    $\operatorname{Proj}(\mathcal{S}')\to X$에 결부된 부분준스킴이다
    (\hyperref[I.4.2.1-ko]{I, 4.2.1}).
\end{enumerate}
\end{proposition}

\begin{proof}
\begin{enumerate}
  \item[\rm(i)] $\varphi$가 전사라고 가정해도 된다
    (\hyperref[II.3.6.1-ko]{3.6.1}). 그러면 $Y$의 모든 아핀 열린집합
    $U$에 대하여 $\Gamma(U,\mathcal{S})\to
    \Gamma(U,\mathcal{S}')$은 전사이다
    (\hyperref[I.1.3.9-ko]{I, 1.3.9}). 따라서
    (\hyperref[II.3.5.1-ko]{3.5.1})에 의해 $G(\varphi)=X$이다. 문제는
    곧바로 $Y$가 아핀인 경우의 명제를 증명하는 것으로 환원되며, 이
    경우에는 (\hyperref[II.2.9.2-ko]{2.9.2, (i)})에서 따른다.
  \item[\rm(ii)] 표준 단사 준동형 $\mathcal{J}\to\mathcal{S}$에서
    유도되는 준동형 $\widetilde{\mathcal{J}}\to\mathcal{O}_X$가
    $\widetilde{\mathcal{J}}$에서 $\mathcal{I}$로 가는 동형임을 보이면
    된다. 이 문제는 $Y$ 위에서 국소적이므로, $Y$가 환 $A$의 아핀
    준스킴이라고 가정할 수 있다. 그러면 $\mathcal{S}=\widetilde{S}$이고,
    여기서 $S$는 $S_1$로 생성되는 등급 $A$-대수이며 $S_1$은 $A$ 위에서
    유한형이다. 이제 (\hyperref[II.2.9.2-ko]{2.9.2, (ii)})를 적용하면
    충분하다.
\end{enumerate}
\end{proof}

\begin{corollary}[3.6.3]
\label{II.3.6.3-ko}
(\hyperref[II.3.6.2-ko]{3.6.2, (i)})의 조건 아래, 더 나아가
$\mathcal{S}$가 $\mathcal{S}_1$로 생성된다고 가정하면, 모든
$n\in\mathbf{Z}$에 대하여 $\Phi^*(\mathcal{O}_X(n))$은
$\mathcal{O}_{X'}(n)$과 표준적으로 동일시된다.
\end{corollary}

\begin{proof}
$Y$가 아핀일 때에는 이러한 표준 동형을 이미 정의하였다
(\hyperref[II.2.9.3-ko]{2.9.3}). 일반적인 경우에는 $Y$의 각 아핀
열린집합 $U$에 대하여 이렇게 정의한 동형들이 아핀 열린집합
$U'\subset U$로 옮겨갈 때 서로 호환됨을 확인하면 충분하며, 이는
명백하다.
\end{proof}

\begin{corollary}[3.6.4]
\label{II.3.6.4-ko}
$Y$를 준스킴, $\mathcal{S}$를 $\mathcal{S}_1$로 생성되는 준연접 등급
$\mathcal{O}_Y$-대수, $\mathcal{M}$을 준연접 $\mathcal{O}_Y$-가군,
$u:\mathcal{M}\to\mathcal{S}_1$을 전사 $\mathcal{O}_Y$-준동형이라
하자. 또한
$\overline{u}:\mathbf{S}_{\mathcal{O}_Y}(\mathcal{M})\to\mathcal{S}$를
$u$를 연장하는 등급 $\mathcal{O}_Y$-대수의 준동형이라 하자
(\hyperref[II.1.7.4-ko]{1.7.4}). 그러면 $\overline{u}$에 대응하는
사상은 $\operatorname{Proj}(\mathcal{S})$에서
$\operatorname{Proj}(\mathbf{S}_{\mathcal{O}_Y}(\mathcal{M}))$로 가는
닫힌 몰입 사상이다.
\end{corollary}

\begin{proof}
실제로 가정에 따라 $\overline{u}$는 전사이므로
(\hyperref[II.3.6.2-ko]{3.6.2, (i)})을 적용한다.
\end{proof}

\subsection{준스킴에서 동차 스펙트럼으로 가는 사상.}
\label{subsection:II.3.7-ko}

\begin{env}[3.7.1]
\label{II.3.7.1-ko}
$q:X\to Y$를 준스킴의 사상, $\mathcal{L}$을 가역인
$\mathcal{O}_X$-가군, $\mathcal{S}$를 음이 아닌 차수로 등급이 주어진
준연접 $\mathcal{O}_Y$-대수라 하자. 그러면 $q^*(\mathcal{S})$는 음이
아닌 차수로 등급이 주어진 준연접 $\mathcal{O}_X$-대수이다.
준연접 등급 $\mathcal{O}_X$-대수
$\mathcal{S}'=\bigoplus_{n\geq0}\mathcal{L}^{\otimes n}$을 생각하고,
등급 $\mathcal{O}_X$-대수의 준동형
\[
  \psi:q^*(\mathcal{S})\longrightarrow
  \mathcal{S}'=\bigoplus_{n\geq0}\mathcal{L}^{\otimes n}
\]
이 주어졌다고 가정하자. 이는 또한 등급 대수의 $q$-사상
\[
  \psi^\flat:\mathcal{S}\longrightarrow q_*(\mathcal{S}').
\]
을 주는 것과 같다. $\operatorname{Proj}(\mathcal{S}')$가 $X$와
표준적으로 동일시됨을 알고 있다
(\hyperref[II.3.1.7-ko]{3.1.7} 및
\hyperref[II.3.1.8-ko]{3.1.8, (iii)}). 따라서 $\psi$로부터 $X$의
열린부분집합 $G(\psi)$와 $Y$-사상
\[
\label{II.3.7.1.1-ko}
  r_{\mathcal{L},\psi}:G(\psi)\longrightarrow
  \operatorname{Proj}(\mathcal{S})=P
\tag{3.7.1.1}
\]

\oldpage[II]{66}

을 표준적으로 얻는다. 이 사상을 \emph{$\mathcal{L}$과 $\psi$에 대응하는
사상}이라 부르겠다. 이 사상은 다음 $Y$-사상에 이어 첫째 사영
$\pi:\operatorname{Proj}(q^*(\mathcal{S}))=P\times_YX\to P$를 취한
합성으로 얻어짐을 상기하자(\hyperref[II.3.5.6-ko]{3.5.6}).
\[
  \tau=\operatorname{Proj}(\psi):G(\psi)\longrightarrow
  \operatorname{Proj}(q^*(\mathcal{S})).
\]
\end{env}

\begin{env}[3.7.2]
\label{II.3.7.2-ko}
$Y=\operatorname{Spec}(A)$가 아핀인 경우의
$r=r_{\mathcal{L},\psi}$를 구체적으로 기술하자. 이때
$\mathcal{S}=\widetilde{S}$이고, 여기서 $S$는 음이 아닌 차수로 등급이
주어진 $A$-대수이다. 먼저 $X=\operatorname{Spec}(B)$도 아핀이며
$\mathcal{L}=\widetilde{L}$이라고 가정하자. 여기서 $L$은 계수가 $1$인
자유 $B$-가군이다. 그러면
$q^*(\mathcal{S})=(S\otimes_AB)^{\sim}$이다
(\hyperref[I.1.6.5-ko]{I, 1.6.5}). $c$가 $L$의 생성원이면,
$\psi_n:q^*(\mathcal{S}_n)\to\mathcal{L}^{\otimes n}$에는
$S_n\otimes_AB$에서 $L^{\otimes n}$으로 가는 준동형
$w_n:s\otimes b\mapsto bv_n(s)c^{\otimes n}$이 대응한다. 여기서
$v_n:S_n\to B$는 $A$-가군의 준동형이고, $v_n$들은 모여 대수 준동형
$S\to B$를 이룬다. $f\in S_d$ ($d>0$)라 하고 $g=v_d(f)$로 놓자.
(\hyperref[II.2.8.10-ko]{2.8.10})과 $D_+(f)$의
$\operatorname{Spec}(S_{(f)})$와의 동일시
(\hyperref[II.2.3.6-ko]{2.3.6})에 따라
$\pi^{-1}(D_+(f))=D_+(f\otimes1)$이다. 한편 식
(\hyperref[II.2.8.1.1-ko]{2.8.1.1})은 $X$와
$\operatorname{Proj}(\mathcal{S}')$의 표준 동일시를 고려하면
\[
  \tau^{-1}(D_+(f\otimes1))=D(g)
\]
임을 보인다. 따라서
\[
\label{II.3.7.2.1-ko}
  r^{-1}(D_+(f))=D(g).
\tag{3.7.2.1}
\]

더 나아가 $D(g)$에 제한한 사상 $\tau=\operatorname{Proj}(\psi)$는
$(s\otimes1)/(f\otimes1)^n$을 $v_{nd}(s)/g^n$으로 보내는 준동형에
대응한다($s\in S_{nd}$인 경우)
(\hyperref[II.2.8.1-ko]{2.8.1}). 또한 $D_+(f\otimes1)$에 제한한 사영
$\pi$는 준동형
$s/f^n\mapsto(s\otimes1)/(f\otimes1)^n$에 대응한다. 이로부터
$D(g)$에 제한한 $r$은 다음 조건을 만족하는 $A$-대수 준동형
$\omega:S_{(f)}\to B_g$에 대응한다.
$s\in S_{nd}$ ($n>0$)이면
$\omega(s/f^n)=v_{nd}(s)/g^n$이다. 이제 $X$가 임의인 경우로 넘어가면
($Y$는 여전히 아핀이다), (\hyperref[II.2.8.1-ko]{2.8.1})을 고려하여
다음을 얻는다:
\end{env}

\begin{proposition}[3.7.3]
\label{II.3.7.3-ko}
$Y=\operatorname{Spec}(A)$가 아핀이고 $\mathcal{S}=\widetilde{S}$라 하자.
여기서 $S$는 등급 $A$-대수이다. 모든
$f\in S_d=\Gamma(Y,\mathcal{S}_d)$에 대하여
\[
\label{II.3.7.3.1-ko}
  r_{\mathcal{L},\psi}^{-1}(D_+(f))=X_{\psi^\flat(f)}
  \qquad\bigl(\text{여기서 }\psi^\flat(f)\in
  \Gamma(X,\mathcal{L}^{\otimes d})\bigr)
\tag{3.7.3.1}
\]
이고, $r_{\mathcal{L},\psi}$의 제한
$X_{\psi^\flat(f)}\to D_+(f)=\operatorname{Spec}(S_{(f)})$은
(\hyperref[I.2.2.4-ko]{I, 2.2.4})에 따라 대수의 준동형
\[
\label{II.3.7.3.2-ko}
  \psi_{(f)}^\flat:S_{(f)}\longrightarrow
  \Gamma(X_{\psi^\flat(f)},\mathcal{O}_X)
\tag{3.7.3.2}
\]
에 대응한다. 이 준동형은
$s\in S_{nd}=\Gamma(Y,\mathcal{S}_{nd})$에 대하여
\[
\label{II.3.7.3.3-ko}
  \psi_{(f)}^\flat(s/f^n)
  =(\psi^\flat(s)|X_{\psi^\flat(f)})
   (\psi^\flat(f)|X_{\psi^\flat(f)})^{-n}.
\tag{3.7.3.3}
\]
을 만족한다.
\end{proposition}

$G(\psi)=X$이면 $r_{\mathcal{L},\psi}$가 \emph{전체에서 정의된다}고
하겠다. 이것이 성립할 필요충분조건은 명백히 모든 아핀 열린집합
$U\subset Y$에 대하여
$G(\psi)\cap q^{-1}(U)=q^{-1}(U)$인 것이다. 다시 말해 이 문제는
\emph{$Y$ 위에서 국소적}이다. $Y$가 아핀이면 $G(\psi)$는 $S_+$의
동차 원소 $f$들을 모두 달리할 때 $r^{-1}(D_+(f))$들의 합집합이다
(\hyperref[II.2.8.1-ko]{2.8.1}). 따라서
(\hyperref[II.3.7.3.1-ko]{3.7.3.1})에 의해
$X_{\psi^\flat(f)}$들은 $X$의 \emph{덮개}를 이루어야 한다. 즉 다음을
얻는다.

\begin{corollary}[3.7.4]
\label{II.3.7.4-ko}
(\hyperref[II.3.7.3-ko]{3.7.3})의 가정 아래에서
$r_{\mathcal{L},\psi}$가 전체에서 정의될 필요충분조건은 모든
$x\in X$에 대하여 정수 $n>0$와 $Y$ 위의 $\mathcal{S}_n$의 단면
$s$가 존재하여, 이때
$t=\psi^\flat(s)\in\Gamma(X,\mathcal{L}^{\otimes n})$로 놓으면
$t(x)\neq0$인 것이다.
\end{corollary}

\oldpage[II]{67}

$\psi$가 (\textbf{TN})-전사이면 이 조건이 항상 만족됨에 유의하자.

마찬가지로 $r_{\mathcal{L},\psi}$가 \emph{우세 사상}인지의 문제도
$Y$ 위에서 국소적이며, 다음이 성립한다.

\begin{corollary}[3.7.5]
\label{II.3.7.5-ko}
(\hyperref[II.3.7.3-ko]{3.7.3})의 가정 아래에서
$r_{\mathcal{L},\psi}$가 우세 사상일 필요충분조건은 모든 정수
$n>0$에 대하여, $\psi^\flat(s)\in
\Gamma(X,\mathcal{L}^{\otimes n})$가 국소적으로 멱영인 모든 단면
$s\in S_n$이 그 자체로 멱영인 것이다.
\end{corollary}

\begin{proof}
실제로 $r_{\mathcal{L},\psi}^{-1}(D_+(s))$가 공집합이면
$D_+(s)$도 공집합이어야 함을 나타내면 된다. 따라서 따름정리는
(\hyperref[II.3.7.3.1-ko]{3.7.3.1})과
(\hyperref[II.2.3.7-ko]{2.3.7})에서 따른다.
\end{proof}

\begin{proposition}[3.7.6]
\label{II.3.7.6-ko}
$q:X\to Y$를 사상, $\mathcal{L}$을 가역 $\mathcal{O}_X$-가군,
$\mathcal{S}$, $\mathcal{S}'$을 두 준연접 등급
$\mathcal{O}_Y$-대수,
$u:\mathcal{S}'\to\mathcal{S}$를 등급 대수의 준동형,
$\psi:q^*(\mathcal{S})\to\bigoplus_{n\geq0}\mathcal{L}^{\otimes n}$를
등급 대수의 준동형,
$\psi'=\psi\circ q^*(u)$를 그 합성 준동형이라 하자.
$r_{\mathcal{L},\psi'}$이 전체에서 정의되면
$r_{\mathcal{L},\psi}$도 그러하다. $u$가 (\textbf{TN})-전사이고
$r_{\mathcal{L},\psi'}$이 우세 사상이면
$r_{\mathcal{L},\psi}$도 우세 사상이다. 역으로 $u$가
(\textbf{TN})-단사이고 $r_{\mathcal{L},\psi}$가 우세 사상이면
$r_{\mathcal{L},\psi'}$도 우세 사상이다.
\end{proposition}

\begin{proof}
실제로 $G(\psi')\subset G(\psi)$이다
(\hyperref[II.2.8.4-ko]{2.8.4}). 이로부터 첫째 주장이 따른다. $u$가
(\textbf{TN})-전사이면
$\operatorname{Proj}(u):\operatorname{Proj}(\mathcal{S})\to
\operatorname{Proj}(\mathcal{S}')$은 전체에서 정의되며 닫힌 몰입
사상이다. 또한 $r_{\mathcal{L},\psi'}$은
$r_{\mathcal{L},\psi}$의 $G(\psi')$에 대한 제한에
$\operatorname{Proj}(u)$를 합성하여 얻어지므로,
$r_{\mathcal{L},\psi'}$이 우세 사상이면
$r_{\mathcal{L},\psi}$도 우세 사상임을 알 수 있다. 마지막으로 $u$가
(\textbf{TN})-단사이면 $\operatorname{Proj}(u)$는 $G(u)$에서
$\operatorname{Proj}(\mathcal{S}')$로 가는 우세 사상임을 알고 있다
(\hyperref[II.2.8.3-ko]{2.8.3}). $G(\psi')$은
$r_{\mathcal{L},\psi}$에 의한 $G(u)$의 역상이므로,
$r_{\mathcal{L},\psi}$가 우세 사상이면
$r_{\mathcal{L},\psi'}$도 우세 사상임을 알 수 있다.
\end{proof}

\begin{proposition}[3.7.7]
\phantomsection
\label{II.3.7.7-ko}
$Y$를 준콤팩트 준스킴, $q:X\to Y$를 준콤팩트 사상,
$\mathcal{L}$을 가역 $\mathcal{O}_X$-가군, $\mathcal{S}$를 준연접 등급
$\mathcal{O}_Y$-대수라 하되, 준연접 $\mathcal{O}_Y$-대수들의 귀납계
$(\mathcal{S}^\lambda)$의 여과 귀납극한이라 하자.
$\varphi_\lambda:\mathcal{S}^\lambda\to\mathcal{S}$를 표준 준동형,
$\psi:q^*(\mathcal{S})\to\bigoplus_{n\geq0}\mathcal{L}^{\otimes n}$를
등급 대수의 준동형이라 하고,
$\psi_\lambda=\psi\circ q^*(\varphi_\lambda)$로 놓자.
$r_{\mathcal{L},\psi}$가 모든 곳에서 정의되기 위한 필요충분조건은
$r_{\mathcal{L},\psi_\lambda}$가 모든 곳에서 정의되는 $\lambda$가
존재하는 것이다~; 이때 $\mu\geq\lambda$이면
$r_{\mathcal{L},\psi_\mu}$도 모든 곳에서 정의된다.
\end{proposition}

\begin{proof}
충분성은 (\hyperref[II.3.7.6-ko]{3.7.6})에서 따른다. 역으로,
$r_{\mathcal{L},\psi}$가 모든 곳에서 정의된다고 하자~; $Y$가 아핀인
경우로 환원할 수 있다. 실제로 모든 아핀 열린집합 $U\subset Y$에 대하여
$r_{\mathcal{L},\psi_{\lambda(U)}}$의 $q^{-1}(U)$에 대한 제한이 모든
곳에서 정의되는 $\lambda(U)$가 존재한다면, $Y$가 준콤팩트이므로 $Y$를
유한 개의 아핀 열린집합 $U_i$로 덮고 모든 지표 $i$에 대하여
$\lambda\geq\lambda(U_i)$인 $\lambda$를 택하면
(\hyperref[II.3.7.6-ko]{3.7.6})에 따라 충분하다. $Y$가 아핀이면 가정에
따라 모든 $x\in X$에 대하여 어떤 $S_n$의 단면 $s^{(x)}$가 존재하여,
$t^{(x)}=\psi^\flat(s^{(x)})$로 놓을 때 $t^{(x)}(x)\neq0$이다
($t^{(x)}$는 $X$ 위의 $\mathcal{L}^{\otimes n}$의 단면으로 본다).
따라서 $x$의 어떤 근방 $V(x)$에 속하는 모든 $z$에 대하여
$t^{(x)}(z)\neq0$이다. $X$를 유한 개의 $V(x_i)$로 덮고 $s^{(i)}$를
이에 대응하는 $S$의 단면이라 하자~; 그러면 어떤 $\lambda$가 존재하여
모든 $i$에 대하여 $s^{(i)}=\varphi_\lambda(s_\lambda^{(i)})$이고
$s_\lambda^{(i)}\in S^\lambda$이다~; 따라서
(\hyperref[II.3.7.4-ko]{3.7.4})에 따라
$r_{\mathcal{L},\psi_\lambda}$는 모든 곳에서 정의된다. 마지막 주장은
(\hyperref[II.3.7.6-ko]{3.7.6})의 자명한 귀결이다.
\end{proof}

\begin{corollary}[3.7.8]
\phantomsection
\label{II.3.7.8-ko}
(\hyperref[II.3.7.7-ko]{3.7.7})의 가정 아래에서 각
$r_{\mathcal{L},\psi_\lambda}$가 우세 사상이면
$r_{\mathcal{L},\psi}$도 우세 사상이다~; $\varphi_\lambda$들이 단사이면
그 역도 성립한다.
\end{corollary}

\oldpage[II]{68}

\begin{proof}
둘째 주장은 (\hyperref[II.3.7.6-ko]{3.7.6})의 특수한 경우이다~; 한편
$r_{\mathcal{L},\psi}$가 우세 사상임을 보이기 위해서는 $Y$가 아핀인
경우만 생각하면 충분하다~; $s\in S$이고 $\psi^\flat(s)$가 국소적으로
멱영이라 하자. 적어도 하나의 $\lambda$에 대하여
$s=\varphi_\lambda(s_\lambda)$로 쓸 수 있으므로, 가정과
(\hyperref[II.3.7.5-ko]{3.7.5})에서 $s_\lambda$가 멱영원임이 따른다.
따라서 $s$도 멱영원이고, (\hyperref[II.3.7.5-ko]{3.7.5})의 판정법을
적용할 수 있다.
\end{proof}

\begin{namedtheorem}[3.7.9]{비고}
\label{II.3.7.9-ko}
\begin{enumerate}
  \item[(i)] (\hyperref[II.3.7.1-ko]{3.7.1})의 표기법을 쓰고
    (\hyperref[II.3.2.10-ko]{3.2.10})을 고려하면, 모든
    $n\in\mathbf{Z}$에 대하여 다음 표준 준동형이 있다.
    \[
    \label{II.3.7.9.1-ko}
      \theta:r_{\mathcal{L},\psi}^*(\mathcal{O}_P(n))\longrightarrow
      \mathcal{L}^{\otimes n}
    \tag{3.7.9.1}
    \]
    이는 (\hyperref[II.3.5.6.2-ko]{3.5.6.2})에서 일반적으로 정의되어
    있다. (\hyperref[II.3.7.3-ko]{3.7.3})의 조건 아래에서 이 준동형의
    $X_{\psi^\flat(f)}$에 대한 제한은 $s/f^k$ ($s\in S_{n+kd}$)에 원소
    $(\psi^\flat(s)|X_{\psi^\flat(f)})
    (\psi^\flat(f)|X_{\psi^\flat(f)})^{-k}$를 대응시키는 것으로
    구체적으로 주어짐을 곧바로 알 수 있다.
  \item[(ii)] $\mathcal{F}$를 준연접 $\mathcal{O}_X$-가군이라 하고,
    $q$가 분리 사상이고 준콤팩트라고 가정하자. 그러면 모든
    $n\geq0$에 대하여 $q_*(\mathcal{F}\otimes\mathcal{L}^{\otimes n})$은
    준연접 $\mathcal{O}_Y$-가군이다
    (\hyperref[I.9.2.2-ko]{I, 9.2.2}).
    $\mathcal{M}'=\bigoplus_{n\geq0}\mathcal{F}\otimes
    \mathcal{L}^{\otimes n}$이라 놓자. 이는 준연접 등급
    $\mathcal{S}'$-가군이다. 그리고 그 직상
    $\mathcal{M}=q_*(\mathcal{M}')=
    \bigoplus_{n\geq0}q_*(\mathcal{F}\otimes\mathcal{L}^{\otimes n})$을
    생각하자(준동형 $\psi^\flat$을 통해 이는 준연접
    $\mathcal{S}$-가군이다). 다음과 같은 표준 $\mathcal{O}_X$-가군
    준동형이 있음을 보이겠다.
    \[
    \label{II.3.7.9.2-ko}
      \xi:r_{\mathcal{L},\psi}^*(\widetilde{\mathcal{M}})
      \longrightarrow\mathcal{F}|G(\psi).
    \tag{3.7.9.2}
    \]

    실제로 (\hyperref[II.3.5.6.1-ko]{3.5.6.1})에서 다음 표준 준동형을
    정의하였다.
    \[
    \label{II.3.7.9.3-ko}
      r_{\mathcal{L},\psi}^*(\widetilde{\mathcal{M}})
      \longrightarrow
      (q^*(\mathcal{M})\otimes_{q^*(\mathcal{S})}\mathcal{S}')^{\sim}
      |G(\psi),
    \tag{3.7.9.3}
    \]
    여기서 우변은 $\operatorname{Proj}(\mathcal{S}')$ 위의 준연접층으로
    본다. 한편 모든 준연접 등급 $\mathcal{S}'$-가군 $\mathcal{M}'$에
    대하여 다음 표준 준동형이 있다~:
    \[
    \label{II.3.7.9.4-ko}
      q^*(q_*(\mathcal{M}'))\otimes_{q^*(\mathcal{S})}\mathcal{S}'
      \longrightarrow\mathcal{M}'
    \tag{3.7.9.4}
    \]
    즉 $X$의 모든 열린집합 $U$, $U$ 위의
    $q^*(q_*(\mathcal{M}'_h))$의 모든 단면 $t'$, 그리고 $U$ 위의
    $\mathcal{S}'_k$의 모든 단면 $b'$에 대하여, $t'\otimes b'$에는
    $\mathcal{M}'_{h+k}$의 단면 $b'\sigma(t')$를 대응시킨다. 여기서
    $\sigma(t')$는 (\hyperref[0.4.4.3-ko]{0, 4.4.3})에 따라 $t'$에
    표준적으로 대응하는 $U$ 위의 $\mathcal{M}'_h$의 단면이다. 이로부터
    다음 표준 준동형을 얻는다.
    \[
    \label{II.3.7.9.5-ko}
      (q^*(q_*(\mathcal{M}'))\otimes_{q^*(\mathcal{S})}\mathcal{S}')^{\sim}
      |G(\psi)\longrightarrow\widetilde{\mathcal{M}'}|G(\psi)
    \tag{3.7.9.5}
    \]
    마지막으로 $\widetilde{\mathcal{M}'}$은 $\mathcal{F}$와 표준적으로
    동일시되므로 (\hyperref[II.3.2.9-ko]{3.2.9, (i)}), 찾던 표준
    준동형을 얻는다.

    (\hyperref[II.3.7.3-ko]{3.7.3})의 조건 아래에서 이 준동형의
    $X_{\psi^\flat(f)}$에 대한 제한은 다음과 같이 구체적으로 주어진다~:
    $X$ 위의 $\mathcal{F}\otimes\mathcal{L}^{\otimes nd}$의 단면
    $t_{nd}$가 주어졌다고 하자. $t'_{nd}$로 $t_{nd}$를 $Y$ 위의
    $q_*(\mathcal{F}\otimes\mathcal{L}^{\otimes nd})$의 단면으로 본 것을
    나타내면, 원소 $t'_{nd}/f^n$에는 $X_{\psi^\flat(f)}$ 위의
    $\mathcal{F}$의 단면
    $(t_{nd}|X_{\psi^\flat(f)})
    (\psi^\flat(f)|X_{\psi^\flat(f)})^{-n}$이 대응한다.
\end{enumerate}
\end{namedtheorem}

\oldpage[II]{69}

\subsection{동차 스펙트럼으로의 몰입 사상 판정법.}
\label{subsection:II.3.8-ko}

\begin{env}[3.8.1]
\label{II.3.8.1-ko}
(\hyperref[II.3.7.1-ko]{3.7.1})의 표기법을 쓰면,
$r_{\mathcal{L},\psi}$가 몰입 사상(각각 열린 몰입 사상, 닫힌 몰입 사상)인지의
문제는 명백히 \emph{$Y$ 위에서 국소적이다}.
\end{env}

\begin{proposition}[3.8.2]
\label{II.3.8.2-ko}
(\hyperref[II.3.7.3-ko]{3.7.3})의 가정 아래에서,
$r_{\mathcal{L},\psi}$가 모든 점에서 정의되고 몰입 사상이기 위한
필요충분조건은 다음과 같은 단면들의 족 $s_\alpha\in S_{n_\alpha}$
($n_\alpha>0$)가 존재하는 것이다. 곧,
$f_\alpha=\psi^\flat(s_\alpha)$라 놓으면 다음 조건들이 성립한다.
\begin{enumerate}
  \item[(i)] $X_{f_\alpha}$들이 $X$의 덮개를 이룬다.
  \item[(ii)] $X_{f_\alpha}$들은 아핀 열린집합이다.
  \item[(iii)] 모든 $\alpha$와 모든
    $t\in\Gamma(X_{f_\alpha},\mathcal{O}_X)$에 대하여, 다음을 만족하는
    정수 $m>0$과 $s\in S_{mn_\alpha}$가 존재한다.
    $t=(\psi^\flat(s)|X_{f_\alpha})(f_\alpha|X_{f_\alpha})^{-m}$.
\end{enumerate}

$r_{\mathcal{L},\psi}$가 모든 점에서 정의되고 열린 몰입 사상이기 위한
필요충분조건은 조건 (i), (ii), (iii)과 다음 조건을 만족하는 족
$(s_\alpha)$가 존재하는 것이다.
\begin{enumerate}
  \item[(iv)] 모든 $n>0$과
    $\psi^\flat(s)|X_{f_\alpha}=0$을 만족하는 모든
    $s\in S_{nn_\alpha}$에 대하여, $s_\alpha^ks=0$을 만족하는 정수
    $k>0$이 존재한다.
\end{enumerate}

$r_{\mathcal{L},\psi}$가 모든 점에서 정의되고 닫힌 몰입 사상이기 위한
필요충분조건은 (i), (ii), (iii)과 다음 조건을 만족하는 족
$(s_\alpha)$가 존재하는 것이다.
\begin{enumerate}
  \item[(v)] $D_+(s_\alpha)$들이
    $P=\operatorname{Proj}(S)$의 덮개를 이룬다.
\end{enumerate}
\end{proposition}

\begin{proof}
실제로 $r$가 몰입 사상(각각 닫힌 몰입 사상)이기 위한 필요충분조건은
$D_+(s_\alpha)$들로 이루어진 $r(G(\psi))$의 덮개(각각 $P$의 덮개)가
존재하여, $V_\alpha=r^{-1}(D_+(s_\alpha))$라 놓으면 $V_\alpha$에 대한
$r$의 제한이 $V_\alpha$에서 $D_+(s_\alpha)$로 가는 닫힌 몰입 사상이
되는 것이다(\hyperref[I.4.2.4-ko]{I, 4.2.4}). 조건 (i)는
$r$가 모든 점에서 정의된다는 것과 $D_+(s_\alpha)$들이 $r(X)$를
덮는다는 것을 동시에 나타낸다
(\hyperref[II.3.7.3.1-ko]{3.7.3.1})~; $D_+(s_\alpha)$가 아핀이므로,
조건 (ii)와 (iii)은 $X_{f_\alpha}$에 대한 $r$의 제한이
$D_+(s_\alpha)$로 가는 닫힌 몰입 사상임을 나타낸다
(\hyperref[I.4.2.3-ko]{I, 4.2.3})~; 끝으로 (iii)과 (iv)는
$\psi_{(s_\alpha)}^\flat$가 전단사임을 나타내므로
(\hyperref[II.3.7.3.2-ko]{3.7.3.2}의 표기법), (ii), (iii), (iv)는 모든
$\alpha$에 대하여 $X_{f_\alpha}$에 대한 $r$의 제한이
$D_+(s_\alpha)$ 위로 가는 동형사상임을 나타낸다. 따라서 (i), (ii),
(iii), (iv)는 $r$가 열린 몰입 사상임을 나타낸다.
\end{proof}

\begin{corollary}[3.8.3]
\label{II.3.8.3-ko}
(\hyperref[II.3.7.6-ko]{3.7.6})의 가정 아래에서,
$r_{\mathcal{L},\psi'}$가 모든 점에서 정의되고 몰입 사상이면
$r_{\mathcal{L},\psi}$도 그러하다. 더 나아가 $u$가
(\textbf{TN})-전사이고 $r_{\mathcal{L},\psi'}$가 열린 몰입 사상(각각
닫힌 몰입 사상)이면 $r_{\mathcal{L},\psi}$도 그러하다.
\end{corollary}

\begin{proof}
실제로 (\hyperref[II.3.8.2-ko]{3.8.2})에 따라,
$s'_\alpha\in S'_{n_\alpha}$인 족이 존재하여
$f_\alpha=\psi'^\flat(s'_\alpha)$라 놓으면 조건 (i), (ii), (iii)이
성립한다. 그런데 $s_\alpha=u(s'_\alpha)$라 놓으면 역시
$f_\alpha=\psi^\flat(s_\alpha)$이고,
$t=(\psi'^\flat(s')|X_{f_\alpha})(f_\alpha|X_{f_\alpha})^{-m}$이면
$s=u(s')$라 놓음으로써
$t=(\psi^\flat(s)|X_{f_\alpha})(f_\alpha|X_{f_\alpha})^{-m}$도 성립한다.
따라서 첫째 주장을 얻는다. 둘째 주장은
$\operatorname{Proj}(u)$가 닫힌 몰입 사상이라는 사실에서 곧바로 따른다.
\end{proof}

\begin{proposition}[3.8.4]
\label{II.3.8.4-ko}
(\hyperref[II.3.7.7-ko]{3.7.7})의 가정들이 성립하고, 더 나아가
$q:X\to Y$가 유한형 사상이라고 하자. 그러면
$r_{\mathcal{L},\psi}$가 전체에서 정의되는 몰입 사상일 필요충분조건은
어떤 $\lambda$가 존재하여 $r_{\mathcal{L},\psi_\lambda}$가 전체에서
정의되는 몰입 사상인 것이다. 이 경우 $\mu\geq\lambda$이면
$r_{\mathcal{L},\psi_\mu}$도 전체에서 정의되는 몰입 사상이다.
\end{proposition}

\oldpage[II]{70}

\begin{proof}
(\hyperref[II.3.8.3-ko]{3.8.3})을 고려하면,
$r_{\mathcal{L},\psi}$가 전체에서 정의되는 몰입 사상일 때 적어도 하나의
$\lambda$에 대하여 $r_{\mathcal{L},\psi_\lambda}$도 그러함을 보이면
충분하다. $Y$의 준콤팩트성을 사용하는
(\hyperref[II.3.7.7-ko]{3.7.7})에서와 같은 논법으로, 먼저 $Y$가
아핀인 경우로 환원된다. $X$가 준콤팩트이므로
(\hyperref[II.3.8.2-ko]{3.8.2})에 따라 성질 (i), (ii), (iii)을 갖는
$S$의 원소들의 \emph{유한}족 $(s_i)$ ($s_i\in S_{n_i}$)이 존재한다.
사상 $X_{f_i}\to Y$ (여기서 $f_i=\psi^\flat(s_i)$)는 유한형이다.
실제로 이는 아핀 스킴의 사상이므로 준콤팩트하고
(\hyperref[I.6.6.1-ko]{I, 6.6.1}), $q$가 유한형이므로 국소 유한형이며
(\hyperref[I.6.3.2-ko]{I, 6.3.2}), 결론은
(\hyperref[I.6.6.3-ko]{I, 6.6.3})에서 따른다. 따라서 $X_{f_i}$의 환
$B_i$는 유한형 $A$-대수이다
(\hyperref[I.6.3.3-ko]{I, 6.3.3}). 이 대수의 생성원족을 $(t_{ij})$라
하자. 가정에 따라 다음을 만족하는 원소
$s'_{ij}\in S_{m_{ij}n_i}$가 존재한다.
\[
  t_{ij}=(\psi^\flat(s'_{ij})|X_{f_i})
  (\psi^\flat(s_i)|X_{f_i})^{-m_{ij}}.
\]
가정에 따라 어떤 $\lambda$와 원소
$s_{i\lambda}\in S_{n_i}^\lambda$,
$s'_{ij\lambda}\in S_{m_{ij}n_i}^\lambda$가 존재하여, 이들의
$\varphi_\lambda$에 의한 상은 각각 $s_i$, $s'_{ij}$이다. 그러면
$r_{\mathcal{L},\psi_\lambda}$가
(\hyperref[II.3.8.2-ko]{3.8.2})의 조건 (i), (ii), (iii)을 만족함은
명백하다.
\end{proof}

\begin{proposition}[3.8.5]
\label{II.3.8.5-ko}
$Y$를 준콤팩트 스킴 또는 그 기저공간이 뇌터인 준스킴,
$q:X\to Y$를 유한형 사상, $\mathcal{L}$을 가역
$\mathcal{O}_X$-가군, $\mathcal{S}$를 준연접 등급
$\mathcal{O}_Y$-대수,
$\psi:q^*(\mathcal{S})\to\bigoplus_{n\geq0}\mathcal{L}^{\otimes n}$을
등급 대수의 준동형이라 하자. $r_{\mathcal{L},\psi}$가 전체에서
정의되는 몰입 사상일 필요충분조건은 다음 조건들을 만족하는 정수
$n>0$과 $\mathcal{S}_n$의 유한형 준연접 부분-$\mathcal{O}_Y$-가군
$\mathcal{E}$가 존재하는 것이다.
\begin{enumerate}
  \item[a)] 준동형
    $\psi_n\circ q^*(j_n):q^*(\mathcal{E})\to\mathcal{L}^{\otimes n}$
    (여기서 $j_n:\mathcal{E}\to\mathcal{S}$은 표준 단사 준동형이다)는
    전사이다.
  \item[b)] $\mathcal{S}'$을 $\mathcal{E}$로 생성되는
    $\mathcal{S}$의 (등급) 부분-$\mathcal{O}_Y$-대수,
    $\psi'$을 준동형 $\psi\circ q^*(j')$ (여기서 $j'$은
    $\mathcal{S}'$에서 $\mathcal{S}$로 가는 단사 준동형이다)이라 하면,
    $r_{\mathcal{L},\psi'}$은 전체에서 정의되는 몰입 사상이다.
\end{enumerate}

이 조건들이 성립할 때, $\mathcal{E}$를 포함하는 $\mathcal{S}_n$의
모든 준연접 부분-$\mathcal{O}_Y$-가군 $\mathcal{E}'$도 같은 성질을
가지며, 모든 $k>0$에 대하여 $\mathcal{S}_{kn}$ 안에서
$\bigotimes^k\mathcal{E}$의 상도 마찬가지이다.
\end{proposition}

\begin{proof}
조건의 충분성과 마지막 두 주장은
$\operatorname{Proj}(\mathcal{S})$와
$\operatorname{Proj}(\mathcal{S}^{(k)})$ 사이의 표준 동형
(\hyperref[II.3.1.8-ko]{3.1.8})을 고려하면
(\hyperref[II.3.8.3-ko]{3.8.3})의 특수한 경우이다.

$(U_i)$를 $Y$의 아핀 열린집합들로 이루어진 유한 열린 덮개라 하고
$A_i=A(U_i)$라 놓자. $q^{-1}(U_i)$가 콤팩트이므로,
$r_{\mathcal{L},\psi}$가 전체에서 정의되는 몰입 사상이라는 가정과
(\hyperref[II.3.8.2-ko]{3.8.2})에 따라 성질 (i), (ii), (iii)을 갖는
$S^{(i)}=\Gamma(U_i,\mathcal{S})$의 원소들의 유한족 $(s_{ij})$
($s_{ij}\in S_{n_{ij}}^{(i)}$)이 존재한다.
(\hyperref[II.3.8.4-ko]{3.8.4})의 증명에서와 같이, 사상
$X_{f_{ij}}\to U_i$ (여기서 $f_{ij}=\psi^\flat(s_{ij})$)가 유한형이고,
따라서 $X_{f_{ij}}$의 환 $B_{ij}$가 유한형 $A_i$-대수임을 알 수 있다.
이 대수에는 다음 꼴의 생성원계가 존재한다.
$(\psi^\flat(t_{ijk})|X_{f_{ij}})(f_{ij}|X_{f_{ij}})^{-m_{ijk}}$,
여기서 $t_{ijk}\in S_{m_{ijk}n_{ij}}^{(i)}$이다. $n$을 모든
$m_{ijk}n_{ij}$의 공배수라 하자. 각 $(i,j,k)$에 대하여 $s_{ij}$를
$\rho m_{ijk}n_{ij}=n$을 만족하는 거듭제곱 $s_{ij}^\rho$로 바꾸고,
$t_{ijk}$에 $s_{ij}^{\rho-m_{ijk}}$를 곱하면, 각 $i$에 대하여 모든
$s_{ij}$와 $t_{ijk}$가 $S_n^{(i)}$에 속하고 $m_{ijk}=1$이라고 가정할
수 있다. $E_i$를 이 원소들로 생성되는 $S_n^{(i)}$의 부분-$A_i$-가군이라
하자($i$는 고정한다). 다음을 만족하는 $\mathcal{S}_n$의 유한형 연접
부분-$\mathcal{O}_Y$-가군 $\mathcal{E}_i$가 존재한다.

\oldpage[II]{71}

$\mathcal{E}_i|U_i=(E_i)^{\sim}$
(\hyperref[I.9.4.7-ko]{I, 9.4.7}). $\mathcal{E}_i$들의 합인
$\mathcal{S}_n$의 부분-$\mathcal{O}_Y$-가군 $\mathcal{E}$가 요구된
조건을 만족함은 명백하다(각 단면 $f_{ij}$는 다음 성질을 갖는다. 모든
$x\in X_{f_{ij}}$에 대하여 $x$의 아핀 근방
$V\subset X_{f_{ij}}$가 존재하여 $f_{ij}|V$가
$\Gamma(V,\mathcal{L}^{\otimes n})$의 기저가 된다).
\end{proof}

\section{사영 다발. 풍부한 층}
\label{section:II.4-ko}

\subsection{사영 다발의 정의.}
\label{subsection:II.4.1-ko}

\begin{definition}[4.1.1]
\label{II.4.1.1-ko}
$Y$를 준스킴, $\mathcal{E}$를 준연접 $\mathcal{O}_Y$-가군,
$\mathbf{S}_{\mathcal{O}_Y}(\mathcal{E})$를 $\mathcal{E}$의 대칭
$\mathcal{O}_Y$-대수라 하자
(\hyperref[II.1.7.4-ko]{1.7.4}). 이 대수는 준연접이다
(\hyperref[II.1.7.7-ko]{1.7.7}). $Y$-스킴
$P=\operatorname{Proj}(\mathbf{S}_{\mathcal{O}_Y}(\mathcal{E}))$를
\emph{$\mathcal{E}$로 정의되는 $Y$ 위의 사영 다발}이라 하고
$\mathbf{P}(\mathcal{E})$로 나타낸다. $\mathcal{O}_P$-가군
$\mathcal{O}_P(1)$을 \emph{$P$ 위의 기본 층}이라 한다.
\end{definition}

$Y$가 환 $A$의 아핀 스킴이면 $\mathcal{E}=\widetilde{E}$이며, 여기서
$E$는 $A$-가군이다. 이때 $\mathbf{P}(\widetilde{E})$ 대신
$\mathbf{P}(E)$로 쓴다.

$\mathcal{E}=\mathcal{O}_Y^n$으로 둘 때에는
$\mathbf{P}(\mathcal{E})$ 대신 $\mathbf{P}_Y^{n-1}$로 쓴다~; 더 나아가
$Y$가 환 $A$의 아핀 스킴이면 $\mathbf{P}_Y^{n-1}$ 대신
$\mathbf{P}_A^{n-1}$로도 쓴다.
$\mathbf{S}_{\mathcal{O}_Y}(\mathcal{O}_Y)$는 $\mathcal{O}_Y[T]$와
표준적으로 동일시되므로 (\hyperref[II.1.7.4-ko]{1.7.4}),
$\mathbf{P}_Y^0$은 $Y$와 표준적으로 동일시된다
(\hyperref[II.3.1.7-ko]{3.1.7})~; 예
(\hyperref[II.2.4.3-ko]{2.4.3})은 다름 아닌 $\mathbf{P}_K^1$이다.

\begin{env}[4.1.2]
\label{II.4.1.2-ko}
$\mathcal{E}$와 $\mathcal{F}$를 두 준연접 $\mathcal{O}_Y$-가군,
$u:\mathcal{E}\to\mathcal{F}$를 $\mathcal{O}_Y$-준동형이라 하자. 이에
등급 $\mathcal{O}_Y$-대수의 준동형
$\mathbf{S}(u):\mathbf{S}_{\mathcal{O}_Y}(\mathcal{E})\to
\mathbf{S}_{\mathcal{O}_Y}(\mathcal{F})$가 표준적으로 대응한다
(\hyperref[II.1.7.4-ko]{1.7.4}). $u$가 \emph{전사}이면
$\mathbf{S}(u)$도 전사이고, 따라서
(\hyperref[II.3.6.2-ko]{3.6.2, (i)}),
$\operatorname{Proj}(\mathbf{S}(u))$는 닫힌 몰입 사상
$\mathbf{P}(\mathcal{F})\to\mathbf{P}(\mathcal{E})$이다. 이를
$\mathbf{P}(u)$로 나타내겠다. 그러므로 준연접 $\mathcal{O}_Y$-가군들의
사상으로 전사 준동형만을 취한다는 조건 아래, $\mathbf{P}(\mathcal{E})$는
\emph{$\mathcal{E}$에 관한 반변 함자}라고 할 수 있다.

계속해서 $u$가 전사라고 가정하고
$P=\mathbf{P}(\mathcal{E})$, $Q=\mathbf{P}(\mathcal{F})$,
$j=\mathbf{P}(u)$로 놓으면, 동형까지 다음이 성립한다.
\[
\label{II.4.1.2.1-ko}
  j^*(\mathcal{O}_P(n))=\mathcal{O}_Q(n)
  \qquad\text{모든 }n\in\mathbf{Z}\text{에 대하여}
\tag{4.1.2.1}
\]
이는 (\hyperref[II.3.6.3-ko]{3.6.3})에서 따른다.
\end{env}

\begin{env}[4.1.3]
\label{II.4.1.3-ko}
이제 $\psi:Y'\to Y$를 사상이라 하고
$\mathcal{E}'=\psi^*(\mathcal{E})$로 놓자. 그러면
$\mathbf{S}_{\mathcal{O}_{Y'}}(\mathcal{E}')=
\psi^*(\mathbf{S}_{\mathcal{O}_Y}(\mathcal{E}))$
이다 (\hyperref[II.1.7.5-ko]{1.7.5})~; 따라서
(\hyperref[II.3.5.3-ko]{3.5.3}), 표준 동형까지 다음이 성립한다.
\[
\label{II.4.1.3.1-ko}
  \mathbf{P}(\psi^*(\mathcal{E}))=
  \mathbf{P}(\mathcal{E})\times_YY'
\tag{4.1.3.1}
\]
또한 모든 $n\in\mathbf{Z}$에 대하여 명백히
\[
  \psi^*((\mathbf{S}_{\mathcal{O}_Y}(\mathcal{E}))(n))
  =(\mathbf{S}_{\mathcal{O}_{Y'}}(\mathcal{E}'))(n)
\]
이다. 따라서 $P=\mathbf{P}(\mathcal{E})$,
$P'=\mathbf{P}(\mathcal{E}')$로 놓으면
(\hyperref[II.3.5.4-ko]{3.5.4}), 동형까지 다음이 성립한다.
\[
\label{II.4.1.3.2-ko}
  \mathcal{O}_{P'}(n)=\mathcal{O}_P(n)\otimes_Y\mathcal{O}_{Y'}
  \qquad\text{모든 }n\in\mathbf{Z}\text{에 대하여}.
\tag{4.1.3.2}
\]
\end{env}

\oldpage[II]{72}

\begin{proposition}[4.1.4]
\label{II.4.1.4-ko}
$\mathcal{L}$을 가역 $\mathcal{O}_Y$-가군이라 하자. 모든 준연접
$\mathcal{O}_Y$-가군 $\mathcal{E}$에 대하여 표준 $Y$-동형사상
$i_{\mathcal{L}}:\mathbf{P}(\mathcal{E})\xrightarrow{\sim}
\mathbf{P}(\mathcal{E}\otimes\mathcal{L})$이 존재한다~; 더 나아가
$P=\mathbf{P}(\mathcal{E})$, $Q=\mathbf{P}(\mathcal{E}\otimes\mathcal{L})$로
놓으면, 모든 $n\in\mathbf{Z}$에 대하여
$i_{\mathcal{L}}^*(\mathcal{O}_Q(n))$은
$\mathcal{O}_P(n)\otimes_Y\mathcal{L}^{\otimes n}$과 표준적으로 동형이다.
\end{proposition}

\begin{proof}
먼저 $A$가 환, $E$가 $A$-가군, $L$이 \emph{생성원 하나를 갖는 자유}
$A$-가군이면, 다음 $A$-가군의 준동형이 표준적으로 정의됨을 지적하자.
\[
  \mathbf{S}_n(E\otimes L)\longrightarrow
  \mathbf{S}_n(E)\otimes L^{\otimes n}
\]
즉 $(x_1\otimes y_1)\cdots(x_n\otimes y_n)$에 원소
\[
  (x_1x_2\cdots x_n)\otimes(y_1\otimes y_2\otimes\cdots\otimes y_n)
  \qquad(x_i\in E,\ y_i\in L\text{이고 }1\leq i\leq n)~;
\]
를 대응시킨다. 이 준동형은 ($L=A$인 경우로 환원하여) 실제로
동형사상임을 곧바로 확인할 수 있다. 따라서 등급 $A$-대수의 표준 동형
$\mathbf{S}_A(E\otimes L)\xrightarrow{\sim}
\bigoplus_{n\geq0}\mathbf{S}_n(E)\otimes L^{\otimes n}$을 얻는다. 이제
(4.1.4)의 조건들로 돌아가면, 앞의 고찰에 따라 등급
$\mathcal{O}_Y$-대수의 표준 동형
\[
\label{II.4.1.4.1-ko}
  \mathbf{S}_{\mathcal{O}_Y}(\mathcal{E}\otimes_{\mathcal{O}_Y}\mathcal{L})
  \xrightarrow{\sim}
  \bigoplus_{n\geq0}\mathbf{S}_n(\mathcal{E})
  \otimes_{\mathcal{O}_Y}\mathcal{L}^{\otimes n}
\tag{4.1.4.1}
\]
을 정의할 수 있다. 이는 이 동형을 준층들의 동형으로 정의하고
(\hyperref[II.1.7.4-ko]{1.7.4}),
(\hyperref[I.1.3.9-ko]{I, 1.3.9}) 및
(\hyperref[I.1.3.12-ko]{I, 1.3.12})을 고려하여 얻는다. 이제 명제는
(\hyperref[II.3.1.8-ko]{3.1.8, (iii)})과
(\hyperref[II.3.2.10-ko]{3.2.10})에서 따른다.
\end{proof}

\begin{env}[4.1.5]
\label{II.4.1.5-ko}
(\hyperref[II.4.1.1-ko]{4.1.1})의 가정 아래
$P=\mathbf{P}(\mathcal{E})$로 놓고, $p$로 구조 사상 $P\to Y$를 나타내자.
정의에 따라 $\mathcal{E}=(\mathbf{S}_{\mathcal{O}_Y}(\mathcal{E}))_1$이므로
표준 준동형 $\alpha_1:\mathcal{E}\to p_*(\mathcal{O}_P(1))$이 있고
(\hyperref[II.3.3.2.2-ko]{3.3.2.2}), 따라서 또한
(\hyperref[0.4.4.3-ko]{0, 4.4.3})에 따라 표준 준동형
\[
\label{II.4.1.5.1-ko}
  \alpha_1^\sharp:p^*(\mathcal{E})\longrightarrow\mathcal{O}_P(1).
\tag{4.1.5.1}
\]
이 있다.
\end{env}

\begin{proposition}[4.1.6]
\label{II.4.1.6-ko}
표준 준동형 (\hyperref[II.4.1.5.1-ko]{4.1.5.1})은 전사이다.
\end{proposition}

\begin{proof}
실제로 (\hyperref[II.3.3.2-ko]{3.3.2})에서
$\alpha_1^\sharp$이 다음 표준 준동형에 함자적으로 대응함을 보았다.
\[
  \mathcal{E}\otimes_{\mathcal{O}_Y}
  \mathbf{S}_{\mathcal{O}_Y}(\mathcal{E})
  \longrightarrow
  (\mathbf{S}_{\mathcal{O}_Y}(\mathcal{E}))(1)~;
\]
정의상 $\mathcal{E}$가 $\mathbf{S}_{\mathcal{O}_Y}(\mathcal{E})$를 생성하므로
이 준동형은 전사이며, 따라서 (\hyperref[II.3.2.4-ko]{3.2.4})에 의해
결론을 얻는다.
\end{proof}

\subsection{준스킴에서 사영 다발로 가는 사상.}
\label{subsection:II.4.2-ko}

\begin{env}[4.2.1]
\label{II.4.2.1-ko}
(\hyperref[II.4.1.5-ko]{4.1.5})의 표기법을 유지하자. $X$를
$Y$-준스킴, $q:X\to Y$를 그 구조 사상이라 하고, $r:X\to P$를
$Y$-\emph{사상}이라 하자. 그러면 다음 도식은 가환한다.
\[
  \xymatrix{
    P \ar[d]_p & X \ar[l]_r \ar[dl]^q
    \\Y
  }
\]

\oldpage[II]{73}

함자 $r^*$는 오른쪽 완전하므로
(\hyperref[0.4.3.1-ko]{0, 4.3.1}),
(\hyperref[II.4.1.5.1-ko]{4.1.5.1})의 준동형이 전사라는 사실
(\hyperref[II.4.1.6-ko]{4.1.6})에서 다음 전사 준동형을 얻는다.
\[
  r^*(\alpha_1^\sharp):r^*(p^*(\mathcal{E}))\longrightarrow
  r^*(\mathcal{O}_P(1)).
\]
그런데 $r^*(p^*(\mathcal{E}))=q^*(\mathcal{E})$이고,
$r^*(\mathcal{O}_P(1))$은
$r^*(\mathcal{O}_P)=\mathcal{O}_X$와 국소적으로 동형이다. 다시 말해 이는
가역 $\mathcal{O}_X$-가군 $\mathcal{L}_r$이다. 따라서 $r$로부터 다음
표준 전사 $\mathcal{O}_X$-준동형을 정의하였다.
\[
\label{II.4.2.1.1-ko}
  \varphi_r:q^*(\mathcal{E})\longrightarrow\mathcal{L}_r.
\tag{4.2.1.1}
\]

$Y=\operatorname{Spec}(A)$가 아핀이고
$\mathcal{E}=\widetilde{E}$일 때에는 이 준동형을 다음과 같이 더
구체적으로 기술할 수 있다. $f\in E$가 주어지면, 먼저
(\hyperref[II.2.6.3-ko]{2.6.3})에서 다음을 얻는다.
\[
\label{II.4.2.1.2-ko}
  r^{-1}(D_+(f))=X_{\varphi_r^\flat(f)}.
\tag{4.2.1.2}
\]
다른 한편, $V$를 $r^{-1}(D_+(f))$에 포함되는 $X$의 아핀 열린집합이라
하고, $B$를 그 환이라 하자. 그러면 $B$는 $A$-대수이다. 이제
$S=\mathbf{S}_A(E)$로 놓자. $r$의 $V$에 대한 제한은
$A$-준동형 $\omega:S_{(f)}\to B$에 대응하며,
$q^*(\mathcal{E})|V=\widetilde{E\otimes_A B}$이고
$\mathcal{L}_r|V=\widetilde{L_r}$이다. 여기서
$L_r=(S(1))_{(f)}\otimes_{S_{(f)}}B_{[\omega]}$이다
(\hyperref[I.1.6.5-ko]{I, 1.6.5}). $\varphi_r$의
$q^*(\mathcal{E})|V$에 대한 제한은 $x\otimes1$을
$(x/1)\otimes f=(f/1)\otimes\omega(x/f)$로 보내는 $B$-준동형
$u:E\otimes_A B\to L_r$에 대응한다. 그러므로 $\varphi_r$를
표준적으로 연장한 다음 $\mathcal{O}_X$-대수 준동형
\[
  \psi_r:q^*(\mathbf{S}(\mathcal{E}))=
  \mathbf{S}(q^*(\mathcal{E}))\longrightarrow
  \mathbf{S}(\mathcal{L}_r)=\bigoplus_{n\geq0}\mathcal{L}_r^{\otimes n}
\]
은 $q^*(\mathbf{S}_n(\mathcal{E}))|V$에 제한했을 때 준동형
$\mathbf{S}_n(E\otimes_A B)=\mathbf{S}_n(E)\otimes_A B\to L_r^{\otimes n}$
에 대응하며, 이 준동형은 $s\otimes1$을
$(f/1)^{\otimes n}\otimes\omega(s/f^n)$으로 보낸다.
\end{env}

\begin{env}[4.2.2]
\label{II.4.2.2-ko}
거꾸로, 사상 $q:X\to Y$, 가역 $\mathcal{O}_X$-가군 $\mathcal{L}$,
준연접 $\mathcal{O}_Y$-가군 $\mathcal{E}$가 주어졌다고 하자. 각 준동형
$\varphi:q^*(\mathcal{E})\to\mathcal{L}$에는 다음 준연접
$\mathcal{O}_X$-대수들의 준동형이 표준적으로 대응한다.
\[
  \psi:\mathbf{S}(q^*(\mathcal{E}))=q^*(\mathbf{S}(\mathcal{E}))
  \longrightarrow\bigoplus_{n\geq0}\mathcal{L}^{\otimes n},
\]
따라서 (\hyperref[II.3.7.1-ko]{3.7.1})에 의해 $Y$-사상
$r_{\mathcal{L},\psi}:G(\psi)\to
\operatorname{Proj}(\mathbf{S}(\mathcal{E}))=\mathbf{P}(\mathcal{E})$가
대응하며, 이를 $r_{\mathcal{L},\varphi}$라 쓰겠다. $\varphi$가 전사이면
$\psi$도 전사이므로 (\hyperref[II.3.7.4-ko]{3.7.4})에 의해
$r_{\mathcal{L},\varphi}$는 전체에서 정의된다. 더 나아가
(\hyperref[II.4.2.1-ko]{4.2.1})과 (4.2.2)의 표기법을 쓰면 다음이
성립한다:
\end{env}

\begin{proposition}[4.2.3]
\label{II.4.2.3-ko}
사상 $q:X\to Y$와 준연접 $\mathcal{O}_Y$-가군
$\mathcal{E}$가 주어졌다고 하자. 그러면 대응
$r\mapsto(\mathcal{L}_r,\varphi_r)$와
$(\mathcal{L},\varphi)\mapsto r_{\mathcal{L},\varphi}$는
$Y$-사상 $r:X\to\mathbf{P}(\mathcal{E})$들의 집합과, 가역
$\mathcal{O}_X$-가군 $\mathcal{L}$ 및 전사 준동형
$\varphi:q^*(\mathcal{E})\to\mathcal{L}$로 이루어진 쌍
$(\mathcal{L},\varphi)$들의 동치류의 집합 사이에 일대일 대응을 이룬다.
여기서 두 쌍 $(\mathcal{L},\varphi)$,
$(\mathcal{L}',\varphi')$가 동치라는 것은
$\varphi'=\tau\circ\varphi$를 만족하는 $\mathcal{O}_X$-동형사상
$\tau:\mathcal{L}\xrightarrow{\sim}\mathcal{L}'$가 존재한다는 뜻이다.
\end{proposition}

\begin{proof}
먼저 $Y$-사상 $r:X\to\mathbf{P}(\mathcal{E})$에서 출발하여
$\mathcal{L}_r$와 $\varphi_r$를 만들고
(\hyperref[II.4.2.1-ko]{4.2.1}), 이어서
$r'=r_{\mathcal{L}_r,\varphi_r}$를 만들자~;
(\hyperref[II.4.2.1-ko]{4.2.1})과
(\hyperref[II.3.7.2-ko]{3.7.2})에서 사상 $r$와 $r'$가 동일함이 곧바로
따른다(\hyperref[II.3.7.2-ko]{3.7.2}의 준동형 $v_n$들을 정의할 때
$\mathcal{L}_r$의 생성원으로 원소 $(f/1)\otimes1$을 택한다).
역으로 쌍 $(\mathcal{L},\varphi)$에서 출발하여

\oldpage[II]{74}

$r''=r_{\mathcal{L},\varphi}$를 만들고, 이어서 $\mathcal{L}_{r''}$와
$\varphi_{r''}$를 만들자~; 이제
$\varphi=\tau\circ\varphi_{r''}$를 만족하는 표준 동형사상
$\tau:\mathcal{L}_{r''}\xrightarrow{\sim}\mathcal{L}$가 있음을 보이자.
이를 정의하기 위해 $Y=\operatorname{Spec}(A)$와
$X=\operatorname{Spec}(B)$가 아핀인 경우로 국소화할 수 있으며,
(\hyperref[II.4.2.1-ko]{4.2.1})과
(\hyperref[II.3.7.2-ko]{3.7.2})의 표기법으로
$L_{r''}$의 각 원소 $(x/1)\otimes1$ ($x\in E$)에 $L$의 원소
$v_1(x)c$를 대응시키면 된다. $\tau$가 $L$의 생성원으로 택한 $c$와
무관함은 곧바로 확인된다~; 가정에 따라 $v_1$은 전사이므로,
$\tau$가 동형사상임을 보이려면 $x/1=0$이 $(S(1))_{(f)}$에서 성립할 때
$v_1(x)/1=0$이 $B_g$에서 성립함을 보이면 충분하다~; 그런데 첫째
관계는 어떤 $n$에 대하여 $f^nx=0$이 $\mathbf{S}_{n+1}(E)$에서
성립한다는 뜻이고, 이로부터
$v_{n+1}(f^nx)=g^nv_1(x)=0$이 $B$에서 성립하므로 결론을 얻는다.
끝으로 동치인 두 쌍 $(\mathcal{L},\varphi)$,
$(\mathcal{L}',\varphi')$에 대하여
$r_{\mathcal{L},\varphi}=r_{\mathcal{L}',\varphi'}$임은 명백하다.
\end{proof}

특히 $X=Y$이면~:

\begin{theorem}[4.2.4]
\label{II.4.2.4-ko}
$Y$ 위의 $\mathbf{P}(\mathcal{E})$의 단면들의 집합은, 준연접
$\mathcal{O}_Y$-가군 $\mathcal{F}$들로서 $\mathcal{E}$의 부분가군이고
$\mathcal{E}/\mathcal{F}$가 가역인 것들의 집합과 표준적으로 전단사
대응한다.
\end{theorem}

이 성질은 벡터 공간의 초평면들의 집합으로서의 고전적인
‘사영공간’의 정의에 대응함을 유의하자(고전적인 경우는
$Y=\operatorname{Spec}(K)$인 경우에 해당하는데, 여기서 $K$는 체이고
$\mathcal{E}=\widetilde{E}$이며 $E$는 유한 차원 $K$-벡터 공간이다~;
$\mathcal{F}$들 가운데 (\hyperref[II.4.2.4-ko]{4.2.4})에서 서술한
성질을 갖는 것들은 이때 $E$의 초평면들에 대응하고, 한편 $Y$ 위의
$\mathbf{P}(\mathcal{E})$의 단면들은 $\mathbf{P}(\mathcal{E})$의
$K$-유리점들이다
(\hyperref[I.3.4.5-ko]{I, 3.4.5})).

\begin{remark}[4.2.5]
\label{II.4.2.5-ko}
$Y$ 위에서 $X$로부터 $P$로 가는 사상들과 그 그래프 사상, 곧
$P\times_YX$의 $X$-단면들 사이에는 표준 전단사 대응이 있으므로
(\hyperref[I.3.3.14-ko]{I, 3.3.14}), 역으로
(\hyperref[II.4.2.3-ko]{4.2.3})은
(\hyperref[II.4.2.4-ko]{4.2.4})에서 도출할 수 있음을 알 수 있다.
$\operatorname{Hyp}_Y(X,\mathcal{E})$로 준연접
$\mathcal{O}_X$-가군 $\mathcal{F}$들 중에서
$\mathcal{E}\otimes_Y\mathcal{O}_X=q^*(\mathcal{E})$의 부분가군이고
$q^*(\mathcal{E})/\mathcal{F}$가 가역 $\mathcal{O}_X$-가군인 것들의
집합을 나타내자. $g:X'\to X$가 $Y$-사상이면 $g^*$이 오른쪽
완전하다는 사실에서
\[
  g^*(q^*(\mathcal{E})/\mathcal{F})=
  g^*(q^*(\mathcal{E}))/g^*(\mathcal{F}),
\]
를 얻으므로, 후자의 층은 가역이다. 따라서
$\operatorname{Hyp}_Y(X,\mathcal{E})$는 $Y$-준스킴의 범주 위의 반변
함자이다. 그러면 정리 (\hyperref[II.4.2.4-ko]{4.2.4})는
$\operatorname{Hom}_Y(X,\mathbf{P}(\mathcal{E}))$와
$\operatorname{Hyp}_Y(X,\mathcal{E})$라는 두 함자 사이의 표준 동형을
정의하는 것으로 해석할 수 있는데, 이 함자들은 $Y$-준스킴의 범주에서
변하는 $X$에 관하여 반변이다. 이는 또한 사영 다발
$P=\mathbf{P}(\mathcal{E})$를 다음 보편 성질로 특징짓는다. 이 성질은
\S\S~2와 3의 구성들보다 기하학적 직관에 더 가깝다~: 모든 사상
$q:X\to Y$와, 모든 가역 $\mathcal{O}_X$-가군 $\mathcal{L}$로서
$\mathcal{E}\otimes_{\mathcal{O}_Y}\mathcal{O}_X$의 몫인 것에 대하여,
다음을 만족하는
유일한 $Y$-사상 $r:X\to P$가 존재한다.
$\mathcal{L}=r^*(\mathcal{O}_P(1))$.

나중에 같은 방식으로, 특히 ‘그라스만’ 스킴들을 어떻게 정의할
수 있는지 살펴볼 것이다.
\end{remark}

\begin{corollary}[4.2.6]
\label{II.4.2.6-ko}
모든 가역 $\mathcal{O}_Y$-가군이 자명하다고 가정하자
(\hyperref[I.2.4.8-ko]{I, 2.4.8}).
$V$를 군 $\operatorname{Hom}_{\mathcal{O}_Y}(\mathcal{E},\mathcal{O}_Y)$라
하고, 이를 환 $A=\Gamma(Y,\mathcal{O}_Y)$ 위의 가군으로 보자. 또한
$V^*$를 전사 준동형들로 이루어진 $V$의 부분집합이라 하자. 그러면
$\mathbf{P}(\mathcal{E})$의 $Y$-단면들의 집합은 $V^*/A^*$와
표준적으로 동일시된다. 여기서 $A^*$는 $A$의 단원군이다.
\end{corollary}

\oldpage[II]{75}

특히 다음이 성립한다~:
\begin{enumerate}
  \item[1\textsuperscript{o}] $Y$가 \emph{국소 스킴}이면 따름정리
  (\hyperref[II.4.2.6-ko]{4.2.6})을 적용할 수 있다
  (\hyperref[I.2.4.8-ko]{I, 2.4.8}). $Y$를 임의의 준스킴, $y$를 $Y$의
  한 점, $Y'=\operatorname{Spec}(k(y))$라 하자~;
  (\hyperref[II.4.1.3.1-ko]{4.1.3.1})에 따라
  $\mathbf{P}(\mathcal{E})$의 섬유 $p^{-1}(y)$는
  $\mathbf{P}(\mathcal{E}^y)$와 동일시된다. 여기서
  $\mathcal{E}^y=\mathcal{E}_y\otimes_{\mathcal{O}_y}k(y)
  =\mathcal{E}_y/\mathfrak{m}_y\mathcal{E}_y$는 $k(y)$-벡터 공간으로
  본다. 더 일반적으로 $K$가 $k(y)$의 확대이면,
  $p^{-1}(y)\otimes_{k(y)}K$는
  $\mathbf{P}(\mathcal{E}^y\otimes_{k(y)}K)$와 동일시된다. 따라서
  따름정리 (\hyperref[II.4.2.6-ko]{4.2.6})에 의해, $k(y)$의 확대
  $K$를 값체로 하는 $\mathbf{P}(\mathcal{E})$의 \emph{기하적 점들}
  (\hyperref[I.3.4.5-ko]{I, 3.4.5})의 집합, 달리 말해 점 $y$ 위의
  $\mathbf{P}(\mathcal{E})$의 \emph{$K$-유리 기하적 섬유}는
  $K$-벡터 공간 $\mathcal{E}^y\otimes_{k(y)}K$의 \emph{쌍대}에
  대응하는 \emph{사영공간}과 동일시된다.
  \item[2\textsuperscript{o}] $Y$가 환 $A$의 아핀 스킴이고 모든 가역
  $\mathcal{O}_Y$-가군이 자명하다고 가정하자~; 또한
  $\mathcal{E}=\mathcal{O}_Y^n$으로 놓자. 그러면
  (\hyperref[II.4.2.6-ko]{4.2.6})에서 $V$는 $A^n$과 동일시되고
  (\hyperref[I.1.3.8-ko]{I, 1.3.8}), $V^*$는 $A$의 원소들의 계
  $(f_i)_{1\leq i\leq n}$ 가운데 그 원소들이 생성하는 아이디얼이
  $A$인 것들의 집합과 동일시된다~; 이와 같은 두 계가
  $\mathbf{P}_Y^{n-1}=\mathbf{P}_A^{n-1}$의 같은 $Y$-단면, 다시 말해
  $A$를 값으로 하는 $\mathbf{P}_A^{n-1}$의 같은 점을 정의할
  필요충분조건은 한 계가 다른 계에 $A$의 한 단원을 곱하여 얻어지는
  것이다.
\end{enumerate}

이 성질들은 $\mathbf{P}(\mathcal{E})$에 대하여 ‘사영 다발’이라는
용어를 쓰는 것을 정당화한다. 이렇게 얻은 ‘사영공간’의 정의는
실제로 고전적 정의의 쌍대임에 유의하자~; 이는 반드시 국소 자유일
필요는 없는 임의의 준연접 $\mathcal{O}_Y$-가군 $\mathcal{E}$에 대하여
$\mathbf{P}(\mathcal{E})$를 정의해야 한다는 필요성 때문에 강제된다.

\begin{remark}[4.2.7]
\label{II.4.2.7-ko}
제V장에서 다음을 보일 것이다. $Y$가 국소 뇌터이고 연결이며
$\mathcal{E}$가 국소 자유라고 하자. $P=\mathbf{P}(\mathcal{E})$로 놓으면,
모든 가역 $\mathcal{O}_P$-가군 $\mathcal{L}$은
$\mathcal{L}'\otimes_{\mathcal{O}_Y}\mathcal{O}_P(m)$ 꼴의
$\mathcal{O}_P$-가군과 동형이다. 여기서 $\mathcal{L}'$은 동형을
제외하고 유일하게 결정되는 가역 $\mathcal{O}_Y$-가군이고, $m$은
유일하게 결정되는 정수이다. 다시 말해
$\mathrm{H}^1(P,\mathcal{O}_P^*)$는
$\mathbf{Z}\times\mathrm{H}^1(Y,\mathcal{O}_Y^*)$와 동형이다
(\hyperref[0.5.4.7-ko]{0, 5.4.7}). 또한
(\agref{III.2.1.14-ko}{III, 2.1.14};
(\hyperref[0.5.4.10-ko]{0, 5.4.10})을 고려하여),
$m<0$이면 $p_*(\mathcal{L}^{\otimes m})=0$이고 $m\geq0$이면
$p_*(\mathcal{L}^{\otimes m})$이
$\mathcal{L}'\otimes_{\mathcal{O}_Y}
(\mathbf{S}_{\mathcal{O}_Y}(\mathcal{E}))_m$과 동형임을 보일 것이다.
$\mathcal{F}$가 준연접 $\mathcal{O}_Y$-가군이면, 모든 $Y$-사상
$\mathbf{P}(\mathcal{E})\to\mathbf{P}(\mathcal{F})$은 가역
$\mathcal{O}_Y$-가군 $\mathcal{L}'$, 정수 $m\geq0$, 그리고 대응하는
$\mathcal{O}_{\mathbf{P}(\mathcal{F})}$-가군 준동형 $\psi^\sharp$이
전사인 $\mathcal{O}_Y$-준동형
$\psi:\mathcal{F}\to\mathcal{L}'\otimes_{\mathcal{O}_Y}
(\mathbf{S}_{\mathcal{O}_Y}(\mathcal{E}))_m$의 데이터로 결정된다. 또한
고려하는 $Y$-사상이 동형사상이면 $m=1$이고 $\mathcal{F}$는
$\mathcal{E}\otimes_{\mathcal{O}_Y}\mathcal{L}'$과 동형임을 보일 것이다
((\hyperref[II.4.1.4-ko]{4.1.4})의 역). 이를 통해
$\mathbf{P}(\mathcal{E})$의 자동사상들의 싹으로 이루어진 층을 군의 층
$\mathcal{A}ut(\mathcal{E})$ ($\mathcal{E}$의 계수가 $n$이면 국소적으로
$\operatorname{GL}(n,\mathcal{O}_Y)$와 동형이다)를
$\mathcal{O}_Y^*$로 나눈 몫으로 결정할 수 있을 것이다.
\end{remark}

\begin{env}[4.2.8]
\label{II.4.2.8-ko}
(\hyperref[II.4.2.1-ko]{4.2.1})의 표기법을 유지하자. $u:X'\to X$를
사상이라 하자~; $Y$-사상 $r:X\to P$가 준동형
$\varphi:q^*(\mathcal{E})\to\mathcal{L}$에 대응하면, $Y$-사상
$r\circ u$는
$u^*(\varphi):u^*(q^*(\mathcal{E}))\to u^*(\mathcal{L})$에 대응한다.
이는 정의에서 곧바로 따른다.
\end{env}

\begin{env}[4.2.9]
\label{II.4.2.9-ko}
$\mathcal{E}$, $\mathcal{F}$를 두 준연접 $\mathcal{O}_Y$-가군,
$v:\mathcal{E}\to\mathcal{F}$를 전사 준동형, $j=\mathbf{P}(v)$를 이에
대응하는 닫힌 몰입
$\mathbf{P}(\mathcal{F})\to\mathbf{P}(\mathcal{E})$
(\hyperref[II.4.1.2-ko]{4.1.2})이라 하자. $Y$-사상
$r:X\to\mathbf{P}(\mathcal{F})$가 준동형
$\varphi:q^*(\mathcal{F})\to\mathcal{L}$에 대응하면,

\oldpage[II]{76}

$Y$-사상 $j\circ r$는 다음 합성에 대응한다.
\[
  q^*(\mathcal{E})\xrightarrow{q^*(v)}q^*(\mathcal{F})
  \xrightarrow{\varphi}\mathcal{L}~;
\]
이 역시 (\hyperref[II.4.2.1-ko]{4.2.1})에서 준 정의로부터 따른다.
\end{env}

\begin{env}[4.2.10]
\label{II.4.2.10-ko}
$\psi:Y'\to Y$를 사상이라 하고,
$\mathcal{E}'=\psi^*(\mathcal{E})$로 놓자. $Y$-사상 $r:X\to P$가
준동형 $\varphi:q^*(\mathcal{E})\to\mathcal{L}$에 대응하면,
$Y'$-사상
\[
  r_{(Y')}:X_{(Y')}\longrightarrow P'=\mathbf{P}(\mathcal{E}')
\]
은 다음 준동형에 대응한다.
$\varphi_{(Y')}:q_{(Y')}^*(\mathcal{E}')=
  q^*(\mathcal{E})\otimes_Y\mathcal{O}_{Y'}\to
\mathcal{L}\otimes_Y\mathcal{O}_{Y'}$. 실제로
(\hyperref[II.4.1.3.1-ko]{4.1.3.1})에 의해 다음 가환 도식을 얻는다.
\[
  \xymatrix{
    Y' \ar[d]
    & P'=P_{(Y')} \ar[l]_{p_{(Y')}} \ar[d]^{u}
    & X_{(Y')} \ar[l]_{r_{(Y')}} \ar[d]^{v}
    \\Y
    & P \ar[l]_{p}
    & X \ar[l]_{r}
  }
\]
(\hyperref[II.4.1.3.2-ko]{4.1.3.2})에 따르면
\[
  (r_{(Y')})^*(\mathcal{O}_{P'}(1))
  =(r_{(Y')})^*(u^*(\mathcal{O}_P(1)))
  =v^*(r^*(\mathcal{O}_P(1)))
  =v^*(\mathcal{L})=\mathcal{L}\otimes_Y\mathcal{O}_{Y'}~;
\]
다른 한편, $u^*(\alpha_1^\sharp)$은 바로 표준 준동형
$\alpha_1^\sharp:(p_{(Y')})^*(\mathcal{E}')\to\mathcal{O}_{P'}(1)$이다~;
이는 (\hyperref[II.4.1.6-ko]{4.1.6})에서처럼 $P$와 $P'$의 표준
준동형 $\alpha_1^\sharp$을 구체적으로 적어 보면 알 수 있다. 따라서
주장이 따른다.
\end{env}

\subsection{세그레 사상.}
\label{subsection:II.4.3-ko}

\begin{env}[4.3.1]
\label{II.4.3.1-ko}
$Y$를 준스킴, $\mathcal{E}$와 $\mathcal{F}$를 두 준연접
$\mathcal{O}_Y$-가군이라 하자~; $P_1=\mathbf{P}(\mathcal{E})$,
$P_2=\mathbf{P}(\mathcal{F})$로 놓고,
$p_1:P_1\to Y$, $p_2:P_2\to Y$로 구조 사상들을 나타내자.
$Q=P_1\times_YP_2$라 하고, $q_1:Q\to P_1$, $q_2:Q\to P_2$를
표준 사영이라 하자~; $\mathcal{O}_Q$-가군
\[
  \mathcal{L}=\mathcal{O}_{P_1}(1)\otimes_Y\mathcal{O}_{P_2}(1)
  :=q_1^*(\mathcal{O}_{P_1}(1))\otimes_{\mathcal{O}_Q}
  q_2^*(\mathcal{O}_{P_2}(1))
\]
은 두 가역 $\mathcal{O}_Q$-가군의 텐서곱이므로 가역이다
(\hyperref[0.5.4.4-ko]{0, 5.4.4}). 다른 한편,
$r=p_1\circ q_1=p_2\circ q_2$가 구조 사상 $Q\to Y$이면
\[
  r^*(\mathcal{E}\otimes_{\mathcal{O}_Y}\mathcal{F})
  =q_1^*(p_1^*(\mathcal{E}))\otimes_{\mathcal{O}_Q}
  q_2^*(p_2^*(\mathcal{F}))
\]
이다 (\hyperref[0.4.3.3-ko]{0, 4.3.3})~; 따라서 표준 전사 준동형
(\hyperref[II.4.1.5.1-ko]{4.1.5.1})
$p_1^*(\mathcal{E})\to\mathcal{O}_{P_1}(1)$,
$p_2^*(\mathcal{F})\to\mathcal{O}_{P_2}(1)$은 텐서곱에 의해 표준 준동형
\[
\label{II.4.3.1.1-ko}
  s:r^*(\mathcal{E}\otimes_{\mathcal{O}_Y}\mathcal{F})\longrightarrow
  \mathcal{L}
\tag{4.3.1.1}
\]
을 주며, 이는 명백히 전사이다~; 그러므로
(\hyperref[II.4.2.2-ko]{4.2.2})에서 \emph{세그레 사상}이라 하는 표준 사상
\[
\label{II.4.3.1.2-ko}
  \varsigma:\mathbf{P}(\mathcal{E})\times_Y\mathbf{P}(\mathcal{F})
  \longrightarrow
  \mathbf{P}(\mathcal{E}\otimes_{\mathcal{O}_Y}\mathcal{F}).
\tag{4.3.1.2}
\]
을 얻는다.
$Y=\operatorname{Spec}(A)$가 아핀이고,
$\mathcal{E}=\widetilde{E}$, $\mathcal{F}=\widetilde{F}$이며 $E$와 $F$가
두 $A$-가군인 경우에 사상 $\varsigma$를 구체적으로 기술하자. 이때
$\mathcal{E}\otimes_{\mathcal{O}_Y}\mathcal{F}=\widetilde{E\otimes_A F}$
이다 (\hyperref[I.1.3.12-ko]{I, 1.3.12})~; 다음과 같이 놓자.
$R=\mathbf{S}_A(E)$, $S=\mathbf{S}_A(F)$,
$T=\mathbf{S}_A(E\otimes_A F)$~; $f\in E$, $g\in F$라 하고 $Q$의
아핀 열린집합
\[
  D_+(f)\times_YD_+(g)=\operatorname{Spec}(B)
\]

\oldpage[II]{77}

을 생각하자. 여기서 $B=R_{(f)}\otimes_AS_{(g)}$이다~; 이 아핀
열린집합에 대한 $\mathcal{L}$의 제한은 $\widetilde{L}$이며, 여기서
\[
  L=(R(1))_{(f)}\otimes_A(S(1))_{(g)}
\]
이고, 원소 $c=(f/1)\otimes(g/1)$은 $L$을 자유 $B$-가군으로 보았을
때의 생성원이다 (\hyperref[II.2.5.7-ko]{2.5.7}). 준동형
(\hyperref[II.4.3.1.1-ko]{4.3.1.1})은
\[
  (x\otimes y)\otimes b\longmapsto b((x/1)\otimes(y/1))
\]
로 주어지는 $(E\otimes_A F)\otimes_AB$에서 $L$로 가는 준동형에
대응한다. (\hyperref[II.3.7.2-ko]{3.7.2})의 표기법을 쓰면, 따라서
$v_1(x\otimes y)=(x/f)\otimes(y/g)$이다~; $D_+(f)\times_YD_+(g)$에
대한 $\varsigma$의 제한은 이 아핀 스킴에서 $D_+(f\otimes g)$로 가는
사상이며, 다음을 만족하는 환 준동형
$\omega:T_{(f\otimes g)}\to R_{(f)}\otimes_AS_{(g)}$에 대응한다.
\[
\label{II.4.3.1.3-ko}
  \omega((x\otimes y)/(f\otimes g))=(x/f)\otimes(y/g)
\tag{4.3.1.3}
\]
여기서 $x\in E$이고 $y\in F$이다.
\end{env}

\begin{env}[4.3.2]
\label{II.4.3.2-ko}
(\hyperref[II.4.2.3-ko]{4.2.3})에 따라 표준 동형사상
\[
\label{II.4.3.2.1-ko}
  \tau:\varsigma^*(\mathcal{O}_P(1))\xrightarrow{\sim}
  \mathcal{O}_{P_1}(1)\otimes_Y\mathcal{O}_{P_2}(1)
\tag{4.3.2.1}
\]
이 있다. 여기서
$P=\mathbf{P}(\mathcal{E}\otimes_{\mathcal{O}_Y}\mathcal{F})$로 놓았다.
$x\in\Gamma(Y,\mathcal{E})$와 $y\in\Gamma(Y,\mathcal{F})$에 대하여
\[
\label{II.4.3.2.2-ko}
  \tau(\alpha_1(x\otimes y))=\alpha_1(x)\otimes\alpha_1(y).
\tag{4.3.2.2}
\]
임을 보이자. 실제로 $Y$가 아핀인 경우로 환원되며, 이때
(\hyperref[II.4.3.1-ko]{4.3.1})과
(\hyperref[II.2.6.2-ko]{2.6.2})의 표기법을 쓰면
$\alpha_1^{f\otimes g}(x\otimes y)=(x\otimes y)/1$은
$(T(1))_{(f\otimes g)}$에 속하고, $\alpha_1^f(x)=x/1$은
$(R(1))_{(f)}$에 속하며, $\alpha_1^g(y)=y/1$은 $(S(1))_{(g)}$에
속한다. (\hyperref[II.4.2.3-ko]{4.2.3})에서 주어진 $\tau$의 정의와
(\hyperref[II.4.3.1-ko]{4.3.1})에서 계산한 $v_1$은 곧바로
(\hyperref[II.4.3.2.2-ko]{4.3.2.2})을 증명한다. 이로부터
(\hyperref[II.3.1.4-ko]{3.1.4})의 표기법으로 다음 공식을 얻는다.
\[
\label{II.4.3.2.3-ko}
  \varsigma^{-1}(P_{x\otimes y})=(P_1)_x\times_Y(P_2)_y
\tag{4.3.2.3}
\]
실제로 (\hyperref[II.3.3.3-ko]{3.3.3})을 고려하면, 공식
(\hyperref[II.4.3.2.2-ko]{4.3.2.2})에 의해
((\hyperref[I.3.2.7-ko]{I, 3.2.7})과
(\hyperref[I.3.2.3-ko]{3.2.3})을 사용하여 다시 아핀인 경우로
돌아감으로써) 다음 보조정리를 증명하는 것으로 환원된다~:

\begin{lemma}[4.3.2.4]
\label{II.4.3.2.4-ko}
$B$, $B'$를 두 $A$-대수라 하고,
$Y=\operatorname{Spec}(A)$, $Z=\operatorname{Spec}(B)$,
$Z'=\operatorname{Spec}(B')$로 놓자~; 임의의 $t\in B$, $t'\in B'$에
대하여
$D(t\otimes t')=D(t)\times_YD(t')$이다.
\end{lemma}

\begin{proof}
실제로 $p:Z\times_YZ'\to Z$와 $p':Z\times_YZ'\to Z'$가 표준
사영이면, (\hyperref[I.1.2.2.2-ko]{I, 1.2.2.2})에서
$p^{-1}(D(t))=D(t\otimes1)$과
$p'^{-1}(D(t'))=D(1\otimes t')$임을 얻는다~; 이제
(\hyperref[I.3.2.7-ko]{I, 3.2.7})과
(\hyperref[I.1.1.9.1-ko]{I, 1.1.9.1})에서 결론이 따른다. 실제로
$(t\otimes1)(1\otimes t')=t\otimes t'$이다.
\end{proof}
\end{env}

\begin{proposition}[4.3.3]
\label{II.4.3.3-ko}
세그레 사상은 닫힌 몰입 사상이다.
\end{proposition}

\begin{proof}
실제로 문제는 $Y$ 위에서 국소적이므로 $Y$가 아핀인 경우로 환원된다.
(\hyperref[II.4.3.1-ko]{4.3.1})과
(\hyperref[II.4.3.2-ko]{4.3.2})의 표기법으로, $f$가 $E$를 달리고
$g$가 $F$를 달릴 때 $f\otimes g$들이 $T$를 생성하므로
$D_+(f\otimes g)$들은 $P$ 위상의 기저를 이룬다. 다른 한편
(\hyperref[II.4.3.2.3-ko]{4.3.2.3})에 따라
$\varsigma^{-1}(D_+(f\otimes g))=D_+(f)\times_YD_+(g)$이다. 따라서
(\hyperref[I.4.2.4-ko]{I, 4.2.4})에 의해
$D_+(f)\times_YD_+(g)$에 대한 $\varsigma$의 제한이
$D_+(f\otimes g)$로 가는 닫힌 몰입 사상임을 보이면 충분하다. 그런데
같은 표기법으로 공식 (\hyperref[II.4.3.1.3-ko]{4.3.1.3})은
$\omega$가 전사임을 보이며, 이로써 증명이 끝난다.
\end{proof}

\begin{env}[4.3.4]
\label{II.4.3.4-ko}
준연접 $\mathcal{O}_Y$-가군의

\oldpage[II]{78}

준동형을 전사 준동형으로 제한하면, 세그레 사상은 $\mathcal{E}$와
$\mathcal{F}$에 관해 함자적이다. 실제로
$\mathcal{E}\to\mathcal{E}'$이 전사 $\mathcal{O}_Y$-준동형일 때 도식
\[
  \xymatrix{
    \mathbf{P}(\mathcal{E}')\times_Y\mathbf{P}(\mathcal{F})
      \ar[r]^{j\times1} \ar[d]_{\varsigma}
    & \mathbf{P}(\mathcal{E})\times_Y\mathbf{P}(\mathcal{F})
      \ar[d]^{\varsigma}
    \\\mathbf{P}(\mathcal{E}'\otimes\mathcal{F}) \ar[r]
    & \mathbf{P}(\mathcal{E}\otimes\mathcal{F})
  }
\]
이 가환임을 보이면 된다. 여기서 $j$는 표준 닫힌 몰입 사상
$\mathbf{P}(\mathcal{E}')\to\mathbf{P}(\mathcal{E})$을 나타낸다.
$P_1'=\mathbf{P}(\mathcal{E}')$로 놓고 다른 표기는
(\hyperref[II.4.3.1-ko]{4.3.1})의 것을 그대로 쓰자~; $j\times1$은 닫힌
몰입 사상이며 (\hyperref[I.4.3.1-ko]{I, 4.3.1}), 동형까지 다음이
성립한다.
\[
  (j\times1)^*(\mathcal{O}_{P_1}(1)\otimes\mathcal{O}_{P_2}(1))
  =j^*(\mathcal{O}_{P_1}(1))\otimes\mathcal{O}_{P_2}(1)
  =\mathcal{O}_{P_1'}(1)\otimes\mathcal{O}_{P_2}(1)
\]
이는 (\hyperref[II.4.1.2.1-ko]{4.1.2.1})과
(\hyperref[I.9.1.5-ko]{I, 9.1.5})에서 따르므로, 우리의 주장은
(\hyperref[II.4.2.8-ko]{4.2.8})과
(\hyperref[II.4.2.9-ko]{4.2.9})에서 따른다.
\end{env}

\begin{env}[4.3.5]
\label{II.4.3.5-ko}
(\hyperref[II.4.3.1-ko]{4.3.1})의 표기법 아래에서,
$\psi:Y'\to Y$를 사상이라 하고
$\mathcal{E}'=\psi^*(\mathcal{E})$, $\mathcal{F}'=\psi^*(\mathcal{F})$로 놓자~;
그러면 세그레 사상
$\mathbf{P}(\mathcal{E}')\times\mathbf{P}(\mathcal{F}')
\to\mathbf{P}(\mathcal{E}'\otimes\mathcal{F}')$은
$\varsigma_{(Y')}$와 동일시된다. 실제로
(\hyperref[II.4.3.1-ko]{4.3.1})의 표기법을 유지하면서 또한
$P_1'=\mathbf{P}(\mathcal{E}')$, $P_2'=\mathbf{P}(\mathcal{F}')$로
놓자~; (\hyperref[II.4.1.3.1-ko]{4.1.3.1})에 의해 $P_i'$은
$(P_i)_{(Y')}$와 동일시되므로($i=1,2$), 구조 사상
$P_1'\times P_2'\to Y'$은 $r_{(Y')}$와 동일시된다. 다른 한편,
$\mathcal{E}'\otimes\mathcal{F}'$은
$\psi^*(\mathcal{E}\otimes\mathcal{F})$와 동일시되므로
$\mathbf{P}(\mathcal{E}'\otimes\mathcal{F}')$은
$(\mathbf{P}(\mathcal{E}\otimes\mathcal{F}))_{(Y')}$와 동일시된다
(\hyperref[II.4.1.3.1-ko]{4.1.3.1}). 마지막으로,
$\mathcal{O}_{P_1'}(1)\otimes_{Y'}\mathcal{O}_{P_2'}(1)=\mathcal{L}'$은
(\hyperref[II.4.1.3.2-ko]{4.1.3.2})와
(\hyperref[I.9.1.11-ko]{I, 9.1.11})에 의해
$\mathcal{L}\otimes_Y\mathcal{O}_{Y'}$와 동일시된다. 그러면 표준 준동형
$(r_{(Y')})^*(\mathcal{E}'\otimes\mathcal{F}')\to\mathcal{L}'$은
$s_{(Y')}$와 동일시되며, 우리의 주장은
(\hyperref[II.4.2.10-ko]{4.2.10})에서 따른다.
\end{env}

\begin{remark}[4.3.6]
\label{II.4.3.6-ko}
$\mathbf{P}(\mathcal{E})$와 $\mathbf{P}(\mathcal{F})$의 합 준스킴도
마찬가지로 $\mathbf{P}(\mathcal{E}\oplus\mathcal{F})$의 닫힌
부분준스킴과 표준적으로 동형이다. 실제로 전사 준동형
$\mathcal{E}\oplus\mathcal{F}\to\mathcal{E}$,
$\mathcal{E}\oplus\mathcal{F}\to\mathcal{F}$에는 닫힌 몰입 사상
$\mathbf{P}(\mathcal{E})\to\mathbf{P}(\mathcal{E}\oplus\mathcal{F})$,
$\mathbf{P}(\mathcal{F})\to\mathbf{P}(\mathcal{E}\oplus\mathcal{F})$이
대응한다~; 따라서 이렇게 얻은 닫힌 부분준스킴들의 바탕공간에 공통점이
없음을 보이면 충분하다. 이 문제는 $Y$ 위에서 국소적이므로
(\hyperref[II.4.3.1-ko]{4.3.1})의 표기법을 채택할 수 있다~; 그런데
$\mathbf{S}_n(E)$와 $\mathbf{S}_n(F)$는
$\mathbf{S}_n(E\oplus F)$의 교집합이 $0$인 부분가군들과 동일시된다~;
또 $\mathfrak{p}$가 $\mathbf{S}(E)$의 등급 소아이디얼이고 모든
$n\geq0$에 대하여
$\mathfrak{p}\cap\mathbf{S}_n(E)\neq\mathbf{S}_n(E)$라면, 그에
대응하는 $\mathbf{S}(E\oplus F)$의 등급 소아이디얼은
$\mathbf{S}_n(E)$와의 교집합이
$\mathfrak{p}\cap\mathbf{S}_n(E)$이고 $\mathbf{S}_+(F)$를 포함함을
곧 알 수 있다~; 따라서 각각
$\operatorname{Proj}(\mathbf{S}(E))$와
$\operatorname{Proj}(\mathbf{S}(F))$의 점인 두 점은
$\operatorname{Proj}(\mathbf{S}(E\oplus F))$에서 같은 상을 가질 수 없다.
\end{remark}

\subsection{사영 다발로의 몰입 사상. 매우 풍부한 층}
\label{subsection:II.4.4-ko}

\begin{proposition}[4.4.1]
\label{II.4.4.1-ko}
$Y$를 준콤팩트 스킴 또는 그 기저공간이 뇌터인 준스킴,
$q:X\to Y$를 유한형 사상, $\mathcal{L}$을 가역
$\mathcal{O}_X$-가군이라 하자.
\begin{enumerate}
  \item[\rm{(i)}] $\mathcal{S}$를 양의 차수로 등급화된 준연접
  $\mathcal{O}_Y$-대수,
  $\psi:q^*(\mathcal{S})\to\bigoplus_{n\geq0}\mathcal{L}^{\otimes n}$을
  등급 대수의 준동형이라 하자. $r_{\mathcal{L},\psi}$가 전체에서
  정의되는 몰입 사상일 필요충분조건은, 어떤 정수 $n\geq0$과

\oldpage[II]{79}

  $\mathcal{S}_n$의 유한형 준연접 부분-$\mathcal{O}_Y$-가군
  $\mathcal{E}$가 존재하여 준동형
  \[
    \psi'=\psi_n\circ q^*(j):q^*(\mathcal{E})
      \longrightarrow\mathcal{L}^{\otimes n}=\mathcal{L}'
  \]
  (여기서 $j$는 단사 준동형 $\mathcal{E}\to\mathcal{S}_n$이다)가
  전사이고, 사상
  $r_{\mathcal{L}',\psi'}:X\to\mathbf{P}(\mathcal{E})$가 몰입
  사상인 것이다.
  \item[\rm{(ii)}] $\mathcal{F}$를 준연접 $\mathcal{O}_Y$-가군,
  $\varphi:q^*(\mathcal{F})\to\mathcal{L}$을 전사 준동형이라 하자.
  사상 $r_{\mathcal{L},\varphi}$가 몰입 사상
  $X\to\mathbf{P}(\mathcal{F})$일 필요충분조건은,
  $\mathcal{F}$의 유한형 준연접 부분-$\mathcal{O}_Y$-가군
  $\mathcal{E}$가 존재하여 준동형
  $\varphi'=\varphi\circ q^*(j):q^*(\mathcal{E})\to\mathcal{L}$
  (여기서 $j:\mathcal{E}\to\mathcal{F}$는 표준 단사 준동형이다)가
  전사이고, $r_{\mathcal{L},\varphi'}:X\to\mathbf{P}(\mathcal{E})$가
  몰입 사상인 것이다.
\end{enumerate}
\end{proposition}

\begin{proof}
\begin{enumerate}
  \item[\rm{(i)}] $r_{\mathcal{L},\psi}$가 전체에서 정의되는 몰입
  사상이라는 것은
  (\hyperref[II.3.8.5-ko]{3.8.5})에 따라 다음 조건을 만족하는
  $n\geq0$과 $\mathcal{E}$가 존재한다는 것과 동치이다.
  $\mathcal{S}'$을 $\mathcal{E}$로 생성되는 $\mathcal{S}$의
  부분대수라 할 때, 준동형
  $q^*(\mathcal{E})\to\mathcal{L}^{\otimes n}$은 전사이고
  $r_{\mathcal{L}',\psi'}:X\to\operatorname{Proj}(\mathcal{S}')$은
  전체에서 정의되는 몰입 사상이다. 다른 한편, 표준 전사 준동형
  $\mathbf{S}(\mathcal{E})\to\mathcal{S}'$이 있고, 이에 닫힌 몰입 사상
  $\operatorname{Proj}(\mathcal{S}')\to\mathbf{P}(\mathcal{E})$이
  대응한다 (\hyperref[II.3.6.2-ko]{3.6.2})~; 이로부터 결론이 따른다.
  \item[\rm{(ii)}] $\mathcal{F}$는 그 유한형 준연접 부분가군
  $\mathcal{E}_\lambda$들의 귀납적 극한이므로
  (\hyperref[I.9.4.9-ko]{I, 9.4.9}), $\mathbf{S}(\mathcal{F})$는
  $\mathbf{S}(\mathcal{E}_\lambda)$들의 귀납적 극한이다~; 따라서
  (\hyperref[II.3.8.4-ko]{3.8.4})에서 결론이 따른다. 이때
  (\hyperref[II.3.8.4-ko]{3.8.4})의 증명에서 모든 $n_i$를 $1$로
  택할 수 있음을 관찰하면 된다~: 실제로 ($Y$가 아핀이라고 가정할 때)
  $r=r_{\mathcal{L},\varphi}$가 몰입 사상이면, $r(X)$는
  $\mathbf{P}(\mathcal{F})$의 준콤팩트 부분공간이며, 이를
  $f\in F$이고 $D_+(f)\cap r(X)$가 닫힌집합인 꼴의
  $\mathbf{P}(\mathcal{F})$의 열린집합 $D_+(f)$ 유한 개로 덮을 수 있다.
\end{enumerate}
\end{proof}

\begin{definition}[4.4.2]
\label{II.4.4.2-ko}
$Y$를 준스킴, $q:X\to Y$를 사상이라 하자. 가역
$\mathcal{O}_X$-가군 $\mathcal{L}$에 대하여, 준연접
$\mathcal{O}_Y$-가군 $\mathcal{E}$와 $X$에서
$P=\mathbf{P}(\mathcal{E})$로 가는 $Y$-몰입 사상 $i$가 존재하여
$\mathcal{L}$이 $i^*(\mathcal{O}_P(1))$과 동형이면,
$\mathcal{L}$을 \emph{$q$에 관해 매우 풍부하다} (또는
\emph{$Y$에 관해 매우 풍부하다}, 혼동의 우려가 없으면 단순히
\emph{매우 풍부하다})고 한다.
\end{definition}

동치로 (\hyperref[II.4.2.3-ko]{4.2.3}), 준연접
$\mathcal{O}_Y$-가군 $\mathcal{E}$와 전사 준동형
$\varphi:q^*(\mathcal{E})\to\mathcal{L}$가 존재하여
$r_{\mathcal{L},\varphi}:X\to\mathbf{P}(\mathcal{E})$가 몰입
사상이라고 말할 수도 있다.

$Y$에 관해 매우 풍부한 $\mathcal{O}_X$-가군이 존재하면 $q$가
분리 사상임을 유의하자
((\hyperref[II.3.1.3-ko]{3.1.3}) 및
(\hyperref[I.5.5.1-ko]{I, 5.5.1, (i)와 (ii)})).

\begin{corollary}[4.4.3]
\label{II.4.4.3-ko}
$\mathcal{S}_1$로 생성되는 준연접 등급 $\mathcal{O}_Y$-대수
$\mathcal{S}$와 $Y$-몰입 사상
$i:X\to P=\operatorname{Proj}(\mathcal{S})$가 존재하여
$\mathcal{L}$이 $i^*(\mathcal{O}_P(1))$과 동형이라고 가정하자~;
그러면 $\mathcal{L}$은 $q$에 관해 매우 풍부하다.
\end{corollary}

\begin{proof}
실제로 $\mathcal{F}=\mathcal{S}_1$이면 표준 준동형
$\mathbf{S}(\mathcal{F})\to\mathcal{S}$는 전사이다. 따라서 이에
대응하는 닫힌 몰입 사상
$\operatorname{Proj}(\mathcal{S})\to\mathbf{P}(\mathcal{F})$
(\hyperref[II.3.6.2-ko]{3.6.2})과 몰입 사상 $i$를 합성하면,
$\mathcal{L}$이 $j^*(\mathcal{O}_{P'}(1))$과 동형인 몰입 사상
$j:X\to\mathbf{P}(\mathcal{F})=P'$을 얻는다
(\hyperref[II.3.6.3-ko]{3.6.3}).
\end{proof}

\begin{proposition}[4.4.4]
\label{II.4.4.4-ko}
$q:X\to Y$를 준콤팩트 사상, $\mathcal{L}$을 가역
$\mathcal{O}_X$-가군이라 하자. 다음 성질들은 동치이다~:
\begin{enumerate}
  \item[\rm{a)}] $\mathcal{L}$은 $q$에 관해 매우 풍부하다.
  \item[\rm{b)}] $q_*(\mathcal{L})$은 준연접이고, 표준 준동형
  $\sigma:q^*(q_*(\mathcal{L}))\to\mathcal{L}$은 전사이며, 사상
  $r_{\mathcal{L},\sigma}:X\to\mathbf{P}(q_*(\mathcal{L}))$은 몰입
  사상이다.
\end{enumerate}
\end{proposition}

\begin{proof}
$q$가 준콤팩트이므로, $q$가 분리 사상이면 $q_*(\mathcal{L})$이
준연접임을 안다 (\hyperref[I.9.2.2-ko]{I, 9.2.2}).

\oldpage[II]{80}

전사 준동형 $\varphi:q^*(\mathcal{E})\to\mathcal{L}$
($\mathcal{E}$는 준연접 $\mathcal{O}_Y$-가군이다)의 존재로부터
$\sigma$가 전사임을 안다 (\hyperref[II.3.4.7-ko]{3.4.7})~; 더 나아가
$\varphi$의 인수분해
$q^*(\mathcal{E})\to q^*(q_*(\mathcal{L}))\xrightarrow{\sigma}\mathcal{L}$
에는 다음 인수분해가 표준적으로 대응한다.
\[
  q^*(\mathbf{S}(\mathcal{E}))
    \longrightarrow q^*(\mathbf{S}(q_*(\mathcal{L})))
    \longrightarrow\bigoplus_{n\geq0}\mathcal{L}^{\otimes n}
\]
따라서 (\hyperref[II.3.8.3-ko]{3.8.3})에 의해
$r_{\mathcal{L},\varphi}$가 몰입 사상이라는 가정으로부터
$j=r_{\mathcal{L},\sigma}$도 그러함이 따른다~; 또한
(\hyperref[II.4.2.4-ko]{4.2.4})에 의해,
$P'=\mathbf{P}(q_*(\mathcal{L}))$로 놓으면 $\mathcal{L}$은
$j^*(\mathcal{O}_{P'}(1))$과 동형이다. 따라서 a)와 b)는 동치이다.
\end{proof}

\begin{corollary}[4.4.5]
\label{II.4.4.5-ko}
$q$가 준콤팩트라고 가정하자. $\mathcal{L}$이 $Y$에 관해 매우 풍부할
필요충분조건은, $Y$의 열린 덮개 $(U_\alpha)$가 존재하여 모든
$\alpha$에 대하여 $\mathcal{L}|q^{-1}(U_\alpha)$가 $U_\alpha$에 관해
매우 풍부한 것이다.
\end{corollary}

\begin{proof}
실제로 (\hyperref[II.4.4.4-ko]{4.4.4})의 조건 b)는 $Y$ 위에서
국소적이다.
\end{proof}

\begin{proposition}[4.4.6]
\label{II.4.4.6-ko}
$Y$를 준콤팩트 스킴 또는 그 기저공간이 뇌터인 준스킴,
$q:X\to Y$를 유한형 사상, $\mathcal{L}$을 가역
$\mathcal{O}_X$-가군이라 하자. 그러면
(\hyperref[II.4.4.4-ko]{4.4.4})의 조건 a), b)는 다음 조건들과
동치이다~:
\begin{enumerate}
  \item[\rm{a')}] 유한형 준연접 $\mathcal{O}_Y$-가군 $\mathcal{E}$와
  전사 준동형 $\varphi:q^*(\mathcal{E})\to\mathcal{L}$이 존재하여
  $r_{\mathcal{L},\varphi}$가 몰입 사상이다.
  \item[\rm{b')}] $q_*(\mathcal{L})$의 유한형 연접
  부분-$\mathcal{O}_Y$-가군 $\mathcal{E}$가 존재하여 a')에서 서술한
  성질들을 갖는다.
\end{enumerate}
\end{proposition}

\begin{proof}
a') 또는 b')가 a)를 함의함은 명백하다~; 다른 한편,
(\hyperref[II.4.4.1-ko]{4.4.1})에 의해 a)는 a')를 함의하고, 마찬가지로
b)는 b')를 함의한다.
\end{proof}

\begin{corollary}[4.4.7]
\label{II.4.4.7-ko}
$Y$가 준콤팩트 스킴 또는 뇌터 준스킴이라고 가정하자.
$\mathcal{L}$이 $q$에 관해 매우 풍부하면, $\mathcal{S}_1$이 유한형이고
$\mathcal{S}$를 생성하는 준연접 등급 $\mathcal{O}_Y$-대수
$\mathcal{S}$와 우세인 열린 $Y$-몰입 사상
$i:X\to P=\operatorname{Proj}(\mathcal{S})$가 존재하여
$\mathcal{L}$은 $i^*(\mathcal{O}_P(1))$과 동형이다.
\end{corollary}

\begin{proof}
실제로 (\hyperref[II.4.4.6-ko]{4.4.6})의 조건 b')가 성립한다~;
구조 사상 $p:\mathbf{P}(\mathcal{E})=P'\to Y$는 이때 분리 사상이자
유한형이므로 (\hyperref[II.3.1.3-ko]{3.1.3}), $Y$가 준콤팩트 스킴
(각각 뇌터 준스킴)이면 $P'$도 준콤팩트 스킴(각각 뇌터 준스킴)이다.
몰입 사상 $j=r_{\mathcal{L},\varphi}$에 대응하는 $P'$의 부분준스킴
$X'$의 폐포 (\hyperref[I.9.5.11-ko]{I, 9.5.11})를 $Z$라 하자~;
$j$는 우세인 열린 몰입 사상 $i:X\to Z$와 표준 포함 사상
$Z\to P'$의 합성으로 인수분해됨이 명백하다. 그런데 $Z$는 준스킴
$\operatorname{Proj}(\mathcal{S})$와 동일시되며, 여기서 $\mathcal{S}$는
준연접 등급 아이디얼층으로 $\mathbf{S}(\mathcal{E})$을 나눈 몫인 등급
$\mathcal{O}_Y$-대수이다 (\hyperref[II.3.6.2-ko]{3.6.2}). 또한
$\mathcal{S}_1$이 유한형이고 $\mathcal{S}$를 생성함은 명백하다~;
더 나아가 $\mathcal{O}_Z(1)$은 표준 포함 사상에 의한
$\mathcal{O}_{P'}(1)$의 역상이므로
(\hyperref[II.3.6.3-ko]{3.6.3}),
$\mathcal{L}=i^*(\mathcal{O}_Z(1))$이다.
\end{proof}

\begin{proposition}[4.4.8]
\label{II.4.4.8-ko}
$q:X\to Y$를 사상, $\mathcal{L}$을 $q$에 관해 매우 풍부한
$\mathcal{O}_X$-가군, $\mathcal{L}'$을 다음 조건을 만족하는 가역
$\mathcal{O}_X$-가군이라 하자. 즉, 준연접 $\mathcal{O}_Y$-가군
$\mathcal{E}'$과 전사 준동형
$q^*(\mathcal{E}')\to\mathcal{L}'$이 존재한다고 하자. 그러면
$\mathcal{L}\otimes_{\mathcal{O}_X}\mathcal{L}'$은 $q$에 관해 매우
풍부하다.
\end{proposition}

\begin{proof}
실제로 가정에 따라
$\mathcal{L}'=r'^*(\mathcal{O}_{P'}(1))$을 만족하는 $Y$-사상
$r':X\to P'=\mathbf{P}(\mathcal{E}')$이 존재한다
(\hyperref[II.4.2.1-ko]{4.2.1}). 또한 가정에 따라 준연접
$\mathcal{O}_Y$-가군 $\mathcal{E}$와

\oldpage[II]{81}

$\mathcal{L}=r^*(\mathcal{O}_P(1))$을 만족하는 $Y$-몰입 사상
$r:X\to P=\mathbf{P}(\mathcal{E})$가 존재한다.
$Q=\mathbf{P}(\mathcal{E}\otimes\mathcal{E}')$라 놓고 세그레 사상
$\varsigma:P\times_YP'\to Q$를 생각하자
(\hyperref[II.4.3.1-ko]{4.3.1}). $r$이 몰입 사상이므로
$(r,r')_Y:X\to P\times_YP'$도 몰입 사상이다
(\hyperref[I.5.3.14-ko]{I, 5.3.14})~; 한편 $\varsigma$도 몰입
사상이므로 (\hyperref[II.4.3.3-ko]{4.3.3}),
마찬가지로
$r'':X\xrightarrow{(r,r')_Y}P\times_YP'\xrightarrow{\varsigma}Q$도
몰입 사상이다. 다른 한편 (\hyperref[II.4.3.2.1-ko]{4.3.2.1}),
$\varsigma^*(\mathcal{O}_Q(1))$은
$\mathcal{O}_P(1)\otimes_Y\mathcal{O}_{P'}(1)$과 동형이므로,
(\hyperref[I.9.1.4-ko]{I, 9.1.4})에 따라
$r''^*(\mathcal{O}_Q(1))$은 $\mathcal{L}\otimes\mathcal{L}'$과
동형이다. 이로써 명제가 증명된다.
\end{proof}

\begin{corollary}[4.4.9]
\label{II.4.4.9-ko}
$q:X\to Y$를 사상이라 하자.
\begin{enumerate}
  \item[\rm{(i)}] $\mathcal{L}$을 가역 $\mathcal{O}_X$-가군,
  $\mathcal{K}$를 가역 $\mathcal{O}_Y$-가군이라 하자.
  $\mathcal{L}$이 $q$에 관해 매우 풍부할 필요충분조건은
  $\mathcal{L}\otimes q^*(\mathcal{K})$가 그러한 것이다.
  \item[\rm{(ii)}] $\mathcal{L}$과 $\mathcal{L}'$이 $q$에 관해
  매우 풍부한 두 $\mathcal{O}_X$-가군이면,
  $\mathcal{L}\otimes\mathcal{L}'$도 그러하다~; 특히 모든
  $n>0$에 대하여 $\mathcal{L}^{\otimes n}$은 $q$에 관해 매우
  풍부하다.
\end{enumerate}
\end{corollary}

\begin{proof}
(ii)는 (\hyperref[II.4.4.8-ko]{4.4.8})의 직접적인 귀결이고,
(i)의 필요조건 방향도 마찬가지이다~; 다른 한편
$\mathcal{L}\otimes q^*(\mathcal{K})$가 매우 풍부하면 앞의 결과에
따라
$(\mathcal{L}\otimes q^*(\mathcal{K}))\otimes q^*(\mathcal{K}^{-1})$
도 그러하며, 이 마지막 $\mathcal{O}_X$-가군은 $\mathcal{L}$과
동형이다 (\hyperref[0.5.4.3-ko]{0, 5.4.3} 및
\hyperref[0.5.4.5-ko]{5.4.5}).
\end{proof}

\begin{proposition}[4.4.10]
\label{II.4.4.10-ko}
\begin{enumerate}
  \item[\rm{(i)}] 임의의 준스킴 $Y$에 대하여, 모든 가역
  $\mathcal{O}_Y$-가군 $\mathcal{L}$은 항등사상 $1_Y$에 관해 매우
  풍부하다.
  \item[\rm{(i bis)}] $f:X\to Y$를 사상, $j:X'\to X$를 몰입이라 하자.
  $\mathcal{L}$이 $f$에 관해 매우 풍부한 $\mathcal{O}_X$-가군이면,
  $j^*(\mathcal{L})$은 $f\circ j$에 관해 매우 풍부하다.
  \item[\rm{(ii)}] $Z$를 준콤팩트 준스킴, $f:X\to Y$를 유한형 사상,
  $g:Y\to Z$를 준콤팩트 사상, $\mathcal{L}$을 $f$에 관해 매우 풍부한
  $\mathcal{O}_X$-가군, $\mathcal{K}$를 $g$에 관해 매우 풍부한
  $\mathcal{O}_Y$-가군이라 하자. 그러면 모든 $n\geq n_0$에 대하여
  $\mathcal{L}\otimes f^*(\mathcal{K}^{\otimes n})$이 $g\circ f$에 관해
  매우 풍부하도록 하는 정수 $n_0>0$이 존재한다.
  \item[\rm{(iii)}] $f:X\to Y$, $g:Y'\to Y$를 두 사상이라 하고
  $X'=X_{(Y')}$라 놓자. $\mathcal{L}$이 $f$에 관해 매우 풍부한
  $\mathcal{O}_X$-가군이면,
  $\mathcal{L}'=\mathcal{L}\otimes_Y\mathcal{O}_{Y'}$은
  $f_{(Y')}$에 관해 매우 풍부한 $\mathcal{O}_{X'}$-가군이다.
  \item[\rm{(iv)}] $f_i:X_i\to Y_i$ ($i=1,2$)를 두 $S$-사상이라 하자.
  $\mathcal{L}_i$가 $f_i$에 관해 매우 풍부한
  $\mathcal{O}_{X_i}$-가군이면 ($i=1,2$),
  $\mathcal{L}_1\otimes_S\mathcal{L}_2$는 $f_1\times_Sf_2$에 관해
  매우 풍부하다.
  \item[\rm{(v)}] $f:X\to Y$, $g:Y\to Z$를 두 사상이라 하자.
  $\mathcal{O}_X$-가군 $\mathcal{L}$이 $g\circ f$에 관해 매우
  풍부하면, $\mathcal{L}$은 $f$에 관해 매우 풍부하다.
  \item[\rm{(vi)}] $f:X\to Y$를 사상, $j$를 표준 포함 사상
  $X_{\mathrm{red}}\to X$라 하자. $\mathcal{L}$이 $f$에 관해 매우
  풍부한 $\mathcal{O}_X$-가군이면, $j^*(\mathcal{L})$은
  $f_{\mathrm{red}}$에 관해 매우 풍부하다.
\end{enumerate}
\end{proposition}

\begin{proof}
성질 (i bis)는 정의
(\hyperref[II.4.4.2-ko]{4.4.2})에서 곧바로 따르며, (vi)는
(\hyperref[I.5.5.12-ko]{I, 5.5.12})의 논증을 그대로 본뜬 형식적
논증에 의해 (i bis)와 (v)에서 도출됨이 명백하므로, 그 논증은 독자에게
맡긴다. (v)를 증명하기 위하여
(\hyperref[I.5.5.12-ko]{I, 5.5.12})에서와 같이 분해
$X\xrightarrow{\Gamma_f}X\times_ZY\xrightarrow{p_2}Y$를 생각하자. 이때
$p_2=(g\circ f)\times1_Y$임에 유의한다. 가정과 (i), (iv)에서
$\mathcal{L}\otimes_{\mathcal{O}_Z}\mathcal{O}_Y$가 $p_2$에 관해 매우
풍부함이 따른다~; 다른 한편
$\mathcal{L}=\Gamma_f^*(\mathcal{L}\otimes_{\mathcal{O}_Z}\mathcal{O}_Y)$
(\hyperref[I.9.1.4-ko]{I, 9.1.4})이고, $\Gamma_f$는 몰입이다
(\hyperref[I.5.3.11-ko]{I, 5.3.11})~; 따라서 (i bis)를 적용할 수 있다.

\oldpage[II]{82}

(i)를 증명하려면 $\mathcal{E}=\mathcal{L}$로 놓고 정의
(\hyperref[II.4.4.2-ko]{4.4.2})를 적용하면 된다. 실제로 이때
$\mathbf{P}(\mathcal{E})$는 $Y$와 동일시된다
(\hyperref[II.4.1.4-ko]{4.1.4}).

(iii)을 증명하자. 준연접 $\mathcal{O}_Y$-가군 $\mathcal{E}$와
$\mathcal{L}=i^*(\mathcal{O}_P(1))$을 만족하는 $Y$-몰입 사상
$i:X\to\mathbf{P}(\mathcal{E})=P$가 존재한다~; 이제
$\mathcal{E}'=g^*(\mathcal{E})$라 놓으면, $\mathcal{E}'$은 준연접
$\mathcal{O}_{Y'}$-가군이고
$P'=\mathbf{P}(\mathcal{E}')=P_{(Y')}$
(\hyperref[II.4.1.3.1-ko]{4.1.3.1})이며, $i_{(Y')}$은
$X_{(Y')}$에서 $P'$로 가는 몰입이고
(\hyperref[I.4.3.2-ko]{I, 4.3.2}), $\mathcal{L}'$은
$(i_{(Y')})^*(\mathcal{O}_{P'}(1))$과 동형이다
(\hyperref[II.4.2.10-ko]{4.2.10}).

(iv)를 증명하기 위하여, 가정에 따라 준연접
$\mathcal{O}_{Y_i}$-가군 $\mathcal{E}_i$와 $Y_i$-몰입 사상
$r_i:X_i\to P_i=\mathbf{P}(\mathcal{E}_i)$가 존재하고
$\mathcal{L}_i=r_i^*(\mathcal{O}_{P_i}(1))$임을 관찰하자 ($i=1,2$)~;
$r_1\times_Sr_2$는 $X_1\times_SX_2$에서 $P_1\times_SP_2$로 가는
$S$-몰입이고 (\hyperref[I.4.3.1-ko]{I, 4.3.1}), 이 몰입에 의한
$\mathcal{O}_{P_1}(1)\otimes_S\mathcal{O}_{P_2}(1)$의 역상은
$\mathcal{L}_1\otimes_S\mathcal{L}_2$이다. 다른 한편
$T=Y_1\times_SY_2$라 놓고, $p_1,p_2$를 각각 $T$에서 $Y_1,Y_2$로 가는
사영이라 하자. $P_i'=\mathbf{P}(p_i^*(\mathcal{E}_i))$ ($i=1,2$)라 놓으면,
(\hyperref[II.4.1.3.1-ko]{4.1.3.1})에 따라
$P_i'=P_i\times_{Y_i}T$이고, 따라서
\[
  P_1'\times_TP_2'
  =(P_1\times_{Y_1}T)\times_T(P_2\times_{Y_2}T)
  =P_1\times_{Y_1}(T\times_{Y_2}P_2)
  =P_1\times_{Y_1}(Y_1\times_SP_2)
  =P_1\times_SP_2
\]
가 표준 동형사상까지 성립한다. 마찬가지로
$\mathcal{O}_{P_i'}(1)=\mathcal{O}_{P_i}(1)\otimes_{Y_i}\mathcal{O}_T$
(\hyperref[II.4.1.3.2-ko]{4.1.3.2})이고, 특히
(\hyperref[I.9.1.9.1-ko]{I, 9.1.9.1}과
\hyperref[I.9.1.2-ko]{9.1.2})에 바탕을 둔 유사한 계산에 의하면, 앞의
동일시에서 $\mathcal{O}_{P_1'}(1)\otimes_T\mathcal{O}_{P_2'}(1)$은
$\mathcal{O}_{P_1}(1)\otimes_S\mathcal{O}_{P_2}(1)$과 동일시된다.
따라서 $r_1\times_Sr_2$를 $X_1\times_SX_2$에서
$P_1'\times_TP_2'$로 가는 $T$-몰입으로 볼 수 있고, 이 몰입에 의한
$\mathcal{O}_{P_1'}(1)\otimes_T\mathcal{O}_{P_2'}(1)$의 역상은
$\mathcal{L}_1\otimes_S\mathcal{L}_2$이다. 이제 세그레 사상을 사용하여
(\hyperref[II.4.4.8-ko]{4.4.8})에서와 같이 논증을 마치면 된다.

(ii)를 증명하는 일만 남았다. 먼저 $Z$가 아핀 스킴인 경우로 제한할 수
있다. 실제로 일반적인 경우에는 $Z$의 유한 아핀 열린 덮개 $(U_i)$가
존재한다~; $\mathcal{K}|g^{-1}(U_i)$,
$\mathcal{L}|f^{-1}(g^{-1}(U_i))$ 및 정수 $n_i$에 대하여 명제가
증명되었다면, $n_0$로 $n_i$들 중 가장 큰 것을 택하는 것만으로
$\mathcal{K}$와 $\mathcal{L}$에 대한 명제를 증명할 수 있다
(\hyperref[II.4.4.5-ko]{4.4.5}). 가정에 따라 $f$와 $g$는 분리
사상이므로, $X$와 $Y$는 준콤팩트 스킴이다.

유한형 준연접 $\mathcal{O}_Y$-가군 $\mathcal{E}$와 몰입 사상
$r:X\to P=\mathbf{P}(\mathcal{E})$가 존재하고
$\mathcal{L}=r^*(\mathcal{O}_P(1))$이다
(\hyperref[II.4.4.6-ko]{4.4.6}). 합성사상
$P\xrightarrow{h}Y\xrightarrow{g}Z$에 관해 매우 풍부한
$\mathcal{O}_P$-가군 $\mathcal{M}$과 정수 $m$이 존재하여
$\mathcal{O}_P(1)$이
$\mathcal{M}\otimes_Y\mathcal{K}^{\otimes(-m)}$과 동형임을 보이자.
$n\geq m+1$이면, 가정과 사상 $h:P\to Y$ 및 $1_Y$에 적용한 (iv)에
따라 $\mathcal{O}_P(1)\otimes_Y\mathcal{K}^{\otimes n}$은 $Z$에 관해
매우 풍부하다~; $r$은 몰입이고
$\mathcal{L}\otimes f^*(\mathcal{K}^{\otimes n})
=r^*(\mathcal{O}_P(1)\otimes_Y\mathcal{K}^{\otimes n})$이므로, 결론은
(i bis)에서 따른다. $\mathcal{O}_P(1)$에 관한 위 주장을 증명하기 위해
다음 보조정리를 사용한다~:

\begin{lemma}[4.4.10.1]
\label{II.4.4.10.1-ko}
$Z$를 준콤팩트 스킴 또는 그 기저공간이 뇌터인 준스킴,
$g:Y\to Z$를 준콤팩트 사상, $\mathcal{K}$를 $g$에 관해 매우 풍부한
가역 $\mathcal{O}_Y$-가군, $\mathcal{E}$를 유한형 준연접
$\mathcal{O}_Y$-가군이라 하자. 그러면 어떤 정수 $m_0$가 존재하여,
모든 $m\geq m_0$에 대하여 $\mathcal{E}$는
$g^*(\mathcal{F})\otimes\mathcal{K}^{\otimes(-m)}$ 꼴의
$\mathcal{O}_Y$-가군의 몫과 동형이다. 여기서 $\mathcal{F}$는
$m$에 의존하는 유한형 준연접 $\mathcal{O}_Z$-가군이다.
\end{lemma}

이 보조정리는 (\hyperref[II.4.5.10.1-ko]{4.5.10.1})에서 증명할 것이다~;
독자는 (\hyperref[II.4.4.10-ko]{4.4.10})이 4.5절 어디에서도 사용되지
않음을 확인할 수 있다.

\oldpage[II]{83}

이 보조정리를 인정하면, $P$에서
\[
  P_1=\mathbf{P}(g^*(\mathcal{F})\otimes\mathcal{K}^{\otimes(-m)})
\]
으로 가는 닫힌 몰입 사상 $j_1$이 존재하여 $\mathcal{O}_P(1)$은
$j_1^*(\mathcal{O}_{P_1}(1))$과 동형이다
(\hyperref[II.4.1.2-ko]{4.1.2}). 다른 한편, $P_1$에서
$P_2=\mathbf{P}(g^*(\mathcal{F}))$ 위로 가는 동형사상이 존재하고,
이 동형사상은 $\mathcal{O}_{P_1}(1)$을
$\mathcal{O}_{P_2}(1)\otimes_Y\mathcal{K}^{\otimes(-m)}$과 동일시한다
(\hyperref[II.4.1.4-ko]{4.1.4})~; 따라서 닫힌 몰입 사상
$j_2:P\to P_2$가 존재하여 $\mathcal{O}_P(1)$은
$j_2^*(\mathcal{O}_{P_2}(1))\otimes_Y\mathcal{K}^{\otimes(-m)}$과
동형이다. 끝으로 $P_2$는 $P_3=\mathbf{P}(\mathcal{F})$인
$P_3\times_ZY$와 동일시되고, $\mathcal{O}_{P_2}(1)$은
$\mathcal{O}_{P_3}(1)\otimes_Z\mathcal{O}_Y$와 동일시된다
(\hyperref[II.4.1.3-ko]{4.1.3}). 정의상 $\mathcal{O}_{P_3}(1)$은
$Z$에 관해 매우 풍부하다~; $\mathcal{K}$도 그러하므로 (iv)에 따라
$\mathcal{O}_{P_2}(1)\otimes_Y\mathcal{K}$는 $Z$에 관해 매우
풍부하다~; (i bis)에 따라
$\mathcal{M}=j_2^*(\mathcal{O}_{P_2}(1)\otimes_Y\mathcal{K})$도
그러하고, $\mathcal{O}_P(1)$은
$\mathcal{M}\otimes_Y\mathcal{K}^{\otimes(-m-1)}$과 동형이다. 이로써
증명이 끝난다.
\end{proof}

\begin{proposition}[4.4.11]
\label{II.4.4.11-ko}
$f:X\to Y$, $f':X'\to Y$를 두 사상, $X''$을 합 준스킴
$X\sqcup X'$, $f''$을 $X$에서 $f$와, $X'$에서 $f'$과 각각 일치하는
사상 $X''\to Y$라 하자. $\mathcal{L}$ (각각 $\mathcal{L}'$)을 가역
$\mathcal{O}_X$-가군 (각각 가역 $\mathcal{O}_{X'}$-가군)이라 하고,
$\mathcal{L}''$을 $X$에서 $\mathcal{L}$과, $X'$에서
$\mathcal{L}'$과 각각 일치하는 가역 $\mathcal{O}_{X''}$-가군이라 하자.
$\mathcal{L}''$이 $f''$에 관해 매우 풍부할 필요충분조건은
$\mathcal{L}$이 $f$에 관해 매우 풍부하고 $\mathcal{L}'$이 $f'$에
관해 매우 풍부한 것이다.
\end{proposition}

\begin{proof}
$Y$가 아핀인 경우로 곧바로 환원된다. $\mathcal{L}''$이 매우
풍부하면 (\hyperref[II.4.4.10-ko]{4.4.10, (i bis)})에 따라
$\mathcal{L}$과 $\mathcal{L}'$도 그러하다. 역으로 $\mathcal{L}$과
$\mathcal{L}'$이 매우 풍부하면, 정의
(\hyperref[II.4.4.2-ko]{4.4.2})와
(\hyperref[II.4.3.6-ko]{4.3.6})에 따라 $\mathcal{L}''$이 매우
풍부하다는 것이 곧바로 따른다.
\end{proof}

\subsection{풍부한 층}
\label{subsection:II.4.5-ko}

\begin{env}[4.5.1]
\label{II.4.5.1-ko}
준스킴 $X$와 가역 $\mathcal{O}_X$-가군 $\mathcal{L}$이 주어졌다고
하자. 모든 $\mathcal{O}_X$-가군 $\mathcal{F}$에 대하여
($\mathcal{L}$에 관해 혼동의 여지가 없을 때)
$\mathcal{F}(n)=\mathcal{F}\otimes_{\mathcal{O}_X}\mathcal{L}^{\otimes n}$
($n\in\mathbf{Z}$)으로 놓겠다~; 한편
\[
  S=\bigoplus_{n\geq0}\Gamma(X,\mathcal{L}^{\otimes n})
\]
으로 놓겠다. 이는 (\hyperref[0.5.4.6-ko]{0, 5.4.6})에서 정의한 환
$\Gamma_\bullet(\mathcal{L})$의 등급 부분환이다. $X$를
$\mathbf{Z}$-준스킴으로 보고 $p$로 구조 사상
$X\to\operatorname{Spec}(\mathbf{Z})$를 나타내면, 등급
$\mathcal{O}_X$-대수의 준동형
\[
  p^*(\widetilde{S})\longrightarrow
  \bigoplus_{n\geq0}\mathcal{L}^{\otimes n}
\]
들과 등급환 $S$의 자기준동형들 사이에는 전단사 대응이 있다
(\hyperref[I.2.2.5-ko]{I, 2.2.5})~; 이 대응에서 $S$의 항등
자기동형사상에 대응하는 준동형
\[
  \varepsilon:p^*(\widetilde{S})\longrightarrow
  \bigoplus_{n\geq0}\mathcal{L}^{\otimes n}
\]
을 표준 준동형이라 하겠다. 이에 대응하는
(\hyperref[II.3.7.1-ko]{3.7.1}) 사상
$G(\varepsilon)\to\operatorname{Proj}(S)$도 표준 사상이라 하겠다.
\end{env}

\begin{theorem}[4.5.2]
\label{II.4.5.2-ko}
$X$를 준콤팩트 스킴 또는 그 바탕공간이 뇌터인 준스킴,
$\mathcal{L}$을 가역 $\mathcal{O}_X$-가군, $S$를 등급환
$\bigoplus_{n\geq0}\Gamma(X,\mathcal{L}^{\otimes n})$이라 하자. 다음
조건들은 동치이다~:
\begin{enumerate}
  \item[\rm{a)}] $f$가 $S_+$의 동차 원소 전체를 움직일 때,
  $X_f$들은 $X$의 위상의 기저를 이룬다.
  \item[\rm{a')}] $f$가 $S_+$의 동차 원소 전체를 움직일 때,
  $X_f$들 가운데 아핀인 것들이 $X$의 덮개를 이룬다.
  \item[\rm{b)}] 표준 사상
  $G(\varepsilon)\to\operatorname{Proj}(S)$
  (\hyperref[II.4.5.1-ko]{4.5.1})은 전체에서 정의되고 우세한 열린 몰입
  사상이다.

\oldpage[II]{84}

  \item[\rm{b')}] 표준 사상
  $G(\varepsilon)\to\operatorname{Proj}(S)$은 전체에서 정의되고,
  $X$의 바탕공간에서 $\operatorname{Proj}(S)$의 한 부분공간 위로 가는
  위상동형이다.
  \item[\rm{c)}] 모든 준연접 $\mathcal{O}_X$-가군 $\mathcal{F}$에
  대하여, $\mathcal{F}_n$을 $X$ 위의 $\mathcal{F}(n)$의 단면들로
  생성되는 $\mathcal{F}(n)$의 부분-$\mathcal{O}_X$-가군이라 하면,
  $\mathcal{F}$는 정수 $n>0$에 대한 부분-$\mathcal{O}_X$-가군
  $\mathcal{F}_n(-n)$들의 합이다.
  \item[\rm{c')}] 성질 c)는 $\mathcal{O}_X$의 모든 준연접 아이디얼층에
  대하여 성립한다.
\end{enumerate}
더 나아가 $(f_\alpha)$가 $S_+$의 동차 원소들의 족이고
$X_{f_\alpha}$들이 아핀이면, 표준 사상
$X\to\operatorname{Proj}(S)$의 $\bigcup_\alpha X_{f_\alpha}$에 대한
제한은 $\bigcup_\alpha X_{f_\alpha}$에서
$\bigcup_\alpha(\operatorname{Proj}(S))_{f_\alpha}$ 위로 가는
동형사상이다.
\end{theorem}

\begin{proof}
b)가 b')을 함의함은 명백하고, b')은 공식
(\hyperref[II.3.7.3.1-ko]{3.7.3.1})에 따라 a)를 함의한다
($\varepsilon^\flat$가 항등사상임을 고려한다). 조건 a)는 a')을
함의한다. 실제로 모든 $x\in X$에는 $\mathcal{L}|U$가
$\mathcal{O}_X|U$와 동형인 아핀 근방 $U$가 있다~; 만일
$f\in\Gamma(X,\mathcal{L}^{\otimes n})$가
$x\in X_f\subset U$를 만족하면, $X_f$는
$(f|U)(x')\neq0$인 $x'\in U$들의 집합이기도 하며, 따라서 아핀
열린집합이다 (\hyperref[I.1.3.6-ko]{I, 1.3.6}). a')이 b)를
함의함을 증명하려면 정리의 마지막 주장을 증명하고, 더 나아가
$X=\bigcup_\alpha X_{f_\alpha}$이면
(\hyperref[II.3.8.2-ko]{3.8.2})의 조건 (iv)가 만족됨을 확인하는 것으로
충분하다. 이 마지막 사실은 (\hyperref[I.9.3.1-ko]{I, 9.3.1, (i)})에서
곧바로 따른다. 한편 (\hyperref[II.4.5.2-ko]{4.5.2})의 마지막 주장에
대해서는, $X_{f_\alpha}$가 사상
$G(\varepsilon)\to\operatorname{Proj}(S)$에 의한
$(\operatorname{Proj}(S))_{f_\alpha}$의 역상이므로
(\hyperref[I.9.3.2-ko]{I, 9.3.2})를 적용하면 충분하다. 따라서 a),
a'), b), b')은 동치이다.

a')이 c)를 함의함을 보이기 위해, $X_f$가 아핀이고($f\in S_k$),
$\mathcal{F}|X_f$가 $X_f$ 위의 단면들로 생성됨을 유의하자
(\hyperref[I.1.3.9-ko]{I, 1.3.9})~; 다른 한편
(\hyperref[I.9.3.1-ko]{I, 9.3.1, (ii)}) 그러한 단면 $s$는
$(t|X_f)\otimes(f|X_f)^{-m}$의 꼴이며, 여기서
$t\in\Gamma(X,\mathcal{F}(km))$이다~; 정의에 따라
$t$는 $\mathcal{F}_{km}$의 단면이기도 하므로, $s$는 실제로 $X_f$
위의 $\mathcal{F}_{km}(-km)$의 단면이고, 이로써 c)가 증명된다. c)가
c')을 함의함은 명백하며, 이제 c')이 a)를 함의함을 보이는 일만 남는다.
$U$를 $x\in X$의 열린 근방이라 하고, $\mathcal{J}$를 바탕공간이
$X-U$인 $X$의 닫힌 부분준스킴을 정의하는 $\mathcal{O}_X$의 준연접
아이디얼층이라 하자 (\hyperref[I.5.2.1-ko]{I, 5.2.1}). 가정 c')에 따라
어떤 정수 $n>0$과 $X$ 위의 $\mathcal{J}(n)$의 단면 $f$가 존재하여
$f(x)\neq0$이다. 그런데 명백히 $f\in S_n$이고
$x\in X_f\subset U$이므로, a)가 증명된다.
\end{proof}

$X$가 그 바탕공간이 뇌터인 준스킴이면,
(\hyperref[II.4.5.2-ko]{4.5.2})의 동치인 조건들은 $X$가 스킴임을
함의한다. 실제로 (\hyperref[II.4.5.2-ko]{4.5.2, b)})에 따라 $X$는
스킴 $S=\operatorname{Proj}(A)$의 한 부분준스킴과 동형이다.

\begin{definition}[4.5.3]
\label{II.4.5.3-ko}
$X$가 준콤팩트 스킴이고 (\hyperref[II.4.5.2-ko]{4.5.2})의 동치인
조건들이 성립할 때, 가역 $\mathcal{O}_X$-가군 $\mathcal{L}$을
풍부하다고 한다.
\end{definition}

판정법 (\hyperref[II.4.5.2-ko]{4.5.2, a)})에서 명백히 다음이 따른다.
$\mathcal{L}$이 풍부한 $\mathcal{O}_X$-가군이면, $X$의 모든 열린집합
$U$에 대하여 $\mathcal{L}|U$는 풍부한
$(\mathcal{O}_X|U)$-가군이다.

(\hyperref[II.4.5.2-ko]{4.5.2})의 증명에 따라, 아핀인 $X_f$들은 이미
$X$의 위상의 기저를 이룬다. 더 나아가 다음이 성립한다~:

\begin{corollary}[4.5.4]
\label{II.4.5.4-ko}
$\mathcal{L}$을 풍부한 $\mathcal{O}_X$-가군이라 하자. $X$의 임의의 유한
부분공간 $Z$와 $Z$의 임의의 근방 $U$에 대하여, $X_f$가 $U$에 포함되는
$Z$의 아핀 근방이 되도록 하는 정수 $n$과
$f\in\Gamma(X,\mathcal{L}^{\otimes n})$이 존재한다.
\end{corollary}

\begin{proof}
\oldpage[II]{85}

(\hyperref[II.4.5.2-ko]{4.5.2, b)})에 따라, $\operatorname{Proj}(S)$의
임의의 유한 부분집합 $Z'$과 $Z'$의 임의의 열린 근방 $U$에 대하여
$Z'\subset(\operatorname{Proj}(S))_f\subset U$를 만족하는 동차 원소
$f\in S_+$가 존재함을 증명하면 충분하다.
(\hyperref[II.2.4.1-ko]{2.4.1}). 그런데 정의에 따라, $U$의 여집합인
닫힌집합 $Y$는 $V_+(\mathfrak{I})$ 꼴이다. 여기서 $\mathfrak{I}$는
$S_+$를 포함하지 않는 $S$의 등급 아이디얼이다
(\hyperref[II.2.3.2-ko]{2.3.2})~; 다른 한편, $Z'$의 점들은 정의상
$\mathfrak{I}$를 포함하지 않는 $S_+$의 등급 소아이디얼
$\mathfrak{p}_i$이다 (\hyperref[II.2.3.1-ko]{2.3.1}). 따라서 어느
$\mathfrak{p}_i$에도 속하지 않는 원소 $f\in\mathfrak{I}$가 존재한다
(Bourbaki, \emph{Alg. comm.}, chap.~II, \S~1,
n\textsuperscript{o}~1, prop.~2). $\mathfrak{p}_i$들이 등급
아이디얼이므로, \emph{같은 곳}에서 한 논증은 $f$를 동차라고까지
가정할 수 있음을 보여 준다~; 그러면 이 원소가 요구를 충족한다.
\end{proof}

\begin{proposition}[4.5.5]
\label{II.4.5.5-ko}
$X$가 준콤팩트 스킴이거나 그 기저공간이 뇌터인 준스킴이라고 가정하자.
(\hyperref[II.4.5.2-ko]{4.5.2})의 조건 a)에서 c')까지는 다음 조건들과도
동치이다~:
\begin{enumerate}
  \item[\rm{d)}] 임의의 유한형 준연접 $\mathcal{O}_X$-가군
  $\mathcal{F}$에 대하여, 모든 $n\geq n_0$에 대해 $\mathcal{F}(n)$이
  $X$ 위의 단면들로 생성되도록 하는 정수 $n_0$가 존재한다.
  \item[\rm{d')}] 임의의 유한형 준연접 $\mathcal{O}_X$-가군
  $\mathcal{F}$에 대하여, $\mathcal{F}$가 $\mathcal{O}_X$-가군
  $\mathcal{L}^{\otimes(-n)}\otimes\mathcal{O}_X^k$의 몫과 동형이
  되도록 하는 두 정수 $n>0$, $k>0$가 존재한다.
  \item[\rm{d'')}] 성질 d')는 $\mathcal{O}_X$의 임의의 유한형 준연접
  아이디얼층에 대하여 성립한다.
\end{enumerate}
\end{proposition}

\begin{proof}
$X$가 준콤팩트이므로, 유한형 준연접 $\mathcal{O}_X$-가군
$\mathcal{F}$에 대하여 (그 자체도 유한형인) $\mathcal{F}(n)$이 $X$
위의 단면들로 생성되면, $\mathcal{F}(n)$은 그 단면들 가운데 유한 개로
생성된다 (\hyperref[0.5.2.3-ko]{0, 5.2.3}). 따라서 d)는 d')를
함의하고, d')가 d'')를 함의함은 명백하다. 임의의 준연접
$\mathcal{O}_X$-가군 $\mathcal{G}$는 그 유한형 부분-$\mathcal{O}_X$-가군들의
귀납적 극한이므로 (\hyperref[I.9.4.9-ko]{I, 9.4.9}),
(\hyperref[II.4.5.2-ko]{4.5.2})의 조건 c')를 확인하려면 유한형인
$\mathcal{O}_X$의 준연접 아이디얼층에 대해서만 확인하면 충분하다.
따라서 d'')는 c')를 함의한다. 이제 $\mathcal{L}$이 풍부하면 성질 d)가
성립함을 보이는 일만 남았다. $X$의 유한 열린 덮개
$X_{f_i}$ ($f_i\in S_{n_i}$)를 생각하자. 이들은 아핀이라고 가정할 수
있다~; $f_i$들을 적당한 거듭제곱으로 바꾸면(이는 $X_{f_i}$들을
바꾸지 않는다), 모든 $n_i$가 하나의 같은 정수 $m$과 같다고 가정할
수 있다. 가정에 따라 유한형인 층 $\mathcal{F}|X_{f_i}$는
$X_{f_i}$ 위의 유한 개 단면 $h_{ij}$로 생성된다
(\hyperref[I.1.3.13-ko]{I, 1.3.13})~; 따라서 모든 쌍 $(i,j)$에 대하여
단면 $h_{ij}\otimes f_i^{\otimes k_0}$가 $X$ 위의
$\mathcal{F}(k_0m)$의 단면으로 연장되도록 하는 정수 $k_0$가 존재한다
(\hyperref[I.9.3.1-ko]{I, 9.3.1}). \emph{더 나아가}, 모든
$k\geq k_0$에 대하여 $h_{ij}\otimes f_i^{\otimes k}$는 $X$ 위의
$\mathcal{F}(km)$의 단면으로 연장되고, 따라서 이러한 $k$에 대하여
$\mathcal{F}(km)$은 $X$ 위의 단면들로 생성된다. $0<p<m$인 임의의
$p$에 대하여 $\mathcal{F}(p)$ 역시 유한형이므로,
$k\geq k_p$일 때 $\mathcal{F}(p)(km)=\mathcal{F}(p+km)$이 $X$ 위의
단면들로 생성되도록 하는 정수 $k_p$가 존재한다. $n_0$를 모든
$k_pm$보다 크게 택하면, 모든 $n\geq n_0$에 대하여 $\mathcal{F}(n)$은
$X$ 위의 단면들로 생성된다는 결론을 얻는다. 실제로 그러한 $n$은
$k\geq k_p$ 및 $0\leq p<m$을 만족하도록 $n=km+p$로 쓸 수 있다.
\end{proof}

\begin{proposition}[4.5.6]
\label{II.4.5.6-ko}
$X$를 준콤팩트 스킴, $\mathcal{L}$을 가역 $\mathcal{O}_X$-가군이라
하자.
\begin{enumerate}
  \item[\rm{(i)}] $n$을 $>0$인 정수라 하자. $\mathcal{L}$이
  풍부할 필요충분조건은 $\mathcal{L}^{\otimes n}$이 풍부한 것이다.
  \item[\rm{(ii)}] $\mathcal{L}'$을 다음 성질을 만족하는 가역
  $\mathcal{O}_X$-가군이라 하자. 모든 $x\in X$에 대하여 어떤 정수
  $n>0$과 $X$ 위의 $\mathcal{L}'^{\otimes n}$의 단면 $s'$이 존재하여

\oldpage[II]{86}

  $s'(x)\neq0$이라고 하자. 그러면 $\mathcal{L}$이 풍부하면
  $\mathcal{L}\otimes\mathcal{L}'$도 풍부하다.
\end{enumerate}
\end{proposition}

\begin{proof}
성질 (i)는 $X_{f^{\otimes n}}=X_f$이므로
(\hyperref[II.4.5.2-ko]{4.5.2})의 판정법 a)의 명백한 귀결이다. 한편
$\mathcal{L}$이 풍부하면, 모든 $x\in X$와 $x$의 모든 근방 $U$에 대하여
어떤 $m>0$과 $f\in\Gamma(X,\mathcal{L}^{\otimes m})$이 존재하여
$x\in X_f\subset U$이다 (\hyperref[II.4.5.2-ko]{4.5.2, a)})~;
$f'\in\Gamma(X,\mathcal{L}'^{\otimes n})$이고 $f'(x)\neq0$이면,
다음과 같이 놓은 $s$에 대하여 $s(x)\neq0$이다.
$s=f^{\otimes n}\otimes f'^{\otimes m}\in
\Gamma(X,(\mathcal{L}\otimes\mathcal{L}')^{\otimes mn})$ 따라서
$x\in X_s\subset X_f\subset U$이고, 이는
$\mathcal{L}\otimes\mathcal{L}'$이 풍부함을 증명한다
(\hyperref[II.4.5.2-ko]{4.5.2, a)}).
\end{proof}

\begin{corollary}[4.5.7]
\label{II.4.5.7-ko}
풍부한 두 $\mathcal{O}_X$-가군의 텐서곱은 풍부하다.
\end{corollary}

\begin{corollary}[4.5.8]
\label{II.4.5.8-ko}
$\mathcal{L}$을 풍부한 $\mathcal{O}_X$-가군,
$\mathcal{L}'$을 가역 $\mathcal{O}_X$-가군이라 하자. 그러면 어떤 정수
$n_0>0$이 존재하여, $n\geq n_0$이면
$\mathcal{L}^{\otimes n}\otimes\mathcal{L}'$은 풍부하고 $X$ 위의
단면들에 의해 생성된다.
\end{corollary}

\begin{proof}
실제로 (\hyperref[II.4.5.5-ko]{4.5.5})에 따라 어떤 정수 $m_0$이
존재하여, $m\geq m_0$이면
$\mathcal{L}^{\otimes m}\otimes\mathcal{L}'$은 $X$ 위의 단면들에
의해 생성된다~; 따라서 (\hyperref[II.4.5.6-ko]{4.5.6})에 따라
$n_0=m_0+1$로 둘 수 있다.
\end{proof}

\begin{remark}[4.5.9]
\label{II.4.5.9-ko}
$P=\mathrm{H}^1(X,\mathcal{O}_X^*)$를 가역
$\mathcal{O}_X$-가군들의 동치류의 군이라 하고
(\hyperref[0.5.4.7-ko]{0, 5.4.7}), $P^+$를 풍부한 층들의 동치류로
이루어진 $P$의 부분집합이라 하자. $P^+$가 공집합이 아니라고 가정하자.
그러면 (\hyperref[II.4.5.7-ko]{4.5.7})과
(\hyperref[II.4.5.8-ko]{4.5.8})에 따라
\[
  P^++P^+\subset P^+
  \qquad\text{그리고}\qquad
  P^+-P^+=P
\]
를 얻는다. 다시 말해 $P^+\cup\{0\}$은 $P$의 군 구조와 양립하는
준순서 구조에서 $P$의 양의 원소들의 집합이며, 이 준순서는
(\hyperref[II.4.5.8-ko]{4.5.8})에 따라 아르키메데스적이기까지 하다.
이 때문에 때로는 풍부한 층 대신 «~양의 층~»이라 하고, 그런 층의
역원에 대하여 «~음의 층~»이라 한다(우리는 이 용어법을 따르지 않는다).
\end{remark}

\begin{proposition}[4.5.10]
\label{II.4.5.10-ko}
$Y$를 아핀 스킴, $q:X\to Y$를 분리 준콤팩트 사상,
$\mathcal{L}$을 가역 $\mathcal{O}_X$-가군이라 하자.
\begin{enumerate}
  \item[\rm{(i)}] $\mathcal{L}$이 $q$에 관해 매우 풍부하면
  $\mathcal{L}$은 풍부하다.
  \item[\rm{(ii)}] 더 나아가 사상 $q$가 유한형이라고 가정하자.
  $\mathcal{L}$이 풍부할 필요충분조건은 다음 동치인 성질들 가운데
  하나를 갖는 것이다~:
  \begin{enumerate}
    \item[\rm{e)}] 어떤 $n_0>0$가 존재하여, 모든 정수 $n\geq n_0$에
    대하여 $\mathcal{L}^{\otimes n}$은 $q$에 관해 매우 풍부하다.
    \item[\rm{e')}] 어떤 $n>0$가 존재하여
    $\mathcal{L}^{\otimes n}$은 $q$에 관해 매우 풍부하다.
  \end{enumerate}
\end{enumerate}
\end{proposition}

\begin{proof}
첫째 주장은 매우 풍부한 $\mathcal{O}_X$-가군의 정의
(\hyperref[II.4.4.2-ko]{4.4.2})에서 따른다~: $A$가 $Y$의 환이면,
$A$-가군 $E$와 전사 준동형
\[
  \psi:q^*((\mathbf{S}(E))^{\sim})\longrightarrow
  \bigoplus_{n\geq0}\mathcal{L}^{\otimes n}
\]
가 존재하여 $i=r_{\mathcal{L},\psi}$는 전체에서 정의되는 몰입 사상
$X\to P=\mathbf{P}(\widetilde{E})$이고
$\mathcal{L}=i^*(\mathcal{O}_P(1))$이다~; $(\mathbf{S}(E))_+$에 속하는
동차 원소 $f$에 대한 $D_+(f)$들이 $P$의 위상의 기저를 이루고,
(\hyperref[II.3.7.3.1-ko]{3.7.3.1})에 따라
$i^{-1}(D_+(f))=X_{\psi^\flat(f)}$이므로,
(\hyperref[II.4.5.2-ko]{4.5.2})의 조건 a)가 성립함을 알 수 있다.
따라서 $\mathcal{L}$은 풍부하다.

이제 $q$가 유한형이고 $\mathcal{L}$이 풍부하면 조건 e)를 만족함을
증명하자. 먼저 (\hyperref[II.4.5.2-ko]{4.5.2})의 판정법 b)와
(\hyperref[II.4.4.1-ko]{4.4.1, (i)})에서 어떤

\oldpage[II]{87}

정수 $k_0>0$가
존재하여 $\mathcal{L}^{\otimes k_0}$이 $q$에 관해 매우 풍부함이
따른다. 한편 (\hyperref[II.4.5.5-ko]{4.5.5})에 따라 어떤 정수
$m_0$가 존재하여, $m\geq m_0$이면 $\mathcal{L}^{\otimes m}$은
$X$ 위의 단면들로 생성된다. $n_0=k_0+m_0$라 놓자~; $n\geq n_0$이면
$m\geq m_0$인 $m$을 써서 $n=k_0+m$으로 나타낼 수 있고, 따라서
$\mathcal{L}^{\otimes n}
=\mathcal{L}^{\otimes k_0}\otimes\mathcal{L}^{\otimes m}$이다.
$\mathcal{L}^{\otimes m}$은 $X$ 위의 단면들로 생성되므로,
(\hyperref[II.4.4.8-ko]{4.4.8})의 판정법과
(\hyperref[II.3.4.7-ko]{3.4.7})에서 $\mathcal{L}^{\otimes n}$이
$q$에 관해 매우 풍부함이 따른다. 마지막으로 e)가 e')를 함의함은
명백하고, e')는 (i)와 (\hyperref[II.4.5.6-ko]{4.5.6, (i)})에 따라
$\mathcal{L}$이 풍부함을 함의한다.

\begin{env}[4.5.10.1]
\label{II.4.5.10.1-ko}
\emph{보조정리 (\hyperref[II.4.4.10.1-ko]{4.4.10.1})의 증명.}
$\mathcal{E}(n)=\mathcal{E}\otimes\mathcal{K}^{\otimes n}$이라 놓자~;
$g$가 분리 사상이므로 (\hyperref[II.4.4.2-ko]{4.4.2}),
(\hyperref[II.3.4.8-ko]{3.4.8})의 논법을 적용할 수 있고, 문제는 충분히
큰 $n$에 대하여 표준 준동형
$g^*(g_*(\mathcal{E}(n)))\to\mathcal{E}(n)$이 전사임을 보이는 것으로
환원된다. 더 나아가 $Z$가 준콤팩트이므로
(\hyperref[II.3.4.6-ko]{3.4.6})의 논법에 따라 $Z$가 아핀인 경우에
명제를 증명하면 충분하다. 그런데 이 경우 $\mathcal{K}$는
(\hyperref[II.4.5.10-ko]{4.5.10, (i)})에 따라 풍부하고, 결론은
(\hyperref[II.4.5.5-ko]{4.5.5, d')})에서 따른다.
\end{env}
\end{proof}

\begin{corollary}[4.5.11]
\label{II.4.5.11-ko}
$Y$를 아핀 스킴, $q:X\to Y$를 유한형 분리 사상,
$\mathcal{L}$을 풍부한 $\mathcal{O}_X$-가군,
$\mathcal{L}'$을 가역 $\mathcal{O}_X$-가군이라 하자. 어떤 정수
$n_0$가 존재하여, $n\geq n_0$이면
$\mathcal{L}^{\otimes n}\otimes\mathcal{L}'$은 $q$에 관해 매우 풍부하다.
\end{corollary}

\begin{proof}
실제로 어떤 $m_0$가 존재하여, $m\geq m_0$이면
$\mathcal{L}^{\otimes m}\otimes\mathcal{L}'$은 $X$ 위의 단면들로
생성된다 (\hyperref[II.4.5.8-ko]{4.5.8})~; 한편 어떤 $k_0$가 존재하여,
$k\geq k_0$이면 $\mathcal{L}^{\otimes k}$은 $q$에 관해 매우 풍부하다.
따라서 $k\geq k_0$이면
$\mathcal{L}^{\otimes(k+m_0)}\otimes\mathcal{L}'$은 매우 풍부하다
((\hyperref[II.4.4.8-ko]{4.4.8}) 및
(\hyperref[II.3.4.7-ko]{3.4.7})).
\end{proof}

\begin{remark}[4.5.12]
\label{II.4.5.12-ko}
어떤 $n$에 대하여 $\mathcal{L}^{\otimes n}$이 ($q$에 관해) 매우
풍부한 $\mathcal{O}_X$-가군 $\mathcal{L}$이면
$\mathcal{L}^{\otimes(n+1)}$도 그렇다는 결론이 따르는지는 알려져 있지
않다.
\end{remark}

\begin{proposition}[4.5.13]
\label{II.4.5.13-ko}
$X$를 준콤팩트 준스킴, $Z$를 $\mathcal{O}_X$의 준연접 멱영 아이디얼층
$\mathcal{J}$로 정의되는 $X$의 닫힌 부분준스킴, $j$를 표준 포함 사상
$Z\to X$라 하자. 가역 $\mathcal{O}_X$-가군 $\mathcal{L}$이 풍부할
필요충분조건은 $\mathcal{L}'=j^*(\mathcal{L})$이 풍부한
$\mathcal{O}_Z$-가군인 것이다.
\end{proposition}

\begin{proof}
이 조건은 필요하다. 실제로 $X$ 위의 $\mathcal{L}^{\otimes n}$의 임의의
단면 $f$에 대하여, $f'$을 그 표준상 $f\otimes1$, 곧 $X$와 동일한
기저공간을 갖는 $Z$ 위의
$\mathcal{L}'^{\otimes n}=\mathcal{L}^{\otimes n}
\otimes_{\mathcal{O}_X}(\mathcal{O}_X/\mathcal{J})$의 단면이라 하자.
그러면 $X_f=Z_{f'}$임은 명백하므로
(\hyperref[II.4.5.2-ko]{4.5.2})의 판정법 a)에 따라
$\mathcal{L}'$은 풍부하다.

이 조건이 충분함을 보이기 위해, 먼저 $\mathcal{J}^2=0$인 경우로
환원할 수 있음을 주목하자. 실제로 준스킴들의 유한 열
$X_k=(X,\mathcal{O}_X/\mathcal{J}^{k+1})$을 생각하면, 각각은 그 다음
준스킴의 닫힌 부분준스킴이며 제곱이 영인 아이디얼층으로 정의된다.
또한 가정에 의해 $X_{\mathrm{red}}$가 스킴이므로 $X$도 스킴이다
((\hyperref[II.4.5.3-ko]{4.5.3}) 및
(\hyperref[I.5.5.1-ko]{I, 5.5.1})).
(\hyperref[II.4.5.2-ko]{4.5.2})의 판정법 a)에 따라 다음 보조정리를
증명하면 충분하다.

\begin{lemma}[4.5.13.1]
\label{II.4.5.13.1-ko}
(\hyperref[II.4.5.13-ko]{4.5.13})의 가정하에서, 더 나아가
$\mathcal{J}$가 제곱영이라고 가정하자~; $\mathcal{L}$을 가역
$\mathcal{O}_X$-가군이라 하고, $g$를 $Z_g$가 아핀이 되도록 하는
$Z$ 위의 $\mathcal{L}'^{\otimes n}$의 단면이라 하자. 그러면 어떤
정수 $m>0$가 존재하여 $g^{\otimes m}$은 $X$ 위의
$\mathcal{L}^{\otimes nm}$의 한 단면 $f$의 표준상이 된다.
\end{lemma}

$\mathcal{O}_X$-가군의 완전열
\[
  0\longrightarrow\mathcal{J}(n)\longrightarrow
  \mathcal{O}_X(n)=\mathcal{L}^{\otimes n}\longrightarrow
  \mathcal{O}_Z(n)=\mathcal{L}'^{\otimes n}\longrightarrow0
\]

\oldpage[II]{88}

이 있다. 실제로 $\mathcal{F}(n)$은 $\mathcal{F}$에 관한 완전
함자이다~; 따라서 코호몰로지 완전열
\[
  0\longrightarrow\Gamma(X,\mathcal{J}(n))\longrightarrow
  \Gamma(X,\mathcal{L}^{\otimes n})\longrightarrow
  \Gamma(X,\mathcal{L}'^{\otimes n})\xrightarrow{\partial}
  \mathrm{H}^1(X,\mathcal{J}(n))
\]
을 얻으며, 이 열은 특히 $g$에 원소
$\partial g\in\mathrm{H}^1(X,\mathcal{J}(n))$을 대응시킨다.

$\mathcal{J}^2=0$이므로 $\mathcal{J}$를 준연접
$\mathcal{O}_Z$-가군으로 볼 수 있고, 모든 $k$에 대하여
$\mathcal{L}'^{\otimes k}\otimes_{\mathcal{O}_Z}\mathcal{J}(n)
=\mathcal{J}(n+k)$임에 유의하자~; 따라서 모든 단면
$s\in\Gamma(X,\mathcal{L}'^{\otimes k})$에 대하여 $s$와의 텐서곱은
$\mathcal{O}_Z$-가군의 준동형
$\mathcal{J}(n)\to\mathcal{J}(n+k)$이고, 그 결과 코호몰로지군의
준동형
$\mathrm{H}^i(X,\mathcal{J}(n))\to
\mathrm{H}^i(X,\mathcal{J}(n+k))$을 준다.

이제 충분히 큰 어떤 $m>0$에 대하여
\begin{equation}
\label{II.4.5.13.2-ko}
  g^{\otimes m}\otimes\partial g=0
\tag{4.5.13.2}
\end{equation}
임을 보이겠다. 실제로 $Z_g$는 $Z$의 아핀 열린집합이고,
$\mathcal{J}(n)$을 $\mathcal{O}_Z$-가군으로 볼 때
$\mathrm{H}^1(Z_g,\mathcal{J}(n))=0$이다
(\hyperref[I.5.1.9.2-ko]{I, 5.1.9.2}). 특히 $g'=g|Z_g$로 놓고 그 상을
사상
$\partial:\Gamma(Z_g,\mathcal{L}'^{\otimes n})\to
\mathrm{H}^1(Z_g,\mathcal{J}(n))$에 의해 생각하면 $\partial g'=0$이다.
이 관계를 명시하기 위하여, 차수 $1$에서 아벨군층의 코호몰로지는
그 체흐 코호몰로지와 같음을 관찰하자 (G, II, 5.9)~; 따라서
$\partial g$를 구성하려면 $X$의 충분히 미세한 열린 덮개
$(U_\alpha)$를 생각해야 하는데, 이를 유한하고 아핀 열린집합들로
이루어졌다고 가정할 수 있다. 각 $\alpha$에 대하여 표준상이
$\Gamma(U_\alpha,\mathcal{L}'^{\otimes n})$에서 $g|U_\alpha$인 단면
$g_\alpha\in\Gamma(U_\alpha,\mathcal{L}^{\otimes n})$을 택하고,
공순환 $(g_{\alpha|\beta}-g_{\beta|\alpha})$의 코호몰로지류를
생각한다. 여기서 $g_{\alpha|\beta}$는 $g_\alpha$의
$U_\alpha\cap U_\beta$로의 제한이고, 이 공순환은
$\mathcal{J}(n)$에 값을 갖는다. 또한 필요하다면 $(U_\alpha)$를 더
미세한 덮개로 바꾸어, $\partial g'$도 $U_\alpha\cap Z_g$들로 이루어진
덮개와 제한 $g_\alpha|(U_\alpha\cap Z_g)$을 사용하여 같은 방식으로
계산된다고 가정할 수 있다~; 그러면 관계 $\partial g'=0$은 모든
$\alpha$에 대하여 단면
$h_\alpha\in\Gamma(U_\alpha\cap Z_g,\mathcal{J}(n))$이 존재하여
\[
  (g_{\alpha|\beta}-g_{\beta|\alpha})
  |(U_\alpha\cap U_\beta\cap Z_g)
  =h_{\alpha|\beta}-h_{\beta|\alpha}
\]
임을 뜻한다. 여기서 $h_{\alpha|\beta}$는 $h_\alpha$의
$U_\alpha\cap U_\beta\cap Z_g$로의 제한을 나타낸다 (G, II, 5.11).
이제 어떤 정수 $m>0$가 존재하여 모든 $\alpha$에 대해
$g^{\otimes m}\otimes h_\alpha$는 어떤 단면
$t_\alpha\in\Gamma(U_\alpha,\mathcal{J}(n+nm))$의
$U_\alpha\cap Z_g$로의 제한이 된다
(\hyperref[I.9.3.1-ko]{I, 9.3.1})~; 따라서 모든 지표의 쌍에 대하여
\[
  g^{\otimes m}\otimes(g_{\alpha|\beta}-g_{\beta|\alpha})
  =t_{\alpha|\beta}-t_{\beta|\alpha}
\]
이고, 이는 (\hyperref[II.4.5.13.2-ko]{4.5.13.2})를 증명한다.

다른 한편
$s\in\Gamma(X,\mathcal{O}_Z(p))$,
$t\in\Gamma(X,\mathcal{O}_Z(q))$이면, 군
$\mathrm{H}^1(X,\mathcal{J}(p+q))$에서
\begin{equation}
\label{II.4.5.13.3-ko}
  \partial(s\otimes t)=(\partial s)\otimes t+s\otimes(\partial t).
\tag{4.5.13.3}
\end{equation}
이다.

실제로 양변을 계산하기 위하여 다시 $X$의 열린 덮개
$(U_\alpha)$를 생각할 수 있고, 각 $\alpha$에 대하여 표준상이
$\Gamma(U_\alpha,\mathcal{O}_Z(p))$ (각각
$\Gamma(U_\alpha,\mathcal{O}_Z(q))$)에서 $s|U_\alpha$ (각각
$t|U_\alpha$)인 단면
$s_\alpha\in\Gamma(U_\alpha,\mathcal{O}_X(p))$ (각각
$t_\alpha\in\Gamma(U_\alpha,\mathcal{O}_X(q))$)를 생각할 수 있다~;
관계 (\hyperref[II.4.5.13.3-ko]{4.5.13.3})은 위와 같은 표기 아래
관계
\[
  (s_{\alpha|\beta}\otimes t_{\alpha|\beta})
  -(s_{\beta|\alpha}\otimes t_{\beta|\alpha})
  =(s_{\alpha|\beta}-s_{\beta|\alpha})\otimes t_{\alpha|\beta}
  +s_{\beta|\alpha}\otimes(t_{\alpha|\beta}-t_{\beta|\alpha})
\]
에서 따른다. 따라서 $k$에 관한 귀납법으로
\begin{equation}
\label{II.4.5.13.4-ko}
  \partial(g^{\otimes k})=(kg^{\otimes(k-1)})\otimes(\partial g)
\tag{4.5.13.4}
\end{equation}
를 얻는다.

\oldpage[II]{89}

(\hyperref[II.4.5.13.2-ko]{4.5.13.2})와
(\hyperref[II.4.5.13.4-ko]{4.5.13.4})로부터
$\partial(g^{\otimes(m+1)})=0$이라고 결론짓는다~; 따라서
$g^{\otimes(m+1)}$은 $X$ 위의 $\mathcal{L}^{\otimes n(m+1)}$의
한 단면 $f$의 표준상이고, 이로써
(\hyperref[II.4.5.13-ko]{4.5.13})의 증명이 끝난다.
\end{proof}

\begin{corollary}[4.5.14]
\label{II.4.5.14-ko}
$X$를 뇌터 스킴, $j$를 표준 포함 사상
$X_{\mathrm{red}}\to X$라 하자. 가역 $\mathcal{O}_X$-가군
$\mathcal{L}$이 풍부할 필요충분조건은 $j^*(\mathcal{L})$이 풍부한
$\mathcal{O}_{X_{\mathrm{red}}}$-가군인 것이다.
\end{corollary}

이는 (\hyperref[I.6.1.6-ko]{I, 6.1.6})에서 따른다.

\subsection{상대적으로 풍부한 층}
\label{II.4.6-ko}

\begin{definition}[4.6.1]
\label{II.4.6.1-ko}
$f:X\to Y$를 준콤팩트 사상, $\mathcal{L}$을 가역
$\mathcal{O}_X$-가군이라 하자. $Y$에 아핀 열린집합들로 이루어진 덮개
$(U_\alpha)$가 존재하여, $X_\alpha=f^{-1}(U_\alpha)$라 놓으면 모든
$\alpha$에 대하여 $\mathcal{L}|X_\alpha$가 풍부한
$\mathcal{O}_{X_\alpha}$-가군일 때, $\mathcal{L}$을 $f$에 관해
풍부하다, 또는 $Y$에 관해 풍부하다, 또는 $f$-풍부하다, 또는
$Y$-풍부하다라고 한다(또는 (\hyperref[II.4.5.3-ko]{4.5.3})에서 정의한
개념과 혼동할 염려가 없으면 단순히 풍부하다고도 한다).
\end{definition}

$f$-풍부한 $\mathcal{O}_X$-가군이 존재하면 $f$는 반드시 분리
사상이다((\hyperref[II.4.5.3-ko]{4.5.3}) 및
(\hyperref[I.5.5.5-ko]{I, 5.5.5})).

\begin{proposition}[4.6.2]
\label{II.4.6.2-ko}
$f:X\to Y$를 준콤팩트 사상, $\mathcal{L}$을 가역
$\mathcal{O}_X$-가군이라 하자. $\mathcal{L}$이 $f$에 관해 매우
풍부하면, $f$에 관해 풍부하다.
\end{proposition}

이는 매우 풍부한 층의 개념이 $Y$ 위에서 국소적이라는 사실
(\hyperref[II.4.4.5-ko]{4.4.5}), 정의
(\hyperref[II.4.6.1-ko]{4.6.1}), 그리고 판정법
(\hyperref[II.4.5.10-ko]{4.5.10, (i)})에서 따른다.

\begin{proposition}[4.6.3]
\label{II.4.6.3-ko}
$f:X\to Y$를 준콤팩트 사상, $\mathcal{L}$을 가역
$\mathcal{O}_X$-가군이라 하고, $\mathcal{S}$를 등급
$\mathcal{O}_Y$-대수
$\bigoplus_{n\geq0}f_*(\mathcal{L}^{\otimes n})$라 하자. 다음 조건들은
동치이다~:
\begin{enumerate}
  \item[\rm{a)}] $\mathcal{L}$은 $f$-풍부하다.
  \item[\rm{b)}] $\mathcal{S}$는 준연접이고 표준 준동형
  $\sigma:f^*(\mathcal{S})\to
  \bigoplus_{n\geq0}\mathcal{L}^{\otimes n}$
  (\hyperref[0.4.4.3-ko]{0, 4.4.3})은 $Y$-사상
  $r_{\mathcal{L},\sigma}:G(\sigma)\to\operatorname{Proj}(\mathcal{S})=P$가
  모든 점에서 정의되고 우세한 열린 몰입 사상이 되도록 한다.
  \item[\rm{b')}] 사상 $f$는 분리이고, $Y$-사상
  $r_{\mathcal{L},\sigma}$는 모든 점에서 정의되며 $X$의 기저공간을
  $\operatorname{Proj}(\mathcal{S})$의 한 부분공간 위로 보내는
  위상동형사상이다.
\end{enumerate}
더 나아가 이 조건들이 성립할 때, 모든 $n\in\mathbf{Z}$에 대하여
(\hyperref[II.3.7.9.1-ko]{3.7.9.1})에서 정의한 표준 준동형
\begin{equation}
\label{II.4.6.3.1-ko}
  r_{\mathcal{L},\sigma}^*(\mathcal{O}_P(n))
  \longrightarrow\mathcal{L}^{\otimes n}
\tag{4.6.3.1}
\end{equation}
은 동형사상이다.

끝으로, 모든 준연접 $\mathcal{O}_X$-가군 $\mathcal{F}$에 대하여
$\mathcal{M}=\bigoplus_{n\geq0}f_*(\mathcal{F}\otimes
\mathcal{L}^{\otimes n})$라 놓으면,
(\hyperref[II.3.7.9.2-ko]{3.7.9.2})에서 정의한 표준 준동형
\begin{equation}
\label{II.4.6.3.2-ko}
  r_{\mathcal{L},\sigma}^*(\widetilde{\mathcal{M}})
  \longrightarrow\mathcal{F}
\tag{4.6.3.2}
\end{equation}
은 동형사상이다.
\end{proposition}

a)는 $f$가 분리 사상임을 함의한다고 이미 지적했으므로,
$\mathcal{S}$는 준연접이다
(\hyperref[I.9.2.2-ko]{I, 9.2.2, a)}). $r_{\mathcal{L},\sigma}$가
모든 점에서 정의된 몰입 사상이라는 것은 $Y$ 위에서 국소적인
성질이므로, a)가 b)를 함의함을 증명하기 위해서는 $Y$가 아핀이고
$\mathcal{L}$이 풍부한 경우만 생각하면 충분하다~;

\oldpage[II]{90}

이때 주장은 (\hyperref[II.4.5.2-ko]{4.5.2, b)})에서 따른다. b)가
b')을 함의함은 명백하다~; 끝으로 b')이 a)를 함의함을 증명하려면,
$Y$의 아핀 열린 덮개 $(U_\alpha)$를 택하고 각 층
$\mathcal{L}|f^{-1}(U_\alpha)$에 판정법
(\hyperref[II.4.5.2-ko]{4.5.2, b')})을 적용하면 충분하다.

마지막 두 주장에 대해서는 여기서 $\sigma^\flat$이 항등사상이라는
사실과 준동형
(\hyperref[II.3.7.9.1-ko]{3.7.9.1}) 및
(\hyperref[II.3.7.9.2-ko]{3.7.9.2})의 구체적 표현을 사용한다~;
그러면 (\hyperref[II.4.6.3.1-ko]{4.6.3.1})이 동형사상임이 곧바로
따른다. (\hyperref[II.4.6.3.2-ko]{4.6.3.2})에 대해서는 $Y$가
아핀이고 따라서 $\mathcal{L}$이 풍부한 경우로 환원할 수 있다~;
준동형 (\hyperref[II.4.6.3.2-ko]{4.6.3.2})이 단사임은 명백하고,
판정법 (\hyperref[II.4.5.2-ko]{4.5.2, c)})에 따라 전사이므로 결론을
얻는다.

\begin{corollary}[4.6.4]
\label{II.4.6.4-ko}
$(U_\alpha)$를 $Y$의 열린 덮개라 하자. $\mathcal{L}$이 $Y$에 관해
풍부할 필요충분조건은 모든 $\alpha$에 대하여
$\mathcal{L}|f^{-1}(U_\alpha)$가 $U_\alpha$에 관해 풍부한 것이다.
\end{corollary}

실제로 조건 b)는 $Y$ 위에서 국소적이다.

\begin{corollary}[4.6.5]
\label{II.4.6.5-ko}
$\mathcal{K}$를 가역 $\mathcal{O}_Y$-가군이라 하자.
$\mathcal{L}$이 $Y$ 위에서 풍부하기 위한 필요충분조건은
$\mathcal{L}\otimes f^*(\mathcal{K})$가 $Y$ 위에서 풍부한 것이다.
\end{corollary}

실제로 모든 $\alpha$에 대하여 $\mathcal{K}|U_\alpha$가
$\mathcal{O}_Y|U_\alpha$와 동형이 되도록 $U_\alpha$들을 잡으면, 이는
(\hyperref[II.4.6.4-ko]{4.6.4})의 명백한 귀결이다.

\begin{corollary}[4.6.6]
\label{II.4.6.6-ko}
$Y$가 아핀이라고 가정하자. $\mathcal{L}$이 $Y$ 위에서 풍부하기 위한
필요충분조건은 $\mathcal{L}$이 풍부한 것이다.
\end{corollary}

이는 정의 (\hyperref[II.4.6.1-ko]{4.6.1})와 판정법
(\hyperref[II.4.6.3-ko]{4.6.3, b)}) 및
(\hyperref[II.4.5.2-ko]{4.5.2, b)})의 직접적인 귀결이다. 실제로 여기서는
정의에 따라
$\operatorname{Proj}(\mathcal{S})
=\operatorname{Proj}(\Gamma(Y,\mathcal{S}))$이다.

\begin{corollary}[4.6.7]
\label{II.4.6.7-ko}
$f:X\to Y$를 준콤팩트 사상이라 하자. 준연접 $\mathcal{O}_Y$-가군
$\mathcal{E}$와 $Y$-사상 $g:X\to P=\mathbf{P}(\mathcal{E})$가 존재하고,
$g$가 $X$의 바탕공간에서 $P$의 한 부분공간으로의 위상동형이라고
가정하자. 그러면
$\mathcal{L}=g^*(\mathcal{O}_P(1))$은 $Y$ 위에서 풍부하다.
\end{corollary}

실제로 $Y$가 아핀이라고 가정해도 된다. 그러면 이 따름정리는 판정법
(\hyperref[II.4.5.2-ko]{4.5.2, a)}), 공식
(\hyperref[II.3.7.3.1-ko]{3.7.3.1}), 그리고
(\hyperref[II.4.2.3-ko]{4.2.3})에서 따른다.

\begin{proposition}[4.6.8]
\label{II.4.6.8-ko}
$X$를 준콤팩트 스킴 또는 바탕공간이 뇌터인 준스킴이라 하고,
$f:X\to Y$를 분리된 준콤팩트 사상이라 하자. 가역
$\mathcal{O}_X$-가군 $\mathcal{L}$이 $f$에 대해 풍부하기 위한
필요충분조건은 다음 동치 조건들 중 하나가 성립하는 것이다.
\begin{enumerate}
  \item[\rm{c)}] 모든 유한형 준연접 $\mathcal{O}_X$-가군
  $\mathcal{F}$에 대하여, 정수 $n_0>0$이 존재하여 모든 $n\geq n_0$에
  대해 표준 준동형
  $\sigma:f^*(f_*(\mathcal{F}\otimes\mathcal{L}^{\otimes n}))\to
  \mathcal{F}\otimes\mathcal{L}^{\otimes n}$이 전사이다.
  \item[\rm{c')}] $\mathcal{O}_X$의 모든 유한형 준연접 아이디얼층
  $\mathcal{J}$에 대하여, 정수 $n>0$이 존재하여 표준 준동형
  $\sigma:f^*(f_*(\mathcal{J}\otimes\mathcal{L}^{\otimes n}))\to
  \mathcal{J}\otimes\mathcal{L}^{\otimes n}$이 전사이다.
\end{enumerate}
\end{proposition}

$X$가 준콤팩트이므로 $f(X)$도 준콤팩트이고, 따라서 $Y$의 아핀
열린집합들로 이루어진 $f(X)$의 유한 덮개 $(U_i)$가 존재한다.
$\mathcal{L}$이 $f$에 대해 풍부할 때 조건 c)를 증명하기 위해서는
$Y$를 각 $U_i$로, $X$를 각 $f^{-1}(U_i)$로 바꾸어도 된다. 실제로 각
$i$에 대하여 $n\geq n_i$이면 ($U_i$, $f^{-1}(U_i)$,
$\mathcal{L}|f^{-1}(U_i)$에 대해) c)가 성립하도록 하는 정수 $n_i$를
얻는다면, $n_i$들 가운데 가장 큰 것을 $n_0$로 잡기만 하면 $Y$, $X$,
$\mathcal{L}$에 대해 c)를 얻는다. 그런데 $Y$가 아핀일 때에는
(\hyperref[II.4.6.6-ko]{4.6.6})을 고려하면 조건 c)는
(\hyperref[II.4.5.5-ko]{4.5.5, d)})에서 따른다. c)가 c')를 함의하는
것은 자명하다. 마지막으로 c')가 $\mathcal{L}$이 $f$에 대해 풍부함을
함의한다는 것을 증명할 때에도 $Y$가 아핀인 경우로 한정할 수 있다.
실제로 $\mathcal{O}_X|f^{-1}(U_i)$의 모든 유한형 준연접 아이디얼층
$\mathcal{J}_i$는 $\mathcal{O}_X$의 유한형 연접 아이디얼층의 제한이다
(\hyperref[I.9.4.7-ko]{I, 9.4.7}). 그리고 가정 c')에 의하면

\oldpage[II]{91}

$\mathcal{J}_i\otimes(\mathcal{L}^{\otimes n}|f^{-1}(U_i))$은 그
단면들로 생성된다((\hyperref[I.9.2.2-ko]{I, 9.2.2})와
(\hyperref[II.3.4.7-ko]{3.4.7})을 고려한다). 따라서 판정법
(\hyperref[II.4.5.5-ko]{4.5.5, d'')})을 적용하기만 하면 된다.

\begin{proposition}[4.6.9]
\label{II.4.6.9-ko}
$f:X\to Y$를 준콤팩트 사상, $\mathcal{L}$을 가역
$\mathcal{O}_X$-가군이라 하자.
\begin{enumerate}
  \item[\rm{(i)}] $n$을 $>0$인 정수라 하자. $\mathcal{L}$이 $f$에 대해
  풍부하기 위한 필요충분조건은 $\mathcal{L}^{\otimes n}$이 $f$에 대해
  풍부한 것이다.
  \item[\rm{(ii)}] $\mathcal{L}'$을 가역 $\mathcal{O}_X$-가군이라 하고,
  정수 $n>0$이 존재하여 표준 준동형
  $\sigma:f^*(f_*(\mathcal{L}'^{\otimes n}))\to
  \mathcal{L}'^{\otimes n}$이 전사라고 가정하자. 그러면
  $\mathcal{L}$이 $f$에 대해 풍부하면
  $\mathcal{L}\otimes\mathcal{L}'$도 그러하다.
\end{enumerate}
\end{proposition}

실제로 곧바로 $Y$가 아핀인 경우로 환원되며, 이때 명제는
(\hyperref[II.4.5.6-ko]{4.5.6})의 직접적인 귀결이다.

\begin{corollary}[4.6.10]
\label{II.4.6.10-ko}
두 $f$-풍부 $\mathcal{O}_X$-가군의 텐서곱은 $f$-풍부하다.
\end{corollary}

\begin{proposition}[4.6.11]
\label{II.4.6.11-ko}
$Y$를 준콤팩트 준스킴, $f:X\to Y$를 유한형 사상,
$\mathcal{L}$을 가역 $\mathcal{O}_X$-가군이라 하자.
$\mathcal{L}$이 $f$에 관해 풍부할 필요충분조건은 다음 동치인
성질들 가운데 하나를 갖는 것이다~:
\begin{enumerate}
  \item[\rm{d)}] 어떤 $n_0>0$가 존재하여, 모든 정수 $n\geq n_0$에
  대하여 $\mathcal{L}^{\otimes n}$은 $f$에 관해 매우 풍부하다.
  \item[\rm{d')}] 어떤 $n>0$가 존재하여
  $\mathcal{L}^{\otimes n}$은 $f$에 관해 매우 풍부하다.
\end{enumerate}
\end{proposition}

$\mathcal{L}$이 $f$에 관해 풍부하면, $Y$에는 유한 아핀 열린 덮개
$(U_i)$가 존재하여 $\mathcal{L}|f^{-1}(U_i)$들은 풍부하다.
그러므로 (\hyperref[II.4.5.10-ko]{4.5.10})에 따라 어떤 정수 $n_0$가
존재하여, 모든 $n\geq n_0$와 모든 $i$에 대하여
$\mathcal{L}^{\otimes n}|f^{-1}(U_i)$는
$f^{-1}(U_i)\to U_i$에 관해 매우 풍부하고, 따라서
$\mathcal{L}^{\otimes n}$은 $f$에 관해 매우 풍부하다
((\hyperref[II.4.4.5-ko]{4.4.5})). 역으로 d')에서 이미
$\mathcal{L}^{\otimes n}$이 $f$-풍부함이 따르고
((\hyperref[II.4.6.2-ko]{4.6.2})), 따라서 $\mathcal{L}$도 그러하다
((\hyperref[II.4.6.9-ko]{4.6.9, (i)})).

\begin{corollary}[4.6.12]
\label{II.4.6.12-ko}
$Y$를 준콤팩트 준스킴, $f:X\to Y$를 유한형 사상,
$\mathcal{L}$과 $\mathcal{L}'$을 두 가역 $\mathcal{O}_X$-가군이라
하자. $\mathcal{L}$이 $f$에 관해 풍부하면, 어떤 $n_0$가 존재하여
모든 $n\geq n_0$에 대하여
$\mathcal{L}^{\otimes n}\otimes\mathcal{L}'$은 $f$에 관해 매우
풍부하다.
\end{corollary}

$Y$의 유한 아핀 열린 덮개와
(\hyperref[II.4.5.11-ko]{4.5.11})을 사용하여
(\hyperref[II.4.6.11-ko]{4.6.11})에서와 같이 논증한다.

\begin{proposition}[4.6.13]
\label{II.4.6.13-ko}
\begin{enumerate}
  \item[\rm{(i)}] 모든 준스킴 $Y$에 대하여, 모든 가역
  $\mathcal{O}_Y$-가군 $\mathcal{L}$은 항등 사상 $1_Y$에 관해
  풍부하다.
  \item[\rm{(i bis)}] $f:X\to Y$를 준콤팩트 사상,
  $j:X'\to X$를 $X'$의 바탕공간에서 $X$의 한 부분공간으로의
  위상동형인 준콤팩트 사상이라 하자. $\mathcal{L}$이 $f$에 관해
  풍부한 $\mathcal{O}_X$-가군이면, $j^*(\mathcal{L})$은
  $f\circ j$에 관해 풍부하다.
  \item[\rm{(ii)}] $Z$를 준콤팩트 준스킴, $f:X\to Y$와
  $g:Y\to Z$를 두 준콤팩트 사상, $\mathcal{L}$을 $f$에 관해 풍부한
  $\mathcal{O}_X$-가군, $\mathcal{K}$를 $g$에 관해 풍부한
  $\mathcal{O}_Y$-가군이라 하자. 어떤 정수 $n_0>0$가 존재하여,
  모든 $n\geq n_0$에 대하여
  $\mathcal{L}\otimes f^*(\mathcal{K}^{\otimes n})$은
  $g\circ f$에 관해 풍부하다.
  \item[\rm{(iii)}] $f:X\to Y$를 준콤팩트 사상,
  $g:Y'\to Y$를 사상이라 하고, $X'=X_{(Y')}$라 놓자.
  $\mathcal{L}$이 $f$에 관해 풍부한 $\mathcal{O}_X$-가군이면,
  $\mathcal{L}'=\mathcal{L}\otimes_Y\mathcal{O}_{Y'}$은
  $f_{(Y')}$에 관해 풍부한 $\mathcal{O}_{X'}$-가군이다.
  \item[\rm{(iv)}] $f_i:X_i\to Y_i$ ($i=1,2$)를 두 준콤팩트
  $S$-사상이라 하자. $\mathcal{L}_i$가 $f_i$에 관해 풍부한
  $\mathcal{O}_{X_i}$-가군이면 ($i=1,2$),
  $\mathcal{L}_1\otimes_S\mathcal{L}_2$는
  $f_1\times_S f_2$에 관해 풍부하다.
  \item[\rm{(v)}] $f:X\to Y$와 $g:Y\to Z$를 $g\circ f$가
  준콤팩트인 두 사상이라 하자. $\mathcal{O}_X$-가군
  $\mathcal{L}$이 $g\circ f$에 관해 풍부하고, $g$가 분리 사상이거나
  $X$의 바탕공간이 국소 뇌터이면, $\mathcal{L}$은 $f$에 관해
  풍부하다.

  \oldpage[II]{92}

  \item[\rm{(vi)}] $f:X\to Y$를 준콤팩트 사상,
  $j$를 표준 주입 사상 $X_{\mathrm{red}}\to X$라 하자.
  $\mathcal{L}$이 $f$에 관해 풍부한 $\mathcal{O}_X$-가군이면,
  $j^*(\mathcal{L})$은 $f_{\mathrm{red}}$에 관해 풍부하다.
\end{enumerate}
\end{proposition}

먼저 (v)와 (vi)는
(\hyperref[II.4.6.4-ko]{4.6.4})를
(\hyperref[II.4.4.5-ko]{4.4.5}) 대신 사용하여
(\hyperref[II.4.4.10-ko]{4.4.10})에서와 같은 논법으로 (i), (i bis),
(iv)에서 도출되며, 논증의 세부는 독자에게 맡긴다. 주장 (i)는
(\hyperref[II.4.4.10-ko]{4.4.10, (i)})와
(\hyperref[II.4.6.2-ko]{4.6.2})의 자명한 귀결이다. (i bis), (iii),
(iv)를 증명하기 위해 다음 보조정리를 사용할 것이다~:

\begin{lemma}[4.6.13.1]
\label{II.4.6.13.1-ko}
\begin{enumerate}
  \item[\rm{(i)}] $u:Z\to S$를 사상, $\mathcal{L}$을 가역
  $\mathcal{O}_S$-가군, $s$를 $S$ 위의 $\mathcal{L}$의 단면,
  $s'$을 이에 표준적으로 대응하는 $Z$ 위의
  $u^*(\mathcal{L})=\mathcal{L}'$의 단면이라 하자. 그러면
  $Z_{s'}=u^{-1}(S_s)$이다.
  \item[\rm{(ii)}] $Z$, $Z'$을 두 $S$-준스킴, $p$, $p'$을
  $T=Z\times_S Z'$의 두 사영, $\mathcal{L}$ (각각
  $\mathcal{L}'$)을 가역 $\mathcal{O}_Z$-가군 (각각 가역
  $\mathcal{O}_{Z'}$-가군), $t$ (각각 $t'$)를 $Z$ (각각 $Z'$)
  위의 $\mathcal{L}$ (각각 $\mathcal{L}'$)의 단면, $s$ (각각
  $s'$)를 이에 대응하는 $Z\times_S Z'$ 위의 $p^*(\mathcal{L})$
  (각각 $p'^*(\mathcal{L}')$)의 단면이라 하자. 그러면
  $T_{s\otimes s'}=Z_t\times_S Z'_{t'}$이다.
\end{enumerate}
\end{lemma}

정의들로부터, 고려하는 모든 준스킴이 아핀인 경우로 환원할 수 있다.
더 나아가 (i)에서는 $\mathcal{L}=\mathcal{O}_S$라 가정할 수 있다~;
그러면 주장 (i)는 (\hyperref[I.1.2.2.2-ko]{I, 1.2.2.2})에서 곧바로
따른다. 마찬가지로 (ii)에서는
$\mathcal{L}=\mathcal{O}_Z$, $\mathcal{L}'=\mathcal{O}_{Z'}$인
경우로 한정할 수 있고, 이때 주장은 보조정리
(\hyperref[II.4.3.2.4-ko]{4.3.2.4})로 환원된다.

이제 (i bis)를 증명하자. $Y$가 아핀이라고 가정할 수 있으므로
(\hyperref[II.4.6.4-ko]{4.6.4}), $\mathcal{L}$은 풍부하다
(\hyperref[II.4.6.6-ko]{4.6.6})~; $s$가
$\Gamma(X,\mathcal{L}^{\otimes n})$ ($n>0$)들의 합집합을 달릴 때,
$X_s$들은 $X$의 위상의 기저를 이루므로
(\hyperref[II.4.5.2-ko]{4.5.2, a)}), 가정에 따라
$j^{-1}(X_s)$들은 $X'$의 위상의 기저를 이룬다~; 따라서 보조정리
(\hyperref[II.4.6.13.1-ko]{4.6.13.1, (i)})와
(\hyperref[II.4.5.2-ko]{4.5.2, a)})에 의해
$j^*(\mathcal{L})$이 풍부하다는 결론을 얻는다.

다음으로 (iii)을 증명하자. 다시 $Y$와 $Y'$이 아핀이라고 가정할 수
있고 (\hyperref[II.4.6.4-ko]{4.6.4}), 이로부터 사영
$h:X'\to X$가 아핀임이 따른다
(\hyperref[II.1.5.5-ko]{1.5.5}). $\mathcal{L}$은 풍부하므로
(\hyperref[II.4.6.6-ko]{4.6.6}), $s$가 $X$ 위의
$\mathcal{L}^{\otimes n}$ ($n>0$)의 단면들 가운데 $X_s$가 아핀인
것들을 달릴 때 $X_s$들은 $X$를 덮는다
(\hyperref[II.4.5.2-ko]{4.5.2, a')})~; 따라서 $h^{-1}(X_s)$들은
아핀이고 (\hyperref[II.1.2.5-ko]{1.2.5}) $X'$을 덮는다~; 그러므로
다시 보조정리 (\hyperref[II.4.6.13.1-ko]{4.6.13.1, (i)})와
(\hyperref[II.4.5.2-ko]{4.5.2, a')})에 의해 $\mathcal{L}'$이
풍부하다는 결론을 얻는다. 여기서 사상 $f_{(Y')}$은 준콤팩트이다
(\hyperref[I.6.6.4-ko]{I, 6.6.4, (iii)}).

(iv)를 증명하기 위해, 먼저 $f_1\times_S f_2$가 준콤팩트임을
주목하자 (\hyperref[I.6.6.4-ko]{I, 6.6.4, (iv)}). 더 나아가
$S$, $Y_1$, $Y_2$가 아핀이라고 가정할 수 있으므로
((\hyperref[II.4.6.4-ko]{4.6.4}) 및
(\hyperref[I.3.2.7-ko]{I, 3.2.7})), $\mathcal{L}_i$는 풍부하다
($i=1,2$) (\hyperref[II.4.6.6-ko]{4.6.6}). $s_i$가
$\mathcal{L}_i^{\otimes n_i}$의 단면들 가운데 $(X_i)_{s_i}$가 아핀인
것들을 달릴 때, 열린집합
$(X_1)_{s_1}\times_S(X_2)_{s_2}$들은 $X_1\times_S X_2$의 덮개를
이룬다 (\hyperref[II.4.5.2-ko]{4.5.2, a')}). 더욱이 $s_1$과
$s_2$를 적절한 거듭제곱으로 바꾸어도 $(X_i)_{s_i}$는 변하지 않으므로,
언제나 $n_1=n_2$라고 가정할 수 있다. 이제
(\hyperref[II.4.6.13.1-ko]{4.6.13.1, (ii)})와
(\hyperref[II.4.5.2-ko]{4.5.2, a')})로부터
$\mathcal{L}_1\otimes_S\mathcal{L}_2$가 풍부함이 따르고, 따라서
$Y_1\times_S Y_2$가 아핀이므로
(\hyperref[II.4.6.6-ko]{4.6.6}) 주장을 얻는다.

(ii)를 증명하는 일만 남았다.
(\hyperref[II.4.4.10-ko]{4.4.10})에서와 같은 논법을 쓰되 여기서는
(\hyperref[II.4.6.4-ko]{4.6.4})를 사용하면, $Z$가 아핀인 경우로
한정할 수 있다. 그러면 $\mathcal{K}$은 풍부하고 $Y$는 준콤팩트이므로,
유한 개의 단면
$s_i\in\Gamma(Y,\mathcal{K}^{\otimes k_i})$가 존재하여 $Y_{s_i}$들은

\oldpage[II]{93}

아핀이고 $Y$를 덮는다
(\hyperref[II.4.5.2-ko]{4.5.2, a')})~; $s_i$들을 적절한
거듭제곱으로 바꾸면, 더 나아가 모든 $k_i$가 하나의 같은 정수 $k$와
같다고 가정할 수 있다. $s'_i$를 $s_i$에 표준적으로 대응하는 $X$ 위의
$f^*(\mathcal{K}^{\otimes k})$의 단면이라 하자. 그러면
$X_{s'_i}=f^{-1}(Y_{s_i})$ (4.6.1.13, (i))이고 이들은 $X$를 덮는다.
$\mathcal{L}|X_{s'_i}$가 풍부하므로
((\hyperref[II.4.6.4-ko]{4.6.4}) 및
(\hyperref[II.4.6.6-ko]{4.6.6})), 각 $i$에 대하여 유한 개의 단면
$t_{ij}\in\Gamma(X,\mathcal{L}^{\otimes n_{ij}})$가 존재하여,
$X_{t_{ij}}$들은 아핀이고 $X_{s'_i}$에 포함되며 $X_{s'_i}$를 덮는다
(\hyperref[II.4.5.2-ko]{4.5.2, a')})~; 더 나아가 모든
$n_{ij}$가 하나의 같은 수 $n$과 같다고 가정할 수 있다. 그런데 $X$는
분리이고 준콤팩트이므로, 어떤 정수 $m>0$이 존재하고 각 $(i,j)$에
대하여 단면
\[
  u_{ij}\in\Gamma(X,\mathcal{L}^{\otimes n}\otimes
  f^*(\mathcal{K}^{\otimes mk}))
\]
이 존재하여 $t_{ij}\otimes s_i'^{\otimes m}$은 $u_{ij}$의
$X_{s'_i}$에 대한 제한이다 (\hyperref[I.9.3.1-ko]{I, 9.3.1})~;
더 나아가 $X_{u_{ij}}=X_{t_{ij}}$이므로, $X_{u_{ij}}$들은 아핀이고
$X$를 덮는다. 또한 $m$이 $nr$ 꼴이라고 가정할 수 있다~;
$n_0=rk$라 놓으면
(\hyperref[II.4.5.2-ko]{4.5.2, a')})에 따라
$\mathcal{L}\otimes_{\mathcal{O}_X}f^*(\mathcal{K}^{\otimes n_0})$이
풍부함을 알 수 있다. 더 나아가 어떤 $h_0>0$이 존재하여,
$h\geq h_0$이면 $\mathcal{K}^{\otimes h}$는 $Y$ 위의 단면들로
생성된다 (\hyperref[II.4.5.5-ko]{4.5.5})~; 따라서 더욱이
$h\geq h_0$이면 $f^*(\mathcal{K}^{\otimes h})$는 정의상 $X$ 위의
단면들로 생성된다 ((\hyperref[0.3.7.1-ko]{0, 3.7.1}) 및
(\hyperref[II.4.4.1-ko]{4.4.1})). 이로부터 $h\geq h_0$이면
$\mathcal{L}\otimes f^*(\mathcal{K}^{\otimes n_0+h})$가 풍부함이
따르고 (\hyperref[II.4.5.6-ko]{4.5.6}), 이로써 증명이 끝난다.

\begin{remark}[4.6.14]
\label{II.4.6.14-ko}
(ii)의 조건 아래에서
$\mathcal{L}\otimes f^*(\mathcal{K})$가 $g\circ f$에 관해 풍부하다고
생각해서는 안 된다~; 실제로
$\mathcal{L}\otimes f^*(\mathcal{K}^{-1})$도 $f$에 관해 풍부하므로
(\hyperref[II.4.6.5-ko]{4.6.5}), 그렇게 생각한다면
$\mathcal{L}$이 $g\circ f$에 관해 풍부하다는 결론을 얻게 된다~;
특히 $g$로 항등 사상을 취하면 모든 가역 $\mathcal{O}_X$-가군이
$f$에 관해 풍부하게 되지만, 일반적으로 그렇지 않다
((\hyperref[II.5.1.6-ko]{5.1.6}),
(\hyperref[II.5.3.4-ko]{5.3.4, (i)}) 및
(\hyperref[II.5.3.1-ko]{5.3.1}) 참조).
\end{remark}

\begin{proposition}[4.6.15]
\label{II.4.6.15-ko}
$f:X\to Y$를 준콤팩트 사상, $\mathcal{J}$를 $\mathcal{O}_X$의
국소적으로 멱영인 준연접 아이디얼층, $Z$를 $\mathcal{J}$로 정의되는
$X$의 닫힌 부분준스킴, $j:Z\to X$를 표준 포함 사상이라 하자. 가역
$\mathcal{O}_X$-가군 $\mathcal{L}$이 $f$에 관해 풍부할 필요충분조건은
$j^*(\mathcal{L})$이 $f\circ j$에 관해 풍부한 것이다.
\end{proposition}

실제로 문제는 $Y$ 위에서 국소적이므로
(\hyperref[II.4.6.4-ko]{4.6.4}), $Y$가 아핀이라고 가정할 수 있다~;
그러면 $X$가 준콤팩트이므로 $\mathcal{J}$가 멱영이라고 가정할 수
있다. (\hyperref[II.4.6.6-ko]{4.6.6})을 고려하면, 이 명제는 곧
(\hyperref[II.4.5.13-ko]{4.5.13})이다.

\begin{corollary}[4.6.16]
\label{II.4.6.16-ko}
$X$를 국소 뇌터 준스킴, $f:X\to Y$를 준콤팩트 사상이라 하자. 가역
$\mathcal{O}_X$-가군 $\mathcal{L}$이 $f$에 관해 풍부할 필요충분조건은
표준 포함 사상 $X_{\mathrm{red}}\to X$에 의한 그 역상
$\mathcal{L}'$이 $f_{\mathrm{red}}$에 관해 풍부한 것이다.
\end{corollary}

이 조건이 필요함은 이미 보았다
(\hyperref[II.4.6.13-ko]{4.6.13, (vi)})~; 역으로 이 조건이 성립한다고
하자. $\mathcal{L}$이 $f$에 관해 풍부함을 증명하기 위해서는 $Y$가
아핀인 경우만 생각하면 된다 (\hyperref[II.4.6.4-ko]{4.6.4})~;
그러면 $Y_{\mathrm{red}}$도 아핀이므로 $\mathcal{L}'$은 풍부하고
(\hyperref[II.4.6.6-ko]{4.6.6}), 따라서
(\hyperref[II.4.5.13-ko]{4.5.13})에 의해 $\mathcal{L}$도 풍부하다.
실제로 이때 $X$는 뇌터이고 $X_{\mathrm{red}}$는 준연접 멱영
아이디얼층으로 정의되는 $X$의 닫힌 부분준스킴이다
(\hyperref[I.6.1.6-ko]{I, 6.1.6}).

\begin{proposition}[4.6.17]
\label{II.4.6.17-ko}
(\hyperref[II.4.4.11-ko]{4.4.11})의 표기와 가정 아래에서,
$\mathcal{L}''$이 $f''$에 관해 풍부할 필요충분조건은
$\mathcal{L}$이 $f$에 관해 풍부하고 $\mathcal{L}'$이 $f'$에 관해
풍부한 것이다.
\end{proposition}

\oldpage[II]{94}

조건의 필요성은
(\hyperref[II.4.6.13-ko]{4.6.13, (i bis)})에서 따른다. 조건이
충분함을 보이기 위해서는 $Y$가 아핀인 경우만 생각하면 된다. 그러면
$\mathcal{L}$, $\mathcal{L}'$, $\mathcal{L}''$에 판정법
(\hyperref[II.4.5.2-ko]{4.5.2, a)})을 적용하고, $X$ 위의
$\mathcal{L}$의 한 단면은 $0$으로 연장하여 $X''$ 위의
$\mathcal{L}''$의 한 단면으로 만들 수 있음을 관찰하면,
$\mathcal{L}''$이 풍부하다는 결론을 얻는다.

\begin{proposition}[4.6.18]
\label{II.4.6.18-ko}
$Y$를 준콤팩트 준스킴, $\mathcal{S}$를 유한 생성 준연접 등급
$\mathcal{O}_Y$-대수, $X=\operatorname{Proj}(\mathcal{S})$라 하고,
$f:X\to Y$를 구조 사상이라 하자. 그러면 $f$는 유한형이고, 어떤
정수 $d>0$가 존재하여 $\mathcal{O}_X(d)$는 가역이고 $f$-풍부하다.
\end{proposition}

(\hyperref[II.3.1.10-ko]{3.1.10})에 따라 어떤 정수 $d>0$가 존재하여
$\mathcal{S}^{(d)}$는 $\mathcal{S}_d$에 의해 생성된다.
$X$와 $X^{(d)}=\operatorname{Proj}(\mathcal{S}^{(d)})$ 사이의 표준
동형사상에서 $\mathcal{O}_X(d)$는 $\mathcal{O}_{X^{(d)}}(1)$과
동일시됨을 알고 있다
(\hyperref[II.3.2.9-ko]{3.2.9, (ii)}). 따라서 $\mathcal{S}$가
$\mathcal{S}_1$에 의해 생성되는 경우로 환원된다~; 그러면 명제는
(\hyperref[II.4.4.3-ko]{4.4.3})과
(\hyperref[II.4.6.2-ko]{4.6.2})에서 따른다($f$가 유한형 사상이라는
사실 (\hyperref[II.3.4.1-ko]{3.4.1})을 고려한다).

\section{준아핀 사상~; 준사영 사상. 고유 사상~; 사영 사상}
\label{section:II.5-ko}

\subsection{준아핀 사상.}
\label{subsection:II.5.1-ko}

\begin{definition}[5.1.1]
\label{II.5.1.1-ko}
아핀 스킴의 어떤 준콤팩트 열린집합 위에 유도된 부분스킴과 동형인 스킴을
\emph{준아핀 스킴}이라 한다. 사상 $f:X\to Y$에 대하여, $f$를 통해
$X$를 $Y$-준스킴으로 보았을 때, $Y$의 아핀 열린집합들로 이루어진 덮개
$(U_\alpha)$가 존재하여 $f^{-1}(U_\alpha)$들이 준아핀 스킴이면, $f$를
\emph{준아핀}이라 하며, 또는 $X$를 \emph{준아핀 $Y$-스킴}이라 한다.
\end{definition}

준아핀 사상이 \emph{분리 사상}
(\hyperref[I.5.5.5-ko]{I, 5.5.5}와 \hyperref[I.5.5.8-ko]{5.5.8})이고
\emph{준콤팩트 사상}(\hyperref[I.6.6.1-ko]{I, 6.6.1})임은 명백하다~;
모든 아핀 사상은 분명히 준아핀이다.

모든 준스킴 $X$에 대하여 $A=\Gamma(X,\mathcal{O}_X)$라 놓으면, 항등
준동형 $A\to A=\Gamma(X,\mathcal{O}_X)$는 \emph{표준}이라 불리는 사상
$X\to\operatorname{Spec}(A)$를 정의함을 상기하자
(\hyperref[I.2.2.4-ko]{I, 2.2.4})~; 이는 사실 $\mathcal{L}=\mathcal{O}_X$인
특별한 경우에 (\hyperref[II.4.5.1-ko]{4.5.1})에서 정의된 표준 사상에
다름 아니며, 여기서는 $\operatorname{Proj}(A[T])$가
$\operatorname{Spec}(A)$와 표준적으로 동일시됨을 사용하였다
(\hyperref[II.3.1.7-ko]{3.1.7}).

\begin{proposition}[5.1.2]
\label{II.5.1.2-ko}
$X$를 준콤팩트 스킴 또는 바탕공간이 뇌터인 준스킴이라 하고,
$A=\Gamma(X,\mathcal{O}_X)$라 하자. 다음 조건들은 서로 동치이다~:
\begin{enumerate}
  \item[\rm{a)}] $X$는 준아핀 스킴이다.
  \item[\rm{b)}] 표준 사상 $u:X\to\operatorname{Spec}(A)$는 열린 몰입
  사상이다.
  \item[\rm{b')}] 표준 사상 $u:X\to\operatorname{Spec}(A)$는 $X$에서
  $\operatorname{Spec}(A)$의 바탕공간의 한 부분공간 위로 가는
  위상동형이다.
  \item[\rm{c)}] $\mathcal{O}_X$-가군 $\mathcal{O}_X$는 $u$에 관해 매우
  풍부하다 (\hyperref[II.4.4.2-ko]{4.4.2}).
  \item[\rm{c')}] $\mathcal{O}_X$-가군 $\mathcal{O}_X$는 풍부하다
  (\hyperref[II.4.5.1-ko]{4.5.1}).
  \item[\rm{d)}] $f$가 $A$의 원소 전체를 달릴 때, $X_f$들은 $X$의
  위상의 기저를 이룬다.
  \item[\rm{d')}] $f$가 $A$의 원소 전체를 달릴 때, 그중 아핀인
  $X_f$들은 $X$의 덮개를 이룬다.

\oldpage[II]{95}

  \item[\rm{e)}] 모든 준연접 $\mathcal{O}_X$-가군은 $X$ 위의 단면들로
  생성된다.
  \item[\rm{e')}] $\mathcal{O}_X$의 모든 유한형 준연접 아이디얼층은
  $X$ 위의 단면들로 생성된다.
\end{enumerate}
\end{proposition}

b)가 a)를 함의함은 명백하고, a)는 항등 사상에
(\hyperref[II.4.4.4-ko]{4.4.4})의 판정법 b)를 적용하면 c)를 함의한다
(명제 앞의 주의도 고려한다)~; 한편 c)는 c')을 함의하며
(\hyperref[II.4.5.10-ko]{4.5.10, (i)}),
(\hyperref[II.4.5.2-ko]{4.5.2, 판정법 b)와 b')})에 따라 c'), b), b')은
서로 동치이다. 끝으로 (\hyperref[II.4.5.2-ko]{4.5.2, 판정법 a), a'), c)})와
(\hyperref[II.4.5.5-ko]{4.5.5, 판정법 d'')})에 따라 c')은 판정법 d), d'),
e), e') 각각과 동치이다.

더 나아가 위의 표기 아래에서, 그중 아핀인 $X_f$들이 $X$의 위상의
\emph{기저}를 이루며, 표준 사상 $u$가 \emph{우세 사상}임을 주목하자
(\hyperref[II.4.5.2-ko]{4.5.2}).

\begin{corollary}[5.1.3]
\label{II.5.1.3-ko}
$X$를 준콤팩트 준스킴이라 하자. $X$에서 아핀 스킴 $Y$로 가는 사상
$v:X\to Y$가 존재하고, 이것이 $X$에서 $Y$의 한 열린 부분공간 위로
가는 위상동형이면, $X$는 준아핀이다.
\end{corollary}

실제로 $Y$ 위의 $\mathcal{O}_Y$의 단면들의 족 $(g_\alpha)$가 존재하여
$D(g_\alpha)$들이 $v(X)$의 위상의 기저를 이룬다~; $v=(\psi,\theta)$이고
$f_\alpha=\theta(g_\alpha)$라 놓으면
$X_{f_\alpha}=\psi^{-1}(D(g_\alpha))$
(\hyperref[I.2.2.4.1-ko]{I, 2.2.4.1})이므로, $X_{f_\alpha}$들은 $X$의
위상의 기저를 이루며 (\hyperref[II.5.1.2-ko]{5.1.2})의 판정법 d)가
성립한다.

\begin{corollary}[5.1.4]
\label{II.5.1.4-ko}
$X$가 준아핀 스킴이면, 모든 가역 $\mathcal{O}_X$-가군은 매우 풍부하고
(표준 사상에 관하여), 따라서 특히 풍부하다.
\end{corollary}

실제로 그러한 가군 $\mathcal{L}$은 $X$ 위의 단면들로 생성되므로
(\hyperref[II.5.1.2-ko]{5.1.2, e)}),
$\mathcal{L}\otimes\mathcal{O}_X=\mathcal{L}$은 매우 풍부하다
(\hyperref[II.4.4.8-ko]{4.4.8}).

\begin{corollary}[5.1.5]
\label{II.5.1.5-ko}
$X$를 준콤팩트 준스킴이라 하자. $\mathcal{L}$과
$\mathcal{L}^{-1}$이 모두 풍부한 가역 $\mathcal{O}_X$-가군
$\mathcal{L}$이 존재하면, $X$는 준아핀 스킴이다.
\end{corollary}

실제로 이 경우 $\mathcal{O}_X=\mathcal{L}\otimes\mathcal{L}^{-1}$은
풍부하다 (\hyperref[II.4.5.7-ko]{4.5.7}).

\begin{proposition}[5.1.6]
\label{II.5.1.6-ko}
$f:X\to Y$를 준콤팩트 사상이라 하자. 다음 조건들은 동치이다~:
\begin{enumerate}
  \item[\rm{a)}] $f$는 준아핀 사상이다.
  \item[\rm{b)}] $\mathcal{O}_Y$-대수
  $f_*(\mathcal{O}_X)=\mathcal{A}(X)$는 준연접이고, 항등 준동형
  $\mathcal{A}(X)\to\mathcal{A}(X)$에 대응하는 표준 사상
  $X\to\operatorname{Spec}(\mathcal{A}(X))$
  (\hyperref[II.1.2.7-ko]{1.2.7})는 열린 몰입이다.
  \item[\rm{b')}] $\mathcal{O}_Y$-대수 $\mathcal{A}(X)$는
  준연접이고, 표준 사상 $X\to\operatorname{Spec}(\mathcal{A}(X))$는
  $X$에서 $\operatorname{Spec}(\mathcal{A}(X))$의 한 부분공간 위로
  가는 위상동형사상이다.
  \item[\rm{c)}] $\mathcal{O}_X$-가군 $\mathcal{O}_X$는
  $f$에 관해 매우 풍부하다.
  \item[\rm{c')}] $\mathcal{O}_X$-가군 $\mathcal{O}_X$는
  $f$에 관해 풍부하다.
  \item[\rm{d)}] 사상 $f$는 분리 사상이고, 모든 준연접
  $\mathcal{O}_X$-가군 $\mathcal{F}$에 대하여 표준 준동형
  $\sigma:f^*(f_*(\mathcal{F}))\to\mathcal{F}$
  (\hyperref[0.4.4.3-ko]{0, 4.4.3})는 전사이다.
\end{enumerate}
더 나아가 $f$가 준아핀 사상이면, 모든 가역 $\mathcal{O}_X$-가군
$\mathcal{L}$은 $f$에 관해 매우 풍부하다.
\end{proposition}

a)와 c')의 동치는 $f$-풍부성이 $Y$ 위에서 국소적이라는 사실
(\hyperref[II.4.6.4-ko]{4.6.4}), 정의
(\hyperref[II.5.1.1-ko]{5.1.1}), 그리고 판정법
(\hyperref[II.5.1.2-ko]{5.1.2, c')})에서 따른다. 나머지 성질들도
$Y$ 위에서 국소적이며

\oldpage[II]{96}

따라서 (\hyperref[II.5.1.2-ko]{5.1.2})와
(\hyperref[II.5.1.4-ko]{5.1.4})에서 곧바로 따른다. 여기서는 $f$가
분리 사상일 때 $f_*(\mathcal{F})$가 준연접이라는 사실을 사용하였다
(\hyperref[I.9.2.2-ko]{I, 9.2.2, a)}).

\begin{corollary}[5.1.7]
\label{II.5.1.7-ko}
$f:X\to Y$를 준아핀 사상이라 하자. $Y$의 모든 열린집합 $U$에
대하여 $f$의 제한 사상 $f^{-1}(U)\to U$는 준아핀 사상이다.
\end{corollary}

\begin{corollary}[5.1.8]
\label{II.5.1.8-ko}
$Y$를 아핀 스킴, $f:X\to Y$를 준콤팩트 사상이라 하자. $f$가
준아핀 사상일 필요충분조건은 $X$가 준아핀 스킴인 것이다.
\end{corollary}

이는 (\hyperref[II.5.1.6-ko]{5.1.6})과
(\hyperref[II.4.6.6-ko]{4.6.6})의 직접적인 귀결이다.

\begin{corollary}[5.1.9]
\label{II.5.1.9-ko}
$Y$를 준콤팩트 스킴 또는 그 바탕공간이 뇌터인 준스킴,
$f:X\to Y$를 유한형 사상이라 하자. $f$가 준아핀 사상이면,
$\mathcal{A}(X)=f_*(\mathcal{O}_X)$의 준연접
부분-$\mathcal{O}_Y$-대수 $\mathcal{B}$로서 유한형이고
(\hyperref[I.9.6.2-ko]{I, 9.6.2}), 표준 포함 준동형
$\mathcal{B}\to\mathcal{A}(X)$에 대응하는 사상
$X\to\operatorname{Spec}(\mathcal{B})$
(\hyperref[II.1.2.7-ko]{1.2.7})가 몰입이 되는 것이 존재한다.
더 나아가 $\mathcal{B}$를 포함하는 $\mathcal{A}(X)$의 모든 유한형
준연접 부분-$\mathcal{O}_Y$-대수 $\mathcal{B}'$도 같은 성질을 갖는다.
\end{corollary}

실제로 $\mathcal{A}(X)$는 그 유한형 준연접
부분-$\mathcal{O}_Y$-대수들의 귀납극한이다
(\hyperref[I.9.6.5-ko]{I, 9.6.5})~; 따라서 이 결과는
$\operatorname{Spec}(\mathcal{A}(X))$와
$\operatorname{Proj}(\mathcal{A}(X)[T])$의 동일시를 고려하면
(\hyperref[II.3.8.4-ko]{3.8.4})의 특별한 경우이다
(\hyperref[II.3.1.7-ko]{3.1.7}).

\begin{proposition}[5.1.10]
\label{II.5.1.10-ko}
\begin{enumerate}
  \item[\rm{(i)}] 준콤팩트 사상 $X\to Y$가 $X$의 바탕공간에서
  $Y$의 바탕공간의 한 부분공간으로의 위상동형이면(특히 닫힌 몰입
  사상이면), 이 사상은 준아핀이다.
  \item[\rm{(ii)}] 두 준아핀 사상의 합성은 준아핀이다.
  \item[\rm{(iii)}] $f:X\to Y$가 준아핀 $S$-사상이면, 기저 준스킴의
  모든 확장 $S'\to S$에 대하여
  $f_{(S')}:X_{(S')}\to Y_{(S')}$는 준아핀 사상이다.
  \item[\rm{(iv)}] $f:X\to Y$와 $g:X'\to Y'$이 두 준아핀
  $S$-사상이면, $f\times_S g$는 준아핀이다.
  \item[\rm{(v)}] $f:X\to Y$, $g:Y\to Z$가 $g\circ f$가 준아핀인 두
  사상이고, $g$가 분리 사상이거나 $X$의 바탕공간이 국소 뇌터이면,
  $f$는 준아핀이다.
  \item[\rm{(vi)}] $f$가 준아핀 사상이면 $f_{\mathrm{red}}$도
  준아핀 사상이다.
\end{enumerate}
\end{proposition}

판정법 (\hyperref[II.5.1.6-ko]{5.1.6, c')})을 고려하면, (i), (iii),
(iv), (v), (vi)는 각각
(\hyperref[II.4.6.13-ko]{4.6.13, (i bis), (iii), (iv), (v), (vi)})에서
즉시 따른다. (ii)를 증명하기 위해서는 $Z$가 아핀인 경우로 한정할 수
있고, 그러면 주장은 $\mathcal{L}=\mathcal{O}_X$와
$\mathcal{K}=\mathcal{O}_Y$에 적용한
(\hyperref[II.4.6.13-ko]{4.6.13, (ii)})에서 곧바로 따른다.

\begin{remark}[5.1.11]
\label{II.5.1.11-ko}
$f:X\to Y$, $g:Y\to Z$를 $X\times_ZY$가 국소 뇌터인 두 사상이라
하자. 그래프 몰입 사상 $\Gamma_f:X\to X\times_ZY$는 준콤팩트이므로
(\hyperref[I.6.3.5-ko]{I, 6.3.5}) 준아핀이고,
(\hyperref[I.5.5.12-ko]{I, 5.5.12})에 따라 (v)에서 $g$가 분리
사상이라는 가정을 없애도 결론은 여전히 성립한다.
\end{remark}

\begin{proposition}[5.1.12]
\label{II.5.1.12-ko}
$f:X\to Y$를 준콤팩트 사상, $g:X'\to X$를 준아핀 사상이라 하자.
$\mathcal{L}$이 $f$에 관해 풍부한 $\mathcal{O}_X$-가군이면,
$g^*(\mathcal{L})$은 $f\circ g$에 관해 풍부한
$\mathcal{O}_{X'}$-가군이다.
\end{proposition}

실제로 $\mathcal{O}_{X'}$은 $g$에 관해 매우 풍부하고, 문제는 $Y$
위에서 국소적이므로 (\hyperref[II.4.6.4-ko]{4.6.4}),
(\hyperref[II.4.6.13-ko]{4.6.13, (ii)})에 따라 ($Y$가 아핀일 때)
어떤 정수 $n$이 존재하여 $g^*(\mathcal{L})^{\otimes n}$은
$f\circ g$에 관해 풍부하다. 따라서
$g^*(\mathcal{L})$은 $f\circ g$에 관해 풍부하다
(\hyperref[II.4.6.9-ko]{4.6.9}).

\oldpage[II]{97}

\subsection{세르 판정법.}
\label{subsection:II.5.2-ko}

\begin{theorem}[5.2.1]
\label{II.5.2.1-ko}
\emph{(세르 판정법.)} $X$를 준콤팩트 스킴 또는 바탕공간이 뇌터인
준스킴이라 하자. 다음 조건들은 서로 동치이다~:
\begin{enumerate}
  \item[\rm{a)}] $X$는 아핀 스킴이다.
  \item[\rm{b)}] 원소족
  $f_\alpha\in A=\Gamma(X,\mathcal{O}_X)$가 존재하여, $X_{f_\alpha}$들은
  아핀이고 $A$에서 $f_\alpha$들로 생성되는 아이디얼은 $A$와 같다.
  \item[\rm{c)}] 함자 $\Gamma(X,\mathcal{F})$는 준연접
  $\mathcal{O}_X$-가군의 범주에서 $\mathcal{F}$에 관해 완전하다~; 다시
  말해
  \[
    0\to\mathcal{F}'\to\mathcal{F}\to\mathcal{F}''\to0 \tag{*}
  \]
  이 준연접 $\mathcal{O}_X$-가군들의 완전열이면, 열
  \[
    0\to\Gamma(X,\mathcal{F}')\to\Gamma(X,\mathcal{F})
      \to\Gamma(X,\mathcal{F}'')\to0
  \]
  은 완전하다.
  \item[\rm{c')}] $\mathcal{F}$가 유한곱 $\mathcal{O}_X^n$의 한
  부분-$\mathcal{O}_X$-가군과 동형인 모든 준연접
  $\mathcal{O}_X$-가군의 완전열 (*)에 대하여 성질 c)가 성립한다.
  \item[\rm{d)}] 모든 준연접 $\mathcal{O}_X$-가군 $\mathcal{F}$에
  대하여 $\mathrm{H}^1(X,\mathcal{F})=0$이다.
  \item[\rm{d')}] $\mathcal{O}_X$의 모든 준연접 아이디얼층
  $\mathcal{J}$에 대하여 $\mathrm{H}^1(X,\mathcal{J})=0$이다.
\end{enumerate}
\end{theorem}

a)가 b)를 함의함은 명백하다~; 한편 b)는 $X_{f_\alpha}$들이 $X$를
덮음을 함의한다. 실제로 가정에 따라 단면 $1$은 $f_\alpha$들의
선형결합이고, 따라서 $D(f_\alpha)$들은 $\operatorname{Spec}(A)$를
덮는다. 그러므로 (\hyperref[II.4.5.2-ko]{4.5.2})의 마지막 주장에
따라 $X\to\operatorname{Spec}(A)$는 동형사상이다.

a)가 c)를 함의함을 알고 있고 (\hyperref[I.1.3.11-ko]{I, 1.3.11}),
c)는 자명하게 c')을 함의한다. c')이 b)를 함의함을 증명하자. 먼저
c')에 따라, 모든 닫힌점 $x\in X$와 $x$의 모든 열린 근방 $U$에
대하여 $x\in X_f\subset X-U$를 만족하는 $f\in A$가 존재한다.
실제로 $\mathcal{J}$ (각각 $\mathcal{J}'$)를 바탕공간이 $X-U$
(각각 $(X-U)\cup\{x\}$)인 $X$의 축소 닫힌 부분준스킴을 정의하는
$\mathcal{O}_X$의 준연접 아이디얼층이라 하자
(\hyperref[I.5.2.1-ko]{I, 5.2.1})~; 그러면 분명히
$\mathcal{J}'\subset\mathcal{J}$이고,
$\mathcal{J}''=\mathcal{J}/\mathcal{J}'$은 지지집합이 $\{x\}$이며
$\mathcal{J}''_x=k(x)$인 준연접 $\mathcal{O}_X$-가군이다. 완전열
$0\to\mathcal{J}'\to\mathcal{J}\to\mathcal{J}''\to0$에 c')을
적용하면 $\Gamma(X,\mathcal{J})\to\Gamma(X,\mathcal{J}'')$가
전사임을 알 수 있다. 따라서 $x$에서의 싹이 $1_x$인
$\mathcal{J}''$의 단면은 어떤 단면
$f\in\Gamma(X,\mathcal{J})\subset\Gamma(X,\mathcal{O}_X)$의 상이고,
정의에 따라 $f(x)=1_x$이며 $X-U$에서 $f(y)=0$이다. 이로써 우리의
주장이 증명된다. 더 나아가 $U$가 아핀이면 $X_f$도 아핀이다
(\hyperref[I.1.3.6-ko]{I, 1.3.6}). 따라서 아핀인 $X_f$
($f\in A$)들의 합집합은 $X$의 모든 닫힌점을 포함하는 열린집합
$Z$이다~; $X$가 준콤팩트 콜모고로프 공간이므로 반드시 $Z=X$이다
(\hyperref[0.2.1.3-ko]{0, 2.1.3}). $X$가 준콤팩트이므로, $X_{f_i}$들이
아핀이고 $X$를 덮도록 하는 유한 개의 원소 $f_i\in A$
($1\leq i\leq n$)가 존재한다. 이제 단면 $f_i$들로 정의되는 준동형
$\mathcal{O}_X^n\to\mathcal{O}_X$를 생각하자
(\hyperref[0.5.1.1-ko]{0, 5.1.1})~; 모든 $x\in X$에 대하여
$(f_i)_x$ 가운데 적어도 하나가 가역이므로 이 준동형은 전사이고,
따라서 완전열
$0\to\mathcal{R}\to\mathcal{O}_X^n\to\mathcal{O}_X\to0$이 있다.
여기서 $\mathcal{R}$은 $\mathcal{O}_X^n$의 준연접
부분-$\mathcal{O}_X$-가군이다. 따라서

\oldpage[II]{98}

c')에 의해 대응하는 준동형
$\Gamma(X,\mathcal{O}_X^n)\to\Gamma(X,\mathcal{O}_X)$는 전사이고,
이것이 b)를 증명한다.

끝으로 a)는 d)를 함의하고 (\hyperref[I.5.1.9.2-ko]{I, 5.1.9.2}),
d)는 자명하게 d')을 함의한다. d')이 c')을 함의함을 보이는 일만
남았다. $\mathcal{F}'$이 $\mathcal{O}_X^n$의 준연접
부분-$\mathcal{O}_X$-가군이면, 여과
$0\subset\mathcal{O}_X\subset\mathcal{O}_X^2\subset\cdots\subset
\mathcal{O}_X^n$은 $\mathcal{F}'$ 위에
$\mathcal{F}'_k=\mathcal{O}_X^k\cap\mathcal{F}'$ ($0\leq k\leq n$)로
이루어진 여과를 정의한다. 이들은 준연접 $\mathcal{O}_X$-가군이고
(\hyperref[I.4.1.1-ko]{I, 4.1.1}),
$\mathcal{F}'_{k+1}/\mathcal{F}'_k$는
$\mathcal{O}_X^{k+1}/\mathcal{O}_X^k=\mathcal{O}_X$의 준연접
부분-$\mathcal{O}_X$-가군, 즉 $\mathcal{O}_X$의 준연접 아이디얼층과
동형이다. 따라서 가정 d')은
$\mathrm{H}^1(X,\mathcal{F}'_{k+1}/\mathcal{F}'_k)=0$을 함의한다~;
그러면 코호몰로지 완전열
$\mathrm{H}^1(X,\mathcal{F}'_k)\to\mathrm{H}^1(X,\mathcal{F}'_{k+1})
\to\mathrm{H}^1(X,\mathcal{F}'_{k+1}/\mathcal{F}'_k)=0$을 사용하여,
$k$에 관한 귀납법으로 모든 $k$에 대하여
$\mathrm{H}^1(X,\mathcal{F}'_k)=0$임을 증명할 수 있다. C.Q.F.D.

\begin{remark}[5.2.1.1]
\label{II.5.2.1.1-ko}
$X$가 뇌터 준스킴이면, c')과 d')의 서술에서 ``준연접''을 ``연접''으로
바꿀 수 있다. 실제로 c')이 b)를 함의한다는 증명에서
$\mathcal{J}$와 $\mathcal{J}'$은 연접 아이디얼층이고, 한편 연접 가군의
모든 준연접 부분가군은 연접이다
(\hyperref[I.6.1.1-ko]{I, 6.1.1})~; 이로부터 결론이 따른다.
\end{remark}

\begin{corollary}[5.2.2]
\label{II.5.2.2-ko}
$f:X\to Y$를 분리된 준콤팩트 사상이라 하자. 다음 조건들은
동치이다~:
\begin{enumerate}
  \item[\rm{a)}] $f$는 아핀 사상이다.
  \item[\rm{b)}] 함자 $f_*$는 준연접
  $\mathcal{O}_X$-가군들의 범주에서 완전하다.
  \item[\rm{c)}] 모든 준연접 $\mathcal{O}_X$-가군
  $\mathcal{F}$에 대하여 $\mathrm{R}^1f_*(\mathcal{F})=0$이다.
  \item[\rm{c')}] $\mathcal{O}_X$의 모든 준연접 아이디얼층
  $\mathcal{J}$에 대하여 $\mathrm{R}^1f_*(\mathcal{J})=0$이다.
\end{enumerate}
\end{corollary}

이 조건들은 모두 $Y$ 위에서 국소적이므로, 함자
$\mathrm{R}^1f_*$의 정의 (T, 3.7.3)에 따라 $Y$가 아핀이라고
가정할 수 있다. $f$가 아핀이면 $X$도 아핀이고, 성질 b)는 바로
(\hyperref[I.1.6.4-ko]{I, 1.6.4})이다. 반대로 b)가 a)를 함의함을
보이자~: 모든 준연접 $\mathcal{O}_X$-가군 $\mathcal{F}$에 대하여
$f_*(\mathcal{F})$는 준연접 $\mathcal{O}_Y$-가군이다
(\hyperref[I.9.2.2-ko]{I, 9.2.2, a)}). 가정에 따라 함자
$\mathcal{F}\mapsto f_*(\mathcal{F})$는 $\mathcal{F}$에 관해 완전하고,
$Y$가 아핀이므로 함자 $\mathcal{G}\mapsto\Gamma(Y,\mathcal{G})$는
$\mathcal{G}$에 관해 (준연접 $\mathcal{O}_Y$-가군들의 범주에서)
완전하다 (\hyperref[I.1.3.11-ko]{I, 1.3.11})~; 따라서
$\Gamma(Y,f_*(\mathcal{F}))=\Gamma(X,\mathcal{F})$는
$\mathcal{F}$에 관해 완전하고, 이는
(\hyperref[II.5.2.1-ko]{5.2.1, c)})에 의해 우리의 주장을 증명한다.

$f$가 아핀이면 $Y$의 모든 아핀 열린집합 $U$에 대하여
$f^{-1}(U)$는 아핀이다 (\hyperref[II.1.3.2-ko]{1.3.2}). 따라서
$\mathrm{H}^1(f^{-1}(U),\mathcal{F})=0$
(\hyperref[II.5.2.1-ko]{5.2.1, d)})이고, 이는 정의에 따라
$\mathrm{R}^1f_*(\mathcal{F})=0$을 함의한다. 끝으로 조건 c')이
성립한다고 가정하자~; 르레 스펙트럼 열의 낮은 차수 항들로 이루어진
완전열 (G, II, 4.17.1 및 I, 4.5.1)은 특히 다음 완전열을 준다.
\[
  0\to\mathrm{H}^1(Y,f_*(\mathcal{J}))\to
  \mathrm{H}^1(X,\mathcal{J})\to
  \mathrm{H}^0(Y,\mathrm{R}^1f_*(\mathcal{J})).
\]
$Y$는 아핀이고 $f_*(\mathcal{J})$는 준연접이므로
(\hyperref[I.9.2.2-ko]{I, 9.2.2, a)}),
$\mathrm{H}^1(Y,f_*(\mathcal{J}))=0$이다
(\hyperref[II.5.2.1-ko]{5.2.1})~; 따라서 가정 c')에 의해
$\mathrm{H}^1(X,\mathcal{J})=0$이고,
(\hyperref[II.5.2.1-ko]{5.2.1})에 따라 $X$가 아핀 스킴이라는
결론을 얻는다.

\begin{corollary}[5.2.3]
\label{II.5.2.3-ko}
$f:X\to Y$가 아핀 사상이면, 모든 준연접
$\mathcal{O}_X$-가군 $\mathcal{F}$에 대하여 표준 준동형
$\mathrm{H}^1(Y,f_*(\mathcal{F}))\to\mathrm{H}^1(X,\mathcal{F})$은
전단사이다.
\end{corollary}

\oldpage[II]{99}

실제로 르레 스펙트럼 열의 낮은 차수 항들로 이루어진 완전열
\[
  0\to\mathrm{H}^1(Y,f_*(\mathcal{F}))\to
  \mathrm{H}^1(X,\mathcal{F})\to
  \mathrm{H}^0(Y,\mathrm{R}^1f_*(\mathcal{F}))
\]
이 있고, 결론은 (\hyperref[II.5.2.2-ko]{5.2.2})에서 따른다.

\begin{remark}[5.2.4]
\label{II.5.2.4-ko}
제III장 § 1에서 우리는 $X$가 아핀이면 모든 $i>0$과 모든 준연접
$\mathcal{O}_X$-가군 $\mathcal{F}$에 대하여
$\mathrm{H}^i(X,\mathcal{F})=0$임을 증명할 것이다.
\end{remark}

\subsection{준사영 사상}
\label{subsection:II.5.3-ko}

\begin{definition}[5.3.1]
\label{II.5.3.1-ko}
사상 $f:X\to Y$가 \emph{준사영}이라는 것은, 또는 $X$를 $f$를
통하여 $Y$-준스킴으로 보았을 때 $X$가 \emph{$Y$ 위에서 준사영}이라는
것은, 또는 $X$가 \emph{준사영 $Y$-스킴}이라는 것은, $f$가 유한형이고
$f$-풍부한 가역 $\mathcal{O}_X$-가군이 존재한다는 뜻이다.
\end{definition}

이 개념은 $Y$ 위에서 국소적이지 않음에 주의해야 한다~: Nagata [26]와
Hironaka의 반례들은, $X$와 $Y$가 대수적으로 닫힌 체 위의 비특이 대수
스킴인 경우에도, $Y$의 모든 점이 아핀 근방 $U$를 가져
$f^{-1}(U)$가 $U$ 위에서 준사영이지만 $f$ 자체는 준사영이 아닐 수
있음을 보여 준다.

준사영 사상은 반드시 분리 사상임에 주의하자
(\hyperref[II.4.6.1-ko]{4.6.1}). $Y$가 준콤팩트일 때에는 $f$가
준사영이라고 말하는 것과, $f$가 유한형이고 $f$에 관해 매우 풍부한
$\mathcal{O}_X$-가군이 존재한다고 말하는 것이 같다
(\hyperref[II.4.6.2-ko]{4.6.2} 및
\hyperref[II.4.6.11-ko]{4.6.11}). 더 나아가~:

\begin{proposition}[5.3.2]
\label{II.5.3.2-ko}
$Y$를 준콤팩트 스킴 또는 그 바탕공간이 뇌터인 준스킴이라 하고,
$X$를 $Y$-준스킴이라 하자. 다음 조건들은 동치이다~:
\begin{enumerate}
  \item[\rm{a)}] $X$는 준사영 $Y$-스킴이다.
  \item[\rm{b)}] $X$는 $Y$ 위에서 유한형이고, 유한형 준연접
  $\mathcal{O}_Y$-가군 $\mathcal{E}$가 존재하여 $X$는
  $\mathbf{P}(\mathcal{E})$의 한 부분준스킴과 $Y$-동형이다.
  \item[\rm{c)}] $X$는 $Y$ 위에서 유한형이고, 준연접 등급
  $\mathcal{O}_Y$-대수 $\mathcal{S}$가 존재하여 $\mathcal{S}_1$은
  유한형이고 $\mathcal{S}$를 생성하며, $X$는
  $\operatorname{Proj}(\mathcal{S})$의 모든 곳에서 조밀한 한 열린집합
  위에 유도된 부분준스킴과 $Y$-동형이다.
\end{enumerate}
\end{proposition}

이는 앞의 주의와 (\hyperref[II.4.4.3-ko]{4.4.3}),
(\hyperref[II.4.4.6-ko]{4.4.6}) 및
(\hyperref[II.4.4.7-ko]{4.4.7})에서 곧바로 따른다.

$Y$가 뇌터 준스킴이면, (\hyperref[II.5.3.2-ko]{5.3.2})의 조건 b)와
c)에서 $X$가 $Y$ 위에서 유한형이라는 가정을 생략할 수 있음을
주목하자. 이 가정은 이때 자동으로 성립한다
(\hyperref[I.6.3.5-ko]{I, 6.3.5}).

\begin{corollary}[5.3.3]
\label{II.5.3.3-ko}
$Y$를 풍부한 $\mathcal{O}_Y$-가군 $\mathcal{L}$이 존재하는 준콤팩트
스킴이라 하자 (\hyperref[II.4.5.3-ko]{4.5.3}). $Y$-스킴 $X$가
준사영일 필요충분조건은 $X$가 $Y$ 위에서 유한형이고
$\mathbf{P}_Y^r$ 꼴의 사영 다발의 한 부분-$Y$-스킴과 동형인 것이다.
\end{corollary}

실제로 $\mathcal{E}$가 유한형 준연접 $\mathcal{O}_Y$-가군이면,
$\mathcal{E}$는 $\mathcal{O}_Y$-가군
$\mathcal{L}^{\otimes(-n)}\otimes_{\mathcal{O}_Y}\mathcal{O}_Y^k$의
한 몫과 동형이고 (\hyperref[II.4.5.5-ko]{4.5.5}), 따라서
$\mathbf{P}(\mathcal{E})$는 $\mathbf{P}_Y^{k-1}$의 닫힌 부분스킴과
동형이다 (\hyperref[II.4.1.2-ko]{4.1.2}와
\hyperref[II.4.1.4-ko]{4.1.4}).

\begin{proposition}[5.3.4]
\label{II.5.3.4-ko}
\begin{enumerate}
  \item[\rm{(i)}] 유한형 준아핀 사상(특히 준콤팩트 몰입 사상 또는
  유한형 아핀 사상)은 준사영이다.
  \item[\rm{(ii)}] $f:X\to Y$와 $g:Y\to Z$가 준사영 사상이고
  $Z$가 준콤팩트이면, $g\circ f$는 준사영이다.

\oldpage[II]{100}

  \item[\rm{(iii)}] $f:X\to Y$가 준사영 $S$-사상이면, 기저 준스킴의
  모든 확장 $S'\to S$에 대하여
  $f_{(S')}:X_{(S')}\to Y_{(S')}$는 준사영이다.
  \item[\rm{(iv)}] $f:X\to Y$와 $g:X'\to Y'$이 두 준사영
  $S$-사상이면, $f\times_Sg$는 준사영이다.
  \item[\rm{(v)}] $f:X\to Y$, $g:Y\to Z$가 $g\circ f$가 준사영인
  두 사상이고, $g$가 분리 사상이거나 $X$가 국소 뇌터이면, $f$는
  준사영이다.
  \item[\rm{(vi)}] $f$가 준사영 사상이면 $f_{\mathrm{red}}$도
  준사영 사상이다.
\end{enumerate}
\end{proposition}

(i)은 (\hyperref[II.5.1.6-ko]{5.1.6})과
(\hyperref[II.5.1.10-ko]{5.1.10, (i)})에서 따른다. 나머지 주장들은
정의 (\hyperref[II.5.3.1-ko]{5.3.1}), 유한형 사상의 성질들
(\hyperref[I.6.3.4-ko]{I, 6.3.4}) 및
(\hyperref[II.4.6.13-ko]{4.6.13})의 직접적인 귀결이다.

\begin{remark}[5.3.5]
\label{II.5.3.5-ko}
$Y$가 유한 계수 $\mathbf{C}$-대수의 스펙트럼이고 $f$가 고유하다고
가정하더라도, $f_{\mathrm{red}}$는 준사영이지만 $f$는 그렇지 않은
경우가 있을 수 있음을 주목하자.
\end{remark}

\begin{corollary}[5.3.6]
\label{II.5.3.6-ko}
$X$와 $X'$이 두 준사영 $Y$-스킴이면, $X\amalg X'$은 준사영
$Y$-스킴이다.
\end{corollary}

이는 (\hyperref[II.4.6.18-ko]{4.6.18})에서 따른다.

\subsection{고유 사상과 보편적으로 닫힌 사상.}
\label{subsection:II.5.4-ko}
\label{II.5.4-ko}

\begin{definition}[5.4.1]
\label{II.5.4.1-ko}
준스킴의 사상 $f:X\to Y$가 다음 두 조건을 만족할 때 $f$를
\emph{고유} 사상이라고 한다.
\begin{enumerate}
  \item[\rm{a)}] $f$는 분리 사상이고 유한형이다.
  \item[\rm{b)}] 모든 준스킴 $Y'$과 모든 사상 $Y'\to Y$에 대하여,
  사영 $f_{(Y')}:X\times_YY'\to Y'$은 닫힌 사상이다
  (\hyperref[I.2.2.6-ko]{I, 2.2.6}).
\end{enumerate}
이 조건들이 성립할 때에는 $X$도 (구조 사상이 $f$인 $Y$-준스킴으로
보아) $Y$ 위에서 고유라고 한다.
\end{definition}

조건 a)와 b)가 $Y$ 위에서 국소적임은 바로 알 수 있다. $X\times_YY'$의
닫힌 부분집합 $Z$의 사영
$q:X\times_YY'\to Y'$에 의한 상이 $Y'$에서 닫혀 있음을 확인하려면,
$Y'$의 모든 아핀 열린집합 $U'$에 대하여 $q(Z)\cap U'$이 $U'$에서 닫혀
있음을 보이면 충분하다. 그런데
$q(Z)\cap U'=q(Z\cap q^{-1}(U'))$이고 $q^{-1}(U')$은
$X\times_YU'$과 동일시되므로
(\hyperref[I.4.4.1-ko]{I, 4.4.1}), 정의
(\hyperref[II.5.4.1-ko]{5.4.1})의 조건 b)를 확인할 때에는 $Y'$이 아핀
스킴인 경우로 한정할 수 있다. 뒤에서
(\hyperref[II.5.6.3-ko]{5.6.3}) $Y$가 국소 뇌터이면 $Y'$이 $Y$ 위에서
유한형인 경우로까지 한정할 수 있음을 보일 것이다.

모든 고유 사상이 \emph{닫힌} 사상임은 명백하다.

\begin{proposition}[5.4.2]
\label{II.5.4.2-ko}
\begin{enumerate}
  \item[\rm{(i)}] 닫힌 몰입 사상은 고유 사상이다.
  \item[\rm{(ii)}] 두 고유 사상의 합성은 고유 사상이다.
  \item[\rm{(iii)}] $X$, $Y$가 $S$-준스킴이고
  $f:X\to Y$가 고유 $S$-사상이면, 밑준스킴의 모든 확장
  $S'\to S$에 대하여 $f_{(S')}:X_{(S')}\to Y_{(S')}$도 고유 사상이다.
  \item[\rm{(iv)}] $f:X\to Y$와 $g:X'\to Y'$가 두 고유 $S$-사상이면,
  $S$-사상
  $f\times_Sg:X\times_SX'\to Y\times_SY'$도 고유 사상이다.
\end{enumerate}
\end{proposition}

\oldpage[II]{101}

(i), (ii), (iii)만 증명하면 충분하다
(\hyperref[I.3.5.1-ko]{I, 3.5.1}). 이 세 경우 각각에서
(\hyperref[II.5.4.1-ko]{5.4.1})의 조건 a)는 앞에서 얻은 결과들
(\hyperref[I.5.5.1-ko]{I, 5.5.1}과
\hyperref[I.6.3.4-ko]{6.3.4})로부터 따르므로, 조건 b)만 확인하면 된다.
(i)의 경우에는 이것이 곧바로 성립한다. 실제로 $X\to Y$가 닫힌 몰입
사상이면 $X\times_YY'\to Y\times_YY'=Y'$도 닫힌 몰입 사상이기 때문이다
(\hyperref[I.4.3.2-ko]{I, 4.3.2}와
\hyperref[II.3.3.3-ko]{3.3.3}). (ii)를 증명하기 위해 두 고유 사상
$X\to Y$, $Y\to Z$와 사상 $Z'\to Z$를 생각하자. 다음과 같이 쓸 수 있다.
$X\times_ZZ'=X\times_Y(Y\times_ZZ')$
(\hyperref[I.3.3.9.1-ko]{I, 3.3.9.1}). 따라서 사영
$X\times_ZZ'\to Z'$은
$X\times_Y(Y\times_ZZ')\to Y\times_ZZ'\to Z'$로 분해된다. 처음의
관찰과 두 닫힌 사상의 합성이 닫힌 사상이라는 사실로부터 (ii)가 따른다.
마지막으로 모든 사상 $S'\to S$에 대하여 $X_{(S')}$은
$X\times_YY_{(S')}$과 동일시된다
(\hyperref[I.3.3.11-ko]{I, 3.3.11}). 모든 사상
$Z\to Y_{(S')}$에 대하여 다음과 같이 쓸 수 있다.
\[
  X_{(S')}\times_{Y_{(S')}}Z
  =(X\times_YY_{(S')})\times_{Y_{(S')}}Z=X\times_YZ~;
\]
가정에 따라 $X\times_YZ\to Z$는 닫힌 사상이므로, 이것으로 (iii)가
증명된다.

\begin{corollary}[5.4.3]
\label{II.5.4.3-ko}
$f:X\to Y$, $g:Y\to Z$를 $g\circ f$가 고유 사상이 되는 두 사상이라
하자.
\begin{enumerate}
  \item[\rm{(i)}] $g$가 분리 사상이면 $f$는 고유 사상이다.
  \item[\rm{(ii)}] $g$가 분리 사상이고 유한형이며 $f$가 전사이면,
  $g$는 고유 사상이다.
\end{enumerate}
\end{corollary}

(i)는 일반 절차
(\hyperref[I.5.5.12-ko]{I, 5.5.12})를
(\hyperref[II.5.4.2-ko]{5.4.2})에 적용하면 따른다. (ii)를 증명할
때에는 정의 (\hyperref[II.5.4.1-ko]{5.4.1})의 조건 b)만 확인하면 된다.
모든 사상 $Z'\to Z$에 대하여 도식
\[
  \xymatrix{
    X\times_ZZ'\ar[r]^{f\times1_{Z'}}\ar[dr]_p &
    Y\times_ZZ'\ar[d]^{p'}\\
    & Z'
  }
\]
은 가환한다(여기서 $p$와 $p'$은 사영이다)
(\hyperref[I.3.2.1-ko]{I, 3.2.1}). 또한 $f$가 전사이므로
$f\times1_{Z'}$도 전사이고
(\hyperref[I.3.5.2-ko]{I, 3.5.2}), 가정에 따라 $p$는 닫힌 사상이다.
그러므로 $Y\times_ZZ'$의 모든 닫힌 부분집합 $F$는
$X\times_ZZ'$의 어떤 닫힌 부분집합 $E$의 $f\times1_{Z'}$에 의한 상이고,
따라서 $p'(F)=p(E)$는 가정에 따라 $Z'$에서 닫혀 있다. 이로써 따름정리가
증명된다.

\begin{corollary}[5.4.4]
\label{II.5.4.4-ko}
$X$가 $Y$ 위에서 고유인 준스킴이고 $mathcal{S}$가 준연접
$\mathcal{O}_Y$-대수이면, 모든 $Y$-사상
$f:X\to\operatorname{Proj}(\mathcal{S})$은 고유 사상이다(따라서 특히
닫힌 사상이다).
\end{corollary}

실제로 구조 사상 $p:\operatorname{Proj}(\mathcal{S})\to Y$는 분리
사상이고, 가정에 따라 $p\circ f$는 고유 사상이다.

\begin{corollary}[5.4.5]
\label{II.5.4.5-ko}
$f:X\to Y$를 분리 사상이고 유한형인 사상이라 하자.
$(X_i)_{1\leq i\leq n}$ (각각 $(Y_i)_{1\leq i\leq n}$)을 $X$ (각각
$Y$)의 닫힌 부분준스킴들의 유한족, $j_i$ (각각 $h_i$)를 표준 포함 사상
$X_i\to X$ (각각 $Y_i\to Y$)라 하자. $X$의 바탕공간이 $X_i$들의
합집합이고, 모든 $i$에 대하여 도식
\[
  \xymatrix{
    X_i\ar[r]^{f_i}\ar[d]_{j_i} & Y_i\ar[d]^{h_i}\\
    X\ar[r]_f & Y
  }
\]
을 가환이 되게 하는 사상 $f_i:X_i\to Y_i$가 존재한다고 가정하자.
그러면 $f$가 고유 사상일 필요충분조건은 각 $f_i$가 고유 사상인 것이다.
\end{corollary}

\oldpage[II]{102}

$f$가 고유 사상이면 $j_i$가 닫힌 몰입 사상이므로
(\hyperref[II.5.4.2-ko]{5.4.2}) $f\circ j_i$도 고유 사상이다~; $h_i$는
닫힌 몰입 사상이고 따라서 분리 사상이므로, $f_i$는
(\hyperref[II.5.4.3-ko]{5.4.3})에 따라 고유 사상이다. 역으로 각
$f_i$가 고유 사상이라고 가정하고, $X_i$들의 합인 준스킴 $Z$와 각
$X_i$에서 $j_i$와 같은 사상 $u:Z\to X$를 생각하자. $f\circ u$의 각
$X_i$에 대한 제한은 $f\circ j_i=h_i\circ f_i$와 같으므로, $h_i$와
$f_i$가 모두 고유 사상이라는 사실에서
(\hyperref[II.5.4.2-ko]{5.4.2}) 고유 사상이다~; 그러면 정의
(\hyperref[II.5.4.1-ko]{5.4.1})에서 곧바로 $f\circ u$가 고유 사상임이
따른다. 그런데 가정에 따라 $u$는 전사이므로,
(\hyperref[II.5.4.3-ko]{5.4.3})에 따라 $f$가 고유 사상이라는 결론을
얻는다.

\begin{corollary}[5.4.6]
\label{II.5.4.6-ko}
$f:X\to Y$를 분리 사상이고 유한형인 사상이라 하자~; $f$가 고유
사상일 필요충분조건은
$f_{\mathrm{red}}:X_{\mathrm{red}}\to Y_{\mathrm{red}}$가 고유 사상인
것이다.
\end{corollary}

이는 $n=1$, $X_1=X_{\mathrm{red}}$, $Y_1=Y_{\mathrm{red}}$인
(\hyperref[II.5.4.5-ko]{5.4.5})의 특별한 경우이다
(\hyperref[I.5.1.5-ko]{I, 5.1.5}).

\begin{env}[5.4.7]
\label{II.5.4.7-ko}
$X$와 $Y$가 뇌터 준스킴이고 $f:X\to Y$가 분리 사상이면서 유한형이면,
$f$가 고유 사상임을 확인하기 위해 정역 준스킴들의 우세 사상의 경우로
환원할 수 있다. 실제로 $X_i$ ($1\leq i\leq n$)를 $X$의 유한 개의
기약성분이라 하고, 각 $i$에 대하여 바탕공간이 $X_i$인 $X$의 유일한
축소 닫힌 부분준스킴을 생각하여, 이것도 $X_i$로 나타내자
(\hyperref[I.5.2.1-ko]{I, 5.2.1}). 한편 $Y_i$를 바탕공간이
$\overline{f(X_i)}$인 $Y$의 유일한 축소 닫힌 부분준스킴이라 하자.
$g_i$ (각각 $h_i$)가 포함 사상 $X_i\to X$ (각각 $Y_i\to Y$)이면,
$f\circ g_i=h_i\circ f_i$임을 얻는데, 여기서 $f_i$는 우세 사상
$X_i\to Y_i$이다 (\hyperref[I.5.2.2-ko]{I, 5.2.2})~; 따라서
(\hyperref[II.5.4.5-ko]{5.4.5})를 적용할 수 있는 조건들이 성립하고,
$f$가 고유 사상일 필요충분조건은 $f_i$들이 고유 사상인 것이다.
\end{env}

\begin{corollary}[5.4.8]
\label{II.5.4.8-ko}
$X$, $Y$를 $S$ 위의 분리 유한형 $S$-준스킴이라 하고,
$f:X\to Y$를 $S$-사상이라 하자. $f$가 고유 사상일 필요충분조건은
모든 $S$-준스킴 $S'$에 대하여 사상
$f\times_S1_{S'}:X\times_SS'\to Y\times_SS'$가 닫힌 사상인 것이다.
\end{corollary}

먼저 $g:X\to S$와 $h:Y\to S$가 구조 사상이면 정의에 따라
$g=h\circ f$이므로, $f$는 분리 사상이고 유한형임을 주목하자
(\hyperref[I.5.5.1-ko]{I, 5.5.1}과
\hyperref[I.6.3.4-ko]{6.3.4}). $f$가 고유 사상이면
$f\times_S1_{S'}$도 그러하다 (\hyperref[II.5.4.2-ko]{5.4.2})~;
특히 $f\times_S1_{S'}$는 닫힌 사상이다. 역으로 명제의 조건이
성립한다고 가정하고 $Y'$을 $Y$-준스킴이라 하자~; $Y'$은
$S$-준스킴으로도 볼 수 있고, $Y\to S$가 분리 사상이므로
$X\times_YY'$은 $X\times_SY'$의 닫힌 부분준스킴과 동일시된다
(\hyperref[I.5.4.2-ko]{I, 5.4.2}). 가환 도식
\[
  \xymatrix{
    X\times_YY'\ar[r]^{f\times1_{Y'}}\ar[d] &
    Y\times_YY'=Y'\ar[d]\\
    X\times_SY'\ar[r]_{f\times_S1_{Y'}} & Y\times_SY'
  }
\]
에서 수직 화살표들은 닫힌 몰입 사상이다~; 따라서
$f\times_S1_{Y'}$가 닫힌 사상이면 $f\times_Y1_{Y'}$도 닫힌 사상임이
곧바로 따른다.

\begin{remark}[5.4.9]
\label{II.5.4.9-ko}
사상 $f:X\to Y$가 정의 (\hyperref[II.5.4.1-ko]{5.4.1})의 조건 b)를
만족하면 $f$를 \emph{보편적으로 닫힌} 사상이라 하겠다. 독자는

\oldpage[II]{103}

명제 (\hyperref[II.5.4.2-ko]{5.4.2})부터
(\hyperref[II.5.4.8-ko]{5.4.8})까지 어디서나 결과의 타당성을 바꾸지
않고 ``고유''를 ``보편적으로 닫힌''으로 바꿀 수 있음을 알 수 있을
것이다(그리고 (\hyperref[II.5.4.3-ko]{5.4.3}),
(\hyperref[II.5.4.5-ko]{5.4.5}),
(\hyperref[II.5.4.6-ko]{5.4.6}),
(\hyperref[II.5.4.8-ko]{5.4.8})의 가정에서는 유한성 조건들을 생략할
수 있다).
\end{remark}

\begin{env}[5.4.10]
\label{II.5.4.10-ko}
$f:X\to Y$를 유한형 사상이라 하자. $X$의 닫힌 부분집합 $Z$에 대하여,
바탕공간이 $Z$인 $X$의 닫힌 부분준스킴으로 제한한 $f$가 고유 사상이면
$Z$는 \emph{$Y$ 위에서 고유하다} (또는 \emph{$Y$-고유하다}, 또는
\emph{$f$에 관해 고유하다})고 한다
(\hyperref[I.5.2.1-ko]{I, 5.2.1}). 이 제한 사상은 이때 분리 사상이므로,
(\hyperref[II.5.4.6-ko]{5.4.6})과
(\hyperref[I.5.5.1-ko]{I, 5.5.1, (vi)})에서 앞의 성질은 바탕공간이
$Z$인 $X$의 닫힌 부분준스킴의 선택에 의존하지 않음이 따른다.
$g:X'\to X$가 고유 사상이면, $g^{-1}(Z)$는 $X'$의 $Y$-고유
부분집합이다~: 실제로 바탕공간이 $Z$인 $X$의 부분준스킴을 $T$라 하면,
$X'$의 닫힌 부분준스킴 $g^{-1}(T)$로 제한한 $g$는
(\hyperref[II.5.4.2-ko]{5.4.2, (iii)})에 따라 고유 사상
$g^{-1}(T)\to T$이고, 이어서
(\hyperref[II.5.4.2-ko]{5.4.2, (ii)})를 적용하면 충분하다. 한편 $X''$이
유한형 $Y$-스킴이고 $u:X\to X''$이 $Y$-사상이면, $u(Z)$는 $X''$의
$Y$-고유 부분집합이다~; 실제로 바탕공간이 $Z$인 $X$의 축약 닫힌
부분준스킴을 $T$라 하자~; $T$로 제한한 $f$가 고유 사상이므로,
(\hyperref[II.5.4.3-ko]{5.4.3, (i)})에 따라 $T$로 제한한 $u$도 고유
사상이고, 따라서 $u(Z)$는 $X''$에서 닫혀 있다~; 바탕공간이 $u(Z)$인
$X''$의 닫힌 부분준스킴을 $T''$이라 하면
(\hyperref[I.5.2.1-ko]{I, 5.2.1}), $u|T$는 다음과 같이 분해된다.
\[
  T\xrightarrow{v}T''\xrightarrow{j}X'',
\]
여기서 $j$는 표준 포함 사상이고
(\hyperref[I.5.2.2-ko]{I, 5.2.2}), 따라서 $v$는 고유 전사 사상이다
(\hyperref[II.5.4.5-ko]{5.4.5})~; $g$를 구조 사상 $X''\to Y$의
$T''$로의 제한이라 하면, $g$는 분리 사상이고 유한형이며
$f|T=g\circ v$이다~; 그러므로
(\hyperref[II.5.4.3-ko]{5.4.3, (ii)})에서 $g$가 고유 사상임이 따르고,
이는 주장을 증명한다.
\end{env}

특히 이들 주의로부터, $Z$가 $X$의 $Y$-고유 부분집합이면 다음이
따른다~:
\begin{enumerate}
  \item[\rm{1\textsuperscript{o}}] $X$의 모든 닫힌 부분준스킴 $X'$에
  대하여, $Z\cap X'$은 $X'$의 $Y$-고유 부분집합이다.
  \item[\rm{2\textsuperscript{o}}] $X$가 유한형 $Y$-스킴 $X''$의
  부분준스킴이면, $Z$는 $X''$의 $Y$-고유 부분집합이기도 하다
  (특히 $X''$에서 닫혀 있다).
\end{enumerate}

\subsection{사영 사상.}
\label{subsection:II.5.5-ko}

\begin{proposition}[5.5.1]
\label{II.5.5.1-ko}
$X$를 $Y$-준스킴이라 하자. 다음 조건들은 서로 동치이다~:
\begin{enumerate}
  \item[\rm{a)}] $X$는 어떤 사영 다발 $\mathbf{P}(\mathcal{E})$의 닫힌
  부분준스킴과 $Y$-동형이다. 여기서 $\mathcal{E}$는 유한형 준연접
  $\mathcal{O}_Y$-가군이다.
  \item[\rm{b)}] $\mathcal{S}_1$이 유한형이고 $\mathcal{S}$를 생성하며,
  $X$가 $\operatorname{Proj}(\mathcal{S})$와 $Y$-동형이 되게 하는 준연접
  등급 $\mathcal{O}_Y$-대수 $\mathcal{S}$가 존재한다.
\end{enumerate}
\end{proposition}

(\hyperref[II.3.6.2-ko]{3.6.2, (ii)})에 따라 조건 a)는 b)를 함의한다~:
$\mathcal{J}$가 $\mathbf{S}(\mathcal{E})$의 등급 준연접 아이디얼층이면,
준연접 등급 $\mathcal{O}_Y$-대수
$\mathcal{S}=\mathbf{S}(\mathcal{E})/\mathcal{J}$는
$\mathcal{S}_1$에 의해 생성되고, $\mathcal{E}$의 표준 상인
$\mathcal{S}_1$은 유한형 $\mathcal{O}_Y$-가군이다. 조건 b)는
$\mathcal{M}\to\mathcal{S}_1$이 항등 사상인 경우에
(\hyperref[II.3.6.2-ko]{3.6.2})를 적용하면 a)를 함의한다.

\oldpage[II]{104}

\begin{definition}[5.5.2]
\label{II.5.5.2-ko}
$Y$-준스킴 $X$가 (\hyperref[II.5.5.1-ko]{5.5.1})의 서로 동치인 조건
a)와 b)를 만족하면, $X$는 \emph{$Y$ 위에서 사영적이다} 또는
\emph{사영 $Y$-스킴이다}라고 한다. 사상 $f:X\to Y$가 $X$를 사영
$Y$-스킴으로 만들면 $f$를 \emph{사영 사상}이라 한다.
\end{definition}

$f:X\to Y$가 사영 사상이면 $f$에 관해 \emph{매우 풍부한}
$\mathcal{O}_X$-가군이 존재함은 명백하다
(\hyperref[II.4.4.2-ko]{4.4.2}).

\begin{theorem}[5.5.3]
\label{II.5.5.3-ko}
\begin{enumerate}
  \item[\rm{(i)}] 모든 사영 사상은 준사영 사상이고 고유 사상이다.
  \item[\rm{(ii)}] 역으로, $Y$를 준콤팩트 스킴 또는 그 바탕공간이
  뇌터인 준스킴이라 하자~; 준사영 사상이면서 고유 사상인 모든 사상
  $f:X\to Y$는 사영 사상이다.
\end{enumerate}
\end{theorem}

(i) $f:X\to Y$가 사영 사상이면 유한형이고 준사영 사상이며(따라서 특히
분리 사상이며), 한편 (\hyperref[II.5.5.1-ko]{5.5.1, b)})와
(\hyperref[II.3.5.3-ko]{3.5.3})에서 모든 사상 $Y'\to Y$에 대하여
$f$가 사영 사상이면
$f\times_Y1_{Y'}:X\times_YY'\to Y'$도 사영 사상임이 곧바로 따른다.
따라서 $f$가 보편적으로 닫힌 사상임을 보이기 위해서는 사영 사상 $f$가
\emph{닫힌 사상}임을 증명하면 충분하다. 문제는 $Y$ 위에서 국소적이므로
$Y=\operatorname{Spec}(A)$라고 가정할 수 있고, 따라서
(\hyperref[II.5.5.1-ko]{5.5.1})에 따라
$X=\operatorname{Proj}(S)$이다. 여기서 $S$는 $S_1$의 유한 개의
원소로 생성되는 등급 $A$-대수이다. 모든 $y\in Y$에 대하여 올
$f^{-1}(y)$는
$\operatorname{Proj}(S)\times_Y\operatorname{Spec}(k(y))$와 동일시되고
(\hyperref[I.3.6.1-ko]{I, 3.6.1}), 따라서
$\operatorname{Proj}(S\otimes_Ak(y))$와 동일시된다
(\hyperref[II.2.8.10-ko]{2.8.10})~; 그러므로 $f^{-1}(y)$가 공집합일
필요충분조건은 $S\otimes_Ak(y)$가 조건 (TN)을 만족하는 것이고
(\hyperref[II.2.7.4-ko]{2.7.4}), 다시 말해 충분히 큰 $n$에 대하여
$S_n\otimes_Ak(y)=0$인 것이다. 그런데 $(S_n)_y$는 유한형
$\mathcal{O}_y$-가군이므로, 나카야마 보조정리에 따라 앞의 조건은
충분히 큰 $n$에 대하여 $(S_n)_y=0$이라는 뜻이다. $A$에서 $A$-가군
$S_n$의 소멸자를 $\mathfrak a_n$이라 하면, 앞의 조건은 다시 충분히 큰
$n$에 대하여 $\mathfrak a_n\not\subset\mathfrak j_y$라는 뜻이다
(\hyperref[0.1.7.4-ko]{0, 1.7.4}). 가정에 따라
$S_nS_1=S_{n+1}$이므로 $\mathfrak a_n\subset\mathfrak a_{n+1}$이고,
$\mathfrak a$를 $\mathfrak a_n$들의 합집합이라 하면
$f(X)=V(\mathfrak a)$임을 알 수 있다. 따라서 $f(X)$는 $Y$에서 닫혀
있다. 이제 $X'$이 $X$의 임의의 닫힌 부분집합이면, 바탕공간이 $X'$인
$X$의 닫힌 부분준스킴이 존재하고
(\hyperref[I.5.2.1-ko]{I, 5.2.1}), 사상
$X'\to X\xrightarrow{f}Y$가 사영 사상임은 명백하다
(\hyperref[II.5.5.1-ko]{5.5.1, a)})~; 따라서 $f(X')$은 $Y$에서 닫혀
있다.

(ii) $Y$에 대한 가정과 $f$가 준사영 사상이라는 사실에 따라, 유한형
준연접 $\mathcal{O}_Y$-가군 $\mathcal{E}$와 $Y$-몰입 사상
$j:X\to\mathbf{P}(\mathcal{E})$가 존재한다
(\hyperref[II.5.3.2-ko]{5.3.2}). 그런데 $f$가 고유 사상이므로
(\hyperref[II.5.4.4-ko]{5.4.4})에 따라 $j$는 닫힌 몰입 사상이고,
따라서 $f$는 사영 사상이다.

\par\medskip\noindent\textit{주의 (5.5.4). ---\ }
\label{II.5.5.4-ko}
\begin{enumerate}
  \item[\rm{(i)}] $f:X\to Y$를 다음 조건들을 만족하는 사상이라 하자~:
  1\textsuperscript{o} $f$는 고유 사상이다~; 2\textsuperscript{o}
  $f$에 관해 매우 풍부한 $\mathcal{O}_X$-가군 $\mathcal{L}$이
  존재한다~; 3\textsuperscript{o} 준연접 $\mathcal{O}_Y$-가군
  $\mathcal{E}=f_*(\mathcal{L})$은 유한형이다. 그러면 $f$는 사영
  사상이다~: 실제로 (\hyperref[II.4.4.4-ko]{4.4.4})에 따라 이때
  $Y$-몰입 사상 $r:X\to\mathbf{P}(\mathcal{E})$가 존재하고, $f$가
  고유 사상이므로 $r$는 닫힌 몰입 사상이다
  (\hyperref[II.5.4.4-ko]{5.4.4}). 제III장, \S~3에서 보겠듯이, $Y$가
  국소 뇌터이면 위의 조건 3\textsuperscript{o}은 나머지 두 조건의
  귀결이다. 따라서 이 경우 조건 1\textsuperscript{o}과
  2\textsuperscript{o}은 사영 사상들을 특징짓고, $Y$가 준콤팩트이면
  조건 2\textsuperscript{o}을 $f$에 관해 풍부한
  $\mathcal{O}_X$-가군이 존재한다는 가정으로 바꿀 수 있다
  (\hyperref[II.4.6.11-ko]{4.6.11}).
  \item[\rm{(ii)}] 풍부한 $\mathcal{O}_Y$-가군이 존재하는 준콤팩트
  스킴 $Y$를 생각하자. $Y$-스킴 $X$가 사영적일 필요충분조건은 어떤
  $\mathbf{P}_Y^r$ 꼴의 사영 다발의 닫힌 부분-$Y$-스킴과
  $Y$-동형인 것이다. 이 조건이 충분함은 명백하다.

\oldpage[II]{105}

  역으로, $X$가 $Y$ 위에서 사영적이면 준사영적이고, 따라서 $X$에서
  어떤 $\mathbf{P}_Y^r$로 가는 $Y$-몰입 사상 $j$가 존재하며
  (\hyperref[II.5.3.3-ko]{5.3.3}), 이는
  (\hyperref[II.5.4.4-ko]{5.4.4})와
  (\hyperref[II.5.5.3-ko]{5.5.3})에 따라 닫힌 몰입 사상이다.
  \item[\rm{(iii)}] (\hyperref[II.5.5.3-ko]{5.5.3})의 논법은 모든
  준스킴 $Y$와 모든 정수 $r\geq0$에 대하여 구조 사상
  $\mathbf{P}_Y^r\to Y$가 전사임을 보여 준다. 실제로
  $\mathcal{S}=\mathbf{S}_{\mathcal{O}_Y}(\mathcal{O}_Y^{r+1})$이라
  놓으면
  $\mathcal{S}_y=\mathbf{S}_{k(y)}((k(y))^{r+1})$임은 명백하고
  (\hyperref[II.1.7.3-ko]{1.7.3}), 따라서 모든 $y\in Y$와 모든
  $n\geq0$에 대하여 $(\mathcal{S}_n)_y\neq0$이다.
  \item[\rm{(iv)}] 나가타의 예들 \cite{II-26}에서 준사영 사상이 아닌
  고유 사상들이 존재함이 따른다.
\end{enumerate}

\begin{proposition}[5.5.5]
\label{II.5.5.5-ko}
\begin{enumerate}
  \item[\rm{(i)}] 닫힌 몰입 사상은 사영 사상이다.
  \item[\rm{(ii)}] $f:X\to Y$와 $g:Y\to Z$가 사영 사상이고, $Z$가
  준콤팩트 스킴이거나 그 바탕공간이 뇌터인 준스킴이면 $g\circ f$는
  사영 사상이다.
  \item[\rm{(iii)}] $f:X\to Y$가 사영 $S$-사상이면, 밑준스킴의 모든
  확장 $S'\to S$에 대하여
  $f_{(S')}:X_{(S')}\to Y_{(S')}$는 사영 사상이다.
  \item[\rm{(iv)}] $f:X\to Y$와 $g:X'\to Y'$가 사영 $S$-사상이면,
  $f\times_Sg$도 사영 사상이다.
  \item[\rm{(v)}] $g\circ f$가 사영 사상이고 $g$가 분리 사상이면,
  $f$는 사영 사상이다.
  \item[\rm{(vi)}] $f$가 사영 사상이면 $f_{\mathrm{red}}$도 사영
  사상이다.
\end{enumerate}
\end{proposition}

(i)는 (\hyperref[II.3.1.7-ko]{3.1.7})에서 곧바로 따른다. 여기서는
(ii)의 $Z$에 둔 제약 때문에 (iii)과 (iv)를 따로 증명해야 한다
(\hyperref[I.3.5.1-ko]{I, 3.5.1} 참조). (iii)을 증명할 때에는
$S=Y$인 경우로 환원할 수 있고
(\hyperref[I.3.3.11-ko]{I, 3.3.11}), 그러면 주장은
(\hyperref[II.5.5.1-ko]{5.5.1, b)})와
(\hyperref[II.3.5.3-ko]{3.5.3})에서 곧바로 따른다. (iv)를 증명할
때에는 곧바로 $X=\mathbf{P}(\mathcal{E})$,
$X'=\mathbf{P}(\mathcal{E}')$인 경우로 환원된다. 여기서
$\mathcal{E}$ (각각 $\mathcal{E}'$)는 유한형 준연접
$\mathcal{O}_Y$-가군 (각각 $\mathcal{O}_{Y'}$-가군)이다. $p$, $p'$을
각각 $T=Y\times_SY'$에서 $Y$, $Y'$로 가는 표준 사영이라 하자~;
(\hyperref[II.4.1.3.1-ko]{4.1.3.1})에 따라
$\mathbf{P}(p^*(\mathcal{E}))=\mathbf{P}(\mathcal{E})\times_YT$,
$\mathbf{P}(p'^*(\mathcal{E}'))=\mathbf{P}(\mathcal{E}')\times_{Y'}T$이고,
따라서
\begin{align*}
\mathbf{P}(p^*(\mathcal{E}))\times_T\mathbf{P}(p'^*(\mathcal{E}'))
&=(\mathbf{P}(\mathcal{E})\times_YT)\times_T
(T\times_{Y'}\mathbf{P}(\mathcal{E}'))\\
&=\mathbf{P}(\mathcal{E})\times_Y
(T\times_{Y'}\mathbf{P}(\mathcal{E}'))
=\mathbf{P}(\mathcal{E})\times_S\mathbf{P}(\mathcal{E}')
\end{align*}
를 얻는다. 여기서 $T$를 $Y\times_SY'$로 바꾸고
(\hyperref[I.3.3.9.1-ko]{I, 3.3.9.1})을 사용하였다. 그런데
$p^*(\mathcal{E})$와 $p'^*(\mathcal{E}')$는 $T$ 위에서 유한형이므로
(\hyperref[0.5.2.4-ko]{0, 5.2.4}),
$p^*(\mathcal{E})\otimes_{\mathcal{O}_T}p'^*(\mathcal{E}')$도 유한형이다~;
$\mathbf{P}(p^*(\mathcal{E}))\times_T\mathbf{P}(p'^*(\mathcal{E}'))$은
$\mathbf{P}(p^*(\mathcal{E})\otimes_{\mathcal{O}_T}p'^*(\mathcal{E}'))$의
닫힌 부분준스킴과 동일시되므로
(\hyperref[II.4.3.3-ko]{4.3.3}), 이로써 (iv)의 증명이 끝난다. (v)와
(vi)을 얻기 위해서는 (\hyperref[I.5.5.13-ko]{I, 5.5.13})을 적용할 수
있다. 실제로 사영 $Y$-스킴의 모든 닫힌 부분준스킴은
(\hyperref[II.5.5.1-ko]{5.5.1, a)})에 따라 사영 $Y$-스킴이다.

(ii)만 증명하면 된다~; $Z$에 대한 가정에 비추어 이는
(\hyperref[II.5.5.3-ko]{5.5.3}),
(\hyperref[II.5.3.4-ko]{5.3.4, (ii)}) 및
(\hyperref[II.5.4.2-ko]{5.4.2, (ii)})에서 따른다.

\begin{proposition}[5.5.6]
\label{II.5.5.6-ko}
$X$와 $X'$이 두 사영 $Y$-스킴이면, $X\amalg X'$도 사영
$Y$-스킴이다.
\end{proposition}

이는 (\hyperref[II.5.5.2-ko]{5.5.2})와
(\hyperref[II.4.3.6-ko]{4.3.6})의 명백한 귀결이다.

\begin{proposition}[5.5.7]
\label{II.5.5.7-ko}
$X$를 사영 $Y$-스킴, $\mathcal{L}$을 $Y$에 관해 풍부한
$\mathcal{O}_X$-가군이라 하자~; $X$ 위의 $\mathcal{L}$의 모든 단면
$f$에 대하여 $X_f$는 $Y$ 위에서 아핀이다.
\end{proposition}

\oldpage[II]{106}

문제는 $Y$ 위에서 국소적이므로 $Y=\operatorname{Spec}(A)$라고 가정할
수 있다~; 또한 $X_{f^{\otimes n}}=X_f$이므로 $\mathcal{L}$을 적당한
$\mathcal{L}^{\otimes n}$으로 바꾸어, $\mathcal{L}$이 구조 사상
$q:X\to Y$에 관해 매우 풍부하다고 가정할 수 있다
(\hyperref[II.4.6.11-ko]{4.6.11}). 따라서 표준 준동형
$\sigma:q^*(q_*(\mathcal{L}))\to\mathcal{L}$은 전사이고, 이에 대응하는
사상
\[
r=r_{\mathcal{L},\sigma}:X\to P=\mathbf{P}(q_*(\mathcal{L}))
\]
은 $\mathcal{L}=r^*(\mathcal{O}_P(1))$을 만족하는 몰입 사상이다
(\hyperref[II.4.4.4-ko]{4.4.4})~; 또한 $X$가 $Y$ 위에서 고유하므로
몰입 사상 $r$는 닫혀 있다 (\hyperref[II.5.4.4-ko]{5.4.4}). 한편
정의에 따라 $f\in\Gamma(Y,q_*(\mathcal{L}))$이고
$\sigma^\flat$는 $q_*(\mathcal{L})$의 항등 사상이다~; 그러므로 공식
(\hyperref[II.3.7.3.1-ko]{3.7.3.1})에서
$X_f=r^{-1}(D_+(f))$를 얻는다~; 따라서 $X_f$는 아핀 스킴
$D_+(f)$의 닫힌 부분준스킴이고, 따라서 실제로 아핀 스킴이다.

특히 $Y=X$로 잡으면 (\hyperref[II.4.6.13-ko]{4.6.13, (i)})을
고려하여 다음 따름정리를 얻는다. 한편 이 따름정리의 직접적인 증명도
곧바로 얻어진다~:

\begin{corollary}[5.5.8]
\label{II.5.5.8-ko}
$X$를 준스킴, $\mathcal{L}$을 가역 $\mathcal{O}_X$-가군이라 하자.
$X$ 위의 $\mathcal{L}$의 모든 단면 $f$에 대하여 $X_f$는 $X$ 위에서
아핀이다 (따라서 $X$가 아핀 스킴이면 $X_f$도 아핀 스킴이다).
\end{corollary}

\subsection{차우 보조정리.}
\label{subsection:II.5.6-ko}

\begin{theorem}[5.6.1]
\label{II.5.6.1-ko}
(저우의 보조정리.) $S$를 준스킴, $X$를 유한형 $S$-스킴이라 하자.
다음 가정들 가운데 하나가 성립한다고 하자~:
\begin{enumerate}
  \item[\rm{a)}] $S$는 뇌터이다.
  \item[\rm{b)}] $S$는 준콤팩트 스킴이고, $X$의 기약성분은 유한 개이다.
\end{enumerate}
이 조건들 아래에서 다음이 성립한다~:
\begin{enumerate}
  \item[\rm{(i)}] $S$ 위에서 준사영인 $S$-스킴 $X'$과 사영 전사
  $S$-사상 $f:X'\to X$가 존재한다.
  \item[\rm{(ii)}] 어떤 열린집합 $U\subset X$에 대하여
  $U'=f^{-1}(U)$가 $X'$에서 조밀하고 $f$의 $U'$로의 제한이 동형사상
  $U'\simeq U$이 되도록 $X'$과 $f$를 택할 수 있다.
  \item[\rm{(iii)}] $X$가 축소 스킴(각각 기약 스킴, 정역 스킴)이면,
  $X'$도 축소 스킴(각각 기약 스킴, 정역 스킴)이라고 가정할 수 있다.
\end{enumerate}
\end{theorem}

증명은 여러 단계로 이루어진다.

A) 먼저 $X$가 기약인 경우로 환원할 수 있다. 실제로 가정 a)에서
$X$는 뇌터이므로, 두 가정 어느 쪽에서도 $X$의 기약성분 $X_i$들은
유한 개이다. 바탕공간이 $X_i$인 $X$의 축소 닫힌 부분준스킴 각각에
대하여 정리가 증명되었고, $X'_i$와 $f_i:X'_i\to X_i$가 $X_i$에
대응하는 준스킴과 사상이라고 하자. $X'_i$들의 합인 준스킴 $X'$과,
각 $X'_i$에 대한 제한이 $j_i\circ f_i$인 사상 $f:X'\to X$
($j_i$는 표준 단사 사상 $X_i\to X$이다)은 요구를 만족한다. 실제로
$X'_i$들이 축소이면 $X'$도 축소임은 자명하다~; 한편 $U$를 집합들
$U_i\cap(\bigcup_{j\neq i}X_j)^{\complement}$의 합집합으로 택하면 (ii)가
성립한다. 끝으로 $X'_i$들이 $S$ 위에서 준사영이므로 $X'$도 그러하다

\oldpage[II]{107}

(\hyperref[II.5.3.6-ko]{5.3.6})~; 마찬가지로 사상 $X'_i\to X$들은
(\hyperref[II.5.5.5-ko]{5.5.5, (i)와 (ii)})에 따라 사영이고, 따라서
$f$는 사영이며 (\hyperref[II.5.5.6-ko]{5.5.6}), 정의에 따라 명백히
전사이다.

B) 이제 $X$가 기약이라고 하자. 구조 사상 $r:X\to S$가 유한형이므로,
$S$의 아핀 열린집합들로 이루어진 유한 덮개 $(S_i)$가 있고, 각 $i$에
대하여 $r^{-1}(S_i)$의 아핀 열린집합들로 이루어진 유한 덮개
$(T_{ij})$가 있다. 사상 $T_{ij}\to S_i$는 아핀이고 유한형이므로
준사영이다 (\hyperref[II.5.3.4-ko]{5.3.4, (i)})~; 가정 a), b) 어느
쪽에서도 몰입 사상 $S_i\to S$는 준콤팩트이므로 준사영 사상이고
(\hyperref[II.5.3.4-ko]{5.3.4, (ii)}), 따라서 $r$의 $T_{ij}$로의 제한은
준사영 사상이다 (\hyperref[II.5.3.4-ko]{5.3.4, (iii)}). $T_{ij}$들을
$U_k$ ($1\leq k\leq n$)로 나타내자. 각 지표 $k$에 대하여 열린 몰입
사상 $\varphi_k:U_k\to P_k$가 존재하는데, 여기서 $P_k$는 $S$ 위에서
사영이다 (\hyperref[II.5.3.2-ko]{5.3.2}와
\hyperref[II.5.5.2-ko]{5.5.2}). $U=\bigcap_kU_k$라 하자~; $X$가
기약이고 $U_k$들이 공집합이 아니므로 $U$도 공집합이 아니며, 따라서
$X$에서 어디서나 조밀하다~; $\varphi_k$들의 $U$로의 제한들은 사상
\[
\varphi:U\to P=P_1\times_SP_2\times_S\cdots\times_SP_n
\]
을 정의하여 도식들
\[
\tag{5.6.1.1}
\label{II.5.6.1.1-ko}
\xymatrix{
U\ar[r]^\varphi\ar[d]_{j_k} & P\ar[d]^{p_k}\\
U_k\ar[r]_{\varphi_k} & P_k
}
\]
이 가환하게 한다. 여기서 $j_k$는 표준 단사 사상 $U\to U_k$이고
$p_k$는 표준 사영 $P\to P_k$이다. $j$가 표준 단사 사상 $U\to X$이면,
사상 $\psi=(j,\varphi)_S:U\to X\times_SP$는 몰입 사상이다
(\hyperref[I.5.3.14-ko]{I, 5.3.14}). 가정 a)에서 $X\times_SP$는 국소
뇌터이고 (\hyperref[II.3.4.1-ko]{3.4.1}과
\hyperref[I.6.3.7-ko]{I, 6.3.7과 6.3.8}), 가정 b)에서
$X\times_SP$는 준콤팩트 스킴이다
(\hyperref[I.5.5.1-ko]{I, 5.5.1과 6.6.4})~; 두 경우 모두 $\psi$에
결부된 부분준스킴 $Z$ (그 바탕공간은 $\psi(U)$이다)의
$X\times_SP$ 안에서의 폐포 $X'$이 존재하고, $\psi$는 다음과 같이
분해된다.
\[
\tag{5.6.1.2}
\label{II.5.6.1.2-ko}
\psi:U\xrightarrow{\psi'}X'\xrightarrow{h}X\times_SP
\]
여기서 $\psi'$은 열린 몰입 사상이고 $h$는 닫힌 몰입 사상이다
(\hyperref[I.9.5.10-ko]{I, 9.5.10}). $q_1:X\times_SP\to X$와
$q_2:X\times_SP\to P$를 표준 사영이라 하자~; 다음과 같이 놓겠다.
\[
\tag{5.6.1.3}
\label{II.5.6.1.3-ko}
f:X'\xrightarrow{h}X\times_SP\xrightarrow{q_1}X
\]
그리고
\[
\tag{5.6.1.4}
\label{II.5.6.1.4-ko}
g:X'\xrightarrow{h}X\times_SP\xrightarrow{q_2}P.
\]
$X'$과 $f$가 요구를 만족함을 보이겠다.

C) 먼저 $f$가 사영 전사 사상이고 $f$의 $U'=f^{-1}(U)$로의 제한이
$U'$에서 $U$로 가는 동형사상임을 보이자. $P_k$들이 $S$ 위에서
사영이므로 $P$도 그러하고 (\hyperref[II.5.5.5-ko]{5.5.5, (iv)}), 따라서
$X\times_SP$는 $X$ 위에서 사영이며
(\hyperref[II.5.5.5-ko]{5.5.5, (iii)}), $X\times_SP$의 닫힌
부분준스킴인 $X'$도 그러하다. 한편
$f\circ\psi'=q_1\circ(h\circ\psi')=q_1\circ\psi=j$이므로 $f(X')$은
$X$의 어디서나 조밀한 열린집합 $U$를 포함한다~; 그런데 $f$는 닫힌
사상이므로 (\hyperref[II.5.5.3-ko]{5.5.3}) $f(X')=X$이다. 이제
$q_1^{-1}(U)=U\times_SP$는 $X\times_SP$의 한 열린집합 위에 유도되고,
정의에 따라 준스킴 $U'=h^{-1}(U\times_SP)$은 $X'$에서 열린집합
$U'$ 위에 유도된다~; 따라서 이는 준스킴 $Z$의 $U\times_SP$에 대한
폐포이다

\oldpage[II]{108}

(\hyperref[I.9.5.8-ko]{I, 9.5.8}). 그런데 몰입 사상 $\psi$는 다음과
같이 분해된다.
\[
U\xrightarrow{\Gamma_\varphi}U\times_SP\xrightarrow{j\times1}X\times_SP,
\]
$P$가 $S$ 위에서 분리되어 있으므로 그래프 사상 $\Gamma_\varphi$는
닫힌 몰입 사상이다 (\hyperref[I.5.4.3-ko]{I, 5.4.3}). 따라서 $Z$는
$U\times_SP$의 닫힌 부분준스킴이고, 이로부터 $U'=Z$임을 얻는다. 또한
$\psi$가 몰입 사상이므로 $f$의 $U'$로의 제한은 $U$ 위로 가는
동형사상이고 그 역은 $\psi'$이다~; 끝으로 $X'$의 정의에 따라 $U'$은
$X'$에서 조밀하다.

D) 이제 $g$가 몰입 사상임을 증명하자. 그러면 $P$가 $S$ 위에서
사영이므로 $X'$이 $S$ 위에서 준사영임이 확립된다. 다음과 같이 놓자.
\begin{align*}
V_k&=\varphi_k(U_k) &&\text{($P_k$의 열린집합)},\\
W_k&=p_k^{-1}(V_k) &&\text{($P$의 열린집합)},\\
U'_k&=f^{-1}(U_k),\qquad U''_k=g^{-1}(W_k)
&&\text{($X'$의 열린집합)}.
\end{align*}
$U'_k$들이 $X'$의 열린 덮개를 이룸은 명백하다~; 먼저
$U'_k\subset U''_k$임을 증명하여 $U''_k$들도 그러함을 보이겠다. 이를
위해서는 도식
\[
\tag{5.6.1.5}
\label{II.5.6.1.5-ko}
\xymatrix{
U'_k\ar[r]^{g|U'_k}\ar[d]_{f|U'_k} & P\ar[d]^{p_k}\\
U_k\ar[r]_{\varphi_k} & P_k
}
\]
이 가환함을 보이면 충분하다. 그런데 준스킴
$U'_k=h^{-1}(U_k\times_SP)$는 $X'$에서 열린집합 $U'_k$ 위에 유도되며,
따라서 $U'_k$에 대한 $Z=U'\subset U'_k$의 폐포이다
(\hyperref[I.9.5.8-ko]{I, 9.5.8}). 도식
(\hyperref[II.5.6.1.5-ko]{5.6.1.5})이 가환함을 보이기 위해서는,
$P_k$가 $S$-스킴이므로, 이 도식을 표준 단사 사상 $U'\to U'_k$와
합성하여 (또는 $U'$과 $U$의 동형사상으로 인해 마찬가지인 $\psi$와
합성하여) 가환 도식을 얻음을 보이면 충분하다
(\hyperref[I.9.5.6-ko]{I, 9.5.6}). 그러나 정의에 따라 이렇게 얻는
도식은 바로 (\hyperref[II.5.6.1.1-ko]{5.6.1.1})이므로 주장이 따른다.

$W_k$들은 $g(X')$의 열린 덮개를 이룬다~; $g$가 몰입 사상임을
증명하려면 각 제한 사상 $g|U''_k$가 $W_k$ 안으로 가는 몰입 사상임을
보이면 충분하다 (\hyperref[I.4.2.4-ko]{I, 4.2.4}). 이를 위해 사상
$u_k:W_k\xrightarrow{p_k}V_k\xrightarrow{\varphi_k^{-1}}U_k\to X$를
생각하자~; $X$가 $S$ 위에서 분리되어 있으므로 그래프 사상
$\Gamma_{u_k}:W_k\to X\times_SW_k$는 닫힌 몰입 사상이고
(\hyperref[I.5.4.3-ko]{I, 5.4.3}), 따라서 그 그래프
$T_k=\Gamma_{u_k}(W_k)$는 $X\times_SW_k$의 닫힌 부분준스킴이다~;
이 부분준스킴이 $U'$을 포함함을 보이면,
(\hyperref[I.9.5.8-ko]{I, 9.5.8})에 따라 $X'$에서 열린집합 $X''_k$ 위에
유도된 부분준스킴도 포함한다. $q_2$의 $T_k$로의 제한은 $W_k$ 위로
가는 동형사상이므로, $g$의 $X''_k$로의 제한은 $W_k$ 안으로 가는 몰입
사상이며, 주장이 증명된다. $v_k$를 표준 단사 사상
$U'\to X\times_SW_k$라 하자~; $v_k=\Gamma_{u_k}\circ w_k$를 만족하는
사상 $w_k:U'\to W_k$가 존재함을 보여야 한다. 곱의 정의에 따라 다음을
증명하면 충분하다.
$q_1\circ v_k=u_k\circ q_2\circ v_k$
(\hyperref[I.3.2.1-ko]{I, 3.2.1}), 또는 오른쪽에서

\oldpage[II]{109}

동형사상 $\psi':U\to U'$을 합성하면,
$q_1\circ\psi=u_k\circ q_2\circ\psi$임을 증명하면 충분하다. 그런데
$q_1\circ\psi=j$, $q_2\circ\psi=\varphi$이므로, 주장은 $u_k$의 정의를
고려할 때 도식 (\hyperref[II.5.6.1.1-ko]{5.6.1.1})의 가환성에서 따른다.

E) 위의 구성에서 $U$, 따라서 $U'$이 기약이므로 $X'$도 기약이고,
따라서 사상 $f$는 쌍유리 사상임이 명백하다
(\hyperref[I.2.2.9-ko]{I, 2.2.9}). 더 나아가 $X$가 축소이면 $U'$도
축소이고, 따라서 $X'$도 축소이다
(\hyperref[I.9.5.9-ko]{I, 9.5.9}). 이로써 증명이 끝난다.

\begin{corollary}[5.6.2]
\label{II.5.6.2-ko}
(\hyperref[II.5.6.1-ko]{5.6.1})의 가정 a), b) 가운데 하나가 성립한다고
가정하자. $X$가 $S$ 위에서 고유일 필요충분조건은 $S$ 위에서 사영적인
스킴 $X'$과 전사 $S$-사상 $f:X'\to X$가 존재하는 것이다(이때 $f$는
(\hyperref[II.5.5.5-ko]{5.5.5, (v)})에 따라 사영 사상이다). 이 조건이
성립할 때에는, $X$의 조밀 열린집합 $U$가 존재하여 $f$를
$f^{-1}(U)$에 제한한 사상이 동형사상 $f^{-1}(U)\simeq U$이고
$f^{-1}(U)$가 $X'$에서 조밀하도록 $f$를 택할 수도 있다. 더 나아가
$X$가 기약이면(각각 축소이면) $X'$도 그러하다고 가정할 수 있다~;
$X$와 $X'$이 기약일 때 $f$는 쌍유리 사상이다.
\end{corollary}

조건의 충분성은 (\hyperref[II.5.5.3-ko]{5.5.3})과
(\hyperref[II.5.4.3-ko]{5.4.3, (ii)})에서 따른다. 필요성을 보이기 위해
(\hyperref[II.5.6.1-ko]{5.6.1})의 표기를 사용하자. $X$가 $S$ 위에서
고유이면 $X'$도 $S$ 위에서 고유이다. 실제로 $X'$은 $X$ 위에서
사영적이므로 $X$ 위에서 고유이고
(\hyperref[II.5.5.3-ko]{5.5.3}), 우리의 주장은
(\hyperref[II.5.4.2-ko]{5.4.2, (ii)})에서 따른다~; 또한 $X'$은 $S$
위에서 준사영이므로 (\hyperref[II.5.5.3-ko]{5.5.3})에 따라 $S$ 위에서
사영적이다.

\begin{corollary}[5.6.3]
\label{II.5.6.3-ko}
$S$를 국소 뇌터 준스킴, $X$를 $S$ 위에서 유한형인 $S$-스킴,
$f_0:X\to S$를 그 구조 사상이라 하자. $X$가 $S$ 위에서 고유일
필요충분조건은 모든 유한형 사상 $S'\to S$에 대하여
$(f_0)_{(S')}:X_{(S')}\to S'$이 닫힌 사상인 것이다. 이 조건은 심지어
$S'=S\otimes_{\mathbf Z}\mathbf Z[T_1,\ldots,T_n]$ ($T_i$들은
부정원) 꼴의 모든 $S$-준스킴에 대하여 성립하기만 하면 충분하다.
\end{corollary}

이 조건이 필요함은 명백하므로, 충분함을 보이자. 문제는 $S$와 $S'$ 위에서
국소적이므로 (\hyperref[II.5.4.1-ko]{5.4.1}), $S$와 $S'$이 아핀이고
뇌터라고 가정할 수 있다. 저우 보조정리에 따라 $S$ 위에서 사영적인 스킴
$P$, 몰입 사상 $j:X'\to P$, 사영 전사 사상 $f:X'\to X$가 존재하여
다음 도식이 가환한다.
\[
\xymatrix{
X\ar[d]_{f_0} & X'\ar[l]_f\ar[d]^j\\
S & P\ar[l]_r
}
\]
$P$는 $S$ 위에서 유한형이므로 첫 번째 가정에 따라 사영
$q_2:X\times_SP\to P$는 닫힌 사상이다. 그런데 몰입 사상 $j$는
$q_2$와 $X'\times_SP$에서 $X\times_SP$로 가는 사상 $f\times1$의
합성이다~; 한편 $f$는 사영 사상이므로 고유이고
(\hyperref[II.5.5.3-ko]{5.5.3}), 따라서 $f\times1$은 닫힌 사상이다.
그러므로 $j$는 닫힌 몰입 사상이고, 따라서 고유이다
(\hyperref[II.5.4.2-ko]{5.4.2, (ii)}). 또한 구조 사상 $r:P\to S$는
사영 사상이므로 고유이고 (\hyperref[II.5.5.3-ko]{5.5.3}), 따라서
$f_0\circ f=r\circ j$는 고유이다
(\hyperref[II.5.4.2-ko]{5.4.2, (ii)})~; 끝으로 $f$가 전사이므로
(\hyperref[II.5.4.3-ko]{5.4.3})에 따라 $f_0$는 고유이다.

두 번째 가정만을 사용하여 이 명제를 증명하려면, 이 가정이 첫 번째
가정을 함의함을 보이면 충분하다. 실제로 $S'$이
$S=\operatorname{Spec}(A)$ 위에서 유한형인 아핀 스킴이면,

\oldpage[II]{110}

$S'=\operatorname{Spec}(A[c_1,\ldots,c_n])$이고
(\hyperref[I.6.3.3-ko]{I, 6.3.3}), 따라서 $S'$은
$S''=\operatorname{Spec}(A[T_1,\ldots,T_n])$ ($T_i$들은 부정원)의
닫힌 부분준스킴과 동형이다. 가환 도식
\[
\xymatrix{
X\times_SS'\ar[r]^{1_X\times j}\ar[d]_{(f_0)_{(S')}} &
X\times_SS''\ar[d]^{(f_0)_{(S'')}}\\
S'\ar[r]_j & S''
}
\]
에서 $j$와 $1_X\times j$는 닫힌 몰입 사상이고
(\hyperref[I.4.3.1-ko]{I, 4.3.1}), $(f_0)_{(S'')}$는 가정에 따라 닫힌
사상이다~; 따라서 $(f_0)_{(S')}$도 닫힌 사상이다.

\section{정수적 사상과 유한 사상}
\label{section:II.6-ko}

\subsection{다른 준스킴 위에서 정수적인 준스킴.}
\label{subsection:II.6.1-ko}

\begin{definition}[6.1.1]
\label{II.6.1.1-ko}
$X$를 $S$-준스킴, $f:X\to S$를 그 구조 사상이라 하자. $S$의 아핀
열린집합들로 이루어진 덮개 $(S_\alpha)$가 존재하여, 모든 $\alpha$에
대하여 유도된 준스킴 $f^{-1}(S_\alpha)$가 아핀 스킴이고 그 환
$B_\alpha$가 $S_\alpha$의 환 $A_\alpha$ 위에서 정수적인 대수이면,
$X$가 \emph{$S$ 위에서 정수적}이라고, 또는 $f$가 \emph{정수적 사상}이라고
한다. $X$가 $S$ 위에서 정수적이고 $S$ 위에서 유한형이면, $X$가
\emph{$S$ 위에서 유한}이라고, 또는 $f$가 \emph{유한 사상}이라고 한다.
\end{definition}

$S$가 환 $A$의 아핀 스킴이면, ``$S$ 위에서 정수적(각각 유한)''이라고
하는 대신 ``$A$ 위에서 정수적(각각 유한)''이라고도 할 것이다.

\begin{env}[6.1.2]
\label{II.6.1.2-ko}
$X$가 $S$ 위에서 정수적이면 $S$ 위에서 아핀임은 명백하다. $S$ 위에서
아핀인 준스킴 $X$가 $S$ 위에서 정수적(각각 유한)일 필요충분조건은,
대응하는 준연접 $\mathcal{O}_S$-대수 $\mathcal{A}(X)$에 대하여 다음
성질을 갖는 $S$의 아핀 열린집합들로 이루어진 덮개 $(S_\alpha)$가
존재하는 것이다~: 모든 $\alpha$에 대하여
$\Gamma(S_\alpha,\mathcal{A}(X))$는
$\Gamma(S_\alpha,\mathcal{O}_S)$ 위에서 정수적(각각 정수적이고 유한
생성인) 대수이다. 이 성질을 갖는 준연접 $\mathcal{O}_S$-대수를
$\mathcal{O}_S$ 위에서 \emph{정수적}(각각 \emph{유한})이라고 한다.
따라서 $S$ 위에서 정수적(각각 유한)인 준스킴을 정하는 것은
(\hyperref[II.1.3.1-ko]{1.3.1}) $\mathcal{O}_S$ 위에서 정수적(각각
유한)인 준연접 $\mathcal{O}_S$-대수를 정하는 것과 같다. 준연접
$\mathcal{O}_S$-대수 $\mathcal{B}$가 유한일 필요충분조건은 이것이
유한 생성 $\mathcal{O}_S$-가군인 것이다
(\hyperref[I.1.3.9-ko]{I, 1.3.9})~; 이는 $\mathcal{B}$가
$\mathcal{O}_S$ 위에서 정수적이고 유한 생성인 $\mathcal{O}_S$-대수라고
하는 것과 같다. 실제로 환 $A$ 위에서 정수적이고 유한 생성인 대수는
유한 생성 $A$-가군이다.
\end{env}

\begin{proposition}[6.1.3]
\label{II.6.1.3-ko}
$S$를 국소 뇌터 준스킴이라 하자. $S$ 위에서 아핀인 준스킴 $X$가
$S$ 위에서 유한일 필요충분조건은 $\mathcal{O}_S$-대수
$\mathcal{A}(X)$가 연접인 것이다.
\end{proposition}

앞의 주의를 고려하면, 이는 $S$가 국소 뇌터일 때 유한 생성 준연접
$\mathcal{O}_S$-가군들이 바로 연접 $\mathcal{O}_S$-가군들임을 주목하는
것과 같다 (\hyperref[I.1.5.1-ko]{I, 1.5.1}).

\begin{proposition}[6.1.4]
\label{II.6.1.4-ko}
$X$를 $S$ 위에서 정수적(각각 유한)인 준스킴, $f:X\to S$를 그 구조
사상이라 하자. 그러면 환이 $A$인 모든 아핀 열린집합 $U\subset S$에
대하여, $f^{-1}(U)$는 아핀 스킴이고 그 환 $B$는 $A$ 위에서 정수적
(각각 유한)인 대수이다.
\end{proposition}

\oldpage[II]{111}

먼저 다음 보조정리를 증명하겠다~:

\begin{lemma}[6.1.4.1]
\label{II.6.1.4.1-ko}
$A$를 환, $M$을 $A$-가군, $(g_i)_{1\leq i\leq m}$를 $A$의 원소들로
이루어진 유한족이라 하고, $D(g_i)$들 ($1\leq i\leq m$)이
$\operatorname{Spec}(A)$를 덮는다고 하자. 모든 $i$에 대하여 $M_{g_i}$가
유한 생성 $A_{g_i}$-가군이면, $M$은 유한 생성 $A$-가군이다.
\end{lemma}

실제로 각 $M_{g_i}$가 $(m_{ij}/g_i^n)$ 꼴의 유한 생성계를 가지며,
$m_{ij}\in M$이고 $n$은 모든 지표 $i$에 대하여 같다고 가정할 수 있다.
$m_{ij}$들이 $M$의 생성계를 이룸을 보이자. 이 원소들이 생성하는 $M$의
$A$-부분가군을 $M'$이라 하고, $m$을 $M$의 한 원소라 하자. 가정에 따라
각 $i$에 대하여 $a_{ij}\in A$인 원소들과 $i$와 무관한 정수 $p$가 존재하여
$M_{g_i}$에서
$m/1=(\sum_j a_{ij}m_{ij})/g_i^p$이다~; 이로부터 모든 $i$에 대하여
$g_i^r m\in M'$이 되는 정수 $r\geq p$가 존재한다. 그런데
$D(g_i^r)=D(g_i)$들이 $\operatorname{Spec}(A)$를 덮으므로, $g_i^r$들이
생성하는 $A$의 아이디얼은 $A$와 같다. 다시 말해
$\sum_i a_i g_i^r=1$을 만족하는 원소 $a_i\in A$들이 존재한다~; 따라서
$m=(\sum_i a_i g_i^r)m\in M'$이고, 보조정리가 증명된다.

이제 $f^{-1}(U)$가 아핀임은 이미 알고 있다
(\hyperref[II.1.3.2-ko]{1.3.2}). $f$에 대응하는 준동형이
$\varphi:A\to B$이면, $U$에는 $D(g_i)$ ($g_i\in A$)들로 이루어진 유한
덮개가 존재하여, $h_i=\varphi(g_i)$라 놓을 때 $B_{h_i}$는
$A_{g_i}$ 위에서 정수적(각각 유한)인 대수가 된다. 실제로 $U$에는 아핀
열린집합 $V_\alpha\subset U$들로 이루어진 덮개가 존재하여,
$A_\alpha=A(V_\alpha)$,
$B_\alpha=A(f^{-1}(V_\alpha))$라 놓으면 $B_\alpha$는 $A_\alpha$ 위에서
정수적(각각 유한)인 대수가 된다. 모든 $x\in U$는 어떤 $V_\alpha$에
속하므로, $x\in D(g)\subset V_\alpha$를 만족하는 $g\in A$가 존재한다~;
$g_\alpha$가 $A_\alpha$에서 $g$의 상이면
$A(D(g))=A_g=(A_\alpha)_{g_\alpha}$이다~; $h=\varphi(g)$라 하고,
$h_\alpha$를 $B_\alpha$에서 $g_\alpha$의 상이라 하자~; 그러면
\[
A(D(h))=B_h=(B_\alpha)_{h_\alpha}
\]
이고, $B_\alpha$가 $A_\alpha$ 위에서 정수적(각각 유한)이므로
$(B_\alpha)_{h_\alpha}$는 $(A_\alpha)_{g_\alpha}$ 위에서 정수적
(각각 유한)이다. 이제 $U$가 준콤팩트라는 사실을 사용하면 원하는 덮개를
얻는다.

먼저 $B_{h_i}$들이 $A_{g_i}$ 위에서 유한이라고 가정하자.
$A_{g_i}$-가군으로서 $B_{h_i}$를 $B_{g_i}$라고도 쓸 수 있으므로,
보조정리 (\hyperref[II.6.1.4.1-ko]{6.1.4.1})에 따라 이 경우 $B$는 유한
생성 $A$-가군이다.

이제 각 $B_{h_i}$가 $A_{g_i}$ 위에서 정수적이라고만 가정하자~;
$b\in B$라 하고, $C$를 $b$가 생성하는 $B$의 $A$-부분대수라 하자. 모든
$i$에 대하여 $C_{h_i}$는 $B_{h_i}$ 안에서 $b/1$이 생성하는
$A_{g_i}$-대수이다~; 가정으로부터 각 $C_{h_i}$가 유한 생성
$A_{g_i}$-가군임이 따르므로,
(\hyperref[II.6.1.4.1-ko]{6.1.4.1})에 따라 $C$는 유한 생성 $A$-가군이다.
이는 $B$가 $A$ 위에서 정수적임을 증명한다.

\begin{proposition}[6.1.5]
\label{II.6.1.5-ko}
\begin{enumerate}
\item[(i)] 닫힌 몰입 사상은 유한 사상이다(따라서 특히 정수적 사상이다).
\item[(ii)] 두 유한 사상(각각 정수적 사상)의 합성은 유한 사상(각각
정수적 사상)이다.
\item[(iii)] $f:X\to Y$가 유한(각각 정수적) $S$-사상이면, 밑준스킴의
모든 확장 $S'\to S$에 대하여
$f_{(S')}:X_{(S')}\to Y_{(S')}$도 유한 사상(각각 정수적 사상)이다.
\item[(iv)] $f:X\to Y$와 $g:X'\to Y'$이 두 유한(각각 정수적)
$S$-사상이면,
$f\times_S g:X\times_S X'\to Y\times_S Y'$도 그러하다.
\item[(v)] $f:X\to Y$와 $g:Y\to Z$가 두 사상이고 $g\circ f$가 유한
사상(각각 정수적 사상)이며 $g$가 분리 사상이면, $f$는 유한 사상
(각각 정수적 사상)이다.
\item[(vi)] $f:X\to Y$가 유한 사상(각각 정수적 사상)이면
$f_{\mathrm{red}}$도 그러하다.
\end{enumerate}
\end{proposition}

\oldpage[II]{112}

(\hyperref[I.5.5.12-ko]{I, 5.5.12})에 따라 (i), (ii), (iii)을 증명하면
충분하다. 닫힌 몰입 사상 $X\to S$가 유한임을 증명할 때에는
$S=\operatorname{Spec}(A)$인 경우로 한정할 수 있고, 그러면 몫환
$A/\mathfrak{J}$가 한 원소로 생성되는 $A$-가군임을 주목하면 된다.
두 유한(각각 정수적) 사상 $X\to Y$, $Y\to Z$의 합성이 유한(각각
정수적)임을 증명할 때에도 $Z$가 아핀이라고, 따라서 $X$, $Y$도
(\hyperref[II.1.3.4-ko]{1.3.4}) 아핀이라고 가정할 수 있다. 그러면 이
명제는 $B$가 $A$ 위에서 유한인(각각 정수적인) 대수이고 $C$가 $B$
위에서 유한인(각각 정수적인) 대수이면, $C$는 $A$ 위에서 유한인
(각각 정수적인) 대수라는 것과 동치이며, 이는 자명하다. 끝으로
(iii)을 증명할 때에는 $X_{(S')}$이 $X\times_Y Y_{(S')}$과 동일시되므로
(\hyperref[I.3.3.11-ko]{I, 3.3.11}), $S=Y$인 경우로 한정할 수 있다~;
더 나아가 $S=\operatorname{Spec}(A)$,
$S'=\operatorname{Spec}(A')$라고 가정할 수 있다. 이때 $X$는 환이
$B$인 아핀 스킴이고 (\hyperref[II.1.3.4-ko]{1.3.4}), $X_{(S')}$은 환이
$A'\otimes_A B$인 아핀 스킴이므로, $B$가 $A$ 위에서 유한인(각각
정수적인) 대수이면 $A'\otimes_A B$가 $A'$ 위에서 유한인(각각
정수적인) 대수임을 주목하면 충분하다.

또한 $X$와 $Y$가 $S$ 위에서 유한인(각각 정수적인) 두 $S$-준스킴이면,
그들의 \emph{합} $X\amalg Y$는 $S$ 위에서 유한인(각각 정수적인)
준스킴이다. 실제로 이는 $B$와 $C$가 $A$ 위에서 유한인(각각 정수적인)
두 대수이면 $B\times C$도 그러함을 확인하는 것으로 귀착된다.

\begin{corollary}[6.1.6]
\label{II.6.1.6-ko}
$X$가 $S$ 위에서 정수적(각각 유한)이면, 모든 열린집합 $U\subset S$에
대하여 $f^{-1}(U)$는 $U$ 위에서 정수적(각각 유한)이다.
\end{corollary}

이는 (\hyperref[II.6.1.5-ko]{6.1.5, (iii)})의 특수한 경우이다.

\begin{corollary}[6.1.7]
\label{II.6.1.7-ko}
$f:X\to Y$를 유한 사상이라 하자. 그러면 모든 $y\in Y$에 대하여 섬유
$f^{-1}(y)$는 $k(y)$ 위의 유한 대수적 스킴이고, 따라서 특히 그
바탕공간은 이산이고 유한하다.
\end{corollary}

실제로 $k(y)$-준스킴으로서 $f^{-1}(y)$는
$X\times_Y\operatorname{Spec}(k(y))$와 동일시되고
(\hyperref[I.3.6.1-ko]{I, 3.6.1}), 이는
$\operatorname{Spec}(k(y))$ 위에서 유한이다
(\hyperref[II.6.1.5-ko]{6.1.5, (iii)})~; 따라서 이는 환이 $k(y)$ 위의
유한 랭크 대수인 아핀 스킴이다
(\hyperref[II.6.1.4-ko]{6.1.4}). 이제 명제는
(\hyperref[I.6.4.4-ko]{I, 6.4.4})에서 따른다.

\begin{corollary}[6.1.8]
\label{II.6.1.8-ko}
$X$, $S$를 두 정역 준스킴, $f:X\to S$를 우세 사상이라 하자. $f$가
정수적(각각 유한)이면, $X$의 유리함수체 $R(X)$는 $S$의 유리함수체
$R(S)$의 대수적 확대체(각각 유한 차수의 대수적 확대체)이다.
\end{corollary}

실제로 $s$를 $S$의 일반점이라 하자~; $k(s)$-준스킴 $f^{-1}(s)$는
$\operatorname{Spec}(k(s))$ 위에서 정수적(각각 유한)이고
(\hyperref[II.6.1.5-ko]{6.1.5, (iii)}), 가정상 $X$의 일반점 $x$를
포함한다~; $f^{-1}(s)$에서 $x$의 국소환은 $k(x)$와 같고
(\hyperref[I.3.6.5-ko]{I, 3.6.5}), $k(s)$ 위에서 정수적인(각각 유한인)
대수의 국소환이므로 (\hyperref[II.6.1.4-ko]{6.1.4}) 따름정리를 얻는다.

\begin{remark}[6.1.9]
\label{II.6.1.9-ko}
(\hyperref[II.6.1.5-ko]{6.1.5, (v)})가 성립하려면 $g$가 \emph{분리
사상}이라는 가정이 필수적이다~: 실제로 $Y$가 $Z$ 위에서 분리되지
않으면 항등사상 $1_Y$는 합성 사상
$Y\xrightarrow{\Delta_Y}Y\times_ZY\xrightarrow{p_1}Y$이지만,
(\hyperref[II.6.1.10-ko]{6.1.10})에서 알 수 있듯이 $\Delta_Y$는 정수적
사상이 아니다~:
\end{remark}

\begin{proposition}[6.1.10]
\label{II.6.1.10-ko}
모든 정수적 사상은 보편적으로 닫혀 있다.
\end{proposition}

$f:X\to Y$를 정수적 사상이라 하자~; 
(\hyperref[II.6.1.5-ko]{6.1.5, (iii)})에 따라 $f$가 닫혀 있음을 보이면
충분하다. $Z$를 $X$의 닫힌 부분집합이라 하자~; $X$의 한 부분준스킴이
존재하여

\oldpage[II]{113}

그 바탕공간이 $Z$이고 (\hyperref[I.5.2.1-ko]{I, 5.2.1}), 따라서
(\hyperref[II.6.1.5-ko]{6.1.5, (i)와 (ii)})에서 $f(X)$가 $Y$에서
닫혀 있음을 증명하는 경우로 귀착시킬 수 있음을 얻는다.
(\hyperref[II.6.1.5-ko]{6.1.5, (vii)})에 따라 $X$와 $Y$가 축소라고
가정할 수 있다~; 또한 $T$를 바탕공간이 $\overline{f(X)}$인 $Y$의 축소
닫힌 부분준스킴이라 하면 (\hyperref[I.5.2.1-ko]{I, 5.2.1}), $f$는
$X\xrightarrow{g}T\xrightarrow{j}Y$로 분해됨을 안다. 여기서 $j$는 단사
사상이고 (\hyperref[I.5.2.2-ko]{I, 5.2.2}), $j$가 분리되어 있으므로
(\hyperref[I.5.5.1-ko]{I, 5.5.1, (i)}),
(\hyperref[II.6.1.5-ko]{6.1.5, (v)})에서 $g$가 정수적 사상임을 얻는다.
따라서 $f(X)$가 $Y$에서 조밀하다고 가정할 수 있다. 끝으로 문제는
$Y$ 위에서 국소적이므로 $Y=\operatorname{Spec}(A)$인 경우로 귀착시킬
수 있다. 이때 $X=\operatorname{Spec}(B)$이고, $B$는 $A$ 위에서
정수적인 $A$-대수이다 (\hyperref[II.6.1.4-ko]{6.1.4})~; 또한 $A$는
축소이고 (\hyperref[I.5.1.4-ko]{I, 5.1.4}), $f(X)$가 $Y$에서
조밀하다는 가정으로부터 $f$에 대응하는 준동형 $\varphi:A\to B$가
단사임을 얻는다 (\hyperref[I.1.2.7-ko]{I, 1.2.7}). 이 조건 아래에서
$f(X)=Y$라는 것은 $A$의 모든 소아이디얼이 $B$의 어떤 소아이디얼을
$A$로 축약한 것이라는 뜻이며, 이는 바로 제1 코언--자이덴베르크 정리이다
(\cite[t.~I, p.~257, th.~3]{I-13}).

\begin{corollary}[6.1.11]
\label{II.6.1.11-ko}
모든 유한 사상 $f:X\to Y$는 사영적이다.
\end{corollary}

실제로 $f$가 아핀이므로 $\mathcal{O}_X$는 $f$에 관해 매우 풍부한
$\mathcal{O}_X$-가군이다 (\hyperref[II.5.1.2-ko]{5.1.2})~; 또한
$f_*(\mathcal{O}_X)$는 유한형 준연접 $\mathcal{O}_Y$-가군이고
(\hyperref[II.6.1.2-ko]{6.1.2})~; 끝으로 $f$는 분리된 유한형 사상이며
보편적으로 닫혀 있으므로 (\hyperref[II.6.1.10-ko]{6.1.10}), 판정법
(\hyperref[II.5.5.4-ko]{5.5.4, (ii)})을 적용할 수 있는 조건에 놓여 있다.

\begin{proposition}[6.1.12]
\label{II.6.1.12-ko}
$f:X'\to X$를 유한 사상이라 하고,
$\mathcal{B}=f_*(\mathcal{O}_{X'})$라 하자(이는 준연접
$\mathcal{O}_X$-대수이자 유한형 $\mathcal{O}_X$-가군이다).
$\mathcal{F}'$을 준연접 $\mathcal{O}_{X'}$-가군이라 하자~;
$\mathcal{F}'$이 계수 $r$인 국소 자유 가군일 필요충분조건은
$f_*(\mathcal{F}')$이 계수 $r$인 국소 자유 $\mathcal{B}$-가군인 것이다.
\end{proposition}

$f_*(\mathcal{F}')|U$가 $\mathcal{B}^r|U$와 동형이면($U$는 $X$의
열린집합), $\mathcal{F}'|f^{-1}(U)$가
$\mathcal{O}_{X'}^r|f^{-1}(U)$와 동형임은 명백하다
(\hyperref[II.1.4.2-ko]{1.4.2}). 역으로 $\mathcal{F}'$이 계수 $r$인
국소 자유 가군이라고 가정하고, $f_*(\mathcal{F}')$이
$\mathcal{B}$-가군으로서 $\mathcal{B}^r$과 국소적으로 동형임을
보이자. $x$를 $X$의 한 점이라 하자~; $U$가 $x$의 아핀 근방들로
이루어진 기본계를 달릴 때, $f^{-1}(U)$는 유한집합 $f^{-1}(x)$의 아핀
근방들로 이루어진 기본계를 달린다
(\hyperref[II.1.2.5-ko]{1.2.5}). 실제로 $f$가 닫힌 사상이기 때문이다
(\hyperref[II.6.1.10-ko]{6.1.10}). 이제 명제는 다음 보조정리에서
따른다~:

\begin{lemma}[6.1.12.1]
\label{II.6.1.12.1-ko}
$Y$를 준스킴, $\mathcal{E}$를 계수 $r$의 국소 자유
$\mathcal{O}_Y$-가군, $Z$를 아핀 열린집합 $V$에 포함되는 $Y$의 유한
부분집합이라 하자. 그러면 $Z$의 근방 $U\subset V$가 존재하여
$\mathcal{E}|U$는 $\mathcal{O}_Y^r|U$와 동형이다.
\end{lemma}

$Y$가 아핀이라고 가정해도 됨은 명백하다~; 모든 $z_i\in Z$에 대하여
폐포 $\overline{\{z_i\}}$ 안에는 적어도 하나의 닫힌점 $z'_i$가 존재한다
(\hyperref[0.2.1.3-ko]{0, 2.1.3})~; $Z'$을 $z'_i$들의 집합이라 하면
$Z'$의 모든 근방은 $Z$의 근방이므로, $Z$가 $Y$에서 닫혀 있고 이산이라고
가정할 수 있다. 바탕공간이 $Z$인 $Y$의 축소 닫힌 부분준스킴을 생각하고
(\hyperref[I.5.2.1-ko]{I, 5.2.1}), $j:Z\to Y$를 표준 포함 사상이라 하자~;
$j^*(\mathcal{E})=\mathcal{E}\otimes_Y
\mathcal{O}_Z$는 이산 스킴 $Z$ 위에서 계수 $r$의 국소 자유이고,
따라서 $\mathcal{O}_Z^r$와 동형이다~; 다시 말해 $Z$ 위의
$\mathcal{E}\otimes_Y\mathcal{O}_Z$의 단면 $s_i$ ($1\leq i\leq r$)가
$r$개 존재하여, 이 단면들이 정의하는 준동형
$\mathcal{O}_Z^r\to\mathcal{E}\otimes_Y\mathcal{O}_Z$는 전단사이다.
그런데 $Y=\operatorname{Spec}(A)$는 아핀이고, $Z$는 $A$의 아이디얼
$\mathfrak{J}$로 정의되며, $\mathcal{E}=\widetilde{M}$이다. 여기서 $M$은
$A$-가군이다~; $s_i$들은

\oldpage[II]{114}

$M\otimes_A(A/\mathfrak{J})$의 원소들이므로, $t_i\in M=\Gamma(Y,\mathcal{E})$인
$r$개의 원소의 상이다. 모든 $z_j\in Z$에
대하여 $z_j$의 근방 $V_j$가 존재하여, $t_i$들을 $V_j$에 제한한 것들이
동형사상
$\mathcal{O}_Y^r|V_j\to\mathcal{E}|V_j$
을 정의한다(\hyperref[0.5.5.4-ko]{0, 5.5.4})~; 따라서 $V_j$들의 합집합인
근방 $U$가 원하는 조건을 만족한다.

\begin{proposition}[6.1.13]
\label{II.6.1.13-ko}
$g:X'\to X$를 준스킴들의 정수적 사상, $Y$를 정규인 국소 정역 준스킴,
$f$를 $Y$에서 $X'$으로 가는 유리 사상이라 하자. $g\circ f$가 모든 곳에서
정의된 유리 사상이라면(\hyperref[I.7.2.1-ko]{I, 7.2.1}), $f$도 모든
곳에서 정의된다.
\end{proposition}

$f_1$과 $f_2$를 두 사상(즉 $Y$의 조밀한 열린집합들에서 $X'$으로
가는 사상)으로서 동치류 $f$에 속하는 것들이라 하면, $g\circ f_1$과
$g\circ f_2$가 동치인
사상임은 명백하다. 이로써 그 동치류에 $g\circ f$라는 표기를 쓰는 것이
정당화된다. 또한 $Y$가 국소 뇌터라고 더 가정할 때에는, $Y$가 정규라는
가정만으로도 이미 $Y$가 국소 정역 준스킴임이 따른다는 것을 상기하자
(\hyperref[I.6.1.13-ko]{I, 6.1.13}).

(\hyperref[II.6.1.13-ko]{6.1.13})을 증명하기 위해, 먼저 이 문제는 $Y$
위에서 국소적이므로 $g\circ f$의 동치류에 사상 $h:Y\to X$가 속한다고
가정할 수 있음에 유의하자. 역상
$Y'=X'_{(h)}=X'_{(Y)}$을 생각하면, 사상
$g'=g_{(Y)}:Y'\to Y$는 정수적이다
(\hyperref[II.6.1.5-ko]{6.1.5, (iii)}). $Y$에서 $X'$으로 가는 유리
사상과 $Y$-유리 단면으로 본 $Y'$의 단면 사이의 대응
(\hyperref[I.7.1.2-ko]{I, 7.1.2})을 고려하면,
(\hyperref[II.6.1.13-ko]{6.1.13})에서 $X=Y$인 특수한 경우, 다시 말해 다음으로
환원됨을 알 수 있다.

\begin{corollary}[6.1.14]
\label{II.6.1.14-ko}
$X$를 정규인 국소 정역 준스킴, $g:X'\to X$를 정수적 사상, $f$를
$X'$의 $X$-유리 단면이라 하자. 그러면 $f$는 모든 곳에서 정의된다.
\end{corollary}

이 문제는 $X$ 위에서 국소적이므로 $X$가 정역 준스킴이라고 가정할 수
있고, 그러면 $f$는 $X$의 열린집합 $U$에서 $X'$으로 가는 사상과
동일시된다(\hyperref[I.7.2.2-ko]{I, 7.2.2}). 이 사상은
$g^{-1}(U)$의 $U$-단면이다. $g$가 분리 사상이므로, $f$는 $U$에서
$g^{-1}(U)$로 가는 닫힌 몰입 사상이다
(\hyperref[I.5.4.6-ko]{I, 5.4.6})~; $Z$를 $g^{-1}(U)$에서 $f$에 대응하는
닫힌 부분준스킴이라 하자
(\hyperref[I.4.2.1-ko]{I, 4.2.1}). 이는 $U$와 동형이므로 정역
준스킴이다~; $X_1$을 $X'$의 축소 부분준스킴이라 하되, 그 바탕공간은
$\overline{Z}$, 곧 $Z$의 $X'$ 안에서의 폐포라 하자
(\hyperref[I.5.2.1-ko]{I, 5.2.1})~; 그러면 $Z$는 $X_1$의 어떤
열린집합 위에 유도된 부분준스킴이고
(\hyperref[I.5.2.3-ko]{I, 5.2.3}), 기약이므로 $X_1$도 기약이다. 따라서
$X_1$은 정역 준스킴이다. 이제 사상 $f$를 $X$-유리 단면으로서 $X_1$에
값을 갖는 것으로 볼 수 있다~; $g$를 $X_1$에 제한한 사상은 정수적이므로
(\hyperref[II.6.1.5-ko]{6.1.5, (i) et (ii)}), 결국
(\hyperref[II.6.1.14-ko]{6.1.14})을 $X'=X_1$인 특수한 경우, 다시 말해
다음을 증명하는 것으로 환원된다.

\begin{corollary}[6.1.15]
\label{II.6.1.15-ko}
$X$를 정규인 정역 준스킴, $X'$을 정역 준스킴,
$g:X'\to X$를 정수적 사상이라 하자. $X'$의 $X$-유리 단면 $f$가
존재하면 $g$는 동형사상이다.
\end{corollary}

문제는 $X$ 위에서 국소적이므로, $X$가 환이 정역 $A$인 아핀이고
$X'$이 환이 $A$ 위에서 정수적이고 정역인 $A'$인 아핀이라고 가정할 수
있다 (\hyperref[II.6.1.4-ko]{6.1.4})~; 더구나
(\hyperref[II.6.1.14-ko]{6.1.14})의 논증은 $X$의 한 조밀 열린집합과
$X'$의 한 조밀 열린집합이 서로 동형임을 보이므로, $A$와 $A'$은 같은
분수체를 갖는다. 한편 (\hyperref[I.8.2.1.1-ko]{I, 8.2.1.1})과
$\mathcal{O}_x$들이 정수적으로 닫혀 있다는 가정에 따라 환 $A$는
정수적으로 닫혀 있다. 따라서 $A'=A$이고, 이로써
(\hyperref[II.6.1.13-ko]{6.1.13})의 증명이 끝난다.

\subsection{준유한 사상.}
\label{subsection:II.6.2-ko}

\begin{proposition}[6.2.1]
\label{II.6.2.1-ko}
$f:X\to Y$를 국소 유한형 사상, $x$를 $X$의 한 점이라 하자. 다음
조건들은 동치이다~:

\oldpage[II]{115}

\begin{enumerate}
\item[a)] 점 $x$는 자신의 올 $f^{-1}(f(x))$에서 고립점이다.
\item[b)] 환 $\mathcal{O}_x$는 준유한 $\mathcal{O}_{f(x)}$-가군이다
(\hyperref[0.7.4.1-ko]{0, 7.4.1}).
\end{enumerate}
\end{proposition}

문제는 명백히 $X$와 $Y$ 위에서 국소적이므로,
$X=\operatorname{Spec}(A)$와 $Y=\operatorname{Spec}(B)$가 아핀이고
$A$가 유한형 $B$-대수라고 가정할 수 있다
(\hyperref[I.6.3.3-ko]{I, 6.3.3}). 더 나아가 $X$를
$X\times_Y\operatorname{Spec}(\mathcal{O}_{f(x)})$로 바꾸어도 올
$f^{-1}(f(x))$과 국소환 $\mathcal{O}_x$는 바뀌지 않는다
(\hyperref[I.3.6.5-ko]{I, 3.6.5})~; 따라서 $B$가
$\mathcal{O}_{f(x)}$와 같은 국소환이라고 가정할 수 있다. $\mathfrak{n}$을
$B$의 극대 아이디얼이라 하면, $f^{-1}(f(x))$은 환이
$A/\mathfrak{n}A$인 아핀 스킴이고 $k(f(x))=B/\mathfrak{n}$ 위에서
유한형이다 (\hyperref[I.6.4.11-ko]{I, 6.4.11}). 이제 a)가 성립하면
$f^{-1}(f(x))$이 $x$ 한 점으로만 이루어져 있다고 더 가정할 수 있다~;
따라서 $A/\mathfrak{n}A$는 $B/\mathfrak{n}$ 위에서 유한 차원이다
(\hyperref[I.6.4.4-ko]{I, 6.4.4}). 다시 말해 $A$는 준유한
$B$-가군이다. 역으로 이 조건이 성립하면 $f^{-1}(f(x))$은 아르틴 아핀
스킴이고, 따라서 이산적이다 (\hyperref[I.6.4.4-ko]{I, 6.4.4})~;
그러므로 $x$는 자신의 올에서 고립점이며, 이는 b)가 a)를 함의함을
보인다.

\begin{corollary}[6.2.2]
\label{II.6.2.2-ko}
$f:X\to Y$를 유한형 사상이라 하자. 다음 조건들은 동치이다~:
\begin{enumerate}
\item[a)] 모든 점 $x\in X$는 자신의 올 $f^{-1}(f(x))$에서 고립점이다
(다시 말해 부분공간 $f^{-1}(f(x))$은 이산적이다).
\item[b)] 모든 $x\in X$에 대하여 준스킴 $f^{-1}(f(x))$은
$k(f(x))$ 위에서 유한이다.
\item[c)] 모든 $x\in X$에 대하여 환 $\mathcal{O}_x$는 준유한
$\mathcal{O}_{f(x)}$-가군이다.
\end{enumerate}
\end{corollary}

a)와 c)의 동치는 (\hyperref[II.6.2.1-ko]{6.2.1})에서 따른다. 한편
$f^{-1}(f(x))$은 $k(f(x))$-대수적 준스킴이므로
(\hyperref[I.6.4.11-ko]{I, 6.4.11}), a)와 b)의 동치는
(\hyperref[I.6.4.4-ko]{I, 6.4.4})에서 따른다.

\begin{definition}[6.2.3]
\label{II.6.2.3-ko}
$f:X\to Y$가 (\hyperref[II.6.2.2-ko]{6.2.2})의 동치인 조건들을
만족하는 유한형 사상일 때, $f$가 \emph{준유한}이라고, 또는 $X$가
$Y$ 위에서 준유한이라고 한다.
\end{definition}

모든 유한 사상이 준유한임은 명백하다
(\hyperref[II.6.1.8-ko]{6.1.8}).

\begin{proposition}[6.2.4]
\label{II.6.2.4-ko}
\begin{enumerate}
\item[(i)] 닫힌 몰입 사상이거나 $X$가 뇌터인 몰입 사상 $X\to Y$는
준유한 사상이다.
\item[(ii)] $f:X\to Y$와 $g:Y\to Z$가 준유한 사상이면
$g\circ f$도 준유한이다.
\item[(iii)] $X$와 $Y$가 $S$-준스킴이고 $f:X\to Y$가 준유한
$S$-사상이면, 기저의 모든 확장 $g:S'\to S$에 대하여
$f_{(S')}:X_{(S')}\to Y_{(S')}$는 준유한이다.
\item[(iv)] $f:X\to Y$와 $g:X'\to Y'$이 두 준유한 $S$-사상이면,
\[
f\times_S g:X\times_S Y\to X'\times_S Y'
\]
는 준유한이다.
\item[(v)] $f:X\to Y$와 $g:Y\to Z$가 $g\circ f$가 준유한인
두 사상이라 하자~; 여기에 더하여 $g$가 분리 사상이거나, $X$가
뇌터이거나, $X\times_ZY$가 국소 뇌터이면 $f$는 준유한이다.
\item[(vi)] $f$가 준유한이면 $f_{\mathrm{red}}$도 그러하다.
\end{enumerate}
\end{proposition}

$f:X\to Y$가 몰입 사상이면 각 올은 한 점으로만 이루어지므로,
(i)은 (\hyperref[I.6.3.4-ko]{I, 6.3.4 (i)와 6.3.5})에서 따른다.
(ii)를 증명하기 위해 먼저 $h=g\circ f$가 유한형임을 주목하자
(\hyperref[I.6.3.4-ko]{I, 6.3.4 (ii)})~; 더 나아가 $z=h(x)$ 및
$y=f(x)$이면 $y$는 $g^{-1}(z)$에서 고립되어 있으므로, $Y$ 안에서
$y$의 열린 이웃 $V$로서 $g^{-1}(z)$의 $y$와 다른 점을 하나도 만나지
않는 것이 존재한다~; 따라서 $f^{-1}(V)$는 $x$의 열린 이웃이고,
$g^{-1}(z)$에 속하는 $y'\neq y$에 대한 어느

\oldpage[II]{116}

$f^{-1}(y')$와도 만나지 않는다~; $x$가 $f^{-1}(y)$에서 고립되어
있으므로, $x$는 $h^{-1}(z)=f^{-1}(g^{-1}(z))$에서도 고립되어 있다.
(iii)을 증명할 때에는 $Y=S$인 경우로 한정할 수 있다
(\hyperref[I.3.3.11-ko]{I, 3.3.11})~; 다시 먼저
$f'=f_{(S')}$가 유한형임을 주목하자
(\hyperref[I.6.3.4-ko]{I, 6.3.4, (iii)})~; 다른 한편
$x'\in X'=X_{(S')}$이고 $y'=f'(x')$, $y=g(y')$로 놓으면,
$f'^{-1}(y')$는 $f^{-1}(y)\otimes_{k(y)}k(y')$와 동일시된다
(\hyperref[I.3.6.5-ko]{I, 3.6.5})~; 가정에 따라 $f^{-1}(y)$는
$k(y)$ 위에서 유한 랭크이므로 $f'^{-1}(y')$는 $k(y')$ 위에서
유한 랭크이고, 따라서 이산적이다. (iv), (v), (vi)은 일반적인 방법
(\hyperref[I.5.5.12-ko]{I, 5.5.12})에 따라 (i), (ii), (iii)에서
도출된다. 다만 (v)의 가정이 ``$g$가 분리 사상이다''라는 가정 이외의
것인 경우는 예외이다~; 이 경우를 다루기 위해 먼저 $x$가
$f^{-1}(g^{-1}(g(f(x))))$에서 고립되어 있으면 하물며
$f^{-1}(f(x))$에서도 고립되어 있음을 주목하자~; 그러면 $f$가
유한형이라는 사실은 (\hyperref[I.6.3.6-ko]{I, 6.3.6})에서 따른다.

\begin{proposition}[6.2.5]
\label{II.6.2.5-ko}
$A$를 완비 뇌터 국소환, $X$를 $A$ 위에서 국소 유한형인 스킴,
$x$를 $Y=\operatorname{Spec}(A)$의 닫힌점 $y$ 위에 놓인 $X$의
점이라 하고, $x$가 자신의 올 $f^{-1}(y)$에서 고립되어 있다고
가정하자(여기서 $f$는 구조 사상 $X\to Y$이다). 그러면
$\mathcal{O}_x$는 유한형 $A$-가군이고, $X$는
$X'=\operatorname{Spec}(\mathcal{O}_x)$($Y$-유한 스킴)와 어떤
$A$-스킴 $X''$의 합 (I, 3.1)과 $Y$-동형이다.
\end{proposition}

(\hyperref[II.6.2.1-ko]{6.2.1})에 따라 $\mathcal{O}_x$는 준유한
$A$-가군이다. $\mathcal{O}_x$가 뇌터이고
(\hyperref[I.6.3.7-ko]{I, 6.3.7}) 준동형
$A\to\mathcal{O}_x$가 국소 준동형이므로, $A$가 완비라는 가정에서
$\mathcal{O}_x$가 유한형 $A$-가군임이 따른다
(\hyperref[0.7.4.3-ko]{0, 7.4.3}).
$X'=\operatorname{Spec}(\mathcal{O}_x)$를 점 $x$에서의 $X$의
국소 스킴이라 하고 (\hyperref[I.2.4.1-ko]{I, 2.4.1}),
$g:X'\to X$를 표준 사상이라 하자. 합성
$X'\xrightarrow{g}X\xrightarrow{f}Y$가 유한하고
(\hyperref[II.6.1.1-ko]{6.1.1}) $f$가 분리 사상이므로 $g$는
유한하며 (\hyperref[II.6.1.5-ko]{6.1.5, (v)}), 따라서 $g(X')$는
$X$에서 닫혀 있다 (\hyperref[II.6.1.10-ko]{6.1.10})~; 다른 한편
$g$가 유한형이므로, $g$의 정의와
(\hyperref[I.6.5.4-ko]{I, 6.5.4})에 따라 $g$는 $X'$의 닫힌점
$x'$에서 국소 동형사상이다~; 그런데 $X'$ 자체가 $x'$의 유일한
열린 이웃이므로, 이는 $g$가 열린 몰입 사상이라는 뜻이다. 따라서
$g(X')$는 $X$에서 열려 있기도 하며, 이것으로 증명이 끝난다.

\begin{corollary}[6.2.6]
\label{II.6.2.6-ko}
$A$를 완비 뇌터 국소환, $Y=\operatorname{Spec}(A)$라 하고,
$f:X\to Y$를 분리 준유한 사상이라 하자. 그러면 $X$는
$X'\amalg X''$ 꼴의 합과 $Y$-동형이다. 여기서 $X'$은 $Y$-유한
스킴이고 $X''$은 준유한 $Y$-스킴으로서, $y$가 $Y$의 닫힌점이면
$X''\cap f^{-1}(y)=\varnothing$이다.
\end{corollary}

실제로 가정에 따라 올 $f^{-1}(y)$는 유한 이산 공간이고, 따라서
따름정리는 이 올의 점 개수에 관한 귀납법과
(\hyperref[II.6.2.5-ko]{6.2.5})에서 따른다.

\begin{remark}[6.2.7]
\label{II.6.2.7-ko}
제V장에서 $Y$가 국소 뇌터이면 분리 준유한 사상 $X\to Y$가
반드시 준아핀 사상임을 보게 될 것이다.
\end{remark}

\subsection{준스킴의 정수적 폐포.}
\label{subsection:II.6.3-ko}

\begin{proposition}[6.3.1]
\label{II.6.3.1-ko}
$(X,\mathcal{A})$를 환 달린 공간, $\mathcal{B}$를 (가환)
$\mathcal{A}$-대수, $f$를 $X$ 위의 $\mathcal{B}$의 단면이라 하자.
다음 성질들은 서로 동치이다~:
\begin{enumerate}
\item[a)] $n\geq 0$인 $f^n$들로 생성되는 $\mathcal{B}$의
부분-$\mathcal{A}$-가군 (\hyperref[0.5.1.1-ko]{0, 5.1.1})은 유한
생성이다.
\item[b)] $\mathcal{B}$의 부분-$\mathcal{A}$-대수 $\mathcal{C}$가
존재하여, $\mathcal{C}$는 유한 생성 $\mathcal{A}$-가군이고
$f\in\Gamma(X,\mathcal{C})$이다.
\item[c)] 모든 $x\in X$에 대하여 $f_x$는 줄기 $\mathcal{A}_x$ 위에서
정수적이다.
\end{enumerate}
\end{proposition}

\oldpage[II]{117}

$f^n$들로 생성되는 $\mathcal{B}$의 부분-$\mathcal{A}$-가군은
$\mathcal{A}$-대수이므로 a)가 b)를 함의함은 명백하다. 한편 b)는 모든
$x\in X$에 대하여 $\mathcal{A}_x$-가군 $\mathcal{C}_x$가 유한
생성임을 함의하고, 따라서 대수 $\mathcal{C}_x$의 모든 원소, 특히
$f_x$는 $\mathcal{A}_x$ 위에서 정수적이다. 끝으로 어떤 $x\in X$에
대하여 다음 꼴의 관계가 성립한다고 하자.
\[
f_x^n+(a_1)_xf_x^{n-1}+\ldots+(a_n)_x=0
\]
여기서 $a_i$ ($1\leq i\leq n$)들은 $x$의 열린 근방 $U$ 위의
$\mathcal{A}$의 단면들이다. 그러면 단면
$f^n|U+a_1\cdot f^{n-1}|U+\ldots+a_n$은 $x$의 어떤 근방
$V\subset U$ 위에서 영이다. 따라서 모든 $f^k|V$ ($k\geq0$)는
$0\leq j\leq n-1$인 $f^j|V$들의 $\Gamma(V,\mathcal{A})$-계수 선형결합이고,
이로부터 c)가 a)를 함의함을 얻는다.

(\hyperref[II.6.3.1-ko]{6.3.1})의 동치 조건들이 성립할 때 단면 $f$가
\emph{$\mathcal{A}$ 위에서 정수적}이라고 한다.

\begin{corollary}[6.3.2]
\label{II.6.3.2-ko}
(\hyperref[II.6.3.1-ko]{6.3.1})의 가정 아래에서 $\mathcal{B}$의
부분-$\mathcal{A}$-가군 $\mathcal{A}'$이 (유일하게) 존재하여, 모든
$x\in X$에 대하여 $\mathcal{A}'_x$은 $\mathcal{A}_x$ 위에서 정수적인
싹 $f_x\in\mathcal{B}_x$들의 집합이다. 모든 열린집합 $U\subset X$에
대하여 $U$ 위의 $\mathcal{A}'$의 단면들은 $\mathcal{A}|U$ 위에서
정수적인 단면 $f\in\Gamma(U,\mathcal{B})$들이다.
\end{corollary}

$\Gamma(U,\mathcal{A}')$을 모든 $x\in U$에 대하여 $f_x$가
$\mathcal{A}_x$ 위에서 정수적인 $f\in\Gamma(U,\mathcal{B})$들의 집합으로
정하면 $\mathcal{A}'$의 존재는 즉시 따른다. 두 번째 주장은
(\hyperref[II.6.3.1-ko]{6.3.1})에서 곧바로 따른다.

$\mathcal{A}'$은 명백히 $\mathcal{B}$의 부분-$\mathcal{A}$-대수이다~;
이를 $\mathcal{B}$ 안에서 $\mathcal{A}$의 \emph{정수적 폐포}라고 한다.

\begin{env}[6.3.3]
\label{II.6.3.3-ko}
$(X,\mathcal{A})$, $(Y,\mathcal{B})$를 두 환 달린 공간,
\[
g=(\psi,\theta):X\to Y
\]
를 사상이라 하자. $\mathcal{C}$ (각각 $\mathcal{D}$)를
$\mathcal{A}$-대수 (각각 $\mathcal{B}$-대수)라 하고,
\[
u:\mathcal{D}\to\mathcal{C}
\]
를 $g$-사상이라 하자 (\hyperref[0.4.4.1-ko]{0, 4.4.1}). 이때
$\mathcal{A}'$ (각각 $\mathcal{B}'$)이 $\mathcal{C}$ (각각
$\mathcal{D}$) 안에서 $\mathcal{A}$ (각각 $\mathcal{B}$)의 정수적
폐포이면, $u$의 $\mathcal{B}'$로의 \emph{제한}은 $g$-사상
\[
u':\mathcal{B}'\to\mathcal{A}'.
\]
이다.

실제로 $j$가 표준 단사 사상 $\mathcal{B}'\to\mathcal{D}$이면, 다음 사상이
\[
v=u^\sharp\circ g^*(j):g^*(\mathcal{B}')\to\mathcal{C}
\]
$g^*(\mathcal{B}')$을 $\mathcal{A}'$ 안으로 보냄을 보이면 충분하다. 그런데
$(g^*(\mathcal{B}'))_x=\mathcal{B}'_{\psi(x)}\otimes_{\mathcal{B}_{\psi(x)}}
\mathcal{A}_x$의 원소는 $\mathcal{B}'$의 정의에 따라
$\mathcal{A}_x$ 위에서 정수적이고, 따라서 그 $v_x$에 의한 상도 그러하다.
이로써 주장이 증명된다.
\end{env}

\begin{proposition}[6.3.4]
\label{II.6.3.4-ko}
$X$를 준스킴, $\mathcal{A}$를 준연접 $\mathcal{O}_X$-대수라 하자.
그러면 $\mathcal{A}$ 안에서 $\mathcal{O}_X$의 정수적 폐포
$\mathcal{O}'_X$은 준연접 $\mathcal{O}_X$-대수이고, $X$의 모든 아핀
열린집합 $U$에 대하여 $\Gamma(U,\mathcal{O}'_X)$는
$\Gamma(U,\mathcal{A})$ 안에서 $\Gamma(U,\mathcal{O}_X)$의 정수적
폐포이다.
\end{proposition}

$X=\operatorname{Spec}(B)$가 아핀이고 $\mathcal{A}=\widetilde{A}$이며,
$A$가

\oldpage[II]{118}

$B$-대수인 경우만 생각하면 된다~; $B'$을 $A$ 안에서 $B$의 정수적
폐포라 하자. 모든 $x\in X$에 대하여 $B_x$ 위에서 정수적인 $A_x$의
원소가 반드시 $B'_x$에 속함을 보이면 충분한데, 이는 가환환 $C$에
대하여 $C$-대수 안에서 정수적 폐포를 취하는 연산과 ($C$의 곱셈적
부분집합에 대하여) 분수환으로 이행하는 연산이 가환한다는 사실
(\cite[t.~I, p.~261 et 257]{I-13})에서 따른다.

이때 $X$-스킴 $X'=\operatorname{Spec}(\mathcal{O}'_X)$을
\emph{$\mathcal{A}$에 대한 $X$의 정수적 폐포}(또는
\emph{$\operatorname{Spec}(\mathcal{A})$에 대한 정수적 폐포})라 부른다~;
$X'$이 $X$ 위에서 \emph{정수적}임은 명백하다
(\hyperref[II.6.1.2-ko]{6.1.2}).

(\hyperref[II.6.3.4-ko]{6.3.4})에서 곧바로 다음을 얻는다. $f:X'\to X$가
구조 사상이면, $X$의 모든 열린집합 $U$에 대하여 $f^{-1}(U)$는
\emph{$\mathcal{A}|U$에 대한, $U$ 위에 $X$가 유도하는 준스킴의 정수적
폐포}이다.

\begin{env}[6.3.5]
\label{II.6.3.5-ko}
$X$, $Y$를 두 준스킴, $f:X\to Y$를 사상, $\mathcal{A}$(각각
$\mathcal{B}$)를 준연접 $\mathcal{O}_X$-대수(각각
$\mathcal{O}_Y$-대수), $u:\mathcal{B}\to\mathcal{A}$를 $f$-사상이라
하자. 우리는 (\hyperref[II.6.3.3-ko]{6.3.3})에서 이로부터
$f$-사상 $u':\mathcal{O}'_Y\to\mathcal{O}'_X$을 얻음을 보았다. 여기서
$\mathcal{O}'_X$(각각 $\mathcal{O}'_Y$)는 $\mathcal{A}$(각각
$\mathcal{B}$) 안에서 $\mathcal{O}_X$(각각 $\mathcal{O}_Y$)의 정수적
폐포이다. 따라서 $X'$(각각 $Y'$)이 $\mathcal{A}$(각각
$\mathcal{B}$)에 대한 $X$(각각 $Y$)의 정수적 폐포이면, $u$로부터
도식
\[
\label{II.6.3.5.1-ko}
\xymatrix{
X'\ar[r]^{f'}\ar[d]&Y'\ar[d]\\
X\ar[r]_{f}&Y
}
\tag{6.3.5.1}
\]
을 가환하게 하는 사상 $f'=\operatorname{Spec}(u'):X'\to Y'$을 표준적으로
얻는다 (\hyperref[II.1.5.6-ko]{1.5.6}).
\end{env}

\begin{env}[6.3.6]
\label{II.6.3.6-ko}
$X$가 \emph{유한} 개의 기약 성분 $X_i$ ($1\leq i\leq r$)만을 가지며, 그
일반점들을 $\xi_i$라 하자. 특히 준연접 $\mathcal{R}(X)$-대수
$\mathcal{A}$에 대한 $X$의 정수적 폐포를 생각하자
($\mathcal{O}_X$-대수로 보든 $\mathcal{R}(X)$-대수로 보든 같은
뜻이다). (\hyperref[I.7.3.5-ko]{I, 7.3.5})에서 알 수 있듯이
$\mathcal{A}$는 $r$개의 준연접 $\mathcal{O}_X$-대수
$\mathcal{A}_i$의 직접곱이고, $\mathcal{A}_i$의 지지집합은 $X_i$에
포함되며, $\mathcal{A}_i$가 $X_i$ 위에 유도하는 층은 그 섬유 $A_i$가
$\mathcal{O}_{\xi_i}$ 위의 대수인 상수층이다. 따라서
(\hyperref[II.6.3.4-ko]{6.3.4}) $\mathcal{A}$ 안에서
$\mathcal{O}_X$의 정수적 폐포 $\mathcal{O}'_X$은 각
$\mathcal{A}_i$ 안에서 $\mathcal{O}_X$의 정수적 폐포
$\mathcal{O}^{(i)}_X$들의 \emph{직접곱}임이 명백하다. 그러므로
$\mathcal{A}$에 대한 $X$의 정수적 폐포
$X'=\operatorname{Spec}(\mathcal{O}'_X)$은
$\operatorname{Spec}(\mathcal{O}^{(i)}_X)=X'_i$ ($1\leq i\leq r$)들의
\emph{합}인 $X$-스킴이다.

더 나아가 $\mathcal{O}_X$-대수 $\mathcal{A}$가 \emph{축소}라고
가정하자. 이는 각 대수 $A_i$가 축소라는 것과 동치이며, 따라서
$A_i$를 \emph{체} $k(\xi_i)$ 위의 대수로 볼 수 있다(이 체는
바탕공간이 $X_i$인 $X$의 축소 부분준스킴의 유리함수체와 같다)~;
그러면 (\hyperref[II.1.3.8-ko]{1.3.8}) 각 $X'_i$는 \emph{축소}
$X$-스킴이고, $X'$은 또한 $X_{\mathrm{red}}$의 정수적 폐포이다. 더욱이
각 대수 $A_i$가 \emph{유한 개의 체 $K_{ij}$의 직접곱}
($1\leq j\leq s_i$)이라고 가정하자~; $\mathcal{K}_{ij}$가
$K_{ij}$에 대응하는 $\mathcal{A}_i$의 부분대수이면,
$\mathcal{O}^{(i)}_X$은 각 $\mathcal{K}_{ij}$ 안에서
$\mathcal{O}_X$의 정수적 폐포 $\mathcal{O}^{(ij)}_X$들의
\emph{직접곱}임이 명백하다. 따라서 $X'_i$는 $X$-스킴
$X'_{ij}=\operatorname{Spec}(\mathcal{O}^{(ij)}_X)$
($1\leq j\leq s_i$)들의 \emph{합}이다. 또한 이 가정들과 표기 아래에서
다음이 성립한다~:
\end{env}

\oldpage[II]{119}

\begin{proposition}[6.3.7]
\label{II.6.3.7-ko}
각 $X'_{ij}$는 정역이고 정규인 $X$-준스킴이며, 그 유리함수체
$R(X'_{ij})$는 $K_{ij}$ 안에서 $k(\xi_i)$의 대수적 폐포 $K'_{ij}$와
표준적으로 동일시된다.
\end{proposition}

앞의 논의에 따라 $X$가 정역이라고 가정할 수 있고, 따라서 $r=1$이며
$s_1=1$이므로 유일한 대수 $A_1$은 체 $K$이다~; $\xi$를 $X$의 일반점,
$f:X'\to X$를 구조 사상이라 하자. $X$의 공집합이 아닌 모든 아핀
열린집합 $U$에 대하여, $f^{-1}(U)$는 정역
$B_U=\Gamma(U,\mathcal{O}_X)$의 체 $K$ 안에서의 정수적 폐포 $B'_U$의
스펙트럼과 동일시된다 (\hyperref[II.6.3.4-ko]{6.3.4})~; 환 $B'_U$가
정역이고 정수적으로 닫혀 있으므로 그 스펙트럼의 점들의 국소환도
그러하며, 따라서 정의에 따라 $f^{-1}(U)$는 \emph{정역이고 정규인}
스킴이다 ((\hyperref[0.4.1.4-ko]{0, 4.1.4}) 및
(\hyperref[I.5.1.4-ko]{I, 5.1.4})). 또한 $B_U$의 소아이디얼 $(0)$ 위에
놓인 $B'_U$의 소아이디얼은 $(0)$뿐이므로
(\cite[t.~I, p.~259]{I-13}), $f^{-1}(\xi)$는 한 점 $\xi'$만으로
이루어지고, $k(\xi')$는 $B'_U$의 분수체 $K'$이며, 이는 다름 아닌
$K$ 안에서 $k(\xi)$의 대수적 폐포이다. 끝으로 $X'$은 기약이다. 실제로
$U$가 $X$의 공집합이 아닌 아핀 열린집합들을 달릴 때, $f^{-1}(U)$들은
기약 열린집합들로 이루어진 $X'$의 열린 덮개를 이룬다~; 또한 이 열린집합
둘의 교집합 $f^{-1}(U\cap V)$는 $\xi'$을 포함하므로 공집합이 아니며,
(\hyperref[0.2.1.4-ko]{0, 2.1.4})로 결론을 얻는다.

\begin{corollary}[6.3.8]
\label{II.6.3.8-ko}
$X$를 유한개의 기약성분 $X_i$ ($1\leq i\leq r$)만을 갖는 축소
준스킴이라 하고, $\xi_i$를 $X_i$의 일반점이라 하자. $\mathcal{R}(X)$에
대한 $X$의 정수적 폐포 $X'$은 정역이고 정규인 $r$개의 $X$-스킴
$X'_i$의 합이다. $f:X'\to X$가 구조 사상이면, $f^{-1}(\xi_i)$는
$X'_i$의 일반점 $\xi'_i$ 하나만으로 이루어지고
$k(\xi'_i)=k(\xi_i)$이며, 다시 말해 $f$는 쌍유리 사상이다.
\end{corollary}

이 경우 $X'$을 축소 준스킴 $X$의 \emph{정규화}라고 한다~; $f$는
쌍유리이고 정수적이므로 \emph{전사}임에 유의하자
(\hyperref[II.6.1.10-ko]{6.1.10}). 이때 $X'=X$가 성립할 필요충분조건은
$X$가 정규인 것이다. $X$가 정역 준스킴이면
(\hyperref[II.6.3.8-ko]{6.3.8})로부터 그 정규화 $X'$이 정역임을 얻는다.

\begin{env}[6.3.9]
\label{II.6.3.9-ko}
$X$, $Y$를 두 정역 준스킴, $f:X\to Y$를 우세 사상이라 하고,
$L=R(X)$, $K=R(Y)$를 각각 $X$와 $Y$의 유리함수체라 하자~; $f$에는
표준적으로 단사 준동형 $K\to L$이 대응하며, $K$ (각각 $L$)를 상수층
$\mathcal{R}(Y)$ (각각 $\mathcal{R}(X)$)와 동일시하면 이 단사 준동형은
$f$-사상이다. $K_1$ (각각 $L_1$)을 $K$ (각각 $L$)의 확대체라 하고,
도식
\[
\xymatrix{
K_1\ar[r]&L_1\\
K\ar[u]\ar[r]&L\ar[u]
}
\]
이 가환이 되게 하는 단사 준동형 $K_1\to L_1$이 주어졌다고 하자~;
$K_1$ (각각 $L_1$)을 $Y$ (각각 $X$) 위의 상수층, 따라서
$\mathcal{R}(Y)$-대수 (각각 $\mathcal{R}(X)$-대수)로 보면, 이는
$K_1\to L_1$이 $f$-사상이라는 뜻이다. 이제 $X'$ (각각 $Y'$)을
$L_1$ (각각 $K_1$)에 대한 $X$ (각각 $Y$)의 정수적 폐포라 하자.
$X'$ (각각 $Y'$)은 정역이고 정규인 준스킴이며
(\hyperref[II.6.3.6-ko]{6.3.6}), 그 유리함수체는 $L_1$ (각각 $K_1$)
안에서 $L$ (각각 $K$)의 대수적 폐포

\oldpage[II]{120}

$L'$ (각각 $K'$)와 표준적으로 동일시된다. 또한 도식
(\hyperref[II.6.3.5.1-ko]{6.3.5.1})을 가환이 되게 하는 표준 사상
$f':X'\to Y'$이 존재하며, 이 사상은 반드시 우세하다.

가장 중요한 경우는 $L_1=L$로 잡고, 따라서 $K_1$은 $L$에 포함되는
$K$의 확대체이며, $X$가 정역이고 \emph{정규}라고 가정하여 $X'=X$가
되는 경우이다. 앞의 결과에 따르면 $X$가 정규이고 $Y'$이
$K_1\subset L=R(X)$에 대한 $Y$의 정수적 폐포일 때, 모든 우세 사상
$f:X\to Y$는
\[
f:X\xrightarrow{f'}Y'\to Y
\]
로 인수분해되며, 여기서 $f'$은 우세하다~; 더구나 단사 준동형
$K_1\to L$이 주어지면 $f'$은 반드시 유일하다. 이는 $X$와 $Y$가
아핀인 경우로 귀착시키면 알 수 있다. 따라서 $Y$, $L$, 그리고
$K$-단사 준동형 $K_1\to L$이 주어졌을 때, $K_1$에 대한 $Y$의 정수적
폐포 $Y'$은 \emph{보편 문제의 해}이다.
\end{env}

\begin{remark}[6.3.10]
\label{II.6.3.10-ko}
(\hyperref[II.6.3.6-ko]{6.3.6})의 가정으로 돌아가, 각 대수 $A_i$가
$k(\xi_i)$ 위에서 \emph{유한 차원}이라고 더 가정하자(따라서 $A_i$는
유한 개의 체의 직접곱으로 분해된다)~; 이때 어떤 경우에는 구조 사상
$X'\to X$가 정수적일 뿐 아니라 \emph{유한}하기까지 하다고 말할 수
있다. $X$가 \emph{축소}인 경우만 다루자~; 문제는 $X$ 위에서
국소적이므로, 더 나아가 $X$가 환 $C$의 아핀 스킴이고 $C$는 유한 개의
극소 소아이디얼 $\mathfrak{p}_i$ ($1\leq i\leq r$)만을 가지며
$C_i=C/\mathfrak{p}_i$들이 정역이라고 가정할 수 있다~; 각 $C_i$의
분수체의 모든 유한 차수 확대체 안에서 $C_i$의 정수적 폐포가 유한
생성 $C$-가군이면, $X'$은 $X$ 위에서 유한하다
(\hyperref[II.6.3.4-ko]{6.3.4}). 이 조건은 $C$가 \emph{체 위의 유한
생성 대수}일 때 (\cite[t.~I, p.~267, th.~9]{I-13}), 또는
\emph{$\mathbf{Z}$ 위의 유한 생성 대수}일 때
(\cite[t.~I, p.~93, th.~3]{I-9}), 또는 \emph{완비 뇌터 국소환 위의
유한 생성 대수}일 때 ([25], p.~298) 언제나 성립한다고 알려져 있다.
따라서 $X$가 체 위의 \emph{유한형 스킴}, 또는 $\mathbf{Z}$ 위의
유한형 스킴, 또는 완비 뇌터 국소환 위의 유한형 스킴이면
$X'\to X$는 유한 사상이다.
\end{remark}

\subsection{$\mathcal{O}_X$-가군의 자기준동형의 행렬식.}
\label{subsection:II.6.4-ko}

\begin{env}[6.4.1]
\label{II.6.4.1-ko}
$A$를 (가환)환, $E$를 계수 $n$인 자유 $A$-가군, $u$를 $E$의
자기준동형이라 하자~; $u$의 \emph{특성다항식}을 정의하기 위해서는
계수 $n$인 자유 $A[T]$-가군 $E\otimes_AA[T]$ ($T$는 부정원)의
자기준동형 $u\otimes1$을 생각하고 다음과 같이 놓는다는 것을 상기하자.
\[
\label{II.6.4.1.1-ko}
P(u,T)=\det(T\cdot I-(u\otimes1))
\tag{6.4.1.1}
\]
(여기서 $I$는 $E\otimes_AA[T]$의 항등 자기동형사상이다). 다음이
성립한다.
\[
\label{II.6.4.1.2-ko}
P(u,T)=T^n-\sigma_1(u)T^{n-1}+\ldots+(-1)^n\sigma_n(u)
\tag{6.4.1.2}
\]
여기서 $\sigma_i(u)$는 $A$의 원소이며, $E$의 임의의 기저에 대한
$u$의 행렬 원소들에 관하여 차수 $i$인 (정수 계수) 동차다항식과 같다~;
$\sigma_i(u)$들을 $u$의 \emph{기본대칭함수}라 한다. 특히
$\sigma_1(u)=\operatorname{Tr}u$이고 $\sigma_n(u)=\det u$이다.
케일리--해밀턴 정리에 따라 다음이 성립함을 상기하자.
\[
\label{II.6.4.1.3-ko}
P(u,u)=u^n-\sigma_1(u)u^{n-1}+\ldots+(-1)^n\sigma_n(u)=0
\tag{6.4.1.3}
\]

\oldpage[II]{121}

이는 다음과 같이도 쓸 수 있다.
\[
\label{II.6.4.1.4-ko}
(\det u)\cdot I_E=uQ(u)
\tag{6.4.1.4}
\]
($I_E$는 $E$의 항등 자기동형사상이다). 여기서
\[
\label{II.6.4.1.5-ko}
Q(u)=(-1)^{n+1}(u^{n-1}-\sigma_1(u)u^{n-2}+\ldots+(-1)^{n-1}\sigma_{n-1}(u)).
\tag{6.4.1.5}
\]

$\varphi:A\to B$를 환의 준동형이라 하여 $B$를 $A$-대수로 만들자~;
계수 $n$인 자유 $B$-가군 $E_{(B)}=E\otimes_AB$와 $u$를
$E_{(B)}$의 자기준동형으로 확장한 $u\otimes1$을 생각하자. 모든 지표
$i$에 대하여 $\sigma_i(u\otimes1)=\varphi(\sigma_i(u))$임은 명백하다.
\end{env}

\begin{env}[6.4.2]
\label{II.6.4.2-ko}
이제 $A$가 분수체 $K$를 갖는 \emph{정역}이고, $E$가 (반드시 자유일
필요는 없는) \emph{유한형} $A$-가군이라고 가정하자. $n$을 $E$의
\emph{계수}, 즉 자유 $K$-가군 $E\otimes_AK$의 계수라 하자~; $E$의
각 자기준동형 $u$에는 $E\otimes_AK$의 자기준동형 $u\otimes1$이
표준적으로 대응한다~; 용어를 남용하여 $u$의 \emph{특성다항식}이라
부르고 $P(u,T)$로 나타내는 것은 다항식 $P(u\otimes1,T)$이며, 그
계수 $\sigma_i(u\otimes1)$ ($K$에 속한다)도 여전히 $\sigma_i(u)$로
나타내고 $u$의 \emph{기본대칭함수}라 부른다~; 특히 정의에 따라
$\det u=\det(u\otimes1)$이다. 이 표기 아래에서 공식
(\hyperref[II.6.4.1.3-ko]{6.4.1.3})부터
(\hyperref[II.6.4.1.5-ko]{6.4.1.5})까지는 의미가 있으며 여전히
성립한다. 여기서 $u^j$는 표준 준동형 $x\mapsto x\otimes1$과
$E\otimes_AK$의 자기준동형
$u^j\otimes1=(u\otimes1)^j$을 합성하여 얻는 준동형
$E\to E\otimes_AK$로 해석한다.

$F$가 $E$의 꼬임 부분가군이고 $E_0=E/F$이면 $u(F)\subset F$이므로,
몫으로 내려가면 $u$는 $E_0$의 자기준동형 $u_0$를 준다~; 더 나아가
$E\otimes_AK$는 $E_0\otimes_AK$와 동일시되고 $u\otimes1$은
$u_0\otimes1$과 동일시되므로, $1\leq i\leq n$에 대하여
$\sigma_i(u)=\sigma_i(u_0)$이다.

$E$에 \emph{꼬임이 없을} 때 $E$는 $E\otimes_AK$의 $A$-부분가군과
표준적으로 동일시되며, 관계 $u\otimes1=0$은 $u=0$과 동치이다.
$E$가 자유 $A$-가군일 때
(\hyperref[II.6.4.1-ko]{6.4.1})과
(\hyperref[II.6.4.2-ko]{6.4.2})에서 주어진 $\sigma_i(u)$의 두 정의는
앞의 사실에 따라 일치하며, 이는 채택한 표기를 정당화한다.

$E$가 \emph{꼬임가군}, 다시 말해 $E_0=\{0\}$일 때 $E_0$의 외대수는
$K$로 축약되고, $E_0$의 유일한 자기준동형 $u_0$의 행렬식은
\emph{$1$과 같다}는 것도 주목하자.
\end{env}

\begin{proposition}[6.4.3]
\label{II.6.4.3-ko}
$A$를 정역, $E$를 유한형 $A$-가군, $u$를 $E$의 자기준동형이라 하자~;
그러면 $u$의 기본대칭함수 $\sigma_i(u)$ (특히 $\det u$)는 $A$ 위에서
정수적인 $K$의 원소이다.
\end{proposition}

$A'$을 $A$의 정수적 폐포라 하자~; $A'[T]$는 정수적으로 닫힌 환이므로
([24], p.~99), 분수체 $K(T)$ 안에서 $A[T]$의 정수적 폐포이다.
$u$를 $T\cdot I-u\otimes1$로, $A$를 $A[T]$로 바꾸면, $\det u$가
$A$ 위에서 정수적임을 증명하는 것으로 귀착됨을 알 수 있다. $n$이
$E$의 계수이면 $\det u=\det(\bigwedge^n u)$이고
$(\bigwedge^n u)\otimes1=\bigwedge^n(u\otimes1)$이므로, $n=1$이라고
가정할 수 있다. 그런데 이때 사상 $u\mapsto\det u$는 $A$-가군
$\operatorname{Hom}_A(E,E)$에서 $K$로 가는 준동형이다~; $E$가
유한형이므로 $\operatorname{Hom}_A(E,E)$는 유한형 $A$-가군 $E^n$의
한 부분가군과 동형이다($E$가 $n$개의 생성원으로 된 생성계를 갖는
경우)~; 따라서 원소 $\det u$들은 $K$의 한 유한형 $A$-부분가군에
속하고, 그러므로 $A$ 위에서 정수적이다.

\oldpage[II]{122}

\begin{corollary}[6.4.4]
\label{II.6.4.4-ko}
(\hyperref[II.6.4.3-ko]{6.4.3})의 가정 아래에서, 더 나아가 $A$가
정규환이라고 가정하면 $\sigma_i(u)$ (특히 $\operatorname{Tr}u$와
$\det u$)는 $A$에 속한다.
\end{corollary}

\begin{proposition}[6.4.5]
\label{II.6.4.5-ko}
$A$를 정역, $E$를 계수 $n$인 유한형 $A$-가군, $u$를 $E$의
자기준동형이라 하고, $\sigma_i(u)$들이 $A$에 속한다고 하자. $u$가
$E$의 자기동형사상이기 위해서는 $\det u$가 $A$에서 가역이어야 한다~;
$E$에 꼬임이 없을 때에는 이 조건이 충분하다.
\end{proposition}

이 조건은 \emph{충분하다}. 실제로 가정과 공식
(\hyperref[II.6.4.1.4-ko]{6.4.1.4}) 및
(\hyperref[II.6.4.1.5-ko]{6.4.1.5})는 ($E$에 꼬임이 없으므로
$E\otimes_AK$에서뿐 아니라 $E$에서도 유효하다)
$(\det u)^{-1}Q(u)$가 $u$의 역임을 증명한다.

이 조건은 \emph{필요하다}. 실제로 $u$가 가역이면
(\hyperref[II.6.4.3-ko]{6.4.3})에 따라 $\det(u^{-1})$은 $A$의
분수체 $K$ 안에서 $A$의 정수적 폐포 $A'$에 속하고, 명백히 $A'$에서
$\det u$의 역이다. 이제 우리의 주장은 다음 보조정리에서 따른다.

\begin{lemma}[6.4.5.1]
\label{II.6.4.5.1-ko}
$A$를 환 $A'$의 부분환이라 하고, $A'$가 $A$ 위에서 정수적이라고 하자.
원소 $x\in A$가 $A'$에서 가역이면, $A$에서도 가역이다.
\end{lemma}

그렇지 않으면 $x$는 $A$의 어떤 극대 아이디얼 $\mathfrak{m}$에 속할
것이고, 제1 코언--자이덴베르크 정리
(\cite[t.~I, p.~257, th.~3]{I-13})에 따라
$\mathfrak{m}=\mathfrak{m}'\cap A$를 만족하는 $A'$의 극대 아이디얼
$\mathfrak{m}'$이 존재할 것이다~; 그러면 $x\in\mathfrak{m}'$이 되어
모순이다.

\begin{corollary}[6.4.6]
\label{II.6.4.6-ko}
$A$를 정수적으로 닫힌 정역, $E$를 꼬임 없는 유한형 $A$-가군,
$u$를 $E$의 자기준동형이라 하자. $u$가 $E$의 자기동형사상일
필요충분조건은 $\det u$가 $A$에서 가역인 것이다.
\end{corollary}

이는 (\hyperref[II.6.4.4-ko]{6.4.4})와
(\hyperref[II.6.4.5-ko]{6.4.5})에서 따른다.

\begin{remark}[6.4.7]
\label{II.6.4.7-ko}
뒤에서 앞의 결과들을 일반화한 것이 필요할 것이다. 뇌터
\emph{축소}환 $A$를 생각하자~; $\mathfrak{p}_\alpha$
($1\leq\alpha\leq r$)를 그 극소 소아이디얼들, $K_\alpha$를 정역
$A/\mathfrak{p}_\alpha$의 분수체, $K$를 $A$의 전분수환, 즉 체
$K_\alpha$들의 \emph{직접곱}이라 하자. $E$를 유한형 $A$-가군이라
하고, \emph{$E\otimes_AK$가 계수 $n$인 자유 $K$-가군이라고 가정하자}
(여기서는 이것이 더 이상 다른 가정들의 귀결이 아니다)~; 이는 모든
$K_\alpha$-벡터공간 $E\otimes_AK_\alpha=E_\alpha$가 \emph{같은 차원
$n$}을 갖는다는 것과 동치이다~; 이때 $u$가 $E$의 자기준동형이면 다시
$P(u,T)=P(u\otimes1,T)$ 및 $\sigma_j(u)=\sigma_j(u\otimes1)$로 놓고,
특히 $\det u=\det(u\otimes1)$로 놓는다~; 따라서 $\sigma_j(u)$들은
$K$의 원소이다. $E\otimes_AK$가 $E_\alpha$들의 직합이고 이들이
$u\otimes1$에 대하여 안정적임은 즉시 알 수 있다~; 더욱이
$u\otimes1$을 $E_\alpha$에 제한한 것은 바로 $u$를 이
$K_\alpha$-벡터공간으로 확장한 $u_\alpha$이다~; 이로부터
$\sigma_j(u)$는 각 $K_\alpha$에서의 성분이 $\sigma_j(u_\alpha)$인
$K$의 원소임을 얻는다. $K$ 안에서 $A$의 정수적 폐포는 각
$K_\alpha$ 안에서 $A$의 정수적 폐포들의 \emph{직접곱}이므로,
$\sigma_j(u)$들은 $A$ 위에서 \emph{정수적}이다.

\begin{lemma}[6.4.7.1]
\label{II.6.4.7.1-ko}
$u$가 $\operatorname{Hom}_A(E,E)$를 달릴 때 모든 원소 $\sigma_j(u)$
($1\leq j\leq n$)에 의해 생성되는 $K$의 부분-$A$-대수는 유한 생성
$A$-가군이다.
\end{lemma}

$P(u,T)$들에 의해 생성되는 $K[T]$의 부분-$A[T]$-대수가 유한 생성
$A[T]$-가군임을 보이면 충분하다. 실제로 $F_i(T)$ ($1\leq i\leq m$)가
이 $A[T]$-가군의 생성계를 이루면, $F_i(T)$들의 계수들은 $A$ 위에서
정수적이므로 이 계수들이 생성하는 $A$-대수는 유한 생성 $A$-가군이다
(\cite[t.~I, p.~255, th.~1]{I-13}). 따라서

\oldpage[II]{123}

$A$를 $A[T]$로(이는 뇌터 환이다), $E$를
$E\otimes_AA[T]=E'$으로 바꿀 수 있다. 이때
$E'\otimes_{A[T]}K[T]=E\otimes_AK[T]$는 계수 $n$인 자유
$K[T]$-가군이다.

처음 표기로 돌아가면, $u$가 $\operatorname{Hom}_A(E,E)$를 달릴 때
원소 $\det u$들이 생성하는 $A$-가군이 유한 생성임을 증명하면 충분함을
알 수 있다~; \emph{더 강하게} (유한 생성 $A$-가군의 모든 부분가군은
유한 생성이므로), $v$가 $\bigwedge^nE$의 자기준동형 전체를 달릴 때
$\det v$들이 생성하는 $A$-가군이 유한 생성임을 증명하면 충분하다~;
다시 말해 우리는 또다시 $n=1$인 경우로 귀착된다. 그런데 이 경우 주장은
$\operatorname{Hom}_A(E,E)$가 유한 생성 $A$-가군이고
$v\mapsto\det v$가 이 $A$-가군에서 $K$로 가는 준동형이라는 사실에서
따른다.

$F$를 표준 준동형 $E\to E\otimes_AK$의 핵이라 하고 $E_0=E/F$라 하자~;
위와 같이 $E\otimes_AK$가 $E_0\otimes_AK$와 동일시되고,
$u(F)\subset F$이며, $u_0$가 몫으로 내려가 $u$에서 얻어지는 $E_0$의
자기준동형이면 $u\otimes1$이 $u_0\otimes1$과 동일시되고 모든 $j$에
대하여 $\sigma_j(u)=\sigma_j(u_0)$임을 알 수 있다. \emph{$F=0$이면},
$E$를 $E\otimes_AK$의 부분가군과 동일시하고 $u^j$들을
준동형 $E\to E\otimes_AK$와 동일시할 때 공식
(\hyperref[II.6.4.1.3-ko]{6.4.1.3})과
(\hyperref[II.6.4.1.5-ko]{6.4.1.5})는 의미가 있으며 성립한다~; 따라서
명제 (\hyperref[II.6.4.5-ko]{6.4.5})도 그 증명과 함께 이 경우로
확장된다.
\end{remark}

\begin{env}[6.4.8]
\label{II.6.4.8-ko}
$(X,\mathcal{A})$를 환 달린 공간, $\mathcal{E}$를 (유한 계수의)
\emph{국소 자유} $\mathcal{A}$-가군, $u$를 $\mathcal{E}$의
자기준동형이라 하자. 가정에 따라 $X$의 위상의 기저 $\mathfrak{B}$가
존재하여, 모든 $V\in\mathfrak{B}$에 대하여 $\mathcal{E}|V$는
$\mathcal{A}^n|V$와 동형이다(여기서 $n$은 $V$에 따라 달라질 수
있다). $u$를 $\mathcal{E}$의 자기준동형이라 하자~; 그러면 모든
$V\in\mathfrak{B}$에 대하여 $u_V$는
$\Gamma(V,\mathcal{A})$-가군 $\Gamma(V,\mathcal{E})$의
자기준동형이고, 이 가군은 가정에 따라 자유이다~; 따라서 결정식
$\det u_V$가 정의되며 $\Gamma(V,\mathcal{A})$에 속한다. 더욱이
$e_1,\ldots,e_n$이 $\Gamma(V,\mathcal{E})$의 기저를 이루면, 임의의
열린집합 $W\subset V$로의 그 제한들은 $\Gamma(W,\mathcal{A})$ 위에서
$\Gamma(W,\mathcal{E})$의 기저를 이루므로, $\det u_W$는
$\det u_V$의 $W$로의 제한이다. 따라서 $X$ 위의 $\mathcal{A}$의
유일한 절단이 존재하여, 이를 $\det u$로 표기하고
\emph{$u$의 결정식}이라 부르겠다. 이 절단은 모든
$V\in\mathfrak{B}$로의 $\det u$의 제한이 $\det u_V$가 되는 절단이다. 모든
$x\in X$에 대하여 $(\det u)_x=\det u_x$임은 명백하다~; $\mathcal{E}$의
두 자기준동형 $u$, $v$에 대하여
\[
\label{II.6.4.8.1-ko}
\det(u\circ v)=(\det u)(\det v)
\tag{6.4.8.1}
\]
이며, 또한
\[
\label{II.6.4.8.2-ko}
\det(1_{\mathcal{E}})=1_{\mathcal{A}}
\tag{6.4.8.2}
\]
이다. 그리고 $\mathcal{E}$의 계수가 일정하고(이는 $X$가 연결이면
(\hyperref[0.5.4.1-ko]{0, 5.4.1})에 따라 성립한다) 그 값이 $n$이면,
\[
\label{II.6.4.8.3-ko}
\det(s\cdot u)=s^n\det u
\tag{6.4.8.3}
\]
가 모든 $s\in\Gamma(X,\mathcal{A})$에 대하여 성립한다
($n\geq1$이면 $\det(0)=0_{\mathcal{A}}$이지만, $n=0$이면
$\det(0)=1_{\mathcal{A}}$임에 유의하자). 더욱이 $u$가
$\mathcal{E}$의 \emph{자기동형}일 필요충분조건은 $\det u$가
$\Gamma(X,\mathcal{A})$의 \emph{가역원}인 것이다.

$\mathcal{E}$의 계수가 일정하면 같은 방식으로 기본 대칭함수
$\sigma_i(u)$를 정의할 수 있으며, 이들은 $\Gamma(X,\mathcal{A})$의
원소이다~; 이때도 (\hyperref[II.6.4.1.3-ko]{6.4.1.3})부터
(\hyperref[II.6.4.1.5-ko]{6.4.1.5})까지의 관계들이 성립한다.

\oldpage[II]{124}

이로써 곱셈 모노이드의 준동형
$\det:\operatorname{Hom}_{\mathcal{A}}(\mathcal{E},\mathcal{E})
\to\Gamma(X,\mathcal{A})$를 정의하였다~; 정의에 따라
$\operatorname{Hom}_{\mathcal{A}}(\mathcal{E},\mathcal{E})=
\Gamma(X,\mathcal{H}om_{\mathcal{A}}(\mathcal{E},\mathcal{E}))$임에
유의하면, 이 정의에서 $X$를 $X$의 임의의 열린집합 $U$로 바꿀 수
있음을 알 수 있다. 이로부터 곱셈 모노이드의 층의 준동형
$\det:\mathcal{H}om_{\mathcal{A}}(\mathcal{E},\mathcal{E})\to\mathcal{A}$가
곧바로 정의된다. $\mathcal{E}$의 계수가 일정할 때에는 마찬가지로
집합의 층의 준동형
$\sigma_i:\mathcal{H}om_{\mathcal{A}}(\mathcal{E},\mathcal{E})
\to\mathcal{A}$를 정의한다~; $i=1$일 때 준동형
$\sigma_1=\operatorname{Tr}$은 $\mathcal{A}$-가군의 준동형이다.

$(Y,\mathcal{B})$를 또 하나의 환 달린 공간,
$f:(X,\mathcal{A})\to(Y,\mathcal{B})$를 환 달린 공간의 사상이라 하자~;
$\mathcal{F}$가 국소 자유 $\mathcal{B}$-가군이면,
$f^*(\mathcal{F})$는 국소 자유 $\mathcal{A}$-가군이다(더 나아가
$\mathcal{F}$의 계수가 일정하면 이와 같은 계수를 갖는다)
(\hyperref[0.5.4.5-ko]{0, 5.4.5}). $\mathcal{F}$의 임의의 자기준동형
$v$에 대하여 $f^*(v)$는 $f^*(\mathcal{F})$의 자기준동형이며, 정의에서
곧바로 $\det f^*(v)$는 $\det v\in\Gamma(Y,\mathcal{B})$에 표준적으로
대응하는 $X$ 위의 $\mathcal{A}=f^*(\mathcal{B})$의 절단임을 얻는다.
또한 준동형
$f^*(\det):f^*(\mathcal{H}om_{\mathcal{B}}(\mathcal{F},\mathcal{F}))
\to f^*(\mathcal{B})=\mathcal{A}$가 다음 합성이라고 말할 수도 있다.
\[
f^*(\mathcal{H}om_{\mathcal{B}}(\mathcal{F},\mathcal{F}))
\xrightarrow{\gamma^\sharp}
\mathcal{H}om_{\mathcal{A}}(f^*(\mathcal{F}),f^*(\mathcal{F}))
\xrightarrow{\det}\mathcal{A}
\]
(\hyperref[0.4.4.6-ko]{0, 4.4.6}). $\sigma_i$들에 대해서도 이와 같은
결과들이 성립한다.
\end{env}

\begin{env}[6.4.9]
\label{II.6.4.9-ko}
이제 $X$가 \emph{국소 정역 준스킴}이라고 가정하자. 그러면 $X$ 위의
유리함수층 $\mathcal{R}(X)$는 국소 상수인 체층이고
(\hyperref[I.7.3.4-ko]{I, 7.3.4}), $\mathcal{O}_X$-가군으로서는
준연접이다. $\mathcal{E}$가 \emph{유한형} 준연접
$\mathcal{O}_X$-가군이면,
$\mathcal{E}'=\mathcal{E}\otimes_{\mathcal{O}_X}\mathcal{R}(X)$는
국소 자유 $\mathcal{R}(X)$-가군이다
(\hyperref[I.7.3.6-ko]{I, 7.3.6})~; $\mathcal{E}$의 모든 자기준동형
$u$에 대하여 $u\otimes1_{\mathcal{R}(X)}$는 $\mathcal{E}'$의
자기준동형이고, $\det(u\otimes1)$은 $X$ 위의 $\mathcal{R}(X)$의
단면이다. 이 단면을 역시 \emph{$u$의 행렬식}이라 하고 여전히
$\det u$로 나타낸다. (\hyperref[II.6.4.3-ko]{6.4.3})에 따라
$\det u$는 $\mathcal{R}(X)$ 안에서 $\mathcal{O}_X$의
\emph{정수적 폐포}의 단면이다
(\hyperref[II.6.3.2-ko]{6.3.2})~; 더 나아가 $X$가 \emph{정규}이면
$\det u$는 $X$ 위의 \emph{$\mathcal{O}_X$의 단면}이고, 또한
$\mathcal{E}$가 \emph{꼬임이 없다}고 가정하면
(\hyperref[II.6.4.6-ko]{6.4.6})에 따라 $u$가 $\mathcal{E}$의
자동사상일 필요충분조건은 $\det u$가 \emph{가역}인 것이다. 공식
(\hyperref[II.6.4.8.1-ko]{6.4.8.1})부터
(\hyperref[II.6.4.8.3-ko]{6.4.8.3})까지도 여전히 성립한다~;
$u\mapsto\det u$인 준동형을
$\mathcal{H}om_{\mathcal{O}_X}(\mathcal{E},\mathcal{E})$의
단면가군들에 적용하면 층의 준동형
$\det:\mathcal{H}om_{\mathcal{O}_X}(\mathcal{E},\mathcal{E})
\to\mathcal{R}(X)$를 얻으며, $X$가 정규이면 그 값은
$\mathcal{O}_X$에 속한다. $\mathcal{E}'$이 \emph{일정한 계수}를 가질
때에는 다른 기본대칭함수 $\sigma_j(u)$에 대해서도 이와 유사한 정의와
결과가 성립한다~; 더 나아가 $X$가 \emph{정규}이면 $\sigma_j(u)$들은
$X$ 위의 \emph{$\mathcal{O}_X$의 단면}이다.

마지막으로 $X$, $Y$를 두 정역 준스킴, $f:X\to Y$를 \emph{우세 사상}이라
하자. 이때 표준 준동형 $f^*(\mathcal{R}(Y))\to\mathcal{R}(X)$가
존재한다는 것을 알고 있다
(\hyperref[I.7.3.8-ko]{I, 7.3.8}\footnote{\emph{정오표 및 추가 사항},
목록 1 참조.}). 따라서 모든 유한형 준연접 $\mathcal{O}_Y$-가군
$\mathcal{F}$에 대하여 표준 준동형
$\theta:f^*(\mathcal{F}\otimes_{\mathcal{O}_Y}\mathcal{R}(Y))
\to f^*(\mathcal{F})\otimes_{\mathcal{O}_X}\mathcal{R}(X)$를 얻는다.
$v$가 $\mathcal{F}$의

\oldpage[II]{125}

자기준동형이면,
$f^*(v\otimes1_{\mathcal{R}(Y)})$는
$f^*(\mathcal{F}\otimes_{\mathcal{O}_Y}\mathcal{R}(Y))$의
자기준동형이고 다음 가환 그림을 얻는다.
\[
\xymatrix{
f^*(\mathcal{F}\otimes_{\mathcal{O}_Y}\mathcal{R}(Y))
  \ar[r]^{f^*(v\otimes1)}\ar[d]_{\theta}
& f^*(\mathcal{F}\otimes_{\mathcal{O}_Y}\mathcal{R}(Y))\ar[d]^{\theta}\\
f^*(\mathcal{F})\otimes_{\mathcal{O}_X}\mathcal{R}(X)
  \ar[r]_{f^*(v)\otimes1}
& f^*(\mathcal{F})\otimes_{\mathcal{O}_X}\mathcal{R}(X).
}
\]
이로부터 $\det f^*(v)$는 $\mathcal{R}(Y)$의 단면 $\det v$를
준동형 $f^*(\mathcal{R}(Y))\to\mathcal{R}(X)$로 보낸 표준 상임을 쉽게
알 수 있다~; 실제로 $X=\operatorname{Spec}(A)$와
$Y=\operatorname{Spec}(B)$가 아핀인 경우로 즉시 환원된다. 여기서 $A$와
$B$는 각각 분수체 $K$와 $L$을 갖는 정역이고, 준동형 $B\to A$는
단사이며 따라서 단사 준동형 $L\to K$로 확장된다~;
$\mathcal{F}=\widetilde{M}$이고 $M$이 유한형 $B$-가군이면,
$M\otimes_B L$의 $L$ 위의 계수는
$(M\otimes_B A)\otimes_A K$의 $K$ 위의 계수와 \emph{같고}, $M$의 모든
자기준동형 $u$에 대하여 $\det((u\otimes1)\otimes1)$은
$\det(u\otimes1)$의 $K$에서의 상이다. 이로써 결론이 따른다.
\end{env}

\begin{env}[6.4.10]
\label{II.6.4.10-ko}
끝으로 $X$가 \emph{국소 뇌터 축소 준스킴}이라고 가정하자. 그러면
$X$ 위의 유리함수층 $\mathcal{R}(X)$는 여전히 준연접
$\mathcal{O}_X$-가군이다
(\hyperref[I.7.3.4-ko]{I, 7.3.4})~; 한편 $\mathcal{E}$를
$\mathcal{E}'=\mathcal{E}\otimes_{\mathcal{O}_X}\mathcal{R}(X)$가
\emph{국소 자유} (유한 계수)인 \emph{연접} $\mathcal{O}_X$-가군이라
하자. (\hyperref[II.6.4.7-ko]{6.4.7})에 따라 $\mathcal{E}'$이 일정한
계수를 가지면, $\mathcal{E}$의 모든 자기준동형 $u$에 대하여
$X$ 위의 $\mathcal{R}(X)$의 단면인 $\sigma_j(u)$들을 정의할 수 있다.
$\mathcal{E}'$이 일정한 계수를 갖는다고 더 이상 가정하지 않을 때에도
준동형
$\det:\mathcal{H}om_{\mathcal{O}_X}(\mathcal{E},\mathcal{E})
\to\mathcal{R}(X)$를 정의할 수 있다.
\end{env}

\subsection{가역층의 노름.}
\label{subsection:II.6.5-ko}

\begin{env}[6.5.1]
\label{II.6.5.1-ko}
$(X,\mathcal{A})$를 환 달린 공간이라 하고, $\mathcal{B}$를 (교환)
$\mathcal{A}$-대수라 하자. $\mathcal{A}$-가군 $\mathcal{B}$는
$\mathcal{H}om_{\mathcal{A}}(\mathcal{B},\mathcal{B})$의 한 부분
$\mathcal{A}$-가군과 표준적으로 동일시된다. 이때 $X$의 열린집합 $U$
위의 $\mathcal{B}$의 단면 $f$는 그 단면을 곱하는 사상과 동일시된다.
$(X,\mathcal{A})$와 $\mathcal{B}$가
(\hyperref[II.6.4.8-ko]{6.4.8}), (\hyperref[II.6.4.9-ko]{6.4.9}) 또는
(\hyperref[II.6.4.10-ko]{6.4.10})에 열거된 조건 가운데 하나를
만족하면, 따라서 $\det(f)$를 (그리고 어떤 경우에는 $\sigma_j(f)$들을)
정의할 수 있다. 이들은 $U$ 위의 $\mathcal{A}$ 또는
$\mathcal{R}(X)$의 단면이며, 이를 $f$의 \emph{노름}(각각 $f$의
\emph{기본대칭함수})이라 하고 $N_{\mathcal{B}/\mathcal{A}}(f)$로
나타내겠다. 이제 다음 조건 가운데 \emph{하나}가 성립한다고 가정하자~:
\begin{enumerate}
\item[(I)] \emph{$\mathcal{B}$는 (유한 계수의) 국소 자유
$\mathcal{A}$-가군이다.}
\item[(II)] \emph{$(X,\mathcal{A})$는 국소 뇌터 축소 준스킴이고,
$\mathcal{B}$는 $\mathcal{B}\otimes_{\mathcal{A}}\mathcal{R}(X)$가
국소 자유 $\mathcal{R}(X)$-가군이 되도록 하는 연접
$\mathcal{A}$-가군이며, 열린집합 $U\subset X$ 위의 모든 단면
$f\in\Gamma(U,\mathcal{B})$에 대하여
$N_{\mathcal{B}/\mathcal{A}}(f)$는 $U$ 위의 $\mathcal{A}$의 단면이다.}
\end{enumerate}

\oldpage[II]{126}

준스킴 $X$가 국소 뇌터이고 \emph{정규}이면 이 마지막 조건은 자동으로
성립한다는 점에 유의하자(\hyperref[II.6.4.9-ko]{6.4.9}).

한편 $\mathcal{B}\otimes_{\mathcal{A}}\mathcal{R}(X)$가 국소 자유라는
가정은 다음과 같이 표현할 수도 있다. $X$의 기약 성분을 바탕공간으로
갖는 $X$의 축소 닫힌 부분준스킴들을 $X_\alpha$라 하자
(\hyperref[I.5.2.1-ko]{I, 5.2.1}). 이들은 국소 뇌터 정역 준스킴이다.
각 $x\in X$는 유한 개의 부분공간 $X_\alpha$에만 속한다. 한편
$\mathcal{B}\otimes_{\mathcal{A}}\mathcal{R}(X_\alpha)$는 상수 계수
$k_\alpha$를 갖는 국소 자유 $\mathcal{R}(X_\alpha)$-가군이다
(\hyperref[I.7.3.6-ko]{I, 7.3.6}).
$\mathcal{B}\otimes_{\mathcal{A}}\mathcal{R}(X)$가 국소 자유
$\mathcal{R}(X)$-가군이라는 것은 모든 $x\in X$에 대하여
\emph{$x\in X_\alpha$인 지표들에 대응하는 계수 $k_\alpha$들이
서로 같다}는 뜻이다. 실제로 이 문제는 국소적이므로
$X=\operatorname{Spec}(C)$인 경우로 귀착된다. 여기서 $C$는 뇌터 축소
환이고 $\mathcal{B}=\widetilde{D}$이며, $D$는 유한형 $C$-가군인
$C$-대수이다. $\mathfrak{p}_i$ ($1\leq i\leq m$)가 $C$의 극소 소
아이디얼들이라 하면, $C$의 전체 분수환 $L$은 정역
$C/\mathfrak{p}_i$들의 분수체 $K_i$들의 직접곱이고,
$D\otimes_C L$은 $D\otimes_C K_i$들의 직합이므로 결론이 따른다.

(I) 또는 (II)의 가정 아래에서는 이와 같이 \emph{곱셈 모노이드층의
준동형} $N_{\mathcal{B}/\mathcal{A}}:\mathcal{B}\to\mathcal{A}$를
정의했음이 분명하다. 혼동의 우려가 없으면 이를 $N$으로도 나타내고,
\emph{노름} 준동형이라 한다. 같은 열린집합 $U$ 위의
$\mathcal{B}$의 두 단면 $f$, $g$에 대하여
\[
\label{II.6.5.1.1-ko}
N_{\mathcal{B}/\mathcal{A}}(fg)
=N_{\mathcal{B}/\mathcal{A}}(f)N_{\mathcal{B}/\mathcal{A}}(g)
\tag{6.5.1.1}
\]
가 $U$ 위의 $\mathcal{A}$의 대응하는 단면들에 대하여 성립하고,
\[
\label{II.6.5.1.2-ko}
N_{\mathcal{B}/\mathcal{A}}(1_{\mathcal{B}})=1_{\mathcal{A}}~;
\tag{6.5.1.2}
\]
끝으로 $U$ 위의 $\mathcal{A}$의 모든 단면 $s$에 대하여
\[
\label{II.6.5.1.3-ko}
N_{\mathcal{B}/\mathcal{A}}(s\cdot1_{\mathcal{B}})=s^n
\tag{6.5.1.3}
\]
이 성립한다. 여기서 $\mathcal{B}$의 계수는 상수이고 $n$과 같다
($s=0_{\mathcal{A}}$이면 이 공식은 $n\geq1$일 때
$N(0_{\mathcal{B}})=0_{\mathcal{A}}$를, $n=0$일 때
$N(0_{\mathcal{B}})=N(1_{\mathcal{B}})=1_{\mathcal{A}}$를 준다).

가정 (I)에서 $f\in\Gamma(U,\mathcal{B})$가 가역이기 위한
필요충분조건은 $N(f)\in\Gamma(U,\mathcal{A})$가 가역인 것이다.
가정 (II)에서 이 조건은 필요조건이다. 또한
$\mathcal{B}\to\mathcal{B}\otimes_{\mathcal{A}}\mathcal{R}(X)$가
\emph{단사}이고 다음의 더 강한 가정이 성립하면 이 조건은
(\hyperref[II.6.4.7-ko]{6.4.7})에 따라 충분조건이기도 하다~:
\begin{enumerate}
\item[(II \emph{bis})] \emph{$(X,\mathcal{A})$는 국소 뇌터 축소
준스킴이고, $\mathcal{B}$는
$\mathcal{B}\otimes_{\mathcal{A}}\mathcal{R}(X)$가 국소 자유
$\mathcal{R}(X)$-가군이 되도록 하는 연접 $\mathcal{A}$-가군이며,
$\mathcal{B}\otimes_{\mathcal{A}}\mathcal{R}(X)|U$가
$\mathcal{R}(X)|U$ 위에서 상수 계수 $n$을 갖는 열린집합 위의 모든
단면 $f\in\Gamma(U,\mathcal{B})$에 대하여 $\sigma_j(f)$
($1\leq j\leq n$)는 $U$ 위의 $\mathcal{A}$의 단면이다.}
\end{enumerate}
($X$가 \emph{정규}이면 이 조건도 성립한다는 점에 다시 유의하자.)
\end{env}

\begin{env}[6.5.2]
\label{II.6.5.2-ko}
(\hyperref[II.6.5.1-ko]{6.5.1})의 가정 (I), (II) 가운데 하나가 성립한다고
가정하고, $\mathcal{L}'$을 \emph{가역} $\mathcal{B}$-가군이라 하자. 다음과
같이 이에 (유일한 동형사상까지) 표준적으로 \emph{가역}
$\mathcal{A}$-가군 $\mathcal{L}$을 대응시키겠다. $\mathcal{A}^*$ (각각
$\mathcal{B}^*$)를, 모든 열린집합 $U\subset X$에 대하여
$\Gamma(U,\mathcal{A}^*)$ (각각 $\Gamma(U,\mathcal{B}^*)$)가
$\Gamma(U,\mathcal{A})$ (각각 $\Gamma(U,\mathcal{B})$)의

\oldpage[II]{127}

가역원들의 집합이 되도록 하는 $\mathcal{A}$ (각각 $\mathcal{B}$)의
부분층이라 하자~; 이들은 곱셈군의 층이며, $\mathcal{B}^*$로 제한한
$N_{\mathcal{B}/\mathcal{A}}$는 군의 층의 \emph{준동형}
$\mathcal{B}^*\to\mathcal{A}^*$이다
(\hyperref[II.6.5.1-ko]{6.5.1}). 다음 성질을 갖는 쌍
$(U_\lambda,\eta_\lambda)$들의 집합을 $\mathfrak{L}$이라 하자~:
$U_\lambda$는 $X$의 열린집합이고, $\eta_\lambda$는
$(\mathcal{B}|U_\lambda)$-가군의 동형사상
$\mathcal{L}'|U_\lambda\simeq\mathcal{B}|U_\lambda$이다. 가정에 따라
$U_\lambda$들은 $X$의 덮개를 이룬다~; 임의의 두 지표 $\lambda$, $\mu$에
대하여
$\omega_{\lambda\mu}=(\eta_\lambda|U_\lambda\cap U_\mu)
\circ(\eta_\mu|U_\lambda\cap U_\mu)^{-1}$로 놓자. 이는
$\mathcal{B}|U_\lambda\cap U_\mu$의 자동사상으로서
$U_\lambda\cap U_\mu$ 위의 $\mathcal{B}^*$의 절단과 표준적으로
동일시되며, $(\omega_{\lambda\mu})$는 $\mathcal{B}^*$에 값을 갖는 덮개
$\mathfrak{U}=(U_\lambda)$의 $1$-공순환이다
(\hyperref[0.5.4.7-ko]{0, 5.4.7}).
$N_{\mathcal{B}/\mathcal{A}}:\mathcal{B}^*\to\mathcal{A}^*$가
준동형이라는 사실로부터
$(N_{\mathcal{B}/\mathcal{A}}\omega_{\lambda\mu})$는
$\mathcal{A}^*$에 값을 갖는 $\mathfrak{U}$의 $1$-공순환이고, 따라서
이에 (유일한 동형사상까지) 가역 $\mathcal{A}$-가군 $\mathcal{L}$이
대응한다 (\hyperref[0.5.4.7-ko]{0, 5.4.7}). 이를
$N_{\mathcal{B}/\mathcal{A}}(\mathcal{L}')$로 나타내고 \emph{가역
$\mathcal{B}$-가군 $\mathcal{L}'$의 노름}이라 부르겠다.

$\mathfrak{M}$을, 이에 대응하는 $U_\lambda$들이 여전히 $X$의 덮개를
이루도록 하는 $\mathfrak{L}$의 부분집합이라 하고, 이 덮개를
$\mathfrak{B}$라 하자~; 공순환 $(\omega_{\lambda\mu})$를
$\mathfrak{B}$에 제한하면 $\mathfrak{B}$의 $1$-공순환
$(N_{\mathcal{B}/\mathcal{A}}\omega_{\lambda\mu})$가 정의되며, 이는
$\mathfrak{U}$의 $1$-공순환
$(N_{\mathcal{B}/\mathcal{A}}\omega_{\lambda\mu})$의 제한으로서
$\mathcal{A}^*$에 값을 갖는다~; 이 $\mathfrak{B}$의 $1$-공순환으로
정의되는 가역 $\mathcal{A}$-가군과
$N_{\mathcal{B}/\mathcal{A}}(\mathcal{L}')$ 사이에 표준 동형사상이
있음은 명백하고, 따라서 임의의 $\mathfrak{U}$의 부분덮개를 사용하여 이
가역 $\mathcal{A}$-가군을 정의할 수 있다. 이 가능성으로부터,
$\mathcal{L}'_1$, $\mathcal{L}'_2$가 두 가역 $\mathcal{B}$-가군이면
(\hyperref[II.6.5.1.1-ko]{6.5.1.1})과
(\hyperref[II.6.5.1.2-ko]{6.5.1.2})에 따라 곧바로
\[
\label{II.6.5.2.1-ko}
N(\mathcal{L}'_1\otimes_{\mathcal{B}}\mathcal{L}'_2)
=N(\mathcal{L}'_1)\otimes_{\mathcal{A}}N(\mathcal{L}'_2)
\tag{6.5.2.1}
\]
이고
\[
\label{II.6.5.2.2-ko}
N_{\mathcal{B}/\mathcal{A}}(\mathcal{B})=\mathcal{A}
\tag{6.5.2.2}
\]
이며
\[
\label{II.6.5.2.3-ko}
N((\mathcal{L}')^{-1})=(N(\mathcal{L}'))^{-1}
\tag{6.5.2.3}
\]
임을 표준 동형사상까지 얻는다. 더 나아가
(\hyperref[II.6.5.1.3-ko]{6.5.1.3})으로부터, $\mathcal{L}$이 가역
$\mathcal{A}$-가군이고 경우 (I)에서 $\mathcal{B}$가
$\mathcal{A}$ 위에서 상수 계수 $n$이면 (각각 경우 (II)에서
$\mathcal{B}\otimes_{\mathcal{A}}\mathcal{R}(X)$가
$\mathcal{R}(X)$ 위에서 상수 계수 $n$이면), 표준 동형사상까지
\[
\label{II.6.5.2.4-ko}
N_{\mathcal{B}/\mathcal{A}}(\mathcal{L}\otimes_{\mathcal{A}}\mathcal{B})
=\mathcal{L}^{\otimes n}.
\tag{6.5.2.4}
\]
이다.
\end{env}

\begin{env}[6.5.3]
\label{II.6.5.3-ko}
이제 $N_{\mathcal{B}/\mathcal{A}}(\mathcal{L}')$가 가역
$\mathcal{B}$-가군들의 범주에서 \emph{공변 함자}임을 보이자. 실제로
$h':\mathcal{L}'_1\to\mathcal{L}'_2$를 가역 $\mathcal{B}$-가군들의
준동형이라 하고, $\mathfrak{B}=(U_\lambda)$를 모든 $\lambda$에 대하여
$(\mathcal{B}|U_\lambda)$-동형사상
$\eta^{(1)}_\lambda:\mathcal{L}'_1|U_\lambda\simeq\mathcal{B}|U_\lambda$와
$\eta^{(2)}_\lambda:\mathcal{L}'_2|U_\lambda\simeq\mathcal{B}|U_\lambda$가
존재하는 $X$의 열린 덮개라 하자~; 그러면 각 $\lambda$에 대하여
$\mathcal{B}|U_\lambda$의 자기준동형 $h'_\lambda$가 존재하여
$h'_\lambda\circ\eta^{(1)}_\lambda
=\eta^{(2)}_\lambda\circ(h'|U_\lambda)$이고, $h'_\lambda$를
$U_\lambda$ 위의 $\mathcal{B}$의 한 단면과 분명히 동일시할 수 있다
(\hyperref[0.5.1.1-ko]{0, 5.1.1}). 따라서 모든 지표쌍 $(\lambda,\mu)$에
대하여 $(\eta^{(2)}_\lambda)^{-1}\circ h'_\lambda\circ
\eta^{(1)}_\lambda$와 $(\eta^{(2)}_\mu)^{-1}\circ h'_\mu\circ
\eta^{(1)}_\mu$의 $U_\lambda\cap U_\mu$에 대한 제한은 일치한다. 이로부터
$\mathcal{L}'_1$과 $\mathcal{L}'_2$에 각각 대응하는 $\mathcal{B}^*$ 값을
갖는 $1$-코사이클 $(\omega^{(1)}_{\lambda\mu})$와
$(\omega^{(2)}_{\lambda\mu})$에 대하여 다음 관계를 얻는다.
\[
\omega^{(2)}_{\lambda\mu}h'_\mu
=h'_\lambda\omega^{(1)}_{\lambda\mu}.
\]

\oldpage[II]{128}

$h_\lambda=N(h'_\lambda)$로 놓으면, 따라서 이와 유사한 관계
\[
N(\omega^{(2)}_{\lambda\mu})h_\mu
=h_\lambda N(\omega^{(1)}_{\lambda\mu})
\]
가 성립하며, 그러므로 $h_\lambda$들은 준동형
$N(\mathcal{L}'_1)\to N(\mathcal{L}'_2)$를 정의한다. 이를
$N_{\mathcal{B}/\mathcal{A}}(h')$ 또는 $N(h')$로 표기하겠다. 가정 (I)에서
$h'$가 \emph{동형사상}이기 위한 필요충분조건은 (이는 국소적인 문제이므로)
$N_{\mathcal{B}/\mathcal{A}}(h')$가 동형사상인 것이다. 가정 (II)에서는
이 조건이 여전히 필요하다~; 가정 (II \emph{bis})가 성립하고
$\mathcal{B}\to\mathcal{B}\otimes_{\mathcal{A}}\mathcal{R}(X)$가 단사이면
그 조건은 충분하기도 하다.

특히 $\mathcal{L}'_1=\mathcal{B}$로 취하자~; 그러면 준동형
$\mathcal{B}\to\mathcal{L}'$은
(\hyperref[0.5.1.1-ko]{0, 5.1.1})에 따라 $X$ 위의 $\mathcal{L}'$의
단면과 동일시되므로, 표준 사상
\[
N_{\mathcal{B}/\mathcal{A}}:\Gamma(X,\mathcal{L}')
\to\Gamma(X,N_{\mathcal{B}/\mathcal{A}}(\mathcal{L}')).
\]
을 얻는다. 또한 (\hyperref[II.6.5.1.1-ko]{6.5.1.1})로부터
$f'_1\in\Gamma(X,\mathcal{L}'_1)$,
$f'_2\in\Gamma(X,\mathcal{L}'_2)$이면 다음이 성립한다.
\[
\label{II.6.5.3.1-ko}
N(f'_1\otimes f'_2)=N(f'_1)\otimes N(f'_2).
\tag{6.5.3.1}
\]
모든 가역 $\mathcal{A}$-가군 $\mathcal{L}$과 모든 단면
$f\in\Gamma(X,\mathcal{L})$에 대하여,
(\hyperref[II.6.5.2.4-ko]{6.5.2.4})를 고려하면
\[
\label{II.6.5.3.2-ko}
N_{\mathcal{B}/\mathcal{A}}(f\otimes1_{\mathcal{B}})=f^{\otimes n}
\tag{6.5.3.2}
\]
가 성립한다. 여기서 가정 (I)에서는 $\mathcal{B}$의 계수가 일정하게
$n$이고 (각각 가정 (II)에서는
$\mathcal{B}\otimes_{\mathcal{A}}\mathcal{R}(X)$의 계수가 일정하게
$n$이다). 마지막으로, $X$ 위의 $\mathcal{L}'$의 단면 $f'$에 대응하는
준동형 $\mathcal{B}\to\mathcal{L}'$이 동형사상이기 위한 필요충분조건은
모든 $x\in X$에 대하여 $f'_x$가 $\mathcal{L}'_x$의 기저인 것이다~;
따라서 가정 (I)에서 이 조건은 모든 $x$에 대하여 $(N(f'))_x$가
$(N(\mathcal{L}'))_x$의 기저라는 것과 동치이다~; 가정 (II)에서 이 조건은
여전히 필요하고, $\mathcal{B}$가 (II \emph{bis})를 만족하며
$\mathcal{B}\to\mathcal{B}\otimes_{\mathcal{A}}\mathcal{R}(X)$가 단사이면
충분하다.
\end{env}

\begin{env}[6.5.4]
\label{II.6.5.4-ko}
$(X,\mathcal{A})$, $(X',\mathcal{A}')$를 두 환 달린 공간이라 하고,
$f:X'\to X$를 사상, $\mathcal{B}$를 $\mathcal{A}$-대수,
$\mathcal{B}'=f^*(\mathcal{B})$를 역상 $\mathcal{A}'$-대수라 하자. 다음
가정 가운데 하나가 성립한다고 가정한다~:
\begin{enumerate}
\item[$1^\circ$] $\mathcal{B}$는
(\hyperref[II.6.5.1-ko]{6.5.1})의 가정 (I)을 만족한다.
\item[$2^\circ$] $(X,\mathcal{A})$와 $\mathcal{B}$는
(\hyperref[II.6.5.1-ko]{6.5.1})의 가정 (II)를 만족하고,
$(X',\mathcal{A}')$는 국소 뇌터 축소 준스킴이다. 또한 각각 $X$와
$X'$의 기약 성분들을 바탕공간으로 갖는 축소 닫힌 부분준스킴들을
$X_\alpha$와 $X'_\beta$로 나타낼 때, 각 $X'_\beta$에 대한 $f$의
제한은 $X'_\beta$에서 어떤 $X_\alpha$로 가는 \emph{우세} 사상이다.
\end{enumerate}
이 조건들 아래에서 $\mathcal{B}'$는 각각
(\hyperref[II.6.5.1-ko]{6.5.1})의 가정 (I) 또는 가정 (II)를 만족한다~;
첫 번째 주장은 자명하다. 두 번째 주장을 보이려면 모든 $x'\in X'$에
대하여, $x'\in X'_\beta$인 $\beta$들에 대응하는
$\mathcal{B}'\otimes_{\mathcal{O}_{X'}}\mathcal{R}(X'_\beta)$들의
계수가 \emph{모두 같음}을 보이면 충분하다. 그런데 $X'_\beta$에 대한
$f$의 제한이 어떤 $X_\alpha$로 가는 우세 사상이면,
$\mathcal{B}'\otimes_{\mathcal{O}_{X'}}\mathcal{R}(X'_\beta)$의 계수는
$\mathcal{B}\otimes_{\mathcal{O}_X}\mathcal{R}(X_\alpha)$의 계수와 같다
((\hyperref[II.6.4.9-ko]{6.4.9})에서처럼 아핀인 경우로 환원하면 곧바로
알 수 있다). 따라서 가정 (II)와
(\hyperref[II.6.5.1-ko]{6.5.1})에 의해 주장이 성립한다.

\oldpage[II]{129}

이제 (\hyperref[II.6.4.8-ko]{6.4.8}),
(\hyperref[II.6.4.9-ko]{6.4.9}),
(\hyperref[II.6.4.10-ko]{6.4.10})으로부터 다음을 얻는다. $s$가 열린집합
$U\subset X$ 위의 $\mathcal{B}$의 단면이고, $s'$이 $f^{-1}(U)$ 위의
$\mathcal{B}'$의 대응하는 단면이면,
$N_{\mathcal{B}'/\mathcal{A}'}(s')$은 $U$ 위의 $\mathcal{A}$의 단면
$N_{\mathcal{B}/\mathcal{A}}(s)$에 대응하는 $f^{-1}(U)$ 위의
$\mathcal{A}'$의 단면이다.

$\mathcal{M}$이 가역 $\mathcal{B}$-가군이면, 앞의 결과로부터
$\mathcal{M}'=f^*(\mathcal{M})$ (이는 가역 $\mathcal{B}'$-가군이다)에
대하여 표준 동형사상까지
$N_{\mathcal{B}'/\mathcal{A}'}(\mathcal{M}')
=f^*(N_{\mathcal{B}/\mathcal{A}}(\mathcal{M}))$임을 얻는다.
\end{env}

\begin{env}[6.5.5]
\label{II.6.5.5-ko}
이제부터 $(X,\mathcal{A})$가 \emph{준스킴}이라고 가정하자. 잘 알려진
대로, 준연접 $\mathcal{A}$-대수 $\mathcal{B}$이면서 \emph{유한 생성}
$\mathcal{A}$-가군인 $\mathcal{B}$를 주는 것은
$g_*(\mathcal{O}_{X'})=\mathcal{B}$를 만족하는 \emph{유한} 사상
$g:X'\to X$를 주는 것과 동치이고, 이 사상은 $X$-동형까지 정의된다
(\hyperref[II.6.1.2-ko]{6.1.2}와
\hyperref[II.1.3.1-ko]{1.3.1}). 더욱이 준연접
$\mathcal{O}_{X'}$-가군 $\mathcal{F}'$을 주는 것은
$g_*(\mathcal{F}')=\mathcal{F}$를 만족하는 준연접 $\mathcal{B}$-가군
$\mathcal{F}$를 주는 것과 동치이고
(\hyperref[II.1.4.3-ko]{1.4.3}), $\mathcal{F}'$이 가역이기 위한
필요충분조건은 $\mathcal{F}$가 가역인 것이다
(\hyperref[II.6.1.12-ko]{6.1.12}). 앞의 결과들을 유한 사상 $g$의
용어로 옮기려면, 따라서 $g_*(\mathcal{O}_{X'})$가 (유한 계수의) \emph{국소
자유} $\mathcal{O}_X$-가군이라고 가정하거나,
$(X,\mathcal{O}_X)$와 $g_*(\mathcal{O}_{X'})$가 가정 (II)를 만족한다고
가정해야 한다. 임의의 가역 $\mathcal{O}_{X'}$-가군 $\mathcal{L}'$에
대하여 이제
\[
\label{II.6.5.5.1-ko}
N_{X'/X}(\mathcal{L}')
=N_{g_*(\mathcal{O}_{X'})/\mathcal{O}_X}(g_*(\mathcal{L}'))
\tag{6.5.5.1}
\]
로 놓고, 이를 $\mathcal{L}'$의 ($g$에 대한) \emph{노름}이라고 한다.
마찬가지로 $h':\mathcal{L}'_1\to\mathcal{L}'_2$가 가역
$\mathcal{O}_{X'}$-가군들의 준동형이면
\[
\label{II.6.5.5.2-ko}
N_{X'/X}(h')
=N_{g_*(\mathcal{O}_{X'})/\mathcal{O}_X}(g_*(h')):
N_{X'/X}(\mathcal{L}'_1)\to N_{X'/X}(\mathcal{L}'_2).
\tag{6.5.5.2}
\]
로 놓는다. 특히 $\mathcal{L}'_1=\mathcal{O}_{X'}$로 놓으면 표준 사상
\[
\label{II.6.5.5.3-ko}
N_{X'/X}:\Gamma(X',\mathcal{L}')
\to\Gamma(X,N_{X'/X}(\mathcal{L}')).
\tag{6.5.5.3}
\]
을 얻는다. 이러한 번역의 대부분은 독자에게 맡기고, 다음 것들을
명시하는 데 그치겠다~:
\end{env}

\begin{proposition}[6.5.6]
\label{II.6.5.6-ko}
$g:X'\to X$를 유한 사상이라 하자. $g_*(\mathcal{O}_{X'})$가
국소 자유 $\mathcal{O}_X$-가군이거나, 또는 $(X,\mathcal{O}_X)$와
$g_*(\mathcal{O}_{X'})$가 (II \emph{bis})를 만족한다고 가정하자
(특히 $X$가 국소 뇌터 정규 준스킴이면 후자의 경우가 성립한다).
가역 $\mathcal{O}_{X'}$-가군의 준동형
$h':\mathcal{L}'_1\to\mathcal{L}'_2$가 동형사상일 필요충분조건은,
첫 번째 가정에서는 $N_{X'/X}(h')$가 동형사상인 것이다~; 두 번째
가정에서는 이 조건이 여전히 필요하며, 준동형
$g_*(\mathcal{O}_{X'})\to
g_*(\mathcal{O}_{X'})\otimes_{\mathcal{O}_X}\mathcal{R}(X)$가
단사이면 충분하기도 하다.
\end{proposition}

여기서는 $g_*(h')$가 동형사상일 필요충분조건이 $h'$가 동형사상인
것이라는 사실을 사용했음에 유의하자
(\hyperref[II.1.4.2-ko]{1.4.2}).

\begin{corollary}[6.5.7]
\label{II.6.5.7-ko}
$g:X'\to X$를 유한 사상이라 하자. $g_*(\mathcal{O}_{X'})$가
국소 자유 $\mathcal{O}_X$-가군이거나, 또는 $(X,\mathcal{O}_X)$와
$g_*(\mathcal{O}_{X'})$가 (II \emph{bis})를 만족하고 준동형
$g_*(\mathcal{O}_{X'})\to
g_*(\mathcal{O}_{X'})\otimes_{\mathcal{O}_X}\mathcal{R}(X)$가
단사라고 가정하자. $\mathcal{L}'$을 가역 $\mathcal{O}_{X'}$-가군,
$f'$를 $X'$ 위의 $\mathcal{L}'$의 절단이라 하고,
$f=N_{X'/X}(f')$를 이에 대응하는 $X$ 위의
$\mathcal{L}=N_{X'/X}(\mathcal{L}')$의 절단이라 하자
(\hyperref[II.6.5.5.3-ko]{6.5.5.3}).

\oldpage[II]{130}

그러면 $g(X'-X'_{f'})=X-X_f$이고, $X_f$는
$g^{-1}(U)\subset X'_{f'}$를 만족하는 $X$의 가장 큰 열린집합 $U$이다.
\end{corollary}

실제로 $g(X'-X'_{f'})$는 $X$에서 닫혀 있다
(\hyperref[II.6.1.10-ko]{6.1.10})~; 따라서 마지막 주장을 증명하면
충분하다. 그런데 관계 $U\subset X_f$는 $f|U$로 정의되는 준동형
$\mathcal{O}_X|U\to\mathcal{L}|U$가 동형사상이라는 것과 동치이다.
(\hyperref[II.6.5.6-ko]{6.5.6})에 따라 이는 $f'|g^{-1}(U)$로 정의되는
준동형
$\mathcal{O}_{X'}|g^{-1}(U)\to\mathcal{L}'|g^{-1}(U)$가
동형사상이라는 것, 즉 관계 $g^{-1}(U)\subset X'_{f'}$와 동치이다.

\begin{proposition}[6.5.8]
\label{II.6.5.8-ko}
$g:X'\to X$를 유한 사상, $f:Y\to X$를 사상이라 하자~;
$Y'=X'_{(Y)}$, $g'=g_{(Y)}$, $f'=f_{(X')}$라 놓으면 다음 가환 도식을
얻는다.
\[
\xymatrix{
X'\ar[d]_g & Y'\ar[l]_{f'}\ar[d]^{g'}\\
X & Y\ar[l]^f
}
\]
$g_*(\mathcal{O}_{X'})$가 국소 자유이거나, 또는
$(X,\mathcal{O}_X)$와 $g_*(\mathcal{O}_{X'})$가 (II)를 만족하고
$Y$가 국소 뇌터 축소 준스킴이며 $Y$의 각 기약 성분에 대한 $f$의
제한이 $X$의 한 기약 성분으로 가는 우세 사상이라고 가정하자. 그러면
모든 가역 $\mathcal{O}_{X'}$-가군 $\mathcal{L}'$에 대하여
\[
N_{Y'/Y}({f'}^*(\mathcal{L}'))=f^*(N_{X'/X}(\mathcal{L}'))
\]
가 표준 동형사상까지 성립한다.
\end{proposition}

(\hyperref[II.1.5.2-ko]{1.5.2})에 따라
$f^*(g_*(\mathcal{L}'))=g'_*({f'}^*(\mathcal{L}'))$이고, 특히
$g'_*(\mathcal{O}_{Y'})=f^*(g_*(\mathcal{O}_{X'}))$임에 유의하자~;
$g_*(\mathcal{O}_{X'})$가 국소 자유이면
$g'_*(\mathcal{O}_{Y'})$도 그러하다. 이제 결론은 정의와
(\hyperref[II.6.5.4-ko]{6.5.4})에서 따른다.

\begin{remark}[6.5.9]
\label{II.6.5.9-ko}
위에서 전개한 노름의 개념은 뒤에서 일반화하고, 약수의 직상이라는
개념과 관련지을 것이다.
\end{remark}

\subsection{응용: 풍부성 판정법.}
\label{subsection:II.6.6-ko}

\begin{proposition}[6.6.1]
\label{II.6.6.1-ko}
$Y$를 준스킴, $f:X\to Y$를 준콤팩트 사상,
$g:X'\to X$를 유한 전사 사상이라 하자.
$g_*(\mathcal{O}_{X'})$가 국소 자유 $\mathcal{O}_X$-가군이거나, 또는
$(X,\mathcal{O}_X)$와 $g_*(\mathcal{O}_{X'})$가 (II
\emph{bis})를 만족한다고 가정하자. 그러면 $f\circ g$에 관해
풍부한 모든 가역 $\mathcal{O}_{X'}$-가군 $\mathcal{L}'$에 대하여
$N_{X'/X}(\mathcal{L}')=\mathcal{L}$은 $f$에 관해 풍부하다.
\end{proposition}

$Y$가 아핀이라고 가정할 수 있으며
(\hyperref[II.4.6.4-ko]{4.6.4}), 이 경우
(\hyperref[II.4.6.6-ko]{4.6.6})에 따라 이 명제는 다음 따름정리와
동치이다.

\begin{corollary}[6.6.2]
\label{II.6.6.2-ko}
$X$를 준콤팩트 준스킴, $g:X'\to X$를 유한 전사 사상이라 하자.
$g_*(\mathcal{O}_{X'})$가 국소 자유 $\mathcal{O}_X$-가군이거나,
$(X,\mathcal{O}_X)$와 $g_*(\mathcal{O}_{X'})$가 (II \emph{bis})를
만족한다고 하자. 그러면 모든 풍부한 $\mathcal{O}_{X'}$-가군
$\mathcal{L}'$에 대하여
$\mathcal{L}=N_{X'/X}(\mathcal{L}')$은 풍부하다.
\end{corollary}

둘째 가정에서는 표준 준동형
$g_*(\mathcal{O}_{X'})\to
g_*(\mathcal{O}_{X'})\otimes_{\mathcal{O}_X}\mathcal{R}(X)$가
\emph{단사}라고 더 가정할 수 있다. 실제로 그렇지 않으면
$\mathcal{T}$를

\oldpage[II]{131}

이 준동형의 핵이라 하자. 이는
$g_*(\mathcal{O}_{X'})=\mathcal{B}$의 연접 아이디얼이고
(\hyperref[I.6.1.1-ko]{I, 6.1.1}),
$X''=\operatorname{Spec}(\mathcal{B}/\mathcal{T})$로 놓으면 가환 도식
\[
\xymatrix{
X''\ar[rr]^h\ar[dr]_{g'} && X'\ar[dl]^g\\
& X
}
\]
을 얻는다. 여기서 $h$는 닫힌 몰입 사상이다
(\hyperref[II.1.4.10-ko]{1.4.10}). 또한 $\mathcal{T}$의 지지집합은
$X$에서 희박한 닫힌집합임을 안다
(\hyperref[0.5.2.2-ko]{0, 5.2.2})
(\hyperref[I.7.4.6-ko]{I, 7.4.6}). 따라서 $X$의 한 기약 성분의 일반점
$x$마다 $x$의 아핀 열린 근방 $U$가 존재하여
$\mathcal{B}|U=(\mathcal{B}/\mathcal{T})|U$이다. 가정에 따라 $g$가
전사이므로 $x\in g'(X'')$임이 따른다~; 따라서 $g'$은 우세 사상이고,
유한 사상이므로 \emph{전사}이다
(\hyperref[II.6.1.10-ko]{6.1.10}). 정의에 따라
\[
g'_*(\mathcal{O}_{X''})\otimes_{\mathcal{O}_X}\mathcal{R}(X)
=(\mathcal{B}/\mathcal{T})\otimes_{\mathcal{O}_X}\mathcal{R}(X)
=g_*(\mathcal{O}_{X'})\otimes_{\mathcal{O}_X}\mathcal{R}(X),
\]
이므로 $(X,\mathcal{O}_X)$와 $g'_*(\mathcal{O}_{X''})$는
(II \emph{bis})를 만족하고, 더 나아가
\[
g'_*(\mathcal{O}_{X''})\longrightarrow
g'_*(\mathcal{O}_{X''})\otimes_{\mathcal{O}_X}\mathcal{R}(X)
\]
는 단사이다. 끝으로 $h^*(\mathcal{L}')=\mathcal{L}''$은 풍부한
$\mathcal{O}_{X''}$-가군이고
(\hyperref[II.4.6.13-ko]{4.6.13, (i \emph{bis})}),
$N_{X''/X}(\mathcal{L}'')=N_{X'/X}(\mathcal{L}')$이다. 실제로 이 두
가역 $\mathcal{O}_X$-가군을 정의할 때, $g_*(\mathcal{L}')$과
$g'_*(\mathcal{L}'')$을 $U_\lambda$에 제한한 것이 각각
$\mathcal{B}|U_\lambda$와
$(\mathcal{B}/\mathcal{T})|U_\lambda$에 동형이 되도록 하는 $X$의 같은
아핀 열린 덮개 $(U_\lambda)$를 사용할 수 있다. 모든 동형사상
\[
\eta_\lambda:g_*(\mathcal{L}')|U_\lambda\longrightarrow
\mathcal{B}|U_\lambda
\]
에는 표준적으로 동형사상
\[
\eta'_\lambda:g'_*(\mathcal{L}'')|U_\lambda\longrightarrow
(\mathcal{B}/\mathcal{T})|U_\lambda,
\]
이 대응한다. 따라서 동형사상계 $(\eta_\lambda)$와
$(\eta'_\lambda)$에 대응하는 $1$-공순환들을 각각
$(\omega_{\lambda\mu})$와 $(\omega'_{\lambda\mu})$라 하면
(\hyperref[II.6.5.2-ko]{6.5.2}), $\omega'_{\lambda\mu}$는
$\omega_{\lambda\mu}\in
\Gamma(U_\lambda\cap U_\mu,\mathcal{B})$의
$\Gamma(U_\lambda\cap U_\mu,\mathcal{B}/\mathcal{T})$ 안에서의 표준
상이다. $\mathcal{T}$의 정의에 따라
\[
N_{\mathcal{B}/\mathcal{A}}\omega_{\lambda\mu}
=N_{(\mathcal{B}/\mathcal{T})/\mathcal{A}}\omega'_{\lambda\mu}
\]
(여기서 $\mathcal{A}=\mathcal{O}_X$)임이 따르고, 이로부터 말한 등식을
얻는다.

그러므로 (II \emph{bis})의 가정 아래에서는 준동형
$g_*(\mathcal{O}_{X'})\to
g_*(\mathcal{O}_{X'})\otimes_{\mathcal{O}_X}\mathcal{R}(X)$가
단사라고 가정하자. $f$가 $X$ 위의 $\mathcal{L}^{\otimes n}$의 단면들을
달릴 때($n>0$), $X_f$들이 $X$의 위상의 기저를 이룸을 증명하면
충분하다 (\hyperref[II.4.5.2-ko]{4.5.2}). 이제 $x\in X$라 하고,
$U$를 $x$의 임의의 근방이라 하자~; $g^{-1}(x)$는 유한하고
(\hyperref[II.6.1.7-ko]{6.1.7}) $\mathcal{L}'$은 풍부하므로, 어떤 정수
$n>0$과 $X'$ 위의 ${\mathcal{L}'}^{\otimes n}$의 단면 $f'$이 존재하여
$X'_{f'}$은 $g^{-1}(U)$에 포함되는 $g^{-1}(x)$의 근방이다
(\hyperref[II.4.5.4-ko]{4.5.4}). 그런데
\[
\mathcal{L}^{\otimes n}=N_{X'/X}({\mathcal{L}'}^{\otimes n}),
\]
이므로 $f=N_{X'/X}(f')$로 택하면 충분하다~; 실제로 이때
$X-X_f=g(X'-X'_{f'})$이고
(\hyperref[II.6.5.7-ko]{6.5.7}), 따라서 $x\in X_f\subset U$이다.

\begin{corollary}[6.6.3]
\label{II.6.6.3-ko}
(\hyperref[II.6.6.1-ko]{6.6.1})의 가정 아래, 가역
$\mathcal{O}_X$-가군 $\mathcal{L}$이 $f$에 관해 풍부하기 위한
필요충분조건은 $\mathcal{L}'=g^*(\mathcal{L})$이 $f\circ g$에 관해
풍부한 것이다.
\end{corollary}

이 조건은 필요하다. 실제로 $g$가 아핀이기 때문이다
(\hyperref[II.5.1.12-ko]{5.1.12}). 조건이 충분함을 보이기 위해서는
$Y$가 아핀이라고 가정할 수 있다 (\hyperref[II.4.6.4-ko]{4.6.4}).
따라서 $X$와 $X'$은 준콤팩트이고 $\mathcal{L}'$은 풍부하며
(\hyperref[II.4.6.6-ko]{4.6.6}), $\mathcal{L}$이 풍부함을 보이면 된다.
한편, 첫째 (각각 둘째) 가정에서
$g_*(\mathcal{O}_{X'})$ (각각
$g_*(\mathcal{O}_{X'})\otimes_{\mathcal{O}_X}\mathcal{R}(X)$)가 그
근방에서 주어진 계수 $n$을 갖는 점

\oldpage[II]{132}

$x\in X$들의 집합은 $X$에서 열리고 닫혀 있다. 따라서 $X$는 이러한
열린집합들 가운데 유한 개의 합인 준스킴이고, 그러므로 그 가운데 하나와
같다고 가정할 수 있다 (\hyperref[II.4.6.17-ko]{4.6.17}). 그런데 이때
$N_{X'/X}(\mathcal{L}')=\mathcal{L}^{\otimes n}$이므로,
(\hyperref[II.6.6.2-ko]{6.6.2})에 따라
$\mathcal{L}^{\otimes n}$은 풍부하고, 따라서 $\mathcal{L}$도 풍부하다
(\hyperref[II.4.5.6-ko]{4.5.6}).

\begin{corollary}[6.6.4]
\label{II.6.6.4-ko}
(\hyperref[II.6.6.1-ko]{6.6.1})의 가정들이 성립하고, 더 나아가
$f:X\to Y$가 유한형이라고 가정하자. 그러면 $f$가 준사영적이기 위한
필요충분조건은 $f\circ g$가 준사영적인 것이다. 또한 $Y$가 준콤팩트
스킴이거나 바탕공간이 뇌터인 준스킴이라고 가정하면, $f$가 사영적이기
위한 필요충분조건은 $f\circ g$가 사영적인 것이다.
\end{corollary}

가정으로부터 $f\circ g$가 유한형임을 알 수 있다. 준사영 사상의 정의
(\hyperref[II.5.3.1-ko]{5.3.1})를 고려하면, 첫째 명제는
(\hyperref[II.6.6.1-ko]{6.6.1})과
(\hyperref[II.6.6.3-ko]{6.6.3})에서 나온다. 이 결과와
(\hyperref[II.5.5.3-ko]{5.5.3, (ii)})를 고려하면, $f$가
준사영적이라고 가정할 때 $f$가 고유이기 위한 필요충분조건이
$f\circ g$가 고유인 것임을 보이는 일만 남는다. 그런데 이때 $f$는
분리되어 있고 (\hyperref[II.5.3.1-ko]{5.3.1}) 유한형이다~; $g$가
전사이므로, 이 명제는 (\hyperref[II.5.4.2-ko]{5.4.2, (ii)})와
(\hyperref[II.5.4.3-ko]{5.4.3, (ii)})에서 나온다.

특히 다음을 얻는다~:

\begin{corollary}[6.6.5]
\label{II.6.6.5-ko}
$X$를 체 $K$ 위의 유한형 준스킴, $K'$을 $K$의 유한 차수 확대라 하자.
$X$가 $K$ 위에서 사영적 (각각 준사영적)이기 위한 필요충분조건은
$X'=X\otimes_K K'$이 $K'$ 위에서 사영적 (각각 준사영적)인 것이다.
\end{corollary}

실제로 이 조건은 필요하다
(\hyperref[II.5.3.4-ko]{5.3.4, (iii)} 및
\hyperref[II.5.5.5-ko]{5.5.5, (iii)}). 역으로 이 조건이 성립한다고
가정하고, $g:X'\to X$를 표준 사영이라 하자. $g$가 유한 사상
(\hyperref[II.6.1.5-ko]{6.1.5, (iii)})이자 전사 사상
(\hyperref[I.3.5.2-ko]{I, 3.5.2, (ii)})임은 명백하다. 더 나아가
$g_*(\mathcal{O}_{X'})$은 $\mathcal{O}_X\otimes_K K'$과 동형이므로
(\hyperref[II.1.5.2-ko]{1.5.2}) 국소 자유 $\mathcal{O}_X$-가군이다.
그러므로 가정과 (\hyperref[II.6.1.11-ko]{6.1.11}) 및
(\hyperref[II.5.5.5-ko]{5.5.5, (ii)})에 따라 $X'$은
\emph{$K$ 위에서} 사영적 (각각 준사영적)이다~; 이제
(\hyperref[II.6.6.4-ko]{6.6.4})로부터 $X$가 $K$ 위에서 사영적
(각각 준사영적)임을 얻는다.

제V장에서 우리는 명제 (\hyperref[II.6.6.5-ko]{6.6.5})가 $K'$이
$K$의 \emph{임의의} 확대일 때에도 성립함을 보일 것이다.

이 항의 나머지는 (\hyperref[II.6.6.1-ko]{6.6.1})을 상당히 기술적으로
정밀화한 판정법 (\hyperref[II.6.6.11-ko]{6.6.11})의 증명에 할애한다~;
처음 읽을 때에는 이 부분을 생략해도 된다.

\begin{lemma}[6.6.6]
\label{II.6.6.6-ko}
$X$를 뇌터 축소 준스킴, $\mathcal{E}$를
$\mathcal{E}\otimes_{\mathcal{O}_X}\mathcal{R}(X)$가 상수 계수 $n$을
갖는 국소 자유 $\mathcal{R}(X)$-가군이 되도록 하는 연접
$\mathcal{O}_X$-가군이라 하자. 그러면 다음 성질을 갖는 뇌터 축소 준스킴
$Z$와 유한 쌍유리 사상 $h:Z\to X$가 존재한다~: 집합층의 사상
$\sigma_i:\mathcal{H}om_{\mathcal{O}_X}(\mathcal{E},\mathcal{E})
\to\mathcal{R}(X)$ ($1\leq i\leq n$)
(\hyperref[II.6.4.10-ko]{6.4.10})은
$\mathcal{H}om_{\mathcal{O}_X}(\mathcal{E},\mathcal{E})$를 연접
$\mathcal{O}_X$-대수 $h_*(\mathcal{O}_Z)$ 안으로 보낸다.
\end{lemma}

실제로 $X$의 아핀 열린집합 $U$를 생각하고 그 환을 $A(U)=A$라 하자~;
$E=\Gamma(U,\mathcal{E})$라 하고, $u$가
$\operatorname{Hom}_A(E,E)$를 움직일 때 $\sigma_i(u)$들이 생성하는
$R(U)$의 부분대수를 $C_U$라 하자~; 이 $A$-대수가 \emph{유한 계수}임을
앞에서 보았다 (\hyperref[II.6.4.7.1-ko]{6.4.7.1}). 더 나아가 대수
$C_U$의 형성은 아핀 열린집합 $U$에서 그 안의 아핀 열린집합
$U'\subset U$로 가는 제한 연산과 가환함이 명백하다. 따라서
$\mathcal{R}(X)$의 \emph{유한} 부분 $\mathcal{O}_X$-대수
$\mathcal{C}$를 이렇게 정의하였으며, 이는 $X$의 모든 아핀 열린집합
$U$에 대하여 $\Gamma(U,\mathcal{C})=C_U$를 만족한다. 이제

\oldpage[II]{133}

$Z=\operatorname{Spec}(\mathcal{C})$로 잡고 $h$를 구조 사상으로 잡자.
그러면 $h$는 유한이다 (\hyperref[II.6.1.2-ko]{6.1.2})~;
$\mathcal{C}$가 축소이므로 $Z$는 뇌터 축소 준스킴이다
(\hyperref[II.1.3.8-ko]{1.3.8}). 끝으로 정의에 따라 $C_U$의 전체
분수환은 $R(U)$이고, $C_U$는 $R(U)$ 안에서 $A(U)$의 정수적 폐포에
포함되므로 $A(U)$의 극소 소 아이디얼들과 $C_U$의 극소 소 아이디얼들
사이에 일대일 대응이 있다 (
\cite[t.~I, p.~259]{I-13}). 이로부터 $h$가 쌍유리 사상임이 증명되고
증명이 끝난다.

\begin{corollary}[6.6.7]
\label{II.6.6.7-ko}
(\hyperref[II.6.6.6-ko]{6.6.6})의 가정 아래에서, $W$를 모든
$x\in W$에 대하여 $X$가 점 $x$에서 정규이거나
$\mathcal{E}_x$가 자유 $\mathcal{O}_x$-가군인 $X$의 열린집합이라 하자.
그러면 $h$를 $h^{-1}(W)$에 제한한 사상이 $h^{-1}(W)$에서 $W$ 위로
가는 동형사상이 되도록 $h$를 정의할 수 있다.
\end{corollary}

실제로 어느 가정이든, $U\subset W$가 아핀 열린집합이면
(\hyperref[II.6.6.6-ko]{6.6.6})의 표기 아래 모든 $x\in U$에 대하여
$(\sigma_i(u))_x\in A_x$이고
(\hyperref[II.6.4.3-ko]{6.4.3}), 따라서 $\sigma_i(u)\in A$임을
함의한다. 결론은 (\hyperref[II.6.6.6-ko]{6.6.6})에서 주어진 $h$의
정의로부터 따른다.

\begin{env}[6.6.8]
\label{II.6.6.8-ko}
$X$를 뇌터 축소 준스킴, $g:X'\to X$를 유한 전사 사상이라 하자. 그러면
$\mathcal{B}=g_*(\mathcal{O}_{X'})$는 \emph{연접}
$\mathcal{O}_X$-대수이다~; 더 나아가
$\mathcal{B}\otimes_{\mathcal{O}_X}\mathcal{R}(X)$가
$\mathcal{R}(X)$-가군으로서 \emph{상수 계수 $n$을 갖는 국소 자유}라고
가정하자. 이제 $\mathcal{E}=\mathcal{B}$로 두어 보조정리
(\hyperref[II.6.6.6-ko]{6.6.6})을 적용할 수 있고, 따라서
(\hyperref[II.6.6.6-ko]{6.6.6})의 표기로 곱셈 모노이드층의 준동형
$\sigma_n:\mathcal{H}om_{\mathcal{O}_X}(\mathcal{B},\mathcal{B})
\to h_*(\mathcal{O}_Z)$를 얻는다. 이 준동형과 표준 준동형
$\mathcal{B}\to
\mathcal{H}om_{\mathcal{O}_X}(\mathcal{B},\mathcal{B})$
(\hyperref[II.6.5.1-ko]{6.5.1})을 합성하면 다음과 같은 곱셈
모노이드층의 준동형을 얻는다~:
\[
\label{II.6.6.8.1-ko}
N':\mathcal{B}=g_*(\mathcal{O}_{X'})\longrightarrow
h_*(\mathcal{O}_Z)=\mathcal{C}.
\tag{6.6.8.1}
\]

이제 모든 가역 $\mathcal{O}_{X'}$-가군 $\mathcal{L}'$에 대하여
$g_*(\mathcal{L}')$은 가역 $\mathcal{B}$-가군이고
(\hyperref[II.6.1.12-ko]{6.1.12}),
(\hyperref[II.6.5.2-ko]{6.5.2})의 방법에 따라 $\mathcal{L}'$에 가역
$\mathcal{C}$-가군을 함자적으로 대응시킬 수 있다. 이를
$N'(g_*(\mathcal{L}'))$로 나타내겠다.
\end{env}

\begin{lemma}[6.6.9]
\label{II.6.6.9-ko}
$X$를 뇌터 축소 준스킴, $g:X'\to X$를 유한 전사 사상이라 하고,
$g_*(\mathcal{O}_{X'})\otimes_{\mathcal{O}_X}\mathcal{R}(X)$가 계수가
일정한 국소 자유 $\mathcal{R}(X)$-가군이라고 하자. 그러면 뇌터 축소
준스킴 $Z$와 다음 성질을 갖는 유한 쌍유리 사상 $h:Z\to X$가 존재한다~:
모든 풍부한 $\mathcal{O}_{X'}$-가군 $\mathcal{L}'$에 대하여,
$h_*(\mathcal{M})=N'(g_*(\mathcal{L}'))$인 가역
$\mathcal{O}_Z$-가군 $\mathcal{M}$
((\hyperref[II.6.6.8-ko]{6.6.8})의 표기법)은 풍부하다.
\end{lemma}

먼저 준동형
$\mathcal{B}\to
\mathcal{B}\otimes_{\mathcal{O}_X}\mathcal{R}(X)$가 \emph{단사}라고
가정하자. $Z$와 $h$를 (\hyperref[II.6.6.6-ko]{6.6.6})에서와 같이
(단, $\mathcal{E}=g_*(\mathcal{O}_{X'})$로 하여) 정의하자.
$z\in Z$라 하자~; $z$를 포함하는 아핀 열린집합 $Z_t$를 갖도록 하는
정수 $m>0$과 $Z$ 위의 $\mathcal{M}^{\otimes m}$의 단면 $t$가 존재함을
보여야 한다 (\hyperref[II.4.5.2-ko]{4.5.2}). $x=h(z)$라 하고, $U$를
$x$를 포함하는 $X$의 아핀 열린집합이라 하자~; 그러면 $h^{-1}(U)$는
$z$의 아핀 열린 근방이고, $z\in Z_t\subset h^{-1}(U)$인 $t$를 찾으면
충분하다. 실제로 이때 $Z_t$는 반드시 아핀이다
(\hyperref[II.5.5.8-ko]{5.5.8}). 가정에 따라 정수 $n>0$과 $X'$ 위의
${\mathcal{L}'}^{\otimes n}$의 단면 $s'$가 존재하여
\[
\label{II.6.6.9.1-ko}
g^{-1}(x)\subset X'_{s'}\subset g^{-1}(U)
\tag{6.6.9.1}
\]
가 성립한다 (\hyperref[II.4.5.4-ko]{4.5.4}). 정의에 따라 $s'$은 또한
$X$ 위의 $g_*({\mathcal{L}'}^{\otimes n})$의 단면이고,
(\hyperref[II.6.5.2-ko]{6.5.2})에서와 같이 이에 대응하는 $X$ 위의
$N'(g_*({\mathcal{L}'}^{\otimes n}))$의 단면은
$s=N'(s')$이다.

\oldpage[II]{134}

$s$를 $Z$ 위의 $\mathcal{M}^{\otimes n}$의 단면으로 여긴 것을 $t$라
하면, 이 $t$가 조건을 만족함을 보이겠다. 실제로
\[
\label{II.6.6.9.2-ko}
V=X-g(X'-X'_{s'})
\tag{6.6.9.2}
\]
라 놓자. 이는 (\hyperref[II.6.6.9.1-ko]{6.6.9.1})과
(\hyperref[II.6.1.10-ko]{6.1.10})에 따라 $x$를 포함하고 $U$에 포함되는
$X$의 열린집합이다. 이제
\[
\label{II.6.6.9.3-ko}
h^{-1}(V)\subset Z_t\subset h^{-1}(U),
\tag{6.6.9.3}
\]
임을 보이면 증명이 끝난다. 이는 $s_y$가 가역인 $y\in X$들의 집합
$T$가 $V$를 포함하고 $U$에 포함된다고 말하는 것과 같다. 이를 위해
먼저 $V$에 포함되는 아핀 열린집합 $W$를 생각하자~; $g^{-1}(W)$는
$X'$의 아핀 열린집합이고, (\hyperref[II.6.6.9.2-ko]{6.6.9.2})에 따라
모든 $y'\in g^{-1}(W)$에 대하여 $s'_{y'}$은 가역이다~; $X$와
$\mathcal{B}$에 관한 가정에 따라 (\hyperref[II.6.4.7-ko]{6.4.7})의
결과들을 적용할 수 있으므로, $y=g(y')$이면 $s_y$가 가역임을 알 수
있다. 다시 말해 $V\subset T$이다. 다른 한편,
(\hyperref[II.6.4.7-ko]{6.4.7})로부터 역으로 $s_y$가 가역이면
$s'_{y'}$도 가역임이 또한 따른다. 그러면
(\hyperref[II.6.6.9.1-ko]{6.6.9.1})에 따라 $y'\in g^{-1}(U)$이고,
따라서 $y\in U$이다. 이 경우 $T\subset U$임을 얻는다.

이제 (\hyperref[II.6.6.2-ko]{6.6.2})에서와 같은 논법으로 일반적인
경우로 넘어간다. 즉 $X'$을
\[
g'_*(\mathcal{O}_{X''})\longrightarrow
g'_*(\mathcal{O}_{X''})\otimes_{\mathcal{O}_X}\mathcal{R}(X)
\]
가 단사인 $X''$으로 바꾸고, $\mathcal{L}'$을
$N'(g_*(\mathcal{L}'))=N'(g'_*(\mathcal{L}''))$인 풍부한
$\mathcal{O}_{X''}$-가군 $\mathcal{L}''$으로 바꾼다. 따라서 보조정리
(\hyperref[II.6.6.9-ko]{6.6.9})는 모든 경우에 증명되었다
($h$를 적절히 선택한다).

\begin{corollary}[6.6.10]
\label{II.6.6.10-ko}
(\hyperref[II.6.6.9-ko]{6.6.9})의 가정들이 성립한다고 하자~;
$g^*(\mathcal{L})$이 풍부한 모든 가역 $\mathcal{O}_X$-가군
$\mathcal{L}$에 대하여 $h^*(\mathcal{L})$은 풍부하다.
\end{corollary}

실제로 $\mathcal{L}'=g^*(\mathcal{L})$로 놓으면
$g_*(\mathcal{L}')=\mathcal{L}\otimes_{\mathcal{O}_X}\mathcal{B}$이다
(\hyperref[0.5.4.10-ko]{0, 5.4.10}). 따라서
\[
N'(g_*(\mathcal{L}'))=
(\mathcal{L}\otimes_{\mathcal{O}_X}\mathcal{C})^{\otimes n}
\]
이다 ((\hyperref[II.6.5.2.4-ko]{6.5.2.4})에서와 같은 논법을 썼다).
그러므로 $\mathcal{M}=(h^*(\mathcal{L}))^{\otimes n}$이고,
$\mathcal{M}$이 풍부하므로 $h^*(\mathcal{L})$도 풍부하다
(\hyperref[II.4.5.6-ko]{4.5.6}).

\begin{proposition}[6.6.11]
\label{II.6.6.11-ko}
$Y$를 아핀 스킴, $X$를 축소 뇌터 준스킴,
$f:X\to Y$를 준콤팩트 사상, $g:X'\to X$를 유한 전사 사상이라 하자.
$W$를 $X$의 열린 부분집합이라 하고, 모든 $x\in W$에 대하여 $X$가
$x$에서 정규이거나, 아니면 $x$의 열린 근방 $T\subset W$가 존재하여
$(g_*(\mathcal{O}_{X'}))|T$가 국소 자유
$(\mathcal{O}_X|T)$-가군이라고 하자. 그러면 축소 $Y$-준스킴 $Z$와
유한 쌍유리 $Y$-사상 $h:Z\to X$가 존재하여, $h^{-1}(W)$에 대한
$h$의 제한은 동형사상 $h^{-1}(W)\xrightarrow{\sim}W$이고 다음 성질을
갖는다~: $g^*(\mathcal{L})$이 $f\circ g$에 대해 풍부한 모든 가역
$\mathcal{O}_X$-가군 $\mathcal{L}$에 대하여 $h^*(\mathcal{L})$은
$f\circ h$에 대해 풍부하다.
\end{proposition}

$Y$가 아핀이므로 $g^*(\mathcal{L})$은 풍부하고, 따라서 적절한 $h$를
택하여 $h^*(\mathcal{L})$이 풍부함을 보이면 된다
(\hyperref[II.4.6.6-ko]{4.6.6}). 이제 $g$를 다음을 만족하는 유한 전사
사상 $g':X''\to X$로 바꿀 수 있음을 보이겠다~: ${g'}^*(\mathcal{L})$은
풍부하고,
$g'_*(\mathcal{O}_{X''})\otimes_{\mathcal{O}_X}\mathcal{R}(X)$는
\emph{상수 계수의 국소 자유} $\mathcal{R}(X)$-가군이다. 그러면
(\hyperref[II.6.6.10-ko]{6.6.10})의 조건으로 환원되어 명제가 증명된다.

이를 위해 $\mathcal{B}=g_*(\mathcal{O}_{X'})$라 하자. $X$의 기약
성분들을 바탕공간으로 갖는 $X$의 축소 닫힌 부분준스킴들을
$X_i$ ($1\leq i\leq n$)로 나타내자
(\hyperref[I.5.2.1-ko]{I, 5.2.1})~; 가정에 따라 이들은 정역
준스킴이다. $X'_i$를 $X'$의 닫힌 부분준스킴 $g^{-1}(X_i)$라 하고,

\oldpage[II]{135}

$g_i:X'_i\to X_i$를 $X'_i$에 대한 $g$의 제한이라 하자. 이는 유한
(\hyperref[II.6.1.5-ko]{6.1.5, (iii)}) 전사 사상이다. 또한
$\mathcal{B}_i=(g_i)_*(\mathcal{O}_{X'_i})$인
$\mathcal{O}_{X_i}$-대수 $\mathcal{B}_i$의 \emph{계수}를 $k_i$라 하자.
$\mathcal{B}_i\otimes_{\mathcal{O}_{X_i}}\mathcal{R}(X_i)$가 상수
준층이므로 (\hyperref[I.7.3.5-ko]{I, 7.3.5}), 계수 $k_i$는 기약 성분
가운데 \emph{오직} $X_i$와만 만나는 $X$의 모든 열린집합 $U$에 대해
$(\mathcal{O}_X|U)$-대수 $\mathcal{B}|U$의 계수이기도 하다. $T$가
$\mathcal{B}|T\simeq\mathcal{O}_X^m|T$를 만족하는 $X$의 열린집합이면,
앞의 주의로부터 $T\cap X_i\neq\varnothing$인 모든 지표 $i$에 대해
$k_i=m$임을 얻는다. 이제 $U$를 그 근방에서 $\mathcal{B}$가 국소 자유
$\mathcal{O}_X$-가군인 $X$의 점들로 이루어진 열린집합이라 하고,
$U_j$ ($1\leq j\leq s$)를 그 연결 성분들이라 하자. 이들은 $X$에서
열려 있고 유한 개이다($U$가 뇌터이기 때문이다). $V_j$를 열린집합
$g^{-1}(U_j)$ 위에 유도된 부분준스킴의 \emph{폐포}인 $X'$의 닫힌
부분준스킴이라 하자 (\hyperref[I.9.5.11-ko]{I, 9.5.11}). 앞에서 본
바에 따라 $X_i\cap U_j\neq\varnothing$인 모든 지표 $i$에 대해 계수
$k_i$들은 모두 같은 정수 $m_j$와 같다. 또한 같은 $X_i$가 지표가 서로
다른 두 $U_j$와 만날 수 없음에 주의하자. $X_i\cap U=\varnothing$인
지표 $i$들을 $i_\lambda$라 하자. 모든 $k_i$의 곱을 $k$라 하고,
$n_i=k/k_i$로 놓은 뒤, 준스킴 $X''$을 다음과 같이 정의한다. 각
$j$ ($1\leq j\leq s$)에 대해 $V_j$와 동형인 준스킴 $k/m_j$개를
취하고, 각 $\lambda$에 대해 $X'_{i_\lambda}$와 동형인 준스킴
$k/k_{i_\lambda}$개를 취한다~; $X''$은 이 모든 준스킴의 \emph{합}이다.
$X''$의 각 합성분 위에서 표준 주입으로 주어지는 사상
$g'':X''\to X'$을 정의한다. $g''$은 유한 우세 사상이므로 전사임이
명백하다(유한 사상은 닫힌 사상이기 때문이다
(\hyperref[II.6.1.10-ko]{6.1.10})). $g'=g\circ g''$로 놓으면 이는
유한 전사 사상 $X''\to X$이다. 또한
${g'}^*(\mathcal{L})={g''}^*(g^*(\mathcal{L}))$이므로
${g'}^*(\mathcal{L})$은 풍부한 $\mathcal{O}_{X''}$-가군이다
(\hyperref[II.5.1.12-ko]{5.1.12}). 이제 이 새로운 준스킴 $X''$에
대하여 $X'$의 $k_i$들과 같은 방식으로 정의한 계수들은 모두
\emph{$k$와 같음}이 명백하다. 따라서
(\hyperref[I.7.3.3-ko]{I, 7.3.3})을 고려하면, $X$의 모든 아핀
열린집합 $T$에 대하여
\[
\bigl(g'_*(\mathcal{O}_{X''})
\otimes_{\mathcal{O}_X}\mathcal{R}(X)\bigr)|T
\]
는 $(\mathcal{R}(X)|T)^k$와 동형인
$(\mathcal{R}(X)|T)$-가군임을 곧바로 얻는다.

\begin{corollary}[6.6.12]
\label{II.6.6.12-ko}
(\hyperref[II.6.6.11-ko]{6.6.11})의 명제에서 $W=X$이면, 가역
$\mathcal{O}_X$-가군 $\mathcal{L}$이 $f$에 관해 풍부하기 위한
필요충분조건은 $g^*(\mathcal{L})$이 $f\circ g$에 관해 풍부한 것이다.
\end{corollary}

\begin{remark}[6.6.13]
\label{II.6.6.13-ko}
제III장에서 우리는 $Y$가 뇌터이고 $f$가 유한형이며, 바탕공간이
$X-W$인 $X$의 축소 닫힌 부분준스킴에 제한한 $f$가 \emph{고유}이면
(\hyperref[II.6.6.12-ko]{6.6.12})의 결론이 여전히 성립함을 볼 것이다.
그러나 제V장에서는 체 $K$ 위의 \emph{대수적} 스킴 $X$ 가운데 (구조 사상
$X\to\operatorname{Spec}(K)$가 고유하지 않고) 그 정규화 $X'$은
준아핀이지만 $X$ 자체는 준아핀이 아닌 예들을 제시할 것이다(따라서
$\mathcal{O}_{X'}$은 풍부하고
(\hyperref[II.5.1.12-ko]{5.1.12}) 사상 $X'\to X$는 유한 전사이지만
(\hyperref[II.6.3.10-ko]{6.3.10}), $\mathcal{O}_X$는 풍부하지 않다).
다음 항에서 우리는 «~준아핀~»을 «~아핀~»으로 바꾸면 이러한 상황이
일어날 수 없음을 볼 것이다.
\end{remark}

\subsection{슈발레의 정리.}
\label{subsection:II.6.7-ko}

우리는 세르 판정법 (\hyperref[II.5.2.1-ko]{5.2.1})을 사용하여 다음
정리를 증명할 것이다. 이 정리는 대수적 스킴의 경우 C.~슈발레가 다른
방법으로 증명하였다.

\oldpage[II]{136}

\begin{theorem}[6.7.1]
\label{II.6.7.1-ko}
$X$를 아핀 스킴, $Y$를 뇌터 준스킴, $f:X\to Y$를 유한 전사 사상이라
하자. 그러면 $Y$는 아핀 스킴이다.
\end{theorem}

$f_{\mathrm{red}}:X_{\mathrm{red}}\to Y_{\mathrm{red}}$가 유한 사상임은
명백하다 (\hyperref[II.6.1.5-ko]{6.1.5, (vi)})~; $X_{\mathrm{red}}$는 아핀
스킴이고, $Y$가 뇌터이므로 $Y$가 아핀인 것과 $Y_{\mathrm{red}}$가 아핀인
것은 동치이므로 (\hyperref[I.6.1.7-ko]{I, 6.1.7}), $Y$가
\emph{축소}라고 가정할 수 있다. $Y$의 임의의 닫힌 부분집합 $Y'$에
대하여, 바탕공간이 $Y'$인 $Y$의 축소 부분준스킴은 유일하다
(\hyperref[I.5.1.2-ko]{I, 5.1.2})~; 그 역상 $f^{-1}(Y')$은
$X\times_Y Y'$과 표준적으로 동형이고
(\hyperref[I.4.4.1-ko]{I, 4.4.1}), $X$의 닫힌 부분준스킴이므로
아핀이며, $f^{-1}(Y')$에 대한 $f$의 제한은 $f\times_Y 1_{Y'}$과
동일시되고 유한 전사 사상이다
(\hyperref[II.6.1.5-ko]{6.1.5, (iii)}). 따라서 뇌터 귀납법 원리
(\hyperref[0.2.2.2-ko]{0, 2.2.2})에 따라
(\hyperref[I.6.1.7-ko]{I, 6.1.7})을 고려하면, \emph{모든 닫힌
부분집합 $Y'\neq Y$}에 대하여 바탕공간이 $Y'$인 $Y$의 모든 닫힌
부분준스킴이 아핀이라는 가정 아래에서 정리를 증명하는 것으로 환원할 수
있다. 이로부터 \emph{(닫힌) 지지집합 $Z$가 $Y$와 다른 모든 연접
$\mathcal{O}_Y$-가군 $\mathcal{F}$에 대하여
$\mathrm{H}^1(Y,\mathcal{F})=0$이다}. 실제로 바탕공간이 $Z$인 $Y$의
닫힌 부분준스킴 $Y'$이 존재하여, $j:Y'\to Y$가 표준 포함 사상이면
$\mathcal{F}=j_*(j^*(\mathcal{F}))$이다
(\hyperref[I.9.3.5-ko]{I, 9.3.5})~; 따라서
(\hyperref[II.5.2.3-ko]{5.2.3})에 의해
$\mathrm{H}^1(Y,\mathcal{F})
=\mathrm{H}^1(Y',j^*(\mathcal{F}))=0$이다
(\hyperref[I.5.1.9.2-ko]{I, 5.1.9.2}).

먼저 $Y$가 기약이 아니라고 가정하고, $Y'$을 $Y$의 한 기약 성분,
$Y''=Y-Y'$이라 하자~; 바탕공간이 $Y'$인 $Y$의 축소 닫힌
부분준스킴도 다시 $Y'$으로 나타내고, 표준 포함 사상 $Y'\to Y$를 $j$로
나타내자. $\mathcal{F}$를 연접 $\mathcal{O}_Y$-가군이라 하고, 표준
준동형
\[
\rho:\mathcal{F}\longrightarrow
\mathcal{F}'=j_*(j^*(\mathcal{F}))
\]
을 생각하자 (\hyperref[0.4.4.3-ko]{0, 4.4.3})~;
$\mathcal{F}'$은 (\hyperref[0.5.3.10-ko]{0, 5.3.10})과
(\hyperref[0.5.3.12-ko]{0, 5.3.12})에 따라 연접
$\mathcal{O}_Y$-가군이다. 실제로 $Y'$을 정의하는
$\mathcal{O}_Y$의 아이디얼층을 $\mathcal{J}$로 나타내면
$j_*(\mathcal{O}_{Y'})=\mathcal{O}_Y/\mathcal{J}$이다. 따라서
$\mathcal{G}=\operatorname{Ker}\rho$와
$\mathcal{K}=\operatorname{Im}\rho$도 연접 $\mathcal{O}_Y$-가군이다
(\hyperref[0.5.3.4-ko]{0, 5.3.4})~; 그런데 정의에 따라 $Y'$의 일반점
$y$에서 $\mathcal{F}'$의 섬유 $\mathcal{F}'_y$는
$\mathcal{F}_y$와 같다. 실제로 $y$는 $Y'$의 내부점이고, $Y$가
축소이므로 $\mathcal{J}_y=0$이다. 따라서 $y$는 $\mathcal{G}$의
(닫힌) 지지집합에 포함되지 않는다~; 한편 $\mathcal{F}'$의 지지집합은
($\mathcal{K}$의 지지집합은 \emph{더욱이}) $Y'$에 포함된다~; 다시 말해
$\mathcal{G}$와 $\mathcal{K}$의 지지집합은 \emph{$Y$와 다르다}.
그러므로
$\mathrm{H}^1(Y,\mathcal{G})=\mathrm{H}^1(Y,\mathcal{K})=0$이고,
완전열 $0\to\mathcal{G}\to\mathcal{F}\to\mathcal{K}\to0$에 적용한
코호몰로지 완전열로부터 $\mathrm{H}^1(Y,\mathcal{F})=0$이다. 이제
세르 판정법 (\hyperref[II.5.2.1-ko]{5.2.1})으로 결론을 얻는다.

이제 $Y$가 기약이고, 따라서 \emph{정역}이라고 가정하자. $X$도
\emph{정역}이라고 가정할 수 있다~: 실제로 $X_i$들을 바탕공간이 $X$의
기약 성분인 $X$의 축소 닫힌 부분준스킴들이라 하고
(\hyperref[I.5.2.1-ko]{I, 5.2.1}), $f_i$를 $X_i$에 대한 $f$의
제한이라 하자. 이때 $f_i$들 가운데 적어도 하나는 우세 사상이고, 이는
유한 사상이므로 (\hyperref[II.6.1.5-ko]{6.1.5}) 전사 사상이다
(\hyperref[II.6.1.10-ko]{6.1.10})~; 한편 $X_i$는 아핀 스킴이므로,
정리의 서술에서 $X$를 $X_i$로 대체할 수 있음을 알 수 있다.

\begin{lemma}[6.7.1.1]
\label{II.6.7.1.1-ko}
$X$, $Y$를 두 뇌터 정역 준스킴, $x$ (각각 $y$)를 $X$ (각각 $Y$)의
일반점, $f:X\to Y$를 유한 전사 사상이라 하자. $y$의 아핀 열린 근방
$U$와 단면 $g\in\Gamma(X,\mathcal{L})$가 존재하여

\oldpage[II]{137}

$x\in X_g\subset f^{-1}(U)$가 되도록 하는 가역 $\mathcal{O}_X$-가군
$\mathcal{L}$이 주어졌다고 하자. 그러면 두 정수 $m>0$, $n>0$, 준동형
$u:\mathcal{O}_Y^m\to f_*(\mathcal{L}^{\otimes n})$, 그리고 $y$의 열린
근방 $V$가 존재하여, 제한 $u|V$는 $\mathcal{O}_Y^m|V$에서
$f_*(\mathcal{L}^{\otimes n})|V$ 위로 가는 동형사상이다.
\end{lemma}

$C$를 $U$의 (정역인) 환, $k=\mathcal{O}_y$를 그 분수체라 하자~;
$f$가 유한이므로 $U'=f^{-1}(U)$는 아핀이다
(\hyperref[II.1.3.2-ko]{1.3.2})~; $D$를 그 (정역인) 환이라 하고,
그 분수체를 $K=\mathcal{O}_x$라 하자~; 가정에 따라 $D$는 유한형
$C$-가군이므로 (\hyperref[II.6.1.4-ko]{6.1.4}), $K$는 $k$의 유한
차수 확대이다. 섬유
\[
f^{-1}(y)=X\times_Y\operatorname{Spec}(k(y))
=X\times_Y\operatorname{Spec}(\mathcal{O}_y)
\]
는 $\operatorname{Spec}(K)$와 동일시된다
(\hyperref[I.3.6.5-ko]{I, 3.6.6})~; $K$의 $k$ 위의 기저를 이루는
$D$의 원소들을 $s_i$ ($1\leq i\leq m$)라 하자. 어떤 $n>0$가
존재하여, $X_g$ 위의 $\mathcal{L}^{\otimes n}$의 단면
$(s_i|X_g)g^{\otimes n}$은 $X$ 위의 $\mathcal{L}^{\otimes n}$의 단면
$b_i$ ($1\leq i\leq m$)로 연장된다
(\hyperref[I.9.3.1-ko]{I, 9.3.1}). 정의에 따라 $b_i$들은 또한 $Y$ 위의
$f_*(\mathcal{L}^{\otimes n})$의 단면들이므로 준동형
$u:\mathcal{O}_Y^m\to f_*(\mathcal{L}^{\otimes n})$을 정의한다
(\hyperref[0.5.1.1-ko]{0, 5.1.1})~; 이 $u$가 요구를 만족함을 보이겠다.
$\mathcal{L}^{\otimes n}|U'=\widetilde{M}$이며, 여기서 $M$은 유한형
$D$-가군이다. 따라서 $\varphi$를 $f$의 제한
$f^{-1}(U)\to U$에 대응하는 단사 준동형 $C\to D$라 하면,
$M_{[\varphi]}$는 유한형 $C$-가군이다~; 그러므로
\[
f_*(\mathcal{L}^{\otimes n})|U=(M_{[\varphi]})^\sim
\]
(\hyperref[I.1.6.3-ko]{I, 1.6.3})는 연접이고, $U$는 $Y$의 임의의 아핀
열린집합이므로 $f_*(\mathcal{L}^{\otimes n})$은 연접이다~; 더 나아가
$u|U=\widetilde{\theta}$이며, 여기서 $\theta$는 $C$-준동형
$C^m\to M_{[\varphi]}$이고, $u_y=\theta_y$는 준동형
\[
\theta\otimes1:K^m=C^m\otimes K\longrightarrow
M_{[\varphi]}\otimes K,
\]
이다. 그런데 후자는 정의에 따라 \emph{동형사상}이다. 실제로
$(b_i)_x$들은 $M_{[\varphi]}\otimes K$의 $k$ 위의 기저를 이루고,
가정에 따라 $g_x$는 $(\mathcal{L}^{\otimes n})_x$의 생성원이다.
따라서 $\mathcal{O}_Y$-가군 $\operatorname{Ker}u$와
$\operatorname{Coker}u$의 지지집합은 $y$를 포함하지 않는다~;
이 $\mathcal{O}_Y$-가군들은 연접이므로
(\hyperref[0.5.3.3-ko]{0, 5.3.3}), 그 지지집합들은 닫혀 있고
(\hyperref[0.5.2.2-ko]{0, 5.2.2}), 이로써 보조정리가 증명된다.

이제 우리가 생각하는 경우에는 $X$가 아핀이므로
(\hyperref[I.1.1.10-ko]{I, 1.1.10}), $\mathcal{L}=\mathcal{O}_X$로 놓으면
보조정리 (\hyperref[II.6.7.1.1-ko]{6.7.1.1})의 가정들이 성립한다~;
$\mathcal{A}=\mathcal{O}_Y$, $\mathcal{B}=f_*(\mathcal{O}_X)$로 놓겠다.
세르 판정법 (\hyperref[II.5.2.1.1-ko]{5.2.1.1})에 따라, 모든 연접
$\mathcal{O}_Y$-가군 $\mathcal{F}$에 대하여
$\mathrm{H}^1(Y,\mathcal{F})=0$임을 증명하면 충분하다~; 심지어
$\mathcal{F}\subset\mathcal{O}_Y$인 경우에 이를 증명하면 충분하고,
이 포함 관계와 $Y$가 정역이라는 사실로부터 $\mathcal{F}$에는 꼬임이
없음이 따른다~; 실제로 우리는 \emph{꼬임이 없는} \emph{모든} 연접
$\mathcal{O}_Y$-가군 $\mathcal{F}$에 대하여
$\mathrm{H}^1(Y,\mathcal{F})=0$임을 보이겠다. 이제 준동형
$u:\mathcal{A}^m\to\mathcal{B}$는 준동형
\[
v:\mathcal{G}
=\mathcal{H}om_{\mathcal{A}}(\mathcal{B},\mathcal{F})
\longrightarrow
\mathcal{H}om_{\mathcal{A}}(\mathcal{A}^m,\mathcal{F})
=\mathcal{F}^m.
\]
을 정의한다. 먼저 $v$가 \emph{단사}임을 보이자~: 가정에 따라
$\mathcal{T}=\operatorname{Coker}u$의 지지집합은 $V$와 만나지 않으므로,
$\mathcal{T}$는 꼬임 $\mathcal{O}_Y$-가군이다
(\hyperref[I.7.4.6-ko]{I, 7.4.6})~; 완전열
\[
\mathcal{A}^m\longrightarrow\mathcal{B}\longrightarrow\mathcal{T}
\longrightarrow0
\]
과 함자 $\mathcal{H}om_{\mathcal{A}}$의 왼쪽 완전성으로부터 완전열
\[
0\longrightarrow
\mathcal{H}om_{\mathcal{A}}(\mathcal{T},\mathcal{F})
\longrightarrow\mathcal{G}\xrightarrow{v}\mathcal{F}^m.
\]
을 얻는다. 그런데 $\mathcal{F}$에는 꼬임이 없으므로
$\mathcal{H}om_{\mathcal{A}}(\mathcal{T},\mathcal{F})=0$이고
(\hyperref[0.5.2.6-ko]{0, 5.2.6}), 따라서 주장이 성립한다. 그러므로
완전열
\[
0\longrightarrow\mathcal{G}\longrightarrow\mathcal{F}^m
\longrightarrow\operatorname{Coker}v\longrightarrow0
\]

\oldpage[II]{138}

을 얻는다. 여기서 $\mathcal{G}$와 $\operatorname{Coker}v$는 연접
$\mathcal{O}_Y$-가군이다
(\hyperref[0.5.3.4-ko]{0, 5.3.4} 및
\hyperref[0.5.3.5-ko]{5.3.5})~; 코호몰로지 완전열에 따라
$\mathrm{H}^1(Y,\mathcal{G})
=\mathrm{H}^1(Y,\operatorname{Coker}v)=0$임을 보이면
$\mathrm{H}^1(Y,\mathcal{F}^m)
=(\mathrm{H}^1(Y,\mathcal{F}))^m=0$, 따라서
$\mathrm{H}^1(Y,\mathcal{F})=0$임을 얻기에 충분하다. 그런데 $v|V$는
동형사상이므로 $\operatorname{Coker}v$의 지지집합은 $Y$와 다르고,
따라서 처음의 가정에 의해
$\mathrm{H}^1(Y,\operatorname{Coker}v)=0$이다. 한편 $\mathcal{G}$는 연접
$\mathcal{B}$-가군이고 (\hyperref[I.9.6.4-ko]{I, 9.6.4}), $X$는 $Y$
위에서 아핀이므로 $\mathcal{G}$와 동형인 $f_*(\mathcal{K})$를 갖는
준연접 $\mathcal{O}_X$-가군 $\mathcal{K}$가 존재한다
(\hyperref[II.1.4.3-ko]{1.4.3})~; $X$가 아핀이므로
$\mathrm{H}^1(X,\mathcal{K})=0$이고
(\hyperref[I.5.1.9.2-ko]{I, 5.1.9.2}), 따라서
(\hyperref[II.5.2.3-ko]{5.2.3})에 의해
$\mathrm{H}^1(Y,\mathcal{G})=0$이다. 이로써 정리
(\hyperref[II.6.7.1-ko]{6.7.1})의 증명이 끝난다.

\begin{corollary}[6.7.2]
\label{II.6.7.2-ko}
$X$를 뇌터 준스킴, $(X_i)_{1\leq i\leq n}$를 닫힌집합들로 이루어진
공간 $X$의 유한 덮개라 하자. $X$가 아핀이기 위한 필요충분조건은 각
$i$에 대하여 바탕공간이 $X_i$이고 아핀인 $X$의 닫힌 부분준스킴이
존재하는 것이다.
\end{corollary}

실제로 이 조건이 성립한다고 하고, $X'$을 $X_i$들의 \emph{합} 스킴이라
하자~; $X'$이 아핀임은 분명하며, $f$의 $X_i$에 대한 제한을 표준 주입
사상으로 잡아 전사 사상 $f:X'\to X$를 정의한다.
(\hyperref[II.6.7.1-ko]{6.7.1})에 따라 이제 $f$가 \emph{유한}임을
확인하면 충분하며, 이는 (\hyperref[II.6.1.5-ko]{6.1.5})에서
보였다.

\begin{corollary}[6.7.3]
\label{II.6.7.3-ko}
뇌터 준스킴 $X$가 아핀이기 위한 필요충분조건은 $X$의 기약 성분들을
바탕공간으로 갖는 축소 닫힌 부분준스킴들이 아핀인 것이다.
\end{corollary}

\section{값매김 판정법}
\label{section:II.7-ko}

이 절에서는 사상의 분리성과 고유성에 대한 값매김 판정법, 즉 $A$가
값매김환일 때 $\operatorname{Spec}(A)$ 꼴의 변하는 보조 스킴을 사용하는
판정법을 제시한다. 적절한 ``뇌터'' 가정 아래에서는 이 판정법들에서
$A$가 \emph{이산} 값매김환인 경우로 한정할 수 있다. 아마 이후에
적용할 경우는 이것뿐일 것이며, 일반적인 경우에 임의의 값매김환을
도입한 것은 고전적인 전개와 연결하기 위해서일 뿐이다.

\subsection{값매김환에 대한 복습.}
\label{subsection:II.7.1-ko}

\begin{env}[7.1.1]
\label{II.7.1.1-ko}
값매김환을 특징짓는 여러 동치인 성질 가운데 다음 것을 사용하겠다~:
환 $A$가 체가 아닌 정역이고, $A$의 분수체 $K$에 포함되며 $K$와는
다른 국소환들의 집합에서 지배관계에 관하여 \emph{극대}이면 $A$를
\emph{값매김환}이라고 한다
(\hyperref[I.8.1.1-ko]{I, 8.1.1}). 값매김환은 \emph{정수적으로
닫혀 있음}을 상기하자. $A$가 값매김환이면, $A$의 모든 소 아이디얼
$\mathfrak{p}\neq0$에 대하여 $A_{\mathfrak{p}}$도 값매김환이다.
\end{env}

\begin{env}[7.1.2]
\label{II.7.1.2-ko}
$K$를 체, $A$를 체가 아닌 $K$의 국소 부분환이라 하자~;

\oldpage[II]{139}

그러면 분수체가 $K$이고 $A$를 지배하는 값매김환이 존재한다
(\cite[p.~1-07, lemme~2]{I-1}).

한편 $B$를 값매김환, $k$를 그 잉여체, $K$를 그 분수체, $L$을 $k$의
확대체라 하자. 그러면 $B$를 지배하고 잉여체가 $L$인 \emph{완비}
값매김환 $C$가 존재한다. 실제로 $L$은 순수 초월 확대체
$L'=k(T_\mu)_{\mu\in M}$의 대수적 확대체이다~; $B$에 대응하는 $K$의
값매김을 $K'=K(T_\mu)_{\mu\in M}$의 값매김으로 연장하여 $L'$이 이
값매김의 잉여체가 되게 할 수 있음이 알려져 있다
(\cite[p.~98]{II-24})~; $B$를 이 연장된 값매김의 값매김환의 완비화로
바꾸면, $B$가 완비이고 $L$이 $k$의 대수적 폐포인 경우로 한정할 수
있음을 알 수 있다. $\overline{K}$가 $K$의 대수적 폐포이면 $B$를
정의하는 값매김을 $\overline{K}$로 연장할 수 있고, 이에 대응하는
잉여체는 $k$의 대수적 폐포이다. 이는 $k[T]$의 일계수 다항식의
계수들을 $\overline{K}$로 들어 올려 보면 알 수 있다. 따라서 결국
$L=k$인 경우로 귀착되며, 이 경우에는 $B$의 완비화를 $C$로 택하면
충분하다.
\end{env}

\begin{env}[7.1.3]
\label{II.7.1.3-ko}
$K$를 체, $A$를 $K$의 부분환이라 하자~; $K$에서 $A$의 정수적 폐포
$A'$은 $A$를 포함하고 분수체가 $K$인 값매김환들의 교집합이다
(\cite[p.~51, th.~2]{I-11}). 명제
(\hyperref[II.7.1.2-ko]{7.1.2})는 다음과 같은 동치인 기하학적 형태로
표현된다~:
\end{env}

\begin{proposition}[7.1.4]
\label{II.7.1.4-ko}
$Y$를 준스킴, $p:X\to Y$를 사상, $x$를 $X$의 점, $y=p(x)$라 하고,
$y'\neq y$를 $y$의 특수화
(\hyperref[0.2.1.2-ko]{0, 2.1.2})라 하자. 그러면 값매김환의 스펙트럼인
국소 스킴 $Y'$과 분리 사상 $f:Y'\to Y$가 존재하여, $a$를 $Y'$의
유일한 닫힌점, $b$를 $Y'$의 일반점이라 하면 $f(a)=y'$이고
$f(b)=y$이다. 더 나아가 다음 두 부가 성질 중 하나가 성립한다고
가정할 수 있다~:
\begin{enumerate}
\item[\rm{(i)}] $Y'$은 잉여체가 대수적으로 닫힌 \emph{완비}
값매김환의 스펙트럼이고, $k(y)$-준동형 $k(x)\to k(b)$가 존재한다.
\item[\rm{(ii)}] $k(y)$-동형사상
$k(x)\xrightarrow{\sim}k(b)$가 존재한다.
\end{enumerate}
\end{proposition}

$Y_1$을 바탕공간이 $\overline{\{y\}}$인 $Y$의 축소 닫힌 부분준스킴
(\hyperref[I.5.2.1-ko]{I, 5.2.1})이라 하고, $X_1$을 그 역상인 닫힌
부분준스킴 $p^{-1}(Y_1)$이라 하자~; 가정에 따라
$y'\in\overline{\{y\}}$이고 $k(x)$는 $X$에서나 $X_1$에서나 같으므로,
$Y$가 일반점 $y$를 갖는 \emph{정역 준스킴}이라고 가정할 수 있다~;
$\mathcal{O}_{y'}$는 이때 체가 아닌 국소 정역이고 그 분수체는
$\mathcal{O}_y=k(y)$이며, $k(x)$는 $k(y)$의 확대체이다. 조건
$f(a)=y'$, $f(b)=y$와 부가 조건 (i) (각각 (ii))를 실현하기 위해
$Y'=\operatorname{Spec}(A')$로 잡는다. 여기서 $A'$은
$\mathcal{O}_{y'}$를 지배하고, 완비이며, 잉여체가 $k(x)$의 대수적으로
닫힌 확대체인 값매김환이다(각각 $\mathcal{O}_{y'}$를 지배하고
분수체가 $k(x)$인 값매김환이다)~; $A'$의 존재는
(\hyperref[II.7.1.2-ko]{7.1.2})에 의해 보장된다.

\begin{env}[7.1.5]
\label{II.7.1.5-ko}
국소환 $A$에 대하여, 극대 아이디얼 $\mathfrak{m}$과 다른 소 아이디얼이
존재하고 $\mathfrak{m}$과 다른 $A$의 모든 소 아이디얼이
\emph{극소} 소 아이디얼이면 $A$는 \emph{차원이 $1$}이라고 한다~;
$A$가 \emph{정역}일 때에는 $\mathfrak{m}$과 $(0)$이 유일한 소 아이디얼이고
$\mathfrak{m}\neq(0)$이라고 하는 것과 같다~; 다시 말해,
$Y=\operatorname{Spec}(A)$는 두

\oldpage[II]{140}

점 $a$, $b$로만 이루어진다~: $a$는 유일한 \emph{닫힌점}이고,
$\mathfrak{j}_a=\mathfrak{m}$이며 $k(a)=k$는 \emph{잉여체}
$k=A/\mathfrak{m}$이다~; $b$는 $Y$의 \emph{일반점}이고,
$\mathfrak{j}_b=(0)$이며, 집합 $\{b\}$는 $\varnothing$ 및 $Y$와 다른
$Y$의 유일한 열린집합(따라서 \emph{모든 곳에서 조밀한} 열린집합)이고,
$k(b)=K$는 $A$의 \emph{분수체}이다.
\end{env}

\begin{env}[7.1.6]
\label{II.7.1.6-ko}
차원이 $1$인 뇌터 국소환 $A$에 대하여 다음 조건들이 서로 동치임이
알려져 있다
(\cite[p.~2-08 et 17-01]{I-1})~:
\begin{enumerate}
\item[\rm{a)}] $A$는 정규이다~;
\item[\rm{b)}] $A$는 정칙이다~;
\item[\rm{c)}] $A$는 값매김환이다~;
\end{enumerate}
더욱이 이때 $A$는 \emph{이산 값매김환}이다. 이산 값매김환에 대해서는
명제 (\hyperref[II.7.1.2-ko]{7.1.2})와
(\hyperref[II.7.1.3-ko]{7.1.3})에 다음과 같은 대응 명제가 있다~:
\end{env}

\begin{proposition}[7.1.7]
\label{II.7.1.7-ko}
$A$를 체가 아닌 뇌터 국소 정역, $K$를 그 분수체, $L$을 $K$의
유한 생성 확대체라 하자~; 그러면 분수체가 $L$이고 $A$를 지배하는
이산 값매김환이 존재한다.
\end{proposition}

먼저 $L=K$라고 하자. $\mathfrak{m}$을 $A$의 극대 아이디얼,
$(x_1,\ldots,x_n)$을 $\mathfrak{m}$의 영이 아닌 생성원계라 하고,
$B$를 $K$의 뇌터 부분환
$A[x_2/x_1,\ldots,x_n/x_1]$이라 하자. $B$의 아이디얼
$\mathfrak{m}B$가 주 아이디얼 $x_1B$와 같음은 자명하다~;
$\mathfrak{p}$가 $x_1B$의 극소 소 아이디얼이면 $\mathfrak{p}$의 높이는
$1$이다 (\cite[t.~I, p.~277]{I-13})~; 다시 말해
$B_{\mathfrak{p}}$는 \emph{차원이 $1$}인 뇌터 국소환이다~;
$\mathfrak{p}B_{\mathfrak{p}}\cap A$가 $\mathfrak{m}$을 포함하지만
$1$은 포함하지 않는 $A$의 아이디얼임은 분명하므로 이는
$\mathfrak{m}$과 같고, 따라서 $B_{\mathfrak{p}}$는 $A$를
\emph{지배한다} (\hyperref[I.8.1.1-ko]{I, 8.1.1}). Krull--Akizuki 정리
(\cite[p.~293]{II-25})에 따라 $B_{\mathfrak{p}}$의 정수적 폐포 $C$는
뇌터환이다($C$가 반드시 유한 생성 $B_{\mathfrak{p}}$-가군인 것은
아니다)~; $\mathfrak{n}$이 $C$의 극대 아이디얼이면
$C_{\mathfrak{n}}$은 차원이 $1$인 정규 뇌터 국소환
(\cite[p.~295]{II-25})이고, 따라서 $B_{\mathfrak{p}}$를 지배하고
\emph{더욱이} $A$를 지배하는 이산 값매김환이다.

이제 $L$이 $K$의 유한 생성 확대체이면, 앞서 보인 바에 따라 $A$가
이미 이산 값매김환인 경우로 한정할 수 있다. $w$를 $A$에 대응하는
$K$의 값매김이라 하자~; $w$를 \emph{연장하는} $L$의 이산 값매김
$w'$이 존재한다~: 실제로 $L$의 생성원 수에 관한 귀납법으로
$L=K(\alpha)$인 경우로 환원되며, 이때 명제는 고전적인 결과이다
(\cite[p.~106]{II-24}).

\begin{corollary}[7.1.8]
\label{II.7.1.8-ko}
$A$를 뇌터 정역, $K$를 그 분수체, $L$을 $K$의 유한 생성 확대체라
하자. 그러면 $L$ 안에서 $A$의 정수적 폐포는 분수체가 $L$이고
$A$를 포함하는 모든 이산 값매김환의 교집합이다.
\end{corollary}

실제로 그러한 이산 값매김환은 정규이므로 $A$ 위에서 정수적인
$L$의 모든 원소를 \emph{더욱이} 포함한다. 따라서 $x\in L$이 $A$
위에서 정수적이지 않을 때, 분수체가 $L$이고 $A$를 포함하지만 $x$는
포함하지 않는 이산 값매김환 $C$가 존재함을 보이면 충분하다. 실제로
$x$에 대한 가정은 $x\notin B=A[1/x]$임을 뜻한다. 다시 말해,
$1/x$는 뇌터환 $B$에서 가역원이 아니다. 따라서 $1/x$를 포함하는
$B$의 소 아이디얼 $\mathfrak{p}$가 존재한다. 국소 정역
$B_{\mathfrak{p}}$는 뇌터이고 $L$에 포함되며, $L$은
$B_{\mathfrak{p}}$의 분수체(이는 $K$를 포함한다)의 유한 생성
확대체이다. (\hyperref[II.7.1.7-ko]{7.1.7})에 따라 분수체가 $L$이고
$B_{\mathfrak{p}}$를 지배하는 이산 값매김환 $C$가 존재한다~;
$1/x\in\mathfrak{p}B_{\mathfrak{p}}$는 $C$의 극대 아이디얼에
속하므로 $x\notin C$이고, 이로써 증명이 끝난다.

(\hyperref[II.7.1.7-ko]{7.1.7})의 기하학적 형태는 다음과 같다~:

\oldpage[II]{141}

\begin{proposition}[7.1.9]
\label{II.7.1.9-ko}
$Y$를 국소 뇌터 준스킴, $p:X\to Y$를 국소 유한형 사상,
$x$를 $X$의 한 점, $y=p(x)$라 하고, $y'\neq y$를 $y$의
특수화라 하자. 그러면 이산 값매김환의 스펙트럼인 국소 스킴 $Y'$,
분리 사상 $f:Y'\to Y$, 그리고 $Y'$에서 $X$로 가는 $Y$-유리 사상
$g$가 존재하여, $a$를 $Y'$의 닫힌점, $b$를 그 일반점이라 할 때
$f(a)=y'$, $f(b)=y$, $g(b)=x$이고 다음 가환 그림에서
\[
\tag{7.1.9.1}
\label{II.7.1.9.1-ko}
\begin{tikzcd}[column sep=large,row sep=large]
& k(x) \arrow[dl,"\gamma"'] \\
k(b) & k(y) \arrow[u,"\pi"] \arrow[l,"\varphi"']
\end{tikzcd}
\]
($\pi$, $\varphi$, $\gamma$는 각각 $p$, $f$, $g$에 대응하는
준동형이다) $\gamma$는 전단사이다.
\end{proposition}

(\hyperref[II.7.1.4-ko]{7.1.4})에서와 같이
(\hyperref[I.6.3.4-ko]{I, 6.3.4, (iv)})를 고려하면 $Y$가 일반점 $y$를
갖는 정역 준스킴이라고 가정할 수 있고, 문제는 $X$와 $Y$ 위에서 국소적이므로
$p$가 유한형이라고 가정할 수 있다~; 그러면
(\hyperref[II.7.1.4-ko]{7.1.4})의 상황에 있으며, 여기에 $k(x)$가
$k(y)$의 \emph{유한형} 확대체이고
(\hyperref[I.6.4.11-ko]{I, 6.4.11}) $\mathcal{O}_{y'}$가 뇌터라는
성질이 더해진다~; 따라서 (\hyperref[II.7.1.7-ko]{7.1.7})을 적용하여
$Y'=\operatorname{Spec}(A')$로 잡을 수 있다. 여기서 $A'$은
$\mathcal{O}_{y'}$를 지배하고 $k(x)$를 분수체로 갖는 이산 값매김환이다.
이로써 $\gamma$가 전단사이고 $\pi$와 $\varphi$가 사상 $p$와 $f$에
대응하는 가환 그림 (\hyperref[II.7.1.9.1-ko]{7.1.9.1})을 얻는다.
또한 $X$와 $Y$는 국소 뇌터이고
(\hyperref[I.6.6.2-ko]{I, 6.6.2}) $Y'$은 정역 준스킴이므로, 동형사상
$\gamma$에 대응하는 $Y'$에서 $X$로 가는 $Y$-유리 사상 $g$가
유일하게 존재한다 (\hyperref[I.7.1.15-ko]{I, 7.1.15}). 이로써 증명이
끝난다.

\subsection{분리성의 값매김 판정법}
\label{subsection:II.7.2-ko}

\begin{proposition}[7.2.1]
\label{II.7.2.1-ko}
$X$, $Y$를 두 준스킴, $f:X\to Y$를 준콤팩트 사상이라 하자. 다음 두
조건은 서로 동치이다~:
\begin{enumerate}
\item[\rm a)] 사상 $f$는 닫힌 사상이다.
\item[\rm b)] 모든 $x\in X$와 $y=f(x)$의 특수화 $y'\neq y$에 대하여,
$f(x')=x$를 만족하는 $x$의 특수화 $x'$이 존재한다.
\end{enumerate}
\end{proposition}

조건 b)는 $f(\overline{\{x\}})=\overline{\{y\}}$임을 나타내므로 a)의
결과이다. b)가 a)를 함의함을 보이기 위해 바탕공간 $X$의 닫힌 부분집합
$X'$을 생각하자~; $Y'=\overline{f(X')}$라 놓고 $Y'=f(X')$임을
보이자. 각각 $X'$과 $Y'$을 바탕공간으로 갖는 $X$와 $Y$의 축소 닫힌
부분준스킴을 생각하자
(\hyperref[I.5.2.1-ko]{I, 5.2.1})~; 그러면 다음 그림을 가환이게 하는 사상
$f':X'\to Y'$이 존재한다
\[
\begin{tikzcd}[column sep=large,row sep=large]
X' \arrow[r,"f'"] \arrow[d] & Y' \arrow[d] \\
X \arrow[r,"f"'] & Y
\end{tikzcd}
\]
(\hyperref[I.5.2.2-ko]{I, 5.2.2}). $f$가 준콤팩트이므로 $f'$도
준콤팩트이다. 따라서 $f$가 준콤팩트 \emph{우세 사상}이면 조건 b)가
\oldpage[II]{142}
$f(X)=Y$를 함의함을 보이는 것으로 환원된다. 이제 $y'$을 $Y$의 한 점,
$y$를 $y'$을 포함하는 $Y$의 한 기약 성분의 일반점이라 하자~; b)에 따라
$f^{-1}(y)$가 공집합이 아님을 보이면 충분하다. 그런데 이 성질은 $f$가
준콤팩트 우세 사상이라는 사실의 결과임을 알고 있다
(\hyperref[I.6.6.5-ko]{I, 6.6.5}).

\begin{corollary}[7.2.2]
\label{II.7.2.2-ko}
$f:X\to Y$를 준콤팩트 몰입 사상이라 하자. 바탕공간 $X$가 $Y$에서
닫혀 있기 위한 필요충분조건은 $X$가 그 각 점의 모든 특수화($Y$에서)를
포함하는 것이다.
\end{corollary}

\begin{proposition}[7.2.3]
\label{II.7.2.3-ko}
$Y$를 준스킴(각각, 국소 뇌터 준스킴), $f:X\to Y$를
사상(각각, 국소 유한형 사상)이라 하자. 다음 조건들은
서로 동치이다~:
\begin{enumerate}
\item[\rm a)] $f$는 분리 사상이다.
\item[\rm b)] 대각 사상 $X\to X\times_YX$는 준콤팩트이고,
$A$가 값매김환(각각, 이산 값매김환)인 꼴의 모든
$Y$-준스킴 $Y'=\operatorname{Spec}(A)$에 대하여, $Y'$의 일반점에서
일치하는 $Y'$에서 $X$로 가는 두 $Y$-사상은 서로 같다.
\item[\rm c)] 대각 사상 $X\to X\times_YX$는 준콤팩트이고,
$A$가 값매김환(각각, 이산 값매김환)인 꼴의 모든
$Y$-준스킴 $Y'=\operatorname{Spec}(A)$에 대하여, $Y'$의 일반점에서
일치하는 $X'=X_{(Y')}$의 두 $Y'$-단면은 서로 같다.
\end{enumerate}
\end{proposition}

b)와 c)의 동치는 $Y'$에서 $X$로 가는 $Y$-사상과 $X'$의
$Y'$-단면 사이의 일대일 대응
(\hyperref[I.3.3.14-ko]{I, 3.3.14})에서 따른다. $X$가 $Y$ 위에서 분리이면
$Y'$이 정역 준스킴이므로 (\hyperref[I.7.2.2.1-ko]{I, 7.2.2.1})에 따라 조건 b)가
성립한다. 이제 b)가 대각 사상
$\Delta:X\to X\times_YX$가 닫힌 사상임을 함의한다는 것을 증명하면 되고, 이는
$\Delta$가 (\hyperref[II.7.2.2-ko]{7.2.2})의 판정법을 만족함을 보이는 것과
같다. 실제로 $z$를 대각선 $\Delta(X)$의 한 점, $z'\neq z$를
$X\times_YX$에서 $z$의 한 특수화라 하자. 그러면
(\hyperref[II.7.1.4-ko]{7.1.4})에 따라 값매김환 $A$와
$Y'=\operatorname{Spec}(A)$에서 $X\times_YX$로 가는 사상 $f$가 존재하여,
$Y'$의 닫힌점 $a$를 $z'$에, $Y'$의 일반점 $b$를 $z$에 보낸다~; 이 사상은
$Y'$을 $(X\times_YX)$-준스킴으로 만들며, 따라서 \emph{특히} $Y$-준스킴으로
만든다. $X\times_YX$의 두 사영을 $f$와 합성하면 $Y'$에서 $X$로 가는
두 $Y$-사상 $g_1$, $g_2$를 얻는데, 가정에 따라 이들은 점 $b$에서
일치한다~; 따라서 이들은 하나의 동일한 사상 $g$와 같고, 이는
(\hyperref[I.5.3.1-ko]{I, 5.3.1})에 따라 $f$가 $f=\Delta\circ g$로
분해됨을 뜻하므로 $z'\in\Delta(X)$이다. $Y$가 국소 뇌터이고 $f$가
국소 유한형이라고 가정하면 $X\times_YX$는 국소 뇌터이다
(\hyperref[I.6.6.7-ko]{I, 6.6.7})~; 그러면
(\hyperref[II.7.1.9-ko]{7.1.9})에 따라 앞의 논증에서 $A$를
이산 값매김환으로 가정할 수 있다.

\begin{remark}[7.2.4]
\label{II.7.2.4-ko}
\begin{enumerate}
\item[\rm (i)] $\Delta$가 준콤팩트라는 가정은 $Y$가 국소 뇌터이고 $f$가
국소 유한형일 때에는 항상 성립한다. 실제로 이때 $X\times_YX$는 국소 뇌터이다
(\hyperref[I.6.6.4-ko]{I, 6.6.4, (i)}). 일반적인 경우에는 이 가정은 또한
$X$의 아핀 열린집합들로 이루어진 모든 덮개 $(U_\alpha)$에 대하여
집합 $U_\alpha\cap U_\beta$가 \emph{준콤팩트}라는 것을 뜻한다.
\item[\rm (ii)] $f$가 분리 사상이기 위해서는 조건 b) 또는 조건 c)가
\emph{완비}이고 잉여체가 \emph{대수적으로 닫힌} 값매김환 $A$에 대하여
성립하는 것으로 \emph{충분하다}~; 이는
(\hyperref[II.7.2.3-ko]{7.2.3})과
(\hyperref[II.7.1.4-ko]{7.1.4})의 증명에서 따른다.
\end{enumerate}
\end{remark}

\oldpage[II]{143}
\subsection{고유성의 값매김 판정법}
\label{subsection:II.7.3-ko}

\begin{proposition}[7.3.1]
\label{II.7.3.1-ko}
$A$를 값매김환, $Y=\operatorname{Spec}(A)$, $b$를 $Y$의 일반점,
$X$를 정역 스킴이라 하고, $f:X\to Y$를 닫힌 사상이라 하자. $f^{-1}(b)$가
한 점 $x$만으로 이루어지고 이에 대응하는 준동형
$k(b)\to k(x)$가 전단사라고 하자. 그러면 $f$는 동형 사상이다.
\end{proposition}

$f$가 닫힌 우세 사상이므로 $f(X)=Y$이다~;
(\hyperref[I.4.2.2-ko]{I, 4.2.2})에 따라, $Y$의 모든 $y'\neq b$에 대하여
$f(x')=y'$인 점 $x'$가 \emph{유일하게} 존재하고 이에 대응하는 준동형
$\mathcal{O}_{y'}\to\mathcal{O}_{x'}$가 전단사임을 증명하면 충분하다. 그러면
$f$는 위상동형이기 때문이다. 이제 $f(x')=y'$이면 $\mathcal{O}_{x'}$는
$K=k(x)=k(b)$에 포함되고 $\mathcal{O}_{y'}$를 지배하는 국소환이다~; 후자는
국소환 $A_{y'}$이므로, $K$를 분수체로 갖는 값매김환이다
(\hyperref[II.7.1.1-ko]{7.1.1}). 한편 $x'$는 $X$의 일반점이 아니므로
(\hyperref[0.2.1.3-ko]{0, 2.1.3}) $\mathcal{O}_{x'}\neq K$이다~; 따라서
$\mathcal{O}_{x'}=\mathcal{O}_{y'}$이다. $X$가 정역 스킴이므로
$\mathcal{O}_{x'}=\mathcal{O}_{x''}$이라는 관계는 $x'=x''$을 함의하며
(\hyperref[I.8.2.2-ko]{I, 8.2.2}), 이로써 증명이 끝난다.

\begin{env}[7.3.2]
\label{II.7.3.2-ko}
$A$를 값매김환, $Y=\operatorname{Spec}(A)$, $b$를 $Y$의 일반점이라 하자.
그러면 $\mathcal{O}_b=k(b)$는 $A$의 분수체 $K$와 같다~; 또한
$f:X\to Y$를 사상이라 하자. $X$의 $Y$-\emph{유리 단면}들은 점 $b$에서
($b$의 근방들에 정의된) $Y$-단면들의 \emph{싹}들과 일대일 대응함이
알려져 있으므로 (\hyperref[I.7.1.4-ko]{I, 7.1.4}), 표준 사상
\[
\tag{7.3.2.1}
\label{II.7.3.2.1-ko}
\Gamma_{\mathrm{rat}}(X/Y)\longrightarrow
\Gamma(f^{-1}(b)/\operatorname{Spec}(K))
\]
을 얻는다. 한편 정의에 따라 (\hyperref[I.3.4.5-ko]{I, 3.4.5}),
$\Gamma(f^{-1}(b)/\operatorname{Spec}(K))$의 원소들은
\emph{$f^{-1}(b)=X\otimes_AK$의 $K$-유리점들}과 동일시된다.
$f$가 \emph{분리}이면 $Y$가 정역 스킴이므로
(\hyperref[I.5.4.7-ko]{I, 5.4.7}), 사상
(\hyperref[II.7.3.2.1-ko]{7.3.2.1})은 \emph{단사}이다.

(\hyperref[II.7.3.2.1-ko]{7.3.2.1})과 표준 사상
$\Gamma(X/Y)\to\Gamma_{\mathrm{rat}}(X/Y)$
(\hyperref[I.7.1.2-ko]{I, 7.1.2})을 합성하면 표준 사상
\[
\tag{7.3.2.2}
\label{II.7.3.2.2-ko}
\Gamma(X/Y)\longrightarrow \Gamma(f^{-1}(b)/\operatorname{Spec}(K)).
\]
을 얻는다. $f$가 \emph{분리}이면 이 사상 역시 \emph{단사}이다
(\hyperref[I.5.4.7-ko]{I, 5.4.7}).
\end{env}

\begin{proposition}[7.3.3]
\label{II.7.3.3-ko}
$A$를 분수체가 $K$인 값매김환, $Y=\operatorname{Spec}(A)$, $b$를 $Y$의
일반점, $f:X\to Y$를 \emph{분리}이고 \emph{닫힌} 사상이라 하자. 그러면
표준 사상 (\hyperref[II.7.3.2.2-ko]{7.3.2.2})은 전단사이다 (이는 이
사상이 \emph{전사}라는 것과 동치이고, $X$의 $Y$-유리 단면들이
\emph{모든 곳에서 정의됨}을 함의한다).
\end{proposition}

따라서 $x$를 $f^{-1}(b)$의 점으로서 \emph{$K$ 위에서 유리}인 점이라
하자. $f$가 분리이므로 $f$에 대응하는 사상
$f^{-1}(b)\to\operatorname{Spec}(K)$도 분리이고
(\hyperref[I.5.5.1-ko]{I, 5.5.1, (iv)}), $f^{-1}(b)$의 모든 단면은 닫힌
몰입 사상이므로 (\hyperref[I.5.4.6-ko]{I, 5.4.6}), $\{x\}$는
\emph{$f^{-1}(b)$ 안에서} 닫혀 있다. $X$에서 $\{x\}$의 폐포
$\overline{\{x\}}$를 바탕공간으로 갖는 $X$의 축소 닫힌 부분준스킴
$X'$을 생각하자. $f$를 $X'$에 제한한 사상이
(\hyperref[II.7.3.1-ko]{7.3.1})의 가정들을 만족함은 분명하다. 따라서
이 제한은 $X'$에서 $Y$로 가는 \emph{동형사상}이고, 그 역동형사상이
구하는 $X$의 $Y$-단면이다.

\begin{env}[7.3.4]
\label{II.7.3.4-ko}
다음 두 결과를 서술하기 위해, 제IV장에서 정당화하고 논의할 용어를
사용하겠다~: $F$를 준스킴 $Y$의 부분집합
\oldpage[II]{144}
이라 하면, $Y$ 안에서 $F$의 \emph{여차원}이라 부르고
$\operatorname{codim}_YF$로 나타내는 것은 $z$가 $F$를 움직일 때의 정수
$\dim(\mathcal{O}_z)$들의 하한이다.
\end{env}

\begin{corollary}[7.3.5]
\label{II.7.3.5-ko}
$Y$를 축소 국소 뇌터 준스킴, $N$을 $Y$가 정칙이지 않은 점 $y\in Y$들의
집합이라 하자 (\hyperref[0.4.1.4-ko]{0, 4.1.4})~; 이때
$\operatorname{codim}_YN\geq2$라고 가정한다. $f:X\to Y$를 유한형이고
분리이며 닫힌 사상, $g$를 $X$의 $Y$-유리 단면이라 하자~; $Y'$을 $g$가
정의되지 않는 $Y$의 점들의 집합이라 하면 (이 집합은 \emph{닫혀 있다}
(\hyperref[I.7.2.1-ko]{I, 7.2.1})),
$\operatorname{codim}_YY'\geq2$이다.
\end{corollary}

$\dim\mathcal{O}_z\leq1$인 모든 점 $z\in Y$에서 $g$가 정의됨을 증명하면
충분하다. $\dim\mathcal{O}_z=0$이면 $z$는 $Y$의 한 기약 성분의
일반점이고 (\hyperref[I.1.1.14-ko]{I, 1.1.14}), 따라서 $Y$의 모든
곳에서 조밀한 모든 열린집합에, 특히 $g$의 정의역에 속한다. 이제
$\dim\mathcal{O}_z=1$이라고 하자~; 그러면 가정에 따라
$\mathcal{O}_z$는 뇌터 정칙 국소환이고, 따라서
(\hyperref[II.7.1.6-ko]{7.1.6}) \emph{이산 값매김환}이다.
$Z=\operatorname{Spec}(\mathcal{O}_z)$라 하자~; $U=Y-Y'$은 모든 곳에서
조밀하므로 $U\cap Z$는 공집합이 아니며
(\hyperref[I.2.4.2-ko]{I, 2.4.2}), $g'$을 $g$가 유도하는 $Z$에서 $X$로
가는 유리 사상이라 하자 (\hyperref[I.7.2.8-ko]{I, 7.2.8})~; $g'$이
\emph{사상}임을 보이면 충분하다 (\hyperref[I.7.2.9-ko]{I, 7.2.9}).
그런데 $g'$은 $Z$-준스킴 $f^{-1}(Z)=X\times_YZ$의 $Z$-유리 단면으로
볼 수 있다~; $f$에 대응하는 사상 $f^{-1}(Z)\to Z$가 닫힌 사상임은
분명하고, (\hyperref[I.5.5.1-ko]{I, 5.5.1, (i)})에 따라 분리이다~;
따라서 (\hyperref[II.7.3.3-ko]{7.3.3})으로부터 $g'$이 모든 곳에서
정의됨을 얻는다~; $Z$가 축소이고 $X$가 $Y$ 위에서 분리이므로 $g'$은
사상이다 (\hyperref[I.7.2.2-ko]{I, 7.2.2}).

\begin{corollary}[7.3.6]
\label{II.7.3.6-ko}
$S$를 국소 뇌터 준스킴, $X$와 $Y$를 두 $S$-준스킴이라 하자~; $Y$가
축소라고 가정하고, 더 나아가 $Y$가 정칙이지 않은 점 $y\in Y$들의 집합
$N$이 $\operatorname{codim}_YN\geq2$를 만족한다고 가정한다~; 마지막으로
구조 사상 $X\to S$가 고유라고 가정한다. $f$를 $Y$에서 $X$로 가는
$S$-유리 사상, $Y'$을 $f$가 정의되지 않는 $Y$의 점들의 집합이라 하자~;
그러면 $\operatorname{codim}_YY'\geq2$이다.
\end{corollary}

$Y$에서 $X$로 가는 $S$-유리 사상들을 $X\times_SY$의 $Y$-유리
단면들과 동일시할 수 있음이 알려져 있다
(\hyperref[I.7.1.2-ko]{I, 7.1.2})~; 구조 사상 $X\times_SY\to Y$는 닫힌
사상이므로 (\hyperref[II.5.4.1-ko]{5.4.1}),
(\hyperref[II.7.3.5-ko]{7.3.5})를 적용하면 따름정리를 얻는다.

\begin{remark}[7.3.7]
\label{II.7.3.7-ko}
(\hyperref[II.7.3.5-ko]{7.3.5})와
(\hyperref[II.7.3.6-ko]{7.3.6})에서 $Y$에 둔 가정들은 특히 $Y$가
\emph{정규}일 때 (\hyperref[0.4.1.4-ko]{0, 4.1.4}),
(\hyperref[II.7.1.6-ko]{7.1.6})에 따라 성립한다.
\end{remark}

보편적으로 닫힌 사상 (각각 고유 사상)을
(\hyperref[II.7.3.3-ko]{7.3.3})의 역에 해당하는 명제로 다음과 같이
특징지을 수 있다~:

\begin{theorem}[7.3.8]
\label{II.7.3.8-ko}
$Y$를 준스킴(각각 국소 뇌터 준스킴), $f:X\to Y$를 분리 준콤팩트
사상(각각 유한형 사상)이라 하자. 다음 조건들은 동치이다~:
\begin{enumerate}
\item[\rm a)] $f$는 보편적으로 닫혀 있다(각각 고유이다).
\item[\rm b)] $A$가 분수체 $K$를 갖는 값매김환(각각 이산 값매김환)인
$Y$-스킴 $Y'=\operatorname{Spec}(A)$ 각각에 대하여, 표준 단사 준동형
$A\to K$에 대응하는 표준 사상
\[
\operatorname{Hom}_Y(Y',X)\longrightarrow
\operatorname{Hom}_Y(\operatorname{Spec}(K),X)
\]
은 전사이다(각각 전단사이다).
\oldpage[II]{145}
\item[\rm c)] $A$가 값매김환(각각 이산 값매김환)인 꼴
$Y'=\operatorname{Spec}(A)$의 모든 $Y$-스킴에 대하여, $Y'$-준스킴
$X_{(Y')}$에 관한 표준 사상
(\hyperref[II.7.3.2.2-ko]{7.3.2.2})은 전사이다(각각 전단사이다).
\end{enumerate}
\end{theorem}

b)와 c)의 동치는 (\hyperref[I.3.3.14-ko]{I, 3.3.14})에서 곧바로
따른다~; a)는 c)를 함의한다. 실제로 두 경우 모두 a)에 따라
$f_{(Y')}$는 분리되고
(\hyperref[I.5.5.1-ko]{I, 5.5.1, (iv)}) 닫혀 있으므로
(\hyperref[II.7.3.3-ko]{7.3.3})을 적용하면 충분하다. 이제 b)가 a)를
함의함을 증명하면 된다. 먼저 $Y$가 임의의 준스킴이고 $f$가 분리
준콤팩트 사상인 경우를 생각하자. $f$가 조건 b)를 만족하면 임의의
$Y$-준스킴 $Y''$에 대하여 $f_{(Y'')}:X_{(Y'')}\to Y''$도 조건 b)를
만족한다. 이는 b)와 c)의 동치 및 모든 사상 $Y'\to Y''$에 대한 등식
$X_{(Y'')}\times_{Y''}Y'=X\times_YY'$
(\hyperref[I.3.3.9.1-ko]{I, 3.3.9.1})에서 따른다~; 또한 $f$가 분리되고
준콤팩트이면 $f_{(Y'')}$도 그러하므로
(\hyperref[I.5.5.1-ko]{I, 5.5.1, (iv)} 및
\hyperref[I.6.6.4-ko]{6.6.4, (iii)}), b)가 $f$가 \emph{닫힌}
사상임을 함의하는지만 증명하면 된다. 이를 위해서는
(\hyperref[II.7.2.1-ko]{7.2.1})의 조건 b)를 확인하면 충분하다. 따라서
$x\in X$라 하고, $y'$를 $y=f(x)$의 특수화로서 $y$와 다르다고 하자~;
(\hyperref[II.7.1.4-ko]{7.1.4})에 따라 값매김환의 스펙트럼인 스킴
$Y'$과 분리 사상 $g:Y'\to Y$가 존재하여, $a$를 $Y'$의 닫힌점, $b$를
일반점이라 하면 $g(a)=y'$, $g(b)=y$이고 $k(y)$-준동형
$k(x)\to k(b)$가 존재한다. 이 준동형은 표준적으로 $Y$-사상
$\operatorname{Spec}(k(b))\to X$에 대응하고
(\hyperref[I.2.4.6-ko]{I, 2.4.6}), 따라서 b)에 의해 앞의 사상에
대응하는 $Y$-사상 $h:Y'\to X$가 존재한다. 그러면 $h(b)=x$이다~;
$h(a)=x'$로 놓으면 $x'$은 $x$의 특수화이고
$f(x')=f(h(a))=g(a)=y'$이다.

이제 $Y$가 국소 뇌터이고 $f$가 유한형이면, 대각 사상
$X\to X\times_YX$가 준콤팩트이므로
(\hyperref[II.7.2.4-ko]{7.2.4}), 가정 b)는 우선
(\hyperref[II.7.2.3-ko]{7.2.3})에 따라 $f$가 \emph{분리} 사상임을
함의한다. 또한 $f$가 고유임을 확인하기 위해서는
(\hyperref[II.5.6.3-ko]{5.6.3})에 따라 모든 \emph{유한형}
$Y$-준스킴 $Y''$에 대하여 $f_{(Y'')}:X_{(Y'')}\to Y''$가
\emph{닫힌} 사상임을 보이면 충분하다. 이때 $Y''$은 국소 뇌터이므로,
$Y'$을 이산 값매김환의 스펙트럼으로 취하고
(\hyperref[II.7.1.4-ko]{7.1.4}) 대신
(\hyperref[II.7.1.9-ko]{7.1.9})를 적용하여 첫 번째 경우의 논증을
되풀이할 수 있다.

\begin{remark}[7.3.9]
\label{II.7.3.9-ko}
\begin{enumerate}
\item[\rm (i)] $Y$가 임의의 준스킴이고 $f$가 분리 사상일 때, $f$가
보편적으로 닫혀 있기 위해서는 값매김환 $A$가 \emph{완비}이고 그
잉여체가 \emph{대수적으로 닫혀 있는} 경우에 조건 b) 또는 조건 c)가
성립하는 것으로 \emph{충분하다}~; 이는 실제로 위의 증명과
(\hyperref[II.7.1.4-ko]{7.1.4})에서 따른다.
\item[\rm (ii)] (\hyperref[II.7.3.8-ko]{7.3.8})의 판정법 c)에서 사영
사상 $X\to Y$가 닫힌 사상이라는 사실
(\hyperref[II.5.5.3-ko]{5.5.3})의 고전적 방법에 더 가까운 새로운
증명이 나온다. 실제로 $Y$가 아핀이라고 가정할 수 있고, 따라서
$X$를 사영 다발 $\mathbf{P}_Y^n$의 닫힌 부분준스킴과 동일시할 수 있다
(\hyperref[II.5.3.3-ko]{5.3.3})~; $X\to Y$가 닫힌 사상임을 증명하려면
구조 사상 $\mathbf{P}_Y^n\to Y$가 그러함을 확인하면 충분하며,
(\hyperref[II.7.3.8-ko]{7.3.8})의 판정법 c)와
(\hyperref[II.4.1.3.1-ko]{4.1.3.1})을 함께 적용하면 다음 사실을
증명하는 것으로 환원된다~: \emph{$Y$가 분수체 $K$를 갖는 값매김환
$A$의 스펙트럼이면, $K$에 값을 갖는 $\mathbf{P}_Y^n$의 모든 점은
($Y$의 일반점으로 제한하여) $A$에 값을 갖는 $\mathbf{P}_Y^n$의 한
점에서 유도된다}. 그런데 모든 가역 $\mathcal{O}_Y$-가군은 자명하므로
(\hyperref[I.2.4.8-ko]{I, 2.4.8}),
(\hyperref[II.4.2.6-ko]{4.2.6})에 따라 $K$에 값을 갖는
$\mathbf{P}_Y^n$의 한 점은 $K$의 원소들
$(\zeta c_0,\zeta c_1,\ldots,\zeta c_n)$의 한 동치류와 동일시된다.
여기서 $\zeta\neq0$이고 $c_i$들은 모두 영은 아닌 $K$의 원소들이다.
$c_i$들에 $A$의 한 원소를 곱하되 그
\oldpage[II]{146}
값매김을 적당히 택하면, 모든 $c_i$가 $A$에 속하고 그 가운데 적어도
하나가 가역이라고 가정할 수 있다.
그러면 (\hyperref[II.4.2.6-ko]{4.2.6})에 따라 계
$(c_0,\ldots,c_n)$은 $A$에 값을 갖는 $\mathbf{P}_Y^n$의 한 점도
정의하며, 이로써 주장이 증명된다.
\item[\rm (iii)] $Y$-준스킴 $X$를 주는 것을 $Y$-준스킴 $Y'$에
함자
\[
X(Y')=\operatorname{Hom}_Y(Y',X)
\]
을 주는 것과 동치로 볼 때 판정법
(\hyperref[II.7.2.3-ko]{7.2.3})과
(\hyperref[II.7.3.8-ko]{7.3.8})이 특히 편리하다~; 예를 들어 이
판정법들을 사용하여 특정 조건 아래에서 ``피카르 스킴''이 고유임을
증명할 것이다.
\end{enumerate}
\end{remark}

\begin{corollary}[7.3.10]
\label{II.7.3.10-ko}
$Y$를 정역 스킴(각각 국소 뇌터 정역 스킴), $X$를 정역 스킴,
$f:X\to Y$를 우세 사상이라 하자.
\begin{enumerate}
\item[\rm (i)] $f$가 준콤팩트이고 보편적으로 닫혀 있으면, 분수체가
$X$ 위의 유리함수체 $R(X)$이고 $Y$의 한 국소환을 지배하는 모든
값매김환은 $X$의 한 국소환도 지배한다.
\item[\rm (ii)] 역으로 $f$가 유한형이고, 분수체가 $R(X)$인 모든
값매김환(각각 모든 이산 값매김환)이 (i)의 성질을 만족한다고
가정하자. 그러면 $f$는 고유이다.
\end{enumerate}
\end{corollary}

먼저 가정들은 모든 경우에 $f$가 분리 사상임을 함의함을 주목하자
(\hyperref[I.5.5.9-ko]{I, 5.5.9}).
\begin{enumerate}
\item[\rm (i)] $K=R(Y)$, $L=R(X)$라 하고, $y$를 $Y$의 한 점,
$A$를 분수체가 $L$이고 $\mathcal{O}_y$를 지배하는 값매김환이라 하자~;
그러면 단사 준동형 $\mathcal{O}_y\to A$는 $Y'=\operatorname{Spec}(A)$에서
$Y$로 가는 사상 $h$를 정의한다
(\hyperref[I.2.4.4-ko]{I, 2.4.4}). 여기서 $a$를 $Y'$의 닫힌점이라 하면
$h(a)=y$이다~; 또한 $Y$의 일반점 $\eta$는
$\operatorname{Spec}(\mathcal{O}_y)$의 일반점이기도 하므로, $b$를
$Y'$의 일반점이라 할 때 $h(b)=\eta$이다(가정상 $K\subset L$이기
때문이다). $\xi$를 $X$의 일반점이라 하면 가정상
$k(\xi)=k(b)=L$이고, 따라서 $g(b)=\xi$인 $Y$-사상
$g:\operatorname{Spec}(L)\to X$가 존재한다~;
(\hyperref[II.7.3.8-ko]{7.3.8})에 따라 $g$는 $Y$-사상
$g':Y'\to X$에서 유도된다. $x=g'(a)$로 놓으면 $A$가
$\mathcal{O}_x$를 지배함은 명백하다.
\item[\rm (ii)] 문제는 $Y$ 위에서 국소적이므로 언제나 $Y$가
아핀(각각 아핀이고 뇌터)이라고 가정할 수 있다. $f$가 유한형이므로
두 경우 모두 차우 보조정리 (\hyperref[II.5.6.1-ko]{5.6.1})를 적용할
수 있다. 따라서 사영 사상 $p:P\to Y$, 몰입 사상 $j:X'\to P$, 그리고
사영적이고 전사이며 쌍유리인 사상 $g:X'\to X$가 존재하여($X'$은
정역이다) 다음 도식이 가환한다~:
\[
\begin{tikzcd}[column sep=large,row sep=large]
P \arrow[d,"p"'] & X' \arrow[l,"j"'] \arrow[d,"g"] \\
Y & X \arrow[l,"f"]
\end{tikzcd}
\]
$j$가 \emph{닫힌} 몰입 사상임을 증명하면 충분하다. 실제로 이때
$f\circ g=p\circ j$는 사영 사상이고 따라서 고유이며, $g$가
전사이므로 $f$도 고유이다
(\hyperref[II.5.4.3-ko]{5.4.3}). $Z$를 바탕공간이
$\overline{j(X')}$인 $P$의 축소 닫힌 부분준스킴이라 하자
(\hyperref[I.5.2.1-ko]{I, 5.2.1})~; $X'$이 정역이므로 $j$는
$i\circ h$로 인수분해된다. 여기서 $i:Z\to P$는 표준 주입 사상,
$h:X'\to Z$는 우세인 열린 몰입 사상이고
(\hyperref[I.5.2.3-ko]{I, 5.2.3}), $Z$는 정역이다~;

\oldpage[II]{147}
또한 $Z$는 $Y$ 위에서 사영적이므로, $P$가 \emph{정역}이고 $j$가
\emph{우세이며 쌍유리}인 경우로 국한할 수 있음을 알 수 있다. 이제
모든 것은 $j$가 \emph{전사}임을 확인하는 일로 귀착된다. $z\in P$라
하자~; $\mathcal{O}_z$는 분수체가
\[
L=R(P)=R(X')=R(X).
\]
인 정역 국소환(각각 뇌터 정역 국소환)이다. $z$가 $P$의 일반점이 아닌
경우로 국한할 수 있다. 그러면
(\hyperref[II.7.1.2-ko]{7.1.2} 및
\hyperref[II.7.1.7-ko]{7.1.7})에 따라 분수체가 $L$이고
$\mathcal{O}_z$를 지배하는 값매김환(각각 이산 값매김환) $A$가
존재한다. $y=p(z)$로 놓으면 \emph{더욱이} $A$는 $\mathcal{O}_y$를
지배하며, 따라서 가정에 의해 $A$가 $\mathcal{O}_x$를 지배하는 점
$x\in X$가 존재한다. $g$가 고유이므로 증명의 첫째 부분에 따라 어떤
$x'\in X'$에 대하여 $A$는 $\mathcal{O}_{x'}$를 지배한다~; 따라서
$\mathcal{O}_z$와 $\mathcal{O}_{j(x')}=\mathcal{O}_{x'}$는 서로
연관되어 있고 (\hyperref[I.8.1.4-ko]{I, 8.1.4}), $P$가 스킴이므로
이는 $z=j(x')$를 함의한다 (\hyperref[I.8.2.2-ko]{I, 8.2.2}). 이로써
증명이 끝난다.
\end{enumerate}

\begin{corollary}[7.3.11]
\label{II.7.3.11-ko}
$X$, $Y$를 두 정역 스킴, $f:X\to Y$를 준콤팩트이고 보편적으로 닫힌
우세 사상이라 하자. 또한 $Y$가 (정역인) 환 $B$의 아핀 스킴이라고
가정하자. 그러면 $\Gamma(X,\mathcal{O}_X)$는 $R(X)$ 안에서 $B$의
정수적 폐포의 한 부분환과 표준적으로 동형이다.
\end{corollary}

실제로 (\hyperref[I.8.2.1.1-ko]{I, 8.2.1.1})에 따라
$\Gamma(X,\mathcal{O}_X)$는 $x\in X$에 대한 $\mathcal{O}_x$들의
교집합과 동일시된다~; (\hyperref[II.7.3.10-ko]{7.3.10}),
(\hyperref[II.7.1.2-ko]{7.1.2}), (\hyperref[II.7.1.3-ko]{7.1.3})에 따라
$\Gamma(X,\mathcal{O}_X)$는 따라서 $B$를 포함하고 $R(X)$를 분수체로
갖는 값매김환들의 교집합에 포함된다~; 이제 결론은
(\hyperref[II.7.1.3-ko]{7.1.3})에서 따른다.

\begin{remark}[7.3.12]
\label{II.7.3.12-ko}
(\hyperref[II.7.3.11-ko]{7.3.11})의 가정 아래에서 $R(X)$가 $R(Y)$의
유한 생성 확대체라고 가정하면, 많은 경우에 $\Gamma(X,\mathcal{O}_X)$가
환 $B=\Gamma(Y,\mathcal{O}_Y)$ 위의 \emph{유한형 가군}이라고 결론지을
수 있다. 예를 들어 $B$가 \emph{체 위의 유한형 대수}인 경우가 그러하다.
실제로 이때 $B$의 분수체의 유한 생성 확대체 안에서 $B$의 정수적 폐포가
유한형 $B$-가군임이 알려져 있기 때문이다([13], t.~I, p.~267, th.~9)~;
따라서 결론은 (\hyperref[II.7.3.11-ko]{7.3.11})과 $B$가 뇌터라는
사실에서 따른다.

특히 \emph{체 $K$ 위에서 고유하고 아핀인 스킴 $X$는 유한 스킴이다}. 실제로
(\hyperref[II.1.6.4-ko]{1.6.4}), (\hyperref[II.5.4.6-ko]{5.4.6}),
(\hyperref[I.6.4.4-ko]{I, 6.4.4, c)})에 따라 $X$가 \emph{축소}인 경우로
환원할 수 있다. 또한 $X$의 (유한 개의) 각 기약 성분을 바탕공간으로
갖는 닫힌 부분준스킴이 $K$ 위에서 유한함을 증명하면 충분하다~;
그러므로 (\hyperref[II.5.4.5-ko]{5.4.5})를 고려하면 결국 $X$가
\emph{정역}인 경우로 환원된다. 그러나 이때 결과는 위의 고찰에서
따른다.

제III장에서는 이 마지막 명제를 다른 방법으로 다시 얻을 것이며, 더
일반적으로 $f:X\to Y$가 고유 사상이고 $Y$가 국소 뇌터이면 모든 연접
$\mathcal{O}_X$-가군 $\mathcal{F}$에 대하여 $f_*(\mathcal{F})$가
연접임을 보이는 결과의 따름정리로 얻을 것이다
(\agref{III.4.4.2-ko}{III, 4.4.2}).

끝으로 고전 대수기하학에서는 판정법
(\hyperref[II.7.3.10-ko]{7.3.10})을 고유 사상의 \emph{정의}로 삼는다는
것을 지적해 두자. 우리가 이를 언급한 것은 오직 이 이유에서이며, 우리가
아는 모든 응용에서는 판정법 (\hyperref[II.7.3.8-ko]{7.3.8})이 더
다루기 쉬워 보인다.
\end{remark}

\oldpage[II]{148}
\subsection{대수곡선과 차원 1인 함수체}
\label{subsection:II.7.4-ko}

이 번호의 목적은 대수곡선의 고전적 개념(예를 들어 C.~Chevalley의 책
[23]에 서술된 개념)이 스킴의 언어로 어떻게 정식화되는지를 보이는
것이다. 이 번호 전체에서 \emph{$k$는 체를 나타내고, 고려하는 모든
스킴은 유한형 $k$-스킴이며, 모든 사상은 $k$-사상이다}.

\begin{proposition}[7.4.1]
\label{II.7.4.1-ko}
$X$를 $k$ 위의 유한형 준스킴(따라서 뇌터 준스킴), $x_i$
$(1\leq i\leq n)$를 $X$의 기약 성분 $X_i$들의 일반점, 그리고
$K_i=k(x_i)$ $(1\leq i\leq n)$라 하자. 다음 조건들은 동치이다~:
\begin{enumerate}
\item[\rm a)] 각 $K_i$는 초월차수가 1인 $k$의 확대체이다.
\item[\rm b)] $X$의 모든 닫힌점 $x$에 대하여 국소환
$\mathcal{O}_x$의 차원은 1이다(\hyperref[II.7.1.5-ko]{7.1.5}).
\item[\rm c)] $X_i$들과 다른 $X$의 닫힌 기약 부분집합들은
$X$의 닫힌점들이다.
\end{enumerate}
\end{proposition}

$X$가 준콤팩트이므로 $X$의 모든 닫힌 기약 부분집합 $F$는 닫힌점을
포함한다(\hyperref[0.2.1.3-ko]{0, 2.1.3}).
(\hyperref[I.2.4.2-ko]{I, 2.4.2})에 따라 $\mathcal{O}_x$의 소 아이디얼들과
$x$를 포함하는 $X$의 닫힌 기약 부분집합들 사이에는 전단사 대응이 있다
(\hyperref[I.1.1.14-ko]{I, 1.1.14})~; 따라서 b)와 c)의 동치는 곧바로
따른다. 한편 $\mathfrak{p}_\alpha$ $(1\leq\alpha\leq r)$가 뇌터 국소환
$\mathcal{O}_x$의 극소 소 아이디얼들이면, 국소환
$\mathcal{O}_x/\mathfrak{p}_\alpha$들은 정역이고 그 분수체들은
$x\in X_i$인 $K_i$들이다. 또한 유한형 정역 국소 $k$-대수의 차원은 그
분수체의 $k$ 위의 초월차수와 같다는 것이 알려져 있다
([1], p.~4-06, th.~2). 끝으로 $\mathcal{O}_x$의 차원은
$\mathcal{O}_x/\mathfrak{p}_\alpha$들의 차원의 상한이다~; 그런데 조건
a)는 이 차원들이 1임을 함의하므로 a)는 b)를 함의한다~; 역으로
$\mathcal{O}_x$의 차원이 1이면 어느 $\mathfrak{p}_\alpha$도
$\mathcal{O}_x$의 극대 아이디얼과 같을 수 없다. 그렇지 않으면
$\mathcal{O}_x$의 차원이 0일 것이기 때문이다~; 따라서 각
$\mathcal{O}_x/\mathfrak{p}_\alpha$의 차원은 1이고, 이는 b)가 a)를
함의함을 보인다.

(\hyperref[II.7.4.1-ko]{7.4.1})의 조건 아래에서 집합 $X$는
\emph{공집합이거나 무한집합}임을 주목하자. 이는
(\hyperref[I.6.4.4-ko]{I, 6.4.4})에서 곧바로 따른다.

\begin{definition}[7.4.2]
\label{II.7.4.2-ko}
(\hyperref[II.7.4.1-ko]{7.4.1})의 조건을 만족하는 공집합이 아닌
$k$-대수적 스킴을 $k$ 위의 대수곡선이라고 한다.
\end{definition}

제IV장에서 도입할 차원의 언어로는, $k$ 위의 대수곡선이란
\emph{모든 기약 성분의 차원이 1인} 공집합이 아닌 $k$-대수적 스킴이라고
표현할 수 있다.

$X$가 $k$ 위의 대수곡선이면, $X$의 기약 성분들을 바탕공간으로 갖는
$X$의 축소 닫힌 부분준스킴 $X_i$ $(1\leq i\leq n)$들도 $k$ 위의
대수곡선임을 주목하자.

\begin{corollary}[7.4.3]
\label{II.7.4.3-ko}
$X$를 기약 대수곡선이라 하자. $X$의 닫힌점이 아닌 유일한 점은 그
일반점이다. $X$와 다른 $X$의 닫힌 부분집합들은 닫힌점들의 유한집합이며,
이들은 또한 $X$에서 모든 곳에서 조밀하지 않은 유일한 부분집합들이다.
\end{corollary}

점 $x\in X$가 닫힌점이 아니면 $X$에서의 그 폐포는 $X$의 닫힌 기약
부분집합이므로, (\hyperref[II.7.4.1-ko]{7.4.1})에 따라 반드시 $X$
전체이다. 따라서 $x$는 $X$의 일반점이다. $X$와 다른 $X$의 닫힌
부분집합 $F$는 $X$의 일반점을 포함할 수 없으므로, 그 모든 점은
$X$에서 닫혀 있고(\emph{더욱이} $F$에서도 닫혀 있다)~; $F$를
바탕공간으로 갖는 $X$의 축소 닫힌 부분준스킴을 고려하면
(\hyperref[I.5.2.1-ko]{I, 5.2.1}),

\oldpage[II]{149}

(\hyperref[I.6.2.2-ko]{I, 6.2.2})에 따라 $F$가 유한 이산집합임이
따른다. 그러므로 $X$의 무한 부분집합의 $X$에서의 폐포는 반드시
$X$와 같다.

$X$가 임의의 대수곡선인 경우, $X$의 기약 성분들에
(\hyperref[II.7.4.3-ko]{7.4.3})을 적용하면 $X$의 닫힌점이 아닌 유일한
점들은 이 성분들의 일반점들임을 알 수 있다.

\begin{corollary}[7.4.4]
\label{II.7.4.4-ko}
$X$와 $Y$를 $k$ 위의 두 기약 대수곡선, $f:X\to Y$를
$k$-사상이라 하자. $f$가 우세이기 위한 필요충분조건은 모든 $y\in Y$에
대하여 $f^{-1}(y)$가 유한인 것이다.
\end{corollary}

실제로 $f$가 우세가 아니면, (\hyperref[II.7.4.3-ko]{7.4.3})에 따라
$f(X)$는 반드시 $Y$의 \emph{유한} 부분집합이다. 따라서 $Y$의 모든 점에
대하여 $f^{-1}(y)$가 유한일 수는 없다. 그렇다면 $X$가 유한이 되어
모순이기 때문이다 (\hyperref[II.7.4.1-ko]{7.4.1}). 역으로 $f$가 우세이면,
$Y$의 일반점 $\eta$와 다른 모든 $y\in Y$에 대하여 $\{y\}$가 $Y$에서
닫혀 있으므로 (\hyperref[II.7.4.3-ko]{7.4.3}), $f^{-1}(y)$는 $X$에서
닫혀 있다~; 한편 가정에 따라 $f^{-1}(y)$는 $X$의 일반점 $\xi$를
포함하지 않으므로, (\hyperref[II.7.4.3-ko]{7.4.3})에 따라 유한이다.
끝으로 $f$가 우세일 때 $f^{-1}(\eta)$가 유한임을 보이기 위해, 섬유
$f^{-1}(\eta)$가 $k(\eta)$ 위의 유한형 기약 스킴이고 그 일반점이
$\xi$임에 주의하자 (\hyperref[I.6.3.9-ko]{I, 6.3.9} 및
\hyperref[I.6.4.11-ko]{6.4.11}). $k(\xi)$와 $k(\eta)$는 모두 $k$의
유한형 확대이고 초월차수가 1로 같으므로, $k(\xi)$는 반드시
$k(\eta)$의 유한 차수 확대이다. 따라서 $\xi$는 $f^{-1}(\eta)$에서
닫혀 있고 (\hyperref[I.6.4.2-ko]{I, 6.4.2}), 그 결과
$f^{-1}(\eta)$는 한 점 $\xi$로 이루어진다.

제III장에서 우리는 뇌터 준스킴의 \emph{고유} 사상 $f:X\to Y$가 모든
$y\in Y$에 대하여 $f^{-1}(y)$가 유한이면 반드시 \emph{유한} 사상임을
볼 것이다~; 따라서 (\hyperref[II.7.4.4-ko]{7.4.4})로부터 기약
대수곡선에서 대수곡선으로 가는 우세 고유 사상은 \emph{유한} 사상임이
따른다.

\begin{corollary}[7.4.5]
\label{II.7.4.5-ko}
$X$를 $k$ 위의 대수곡선이라 하자. $X$가 정칙이기 위한 필요충분조건은
$X$가 정규인 것이며, 이는 다시 닫힌점들의 국소환이 이산 값매김환인 것과
동치이다.
\end{corollary}

이는 (\hyperref[II.7.4.1-ko]{7.4.1})의 조건 b)와
(\hyperref[II.7.1.6-ko]{7.1.6})에서 곧바로 따른다.

\begin{corollary}[7.4.6]
\label{II.7.4.6-ko}
$X$를 축소 대수곡선, $\mathcal{A}$를 축소 연접
$\mathcal{R}(X)$-대수라 하자~; 그러면 $\mathcal{A}$에 관한 $X$의
적분폐포 $X'$ (\hyperref[II.6.3.4-ko]{6.3.4})은 정규 대수곡선이고,
표준 사상 $X'\to X$는 유한이다.
\end{corollary}

$X'\to X$가 유한이라는 사실은 (\hyperref[II.6.3.10-ko]{6.3.10})에서
따른다~; 따라서 $X'$은 $k$-대수적 스킴이다~; 더 나아가 $x_i$
$(1\leq i\leq n)$를 $X$의 기약 성분들의 일반점, $x'_j$
$(1\leq j\leq m)$를 $X'$의 기약 성분들의 일반점이라 하면, 각
$k(x'_j)$는 어느 한 $k(x_i)$의 유한 대수적 확대이므로
(\hyperref[II.6.3.6-ko]{6.3.6}), $k$ 위의 초월차수가 1이다. 따라서
$X'$은 실제로 $k$ 위의 대수곡선이며, 더욱이 $X'$이 유한 개의 정역
정규 스킴들의 합임을 알고 있다
(\hyperref[II.6.3.6-ko]{6.3.6} 및
\hyperref[II.6.3.7-ko]{6.3.7}).

\begin{env}[7.4.7]
\label{II.7.4.7-ko}
$k$ 위의 대수곡선 $X$가 $k$ 위에서 \emph{고유}이면 $X$를
\emph{완비}라고 한다.
\end{env}

\begin{corollary}[7.4.8]
\label{II.7.4.8-ko}
$k$ 위의 축소 대수곡선 $X$가 완비이기 위한 필요충분조건은 그 정규화
$X'$이 완비인 것이다.
\end{corollary}

\oldpage[II]{150}
실제로 표준 사상 $f:X'\to X$는 이때 유한이고
(\hyperref[II.7.4.6-ko]{7.4.6}), 따라서 고유이며
(\hyperref[II.6.1.11-ko]{6.1.11}) 전사이다
(\hyperref[II.6.3.8-ko]{6.3.8})~; $g:X\to\operatorname{Spec}(k)$를
구조 사상이라 하면, $g$가 가정에 따라 분리되어 있으므로
(\hyperref[II.5.4.2-ko]{5.4.2, (iii)})와
(\hyperref[II.5.4.3-ko]{5.4.3, (ii)})에 따라 $g$와 $g\circ f$는
동시에 고유이다.

\begin{proposition}[7.4.9]
\label{II.7.4.9-ko}
$X$를 $k$ 위의 정규 대수곡선, $Y$를 $k$ 위에서 고유인 $k$-대수적
스킴이라 하자. 그러면 $X$에서 $Y$로 가는 모든 $k$-유리 사상은 모든
곳에서 정의되며, 다시 말해 사상이다.
\end{proposition}

실제로 (\hyperref[II.7.3.7-ko]{7.3.7})에 따라 그러한 사상이 정의되지
않는 점 $x\in X$에서는 $\mathcal{O}_x$의 차원이 $\geq2$이어야 하므로,
그러한 점들의 집합은 공집합이다~; 마지막 주장은
(\hyperref[I.7.2.3-ko]{I, 7.2.3})에서 따른다.

\begin{corollary}[7.4.10]
\label{II.7.4.10-ko}
$k$ 위의 정규 대수곡선은 $k$ 위에서 준사영적이다.
\end{corollary}

$X$는 유한 개의 정역 정규 대수곡선들의 합이므로
(\hyperref[II.6.3.8-ko]{6.3.8}), $X$가 정역인 경우로 제한할 수 있다
(\hyperref[II.5.3.6-ko]{5.3.6}). $X$는 준콤팩트이므로 유한 개의 아핀
열린집합 $U_i$ $(1\leq i\leq n)$로 덮이고, 각 $U_i$는 $k$ 위에서
유한형이므로 어떤 정수 $n_i$와 $k$-몰입 사상
$f_i:U_i\to\mathbf{P}_k^{n_i}$가 존재한다
(\hyperref[II.5.3.3-ko]{5.3.3} 및
\hyperref[II.5.3.4-ko]{5.3.4, (i)}). $U_i$가 $X$에서 조밀하므로
(\hyperref[II.7.4.9-ko]{7.4.9})에 따라 $f_i$는 $k$-사상
$g_i:X\to\mathbf{P}_k^{n_i}$로 연장된다. 이에 따라 $X$에서
$\mathbf{P}_k^{n_i}$들의 $k$ 위의 곱 $P$로 가는 $k$-사상
$g=(g_1,\ldots,g_n)_k$를 얻는다. 또한 각 지표 $i$에 대하여 $g_i$의
$U_i$에 대한 제한이 몰입 사상이므로, $g$의 $U_i$에 대한 제한도
몰입 사상이다 (\hyperref[I.5.3.14-ko]{I, 5.3.14}). $U_i$들이 $X$를
덮고 $g$가 분리 사상이므로 (\hyperref[I.5.5.1-ko]{I, 5.5.1, (v)}),
$g$는 $X$에서 $P$로 가는 몰입 사상이다
(\hyperref[I.8.2.8-ko]{I, 8.2.8}). 세그레 사상
(\hyperref[II.4.3.3-ko]{4.3.3})이 $P$에서 어떤
$\mathbf{P}_k^N$으로 가는 몰입 사상을 주므로, 이로써 $X$가
준사영적임이 증명된다.

\begin{corollary}[7.4.11]
\label{II.7.4.11-ko}
정규 대수곡선 $X$는 정규 완비 대수곡선 $\widehat{X}$의 조밀한
열린집합 위에 유도된 스킴과 동형이며, $\widehat{X}$은 유일한
동형사상까지 유일하게 결정된다.
\end{corollary}

$X_1$, $X_2$를 두 정규 완비 곡선이라 하면,
(\hyperref[II.7.4.9-ko]{7.4.9})에서 $X_1$의 조밀한 열린집합 $U_1$에서
$X_2$의 조밀한 열린집합 $U_2$ 위로 가는 모든 동형사상이 $X_1$에서
$X_2$ 위로 가는 동형사상으로 유일하게 연장됨이 곧바로 따른다~;
이로써 유일성 주장이 성립한다. $\widehat{X}$의 존재를 증명하기 위해서는
$X$를 사영공간 $\mathbf{P}_k^n$의 부분스킴으로 볼 수 있음에 주의하면
충분하다 (\hyperref[II.7.4.10-ko]{7.4.10}). $\overline{X}$를
$\mathbf{P}_k^n$에서 $X$의 폐포라 하자
(\hyperref[I.9.5.11-ko]{I, 9.5.11})~; $X$는 $\overline{X}$의 조밀한
열린집합 위에 유도되므로 (\hyperref[I.9.5.10-ko]{I, 9.5.10}), $X$의
기약 성분들의 일반점 $x_i$들은 $\overline{X}$의 기약 성분들의
일반점이기도 하고, 두 스킴에서 $k(x_i)$들은 서로 같다. 따라서
(\hyperref[II.7.4.1-ko]{7.4.1}) $\overline{X}$는 $k$ 위의
대수곡선이며, 축소이고 (\hyperref[I.9.5.9-ko]{I, 9.5.9}) $k$ 위에서
사영적이므로 (\hyperref[II.5.5.1-ko]{5.5.1}) 완비이다
(\hyperref[II.5.5.3-ko]{5.5.3}). 이제 $\widehat{X}$을
$\overline{X}$의 \emph{정규화}로 잡으면, 이는 역시 완비이다
(\hyperref[II.7.4.8-ko]{7.4.8})~; 더 나아가
$h:\widehat{X}\to\overline{X}$가 표준 사상이면, $X$가 정규이므로
$h^{-1}(X)$에 대한 $h$의 제한은 $X$ 위로 가는 동형사상이다
(\hyperref[II.6.3.4-ko]{6.3.4}). 그리고 $h^{-1}(X)$는
$\widehat{X}$의 기약 성분들의 일반점들을 포함하므로
(\hyperref[II.6.3.8-ko]{6.3.8}) $\widehat{X}$에서 조밀하다. 이로써
증명이 끝난다.

\oldpage[II]{151}

\begin{remark}[7.4.12]
\label{II.7.4.12-ko}
제V장에서 우리는 곡선이 정규라는 가정(더 나아가 축소라는 가정) 없이도
(\hyperref[II.7.4.10-ko]{7.4.10})의 결론이 여전히 성립함을 보일
것이다~; 또한 대수곡선(축소이든 아니든)이 아핀이기 위한 필요충분조건은
그 (축소) 기약 성분들이 완비가 아닌 것임을 보일 것이다.
\end{remark}

\begin{corollary}[7.4.13]
\label{II.7.4.13-ko}
$X$를 체 $R(X)=K$인 기약 정규 곡선, $Y$를 체 $L=R(Y)$인 정역 완비
곡선이라 하자. 우세 $k$-사상 $X\to Y$와 $k$-단사 준동형 $L\to K$
사이에는 표준 일대일 대응이 있다.
\end{corollary}

(\hyperref[II.7.4.9-ko]{7.4.9})에 따라 $X$에서 $Y$로 가는
$k$-유리 사상들은 $k$-사상 $u:X\to Y$와 동일시된다. 우세 사상 $u$는
$u(x)=y$라는 사실로 특징지어지므로(여기서 $x$와 $y$는 각각 $X$와
$Y$의 일반점이다), 따름정리는 이들 주의와
(\hyperref[I.7.1.13-ko]{I, 7.1.13})에서 따른다.

\begin{env}[7.4.14]
\label{II.7.4.14-ko}
(\hyperref[II.7.4.13-ko]{7.4.13})의 결과는 $Y$를 \emph{사영직선}
$\mathbf{P}_k^1=\operatorname{Proj}(k[T_0,T])$으로 취할 때 더 구체화할 수
있다. 여기서 $T_0$와 $T$는 두 부정원이다. 이는 실제로 정역 스킴이고
(\hyperref[II.2.4.4-ko]{2.4.4}), $Y$의 열린집합 $D_+(T_0)$ 위에 유도된
스킴은 $\operatorname{Spec}(k[T])$과 동형이다
(\hyperref[II.2.3.6-ko]{2.3.6}). 따라서 $Y$의 일반점은 $k[T]$의
아이디얼 $(0)$이고 그 유리함수체는 $k(T)$이므로, $Y$는 $k$ 위의 완전
대수곡선이다. 또한 $S=k[T_0,T]$에서 $T_0$를 포함하고 $S_+$와 다른
등급 소 아이디얼은 주 아이디얼 $(T_0)$뿐이다. 따라서
$Y=\mathbf{P}_k^1$에서 $D_+(T_0)$의 여집합은 \emph{하나의 닫힌점}으로만
이루어지며, 이를 ``무한원점''이라 하고 $\infty$로 나타내겠다(벡터다발과
사영다발 사이의 관계에 관한 일반적인 연구는
\hyperref[subsection:II.8.4-ko]{8.4}를 보라). 이 표기 아래 다음이
성립한다~:
\end{env}

\begin{corollary}[7.4.15]
\label{II.7.4.15-ko}
$X$를 유리함수체가 $R(X)=K$인 정규 기약 곡선이라 하자. $K$에서
$X$로부터 $\mathbf{P}_k^1$로 가는 사상 $u$ 가운데 값이 $\infty$인
상수 사상과 다른 것들의 집합으로 가는 표준적인 전단사 사상이 존재한다.
$u$가 우세하기 위한 필요충분조건은 $K$의 대응하는 원소가 $k$ 위에서
초월적인 것이다.
\end{corollary}

이 명제는 (\hyperref[II.7.4.9-ko]{7.4.9})와 다음 보조정리에서 곧바로
따른다.

\begin{lemma}[7.4.15.1]
\label{II.7.4.15.1-ko}
$X$를 $k$ 위의 정역 준스킴이라 하고, $K=R(X)$를 그 유리함수체라 하자.
$K$에서 $X$로부터 $\mathbf{P}_k^1$로 가는 유리사상 $u$ 가운데 값이
$\infty$인 상수 사상과 다른 것들의 집합으로 가는 표준적인 전단사 사상이
존재한다. 그러한 유리사상이 우세하기 위한 필요충분조건은 $K$의 대응하는
원소가 $k$ 위에서 초월적인 것이다.
\end{lemma}

먼저 $X$에서 $\mathbf{P}_k^1$로 가는 유리사상들은 $k$의 확대체 $K$에
값을 갖는 $\mathbf{P}_k^1$의 점들과 일대일로 대응한다
(\hyperref[I.7.1.12-ko]{I, 7.1.12}). 그러한 점이 $\mathbf{P}_k^1$의
일반점에 위치하면
(\hyperref[I.3.4.5-ko]{I, 3.4.5}), 대응하는 유리사상은 명백히
우세하다. 그렇지 않으면 $\mathbf{P}_k^1$의 일반점과 다른 모든 점은
닫혀 있으므로 (\hyperref[II.7.4.3-ko]{7.4.3}), $u$의 정의역 $U$에서
그 동치류 $u$의 유일한 사상 $U\to\mathbf{P}_k^1$로 얻는 상은
(\hyperref[I.7.2.2-ko]{I, 7.2.2}) $\mathbf{P}_k^1$의 한 닫힌점 $y$로만
이루어진다. 따라서 이 사상은 ($X$의 모든 곳에서 정의될 필요는 없으며)
우세하지 않다~; 용어를 남용하여 이때 유리사상 $u$를 ``값이 $y$인
상수''라고 한다. 이제 소재점이
(\hyperref[I.3.4.5-ko]{I, 3.4.5}) $\infty$와 다른, $K$에 값을 갖는
$\mathbf{P}_k^1$의 점들과 $K$의 원소들을 일대일로 대응시키고, 그러한
점의 소재점이 $\mathbf{P}_k^1$의 일반점인 것은 그 점에 대응하는 원소가

\oldpage[II]{152}

$k$ 위에서 초월적일 필요충분조건임을 확인하기만 하면 된다. 그런데 이 확인은
(\hyperref[II.4.2.6-ko]{4.2.6, 예 1\textsuperscript{o}})에서
즉시 이루어진다.

\begin{corollary}[7.4.16]
\label{II.7.4.16-ko}
$X$, $Y$를 $k$ 위의 정규이고 완전하며 기약인 두 대수곡선이라 하고,
$K=R(X)$, $L=R(Y)$을 각각 그 유리함수체라 하자. $k$-동형사상
$X\xrightarrow{\sim}Y$들의 집합과 $k$-동형사상
$L\xrightarrow{\sim}K$들의 집합 사이에는 표준적인 일대일 대응이
존재한다.
\end{corollary}

이는 (\hyperref[II.7.4.13-ko]{7.4.13})의 자명한 귀결이다.

\begin{env}[7.4.17]
\label{II.7.4.17-ko}
따름정리 (\hyperref[II.7.4.16-ko]{7.4.16})에 따르면 $k$ 위의 정규이고
완전하며 기약인 대수곡선은 \emph{유일한 동형사상까지 그 유리함수체
$K$에 의해 결정된다}~; 정의에 따라 $K$는 $k$의 유한 생성 확대체이고
초월차수가 $1$이다(이를 고전적으로 \emph{한 변수의 대수함수체}라고
한다). 더욱이 다음이 성립한다~:
\end{env}

\begin{proposition}[7.4.18]
\label{II.7.4.18-ko}
$k$의 유한 생성 확대체 $K$로서 초월차수가 $1$인 것마다, $R(X)=K$인
정규이고 완전하며 기약인 대수곡선 $X$가 존재한다(유일한 동형사상까지
결정된다). $X$의 국소환들의 집합은 (\hyperref[I.8.2.1-ko]{I, 8.2.1})
$K$ 및 $k$를 포함하고 $K$를 분수체로 갖는 값매김환들로 이루어진
집합과 동일시된다.
\end{proposition}

실제로 $K$는 $k$의 순수 초월 확대체 $k(T)$의 유한차수 확대체이고,
앞에서 본 바와 같이 $k(T)$는 사영직선 $Y=\mathbf{P}_k^1$의
유리함수체와 동일시된다. $X$를 $K$에 관한 $Y$의 \emph{정수적 폐포}라
하자 (\hyperref[II.6.3.4-ko]{6.3.4})~; $X$는 유리함수체가 $K$인 정규
대수곡선이고 (\hyperref[II.6.3.7-ko]{6.3.7}), 사상 $X\to Y$가
유한이므로 완전하다 (\hyperref[II.7.4.6-ko]{7.4.6}). $X$의 국소환
$\mathcal{O}_x$들은 다음과 같다~: $x$가 일반점이면 체 $K$이고, $x$가
일반점과 다르면 $\mathcal{O}_x$는 $k$를 포함하고 $K$를 분수체로 갖는
이산 값매김환이다 (\hyperref[II.7.4.5-ko]{7.4.5}). 역으로 $A$를 그러한
값매김환이라 하자~; 사상 $X\to\operatorname{Spec}(k)$가 고유이고
$A$가 $k$를 지배하므로, $A$는 $X$의 어떤 국소환 $\mathcal{O}_x$를
지배한다 (\hyperref[II.7.3.10-ko]{7.3.10})~; 후자도 $K$를 분수체로 갖는
값매김환이므로 반드시 $A$와 같다.

\begin{namedtheorem}[7.4.19]{비고}
\label{II.7.4.19-ko}
(\hyperref[II.7.4.16-ko]{7.4.16})과
(\hyperref[II.7.4.18-ko]{7.4.18})에서 \emph{$k$ 위의 정규이고 완전하며
기약인 대수곡선을 주는 것은 본질적으로 $k$의 유한 생성 확대체이면서
초월차수가 $1$인 $K$를 주는 것과 동치이다}라는 사실이 따른다. $k'$이
밑체 $k$의 확대체이면 $X\otimes_k k'$은 여전히 $k'$ 위의 완전
대수곡선이지만 (\hyperref[II.5.4.2-ko]{5.4.2, (iii)}), 일반적으로는
축소이지도 기약이지도 않다는 점에 유의하자. 그러나 $K$가 $k$의
분리가능 확대체이고 $k$가 $K$ 안에서 대수적으로 닫혀 있으면 축소이고
기약이다(우리가 따르지 않을 고전적인 용어로는 이를 $K$가 $k$의
``정칙 확대체''라고 표현한다). 하지만 이 경우에도
$X\otimes_k k'$이 정규가 아닐 수 있다. 이 문제들에 관한 자세한 내용은
제IV장에서 찾을 수 있다.
\end{namedtheorem}

\section{블로업 스킴; 사영 원뿔; 사영 폐포}
\label{section:II.8-ko}

\subsection{블로업 준스킴}
\label{subsection:II.8.1-ko}

\begin{env}[8.1.1]
\label{II.8.1.1-ko}
$Y$를 준스킴이라 하고, 각 정수 $n\geq0$에 대하여 $\mathcal{I}_n$을
$\mathcal{O}_Y$의 준연접 아이디얼이라 하자. 다음 조건들이 성립한다고
가정한다~:
\begin{equation}
\label{II.8.1.1.1-ko}
  \mathcal{I}_0=\mathcal{O}_Y,\qquad
  \mathcal{I}_n\subset\mathcal{I}_m\quad\text{단 }m\leq n
\tag{8.1.1.1}
\end{equation}
\begin{equation}
\label{II.8.1.1.2-ko}
  \mathcal{I}_m\mathcal{I}_n\subset\mathcal{I}_{m+n}
  \quad\text{임의의 }m,n\text{에 대하여}.
\tag{8.1.1.2}
\end{equation}

\oldpage[II]{153}
이 가정들로부터 다음이 따름을 주목하자.
\begin{equation}
\label{II.8.1.1.3-ko}
  \mathcal{I}_1^n\subset\mathcal{I}_n.
\tag{8.1.1.3}
\end{equation}

다음과 같이 놓자.
\begin{equation}
\label{II.8.1.1.4-ko}
  \mathcal{S}=\bigoplus_{n\geq0}\mathcal{I}_n.
\tag{8.1.1.4}
\end{equation}

(\hyperref[II.8.1.1.1-ko]{8.1.1.1})과
(\hyperref[II.8.1.1.2-ko]{8.1.1.2})로부터 $\mathcal{S}$는 준연접 등급
$\mathcal{O}_Y$-대수이고, 따라서 $Y$-스킴
$X=\operatorname{Proj}(\mathcal{S})$를 정의한다. $\mathcal{J}$가
$\mathcal{O}_Y$의 \emph{가역} 아이디얼이면
$\mathcal{I}_n\otimes_{\mathcal{O}_Y}\mathcal{J}^{\otimes n}$은
$\mathcal{I}_n\mathcal{J}^n$과 표준적으로 동일시된다. 따라서
$\mathcal{I}_n$들을 $\mathcal{I}_n\mathcal{J}^n$들로 바꾸어
$\mathcal{S}$를 준연접 $\mathcal{O}_Y$-대수
$\mathcal{S}_{(\mathcal{J})}$로 바꾸면,
$X_{(\mathcal{J})}=\operatorname{Proj}(\mathcal{S}_{(\mathcal{J})})$는
$X$와 표준적으로 동형이다 (\hyperref[II.3.1.8-ko]{3.1.8}).
\end{env}

\begin{env}[8.1.2]
\label{II.8.1.2-ko}
$Y$가 \emph{국소 정역 준스킴}이라고 가정하자. 그러면 유리함수층
$\mathcal{R}(Y)$는 준연접 $\mathcal{O}_Y$-대수이다
(\hyperref[I.7.3.7-ko]{I, 7.3.7}). $\mathcal{R}(Y)$의
$\mathcal{O}_Y$-부분가군 $\mathcal{I}$가 \emph{유한형}이면 이를
$\mathcal{R}(Y)$의 \emph{분수 아이디얼}이라고 한다
(\hyperref[0.5.2.1-ko]{0, 5.2.1}). 각 $n\geq0$에 대하여
$\mathcal{R}(Y)$의 준연접 분수 아이디얼 $\mathcal{I}_n$이 주어지고,
$\mathcal{I}_0=\mathcal{O}_Y$이며 조건
(\hyperref[II.8.1.1.2-ko]{8.1.1.2})가 성립한다고 하자(그러나
(\hyperref[II.8.1.1.1-ko]{8.1.1.1})의 두 번째 조건은 반드시 성립할
필요가 없다). 그러면 여전히 식 (\hyperref[II.8.1.1.4-ko]{8.1.1.4})로
준연접 등급 $\mathcal{O}_Y$-대수와 이에 대응하는 $Y$-스킴
$X=\operatorname{Proj}(\mathcal{S})$를 정의할 수 있다. 또한
$\mathcal{R}(Y)$의 임의의 가역 분수 아이디얼 $\mathcal{J}$에 대하여
$X$에서 $X_{(\mathcal{J})}$로 가는 표준 동형사상이 여전히 존재한다.
\end{env}

\begin{definition}[8.1.3]
\label{II.8.1.3-ko}
$Y$를 준스킴(각각 국소 정역 준스킴)이라 하고, $\mathcal{I}$를
$\mathcal{O}_Y$의 준연접 아이디얼(각각 $\mathcal{R}(Y)$의 준연접 분수
아이디얼)이라 하자. $Y$-스킴
$X=\operatorname{Proj}(\bigoplus_{n\geq0}\mathcal{I}^n)$을
\emph{아이디얼 $\mathcal{I}$을 따라 블로업하여 얻었다}고 하며, 또는
이를 \emph{$\mathcal{I}$에 대한 $Y$의 블로업 준스킴}이라고 한다.
$\mathcal{I}$가 $\mathcal{O}_Y$의 준연접 아이디얼이고 $Y'$이
$\mathcal{I}$로 정의되는 $Y$의 닫힌 부분준스킴일 때에는, $X$를
$Y'$을 블로업하여 얻은 $Y$-스킴이라고도 한다.
\end{definition}

정의에 따라 이때
$\mathcal{S}=\bigoplus_{n\geq0}\mathcal{I}^n$은
$\mathcal{S}_1=\mathcal{I}$로 생성된다. $\mathcal{I}$가 
$\mathcal{O}_Y$-가군으로서 \emph{유한형}이면 $X$는 $Y$ 위에서
\emph{사영적}이다 (\hyperref[II.5.5.2-ko]{5.5.2}). $\mathcal{I}$에
아무 가정도 하지 않더라도 $\mathcal{O}_X$-가군 $\mathcal{O}_X(1)$은
\emph{가역}이고 (\hyperref[II.3.2.5-ko]{3.2.5}), 구조 사상
$X\to Y$에 대하여 (\hyperref[II.4.4.3-ko]{4.4.3})에 의해
\emph{매우 풍부}하다.

$f:X\to Y$가 구조 사상이면, $\mathcal{I}$가 $\mathcal{O}_Y$의
아이디얼이고 $Y'$이 그것으로 정의되는 닫힌 부분준스킴일 때
$f^{-1}(Y-Y')$에 대한 $f$의 제한은 $Y-Y'$ 위의
\emph{동형사상}임을 주목하자. 실제로 이 문제는 $Y$에서 국소적이므로
$\mathcal{I}=\mathcal{O}_Y$라고 가정하면 충분하고, 주장은
(\hyperref[II.3.1.7-ko]{3.1.7})에서 따른다.

$\mathcal{I}$를 $\mathcal{I}^d$ ($d>0$)로 바꾸면 블로업 $Y$-스킴
$X$는 표준적으로 동형인 $Y$-스킴 $X'$으로 바뀐다
(\hyperref[II.8.1.1-ko]{8.1.1}). 마찬가지로 임의의 가역 아이디얼
(각각 가역 분수 아이디얼) $\mathcal{J}$에 대하여 아이디얼
$\mathcal{I}\mathcal{J}$에 대한 블로업 준스킴
$X_{(\mathcal{J})}$는 $X$와 표준적으로 동형이다
(\hyperref[II.8.1.1-ko]{8.1.1}).

특히 $\mathcal{I}$가 \emph{가역} 아이디얼(각각 분수 아이디얼)이면,
$\mathcal{I}$을 따라 블로업하여 얻은 $Y$-스킴은 $Y$와
\emph{동형}이다 (\hyperref[II.3.1.7-ko]{3.1.7}).

\begin{proposition}[8.1.4]
\label{II.8.1.4-ko}
$Y$를 정역 준스킴이라 하자.
\begin{enumerate}
\item[\rm(i)] $\mathcal{R}(Y)$의 임의의 준연접 분수 아이디얼열
$(\mathcal{I}_n)$이 (\hyperref[II.8.1.1.2-ko]{8.1.1.2})를
만족한다고 하자.

\oldpage[II]{154}
더 나아가 $\mathcal{I}_0=\mathcal{O}_Y$라고 하자. 그러면 $Y$-스킴
$X=\operatorname{Proj}(\bigoplus_{n\geq0}\mathcal{I}_n)$은
정역이고 구조 사상 $f:X\to Y$는 우세 사상이다.
\item[\rm(ii)] $\mathcal{I}$를 $\mathcal{R}(Y)$의 준연접 분수
아이디얼이라 하고, $X$를 $\mathcal{I}$에 대한 $Y$의 블로업
$Y$-스킴이라 하자. $\mathcal{I}\neq0$이면 구조 사상 $f:X\to Y$는
쌍유리이며 전사이다.
\end{enumerate}
\end{proposition}

\begin{enumerate}
\item[\rm(i)] 모든 $y\in Y$에 대하여 $\mathcal{O}_y$가 정역이므로
(\hyperref[I.5.1.4-ko]{I, 5.1.4}),
$\mathcal{S}=\bigoplus_{n\geq0}\mathcal{I}_n$이 \emph{정역}
$\mathcal{O}_Y$-대수라는 사실
(\hyperref[II.3.1.12-ko]{3.1.12} 및
\hyperref[II.3.1.14-ko]{3.1.14})에서 따른다.
\item[\rm(ii)] (i)에 따라 $X$는 정역이다. 더 나아가 $x$와 $y$가
각각 $X$와 $Y$의 일반점이면 $f(x)=y$이고, $k(x)$가 $k(y)$ 위에서
랭크 $1$임을 증명해야 한다. 그런데 $x$는 올 $f^{-1}(y)$의
일반점이기도 하다. $\psi$가 표준 사상 $Z\to Y$이고
$Z=\operatorname{Spec}(k(y))$이면, 준스킴 $f^{-1}(y)$은
$\operatorname{Proj}(\mathcal{S}')$과 동일시된다. 여기서
$\mathcal{S}'=\psi^*(\mathcal{S})$이다
(\hyperref[II.3.5.3-ko]{3.5.3}). 한편
$\mathcal{S}'=\bigoplus_{n\geq0}(\mathcal{I}_y)^n$임은 명백하고,
$\mathcal{I}$가 $\mathcal{R}(Y)$의 영이 아닌 준연접 분수
아이디얼이므로 $\mathcal{I}_y\neq0$
(\hyperref[I.7.3.6-ko]{I, 7.3.6})이다. 따라서
$\mathcal{I}_y=k(y)$이고, 이에 따라
$\operatorname{Proj}(\mathcal{S}')$은
$\operatorname{Spec}(k(y))$과 동일시된다
(\hyperref[II.3.1.7-ko]{3.1.7}). 이로부터 결론이 따른다.
\end{enumerate}

(\hyperref[II.8.1.4-ko]{8.1.4})의 \emph{역}을
(\agref{III.2.3.8-ko}{III, 2.3.8})에서 증명할 것이다.

\begin{env}[8.1.5]
\label{II.8.1.5-ko}
(\hyperref[II.8.1.1-ko]{8.1.1})의 상황과 표기로 돌아가자. 정의에 따라
단사 준동형 $\mathcal{I}_{n+1}\to\mathcal{I}_n$
(\hyperref[II.8.1.1.1-ko]{8.1.1.1})은 각 $k\in\mathbf{Z}$에 대하여
차수가 $0$인 다음 등급 $\mathcal{S}$-가군의 단사 준동형을 정의한다.
\begin{equation}
\label{II.8.1.5.1-ko}
  u_k:\mathcal{S}_+(k+1)\to\mathcal{S}(k)~;
\tag{8.1.5.1}
\end{equation}
$\mathcal{S}_+(k+1)$과 $\mathcal{S}(k+1)$은 표준적으로
(TN)-동형이므로, $u_k$에는 다음 단사 $\mathcal{O}_X$-가군 준동형이
표준적으로 대응한다 (\hyperref[II.3.4.2-ko]{3.4.2})~:
\begin{equation}
\label{II.8.1.5.2-ko}
  \widetilde{u}_k:\mathcal{O}_X(k+1)\to\mathcal{O}_X(k).
\tag{8.1.5.2}
\end{equation}

한편 (\hyperref[II.3.2.6-ko]{3.2.6})에서 표준 준동형들을
정의했음을 상기하자.
\begin{equation}
\label{II.8.1.5.3-ko}
  \lambda:\mathcal{O}_X(h)\otimes_{\mathcal{O}_X}\mathcal{O}_X(k)
  \to\mathcal{O}_X(h+k)
\tag{8.1.5.3}
\end{equation}
그리고 다음 도식은 가환이다.
\[
\xymatrix{
  \mathcal{S}(h)\otimes_{\mathcal{S}}\mathcal{S}(k)
    \otimes_{\mathcal{S}}\mathcal{S}(l)\ar[r]\ar[d]
  & \mathcal{S}(h+k)\otimes_{\mathcal{S}}\mathcal{S}(l)\ar[d]\\
  \mathcal{S}(h)\otimes_{\mathcal{S}}\mathcal{S}(k+l)\ar[r]
  & \mathcal{S}(h+k+l)
}
\]
따라서 $\lambda$의 함자성(\hyperref[II.3.2.6-ko]{3.2.6})에 의해
준동형들 (\hyperref[II.8.1.5.3-ko]{8.1.5.3})은
\begin{equation}
\label{II.8.1.5.4-ko}
  \mathcal{S}_X=\bigoplus_{n\in\mathbf{Z}}\mathcal{O}_X(n)
\tag{8.1.5.4}
\end{equation}

\oldpage[II]{155}
위에 준연접 등급 $\mathcal{O}_X$-대수 구조를 정의한다. 또한 다음 도식은
가환이다.
\[
\xymatrix{
  \mathcal{S}(h)\otimes_{\mathcal{S}}\mathcal{S}_+(k+1)\ar[r]\ar[d]_{1\otimes u_k}
  & \mathcal{S}_+(h+k+1)\ar[d]^{u_{k+h}}\\
  \mathcal{S}(h)\otimes_{\mathcal{S}}\mathcal{S}(k)\ar[r]
  & \mathcal{S}(h+k)
}
\]
$\lambda$의 함자성에 의해 다음 가환 도식을 얻는다.
\begin{equation}
\label{II.8.1.5.5-ko}
\xymatrix{
  \mathcal{O}_X(h)\otimes_{\mathcal{O}_X}\mathcal{O}_X(k+1)
    \ar[r]^{\lambda}\ar[d]_{1\otimes\widetilde{u}_k}
  & \mathcal{O}_X(h+k+1)\ar[d]^{\widetilde{u}_{k+h}}\\
  \mathcal{O}_X(h)\otimes_{\mathcal{O}_X}\mathcal{O}_X(k)
    \ar[r]^{\lambda}
  & \mathcal{O}_X(h+k)
}
\tag{8.1.5.5}
\end{equation}
여기서 수평 화살표들은 표준 준동형들이다. 따라서 $\widetilde{u}_k$들은
\emph{단사 준동형}, 즉 (차수 $0$인) \emph{등급 $\mathcal{S}_X$-가군의 준동형}
\begin{equation}
\label{II.8.1.5.6-ko}
  \widetilde{u}:\mathcal{S}_X(1)\to\mathcal{S}_X.
\tag{8.1.5.6}
\end{equation}
을 정의한다고 말할 수 있다.
\end{env}

\begin{env}[8.1.6]
\label{II.8.1.6-ko}
(\hyperref[II.8.1.5-ko]{8.1.5})의 표기를 유지하자. $n\geq0$일 때 합성 준동형
$\widetilde{v}_n=\widetilde{u}_{n-1}\widetilde{u}_{n-2}\cdots\widetilde{u}_0$
은 \emph{단사} 준동형 $\mathcal{O}_X(n)\to\mathcal{O}_X$이다. 그 상을
$\mathcal{I}_{n,X}$로 나타내자. 이는 $\mathcal{O}_X$의 준연접 아이디얼층이며
$\mathcal{O}_X(n)$과 \emph{동형}이다. 또한 다음 도식은
\[
\xymatrix{
  \mathcal{O}_X(m)\otimes_{\mathcal{O}_X}\mathcal{O}_X(n)
    \ar[r]^{\lambda}\ar[d]_{\widetilde{v}_m\otimes\widetilde{v}_n}
  & \mathcal{O}_X(m+n)\ar[d]^{\widetilde{v}_{m+n}}\\
  \mathcal{O}_X\ar[r]_{\mathrm{id}}
  & \mathcal{O}_X
}
\]
$m\geq0$, $n\geq0$일 때 가환이다. 이로부터 다음 포함관계들을 얻는다.
\begin{equation}
\label{II.8.1.6.1-ko}
  \mathcal{I}_{0,X}=\mathcal{O}_X,\qquad
  \mathcal{I}_{n,X}\subset\mathcal{I}_{m,X}
  \quad\text{단, }0\leq m\leq n
\tag{8.1.6.1}
\end{equation}
\begin{equation}
\label{II.8.1.6.2-ko}
  \mathcal{I}_{m,X}\mathcal{I}_{n,X}\subset\mathcal{I}_{m+n,X}
  \quad\text{단, }m\geq0,\ n\geq0.
\tag{8.1.6.2}
\end{equation}
\end{env}

\oldpage[II]{156}

\begin{proposition}[8.1.7]
\label{II.8.1.7-ko}
$Y$를 준스킴, $\mathcal{I}$를 $\mathcal{O}_Y$의 준연접 아이디얼층이라 하고,
$X=\operatorname{Proj}(\bigoplus_{n\geq0}\mathcal{I}^n)$를
$\mathcal{I}$를 블로업하여 얻은 $Y$-스킴이라 하자. 그러면 모든
$n>0$에 대하여 표준 동형사상
\begin{equation}
\label{II.8.1.7.1-ko}
  \mathcal{O}_X(n)\xrightarrow{\sim}\mathcal{I}^n\mathcal{O}_X
  =\mathcal{I}_{n,X}
\tag{8.1.7.1}
\end{equation}
이 존재한다((\hyperref[0.4.3.5-ko]{0, 4.3.5}) 참조). 따라서 $n>0$이면
$\mathcal{I}^n\mathcal{O}_X$는 가역이고 매우 풍부한 $\mathcal{O}_X$-가군이다.
\end{proposition}

마지막 단언은 정의상 $\mathcal{O}_X(1)$이 가역이고 $Y$에 관해 매우
풍부하므로 자명하다
(\hyperref[II.3.2.5-ko]{3.2.5} 및 \hyperref[II.4.4.3-ko]{4.4.3}과
\hyperref[II.4.4.9-ko]{4.4.9}). 한편 정의상 $v_n$의 상은
$\mathcal{I}^n\mathcal{S}$와 다르지 않으며,
(\hyperref[II.8.1.7.1-ko]{8.1.7.1})은 함자
$\widetilde{\mathcal{M}}$의 완전성
(\hyperref[II.3.2.4-ko]{3.2.4})과 공식
(\hyperref[II.3.2.4.1-ko]{3.2.4.1})에서 따른다.

\begin{corollary}[8.1.8]
\label{II.8.1.8-ko}
(\hyperref[II.8.1.7-ko]{8.1.7})의 가정 아래에서 $f:X\to Y$를 구조
사상이라 하고 $Y'$을 $\mathcal{I}$로 정의되는 $Y$의 닫힌
부분준스킴이라 하자. 그러면 $X$의 닫힌 부분준스킴
$X'=f^{-1}(Y')$은 $\mathcal{I}\mathcal{O}_X$로 정의되고
(이는 $\mathcal{O}_X(1)$과 표준적으로
동형이다), 따라서 표준 완전열
\begin{equation}
\label{II.8.1.8.1-ko}
  0\to\mathcal{O}_X(1)\to\mathcal{O}_X\to\mathcal{O}_{X'}\to0.
\tag{8.1.8.1}
\end{equation}
을 얻는다.
\end{corollary}

이는 (\hyperref[II.8.1.7.1-ko]{8.1.7.1})과
(\hyperref[I.4.4.5-ko]{I, 4.4.5})에서 따른다.

\begin{env}[8.1.9]
\label{II.8.1.9-ko}
(\hyperref[II.8.1.7-ko]{8.1.7})의 가정 아래에서
$\mathcal{I}_{n,X}$의 구조를 더 구체적으로 기술할 수 있다. 실제로
준동형
\[
  \widetilde{u}_{-1}:\mathcal{O}_X\to\mathcal{O}_X(-1)
\]
은 $X$ 위의 $\mathcal{O}_X(-1)$의 한 단면 $s$와 표준적으로 대응하며,
이를 ($\mathcal{I}$에 대한) \emph{표준 단면}이라 하겠다
(\hyperref[0.5.1.1-ko]{0, 5.1.1}). 한편 도표
(\hyperref[II.8.1.5.5-ko]{8.1.5.5})의 수평 화살표들은 동형사상이다
(\hyperref[II.3.2.7-ko]{3.2.7})~; 이 도표에서 $h$를 $k$로, $k$를
$-1$로 바꾸면
$\widetilde{u}_k=1_k\otimes\widetilde{u}_{-1}$을 얻는다
($1_h$는 $\mathcal{O}_X(h)$의 항등사상을 나타낸다). 다시 말해 모든
$k\in\mathbf{Z}$에 대하여 준동형 $\widetilde{u}_k$는 \emph{표준 단면
$s$에 의한 텐서 곱셈}과 다르지 않다. 따라서 준동형
$\widetilde{u}$ (\hyperref[II.8.1.5.6-ko]{8.1.5.6})도 같은 방식으로
해석된다.

따라서 모든 $n\geq0$에 대하여 준동형
$\widetilde{v}_n:\mathcal{O}_X(n)\to\mathcal{O}_X$은
$s^{\otimes n}$에 의한 텐서 곱셈과 다르지 않으며, 이로부터 다음을
얻는다~:
\end{env}

\begin{corollary}[8.1.10]
\label{II.8.1.10-ko}
(\hyperref[II.8.1.8-ko]{8.1.8})의 표기에서 $X'$의 바탕공간은
$s(x)=0$인 $x\in X$들의 집합이다. 여기서 $s$는
$\mathcal{O}_X(-1)$의 표준 단면을 나타낸다.
\end{corollary}

실제로 $c_x$가 점 $x$에서 줄기 $(\mathcal{O}_X(1))_x$의 생성원이면,
$s_x\otimes c_x$는 점 $x$에서 $\mathcal{I}_{1,X}$의 줄기의 한
생성원과 표준적으로 동일시되며, 따라서
$s_x\notin\mathfrak{m}_x(\mathcal{O}_X(-1))_x$일 때 그리고 그때에만
가역원이다. 다시 말해 이는 $s(x)\neq0$일 때 그리고 그때에만
성립한다.

\begin{proposition}[8.1.11]
\label{II.8.1.11-ko}
$Y$를 정역 준스킴, $\mathcal{I}$을 $\mathcal{R}(Y)$의 준연접 분수
아이디얼, $X$를 $\mathcal{I}$을 블로업하여 얻은 $Y$-스킴이라 하자.
그러면 $\mathcal{I}\mathcal{O}_X$는 가역이고 $Y$에 관해 매우 풍부한
$\mathcal{O}_X$-가군이다.
\end{proposition}

문제는 $Y$ 위에서 국소적이므로
(\hyperref[II.4.4.5-ko]{4.4.5}), $Y=\operatorname{Spec}(A)$이고
$A$가 분수체 $K$를 갖는 정역이며
$\mathcal{I}=\widetilde{\mathfrak{I}}$인 경우로 한정할 수 있다. 여기서
$\mathfrak{I}$은 $K$의 분수 아이디얼이다~; 따라서
$a\neq0$인 $A$의 원소가 존재하여 $a\mathfrak{I}\subset A$를
만족한다. $S=\bigoplus_{n\geq0}\mathfrak{I}^n$이라 놓자~; 사상
$x\mapsto ax$는 $\mathfrak{I}^{n+1}=(S(1))_n$에서
$a\mathfrak{I}^{n+1}=a\mathfrak{I}S_n\subset\mathfrak{I}^n=S_n$으로
가는 $A$-동형사상이고,

\oldpage[II]{157}
따라서 차수 $0$인 등급 $S$-가군의 (TN)-동형사상
$S_+(1)\to a\mathfrak{I}S$를 정의한다. 한편 $x\mapsto a^{-1}x$는 차수
$0$인 등급 $S$-가군의 동형사상
$a\mathfrak{I}S\xrightarrow{\sim}\mathfrak{I}S$이다. 그러므로 이를
합성하면(\hyperref[II.3.2.4-ko]{3.2.4}) $\mathcal{O}_X$-가군의 동형사상
$\mathcal{O}_X(1)\xrightarrow{\sim}\mathcal{I}\mathcal{O}_X$를 얻는다. 또한
$S$는 $S_1=\mathfrak{I}$에 의해 생성되므로 $\mathcal{O}_X(1)$은 가역이고
(\hyperref[II.3.2.5-ko]{3.2.5}) 매우 풍부하다
(\hyperref[II.4.4.3-ko]{4.4.3}과 \hyperref[II.4.4.9-ko]{4.4.9}). 따라서
주장이 증명된다.

\subsection{등급환에서의 국소화에 관한 예비 결과}
\label{subsection:II.8.2-ko}

\begin{env}[8.2.1]
\label{II.8.2.1-ko}
$S$를 등급환이라 하자. 여기서는 아직 차수들이 양수라고 가정하지 않는다.
다음과 같이 놓자.
\begin{equation}
\label{II.8.2.1.1-ko}
  S^{\geq}=\bigoplus_{n\geq0}S_n,
  \qquad
  S^{\leq}=\bigoplus_{n\leq0}S_n,
\tag{8.2.1.1}
\end{equation}
이들은 각각 음이 아닌 차수와 양이 아닌 차수로만 등급이 주어진 $S$의
등급 부분환이다. $f$가 $S$의 차수 $d$(양수이거나 음수)의 동차
원소이면, 분수환 $S_f=S'$에도 여전히 등급환 구조가 주어진다. 즉
$S'_n$($n\in\mathbf{Z}$)을 $x\in S_{n+kd}$($k\geq0$)인 $x/f^k$들의
집합으로 잡는다. 또한 $S_{(f)}=S'_0$로 놓고, $S'^{\geq}$와
$S'^{\leq}$를 각각 $S_f^{\geq}$와 $S_f^{\leq}$로 쓰겠다. $d>0$이면
\begin{equation}
\label{II.8.2.1.2-ko}
  (S^{\geq})_f=S_f
\tag{8.2.1.2}
\end{equation}
이다. 실제로 $x\in S_{n+kd}$이고 $n+kd<0$이면
$x/f^k=xf^h/f^{h+k}$로 쓸 수 있으며, $h$를 충분히 크게, 그리고 $>0$이 되게 잡으면
$n+(h+k)d>0$이기 때문이다. 따라서 정의에 의해
\begin{equation}
\label{II.8.2.1.3-ko}
  (S^{\geq})_{(f)}=(S_f^{\geq})_0=S_{(f)}.
\tag{8.2.1.3}
\end{equation}

$M$이 등급 $S$-가군이면 마찬가지로
\begin{equation}
\label{II.8.2.1.4-ko}
  M^{\geq}=\bigoplus_{n\geq0}M_n,
  \qquad
  M^{\leq}=\bigoplus_{n\leq0}M_n
\tag{8.2.1.4}
\end{equation}
로 놓는다. 이들은 각각 등급 $S^{\geq}$-가군과 등급
$S^{\leq}$-가군이고, 그 교집합은 $S_0$-가군 $M_0$이다. $f\in S_d$이면
$M_f$에도 등급 $S_f$-가군 구조를 정의한다. 즉 차수 $n$인 원소를
$z\in M_{n+kd}$($k\geq0$)인 $z/f^k$들로 잡는다. $M_f$의 차수 $0$인
원소들의 집합을 $M_{(f)}$로 나타내기로 하며, 이는 $S_{(f)}$-가군이다.
또한 $(M_f)^{\geq}$와 $(M_f)^{\leq}$ 대신 각각 $M_f^{\geq}$와
$M_f^{\leq}$로 쓰겠다. $d>0$이면 위와 마찬가지로
\begin{equation}
\label{II.8.2.1.5-ko}
  (M^{\geq})_f=M_f
\tag{8.2.1.5}
\end{equation}
이고
\begin{equation}
\label{II.8.2.1.6-ko}
  (M^{\geq})_{(f)}=(M_f^{\geq})_0=M_{(f)}.
\tag{8.2.1.6}
\end{equation}
\end{env}

\begin{env}[8.2.2]
\label{II.8.2.2-ko}
$z$를 부정원이라 하고 이를 \emph{동차화 변수}라고 부르자. $S$가
등급환이면(양의 차수이든 음의 차수이든), 다항식 대수\footnote{여기서
환의 분리 완비화를 나타내는 표기 $\widehat{S}$와 혼동할 여지는 없다.}
\begin{equation}
\label{II.8.2.2.1-ko}
  \widehat{S}=S[z]
\tag{8.2.2.1}
\end{equation}

\oldpage[II]{158}
은 등급 $S$-대수이다. 여기서 $f$가 동차이고 $n\geq0$일 때
$fz^n$의 차수를
\begin{equation}
\label{II.8.2.2.2-ko}
  \deg(fz^n)=n+\deg f.
\tag{8.2.2.2}
\end{equation}
로 잡는다.
\end{env}

\begin{lemma}[8.2.3]
\label{II.8.2.3-ko}
\begin{enumerate}
\item[\rm(i)] (등급이 주어지지 않은) 환의 표준 동형사상들이 있다.
\begin{equation}
\label{II.8.2.3.1-ko}
  \widehat{S}_{(z)}\xrightarrow{\sim}
  \widehat{S}/(z-1)\widehat{S}\xrightarrow{\sim}S.
\tag{8.2.3.1}
\end{equation}
\item[\rm(ii)] 모든 $f\in S_d$, $d>0$에 대하여 (등급이 주어지지 않은)
환의 표준 동형사상이 있다.
\begin{equation}
\label{II.8.2.3.2-ko}
  \widehat{S}_{(f)}\xrightarrow{\sim}S_f^{\leq}
\tag{8.2.3.2}
\end{equation}
\end{enumerate}
\end{lemma}

(\hyperref[II.8.2.3.1-ko]{8.2.3.1})의 첫 번째 동형사상은
(\hyperref[II.2.2.5-ko]{2.2.5})에서 정의했으며 두 번째 것은 자명하다.
따라서 이렇게 정의된 동형사상
$\widehat{S}_{(z)}\xrightarrow{\sim}S$는 $\deg(x)=k$($k\geq-n$)인
$xz^n/z^{n+k}$를 원소 $x$에 대응시킨다. 준동형
(\hyperref[II.8.2.3.2-ko]{8.2.3.2})는 $\deg(x)=kd-n$인 $xz^n/f^k$를
$S_f^{\leq}$에서 차수가 $-n$인 원소 $x/f^k$에 대응시키며, 이것이
실제로 동형사상이라는 것도 자명하다.

\begin{env}[8.2.4]
\label{II.8.2.4-ko}
$M$을 등급 $S$-가군이라 하자. $S$-가군
\begin{equation}
\label{II.8.2.4.1-ko}
  \widehat{M}=M\otimes_S\widehat{S}=M\otimes_S S[z]
\tag{8.2.4.1}
\end{equation}
은 $S$-가군 $M\otimes_S Sz^n$들의 직합이고, 따라서 아벨군
$M_k\otimes_S Sz^n$($k\in\mathbf{Z}$, $n\geq0$)들의 직합임이 명백하다.
$\widehat{M}$에 다음과 같이 등급 $\widehat{S}$-가군 구조를 정의한다.
\begin{equation}
\label{II.8.2.4.2-ko}
  \deg(x\otimes z^n)=n+\deg x
\tag{8.2.4.2}
\end{equation}
여기서 $x$는 $M$의 임의의 동차 원소이다. 다음과 같은
(\hyperref[II.8.2.3-ko]{8.2.3})의 대응 명제는 독자에게 증명을 맡긴다.
\end{env}

\begin{lemma}[8.2.5]
\label{II.8.2.5-ko}
\begin{enumerate}
\item[\rm(i)] (등급이 주어지지 않은) 가군의 표준 디-동형사상이 있다.
\begin{equation}
\label{II.8.2.5.1-ko}
  \widehat{M}_{(z)}\xrightarrow{\sim}M.
\tag{8.2.5.1}
\end{equation}
\item[\rm(ii)] 모든 $f\in S_d$($d>0$)에 대하여 (등급이 주어지지 않은)
가군의 디-동형사상이 있다.
\begin{equation}
\label{II.8.2.5.2-ko}
  \widehat{M}_{(f)}\xrightarrow{\sim}M_f^{\leq}.
\tag{8.2.5.2}
\end{equation}
\end{enumerate}
\end{lemma}

\begin{env}[8.2.6]
\label{II.8.2.6-ko}
$S$를 \emph{음이 아닌 차수}로 등급이 주어진 환이라 하고, $S$에서
등급 아이디얼들의 감소열
\begin{equation}
\label{II.8.2.6.1-ko}
  S_{[n]}=\bigoplus_{m\geq n}S_m\qquad(n\geq0)
\tag{8.2.6.1}
\end{equation}
을 생각하자(특히 $S_{[0]}=S$, $S_{[1]}=S_+$이다). 명백히
$S_{[m]}S_{[n]}\subset S_{[m+n]}$이므로, 다음과 같이 잡아
\emph{등급환} $S^\natural$을 정의할 수 있다.
\begin{equation}
\label{II.8.2.6.2-ko}
  S^\natural=\bigoplus_{n\geq0}S_n^\natural
  \quad\text{이고}\quad S_n^\natural=S_{[n]}.
\tag{8.2.6.2}
\end{equation}
따라서 $S_0^\natural$은 환 $S$를 \emph{등급이 주어지지 않은 환}으로
본 것이고, 그 결과 $S^\natural$은 $S_0^\natural$-대수이다. 차수
$d>0$인 모든 동차 원소 $f\in S_d$에 대하여, $S_{[d]}=S_d^\natural$에
속하는 것으로 본 원소 $f$를 $f^\natural$로 나타내겠다. 이 표기 아래
다음이 성립한다.
\end{env}

\oldpage[II]{159}

\begin{lemma}[8.2.7]
\label{II.8.2.7-ko}
$S$를 음이 아닌 차수로 등급이 주어진 환, $f$를 $S_d$의 차수
$d>0$인 동차 원소라 하자. 다음과 같은 표준 환 동형사상들이 있다.
\begin{equation}
\label{II.8.2.7.1-ko}
  S_f\xrightarrow{\sim}\bigoplus_{n\in\mathbf{Z}}S(n)_{(f)}
\tag{8.2.7.1}
\end{equation}
\begin{equation}
\label{II.8.2.7.2-ko}
  (S_f^{\geq})_{f/1}\xrightarrow{\sim}S_f
\tag{8.2.7.2}
\end{equation}
\begin{equation}
\label{II.8.2.7.3-ko}
  S^\natural_{(f^\natural)}\xrightarrow{\sim}S_f^{\geq}
\tag{8.2.7.3}
\end{equation}
이 가운데 처음 두 개는 등급환의 동형사상이다.
\end{lemma}

정의에서 곧바로 $(S_f)_n=(S(n)_f)_0$을 얻으므로 동형사상
(\hyperref[II.8.2.7.1-ko]{8.2.7.1})을 얻으며, 이는 항등사상에 지나지
않는다. 한편 $f/1$은 $S_f$에서 가역이므로, $S_f$에
(\hyperref[II.8.2.1.2-ko]{8.2.1.2})를 적용하면 표준 동형사상
$S_f\xrightarrow{\sim}(S_f^{\geq})_{f/1}=(S_f)_{f/1}$이 있다. 그 역은
정의상 동형사상 (\hyperref[II.8.2.7.2-ko]{8.2.7.2})이다. 끝으로
$x=\sum_{m\geq n}y_m$이 $n=kd$인 $S_{[n]}$의 원소이면, 원소
$x/(f^\natural)^k$를 $S_f^{\geq}$의 원소
$\sum_m y_m/f^k$에 대응시킨다. 이것이 동형사상
(\hyperref[II.8.2.7.3-ko]{8.2.7.3})을 정의함을 곧바로 확인할 수 있다.

\begin{env}[8.2.8]
\label{II.8.2.8-ko}
$M$이 등급 $S$-가군이면, 마찬가지로 모든 $n\in\mathbf{Z}$에 대하여
\begin{equation}
\label{II.8.2.8.1-ko}
  M_{[n]}=\bigoplus_{m\geq n}M_m
\tag{8.2.8.1}
\end{equation}
로 놓는다. $S_{[m]}M_{[n]}\subset M_{[m+n]}$ ($m\geq0$)이므로,
다음과 같이 놓아 등급 $S^\natural$-가군 $M^\natural$을 정의할 수 있다.
\begin{equation}
\label{II.8.2.8.2-ko}
  M^\natural=\bigoplus_{n\in\mathbf{Z}}M_n^\natural,
  \quad\text{여기서}\quad M_n^\natural=M_{[n]}.
\tag{8.2.8.2}
\end{equation}
\end{env}

다음 보조정리의 증명도 독자에게 맡긴다~:

\begin{lemma}[8.2.9]
\label{II.8.2.9-ko}
(\hyperref[II.8.2.7-ko]{8.2.7})과
(\hyperref[II.8.2.8-ko]{8.2.8})의 표기 아래, 다음과 같은 가군의 표준
디-동형사상들이 존재한다.
\begin{equation}
\label{II.8.2.9.1-ko}
  M_f\xrightarrow{\sim}\bigoplus_{n\in\mathbf{Z}}M(n)_{(f)}
\tag{8.2.9.1}
\end{equation}
\begin{equation}
\label{II.8.2.9.2-ko}
  (M_f^{\geq})_{f/1}\xrightarrow{\sim}M_f
\tag{8.2.9.2}
\end{equation}
\begin{equation}
\label{II.8.2.9.3-ko}
  M^\natural_{(f^\natural)}\xrightarrow{\sim}M_f^{\geq}
\tag{8.2.9.3}
\end{equation}
이 가운데 처음 두 개는 등급가군의 디-동형사상이다.
\end{lemma}

\begin{lemma}[8.2.10]
\label{II.8.2.10-ko}
$S$를 양의 차수로 등급화된 환이라 하자.
\begin{enumerate}
\item[\rm(i)] $S^\natural$이 유한형(각각, 뇌터) $S_0^\natural$-대수이기
위한 필요충분조건은 $S$가 유한형(각각, 뇌터) $S_0$-대수인 것이다.
\item[\rm(ii)] $n\geq n_0$에 대하여
$S_{n+1}^\natural=S_1^\natural S_n^\natural$이기 위한 필요충분조건은
$n\geq n_0$에 대하여 $S_{n+1}=S_1S_n$인 것이다.
\item[\rm(iii)] $n\geq n_0$에 대하여 $S_n^\natural=S_1^{\natural n}$이기
위한 필요충분조건은 $n\geq n_0$에 대하여 $S_n=S_1^n$인 것이다.
\item[\rm(iv)] $(f_\alpha)$가 $S_+$의 동차원소들의 집합이고, $S_+$가
$f_\alpha$들이 생성하는 $S_+$의 아이디얼의 $S_+$ 안에서의 근기라면,
$S_+^\natural$은 $f_\alpha^\natural$들이 생성하는
$S_+^\natural$의 아이디얼의 $S_+^\natural$ 안에서의 근기이다.
\end{enumerate}
\end{lemma}

\begin{enumerate}
\item[\rm(i)] $S^\natural$이 유한형 $S_0^\natural$-대수이면,
(\hyperref[II.2.1.6-ko]{2.1.6}, (i))에 따라
$S_+=S_1^\natural$은 $S=S_0^\natural$ 위의 유한 생성 가군이므로,
$S$는 유한형 $S_0$-대수이다(\hyperref[II.2.1.4-ko]{2.1.4})~;
$S^\natural$이 뇌터환이면 $S_0^\natural=S$도 그러하다
(\hyperref[II.2.1.5-ko]{2.1.5}). 역으로 $S$가 유한형 $S_0$-대수이면,

\oldpage[II]{160}
(\hyperref[II.2.1.6-ko]{2.1.6}, (ii))에 따라 $n\geq m_0$일 때
$S_{n+h}=S_hS_n$이 되도록 하는 $h>0$과 $m_0>0$이 존재함을 안다~;
명백히 $m_0\geq h$라고 가정할 수 있다. 또한 $S_m$들은 유한 생성
$S_0$-가군이다(\hyperref[II.2.1.6-ko]{2.1.6}, (i)). 이제
$n\geq m_0+h$이면 $S_n^\natural=S_hS_{n-h}^\natural
=S_h^\natural S_{n-h}^\natural$이고~; $m<m_0+h$이면
$E=S_{m_0}+\cdots+S_{m_0+h-1}$로 놓을 때
\[
  S_m^\natural=S_m+\cdots+S_{m_0+h-1}+S_hE+S_h^2E+\cdots.
\]
$1\leq m\leq m_0$에 대하여, $m\leq i\leq m_0+h-1$인 $S_i$들의
$S_0$-가군으로서의 유한 생성계들을 합친 집합을 $G_m$이라 하고,
이를 $S_{[m]}$의 부분집합으로 본다. $m_0+1\leq m\leq m_0+h-1$에
대하여도 마찬가지로, $m\leq i\leq m_0+h-1$인 $S_i$들의
$S_0$-가군으로서의 유한 생성계들과 $S_hE$의 생성계를 합친 집합을
$G_m$이라 하고, 이를 $S_{[m]}$의 부분집합으로 본다.
$1\leq m\leq m_0+h-1$에 대하여
$S_m^\natural=S_0^\natural G_m$임은 명백하고, 따라서
$1\leq m\leq m_0+h-1$인 $G_m$들을 합친 집합 $G$는
$S_0^\natural$-대수 $S^\natural$의
생성계이다. 이로부터 $S=S_0^\natural$이 뇌터환이면
$S^\natural$도 뇌터환이라는 결론을 얻는다.

\item[\rm(ii)] $n\geq n_0$에 대하여 $S_{n+1}=S_1S_n$이면
$S_{n+1}^\natural=S_1S_n^\natural$이고, 
\emph{더욱이} $n\geq n_0$에 대하여
$S_{n+1}^\natural=S_1^\natural S_n^\natural$임은 명백하다.
역으로 마지막 관계는
\[
  S_{n+1}+S_{n+2}+\cdots=(S_1+S_2+\cdots)(S_n+S_{n+1}+\cdots)
\]
로 쓰이며, 양변에서 ($S$ 안의) 차수 $n+1$인 항들을 비교하면
$S_{n+1}=S_1S_n$을 얻는다.

\item[\rm(iii)] $n\geq n_0$에 대하여 $S_n=S_1^n$이면
$S_n^\natural=S_1^n+S_1^{n+1}+\cdots$이다~;
$S_1^\natural$은 $S_1+S_1^2+\cdots$를 포함하므로
$S_n^\natural\subset S_1^{\natural n}$이고, 따라서 $n\geq n_0$에
대하여 $S_n^\natural=S_1^{\natural n}$이다. 역으로
$S_1^{\natural n}=(S_1+S_2+\cdots)^n$에서 $S$ 안의 차수가 $n$인
항들은 $S_1^n$의 항들뿐이다~; 따라서
$S_n^\natural=S_1^{\natural n}$이라는 관계는 $S_n=S_1^n$을 함의한다.

\item[\rm(iv)] 원소 $g\in S_{k+h}$를 $S_k^\natural$의 원소로 볼 때
($k>0$, $h\geq0$), $S_{kn}^\natural$ 안에서 $g^n$이
$S^\natural$의 원소들을 계수로 하는 $f_\alpha^\natural$들의
선형결합이 되도록 하는 정수 $n>0$이 존재함을 증명하면 충분하다.
가정에 따라, $m\geq m_0$이면 
\emph{$S$ 안에서} $g^m=\sum_\alpha c_{\alpha m}f_\alpha$가 되도록 하는
정수 $m_0$이 존재하며, 이 식에 나타나는 첨자 $\alpha$들은
\emph{$m$과 무관하다}~; 또한 $c_{\alpha m}$들은 다음을 만족하는
동차원소라고 명백히 가정할 수 있다.
\[
  \deg(c_{\alpha m})=m(k+h)-\deg f_\alpha
\]
이는 $S$에서의 차수이다. 이제 $g^{m_0}$에 나타나는 모든
$f_\alpha$에 대하여 $km_0>\deg f_\alpha$가 되도록 $m_0$을 충분히
크게 잡자~; 각 $\alpha$에 대하여 $c'_{\alpha m}$을
$S^\natural$에서 차수 $km-\deg(f_\alpha)$인 원소로 본
$c_{\alpha m}$이라 하자~; 그러면 $S^\natural$ 안에서
$g^m=\sum_\alpha c'_{\alpha m}f_\alpha^\natural$이고, 이로써 증명이 끝난다.
\end{enumerate}

\begin{env}[8.2.11]
\label{II.8.2.11-ko}
다음 등급 $S_0$-대수를 생각하자.
\begin{equation}
\label{II.8.2.11.1-ko}
  S^\natural\otimes_S S_0=S^\natural/S_+S^\natural
  =\bigoplus_{n\geq0}S_{[n]}/S_+S_{[n]}.
\tag{8.2.11.1}
\end{equation}
$S_n$은 $S_{[n]}/S_+S_{[n]}$의 몫 $S_0$-가군이므로, 등급
$S_0$-대수의 표준 준동형
\begin{equation}
\label{II.8.2.11.2-ko}
  S^\natural\otimes_S S_0\to S
\tag{8.2.11.2}
\end{equation}
를 얻는다. 이는 명백히 
\emph{전사}이고, 따라서 (\hyperref[II.2.9.2-ko]{2.9.2})에 의해 표준
\emph{닫힌 몰입 사상}
\begin{equation}
\label{II.8.2.11.3-ko}
  \operatorname{Proj}(S)\to
  \operatorname{Proj}(S^\natural\otimes_S S_0)
\tag{8.2.11.3}
\end{equation}
에 대응한다.
\end{env}

\oldpage[II]{161}

\begin{proposition}[8.2.12]
\label{II.8.2.12-ko}
표준 사상 (\hyperref[II.8.2.11.3-ko]{8.2.11.3})은 전단사이다.
준동형 (\hyperref[II.8.2.11.2-ko]{8.2.11.2})가 (TN)-전단사이기 위한
필요충분조건은 $n\geq n_0$에 대하여 $S_{n+1}=S_1S_n$이 되도록 하는
$n_0$가 존재하는 것이다. 이 마지막 조건이 성립하면
(\hyperref[II.8.2.11.3-ko]{8.2.11.3})은 동형사상이다~;
$S$가 뇌터환일 때에는 그 역도 성립한다.
\end{proposition}

첫 번째 주장을 증명하려면, (\hyperref[II.2.8.3-ko]{2.8.3})에 따라
준동형 (\hyperref[II.8.2.11.2-ko]{8.2.11.2})의 핵 $\mathfrak{I}$가
\emph{멱영원소}들로 이루어짐을 증명하면 충분하다. 실제로
$f\in S_{[n]}$이고 $S_+S_{[n]}$을 법으로 한 그 잉여류가 이 핵에
속한다면, 이는 $f\in S_{[n+1]}$임을 뜻한다~;
$f^{n+1}$을 $S_{[n(n+1)]}$에 속하는 원소로 보면
$f\cdot f^n$으로 쓰이므로 $S_+S_{[n(n+1)]}$에 속한다~; 따라서
$S_+S_{[n(n+1)]}$을 법으로 한 $f^{n+1}$의 잉여류는 영이고, 이로써
주장이 증명된다. $n\geq n_0$에 대한 가정
$S_{n+1}=S_1S_n$은 $n\geq n_0$에 대한
$S_{n+1}^{\natural}=S_1^{\natural}S_n^{\natural}$과 동치이므로
(\hyperref[II.8.2.10-ko]{8.2.10}, (ii)), 이 가정은 정의에 따라
(\hyperref[II.8.2.11.2-ko]{8.2.11.2})가 (TN)-단사라는 사실과 동치이고,
따라서 (TN)-전단사라는 사실과 동치이다. 그러면
(\hyperref[II.2.9.1-ko]{2.9.1})에 따라
(\hyperref[II.8.2.11.3-ko]{8.2.11.3})은 동형사상이다. 역으로
(\hyperref[II.8.2.11.3-ko]{8.2.11.3})이 동형사상이면,
$\operatorname{Proj}(S^{\natural}\otimes_S S_0)$ 위의 층
$\widetilde{\mathfrak{I}}$은 영이다
(\hyperref[II.2.9.2-ko]{2.9.2}, (i))~;
$S^{\natural}\otimes_S S_0$는 $S^{\natural}$의 몫으로서 뇌터환이므로
(\hyperref[II.8.2.10-ko]{8.2.10}, (i)),
(\hyperref[II.2.7.3-ko]{2.7.3})으로부터 $\mathfrak{I}$가 조건 (TN)을
만족한다는 결론을 얻는다. 따라서 $n\geq n_0$에 대하여
$S_{n+1}^{\natural}=S_1^{\natural}S_n^{\natural}$이고,

이로써 (\hyperref[II.8.2.10-ko]{8.2.10}, (ii))에 의해 증명이 끝난다.

\begin{env}[8.2.13]
\label{II.8.2.13-ko}
이제 표준 단사사상 $(S_+)^n\to S_{[n]}$을 생각하자. 이들은 차수 $0$인
등급환의 단사 준동형
\begin{equation}
\label{II.8.2.13.1-ko}
  \bigoplus_{n\geq0}(S_+)^n\to S^{\natural}.
\tag{8.2.13.1}
\end{equation}
을 정의한다.
\end{env}

\begin{proposition}[8.2.14]
\label{II.8.2.14-ko}
준동형 (\hyperref[II.8.2.13.1-ko]{8.2.13.1})이 (TN)-동형사상이기 위한
필요충분조건은 $n_0$가 존재하여 $S_n=S_1^n$이 모든 $n\geq n_0$에
대하여 성립하는 것이다. 이 조건이 성립하면
(\hyperref[II.8.2.13.1-ko]{8.2.13.1})에 대응하는 사상은 전체에서
정의되고 다음 동형사상이 된다.
\[
  \operatorname{Proj}(S^{\natural})\xrightarrow{\sim}
  \operatorname{Proj}\left(\bigoplus_{n\geq0}(S_+)^n\right)~;
\]
$S$가 뇌터환이면 그 역도 성립한다.
\end{proposition}

첫 두 주장은 (\hyperref[II.8.2.10-ko]{8.2.10}, (iii))과
(\hyperref[II.2.9.1-ko]{2.9.1})에 비추어 명백하다. 셋째 주장은
(\hyperref[II.8.2.10-ko]{8.2.10}, (i), (iii))과 다음 보조정리에서
따른다~:

\begin{lemma}[8.2.14.1]
\label{II.8.2.14.1-ko}
$T$를 양의 차수로 등급화되고 유한형 $T_0$-대수인 등급환이라 하자.
단사 준동형 $\bigoplus_{n\geq0}T_1^n\to T$에 대응하는 사상이 전체에서
정의되고 다음 동형사상이라고 하자.
\[
  \operatorname{Proj}(T)\to
  \operatorname{Proj}\left(\bigoplus_{n\geq0}T_1^n\right),
\]
그러면 $n_0$가 존재하여 $T_n=T_1^n$이 모든 $n\geq n_0$에 대하여
성립한다.
\end{lemma}

실제로 $g_i$ ($1\leq i\leq r$)를 $T_0$-가군 $T_1$의 생성원들이라 하자.
가정에 따르면 우선 $D_+(g_i)$들이 $\operatorname{Proj}(T)$를 덮는다
(\hyperref[II.2.8.1-ko]{2.8.1}). 족 $(h_j)_{1\leq j\leq s}$를
$T_+$의 동차 원소들로 택하되, $\deg(h_j)=n_j$이고 이 원소들이
$g_i$들과 함께 아이디얼 $T_+$의 생성계, 또는 (동치인 말로
(\hyperref[II.2.1.3-ko]{2.1.3})) $T$의 $T_0$-대수로서의 생성계를 이루게
하자. $T'=\bigoplus_{n\geq0}T_1^n$이라 놓으면, 가정에 따라 환
$h_j/g_i^{n_j}$는 환 $T_{(g_i)}$의 원소이며 부분환 $T'_{(g_i)}$에
속해야 한다. 따라서 정수 $k$가 존재하여
$T_1^k h_j\subset T_1^{k+n_j}$가 모든 $j$에 대하여 성립한다.
거듭제곱 지수 $r$에 대해 귀납하면 $T_1^k h_j^r\subset T'$가 모든
$r\geq1$에 대하여 성립한다. $h_j$들의 정의에
따라 $T_1^kT\subset T'$이다. 한편 각 $j$에 대하여 다음 성질을 갖는
정수 $m_j$가 있다~: $h_j^{m_j}$는 $T$에서
$g_i$들로 생성되는 아이디얼에 속한다
(\hyperref[II.2.3.14-ko]{2.3.14}), $h_j^{m_j}\in T_1T$이고,

\oldpage[II]{162}
$h_j^{m_jk}\in T_1^kT\subset T'$이다. 따라서 정수 $m_0\geq k$가
존재하여 $h_j^m\in T_1^{mn_j}$가 모든 $m\geq m_0$에 대하여 성립한다. 이제
$q$를 정수 $n_j$들 가운데 가장 큰 것으로 놓으면 $n_0=qsm_0+k$가
요구를 만족한다. 실제로 $T_n$의 원소는 $n\geq n_0$일 때
$T_1^\alpha u$에 속하는 단항식들의 합이다. 여기서 $u$는 $h_j$들의
거듭제곱들의 곱이다. $\alpha\geq k$이면 앞의 결과에서
$T_1^\alpha u\subset T_1^n$이 따른다. 반대의 경우에는 $h_j$ 가운데
하나의 지수가 $\geq m_0$이므로 $u\in T_1^\beta v$라 쓸 수 있다.
여기서 $\beta\geq k$이고 $v$는 여전히 $h_j$들의 거듭제곱들의 곱이다.
그러면 앞의 경우로 귀착되며, 결국 모든 경우에
$T_1^\alpha u\subset T_1^n$임을 알 수 있다.

\begin{remark}[8.2.15]
\label{II.8.2.15-ko}
조건 $S_n=S_1^n$ ($n\geq n_0$)은 분명히
$S_{n+1}=S_1S_n$ ($n\geq n_0$)을 함의하지만, $S$를 뇌터환이라고
가정해도 그 역은
성립하지 않는다. 예를 들어 $K$를 체라 하고
$A=K[x]$, $B=K[y]/y^2K[y]$라 하자. 여기서 $x$와 $y$는 두 부정원이고
$x$의 차수는 $1$, $y$의 차수는 $2$이다. $S=A\otimes_KB$라 놓으면
$S$는 $K$ 위의 등급 대수이고, 원소 $1$, $x^n$ ($n\geq1$),
$x^ny$ ($n\geq0$)로 이루어진 기저를 갖는다.
$S_{n+1}=S_1S_n$ ($n\geq2$)임은 즉시 알 수 있지만,
$S_1^n=Kx^n$인 반면 $S_n=Kx^n+Kx^{n-2}y$ ($n\geq2$)이다.
\end{remark}

\subsection{사영 원뿔}
\label{subsection:II.8.3-ko}

\begin{env}[8.3.1]
\label{II.8.3.1-ko}
$Y$를 준스킴이라 하자. 이 항 전체에서는 \emph{$Y$-준스킴}과
\emph{$Y$-사상}만을 다룬다. $\mathcal{S}$를 \emph{양의 차수로} 등급화된 준연접
$\mathcal{O}_Y$-대수라 하자. \emph{또한 $\mathcal{S}_0=\mathcal{O}_Y$라고
가정한다}. (\hyperref[II.8.2.2-ko]{8.2.2})에서 도입한 표기법에 따라
다음과 같이 놓는다.
\begin{equation}
\label{II.8.3.1.1-ko}
  \widehat{\mathcal{S}}=\mathcal{S}[z]
  =\mathcal{S}\otimes_{\mathcal{O}_Y}\mathcal{O}_Y[z]
\tag{8.3.1.1}
\end{equation}
이를 양의 차수로 등급화된 $\mathcal{O}_Y$-대수로 보되, 차수는
(\hyperref[II.8.2.2.2-ko]{8.2.2.2})의 공식으로 정의한다. 그러면 아핀
열린집합 $U$가 $Y$ 안에 있을 때마다
\[
  \Gamma(U,\widehat{\mathcal{S}})=(\Gamma(U,\mathcal{S}))[z].
\]
이다. 이하에서는
\begin{equation}
\label{II.8.3.1.2-ko}
  X=\operatorname{Proj}(\mathcal{S}),\qquad
  C=\operatorname{Spec}(\mathcal{S}),\qquad
  \widehat{C}=\operatorname{Proj}(\widehat{\mathcal{S}})
\tag{8.3.1.2}
\end{equation}
로 놓고 ($C$의 정의에서는 $\mathcal{S}$를 등급화되지 않은
$\mathcal{O}_Y$-대수로 본다), $C$ (각각 $\widehat{C}$)를
$\mathcal{S}$로 정의되는 \emph{아핀 원뿔} (각각 \emph{사영 원뿔})이라 한다.
또한 “아핀 원뿔” 대신 단순히 “원뿔”이라고도 한다. 용어를
남용하여, $C$ (각각 $\widehat{C}$)를 \emph{$X$를 사영하는 아핀 원뿔}
(각각 \emph{$X$를 사영하는 사영 원뿔})이라고도 한다. 여기서는 준스킴
$X$가 $\operatorname{Proj}(\mathcal{S})$의 꼴로 주어져 있다고 이해한다.
끝으로 $\widehat{C}$를 $C$의 \emph{사영 폐포}라고 한다
($\mathcal{S}$의 데이터가 $C$의 구조에 포함되어 있다고 이해한다).
\end{env}

\begin{proposition}[8.3.2]
\label{II.8.3.2-ko}
다음과 같은 표준 $Y$-사상들이 존재한다.
\begin{equation}
\label{II.8.3.2.1-ko}
  Y\xrightarrow{\varepsilon}C\xrightarrow{i}\widehat{C}
\tag{8.3.2.1}
\end{equation}
\begin{equation}
\label{II.8.3.2.2-ko}
  X\xrightarrow{j}\widehat{C}
\tag{8.3.2.2}
\end{equation}

\oldpage[II]{163}
$\varepsilon$과 $j$는 닫힌 몰입이고, $i$는 아핀 사상이며 우세 열린
몰입이다. 또한
\begin{equation}
\label{II.8.3.2.3-ko}
  i(C)=\widehat{C}-j(X)~;
\tag{8.3.2.3}
\end{equation}
이고, $\widehat{C}$는 $\widehat{C}$의 닫힌 부분준스킴 가운데
$i(C)$를 포함하는 가장 작은 것이다.
\end{proposition}

$i$를 정의하기 위해 $\widehat{C}$의 열린집합
\begin{equation}
\label{II.8.3.2.4-ko}
  \widehat{C}_z=
  \operatorname{Spec}(\widehat{\mathcal{S}}/(z-1)\widehat{\mathcal{S}})
\tag{8.3.2.4}
\end{equation}
을 생각한다 (\hyperref[II.3.1.4-ko]{3.1.4}). 여기서 $z$는 표준적으로
$\widehat{\mathcal{S}}$의 $Y$ 위의 한 단면과 동일시된다. 그러면 동형사상
$i:C\xrightarrow{\sim}\widehat{C}_z$는 표준 동형사상
(\hyperref[II.8.2.3.1-ko]{8.2.3.1})
\[
  \widehat{\mathcal{S}}/(z-1)\widehat{\mathcal{S}}
  \xrightarrow{\sim}\mathcal{S}.
\]
에 대응한다. 사상 $\varepsilon$은 증대 준동형
$\mathcal{S}\to\mathcal{S}_0=\mathcal{O}_Y$에 대응하고 그 핵은
$\mathcal{S}_+$이며 (\hyperref[II.1.2.7-ko]{1.2.7}), 이 준동형은 전사이므로
$\varepsilon$은 닫힌 몰입이다
(\hyperref[II.1.4.10-ko]{1.4.10}). 끝으로 $j$도 마찬가지로
(\hyperref[II.3.5.1-ko]{3.5.1}) 차수 $0$인 전사 준동형
$\widehat{\mathcal{S}}\to\mathcal{S}$에 대응한다. 이 준동형은
$\mathcal{S}$에서는 항등사상이고 그 핵인 $z\widehat{\mathcal{S}}$에서는
영이다. $j$는 전체에서 정의되고 (\hyperref[II.3.6.2-ko]{3.6.2})에
의해 닫힌 몰입이다.

(\hyperref[II.8.3.2-ko]{8.3.2})의 나머지 주장들을 증명하기 위해서는
분명히 $Y=\operatorname{Spec}(A)$가 아핀이고
$\mathcal{S}=\widetilde{S}$인 경우로 한정할 수 있다. 여기서 $S$는
등급 $A$-대수이고, 따라서
$\widehat{\mathcal{S}}=(\widehat{S})^{\sim}$이다~; 이때 $S_+$의 동차
원소 $f$는 $Y$ 위의 $\widehat{\mathcal{S}}$의 단면과 동일시되고,
(\hyperref[II.2.3.3-ko]{2.3.3})에서 $D_+(f)$로 나타낸
$\widehat{C}$의 열린집합은 $\widehat{C}_f$로도 쓸 수 있다
(\hyperref[II.3.1.4-ko]{3.1.4})~; 마찬가지로
(\hyperref[I.1.1.1-ko]{I, 1.1.1})에서 $D(f)$로 나타낸 $C$의 열린집합은
$C_f$로도 쓸 수 있다 (\hyperref[0.5.5.2-ko]{0, 5.5.2}). 이제
(\hyperref[II.2.3.14-ko]{2.3.14})와 $\widehat{S}$의 정의로부터, 현재의
경우 열린집합 $\widehat{C}_z=i(C)$와 $\widehat{C}_f$ ($f$는 $S_+$의
동차 원소)가 $\widehat{C}$의 \emph{덮개}를 이룬다. 또한 이 표기 아래
다음이 성립한다.
\begin{equation}
\label{II.8.3.2.5-ko}
  i^{-1}(\widehat{C}_f)=C_f~;
\tag{8.3.2.5}
\end{equation}
실제로
$\widehat{C}_f\cap i(C)=\widehat{C}_f\cap\widehat{C}_z
=\widehat{C}_{fz}=\operatorname{Spec}(\widehat{S}_{(fz)})$이다. 그런데
$d=\deg(f)$라 하면, $\widehat{S}_{(fz)}$는
$(\widehat{S}_{(z)})_{f/z^d}$와 표준적으로 동형이고
(\hyperref[II.2.2.2-ko]{2.2.2}), 동형사상
(\hyperref[II.8.2.3.1-ko]{8.2.3.1})의 정의로부터 분수환들의 대응하는
동형사상에 의한 $(\widehat{S}_{(z)})_{f/z^d}$의 상은 정확히 $S_f$이다.
$C_f=\operatorname{Spec}(S_f)$이므로, 이는
(\hyperref[II.8.3.2.5-ko]{8.3.2.5})를 증명하는 동시에 사상 $i$가
아핀임을 보인다~; 더 나아가 $\widehat{C}_f$ 안으로 가는 사상으로 본
$i$의 $C_f$에 대한 제한은
(\hyperref[I.1.7.3-ko]{I, 1.7.3})에 따라 표준 준동형
$\widehat{S}_{(f)}\to\widehat{S}_{(fz)}$에 대응한다. 앞서 본 것과
(\hyperref[II.8.2.3.2-ko]{8.2.3.2})에 따라 다음 결과를 서술할 수 있다.

\begin{env}[8.3.2.6]
\label{II.8.3.2.6-ko}
$Y=\operatorname{Spec}(A)$가 아핀이고 $\mathcal{S}=\widetilde{S}$라 하자.
그러면 $S_+$의 모든 동차 원소 $f$에 대하여 $\widehat{C}_f$는
$\operatorname{Spec}(S_f^{\leq})$와 표준적으로 동일시되고, $i$의
제한인 사상 $C_f\to\widehat{C}_f$는 표준 단사 준동형
$S_f^{\leq}\to S_f$에 대응한다.
\end{env}

이제 ($Y$가 아핀일 때) $\widehat{C}=\operatorname{Proj}(\widehat{S})$에서
$\widehat{C}_z$의 여집합은

\oldpage[II]{164}
$z$를 포함하는 $\widehat{S}$의 등급 소아이디얼들의 집합이며, 이는
$j$의 정의에 따라 $j(X)$이다. 이로써
(\hyperref[II.8.3.2.3-ko]{8.3.2.3})이 증명된다.

마지막으로 (\hyperref[II.8.3.2-ko]{8.3.2})의 마지막 주장을 증명하기
위해서도 언제나 $Y$가 아핀이라고 가정할 수 있다. 위의 표기 아래
$\widehat{S}$에서 $z$는 $0$의 약수(영인자)가 아님에 유의하자~;
$\overline{i(C)}=\widehat{C}$이므로 다음 보조정리를 증명하면 충분하다.

\begin{lemma}[8.3.2.7]
\label{II.8.3.2.7-ko}
$T$를 양의 차수들로 등급화된 환, $Z=\operatorname{Proj}(T)$,
$g$를 차수 $d>0$인 $T$의 동차 원소라 하자. $g$가 $T$에서 $0$의
약수(영인자)가
아니면, $Z$는 $Z_g=D_+(g)$를 포함하는 $Z$의 가장 작은 닫힌
부분준스킴이다.
\end{lemma}

(\hyperref[I.4.1.9-ko]{I, 4.1.9})에 따라 문제는 $Z$ 위에서 국소적이다~;
따라서 모든 동차 원소 $h\in T_e$ ($e>0$)에 대하여 $Z_h$가
$Z_{gh}$를 포함하는 $Z_h$의 가장 작은 닫힌 부분준스킴임을 증명하면
충분하다~; 정의와 (\hyperref[I.4.3.2-ko]{I, 4.3.2})로부터 이 조건은
표준 준동형 $T_{(h)}\to T_{(gh)}$가 \emph{단사}라는 것과 동치이다.
그런데 이 준동형은 표준 준동형
$T_{(h)}\to(T_{(h)})_{g^e/h^d}$와 동일시된다
(\hyperref[II.2.2.3-ko]{2.2.3}). 한편 $g^e$는 $T$에서 $0$의
약수(영인자)가 아니므로,
$g^e/h^d$는 $T_h$에서 (따라서 \emph{더욱이} $T_{(h)}$에서도) 영인자가
아니다. 실제로 $t\in T$, $m>0$에 대하여
$(g^e/h^d)(t/h^m)=0$이면 어떤 $n>0$에 대하여 $h^ng^et=0$이고, 따라서
$h^nt=0$이며, 결국 $T_h$에서 $t/h^m=0$이다. 이로써 증명이 끝난다
(\hyperref[0.1.2.2-ko]{0, 1.2.2}).

\begin{env}[8.3.3]
\label{II.8.3.3-ko}
열린 몰입 사상 $i$를 통하여 아핀 원뿔 $C$를 열린집합 $i(C)$ 위에서
사영 원뿔 $\widehat{C}$가 유도하는 부분준스킴과 흔히 동일시한다.
닫힌 몰입 사상 $\varepsilon$에 결부된 $C$의 닫힌 부분준스킴을
\emph{$C$의 꼭짓점 준스킴}이라 한다~; $C$의 $Y$-단면인
$\varepsilon$을 \emph{꼭짓점 단면} 또는 \emph{$C$의 영단면}이라고도
한다~; $\varepsilon$을 통하여 $Y$를 $C$의 꼭짓점 준스킴과 동일시할
수 있다. 또한 $i\circ\varepsilon$은 $\widehat{C}$의 $Y$-단면이고,
따라서 역시 닫힌 몰입 사상이며
(\hyperref[I.5.4.6-ko]{I, 5.4.6}), 이는 핵이
$\mathcal{S}_+[z]=\mathcal{S}_+\widehat{\mathcal{S}}$인 차수 $0$의 표준
전사 준동형
$\widehat{\mathcal{S}}=\mathcal{S}[z]\to\mathcal{O}_Y[z]$
(\hyperref[II.3.1.7-ko]{3.1.7} 참조)에 대응한다~; 이 닫힌 몰입 사상에
결부된 $\widehat{C}$의 부분준스킴도 \emph{$\widehat{C}$의 꼭짓점
준스킴}이라 하고, $i\circ\varepsilon$을 \emph{$\widehat{C}$의 꼭짓점
단면}이라 한다~; $i\circ\varepsilon$을 통하여 이를 $Y$와 동일시할
수 있다. 끝으로 $j$에 결부된 $\widehat{C}$의 닫힌 부분준스킴을
\emph{$\widehat{C}$의 무한원 자취}라 하며, $j$를 통하여 $X$와
동일시할 수 있다.
\end{env}

\begin{env}[8.3.4]
\label{II.8.3.4-ko}
각각의 \emph{열린집합}
\begin{equation}
\label{II.8.3.4.1-ko}
  E=C-\varepsilon(Y),\qquad
  \widehat{E}=\widehat{C}-i(\varepsilon(Y))
\tag{8.3.4.1}
\end{equation}
위에서 각각 $C$와 $\widehat{C}$로부터 유도된 부분준스킴을 (용어의
남용에 따라) $\mathcal{S}$가 정의하는 \emph{천공 아핀 원뿔}과
\emph{천공 사영 원뿔}이라 한다~; 이 용어에도 불구하고 $E$는 $Y$
위에서 \emph{반드시 아핀인 것은 아니며}, $\widehat{E}$도 $Y$ 위에서
사영인 것은 아님에 유의해야 한다
(\hyperref[II.8.4.3-ko]{8.4.3} 참조). $C$를 $i(C)$와 동일시하면
바탕공간들에 대하여 다음을 얻는다.
\begin{equation}
\label{II.8.3.4.2-ko}
  C\cup\widehat{E}=\widehat{C},\qquad C\cap\widehat{E}=E
\tag{8.3.4.2}
\end{equation}
따라서 $\widehat{C}$는 열린 부분준스킴 $C$와 $\widehat{E}$를
\emph{붙여서} 얻은 것으로 볼 수 있다~; 더 나아가
(\hyperref[II.8.3.2.3-ko]{8.3.2.3})에 따라
\begin{equation}
\label{II.8.3.4.3-ko}
  E=\widehat{E}-j(X).
\tag{8.3.4.3}
\end{equation}

\oldpage[II]{165}
$Y=\operatorname{Spec}(A)$가 아핀일 때,
(\hyperref[II.8.3.2-ko]{8.3.2})의 표기 아래
\begin{equation}
\label{II.8.3.4.4-ko}
  E=\bigcup C_f,\qquad
  \widehat{E}=\bigcup\widehat{C}_f,\qquad
  C_f=C\cap\widehat{C}_f
\tag{8.3.4.4}
\end{equation}
이다. 여기서 $f$는 $S_+$의 동차 원소들의 집합 전체를 돈다 (또는 이
집합의 부분집합 $M$만을 도는데, 이 부분집합이 생성하는 $S_+$의 아이디얼의
$S_+$ 안에서의 근기는 $S_+$ 자체이다. 다시 말해 $f\in M$에 대한
$X_f$들이 $X$를 덮는다
(\hyperref[II.2.3.14-ko]{2.3.14})). $C_f$를 따라 $C$와
$\widehat{C}_f$를 붙이는 것은 단사 사상 $C_f\to C$와
$C_f\to\widehat{C}_f$에 의해 결정되며, 앞서 본 것처럼
(\hyperref[II.8.3.2.6-ko]{8.3.2.6}) 이들은 각각 표준 준동형
$S\to S_f$와 $S_f^{\leq}\to S_f$에 대응한다.
\end{env}

\begin{proposition}[8.3.5]
\label{II.8.3.5-ko}
(\hyperref[II.8.3.1-ko]{8.3.1})과
(\hyperref[II.8.3.4-ko]{8.3.4})의 표기 아래, 표준 단사 준동형
$\varphi:\mathcal{S}\to\widehat{\mathcal{S}}=\mathcal{S}[z]$에 결부된
사상 (\hyperref[II.3.5.1-ko]{3.5.1})은 다음을 만족하는 아핀 전사 사상
(\emph{표준 수축}이라 한다)이다.
\begin{equation}
\label{II.8.3.5.1-ko}
  p:\widehat{E}\to X
\tag{8.3.5.1}
\end{equation}
즉,
\begin{equation}
\label{II.8.3.5.2-ko}
  p\circ j=1_X.
\tag{8.3.5.2}
\end{equation}
\end{proposition}

명제를 증명하기 위해서는 $Y$가 아핀인 경우로 한정할 수 있다.
$\widehat{E}$의 표현 (\hyperref[II.8.3.4.4-ko]{8.3.4.4})를 고려하면,
$p$의 정의역 $G(\varphi)$가 $\widehat{E}$와 같다는 사실은 다음 주장들
가운데 첫 번째 주장으로부터 따를 것이다.

\begin{env}[8.3.5.3]
\label{II.8.3.5.3-ko}
$Y=\operatorname{Spec}(A)$가 아핀이고 $\mathcal{S}=\widetilde{S}$이면,
모든 동차원소 $f\in S_+$에 대하여
\begin{equation}
\label{II.8.3.5.4-ko}
  p^{-1}(X_f)=\widehat{C}_f
\tag{8.3.5.4}
\end{equation}
이다. 또한 $\widehat{C}_f=\operatorname{Spec}(S_f^{\leq})$에 대한
$p$의 제한을 $\widehat{C}_f$에서 $X_f$로 가는 사상으로 보면, 이는
표준 단사 준동형 $S_{(f)}\to S_f^{\leq}$에 대응한다. 더 나아가
$f\in S_1$이면 $\widehat{C}_f$는
$X_f\otimes_{\mathbf{Z}}\mathbf{Z}[T]$ ($T$는 부정원)와 동형이다.
\end{env}

실제로 공식 (\hyperref[II.8.3.5.4-ko]{8.3.5.4})는
(\hyperref[II.2.8.1.1-ko]{2.8.1.1})의 특수한 경우일 뿐이며, 두 번째
명제는 $Y$가 아핀일 때 $\operatorname{Proj}(\varphi)$의 정의에
지나지 않는다 (\hyperref[II.2.8.1-ko]{2.8.1}). 이제 공식
(\hyperref[II.8.3.5.2-ko]{8.3.5.2})와 $p$의 전사성은 표준 준동형들의
합성 $\mathcal{S}\to\widehat{\mathcal{S}}\to\mathcal{S}$가
$\mathcal{S}$의 항등 준동형이라는 사실에서 따른다. 끝으로
(\hyperref[II.8.3.5.3-ko]{8.3.5.3})의 마지막 명제는 $f\in S_1$일 때
$S_f^{\leq}$가 $S_{(f)}[T]$와 동형이라는 사실에서 나온다
(\hyperref[II.2.2.1-ko]{2.2.1}).

\begin{corollary}[8.3.6]
\label{II.8.3.6-ko}
$p$를 $E$에 제한한 사상
\begin{equation}
\label{II.8.3.6.1-ko}
  \pi:E\to X
\tag{8.3.6.1}
\end{equation}
은 아핀 전사 사상이다. $Y$가 아핀이고 $f$가 $S_+$의 동차원소이면
\begin{equation}
\label{II.8.3.6.2-ko}
  \pi^{-1}(X_f)=C_f
\tag{8.3.6.2}
\end{equation}
이고, $\pi$를 $C_f$에 제한한 사상은 표준 단사 준동형
$S_{(f)}\to S_f$에 대응한다. 또한 $f\in S_1$이면 $C_f$는
$X_f\otimes_{\mathbf{Z}}\mathbf{Z}[T,T^{-1}]$ ($T$는 부정원)와
동형이다.
\end{corollary}

공식 (\hyperref[II.8.3.6.2-ko]{8.3.6.2})는
(\hyperref[II.8.3.5.3-ko]{8.3.5.3})과
(\hyperref[II.8.3.2.5-ko]{8.3.2.5})에서 곧바로 따르며, $\pi$의
전사성을 보인다~; 한편 몰입 사상 $i$를 $C_f$에 제한하면

\oldpage[II]{166}
$S_f^{\leq}\to S_f$라는 단사 준동형에 대응한다는 것을 앞에서 보았다
(\hyperref[II.8.3.2-ko]{8.3.2}). 끝으로 마지막 명제는 $f\in S_1$일
때 $S_f$가 $S_{(f)}[T,T^{-1}]$와 동형이라는 사실에서 따른다
(\hyperref[II.2.2.1-ko]{2.2.1}).

\begin{remark}[8.3.7]
\label{II.8.3.7-ko}
$Y$가 아핀이면 $E$의 바탕공간의 원소들은 $S_+$를 포함하지 않는
$S$의 소 아이디얼 $\mathfrak{p}$ (반드시 등급 아이디얼일 필요는 없다)이다.
이는 몰입 사상 $\varepsilon$의 정의에서 따른다
(\hyperref[II.8.3.2-ko]{8.3.2}). 그러한 아이디얼 $\mathfrak{p}$에
대하여 $\mathfrak{p}\cap S_n$은
(\hyperref[II.2.1.9-ko]{2.1.9})의 조건을 자명하게 만족하므로, 모든
$n$에 대하여 $\mathfrak{q}\cap S_n=\mathfrak{p}\cap S_n$인 $S$의
유일한 \emph{등급} 소 아이디얼 $\mathfrak{q}$가 존재한다~; 바탕공간
사이의 사상 $\pi:E\to X$는 이때 다음 관계로 해석된다.
\begin{equation}
\label{II.8.3.7.1-ko}
  \pi(\mathfrak{p})=\mathfrak{q}.
\tag{8.3.7.1}
\end{equation}
실제로 이 관계를 확인하려면 $\mathfrak{p}\in D(f)$인 $S_+$의 동차원소
$f$를 하나 택하고, 단사 준동형 $S_{(f)}\to S_f$에 의한
$\mathfrak{p}_f$의 역상이 $\mathfrak{q}_{(f)}$임을 관찰하면 충분하다.
\end{remark}

\begin{corollary}[8.3.8]
\label{II.8.3.8-ko}
$\mathcal{S}$가 $\mathcal{S}_1$로 생성되면 사상 $p$와 $\pi$는
유한형이다~; 모든 $x\in X$에 대하여 섬유 $p^{-1}(x)$는
$\operatorname{Spec}(k(x)[T])$와 동형이고, 섬유 $\pi^{-1}(x)$는
$\operatorname{Spec}(k(x)[T,T^{-1}])$와 동형이다.
\end{corollary}

이는 (\hyperref[II.8.3.5-ko]{8.3.5})와
(\hyperref[II.8.3.6-ko]{8.3.6})에서 곧바로 따른다. 실제로 $Y$가
아핀이고 $S$가 $S_1$로 생성되면 $f\in S_1$에 대한 $X_f$들이 $X$를
덮는다 (\hyperref[II.2.3.14-ko]{2.3.14}).

\begin{remark}[8.3.9]
\label{II.8.3.9-ko}
등급 $\mathcal{O}_Y$-대수 $\mathcal{O}_Y[T]$ ($T$는 부정원)에 대응하는
꼭짓점이 제거된 아핀 원뿔은
$G_m=\operatorname{Spec}(\mathcal{O}_Y[T,T^{-1}])$와 동일시된다. 실제로
이는 (\hyperref[II.8.3.2-ko]{8.3.2})에서 본 $C_T$에 지나지 않는다
(더 일반적인 결과는 (\hyperref[II.8.4.4-ko]{8.4.4})를 보라). 이
준스킴에는 표준적으로 «~\emph{가환 $Y$-군 스킴}~»의 구조가 주어진다.
이 개념은 뒤에서 자세히 전개할 것이며, 여기서는 다음과 같이 간단히
요약할 수 있다. $Y$-군 스킴이란 두 $Y$-사상
$p:G\times_YG\to G$와 $s:G\to G$가 주어진 $Y$-스킴 $G$로서, 이 두
사상이 군의 합성법칙과 역원 사상의 공리와 형식적으로 유사한 조건을
만족하는 것이다~: 도식
\[
  \xymatrix{
    G\times G\times G \ar[r]^{p\times1}\ar[d]_{1\times p}
      & G\times G \ar[d]^{p} \\
    G\times G \ar[r]_{p} & G
  }
\]
은 가환이어야 하고(«~결합법칙~»), 군의 경우 다음 두 사상
\[
  (x,y)\mapsto(x,x^{-1},y)\mapsto(x,x^{-1}y)\mapsto x(x^{-1}y)
\]
과
\[
  (x,y)\mapsto(x,x^{-1},y)\mapsto(x,yx^{-1})\mapsto(yx^{-1})x
\]
이 모두 $(x,y)\mapsto y$로 귀착한다는 사실에 대응하는 조건도
만족해야 한다~; 예를 들어 첫 번째 합성 사상에 대응하는 사상열은
\[
  G\times G\xrightarrow{(1,s)\times1}G\times G\times G
  \xrightarrow{1\times p}G\times G\xrightarrow{p}G
\]
이며, 두 번째 사상열도 같은 방식으로 쓸 수 있다.

\oldpage[II]{167}
(\hyperref[I.3.4.3-ko]{I, 3.4.3})에 따라, $Y$-스킴 $G$에 $Y$-군 스킴
구조를 주는 것은 \emph{모든} $Y$-준스킴 $Z$에 대하여 집합
$\operatorname{Hom}_Y(Z,G)$에 \emph{군} 구조를 주는 것과 동치임은
명백하다. 이 구조들은 모든 $Y$-사상 $Z\to Z'$에 대하여 대응하는 사상
$\operatorname{Hom}_Y(Z',G)\to\operatorname{Hom}_Y(Z,G)$가 군 준동형이
되도록 해야 한다. 여기서 다루는 $G_m$의 경우,
$\operatorname{Hom}_Y(Z,G)$는 $Z\times_YG_m$의 $Z$-단면들의 집합과
동일시되고 (\hyperref[I.3.3.14-ko]{I, 3.3.14}), 따라서
$\operatorname{Spec}(\mathcal{O}_Z[T,T^{-1}])$의 $Z$-단면들의 집합과
동일시된다~; 끝으로 (\hyperref[I.3.3.15-ko]{I, 3.3.15})에서와 같은
논증에 따라 이 집합은 환 $\Gamma(Z,\mathcal{O}_Z)$의 \emph{가역원}들의
집합과 표준적으로 동일시되며, 이 집합의 군 구조는 환
$\Gamma(Z,\mathcal{O}_Z)$의 곱셈에서 유도된다. 위에서 다룬 사상 $p$와
$s$는 다음과 같이 얻어진다는 것을 확인할 수 있다~: 이들은
(\hyperref[II.1.2.7-ko]{1.2.7})과
(\hyperref[II.1.4.6-ko]{1.4.6})에 따라 다음 $\mathcal{O}_Y$-대수
준동형들에 대응한다.
\begin{align*}
  \pi &: \mathcal{O}_Y[T,T^{-1}]
  \to\mathcal{O}_Y[T,T^{-1},T',T'^{-1}],\\
  \sigma &: \mathcal{O}_Y[T,T^{-1}]\to\mathcal{O}_Y[T,T^{-1}]
\end{align*}
이 준동형들은 $\pi(T)=TT'$와 $\sigma(T)=T^{-1}$라는 값으로 완전히
정해진다.

이제 $G_m$은 양의 차수로 등급화된 준연접 $\mathcal{O}_Y$-대수
$\mathcal{S}$에 대한 임의의 \emph{아핀 원뿔}
$C=\operatorname{Spec}(\mathcal{S})$의 «~\emph{보편 작용소 집합}~»으로
볼 수 있다. 이는 작용소 군이 주어진 집합의 외부 연산과 같은 형식적
성질을 갖는 $Y$-사상 $G_m\times_YC\to C$를 표준적으로 정의할 수
있다는 뜻이다~; 또는 위의 군 스킴 경우와 마찬가지로, 모든
$Y$-준스킴 $Z$에 대하여 군 $\operatorname{Hom}_Y(Z,G_m)$을 작용소
집합으로 하는 $\operatorname{Hom}_Y(Z,C)$ 위의 외부 연산을 주되,
작용소 군이 주어진 집합의 통상적인 공리와 $Y$-사상 $Z\to Z'$에 대한
호환 조건을 만족하게 하는 것이다. 여기서 사상 $G_m\times_YC\to C$는
$\mathcal{O}_Y$-대수 준동형
\[
  \mathcal{S}\to
  \mathcal{S}\otimes_{\mathcal{O}_Y}\mathcal{O}_Y[T,T^{-1}]
  =\mathcal{S}[T,T^{-1}]
\]
으로 정의되며, 이 준동형은 각 단면
$s_n\in\Gamma(U,\mathcal{S}_n)$ ($U$는 $Y$의 열린집합)에

단면
$s_nT^n\in\Gamma(U,\mathcal{S}\otimes_{\mathcal{O}_Y}
\mathcal{O}_Y[T,T^{-1}])$을 대응시킨다.

거꾸로, 선험적으로 \emph{등급이 주어지지 않은} 준연접
$\mathcal{O}_Y$-대수 $\mathcal{S}$가 주어지고,
$C=\operatorname{Spec}(\mathcal{S})$ 위에 작용자 영역이 군
$Y$-스킴 $G_m$인 ``\emph{작용자군을 갖춘} 집합으로 이루어진
$Y$-스킴'' 구조가 주어졌다고 하자. 그러면 이로부터
$\mathcal{S}$ 위에 $\mathcal{O}_Y$-대수의 \emph{등급}을 표준적으로
얻는다. 실제로 $Y$-사상 $G_m\times_YC\to C$를 주는 것은
$\mathcal{O}_Y$-대수의 준동형
$\psi:\mathcal{S}\to\mathcal{S}[T,T^{-1}]$을 주는 것과 동치이므로,
이는
$\psi=\sum_{n\in\mathbf{Z}}\psi_nT^n$로 쓸 수 있다. 여기서
$\psi_n:\mathcal{S}\to\mathcal{S}$는 $\mathcal{O}_Y$-가군의
준동형이며($Y$의 열린집합 $U$와 모든 단면
$s\in\Gamma(U,\mathcal{S})$에 대하여 유한 개의 지표를 제외하면
$\psi_n(s)=0$이다), 작용자군을 갖춘 집합의 공리들이
$\psi_n(\mathcal{S})=\mathcal{S}_n$으로 하여
$\mathcal{S}$ 위에 $\mathcal{O}_Y$-대수의 등급(양수 또는 음수
차수를 허용한다)을 정의한다는 것을 확인할 수 있다. 이때
$\psi_n$들은 대응하는 사영상이다. 이로써 임의의 아핀
$Y$-스킴 위에 ``기하적인'' 방법으로, 미리 등급을 주지 않고도
\emph{아핀 원뽔} 구조를 정의할 수 있다.

\oldpage[II]{168}
여기서는 이 관점을 더 전개하지 않고, 앞의 설명에 해당하는
정의와 결과를 엄밀하게 정식화하는 일은 독자에게 맡긴다.
\end{remark}

\subsection{벡터 다발의 사영 폐포}
\label{subsection:II.8.4-ko}

\begin{env}[8.4.1]
\label{II.8.4.1-ko}
$Y$를 준스킴, $\mathcal{E}$를 준연접 $\mathcal{O}_Y$-가군이라
하자. $\mathcal{S}$로 등급 $\mathcal{O}_Y$-대수
$\mathbf{S}_{\mathcal{O}_Y}(\mathcal{E})$를 취하면, 정의
(\hyperref[II.8.3.1.1-ko]{8.3.1.1})에 따라
$\widehat{\mathcal{S}}$는
$\mathbf{S}_{\mathcal{O}_Y}(\mathcal{E}\oplus\mathcal{O}_Y)$와 동일시된다.
$\mathcal{S}$로 정의되는 아핀 원뽔
$\operatorname{Spec}(\mathcal{S})$는 정의상
$\mathbf{V}(\mathcal{E})$이고, $\operatorname{Proj}(\mathcal{S})$는 정의상
$\mathbf{P}(\mathcal{E})$이므로, 다음을 알 수 있다:
\end{env}

\begin{proposition}[8.4.2]
\label{II.8.4.2-ko}
$Y$ 위의 벡터 다발 $\mathbf{V}(\mathcal{E})$의 사영 폐포는
$\mathbf{P}(\mathcal{E}\oplus\mathcal{O}_Y)$와 표준적으로 동형이고,
후자의 무한원 자취는 $\mathbf{P}(\mathcal{E})$와 표준적으로 동형이다.
\end{proposition}

\begin{remark}[8.4.3]
\label{II.8.4.3-ko}
예를 들어 $r\geq2$인 $\mathcal{E}=\mathcal{O}_Y^r$을 취하자. 그러면
$\mathcal{S}$로 정의되는 꼭지점을 제거한 원뽔 $E$, $\widehat{E}$는
$Y\neq\varnothing$일 때 $Y$ 위에서 아핀이지도, 사영적이지도 않다.
둘째 주장은 즉시 따른다. 실제로
$\widehat{C}=\mathbf{P}(\mathcal{O}_Y^{r+1})$은 $Y$ 위에서 사영적이고,
$E$와 $\widehat{E}$의 바탕공간은 $\widehat{C}$의 닫히지 않은
열린집합이므로, 표준 몰입사상
$E\to\widehat{E}$와 $\widehat{E}\to\widehat{C}$는 사영적이지 않다
(\hyperref[II.5.5.3-ko]{5.5.3}). 따라서
(\hyperref[II.5.5.5-ko]{5.5.5, (v)})를 적용하면 결론을 얻는다.
다른 한편 $Y=\operatorname{Spec}(A)$가 아핀이고, 예를 들어
$r=2$라고 하자. 그러면
$C=\operatorname{Spec}(A[T_1,T_2])$이고 $E$는 열린집합
$D(T_1)\cup D(T_2)$ 위에서 $C$가 유도하는 준스킴이다. 그런데
이 열린집합은 아핀이 아님을 이미 보았다
(\hyperref[I.5.5.11-ko]{I, 5.5.11}). \emph{더구나} $E$는
$\widehat{E}$ 위의 단면 $z$가 영이 아닌 열린집합이므로
(\hyperref[II.8.3.2-ko]{8.3.2}), $\widehat{E}$는 아핀일 수 없다.

그러나 다음이 성립한다:
\end{remark}

\begin{proposition}[8.4.4]
\label{II.8.4.4-ko}
$\mathcal{L}$이 가역 $\mathcal{O}_Y$-가군이면,
$C=\mathbf{V}(\mathcal{L})$에 대응하는 꼭지점을 제거한 원뽔
$E$와 $\widehat{E}$에 대하여 다음 표준 동형들이 성립한다.
\begin{equation}
\label{II.8.4.4.1-ko}
  \operatorname{Spec}\left(
    \bigoplus_{n\in\mathbf{Z}}\mathcal{L}^{\otimes n}
  \right)\xrightarrow{\sim}E
\tag{8.4.4.1}
\end{equation}
\begin{equation}
\label{II.8.4.4.2-ko}
  \mathbf{V}(\mathcal{L}^{-1})\xrightarrow{\sim}\widehat{E}.
\tag{8.4.4.2}
\end{equation}

또한 $\mathbf{V}(\mathcal{L})$의 사영 폐포에서
$\mathbf{V}(\mathcal{L}^{-1})$의 사영 폐포로 가는 표준 동형사상이
존재하여, 전자의 영 단면(각각 무한원 자취)을 후자의 무한원
자취(각각 영 단면)로 보낸다.
\end{proposition}

\begin{proof}
여기서는
$\mathcal{S}=\bigoplus_{n\geq0}\mathcal{L}^{\otimes n}$이다. 표준 단사 준동형
\[
  \mathcal{S}\to\bigoplus_{n\in\mathbf{Z}}\mathcal{L}^{\otimes n}
\]
은 표준 우세 사상
\begin{equation}
\label{II.8.4.4.3-ko}
  \operatorname{Spec}\left(
    \bigoplus_{n\in\mathbf{Z}}\mathcal{L}^{\otimes n}
  \right)
  \to\mathbf{V}(\mathcal{L})
  =\operatorname{Spec}\left(
    \bigoplus_{n\geq0}\mathcal{L}^{\otimes n}
  \right)
\tag{8.4.4.3}
\end{equation}
을 정의한다. 이 사상이 스킴
$\operatorname{Spec}(\bigoplus_{n\in\mathbf{Z}}\mathcal{L}^{\otimes n})$에서
$E$ 위로 가는 동형사상임을 보이면 충분하다. 문제는
$Y$ 위에서 국소적이므로 $Y=\operatorname{Spec}(A)$가 아핀이고

\oldpage[II]{169}
$\mathcal{L}=\mathcal{O}_Y$라고 가정할 수 있다. 그러면
$\mathcal{S}=(A[T])^{\sim}$이고
$\bigoplus_{n\in\mathbf{Z}}\mathcal{L}^{\otimes n}
=(A[T,T^{-1}])^{\sim}$이다. 그런데 $A[T,T^{-1}]$는 $A[T]$의
분수환 $A[T]_T$이므로, (\hyperref[II.8.4.4.3-ko]{8.4.4.3})은
좌변을 열린집합 $D(T)$ 위에서
$C=\mathbf{V}(\mathcal{L})$가 유도하는 준스킴과 동일시한다. $C$에서
이 열린집합의 여집합 $V(T)$는 아이디얼 $TA[T]$로 정의되는
$C$의 닫힌 부분준스킴의 바탕공간이고, 이 부분준스킴은
$C$의 영 단면이다. 따라서 $E=D(T)$이다.

다른 한편 마지막 주장에서 동형
(\hyperref[II.8.4.4.2-ko]{8.4.4.2})가 따른다. 실제로
$\mathbf{V}(\mathcal{L}^{-1})$은 그 사영 폐포에서 무한원 자취의
여집합이고, $\widehat{E}$는
$C=\mathbf{V}(\mathcal{L})$의 사영 폐포에서 영 단면의 여집합이다.
이 사영 폐포들은 각각
$\mathbf{P}(\mathcal{L}^{-1}\oplus\mathcal{O}_Y)$와
$\mathbf{P}(\mathcal{L}\oplus\mathcal{O}_Y)$이다. 그런데
$\mathcal{L}\oplus\mathcal{O}_Y
=\mathcal{L}\otimes(\mathcal{L}^{-1}\oplus\mathcal{O}_Y)$로 쓸 수 있다.
따라서 원하는 표준 동형사상의 존재는
(\hyperref[II.4.1.4-ko]{4.1.4})에서 따른다. 이제 이 동형사상이
영 단면과 무한원 자취를 서로 바꾸어 보냄을 보이면 된다. 이를 위해
$Y=\operatorname{Spec}(A)$가 아핀이고,
$L=Ac$, $L^{-1}=Ac'$이며, 표준 동형사상
$L\otimes L^{-1}\to A$가 $c\otimes c'$을 $A$의 원소 $1$에
대응시키는 경우로 한정할 수 있다. 그러면
$\mathbf{S}(L\oplus A)$는 $A[z]$과
$\bigoplus_{n\geq0}Ac^{\otimes n}$의 텐서곱이고,
$\mathbf{S}(L^{-1}\oplus A)$는 $A[z]$과
$\bigoplus_{n\geq0}Ac'^{\otimes n}$의 텐서곱이며,
(\hyperref[II.4.1.4-ko]{4.1.4})에서 정의한 동형사상은
$z^h\otimes c'^{\otimes(n-h)}$을
$z^{n-h}\otimes c^{\otimes h}$에 대응시킨다. 그런데
$\mathbf{P}(\mathcal{L}^{-1}\oplus\mathcal{O}_Y)$에서 무한원 자취는
단면 $z$가 영이 되는 점들의 집합이고, 영 단면은 단면
$c'$가 영이 되는 점들의 집합이다.
$\mathbf{P}(\mathcal{L}\oplus\mathcal{O}_Y)$에 대해서도 동일하게 정의되므로,
앞의 구체적 설명에서 결론이 즉시 따른다.
\end{proof}

\subsection{함자적 성질}
\label{subsection:II.8.5-ko}

\begin{env}[8.5.1]
\label{II.8.5.1-ko}
$Y$, $Y'$을 두 준스킴, $q:Y'\to Y$를 사상,
$\mathcal{S}$ (각각 $\mathcal{S}'$)를 \emph{양의} 차수인 등급 준연접
$\mathcal{O}_Y$-대수(각각 $\mathcal{O}_{Y'}$-대수)라 하자.
등급 대수의 $q$-사상
\begin{equation}
\label{II.8.5.1.1-ko}
  \varphi:\mathcal{S}\to\mathcal{S}'.
\tag{8.5.1.1}
\end{equation}
을 생각하자. 이에 표준적으로 대응하는 사상
(\hyperref[II.1.5.6-ko]{1.5.6})
\[
  \Phi=\operatorname{Spec}(\varphi):
  \operatorname{Spec}(\mathcal{S}')\to\operatorname{Spec}(\mathcal{S})
\]
이 있어서, 도식
\begin{equation}
\label{II.8.5.1.2-ko}
  \xymatrix{
    C'\ar[r]^{\Phi}\ar[d] & C\ar[d]\\
    Y'\ar[r]_{q} & Y
  }
\tag{8.5.1.2}
\end{equation}
은 가환이다. 여기서
$C=\operatorname{Spec}(\mathcal{S})$,
$C'=\operatorname{Spec}(\mathcal{S}')$로 두었다.
\emph{또한
$\mathcal{S}_0=\mathcal{O}_Y$, $\mathcal{S}'_0=\mathcal{O}_{Y'}$라고
가정하자}. $\varepsilon:Y\to C$와 $\varepsilon':Y'\to C'$을 표준
몰입사상(\hyperref[II.8.3.2-ko]{8.3.2})이라 하면, 가환 도식
\begin{equation}
\label{II.8.5.1.3-ko}
  \xymatrix{
    Y'\ar[r]^{q}\ar[d]_{\varepsilon'} & Y\ar[d]^{\varepsilon}\\
    C'\ar[r]_{\Phi} & C
  }
\tag{8.5.1.3}
\end{equation}
을 얻는다.

\oldpage[II]{170}
이는 도식
\[
  \xymatrix{
    \mathcal{S}\ar[r]^{\varphi}\ar[d] & \mathcal{S}'\ar[d]\\
    \mathcal{O}_Y\ar[r] & \mathcal{O}_{Y'}
  }
\]
에 대응한다. 여기서 수직 화살표들은 증대 준동형이고, 이 도식이
가환인 것은 $\varphi$가 \emph{등급} 대수의 준동형이라는 가정에서
따른다.
\end{env}

\begin{proposition}[8.5.2]
\label{II.8.5.2-ko}
$E$ (각각 $E'$)가 $\mathcal{S}$ (각각 $\mathcal{S}'$)로 정의되는
꼭지점을 제거한 아핀 원뿔이면 $\Phi^{-1}(E)\subset E'$이다. 또한
$\operatorname{Proj}(\varphi):G(\varphi)\to\operatorname{Proj}(\mathcal{S})$
가 모든 곳에서 정의되면(다시 말해
$G(\varphi)=\operatorname{Proj}(\mathcal{S}')$이면),
$\Phi^{-1}(E)=E'$이고, 그 역도 성립한다.
\end{proposition}

\begin{proof}
첫째 주장은 (\hyperref[II.8.5.1.3-ko]{8.5.1.3})의 가환성에서 따른다.
둘째 주장을 증명하기 위해서는
$Y=\operatorname{Spec}(A)$와 $Y'=\operatorname{Spec}(A')$가 아핀이고
$\mathcal{S}=\widetilde{S}$,
$\mathcal{S}'=\widetilde{S'}$인 경우로 한정할 수 있다. $S_+$의 임의의
동차 원소 $f$에 대하여 $f'=\varphi(f)$로 두면
$\Phi^{-1}(C_f)=C'_{f'}$이다
(\hyperref[I.2.2.4.1-ko]{I, 2.2.4.1}).
$G(\varphi)=\operatorname{Proj}(S')$라는 것은 $S'_+$ 안에서
$f'=\varphi(f)$들이 생성하는 아이디얼의 근기가 $S'_+$ 자체라는
뜻이고 (\hyperref[II.2.8.1-ko]{2.8.1}) 및
(\hyperref[II.2.3.14-ko]{2.3.14}), 이는 다시 $C'_{f'}$들이 $E'$을
덮는다는 것과 동치이다
(\hyperref[II.8.3.4.4-ko]{8.3.4.4}).
\end{proof}

\begin{env}[8.5.3]
\label{II.8.5.3-ko}
$q$-사상 $\varphi$는 $\widehat{\varphi}(z)=z$로 둠으로써 등급 대수의
$q$-사상
\begin{equation}
\label{II.8.5.3.1-ko}
  \widehat{\varphi}:\widehat{\mathcal{S}}
  \to\widehat{\mathcal{S}'}
\tag{8.5.3.1}
\end{equation}
으로 표준적으로 연장된다. 따라서 사상
\[
  \widehat{\Phi}=\operatorname{Proj}(\widehat{\varphi}):
  G(\widehat{\varphi})\to\widehat{C}
  =\operatorname{Proj}(\widehat{\mathcal{S}})
\]
을 얻으며, 도식
\[
  \xymatrix{
    G(\widehat{\varphi})\ar[r]^{\widehat{\Phi}}\ar[d] &
    \widehat{C}\ar[d]\\
    Y'\ar[r]_{q} & Y
  }
\]
은 가환이다 (\hyperref[II.3.5.6-ko]{3.5.6}).
$i:C\to\widehat{C}$와 $i':C'\to\widehat{C'}$을 표준 열린 몰입사상
(\hyperref[II.8.3.2-ko]{8.3.2})이라 하면, 정의에서 곧바로
$i'(C')\subset G(\widehat{\varphi})$이고 도식
\begin{equation}
\label{II.8.5.3.2-ko}
  \xymatrix{
    C'\ar[r]^{\Phi}\ar[d]_{i'} & C\ar[d]^{i}\\
    G(\widehat{\varphi})\ar[r]_{\widehat{\Phi}} & \widehat{C}
  }
\tag{8.5.3.2}
\end{equation}
은 가환임을 알 수 있다. 끝으로
$X=\operatorname{Proj}(\mathcal{S})$,
$X'=\operatorname{Proj}(\mathcal{S}')$로 두고,
$j:X\to\widehat{C}$, $j':X'\to\widehat{C'}$을 표준 닫힌 몰입사상
(\hyperref[II.8.3.2-ko]{8.3.2})이라 하면, 이 몰입사상들의

정의에서 다시 $j'(G(\varphi))\subset G(\widehat{\varphi})$이고 도식

\oldpage[II]{171}
\begin{equation}
\label{II.8.5.3.3-ko}
  \xymatrix{
    G(\varphi)\ar[r]^{\operatorname{Proj}(\varphi)}\ar[d]_{j'} &
    X\ar[d]^{j}\\
    G(\widehat{\varphi})\ar[r]_{\widehat{\Phi}} & \widehat{C}
  }
\tag{8.5.3.3}
\end{equation}
은 가환임을 알 수 있다.
\end{env}

\begin{proposition}[8.5.4]
\label{II.8.5.4-ko}
$\widehat{E}$ (각각 $\widehat{E'}$)가 $\mathcal{S}$
(각각 $\mathcal{S}'$)로 정의되는 꼭지점을 제거한 사영 원뿔이면
$\widehat{\Phi}^{-1}(\widehat{E})\subset\widehat{E'}$이다. 또한
$p:\widehat{E}\to X$와 $p':\widehat{E'}\to X'$이 표준 수축사상이면
$p'(\widehat{\Phi}^{-1}(\widehat{E}))\subset G(\varphi)$이고, 도식
\begin{equation}
\label{II.8.5.4.1-ko}
  \xymatrix{
    \widehat{\Phi}^{-1}(\widehat{E})
      \ar[r]^{\widehat{\Phi}}\ar[d]_{p'} &
    \widehat{E}\ar[d]^{p}\\
    G(\varphi)\ar[r]_{\operatorname{Proj}(\varphi)} & X
  }
\tag{8.5.4.1}
\end{equation}
은 가환이다. $\operatorname{Proj}(\varphi)$가 모든 곳에서 정의되면
$\widehat{\Phi}$도 모든 곳에서 정의되고
$\widehat{\Phi}^{-1}(\widehat{E})=\widehat{E'}$이다.
\end{proposition}

\begin{proof}
첫째 주장은 도식 (\hyperref[II.8.5.1.3-ko]{8.5.1.3})과
(\hyperref[II.8.5.3.2-ko]{8.5.3.2})의 가환성에서 따르고, 그다음 두
주장은 표준 수축사상의 정의 (\hyperref[II.8.3.5-ko]{8.3.5})와
$\widehat{\varphi}$의 정의에서 따른다. 다른 한편
$\operatorname{Proj}(\varphi)$가 모든 곳에서 정의될 때
$\widehat{\Phi}$도 그러함을 보이기 위해서는
$Y=\operatorname{Spec}(A)$와 $Y'=\operatorname{Spec}(A')$가 아핀이고
$\mathcal{S}=\widetilde{S}$,
$\mathcal{S}'=\widetilde{S'}$인 경우만 생각하면 된다. 가정은 $f$가
$S_+$의 동차 원소 전체를 움직일 때 $\varphi(f)$들이 $S'_+$ 안에서
생성하는 아이디얼의 $S'_+$ 안에서의 근기가 아이디얼 $S'_+$ 자체라는
것이다. 이로부터 $z$와 $\varphi(f)$들이 생성하는 아이디얼의
$(S'[z])_+$ 안에서의 근기가 $(S'[z])_+$ 자체임을 곧바로 알 수 있고,
따라서 주장이 성립한다. 이 논증은 동시에 $\widehat{E'}$이
$\widehat{C'}_{\varphi(f)}$들의 합집합이고, 따라서
$\widehat{\Phi}^{-1}(\widehat{E})$과 같음을 증명한다.
\end{proof}

\begin{corollary}[8.5.5]
\label{II.8.5.5-ko}
$\operatorname{Proj}(\varphi)$가 모든 곳에서 정의될 때,
$\widehat{\Phi}$에 의한 $\widehat{C}$의 무한원 자취(각각 꼭지점
준스킴)의 바탕공간의 역상은 $\widehat{C'}$의 무한원 자취(각각
꼭지점 준스킴)의 바탕공간이다.
\end{corollary}

\begin{proof}
이는 (\hyperref[II.8.5.4-ko]{8.5.4})와
(\hyperref[II.8.5.2-ko]{8.5.2})에서 곧바로 따르며, 이때 관계식
(\hyperref[II.8.3.4.1-ko]{8.3.4.1})과
(\hyperref[II.8.3.4.2-ko]{8.3.4.2})를 사용한다.
\end{proof}

\subsection{꼭지점을 제거한 원뿔에 대한 표준 동형사상}
\label{subsection:II.8.6-ko}

\begin{env}[8.6.1]
\label{II.8.6.1-ko}
$Y$를 준스킴, $\mathcal{S}$를 양의 차수인 등급 준연접
$\mathcal{O}_Y$-대수로서
\emph{$\mathcal{S}_0=\mathcal{O}_Y$인 것}, $X$를 $Y$-스킴
$\operatorname{Proj}(\mathcal{S})$라 하자.
(\hyperref[subsection:II.8.5-ko]{8.5})의 결과를 $Y'=X$로 취하고
$q:X\to Y$를 구조사상으로 취한 경우에 적용하고자 한다. 다음과 같이
두자.
\begin{equation}
\label{II.8.6.1.1-ko}
  \mathcal{S}_X=\bigoplus_{n\in\mathbf{Z}}\mathcal{O}_X(n)
\tag{8.6.1.1}
\end{equation}

\oldpage[II]{172}
이는 등급 준연접 $\mathcal{O}_X$-대수이며, 그 곱셈은 표준 준동형
(\hyperref[II.3.2.6.1-ko]{3.2.6.1})
\[
  \mathcal{O}_X(m)\otimes_{\mathcal{O}_X}\mathcal{O}_X(n)
  \to\mathcal{O}_X(m+n)
\]
을 사용하여 정의된다. 이 곱셈의 결합법칙은 가환 도식
(\hyperref[II.2.5.11.4-ko]{2.5.11.4})으로 보장된다.
$\mathcal{S}'$로 $\mathcal{S}_X$의 양의 차수인 등급 준연접
부분-$\mathcal{O}_X$-대수
$\mathcal{S}_X^{\geq}=\bigoplus_{n\geq0}\mathcal{O}_X(n)$을 취하자.

마지막으로, 우리는 표준 $q$-사상
\begin{equation}
\label{II.8.6.1.2-ko}
  \alpha:\mathcal{S}\to\mathcal{S}_X^{\geq}
\tag{8.6.1.2}
\end{equation}
을 생각한다. 이는 (\hyperref[II.3.3.2.3-ko]{3.3.2.3})에서
$\mathcal{S}\to q_*(\mathcal{S}_X)$인 준동형으로 정의되지만, 명백히 $\mathcal{S}$를
$q_*(\mathcal{S}_X^{\geq})$ 안으로 보낸다. 다음과 같이 놓는다.
\begin{equation}
\label{II.8.6.1.3-ko}
  C_X=\operatorname{Spec}(\mathcal{S}_X^{\geq}),\qquad
  \widehat{C}_X=\operatorname{Proj}(\mathcal{S}_X^{\geq}[z]),\qquad
  X'=\operatorname{Proj}(\mathcal{S}_X^{\geq})
\tag{8.6.1.3}
\end{equation}
또한 대응하는 천공 아핀 원뿔과 천공 사영 원뿔을 각각 $E_X$와
$\widehat{E}_X$로 나타낸다~; (\hyperref[subsection:II.8.3-ko]{8.3})에서
정의된 표준 사상들을
$\varepsilon_X:X\to C_X$, $i_X:C_X\to\widehat{C}_X$,
$j_X:X'\to\widehat{C}_X$, $p_X:\widehat{E}_X\to X'$,
$\pi_X:E_X\to X'$로 나타내기로 한다.
\end{env}

\begin{proposition}[8.6.2]
\label{II.8.6.2-ko}
구조 사상 $u:X'\to X$는 동형사상이고,
$\operatorname{Proj}(\alpha)$는 모든 점에서 정의되며 $u$와 같다.
사상
$\operatorname{Proj}(\widehat{\alpha}):\widehat{C}_X\to\widehat{C}$는
모든 점에서 정의되고, 이를 $\widehat{E}_X$와 $E_X$에 제한하면 각각
$\widehat{E}$와 $E$ 위로 가는 동형사상이 된다. 마지막으로 $u$를
통하여 $X'$과 $X$를 동일시하면, 사상 $p_X$와 $\pi_X$는 각각
$X$-준스킴 $\widehat{E}_X$와 $E_X$의 구조 사상과 동일시된다.
\end{proposition}

\begin{proof}
명백히 $Y=\operatorname{Spec}(A)$가 아핀이고
$\mathcal{S}=\widetilde{S}$인 경우만 보면 된다~; 이때 $X$는 $S_+$의
동차 원소 $f$ 전체를 따라 움직이는 아핀 열린집합 $X_f$들의 합집합이고,
$X_f$의 환은 $S_{(f)}$이다. 더욱이
(\hyperref[II.8.2.7.1-ko]{8.2.7.1})에서 다음을 얻는다.
\begin{equation}
\label{II.8.6.2.1-ko}
  \Gamma(X_f,\mathcal{S}_X^{\geq})=S_f^{\geq}.
\tag{8.6.2.1}
\end{equation}

따라서 $u^{-1}(X_f)=\operatorname{Proj}(S_f^{\geq})$이다. 그런데
$f\in S_d$ ($d>0$)이면 $\operatorname{Proj}(S_f^{\geq})$는
$\operatorname{Proj}((S_f^{\geq})^{(d)})$와 표준적으로 동형이고
(\hyperref[II.2.4.7-ko]{2.4.7}), 한편
$(S_f^{\geq})^{(d)}=(S^{(d)})_f^{\geq}$는 $T\to f/1$인 사상에 의하여
$S_{(f)}[T]$와 동일시된다
(\hyperref[II.2.2.1-ko]{2.2.1})~; 그러므로
(\hyperref[II.3.1.7-ko]{3.1.7})에서 구조 사상
$u^{-1}(X_f)\to X_f$가 동형사상임을 얻으며, 첫 번째 주장이 따른다.
두 번째 주장을 보이기 위하여, $\operatorname{Proj}(\alpha)$의 제한
$u^{-1}(X_f)\cap G(\alpha)\to X=\operatorname{Proj}(S)$가 $S$에서
$S_f^{\geq}$로 가는 표준 사상 $x\to x/1$에 대응함을 주목하자
(\hyperref[II.2.6.2-ko]{2.6.2})~; 우선 여기서 $G(\alpha)=X'$을 얻는다.
이제 $u^{-1}(X_f)=(u^{-1}(X_f))_{f/1}$임을 고려하면,
(\hyperref[II.2.8.1.1-ko]{2.8.1.1})에 따라
$\operatorname{Proj}(\alpha)$에 의한 $u^{-1}(X_f)$의 상은 $X_f$에
포함된다. 또한 $\operatorname{Proj}(\alpha)$를 $u^{-1}(X_f)$에
제한하여 $X_f=\operatorname{Spec}(S_{(f)})$로 가는 사상으로 보면,
이는 실제로 $u$의 제한과 같다.
마지막으로 $p$ 대신 $p_X$에 공식
(\hyperref[II.8.3.5.4-ko]{8.3.5.4})를 적용하면
$p_X^{-1}(u^{-1}(X_f))=
\operatorname{Spec}((S_f^{\geq})_{f/1}^{\leq})$임을 알 수 있다.
(\hyperref[II.8.5.4.1-ko]{8.5.4.1})에 따르면 이 열린집합은
$\operatorname{Proj}(\widehat{\alpha})$에 의한
$p^{-1}(X_f)=\operatorname{Spec}(S_f^{\leq})$의 역상이다
(\hyperref[II.8.3.5.3-ko]{8.3.5.3}).

(\hyperref[II.8.2.3.2-ko]{8.2.3.2})를 고려하면,
$\operatorname{Proj}(\widehat{\alpha})$를
$p_X^{-1}(u^{-1}(X_f))$에 제한한 사상은
(\hyperref[II.8.2.7.2-ko]{8.2.7.2})의 동형사상의 역을
$S_f^{\leq}$에 제한한 것에 대응한다. 이로부터 세 번째 주장이
따르며, 마지막 주장은 정의에 의하여 명백하다.

\oldpage[II]{173}
또한 가환 그림 (\hyperref[II.8.5.3.2-ko]{8.5.3.2})에서
\emph{$\operatorname{Proj}(\widehat{\alpha})$를 $C_X$에 제한한 사상은
다름 아닌 $\operatorname{Spec}(\alpha)$임}을 알 수 있다.
\end{proof}

\begin{corollary}[8.6.3]
\label{II.8.6.3-ko}
$X$-스킴으로 보면, $\widehat{E}_X$는
$\operatorname{Spec}(\mathcal{S}_X^{\leq})$와 표준적으로 동형이고,
$E_X$는 $\operatorname{Spec}(\mathcal{S}_X)$와 표준적으로 동형이다.
\end{corollary}

\begin{proof}
사상 $p_X$와 $\pi_X$가 아핀임을 알고 있으므로
(\hyperref[II.8.3.5-ko]{8.3.5} 및
\hyperref[II.8.3.6-ko]{8.3.6}),
(\hyperref[II.1.3.1-ko]{1.3.1})에 따라
$Y=\operatorname{Spec}(A)$가 아핀이고 $\mathcal{S}=\widetilde{S}$인
경우에 따름정리를 확인하면 충분하다. 첫 번째 주장은 표준 동형사상
(\hyperref[II.8.2.7.2-ko]{8.2.7.2})
$(S_f^{\geq})_{f/1}^{\leq}\xrightarrow{\sim}S_f^{\leq}$가 존재하고,
이 동형사상들이 $f$에서 $fg$로 옮겨갈 때 서로 양립한다는 사실에서
따른다($f$와 $g$는 $S_+$의 동차 원소이다). 마찬가지로 $\pi$ 대신
$\pi_X$에 공식 (\hyperref[II.8.3.6.2-ko]{8.3.6.2})를 적용하면,
$S_+$의 동차 원소 $f$에 대하여
$\pi_X^{-1}(u^{-1}(X_f))=
\operatorname{Spec}((S_f^{\geq})_{f/1})$임을 알 수 있다. 두 번째 주장은
표준 동형사상 (\hyperref[II.8.2.7.2-ko]{8.2.7.2})
$(S_f^{\geq})_{f/1}\xrightarrow{\sim}S_f$가 존재한다는 사실에서 따른다.

따라서 $\widehat{C}_X$를 $X$-스킴으로 보면, $X$ 위의 두 아핀
$X$-스킴
$C_X=\operatorname{Spec}(\mathcal{S}_X^{\geq})$와
$\widehat{E}_X=\operatorname{Spec}(\mathcal{S}_X^{\leq})$를
\emph{붙이기}로 얻는다고 말할 수 있다. 이 둘의 교집합은 열린집합
$E_X=\operatorname{Spec}(\mathcal{S}_X)$이다.
\end{proof}

\begin{corollary}[8.6.4]
\label{II.8.6.4-ko}
$\mathcal{O}_X(1)$이 가역 $\mathcal{O}_X$-가군이고
$\mathcal{S}_X$가
$\bigoplus_{n\in\mathbf{Z}}(\mathcal{O}_X(1))^{\otimes n}$과
동형이라고 가정하자(특히 $\mathcal{S}$가 $\mathcal{S}_1$로 생성되면
이 조건이 성립한다
(\hyperref[II.3.2.5-ko]{3.2.5} 및
\hyperref[II.3.2.7-ko]{3.2.7})). 그러면 천공 사영 원뿔
$\widehat{E}$는 $X$ 위의 랭크 1 벡터다발
$\mathbf{V}(\mathcal{O}_X(-1))$과 동일시되고, 천공 아핀 원뿔 $E$는
이 벡터다발의 영 단면의 여집합에 유도된 부분준스킴과 동일시된다.
이 동일시 아래에서 표준 리트랙션 $\widehat{E}\to X$는
$X$-스킴 $\mathbf{V}(\mathcal{O}_X(-1))$의 구조 사상과 동일시된다.
마지막으로 표준 $Y$-사상
$\mathbf{V}(\mathcal{O}_X(1))\to C$가 존재하며, 이 사상을
$\mathbf{V}(\mathcal{O}_X(1))$의 영 단면의 여집합에 제한하면 그
여집합에서 천공 아핀 원뿔 $E$ 위로 가는 동형사상이 된다.
\end{corollary}

\begin{proof}
실제로 $\mathcal{L}=\mathcal{O}_X(1)$로 놓으면,
$\mathcal{S}_X^{\geq}$는
$\mathbf{S}_{\mathcal{O}_X}(\mathcal{L})$과 같으므로,
(\hyperref[II.8.6.3-ko]{8.6.3})에 따라 $\widehat{E}_X$는
$\mathbf{V}(\mathcal{L}^{-1})$과, $C_X$는
$\mathbf{V}(\mathcal{L})$과 표준적으로 동일시된다. 사상
$\mathbf{V}(\mathcal{L})\to C$는
$\operatorname{Proj}(\widehat{\alpha})$의 제한이고, 따라서 따름정리의
주장들은 (\hyperref[II.8.6.2-ko]{8.6.2})의 특수한 경우이다.
\end{proof}

사상 $\mathbf{V}(\mathcal{O}_X(1))\to C$에 의한 $C$의 꼭짓점
준스킴의 바탕공간의 역상은 $\mathbf{V}(\mathcal{O}_X(1))$의 영
단면의 바탕공간이다
(\hyperref[II.8.5.5-ko]{8.5.5})~; 그러나 일반적으로 $C$와
$\mathbf{V}(\mathcal{O}_X(1))$에서 이에 대응하는 부분준스킴들은
동형이 아니다. 이 문제를 아래에서 연구한다.

\subsection{사영 원뿔의 블로업}
\label{subsection:II.8.7-ko}

\begin{env}[8.7.1]
\label{II.8.7.1-ko}
(\hyperref[II.8.6.1-ko]{8.6.1})의 조건 아래에서
$r=\operatorname{Proj}(\widehat{\alpha})$로 놓으면, 다음 가환 그림을
얻는다.
\begin{equation}
\label{II.8.7.1.1-ko}
  \xymatrix{
    X\ar[r]^{i_X\circ\varepsilon_X}\ar[d]_{q} &
    \widehat{C}_X\ar[d]^{r}\\
    Y\ar[r]_{i\circ\varepsilon} & \widehat{C}
  }
\tag{8.7.1.1}
\end{equation}

\oldpage[II]{174}
(\hyperref[II.8.5.1.3-ko]{8.5.1.3})과
(\hyperref[II.8.5.3.2-ko]{8.5.3.2})에 따른 것이다~; 또한 영 단면의
여집합 $\widehat{C}_X-i_X(\varepsilon_X(X))$에 대한 $r$의 제한은
(\hyperref[II.8.6.2-ko]{8.6.2})에 따라 영 단면의 여집합
$\widehat{C}-i(\varepsilon(Y))$ 위의 \emph{동형사상}이다. 간단히 하기
위해 $Y$가 아핀이고, $\mathcal{S}$가 유한형이며 $\mathcal{S}_1$로
생성된다고 가정하면, $X$는 $Y$ 위에서 사영이고 $\widehat{C}_X$는
$X$ 위에서 사영이므로 (\hyperref[II.5.5.1-ko]{5.5.1}),
$\widehat{C}_X$는 $Y$ 위에서 사영이며
(\hyperref[II.5.5.5-ko]{5.5.5, (ii)}), 따라서 \emph{더구나}
$\widehat{C}$ 위에서도 사영이다
(\hyperref[II.5.5.5-ko]{5.5.5, (v)}). 이와 같이 사영 $Y$-사상
$r:\widehat{C}_X\to\widehat{C}$를 얻는다 (그 제한
$C_X\to C$는 사영 $Y$-사상이다). 이 사상은 \emph{$X$를 $Y$로
수축하고}, 이를 \emph{$X$와 $Y$의 여집합들}에 제한하면
\emph{동형사상}을 유도한다. 따라서 $C_X$와 $C$ 사이에는 블로업
준스킴과 그 출발 준스킴 사이의 관계
(\hyperref[II.8.1.3-ko]{8.1.3})와 유사한 관계가 있다. 실제로
$C_X$를 어떤 등급 $\mathcal{O}_C$-대수의 동차 스펙트럼과 동일시할
수 있음을 보이겠다.
\end{env}

\begin{env}[8.7.2]
\label{II.8.7.2-ko}
(\hyperref[II.8.6.1-ko]{8.6.1})의 표기를 그대로 두고, 각
$n\geqslant 0$에 대하여 등급 $\mathcal{O}_Y$-대수 $\mathcal{S}$의
준연접 아이디얼
\begin{equation}
\label{II.8.7.2.1-ko}
  \mathcal{S}_{[n]}=\bigoplus_{m\geqslant n}\mathcal{S}_m
\tag{8.7.2.1}
\end{equation}
을 생각하자. 다음은 명백하다.
\begin{equation}
\label{II.8.7.2.2-ko}
  \mathcal{S}_{[0]}=\mathcal{S},\qquad
  \mathcal{S}_{[n]}\subset\mathcal{S}_{[m]}
  \qquad\text{$m\leqslant n$일 때}
\tag{8.7.2.2}
\end{equation}
\begin{equation}
\label{II.8.7.2.3-ko}
  \mathcal{S}_{[n]}\mathcal{S}_{[m]}\subset\mathcal{S}_{[m+n]}.
\tag{8.7.2.3}
\end{equation}

$\mathcal{S}_{[n]}$에 결부된 $\mathcal{O}_C$-가군을 생각하자. 이는
$\mathcal{O}_C=\widetilde{\mathcal{S}}$의 준연접 아이디얼이며
(\hyperref[II.1.4.4-ko]{1.4.4}), 다음과 같이 쓴다.
\begin{equation}
\label{II.8.7.2.4-ko}
  \mathcal{I}_n=(\mathcal{S}_{[n]})^{\sim}.
\tag{8.7.2.4}
\end{equation}

그러면 (\hyperref[II.8.7.2.2-ko]{8.7.2.2})와
(\hyperref[II.8.7.2.3-ko]{8.7.2.3})로부터,
(\hyperref[II.1.4.4-ko]{1.4.4})와
(\hyperref[II.1.4.8.1-ko]{1.4.8.1})을 사용하여 다음과 같은 유사한
공식들을 얻는다.
\begin{equation}
\label{II.8.7.2.5-ko}
  \mathcal{I}_0=\mathcal{O}_C,\qquad
  \mathcal{I}_n\subset\mathcal{I}_m
  \qquad\text{$m\leqslant n$일 때}
\tag{8.7.2.5}
\end{equation}
\begin{equation}
\label{II.8.7.2.6-ko}
  \mathcal{I}_n\mathcal{I}_m\subset\mathcal{I}_{n+m}.
\tag{8.7.2.6}
\end{equation}

따라서 (\hyperref[II.8.1.1-ko]{8.1.1})의 조건들이 성립하며, 이에
따라 다음 준연접 등급 $\mathcal{O}_C$-대수를 도입한다.
\begin{equation}
\label{II.8.7.2.7-ko}
  \mathcal{S}^{\natural}
  =\bigoplus_{n\geqslant 0}\mathcal{I}_n
  =\left(\bigoplus_{n\geqslant 0}\mathcal{S}_{[n]}\right)^{\sim}.
\tag{8.7.2.7}
\end{equation}
\end{env}

\begin{proposition}[8.7.3]
\label{II.8.7.3-ko}
다음 표준 $C$-동형사상이 존재한다.
\begin{equation}
\label{II.8.7.3.1-ko}
  h:C_X\xrightarrow{\sim}\operatorname{Proj}(\mathcal{S}^{\natural}).
\tag{8.7.3.1}
\end{equation}
\end{proposition}

\begin{proof}
먼저 $Y=\operatorname{Spec}(A)$가 아핀이라고 가정하자. 그러면
$\mathcal{S}=\widetilde{S}$이며, 여기서 $S$는 양의 차수로 등급화된
$A$-대수이고 $C=\operatorname{Spec}(S)$이다. 정의
(\hyperref[II.8.2.6.2-ko]{8.2.7.4})에 따라
(\hyperref[II.8.2.6-ko]{8.2.6})의 표기 아래
$\mathcal{S}^{\natural}=(S^{\natural})^{\sim}$이다.
(\hyperref[II.8.7.3.1-ko]{8.7.3.1})을 정의하기 위해 동차 원소
$f\in S_d$ ($d>0$)와 이에 대응하는 원소
$f^{\natural}\in S^{\natural}$
(\hyperref[II.8.2.6-ko]{8.2.6})을 생각하자~; 그러면 $S$-동형사상
(\hyperref[II.8.2.7.3-ko]{8.2.7.3})은 다음 $C$-동형사상을 정의한다.
\begin{equation}
\label{II.8.7.3.2-ko}
  \operatorname{Spec}(S_f^{\geqslant})
  \xrightarrow{\sim}
  \operatorname{Spec}(S_{(f^{\natural})}^{\natural}).
\tag{8.7.3.2}
\end{equation}

\oldpage[II]{175}
그런데 (\hyperref[II.8.6.2-ko]{8.6.2})의 표기에서
$v:C_X\to X$가 구조 사상이면,
(\hyperref[II.8.6.2.1-ko]{8.6.2.1})에 의해
$v^{-1}(X_f)=\operatorname{Spec}(S_f^{\geqslant})$이다. 한편
$\operatorname{Spec}(S_{(f^{\natural})}^{\natural})=D_+(f^{\natural})$이므로,
(\hyperref[II.8.7.3.2-ko]{8.7.3.2})는 동형
$v^{-1}(X_f)\to D_+(f^{\natural})$을 정의한다. 또한 $g\in S_e$ ($e>0$)이면
도식
\[
  \xymatrix{
    v^{-1}(X_{fg})\ar[r]^{\sim}\ar[d] &
    D_+(f^{\natural}g^{\natural})\ar[d]\\
    v^{-1}(X_f)\ar[r]^{\sim} & D_+(f^{\natural})
  }
\]
은 가환하는데, 이는 동형
(\hyperref[II.8.2.7.3-ko]{8.2.7.3})의 정의에서 즉시 따른다. 끝으로 정의에 의해
$S_+$는 동차 원소 $f$들로 생성되므로,
(\hyperref[II.8.2.10-ko]{8.2.10, (iv)})와
(\hyperref[II.2.3.14-ko]{2.3.14})에 의해 $D_+(f^{\natural})$들은
$\operatorname{Proj}(S^{\natural})$의 덮개를 이루고 $v^{-1}(X_f)$들은
$C_X$의 덮개를 이룬다. 실제로 $X_f$들이 $X$의 덮개를 이루기 때문이다~;
따라서 이 경우의 동형 (\hyperref[II.8.7.3.1-ko]{8.7.3.1})의 정의가
완성된다.

(\hyperref[II.8.7.3-ko]{8.7.3})을 일반적인 경우에 증명하기 위해서는,
$U$, $U'$가 $Y$의 아핀 열린집합이고 $U'\subset U$이며 그 환들이 각각
$A$, $A'$일 때, 그리고
$\mathcal{S}|U=\widetilde{S}$, $\mathcal{S}|U'=\widetilde{S'}$로 놓을 때 도식
\begin{equation}
\label{II.8.7.3.3-ko}
  \xymatrix{
    C_{U'}\ar[r]\ar[d] & \operatorname{Proj}(S'^{\natural})\ar[d]\\
    C_U\ar[r] & \operatorname{Proj}(S^{\natural})
  }
\tag{8.7.3.3}
\end{equation}
이 가환함을 보이면 충분하다. 그런데 $S'$는 $S\otimes_A A'$와 표준적으로
동일시되고, 따라서 $S'^{\natural}$은
\[
  S^{\natural}\otimes_S S'=S^{\natural}\otimes_A A'~;
\]
와 표준적으로 동일시된다. 그러므로
$\operatorname{Proj}(S'^{\natural})=\operatorname{Proj}(S^{\natural})
\times_U U'$이다 (\hyperref[II.2.8.10-ko]{2.8.10})~; 마찬가지로
$X=\operatorname{Proj}(S)$, $X'=\operatorname{Proj}(S')$이면,
$X'=X\times_U U'$이고
$\mathcal{S}_{X'}=\mathcal{S}_X\otimes_{\mathcal{O}_U}\mathcal{O}_{U'}$
이다 (\hyperref[II.3.5.4-ko]{3.5.4}). 다시 말해
$\mathcal{S}_{X'}=j^*(\mathcal{S}_X)$이며, 여기서 $j$는 사영 $X'\to X$이다.
따라서 (\hyperref[II.1.5.2-ko]{1.5.2})에 의해
$C_{U'}=C_U\times_X X'=C_U\times_U U'$이고,
(\hyperref[II.8.7.3.3-ko]{8.7.3.3})의 가환성은 이제 자명하다.
\end{proof}

\begin{remark}[8.7.4]
\label{II.8.7.4-ko}
\begin{enumerate}
  \item[\rm{(i)}] (\hyperref[II.8.7.3-ko]{8.7.3})의 논증 마지막 부분은
  다음과 같이 즉시 일반화된다. $g:Y'\to Y$를 사상이라 하고,
  $\mathcal{S}'=g^*(\mathcal{S})$, $X'=\operatorname{Proj}(\mathcal{S}')$로
  놓자~; 그러면 가환 도식
  \begin{equation}
  \label{II.8.7.4.1-ko}
    \xymatrix{
      C_{X'}\ar[r]\ar[d] &
      \operatorname{Proj}(\mathcal{S}'^{\natural})\ar[d]\\
      C_X\ar[r] & \operatorname{Proj}(\mathcal{S}^{\natural})
    }
  \tag{8.7.4.1}
  \end{equation}
  을 얻는다.

  한편 $\varphi:\mathcal{S}''\to\mathcal{S}$를 등급
  $\mathcal{O}_Y$-대수의 준동형이라 하자. 이때
  $X''=\operatorname{Proj}(\mathcal{S}'')$로 놓으면
  $u=\operatorname{Proj}(\varphi):X\to X''$가 모든 점에서 정의된다고
  가정한다~; 다른 한편

  \oldpage[II]{176}
  $\mathcal{A}(v)=\varphi$를 만족하는 $Y$-사상 $v:C\to C''$
  ($C''=\operatorname{Spec}(\mathcal{S}'')$)를 얻는다. 또한
  $\varphi$는 등급 대수의 준동형이므로, $\varphi$로부터 등급 대수의
  $v$-준동형
  $\psi:\mathcal{S}''^{\natural}\to\mathcal{S}^{\natural}$
  을 얻는다(\hyperref[II.1.4.1-ko]{1.4.1}). 더구나
  (\hyperref[II.8.2.10-ko]{8.2.10, (iv)})과 $\varphi$에 대한 가정으로부터
  $\operatorname{Proj}(\psi)$가 모든 곳에서 정의된다는 것이 쉽게 따른다.
  마지막으로 (\hyperref[II.3.5.6.1-ko]{3.5.6.1})을 고려하면 표준
  $u$-준동형 $\mathcal{S}_{X''}\to\mathcal{S}_X$가 있으므로,
  (\hyperref[II.1.5.6-ko]{1.5.6})에 의해 사상 $w:C_{X''}\to C_X$를 얻는다.
  이때 도표
  \begin{equation}
  \label{II.8.7.4.2-ko}
    \xymatrix{
      C_{X''}\ar[r]^{\sim}\ar[d]_{w} &
      \operatorname{Proj}(\mathcal{S}''^{\natural})
        \ar[d]^{\operatorname{Proj}(\psi)}\\
      C_X\ar[r]^{\sim} & \operatorname{Proj}(\mathcal{S}^{\natural})
    }
  \tag{8.7.4.2}
  \end{equation}
  는 가환이다. 이는 $Y$가 아핀인 경우로 환원하여 곧바로 확인할 수 있다.

  \item[\rm{(ii)}] (\hyperref[II.8.7.2.5-ko]{8.7.2.5})와
  (\hyperref[II.8.7.2.6-ko]{8.7.2.6})에 의해 모든 $m>0$에 대하여
  $\mathcal{I}_1^m\subset\mathcal{I}_m\subset\mathcal{I}_1$임을 주목하자.
  그런데 정의에 따라 $\mathcal{I}_1=(\mathcal{S}_+)^{\sim}$이므로,
  $\mathcal{I}_1$은 $C$에서 닫힌 부분준스킴 $\varepsilon(Y)$를 정의한다
  ((\hyperref[II.1.4.10-ko]{1.4.10}) 및
  (\hyperref[II.8.3.2-ko]{8.3.2})). 따라서 모든 $m>0$에 대하여
  \emph{$\mathcal{O}_C/\mathcal{I}_m$의 지지집합은 꼭지점 부분준스킴
  $\varepsilon(Y)$의 바탕공간에 포함된다}. 그러므로 꼭지점을 제거한
  아핀 원뿔 $E$의 역상에서는 구조 사상
  $\operatorname{Proj}(\mathcal{S}^{\natural})\to C$가
  \emph{동형사상}으로 제한된다(이는
  (\hyperref[II.8.7.3-ko]{8.7.3})과
  (\hyperref[II.8.7.1-ko]{8.7.1})에서도 따른다). 또한 $C$를
  $\widehat{C}$의 열린집합과 표준적으로 동일시하면
  (\hyperref[II.8.3.3-ko]{8.3.3}), $\mathcal{O}_C$의 아이디얼
  $\mathcal{I}_m$을 $\mathcal{O}_{\widehat{C}}$의 아이디얼
  $\mathcal{J}_m$으로 당연히 연장할 수 있다. 그 조건은
  $\widehat{C}$의 열린집합 $\widehat{E}$에서
  $\mathcal{O}_{\widehat{C}}$와 일치하는 것이다. 다음과 같이 놓자.
  $\mathcal{T}=\bigoplus_{n\geqslant 0}\mathcal{J}_m$.
  이는 준연접 등급 $\mathcal{O}_{\widehat{C}}$-대수이며, 동형사상
  (\hyperref[II.8.7.3.1-ko]{8.7.3.1})을 다음
  $\widehat{C}$-동형사상으로 연장할 수 있다.
  \begin{equation}
  \label{II.8.7.4.3-ko}
    \widehat{C}_X\xrightarrow{\sim}\operatorname{Proj}(\mathcal{T}).
  \tag{8.7.4.3}
  \end{equation}

  실제로 $\widehat{E}$ 위에서 앞의 결과에 의해
  $\operatorname{Proj}(\mathcal{T})$는 $\widehat{E}$와 표준적으로
  동일시된다. 따라서 $\widehat{E}$ 위의 동형사상
  (\hyperref[II.8.7.4.3-ko]{8.7.4.3})은 표준 동형사상
  $\widehat{E}_X\to\widehat{E}$
  (\hyperref[II.8.6.2-ko]{8.6.2})과 일치하도록 정의한다. 그러면 이
  동형사상과 (\hyperref[II.8.7.3.1-ko]{8.7.3.1})이 $E$ 위에서
  일치함은 명백하다.
\end{enumerate}
\end{remark}

\begin{corollary}[8.7.5]
\label{II.8.7.5-ko}
어떤 $n_0>0$가 존재하여 다음이 성립한다고 가정하자.
\begin{equation}
\label{II.8.7.5.1-ko}
  \mathcal{S}_{n+1}=\mathcal{S}_1\mathcal{S}_n
  \qquad\text{$n\geqslant n_0$일 때}.
\tag{8.7.5.1}
\end{equation}
그러면 $C_X$의 꼭지점 부분준스킴($X$와 동형)은 표준 사상
$r:C_X\to C$에 의한 $C$의 꼭지점 부분준스킴($Y$와 동형)의 역상이다.
역으로 이 성질이 성립하고, 더 나아가 $Y$가 뇌터이고 $\mathcal{S}$가
유한형이라고 가정하면, (\hyperref[II.8.7.5.1-ko]{8.7.5.1})이 성립하는
어떤 $n_0>0$가 존재한다.
\end{corollary}

\begin{proof}
첫째 주장은 $Y$ 위에서 국소적인 명제이므로, 이를 증명할 때에는
$Y=\operatorname{Spec}(A)$가 아핀이고 $\mathcal{S}=\widetilde{S}$이며,
$S$가 양의 차수로 등급이 주어진 $A$-대수라고 가정할 수 있다. 그러면
(\hyperref[II.8.2.12-ko]{8.2.12})에서 주장이 따른다. 실제로
$\operatorname{Proj}(S^{\natural}\otimes_S S_0)=
C_X\times_C\varepsilon(Y)$이고(동일시
(\hyperref[II.8.7.3.1-ko]{8.7.3.1})에 의하여), 다시 말해 이 준스킴은
$C_X$에서 $\varepsilon(Y)$의 역상이다
(\hyperref[I.4.4.1-ko]{I, 4.4.1}). $Y$가 아핀 뇌터이고 $S$가
유한형인 경우에는 그 역도 (\hyperref[II.8.2.12-ko]{8.2.12})에서 따른다.

\oldpage[II]{177}
$Y$가 뇌터이고(반드시 아핀일 필요는 없다) $\mathcal{S}$가 유한형이면,
$Y$는 유한 개의 아핀 뇌터 열린집합 $U_i$로 덮인다. 따라서 앞의
결과로부터 각 $i$에 대하여 정수 $n_i$가 존재하여, $n\geqslant n_i$이면
$\mathcal{S}_{n+1}|U_i=(\mathcal{S}_1|U_i)(\mathcal{S}_n|U_i)$임을 얻는다.
그러므로 $n_i$들 가운데 가장 큰 것은
(\hyperref[II.8.7.5.1-ko]{8.7.5.1})을 만족한다.
\end{proof}

\begin{env}[8.7.6]
\label{II.8.7.6-ko}
이제 아핀 원뿔 $C$에서 \emph{꼭지점 부분준스킴}
$\varepsilon(Y)$을 따라 \emph{블로업하여} 얻은 $C$-준스킴 $Z$를
생각하자. 정의에 따라 (\hyperref[II.8.1.3-ko]{8.1.3}) 이는 준스킴
$\operatorname{Proj}(\bigoplus_{n\geqslant 0}\mathcal{S}_+^n)$이다.
표준 단사 준동형
\begin{equation}
\label{II.8.7.6.1-ko}
  \iota:\bigoplus_{n\geqslant 0}\mathcal{S}_+^n
  \longrightarrow\mathcal{S}^{\natural}
\tag{8.7.6.1}
\end{equation}
은 동일시 (\hyperref[II.8.7.3-ko]{8.7.3})에 의해 표준 우세
$C$-사상
\begin{equation}
\label{II.8.7.6.2-ko}
  G(\iota)\longrightarrow Z
\tag{8.7.6.2}
\end{equation}
을 정의한다. 여기서 $G(\iota)$는 $C_X$의 열린집합이다
(\hyperref[II.3.5.1-ko]{3.5.1}). $G(\iota)\ne C_X$일 수도 있음에
유의하자. 이는 다음 예에서 드러난다. $Y=\operatorname{Spec}(K)$이고
$K$가 체이며, $\mathcal{S}=\widetilde{S}$이고 $S=K[\mathbf{y}]$라 하자.
여기서 $\mathbf{y}$는 \emph{차수가 2인} 부정원이다. $R_n$이 집합
$(S_+)^n$을 나타내고 이를 $S_{[n]}=S_n^{\natural}$의 부분집합으로
간주하면, $S_+^{\natural}$은 $R_n$들의 합집합이 생성하는 아이디얼의
$S^{\natural}$에서의 근기가 아니다(\hyperref[II.2.3.14-ko]{2.3.14} 참조).
\end{env}

\begin{corollary}[8.7.7]
\label{II.8.7.7-ko}
어떤 $n_0>0$가 존재하여 다음이 성립한다고 가정하자.
\begin{equation}
\label{II.8.7.7.1-ko}
  \mathcal{S}_n=\mathcal{S}_1^n
  \qquad\text{$n\geqslant n_0$일 때}.
\tag{8.7.7.1}
\end{equation}
그러면 표준 사상 (\hyperref[II.8.7.6.2-ko]{8.7.6.2})은 모든 곳에서
정의되고 동형사상 $C_X\xrightarrow{\sim}Z$이다. 역으로 이 성질이
성립하고, 더 나아가 $Y$가 뇌터이고 $\mathcal{S}$가 유한형이라고
가정하면, (\hyperref[II.8.7.7.1-ko]{8.7.7.1})이 성립하는 $n_0$가
존재한다.
\end{corollary}

\begin{proof}
실제로 첫째 주장은 $Y$ 위에서 국소적이므로
(\hyperref[II.8.2.14-ko]{8.2.14})에서 따른다. 같은 방식으로
(\hyperref[II.8.7.5-ko]{8.7.5})에서처럼 논증하면 그 역도 따른다.
\end{proof}

\begin{remark}[8.7.8]
\label{II.8.7.8-ko}
조건 (\hyperref[II.8.7.7.1-ko]{8.7.7.1})은
(\hyperref[II.8.7.5.1-ko]{8.7.5.1})을 함의한다. 따라서 이 조건이
성립할 때 $C_X$는 아핀 원뿔 $C$의 꼭지점($Y$와 동일시된다)을
블로업하여 얻은 준스킴과 동일시될 뿐 아니라, $C_X$의 꼭지점
($X$와 동일시된다)은 $C$의 꼭지점 $Y$의 역상인 닫힌 부분준스킴과
동일시됨을 알 수 있다. 더구나 가정
(\hyperref[II.8.7.7.1-ko]{8.7.7.1})은
$X=\operatorname{Proj}(\mathcal{S})$ 위의 $\mathcal{O}_X$-가군
$\mathcal{O}_X(n)$들이 가역임을 함의하며
(\hyperref[II.3.2.5-ko]{3.2.5} 및
\hyperref[II.3.2.9-ko]{3.2.9}), 다음이 성립한다.
$\mathcal{O}_X(n)=\mathcal{L}^{\otimes n}$, 여기서
$\mathcal{L}=\mathcal{O}_X(1)$이다
(\hyperref[II.3.2.7-ko]{3.2.7} 및
\hyperref[II.3.2.9-ko]{3.2.9}). 정의
(\hyperref[II.8.6.1.1-ko]{8.6.1.1})에 따라 $C_X$는 $X$ 위의
\emph{벡터 다발} $\mathbf{V}(\mathcal{L})$이고, 그 꼭지점은 이
벡터 다발의 \emph{영 단면}이다.
\end{remark}

\subsection{풍부한 층과 수축}
\label{subsection:II.8.8-ko}

\begin{env}[8.8.1]
\label{II.8.8.1-ko}
$Y$를 준스킴, $f:X\to Y$를 \emph{분리된 준콤팩트} 사상,
$\mathcal{L}$을 $f$에 관해 \emph{풍부한} 가역
$\mathcal{O}_X$-가군이라 하자. 양의 차수로 등급화된 다음
$\mathcal{O}_Y$-대수를 생각하자.
\begin{equation}
\label{II.8.8.1.1-ko}
  \mathcal{S}=\mathcal{O}_Y\oplus
  \bigoplus_{n\geqslant 1}f_*(\mathcal{L}^{\otimes n})
\tag{8.8.1.1}
\end{equation}

\oldpage[II]{178}
이는 준연접이다 (\hyperref[I.9.2.2-ko]{I, 9.2.2, a)}). 등급
$\mathcal{O}_X$-대수의 표준 준동형
\begin{equation}
\label{II.8.8.1.2-ko}
  \tau:f^*(\mathcal{S})\longrightarrow
  \bigoplus_{n\geqslant 0}\mathcal{L}^{\otimes n}
\tag{8.8.1.2}
\end{equation}
이 존재한다. 이는 차수 $\geqslant 1$에서는 표준 준동형
$\sigma:f^*(f_*(\mathcal{L}^{\otimes n}))\to\mathcal{L}^{\otimes n}$
(\hyperref[0.4.4.3-ko]{0, 4.4.3})과 일치하고, 차수 $0$에서는 항등사상이다.
$\mathcal{L}$이 $f$에 관해 풍부하다는 가정으로부터
(\hyperref[II.4.6.3-ko]{4.6.3} 및 \hyperref[II.3.6.1-ko]{3.6.1})
이에 대응하는 $Y$-사상
\begin{equation}
\label{II.8.8.1.3-ko}
  r=r_{\mathcal{L},\tau}:X\longrightarrow
  P=\operatorname{Proj}(\mathcal{S})
\tag{8.8.1.3}
\end{equation}
은 모든 곳에서 정의되며 \emph{우세 열린 몰입 사상}이고, 다음이 성립한다.
\begin{equation}
\label{II.8.8.1.4-ko}
  r^*(\mathcal{O}_P(n))=\mathcal{L}^{\otimes n}
  \qquad\text{모든 }n\in\mathbf{Z}\text{에 대하여}.
\tag{8.8.1.4}
\end{equation}
\end{env}

\begin{proposition}[8.8.2]
\label{II.8.8.2-ko}
$C=\operatorname{Spec}(\mathcal{S})$를 \emph{$\mathcal{S}$로 정의되는
아핀 원뿔}이라 하자. $\mathcal{L}$이 $f$에 관해 풍부하면 표준
$Y$-사상
\begin{equation}
\label{II.8.8.2.1-ko}
  g:V=\mathbf{V}(\mathcal{L})\longrightarrow C
\tag{8.8.2.1}
\end{equation}
이 존재하여 도식
\begin{equation}
\label{II.8.8.2.2-ko}
  \xymatrix{
    X\ar[r]^{j}\ar[d]_{f} &
    \mathbf{V}(\mathcal{L})\ar[r]^{\pi}\ar[d]^{g} & X\ar[d]^{f}\\
    Y\ar[r]_{\varepsilon} & C\ar[r]_{\psi} & Y
  }
\tag{8.8.2.2}
\end{equation}
은 가환한다. 여기서 $\psi$와 $\pi$는 구조 사상이고, $j$와
$\varepsilon$은 각각 $X$와 $Y$를 $\mathbf{V}(\mathcal{L})$의 영 단면과
$C$의 꼭짓점 준스킴 위로 보내는 표준 몰입 사상이다. 또한
$\mathbf{V}(\mathcal{L})-j(X)$에 대한 $g$의 제한은
\begin{equation}
\label{II.8.8.2.3-ko}
  \mathbf{V}(\mathcal{L})-j(X)\longrightarrow
  E=C-\varepsilon(Y)
\tag{8.8.2.3}
\end{equation}
으로 주어지는 열린 몰입 사상이며, 여기서 $E$는 $\mathcal{S}$에
대응하는 천공 아핀 원뿔이다.
\end{proposition}

\begin{proof}
(\hyperref[II.8.8.1-ko]{8.8.1})의 표기법을 사용하여
$\mathcal{S}_P^{\geqslant}=\bigoplus_{n\geqslant 0}\mathcal{O}_P(n)$,
$C_P=\operatorname{Spec}(\mathcal{S}_P^{\geqslant})$로 놓자.
(\hyperref[II.8.6.2-ko]{8.6.2})에 의해 가환 도식
\begin{equation}
\label{II.8.8.2.4-ko}
  \xymatrix{
    C_P\ar[r]\ar[d]_{h} & P\ar[d]^{p}\\
    C\ar[r]_{\psi} & Y
  }
\tag{8.8.2.4}
\end{equation}
을 이루는 표준 사상 $h=\operatorname{Spec}(\alpha):C_P\to C$가 존재한다.
또한 $\varepsilon_P:P\to C_P$가 표준 몰입 사상이면 도식
\begin{equation}
\label{II.8.8.2.5-ko}
  \xymatrix{
    P\ar[r]^{\varepsilon_P}\ar[d]_{p} & C_P\ar[d]^{h}\\
    Y\ar[r]_{\varepsilon} & C
  }
\tag{8.8.2.5}
\end{equation}
은 가환하고 (\hyperref[II.8.7.1.1-ko]{8.7.1.1}), $h$를 천공 아핀
원뿔 $E_P$에 제한한 사상은 \emph{동형사상}
$E_P\xrightarrow{\sim}E$이다 (\hyperref[II.8.6.2-ko]{8.6.2}). 한편
(\hyperref[II.8.8.1.4-ko]{8.8.1.4})로부터 다음을 얻는다.
\[
  r^*(\mathcal{S}_P^{\geqslant})=
  \mathbf{S}_{\mathcal{O}_X}(\mathcal{L})
\]

\oldpage[II]{179}
따라서 표준 $P$-사상 $q:\mathbf{V}(\mathcal{L})\to C_P$가 존재하고,
가환 도식
\begin{equation}
\label{II.8.8.2.6-ko}
  \xymatrix{
    \mathbf{V}(\mathcal{L})\ar[r]^{\pi}\ar[d]_{q} & X\ar[d]^{r}\\
    C_P\ar[r] & P
  }
\tag{8.8.2.6}
\end{equation}
은 $\mathbf{V}(\mathcal{L})$을 곱 $C_P\times_P X$와 동일시한다
(\hyperref[II.1.5.2-ko]{1.5.2}). $r$이 열린 몰입 사상이므로 $q$도
열린 몰입 사상이다 (\hyperref[I.4.3.2-ko]{I, 4.3.2}). 또한
(\hyperref[II.8.5.2-ko]{8.5.2})에 의해
$\mathbf{V}(\mathcal{L})-j(X)$에 대한 $q$의 제한은 이 준스킴을
$E_P$ 안으로 보내고, 도식
\begin{equation}
\label{II.8.8.2.7-ko}
  \xymatrix{
    X\ar[r]^{j}\ar[d]_{r} & \mathbf{V}(\mathcal{L})\ar[d]^{q}\\
    P\ar[r]_{\varepsilon_P} & C_P
  }
\tag{8.8.2.7}
\end{equation}
은 가환한다(이는
(\hyperref[II.8.5.1.3-ko]{8.5.1.3})의 특수한 경우이다).
$g$를 합성 사상 $h\circ q$로 취하면, 이 사실들로부터
(\hyperref[II.8.8.2-ko]{8.8.2})의 주장들이 곧바로 따른다.
\end{proof}

\begin{remark}[8.8.3]
\label{II.8.8.3-ko}
더 나아가 $Y$가 \emph{뇌터} 준스킴이고 $f$가 \emph{고유} 사상이라고
가정하자. 그러면 $r$도 \emph{고유}이고
(\hyperref[II.5.4.4-ko]{5.4.4}), 따라서 닫힌 사상이다. 한편 $r$은
우세 열린 몰입 사상이므로 반드시 \emph{동형사상}
$X\xrightarrow{\sim}P$이다. 또한 제III장에서
(\agref{III.2.3.5.1-ko}{III, 2.3.5.1}) $\mathcal{S}$가 반드시
\emph{유한형} $\mathcal{O}_Y$-대수임을 보일 것이다. 따라서
$\mathcal{S}^{\natural}$은 $\mathcal{S}_0^{\natural}$ 위의
\emph{유한형} 대수이다
(\hyperref[II.8.2.10-ko]{8.2.10, (i)} 및
\hyperref[II.8.7.2.7-ko]{8.7.2.7}). $C_P$는
$\operatorname{Proj}(\mathcal{S}^{\natural})$과 $C$ 위에서 동형이므로
(\hyperref[II.8.7.3-ko]{8.7.3}), 사상 $h:C_P\to C$는
\emph{사영적}이다. 사상 $r$이 동형사상이므로
$q:\mathbf{V}(\mathcal{L})\to C_P$도 동형사상이고, 이에 따라
$g:\mathbf{V}(\mathcal{L})\to C$가 \emph{사영적}임을 얻는다. 또한
$E_P$에 대한 $h$의 제한은 $E$ 위로 가는 동형사상이고 $q$도
동형사상이므로, $g$의 제한
(\hyperref[II.8.8.2.3-ko]{8.8.2.3})은 \emph{동형사상}
$\mathbf{V}(\mathcal{L})-j(X)\xrightarrow{\sim}E$이다.

더 나아가 $\mathcal{L}$이 $f$에 관해 \emph{매우 풍부}하다고 가정하자.
제III장에서 (\agref{III.2.3.5.1-ko}{III, 2.3.5.1}) 어떤 정수
$n_0>0$이 존재하여 $n\geqslant n_0$이면
$\mathcal{S}_n=\mathcal{S}_1^n$임도 보일 것이다. 이로부터
(\hyperref[II.8.7.7-ko]{8.7.7})에 의해 $\mathbf{V}(\mathcal{L})$은
\emph{아핀 원뿔 $C$ 안의 꼭짓점 부분준스킴을 블로업하여} 얻은
준스킴 $Z$와 동일시되고(이 부분준스킴은 $Y$와 동일시된다),
$\mathbf{V}(\mathcal{L})$의 \emph{영 단면}($Y$와 동일시된다)은
$C$의 꼭짓점 부분준스킴 $Y$의 \emph{역상}임을 얻을 것이다. 앞의
결과 중 일부는 실제로 뇌터 가정 없이도 확립할 수 있다~:
\end{remark}

\begin{corollary}[8.8.4]
\label{II.8.8.4-ko}
$Y$를 준스킴(각각 준콤팩트 스킴), $f:X\to Y$를 고유 사상,
$\mathcal{L}$을 $f$에 관해 풍부한 가역 $\mathcal{O}_X$-가군이라 하자.
그러면 사상 (\hyperref[II.8.8.2.1-ko]{8.8.2.1})은 고유하고(각각
사영적이고), 그 제한 (\hyperref[II.8.8.2.3-ko]{8.8.2.3})은
동형사상이다.
\end{corollary}

\begin{proof}
$g$가 고유임을 증명하기 위해 $Y$가 아핀인 경우로 환원할 수 있으며,
따라서 $Y$가 준콤팩트 스킴인 경우만 생각하면 충분하다.
(\hyperref[II.8.8.3-ko]{8.8.3})과 같은 논증에 의해 먼저 $r$이
\emph{동형사상} $X\xrightarrow{\sim}P$임을 알 수 있다. 따라서 $q$도
동형사상이고, $E_P$에 대한 $h$의 제한이 동형사상
$E_P\xrightarrow{\sim}E$이므로
(\hyperref[II.8.8.2.3-ko]{8.8.2.3})이 동형사상임을 이미 알 수 있다.
이제 $g$가 \emph{사영적}임을 보이면 된다.

$f$는 가정에 의해 유한형이므로, (\hyperref[II.3.8.5-ko]{3.8.5})를
다음 준동형에 적용할 수 있다~:

\oldpage[II]{180}
(\hyperref[II.8.8.1.2-ko]{8.8.1.2})의 $\tau$에 적용할 수 있다~: 따라서
$d>0$인 정수와 $\mathcal{S}_d$의 유한형 준연접
$\mathcal{O}_Y$-부분가군 $\mathcal{E}$가 존재하여, $\mathcal{S}'$을
$\mathcal{E}$가 생성하는 $\mathcal{S}$의 $\mathcal{O}_Y$-부분대수라 하고
$\tau'=\tau\circ f^*(\varphi)$ ($\varphi$는 표준 단사 준동형
$\mathcal{S}'\to\mathcal{S}$)라 하면, $r'=r_{\mathcal{L},\tau'}$은 몰입사상
\[
  X\longrightarrow P'=\operatorname{Proj}(\mathcal{S}').
\]
이 된다. 더구나 $\varphi$가 단사이므로 $r'$은 또한 \emph{우세
몰입사상}이다 (\hyperref[II.3.7.6-ko]{3.7.6})~; 따라서 $r$에 대한 것과
같은 논법으로 $r'$이 \emph{전사 닫힌 몰입사상}임을 알 수 있다. $r'$은
$X\xrightarrow{r}\operatorname{Proj}(\mathcal{S})
\xrightarrow{\Phi}\operatorname{Proj}(\mathcal{S}')$로 분해되고, 여기서
$\Phi=\operatorname{Proj}(\varphi)$이므로, $\Phi$도 \emph{전사 닫힌
몰입사상}이라고 결론 내릴 수 있다. 그런데 이로부터 $\Phi$가
\emph{동형사상}임이 따른다~; 실제로
$Y=\operatorname{Spec}(A)$가 아핀이고, $\mathcal{S}=\widetilde{S}$,
$\mathcal{S}'=\widetilde{S'}$이며, $S$는 등급 $A$-대수, $S'$은 $S$의
등급 부분대수인 경우로 귀착시킬 수 있다. 임의의 동차 원소 $t\in S'$에
대하여 $S'_{(t)}$는 $S_{(t)}$의 부분환이다~; 따라서
$\operatorname{Proj}(\varphi)$의 정의
(\hyperref[II.2.8.1-ko]{2.8.1})로 돌아가면, 환 $B$의 부분환을 $B'$이라
할 때 표준 단사 준동형 $B'\to B$에 대응하는 사상
$\operatorname{Spec}(B)\to\operatorname{Spec}(B')$가 닫힌 몰입사상이면
이 사상이 반드시 \emph{동형사상}임을 증명하는 문제로 귀착된다~; 그런데
이는 (\hyperref[I.4.2.3-ko]{I, 4.2.3})에서 따른다.

또한
$\Phi^*(\mathcal{O}_{P'}(n))=\mathcal{O}_P(n)$
(\hyperref[II.3.5.2-ko]{3.5.2, (ii)} 및
\hyperref[II.3.5.4-ko]{3.5.4})이므로,
$r'^*(\mathcal{O}_{P'}(n))$은 $\mathcal{L}^{\otimes n}$과 동형이다
(\hyperref[II.4.6.3-ko]{4.6.3}). $\mathcal{S}''=\mathcal{S}'^{(d)}$라
놓자. 그러면 (\hyperref[II.3.1.8-ko]{3.1.8, (i)})에 따라 $X$는
$P''=\operatorname{Proj}(\mathcal{S}'')$와 표준적으로 동일시되고,
$\mathcal{L}''=\mathcal{L}^{\otimes d}$는 $\mathcal{O}_{P''}(1)$과
표준적으로 동일시된다 (\hyperref[II.3.2.9-ko]{3.2.9, (ii)}).

이제 $C''=\operatorname{Spec}(\mathcal{S}'')$라 하면,
$\mathcal{S}_{P''}^{\geqslant}=
\bigoplus_{n\geqslant 0}\mathcal{O}_{P''}(n)$은
$\bigoplus_{n\geqslant 0}\mathcal{L}''^{\otimes n}$과 동일시되므로,
$C_{P''}=\operatorname{Spec}(\mathcal{S}_{P''}^{\geqslant})$은
$\mathbf{V}(\mathcal{L}'')$와 동일시된다~; 다른 한편
(\hyperref[II.8.7.3-ko]{8.7.3})에 따라 $C_{P''}$은
$\operatorname{Proj}(\mathcal{S}''^{\natural})$과 $C''$-동형이다.
$\mathcal{S}''$의 정의상 $\mathcal{S}''^{\natural}$은
$\mathcal{S}_1''^{\natural}$에 의해 생성되고,
$\mathcal{S}_1''^{\natural}$은
$\mathcal{S}_0''^{\natural}=\mathcal{S}''$ 위에서 유한형이다
(\hyperref[II.8.2.10-ko]{8.2.10, (i) 및 (iii)}). 따라서
$\operatorname{Proj}(\mathcal{S}''^{\natural})$은 $C''$ 위에서
\emph{사영적}이다 (\hyperref[II.5.5.1-ko]{5.5.1}). 이제 다음 도식을
고려하자.
\begin{equation}
\label{II.8.8.4.1-ko}
  \xymatrix{
    \mathbf{V}(\mathcal{L})\ar[r]^{g}\ar[d]_{u} &
    \operatorname{Spec}(\mathcal{S})=C\ar[d]^{v}\\
    \mathbf{V}(\mathcal{L}'')\ar[r]_{g''} &
    \operatorname{Spec}(\mathcal{S}'')=C''
  }
\tag{8.8.4.1}
\end{equation}
여기서 $g$와 $g''$은 (\hyperref[II.1.5.6-ko]{1.5.6})에 의하여 각각
표준 $f$-사상
\[
  \mathcal{S}\longrightarrow
  \bigoplus_{n\geqslant 0}\mathcal{L}^{\otimes n}
  \quad\text{및}\quad
  \mathcal{S}''\longrightarrow
  \bigoplus_{n\geqslant 0}\mathcal{L}''^{\otimes n}
\]
(\hyperref[II.3.3.2.3-ko]{3.3.2.3})에 대응하고(아래의
(\hyperref[II.8.8.5-ko]{8.8.5})를 보라), $v$는 포함 사상
$\mathcal{S}''\to\mathcal{S}$에 대응하며, $u$는 포함 사상
$\bigoplus_{n\geqslant 0}\mathcal{L}^{\otimes nd}\to
\bigoplus_{n\geqslant 0}\mathcal{L}^{\otimes n}$에 대응한다~;
(\hyperref[II.3.3.2-ko]{3.3.2})에 따라 이 도식이 가환임은 자명하다.
앞에서 $g''$이 사영적 사상임을 보았다~; 다른 한편 $u$는 \emph{유한}
사상이다. 실제로 문제는 $X$ 위에서 국소적이므로,

$X$가 환 $A$의 아핀 스킴이고 $\mathcal{L}=\mathcal{O}_X$라고 가정할 수
있다~; 그러면 모든 것은 환 $A[T]$가 그 부분환 $A[T^d]$ 위의
유한형 가군임($T$는 부정원)에 유의하는 것으로 귀착된다. $Y$는
준콤팩트 스킴이고 $C''$은 $Y$ 위에서 아핀이므로, $C''$도 준콤팩트
스킴이다.

\oldpage[II]{181}
따라서 $g''\circ u$는 사영적 사상이다
(\hyperref[II.5.5.5-ko]{5.5.5, (ii)})~; 도식
(\hyperref[II.8.8.4.1-ko]{8.8.4.1})의 가환성에 의하여 $v\circ g$도
그러하며, $v$는 아핀이고 따라서 분리이므로, 마침내 $g$가 사영적임을
얻는다 (\hyperref[II.5.5.5-ko]{5.5.5, (v)}).
\end{proof}

\begin{env}[8.8.5]
\label{II.8.8.5-ko}
(\hyperref[II.8.8.1-ko]{8.8.1})의 상황을 다시 취하자. 사상
$g:\mathbf{V}(\mathcal{L})\to C$에 대하여, 임의의 가역
$\mathcal{O}_X$-가군 $\mathcal{L}$ (반드시 풍부할 필요는 없다)에
유효한 다른 정의를 줄 수 있음을 보이겠다. 이를 위해
(\hyperref[II.8.8.1.2-ko]{8.8.1.2})의 사상 $\tau$에 대응하는
$f$-사상
\begin{equation}
\label{II.8.8.5.1-ko}
  \tau^\flat:\mathcal{S}\longrightarrow
  \bigoplus_{n\geqslant 0}\mathcal{L}^{\otimes n}
\tag{8.8.5.1}
\end{equation}
을 고려하자. 이로부터 (\hyperref[II.1.5.6-ko]{1.5.6})에 따라 사상
$g':V\to C$를 얻으며, $\pi:V\to X$와 $\psi:C\to Y$가 구조 사상이면
다음 도식들이 가환이다.
\begin{equation}
\label{II.8.8.5.2-ko}
  \xymatrix{
    X\ar[d]_{f} & V\ar[l]^{\pi}\ar[d]^{g'}\\
    Y & C\ar[l]_{\psi}
  }
  \qquad
  \xymatrix{
    X\ar[r]^{j}\ar[d]_{f} & V\ar[d]^{g'}\\
    Y\ar[r]_{\varepsilon} & C
  }
\tag{8.8.5.2}
\end{equation}
(\hyperref[II.8.5.1.2-ko]{8.5.1.2} 및
\hyperref[II.8.5.1.3-ko]{8.5.1.3}). $\mathcal{L}$이 $f$에 대하여
풍부하다고 가정하면 사상 $g$와 $g'$이 동일함을 보이겠다.

실제로 문제는 $Y$ 위에서 국소적이므로,
$Y=\operatorname{Spec}(A)$가 아핀이라고 가정할 수 있고,
(\hyperref[II.8.8.1.3-ko]{8.8.1.3})에 따라 $X$를
$P=\operatorname{Proj}(S)$의 열린집합과 동일시할 수 있다. 여기서
\[
  S=A\oplus\bigoplus_{n\geqslant 1}
  \Gamma(X,\mathcal{L}^{\otimes n})~;
\]
그러면 (\hyperref[II.8.8.1.4-ko]{8.8.1.4})로부터 모든
$n\in\mathbf{Z}$에 대하여
$\Gamma(X,\mathcal{O}_P(n))=\Gamma(X,\mathcal{L}^{\otimes n})$임을 얻는다.
$\alpha$가 표준 $p$-사상
$\widetilde{S}\to\mathcal{S}_P^{\geqslant}$이고
(\hyperref[II.8.6.1.2-ko]{8.6.1.2}),
$h=\operatorname{Spec}(\alpha)$라는 정의를 고려하면, $X$에 제한한
$\alpha^\sharp:p^*(\widetilde{S})\to\mathcal{S}_P^{\geqslant}$이
$\tau$와 동일함을 확인하면 된다. (\hyperref[0.4.4.3-ko]{0, 4.4.3})을
고려하면, 모든 $n>0$에 대하여 표준 준동형
$\alpha_n:S_n\to\Gamma(P,\mathcal{O}_P(n))$과 제한 준동형
\[
  \Gamma(P,\mathcal{O}_P(n))\longrightarrow
  \Gamma(X,\mathcal{O}_P(n))=
  \Gamma(X,\mathcal{L}^{\otimes n}),
\]
의 합성이 항등사상임을 보이는 문제로 귀착된다~; 그런데 이는 대수
$S$의 정의와 $\alpha_n$의 정의
(\hyperref[II.2.6.2-ko]{2.6.2})에서 즉시 따른다.
\end{env}

\begin{proposition}[8.8.6]
\label{II.8.8.6-ko}
(\hyperref[II.8.8.5-ko]{8.8.5})의 표기에서 $f=(f_0,\lambda)$라 놓을 때
준동형 $\lambda:\mathcal{O}_Y\to f_*(\mathcal{O}_X)$가 전단사라고
가정하자~; 그러면 다음이 성립한다~:
\begin{enumerate}
  \item[\rm{(i)}] $g=(g_0,\mu)$라 놓으면,
  $\mu:\mathcal{O}_C\to g_*(\mathcal{O}_V)$는 동형사상이다.
  \item[\rm{(ii)}] $X$가 정역 스킴이면(각각 국소 정역이고 정규이면),
  $C$는 정역 스킴이다(각각 정규이다).
\end{enumerate}
\end{proposition}

\begin{proof}
실제로 이때 $f$-사상 $\tau^\flat$은 \emph{동형사상}
\[
  \tau^\flat:\mathcal{S}=\psi_*(\mathcal{O}_C)
  \longrightarrow f_*(\pi_*(\mathcal{O}_V))
  =\psi_*(g_*(\mathcal{O}_V))
\]
이고, $Y$-사상 $g$는 준동형 $\mathcal{A}(g)$
(\hyperref[II.1.1.2-ko]{1.1.2})가 $\tau^\flat$과 같은 사상으로 볼 수
있다. $\mu$가 $\mathcal{O}_C$-가군의 동형사상임을 보이려면
(\hyperref[II.1.4.2-ko]{1.4.2})
\[
  \mathcal{A}(\mu):\psi_*(\mathcal{O}_C)
  \longrightarrow\psi_*(g_*(\mathcal{O}_V))
\]
가 동형사상임을 보이면 충분하다. 그런데 정의에 따라
(\hyperref[II.1.1.2-ko]{1.1.2})
$\mathcal{A}(\mu)=\mathcal{A}(g)$이므로 (i)의 결론을 얻는다.

(ii)를 증명하기 위해서는 $Y$가 아핀인 경우로 환원할 수 있고, 따라서
$\mathcal{S}=\widetilde{S}$이며

\oldpage[II]{182}
$S=\bigoplus_{n\geqslant 0}\Gamma(X,\mathcal{L}^{\otimes n})$이다~;
$X$가 정역이라는 가정으로부터 환 $S$가 정역임이 따르고
(\hyperref[I.7.4.4-ko]{I, 7.4.4}), 따라서 $C$도 정역이다
(\hyperref[I.5.1.4-ko]{I, 5.1.4}). $C$가 정규임을 증명하기 위해 다음
보조정리를 사용하겠다~:

\begin{lemma}[8.8.6.1]
\label{II.8.8.6.1-ko}
$Z$를 정규 정역 준스킴이라 하자. 그러면 환
$\Gamma(Z,\mathcal{O}_Z)$는 정역이고 정수적으로 닫혀 있다.
\end{lemma}

\begin{proof}
실제로 (\hyperref[I.8.2.1.1-ko]{I, 8.2.1.1})에 따라
$\Gamma(Z,\mathcal{O}_Z)$는 유리함수체 $R(Z)$ 안에서 $z\in Z$에 대한
정수적으로 닫힌 환 $\mathcal{O}_z$들의 교집합이다.
\end{proof}

이제 먼저 $V$가 \emph{국소적으로 정역이고 정규}임을 보이자~; 이를 위해서는
$X=\operatorname{Spec}(A)$가 아핀이고 그 환 $A$가 정역이며 정수적으로
닫혀 있고 (\hyperref[II.6.3.8-ko]{6.3.8}),
$\mathcal{L}=\mathcal{O}_X$인 경우만 생각하면 충분하다. 이때
$V=\operatorname{Spec}(A[T])$이고 $A[T]$는 정역이며 정수적으로 닫혀
있으므로 ([24], p.~99), 주장이 성립한다. 한편 $C$의 임의의 아핀
열린집합 $U$에 대하여 $g$가 준콤팩트이므로 $g^{-1}(U)$는
준콤팩트이다~; $V$가 국소적으로 정역이므로 $g^{-1}(U)$의 연결
성분들은 $g^{-1}(U)$ 안의 열린 정역 준스킴들이고, 따라서 그 수는
유한하다. 또한 $V$가 정규이므로 이 준스킴들은 정규이다
(\hyperref[II.6.3.8-ko]{6.3.8}). 이제 (i)에 따라
$\Gamma(U,\mathcal{O}_C)$는 $\Gamma(g^{-1}(U),\mathcal{O}_V)$와 같고,
이는 유한 개의 정수적으로 닫힌 정역들의 직접곱이다
(\hyperref[II.8.8.6.1-ko]{8.8.6.1}). 그러므로 $C$는 정규이다
(\hyperref[II.6.3.4-ko]{6.3.4}).
\end{proof}

\subsection{그라우어트의 풍부성 판정법: 진술}
\label{subsection:II.8.9-ko}

(\hyperref[II.8.8.2-ko]{8.8.2})에서 증명한 성질들이 $f$에 관해 풍부한
$\mathcal{O}_X$-가군을 \emph{특징짓는다}는 것을 보이고자 한다. 더 정확히는
다음 판정법을 증명하고자 한다~:

\begin{theorem}[8.9.1]
\label{II.8.9.1-ko}
\emph{(그라우어트 판정법).} $Y$를 준스킴,
$p:X\to Y$를 분리 준콤팩트 사상, $\mathcal{L}$을 가역
$\mathcal{O}_X$-가군이라 하자. $\mathcal{L}$이 $p$에 관해 풍부할
필요충분조건은 다음 성질들을 갖는 $Y$-준스킴 $C$, $C$의 $Y$-단면
$\varepsilon:Y\to C$, 그리고 $Y$-사상
$q:\mathbf{V}(\mathcal{L})\to C$가 존재하는 것이다~:
\begin{enumerate}
  \item[\rm{(i)}] 다음 그림은
  \begin{equation}
  \label{II.8.9.1.1-ko}
    \xymatrix{
      X\ar[r]^{j}\ar[d]_{p} & \mathbf{V}(\mathcal{L})\ar[d]^{q}\\
      Y\ar[r]_{\varepsilon} & C
    }
  \tag{8.9.1.1}
  \end{equation}
  가환한다. 여기서 $j$는 벡터 다발 $\mathbf{V}(\mathcal{L})$의 영
  단면이다.
  \item[\rm{(ii)}] $q$를 $\mathbf{V}(\mathcal{L})-j(X)$에 제한한
  사상은 준콤팩트 열린 몰입
  \[
    \mathbf{V}(\mathcal{L})-j(X)\longrightarrow C
  \]
  이고, 그 상은 $\varepsilon(Y)$와 만나지 않는다.
\end{enumerate}
\end{theorem}

$C$가 $Y$ 위에서 \emph{분리}이면, 조건 (ii)에서 열린 몰입이
준콤팩트이라는 가정을 생략할 수 있다~; 실제로 이 마지막 성질이 다른
조건들로부터 따름을 보이기 위해서는 $Y$가 아핀인 경우만 생각하면
되고, 이때 주장은 (\hyperref[I.5.5.1-ko]{I, 5.5.1, (i)} 및
\hyperref[I.5.5.10-ko]{5.5.10})에서 따른다. $X$가 뇌터라고 가정할
때에도

\oldpage[II]{183}
같은 가정을 생략할 수 있다. 실제로 이때 $V$도 뇌터이고, 이번에는
(\hyperref[I.6.3.5-ko]{I, 6.3.5})에서 주장이 따른다.

\begin{corollary}[8.9.2]
\label{II.8.9.2-ko}
사상 $p:X\to Y$가 고유이면 (\hyperref[II.8.9.1-ko]{8.9.1})의 진술에서
$q$가 고유라고 가정할 수 있고, «~열린 몰입~»을 «~동형사상~»으로
바꿀 수 있다.
\end{corollary}

기하학적으로 말하면 ($p:X\to Y$가 고유일 때) $\mathcal{L}$이 $p$에
관해 풍부할 필요충분조건은 벡터 다발 $\mathbf{V}(\mathcal{L})$의 영
단면을 밑준스킴 $Y$로 «~수축~»할 수 있는 것이다. 특히 중요한 경우는
$Y$가 체의 스펙트럼인 경우이며, 이때 «~수축~» 연산은
$\mathbf{V}(\mathcal{L})$의 영 단면을 \emph{하나의 점으로} 수축하는 것이다.

\begin{env}[8.9.3]
\label{II.8.9.3-ko}
(\hyperref[II.8.9.1-ko]{8.9.1})의 조건들이 필요함과 따름정리
(\hyperref[II.8.9.2-ko]{8.9.2})는
(\hyperref[II.8.8.2-ko]{8.8.2})와
(\hyperref[II.8.8.4-ko]{8.8.4})에서 곧바로 따른다.

(\hyperref[II.8.9.1-ko]{8.9.1})의 충분성을 증명하기 위해 조금 더
일반적인 상황을 생각하겠다. 이를 위해
(\hyperref[II.8.8.2-ko]{8.8.2})의 표기법으로
\[
  \mathcal{S}'=\bigoplus_{n\geqslant 0}\mathcal{L}^{\otimes n}
\]
이라 놓고,
\[
  V=\mathbf{V}(\mathcal{L})=\operatorname{Spec}(\mathcal{S}').
\]
이라 하자. $\mathbf{V}(\mathcal{L})$의 영 단면인 닫힌 부분준스킴
$j(X)$는 $\mathcal{O}_V$의 준연접 아이디얼 층
$\mathcal{J}=(\mathcal{S}'_+)^{\sim}$로 정의된다
(\hyperref[II.1.4.10-ko]{1.4.10}). 이 $\mathcal{O}_V$-가군은
\emph{가역}이다. 실제로 문제는 $X$ 위에서 국소적이고, 이는 다항식환
$A[T]$의 아이디얼 $TA[T]$가 단일 생성 자유 $A[T]$-가군임을 확인하는
것으로 귀착된다. 또한 (역시 문제가 $X$ 위에서 국소적이므로)
\[
  \mathcal{L}=j^*(\mathcal{J})
\]
이고
\[
  j_*(\mathcal{L})=\mathcal{J}/\mathcal{J}^2.
\]
임은 즉시 알 수 있다. 한편
\[
  \pi:\mathbf{V}(\mathcal{L})\longrightarrow X
\]
가 구조사상이면 $\pi_*(\mathcal{J})=\mathcal{S}'_+$이고
$\pi_*(\mathcal{J}/\mathcal{J}^2)=\mathcal{L}$이다~; 따라서 두 표준
준동형 $\mathcal{L}\to\pi_*(\mathcal{J})\to\mathcal{L}$이 존재한다.
첫째는 표준 단사 준동형 $\mathcal{L}\to\mathcal{S}'_+$이고, 둘째는
$\mathcal{S}'_+$에서 $\mathcal{S}'_1=\mathcal{L}$ 위로 가는 표준
사영이며, 이들의 합성은 항등사상이다. 더 나아가
\[
  \pi_*(\mathcal{J})=\mathcal{S}'_+
  =\bigoplus_{n\geqslant 1}\mathcal{L}^{\otimes n}
\]
을 \emph{직접곱}
\[
  \prod_{n\geqslant 1}\mathcal{L}^{\otimes n}
  =\varprojlim_n\pi_*(\mathcal{J}/\mathcal{J}^{n+1})
\]
에 표준적으로 매장할 수 있다(실제로
$\pi_*(\mathcal{J}/\mathcal{J}^{n+1})=
\mathcal{L}\oplus\mathcal{L}^{\otimes2}\oplus\cdots\oplus
\mathcal{L}^{\otimes n}$이다). 따라서 합성이 항등사상인 두 표준
준동형
\begin{equation}
\label{II.8.9.3.1-ko}
  \mathcal{L}\longrightarrow
  \varprojlim_n\pi_*(\mathcal{J}/\mathcal{J}^{n+1})
  \longrightarrow\mathcal{L}
\tag{8.9.3.1}
\end{equation}
을 다시 얻는다.

이제 우리가 증명할 (\hyperref[II.8.9.1-ko]{8.9.1})의 일반화는 다음과
같다~:
\end{env}

\begin{proposition}[8.9.4]
\label{II.8.9.4-ko}
$Y$를 준스킴, $V$를 $Y$-준스킴, $X$를 $V$의 닫힌 부분준스킴이라
하자. $X$는 가역 $\mathcal{O}_V$-가군인 $\mathcal{O}_V$의 아이디얼
$\mathcal{J}$로 정의된다고 하자~; $j:X\to V$를

\oldpage[II]{184}
$j:X\to V$를 표준 단사 사상이라 하고,
$\mathcal{L}=j^*(\mathcal{J})=
\mathcal{J}\otimes_{\mathcal{O}_V}\mathcal{O}_X$로 놓자. 그러면
$j_*(\mathcal{L})=\mathcal{J}/\mathcal{J}^2$이다.
구조 사상 $p:X\to Y$가 분리이고 준콤팩트이며, 다음 조건들이
성립한다고 가정한다.
\begin{enumerate}
  \item[\rm{(i)}] $\pi\circ j=1_X$를 만족하는 유한형
  $Y$-사상 $\pi:V\to X$가 존재한다. 따라서
  $\pi_*(\mathcal{J}/\mathcal{J}^2)=\mathcal{L}$이다.
  \item[\rm{(ii)}] $\mathcal{O}_X$-가군 준동형
  \[
    \varphi:\mathcal{L}\longrightarrow
    \varprojlim_n\pi_*(\mathcal{J}/\mathcal{J}^{n+1})
  \]
  이 존재하여, 합성
  \[
    \mathcal{L}\xrightarrow{\varphi}
    \varprojlim_n\pi_*(\mathcal{J}/\mathcal{J}^{n+1})
    \xrightarrow{\alpha}\pi_*(\mathcal{J}/\mathcal{J}^2)=\mathcal{L}
  \]
  이 항등사상이 된다. 여기서 $\alpha$는 표준 준동형이다.
  \item[\rm{(iii)}] $Y$-준스킴 $C$, $C$의 $Y$-단면
  $\varepsilon$, 그리고 다음 도표를 가환하게 하는 $Y$-사상
  $q:V\to C$가 존재한다.
  \begin{equation}
  \label{II.8.9.4.1-ko}
    \xymatrix{
      X\ar[r]^{j}\ar[d]_{p} & V\ar[d]^{q}\\
      Y\ar[r]_{\varepsilon} & C
    }
  \tag{8.9.4.1}
  \end{equation}
  \item[\rm{(iv)}] $q$의 $W=V-j(X)$에 대한 제한은 $C$로 가는
  준콤팩트 열린 몰입 사상이고, 그 상은 $\varepsilon(Y)$와 만나지
  않는다.
\end{enumerate}
이 조건들 아래에서 $\mathcal{L}$은 $p$에 대해 풍부하다.
\end{proposition}

\subsection{그라우에르트의 풍부성 판정법: 증명}
\label{subsection:II.8.10-ko}

\begin{lemma}[8.10.1]
\label{II.8.10.1-ko}
$\pi:V\to X$를 사상, $j:X\to V$를 닫힌 몰입 사상인 $V$의
$X$-단면이라 하고, $\mathcal{J}$를 $j$에 대응하는 $V$의 닫힌
부분준스킴을 정의하는 $\mathcal{O}_V$의 준연접 아이디얼이라 하자.
\begin{enumerate}
  \item[\rm{(i)}] 모든 $n\geqslant 0$에 대하여
  $\pi_*(\mathcal{O}_V/\mathcal{J}^{n+1})$과
  $\pi_*(\mathcal{J}/\mathcal{J}^{n+1})$은 준연접
  $\mathcal{O}_X$-가군이고,
  $\pi_*(\mathcal{O}_V/\mathcal{J})=\mathcal{O}_X$,
  $\pi_*(\mathcal{J}/\mathcal{J}^2)=j^*(\mathcal{J})$이다.
  \item[\rm{(ii)}] $X=\{\xi\}=\operatorname{Spec}(k)$이고 $k$가
  체이면, $\varprojlim_n\pi_*(\mathcal{O}_V/\mathcal{J}^{n+1})$은
  국소환 $\mathcal{O}_{j(\xi)}$의
  $\mathfrak{m}_{j(\xi)}$-진 위상에 대한 분리 완비화와 동형이다.
  \item[\rm{(iii)}] $\mathcal{J}$가 가역 $\mathcal{O}_V$-가군이라고
  가정하자. 이로부터
  \[
    \mathcal{L}=j^*(\mathcal{J})=
    \pi_*(\mathcal{J}/\mathcal{J}^2)
  \]
  은 가역 $\mathcal{O}_X$-가군이다. 또한 준동형
  \[
    \varphi:\mathcal{L}\longrightarrow
    \varprojlim_n\pi_*(\mathcal{J}/\mathcal{J}^{n+1})
  \]
  이 존재하여, 합성
  \[
    \mathcal{L}\xrightarrow{\varphi}
    \varprojlim_n\pi_*(\mathcal{J}/\mathcal{J}^{n+1})
    \xrightarrow{\alpha}\pi_*(\mathcal{J}/\mathcal{J}^2)
  \]
  이 항등사상이 된다고 하자. 여기서 $\alpha$는 표준 준동형이다.
  $\mathcal{S}=\bigoplus_{n\geqslant0}\mathcal{L}^{\otimes n}$로 놓으면,
  $\varphi$로부터 표준적으로 $\mathcal{S}$의 표준 여과에 대한 완비화
  $\widehat{\mathcal{S}}$에서
  $\varprojlim_n\pi_*(\mathcal{O}_V/\mathcal{J}^{n+1})$ 위로 가는
  $\mathcal{O}_X$-대수 동형사상이 유도된다. 여기서 이 완비화는
  $\prod_{n\geqslant0}\mathcal{L}^{\otimes n}$과 동형이다.
\end{enumerate}
\end{lemma}

\begin{proof}
먼저 $\mathcal{O}_V$-가군
$\mathcal{O}_V/\mathcal{J}^{n+1}$의 지지는 $j(X)$이고,
$\mathcal{J}/\mathcal{J}^{n+1}$의 지지는 $j(X)$에 포함됨을 주목하자.
(ii)의 경우 $j(X)$는 $V$의 닫힌점 $j(\xi)$이다.

\oldpage[II]{185}
정의에 따라 $\pi_*(\mathcal{O}_V/\mathcal{J}^{n+1})$은
$j(\xi)$에서 $\mathcal{O}_V/\mathcal{J}^{n+1}$의 섬유이다. 즉
$C=\mathcal{O}_{j(\xi)}$로 놓고 $C$의 극대 아이디얼을
$\mathfrak{m}$으로 나타내면, 이는 $C$-가군
$C/\mathfrak{m}^{n+1}$이다. 따라서 (ii)의 주장은 자명하다.

(i)를 증명하기 위해 이 문제는 $X$ 위에서 국소적임을 주목하자. 따라서
$X$가 아핀인 경우로 한정할 수 있다. $U$를 $V$의 아핀 열린집합이라
하자. $j(X)\cap U$는 $j(X)$의 아핀 열린집합이므로, 이와 동형인
$U_0=\pi(j(X)\cap U)$는 $X$의 아핀 열린집합이다. $X$에서
$W_0\subset U_0$인 모든 아핀 열린집합 $W_0$에 대하여
$W=\pi^{-1}(W_0)\cap U$는 $V$의 아핀 열린집합이다. 이는 $X$가
스킴이기 때문이다(\hyperref[I.5.5.10-ko]{I, 5.5.10}). 특히
$U'=U\cap\pi^{-1}(U_0)$는 $V$의 아핀 열린집합이고, 명백히
$\pi(U')=U_0$ 및 $j(U_0)=j(X)\cap U$이다. 한편 정의에 따라
\[
  \Gamma(W_0,\pi_*(\mathcal{O}_V/\mathcal{J}^{n+1}))
  =\Gamma(\pi^{-1}(W_0),\mathcal{O}_V/\mathcal{J}^{n+1})~;
\]
그런데 $\pi^{-1}(W_0)$에서 $j(W_0)$에 속하지 않는 각 점은
$\pi^{-1}(W_0)$ 안에서 $j(X)$와 만나지 않는 열린 근방을 가지며,
그 근방에서는 $\mathcal{O}_V/\mathcal{J}^{n+1}$이 영이다. 따라서
$\pi^{-1}(W_0)$ 위의 $\mathcal{O}_V/\mathcal{J}^{n+1}$의 단면들과
$W$ 위의 단면들은 일대일로 대응한다. 달리 말해 $\pi'$을 $\pi$의
$U'$에 대한 제한이라 하면, 두 $(\mathcal{O}_X|U_0)$-가군
\[
  \pi_*(\mathcal{O}_V/\mathcal{J}^{n+1})|U_0
  \quad\hbox{와}\quad
  \pi'_*\bigl((\mathcal{O}_V/\mathcal{J}^{n+1})|U'\bigr)
\]
은 동일하다. $U'$과 $U_0$는 아핀이고 이러한 $U_0$들이 $X$를
덮으므로, (\hyperref[I.1.6.3-ko]{I, 1.6.3})에서
$\pi_*(\mathcal{O}_V/\mathcal{J}^{n+1})$이 준연접임을 얻는다.
$\pi_*(\mathcal{J}/\mathcal{J}^{n+1})$에 대한 증명도 동일하다.

마지막으로 (iii)을 증명하자. $\mathcal{S}$는 바로
$\mathbf{S}_{\mathcal{O}_X}(\mathcal{L})$이므로, $\varphi$로부터
표준적으로 $\mathcal{O}_X$-대수 준동형
\[
  \psi:\mathcal{S}\longrightarrow
  \varprojlim_n\pi_*(\mathcal{O}_V/\mathcal{J}^{n+1}).
\]
을 얻는다(\hyperref[II.1.7.4-ko]{1.7.4}). 또한 이 준동형은
$\mathcal{L}^{\otimes n}$을
$\varprojlim_m\pi_*(\mathcal{J}^n/\mathcal{J}^{m+1})$ 안으로 보내므로,
고려하는 위상들에 대하여 연속이고, 따라서 실제로 준동형
\[
  \widehat{\psi}:\widehat{\mathcal{S}}\longrightarrow
  \varprojlim_n\pi_*(\mathcal{O}_V/\mathcal{J}^{n+1})
\]
으로 연장된다. 이것이 동형사상임을 보이기 위해 (i)의 증명에서처럼
$X=\operatorname{Spec}(A)$와 $V=\operatorname{Spec}(B)$가 아핀이고
$\mathcal{J}=\widetilde{\mathfrak{J}}$이며 $\mathfrak{J}$가 $B$의
아이디얼인 경우로 한정할 수 있다. $\pi$에는 단사 준동형
$A\to B$가 대응하며, 이는 $A$를 $\mathfrak{J}$와 직합을 이루는
$B$의 부분환으로 동일시한다. 그리고 $\mathcal{L}$은
$A$-가군 $L=\mathfrak{J}/\mathfrak{J}^2$에 대응하는 준연접
$\mathcal{O}_X$-가군이고,
$\pi_*(\mathcal{O}_V/\mathcal{J}^{n+1})$은
$A$-가군 $B/\mathfrak{J}^{n+1}$에 대응하는 준연접
$\mathcal{O}_X$-가군이다. $\mathcal{J}$는 \emph{가역}
$\mathcal{O}_V$-가군이므로, 더 나아가
$\mathfrak{J}=Bt$이고 $t$가 $B$의 영인자가 아닌 경우로 가정할 수
있다. 관계 $B=A\oplus Bt$에서 모든 $n>0$에 대하여
\[
  B=A\oplus At\oplus At^2\oplus\cdots\oplus At^n\oplus Bt^{n+1}
\]
을 얻는다. 따라서 $T$와 $t$를 대응시키는, 형식적 멱급수환
$A[[T]]$에서 $C=\varprojlim B/\mathfrak{J}^{n+1}$ 위로 가는 표준
$A$-동형사상이 존재한다. 다른 한편
$L=A\bar t$이고, 여기서 $\bar t$는 $Bt^2$를 법으로 한 $t$의
동치류이다. 가정에 따라 준동형 $\varphi$는 $\bar t$를 $Ct^2$를
법으로 하여 $t$와 합동인 원소 $t'\in C$로 보낸다. $n$에 대한
귀납법으로
\[
  A\oplus At'\oplus\cdots\oplus At'{}^n\oplus Ct^{n+1}
  =A\oplus At\oplus\cdots\oplus At^n\oplus Ct^{n+1}
\]
을 얻으며, 이는 $\widehat{\psi}$가
$\prod_{n\geqslant0}L^{\otimes n}$에서 $C$ 위로 가는 동형사상에
실제로 대응함을 증명한다.
\end{proof}

\begin{lemma}[8.10.2]
\label{II.8.10.2-ko}
(\hyperref[II.8.10.1-ko]{8.10.1})의 가정 아래에서,
$g:X'\to X$를 사상이라 하고,

\oldpage[II]{186}
$V'=V\times_X X'$라 하고, $\pi':V'\to X'$와
$g':V'\to V$를 표준 사영이라 하자. 그러면 다음 가환 도표를 얻는다.
\[
  \xymatrix{
    V\ar[d]_{\pi} & V'\ar[l]^{g'}\ar[d]^{\pi'}\\
    X & X'\ar[l]_{g}
  }
\]
그러면 $j'=j\times1_{X'}$은 닫힌 몰입 사상인 $V'$의 $X'$-단면이고,
$\mathcal{J}'=(g')^*(\mathcal{J})\mathcal{O}_{V'}$은 $j'$에 대응하는
$V'$의 닫힌 부분준스킴을 정의하는 $\mathcal{O}_{V'}$의 준연접
아이디얼층이다. 또한
\[
  \pi'_*(\mathcal{O}_{V'}/\mathcal{J}'^{n+1})
  =g^*(\pi_*(\mathcal{O}_V/\mathcal{J}^{n+1})).
\]
마지막으로 $\mathcal{J}'$은 $(g')^*(\mathcal{J})$와 표준적으로 동형인
$\mathcal{O}_{V'}$-가군이며, 특히 $\mathcal{J}$가 가역
$\mathcal{O}_V$-가군이면 가역이다.
\end{lemma}

\begin{proof}
$j'$이 닫힌 몰입 사상이라는 사실은
(\hyperref[I.4.3.1-ko]{I, 4.3.1})에서 따르며, 밑준스킴 확장의
함자성에 의하여 그것은 $V'$의 $X'$-단면이다. 또한 $Z$(각각 $Z'$)를
$j$(각각 $j'$)에 대응하는 $V$(각각 $V'$)의 닫힌 부분준스킴이라 하면
$Z'=g'{}^{-1}(Z)$이다
(\hyperref[I.4.3.1-ko]{I, 4.3.1}). 따라서 두 번째 주장은
(\hyperref[I.4.4.5-ko]{I, 4.4.5})에서 따른다. 나머지 주장들을
증명하기 위해 (\hyperref[II.8.10.1-ko]{8.10.1})에서와 같이 $X$, $V$,
$X'$가 아핀인 경우로(따라서 $V'$도 아핀인 경우로) 한정할 수 있다.
(\hyperref[II.8.10.1-ko]{8.10.1})의 증명에서 쓴 표기를 유지하고
$X'=\operatorname{Spec}(A')$라 하자. 그러면
$V'=\operatorname{Spec}(B')$이고 $B'=B\otimes_A A'$이며,
$\mathcal{J}'=\widetilde{\mathfrak{J}'}$이고
$\mathfrak{J}'=\operatorname{Im}(\mathfrak{J}\otimes_A A')$이다. 따라서
\[
  B'/\mathfrak{J}'^{n+1}=(B/\mathfrak{J}^{n+1})\otimes_A A'~;
\]
또한 $\mathfrak{J}$는 $B$의 직합인자($A$-가군으로서)이므로,
$\mathfrak{J}\otimes_A A'$은 $B'$의 직합인자($A'$-가군으로서)이며
따라서 $\mathfrak{J}'$과 표준적으로 동일시된다.
\end{proof}

\begin{corollary}[8.10.3]
\label{II.8.10.3-ko}
(\hyperref[II.8.10.1-ko]{8.10.1})의 가정들이 성립하고, 더 나아가
$\pi$가 유한형이고 $\mathcal{J}$가 가역 $\mathcal{O}_V$-가군이라고
가정하자. 그러면 모든 $x\in X$에 대하여 섬유 $\pi^{-1}(x)$의 점
$j(x)$에서의 국소환은 1차원 정칙환(따라서 정역)이고, 그 완비화는
형식적 멱급수환 $k(x)[[T]]$($T$는 부정원)과 동형이다. 또한
$j(x)$를 포함하는 $\pi^{-1}(x)$의 기약성분은 단 하나뿐이다.
\end{corollary}

\begin{proof}
$\pi^{-1}(x)=V\times_X\operatorname{Spec}(k(x))$이므로,
(\hyperref[II.8.10.2-ko]{8.10.2})에 의하여 $X$가 체 $K$의 스펙트럼인
경우로 한정된다. $\pi$가 유한형이므로
(\hyperref[I.6.3.4-ko]{I, 6.3.4, (iv)}), $\mathcal{O}_{j(x)}$는 뇌터
국소환이고, 따라서 $\mathfrak{m}_{j(x)}$-진 위상에 관하여 분리되어
있다(\hyperref[0.7.3.5-ko]{0, 7.3.5}).
(\hyperref[II.8.10.1-ko]{8.10.1, (ii)와 (iii)})에서 이 환의 완비화가
$K[[T]]$와 동형임이 따른다. 따라서 $\mathcal{O}_{j(x)}$는 1차원
정칙환이다([1], p.~17-01, th.~1). 마지막으로 $\mathcal{O}_{j(x)}$가
정역이므로, $j(x)$는 $V$의 (유한 개의) 기약성분 중 단 하나에만
속한다(\hyperref[I.5.1.4-ko]{I, 5.1.4}).
\end{proof}

\begin{corollary}[8.10.4]
\label{II.8.10.4-ko}
(\hyperref[II.8.10.1-ko]{8.10.1})의 가정들이 성립하고, 더 나아가
$\mathcal{J}$가 가역 $\mathcal{O}_V$-가군이라고 가정하자.
$W=V-j(X)$라 하자. $\mathcal{O}_X$의 모든 준연접 아이디얼층
$\mathcal{K}$에 대하여
$\mathcal{K}_V=\pi^*(\mathcal{K})\mathcal{O}_V$ 및
$\mathcal{K}_W=\mathcal{K}_V|W$라 놓자. 그러면 $\mathcal{K}_V$는
$W$에 대한 제한이 $\mathcal{K}_W$인 $\mathcal{O}_V$의 준연접
아이디얼층 중 가장 큰 것이다.
\end{corollary}

\begin{proof}
실제로 (\hyperref[II.8.10.1-ko]{8.10.1})에서와 같이 문제는 $X$와
$V$에서 국소적임을 알 수 있다. 따라서
(\hyperref[II.8.10.1-ko]{8.10.1})의 증명에서 쓴 표기를 유지하여
$\mathfrak{J}=Bt$라 할 수 있으며, 여기서 $t$는 $B$의 영인자가 아니다.
또한 $W=\operatorname{Spec}(B_t)$이고
$\mathcal{K}=\widetilde{\mathfrak{K}}$이며, 여기서 $\mathfrak{K}$는
$A$의 아이디얼이다. 따라서
$\pi^*(\mathcal{K})\mathcal{O}_V=(\mathfrak{K}.B)^{\sim}$
(\hyperref[I.1.6.9-ko]{I, 1.6.9})이고,
$\mathcal{K}_W=(\mathfrak{K}.B_t)^{\sim}$이다. $B_t$에서의 표준상이
$\mathfrak{K}.B_t$인 $B$의 가장 큰 아이디얼은

\oldpage[II]{187}
$\mathfrak{K}.B_t$의 역상, 즉 어떤 정수 $n>0$에 대하여
$t^ns\in\mathfrak{K}.B$인 $s\in B$들의 집합이다. 이 마지막 관계로부터
$s\in\mathfrak{K}.B$가 따라옴을 보여야 한다. 이는 동치로 $t$의
표준상이
$B/\mathfrak{K}B=(A/\mathfrak{K})\otimes_A B$에서 영인자가 아님을
보이는 것이며, 이 사실은 $X'=\operatorname{Spec}(A/\mathfrak{K})$에
(\hyperref[II.8.10.2-ko]{8.10.2})를 적용하면 따른다.
\end{proof}

\begin{corollary}[8.10.5]
\label{II.8.10.5-ko}
(\hyperref[II.8.10.3-ko]{8.10.3})의 가정들이 성립한다고 하자.
$W=V-j(X)$라 하고, $x$를 $X$의 한 점, $\mathcal{K}$를
$\mathcal{O}_X$의 준연접 아이디얼층, $z$를 $j(x)$를 포함하는
$\pi^{-1}(x)$의 기약성분의 일반점이라 하자
(\hyperref[II.8.10.3-ko]{8.10.3}).
\begin{enumerate}
  \item[\rm{(i)}] $g$를 $V$ 위의 $\mathcal{O}_V$의 단면이라 하고,
  $g|W$가 $W$ 위의 $\mathcal{K}_W$의 단면이라고 하자
  ((\hyperref[II.8.10.4-ko]{8.10.4})의 표기). 그러면 $g$는
  $\mathcal{K}_V$의 단면이다. 더 나아가 $g(z)\neq0$이고, 모든 정수
  $m>0$에 대하여 $g_m^x$가
  $\Gamma(X,\pi_*(\mathcal{O}_V/\mathcal{J}^{m+1}))$에서 $g$의
  표준상 $g_m$의 점 $x$에서의 싹을 나타낸다면, 어떤 정수 $m>0$에
  대하여 $g_m^x$의 다음 공간에서의 상은 $\neq0$이다.
  \[
    (\pi_*(\mathcal{O}_V/\mathcal{J}^{m+1}))_x
    \otimes_{\mathcal{O}_x}k(x)
  \]
  \item[\rm{(ii)}] 더 나아가
  (\hyperref[II.8.10.1-ko]{8.10.1, (iii)})의 조건들이 성립한다고 하자.
  $g(z)\neq0$인 $V$ 위의 $\mathcal{K}_V$의 단면 $g$가 존재하면,
  어떤 정수 $n\geqslant0$과
  $\mathcal{K}.\mathcal{L}^{\otimes n}=
  \mathcal{K}\otimes\mathcal{L}^{\otimes n}\subset\mathcal{L}^{\otimes n}$의
  단면 $f$가 존재하여 $f(x)\neq0$이다. $g$가 $\mathcal{J}$의
  단면이면 $n>0$으로 택할 수 있다.
\end{enumerate}
\end{corollary}

\begin{proof}
\begin{enumerate}
  \item[\rm{(i)}] 가정에 의하여 $g|W$가 생성하는 $\mathcal{O}_W$의
  아이디얼층은 $\mathcal{K}_W$에 포함된다. 그러므로
  (\hyperref[II.8.10.4-ko]{8.10.4})에 의하여 $g$가 생성하는
  $\mathcal{O}_V$의 아이디얼층은 $\mathcal{K}_V$에 포함된다. 다시 말해
  $g$는 $\mathcal{K}_V$의 단면이다. (i)의 두 번째 주장을 증명하기
  위해 다시 $X$와 $V$가 아핀이라고 가정하고
  (\hyperref[II.8.10.1-ko]{8.10.1})의 표기를 유지할 수 있다. 그러면
  섬유 $\pi^{-1}(x)$는 환 $B'=B\otimes_A k(x)$을 갖는 아핀
  준스킴이고, $B'=k(x)\oplus B't'$이 되게 하는 영인자가 아닌 원소
  $t'$이 $B'$ 안에 존재한다. $j(x)$는 $z$의 특수화이고 $g(z)\neq0$이므로
  반드시 $g_{j(x)}\neq0$이다. 그런데 $\mathcal{O}_{j(x)}$는 분리된
  국소환이므로(\hyperref[II.8.10.3-ko]{8.10.3}) 자신의 완비화에
  단사로 들어가며, 따라서 이 완비화에서 $g$의 상은 영이 아니다.
  이 완비화는
  $\varprojlim_n(B'/B't'^{n+1})$와 동형이다
  (\hyperref[II.8.10.3-ko]{8.10.3}). $g'=g\otimes1\in B'$라 하면,
  따라서 어떤 정수 $m$에 대하여 $g'\notin B't'^{m+1}$이다. 다시 말해
  $B'/B't'^{m+1}$에서 $g'$의 상 $g'_m$은 영이 아니다. 그런데
  $g'_m$은 바로 $g_m^x$의 상이므로 주장이 증명되었다.
  \item[\rm{(ii)}]
  (\hyperref[II.8.10.1-ko]{8.10.1, (iii)})에 의하여
  $\pi_*(\mathcal{O}_V/\mathcal{J}^{m+1})$은
  $0\leqslant k\leqslant m$인 $\mathcal{L}^{\otimes k}$들의 직합과
  동형이다. 각 $k$에 대하여 $f_k$를 $X$ 위의
  $\mathcal{L}^{\otimes k}$의 단면, 즉 다음 공간의 한 원소의 성분이라
  하자.

  $\bigoplus_{k=0}^m\Gamma(X,\mathcal{L}^{\otimes k})$.
  이 원소는 이 동형에 의해 $g_m$에 대응한다. (i)에서와 같이 $m$을 택하면,
  (i)에 의해 $f_k(x)\neq0$인 지표 $k$가 존재한다. $f_k$가
  $\mathcal{K}\mathcal{L}^{\otimes k}$의 단면임을 보이려면 위에서처럼
  $X$와 $V$가 아핀인 경우를 생각하면 충분하며, 그 경우 이는
  $g\in\mathfrak{K}.B$라는 사실
  ((\hyperref[II.8.10.4-ko]{8.10.4})의 표기)에서 곧바로 따른다.
  마지막 주장은 $g\in\Gamma(V,\mathcal{J})$라는 가정이 $f_0=0$을
  함의한다는 사실에서 따른다.
\end{enumerate}
\end{proof}

\begin{env}[8.10.6]
\label{II.8.10.6-ko}
\emph{(\hyperref[II.8.9.4-ko]{8.9.4})의 증명.} 문제는 $Y$ 위에서
국소적이다 (\hyperref[II.4.6.4-ko]{4.6.4})~; $\varepsilon$은
$Y$-단면이므로, $C$를 $\varepsilon(Y)$의 한 점의 아핀 열린 근방 $U$로
바꾸어 $\varepsilon(Y)\cap U$가 $U$에서 닫혀 있다고 할 수 있다.
다시 말해 $C$가 아핀이고, $Y$가 $C$의 닫힌 부분준스킴(따라서 아핀)으로서
$\mathcal{O}_C$의

\oldpage[II]{188}

준연접 아이디얼층 $\mathcal{I}$로 정의된다고 가정할 수 있다. $p$가 분리
사상이자 준콤팩트 사상이므로 $X$는 준콤팩트 스킴이고,
문제는 $\mathcal{L}$이 풍부함을 증명하는 것으로 환원된다
(\hyperref[II.4.6.4-ko]{4.6.4}). 판정법
(\hyperref[II.4.5.2-ko]{4.5.2, a)})에 따라, 따라서 다음을 증명해야 한다~:
$\mathcal{O}_X$의 모든 준연접 아이디얼층 $\mathcal{K}$와
$\mathcal{O}_X/\mathcal{K}$의 지지집합에 속하지 않는 모든 점 $x\in X$에
대하여, 정수 $n>0$과 $X$ 위의
$\mathcal{K}\otimes\mathcal{L}^{\otimes n}$의 단면 $f$로서
$f(x)\neq0$인 것이 존재한다.

이를 위해 다음과 같이 놓자.
\[
  \mathcal{K}_V=\pi^*(\mathcal{K})\mathcal{O}_V
\]
그리고
\[
  \mathcal{K}_W=\mathcal{K}_V|W,\qquad\hbox{여기서}\qquad W=V-j(X)~;
\]
$q$의 $W$로의 제한은 $C$로 가는 준콤팩트 몰입 사상이므로,
(\hyperref[I.9.4.2-ko]{I, 9.4.2})에 따라 $\mathcal{K}_W$는 다음 꼴인
$\mathcal{O}_V$의 준연접 아이디얼층 $\mathcal{K}'_V$의 $W$로의 제한이다.
\[
  \mathcal{K}'_V=q^*(\mathcal{K}_C)\mathcal{O}_V
\]
여기서 $\mathcal{K}_C$는 $\mathcal{O}_C$의 준연접 아이디얼층이다.
또한 가정에 따라 $q^{-1}(Y)\subset j(X)$이고 $Y$는 아이디얼층
$\mathcal{I}$로 정의되므로, $q^*(\mathcal{I})\mathcal{O}_V$의 $W$로의
제한은 $\mathcal{O}_V$의 제한과 동일하다. 따라서 $\mathcal{K}_W$는
$q^*(\mathcal{I}\mathcal{K}_C)\mathcal{O}_V$의 $W$로의 제한이기도 하며,
이에 따라 $\mathcal{K}_C\subset\mathcal{I}$라고 가정할 수 있다. 그러므로
\begin{equation}
\label{II.8.10.6.1-ko}
  \mathcal{K}'_V\subset q^*(\mathcal{I})\mathcal{O}_V\subset\mathcal{J}
\tag{8.10.6.1}
\end{equation}
이다. 여기서는 (\hyperref[I.4.4.6-ko]{I, 4.4.6})과
(\hyperref[II.8.9.4.1-ko]{8.9.4.1})이 가환한다는 사실을 함께 사용했다.
또한 (\hyperref[II.8.10.4-ko]{8.10.4})로부터
\begin{equation}
\label{II.8.10.6.2-ko}
  \mathcal{K}'_V\subset\mathcal{K}_V
\tag{8.10.6.2}
\end{equation}
를 얻는다.

이제 (\hyperref[II.8.10.3-ko]{8.10.3})에 따라 $j(x)$는
$\pi^{-1}(x)$의 유일한 기약 성분에 속한다~; 이 성분의 일반점을 $z$라
하고 $z'=q(z)$로 놓자. (\hyperref[II.8.10.5-ko]{8.10.5})에 따라,
(\hyperref[II.8.10.6.1-ko]{8.10.6.1})과
(\hyperref[II.8.10.6.2-ko]{8.10.6.2})을 고려하면 $V$ 위의
$\mathcal{K}'_V$의 단면 $g$로서 $g(z)\neq0$인 것이 존재함을 증명하는
것으로 증명이 완결된다. 그런데 가정에 따라 $x$의 어떤 열린 근방에서
$\mathcal{K}$의 제한은 $\mathcal{O}_X$의 제한과 같다~; 다른 한편
(\hyperref[II.8.10.3-ko]{8.10.3})에 따라 $z\neq j(x)$이므로 $z\in W$이고,
따라서 $(\mathcal{K}_W)_z=\mathcal{O}_{V,z}$이다. 그러므로 정의에 의해
$(\mathcal{K}_C)_{z'}=\mathcal{O}_{C,z'}$이다. $C$가 아핀이므로,
$C$ 위의 $\mathcal{K}_C$의 단면 $g'$로서 $g'(z')\neq0$인 것이 존재한다.
$g'$에 표준적으로 대응하는 $\mathcal{K}'_V$의 단면을 $g$로 택하면 실제로
$g(z)\neq0$이고, 이로써 증명이 끝난다.
\end{env}

\begin{remark}[8.10.7]
\label{II.8.10.7-ko}
(\hyperref[II.8.9.4-ko]{8.9.4})의 명제에서 조건 (ii)가 불필요한지
아닌지는 알려져 있지 않다. 어쨌든 $\pi\circ j=1_X$를 만족하는
$Y$-사상 $\pi:V\to X$의 존재를 가정하지 않으면 결론은 더 이상
성립하지 않는다~; 실제로 반례를 구성할 수 있는 방법을 간략히 설명하자.
그 세부 사항은 뒤에서야 전개할 수 있다. $k$를 체라 하고
$Y=\operatorname{Spec}(k)$, $A=k[T_1,T_2]$에 대하여
$C=\operatorname{Spec}(A)$로 잡으며, $Y$-단면 $\varepsilon$은 증대
준동형 $A\to k$에 대응한다고 하자. $C$의 닫힌 점
$a=\varepsilon(Y)$를 $C$에서 블로업하여 얻은 스킴을 $C'$이라 하자~; $D$를
$C'$에서 $a$의 역상이라 하고, $D$의 닫힌 점 $b$를 하나 택한다.

\oldpage[II]{189}
$b$를 블로업하여 $C'$에서 얻은 스킴을 $V$라 하자~; $X$는 구조 사상
$q:V\to C$에 의한 $a$의 역상인 $V$의 닫힌 부분준스킴이다. $X$가 두
기약 성분 $X_1$, $X_2$의 합집합이고, 여기서 $X_1$은 $V$에서 $b$의
역상임을 보일 수 있다. $X$를 정의하는 $\mathcal{O}_V$의 아이디얼층
$\mathcal{J}$가 여전히 가역이라는 것은 자명하지만,
$j^*(\mathcal{J})=\mathcal{L}$ ($j$는 표준 포함 사상 $X\to V$이다)이
풍부하지 않음을 증명할 수 있다. 실제로 $X_1$ 위로 끌어올린
$\mathcal{L}$의 «~차수~»를 생각하면, $\mathcal{L}$이 풍부할 경우 이는
$>0$이어야 하지만 (교차에 관한 초등적 계산으로) 실제로 0임을 보일 수
있다.
\end{remark}

\subsection{수축의 유일성}
\label{subsection:II.8.11-ko}

\begin{lemma}[8.11.1]
\label{II.8.11.1-ko}
$U$, $V$를 두 준스킴, $h=(h_0,\lambda):U\to V$를 전사 사상이라 하자.
다음을 가정한다~:
\begin{enumerate}
  \item[$1^\circ$]
  $\lambda:\mathcal{O}_V\to h_*(\mathcal{O}_U)
  =(h_0)_*(\mathcal{O}_U)$는 동형사상이다.
  \item[$2^\circ$] $V$의 바탕공간은 $U$의 바탕공간을 관계
  $R:h_0(x)=h_0(y)$로 나눈 몫공간과 동일시된다 (사상 $h$가 열린
  사상이거나 닫힌 사상일 때, 특히 $h$가 고유 사상일 때에는 이 조건이
  항상 성립한다).
\end{enumerate}
그러면 모든 준스킴 $W$에 대하여 사상
\begin{equation}
\label{II.8.11.1.1-ko}
  \operatorname{Hom}(V,W)\longrightarrow\operatorname{Hom}(U,W)
\tag{8.11.1.1}
\end{equation}
은, $V$에서 $W$로 가는 각 사상 $v=(v_0,\nu)$에 사상
$u=v\circ h=(u_0,\mu)$를 대응시키며, $\operatorname{Hom}(V,W)$에서
$u_0$가 모든 올 $h_0^{-1}(x)$ 위에서 상수인 $u$들의 집합으로 가는
전단사이다.
\end{lemma}

\begin{proof}
$u=v\circ h$, 따라서 $u_0=v_0\circ h_0$이면 $u_0$가 모든 집합
$h_0^{-1}(x)$ 위에서 상수임은 명백하다. 역으로 $u$가 이 성질을
갖는다고 하자. $u=v\circ h$가 되게 하는
$v\in\operatorname{Hom}(V,W)$가 유일하게 존재함을 보이자. $h_0$는
$U$에서 $U/R$로 가는 표준 사상과 동일시되므로, 가정에 의해
$u_0=v_0\circ h_0$를 만족하는 연속 사상 $v_0:V\to W$의 존재성과
유일성이 따른다. 다른 한편 필요하면 $V$를 그와 동형인 준스킴으로
바꾸어 $\lambda$가 항등사상이라고 가정할 수 있다~; 가정에 따라 이때
$\mu$는 다음 준동형이다.
\[
  \mu:\mathcal{O}_W\longrightarrow(u_0)_*(\mathcal{O}_U)
  =(v_0)_*((h_0)_*(\mathcal{O}_U))
\]
그리고 이에 대응하는 준동형
$\mu^\sharp:u_0^*(\mathcal{O}_W)\to\mathcal{O}_U$는 모든 점에서 국소
준동형이다. 그런데
\[
  (v_0)_*((h_0)_*(\mathcal{O}_U))=(v_0)_*(\mathcal{O}_V)
\]
이므로 반드시
$\nu=\mu$여야 하며, 이제 대응하는 준동형
$\nu^\sharp:v_0^*(\mathcal{O}_W)\to\mathcal{O}_V$가 모든 점에서 국소
준동형임을 보이면 된다. 모든 $y\in V$는 어떤 $x\in U$에 대하여
$h_0(x)$의 꼴이다~; $z=v_0(y)=u_0(x)$로 놓자. 그러면
(\hyperref[0.3.5.5-ko]{0, 3.5.5})에 따라 준동형 $\mu_x^\sharp$는 다음과
같이 분해된다.
\[
  \mu_x^\sharp:\mathcal{O}_z
  \xrightarrow{\nu_y^\sharp}\mathcal{O}_y
  \xrightarrow{\lambda_x^\sharp}\mathcal{O}_x.
\]
가정에 따라 $\lambda_x^\sharp$와 $\mu_x^\sharp$는 국소 준동형이다~;
따라서 $\lambda_x^\sharp$는 $\mathcal{O}_y$의 모든 가역원을
$\mathcal{O}_x$의 가역원으로 보낸다~; 만일 $\nu_y^\sharp$가
$\mathcal{O}_z$의 비가역원을

$\mathcal{O}_y$의 가역원으로 보낸다면, $\mu_x^\sharp$는 그
$\mathcal{O}_z$의 원소를 $\mathcal{O}_x$의 가역원으로 보낼 것이므로
가정에 모순이다. 이로부터 보조정리가 따른다.
\end{proof}

\begin{corollary}[8.11.2]
\label{II.8.11.2-ko}
$U$를 정역 준스킴, $V$를 정규 준스킴이라 하자. 보편적으로 닫혀 있고
쌍유리이며 라디시엘인 모든 사상 $h:U\to V$는 동형사상이다.
\end{corollary}

\begin{proof}
$h=(h_0,\lambda)$이면 가정들로부터 $h_0$가 단사인 닫힌 사상이고
$h_0(U)$가 $V$에서 조밀함이 따른다. 따라서 $h_0$는 $U$에서 $V$로 가는

\oldpage[II]{190}
\emph{위상동형}이다. 따름정리를 증명하려면
$\lambda:\mathcal{O}_V\to(h_0)_*(\mathcal{O}_U)$가 동형사상임을 보이면
충분하다~: 그러면 (\hyperref[II.8.11.1-ko]{8.11.1})을 적용할 수 있고,
각 섬유 $h_0^{-1}(x)$가 한 점으로만 이루어져 있으므로 이 보조정리에 의해
사상 (\hyperref[II.8.11.1.1-ko]{8.11.1.1})은 전단사이다~; 따라서 $h$는
동형사상이다. 문제는 명백히 $V$ 위에서 국소적이므로, 정수적으로 닫힌
정역 $A$에 대하여 $V=\operatorname{Spec}(A)$라고 가정할 수 있다
(\hyperref[II.8.8.6.1-ko]{8.8.6.1})~; 그러면 $h$는 준동형
$\varphi:A\to\Gamma(U,\mathcal{O}_U)$에 대응하고
(\hyperref[I.2.2.4-ko]{I, 2.2.4}), 모든 것은 $\varphi$가 동형사상임을
보이는 것으로 환원된다. 그런데 $K$가 $A$의 분수체이면, 가정에 따라
$\Gamma(U,\mathcal{O}_U)$도 $K$를 분수체로 가지며 $A$는
$\Gamma(U,\mathcal{O}_U)$의 부분환이고, $\varphi$는 표준 단사 준동형이다
(\hyperref[I.8.2.7-ko]{I, 8.2.7}). 사상 $h$가
(\hyperref[II.7.3.11-ko]{7.3.11})의 가정들을 만족하므로,
$\Gamma(U,\mathcal{O}_U)$는 $K$ 안에서 $A$의 정수적 폐포의 부분환이다.
따라서 가정에 의해 이는 $A$와 같다.
\end{proof}

\begin{remark}[8.11.3]
\label{II.8.11.3-ko}
제III장에서 (\agref{III.4.4.11-ko}{III, 4.4.11}), $V$가
\emph{국소 뇌터} 준스킴이면 고유이고 준유한인 모든 사상 $h:U\to V$
(특히 (\hyperref[II.8.11.2-ko]{8.11.2})의 가정들을 만족하는 모든 사상)는
반드시 \emph{유한}임을 보일 것이다. 그러므로 이 경우
(\hyperref[II.8.11.2-ko]{8.11.2})의 결론은
(\hyperref[II.6.1.15-ko]{6.1.15})에서 따른다.
\end{remark}

\begin{env}[8.11.4]
\label{II.8.11.4-ko}
이제 그라우어트 판정법 (\hyperref[II.8.9.1-ko]{8.9.1})에서 준스킴 $C$와
«~수축~» $q$가 본질적으로 유일하게 결정된다고 흔히 단언할 수 있음을
보이겠다.
\end{env}

\begin{lemma}[8.11.5]
\label{II.8.11.5-ko}
$Y$를 준스킴, $p:X\to Y$를 고유 사상, $\mathcal{L}$을 $p$-풍부한 가역
$\mathcal{O}_X$-가군, $C$를 $Y$-준스킴, $\varepsilon:Y\to C$를
$Y$-단면, $q:V=\mathbf{V}(\mathcal{L})\to C$를 $Y$-사상이라 하고, 도식
(\hyperref[II.8.9.1.1-ko]{8.9.1.1})이 가환한다고 하자. 또한
$p=(p_0,\theta)$일 때
$\theta:\mathcal{O}_Y\to p_*(\mathcal{O}_X)$가 동형사상이라고 가정하자.
이때
$\mathcal{S}'=\bigoplus_{n\geqslant0}p_*(\mathcal{L}^{\otimes n})$,
$C'=\operatorname{Spec}(\mathcal{S}')$라 하고,
$q':\mathbf{V}(\mathcal{L})\to C'$을 표준 $Y$-사상
(\hyperref[II.8.8.5-ko]{8.8.5})이라 하자. 그러면 $q=u\circ q'$을
만족하는 $Y$-사상 $u:C'\to C$가 유일하게 존재한다.
\end{lemma}

\begin{proof}
$\theta$에 대한 가정은 특히 $p$가 전사임을 함의한다~;
(\hyperref[II.8.8.4-ko]{8.8.4})에 따라 $q'$을
$\mathbf{V}(\mathcal{L})-j(X)$에 제한하면
$C'-\varepsilon'(Y)$ 위의 \emph{동형사상}이고($\varepsilon'$은 $C'$의
꼭짓점 단면이다), (\hyperref[II.8.8.4-ko]{8.8.4})에서 $q'$이
\emph{고유이고 전사}임이 따른다~; 또한
(\hyperref[II.8.8.6-ko]{8.8.6})에 따라 $q'=(q'_0,\tau)$로 놓으면
$\tau:\mathcal{O}_{C'}\to q'_*(\mathcal{O}_V)$는 동형사상이다. 따라서
보조정리 (\hyperref[II.8.11.1-ko]{8.11.1})을 적용할 수 있으며, 모든
$z'\in C'$에 대하여 $q$가 섬유 $q'{}^{-1}(z')$ 위에서 상수임을 보이면
보조정리가 증명된다. $z'\notin\varepsilon'(Y)$일 때에는 이 조건이
자명하다. 한편 $z'\in\varepsilon'(Y)$이면
$z'=\varepsilon'(y)$를 만족하는 $y\in Y$가 유일하게 존재한다~;
(\hyperref[II.8.8.5.2-ko]{8.8.5.2})의 가환성과 $q'$이
$\mathbf{V}(\mathcal{L})-j(X)$를 $C'-\varepsilon'(Y)$ 안으로 보낸다는
사실에 따라 $q'{}^{-1}(z')=j(p^{-1}(y))$이다~; 따라서 도식
(\hyperref[II.8.9.1.1-ko]{8.9.1.1})의 가환성이 우리의 주장을 증명한다.
\end{proof}

\begin{corollary}[8.11.6]
\label{II.8.11.6-ko}
(\hyperref[II.8.11.5-ko]{8.11.5})의 가정 아래에서, 더 나아가 $q$가
고유이고 $q$을 $\mathbf{V}(\mathcal{L})-j(X)$에 제한하면
$C-\varepsilon(Y)$ 위의 동형사상이라고 하자. 그러면 사상 $u$는
보편적으로 닫혀 있고 전사이며 라디시엘이고, $u$를
$C'-\varepsilon'(Y)$에 제한하면 $C-\varepsilon(Y)$ 위의 동형사상이다.
\end{corollary}

\begin{proof}
$q'$은 $\mathbf{V}(\mathcal{L})-j(X)$에서 $C'-\varepsilon'(Y)$로 가는
동형사상이므로 (\hyperref[II.8.8.4-ko]{8.8.4}), 마지막 주장은
$q=u\circ q'$에서 곧바로 따른다. 또한 도식들

\oldpage[II]{191}
(\hyperref[II.8.8.5.2-ko]{8.8.5.2})와
(\hyperref[II.8.9.1.1-ko]{8.9.1.1})의 가환성으로부터, $u$를 $C'$의
닫힌 부분준스킴 $\varepsilon'(Y)$에 제한하면 $C$의 닫힌 부분준스킴
$\varepsilon(Y)$ 위의 동형사상임이 따른다. 따라서 모든
$z'\in\varepsilon'(Y)$에 대하여 $z=u(z')$라 놓으면, $u$는 $k(z)$에서
$k(z')$로 가는 동형사상을 정의한다. 이 관찰들은 $u$가 전단사이고
라디시엘임을 증명한다~; 또한 $\psi$와 $\psi'$이 구조 사상
$C\to Y$, $C'\to Y$이면 $\psi'=\psi\circ u$이고, $\psi'$이 분리 사상이므로
(\hyperref[II.1.2.4-ko]{1.2.4}), $u$도 분리 사상이다
(\hyperref[I.5.5.1-ko]{I, 5.5.1, (v)}). 한편 보조정리
(\hyperref[II.8.11.5-ko]{8.11.5})의 증명에서 $q'$이 전사임을 보였다~;
$q=u\circ q'$이 고유이므로, 끝으로 (\hyperref[II.5.4.3-ko]{5.4.3})과
(\hyperref[II.5.4.9-ko]{5.4.9})에서 $u$가 보편적으로 닫혀 있음이 따른다.
\end{proof}

\begin{proposition}[8.11.7]
\label{II.8.11.7-ko}
$Y$를 준스킴, $X$를 정역 준스킴, $p:X\to Y$를 고유 사상,
$\mathcal{L}$을 $p$-풍부한 가역 $\mathcal{O}_X$-가군, $C$를 정규
$Y$-준스킴, $\varepsilon:Y\to C$를 $Y$-단면,
$q:V=\mathbf{V}(\mathcal{L})\to C$를 $Y$-사상이라 하고, 도식
(\hyperref[II.8.9.1.1-ko]{8.9.1.1})이 가환한다고 하자. 또한
$p=(p_0,\theta)$일 때 $\theta:\mathcal{O}_Y\to p_*(\mathcal{O}_X)$가
동형사상이라고 가정하자. 이때
$\mathcal{S}'=\bigoplus_{n\geqslant0}p_*(\mathcal{L}^{\otimes n})$,
$C'=\operatorname{Spec}(\mathcal{S}')$라 하고,
$q':\mathbf{V}(\mathcal{L})\to C'$을 표준 $Y$-사상
(\hyperref[II.8.8.5-ko]{8.8.5})이라 하자. 그러면 $q=u\circ q'$을
만족하는 유일한 $Y$-사상 $u:C'\to C$는 동형사상이다.
\end{proposition}

\begin{proof}
(\hyperref[II.8.8.6-ko]{8.8.6})에서 $C'$이 정역 준스킴임이 따른다~;
$u$는 바탕공간의 위상동형사상 $C'\to C$이므로
((\hyperref[II.8.11.6-ko]{8.11.6})에 따라 $u$는 전단사이고 닫힌
사상이다), $C$는 기약이고 따라서 정역 준스킴이다. 또한 $C'$의
공집합이 아닌 열린집합에 제한한 $u$가 $C$의 열린집합 위의
동형사상이므로 $u$는 쌍유리 사상이다. $C$가 정규라고 가정했으므로,
(\hyperref[II.8.11.2-ko]{8.11.2})을 적용하면 결론을 얻는다.
\end{proof}

\begin{remark}[8.11.8]
\label{II.8.11.8-ko}
\begin{enumerate}
  \item[\rm{(i)}] $C$가 정규라는 가정은 $X$도 정규임을 함의한다는 점에
  유의하자. 실제로 $C'=\operatorname{Spec}(\mathcal{S}')$은 $C$와
  동형이므로 정규이고, (\hyperref[II.8.8.6-ko]{8.8.6})에 따라
  정역이다~; 이로부터 $\operatorname{Proj}(\mathcal{S}')$가
  \emph{정규}임을 얻는다. 실제로 이 문제는 $Y$ 위에서 국소적이다~;
  $Y$가 아핀이면 $\mathcal{S}'=\widetilde{S'}$이고, 환
  $S'=\Gamma(C',\mathcal{O}_{C'})$은 정역이며 정수적으로 닫혀 있다
  (\hyperref[II.8.8.6.1-ko]{8.8.6.1}). 따라서 모든 동차 원소
  $f\in S'_+$에 대하여 등급환 $S'_f$는 정역이며 정수적으로 닫혀 있고
  ([13], t.~I, p.~257 및 261), 그러므로 그 차수 0인 항들로 이루어진 환
  $S'_{(f)}$도 그러하다. 실제로 $S'_f$와 $S'_{(f)}$의 분수체의 교집합은
  $S'_{(f)}$와 같다~; 이는 우리의 주장을 증명한다
  (\hyperref[II.6.3.4-ko]{6.3.4}). 끝으로 $X$는
  $\operatorname{Proj}(\mathcal{S}')$의 열린 부분준스킴과 동형이므로
  (\hyperref[II.8.8.1-ko]{8.8.1}), $X$는 실제로 정규이다. 따라서 명제
  (\hyperref[II.8.11.7-ko]{8.11.7})을 다음과 같은 형태로 나타낼 수 있다~:
  \emph{$X$가 정역이고 정규이며,
  $p=(p_0,\theta):X\to Y$가
  $\theta:\mathcal{O}_Y\to p_*(\mathcal{O}_X)$가 동형사상이 되게 하는
  고유 사상이면, 모든 $p$-풍부한 $\mathcal{O}_X$-가군
  $\mathcal{L}$에 대하여 $V=\mathbf{V}(\mathcal{L})$의 영 단면을
  수축하여 정규 $Y$-스킴 $C$와 고유 $Y$-사상 $q:V\to C$를 얻는 방법이
  하나 존재하고 오직 하나뿐이다.}
  \item[\rm{(ii)}] $p$가 고유일 때
  $p_*(\mathcal{O}_X)=\mathcal{O}_Y$라는 가정은 보조적인 가정으로
  볼 수 있으며, 결과의 일반성을 실제로 제한하지 않는다. 실제로 이
  가정이 성립하지 않으면 $Y$를 $Y$-스킴
  $Y'=\operatorname{Spec}(p_*(\mathcal{O}_X))$으로 바꾸고 $X$를
  $Y'$-스킴으로 보면 충분하다. 이 일반적인 방법은 제III장 \S~4에서
  다시 다룬다.
\end{enumerate}
\end{remark}
\subsection{사영하는 원뿔 위의 준연접층}
\label{subsection:II.8.12-ko}

\begin{env}[8.12.1]
\label{II.8.12.1-ko}
(\hyperref[II.8.3.1-ko]{8.3.1})의 가정과 표기를 다시 사용하자.
$\mathcal{M}$을 준연접 등급 $\mathcal{S}$-가군이라 하자~; 혼동을
피하기 위하여 $\mathcal{M}$을 \emph{등급화되지 않은}
$\mathcal{S}$-가군으로 볼 때 $\mathcal{M}$에 결부된 준연접
$\mathcal{O}_C$-가군을 $\widetilde{\mathcal{M}}$으로 나타내고
(\hyperref[II.1.4.3-ko]{1.4.3}),

\oldpage[II]{192}
$\mathcal{M}$을 등급 $\mathcal{S}$-가군으로 볼 때 $\mathcal{M}$에
결부된 준연접 $\mathcal{O}_X$-가군을
$\operatorname{Proj}_0(\mathcal{M})$으로 나타낸다 (다시 말해 이는
(\hyperref[II.3.2.2-ko]{3.2.2})에서
$\widetilde{\mathcal{M}}$으로 표기한 $\mathcal{O}_X$-가군이다).
또한 다음과 같이 놓는다.
\begin{equation}
\label{II.8.12.1.1-ko}
  \mathcal{M}_X=\operatorname{Proj}(\mathcal{M})
  =\bigoplus_{n\in\mathbf{Z}}\operatorname{Proj}_0(\mathcal{M}(n))~;
\tag{8.12.1.1}
\end{equation}
준연접 등급 $\mathcal{O}_X$-대수 $\mathcal{S}_X$는
(\hyperref[II.8.6.1.1-ko]{8.6.1.1})로 정의되어 있으므로,
$\operatorname{Proj}(\mathcal{M})$에는 표준 준동형사상
(\hyperref[II.3.2.6.1-ko]{3.2.6.1})
\begin{equation}
\label{II.8.12.1.2-ko}
  \mathcal{O}_X(m)\otimes_{\mathcal{O}_X}
  \operatorname{Proj}_0(\mathcal{M}(n))
  \to \operatorname{Proj}_0(\mathcal{S}(m)\otimes_{\mathcal{S}}\mathcal{M}(n))
  \to \operatorname{Proj}_0(\mathcal{M}(m+n))
\tag{8.12.1.2}
\end{equation}
을 통하여 등급 $\mathcal{S}_X$-가군의 구조가 주어진다 (이 구조는
준연접이다). 가군 공리의 확인에는 가환 도식
(\hyperref[II.2.5.11.4-ko]{2.5.11.4})을 사용한다.

$Y=\operatorname{Spec}(A)$가 아핀이고 $\mathcal{S}=\widetilde{S}$,
$\mathcal{M}=\widetilde{M}$이라 하자. 여기서 $S$는 등급 $A$-대수이고
$M$은 등급 $S$-가군이다. 그러면 모든 동차 원소 $f\in S_+$에 대하여
정의와 (\hyperref[II.8.2.9.1-ko]{8.2.9.1})에 따라 다음이 성립한다.
\begin{equation}
\label{II.8.12.1.3-ko}
  \Gamma(X_f,\operatorname{Proj}(\widetilde{M}))=M_f
\tag{8.12.1.3}
\end{equation}

이제 다음 준연접 등급 $\widehat{\mathcal{S}}$-가군을 생각하자.
\begin{equation}
\label{II.8.12.1.4-ko}
  \widehat{\mathcal{M}}
  =\mathcal{M}\otimes_{\mathcal{S}}\widehat{\mathcal{S}}
\tag{8.12.1.4}
\end{equation}
(여기서 $\widehat{\mathcal{S}}$는
(\hyperref[II.8.3.1.1-ko]{8.3.1.1})로 정의된다)~; 이로부터 준연접
등급 $\mathcal{O}_{\widehat{C}}$-가군
$\operatorname{Proj}_0(\widehat{\mathcal{M}})$을 얻으며, 이를 다음과
같이 쓰기도 한다.
\begin{equation}
\label{II.8.12.1.5-ko}
  \mathcal{M}^{\square}=\operatorname{Proj}_0(\widehat{\mathcal{M}}).
\tag{8.12.1.5}
\end{equation}
(\hyperref[II.3.2.4-ko]{3.2.4})에 따라 $\mathcal{M}^{\square}$은
$\mathcal{M}$에 대한 가법 완전 함자이고, 직합과
귀납극한에 가환함이 명백하다.
\end{env}

\begin{proposition}[8.12.2]
\label{II.8.12.2-ko}
(\hyperref[II.8.3.2-ko]{8.3.2})의 표기 아래 다음 함자적인 표준
동형사상들이 성립한다.
\begin{equation}
\label{II.8.12.2.1-ko}
  i^*(\mathcal{M}^{\square})\xrightarrow{\sim}\widetilde{\mathcal{M}},
  \qquad
  j^*(\mathcal{M}^{\square})\xrightarrow{\sim}
  \operatorname{Proj}_0(\mathcal{M}).
\tag{8.12.2.1}
\end{equation}
\end{proposition}

\begin{proof}
실제로 $\operatorname{Spec}(\widehat{\mathcal{S}}/(z-1)\widehat{\mathcal{S}})$
위에서 $i^*(\mathcal{M}^{\square})$은
$\left(\widehat{\mathcal{M}}/(z-1)\widehat{\mathcal{M}}\right)^{\sim}$과
표준적으로 동일시된다
(\hyperref[II.3.2.3-ko]{3.2.3})~; 따라서 표준 동형사상
$\widehat{\mathcal{M}}/(z-1)\widehat{\mathcal{M}}
\xrightarrow{\sim}\mathcal{M}$으로부터
(\hyperref[II.1.4.1-ko]{1.4.1})에 의해 표준 동형사상
(\hyperref[II.8.12.2.1-ko]{8.12.2.1}) 가운데 첫 번째 것이 곧바로
따른다. 한편 표준 몰입 사상 $j:X\to\widehat{C}$는 핵이
$z\widehat{\mathcal{S}}$인 표준 준동형사상
$\widehat{\mathcal{S}}\to\mathcal{S}$에 대응한다
(\hyperref[II.8.3.2-ko]{8.3.2})~; 준동형사상
(\hyperref[II.8.12.2.1-ko]{8.12.2.1}) 가운데 두 번째 것은 표준
준동형사상 (\hyperref[II.3.5.2-ko]{3.5.2, (ii)})의 특수한 경우이다.
실제로 여기서는
$\widehat{\mathcal{M}}\otimes_{\widehat{\mathcal{S}}}\mathcal{S}
=\mathcal{M}$이다~; 이것이 동형사상임을 확인하기 위해서는
$Y=\operatorname{Spec}(A)$가 아핀이고
$\mathcal{S}=\widetilde{S}$, $\mathcal{M}=\widetilde{M}$인 경우로
한정할 수 있다~; 이때 (\hyperref[II.2.8.8-ko]{2.8.8})을 참조하면,
모든 동차 원소 $f\in S_+$에 대하여 앞의 준동형사상을 $X_f$에 제한한
것이 동형사상으로 귀착됨은 즉시 확인된다.
\end{proof}

\oldpage[II]{193}
용어를 남용하여, 첫 번째 동형사상
(\hyperref[II.8.12.2.1-ko]{8.12.2.1})이 존재한다는 이유로
$\mathcal{M}^{\square}$을 $\mathcal{O}_C$-가군
$\widetilde{\mathcal{M}}$의 \emph{사영 폐포}라고도 한다 (이때
$\mathcal{O}_C$-가군 $\widetilde{\mathcal{M}}$의 데이터에는
$\mathcal{S}$-가군 $\mathcal{M}$의 등급이 포함되어 있다고
암묵적으로 전제한다).

\begin{env}[8.12.3]
\label{II.8.12.3-ko}
(\hyperref[II.8.3.5-ko]{8.3.5})의 표기 아래에서 함자적인 표준 준동형
\begin{equation}
\label{II.8.12.3.1-ko}
  p^*(\operatorname{Proj}_0(\mathcal{M}))
  \to\mathcal{M}^{\square}|\widehat{E}.
\tag{8.12.3.1}
\end{equation}
이 존재한다. 실제로 이는 (\hyperref[II.3.5.6-ko]{3.5.6})에서 일반적으로
정의한 준동형 $u^{\sharp}$의 특수한 경우이다.
$Y=\operatorname{Spec}(A)$가 아핀이고, $\mathcal{S}=\widetilde{S}$,
$\mathcal{M}=\widetilde{M}$이면, (\hyperref[II.2.8.8-ko]{2.8.8})을
참조하여 (\hyperref[II.8.12.3.1-ko]{8.12.3.1})을
$p^{-1}(X_f)=\widehat{C}_f$에 제한한 사상은(여기서 $f$는 $S_+$의
동차 원소이다) 표준 준동형
\begin{equation}
\label{II.8.12.3.2-ko}
  M_{(f)}\otimes_{S_{(f)}}S_f^{\leq}\to M_f^{\leq}
\tag{8.12.3.2}
\end{equation}
에 대응함을 알 수 있다. 여기서는
(\hyperref[II.8.2.3.2-ko]{8.2.3.2})와
(\hyperref[II.8.2.5.2-ko]{8.2.5.2})를 고려하였다.
\end{env}

\begin{env}[8.12.4]
\label{II.8.12.4-ko}
(\hyperref[II.8.5.1-ko]{8.5.1})의 가정 아래에 놓이고 그 표기를
유지하자. (\hyperref[II.1.5.6-ko]{1.5.6})에 따르면 모든 준연접 등급
$\mathcal{S}$-가군 $\mathcal{M}$에 대하여 한편으로는
$\mathcal{O}_{C'}$-가군의 표준 동형사상
\begin{equation}
\label{II.8.12.4.1-ko}
  \Phi^*(\widetilde{\mathcal{M}})
  \xrightarrow{\sim}
  \left(q^*(\mathcal{M})\otimes_{q^*(\mathcal{S})}\mathcal{S}'\right)^{\sim}
\tag{8.12.4.1}
\end{equation}
이 존재한다. 다른 한편 (\hyperref[II.3.5.6-ko]{3.5.6})에 따르면
표준 $\operatorname{Proj}(\varphi)$-사상
\begin{equation}
\label{II.8.12.4.2-ko}
  \operatorname{Proj}_0(\mathcal{M})\to
  \left(\operatorname{Proj}_0
  (q^*(\mathcal{M})\otimes_{q^*(\mathcal{S})}\mathcal{S}')\right)
  |G(\varphi)
\tag{8.12.4.2}
\end{equation}
과 표준 $\widehat{\Phi}$-사상
\begin{equation}
\label{II.8.12.4.3-ko}
  \operatorname{Proj}_0(\widehat{\mathcal{M}})\to
  \left(\operatorname{Proj}_0
  (q^*(\widehat{\mathcal{M}})
  \otimes_{q^*(\widehat{\mathcal{S}})}\widehat{\mathcal{S}'})\right)
  |G(\widehat{\varphi}).
\tag{8.12.4.3}
\end{equation}
도 존재한다.
\end{env}

\begin{env}[8.12.5]
\label{II.8.12.5-ko}
이제 같은 표기로 (\hyperref[II.8.6.1-ko]{8.6.1})의 상황을 생각하자.
그러므로 앞에서 $Y'=X$로 취하고, $q:X\to Y$는 구조 사상이며,
$\varphi$는 표준 $q$-사상
(\hyperref[II.8.6.1.2-ko]{8.6.1.2})이다. 그러면 표준 동형사상
\begin{equation}
\label{II.8.12.5.1-ko}
  q^*(\mathcal{M})\otimes_{q^*(\mathcal{S})}\mathcal{S}_X^{\geq}
  \xrightarrow{\sim}\mathcal{M}_X^{\geq}
\tag{8.12.5.1}
\end{equation}
이 존재한다. 여기서
$\mathcal{M}_X^{\geq}=\bigoplus_{n\geq0}\operatorname{Proj}_0(\mathcal{M}(n))$로
놓는다. 실제로 $Y=\operatorname{Spec}(A)$가 아핀이고,
$\mathcal{S}=\widetilde{S}$이며 $\mathcal{M}=\widetilde{M}$인 경우로
한정할 수 있다. 각 아핀 열린집합 $X_f$에서($f$는 $S_+$의 동차
원소이다) 동형사상 (\hyperref[II.8.12.5.1-ko]{8.12.5.1})을 정의하고,
$f$를 그 동차 배수로 바꾸는 것과 양립함을 확인하면 된다. 그런데
(\hyperref[II.8.12.5.1-ko]{8.12.5.1})의 좌변을 $X_f$에 제한한 것은
$\widetilde{M}'=
\left((M\otimes_A S_{(f)})
\otimes_{S\otimes_A S_{(f)}}S_f^{\geq}\right)^{\sim}$이다
(\hyperref[II.8.6.2.1-ko]{8.6.2.1}).
$M\otimes_A S_{(f)}$에서
$M\otimes_S(S\otimes_A S_{(f)})$ 위로 가는 표준 동형사상이 있으므로,
이로부터 $\widetilde{M}'$에서
$\left(M\otimes_S S_f^{\geq}\right)^{\sim}$ 위로 가는 동형사상을
얻는다. 후자는 (\hyperref[II.8.2.9.1-ko]{8.2.9.1})에 의해
(\hyperref[II.8.12.5.1-ko]{8.12.5.1})의 우변을 $X_f$에 제한한 것과
표준적으로 동형이고, 요구되는 양립 조건을 만족한다.

\oldpage[II]{194}
$\mathcal{M}$을 $\widehat{\mathcal{M}}$으로, $\mathcal{S}$를
$\widehat{\mathcal{S}}$으로, $\mathcal{S}_X$를
$(\mathcal{S}_X^{\geq})^{\wedge}$로 바꾸어 앞의 논증을 적용하면,
마찬가지로 표준 동형사상
\begin{equation}
\label{II.8.12.5.2-ko}
  q^*(\widehat{\mathcal{M}})
  \otimes_{q^*(\widehat{\mathcal{S}})}(\mathcal{S}_X^{\geq})^{\wedge}
  \xrightarrow{\sim}(\mathcal{M}_X^{\geq})^{\wedge}.
\tag{8.12.5.2}
\end{equation}
을 얻는다. 구조 사상
$u:\operatorname{Proj}(\mathcal{S}_X^{\geq})\to X$가 동형사상이라는
(\hyperref[II.8.6.2-ko]{8.6.2})의 사실을 상기하면, 먼저 앞의 결과에서
표준 $u$-동형사상
\begin{equation}
\label{II.8.12.5.3-ko}
  \operatorname{Proj}_0\mathcal{M}
  \xrightarrow{\sim}\operatorname{Proj}_0(\mathcal{M}_X^{\geq})
\tag{8.12.5.3}
\end{equation}
을 (\hyperref[II.8.12.4.2-ko]{8.12.4.2})의 특수한 경우로 얻는다.
실제로 (\hyperref[II.8.6.2-ko]{8.6.2})의 증명에 사용한 표기 아래에서,
이는 표준 준동형
$M_{(f)}\otimes_{S_{(f)}}(S_f^{\geq})^{(d)}\to(M_f^{\geq})^{(d)}$가
$f\in S_d$일 때
동형사상임을 확인하는 것과 같은데, 이는 자명하다.

둘째로 동형사상 (\hyperref[II.8.12.5.2-ko]{8.12.5.2})을 이용하여,
이번에는 (\hyperref[II.8.12.4.3-ko]{8.12.4.3})을 표준 사상
$r=\operatorname{Proj}(\widehat{\alpha}):\widehat{C}_X\to
\widehat{C}$에 적용하면 표준 $r$-사상
\begin{equation}
\label{II.8.12.5.4-ko}
  \mathcal{M}^{\square}\to(\mathcal{M}_X^{\geq})^{\square}.
\tag{8.12.5.4}
\end{equation}
을 얻는다. 이제 (\hyperref[II.8.6.2-ko]{8.6.2})에 따라 $r$을
꼭지점을 제거한 원뿔 $\widehat{E}_X$와 $E_X$에 제한한 사상은 각각
$\widehat{E}$와 $E$ 위로 가는 \emph{동형사상}임을 상기하자. 또한
다음이 성립한다.
\end{env}

\begin{proposition}[8.12.6]
\label{II.8.12.6-ko}
$\widehat{E}_X$와 $E_X$에 대한 표준 $r$-사상
(\hyperref[II.8.12.5.4-ko]{8.12.5.4})의 제한은 다음 동형사상들이다.
\begin{equation}
\label{II.8.12.6.1-ko}
  \mathcal{M}^{\square}|\widehat{E}
  \xrightarrow{\sim}(\mathcal{M}_X^{\geq})^{\square}|\widehat{E}_X
\tag{8.12.6.1}
\end{equation}
\begin{equation}
\label{II.8.12.6.2-ko}
  \widetilde{\mathcal{M}}|E
  \xrightarrow{\sim}(\mathcal{M}_X^{\geq})^{\sim}|E_X.
\tag{8.12.6.2}
\end{equation}
\end{proposition}

\begin{proof}
(\hyperref[II.8.6.2-ko]{8.6.2})의 증명에서와 같이 $Y$가 아핀인
경우로 환원한다. 그 증명의 표기를 사용하고 정의
(\hyperref[II.2.8.8-ko]{2.8.8})로 돌아가면, 표준 준동형
\[
  \widehat{M}_{(f)}\otimes_{\widehat{S}_{(f)}}
  (S_f^{\geq})_{(f/1)}^{\wedge}
  \to(M\otimes_S S_f^{\geq})_{(f/1)}^{\wedge}
\]
이 동형사상임을 보여야 한다. 그런데
(\hyperref[II.8.2.3.2-ko]{8.2.3.2})와
(\hyperref[II.8.2.5.2-ko]{8.2.5.2})에 따라 좌변은 표준적으로
$M_f^{\leq}\otimes_{S_f^{\leq}}(S_f^{\geq})_{f/1}^{\leq}$와 동일시되고,
따라서 (\hyperref[II.8.2.7.2-ko]{8.2.7.2})에 의해 $M_f^{\leq}$와
동일시된다. 우변은 $(M_f^{\geq})_{f/1}^{\leq}$와 동일시되고,
따라서 (\hyperref[II.8.2.9.2-ko]{8.2.9.2})에 의해 역시
$M_f^{\leq}$와 동일시된다. 이로부터
(\hyperref[II.8.12.6.1-ko]{8.12.6.1})에 관한 결론이 따른다.
그러면 식 (\hyperref[II.8.12.6.2-ko]{8.12.6.2})는
(\hyperref[II.8.12.6.1-ko]{8.12.6.1})과
(\hyperref[II.8.12.2.1-ko]{8.12.2.1})에서 따른다.
\end{proof}

\begin{corollary}[8.12.7]
\label{II.8.12.7-ko}
(\hyperref[II.8.6.3-ko]{8.6.3})의 동일시 아래에서,
$(\mathcal{M}_X^{\geq})^{\square}$을 $\widehat{E}_X$에 제한한 것은
$(\mathcal{M}_X^{\leq})^{\sim}$와 동일시되고,
$(\mathcal{M}_X^{\geq})^{\square}$을 $E_X$에 제한한 것은
$\widetilde{\mathcal{M}}_X$와 동일시된다.
\end{corollary}

\begin{proof}
실제로 아핀인 경우로 환원되며, 이는 $(M_f^{\geq})_{f/1}^{\leq}$와
$M_f^{\leq}$의 동일시 및 $(M_f^{\geq})_{f/1}$와 $M_f$의 동일시에서
따른다 (\hyperref[II.8.2.9.2-ko]{8.2.9.2}).
\end{proof}

\begin{proposition}[8.12.8]
\label{II.8.12.8-ko}
(\hyperref[II.8.6.4-ko]{8.6.4})의 가정 아래에서 표준 준동형
(\hyperref[II.8.12.3.1-ko]{8.12.3.1})은 동형사상이다.
\end{proposition}

\begin{proof}
$\operatorname{Proj}(\mathcal{S}_X^{\geq})\to X$가 동형사상이라는 사실
(\hyperref[II.8.6.2-ko]{8.6.2})과

\oldpage[II]{195}
동형사상들 (\hyperref[II.8.12.5.4-ko]{8.12.5.4}) 및
(\hyperref[II.8.12.6.1-ko]{8.12.6.1})을 고려하면, 표준 준동형
$p_X^*(\operatorname{Proj}_0(\mathcal{M}_X^{\geq}))
\to(\mathcal{M}_X^{\geq})^{\square}|\widehat{E}_X$에 대한 대응하는
명제를 증명하는 것으로 귀착된다. 다시 말해, $\mathcal{S}_1$이 가역
$\mathcal{O}_Y$-가군이고 $\mathcal{S}$가 $\mathcal{S}_1$에 의해
생성되는 경우로 귀착된다. (\hyperref[II.8.12.3-ko]{8.12.3})의 표기법을
쓰면, 이 경우 $f\in S_1$에 대하여
$S_f^{\leq}=S_{(f)}[1/f]$이고, 표준 준동형
$M_{(f)}\otimes_{S_{(f)}}S_f^{\leq}\to M_f^{\leq}$는
$M_f^{\leq}$의 정의에 의해 동형사상이다.
\end{proof}

\begin{env}[8.12.9]
\label{II.8.12.9-ko}
이제 준연접 $\mathcal{S}$-가군
\[
  \mathcal{M}_{[n]}=\bigoplus_{m\geqslant n}\mathcal{M}_m
\]
들과 (\hyperref[II.8.7.2-ko]{8.7.2})의 표기법을 사용하여 준연접 등급
$\mathcal{S}^{\natural}$-가군
\begin{equation}
\label{II.8.12.9.1-ko}
  \mathcal{M}^{\natural}
  =\left(\bigoplus_{n\geqslant0}\mathcal{M}_{[n]}\right)^{\sim}.
\tag{8.12.9.1}
\end{equation}
을 생각하자. 표준 $C$-동형사상
$h:C_X\xrightarrow{\sim}\operatorname{Proj}(\mathcal{S}^{\natural})$이
존재함은 이미 보았다 (\hyperref[II.8.7.3-ko]{8.7.3}). 더 나아가 다음이
성립한다:
\end{env}

\begin{proposition}[8.12.10]
\label{II.8.12.10-ko}
표준 $h$-동형사상
\begin{equation}
\label{II.8.12.10.1-ko}
  \operatorname{Proj}_0(\mathcal{M}^{\natural})
  \xrightarrow{\sim}\widetilde{\mathcal{M}}_X.
\tag{8.12.10.1}
\end{equation}
이 존재한다.
\end{proposition}

\begin{proof}
(\hyperref[II.8.7.3-ko]{8.7.3})에서와 같이 논증하되, 이번에는
(\hyperref[II.8.2.7.3-ko]{8.2.7.3}) 대신 di-동형사상
(\hyperref[II.8.2.9.3-ko]{8.2.9.3})의 존재를 사용한다. 세부 사항은
독자에게 맡긴다.
\end{proof}

\subsection{부분층과 닫힌 부분스킴의 사영 폐포}
\label{subsection:II.8.13-ko}

\begin{env}[8.13.1]
\label{II.8.13.1-ko}
가정과 표기법은 (\hyperref[II.8.12.1-ko]{8.12.1})에서와 같다고 하고,
$\mathcal{M}$의 준연접 부분-$\mathcal{S}$-가군 $\mathcal{N}$을
생각하자. 이때 $\mathcal{N}$에는 \emph{반드시 등급이 주어져 있을
필요는 없다}. 따라서 $\mathcal{N}$에 연관된 준연접
$\mathcal{O}_C$-가군 $\widetilde{\mathcal{N}}$을 생각할 수 있으며,
이는 $\widetilde{\mathcal{M}}$의 부분-$\mathcal{O}_C$-가군이다. 한편
(\hyperref[II.8.12.2.1-ko]{8.12.2.1})에서
$\widetilde{\mathcal{M}}$이 $\mathcal{M}^{\square}$의 $C$로의 제한과
동일시됨을 보았다. 표준 단사 사상 $i:C\to\widehat{C}$는 아핀 사상이고
(\hyperref[II.8.3.2-ko]{8.3.2}), \emph{더구나} 준콤팩트하므로,
$\mathcal{M}^{\square}$에 포함되고 $C$ 위에서
$\widetilde{\mathcal{N}}$을 유도하는 가장 큰
$\mathcal{O}_{\widehat{C}}$-부분가군인
\emph{표준 연장} $(\widetilde{\mathcal{N}})^{-}$는 준연접
$\mathcal{O}_{\widehat{C}}$-가군이다
(\hyperref[I.9.4.2-ko]{I, 9.4.2}). 이제 등급
$\widehat{\mathcal{S}}$-가군을 이용하여 이것을 명시적으로
기술하겠다.
\end{env}

\begin{env}[8.13.2]
\label{II.8.13.2-ko}
이를 위해 각 정수 $n\geqslant0$에 대하여 준동형
$\bigoplus_{i\leqslant n}\mathcal{M}_i\to\mathcal{M}$을 생각하자.
이 준동형은 $Y$의 각 열린집합 $U$에 대하여 족
\[
  (s_i)\in\bigoplus_{i\leqslant n}\Gamma(U,\mathcal{M}_i)
\]
에 단면 $\sum_i s_i\in\Gamma(U,\mathcal{M})$을 대응시킨다.
이 준동형에 의한 $\mathcal{N}$의 역상을 $\mathcal{N}'_n$이라 하자.
이는 $\bigoplus_{i\leqslant n}\mathcal{M}_i$의 준연접
부분-$\mathcal{O}_Y$-가군이다. 이제 $(s_i)$에 단면
$\sum_{i\leqslant n}s_i z^{n-i}\in
\Gamma(U,\widehat{\mathcal{M}}_n)$을 대응시키는 준동형
$\bigoplus_{i\leqslant n}\mathcal{M}_i\to
\widehat{\mathcal{M}}=\mathcal{M}[z]$을 생각하고, 이 준동형에 의한
$\mathcal{N}'_n$의 상을 $\mathcal{N}_n$이라 하자. 그러면
$\overline{\mathcal{N}}=\bigoplus_{n\geqslant0}\mathcal{N}_n$이
$\widehat{\mathcal{M}}$의 (준연접) 등급
부분-$\widehat{\mathcal{S}}$-가군임을 즉시 확인할 수 있다.
$\overline{\mathcal{N}}$은 “동차화 변수” $z$를 이용한
\emph{동차화}에 의해 $\mathcal{N}$으로부터 얻어진다고 말한다.

\oldpage[II]{196}
$\mathcal{N}$이 이미 $\mathcal{M}$의 \emph{등급}
부분-$\mathcal{S}$-가군이면, $\overline{\mathcal{N}}$은
$\widehat{\mathcal{N}}=\mathcal{N}[z]$에서 차수가 $n\geqslant0$인
성분 $\widehat{\mathcal{N}}_n$들의 직합과 동일시된다는 점에
유의하자.
\end{env}

\begin{proposition}[8.13.3]
\label{II.8.13.3-ko}
$\mathcal{O}_{\widehat{C}}$-가군
$\operatorname{Proj}_0(\overline{\mathcal{N}})$은
$\widetilde{\mathcal{N}}$의 $\widehat{C}$로의 표준 연장
$(\widetilde{\mathcal{N}})^{-}$이다.
\end{proposition}

\begin{proof}
표준 연장의 정의 (\hyperref[I.9.4.1-ko]{I, 9.4.1})에 따라 문제는
$Y$와 $\widehat{C}$ 위에서 국소적이다. 따라서 먼저
$Y=\operatorname{Spec}(A)$가 아핀이고
$\mathcal{S}=\widetilde{S}$, $\mathcal{M}=\widetilde{M}$,
$\mathcal{N}=\widetilde{N}$이라고 가정할 수 있다. 여기서 $N$은
$M$의 부분-$S$-가군이며 반드시 등급이 주어져 있을 필요는 없다.
또한 (\hyperref[II.8.3.2.6-ko]{8.3.2.6})에 따라 $\widehat{C}$는
아핀 열린집합 $\widehat{C}_z=C$와
$\widehat{C}_f=\operatorname{Spec}(S_f^{\leq})$
($f$는 $S_+$의 동차 원소)의 합집합이다. 그러므로 다음 두 가지를
보이면 충분하다~: $1^\circ$
$\operatorname{Proj}_0(\overline{\mathcal{N}})$의 $C$로의 제한은
$\widetilde{\mathcal{N}}$이다~; $2^\circ$
$\operatorname{Proj}_0(\overline{\mathcal{N}})$의 각
$\widehat{C}_f$로의 제한은 $\mathcal{N}$의
$C\cap\widehat{C}_f=\operatorname{Spec}(S_f)$로의 제한의 표준
연장이다 (\hyperref[II.8.3.2.6-ko]{8.3.2.6}). 첫 번째 항목을 위해
$\operatorname{Proj}_0(\overline{\mathcal{N}})|C$가
$(\overline{N}_{(z)})^{\sim}$과 동일시된다는 점에 유의하자
(\hyperref[II.8.3.2.4-ko]{8.3.2.4}). 그런데
$\overline{N}_{(z)}$는 $\widehat{M}/(z-1)\widehat{M}$ 안에서
$\overline{N}$의 상과 표준적으로 동일시되고
(\hyperref[II.2.2.5-ko]{2.2.5}), 후자에서 $M$으로 가는 표준
동형사상 (\hyperref[II.8.2.5-ko]{8.2.5})에 의해 이 상은
(\hyperref[II.8.13.2-ko]{8.13.2})에서 주어진
$\overline{N}$의 정의에 따라 $N$과 동일시된다.

둘째 주장을 증명하기 위해, 단사사상
$i:C\cap\widehat{C}_f\to\widehat{C}_f$는 표준 단사 준동형
$S_f^{\leq}\to S_f$에 대응함을 주목하자
(\hyperref[II.8.3.2.6-ko]{8.3.2.6})~; 한편
$\Gamma(\widehat{C}_f,\mathcal{M}^{\square})=M_f^{\leq}$,
$\Gamma(\widehat{C}_f,i_*(\widetilde{\mathcal{N}}))=N_f$이고
((\hyperref[II.8.12.2.1-ko]{8.12.2.1})에 의해)
$\Gamma(\widehat{C}_f,i_*(i^*(\mathcal{M}^{\square})))=M_f$이다.
(\hyperref[I.9.4.2-ko]{I, 9.4.2})를 고려하면, 따라서
$\overline{N}_{(f)}\subset\widehat{M}_{(f)}=M_f^{\leq}$가 $N_f$의 역상,
즉 표준 단사 준동형 $M_f^{\leq}\to M_f$에 의한 역상과 표준적으로
동일시됨을 보이는 것으로 환원된다. 실제로 $d=\deg(f)>0$라 하고,
$(\sum_{k\leqslant md}x_k)/f^m$라는 $M_f$의 원소
($x_k\in M_k$)가 $y/f^m$ 꼴이며 $y\in N$이라고 가정하자.
$y$와 $x_k$들을 같은 적당한 $f^h$로
곱함으로써, 이미 $\sum_{k\leqslant md}x_k=y$라고 가정할 수 있다.
그런데 동일시 (\hyperref[II.8.2.5.2-ko]{8.2.5.2})에서
$(\sum_{k\leqslant md}x_k)/f^m$는
$\sum_{k\leqslant md}x_kz^{md-k}/f^m$에 대응하고, 이 원소는 실제로
$\overline{N}_{(f)}$에 속한다. 왜냐하면
$\sum_{k\leqslant md}x_k\in N$이기 때문이다~; 역은 자명하다.
\end{proof}

\begin{remark}[8.13.4]
\label{II.8.13.4-ko}
\begin{enumerate}
  \item[\rm{(i)}] (\hyperref[II.8.13.3-ko]{8.13.3})의 가장 중요한 적용
  사례는 $\mathcal{M}=\mathcal{S}$인 경우이다. 이때
  $\widetilde{\mathcal{N}}$은 \emph{임의의} 준연접층 $\mathcal{J}$인
  $\mathcal{O}_C$의 아이디얼이고 (\hyperref[II.1.4.3-ko]{1.4.3}),
  \emph{닫힌 부분준스킴} $Z$라는 $C$의 부분준스킴과 일대일로 대응한다.
  그러면 $\overline{\mathcal{J}}$은 $\mathcal{J}$의 표준 연장이고,
  $\mathcal{O}_{\widehat{C}}$의 준연접 아이디얼층으로서 \emph{폐포}
  $\overline{Z}$, 곧 $Z$의 $\widehat{C}$에서의 폐포를 정의한다
  (\hyperref[I.9.5.10-ko]{I, 9.5.10})~; 명제
  (\hyperref[II.8.13.3-ko]{8.13.3})은 $\overline{Z}$를
  $\widehat{\mathcal{S}}=\mathcal{S}[z]$ 안의 등급 아이디얼로 정의하는
  표준적인 방법을 준다.
  \item[\rm{(ii)}] 간단히 하기 위해 $Y$가 아핀이라고 가정하고,
  (\hyperref[II.8.13.3-ko]{8.13.3})의 증명에서 사용한 표기를 유지하자.
  영이 아닌 모든 $x\in N$에 대하여, $d(x)$를 $x_i$가 $x$의,
  $M$ 안에서의 동차 성분일 때 그 차수들 가운데 가장 큰 것으로 놓자~;
  정의에 따라 $\overline{N}$은 $\widehat{M}$의 부분가군으로서 $0$과
  다음 꼴의 원소들로 이루어진다.
  $h(x,k)=z^k\sum_{i\leqslant d(x)}x_i z^{d(x)-i}$ ($k$는
  $\geqslant0$인 정수)~; 따라서 이는 다음 꼴의 원소들로
  $\widehat{S}=S[z]$ 위에서 생성된다.
  \[
    h(x,0)=\sum_{i\leqslant d(x)}x_i z^{d(x)-i}.
  \]

  \oldpage[II]{197}
  $h(x,0)$은 $x$에서 \emph{동차화}를 하여 ``동차화 변수'' $z$를
  사용해 얻어진다고 말한다. 그러나 $h(x,0)$은 $x$에 대하여 가법적이지
  않고 (따라서 \emph{더더욱} $S$-선형이지 않으므로), 심지어 $M=S$일
  때조차 $h(x,0)$들이 등급 $\widehat{S}$-가군 $\overline{N}$의 한
  \emph{생성계}를 이룬다고, $x$가 $S$-가군 $N$의 한 \emph{생성계}를
  달리게 하면서 \emph{믿지 않도록 주의해야 한다}. 하지만 $N$이
  \emph{하나의 생성원을 갖는 자유} $S$-가군인 경우(초등 대수기하학에서
  유일하게 고려되는 경우)에는 실제로 그러하다. 왜냐하면 $t$가 $N$의
  기저이면 $h(t,0)$이 $\widehat{S}$-가군 $\overline{N}$을 생성하기
  때문이다.
\end{enumerate}
\end{remark}

\subsection{등급 $\mathcal{S}$-가군에 수반되는 층에 대한 보충}
\label{subsection:II.8.14-ko}

\begin{env}[8.14.1]
\label{II.8.14.1-ko}
$Y$를 준스킴, $\mathcal{S}$를 \emph{양의 차수}를 갖는 준연접 등급
$\mathcal{O}_Y$-대수,
$X=\operatorname{Proj}(\mathcal{S})$라 하고, $q:X\to Y$를 구조 사상이라
하자 ((\hyperref[II.3.1.3-ko]{3.1.3})에 의해 분리 사상이다).
(\hyperref[II.8.12.1-ko]{8.12.1})의 표기를 유지하자. 그에 따라
$\mathcal{M}_X=\operatorname{Proj}(\mathcal{M})$을 $\mathcal{M}$에
대응시키며, 준연접 등급 $\mathcal{S}$-가군의 범주에서 준연접 등급
$\mathcal{S}_X$-가군의 범주로 가는 함자를 정의하였다~; 더구나
(\hyperref[II.3.2.4-ko]{3.2.4})에 의해 이것이
귀납극한과 가환하는 \emph{가법 완전 함자}임은 명백하다.

또한 정의 (\hyperref[II.8.12.1.1-ko]{8.12.1.1})에서 즉시 다음이
따름을 주목하자.
\begin{equation}
\label{II.8.14.1.1-ko}
  \operatorname{Proj}(\mathcal{M}(n))
  =(\operatorname{Proj}(\mathcal{M}))(n)
  \qquad\text{pour tout }n\in\mathbf{Z}.
\tag{8.14.1.1}
\end{equation}
\end{env}

\begin{env}[8.14.2]
\label{II.8.14.2-ko}
먼저 $\mathcal{S}_X$-가군 가운데 $\operatorname{Proj}(\mathcal{M})$
꼴인 것들에 (\hyperref[II.3.2.6-ko]{3.2.6})에서 정의한 표준 준동형
$\lambda$와 $\mu$를 확장하겠다. 이를 위해
$m\in\mathbf{Z}$와 $n\in\mathbf{Z}$가 무엇이든
(\hyperref[II.2.1.2.1-ko]{2.1.2.1})에 의해 다음과 같은 표준
$\mathcal{O}_X$-가군 준동형이 있음을 주목하자.
\begin{equation}
\label{II.8.14.2.1-ko}
  \lambda_{mn}:\operatorname{Proj}_0(\mathcal{M}(m))
  \otimes_{\mathcal{O}_X}\operatorname{Proj}_0(\mathcal{N}(n))
  \to\operatorname{Proj}_0
  ((\mathcal{M}\otimes_{\mathcal{S}}\mathcal{N})(m+n))
\tag{8.14.2.1}
\end{equation}
또한 마찬가지로 (\hyperref[II.2.1.2.2-ko]{2.1.2.2})에 의해 다음과 같은
표준 $\mathcal{O}_X$-가군 준동형이 있다.
\begin{equation}
\label{II.8.14.2.2-ko}
  \mu_{mn}:\operatorname{Proj}_0
  ((\mathcal{H}om_{\mathcal{S}}(\mathcal{M},\mathcal{N}))(n-m))
  \to\mathcal{H}om_{\mathcal{O}_X}
  (\operatorname{Proj}_0(\mathcal{M}(m)),
   \operatorname{Proj}_0(\mathcal{N}(n)))
\tag{8.14.2.2}
\end{equation}
여기서 임의의 준연접 등급 $\mathcal{S}$-가군 $\mathcal{M}$,
$\mathcal{N}$을 취할 수 있다. 이로부터 다음 준동형을 얻는다.
\[
  \mu_k:\operatorname{Proj}_0
  ((\mathcal{H}om_{\mathcal{S}}(\mathcal{M},\mathcal{N}))(k))
  \to
  (\mathcal{H}om_{\mathcal{S}_X}
  (\operatorname{Proj}(\mathcal{M}),\operatorname{Proj}(\mathcal{N})))_k
\]
이는 모든
$u\in\Gamma(U,\operatorname{Proj}_0
((\mathcal{H}om_{\mathcal{S}}(\mathcal{M},\mathcal{N}))(k)))$에 준동형
$\mu_k(u)$를 대응시켜 얻는다. 이것은 차수 $k$인 등급
$\mathbf{Z}$-가군의 준동형
$\Gamma(U,\operatorname{Proj}(\mathcal{M}))\to
\Gamma(U,\operatorname{Proj}(\mathcal{N}))$이다. 여기서 $U$는 $X$의
열린집합이고, 이 준동형은 각각의
$\Gamma(U,\operatorname{Proj}_0(\mathcal{M}(m)))$에서
$\mu_{m,m+k}(u)$와 일치한다~; 더구나 $\mu_{mn}$의 정의
(\hyperref[II.2.5.12.1-ko]{2.5.12.1})로 돌아가 보면,
$\mu_k(u)$가 실제로 차수 $k$인 등급 $\Gamma(U,\mathcal{S}_X)$-가군
사이의 준동형이고, 나아가 $\mu_k$들이 다음과 같은 등급
$\mathcal{S}_X$-가군의 준동형을 정의함을 즉시 알 수 있다.
\begin{equation}
\label{II.8.14.2.3-ko}
  \operatorname{Proj}
  (\mathcal{H}om_{\mathcal{S}}(\mathcal{M},\mathcal{N}))
  \to\mathcal{H}om_{\mathcal{S}_X}
  (\operatorname{Proj}(\mathcal{M}),\operatorname{Proj}(\mathcal{N})).
\tag{8.14.2.3}
\end{equation}

마찬가지로 결합법칙 도식
(\hyperref[II.2.5.11.4-ko]{2.5.11.4})을 고려하면, 준동형
(\hyperref[II.8.14.2.1-ko]{8.14.2.1})은 다음과 같은 등급
$\mathcal{S}_X$-가군의 준동형을 준다.
\begin{equation}
\label{II.8.14.2.4-ko}
  \lambda:\operatorname{Proj}(\mathcal{M})
  \otimes_{\mathcal{S}_X}\operatorname{Proj}(\mathcal{N})
  \to\operatorname{Proj}(\mathcal{M}\otimes_{\mathcal{S}}\mathcal{N}).
\tag{8.14.2.4}
\end{equation}
\end{env}

\oldpage[II]{198}
\begin{proposition}[8.14.3]
\label{II.8.14.3-ko}
준동형 (\hyperref[II.8.14.2.4-ko]{8.14.2.4})는 전단사이다~; 준동형
(\hyperref[II.8.14.2.3-ko]{8.14.2.3})도, 등급 $\mathcal{S}$-가군
$\mathcal{M}$이 유한 표시를 갖는 경우
(\hyperref[II.3.1.1-ko]{3.1.1})에는 마찬가지이다.
\end{proposition}

\begin{proof}
문제는 자명하게 $X$와 $Y$에 대해 국소적이다~; 따라서
$Y=\operatorname{Spec}(A)$가 아핀이고, $\mathcal{S}=\widetilde{S}$,
$\mathcal{M}=\widetilde{M}$, $\mathcal{N}=\widetilde{N}$이라고 가정할 수
있다. 여기서 $S$는 양의 차수를 갖는 등급 $A$-대수이고, $M$과 $N$은
두 등급 $S$-가군이다. $f$가 $S_+$의 동차 원소이면, 준동형
(\hyperref[II.8.14.2.1-ko]{8.14.2.1})과
(\hyperref[II.8.14.2.2-ko]{8.14.2.2})를 아핀 열린집합 $D_+(f)$에
제한한 것은 다음 표준 준동형들
(\hyperref[II.2.5.11.1-ko]{2.5.11.1})과
(\hyperref[II.2.5.12.1-ko]{2.5.12.1})에 대응한다~:
\begin{align*}
  M(m)_{(f)}\otimes_{S_{(f)}}N(n)_{(f)}
  &\to(M\otimes_S N)(m+n)_{(f)},\\
  (\operatorname{Hom}_S(M,N))(n-m)_{(f)}
  &\to\operatorname{Hom}_{S_{(f)}}(M(m)_{(f)},N(n)_{(f)}).
\end{align*}
이 준동형들의 정의로 돌아가 보면, 따라서
(\hyperref[II.8.2.9.1-ko]{8.2.9.1})을 고려할 때
(\hyperref[II.8.14.2.4-ko]{8.14.2.4})를 $D_+(f)$에 제한한 것은
다음 표준 준동형에 대응함을 알 수 있다.
\[
  M_f\otimes_{S_f}N_f\to(M\otimes_S N)_f
\]
이는 (\hyperref[0.1.3.4-ko]{0, 1.3.4})에서 정의되었고, 동형사상임이
알려져 있다. 마찬가지로
(\hyperref[II.8.14.2.3-ko]{8.14.2.3})을 $D_+(f)$에 제한한 것은
표준 준동형 (\hyperref[0.1.3.5-ko]{0, 1.3.5})
\[
  (\operatorname{Hom}_S(M,N))_f\to
  \operatorname{Hom}_{S_f}(M_f,N_f)
\]
에 대응한다. 여기서는 $M$이 유한형이므로
$\operatorname{Hom}_S(M,N)$, 즉 $S$-가군의 \emph{동차} 준동형들로
이루어진 부분군들의 직합
(\hyperref[II.2.1.2-ko]{2.1.2})이 $S$-가군의 \emph{모든} 준동형
$M\to N$의 집합과 일치한다는 사실을 고려하였다. 이제 $M$이 유한
표시를 갖는다는 가정과 (\hyperref[0.1.3.5-ko]{0, 1.3.5})로부터, 문제의
표준 준동형이 실제로 동형사상임이 따른다.
\end{proof}

\begin{proposition}[8.14.4]
\label{II.8.14.4-ko}
$U$가 $X$의 준콤팩트 열린집합이면, 어떤 정수 $d$가 존재하여 $d$의
배수인 모든 정수 $n$에 대해 $\mathcal{O}_X(n)|U$가 가역이고, 그 역은
$\mathcal{O}_X(-n)|U$이다.
\end{proposition}

\begin{proof}
$q(U)$는 준콤팩트이므로 유한 개의 아핀 열린집합 $V_i$로 덮인다.
따라서 모든 $x\in U$는 $D_+(f)$ 꼴의 아핀 열린집합에 포함되는데,
여기서 $f$는 환들 $\Gamma(V_i,\mathcal{S})$ 가운데 하나에 속하는 차수
$>0$인 동차 원소이다. $U$가 준콤팩트이므로 이와 같은 유한 개의
열린집합 $D_+(f_j)$로 덮을 수 있다~; $d$를 $f_j$들의 차수의
공배수라 하자. 그러면 (\hyperref[II.2.5.17-ko]{2.5.17})에 의해 이 $d$가
요구되는 조건을 만족한다.
\end{proof}

\begin{env}[8.14.5]
\label{II.8.14.5-ko}
(\hyperref[II.8.14.1-ko]{8.14.1})의 가정과 표기를 사용하자.
(\hyperref[II.3.3.2-ko]{3.3.2})에서 다음과 같은 표준
$\mathcal{O}_Y$-가군 준동형들을 정의하였다.
\begin{equation}
\label{II.8.14.5.1-ko}
  \alpha_n:\mathcal{M}_n\to
  q_*(\operatorname{Proj}_0(\mathcal{M}(n)))
  \qquad(n\in\mathbf{Z})
\tag{8.14.5.1}
\end{equation}
(\hyperref[II.3.3.1-ko]{3.3.1})의 표기를 일반화하여, 모든 등급
$\mathcal{S}_X$-가군 $\mathcal{F}$에 대해 다음과 같이 놓겠다.
\begin{equation}
\label{II.8.14.5.2-ko}
  \Gamma_*(\mathcal{F})=
  \bigoplus_{n\in\mathbf{Z}}q_*(\mathcal{F}_n).
\tag{8.14.5.2}
\end{equation}
특히
$\Gamma_*(\mathcal{S}_X)=\bigoplus_{n\in\mathbf{Z}}q_*(\mathcal{O}_X(n))$은
(\hyperref[II.3.3.1.2-ko]{3.3.1.2})에서 $\Gamma_*(\mathcal{O}_X)$로
표기한 등급 $\mathcal{O}_Y$-대수이다~; $\Gamma_*(\mathcal{F})$가 등급
$\Gamma_*(\mathcal{S}_X)$-가군임은 명백하다
(\hyperref[0.4.2.2-ko]{0, 4.2.2}). 다음의 경우,

\oldpage[II]{199}
준동형들 (\hyperref[II.8.14.5.1-ko]{8.14.5.1})에서
$\mathcal{M}=\mathcal{S}$로 놓으면, 등급 $\mathcal{O}_Y$-대수의 준동형
\begin{equation}
\label{II.8.14.5.3-ko}
  \alpha:\mathcal{S}\to\Gamma_*(\mathcal{S}_X)
\tag{8.14.5.3}
\end{equation}
을 얻는데, 이는 (\hyperref[II.3.3.2-ko]{3.3.2})에서 이미 정의한 것이며,
따라서 $\Gamma_*(\mathcal{F})$에 등급 $\mathcal{S}$-가군의 구조를
부여한다~; 그러면 준동형들
(\hyperref[II.8.14.5.1-ko]{8.14.5.1})은 (차수 $0$인) 등급
$\mathcal{S}$-가군의 준동형
\begin{equation}
\label{II.8.14.5.4-ko}
  \alpha:\mathcal{M}\to\Gamma_*(\operatorname{Proj}(\mathcal{M})).
\tag{8.14.5.4}
\end{equation}
을 정의한다.
\end{env}

\begin{env}[8.14.6]
\label{II.8.14.6-ko}
일반적으로 준연접 등급 $\mathcal{S}_X$-가군 $\mathcal{F}$에 대하여,
등급 $\mathcal{S}$-가군 $\Gamma_*(\mathcal{F})$가 반드시 준연접인지는
확실하지 않다. $X'$을 $X$의 열린집합이라 하되, $q':X'\to Y$는 $q$를
$X'$에 제한한 \emph{준콤팩트} 사상이라 하자. 더구나 $q'$은 분리이므로,
$q'_*(\mathcal{F}')$은 준연접 $\mathcal{O}_Y$-가군이다. 이는 모든 준연접
$\mathcal{O}_{X'}$-가군 $\mathcal{F}'$에 대하여 성립한다
(\hyperref[I.9.2.2-ko]{I, 9.2.2, b)}). 다음과 같이 놓겠다.
\begin{equation}
\label{II.8.14.6.1-ko}
  \mathcal{S}_{X'}=\mathcal{S}_X|X'
  =\bigoplus_{n\in\mathbf{Z}}\mathcal{O}_X(n)|X'
\tag{8.14.6.1}
\end{equation}
또한 모든 등급 $\mathcal{S}_{X'}$-가군 $\mathcal{F}'$에 대하여
\begin{equation}
\label{II.8.14.6.2-ko}
  \Gamma'_*(\mathcal{F}')=
  \bigoplus_{n\in\mathbf{Z}}q'_*(\mathcal{F}'_n).
\tag{8.14.6.2}
\end{equation}
그러면 앞의 주의에 따라, $\mathcal{F}'$이 준연접
$\mathcal{S}_{X'}$-가군이면 $\Gamma'_*(\mathcal{F}')$은 준연접 등급
$\mathcal{S}$-가군이다
(\hyperref[I.9.6.1-ko]{I, 9.6.1}).

한편 표준 포함 사상 $j:X'\to X$는 \emph{준콤팩트}임을 주목하자. 실제로
$q'=q\circ j$가 준콤팩트이고 $q$가 분리이기 때문이다
(\hyperref[I.6.6.4-ko]{I, 6.6.4, (v)}). 따라서
$\mathcal{F}=j_*(\mathcal{F}')$은 준연접 등급 $\mathcal{S}_X$-가군이다.
이는 모든 준연접 등급 $\mathcal{S}_{X'}$-가군 $\mathcal{F}'$에 대하여
성립하며, 앞의 정의들로부터 다음을 얻는다.
\begin{equation}
\label{II.8.14.6.3-ko}
  \Gamma'_*(\mathcal{F}')=\Gamma_*(\mathcal{F}).
\tag{8.14.6.3}
\end{equation}
$X'$에 대해 같은 가정을 두고, 모든 준연접 등급 $\mathcal{S}$-가군
$\mathcal{M}$에 대하여 다음과 같이 놓겠다.
\begin{equation}
\label{II.8.14.6.4-ko}
  \operatorname{Proj}'(\mathcal{M})
  =\operatorname{Proj}(\mathcal{M})|X'
\tag{8.14.6.4}
\end{equation}
이는 준연접 등급 $\mathcal{S}_{X'}$-가군이다. 표준 준동형
\[
  \operatorname{Proj}(\mathcal{M})\to
  j_*(\operatorname{Proj}'(\mathcal{M}))
\]
(\hyperref[0.4.4.3-ko]{0, 4.4.3})은 표준 준동형
$\Gamma_*(\operatorname{Proj}(\mathcal{M}))\to
\Gamma'_*(\operatorname{Proj}'(\mathcal{M}))$을 주는데, 이는 등급
$\mathcal{S}$-가군의 준동형이다. 따라서 이를
(\hyperref[II.8.14.5.4-ko]{8.14.5.4})와 합성하면 함자적인 표준 준동형
(차수 $0$), 즉 준연접 등급 $\mathcal{S}$-가군의 준동형
\begin{equation}
\label{II.8.14.6.5-ko}
  \alpha':\mathcal{M}\to
  \Gamma'_*(\operatorname{Proj}'(\mathcal{M})).
\tag{8.14.6.5}
\end{equation}
을 얻는다.
\end{env}

\begin{env}[8.14.7]
\label{II.8.14.7-ko}
$X'$에 대해 (\hyperref[II.8.14.6-ko]{8.14.6})의 가정을 그대로 두고,
$\mathcal{F}'$을 준연접 등급 $\mathcal{S}_{X'}$-가군이라 하자. 그러면
$\operatorname{Proj}'(\Gamma'_*(\mathcal{F}'))$도 준연접 등급
$\mathcal{S}_{X'}$-가군이다. 다음과 같은 함자적인 표준 준동형
(차수 $0$), 즉 등급 $\mathcal{S}_{X'}$-가군의 준동형을 정의하겠다.
\begin{equation}
\label{II.8.14.7.1-ko}
  \beta':\operatorname{Proj}'(\Gamma'_*(\mathcal{F}'))\to\mathcal{F}'.
\tag{8.14.7.1}
\end{equation}

\oldpage[II]{200}
먼저 $Y=\operatorname{Spec}(A)$가 아핀이고,
$\mathcal{S}=\widetilde{S}$이며, 여기서 $S$는 양의 차수로 등급화된
$A$-대수라고 가정하자~; 그러면 $\Gamma'_*(\mathcal{F}')=\widetilde{M}$이며,
여기서 $M=\bigoplus_{n\in\mathbf{Z}}\Gamma(X',\mathcal{F}'_n)$은 등급
$S$-가군이다. $D_+(f)\subset X'$인 $f\in S_d$를 택하자~; 정의
(\hyperref[II.2.6.2-ko]{2.6.2})에 따라, $D_+(f)$에 제한한
$\alpha_d(f)$는 $(S(d))_{(f)}$의 원소 $f/1$에 대응하는 $D_+(f)$ 위의
$\mathcal{O}_X(d)$의 단면이며, 따라서 가역이다~; 그러므로 모든 $n>0$에
대하여 $\alpha_d(f^n)$도 마찬가지이다. 따라서 각 원소
$z/f^n\in M_f$ (여기서 $z\in M$)에 $D_+(f)$ 위의 $\mathcal{F}'$의 단면
$(z|D_+(f))(\alpha_d(f^n)|D_+(f))^{-1}$을 대응시킴으로써, 등급 가군의
(차수 $0$인) $S_f$-준동형
$\beta_f:M_f\to\Gamma(D_+(f),\mathcal{F}')$을 정의할 수 있음을 즉시
알 수 있다. 또한 (\hyperref[II.2.6.4.1-ko]{2.6.4.1})에 대응하는 가환
도식이 있으며, 이로부터 이 경우의 $\beta'$가 정의된다. 일반적인 경우로
넘어가려면 $A$-대수 $A'$, 등급 $A'$-대수 $S'=S\otimes_A A'$을
고려하고 (\hyperref[II.2.8.13.2-ko]{2.8.13.2})와 유사한 가환 도식을
사용해야 한다~; 자세한 내용은 독자에게 맡긴다.
\end{env}

\begin{proposition}[8.14.8]
\label{II.8.14.8-ko}
$X'$이 $X=\operatorname{Proj}(\mathcal{S})$의 열린집합이고
$q':X'\to Y$가 준콤팩트이면, (\hyperref[II.8.14.7-ko]{8.14.7})에서
정의한 준동형 $\beta'$는 전단사이다.
\end{proposition}

\begin{proof}
명백히 $Y$가 아핀인 경우로 한정할 수 있으며, 모든 것은
(\hyperref[II.8.14.7-ko]{8.14.7})의 표기 아래에서 준동형
$\beta_f:M_f\to\Gamma(D_+(f),\mathcal{F}')$이 동형사상임을 보이는 것으로
귀착된다. 그런데 $f$를 그 거듭제곱 가운데 하나로 바꾸어도 $D_+(f)$와
$\beta_f$는 바뀌지 않는다~; 가정에 따라 $X'$은 \emph{준콤팩트}이므로,
(\hyperref[II.8.14.4-ko]{8.14.4})에 의해 층 $\mathcal{O}_X(d)$가
\emph{가역}이라고 언제나 가정할 수 있다. $X'$은 스킴이므로
($q'$가 분리 사상이기 때문이다), 그때 이 명제는 다름 아닌
(\hyperref[I.9.3.1-ko]{I, 9.3.1})이다.
\end{proof}

\begin{corollary}[8.14.9]
\label{II.8.14.9-ko}
(\hyperref[II.8.14.8-ko]{8.14.8})의 가정 아래에서, 모든 준연접 등급
$\mathcal{S}_{X'}$-가군 $\mathcal{F}'$은
$\operatorname{Proj}'(\mathcal{M})$ 꼴의 등급 $\mathcal{S}_{X'}$-가군과
동형이며, 여기서 $\mathcal{M}$은 준연접 등급 $\mathcal{S}$-가군이다.
더 나아가 $\mathcal{F}'$이 유한형이고, $Y$가 준콤팩트 스킴이거나 그
바탕공간이 뇌터인 준스킴이라고 가정하면, $\mathcal{M}$을 유한형으로
택할 수 있다.
\end{corollary}

\begin{proof}
(\hyperref[II.8.14.8-ko]{8.14.8})에서 출발하는 증명은
(\hyperref[II.3.4.5-ko]{3.4.5})의 증명이
(\hyperref[II.3.4.4-ko]{3.4.4})에서 출발하여 따르는 절차를 정확히 따르므로,
자세한 내용은 독자에게 맡긴다.
\end{proof}

\begin{proposition}[8.14.10]
\label{II.8.14.10-ko}
(\hyperref[II.8.14.7-ko]{8.14.7})의 가정 아래에서,
$\mathcal{M}$을 준연접 등급 $\mathcal{S}$-가군,
$\mathcal{F}'$을 준연접 등급 $\mathcal{S}_{X'}$-가군이라 하자~; 합성
준동형
\begin{equation}
\label{II.8.14.10.1-ko}
  \operatorname{Proj}'(\mathcal{M})
  \xrightarrow{\operatorname{Proj}'(\alpha')}
  \operatorname{Proj}'(\Gamma'_*(\operatorname{Proj}'(\mathcal{M})))
  \xrightarrow{\beta'}\operatorname{Proj}'(\mathcal{M})
\tag{8.14.10.1}
\end{equation}
\begin{equation}
\label{II.8.14.10.2-ko}
  \Gamma'_*(\mathcal{F}')
  \xrightarrow{\alpha'}
  \Gamma'_*(\operatorname{Proj}'(\Gamma'_*(\mathcal{F}')))
  \xrightarrow{\Gamma'_*(\beta')}\Gamma'_*(\mathcal{F}')
\tag{8.14.10.2}
\end{equation}
은 항등 동형사상이다.
\end{proposition}

\begin{proof}
문제는 $Y$ 위에서 국소적이며, 검증은
(\hyperref[II.2.6.5-ko]{2.6.5})에서와 같이 이루어진다~; 자세한 내용은
독자에게 맡긴다.
\end{proof}

\begin{remark}[8.14.11]
\label{II.8.14.11-ko}
제III장 (\agref{III.2.3.1-ko}{III, 2.3.1})에서 다음을 보게 될 것이다.
$Y$가 \emph{국소 뇌터}이고 $\mathcal{S}$가 준연접 등급
$\mathcal{O}_Y$-대수이며 \emph{유한형}일 때(이
경우에는

\oldpage[II]{201}
$X'=X$로 취할 수 있다), 준동형 $\alpha$
(\hyperref[II.8.14.5.4-ko]{8.14.5.4})는 \textbf{(TN)}-\emph{전단사}이고,
이는 모든 준연접 등급 $\mathcal{S}$-가군 $\mathcal{M}$ 가운데
조건 \textbf{(TF)}를 만족하는 것에 대해 성립한다.
\end{remark}

\begin{remark}[8.14.12]
\label{II.8.14.12-ko}
(\hyperref[II.8.14.4-ko]{8.14.4})에서 서술한 상황은 다음 상황의
특수한 경우이다. $X$를 환 달린 공간, $\mathcal{S}$를 (양의 차수와
음의 차수를 갖는) 등급 $\mathcal{O}_X$-대수라 하자~; 어떤 정수
$d>0$가 존재하여 $\mathcal{S}_d$와 $\mathcal{S}_{-d}$가
\emph{가역}이고, 표준 준동형
\begin{equation}
\label{II.8.14.12.1-ko}
  \mathcal{S}_d\otimes_{\mathcal{O}_X}\mathcal{S}_{-d}
  \to\mathcal{O}_X
\tag{8.14.12.1}
\end{equation}
이 \emph{동형사상}이라고 가정하자(따라서 $\mathcal{S}_{-d}$는
$\mathcal{S}_d^{-1}$과 동일시된다). 그러면 등급
$\mathcal{O}_X$-대수 $\mathcal{S}$를 \emph{주기가 $d$인 주기적 대수}라
한다. 이 명칭은 다음 성질에서 유래한다~: \emph{앞의 가정 아래에서,
모든 등급 $\mathcal{S}$-가군 $\mathcal{F}$에 대하여 표준 준동형}
\begin{equation}
\label{II.8.14.12.2-ko}
  \mathcal{S}_d\otimes\mathcal{F}_n\to\mathcal{F}_{n+d}
\tag{8.14.12.2}
\end{equation}
\emph{은 모든 $n\in\mathbf{Z}$에 대하여 동형사상이다.} 실제로 이 문제는
$X$ 위에서 국소적이므로, $\mathcal{S}_d$가 \emph{가역인} 단면 $s$를
$X$ 위에서 갖는다고 가정할 수 있고, 그 역원 $s'$은
$\mathcal{S}_{-d}$의 단면이다. 준동형
$\mathcal{F}_{n+d}\to\mathcal{S}_d\otimes\mathcal{F}_n$은 각 단면
$z\in\Gamma(U,\mathcal{F}_{n+d})$에
$(s|U)\otimes(s'|U)z$를 대응시키는데, 이것은
$\mathcal{S}_d\otimes\mathcal{F}_n$의 $U$ 위 단면이다. 이 준동형은
(\hyperref[II.8.14.12.2-ko]{8.14.12.2})의 역준동형이므로 주장이
따른다. 이로부터 모든 $k\in\mathbf{Z}$에 대하여 다음 표준 동형사상을
얻는다.
\[
  (\mathcal{S}_d)^{\otimes k}\otimes\mathcal{F}_n
  \xrightarrow{\sim}\mathcal{F}_{n+kd}.
\]
따라서 \emph{등급 $\mathcal{S}$-가군 $\mathcal{F}$는
$\mathcal{S}_0$-가군 $\mathcal{F}_i$들
$(0\leqslant i\leqslant d-1)$과 다음 표준 준동형들을 알면 결정된다}
\[
  \mathcal{S}_i\otimes\mathcal{F}_j\to\mathcal{F}_{i+j}
  \qquad\text{모든 }0\leqslant i,j\leqslant d-1
\]
(여기서
$\mathcal{F}_{i+j}=\mathcal{S}_d\otimes_{\mathcal{S}_0}
\mathcal{F}_{i+j-d}$로 두며, 이는 $i+j\geqslant d$일 때이다). 물론 이
준동형들이 $\mathcal{S}$-가군 구조를
$(\mathcal{S}_d)^{\otimes k}\otimes\mathcal{F}_i$들의 직합
$(k\in\mathbf{Z},\ 0\leqslant i\leqslant d-1)$ 위에 올바르게 정의하려면,
여기서는 명시하지 않을 결합법칙 조건들을 만족해야 한다.

$d=1$인 경우(이는
(\hyperref[subsection:II.3.3-ko]{3.3})에서 고찰한 경우이다)에는 등급
$\mathcal{S}$-가군(각각 $X$가 준스킴이고 $\mathcal{S}$가 준연접일 때의
준연접 등급 가군)의 범주가 임의의 $\mathcal{S}_0$-가군(각각 준연접
가군)의 범주와 \emph{동치}라고 할 수 있다~; 이러한 의미에서 이 번호의
전개를 \S~\hyperref[section:II.3-ko]{3}의 전개를 일반화한 것으로 볼 수
있다. 더 나아가 적절한 유한성 조건 아래에서, 이 번호의 결과들은
((\hyperref[II.8.14.11-ko]{8.14.11})과 함께) 준스킴 위의 준연접 등급
대수와 그러한 대수 위의 등급 가군 «~mod. \textbf{(TN)}~»의 연구를,
고려하는 대수들이 \emph{주기적}인 특수한 경우의 연구로 어느 정도
환원함을 알 수 있다(그리고 이 경우에는 $\mathcal{M}$에 대한 조건
\textbf{(TN)} (\hyperref[II.3.4.2-ko]{3.4.2})가 따라서
$\mathcal{M}=0$을 함의한다).
\end{remark}

\begin{remark}[8.14.13]
\label{II.8.14.13-ko}
(\hyperref[II.8.14.1-ko]{8.14.1})의 가정 아래에서 $d$를 정수
$>0$이라 하자~; 표준 $Y$-동형사상 $h$로 $X$를
$X^{(d)}=\operatorname{Proj}(\mathcal{S}^{(d)})$에 대응시켰다
(\hyperref[II.3.1.8-ko]{3.1.8}). 모든 준연접 등급

\oldpage[II]{202}
$\mathcal{S}$-가군 $\mathcal{M}$과 모든 정수 $k$에 대하여,
$0\leqslant k\leqslant d-1$이면
(\hyperref[II.3.1.1-ko]{3.1.1})의 표기 아래 다음 표준
$h$-동형사상도 존재한다.
\begin{equation}
\label{II.8.14.13.1-ko}
  (\operatorname{Proj}(\mathcal{M}))^{(d,k)}
  \xleftarrow{\sim}\operatorname{Proj}(\mathcal{M}^{(d,k)}).
\tag{8.14.13.1}
\end{equation}

먼저 $Y=\operatorname{Spec}(A)$가 아핀이고,
$\mathcal{S}=\widetilde{S}$, $\mathcal{M}=\widetilde{M}$이라 가정하자.
여기서 $S$는 양의 차수로 등급이 주어진 $A$-대수이고 $M$은 등급
$S$-가군이다. 모든 $f\in S_e$ ($e>0$)에 대하여 $h$는 $D_+(f)$를
$D_+(f^d)$ 위로 보내며, 표준 동형사상
$S_{(f^d)}\xrightarrow{\sim}S_{(f)}$에 대응함을 안다
(\hyperref[II.2.2.2-ko]{2.2.2}). 그러면
(\hyperref[II.8.14.13.1-ko]{8.14.13.1})을 $D_+(f^d)$에 제한한 것은
표준 디-동형사상 $M_{f^d}\xrightarrow{\sim}M_f$를 $M_{f^d}$의 원소
중 차수가 $k$와 합동인 것들(mod. $d$)로 제한한 것에 대응한다. 이
동형사상들이 $f$를 그 동차 배수 $fg$로 바꾸는 것과 양립함을
확인하고, 이어서 $S$를 등급 $A'$-대수 $S'=S\otimes_A A'$로 바꾸는
것과도 마찬가지로 양립함을 확인하는 일은 독자에게 맡긴다. 여기서
$A'$은 $A$-대수이다.

특히 이에 따라 다음과 같은 $h$-동형사상이 있다.
\begin{equation}
\label{II.8.14.13.2-ko}
  (\mathcal{S}^{(d)})_{X^{(d)}}
  \xrightarrow{\sim}(\mathcal{S}_X)^{(d)}
\tag{8.14.13.2}
\end{equation}
이 동형사상은 양변의 곱셈 구조를 보존하며, 이를 통해
(\hyperref[II.8.14.13.1-ko]{8.14.13.1})은 등급
$(\mathcal{S}^{(d)})_{X^{(d)}}$-가군에서 등급
$(\mathcal{S}_X)^{(d)}$-가군으로 가는 $h$-디-동형사상이 된다.
마찬가지로 다음과 같은 $h$-동형사상이 있다.
\begin{equation}
\label{II.8.14.13.3-ko}
  \operatorname{Proj}_0(\mathcal{S}^{(d,k)}(n))
  \xrightarrow{\sim}\mathcal{O}_X(nd+k)
\tag{8.14.13.3}
\end{equation}
이는 (\hyperref[II.3.2.9-ko]{3.2.9, (ii)})의 결과를 보완한다.

동형사상 (\hyperref[II.8.14.13.1-ko]{8.14.13.1})에서 다음과 같은 등급
$\mathcal{S}^{(d)}$-가군의 동형사상을 즉시 얻는다.
\begin{equation}
\label{II.8.14.13.4-ko}
  \Gamma_*^{(d)}(\operatorname{Proj}(\mathcal{M}^{(d,k)}))
  \xrightarrow{\sim}
  \Gamma_*((\operatorname{Proj}(\mathcal{M}))^{(d,k)})
\tag{8.14.13.4}
\end{equation}
여기서 $\Gamma_*^{(d)}$는 구조 사상 $q^{(d)}:X^{(d)}\to Y$에
대응한다~; 표준 준동형 $\alpha$
(\hyperref[II.8.14.5.4-ko]{8.14.5.4})와 유사한 준동형
$\alpha^{(d)}$가 $X^{(d)}$에 대해 다음 도식을 가환하게 함은 곧바로
확인된다.
\begin{equation}
\label{II.8.14.13.5-ko}
  \xymatrix{
    & \mathcal{M}^{(d,k)}\ar[dl]_{\alpha^{(d)}}\ar[dr]^{\alpha} & \\
    \Gamma_*^{(d)}(\operatorname{Proj}(\mathcal{M}^{(d,k)}))
      \ar[rr]^{\sim}
      && \Gamma_*((\operatorname{Proj}(\mathcal{M}))^{(d,k)})
  }
\tag{8.14.13.5}
\end{equation}
이를 확인하려면 $Y$가 아핀이라고 가정하고, $\alpha^{(d)}$와
$\alpha$에 의한 $M^{(d,k)}$의 같은 원소의 상들을 각각 열린집합
$D_+(f^d)$와 $D_+(f)$에 제한한 것을 계산하면 된다(표기는 위와 같다).
\end{remark}

\begin{proposition}[8.14.14]
\label{II.8.14.14-ko}
$Y$를 준콤팩트 준스킴, $\mathcal{S}$를 유한형 준연접 등급
$\mathcal{O}_Y$-대수, $\mathcal{M}$을 조건
\textbf{(TF)}를 만족하는 준연접 등급 $\mathcal{S}$-가군이라 하자~;
$X=\operatorname{Proj}(\mathcal{S})$라 하자. 그러면
$\mathcal{S}_X$는 주기적 등급 $\mathcal{O}_X$-대수이고
(\hyperref[II.8.14.12-ko]{8.14.12}), 다음을 만족하는

\oldpage[II]{203}
$\mathcal{S}_X$의 주기 $d$가 존재하여 $(0\leqslant k\leqslant d-1)$인
$(\operatorname{Proj}(\mathcal{M}))^{(d,k)}$들은
유한형 $(\mathcal{S}_X)^{(d)}$-가군이다.
\end{proposition}

\begin{proof}
실제로 (\hyperref[II.3.1.10-ko]{3.1.10})에 따르면 $d$가 존재하여,
$\mathcal{S}^{(d)}$는 $\mathcal{S}_d=(\mathcal{S}^{(d)})_1$에 의해
생성되고 후자는 유한형 $\mathcal{S}_0$-가군이다. 따라서 첫째 주장을
증명할 때에는
(\hyperref[II.8.14.13.2-ko]{8.14.13.2})에 의해 $d=1$인 경우로
한정할 수 있고, 이 경우 그 주장은
(\hyperref[II.3.2.7-ko]{3.2.7})에서 따른다. 더구나
(\hyperref[II.8.14.13.1-ko]{8.14.13.1})을 고려하면, 둘째 주장은
(\hyperref[II.2.1.6-ko]{2.1.6, (iii)})과
(\hyperref[II.3.4.3-ko]{3.4.3})의 귀결이다.
\end{proof}
% ===== END EXACT INLINE INPUT: c2s1.tex =====


% ===== BEGIN EXACT INLINE INPUT: c2bibliography.tex =====
\clearpage
\phantomsection
\addcontentsline{toc}{section}{참고문헌 (계속)}
\renewcommand{\bibname}{참고문헌 \textit{(계속)}}

\begin{thebibliography}{26}

\item[]\noindent\textit{편집자 주. 독립형 누적 독자를 위해, 이 권에서
인용되는 [1], [9], [11], [13]을 앞선 시리즈의 참고문헌에서 재수록하였다.
EGA~II 권위본 자체에는 인쇄본 205쪽의 계속 목록 [23]--[26]만 수록되어
있다.}
\par\medskip

\bibitem[1]{II-reprint-I-1}
H.~\textsc{Cartan} et C.~\textsc{Chevalley},
\emph{Séminaire de l'École Normale Supérieure},
8\textsuperscript{e} année (1955-56), \emph{Géométrie algébrique}.

\bibitem[9]{II-reprint-I-9}
M.~\textsc{Nagata},
\emph{A general theory of algebraic geometry over Dedekind domains},
\emph{Amer. Math. Journ.}~: I, t.~LXXVIII, p.~78-116 (1956)~; II, t.~LXXX,
p.~382-420 (1958).

\bibitem[11]{II-reprint-I-11}
P.~\textsc{Samuel}, \emph{Commutative algebra} (Notes by D.~Herzig), Cornell
Univ., 1953.

\bibitem[13]{II-reprint-I-13}
P.~\textsc{Samuel} and O.~\textsc{Zariski},
\emph{Commutative algebra}, 2 vol., New York (Van Nostrand), 1958-60.

\oldpage[II]{205}
\bibitem[23]{II-23}
C.~\textsc{Chevalley},
\emph{Introduction to the theory of algebraic functions of one variable}.
Math. Surveys, n\textsuperscript{o}~6, Amer. Math. Soc., New York, 1951.

\bibitem[24]{II-24}
P.~\textsc{Jaffard},
\emph{Les systèmes d'idéaux}, Paris (Dunod), 1960.

\bibitem[25]{II-25}
M.~\textsc{Nagata},
\emph{On the derived normal rings of Noetherian integral domains},
Mem. Coll. Sci. Kyoto, sér.~A, t.~XXIX (1955), p.~293--303.

\bibitem[26]{II-26}
M.~\textsc{Nagata},
\emph{Existence theorems for non projective complete algebraic varieties},
Ill. J. Math., t.~II (1958), p.~490--498.

\end{thebibliography}
% ===== END EXACT INLINE INPUT: c2bibliography.tex =====


% ===== BEGIN EXACT INLINE INPUT: c2indexes-and-contents.tex =====
\clearpage
\oldpage[II]{207}
\markboth{제II장}{제II장}
\phantomsection
\addcontentsline{toc}{section}{기호 색인}
\section*{기호 색인}
\thispagestyle{plain}

\begingroup
\small
\setlength{\parindent}{0pt}
\setlength{\parskip}{2pt}
\newcommand{\egaShProj}{\mathcal{P}\!\operatorname{roj}}

$\mathcal{A}(X),\mathcal{A}(\mathcal{F}),\mathcal{A}(\mathcal{B})$
($X$는 $S$-준스킴, $\mathcal{F}$는 $\mathcal{O}_X$-가군,
$\mathcal{B}$는 $\mathcal{O}_X$-대수): \hyperref[II.1.1.1-ko]{1.1.1}.

$\mathcal{A}(h)$ ($h$는 $S$-사상): \hyperref[II.1.1.2-ko]{1.1.2}.

$\Spec(\mathcal{B})$ ($\mathcal{B}$는 $\mathcal{O}_S$-대수): \hyperref[II.1.3.1-ko]{1.3.1}.

$\widetilde{\mathcal{M}}$ ($\mathcal{M}$은 $\mathcal{B}$-가군,
$\mathcal{B}$는 $\mathcal{O}_S$-대수): \hyperref[II.1.4.3-ko]{1.4.3}.

$\mathbf{T}(E),\mathbf{S}(E),\mathbf{S}_A(E)$ ($E$는 $A$-가군): \hyperref[II.1.7.1-ko]{1.7.1}.

$\mathbf{S}(\mathcal{E}),\mathbf{S}_{\mathcal{A}}(\mathcal{E})$
($\mathcal{E}$는 $\mathcal{A}$-가군): \hyperref[II.1.7.4-ko]{1.7.4}.

$\mathbf{V}(\mathcal{E})$ ($\mathcal{E}$는 준연접 $\mathcal{O}_S$-가군): \hyperref[II.1.7.8-ko]{1.7.8}.

$S_n,S_+,M_n,S^{(d)},M^{(d,k)},M^{(d)}$
($S$는 등급환, $M$은 등급 $S$-가군): \hyperref[II.2.1.1-ko]{2.1.1}.

$M(n),S(n)$ ($S$는 등급환, $M$은 등급 $S$-가군): \hyperref[II.2.1.1-ko]{2.1.1}.

$\Hom_S(M,N)$ ($S$는 등급환, $M,N$은 등급 $S$-가군): \hyperref[II.2.1.2-ko]{2.1.2}.

$\mathfrak N_+,\mathfrak r_+(\mathfrak J)$
($\mathfrak J$는 $S_+$의 아이디얼): \hyperref[II.2.1.10-ko]{2.1.10}.

$S_{(f)},M_{(f)}$ ($f$는 $S_+$의 동차 원소): \hyperref[II.2.2.1-ko]{2.2.1}.

$S_{(T)},M_{(T)},S_{(\mathfrak p)},M_{(\mathfrak p)}$
($T$는 동차 원소들로 이루어진 $S_+$의 곱셈적 부분집합,
$\mathfrak p$는 $S_+$의 동차 소아이디얼): \hyperref[II.2.2.7-ko]{2.2.7}.

$\Proj(S)$ ($S$는 등급환): \hyperref[II.2.3.1-ko]{2.3.1}.

$V_+(E)$ ($E$는 $S_+$의 부분집합): \hyperref[II.2.3.2-ko]{2.3.2}.

$D_+(f)$ ($f$는 $S$의 원소): \hyperref[II.2.3.3-ko]{2.3.3}.

$\mathfrak j_+(Y)$ ($Y$는 $\Proj(S)$의 부분집합): \hyperref[II.2.3.10-ko]{2.3.10}.

$\widetilde M$ ($M$은 등급 $S$-가군): \hyperref[II.2.5.2-ko]{2.5.2}.

$\widetilde u$ ($u$는 차수 $0$인 준동형): \hyperref[II.2.5.4-ko]{2.5.4}.

$\mathcal{O}_X(n),\mathcal{F}(n)$
($X=\Proj(S)$, $\mathcal{F}$는 $\mathcal{O}_X$-가군): \hyperref[II.2.5.10-ko]{2.5.10}.

$\lambda$ (표준 준동형): \hyperref[II.2.5.11-ko]{2.5.11}.

$\mu$ (표준 준동형): \hyperref[II.2.5.12-ko]{2.5.12}.

$X^{(d)}$ ($X=\Proj(S)$, $d$는 $>0$인 정수): \hyperref[II.2.5.16-ko]{2.5.16}.

$\Gamma_*(\mathcal{F})$ ($\mathcal{F}$는 $\mathcal{O}_X$-가군,
$X=\Proj(S)$): \hyperref[II.2.6.1-ko]{2.6.1}.

$\alpha_n,\alpha,\alpha_{\mathcal{M}}$ (표준 준동형들): \hyperref[II.2.6.2-ko]{2.6.2}.

$\beta,\beta_{\mathcal{F}}$ (표준 준동형들): \hyperref[II.2.6.4-ko]{2.6.4}.

$G(\vphi),{}^a\vphi$ (표기의 남용)
($\vphi$는 등급환의 준동형): \hyperref[II.2.8.1-ko]{2.8.1}.

$\Proj(\vphi)$ ($\vphi$는 등급환의 준동형): \hyperref[II.2.8.2-ko]{2.8.2}.

$\mathcal{S}^{(d)},\mathcal{M}^{(d,k)},\mathcal{M}^{(d)},\mathcal{M}(n)$
($\mathcal{S}$는 등급 $\mathcal{O}_Y$-대수,
$\mathcal{M}$은 등급 $\mathcal{S}$-가군): \hyperref[II.3.1.1-ko]{3.1.1}.

$\Proj(\mathcal{S})$ ($\mathcal{S}$는 준연접 등급 $\mathcal{O}_Y$-대수): \hyperref[II.3.1.3-ko]{3.1.3}.

$X_f$ ($X=\Proj(\mathcal{S})$, $f$는 $Y$ 위의 $\mathcal{S}_d$의 단면): \hyperref[II.3.1.4-ko]{3.1.4}.

$\widetilde{\mathcal{M}}$ ($\mathcal{M}$은 등급 $\mathcal{S}$-가군): \hyperref[II.3.2.2-ko]{3.2.2}.

$\mathcal{O}_X(n),\mathcal{F}(n)$
($X=\Proj(\mathcal{S})$, $\mathcal{F}$는 $\mathcal{O}_X$-가군): \hyperref[II.3.2.5-ko]{3.2.5}.

$\lambda,\mu$ (표준 준동형들): \hyperref[II.3.2.6-ko]{3.2.6}.

$\Gamma_*(\mathcal{F})$ ($\mathcal{F}$는 $\mathcal{O}_X$-가군,
$X=\Proj(\mathcal{S})$): \hyperref[II.3.3.1-ko]{3.3.1}.

$\alpha_n,\alpha,\alpha_{\mathcal{M}}$ (표준 준동형들): \hyperref[II.3.3.2-ko]{3.3.2}.

\oldpage[II]{208}
\markboth{제II장}{제II장}
$\beta,\beta_{\mathcal{F}}$ (표준 준동형들): \hyperref[II.3.3.4-ko]{3.3.4}.

$G(\vphi),\Proj(\vphi)$
($\vphi$는 등급 $\mathcal{O}_Y$-대수의 준동형): \hyperref[II.3.5.1-ko]{3.5.1}.

$\Proj(u)$ ($u$는 등급 $\mathcal{O}_Y$-대수에서
등급 $\mathcal{O}_{Y'}$-대수로 가는 $g$-사상): \hyperref[II.3.5.6-ko]{3.5.6}.

$r_{\mathcal{L},\psi}$ ($\mathcal{L}$은 가역 $\mathcal{O}_X$-가군): \hyperref[II.3.7.1-ko]{3.7.1}.

$\mathbf{P}(\mathcal{E}),\mathbf{P}(E),\mathbf{P}_Y^n,\mathbf{P}_A^n$ : \hyperref[II.4.1.1-ko]{4.1.1}.

$\mathbf{P}(u)$ ($u$는 $\mathcal{O}_Y$-가군의 전사 준동형): \hyperref[II.4.1.2-ko]{4.1.2}.

$\Hyp_Y(X,\mathcal{E})$ ($\mathcal{E}$는 준연접 $\mathcal{O}_Y$-가군): \hyperref[II.4.2.5-ko]{4.2.5}.

$\varsigma$ (세그레 사상): \hyperref[II.4.3.1-ko]{4.3.1}.

$\mathcal{F}(n)$ ($\mathcal{F}$는 $\mathcal{O}_X$-가군, $X$에는
가역 $\mathcal{O}_X$-가군이 주어져 있음): \hyperref[II.4.5.1-ko]{4.5.1}.

$\varepsilon$ (항등 자기동형사상): \hyperref[II.4.5.1-ko]{4.5.1}.

$P(u,T),\sigma_i(u)$ ($u$는 유한 기저를 갖는 자유 가군의 자기준동형): \hyperref[II.6.4.1-ko]{6.4.1}.

$P(u,T),\sigma_i(u),\det u$
($u$는 정역 또는 뇌터 축소환 위의 유한형 가군의 자기준동형):
\hyperref[II.6.4.2-ko]{6.4.2} 및 \hyperref[II.6.4.7-ko]{6.4.7}.

$\sigma_i(u),\det u$ ($u$는 국소 자유 $\mathcal{A}$-가군의 자기준동형): \hyperref[II.6.4.8-ko]{6.4.8}.

$\sigma_i,\det$ (층의 준동형들): \hyperref[II.6.4.8-ko]{6.4.8}, \hyperref[II.6.4.9-ko]{6.4.9} 및 \hyperref[II.6.4.10-ko]{6.4.10}.

$\sigma_i(u),\det u$
($u$는 국소 정역 준스킴 또는 국소 뇌터 축소 준스킴 위의
유한형 $\mathcal{O}_X$-가군의 자기준동형):
\hyperref[II.6.4.9-ko]{6.4.9} 및 \hyperref[II.6.4.10-ko]{6.4.10}.

$N_{\mathcal{B}/\mathcal{A}}(f)$
($f$는 $\mathcal{A}$-대수 $\mathcal{B}$의 단면): \hyperref[II.6.5.1-ko]{6.5.1}.

$N_{\mathcal{B}/\mathcal{A}}(\mathcal{L}')$
($\mathcal{L}'$은 가역 $\mathcal{B}$-가군): \hyperref[II.6.5.2-ko]{6.5.2}.

$N_{\mathcal{B}/\mathcal{A}}(h')$
($h'$은 가역 $\mathcal{B}$-가군들의 준동형): \hyperref[II.6.5.3-ko]{6.5.3}.

$N_{X'/X}(\mathcal{L}')$
($X'$은 $X$ 위에서 유한이고, $\mathcal{L}'$은 가역 $\mathcal{O}_{X'}$-가군): \hyperref[II.6.5.5-ko]{6.5.5}.

$N_{X'/X}(h')$
($h'$은 가역 $\mathcal{O}_{X'}$-가군들의 준동형): \hyperref[II.6.5.5-ko]{6.5.5}.

$\mathcal{S}_X$ ($X=\Proj(\mathcal{S})$, $\mathcal{S}$는
$\mathcal{O}_Y$의 아이디얼들의 합): \hyperref[II.8.1.5-ko]{8.1.5}.

$S^\geq,S^\leq,S_f^\geq,S_f^\leq,
M^\geq,M^\leq,M_f^\geq,M_f^\leq$
($S$는 등급환, $M$은 등급 $S$-가군,
$f$는 $S_+$의 동차 원소): \hyperref[II.8.2.1-ko]{8.2.1}.

$\widehat S$ ($S$는 등급환): \hyperref[II.8.2.2-ko]{8.2.2}.

$\widehat M$ ($M$은 등급 $S$-가군): \hyperref[II.8.2.4-ko]{8.2.4}.

$S_{[n]},S^\natural,f^\natural$
($S$는 등급환, $f$는 $S$의 동차 원소): \hyperref[II.8.2.6-ko]{8.2.6}.

$M_{[n]},M^\natural$ ($M$은 등급 $S$-가군): \hyperref[II.8.2.8-ko]{8.2.8}.

$\widehat{\mathcal{S}}$ ($\mathcal{S}$는 등급 $\mathcal{O}_Y$-대수):
\hyperref[II.8.3.1-ko]{8.3.1}\footnote{편집자 주: 프랑스어 원문의 이 항목에는 참조 번호가 3.8.1로 인쇄되어 있으나, 연결되는 정의는 8.3.1이다.}.

$G_m$ (군 스킴): \hyperref[II.8.3.9-ko]{8.3.9}.

$\mathcal{S}_X$ ($X=\Proj(\mathcal{S})$, $\mathcal{S}$는 등급 $\mathcal{O}_Y$-대수): \hyperref[II.8.6.1-ko]{8.6.1}.

$\mathcal{S}_X,C_X,\widehat C_X,E_X,\widehat E_X$
($X=\Proj(\mathcal{S})$, $\mathcal{S}$는 $\mathcal{S}_0=\mathcal{O}_Y$인
등급 $\mathcal{O}_Y$-대수): \hyperref[II.8.6.1-ko]{8.6.1}.

$\mathcal{S}_{[n]},\mathcal{S}^{\natural}$
($\mathcal{S}$는 등급 $\mathcal{O}_Y$-대수): \hyperref[II.8.7.2-ko]{8.7.2}.

$\egaShProj_0(\mathcal{M}),\egaShProj(\mathcal{M}),\mathcal{M}_X,
\widehat{\mathcal{M}},\mathcal{M}^{\square}$
($\mathcal{M}$은 등급 $\mathcal{S}$-가군): \hyperref[II.8.12.1-ko]{8.12.1}.

$\mathcal{M}_X^\geq$
($X=\Proj(\mathcal{S})$, $\mathcal{M}$은 등급 $\mathcal{S}$-가군): \hyperref[II.8.12.5-ko]{8.12.5}.

$\mathcal{M}_{[n]},\mathcal{M}^{\natural}$
($\mathcal{M}$은 등급 $\mathcal{S}$-가군): \hyperref[II.8.12.9-ko]{8.12.9}.

$(\widetilde{\mathcal{N}})^-$
($\mathcal{N}$은 어떤 등급 $\mathcal{S}$-가군의 부분-$\mathcal{S}$-가군): \hyperref[II.8.13.1-ko]{8.13.1}.

$\overline{\mathcal{N}}$
($\mathcal{N}$은 어떤 등급 $\mathcal{S}$-가군의 부분-$\mathcal{S}$-가군): \hyperref[II.8.13.2-ko]{8.13.2}.

$\Gamma_*(\mathcal{F})$
($X=\Proj(\mathcal{S})$, $\mathcal{F}$는 등급 $\mathcal{S}_X$-가군): \hyperref[II.8.14.5-ko]{8.14.5}.

$\mathcal{S}_{X'},\Gamma'_*(\mathcal{F}'),\egaShProj'(\mathcal{M}),
\alpha',\beta'$
($X'$은 $X=\Proj(\mathcal{S})$의 열린집합, $\mathcal{F}'$은 등급
$\mathcal{S}_{X'}$-가군, $\mathcal{M}$은 등급 $\mathcal{S}$-가군): \hyperref[II.8.14.6-ko]{8.14.6}.

\endgroup

\clearpage
\oldpage[II]{209}
\markboth{제II장}{제II장}
\thispagestyle{plain}
\phantomsection
\addcontentsline{toc}{section}{용어 색인}
\section*{용어 색인}

\begingroup
\small
\setlength{\parindent}{0pt}
\setlength{\parskip}{2pt}
\newcommand{\egaTerme}[2]{\textbf{#1} : #2.\par}

\egaTerme{아핀 (사상)}{\hyperref[II.1.6.1-ko]{1.6.1}}
\egaTerme{준스킴 위에서 아핀인 준스킴}{\hyperref[II.1.2.1-ko]{1.2.1}}
\egaTerme{등급환 위의 등급대수}{\hyperref[II.2.1.2-ko]{2.1.2}}
\egaTerme{$A$-가군의 대칭 대수}{\hyperref[II.1.7.1-ko]{1.7.1}}
\egaTerme{$\mathcal{A}$-가군의 대칭 $\mathcal{A}$-대수}{\hyperref[II.1.7.4-ko]{1.7.4}}
\egaTerme{본질적으로 축소인 등급 $\mathcal{O}_Y$-대수}{\hyperref[II.3.1.12-ko]{3.1.12}}
\egaTerme{정역인 등급 $\mathcal{O}_Y$-대수}{\hyperref[II.3.1.12-ko]{3.1.12}}
\egaTerme{풍부한 가역 $\mathcal{O}_X$-가군}{\hyperref[II.4.5.3-ko]{4.5.3}}
\egaTerme{$f$에 관해 풍부한, $f$-풍부한, $Y$-풍부한 가역 $\mathcal{O}_X$-가군}{\hyperref[II.4.6.1-ko]{4.6.1}}
\egaTerme{본질적으로 정역인 등급환}{\hyperref[II.2.1.11-ko]{2.1.11}}
\egaTerme{본질적으로 축소인 등급환}{\hyperref[II.2.1.10-ko]{2.1.10}}
\egaTerme{$\mathcal{B}$-가군에 대응하는 층}{\hyperref[II.1.4.3-ko]{1.4.3}}
\egaTerme{등급 $S$-가군에 대응하는 층}{\hyperref[II.2.5.3-ko]{2.5.3}}
\egaTerme{등급 $\mathcal{S}$-가군에 대응하는 층}{\hyperref[II.3.2.2-ko]{3.2.2}}
\egaTerme{$\mathcal{O}_S$-대수에 대응하는 $S$-스킴}{\hyperref[II.1.3.1-ko]{1.3.1}}

\egaTerme{조건 (TF), 조건 (TN)}{\hyperref[II.2.7.2-ko]{2.7.2}}
\egaTerme{조건 (TF), 조건 (TN)}{\hyperref[II.3.4.2-ko]{3.4.2}}
\egaTerme{등급 $\mathcal{O}_Y$-대수로 정의되는 아핀 원뿔과 사영 원뿔}{\hyperref[II.8.3.1-ko]{8.3.1}}
\egaTerme{등급 $\mathcal{O}_Y$-대수로 정의되는 천공 아핀 원뿔과 천공 사영 원뿔}{\hyperref[II.8.3.4-ko]{8.3.4}}
\egaTerme{준스킴 $\Proj(\mathcal{S})$를 사영하는 아핀 원뿔과 사영 원뿔}{\hyperref[II.8.3.1-ko]{8.3.1}}
\egaTerme{체 위의 대수곡선}{\hyperref[II.7.4.2-ko]{7.4.2}}
\egaTerme{완비 대수곡선}{\hyperref[II.7.4.7-ko]{7.4.7}}
\egaTerme{세르 판정법}{\hyperref[II.5.2.1-ko]{5.2.1}}

\egaTerme{정역 위의 유한형 가군 자기준동형의 결정식}{\hyperref[II.6.4.2-ko]{6.4.2}}
\egaTerme{뇌터 축소환 위의 유한형 가군 자기준동형의 결정식}{\hyperref[II.6.4.7-ko]{6.4.7}}
\egaTerme{국소 자유 $\mathcal{A}$-가군 자기준동형의 결정식}{\hyperref[II.6.4.8-ko]{6.4.8}}
\egaTerme{국소 정역 준스킴 $X$ 위의 유한형 $\mathcal{O}_X$-가군 자기준동형의 결정식}{\hyperref[II.6.4.9-ko]{6.4.9}}
\egaTerme{국소 뇌터 축소 준스킴 $X$ 위의 유한형 $\mathcal{O}_X$-가군 자기준동형의 결정식}{\hyperref[II.6.4.10-ko]{6.4.10}}

\egaTerme{블로업 (준스킴)}{\hyperref[II.8.1.3-ko]{8.1.3}}
\egaTerme{정수적 (사상)}{\hyperref[II.6.1.1-ko]{6.1.1}}
\egaTerme{준스킴 위에서 정수적인 준스킴}{\hyperref[II.6.1.1-ko]{6.1.1}}
\egaTerme{정수적인 준연접 $\mathcal{O}_S$-대수}{\hyperref[II.6.1.2-ko]{6.1.2}}
\egaTerme{$\mathcal{A}$ 위에서 정수적인 $\mathcal{A}$-대수의 단면}{\hyperref[II.6.3.1-ko]{6.3.1}}

\egaTerme{사영 다발의 기본 층}{\hyperref[II.4.1.1-ko]{4.1.1}}
\egaTerme{$\mathcal{A}$-대수 안에서 환의 층 $\mathcal{A}$의 정수적 폐포}{\hyperref[II.6.3.2-ko]{6.3.2}}
\egaTerme{$\mathcal{O}_X$-대수에 대한 준스킴 $X$의 정수적 폐포}{\hyperref[II.6.3.4-ko]{6.3.4}}
\egaTerme{아핀 원뿔의 사영 폐포}{\hyperref[II.8.3.1-ko]{8.3.1}}

\oldpage[II]{210}
\markboth{A. GROTHENDIECK --- 제II장}{A. GROTHENDIECK --- 제II장}
\egaTerme{$k(y)$의 확장체 $K$에 대한 사영 다발의 $K$-유리 기하적 섬유}{\hyperref[II.4.2.6-ko]{4.2.6}}
\egaTerme{$\mathcal{O}_Y$-가군으로 정의되는 사영 다발}{\hyperref[II.4.1.1-ko]{4.1.1}}
\egaTerme{$\mathcal{O}_S$-가군으로 정의되는 벡터 다발}{\hyperref[II.1.7.8-ko]{1.7.8}}
\egaTerme{유한 (사상)}{\hyperref[II.6.1.1-ko]{6.1.1}}
\egaTerme{준스킴 위에서 유한인 준스킴}{\hyperref[II.6.1.1-ko]{6.1.1}}
\egaTerme{유한인 준연접 $\mathcal{O}_S$-대수}{\hyperref[II.6.1.2-ko]{6.1.2}}
\egaTerme{국소 자유 가군 자기준동형의 기본대칭함수}{\hyperref[II.6.4.1-ko]{6.4.1}}
\egaTerme{정역 위의 유한형 가군 자기준동형의 기본대칭함수}{\hyperref[II.6.4.2-ko]{6.4.2}}
\egaTerme{뇌터 축소환 위의 유한형 가군 자기준동형의 기본대칭함수}{\hyperref[II.6.4.7-ko]{6.4.7}}
\egaTerme{일정한 계수의 국소 자유 $\mathcal{A}$-가군 자기준동형의 기본대칭함수}{\hyperref[II.6.4.8-ko]{6.4.8}}
\egaTerme{국소 정역 준스킴 $X$ 위의 유한형 $\mathcal{O}_X$-가군 자기준동형의 기본대칭함수}{\hyperref[II.6.4.9-ko]{6.4.9}}
\egaTerme{국소 뇌터 축소 준스킴 $X$ 위의 유한형 $\mathcal{O}_X$-가군 자기준동형의 기본대칭함수}{\hyperref[II.6.4.10-ko]{6.4.10}}
\egaTerme{$\mathcal{A}$-대수의 단면의 기본대칭함수}{\hyperref[II.6.5.1-ko]{6.5.1}}

\egaTerme{등급 $\mathcal{S}$-가군의 부분-$\mathcal{S}$-가군의 동차화}{\hyperref[II.8.13.2-ko]{8.13.2}}
\egaTerme{등급환의 준동형}{\hyperref[II.2.1.2-ko]{2.1.2}}
\egaTerme{등급가군의 차수 $k$인 준동형}{\hyperref[II.2.1.2-ko]{2.1.2}}

\egaTerme{$S_+$의 아이디얼, $S_+$의 동차 소아이디얼 ($S$는 등급환)}{\hyperref[II.2.1.10-ko]{2.1.10}}

\egaTerme{저우의 보조정리}{\hyperref[II.5.6.1-ko]{5.6.1}}
\egaTerme{사영 원뿔의 무한원 자취}{\hyperref[II.8.3.3-ko]{8.3.3}}

\egaTerme{표준 사상 $G(\mathcal{E})\to\Proj(S)$, $S=\bigoplus_{n\geq0}\Gamma(X,\mathcal{L}^{\otimes n})$}{\hyperref[II.4.5.1-ko]{4.5.1}}
\egaTerme{표준 사상 $X\to\Spec(\Gamma(X,\mathcal{O}_X))$}{\hyperref[II.5.1.1-ko]{5.1.1}}
\egaTerme{세그레 사상}{\hyperref[II.4.3.1-ko]{4.3.1}}

\egaTerme{$S_+$의 멱영근기 ($S$는 등급환)}{\hyperref[II.2.1.10-ko]{2.1.10}}
\egaTerme{$\mathcal{S}_+$의 멱영근기 ($\mathcal{S}$는 준연접 등급 $\mathcal{O}_Y$-대수)}{\hyperref[II.3.1.12-ko]{3.1.12}}
\egaTerme{$\mathcal{A}$-대수의 단면의 노름}{\hyperref[II.6.5.1-ko]{6.5.1}}
\egaTerme{가역 $\mathcal{B}$-가군의 노름}{\hyperref[II.6.5.2-ko]{6.5.2}}
\egaTerme{가역 $\mathcal{B}$-가군의 단면의 노름}{\hyperref[II.6.5.3-ko]{6.5.3}}
\egaTerme{가역 $\mathcal{O}_{X'}$-가군의 노름}{\hyperref[II.6.5.5-ko]{6.5.5}}

\egaTerme{주기적 (등급대수)}{\hyperref[II.8.14.12-ko]{8.14.12}}
\egaTerme{정역 위의 유한형 가군 자기준동형의 특성다항식}{\hyperref[II.6.4.2-ko]{6.4.2}}
\egaTerme{뇌터 축소환 위의 유한형 가군 자기준동형의 특성다항식}{\hyperref[II.6.4.7-ko]{6.4.7}}
\egaTerme{아핀 (사영) 원뿔의 꼭짓점 준스킴}{\hyperref[II.8.3.3-ko]{8.3.3}}
\egaTerme{유한 표시를 갖는 등급가군}{\hyperref[II.2.1.1-ko]{2.1.1}}
\egaTerme{유한 표시를 갖는 등급 $\mathcal{S}$-가군}{\hyperref[II.3.1.1-ko]{3.1.1}}
\egaTerme{두 등급가군의 등급 텐서곱}{\hyperref[II.2.1.2-ko]{2.1.2}}
\egaTerme{사영 (사상)}{\hyperref[II.5.5.2-ko]{5.5.2}}
\egaTerme{준스킴 위에서 사영적인 준스킴}{\hyperref[II.5.5.2-ko]{5.5.2}}
\egaTerme{고유 (사상)}{\hyperref[II.5.4.1-ko]{5.4.1}}
\egaTerme{고유 (부분집합)}{\hyperref[II.5.4.10-ko]{5.4.10}}
\egaTerme{준스킴 위에서 고유한 준스킴}{\hyperref[II.5.4.1-ko]{5.4.1}}

\egaTerme{준아핀 (사상)}{\hyperref[II.5.1.1-ko]{5.1.1}}
\egaTerme{준스킴 위에서 준아핀인 준스킴}{\hyperref[II.5.1.1-ko]{5.1.1}}
\egaTerme{준아핀 (스킴)}{\hyperref[II.5.1.1-ko]{5.1.1}}

\oldpage[II]{211}
\markboth{제II장}{제II장}
\egaTerme{준유한 (사상)}{\hyperref[II.6.2.3-ko]{6.2.3}}
\egaTerme{준스킴 위에서 준유한인 준스킴}{\hyperref[II.6.2.3-ko]{6.2.3}}
\egaTerme{준사영 (사상)}{\hyperref[II.5.3.1-ko]{5.3.1}}
\egaTerme{준스킴 위에서 준사영인 준스킴}{\hyperref[II.5.3.1-ko]{5.3.1}}

\egaTerme{$S_+$의 아이디얼의 근기 ($S$는 등급환)}{\hyperref[II.2.1.10-ko]{2.1.10}}
\egaTerme{천공 사영 원뿔의 무한원 자취 위로의 표준 수축}{\hyperref[II.8.3.5-ko]{8.3.5}}

\egaTerme{$\mathcal{O}_X(-1)$의 표준 단면 ($X$는 블로업 준스킴)}{\hyperref[II.8.1.9-ko]{8.1.9}}
\egaTerme{아핀 원뿔의 영단면}{\hyperref[II.8.3.3-ko]{8.3.3}}
\egaTerme{벡터 다발의 영단면}{\hyperref[II.1.7.9-ko]{1.7.9}}
\egaTerme{아핀 (사영) 원뿔의 꼭짓점 단면}{\hyperref[II.8.3.3-ko]{8.3.3}}
\egaTerme{$\mathcal{O}_S$-대수의 스펙트럼}{\hyperref[II.1.3.1-ko]{1.3.1}}
\egaTerme{등급환의 동차 소 스펙트럼}{\hyperref[II.2.3.1-ko]{2.3.1}}
\egaTerme{준연접 등급 $\mathcal{O}_S$-대수의 동차 스펙트럼}{\hyperref[II.3.1.3-ko]{3.1.3}}

\egaTerme{(TN)-단사, (TN)-전사, (TN)-전단사 (등급가군의 준동형)}{\hyperref[II.2.7.2-ko]{2.7.2}}
\egaTerme{등급가군의 (TN)-동형사상}{\hyperref[II.2.7.2-ko]{2.7.2}}
\egaTerme{(TN)-단사, (TN)-전사, (TN)-전단사 (등급환의 준동형)}{\hyperref[II.2.9.1-ko]{2.9.1}}
\egaTerme{등급환의 (TN)-동형사상}{\hyperref[II.2.9.1-ko]{2.9.1}}
\egaTerme{(TN)-단사, (TN)-전사, (TN)-전단사 (등급 $\mathcal{S}$-가군의 준동형)}{\hyperref[II.3.4.2-ko]{3.4.2}}
\egaTerme{등급 $\mathcal{S}$-가군의 (TN)-동형사상}{\hyperref[II.3.4.2-ko]{3.4.2}}
\egaTerme{(TN)-단사, (TN)-전사, (TN)-전단사 (등급 $\mathcal{O}_Y$-대수의 준동형)}{\hyperref[II.3.6.1-ko]{3.6.1}}
\egaTerme{등급 $\mathcal{O}_Y$-대수의 (TN)-동형사상}{\hyperref[II.3.6.1-ko]{3.6.1}}
\egaTerme{$\Proj(S)$ 위의 스펙트럼 위상}{\hyperref[II.2.3.3-ko]{2.3.3}}
\egaTerme{$q$에 관해 매우 풍부한, $Y$에 관해 매우 풍부한 가역 $\mathcal{O}_X$-가군}{\hyperref[II.4.4.2-ko]{4.4.2}}
\egaTerme{유한형 (등급 $\mathcal{S}$-가군)}{\hyperref[II.3.1.1-ko]{3.1.1}}

\egaTerme{보편적으로 닫힌 (사상)}{\hyperref[II.5.4.9-ko]{5.4.9}}

\endgroup
\clearpage
\oldpage[II]{212}
\thispagestyle{empty}
\mbox{}

\clearpage
\oldpage[II]{213}
\markboth{목차}{목차}
\phantomsection
\addcontentsline{toc}{section}{원문 목차}
\section*{목차}
\thispagestyle{plain}

\begingroup
\small
\renewcommand{\arraystretch}{1.04}
\begin{longtable}{@{}p{0.09\textwidth}p{0.69\textwidth}r@{}}
\textbf{단위} & \textbf{제목} & \textbf{원문 쪽}\\
\hline
\textbf{II} & \textbf{몇몇 사상 부류에 대한 초보적 전역 연구} & \textbf{5}\\
\hyperref[section:II.1-ko]{\textsection1} & 아핀 사상 & 5\\
\hyperref[subsection:II.1.1-ko]{1.1} & $S$-준스킴과 $\mathcal{O}_S$-대수 & 5\\
\hyperref[subsection:II.1.2-ko]{1.2} & 준스킴 위에서 아핀인 준스킴 & 6\\
\hyperref[subsection:II.1.3-ko]{1.3} & $\mathcal{O}_S$-대수에 대응하는 $S$ 위에서 아핀인 준스킴 & 8\\
\hyperref[subsection:II.1.4-ko]{1.4} & $S$ 위에서 아핀인 준스킴 위의 준연접층 & 9\\
\hyperref[II.1.5-ko]{1.5} & 밑준스킴의 변경 & 12\\
\hyperref[II.1.6-ko]{1.6} & 아핀 사상 & 14\\
\hyperref[II.1.7-ko]{1.7} & 가군의 층에 대응하는 벡터 다발 & 14\\
\hyperref[section:II.2-ko]{\textsection2} & 동차 소 스펙트럼 & 19\\
\hyperref[subsection:II.2.1-ko]{2.1} & 등급환과 등급가군에 관한 일반 사항. & 19\\
\hyperref[subsection:II.2.2-ko]{2.2} & 등급환의 분수환. & 23\\
\hyperref[subsection:II.2.3-ko]{2.3} & 등급환의 동차 소 스펙트럼. & 25\\
\hyperref[subsection:II.2.4-ko]{2.4} & $\operatorname{Proj}(S)$의 스킴 구조 & 28\\
\hyperref[subsection:II.2.5-ko]{2.5} & 등급가군에 대응하는 층. & 30\\
\hyperref[subsection:II.2.6-ko]{2.6} & $\operatorname{Proj}(S)$ 위의 층에 대응하는 등급 $S$-가군. & 36\\
\hyperref[subsection:II.2.7-ko]{2.7} & 유한성 조건. & 38\\
\hyperref[subsection:II.2.8-ko]{2.8} & 함자적 성질. & 41\\
\hyperref[subsection:II.2.9-ko]{2.9} & $\operatorname{Proj}(S)$ 스킴의 닫힌 부분스킴. & 48\\
\hyperref[section:II.3-ko]{\textsection3} & 등급 대수의 층의 동차 스펙트럼 & 49\\
\hyperref[subsection:II.3.1-ko]{3.1} & 준연접 등급 $\mathcal{O}_Y$-대수의 동차 스펙트럼. & 49\\
\hyperref[subsection:II.3.2-ko]{3.2} & 등급 $\mathcal{S}$-가군에 대응하는 $\operatorname{Proj}(\mathcal{S})$ 위의 층. & 54\\
\hyperref[subsection:II.3.3-ko]{3.3} & $\operatorname{Proj}(\mathcal{S})$ 위의 층에 대응하는 등급 $\mathcal{S}$-가군. & 56\\
\hyperref[subsection:II.3.4-ko]{3.4} & 유한성 조건. & 59\\
\hyperref[subsection:II.3.5-ko]{3.5} & 함자적 성질. & 61\\
\hyperref[subsection:II.3.6-ko]{3.6} & 준스킴 $\operatorname{Proj}(\mathcal{S})$의 닫힌 부분준스킴. & 64\\
\hyperref[subsection:II.3.7-ko]{3.7} & 준스킴에서 동차 스펙트럼으로 가는 사상. & 65\\
\hyperref[subsection:II.3.8-ko]{3.8} & 동차 스펙트럼으로의 몰입 사상 판정법. & 69\\
\end{longtable}
\oldpage[II]{214}
\markboth{A. 그로텐디크 --- 제II장}{A. 그로텐디크 --- 제II장}
\begin{longtable}{@{}p{0.09\textwidth}p{0.69\textwidth}r@{}}
\hyperref[section:II.4-ko]{\textsection4} & 사영 다발. 풍부한 층 & 71\\
\hyperref[subsection:II.4.1-ko]{4.1} & 사영 다발의 정의. & 71\\
\hyperref[subsection:II.4.2-ko]{4.2} & 준스킴에서 사영 다발로 가는 사상. & 72\\
\hyperref[subsection:II.4.3-ko]{4.3} & 세그레 사상. & 76\\
\hyperref[subsection:II.4.4-ko]{4.4} & 사영 다발로의 몰입 사상. 매우 풍부한 층 & 78\\
\hyperref[subsection:II.4.5-ko]{4.5} & 풍부한 층 & 83\\
\hyperref[II.4.6-ko]{4.6} & 상대적으로 풍부한 층 & 89\\
\hyperref[section:II.5-ko]{\textsection5} & 준아핀 사상~; 준사영 사상. 고유 사상~; 사영 사상 & 94\\
\hyperref[subsection:II.5.1-ko]{5.1} & 준아핀 사상. & 94\\
\hyperref[subsection:II.5.2-ko]{5.2} & 세르 판정법. & 97\\
\hyperref[subsection:II.5.3-ko]{5.3} & 준사영 사상 & 99\\
\hyperref[subsection:II.5.4-ko]{5.4} & 고유 사상과 보편적으로 닫힌 사상. & 100\\
\hyperref[subsection:II.5.5-ko]{5.5} & 사영 사상. & 103\\
\hyperref[subsection:II.5.6-ko]{5.6} & 차우 보조정리. & 106\\
\hyperref[section:II.6-ko]{\textsection6} & 정수적 사상과 유한 사상 & 110\\
\hyperref[subsection:II.6.1-ko]{6.1} & 다른 준스킴 위에서 정수적인 준스킴. & 110\\
\hyperref[subsection:II.6.2-ko]{6.2} & 준유한 사상. & 114\\
\hyperref[subsection:II.6.3-ko]{6.3} & 준스킴의 정수적 폐포. & 116\\
\hyperref[subsection:II.6.4-ko]{6.4} & $\mathcal{O}_X$-가군의 자기준동형의 행렬식. & 120\\
\hyperref[subsection:II.6.5-ko]{6.5} & 가역층의 노름. & 125\\
\hyperref[subsection:II.6.6-ko]{6.6} & 응용: 풍부성 판정법. & 130\\
\hyperref[subsection:II.6.7-ko]{6.7} & 슈발레의 정리. & 135\\
\hyperref[section:II.7-ko]{\textsection7} & 값매김 판정법 & 138\\
\hyperref[subsection:II.7.1-ko]{7.1} & 값매김환에 대한 복습. & 138\\
\hyperref[subsection:II.7.2-ko]{7.2} & 분리성의 값매김 판정법 & 141\\
\hyperref[subsection:II.7.3-ko]{7.3} & 고유성의 값매김 판정법 & 143\\
\hyperref[subsection:II.7.4-ko]{7.4} & 대수곡선과 차원 1인 함수체 & 148\\
\hyperref[section:II.8-ko]{\textsection8} & 블로업 스킴; 사영 원뿔; 사영 폐포 & 152\\
\hyperref[subsection:II.8.1-ko]{8.1} & 블로업 준스킴 & 152\\
\hyperref[subsection:II.8.2-ko]{8.2} & 등급환에서의 국소화에 관한 예비 결과 & 157\\
\hyperref[subsection:II.8.3-ko]{8.3} & 사영 원뿔 & 162\\
\hyperref[subsection:II.8.4-ko]{8.4} & 벡터 다발의 사영 폐포 & 168\\
\hyperref[subsection:II.8.5-ko]{8.5} & 함자적 성질 & 169\\
\hyperref[subsection:II.8.6-ko]{8.6} & 꼭지점을 제거한 원뿔에 대한 표준 동형사상 & 171\\
\end{longtable}
\oldpage[II]{215}
\markboth{목차}{목차}
\begin{longtable}{@{}p{0.09\textwidth}p{0.69\textwidth}r@{}}
\hyperref[subsection:II.8.7-ko]{8.7} & 사영 원뿔의 블로업 & 173\\
\hyperref[subsection:II.8.8-ko]{8.8} & 풍부한 층과 수축 & 177\\
\hyperref[subsection:II.8.9-ko]{8.9} & 그라우어트의 풍부성 판정법: 진술 & 182\\
\hyperref[subsection:II.8.10-ko]{8.10} & 그라우에르트의 풍부성 판정법: 증명 & 184\\
\hyperref[subsection:II.8.11-ko]{8.11} & 수축의 유일성 & 189\\
\hyperref[subsection:II.8.12-ko]{8.12} & 사영하는 원뿔 위의 준연접층 & 191\\
\hyperref[subsection:II.8.13-ko]{8.13} & 부분층과 닫힌 부분스킴의 사영 폐포 & 195\\
\hyperref[subsection:II.8.14-ko]{8.14} & 등급 $\mathcal{S}$-가군에 수반되는 층에 대한 보충 & 197\\
\multicolumn{2}{@{}l}{참고문헌(계속)} & 205\\
\multicolumn{2}{@{}l}{기호 색인} & 207\\
\multicolumn{2}{@{}l}{용어 색인} & 209\\
\multicolumn{2}{@{}l}{정오표 및 추록 (목록 1)} & 217\\
\end{longtable}

\begin{flushright}
\emph{1960년 7월 22일 접수.}
\end{flushright}
\endgroup
% ===== END EXACT INLINE INPUT: c2indexes-and-contents.tex =====


% ===== BEGIN EXACT INLINE INPUT: c2errata-addenda.tex =====
\clearpage
\oldpage[II]{217}
\markboth{정오표 및 추록, 목록 1}{정오표 및 추록, 목록 1}
\phantomsection
\addcontentsline{toc}{section}{정오표 및 추록 (목록 1)}
\section*{정오표 및 추록}
\thispagestyle{plain}
\begin{center}
\large (목록 1)
\end{center}

\subsection*{제0장}

\paragraph{(0, 2.1.3).}
p.~21의 아래에서 15번째 줄에서 «~inter-section~»을
«~intersection~»으로 바꾼다.

\paragraph{(0, 2.1.6).}
p.~22의 13번째 줄에서 다음을
\[
  Z_i\cap C\left(\bigcup_{j\ne i}Z_j\right)
  \quad\hbox{를}\quad
  \bigcup_{j\ne i}Z_j\quad\hbox{로 바꾼다}.
\]

\paragraph{(0, 3.1.6).}
p.~25의 아래에서 17번째 줄에서
$\Gamma(U,\sh{F})$를 $\Gamma(V,\sh{F})$로 바꾼다.

\paragraph{(0, 4.1.1).}
p.~36의 12번째 줄에서 $\sh{G}\to\sh{F}$를
$\sh{B}\to\sh{A}$로 바꾼다.

\paragraph{(0, 4.2.1).}
p.~39의 아래에서 11번째 줄에서 «~faisceaux~»를 «~faisceau~»로 바꾼다.

\paragraph{(0, 5.1.2).}
p.~45의 14번째 줄 뒤에 다음을 추가한다:

$\sh{G}$가 $Y$ 위의 단면들로 생성되는 $\sh{O}_Y$-가군이면,
$f^*(\sh{G})$는 $X$ 위의 단면들로 생성된다. 이는 $f^*$가
오른쪽 완전한 함자이기 때문이다.

\paragraph{(0, 5.2.4).}
p.~46의 아래에서 18번째 줄에서 $f_U$를 $f_U^*$로 바꾼다.

\paragraph{(0, 5.3.4).}
p.~47의 아래에서 16번째 줄 앞에 다음을 추가한다:

\[
  \sh{A}\to\sh{B}\to\sh{C}\to\sh{D}\to\sh{E}
\]
가 $\sh{O}_X$-가군의 완전열이고
$\sh{A},\sh{B},\sh{D},\sh{E}$가 연접이면, $\sh{C}$도
연접이다.

\paragraph{(0, 5.4.1).}
p.~49의 2번째 줄 뒤에 다음을 추가한다:

$\sh{O}_X$가 연접인 환의 층이라고 가정하고, $\sh{F}$를 연접
$\sh{O}_X$-가군이라 하자. 어떤 점 $x\in X$에서 $\sh{F}_x$가
계수 $n$인 자유 $\sh{O}_x$-가군이면, $x$의 근방 $U$가 존재하여
$\sh{F}|U$는 계수 $n$인 국소 자유 $(\sh{O}_X|U)$-가군이다. 실제로
$\sh{F}_x$는 $\sh{O}_x^n$과 동형이고, 이 명제는
(0, 5.2.7)에서 따른다.

\paragraph{(0, 6.6.2).}
p.~59의 10번째 줄과 11번째 줄을 다음으로 바꾼다:

실제로 이는 (0, 6.4.1, d))에서 따른다.

\paragraph{(0, 6.7.1).}
p.~59의 아래에서 10번째 줄 앞에 다음을 추가한다:

$f$가 평탄 사상이면, $X$가 \emph{$Y$ 위에서 평탄} 또는
\emph{$Y$-평탄}이라고도 한다.

\paragraph{(0, 7.2.9).}
p.~65의 아래에서 11번째 줄에서 $M_n$을 $M_{n-1}$로 바꾼다.

\subsection*{제I장}

\paragraph{(I, 1.2.1).}
p.~83의 18번째 줄에서 $K$를 $k$로 바꾼다(두 곳).

\paragraph{(I, 1.2.7).}
p.~84의 아래에서 16번째 줄에서 $A$를 $A'$로 바꾸고(두 곳),
$\mathfrak N$을 $\mathfrak N'$으로 바꾼다. 아래에서 17번째 줄에서는
$X$를 $X'$으로 바꾼다.

\paragraph{(I, 1.7.3) 및 (I, 2.2.4).}
이 번호들에 제시된 명제들은 J.~Tate의 지적에 따라 다음과 같이
일반화할 수 있다:
\subsection*{1.8. 국소환 달린 공간에서 아핀 스킴으로 가는 사상}
\label{ega2:addendum:subsection:I.1.8-ko}

\begin{proposition}[1.8.1]
\label{ega2:addendum:I.1.8.1-ko}
$(S,\sh{O}_S)$를 아핀 스킴, $(X,\sh{O}_X)$를 국소환 달린 공간이라 하자.
그러면 환 준동형들의 집합
\oldpage[II]{218}
\markboth{A. 그로텐디크 --- 제II장}{A. 그로텐디크 --- 제II장}
\[
  \Gamma(S,\sh{O}_S)\to\Gamma(X,\sh{O}_X)
\]
에서, 모든 $x\in X$에 대하여
$\theta_x^\sharp:\sh{O}_{\psi(x)}\to\sh{O}_x$가 국소 준동형인 환 달린
공간의 사상
$(\psi,\theta):(X,\sh{O}_X)\to(S,\sh{O}_S)$들의 집합으로 가는 표준
전단사가 존재한다.
\end{proposition}

먼저 $(X,\sh{O}_X)$와 $(S,\sh{O}_S)$가 임의의 두 환 달린 공간이면,
$(X,\sh{O}_X)$에서 $(S,\sh{O}_S)$로 가는 사상 $(\psi,\theta)$가
표준적으로 환 준동형
$\Gamma(\theta):\Gamma(S,\sh{O}_S)\to\Gamma(X,\sh{O}_X)$를 정의함을
유의하자. 따라서 첫 번째 사상
\begin{equation}
\label{ega2:addendum:I.1.8.1.1-ko}
  \rho:\Hom((X,\sh{O}_X),(S,\sh{O}_S))
  \to\Hom(\Gamma(S,\sh{O}_S),\Gamma(X,\sh{O}_X)).
\tag{1.8.1.1}
\end{equation}
을 얻는다. 역으로, 명제의 가정 아래 $A=\Gamma(S,\sh{O}_S)$로 놓고
환 준동형 $\vphi:A\to\Gamma(X,\sh{O}_X)$를 생각하자. 모든 $x\in X$에
대하여 $\vphi(f)(x)=0$을 만족하는 $f\in A$들의 집합은 $A$의 소
아이디얼임이 분명하다. 실제로 $\sh{O}_x/\mathfrak m_x=\kres(x)$는 체이다.
따라서 이 소 아이디얼은 $S=\Spec(A)$의 한 원소이며, 이를 다시
${}^a\vphi(x)$로 나타내겠다. 또한 모든 $f\in A$에 대하여 정의상
(\hyperref[0.5.5.2-ko]{0, 5.5.2})
\[
  {}^a\vphi^{-1}(D(f))=X_f,
\]
이므로 ${}^a\vphi$는 연속사상 $X\to S$이다. 이제 $\sh{O}_S$-가군의
준동형
\[
  \widetilde{\vphi}:\sh{O}_S\to{}^a\vphi_*(\sh{O}_X)
\]
을 정의하자. 모든 $f\in A$에 대하여
$\Gamma(D(f),\sh{O}_S)=A_f$이다
(\hyperref[I.1.3.6-ko]{I, 1.3.6}). 모든 $s\in A$에 대하여
$s/f\in A_f$에는 원소
\[
  (\vphi(s)|X_f)(\vphi(f)|X_f)^{-1}
  \in\Gamma(X_f,\sh{O}_X)
  =\Gamma(D(f),{}^a\vphi_*(\sh{O}_X))
\]
를 대응시키자. 그러면 $D(f)$에서 $D(fg)$로 넘어감으로써 이것이 실제로
$\sh{O}_S$-가군의 준동형을 정의함을 곧바로 확인할 수 있고, 따라서 환
달린 공간의 사상 $({}^a\vphi,\widetilde{\vphi})$를 얻는다. 또한 같은
표기를 쓰고 간단히 $y={}^a\vphi(x)$로 놓으면
(\hyperref[0.3.7.1-ko]{0, 3.7.1})
\[
  \widetilde{\vphi}_x^\sharp(s_y/f_y)
  =(\vphi(s)_x)(\vphi(f)_x)^{-1};
\]
이다. 관계 $s_y\in\mathfrak m_y$는 정의상
$\vphi(s)_x\in\mathfrak m_x$와 동치이므로
$\widetilde{\vphi}_x^\sharp$는 국소 준동형
$\sh{O}_y\to\sh{O}_x$이다. 이로써 두 번째 사상
\[
  \sigma:\Hom(\Gamma(S,\sh{O}_S),\Gamma(X,\sh{O}_X))\to\mathfrak L
\]
을 정의하였다. 여기서 $\mathfrak L$은 모든 $x\in X$에 대하여
$\theta_x^\sharp$이 국소 준동형인 사상
$(\psi,\theta):(X,\sh{O}_X)\to(S,\sh{O}_S)$들의 집합이다.

이제 $\sigma$와 $\rho$의 $\mathfrak L$에 대한 제한이 서로 역임을
보이면 된다. 그런데 $\widetilde{\vphi}$의 정의에서 곧바로
$\Gamma(\widetilde{\vphi})=\vphi$를 얻으므로 $\rho\circ\sigma$는
항등사상이다. $\sigma\circ\rho$가 항등사상임을 보이기 위하여 사상
$(\psi,\theta)\in\mathfrak L$에서 출발하여
$\vphi=\Gamma(\theta)$로 놓자. $\theta_x^\sharp$에 대한 가정으로부터
몫으로 내려가 단사 준동형
\[
  \theta^x:\kres(\psi(x))\to\kres(x)
\]
을 얻는다. 모든 단면 $f\in A=\Gamma(S,\sh{O}_S)$에 대하여
$\theta^x(f(\psi(x)))=\vphi(f)(x)$이고, 따라서 관계
$f(\psi(x))=0$은 $\vphi(f)(x)=0$과 동치이다. 이는 우선
${}^a\vphi=\psi$임을 보인다. 한편 정의들에 따라 도식
\[
  \xymatrix{
    A\ar[r]^\vphi\ar[d] &
    \Gamma(X,\sh{O}_X)\ar[d]\\
    A_{\psi(x)}\ar[r]^{\theta_x^\sharp} &
    \sh{O}_x
  }
\]
은 가환이다. $\theta_x^\sharp$을
$\widetilde{\vphi}_x^\sharp$으로 바꾼 유사한 도식도 가환하므로
$\widetilde{\vphi}_x^\sharp=\theta_x^\sharp$이다
(\hyperref[0.1.2.4-ko]{0, 1.2.4}). 따라서
$\widetilde{\vphi}=\theta$이다.

\begin{env}[1.8.2]
\label{ega2:addendum:I.1.8.2-ko}
$(X,\sh{O}_X)$와 $(Y,\sh{O}_Y)$가 국소환 달린 공간일 때에는, 모든
$x\in X$에 대하여 $\theta_x^\sharp$이 국소 준동형
$\sh{O}_{\psi(x)}\to\sh{O}_x$인 사상
$(\psi,\theta):(X,\sh{O}_X)\to(Y,\sh{O}_Y)$를 고려할 것이다.

이제부터
\oldpage[II]{219}
\markboth{정오표 및 추록, 목록 1}{정오표 및 추록, 목록 1}
\emph{국소환 달린 공간의 사상}이라고 할 때에는 언제나 이러한 사상을
뜻한다. 이러한 사상의 정의 아래 국소환 달린 공간들이 하나의 범주를
이룸은 분명하다. 따라서 이 범주의 두 대상 $X,Y$에 대하여 $\Hom(X,Y)$는
$X$에서 $Y$로 가는 국소환 달린 공간의 사상들의 집합
(\hyperref[ega2:addendum:I.1.8.1-ko]{1.8.1}에서 $\mathfrak L$로 표기한
집합)을 뜻한다. 혼동을 피하기 위하여 $X$에서 $Y$로 가는 환 달린 공간의
사상들의 집합을 고려할 때에는 이를 $\Hom_{\mathrm{an}}(X,Y)$로
나타내겠다. 따라서 사상 (\hyperref[ega2:addendum:I.1.8.1.1-ko]{1.8.1.1})은
\begin{equation}
\label{ega2:addendum:I.1.8.2.1-ko}
  \rho:\Hom_{\mathrm{an}}(X,Y)
  \to\Hom(\Gamma(Y,\sh{O}_Y),\Gamma(X,\sh{O}_X)),
\tag{1.8.2.1}
\end{equation}
로 쓰이며, 그 제한
\begin{equation}
\label{ega2:addendum:I.1.8.2.2-ko}
  \rho':\Hom(X,Y)
  \to\Hom(\Gamma(Y,\sh{O}_Y),\Gamma(X,\sh{O}_X))
\tag{1.8.2.2}
\end{equation}
은 국소환 달린 공간의 범주에서 $X$와 $Y$에 관하여 함자적인 사상이다.
\end{env}
\begin{corollary}[1.8.3]
\label{ega2:addendum:I.1.8.3-ko}
$(Y,\sh{O}_Y)$를 국소환 달린 공간이라 하자. $Y$가 아핀 스킴이기 위한
필요충분조건은 모든 국소환 달린 공간 $(X,\sh{O}_X)$에 대하여 사상
(\hyperref[ega2:addendum:I.1.8.2.2-ko]{1.8.2.2})가 전단사인 것이다.
\end{corollary}

명제 (\hyperref[ega2:addendum:I.1.8.1-ko]{1.8.1})은 이 조건이 필요함을
보여 준다. 역으로 이 조건이 성립한다고 가정하고
$A=\Gamma(Y,\sh{O}_Y)$라 놓자. 가정과
(\hyperref[ega2:addendum:I.1.8.1-ko]{1.8.1})에 따라, 국소환 달린 공간의
범주에서 집합의 범주로 가는 두 함자 $X\mapsto\Hom(X,Y)$와
$X\mapsto\Hom(X,\Spec(A))$는 서로 동형이다. 이로부터 표준 동형
$Y\to\Spec(A)$가 존재한다\footnote{\textbf{번역 주.} 인쇄된 프랑스어
원문은 $X\to\Spec(A)$라고 쓰지만, 바로 앞에서 비교한 두 함자의 표현
대상은 $Y$와 $\Spec(A)$이다. 따라서 EGA2-CANON-DISP-0026에 따라
$Y\to\Spec(A)$로 교정했으며, 프랑스어 외교적 전사는 인쇄형을
보존한다.} (0, 8 참조).

\begin{env}[1.8.4]
\label{ega2:addendum:I.1.8.4-ko}
$S=\Spec(A)$를 아핀 스킴이라 하자. 바탕공간이 한 점으로 이루어지고
구조층 $A'$가 환 $A$로 정의되는 $S'$ 위의 층(원문의 용어로 ``단순층'',
이 한 점 공간에서는 필연적으로 상수층)인
환 달린 공간을 $(S',A')$로 나타내자. $\pi:S\to S'$를 $S$에서 $S'$로
가는 유일한 사상이라 하자. 한편 $S$의 모든 열린집합 $U$에 대하여
표준 사상
$\Gamma(S',A')=\Gamma(S,\sh{O}_S)\to\Gamma(U,\sh{O}_S)$가 있으며,
이는 환의 층의 $\pi$-사상 $\iota:A'\to\sh{O}_S$를 정의한다. 따라서
환 달린 공간의 사상
$i=(\pi,\iota):(S,\sh{O}_S)\to(S',A')$를 표준적으로 정의하였다.
모든 $A$-가군 $M$에 대하여, $M$으로 정의되는 $S'$ 위의 상수층을
$M'$으로 나타내겠다. 이는 명백히 $A'$-가군이다. 그리고
$i_*(\widetilde M)=M'$이다
(\hyperref[I.1.3.7-ko]{I, 1.3.7}).
\end{env}

\begin{lemma}[1.8.5]
\label{ega2:addendum:I.1.8.5-ko}
(\hyperref[ega2:addendum:I.1.8.4-ko]{1.8.4})의 표기에서, 모든
$A$-가군 $M$에 대하여 함자적인 표준 $\sh{O}_S$-준동형
(\hyperref[0.4.4.3.3-ko]{0, 4.4.3.3})
\begin{equation}
\label{ega2:addendum:I.1.8.5.1-ko}
  i^*(i_*(\widetilde M))\to\widetilde M
\tag{1.8.5.1}
\end{equation}
은 동형이다.
\end{lemma}

실제로 (\hyperref[ega2:addendum:I.1.8.5.1-ko]{1.8.5.1})의 양변은
$M$에 관하여 오른쪽 완전하고(함자 $M\mapsto i_*(\widetilde M)$는
명백히 완전하다) 직합과 가환한다. $M$을 준동형
$A^{(I)}\to A^{(J)}$의 여핵으로 보면, 이 보조정리를 $M=A$인 경우에
증명하는 것으로 환원할 수 있으며, 이 경우에는 자명하다.

\begin{corollary}[1.8.6]
\label{ega2:addendum:I.1.8.6-ko}
$(X,\sh{O}_X)$를 환 달린 공간이라 하고, $u:X\to S$를 환 달린 공간의
\oldpage[II]{220}
\markboth{A. 그로텐디크 --- 제II장}{A. 그로텐디크 --- 제II장}
사상이라 하자. 모든 $A$-가군 $M$에 대하여
(\hyperref[ega2:addendum:I.1.8.4-ko]{1.8.4})의 표기 아래
$\sh{O}_X$-가군의 함자적인 표준 동형
\begin{equation}
\label{ega2:addendum:I.1.8.6.1-ko}
  u^*(\widetilde M)\isoto u^*(i^*(M')).
\tag{1.8.6.1}
\end{equation}
이 있다.
\end{corollary}

\begin{corollary}[1.8.7]
\label{ega2:addendum:I.1.8.7-ko}
(\hyperref[ega2:addendum:I.1.8.6-ko]{1.8.6})의 가정 아래, 모든
$A$-가군 $M$과 모든 $\sh{O}_X$-가군 $\sh{F}$에 대하여 $M$과
$\sh{F}$에 관한 함자적인 표준 동형
\begin{equation}
\label{ega2:addendum:I.1.8.7.1-ko}
  \Hom_{\sh{O}_S}(\widetilde M,u_*(\sh{F}))
  \isoto\Hom_A(M,\Gamma(X,\sh{F})).
\tag{1.8.7.1}
\end{equation}
이 있다.
\end{corollary}

실제로 (\hyperref[0.4.4.3-ko]{0, 4.4.3})과 보조정리
(\hyperref[ega2:addendum:I.1.8.5-ko]{1.8.5})에 따라 쌍함자들의 표준
동형
\[
  \Hom_{\sh{O}_S}(\widetilde M,u_*(\sh{F}))
  \isoto\Hom_{A'}(M',i_*(u_*(\sh{F})))
\]
이 있고, 오른쪽 변은 바로 $\Hom_A(M,\Gamma(X,\sh{F}))$이다. 표준
준동형 (\hyperref[ega2:addendum:I.1.8.7.1-ko]{1.8.7.1})은 모든
$\sh{O}_S$-준동형 $h:\widetilde M\to u_*(\sh{F})$(다시 말해 모든
$u$-사상 $\widetilde M\to\sh{F}$)에 $A$-준동형
\[
  \Gamma(h):M\to\Gamma(S,u_*(\sh{F}))=\Gamma(X,\sh{F})
\]
을 대응시킨다는 점에 유의하자.

\begin{env}[1.8.8]
\label{ega2:addendum:I.1.8.8-ko}
(\hyperref[ega2:addendum:I.1.8.4-ko]{1.8.4})의 표기에서, 환 달린 공간의
사상 $(\psi,\theta):X\to S'$를 주는 것은 환 준동형
$A\to\Gamma(X,\sh{O}_X)$를 주는 것과 동치임이 명백하다
(\hyperref[0.4.1.1-ko]{0, 4.1.1}). 따라서
(\hyperref[ega2:addendum:I.1.8.1-ko]{1.8.1})을 표준 전단사
\[
  \Hom(X,S)\isoto\Hom(X,S')
\]
를 정의하는 것으로도 해석할 수 있다. 물론 오른쪽 변에서는 환 달린
공간의 사상을 뜻한다. 일반적으로 $A$가 국소환은 아니기 때문이다.
더 일반적으로, $X,Y$를 두 국소환 달린 공간이라 하고, 바탕공간이 한
점으로 이루어지며 환의 층 $A'$가 환 $\Gamma(Y,\sh{O}_Y)$로 정의되는
상수층인 환 달린 공간을 $(Y',A')$라 하자. 그러면
(\hyperref[ega2:addendum:I.1.8.2.1-ko]{1.8.2.1})을 사상
\begin{equation}
\label{ega2:addendum:I.1.8.8.1-ko}
  \rho:\Hom_{\mathrm{an}}(X,Y)\to\Hom(X,Y')
\tag{1.8.8.1}
\end{equation}
으로 해석할 수 있다. 따라서 따름정리
(\hyperref[ega2:addendum:I.1.8.3-ko]{1.8.3})의 결과는, 국소환 달린
공간들 가운데 아핀 스킴이 다음 성질을 갖는 것들로 특징지어진다는
뜻이다. 모든 국소환 달린 공간 $X$에 대하여 $\rho$를
$\Hom(X,Y)$에 제한한 사상
\begin{equation}
\label{ega2:addendum:I.1.8.8.2-ko}
  \rho':\Hom(X,Y)\to\Hom(X,Y')
\tag{1.8.8.2}
\end{equation}
이 전단사이다. 뒤의 한 장에서 이 정의를 일반화하여, 임의의 환 달린
공간 $Z$(바탕공간이 한 점으로 이루어진 환 달린 공간만이 아니라)에
대해 다시 $\Spec(Z)$라 부를 국소환 달린 공간을 대응시킬 것이다. 이는
임의의 환 달린 공간 위의 준스킴에 관한 «~상대적~» 이론의 출발점이
되며, 제I장의 결과들을 확장할 것이다.
\end{env}
\begin{env}[1.8.9]
\label{ega2:addendum:I.1.8.9-ko}
국소환 달린 공간 $X$와 $\sh{O}_X$-가군 $\sh{F}$로 이루어진 쌍
$(X,\sh{F})$들을 하나의 범주의 대상들로 볼 수 있다. 이 범주에서
$(X,\sh{F})\to(Y,\sh{G})$인 사상은 쌍 $(u,h)$이다. 여기서 $u:X\to Y$는
국소환 달린 공간의
\oldpage[II]{221}
\markboth{정오표 및 추록, 목록 1}{정오표 및 추록, 목록 1}
사상이고 $h:\sh{G}\to\sh{F}$는 가군의 $u$-사상이다. 고정된
$(X,\sh{F})$와 $(Y,\sh{G})$에 대하여 이러한 사상들은 하나의 집합을
이루며, 이를 $\Hom((X,\sh{F}),(Y,\sh{G}))$으로 나타내겠다. 대응
$(u,h)\mapsto(\rho'(u),\Gamma(h))$는 표준 사상
\begin{equation}
\label{ega2:addendum:I.1.8.9.1-ko}
\begin{split}
  \Hom((X,\sh{F}),(Y,\sh{G}))\to
  \Hom\bigl(
    &(\Gamma(Y,\sh{O}_Y),\Gamma(Y,\sh{G})),\\
    &(\Gamma(X,\sh{O}_X),\Gamma(X,\sh{F}))
  \bigr),
\end{split}
\tag{1.8.9.1}
\end{equation}
을 주며, 이는 $(X,\sh{F})$와 $(Y,\sh{G})$에 관하여 함자적이다. 이때
우변은 고려 중인 환들과 가군들에 대응하는 디-준동형들의 집합이다
(0, 1.0.2).
\end{env}

\begin{corollary}[1.8.10]
\label{ega2:addendum:I.1.8.10-ko}
$Y$를 국소환 달린 공간, $\sh{G}$를 $\sh{O}_Y$-가군이라 하자. $Y$가
아핀 스킴이고 $\sh{G}$가 준연접 $\sh{O}_Y$-가군이기 위한 필요충분조건은
국소환 달린 공간 $X$와 $\sh{O}_X$-가군 $\sh{F}$로 이루어진 모든 쌍
$(X,\sh{F})$에 대하여 표준 사상
\hyperref[ega2:addendum:I.1.8.9.1-ko]{(1.8.9.1)}\footnote{\emph{[번역 주]
원문에 인쇄된 (1.8.8.1)을 (1.8.9.1)로 바로잡았다.}}이 전단사인 것이다.
\end{corollary}

이 추론을 자세히 전개하는 일은 독자에게 맡긴다. 그 추론은
\hyperref[ega2:addendum:I.1.8.3-ko]{(1.8.3)}의 추론을 본뜬 것으로,
\hyperref[ega2:addendum:I.1.8.1-ko]{(1.8.1)}과
\hyperref[ega2:addendum:I.1.8.7-ko]{(1.8.7)}을 사용한다.

\begin{remark}[1.8.11]
\label{ega2:addendum:I.1.8.11-ko}
명제 (I, 1.7.3), (I, 1.7.4), (I, 2.2.4)와 정의 (I, 1.6.1)은
\hyperref[ega2:addendum:I.1.8.1-ko]{(1.8.1)}의 특수한 경우이고, 마찬가지로
(I, 2.2.5)는 \hyperref[ega2:addendum:I.1.8.7-ko]{(1.8.7)}에서 따른다.
명제 \hyperref[ega2:addendum:I.1.8.7-ko]{(1.8.7)}은 또한 (I, 1.6.3)을
특수한 경우로서 함의하고, 따라서 (I, 1.6.4)도 함의한다. 실제로 $X$가
아핀이고 $\Gamma(X,\sh{F})=N$이면, 두 함자
\[
  M\longmapsto
  \Hom_{\sh{O}_S}(\widetilde M,u_*(\widetilde N))
  \quad\hbox{와}\quad
  M\longmapsto
  \Hom_{\sh{O}_S}(\widetilde M,(N_{[\vphi]})\supertilde)
\]
는 \hyperref[ega2:addendum:I.1.8.7-ko]{(1.8.7)}과 (I, 1.3.8)에 따라 서로
동형이다. 여기서 $\vphi:A\to\Gamma(X,\sh{O}_X)$는 $u$에 대응한다.
끝으로 (I, 1.6.5), 따라서 (I, 1.6.6)은
\hyperref[ega2:addendum:I.1.8.6-ko]{(1.8.6)}과, 모든 $f\in A$에 대하여
두 $A_f$-가군
\[
  N'\otimes_{A'}A_f
  \quad\hbox{와}\quad
  (N'\otimes_{A'}A)_f
\]
가 표준적으로 동형이라는 사실에서 따른다. 여기서는 (I, 1.6.5)의 표기를
사용하였다.
\end{remark}

\paragraph{(I, 3.2.8) 뒤에 다음을 삽입한다:}

\begin{env}[3.2.9]
\label{ega2:addendum:I.3.2.9-ko}
\hyperref[ega2:addendum:I.1.8.1-ko]{(1.8.1)}에서 (I, 3.2.2)를 다음과 같이
더 정확히 할 수 있음이 따른다. $Z=\Spec(B\otimes_A C)$는
$S$-준스킴들의 범주에서 $X=\Spec(B)$와 $Y=\Spec(C)$의 곱일 뿐만 아니라,
$S$ 위의 국소환 달린 공간들의 범주에서도 그 곱이다. 여기서
$S$-사상의 정의는 (I, 2.5.2)의 정의를 본뜬다. 그러면 (I, 3.2.6)의
증명은 실제로 임의의 두 $S$-준스킴 $X,Y$에 대하여 준스킴
$X\times_S Y$가 $S$-준스킴들의 범주에서 $X$와 $Y$의 곱일 뿐만 아니라,
준스킴 $S$ 위의 국소환 달린 공간들의 범주에서도 그 곱임을 증명한다.
\end{env}

\paragraph{(I, 4.4.5).}
126쪽 아래에서 14번째 줄의 $B$를 $A$로 바꾸고, 아래에서 13번째 줄의
``$A$-대수''를 ``$B$-대수''로 바꾼다.

\paragraph{(I, 5.3.5).}
132쪽 아래에서 2번째 줄의 $Y\to S$를 $g:Y\to S$로 바꾼다.

\paragraph{(I, 6.3.2.1).}
144쪽 아래에서 6번째 줄은 다음과 같이 읽는다.
\[
  D(g_i)\subset W.
\]

\paragraph{(I, 7.3.8).}
오류가 있는 현재의 본문을 다음 본문으로 바꾼다.

$X,Y$를 두 정역 준스킴이라 하자. 그러면 $\sh{R}(X)$는 준연접
$\sh{O}_X$-가군이고 $\sh{R}(Y)$는 준연접 $\sh{O}_Y$-가군이다
(I, 7.3.3). $f:X\to Y$를 우세 사상이라 하자. 그러면 표준
$\sh{O}_X$-가군 준동형
\oldpage[II]{222}
\markboth{A. 그로텐디크 --- 제II장}{A. 그로텐디크 --- 제II장}
\begin{equation}
\label{ega2:addendum:I.7.3.8.1-ko}
  \tau:f^*(\sh{R}(Y))\to\sh{R}(X).
\tag{7.3.8.1}
\end{equation}
이 존재한다.

먼저 $X=\Spec(A)$와 $Y=\Spec(B)$가 각각 정역 $A$와 $B$의 아핀
준스킴이라고 하자. 그러면 $f$는 단사 준동형 $B\to A$에 대응하고, 이
준동형은 $B$의 분수체 $L$에서 $A$의 분수체 $K$로 가는 단사 준동형
$L\to K$로 연장된다. 이때
\hyperref[ega2:addendum:I.7.3.8.1-ko]{(7.3.8.1)}의 준동형은 표준 준동형
\[
  L\otimes_B A\to K
\]
에 대응한다 (I, 1.6.5). 일반적인 경우에는 $f(U)\subset V$인 공집합이
아닌 아핀 열린집합 $U\subset X$, $V\subset Y$의 모든 쌍에 대하여 앞과
같이 준동형 $\tau_{U,V}$를 정의한다. $U'\subset U$, $V'\subset V$이고
$f(U')\subset V'$이면 $\tau_{U,V}$가 $\tau_{U',V'}$의 연장임을 곧바로
알 수 있고, 이로부터 주장이 따른다. $x,y$가 각각 $X,Y$의 일반점이면
$f(x)=y$이고,
\[
  \bigl(f^*(\sh{R}(Y))\bigr)_x
  =\sh{O}_y\otimes_{\sh{O}_y}\sh{O}_x
  =\sh{O}_x
\]
이다 (0, 4.3.1). 따라서 $\tau_x$는 동형사상이다.

\paragraph{(I, 9.5.2).}
176쪽의 마지막 세 줄과 177쪽의 처음 다섯 줄에서 $X'$은 모두 $Y'$으로,
$X''$은 모두 $Y''$으로 바꾼다.

\paragraph{(I, 10.5.4).}
187쪽 아래에서 11번째 줄의 (10.3.4)를 (10.3.5)로 바꾼다.

\paragraph{(I, 10.6.3).}
189쪽 아래에서 3번째 줄의 $X$를 $\mathfrak X$로 바꾼다.

\paragraph{(I, 10.14.2).}
210쪽의 5번째 줄에서 ``연접 $\sh{O}_{\mathfrak X}$-가군''을
``$\sh{O}_{\mathfrak X}$의 연접 아이디얼층''으로 바꾸고, 7번째 줄에서
``연접 $\sh{O}_{\mathfrak X}$-가군들''을
``$\sh{O}_{\mathfrak X}$의 연접 아이디얼층들''로 바꾼다.
% ===== END EXACT INLINE INPUT: c2errata-addenda.tex =====

\clearpage

% ===== BEGIN EXACT INLINE INPUT: c3-1.tex =====
\clearpage
\thispagestyle{plain}
\markboth{예비사항}{예비사항}
\oldpage[0\textsubscript{III}]{5}
\phantomsection
\label{chapter:0III-fr}

\begin{center}
{\large 제0장 (계속)\footnotemark\par}
\vspace{1.5em}
{\Large\bfseries 예비사항\par}
\end{center}

\footnotetext{참조를 쉽게 찾을 수 있도록, 이제부터 제I장과 함께 출판된
제0장의 절은 기호 $0_{\mathrm I}$로 인용한다.}

\setcounter{section}{7}
\section{표현 가능 함자}
\label{section:0.8-fr}

\setcounter{subsection}{0}
\subsection{표현 가능 함자}
\label{subsection:0.8.1-fr}

\begin{env}[8.1.1]
\label{0.8.1.1-fr}
집합의 범주를 $\mathrm{Ens}$로 나타내기로 한다. $C$를 하나의
범주라 하자~; $C$의 두 대상 $X$, $Y$에 대하여
$h_X(Y)=\operatorname{Hom}(Y,X)$로 놓고~; $C$의 임의의 사상
$u:Y\to Y'$에 대하여 $h_X(u)$를 $\operatorname{Hom}(Y',X)$에서
$\operatorname{Hom}(Y,X)$로 가는 함수 $v\mapsto vu$로 나타낸다. 이
정의들에 따르면 $h_X:C\to\mathrm{Ens}$가 \emph{반변 함자}, 즉 범주
$C$의 쌍대범주인 $C^0$에서 범주 $\mathrm{Ens}$로 가는 공변 함자들의
범주 $\operatorname{Hom}(C^0,\mathrm{Ens})$의 대상임은 자명하다
(T, 1.7, d) 및 [29]).
\end{env}

\begin{env}[8.1.2]
\label{0.8.1.2-fr}
이제 $w:X\to X'$을 $C$의 사상이라 하자~; 임의의 $Y\in C$와 임의의
$v\in\operatorname{Hom}(Y,X)=h_X(Y)$에 대하여
$wv\in\operatorname{Hom}(Y,X')=h_{X'}(Y)$이다~; $h_w(Y)$를
$h_X(Y)$에서 $h_{X'}(Y)$로 가는 함수 $v\mapsto wv$로 나타내자.
$C$의 임의의 사상 $u:Y\to Y'$에 대하여 다음 도식은
\[
  \xymatrix{
    h_X(Y')\ar[r]^{h_X(u)}\ar[d]_{h_w(Y')} &
    h_X(Y)\ar[d]^{h_w(Y)}\\
    h_{X'}(Y')\ar[r]_{h_{X'}(u)} &
    h_{X'}(Y)
  }
\]
가환한다~; 다시 말해 $h_w$는 \emph{자연변환}
$h_X\to h_{X'}$이다 (T, 1.2).
또는 범주 $\operatorname{Hom}(C^0,\mathrm{Ens})$의 사상이다
(T, 1.7, d)). 따라서 $h_X$와 $h_w$의 정의는 \emph{표준 공변 함자}
\[
  h:C\to\operatorname{Hom}(C^0,\mathrm{Ens}),\qquad X\mapsto h_X.
  \tag{8.1.2.1}
\]
를 정의한다.
\end{env}

\begin{env}[8.1.3]
\label{0.8.1.3-fr}
$X$를 $C$의 대상, $F$를 $C$에서 $\mathrm{Ens}$로 가는 반변 함자
($\operatorname{Hom}(C^0,\mathrm{Ens})$의 대상)라 하자.
$g:h_X\to F$를 \emph{자연변환}이라 하자~: 임의의 $Y\in C$에 대하여
\oldpage[0\textsubscript{III}]{6}
$g(Y)$는 $h_X(Y)$에서 $F(Y)$로 가는 함수이고, $C$의 임의의 사상
$u:Y\to Y'$에 대하여 다음 도식이 가환한다.
\[
  \xymatrix{
    h_X(Y')\ar[r]^{h_X(u)}\ar[d]_{g(Y')} &
    h_X(Y)\ar[d]^{g(Y)}\\
    F(Y')\ar[r]_{F(u)} & F(Y)
  }
  \tag{8.1.3.1}\label{0.8.1.3.1-fr}
\]
특히 함수 $g(X):h_X(X)=\operatorname{Hom}(X,X)\to F(X)$가 있으므로,
다음 원소를 얻는다.
\[
  \alpha(g)=(g(X))(1_X)\in F(X)
  \tag{8.1.3.2}
\]
따라서 표준 함수
\[
  \alpha:\operatorname{Hom}(h_X,F)\to F(X).
  \tag{8.1.3.3}
\]
가 정의된다.

반대로 원소 $\xi\in F(X)$를 생각하자~; $C$의 임의의 사상
$v:Y\to X$에 대하여 $F(v)$는 $F(X)$에서 $F(Y)$로 가는 함수이다~;
다음 함수를 생각하자.
\[
  v\mapsto(F(v))(\xi)
  \tag{8.1.3.4}
\]
이는 $h_X(Y)$에서 $F(Y)$로 가는 함수이다~; 이 함수를
$(\beta(\xi))(Y)$로 나타내면,
\[
  \beta(\xi):h_X\to F
  \tag{8.1.3.5}
\]
는 \emph{자연변환}이다. 실제로 $C$의 임의의 사상 $u:Y\to Y'$에
대하여 $(F(vu))(\xi)=(F(u)\circ F(v))(\xi)$이고, 이는
$g=\beta(\xi)$일 때 (\hyperref[0.8.1.3.1-fr]{8.1.3.1})의 가환성을
확인해 준다. 이로써 표준 함수
\[
  \beta:F(X)\to\operatorname{Hom}(h_X,F).
  \tag{8.1.3.6}
\]
를 정의하였다.
\end{env}

\begin{proposition}[8.1.4]
\label{0.8.1.4-fr}
함수 $\alpha$와 $\beta$는 서로 역인 전단사이다.
\end{proposition}

$\xi\in F(X)$에 대하여 $\alpha(\beta(\xi))$를 계산하자~; 임의의
$Y\in C$에 대하여 $(\beta(\xi))(Y)$는 $h_X(Y)$에서 $F(Y)$로 가는
함수 $g_1(Y):v\mapsto(F(v))(\xi)$이다. 따라서
\[
  \alpha(\beta(\xi))=(g_1(X))(1_X)=(F(1_X))(\xi)
  =1_{F(X)}(\xi)=\xi.
\]
이제 $g\in\operatorname{Hom}(h_X,F)$에 대하여
$\beta(\alpha(g))$를 계산하자~; 임의의 $Y\in C$에 대하여
$(\beta(\alpha(g)))(Y)$는 함수
$v\mapsto(F(v))((g(X))(1_X))$이다~; (\hyperref[0.8.1.3.1-fr]{8.1.3.1})의
가환성에 따라, 이 함수는 $h_X(v)$의 정의에 의해
$v\mapsto(g(Y))((h_X(v))(1_X))=(g(Y))(v)$와 다르지 않다. 다시 말해
$g(Y)$와 같으며, 이로써 명제가 증명된다.

\begin{env}[8.1.5]
\label{0.8.1.5-fr}
범주 $C$의 \emph{부분범주} $C'$은 다음 조건으로 정의됨을 상기하자.
$C'$의 대상들은 $C$의 대상들이고, $X'$, $Y'$이 $C'$의 두 대상이면
\emph{$C'$ 안에서} $X'\to Y'$인 사상들의 집합
$\operatorname{Hom}_{C'}(X',Y')$은 \emph{$C$ 안에서} $X'\to Y'$인
사상들의 집합 $\operatorname{Hom}_C(X',Y')$의 부분집합이며, 표준적인
<<~사상의 합성~>> 함수
\[
  \operatorname{Hom}_{C'}(X',Y')\times\operatorname{Hom}_{C'}(Y',Z')
  \to\operatorname{Hom}_{C'}(X',Z')
\]
\oldpage[0\textsubscript{III}]{7}
는 다음 표준 함수의 제한이다.
\[
  \operatorname{Hom}_C(X',Y')\times\operatorname{Hom}_C(Y',Z')
  \to\operatorname{Hom}_C(X',Z').
\]

$C'$의 임의의 두 대상에 대하여
$\operatorname{Hom}_{C'}(X',Y')=\operatorname{Hom}_C(X',Y')$이면,
$C'$을 $C$의 \emph{충만한 부분범주}라고 한다. 이때 $C'$의 대상들과
동형인 $C$의 대상들로 이루어진 $C$의 부분범주 $C''$도 $C$의 충만한
부분범주이며, 쉽게 확인할 수 있듯이 $C'$과 \emph{동치}이다 (T, 1.2).

공변 함자 $F:C_1\to C_2$에 대하여, $C_1$의 임의의 두 대상 $X_1$,
$Y_1$에 대해 $\operatorname{Hom}(X_1,Y_1)$에서
$\operatorname{Hom}(F(X_1),F(Y_1))$로 가는 함수 $u\mapsto F(u)$가
\emph{전단사}이면, $F$가 \emph{충만하고 충실하다}고 한다~; 이 조건은
$C_2$의 부분범주 $F(C_1)$이 \emph{충만함}을 함의한다. 더욱이 두 대상
$X_1$, $X'_1$의 상이 같은 대상 $X_2$이면, $F(u)=1_{X_2}$를 만족하는
동형사상 $u:X_1\to X'_1$가 유일하게 존재한다. 이제 $F(C_1)$의 임의의
대상 $X_2$에 대하여 $F(X_1)=X_2$를 만족하는 $C_1$의 대상 $X_1$ 중
하나를 $G(X_2)$로 택하자 ($G$는 선택공리를 사용하여 정의한다)~;
$F(C_1)$의 임의의 사상 $v:X_2\to Y_2$에 대하여 $G(v)$를
$F(u)=v$를 만족하는 유일한 사상 $u:G(X_2)\to G(Y_2)$로 정한다~;
그러면 $G$는 $F(C_1)$에서 $C_1$로 가는 \emph{함자}이다~;
$FG$는 $F(C_1)$의 항등함자이고, 앞의 논의에 따라 함자의 동형사상
$\varphi:1_{C_1}\to GF$가 존재하여 $F$, $G$, $\varphi$ 및 항등사상
$1_{F(C_1)}\to FG$가 범주 $C_1$과 $C_2$의 충만한 부분범주
$F(C_1)$ 사이의 \emph{동치}를 정의한다 (T, 1.2).
\end{env}

\begin{env}[8.1.6]
\label{0.8.1.6-fr}
함자 $F$가 $h_{X'}$이고 $X'$이 $C$의 임의의 대상인 경우에 명제
(\hyperref[0.8.1.4-fr]{8.1.4})를 적용하자~; 이때 함수
$\beta:\operatorname{Hom}(X,X')\to\operatorname{Hom}(h_X,h_{X'})$는
(\hyperref[0.8.1.2-fr]{8.1.2})에서 정의한 함수 $w\mapsto h_w$와
다르지 않다~; 이 함수가 \emph{전단사}이므로,
(\hyperref[0.8.1.5-fr]{8.1.5})의 용어를 사용하면 다음을 얻는다~:
\end{env}

\begin{proposition}[8.1.7]
\label{0.8.1.7-fr}
표준 함자 $h:C\to\operatorname{Hom}(C^0,\mathrm{Ens})$는 충만하고
충실하다.
\end{proposition}

\begin{env}[8.1.8]
\label{0.8.1.8-fr}
$F$를 $C$에서 $\mathrm{Ens}$로 가는 반변 함자라 하자~; 어떤 대상
$X\in C$가 존재하여 $F$가 $h_X$와 \emph{동형}이면, $F$가
\emph{표현 가능하다}고 한다~; (\hyperref[0.8.1.7-fr]{8.1.7})에 따라
$X\in C$와 함자의 동형사상 $g:h_X\to F$를 주면 $X$는 유일한
동형사상까지 결정된다. 명제 (\hyperref[0.8.1.7-fr]{8.1.7})는 또한
$h$가 $C$와, $\operatorname{Hom}(C^0,\mathrm{Ens})$에서 \emph{표현 가능한
반변 함자}들로 이루어진 충만한 부분범주 사이의 \emph{동치}를
정의한다는 뜻이다. 한편 (\hyperref[0.8.1.4-fr]{8.1.4})에 따라
자연변환 $g:h_X\to F$를 주는 것은 원소 $\xi\in F(X)$를 주는 것과
동치이다~; $g$가 \emph{동형사상}이라는 것은 $\xi$에 대하여 다음
조건과 동치이다~: \emph{$C$의 임의의 대상 $Y$에 대하여
$\operatorname{Hom}(Y,X)$에서 $F(Y)$로 가는 함수
$v\mapsto(F(v))(\xi)$가 전단사이다}. $\xi$가 이 조건을 만족하면,
쌍 $(X,\xi)$가 표현 가능 함자 $F$를 \emph{표현한다}고 한다. 용어를
다소 남용하여, 어떤 $\xi\in F(X)$에 대하여 $(X,\xi)$가 $F$를
표현할 때, 즉 $h_X$가 $F$와 동형일 때, 대상 $X\in C$가 $F$를
표현한다고도 한다.

$F$, $F'$을 $C$에서 $\mathrm{Ens}$로 가는 표현 가능한 두 반변
함자라 하고, $h_X\to F$와 $h_{X'}\to F'$을 함자의 두 동형사상이라
하자. 그러면 (\hyperref[0.8.1.6-fr]{8.1.6})에 따라
$\operatorname{Hom}(X,X')$과, 자연변환 $F\to F'$들의 집합
$\operatorname{Hom}(F,F')$ 사이에 표준적인 일대일 대응이 있다.
\end{env}

\begin{env}[8.1.9]
\label{0.8.1.9-fr}
\textit{예. I: 사영극한.} 표현 가능한 반변 함자의 개념은 특히
통상적인 <<~보편 문제의
\oldpage[0\textsubscript{III}]{8}
해~>>라는 개념의 <<~쌍대~>> 개념을 포괄한다. 더 일반적으로,
\emph{사영극한}의 개념이 표현 가능 함자의 개념의 특수한 경우임을
살펴보겠다. 범주 $C$에서 \emph{사영계}는 준순서 집합 $I$, $C$의
대상들의 족 $(A_\alpha)_{\alpha\in I}$, 그리고
$\alpha\leq\beta$인 임의의 두 지표 $(\alpha,\beta)$에 대한 사상
$u_{\alpha\beta}:A_\beta\to A_\alpha$를 주어 정의함을 상기하자.
$C$에서 이 사영계의 \emph{사영극한}은 $C$의 대상 $B$
($\varprojlim A_\alpha$로 표기한다)와 각 $\alpha\in I$에 대한 사상
$u_\alpha:B\to A_\alpha$로 이루어지며, 다음 조건들을 만족한다~:
1\textsuperscript{o} $\alpha\leq\beta$일 때
$u_\alpha=u_{\alpha\beta}u_\beta$이다~;
2\textsuperscript{o} $C$의 임의의 대상 $X$와, $\alpha\leq\beta$일 때
$v_\alpha=u_{\alpha\beta}v_\beta$를 만족하는 사상
$v_\alpha:X\to A_\alpha$들의 임의의 족 $(v_\alpha)_{\alpha\in I}$에
대하여, 모든 $\alpha\in I$에 대해 $v_\alpha=u_\alpha v$를 만족하는
유일한 사상 $v:X\to B$가 존재한다 ($\varprojlim v_\alpha$로 표기한다)
(T, 1.8). 이는 다음과 같이 해석된다~: $u_{\alpha\beta}$들은 표준적으로
함수
\[
  \overline{u}_{\alpha\beta}:\operatorname{Hom}(X,A_\beta)
  \to\operatorname{Hom}(X,A_\alpha)
\]
들을 정의하며, 이 함수들은 집합들의 \emph{사영계}
$(\operatorname{Hom}(X,A_\alpha),\overline{u}_{\alpha\beta})$를
정의한다. 그리고 $(v_\alpha)$는 정의에 따라 집합
$\varprojlim_\alpha\operatorname{Hom}(X,A_\alpha)$의 원소이다~;
$X\mapsto\varprojlim_\alpha\operatorname{Hom}(X,A_\alpha)$가 $C$에서
$\mathrm{Ens}$로 가는 \emph{반변 함자}임은 명백하며, 사영극한 $B$가
존재한다는 것은 $(v_\alpha)\mapsto\varprojlim v_\alpha$가 $X$에 관한
함자들의 \emph{동형사상}
\[
  \varprojlim_\alpha\operatorname{Hom}(X,A_\alpha)
  \xrightarrow{\sim}\operatorname{Hom}(X,B)
  \tag{8.1.9.1}
\]
이라는 것, 다시 말해 함자
$X\mapsto\varprojlim_\alpha\operatorname{Hom}(X,A_\alpha)$가
\emph{표현 가능하다}는 것과 동치이다.
\end{env}

\begin{env}[8.1.10]
\label{0.8.1.10-fr}
\textit{예. II: 끝 대상.} $C$를 범주라 하고, $\{a\}$를 한 원소로만
이루어진 집합이라 하자. 임의의 대상 $X$에 집합 $\{a\}$를 대응시키고,
$C$의 임의의 사상 $X\to X'$에 유일한 함수 $\{a\}\to\{a\}$를
대응시키는 반변 함자 $F:C\to\mathrm{Ens}$를 생각하자. 이 함자가
\emph{표현 가능하다}는 것은, 임의의 $Y\in C$에 대하여
$\operatorname{Hom}(Y,e)=h_e(Y)$가 \emph{한 원소로만 이루어지는} 대상
$e\in C$가 존재한다는 뜻이다~; 이러한 $e$를 $C$의 \emph{끝 대상}이라
한다. $C$의 두 끝 대상이 동형임은 명백하다 (이에 따라, 일반적으로
선택공리를 사용하여, $C$의 끝 대상을 정하고 이를 $e_C$로 나타낼 수
있다). 예를 들어 범주 $\mathrm{Ens}$에서 끝 대상은 한 원소로만
이루어진 집합이다~; 체 $K$ 위의 \emph{증대 대수}의 범주(사상은
증대와 양립하는 대수 준동형이다)에서는 $K$가 끝 대상이다~;
\emph{$S$-준스킴}의 범주 (I, 2.5.1)에서는 $S$가 끝 대상이다.
\end{env}

\begin{env}[8.1.11]
\label{0.8.1.11-fr}
범주 $C$의 두 대상 $X$, $Y$에 대하여
$h'_X(Y)=\operatorname{Hom}(X,Y)$로 놓고, 임의의 사상 $u:Y\to Y'$에
대하여 $h'_X(u)$를 $\operatorname{Hom}(X,Y)$에서
$\operatorname{Hom}(X,Y')$로 가는 함수 $v\mapsto uv$로 두자~;
그러면 $h'_X$는 \emph{공변 함자} $C\to\mathrm{Ens}$이다. 따라서
(\hyperref[0.8.1.2-fr]{8.1.2})에서와 같이 표준 공변 함자
$h':C^0\to\operatorname{Hom}(C,\mathrm{Ens})$의 정의를 얻는다~;
$C$에서 $\mathrm{Ens}$로 가는 \emph{공변 함자} $F$, 즉
$\operatorname{Hom}(C,\mathrm{Ens})$의 대상은, 어떤 대상 $X\in C$
(필연적으로 유일한 동형사상까지 유일하다)가 존재하여 $F$가
\emph{$h'_X$와 동형}이면 \emph{표현 가능하다}고 한다~; 이 개념에
대하여 앞의 논의와 <<~쌍대인~>> 논의를 전개하는 일은 독자에게
맡긴다. 이 개념은 \emph{귀납극한}의 개념, 특히 통상적인
<<~보편 문제의 해~>>라는 개념을 포괄한다.
\end{env}

\subsection{범주에서의 대수적 구조}
\label{subsection:0.8.2-ko}

\begin{env}[8.2.1]
\label{0.8.2.1-ko}
\oldpage[0\textsubscript{III}]{9}
$C$에서 $\mathrm{Ens}$로 가는 두 반변 함자 $F$, $F'$이 주어졌다고 하자.
모든 대상 $Y\in C$에 대하여
$(F\times F')(Y)=F(Y)\times F'(Y)$로 놓고, $C$의 모든 사상
$u:Y\to Y'$에 대하여 $(F\times F')(u)=F(u)\times F'(u)$로 놓음을
상기하자. 뒤의 사상은 $F(Y')\times F'(Y')$에서 $F(Y)\times F'(Y)$로
가는 사상
$(t,t')\mapsto((F(u))(t),(F'(u))(t'))$이다. 따라서
$F\times F':C\to\mathrm{Ens}$는 
\emph{반변 함자}이다(이는 사실 범주
$\operatorname{Hom}(C^0,\mathrm{Ens})$에서 대상 $F$, $F'$의
\emph{곱}과 다르지 않다). 대상 $X\in C$가 주어졌을 때, 함자적 사상
\[
  \gamma_X:h_X\times h_X\to h_X
  \tag{8.2.1.1}
\]
을 
\emph{$X$ 위의 내부 합성법칙}이라 하겠다.

다시 말해 (T, 1.2), 모든 대상 $Y\in C$에 대하여 $\gamma_X(Y)$는
사상 $h_X(Y)\times h_X(Y)\to h_X(Y)$이고(따라서 정의상 집합
$h_X(Y)$ 위의 \emph{내부 합성법칙}이다), $C$의 모든 사상
$u:Y\to Y'$에 대하여 도식
\[
  \xymatrix{
    h_X(Y')\times h_X(Y')
      \ar[rr]^{h_X(u)\times h_X(u)}\ar[dd]_{\gamma_X(Y')} &&
    h_X(Y)\times h_X(Y)\ar[dd]^{\gamma_X(Y)}\\\\
    h_X(Y')\ar[rr]_{h_X(u)} && h_X(Y)
  }
\]
이 가환이라는 조건을 만족한다. 이는 합성법칙 $\gamma_X(Y)$와
$\gamma_X(Y')$에 대하여 $h_X(u)$가 $h_X(Y')$에서 $h_X(Y)$로 가는
\emph{준동형}임을 뜻한다.

마찬가지로 $C$의 두 대상 $Z$, $X$가 주어졌을 때, 함자적 사상
\[
  \omega_{X,Z}:h_Z\times h_X\to h_X
  \tag{8.2.1.2}
\]
을 
\emph{$Z$를 작용자 영역으로 하는 $X$ 위의 외부 합성법칙}이라 한다.

위와 같이 모든 $Y\in C$에 대하여 $\omega_{X,Z}(Y)$는
$h_Z(Y)$를 작용자 영역으로 하는 $h_X(Y)$ 위의 외부 합성법칙이고,
모든 사상 $u:Y\to Y'$에 대하여 $h_X(u)$와 $h_Z(u)$는
$(h_Z(Y'),h_X(Y'))$에서 $(h_Z(Y),h_X(Y))$로 가는
\emph{디-준동형}을 이룸을 알 수 있다.
\end{env}

\begin{env}[8.2.2]
\label{0.8.2.2-ko}
$C$의 두 번째 대상 $X'$에 내부 합성법칙 $\gamma_{X'}$가 주어졌다고
하자. $C$의 사상 $w:X\to X'$가 모든 $Y\in C$에 대하여
$h_w(Y):h_X(Y)\to h_{X'}(Y)$가 합성법칙 $\gamma_X(Y)$와
$\gamma_{X'}(Y)$에 관한 준동형이면, $w$를 이 합성법칙들에 관한
\emph{준동형}이라 하겠다. $X''$이 내부 합성법칙 $\gamma_{X''}$를 갖춘
$C$의 세 번째
\oldpage[0\textsubscript{III}]{10}
대상이고, $w':X'\to X''$이 $\gamma_{X'}$와 $\gamma_{X''}$에 관한
준동형인 $C$의 사상이면, 사상 $w'w:X\to X''$이 합성법칙
$\gamma_X$와 $\gamma_{X''}$에 관한 준동형임은 명백하다. $C$의 동형
$w:X\xrightarrow{\sim}X'$가 이 합성법칙들에 관한 준동형이고 그 역 사상
$w^{-1}$가 합성법칙 $\gamma_{X'}$와 $\gamma_X$에 관한 준동형이면,
$w$를 \emph{합성법칙 $\gamma_X$와 $\gamma_{X'}$에 관한 동형}이라 한다.

외부 합성법칙을 갖춘 $C$의 대상쌍들 사이의 \emph{디-준동형}도 같은
방식으로 정의한다.
\end{env}

\begin{env}[8.2.3]
\label{0.8.2.3-ko}
대상 $X\in C$ 위의 내부 합성법칙 $\gamma_X$가 \emph{모든}
$Y\in C$에 대하여 $h_X(Y)$ 위의 \emph{군법칙} $\gamma_X(Y)$를 주면,
이 법칙을 갖춘 $X$를 \emph{$C$-군} 또는 \emph{범주 $C$의 군 대상}이라
한다. \emph{$C$-환}, \emph{$C$-가군} 등도 마찬가지로 정의한다.
\end{env}

\begin{env}[8.2.4]
\label{0.8.2.4-ko}
대상 $X\in C$와 자기 자신의 \emph{곱} $X\times X$가 $C$에서
존재한다고 하자. 이는 사영극한의 특수한 경우이므로
(\hyperref[0.8.1.9-fr]{8.1.9}), 정의에 따라 표준 동형을 제외하고
$h_{X\times X}=h_X\times h_X$이다. 따라서 $X$ 위의 내부 합성법칙은
함자적 사상 $\gamma_X:h_{X\times X}\to h_X$로 볼 수 있고, 이에 따라
(\hyperref[0.8.1.6-fr]{8.1.6}) $h_{c_X}=\gamma_X$인 원소
$c_X\in\operatorname{Hom}(X\times X,X)$를 표준적으로 정한다. 그러므로
이 경우 $X$ 위의 내부 합성법칙을 주는 것은 사상 $X\times X\to X$를
주는 것과 동치이다. $C$가 범주 $\mathrm{Ens}$일 때에는 집합 위의 내부
합성법칙이라는 고전적 개념을 다시 얻는다. 곱 $Z\times X$가 $C$에서
존재하는 경우의 외부 합성법칙에 대해서도 같은 결과가 성립한다.
\end{env}

\begin{env}[8.2.5]
\label{0.8.2.5-ko}
앞의 표기 아래에서 $X\times X\times X$도 $C$에서 존재한다고 하자.
함자를 표현하는 대상으로서 곱의 특징짓기
(\hyperref[0.8.1.9-fr]{8.1.9})로부터 표준 동형
\[
  (X\times X)\times X\xrightarrow{\sim}X\times X\times X
  \xrightarrow{\sim}X\times(X\times X)
\]
들이 존재한다. $X\times X\times X$를 $(X\times X)\times X$와
표준적으로 동일시하면, 모든 $Y\in C$에 대하여 사상
$\gamma_X(Y)\times 1_{h_X(Y)}$는 $h_{c_X\times 1_X}(Y)$와 동일시된다.
따라서 모든 $Y\in C$에 대하여 내부 합성법칙 $\gamma_X(Y)$가
\emph{결합적}이라고 하는 것과 사상들의 도식
\[
  \xymatrix{
    h_X(Y)\times h_X(Y)\times h_X(Y)
      \ar[rr]^{\gamma_X(Y)\times 1}\ar[dd]_{1\times\gamma_X(Y)} &&
    h_X(Y)\times h_X(Y)\ar[dd]^{\gamma_X(Y)}\\\\
    h_X(Y)\times h_X(Y)\ar[rr]_{\gamma_X(Y)} && h_X(Y)
  }
\]
이 가환이라고 하는 것, 그리고 사상들의 도식
\oldpage[0\textsubscript{III}]{11}
\[
  \xymatrix{
    X\times X\times X
      \ar[rr]^{c_X\times 1_X}\ar[dd]_{1_X\times c_X} &&
    X\times X\ar[dd]^{c_X}\\\\
    X\times X\ar[rr]_{c_X} && X
  }
\]
이 가환이라고 하는 것은 서로 동치이다.
\end{env}

\begin{env}[8.2.6]
\label{0.8.2.6-ko}
(\hyperref[0.8.2.5-ko]{8.2.5})의 가정 아래에서 모든 $Y\in C$에 대하여
내부 합성법칙 $\gamma_X(Y)$가 \emph{군법칙}임을 나타내려면, 한편으로
그것이 결합적임을 나타내고 다른 한편으로 군에서의 \emph{역원}의 성질을
갖는 사상 $\alpha_X(Y):h_X(Y)\to h_X(Y)$가 존재함을 나타내야 한다.
$C$의 모든 사상 $u:Y\to Y'$에 대하여 $h_X(u)$가 군 $h_X(Y')$에서
$h_X(Y)$로 가는 준동형이어야 함을 앞에서 보았으므로, 먼저
$\alpha_X:h_X\to h_X$는 \emph{함자적 사상}이어야 한다. 한편 군 $G$에서
역원을 취하는 사상 $s\mapsto s^{-1}$의 특징적인 성질은 항등원을 쓰지
않고도 나타낼 수 있다. 즉 두 합성사상
\[
  (s,t)\mapsto(s,s^{-1},t)\mapsto(s,s^{-1}t)\mapsto s(s^{-1}t),
\]
\[
  (s,t)\mapsto(s,s^{-1},t)\mapsto(s,ts^{-1})\mapsto(ts^{-1})s
\]
이 $G\times G$에서 $G$로 가는 두 번째 사영 $(s,t)\mapsto t$와 같다고
쓰면 충분하다. (\hyperref[0.8.1.3-fr]{8.1.3})에 따라
$a_X\in\operatorname{Hom}(X,X)$인 $\alpha_X=h_{a_X}$이다. 그러면 앞의
첫째 조건은 합성사상
\[
  X\times X\xrightarrow{(1_X,a_X)\times 1_X}X\times X\times X
  \xrightarrow{1_X\times c_X}X\times X\xrightarrow{c_X}X
\]
이 $C$에서의 두 번째 사영 $X\times X\to X$라는 것을 나타내며,
둘째 조건도 같은 방식으로 옮겨진다.
\end{env}

\begin{env}[8.2.7]
\label{0.8.2.7-ko}
이제 $C$에 \emph{끝 대상} $e$가 존재한다고 하자
(\hyperref[0.8.1.10-fr]{8.1.10}). 여전히 모든 $Y\in C$에 대하여
$\gamma_X(Y)$가 $h_X(Y)$ 위의 군법칙이라고 하고, $\eta_X(Y)$로
$\gamma_X(Y)$의 항등원을 나타내자. $C$의 모든 사상 $u:Y\to Y'$에
대하여 $h_X(u)$가 군 준동형이므로
$\eta_X(Y)=(h_X(u))(\eta_X(Y'))$이다. 특히 $Y'=e$로 취하면 $u$는
$\operatorname{Hom}(Y,e)$의 유일한 원소 $\varepsilon$이므로,
$\eta_X(e)$가 모든 $Y\in C$에 대하여 $\eta_X(Y)$를 완전히 정함을 알 수
있다. $h_X(X)=\operatorname{Hom}(X,X)$인 군의 항등원
$e_X=\eta_X(X)$로 놓자. 도식
\[
  \xymatrix{
    h_X(e)\ar[r]^{h_X(\varepsilon)}\ar[d]_{h_{e_X}(e)} &
    h_X(Y)\ar[d]^{h_{e_X}(Y)}\\
    h_X(e)\ar[r]_{h_X(\varepsilon)} & h_X(Y)
  }
\]
의 가환성((\hyperref[0.8.1.2-fr]{8.1.2}) 참조)은 집합 $h_X(Y)$에서
사상 $h_{e_X}(Y)$가 모든 원소를 항등원으로 보내는
\oldpage[0\textsubscript{III}]{12}
$s\mapsto\eta_X(Y)$와 다르지 않음을 보인다. 그러면 모든 $Y\in C$에
대하여 $\eta_X(Y)$가 $\gamma_X(Y)$의 항등원이라는 것은 합성사상
\[
  X\xrightarrow{(1_X,1_X)}X\times X
  \xrightarrow{1_X\times e_X}X\times X\xrightarrow{c_X}X,
\]
그리고 $1_X$와 $e_X$의 자리를 바꾼 유사한 사상이 둘 다 $1_X$와
\emph{같다}고 하는 것과 동치임을 확인할 수 있다.
\end{env}

\begin{env}[8.2.8]
\label{0.8.2.8-ko}
물론 범주에서의 대수적 구조에 관한 예는 어렵지 않게 얼마든지 더 들 수
있다. 군의 예는 충분히 상세히 다루었지만, 앞으로 만나게 될 대수적 구조의
예에서는 이와 비슷한 논의를 전개하는 일을 대체로 독자에게 맡기겠다.
\end{env}

\section{구성가능 집합}
\label{section:0.9-ko}

\subsection{구성가능 집합}
\label{subsection:0.9.1-ko}

\begin{definition}[9.1.1]
\label{0.9.1.1-ko}
연속사상 $f:X\to Y$가 $Y$의 모든 준콤팩트 열린집합 $U$에 대하여
$f^{-1}(U)$가 준콤팩트이면, $f$를 \emph{준콤팩트}라 한다. 위상공간
$X$의 부분집합 $Z$에 대한 표준 단사사상 $Z\to X$가 준콤팩트이면,
$Z$를 \emph{$X$에서 역콤팩트}라 한다. 다시 말해 $X$의 모든 준콤팩트
열린집합 $U$에 대하여 $U\cap Z$가 준콤팩트인 경우이다.
\end{definition}

$X$의 \emph{닫힌} 부분집합은 $X$에서 역콤팩트이지만, $X$의 준콤팩트
부분집합은 반드시 $X$에서 역콤팩트인 것은 아니다. $X$가 준콤팩트이면
$X$에서 역콤팩트인 모든 열린 부분집합은 준콤팩트이다. $X$에서
역콤팩트인 집합들의 모든 \emph{유한} 합집합은 $X$에서 역콤팩트임이
명백하다. 이는 준콤팩트 집합들의 모든 유한 합집합이 준콤팩트이기
때문이다. $X$에서 역콤팩트인 열린집합들의 모든 유한 교집합은
$X$에서 역콤팩트인 열린집합이다. \emph{국소 뇌터} 공간 $X$에서는
모든 준콤팩트 집합이 뇌터 부분공간이고, 따라서 $X$의 \emph{모든
부분집합}은 $X$에서 역콤팩트이다.

\begin{definition}[9.1.2]
\label{0.9.1.2-ko}
위상공간 $X$가 주어졌다고 하자. $X$에서 역콤팩트인 모든 열린
부분집합을 포함하고 유한 교집합과 여집합을 취하는 것에 관하여 닫혀
있는 $X$의 부분집합들의 최소 집합 $\mathfrak{F}$에 속하는 $X$의
부분집합을 \emph{구성가능}이라 한다(따라서 $\mathfrak{F}$는 유한
합집합에 관해서도 닫혀 있다).
\end{definition}

\begin{proposition}[9.1.3]
\label{0.9.1.3-ko}
$X$의 부분집합이 구성가능이기 위한 필요충분조건은 $X$에서 역콤팩트인
열린집합 $U$, $V$에 대한 $U\cap\complement{V}$ 꼴의 집합들의 유한
합집합인 것이다.
\end{proposition}

이 조건이 충분함은 명백하다. 필요함을 보이기 위하여, $X$에서
역콤팩트인 열린집합 $U$, $V$에 대한 $U\cap\complement{V}$ 꼴의
집합들의 유한 합집합으로 이루어진 집합 $\mathfrak{G}$를 생각하자.
$\mathfrak{G}$의 모든 원소의 여집합이 $\mathfrak{G}$에 속함을 보이면
충분하다. 따라서 $I$가 유한하고 $U_i$, $V_i$가 $X$에서 역콤팩트인
열린집합인
$Z=\bigcup_{i\in I}(U_i\cap\complement{V_i})$를 취하자. 이때
$\complement{Z}=\bigcap_{i\in I}(V_i\cup\complement{U_i})$이므로,
$\complement{Z}$는 몇 개의 $V_i$와 몇 개의 $\complement{U_i}$를
교차하여 얻는 집합들의 유한 합집합이다. 따라서 각 집합은
\oldpage[0\textsubscript{III}]{13}
$V\cap\complement{U}$ 꼴이며, 여기서 $U$는 몇 개의 $U_i$의
합집합이고 $V$는 몇 개의 $V_i$의 교집합이다. 그런데 $X$에서
역콤팩트인 열린집합들의 유한 합집합과 유한 교집합은 역콤팩트인
열린집합임을 앞에서 보았으므로 결론이 따른다.

\begin{corollary}[9.1.4]
\label{0.9.1.4-ko}
$X$의 모든 구성가능 부분집합은 $X$에서 역콤팩트이다.
\end{corollary}

$U$, $V$가 $X$에서 역콤팩트인 열린집합일 때
$U\cap\complement{V}$가 $X$에서 역콤팩트임을 보이면 충분하다.
그런데 $W$가 $X$의 준콤팩트 열린집합이면
$W\cap U\cap\complement{V}$는 준콤팩트 공간 $W\cap U$의 닫힌
부분집합이므로 준콤팩트이다.

특히 다음을 얻는다.

\begin{corollary}[9.1.5]
\label{0.9.1.5-ko}
$X$의 열린 부분집합 $U$가 구성가능이기 위한 필요충분조건은 $X$에서
역콤팩트인 것이다. $X$의 닫힌 부분집합 $F$가 구성가능이기 위한
필요충분조건은 열린집합 $\complement{F}$가 역콤팩트인 것이다.
\end{corollary}

\begin{env}[9.1.6]
\label{0.9.1.6-ko}
$X$의 모든 준콤팩트 열린 부분집합이 역콤팩트인 경우, 다시 말해 $X$의
두 준콤팩트 열린 부분집합의 교집합이 준콤팩트인 경우가 중요하다
((\hyperref[I.5.5.6-ko]{I, 5.5.6}) 참조). $X$ 자체가 준콤팩트이면,
이는 $X$에서 역콤팩트인 열린 부분집합들이 $X$의 준콤팩트 열린
부분집합들과 같고, $X$의 구성가능 부분집합들이 준콤팩트 열린집합
$U$, $V$에 대한 $U\cap\complement{V}$ 꼴의 집합들의 유한 합집합과
같다는 뜻이다.
\end{env}

\begin{corollary}[9.1.7]
\label{0.9.1.7-ko}
뇌터 공간 $X$의 부분집합이 구성가능이기 위한 필요충분조건은 $X$의
국소 닫힌 부분집합들의 유한 합집합인 것이다.
\end{corollary}

\begin{proposition}[9.1.8]
\label{0.9.1.8-ko}
$X$를 위상공간, $U$를 $X$의 열린 부분집합이라 하자.
\begin{enumerate}
  \item[(i)] $T$가 $X$의 구성가능 부분집합이면 $T\cap U$는 $U$의
    구성가능 부분집합이다.
  \item[(ii)] 더 나아가 $U$가 $X$에서 역콤팩트라고 하자. $U$의
    부분집합 $Z$가 $X$에서 구성가능이기 위한 필요충분조건은 $U$에서
    구성가능인 것이다.
\end{enumerate}
\end{proposition}

\begin{enumerate}
  \item[(i)] (\hyperref[0.9.1.3-ko]{9.1.3})을 사용하면, $T$가 $X$에서
    역콤팩트인 열린집합일 때 $T\cap U$가 $U$에서 역콤팩트인
    열린집합임을 보이는 것으로 환원된다. 다시 말해 모든 준콤팩트
    열린집합 $W\subset U$에 대하여 $T\cap U\cap W=T\cap W$가
    준콤팩트임을 보이면 되는데, 이는 가정에서 곧바로 따른다.
  \item[(ii)] (i)에 따라 조건은 필요하므로, 충분함만 보이면 된다.
    (\hyperref[0.9.1.3-ko]{9.1.3})을 고려하면 $Z$가 \emph{$U$에서}
    역콤팩트인 열린집합인 경우만 생각하면 충분하다. 실제로 그러면
    $U\setminus Z$가 $X$에서 구성가능이고, $Z$, $Z'$이 \emph{$U$에서}
    역콤팩트인 두 열린집합이면 $Z\cap(U\setminus Z')$가 참으로
    $X$에서 구성가능임이 따른다. 이제 $W$가 $X$의 준콤팩트
    열린집합이고 $Z$가 $U$에서 역콤팩트인 열린집합이면
    $Z\cap W=Z\cap(W\cap U)$이고, 가정에 따라 $W\cap U$는 $U$의
    준콤팩트 열린집합이다. 따라서 $W\cap Z$는 준콤팩트이고, 이에 따라
    $Z$는 $X$에서 역콤팩트인 열린집합이며 \emph{더구나} $X$에서
    구성가능이다.
\end{enumerate}

\begin{corollary}[9.1.9]
\label{0.9.1.9-ko}
$X$를 위상공간, $(U_i)_{i\in I}$를 $X$에서 역콤팩트인 열린집합들로
이루어진 $X$의 유한 덮개라 하자. $X$의 부분집합 $Z$가 $X$에서
구성가능이기 위한 필요충분조건은 모든 $i\in I$에 대하여 $Z\cap U_i$가
$U_i$에서 구성가능인 것이다.
\end{corollary}

\begin{env}[9.1.10]
\label{0.9.1.10-ko}
특히 $X$가 준콤팩트이고 $X$의 모든 점이
\oldpage[0\textsubscript{III}]{14}
$X$에서 역콤팩트인 열린
근방들(따라서 \emph{더구나} 준콤팩트 열린 근방들)로 이루어진 기본계를
갖는다고 하자. 그러면 $X$의 부분집합 $Z$가 $X$에서 구성가능이라는
조건은 \emph{국소적}이다. 다시 말해 모든 $x\in X$에 대하여 $x$의
열린 근방 $V$가 존재하여 $V\cap Z$가 $V$에서 구성가능인 것이
필요충분하다. 실제로 이 조건이 성립하면, 모든 $x\in X$에 대하여
$X$에서 \emph{역콤팩트}이고 $V\cap Z$가 $V$에서 구성가능인 열린 근방
$V$가 존재한다. 이는 $X$에 관한 가정과
(\hyperref[0.9.1.8-ko]{9.1.8, (i)})에서 따른다. 이제 이 근방들의 유한
개로 $X$를 덮고 (\hyperref[0.9.1.9-ko]{9.1.9})를 적용하면 충분하다.
\end{env}

\begin{definition}[9.1.11]
\label{0.9.1.11-ko}
$X$를 위상공간이라 하자. 모든 $x\in X$에 대하여 $x$의 열린 근방 $V$가
존재하여 $T\cap V$가 $V$에서 구성가능이면, $X$의 부분집합 $T$를
\emph{$X$에서 국소 구성가능}이라 한다.
\end{definition}

$V\cap T$가 $V$에서 구성가능이면 모든 열린집합 $W\subset V$에 대하여
$W\cap T$가 $W$에서 구성가능임이
(\hyperref[0.9.1.8-ko]{9.1.8, (i)})에서 곧바로 따른다. $T$가 $X$에서
국소 구성가능이면, 앞의 주석에 따라 $X$의 모든 열린집합 $U$에 대하여
$T\cap U$는 $U$에서 국소 구성가능이다. 같은 주석으로부터 $X$에서
국소 구성가능인 부분집합들의 집합은 유한 합집합과 유한 교집합에
관하여 닫혀 있음을 알 수 있다. 한편 이 집합이 여집합을 취하는 것에
관해서도 닫혀 있음은 명백하다.

\begin{proposition}[9.1.12]
\label{0.9.1.12-ko}
$X$를 위상공간이라 하자. $X$에서 구성가능인 모든 집합은 $X$에서
국소 구성가능이다. $X$가 준콤팩트이고 $X$의 위상이 $X$에서 역콤팩트인
집합들로 이루어진 기저를 가지면 그 역도 성립한다.
\end{proposition}

첫째 주장은 정의 (\hyperref[0.9.1.11-ko]{9.1.11})에서 따르고, 둘째
주장은 (\hyperref[0.9.1.10-ko]{9.1.10})에서 따른다.

\begin{corollary}[9.1.13]
\label{0.9.1.13-ko}
$X$를 그 위상이 $X$에서 역콤팩트인 집합들로 이루어진 기저를 갖는
위상공간이라 하자. $X$에서 국소 구성가능인 모든 부분집합 $T$는
$X$에서 역콤팩트이다.
\end{corollary}

실제로 $U$를 $X$의 준콤팩트 열린집합이라 하자. $T\cap U$는 $U$에서
국소 구성가능이므로 (\hyperref[0.9.1.12-ko]{9.1.12})에 따라 $U$에서
구성가능이고, 따라서 (\hyperref[0.9.1.4-ko]{9.1.4})에 따라
준콤팩트이다.

\subsection{뇌터 공간에서의 구성가능 집합}
\label{subsection:0.9.2-ko}

\begin{env}[9.2.1]
\label{0.9.2.1-ko}
뇌터 공간 $X$에서 $X$의 구성가능 부분집합들은 \emph{$X$의 국소 닫힌
부분집합들의 유한 합집합}임을
(\hyperref[0.9.1.7-ko]{9.1.7})에서 보았다.

뇌터 공간 $X'$에서 $X$로 가는 연속사상에 의한 $X$의 구성가능 집합의
역상은 $X'$에서 구성가능이다. $Y$가 뇌터 공간 $X$의 구성가능
부분집합이면, $Y$의 부분공간으로서 구성가능인 $Y$의 부분집합들과
$X$의 부분공간으로서 구성가능인 것들은 서로 같다.
\end{env}

\begin{proposition}[9.2.2]
\label{0.9.2.2-ko}
$X$를 뇌터 기약공간, $E$를 $X$의 구성가능 부분집합이라 하자. $E$가
$X$에서 모든 곳에서 조밀하기 위한 필요충분조건은 $E$가 $X$의 공집합이
아닌 열린 부분집합 하나를 포함하는 것이다.
\end{proposition}

\oldpage[0\textsubscript{III}]{15}
공집합이 아닌 모든 열린집합은 $X$에서 조밀하므로 이 조건이 충분함은
명백하다. 역으로 $E=\bigcup_{i=1}^n(U_i\cap F_i)$를 $X$의 구성가능
부분집합이라 하자. 여기서 $U_i$들은 공집합이 아닌 열린집합이고
$F_i$들은 $X$의 닫힌집합이다. 그러면
$\overline{E}\subset\bigcup_i F_i$이다. 따라서 $\overline{E}=X$이면
$X$는 $F_i$ 중 하나와 같고, 이에 따라 $E\supset U_i$이다. 이로써
증명이 끝난다.

$X$가 일반점 $x$를 가지면
(\hyperref[0.2.1.2-ko]{$0_{\mathrm I}$, 2.1.2}),
(\hyperref[0.9.2.2-ko]{9.2.2})의 조건은 관계 $x\in E$와 동치이다.

\begin{proposition}[9.2.3]
\label{0.9.2.3-ko}
$X$를 뇌터 공간이라 하자. $X$의 부분집합 $E$가 구성가능이기 위한
필요충분조건은 $X$의 모든 닫힌 기약 부분집합 $Y$에 대하여 $E\cap Y$가
$Y$에서 어디에서도 조밀하지 않거나 $Y$의 공집합이 아닌 열린
부분집합 하나를 포함하는 것이다.
\end{proposition}

조건의 필요성은 $E\cap Y$가 $Y$의 구성가능 부분집합이어야 한다는
사실과 (\hyperref[0.9.2.2-ko]{9.2.2})에서 따른다. 실제로 $Y$에서 조밀하지
않은 부분집합은 기약공간 $Y$에서 반드시 어디에서도 조밀하지 않다
(\hyperref[0.2.1.1-ko]{$0_{\mathrm I}$, 2.1.1}). 조건이 충분함을 보이기
위하여, $Y\cap E$가 구성가능인(여기서 $Y$에 관하여 구성가능이든
$X$에 관하여 구성가능이든 같은 뜻이다) $X$의 닫힌 부분집합 $Y$들의
집합 $\mathfrak{F}$에 뇌터 귀납법의 원리
(\hyperref[0.2.2.2-ko]{$0_{\mathrm I}$, 2.2.2})를 적용하자. 따라서
$X$의 모든 닫힌 부분집합 $Y\ne X$에 대하여 $E\cap Y$가 구성가능이라고
가정할 수 있다. 먼저 $X$가 기약이 아니라고 하고, 그 기약 성분들을
$X_i$ ($1\leq i\leq m$)라 하자. 이들은 반드시 유한 개이다
(\hyperref[0.2.2.5-ko]{$0_{\mathrm I}$, 2.2.5}). 가정에 따라
$E\cap X_i$들은 구성가능이고, 따라서 그 합집합 $E$도 구성가능이다.
이제 $X$가 기약이라고 하자. 이때 가정에 따라 $E$가 어디에서도
조밀하지 않아서 $\overline{E}\ne X$이고
$E=E\cap\overline{E}$가 구성가능이거나, 또는 $E$가 $X$의 공집합이
아닌 열린집합 $U$를 포함하여 $U$와 $E\cap(X\setminus U)$의
합집합이다. 그런데 $X\setminus U$는 $X$와 다른 닫힌집합이므로
$E\cap(X\setminus U)$는 구성가능이다. 따라서 $E$ 자체도 구성가능이고,
이로써 증명이 끝난다.

\begin{corollary}[9.2.4]
\label{0.9.2.4-ko}
$X$를 뇌터 공간, $(E_\alpha)$를 다음 조건을 만족하는 $X$의 구성가능
부분집합들의 증가 여과족이라 하자.
\begin{enumerate}
  \item[$1^\circ$] $X$는 족 $(E_\alpha)$의 합집합이다.
  \item[$2^\circ$] $X$의 모든 닫힌 기약 부분집합은 어떤 $E_\alpha$의
    폐포에 포함된다.
\end{enumerate}
그러면 $X=E_\alpha$인 지표 $\alpha$가 존재한다.

$X$의 모든 닫힌 기약 부분집합이 일반점을 가지면 가정 $2^\circ$를
생략할 수 있다.
\end{corollary}

$X$의 닫힌 부분집합들 가운데 적어도 하나의 $E_\alpha$에 포함되는
것들의 집합 $\mathfrak{M}$에 뇌터 귀납법의 원리
(\hyperref[0.2.2.2-ko]{$0_{\mathrm I}$, 2.2.2})를 적용하자. 따라서
$X$의 모든 닫힌 부분집합 $Y\ne X$가 어떤 $E_\alpha$에 포함된다고
가정할 수 있다. $X$가 기약이 아니면 명제는 명백하다. 실제로 $X$의
각 기약 성분 $X_i$ ($1\leq i\leq m$)는 어떤 $E_{\alpha_i}$에 포함되고,
모든 $E_{\alpha_i}$를 포함하는 $E_\alpha$가 존재한다. 따라서 $X$가
기약인 경우를 생각하자. 가정에 따라 $X=\overline{E_\beta}$인 $\beta$가
존재하고, 따라서 (\hyperref[0.9.2.2-ko]{9.2.2})에 따라 $E_\beta$는
$X$의 공집합이 아닌 열린집합 $U$를 포함한다. 그런데 닫힌집합
$X\setminus U$는 어떤 $E_\gamma$에 포함되므로, $E_\beta$와
$E_\gamma$를 포함하는 $E_\alpha$를 취하면 충분하다. $X$의 모든 닫힌
기약 부분집합 $Y$가
\oldpage[0\textsubscript{III}]{16}
일반점 $y$를 가지면 $y\in E_\alpha$인 $\alpha$가 존재한다. 따라서
$Y=\overline{\{y\}}\subset\overline{E_\alpha}$이고, 조건 $2^\circ$는
$1^\circ$의 귀결이다.

\begin{proposition}[9.2.5]
\label{0.9.2.5-ko}
$X$를 뇌터 공간, $x$를 $X$의 점, $E$를 $X$의 구성가능 부분집합이라
하자. $E$가 $x$의 근방이기 위한 필요충분조건은 $x$를 포함하는 $X$의
모든 닫힌 기약 부분집합 $Y$에 대하여 $E\cap Y$가 $Y$에서 조밀한
것이다($Y$에 일반점 $y$가 존재하면, 이는
(\hyperref[0.9.2.2-ko]{9.2.2})에 따라 $y\in E$라는 뜻이기도 하다).
\end{proposition}

조건이 필요함은 명백하다. 충분함을 보이자. $x$를 포함하며 $E\cap Y$가
$Y$에서 $x$의 근방인 $X$의 닫힌 부분집합 $Y$들의 집합
$\mathfrak{M}$에 뇌터 귀납법의 원리를 적용하면, $x$를 포함하는
$X$의 모든 닫힌 부분집합 $Y\ne X$가 $\mathfrak{M}$에 속한다고 가정할
수 있다. $X$가 기약이 아니면 $x$를 포함하는 $X$의 각 기약 성분
$X_i$는 $X$와 다르므로 $E\cap X_i$는 $X_i$에서 $x$의 근방이다.
따라서 $E$는 $x$를 포함하는 $X$의 기약 성분들의 합집합에서 $x$의
근방이고, 이 합집합 자체가 $X$에서 $x$의 근방이므로 $E$도 그러하다.
$X$가 기약이면 가정에 따라 $E$는 $X$에서 조밀하고, 따라서
(\hyperref[0.9.2.2-ko]{9.2.2})에 따라 $X$의 공집합이 아닌 열린
부분집합 $U$를 포함한다. $x\in U$이면 명제는 명백하다. 그렇지 않으면
가정에 따라 $x$는 $X\setminus U$에 관한 $E\cap(X\setminus U)$의
내점이다. 따라서 $X$에서 $X\setminus E$의 폐포는 $x$를 포함하지 않고,
이 폐포의 여집합은 $E$에 포함되는 $x$의 근방이다. 이로써 증명이
끝난다.

\begin{corollary}[9.2.6]
\label{0.9.2.6-ko}
$X$를 뇌터 공간, $E$를 $X$의 부분집합이라 하자. $E$가 $X$의
열린집합이기 위한 필요충분조건은 $E$와 만나는 $X$의 모든 닫힌 기약
부분집합 $Y$에 대하여 $E\cap Y$가 $Y$의 공집합이 아닌 열린 부분집합
하나를 포함하는 것이다.
\end{corollary}

조건이 필요함은 명백하다. 역으로 조건이 성립하면
(\hyperref[0.9.2.3-ko]{9.2.3})에 따라 $E$는 구성가능이다. 더 나아가
(\hyperref[0.9.2.5-ko]{9.2.5})에 따라 $E$는 자신의 모든 점의 근방이고,
따라서 결론이 따른다.

\subsection{구성가능 함수}
\label{subsection:0.9.3-ko}

\begin{definition}[9.3.1]
\label{0.9.3.1-ko}
$h$를 위상공간 $X$에서 집합 $T$로 가는 사상이라 하자. 모든 $t\in T$에
대하여 $h^{-1}(t)$가 구성가능이고 유한 개의 $t$를 제외하면 공집합이면,
$h$를 \emph{구성가능}이라 한다. 이때 $T$의 모든 부분집합 $S$에 대하여
$h^{-1}(S)$는 구성가능이다. 모든 $x\in X$가 $h|V$가 구성가능인 열린
근방 $V$를 가지면 $h$를 \emph{국소 구성가능}이라 한다.
\end{definition}

모든 구성가능 함수는 국소 구성가능이다. $X$가 준콤팩트이고 $X$에서
역콤팩트인 열린집합들로 이루어진 기저를 가지면(특히 $X$가 뇌터이면)
그 역도 성립한다.

\begin{proposition}[9.3.2]
\label{0.9.3.2-ko}
$h$를 뇌터 공간 $X$에서 집합 $T$로 가는 사상이라 하자. $h$가
구성가능이기 위한 필요충분조건은 $X$의 모든 닫힌 기약 부분집합 $Y$에
대하여 $Y$에서 상대적으로 열려 있는 공집합이 아닌 부분집합 $U$가
존재하여 $U$에서 $h$가 상수인 것이다.
\end{proposition}

조건은 필요하다. 실제로 가정에 따라 $h$가 $Y$에서 취하는 값들은
유한 개의 $t_i$뿐이고, 각 집합 $h^{-1}(t_i)\cap Y$는 $Y$에서
구성가능이다(\hyperref[0.9.2.1-ko]{9.2.1}). 이 집합들이 모두 $Y$에서
어디에서도 조밀하지 않을 수는 없으므로, 그중 적어도 하나는 공집합이
아닌 열린집합을 포함한다
(\hyperref[0.9.2.3-ko]{9.2.3}).

\oldpage[0\textsubscript{III}]{17}
조건이 충분함을 보이기 위하여, $h|Y$가 구성가능인 $X$의 닫힌 부분집합
$Y$들의 집합 $\mathfrak{M}$에 뇌터 귀납법의 원리를 적용하자. 따라서
$X$의 모든 닫힌 부분집합 $Y\ne X$에 대하여 $h|Y$가 구성가능이라고
가정할 수 있다. $X$가 기약이 아니면 $X$의 각 기약 성분 $X_i$로 제한한
$h$는 구성가능이다. 이 성분들은 유한 개이므로 정의
(\hyperref[0.9.3.1-ko]{9.3.1})에서 $h$가 구성가능임이 곧바로 따른다.
$X$가 기약이면 가정에 따라 $h$가 상수인 $X$의 공집합이 아닌 열린
부분집합 $U$가 존재한다. 한편 가정에 따라 $X\setminus U$로 제한한
$h$는 구성가능이므로, 이로부터 $h$가 구성가능임이 곧바로 따른다.

\begin{corollary}[9.3.3]
\label{0.9.3.3-ko}
$X$를 모든 닫힌 기약 부분집합이 일반점을 갖는 뇌터 공간이라 하자.
$h$가 $X$에서 집합 $T$로 가는 사상이고 모든 $t\in T$에 대하여
$h^{-1}(t)$가 구성가능이면, $h$는 구성가능이다.
\end{corollary}

실제로 $Y$가 $X$의 닫힌 기약 부분집합이고 $y$가 그 일반점이면,
$Y\cap h^{-1}(h(y))$는 구성가능이고 $y$를 포함한다. 따라서
(\hyperref[0.9.2.2-ko]{9.2.2})에 따라 이 집합은 $Y$의 공집합이 아닌
열린 부분집합 하나를 포함하며,
(\hyperref[0.9.3.2-ko]{9.3.2})를 적용하면 충분하다.

\begin{proposition}[9.3.4]
\label{0.9.3.4-ko}
$X$를 모든 닫힌 기약 부분집합이 일반점을 갖는 뇌터 공간, $h$를
$X$에서 순서집합으로 가는 구성가능 사상이라 하자. $h$가 $X$에서
상반연속이기 위한 필요충분조건은 모든 $x\in X$와 $x$의 모든 일반화
$x'$에 대하여
(\hyperref[0.2.1.2-ko]{$0_{\mathrm I}$, 2.1.2})
$h(x')\leq h(x)$인 것이다.
\end{proposition}

함수 $h$는 유한 개의 값만 취한다. 따라서 $h$가 상반연속이라는 것은
모든 $x\in X$에 대하여 $h(y)\leq h(x)$인 $y\in X$들의 집합 $E$가
$x$의 근방이라는 뜻이다. 가정에 따라 $E$는 $X$의 구성가능
부분집합이다. 한편 $X$의 닫힌 기약 부분집합 $Y$가 $x$를 포함한다는
것은 그 일반점 $y$가 $x$의 일반화라는 뜻이다. 따라서 결론은
(\hyperref[0.9.2.5-ko]{9.2.5})에서 따른다.

\section{평탄 가군에 관한 보충}
\label{section:0.10-fr}

(\hyperref[subsection:0.10.1-fr]{10.1})과
(\hyperref[subsection:0.10.2-fr]{10.2})에서 증명 없이 서술한 성질들의
증명은 부르바키, \emph{Alg. comm.}, 제II장과 제III장을 참조한다.

\subsection{평탄 가군과 자유 가군 사이의 관계}
\label{subsection:0.10.1-fr}

\begin{env}[10.1.1]
\label{0.10.1.1-fr}
$A$를 환, $\mathfrak{J}$를 $A$의 아이디얼, $M$을 $A$-가군이라 하자. 모든
정수 $p\geq0$에 대하여 다음과 같은 표준 $(A/\mathfrak{J})$-가군
준동형이 있다.
\[
\label{0.10.1.1.1-fr}
  \varphi_p:(M/\mathfrak{J}M)\otimes_{A/\mathfrak{J}}
  (\mathfrak{J}^p/\mathfrak{J}^{p+1})
  \longrightarrow \mathfrak{J}^pM/\mathfrak{J}^{p+1}M
  \tag{10.1.1.1}
\]
이 준동형은 명백히 \emph{전사}이다. $\mathfrak{J}^p$들로 여과된 $A$에
연관된 등급환을
$\operatorname{gr}(A)=\bigoplus_{p\geq0}\mathfrak{J}^p/\mathfrak{J}^{p+1}$로,
$\mathfrak{J}^pM$들로 여과된 $M$에 연관된 등급 $\operatorname{gr}(A)$-가군을
$\operatorname{gr}(M)=\bigoplus_{p\geq0}\mathfrak{J}^pM/\mathfrak{J}^{p+1}M$로
나타내자. 그러면
$\operatorname{gr}_p(A)=\mathfrak{J}^p/\mathfrak{J}^{p+1}$,
$\operatorname{gr}_p(M)=\mathfrak{J}^pM/\mathfrak{J}^{p+1}M$이고,
$\varphi_p$들은 다음의 \emph{전사}인 등급 $\operatorname{gr}(A)$-가군
준동형을 정의한다.
\[
\label{0.10.1.1.2-fr}
  \varphi:\operatorname{gr}_0(M)\otimes_{\operatorname{gr}_0(A)}
  \operatorname{gr}(A)\longrightarrow\operatorname{gr}(M).
  \tag{10.1.1.2}
\]
\end{env}

\oldpage[0\textsubscript{III}]{18}
\begin{env}[10.1.2]
\label{0.10.1.2-fr}
다음 가정 가운데 하나가 성립한다고 하자.
\begin{enumerate}
  \item[(i)] $\mathfrak{J}$는 멱영이다.
  \item[(ii)] $A$는 뇌터이고, $\mathfrak{J}$는 $A$의 근기에 포함되며,
    $M$은 유한 생성이다.
\end{enumerate}
그러면 다음 성질들은 서로 동치이다.
\begin{enumerate}
  \item[$a)$] $M$은 자유 $A$-가군이다.
  \item[$b)$] $M/\mathfrak{J}M=M\otimes_A(A/\mathfrak{J})$는 자유
    $(A/\mathfrak{J})$-가군이고,
    $\operatorname{Tor}_1^A(M,A/\mathfrak{J})=0$이다.
  \item[$c)$] $M/\mathfrak{J}M$은 자유 $(A/\mathfrak{J})$-가군이고
    표준 준동형 (\hyperref[0.10.1.1.2-fr]{10.1.1.2})는 단사이다
    (따라서 전단사이다).
\end{enumerate}
\end{env}

\begin{env}[10.1.3]
\label{0.10.1.3-fr}
$A/\mathfrak{J}$가 \emph{체}라고(다시 말해 $\mathfrak{J}$가 극대
아이디얼이라고) 하고, (\hyperref[0.10.1.2-fr]{10.1.2})의 가정 (i),
(ii) 가운데 하나가 성립한다고 하자. 그러면 다음 성질들은 서로 동치이다.
\begin{enumerate}
  \item[$a)$] $M$은 자유 $A$-가군이다.
  \item[$b)$] $M$은 사영 $A$-가군이다.
  \item[$c)$] $M$은 평탄 $A$-가군이다.
  \item[$d)$] $\operatorname{Tor}_1^A(M,A/\mathfrak{J})=0$이다.
  \item[$e)$] 표준 준동형
    (\hyperref[0.10.1.1.2-fr]{10.1.1.2})는 전단사이다.
\end{enumerate}
이 결과는 특히 다음 두 경우에 적용된다.
\begin{enumerate}
  \item[(i)] 극대 아이디얼 $\mathfrak{J}$가 \emph{멱영}인 국소환 $A$ 위의
    \emph{임의의} 가군 $M$(예를 들어 $A$가 아르틴 국소환인 경우).
  \item[(ii)] \emph{뇌터} 국소환 위의 \emph{유한 생성} 가군 $M$.
\end{enumerate}
\end{env}

\subsection{평탄성의 국소 판정법}
\label{subsection:0.10.2-fr}

\begin{env}[10.2.1]
\label{0.10.2.1-fr}
가정과 표기는 (\hyperref[0.10.1.1-fr]{10.1.1})에서와 같다고 하고, 다음
조건들을 생각하자.
\begin{enumerate}
  \item[$a)$] $M$은 평탄 $A$-가군이다.
  \item[$b)$] $M/\mathfrak{J}M$은 평탄 $(A/\mathfrak{J})$-가군이고
    $\operatorname{Tor}_1^A(M,A/\mathfrak{J})=0$이다.
  \item[$c)$] $M/\mathfrak{J}M$은 평탄 $(A/\mathfrak{J})$-가군이고
    표준 준동형 (\hyperref[0.10.1.1.2-fr]{10.1.1.2})는 전단사이다.
  \item[$d)$] 모든 $n\geq1$에 대하여 $M/\mathfrak{J}^nM$은 평탄
    $(A/\mathfrak{J}^n)$-가군이다.
\end{enumerate}
그러면 다음 함의들이 성립한다.
\[
  a\Rightarrow b\Rightarrow c\Rightarrow d
\]
또한 $\mathfrak{J}$가 \emph{멱영}이면 네 조건 $a)$, $b)$, $c)$, $d)$는
\emph{서로 동치}이다. $A$가 뇌터이고, 더 나아가 $M$이
\emph{아이디얼적으로 분리}되어 있을 때에도 마찬가지이다. 여기서
아이디얼적으로 분리되어 있다는 것은 $A$의 모든 아이디얼
$\mathfrak{a}$에 대하여 $A$-가군 $\mathfrak{a}\otimes_A M$이
$\mathfrak{J}$-준아딕 위상에 관해 \emph{분리}되어 있다는 뜻이다.
\end{env}

\begin{env}[10.2.2]
\label{0.10.2.2-fr}
$A$를 뇌터환, $B$를 가환 뇌터 $A$-대수, $\mathfrak{J}$를
$\mathfrak{J}B$가 $B$의 근기에 포함되도록 하는 $A$의 아이디얼,
$M$을 유한 생성 $B$-가군이라 하자. 그러면 $M$을 $A$-가군으로 볼 때
(\hyperref[0.10.2.1-fr]{10.2.1})의 조건 $a)$, $b)$, $c)$, $d)$는 서로
동치이다.

이 결과는 특히 $A$와 $B$가 뇌터 \emph{국소환}이고 준동형 $A\to B$가
\emph{국소 준동형}일 때 적용된다. 더 구체적으로 이때
$\mathfrak{J}$가 $A$의 \emph{극대} 아이디얼이면, 조건 $b)$와 $c)$에서
$M/\mathfrak{J}M$이 평탄이라는 자동으로 성립하는 가정을 생략할 수
있고, 조건 $d)$는 가군 $M/\mathfrak{J}^nM$이
$A/\mathfrak{J}^n$ 위에서 \emph{자유}라는 뜻이다.
\end{env}

\oldpage[0\textsubscript{III}]{19}
\begin{env}[10.2.3]
\label{0.10.2.3-fr}
$A$, $B$, $\mathfrak{J}$, $M$에 대한 가정은
(\hyperref[0.10.2.2-fr]{10.2.2})의 처음에 서술한 것과 같다고 하자.
$\widehat A$를 $A$의 $\mathfrak{J}$-준아딕 위상에 대한 분리 완비화,
$\widehat M$을 $M$의 $\mathfrak{J}B$-준아딕 위상에 대한 분리
완비화라 하자. 그러면 $M$이 평탄 $A$-가군이기 위한 필요충분조건은
$\widehat M$이 평탄 $\widehat A$-가군인 것이다.
\end{env}

\begin{env}[10.2.4]
\label{0.10.2.4-fr}
$\rho:A\to B$를 뇌터 국소환들의 국소 준동형, $k$를 $A$의 잉여체,
$M$, $N$을 두 유한 생성 $B$-가군이라 하고, $N$은 $A$-\emph{평탄}이라고
가정하자. $u:M\to N$을 $B$-준동형이라 하자. 그러면 다음 조건들은 서로
동치이다.
\begin{enumerate}
  \item[$a)$] $u$는 단사이고 $\operatorname{Coker}(u)$는 평탄
    $A$-가군이다.
  \item[$b)$] $u\otimes1:M\otimes_A k\to N\otimes_A k$는 단사이다.
\end{enumerate}
\end{env}

\begin{env}[10.2.5]
\label{0.10.2.5-fr}
$\rho:A\to B$, $\sigma:B\to C$를 뇌터 국소환들의 국소 준동형,
$k$를 $A$의 잉여체, $M$을 유한 생성 $C$-가군이라 하자. $B$가 평탄
$A$-가군이라고 가정한다. 그러면 다음 조건들은 서로 동치이다.
\begin{enumerate}
  \item[$a)$] $M$은 평탄 $B$-가군이다.
  \item[$b)$] $M$은 평탄 $A$-가군이고, $M\otimes_A k$는 평탄
    $(B\otimes_A k)$-가군이다.
\end{enumerate}
\end{env}

\begin{proposition}[10.2.6]
\label{0.10.2.6-fr}
$A$, $B$를 두 뇌터 국소환, $\rho:A\to B$를 국소 준동형,
$\mathfrak{J}$를 극대 아이디얼에 포함되는 $B$의 아이디얼, $M$을
유한 생성 $B$-가군이라 하자. 모든 $n\geq0$에 대하여
$M_n=M/\mathfrak{J}^{n+1}M$이 평탄 $A$-가군이라고 가정한다. 그러면
$M$은 평탄 $A$-가군이다.
\end{proposition}

유한 생성 $A$-가군들의 모든 단사 준동형 $u:N'\to N$에 대하여
$v=1\otimes u:M\otimes_A N'\to M\otimes_A N$이 단사임을 증명해야 한다.
그런데 $M\otimes_A N'$과 $M\otimes_A N$은 유한 생성 $B$-가군이므로
$\mathfrak{J}$-준아딕 위상에 관해 분리되어 있다
(\hyperref[0.7.3.5-ko]{$0_{\mathrm I}$, 7.3.5}). 따라서 분리 완비화들
사이의 준동형
$\widehat v:\widehat{(M\otimes_A N')}\to\widehat{(M\otimes_A N)}$이
단사임을 증명하면 충분하다. 한편 $\widehat v=\varprojlim v_n$이고,
여기서 $v_n$은 준동형
$1\otimes u:M_n\otimes_A N'\to M_n\otimes_A N$이다. 가정에 따라
$M_n$은 $A$-평탄이므로 모든 $n$에 대해 $v_n$은 단사이고, 따라서
$\widehat v$도 단사이다. 실제로 함자 $\varprojlim$은 왼쪽 완전이다.

\begin{corollary}[10.2.7]
\label{0.10.2.7-fr}
$A$를 뇌터환, $B$를 뇌터 국소환, $\rho:A\to B$를 준동형,
$f$를 $B$의 극대 아이디얼의 원소, $M$을 유한 생성 $B$-가군이라 하자.
$M$의 곱셈 준동형 $f_M:x\mapsto fx$가 단사이고 $M/fM$이 평탄
$A$-가군이라고 가정한다. 그러면 $M$은 평탄 $A$-가군이다.
\end{corollary}

$i\geq0$에 대하여 $M_i=f^iM$이라 두자. $f_M$이 단사이므로
$M_i/M_{i+1}$은 $M/fM$과 동형이고, 따라서 모든 $i\geq0$에 대하여
$A$-평탄이다. 완전열
\[
  0\longrightarrow M_i/M_{i+1}\longrightarrow M/M_{i+1}
  \longrightarrow M/M_i\longrightarrow0
\]
에서 $i$에 대한 귀납법으로 모든 $i\geq0$에 대해 $M/M_i$가
$A$-\emph{평탄}임을 얻는다
(\hyperref[0.6.1.2-ko]{$0_{\mathrm I}$, 6.1.2}). 따라서
(\hyperref[0.10.2.6-fr]{10.2.6})을 적용할 수 있다. 다음과 같이 직접
논증할 수도 있다. 모든 유한 생성 $A$-가군 $N$에 대하여
$M\otimes_A N$은 유한 생성 $B$-가군이다. $f$가 $B$의 근기
$\mathfrak n$에 속하므로 $M\otimes_A N$ 위의 $(f)$-아딕 위상은
\oldpage[0\textsubscript{III}]{20}
$\mathfrak n$-아딕 위상보다 더 섬세하며, 후자는 \emph{분리}되어 있음을
알고 있다 (\hyperref[0.7.3.5-ko]{$0_{\mathrm I}$, 7.3.5}). 또한
$M/M_i$가 $A$-평탄이므로
\[
  f^i(M\otimes_A N)=\operatorname{Im}(M_i\otimes_A N\to M\otimes_A N)
  =\operatorname{Ker}(M\otimes_A N\to(M/M_i)\otimes_A N)
\]
이다 (\hyperref[0.6.1.2-ko]{$0_{\mathrm I}$, 6.1.2}). 이제 $N$을 유한
생성 $A$-가군, $N'$을 $N$의 부분가군, $j:N'\to N$을 표준 포함
준동형이라 하자. 가환 도식
\[
\xymatrix{
  M\otimes_A N'\ar[r]\ar[d]_{1_M\otimes j} &
  (M/M_i)\otimes_A N'\ar[d]^{1_{M/M_i}\otimes j}\\
  M\otimes_A N\ar[r] & (M/M_i)\otimes_A N
}
\]
에서 $M/M_i$가 $A$-평탄이므로 $1_{M/M_i}\otimes j$는 단사이다.
따라서 모든 $i$에 대하여
\[
  \operatorname{Ker}(M\otimes_A N'\to M\otimes_A N)
  \subset\operatorname{Ker}(M\otimes_A N'\to(M/M_i)\otimes_A N')
\]
이다. 위에서 본 대로 우변들의 교집합은 $0$만으로 이루어져 있으므로
좌변도 마찬가지이고, 따라서 $M$은 $A$-평탄이다.

\begin{proposition}[10.2.8]
\label{0.10.2.8-fr}
$A$를 축소 뇌터환, $M$을 유한 생성 $A$-가군이라 하자. 이산
값매김환인 모든 $A$-대수 $B$에 대하여 $M\otimes_A B$가 평탄
$B$-가군이라고(따라서 자유 가군이라고
(\hyperref[0.10.1.3-fr]{10.1.3})) 가정한다. 그러면 $M$은 평탄
$A$-가군이다.
\end{proposition}

$M$이 평탄이기 위한 필요충분조건은 $A$의 모든 극대 아이디얼
$\mathfrak m$에 대하여 $M_{\mathfrak m}$이 평탄
$A_{\mathfrak m}$-가군인 것임을 알고 있다
(\hyperref[0.6.3.3-ko]{$0_{\mathrm I}$, 6.3.3}). 따라서 $A$가
\emph{국소}인 경우로 한정해도 된다
(\hyperref[0.1.2.8-ko]{$0_{\mathrm I}$, 1.2.8}). 이제
$\mathfrak m$을 $A$의 극대 아이디얼, $\mathfrak p_i$
$(1\leq i\leq r)$를 그 극소 소아이디얼들, $k$를 잉여체
$A/\mathfrak m$이라 하자. 각 $i$에 대하여 정역
$A/\mathfrak p_i$와 같은 분수체 $K_i$를 가지며 이 정역을 지배하는
이산 값매김환 $B_i$가 존재함을 알고 있다
(\hyperref[II.7.1.7-ko]{II, 7.1.7}). $M_i=M\otimes_A B_i$라 두자.
가정에 따라 $M_i$는 $B_i$ 위에서 자유이므로, $k_i$를 $B_i$의
잉여체라 하면
\[
\label{0.10.2.8.1-fr}
  \operatorname{rg}_{k_i}(M_i\otimes_{B_i}k_i)
  =\operatorname{rg}_{K_i}(M_i\otimes_{B_i}K_i).
  \tag{10.2.8.1}
\]
이다. 그런데 합성 준동형 $A\to A/\mathfrak p_i\to B_i$는 명백히
국소 준동형이므로 $k_i$는 $k$의 확대체이고,
$M_i\otimes_{B_i}k_i=M\otimes_A k_i=(M\otimes_A k)\otimes_k k_i$이며,
한편 $M_i\otimes_{B_i}K_i=M\otimes_A K_i$이다. 따라서 등식
(\hyperref[0.10.2.8.1-fr]{10.2.8.1})로부터
\[
  \operatorname{rg}_k(M\otimes_A k)=\operatorname{rg}_{K_i}(M\otimes_A K_i)
  \qquad(1\leq i\leq r\text{에 대하여})
\]
를 얻는다. $A$는 축소환이므로 이 조건으로부터 $M$이 자유
$A$-가군임을 알고 있다(부르바키, \emph{Alg. comm.}, 제II장,
\S~3, n\textsuperscript{o}~2, 명제 7).

\subsection{국소환의 평탄 확대의 존재}
\label{subsection:0.10.3-fr}

\begin{proposition}[10.3.1]
\label{0.10.3.1-fr}
$A$를 뇌터 국소환, $\mathfrak J$를 그 극대 아이디얼,
$k=A/\mathfrak J$를 그 잉여체라 하자. $K$를 $k$의 확대체라 하자.
$B/\mathfrak J B$가 $K$와 $k$-동형이고 $B$가 평탄 $A$-가군이 되도록
하는, $A$에서 뇌터 국소환 $B$로의 국소 준동형이 존재한다.
\end{proposition}

이 명제를 여러 단계로 증명하겠다.

\begin{env}[10.3.1.1]
\label{0.10.3.1.1-fr}
먼저 $T$가 부정원이고 $K=k(T)$라고 하자. 다항식환 $A'=A[T]$에서
계수들이 아이디얼 $\mathfrak J$에 속하는 다항식들로 이루어진
소아이디얼 $\mathfrak J'=\mathfrak J A'$을 생각하자.
\oldpage[0\textsubscript{III}]{21}
$A'/\mathfrak J'$은 $k[T]$와 표준적으로 동형임이 명백하다. 분수환
$B=A'_{\mathfrak J'}$이 조건을 충족함을 보이자. 이는 극대 아이디얼이
$\mathfrak L=\mathfrak J B$인 뇌터 국소환임이 명백하다. 또한
$B/\mathfrak L=(A'/\mathfrak J')_{\mathfrak J'}=(k[T])_{\mathfrak J'}$은
$k[T]$의 분수체 $K$에 다름 아니다. 마지막으로 $B$는 평탄
$A'$-가군이고 $A'$은 자유 $A$-가군이므로, $B$는 평탄 $A$-가군이다
(\hyperref[0.6.2.1-ko]{$0_{\mathrm I}$, 6.2.1}).
\end{env}

\begin{env}[10.3.1.2]
\label{0.10.3.1.2-fr}
다음으로 $t$가 $k$ 위에서 대수적이고 $K=k(t)=k[t]$라고 하자.
$f\in k[T]$를 $t$의 최소다항식이라 하자. $k[T]$에서의 표준 상이
$f$인 일계수 다항식 $F\in A[T]$가 존재한다. 다시 $A'=A[T]$라 두고,
$\mathfrak J'$을 $A'$의 아이디얼 $\mathfrak J A'+(F)$라 하자. 이번에는
몫환 $B=A'/(F)$가 조건을 충족함을 보이겠다. 우선 $B$가 자유
$A$-가군이고, 따라서 평탄임이 명백하다. 환 $A'/\mathfrak J'$은
\[
  (A'/\mathfrak J A')/((\mathfrak J A'+(F))/\mathfrak J A')
  =k[T]/(f)=K
\]
와 동형이다. 따라서 $\mathfrak J'$의 $B$에서의 상 $\mathfrak L$은
극대이고, 명백히 $\mathfrak L=\mathfrak J B$이다. 마지막으로 $B$는
유한 생성 $A$-가군이므로 반국소환이다(부르바키,
\emph{Alg. comm.}, 제IV장, \S~2, n\textsuperscript{o}~5, 명제 9의
따름정리 3). 그리고 $B$의 극대 아이디얼들은 $B/\mathfrak J B$의
극대 아이디얼들과 일대일 대응한다([13], 제I권, 259쪽). 따라서 앞의
결과로부터 $B$가 국소환임을 알 수 있다.
\end{env}

\begin{lemma}[10.3.1.3]
\label{0.10.3.1.3-fr}
$(A_\lambda,f_{\mu\lambda})$를 국소환들의 여과 귀납계라 하고,
$f_{\mu\lambda}$들은 국소 준동형이라고 하자. $\mathfrak m_\lambda$를
$A_\lambda$의 극대 아이디얼, $K_\lambda=A_\lambda/\mathfrak m_\lambda$라
하자. 그러면 $A'=\varinjlim A_\lambda$는 국소환이고,
$\mathfrak m'=\varinjlim\mathfrak m_\lambda$는 그 극대 아이디얼이며,
$K=\varinjlim K_\lambda$는 그 잉여체이다. 또한 $\lambda<\mu$일 때
$\mathfrak m_\mu=\mathfrak m_\lambda A_\mu$이면, 모든 $\lambda$에 대해
$\mathfrak m'=\mathfrak m_\lambda A'$이다. 더 나아가
$\lambda<\mu$일 때 $A_\mu$가 평탄 $A_\lambda$-가군이고 모든
$A_\lambda$가 뇌터이면, $A'$은 뇌터이며 모든 $\lambda$에 대하여
평탄 $A_\lambda$-가군이다.
\end{lemma}

가정에 따라 $\lambda<\mu$이면
$f_{\mu\lambda}(\mathfrak m_\lambda)\subset\mathfrak m_\mu$이므로,
$\mathfrak m_\lambda$들은 귀납계를 이루고 그 극한 $\mathfrak m'$은
명백히 $A'$의 아이디얼이다. 또한 $x'\notin\mathfrak m'$이면 어떤
$\lambda$와 $x_\lambda\in A_\lambda$에 대하여
$x'=f_\lambda(x_\lambda)$이다. 여기서
$f_\lambda:A_\lambda\to A'$은 표준 준동형이다. $x'$가
$\mathfrak m'$에 속하지 않으므로 필연적으로
$x_\lambda\notin\mathfrak m_\lambda$이고, 따라서 $x_\lambda$는
$A_\lambda$에서 역원 $y_\lambda$를 가진다. 그러면
$y'=f_\lambda(y_\lambda)$는 $A'$에서 $x'$의 역원이다. 이로써 $A'$이
극대 아이디얼 $\mathfrak m'$을 갖는 국소환임이 증명된다. $K$에 관한
주장은 $\varinjlim$이 완전 함자라는 사실에서 곧바로 따른다. 가정
$\mathfrak m_\mu=\mathfrak m_\lambda A_\mu$는 표준 사상
$\mathfrak m_\lambda\otimes_{A_\lambda}A_\mu\to\mathfrak m_\mu$가
전사라는 뜻이다. 따라서 관계
$\mathfrak m'=\mathfrak m_\lambda A'$도 함자 $\varinjlim$의 완전성과
이 함자가 텐서곱과 가환한다는 사실에서 따른다.

이제 $\lambda<\mu$일 때
$\mathfrak m_\mu=\mathfrak m_\lambda A_\mu$이고 $A_\mu$가 평탄
$A_\lambda$-가군이라고 하자. 그러면
(\hyperref[0.6.2.3-ko]{$0_{\mathrm I}$, 6.2.3})에 따라 모든
$\lambda$에 대하여 $A'$은 평탄 $A_\lambda$-가군이다. $A'$과
$A_\lambda$가 국소환이고 $\mathfrak m'=\mathfrak m_\lambda A'$이므로,
$A'$은 실제로 충실 평탄 $A_\lambda$-가군이다
(\hyperref[0.6.6.2-ko]{$0_{\mathrm I}$, 6.6.2}). 끝으로 모든
$A_\lambda$가 \emph{뇌터}라고 하자. 그러면
$\mathfrak m_\lambda$-준아딕 위상들은 분리되어 있다
(\hyperref[0.7.3.5-ko]{$0_{\mathrm I}$, 7.3.5}). 이로부터 먼저 $A'$의
$\mathfrak m'$-아딕 위상이 \emph{분리}되어 있음을 보이자. 실제로
$x'\in A'$이 모든 $\mathfrak m'^n$ ($n>0$)에 속하면, 어떤 지표
$\mu$와 $x_\mu\in A_\mu$에 대하여 $x'$은 $x_\mu$의 상이다. 그런데
$A_\mu$에서 $\mathfrak m'^n=\mathfrak m_\mu^nA'$의 역상은
$\mathfrak m_\mu^n$이므로
(\hyperref[0.6.6.1-ko]{$0_{\mathrm I}$, 6.6.1}), $x_\mu$는 모든
$\mathfrak m_\mu^n$에 속한다. 따라서 가정에 의해 $x_\mu=0$이고,
결국 $x'=0$이다. $\widehat A'$을 $A'$의 $\mathfrak m'$-아딕 위상에
대한 완비화라 하자. 앞의 결과로부터 $A'\subset\widehat A'$이다.
$\widehat A'$이 \emph{뇌터}이고 모든 $\lambda$에 대하여
$A_\lambda$-\emph{평탄}임을 보이겠다. 그러면
\oldpage[0\textsubscript{III}]{22}
$\widehat A'$은 $A'$-평탄이고
(\hyperref[0.6.2.3-ko]{$0_{\mathrm I}$, 6.2.3}),
$\mathfrak m'\widehat A'\ne\widehat A'$이므로 $\widehat A'$은 충실 평탄
$A'$-가군이다 (\hyperref[0.6.6.2-ko]{$0_{\mathrm I}$, 6.6.2}). 따라서
마침내 $A'$이 \emph{뇌터}임을 얻으며
(\hyperref[0.6.5.2-ko]{$0_{\mathrm I}$, 6.5.2}), 이로써 보조정리의
증명이 끝난다.

$\widehat A'=\varprojlim_n A'/\mathfrak m'^n$이다. $A'$이
$A_\lambda$-평탄이므로
\[
  \mathfrak m'^{n}/\mathfrak m'^{n+1}
  =(\mathfrak m_\lambda^n/\mathfrak m_\lambda^{n+1})
    \otimes_{A_\lambda}A'
  =(\mathfrak m_\lambda^n/\mathfrak m_\lambda^{n+1})
    \otimes_{K_\lambda}(K_\lambda\otimes_{A_\lambda}A')
  =(\mathfrak m_\lambda^n/\mathfrak m_\lambda^{n+1})
    \otimes_{K_\lambda}K
\]
이다. $\mathfrak m_\lambda^n/\mathfrak m_\lambda^{n+1}$은 유한 차원
$K_\lambda$-벡터공간이므로, 모든 $n\geq0$에 대하여
$\mathfrak m'^n/\mathfrak m'^{n+1}$은 유한 차원 $K$-벡터공간이다.
따라서 (\hyperref[0.7.2.12-ko]{$0_{\mathrm I}$, 7.2.12})와
(\hyperref[0.7.2.8-ko]{$0_{\mathrm I}$, 7.2.8})에 의해
$\widehat A'$은 뇌터이다. 또한 $\widehat A'$의 극대 아이디얼은
$\mathfrak m'\widehat A'$이고,
$\widehat A'/\mathfrak m'^n\widehat A'$은 $A'/\mathfrak m'^n$과
동형임을 알고 있다. 그리고
$A'/\mathfrak m'^n=(A_\lambda/\mathfrak m_\lambda^n)
\otimes_{A_\lambda}A'$이므로, $A'/\mathfrak m'^n$은 평탄
$(A_\lambda/\mathfrak m_\lambda^n)$-가군이다
(\hyperref[0.6.2.1-ko]{$0_{\mathrm I}$, 6.2.1}). 따라서 판정법
(\hyperref[0.10.2.2-fr]{10.2.2})을 뇌터 $A_\lambda$-대수
$\widehat A'$에 적용할 수 있고, 이로부터 $\widehat A'$이
$A_\lambda$-평탄임을 얻는다. C.Q.F.D.

\begin{env}[10.3.1.4]
\label{0.10.3.1.4-fr}
이제 일반적인 경우를 다루자. 순서수 $\gamma$와, 모든 순서수
$\lambda\leq\gamma$에 대하여 $k$를 포함하는 $K$의 부분체
$k_\lambda$가 다음 조건들을 만족하도록 존재한다. $1^\circ$ 모든
$\lambda<\gamma$에 대하여 $k_{\lambda+1}$은 단 하나의 원소로 생성되는
$k_\lambda$의 확대체이다. $2^\circ$ 바로 앞 원소가 없는 모든 순서수
$\mu$에 대하여
$k_\mu=\bigcup_{\lambda<\mu}k_\lambda$이다. $3^\circ$ $K=k_\gamma$이다.
실제로 적당한 $\beta$에 대하여 순서수 $\xi\leq\beta$들의 집합에서
$K$로 가는 전단사 $\xi\mapsto t_\xi$를 택하면 충분하다.
$\lambda\leq\beta$에 대하여 $k_\lambda$를 초한 귀납법으로 다음과 같이
정의한다. $\lambda$에 바로 앞 원소가 없으면 $\mu<\lambda$인
$k_\mu$들의 합집합으로 두고, $\lambda=\nu+1$이면
$t_\xi\notin k_\nu$가 되는 가장 작은 순서수 $\xi$를 택하여
$k_\nu(t_\xi)$로 둔다. 그러면 $\gamma$는 정의상
$k_\gamma=K$가 되는 $\beta$ 이하의 가장 작은 순서수이다.

이제 초한 귀납법으로 $\lambda\leq\gamma$에 대한 뇌터 국소환
$A_\lambda$의 족과, $\lambda\leq\mu$에 대한 국소 준동형
$f_{\mu\lambda}:A_\lambda\to A_\mu$를 다음 조건들을 만족하도록
정의하겠다.
\begin{enumerate}
  \item[(i)] $(A_\lambda,f_{\mu\lambda})$는 귀납계이고 $A_0=A$이다.
  \item[(ii)] 모든 $\lambda$에 대하여 $k$-동형
    $A_\lambda/\mathfrak J A_\lambda\xrightarrow{\sim}k_\lambda$가 있다.
  \item[(iii)] $\lambda\leq\mu$이면 $A_\mu$는 평탄
    $A_\lambda$-가군이다.
\end{enumerate}

$\lambda<\mu<\xi$에 대해 $A_\lambda$와 $f_{\mu\lambda}$가 정의되어
있다고 하자. 먼저 $\xi=\zeta+1$이어서
$k_\xi=k_\zeta(t)$인 경우를 생각하자. $t$가 $k_\zeta$ 위에서
초월적이면, (\hyperref[0.10.3.1.1-fr]{10.3.1.1})의 절차에 따라
$A_\xi=(A_\zeta[t])_{\mathfrak J A_\zeta[t]}$로 정의한다.
$f_{\xi\zeta}$는 표준 사상이고, $\lambda<\zeta$이면
$f_{\xi\lambda}=f_{\xi\zeta}\circ f_{\zeta\lambda}$로 둔다.
(\hyperref[0.10.3.1.1-fr]{10.3.1.1})에서 증명한 것에 비추어 조건
(i)--(iii)의 확인은 즉시 이루어진다. 다음으로 $t$가 대수적이라고 하고,
$h$를 $k_\zeta[T]$에서의 그 최소다항식, $H$를 $k_\zeta[T]$에서의 상이
$h$인 $A_\zeta[T]$의 일계수 다항식이라 하자. 이때
$A_\xi=A_\zeta[T]/(H)$로 두고, $f_{\xi\lambda}$들은 앞에서와 같이
정의한다. 그러면 조건 (i)--(iii)은
(\hyperref[0.10.3.1.2-fr]{10.3.1.2})에서 본 결과로부터 따른다.

이제 $\xi$에 바로 앞 원소가 없다고 하자. $A_\xi$를
$\lambda<\xi$에 대한 국소환들의 귀납계
$(A_\lambda,f_{\mu\lambda})$의 귀납극한으로 두고,
$\lambda<\xi$에 대한 $f_{\xi\lambda}$는 표준 사상으로 정의한다.
$A_\xi$가 뇌터 국소환이라는 사실, $f_{\xi\lambda}$들이 국소
준동형이라는 사실, 그리고 $\lambda\leq\xi$에 대한 조건 (i)--(iii)은
\oldpage[0\textsubscript{III}]{23}
귀납가정과 보조정리
(\hyperref[0.10.3.1.3-fr]{10.3.1.3})에서 따른다. 이 구성을 마치면
$B=A_\gamma$가 (\hyperref[0.10.3.1-fr]{10.3.1})의 명제를 만족함이
명백하다.
\end{env}

(\hyperref[0.10.2.1-fr]{10.2.1, $c)$})에 따라 표준 동형
\[
\label{0.10.3.1.5-fr}
  \operatorname{gr}(A)\otimes_k K\xrightarrow{\sim}\operatorname{gr}(B).
  \tag{10.3.1.5}
\]
이 있음을 유의하자.

한편 (\hyperref[0.10.3.1-fr]{10.3.1})의 결론을 바꾸지 않고 $B$를 그
$\mathfrak J B$-아딕 완비화 $\widehat B$로 바꿀 수 있다. 실제로
$\widehat B$는 평탄 $B$-가군이고
(\hyperref[0.7.3.3-ko]{$0_{\mathrm I}$, 7.3.3}), 따라서 평탄
$A$-가군이다 (\hyperref[0.6.2.1-ko]{$0_{\mathrm I}$, 6.2.1}).

또한 다음 따름정리가 증명되었다.

\begin{corollary}[10.3.2]
\label{0.10.3.2-fr}
$K$가 유한 차수 확대체이면, $B$가 유한 $A$-대수라고 가정할 수 있다.
\end{corollary}

\section{호몰로지 대수학에 관한 보충}
\label{section:0.11-fr}

\subsection{스펙트럼 열에 관한 복습}
\label{subsection:0.11.1-fr}

\begin{env}[11.1.1]
\label{0.11.1.1-fr}
이후에는 (T, 2.4)에서 정의한 것보다 더 일반적인 스펙트럼 열의 개념을
사용한다. (T, 2.4)의 표기를 유지하면서, 아벨 범주 $\mathcal C$에서
다음 자료들로 이루어진 계 E를 \emph{스펙트럼 열}이라 부르겠다.
\begin{enumerate}
  \item[(a)] $p\in\mathbb Z$, $q\in\mathbb Z$, $r\geq2$에 대하여
    정의되는 $\mathcal C$의 대상들의 족 $(E_r^{pq})$.
  \item[(b)]
    $d_r^{p+r,q-r+1}d_r^{pq}=0$을 만족하는 사상들의 족
    $d_r^{pq}:E_r^{pq}\to E_r^{p+r,q-r+1}$. 다음과 같이 둔다.
    $Z_{r+1}(E_r^{pq})=\operatorname{Ker}(d_r^{pq})$,
    $B_{r+1}(E_r^{pq})=\operatorname{Im}(d_r^{p-r,q+r-1})$. 그러면
    \[
      B_{r+1}(E_r^{pq})\subset Z_{r+1}(E_r^{pq})\subset E_r^{pq}.
    \]
  \item[(c)] 동형사상들의 족
    $\alpha_r^{pq}:Z_{r+1}(E_r^{pq})/B_{r+1}(E_r^{pq})
    \xrightarrow{\sim}E_{r+1}^{pq}$.
\end{enumerate}

이제 $k\geq r+1$에 대하여 $k$에 대한 귀납법으로 부분대상
$B_k(E_r^{pq})$와 $Z_k(E_r^{pq})$를 다음과 같이 정의한다. 이들은
표준 사상 $E_r^{pq}\to E_r^{pq}/B_{r+1}(E_r^{pq})$에 의한, 각각
$\alpha_r^{pq}$를 통해 부분대상 $B_k(E_{r+1}^{pq})$와
$Z_k(E_{r+1}^{pq})$로 식별되는 이 몫의 부분대상들의 역상이다. 그러면
동형을 제외하면 명백히
\[
\label{0.11.1.1.1-fr}
  Z_k(E_r^{pq})/B_k(E_r^{pq})=E_k^{pq}
  \qquad(k\geq r+1\text{에 대하여})
  \tag{11.1.1.1}
\]
이고, 다시 $B_r(E_r^{pq})=0$, $Z_r(E_r^{pq})=E_r^{pq}$로 두면 다음
포함관계들이 성립한다.
\[
\label{0.11.1.1.2-fr}
  0=B_r(E_r^{pq})\subset B_{r+1}(E_r^{pq})\subset B_{r+2}(E_r^{pq})
  \subset\cdots\subset Z_{r+2}(E_r^{pq})\subset Z_{r+1}(E_r^{pq})
  \subset Z_r(E_r^{pq})=E_r^{pq}.
  \tag{11.1.1.2}
\]

E의 나머지 자료는 다음과 같다.
\begin{enumerate}
  \item[(d)] 다음을 만족하는 $E_2^{pq}$의 두 부분대상
    $B_\infty(E_2^{pq})$와 $Z_\infty(E_2^{pq})$:
    $B_\infty(E_2^{pq})\subset Z_\infty(E_2^{pq})$이고, 모든
    $k\geq2$에 대하여
    \[
      B_k(E_2^{pq})\subset B_\infty(E_2^{pq})
      \qquad\text{및}\qquad
      Z_\infty(E_2^{pq})\subset Z_k(E_2^{pq}).
    \]
    다음과 같이 둔다.
    \[
    \label{0.11.1.1.3-fr}
      E_\infty^{pq}=Z_\infty(E_2^{pq})/B_\infty(E_2^{pq}).
      \tag{11.1.1.3}
    \]
  \item[(e)]
    \oldpage[0\textsubscript{III}]{24}
    $\mathcal C$의 대상들의 족 $(E^n)$. 각 대상에는 \emph{감소 여과}
    $(F^p(E^n))_{p\in\mathbb Z}$가 주어져 있다. 통상과 같이
    $\operatorname{gr}(E^n)$으로 여과 대상 $E^n$에 연관된 등급 대상을
    나타낸다. 이는
    $\operatorname{gr}_p(E^n)=F^p(E^n)/F^{p+1}(E^n)$들의 직합이다.
  \item[(f)] 모든 쌍 $(p,q)\in\mathbb Z\times\mathbb Z$에 대한 동형사상
    $\beta^{pq}:E_\infty^{pq}\xrightarrow{\sim}
    \operatorname{gr}_p(E^{p+q})$.
\end{enumerate}

여과를 제외한 족 $(E^n)$을 스펙트럼 열 E의 \emph{수렴 대상}이라
부른다.

범주 $\mathcal C$가 무한 직합을 가지거나, 모든 $r\geq2$와 모든
$n\in\mathbb Z$에 대하여 $p+q=n$이고 $E_r^{pq}\ne0$인 쌍 $(p,q)$이
유한 개라고 하자($r=2$에 대해서만 그러하다고 해도 충분하다). 그러면
\[
  E_r^{(n)}=\sum_{p+q=n}E_r^{pq}
\]
이 정의된다. $p+q=n$인 각 쌍 $(p,q)$에 대하여 $E_r^{pq}$로의 제한이
$d_r^{pq}$인 사상 $E_r^{(n)}\to E_r^{(n+1)}$을 $d_r^{(n)}$으로
나타내면 $d_r^{(n+1)}\circ d_r^{(n)}=0$이다. 다시 말해
$(E_r^{(n)})_{n\in\mathbb Z}$는 $\mathcal C$에서 차수 $+1$인 미분
연산자를 갖는 \emph{복합체 $E_r^{(\bullet)}$}이다. 그리고 c)로부터
모든 $r\geq2$에 대하여 $\mathrm H^n(E_r^{(\bullet)})$가
$E_{r+1}^{(n)}$과 동형임을 얻는다.
\end{env}

\begin{env}[11.1.2]
\label{0.11.1.2-fr}
스펙트럼 열 E에서 스펙트럼 열
$E'=(E_r^{\prime pq},E^{\prime n})$로 가는 \emph{사상 $u:E\to E'$}은
사상들의 계
$u_r^{pq}:E_r^{pq}\to E_r^{\prime pq}$,
$u^n:E^n\to E^{\prime n}$으로 이루어진다. 여기서 $u^n$들은
$E^n$과 $E^{\prime n}$의 여과와 양립하고, 도식
\[
  \xymatrix{
    E_r^{pq}\ar[r]^-{d_r^{pq}}\ar[d]_{u_r^{pq}} &
    E_r^{p+r,q-r+1}\ar[d]^{u_r^{p+r,q-r+1}}\\
    E_r^{\prime pq}\ar[r]_{d_r^{\prime pq}} &
    E_r^{\prime p+r,q-r+1}
  }
\]
은 가환한다. 또한 몫을 취하면 $u_r^{pq}$는 사상
\[
  \overline u_r^{pq}:
  Z_{r+1}(E_r^{pq})/B_{r+1}(E_r^{pq})\longrightarrow
  Z_{r+1}(E_r^{\prime pq})/B_{r+1}(E_r^{\prime pq})
\]
을 주며,
$\alpha_r^{\prime pq}\circ\overline u_r^{pq}
=u_{r+1}^{pq}\circ\alpha_r^{pq}$이어야 한다. 끝으로
$u_2^{pq}(B_\infty(E_2^{pq}))\subset B_\infty(E_2^{\prime pq})$,
$u_2^{pq}(Z_\infty(E_2^{pq}))\subset Z_\infty(E_2^{\prime pq})$이어야
한다. 몫을 취하면 $u_2^{pq}$는 사상
$u_\infty^{pq}:E_\infty^{pq}\to E_\infty^{\prime pq}$을 주며, 도식
\[
  \xymatrix{
    E_\infty^{pq}\ar[r]^-{u_\infty^{pq}}\ar[d]_{\beta^{pq}} &
    E_\infty^{\prime pq}\ar[d]^{\beta^{\prime pq}}\\
    \operatorname{gr}_p(E^{p+q})
      \ar[r]_{\operatorname{gr}_p(u^{p+q})} &
    \operatorname{gr}_p(E^{\prime p+q})
  }
\]
은 가환해야 한다.

앞의 정의들로부터 $r$에 대한 귀납법에 의해 다음을 알 수 있다.
$u_2^{pq}$들이 \emph{동형사상}이면 $r\geq2$에 대한
$u_r^{pq}$들도 동형사상이다. \emph{더 나아가}
$u_2^{pq}(B_\infty(E_2^{pq}))=B_\infty(E_2^{\prime pq})$,
$u_2^{pq}(Z_\infty(E_2^{pq}))=Z_\infty(E_2^{\prime pq})$이고
$u^n$들이 동형사상임을 \emph{알고 있다면}, $u$가
\emph{동형사상}이라고 결론지을 수 있다.
\end{env}

\begin{env}[11.1.3]
\label{0.11.1.3-fr}
\oldpage[0\textsubscript{III}]{25}
$(F^p(X))_{p\in\mathbb Z}$가 대상 $X\in\mathcal C$의 (감소) 여과라고
하자. $\inf(F^p(X))=0$이면 이 여과를 \emph{분리}, 어떤 $p$에 대해
$F^p(X)=0$이면 \emph{이산}, $\sup(F^p(X))=X$이면 \emph{소진적}(또는
\emph{공분리}), 어떤 $p$에 대해 $F^p(X)=X$이면 \emph{공이산}이라고
함을 상기하자.

스펙트럼 열 $E=(E_r^{pq},E^n)$가
\[
  B_\infty(E_2^{pq})=\sup_k(B_k(E_2^{pq})),\qquad
  Z_\infty(E_2^{pq})=\inf_k(Z_k(E_2^{pq}))
\]
를 만족하면 \emph{약수렴}한다고 한다. 다시 말해 대상
$B_\infty(E_2^{pq})$와 $Z_\infty(E_2^{pq})$가 스펙트럼 열 E의 자료
a)--c)에 의해 결정된다는 뜻이다. 스펙트럼 열 E가 약수렴하고, 더
나아가 다음을 만족하면 \emph{정칙}이라고 한다.
\begin{enumerate}
  \item[$1^\circ$] 모든 쌍 $(p,q)$에 대하여 감소열
    $(Z_k(E_2^{pq}))_{k\geq2}$는 \emph{어떤 단계 이후로 일정하다}.
    그러면 E가 약수렴한다는 가정에 따라, $p$와 $q$에 의존하는 충분히
    큰 $k$에 대하여
    $Z_\infty(E_2^{pq})=Z_k(E_2^{pq})$이다.
  \item[$2^\circ$] 모든 $n$에 대하여 $E^n$의 여과
    $(F^p(E^n))_{p\in\mathbb Z}$는 \emph{이산}이고 \emph{소진적}이다.
\end{enumerate}

스펙트럼 열 E가 약수렴하고, 더 나아가 다음을 만족하면
\emph{공정칙}이라고 한다.
\begin{enumerate}
  \item[$3^\circ$] 모든 쌍 $(p,q)$에 대하여 증가열
    $(B_k(E_2^{pq}))_{k\geq2}$는 \emph{어떤 단계 이후로 일정하다}.
    따라서 $B_\infty(E_2^{pq})=B_k(E_2^{pq})$이고, 이어서
    $E_\infty^{pq}=\inf E_k^{pq}$이다.
  \item[$4^\circ$] 모든 $n$에 대하여 $E^n$의 여과는
    \emph{공이산}이다.
\end{enumerate}

끝으로 E가 정칙인 동시에 공정칙이면 \emph{쌍정칙}이라고 한다. 다시
말해 다음 조건들이 성립한다.
\begin{enumerate}
  \item[(a)] 모든 쌍 $(p,q)$에 대하여 열
    $(B_k(E_2^{pq}))_{k\geq2}$과 $(Z_k(E_2^{pq}))_{k\geq2}$은
    \emph{어떤 단계 이후로 일정}하고, 충분히 큰 $k$에 대하여
    $B_\infty(E_2^{pq})=B_k(E_2^{pq})$ 및
    $Z_\infty(E_2^{pq})=Z_k(E_2^{pq})$이다. 따라서
    $E_\infty^{pq}=E_k^{pq}$이다.
  \item[(b)] 모든 $n$에 대하여 여과
    $(F^p(E^n))_{p\in\mathbb Z}$는 \emph{이산}이고 \emph{공이산}이다.
    이는 이 여과가 \emph{유한}이라고도 표현한다.
\end{enumerate}

따라서 (T, 2.4)에서 정의한 스펙트럼 열들은 쌍정칙 스펙트럼 열이다.
\end{env}

\begin{env}[11.1.4]
\label{0.11.1.4-fr}
범주 $\mathcal C$에서 여과 귀납극한들이 존재하고 함자 $\varinjlim$이
\emph{완전}하다고 하자. 이는 (T, 1.5)의 공리 AB 5)가 성립한다는 것과
동치이다(T, 1.8 참조). 그러면 대상 $X\in\mathcal C$의 여과
$(F^p(X))_{p\in\mathbb Z}$가 소진적이라는 조건은
$\varinjlim_{p\to-\infty}F^p(X)=X$로도 쓸 수 있다. 스펙트럼 열 E가
약수렴하면
$B_\infty(E_2^{pq})=\varinjlim_{k\to\infty}B_k(E_2^{pq})$이다. 더
나아가 $u:E\to E'$가 E에서 $\mathcal C$의 약수렴 스펙트럼 열 $E'$로
가는 사상이면, $\varinjlim$의 완전성에 의해
$u_2^{pq}(B_\infty(E_2^{pq}))=B_\infty(E_2^{\prime pq})$이다. 또한
다음이 성립한다.
\end{env}

\begin{proposition}[11.1.5]
\label{0.11.1.5-fr}
$\mathcal C$를 여과 귀납극한들이 완전한 아벨 범주, E와 $E'$을
$\mathcal C$의 두 정칙 스펙트럼 열, $u:E\to E'$을 스펙트럼 열들의
사상이라 하자. $u_2^{pq}$들이 동형사상이면 $u$도 동형사상이다.
\end{proposition}

\begin{proof}
$u_r^{pq}$들이 동형사상이고
\[
  u_2^{pq}(B_\infty(E_2^{pq}))=B_\infty(E_2^{\prime pq})
\]
임을 이미 알고 있다 (\hyperref[0.11.1.2-fr]{11.1.2}).
\oldpage[0\textsubscript{III}]{26}
E와 $E'$이 정칙이라는 가정으로부터
$u_2^{pq}(Z_\infty(E_2^{pq}))=Z_\infty(E_2^{\prime pq})$도 따른다.
$u_2^{pq}$가 동형사상이므로 $u_\infty^{pq}$도 동형사상이고, 따라서
$\operatorname{gr}_p(u^{p+q})$도 동형사상이라고 결론지을 수 있다.
그런데 $E^n$과 $E^{\prime n}$의 여과들은 이산이고 소진적이므로,
이로부터 $u^n$들도 동형사상임이 따른다(부르바키,
\emph{Alg. comm.}, 제III장, \S~2, n\textsuperscript{o} 8, 정리 1).
\end{proof}

\begin{env}[11.1.6]
\label{0.11.1.6-fr}
(\hyperref[0.11.1.1.2-fr]{11.1.1.2})와 정의
(\hyperref[0.11.1.1.3-fr]{11.1.1.3})로부터, 스펙트럼 열 E에서
$E_r^{pq}=0$이면 $k\geq r$에 대하여 $E_k^{pq}=0$이고
$E_\infty^{pq}=0$임이 따른다. 어떤 정수 $r\geq2$와, 모든 정수
$n\in\mathbb Z$에 대한 정수 $q(n)$이 존재하여 $q\ne q(n)$이면
$E_r^{n-q,q}=0$일 때, 스펙트럼 열을 \emph{퇴화}한다고 한다. 앞의
관찰로부터 우선 $k\geq r$($k=\infty$ 포함)이고 $q\ne q(n)$이면
$E_k^{n-q,q}=0$임을 얻는다. 또한 $E_{r+1}^{pq}$의 정의로부터
$E_{r+1}^{n-q(n),q(n)}=E_r^{n-q(n),q(n)}$이다. E가 약수렴하면 따라서
$E_\infty^{n-q(n),q(n)}=E_r^{n-q(n),q(n)}$도 성립한다. 다시 말해 모든
$n\in\mathbb Z$에 대하여 $p\ne n-q(n)$이면
$\operatorname{gr}_p(E^n)=0$이고,
$\operatorname{gr}_{n-q(n)}(E^n)=E_r^{n-q(n),q(n)}$이다. 더 나아가
$E^n$의 여과가 이산이고 소진적이면, 열 E는 정칙이고 동형을 제외하면
$E^n=E_r^{n-q(n),q(n)}$이다.
\end{env}

\begin{env}[11.1.7]
\label{0.11.1.7-fr}
범주 $\mathcal C$에서 여과 귀납극한들이 존재하고 완전하다고 하자.
$(E_\lambda,u_{\mu\lambda})$를 여과 지표집합에 따른 $\mathcal C$의
스펙트럼 열들의 귀납계라 하자. 그러면 이 귀납계의 \emph{귀납극한}은
$\mathcal C$의 대상들로 이루어진 스펙트럼 열들의 가법 범주에서
존재한다. 이를 보려면 $E_r^{pq}$, $d_r^{pq}$, $\alpha_r^{pq}$,
$B_\infty(E_2^{pq})$, $Z_\infty(E_2^{pq})$, $E^n$, $F^p(E^n)$,
$\beta^{pq}$를 각각 $E_{r,\lambda}^{pq}$, $d_{r,\lambda}^{pq}$,
$\alpha_{r,\lambda}^{pq}$, $B_\infty(E_{2,\lambda}^{pq})$,
$Z_\infty(E_{2,\lambda}^{pq})$, $E_\lambda^n$,
$F^p(E_\lambda^n)$, $\beta_\lambda^{pq}$의 귀납극한으로 정의하면
충분하다. (\hyperref[0.11.1.1-fr]{11.1.1})의 조건들은 $\mathcal C$에서
함자 $\varinjlim$이 완전하다는 사실로부터 확인된다.
\end{env}

\begin{remark}[11.1.8]
\label{0.11.1.8-fr}
범주 $\mathcal C$가 \emph{뇌터}환 $A$ 위의 $A$-가군들의 범주(각각,
환 $A$ 위의 $A$-가군들의 범주)라고 하자. 그러면
(\hyperref[0.11.1.1-fr]{11.1.1})의 정의들로부터, 주어진 $r$에 대하여
$E_r^{pq}$들이 \emph{유한 생성} $A$-가군(각각, \emph{유한 길이}
$A$-가군)이면 모든 $s\geq r$에 대한 $E_s^{pq}$와
$E_\infty^{pq}$도 그러함을 알 수 있다. 더 나아가 모든 $n$에 대해
수렴 대상 $(E^n)$의 여과가 이산이고 공이산이면, 각 $E^n$도
\emph{유한 생성} $A$-가군(각각, \emph{유한 길이} $A$-가군)이라고
결론지을 수 있다.
\end{remark}

\begin{env}[11.1.9]
\label{0.11.1.9-fr}
스펙트럼 열 E가 $p+q=n$에 관하여 ``균일하게'' 쌍정칙임을 보장하는
조건들을 살펴볼 것이다. 이때 다음 보조정리를 사용한다.
\end{env}

\begin{lemma}[11.1.10]
\label{0.11.1.10-fr}
$(E_r^{pq})$를 (\hyperref[0.11.1.1-fr]{11.1.1})의 자료 a), b), c)로
연결된 $\mathcal C$의 대상들의 족이라 하자. 고정된 정수 $n$에 대하여
다음 성질들은 서로 동치이다.
\begin{enumerate}
  \item[(a)] 정수 $r(n)$이 존재하여, $r\geq r(n)$이고 $p+q=n$ 또는
    $p+q=n-1$이면 모든 사상 $d_r^{pq}$가 0이다.
  \item[(b)] 정수 $r(n)$이 존재하여, $p+q=n$ 또는 $p+q=n+1$이면
    $s\geq r\geq r(n)$에 대하여
    $B_r(E_2^{pq})=B_s(E_2^{pq})$이다.
  \item[(c)] 정수 $r(n)$이 존재하여, $p+q=n$ 또는 $p+q=n-1$이면
    $s\geq r\geq r(n)$에 대하여
    $Z_r(E_2^{pq})=Z_s(E_2^{pq})$이다.
  \item[(d)] 정수 $r(n)$이 존재하여, $p+q=n$이면
    $s\geq r\geq r(n)$에 대하여
    $B_r(E_2^{pq})=B_s(E_2^{pq})$ 및
    $Z_r(E_2^{pq})=Z_s(E_2^{pq})$이다.
\end{enumerate}
\end{lemma}

\begin{proof}
\oldpage[0\textsubscript{III}]{27}
실제로 (\hyperref[0.11.1.1-fr]{11.1.1})의 조건 a), b), c)에 따라
$Z_{r+1}(E_2^{pq})=Z_r(E_2^{pq})$이라는 것은 $d_r^{pq}=0$이라는
것과 동치이고,
$B_r(E_2^{p+r,q-r+1})=B_{r+1}(E_2^{p+r,q-r+1})$이라는 것도
$d_r^{pq}=0$이라는 것과 동치이다. 보조정리는 이 관찰에서 즉시 따른다.
\end{proof}

\subsection{여과 복합체의 스펙트럼 열}
\label{subsection:0.11.2-fr}

\begin{env}[11.2.1]
\label{0.11.2.1-fr}
아벨 범주 $\mathcal C$가 주어졌다고 하자. 미분 연산자의 차수가 $+1$인
$\mathcal C$의 대상들의 \emph{복합체} $(K^i)_{i\in\mathbb Z}$는
$K^\bullet$과 같은 표기로 나타내고, 미분 연산자의 차수가 $-1$인
$\mathcal C$의 대상들의 복합체 $(K_i)_{i\in\mathbb Z}$는 $K_\bullet$과 같은
표기로 나타내기로 한다. 미분 연산자 $d$의 차수가 $+1$인 모든 복합체
$K^\bullet=(K^i)$에 대하여 $K'_i=K^{-i}$로 놓으면 복합체
$K'_\bullet=(K'_i)$를 대응시킬 수 있으며, 이때 미분 연산자
$K'_i\to K'_{i-1}$는 연산자 $d:K^{-i}\to K^{-i+1}$이다. 그 역도
마찬가지이다. 따라서 상황에 따라 어느 한 종류의 복합체를 고찰하고, 한
종류에 대한 모든 결과를 다른 종류에 대한 결과로 옮길 수 있다. 마찬가지로
$K^{\bullet\bullet}=(K^{ij})$ (각각 $K_{\bullet\bullet}=(K_{ij})$)와 같은
표기로 두 미분 연산자의 차수가 모두 $+1$ (각각 $-1$)인 $\mathcal C$의
대상들의 \emph{이중복합체}를 나타낸다. 첨자의 부호를 바꾸면 다시 한
종류에서 다른 종류로 넘어가며, 임의의 다중복합체에도 이와 유사한 표기와
고찰을 적용한다. $K^\bullet$ 또는 $K_\bullet$이라는 표기는 반드시 복합체일
필요는 없는 $\mathbb Z$형의 $\mathcal C$의 \emph{등급 대상}에도 사용한다
(또는 영인 미분 연산자를 주어 이를 복합체로 볼 수 있다). 예를 들어 미분
연산자의 차수가 $+1$인 복합체 $K^\bullet$의 \emph{코호몰로지}를
$\mathrm H^\bullet(K^\bullet)=(\mathrm H^i(K^\bullet))_{i\in\mathbb Z}$로
쓰고, 미분 연산자의 차수가 $-1$인 복합체 $K_\bullet$의
\emph{호몰로지}를
$\mathrm H_\bullet(K_\bullet)=(\mathrm H_i(K_\bullet))_{i\in\mathbb Z}$로
쓴다. 위에서 설명한 연산으로 $K^\bullet$에서 $K'_\bullet$로 넘어가면
$\mathrm H_i(K'_\bullet)=\mathrm H^{-i}(K^\bullet)$이다.

이와 관련하여, 복합체 $K^\bullet$ (각각 $K_\bullet$)에 대하여 일반적으로
$Z^i(K^\bullet)=\operatorname{Ker}(K^i\to K^{i+1})$ (``공순환 대상'')와
$B^i(K^\bullet)=\operatorname{Im}(K^{i-1}\to K^i)$ (``공경계 대상'')
(각각 $Z_i(K_\bullet)=\operatorname{Ker}(K_i\to K_{i-1})$ (``순환 대상'')와
$B_i(K_\bullet)=\operatorname{Im}(K_{i+1}\to K_i)$ (``경계 대상''))로
쓰기로 한다. 따라서
$\mathrm H^i(K^\bullet)=Z^i(K^\bullet)/B^i(K^\bullet)$ (각각
$\mathrm H_i(K_\bullet)=Z_i(K_\bullet)/B_i(K_\bullet)$)이다.

$K^\bullet=(K^i)$ (각각 $K_\bullet=(K_i)$)가 $\mathcal C$의 복합체이고
$T:\mathcal C\to\mathcal C'$가 $\mathcal C$에서 아벨 범주 $\mathcal C'$로
가는 함자이면, $\mathcal C'$의 복합체 $(T(K^i))$ (각각 $(T(K_i))$)를
$T(K^\bullet)$ (각각 $T(K_\bullet)$)로 나타낸다.

$\partial$-함자의 정의(T, 2.1)는 되풀이하지 않는다. 다만 사상
$\partial$가 차수를 하나 낮출 때에는 $\partial^*$-함자 대신 역시
$\partial$-함자라고 부르며, 혼동의 여지가 있을 때에는 매번 문맥으로 이를
명시한다.

마지막으로, $\mathcal C$의 \emph{등급 대상} $(A_i)_{i\in\mathbb Z}$에
대하여 어떤 $i_0$가 존재하여 $i<i_0$ (각각 $i>i_0$)이면 $A_i=0$일 때,
이를 \emph{아래로 유계} (각각 \emph{위로 유계})라고 한다.
\end{env}

\begin{env}[11.2.2]
\label{0.11.2.2-fr}
$K^\bullet$을 미분 연산자 $d$의 차수가 $+1$인 $\mathcal C$의 복합체라
하자. 또한 $K^\bullet$에 그 등급 부분대상들로 이루어진 \emph{여과}
$F(K^\bullet)=(F^p(K^\bullet))_{p\in\mathbb Z}$가 주어졌다고 하자. 다시
말해
\oldpage[0\textsubscript{III}]{28}
$F^p(K^\bullet)=(K^i\cap F^p(K^\bullet))_{i\in\mathbb Z}$이다. 더 나아가
모든 $p\in\mathbb Z$에 대하여
$d(F^p(K^\bullet))\subset F^p(K^\bullet)$라고 가정한다. 이제
$K^\bullet$으로부터 \emph{함자적으로} 스펙트럼 열 $E(K^\bullet)$을
정의하는 방법을 간단히 상기하자(M, XV, 4 및 G, I, 4.3). $r\geq2$일 때,
표준 사상
$F^p(K^\bullet)/F^{p+r}(K^\bullet)\to
F^p(K^\bullet)/F^{p+1}(K^\bullet)$은 코호몰로지에 대한 사상
\[
  \mathrm H^{p+q}(F^p(K^\bullet)/F^{p+r}(K^\bullet))
  \longrightarrow
  \mathrm H^{p+q}(F^p(K^\bullet)/F^{p+1}(K^\bullet))
\]
을 정의한다. 이 사상의 상을 $Z_r^{pq}(K^\bullet)$로 나타낸다. 마찬가지로
완전열
\[
  0\to F^p(K^\bullet)/F^{p+1}(K^\bullet)
  \to F^{p-r+1}(K^\bullet)/F^{p+1}(K^\bullet)
  \to F^{p-r+1}(K^\bullet)/F^p(K^\bullet)\to0
\]
로부터 코호몰로지 완전열에 의해 사상
\[
  \mathrm H^{p+q-1}(F^{p-r+1}(K^\bullet)/F^p(K^\bullet))
  \longrightarrow
  \mathrm H^{p+q}(F^p(K^\bullet)/F^{p+1}(K^\bullet))
\]
을 얻으며, 이 사상의 상을 $B_r^{pq}(K^\bullet)$로 나타낸다. 이때
$B_r^{pq}(K^\bullet)\subset Z_r^{pq}(K^\bullet)$임을 보일 수 있고,
$E_r^{pq}(K^\bullet)=Z_r^{pq}(K^\bullet)/B_r^{pq}(K^\bullet)$로 놓는다.
여기서는 $d_r^{pq}$와 $\alpha_r^{pq}$의 정의를 명시하지 않는다.

여기서 $p$, $q$를 고정하면 모든 $Z_r^{pq}(K^\bullet)$와
$B_r^{pq}(K^\bullet)$는 같은 대상
$\mathrm H^{p+q}(F^p(K^\bullet)/F^{p+1}(K^\bullet))$의 부분대상이며, 이
대상을 $Z_1^{pq}(K^\bullet)$로 쓴다는 점에 유의하자.
$B_1^{pq}(K^\bullet)=0$으로 놓으면 위의 $Z_r^{pq}(K^\bullet)$와
$B_r^{pq}(K^\bullet)$의 정의는 $r=1$에도 적용된다. 또한
$E_1^{pq}(K^\bullet)=Z_1^{pq}(K^\bullet)$로 놓는다.
(\hyperref[0.11.1.1-fr]{11.1.1})의 조건들이 $r=1$에 대해서도 성립하도록
$d_1^{pq}$와 $\alpha_1^{pq}$를 정의한다. 한편 부분대상
$Z_\infty^{pq}(K^\bullet)$는 사상
\[
  \mathrm H^{p+q}(F^p(K^\bullet))\longrightarrow
  \mathrm H^{p+q}(F^p(K^\bullet)/F^{p+1}(K^\bullet))
  =E_1^{pq}(K^\bullet)
\]
의 상으로 정의하고, $B_\infty^{pq}(K^\bullet)$는 사상
\[
  \mathrm H^{p+q-1}(K^\bullet/F^p(K^\bullet))\longrightarrow
  \mathrm H^{p+q}(F^p(K^\bullet)/F^{p+1}(K^\bullet))
  =E_1^{pq}(K^\bullet)
\]
의 상으로 정의한다. 이 사상은 위와 같이 코호몰로지 완전열에서 얻는다.
$Z_\infty(E_2^{pq}(K^\bullet))$와 $B_\infty(E_2^{pq}(K^\bullet))$로는
$Z_\infty^{pq}(K^\bullet)$와 $B_\infty^{pq}(K^\bullet)$의
$E_2^{pq}(K^\bullet)$ 안에서의 표준 상을 취한다.

마지막으로 $F^p(\mathrm H^n(K^\bullet))$는 표준 포함 사상
$F^p(K^\bullet)\to K^\bullet$에서 오는 사상
$\mathrm H^n(F^p(K^\bullet))\to\mathrm H^n(K^\bullet)$의
$\mathrm H^n(K^\bullet)$ 안에서의 상으로 나타낸다. 코호몰로지 완전열에
의하면 이는 또한 사상
$\mathrm H^n(K^\bullet)\to\mathrm H^n(K^\bullet/F^p(K^\bullet))$의
핵이다. 이로써 $E^n(K^\bullet)=\mathrm H^n(K^\bullet)$ 위의 여과를
정의한다. 여기서는 동형사상 $\beta^{pq}$의 정의도 제시하지 않는다.
\end{env}

\begin{env}[11.2.3]
\label{0.11.2.3-fr}
$E(K^\bullet)$의 \emph{함자성}은 다음과 같은 뜻으로 이해해야 한다.
$\mathcal C$의 두 \emph{여과 복합체} $K^\bullet$,
$K^{\prime\bullet}$과 \emph{여과와 양립하는} 복합체 사상
$u:K^\bullet\to K^{\prime\bullet}$이 주어지면, 이로부터 자명한 방식으로
사상 $u_r^{pq}$ ($r\geq1$)와 $u^n$을 얻는다. 이 사상들이
(\hyperref[0.11.1.2-fr]{11.1.2})의 의미에서 $d_r^{pq}$,
$\alpha_r^{pq}$ 및 $\beta^{pq}$와 양립함을 보일 수 있으므로, 이들은
실제로 스펙트럼 열의 사상
$E(u):E(K^\bullet)\to E(K^{\prime\bullet})$을 정의한다. 더 나아가 $u$와
$v$가 위와 같은 종류의 $K^\bullet\to K^{\prime\bullet}$ 사상들이고
\emph{차수 $\leq k$로 호모토픽}이면, $r>k$에 대하여
$u_r^{pq}=v_r^{pq}$이고 모든 $n$에 대하여 $u^n=v^n$이다(M, XV, 3.1).
\end{env}

\begin{env}[11.2.4]
\label{0.11.2.4-fr}
\oldpage[0\textsubscript{III}]{29}
$\mathcal C$에서 여과 귀납극한들이 완전하다고 가정하자. 그러면
$K^\bullet$의 여과 $(F^p(K^\bullet))$가 \emph{소진적}이면 모든 $n$에
대하여 여과 $(F^p(\mathrm H^n(K^\bullet)))$도 소진적이다. 실제로 가정에
의해 $K^\bullet=\varinjlim_{p\to-\infty}F^p(K^\bullet)$이고,
$\mathcal C$에 대한 가정에 의해 코호몰로지는 귀납극한과 교환한다. 같은
이유로
$B_\infty(E_2^{pq}(K^\bullet))=\sup_k B_k(E_2^{pq}(K^\bullet))$이다.
모든 $n$에 대하여 어떤 정수 $u(n)$가 존재하여 $p>u(n)$이면
$\mathrm H^n(F^p(K^\bullet))=0$일 때, $K^\bullet$의 여과
$(F^p(K^\bullet))$를 \emph{정칙}이라고 한다. 특히 $K^\bullet$의 여과가
\emph{이산}이면 그러하다. $K^\bullet$의 여과가 정칙이고 소진적이며
$\mathcal C$에서 여과 귀납극한들이 완전하면, 스펙트럼 열
$E(K^\bullet)$이 \emph{정칙}임을 보일 수 있다(M, XV, 4).
\end{env}

\subsection{이중복합체의 스펙트럼 열}
\label{subsection:0.11.3-fr}

\begin{env}[11.3.1]
\label{0.11.3.1-fr}
이중복합체에 관한 규약으로는 (M)의 규약보다 (T, 2.4)의 규약을 따른다.
따라서 그러한 이중복합체 $K^{\bullet\bullet}=(K^{ij})$의 두 미분 연산자
$d'$, $d''$ (차수 $+1$)는 \emph{서로 교환한다}고 가정한다. 다음 두 조건
중 하나가 성립한다고 하자.
\begin{enumerate}
  \item[$1^\circ$] $\mathcal C$에서 무한 직합들이 존재한다.
  \item[$2^\circ$] 모든 $n\in\mathbb Z$에 대하여 $p+q=n$이고
    $K^{pq}\ne0$인 쌍 $(p,q)$는 \emph{유한}개뿐이다.
\end{enumerate}
그러면 이중복합체 $K^{\bullet\bullet}$는
$K^{\prime n}=\sum_{i+j=n}K^{ij}$인 (전체) 복합체
$(K^{\prime n})_{n\in\mathbb Z}$를 정의한다. 이 복합체의 미분 연산자
$d$ (차수 $+1$)는 $x\in K^{ij}$에 대하여
$dx=d'x+(-1)^i d''x$로 주어진다. 이하에서 \emph{이중복합체
$K^{\bullet\bullet}$가 정의하는 (전체) 복합체}를 말할 때에는 언제나 위
조건들 중 하나가 성립한다고 암묵적으로 가정한다. 다중복합체에 대해서도
이와 유사한 규약을 채택한다.

$K^{i,\bullet}$ (각각 $K^{\bullet,j}$)로 전체 복합체
$(K^{ij})_{j\in\mathbb Z}$ (각각 $(K^{ij})_{i\in\mathbb Z}$)를 나타내고,
$Z_{\mathrm{II}}^p(K^{i,\bullet})$, $B_{\mathrm{II}}^p(K^{i,\bullet})$,
$\mathrm H_{\mathrm{II}}^p(K^{i,\bullet})$ (각각
$Z_{\mathrm I}^p(K^{\bullet,j})$, $B_{\mathrm I}^p(K^{\bullet,j})$,
$\mathrm H_{\mathrm I}^p(K^{\bullet,j})$)로 각각 그 $p$번째 공순환 대상,
공경계 대상 및 코호몰로지 대상을 나타낸다. 미분
$d':K^{i,\bullet}\to K^{i+1,\bullet}$은 복합체의 사상이므로 공순환,
공경계 및 코호몰로지 위에 다음 연산자들을 준다.
\[
\begin{aligned}
  d'&:Z_{\mathrm{II}}^p(K^{i,\bullet})
      \longrightarrow Z_{\mathrm{II}}^p(K^{i+1,\bullet}),\\
  d'&:B_{\mathrm{II}}^p(K^{i,\bullet})
      \longrightarrow B_{\mathrm{II}}^p(K^{i+1,\bullet}),\\
  d'&:\mathrm H_{\mathrm{II}}^p(K^{i,\bullet})
      \longrightarrow \mathrm H_{\mathrm{II}}^p(K^{i+1,\bullet}).
\end{aligned}
\]
이 연산자들에 대하여
$(Z_{\mathrm{II}}^p(K^{i,\bullet}))_{i\in\mathbb Z}$,
$(B_{\mathrm{II}}^p(K^{i,\bullet}))_{i\in\mathbb Z}$ 및
$(\mathrm H_{\mathrm{II}}^p(K^{i,\bullet}))_{i\in\mathbb Z}$가
복합체임은 자명하다. 복합체
$(\mathrm H_{\mathrm{II}}^p(K^{i,\bullet}))_{i\in\mathbb Z}$를
$\mathrm H_{\mathrm{II}}^p(K^{\bullet\bullet})$로 나타내고, 그 $q$번째
공순환 대상, 공경계 대상 및 코호몰로지 대상을 각각
$Z_{\mathrm I}^q(\mathrm H_{\mathrm{II}}^p(K^{\bullet\bullet}))$,
$B_{\mathrm I}^q(\mathrm H_{\mathrm{II}}^p(K^{\bullet\bullet}))$ 및
$\mathrm H_{\mathrm I}^q(\mathrm H_{\mathrm{II}}^p(K^{\bullet\bullet}))$로
나타낸다. 마찬가지로 복합체
$\mathrm H_{\mathrm I}^p(K^{\bullet\bullet})$와 그 코호몰로지 대상
$\mathrm H_{\mathrm{II}}^q(\mathrm H_{\mathrm I}^p(K^{\bullet\bullet}))$를
정의한다. 한편 $\mathrm H^n(K^{\bullet\bullet})$는
$K^{\bullet\bullet}$가 정의하는 (전체) 복합체의 $n$번째 코호몰로지
대상을 나타냄을 상기하자.
\end{env}

\begin{env}[11.3.2]
\label{0.11.3.2-fr}
이중복합체 $K^{\bullet\bullet}$가 정의하는 복합체 위에는 다음과 같이
주어지는 두 표준 여과 $(F_{\mathrm I}^p(K^{\bullet\bullet}))$와
$(F_{\mathrm{II}}^p(K^{\bullet\bullet}))$를 고찰할 수 있다.
\[
\label{0.11.3.2.1-fr}
  F_{\mathrm I}^p(K^{\bullet\bullet})
  =\left(\sum_{\substack{i+j=n\\i\geq p}}K^{ij}\right)_{n\in\mathbb Z},
  \qquad
  F_{\mathrm{II}}^p(K^{\bullet\bullet})
  =\left(\sum_{\substack{i+j=n\\j\geq p}}K^{ij}\right)_{n\in\mathbb Z}.
  \tag{11.3.2.1}
\]
\oldpage[0\textsubscript{III}]{30}
정의에 따라 이들은 실제로 $K^{\bullet\bullet}$가 정의하는 (전체)
복합체의 등급 부분대상들이므로 이 복합체를 여과 복합체로 만든다. 또한 이
여과들이 소진적이고 분리되어 있음은 자명하다.

각 여과에는 스펙트럼 열
(\hyperref[0.11.2.2-fr]{11.2.2})이 대응한다.
$(F_{\mathrm I}^p(K^{\bullet\bullet}))$와
$(F_{\mathrm{II}}^p(K^{\bullet\bullet}))$에 각각 대응하는 스펙트럼 열을
$'E(K^{\bullet\bullet})$와 $''E(K^{\bullet\bullet})$로 나타내며, 이들을
$K^{\bullet\bullet}$의 \emph{스펙트럼 열}이라 한다. 두 열은 모두
코호몰로지 $(\mathrm H^n(K^{\bullet\bullet}))$를 수렴 대상으로 갖는다.
더 나아가 다음이 성립함을 보일 수 있다(M, XV, 6).
\[
\label{0.11.3.2.2-fr}
  'E_2^{pq}(K^{\bullet\bullet})
  =\mathrm H_{\mathrm I}^p(\mathrm H_{\mathrm{II}}^q(K^{\bullet\bullet})),
  \qquad
  ''E_2^{pq}(K^{\bullet\bullet})
  =\mathrm H_{\mathrm{II}}^q(\mathrm H_{\mathrm I}^p(K^{\bullet\bullet})).
  \tag{11.3.2.2}
\]

이중복합체의 모든 사상
$u:K^{\bullet\bullet}\to K^{\prime\bullet\bullet}$은 그 자체로
$K^{\bullet\bullet}$와 $K^{\prime\bullet\bullet}$의 같은 종류의 여과와
양립하므로, 두 스펙트럼 열 각각에 대한 사상을 정의한다. 더 나아가 호모토픽인
두 사상은 대응하는 여과 (전체) 복합체들의 \emph{차수 $\leq1$인
호모토피}를 정의하므로, 두 스펙트럼 열 각각에 대하여 \emph{같은} 사상을
정의한다(M, XV, 6.1).
\end{env}

\begin{proposition}[11.3.3]
\label{0.11.3.3-fr}
$K^{\bullet\bullet}=(K^{ij})$를 아벨 범주 $\mathcal C$의 이중복합체라
하자.
\begin{enumerate}
  \item[(i)] 어떤 $i_0$와 $j_0$가 존재하여 $i<i_0$ 또는 $j<j_0$이면
    $K^{ij}=0$일 때(각각 $i>i_0$ 또는 $j>j_0$이면 $K^{ij}=0$일 때), 두
    스펙트럼 열 $'E(K^{\bullet\bullet})$와
    $''E(K^{\bullet\bullet})$는 쌍정칙이다.
  \item[(ii)] 어떤 $i_0$와 $i_1$이 존재하여 $i<i_0$ 또는 $i>i_1$이면
    $K^{ij}=0$일 때(각각 어떤 $j_0$와 $j_1$이 존재하여 $j<j_0$ 또는
    $j>j_1$이면 $K^{ij}=0$일 때), 두 스펙트럼 열
    $'E(K^{\bullet\bullet})$와 $''E(K^{\bullet\bullet})$는 쌍정칙이다.
\end{enumerate}
더 나아가 $\mathcal C$에서 여과 귀납극한들이 존재하고 완전하다고 하자.
그러면 다음이 성립한다.
\begin{enumerate}
  \item[(iii)] 어떤 $i_0$가 존재하여 $i>i_0$이면 $K^{ij}=0$일
    때(각각 어떤 $j_0$가 존재하여 $j<j_0$이면 $K^{ij}=0$일 때), 열
    $'E(K^{\bullet\bullet})$은 정칙이다.
  \item[(iv)] 어떤 $i_0$가 존재하여 $i<i_0$이면 $K^{ij}=0$일
    때(각각 어떤 $j_0$가 존재하여 $j>j_0$이면 $K^{ij}=0$일 때), 열
    $''E(K^{\bullet\bullet})$은 정칙이다.
\end{enumerate}
\end{proposition}

\begin{proof}
이 명제는 정의 (\hyperref[0.11.1.3-fr]{11.1.3})과
(\hyperref[0.11.2.4-fr]{11.2.4}) 및 여과 $F_{\mathrm I}$에 관한 다음
관찰들에서 곧바로 따른다($K^{\bullet\bullet}$의 두 첨자의 역할을 바꾸면
$F_{\mathrm{II}}$에 관한 유사한 관찰들을 얻는다).
\begin{enumerate}
  \item[$1^\circ$] 어떤 $i_0$가 존재하여 $i>i_0$이면 $K^{ij}=0$일 때,
    여과 $F_{\mathrm I}(K^{\bullet\bullet})$는 이산이다.
  \item[$2^\circ$] 어떤 $i_0$가 존재하여 $i<i_0$이면 $K^{ij}=0$일 때,
    여과 $F_{\mathrm I}(K^{\bullet\bullet})$는 쌍대 이산이다. 이로부터 모든
    $n$에 대하여 대응하는 여과
    $F_{\mathrm I}(\mathrm H^n(K^{\bullet\bullet}))$도 쌍대 이산임을
    곧바로 얻는다. 더 나아가 여과
    $F_{\mathrm I}(K^{\bullet\bullet})$에 대응하는 $B_r^{pq}$의 정의
    (\hyperref[0.11.2.2-fr]{11.2.2})에 의하면 모든 쌍 $(p,q)$에 대하여
    열 $(B_r^{pq})_{r\geq2}$는 정상적이다.
  \item[$3^\circ$] 어떤 $j_0$가 존재하여 $j<j_0$이면 $K^{ij}=0$일 때,
    $p+r+j_0>n$이면
    \[
      F_{\mathrm I}^{p+r}(K^{\bullet\bullet})
      \cap\left(\sum_{i+j=n}K^{ij}\right)=0
    \]
    이다. 따라서 $r>q-j_0+1$이면
    $Z_r^{pq}=Z_\infty(E_2^{pq})$이다. 한편 $p>n-j_0+1$이면
    $\mathrm H^n(F_{\mathrm I}^p(K^{\bullet\bullet}))=0$이다.
  \item[$4^\circ$] 어떤 $j_0$가 존재하여 $j>j_0$이면 $K^{ij}=0$일 때,
    $p-r+1+j_0<n$이면
    \[
      F_{\mathrm I}^{p-r+1}(K^{\bullet\bullet})
      \cap\left(\sum_{i+j=n}K^{ij}\right)=\sum_{i+j=n}K^{ij}.
    \]
    \oldpage[0\textsubscript{III}]{31}
    따라서 $r>j_0-q+1$이면
    $B_r^{pq}=B_\infty(E_2^{pq})$이다. 한편 $p+j_0<n-1$이면
    $\mathrm H^n(F_{\mathrm I}^p(K^{\bullet\bullet}))
    =\mathrm H^n(K^{\bullet\bullet})$이다.
\end{enumerate}
\end{proof}

\begin{env}[11.3.4]
\label{0.11.3.4-fr}
이중복합체 $K^{\bullet\bullet}=(K^{ij})$가 $i<0$ 또는 $j<0$이면
$K^{ij}=0$이라고 하자. 그러면 모든 $p\in\mathbb Z$에 대하여 표준
«~가장자리 준동형~»
\[
\label{0.11.3.4.1-fr}
  'E_2^{p0}(K^{\bullet\bullet})\longrightarrow
  \mathrm H^p(K^{\bullet\bullet})
  \tag{11.3.4.1}
\]
을 정의할 수 있음이 알려져 있다(M, XV, 6). 그 이유를 간단히 상기하자.
한편으로 스펙트럼 열 $'E(K^{\bullet\bullet})$에서
$2\leq r\leq+\infty$일 때
$Z_r^{p0}=Z_{\mathrm I}^p(Z_{\mathrm{II}}^0(K^{\bullet\bullet}))$이고,
다른 한편으로
$\mathrm H^p(F_{\mathrm I}^{p+1}(K^{\bullet\bullet}))=0$이다. 따라서
동형사상
$\beta^{p0}:'E_\infty^{p0}\xrightarrow{\sim}
\mathrm H^p(F_{\mathrm I}^p)/\mathrm H^p(F_{\mathrm I}^{p+1})$은
준동형
$'E_\infty^{p0}\to\mathrm H^p(F_{\mathrm I}^p(K^{\bullet\bullet}))
\to\mathrm H^p(K^{\bullet\bullet})$을 준다. 모든 $Z_r^{p0}$가 같으므로
$r\leq s$에 대하여 표준 준동형
$'E_r^{p0}\to'E_s^{p0}$을 정의할 수 있고, 특히 준동형
$'E_2^{p0}\to'E_\infty^{p0}$을 정의할 수 있다. 이를 합성하여 가장자리
준동형 $'E_2^{p0}\to\mathrm H^p(K^{\bullet\bullet})$을 얻는다. 더 나아가
$d'z=0$을 만족하는 원소
$z\in Z_{\mathrm{II}}^0(K^{\bullet\bullet})\subset K^{p0}$의
법 $B_2^{p0}$인 유에 위에서 정의한 가장자리 준동형을 적용하면,
$'E_\infty^{p0}$에서 $z$의 법 $B_\infty^{p0}$인 유를 얻고, 이어서
$\mathrm H^p(K^{\bullet\bullet})$에서 $z$의 코호몰로지류를 얻음을 곧바로
확인할 수 있다. 그러므로 끝으로 \emph{가장자리 준동형
(\hyperref[0.11.3.4.1-fr]{11.3.4.1})은 표준 포함 사상}
$Z_{\mathrm{II}}^0(K^{\bullet\bullet})\to K^{\bullet\bullet}$
\emph{에 코호몰로지를 취하여 얻어진다}($K^{\bullet\bullet}$는 전체
복합체로 본다). 마찬가지로 가장자리 준동형
\[
\label{0.11.3.4.2-fr}
  ''E_2^{p0}(K^{\bullet\bullet})\longrightarrow
  \mathrm H^p(K^{\bullet\bullet})
  \tag{11.3.4.2}
\]
은 자연스럽게 표준 포함 사상
$Z_{\mathrm I}^0(K^{\bullet\bullet})\to K^{\bullet\bullet}$에서
얻어지는 것으로 해석된다.
\end{env}

\begin{env}[11.3.5]
\label{0.11.3.5-fr}
이제 $K_{\bullet\bullet}=(K_{ij})$를 두 미분 연산자의 차수가 모두
$-1$인 $\mathcal C$의 이중복합체라 하자. 이때 전체 복합체
$(K_{ij})_{j\in\mathbb Z}$ (각각 $(K_{ij})_{i\in\mathbb Z}$)를
$K_{i,\bullet}$ (각각 $K_{\bullet,j}$)로 쓰고, 그 $p$번째 호몰로지 대상을
$\mathrm H_p^{\mathrm{II}}(K_{i,\bullet})$ (각각
$\mathrm H_p^{\mathrm I}(K_{\bullet,j})$)로 쓴다. 복합체
$(\mathrm H_p^{\mathrm{II}}(K_{i,\bullet}))_{i\in\mathbb Z}$ (각각
$(\mathrm H_p^{\mathrm I}(K_{\bullet,j}))_{j\in\mathbb Z}$)은
$\mathrm H_p^{\mathrm{II}}(K_{\bullet\bullet})$ (각각
$\mathrm H_p^{\mathrm I}(K_{\bullet\bullet})$)로 쓰고, 그 $q$번째
호몰로지 대상은
$\mathrm H_q^{\mathrm I}(\mathrm H_p^{\mathrm{II}}(K_{\bullet\bullet}))$
(각각
$\mathrm H_q^{\mathrm{II}}(\mathrm H_p^{\mathrm I}(K_{\bullet\bullet}))$)
로 쓴다. 순환 대상과 경계 대상에도 이와 유사한 표기를 쓴다. 마지막으로
$\mathrm H_n(K_{\bullet\bullet})$는 $K_{\bullet\bullet}$가 정의하는
(차수가 $-1$인 미분 연산자를 갖는) 전체 복합체의 $n$번째 호몰로지 대상이
존재할 때 이를 나타낸다.

$K^{\prime ij}=K_{-i,-j}$로 놓고, $K_{\bullet\bullet}$에 대응하는 차수
$+1$의 미분 연산자들을 갖는 이중복합체를
$K^{\prime\bullet\bullet}=(K^{\prime ij})$라 하자. 정의에 따라
$K_{\bullet\bullet}$의 \emph{스펙트럼 열}은
$K^{\prime\bullet\bullet}$의 스펙트럼 열이며, 이를
$'E(K_{\bullet\bullet})$와 $''E(K_{\bullet\bullet})$로 쓴다. 다만 표기는
다음과 같이 바꾼다.
\[
  'E_{pq}^r(K_{\bullet\bullet})
  ='E_r^{-p,-q}(K^{\prime\bullet\bullet}),\qquad
  ''E_{pq}^r(K_{\bullet\bullet})
  =''E_r^{-p,-q}(K^{\prime\bullet\bullet})
\]
$2\leq r\leq\infty$에 대하여 이 표기를 쓰면 다음을 얻는다.
\[
  'E_{pq}^2(K_{\bullet\bullet})
  =\mathrm H_p^{\mathrm I}(\mathrm H_q^{\mathrm{II}}(K_{\bullet\bullet})),
  \qquad
  ''E_{pq}^2(K_{\bullet\bullet})
  =\mathrm H_q^{\mathrm{II}}(\mathrm H_p^{\mathrm I}(K_{\bullet\bullet})).
\]
부호 오류를 피하려면 일반적으로 이 스펙트럼 열들과 그 수렴 대상 사이의
관계는 복합체 $K^{\prime\bullet\bullet}$로 돌아가 고찰하는 편이 낫다.
다만 (\hyperref[0.11.3.3-fr]{11.3.3})에 대응하는 판정법은 다음과 같다.
\end{env}

\begin{env}[11.3.6]
\label{0.11.3.6-fr}
스펙트럼 열 $'E(K_{\bullet\bullet})$와
$''E(K_{\bullet\bullet})$는 다음 경우에 \emph{쌍정칙}이다.
\begin{enumerate}
  \item[(a)] 어떤 $i_0$와 $j_0$가 존재하여 $i>i_0$ 또는 $j>j_0$이면
    $K_{ij}=0$이다(각각 $i<i_0$
    \oldpage[0\textsubscript{III}]{32}
    또는 $j<j_0$이면 $K_{ij}=0$이다).
  \item[(b)] 어떤 $i_0$와 $i_1$이 존재하여 $i<i_0$이고 $i>i_1$이면
    $K_{ij}=0$이다.
  \item[(c)] 어떤 $j_0$와 $j_1$이 존재하여 $j<j_0$이고 $j>j_1$이면
    $K_{ij}=0$이다.
\end{enumerate}

어떤 $i_0$가 존재하여 $i<i_0$이면 $K_{ij}=0$이거나, 어떤 $j_0$가
존재하여 $j>j_0$이면 $K_{ij}=0$일 때, 열
$'E(K_{\bullet\bullet})$은 정칙이다.

어떤 $i_0$가 존재하여 $i>i_0$이면 $K_{ij}=0$이거나, 어떤 $j_0$가
존재하여 $j<j_0$이면 $K_{ij}=0$일 때, 열
$''E(K_{\bullet\bullet})$은 정칙이다.
\end{env}

\subsection{복합체 $K^\bullet$에 관한 함자의 초코호몰로지}
\label{subsection:0.11.4-fr}

\begin{env}[11.4.1]
\label{0.11.4.1-fr}
$\mathcal C$를 아벨 범주라 하자. $\mathcal C$의 대상 $A$의
\emph{오른쪽 분해}(또는 \emph{코호몰로지 분해})란, 미분작용소의 차수가
$+1$인 $\mathcal C$의 대상들의 복합체
\[
  0\longrightarrow L^0\longrightarrow L^1\longrightarrow L^2
  \longrightarrow\cdots
\]
와, 분해의 \emph{증대}라 불리는 사상 $\varepsilon:A\to L^0$로
이루어지며, 이 증대는 복합체의 사상
\[
\xymatrix{
  0\ar[r] & A\ar[r]\ar[d]_{\varepsilon} & 0\ar[r]\ar[d]
    & 0\ar[r]\ar[d] & \cdots\\
  0\ar[r] & L^0\ar[r] & L^1\ar[r] & L^2\ar[r] & \cdots
})
\]
로 볼 수 있고, 열
\[
  0\longrightarrow A\xrightarrow{\varepsilon}L^0\longrightarrow L^1
  \longrightarrow\cdots
\]
이 완전한 것을 말한다. 마찬가지로 $A$의 \emph{왼쪽 분해}(또는
\emph{호몰로지 분해})란, 미분작용소의 차수가 $-1$인 $\mathcal C$의
대상들의 복합체 $0\leftarrow L_0\leftarrow L_1\leftarrow\cdots$와
증대 $\varepsilon:L_0\to A$로 이루어지며, 열
\[
  0\longleftarrow A\xleftarrow{\varepsilon}L_0\longleftarrow L_1
  \longleftarrow\cdots
\]
이 완전한 것을 말한다.

대상 $A$의 오른쪽 분해 $(L^i)_{i\geq0}$가 $i\geq n+1$일 때
$L^i=0$을 만족하면, 이 분해의 \emph{길이는 $\leq n$}이라고 한다.
길이가 $\leq n$인 왼쪽 분해도 같은 방식으로 정의한다. 어떤 정수
$n$에 대하여 길이가 $\leq n$인 분해를 \emph{유한} 분해라 한다.

$A$의 분해를 이루는 $A$ 이외의 $\mathcal C$의 대상들이
\emph{사영 대상}(각각 \emph{단사 대상})이면 그 분해를
\emph{사영 분해}(각각 \emph{단사 분해})라 한다. $\mathcal C$가 한
환 위의 가군들(예를 들어 왼쪽 가군들)의 범주인 경우, 분해를 이루는
$A$ 이외의 가군들이 \emph{평탄}(각각 \emph{자유})이면 마찬가지로
$A$의 분해를 \emph{평탄 분해}(각각 \emph{자유 분해})라 한다.
\end{env}

\begin{env}[11.4.2]
\label{0.11.4.2-fr}
$K^\bullet=(K^i)_{i\in\mathbb Z}$를 미분작용소의 차수가 $+1$인
$\mathcal C$의 대상들의 복합체라 하자.

$K^\bullet$의 \emph{오른쪽 카르탕--아일렌베르크 분해}란,
미분작용소들의 차수가 $+1$이고 $j<0$일 때 $L^{ij}=0$인 이중복합체
$L^{\bullet\bullet}=(L^{ij})$와 단순복합체의 사상
$\varepsilon:K^\bullet\to L^{\bullet,0}$의 쌍으로서, 다음 조건들을
만족하는 것을 말한다.
\oldpage[0\textsubscript{III}]{33}
\begin{enumerate}
  \item[(i)] 각 지표 $i$에 대하여 열들
    \[
    \begin{aligned}
      0\longrightarrow{}&K^i\xrightarrow{\varepsilon}L^{i0}
        \longrightarrow L^{i1}\longrightarrow\cdots\\
      0\longrightarrow{}&\mathrm B^i(K^\bullet)\xrightarrow{\varepsilon}
        \mathrm B_{\mathrm I}^i(L^{\bullet,0})
        \longrightarrow\mathrm B_{\mathrm I}^i(L^{\bullet,1})
        \longrightarrow\cdots\\
      0\longrightarrow{}&\mathrm Z^i(K^\bullet)\xrightarrow{\varepsilon}
        \mathrm Z_{\mathrm I}^i(L^{\bullet,0})
        \longrightarrow\mathrm Z_{\mathrm I}^i(L^{\bullet,1})
        \longrightarrow\cdots\\
      0\longrightarrow{}&\mathrm H^i(K^\bullet)\xrightarrow{\varepsilon}
        \mathrm H_{\mathrm I}^i(L^{\bullet,0})
        \longrightarrow\mathrm H_{\mathrm I}^i(L^{\bullet,1})
        \longrightarrow\cdots
    \end{aligned}
    \]
    은 완전하다. 다시 말해 $(L^{i,\bullet})$,
    $(\mathrm B_{\mathrm I}^i(L^{\bullet,\bullet}))$,
    $(\mathrm Z_{\mathrm I}^i(L^{\bullet,\bullet}))$,
    $(\mathrm H_{\mathrm I}^i(L^{\bullet,\bullet}))$은 각각 $K^i$,
    $\mathrm B^i(K^\bullet)$, $\mathrm Z^i(K^\bullet)$,
    $\mathrm H^i(K^\bullet)$의 \emph{분해}이다.
  \item[(ii)] 각 $j$에 대하여 단순복합체 $L^{\bullet,j}$는
    \emph{분할}된다. 다시 말해 완전열들
    \[
      \label{0.11.4.2.1-fr}
      0\longrightarrow\mathrm B_{\mathrm I}^i(L^{\bullet,j})
      \longrightarrow\mathrm Z_{\mathrm I}^i(L^{\bullet,j})
      \longrightarrow\mathrm H_{\mathrm I}^i(L^{\bullet,j})
      \longrightarrow0
      \tag{11.4.2.1}
    \]
    \[
      \label{0.11.4.2.2-fr}
      0\longrightarrow\mathrm Z_{\mathrm I}^i(L^{\bullet,j})
      \longrightarrow L^{ij}
      \longrightarrow\mathrm B_{\mathrm I}^{i+1}(L^{\bullet,j})
      \longrightarrow0
      \tag{11.4.2.2}
    \]
    은 분할된다.
\end{enumerate}

$\mathcal C$의 모든 대상이 어떤 단사 대상의 부분대상이면,
$\mathcal C$의 모든 복합체 $K^\bullet$은 \emph{단사} 카르탕--아일렌베르크
분해, 즉 단사 대상들 $L^{ij}$로 이루어진 분해를 가짐이 증명된다
(M, XVII, 1.2). 이때 위의 조건 (ii)에 따라
$\mathrm B_{\mathrm I}^i(L^{\bullet,j})$,
$\mathrm Z_{\mathrm I}^i(L^{\bullet,j})$,
$\mathrm H_{\mathrm I}^i(L^{\bullet,j})$도 단사 대상이다. 더 나아가,
$\mathcal C$의 복합체의 임의의 사상 $f:K^\bullet\to K'^\bullet$,
$K^\bullet$의 임의의 카르탕--아일렌베르크 분해
$L^{\bullet\bullet}$, 그리고 $K'^\bullet$의 임의의 단사
카르탕--아일렌베르크 분해 $L'^{\bullet\bullet}$에 대하여, $f$ 및
증대들과 양립하는 이중복합체의 사상
$F:L^{\bullet\bullet}\to L'^{\bullet\bullet}$가 존재한다. 또한 $f$와
$g$가 $K^\bullet$에서 $K'^\bullet$로 가는 서로 호모토픽한 두 사상이면,
이에 대응하는 $L^{\bullet\bullet}$에서 $L'^{\bullet\bullet}$로 가는
사상들도 서로 호모토픽하다(\emph{같은 곳}).

$K^\bullet$이 \emph{아래로 유계}(각각 \emph{위로 유계})이고 $i<i_0$
(각각 $i>i_0$)일 때 $K^i=0$이면, $i<i_0$ (각각 $i>i_0$)일 때
$L^{ij}=0$이 되도록 $L^{\bullet\bullet}$을 택할 수 있다
(M, XVII, 1.3).

한편 $\mathcal C$의 모든 대상이 길이 $\leq n$인 \emph{단사 분해}를
갖게 하는 정수 $n$이 존재한다고 하자. 그러면 $j>n$일 때
$L^{ij}=0$이라고 가정할 수 있다 (M, XVII, 1.4).
\end{env}

\begin{env}[11.4.3]
\label{0.11.4.3-fr}
이제 $T$를 $\mathcal C$에서 아벨 범주 $\mathcal C'$으로 가는
\emph{가법 공변 함자}라 하자. $\mathcal C$의 복합체 $K^\bullet$과
그 단사 카르탕--아일렌베르크 분해 $L^{\bullet\bullet}$이 주어졌다고
하고, 이중복합체 $T(L^{\bullet\bullet})$이 정의하는 (단순)복합체가
존재한다고 가정하자
(\hyperref[0.11.3.1-fr]{11.3.1} 참조). 이 이중복합체의 두 스펙트럼 열
$'E(T(L^{\bullet\bullet}))$와 $''E(T(L^{\bullet\bullet}))$을
\emph{복합체 $K^\bullet$에 관한 $T$의 초코호몰로지 스펙트럼 열}이라
한다. (\hyperref[0.11.4.2-fr]{11.4.2})와
(\hyperref[0.11.3.2-fr]{11.3.2})에 따라 이들은 실제로 선택한 단사
카르탕--아일렌베르크 분해 $L^{\bullet\bullet}$이 아니라
$K^\bullet$에만 의존하며, 더 나아가 $K^\bullet$에 \emph{함자적으로}
의존한다. 두 열은 공통 수렴대상
$\mathrm H^\bullet(T(L^{\bullet\bullet}))$을 가지며, 이 또한
\emph{복합체 $K^\bullet$에 관한 $T$의 초코호몰로지}라 하고
$\mathrm R^\bullet T(K^\bullet)$로 나타낸다. 앞의 두 스펙트럼 열의
$E_2$항은 다음과 같음이 증명된다.
\[
  \label{0.11.4.3.1-fr}
  'E_2^{pq}=\mathrm H^p(\mathrm R^qT(K^\bullet))
  \tag{11.4.3.1}
\]
\[
  \label{0.11.4.3.2-fr}
  ''E_2^{pq}=\mathrm R^pT(\mathrm H^q(K^\bullet))
  \tag{11.4.3.2}
\]
\oldpage[0\textsubscript{III}]{34}
여기서 $p\in\mathbb Z$에 대하여 $\mathrm R^pT$는 통상과 같이 $T$의
$p$번째 \emph{유도함자}이고, $\mathrm R^qT(K^\bullet)$는 복합체
$(\mathrm R^qT(K^i))_{i\in\mathbb Z}$를 뜻한다.

\phantomsection\label{0.11.4.4-fr}
달리 명시하지 않는 한, 이제부터 $\mathcal C$의 모든 대상이
$\mathcal C$의 어떤 단사 대상의 부분대상이라고 가정한다. 따라서
$\mathcal C$의 모든 복합체에는 단사 카르탕--아일렌베르크 분해가
존재한다. $j<0$일 때 $L^{ij}=0$이므로
(\hyperref[0.11.3.3-fr]{11.3.3})의 판정법에 따라 다음 두 경우에는
$K^\bullet$에 관한 $T$의 초코호몰로지 스펙트럼 열 \emph{둘 다} 존재하며
\emph{쌍정칙}이다.
\begin{enumerate}
  \item[$1^\circ$] $K^\bullet$이 아래로 유계이다.
  \item[$2^\circ$] $\mathcal C$의 모든 대상이, 고려하는 대상과 무관한
    어떤 정수 $n$ 이하의 길이를 갖는 단사 분해를 가진다.
\end{enumerate}
실제로 첫째 경우에는 (\hyperref[0.11.4.2-fr]{11.4.2})에 따라
$i<i_0$일 때 $L^{ij}=0$이 되게 하는 $i_0$가 존재한다고 가정할 수
있고, 둘째 경우에는 $j>j_1$일 때 $L^{ij}=0$이 되게 하는 $j_1$이
존재한다고 가정할 수 있다. 두 경우 모두 주어진 $n$에 대하여
$L^{ij}\neq0$이고 $i+j=n$인 쌍 $(i,j)$는 유한 개뿐임도 명백하므로,
주장이 성립한다.

$\mathcal C'$에서 여과 귀납극한들이 존재하고 완전하다고 가정하면
(특히 $\mathcal C'$에서 무한 직합들이 존재한다), 이중복합체
$T(L^{\bullet\bullet})$이 정의하는 복합체가 존재하고,
(\hyperref[0.11.3.3-fr]{11.3.3})의 판정법에 따라 스펙트럼 열
$'E(T(L^{\bullet\bullet}))$은 항상 \emph{정칙}이다.

$K^\bullet$이 단 하나의 항 $K^{i_0}$을 제외하고 모든 항 $K^i$가
0인 복합체이면, $\mathrm R^nT(K^\bullet)$은
$\mathrm R^{n-i_0}T(K^{i_0})$와 동형이다. 실제로 $i\neq i_0$일 때
$L^{ij}=0$인 카르탕--아일렌베르크 분해 $L^{\bullet\bullet}$을 취하면
정의에서 곧바로 따른다.

$K^\bullet$, $K'^\bullet$이 $\mathcal C$의 두 복합체이고 $f$, $g$가
$K^\bullet$에서 $K'^\bullet$로 가는 서로 \emph{호모토픽한} 두
사상이면, $f$와 $g$로부터 얻는 사상
$\mathrm R^\bullet T(K^\bullet)\to
\mathrm R^\bullet T(K'^\bullet)$은 서로 같으며, 코호몰로지 스펙트럼
열의 사상들도 마찬가지이다.
\end{env}

\begin{proposition}[11.4.5]
\label{0.11.4.5-fr}
$\mathcal C'$에서 여과 귀납극한들이 존재하고 완전하다고 가정하자.
모든 $n>0$과 모든 $i\in\mathbb Z$에 대하여
$\mathrm R^nT(K^i)=0$이면, 함자적 동형
\[
  \label{0.11.4.5.1-fr}
  \mathrm R^iT(K^\bullet)\xrightarrow{\sim}
  \mathrm H^i(T(K^\bullet))\qquad(i\in\mathbb Z)
  \tag{11.4.5.1}
\]
이 성립한다.
\end{proposition}

\begin{proof}
실제로 첫째 스펙트럼 열 (\hyperref[0.11.4.3.1-fr]{11.4.3.1})에서
0이 아닌 $E_2$항은 $'E_2^{p0}=\mathrm H^p(T(K^\bullet))$뿐이다.
다시 말해 이 스펙트럼 열은 \emph{퇴화}한다. 이 열은 정칙이므로
(\hyperref[0.11.4.4-fr]{11.4.4}), 결론은
(\hyperref[0.11.1.6-fr]{11.1.6})에서 따른다.
\end{proof}

\begin{env}[11.4.6]
\label{0.11.4.6-fr}
이제 예를 들어 $\mathcal C\times\mathcal C'$에서 $\mathcal C''$으로
가는 \emph{공변 쌍함자} $(M,N)\mapsto T(M,N)$을 생각하자. 여기서
$\mathcal C$, $\mathcal C'$, $\mathcal C''$은 세 아벨 범주이다. 간단히
하기 위하여 $T$는 각 변수에 관하여 가법이고, $\mathcal C$와
$\mathcal C'$의 모든 대상은 각각 어떤 단사 대상의 부분대상이며,
$\mathcal C''$에서 여과 귀납극한들이 존재하고 완전하다고 가정한다.
그러면 미분작용소의 차수가 $+1$인 $\mathcal C$와 $\mathcal C'$의 두
복합체 $K^\bullet$, $K'^\bullet$에 관한 $T$의 \emph{초코호몰로지}를
다음과 같이 정의한다. $K^\bullet$ (각각 $K'^\bullet$)의 단사
카르탕--아일렌베르크 분해 $L^{\bullet\bullet}$ (각각
$L'^{\bullet\bullet}$)을 취한다. 그러면
$T(L^{\bullet\bullet},L'^{\bullet\bullet})$은 $\mathcal C''$의
사중복합체이며, $T(L^{ij},L'^{hk})$의 차수를 $i+h$와 $j+k$로 취하여
$\mathcal C''$의 \emph{이중복합체}로 본다. \emph{$K^\bullet$과
$K'^\bullet$에 관한 $T$의 초코호몰로지}는 정의에 따라 이
이중복합체의 코호몰로지(즉, 연관된 단순복합체의 코호몰로지)
$\mathrm H^\bullet(T(L^{\bullet\bullet},L'^{\bullet\bullet}))$이며,
$\mathrm R^\bullet T(K^\bullet,K'^\bullet)$로 나타낸다. 이는 두
스펙트럼 열의 수렴대상이고, 그 $E_2$항들은 다음과 같다.
\oldpage[0\textsubscript{III}]{35}
\[
  \label{0.11.4.6.1-fr}
  'E_2^{pq}=\mathrm H^p(\mathrm R^qT(K^\bullet,K'^\bullet))
  \tag{11.4.6.1}
\]
\[
  \label{0.11.4.6.2-fr}
  ''E_2^{pq}=\sum_{q'+q''=q}
  \mathrm R^pT(\mathrm H^{q'}(K^\bullet),\mathrm H^{q''}(K'^\bullet))
  \tag{11.4.6.2}
\]
(M, XVII, 2 참조).

여기서 $\mathrm R^qT(K^\bullet,K'^\bullet)$은 이중복합체
$(\mathrm R^qT(K^i,K'^j))_{(i,j)\in\mathbb Z\times\mathbb Z}$이고,
(\hyperref[0.11.4.6.1-fr]{11.4.6.1})의 우변은 이를 단순복합체로 볼 때의
코호몰로지이다.

더 나아가 첫째 스펙트럼 열은 항상 정칙이다. $\mathcal C$와
$\mathcal C'$의 모든 대상이 길이 $\leq n$인 단사 분해를 갖게 하는
$n$이 존재하거나, $K^\bullet$과 $K'^\bullet$이 아래로 유계이면 두
스펙트럼 열은 쌍정칙이다. 뒤의 경우에는 $\mathcal C''$에서
귀납극한들이 존재한다는 가정도 생략할 수 있다.

$K_1^\bullet$, $K_1'^\bullet$을 각각 $\mathcal C$와 $\mathcal C'$의
다른 두 복합체라 하고, $f$, $g$를 $K^\bullet$에서 $K_1^\bullet$로
가는 서로 호모토픽한 두 사상, $f'$, $g'$를 $K'^\bullet$에서
$K_1'^\bullet$로 가는 서로 호모토픽한 두 사상이라 하자. 그러면
$f$, $f'$에서 얻는 사상과 $g$, $g'$에서 얻는 사상
\[
\mathrm R^\bullet T(K^\bullet,K'^\bullet)\to
\mathrm R^\bullet T(K_1^\bullet,K_1'^\bullet)
\]
은 서로 같고, 초코호몰로지 스펙트럼 열의 사상들도 마찬가지이다.

이는 임의의 가법 공변 다중함자로 쉽게 일반화된다.
\end{env}

\begin{proposition}[11.4.7]
\label{0.11.4.7-fr}
$\mathcal C$의 모든 단사 대상 $I$ (각각 $\mathcal C'$의 단사 대상
$I'$)에 대하여 $A'\mapsto T(I,A')$ (각각 $A\mapsto T(A,I')$)가
완전함자라고 가정하자. 그러면 (\hyperref[0.11.4.6-fr]{11.4.6})의
표기 아래 표준 동형
\[
  \label{0.11.4.7.1-fr}
  \mathrm R^\bullet T(K^\bullet,K'^\bullet)
  \xrightarrow{\sim}\mathrm H^\bullet(T(L^{\bullet\bullet},K'^\bullet))
  \xrightarrow{\sim}\mathrm H^\bullet(T(K^\bullet,L'^{\bullet\bullet}))
  \tag{11.4.7.1}
\]
이 성립한다. 여기서 마지막 두 항은 각각 삼중복합체
$T(L^{\bullet\bullet},K'^\bullet)$과
$T(K^\bullet,L'^{\bullet\bullet})$이 정의하는 단순복합체의
코호몰로지이다.
\end{proposition}

\begin{proof}
예를 들어 이 동형들 중 첫째 것을 정의하자. 사중복합체
$T(L^{\bullet\bullet},L'^{\bullet\bullet})$은
$T(L^{ij},L'^{hk})$의 차수를 $i+j+h$와 $k$로 취하여
\emph{이중복합체}로 볼 수 있다. 각 $h$에 대하여
$L'^{h,\bullet}$은 $K'^h$의 \emph{분해}이므로, $T$에 관한 가정에 따라
이 이중복합체에 대하여
$\mathrm H_{\mathrm{II}}^q(T(L^{\bullet\bullet},L'^{\bullet\bullet}))=0$
($q\neq0$)이고
\[
\mathrm H_{\mathrm{II}}^0(T(L^{\bullet\bullet},L'^{\bullet\bullet}))
=T(L^{\bullet\bullet},K'^\bullet).
\]
따라서 이 이중복합체의 첫째 스펙트럼 열은 \emph{퇴화}한다. 또한
$k<0$일 때 $L'^{hk}=0$이므로 이 열은 정칙이고
(\hyperref[0.11.3.3-fr]{11.3.3}), 결론은
(\hyperref[0.11.1.6-fr]{11.1.6})에서 따른다.
\end{proof}

임의의 수 $n$개의 변수를 갖는 공변 다중함자에 대해서도 비슷한 결과가
성립한다. 초코호몰로지를 계산할 때 모든 복합체를 카르탕--아일렌베르크
분해로 바꿀 필요는 없고, 그중 $n-1$개만 바꾸면 된다. 다만 임의의
$n-1$개 변수를 단사 대상들로 고정했을 때 나머지 변수에 관한 공변
함자는 완전해야 한다.

\subsection{초코호몰로지에서 귀납극한으로의 이행}
\label{subsection:0.11.5-ko}

\begin{lemma}[11.5.1]
\label{0.11.5.1-ko}
$K^\bullet=(K^i)_{i\in\mathbb Z}$를 $\mathcal C$의 복합체라 하고, 각 정수
$r\in\mathbb Z$에 대하여 $K_{(r)}^i=0$ ($i<r$),
$K_{(r)}^i=K^i$ ($i\geq r$)인 복합체를 $K_{(r)}^\bullet$이라 하자.
$T$를 $\mathcal C$에서 $\mathcal C'$로 가는 공변 가법 함자로서
\oldpage[0\textsubscript{III}]{36}
귀납극한과 가환하는 함자라 하자(여과 귀납극한은 $\mathcal C$와
$\mathcal C'$에서 존재하고 완전하다고 가정한다). 그러면 $r$이
$-\infty$로 갈 때 $\mathrm R^\bullet T(K^\bullet)$는 귀납극한
\[
  \varinjlim_{r\to-\infty}\mathrm R^\bullet T(K_{(r)}^\bullet)
\]
과 동형이다.
\end{lemma}

$K^\bullet$의 단사 카르탕--아일렌베르크 분해는 각 $i$에 대하여
$\mathrm B^i(K^\bullet)$의 단사 분해
$(X_{\mathrm B}^{ij})_{j\geq0}$와 $\mathrm H^i(K^\bullet)$의 단사 분해
$(X_{\mathrm H}^{ij})_{j\geq0}$를 임의로 택하여 구성한다. 이렇게 하면
구성법으로부터 $K^i$의 단사 분해 $(L^{ij})_{j\geq0}$와 미분작용소
$L^{i,j}\to L^{i+1,j}$는 \emph{오직 분해들}
$(X_{\mathrm B}^{i,\bullet})$, $(X_{\mathrm H}^{i,\bullet})$ 및
$(X_{\mathrm B}^{i+1,\bullet})$에만 의존함을 알 수 있다
(M, XVII, 1.2). 한편 $i\geq r+1$이면 분명히
$\mathrm B^i(K_{(r)}^\bullet)=\mathrm B^i(K^\bullet)$이고
$\mathrm H^i(K_{(r)}^\bullet)=\mathrm H^i(K^\bullet)$이다. 또 각 $r$에
대하여 표준 단사사상
$\varphi_{r-1,r}^\bullet:K_{(r)}^\bullet\to K_{(r-1)}^\bullet$이 있으며,
$\varphi_{r-1,r}^i$는 $i\neq r-1$일 때 항등사상이다. 앞의 관찰에
따르면 $L^{\bullet\bullet}=(L^{ij})$가 $K^\bullet$의 단사
카르탕--아일렌베르크 분해일 때, 각 $r$에 대하여 $K_{(r)}^\bullet$의
단사 카르탕--아일렌베르크 분해
$L_{(r)}^{\bullet\bullet}=(L_{(r)}^{ij})$를
$i<r$이면 $L_{(r)}^{ij}=0$, $i\geq r+1$이면
$L_{(r)}^{ij}=L^{ij}$이 되도록 정의할 수 있다. 또한
$\varphi_{r-1,r}^\bullet$에 대응하는 이중복합체의 사상
$\Phi_{r-1,r}^{\bullet\bullet}:L_{(r)}^{\bullet\bullet}\to
L_{(r-1)}^{\bullet\bullet}$을 정의할 수 있다. 이 사상의 정의법
(\emph{같은 곳})을 다시 적용하면 $i\neq r$이고 $i\neq r-1$일 때
$\Phi_{r-1,r}^{ij}$가 항등사상이 되도록 구성할 수 있다. 이와 같이
$\mathcal C$의 이중복합체들로 이루어진 귀납계
$(L_{(r)}^{\bullet\bullet})$를 정의하였다. $r$이 $-\infty$로 갈 때
그 귀납극한은 분명히 $L^{\bullet\bullet}$이다. 직합과 귀납극한이
가환하므로, $L^{\bullet\bullet}$에 결부된 단순복합체도
$L_{(r)}^{\bullet\bullet}$에 결부된 단순복합체들의 귀납극한이다.
$T$는 가정에 따라 귀납극한과 가환하고, 코호몰로지도 그러하므로
(함자 $\varinjlim$의 완전성에 의하여), 동형을 제외하면
\[
  \mathrm H^\bullet(T(L^{\bullet\bullet}))=
  \varinjlim\mathrm H^\bullet(T(L_{(r)}^{\bullet\bullet}))
\]
이다.

보조정리 (\hyperref[0.11.5.1-ko]{11.5.1})을 쓰면 아래로 유계인
복합체들에 대하여 성립하는 성질을 귀납극한으로 이행시켜 임의의
복합체 $K^\bullet$로 확장할 수 있다. 첫 번째 예로 다음을 증명한다.

\begin{proposition}[11.5.2]
\label{0.11.5.2-ko}
$\mathcal C$, $\mathcal C'$ 및 $T$에 관한
(\hyperref[0.11.5.1-ko]{11.5.1})의 가정 아래에서
$\mathrm R^\bullet T(K^\bullet)$는 $\mathcal C$의 복합체들의 아벨
범주에서 코호몰로지 함자이다.
\end{proposition}

\begin{proof}
아래로 유계인 복합체들의 경우로 환원할 수 있음을 보이자. 복합체들의
완전열
\[
  0\longrightarrow K'^\bullet\longrightarrow K^\bullet
  \longrightarrow K''^\bullet\longrightarrow0
\]
이 주어지면, 각 $r$에 대하여 분명히 완전열
\[
  0\longrightarrow K_{(r)}'^\bullet\longrightarrow K_{(r)}^\bullet
  \longrightarrow K_{(r)}''^\bullet\longrightarrow0
\]
을 얻으며, 따라서 가정에 의하여 완전열
\[
  \cdots\longrightarrow\mathrm R^nT(K_{(r)}'^\bullet)
  \longrightarrow\mathrm R^nT(K_{(r)}^\bullet)
  \longrightarrow\mathrm R^nT(K_{(r)}''^\bullet)
  \xrightarrow{\partial}\mathrm R^{n+1}T(K_{(r)}'^\bullet)
  \longrightarrow\cdots
\]
을 얻는다. 이 완전열들은 귀납계를 이룬다. 보조정리
(\hyperref[0.11.5.1-ko]{11.5.1})과 함자 $\varinjlim$의 완전성으로부터
완전열
\[
  \cdots\longrightarrow\mathrm R^nT(K'^\bullet)
  \longrightarrow\mathrm R^nT(K^\bullet)
  \longrightarrow\mathrm R^nT(K''^\bullet)
  \xrightarrow{\partial}\mathrm R^{n+1}T(K'^\bullet)
  \longrightarrow\cdots
\]
을 얻는다.
\end{proof}

아래로 유계인 복합체들을 다룰 때에는 $i<0$이면 $K^i=0$인 것들만
고려해도 충분하다. 이들은 분명히 아벨 범주 $\mathcal K$를 이룬다.

\begin{lemma}[11.5.2.1]
\label{0.11.5.2.1-ko}
$\mathcal K$에서 다음 성질들을 갖는 복합체
$Q^\bullet=(Q^i)_{i\geq0}$의 집합을 $\mathfrak J$라 하자.
\oldpage[0\textsubscript{III}]{37}
\begin{enumerate}
  \item[$1^\circ$] 각 $Q^i$는 $\mathcal C$의 단사 대상이다.
  \item[$2^\circ$] 각 $i\geq0$에 대하여
    $\mathrm Z^i(Q^\bullet)=\mathrm B^i(Q^\bullet)$이고,
    $\mathrm Z^i(Q^\bullet)$는 $Q^i$의 직합인자이다.
\end{enumerate}
그러면 다음이 성립한다.
\begin{enumerate}
  \item[(i)] 모든 $Q^\bullet\in\mathfrak J$는 $\mathcal K$의 단사
    대상이다.
  \item[(ii)] $\mathcal K$의 모든 대상은 $\mathfrak J$에 속하는 어떤
    복합체의 부분복합체와 동형이다.
\end{enumerate}
\end{lemma}

\begin{proof}
\textup{(i)} $A^\bullet=(A^i)$를 $\mathcal K$의 대상,
$A'^\bullet=(A'^i)$를 $A^\bullet$의 부분대상,
$Q^\bullet=(Q^i)$를 $\mathfrak J$의 대상이라 하자. 사상
$f=(f^i):A'^\bullet\to Q^\bullet$이 주어졌다고 하고, 이를 사상
$g=(g^i):A^\bullet\to Q^\bullet$으로 연장하려 한다. 표기를
간단히 하기 위하여 가군 범주의 언어를 쓰겠다([27] 참조).

$Q^i$를
$\mathrm B^i(Q^\bullet)\oplus\mathrm B^{i+1}(Q^\bullet)$와
동일시하자. $i$에 관한 귀납법을 쓰며, 따라서 $j<i$에 대하여
$g^j$가 정의되어 있고 $j<i-1$에 대하여 미분작용소
$d^j:A^j\to A^{j+1}$ 및 $d'^j:Q^j\to Q^{j+1}$과 양립하며, 더 나아가
다음을 만족한다고 가정하자.
\begin{enumerate}
  \item[$1^\circ$] $g^{i-1}(\mathrm Z^{i-1}(A^\bullet))\subset
    \mathrm Z^{i-1}(Q^\bullet)$.
  \item[$2^\circ$] 모든 $j$에 대하여
    $C^j=(d^j)^{-1}(A'^{j+1})$라 놓으면,
    $d'^{i-1}\circ g^{i-1}$과 $f^i\circ d^{i-1}$은 $C^{i-1}$ 위에서
    일치한다.
\end{enumerate}
사상 $f^i:A'^i\to Q^i$를 두 사영과 합성하면 두 사상
$f'^i:A'^i\to\mathrm B^i(Q^\bullet)$ 및
$f''^i:A'^i\to\mathrm B^{i+1}(Q^\bullet)$을 얻는다.
$d'^{i-1}\circ g^{i-1}$은 $A^{i-1}$을 $\mathrm B^i(Q^\bullet)$로
보내고 $\mathrm Z^{i-1}(A^\bullet)$ 위에서 0이므로 사상
$h^i:\mathrm B^i(A^\bullet)\to\mathrm B^i(Q^\bullet)$을 정의한다.
또 $d'^{i-1}\circ g^{i-1}$과 $f^i\circ d^{i-1}$은 $C^{i-1}$ 위에서
일치하므로, $h^i$와 $f'^i$는
$\mathrm B^i(A^\bullet)\cap A'^i$ 위에서 일치한다.
$\mathrm B^i(Q^\bullet)$은 $Q^i$의 직합인자이므로 단사 대상이다.
따라서 $\mathrm B^i(A^\bullet)$ 위에서는 $h^i$와, $A'^i$ 위에서는
$f'^i$와 일치하는 사상
$g'^i:A^i\to\mathrm B^i(Q^\bullet)$가 존재한다. 한편 사상
$f'^{i+1}\circ d^i:C^i\to\mathrm B^{i+1}(Q^\bullet)$을 생각하자.
이는 $\mathrm Z^i(A^\bullet)$ 위에서 0이다. 그런데
$\mathrm B^{i+1}(Q^\bullet)$은 단사 대상이므로, $C^i$ 위에서는
$f'^{i+1}\circ d^i$와 일치하고 $\mathrm Z^i(A^\bullet)$ 위에서는
$0$과 일치하는 사상
$g''^i:A^i\to\mathrm B^{i+1}(Q^\bullet)$가 존재한다. 이제
$g^i=g'^i+g''^i$로 놓으면 귀납을 계속할 수 있다.

\textup{(ii)} $A^\bullet=(A^i)$를 $\mathfrak J$에 속하는 복합체에
매입하기 위하여, 각 $i\geq1$에 대하여 단사사상
$f'^i:A^i\to Q'^i$가 존재하도록 $\mathcal C$의 단사 대상 $Q'^i$를
택한다. 이제
\[
  Q^i=0\quad(i<0),\qquad Q^0=Q'^1,\qquad
  Q^i=Q'^i\oplus Q'^{i+1}\quad(i\geq1)
\]
로 놓고, 자명한 미분작용소를 준다. 모든 $i$에 대하여
$f^i=f'^i+(f'^{i+1}\circ d^i)$로 놓자($i\leq0$이면
$f'^i=0$으로 한다). 각 $f^i$가 단사이고 $f^i$들이 미분작용소와
양립함은 곧바로 알 수 있다.
\end{proof}

\begin{corollary}[11.5.2.2]
\label{0.11.5.2.2-ko}
$\mathcal K$의 모든 대상 $K^\bullet$은 $\mathfrak J$의 대상들로
이루어진 오른쪽 분해를 갖는다. $L^{\bullet\bullet}$이 그러한
분해이고, $L'^{\bullet\bullet}$이 $\mathcal K$의 대상들로 이루어진
$K^\bullet$의 임의의 분해이면, 확대사상들과 양립하는 이중복합체의
사상 $F:L'^{\bullet\bullet}\to L^{\bullet\bullet}$이 존재하고, 그러한
두 사상 $F$, $F'$은 서로 호모토픽이다.
\end{corollary}

\begin{proof}
이는 아벨 범주 $\mathcal K$에 (M, V, 1.1 a))를 적용한 것에 지나지
않는다.
\end{proof}

\begin{env}[11.5.2.3]
\label{0.11.5.2.3-ko}
이 준비들을 갖추었으므로, $K^\bullet$의 단사 카르탕--아일렌베르크
분해 $L'^{\bullet\bullet}$과 $\mathfrak J$의 대상들로 이루어진
$K^\bullet$의 분해 $L^{\bullet\bullet}$을 생각하고, 다음 동형이
성립함을 보이자.
\[
  \mathrm H^\bullet(T(L'^{\bullet\bullet}))\xrightarrow{\sim}
  \mathrm H^\bullet(T(L^{\bullet\bullet})).
\]
실제로 (\hyperref[0.11.5.2.2-ko]{11.5.2.2})로부터 이중복합체의 사상
$T(L'^{\bullet\bullet})\to T(L^{\bullet\bullet})$을 얻고, 따라서 이
이중복합체들의 첫째 분광수열 사이의 사상
$'E(T(L'^{\bullet\bullet}))\to{}'E(T(L^{\bullet\bullet}))$을 얻는다.
(\hyperref[0.11.3.3-fr]{11.3.3})에 의하여 이 분광수열들은 정칙이므로,
앞의 사상이 $E_2$ 항들에서 동형임을 보이면 충분하다
(\hyperref[0.11.1.5-fr]{11.1.5}). 다시 말해
$\mathrm H_{\mathrm{II}}^q(T(L^{i,\bullet}))$가
$\mathrm R^qT(K^i)$와 같음을 보이면 충분하다. 그런데
$L^{i,\bullet}$은 $K^i$의 오른쪽 분해이므로 다음 보조정리를
증명하는 문제로 환원된다.
\end{env}

\oldpage[0\textsubscript{III}]{38}
\begin{lemma}[11.5.2.4]
\label{0.11.5.2.4-ko}
$L^\bullet=(L^i)_{i\in\mathbb Z}$가 $\mathcal C$의 대상 $A$의 오른쪽
분해이고, 모든 $i\in\mathbb Z$와 모든 $n>0$에 대하여
$\mathrm R^nT(L^i)=0$이라 하자. 그러면
$\mathrm H^\bullet T(L^\bullet)=\mathrm R^\bullet T(A)$이다.
\end{lemma}

\begin{proof}
이는 (T, 2.5.2)의 특수한 경우이다.
\end{proof}

\begin{env}[11.5.2.5]
\label{0.11.5.2.5-ko}
이제 (\hyperref[0.11.5.2-ko]{11.5.2})의 증명은 곧바로 끝난다.
실제로 $L'^{\bullet\bullet}$은 아벨 범주 $\mathcal K$에서
$K^\bullet$의 단사 분해이다. 다시 말해
$K^\bullet\mapsto\mathrm H^\bullet(T(L'^{\bullet\bullet}))$은 바로
범주 $\mathcal K$에서 $T$의 오른쪽 유도 코호몰로지 함자이다
(T, 2.3).
\end{env}

\begin{proposition}[11.5.3]
\label{0.11.5.3-ko}
$\mathcal C$, $\mathcal C'$ 및 $T$에 관한
(\hyperref[0.11.5.1-ko]{11.5.1})의 가정 아래에서,
$L^{\bullet\bullet}=(L^{ij})$를 $j<0$이면 $L^{ij}=0$이고 각 $i$에
대하여 $L^{i,\bullet}$이 $K^i$의 분해인 이중복합체라 하자. 끝으로
모든 쌍 $(i,j)$와 모든 $n>0$에 대하여
$\mathrm R^nT(L^{ij})=0$이라고 가정하자. 그러면 함자적 동형
\[
  \label{0.11.5.3.1-ko}
  \mathrm R^\bullet T(K^\bullet)\xrightarrow{\sim}
  \mathrm H^\bullet(T(L^{\bullet\bullet})).
  \tag{11.5.3.1}
\]
이 존재한다.
\end{proposition}

\begin{proof}
$i<r$이면 $L_{(r)}^{ij}=0$, $i\geq r$이면
$L_{(r)}^{ij}=L^{ij}$인 이중복합체를
$L_{(r)}^{\bullet\bullet}=(L_{(r)}^{ij})$라 하자. $r$이
$-\infty$로 갈 때 $L^{\bullet\bullet}$이
$L_{(r)}^{\bullet\bullet}$의 귀납극한임은 곧바로 알 수 있다.
$T$에 대한 가정과 (\hyperref[0.11.5.1-ko]{11.5.1})에 의하여, 따라서
$K^\bullet$이 아래로 유계인 경우만 증명하면 충분하다. 예를 들어
$i<0$이면 $K^i=0$이고 $L^{ij}=0$인 경우만 다루면 된다. 이제
$L'^{\bullet\bullet}=(L'^{ij})$을 $\mathfrak J$의 대상들로 이루어진
$K^\bullet$의 오른쪽 분해라 하자
(\hyperref[0.11.5.2.2-ko]{11.5.2.2}). 확대사상들과 양립하는
이중복합체의 사상 $L^{\bullet\bullet}\to L'^{\bullet\bullet}$이
존재하므로, 첫째 분광수열 사이의 사상
$'E(T(L^{\bullet\bullet}))\to{}'E(T(L'^{\bullet\bullet}))$을 얻는다.
보조정리 (\hyperref[0.11.5.2.4-ko]{11.5.2.4})에 의하여,
(\hyperref[0.11.5.2.3-ko]{11.5.2.3})에서와 같이 이 사상은
\emph{동형}이고, 이로부터 결론이 따른다.
\end{proof}

\begin{remark}[11.5.4]
\label{0.11.5.4-ko}
앞의 논증은 아래로 유계인 복합체 $K^\bullet$의 범주에서는
\emph{$T$가 여과 귀납극한과 가환한다고 가정하지 않아도}
(\hyperref[0.11.5.2-ko]{11.5.2})와
(\hyperref[0.11.5.3-ko]{11.5.3})의 결론들이 성립함을 증명한다.
더욱이 $i<0$이면 $K^i=0$인 복합체 $K^\bullet$들의 범주
$\mathcal K$만을 생각할 때에는, $\mathrm R^\bullet T(K^\bullet)$가
$\mathcal K$에서 $T$의 \emph{오른쪽 유도 함자}들의 계라는 특징으로부터
이 코호몰로지 함자가 \emph{보편적}임을 알 수 있다(T, 2.3).

$T$에 관하여 추가 가정을 하지 않고도
(\hyperref[0.11.5.2-ko]{11.5.2})가 성립하는 또 하나의 경우는,
$\mathcal C$의 모든 대상이 길이 $\leq m$인 단사 분해를 갖게 하는
정수 $m>0$이 존재하는 경우이다. 실제로
(\hyperref[0.11.5.1-ko]{11.5.1})의 증명에 나타나는
$\mathcal C$의 대상들의 단사 분해를 모두 길이 $\leq m$인 것으로
택할 수 있다. 따라서 $-r$이 충분히 크면 이중복합체
$T(L_{(r)}^{\bullet\bullet})$의 전체차수 $n$인 항들은
$T(L^{\bullet\bullet})$의 항들과 같고 그 수가 유한하다. 이에 따라
각 $n$에 대하여 $-r$이 충분히 크면
$\mathrm H^n(T(L^{\bullet\bullet}))=
\mathrm H^n(T(L_{(r)}^{\bullet\bullet}))$이다.
(\hyperref[0.11.5.2-ko]{11.5.2})의 표기를 쓰면, 따라서 $-r$이 충분히
클 때($n$에 의존한다)
$\mathrm R^nT(K_{(r)}^\bullet)=\mathrm R^nT(K^\bullet)$이고,
$K'^\bullet$과 $K''^\bullet$에 대해서도 마찬가지이므로 결론이 따른다.
같은 방법으로, $\mathcal C$가 앞의 가정을 만족하고 분해들
$L^{i,\bullet}$의 길이가 $\leq m$이라고 가정하면
(\hyperref[0.11.5.3-ko]{11.5.3})도 $T$에 대한 추가 조건 없이
성립한다.
\end{remark}

\begin{env}[11.5.5]
\label{0.11.5.5-ko}
(\hyperref[0.11.5.2-ko]{11.5.2})의 결과는 공변 다변함자로
일반화된다. 예를 들어 (\hyperref[0.11.4.6-fr]{11.4.6})의 상황을
생각하자. 여기서는 $\mathcal C$, $\mathcal C'$ 및 $\mathcal C''$에서
\oldpage[0\textsubscript{III}]{39}
여과 귀납극한이 존재하고 완전하며, $T$가 귀납극한과 가환한다고
가정한다. 그러면 $\mathrm R^\bullet T(K^\bullet,K'^\bullet)$은
각 복합체 $K^\bullet$, $K'^\bullet$에 관하여 코호몰로지 쌍함자이다.
이를 보이기 위하여 (\hyperref[0.11.5.2-ko]{11.5.2})에서와 같이
$K^\bullet$과 $K'^\bullet$이 아래로 유계인 경우로 환원한다. 그런 다음
$K^\bullet$과 $K'^\bullet$에 대하여
(\hyperref[0.11.5.2.2-ko]{11.5.2.2})에 기술된 유형의 단사 분해를
택하면, (M, V, 4.1)에서 증명한 일반 성질로 환원된다.
\end{env}

\begin{env}[11.5.6]
\label{0.11.5.6-ko}
마찬가지로 (\hyperref[0.11.4.7-fr]{11.4.7})과
(\hyperref[0.11.5.3-ko]{11.5.3})의 결과들은 다음과 같이 일반화된다.
(\hyperref[0.11.5.5-ko]{11.5.5})의 가정 아래에서 두 이중복합체
$L^{\bullet\bullet}=(L^{ij})$,
$L'^{\bullet\bullet}=(L'^{ij})$가 주어졌다고 하자. $j<0$이면
$L^{ij}=0$이고 $L'^{ij}=0$이며, 각 $i$에 대하여
$L^{i,\bullet}$은 $K^i$의 분해이고 $L'^{i,\bullet}$은 $K'^i$의
분해라고 하자. 끝으로 모든 지표계 $(i,j,h,k)$와 $n>0$에 대하여
$\mathrm R^nT(L^{ij},L'^{hk})=0$이라고 하자. 그러면
$K^\bullet$과 $K'^\bullet$에 관한 함자적 동형
\[
  \label{0.11.5.6.1-ko}
  \mathrm R^\bullet T(K^\bullet,K'^\bullet)\xrightarrow{\sim}
  \mathrm H^\bullet(T(L^{\bullet\bullet},L'^{\bullet\bullet})).
  \tag{11.5.6.1}
\]
이 성립한다. 이는 (\hyperref[0.11.5.3-ko]{11.5.3})에서와 같이
$K^\bullet$과 $K'^\bullet$이 아래로 유계인 경우로 환원하여 증명한다.

더 나아가 각 쌍 $(i,j)$와 각 쌍 $(h,k)$에 대하여 함자
$A\mapsto T(A,L'^{hk})$와 $A'\mapsto T(L^{ij},A')$가 각각
$\mathcal C$와 $\mathcal C'$에서 완전하다고 가정하자. 그러면 또한
함자적 동형
\[
  \label{0.11.5.6.2-ko}
  \mathrm R^\bullet T(K^\bullet,K'^\bullet)
  \xrightarrow{\sim}\mathrm H^\bullet(T(L^{\bullet\bullet},K'^\bullet))
  \xrightarrow{\sim}\mathrm H^\bullet(T(K^\bullet,L'^{\bullet\bullet})).
  \tag{11.5.6.2}
\]
이 성립한다. 증명은 (\hyperref[0.11.4.7-fr]{11.4.7})의 증명과
비슷하다.
\end{env}

\begin{env}[11.5.7]
\label{0.11.5.7-ko}
끝으로 (\hyperref[0.11.5.5-ko]{11.5.5})와
(\hyperref[0.11.5.6-ko]{11.5.6})의 결과들은 다음 두 경우에 $T$가
귀납극한과 가환한다고 가정하지 않아도 성립함에 유의하자. 하나는
아래로 유계인 복합체 $K^\bullet$, $K'^\bullet$에 한정하는 경우이고,
다른 하나는 $\mathcal C$(각각 $\mathcal C'$)의 모든 대상이 길이가
유계인 단사 분해를 갖는다고 가정하고
(\hyperref[0.11.5.6-ko]{11.5.6})에서 이중복합체
$L^{\bullet\bullet}$과 $L'^{\bullet\bullet}$의 둘째 차수가 위로
유계라고 가정하는 경우이다.
\end{env}

\subsection{복합체 $K_\bullet$에 관한 함자의 초호몰로지}
\label{subsection:0.11.6-fr}

\begin{env}[11.6.1]
\label{0.11.6.1-fr}
$\mathcal C$를 아벨 범주,
$K_\bullet=(K_i)_{i\in\mathbb Z}$를 미분작용소의 차수가 $-1$인
$\mathcal C$의 대상들의 복합체라 하자. $K_\bullet$의
\emph{왼쪽 카르탕--아일렌베르크 분해}란, 미분작용소들의 차수가
$-1$이고 $j<0$이면 $L_{ij}=0$인 이중복합체
$L_{\bullet\bullet}=(L_{ij})$와 단순복합체들의 사상
\[
\varepsilon:L_{\bullet,0}\to K_\bullet
\]
로 이루어지며, (\hyperref[0.11.4.2-fr]{11.4.2})의 조건들에서
``화살표의 방향을 뒤집어서'' 얻은 조건들을 만족하는 것이다.
$\mathcal C$의 모든 대상이 어떤 사영 대상의 몫이면,
$\mathcal C$의 모든 복합체 $K_\bullet$은
\emph{사영 카르탕--아일렌베르크 분해}, 즉 사영 대상들 $L_{ij}$로
이루어진 분해를 가지며, 이는
(\hyperref[0.11.4.2-fr]{11.4.2})와 비슷한 함자적 성질들을 갖는다.
또한 $K_\bullet$이 아래로 유계이면 (각각 위로 유계이면),
$i<i_0$일 때 $K_i=0$인 경우 (각각 $i>i_0$일 때 $K_i=0$인 경우)
$i<i_0$이면 $L_{ij}=0$이라고 (각각 $i>i_0$이면
$L_{ij}=0$이라고) 가정할 수 있다. $\mathcal C$의 모든 대상이
길이 $\leq n$인 사영 분해를 가지면 $j>n$일 때 $L_{ij}=0$이라고
가정할 수 있다.
\end{env}

\begin{env}[11.6.2]
\label{0.11.6.2-fr}
\oldpage[0\textsubscript{III}]{40}
$T$를 $\mathcal C$에서 아벨 범주 $\mathcal C'$으로 가는 공변 가법
함자라고 가정하자. $\mathcal C$의 복합체 $K_\bullet$에 관한
$T$의 \emph{초호몰로지} $L_\bullet T(K_\bullet)$와
\emph{초호몰로지 스펙트럼 열}들이 존재할 때의 정의는 여전히
(\hyperref[0.11.4.3-fr]{11.4.3})에서 ``화살표의 방향을
뒤집어서'' 얻는다. 이렇게 얻은 두 스펙트럼 열의 $E_2$ 항은
\[
\label{0.11.6.2.1-fr}
\tag{11.6.2.1}
{}'E^2_{pq}=H_p(L_qT(K_\bullet))
\]
및
\[
\label{0.11.6.2.2-fr}
\tag{11.6.2.2}
{}''E^2_{pq}=L_pT(H_q(K_\bullet))
\]
이다. 여기서 $L_pT$는 $p\geq0$이면 $T$의 $p$번째 왼쪽 유도
함자이고 $p<0$이면 $0$이며, $L_qT(K_\bullet)$는 복합체
$(L_qT(K_i))_{i\in\mathbb Z}$이다.

초호몰로지의 성질들이 모두 초코호몰로지의 성질들에서 단순히
``화살표의 방향을 뒤집어서'' 따라오는 것은 아니다. 범주
$\mathcal C'$에 (T, 1.5)의 (AB~5*) 유형의 추가 가정을 두지 않으면,
앞의 두 스펙트럼 열의 정칙성 조건 때문에 이번에는
(\hyperref[0.11.3.4-fr]{11.3.4})의 판정법들을 적용해야 하기
때문이다. 이 판정법들에 의하면 $\mathcal C'$에서 여과 귀납극한들이
존재하고 완전하다고 가정할 때, 이중복합체
$T(L_{\bullet\bullet})$가 정하는 복합체가 존재하고 둘째 스펙트럼 열
${}''E(T(L_{\bullet\bullet}))$은 정칙이다. $K_\bullet$이 아래로
유계라고 가정하거나, $\mathcal C$의 모든 대상이 길이 $\leq n$인
사영 분해를 갖게 하는 어떤 정수 $n$이 존재한다고 가정하면, 두
초호몰로지 스펙트럼 열은 $\mathcal C'$에 대한 아무 가정 없이도
존재하며 쌍정칙이다.
\end{env}

\begin{proposition}[11.6.3]
\label{0.11.6.3-fr}
$\mathcal C$, $\mathcal C'$을 두 아벨 범주, $T$를 $\mathcal C$에서
$\mathcal C'$으로 가는 공변 가법 함자라 하자. 그러면 다음이
성립한다.
\begin{enumerate}
\item[\rm(i)]
\emph{초호몰로지 $L_\bullet T(K_\bullet)$은 $\mathcal C$의 아래로
유계인 복합체들의 아벨 범주에서 호몰로지 함자이다.}
\item[\rm(ii)]
\emph{$K_\bullet$을 $\mathcal C$의 아래로 유계인 복합체라 하자.
모든 $n>0$와 모든 $i\in\mathbb Z$에 대해 $L_nT(K_i)=0$이면
함자적 동형}
\[
\label{0.11.6.3.1-fr}
\tag{11.6.3.1}
L_iT(K_\bullet)\simeq H_i(T(K_\bullet))\qquad(i\in\mathbb Z)
\]
\emph{이 성립한다.}
\item[\rm(iii)]
\emph{$K_\bullet$을 $\mathcal C$의 아래로 유계인 복합체라 하자.
$L_{\bullet\bullet}=(L_{ij})$를 $j<0$이면 $L_{ij}=0$이고, 모든
$i$에 대해 $L_{i,\bullet}$이 $K_i$의 분해인 이중복합체라 하자.
끝으로 모든 쌍 $(i,j)$와 모든 $n>0$에 대해
$L_nT(L_{ij})=0$이라고 가정하자. 그러면 함자적 동형}
\[
\label{0.11.6.3.2-fr}
\tag{11.6.3.2}
L_\bullet T(K_\bullet)\simeq H_\bullet(T(L_{\bullet\bullet}))
\]
\emph{이 존재한다.}
\end{enumerate}
아래로 유계인 복합체의 경우 증명은 각각
(\hyperref[0.11.5.2-ko]{11.5.2}),
(\hyperref[0.11.4.5-fr]{11.4.5}),
(\hyperref[0.11.5.3-ko]{11.5.3})의 증명과 같은 방식으로 한다.
이 논증들의 세부사항은 독자에게 맡긴다.
\end{proposition}

\begin{env}[11.6.4]
\label{0.11.6.4-fr}
각 변수에 관해 공변이고 가법인 다변함자에 대해서도 완전히 비슷한
결과들이 성립한다. 예를 들어 쌍함자 $T$에 대한 두 초호몰로지
스펙트럼 열의 $E_2$ 항은 다음과 같다.
\oldpage[0\textsubscript{III}]{41}
\[
\label{0.11.6.4.1-fr}
\tag{11.6.4.1}
{}'E^2_{pq}=H_p(L_qT(K_\bullet,K'_\bullet))
\]
\[
\label{0.11.6.4.2-fr}
\tag{11.6.4.2}
{}''E^2_{pq}=
\sum_{q'+q''=q}L_pT(H_{q'}(K_\bullet),H_{q''}(K'_\bullet)).
\]
여기에서도 둘째 스펙트럼 열이 정칙이다. 복합체 $K_\bullet$과
$K'_\bullet$이 아래로 유계이거나, 고려하는 아벨 범주들의 대상들이
고정된 길이의 사영 분해를 가질 때에는 두 스펙트럼 열이 쌍정칙이다.

또한 다음이 성립한다.
\end{env}

\begin{proposition}[11.6.5]
\label{0.11.6.5-fr}
$\mathcal C$, $\mathcal C'$, $\mathcal C''$을 세 아벨 범주,
$T$를 $\mathcal C\times\mathcal C'$에서 $\mathcal C''$으로 가는
공변 쌍가법 쌍함자라 하자.
\begin{enumerate}
\item[\rm(i)]
\emph{$L_\bullet T(K_\bullet,K'_\bullet)$은 아래로 유계인 각 복합체
$K_\bullet$, $K'_\bullet$에 관해 호몰로지 쌍함자이다. 여기서 이
복합체들은 각각 $\mathcal C$의 대상들과 $\mathcal C'$의 대상들로
이루어진다.}
\item[\rm(ii)]
\emph{$K_\bullet$과 $K'_\bullet$이 아래로 유계라고 가정하자.
$L_{\bullet\bullet}=(L_{ij})$,
$L'_{\bullet\bullet}=(L'_{ij})$를 $j<0$이면
$L_{ij}=L'_{ij}=0$이고, 모든 $i$에 대해 $L_{i,\bullet}$은
$K_i$의 분해이며 $L'_{i,\bullet}$은 $K'_i$의 분해인 두
이중복합체라 하자. 끝으로 모든 $n>0$와 모든 지표계
$(i,j,h,k)$에 대해 $L_nT(L_{ij},L'_{hk})=0$이라고 하자. 그러면
함자적 동형}
\[
\label{0.11.6.5.1-fr}
\tag{11.6.5.1}
L_\bullet T(K_\bullet,K'_\bullet)
\simeq H_\bullet(T(L_{\bullet\bullet},L'_{\bullet\bullet}))
\]
\emph{이 성립한다.}
\item[\rm(iii)]
\emph{더 나아가 모든 쌍 $(i,j)$와 모든 쌍 $(h,k)$에 대해 함자}
\[
A\longmapsto T(A,L'_{hk}),\qquad
A'\longmapsto T(L_{ij},A')
\]
\emph{가 각각 $\mathcal C$와 $\mathcal C'$에서 완전하다고 하자.
그러면 함자적 동형}
\[
\label{0.11.6.5.2-fr}
\tag{11.6.5.2}
L_\bullet T(K_\bullet,K'_\bullet)
\simeq H_\bullet(T(L_{\bullet\bullet},K'_\bullet))
\simeq H_\bullet(T(K_\bullet,L'_{\bullet\bullet}))
\]
\emph{들이 성립한다.}
\end{enumerate}
증명은 (\hyperref[0.11.5.5-ko]{11.5.5})와
(\hyperref[0.11.5.6-ko]{11.5.6})의 증명과 비슷하다.
\end{proposition}

\subsection{이중복합체 $K_{\bullet\bullet}$에 관한 함자의 초호몰로지}
\label{subsection:0.11.7-fr}

\begin{env}[11.7.1]
\label{0.11.7.1-fr}
$\mathcal C$를 모든 대상이 어떤 사영 대상의 몫인 아벨 범주라 하자.
$\mathcal C$의 대상들로 이루어지고 두 차수가 모두 아래로 유계인
이중복합체 $K_{\bullet\bullet}=(K_{ij})$를 생각하자. 언제나
$i<0$ 또는 $j<0$이면 $K_{ij}=0$인 경우로 한정할 수 있으므로,
이제부터 그렇게 하겠다. $K_{\bullet\bullet}$를 $\mathcal C$의
대상들로 이루어진 양의 차수 복합체들의 아벨 범주 $\mathcal K$의 대상
\[
K_{i,\bullet}=(K_{ij})_{j\geq0}
\]
들로 이루어진 (단순)복합체로 볼 수 있다. 보조정리
(\hyperref[0.11.5.2.1-ko]{11.5.2.1})(보다 정확히는 ``화살표의
방향을 뒤집어서'' 얻은 ``쌍대'' 보조정리)와 (M, V, 2.2)에 의해
$K_{\bullet\bullet}$은 범주 $\mathcal K$에서
\emph{사영 카르탕--아일렌베르크 분해}를 갖는다. 그러한 분해는 모든
차수가 $\geq0$이고 사영 대상들로 이루어진 $\mathcal C$의 삼중복합체
\[
M_{\bullet\bullet\bullet}=(M_{ijk})
\]
이다. 모든 $i$에 대해
\[
M_{i,\bullet\bullet},\quad
\mathrm B_i^{\mathrm I}(M_{\bullet\bullet\bullet}),\quad
\mathrm Z_i^{\mathrm I}(M_{\bullet\bullet\bullet}),\quad
H_i^{\mathrm I}(M_{\bullet\bullet\bullet})
\]
는 각각 범주 $\mathcal K$에서
\[
K_{i,\bullet},\quad
\mathrm B_i^{\mathrm I}(K_{\bullet\bullet}),\quad
\mathrm Z_i^{\mathrm I}(K_{\bullet\bullet}),\quad
H_i^{\mathrm I}(K_{\bullet\bullet})
\]
의 사영 분해를 이룬다. 특히 모든 쌍 $(i,j)$에 대해
$M_{ij,\bullet}$은 $\mathcal C$에서 $K_{ij}$의 사영 분해이다.
\end{env}

\begin{proposition}[11.7.2]
\label{0.11.7.2-fr}
$T$를 $\mathcal C$에서 아벨 범주 $\mathcal C'$으로 가는 공변 가법
함자라 하자. (\hyperref[0.11.7.1-fr]{11.7.1})의 표기 아래에서,
\emph{삼중복합체 $T(M_{\bullet\bullet\bullet})$에 대응하는
단순복합체의 호몰로지
$H_\bullet(T(M_{\bullet\bullet\bullet}))$는
$K_{\bullet\bullet}$에 대응하는 단순복합체의 초호몰로지
$L_\bullet T(K_{\bullet\bullet})$와 표준적으로 동형이다}
(\hyperref[0.11.6.2-fr]{11.6.2}). \emph{또한 이는}

\oldpage[0\textsubscript{III}]{42}

\emph{$t=a,b,a',b',c,d$에 대해 $\mathop{}^{(t)}E$로 나타내는 여섯
쌍정칙 스펙트럼 열의 공통 수렴 대상이며, 그 $E_2$ 항들은 다음
공식들로 주어진다.}
\[
\begin{aligned}
{}^{(a)}E^2_{pq}&=L_pT(H_q(K_{\bullet\bullet})),&
{}^{(b)}E^2_{pq}&=H_p(L_q^{\mathrm{II}}T(K_{\bullet\bullet})),\\
{}^{(a')}E^2_{pq}&=L_pT(H_q^{\mathrm I}(K_{\bullet\bullet})),&
{}^{(b')}E^2_{pq}&=H_p(L_qT(K_{\bullet\bullet})),\\
{}^{(c)}E^2_{pq}&=L_pT(H_q^{\mathrm{II}}(K_{\bullet\bullet})),&
{}^{(d)}E^2_{pq}&=H_p(L_q^{\mathrm I}T(K_{\bullet\bullet})).
\end{aligned}
\]
\end{proposition}

복합체 $A_\bullet=(A_i)$마다 대상 $F(A_i)$들로 이루어진 복합체를
$F(A_\bullet)$로 나타낸다는 표기를 사용함을 상기하자. 예를 들어
$L_q^{\mathrm{II}}T(K_{\bullet\bullet})$는 복합체
\[
(L_q^{\mathrm{II}}T(K_{i,\bullet}))_{i\geq0}
\]
를 나타내며, 여기서 $L_q^{\mathrm{II}}T(K_{i,\bullet})$는
단순복합체 $K_{i,\bullet}$에 관한 함자 $T$의 지표 $q$인
초호몰로지이다.

\begin{proof}
$L_\bullet$을 $K_{\bullet\bullet}$에 대응하는 단순복합체라 하자.
그러면
\[
L_i=\bigoplus_{r+s=i}K_{rs}.
\]
또한
\[
N_{ij}=\bigoplus_{r+s=i}M_{rsj}
\]
로 놓자. 모든 $i$에 대해 $N_{i,\bullet}$은 $\mathcal C$에서
$L_i$의 사영 분해임이 분명하다. 따라서
(\hyperref[0.11.6.3-fr]{11.6.3})과
(\hyperref[0.11.6.4-fr]{11.6.4})에 의해 함자적 동형
\[
L_\bullet T(L_\bullet)\simeq H_\bullet(T(N_{\bullet\bullet}))
\]
이 성립한다. 이중복합체 $T(N_{\bullet\bullet})$에 대응하는
단순복합체는 삼중복합체 $T(M_{\bullet\bullet\bullet})$에도
대응하므로, 이는 명제의 첫째 주장을 증명한다.

또한 $L_\bullet T(L_\bullet)$은 단순복합체 $L_\bullet$에 관한
$T$의 두 초호몰로지 스펙트럼 열
(\hyperref[0.11.6.2-fr]{11.6.2})의 수렴 대상이고, 이 두 열은 각각
$\mathop{}^{(b')}E$와 $\mathop{}^{(a)}E$이다.

이제 $M_{\bullet\bullet\bullet}$을
\[
U_{ij}=\bigoplus_{r+s=j}M_{irs}
\]
인 이중복합체 $U_{\bullet\bullet}$로 보자.
$H_\bullet(T(M_{\bullet\bullet\bullet}))$은 여전히 이중복합체
$T(U_{\bullet\bullet})$의 두 스펙트럼 열의 수렴 대상이다. 그런데
모든 $i$에 대해 $M_{i,\bullet\bullet}$은 단순복합체
$K_{i,\bullet}$에 관하여
(\hyperref[0.11.6.3-fr]{11.6.3})의 조건을 만족하는 이중복합체이다.
따라서
\[
H_q^{\mathrm{II}}(T(U_{i,\bullet}))
=L_q^{\mathrm{II}}T(K_{i,\bullet}),
\]
이고 $T(U_{\bullet\bullet})$의 첫째 스펙트럼 열은 바로
$\mathop{}^{(b)}E$이다. 한편 모든 $r$에 대해
$M_{\bullet,r,\bullet}$은 단순복합체 $K_{\bullet,r}$의
카르탕--아일렌베르크 분해이다. (M, XV, 2)의 계산에 의해
\[
H_q^{\mathrm I}(T(M_{\bullet,rs}))
=T(H_q^{\mathrm I}(M_{\bullet,rs})),
\]
따라서
\[
H_q^{\mathrm I}(T(U_{\bullet,j}))
=\bigoplus_{r+s=j}T(H_q^{\mathrm I}(M_{\bullet,rs})).
\]
다시 말해 단순복합체
$H_q^{\mathrm I}(T(U_{\bullet\bullet}))$은 이중복합체
$T(H_q^{\mathrm I}(M_{\bullet\bullet\bullet}))$에 대응하는
단순복합체이다. 그런데 모든 $q$에 대해
$H_q^{\mathrm I}(M_{\bullet\bullet\bullet})$은 범주 $\mathcal K$에서
단순복합체 $H_q^{\mathrm I}(K_{\bullet\bullet})$의 사영 분해이다.
(\hyperref[0.11.6.3-fr]{11.6.3})을 적용하면
\[
H_p^{\mathrm{II}}
\bigl(T(H_q^{\mathrm I}(M_{\bullet\bullet\bullet}))\bigr)
=L_pT(H_q^{\mathrm I}(K_{\bullet\bullet}))
\]
을 얻고, 따라서 스펙트럼 열 $\mathop{}^{(a')}E$를 얻는다. 끝으로
삼중복합체 $M_{\bullet\bullet\bullet}$의 정의에서 지표 $i$와 $j$의
역할을 서로 바꾸고 앞의 논증을 이 새 삼중복합체에 적용하면
$\mathop{}^{(c)}E$와 $\mathop{}^{(d)}E$를 얻는다.
\end{proof}

$L_\bullet T(K_{\bullet\bullet})$를 이중복합체
$K_{\bullet\bullet}$에 관한 $T$의 \emph{초호몰로지}라 한다.

\begin{env}[11.7.3]
\label{0.11.7.3-fr}
\emph{주석.}
\begin{enumerate}
\item[\rm(i)]
(\hyperref[0.11.6.3-fr]{11.6.3})에 의해
$L_\bullet T(K_{\bullet\bullet})$은 두 차수가 각각 아래로 유계인
$\mathcal C$의 이중복합체 $K_{\bullet\bullet}$들의 범주에서
호몰로지 함자이다.

\oldpage[0\textsubscript{III}]{43}

\item[\rm(ii)]
$M_{\bullet\bullet\bullet}$을 모든 쌍 $(r,s)$에 대해
$M_{rs,\bullet}$이 $K_{rs}$의 분해이고, 모든 삼중지표 $(i,j,k)$와
모든 $n>0$에 대해 $L_nT(M_{ijk})=0$인 $\mathcal C$의
삼중복합체라 하자. 그러면 동형
\[
L_\bullet T(K_{\bullet\bullet})
\simeq H_\bullet(T(M_{\bullet\bullet\bullet}))
\]
이 성립한다. 실제로 (\hyperref[0.11.7.2-fr]{11.7.2})의 증명에서
사용한 표기 아래에서 $N_{i,\bullet}$은 $L_i$의 분해이고 모든
$n>0$와 모든 쌍 $(i,j)$에 대해 $L_nT(N_{ij})=0$이므로,
(\hyperref[0.11.6.3-fr]{11.6.3}, (iii))을 적용하면 충분하다.

\item[\rm(iii)]
(\hyperref[0.11.7.2-fr]{11.7.2})의 결과는 공변 다변함자로 곧바로
일반화된다. 예를 들어 $\mathcal C'$을 모든 대상이 어떤 사영 대상의
몫인 또 하나의 아벨 범주, $K'_{\bullet\bullet}$을 두 차수가 아래로
유계인 $\mathcal C'$의 이중복합체, $T$를
$\mathcal C\times\mathcal C'$에서 아벨 범주 $\mathcal C''$으로
가는 공변 가법 쌍함자라 하자. $L_\bullet$과 $L'_\bullet$을 각각
$K_{\bullet\bullet}$과 $K'_{\bullet\bullet}$에 대응하는
단순복합체라 하면, 두 이중복합체 $K_{\bullet\bullet}$,
$K'_{\bullet\bullet}$에 관한 $T$의 초호몰로지를 초호몰로지
\[
L_\bullet T(L_\bullet,L'_\bullet)
\]
로 정의한다. (\hyperref[0.11.6.4-fr]{11.6.4})와
(\hyperref[0.11.6.5-fr]{11.6.5})를 적용하면
(\hyperref[0.11.7.2-fr]{11.7.2})에서와 같이 이 초호몰로지를 공통
수렴 대상으로 갖는 여섯 스펙트럼 열을 얻는다. 이 열들을 써 보는
일은 독자에게 맡긴다.
\end{enumerate}
\end{env}

\subsection{단체 복합체의 코호몰로지에 관한 보충}
\label{subsection:0.11.8-fr}

\begin{env}[11.8.1]
\label{0.11.8.1-fr}
$A$를 유한집합이라 하고, $\Sigma(A)$를 $A$의 원소들로 이루어진 유한열
$\sigma=(\alpha_0,\ldots,\alpha_h)$들의 집합($A$의 ``단체''들의 집합)이라 하자.
$|\sigma|=\{\alpha_0,\ldots,\alpha_h\}$로 놓는다. 사슬 복합체
$\mathrm C_\bullet(A)$는 $\Sigma(A)$의 원소들이 생성하는 등급 자유
아벨군이며, $(\alpha_0,\ldots,\alpha_h)$의 차수는 $h$이고, 그 미분은
다음과 같이 정의됨을 상기하자.
\[
d(\alpha_0,\ldots,\alpha_h)=
\sum_{j=0}^{h}(-1)^j
(\alpha_0,\ldots,\widehat{\alpha_j},\ldots,\alpha_h).
\]
$\mathrm C_\bullet(A)$에서 두 $\alpha_i$가 같은 사슬
$\sigma=(\alpha_0,\ldots,\alpha_h)$들과 다음 사슬들이 생성하는 부분군을
$\mathrm D_\bullet(A)$라 하자.
$\pi(\sigma)-\varepsilon_\pi\sigma$. 여기서 각 순열 $\pi$에 대하여
\[
\pi(\sigma)=(\alpha_{\pi(0)},\ldots,\alpha_{\pi(h)})
\]
이고 $\varepsilon_\pi$는 $\pi$의 부호이다. 그러면
$\mathrm D_\bullet(A)$는 $\mathrm C_\bullet(A)$의 부분복합체이다. 그
원소들을 \emph{퇴화 사슬}이라 부른다. 이제
$\mathrm L_\bullet(A)=\mathrm C_\bullet(A)/\mathrm D_\bullet(A)$로 놓으면
복합체의 자연 준동형
\[
p:\mathrm C_\bullet(A)\longrightarrow\mathrm L_\bullet(A)
\]
이 있다. 한편 복합체의 준동형
\[
j:\mathrm L_\bullet(A)\longrightarrow\mathrm C_\bullet(A)
\]
을 다음과 같이 정의한다. $A$에 전순서를 하나 정한다. 단체
$\sigma=(\alpha_0,\ldots,\alpha_h)$의 mod.~$\mathrm D_\bullet(A)$ 동치류에
대하여, 두 $\alpha_i$가 같으면 $0$을 대응시키고, 그렇지 않으면
$\alpha_i$들을 오름차순으로 배열한 열 $(\beta_0,\ldots,\beta_h)$를
대응시킨다. $p\circ j$가 $\mathrm L_\bullet(A)$의 항등사상임은 분명하다.
\end{env}

\begin{env}[11.8.2]
\label{0.11.8.2-fr}
$B$를 또 하나의 유한집합이라 하자. $d'$, $d''$가 각각
$\mathrm C_\bullet(A)$와 $\mathrm C_\bullet(B)$의 미분이면, 텐서곱 복합체
$\mathrm C_\bullet(A)\otimes\mathrm C_\bullet(B)$는
$\Sigma(A)\times\Sigma(B)$의 원소들이 생성하는 자유 아벨군으로 볼 수
있으며, 그 미분은 다음과 같다.
\[
d(\sigma,\tau)=(d'\sigma,\tau)+(-1)^{h+1}(\sigma,d''\tau)
\quad\text{$\operatorname{Card}|\sigma|=h+1$이면.}
\]
자연 준동형들
$\mathrm C_\bullet(A)\to\mathrm L_\bullet(A)$,
$\mathrm C_\bullet(B)\to\mathrm L_\bullet(B)$는 준동형
\[
p:\mathrm C_\bullet(A)\otimes\mathrm C_\bullet(B)
\longrightarrow\mathrm L_\bullet(A)\otimes\mathrm L_\bullet(B)
\]
을 정의한다. 뒤의 텐서곱은 다음과 동형이다.
\[
\frac{\mathrm C_\bullet(A)\otimes\mathrm C_\bullet(B)}
{\mathrm C_\bullet(A)\otimes\mathrm D_\bullet(B)
+\mathrm D_\bullet(A)\otimes\mathrm C_\bullet(B)}.
\]
마찬가지로, (\hyperref[0.11.8.1-fr]{11.8.1})에서 정의한 준동형
$\mathrm L_\bullet(A)\to\mathrm C_\bullet(A)$와
$\mathrm L_\bullet(B)\to\mathrm C_\bullet(B)$를($A$와 $B$에 전순서를
정하여) 이용하면, $p\circ j$가 항등사상이 되도록 하는 준동형
\[
j:\mathrm L_\bullet(A)\otimes\mathrm L_\bullet(B)
\longrightarrow\mathrm C_\bullet(A)\otimes\mathrm C_\bullet(B)
\]
을 정의할 수 있다.
\end{env}

\oldpage[0\textsubscript{III}]{44}
이 표기 아래에서 다음이 성립한다.
\begin{proposition}[11.8.3]
\label{0.11.8.3-fr}
\emph{호모토피
\[
h:\mathrm C_\bullet(A)\otimes\mathrm C_\bullet(B)
\longrightarrow\mathrm C_\bullet(A)\otimes\mathrm C_\bullet(B)
\]
가 존재하여, $h(\sigma,\tau)$는 $|\sigma_i|\subset|\sigma|$,
$|\tau_i|\subset|\tau|$를 만족하는 단체쌍 $(\sigma_i,\tau_i)$들의
선형결합이고, $f=j\circ p$에 대하여 다음이 성립한다.}
\[
\label{0.11.8.3.1-fr}
\tag{11.8.3.1}
f-1=h\circ d+d\circ h.
\]
\end{proposition}

\begin{proof}
$\sigma$와 $\tau$의 차수의 합에 관한 귀납법으로 각 단체쌍
$(\sigma,\tau)$ 위에서 $h$를 정의하면 충분하다. 그 합이 $0$이면
$h=0$으로 취할 수 있기 때문이다. 다음과 같이 놓자.
\[
\omega=f(\sigma,\tau)-(\sigma,\tau)-h(d(\sigma,\tau)).
\]
귀납 가정과 $d$의 정의에 따라
\[
\omega\in\mathrm C_\bullet(|\sigma|)\otimes\mathrm C_\bullet(|\tau|)
\]
이다. 또한
\[
d\omega=f(d(\sigma,\tau))-d(\sigma,\tau)-d(h(d(\sigma,\tau)))
=h(d(d(\sigma,\tau)))=0
\]
이다. 이는 (\hyperref[0.11.8.3.1-fr]{11.8.3.1})과 귀납 가정에서
따른다. 그런데 $q>0$이면 $H_q(\mathrm C_\bullet(A))=0$이고
(G, I, 3.7.4), 따라서 퀸네트 공식(G, I, 5.5.2)에 의해
\[
H_q(\mathrm C_\bullet(A)\otimes\mathrm C_\bullet(B))=0\quad(q>0)
\]
이기도 하다. 이 사실에서 $A$를 $|\sigma|$로, $B$를 $|\tau|$로
바꾸어 적용하면
\[
\omega'\in\mathrm C_\bullet(|\sigma|)\otimes\mathrm C_\bullet(|\tau|)
\]
이며 $\omega=d\omega'$인 원소가 존재함을 알 수 있다.
$h(\sigma,\tau)=\omega'$으로 취하면 단체쌍 $(\sigma,\tau)$에 대하여
(\hyperref[0.11.8.3.1-fr]{11.8.3.1})을 확인할 수 있으므로 귀납을 계속할
수 있다.
\end{proof}

\begin{env}[11.8.4]
\label{0.11.8.4-fr}
표기는 (\hyperref[0.11.8.2-fr]{11.8.2})에서와 같이 하자.
$|\sigma|\subset|\sigma'|$이고 $|\tau|\subset|\tau'|$일 때
$(\sigma,\tau)\leq(\sigma',\tau')$로 놓는다.
$\Sigma(A)\times\Sigma(B)$ 위의 \emph{계수계} $\Gamma$는 다음 데이터로
이루어진다. 첫째, $\Gamma_{\sigma,\tau}$가 오직 집합 $|\sigma|$와
$|\tau|$에만 의존하는 아벨군들의 족 $(\Gamma_{\sigma,\tau})$가 있고,
둘째,
\[
\Gamma_{\sigma,\tau}\longrightarrow\Gamma_{\sigma',\tau'}
\quad\text{$(\sigma,\tau)\leq(\sigma',\tau')$이면}
\]
인 준동형들의 족이 있어서 이 전순서 관계에 관한 귀납계를 이룬다. 이제
공사슬 복합체 $\mathrm C^\bullet(A,B;\Gamma)$를 다음과 같이 정의한다.
그 원소는 $\Sigma(A)\times\Sigma(B)$를 움직이는
$(\sigma,\tau)$에 대한 족 $\lambda=(\lambda(\sigma,\tau))$이며, 각 쌍에
대하여 $\lambda(\sigma,\tau)\in\Gamma_{\sigma,\tau}$이다. 미분은 다음과
같이 준다. 만일
\[
d(\sigma,\tau)=\sum_i\pm(\sigma_i,\tau_i)
\]
이면, 모든 $i$에 대하여 $|\sigma_i|\subset|\sigma|$이고
$|\tau_i|\subset|\tau|$이며,
\[
d\lambda(\sigma,\tau)=\sum_i\pm\lambda_i(\sigma_i,\tau_i)
\]
로 취한다. 여기서 $\lambda_i(\sigma_i,\tau_i)$는
$\lambda(\sigma_i,\tau_i)$의 $\Gamma_{\sigma,\tau}$ 안에서의 표준 상을
나타낸다.

공사슬 $\lambda\in\mathrm C^\bullet(A,B;\Gamma)$가 다음 조건을 만족하면
\emph{이중 교대적}이라고 하겠다. 두 단체 $\sigma$, $\tau$ 중 하나에
같은 항이 두 개 있으면 $\lambda(\sigma,\tau)=0$이고, 지표들의 임의의
순열 $\pi$, $\pi'$에 대하여
\[
\lambda(\pi(\sigma),\tau)=\varepsilon_\pi\lambda(\sigma,\tau),
\qquad
\lambda(\sigma,\pi'(\tau))=\varepsilon_{\pi'}\lambda(\sigma,\tau)
\]
이다. 이러한 공사슬들이 $\mathrm C^\bullet(A,B;\Gamma)$의 부분복합체
$\mathrm L^\bullet(A,B;\Gamma)$를 이룸은 분명하다.
\end{env}

\begin{proposition}[11.8.5]
\label{0.11.8.5-fr}
\emph{표준 단사 준동형
\[
\mathrm L^\bullet(A,B;\Gamma)\longrightarrow
\mathrm C^\bullet(A,B;\Gamma)
\]
은 이 두 복합체의 코호몰로지 사이의 동형을 정의한다.}
\end{proposition}

\begin{proof}
$p$와 $j$를 (\hyperref[0.11.8.2-fr]{11.8.2})에서 정의한 뜻으로 쓰면,
사상들
\[
{}'p:\lambda\longmapsto\lambda\circ p,\qquad
{}'j:\lambda\longmapsto\lambda\circ j
\]
은 각각 $\mathrm L^\bullet(A,B;\Gamma)$와
$\mathrm C^\bullet(A,B;\Gamma)$에서 정의되며, 첫째 것은 바로 표준 단사
준동형이다. $\mathop{}'j\circ{}'p$가 항등사상이므로
$\mathop{}'p\circ{}'j$가 항등사상과 호모토픽임을 보이면 충분하다. 그런데
(\hyperref[0.11.8.3-fr]{11.8.3})에 따라
$\mathop{}'h:\lambda\mapsto\lambda\circ h$가
$\mathrm C^\bullet(A,B;\Gamma)$에서 정의되므로, 항등식
(\hyperref[0.11.8.3.1-fr]{11.8.3.1})을 전치하여 원하는 결과를 얻는다.
\end{proof}

\begin{env}[11.8.6]
\label{0.11.8.6-fr}
\oldpage[0\textsubscript{III}]{45}
명제 (\hyperref[0.11.8.5-fr]{11.8.5})에 의해
$\mathrm L^\bullet(A,B;\Gamma)$의 코호몰로지 계산은
$\mathrm C^\bullet(A,B;\Gamma)$의 코호몰로지 계산으로 귀착된다. 한편
아일렌베르크--질버 정리(G, I, 3.10.2)에 따라 뒤의 코호몰로지는 다음과
같이 정의되는 공사슬 복합체의 코호몰로지와 표준적으로 동형임을 상기하자.
사슬 복합체 $\mathrm P_\bullet(A,B)$를 구성한다. 이는 $\sigma$와 $\tau$의
차수가 같은 $(\sigma,\tau)\in\Sigma(A)\times\Sigma(B)$들의 선형결합으로
이루어지며, 그 미분은
\[
d(\sigma,\tau)=\sum_{j,k}(-1)^{j+k}(\sigma_j,\tau_k)
\]
이다. 여기서
\[
d\sigma=\sum_j(-1)^j\sigma_j,\qquad
d\tau=\sum_k(-1)^k\tau_k
\]
이다. 이때 복합체의 표준 준동형 두 개
\[
f:\mathrm P_\bullet(A,B)\longrightarrow
\mathrm C_\bullet(A)\otimes\mathrm C_\bullet(B),\qquad
g:\mathrm C_\bullet(A)\otimes\mathrm C_\bullet(B)
\longrightarrow\mathrm P_\bullet(A,B)
\]
가 있으며, 호모토피 $h$, $h'$이 존재하여 다음을 만족함이
증명되어 있다(\emph{인용한 곳}).
\[
f\circ g-1=d\circ h+h\circ d,\qquad
g\circ f-1=d\circ h'+h'\circ d.
\]
더욱이
\[
f(\sigma,\tau)\in\mathrm C_\bullet(|\sigma|)
\otimes\mathrm C_\bullet(|\tau|),\qquad
g(\sigma,\tau)\in\mathrm P_\bullet(|\sigma|,|\tau|)
\]
이고, 호모토피 $h$, $h'$은 다음을 만족하도록 취할 수 있다.
\[
h(\sigma,\tau)\in\mathrm C_\bullet(|\sigma|)
\otimes\mathrm C_\bullet(|\tau|),\qquad
h'(\sigma,\tau)\in\mathrm P_\bullet(|\sigma|,|\tau|).
\]
이는 $h(\sigma,\tau)$와 $h'(\sigma,\tau)$를 $\sigma$와 $\tau$의 차수의
합에 관한 귀납법으로 정의할 수 있고,
\[
H_q(\mathrm C_\bullet(|\sigma|)\otimes\mathrm C_\bullet(|\tau|))
\quad\text{와}\quad
H_q(\mathrm P_\bullet(|\sigma|,|\tau|))
\]
가 $q>0$일 때 $0$이기 때문이다(\emph{인용한 곳}). 그러면
(\hyperref[0.11.8.3-fr]{11.8.3})에서와 같이 논증하여 결론을 얻는다.

이제 $\mathrm P^\bullet(A,B;\Gamma)$를 다음 족들의 집합으로 정의한다.
$(\sigma,\tau)$는 두 항의 차수가 같은 쌍들을 움직이고,
$\lambda=(\lambda(\sigma,\tau))$이며
$\lambda(\sigma,\tau)\in\Gamma_{\sigma,\tau}$이다. 또한
\[
d\sigma=\sum_j(-1)^j\sigma_j,\qquad
d\tau=\sum_k(-1)^k\tau_k
\]
이므로
\[
d\lambda(\sigma,\tau)=
\sum_{j,k}(-1)^{j+k}\lambda(\sigma_j,\tau_k)
\]
가 정의되며, 이것이 복합체 $\mathrm P^\bullet(A,B;\Gamma)$의 미분을
준다. 이때 앞의 고찰에 따라 사상들
\[
{}'f:\lambda\mapsto\lambda\circ f,\quad
{}'g:\lambda\mapsto\lambda\circ g,\quad
{}'h:\lambda\mapsto\lambda\circ h,\quad
{}'h':\lambda\mapsto\lambda\circ h'
\]
은 모두 정의된다. 따라서 $\mathop{}'f\circ{}'g$와
$\mathop{}'g\circ{}'f$는 항등사상과 호모토픽이고, 이로부터
$\mathrm C^\bullet(A,B;\Gamma)$의 코호몰로지와
$\mathrm P^\bullet(A,B;\Gamma)$의 코호몰로지 사이의 원하는 동형을
얻는다.
\end{env}

\begin{env}[11.8.7]
\label{0.11.8.7-fr}
\emph{주의}. (\hyperref[0.11.8.3-fr]{11.8.3})에서와 같은 논증을
$\mathrm C_\bullet(A)$와 $\mathrm L_\bullet(A)$에 적용하면 다음을 알 수
있다. $j$와 $p$를 (\hyperref[0.11.8.1-fr]{11.8.1})에서와 같이 정의할 때
$f=j\circ p$는 다시 (\hyperref[0.11.8.3.1-fr]{11.8.3.1})의 관계식을
만족하고, $|h(\sigma)|\subset|\sigma|$이다. 따라서
(\hyperref[0.11.8.5-fr]{11.8.5})에서와 같이, 자명한 방식으로 정의되는 두
복합체 $\mathrm L^\bullet(A;\Gamma)$와
$\mathrm C^\bullet(A;\Gamma)$의 코호몰로지 사이의 동형을 얻는다. 이것이
(G, I, 3.8.1)에서 증명이 개략적으로 제시된 결과이다.
\end{env}

\begin{env}[11.8.8]
\label{0.11.8.8-fr}
이제 (\hyperref[0.11.8.2-fr]{11.8.2})의 표기와 가정으로 돌아가서,
$\Sigma(A)\times\Sigma(B)$ 위의 계수계들의 복합체
$\Gamma^\bullet=(\Gamma^k)$를 생각하자. 따라서 각 $(\sigma,\tau)$에
대하여 $\Gamma^k_{\sigma,\tau}$들은 아벨군의 복합체를 이루며
$(k\in\mathbb Z)$, $(\sigma,\tau)\leq(\sigma',\tau')$일 때 다음 가환
도표들이 있다.
\[
\xymatrix{
\Gamma^k_{\sigma,\tau}\ar[r]\ar[d] &
\Gamma^{k+1}_{\sigma,\tau}\ar[d]\\
\Gamma^k_{\sigma',\tau'}\ar[r] &
\Gamma^{k+1}_{\sigma',\tau'}
}
\]
\oldpage[0\textsubscript{III}]{46}
그러면
\[
\mathrm C^\bullet(A,B;\Gamma^\bullet)
=\bigl(\mathrm C^h(A,B;\Gamma^k)\bigr)_{(h,k)\in\mathbb Z\times\mathbb Z}
\]
가 아벨군의 이중복합체이고,
\[
\mathrm L^\bullet(A,B;\Gamma^\bullet)
=\bigl(\mathrm L^h(A,B;\Gamma^k)\bigr)
\]
가 그 부분 이중복합체임을 곧바로 확인할 수 있다.
\end{env}

\begin{proposition}[11.8.9]
\label{0.11.8.9-fr}
\emph{표준 단사 준동형
\[
\mathrm L^\bullet(A,B;\Gamma^\bullet)\longrightarrow
\mathrm C^\bullet(A,B;\Gamma^\bullet)
\]
은 이 두 이중복합체의 코호몰로지 사이의 동형을 정의한다.}
\end{proposition}

\begin{proof}
간단히 쓰기 위해 $\mathrm C^{\bullet\bullet}=
\mathrm C^\bullet(A,B;\Gamma^\bullet)$,
$\mathrm L^{\bullet\bullet}=
\mathrm L^\bullet(A,B;\Gamma^\bullet)$로 놓자.
$h<0$이면 $\mathrm C^{hk}=\mathrm L^{hk}=0$이므로 이 이중복합체들의 둘째
스펙트럼 열들은 정칙이다
(\hyperref[0.11.3.3-fr]{11.3.3}). 따라서 준동형
$\mathrm L^{\bullet\bullet}\to\mathrm C^{\bullet\bullet}$은 스펙트럼 열의
사상
\[
{}''E(\mathrm L^{\bullet\bullet})\longrightarrow
{}''E(\mathrm C^{\bullet\bullet})
\]
을 주고, $E_2$ 항에서는 이것이 다음으로 귀착된다.
\[
\label{0.11.8.9.1-fr}
\tag{11.8.9.1}
H^p_{\mathrm{II}}(H^q_{\mathrm I}(\mathrm L^{\bullet\bullet}))
\longrightarrow
H^p_{\mathrm{II}}(H^q_{\mathrm I}(\mathrm C^{\bullet\bullet})).
\]
그런데 모든 $k\in\mathbb Z$에 대하여
(\hyperref[0.11.8.5-fr]{11.8.5})로부터 준동형
$H^q_{\mathrm I}(\mathrm L^{\bullet,k})\to
H^q_{\mathrm I}(\mathrm C^{\bullet,k})$가 전단사임을 알 수 있다. 따라서
(\hyperref[0.11.1.5-fr]{11.1.5})로부터 결론이 따른다.
\end{proof}

\begin{env}[11.8.10]
\label{0.11.8.10-fr}
마찬가지로 (\hyperref[0.11.8.6-fr]{11.8.6})의 표기 아래에서 이중복합체의
표준 준동형들
\[
\mathrm C^\bullet(A,B;\Gamma^\bullet)\longrightarrow
\mathrm P^\bullet(A,B;\Gamma^\bullet)
\]
이 있다(표기는 자명하다). 이번에는
(\hyperref[0.11.8.6-fr]{11.8.6})에 근거하여
(\hyperref[0.11.8.9-fr]{11.8.9})에서와 같이 논증하면, 이 준동형도
코호몰로지 사이의 동형을 줌을 알 수 있다.
\end{env}

\subsection{유한 생성 복합체에 관한 보조정리}
\label{subsection:0.11.9-fr}

\begin{proposition}[11.9.1]
\label{0.11.9.1-fr}
$\mathcal C$를 아벨 범주라 하고, $K'$과 $K''$을 $\mathcal C$의 대상들의
집합의 부분집합이라 하자. $K''\subset K'$이고 다음 조건들을 만족한다고
하자.
\begin{enumerate}
\item[{\rm(i)}] \emph{모든 대상 $A'\in K'$과 $\mathcal C$ 안의 모든
에피사상 $u:A\to A'$에 대하여, 어떤 대상 $B\in K''$과 사상
$v:B\to A$가 존재하여 $u\circ v$가 에피사상이다.}
\item[{\rm(ii)}] \emph{모든 대상쌍 $A\in K'$, $B\in K'$과 모든
에피사상 $u:A\to B$에 대하여 $\operatorname{Ker}(u)$는 $K'$에 속한다.}
\item[{\rm(iii)}] \emph{$K''$의 두 대상의 곱은 $K''$에 속한다.}
\end{enumerate}
\emph{$P_\bullet=(P_i)_{i\in\mathbb Z}$를 $\mathcal C$ 안의 복합체라 하자.
모든 $i$에 대하여 $H_i(P_\bullet)\in K'$이고, 어떤 $d$가 존재하여
$i<d$일 때 $H_i(P_\bullet)=0$이라고 하자. 그러면 $\mathcal C$ 안의
복합체 $Q_\bullet=(Q_i)_{i\in\mathbb Z}$와 복합체의 사상
$u:Q_\bullet\to P_\bullet$가 존재하여, 모든 $i$에 대하여 $Q_i\in K''$이고
$i<d$일 때 $Q_i=0$이며, 이에 대응하는 사상
$H_\bullet(Q_\bullet)\to H_\bullet(P_\bullet)$은 동형이다.}
\end{proposition}

\begin{proof}
먼저 성질 (i)에서 따르는 다음 결과를 증명하자.

\emph{(i bis) $u:C\to B$를 $\mathcal C$ 안의 에피사상, $A$를 $K'$의
대상, $v:A\to B$를 $\mathcal C$ 안의 사상이라 하자. 그러면 대상
$D\in K''$, 에피사상 $u':D\to A$, 사상 $v':D\to C$가 존재하여 도표}
\[
\xymatrix{
D\ar[r]^{u'}\ar[d]_{v'} & A\ar[d]^{v}\\
C\ar[r]_{u} & B
}
\tag{11.9.1.1}\label{0.11.9.1.1-fr}
\]
\emph{가 가환이다.}

실제로 $\mathcal C$ 안의 섬유곱 $C\times_B A$를 생각하고, 도표
\oldpage[0\textsubscript{III}]{47}
\[
\xymatrix{
C\times_B A\ar[r]^{q}\ar[d]_{p} & A\ar[d]^{v}\\
C\ar[r]_{u} & B
}
\tag{11.9.1.2}\label{0.11.9.1.2-fr}
\]
가 가환이 되게 하는 표준 사영 $p:C\times_B A\to C$,
$q:C\times_B A\to A$를 취하자. $q$의 여핵은
$v^{-1}(u(C))$에 의한 $A$의 몫임이 알려져 있다([27], p.~1-12).
$u$가 에피사상이므로 $u(C)=B$이고 $v^{-1}(u(C))=A$이다. 따라서 $q$는
에피사상이다. 이제 에피사상 $q:C\times_B A\to A$에 (i)를 적용하면,
어떤 대상 $D\in K''$과 사상 $w:D\to C\times_B A$가 존재하여
$q\circ w$가 에피사상이다. $u'=q\circ w$, $v'=p\circ w$로 취하면 된다.

이제 명제를 증명하기 위해 귀납법을 쓰자. 어떤 $i\geq d-1$에 대하여,
$j\leq i$인 대상 $Q_j$, 사상 $d_j:Q_j\to Q_{j-1}$, 사상
$u_j:Q_j\to P_j$가 이미 구성되어 다음을 만족한다고 하자. $j<d$이면
$Q_j=0$이고, $j\leq i$에 대하여 $d_{j-1}\circ d_j=0$이며
$d_j\circ u_j=u_{j-1}\circ d_j$이다. 더욱이 다음 조건들이 성립한다고
가정한다.
\begin{enumerate}
\item[$(\mathrm I_i)$] $j\leq i$이면 $Q_j\in K''$이고, $j<i$이면
$B_j(Q_\bullet)\in K'$이다.
\item[$(\mathrm{II}_i)$] $j<i$이면, 족 $(u_k)_{k\leq i}$에서 유도되는
준동형 $H_j(Q_\bullet)\to H_j(P_\bullet)$은 동형이다.
\item[$(\mathrm{III}_i)$] 합성사상
\[
v_i:Z_i(Q_\bullet)\longrightarrow Z_i(P_\bullet)
\longrightarrow H_i(P_\bullet)
\]
은 에피사상이다. 여기서 왼쪽 화살표는 $u_i$의 제한이고 오른쪽 화살표는
표준 사상이다.
\end{enumerate}
(ii)에 따라, 에피사상 $Q_i\to B_{i-1}(Q_\bullet)$의 핵
$Z_i(Q_\bullet)$은 가정 $(\mathrm I_i)$에 의해 $K'$에 속한다. 또한 (ii)와
가정 $(\mathrm{III}_i)$에 따라 $N_i=\operatorname{Ker}(v_i)$도 $K'$에
속한다. (i bis)에 의해 $Q'_{i+1}\in K''$, 에피사상
$d'_{i+1}:Q'_{i+1}\to N_i$, 사상 $u'_{i+1}:Q'_{i+1}\to P_{i+1}$이
존재하여 도표
\[
\xymatrix{
Q'_{i+1}\ar[r]^{d'_{i+1}}\ar[d]_{u'_{i+1}} &
N_i\ar[d]^{u_i}\\
P_{i+1}\ar[r]_{d_{i+1}} & B_i(P_\bullet)
}
\tag{11.9.1.3}\label{0.11.9.1.3-fr}
\]
가 가환이다.

표준 사상 $Z_{i+1}(P_\bullet)\to H_{i+1}(P_\bullet)$은 에피사상이고,
가정에 따라 $H_{i+1}(P_\bullet)\in K'$이다. 따라서 (i)에 의해 대상
$Q''_{i+1}\in K''$과 사상
$u''_{i+1}:Q''_{i+1}\to Z_{i+1}(P_\bullet)$이 존재하여 합성사상
$Q''_{i+1}\to Z_{i+1}(P_\bullet)\to H_{i+1}(P_\bullet)$이
에피사상이다. $d''_{i+1}:Q''_{i+1}\to N_i$를 $0$으로 취하면 도표
\[
\xymatrix{
Q''_{i+1}\ar[r]^{d''_{i+1}}\ar[d]_{u''_{i+1}} &
N_i\ar[d]^{u_i}\\
Z_{i+1}(P_\bullet)\ar[r]_{d_{i+1}} & P_i
}
\tag{11.9.1.4}\label{0.11.9.1.4-fr}
\]
는 가환이다. 아래쪽 가로 화살표가 $0$이기 때문이다.
\oldpage[0\textsubscript{III}]{48}
이제 다음과 같이 취하자.
\[
Q_{i+1}=Q'_{i+1}\times Q''_{i+1},\qquad
d_{i+1}=d'_{i+1}+d''_{i+1},\qquad
u_{i+1}=u'_{i+1}+u''_{i+1}.
\]
$d_{i+1}(Q_{i+1})=d'_{i+1}(Q'_{i+1})=N_i\subset Z_i(Q_\bullet)$이므로
$d_i\circ d_{i+1}=0$이고, 통상적인 표기 아래에서
$B_i(Q_\bullet)=N_i$이다. 따라서 $(\mathrm I_{i+1})$이 확인된다. 도표
(\hyperref[0.11.9.1.3-fr]{11.9.1.3})과
(\hyperref[0.11.9.1.4-fr]{11.9.1.4})의 가환성에 따라
$d_{i+1}\circ u_{i+1}=u_i\circ d_{i+1}$이다. $N_i$의 정의에 의해,
$u_k$들의 계($k\leq i+1$)에서 유도되는 사상
\[
H_i(Q_\bullet)=Z_i(Q_\bullet)/N_i\longrightarrow H_i(P_\bullet)
\]
은 몫으로 내려서 $v_i$에서 유도되는 사상이다. $v_i$가 에피사상이므로
이 사상은 동형이고, 따라서 $(\mathrm{II}_{i+1})$이 성립한다. 마지막으로
정의에 의해 $Q''_{i+1}\subset Z_{i+1}(Q_\bullet)$이다.
$u''_{i+1}$의 선택에 따라 사상
\[
v_{i+1}:Z_{i+1}(Q_\bullet)\longrightarrow Z_{i+1}(P_\bullet)
\longrightarrow H_{i+1}(P_\bullet)
\]
은 에피사상이다. 실제로 $Q''_{i+1}$에 대한 그 제한이 이미
에피사상이다. 따라서 $(\mathrm{III}_{i+1})$이 성립한다. 그러므로 귀납을
계속할 수 있고, 명제가 증명된다.
\end{proof}

\begin{corollary}[11.9.2]
\label{0.11.9.2-fr}
$A$를 뇌터 환이라 하자(가환일 필요는 없다). $P_\bullet=(P_i)_{i\in\mathbb Z}$를
오른쪽 $A$-가군들의 복합체라 하자. $H_i(P_\bullet)$들이 유한 생성
$A$-가군이고, 어떤 $d$가 존재하여 $i<d$일 때 $H_i(P_\bullet)=0$이라고
가정하자. 그러면 유한 계수의 자유 오른쪽 $A$-가군들로 이루어진 복합체
$Q_\bullet=(Q_i)_{i\in\mathbb Z}$가 존재하여 $i<d$일 때 $Q_i=0$이고,
복합체의 준동형 $u:Q_\bullet\to P_\bullet$가 존재하여 $u$에 대응하는
준동형 $H_\bullet(Q_\bullet)\to H_\bullet(P_\bullet)$은 전단사이다.
\end{corollary}

\begin{proof}
(\hyperref[0.11.9.1-fr]{11.9.1})에서 $\mathcal C$를 오른쪽
$A$-가군들의 범주로, $K'$을 유한 생성 $A$-가군들의 집합으로, $K''$을
유한 계수의 자유 $A$-가군들의 집합으로 취하여 적용한다.
$A$가 뇌터라는 가정에 유의하면
(\hyperref[0.11.9.1-fr]{11.9.1})의 조건 (i), (ii), (iii)은 곧바로
확인된다.
\end{proof}

\begin{env}[11.9.3]
\label{0.11.9.3-fr}
\emph{주의}.
\begin{enumerate}
\item[{\rm(i)}] (\hyperref[0.11.9.2-fr]{11.9.2})의 조건 아래에서
$P_i$들이 평탄 오른쪽 $A$-가군이라고 더 가정하자. 그러면 모든 왼쪽
$A$-가군 $M$에 대하여 복합체의 준동형
\[
u\otimes 1:Q_\bullet\otimes_A M\longrightarrow P_\bullet\otimes_A M
\]
은 여전히 호몰로지의 동형
$H_\bullet(Q_\bullet\otimes_A M)\simeq
H_\bullet(P_\bullet\otimes_A M)$을 정의한다. 이는 제III장에서 보게 될
것이다.
\item[{\rm(ii)}] $A$가 뇌터라고 가정하지 않으면
(\hyperref[0.11.9.2-fr]{11.9.2})의 결론은 더 이상 반드시 참인 것은
아니다. 실제로 한 항을 제외하고는 모두 $0$인 복합체에 이를 적용하면,
모든 유한 생성 오른쪽 $A$-가군이 유한 생성 자유 가군들에 의한 분해를
가진다는 결론을 얻게 되지만, 이는 일반적으로 참이 아니다
(cf. Bourbaki, \emph{Alg. comm.}, chap.~I, \S~2, exerc.~6).

그러나 $A$가 뇌터라고 가정하는 대신 $H_i(P_\bullet)$들이 유한
$\infty$-표시를 갖는다고만 가정할 수도 있다(cf. chap.~IV).
\end{enumerate}
\end{env}

\subsection{유한 길이 가군 복합체의 오일러--푸앵카레 특성}
\label{subsection:0.11.10-fr}

\begin{env}[11.10.1]
\label{0.11.10.1-fr}
$A$를 환이라 하자(가환일 필요는 없다). 유한 길이의 왼쪽
$A$-가군들로 이루어진 유한 복합체
\[
\label{0.11.10.1.1-fr}
\tag{11.10.1.1}
M^\bullet:0\longrightarrow M^0\longrightarrow M^1
\longrightarrow\cdots\longrightarrow M^n\longrightarrow0
\]
를 생각하자. 이 복합체의 \emph{오일러--푸앵카레 특성}은 다음 수로
정의된다.
\oldpage[0\textsubscript{III}]{49}
\[
\label{0.11.10.1.2-fr}
\tag{11.10.1.2}
\chi(M^\bullet)=\sum_{i=0}^{n}(-1)^i\operatorname{long}(M^i).
\]
\end{env}

\begin{proposition}[11.10.2]
\label{0.11.10.2-fr}
\emph{유한 길이의 왼쪽 $A$-가군들로 이루어진 모든 유한 복합체
$M^\bullet$에 대하여
\[
\chi(M^\bullet)=\chi(H^\bullet(M^\bullet))
\]
이다. 여기서 $H^\bullet(M^\bullet)$은 영 미분을 갖는 복합체로 본다.
특히 열 (\hyperref[0.11.10.1.1-fr]{11.10.1.1})이 완전하면
$\chi(M^\bullet)=0$이다.}
\end{proposition}

\begin{proof}
간단히 $B^i=B^i(M^\bullet)$, $Z^i=Z^i(M^\bullet)$,
$H^i=H^i(M^\bullet)=Z^i/B^i$로 놓자. $B^i$, $Z^i$, $H^i$는 유한
길이이다. 완전열
\[
0\longrightarrow B^i\longrightarrow Z^i\longrightarrow H^i
\longrightarrow0,
\qquad
0\longrightarrow Z^i\longrightarrow M^i\longrightarrow B^{i+1}
\longrightarrow0
\]
로부터 관계식
\[
\operatorname{long}(Z^i)=\operatorname{long}(H^i)+
\operatorname{long}(B^i),
\qquad
\operatorname{long}(M^i)=\operatorname{long}(Z^i)+
\operatorname{long}(B^{i+1})
\]
을 얻는다. 따라서
\[
\operatorname{long}(M^i)-\operatorname{long}(H^i)=
\operatorname{long}(B^{i+1})+\operatorname{long}(B^i).
\]
이 관계식에 $(-1)^i$를 곱하고 $0\leq i\leq n$에 대하여 합하자.
$B^0=B^{n+1}=0$임에 유의하면 원하는 등식을 얻는다.
\end{proof}

\begin{corollary}[11.10.3]
\label{0.11.10.3-fr}
$E=(E_r^{pq})$를 환 $A$ 위의 가군들의 범주에서의 스펙트럼 열이라 하자.
$E_2^{pq}$들이 유한 길이 $A$-가군이고, $E_2^{pq}\neq0$인 쌍 $(p,q)$는
유한 개뿐이라고 가정하자. 그러면 모든 복합체
$E_r^\bullet=(E_r^{(n)})_{n\in\mathbb Z}$
(\hyperref[0.11.1.1-fr]{11.1.1})의 오일러--푸앵카레 특성
$\chi(E_r^\bullet)$들은 모두 같다. 더 나아가 스펙트럼 열 $E$가
약수렴하고
\[
E_\infty^{(n)}=\bigoplus_{p+q=n}E_\infty^{pq}\quad(n\in\mathbb Z)
\]
로 놓으면 $\chi(E_\infty^\bullet)=\chi(E_2^\bullet)$이기도 하다. 여기서
$E_\infty^\bullet=(E_\infty^{(n)})_{n\in\mathbb Z}$는 영 미분을 갖는
복합체로 본다.
\end{corollary}

\begin{proof}
먼저 $E_2^{pq}=0$이면 $2\leq r<+\infty$에 대하여 $E_r^{pq}=0$임에
유의하자. 따라서 모든 복합체 $E_r^\bullet$은 유한하고 유한 길이
$A$-가군들로 이루어진다. 모든 유한한 $r$에 대한 관계식
$\chi(E_r^\bullet)=\chi(E_{r+1}^\bullet)$은
(\hyperref[0.11.10.2-fr]{11.10.2})와, 영 미분을 갖는 복합체로서의
$H^\bullet(E_r^\bullet)$과 $E_{r+1}^\bullet$ 사이의 동형으로부터
따른다. $E_r^{pq}$들이 유한 길이라는 가정에 따라 모든 쌍 $(p,q)$에
대하여 열들
$(B_k(E_2^{pq}))_{k\geq2}$와 $(Z_k(E_2^{pq}))_{k\geq2}$는 각각 어느
단계부터 일정하다. 따라서 $E$가 약수렴한다는 가정과 유한 개의 쌍
$(p,q)$를 제외하면 $E_2^{pq}=0$이라는 가정에 의해, 모든 쌍 $(p,q)$에
대하여 $E_\infty^{pq}=E_r^{pq}$인 어떤 정수 $r\geq2$가 존재한다. 이로부터
$\chi(E_\infty^\bullet)$에 관한 주장이 따른다.
\end{proof}

\section{층의 코호몰로지에 관한 보충}
\label{section:0.12-fr}

\subsection{환 달린 공간 위의 가군층의 코호몰로지}
\label{subsection:0.12.1-fr}

\begin{env}[12.1.1]
\label{0.12.1.1-fr}
$(X,\mathcal O_X)$를 환 달린 공간이라 하자. 임의의
$\mathcal O_X$-가군 $\mathcal F$에 대하여 코호몰로지
$H^\bullet(X,\mathcal F)$를 정의함을 상기하자. 이는
$\mathcal O_X$-가군들의 범주 $C(X)$에서 아벨군들의 범주로 가는
보편 코호몰로지 함자 (T, 2.2)이며, 왼쪽 완전 함자
$\mathcal F\to\Gamma(X,\mathcal F)$의 유도함자이다. 함자
$H^\bullet(X,\mathcal F)$는
\oldpage[0\textsubscript{III}]{50}
$X$ 위의 \emph{아벨군의 층들}의 범주에서 같은 방식으로 정의되는
코호몰로지 함자 (G, II, 7.2.1)를 범주 $C(X)$로 제한한 것과 동형이다.
\end{env}

\begin{env}[12.1.2]
\label{0.12.1.2-fr}
$A=\Gamma(X,\mathcal O_X)$로 놓자. $A$의 각 원소는 아벨군
$\Gamma(X,\mathcal F)$의 자기준동형을 정의하므로, 함자성에 의해
$H^\bullet(X,\mathcal F)$의 $\delta$-함자의 자기준동형을 정의한다.
이 자기준동형들은 각 $H^p(X,\mathcal F)$에 $A$-가군 구조를 정의하며,
연산자 $\partial$는 $A$-선형이다. 또한 임의의 양의 정수 $p$, $q$와
두 $\mathcal O_X$-가군 $\mathcal F$, $\mathcal G$에 대하여
\emph{컵곱}이라 하는 $A$-가군 준동형
\[
\label{0.12.1.2.1-fr}
\tag{12.1.2.1}
H^p(X,\mathcal F)\otimes_A H^q(X,\mathcal G)
\longrightarrow H^{p+q}(X,\mathcal F\otimes_{\mathcal O_X}\mathcal G)
\]
이 있다 (G, II, 6.6). 이 준동형들에 의해 $H^p(X,\mathcal O_X)$
($p\geq0$)들의 직합 $S$는 반가환 등급 $A$-대수가 되고,
$H^p(X,\mathcal F)$들의 직합은 등급 $S$-가군이 된다.

$X$의 임의의 \emph{열린 덮개} $\mathfrak U$에 대하여, 앞으로
$\check C^\bullet(\mathfrak U,\mathcal F)$는 (G, II, 5.1)과 달리 항상
$\Gamma(U_\alpha,\mathcal F)$를 계수계로 하는 $\mathfrak U$의 신경의
\emph{교대} 공사슬 복합체를 나타낸다. 분명
$\check C^\bullet(\mathfrak U,\mathcal F)$는 등급 $A$-가군이므로,
이 복합체의 코호몰로지군 $\check H^p(\mathfrak U,\mathcal F)$에는
$A$-가군 구조가 주어진다. 또한 표준 사상
\[
\check H^p(\mathfrak U,\mathcal F)\longrightarrow H^p(X,\mathcal F)
\]
(G, II, 5.4)는 $A$-선형이다.
\end{env}

\begin{env}[12.1.3]
\label{0.12.1.3-fr}
$(X',\mathcal O_{X'})$를 두 번째 환 달린 공간이라 하고,
$f=(\psi,\theta)$를 $X'$에서 $X$로 가는 사상이라 하자.

$A'=\Gamma(X',\mathcal O_{X'})$로 놓자. $\psi$-사상 $\theta$는 표준적으로
환 준동형 $A\to A'$을 정의한다. $\mathcal F$를 $\mathcal O_X$-가군,
$\mathcal F'$을 $\mathcal O_{X'}$-가군이라 하자. 임의의 $f$-사상
$u:\mathcal F\to\mathcal F'$
(\hyperref[0.4.4.1-ko]{0\textsubscript{I}, 4.4.1})에 대하여, 모든
$p\geq0$에 대해 \emph{디-준동형}
\[
\label{0.12.1.3.1-fr}
\tag{12.1.3.1}
u_p:H^p(X,\mathcal F)\longrightarrow H^p(X',\mathcal F')
\]
을 정의할 수 있음을 보이겠다. 실제로 $\psi^*$는 $X$ 위의 아벨군의
층들의 범주에서 완전하므로,
$\mathcal F\to H^\bullet(X',\psi^*(\mathcal F))$는 이 범주에서
$\delta$-함자이며, 표준 $\delta$-함자 준동형
\[
\label{0.12.1.3.2-fr}
\tag{12.1.3.2}
H^\bullet(X,\mathcal F)\longrightarrow H^\bullet(X',\psi^*(\mathcal F))
\]
이 있음을 안다. 이는 차수 0에서 표준 준동형
$\Gamma(X,\mathcal F)\to\Gamma(X',\psi^*(\mathcal F))$로 귀착한다는
조건에 의하여 유일하게 정해진다 (T, 3.2.2). 또한 $A$의 각 원소는
$\Gamma(X,\mathcal F)$의 자기준동형 $\mu$와
$\Gamma(X',\psi^*(\mathcal F))$의 자기준동형 $\mu'$을 정하며, 이때 도식
\[
\label{0.12.1.3.3-fr}
\tag{12.1.3.3}
\xymatrix{
\Gamma(X,\mathcal F)\ar[r]\ar[d]_\mu
&\Gamma(X',\psi^*(\mathcal F))\ar[d]^{\mu'}\\
\Gamma(X,\mathcal F)\ar[r]
&\Gamma(X',\psi^*(\mathcal F))
}
\]
은 가환한다. 보편 코호몰로지 함자에서 사상 연장의 유일성
성질 (T, 2.2)에 의해, $\mu$와 $\mu'$에는
(\hyperref[0.12.1.3.3-fr]{12.1.3.3})과 같은 도식들을 가환하게
하는 코호몰로지로의 유일한 연장이 있다. 이는
(\hyperref[0.12.1.3.2-fr]{12.1.3.2})가 $A$-가군 준동형임을 뜻한다.
이제
\[
f^*(\mathcal F)=\psi^*(\mathcal F)
\otimes_{\psi^*(\mathcal O_X)}\mathcal O_{X'}
\]
이고, 따라서 $\psi^*(\mathcal O_X)$-가군 $\psi^*(\mathcal F)$에서
$\mathcal O_{X'}$-가군 $f^*(\mathcal F)$로 가는 표준 디-준동형
$\psi^*(\mathcal F)\to f^*(\mathcal F)$가 있음을 유의하자. 함자성에
의해 이에 대응하는 환들이 $A$와 $A'$인 함자적 디-준동형
\[
\label{0.12.1.3.4-fr}
\tag{12.1.3.4}
H^p(X',\psi^*(\mathcal F))\longrightarrow H^p(X',f^*(\mathcal F))
\]
을 얻는다. 이 디-준동형과
(\hyperref[0.12.1.3.2-fr]{12.1.3.2})를 합성하면 $\mathcal F$에 대해
함자적인 표준 디-준동형
\[
\label{0.12.1.3.5-fr}
\tag{12.1.3.5}
\theta_p:H^p(X,\mathcal F)\longrightarrow H^p(X',f^*(\mathcal F))
\]
을 얻는다. 끝으로 함자성에 의해 준동형
$u^\sharp:f^*(\mathcal F)\to\mathcal F'$에서 $A'$-가군 준동형
$H^p(X',f^*(\mathcal F))\to H^p(X',\mathcal F')$을 얻고, 이를
(\hyperref[0.12.1.3.5-fr]{12.1.3.5})와 합성하면
(\hyperref[0.12.1.3.1-fr]{12.1.3.1})을 얻는다.

$f'=(\psi',\theta'):X''\to X'$을 환 달린 공간의 두 번째 사상이라 하고,
$f''=ff'$을 합성사상이라 하자. 함자 $\psi'^*$와 텐서곱의
교환가능성 (\hyperref[0.4.3.3-ko]{0\textsubscript{I}, 4.3.3})을 고려하면,
디-준동형
$H^p(X',f^*(\mathcal F))\to H^p(X'',f'^*(f^*(\mathcal F)))$와
(\hyperref[0.12.1.3.5-fr]{12.1.3.5})의 합성은 이에 대응하는 디-준동형
$H^p(X,\mathcal F)\to H^p(X'',f''^*(\mathcal F))$임을 곧바로 확인한다.
\end{env}

\begin{env}[12.1.4]
\label{0.12.1.4-fr}
준동형 (\hyperref[0.12.1.3.2-fr]{12.1.3.2})의 직접적인 정의는 다음과
같이 얻을 수 있다. $X$ 위의 아벨군의 층들로 이루어진 $\mathcal F$의
단사 분해 $\mathcal L^\bullet=(\mathcal L^i)$를 생각하자. 함자
$\psi^*$가 완전하므로, $\psi^*(\mathcal L^\bullet)$은 $X'$ 위의 층들로
이루어진 $\psi^*(\mathcal F)$의 분해이다. $\mathcal L'^\bullet=
(\mathcal L'^i)$가 $X'$ 위의 아벨군의 층들의 범주에서
$\psi^*(\mathcal F)$의 단사 분해이면, 증대와 양립하는 아벨군의
층들의 복합체 사상
$\psi^*(\mathcal L^\bullet)\to\mathcal L'^\bullet$이 있으며
(M, V, 1.1.a)), 이는 호모토피를 제외하면 잘 정해진다. 이로부터
아벨군의 복합체 준동형
\[
\Gamma(X,\mathcal L^\bullet)\longrightarrow
\Gamma(X',\psi^*(\mathcal L^\bullet))\longrightarrow
\Gamma(X',\mathcal L'^\bullet)
\]
을 얻는다. 그 합성은 코호몰로지를 취함으로써 $\delta$-함자의 사상
$H^\bullet(X,\mathcal F)\to H^\bullet(X',\psi^*(\mathcal F))$을 준다.
이는 차수 0에서 (\hyperref[0.12.1.3.2-fr]{12.1.3.2})와 일치하므로,
그 사상과 같다 (T, 2.2).

이제 $X$의 \emph{열린 덮개} $\mathfrak U=(U_\alpha)$를 생각하고,
$\mathfrak U$의 역상인 $X'$의 열린 덮개를
$\mathfrak U'=(f^{-1}(U_\alpha))$라 하자. $X$의 모든 열린집합 $V$에
대한 표준 준동형
$\Gamma(V,\mathcal F)\to\Gamma(f^{-1}(V),f^*(\mathcal F))$들은
곧바로 (G, II, 5.1 참조) 복합체 준동형
$\check C^\bullet(\mathfrak U,\mathcal F)\to
\check C^\bullet(\mathfrak U',f^*(\mathcal F))$을 정의하며, 따라서
표준 준동형
\[
\label{0.12.1.4.1-fr}
\tag{12.1.4.1}
\theta_p:\check H^p(\mathfrak U,\mathcal F)
\longrightarrow\check H^p(\mathfrak U',f^*(\mathcal F))
\]
을 얻는다.

\oldpage[0\textsubscript{III}]{52}
또한 다음 도식들은 가환한다.
\[
\label{0.12.1.4.2-fr}
\tag{12.1.4.2}
\xymatrix{
\check H^p(\mathfrak U,\mathcal F)\ar[r]^{\theta_p}\ar[d]
&\check H^p(\mathfrak U',f^*(\mathcal F))\ar[d]\\
H^p(X,\mathcal F)\ar[r]_{\theta_p}
&H^p(X',f^*(\mathcal F))
}
\]
여기서 수직 화살표들은 (G, II, 5.2)의 표준 준동형들이다.
(\hyperref[0.12.1.4.2-fr]{12.1.4.2})의 가환성을 보이기 위하여,
$\Gamma(X,\mathcal C^\bullet(\mathfrak U,\mathcal F))=
\check C^\bullet(\mathfrak U,\mathcal F)$인 (G, II, 5.2)
$\mathfrak U$에 대한 $\mathcal F$의 (교대) 공사슬층의 복합체
$\mathcal C^\bullet(\mathfrak U,\mathcal F)$를 생각하자. 표준 준동형
$\Gamma(V,\mathcal F)\to
\Gamma(f^{-1}(V),\psi^*(\mathcal F))$들은 이때 $\psi$-사상
$\mathcal C^\bullet(\mathfrak U,\mathcal F)\to
\mathcal C^\bullet(\mathfrak U',\psi^*(\mathcal F))$을 정의하며,
위의 표기법으로 다음 가환도식을 얻는다.
\[
\xymatrix{
\Gamma(X,\mathcal C^\bullet(\mathfrak U,\mathcal F))\ar[r]\ar[d]
&\Gamma(X',\mathcal C^\bullet(\mathfrak U',\psi^*(\mathcal F)))\ar[d]\\
\Gamma(X,\mathcal L^\bullet)\ar[r]
&\Gamma(X',\psi^*(\mathcal L^\bullet))\ar[r]
&\Gamma(X',\mathcal L'^\bullet)
}
\]
코호몰로지를 취하면 이는 가환도식
\[
\xymatrix{
\check H^p(\mathfrak U,\mathcal F)\ar[r]\ar[d]
&\check H^p(\mathfrak U',\psi^*(\mathcal F))\ar[d]\\
H^p(X,\mathcal F)\ar[r]
&H^p(X',\psi^*(\mathcal F))
}
\]
을 준다. 여기서 수직 화살표들은 (G, II, 5.2)의 표준 준동형들이다.
이 도식들과 다음 가환도식들을 결합하면 충분하다.
\[
\xymatrix{
\check H^p(\mathfrak U',\psi^*(\mathcal F))\ar[r]\ar[d]
&\check H^p(\mathfrak U',f^*(\mathcal F))\ar[d]\\
H^p(X',\psi^*(\mathcal F))\ar[r]
&H^p(X',f^*(\mathcal F))
}
\]
\oldpage[0\textsubscript{III}]{53}
이 도식들은 준동형 $\psi^*(\mathcal F)\to f^*(\mathcal F)$와
(G, II, 5.2)의 표준 준동형들의 함자성에서 나오며, 이들을 결합하면
(\hyperref[0.12.1.4.2-fr]{12.1.4.2})의 가환성을 얻는다.

$A=\Gamma(X,\mathcal O_X)$, $A'=\Gamma(X',\mathcal O_{X'})$일 때
준동형 (12.1.4.1)은 환 $A$와 $A'$에 대응하는 가군의 디-준동형임을
유의하자. 두 사상의 합성에 대하여
(\hyperref[0.12.1.4.1-fr]{12.1.4.1})은
(\hyperref[0.12.1.3.5-fr]{12.1.3.5})의 추이성과 비슷한
\emph{추이성}을 갖는다. 끝으로 앞의 정의들에서 $\mathcal F$의 단사
분해 $\mathcal L^\bullet$ 대신, 모든 $i$와 모든 $p>0$에 대하여
$H^p(X,\mathcal L^i)=0$인 분해로 시작해도 똑같았을 것임을 유의하자
(G, II, 4.7.1).
\end{env}

\begin{env}[12.1.5]
\label{0.12.1.5-fr}
$\mathcal F$, $\mathcal G$, $\mathcal H$를 세 $\mathcal O_X$-가군이라 하고,
코호몰로지 위에 컵곱 (\hyperref[0.12.1.2.1-fr]{12.1.2.1})에서 유도되는
준동형
\[
\label{0.12.1.5.1-fr}
\tag{12.1.5.1}
H^p(X,\mathcal F)\otimes_A H^q(X,\mathcal G)
\longrightarrow H^{p+q}(X,\mathcal H)
\]
을 주는 $\mathcal O_X$-준동형
$u:\mathcal F\otimes_{\mathcal O_X}\mathcal G\to\mathcal H$를 생각하자.
(\hyperref[0.12.1.3-fr]{12.1.3})의 가정과 표기법 아래 다음 도식들이
가환함을 보이자.
\[
\label{0.12.1.5.2-fr}
\tag{12.1.5.2}
\xymatrix@C=10pt{
H^p(X,\mathcal F)\otimes_A H^q(X,\mathcal G)\ar[r]\ar[d]
&H^{p+q}(X,\mathcal H)\ar[d]\\
H^p(X',f^*(\mathcal F))\otimes_{A'}H^q(X',f^*(\mathcal G))\ar[r]
&H^{p+q}(X',f^*(\mathcal H))
}
\]
여기서 수직 화살표들은 표준 준동형
(\hyperref[0.12.1.3.5-fr]{12.1.3.5})에서 나온다. 이를 위해
(\hyperref[0.12.1.5.1-fr]{12.1.5.1})을 얻는 방법을 상기하자.
$\mathcal F$, $\mathcal G$, $\mathcal H$ 각각의 \emph{표준} 분해
(G, II, 4.3) $\mathcal L^\bullet$, $\mathcal M^\bullet$,
$\mathcal N^\bullet$에서 시작한다. 이들은 $\mathcal O_X$-가군들로
이루어져 있다. $u$에 대응하는 $\mathcal O_X$-가군의 복합체의 선형 사상
$\mathcal L^\bullet\otimes_{\mathcal O_X}\mathcal M^\bullet
\to\mathcal N^\bullet$은 $A$-가군의 복합체 준동형
$\Gamma(X,\mathcal L^\bullet)\otimes_A\Gamma(X,\mathcal M^\bullet)
\to\Gamma(X,\mathcal N^\bullet)$을 주며, 코호몰로지를 취하면 준동형
\[
H^p(\Gamma(X,\mathcal L^\bullet))\otimes_A
H^q(\Gamma(X,\mathcal M^\bullet))
\longrightarrow H^{p+q}(\Gamma(X,\mathcal N^\bullet))
\]
을 얻는다 (G, II, 6.6). 그런데 다음 도식은 분명 가환한다.
\[
\label{0.12.1.5.3-fr}
\tag{12.1.5.3}
\xymatrix@C=8pt{
\Gamma(X,\mathcal L^\bullet)\otimes_A\Gamma(X,\mathcal M^\bullet)
\ar[r]\ar[d]
&\Gamma(X,\mathcal N^\bullet)\ar[d]\\
\Gamma(X',\psi^*(\mathcal L^\bullet))
\otimes_{\Gamma(X',\psi^*(\mathcal O_X))}
\Gamma(X',\psi^*(\mathcal M^\bullet))\ar[r]
&\Gamma(X',\psi^*(\mathcal N^\bullet))
}
\]
\oldpage[0\textsubscript{III}]{54}
코호몰로지를 취하면 이 도식은 다음 가환도식들을 준다.
\[
\label{0.12.1.5.4-fr}
\tag{12.1.5.4}
\xymatrix@C=7pt{
H^p(X,\mathcal F)\otimes_A H^q(X,\mathcal G)\ar[r]\ar[d]
&H^{p+q}(X,\mathcal H)\ar[d]\\
H^p(\Gamma(X',\psi^*(\mathcal L^\bullet)))
\otimes_{\Gamma(X',\psi^*(\mathcal O_X))}
H^q(\Gamma(X',\psi^*(\mathcal M^\bullet)))\ar[r]
&H^{p+q}(\Gamma(X',\psi^*(\mathcal N^\bullet)))
}
\]
한편 $\psi^*(\mathcal L^\bullet)$, $\psi^*(\mathcal M^\bullet)$,
$\psi^*(\mathcal N^\bullet)$은 각각 $\psi^*(\mathcal F)$,
$\psi^*(\mathcal G)$, $\psi^*(\mathcal H)$의 \emph{분해}이므로,
다음 가환도식이 있다 (G, II, 6.6.1).
\[
\label{0.12.1.5.5-fr}
\tag{12.1.5.5}
\xymatrix@C=7pt{
H^p(\Gamma(X',\psi^*(\mathcal L^\bullet)))
\otimes_{\Gamma(X',\psi^*(\mathcal O_X))}
H^q(\Gamma(X',\psi^*(\mathcal M^\bullet)))\ar[r]\ar[d]
&H^{p+q}(\Gamma(X',\psi^*(\mathcal N^\bullet)))\ar[d]\\
H^p(X',\psi^*(\mathcal F))
\otimes_{\Gamma(X',\psi^*(\mathcal O_X))}
H^q(X',\psi^*(\mathcal G))\ar[r]
&H^{p+q}(X',\psi^*(\mathcal H))
}
\]
끝으로 함자성에 의해 다음 가환도식을 얻는다.
\[
\label{0.12.1.5.6-fr}
\tag{12.1.5.6}
\xymatrix@C=7pt{
H^p(X',\psi^*(\mathcal F))
\otimes_{\Gamma(X',\psi^*(\mathcal O_X))}
H^q(X',\psi^*(\mathcal G))\ar[r]\ar[d]
&H^{p+q}(X',\psi^*(\mathcal H))\ar[d]\\
H^p(X',f^*(\mathcal F))\otimes_{A'}H^q(X',f^*(\mathcal G))\ar[r]
&H^{p+q}(X',f^*(\mathcal H))
}
\]
세 도식 (\hyperref[0.12.1.5.4-fr]{12.1.5.4}),
(\hyperref[0.12.1.5.5-fr]{12.1.5.5}),
(\hyperref[0.12.1.5.6-fr]{12.1.5.6})을 결합하면 구하는 가환도식
(\hyperref[0.12.1.5.2-fr]{12.1.5.2})를 얻는다.
\end{env}

\begin{remark}[12.1.6]
\label{0.12.1.6-fr}
\oldpage[0\textsubscript{III}]{55}
(\hyperref[0.12.1.3-fr]{12.1.3})의 표기법 아래 다음 가환도식이 있다고
가정하자.
\[
\label{0.12.1.6.1-fr}
\tag{12.1.6.1}
\xymatrix{
0\ar[r]&\mathcal F\ar[r]^r\ar[d]_u
&\mathcal G\ar[r]^s\ar[d]_v
&\mathcal H\ar[r]\ar[d]_w&0\\
0\ar[r]&\mathcal F'\ar[r]_{r'}
&\mathcal G'\ar[r]_{s'}
&\mathcal H'\ar[r]&0
}
\]
여기서 $r$, $s$는 $\mathcal O_X$-가군 준동형, $r'$, $s'$은
$\mathcal O_{X'}$-가군 준동형, $u$, $v$, $w$는 $f$-사상이고 두 행은
\emph{완전}하다. 그러면 가환도식
\[
\label{0.12.1.6.2-fr}
\tag{12.1.6.2}
\xymatrix@C=6pt{
\cdots\ar[r]&H^p(X,\mathcal F)\ar[r]\ar[d]_{u_p}
&H^p(X,\mathcal G)\ar[r]\ar[d]_{v_p}
&H^p(X,\mathcal H)\ar[r]^\partial\ar[d]_{w_p}
&H^{p+1}(X,\mathcal F)\ar[r]\ar[d]_{u_{p+1}}&\cdots\\
\cdots\ar[r]&H^p(X',\mathcal F')\ar[r]
&H^p(X',\mathcal G')\ar[r]
&H^p(X',\mathcal H')\ar[r]_\partial
&H^{p+1}(X',\mathcal F')\ar[r]&\cdots
}
\]
을 얻는다. 실제로 (\hyperref[0.12.1.6.1-fr]{12.1.6.1})은 다음과 같이
인수분해된다.
\[
\xymatrix@R=14pt{
0\ar[r]&\mathcal F\ar[r]\ar[d]
&\mathcal G\ar[r]\ar[d]
&\mathcal H\ar[r]\ar[d]&0\\
0\ar[r]&\psi^*(\mathcal F)\ar[r]\ar[d]
&\psi^*(\mathcal G)\ar[r]\ar[d]
&\psi^*(\mathcal H)\ar[r]\ar[d]&0\\
0\ar[r]&\mathcal F'\ar[r]
&\mathcal G'\ar[r]
&\mathcal H'\ar[r]&0
}
\]
여기서 가운데 행은 \emph{완전}하다
(\hyperref[0.3.7.2-ko]{0\textsubscript{I}, 3.7.2}). 이제
(\hyperref[0.12.1.3.2-fr]{12.1.3.2})가 $\delta$-함자의 준동형이라는
사실과 $H^p(X',\mathcal F')$들이 $\mathcal F'$에 대한
$\delta$-함자를 이룬다는 사실을 쓰면 충분하다.
\end{remark}

\begin{env}[12.1.7]
\label{0.12.1.7-fr}
가정과 표기법을 (\hyperref[0.12.1.3-fr]{12.1.3})과 같이 하고,
이제 $\mathcal F=f_*(\mathcal F')=\psi_*(\mathcal F')$인 경우를 생각하자.
(\hyperref[0.12.1.3-fr]{12.1.3})에서 정의한 디-준동형
\[
\label{0.12.1.7.1-fr}
\tag{12.1.7.1}
H^p(X,f_*(\mathcal F'))\longrightarrow H^p(X',\mathcal F')
\]
은 $H^p(X',\mathcal F')$의 어떤 자기동형의 차이를 제외하면 합성함자
$\mathcal F'\leadsto\Gamma(X,\psi_*(\mathcal F'))$의 스펙트럼 열의
\emph{가장자리 준동형}으로 얻을 수 있음을 보이겠다 (T, 2.4).
(\hyperref[0.12.1.4-fr]{12.1.4})에서 준동형
(\hyperref[0.12.1.7.1-fr]{12.1.7.1})에 대해 준 설명은 여기서 이 준동형을
다음과 같이 얻을 수 있음을 보인다. $\psi_*(\mathcal F')$와
$\mathcal F'$의 단사 분해 $\mathcal L^\bullet$과
$\mathcal L'^\bullet$을 각각 택한 뒤, 표준 준동형
$\psi^*(\psi_*(\mathcal F'))\to\mathcal F'$의 ``위에 놓인'' 복합체 준동형
$v:\psi^*(\mathcal L^\bullet)\to\mathcal L'^\bullet$을 택한다.
\oldpage[0\textsubscript{III}]{56}
이어서 $\Gamma(X',\mathcal L'^\bullet)=
\Gamma(X,\psi_*(\mathcal L'^\bullet))$임을 유의하면, 합성 준동형
\[
\Gamma(X,\mathcal L^\bullet)\longrightarrow
\Gamma(X',\psi^*(\mathcal L^\bullet))
\xrightarrow{\Gamma(v)}\Gamma(X',\mathcal L'^\bullet)
\]
은 다름 아닌
\[
\label{0.12.1.7.2-fr}
\tag{12.1.7.2}
\Gamma(v^\flat):\Gamma(X,\mathcal L^\bullet)
\longrightarrow\Gamma(X,\psi_*(\mathcal L'^\bullet))
\]
이다 (\hyperref[0.3.7.1-ko]{0\textsubscript{I}, 3.7.1}). 그리고
(\hyperref[0.12.1.7.2-fr]{12.1.7.2})에서 코호몰로지를 취하면
(\hyperref[0.12.1.7.1-fr]{12.1.7.1})을 얻는다. 한편 합성함자
$\mathcal F'\leadsto\Gamma(X,\psi_*(\mathcal F'))$의 스펙트럼 열들은
$X$ 위의 아벨군의 층들의 범주에서 복합체
$\psi_*(\mathcal L'^\bullet)$의 단사 카르탕--아일렌베르크 분해
$\mathcal M^{\bullet\bullet}=(\mathcal M^{ij})$를 생각하여 얻는다.
문제의 스펙트럼 열들은 이중복합체
$\Gamma(X,\mathcal M^{\bullet\bullet})$의 스펙트럼 열들이다
($i<0$ 또는 $j<0$이면 $\mathcal M^{ij}=0$이므로 이들은 쌍정칙이다).
그런데 이 이중복합체의 첫 번째 스펙트럼 열은 \emph{퇴화}한다.
실제로 층 $\psi_*(\mathcal L'^i)$들이 플라스크이므로 (G, II, 3.1.1),
$q>0$이면
$H_{\mathrm{II}}^q(\Gamma(X,\mathcal M^{i,\bullet}))
=H^q(\psi_*(\mathcal L'^i))=0$이다 (G, II, 4.4.3). 따라서 전단사인
가장자리 준동형 (\hyperref[0.11.1.6-fr]{11.1.6})
\[
\label{0.12.1.7.3-fr}
\tag{12.1.7.3}
'E_2^{i0}=H^i(H_{\mathrm{II}}^0(\Gamma(X,\mathcal M^{\bullet\bullet})))
\longrightarrow H^i(\Gamma(X,\mathcal M^{\bullet\bullet}))
\]
이 있다. 이 준동형은 증대
\[
\label{0.12.1.7.4-fr}
\tag{12.1.7.4}
\Gamma(X,\psi_*(\mathcal L'^\bullet))
\longrightarrow\Gamma(X,\mathcal M^{\bullet\bullet})
\]
에서 코호몰로지를 취하여 나온다는 것을 안다
(\hyperref[0.11.3.4-fr]{11.3.4}). 이 증대 자체는 증대
$\eta:\psi_*(\mathcal L'^\bullet)\to\mathcal M^{\bullet\bullet}$에서 나온다.
한편 두 번째 스펙트럼 열에는 가장자리 준동형
\[
\label{0.12.1.7.5-fr}
\tag{12.1.7.5}
''E_2^{i0}=H^i(H_{\mathrm I}^0(\Gamma(X,\mathcal M^{\bullet\bullet})))
\longrightarrow H^i(\Gamma(X,\mathcal M^{\bullet\bullet}))
\]
이 있다. 이는 (\hyperref[0.11.3.4-fr]{11.3.4}) 복합체 준동형
$Z_{\mathrm I}^0(\Gamma(X,\mathcal M^{\bullet\bullet}))
\to\Gamma(X,\mathcal M^{\bullet\bullet})$에서 코호몰로지를 취하여 나온다.
그런데 $\psi_*$가 왼쪽 완전하므로 열
\[
0\longrightarrow\psi_*(\mathcal F')\longrightarrow
\psi_*(\mathcal L'^0)\longrightarrow\psi_*(\mathcal L'^1)
\]
은 완전하다. 따라서 카르탕--아일렌베르크 분해의 정의
(\hyperref[0.11.4.2-fr]{11.4.2})에 의해
$B_{\mathrm I}^0(\mathcal M^{\bullet\bullet})=0$,
$Z_{\mathrm I}^0(\mathcal M^{\bullet\bullet})=\mathcal L^\bullet$로
택할 수 있다. 도식
\[
\xymatrix@C=50pt@R=24pt{
\mathcal L^0\ar[r]^{i^0}&\mathcal M^{00}\\
\psi_*(\mathcal F')\ar[u]^\varepsilon
\ar[r]_{\varepsilon'}\ar[ur]^{\varepsilon''}
&\psi_*(\mathcal L'^0)\ar[u]_{\eta^0}
}
\]
은 가환하므로, 복합체 단사사상
$i:\mathcal L^\bullet\to\mathcal M^{\bullet\bullet}$은 증대
$\varepsilon$ 및 $\varepsilon''$과 양립한다. 이로써
$\mathcal L^\bullet$에서 $\mathcal M^{\bullet\bullet}$로 가는 두 복합체
준동형
\[
\xymatrix@C=52pt@R=22pt{
\mathcal L^\bullet\ar[r]^i\ar[dr]_{v^\flat}
&\mathcal M^{\bullet\bullet}\\
&\psi_*(\mathcal L'^\bullet)\ar[u]_\eta
}
\]
\oldpage[0\textsubscript{III}]{57}
을 얻으며, 이들은 증대 $\varepsilon$ 및 $\varepsilon''$과 양립한다.
$\mathcal L^\bullet$은 \emph{단사 분해}이고
$\mathcal M^{\bullet\bullet}$은 단사층들로 이루어져 있으므로,
(M, V, 1.1a))에 의해 이 두 준동형은 \emph{서로 호모토픽}이다.
따라서 이에 대응하는 두 준동형
$\Gamma(X,\mathcal L^\bullet)\to
\Gamma(X,\mathcal M^{\bullet\bullet})$도 서로 호모토픽이고,
코호몰로지를 취하면 \emph{같은} 준동형을 얻는다. 다시 말해
(\hyperref[0.12.1.7.5-fr]{12.1.7.5})의 가장자리 준동형, 즉
$H^p(X,\psi_*(\mathcal F'))\to
H^p(\Gamma(X,\mathcal M^{\bullet\bullet}))$은
(\hyperref[0.12.1.7.1-fr]{12.1.7.1})과
(\hyperref[0.12.1.7.3-fr]{12.1.7.3})의 합성임을 보였다. 여기서 후자는
$H^p(X',\mathcal F')\to
H^p(\Gamma(X,\mathcal M^{\bullet\bullet}))$로 쓰이며 앞에서
\emph{동형}임을 보았다. 이로부터 주장이 따른다.
\end{env}

\subsection{고차 직상}
\label{subsection:0.12.2-fr}

\begin{env}[12.2.1]
\label{0.12.2.1-fr}
$(X,\mathcal O_X)$, $(Y,\mathcal O_Y)$를 두 환 달린 공간이라 하고,
$f=(\psi,\theta)$를 $X$에서 $Y$로 가는 사상이라 하자. 이 사상은
직상 함자
\[
f_*:C(X)\longrightarrow C(Y)
\]
를 정의한다. 이는 사실 $X$ 위의 아벨군층들의 범주에서 정의된 함자
$\psi_*$를 $C(X)$에 제한한 것과 같다. 뒤의 함자는 가법이고 왼쪽
완전하다. $X$ 위의 모든 아벨군층은 어떤 단사 아벨군층의 부분층과
동형이므로, 함자 $\psi_*$의 오른쪽 유도함자
$\mathcal F\leadsto R^p\psi_*(\mathcal F)$가 정의된다. 각
$R^p\psi_*(\mathcal F)$는 $Y$ 위의 아벨군층이고, $R^p\psi_*$들은
보편 코호몰로지 함자를 이룬다 (T, 2.3).

또한 층 $R^p\psi_*(\mathcal F)$는 전층
$V\leadsto H^p(f^{-1}(V),\mathcal F)$에 연관된 층이다 (T, 3.7.2).
이제 $\mathcal F$가 $\mathcal O_X$-가군이라고 가정하면,
$H^p(f^{-1}(V),\mathcal F)$는 자연스럽게
$\Gamma(f^{-1}(V),\mathcal O_X)$-가군, 따라서
$\Gamma(V,\psi_*(\mathcal O_X))$-가군의 구조를 갖는다. 준동형
$\theta:\mathcal O_Y\to\psi_*(\mathcal O_X)$로부터 다시
$\Gamma(V,\mathcal O_Y)$-가군의 구조를 얻는다. 이렇게 정의한 구조에
대하여 열린집합 $V$에서 열린집합 $V'\subset V$로의 제한은
디-준동형을 정의함이 명백하다. 따라서 각 $R^p\psi_*(\mathcal F)$에
$\mathcal O_Y$-가군 구조를 정의할 수 있다. 이 $\mathcal O_Y$-가군을
$R^p f_*(\mathcal F)$로 나타내며, 이로써 $R^p f_*$는 $C(X)$에서
$C(Y)$로 가는 가법 함자로 정의된다. 더 나아가 $R^p f_*$들은
$\partial$-함자를 이룬다. 실제로
\[
0\longrightarrow\mathcal F'\longrightarrow\mathcal F
\longrightarrow\mathcal F''\longrightarrow0
\]
이 $\mathcal O_X$-가군들의 완전열이면, 위에서 준
$R^p\psi_*$와 $R^p\psi_*(\mathcal F)$ 위의 $\mathcal O_Y$-가군
구조의 설명으로부터 준동형
\[
\partial:R^p\psi_*(\mathcal F'')\longrightarrow
R^{p+1}\psi_*(\mathcal F')
\]
이 $\mathcal O_Y$-가군의 준동형임이 곧바로 따른다. 끝으로
$R^p f_*$는 $f_*$의 오른쪽 유도함자와 동일시된다. 실제로 모든
$\mathcal O_X$-가군은 $\mathcal O_X$-가군들로 이루어진 단사 분해를
갖고, 이러한 분해는 아벨군의 플라스크층들로 이루어지므로
(G, II, 7.1), $R^p\psi_*(\mathcal F)$를 계산하는 데 쓸 수 있다.
모든 플라스크층 $\mathcal G$와 모든 $n\geq1$에 대하여
$R^n\psi_*(\mathcal G)=0$이기 때문이다 (T, 2.4.1, 주석 3 및 명제
3.3.2의 따름정리). 따라서 $R^p f_*$들은 $C(X)$에서 $C(Y)$로 가는
보편 코호몰로지 함자를 이룬다 (T, 2.3).
\end{env}

\begin{env}[12.2.2]
\label{0.12.2.2-fr}
$\mathcal F$, $\mathcal G$를 두 $\mathcal O_X$-가군이라 하자.
(\hyperref[0.12.2.1-fr]{12.2.1})의 표기 아래 $Y$의 모든 열린집합
$V$에 대하여 컵곱 준동형 (\hyperref[0.12.1.2.1-fr]{12.1.2.1})
\[
H^p(f^{-1}(V),\mathcal F)
\otimes_{\Gamma(f^{-1}(V),\mathcal O_X)}
H^q(f^{-1}(V),\mathcal G)
\longrightarrow
H^{p+q}(f^{-1}(V),\mathcal F\otimes_{\mathcal O_X}\mathcal G)
\]
이 있다.
\oldpage[0\textsubscript{III}]{58}
컵곱의 정의 (G, II, 6.6)에서 이 준동형들이 $V$에서 그 열린
부분공간 $V'$로의 이행과 가환함이 곧바로 따른다. 한편 $\theta$에서
오는 환 준동형
\[
\Gamma(V,\mathcal O_Y)\longrightarrow\Gamma(V,\psi_*(\mathcal O_X))
=\Gamma(f^{-1}(V),\mathcal O_X)
\]
이 있으므로, 텐서곱의 표준 준동형
\[
\begin{aligned}
&H^p(f^{-1}(V),\mathcal F)\otimes_{\Gamma(V,\mathcal O_Y)}
H^q(f^{-1}(V),\mathcal G)\\
&\qquad\longrightarrow
H^p(f^{-1}(V),\mathcal F)
\otimes_{\Gamma(f^{-1}(V),\mathcal O_X)}
H^q(f^{-1}(V),\mathcal G)
\end{aligned}
\]
을 얻으며, 이것도 $V$에서 $V'$로의 제한과 양립한다. 따라서 합성으로
$\Gamma(V,\mathcal O_Y)$-가군의 준동형을 얻고, 고려하는 전층들에
연관된 층에 대하여 $\mathcal F$, $\mathcal G$에 함자적인 표준 준동형
\[
\label{0.12.2.2.1-fr}
\tag{12.2.2.1}
R^p f_*(\mathcal F)\otimes_{\mathcal O_Y}R^q f_*(\mathcal G)
\longrightarrow
R^{p+q}f_*(\mathcal F\otimes_{\mathcal O_X}\mathcal G)
\]
을 정의한다. $p=q=0$일 때 이 준동형은
(\hyperref[0.4.2.2.1-ko]{0\textsubscript{I}, 4.2.2.1})로 환원된다.
\end{env}

\begin{proposition}[12.2.3]
\label{0.12.2.3-fr}
\emph{모든 $\mathcal O_X$-가군 $\mathcal F$와 유한 계수 국소 자유
$\mathcal O_Y$-가군 $\mathcal L$에 대하여 함자적인 표준 동형
\[
\label{0.12.2.3.1-fr}
\tag{12.2.3.1}
R^p f_*(\mathcal F)\otimes_{\mathcal O_Y}\mathcal L
\xrightarrow{\sim}
R^p f_*(\mathcal F\otimes_{\mathcal O_X}f^*(\mathcal L))
\]
이 성립한다.}
\end{proposition}

\begin{proof}
준동형 (\hyperref[0.12.2.3.1-fr]{12.2.3.1})은
(\hyperref[0.12.2.2.1-fr]{12.2.2.1})의 특수한 경우인 준동형
\[
\label{0.12.2.3.2-fr}
\tag{12.2.3.2}
R^p f_*(\mathcal F)\otimes_{\mathcal O_Y}f_*(f^*(\mathcal L))
\longrightarrow
R^p f_*(\mathcal F\otimes_{\mathcal O_X}f^*(\mathcal L))
\]
과, 표준 준동형
(\hyperref[0.4.4.3.2-ko]{0\textsubscript{I}, 4.4.3.2})에서 오는
(\hyperref[0.12.2.3.1-fr]{12.2.3.1})의 좌변에서
(\hyperref[0.12.2.3.2-fr]{12.2.3.2})의 좌변으로 가는 준동형을
합성하여 얻는다. $\mathcal L$이 국소 자유일 때
(\hyperref[0.12.2.3.1-fr]{12.2.3.1})이 동형임을 확인하려면 곧바로
$\mathcal L=\mathcal O_Y$인 경우로 환원할 수 있다. 문제는 $Y$ 위에서
국소적이고, 고려하는 함자들은 $\mathcal L$에 관하여 가법이기 때문이다.
그런데 이 경우에는 (\hyperref[0.12.2.2.1-fr]{12.2.2.1})의 정의에 따라,
연관된 전층의 준동형이 전단사임을 확인하면 충분하고, 이는 관계
$f^*(\mathcal O_Y)=\mathcal O_X$에서 즉시 따른다.
\end{proof}

\begin{env}[12.2.4]
\label{0.12.2.4-fr}
$(Z,\mathcal O_Z)$를 셋째 환 달린 공간, $g:Y\to Z$를 환 달린 공간의
사상이라 하자. 모든 단사 $\mathcal O_X$-가군 $\mathcal G$에 대하여
$f_*(\mathcal G)$가 아벨군의 플라스크층임을 알고 있다
(G, II, 7.1 및 3.1.1). 따라서
(\hyperref[0.12.2.1-fr]{12.2.1})에 의해 모든 $p>0$에 대하여
$R^p g_*(f_*(\mathcal G))=0$이다. 이에 따라 (T, 2.4.1) 합성함자의
르레 스펙트럼 열을 합성함자 $g_*f_*$에 적용할 수 있다. 즉
$h=g\circ f$일 때 수렴대상이 함자 $R^\bullet h_*$이고 $E_2$항이
\[
\label{0.12.2.4.1-fr}
\tag{12.2.4.1}
E_2^{pq}=R^p g_*(R^q f_*(\mathcal F))
\]
로 주어지는 쌍정칙 스펙트럼 열이 있다.
\end{env}

\begin{env}[12.2.5]
\label{0.12.2.5-fr}
(\hyperref[0.12.2.4-fr]{12.2.4})의 조건 아래 $\mathcal O_Z$-가군의
표준 준동형
\[
\label{0.12.2.5.1-fr}
\tag{12.2.5.1}
R^n g_*(f_*(\mathcal F))\longrightarrow R^n h_*(\mathcal F)
\]
\[
\label{0.12.2.5.2-fr}
\tag{12.2.5.2}
R^n h_*(\mathcal F)\longrightarrow g_*(R^n f_*(\mathcal F))
\]
을 직접 정의하자.
\oldpage[0\textsubscript{III}]{59}
이들은 르레 스펙트럼 열의 ``가장자리 준동형''과 동일시할 수 있다
(\hyperref[0.12.1.7-fr]{12.1.7} 참조). 고차 직상층들에 연관되는
전층들 위에서 작업하면 충분하다
(\hyperref[0.12.2.1-fr]{12.2.1}). 이를 위하여 $Z$의 임의의 열린집합
$W$와 $Y$에서의 그 역상 $g^{-1}(W)$를 생각하자. 대응하는 환들이
$\Gamma(g^{-1}(W),\mathcal O_Y)$와
$\Gamma(h^{-1}(W),\mathcal O_X)$일 때 표준 디-준동형
\[
\label{0.12.2.5.3-fr}
\tag{12.2.5.3}
H^n(g^{-1}(W),f_*(\mathcal F))\longrightarrow
H^n(f^{-1}(g^{-1}(W)),f^*(f_*(\mathcal F)))
\]
이 있다. 한편 표준 준동형
(\hyperref[0.4.4.3.3-ko]{0\textsubscript{I}, 4.4.3.3})은 함자성에 따라
표준 준동형
\[
\label{0.12.2.5.4-fr}
\tag{12.2.5.4}
H^n(h^{-1}(W),f^*(f_*(\mathcal F)))\longrightarrow
H^n(h^{-1}(W),\mathcal F)
\]
을 주며, 이는 $\Gamma(h^{-1}(W),\mathcal O_X)$-가군의 준동형이다.
환 준동형
$\Gamma(W,\mathcal O_Z)\to\Gamma(g^{-1}(W),\mathcal O_Y)$을 고려하면,
(\hyperref[0.12.2.5.4-fr]{12.2.5.4})와
(\hyperref[0.12.2.5.3-fr]{12.2.5.3})의 합성으로 전층의 준동형을 얻고,
이는 층의 준동형 (\hyperref[0.12.2.5.1-fr]{12.2.5.1})을 준다.

(\hyperref[0.12.2.5.2-fr]{12.2.5.2})의 정의는 더 간단하다. 정의에
따라 $R^n h_*(\mathcal F)$는 전층
$W\leadsto H^n(f^{-1}(g^{-1}(W)),\mathcal F)$에 연관되고,
$R^n f_*(\mathcal F)$는 전층
$V\leadsto H^n(f^{-1}(V),\mathcal F)$에 연관된다. 따라서 표준 준동형
\[
H^n(f^{-1}(g^{-1}(W)),\mathcal F)\longrightarrow
\Gamma(g^{-1}(W),R^n f_*(\mathcal F))
\]
이 있고, 이 준동형들이 전층의 준동형을 정의함은 즉시 알 수 있다.
이 전층의 준동형이 다시
(\hyperref[0.12.2.5.2-fr]{12.2.5.2})를 정의한다.
\end{env}

\begin{env}[12.2.6]
\label{0.12.2.6-fr}
(\hyperref[0.12.2.4-fr]{12.2.4})의 가정 아래 $\mathcal F$,
$\mathcal G$, $\mathcal K$를 세 $\mathcal O_X$-가군이라 하고,
$u:\mathcal F\otimes\mathcal G\to\mathcal K$를
$\mathcal O_X$-준동형이라 하자. 그러면 다음 두 도식은 가환한다.
\[
\xymatrix@C=10pt@R=24pt{
R^p g_*(f_*(\mathcal F))\otimes_{\mathcal O_Z}R^q g_*(f_*(\mathcal G))
\ar[r]\ar[d]&R^{p+q}g_*(f_*(\mathcal K))\ar[d]\\
R^p h_*(\mathcal F)\otimes_{\mathcal O_Z}R^q h_*(\mathcal G)
\ar[r]&R^{p+q}h_*(\mathcal K)
}
\tag{12.2.6.1}\label{0.12.2.6.1-fr}
\]
그리고
\[
\xymatrix@C=10pt@R=24pt{
R^p h_*(\mathcal F)\otimes_{\mathcal O_Z}R^q h_*(\mathcal G)
\ar[r]\ar[d]&R^{p+q}h_*(\mathcal K)\ar[d]\\
g_*(R^p f_*(\mathcal F))\otimes_{\mathcal O_Z}g_*(R^q f_*(\mathcal G))
\ar[r]&g_*(R^{p+q}f_*(\mathcal K))
}
\tag{12.2.6.2}\label{0.12.2.6.2-fr}
\]
\oldpage[0\textsubscript{III}]{60}
여기서 수평 화살표들은
(\hyperref[0.12.2.2.1-fr]{12.2.2.1})에서 오며, 마지막 수평 화살표에는
(\hyperref[0.4.2.2.1-ko]{0\textsubscript{I}, 4.2.2.1})도 합성하였다.
수직 화살표들은 각각 준동형
(\hyperref[0.12.2.5.1-fr]{12.2.5.1})과
(\hyperref[0.12.2.5.2-fr]{12.2.5.2})에서 온다.

실제로 대응하는 전층의 준동형들에 대하여 확인하면 충분하다. 이
준동형들에 관하여 (\hyperref[0.12.2.2-fr]{12.2.2})와
(\hyperref[0.12.2.5-fr]{12.2.5})에서 준 정의로 돌아가면,
(\hyperref[0.12.2.6.1-fr]{12.2.6.1})의 경우 곧바로 가환도식들
(\hyperref[0.12.1.5.2-fr]{12.1.5.2})로 환원된다.
(\hyperref[0.12.2.6.2-fr]{12.2.6.2})의 확인은 더욱 간단하다.
\end{env}

\subsection{층의 $\operatorname{Ext}$ 함자에 관한 보충}
\label{subsection:0.12.3-fr}

\begin{env}[12.3.1]
\label{0.12.3.1-fr}
환 달린 공간 $(X,\mathcal O_X)$를 생각하자. $\mathcal O_X$-가군의
범주에서 $\Gamma(X,\mathcal O_X)$-가군의 범주로 가는 쌍함자
$\operatorname{Ext}_{\mathcal O_X}^p(X;\mathcal F,\mathcal G)$와,
$\mathcal O_X$-가군의 범주에서 그 자신으로 가는 쌍함자
$\mathcal{E}xt_{\mathcal O_X}^p(\mathcal F,\mathcal G)$의 정의와 주요
성질들, 그리고 이들을 잇는 쌍정칙 스펙트럼 열
$E(\mathcal F,\mathcal G)$에 대해서는 다시 논하지 않는다
(T, 4.2 및 G, II, 7.3).
\end{env}

\begin{env}[12.3.2]
\label{0.12.3.2-fr}
(M, XIV, 1)에서와 같은 방식으로 $\mathcal O_X$-가군 $\mathcal F$의
$\mathcal O_X$-가군 $\mathcal G$에 의한 확장이라는 개념과 동치인
확장류들 사이의 합성법칙을 정의한다. 실제로 가군에 대하여 한 논증은
임의의 아벨 범주에 명백하게 적용된다. 단사 대상에의 매장이 존재한다는
사실만을 쓰는 (M, XIV, 1.1)의 둘째 증명은 $\mathcal O_X$-가군의
범주에서도 여전히 유효하다. 따라서
$\operatorname{Ext}_{\mathcal O_X}^1(X;\mathcal F,\mathcal G)$는
$\mathcal F$의 $\mathcal G$에 의한 확장류들의 아벨군과 표준적으로
동일시된다.
\end{env}

\begin{proposition}[12.3.3]
\label{0.12.3.3-fr}
\emph{$(X,\mathcal O_X)$를 환층 $\mathcal O_X$가 연접인 환 달린
공간이라 하자. 임의의 두 연접 $\mathcal O_X$-가군 $\mathcal F$,
$\mathcal G$와 모든 $p\geq0$에 대하여
$\mathcal{E}xt_{\mathcal O_X}^p(\mathcal F,\mathcal G)$는 연접
$\mathcal O_X$-가군이다.}
\end{proposition}

\begin{proof}
$\mathcal{E}xt_{\mathcal O_X}^p(\mathcal F,\mathcal G)$들은
$\mathcal F$에 관하여 반변 코호몰로지 함자를 이룸에 유의하자.
$\mathcal F$가 연접이므로, 모든 $p$와 모든 점 $x\in X$에 대하여
$x$의 열린 근방 $U$와 $(\mathcal O_X|U)$-가군들의 완전열
\[
0\longrightarrow\mathcal R\longrightarrow\mathcal L_{p-1}
\longrightarrow\cdots\longrightarrow\mathcal L_0
\longrightarrow\mathcal F|U\longrightarrow0
\]
이 존재하여, 각 $\mathcal L_i$ ($0\leq i\leq p-1$)는 어떤
$\mathcal O_X^{n_i}|U$와 동형이고 $\mathcal R$은 연접이다. 이는
$\mathcal O_X$가 연접이라는 가정 아래
(\hyperref[0.5.3.2-ko]{0\textsubscript{I}, 5.3.2})와
(\hyperref[0.5.3.4-ko]{0\textsubscript{I}, 5.3.4})를 $p$에 관하여
귀납적으로 적용하여 얻는다.

이제 $p\geq1$일 때 $\mathcal L|U$가 어떤 $\mathcal O_X^n|U$와
동형인 모든 $\mathcal O_X$-가군 $\mathcal L$에 대하여
\[
\mathcal{E}xt_{\mathcal O_X|U}^p(\mathcal L|U,\mathcal G|U)=0
\]
임을 유의하자 (T, 4.2.3). 따라서 (M, V, 7.2)의 논증을 반변
코호몰로지 함자
\[
\mathcal F\leadsto
\mathcal{H}om_{\mathcal O_X|U}(\mathcal F|U,\mathcal G|U)
\]
에 적용할 수 있고, 완전열
\[
\mathcal{H}om_{\mathcal O_X|U}(\mathcal L_{p-1},\mathcal G|U)
\longrightarrow
\mathcal{H}om_{\mathcal O_X|U}(\mathcal R,\mathcal G|U)
\longrightarrow
\mathcal{E}xt_{\mathcal O_X|U}^p(\mathcal F|U,\mathcal G|U)
\longrightarrow0
\]
을 얻는다. 이 열의 처음 두 항은 연접 $(\mathcal O_X|U)$-가군이므로
(\hyperref[0.5.3.5-ko]{0\textsubscript{I}, 5.3.5}), 셋째 항도 연접이다
(\hyperref[0.5.3.4-ko]{0\textsubscript{I}, 5.3.4}).
\end{proof}

\oldpage[0\textsubscript{III}]{61}
\begin{proposition}[12.3.4]
\label{0.12.3.4-fr}
\emph{$f:X\to Y$를 환 달린 공간의 평탄 사상이라 하고,
$\mathcal F$, $\mathcal G$를 두 $\mathcal O_Y$-가군이라 하자.}
\begin{enumerate}
\item[{\rm(i)}] \emph{코호몰로지 쌍함자들의 준동형
\[
\label{0.12.3.4.1-fr}
\tag{12.3.4.1}
f^*(\mathcal{E}xt_{\mathcal O_Y}^{\bullet}(\mathcal F,\mathcal G))
\longrightarrow
\mathcal{E}xt_{\mathcal O_X}^{\bullet}(f^*(\mathcal F),f^*(\mathcal G))
\]
이 존재하며, 차수 0에서는 표준 준동형
(\hyperref[0.4.4.6-ko]{0\textsubscript{I}, 4.4.6})으로 환원된다.}

\item[{\rm(ii)}] \emph{스펙트럼 열의 표준 사상
\[
\label{0.12.3.4.2-fr}
\tag{12.3.4.2}
E(\mathcal F,\mathcal G)\longrightarrow
E(f^*(\mathcal F),f^*(\mathcal G))
\]
이 존재하며, $E_2$항에서는
(\hyperref[0.12.3.4.1-fr]{12.3.4.1})과
(\hyperref[0.12.1.3.1-fr]{12.1.3.1})에서 얻는 준동형
\[
\label{0.12.3.4.3-fr}
\tag{12.3.4.3}
H^p(Y,\mathcal{E}xt_{\mathcal O_Y}^q(\mathcal F,\mathcal G))
\longrightarrow
H^p(X,\mathcal{E}xt_{\mathcal O_X}^q(f^*(\mathcal F),f^*(\mathcal G)))
\]
으로 환원된다.}
\end{enumerate}
\end{proposition}

\begin{proof}
\begin{enumerate}
\item[{\rm(i)}]
$f^*$는 $\mathcal O_Y$-가군의 범주에서 완전함자이므로
(\hyperref[0.6.7.2-ko]{0\textsubscript{I}, 6.7.2}), 함자들
\[
\mathcal G\leadsto
f^*(\mathcal{H}om_{\mathcal O_Y}(\mathcal F,\mathcal G))
\quad\hbox{및}\quad
\mathcal G\leadsto
\mathcal{H}om_{\mathcal O_X}(f^*(\mathcal F),f^*(\mathcal G))
\]
은 왼쪽 완전하다. 따라서
(\hyperref[0.4.4.6-ko]{0\textsubscript{I}, 4.4.6})에서 이들의 유도함자
사이를 잇는 준동형을 표준적으로 얻는다.

이를 계산하기 위하여 $\mathcal G$의 단사 분해
$\mathcal L^\bullet=(\mathcal L^i)$을 취하면, 층의 복합체들의
코호몰로지 사상
\[
\mathcal H^p(f^*(\mathcal{H}om_{\mathcal O_Y}
(\mathcal F,\mathcal L^\bullet)))
\longrightarrow
\mathcal H^p(\mathcal{H}om_{\mathcal O_X}
(f^*(\mathcal F),f^*(\mathcal L^\bullet)))
\]
을 얻는다. $f^*$가 완전하므로 정의에 따라
\[
\mathcal H^p(f^*(\mathcal{H}om_{\mathcal O_Y}
(\mathcal F,\mathcal L^\bullet)))
=f^*(\mathcal H^p(\mathcal{H}om_{\mathcal O_Y}
(\mathcal F,\mathcal L^\bullet)))
=f^*(\mathcal{E}xt_{\mathcal O_Y}^p(\mathcal F,\mathcal G)).
\]
한편 $f^*$의 완전성에 따라 $f^*(\mathcal L^\bullet)$은
$f^*(\mathcal G)$의 분해이다. $\mathcal L'^\bullet=(\mathcal L'^i)$을
$\mathcal O_X$-가군 범주에서 $f^*(\mathcal G)$의 단사 분해라 하면,
호모토피를 제외하고 정해지는 복합체의 준동형
$f^*(\mathcal L^\bullet)\to\mathcal L'^\bullet$이 존재하며, 따라서
코호몰로지 위에 잘 정의된 준동형을 준다. 이 준동형과 위에서 정의한
준동형을 합성하여 (\hyperref[0.12.3.4.1-fr]{12.3.4.1})을 얻는다.

\item[{\rm(ii)}]
앞의 표기 아래 $\mathcal O_X$-가군층 복합체의 준동형
\[
f^*(\mathcal{H}om_{\mathcal O_Y}(\mathcal F,\mathcal L^\bullet))
\longrightarrow
\mathcal{H}om_{\mathcal O_X}(f^*(\mathcal F),\mathcal L'^\bullet)
\]
이 있다. $\mathcal M^{\bullet\bullet}$을
$\mathcal O_Y$-가군 범주에서 복합체
$\mathcal{H}om_{\mathcal O_Y}(\mathcal F,\mathcal L^\bullet)$의 단사
카르탕--아일렌베르크 분해라 하자. $f^*$의 완전성에 따라
$f^*(\mathcal M^{\bullet\bullet})$은 복합체
$f^*(\mathcal{H}om_{\mathcal O_Y}(\mathcal F,\mathcal L^\bullet))$의
카르탕--아일렌베르크 분해이다. $\mathcal M'^{\bullet\bullet}$을
복합체
$\mathcal{H}om_{\mathcal O_X}(f^*(\mathcal F),\mathcal L'^\bullet)$의
단사 카르탕--아일렌베르크 분해라 하면,
(\hyperref[0.11.4.2-fr]{11.4.2})에 따라 위의 준동형과 양립하는,
호모토피를 제외하고 정해지는 준동형
\[
f^*(\mathcal M^{\bullet\bullet})\longrightarrow
\mathcal M'^{\bullet\bullet}
\]
이 있다. 다시 말해 이는 층의 이중복합체들의 $f$-사상
$\mathcal M^{\bullet\bullet}\to\mathcal M'^{\bullet\bullet}$이다.
여기서 호모토피를 제외하고 정해지는 가군 이중복합체의 디-준동형
$\Gamma(Y,\mathcal M^{\bullet\bullet})\to
\Gamma(X,\mathcal M'^{\bullet\bullet})$을 얻고, 이어서 잘 정의된
스펙트럼 열의 사상을 얻는다
(\hyperref[0.11.3.2-fr]{11.3.2}). 이것이 구하는 사상
(\hyperref[0.12.3.4.2-fr]{12.3.4.2})이며,
(\hyperref[0.12.3.4.3-fr]{12.3.4.3})의 특징은 정의에서 곧바로 따른다.
\end{enumerate}
\end{proof}

\begin{proposition}[12.3.5]
\label{0.12.3.5-fr}
\emph{(\hyperref[0.12.3.4-fr]{12.3.4})의 가정 아래 환층
$\mathcal O_Y$도 연접이라고 하자. 그러면 모든 연접
$\mathcal O_Y$-가군 $\mathcal F$에 대하여 표준 준동형
(\hyperref[0.12.3.4.1-fr]{12.3.4.1})은 전단사이다.}
\end{proposition}

\begin{proof}
\oldpage[0\textsubscript{III}]{62}
문제는 $Y$ 위에서 국소적이므로 완전열
\[
0\longrightarrow\mathcal R\longrightarrow\mathcal O_Y^n
\longrightarrow\mathcal F\longrightarrow0
\]
이 존재한다고 가정할 수 있다. 이때 $\mathcal R$도 연접
$\mathcal O_Y$-가군이다
(\hyperref[0.5.3.4-ko]{0\textsubscript{I}, 5.3.4}). 준동형
\[
f^*(\mathcal{E}xt_{\mathcal O_Y}^p(\mathcal F,\mathcal G))
\longrightarrow
\mathcal{E}xt_{\mathcal O_X}^p(f^*(\mathcal F),f^*(\mathcal G))
\]
이 전단사임을 보이기 위하여 $p$에 관한 귀납법을 쓴다. $p=0$일 때는
(\hyperref[0.6.7.6.1-ko]{0\textsubscript{I}, 6.7.6.1})에서 명제가
따른다. 이제 가환도식
\[
{\tiny
\xymatrix@C=2pt@R=24pt{
f^*(\mathcal{E}xt_{\mathcal O_Y}^{p-1}(\mathcal O_Y^n,\mathcal G))
\ar[r]\ar[d]&
f^*(\mathcal{E}xt_{\mathcal O_Y}^{p-1}(\mathcal R,\mathcal G))
\ar[r]^\partial\ar[d]&
f^*(\mathcal{E}xt_{\mathcal O_Y}^{p}(\mathcal F,\mathcal G))
\ar[r]\ar[d]&
f^*(\mathcal{E}xt_{\mathcal O_Y}^{p}(\mathcal O_Y^n,\mathcal G))
\ar[d]\\
\mathcal{E}xt_{\mathcal O_X}^{p-1}(\mathcal O_X^n,f^*(\mathcal G))
\ar[r]&
\mathcal{E}xt_{\mathcal O_X}^{p-1}(f^*(\mathcal R),f^*(\mathcal G))
\ar[r]^\partial&
\mathcal{E}xt_{\mathcal O_X}^{p}(f^*(\mathcal F),f^*(\mathcal G))
\ar[r]&
\mathcal{E}xt_{\mathcal O_X}^{p}(\mathcal O_X^n,f^*(\mathcal G)).
}}
\]
을 얻는다. 여기서 $f^*(\mathcal O_Y)=\mathcal O_X$이고, $f^*$가
완전하므로 두 행은 완전하다. 또한
\[
\mathcal{E}xt_{\mathcal O_Y}^p(\mathcal O_Y^n,\mathcal G)=0
\quad\hbox{모든 }p>0\hbox{에 대하여},
\quad\hbox{그리고 마찬가지로}\quad
\mathcal{E}xt_{\mathcal O_X}^p(\mathcal O_X^n,f^*(\mathcal G))=0
\quad\hbox{모든 }p>0\hbox{에 대하여}
\]
이다 (T, 4.2.3). 귀납가정에 따라 앞 도식의 처음 두 수직 화살표는
동형이고 오른쪽 항들은 0이므로,
\[
f^*(\mathcal{E}xt_{\mathcal O_Y}^p(\mathcal F,\mathcal G))
\longrightarrow
\mathcal{E}xt_{\mathcal O_X}^p(f^*(\mathcal F),f^*(\mathcal G))
\]
은 동형이다.
\end{proof}

\subsection{직상 함자의 초코호몰로지}
\label{subsection:0.12.4-fr}

\begin{env}[12.4.1]
\label{0.12.4.1-fr}
$(X,\mathcal O_X)$, $(Y,\mathcal O_Y)$를 두 환 달린 공간이라 하고,
$f:X\to Y$를 환 달린 공간의 사상이라 하자. 임의의
$\mathcal O_X$-가군 복합체
$\mathcal K^\bullet=(\mathcal K^i)_{i\in\mathbb Z}$에 대하여
$f_*$의 \emph{초코호몰로지}를 취할 수 있다
(\hyperref[0.11.4.4-fr]{11.4.4})~; 실제로 $\mathcal O_Y$-가군들의 아벨
범주에서는 여과 귀납극한이 존재하고 완전하다 (T, 3.1.1).
초코호몰로지 $\mathcal O_Y$-가군 $R^p f_*(\mathcal K^\bullet)$는
$\mathcal H^p(f,\mathcal K^\bullet)$ 또는
$\mathcal H_f^p(\mathcal K^\bullet)$로도 나타내기로 한다.
$\mathcal H^\bullet(f,\mathcal K^\bullet)$은
$\mathcal O_Y$-가군의 이중복합체 $f_*(\mathcal L^{\bullet\bullet})$의
코호몰로지임을 상기하자. 여기서 $\mathcal L^{\bullet\bullet}$은
$\mathcal O_X$-가군들의 범주에서 $\mathcal K^\bullet$의
카르탕--아일렌베르크 단사 분해이다~;
$\mathcal H^\bullet(f,\mathcal K^\bullet)$은 두 스펙트럼 열
${}'E(f,\mathcal K^\bullet)$과 ${}''E(f,\mathcal K^\bullet)$의
수렴대상이며, 그 $E_2$항들은 다음과 같다.
\begin{equation}
\label{0.12.4.1.1-fr}
\tag{12.4.1.1}
{}'E_2^{pq}=\mathcal H^p(\mathcal H^q(f,\mathcal K^\bullet))
\end{equation}
\begin{equation}
\label{0.12.4.1.2-fr}
\tag{12.4.1.2}
{}''E_2^{pq}=\mathcal H^p(f,\mathcal H^q(\mathcal K^\bullet))
\quad(=R^p f_*(\mathcal H^q(\mathcal K^\bullet))).
\end{equation}
이 공식들에서는 함자를 복합체에 적용한 것에 대한 일반 표기
$T(A^\bullet)$를 채택하였고
(\hyperref[0.11.2.1-fr]{11.2.1}), $\mathcal O_X$-가군
$\mathcal F$에 대하여 $R^p f_*(\mathcal F)$ 대신
$\mathcal H^p(f,\mathcal F)$를 쓴다. 또한 스펙트럼 열
${}'E(f,\mathcal K^\bullet)$은 언제나 \emph{정칙}임을 상기하자~;
$\mathcal K^\bullet$이 아래로 유계이면 두 스펙트럼 열
${}'E(f,\mathcal K^\bullet)$과 ${}''E(f,\mathcal K^\bullet)$은
\emph{쌍정칙}이고,
\oldpage[0\textsubscript{III}]{63}
또는 모든 $\mathcal O_X$-가군이 길이가 $\leq m$인 플라스크 분해를
갖도록 하는 정수 $m$이 존재할 때에도 그러하다
(\hyperref[0.11.4.4-fr]{11.4.4}).
\end{env}

\begin{env}[12.4.2]
\label{0.12.4.2-fr}
마찬가지로 $\mathbf H^\bullet(X,\mathcal K^\bullet)$으로
$\mathcal O_X$-가군 복합체 $\mathcal K^\bullet$에 대한 함자
$\Gamma$의 초코호몰로지를 나타내기로 한다~; 따라서
$\mathbf H^p(X,\mathcal K^\bullet)$은
$\Gamma(X,\mathcal O_X)$-가군이다. 또한
$\mathbf H^\bullet(X,\mathcal K^\bullet)$을
$\mathcal H^\bullet(f,\mathcal K^\bullet)$의 특수한 경우로 볼 수 있다.
여기서 $f$는 $(X,\mathcal O_X)$에서 환
$\Gamma(X,\mathcal O_X)$을 갖춘 한 점짜리 환 달린 공간으로 가는
사상이다.

$X$의 임의의 열린집합 $V$에 대하여
$\mathbf H^\bullet(V,\mathcal K^\bullet|V)$ 대신
$\mathbf H^\bullet(V,\mathcal K^\bullet)$을 쓰기로 한다.
\end{env}

\begin{proposition}[12.4.3]
\label{0.12.4.3-fr}
\emph{임의의 정수 $p\in\mathbb Z$에 대하여 $\mathcal O_Y$-가군
$\mathcal H^p(f,\mathcal K^\bullet)$은 $Y$ 위의 준층
$U\mapsto\mathbf H^p(f^{-1}(U),\mathcal K^\bullet)$에 대응하는 층과
표준적으로 동형이다.}
\end{proposition}

\begin{proof}
실제로 (\hyperref[0.12.4.1-fr]{12.4.1})의 표기 아래에서 코호몰로지층
$\mathcal H^p(f_*(\mathcal L^{\bullet\bullet}))$은 다음 준층에
대응하는 층이다.
\[
U\longmapsto H^p(\Gamma(U,f_*(\mathcal L^{\bullet\bullet})))
=H^p(\Gamma(f^{-1}(U),\mathcal L^{\bullet\bullet})).
\]
그런데 $\mathcal L^{\bullet\bullet}|f^{-1}(U)$가
$\mathcal K^\bullet|f^{-1}(U)$의 카르탕--아일렌베르크 단사 분해임은
분명하므로 (T, 3.1.3), 정의에 따라
\[
H^p(\Gamma(f^{-1}(U),\mathcal L^{\bullet\bullet}))
=\mathbf H^p(f^{-1}(U),\mathcal K^\bullet)
\]
이다.
\end{proof}

\begin{proposition}[12.4.4]
\label{0.12.4.4-fr}
\emph{초코호몰로지 $\mathcal H^\bullet(f,\mathcal K^\bullet)$은 다음
각 경우에 $\mathcal K^\bullet$에 관한 코호몰로지 함자이다~:}
\begin{enumerate}
\item[\emph{a)}]
\emph{$\mathcal K^\bullet$이 아래로 유계인 복합체들의 범주를 움직인다.}
\item[\emph{b)}]
\emph{모든 $\mathcal O_X$-가군이 길이가 $\leq m$인 플라스크 분해를
갖도록 하는 정수 $m$이 존재한다.}
\item[\emph{c)}]
\emph{$X$는 뇌터 공간이다.}
\end{enumerate}
\end{proposition}

\begin{proof}
경우 a)와 b)는 (\hyperref[0.11.5.4-ko]{11.5.4})의 특수한 경우이다.
한편 이 경우 함자 $f_*$가 귀납극한과 가환함이 알려져 있으므로
(G, II, 3.10.1), 경우 c)는
(\hyperref[0.11.5.2-ko]{11.5.2})에서 따른다.
\end{proof}

\begin{env}[12.4.5]
\label{0.12.4.5-fr}
이제 $X$의 열린 덮개 $\mathfrak U=(U_\alpha)$를 생각하자. $X$ 위의
임의의 \emph{준층} 복합체
$\mathcal K^\bullet=(\mathcal K^j)$에 대하여 이중복합체
$\check C^\bullet(\mathfrak U,\mathcal K^\bullet)$를 생각하자. 이
이중복합체의 $(i,j)$ 성분은
$\check C^i(\mathfrak U,\mathcal K^j)$이며, 이는
$\mathcal K^j$에 값을 갖는 $\mathfrak U$의 신경의 교대
$i$-공사슬군이다 (G, II, 5.1). 이 이중복합체의 코호몰로지를
\emph{$\mathcal K^\bullet$을 계수로 하는 덮개 $\mathfrak U$의
초코호몰로지}라 하고, 다음과 같이 나타내기로 한다.
\[
\mathbf H^\bullet(\mathfrak U,\mathcal K^\bullet)
=H^\bullet(\check C^\bullet(\mathfrak U,\mathcal K^\bullet)).
\]
덮개의 르레 스펙트럼 열 (T, 3.8.1 및 G, II, 5.9.1)은 다음과 같이
초코호몰로지로 일반화된다~:
\end{env}

\begin{proposition}[12.4.6]
\label{0.12.4.6-fr}
\emph{$\mathcal K^\bullet=(\mathcal K^i)$를 $\mathcal O_X$-가군의
복합체라 하자. $\mathcal K^\bullet$에 관하여 정칙인 스펙트럼 함자로서
초코호몰로지 $\mathbf H^\bullet(X,\mathcal K^\bullet)$로 수렴하고,
그 $E_2$항이 다음과 같은 것이 존재한다.}
\begin{equation}
\label{0.12.4.6.1-fr}
\tag{12.4.6.1}
E_2^{pq}=\check H^p(\mathfrak U,h^q(\mathcal K^\bullet)),
\end{equation}
\emph{여기서 $h^q(\mathcal K^\bullet)$은 $X$ 위의 준층 복합체
$V\mapsto H^q(V,\mathcal K^\bullet)$를 나타낸다. 앞의 스펙트럼 열은
$\mathcal K^\bullet$이 아래로 유계이면 쌍정칙이다.}
\end{proposition}

\begin{proof}
$\mathcal K^\bullet$의 카르탕--아일렌베르크 단사 분해
$\mathcal L^{\bullet\bullet}$와 삼중복합체
$\check C^\bullet(\mathfrak U,\mathcal L^{\bullet\bullet})
=(\check C^i(\mathfrak U,\mathcal L^{jk}))$를 생각하자~; 먼저 이
삼중복합체를 차수 $i$와 $j+k$에 관한 이중복합체로 보자. $i$는
$\geq0$인 값만 취하므로 이 이중복합체의 두 번째 스펙트럼 열은
정칙이다 (\hyperref[0.11.3.3-fr]{11.3.3}). 또한 모든 $q>0$에 대하여
$\check H^q(\mathfrak U,\mathcal L^{jk})=0$이므로 퇴화한다. 실제로
$\mathcal O_X$-가군 $\mathcal L^{jk}$들은 플라스크층이다
\oldpage[0\textsubscript{III}]{64}
(G, II, 5.2.3). 따라서 (\hyperref[0.11.1.6-fr]{11.1.6})에 의해 표준
동형
\[
H^n(\check C^\bullet(\mathfrak U,\mathcal L^{\bullet\bullet}))
\xrightarrow{\sim}
H^n(\Gamma(X,\mathcal L^{\bullet\bullet}))
\]
을 얻는다 ((G, II, 5.2.2)에 의함). 그러므로 정의
(\hyperref[0.12.4.2-fr]{12.4.2})에 따라 동형
\[
H^n(\check C^\bullet(\mathfrak U,\mathcal L^{\bullet\bullet}))
\xrightarrow{\sim}\mathbf H^n(X,\mathcal K^\bullet)
\]
을 얻는다. 한편 삼중복합체
$\check C^\bullet(\mathfrak U,\mathcal L^{\bullet\bullet})$를 차수
$i+j$와 $k$에 관한 이중복합체로 보자. $k$는 $\geq0$인 값만 취하므로
이 이중복합체의 첫 번째 스펙트럼 열은 언제나 정칙이다~;
$\mathcal L^{jk}=0$ $(j<j_0)$이면, 즉 $\mathcal K^\bullet$이 아래로
유계이면 쌍정칙이다 (\hyperref[0.11.3.3-fr]{11.3.3}). 이 스펙트럼 열이
구하는 것이다~; 실제로 모든 $j$에 대하여
$\mathcal L^{j,\bullet}$은 $\mathcal K^j$의 단사 분해이다~; 따라서
$H^q(\check C^i(\mathfrak U,\mathcal L^{j,\bullet}))$은 다름 아닌
공사슬 복합체 $\check C^i(\mathfrak U,h^q(\mathcal K^j))$이다. 이로써
증명이 끝난다.
\end{proof}

\begin{corollary}[12.4.7]
\label{0.12.4.7-fr}
\emph{$\mathfrak U$의 신경의 모든 단체 $\sigma$와 모든 정수 $i$에
대하여 $H^q(U_\sigma,\mathcal K^i)=0$ $(q>0)$이면 표준 동형}
\begin{equation}
\label{0.12.4.7.1-fr}
\tag{12.4.7.1}
\mathbf H^\bullet(\mathfrak U,\mathcal K^\bullet)
\xrightarrow{\sim}\mathbf H^\bullet(X,\mathcal K^\bullet).
\end{equation}
\emph{이 존재한다.}
\end{corollary}

\begin{proof}
실제로 가정에 의해 $q>0$이면
$\check C^\bullet(\mathfrak U,h^q(\mathcal K^i))=0$이고, 따라서
$E_2^{pq}=0$이다~; 열 (\hyperref[0.12.4.6.1-fr]{12.4.6.1})이 퇴화하고
정칙이므로, 결론은 $\mathbf H^\bullet(\mathfrak U,\mathcal K^\bullet)$의
정의 (\hyperref[0.12.4.5-fr]{12.4.5})와
(\hyperref[0.11.1.6-fr]{11.1.6})에서 따른다.
\end{proof}

\begin{env}[12.4.8]
\label{0.12.4.8-fr}
$(X',\mathcal O_{X'})$을 두 번째 환 달린 공간이라 하고,
$f=(\psi,\theta)$를 $X'$에서 $X$로 가는 사상이라 하자.
(\hyperref[0.12.1.3-fr]{12.1.3})과
(\hyperref[0.12.1.4-fr]{12.1.4})에서와 같은 방법으로
$\mathcal O_X$-가군 복합체 $\mathcal K^\bullet$의 초코호몰로지에 대한
디-준동형
\begin{equation}
\label{0.12.4.8.1-fr}
\tag{12.4.8.1}
\mathbf H^p(X,\mathcal K^\bullet)
\longrightarrow\mathbf H^p(X',f^*(\mathcal K^\bullet)).
\end{equation}
을 정의한다. $\mathcal K^\bullet$의 카르탕--아일렌베르크 단사 분해
$\mathcal L^{\bullet\bullet}$에서 출발하자. $\psi^*$가 완전하므로
$\psi^*(\mathcal L^{\bullet\bullet})$은 $\psi^*(\mathcal O_X)$-가군들의
범주에서 $\psi^*(\mathcal K^\bullet)$의 카르탕--아일렌베르크 분해이다~;
이때 사상
$\psi^*(\mathcal L^{\bullet\bullet})\to
\mathcal L'^{\bullet\bullet}$이 존재한다. 여기서
$\mathcal L'^{\bullet\bullet}$은 $\psi^*(\mathcal K^\bullet)$의
카르탕--아일렌베르크 단사 분해이다. 이로부터 코호몰로지에 대한 사상
\[
\mathbf H^\bullet(X,\mathcal K^\bullet)
\longrightarrow\mathbf H^\bullet(X',\psi^*(\mathcal K^\bullet))~;
\]
을 얻는다. 이를 함자성으로부터 얻는 사상
$\psi^*(\mathcal K^\bullet)\to f^*(\mathcal K^\bullet)$에서 유도되는
사상과 합성하면, 구하는 사상
(\hyperref[0.12.4.8.1-fr]{12.4.8.1})을 얻는다.

(\hyperref[0.12.4.8.1-fr]{12.4.8.1})과
(\hyperref[0.12.4.3-fr]{12.4.3})에서 출발하여
(\hyperref[0.12.2.5-fr]{12.2.5})에서와 같이 논하면, 환 달린 공간의
두 사상 $f:X\to Y$, $g:Y\to Z$에 대하여 $\mathcal O_X$-가군 복합체
$\mathcal K^\bullet$의 초코호몰로지에 대한 준동형
\begin{equation}
\label{0.12.4.8.2-fr}
\tag{12.4.8.2}
\mathcal H^n(g,f_*(\mathcal K^\bullet))
\longrightarrow\mathcal H^n(h,\mathcal K^\bullet)
\end{equation}
\begin{equation}
\label{0.12.4.8.3-fr}
\tag{12.4.8.3}
\mathcal H^n(h,\mathcal K^\bullet)
\longrightarrow g_*(\mathcal H^n(f,\mathcal K^\bullet)).
\end{equation}
들을 정의할 수 있다. 정의의 세부사항은 독자에게 맡긴다.
\end{env}

\section{호몰로지 대수학에서의 사영극한}
\label{section:0.13-fr}

\subsection{미타그--레플러 조건}
\label{subsection:0.13.1-fr}

\begin{env}[13.1.1]
\label{0.13.1.1-fr}
$\mathcal C$를 무한곱이 존재하는 아벨 범주라 하자 (T, 1.5의 공리
AB~3*). 그러면 $\mathcal C$의 한 대상의 부분대상들로 이루어진 임의의
족의 하한이 존재하고, $\mathcal C$의 대상들의 모든 사영계는
사영극한을 갖는다. 이 사영극한은 고려하는 사영계의 왼쪽 완전함자이다
(T, 1.8). $(A_\alpha,f_{\alpha\beta})$를 지표집합 $I$가 오른쪽
여과인 $\mathcal C$의 대상들의 사영계라 하자.
\oldpage[0\textsubscript{III}]{65}
$A=\varprojlim A_\alpha$라 하고, 모든 $\alpha\in I$에 대하여
$f_\alpha:A\to A_\alpha$를 표준 사상이라 하자. 각 $\alpha\in I$에
대하여 $\alpha\leq\beta$일 때의 $f_{\alpha\beta}(A_\beta)$들은
$A_\alpha$의 부분대상들로 이루어진 감소 여과족을 이룬다. 부분대상
\[
A'_\alpha=\inf_{\beta\geq\alpha}f_{\alpha\beta}(A_\beta)
\]
을 $A_\alpha$ 안의 ``보편상'' 부분대상이라 한다. 명백히
$f_\alpha(A)\subset A'_\alpha$이고, $\alpha\leq\beta$일 때
$f_{\alpha\beta}(A'_\beta)\subset A'_\alpha$이다. 따라서
$(A'_\alpha,f_{\alpha\beta}|A'_\beta)$는 사영계이고
$A=\varprojlim A'_\alpha$이다.
\end{env}

\begin{env}[13.1.2]
\label{0.13.1.2-fr}
$\mathcal C$ 안의 사영계 $(A_\alpha,f_{\alpha\beta})$에 대하여 다음
조건을 \emph{미타그--레플러 조건}이라 한다.
\begin{itemize}
\item[\textbf{(ML)}]
\emph{모든 지표 $\alpha$에 대하여 $\beta\geq\alpha$가 존재하여,
모든 $\gamma\geq\beta$에 대해
$f_{\alpha\gamma}(A_\gamma)=f_{\alpha\beta}(A_\beta)$이다.}
\end{itemize}
$f_{\alpha\beta}$들이 \emph{에피사상}이면 조건 (ML)이 성립함은
명백하다. 역으로 (ML)이 성립하고, 모든 $\alpha\in I$에 대하여
$A'_\alpha$가 $A_\alpha$ 안의 ``보편상'' 부분대상이면,
$\alpha\leq\beta$일 때 $f_{\alpha\beta}$의 $A'_\beta$에 대한 제한은
\emph{에피사상} $A'_\beta\to A'_\alpha$이다. 실제로
$\delta\geq\gamma$일 때
$f_{\beta\delta}(A_\delta)=f_{\beta\gamma}(A_\gamma)$가 되도록
$\gamma\geq\beta$를 택하면
$A'_\beta=f_{\beta\gamma}(A_\gamma)$이다. 한편 이는
$\delta\geq\gamma$일 때
$f_{\alpha\delta}(A_\delta)=f_{\alpha\gamma}(A_\gamma)$를 함의하므로,
$A'_\alpha=f_{\alpha\gamma}(A_\gamma)=f_{\alpha\beta}(A'_\beta)$이다.

또한 각 대상 $A_\alpha$가 $\mathcal C$에서 \emph{아르틴}, 즉
$A_\alpha$의 부분대상들로 이루어진 모든 족이 극소 원소를 가지면
조건 (ML)이 성립한다. 실제로 $A_\alpha$의 부분대상들로 이루어진
감소 여과족 $(f_{\alpha\beta}(A_\beta))$의 극소 원소는 반드시 이
부분대상들 중 가장 작은 것이다.
\end{env}

\begin{remark}[13.1.3]
\label{0.13.1.3-fr}
예를 들어 $\mathcal C$가 집합의 범주일 때에도 조건 (ML)을 같은 방식으로
서술할 수 있다. 이 경우에도 $A_\alpha$의 ``보편상'' 부분집합을 정의할
수 있고, (\hyperref[0.13.1.1-fr]{13.1.1})과
(\hyperref[0.13.1.2-fr]{13.1.2})에서 한 주석들은 여전히 성립한다.
\end{remark}

\subsection{아벨군에 대한 미타그--레플러 조건}
\label{subsection:0.13.2-fr}

\begin{proposition}[13.2.1]
\label{0.13.2.1-fr}
\emph{
\[
0\longrightarrow A_\alpha\xrightarrow{u_\alpha}B_\alpha
\xrightarrow{v_\alpha}C_\alpha\longrightarrow0
\]
을 아벨군들의 사영계들의 완전열이라 하자. 세 사영계는 같은 여과
지표집합 $I$를 갖는다.}
\begin{enumerate}
\item[\emph{(i)}] \emph{$(B_\alpha)$가 (ML)을 만족하면
$(C_\alpha)$도 (ML)을 만족한다.}
\item[\emph{(ii)}] \emph{$(A_\alpha)$와 $(C_\alpha)$가 (ML)을
만족하면 $(B_\alpha)$도 (ML)을 만족한다.}
\end{enumerate}
\end{proposition}

\begin{proof}
$(f_{\alpha\beta})$, $(g_{\alpha\beta})$,
$(h_{\alpha\beta})$를 각각 사영계 $(A_\alpha)$, $(B_\alpha)$,
$(C_\alpha)$를 정의하는 준동형계라 하자.

(i) $\lambda\geq\beta$일 때
$g_{\alpha\beta}(B_\beta)=g_{\alpha\lambda}(B_\lambda)$라고 하자.
$v_\beta$와 $v_\lambda$가 전사이므로,
\[
h_{\alpha\beta}(C_\beta)
=v_\alpha(g_{\alpha\beta}(B_\beta))
=v_\alpha(g_{\alpha\lambda}(B_\lambda))
=h_{\alpha\lambda}(C_\lambda)
\]
가 모든 $\lambda\geq\beta$에 대하여 성립한다.

(ii) $\alpha\in I$라 하고, $\lambda\geq\beta$일 때
$f_{\alpha\beta}(A_\beta)=f_{\alpha\lambda}(A_\lambda)$가 되도록
$\beta\geq\alpha$를 택하자. 또한 $\lambda\geq\gamma$일 때
$h_{\beta\gamma}(C_\gamma)=h_{\beta\lambda}(C_\lambda)$가 되도록
$\gamma\geq\beta$를 택하자. 이제
$y_\alpha\in g_{\alpha\gamma}(B_\gamma)$라 하면,
$y_\gamma\in B_\gamma$가 존재하여
$y_\alpha=g_{\alpha\gamma}(y_\gamma)$이다.
$y_\beta=g_{\beta\gamma}(y_\gamma)$로 놓으면
$v_\beta(y_\beta)=h_{\beta\gamma}(v_\gamma(y_\gamma))$이다. 모든
$\lambda\geq\gamma$에 대하여 가정에 의해 $y_\lambda\in B_\lambda$가
존재하여
\[
h_{\beta\gamma}(v_\gamma(y_\gamma))
=h_{\beta\lambda}(v_\lambda(y_\lambda))
=v_\beta(g_{\beta\lambda}(y_\lambda)).
\]
따라서 $v_\beta(y_\beta-g_{\beta\lambda}(y_\lambda))=0$이고,
$x_\beta\in A_\beta$가 존재하여
\oldpage[0\textsubscript{III}]{66}
\[
y_\beta=g_{\beta\lambda}(y_\lambda)+u_\beta(x_\beta).
\]
그러므로
\[
y_\alpha=g_{\alpha\lambda}(y_\lambda)
+u_\alpha(f_{\alpha\beta}(x_\beta)).
\]
그런데 $\lambda\geq\beta$이므로 $x_\lambda\in A_\lambda$가 존재하여
$f_{\alpha\beta}(x_\beta)=f_{\alpha\lambda}(x_\lambda)$이고, 끝으로
\[
y_\alpha
=g_{\alpha\lambda}(y_\lambda)+u_\alpha(f_{\alpha\lambda}(x_\lambda))
=g_{\alpha\lambda}(y_\lambda+u_\lambda(x_\lambda))
\in g_{\alpha\lambda}(B_\lambda).
\]
이로써 증명이 끝난다.
\end{proof}

\begin{proposition}[13.2.2]
\label{0.13.2.2-fr}
\emph{$I$를 가산 공종 부분집합을 갖는 여과 순서집합이라 하자.
\[
0\longrightarrow A_\alpha\xrightarrow{u_\alpha}B_\alpha
\xrightarrow{v_\alpha}C_\alpha\longrightarrow0
\]
을 지표집합이 $I$인 아벨군들의 사영계들의 완전열이라 하자.
$(A_\alpha)$가 조건 (ML)을 만족하면 열
\[
0\longrightarrow\varprojlim A_\alpha
\longrightarrow\varprojlim B_\alpha
\longrightarrow\varprojlim C_\alpha\longrightarrow0
\]
은 완전하다.}
\end{proposition}

\begin{proof}
준동형
\[
v=\varprojlim v_\alpha:\varprojlim B_\alpha\longrightarrow
\varprojlim C_\alpha
\]
가 전사임을 보이면 충분하다. $z=(z_\alpha)$를
$\varprojlim C_\alpha$의 원소라 하고
$E_\alpha=v_\alpha^{-1}(z_\alpha)$로 놓자. 준동형
$g_{\alpha\beta}:B_\beta\to B_\alpha$의 제한에 대하여
$E_\alpha$들이 공집합이 아닌 집합들의 사영계를 이룸은 명백하다.
이 사영계가 조건 (ML)을 만족함을 보이자. $u_\alpha$를 통하여
$A_\alpha$를 $B_\alpha$의 부분집합으로 동일시한다. 모든
$\alpha\in I$에 대하여 $\beta\geq\alpha$가 존재하여
$\lambda\geq\beta$이면
$g_{\alpha\beta}(A_\beta)=g_{\alpha\lambda}(A_\lambda)$이다. 이때
$\lambda\geq\beta$에 대하여
$g_{\alpha\beta}(E_\beta)=g_{\alpha\lambda}(E_\lambda)$도 성립함을
보이자. 실제로 $y_\lambda\in E_\lambda$를 택하고
$y_\beta=g_{\beta\lambda}(y_\lambda)$,
$y_\alpha=g_{\alpha\lambda}(y_\lambda)$로 놓자. 또한
$y'_\alpha\in g_{\alpha\beta}(E_\beta)$라 하여, 어떤
$y'_\beta\in E_\beta$에 대하여
$y'_\alpha=g_{\alpha\beta}(y'_\beta)$라 하자. 그러면
$y'_\beta-y_\beta=x_\beta\in A_\beta$이고, 가정에 따라
$x_\lambda\in A_\lambda$가 존재하여
$g_{\alpha\beta}(x_\beta)=g_{\alpha\lambda}(x_\lambda)$이다. 따라서
\[
y'_\alpha=g_{\alpha\beta}(y_\beta)+g_{\alpha\beta}(x_\beta)
=g_{\alpha\lambda}(y_\lambda)+g_{\alpha\lambda}(x_\lambda)
=g_{\alpha\lambda}(y_\lambda+x_\lambda)
\in g_{\alpha\lambda}(E_\lambda).
\]
주장이 증명되었다. 이제 $I$에 관한 가정 아래 (ML)을 만족하는 공집합이
아닌 집합들의 사영계는 공집합이 아닌 사영극한을 가짐이 알려져 있다
(Bourbaki, \emph{Top. gén.}, chap.~II, 3\textsuperscript{e} éd., §~3,
th.~1). 따라서 $y=(y_\alpha)\in\varprojlim E_\alpha$가 존재하고,
정의에 의해 모든 $\alpha$에 대하여 $v_\alpha(y_\alpha)=z_\alpha$이므로
$z=v(y)$이다. 증명 끝.
\end{proof}

\begin{proposition}[13.2.3]
\label{0.13.2.3-fr}
\emph{$I$에 관한 가정은 (\hyperref[0.13.2.2-fr]{13.2.2})와 같다고
하자. $(K_\alpha^\bullet)_{\alpha\in I}$를 미분작용소의 차수가 $+1$인
아벨군 복합체
$K_\alpha^\bullet=(K_\alpha^n)_{n\in\mathbb Z}$들의 사영계라 하자.
각 $n$에 대하여 함자적인 표준 준동형}
\begin{equation}
\label{0.13.2.3.1-fr}
\tag{13.2.3.1}
h_n:H^n(\varprojlim K_\alpha^\bullet)
\longrightarrow\varprojlim H^n(K_\alpha^\bullet)
\end{equation}
\emph{이 존재한다. 모든 차수 $n$에 대하여 아벨군들의 사영계
$(K_\alpha^n)_{\alpha\in I}$가 (ML)을 만족하면 모든 준동형 $h_n$은
전사이다. 더 나아가 한 차수 $n$에 대하여 사영계
$(H^{n-1}(K_\alpha^\bullet))_{\alpha\in I}$가 (ML)을 만족하면
$h_n$은 전단사이다.}
\end{proposition}

\begin{proof}
모든 $n$에 대하여 $K^n=\varprojlim K_\alpha^n$로 놓자. 준동형
$h_n$의 정의는 도식
\[
\xymatrix{
\cdots\ar[r]&K^{n-1}\ar[r]\ar[d]&K^n\ar[r]\ar[d]
&K^{n+1}\ar[r]\ar[d]&\cdots\\
\cdots\ar[r]&K_\alpha^{n-1}\ar[r]&K_\alpha^n\ar[r]
&K_\alpha^{n+1}\ar[r]&\cdots
}
\]
의 가환성에서 온다.
\oldpage[0\textsubscript{III}]{67}
여기서 $K^\bullet$의 미분작용소들은 $K_\alpha^\bullet$의 대응하는
미분작용소들의 사영극한이다.

완전열들
\[
(*_n)\qquad
0\longrightarrow \mathrm B^n(K_\alpha^\bullet)
\longrightarrow \mathrm Z^n(K_\alpha^\bullet)
\longrightarrow H^n(K_\alpha^\bullet)\longrightarrow0
\]
\[
(**_n)\qquad
0\longrightarrow \mathrm Z^{n-1}(K_\alpha^\bullet)
\longrightarrow K_\alpha^{n-1}
\longrightarrow \mathrm B^n(K_\alpha^\bullet)\longrightarrow0
\]
을 생각하자. 가정과 명제
(\hyperref[0.13.2.1-fr]{13.2.1}, (i))에 따라 모든 $n$에 대하여
사영계 $(\mathrm B^n(K_\alpha^\bullet))_{\alpha\in I}$가 (ML)을
만족한다. 그러므로 (\hyperref[0.13.2.2-fr]{13.2.2})에 따라 열
\[
(***_n)\qquad
0\longrightarrow\varprojlim_\alpha\mathrm B^n(K_\alpha^\bullet)
\longrightarrow\varprojlim_\alpha\mathrm Z^n(K_\alpha^\bullet)
\longrightarrow\varprojlim_\alpha H^n(K_\alpha^\bullet)
\longrightarrow0
\]
은 완전하다. 그런데
$\varprojlim_\alpha\mathrm B^n(K_\alpha^\bullet)$은
$\mathrm B^n(K^\bullet)$을 포함하는 $K^n$의 부분군과 동일시되고,
$\varprojlim_\alpha\mathrm Z^n(K_\alpha^\bullet)$은
$\mathrm Z^n(K^\bullet)$의 부분군과 동일시된다. 따라서 $h_n$은
전사이다. 이제 사영계
$(H^{n-1}(K_\alpha^\bullet))_{\alpha\in I}$도 (ML)을 만족한다고
가정하면, 완전열 $(*_{n-1})$과 명제
(\hyperref[0.13.2.1-fr]{13.2.1}, (ii))에 따라 사영계
$(\mathrm Z^{n-1}(K_\alpha^\bullet))_{\alpha\in I}$가 (ML)을
만족한다. 그러면 (\hyperref[0.13.2.2-fr]{13.2.2})를 완전열
$(**_n)$에 적용하여 완전열
\[
0\longrightarrow
\varprojlim_\alpha\mathrm Z^{n-1}(K_\alpha^\bullet)
\longrightarrow K^{n-1}\xrightarrow{u}
\varprojlim_\alpha\mathrm B^n(K_\alpha^\bullet)
\longrightarrow0
\]
을 얻는다. 이제
$\varprojlim_\alpha\mathrm B^n(K_\alpha^\bullet)
\supset\mathrm B^n(K^\bullet)$이고, $u$와 포함
$\varprojlim_\alpha\mathrm B^n(K_\alpha^\bullet)\to K^n$의 합성이
미분작용소 $K^{n-1}\to K^n$이다. 따라서 $u$가 전사이면
$\varprojlim_\alpha\mathrm B^n(K_\alpha^\bullet)
=\mathrm B^n(K^\bullet)$이고, 이에 따라 $h_n$은 단사이다. 증명 끝.
\end{proof}

\begin{env}[13.2.4]
\label{0.13.2.4-fr}
\emph{주석} (13.2.4).
\begin{enumerate}
\item[\emph{(i)}]
(\hyperref[0.13.2.2-fr]{13.2.2})의 논증(Bourbaki,
\emph{같은 곳} 참조)에 따르면, 다음 조건들만을 가정해도 그 명제의
결론은 여전히 성립한다. 각 $A_\alpha$에 \emph{완비 거리화 가능}
공간의 구조를 줄 수 있고, 그 구조에서 평행이동은 위상동형이며,
사영계 $(A_\alpha)$를 정의하는 사상
$f_{\alpha\beta}:A_\beta\to A_\alpha$는 고려하는 거리들에 대하여
\emph{균등연속}이고, 끝으로 사영계 $(A_\alpha)$는 다음 조건을
만족한다고 하자.
\begin{itemize}
\item[\textbf{(ML')}]
\emph{모든 지표 $\alpha$에 대하여 $\beta\geq\alpha$가 존재하여,
모든 $\gamma\geq\beta$에 대해 $f_{\alpha\gamma}(A_\gamma)$가
$f_{\alpha\beta}(A_\beta)$에서 조밀하다.}
\end{itemize}

이로부터 (\hyperref[0.13.2.3-fr]{13.2.3})과 비슷한 보충을 얻는다.
모든 $\alpha$와 $n<0$에 대하여 $K_\alpha^n=0$이라고 하자. 또한
$n\geq0$일 때 $(K_\alpha^n)_{\alpha\in I}$가 (ML)을 만족하고,
$A_\alpha=H^0(K_\alpha^\bullet)$에 위의 성질들을 만족하는 거리공간의
구조를 줄 수 있다고 하자. 그러면 $n\geq2$에 대해서는
(\hyperref[0.13.2.3-fr]{13.2.3})의 결론들이 바뀌지 않고, 더 나아가
$h_1$은 \emph{전단사}이다. 실제로
(\hyperref[0.13.2.2-fr]{13.2.2})의 논증은 사영계
$(\mathrm B^1(K_\alpha^\bullet))_{\alpha\in I}$가 (ML)을 만족하고,
열 $(***_1)$이 완전하며, 위의 결과에 따라
$\varprojlim\mathrm B^1(K_\alpha^\bullet)=\mathrm B^1(K^\bullet)$임을
여전히 보여 준다. 이로써 특히 (T, 3.10.2)의 주장들을 얻었다.

\item[\emph{(ii)}]
함자 $\varprojlim$의 오른쪽 유도함자 $\varprojlim^{(n)}$을 도입하여
앞의 명제들보다 더 완전한 명제들을 얻을 수 있다 [28].
\end{enumerate}
\end{env}

\subsection{응용: 층들의 사영극한의 코호몰로지}
\label{subsection:0.13.3-fr}

\begin{proposition}[13.3.1]
\label{0.13.3.1-fr}
\oldpage[0\textsubscript{III}]{68}
\emph{$X$를 위상공간,
$(\mathcal F_k)_{k\in\mathbb N}$를 $X$ 위의 아벨군의 층들로 이루어진
사영계라 하고,
$\mathcal F=\varprojlim_k\mathcal F_k$라 하자. 다음 조건들이 성립한다고
가정한다.}
\begin{enumerate}
\item[\emph{(i)}]
\emph{$X$의 위상의 기저 $\mathfrak B$가 존재하여, 모든
$U\in\mathfrak B$와 모든 $i\geq0$에 대하여 사영계
$(H^i(U,\mathcal F_k))_{k\in\mathbb N}$는 (ML)을 만족한다.}
\item[\emph{(ii)}]
\emph{모든 $x\in X$와 모든 $i>0$에 대하여}
\[
\varinjlim_U\bigl(\varprojlim_k H^i(U,\mathcal F_k)\bigr)=0
\]
\emph{이다. 여기서 $U$는 $\mathfrak B$에 속하는 $x$의 근방들의 집합을
움직인다.}
\item[\emph{(iii)}]
\emph{사영계 $(\mathcal F_k)$를 정의하는 준동형
$u_{hk}:\mathcal F_k\to\mathcal F_h$ $(h\leq k)$는 전사이다.}
\end{enumerate}
\emph{그러면 모든 $i>0$에 대하여 표준 준동형}
\[
h_i:H^i(X,\mathcal F)\longrightarrow
\varprojlim_k H^i(X,\mathcal F_k)
\]
\emph{은 전사이다~; 더 나아가 주어진 $i$에 대하여 사영계
$(H^{i-1}(X,\mathcal F_k))_{k\in\mathbb N}$가 (ML)을 만족하면 $h_i$는
전단사이다.}
\end{proposition}

\begin{proof}
\emph{a)} 먼저 $\mathcal F_k$들과 $u_{hk}$들의 핵 $\mathcal N_{hk}$들이
플라스크라고 가정하자~; 이때 명제의 조건 (iii)로부터
$i>0$이면 $H^i(X,\mathcal F)=0$임을 보이겠다. $X$의 모든 열린집합
$U$와 $U$의 열린집합들로 이루어진 $U$의 모든 덮개 $\mathfrak U$에
대하여 $i>0$이면 $\check H^i(\mathfrak U,\mathcal F)=0$임을 보이면 충분하다.
실제로 이로부터 먼저 체흐 코호몰로지에 대하여 모든 $i>0$에 대해
$\check H^i(U,\mathcal F)=0$을 얻고, 이어서 ($X$의 모든 열린집합의 집합에
(G, II, 5.9.2)를 적용하여) 모든 $i>0$에 대해
$H^i(X,\mathcal F)=0$을 얻는다. $\mathcal F_k$들은 플라스크이므로
$i>0$이면 $\check H^i(\mathfrak U,\mathcal F_k)=0$이다
(G, II, 5.2.3)~; 각 $k$에 대하여 덮개 $\mathfrak U$의 신경의 교대
공사슬 복합체 $\check C^\bullet(\mathfrak U,\mathcal F_k)$를 생각하자
(G, II, 5.1). 이들은 자명하게 아벨군의 복합체들로 이루어진 사영계를
이룬다. 모든 사상
$\check C^\bullet(\mathfrak U,\mathcal F_k)\to
\check C^\bullet(\mathfrak U,\mathcal F_h)$ $(h\leq k)$가 전사임을 보이자.
정의에 따라, $X$의 모든 열린집합 $V$에 대하여 사상
$\Gamma(V,\mathcal F_k)\to\Gamma(V,\mathcal F_h)$가 전사임을 보이면
충분하다~; 그런데 가정에 의해 열
\[
0\longrightarrow\mathcal N_{hk}\longrightarrow\mathcal F_k
\longrightarrow\mathcal F_h\longrightarrow0
\]
은 완전하므로 코호몰로지 완전열
\[
\Gamma(V,\mathcal F_k)\longrightarrow\Gamma(V,\mathcal F_h)
\longrightarrow H^1(V,\mathcal N_{hk})=0
\]
을 준다. 여기서 $\mathcal N_{hk}$가 플라스크임을 썼다. 따라서 사영계
$(\check C^\bullet(\mathfrak U,\mathcal F_k))_{k\in\mathbb N}$는
(ML)을 만족한다~; 모든 $i\geq0$에 대하여
$(\check H^i(\mathfrak U,\mathcal F_k))_{k\in\mathbb N}$도 그러하다.
실제로 $i>0$일 때에는 자명하고,
$\check H^0(\mathfrak U,\mathcal F_k)=\Gamma(U,\mathcal F_k)$이므로
(G, II, 5.2.2), 앞에서 본 바에 의해 $i=0$일 때에도 (ML)이 성립한다.
따라서 (\hyperref[0.13.2.3-fr]{13.2.3})을 적용하면
\[
\check H^i(\mathfrak U,\mathcal F)
=\varprojlim_k\check H^i(\mathfrak U,\mathcal F_k)=0
\quad\hbox{모든 }i>0\hbox{에 대하여}.
\]

\emph{b)} 이제 일반적인 경우로 돌아가자. 각 $k\in\mathbb N$에 대하여
$\mathcal F_k$의 플라스크층들에 의한 표준 분해
$\mathcal C^\bullet(X,\mathcal F_k)
=(\mathcal C^i(X,\mathcal F_k))_{i\geq0}$를 생각하자
(G, II, 4.3). 각 $i\geq0$에 대하여
$(\mathcal C^i(X,\mathcal F_k))_{k\in\mathbb N}$가 플라스크층들로
이루어진 사영계임은 분명하다~; 이것이 \emph{a)}의 조건들을
만족함을 보이자.
\oldpage[0\textsubscript{III}]{69}
실제로 $h\leq k$일 때 $\mathcal N_{hk}$를 $u_{hk}$의 핵이라 하면,
열
\[
0\longrightarrow\mathcal N_{hk}\longrightarrow\mathcal F_k
\longrightarrow\mathcal F_h\longrightarrow0
\]
은 (iii)에 의해 완전하고, 우리의 주장은 함자
$\mathcal A\mapsto\mathcal C^i(X,\mathcal A)$가 완전하다는 사실에서
따른다 (G, II, 4.3). 다음과 같이 놓자.
$\mathcal G^i=\varprojlim_k\mathcal C^i(X,\mathcal F_k)$.
그러면 \emph{a)}에 의해 $j>0$, $i\geq0$이면
$H^j(X,\mathcal G^i)=0$이다. 이제
$\mathcal G^\bullet=(\mathcal G^i)_{i\geq0}$가 층 $\mathcal F$의
분해임을 보이자~; $j>0$이면 $H^j(X,\mathcal G^i)=0$이므로 코호몰로지
$H^\bullet(X,\mathcal F)$는 $H^\bullet(\Gamma(X,\mathcal G^\bullet))$와
같다 (G, II, 4.7.1).

사영극한을 취하면 완전열들
\[
0\longrightarrow\mathcal F_k
\longrightarrow\mathcal C^0(X,\mathcal F_k)
\longrightarrow\mathcal C^1(X,\mathcal F_k)\longrightarrow\cdots
\]
로부터 아벨군의 층들의 복합체
\[
0\longrightarrow\mathcal F\longrightarrow\mathcal G^0
\longrightarrow\mathcal G^1\longrightarrow\cdots
\]
를 얻음은 분명하다. 우리의 주장을 증명하려면 $i>0$일 때
$\mathcal H^i(\mathcal G^\bullet)=0$임을 보여야 한다. 이 층은 준층
$U\mapsto H^i(\Gamma(U,\mathcal G^\bullet))$에 의해 생성된다
(G, II, 4.1)~; 한편 복합체 $\Gamma(U,\mathcal G^\bullet)$는 아벨군의
복합체들의 사영계
\[
(\Gamma(U,\mathcal C^\bullet(X,\mathcal F_k)))_{k\in\mathbb N}
\]
의 사영극한이다 (0\textsubscript{I}, 3.2.6). \emph{a)}에서 각
$i\geq0$에 대하여 사상
$\Gamma(U,\mathcal C^i(X,\mathcal F_k))\to
\Gamma(U,\mathcal C^i(X,\mathcal F_h))$ $(h\leq k)$가 전사임을 보았다~;
다른 한편 표준 분해 $\mathcal C^\bullet(U,\mathcal F_k|U)$는
$\mathcal C^\bullet(X,\mathcal F_k)$가 $U$ 위에 유도하는 것이므로
$H^i(U,\mathcal F_k)=
H^i(\Gamma(U,\mathcal C^\bullet(X,\mathcal F_k)))$이다~; 따라서 가정
(i)에 의해 모든 $U\in\mathfrak B$에 대하여 복합체의 사영계
$(\Gamma(U,\mathcal C^\bullet(X,\mathcal F_k)))_{k\in\mathbb N}$에
(\hyperref[0.13.2.3-fr]{13.2.3})을 적용할 수 있고, 이에 따라
\[
H^i(\Gamma(U,\mathcal G^\bullet))
=\varprojlim_k H^i(U,\mathcal F_k)
\quad\hbox{모든 }i\geq0\hbox{에 대하여}.
\]
이제 가정 (ii)로부터 정의에 의해 $i>0$이면 층
$\mathcal H^i(\mathcal G^\bullet)$가 0임이 바로 따른다.

따라서 모든 $i\geq0$에 대하여
\[
H^i(X,\mathcal F)=H^i(\Gamma(X,\mathcal G^\bullet))
\quad\hbox{이고}\quad
\Gamma(X,\mathcal G^\bullet)
=\varprojlim_k\Gamma(X,\mathcal C^\bullet(X,\mathcal F_k)).
\]
방금 사상들
$\Gamma(X,\mathcal C^i(X,\mathcal F_k))\to
\Gamma(X,\mathcal C^i(X,\mathcal F_h))$ $(h\leq k)$가 모두
전사임을 보았다~; 따라서 결론은 다시
(\hyperref[0.13.2.3-fr]{13.2.3})에서 따른다.
\end{proof}

\begin{env}[13.3.2]
\label{0.13.3.2-fr}
\emph{주의} (13.3.2).
\begin{enumerate}
\item[\emph{(i)}]
(\hyperref[0.13.3.1-fr]{13.3.1})의 명제는 $i>0$일 때에만 내용이 있다.
$i=0$이면 아무 가정 없이도 $h_i$가 언제나 동형이기 때문이다
(0\textsubscript{I}, 3.2.6).
\item[\emph{(ii)}]
모든 $k$, 모든 $i>0$, 모든 $U\in\mathfrak B$에 대하여
$H^i(U,\mathcal F_k)=0$이고, $U\in\mathfrak B$일 때 사상
$\Gamma(U,\mathcal F_k)\to\Gamma(U,\mathcal F_h)$가 전사이면,
(\hyperref[0.13.3.1-fr]{13.3.1})의 조건 (i)와 (ii)는 특히 성립한다.
이 경우가 (\hyperref[0.13.3.1-fr]{13.3.1})의 가장 흔한 적용 사례이다.
\end{enumerate}
\end{env}

\subsection{미타그--레플러 조건과 사영계의 연관 등급 대상}
\label{subsection:0.13.4-fr}
\label{0.13.4-fr}

\begin{env}[13.4.1]
\label{0.13.4.1-fr}
$\mathbf A=(A_k,u_{kh})_{k\in\mathbb Z}$를 아벨 범주 $C$의 사영계라
하자~; 어떤 $k_0$가 존재하여 $k<k_0$이면 $A_k=0$일 때 이 사영계가
\emph{아래로 유계}라고 하겠다. 각 $A_k$ 위에 다음 공식으로
\emph{여과} $(\mathrm F^p(A_k))_{p\in\mathbb Z}$를 정의한다.
\begin{equation}
\label{0.13.4.1.1-fr}
\tag{13.4.1.1}
\mathrm F^p(A_k)=
\begin{cases}
\operatorname{Ker}(A_k\longrightarrow A_{p-1})&\text{$p\leq k+1$일 때},\\
0&\text{$p\geq k+1$일 때}.
\end{cases}
\end{equation}
\oldpage[0\textsubscript{III}]{70}
따라서 가정에 의해 $\mathrm F^{k_0}(A_k)=A_k$이고
$\mathrm F^{k+1}(A_k)=0$이다. 다시 말해 이 여과는 \emph{유한}이다
(\hyperref[0.11.1.3-fr]{11.1.3}). 따라서 이 여과의 연관 등급 대상들은
\[
\operatorname{gr}^p(A_k)=
\operatorname{Ker}(A_k\longrightarrow A_{p-1})
/\operatorname{Ker}(A_k\longrightarrow A_p)
\]
이고, 이에 따라 $\operatorname{gr}^p(A_k)$는
$\operatorname{Ker}(A_k\to A_{p-1})$의 $A_k\to A_p$에 의한 상과
동형이다~; 사영계를 정의하는 사상들의 추이성에 의해
\begin{equation}
\label{0.13.4.1.2-fr}
\tag{13.4.1.2}
\operatorname{gr}^p(A_k)=
\operatorname{Ker}(A_p\longrightarrow A_{p-1})
\cap\operatorname{Im}(A_k\longrightarrow A_p)
\end{equation}
이다. 그런데 (\hyperref[0.13.4.1.1-fr]{13.4.1.1})에 의해
$\operatorname{Ker}(A_p\to A_{p-1})=\operatorname{gr}^p(A_p)$이므로,
또한
\begin{equation}
\label{0.13.4.1.3-fr}
\tag{13.4.1.3}
\operatorname{gr}^p(A_k)=
\operatorname{gr}^p(A_p)\cap\operatorname{Im}(A_k\longrightarrow A_p).
\end{equation}
앞의 정의들로부터 $k\leq h$일 때
\[
u_{kh}(\mathrm F^p(A_h))\subset\mathrm F^p(A_k)
\]
임도 알 수 있다. 따라서 모든 $p\in\mathbb Z$에 대하여
$\operatorname{gr}^p(u_{kh})$들은 \emph{사영계}
$(\operatorname{gr}^p(A_k))_{k\in\mathbb Z}$를 정의한다.
\end{env}

\begin{env}[13.4.2]
\label{0.13.4.2-fr}
충분히 큰 $k$에 대하여 사상 $A_{k+1}\to A_k$가 \emph{동형}이면
사영계 $\mathbf A$가 \emph{본질적으로 상수}라고 하겠다. 사상
$A_i\to A_j$ $(j\leq i)$가 \emph{에피사상}이면 사영계 $\mathbf A$가
\emph{엄밀}하다고 하겠다. $\mathbf A$가 엄밀하면
(\hyperref[0.13.4.1.3-fr]{13.4.1.3})으로부터 $p\leq k\leq h$일 때
표준 사상 $\operatorname{gr}^p(A_h)\to\operatorname{gr}^p(A_k)$가
\emph{동형}임이 따른다. 다시 말해 사영계
$(\operatorname{gr}^p(A_k))_{k\in\mathbb Z}$는 본질적으로 상수이다.
대상들 $\operatorname{gr}^p(A_p)$의 열(각 $p$에 대하여
$\varprojlim_k\operatorname{gr}^p(A_k)$와 동일시한다)을 이때
$\operatorname{gr}^\bullet(\mathbf A)$로 쓰고, 이를 엄밀 사영계
$\mathbf A=(A_k)$의 \emph{연관 등급 대상}이라 한다.

이제 (아래로 유계인) 사영계 $\mathbf A$가 (ML)을 만족한다고
가정하자. 그러면 ``보편상'' 대상들의 사영계
$\mathbf A'=(A'_k)$가 \emph{엄밀}함을 알고 있으며
(\hyperref[0.13.1.2-fr]{13.1.2}), 이 사영계는 또한 아래로 유계이다~;
$\mathbf A'$의 연관 등급 대상 $\operatorname{gr}^\bullet(\mathbf A')$를
이때에도 $\mathbf A$에 \emph{연관된} 등급 대상이라 하고
$\operatorname{gr}^\bullet(\mathbf A)$로 쓴다.
\end{env}

\begin{proposition}[13.4.3]
\label{0.13.4.3-fr}
\emph{$\mathbf A=(A_k)_{k\in\mathbb Z}$를 아벨 범주 $C$의 아래로
유계인 사영계라 하자. 다음 두 조건은 동치이다.}
\begin{enumerate}
\item[\emph{a)}] \emph{$\mathbf A$는 조건 (ML)을 만족한다.}
\item[\emph{b)}] \emph{모든 $p\in\mathbb Z$에 대하여 사영계
$(\operatorname{gr}^p(A_k))_{k\in\mathbb Z}$는 본질적으로 상수이다.}
\end{enumerate}
\emph{더 나아가 이 조건들이 성립하면 모든 $p\in\mathbb Z$에 대하여
표준 동형}
\begin{equation}
\label{0.13.4.3.1-fr}
\tag{13.4.3.1}
\operatorname{gr}^p(\mathbf A)\xrightarrow{\sim}
\varprojlim_k\operatorname{gr}^p(A_k)
\end{equation}
\emph{이 존재한다.}
\end{proposition}

\begin{proof}
(\hyperref[0.13.4.1.2-fr]{13.4.1.2})으로부터 \emph{a)}가 \emph{b)}를
함의함이 곧바로 따른다~; 같은 공식을 사영계 $\mathbf A'$에 적용하면
((\hyperref[0.13.4.2-fr]{13.4.2})의 표기) 정의에 의해 동형
(\hyperref[0.13.4.3.1-fr]{13.4.3.1})을 얻는다. $k\leq h$일 때
$A_{kh}=\operatorname{Im}(A_h\to A_k)$로 놓자~; $k\leq h\leq j$이면
$A_{kj}\subset A_{kh}\subset A_k$이다. $A_{kh}$에
$(\mathrm F^p(A_k))$가 유도하는 여과를 주자~; $\mathbf A$를 정의하는
사상들의 추이성에 의해, 이 여과가 $(\mathrm F^p(A_h))$의 몫 여과이기도
함을 곧바로 확인할 수 있다~; 따라서
\begin{equation}
\label{0.13.4.3.2-fr}
\tag{13.4.3.2}
\operatorname{gr}^p(A_{kh})=
\operatorname{Im}\bigl(\operatorname{gr}^p(A_h)\longrightarrow
\operatorname{gr}^p(A_k)\bigr).
\end{equation}
\oldpage[0\textsubscript{III}]{71}
이제 \emph{b)}가 성립한다고 가정하자~; 모든 $p\in\mathbb Z$와 모든
$k\geq p$에 대하여, (\hyperref[0.13.4.3.2-fr]{13.4.3.2})의 우변이
$h\geq\mathrm L(p,k)$일 때 일정해지는 정수 $\mathrm L(p,k)$가
존재한다~; $p<k_0$이면 $\operatorname{gr}^p(A_k)=0$이므로, 주어진
$k$에 대하여 $p$가 $k$ 이하인 정수들의 집합을 움직일 때 0이 아닌
$\mathrm L(p,k)$는 유한 개뿐이다. 이 정수들의 최댓값을
$\mathrm L(k)=m$이라 하자~; 모든 $h\geq m$에 대하여
$A_{kh}\subset A_{km}$이고, $m$의 정의에 의해 표준 단사
$A_{kh}\to A_{km}$는 \emph{동형}
$\operatorname{gr}^\bullet(A_{kh})\xrightarrow{\sim}
\operatorname{gr}^\bullet(A_{km})$를 정의한다~; 여과들이 \emph{유한}이므로
앞의 단사 자체가 전단사임을 얻는다 (Bourbaki, \emph{Alg. comm.},
chap.~III, §~2, n\textsuperscript{o}~8, th.~1). 이는 $\mathbf A$가
(ML)을 만족함을 증명한다.
\end{proof}

\begin{env}[13.4.4]
\label{0.13.4.4-fr}
$C$에서 사영극한 $A=\varprojlim_k A_k$가 존재한다고 가정하자.
(\hyperref[0.13.4.1-fr]{13.4.1})의 정의에서 $A_k$를 $A$로 바꿀 수
있으며, 이처럼 $A$ 위에 정의한 여과는 여전히
\begin{equation}
\label{0.13.4.4.1-fr}
\tag{13.4.4.1}
\operatorname{gr}^p(A)=\operatorname{gr}^p(A_p)
\cap\operatorname{Im}(A\longrightarrow A_p)
\end{equation}
을 만족한다.
\end{env}

\begin{corollary}[13.4.5]
\label{0.13.4.5-fr}
\emph{$C$가 아벨군의 범주라고 가정하자. 사영계 $\mathbf A$가 (ML)을
만족하고 $A=\varprojlim_k A_k$이면, 모든 $p\in\mathbb Z$에 대하여
표준 동형}
\begin{equation}
\label{0.13.4.5.1-fr}
\tag{13.4.5.1}
\operatorname{gr}^p(A)\xrightarrow{\sim}
\varprojlim_k\operatorname{gr}^p(A_k)
\end{equation}
\emph{이 존재한다.}
\end{corollary}
\begin{proof}
실제로 $k$가 충분히 크면
$\operatorname{Im}(A_k\to A_p)=\operatorname{Im}(A\to A_p)$이다
(Bourbaki, \emph{Top. gén.}, chap.~II,
3\textsuperscript{e} éd., §~3, n\textsuperscript{o}~5, th.~1). 따라서
결론은 (\hyperref[0.13.4.1.3-fr]{13.4.1.3})과
(\hyperref[0.13.4.4.1-fr]{13.4.4.1})에서 따른다.
\end{proof}

\subsection{여과 복합체의 스펙트럼 열들의 사영극한}
\label{subsection:0.13.5-fr}
\begin{env}[13.5.1]
\label{0.13.5.1-fr}
$C$를 아벨 범주라 하고, $X^\bullet$을 $C$의 대상들로 이루어진
복합체라 하자. $X^\bullet$에는 어떤 지표 $p_0$에 대하여
$\mathrm F^{p_0}(X^\bullet)=X^\bullet$을 만족하는 \emph{여과}
$(\mathrm F^p(X^\bullet))_{p\in\mathbb Z}$가 주어져 있다고 하자.
각 $k\in\mathbb Z$에 대하여 복합체
$X_k^\bullet=X^\bullet/\mathrm F^{k+1}(X^\bullet)$을 생각하자~; 이
복합체에는 $p\leq k$일 때의
$\mathrm F^p(X_k^\bullet)=
\mathrm F^p(X^\bullet)/\mathrm F^{k+1}(X^\bullet)$과 $p\geq k+1$일 때의
$\mathrm F^p(X_k^\bullet)=0$으로 이루어진 여과가 표준적으로 주어진다.
더 나아가 표준 사상 $X_{k+1}^\bullet\to X_k^\bullet$들이 있으며,
이들은 $\mathbf X^\bullet=(X_k^\bullet)_{k\in\mathbb Z}$를
\emph{$C$의 대상들의 여과 복합체로 이루어진 사영계}로 만든다. 이
사영계는 \emph{엄밀}하고 $k<p_0$이면 $X_k^\bullet=0$임에 유의하자.
\end{env}
\begin{env}[13.5.2]
\label{0.13.5.2-fr}
더 일반적으로, $C$의 대상들의 복합체로 이루어진 \emph{엄밀}하고
\emph{아래로 유계인} 사영계
$\mathbf X^\bullet=(X_k^\bullet)_{k\in\mathbb Z}$를 생각하자~; 각
$X_k^\bullet$ 위에는 (아래로 유계인 $C$의 복합체들의 아벨 범주에서)
(\hyperref[0.13.4.1-fr]{13.4.1})에 정의한 여과를 생각하자. 사상
$X_k^\bullet\to X_p^\bullet$ $(p\leq k)$는 유한 여과를 갖춘
\emph{여과} 복합체들의 사상이 된다. 여과 복합체의 스펙트럼 열이
갖는 함자성 (\hyperref[0.11.2.3-fr]{11.2.3})에 의해, 사영계
$\mathbf X^\bullet$을 정의하는 사상들은
$\mathrm E(\mathbf X^\bullet)=(\mathrm E(X_k^\bullet))$를
\emph{스펙트럼 열들의 사영계}로 만드는 사상들을 준다.
\end{env}
\begin{lemma}[13.5.3]
\label{0.13.5.3-fr}
\emph{여과 복합체의 사영계
$\mathbf X^\bullet=(X_k^\bullet)_{k\in\mathbb Z}$를
(\hyperref[0.13.5.2-fr]{13.5.2})에서처럼 얻었다고 가정하자. 그러면:}
\begin{enumerate}
\item[\emph{a)}]\emph{$r\geq p-p_0$이면 모든 $k\in\mathbb Z$에 대하여
$\mathrm B_r^{pq}(X_k^\bullet)=\mathrm B_\infty^{pq}(X_k^\bullet)$이다.}
\item[\emph{b)}]\emph{$k+1\leq p+r$이면
$\mathrm Z_r^{pq}(X_k^\bullet)=\mathrm Z_\infty^{pq}(X_k^\bullet)$이다.}
\item[\emph{c)}]\emph{$k+1\geq p+r$이면 모든 $h\geq k$에 대하여 사상
$\mathrm Z_r^{pq}(X_h^\bullet)\to\mathrm Z_r^{pq}(X_k^\bullet)$과
$\mathrm B_r^{pq}(X_h^\bullet)\to\mathrm B_r^{pq}(X_k^\bullet)$은
동형이다.}
\end{enumerate}
\end{lemma}
\oldpage[0\textsubscript{III}]{72}
이 세 성질은 $p-r<p_0$이면
$\mathrm F^{p-r+1}(X_k^\bullet)=X_k^\bullet$이라는 사실을 고려하면
(\hyperref[0.11.2.2-fr]{11.2.2})의 정의들에서 곧바로 따른다.

\begin{env}[13.5.4]
\label{0.13.5.4-fr}
(\hyperref[0.13.5.3-fr]{13.5.3})의 가정들이 성립한다고 하자. 그러면
고정된 $p,q,r$에 대하여($r$은 유한), 사영계들
\[
(\mathrm Z_r^{pq}(X_k^\bullet))_{k\in\mathbb Z},\quad
(\mathrm B_r^{pq}(X_k^\bullet))_{k\in\mathbb Z},\quad
(\mathrm E_r^{pq}(X_k^\bullet))_{k\in\mathbb Z}
\]
은 \emph{본질적으로 상수}이다~; 각각의 사영극한을
$\mathrm Z_r^{pq}(\mathbf X^\bullet)$,
$\mathrm B_r^{pq}(\mathbf X^\bullet)$ 및
$\mathrm E_r^{pq}(\mathbf X^\bullet)=
\mathrm Z_r^{pq}(\mathbf X^\bullet)/
\mathrm B_r^{pq}(\mathbf X^\bullet)$로 쓰겠다.
$\mathrm Z_r^{pq}(\mathbf X^\bullet)$와
$\mathrm B_r^{pq}(\mathbf X^\bullet)$는 표준적으로
$\mathrm E_1^{pq}(\mathbf X^\bullet)$의 부분대상과 동일시된다.
$d_r^{pq}$의 정의 (M, XV, 1)에 따르면 ($X_k^\bullet$들에 관한) 이
사상들 역시 본질적으로 상수이며, 따라서 사상
\begin{equation}
\label{0.13.5.4.1-fr}
\tag{13.5.4.1}
d_r^{pq}:\mathrm E_r^{pq}(\mathbf X^\bullet)
\longrightarrow\mathrm E_r^{p+r,q-r+1}(\mathbf X^\bullet)
\end{equation}
을 정의한다. 이들은 $d_r^{p+r,q-r+1}\circ d_r^{pq}=0$을 만족한다~;
더 나아가 $\operatorname{Ker}(d_r^{pq})$에서
$\mathrm Z_{r+1}^{pq}(\mathbf X^\bullet)/
\mathrm B_r^{pq}(\mathbf X^\bullet)$로 가는 표준 동형과
$\operatorname{Im}(d_r^{pq})$에서
$\mathrm B_{r+1}^{p+r,q-r+1}(\mathbf X^\bullet)/
\mathrm B_r^{p+r,q-r+1}(\mathbf X^\bullet)$로 가는 표준 동형이 있다.
\end{env}

\begin{lemma}[13.5.5]
\label{0.13.5.5-fr}
\emph{(\hyperref[0.13.5.3-fr]{13.5.3})의 가정 아래에서,
$s\geq r>p-p_0$일 때 표준 모노사상}
\begin{equation}
\label{0.13.5.5.1-fr}
\tag{13.5.5.1}
i:\mathrm E_s^{pq}(\mathbf X^\bullet)
\longrightarrow\mathrm E_r^{pq}(\mathbf X^\bullet)
\end{equation}
\emph{과 표준 동형}
\begin{equation}
\label{0.13.5.5.2-fr}
\tag{13.5.5.2}
j_r:\mathrm E_r^{pq}(\mathbf X^\bullet)
\xrightarrow{\sim}\mathrm E_\infty^{pq}(X_{p+r-1}^\bullet)
\end{equation}
\emph{이 존재하여 도식}
\begin{equation}
\label{0.13.5.5.3-fr}
\tag{13.5.5.3}
\xymatrix{
\mathrm E_s^{pq}(\mathbf X^\bullet)\ar[r]^{j_s}\ar[d]_i
&\mathrm E_\infty^{pq}(X_{p+s-1}^\bullet)\ar[d]\\
\mathrm E_r^{pq}(\mathbf X^\bullet)\ar[r]^{j_r}
&\mathrm E_\infty^{pq}(X_{p+r-1}^\bullet)
}
\end{equation}
\emph{은 가환이다} (오른쪽 수직 화살표는 사상
$X_{p+s-1}^\bullet\to X_{p+r-1}^\bullet$에서 유도된다).
\end{lemma}
\begin{proof}
$r>p-p_0$일 때
$\mathrm B_r^{pq}(X_k^\bullet)=\mathrm B_\infty^{pq}(X_k^\bullet)$이므로
(\hyperref[0.13.5.3-fr]{13.5.3}, a), $i$가 존재한다~;
$k+1\leq p+r$이면
$\mathrm Z_r^{pq}(X_k^\bullet)=\mathrm Z_\infty^{pq}(X_k^\bullet)$이므로
(\hyperref[0.13.5.3-fr]{13.5.3}, b), 특히
$\mathrm Z_r^{pq}(X_{p+r-1}^\bullet)=
\mathrm Z_\infty^{pq}(X_{p+r-1}^\bullet)$이다. 다른 한편
$\mathrm Z_r^{pq}(X_{p+r-1}^\bullet)$와
$\mathrm B_r^{pq}(X_{p+r-1}^\bullet)$는
(\hyperref[0.13.5.3-fr]{13.5.3}, c))에 의해 각각
$\mathrm Z_r^{pq}(\mathbf X^\bullet)$와
$\mathrm B_r^{pq}(\mathbf X^\bullet)$에 표준적으로 동일시된다.
이로부터 $j_r$의 존재와
(\hyperref[0.13.5.5.3-fr]{13.5.5.3})의 가환성이 따른다.
\end{proof}

\begin{corollary}[13.5.6]
\label{0.13.5.6-fr}
\emph{(\hyperref[0.13.5.3-fr]{13.5.3})의 가정 아래에서 사영극한
$\varprojlim_r\mathrm E_r^{pq}(\mathbf X^\bullet)$와
$\varprojlim_k\mathrm E_\infty^{pq}(X_k^\bullet)$ 가운데 하나가
존재하면 다른 하나도 존재하며, 표준 동형}
\begin{equation}
\label{0.13.5.6.1-fr}
\tag{13.5.6.1}
j_\infty:\varprojlim_r\mathrm E_r^{pq}(\mathbf X^\bullet)
\xrightarrow{\sim}
\varprojlim_k\mathrm E_\infty^{pq}(X_k^\bullet)
\end{equation}
\emph{이 존재한다. 더 나아가 사영계
$(\mathrm E_r^{pq}(\mathbf X^\bullet))_{r\in\mathbb Z}$가 본질적으로
상수이기 위한 필요충분조건
(\hyperref[0.13.4.2-fr]{13.4.2})은 사영계
$(\mathrm E_\infty^{pq}(X_k^\bullet))_{k\in\mathbb Z}$가
본질적으로 상수인 것이다.}
\end{corollary}

\begin{env}[13.5.7]
\label{0.13.5.7-fr}
$\mathrm E_1^{pq}(\mathbf X^\bullet)$의 부분대상 가운데 각각
$r>p-p_0$일 때의 $\mathrm B_r^{pq}(\mathbf X^\bullet)$ 및
$\inf_r\mathrm Z_r^{pq}(\mathbf X^\bullet)$와 같은 것을
($\inf_r$가 존재할 때)
$\mathrm B_\infty^{pq}(\mathbf X^\bullet)$와
$\mathrm Z_\infty^{pq}(\mathbf X^\bullet)$로 쓴다. 그러면
$\varprojlim_r\mathrm E_r^{pq}(\mathbf X^\bullet)$는 표준적으로
\oldpage[0\textsubscript{III}]{73}
$\mathrm E_\infty^{pq}(\mathbf X^\bullet)=
\mathrm Z_\infty^{pq}(\mathbf X^\bullet)/
\mathrm B_\infty^{pq}(\mathbf X^\bullet)$와 동일시된다. 대상들
$\mathrm Z_r^{pq}(\mathbf X^\bullet)$,
$\mathrm B_r^{pq}(\mathbf X^\bullet)$,
$\mathrm E_r^{pq}(\mathbf X^\bullet)$ $(1\leq r\leq+\infty)$와
$d_r^{pq}$는 (\hyperref[0.13.5.5-fr]{13.5.5})의 제한을 받는 사영계
$\mathbf X^\bullet$에 \emph{함자적으로} 의존하고,
(\hyperref[0.13.5.5-fr]{13.5.5})와
(\hyperref[0.13.5.6-fr]{13.5.6})에서 정의한 사상들도 함자적임에
유의하자.
\end{env}

\subsection{유한 여과를 갖춘 대상에 관한 함자의 스펙트럼 열}
\label{subsection:0.13.6-fr}

\begin{env}[13.6.1]
\label{0.13.6.1-fr}
$C$, $C'$을 두 아벨 범주, $\mathrm T:C\to C'$을 가법 공변 함자라
하자. $C$의 모든 대상이 어떤 단사 대상의 부분대상과 동형이라고
가정한다. 따라서 오른쪽 유도함자 $\mathrm R^p\mathrm T$ ($p\geq0$)가
존재한다.
\end{env}

\begin{lemma}[13.6.2]
\label{0.13.6.2-fr}
\emph{$A$를 유한 여과 $(\mathrm F^i(A))_{i\in\mathbb Z}$를 갖춘
$C$의 대상이라 하자. $A$의 단사 분해
$X^\bullet=(X^j)_{j\geq0}$와 그 위의 유한 여과
$(\mathrm F^i(X^\bullet))_{i\in\mathbb Z}$가 존재하여, 관계
$\mathrm F^i(A)=A$ (각각 $\mathrm F^i(A)=0$)는
$\mathrm F^i(X^\bullet)=X^\bullet$ (각각
$\mathrm F^i(X^\bullet)=0$)를 함의하고, 모든 $i\in\mathbb Z$에
대하여 $\mathrm F^i(X^\bullet)$은 $\mathrm F^i(A)$의 단사 분해이다.}
\end{lemma}
\begin{proof}
$p$를 $\mathrm F^p(A)=A$인 가장 큰 지표, $q>p$를
$\mathrm F^q(A)=0$인 가장 작은 지표라 하자. $q-p$에 관한 귀납법을
쓴다. $q-p=1$일 때 보조정리는 명백하다. 원하는 성질들을 갖는
$A/\mathrm F^{q-1}(A)$의 단사 분해 ${X''}^\bullet$을 구성했다고 하자.
완전열
\[
0\to\mathrm F^{q-1}(A)\to A\to A/\mathrm F^{q-1}(A)\to0
\]
을 생각하고, $\mathrm F^{q-1}(A)$의 단사 분해 ${X'}^\bullet$을
취한다. 이어서 앞의 완전열과 양립하는 완전열
\[
0\to {X'}^\bullet\to X^\bullet\to {X''}^\bullet\to0
\]
을 갖도록 $A$의 단사 분해 $X^\bullet$을 정한다 (M, V, 2.2).
$X^\bullet$이 요구를 만족함은 명백하다.
\end{proof}

\begin{corollary}[13.6.3]
\label{0.13.6.3-fr}
\emph{$B$를 유한 여과 $(\mathrm F^i(B))_{i\in\mathbb Z}$를 갖춘
$C$의 또 다른 대상, $s$를 정수라 하자. $u:A\to B$가 모든
$i\in\mathbb Z$에 대하여
$u(\mathrm F^i(A))\subset\mathrm F^{i+s}(B)$를 만족한다고 하자.
$Y^\bullet=(Y^j)_{j\geq0}$가
(\hyperref[0.13.6.2-fr]{13.6.2})에 열거된 성질들을 갖는 여과
$(\mathrm F^i(Y^\bullet))_{i\in\mathbb Z}$를 갖춘 $B$의 단사
분해이면, $u$와 양립하고 모든 $i\in\mathbb Z$에 대하여
$v(\mathrm F^i(X^\bullet))\subset\mathrm F^{i+s}(Y^\bullet)$를
만족하는 사상 $v:X^\bullet\to Y^\bullet$가 존재한다. 또한 그러한
두 사상 $v$, $v'$는 서로 호모토픽하다.}
\end{corollary}

이는 앞의 구성과 (M, V, 2.3)을 $q-p$에 관하여 귀납적으로 적용하면
곧바로 따른다.

\begin{env}[13.6.4]
\label{0.13.6.4-fr}
(\hyperref[0.13.6.2-fr]{13.6.2})의 가정 아래 이제 $C'$의 복합체
$\mathrm T(X^\bullet)$을 생각하자. $\mathrm F^i(X^\bullet)$이
$X^\bullet$의 직합인자이므로, 이 복합체는 명백히 복합체들
$\mathrm T(\mathrm F^i(X^\bullet))$에 의해 여과된다.
(\hyperref[0.13.6.3-fr]{13.6.3})에 따라 이 여과 복합체의
\emph{스펙트럼 열}은 동형을 제외하면 여과 대상 $A$에만 의존한다.
그 수렴대상은 코호몰로지 $\mathrm R^\bullet\mathrm T(A)$이고, 그 위의
여과는
\begin{equation}
\label{0.13.6.4.1-fr}
\tag{13.6.4.1}
\begin{split}
\mathrm F^p(\mathrm R^n\mathrm T(A))
&=\operatorname{Im}\bigl(\mathrm R^n\mathrm T(\mathrm F^p(A))
\longrightarrow\mathrm R^n\mathrm T(A)\bigr)\\
&=\operatorname{Ker}\bigl(\mathrm R^n\mathrm T(A)
\longrightarrow\mathrm R^n\mathrm T(A/\mathrm F^p(A))\bigr)
\end{split}
\end{equation}
이다 (\hyperref[0.11.2.2-fr]{11.2.2}). 그 $\mathrm E_1$항은
\begin{equation}
\label{0.13.6.4.2-fr}
\tag{13.6.4.2}
\mathrm E_1^{pq}=\mathrm R^{p+q}\mathrm T(\operatorname{gr}^p(A))
\end{equation}
로 주어진다. 여기서 통상과 같이
$\operatorname{gr}^p(A)=\mathrm F^p(A)/\mathrm F^{p+1}(A)$이다.
(\hyperref[0.11.2.2-fr]{11.2.2})에 따라 수렴대상의 여과는 유한이고,
주어진 $p$, $q$에 대하여 열
\oldpage[0\textsubscript{III}]{74}
$\mathrm B_r^{pq}(A)=\mathrm B_r^{pq}(\mathrm T(X^\bullet))$와
$\mathrm Z_r^{pq}(A)=\mathrm Z_r^{pq}(\mathrm T(X^\bullet))$는
정상적이다. 따라서 앞의 스펙트럼 열은 \emph{쌍정칙}이다
(\hyperref[0.11.1.3-fr]{11.1.3}). 이 스펙트럼 열을
$\mathrm E(A)=(\mathrm E_r^{pq}(A))$로 나타내고,
\emph{여과 대상 $A$에 관한 함자 $\mathrm T$의 스펙트럼 열}이라 한다.
\end{env}

\begin{env}[13.6.5]
\label{0.13.6.5-fr}
이제 (\hyperref[0.13.6.3-fr]{13.6.3})의 가정들이 성립한다고 하고 그
표기를 유지하자. $\mathrm F^i(X^\bullet)$ (각각
$\mathrm F^i(Y^\bullet)$)은 $X^\bullet$ (각각 $Y^\bullet$)의
직합인자이므로, 모든 $i\in\mathbb Z$에 대하여
\[
(\mathrm T v)(\mathrm T(\mathrm F^i(X^\bullet)))
\subset\mathrm T(\mathrm F^{i+s}(Y^\bullet)).
\]
그러면 (\hyperref[0.11.2.2-fr]{11.2.2})의 정의에 따라
$1\leq r\leq+\infty$일 때 $\mathrm T v$는 사상
$\mathrm B_r^{pq}(\mathrm T(X^\bullet))\to
\mathrm B_r^{p+s,q-s}(\mathrm T(Y^\bullet))$와 사상
$\mathrm Z_r^{pq}(\mathrm T(X^\bullet))\to
\mathrm Z_r^{p+s,q-s}(\mathrm T(Y^\bullet))$를 정의하며, 따라서 사상
\[
w_r:\mathrm E_r^{pq}(A)\to\mathrm E_r^{p+s,q-s}(B)
\]
을 정의한다. 마찬가지로 수렴대상 위에는
$u_n:\mathrm R^n\mathrm T(A)\to\mathrm R^n\mathrm T(B)$가 있고,
$u_n(\mathrm F^p(\mathrm R^n\mathrm T(A)))\subset
\mathrm F^{p+s}(\mathrm R^n\mathrm T(B))$이다.

$d_r^{pq}$의 정의 (M, XV, 1)에 따라 도식들
\[
\xymatrix{
\mathrm E_r^{pq}(A)\ar[r]^{d_r^{pq}}\ar[d]_{w_r}
&\mathrm E_r^{p+r,q-r+1}(A)\ar[d]^{w_r}\\
\mathrm E_r^{p+s,q-s}(B)\ar[r]_{d_r^{p+s,q-s}}
&\mathrm E_r^{p+r+s,q-r-s+1}(B)
}
\]
도 가환한다. 여기서 동형 $\alpha_r^{pq}$에 대한 비슷한 가환도식을
얻으며, 그 명시는 독자에게 맡긴다. 끝으로 (\emph{같은 곳}), 수렴대상에
대해서도 가환도식
\[
\xymatrix{
\mathrm E_\infty^{pq}(A)\ar[r]^{\beta^{pq}}\ar[d]_{w_\infty}
&\operatorname{gr}^p(\mathrm R^{p+q}\mathrm T(A))
  \ar[d]^{u_{p+q}}\\
\mathrm E_\infty^{p+s,q-s}(B)\ar[r]_{\beta^{p+s,q-s}}
&\operatorname{gr}^{p+s}(\mathrm R^{p+q}\mathrm T(B))
}
\]
을 얻는다.
\end{env}

\begin{env}[13.6.6]
\label{0.13.6.6-fr}
특히 환 $S$가 여과 $(\mathrm F^i(S))_{i\in\mathbb Z}$를 갖고,
환 준동형
\begin{equation}
\label{0.13.6.6.1-fr}
\tag{13.6.6.1}
h:S\to\operatorname{Hom}_C(A,A)
\end{equation}
이 존재하여,
\oldpage[0\textsubscript{III}]{75}
모든 $t\in\mathrm F^i(S)$와 모든 두 지표 $i$, $j$에 대하여
$h_t(\mathrm F^j(A))\subset\mathrm F^{i+j}(A)$라고 하자. 줄여서 이때
$A$는 여과환 $S$ 위의 \emph{여과 $S$-$C$-가군 구조}를 갖는다고 한다.
연관 등급 대상으로 이행하면, 모든 $t\in\mathrm F^j(S)$에 대하여
$h_t$는 $\operatorname{gr}^\bullet(A)$의 차수 $j$인 동차 등급
자기준동형 $\overline h_t$를 정의한다. 이 사상은 $t$의
$\operatorname{gr}^j(S)$에서의 류에만 의존하므로, 등급환 준동형
\[
\overline h:\operatorname{gr}^\bullet(S)\to
\operatorname{Hom}_C^g(\operatorname{gr}^\bullet(A),
\operatorname{gr}^\bullet(A))
\]
을 정의한다. 여기서 우변은 $\operatorname{gr}^\bullet(A)$의 등급
자기준동형들의 환이다. 이때 $\operatorname{gr}^\bullet(A)$가
\emph{등급 $\operatorname{gr}^\bullet(S)$-$C$-가군 구조}를 갖는다고
한다. 이제 (\hyperref[0.13.6.5-fr]{13.6.5})에 따라
$1\leq r\leq+\infty$일 때 모든
$\overline t\in\operatorname{gr}^j(S)$는 이중등급 대상들
$(\mathrm B_r^{pq}(A))_{(p,q)\in\mathbb Z\times\mathbb Z}$,
$(\mathrm Z_r^{pq}(A))_{(p,q)\in\mathbb Z\times\mathbb Z}$,
$\mathrm E_r(A)=(\mathrm E_r^{pq}(A))_{(p,q)\in\mathbb Z\times\mathbb Z}$
위에 이중차수 $(j,-j)$인 이중등급 자기준동형들을 표준적으로 정의한다.
$r$이 유한할 때 $\mathrm E_r(A)$ 위의 이 자기준동형은
$d_r^{pq}$들이 정의하는 이중등급 자기준동형과 가환한다. 이 자기준동형들은
$\operatorname{gr}^\bullet(S)$ 및 고려하는 이중등급 대상들의 덧셈에 대한
통상적인 결합법칙과 분배법칙을 만족하므로, 줄여서 이 대상들을
\emph{이중등급 $\operatorname{gr}^\bullet(S)$-$C'$-가군}이라 한다.
$\alpha_r^{pq}$들이 이 구조에 관한 동형을 정의함은 즉시 알 수 있다.
모든 정수 $n$에 대하여 $\mathrm B_r^{(n)}(A)$ (각각
$\mathrm Z_r^{(n)}(A)$, $\mathrm E_r^{(n)}(A)$)를
$p+q=n$인 $\mathrm B_r^{pq}(A)$ (각각 $\mathrm Z_r^{pq}(A)$,
$\mathrm E_r^{pq}(A)$)들로 이루어진 $\mathrm B_r^\bullet(A)$ (각각
$\mathrm Z_r^\bullet(A)$, $\mathrm E_r^\bullet(A)$)의 등급 부분대상이라
하자 ($1\leq r\leq+\infty$). 이들은 명백히
\emph{등급 $\operatorname{gr}^\bullet(S)$-$C'$-가군}이다. 끝으로 모든
$\overline t\in\operatorname{gr}^j(S)$는 모든 $n$에 대하여 등급 대상
$\operatorname{gr}^\bullet(\mathrm R^n\mathrm T(A))$ 위에 차수 $j$인
등급 자기준동형을 정의한다. 따라서 이 대상도 등급
$\operatorname{gr}^\bullet(S)$-$C'$-가군 구조를 가지며, $p+q=n$일 때
$\beta^{pq}$들은 이 구조에 관한 동형
\[
\mathrm E^{(n)}(A)\xrightarrow{\sim}
\operatorname{gr}^\bullet(\mathrm R^n\mathrm T(A))
\]
을 정의한다.

$C'$이 아벨군의 범주일 때 $S$-$C'$-가군 구조(각각 등급 또는 이중등급
$\operatorname{gr}^\bullet(S)$-$C'$-가군 구조)는 통상적인 $S$-가군
구조(각각 등급 또는 이중등급 $\operatorname{gr}^\bullet(S)$-가군
구조)에 지나지 않음에 유의하자.
\end{env}

\subsection{여러 변수의 사영극한의 유도 함자}
\label{subsection:0.13.7-fr}

\begin{env}[13.7.1]
\label{0.13.7.1-fr}
$C$, $C'$을 두 아벨 범주라 하고, $C$의 모든 대상은 어떤 단사 대상의
부분대상이라고 가정하자. $\mathrm T:C\to C'$을 가법 공변 함자라 하자.
$C$ 안의 \emph{아래로 유계인} 엄밀 사영계
$\mathbf A=(A_k)_{k\in\mathbb Z}$를 생각하자. 구체적으로 어떤 $k_0$가
존재하여 $k<k_0$일 때 $A_k=0$이라고 가정한다. 공식
(\hyperref[0.13.4.1.1-fr]{13.4.1.1})에 따라 각 $A_k$ 위에 여과
$(\mathrm F^p(A_k))_{p\in\mathbb Z}$를 이 사영계에 표준적으로
연관시킨다. 사영계가 엄밀하므로 $i\leq j\leq k+1$ 및 $h\geq k$일 때
표준 사상
\begin{equation}
\label{0.13.7.1.1-fr}
\tag{13.7.1.1}
\mathrm F^i(A_h)/\mathrm F^j(A_h)\to
\mathrm F^i(A_k)/\mathrm F^j(A_k)
\end{equation}
은 동형이다. 또한 모든 $k$에 대하여
$\mathrm F^{k_0}(A_k)=A_k$이고 $\mathrm F^{k+1}(A_k)=0$임을 상기하자.
\end{env}

\begin{env}[13.7.2]
\label{0.13.7.2-fr}
이제 각 $k$에 대하여 (\hyperref[0.13.6.2-fr]{13.6.2})의 성질들을 갖는
$A_k$의 단사 분해 $X_k^\bullet=(X_k^j)_{j\geq0}$을 구성하자. 표준
사상 $A_{k+1}\to A_k$를 사용하면
(\hyperref[0.13.6.3-fr]{13.6.3})에 따라 각 $k$에 대하여 여과들과
양립하는 복합체의 사상 $X_{k+1}^\bullet\to X_k^\bullet$을 정의할 수
있다. 이 사상들은
$\mathbf X^\bullet=(X_k^\bullet)_{k\in\mathbb Z}$를 복합체들의
사영계로 만든다.
\oldpage[0\textsubscript{III}]{76}
더 나아가 이 사영계가 엄밀하다고 가정할 수 있다. 실제로 동형
(\hyperref[0.13.7.1.1-fr]{13.7.1.1})에 따라 $A_k$는
$A_{k+1}/\mathrm F^{k+1}(A_{k+1})$와 동형이다. 따라서
$X_{k+1}^\bullet$을 구성할 때
$A_{k+1}/\mathrm F^{k+1}(A_{k+1})$의 단사 분해를
$X_k^\bullet$로 취할 수 있다. 그러면 (M, V, 2.3)에 따라 복합체의
사상 $X_{k+1}^\bullet\to X_k^\bullet$을 구성할 때 몫으로 이행하여
여과들을 보존하는 동형
\[
X_{k+1}^\bullet/\mathrm F^{k+1}(X_{k+1}^\bullet)
\xrightarrow{\sim}X_k^\bullet
\]
을 주도록 할 수 있다. 이것이
(\hyperref[0.13.5.1-fr]{13.5.1})의 조건이다.
\end{env}

\begin{env}[13.7.3]
\label{0.13.7.3-fr}
구성에 따라 복합체 $\mathrm T(X_k^\bullet)$들의 사영계
$\mathrm T(\mathbf X^\bullet)$은
(\hyperref[0.13.5.3-fr]{13.5.3})의 가정들을 만족한다. 따라서
(\hyperref[0.13.5.4-fr]{13.5.4}),
(\hyperref[0.13.5.5-fr]{13.5.5}),
(\hyperref[0.13.5.6-fr]{13.5.6})의 결과들을 스펙트럼 열
$\mathrm E(\mathrm T(X_k^\bullet))=\mathrm E(A_k)$에 적용할 수 있다.
$1\leq r\leq+\infty$일 때
$\mathrm E_r^{pq}(\mathrm T(\mathbf X^\bullet))$ 대신
$\mathrm E_r^{pq}(\mathbf A)$라 쓰고($r=+\infty$인 경우에는
(\hyperref[0.13.5.7-fr]{13.5.7}) 참조), 비슷한 표기들에도 같은 약속을
쓴다. 특히 (\hyperref[0.13.6.4.2-fr]{13.6.4.2})와 사영계
$(\operatorname{gr}^p(A_k))$가 본질적으로 상수라는 사실에 따라
\begin{equation}
\label{0.13.7.3.1-fr}
\tag{13.7.3.1}
\mathrm E_1^{pq}(\mathbf A)=
\mathrm R^{p+q}\mathrm T(\operatorname{gr}^p(\mathbf A))
\end{equation}
이다.

이 결과들과 (\hyperref[0.13.4.3-fr]{13.4.3})으로부터 다음 명제를
얻는다. 이 명제는 Shih Weishu가 다른 미출판 방법으로 처음 증명하였다.
\end{env}

\begin{proposition}[13.7.4]
\label{0.13.7.4-fr}
\emph{(Shih). --- $n$을 정수라 하자. 다음 두 조건은 동치이다.}
\begin{enumerate}
\item[\emph{a)}] \emph{$p+q=n$인 모든 쌍 $(p,q)$에 대하여 사영계
$(\mathrm E_r^{pq}(\mathbf A))_{r\geq2}$는 본질적으로 상수이다.}
\item[\emph{b)}] \emph{사영계
$\mathrm R^n\mathrm T(\mathbf A)=
(\mathrm R^n\mathrm T(A_k))_{k\in\mathbb Z}$는 (ML)을 만족한다.}
\end{enumerate}
\emph{또한 이 조건들이 성립하면 모든 $p\in\mathbb Z$에 대하여 표준
동형}
\begin{equation}
\label{0.13.7.4.1-fr}
\tag{13.7.4.1}
\operatorname{gr}^p(\mathrm R^n\mathrm T(\mathbf A))
\xrightarrow{\sim}\mathrm E_\infty^{p,n-p}(\mathbf A)
\qquad\emph{모든 $p\in\mathbb Z$에 대하여}
\end{equation}
\emph{가 성립한다.}
\end{proposition}
\begin{proof}
(\hyperref[0.13.5.6-fr]{13.5.6})에 따라 조건 \emph{a)}는 $p+q=n$일
때 사영계 $(\mathrm E_\infty^{pq}(A_k))_{k\in\mathbb Z}$가 본질적으로
상수라는 조건과 동치이다. 한편
$\operatorname{gr}^p(\mathrm R^n\mathrm T(A_k))$는 표준적으로
$\mathrm E_\infty^{p,n-p}(A_k)$와 동형이다. 따라서
(\hyperref[0.13.4.3-fr]{13.4.3})에서 \emph{a)}와 \emph{b)}의 동치가
따른다. 동형 (\hyperref[0.13.7.4.1-fr]{13.7.4.1})은 바로
(\hyperref[0.13.5.6.1-fr]{13.5.6.1})을 지금의 경우에 적용한 것이다.
\end{proof}

\begin{corollary}[13.7.5]
\label{0.13.7.5-fr}
\emph{$\mathcal F=(\mathcal F_k)_{k\in\mathbb N}$를
(\hyperref[0.13.3.1-fr]{13.3.1})의 조건 \emph{(i)}, \emph{(ii)},
\emph{(iii)}을 만족하는 아벨군층들의 사영계라 하고
$\mathcal F=\varprojlim_k\mathcal F_k$라 하자. 함자
$\mathcal G\mapsto\Gamma(X,\mathcal G)$에 대하여, $p+q=n$ 또는
$p+q=n+1$인 모든 쌍 $(p,q)$에서 사영계
$(\mathrm E_r^{pq}(\mathcal F))_{r\in\mathbb Z}$가 본질적으로
상수라고 가정하자. $H^{n+1}(X,\mathcal F)$ 위의 여과를
$\mathrm F^p(H^{n+1}(X,\mathcal F))=
\operatorname{Ker}(H^{n+1}(X,\mathcal F)\to
H^{n+1}(X,\mathcal F_{p-1}))$로 정의하면, 모든 $p\in\mathbb Z$에
대하여 표준 동형}
\begin{equation}
\label{0.13.7.5.1-fr}
\tag{13.7.5.1}
\operatorname{gr}^p(H^{n+1}(X,\mathcal F))
\xrightarrow{\sim}\mathrm E_\infty^{p,n-p+1}(\mathcal F)
\qquad\emph{모든 $p\in\mathbb Z$에 대하여}
\end{equation}
\emph{이 성립한다.}
\end{corollary}
\begin{proof}
\oldpage[0\textsubscript{III}]{77}
$C$가 $X$ 위의 아벨군층들의 범주, $C'$이 아벨군의 범주,
$\mathrm T=\Gamma$인 경우에
(\hyperref[0.13.7.4-fr]{13.7.4})를 적용하면, 모든 $p\in\mathbb Z$에
대하여 표준 동형
$\operatorname{gr}^p(\mathrm R^{n+1}\Gamma(\mathcal F))
\xrightarrow{\sim}\mathrm E_\infty^{p,n-p+1}(\mathcal F)$을 얻는다.
한편 (\hyperref[0.13.7.4-fr]{13.7.4})에 따라 사영계
$(H^n(X,\mathcal F_k))_{k\in\mathbb Z}$가 (ML)을 만족하므로,
(\hyperref[0.13.3.1-fr]{13.3.1})에서 표준 동형
\[
\label{0.13.7.5.1b-fr}
H^{n+1}(X,\mathcal F)\xrightarrow{\sim}
\varprojlim_k H^{n+1}(X,\mathcal F_k)
\]
을 얻는다.

사영계 $\mathrm R^{n+1}\Gamma(\mathcal F)$도
(\hyperref[0.13.7.4-fr]{13.7.4})에 따라 (ML)을 만족하므로 표준 동형
\[
\operatorname{gr}^p(\mathrm R^{n+1}\Gamma(\mathcal F))
\xrightarrow{\sim}
\varprojlim_k\operatorname{gr}^p(H^{n+1}(X,\mathcal F_k))
\]
이 있고 (\hyperref[0.13.4.3-fr]{13.4.3}), 표준 동형
\[
\varprojlim_k\operatorname{gr}^p(H^{n+1}(X,\mathcal F_k))
\xrightarrow{\sim}
\operatorname{gr}^p(\varprojlim_k H^{n+1}(X,\mathcal F_k))
\]
도 있다 (\hyperref[0.13.4.5-fr]{13.4.5}). 따라서 동형
(\hyperref[0.13.7.5.1-fr]{13.7.5.1})이 양변의 여과와 양립함을 보이면
충분하다. 이는 정의들과 모든 $p$에 대한 가환도식
\[
\xymatrix{
H^{n+1}(X,\mathcal F)\ar[r]^{\sim}\ar[dr]
&\varprojlim_k H^{n+1}(X,\mathcal F_k)\ar[d]\\
&H^{n+1}(X,\mathcal F_{p-1})
}
\]
에서 즉시 따른다.
\end{proof}

\begin{env}[13.7.6]
\label{0.13.7.6-fr}
$S$를 $\mathrm F^0(S)=S$인 여과 $(\mathrm F^i(S))_{i\in\mathbb Z}$를
갖는 환이라 하자. 따라서 $i<0$일 때
$\operatorname{gr}^i(S)=0$이다. 각 $A_k$가
(\hyperref[0.13.7.1-fr]{13.7.1})에서 정의한 여과에 관하여 여과
$S$-$C$-가군이고 (\hyperref[0.13.6.6-fr]{13.6.6}), $k\leq h$일 때
사상 $A_h\to A_k$는 여과 $S$-$C$-가군 구조의 사상이라고 가정하자.
줄여서 $\mathbf A$를 \emph{여과 $S$-$C$-가군들의 사영계}라 한다.
그러면 $k\leq h$, $1\leq r\leq+\infty$일 때 사상
$\mathrm B_r^{pq}(A_h)\to\mathrm B_r^{pq}(A_k)$와
$\mathrm Z_r^{pq}(A_h)\to\mathrm Z_r^{pq}(A_k)$는
(\hyperref[0.13.6.5-fr]{13.6.5})의 이중등급
$\operatorname{gr}^\bullet(S)$-$C'$-가군 구조의 사상이고, $r$이
유한할 때 족들 $(\mathrm Z_r^{pq}(\mathbf A))$,
$(\mathrm B_r^{pq}(\mathbf A))$, $(\mathrm E_r^{pq}(\mathbf A))$은
이중등급 $\operatorname{gr}^\bullet(S)$-$C'$-가군이다. 앞의 두 족은
$(\mathrm E_1^{pq}(\mathbf A))$의 부분가군이다. $p+q=n$인 항들만
취하여 얻는 앞의 대상들의 부분대상을 각각
$\mathrm Z_r^{(n)}(\mathbf A)$, $\mathrm B_r^{(n)}(\mathbf A)$,
$\mathrm E_r^{(n)}(\mathbf A)$로 나타낸다. 이들은 등급
$\operatorname{gr}^\bullet(S)$-$C'$-가군이다.

사영계 $(\mathrm E_r^{pq}(\mathbf A))_{r\in\mathbb Z}$가 본질적으로
상수이면, $(\mathrm E_\infty^{pq}(\mathbf A))$도 이중등급
$\operatorname{gr}^\bullet(S)$-$C'$-가군이고, 각
$\mathrm E_\infty^{(n)}(\mathbf A)$는 등급
$\operatorname{gr}^\bullet(S)$-$C'$-가군이다. 또한
$\beta^{p,n-p}:\mathrm E_\infty^{p,n-p}(A_k)\xrightarrow{\sim}
\operatorname{gr}^p(\mathrm R^n\mathrm T(A_k))$들은 각 $k$에 대하여
$\mathrm E_\infty^{(n)}(A_k)$에서
$\operatorname{gr}^\bullet(\mathrm R^n\mathrm T(A_k))$로 가는 등급
$\operatorname{gr}^\bullet(S)$-$C'$-가군의 동형을 이룬다. 앞의
조건들이 성립하면
\[
\beta^{p,n-p}:\mathrm E_\infty^{(n)}(\mathbf A)\xrightarrow{\sim}
\varprojlim_k\operatorname{gr}^\bullet(\mathrm R^n\mathrm T(A_k))
\]
도 이 구조에 관한 동형이다. 표준 동형
\[
\operatorname{gr}^\bullet(\mathrm R^n\mathrm T(\mathbf A))
\xrightarrow{\sim}
\varprojlim_k\operatorname{gr}^\bullet(\mathrm R^n\mathrm T(A_k))
\]
도 명백히 마찬가지이므로, 동형들
(\hyperref[0.13.7.4.1-fr]{13.7.4.1})은 등급
$\operatorname{gr}^\bullet(S)$-$C'$-가군 구조에 관한 동형을 이룬다.
\end{env}

\begin{proposition}[13.7.7]
\label{0.13.7.7-fr}
$S$를 뇌터 $\mathfrak J$-아딕 환이라 하자. $C$를 모든 대상이 어떤
단사 대상의 부분대상과 동형인 아벨 범주라 하고, $\mathrm T$를 $C$에서
아벨군의 범주로 가는 가법 공변 함자라 하자.
$\mathbf A=(A_k)_{k\in\mathbb Z}$를 $S$ 위의 $\mathfrak J$-아딕
여과에 관한 아래로 유계인 여과 $S$-$C$-가군들의 엄밀 사영계라 하자.
\oldpage[0\textsubscript{III}]{78}
주어진 정수 $n$에 대하여 다음 조건이 성립한다고 가정한다.
\begin{itemize}
\item[$(\mathrm F_n)$] (\hyperref[0.13.7.3.1-fr]{13.7.3.1})의 등급
$\operatorname{gr}^\bullet(S)$-가군
\[
\mathrm E_1^{(m)}(\mathbf A)=
\bigl(\mathrm R^m\mathrm T(\operatorname{gr}^p(\mathbf A))\bigr)_{p\in\mathbb Z}
\]
은 $m=n$과 $m=n+1$에 대하여 유한 생성이다.
\end{itemize}
이 조건들 아래 다음이 성립한다.
\begin{enumerate}
\item[\rm(i)] 사영계들
$(\mathrm R^n\mathrm T(A_k))_{k\in\mathbb Z}$와
$(\mathrm R^{n+1}\mathrm T(A_k))_{k\in\mathbb Z}$는 (ML)을 만족한다.
\item[\rm(ii)]
\[
\mathrm R'^{\,n}\mathrm T(\mathbf A)=
\varprojlim_k\mathrm R^n\mathrm T(A_k)
\]
로 놓으면 $\mathrm R'^{\,n}\mathrm T(\mathbf A)$는 유한 생성
$S$-가군이다.
\item[\rm(iii)] $\mathrm R'^{\,n}\mathrm T(\mathbf A)$ 위의 여과를
\[
\mathrm F^p(\mathrm R'^{\,n}\mathrm T(\mathbf A))=
\operatorname{Ker}\bigl(\mathrm R'^{\,n}\mathrm T(\mathbf A)
\longrightarrow\mathrm R^n\mathrm T(A_{p-1})\bigr)
\quad(p\in\mathbb Z)
\]
로 정의하면, 이 여과는 $\mathfrak J$-좋은 여과이다. 즉 모든 $p$에
대하여
$\mathfrak J\mathrm F^p(\mathrm R'^{\,n}\mathrm T(\mathbf A))
\subset\mathrm F^{p+1}(\mathrm R'^{\,n}\mathrm T(\mathbf A))$이고,
$p$가 충분히 크면 양변이 같다. 특히 이 여과가
$\mathrm R'^{\,n}\mathrm T(\mathbf A)$ 위에 정의하는 위상은
$\mathfrak J$-아딕 위상과 같다.
\item[\rm(iv)] $p+q=n$ 또는 $p+q=n+1$일 때 사영계
$(\mathrm E_r^{pq}(\mathbf A))_{r\in\mathbb Z}$는 본질적으로 상수이다.
따라서 $\mathrm E_\infty^{pq}(\mathbf A)$가 정의되고
(\hyperref[0.13.5.7-fr]{13.5.7}), 등급
$\operatorname{gr}^\bullet(S)$-가군의 표준 동형
\begin{equation}
\label{0.13.7.7.1-fr}
\tag{13.7.7.1}
\operatorname{gr}^p(\mathrm R'^{\,n}\mathrm T(\mathbf A))
\xrightarrow{\sim}\mathrm E_\infty^{p,n-p}(\mathbf A)
\quad(p\in\mathbb Z)
\end{equation}
이 성립한다.
\end{enumerate}
\end{proposition}

\begin{proof}
동형 (\hyperref[0.13.7.7.1-fr]{13.7.7.1})과 동형들
(\hyperref[0.13.7.4.1-fr]{13.7.4.1})을 고려하면, 용어를 남용하여
사영계 $\mathrm R^n\mathrm T(A_k)$의 사영극한
$\mathrm R'^{\,n}\mathrm T(\mathbf A)$을
$\mathrm R^n\mathrm T(\mathbf A)$로 나타낼 수 있다.

등급환 $\operatorname{gr}^\bullet(S)$는 뇌터이다
(Bourbaki, \emph{Alg. comm.}, chap.~III, \S~2,
n\textsuperscript{o}~9, cor.~5 du th.~2). 따라서
$\mathrm E_1^{(m)}(\mathbf A)$의 등급
$\operatorname{gr}^\bullet(S)$-부분가군들의 증가열
$\mathrm B_r^{(m)}(\mathbf A)$
(\hyperref[0.13.6.6-fr]{13.6.6})은 $m=n$과 $m=n+1$에 대하여
정상적이다. 그러므로 (\hyperref[0.11.1.10-fr]{11.1.10})의 조건 b)가
성립한다. 이에 따라 (\hyperref[0.13.7.4-fr]{13.7.4})의 조건 a)가
$n$과 $n+1$에 대하여 성립하며, 이것으로 (i)가 증명된다. 또한 동형들
(\hyperref[0.13.7.4.1-fr]{13.7.4.1})과
(\hyperref[0.13.7.6-fr]{13.7.6})의 주석에 따라
$\operatorname{gr}^\bullet(\mathrm R^n\mathrm T(\mathbf A))$는
\[
\mathrm E_\infty^{(n)}(\mathbf A)=
\mathrm Z_\infty^{(n)}(\mathbf A)/\mathrm B_\infty^{(n)}(\mathbf A)
\]
와 동형인 등급 $\operatorname{gr}^\bullet(S)$-가군이다.
$\mathrm Z_\infty^{(n)}(\mathbf A)$는
$\mathrm E_1^{(n)}(\mathbf A)$의 부분가군이므로 유한 생성이고, 따라서
$\mathrm E_\infty^{(n)}(\mathbf A)$도 유한 생성이다. 한편 여과
$(\mathrm F^p(\mathrm R'^{\,n}\mathrm T(\mathbf A)))$에 대하여
(\hyperref[0.13.4.5-fr]{13.4.5})에 따라
$\operatorname{gr}^\bullet(\mathrm R^n\mathrm T(\mathbf A))$와
$\operatorname{gr}^\bullet(\mathrm R'^{\,n}\mathrm T(\mathbf A))$는
동형인 $\operatorname{gr}^\bullet(S)$-가군이다. 이것으로 (iv)가
증명된다. 이제 (ii)와 (iii)은 앞의 결과들과 다음 보조정리에서 따른다.
\end{proof}

\begin{lemma}[13.7.7.2]
\label{0.13.7.7.2-fr}
$S$를 뇌터 $\mathfrak J$-아딕 환이라 하고, $M$을 공이산 여과
$(\mathrm F^p(M))_{p\in\mathbb Z}$를 갖는 $S$-가군이라 하자. 모든
$p$에 대하여
$\mathfrak J\mathrm F^p(M)\subset\mathrm F^{p+1}(M)$이라고 하자.
이는 $M$이 $\mathfrak J$-아딕 여과를 갖춘 여과환 $S$ 위의 여과
가군이라는 뜻이다. 더 나아가 $M$이 여과 $(\mathrm F^p(M))$이 정의하는
위상에 관하여 분리되어 있다고 가정하자. 그러면 다음 조건들은
동치이다.
\begin{enumerate}
\item[\rm a)] $M$은 유한 생성 $S$-가군이고 $(\mathrm F^p(M))$은
$\mathfrak J$-좋은 여과이다.
\item[\rm b)] $\operatorname{gr}^\bullet(M)$은 유한 생성
$\operatorname{gr}^\bullet(S)$-가군이다.
\item[\rm c)] 각 $\operatorname{gr}^p(M)$은 유한 생성 $S$-가군이고,
$p$가 충분히 크면 표준 준동형
\begin{equation}
\label{0.13.7.7.3-fr}
\tag{13.7.7.3}
\mathfrak J\otimes_S\operatorname{gr}^p(M)\longrightarrow
\operatorname{gr}^{p+1}(M)
\end{equation}
은 전사이다.
\oldpage[0\textsubscript{III}]{79}
이 준동형들은 $\mathfrak J\otimes_S\mathrm F^p(M)\to
\mathrm F^{p+1}(M)$에서 얻는다. 이때 합성 준동형
\[
\mathfrak J\otimes_S\mathrm F^{p+1}(M)\longrightarrow
\mathfrak J\otimes_S\mathrm F^p(M)\longrightarrow\mathrm F^{p+1}(M)
\]
의 상이
$\mathfrak J\mathrm F^{p+1}(M)\subset\mathrm F^{p+2}(M)$임을
고려한다.
\end{enumerate}
증명은 Bourbaki, \emph{Alg. comm.}, chap.~III, \S~3,
n\textsuperscript{o}~1, prop.~3을 보라.
\end{lemma}

\begin{env}[13.7.7.4]
\label{0.13.7.7.4-fr}
보조정리 (\hyperref[0.13.7.7.2-fr]{13.7.7.2})를 적용하려면, 고려하는
여과가 $\mathrm R'^{\,n}\mathrm T(\mathbf A)$ 위에 정의하는 위상에서
$\mathrm R'^{\,n}\mathrm T(\mathbf A)$가 분리 완비 $S$-가군임을
관찰하면 된다. 실제로 이 위상은 이산군들
$\mathrm R^n\mathrm T(A_k)$의 사영극한 위상이다. 한편 $k<k_0$일 때
$A_k=0$이면 $k<k_0$일 때 $\mathrm R^n\mathrm T(A_k)=0$도 성립하므로,
\[
\mathrm F^{k_0}(\mathrm R'^{\,n}\mathrm T(\mathbf A))=
\mathrm R'^{\,n}\mathrm T(\mathbf A).
\]
따라서 보조정리의 적용 조건들이 모두 충족된다.
\end{env}

\begin{corollary}[13.7.8]
\label{0.13.7.8-fr}
가정 $(\mathrm F_n)$이 성립하면, 모든 원소 $f\in S$에 대하여 표준 동형
\begin{equation}
\label{0.13.7.8.1-fr}
\tag{13.7.8.1}
\varprojlim_k\bigl((\mathrm R^n\mathrm T(A_k))_f\bigr)
\xrightarrow{\sim}\mathrm R^n\mathrm T(\mathbf A)\otimes_S S_{\{f\}}
\end{equation}
이 성립한다.
\end{corollary}

\begin{proof}
실제로 $\mathrm R^n\mathrm T(\mathbf A)$는 유한 생성 $S$-가군이고,
$S_{\{f\}}$는 뇌터 아딕 $S$-대수이며
(\hyperref[0.7.6.11-ko]{0\textsubscript{I}, 7.6.11}),
$S_f$의 $\mathfrak J$-준아딕 위상에 관한 분리 완비화이다
(\hyperref[0.7.6.2-ko]{0\textsubscript{I}, 7.6.2}).
(\hyperref[0.7.7.8-ko]{0\textsubscript{I}, 7.7.8})과
(\hyperref[0.7.7.1-ko]{0\textsubscript{I}, 7.7.1})에 따라
$\mathrm R^n\mathrm T(\mathbf A)\otimes_S S_{\{f\}}$는
$\mathrm R^n\mathrm T(\mathbf A)\otimes_S S_f$의
$\mathfrak J$-준아딕 위상에 관한 분리 완비화와 동형이다. 이 위상에서
0의 기본 근방계는
$(\mathfrak J^p\mathrm R^n\mathrm T(\mathbf A))\otimes_S S_f$이고,
따라서
$\mathrm F^p(\mathrm R^n\mathrm T(\mathbf A))\otimes_S S_f$도 그러한
기본계이다. 뒤의 가군은 표준 사상
$(\mathrm R^n\mathrm T(\mathbf A))_f\to
(\mathrm R^n\mathrm T(A_{p-1}))_f$의 핵이다. 그러므로
$\mathrm R^n\mathrm T(\mathbf A)\otimes_S S_f$에 연관된 분리군은
$\varprojlim_k((\mathrm R^n\mathrm T(A_k))_f)$의 어떤 부분군 $G$와
동일시된다. 그런데 사영계
$((\mathrm R^n\mathrm T(A_k))_f)$는 명백히 (ML)을 만족하고,
각 $(\mathrm R^n\mathrm T(A_k))_f$ 안에서
$(\mathrm R^n\mathrm T(\mathbf A))_f$의 상은 충분히 큰 $h\geq k$에
대한 $(\mathrm R^n\mathrm T(A_h))_f$들의 공통 상과 같다. 따라서
$G$는 $\varprojlim_k((\mathrm R^n\mathrm T(A_k))_f)$에서 조밀하다.
마지막 군은 분리 완비이므로 따름정리가 증명된다.
\end{proof}

\begin{flushright}

\textit{(계속.)}
\end{flushright}

\clearpage
\thispagestyle{empty}
\null
\clearpage
\thispagestyle{plain}
\markboth{연접층의 코호몰로지적 연구}{제III장}
\oldpage[III]{81}
\phantomsection
\label{chapter:III-fr}

\begin{center}
{\large 제III장\par}
\vspace{1.5em}
{\Large\bfseries 연접층의\par}
{\Large\bfseries 코호몰로지적 연구\par}
\vspace{1.8em}
{\bfseries 목차\par}
\end{center}

\medskip
\begin{tabular}{@{}r@{\ }p{0.82\textwidth}@{}}
\S~1. & 아핀 스킴의 코호몰로지.\\
\S~2. & 사영 사상의 코호몰로지적 연구.\\
\S~3. & 고유 사상에 대한 유한성 정리.\\
\S~4. & 고유 사상의 기본 정리~; 응용.\\
\S~5. & 연접 대수적 층의 존재 정리.\\
\S~6. & 국소 및 전역 Tor, Ext 함자~; 퀸네트 공식.\\
\S~7. & 가군의 공변 코호몰로지 함자에서의 밑변경에 대한
연구.\\
\S~8. & 사영 다발 위의 쌍대성 정리.\\
\S~9. & 상대 코호몰로지와 국소 코호몰로지~; 국소 쌍대성.\\
\S~10. & 사영 코호몰로지와 국소 코호몰로지 사이의 관계.
인자를 따른 형식적 완비화 기법.\\
\S~11. & 전역 및 국소 피카르 군\footnotemark.
\end{tabular}

\medskip
이 장에서는 연접 대수적 층의 코호몰로지에 관한 기본 정리들을 다룬다.
다만 유수 이론에서 나오는 정리들(쌍대성 정리들)은 제외하며, 이들은 뒤의
한 장에서 다룰 것이다. 앞의 정리들 가운데에는 본질적으로 여섯 개의 기본
정리가 있고, 이 장의 처음 여섯 절이 각각 이를 다룬다. 이 결과들은 뒤에서
본래 코호몰로지적인 성격을 갖지 않는 문제들에서도 핵심 도구가 될 것이며,
독자는 이미 \S~4에서 그 첫 사례들을 보게 될 것이다. \S~7은 더 기술적인
성격의 결과들을 제시하지만, 이 결과들은 응용에서 줄곧 사용된다. 마지막으로
\S\S~8부터 11까지에서는 연접층의 쌍대성과 관련되고 응용에 특히 중요한 몇몇
결과를 전개한다. 이 결과들은 일반 유수 이론보다 앞서 설명할 수 있다.

\footnotetext{제IV장은 \S\S~8부터 11까지에 의존하지 않으며, 아마도 이들보다
먼저 출판될 것이다.}

\oldpage[III]{82}
\S\S~1과 2의 내용은 J.-P. Serre에 의한 것이며, 독자는 우리가 (FAC)의
서술을 따르기만 하면 되었음을 알게 될 것이다. \S\S~8과 9 역시 (FAC)에서
착안한 것이다(다만 서로 다른 맥락 때문에 필요한 변환은 그리 자명하지 않다).
마지막으로 서론에서 말했듯이, \S~4는 Zariski의 ``정칙함수론''에서 말하는
기본적인 ``불변성 정리''를 현대 언어로 정식화한 것으로 보아야 한다.

끝으로 3.4번의 결과들(그리고 (0, 13.4부터 13.7까지)의 예비 명제들)은
제III장의 나머지 부분에서 사용되지 않으므로 첫 독해에서는 생략해도 된다.

\setcounter{section}{0}
\section{아핀 스킴의 코호몰로지}
\label{section:III.1-fr}

\subsection{외대수 복합체에 관한 복습}
\label{subsection:III.1.1-fr}

\begin{env}[1.1.1]
\label{III.1.1.1-fr}
$A$를 환이라 하고, $\mathbf f=(f_i)_{1\leq i\leq r}$를 $A$의 $r$개
원소로 이루어진 계라 하자. $\mathbf f$에 대응하는 \emph{외대수 복합체
$K_\bullet(\mathbf f)$}는 다음과 같이 정의되는 사슬 복합체이다
(G, I, 2.2). 등급 $A$-가군 $K_\bullet(\mathbf f)$는 통상적인 방식으로
등급이 매겨진 \emph{외대수 $\bigwedge(A^r)$}이고, 경계 연산자는
쌍대가군 $(A^r)^\vee$의 원소로 본 $\mathbf f$에 의한 \emph{내부곱
$i_{\mathbf f}$}이다. $i_{\mathbf f}$는 $\bigwedge(A^r)$의 차수 $-1$인
\emph{반미분}이고, $(\mathbf e_i)_{1\leq i\leq r}$가 $A^r$의 표준기저이면
$i_{\mathbf f}(\mathbf e_i)=f_i$임을 상기하자. 조건
$i_{\mathbf f}\circ i_{\mathbf f}=0$은 곧바로 확인된다.

동치인 정의는 다음과 같다. 각 $i$에 대하여 다음과 같이 정의되는 사슬
복합체 $K_\bullet(f_i)$를 생각한다.
$K_0(f_i)=K_1(f_i)=A$, $K_n(f_i)=0$ ($n\neq0,1$)로 두고, 경계 연산자는
$d_1:A\to A$가 \emph{$f_i$에 의한 곱셈}이 되도록 정의한다. 그러면
$K_\bullet(\mathbf f)$를 전체 차수를 부여한 \emph{텐서곱
$K_\bullet(f_1)\otimes K_\bullet(f_2)\otimes\cdots\otimes K_\bullet(f_r)$}
(G, I, 2.7)로 잡는다. 이 복합체와 앞에서 정의한 복합체 사이의 동형사상은
곧바로 확인된다.
\end{env}

\begin{env}[1.1.2]
\label{III.1.1.2-fr}
임의의 $A$-가군 $M$에 대하여 \emph{사슬 복합체}
\begin{equation}
\label{III.1.1.2.1-fr}
\tag{1.1.2.1}
K_\bullet(\mathbf f,M)=K_\bullet(\mathbf f)\otimes_A M
\end{equation}
와 \emph{공사슬 복합체} (G, I, 2.2)
\begin{equation}
\label{III.1.1.2.2-fr}
\tag{1.1.2.2}
K^\bullet(\mathbf f,M)=\operatorname{Hom}_A(K_\bullet(\mathbf f),M).
\end{equation}
를 정의한다. $g$가 뒤 복합체의 $k$-공사슬이고
\[
g(i_1,\ldots,i_k)=g(\mathbf e_{i_1}\wedge\cdots\wedge\mathbf e_{i_k})
\]
로 놓으면, $g$는 $[1,r]^k$에서 $M$으로 가는 \emph{교대 사상}과 동일시된다.
앞의 정의들로부터 다음을 얻는다.
\begin{equation}
\label{III.1.1.2.3-fr}
\tag{1.1.2.3}
d^kg(i_1,i_2,\ldots,i_{k+1})=
\sum_{h=1}^{k+1}(-1)^{h-1}f_{i_h}
g(i_1,\ldots,\widehat{i_h},\ldots,i_{k+1}).
\end{equation}
\end{env}

\oldpage[III]{83}
\begin{env}[1.1.3]
\label{III.1.1.3-fr}
앞의 복합체들로부터 통상적인 방식으로 \emph{$A$-호몰로지 가군과
$A$-코호몰로지 가군} (G, I, 2.2)
\begin{align}
\label{III.1.1.3.1-fr}
\tag{1.1.3.1}
H_\bullet(\mathbf f,M)&=H_\bullet(K_\bullet(\mathbf f,M)),\\
\label{III.1.1.3.2-fr}
\tag{1.1.3.2}
H^\bullet(\mathbf f,M)&=H^\bullet(K^\bullet(\mathbf f,M))
\end{align}
을 얻는다. 한편 각 사슬
$z=\sum(\mathbf e_{i_1}\wedge\cdots\wedge\mathbf e_{i_k})\otimes
z_{i_1\ldots i_k}$에 공사슬 $g_z$를 대응시켜 $A$-동형사상
$K_i(\mathbf f,M)\xrightarrow{\sim}K^{r-i}(\mathbf f,M)$을 정의한다. 여기서
$g_z(j_1,\ldots,j_{r-k})=\varepsilon z_{i_1\ldots i_k}$이고,
$(j_h)_{1\leq h\leq r-k}$는 $[1,r]$에서 엄밀히 증가하는 수열
$(i_h)_{1\leq h\leq k}$의 여수열이며,
$\varepsilon=(-1)^\nu$이다. 이때 $\nu$는 $[1,r]$의 순열
$i_1,\ldots,i_k,j_1,\ldots,j_{r-k}$의 역전수이다. $g_{dz}=d(g_z)$임을
확인할 수 있으므로 다음 동형사상을 얻는다.
\begin{equation}
\label{III.1.1.3.3-fr}
\tag{1.1.3.3}
H^i(\mathbf f,M)\xrightarrow{\sim}H_{r-i}(\mathbf f,M)
\qquad(0\leq i\leq r).
\end{equation}
이 장에서는 주로 코호몰로지 가군 $H^\bullet(\mathbf f,M)$을 다룬다.

$\mathbf f$를 고정하면 $M\mapsto H^\bullet(\mathbf f,M)$이 $A$-가군의
범주에서 등급 $A$-가군의 범주로 가며 차수 $<0$ 및 $>r$에서 영인
\emph{코호몰로지 함자} (T, II, 2.1)임은 자명하다 (G, I, 2.1).
또한
\begin{equation}
\label{III.1.1.3.4-fr}
\tag{1.1.3.4}
H^0(\mathbf f,M)=\operatorname{Hom}_A(A/(\mathbf f),M)
\end{equation}
이다. 여기서 $(\mathbf f)$는 $f_1,\ldots,f_r$로 생성되는 $A$의
아이디얼을 나타낸다. 이는 (\hyperref[III.1.1.2.3-fr]{1.1.2.3})에서
곧바로 따르며, $H^0(\mathbf f,M)$은 \emph{$(\mathbf f)$에 의해 소멸되는}
$M$의 부분가군과 동일시된다. 마찬가지로
(\hyperref[III.1.1.2.3-fr]{1.1.2.3})에 의해
\begin{equation}
\label{III.1.1.3.5-fr}
\tag{1.1.3.5}
H^r(\mathbf f,M)=M/\Bigl(\sum_{i=1}^r f_iM\Bigr)
=(A/(\mathbf f))\otimes_A M.
\end{equation}
완전성을 위해 증명을 상기할 다음의 알려진 결과를 사용할 것이다.
\end{env}

\begin{proposition}[1.1.4]
\label{III.1.1.4-fr}
$A$를 환, $\mathbf f=(f_i)_{1\leq i\leq r}$를 $A$의 원소들로 이루어진
유한족, $M$을 $A$-가군이라 하자. $1\leq i\leq r$일 때마다
$M_{i-1}=M/(f_1M+\cdots+f_{i-1}M)$에서의 곱셈 사상
$z\mapsto f_i z$가 단사이면, $i\neq r$인 모든 $i$에 대하여
$H^i(\mathbf f,M)=0$이다.
\end{proposition}

\begin{proof}
(\hyperref[III.1.1.3.3-fr]{1.1.3.3})에 의해 모든 $i>0$에 대하여
$H_i(\mathbf f,M)=0$임을 보이면 충분하다. $r$에 관하여 귀납하자.
$r=0$인 경우는 자명하다. $\mathbf f'=(f_i)_{1\leq i\leq r-1}$로 놓자.
이 족은 명제의 조건을 만족하므로 $L_\bullet=K_\bullet(\mathbf f',M)$로
놓으면 귀납가정에 의해 $i>0$에서 $H_i(L_\bullet)=0$이고,
(\hyperref[III.1.1.3.3-fr]{1.1.3.3})과
(\hyperref[III.1.1.3.5-fr]{1.1.3.5})에 의해
$H_0(L_\bullet)=M_{r-1}$이다. 줄여서
$K_\bullet=K_\bullet(f_r)=K_0\oplus K_1$로 쓰자. 여기서
$K_0=K_1=A$이고 $d_1:K_1\to K_0$는 $f_r$에 의한 곱셈이다.
정의 (\hyperref[III.1.1.1-fr]{1.1.1})에 의해
$K_\bullet(\mathbf f,M)=K_\bullet\otimes_A L_\bullet$이다. 이제 다음 보조정리를
사용한다.

\begin{lemma}[1.1.4.1]
\label{III.1.1.4.1-fr}
$K_\bullet$을 차원 $0$과 $1$ 이외에서는 영인 자유 $A$-가군들로 이루어진
사슬 복합체라 하자. $A$-가군들로 이루어진 임의의 사슬 복합체
$L_\bullet$과 임의의 첨자 $p$에 대하여 완전열
\[
0\longrightarrow H_0(K_\bullet\otimes H_p(L_\bullet))
\longrightarrow H_p(K_\bullet\otimes L_\bullet)
\longrightarrow H_1(K_\bullet\otimes H_{p-1}(L_\bullet))
\longrightarrow0
\]
이 존재한다.
\end{lemma}

\oldpage[III]{84}
이는 퀸네트 스펙트럼 열의 저차항 완전열의 특수한 경우이다
(M, XVII, 5.2 a) 및 G, I, 5.5.2). 다음과 같이 직접 증명할 수도 있다.
$K_0$과 $K_1$을 각각 차원 $\neq0$, $\neq1$에서 영인 사슬 복합체로 보자.
그러면 복합체의 완전열
\[
0\longrightarrow K_0\otimes L_\bullet
\longrightarrow K_\bullet\otimes L_\bullet
\longrightarrow K_1\otimes L_\bullet\longrightarrow0
\]
을 얻으며, 여기에 호몰로지 완전열
\[
\cdots\longrightarrow H_{p+1}(K_1\otimes L_\bullet)
\xrightarrow{\partial}H_p(K_0\otimes L_\bullet)
\longrightarrow H_p(K_\bullet\otimes L_\bullet)
\longrightarrow H_p(K_1\otimes L_\bullet)
\xrightarrow{\partial}H_{p-1}(K_0\otimes L_\bullet)
\longrightarrow\cdots.
\]
을 적용할 수 있다. 그런데 모든 $p$에 대하여
$H_p(K_0\otimes L_\bullet)=K_0\otimes H_p(L_\bullet)$이고
$H_p(K_1\otimes L_\bullet)=K_1\otimes H_{p-1}(L_\bullet)$임은 자명하다.
또한 연산자
$\partial:K_1\otimes H_p(L_\bullet)\to K_0\otimes H_p(L_\bullet)$가 바로
$d_1\otimes1$임을 곧바로 확인한다. 따라서 앞의 완전열과
$H_0(K_\bullet\otimes H_p(L_\bullet))$ 및
$H_1(K_\bullet\otimes H_{p-1}(L_\bullet))$의 정의로부터 보조정리가 따른다.

보조정리가 증명되었으므로 (\hyperref[III.1.1.4-fr]{1.1.4})의 증명의
나머지는 즉각적이다. 귀납가정과 보조정리
(\hyperref[III.1.1.4.1-fr]{1.1.4.1})에 의해 $p\geq2$에서
$H_p(K_\bullet\otimes L_\bullet)=0$이다. 더 나아가
$H_1(K_\bullet,H_0(L_\bullet))=0$임을 보이면 보조정리
(\hyperref[III.1.1.4.1-fr]{1.1.4.1})에서
$H_1(K_\bullet\otimes L_\bullet)=0$도 얻는다. 그런데 정의에 의해
$H_1(K_\bullet,H_0(L_\bullet))$은 $M_{r-1}$에서의 곱셈 사상
$z\mapsto f_rz$의 핵과 다름없고, 가정에 따라 이 핵은 영이다. 이로써
증명이 끝난다.
\end{proof}

\begin{env}[1.1.5]
\label{III.1.1.5-fr}
$\mathbf g=(g_i)_{1\leq i\leq r}$를 $A$의 $r$개 원소로 이루어진 또 다른
수열이라 하고, $\mathbf f\mathbf g=(f_ig_i)_{1\leq i\leq r}$로 놓자.
$A^r$의 $A$-선형 사상
$(x_1,\ldots,x_r)\mapsto(g_1x_1,\ldots,g_rx_r)$을 외대수
$\bigwedge(A^r)$로 표준 연장하여 복합체의 표준 준동형
\begin{equation}
\label{III.1.1.5.1-fr}
\tag{1.1.5.1}
\varphi_{\mathbf g}:K_\bullet(\mathbf f\mathbf g)\longrightarrow
K_\bullet(\mathbf f)
\end{equation}
을 정의할 수 있다. 이것이 실제로 복합체의 준동형임을 보이려면 일반적으로
$u:E\to F$가 $A$-선형 사상이고, $\mathbf x\in F^\vee$이며,
$\mathbf y={}^tu(\mathbf x)\in E^\vee$이면 다음 공식이 성립함을 확인하면
충분하다.
\begin{equation}
\label{III.1.1.5.2-fr}
\tag{1.1.5.2}
(\bigwedge u)\circ i_{\mathbf y}=i_{\mathbf x}\circ(\bigwedge u)~;
\end{equation}
실제로 양변은 $\bigwedge F$의 반미분이므로, 이들이 $F$에서 일치함만
확인하면 충분하고 이는 정의에서 곧바로 따른다.

$K_\bullet(\mathbf f)$를 $K_\bullet(f_i)$들의 텐서곱과 동일시하면
(\hyperref[III.1.1.1-fr]{1.1.1}), $\varphi_{\mathbf g}$는
$\varphi_{g_i}$들의 텐서곱이다. 여기서 $\varphi_{g_i}$는 차수 $0$에서는
항등사상이고 차수 $1$에서는 $g_i$에 의한 곱셈이다.
\end{env}

\begin{env}[1.1.6]
\label{III.1.1.6-fr}
특히 $0\leq n\leq m$인 임의의 두 정수 $m,n$에 대하여 복합체의 준동형
\begin{equation}
\label{III.1.1.6.1-fr}
\tag{1.1.6.1}
\varphi_{\mathbf f^{m-n}}:K_\bullet(\mathbf f^m)\longrightarrow
K_\bullet(\mathbf f^n)
\end{equation}
을 얻고, 따라서 준동형들
\begin{align}
\label{III.1.1.6.2-fr}
\tag{1.1.6.2}
\varphi_{\mathbf f^{m-n}} &:K^\bullet(\mathbf f^n,M)\longrightarrow
K^\bullet(\mathbf f^m,M),\\
\label{III.1.1.6.3-fr}
\tag{1.1.6.3}
\varphi_{\mathbf f^{m-n}} &:H^\bullet(\mathbf f^n,M)\longrightarrow
H^\bullet(\mathbf f^m,M)
\end{align}
을 얻는다.

\oldpage[III]{85}
뒤의 준동형들은 $p\leq n\leq m$일 때의 추이성 조건
$\varphi_{\mathbf f^{m-p}}=
\varphi_{\mathbf f^{m-n}}\circ\varphi_{\mathbf f^{n-p}}$를 자명하게 만족한다.
따라서 이들은 $A$-가군의 두 \emph{귀납계}를 정의한다. 다음과 같이 놓는다.
\begin{align}
\label{III.1.1.6.4-fr}
\tag{1.1.6.4}
C^\bullet((\mathbf f),M)&=\varinjlim_n K^\bullet(\mathbf f^n,M),\\
\label{III.1.1.6.5-fr}
\tag{1.1.6.5}
H^\bullet((\mathbf f),M)&=H^\bullet(C^\bullet((\mathbf f),M))
=\varinjlim_n H^\bullet(\mathbf f^n,M).
\end{align}
마지막 등식은 귀납극한을 취하는 것이 함자 $H^\bullet$과 교환한다는 사실에서
따른다 (G, I, 2.1). 뒤에서 (\hyperref[III.1.4.3-fr]{1.4.3})
$H^\bullet((\mathbf f),M)$이 실제로는 $A$의 \emph{아이디얼
$(\mathbf f)$}에만 의존함(더 나아가 $A$ 위의 $(\mathbf f)$-준아딕 위상에만
의존함)을 볼 것이며, 이것이 이 표기를 정당화한다.

$M\mapsto C^\bullet((\mathbf f),M)$은 완전 $A$-선형 함자이고,
$M\mapsto H^\bullet((\mathbf f),M)$은 코호몰로지 함자임이 자명하다.
\end{env}

\begin{env}[1.1.7]
\label{III.1.1.7-fr}
$\mathbf f=(f_i)\in A^r$, $\mathbf g=(g_i)\in A^r$라 하자. 외대수
$\bigwedge(A^r)$에서 벡터 $\mathbf g\in A^r$에 의한 왼쪽 곱셈을
$e_{\mathbf g}$로 나타내자. $A$-가군 $\bigwedge(A^r)$에서 다음의
\emph{호모토피 공식}이 성립함이 알려져 있다(여기서 $1$은
$\bigwedge(A^r)$의 항등 자기동형사상을 뜻한다).
\begin{equation}
\label{III.1.1.7.1-fr}
\tag{1.1.7.1}
i_{\mathbf f}e_{\mathbf g}+e_{\mathbf g}i_{\mathbf f}
=\langle\mathbf g,\mathbf f\rangle 1
\end{equation}
이 관계는 또한 \emph{복합체 $K_\bullet(\mathbf f)$ 안에서}
\begin{equation}
\label{III.1.1.7.2-fr}
\tag{1.1.7.2}
de_{\mathbf g}+e_{\mathbf g}d=\langle\mathbf g,\mathbf f\rangle 1.
\end{equation}
임을 뜻한다. 아이디얼 $(\mathbf f)$가 $A$와 같으면
$\langle\mathbf g,\mathbf f\rangle=\sum_{i=1}^r g_if_i=1$인
$\mathbf g\in A^r$가 존재한다. 따라서 (G, I, 2.4) 다음을 얻는다.
\end{env}

\begin{proposition}[1.1.8]
\label{III.1.1.8-fr}
$f_i$들이 생성하는 아이디얼 $(\mathbf f)$가 $A$와 같다고 하자. 그러면
복합체 $K_\bullet(\mathbf f)$는 호모토피적으로 자명하고, 따라서 임의의
$A$-가군 $M$에 대하여 복합체 $K_\bullet(\mathbf f,M)$과
$K^\bullet(\mathbf f,M)$도 그러하다.
\end{proposition}

\begin{corollary}[1.1.9]
\label{III.1.1.9-fr}
$(\mathbf f)=A$이면 임의의 $A$-가군 $M$에 대하여
$H^\bullet(\mathbf f,M)=0$이고 $H^\bullet((\mathbf f),M)=0$이다.
\end{corollary}

\begin{proof}
실제로 이 경우 모든 $n$에 대하여 $(\mathbf f^n)=A$이다.
\end{proof}

\begin{remark}[1.1.10]
\label{III.1.1.10-fr}
위와 같은 표기 아래 $X=\operatorname{Spec}(A)$라 하고, $Y$를 아이디얼
$(\mathbf f)$로 정의되는 $X$의 닫힌 부분스킴이라 하자. \S~9에서
$H^\bullet((\mathbf f),M)$이 $Y$의 닫힌 부분집합들의 반필터 $\Phi$에
대응하는 코호몰로지 $H_Y^\bullet(X,\widetilde M)$ (T, 3.2)과 동형임을
증명할 것이다. 또한 $X$와 $\mathcal F=\widetilde M$에 적용한 명제
(\hyperref[III.1.2.3-fr]{1.2.3})가 코호몰로지 완전열
\[
\cdots\longrightarrow H_Y^p(X,\mathcal F)\longrightarrow H^p(X,\mathcal F)
\longrightarrow H^p(X-Y,\mathcal F)\longrightarrow H_Y^{p+1}(X,\mathcal F)
\longrightarrow\cdots.
\]
의 특수한 경우임도 보일 것이다.
\end{remark}

\subsection{열린 덮개의 체흐 코호몰로지}
\label{subsection:III.1.2-fr}

\begin{notation}[1.2.1]
\label{III.1.2.1-fr}
이 번호에서는 다음 표기를 사용한다.
\begin{itemize}
\item $X$는 준스킴이다.
\item $\mathcal F$는 준연접 $\mathcal O_X$-가군이다.
\oldpage[III]{86}
\item $A=\Gamma(X,\mathcal O_X)$, $M=\Gamma(X,\mathcal F)$이다.
\item $\mathbf f=(f_i)_{1\leq i\leq r}$는 $A$의 원소들로 이루어진 유한계이다.
\item $U_i=X_{f_i}$는 $f_i(x)\neq0$인 $x\in X$들로 이루어진 열린집합
(\hyperref[0.5.5.2-ko]{0\textsubscript{I}, 5.5.2})이다.
\item $U=\bigcup_{i=1}^r U_i$이다.
\item $\mathfrak U$는 $U$의 덮개 $(U_i)_{1\leq i\leq r}$이다.
\end{itemize}
\end{notation}

\begin{env}[1.2.2]
\label{III.1.2.2-fr}
$X$가 바탕공간이 \emph{뇌터}인 준스킴이거나, 바탕공간이
\emph{준콤팩트}인 \emph{스킴}이라고 하자. 그러면
(\hyperref[I.9.3.3-ko]{I, 9.3.3}) $\Gamma(U_i,\mathcal F)=M_{f_i}$임이
알려져 있다. 다음과 같이 놓자.
\[
U_{i_0i_1\ldots i_p}=\bigcap_{k=0}^p U_{i_k}
=X_{f_{i_0}f_{i_1}\cdots f_{i_p}}
\]
(\hyperref[0.5.5.3-ko]{0\textsubscript{I}, 5.5.3})이며, 따라서 또한
\begin{equation}
\label{III.1.2.2.1-fr}
\tag{1.2.2.1}
\Gamma(U_{i_0i_1\ldots i_p},\mathcal F)=M_{f_{i_0}f_{i_1}\cdots f_{i_p}}.
\end{equation}
그런데 (\hyperref[0.1.6.1-ko]{0\textsubscript{I}, 1.6.1})
$M_{f_{i_0}f_{i_1}\cdots f_{i_p}}$는 귀납극한
$\varinjlim_n M_{i_0i_1\ldots i_p}^{(n)}$과 동일시된다. 여기서 귀납계는
$M_{i_0i_1\ldots i_p}^{(n)}=M$들로 이루어지고, $m\leq n$일 때 준동형
$\varphi_{nm}:M_{i_0i_1\ldots i_p}^{(m)}\to
M_{i_0i_1\ldots i_p}^{(n)}$은
$(f_{i_0}f_{i_1}\cdots f_{i_p})^{n-m}$에 의한 곱셈이다. $C_n^p(M)$을
모든 $n$에 대하여 $[1,r]^{p+1}$에서 $M$으로 가는 \emph{교대 사상}들의
집합이라 하자. 이 $A$-가군들도 $\varphi_{nm}$에 의해 귀납계를 이룬다.
$C^p(\mathfrak U,\mathcal F)$가 덮개 $\mathfrak U$에 관하여
$\mathcal F$를 계수로 하는 체흐의 \emph{교대} $p$-공사슬군이면
(G, II, 5.1), 앞에서 본 것으로부터
\begin{equation}
\label{III.1.2.2.2-fr}
\tag{1.2.2.2}
C^p(\mathfrak U,\mathcal F)=\varinjlim_n C_n^p(M)
\end{equation}
으로 쓸 수 있다. 그런데 (\hyperref[III.1.1.2-fr]{1.1.2})의 표기 아래
$C_n^p(M)$은 $K^{p+1}(\mathbf f^n,M)$과 동일시되고, 사상
$\varphi_{nm}$은 (\hyperref[III.1.1.6-fr]{1.1.6})에서 정의한 사상
$\varphi_{\mathbf f^{n-m}}$과 동일시된다. 따라서 모든 $p\geq0$에 대하여
$\mathcal F$에 관해 함자적인 표준 동형사상
\begin{equation}
\label{III.1.2.2.3-fr}
\tag{1.2.2.3}
C^p(\mathfrak U,\mathcal F)\xrightarrow{\sim}C^{p+1}((\mathbf f),M)
\end{equation}
을 얻는다. 더 나아가 공식 (\hyperref[III.1.1.2.3-fr]{1.1.2.3})과 덮개의
코호몰로지 정의 (G, II, 5.1)는 동형사상들
(\hyperref[III.1.2.2.3-fr]{1.2.2.3})이 공경계 연산자들과 양립함을 보인다.
\end{env}

\begin{proposition}[1.2.3]
\label{III.1.2.3-fr}
$X$가 바탕공간이 뇌터인 준스킴이거나 바탕공간이 준콤팩트인 스킴이면,
$\mathcal F$에 관해 함자적인 표준 동형사상
\begin{equation}
\label{III.1.2.3.1-fr}
\tag{1.2.3.1}
H^p(\mathfrak U,\mathcal F)\xrightarrow{\sim}H^{p+1}((\mathbf f),M)
\qquad\text{모든 }p\geq1\text{에 대하여}
\end{equation}
이 존재한다. 또한 $\mathcal F$에 관해 함자적인 완전열
\begin{equation}
\label{III.1.2.3.2-fr}
\tag{1.2.3.2}
0\longrightarrow H^0((\mathbf f),M)\longrightarrow M
\longrightarrow H^0(\mathfrak U,\mathcal F)
\longrightarrow H^1((\mathbf f),M)\longrightarrow0
\end{equation}
이 존재한다.
\end{proposition}

\oldpage[III]{87}
\begin{proof}
관계들 (\hyperref[III.1.2.3.1-fr]{1.2.3.1})은 실제로
(\hyperref[III.1.2.2-fr]{1.2.2})에서 본 것으로부터 즉시 따른다. 한편
$C^0(\mathfrak U,\mathcal F)=C^1((\mathbf f),M)$이다. 따라서
$H^0(\mathfrak U,\mathcal F)$는 $C^1((\mathbf f),M)$의 $1$-공순환들로
이루어진 부분군과 동일시된다. $M=C^0((\mathbf f),M)$이므로 완전열
(\hyperref[III.1.2.3.2-fr]{1.2.3.2})은 코호몰로지군
$H^0((\mathbf f),M)$과 $H^1((\mathbf f),M)$의 정의에서 나오는 완전열과
다름없다.
\end{proof}

\begin{corollary}[1.2.4]
\label{III.1.2.4-fr}
$X_{f_i}$들이 준콤팩트이고
$\sum_i g_i(f_i\vert U)=1\vert U$를 만족하는
$g_i\in\Gamma(U,\mathcal O_X)$들이 존재한다고 하자. 그러면 임의의
준연접 $(\mathcal O_X\vert U)$-가군 $\mathcal G$와 모든 $p>0$에 대하여
$H^p(\mathfrak U,\mathcal G)=0$이다. 더 나아가 $U=X$이면 표준 준동형
(\hyperref[III.1.2.3.2-fr]{1.2.3.2})
$\Gamma(X,\mathcal G)\to H^0(\mathfrak U,\mathcal G)$는 전단사이다.
\end{corollary}

\begin{proof}
가정에 따라 $U_i=X_{f_i}$들이 준콤팩트이므로 $U$도 그러하다. 따라서
$U=X$인 경우로 한정해도 된다. 이때 가정으로부터 모든 $p\geq0$에 대하여
$H^p((\mathbf f),M)=0$이다 (\hyperref[III.1.1.9-fr]{1.1.9}). 이제 따름정리는
(\hyperref[III.1.2.3.1-fr]{1.2.3.1})과
(\hyperref[III.1.2.3.2-fr]{1.2.3.2})에서 곧바로 따른다.
\end{proof}

$H^0(\mathfrak U,\mathcal F)=H^0(U,\mathcal F)$ (G, II, 5.2.2)이므로,
이로써 (I, 1.3.7)을 특수한 경우로 다시 증명했음을 주목하자.

\begin{remark}[1.2.5]
\label{III.1.2.5-fr}
$X$가 \emph{아핀 스킴}이라고 하자. 그러면
$U_i=X_{f_i}=D(f_i)$들과 $U_{i_0i_1\ldots i_p}$들은 아핀 열린집합이다
(다만 $U$는 반드시 아핀일 필요는 없다). 이 경우 함자
$\Gamma(X,\mathcal F)$와 $\Gamma(U_{i_0i_1\ldots i_p},\mathcal F)$는
$\mathcal F$에 관해 완전하다 (I, 1.3.11). 준연접
$\mathcal O_X$-가군들의 완전열
$0\to\mathcal F'\to\mathcal F\to\mathcal F''\to0$이 주어지면,
복합체의 열
\[
0\longrightarrow C^\bullet(\mathfrak U,\mathcal F')
\longrightarrow C^\bullet(\mathfrak U,\mathcal F)
\longrightarrow C^\bullet(\mathfrak U,\mathcal F'')
\longrightarrow0
\]
은 완전하고, 따라서 코호몰로지 완전열
\[
\cdots\longrightarrow H^p(\mathfrak U,\mathcal F')
\longrightarrow H^p(\mathfrak U,\mathcal F)
\longrightarrow H^p(\mathfrak U,\mathcal F'')
\xrightarrow{\partial}H^{p+1}(\mathfrak U,\mathcal F')
\longrightarrow\cdots.
\]
을 준다. 한편 $M'=\Gamma(X,\mathcal F')$,
$M''=\Gamma(X,\mathcal F'')$로 놓으면
$0\to M'\to M\to M''\to0$은 완전하다.
$C^\bullet((\mathbf f),M)$이 $M$에 관해 완전한 함자이므로 다음 코호몰로지
완전열도 얻는다.
\[
\cdots\longrightarrow H^p((\mathbf f),M')
\longrightarrow H^p((\mathbf f),M)
\longrightarrow H^p((\mathbf f),M'')
\xrightarrow{\partial}H^{p+1}((\mathbf f),M')
\longrightarrow\cdots.
\]
이제 도식
\[
\xymatrix{
0\ar[r]&C^\bullet(\mathfrak U,\mathcal F')\ar[r]\ar[d]
&C^\bullet(\mathfrak U,\mathcal F)\ar[r]\ar[d]
&C^\bullet(\mathfrak U,\mathcal F'')\ar[r]\ar[d]&0\\
0\ar[r]&C^\bullet((\mathbf f),M')\ar[r]
&C^\bullet((\mathbf f),M)\ar[r]
&C^\bullet((\mathbf f),M'')\ar[r]&0
}
\]
\oldpage[III]{88}
이 가환이므로, 모든 $p$에 대하여 도식들
\begin{equation}
\label{III.1.2.5.1-fr}
\tag{1.2.5.1}
\xymatrix{
H^p(\mathfrak U,\mathcal F'')\ar[r]^{\partial}\ar[d]
&H^{p+1}(\mathfrak U,\mathcal F')\ar[d]\\
H^{p+1}((\mathbf f),M'')\ar[r]^{\partial}
&H^{p+2}((\mathbf f),M')
}
\end{equation}
이 가환임을 얻는다 (G, I, 2.1.1).
\end{remark}

\subsection{아핀 스킴의 코호몰로지}
\label{subsection:III.1.3-fr}

\begin{theorem}[1.3.1]
\label{III.1.3.1-fr}
$X$를 아핀 스킴이라 하자. 임의의 준연접 $\mathcal O_X$-가군
$\mathcal F$와 모든 $p>0$에 대하여 $H^p(X,\mathcal F)=0$이다.
\end{theorem}

\begin{proof}
$\mathfrak U$를 아핀 열린집합 $X_{f_i}=D(f_i)$ ($1\leq i\leq r$)들로
이루어진 $X$의 유한 덮개라 하자. 이때 $f_i$들이
$A=\Gamma(X,\mathcal O_X)$에서 생성하는 아이디얼은 $A$와 같음이 알려져
있다. 따라서 (\hyperref[III.1.2.4-fr]{1.2.4})로부터 $p>0$일 때
$H^p(\mathfrak U,\mathcal F)=0$을 얻는다. 아핀 열린집합들로 이루어진
임의로 세분된 $X$의 유한 덮개들이 존재하므로 (I, 1.1.10), 체흐
코호몰로지의 정의 (G, II, 5.8)로부터 $p>0$일 때
$\check H^p(X,\mathcal F)=0$도 얻는다. 그런데 이는 임의의 $f\in A$에
대한 모든 준스킴 $X_f$에도 적용되므로 (I, 1.3.6), $p>0$일 때
$\check H^p(X_f,\mathcal F)=0$이다. $X_f\cap X_g=X_{fg}$이므로
(G, II, 5.9.2)에 의해 모든 $p>0$에 대하여
$H^p(X,\mathcal F)=0$도 따른다.
\end{proof}

\begin{corollary}[1.3.2]
\label{III.1.3.2-fr}
$Y$를 준스킴, $f:X\to Y$를 아핀 사상이라 하자 (II, 1.6.1).
임의의 준연접 $\mathcal O_X$-가군 $\mathcal F$와 모든 $q>0$에 대하여
$R^qf_*(\mathcal F)=0$이다.
\end{corollary}

\begin{proof}
실제로 정의에 의해 $R^qf_*(\mathcal F)$는 $Y$의 열린집합 $U$에
$H^q(f^{-1}(U),\mathcal F)$를 대응시키는 준층에 대응하는
$\mathcal O_Y$-가군이다. 그런데 아핀 열린집합들은 $Y$의 기저를 이루며,
그러한 열린집합 $U$에 대하여 $f^{-1}(U)$는 아핀이다 (II, 1.3.2).
따라서 (\hyperref[III.1.3.1-fr]{1.3.1})에 의해
$H^q(f^{-1}(U),\mathcal F)=0$이고, 이로써 따름정리가 증명된다.
\end{proof}

\begin{corollary}[1.3.3]
\label{III.1.3.3-fr}
$Y$를 준스킴, $f:X\to Y$를 아핀 사상이라 하자. 임의의 준연접
$\mathcal O_X$-가군 $\mathcal F$에 대하여 표준 준동형
$H^p(Y,f_*(\mathcal F))\to H^p(X,\mathcal F)$ (0, 12.1.3.1)는 모든
$p$에 대하여 전단사이다.
\end{corollary}

\begin{proof}
실제로 (0, 12.1.7)에 의해 합성함자 $\Gamma f_*$의 두 번째 스펙트럼
열의 가장자리 준동형
$''E_2^{p0}=H^p(Y,f_*(\mathcal F))\to H^p(X,\mathcal F)$들이 전단사임을
보이면 충분하다. 이 열의 $E_2$항은
$''E_2^{pq}=H^p(Y,R^qf_*(\mathcal F))$로 주어지므로 (G, II, 4.17.1),
(\hyperref[III.1.3.2-fr]{1.3.2})에 의해 $q>0$이면
$''E_2^{pq}=0$이다. 따라서 스펙트럼 열은 퇴화하며, 이로부터 주장이
따른다 (0, 11.1.6).
\end{proof}

\begin{corollary}[1.3.4]
\label{III.1.3.4-fr}
$f:X\to Y$를 아핀 사상, $g:Y\to Z$를 사상이라 하자.
\oldpage[III]{89}
임의의 준연접 $\mathcal O_X$-가군 $\mathcal F$에 대하여 표준 준동형
$R^pg_*(f_*(\mathcal F))\to R^p(g\circ f)_*(\mathcal F)$ (0, 12.2.5.1)는
모든 $p$에 대하여 전단사이다.
\end{corollary}

\begin{proof}
(\hyperref[III.1.3.3-fr]{1.3.3})에 의해 $Z$의 임의의 아핀 열린집합
$W$에 대하여 표준 준동형
$H^p(g^{-1}(W),f_*(\mathcal F))\to
H^p(f^{-1}(g^{-1}(W)),\mathcal F)$가 전단사임을 주목하면 충분하다.
이는 표준 준동형
$R^pg_*(f_*(\mathcal F))\to R^p(g\circ f)_*(\mathcal F)$를 정의하는
준층의 준동형이 전단사임을 증명한다 (0, 12.2.5).
\end{proof}

\subsection{임의의 준스킴의 코호몰로지에 대한 응용}
\label{subsection:III.1.4-fr}

\begin{proposition}[1.4.1]
\label{III.1.4.1-fr}
$X$를 스킴, $\mathfrak U=(U_\alpha)$를 아핀 열린집합들로 이루어진
$X$의 덮개라 하자. 임의의 준연접 $\mathcal O_X$-가군 $\mathcal F$에
대하여 ($\Gamma(X,\mathcal O_X)$ 위의) 코호몰로지 가군
$H^\bullet(X,\mathcal F)$와 $H^\bullet(\mathfrak U,\mathcal F)$는
표준적으로 동형이다.
\end{proposition}

\begin{proof}
$X$가 스킴이므로 덮개 $\mathfrak U$의 열린집합들의 임의의 유한 교집합
$V$는 아핀이다 (I, 5.5.6). 따라서
(\hyperref[III.1.3.1-fr]{1.3.1})에 의해 $q\geq1$이면
$H^q(V,\mathcal F)=0$이다. 이제 명제는 르레 정리 (G, II, 5.4.1)에서
따른다.
\end{proof}

\begin{remark}[1.4.2]
\label{III.1.4.2-fr}
$X$가 반드시 스킴이라고 가정하지 않더라도 집합 $U_\alpha$들의 유한
교집합들이 아핀이면 (\hyperref[III.1.4.1-fr]{1.4.1})의 결론이 여전히
성립함을 주목하자.
\end{remark}

\begin{corollary}[1.4.3]
\label{III.1.4.3-fr}
$X$를 바탕공간이 준콤팩트인 스킴,
$A=\Gamma(X,\mathcal O_X)$라 하고,
$\mathbf f=(f_i)_{1\leq i\leq r}$를 $X_{f_i}$들
(\hyperref[III.1.2.1-fr]{1.2.1})의 표기)이 아핀인 $A$의 원소들로
이루어진 유한 수열이라 하자. 그러면
(\hyperref[III.1.2.1-fr]{1.2.1})의 표기 아래 임의의 준연접
$\mathcal O_X$-가군 $\mathcal F$에 대하여 $\mathcal F$에 관해 함자적인
표준 동형사상
\begin{equation}
\label{III.1.4.3.1-fr}
\tag{1.4.3.1}
H^q(U,\mathcal F)\xrightarrow{\sim}H^{q+1}((\mathbf f),M)
\qquad(q\geq1\text{에 대하여})
\end{equation}
과 $\mathcal F$에 관해 함자적인 완전열
\begin{equation}
\label{III.1.4.3.2-fr}
\tag{1.4.3.2}
0\longrightarrow H^0((\mathbf f),M)\longrightarrow M
\longrightarrow H^0(U,\mathcal F)
\longrightarrow H^1((\mathbf f),M)\longrightarrow0
\end{equation}
을 얻는다.
\end{corollary}

\begin{proof}
이는 (\hyperref[III.1.4.1-fr]{1.4.1})과
(\hyperref[III.1.2.3-fr]{1.2.3})에서 곧바로 따른다.
\end{proof}

\begin{env}[1.4.4]
\label{III.1.4.4-fr}
$X$가 \emph{아핀 스킴}이면 (\hyperref[III.1.2.5-fr]{1.2.5})와
(\hyperref[III.1.4.1-fr]{1.4.1})에 의해 모든 $q\geq0$에 대하여,
준연접 $\mathcal O_X$-가군들의 완전열
$0\to\mathcal F'\to\mathcal F\to\mathcal F''\to0$에 대응하는
(\hyperref[III.1.2.5-fr]{1.2.5})의 표기 아래의 도식들
\begin{equation}
\label{III.1.4.4.1-fr}
\tag{1.4.4.1}
\xymatrix{
H^q(U,\mathcal F'')\ar[r]^{\partial}\ar[d]
&H^{q+1}(U,\mathcal F')\ar[d]\\
H^{q+1}((\mathbf f),M'')\ar[r]^{\partial}
&H^{q+2}((\mathbf f),M')
}
\end{equation}
은 (\hyperref[III.1.2.5-fr]{1.2.5})의 표기 아래에서 준연접
$\mathcal O_X$-가군들의 완전열
$0\to\mathcal F'\to\mathcal F\to\mathcal F''\to0$에 대응하며,
이 도표들은 가환이다.
\end{env}

\oldpage[III]{90}
\begin{proposition}[1.4.5]
\label{III.1.4.5-fr}
$X$를 준콤팩트 스킴, $\mathcal L$을 가역 $\mathcal O_X$-가군이라 하고,
등급환 $A_\bullet=\Gamma_\bullet(\mathcal L)$을 생각하자
(0\textsubscript{I}, 5.4.6). 그러면
\[
H^\bullet(\mathcal F,\mathcal L)
=\bigoplus_{n\in\mathbb Z}H^\bullet(X,\mathcal F\otimes\mathcal L^{\otimes n})
\]
은 등급 $A_\bullet$-가군이고, 모든 $f\in A_n$에 대하여
$(A_\bullet)_{(f)}$-가군의 표준 동형사상
\begin{equation}
\label{III.1.4.5.1-fr}
\tag{1.4.5.1}
H^\bullet(X_f,\mathcal F)
\xrightarrow{\sim}(H^\bullet(\mathcal F,\mathcal L))_{(f)}
\end{equation}
이 존재한다.
\end{proposition}

\begin{proof}
$X$가 준콤팩트 스킴이므로, 각 $i$에 대하여 제한
$\mathcal L\vert U_i$가 $\mathcal O_X\vert U_i$와 동형인 아핀
열린집합들의 한 유한 덮개 $\mathfrak U=(U_i)$를 사용하여 모든
$\mathcal O_X$-가군 $\mathcal F\otimes\mathcal L^{\otimes n}$의
코호몰로지를 계산할 수 있다 (\hyperref[III.1.4.1-fr]{1.4.1}). 이때
$U_i\cap X_f$들이 아핀 열린집합임은 즉각적이다 (I, 1.3.6). 따라서 제한
덮개 $\mathfrak U\vert X_f=(U_i\cap X_f)$를 사용하여 코호몰로지
$H^\bullet(X_f,\mathcal F\otimes\mathcal L^{\otimes n})$도 계산할 수 있다
(\hyperref[III.1.4.1-fr]{1.4.1}). 임의의 $f\in A_n$에 대하여 $f$에
의한 곱셈이 준동형
\[
C^\bullet(\mathfrak U,\mathcal F\otimes\mathcal L^{\otimes m})
\longrightarrow
C^\bullet(\mathfrak U,\mathcal F\otimes\mathcal L^{\otimes(m+n)})
\]
을 정의함은 자명하고, 이로부터 준동형
\[
H^\bullet(\mathfrak U,\mathcal F\otimes\mathcal L^{\otimes m})
\longrightarrow
H^\bullet(\mathfrak U,\mathcal F\otimes\mathcal L^{\otimes(m+n)})
\]
을 얻는다. 이로써 첫 번째 주장이 증명된다. 한편 주어진
$f\in A_n$에 대하여 (I, 9.3.2)와 (I, 1.3.9, (ii))를 함께 고려하면
$(A_\bullet)_{(f)}$-가군 복합체의 동형사상
\[
C^\bullet(\mathfrak U\vert X_f,\mathcal F)
\xrightarrow{\sim}
\left(C^\bullet\left(\mathfrak U,
\bigoplus_{n\in\mathbb Z}\mathcal F\otimes\mathcal L^{\otimes n}
\right)\right)_{(f)}
\]
을 얻는다. 이 두 복합체의 코호몰로지를 취하고, 등급
$A_\bullet$-가군의 범주에서 함자 $M\mapsto M_{(f)}$가 완전함을
상기하면 동형사상 (\hyperref[III.1.4.5.1-fr]{1.4.5.1})을 얻는다.
\end{proof}

\begin{corollary}[1.4.6]
\label{III.1.4.6-fr}
(\hyperref[III.1.4.5-fr]{1.4.5})의 가정들이 성립하고, 더 나아가
$\mathcal L=\mathcal O_X$라고 하자. $A=\Gamma(X,\mathcal O_X)$로 놓으면,
모든 $f\in A$에 대하여 $A_f$-가군의 표준 동형사상
\[
H^\bullet(X_f,\mathcal F)
\xrightarrow{\sim}(H^\bullet(X,\mathcal F))_f
\]
이 존재한다.
\end{corollary}

\begin{corollary}[1.4.7]
\label{III.1.4.7-fr}
$X$를 준콤팩트 스킴이라 하고, $f$를 $\Gamma(X,\mathcal O_X)$의 원소라
하자.
\begin{enumerate}
\item[(i)] 열린집합 $X_f$가 아핀이라고 하자. 그러면 임의의 준연접
$\mathcal O_X$-가군 $\mathcal F$, 임의의 $i>0$, 임의의
$\xi\in H^i(X,\mathcal F)$에 대하여 $f^n\xi=0$인 정수 $n>0$이 존재한다.
\item[(ii)] 역으로 $X_f$가 준콤팩트이고, $\mathcal O_X$의 임의의 준연접
아이디얼층 $\mathcal J$와 임의의 $\zeta\in H^1(X,\mathcal J)$에 대하여
$f^n\zeta=0$인 $n>0$이 존재한다고 하자. 그러면 $X_f$는 아핀이다.
\end{enumerate}
\end{corollary}

\begin{proof}
\begin{enumerate}
\item[(i)] $X_f$가 아핀이면 모든 $i>0$에 대하여
$H^i(X_f,\mathcal F)=0$이므로
(\hyperref[III.1.3.1-fr]{1.3.1}), 주장은
(\hyperref[III.1.4.6-fr]{1.4.6})에서 직접 따른다.
\item[(ii)] 세르의 판정법 (II, 5.2.1)에 따라
$\mathcal O_X\vert X_f$의 모든 준연접 아이디얼 $\mathcal K$에 대하여
$H^1(X_f,\mathcal K)=0$임을 보이면 충분하다. $X_f$는 준콤팩트 스킴
$X$ 안의 준콤팩트 열린집합이므로, $\mathcal O_X$의 준연접 아이디얼
$\mathcal J$가 존재하여 $\mathcal K=\mathcal J\vert X_f$이다
(I, 9.4.2). (\hyperref[III.1.4.6-fr]{1.4.6})에 따라
$H^1(X_f,\mathcal K)=(H^1(X,\mathcal J))_f$이고, 가정에 따라 우변은
0이다. 이로부터 결론이 따른다.
\end{enumerate}
\end{proof}

\begin{remark}[1.4.8]
\label{III.1.4.8-fr}
(\hyperref[III.1.4.7-fr]{1.4.7}, (i))가 관계 (II, 4.5.13.2)의 더
간단한 증명을 다시 줌에 유의하자.
\end{remark}

\oldpage[III]{91}
\begin{lemma}[1.4.9]
\label{III.1.4.9-fr}
$X$를 준콤팩트 스킴,
$\mathfrak U=(U_i)_{1\leq i\leq n}$를 $X$의 아핀 열린집합들에 의한
유한 덮개, $\mathcal F$를 준연접 $\mathcal O_X$-가군이라 하자.
그러면 덮개 $\mathfrak U$가 정의하는 층의 복합체
$\mathcal C^\bullet(\mathfrak U,\mathcal F)$ (G, II, 5.2)는 준연접
$\mathcal O_X$-가군이다.
\end{lemma}

\begin{proof}
정의 (G, II, 5.2)에 따라
$\mathcal C^p(\mathfrak U,\mathcal F)$는 표준 포함
$U_{i_0\ldots i_p}\to X$에 의한
$\mathcal F\vert U_{i_0\ldots i_p}$의 직상층들의 직합이다. $X$가
스킴이라는 가정에 따라 이 포함들은 아핀 사상이므로 (I, 5.5.6),
$\mathcal C^p(\mathfrak U,\mathcal F)$들은 준연접이다 (II, 1.2.6).
\end{proof}

\begin{proposition}[1.4.10]
\label{III.1.4.10-fr}
$u:X\to Y$를 분리 준콤팩트 사상이라 하자. 모든 준연접
$\mathcal O_X$-가군 $\mathcal F$에 대하여
$R^qu_*(\mathcal F)$들은 준연접 $\mathcal O_Y$-가군이다.
\end{proposition}

\begin{proof}
문제는 $Y$ 위에서 국소적이므로 $Y$가 아핀이라고 가정할 수 있다.
그러면 $X$는 유한 개의 아핀 열린집합 $U_i$ ($1\leq i\leq n$)의
합집합이다. $\mathfrak U=(U_i)$라 하자. 또한 $Y$가 스킴이므로
(I, 5.5.10)에 따라 모든 아핀 열린집합 $V\subset Y$에 대하여 표준
포함 $u^{-1}(V)\to X$는 아핀 사상이다. 따라서
(\hyperref[III.1.4.1-fr]{1.4.1})과 (G, II, 5.2)에 의해 표준 동형
\begin{equation}
\label{III.1.4.10.1-fr}
\tag{1.4.10.1}
H^\bullet(u^{-1}(V),\mathcal F)
\xrightarrow{\sim}H^\bullet(\Gamma(V,\mathcal K^\bullet))
\end{equation}
이 있다. 여기서
$\mathcal K^\bullet=u_*(\mathcal C^\bullet(\mathfrak U,\mathcal F))$로
놓았다. (\hyperref[III.1.4.9-fr]{1.4.9})와 (I, 9.2.2)에 따라
$\mathcal K^\bullet$은 준연접 $\mathcal O_Y$-가군이다. 한편
$\mathcal C^\bullet(\mathfrak U,\mathcal F)$가 층의 복합체이므로
$\mathcal K^\bullet$도 \emph{층의 복합체}이다. 따라서 코호몰로지
$\mathcal H^\bullet(\mathcal K^\bullet)$의 정의 (G, II, 4.1)에 따라
이는 준연접 $\mathcal O_Y$-가군들로 이루어진다 (I, 4.1.1). $Y$의
아핀 열린집합 $V$에 대하여 함자 $\Gamma(V,\mathcal G)$는 준연접
$\mathcal O_Y$-가군의 범주에서 $\mathcal G$에 관하여 완전하므로,
(G, II, 4.1)에 따라
\begin{equation}
\label{III.1.4.10.2-fr}
\tag{1.4.10.2}
H^\bullet(\Gamma(V,\mathcal K^\bullet))
=\Gamma(V,\mathcal H^\bullet(\mathcal K^\bullet)).
\end{equation}
끝으로 (G, II, 5.2)에서 준 표준 준동형
\[
H^\bullet(\mathfrak U,\mathcal F)\longrightarrow H^\bullet(X,\mathcal F)
\]
의 정의에 따라, $V'\subset V$가 $Y$의 또 다른 아핀 열린집합이면
도식
\[
\xymatrix{
H^\bullet(u^{-1}(V),\mathcal F)\ar[r]^{\sim}\ar[d]
&H^\bullet(\Gamma(V,\mathcal K^\bullet))\ar[d]\\
H^\bullet(u^{-1}(V'),\mathcal F)\ar[r]^{\sim}
&H^\bullet(\Gamma(V',\mathcal K^\bullet))
}
\]
은 가환한다. 따라서 앞의 결과들로부터 동형
(\hyperref[III.1.4.10.1-fr]{1.4.10.1})들이 $\mathcal O_Y$-가군의 동형
\begin{equation}
\label{III.1.4.10.3-fr}
\tag{1.4.10.3}
R^\bullet u_*(\mathcal F)
\xrightarrow{\sim}\mathcal H^\bullet(\mathcal K^\bullet)
\end{equation}
을 정의함을 얻는다.
\oldpage[III]{92}
따라서 $R^\bullet u_*(\mathcal F)$는 준연접이다.
\end{proof}

또한 (\hyperref[III.1.4.10.3-fr]{1.4.10.3}),
(\hyperref[III.1.4.10.2-fr]{1.4.10.2}),
(\hyperref[III.1.4.10.1-fr]{1.4.10.1})에서 다음을 얻는다.

\begin{corollary}[1.4.11]
\label{III.1.4.11-fr}
(\hyperref[III.1.4.10-fr]{1.4.10})의 가정 아래, $Y$의 모든 아핀
열린집합 $V$와 모든 $q\geq0$에 대하여 표준 준동형
\begin{equation}
\label{III.1.4.11.1-fr}
\tag{1.4.11.1}
H^q(u^{-1}(V),\mathcal F)\longrightarrow
\Gamma(V,R^qu_*(\mathcal F))
\end{equation}
은 동형이다.
\end{corollary}

\begin{corollary}[1.4.12]
\label{III.1.4.12-fr}
(\hyperref[III.1.4.10-fr]{1.4.10})의 가정들이 성립하고, 더 나아가
$Y$가 준콤팩트라고 하자. 그러면 정수 $r>0$이 존재하여 모든 준연접
$\mathcal O_X$-가군 $\mathcal F$와 모든 정수 $q>r$에 대하여
$R^qu_*(\mathcal F)=0$이다. $Y$가 아핀이면, $X$가 $r$개의 아핀
열린집합으로 이루어진 덮개를 갖도록 $r$을 취할 수 있다.
\end{corollary}

\begin{proof}
$Y$를 유한 개의 아핀 열린집합으로 덮을 수 있으므로
(\hyperref[III.1.4.11-fr]{1.4.11})에 의해 둘째 주장만 증명하면 된다.
$\mathfrak U$가 $X$의 $r$개의 아핀 열린집합에 의한 덮개이면
$q>r$일 때 $H^q(\mathfrak U,\mathcal F)=0$이다. 실제로
$C^q(\mathfrak U,\mathcal F)$의 코사슬은 교대적이다. 따라서 결론은
(\hyperref[III.1.4.1-fr]{1.4.1})에서 따른다.
\end{proof}

\begin{corollary}[1.4.13]
\label{III.1.4.13-fr}
(\hyperref[III.1.4.10-fr]{1.4.10})의 가정들이 성립하고, 더 나아가
$Y=\operatorname{Spec}(A)$가 아핀이라고 하자. 그러면 모든 준연접
$\mathcal O_X$-가군 $\mathcal F$와 모든 $f\in A$에 대하여 표준 동형을
제외하면
\[
\Gamma(Y_f,R^qu_*(\mathcal F))
=(\Gamma(Y,R^qu_*(\mathcal F)))_f
\]
이다.
\end{corollary}

\begin{proof}
이는 $R^qu_*(\mathcal F)$가 준연접 $\mathcal O_Y$-가군이라는 사실에서
따른다 (I, 1.3.7).
\end{proof}

\begin{proposition}[1.4.14]
\label{III.1.4.14-fr}
$f:X\to Y$를 분리 준콤팩트 사상, $g:Y\to Z$를 아핀 사상이라 하자.
모든 준연접 $\mathcal O_X$-가군 $\mathcal F$와 모든 $p$에 대하여
표준 준동형
$R^p(g\circ f)_*(\mathcal F)\to g_*(R^pf_*(\mathcal F))$
(0, 12.2.5.2)은 전단사이다.
\end{proposition}

\begin{proof}
실제로 $Z$의 모든 아핀 열린집합 $W$에 대하여 $g^{-1}(W)$는 $Y$의
아핀 열린집합이다. 그러므로 표준 준동형
\[
R^p(g\circ f)_*(\mathcal F)\longrightarrow g_*(R^pf_*(\mathcal F))
\]
을 정의하는 전층의 준동형 (0, 12.2.5)은
(\hyperref[III.1.4.11-fr]{1.4.11})에 따라 전단사이다.
\end{proof}

\begin{proposition}[1.4.15]
\label{III.1.4.15-fr}
$u:X\to Y$를 분리 준콤팩트 사상, $v:Y'\to Y$를 준스킴의 평탄 사상이라
하자 (0\textsubscript{I}, 6.7.1). $u'=u_{(Y')}$라 하여 가환도식
\begin{equation}
\label{III.1.4.15.1-fr}
\tag{1.4.15.1}
\xymatrix{
X\ar[d]_u & X'=X_{(Y')}\ar[l]_{v'}\ar[d]^{u'}\\
Y & Y'\ar[l]_v
}
\end{equation}
을 얻는다. 모든 준연접 $\mathcal O_X$-가군 $\mathcal F$와 모든
$q\geq0$에 대하여
$\mathcal F'=v'^*(\mathcal F)=\mathcal F\otimes_{\mathcal O_Y}\mathcal O_{Y'}$
로 놓으면, $R^qu_*'(\mathcal F')$은
$R^qu_*(\mathcal F)\otimes_{\mathcal O_Y}\mathcal O_{Y'}
=v^*(R^qu_*(\mathcal F))$와 표준적으로 동형이다.
\end{proposition}

\oldpage[III]{93}
\begin{proof}
표준 준동형
$\rho:\mathcal F\to v_*'(v'^*(\mathcal F))$
(0\textsubscript{I}, 4.4.3.2)는 함자성에 따라 준동형
\begin{equation}
\label{III.1.4.15.2-fr}
\tag{1.4.15.2}
R^qu_*(\mathcal F)\longrightarrow R^qu_*(v_*'(\mathcal F'))
\end{equation}
을 정의한다. 한편 $w=u\circ v'=v\circ u'$로 놓으면 표준 준동형들
(0, 12.2.5.1 및 12.2.5.2)
\begin{equation}
\label{III.1.4.15.3-fr}
\tag{1.4.15.3}
R^qu_*(v_*'(\mathcal F'))\longrightarrow R^qw_*(\mathcal F')
\longrightarrow v_*(R^qu_*'(\mathcal F'))
\end{equation}
이 있다. (\hyperref[III.1.4.15.3-fr]{1.4.15.3})과
(\hyperref[III.1.4.15.2-fr]{1.4.15.2})를 합성하여 준동형
\[
\psi:R^qu_*(\mathcal F)\longrightarrow v_*(R^qu_*'(\mathcal F'))
\]
를 얻고, 끝으로 표준 준동형
\begin{equation}
\label{III.1.4.15.4-fr}
\tag{1.4.15.4}
\psi^\sharp:v^*(R^qu_*(\mathcal F))\longrightarrow R^qu_*'(\mathcal F')
\end{equation}
을 얻는다. 이 정의에는 \emph{$v$에 관한 어떠한 가정도} 들어가지
않는다. $v$가 \emph{평탄}일 때 이것이 동형임을 증명해야 한다.
문제는 $Y$와 $Y'$ 위에서 국소적이므로
$Y=\operatorname{Spec}(A)$, $Y'=\operatorname{Spec}(B)$라고 가정할 수
있다. 다음 보조정리도 사용한다.

\begin{lemma}[1.4.15.5]
\label{III.1.4.15.5-fr}
$\varphi:A\to B$를 환 준동형, $Y=\operatorname{Spec}(A)$,
$X=\operatorname{Spec}(B)$, $f:X\to Y$를 $\varphi$에 대응하는 사상,
$M$을 $B$-가군이라 하자. $\mathcal O_X$-가군 $\widetilde M$이
$f$-평탄이기 위한 필요충분조건은 $M$이 평탄 $A$-가군인 것이다
(0\textsubscript{I}, 6.7.1). 특히 사상 $f$가 평탄이기 위한
필요충분조건은 $B$가 평탄 $A$-가군인 것이다.
\end{lemma}

실제로 이는 정의 (0\textsubscript{I}, 6.7.1)와
(0\textsubscript{I}, 6.3.3)에서 따르며, 이때 (I, 1.3.4)를 고려한다.

이제 (\hyperref[III.1.4.11.1-fr]{1.4.11.1})과 준동형들
(\hyperref[III.1.4.15.3-fr]{1.4.15.3})의 정의(0, 12.2.5 참조)에
따라 $\psi$는 합성사상
\[
H^q(X,\mathcal F)\xrightarrow{\rho_q}
H^q(X,v_*'(v'^*(\mathcal F)))\xrightarrow{\theta_q}
H^q(X',v'^*(v_*'(v'^*(\mathcal F))))\xrightarrow{\sigma_q}
H^q(X',v'^*(\mathcal F))
\]
에 대응한다. 여기서 $\rho_q$와 $\sigma_q$는 표준 준동형 $\rho$와
$\sigma:v'^*(v_*'(\mathcal G'))\to\mathcal G'$가 코호몰로지 위에
정의하는 준동형이고, $\theta_q$는
$\mathcal O_X$-가군 $v_*'(v'^*(\mathcal F))$에 관한
$\varphi$-사상 (0, 12.1.3.1)이다. $\theta_q$의 함자성에 따라 도식
\[
\xymatrix{
H^q(X,\mathcal F)\ar[rr]^{\rho_q}\ar[dd]_{\theta_q}
&&H^q(X,v_*'(v'^*(\mathcal F)))\ar[dd]^{\theta_q}\\\\
H^q(X',v'^*(\mathcal F))\ar[rr]^{v'^*(\rho_q)}
&&H^q(X',v'^*(v_*'(v'^*(\mathcal F))))
}
\]
은 가환한다.
\oldpage[III]{94}
정의 (0\textsubscript{I}, 4.4.3)에 따라 $v'^*(\rho)$는
$\sigma$의 역이므로, 앞의 합성사상은 결국 $\theta_q$이다. 따라서
$\psi^\sharp$은 연관된 $B$-준동형
\[
H^q(X,\mathcal F)\otimes_A B\longrightarrow H^q(X',\mathcal F')
\]
이다. $u$가 준콤팩트이므로 $X$는 유한 개의 아핀 열린집합 $U_i$
($1\leq i\leq r$)의 합집합이다. $\mathfrak U=(U_i)$라 하자. 한편
$v$가 아핀 사상이므로 $v'$도 아핀 사상이고 (II, 1.6.2, (iii)),
$U_i'=v'^{-1}(U_i)$들은 $X'$의 아핀 열린 덮개 $\mathfrak U'$를 이룬다.
이때 (0, 12.1.4.2)에 따라 도식
\[
\xymatrix{
H^q(\mathfrak U,\mathcal F)\ar[r]^{\theta_q}\ar[d]
&H^q(\mathfrak U',\mathcal F')\ar[d]\\
H^q(X,\mathcal F)\ar[r]^{\theta_q}&H^q(X',\mathcal F')
}
\]
은 가환하고, $X$와 $X'$이 스킴이므로 수직 화살표들은 동형이다
(\hyperref[III.1.4.1-fr]{1.4.1}). 따라서 표준 $\varphi$-사상
$\theta_q:H^q(\mathfrak U,\mathcal F)\to
H^q(\mathfrak U',\mathcal F')$에 연관된 $B$-준동형
\[
H^q(\mathfrak U,\mathcal F)\otimes_A B
\longrightarrow H^q(\mathfrak U',\mathcal F')
\]
이 동형임을 보이면 충분하다. $[1,r]$의 $p+1$개 지표로 이루어진 임의의
열 $\mathbf s=(i_k)_{0\leq k\leq p}$에 대하여
\[
U_{\mathbf s}=\bigcap_{k=0}^pU_{i_k},\qquad
U_{\mathbf s}'=\bigcap_{k=0}^pU_{i_k}'=v'^{-1}(U_{\mathbf s}),\qquad
M_{\mathbf s}=\Gamma(U_{\mathbf s},\mathcal F),\qquad
M_{\mathbf s}'=\Gamma(U_{\mathbf s}',\mathcal F')
\]
로 놓자. 표준 사상 $M_{\mathbf s}\otimes_A B\to M_{\mathbf s}'$은
동형이다 (I, 1.6.5). 따라서 표준 사상
$C^p(\mathfrak U,\mathcal F)\otimes_A B\to
C^p(\mathfrak U',\mathcal F')$도 동형이고, 이 동형 아래
$d\otimes1$은 코경계 연산자
$C^p(\mathfrak U',\mathcal F')\to C^{p+1}(\mathfrak U',\mathcal F')$와
동일시된다. $B$가 \emph{평탄} $A$-가군이므로, 코호몰로지 가군의
정의에서 표준 사상
$H^q(\mathfrak U,\mathcal F)\otimes_A B\to
H^q(\mathfrak U',\mathcal F')$이 동형임이 곧바로 따른다
(0\textsubscript{I}, 6.1.1). 이 결과는 뒤에서 일반화할 것이다
(\S~6).
\end{proof}

\begin{corollary}[1.4.16]
\label{III.1.4.16-fr}
$A$를 환, $X$를 유한형 $A$-스킴, $B$를 $A$ 위의 충실 평탄
$A$-대수라 하자. $X$가 아핀이기 위한 필요충분조건은
$X\otimes_A B$가 아핀인 것이다.
\end{corollary}

\begin{proof}
조건의 필요성은 명백하다 (I, 3.2.2). 충분성을 보이자. $X$는 $A$ 위에서
분리이고 사상 $\operatorname{Spec}(B)\to\operatorname{Spec}(A)$는
평탄하므로, (\hyperref[III.1.4.15-fr]{1.4.15})에 따라 모든 $i\geq0$과
모든 준연접 $\mathcal O_X$-가군 $\mathcal F$에 대하여
\begin{equation}
\label{III.1.4.16.1-fr}
\tag{1.4.16.1}
H^i(X\otimes_A B,\mathcal F\otimes_A B)
=H^i(X,\mathcal F)\otimes_A B
\end{equation}
은 모든 $i\geq0$과 모든 준연접 $\mathcal O_X$-가군 $\mathcal F$에
대하여 성립한다. $X\otimes_A B$가 아핀이면
(\hyperref[III.1.4.16.1-fr]{1.4.16.1})의 좌변은 $i=1$일 때 0이다.
$B$가 충실 평탄 $A$-가군이므로 $H^1(X,\mathcal F)=0$도 성립한다.
$X$가 준콤팩트 스킴이므로 세르의 판정법 (II, 5.2.1)으로 결론을 얻는다.
\end{proof}

\begin{proposition}[1.4.17]
\label{III.1.4.17-fr}
$X$를 준스킴,
$0\to\mathcal F\xrightarrow{u}\mathcal G\xrightarrow{v}\mathcal H\to0$
을 $\mathcal O_X$-가군들의 완전열이라 하자. $\mathcal F$와
$\mathcal H$가 준연접이면 $\mathcal G$도 준연접이다.
\end{proposition}

\oldpage[III]{95}
\begin{proof}
문제는 $X$ 위에서 국소적이므로 $X=\operatorname{Spec}(A)$가
아핀이라고 가정할 수 있다. 그러면 $\mathcal G$가 (I, 1.4.1)의 조건
d~1)과 d~2)($V=X$)를 만족함을 보이면 충분하다. d~2)의 확인은
즉시 된다. 실제로 $t\in\Gamma(X,\mathcal G)$의 $D(f)$에 대한 제한이
0이면, 그 상 $v(t)\in\Gamma(X,\mathcal H)$의 제한도 0이다. 따라서
어떤 $m>0$에 대하여 $f^mv(t)=v(f^mt)=0$이다 (I, 1.4.1).
$\Gamma$가 왼쪽 완전하므로 $f^mt=u(s)$인
$s\in\Gamma(X,\mathcal F)$가 존재한다. $u$가 단사이므로 $s$의
$D(f)$에 대한 제한은 0이고, 이에 따라 (I, 1.4.1) 어떤 정수 $n>0$에
대하여 $f^ns=0$이다. 끝으로
$f^{m+n}t=u(f^ns)=0$을 얻는다.

이제 d~1)을 확인하자. $t'\in\Gamma(D(f),\mathcal G)$라 하자.
$\mathcal H$가 준연접이므로 어떤 정수 $m$에 대하여
$f^mv(t')=v(f^mt')$는 어떤 단면 $z\in\Gamma(X,\mathcal H)$로
연장된다 (I, 1.4.1). 그런데 준연접 $\mathcal O_X$-가군
$\mathcal F$에 (\hyperref[III.1.3.1-fr]{1.3.1}) 또는
(I, 5.1.9.2)를 적용하면 열
$\Gamma(X,\mathcal G)\to\Gamma(X,\mathcal H)\to0$은 완전하다.
따라서 $z=v(t)$인 $t\in\Gamma(X,\mathcal G)$가 존재한다.
$t''$를 $t$의 $D(f)$에 대한 제한이라 하면
$v(f^mt'-t'')=0$이다. 그러므로
$f^mt'-t''=u(s')$인 $s'\in\Gamma(D(f),\mathcal F)$가 존재한다.
$\mathcal F$가 준연접이므로 어떤 정수 $n>0$에 대하여 $f^ns'$은
어떤 단면 $s\in\Gamma(X,\mathcal F)$로 연장된다. 관계
$f^{m+n}t'-f^nt''=u(f^ns')$에 따라 $f^{m+n}t'$는 단면
$f^nt+u(s)\in\Gamma(X,\mathcal G)$의 $D(f)$에 대한 제한이다. 이로써
증명이 끝난다.
\end{proof}

\section{사영 사상의 코호몰로지적 연구}
\label{section:III.2-fr}

\subsection{몇몇 코호몰로지 군의 명시적 계산}
\label{subsection:III.2.1-fr}

\begin{env}[2.1.1]
\label{III.2.1.1-fr}
$X$를 준스킴, $\mathcal L$을 가역 $\mathcal O_X$-가군이라 하고, 등급환
(0\textsubscript{I}, 5.4.6)
\begin{equation}
\label{III.2.1.1.1-fr}
\tag{2.1.1.1}
S=\Gamma_\bullet(X,\mathcal L)
=\bigoplus_{n\in\mathbb Z}\Gamma(X,\mathcal L^{\otimes n})
\end{equation}
을 생각하자. $(f_i)_{1\leq i\leq r}$를 $S$의 \emph{동차} 원소들의
유한족이라 하고 $f_i\in S_{d_i}$라 하자. $U_i=X_{f_i}$,
$U=\bigcup_iU_i$라 두고, $\mathfrak U$로 $U$의 덮개 $(U_i)$를
나타내자. 모든 준연접 $\mathcal O_X$-가군 $\mathcal F$에 대하여
다음과 같이 둔다.
\begin{align}
\label{III.2.1.1.2-fr}
\tag{2.1.1.2}
H^\bullet(U,\mathcal F(*))
&=\bigoplus_{n\in\mathbb Z}
H^\bullet(U,\mathcal F\otimes\mathcal L^{\otimes n}),\\
\label{III.2.1.1.3-fr}
\tag{2.1.1.3}
H^\bullet(\mathfrak U,\mathcal F(*))
&=\bigoplus_{n\in\mathbb Z}
H^\bullet(\mathfrak U,\mathcal F\otimes\mathcal L^{\otimes n}).
\end{align}
아벨군 (\hyperref[III.2.1.1.2-fr]{2.1.1.2})와
(\hyperref[III.2.1.1.3-fr]{2.1.1.3})은 다음과 같이 \emph{이중등급}을
갖는다.
\[
(H^\bullet(U,\mathcal F(*)))_{mn}
=H^m(U,\mathcal F\otimes\mathcal L^{\otimes n}).
\]
(\hyperref[III.2.1.1.3-fr]{2.1.1.3})도 마찬가지로 정의한다. 둘째
차수에 관하여 이 군들이 등급 $S$-가군임은 명백하다. 예를 들어 이는
$\mathcal F\mapsto H^m(U,\mathcal F)$와
$\mathcal F\mapsto H^m(\mathfrak U,\mathcal F)$가 함자라는 사실에서
따른다.
\end{env}

\begin{env}[2.1.2]
\label{III.2.1.2-fr}
이제 등급 $S$-가군 (0\textsubscript{I}, 5.4.6)
\begin{equation}
\label{III.2.1.2.1-fr}
\tag{2.1.2.1}
M=\Gamma_\bullet(\mathcal L,\mathcal F)
=H^0(X,\mathcal F(*))
=\bigoplus_{n\in\mathbb Z}
\Gamma(X,\mathcal F\otimes\mathcal L^{\otimes n})
\end{equation}
을 생각하자.
\oldpage[III]{96}
$X$가 바탕공간이 뇌터인 준스킴이거나 준콤팩트 스킴이면, 통상과 같이
$U_{i_0i_1\ldots i_p}=\bigcap_{k=0}^pU_{i_k}$로 둘 때 (I, 9.3.1)로부터
표준 동형을 제외하면
\[
\Gamma(U_{i_0i_1\ldots i_p},\mathcal F(*))
=H^0(U_{i_0i_1\ldots i_p},\mathcal F(*))
=M_{f_{i_0}f_{i_1}\cdots f_{i_p}}
\]
이다. 또한 (\hyperref[III.1.2.2-fr]{1.2.2})의 표기로
$M_{f_{i_0}f_{i_1}\cdots f_{i_p}}$를
$\varinjlim_n M_{i_0i_1\ldots i_p}^{(n)}$와 식별할 수 있다. 동차 원소
$z\in\varinjlim_n M_{i_0i_1\ldots i_p}^{(n)}$의 차수를 다음과 같이
정의하면 이 식별은 등급 $S$-가군의 동형사상이다. $z$가 차수 $m$인
동차 원소
$x\in M_{i_0i_1\ldots i_p}^{(n)}=M$의 표준 상이라고 하자. 이때
$z$의 차수를
$m-n(d_{i_0}+d_{i_1}+\cdots+d_{i_p})$로 둔다. 준동형
$\varphi_{kn}:M_{i_0i_1\ldots i_p}^{(n)}\to
M_{i_0i_1\ldots i_p}^{(k)}$의 정의
(\hyperref[III.1.2.2-fr]{1.2.2})를 고려하면, 이 정의가 택한 $z$의
``대표원'' $x$에 의존하지 않음을 즉시 알 수 있다.
(\hyperref[III.1.2.2-fr]{1.2.2})에서와 같이 $C_n^p(M)$으로
$[1,r]^{p+1}$에서 $M$으로 가는 교대사상들의 집합을 나타내자(모든
$n$에 대하여). 위와 같은 방식으로 $\varinjlim_n C_n^p(M)$에 등급
$S$-가군 구조를 정의한다. 다시
(\hyperref[III.1.2.2-fr]{1.2.2})에서와 같이
\begin{equation}
\label{III.2.1.2.2-fr}
\tag{2.1.2.2}
C^p(\mathfrak U,\mathcal F(*))=\varinjlim_n C_n^p(M)
\end{equation}
이고, 두 변의 동형사상은 \emph{차수를 보존한다}. 그러면
(\hyperref[III.1.2.2-fr]{1.2.2})에서와 같이
\begin{equation}
\label{III.2.1.2.3-fr}
\tag{2.1.2.3}
C^p(\mathfrak U,\mathcal F(*))
=C^{p+1}((\mathbf f),M)
=\varinjlim_n K^{p+1}(\mathbf f^n,M)
\end{equation}
이고, 이 동형사상도 차수를 보존한다. 즉 같은 동차 성분 $M_m$에 값
$g(i_0,\ldots,i_p)$들을 갖는 공사슬
$g\in K^{p+1}(\mathbf f^n,M)$의 표준 상으로 주어진
$\varinjlim_n K^{p+1}(\mathbf f^n,M)$의 원소의 차수는
$m-n(d_{i_0}+\cdots+d_{i_p})$이며, 이 원소의 대표원으로 택한 공사슬에
의존하지 않는다.
\end{env}

앞의 동형사상들이 공경계 연산자들과 양립하므로,
(\hyperref[III.1.2.2-fr]{1.2.2})에서와 같이 다음을 얻는다.

\begin{proposition}[2.1.3]
\label{III.2.1.3-fr}
$X$를 바탕공간이 뇌터인 준스킴이거나 준콤팩트 스킴이라 하자.
$\mathcal F$에 관해 함자적인 표준 동형
\begin{equation}
\label{III.2.1.3.1-fr}
\tag{2.1.3.1}
H^p(\mathfrak U,\mathcal F(*))
\xrightarrow{\sim}H^{p+1}((\mathbf f),M)
\qquad\text{모든 }p\geq1\text{에 대하여}
\end{equation}
이 존재한다. 또한 $\mathcal F$에 관해 함자적인 완전열
\begin{equation}
\label{III.2.1.3.2-fr}
\tag{2.1.3.2}
0\longrightarrow H^0((\mathbf f),M)\longrightarrow M
\longrightarrow H^0(\mathfrak U,\mathcal F(*))
\longrightarrow H^1((\mathbf f),M)\longrightarrow0
\end{equation}
이 있다. 더 나아가, 도입한 모든 준동형은 등급 $S$-가군 구조에 관해
차수 $0$이다. 여기서 $S$는 환
(\hyperref[III.2.1.1.1-fr]{2.1.1.1})이다.
\end{proposition}

\providecommand{\bbGamma}{\boldsymbol{\Gamma}}
\providecommand{\shProj}{\mathcal{P}\mathrm{roj}}

\oldpage[III]{97}
\begin{corollary}[2.1.4]
\label{III.2.1.4-fr}
$X$가 준콤팩트 스킴이고 $U_i=X_{f_i}$들이 아핀이면,
$\mathcal F$에 관해 함자적인 차수 $0$의 표준 동형
\begin{equation}
\label{III.2.1.4.1-fr}
\tag{2.1.4.1}
H^p(U,\mathcal F(*))\xrightarrow{\sim}H^{p+1}((\mathbf f),M)
\qquad(p\geq1\text{에 대하여})
\end{equation}
과 $\mathcal F$에 관해 함자적인 완전열
\begin{equation}
\label{III.2.1.4.2-fr}
\tag{2.1.4.2}
0\longrightarrow H^0((\mathbf f),M)\longrightarrow M
\longrightarrow H^0(U,\mathcal F(*))
\longrightarrow H^1((\mathbf f),M)\longrightarrow0
\end{equation}
이 존재하며, 모든 준동형은 차수 $0$이다.
\end{corollary}

\begin{proof}
실제로 (1.4.1)을 (\hyperref[III.2.1.3-fr]{2.1.3})의 결과에 적용하면
충분하다.
\end{proof}

(\hyperref[III.2.1.3-fr]{2.1.3})과 유사한 ``국소'' 명제는 다음과 같다.

\begin{proposition}[2.1.5]
\label{III.2.1.5-fr}
$S$를 양의 차수를 갖는 등급환, $f_i$ ($1\leq i\leq r$)를 차수
$d_i$인 $S_+$의 동차 원소, $M$을 등급 $S$-가군이라 하자.
$X=\operatorname{Proj}(S)$를 $S$의 동차 소 스펙트럼이라 하고
$U_i=D_+(f_i)$, $U=\bigcup_iU_i$로 두며,
\[
H^\bullet(U,\widetilde M(*))
=\bigoplus_{n\in\mathbb Z}H^\bullet(U,(M(n))^\sim)
\]
로 두자. 그러면 등급 $S$-가군 구조에 관해 차수 $0$이고 $M$에 관해
함자적인 표준 동형
\begin{equation}
\label{III.2.1.5.1-fr}
\tag{2.1.5.1}
H^p(U,\widetilde M(*))\xrightarrow{\sim}H^{p+1}((\mathbf f),M)
\qquad\text{모든 }p\geq1\text{에 대하여}
\end{equation}
들과 $M$에 관해 함자적인 완전열
\begin{equation}
\label{III.2.1.5.2-fr}
\tag{2.1.5.2}
0\longrightarrow H^0((\mathbf f),M)\longrightarrow M
\longrightarrow H^0(U,\widetilde M(*))
\longrightarrow H^1((\mathbf f),M)\longrightarrow0
\end{equation}
이 존재하며, 모든 준동형은 차수 $0$이다.
\end{proposition}

\begin{proof}
실제로 정의에 의해 (II, 2.5.2)
\[
\Gamma(U_{i_0i_1\ldots i_p},(M(n))^\sim)
=(M_{f_{i_0}\cdots f_{i_p}})_n
\]
이므로
$\Gamma(U_{i_0i_1\ldots i_p},\widetilde M(*))
=M_{f_{i_0}\cdots f_{i_p}}$이다. 그러면 $X$가 스킴임을 고려할 때,
나머지 논증은 (\hyperref[III.2.1.4-fr]{2.1.4})의 증명과 같다.
\end{proof}

\begin{env}[2.1.6]
\label{III.2.1.6-fr}
\emph{비고.} ---
(i) (\hyperref[III.2.1.5-fr]{2.1.5})의 조건 아래에서
(0\textsubscript{I}, 1.3.2)에 의해 함자
$\Gamma(U_{i_0i_1\ldots i_p},\widetilde M(*))$는 $M$에 관해 완전하다.
따라서 (1.2.5)에서와 같은 논증으로, 등급 $S$-가군들의 완전열
$0\to M'\to M\to M''\to0$(준동형들은 차수 $0$)이 주어지면 모든
$p\geq0$에 대하여 가환 도식
\begin{equation}
\label{III.2.1.6.1-fr}
\tag{2.1.6.1}
\xymatrix{
H^p(U,\widetilde {M''}(*))\ar[r]^-\partial\ar[d] &
H^{p+1}(U,\widetilde {M'}(*))\ar[d]\\
H^{p+1}((\mathbf f),M'')\ar[r]^-\partial &
H^{p+2}((\mathbf f),M')
}
\end{equation}
을 얻는다.

(ii) 명제 (\hyperref[III.2.1.5-fr]{2.1.5})는 특히 $A$가 뇌터이고
$S$가 차수 $1$의 유한 개 원소로 생성되는 $A$-대수일 때 중요하다.
\oldpage[III]{98}
실제로 이 경우 모든 준연접 $\mathcal O_X$-가군은 $\widetilde M$의
꼴이다 (II, 2.7.7).
\end{env}

\begin{env}[2.1.7]
\label{III.2.1.7-fr}
$A$가 임의의 환이고 $T_i$들이 부정원일 때,
$S=A[T_0,\ldots,T_r]$, $M=S$, $f_i=T_i$인 경우에
(\hyperref[III.2.1.5-fr]{2.1.5})를 적용하겠다. 따라서 본질적으로
$\mathbf T=(T_i)_{0\leq i\leq r}$일 때
$H^\bullet((\mathbf T),S)$를 계산하는 문제로 귀착된다.
\end{env}

\begin{lemma}[2.1.8]
\label{III.2.1.8-fr}
$S=A[T_0,\ldots,T_r]$이고 $\mathbf T=(T_i)_{0\leq i\leq r}$이면
\begin{align}
\label{III.2.1.8.1-fr}
\tag{2.1.8.1}
H^i(\mathbf T^n,S)&=0 &&\text{$i\neq r+1$이면},\\
\label{III.2.1.8.2-fr}
\tag{2.1.8.2}
H^{r+1}(\mathbf T^n,S)&=S/(\mathbf T^n)
\end{align}
이다. 따라서 $A$-가군 $H^{r+1}(\mathbf T^n,S)$는 다음 단항식들의
$(\mathbf T^n)$을 법으로 한 유들로 이루어진 $A$ 위의 기저를 갖는다.
\[
T_{\mathbf p}=T_0^{p_0}\cdots T_r^{p_r},
\qquad \mathbf p=(p_0,\ldots,p_r),\quad 0\leq p_i<n
\quad\text{모든 }i\text{에 대하여}.
\]
\end{lemma}

\begin{proof}
가정들이 자명하게 성립하는 (1.1.3.5)와 명제 (1.1.4)의 즉각적인
귀결이다.
\end{proof}

\begin{env}[2.1.9]
\label{III.2.1.9-fr}
$n$에 관해 귀납극한을 취하자. 관계
(\hyperref[III.2.1.8.1-fr]{2.1.8.1})로부터 $i\neq r+1$이면
$H^i((\mathbf T),S)=0$임을 얻는다. $i=r+1$일 때 귀납계는
$S/(\mathbf T^n)$들로 이루어지고, 준동형
\[
\varphi_{nm}:S/(\mathbf T^n)\longrightarrow S/(\mathbf T^m)
\qquad(0\leq n\leq m)
\]
은 $(T_0\cdots T_r)^{m-n}$을 곱하는 사상이다.
$n\geq\sup(p_i)_{0\leq i\leq r}$일 때,
$\xi_{\mathbf p}^{(n)}=\xi_{p_0,\ldots,p_r}^{(n)}$로
$T_0^{n-p_0}\cdots T_r^{n-p_r}\pmod{(\mathbf T^n)}$의 유를 나타내자.
그러면 $\varphi_{nm}(\xi_{\mathbf p}^{(n)})
=\xi_{\mathbf p}^{(m)}$이고, 따라서 이 원소들은 귀납극한
$H^{r+1}((\mathbf T),S)$에서 같은 표준 상
$\xi_{\mathbf p}=\xi_{p_0,\ldots,p_r}$를 갖는다.
(\hyperref[III.2.1.2-fr]{2.1.2})에서 준 차수의 정의에 따라
$\xi_{\mathbf p}$의 차수는
$-|\mathbf p|=-(p_0+p_1+\cdots+p_r)$이다. $0<p_i\leq n$이고
$0\leq i\leq r$인 $\xi_{\mathbf p}^{(n)}$들이
$S/(\mathbf T^n)$의 기저를 이룸은 명백하다. 따라서
(\hyperref[III.2.1.8-fr]{2.1.8})로부터 즉시 다음을 얻는다.
\end{env}

\begin{corollary}[2.1.10]
\label{III.2.1.10-fr}
(\hyperref[III.2.1.8-fr]{2.1.8})의 표기 아래에서
\begin{equation}
\label{III.2.1.10.1-fr}
\tag{2.1.10.1}
H^i((\mathbf T),S)=0\qquad\text{$i\neq r+1$이면}
\end{equation}
이고, $H^{r+1}((\mathbf T),S)$는 $0\leq i\leq r$에 대하여
$p_i>0$인 원소 $\xi_{p_0,\ldots,p_r}$들로 이루어진 기저를 갖는 자유
$A$-가군이다.
\end{corollary}

\begin{env}[2.1.11]
\label{III.2.1.11-fr}
\emph{비고.} --- $N$을 임의의 $A$-가군이라 하고 $M=S\otimes_A N$이라
하자. (\hyperref[III.2.1.8-fr]{2.1.8})의 논증으로부터 더 일반적으로
\begin{align}
\label{III.2.1.11.1-fr}
\tag{2.1.11.1}
H^i(\mathbf T^n,M)&=0 &&\text{$i\neq r+1$이면},\\
\label{III.2.1.11.2-fr}
\tag{2.1.11.2}
H^{r+1}(\mathbf T^n,M)&=(S/(\mathbf T^n))\otimes_A N
\end{align}
을 얻는다. 실제로 마지막 공식은 (1.1.3.5)에서 직접 따르고, 한편
\[
M/(T_0^nM+\cdots+T_{i-1}^nM)
=\bigl(S/(T_0^nS+\cdots+T_{i-1}^nS)\bigr)\otimes_A N
\]
임이 명백하다. 아이디얼 $T_0^nS+\cdots+T_{i-1}^nS$는 $A$-가군
$S$의 직합인자이므로 (1.1.4)를 $M$에 적용할 수 있고, 이로써
(\hyperref[III.2.1.11.1-fr]{2.1.11.1})을 얻는다.
\end{env}

(\hyperref[III.2.1.10-fr]{2.1.10})과
(\hyperref[III.2.1.5-fr]{2.1.5})를 결합하면 다음을 얻는다.

\begin{proposition}[2.1.12]
\label{III.2.1.12-fr}
$A$를 환, $r$을 양의 정수, $X=\mathbf P_A^r$라 하자 (II, 4.1.1).
그러면 다음이 성립한다.
\begin{enumerate}
\item[{\rm(i)}] $i\neq0,r$이면 $H^i(X,\mathcal O_X(*))=0$이다.
\item[{\rm(ii)}] 표준 준동형
$\alpha:S\to H^0(X,\mathcal O_X(*))$ (II, 2.6.2)는 전단사이다.
\oldpage[III]{99}
\item[{\rm(iii)}] $H^r(X,\mathcal O_X(*))$는 $0\leq i\leq r$에 대해
$p_i>0$인 원소 $\xi_{p_0,\ldots,p_r}$들로 이루어진 기저를 갖는 자유
$A$-가군이다. 여기서 $\xi_{p_0,\ldots,p_r}$의 차수는
$-|\mathbf p|=-(p_0+\cdots+p_r)$이고,
$T_i\xi_{p_0,\ldots,p_r}=\xi_{p_0,\ldots,p_i-1,\ldots,p_r}$이다.
\end{enumerate}
\end{proposition}

\begin{proof}
실제로 $M=S=A[T_0,\ldots,T_r]$에 적용한 완전열
(\hyperref[III.2.1.5.2-fr]{2.1.5.2})에서
(\hyperref[III.2.1.10.1-fr]{2.1.10.1})에 의해
$H^0((\mathbf T),S)=0$이고 $H^1((\mathbf T),S)=0$임을 유의하자.
또한 $X$는 $D_+(T_i)$들의 합집합이므로 (II, 2.3.14), 명제
(\hyperref[III.2.1.5-fr]{2.1.5})를 $U=X$에 적용할 수 있다. 이제 완전열
(\hyperref[III.2.1.5.2-fr]{2.1.5.2})의 사상
$S\to H^0(X,\mathcal O_X(*))$을 표준 사상 $\alpha$와 식별하면 된다.
이는 $H^0(U,\mathcal O_X(*))$와
$H^0(\mathfrak U,\mathcal O_X(*))$의 표준 식별에서 따른다.
\end{proof}

\begin{corollary}[2.1.13]
\label{III.2.1.13-fr}
$H^i(X,\mathcal O_X(n))\neq0$일 수 있는 $(i,n)$의 값은 다음뿐이다.
$i=0$, $n\geq0$인 경우와 $i=r$, $n\leq-(r+1)$인 경우.
\end{corollary}

$A\neq0$이면 (\hyperref[III.2.1.13-fr]{2.1.13})에 열거한 쌍들에 대해
실제로 $H^i(X,\mathcal O_X(n))\neq0$임을 유의하자. 이는 모든
$n\geq0$에 대해 $S_n\neq0$이므로
(\hyperref[III.2.1.12-fr]{2.1.12})에서 따른다.

이 장에서의 응용에는 주로 다음의 덜 정밀한 결과를 사용하겠다.

\begin{corollary}[2.1.14]
\label{III.2.1.14-fr}
$A$-가군 $H^i(X,\mathcal O_X(n))$은 유한 생성 자유 가군이다. $i>0$이면
$n>0$에 대하여 0이다.
\end{corollary}

\begin{proposition}[2.1.15]
\label{III.2.1.15-fr}
$Y$를 준스킴, $\mathcal E$를 계수 $r+1$인 국소 자유
$\mathcal O_Y$-가군, $X=\mathbf P(\mathcal E)$를 $\mathcal E$로 정의되는
사영 다발, $f:X\to Y$를 구조 사상이라 하자.
$R^if_*(\mathcal O_X(n))\neq0$일 수 있는 $i$, $n$의 값은 $i=0$,
$n\geq0$인 경우와 $i=r$, $n\leq-(r+1)$인 경우뿐이다. 또한 표준
준동형 (II, 3.3.2)
\[
\alpha:S_{\mathcal O_Y}(\mathcal E)
\longrightarrow\Gamma_*(\mathcal O_X)
=R^0f_*(\mathcal O_X(*))
=\bigoplus_{n\in\mathbb Z}f_*(\mathcal O_X(n))
\]
은 동형사상이다.
\end{proposition}

\begin{proof}
문제는 $Y$ 위에서 국소적이므로, $Y$가 환 $A$의 아핀이고
$\mathcal E=\widetilde E$, $E=A^{r+1}$이라 가정할 수 있다. 그러면
(1.4.11)을 고려하여 곧바로 (\hyperref[III.2.1.12-fr]{2.1.12})로
귀착된다.
\end{proof}

\begin{env}[2.1.16]
\label{III.2.1.16-fr}
\emph{비고.} --- 나중에 (\hyperref[III.2.1.15-fr]{2.1.15})의 결과를
다음 명제들을 증명함으로써 보완하겠다. 다음과 같이 놓자.
\[
\boldsymbol\omega=f^*\left(\bigwedge^{r+1}\mathcal E\right)(-r-1).
\]
이는 가역 $\mathcal O_X$-가군이다. 그러면 다음이 성립한다.
\begin{enumerate}
\item[{\rm(i)}] 표준 동형사상
\begin{equation}
\label{III.2.1.16.1-fr}
\tag{2.1.16.1}
\rho:R^rf_*(\boldsymbol\omega)\xleftarrow{\sim}\mathcal O_Y
\end{equation}
이 존재한다.
\item[{\rm(ii)}] 컵곱 (0, 12.2.2)에 의한 짝지음
\begin{equation}
\label{III.2.1.16.2-fr}
\tag{2.1.16.2}
R^rf_*(\mathcal O_X(n))\times R^0f_*(\boldsymbol\omega(-n))
\longrightarrow R^rf_*(\boldsymbol\omega)
\end{equation}
을 동형사상 $\rho^{-1}$과 합성하면,
$R^rf_*(\mathcal O_X(n))$에서 국소 자유 $\mathcal O_Y$-가군
\[
R^0f_*(\boldsymbol\omega(-n))
=\left(\bigwedge^{r+1}\mathcal E\right)
\otimes_{\mathcal O_Y}\bigl(S_{\mathcal O_Y}(\mathcal E)\bigr)_{-n}
\]
의 \emph{쌍대}로 가는 \emph{동형사상}이 정의된다.
\end{enumerate}
\end{env}

\subsection{사영 사상의 기본정리}
\label{subsection:III.2.2-fr}
\oldpage[III]{100}

\begin{theorem}[2.2.1]
\label{III.2.2.1-fr}
(세르). $Y$를 뇌터 준스킴, $f:X\to Y$를 고유 사상,
$\mathcal L$을 $f$에 대해 풍부한 가역 $\mathcal O_X$-가군이라 하자.
모든 $\mathcal O_X$-가군 $\mathcal F$에 대하여
\[
\mathcal F(n)=\mathcal F\otimes_{\mathcal O_X}\mathcal L^{\otimes n}
\qquad\text{모든 }n\in\mathbb Z\text{에 대하여}
\]
로 놓는다. 그러면 모든 연접 $\mathcal O_X$-가군 $\mathcal F$에 대해
다음이 성립한다.
\begin{enumerate}
\item[{\rm(i)}] $R^qf_*(\mathcal F)$들은 연접 $\mathcal O_Y$-가군이다.
\item[{\rm(ii)}] 어떤 정수 $N$이 존재하여 $n\geq N$이면
\[
R^qf_*(\mathcal F(n))=0\qquad\text{모든 }q>0\text{에 대하여}
\]
이다.
\item[{\rm(iii)}] 어떤 정수 $N$이 존재하여 $n\geq N$이면 표준 준동형
\[
f^*(f_*(\mathcal F(n)))\longrightarrow\mathcal F(n)
\]
이 전사이다.
\end{enumerate}
\end{theorem}

\begin{proof}
먼저 정리에서 $\mathcal L$을 $\mathcal L^{\otimes d}$ $(d>0)$로
바꾸었을 때 정리가 참이면 원래 형태로도 참임을 유의하자. 실제로
\[
\mathcal F(n)
=(\mathcal F\otimes\mathcal L^{\otimes r})
\otimes\mathcal L^{\otimes hd}
\]
로 쓸 수 있는데, 여기서 $h>0$, $0\leq r<d$이다. 가정에 따라 각
$r$에 대하여 어떤 정수 $N_r$이 존재하여 $h\geq N_r$이면
$\mathcal O_X$-가군 $\mathcal F\otimes\mathcal L^{\otimes r}$에 대해
(ii), (iii)이 성립한다. $dN_r$들 가운데 가장 큰 것을 $N$으로 택하면
$n\geq N$에 대해 (ii), (iii)이 성립한다. 따라서 $\mathcal L$이
$f$에 대해 매우 풍부하다고 가정할 수 있다 (II, 4.6.11). 그러므로
어떤 우세한 $Y$-열린 몰입 $i:X\to P$가 존재한다. 여기서
$P=\operatorname{Proj}(\mathcal S)$이고, $\mathcal S$는 양의 차수들을
갖는 등급 준연접 $\mathcal O_Y$-대수이며, $\mathcal S_1$은 유한
생성이고 $\mathcal S$를 생성한다. 또한 $\mathcal L$은
$i^*(\mathcal O_P(1))$과 동형이다 (II, 4.4.7). 그런데 $f$가 고유이므로
$i$도 고유이고 (II, 5.4.4), 따라서 $i$는 동형사상
$X\xrightarrow{\sim}P$이다. 그러므로
$X=\operatorname{Proj}(\mathcal S)$, $\mathcal L=\mathcal O_X(1)$인
경우로 한정할 수 있다. 이제 정리 (2.2.1)은 다음 명제의 귀결이다.
\end{proof}

\begin{proposition}[2.2.2]
\label{III.2.2.2-fr}
$A$를 뇌터 환, $S$를 양의 차수들을 갖는 등급 $A$-대수라 하자.
$S_1$은 $r+1$개의 생성원계를 갖는 $A$-가군이고 대수 $S$를 생성한다고
하자. $X=\operatorname{Proj}(S)$로 놓는다. 모든 연접
$\mathcal O_X$-가군 $\mathcal F$에 대해 다음이 성립한다.
\begin{enumerate}
\item[{\rm(i)}] $A$-가군 $H^q(X,\mathcal F)$들은 유한 생성이다.
\item[{\rm(ii)}] $q>r$이면 $H^q(X,\mathcal F)=0$이다.
\item[{\rm(iii)}] 어떤 정수 $N$이 존재하여 $n\geq N$이면 모든
$q>0$에 대해 $H^q(X,\mathcal F(n))=0$이다.
\item[{\rm(iv)}] 어떤 정수 $N$이 존재하여 $n\geq N$이면
$\mathcal F(n)$은 $X$ 위의 단면들로 생성된다.
\end{enumerate}
\end{proposition}

\begin{proof}
먼저 (2.2.2)가 (2.2.1)을 함의함을 보이자. 위에서 고려한 특수한 경우
$X=\operatorname{Proj}(\mathcal S)$로 귀착된 (2.2.1)에서 $Y$는
준콤팩트이므로, 각 환이 뇌터인 유한 개의 아핀 열린집합들로 덮인다.
각 열린집합 $U_\alpha$ 위에서 $\mathcal S_1$의 제한은
$U_\alpha$ 위의 유한 개 단면들로 생성된다. (2.2.2)가 증명되었다고
가정하면, (1.4.11)과 (II, 3.4.7)을 고려하여 (2.2.1)의 (ii), (iii)에
나타나는 $N$으로 $U_\alpha$들에 대응하는 정수들 가운데 가장 큰 것을
택하면 충분하다.

(2.2.2)를 증명하기 위해 $X$가 $P=\mathbf P_A^r$의 닫힌 부분스킴과
동일시됨을 유의하자 (II, 3.6.2). 또한 $j:X\to P$가 표준 포함이면
$j_*(\mathcal F)$는 연접 $\mathcal O_P$-가군이고
$j_*(\mathcal F(n))=(j_*(\mathcal F))(n)$이다 (II, 3.4.5 및 3.5.2).
(G, II, 정리 4.9.1의 따름정리)를 고려하면, 따라서
\[
X=\mathbf P_A^r\quad\text{이고}\quad S=A[T_0,\ldots,T_r]
\]
인 특수한 경우에 (2.2.2)를 증명하는 문제로 귀착된다.

\oldpage[III]{101}
$X$는 $r+1$개의 아핀 열린집합 $D_+(T_i)$로 덮이므로 (ii)는
(1.4.12)에서 따른다. 한편 (iv)는 이미 증명되었다 (II, 2.7.9).

(i)와 (iii)을 동시에 증명하겠다. 이 명제들은
$\mathcal F=\mathcal O_X(m)$일 때 참이다
(\hyperref[III.2.1.13-fr]{2.1.13}). 따라서 $\mathcal F$가
$\mathcal O_X(m_j)$ 꼴의 유한 개 $\mathcal O_X$-가군들의 직합일 때도
참이다. 한편 (ii)에 의해 $q>r$이면 (i), (iii)은 자명하게 참이다.
$q$에 대한 \emph{하강 귀납법}을 쓰겠다. $\mathcal F$는 유한 개의
층 $\mathcal O_X(m_j)$의 직합 $\mathcal E$의 몫과 동형임을 안다
(II, 2.7.10). 다시 말해 완전열
\[
0\longrightarrow\mathcal R\longrightarrow\mathcal E
\longrightarrow\mathcal F\longrightarrow0
\]
이 존재한다. 여기서 $\mathcal R$은 연접이고 (0\textsubscript{I}, 5.3.3),
$\mathcal E$는 (i), (iii)을 만족한다. $\mathcal F(n)$은
$\mathcal F$에 관해 완전한 함자이므로, 모든 $n\in\mathbb Z$에 대해
완전열
\[
0\longrightarrow\mathcal R(n)\longrightarrow\mathcal E(n)
\longrightarrow\mathcal F(n)\longrightarrow0
\]
도 존재한다. 이로부터 코호몰로지 완전열
\[
H^{q-1}(X,\mathcal E(n))\longrightarrow
H^{q-1}(X,\mathcal F(n))\longrightarrow
H^q(X,\mathcal R(n))
\]
을 얻는다. $\mathcal E(n)$은 $\mathcal O_X(n+m_j)$들의 직합이므로
(II, 2.5.14), $H^{q-1}(X,\mathcal E(n))$은 유한 생성이다. 귀납가정에
의해 $H^q(X,\mathcal R(n))$도 유한 생성이다. $A$가 뇌터이므로
모든 $n\in\mathbb Z$에 대해 $H^{q-1}(X,\mathcal F(n))$이 유한
생성이고, 특히 $n=0$일 때 그러함을 얻는다. 한편 귀납가정에 의해
어떤 정수 $N$이 존재하여 $n\geq N$이면
$H^q(X,\mathcal R(n))=0$이다. 또한 $\mathcal E$가 (iii)을 만족하므로,
$n\geq N$이면 $H^{q-1}(X,\mathcal E(n))=0$이 되도록 $N$을 택할 수도
있다. 따라서 $n\geq N$이면
$H^{q-1}(X,\mathcal F(n))=0$이고, 이로써 증명이 끝난다.
\end{proof}

\begin{corollary}[2.2.3]
\label{III.2.2.3-fr}
(\hyperref[III.2.2.1-fr]{2.2.1})의 가정 아래에서
$\mathcal F\to\mathcal G\to\mathcal H$를 연접
$\mathcal O_X$-가군들의 완전열이라 하자. 그러면 어떤 정수 $N$이
존재하여 $n\geq N$이면 열
\[
f_*(\mathcal F(n))\longrightarrow f_*(\mathcal G(n))
\longrightarrow f_*(\mathcal H(n))
\]
이 완전하다.
\end{corollary}

\begin{proof}
$\mathcal F\to\mathcal G$의 핵, 상, 여핵을 각각
$\mathcal F'$, $\mathcal G'$, $\mathcal G''$라 하자.
$\mathcal G'$은 $\mathcal G\to\mathcal H$의 핵이고
$\mathcal G''$는 그 상이다. 이 준동형의 여핵을 $\mathcal H''$라 하자.
이 모든 $\mathcal O_X$-가군은 연접이다 (0\textsubscript{I}, 5.3.4).
$\mathcal F(n)$은 $\mathcal F$에 관해 완전한 함자이므로, 충분히 큰
$n$에 대하여 다음 열들이 각각 완전함을 증명하면 충분하다.
\begin{align*}
0&\longrightarrow f_*(\mathcal F'(n))
\longrightarrow f_*(\mathcal F(n))
\longrightarrow f_*(\mathcal G'(n))\longrightarrow0,\\
0&\longrightarrow f_*(\mathcal G'(n))
\longrightarrow f_*(\mathcal G(n))
\longrightarrow f_*(\mathcal G''(n))\longrightarrow0,\\
0&\longrightarrow f_*(\mathcal G''(n))
\longrightarrow f_*(\mathcal H(n))
\longrightarrow f_*(\mathcal H''(n))\longrightarrow0.
\end{align*}
따라서 $0\to\mathcal F\to\mathcal G\to\mathcal H\to0$이
완전하다고 가정할 수 있다. 이때 코호몰로지 완전열
\[
0\longrightarrow f_*(\mathcal F(n))
\longrightarrow f_*(\mathcal G(n))
\longrightarrow f_*(\mathcal H(n))
\longrightarrow R^1f_*(\mathcal F(n))\longrightarrow\cdots
\]
을 얻고, 결론은 (\hyperref[III.2.2.1-fr]{2.2.1}, (ii))에서 따른다.
\end{proof}

\begin{corollary}[2.2.4]
\label{III.2.2.4-fr}
$Y$를 뇌터 준스킴, $f:X\to Y$를 유한형 사상,
$\mathcal L$을 $f$에 대해 풍부한 가역 $\mathcal O_X$-가군이라 하자.
모든 $\mathcal O_X$-가군 $\mathcal F$에 대해
$\mathcal F(n)=\mathcal F\otimes_{\mathcal O_X}\mathcal L^{\otimes n}$
$(n\in\mathbb Z)$으로 놓는다. $\mathcal F\to\mathcal G\to\mathcal H$가
연접 $\mathcal O_X$-가군들의 완전열이고,
$\operatorname{Im}(\mathcal F\to\mathcal G)$의 받침이 $Y$ 위에서
고유라고 하자 (II, 5.4.10). 그러면 어떤 정수 $N$이 존재하여
$n\geq N$이면 열
\[
f_*(\mathcal F(n))\longrightarrow f_*(\mathcal G(n))
\longrightarrow f_*(\mathcal H(n))
\]
이 완전하다.
\end{corollary}

\begin{proof}
(\hyperref[III.2.2.1-fr]{2.2.1})의 증명 첫머리와 같은 논증에 따르면,
$\mathcal L^{\otimes d}$ $(d>0)$에 대해 따름정리가 참이면
$\mathcal L$에 대해서도 참이다. 따라서 $\mathcal L$이 $f$에 대해
매우 풍부한 경우로 한정할 수 있다 (II, 4.6.11). 그러므로 $X$를
$Y$-스킴 $Z=\operatorname{Proj}(\mathcal S)$의 열린집합과 동일시할 수
있다. 여기서 $\mathcal S$는 양의 차수들을 갖는 등급 준연접
$\mathcal O_Y$-대수이고, $\mathcal S_1$은 유한 생성이며
$\mathcal S$를 생성한다. 또한 $i:X\to Z$가 표준 몰입이면
$\mathcal L=i^*(\mathcal O_Z(1))$이다 (II, 4.4.7).
$\mathcal G_1=\operatorname{Im}(\mathcal F\to\mathcal G)
=\operatorname{Ker}(\mathcal G\to\mathcal H)$로 놓으면 이는 연접이다.
완전열 $0\to\mathcal G_1(n)\to\mathcal G(n)\to\mathcal H(n)$이 있고
함자 $f_*$가 왼쪽 완전하므로, 모든 $n$에 대해 열
$0\to f_*(\mathcal G_1(n))\to f_*(\mathcal G(n))
\to f_*(\mathcal H(n))$은 완전하다. 따라서 충분히 큰 $n$에 대하여
$f_*(\mathcal F(n))\to f_*(\mathcal G_1(n))\to0$이 완전함을 보이면
충분하다. 다시 말해 $\mathcal H=0$이고 $\mathcal G$의 받침이
$Y$ 위에서 고유인 경우로 한정할 수 있다. 이 받침은 따라서
$Z$에서 \emph{닫혀 있다} (II, 5.4.10). 그러므로
$\mathcal G'=i_*(\mathcal G)$는 $\mathcal G'|X=\mathcal G$인 연접
$\mathcal O_Z$-가군이다. $\mathcal F'|X=\mathcal F$인 연접
$\mathcal O_Z$-가군 $\mathcal F'$가 존재함을 안다 (I, 9.4.3).
또한 $X$가 $Z$에서 열려 있고
$\operatorname{Supp}(\mathcal G)\subset X$가 $Z$에서 닫혀 있으므로,
$\operatorname{Supp}(\mathcal G)$와 만나지 않는 $Z$의 모든 열린집합
$U$에서는 $u'|U=0$으로, 모든 열린집합 $U\subset X$에서는
$u'|U=u|U$로 놓으면, 주어진 전사 준동형
$u:\mathcal F\to\mathcal G$로 제한되는 전사 층 준동형
$u':\mathcal F'\to\mathcal G'$을 정의함이 곧바로 확인된다.

이제 (2.2.2)의 증명 첫머리에서와 같이 $Y$가 아핀인 경우로 한정할
수 있다. 따라서 충분히 큰 $n$에 대해 준동형
$\Gamma(u):\Gamma(X,\mathcal F(n))\to\Gamma(X,\mathcal G(n))$이
전사임을 보이면 된다. 그런데 가환 도식
\[
\xymatrix@C=35mm@R=12mm{
\Gamma(Z,\mathcal F'(n)) \ar[r]^{\Gamma(u')} \ar[d]_v &
\Gamma(Z,\mathcal G'(n)) \ar[d]^w \\
\Gamma(X,\mathcal F(n)) \ar[r]_{\Gamma(u)} & \Gamma(X,\mathcal G(n))
}
\]
이 있다. 여기서 $v$, $w$는 제한 준동형들이다. $\mathcal G'$의 정의에
의해 $w$는 \emph{전단사}이다. 한편 (2.2.3)에 의해 충분히 큰 $n$에
대해 $\Gamma(u')$은 \emph{전사}이므로, $\Gamma(u)$도 전사이다.
\end{proof}

\begin{env}[2.2.5]
\label{III.2.2.5-fr}
\emph{비고.} ---
(i) (\hyperref[III.2.2.1-fr]{2.2.1})의 명제 (i)는 $Y$가 단지
\emph{국소 뇌터}라고 가정해도 여전히 참이다. 실제로 이 성질은
$Y$ 위에서 명백히 국소적이다. 한편
(\hyperref[III.2.2.1-fr]{2.2.1})의 가정들은 모든 열린집합
$U\subset Y$에 대해 $f$를 $f^{-1}(U)$에 제한한 것이 사영 사상
$f^{-1}(U)\to U$이고 (II, 5.5.5, (iii)),
$\mathcal L|f^{-1}(U)$가 이 사상에 대해 풍부함을 함의한다
(II, 4.6.4).

(ii) 앞에서 보았듯이, (\hyperref[III.2.2.1-fr]{2.2.1})의 명제 (iii)은
$X$가 준콤팩트 스킴이거나 그 바탕공간이 뇌터인 준스킴이고
$f:X\to Y$가 준콤팩트 사상이라고만 가정해도 여전히 성립한다
(II, 4.6.8). 그러나 $Y$가 체 $K$의 \emph{스펙트럼}이고 $f$가
\emph{준사영}이라고 가정해도
(\hyperref[III.2.2.1-fr]{2.2.1})의 명제 (ii)는 더 이상 반드시
성립하지 않음에 유의해야 한다. 예를 들어
$X'=\operatorname{Spec}(K[T_0,\ldots,T_r])$라 하고 $X$를 $X'$의 아핀
열린집합 $D(T_i)$ $(0\leq i\leq r)$들의 합집합이라 하자.
몰입 $X\to X'$가 준콤팩트이므로 구조 사상 $f:X\to Y$는 준아핀이고
(II, 5.1.10), 따라서 $\mathcal O_X$는 $f$에 대해 매우 풍부하다
(II, 5.1.6). 그런데 환 $\Gamma(X,\mathcal O_X)$는
$0\leq i\leq r$에 대한 분수환
$(K[T_0,\ldots,T_r])_{T_i}$들의 교집합과 동일시되며
(I, 8.2.1.1), 이는 $K[T_0,\ldots,T_r]$이다. 따라서 공식 (1.4.3.1)과
(1.1.3.5)로부터 모든 $n$에 대해
$H^r(X,\mathcal O_X^{\otimes n})=H^r(X,\mathcal O_X)\neq0$임을 얻는다.

(iii) $X$, $Y$, $f$, $\mathcal L$에 관한 (2.2.4)의 가정 아래에서,
$\mathcal H$가 받침이 \emph{$Y$ 위에서 고유}인 연접
$\mathcal O_X$-가군이면 (2.2.4)와 유사한 논증에 의해 어떤 정수 $N$이
존재하여 $n\geq N$이면 모든 $i>0$에 대해
$\mathrm R^if_*(\mathcal H(n))=0$임은 즉시 확인된다. 코호몰로지
완전열로부터 다음 결론을 얻는다. $u:\mathcal F\to\mathcal G$가 연접
$\mathcal O_X$-가군들의 준동형이고 $\operatorname{Ker}(u)$와
$\operatorname{Coker}(u)$의 받침들이 \emph{$Y$ 위에서 고유}이면,
어떤 $N$이 존재하여 $n\geq N$일 때 대응하는 준동형
$\mathrm R^if_*(\mathcal F(n))\to\mathrm R^if_*(\mathcal G(n))$은
모든 $i>0$에 대해 \emph{전단사}이다.
\end{env}

\subsection{등급 대수 및 가군의 층에 대한 응용}
\label{subsection:III.2.3-fr}

\begin{theorem}[2.3.1]
\label{III.2.3.1-fr}
$Y$를 뇌터 준스킴, $\mathcal S$를 양의 차수들을 갖는 준연접 유한 생성
등급 $\mathcal O_Y$-대수, $X=\operatorname{Proj}(\mathcal S)$,
$q:X\to Y$를 구조 사상, $\mathcal M$을 조건 \textbf{(TF)}를 만족하는
준연접 등급 $\mathcal S$-가군이라 하자. 그러면 어떤 정수 $N$이
존재하여 $n\geq N$이면 표준 준동형 (II, 8.14.5.1)
\[
\alpha_n:\mathcal M_n\longrightarrow
q_*(\shProj_0(\mathcal M(n)))
=q_*((\shProj(\mathcal M))_n)
\]
\oldpage[III]{103}
이 전단사이다. 다시 말해 표준 준동형
\[
\alpha:\mathcal M\longrightarrow\bbGamma_*(\shProj(\mathcal M))
\]
은 \textbf{(TN)}-동형사상이다.
\end{theorem}

\begin{proof}
$\mathcal M$이 유한 생성 $\mathcal S$-가군인 경우로 한정할 수 있다
(II, 3.4.2). $Y$가 준콤팩트이므로, $\mathcal S^{(d)}$가 준연접
$\mathcal O_Y$-가군 $\mathcal S_d$로 생성되도록 하는 어떤 정수
$d>0$이 존재한다. $\mathcal S_d$는 유한 생성이고 (II, 3.1.10),
$Y$가 뇌터이므로 연접이다. 이제 $\mathcal M$은
$0\leq k<d$에 대한 $\mathcal M^{(d,k)}$들의 직합이며, 문제는
$Y$ 위에서 국소적이므로 (II, 2.1.6, (iii))에 따라 각
$\mathcal M^{(d,k)}$는 유한 생성 준연접 등급
$\mathcal S^{(d)}$-가군이다. 표준 준동형들
\[
\alpha:\mathcal M^{(d,k)}
\longrightarrow\bbGamma_*((\shProj(\mathcal M))^{(d,k)})
\]
이 각각 \textbf{(TN)}-동형사상임을 증명하면 충분함은 명백하다.
(II, 8.14.13), 특히 도식 (8.14.13.4)을 고려하면,
$\mathcal S$가 $\mathcal S_1$로 생성되고 $\mathcal S_1$이 연접
$\mathcal O_Y$-가군인 경우에 정리를 증명하는 문제로 귀착됨을 알 수
있다.

$Y$가 뇌터이므로, (\hyperref[III.2.2.2-fr]{2.2.2})의 증명 첫머리와
같은 논증으로 $Y=\operatorname{Spec}(A)$,
$\mathcal S=\widetilde S$, $\mathcal M=\widetilde M$인 경우로 한정할 수
있다. 여기서 $A$는 뇌터 환, $S_1$은 유한 생성 $A$-가군이고 $M$은
유한 생성 등급 $S$-가군이다. 이때 $M=S$인 경우에 정리를 증명하면
충분함을 보이자. 실제로 일반적인 경우에는 완전열
\[
L'\longrightarrow L\longrightarrow M\longrightarrow0
\]
이 존재하며, 여기서 $L$, $L'$은 $S(m)$ 꼴의 등급가군들의 직합이다.
$M=S$일 때 결과가 참이면 $M=S(m)$일 때도 참이고, 따라서 $L$과
$L'$에 대해서도 참이다. 이제 가환 도식
\[
\xymatrix{
\widetilde {L'_n}\ar[r]\ar[d]_{\alpha_n} &
\widetilde {L_n}\ar[r]\ar[d]^{\alpha_n} &
\widetilde {M_n}\ar[r]\ar[d]^{\alpha_n} & 0\\
q_*(\widetilde {L'}(n))\ar[r] &
q_*(\widetilde L(n))\ar[r] &
q_*(\widetilde M(n))\ar[r] & 0 .
}
\]
을 생각하자. 충분히 큰 $n$에 대해 둘째 줄은
(\hyperref[III.2.2.3-fr]{2.2.3})에 의해 완전하다. 첫째 줄도 그러하며
왼쪽의 두 수직 화살표가 동형사상이므로 셋째 화살표도 동형사상이다.

이제 $M=S$일 때 정리를 증명하기 위해 먼저
$S=A[T_0,\ldots,T_r]$($T_i$들은 부정원)라고 가정하자. 이 경우 우리
명제는 (\hyperref[III.2.1.11-fr]{2.1.11}, (ii))에 지나지 않는다.
일반적인 경우 $S$는 등급 아이디얼에 의한 환
$S'=A[T_0,\ldots,T_r]$의 몫과 동일시되며, 따라서 $X$는
$X'=\mathbf P_A^r$의 닫힌 부분스킴과 동일시된다 (II, 2.9.2).
$j$가 표준 포함 $X\to X'$이면, $j_*(\widetilde S(n))$는 $S$를 등급
$S'$-가군으로 보았을 때의 $\mathcal O_{X'}$-가군
$(\shProj(\widetilde S))(n)$에 지나지 않는다. 이는 (II, 2.8.7)에서
곧바로 따른다. $j_*(\widetilde S(n))$는 조건 \textbf{(TF)}를 만족하는
$\mathcal O_{X'}$-가군이므로, 앞에서 증명한 바에 따라 충분히 큰
$n$에 대해 표준 준동형
$\alpha_n:S_n\to\Gamma(X',j_*(\widetilde S(n)))$은 전단사이다.
그런데
$\Gamma(X',j_*(\widetilde S(n)))=\Gamma(X,\widetilde S(n))$이므로
증명이 끝난다.
\end{proof}

\oldpage[III]{104}
\begin{corollary}[2.3.2]
\label{III.2.3.2-fr}
(\hyperref[III.2.3.1-fr]{2.3.1})의 가정 아래에서
\[
\mathcal S_X=\bigoplus_{n\in\mathbb Z}\mathcal O_X(n)
\]
이라 하고, $\mathcal F$를 유한 생성 준연접 등급
$\mathcal S_X$-가군이라 하자. 그러면 $\bbGamma_*(\mathcal F)$는 조건
\textbf{(TF)}를 만족한다.
\end{corollary}

\begin{proof}
(\hyperref[III.2.3.1-fr]{2.3.1})의 증명에서
$X$가 $\operatorname{Proj}(\mathcal S^{(d)})$와 동형임을 보았다
(II, 3.1.8). 따라서 $X$는 $Y$ 위에서 유한형이다 (II, 3.4.1).
그러면 (II, 8.14.9)에 의해 $\mathcal F$는
$\shProj(\mathcal M)$ 꼴의 등급 $\mathcal S_X$-가군과 동형이다.
여기서 $\mathcal M$은 유한 생성 준연접 등급 $\mathcal S$-가군이다.
(\hyperref[III.2.3.1-fr]{2.3.1})에 의해
$\bbGamma_*(\mathcal F)$는 $\mathcal M$과 \textbf{(TN)}-동형이고,
따라서 \textbf{(TF)}를 만족한다.
\end{proof}

\begin{env}[2.3.3]
\label{III.2.3.3-fr}
\emph{보주.} --- $Y$를 뇌터 준스킴, $\mathcal S$를
(\hyperref[III.2.3.1-fr]{2.3.1})의 조건들을 만족하는 등급
$\mathcal O_Y$-대수, $X=\operatorname{Proj}(\mathcal S)$라 하자.
조건 \textbf{(TF)}를 만족하는 준연접 등급 $\mathcal S$-가군들의
아벨 범주를 $K_{\mathcal S}$라 하고, 그 가운데 조건 \textbf{(TN)}을
만족하는 $\mathcal S$-가군들로 이루어진 부분범주를
$K'_{\mathcal S}$라 하자. 마지막으로, 유한 생성 준연접 등급
$\mathcal S_X$-가군 $\mathcal F$들의 범주를 $K_X$라 하자. 이는
$\mathcal S_X$가 주기적이므로 (II, 8.14.4 및 8.14.12),
$\mathcal F_i$들이 연접 $\mathcal O_X$-가군이라고 말하는 것과 같다.
그러면 함자
\[
\mathcal M\longmapsto\shProj(\mathcal M)
\quad\text{$K_{\mathcal S}$에서},
\qquad
\mathcal F\longmapsto\bbGamma_*(\mathcal F)
\quad\text{$K_X$에서}
\]
는 (II, 8.14.8 및 8.14.10)과
(\hyperref[III.2.3.2-fr]{2.3.2})에 의해 몫범주
$K_{\mathcal S}/K'_{\mathcal S}$ (T, I, 1.11)와 범주 $K_X$ 사이의
동치 (T, I, 1.2)를 정의한다. $\mathcal S$가 $\mathcal S_1$로
생성되면 $K_X$를 연접 $\mathcal O_X$-가군들의 범주로 바꿀 수 있다
(II, 8.14.12).
\end{env}

\begin{proposition}[2.3.4]
\label{III.2.3.4-fr}
$Y$를 뇌터 준스킴이라 하자.
\begin{enumerate}
\item[{\rm(i)}] $\mathcal S$를 양의 차수들을 갖는 준연접 유한 생성
등급 $\mathcal O_Y$-대수라 하자. $X=\operatorname{Proj}(\mathcal S)$,
\[
\mathcal S_X=\shProj(\mathcal S)
=\bigoplus_{n\in\mathbb Z}\mathcal O_X(n)
\]
으로 놓는다. 그러면 $\mathcal S_X$는 동차 성분
$(\mathcal S_X)_n=\mathcal O_X(n)$들이 연접 $\mathcal O_X$-가군인
주기적 등급 $\mathcal O_X$-대수이다 (II, 8.14.12). 또한 $d>0$이
$\mathcal S_X$의 주기이면 $(\mathcal S_X)_d=\mathcal O_X(d)$는
$Y$-풍부한 가역 $\mathcal O_X$-가군이다. 더 나아가 표준 준동형
\[
\alpha:\mathcal S\longrightarrow\bbGamma_*(\mathcal S_X)
\]
은 \textbf{(TN)}-동형사상이다.
\item[{\rm(ii)}] 역으로, $q:X\to Y$를 사영 사상이라 하고,
$\mathcal S'$을 동차 성분 $\mathcal S'_n$ $(n\in\mathbb Z)$들이 연접
$\mathcal O_X$-가군인 등급 $\mathcal O_X$-대수라 하자. 또한
$\mathcal S'_d$가 $q$에 대해 풍부한 가역 $\mathcal O_X$-가군이
되도록 하는 주기 $d>0$을 갖는다고 하자. 그러면
\[
\mathcal S=\bigoplus_{n\geq0}q_*(\mathcal S'_n)
\]
은 양의 차수들을 갖는 준연접 유한 생성 등급 $\mathcal O_Y$-대수이고,
$r^*(\shProj(\mathcal S))$가 등급 $\mathcal O_X$-대수로서
$\mathcal S'$과 동형이 되도록 하는 $Y$-동형사상
\[
r:X\xrightarrow{\sim}\operatorname{Proj}(\mathcal S)
\]
이 존재한다.
\end{enumerate}
\end{proposition}

\begin{proof}
(i) 모든 명제는 사실상 이미 증명되었고, 마지막 것은
(\hyperref[III.2.3.2-fr]{2.3.2})의 특수한 경우에 지나지 않는다.
$\mathcal S_X$가 주기적이라는 사실은 (II, 8.14.14)에서 보았고,
$\mathcal O_X(d)$가 가역이고 $Y$-풍부하도록 하는 주기 $d>0$이
존재한다는 사실은 (II, 4.6.18)에 지나지 않는다. 마지막으로
$0\leq k<d$에 대해 $(\mathcal S_X)^{(d,k)}$는 유한 생성
$(\mathcal S_X)^{(d)}$-가군이고 (II, 8.14.14), 따라서
(II, 2.1.6, (ii))에 의해 각 $(\mathcal S_X)_n$은 유한 생성 준연접
$\mathcal O_X$-가군이다. 문제는 국소적이고 $\mathcal O_X$가
연접이므로 $(\mathcal S_X)_n$들도 연접이다.

(ii) 주기 $d$를 그 배수로 바꾸어 $\mathcal L=\mathcal S'_d$가
$q$에 대해 \emph{매우 풍부한} 가역 $\mathcal O_X$-가군이라고 가정할
수 있다 (II, 4.6.11). 또한 가정에 의해
\[
\mathcal S'^{(d)}=\bigoplus_{n\in\mathbb Z}\mathcal L^{\otimes n}
\quad\text{이고, 따라서}\quad
\mathcal S^{(d)}=\bigoplus_{n\geq0}q_*(\mathcal L^{\otimes n}).
\]
(II, 3.1.8 및 3.2.9)에 의해
$X'=\operatorname{Proj}(\mathcal S)$에서
$X''=\operatorname{Proj}(\mathcal S^{(d)})$로 가는 $Y$-동형사상 $s$가
존재하며
\oldpage[III]{105}
\[
s^*(\mathcal O_{X''}(n))=\mathcal O_{X'}(nd)
\]
이다. 따라서 다음을 증명하면 $Y$-동형사상
$X\xrightarrow{\sim}X'$의 존재가 확립된다.

\begin{env}[2.3.4.1]
\label{III.2.3.4.1-fr}
$Y$를 뇌터 준스킴, $q:X\to Y$를 사영 사상,
$\mathcal L$을 $q$에 대해 매우 풍부한 가역 $\mathcal O_X$-가군이라
하자. 그러면
$\mathcal S=\bigoplus_{n\geq0}q_*(\mathcal L^{\otimes n})$은
$\mathcal S_n=\mathcal S_1^n$이 충분히 큰 $n$에 대해 성립하는 유한
생성 준연접 등급 $\mathcal O_Y$-대수이고, 어떤 $Y$-동형사상
$r:X\xrightarrow{\sim}P=\operatorname{Proj}(\mathcal S)$이 존재하여
$\mathcal L=r^*(\mathcal O_P(1))$이다.
\end{env}

$q$가 사영 사상이므로 (II, 5.4.4 및 4.4.7)에 의해 어떤 $Y$-동형사상
$r':X\xrightarrow{\sim}P'=\operatorname{Proj}(\mathcal T)$가 존재한다.
여기서 $\mathcal T$는 준연접 $\mathcal O_Y$-대수이고,
$\mathcal T_1$은 유한 생성이며 $\mathcal T$를 생성하고,
$\mathcal L=r'^*(\mathcal O_{P'}(1))$이다. $q':P'\to Y$가 구조
사상이면
\[
\mathcal S=\bigoplus_{n\geq0}q'_*(\mathcal O_{P'}(n))
\]
이다. (\hyperref[III.2.3.1-fr]{2.3.1})에 의해 충분히 큰 $n$에 대해
표준 준동형
$\alpha_n:\mathcal T_n\to\mathcal S_n=q'_*(\mathcal O_{P'}(n))$은
전단사이다. $\mathcal T_n=\mathcal T_1^n$이므로, 더구나 충분히 큰
$n$부터 $\mathcal S_n=\mathcal S_1^n$이다. 또한 등급
$\mathcal O_Y$-대수들의 표준 준동형
$\alpha:\mathcal T\to\mathcal S$은 \textbf{(TN)}-동형사상이다.
따라서
\phantomsection\label{III.2.3.5.1-fr}
\[
\Phi=\operatorname{Proj}(\alpha):
\operatorname{Proj}(\mathcal S)\longrightarrow
\operatorname{Proj}(\mathcal T)
\]
는 동형사상이다 (II, 3.6.1). 더 나아가 (II, 3.5.2)와, 스칼라 제한에
의해 $\mathcal S(n)$으로부터 얻은 등급 $\mathcal T$-가군들이
$\mathcal T(n)$과 \textbf{(TN)}-동형이라는 사실에 의해
\[
\Phi_*(\mathcal O_P(n))=\mathcal O_{P'}(n)
\]
이다 (II, 3.4.2). (\hyperref[III.2.3.4.1-fr]{2.3.4.1})의 증명을
끝내려면 $\mathcal S$가 유한 생성 $\mathcal O_Y$-대수임을 보이면
된다. 그런데 (\hyperref[III.2.2.1-fr]{2.2.1})에 의해
$\mathcal S_n=q'_*(\mathcal O_{P'}(n))$은 연접
$\mathcal O_Y$-가군이고, $n\geq n_0$일 때
$\mathcal S_n=\mathcal S_1^n$이면 $\mathcal S$는 연접인
$\bigoplus_{i\leq n_0}\mathcal S_i$로 생성된다. 이로부터 결론이
따른다 (I, 9.6.2).

(\hyperref[III.2.3.4-fr]{2.3.4})의 증명으로 돌아가 그 표기를 다시
사용하자. 모든 $n\in\mathbb Z$에 대해
$r''_*(\mathcal L^{\otimes n})=\mathcal O_{X''}(n)$이 되도록 하는
$Y$-동형사상 $r'':X\xrightarrow{\sim}X''$의 존재를 증명하였다.
$X''\to Y$의 구조 사상을 $q''$라 하자. 이제 $\mathcal S'$은 등급
$\mathcal S'^{(d)}$-가군 $\mathcal S'^{(d,k)}$들의 직합이다. 각
$\mathcal S'^{(d,k)}$는 $\mathcal S'$의 주기성과 $\mathcal S'_n$들이
유한 생성 $\mathcal O_X$-가군이라는 가정에 의해 유한 생성 준연접
등급 $\mathcal S'^{(d)}$-가군이다 (II, 8.14.12). 다음과 같이 놓자.
\[
\mathcal F^{(k)}=r''_*(\mathcal S'^{(d,k)}).
\]
그러면 $\mathcal F^{(k)}$들은 유한 생성 준연접 등급
$\mathcal S_{X''}$-가군이다. 따라서 (II, 8.14.8)에 의해 표준 준동형
\[
\beta:\shProj(\bbGamma_*(\mathcal F^{(k)}))
\longrightarrow\mathcal F^{(k)}
\]
은 $\mathcal S_{X''}$-가군의 동형사상이다. 한편
$q''_*((\mathcal F^{(k)})_n)=q_*((\mathcal S'^{(d,k)})_n)$이고,
$n\geq0$이면 마지막 $\mathcal O_Y$-가군은 정의에 의해
$(\mathcal S^{(d,k)})_n$과 같다. 다시 말해 표준 포함
$\mathcal S^{(d,k)}\to\bbGamma_*(\mathcal F^{(k)})$은
\textbf{(TN)}-동형사상이다. 그러므로
\[
\shProj(\mathcal S^{(d,k)})
=\shProj(\bbGamma_*(\mathcal F^{(k)})),
\]
따라서
$r''{}^*(\shProj(\mathcal S^{(d,k)}))\simeq\mathcal S'^{(d,k)}$이다.
마지막으로 표준 동형을 제외하면
\[
\shProj(\mathcal S^{(d,k)})
=s_*((\shProj(\mathcal S))^{(d,k)})
\]
임을 유의하면 (II, 8.14.13.1), $r^*(\shProj(\mathcal S))$와
$\mathcal S'$의 동형이 증명된다.

끝으로 (\hyperref[III.2.3.2-fr]{2.3.2})에 의해 각
$\bbGamma_*(\mathcal F^{(k)})$는 조건 \textbf{(TF)}를 만족하므로,
각 $\mathcal S^{(d,k)}$도 그러하다. 또한 $\mathcal S'_n$들이
연접이므로 (\hyperref[III.2.2.1-fr]{2.2.1})에 의해
$\mathcal S_n=q_*(\mathcal S'_n)$들은 연접이다. 이로부터 곧바로
$\mathcal S^{(d,k)}$들이 유한 생성 $\mathcal S^{(d)}$-가군임을
얻는다. (\hyperref[III.2.3.4.1-fr]{2.3.4.1})에 의해
$\mathcal S^{(d)}$가 유한 생성 $\mathcal O_Y$-대수이므로
$\mathcal S$도 그러하다.
\end{proof}

\oldpage[III]{106}
\begin{proposition}[2.3.5]
\label{III.2.3.5-fr}
$Y$를 뇌터 정역 준스킴, $X$를 정역 준스킴,
$f:X\to Y$를 사영 쌍유리 사상이라 하자. 그러면 어떤 연접 분수
아이디얼 $\mathcal J\subset\mathcal R(Y)$가 존재하여 (II, 8.1.2),
$X$는 $\mathcal J$를 따라 블로업하여 얻은 준스킴과 $Y$-동형이다
(II, 8.1.3). 또한 $Y$의 어떤 열린집합 $U$가 존재하여, $f$를
$f^{-1}(U)$에 제한한 것은 $f^{-1}(U)$에서 $U$로 가는 동형사상이고
(I, 6.5.5 참조), $\mathcal J|U$는 가역이다.
\end{proposition}

\begin{proof}
$f$에 대해 매우 풍부한 어떤 가역 $\mathcal O_X$-가군 $\mathcal L$이
존재하므로 (II, 4.4.2 및 5.3.2),
(\hyperref[III.2.3.4.1-fr]{2.3.4.1})을 적용할 수 있다. 따라서 $X$는
\[
\operatorname{Proj}(\mathcal S),\qquad
\mathcal S=\bigoplus_{n\geq0}f_*(\mathcal L^{\otimes n})
\]
과 동일시된다. 또한 $f_*(\mathcal L^{\otimes n})$들은 꼬임 없는
$\mathcal O_Y$-가군임을 안다 (I, 7.4.5). 따라서
$\mathcal O_Y$-가군 $\mathcal S$도 꼬임이 없고, 표준적으로
$\mathcal S\otimes_{\mathcal O_Y}\mathcal R(Y)$의 부분
$\mathcal O_Y$-가군과 동일시된다 (I, 7.4.1). 마지막 것은 그 제한을
어떤 공집합이 아닌 열린집합에서 알면 결정되는 \emph{단순한} 층이다
(I, 7.3.6). 예를 들어 $\mathcal L|f^{-1}(U')$가
$\mathcal O_X|f^{-1}(U')$와 동형이 되도록 하는 공집합이 아닌
$U'\subset U$에서 그러하다. 가정에 의해 이때
$f_*(\mathcal L^{\otimes n})|U'$들은 $\mathcal O_Y|U'$와 동형이므로,
$\mathcal S\otimes_{\mathcal O_Y}\mathcal R(Y)$는 부정원 $T$에 대한
$\mathcal R(Y)[T]$와 동형인 $\mathcal R(Y)$-가군이다. 또한
$\mathcal S$는 $f_*(\mathcal L)$의 표준 상이
$\mathcal S\otimes_{\mathcal O_Y}\mathcal R(Y)$ 안에서 생성하는
부분 $\mathcal O_Y$-대수와 \textbf{(TN)}-동형이다
(\hyperref[III.2.3.4.1-fr]{2.3.4.1}). 마지막 층을
$\mathcal R(Y)[T]$와 동일시하면 $f_*(\mathcal L)$의 상은
$\mathcal J T$와 동일시된다. 여기서 $\mathcal J$는
$\mathcal R(Y)$의 연접 부분 $\mathcal O_Y$-가군이고
(\hyperref[III.2.2.1-fr]{2.2.1}), 그 $U'$ 위의 제한은
$\mathcal O_Y|U'$와 동형이며, 따라서 $\mathcal J|U$가 가역이 되도록
한다. 이제 $\mathcal S$가 $\bigoplus_{n\geq0}\mathcal J^n$과
\textbf{(TN)}-동형임을 알 수 있고, 이로써 증명이 끝난다.
\end{proof}

\begin{corollary}[2.3.6]
\label{III.2.3.6-fr}
(\hyperref[III.2.3.5-fr]{2.3.5})의 가정에 더하여,
$\mathcal R(Y)$의 모든 영이 아닌 연접 부분 $\mathcal O_Y$-가군
$\mathcal J$에 대하여 어떤 가역 $\mathcal O_Y$-가군 $\mathcal L$이
존재하여
\[
\Gamma(Y,\mathcal L\otimes_{\mathcal O_Y}
\mathcal{H}om(\mathcal J,\mathcal O_Y))\neq0
\]
이라고 가정하자. 그러면 (\hyperref[III.2.3.5-fr]{2.3.5})의 명제에서
$\mathcal J$가 $\mathcal O_Y$의 아이디얼이라고 가정할 수 있다.
어떤 풍부한 $\mathcal O_Y$-가군이 존재하면 이 추가 조건은 항상
만족된다.
\end{corollary}

\begin{proof}
실제로 (0\textsubscript{I}, 5.4.2)에 의해
\[
\mathcal L\otimes\mathcal{H}om(\mathcal J,\mathcal O_Y)
=\mathcal{H}om(\mathcal L^{-1},\mathcal{H}om(\mathcal J,\mathcal O_Y))
=\mathcal{H}om(\mathcal J\otimes\mathcal L^{-1},\mathcal O_Y)
\]
이다. 따라서 가정은 $\mathcal J\otimes\mathcal L^{-1}$에서
$\mathcal O_Y$로 가는 영이 아닌 준동형 $u$가 존재한다는 뜻이다.
모든 $y\in Y$에 대하여
$(\mathcal J\otimes\mathcal L^{-1})_y$는 $\mathcal O_y$의 분수체
$\mathcal R(Y)_y$의 부분 $\mathcal O_y$-가군과 동일시되므로
(I, 7.1.5), $u_y$는 반드시 단사이다. 따라서 $u$는
$\mathcal J\otimes\mathcal L^{-1}$에서 $\mathcal O_Y$의 어떤 아이디얼
$\mathcal J'$로 가는 동형사상이다. 그런데
\[
\operatorname{Proj}\left(\bigoplus_{n\geq0}\mathcal J^n\right)
\quad\text{과}\quad
\operatorname{Proj}\left(
\bigoplus_{n\geq0}(\mathcal J\otimes\mathcal L^{-1})^n\right)
\]
은 $Y$-동형이므로 (II, 3.1.8), 이는 따름정리의 첫째 명제를
증명한다. 둘째 명제를 증명하기 위해
$\mathcal F=\mathcal{H}om_{\mathcal O_Y}(\mathcal J,\mathcal O_Y)$로
놓자. $\mathcal J|U$가 가역이 되도록 하는 열린집합 $U$가 존재하므로
$\mathcal F$는 연접이고 영이 아니다. $\mathcal L$이 풍부한
$\mathcal O_Y$-가군이면, 어떤 정수 $n$이 존재하여
$\mathcal F(n)=\mathcal F\otimes\mathcal L^{\otimes n}$은 $Y$ 위의
단면들로 생성된다 (II, 4.5.5). 더구나
$\Gamma(Y,\mathcal F(n))\neq0$이고, 이로부터 결론이 따른다.
\end{proof}

\begin{corollary}[2.3.7]
\label{III.2.3.7-fr}
$X$, $Y$를 체 $k$ 위의 두 사영 정역 스킴이라 하고,
$f:X\to Y$를 쌍유리 $k$-사상이라 하자. 그러면 $X$는 $Y$의 어떤
닫힌 부분스킴 $Y'$($Y'$은 반드시 축소일 필요가 없다)을 따라
블로업하여 얻은 $Y$-스킴과 $k$-동형이다.
\end{corollary}

\oldpage[III]{107}
\begin{proof}
실제로 $f$는 사영이고 (II, 5.5.5, (v)), $Y$가 $k$ 위에서
사영이므로 (\hyperref[III.2.3.6-fr]{2.3.6})의 추가 조건이 만족된다.
이제 따름정리 (\hyperref[III.2.3.6-fr]{2.3.6})의 연접 아이디얼
$\mathcal J$로 정의되는 $Y$의 닫힌 부분스킴 $Y'$을 택하면 충분하다.
\end{proof}

\begin{env}[2.3.8]
\label{III.2.3.8-fr}
\emph{비고.} 제IV장에서 약수의 개념을 연구할 때 다음을 보겠다.
(\hyperref[III.2.3.5-fr]{2.3.5})의 명제에서 환
$\mathcal O_y$ $(y\in Y)$들이 유일인수분해환이라고 가정하면
($Y$가 비특이이면 예를 들어 그러하다), $X$는 그 바탕공간이
$Y-U$에 포함되는 $Y$의 어떤 닫힌 부분준스킴 $Y'$을 따라
$Y$를 블로업하여 얻을 수 있다.
\end{env}

\subsection{기본정리의 일반화}
\label{subsection:III.2.4-fr}

\begin{theorem}[2.4.1]
\label{III.2.4.1-fr}
$Y$를 뇌터 준스킴, $\mathcal S$를 유한 생성 준연접
$\mathcal O_Y$-대수라 하자. $f:X\to Y$를 사영 사상,
$\mathcal S'=f^*(\mathcal S)$, $\mathcal M$을 유한 생성 준연접
$\mathcal S'$-가군이라 하자. 그러면 다음이 성립한다.
\begin{enumerate}
\item[(i)] 모든 $p\in\mathbb Z$에 대해 $R^pf_*(\mathcal M)$은 유한
생성 $\mathcal S$-가군이다.
\item[(ii)] 더 나아가 $\mathcal L$을 $f$에 대해 풍부한 가역
$\mathcal O_X$-가군이라 하고, 모든 $n\in\mathbb Z$에 대해
$\mathcal M(n)=\mathcal M\otimes\mathcal L^{\otimes n}$으로 놓자.
그러면 어떤 정수 $N$이 존재하여 $n\geq N$이면 모든 $p>0$에 대해
\[
R^pf_*(\mathcal M(n))=0
\tag{2.4.1.1}\label{III.2.4.1.1-fr}
\]
이고, 표준 준동형
\[
f^*(f_*(\mathcal M(n)))\longrightarrow\mathcal M(n)
\]
(0\textsubscript{I}, 4.4.3)은 전사이다.
\end{enumerate}
\end{theorem}

\begin{proof}
$Y'=\operatorname{Spec}(\mathcal S)$,
$X'=\operatorname{Spec}(\mathcal S')$로 놓자. 그러면
$X'=X\times_Y Y'$이다 (II, 1.5.5). 정의에 의해 아핀인 구조 사상들을
$g:Y'\to Y$, $g':X'\to X$라 하고,
$f'=f_{(Y')}:X'\to Y'$라 하자. 그러면 가환 도식
\[
\xymatrix{
X\ar[d]_f & X'\ar[l]_{g'}\ar[d]^{f'}\ar[dl]^{h}\\
Y & Y'\ar[l]^{g}
}
\]
이 있고, $f'$은 사영 사상이다 (II, 5.5.5, (iii)).
$h=f\circ g'=g\circ f'$로 놓자.

(i) $X'$을 $X$ 위의 아핀 스킴으로 보았을 때 준연접
$\mathcal S'$-가군 $\mathcal M$에 대응하는 $\mathcal O_{X'}$-가군을
$\widetilde{\mathcal M}$이라 하자 (II, 1.4.3). 그러면
$\mathcal M=g'_*(\widetilde{\mathcal M})$이다. $\mathcal M$이 유한
생성 $\mathcal S'$-가군이므로 $\widetilde{\mathcal M}$은 유한 생성
$\mathcal O_{X'}$-가군이다 (II, 1.4.5). 또한 $g$와 $f'$이 유한형이므로
$h$도 유한형이고 (II, 1.3.7 및 I, 6.3.4, (ii)), $X'$은 뇌터이다
(I, 6.3.7). 따라서 $\widetilde{\mathcal M}$은 연접이다.
$g'$이 아핀이므로 표준 준동형
\[
R^pf_*(\mathcal M)\longrightarrow R^ph_*(\widetilde{\mathcal M})
\]
은 전단사이다 (\hyperref[III.1.3.4-fr]{1.3.4}). 더 나아가 이 준동형은
$\mathcal S$-가군의 준동형이다. 실제로 $\mathcal M$의
$\mathcal S'$-가군 구조를 정의하는 표준 준동형
\[
g'_*(\mathcal O_{X'})\otimes_{\mathcal O_X}
g'_*(\widetilde{\mathcal M})\longrightarrow
g'_*(\widetilde{\mathcal M})
\tag{2.4.1.2}\label{III.2.4.1.2-fr}
\]
에서($\mathcal S'=g'_*(\mathcal O_{X'})$임을 상기하라), 표준적으로
준동형
\[
f_*(g'_*(\mathcal O_{X'}))\otimes
R^pf_*(g'_*(\widetilde{\mathcal M}))\longrightarrow
R^pf_*(g'_*(\widetilde{\mathcal M}))
\]
을 얻는다.

\oldpage[III]{108}
(0, 12.2.2). 또한 (\hyperref[III.2.4.1.2-fr]{2.4.1.2}) 자체가
$\widetilde{\mathcal M}$의 $\mathcal O_{X'}$-가군 구조를 정의하는
준동형
\[
\mathcal O_{X'}\otimes\widetilde{\mathcal M}
\longrightarrow\widetilde{\mathcal M}
\]
으로부터 (0\textsubscript{I}, 4.2.2.1)을 적용하여 얻어지므로, 도식
\[
\xymatrix{
f_*(g'_*(\mathcal O_{X'}))\otimes
R^pf_*(g'_*(\widetilde{\mathcal M}))
\ar[r]\ar[d] &
R^pf_*(g'_*(\widetilde{\mathcal M}))\ar[d]\\
h_*(\mathcal O_{X'})\otimes R^ph_*(\widetilde{\mathcal M})
\ar[r] & R^ph_*(\widetilde{\mathcal M})
}
\]
은 가환한다 (0, 12.2.6). 수평 화살표들을 표준 준동형
\[
\mathcal S\longrightarrow f_*(f^*(\mathcal S))=f_*(\mathcal S')
=f_*(g'_*(\mathcal O_{X'}))=h_*(\mathcal O_{X'})
\]
에서 오는 준동형과 합성하면 위의 주장을 얻는다. 한편 $g$가 아핀이고
$f'$이 분리이고 준콤팩트이므로, 표준 준동형
\[
R^ph_*(\widetilde{\mathcal M})\longrightarrow
g_*(R^pf'_*(\widetilde{\mathcal M}))
\]
은 전단사이다 (\hyperref[III.1.4.14-fr]{1.4.14}). 위와 같은 방식으로
이번에는 (0, 12.2.6.2)의 가환성을 사용하면, 이것이
$\mathcal S$-가군의 동형사상임을 증명할 수 있다. 그런데 $f'$은
사영이고 $\widetilde{\mathcal M}$은 연접이므로
(\hyperref[III.2.2.1-fr]{2.2.1})에 의해
$R^pf'_*(\widetilde{\mathcal M})$은 연접 $\mathcal O_{Y'}$-가군이다.
따라서 $g_*(R^pf'_*(\widetilde{\mathcal M}))$은 유한 생성
$\mathcal S$-가군이다 (II, 1.4.5).

(ii) $\mathcal L'=g'^*(\mathcal L)$로 놓자. 이는 가역
$\mathcal O_{X'}$-가군이다. 모든 $n\in\mathbb Z$에 대해 표준 동형을
제외하면
\[
g'_*(\widetilde{\mathcal M}\otimes\mathcal L'^{\otimes n})
=g'_*(\widetilde{\mathcal M})\otimes\mathcal L^{\otimes n}
=\mathcal M(n)
\]
이다 (0\textsubscript{I}, 5.4.10). $\widetilde{\mathcal M}$에 대해
(i)에서 한 논증을
$\widetilde{\mathcal M}\otimes\mathcal L'^{\otimes n}$에 적용하면
\[
R^pf_*(g'_*(\widetilde{\mathcal M}\otimes\mathcal L'^{\otimes n}))
\quad\text{은}\quad
g_*(R^pf'_*(\widetilde{\mathcal M}\otimes\mathcal L'^{\otimes n}))
\]
과 동형이다. $\mathcal L'$은 $f'$에 대해 풍부하므로
(II, 4.6.13, (iii)), (\hyperref[III.2.2.1-fr]{2.2.1})에 의해 어떤
정수 $N$이 존재하여 모든 $p>0$, $n\geq N$에 대해
\[
R^pf'_*(\widetilde{\mathcal M}\otimes\mathcal L'^{\otimes n})=0
\]
이다. 이는 (\hyperref[III.2.4.1.1-fr]{2.4.1.1})을 증명한다.
다시 (\hyperref[III.2.2.1-fr]{2.2.1})에 의해, $n\geq N$이면 표준
준동형
\[
f'^*(f'_*(\widetilde{\mathcal M}\otimes\mathcal L'^{\otimes n}))
\longrightarrow\widetilde{\mathcal M}\otimes\mathcal L'^{\otimes n}
\]
이 전사가 되도록 $N$을 택할 수 있다. $g'_*$은 완전한 함자이므로
(II, 1.4.4), 대응하는 준동형
\[
g'_*(f'^*(f'_*(\widetilde{\mathcal M}\otimes
\mathcal L'^{\otimes n})))
\longrightarrow
g'_*(\widetilde{\mathcal M}\otimes\mathcal L'^{\otimes n})
=\mathcal M(n)
\]
은 전사이다. 한편 (II, 1.5.2)에 의해
\[
g'_*(f'^*(f'_*(\widetilde{\mathcal M}\otimes
\mathcal L'^{\otimes n})))
=f^*(g_*(f'_*(\widetilde{\mathcal M}\otimes
\mathcal L'^{\otimes n})))
\]
이고 $g_*\circ f'_* = f_*\circ g'_*$이므로, 결국
\[
g'_*(f'^*(f'_*(\widetilde{\mathcal M}\otimes\mathcal L'^{\otimes n})))
=f^*(f_*(g'_*(\widetilde{\mathcal M}\otimes
\mathcal L'^{\otimes n})))=f^*(f_*(\mathcal M(n)))
\]
이다. 이로써 증명이 끝난다.
\end{proof}

\begin{env}[2.4.2]
\label{III.2.4.2-fr}
특히 $\mathcal S$가 양의 차수들을 갖는 등급 $\mathcal O_Y$-대수이고
$\mathcal M=\bigoplus_{k\in\mathbb Z}\mathcal M_k$가 등급
$\mathcal S'$-가군일 때 (\hyperref[III.2.4.1-fr]{2.4.1})을 적용하게
된다.

\oldpage[III]{109}
$\mathcal S$, $\mathcal M$에 대해 같은 유한성 가정을 두면,
(\hyperref[III.2.4.1-fr]{2.4.1})로부터 모든 $p$에 대해
\[
\bigoplus_{k\in\mathbb Z}R^pf_*(\mathcal M_k)
\]
가 유한 생성 $\mathcal S$-가군임을 얻는다. 또한
(\hyperref[III.2.4.1-fr]{2.4.1}, (ii))의 추가 가정 아래에서는 어떤
$N$이 존재하여 $n\geq N$이면 모든 $p>0$, $k\in\mathbb Z$에 대해
$R^pf_*(\mathcal M_k(n))=0$이고, 표준 준동형
$f^*(f_*(\mathcal M_k(n)))\to\mathcal M_k(n)$은 모든
$k\in\mathbb Z$에 대해 전사이다.
\end{env}

\subsection{오일러--푸앵카레 특성과 힐베르트 다항식}
\label{subsection:III.2.5-fr}

\begin{env}[2.5.1]
\label{III.2.5.1-fr}
$A$를 아르틴 환, $X$를 $Y=\operatorname{Spec}(A)$ 위의 사영
$A$-스킴이라 하자. 모든 연접 $\mathcal O_X$-가군 $\mathcal F$에 대해
$H^i(X,\mathcal F)$ $(i\geq0)$들은 유한 생성 $A$-가군이다
(\hyperref[III.2.2.1-fr]{2.2.1}). $A$가 아르틴이므로 여기서는 유한
길이를 갖는다. 또한 (\hyperref[III.2.2.1-fr]{2.2.1})에 의해
$H^i(X,\mathcal F)$는 유한 개의 $i\geq0$을 제외하면 0이다. 따라서
정수
\[
\chi_A(\mathcal F)=\sum_{i=0}^{\infty}(-1)^i
\operatorname{long}(H^i(X,\mathcal F))
\tag{2.5.1.1}\label{III.2.5.1.1-fr}
\]
는 모든 연접 $\mathcal O_X$-가군 $\mathcal F$에 대해 정의된다.
$A$가 아르틴 국소환이면 $\chi_A(\mathcal F)$를 $\mathcal F$의
($A$에 대한) 오일러--푸앵카레 특성이라 한다.
$\mathcal F=\mathcal O_X$일 때 $\chi_A(\mathcal O_X)$를 $X$의
($A$에 대한) 산술 종수라고 한다.
\end{env}

\begin{proposition}[2.5.2]
\label{III.2.5.2-fr}
\[
0\longrightarrow\mathcal F'\longrightarrow\mathcal F
\longrightarrow\mathcal F''\longrightarrow0
\]
을 연접 $\mathcal O_X$-가군들의 완전열이라 하자. 그러면
\[
\chi_A(\mathcal F)=\chi_A(\mathcal F')+\chi_A(\mathcal F'').
\tag{2.5.2.1}\label{III.2.5.2.1-fr}
\]
\end{proposition}

\begin{proof}
$\mathcal F$, $\mathcal F'$, $\mathcal F''$의 코호몰로지 가군들은
유한 개를 제외하면 0이므로, 어떤 정수 $r>0$이 존재하여 코호몰로지
완전열은
\[
0\to H^0(X,\mathcal F')\to H^0(X,\mathcal F)
\to H^0(X,\mathcal F'')\to H^1(X,\mathcal F')\to\cdots
\]
\[
\cdots\to H^r(X,\mathcal F')\to H^r(X,\mathcal F)
\to H^r(X,\mathcal F'')\to0
\]
으로 쓰인다. 양 끝이 0인 유한 길이 $A$-가군들의 완전열에서는
길이들의 교대합이 0임을 안다 (0, 11.10.1). 이 결과를 적용하면 공식
(\hyperref[III.2.5.2.1-fr]{2.5.2.1})을 즉시 얻는다.
\end{proof}

$X$가 준콤팩트 $A$-스킴이고 모든 연접 $\mathcal O_X$-가군
$\mathcal F$에 대해 $A$-가군 $H^i(X,\mathcal F)$들이 유한 생성임을
아는 모든 경우에도 (\hyperref[III.2.5.2-fr]{2.5.2})의 결과가 적용됨을
유의하자 (\hyperref[III.1.4.12-fr]{1.4.12}).

\begin{theorem}[2.5.3]
\label{III.2.5.3-fr}
$A$를 아르틴 국소환, $X$를 $Y=\operatorname{Spec}(A)$ 위의 사영 스킴,
$\mathcal L$을 $Y$에 대해 매우 풍부한 가역 $\mathcal O_X$-가군,
$\mathcal F$를 연접 $\mathcal O_X$-가군이라 하자. 모든
$n\in\mathbb Z$에 대해
$\mathcal F(n)=\mathcal F\otimes_{\mathcal O_X}\mathcal L^{\otimes n}$
으로 놓는다.
\begin{enumerate}
\item[(i)] 모든 $n\in\mathbb Z$에 대해
$\chi_A(\mathcal F(n))=P(n)$이 되도록 하는 유일한 다항식
$P\in\mathbb Q[T]$가 존재한다. $P$를 $\mathcal F$의 $A$에 대한
힐베르트 다항식이라 한다.
\item[(ii)] 충분히 큰 $n$에 대해
$\chi_A(\mathcal F(n))=\operatorname{long}_A
\Gamma(X,\mathcal F(n))$이다.
\item[(iii)] $\chi_A(\mathcal F(n))$의 최고차항의 계수는
$\geq0$이다.
\end{enumerate}
\end{theorem}

제IV장의 차원 개념을 다루는 절에서는
$\chi_A(\mathcal F(n))$의 차수가 $\mathcal F$의 받침의 차원과
같음도 증명하겠다.

\begin{proof}
충분히 큰 $n$이면 모든 $i>0$에 대해
$H^i(X,\mathcal F(n))=0$이므로
(\hyperref[III.2.2.1-fr]{2.2.1}),

\oldpage[III]{110}
충분히 큰 $n$에 대해
\[
\chi_A(\mathcal F(n))=\operatorname{long}H^0(X,\mathcal F(n))
=\operatorname{long}\Gamma(X,\mathcal F(n))
\]
이다. 이로써 (ii)가 따르고, 충분히 큰 $n$에 대해
$\chi_A(\mathcal F(n))\geq0$이다. 따라서 (iii)은 (i)에서 따른다.
또한 (i)의 유일성 명제는 명백하므로 다항식 $P$의 존재만 증명하면
된다.

먼저 $\mathfrak m$이 $A$의 극대 아이디얼일 때
$\mathfrak m\mathcal F=0$이라고 가정할 수 있음을 보이자. 실제로
$\mathfrak m^s=0$인 어떤 정수 $s>0$이 존재하므로 $\mathcal F(n)$은
유한 여과
\[
\mathcal F(n)\supset\mathfrak m\mathcal F(n)\supset\cdots
\supset\mathfrak m^{s-1}\mathcal F(n)\supset0
\]
를 갖는다. (\hyperref[III.2.5.2.1-fr]{2.5.2.1})과 귀납법으로
\[
\chi_A(\mathcal F(n))=\sum_{k=1}^{s}
\chi_A(\mathfrak m^{k-1}\mathcal F(n)/\mathfrak m^k\mathcal F(n))
\]
을 얻는다. 그런데
$\mathfrak m^{k-1}\mathcal F(n)/\mathfrak m^k\mathcal F(n)
=\mathcal F'_k(n)$이고
$\mathcal F'_k=\mathfrak m^{k-1}\mathcal F/\mathfrak m^k\mathcal F$이므로,
이는 우리의 주장을 증명한다.

따라서 $\mathfrak m\mathcal F=0$이라고 가정하자. $X'$을 구조 사상
$X\to\operatorname{Spec}(A)$에 의한 $\operatorname{Spec}(A)$의 유일한
닫힌점의 역상인 $X$의 닫힌 부분스킴이라 하고, $j:X'\to X$를 표준
포함이라 하자. 그러면 $\mathcal F=j_*(\mathcal F')$이며,
$\mathcal F'$은 연접 $\mathcal O_{X'}$-가군이다. 또한
$K=A/\mathfrak m$일 때 $X'$은 $\operatorname{Spec}(K)$ 위의 사영
스킴이다. $\mathcal L'=j^*(\mathcal L)$로 놓으면 $\mathcal L'$은
$\operatorname{Spec}(K)$에 대해 매우 풍부하고 (II, 4.4.10),
\[
\mathcal F(n)=j_*(\mathcal F'(n)),\qquad
\mathcal F'(n)=\mathcal F'\otimes_{\mathcal O_{X'}}
\mathcal L'^{\otimes n}
\]
이다 (0\textsubscript{I}, 5.4.10). 따라서
$\chi_A(\mathcal F(n))=\chi_K(\mathcal F'(n))$이고
(G, II, 4.9.1), $A$가 체인 경우로 귀착된다.

이제 어떤 적당한 $r$에 대해 $X$를 $P=\mathbf P_A^r$의 닫힌
부분스킴으로 볼 수 있다 (II, 5.5.4, (ii)). $i:X\to P$가 표준
포함이면 위와 같이
$\chi_A(\mathcal F(n))=\chi_A(i_*(\mathcal F)(n))$임을 알 수 있다.
따라서 $A$가 체이고
$X=\mathbf P_A^r=\operatorname{Proj}(S)$,
$S=A[T_0,\ldots,T_r]$인 경우로 한정할 수 있다.

이제 $\mathcal F=\widetilde M$이며, $M$은 유한 생성 등급
$S$-가군이다 (II, 2.7.8). 힐베르트 시지지 정리 (M, VIII, 6.5)에
의해 $M$은 유한 생성 자유 등급 $S$-가군들에 의한 유한 분해
\[
0\to L_q\to L_{q-1}\to\cdots\to L_1\to M\to0
\]
를 갖는다. $M\mapsto\widetilde M$은 $M$에 관해 완전한 함자이므로
(II, 2.5.4), 완전열
\[
0\to\widetilde L_q\to\widetilde L_{q-1}\to\cdots
\to\widetilde L_1\to\widetilde M\to0
\]
도 존재하고, 따라서 모든 $n\in\mathbb Z$에 대해 열
\[
0\to\widetilde L_q(n)\to\widetilde L_{q-1}(n)\to\cdots
\to\widetilde L_1(n)\to\widetilde M(n)\to0
\]
은 완전하다. 명제 (2.5.2)를 $q$에 대한 귀납법으로 적용하면
\[
\chi_A(\widetilde M(n))=\sum_{j=1}^{q}(-1)^{j+1}
\chi_A(\widetilde L_j(n))
\]
이다. 그러므로 (i)를 증명하기 위해 $M$이 유한 생성 자유 등급
가군인 경우, 나아가 어떤 $h\in\mathbb Z$에 대해 $M=S(h)$인 경우로
귀착된다. 이때
$\widetilde M(n)=(M(n))^\sim=(S(n+h))^\sim$이므로 (II, 2.5.15),
결국 정리는 다음 보조정리에서 따른다.
\end{proof}

\oldpage[III]{111}
\begin{lemma}[2.5.3.1]
\label{III.2.5.3.1-fr}
$A$를 체, $r>0$을 정수, $X=\mathbf P_A^r$라 하자. 그러면 모든
$n\in\mathbb Z$에 대해
\[
\chi_A(\mathcal O_X(n))=\binom{n+r}{r}
\]
이다.
\end{lemma}

\begin{proof}
실제로 $n\geq0$이면
$\chi_A(\mathcal O_X(n))=\operatorname{long}H^0(X,\mathcal O_X(n))$이고,
이는 총차수가 $n$인 $T_i$들의 단항식의 개수, 즉
$\binom{n+r}{r}$이다 (\hyperref[III.2.1.12-fr]{2.1.12}).
$n\leq-r-1$이면 마찬가지로
\[
\chi_A(\mathcal O_X(n))=(-1)^r
\operatorname{long}H^r(X,\mathcal O_X(n)).
\]
$n=-r-h$라 하면 $A$ 위에서 $H^r(X,\mathcal O_X(n))$의 차원은
$p_i>0$, $\sum_{i=0}^{r}p_i=r+h$인 정수열
$(p_i)_{0\leq i\leq r}$의 개수이다
(\hyperref[III.2.1.12-fr]{2.1.12}). 이는 $q_i\geq0$,
$\sum_{i=0}^{r}q_i=h-1$인 정수열 $(q_i)_{0\leq i\leq r}$의 개수와
같으므로
\[
\binom{h+r-1}{r}=(-1)^r\binom{n+r}{r}
\]
이다. 마지막으로 $-r\leq n<0$이면 $\binom{n+r}{r}=0$이고,
한편 모든 $i\geq0$에 대해 $H^i(X,\mathcal O_X(n))=0$이다
(\hyperref[III.2.1.12-fr]{2.1.12}). 이로써 보조정리가 증명된다.
\end{proof}

\begin{corollary}[2.5.4]
\label{III.2.5.4-fr}
$A$를 아르틴 국소환, $S$를 $S_1$로 생성되는 유한 생성 등급
$A$-대수, $M$을 유한 생성 등급 $S$-가군,
$X=\operatorname{Proj}(S)$라 하자. 충분히 큰 $n$에 대해
\[
\chi_A(\widetilde M(n))=\operatorname{long}M_n
\]
이다.
\end{corollary}

\begin{proof}
충분히 큰 $n$에 대해 $M_n$과 $\Gamma(X,\widetilde M(n))$이
동형이라는 사실에서 따른다 (\hyperref[III.2.3.1-fr]{2.3.1}).
\end{proof}

\subsection{응용: 풍부성 판정법}
\label{subsection:III.2.6-fr}

\begin{proposition}[2.6.1]
\label{III.2.6.1-fr}
$Y$를 뇌터 준스킴, $f:X\to Y$를 고유 사상,
$\mathcal L$을 가역 $\mathcal O_X$-가군이라 하자. 다음 조건들은
동치이다.
\begin{enumerate}
\item[a)] $\mathcal L$은 $f$에 대해 풍부하다.
\item[b)] 모든 연접 $\mathcal O_X$-가군 $\mathcal F$에 대해 어떤
정수 $N$이 존재하여 $n\geq N$이면 모든 $q>0$에 대해
\[
R^qf_*(\mathcal F\otimes\mathcal L^{\otimes n})=0
\]
이다.
\item[c)] $\mathcal O_X$의 모든 연접 아이디얼 $\mathcal J$에 대해
어떤 정수 $N$이 존재하여 $n\geq N$이면
\[
R^1f_*(\mathcal J\otimes\mathcal L^{\otimes n})=0
\]
이다.
\end{enumerate}
\end{proposition}

\begin{proof}
a)가 b)를 함의함은 이미 보았다
(\hyperref[III.2.2.1-fr]{2.2.1}, (ii)). b)가 c)를 함의함은 자명하므로,
c)가 a)를 함의함을 증명하면 된다. $Y$가 아핀인 경우로 한정하여
(II, 4.6.4), 이 경우 $\mathcal L$이 풍부함을 증명할 수 있다.
$h$가 $X$ 위의 $\mathcal L^{\otimes n}$ $(n>0)$의 모든 단면을
움직일 때 아핀인 $X_h$들이 $X$를 덮음을 보이면 충분하다
(II, 4.5.2). 이를 위해 $X$의 모든 닫힌점 $x$와 $x$의 모든 아핀
열린 근방 $U$에 대해, 어떤 $n$과
$h\in\Gamma(X,\mathcal L^{\otimes n})$이 존재하여
$x\in X_h\subset U$임을 보이자. 그러면 $X_h$는 반드시 아핀이고
(I, 1.3.6), 이러한 $X_h$들의 합집합은 $X$의 모든 닫힌점을 포함하는
열린집합이다. $X$가 뇌터이므로 이는 $X$ 자체이다
(I, 6.3.7 및 0\textsubscript{I}, 2.1.3).

$X-U$ (각각 $(X-U)\cup\{x\}$)를 바탕공간으로 하는 $X$의 축소 닫힌
부분준스킴을 정의하는 $\mathcal O_X$의 준연접 아이디얼층을
$\mathcal J$ (각각 $\mathcal J'$)라 하자 (I, 5.2.1).
$\mathcal J$, $\mathcal J'$은 연접이고 (I, 6.1.1),
$\mathcal J'\subset\mathcal J$이며,
$\mathcal J''=\mathcal J/\mathcal J'$은 받침이 $\{x\}$이고
$\mathcal J''_x=k(x)$인 연접 $\mathcal O_X$-가군이다
(0\textsubscript{I}, 5.3.3). $\mathcal L$이 국소 자유이므로 모든 $n$에
대해 열
\[
0\to\mathcal J'\otimes\mathcal L^{\otimes n}
\to\mathcal J\otimes\mathcal L^{\otimes n}
\to\mathcal J''\otimes\mathcal L^{\otimes n}\to0
\]
은 완전하다. 가정에 의해 충분히 큰 어떤 $n$에 대해
$H^1(X,\mathcal J'\otimes\mathcal L^{\otimes n})=0$이다. 따라서
코호몰로지 완전열로부터 준동형

\oldpage[III]{112}
\[
\Gamma(X,\mathcal J\otimes\mathcal L^{\otimes n})
\longrightarrow
\Gamma(X,\mathcal J''\otimes\mathcal L^{\otimes n})
\]
이 전사임을 얻는다. $g(x)\neq0$인
$\mathcal J''\otimes\mathcal L^{\otimes n}$의 $X$ 위의 단면 $g$는
따라서 어떤 단면
\[
h\in\Gamma(X,\mathcal J\otimes\mathcal L^{\otimes n})
\subset\Gamma(X,\mathcal L^{\otimes n})
\]
의 상이다. 여기서 포함은 (0\textsubscript{I}, 5.4.1)에 의해
$\mathcal J\otimes\mathcal L^{\otimes n}$이
$\mathcal L^{\otimes n}$의 부분 $\mathcal O_X$-가군이기 때문이다.
정의에 의해 $h(x)\neq0$이고 $z\notin U$이면 $h(z)=0$이다. 이로써
증명이 끝난다.
\end{proof}

\begin{proposition}[2.6.2]
\label{III.2.6.2-fr}
$Y$를 뇌터 준스킴, $f:X\to Y$를 유한형 사상,
$g:X'\to X$를 유한 전사 사상, $\mathcal L$을 가역
$\mathcal O_X$-가군이라 하고 $\mathcal L'=g^*(\mathcal L)$로 놓자.
다음 조건을 가정한다. $Y$ 위에서 고유인 어떤 부분집합 $Z\subset X$가
존재하여 (II, 5.4.10), 모든 $x\in X-Z$에 대해 $X$가 $x$에서
정규이거나 $(g_*(\mathcal O_{X'}))_x$가 자유 $\mathcal O_x$-가군이다.
이 조건 아래에서 $\mathcal L$이 $f$에 대해 풍부일 필요충분조건은
$\mathcal L'$이 $f\circ g$에 대해 풍부인 것이다.
\end{proposition}

\begin{proof}
\begin{env}[2.6.2.1]
\label{III.2.6.2.1-fr}
$g$가 아핀이므로 조건이 필요함은 (II, 5.1.12)에서 따른다. 충분함을
보이기 위해 $Y$가 아핀이라고 가정할 수 있다 (II, 4.6.4). 더 나아가
$X$가 축소인 경우로 한정할 수 있음을 보이자. 실제로
$j:X_{\mathrm{red}}\to X$를 표준 포함이라 하고,
$X_1=X_{\mathrm{red}}$, $X'_1=X'\times_X X_1$으로 놓자. 그러면 가환
도식
\[
\xymatrix{
X'\ar[d]_g & X'_1\ar[l]_{j'}\ar[d]^{g_1\, .}\\
X & X_1\ar[l]^{j}
}
\tag{2.6.2.2}\label{III.2.6.2.2-fr}
\]
을 얻는다. 사상 $f\circ j$는 유한형이고 (I, 6.3.4), $g_1$은 유한
사상이다 (II, 6.1.5, (iii)). $\mathcal L'$이 $f\circ g$에 대해
풍부이면 $j'$이 닫힌 몰입이므로
$j'^*(\mathcal L')$은 $f\circ g\circ j'$에 대해 풍부하다
(II, 5.1.12 및 I, 4.3.2). $Z_1=j^{-1}(Z)$로 놓으면 $Z_1$은 $Y$ 위에서
고유이다 (II, 5.4.10). 한편 $X$가 점 $x$에서 정규이면
$X_{\mathrm{red}}$도 명백히 그러하다. 마지막으로
$(g_*(\mathcal O_{X'}))_x$가 자유 $\mathcal O_x$-가군이면
(II, 1.5.2)에서 곧바로
$((g_1)_*(\mathcal O_{X'_1}))_x$가 자유 $\mathcal O_{X_1,x}$-가군임을
얻는다. 끝으로 $X$가 뇌터이므로 (I, 6.3.7), $j^*(\mathcal L)$이
풍부이면 $\mathcal L$도 풍부하다 (II, 4.5.14). 또한
$j'^*(\mathcal L')=g_1^*(j^*(\mathcal L))$이므로 앞서 말한 귀착이
끝난다. 이제부터 $Y$가 아핀이고 $X$가 축소라고 가정한다.
\end{env}

이제 (II, 6.6.11)의 가정들이 만족되므로, 어떤 축소 $Y$-준스킴
$X_2$와 유한 쌍유리 $Y$-사상 $h:X_2\to X$가 존재하여,
$h$를 $h^{-1}(X-Z)$에 제한한 것은 $X-Z$ 위의 동형사상이고
$h^*(\mathcal L)$은 풍부하다. $X'$을 $X_2$로 바꾸면, $g$가 방금
$h$에 대해 열거한 성질들을 더 만족한다고 가정하여 명제를 증명하는
문제로 귀착된다. 계속해서 $Z$를 바탕공간으로 $Z$를 가지며 $Y$ 위에서
고유인 $X$의 부분준스킴도 $Z$로 나타내겠다 (II, 5.4.10).

\begin{env}[2.6.2.3]
\label{III.2.6.2.3-fr}
이제 $X_1$을 $X$의 닫힌 부분준스킴, $j:X_1\to X$를 표준 포함,
$X'_1=g^{-1}(X_1)=X'\times_X X_1$을 그 역상,
$j':X'_1\to X'$을 표준 포함이라 하자. 그러면 가환 도식
(\hyperref[III.2.6.2.2-fr]{2.6.2.2})가 있다.
$\mathcal L_1=j^*(\mathcal L)$,
$\mathcal L'_1=j'^*(\mathcal L')=g_1^*(\mathcal L_1)$로 놓으면,
$\mathcal L'_1$은 $f\circ g\circ j'$에 대해 풍부하다
(II, 5.1.12). $Z_1=j^{-1}(Z)$로 놓으면 $X_1$의 닫힌 부분준스킴
$Z_1$은 $Y$ 위에서 고유이다 (II, 5.4.2, (ii)). 다시 말해 $X_1$,
$\mathcal L_1$, $g_1$, $Z_1$에 대해
(\hyperref[III.2.6.2-fr]{2.6.2})의 가정들이 만족된다.
\end{env}

\oldpage[III]{113}
이 사실을 이용하여, $g$를 $g^{-1}(X-Z)$에 제한한 것이 $X-Z$ 위의
동형사상인 경우에 뇌터 귀납법 (0\textsubscript{I}, 2.2.2)으로
(\hyperref[III.2.6.2-fr]{2.6.2})를 증명할 수 있다. 앞서
(\hyperref[III.2.6.2.1-fr]{2.6.2.1})에서 본 대로 이것이면 충분하다.
즉 바탕공간이 $X$와 다른 $X$의 모든 닫힌 부분준스킴 $X_1$에 대해
(\hyperref[III.2.6.2-fr]{2.6.2})의 결론이 $\mathcal L_1$에 대해
참이라고 가정하고, 그 결론이 $\mathcal L$에 대해서도 참임을 보이면
된다.

\begin{env}[2.6.2.4]
\label{III.2.6.2.4-fr}
$\mathcal A=\mathcal O_X$, $\mathcal B=g_*(\mathcal O_{X'})$로 놓자.
그러면 $\mathcal B$는 $\mathcal R(X)$의 부분 $\mathcal A$-대수이고,
연접 $\mathcal A$-가군이다. 또한
$\mathcal B|(X-Z)=\mathcal A|(X-Z)$이다. $\mathcal K$를
$\mathcal A$ 위의 $\mathcal B$의 \emph{도체}, 즉
$\mathcal B\mathcal K\subset\mathcal A$를 만족하는
$\mathcal A$의 가장 큰 준연접 부분 $\mathcal A$-가군이라 하자. 이는
$\mathcal A$-가군 $\mathcal B/\mathcal A$의 소멸자와도 같다
(0\textsubscript{I}, 5.3.7). 따라서
$\mathcal B\mathcal K=\mathcal K$이다. $g$가 어떤 근방 $W_x$의 역상
$g^{-1}(W_x)$에서 $W_x$로 가는 동형사상이 되는 모든 점 $x$에서는
$\mathcal K_x=\mathcal A_x$임이 명백하다. 특히 이는 $X-Z$의 모든
점과 $X$의 각 기약 성분의 일반점의 어떤 근방에서 성립한다.
$\mathcal K$가 정의하는 $X$의 닫힌 부분준스킴
$Z_1=\operatorname{Spec}(\mathcal A/\mathcal K)$를 생각하자.
바탕공간 $Z_1$이 $Z$에서 닫혀 있으므로 $Z_1$도 $Y$ 위에서 고유이다
(II, 5.4.10). 또한 $\mathcal K$의 정의로부터
$\mathcal B|(X-Z_1)=\mathcal A|(X-Z_1)$임을 얻는다. 따라서 언제나
$Z=\operatorname{Spec}(\mathcal A/\mathcal K)$인 경우로 귀착할 수
있다. 위에서 $X-Z_1$이 $X$의 공집합이 아닌 열린집합임을 보았으므로,
항상 바탕공간 $Z$가 $X$와 다르다고 가정할 수 있다.
\end{env}

\begin{env}[2.6.2.5]
\label{III.2.6.2.5-fr}
$g$가 아핀이므로 $X'=\operatorname{Spec}(\mathcal B)$로 보고,
$\mathcal K'=\widetilde{\mathcal K}$라 하자. 이는
$g_*(\mathcal K')=\mathcal K$인 $\mathcal O_{X'}$의 연접 아이디얼이다
(II, 1.4.1). $X'$의 닫힌 부분준스킴
$Z'=g^{-1}(Z)=Z\times_X X'$은 $\mathcal K'$로 정의되고
$\operatorname{Spec}(\mathcal B/\mathcal K)$와 같다 (II, 1.4.10).
$h:Z'\to Z$가 유한 사상이므로 (II, 6.1.5, (iii)), $Z'$은 $Y$ 위에서
고유이다 (II, 6.1.11 및 5.4.2, (ii)).

이제 모든 $x\in X$와 모든 열린 근방 $U$에 대해, 어떤 $n>0$과
$X$ 위의 $\mathcal L^{\otimes n}$의 단면 $s$가 존재하여
$x\in X_s\subset U$임을 증명해야 한다 (II, 4.5.2). 두 경우로
나눈다.

$1^\circ$ $x\in X-Z$인 경우. 이때 $U\subset X-Z$라고 가정해도 되고,
따라서 열린집합 $U'=g^{-1}(U)$은 $Z'$과 만나지 않는다.
$\mathcal L'$이 풍부하므로 어떤 $n>0$과 $X'$ 위의
$\mathcal L'^{\otimes n}$의 단면 $s'$이 존재하여
$x'=g^{-1}(x)\in X'_{s'}\subset g^{-1}(U)$이다 (II, 4.5.2).
또한 $\mathcal K'\otimes\mathcal L'^{\otimes n}$이 $X'$ 위의
단면들로 생성된다고 가정할 수 있다 (II, 4.5.5). 따라서
$\mathcal K'_{x'}=\mathcal O_{x'}$이므로 이 단면들 가운데
$s''(x')\neq0$인 $s''$이 존재한다. 이를 $s'$과 곱하면($n$을 $2n$으로
바꾸는 것과 같다), $x'\in X'_{s''}\subset g^{-1}(U)$라고도 가정할 수
있다. 이제 (0\textsubscript{I}, 5.4.10)에 의해 표준 동형사상
\[
\Gamma(X,\mathcal K\otimes\mathcal L^{\otimes n})
\xrightarrow{\sim}
\Gamma(X',\mathcal K'\otimes\mathcal L'^{\otimes n})
\]
이 있다. 이 동형사상에서 $s''$에 대응하는
$\mathcal K\otimes\mathcal L^{\otimes n}$의 단면 $s$는 명백히 원하는
성질들을 갖는다.

$2^\circ$ $x\in Z$인 경우. 바탕공간이 $X-U$인 $X$의 축소 닫힌
부분준스킴을 정의하는 $\mathcal O_X$의 연접 아이디얼을 $\mathcal J$라
하자. $\mathcal B$ 안의 두 연접 아이디얼
$\mathcal J\mathcal B$와
$\mathcal J\mathcal K=\mathcal J(\mathcal K\mathcal B)
=\mathcal K(\mathcal J\mathcal B)$를 생각하면 포함 도식

\oldpage[III]{114}
\[
\xymatrix@R=1.5em{
\mathcal J\mathcal B\ar@{^{(}->}[r] & \mathcal B\\
\mathcal J\ar@{^{(}->}[r]\ar@{^{(}->}[u] &
\mathcal A\ar@{^{(}->}[u]\\
\mathcal J\mathcal K\mathcal B=\mathcal J\mathcal K
\ar@{^{(}->}[r]\ar@{^{(}->}[u] &
\mathcal K\ar@{^{(}->}[u]
}
\tag{2.6.2.6}\label{III.2.6.2.6-fr}
\]
을 얻는다. $\mathcal J'=(\mathcal J\mathcal B)^\sim$를
$\mathcal O_{X'}$의 연접 아이디얼이라 하자. 그러면
$\mathcal J\mathcal B=g_*(\mathcal J')$이고,
$\mathcal J'\mathcal K'=(\mathcal J\mathcal K\mathcal B)^\sim$이므로
\[
\mathcal J'/\mathcal J'\mathcal K'
=(\mathcal J\mathcal B/\mathcal J\mathcal K\mathcal B)^\sim
\]
이다 (II, 1.4.4). $Z$와 만나지 않는 모든 열린집합 $V$에서
$\mathcal J|V=\mathcal J\mathcal K|V$이므로,
$\mathcal J'/\mathcal J'\mathcal K'$의 받침은 $Z'$에 포함된다.
$Z'$이 $Y$ 위에서 고유이므로 (\hyperref[III.2.2.4-fr]{2.2.4})를
적용할 수 있다. 따라서 충분히 큰 $n$에 대해 표준 사상
\[
\Gamma(X',\mathcal J'\otimes\mathcal L'^{\otimes n})
\longrightarrow
\Gamma(X',(\mathcal J'/\mathcal J'\mathcal K')
\otimes\mathcal L'^{\otimes n})
\]
은 전사이다. (0\textsubscript{I}, 5.4.10)에 의해, 이로부터 표준 사상
\[
\Gamma(X,\mathcal J\mathcal B\otimes\mathcal L^{\otimes n})
\longrightarrow
\Gamma(X,(\mathcal J\mathcal B/\mathcal J\mathcal K\mathcal B)
\otimes\mathcal L^{\otimes n})
\]
이 전사임을 얻는다.

$i:Z\to X$, $i':Z'\to X'$을 표준 포함들이라 하자. 그러면 가환 도식
\[
\xymatrix{
X'\ar[d]_g & Z'\ar[l]_{i'}\ar[d]^{h}\\
X & Z\ar[l]^{i}
}
\]
이 있다. $\mathcal M=i^*(\mathcal L)$,
$\mathcal M'=i'^*(\mathcal L')$로 놓자. $\mathcal L'$이 풍부하므로
$\mathcal M'$도 풍부하고 (II, 5.1.12), 한편
$\mathcal M'=h^*(\mathcal M)$이다. 따라서 뇌터 귀납가정에 의해
($Z\neq X$이므로) $\mathcal M$은 풍부하다. 그러므로 충분히 큰
$n$에 대해
$i^*(\mathcal J/\mathcal J\mathcal K)\otimes\mathcal M^{\otimes n}$은
$Z$ 위의 단면들로 생성된다 (II, 4.5.5). 또한
$\mathcal J/\mathcal J\mathcal K
=i_*(i^*(\mathcal J/\mathcal J\mathcal K))$이므로,
(0\textsubscript{I}, 5.4.10)에 의해 충분히 큰 어떤 $n$에 대해
$X$ 위의
$(\mathcal J/\mathcal J\mathcal K)\otimes\mathcal L^{\otimes n}$의
단면 $s$가 존재하여 $s(x)\neq0$이다. 실제로 $\mathcal J$의 정의에
의해 $\mathcal J_x=\mathcal A_x$이고, 가정에 의해
$\mathcal K_x\neq\mathcal A_x$이다. 도식
(\hyperref[III.2.6.2.6-fr]{2.6.2.6})에 따르면 $s$는 또한 $X$ 위의
\[
(\mathcal J\mathcal B/\mathcal J\mathcal K\mathcal B)
\otimes\mathcal L^{\otimes n}
\]
의 단면이다. 따라서 $s$는 $X$ 위의
$(\mathcal J\mathcal B)\otimes\mathcal L^{\otimes n}$의 어떤 단면
$t$의 표준 상이다. 그런데 정의에 의해
$(\mathcal J\mathcal K\mathcal B)\otimes\mathcal L^{\otimes n}$을
법으로 한 $t$의 표준 상 $s$는
\[
(\mathcal J\otimes\mathcal L^{\otimes n})/
(\mathcal J\mathcal K\mathcal B\otimes\mathcal L^{\otimes n})
\]
에 속한다. 따라서 (\hyperref[III.2.6.2.6-fr]{2.6.2.6})에 의해
$t$는 $X$ 위의 $\mathcal J\otimes\mathcal L^{\otimes n}$의 단면이고,
더구나 $\mathcal L^{\otimes n}$의 단면이다. 위에서 $t(x)\neq0$임을
보았으므로 $x\in X_t$이다. 또한 $\mathcal J$의 정의에 의해
$\mathcal O_X/\mathcal J$의 받침인 $X-U$에서 $t(y)=0$이다. 따라서
$X_t\subset U$이고, 이로써 증명이 끝난다.
\end{env}
\end{proof}

\oldpage[III]{115}
\begin{env}[2.6.3]
\label{III.2.6.3-fr}
\emph{비고.} $X$가 $Y$ 위에서 고유이면, 슈발레 정리 (II, 6.7.1)에서와
같이 (\hyperref[III.2.6.1-fr]{2.6.1})과 보조정리 (II, 6.7.1.1)을
사용하여 (\hyperref[III.2.6.2-fr]{2.6.2})를 더 간단히 증명할 수 있다.
\end{env}

\section{고유 사상의 유한성 정리}
\label{section:III.3-fr}

\subsection{데비사주 보조정리}
\label{subsection:III.3.1-fr}

\begin{definition}[3.1.1]
\label{III.3.1.1-fr}
$K$를 아벨 범주라 하자. $K$의 대상들의 집합의 부분집합 $K'$이
$0\in K'$을 만족하고, $K$ 안의 모든 완전열
\[
0\longrightarrow A'\longrightarrow A\longrightarrow A''\longrightarrow0
\]
에서 $A$, $A'$, $A''$ 가운데 둘이 $K'$에 속하면 나머지 하나도
$K'$에 속할 때, $K'$을 \emph{완전}하다고 한다.
\end{definition}

\begin{theorem}[3.1.2]
\label{III.3.1.2-fr}
$X$를 뇌터 준스킴이라 하고, $K$로 연접 $\mathcal O_X$-가군들의 아벨
범주를 나타내자. $K'$을 $K$의 완전한 부분집합, $X'$을 $X$의
바탕공간의 닫힌 부분집합이라 하자. 일반점이 $y$인 $X'$의 모든 기약
닫힌 부분집합 $Y$에 대하여, 받침이 $Y$이고 $\mathcal G_y$가 차원
$1$인 $k(y)$-벡터공간인 어떤 $\mathcal O_X$-가군
$\mathcal G\in K'$이 존재한다고 가정하자. 그러면 받침이 $X'$에
포함되는 모든 연접 $\mathcal O_X$-가군은 $K'$에 속한다. 특히
$X'=X$이면 $K'=K$이다.
\end{theorem}

\begin{proof}
$X'$의 닫힌 부분집합 $Y$에 대한 다음 성질을 $P(Y)$라 하자. 받침이
$Y$에 포함되는 모든 연접 $\mathcal O_X$-가군이 $K'$에 속한다.
뇌터 귀납법의 원리 (0\textsubscript{I}, 2.2.2)에 의해 다음을
증명하면 충분하다. $Y$가 $X'$의 닫힌 부분집합이고, $Y$와 다른
$Y$의 모든 닫힌 부분집합 $Y'$에 대해 $P(Y')$이 참이면 $P(Y)$도
참이다.

따라서 받침이 $Y$에 포함되는 $\mathcal F\in K$를 택하고
$\mathcal F\in K'$임을 증명하자. 계속해서 $Y$로 바탕공간이 $Y$인
$X$의 축소 닫힌 부분준스킴을 나타내자 (I, 5.2.1). 이는
$\mathcal O_X$의 어떤 연접 아이디얼 $\mathcal J$로 정의된다.
$\mathcal J^n\mathcal F=0$이 되도록 하는 어떤 정수 $n>0$이 존재한다
(I, 9.3.4). 따라서 $1\leq k\leq n$에 대해 연접
$\mathcal O_X$-가군들의 완전열
\[
0\longrightarrow
\mathcal J^{k-1}\mathcal F/\mathcal J^k\mathcal F
\longrightarrow \mathcal F/\mathcal J^k\mathcal F
\longrightarrow \mathcal F/\mathcal J^{k-1}\mathcal F
\longrightarrow0
\]
이 있다 (0\textsubscript{I}, 5.3.6 및 5.3.3). $K'$이 완전하므로
$k$에 대한 귀납법에 의해 각
$\mathcal F_k=\mathcal J^{k-1}\mathcal F/\mathcal J^k\mathcal F$가
$K'$에 속함을 보이면 충분하다. 따라서 추가로
$\mathcal J\mathcal F=0$이라고 가정하여 $\mathcal F\in K'$임을
증명하는 문제로 귀착된다. 이는 $j:Y\to X$가 표준 포함일 때
$\mathcal F=j_*(j^*(\mathcal F))$라고 말하는 것과 같다. 이제 두
경우로 나눈다.
\begin{enumerate}
\item[a)] $Y$가 \emph{가약}인 경우. $Y=Y'\cup Y''$라 하자. 여기서
$Y'$, $Y''$은 $Y$와 다른 $Y$의 닫힌 부분집합들이다. 계속해서
$Y'$, $Y''$로 각각 바탕공간이 $Y'$, $Y''$인 $X$의 축소 닫힌
부분준스킴들도 나타내고, 이를 각각 정의하는 $\mathcal O_X$의 연접
아이디얼들을 $\mathcal J'$, $\mathcal J''$라 하자. 다음과 같이 놓자.
\[
\mathcal F'=\mathcal F\otimes_{\mathcal O_X}
(\mathcal O_X/\mathcal J'),\qquad
\mathcal F''=\mathcal F\otimes_{\mathcal O_X}
(\mathcal O_X/\mathcal J'').
\]
표준 준동형 $\mathcal F\to\mathcal F'$,
$\mathcal F\to\mathcal F''$은 준동형
$u:\mathcal F\to\mathcal F'\oplus\mathcal F''$를 정의한다.
$z\notin Y'\cap Y''$이면 준동형
$u_z:\mathcal F_z\to\mathcal F'_z\oplus\mathcal F''_z$가
\emph{전단사}임을 보이자. 실제로
$\mathcal J'\cap\mathcal J''=\mathcal J$이다. 문제는 국소적이고

\oldpage[III]{116}
이 등식은 (I, 5.2.1 및 1.1.5)에서 따르기 때문이다. $z\notin Y''$이면
$\mathcal J'_z=\mathcal J_z$이고, 따라서
$\mathcal F'_z=\mathcal F_z$, $\mathcal F''_z=0$이다. 이는 이 경우의
주장을 증명하며, $z\notin Y'$인 경우도 마찬가지이다. 그러므로 $K$에
속하는 $u$의 핵과 여핵은 (0\textsubscript{I}, 5.3.4) 받침이
$Y'\cap Y''$에 포함되어 귀납가정에 의해 $K'$에 속한다. 같은 이유로
$\mathcal F'$, $\mathcal F''$도 $K'$에 속한다. $K'$이 완전하므로
$\mathcal F'\oplus\mathcal F''$도 $K'$에 속한다. 이제 두 완전열
\[
0\longrightarrow\operatorname{Im}u\longrightarrow
\mathcal F'\oplus\mathcal F''\longrightarrow
\operatorname{Coker}u\longrightarrow0,
\]
\[
0\longrightarrow\operatorname{Ker}u\longrightarrow\mathcal F
\longrightarrow\operatorname{Im}u\longrightarrow0
\]
과 $K'$의 완전성에서 결론이 따른다.

\item[b)] $Y$가 기약인 경우. 따라서 $X$의 부분준스킴 $Y$는
\emph{정역}이다. 그 일반점을 $y$라 하면
$(\mathcal O_Y)_y=k(y)$이다. $j^*(\mathcal F)$가 연접
$\mathcal O_Y$-가군이므로
$\mathcal F_y=(j^*(\mathcal F))_y$는 어떤 유한 차원 $m$의
$k(y)$-벡터공간이다. 가정에 의해 어떤 연접 $\mathcal O_X$-가군
$\mathcal G$가 존재하여($\mathcal G$의 받침은 반드시 $Y$이다),
$\mathcal G_y$는 차원 $1$인 $k(y)$-벡터공간이다. 따라서
$k(y)$-동형사상
$(\mathcal G_y)^m\xrightarrow{\sim}\mathcal F_y$가 존재하며, 이는
$\mathcal O_Y$-동형사상이기도 하다. $\mathcal G^m$과 $\mathcal F$가
연접이므로, $X$에서 $y$의 어떤 열린 근방 $W$와 동형사상
$\mathcal G^m|W\xrightarrow{\sim}\mathcal F|W$가 존재한다
(0\textsubscript{I}, 5.2.7). 이 동형사상의 그래프를 $\mathcal H$라
하자. 이는 $(\mathcal G^m\oplus\mathcal F)|W$의 연접 부분
$(\mathcal O_X|W)$-가군이고, 표준적으로 $\mathcal G^m|W$ 및
$\mathcal F|W$와 동형이다. 따라서 $\mathcal G^m\oplus\mathcal F$의
어떤 연접 부분 $\mathcal O_X$-가군 $\mathcal H_0$가 존재하여,
$W$ 위에서는 $\mathcal H$를 유도하고 $X-Y$ 위에서는 0을 유도한다.
이는 $\mathcal G^m$과 $\mathcal F$의 받침이 $Y$이기 때문이다
(I, 9.4.7).

$\mathcal G^m\oplus\mathcal F$의 표준 사영들을 제한하여 얻는
$v:\mathcal H_0\to\mathcal G^m$과
$w:\mathcal H_0\to\mathcal F$는 연접 $\mathcal O_X$-가군들의
준동형이고, $W$ 및 $X-Y$에서는 동형사상이다. 다시 말해 $v$, $w$의
핵과 여핵들의 받침은 $Y$와 다른 닫힌집합 $Y-(Y\cap W)$에 포함된다.
따라서 이들은 $K'$에 속한다. 한편 $\mathcal G\in K'$이고 $K'$이
완전하므로 $\mathcal G^m\in K'$이다. $K'$의 완전성에 의해 차례로
$\mathcal H_0\in K'$, 이어서 $\mathcal F\in K'$임을 얻는다. 이로써
증명이 끝난다.
\end{enumerate}
\end{proof}

\begin{corollary}[3.1.3]
\label{III.3.1.3-fr}
$K$의 완전한 부분집합 $K'$이 다음 성질도 만족한다고 가정하자.
$\mathcal M\in K'$인 연접 $\mathcal O_X$-가군의 모든 연접 직합인자는
다시 $K'$에 속한다. 이 조건 아래에서
(\hyperref[III.3.1.2-fr]{3.1.2})의 조건
``$\mathcal G_y$는 차원 $1$인 $k(y)$-벡터공간이다''를
$\mathcal G_y\neq0$으로 바꾸어도 같은 결론이 성립한다. 마지막 조건은
$\operatorname{Supp}(\mathcal G)=Y$와 동치이다.
\end{corollary}

\begin{proof}
실제로 (\hyperref[III.3.1.2-fr]{3.1.2})의 논증은 경우 b)에서만
바꾸면 된다. 이번에는 $\mathcal G_y$가 어떤 차원 $q>0$의
$k(y)$-벡터공간이므로
$(\mathcal G_y)^m\xrightarrow{\sim}(\mathcal F_y)^q$인
$\mathcal O_Y$-동형사상이 존재한다. 그러면
(\hyperref[III.3.1.2-fr]{3.1.2})의 논증 끝부분은
$\mathcal F^q\in K'$임을 증명하고, $K'$에 대한 추가 가정으로부터
$\mathcal F\in K'$이 따른다.
\end{proof}

\subsection{유한성 정리: 통상적 스킴의 경우}
\label{subsection:III.3.2-fr}

\begin{theorem}[3.2.1]
\label{III.3.2.1-fr}
$Y$를 국소 뇌터 준스킴, $f:X\to Y$를 고유 사상이라 하자. 모든 연접
$\mathcal O_X$-가군 $\mathcal F$에 대해 $q\geq0$인
$\mathcal O_Y$-가군 $R^qf_*(\mathcal F)$들은 연접이다.
\end{theorem}

\begin{proof}
문제는 $Y$ 위에서 국소적이므로 $Y$가 뇌터라고 가정할 수 있고,
따라서 $X$도 뇌터이다 (I, 6.3.7). 정리
(\hyperref[III.3.2.1-fr]{3.2.1})의 결론이 성립하는 연접
$\mathcal O_X$-가군 $\mathcal F$들은 연접 $\mathcal O_X$-가군들의
범주 $K$의 완전한 부분집합 $K'$을 이룬다.

\oldpage[III]{117}
실제로
$0\to\mathcal F'\to\mathcal F\to\mathcal F''\to0$을 연접
$\mathcal O_X$-가군들의 완전열이라 하자. 예를 들어
$\mathcal F'$, $\mathcal F''$이 $K'$에 속한다고 가정하면 코호몰로지
완전열
\[
R^{q-1}f_*(\mathcal F'')\xrightarrow{\partial}R^qf_*(\mathcal F')
\longrightarrow R^qf_*(\mathcal F)\longrightarrow R^qf_*(\mathcal F'')
\xrightarrow{\partial}R^{q+1}f_*(\mathcal F')
\]
에서 가정에 의해 바깥 네 항이 연접이다. 따라서 가운데 항
$R^qf_*(\mathcal F)$도 연접이다
(0\textsubscript{I}, 5.3.4 및 5.3.3). 같은 방식으로
$\mathcal F$, $\mathcal F'$ (각각 $\mathcal F$, $\mathcal F''$)이
$K'$에 속하면 $\mathcal F''$ (각각 $\mathcal F'$)도 $K'$에 속함을
보인다. 더 나아가 $\mathcal F\in K'$인 $\mathcal O_X$-가군의 모든
연접 \emph{직합인자} $\mathcal F'$도 $K'$에 속한다. 실제로 이때
$R^qf_*(\mathcal F')$은 $R^qf_*(\mathcal F)$의 직합인자이고
(G, II, 4.4.4), 따라서 유한 생성이다. 또한
(\hyperref[III.1.4.10-fr]{1.4.10})에 의해 준연접이고 $Y$가
뇌터이므로 연접이다.

(\hyperref[III.3.1.3-fr]{3.1.3})에 의해 다음을 증명하는 문제로
귀착된다. $X$가 일반점 $x$를 갖는 \emph{기약} 준스킴이면,
$\mathcal F_x\neq0$이고 $K'$에 속하는 어떤 연접
$\mathcal O_X$-가군 $\mathcal F$가 존재한다. 실제로 이 사실이
증명되면 $X$의 모든 기약 닫힌 부분준스킴 $Y$에 적용할 수 있다.
$j:Y\to X$가 표준 포함이면 $f\circ j$는 고유이고 (II, 5.4.2),
$\mathcal G$가 받침이 $Y$인 연접 $\mathcal O_Y$-가군일 때
$j_*(\mathcal G)$는 연접 $\mathcal O_X$-가군이며
\[
R^q(f\circ j)_*(\mathcal G)=R^qf_*(j_*(\mathcal G))
\]
이다 (G, II, 4.9.1). 따라서
(\hyperref[III.3.1.3-fr]{3.1.3})을 적용할 조건을 만족한다.

이제 차우 보조정리 (II, 5.6.2)에 의해 어떤 기약 준스킴 $X'$과
\emph{사영 전사} 사상 $g:X'\to X$가 존재하여
$f\circ g:X'\to Y$는 \emph{사영}이다. $g$에 대해 풍부한 어떤
$\mathcal O_{X'}$-가군 $\mathcal L$이 존재한다 (II, 5.3.1).
$g:X'\to X$와 $\mathcal L$에 사영 사상의 기본정리
(\hyperref[III.2.2.1-fr]{2.2.1})을 적용하자. 그러면 어떤 정수 $n$이
존재하여
\[
\mathcal F=g_*(\mathcal O_{X'}(n))
\]
은 연접 $\mathcal O_X$-가군이고 모든 $q>0$에 대해
$R^qg_*(\mathcal O_{X'}(n))=0$이다. 또한 충분히 큰 $n$에 대해
\[
g^*(g_*(\mathcal O_{X'}(n)))\longrightarrow\mathcal O_{X'}(n)
\]
이 전사이므로 (\hyperref[III.2.2.1-fr]{2.2.1}), $X$의 일반점 $x$에서
$\mathcal F_x\neq0$이라고 가정할 수 있다 (II, 3.4.7). 한편
$f\circ g$가 사영이고 $Y$가 뇌터이므로
$R^p(f\circ g)_*(\mathcal O_{X'}(n))$들은 \emph{연접}이다
(\hyperref[III.2.2.1-fr]{2.2.1}). 이제
$R^*(f\circ g)_*(\mathcal O_{X'}(n))$은 $E_2$항이
\[
E_2^{pq}=R^pf_*(R^qg_*(\mathcal O_{X'}(n)))
\]
인 르레 스펙트럼 열의 수렴 대상이다. 앞의 결과에 의해 이 스펙트럼
열은 퇴화하고, 따라서 (0, 11.1.6)에 의해
$E_2^{p0}=R^pf_*(\mathcal F)$는
$R^p(f\circ g)_*(\mathcal O_{X'}(n))$과 동형이다. 이로써 증명이
끝난다.
\end{proof}

\begin{corollary}[3.2.2]
\label{III.3.2.2-fr}
$Y$를 국소 뇌터 준스킴이라 하자. 모든 고유 사상 $f:X\to Y$에
대하여 연접 $\mathcal O_X$-가군의 $f$에 의한 직접상은 연접
$\mathcal O_Y$-가군이다.
\end{corollary}

\begin{corollary}[3.2.3]
\label{III.3.2.3-fr}
$A$를 뇌터 환, $X$를 $A$ 위의 고유 스킴이라 하자. 모든 연접
$\mathcal O_X$-가군 $\mathcal F$에 대해 $H^p(X,\mathcal F)$는 유한
생성 $A$-가군이다. 또한 어떤 정수 $r>0$이 존재하여 모든 연접
$\mathcal O_X$-가군 $\mathcal F$와 모든 $p>r$에 대해
$H^p(X,\mathcal F)=0$이다.
\end{corollary}

\begin{proof}
둘째 명제는 이미 증명하였다
(\hyperref[III.1.4.12-fr]{1.4.12}). 첫째 명제는
(\hyperref[III.1.4.11-fr]{1.4.11})을 고려하면 유한성 정리
(\hyperref[III.3.2.1-fr]{3.2.1})에서 따른다.
\end{proof}

특히 $X$가 체 $k$ 위의 \emph{고유 대수적 스킴}이면, 모든 연접
$\mathcal O_X$-가군 $\mathcal F$에 대해 $H^p(X,\mathcal F)$들은
\emph{유한 차원} $k$-벡터공간이다.

\begin{corollary}[3.2.4]
\label{III.3.2.4-fr}
$Y$를 국소 뇌터 준스킴, $f:X\to Y$를 유한형 사상이라 하자. 받침이
$Y$ 위에서 고유인 모든 연접 $\mathcal O_X$-가군 $\mathcal F$에 대해
(II, 5.4.10), $\mathcal O_Y$-가군 $R^qf_*(\mathcal F)$들은 연접이다.
\end{corollary}

\oldpage[III]{118}
\begin{proof}
문제는 $Y$ 위에서 국소적이므로 $Y$가 뇌터라고 가정할 수 있고,
따라서 $X$도 뇌터이다 (I, 6.3.7). 가정에 의해 바탕공간이
$\operatorname{Supp}(\mathcal F)$인 $X$의 모든 닫힌 부분준스킴
$Z$는 $Y$ 위에서 고유이다. 다시 말해 $j:Z\to X$가 표준 포함이면
$f\circ j:Z\to Y$는 고유이다. 한편
$\mathcal F=j_*(\mathcal G)$가 되도록 $Z$를 택할 수 있으며, 여기서
$\mathcal G=j^*(\mathcal F)$는 연접 $\mathcal O_Z$-가군이다
(I, 9.3.5). (\hyperref[III.1.3.4-fr]{1.3.4})에 의해
$R^qf_*(\mathcal F)=R^q(f\circ j)_*(\mathcal G)$이므로, 결론은
(\hyperref[III.3.2.1-fr]{3.2.1})에서 즉시 따른다.
\end{proof}

\subsection{유한성 정리의 일반화 (통상적 스킴)}
\label{subsection:III.3.3-fr}

\begin{proposition}[3.3.1]
\label{III.3.3.1-fr}
$Y$를 뇌터 준스킴, $\mathcal S$를 양의 차수들을 갖는 준연접 유한
생성 등급 $\mathcal O_Y$-대수,
$Y'=\operatorname{Proj}(\mathcal S)$라 하고, $g:Y'\to Y$를 구조
사상이라 하자. $f:X\to Y$를 고유 사상,
$\mathcal S'=f^*(\mathcal S)$,
$\mathcal M=\bigoplus_{k\in\mathbb Z}\mathcal M_k$를 유한 생성 준연접
등급 $\mathcal S'$-가군이라 하자. 그러면 모든 $p$에 대해
\[
R^pf_*(\mathcal M)=\bigoplus_{k\in\mathbb Z}R^pf_*(\mathcal M_k)
\]
는 유한 생성 등급 $\mathcal S$-가군이다. 더 나아가 $\mathcal S$가
$\mathcal S_1$로 생성된다고 가정하자. 그러면 각 $p\in\mathbb Z$에
대해 어떤 정수 $k_p$가 존재하여, 모든 $k\geq k_p$, $r\geq0$에 대해
\[
R^pf_*(\mathcal M_{k+r})=\mathcal S_rR^pf_*(\mathcal M_k)
\tag{3.3.1.1}\label{III.3.3.1.1-fr}
\]
이다.
\end{proposition}

\begin{proof}
첫째 명제는 (\hyperref[III.2.4.1-fr]{2.4.1}, (i))의 명제에서 단지
``사영 사상''을 ``고유 사상''으로 바꾼 것이다. 그런데
(\hyperref[III.2.4.1-fr]{2.4.1}, (i))의 증명에서 $f$에 대한 가정은
그 증명의 표기로 $R^pf'_*(\widetilde{\mathcal M})$이 연접
$\mathcal O_{Y'}$-가군임을 보이는 데에만 사용되었다.
(\hyperref[III.3.3.1-fr]{3.3.1})의 가정 아래에서는 $f'$이 고유이므로
(II, 5.4.2, (iii)), 유한성 정리
(\hyperref[III.3.2.1-fr]{3.2.1})에 의해
(\hyperref[III.2.4.1-fr]{2.4.1}, (i))의 증명을 그대로 반복할 수 있다.

둘째 명제에 대해서는, $Y$의 어떤 유한 아핀 열린 덮개 $(U_i)$가
존재하여 모든 $k\geq k_{p,i}$에 대해
(\hyperref[III.3.3.1.1-fr]{3.3.1.1})의 두 변을 $U_i$에 제한한 것들이
같음을 유의하면 충분하다 (II, 2.1.6, (ii)). $k_{p,i}$들 가운데 가장
큰 것을 $k_p$로 택하면 된다.
\end{proof}

\begin{corollary}[3.3.2]
\label{III.3.3.2-fr}
$A$를 뇌터 환, $\mathfrak m$을 $A$의 아이디얼, $X$를 고유
$A$-스킴, $\mathcal F$를 연접 $\mathcal O_X$-가군이라 하자. 그러면
모든 $p\geq0$에 대해 직합
\[
\bigoplus_{k\geq0}H^p(X,\mathfrak m^k\mathcal F)
\]
은 환 $S=\bigoplus_{k\geq0}\mathfrak m^k$ 위의 유한 생성 가군이다.
특히 어떤 정수 $k_p\geq0$이 존재하여 모든 $k\geq k_p$, $r\geq0$에
대해
\[
H^p(X,\mathfrak m^{k+r}\mathcal F)
=\mathfrak m^rH^p(X,\mathfrak m^k\mathcal F)
\tag{3.3.2.1}\label{III.3.3.2.1-fr}
\]
이다.
\end{corollary}

\begin{proof}
$Y=\operatorname{Spec}(A)$, $\mathcal S=\widetilde S$,
$\mathcal M_k=\mathfrak m^k\mathcal F$로 놓고
(\hyperref[III.1.4.11-fr]{1.4.11})을 고려하여
(\hyperref[III.3.3.1-fr]{3.3.1})을 적용하면 된다.
\end{proof}

$\bigoplus_{k\geq0}H^p(X,\mathfrak m^k\mathcal F)$의 $S$-가군 구조는
다음과 같이 얻어짐을 상기하자. 모든 $a\in\mathfrak m^r$에 대해
곱셈
$\mathfrak m^k\mathcal F\to\mathfrak m^{k+r}\mathcal F$에서
코호몰로지로 넘어가 얻는 사상
\[
H^p(X,\mathfrak m^k\mathcal F)
\longrightarrow H^p(X,\mathfrak m^{k+r}\mathcal F)
\]
을 생각한다 (\hyperref[III.2.4.1-fr]{2.4.1}).

\oldpage[III]{119}
\subsection{유한성 정리: 형식적 스킴의 경우}
\label{subsection:III.3.4-fr}

이 절의 결과들 가운데 정의
(\hyperref[III.3.4.1-fr]{3.4.1})를 제외한 것들은 이 장의 나머지
부분에서 사용하지 않는다.

\begin{env}[3.4.1]
\label{III.3.4.1-fr}
$\mathfrak X$와 $\mathfrak S$를 국소 뇌터 형식적 준스킴들
(I, 10.4.2), $f:\mathfrak X\to\mathfrak S$를 형식적 준스킴의 사상이라
하자. 다음 조건들을 만족할 때 $f$를 \emph{고유 사상}이라고 한다.
\begin{enumerate}
\item[$1^\circ$] $f$는 유한형 사상이다 (I, 10.13.3).
\item[$2^\circ$] $\mathcal K$가 $\mathfrak S$의 정의 아이디얼이고
\[
\mathcal J=f^*(\mathcal K)\mathcal O_{\mathfrak X},\qquad
X_0=(\mathfrak X,\mathcal O_{\mathfrak X}/\mathcal J),\qquad
S_0=(\mathfrak S,\mathcal O_{\mathfrak S}/\mathcal K)
\]
로 놓으면, $f$에서 유도되는 사상 $f_0:X_0\to S_0$ (I, 10.5.6)는
고유이다.
\end{enumerate}

이 정의가 고려한 $\mathfrak S$의 정의 아이디얼 $\mathcal K$에 의존하지
않음은 즉시 알 수 있다. 실제로 $\mathcal K'\subset\mathcal K$인 둘째
정의 아이디얼 $\mathcal K'$에 대하여
$\mathcal J'=f^*(\mathcal K')\mathcal O_{\mathfrak X}$,
$X'_0=(\mathfrak X,\mathcal O_{\mathfrak X}/\mathcal J')$,
$S'_0=(\mathfrak S,\mathcal O_{\mathfrak S}/\mathcal K')$로 놓으면,
$f$에서 유도되는 사상 $f'_0:X'_0\to S'_0$에 대한 도표
\[
\xymatrix{
X_0\ar[r]^{f_0}\ar[d]_i & S_0\ar[d]^j\\
X'_0\ar[r]_{f'_0} & S'_0
}
\]
는 가환이고 $i$와 $j$는 전사인 몰입 사상이다. 따라서 (II, 5.4.5)에
의해 $f_0$가 고유라는 것과 $f'_0$가 고유라는 것은 동치이다.

모든 $n\geq0$에 대하여
$X_n=(\mathfrak X,\mathcal O_{\mathfrak X}/\mathcal J^{n+1})$,
$S_n=(\mathfrak S,\mathcal O_{\mathfrak S}/\mathcal K^{n+1})$로 놓자.
$f$에서 유도되는 사상 $f_n:X_n\to S_n$ (I, 10.5.6)는 $n=0$일 때
고유이면 모든 $n$에 대해 고유이다 (II, 5.4.6).

$g:Y\to Z$가 국소 뇌터인 통상적 준스킴들의 고유 사상이고, $Z'$가
$Z$의 닫힌 부분집합, $Y'$가 $g(Y')\subset Z'$를 만족하는 $Y$의 닫힌
부분집합이면, $g$를 완비화들로 연장한
$\widehat g:Y_{/Y'}\to Z_{/Z'}$ (I, 10.9.1)는 정의와 (II, 5.4.5)에
의해 형식적 준스킴들의 고유 사상이다.

$\mathfrak X$와 $\mathfrak S$를 국소 뇌터 형식적 준스킴들,
$f:\mathfrak X\to\mathfrak S$를 유한형 사상이라 하자 (I, 10.13.3).
표기는 위와 같이 하고, $\mathfrak X$의 바탕공간의 부분집합 $Z$를
$X_0$의 부분집합으로 보았을 때 $S_0$ 위에서 고유이면
(II, 5.4.10), $Z$가 \emph{$\mathfrak S$ 위에서 고유하다} (또는
$f$에 대해 고유하다)고 한다. 통상적 준스킴의 고유한 부분집합에 관하여
(II, 5.4.10)에서 서술한 모든 성질은 정의에서 즉시 알 수 있듯이 형식적
준스킴의 고유한 부분집합에도 그대로 성립한다.
\end{env}

\begin{theorem}[3.4.2]
\label{III.3.4.2-fr}
$\mathfrak X$, $\mathfrak Y$를 국소 뇌터 형식적 준스킴들,
$f:\mathfrak X\to\mathfrak Y$를 고유 사상이라 하자. 모든 연접
$\mathcal O_{\mathfrak X}$-가군 $\mathcal F$에 대하여
$\mathcal O_{\mathfrak Y}$-가군 $R^qf_*(\mathcal F)$들은 모든
$q\geq0$에 대해 연접이다.
\end{theorem}

$\mathcal J$를 $\mathfrak Y$의 정의 아이디얼,
$\mathcal K=f^*(\mathcal J)\mathcal O_{\mathfrak X}$라 하고,
$\mathcal O_{\mathfrak X}$-가군
\[
\mathcal F_k=\mathcal F\otimes_{\mathcal O_{\mathfrak Y}}
(\mathcal O_{\mathfrak Y}/\mathcal J^{k+1})
=\mathcal F/\mathcal K^{k+1}\mathcal F\qquad(k\geq0)
\tag{3.4.2.1}\label{III.3.4.2.1-fr}
\]
들을 생각하자. 이들은 명백히 위상 $\mathcal O_{\mathfrak X}$-가군의
사영계를 이루고,

\oldpage[III]{120}
\[
\mathcal F=\varprojlim_k\mathcal F_k
\]
이다 (I, 10.11.3). 한편
(\hyperref[III.3.4.2-fr]{3.4.2})에 의해 각
$R^qf_*(\mathcal F)$는 연접이므로 자연스럽게 위상
$\mathcal O_{\mathfrak Y}$-가군의 구조를 가지며 (I, 10.11.6),
$R^qf_*(\mathcal F_k)$도 마찬가지이다. 표준 준동형사상
$\mathcal F\to\mathcal F_k=\mathcal F/\mathcal K^{k+1}\mathcal F$에는
표준 준동형사상
\[
R^qf_*(\mathcal F)\longrightarrow R^qf_*(\mathcal F_k)
\]
들이 대응한다. 이들은 앞의 위상 $\mathcal O_{\mathfrak Y}$-가군 구조에
관하여 반드시 연속이고 (I, 10.11.6), 사영계를 이룬다. 따라서 극한을
취하면 함자적인 표준 준동형사상
\[
R^qf_*(\mathcal F)\longrightarrow\varprojlim_kR^qf_*(\mathcal F_k)
\tag{3.4.2.2}\label{III.3.4.2.2-fr}
\]
을 얻으며, 이는 위상 $\mathcal O_{\mathfrak Y}$-가군의 연속
준동형사상이다. 우리는 (\hyperref[III.3.4.2-fr]{3.4.2})와 함께 다음을
증명할 것이다.

\begin{corollary}[3.4.3]
\label{III.3.4.3-fr}
각 준동형사상 (\hyperref[III.3.4.2.2-fr]{3.4.2.2})는 위상
동형사상이다. 더 나아가 $\mathfrak Y$가 뇌터이면 사영계
$(R^qf_*(\mathcal F/\mathcal K^{k+1}\mathcal F))_{k\geq0}$는
조건 (ML)을 만족한다 (0, 13.1.1).
\end{corollary}

먼저 $\mathfrak Y$가 뇌터 아핀 형식적 스킴인 경우 (I, 10.4.1)에
(\hyperref[III.3.4.2-fr]{3.4.2})와
(\hyperref[III.3.4.3-fr]{3.4.3})를 증명한다.

\begin{corollary}[3.4.4]
\label{III.3.4.4-fr}
(\hyperref[III.3.4.2-fr]{3.4.2})의 가정들에 더하여
$\mathfrak Y=\operatorname{Spf}(A)$이고 $A$가 뇌터 아딕 환이라고
가정하자. $\mathfrak J$를 $A$의 정의 아이디얼이라 하고 $k\geq0$에
대해 $\mathcal F_k=\mathcal F/\mathfrak J^{k+1}\mathcal F$로 놓자.
그러면 $H^n(\mathfrak X,\mathcal F)$들은 유한 생성 $A$-가군이고,
사영계 $(H^n(\mathfrak X,\mathcal F_k))_{k\geq0}$는 모든 $n$에 대해
조건 (ML)을 만족한다. 또한
\[
N_{n,k}=\operatorname{Ker}\bigl(H^n(\mathfrak X,\mathcal F)
\longrightarrow H^n(\mathfrak X,\mathcal F_k)\bigr)
\tag{3.4.4.1}\label{III.3.4.4.1-fr}
\]
로 놓으면, 이는 코호몰로지 완전열에 의해
$\operatorname{Im}(H^n(\mathfrak X,\mathfrak J^{k+1}\mathcal F)
\to H^n(\mathfrak X,\mathcal F))$와도 같다. $N_{n,k}$들은
$H^n(\mathfrak X,\mathcal F)$ 위에 $\mathfrak J$-좋은 여과를 정의한다
(0, 13.7.7). 끝으로 모든 $n$에 대해 표준 준동형사상
\[
H^n(\mathfrak X,\mathcal F)
\longrightarrow\varprojlim_kH^n(\mathfrak X,\mathcal F_k)
\tag{3.4.4.2}\label{III.3.4.4.2-fr}
\]
은 위상 동형사상이다. 여기서 왼쪽에는 $\mathfrak J$-아딕 위상을,
$H^n(\mathfrak X,\mathcal F_k)$들에는 이산 위상을 준다.
\end{corollary}

다음과 같이 놓자.
\[
S=\operatorname{gr}(A)=\bigoplus_{k\geq0}
\mathfrak J^k/\mathfrak J^{k+1},\qquad
\mathcal M=\operatorname{gr}(\mathcal F)=\bigoplus_{k\geq0}
\mathfrak J^k\mathcal F/\mathfrak J^{k+1}\mathcal F.
\tag{3.4.4.3}\label{III.3.4.4.3-fr}
\]
$\mathfrak J^\Delta$가 $\mathfrak Y$의 정의 아이디얼임은 알려져 있다
(I, 10.3.1). 이제
\[
\mathcal K=f^*(\mathfrak J^\Delta)\mathcal O_{\mathfrak X},\qquad
X_0=(\mathfrak X,\mathcal O_{\mathfrak X}/\mathcal K),\qquad
Y_0=(\mathfrak Y,\mathcal O_{\mathfrak Y}/\mathfrak J^\Delta)
=\operatorname{Spec}(A_0)
\]
라 하자. 여기서 $A_0=A/\mathfrak J$이다. 가군
$\mathcal M_k=\mathfrak J^k\mathcal F/\mathfrak J^{k+1}\mathcal F$들이
연접 $\mathcal O_{X_0}$-가군임은 명백하다 (I, 10.11.3). 한편 준연접
등급 $\mathcal O_{X_0}$-대수
\[
\mathcal S=\mathcal O_{X_0}\otimes_{A_0}S
\tag{3.4.4.4}\label{III.3.4.4.4-fr}
\]
를 생각하자. 곱셈은 등급 대수의 전사 준동형사상
\[
\mathcal S\longrightarrow
\operatorname{gr}(\mathcal O_{\mathfrak X})
=\bigoplus_{k\geq0}\mathcal K^k/\mathcal K^{k+1}
\]
을 정의하며, 이를 통해 $\mathcal M$은 등급 $\mathcal S$-가군이 된다.

$\mathcal F$가 유한 생성 $\mathcal O_{\mathfrak X}$-가군이라는 가정에서
먼저 $\mathcal M$이 유한 생성 등급 $\mathcal S$-가군임이 따른다.

\oldpage[III]{121}
실제로 이 문제는 $\mathfrak X$ 위에서 국소적이므로, 이를 다룰 때
$\mathfrak X=\operatorname{Spf}(B)$라고 가정할 수 있다. 여기서 $B$는
뇌터 아딕 환이고 $\mathcal F=N^\Delta$이며, $N$은 유한 생성
$B$-가군이다 (I, 10.10.5). $B_0=B/\mathfrak J B$,
$N_0=N/\mathfrak JN$으로 놓자. 그러면
$X_0=\operatorname{Spec}(B_0)$이고, 준연접 $\mathcal O_{X_0}$-가군
$\mathcal S$와 $\mathcal M$은 각각
\[
S'=\bigoplus_{k\geq0}
((\mathfrak J^k/\mathfrak J^{k+1})\otimes_{A_0}B_0),\qquad
M'=\bigoplus_{k\geq0}\mathfrak J^kN/\mathfrak J^{k+1}N
\]
에 대응한다. 모든 $k\geq0$에 대해 곱셈은 전사 사상
\[
(\mathfrak J^k/\mathfrak J^{k+1})\otimes_{A_0}N_0
\longrightarrow\mathfrak J^kN/\mathfrak J^{k+1}N,
\qquad \overline a\otimes\overline n\longmapsto\overline{an}
\]
을 유도한다. 이 사상들은 등급 $S'$-가군의 전사 준동형사상
$S'\otimes_{B_0}N_0\to M'$을 정의한다. $N_0$는 유한 생성
$B_0$-가군이므로 $M'$는 유한 생성 $S'$-가군이며, 여기서 주장이 따른다
(I, 1.3.13).

사상 $f_0:X_0\to Y_0$는 가정에 의해 \emph{고유}이므로,
$\mathcal S$, $\mathcal M$, $f_0$에
(\hyperref[III.3.3.1-fr]{3.3.1})을 적용할 수 있다.
(\hyperref[III.1.4.11-fr]{1.4.11})도 고려하면 \emph{모든 $n\geq0$}에
대해 $\bigoplus_{k\geq0}H^n(X_0,\mathcal M_k)$가 유한 생성 등급
$S$-가군임을 얻는다. 이는 아벨 군의 층들로 이루어진 엄밀 사영계
$(\mathcal F/\mathfrak J^k\mathcal F)_{k\geq0}$에 각자의 자연스러운
‘여과된 $A$-가군’ 구조를 줄 때, (0, 13.7.7)의 조건
$(\mathrm F_n)$이 \emph{모든 $n\geq0$}에 대해 성립함을 증명한다.
따라서 (0, 13.7.7)을 적용하면 다음을 얻는다.
\begin{enumerate}
\item[$1^\circ$] 사영계
$(H^n(\mathfrak X,\mathcal F_k))_{k\geq0}$는 조건 (ML)을 만족한다.
\item[$2^\circ$]
$H'^{,n}=\varprojlim_kH^n(\mathfrak X,\mathcal F_k)$로 놓으면,
$H'^{,n}$은 유한 생성 $A$-가군이다.
\item[$3^\circ$] 표준 준동형사상
$H'^{,n}\to H^n(\mathfrak X,\mathcal F_k)$들의 핵들이 $H'^{,n}$ 위에
정의하는 여과는 $\mathfrak J$-좋은 여과이다.
\end{enumerate}

한편 $X_k=(\mathfrak X,\mathcal O_{\mathfrak X}/\mathcal K^{k+1})$로
놓으면, $\mathcal F_k$는 연접 $\mathcal O_{X_k}$-가군이다
(I, 10.11.3). $U$가 $X_0$의 아핀 열린집합이면 $U$는 각 $X_k$에서도
아핀 열린집합이므로 (I, 5.1.9), 모든 $n>0$과 모든 $k$에 대해
$H^n(U,\mathcal F_k)=0$이고
(\hyperref[III.1.3.1-fr]{1.3.1}), $h\leq k$일 때
$H^0(U,\mathcal F_k)\to H^0(U,\mathcal F_h)$는 전사이다 (I, 1.3.9).
그러므로 (0, 13.3.2)의 조건들이 성립하며, (0, 13.3.1)의 사상에 의해
$H'^{,n}$은
\[
H^n(\mathfrak X,\varprojlim_k\mathcal F_k)
=H^n(\mathfrak X,\mathcal F)
\]
와 표준적으로 동일시된다. 이것으로
(\hyperref[III.3.4.4-fr]{3.4.4})의 증명이 끝난다.

\begin{env}[3.4.5]
\label{III.3.4.5-fr}
이제 (\hyperref[III.3.4.2-fr]{3.4.2})와
(\hyperref[III.3.4.3-fr]{3.4.3})의 증명으로 돌아가자. 먼저
(\hyperref[III.3.4.4-fr]{3.4.4})에서 다룬
$\mathfrak Y=\operatorname{Spf}(A)$인 경우에 이 명제들을 증명한다.
이를 위해 모든 $g\in A$에 대하여 $\mathfrak Y$의 열린집합
$\mathfrak Y_g=\mathfrak D(g)$ 위에 유도되는 뇌터 아핀 형식적 스킴과
$f^{-1}(\mathfrak Y_g)$ 위에 유도되는 형식적 준스킴에
(\hyperref[III.3.4.4-fr]{3.4.4})를 적용하자. 여기서
$\mathfrak Y_g=\operatorname{Spf}(A_{\{g\}})$이다. $\mathfrak Y_g$는
준스킴
\[
Y_k=(\mathfrak Y,\mathcal O_{\mathfrak Y}/
(\mathfrak J^\Delta)^{k+1})
\]
에서도 아핀 열린집합이고, $\mathcal F_k$는 연접
$\mathcal O_{X_k}$-가군이므로
(\hyperref[III.1.4.11-fr]{1.4.11})에 의해 모든 $k\geq0$에 대해
\[
H^n(f^{-1}(\mathfrak Y_g),\mathcal F_k)
=\Gamma(\mathfrak Y_g,R^nf_*(\mathcal F_k))
\]
이다. 표준 준동형사상
\[
H^n(f^{-1}(\mathfrak Y_g),\mathcal F)
\longrightarrow\varprojlim_k\Gamma(\mathfrak Y_g,R^nf_*(\mathcal F_k))
\]
은 동형사상이다. 그런데 (0\textsubscript{I}, 3.2.6)에 의해
\[
\varprojlim_k\Gamma(\mathfrak Y_g,R^nf_*(\mathcal F_k))
=\Gamma(\mathfrak Y_g,\varprojlim_kR^nf_*(\mathcal F_k))
\]
이다.

\oldpage[III]{122}
또한 층 $R^nf_*(\mathcal F)$는 $\mathfrak Y_g$들 위의 준층
\[
\mathfrak Y_g\longmapsto
H^n(f^{-1}(\mathfrak Y_g),\mathcal F)
\]
에 대응하므로 (0\textsubscript{I}, 3.2.1), 준동형사상
(\hyperref[III.3.4.2.2-fr]{3.4.2.2})가 전단사임을 보였다. 이제
$R^nf_*(\mathcal F)$가 연접 $\mathcal O_{\mathfrak Y}$-가군임을,
더 정확히는
\[
R^nf_*(\mathcal F)=(H^n(\mathfrak X,\mathcal F))^\Delta
\tag{3.4.5.1}\label{III.3.4.5.1-fr}
\]
임을 증명하자. 앞의 표기를 쓰면 $\mathcal F_k$가 연접
$\mathcal O_{X_k}$-가군이므로
(\hyperref[III.1.4.13-fr]{1.4.13})
\[
\Gamma(\mathfrak Y_g,R^nf_*(\mathcal F_k))
=(\Gamma(\mathfrak Y,R^nf_*(\mathcal F_k)))_g
=(H^n(\mathfrak X,\mathcal F_k))_g.
\]
한편 $H^n(\mathfrak X,\mathcal F_k)$들은 (ML)을 만족하는 사영계를
이루고, 그 사영극한 $H^n(\mathfrak X,\mathcal F)$는 유한 생성
$A$-가군이다. 따라서 (0, 13.7.8)과 $A$, $A_{\{g\}}$에 적용한
(I, 10.10.8)에 의해
\[
\varprojlim_k(H^n(\mathfrak X,\mathcal F_k))_g
=H^n(\mathfrak X,\mathcal F)\otimes_AA_{\{g\}}
=\Gamma(\mathfrak Y_g,(H^n(\mathfrak X,\mathcal F))^\Delta)
\]
이다. 또한
\[
\Gamma(\mathfrak Y_g,R^nf_*(\mathcal F))
=\varprojlim_k\Gamma(\mathfrak Y_g,R^nf_*(\mathcal F_k))
\]
이므로, 이것이 (\hyperref[III.3.4.5.1-fr]{3.4.5.1})을 증명한다.
이제 (\hyperref[III.3.4.2.2-fr]{3.4.2.2})는 연접
$\mathcal O_{\mathfrak Y}$-가군의 동형사상이므로 반드시 \emph{위상}
동형사상이다 (I, 10.11.6). 끝으로
$R^nf_*(\mathcal F_k)=(H^n(\mathfrak X,\mathcal F_k))^\Delta$라는
관계에서 사영계 $(R^nf_*(\mathcal F_k))_{k\geq0}$가 (ML)을 만족함이
따른다 (I, 10.10.2).

(\hyperref[III.3.4.2-fr]{3.4.2})와
(\hyperref[III.3.4.3-fr]{3.4.3})를 $\mathfrak Y$가 뇌터 아핀
형식적 준스킴인 경우에 증명하고 나면,
(\hyperref[III.3.4.2-fr]{3.4.2})와
(\hyperref[III.3.4.3-fr]{3.4.3})의 첫째 명제는 $\mathfrak Y$ 위에서
국소적이므로 일반적인 경우로 넘어가는 것은 즉시 이루어진다.
(\hyperref[III.3.4.3-fr]{3.4.3})의 둘째 명제에 대해서는
$\mathfrak Y$가 뇌터이므로 이를 유한 개의 뇌터 아핀 열린집합
$U_i$들로 덮고, 사영계 $(R^qf_*(\mathcal F_k))$를 각 $U_i$에 제한한
것들이 (ML)을 만족함을 유의하면 충분하다.
\end{env}

더 나아가 위의 증명 과정에서 다음도 증명하였다.

\begin{corollary}[3.4.6]
\label{III.3.4.6-fr}
(\hyperref[III.3.4.4-fr]{3.4.4})의 가정 아래에서 표준 준동형사상
\[
H^q(\mathfrak X,\mathcal F)
\longrightarrow\Gamma(\mathfrak Y,R^qf_*(\mathcal F))
\tag{3.4.6.1}\label{III.3.4.6.1-fr}
\]
은 전단사이다.
\end{corollary}

\section{고유 사상의 기본 정리와 그 응용}
\label{section:III.4-fr}

\subsection{기본 정리}
\label{subsection:III.4.1-fr}

\begin{env}[4.1.1]
\label{III.4.1.1-fr}
$X$, $Y$를 통상적인 뇌터 준스킴들, $f:X\to Y$를 고유 사상,
$Y'$를 $Y$의 닫힌 부분집합, $X'=f^{-1}(Y')$를 그 역상이라 하자.
$X$와 $Y$를 각각 $X'$과 $Y'$을 따라 완비화하여 얻는 형식적 준스킴
$X_{/X'}$, $Y_{/Y'}$ (I, 10.8.5)를 각각 $\widehat X$,
$\widehat Y$로 나타내고, $f$를 이 완비화들로 연장한 형식적 준스킴의
사상 $\widehat f:\widehat X\to\widehat Y$ (I, 10.9.1)를
$\widehat f$로 나타내겠다. 모든 연접 $\mathcal O_X$-가군
$\mathcal F$에 대하여,

\oldpage[III]{123}

$X'$을 따른 그 완비화 $\mathcal F_{/X'}$ (I, 10.8.4)를
$\widehat{\mathcal F}$로 나타내겠다. 이는 연접
$\mathcal O_{\widehat X}$-가군이다 (I, 10.8.8).
\end{env}

\begin{env}[4.1.2]
\label{III.4.1.2-fr}
$\operatorname{Supp}(\mathcal O_Y/\mathcal J)=Y'$인
$\mathcal O_Y$의 연접 아이디얼 $\mathcal J$를 택하자 (I, 5.2.1).
(I, 4.4.5)에 의해
\[
\mathcal K=f^*(\mathcal J)\mathcal O_X
\]
는
\[
\operatorname{Supp}(\mathcal O_X/\mathcal K)=X'
\]
인 $\mathcal O_X$의 연접 아이디얼이다. 모든 $k\geq0$에 대하여 연접
$\mathcal O_X$-가군
\[
\mathcal F_k
=\mathcal F\otimes_{\mathcal O_Y}(\mathcal O_Y/\mathcal J^{k+1})
=\mathcal F/\mathcal K^{k+1}\mathcal F
\]
를 생각하자. 모든 $n$에 대해 $\mathcal O_Y$-가군
$R^nf_*(\mathcal F)$와 $R^nf_*(\mathcal F_k)$는 연접이다
(\hyperref[III.3.2.1-fr]{3.2.1}). 모든 $k\geq0$과 모든 $n$에 대해
표준 준동형사상 $\mathcal F\to\mathcal F_k$는 함자성에 의해
준동형사상
\[
R^nf_*(\mathcal F)\longrightarrow R^nf_*(\mathcal F_k)
\tag{4.1.2.1}\label{III.4.1.2.1-fr}
\]
을 정의한다. 또한 $\mathcal F_k$는
$\mathcal O_X/\mathcal K^{k+1}$-가군이므로
$R^nf_*(\mathcal F_k)$는
$\mathcal O_Y/\mathcal J^{k+1}$-가군이다 (0, 12.2.1). 따라서
(\hyperref[III.4.1.2.1-fr]{4.1.2.1})에서 준동형사상
\[
R^nf_*(\mathcal F)\otimes_{\mathcal O_Y}
(\mathcal O_Y/\mathcal J^{k+1})
\longrightarrow R^nf_*(\mathcal F_k)
\tag{4.1.2.2}\label{III.4.1.2.2-fr}
\]
을 얻는다. (\hyperref[III.4.1.2.2-fr]{4.1.2.2})의 양변은 각각
사영계를 이루며, 왼쪽의 사영극한은 다름 아닌
$(R^nf_*(\mathcal F))_{/Y'}$의 완비화이다. 이를
$\widehat{R^nf_*(\mathcal F)}$로 나타내겠다. 또한
(\hyperref[III.4.1.2.2-fr]{4.1.2.2})의 준동형사상들이 사영계를
이룸은 즉시 알 수 있으므로, 극한을 취하여 표준 준동형사상
\[
\varphi_n:\widehat{R^nf_*(\mathcal F)}
\longrightarrow\varprojlim_kR^nf_*(\mathcal F_k)
\tag{4.1.2.3}\label{III.4.1.2.3-fr}
\]
을 얻는다. 한편 (\hyperref[III.4.1.2.2-fr]{4.1.2.2})는
$(\mathcal O_Y/\mathcal J^{k+1})$-가군의 준동형사상이므로
(I, 10.8.3), 의사이산 위상
$\mathcal O_{\widehat Y}$-가군의 연속 준동형사상으로 볼 수 있다
(0\textsubscript{I}, 3.8.1). 따라서 $\varphi_n$은 위상
$\mathcal O_{\widehat Y}$-가군의 연속 준동형사상이다.
\end{env}

\begin{env}[4.1.3]
\label{III.4.1.3-fr}
$i:\widehat X\to X$를 (I, 10.8.7)에서 정의한 환 달린 공간의 표준
사상이라 하자. 그러면 가환 도표
\[
\xymatrix{
X_k\ar[r]^{h_k}\ar[dr]_{i_k}&\widehat X\ar[d]^i\\
&X
}
\tag{4.1.3.1}\label{III.4.1.3.1-fr}
\]
를 얻는다. 여기서 $X_k$는 아이디얼 $\mathcal K^{k+1}$로 정의되는
$X$의 닫힌 부분준스킴이고, $i_k$는 표준 단사이며, $h_k$는 바탕공간의
항등사상과 표준 준동형사상
$\mathcal O_{\widehat X}\to\mathcal O_X/\mathcal K^{k+1}$에 대응하는
환 달린 공간의 사상이다 (I, 10.5.2). 또한 표준 동형사상을 제외하면
$\widehat{\mathcal F}=i^*(\mathcal F)$이다 (I, 10.8.8). 다음 동형사상이
알려져 있다.
\[
H^n(X_k,i_k^*(\mathcal F_k))=H^n(X,\mathcal F_k)
\tag{4.1.3.2}\label{III.4.1.3.2-fr}
\]

\oldpage[III]{124}

실제로 $\mathcal F_k=(i_k)_*(i_k^*(\mathcal F_k))$이기 때문이다
(G, II, 4.9.1). 그러므로 표준 준동형사상
\[
H^n(\widehat X,\widehat{\mathcal F})
\longrightarrow H^n(X_k,h_k^*(\widehat{\mathcal F}))
\qquad\text{(0, 12.1.3.5)}
\]
은
\[
H^n(\widehat X,\widehat{\mathcal F})\longrightarrow H^n(X,\mathcal F_k)
\tag{4.1.3.3}\label{III.4.1.3.3-fr}
\]
로도 쓸 수 있다. 이 준동형사상들은 명백히 사영계를 이루므로, 극한을
취하여 표준 준동형사상
\[
\psi_{n,X}:H^n(\widehat X,\widehat{\mathcal F})
\longrightarrow\varprojlim_kH^n(X,\mathcal F_k)
\tag{4.1.3.4}\label{III.4.1.3.4-fr}
\]
을 얻는다. $X$를 $Y$의 아핀 열린집합 $V$에 대한
$f^{-1}(V)$ 꼴의 열린집합으로 바꾸고
(\hyperref[III.1.4.11-fr]{1.4.11})을 고려하면 준동형사상
\[
\psi_{n,V}:H^n(\widehat X\cap f^{-1}(V),\widehat{\mathcal F})
\longrightarrow\varprojlim_k\Gamma(V,R^nf_*(\mathcal F_k))
\tag{4.1.3.5}\label{III.4.1.3.5-fr}
\]
을 얻는다. 이 준동형사상들은 $V$를 더 작은 아핀 열린집합으로 제한하는
것과 명백히 가환하므로, 결국 층의 표준 준동형사상
\[
\psi_n:R^n\widehat f_*(\widehat{\mathcal F})
\longrightarrow\varprojlim_kR^nf_*(\mathcal F_k)
\tag{4.1.3.6}\label{III.4.1.3.6-fr}
\]
을 정의한다.
\end{env}

\begin{env}[4.1.4]
\label{III.4.1.4-fr}
끝으로 $j:\widehat Y\to Y$를 환 달린 공간의 표준 사상이라 하자
(I, 10.8.7). $R^nf_*(\mathcal F)$는 연접 $\mathcal O_Y$-가군이므로
(\hyperref[III.3.2.1-fr]{3.2.1}), 표준 동형사상을 제외하면
\[
j^*(R^nf_*(\mathcal F))=\widehat{R^nf_*(\mathcal F)}
\]
이다 (I, 10.8.8). 따라서 환 달린 공간에 대해 일반적으로 정의되는
표준 준동형사상
\[
\rho_n:\widehat{R^nf_*(\mathcal F)}=j^*(R^nf_*(\mathcal F))
\longrightarrow R^n\widehat f_*(i^*(\mathcal F))
=R^n\widehat f_*(\widehat{\mathcal F})
\tag{4.1.4.1}\label{III.4.1.4.1-fr}
\]
을 얻는다
((\hyperref[III.1.4.15-fr]{1.4.15})의 증명을 보라). 도표
\[
\xymatrix{
\widehat{R^nf_*(\mathcal F)}\ar[rr]^{\rho_n}\ar[dr]_{\varphi_n}
&&R^n\widehat f_*(\widehat{\mathcal F})\ar[dl]^{\psi_n}\\
&\displaystyle\varprojlim_kR^nf_*(\mathcal F_k)&
}
\tag{4.1.4.2}\label{III.4.1.4.2-fr}
\]
가 가환임을 보이자. 대응하는 준층 준동형사상의 도표가 가환임을 보이면
충분하므로 $Y$가 아핀인 경우로 한정할 수 있다. 그러면 문제는 도표
\[
\xymatrix{
\widehat{H^n(X,\mathcal F)}\ar[rr]^{\rho_n}\ar[dr]_{\varphi_n}
&&H^n(\widehat X,\widehat{\mathcal F})\ar[dl]^{\psi_{n,X}}\\
&\displaystyle\varprojlim_kH^n(X,\mathcal F_k)&
}
\tag{4.1.4.3}\label{III.4.1.4.3-fr}
\]
가 가환임을 보이는 것으로 귀착된다. 그런데
(\hyperref[III.4.1.3.1-fr]{4.1.3.1})의 가환성과
(\hyperref[III.4.1.3-fr]{4.1.3})에서 본 코호몰로지 군 사이의 관계들은
즉시 가환 도표

\oldpage[III]{125}

\[
\xymatrix{
H^n(X,\mathcal F)\ar[r]\ar[dr]&
H^n(\widehat X,\widehat{\mathcal F})
=H^n(\widehat X,i^*(\mathcal F))\ar[d]\\
&H^n(X,\mathcal F_k)=H^n(X_k,i_k^*(\mathcal F))
}
\]
를 주며, 여기서 (\hyperref[III.4.1.4.3-fr]{4.1.4.3})의 가환성이 즉시
따른다.
\end{env}

\begin{theorem}[4.1.5]
\label{III.4.1.5-fr}
$f:X\to Y$를 뇌터 준스킴들의 고유 사상, $Y'$를 $Y$의 닫힌 부분집합,
$X'=f^{-1}(Y')$라 하자. 그러면 모든 연접 $\mathcal O_X$-가군
$\mathcal F$에 대해 $R^n\widehat f_*(\widehat{\mathcal F})$는 연접
$\mathcal O_{\widehat Y}$-가군이고, 도표
(\hyperref[III.4.1.4.2-fr]{4.1.4.2})의 준동형사상 $\varphi_n$,
$\psi_n$, $\rho_n$은 위상 동형사상이다.
\end{theorem}

\begin{proof}
$\varphi_n$과 $\psi_n$이 동형사상임을 보이면 충분하다. 실제로
$R^nf_*(\mathcal F)$는 연접이므로
(\hyperref[III.3.2.1-fr]{3.2.1}),
$\widehat{R^nf_*(\mathcal F)}$도 연접이고 (I, 10.8.8),
$\varphi_n$, $\psi_n$, $\rho_n$의 양방향 연속성은 자동으로 따른다
(I, 10.11.6).
\end{proof}

\begin{env}[4.1.6]
\label{III.4.1.6-fr}
\emph{주의.}
\begin{enumerate}
\item[(i)]
$\widehat{\mathcal F}_k=
\widehat{\mathcal F}/\widehat{\mathcal K}^{k+1}\widehat{\mathcal F}$로
놓으면 $\mathcal F_k=i_*(\widehat{\mathcal F}_k)$임은 즉시 알 수 있다.
표준 준동형사상 (\hyperref[III.4.1.3.6-fr]{4.1.3.6})은
(\hyperref[III.3.4.2.2-fr]{3.4.2.2})에서 이미 정의한 준동형사상
\[
R^n\widehat f_*(\widehat{\mathcal F})
\longrightarrow\varprojlim_kR^n\widehat f_*(\widehat{\mathcal F}_k)
\tag{4.1.6.1}\label{III.4.1.6.1-fr}
\]
과 다르지 않다. 따라서 $\psi_n$이 동형사상이라는 사실은
(\hyperref[III.3.4.3-fr]{3.4.3})의 특수한 경우이다. 그러나 아래에서는
일반 정리 (\hyperref[III.3.4.3-fr]{3.4.3})가 의존하는 스펙트럼 열들의
사영극한에 관한 섬세한 논의 (0, 13.7)를 피하는 직접 증명을 제시하겠다.

\item[(ii)]
$\psi_n$들이 동형사상이라는 사실을 고려하면, $\varphi_n$들이
동형사상이라는 것과
$\rho_n=\psi_n^{-1}\circ\varphi_n$들이 동형사상이라는 것은 동치이다.
정리 (\hyperref[III.4.1.5-fr]{4.1.5})는 특히 $R^nf_*$의 형성이
완비화와 가환함을 표현하며, 이를 ‘대수적’ 이론과 ‘형식적’ 이론의
제1 비교 정리라고 부를 수 있다.
\end{enumerate}
이제 먼저 (\hyperref[III.4.1.5-fr]{4.1.5})의 아핀 형태를 증명하겠다.
\end{env}

\begin{corollary}[4.1.7]
\label{III.4.1.7-fr}
(\hyperref[III.4.1.5-fr]{4.1.5})의 가정들에 더하여
$Y=\operatorname{Spec}(A)$이고 $A$가 뇌터이며,
$\mathcal J=\widetilde{\mathfrak J}$이고 $\mathfrak J$가 $A$의
아이디얼이라고 가정하자. 따라서
$\mathcal F_k=\mathcal F/\mathfrak J^{k+1}\mathcal F$이다. 표준
준동형사상
\[
\varphi_n:\widehat{H^n(X,\mathcal F)}
\longrightarrow\varprojlim_kH^n(X,\mathcal F_k)
\tag{4.1.7.1}\label{III.4.1.7.1-fr}
\]
은 동형사상이다. 여기서 왼쪽은 $\mathfrak J$-준아딕 위상에 관한
$H^n(X,\mathcal F)$의 분리 완비화이다. 사영계
$(H^n(X,\mathcal F_k))_{k\geq0}$는 모든 $n$에 대해 조건 (ML)을
만족하고, 표준 준동형사상
\[
\psi_n:H^n(\widehat X,\widehat{\mathcal F})
\longrightarrow\varprojlim_kH^n(X,\mathcal F_k)
\tag{4.1.7.2}\label{III.4.1.7.2-fr}
\]

\oldpage[III]{126}

은 동형사상이다. 끝으로 표준 준동형사상
$H^n(X,\mathcal F)\to H^n(X,\mathcal F_k)$들의 핵들이
$H^n(X,\mathcal F)$ 위에 정의하는 여과는 $\mathfrak J$-좋은 여과이고
(0, 13.7.7), $\varphi_n$은 위상 동형사상이다.
\footnote{뒤따르는 증명은 최초의 증명보다 간단하며,
$H^n(X,\mathcal F)$의 여과에 관한 보충과 함께 J.-P. Serre가
알려 주었다.}
\end{corollary}

\begin{proof}
이 증명에서 정수 $n\geq0$을 고정하고, 표기를 간단히 하기 위해
\[
\mathrm H=H^n(X,\mathcal F),\qquad
\mathrm H_k=H^n(X,\mathcal F_k)
\tag{4.1.7.3}\label{III.4.1.7.3-fr}
\]
및
\[
\mathrm R_k=\operatorname{Ker}(\mathrm H\longrightarrow\mathrm H_k),
\qquad\text{$\mathrm H$의 부분 $A$-가군}
\tag{4.1.7.4}\label{III.4.1.7.4-fr}
\]
로 놓겠다. 코호몰로지 완전열
\[
H^n(X,\mathfrak J^{k+1}\mathcal F)\longrightarrow H^n(X,\mathcal F)
\longrightarrow H^n(X,\mathcal F_k)\xrightarrow{\partial}
H^{n+1}(X,\mathfrak J^{k+1}\mathcal F)
\longrightarrow H^{n+1}(X,\mathcal F)
\]
에서 또한
\[
\mathrm R_k=\operatorname{Im}\bigl(
H^n(X,\mathfrak J^{k+1}\mathcal F)\longrightarrow H^n(X,\mathcal F)
\bigr)
\]
임을 알 수 있다. 다음과 같이 놓자.
\[
\begin{aligned}
\mathrm Q_k
&=\operatorname{Ker}\bigl(H^{n+1}(X,\mathfrak J^{k+1}\mathcal F)
\longrightarrow H^{n+1}(X,\mathcal F)\bigr)\\
&=\operatorname{Im}\bigl(H^n(X,\mathcal F_k)
\longrightarrow H^{n+1}(X,\mathfrak J^{k+1}\mathcal F)\bigr).
\end{aligned}
\tag{4.1.7.5}\label{III.4.1.7.5-fr}
\]
따라서 완전열
\[
0\longrightarrow\mathrm R_k\longrightarrow\mathrm H
\longrightarrow\mathrm H_k\longrightarrow\mathrm Q_k\longrightarrow0
\tag{4.1.7.6}\label{III.4.1.7.6-fr}
\]
을 얻는다.

\begin{env}[4.1.7.7]
\label{III.4.1.7.7-fr}
$x$를 $\mathfrak J^m$의 원소라 하자 ($m\geq0$).
$\mathfrak J^k\mathcal F$에서 $x$를 곱하는 것은 준동형사상
$\mathfrak J^k\mathcal F\to\mathfrak J^{k+m}\mathcal F$이며, 따라서
준동형사상
\[
\mu_{x,m}:H^n(X,\mathfrak J^k\mathcal F)
\longrightarrow H^n(X,\mathfrak J^{k+m}\mathcal F)
\tag{4.1.7.8}\label{III.4.1.7.8-fr}
\]
을 준다. $S=\bigoplus_{k\geq0}\mathfrak J^k$라는 등급 $A$-대수를
생각하면, 곱셈 $\mu_{x,m}$들은
\[
\mathrm E=\bigoplus_{k\geq0}H^n(X,\mathfrak J^k\mathcal F)
\]
에 등급환 $S$ 위의 유한 생성 등급 가군 구조를 정의한다
(\hyperref[III.3.3.2-fr]{3.3.2}). 이 등급환은 뇌터이다
(II, 2.1.5).
\end{env}

\begin{lemma}[4.1.7.9]
\label{III.4.1.7.9-fr}
$\mathrm H$의 부분가군 $\mathrm R_k$들은 $\mathrm H$ 위에
$\mathfrak J$-좋은 여과를 정의한다.
\end{lemma}

\begin{proof}
먼저
\[
\mathfrak J^m\mathrm R_k\subset\mathrm R_{k+m}
\tag{4.1.7.10}\label{III.4.1.7.10-fr}
\]
임을 보이자. 여기서 $x\in\mathfrak J^m$을
$\mathrm H=H^n(X,\mathcal F)$에 곱하는 사상은 $\mu_{x,0}$이다. 모든
$x\in\mathfrak J^m$에 대해 도표
\[
\xymatrix{
\mathfrak J^{k+1}\mathcal F\ar[r]^x\ar[d]&
\mathfrak J^{k+m+1}\mathcal F\ar[d]\\
\mathcal F\ar[r]_x&\mathcal F
}
\tag{4.1.7.11}\label{III.4.1.7.11-fr}
\]

\oldpage[III]{127}

는 가환이다. 여기서 수평 화살표들은 $x$의 곱셈이고 수직 화살표들은
표준 단사이다. 따라서 대응하는 도표
\[
\xymatrix{
H^n(X,\mathfrak J^{k+1}\mathcal F)\ar[r]^{\mu_{x,m}}\ar[d]&
H^n(X,\mathfrak J^{k+m+1}\mathcal F)\ar[d]\\
H^n(X,\mathcal F)\ar[r]_{\mu_{x,0}}&H^n(X,\mathcal F)
}
\]
도 가환한다. $\mathrm R_k$를
$H^n(X,\mathfrak J^{k+1}\mathcal F)\to H^n(X,\mathcal F)$의 상으로
해석하면 이것이 (\hyperref[III.4.1.7.10-fr]{4.1.7.10})을 증명한다.
또한 등급 $S$-가군
$\mathrm R=\bigoplus_{k\geq0}\mathrm R_k$가 $\mathrm E$의 부분
$S$-가군
\[
\mathrm M=\bigoplus_{k\geq0}
H^n(X,\mathfrak J^{k+1}\mathcal F)
\]
의 몫임을 보여 준다. 그러므로 위의 주의에 의해 $\mathrm R$은 유한
생성 $S$-가군이다. 이는
(\hyperref[III.4.1.7.9-fr]{4.1.7.9})의 주장과 동치이다
(Bourbaki, \emph{Alg. comm.}, chap. III, §~3,
n\textsuperscript{o}~1, th.~1).
\end{proof}

\begin{env}[4.1.7.12]
\label{III.4.1.7.12-fr}
이제 (4.7.1.8)과 같이 정의되는 등급 $S$-가군
\[
\mathrm N=\bigoplus_{k\geq0}
H^{n+1}(X,\mathfrak J^{k+1}\mathcal F)
\]
을 생각하자. 이것도
(\hyperref[III.3.3.2-fr]{3.3.2})에 의해 유한 생성 $S$-가군이다.
(\hyperref[III.4.1.7.5-fr]{4.1.7.5})에 의해 모든 $k$에 대해
$\mathrm Q_k\subset\mathrm N_k$이고,
(\hyperref[III.4.1.7.11-fr]{4.1.7.11})에서 $n$을 $n+1$로 바꾼 도표는
\[
S_m\mathrm Q_k=\mathfrak J^m\mathrm Q_k\subset\mathrm Q_{k+m}
\]
임을 보여 준다. 다시 말해 $\mathrm Q$는 $\mathrm N$의 등급
부분 $S$-가군이고, 따라서 유한 생성이다.
\end{env}

\begin{env}[4.1.7.13]
\label{III.4.1.7.13-fr}
$S_m\to S_0$로도 쓸 수 있는 표준 단사
$\alpha_m:\mathfrak J^m\to A$를 생각하자.
$\mathfrak J^{k+1}\mathcal F_k=0$이므로 $A$-가군
$H^n(X,\mathcal F_k)$는 $\mathfrak J^{k+1}$에 의해 소멸한다.
$\mathrm Q_k$는 $A$-준동형사상
$H^n(X,\mathcal F_k)\to
H^{n+1}(X,\mathfrak J^{k+1}\mathcal F)$의 상이므로,
$A$-가군 $\mathrm Q_k$도 $\mathfrak J^{k+1}$에 의해 소멸한다.
이는 $S$-가군 $\mathrm Q$에서
\[
\alpha_{k+1}(S_{k+1})\mathrm Q_k=0
\tag{4.1.7.14}\label{III.4.1.7.14-fr}
\]
임을 뜻한다. $\mathrm Q$는 유한 생성 $S$-가군이므로 어떤 정수
$k_0$와 정수 $h$가 존재하여 $k\geq k_0$일 때
$\mathrm Q_{k+h}=S_h\mathrm Q_k$이다 (II, 2.1.6, (ii)).
이 관계와 (\hyperref[III.4.1.7.14-fr]{4.1.7.14})에서 어떤 정수
$r>0$이 존재하여
\[
\alpha_r(S_r)\mathrm Q=0
\tag{4.1.7.15}\label{III.4.1.7.15-fr}
\]
임이 따른다.
\end{env}

\begin{env}[4.1.7.16]
\label{III.4.1.7.16-fr}
표준 단사 $\mathfrak J^{k+m}\mathcal F\to\mathfrak J^k\mathcal F$는
코호몰로지에서 $A$-준동형사상
\[
\nu_m:H^{n+1}(X,\mathfrak J^{k+m}\mathcal F)
\longrightarrow H^{n+1}(X,\mathfrak J^k\mathcal F)
\tag{4.1.7.17}\label{III.4.1.7.17-fr}
\]
을 준다. 모든 $x\in\mathfrak J^m$에 대해 명백한 분해
\[
\xymatrix@C=4.5em{
H^{n+1}(X,\mathfrak J^k\mathcal F)\ar[r]^{\mu_{x,m}}&
H^{n+1}(X,\mathfrak J^{k+m}\mathcal F)\ar[r]^{\nu_m}&
H^{n+1}(X,\mathfrak J^k\mathcal F)
}
\tag{4.1.7.18}\label{III.4.1.7.18-fr}
\]
가 있으며, 그 합성은 $\mu_{x,0}$이다.
\end{env}

\oldpage[III]{128}

따라서 $H^{n+1}(X,\mathfrak J^k\mathcal F)$의 모든 부분
$A$-가군 $P$에 대하여 $S$-가군 $\mathrm N$ 안에서
\[
\nu_m(S_mP)=\alpha_m(S_m)P
\tag{4.1.7.19}\label{III.4.1.7.19-fr}
\]
이다.

\begin{lemma}[4.1.7.20]
\label{III.4.1.7.20-fr}
어떤 정수 $m>0$이 존재하여 모든 $k\geq k_0$에 대해
$\nu_m(\mathrm Q_{k+m})=0$이다.
\end{lemma}

\begin{proof}
$h$의 배수이면서 $m\geq r$인 $m$을 택하자.
$k\geq k_0$일 때 $\mathrm Q_{k+m}=S_m\mathrm Q_k$이므로
(\hyperref[III.4.1.7.19-fr]{4.1.7.19})와
(\hyperref[III.4.1.7.15-fr]{4.1.7.15})에 의해
\[
\nu_m(\mathrm Q_{k+m})=\alpha_m(S_m)\mathrm Q_k
\subset\alpha_r(S_r)\mathrm Q_k=0.
\]
\end{proof}

\begin{env}[4.1.7.21]
\label{III.4.1.7.21-fr}
가환 도표
\[
\xymatrix@C=2.5em{
H^n(X,\mathcal F)\ar[r]&H^n(X,\mathcal F_k)\ar[r]^-\partial&
H^{n+1}(X,\mathfrak J^{k+1}\mathcal F)\ar[r]&H^{n+1}(X,\mathcal F)\\
H^n(X,\mathcal F)\ar[r]\ar[u]&H^n(X,\mathcal F_{k+m})\ar[r]^-\partial\ar[u]&
H^{n+1}(X,\mathfrak J^{k+m+1}\mathcal F)\ar[r]\ar[u]&
H^{n+1}(X,\mathcal F)\ar[u]
}
\]
는 수직 화살표들이 표준 사상인 가환 도표
\[
\xymatrix{
0\ar[r]&\mathfrak J^{k+1}\mathcal F\ar[r]&\mathcal F\ar[r]&
\mathcal F_k\ar[r]&0\\
0\ar[r]&\mathfrak J^{k+m+1}\mathcal F\ar[r]\ar[u]&
\mathcal F\ar[r]\ar[u]&\mathcal F_{k+m}\ar[r]\ar[u]&0
}
\]
에서 유도된다. 따라서 행들이 완전인 가환 도표
\[
\xymatrix@C=2.2em{
0\ar[r]&\mathrm R_k\ar[r]&\mathrm H\ar[r]&\mathrm H_k\ar[r]&
\mathrm Q_k\ar[r]&0\\
0\ar[r]&\mathrm R_{k+m}\ar[r]\ar[u]&\mathrm H\ar[r]\ar@{=}[u]&
\mathrm H_{k+m}\ar[r]\ar[u]&\mathrm Q_{k+m}\ar[r]\ar[u]^{\nu_m}&0
}
\]
를 얻는다. $k\geq k_0$일 때 마지막 수직 화살표는 영이므로
(\hyperref[III.4.1.7.20-fr]{4.1.7.20}), $\mathrm H_{k+m}$의
$\mathrm H_k$ 안의 상은
\[
\operatorname{Ker}(\mathrm H_k\to\mathrm Q_k)
=\operatorname{Im}(\mathrm H\to\mathrm H_k)
\]
에 포함된다. 한편 도표의 가환성에 의해 이 상은
$\operatorname{Im}(\mathrm H\to\mathrm H_k)$를 포함하므로 둘은 같다.
따라서 모든 $k'\geq k+m$에 대해 $\mathrm H_{k'}$의 $\mathrm H_k$
안의 상도 이와 같으며, 이것이 사영계
$(\mathrm H_k)_{k\geq0}$에 대한 조건 (ML)을 증명한다. 또한 $X$의 모든
아핀 열린집합 $U$에 대해 $i>0$이면
$H^i(U,\mathcal F_k)=0$이고
(\hyperref[III.1.3.1-fr]{1.3.1}), $m>0$이면
$H^0(U,\mathcal F_{k+m})\to H^0(U,\mathcal F_k)$가 전사이다
(I, 1.3.9). 따라서 (0, 13.3.1)을 적용할 수 있고,

\oldpage[III]{129}

표준 준동형사상
\[
H^n(\widehat X,\widehat{\mathcal F})
\longrightarrow\varprojlim_kH^n(X,\mathcal F_k)
\]
은 모든 $n\geq0$에 대해 전단사이다.

사영계 $(\mathrm H/\mathrm R_k)_{k\geq0}$는 엄밀하므로 완전열
\[
0\longrightarrow\mathrm H/\mathrm R_k\longrightarrow\mathrm H_k
\longrightarrow\mathrm Q_k\longrightarrow0
\tag{4.1.7.22}\label{III.4.1.7.22-fr}
\]
에서 사영극한을 취할 수 있다 (0, 13.2.2).
$\nu_m(\mathrm Q_{k+m})=0$이므로
$\varprojlim_k\mathrm Q_k=0$이며, 따라서 위상 동형사상
\[
\varprojlim_k(\mathrm H/\mathrm R_k)\xrightarrow{\sim}
\varprojlim_k\mathrm H_k
\]
을 얻는다. 그런데 $\mathrm H$의 여과 $(\mathrm R_k)$는
$\mathfrak J$-좋은 여과이므로 $\mathrm H$ 위에
$\mathfrak J$-준아딕 위상을 정의한다. 따라서
$\varprojlim_k(\mathrm H/\mathrm R_k)$는 이 위상에 관한
$\mathrm H$의 분리 완비화이다. 이것으로
(\hyperref[III.4.1.7-fr]{4.1.7})의 증명이 끝난다.
\end{env}
\end{proof}

\begin{env}[4.1.8]
\label{III.4.1.8-fr}
끝으로 (\hyperref[III.4.1.5-fr]{4.1.5})를 증명하자. $Y$의 모든 아핀
열린집합 $V$에 대하여
$\Gamma(V,\widehat{R^nf_*(\mathcal F)})$는
$\Gamma(V,R^nf_*(\mathcal F))$의 $\mathfrak J$-준아딕 위상에 관한
분리 완비화이다. 여기서
$\mathcal J|V=\widetilde{\mathfrak J}$로 쓴다. 실제로
$R^nf_*(\mathcal F)$는 연접 $\mathcal O_Y$-가군이다 (I, 10.8.4).
또한
\[
\Gamma\bigl(V,\varprojlim_kR^nf_*(\mathcal F_k)\bigr)
=\varprojlim_k\Gamma(V,R^nf_*(\mathcal F_k))
\]
이다 (0\textsubscript{I}, 3.2.6). 따라서 $\varphi_n$이 위상
동형사상이라는 사실은 (\hyperref[III.4.1.7-fr]{4.1.7})과
(\hyperref[III.1.4.11-fr]{1.4.11})에서 따른다. 한편 역시
(\hyperref[III.1.4.11-fr]{1.4.11})에 의해
(\hyperref[III.4.1.7-fr]{4.1.7})에서
(\hyperref[III.4.1.3.3-fr]{4.1.3.3})의 준동형사상 $\psi_{n,V}$가
동형사상임이 따른다. 그러므로
$R^n\widehat f_*(\widehat{\mathcal F})$의 정의에 의해 $\psi_n$도
동형사상이다.
\end{env}

\begin{corollary}[4.1.9]
\label{III.4.1.9-fr}
(\hyperref[III.4.1.4-fr]{4.1.4})의 가정 아래에서 $Y$의 모든 아핀
열린집합 $V$에 대해 표준 준동형사상
\[
H^n(\widehat X\cap f^{-1}(V),\widehat{\mathcal F})
\longrightarrow\Gamma(\widehat Y\cap V,
R^n\widehat f_*(\widehat{\mathcal F}))
\]
은 전단사이다.
\end{corollary}

\begin{env}[4.1.10]
\label{III.4.1.10-fr}
\emph{주의.}
$f:X\to Y$를 통상적인 뇌터 준스킴들의 유한형 사상이라 하고,
$\mathcal F$를 그 받침이 $Y$ 위에서 고유인 연접
$\mathcal O_X$-가군이라 하자 (II, 5.4.10). 그러면 모든 $n\geq0$에
대해 $R^nf_*(\mathcal F)$가 연접 $\mathcal O_Y$-가군임은 알려져 있다
(\hyperref[III.3.2.4-fr]{3.2.4}). 또한 항상
$\mathcal F=u_*(\mathcal G)$라고 가정할 수 있다. 여기서 $Z$는
바탕공간이 $\operatorname{Supp}(\mathcal F)$인 $X$의 적당한 닫힌
부분준스킴이고, $u:Z\to X$는 표준 단사이며,
$\mathcal G=u^*(\mathcal F)$는 연접 $\mathcal O_Z$-가군이다
(I, 9.3.5). 다음과 같이 놓자.
\[
\mathcal G_k=\mathcal G\otimes_{\mathcal O_Y}
(\mathcal O_Y/\mathcal J^{k+1}).
\]
그러면 $\mathcal G_k=u^*(\mathcal F_k)$이고,
\[
R^nf_*(\mathcal F_k)=R^n(f\circ u)_*(\mathcal G_k),\qquad
R^nf_*(\mathcal F)=R^n(f\circ u)_*(\mathcal G)
\]
이다 (\hyperref[III.1.3.4-fr]{1.3.4}). 끝으로 (I, 10.9.5)를 고려하면
\[
R^n\widehat f_*(\widehat{\mathcal F})
=R^n\widehat{(f\circ u)}_*(\widehat{\mathcal G}).
\]
따라서 $\mathcal G$와 고유 사상 $f\circ u$에
(\hyperref[III.4.1.5-fr]{4.1.5})를 적용할 수 있고, 이 가정 아래에서도
(\hyperref[III.4.1.5-fr]{4.1.5})의 결과들이 $\mathcal F$와 $f$에 대해
성립한다.
\end{env}

\subsection{특수한 경우들과 변형들}
\label{subsection:III.4.2-fr}

비교 정리 (\hyperref[III.4.1.5-fr]{4.1.5})의 가장 유용한 형태는
다음과 같다.

\begin{proposition}[4.2.1]
\label{III.4.2.1-fr}
$Y$를 국소 뇌터 준스킴, $f:X\to Y$를 고유 사상, $\mathcal F$를 연접
$\mathcal O_X$-가군이라 하자. 그러면 모든 $y\in Y$와 모든 $p$에 대해
$(R^pf_*(\mathcal F))_y$는

\oldpage[III]{130}

유한 생성 $\mathcal O_y$-가군이고, 따라서 $\mathfrak m_y$-준아딕
위상에 관하여 분리이다. 또한 표준 위상 동형사상
\[
\widehat{(R^pf_*(\mathcal F))_y}
\xrightarrow{\sim}
\varprojlim_kH^p\bigl(f^{-1}(y),
\mathcal F\otimes_{\mathcal O_Y}(\mathcal O_y/\mathfrak m_y^k)\bigr)
\tag{4.2.1.1}\label{III.4.2.1.1-fr}
\]
이 있다. 여기서 왼쪽은 $(R^pf_*(\mathcal F))_y$의
$\mathfrak m_y$-준아딕 위상에 관한 완비화이고, 오른쪽에서는 모든
$k\geq0$에 대해 $f^{-1}(y)$를 준스킴
\[
X\times_Y\operatorname{Spec}(\mathcal O_y/\mathfrak m_y^k)
\]
의 바탕공간으로 본다 (I, 3.6.1).
\end{proposition}

\begin{proof}
$\mathcal O_y$는 뇌터 국소환이고
$(R^pf_*(\mathcal F))_y$는 유한 생성 $\mathcal O_y$-가군이므로
(\hyperref[III.3.2.1-fr]{3.2.1}), 그 위의
$\mathfrak m_y$-준아딕 위상은 분리이다 (0\textsubscript{I}, 7.3.5).
$Y$가 뇌터이고 $y$가 닫힌점인 경우에는 나머지 주장들이
(\hyperref[III.4.1.7-fr]{4.1.7})에서 따른다. 실제로 $Y$를 $y$의
아핀 근방으로 바꾸고 $Y'=\{y\}$로 택하면 된다
(G, II, 4.9.1). 일반적인 경우에는
\[
Y_1=\operatorname{Spec}(\mathcal O_y),\qquad
X_1=X\times_YY_1,\qquad
\mathcal F_1=\mathcal F\otimes_{\mathcal O_Y}\mathcal O_{Y_1}
\]
로 놓고 $f_1=f\times1_{Y_1}:X_1\to Y_1$라 하자. $Y_1$은 뇌터이고,
$f_1$은 고유이며 (II, 5.4.2, (iii)), $\mathcal F_1$은 연접이다
(0\textsubscript{I}, 5.3.11). $Y_1$의 유일한 닫힌점을 $y_1$이라 하자.
명제는 $f_1$, $\mathcal F_1$, $y_1$에 대해 성립한다. 또한
\[
\mathcal O_{y_1}=\mathcal O_y,\qquad
f_1^{-1}(y_1)=f^{-1}(y)
\]
이고 (I, 3.6.5), 준스킴
\[
X\times_Y\operatorname{Spec}(\mathcal O_y/\mathfrak m_y^k)
\quad\text{및}\quad
X_1\times_{Y_1}\operatorname{Spec}(\mathcal O_y/\mathfrak m_y^k)
\]
은 표준적으로 동일시된다 (I, 3.3.9). 더 나아가
\[
\mathcal F_1\otimes_{\mathcal O_{Y_1}}
(\mathcal O_y/\mathfrak m_y^k)
\]
는
$\mathcal F\otimes_{\mathcal O_Y}(\mathcal O_y/\mathfrak m_y^k)$와
동일시된다 (I, 9.1.6). 이제
\[
R^pf_{1*}(\mathcal F_1)\simeq
R^pf_*(\mathcal F)\otimes_{\mathcal O_Y}\mathcal O_{Y_1}
\]
임만 보이면 된다. 이는 국소 사상
$\operatorname{Spec}(\mathcal O_y)\to Y$가 평탄하므로
(0\textsubscript{I}, 6.7.1 및 I, 2.4.2),
(\hyperref[III.1.4.15-fr]{1.4.15})에서 따른다.
\end{proof}

다음 따름정리는 차원 이론의 용어 (제IV장)를 사용하며 제IV장 전에는
적용하지 않는다.

\begin{corollary}[4.2.2]
\label{III.4.2.2-fr}
$Y$를 국소 뇌터 준스킴, $f:X\to Y$를 고유 사상, $y$를 $Y$의 점,
$r$을 $f^{-1}(y)$의 차원이라 하자. 그러면 모든 연접
$\mathcal O_X$-가군 $\mathcal F$에 대해, 모든 $p>r$에서 층
$R^pf_*(\mathcal F)$는 $y$의 어떤 근방에서 영이다.
\end{corollary}

\begin{proof}
실제로 모든 $k$에 대해
\[
H^p\bigl(f^{-1}(y),
\mathcal F\otimes_{\mathcal O_Y}(\mathcal O_y/\mathfrak m_y^k)\bigr)=0
\]
이다 (G, II, 4.15.2). 따라서
(\hyperref[III.4.2.1-fr]{4.2.1})에 의해
$(R^pf_*(\mathcal F))_y$의 $\mathfrak m_y$-준아딕 위상에 관한 분리
완비화는 영이다. 이 위상은 분리이므로
$(R^pf_*(\mathcal F))_y=0$이다. $R^pf_*(\mathcal F)$가 연접이므로
(0\textsubscript{I}, 5.2.2), 결론이 따른다.
\end{proof}

\begin{env}[4.2.3]
\label{III.4.2.3-fr}
결과 (\hyperref[III.4.2.1-fr]{4.2.1})은 주로 $p=0$일 때 사용된다.
따라서 다음 따름정리를 얻는다.
\end{env}

\begin{corollary}[4.2.4]
\label{III.4.2.4-fr}
(\hyperref[III.4.2.1-fr]{4.2.1})의 가정 아래에서 표준 위상 동형사상
\[
\widehat{(f_*(\mathcal F))_y}
\xrightarrow{\sim}
\varprojlim_k\Gamma\bigl(f^{-1}(y),
\mathcal F_y/\mathfrak m_y^k\mathcal F_y\bigr)
\]
이 있다.
\end{corollary}

\subsection{자리스키의 연결성 정리}
\label{subsection:III.4.3-fr}
\label{III.4.3-fr}

이 번호와 다음 번호의 결과들은 자리스키의 잘 알려진 정리들을
일반화하며, 모두 (\hyperref[III.4.2.4-fr]{4.2.4})에서 이끌어낼 수 있다.
이들은 다음 정리의 귀결이다.

\begin{theorem}[4.3.1]
\label{III.4.3.1-fr}
\textnormal{(연결성 정리).}
$Y$를 국소 뇌터 준스킴, $f:X\to Y$를 고유 사상이라 하자. 그러면
\[
\mathcal A(X)=f_*(\mathcal O_X)
\]
는 연접 $\mathcal O_Y$-대수이다.

\oldpage[III]{131}

$\mathcal A(Y')=\mathcal A(X)$이고 $Y$ 위에서 유한인 $Y$-스킴
$Y'$를 택하자. 이는 $Y$-동형사상을 제외하면 유일하게 결정된다
(II, 1.3.1 및 6.1.3). 항등 동형사상
$e:\mathcal A(Y')\to\mathcal A(X)$에서 얻는 $Y$-사상
$f'=\mathcal A(e):X\to Y'$ (II, 1.2.7)는 고유이고,
$f'_*(\mathcal O_X)$는 $\mathcal O_{Y'}$와 동형이며, 사상 $f'$의
모든 올 $f'^{-1}(y')$는 연결이고 공집합이 아니다.
\end{theorem}

\begin{proof}
$g:Y'\to Y$를 구조 사상이라 하자. $f'$의 정의에 들어가는 준동형사상
\[
\theta:\mathcal O_{Y'}\longrightarrow f'_*(\mathcal O_X)
\]
이 전단사임을 보이자. $Y'$는 $Y$ 위에서 아핀이므로,
\[
g_*(\theta):g_*(\mathcal O_{Y'})\longrightarrow
g_*(f'_*(\mathcal O_X))=f_*(\mathcal O_X)
\]
이 항등사상임을 보이면 충분하다 (II, 1.4.2). 그런데 정의에 의해
$g_*(\mathcal O_{Y'})=\mathcal A(Y')$이고
$f_*(\mathcal O_X)=\mathcal A(X)$이므로 그러하다.
$\mathcal A(X)$가 연접이라는 사실은 유한성 정리
(\hyperref[III.3.2.1-fr]{3.2.1})의 특수한 경우이다. $f$는 고유이고
$g$는 분리이므로 $f'$은 고유이다 (II, 5.4.3, (i)). 따라서
(\hyperref[III.4.3.1-fr]{4.3.1})의 증명을 끝내려면 다음 따름정리를
증명하면 충분하다.
\end{proof}

\begin{corollary}[4.3.2]
\label{III.4.3.2-fr}
(\hyperref[III.4.3.1-fr]{4.3.1})의 가정에 더하여
$f_*(\mathcal O_X)\simeq\mathcal O_Y$라고 가정하자. 그러면 모든
$y\in Y$에 대해 $f$의 올 $f^{-1}(y)$는 연결이고 공집합이 아니다.
\end{corollary}

\begin{proof}
$f_*(\mathcal O_X)\simeq\mathcal O_Y$라는 가정에서 이미 $f$가 우세임이
따른다. $f$는 닫힌 사상이므로 전사이다.
(\hyperref[III.4.2.1-fr]{4.2.1})에서와 같이 $y$가 $Y$의 닫힌점인
경우로 귀착시킬 수 있다. $f^{-1}(y)$는 뇌터 공간이므로 유한 개의 연결
성분을 가지며, $f^{-1}(y)$를 따른 완비화 $\widehat X$의 바탕공간이다.
이 연결 성분들을 $Z_i$ ($1\leq i\leq n$)라 하자. 그러면
$\Gamma(\widehat X,\mathcal O_{\widehat X})$는 환들
$\Gamma(Z_i,\mathcal O_{\widehat X})$의 직접곱이고, 각 환은 영환이
아니다. 실제로 $\widehat X$의 모든 점에서 단위원 단면은 영과 다르다.
한편 $\mathcal F=\mathcal O_X$에
(\hyperref[III.4.1.5-fr]{4.1.5})를 적용하면, $f^{-1}(y)$를 따른 그
완비화가 $\mathcal O_{\widehat X}$이므로
$\Gamma(\widehat X,\mathcal O_{\widehat X})$는 국소환
$\mathcal O_y$의 $\mathfrak m_y$-아딕 분리 완비화
$\widehat{\mathcal O}_y$와 동형이다. 따라서 이는 국소환이며, 여러
영이 아닌 환들의 직접곱일 수 없다. 그렇지 않으면 서로 다른 극대
아이디얼들을 여럿 가질 것이기 때문이다. 그러므로 $n=1$이고,
따름정리가 증명되었다.
\end{proof}

\begin{corollary}[4.3.3]
\label{III.4.3.3-fr}
(\hyperref[III.4.3.1-fr]{4.3.1})의 가정 아래에서 모든 $y\in Y$에 대해
올 $f^{-1}(y)$의 연결 성분들의 집합은 구조 사상 $g:Y'\to Y$의 올
$g^{-1}(y)$의 유한한 점 집합과 전단사로 대응한다. 다시 말해 이는
$(f_*(\mathcal O_X))_y$의 극대 아이디얼들의 집합이다.
\end{corollary}

\begin{proof}
$Y'$는 $Y$ 위에서 유한이므로 $g^{-1}(y)$는 유한 이산 공간이다
(II, 6.1.7). 또한 $f^{-1}(y)=f'^{-1}(g^{-1}(y))$이므로, 이 관찰과
(\hyperref[III.4.3.1-fr]{4.3.1})에서 따름정리가 따른다.
\end{proof}

이로써 (\hyperref[III.4.3.1-fr]{4.3.1})에서 정의한 $Y$-준스킴
$Y'$에 대한 주목할 만한 해석을 얻었다. 고유 사상 $f$의 인수분해
$f=g\circ f'$은 해석공간의 정칙 사상에 대해 K.~Stein이 얻은
인수분해와 유사하다. 앞으로 이를 $f$의 \emph{스테인 인자분해}라
부르겠다.

\begin{remark}[4.3.4]
\label{III.4.3.4-fr}
$k$를 체 $k(y)$의 확대라 하자. 준스킴
\[
f^{-1}(y)\otimes_{k(y)}k=X\times_Y\operatorname{Spec}(k)
\]
가 연결이면, 그 사영 사상에 의한 상인 $f^{-1}(y)$도 연결이다
(I, 3.4.7). 준스킴의 사상 $f:X\to Y$와 점 $y\in Y$에 대해, 모든
$k(y)$의 확대 $k$에 대하여 준스킴
\[
f^{-1}(y)\otimes_{k(y)}k=X\times_Y\operatorname{Spec}(k)
\]
가 연결일 때 올 $f^{-1}(y)$가 \emph{기하적으로 연결}이라고 한다.

\oldpage[III]{132}

(\hyperref[III.4.3.2-fr]{4.3.2})의 가정 아래에서는 결론을 강화할 수
있다. 실제로 올 $f^{-1}(y)$들은 \emph{기하적으로 연결}이다. 이를
보이기 위해, 모든 $k(y)$의 확대 $k$에 대하여 어떤 뇌터 국소환 $A$와
국소 준동형사상 $\varphi:\mathcal O_y\to A$가 존재하여 $A$가 평탄
$\mathcal O_y$-가군이고 $A$의 잉여체가 $k$와 $k(y)$-동형임을
유의하자 (0, 10.3.1). $Y_1=\operatorname{Spec}(A)$라 하고,
$h:Y_1\to Y$를 $\varphi$에 대응하며 $Y_1$의 유일한 닫힌점 $y_1$을
$y$로 보내는 국소 사상이라 하자 (I, 2.4.1). 또한
\[
X_1=X\times_YY_1,\qquad f_1=f\times1_{Y_1}
\]
로 놓자. $f_1$은 고유이고 (II, 5.4.2, (iii)),
$f_1^{-1}(y_1)$은 $X\times_Y\operatorname{Spec}(k)$와 동형인
$k(y_1)$-준스킴이다. 따라서
(\hyperref[III.4.3.2-fr]{4.3.2})를 $f_1$에 적용하기 위해
$f_{1*}(\mathcal O_{X_1})=\mathcal O_{Y_1}$임만 보이면 된다.
$h$는 평탄 사상이다 (I, 2.4.2 및
\hyperref[III.1.4.15-fr]{1.4.15.5}). 그러므로 $q=0$에 대해
(\hyperref[III.1.4.15-fr]{1.4.15})를 적용하면
\[
f_{1*}(\mathcal O_{X_1})
=h^*(f_*(\mathcal O_X))
=h^*(\mathcal O_Y)
=\mathcal O_{Y_1}
\]
이다.

일반적인 경우 (\hyperref[III.4.3.1-fr]{4.3.1})에도 같은 논증은 그
정리의 표기로
\[
f_{1*}(\mathcal O_{X_1})=h^*(g_*(\mathcal O_{Y'}))
\]
임을 보여 준다. $f_1$의 스테인 인자분해
$f_1=g_1\circ f'_1$에서
\[
g_1=g\times1_{Y_1}
\]
이고 (II, 1.5.2), 대응하는 유한 $Y_1$-스킴은
$Y'_1=Y'\times_YY_1$이다. 올의 추이성을 고려하면 (I, 3.6.4),
(\hyperref[III.4.3.3-fr]{4.3.3})에 의해 $f_1^{-1}(y_1)$의 연결
성분 수는
\[
g_1^{-1}(y_1)=g^{-1}(y)\otimes_{k(y)}k
\]
의 원소 수와 같다. $k$를 $k(y)$의 대수적으로 닫힌 확대체로 택하면,
이 수는 선택한 대수적으로 닫힌 확대체에 의존하지 않으며
$g^{-1}(y)$의 \emph{기하적 점의 수} (I, 6.4.7), 즉
\emph{분리차수들의 합}
\[
\sum_i[k(y'_i):k(y)]_s
\]
과 같다. 여기서 $y'_i$는 유한집합 $g^{-1}(y)$를 달린다. 이 수를
$f^{-1}(y)$의 \emph{기하적 연결 성분의 수}라고도 한다. 체
$k(y'_i)$들은 반국소환 $(f_*(\mathcal O_X))_y$의 잉여체들과
다르지 않다.
\end{remark}

\begin{proposition}[4.3.5]
\label{III.4.3.5-fr}
$X$, $Y$를 국소 뇌터 정역 준스킴들, $f:X\to Y$를 고유 우세
사상이라 하자. 모든 $y\in Y$에 대해 $f^{-1}(y)$의 연결 성분 수는
유리함수체 $\mathrm R(X)$ 안에서 $\mathcal O_y$의 정수적 폐포
$\mathcal O'_y$의 극대 아이디얼 수 이하이다.
\end{proposition}

\begin{proof}
실제로 $Y$의 모든 열린집합 $U$에 대해
\[
\Gamma(U,f_*(\mathcal O_X))
=\Gamma(f^{-1}(U),\mathcal O_X)
\]
는 $x\in f^{-1}(U)$인 국소환 $\mathcal O_x$들의 교집합이다
(I, 8.2.1.1). 따라서 $(f_*(\mathcal O_X))_y$는
$\mathcal O_y$를 포함하는 $\mathrm R(X)$의 부분환이다. 또한
$f_*(\mathcal O_X)$는 연접 $\mathcal O_Y$-가군이므로
$(f_*(\mathcal O_X))_y$는 유한 생성 $\mathcal O_y$-가군이고,
따라서 $\mathcal O'_y$에 포함된다. 이러한 환 $A$의 모든 극대
아이디얼이 $A$와 $\mathcal O'_y$의 어떤 극대 아이디얼의 교집합임은
알려져 있다 ([13], vol. I, p. 257 및 259). 여기서 명제가 따른다.
\end{proof}

\begin{definition}[4.3.6]
\label{III.4.3.6-fr}
정역 국소환의 정수적 폐포가 국소환이면 그 국소환을
\emph{단일분기(unibranche)}라고 한다. 정역 준스킴 $Y$의 점 $y$에서
국소환 $\mathcal O_y$가 단일분기이면 $y$가 단일분기라고 한다.
특히 $Y$가 $y$에서 정규이면 그러하다.
\end{definition}

$A$를 정역 국소환, $K$를 그 분수체라 하자. $A$가 단일분기이기 위한
필요충분조건은 $A$를 포함하고 유한 $A$-대수인 $K$의 모든 부분환
$A_1$이 국소환인 것이다. 실제로 $A'$을 $A$의 정수적 폐포라 하자.
첫째 Cohen--Seidenberg 정리에 의해 (Bourbaki,
\emph{Alg. comm.}, chap. V, \S~2, n\textsuperscript{o}~1, th.~1),
$A_1$의 모든 극대 아이디얼은 $A'$의 어떤 극대 아이디얼의 자취이다.
따라서 $A'$이 국소이면 $A_1$도 국소이다. 역으로 $A'$은 그 안의 유한
부분 $A$-대수 $A_\alpha$들로 이루어진 증가 여과족의 귀납극한이다.

\oldpage[III]{133}

각 $A_\alpha$가 국소환이면, 위와 같은 논증에 의해
$A_\alpha\subset A_\beta$일 때 $A_\alpha$의 극대 아이디얼은
$A_\beta$의 극대 아이디얼이 $A_\alpha$ 위에 남기는 자취이다. 따라서
$A'$은 국소환이다 (0, 10.3.1.3).

뇌터 국소환 $A$의 완비화가 정역이면, 즉 $A$가
\emph{해석적으로 정역}이면, $A$는 단일분기임도 유의하자. 실제로
$\mathfrak m$을 $A$의 극대 아이디얼, $K$를 그 분수체,
$K'$을 $\widehat A$의 분수체라 하면
\[
K'=K\otimes_A\widehat A.
\]
$A'_F$를 $K$ 안의 유한 부분 $A$-대수라 하자. $\widehat A$와
$A'_F$로 생성되는 $K'$의 부분환 $B_F$는
$A'_F\otimes_A\widehat A$와 동형이다. 이는 유한 생성
$\widehat A$-가군이며, $A'_F$의 $\mathfrak m$-아딕 위상에 관한
완비화이다 (0\textsubscript{I}, 7.3.3 및 7.3.6).
$A'_F$는 반국소환이고 (Bourbaki, \emph{Alg. comm.}, chap. IV,
\S~2, n\textsuperscript{o}~5, prop. 9의 cor. 3), 그 완비화가
정역이므로 $A'_F$는 오직 하나의 극대 아이디얼 $\mathfrak m'_F$만
가질 수 있으며 $\mathfrak m'_F\cap A=\mathfrak m$이다. 여기서
주장이 따른다.

\begin{corollary}[4.3.7]
\label{III.4.3.7-fr}
(\hyperref[III.4.3.5-fr]{4.3.5})의 가정 아래에서
$\mathrm R(X)$ 안의 $\mathrm R(Y)$의 대수적 폐포의 분리차수가
$n$이고 $y\in Y$가 단일분기라고 가정하자. 그러면 올 $f^{-1}(y)$는
연결 성분을 많아야 $n$개 가진다. 특히 $\mathrm R(X)$ 안의
$\mathrm R(Y)$의 대수적 폐포가 $\mathrm R(Y)$ 위에서 순수 비분리이면
$f^{-1}(y)$는 연결이다.
\end{corollary}

\begin{proof}
$\mathcal O'_y$를 $\mathcal O_y$의 정수적 폐포라 하자.
$\mathrm R(X)$ 안의 $\mathcal O_y$의 정수적 폐포
$\mathcal O''_y$는 $\mathcal O'_y$의 정수적 폐포이기도 하다.
$\mathcal O'_y$가 국소환이면 $\mathcal O''_y$는 반국소환이고 그 극대
아이디얼 수는 $n$ 이하임이 알려져 있다
([13], vol. I, p. 289, th. 22).

이 따름정리는 자리스키가 대수적 스킴에 대한 자신의
‘연결성 정리’를 서술한 형태와 본질적으로 같다.
\end{proof}

\begin{remark}[4.3.8]
\label{III.4.3.8-fr}
(\hyperref[III.4.3.7-fr]{4.3.7})의 가정에 더하여 $Y$가 $y$에서
정규라고 가정하면, 올 $f^{-1}(y)$는 \emph{기하적으로 연결}이다.
실제로 (\hyperref[III.4.3.4-fr]{4.3.4})의 표기로 $g^{-1}(y)$는 한
점 $y'$로 이루어지고 $k(y')$는 $k(y)$ 위에서 순수 비분리이다.
\end{remark}

\begin{definition}[4.3.9]
\label{III.4.3.9-fr}
$Y$를 국소 뇌터 준스킴이라 하자. 모든 기약 국소 뇌터 준스킴
$Y'$와 모든 우세 사상 $g:Y'\to Y$에 대해
$X'=X\times_YY'$의 모든 기약 성분이 $Y'$ 위에서 우세이면, 유한형
사상 $f:X\to Y$를 \emph{보편적으로 열린 사상}이라고 한다.
\end{definition}

$Y$가 기약인 경우에는 다음 조건과 동치이다. $Y$, $Y'$의 일반점을
각각 $\eta$, $\eta'$라 하고 $g(\eta')=\eta$라 하며
$f'=f_{(Y')}$로 놓으면, $X'$의 모든 기약 성분이
$f'^{-1}(\eta')$와 만난다 (0\textsubscript{I}, 2.1.8). 따라서
$Y$의 모든 열린집합 $U$에 대해 $f$의 제한
$f^{-1}(U)\to U$도 보편적으로 열린 사상이다.

\begin{corollary}[4.3.10]
\label{III.4.3.10-fr}
$X$, $Y$를 국소 뇌터 정역 준스킴들이라 하고, $f:X\to Y$를 고유하고
우세이며 보편적으로 열린 사상이라 하자. $\mathrm R(X)$ 안의
$\mathrm R(Y)$의 대수적 폐포가 $\mathrm R(Y)$ 위에서 순수
비분리이면, 모든 올 $f^{-1}(y)$ ($y\in Y$)는 기하적으로 연결이다.
\end{corollary}

\begin{proof}
$Y=\operatorname{Spec}(B)$이고 $B$가 뇌터 정역인 경우로 한정할 수
있다. (II, 7.1.7)에 의해 $\mathcal O_y$를 우세 지배하고 분수체가
$\mathrm R(Y)$인 정수적으로 닫힌 뇌터 국소환 $A$가 존재한다.
$Y'=\operatorname{Spec}(A)$라 하고, 표준 단사 $B\to A$에 대응하는
쌍유리 사상 $h:Y'\to Y$를 생각하자. 이는 특히 우세이다. $Y'$의
유일한 닫힌점을 $y'$라 하면 $h(y')=y$이다. 다음과 같이 놓자.
\[
X'=X\times_YY',\qquad f'=f\times1_{Y'}.
\]
$Y$, $Y'$, $X$의 일반점을 각각 $\eta$, $\eta'$, $\xi$라 하자.
그러면 $f(\xi)=\eta$이고 $h^{-1}(\eta)=\{\eta'\}$이다. 또한
$k(\eta)=k(\eta')=\mathrm R(Y)$이므로
$f'^{-1}(\eta')\simeq f^{-1}(\eta)$이다 (I, 3.6.4).
$\xi$는 $f^{-1}(\eta)$의 일반점이므로
(0\textsubscript{I}, 2.1.8), $f'^{-1}(\eta')$는 일반점을 하나만
가진다. 그런데 가정에 의해 $X'$의 모든 기약 성분의 일반점은
$f'^{-1}(\eta')$ 안에 있다. 따라서 $X'$은 기약이고 그 일반점
$\xi'$은 $f'^{-1}(\eta')$의 일반점이며
$k(\xi')=k(\xi)$이다.

\oldpage[III]{134}

$X''=X'_{\mathrm{red}}$, $f''=f'_{\mathrm{red}}$로 놓자. 그러면
$X''$은 뇌터 정역 준스킴이고, $f''$은 고유이며 (II, 5.4.6),
$f'^{-1}(y')$와 $f''^{-1}(y')$의 바탕공간은 같다. 또한
\[
\mathrm R(X'')=k(\xi')=\mathrm R(X)
\]
이므로 $f''$은 (\hyperref[III.4.3.8-fr]{4.3.8})의 가정을 만족하고
$f''^{-1}(y')$는 기하적으로 연결이다. 이제 $k$를 $k(y)$의 임의의
확대라 하자. $k(y')$와 $k$를 모두 부분확대로 갖는 $k(y)$의 확대
$k_1$이 존재한다 (Bourbaki, \emph{Alg.}, chap. V, \S~4, prop. 2).
가정에 의해
\[
f''^{-1}(y')\times_{Y'}\operatorname{Spec}(k_1)
\]
은 연결이고, 이는
$f'^{-1}(y')\times_{Y'}\operatorname{Spec}(k_1)$과 같은 축약
준스킴을 가진다 (I, 5.1.8). 따라서 후자도 연결이다. 후자는
$f^{-1}(y)\times_Y\operatorname{Spec}(k_1)$과 동형이므로
(I, 3.6.4), 이 준스킴도 연결이다. 그러므로
(\hyperref[III.4.3.4-fr]{4.3.4}) 첫머리의 주의에 의해
\emph{더욱이} $f^{-1}(y)\times_Y\operatorname{Spec}(k)$도 연결이다.
이것으로 증명이 끝난다.
\end{proof}

\begin{env}[4.3.11]
\label{III.4.3.11-fr}
\emph{주의} (4.3.11). ---
\begin{enumerate}
\item[(i)]
앞의 논증은 본질적으로 자리스키 [20]의 것이다. 다만 고전 대수기하의
국소환에서는 $\mathcal O_y$의 정수적 폐포가 뇌터환이므로, 자리스키는
$A$로 그 정수적 폐포를 택할 수 있었다. 한편 자리스키는 $Y$가 체
$k$ 위의 사영공간 $\mathbf P_k^r$의 Chow 다양체이고, $X$가
$\mathbf P_k^r$와 $Y$ 사이의 Chow 대응을 정의하는
$\mathbf P_k^r\times_kY$의 닫힌 부분집합이면 사영
$X\to Y$가 보편적으로 열린 사상임을 증명하였다
(\emph{loc. cit.}, p. 82의 lemma). 어떤 응용들에서 쓰인
‘Chow 좌표’의 형식적 성질은 이것뿐인 것으로 보인다. 따라서 이러한
상황에서는 ‘사영공간 안의 순환의 특수화’라는 언어 대신 ‘고유 사상의
올’이라는 언어를 쓰는 것이 유익하다. 필요하면 그 사상이 보편적으로
열려 있다고 가정하거나, 국소 정칙성에 관한 다른 유사한 제한을
부과한다.

\item[(ii)]
제IV장에서 보편적으로 열린 사상 $f:X\to Y$를 다음과 같이 정의할
수 있음도 보겠다. 이 성질이 그 용어를 정당화한다. 모든 사상
$Z\to Y$에 대해 사상 $f_{(Z)}:X_{(Z)}\to Z$가 열린 사상이다. 또한
$f$가 (\hyperref[III.4.3.10-fr]{4.3.10})의 가정들을 만족하고
$y,y'$가 $Y$의 두 점이며 $y$가 $y'$의 특수화이면, $f^{-1}(y)$의
기하적 연결 성분 수는 $f^{-1}(y')$의 연결 성분 수 이하임을 보일 수
있다.
\end{enumerate}
\end{env}

\begin{corollary}[4.3.12]
\label{III.4.3.12-fr}
(\hyperref[III.4.3.5-fr]{4.3.5})의 가정 아래에서
$\mathrm R(Y)$가 $\mathrm R(X)$ 안에서 대수적으로 닫혀 있고 $y$가
$Y$의 정규점이라고 가정하자. 그러면 $f^{-1}(y)$는 기하적으로
연결이고, $Y$ 안의 $y$의 어떤 열린 근방 $U$가 존재하여
\[
f_*(\mathcal O_X|f^{-1}(U))\simeq\mathcal O_Y|U
\]
이다. 특히 $Y$가 정규이고 $\mathrm R(Y)$가 $\mathrm R(X)$ 안에서
대수적으로 닫혀 있으면
$f_*(\mathcal O_X)\simeq\mathcal O_Y$이다.
\end{corollary}

\begin{proof}
$f^{-1}(y)$에 관한 첫째 주장은
(\hyperref[III.4.3.8-fr]{4.3.8})의 특수한 경우이다. 따라서
$X\xrightarrow{f'}Y'\xrightarrow{g}Y$가 $f$의 스테인 인자분해이면
(\hyperref[III.4.3.3-fr]{4.3.3}), $g^{-1}(y)$는 한 점 $y'$로
이루어진다. 또한

\oldpage[III]{135}

\[
\mathcal O_y\subset\mathcal O_{y'}=(f_*(\mathcal O_X))_y
\subset\mathrm R(X).
\]
$\mathcal O_{y'}$는 $\mathcal O_y$ 위에서 유한이므로
$\mathrm R(Y)$ 위에서 대수적이고, 가정에 의해 $\mathrm R(Y)$에
포함된다. $y$는 정규이므로 반드시
$\mathcal O_{y'}=\mathcal O_y$이다. 따라서 $g$는 $y'$에서 국소
동형사상이고 (I, 6.5.4), 이것으로 따름정리의 첫째 부분이 증명된다.
둘째 부분은 첫째 부분에서 따른다. 실제로 추가 가정 아래에서 $g$는
전단사이고 $Y'$의 모든 점의 근방에서 국소 동형사상이므로
동형사상이다.
\end{proof}

(\hyperref[III.4.3.7-fr]{4.3.7})을 스킴의 틀에서 증명했기 때문에
다음과 같은 응용이 가능하다.

\begin{proposition}[4.3.13]
\label{III.4.3.13-fr}
$A$를 단일분기 뇌터 국소환, $\mathfrak a$를 $A$의 정의 아이디얼,
$A_0=A/\mathfrak a$,
\[
S=\operatorname{gr}_{\mathfrak a}(A)
\]
를 $\mathfrak a$-준아딕 여과에 관해 $A$에 대응하는 등급환이라 하자.
$S$는 $S_1$로 생성되는 등급 $A_0$-대수이고, $S_1$은 유한 생성
$A_0$-가군이다. 그러면 $\operatorname{Proj}(S)$는 연결
$A_0$-스킴이다.
\end{proposition}

\begin{proof}
$\mathfrak m$을 $A$의 극대 아이디얼이라 하자.
$Y=\operatorname{Spec}(A)$는 정역 스킴이고 $\mathfrak m$에 대응하는
점 $y$는 그 유일한 닫힌점이다. 가정에 의해 어떤 정수 $p$에 대해
\[
\mathfrak m^p\subset\mathfrak a\subset\mathfrak m
\]
이므로 $V(\mathfrak a)=\{\mathfrak m\}$이다.
$S'=\bigoplus_{n\geq0}\mathfrak a^n$이라 하고,
$X=\operatorname{Proj}(S')$라 하자. 이는 아이디얼 $\mathfrak a$를
부풀려 얻는 $Y$-스킴이다. $X$는 정역이고 구조 사상 $f:X\to Y$는
쌍유리이며 (II, 8.1.4) 명백히 사영적이다. 따라서
(\hyperref[III.4.3.7-fr]{4.3.7})을 적용할 수 있고
$f^{-1}(y)$는 연결이다. 그런데 $f^{-1}(y)$는
\[
\operatorname{Proj}(S'\otimes_AA_0)
\]
의 바탕공간이며 (I, 3.6.1 및 II, 2.8.10),
정의에 의해 $S'\otimes_AA_0=S$이다. 따라서 명제가 증명되었다.
\end{proof}

\subsection{자리스키의 “주정리”}
\label{subsection:III.4.4-fr}

\begin{proposition}[4.4.1]
\label{III.4.4.1-fr}
$Y$를 국소 뇌터 준스킴, $f:X\to Y$를 고유 사상이라 하자.
$X'$을 자신의 올 $f^{-1}(f(x))$ 안에서 고립된 점 $x\in X$들의
집합이라 하자. 그러면 $X'$은 $X$의 열린집합이다. 또한
$f=g\circ f'$이 $f$의 스테인 인자분해이면
(\hyperref[III.4.3.3-fr]{4.3.3}), $f'$의 $X'$에 대한 제한은
$X'$에서 $Y'$의 어떤 열린집합 $U$ 위에 유도된 부분준스킴으로 가는
동형사상이고 $X'=f'^{-1}(U)$이다.
\end{proposition}

\begin{proof}
$g^{-1}(f(x))$는 유한 이산 공간이므로
(\hyperref[III.4.3.3-fr]{4.3.3} 및 II, 6.1.7), $x$가
$f^{-1}(f(x))$ 안에서 고립되기 위한 필요충분조건은
$f'^{-1}(f'(x))$ 안에서 고립되는 것이다. 따라서 $f'=f$, 즉
$f_*(\mathcal O_X)=\mathcal O_Y$인 경우로 한정할 수 있다.
$x\in X'$이면 연결공간 $f^{-1}(f(x))$
(\hyperref[III.4.3.2-fr]{4.3.2})은 반드시 한 점 $x$로 이루어진다.
$f$는 닫힌 사상이므로, $X$ 안의 $x$의 모든 열린 근방 $V$에 대해
$f(X-V)$는 $Y$의 닫힌집합이다. 이는 $y=f(x)$를 포함하지 않는다.
실제로 $f^{-1}(y)=\{x\}$이다. $U$를 $Y$에서 $f(X-V)$의 여집합이라
하면 $f^{-1}(U)\subset V$이다. 따라서 $y$의 열린 근방 기본계의
$f$에 의한 역상들은 $x$의 열린 근방 기본계를 이룬다.

$f_*(\mathcal O_X)=\mathcal O_Y$라는 가정과 층의 직상 정의
(0\textsubscript{I}, 3.4.1 및 4.2.1)에 의해, $f=(\psi,\theta)$로
쓰면 준동형사상
\[
\theta_y^\#:\mathcal O_y\longrightarrow\mathcal O_x
\]
는 동형사상이다. 따라서 $x$의 열린 근방 $V$와 $y$의 열린 근방
$U$가 존재하여 $f$의 $V$에 대한 제한이 $V$에서 $U$로 가는
동형사상이다 (I, 6.5.4). 더 나아가 앞의 관찰에 의해

\oldpage[III]{136}

$f^{-1}(U)=V$라고 가정할 수 있다. 그러면 정의에서 즉시
$V\subset X'$임이 따르며, 증명이 끝난다.
\end{proof}

다음 명제는 대수적 스킴의 경우에 Chevalley가 증명하였다.

\begin{proposition}[4.4.2]
\label{III.4.4.2-fr}
$Y$를 국소 뇌터 준스킴, $f:X\to Y$를 사상이라 하자. 다음 조건들은
서로 동치이다.
\begin{enumerate}
\item[a)] $f$는 유한 사상이다.
\item[b)] $f$는 아핀이고 고유이다.
\item[c)] $f$는 고유이고, 모든 $y\in Y$에 대해 $f^{-1}(y)$는
유한집합이다.
\end{enumerate}
\end{proposition}

\begin{proof}
a)가 b)를 함의함은 알려져 있다 (II, 6.1.2 및 6.1.11). $f$가
고유이고 아핀이면
\[
f^{-1}(y)\longrightarrow\operatorname{Spec}(k(y))
\]
도 그러하다 (II, 1.6.2, (iii) 및 5.4.2, (iii)). $f^{-1}(y)$의
구조층에 유한성 정리 (\hyperref[III.3.2.1-fr]{3.2.1})를 적용하면
$f^{-1}(y)=\operatorname{Spec}(A)$이고 $A$는 유한 $k(y)$-대수이다.
따라서 $f^{-1}(y)$는 유한집합이다 (II, 6.1.7). 이로써 b)가 c)를
함의한다.

끝으로 $f^{-1}(y)$는 $k(y)$ 위의 대수적 준스킴이므로, 집합
$f^{-1}(y)$가 유한이라는 가정에서 공간 $f^{-1}(y)$가 이산임이
따른다 (I, 6.4.4). (\hyperref[III.4.4.1-fr]{4.4.1})의 표기로
$X'=X$이고 $f':X\to Y'$은 동형사상이다. $g$는 유한 사상이므로
c)가 a)를 함의한다.
\end{proof}

\begin{theorem}[4.4.3]
\label{III.4.4.3-fr}
\textnormal{(자리스키의 “주정리”).}
$Y$를 뇌터 준스킴, $f:X\to Y$를 준사영 사상, $X'$을 자신의 올
$f^{-1}(f(x))$ 안에서 고립된 점 $x\in X$들의 집합이라 하자. 그러면
$X'$은 $X$의 열린부분집합이고, 유도 부분준스킴 $X'$은 $Y$ 위에서
유한인 어떤 $Y$-준스킴 $Y'$의 열린부분집합 위에 유도된 준스킴과
동형이다.
\end{theorem}

\begin{proof}
가정에 의해 어떤 사영적 $Y$-준스킴 $Z$가 존재하여 $X$는 $Z$의
어떤 열린집합 위에 유도된 부분준스킴과 $Y$-동형이다
(II, 5.3.2 및 5.5.1). 따라서 $f$가 사영 사상인 경우로 귀착된다.
이때 $f$는 고유이고 (II, 5.5.3), 정리는
(\hyperref[III.4.4.1-fr]{4.4.1})에서 즉시 따른다.
\end{proof}

\begin{remark}[4.4.4]
\label{III.4.4.4-fr}
$X$가 축약이면 (각각 $X$가 기약이고 $X'$이 공집합이 아니면),
(\hyperref[III.4.4.3-fr]{4.4.3})의 명제에서 $Y'$이 축약이라고
(각각 기약이라고) 가정할 수 있다. 실제로 $Y'$을 언제나 $Y'$ 안의
$X'$의 폐포 부분준스킴 $\overline{X'}$로 바꿀 수 있다
(I, 9.5.11 및 II, 6.1.5, (i), (ii)). $X'$이 축약이면
$\overline{X'}$도 축약임이 알려져 있다 (I, 9.5.9, (i)). 또한
$X'$이 공집합이 아니고 $X$가 기약이면 $X'$도 기약이며,
$\overline{X'}$도 기약이다.
\end{remark}

\begin{corollary}[4.4.5]
\label{III.4.4.5-fr}
$Y$를 국소 뇌터 스킴, $f:X\to Y$를 유한형 사상, $x$를 자신의 올
$f^{-1}(f(x))$ 안에서 고립된 $X$의 점이라 하자. 그러면 $X$ 안의
$x$의 어떤 열린 근방은 $Y$ 위에서 유한인 어떤 $Y$-준스킴의
열린부분집합과 동형이다.
\end{corollary}

\begin{proof}
$y=f(x)$라 하고, $U$를 $Y$ 안의 $y$의 아핀 열린 근방,
$V$를 $f^{-1}(U)$에 포함되는 $X$ 안의 $x$의 아핀 열린 근방이라
하자. $Y$는 분리이므로 단사 $U\to Y$는 아핀 사상이다 (II, 1.6.3).
$V$도 $U$ 위에서 아핀이므로 (\emph{ibid.}), $f$의 $V$에 대한
제한은 아핀 사상 $V\to Y$이다 (II, 1.6.2, (ii)). 이 제한은
유한형이므로 \emph{더욱이} 준사영 사상이다
(I, 6.3.5 및 II, 5.3.4, (i)). 이제 이 제한에
(\hyperref[III.4.4.3-fr]{4.4.3})을 적용하면 된다.
\end{proof}

따름정리 (\hyperref[III.4.4.5-fr]{4.4.5})를 가환대수의 언어로
서술하면 다음과 같다.

\oldpage[III]{137}

\begin{corollary}[4.4.6]
\label{III.4.4.6-fr}
$A$를 뇌터환, $B$를 유한 생성 $A$-대수, $\mathfrak q$를 $B$의
소 아이디얼, $\mathfrak p$를 그 $A$ 안의 역상이라 하자.
$\mathfrak q$가 역상이 $\mathfrak p$인 $B$의 소 아이디얼들의
집합에서 동시에 극대이고 극소라고 가정하자. 그러면 어떤
$g\in B-\mathfrak q$, 유한 $A$-대수 $A'$, 원소 $f'\in A'$이
존재하여 $A$-대수 $B_g$와 $A'_{f'}$이 동형이다.
\end{corollary}

\begin{proof}
$Y=\operatorname{Spec}(A)$와 $X=\operatorname{Spec}(B)$에
(\hyperref[III.4.4.5-fr]{4.4.5})를 적용하면 충분하다.
$\mathfrak q$에 대한 가정은 그것이 자신의 올 안에서 고립됨을 정확히
뜻한다 (I, 1.1.7).
\end{proof}

여기서 겉보기에는 덜 일반적인 다음 결과를 얻는다.

\begin{corollary}[4.4.7]
\label{III.4.4.7-fr}
$A$를 뇌터 국소환, $B$를 유한 생성 $A$-대수, $\mathfrak n$을
$A$ 안의 역상이 극대 아이디얼 $\mathfrak m$인 $B$의 소
아이디얼이라 하자. $\mathfrak n$이 $B$에서 극대이고, 역상이
$\mathfrak m$인 $B$의 소 아이디얼들의 집합에서 극소라고 가정하자.
이는 $B\mathfrak m$이 $\mathfrak n$에 대해 일차라는 뜻이기도 하다.
그러면 유한 $A$-대수 $A'$과 $A'$의 극대 아이디얼
$\mathfrak m'$이 존재하여, $\mathfrak m$은 $\mathfrak m'$의
$A$ 안의 역상이고 $B_{\mathfrak n}$은 $A$-대수
$A'_{\mathfrak m'}$과 동형이다.
\end{corollary}

다음은 (\hyperref[III.4.4.7-fr]{4.4.7})의 특수한 경우이며,
때때로 “주정리”라고도 부른다.

\begin{corollary}[4.4.8]
\label{III.4.4.8-fr}
(\hyperref[III.4.4.7-fr]{4.4.7})의 조건 아래에서 $A$, $B$가
정역이고 같은 분수체 $K$를 가진다고 더 가정하자. $A$가 정수적으로
닫혀 있으면 $B=A$이다.
\end{corollary}

\begin{proof}
주의 (\hyperref[III.4.4.4-fr]{4.4.4})에 의해
(\hyperref[III.4.4.7-fr]{4.4.7})을 적용할 때 $A'$이 정역이고
분수체가 $K$라고 가정할 수 있다. $A$에 대한 가정에서
$A'=A$가 따르므로 $B_{\mathfrak n}=A$이다. 이제
\[
A\subset B\subset B_{\mathfrak n}
\]
이므로 $A=B$이다.
\end{proof}

명제 (\hyperref[III.4.4.8-fr]{4.4.8})은 자리스키가 자신의
“주정리”에 부여한 형태이다. 여기서는 임의의 뇌터 정역 국소환으로
확장하였다.

앞의 따름정리들은 전역적 결과인
(\hyperref[III.4.4.3-fr]{4.4.3})의 국소적 변형들이었다. 다음은
또 다른 전역적 귀결이다.

\begin{corollary}[4.4.9]
\label{III.4.4.9-fr}
$X$, $Y$를 국소 뇌터 정역 준스킴들, $f:X\to Y$를 분리이고 국소
유한형인 쌍유리 사상이라 하자. 더 나아가 $Y$가 정규이고 모든
$y\in Y$에 대해 올 $f^{-1}(y)$가 유한이라고 가정하자. 그러면
$f$는 열린 몰입 사상이다. 더 나아가 $f$가 닫힌 사상이면
(특히 $f$가 고유이면) $f$는 동형사상이다.
\end{corollary}

\begin{proof}
$x\in X$라 하고 $y=f(x)$로 놓자. $f^{-1}(y)$는 $k(y)$ 위에서
국소 유한형인 스킴이므로, 유한이라는 가정에서 이산임이 따른다
(I, 6.4.4). 또한 $\mathcal O_y$는 정수적으로 닫혀 있고
$\mathcal O_x$, $\mathcal O_y$는 같은 분수체를 가진다 (I, 7.1.5).
따라서 (\hyperref[III.4.4.8-fr]{4.4.8})을 적용할 수 있다.
$f=(\psi,\theta)$로 쓰면
\[
\theta_y^\#:\mathcal O_y\longrightarrow\mathcal O_x
\]
는 전단사이다. 그러므로 $f$는 국소 동형사상이다 (I, 6.5.4).
그런데 $f$는 분리이고 $X$는 정역이므로 $f$는 열린 몰입 사상이다
(I, 8.2.8). 마지막 주장은 $f$가 우세라는 사실에서 따른다.
\end{proof}

\begin{proposition}[4.4.10]
\label{III.4.4.10-fr}
$Y$를 국소 뇌터 준스킴, $f:X\to Y$를 국소 유한형 사상이라 하자.
자신의 올 $f^{-1}(f(x))$ 안에서 고립된 점 $x\in X$들의 집합 $X'$은
$X$에서 열린집합이다.
\end{proposition}

\begin{proof}
문제는 $X$와 $Y$ 위에서 국소적이므로 $X$, $Y$가 뇌터 아핀이고
$f$가 유한형이라고 가정할 수 있다. 이때 $f$는 유한형 아핀
사상이므로 준사영이다 (II, 5.3.4, (i)). 이제
(\hyperref[III.4.4.3-fr]{4.4.3})을 적용하면 된다.
\end{proof}

\oldpage[III]{138}

\begin{corollary}[4.4.11]
\label{III.4.4.11-fr}
$Y$를 국소 뇌터 준스킴, $f:X\to Y$를 고유 사상이라 하자.
$f^{-1}(y)$가 이산인 점 $y\in Y$들의 집합 $U$는 $Y$에서 열려 있고,
$f$의 제한 $f^{-1}(U)\to U$는 유한 사상이다. 특히 고유 준유한
사상 $X\to Y$는 유한 사상이다.
\end{corollary}

\begin{proof}
$Y$에서 $U$의 여집합은 $X-X'$의 $f$에 의한 상이다.
(\hyperref[III.4.4.10-fr]{4.4.10})에 의해 $X-X'$은 $X$에서
닫혀 있고 $f$는 닫힌 사상이므로 $U$는 열려 있다. 또한
(II, 6.2.2)에 의해 모든 $y\in U$에 대해 $f^{-1}(y)$는 유한하다.
$f$의 제한 $f^{-1}(U)\to U$는 고유이므로 (II, 5.4.1),
(\hyperref[III.4.4.2-fr]{4.4.2})에 의해 유한이다.
\end{proof}

\begin{env}[4.4.12]
\label{III.4.4.12-fr}
\emph{주의} (4.4.12). ---
\begin{enumerate}
\item[(i)]
(II, 6.2.7)에서 예고했듯이, 제V장에서 $Y$가 국소 뇌터이면 모든
준유한 분리 사상 $f:X\to Y$가 준아핀이고 따라서 준사영임을 보일
것이다. 그러면 주정리 (\hyperref[III.4.4.3-fr]{4.4.3})에서 $f$가
분리이고 유한형이라고만 가정해도 결론이 성립한다. 실제로
(\hyperref[III.4.4.10-fr]{4.4.10})에 의해 $X'$은 $X$에서 열려 있다.
$X$가 국소 뇌터이므로 $f$의 $X'$에 대한 제한은 여전히 유한형이고
(I, 6.3.5), $X'$의 정의에 의해 준유한이며 명백히 분리이다. 따라서
이 제한에 (\hyperref[III.4.4.3-fr]{4.4.3})을 적용하면 결론이 따른다.

\item[(ii)]
제IV장에서 차원 이론을 이용하여
(\hyperref[III.4.4.10-fr]{4.4.10})의 더 초등적인 증명을 제시하겠다.
\end{enumerate}
\end{env}

\subsection{준동형사상 가군의 완비화}
\label{subsection:III.4.5-fr}

\begin{proposition}[4.5.1]
\label{III.4.5.1-fr}
$A$를 뇌터환, $\mathfrak J$를 $A$의 아이디얼, $X$를 유한형
$A$-준스킴이라 하고, $\mathcal F$, $\mathcal G$를 그 받침들의
교집합이 $Y=\operatorname{Spec}(A)$ 위에서 고유인 두 연접
$\mathcal O_X$-가군이라 하자 (II, 5.4.10). 그러면 모든 정수
$n\geq0$에 대해
$\operatorname{Ext}^n_{\mathcal O_X}(X;\mathcal F,\mathcal G)$는 유한 생성
$A$-가군이고, 그 $\mathfrak J$-준아딕 위상에 대한 분리 완비화는
(\hyperref[III.4.1.7-fr]{4.1.7})의 표기 아래에서
$\operatorname{Ext}^n_{\mathcal O_{\widehat X}}(\widehat X;
\widehat{\mathcal F},\widehat{\mathcal G})$와 표준적으로 동일시된다.
\end{proposition}

\begin{proof}
(T, 4.2)에 의해, 수렴 대상이
$\operatorname{Ext}_{\mathcal O_X}(X;\mathcal F,\mathcal G)$이고
$\mathrm E_2$ 항들이
\[
\mathrm E_2^{pq}=\mathrm H^p\bigl(X,
\mathcal{E}\!xt^q_{\mathcal O_X}(\mathcal F,\mathcal G)\bigr)
\]
로 주어지는 쌍정칙 스펙트럼 열
$\mathrm E(\mathcal F,\mathcal G)$가 존재한다. 층
$\mathcal{E}\!xt^q_{\mathcal O_X}(\mathcal F,\mathcal G)$는 연접
$\mathcal O_X$-가군이고 (0, 12.3.3), 그 받침은 $\mathcal F$와
$\mathcal G$의 받침들의 교집합에 포함되므로 (T, 4.2.2), $Y$ 위에서
고유이다 (II, 5.4.10). 따라서
(\hyperref[III.3.2.4-fr]{3.2.4})에 의해 $\mathrm E_2^{pq}$들은
유한 생성 $A$-가군이고, 따라서 (0, 11.1.8)에 의해 스펙트럼 열의
모든 항 $\mathrm E_r^{pq}$와 그 수렴 대상도 그러하다.

한편 $i:\widehat X\to X$가 표준 사상이면 $\widehat{\mathcal F}$와
$\widehat{\mathcal G}$는 각각 $i^*(\mathcal F)$와
$i^*(\mathcal G)$로 표준적으로 동일시되고, $i$는 평탄하다
(I, 10.8.8 및 10.8.9). 그러면 (0, 12.3.4)에 의해 모든 $q\geq0$에
대해 표준 $i$-사상
\[
u_q:\mathcal{E}\!xt^q_{\mathcal O_X}(\mathcal F,\mathcal G)
\longrightarrow
\mathcal{E}\!xt^q_{\mathcal O_{\widehat X}}
(\widehat{\mathcal F},\widehat{\mathcal G})
\]
이 존재하고, 이에 대응하는 $\mathcal O_{\widehat X}$-준동형사상
\[
u_q^\#:i^*\bigl(\mathcal{E}\!xt^q_{\mathcal O_X}
(\mathcal F,\mathcal G)\bigr)
\longrightarrow
\mathcal{E}\!xt^q_{\mathcal O_{\widehat X}}
(\widehat{\mathcal F},\widehat{\mathcal G})
\]
은 동형사상이다 (0, 12.3.5). 다시 말해 (I, 10.8.8),
$\mathcal{E}\!xt^q_{\mathcal O_{\widehat X}}
(\widehat{\mathcal F},\widehat{\mathcal G})$는
(\hyperref[III.4.1.7-fr]{4.1.7})의 표기 아래에서 완비화
$(\mathcal{E}\!xt^q_{\mathcal O_X}(\mathcal F,\mathcal G))^\wedge$와
표준적으로 동일시된다. 이제 비교 정리
(\hyperref[III.4.1.10-fr]{4.1.10})에 의해 모든 $p\geq0$에 대해
\[
\mathrm H^p\bigl(\widehat X,
\mathcal{E}\!xt^q_{\mathcal O_{\widehat X}}
(\widehat{\mathcal F},\widehat{\mathcal G})\bigr)
\]
는 $\mathrm H^p\bigl(X,\mathcal{E}\!xt^q_{\mathcal O_X}
(\mathcal F,\mathcal G)\bigr)$의 $\mathfrak J$-준아딕 위상에 대한
분리 완비화와 표준적으로 동일시된다.

\oldpage[III]{139}

$\mathrm E(\widehat{\mathcal F},\widehat{\mathcal G})$를 (T, 4.2)에서
정의된 $\widehat{\mathcal F}$와 $\widehat{\mathcal G}$에 대한 쌍정칙
스펙트럼 열이라 하자. $\widehat A$가 $A$의 $\mathfrak J$-준아딕
위상에 대한 분리 완비화이면, 위에서 표준 동형사상 아래
\[
\mathrm E_2^{pq}(\widehat{\mathcal F},\widehat{\mathcal G})
=\mathrm E_2^{pq}(\mathcal F,\mathcal G)\otimes_A\widehat A
\qquad (0\textsubscript{I}, 7.3.3)
\]
를 얻는다.

이제 평탄 사상 $i$는 스펙트럼 열들의 표준 준동형사상
\[
\varphi:\mathrm E(\mathcal F,\mathcal G)\longrightarrow
\mathrm E(\widehat{\mathcal F},\widehat{\mathcal G})
=\mathrm E(i^*(\mathcal F),i^*(\mathcal G))
\]
을 정한다. 그 $\mathrm E_2$ 항들에서의 사상은
\[
\varphi_2^{pq}:\mathrm H^p\bigl(X,
\mathcal{E}\!xt^q_{\mathcal O_X}(\mathcal F,\mathcal G)\bigr)
\longrightarrow
\mathrm H^p\bigl(\widehat X,
\mathcal{E}\!xt^q_{\mathcal O_{\widehat X}}
(\widehat{\mathcal F},\widehat{\mathcal G})\bigr)
\]
이고, 수렴 대상에서의 사상은
\[
\varphi^n:\operatorname{Ext}^n_{\mathcal O_X}
(X;\mathcal F,\mathcal G)\longrightarrow
\operatorname{Ext}^n_{\mathcal O_{\widehat X}}
(\widehat X;\widehat{\mathcal F},\widehat{\mathcal G})
\]
이며, 이들은 함수성에 의해 각각 $u_q$와 $u_0$에서 얻어진다
(0, 12.3.4). $\widehat A$로 텐서하면 $\varphi_r^{pq}$와
$\varphi^n$에서 $\widehat A$-가군의 준동형사상
\[
\psi_r^{pq}:\mathrm E_r^{pq}(\mathcal F,\mathcal G)
\otimes_A\widehat A\longrightarrow
\mathrm E_r^{pq}(\widehat{\mathcal F},\widehat{\mathcal G})
\]
과
\[
\psi^n:\operatorname{Ext}^n_{\mathcal O_X}(X;\mathcal F,\mathcal G)
\otimes_A\widehat A\longrightarrow
\operatorname{Ext}^n_{\mathcal O_{\widehat X}}
(\widehat X;\widehat{\mathcal F},\widehat{\mathcal G})
\]
을 얻는다. $\widehat A$는 평탄 $A$-가군이므로
(0\textsubscript{I}, 7.3.3), $\widehat A$-가군
$\mathrm E_r^{pq}(\mathcal F,\mathcal G)\otimes_A\widehat A$들은
$\operatorname{Ext}^n_{\mathcal O_X}(X;\mathcal F,\mathcal G)
\otimes_A\widehat A$를 수렴 대상으로 하는 쌍정칙 스펙트럼 열을
이루고, $\psi_r^{pq}$와 $\psi^n$은 스펙트럼 열들의 사상을 이룬다.
$\psi_2^{pq}$들이 동형사상이므로 $\psi^n$들도 동형사상이다
(0, 11.1.5).
\end{proof}

\begin{corollary}[4.5.2]
\label{III.4.5.2-fr}
(\hyperref[III.4.5.1-fr]{4.5.1})의 가정 아래에서 $A$가 뇌터
$\mathfrak J$-아딕 환이라고 더 가정하자. 그러면 모든 정수
$n\geq0$에 대해
$\operatorname{Ext}^n_{\mathcal O_X}(X;\mathcal F,\mathcal G)$는
$\operatorname{Ext}^n_{\mathcal O_{\widehat X}}
(\widehat X;\widehat{\mathcal F},\widehat{\mathcal G})$와 표준적으로
동일시된다.
\end{corollary}

\begin{proof}
$\operatorname{Ext}^n_{\mathcal O_X}(X;\mathcal F,\mathcal G)$는 유한 생성
$A$-가군이므로 $\mathfrak J$-준아딕 위상에 대해 분리이고 완비임을
관찰하면 충분하다 (0\textsubscript{I}, 7.3.6).
\end{proof}

(\hyperref[III.4.5.1-fr]{4.5.1})에서 $n=0$인 특수한 경우는 다음과
같이 서술된다.

\begin{corollary}[4.5.3]
\label{III.4.5.3-fr}
(\hyperref[III.4.5.1-fr]{4.5.1})의 가정 아래에서, 임의의 준동형사상
$u:\mathcal F\to\mathcal G$에 대해 $\widehat u:\widehat{\mathcal F}
\to\widehat{\mathcal G}$를 그 완비화된 준동형사상이라 하자
(I, 10.8.4). 그러면 표준 동형사상
\[
\label{III.4.5.3.1-fr}
\left(\operatorname{Hom}_{\mathcal O_X}(\mathcal F,\mathcal G)\right)^\wedge
\xrightarrow{\sim}
\operatorname{Hom}_{\mathcal O_{\widehat X}}
(\widehat{\mathcal F},\widehat{\mathcal G})
\tag{4.5.3.1}
\]
이 존재한다. 여기서 왼쪽 항은 $A$-가군
$\operatorname{Hom}_{\mathcal O_X}(\mathcal F,\mathcal G)$의
$\mathfrak J$-준아딕 위상에 대한 분리 완비화이고, 이 동형사상은
준동형사상 $u\mapsto\widehat u$에서 분리 완비화로 이행하여 얻어진다.
\end{corollary}

\subsection{형식 사상과 보통 사상의 관계}

\begin{proposition}[4.6.1]
\label{III.4.6.1-fr}
$Y$를 국소 뇌터 준스킴, $f:X\to Y$를 고유 사상,
$\mathcal F$를 연접이고 $f$-평탄인 $\mathcal O_X$-가군, $y$를
$Y$의 점이라 하자. 어떤 정수 $n$에 대해
\[
H^n\bigl(f^{-1}(y),\mathcal F\otimes_{\mathcal O_Y}k(y)\bigr)=0
\]
이라고 가정하자. 그러면 $Y$ 안의 $y$의 어떤 근방 $U$가 존재하여
$R^nf_*(\mathcal F)|U=0$이고, 모든 정수 $p\geq0$에 대해
(\hyperref[III.4.2.1.1-fr]{4.2.1.1})의 표준 준동형사상
\[
\bigl(R^{n-1}f_*(\mathcal F)\bigr)_y
\longrightarrow
H^{n-1}\bigl(f^{-1}(y),
\mathcal F\otimes_{\mathcal O_Y}
(\mathcal O_y/\mathfrak m_y^{p+1})\bigr)
\]
은 전사이다.
\end{proposition}

\begin{proof}
$R^nf_*(\mathcal F)$는 연접 $\mathcal O_Y$-가군이므로
(\hyperref[III.3.2.1-fr]{3.2.1}), 명제의 첫째 주장을 보이려면
$(R^nf_*(\mathcal F))_y=0$임을 보이면 충분하다
(0\textsubscript{I}, 5.2.2). (\hyperref[III.4.2.1-fr]{4.2.1})에
의해, 이를 위해서는 모든 $p$에 대해
\[
H^n\bigl(f^{-1}(y),\mathcal F\otimes_{\mathcal O_Y}
(\mathcal O_y/\mathfrak m_y^{p+1})\bigr)=0
\]
임을 보이면 충분하다. $p=0$일 때는 가정에 의해 참이므로 $p$에
관한 귀납법으로 보이겠다. 다음과 같이 놓자.
\[
X_p=X\times_Y\operatorname{Spec}
(\mathcal O_y/\mathfrak m_y^{p+1}).
\]
그러면 $X_{p-1}$은 $X_p$와 같은 바탕공간을 갖는 닫힌 부분준스킴이다
(I, 3.6.1). 따라서 귀납 가정
\[
H^n\bigl(X_{p-1},\mathcal F\otimes_{\mathcal O_Y}
(\mathcal O_y/\mathfrak m_y^p)\bigr)=0
\]
에서
\[
H^n\bigl(X_p,\mathcal F\otimes_{\mathcal O_Y}
(\mathcal O_y/\mathfrak m_y^p)\bigr)=0
\]
이 따른다. 한편 $\mathcal O_{X_p}$-가군의 완전열
\[
0\longrightarrow
\mathfrak m_y^p\mathcal F/\mathfrak m_y^{p+1}\mathcal F
\longrightarrow
\mathcal F/\mathfrak m_y^{p+1}\mathcal F
\longrightarrow
\mathcal F/\mathfrak m_y^p\mathcal F
\longrightarrow0
\]
에서 얻는 코호몰로지 완전열은
\[
H^n\bigl(X_p,
\mathfrak m_y^p\mathcal F/\mathfrak m_y^{p+1}\mathcal F\bigr)
\longrightarrow
H^n\bigl(X_p,\mathcal F/\mathfrak m_y^{p+1}\mathcal F\bigr)
\longrightarrow
H^n\bigl(X_p,\mathcal F/\mathfrak m_y^p\mathcal F\bigr)
\]
이 완전임을 준다. 따라서
\[
\label{III.4.6.1.1-fr}
H^n\bigl(X_p,
\mathfrak m_y^p\mathcal F/\mathfrak m_y^{p+1}\mathcal F\bigr)=0
\tag{4.6.1.1}
\]
임을 보이면 충분하다. 실제로 그러면
$H^n(X_p,\mathcal F/\mathfrak m_y^{p+1}\mathcal F)$는
$H^n(X_p,\mathcal F/\mathfrak m_y^p\mathcal F)$의 부분가군이고,
귀납 가정에 의해 $0$이다.

\oldpage[III]{140}

이제 올
$Z=f^{-1}(y)=X\times_Y\operatorname{Spec}(k(y))$는 $X_p$의 닫힌
부분준스킴이고,
$\mathfrak m_y^p\mathcal F/\mathfrak m_y^{p+1}\mathcal F$는
$\mathfrak m_y$에 의해 소멸된다. 따라서 이를 $\mathcal O_Z$-가군으로
볼 수 있으며
\[
H^n\bigl(Z,\mathfrak m_y^p\mathcal F/
\mathfrak m_y^{p+1}\mathcal F\bigr)
=H^n\bigl(X_p,\mathfrak m_y^p\mathcal F/
\mathfrak m_y^{p+1}\mathcal F\bigr).
\]
표준 $\mathcal O_Z$-준동형사상
\[
\label{III.4.6.1.2-fr}
(\mathcal F/\mathfrak m_y\mathcal F)
\otimes_{k(y)}(\mathfrak m_y^p/\mathfrak m_y^{p+1})
\longrightarrow
\mathfrak m_y^p\mathcal F/\mathfrak m_y^{p+1}\mathcal F
\tag{4.6.1.2}
\]
이 전단사임을 보이자. 이를 보이면
$\mathfrak m_y^p/\mathfrak m_y^{p+1}$가 자유 $k(y)$-가군이므로
\[
\begin{split}
H^n\bigl(Z,\mathfrak m_y^p\mathcal F/
\mathfrak m_y^{p+1}\mathcal F\bigr)
&=H^n\bigl(Z,\mathcal F/\mathfrak m_y\mathcal F\bigr)
\otimes_{k(y)}(\mathfrak m_y^p/\mathfrak m_y^{p+1})\\
&=0
\end{split}
\]
을 얻는다 (0, 12.2.3). 마지막 등식은 가정에 의해
$H^n(Z,\mathcal F/\mathfrak m_y\mathcal F)=0$이기 때문이고, 이로써
(\hyperref[III.4.6.1.1-fr]{4.6.1.1})을 얻는다. 이제 첫째 주장을
마치려면 (\hyperref[III.4.6.1.2-fr]{4.6.1.2})가 전단사임만 보이면
된다. 이 문제는 $X$ 위에서 점별 문제이고, 가정에 의해 모든
$x\in f^{-1}(y)$에 대해 $\mathcal F_x$는 평탄
$\mathcal O_y$-가군이다. 또한
$\mathcal F_x/\mathfrak m_y\mathcal F_x$는 체
$k(y)=\mathcal O_y/\mathfrak m_y$ 위의 평탄 가군이므로
(0, 10.2.1, c))를 적용하면 된다.

(\hyperref[III.4.6.1-fr]{4.6.1})의 둘째 주장을 보이기 위해서는
(\hyperref[III.4.2.1-fr]{4.2.1})에서와 같이 곧바로 $Y$가 아핀이고
$y$가 닫힌 경우로 귀착된다. 관계
(\hyperref[III.4.6.1.1-fr]{4.6.1.1})은 같은 논증으로 모든 $k>0$과
모든 $p\geq0$에 대해
\[
\label{III.4.6.1.3-fr}
H^n\bigl(X_p,
\mathfrak m_y^k\mathcal F/
\mathfrak m_y^{k+p+1}\mathcal F\bigr)=0
\tag{4.6.1.3}
\]
을 얻는다.

\oldpage[III]{141}

그러므로 (\hyperref[III.4.2.1-fr]{4.2.1})에 의해
\[
\label{III.4.6.1.4-fr}
\bigl(R^nf_*(\mathfrak m_y^k\mathcal F)\bigr)_y=0
\tag{4.6.1.4}
\]
도 성립한다. 이제 코호몰로지 완전열에서
\[
\bigl(R^{n-1}f_*(\mathcal F)\bigr)_y
\longrightarrow
\bigl(R^{n-1}f_*(\mathcal F/\mathfrak m_y^p\mathcal F)\bigr)_y
\longrightarrow
\bigl(R^nf_*(\mathfrak m_y^p\mathcal F)\bigr)_y=0
\]
의 완전성을 얻는다. $y$가 닫혀 있고 $Y$가 아핀이므로
(\hyperref[III.1.4.11-fr]{1.4.11})
\[
R^{n-1}f_*(\mathcal F/\mathfrak m_y^p\mathcal F)
=\bigl(H^{n-1}(X,\mathcal F/\mathfrak m_y^p\mathcal F)\bigr)^\sim
=\bigl(H^{n-1}(f^{-1}(y),
\mathcal F/\mathfrak m_y^p\mathcal F)\bigr)^\sim
\]
이다 (G, II, 4.9.1). 그런데
$H^{n-1}(f^{-1}(y),\mathcal F/\mathfrak m_y^p\mathcal F)$는
$(\mathcal O_y/\mathfrak m_y^p)$-가군이므로
\[
\bigl(R^{n-1}f_*(\mathcal F/\mathfrak m_y^p\mathcal F)\bigr)_y
=H^{n-1}(f^{-1}(y),\mathcal F/\mathfrak m_y^p\mathcal F).
\]
이로써 (\hyperref[III.4.6.1-fr]{4.6.1})의 증명이 끝난다.
\end{proof}

\begin{corollary}[4.6.2]
\label{III.4.6.2-fr}
$Y$를 국소 뇌터 준스킴, $f:X\to Y$를 고유이고 평탄인 사상,
$\mathcal F$, $\mathcal G$를 두 국소 자유 $\mathcal O_X$-가군,
$y$를 $Y$의 점이라 하자. 다음과 같이 놓자.
\[
X_y=f^{-1}(y)=X\otimes_Yk(y),\qquad
\mathcal F_y=\mathcal F\otimes_{\mathcal O_Y}k(y),\qquad
\mathcal G_y=\mathcal G\otimes_{\mathcal O_Y}k(y).
\]
또한
\[
\label{III.4.6.2.1-fr}
H^1\bigl(X_y,
\mathcal{H}\!om_{\mathcal O_{X_y}}(\mathcal F_y,\mathcal G_y)\bigr)=0
\tag{4.6.2.1}
\]
이라고 가정하자. 그러면 모든 준동형사상
$u_0:\mathcal F_y\to\mathcal G_y$에 대해 $y$의 열린 근방 $U$와
준동형사상
\[
u:\mathcal F|f^{-1}(U)\longrightarrow\mathcal G|f^{-1}(U)
\]
가 존재하여 $u_0=u\otimes1$이다.
\end{corollary}

\begin{proof}
실제로 이 가정에 의해 $n=1$, $p=0$으로
(\hyperref[III.4.6.1-fr]{4.6.1})을 연접 $\mathcal O_X$-가군
\[
\mathcal H=\mathcal{H}\!om_{\mathcal O_X}(\mathcal F,\mathcal G)
\]
에 적용할 수 있다. $\mathcal H$는 국소 자유이므로 특히
$f$-평탄이고, 이때 $\mathcal O_{X_y}$-가군
$\mathcal H\otimes_{\mathcal O_Y}k(y)$는
$\mathcal{H}\!om_{\mathcal O_{X_y}}(\mathcal F_y,\mathcal G_y)$와
동일시된다 (0\textsubscript{I}, 6.2.2). $Y=\operatorname{Spec}(A)$가
아핀이라고 가정할 수 있고, 그러면
(\hyperref[III.1.4.11-fr]{1.4.11})
\[
R^0f_*(\mathcal H)=
\bigl(\operatorname{Hom}_{\mathcal O_X}(\mathcal F,\mathcal G)\bigr)^\sim,
\]
따라서
\[
\bigl(R^0f_*(\mathcal H)\bigr)_y
=\operatorname{Hom}_{\mathcal O_X}(\mathcal F,\mathcal G)
\otimes_A\mathcal O_y.
\]
(\hyperref[III.4.6.1-fr]{4.6.1})에 의해 표준 준동형사상
\[
\operatorname{Hom}_{\mathcal O_X}(\mathcal F,\mathcal G)
\otimes_A\mathcal O_y
\longrightarrow
\operatorname{Hom}_{\mathcal O_{X_y}}(\mathcal F_y,\mathcal G_y)
\]
은 전사이다. 또한
$\operatorname{Hom}_{\mathcal O_X}(\mathcal F,\mathcal G)
\otimes_A\mathcal O_y$의 모든 원소는 $s\notin\mathfrak m_y$인
$s\in A$와 어떤 $u$를 써서 $u\otimes(1/s)$의 꼴로 나타낼 수 있다.
이로써 따름정리가 성립한다.
\end{proof}

이 따름정리는 다음 결과로 보충된다.

\begin{corollary}[4.6.3]
\label{III.4.6.3-fr}
(\hyperref[III.4.6.2-fr]{4.6.2})의 가정 아래에서 $u_0$이 단사이면
(각각 전사이면, 전단사이면) $u$도 그러하다고 가정할 수 있다.
\end{corollary}

\begin{proof}
$U=Y$인 경우로 한정할 수 있다. $u_0$이 단사이면 (각각 전사이면)
모든 $x\in f^{-1}(y)$에 대해 $\operatorname{Ker}u_x=0$임을
(각각 $\operatorname{Coker}u_x=0$임을) 보이면 충분하다. 실제로
$\operatorname{Ker}u$와 $\operatorname{Coker}u$는 연접
$\mathcal O_X$-가군이므로 (0\textsubscript{I}, 5.3.4),
$f^{-1}(y)$의 어떤 근방 $V$가 존재하여
$\operatorname{Ker}u$의 (각각 $\operatorname{Coker}u$의) $V$에
대한 제한이 $0$이다 (0\textsubscript{I}, 5.2.2). $f$는 닫힌
사상이므로 $y$의 어떤 근방 $U'\subset U$가 존재하여
$f^{-1}(U')\subset V$이고, 이로써
(\hyperref[III.4.6.3-fr]{4.6.3})이 증명된다.

가정에 의해
\[
u_x\otimes1:
\mathcal F_x\otimes_{\mathcal O_y}k(y)
\longrightarrow
\mathcal G_x\otimes_{\mathcal O_y}k(y)
\]
는 단사이다 (각각 전사이다). $\mathcal F_x$와 $\mathcal G_x$는
유한 생성 자유 $\mathcal O_x$-가군이고, $\mathcal O_x$는 평탄
$\mathcal O_y$-가군이다. $u_x\otimes1$이 단사라고 가정하면
(0, 10.2.4)에 의해 $u_x$도 단사이다. $u_x\otimes1$이 전사라고
가정하면 몫을 취하여 얻는 준동형사상
\[
\mathcal F_x/\mathfrak m_x\mathcal F_x
\longrightarrow
\mathcal G_x/\mathfrak m_x\mathcal G_x
\]
은 더욱이 전사이다. $\mathcal G_x$는 유한 생성
$\mathcal O_x$-가군이고 $\mathcal O_x$는 극대 아이디얼이
$\mathfrak m_x$인 국소환이므로, 나카야마 보조정리에서 결론이
따른다 (Bourbaki, \emph{Alg.}, chap. VIII, \S\,6, no~3,
cor. 4 de la prop. 6).
\end{proof}

특히 (\hyperref[III.4.6.3-fr]{4.6.3})에서 다음이 따른다.

\begin{corollary}[4.6.4]
\label{III.4.6.4-fr}
$Y$를 국소 뇌터 준스킴, $f:X\to Y$를 고유이고 평탄인 사상,
$y$를 $Y$의 점이라 하고 $X_y=X\otimes_Yk(y)$로 놓자.
$\mathcal E_0$을
\[
\label{III.4.6.4.1-fr}
H^1\bigl(X_y,
\mathcal{H}\!om_{\mathcal O_{X_y}}(\mathcal E_0,\mathcal E_0)\bigr)=0
\tag{4.6.4.1}
\]
을 만족하는 국소 자유 $\mathcal O_{X_y}$-가군이라 하자.
$\mathcal F$, $\mathcal G$를 두 국소 자유 $\mathcal O_X$-가군이라
하되, (\hyperref[III.4.6.2-fr]{4.6.2})의 표기에서
$\mathcal F_y$와 $\mathcal G_y$가 $\mathcal E_0$과 동형이라고
하자. 그러면 $y$의 어떤 열린 근방 $U$가 존재하여
$\mathcal F|f^{-1}(U)$와 $\mathcal G|f^{-1}(U)$는 동형이다.
\end{corollary}

더 특별히 다음을 얻는다.

\begin{corollary}[4.6.5]
\label{III.4.6.5-fr}
$f$, $X$, $Y$에 대해 (\hyperref[III.4.6.4-fr]{4.6.4})의 가정이
성립하고 $H^1(X_y,\mathcal O_{X_y})=0$이라고 하자.
$\mathcal F$와 $\mathcal G$가 두 가역 $\mathcal O_X$-가군이고
$\mathcal F_y$와 $\mathcal G_y$가 동형이면, $y$의 어떤 열린 근방
$U$가 존재하여 $\mathcal F|f^{-1}(U)$와
$\mathcal G|f^{-1}(U)$는 동형이다.
\end{corollary}

\begin{proof}
(\hyperref[III.4.6.4-fr]{4.6.4})을
$\mathcal F^{-1}\otimes\mathcal G$와 $\mathcal O_X$에 적용하면
충분하다.
\end{proof}

\begin{remarks}[4.6.6]
\label{III.4.6.6-fr}
\begin{enumerate}
\item[(i)]
(\hyperref[III.4.6.5-fr]{4.6.5})를 이용하여 제V장에서
(II, 4.2.7)에 예고한 사영 다발 위의 가역층의 분류를 확립하겠다.
\item[(ii)]
(\hyperref[III.4.6.1-fr]{4.6.1})의 결과는 \S\,7에서 더 일반적인
명제들의 귀결로 다시 나타날 것이다.
\end{enumerate}
\end{remarks}

\begin{proposition}[4.6.7]
\label{III.4.6.7-fr}
$Z$를 국소 뇌터 준스킴, $X$, $Y$를 두 $Z$-준스킴이라 하고, 구조
사상 $g:X\to Z$, $h:Y\to Z$가 고유라고 하자. $f:X\to Y$를
$Z$-사상, $z$를 $Z$의 점이라 하고
\[
f_z=f\times_Z1:
X\otimes_Zk(z)\longrightarrow Y\otimes_Zk(z)
\]
로 놓자.
\begin{enumerate}
\item[(i)]
$f_z$가 유한 사상이면 (각각 닫힌 몰입 사상이면), $z$의 어떤 열린
근방 $U$가 존재하여 $f$의 제한
$g^{-1}(U)\to h^{-1}(U)$는 유한 사상이다 (각각 닫힌 몰입
사상이다).
\item[(ii)]
더 나아가 $g$가 평탄 사상이라고 가정하자. $f_z$가 동형사상이면
$z$의 어떤 열린 근방 $U$가 존재하여 $f$의 제한
$g^{-1}(U)\to h^{-1}(U)$는 동형사상이다.
\end{enumerate}
\end{proposition}

\begin{proof}
두 경우 모두 모든 $y\in h^{-1}(z)$에 대해 $y$의 어떤 근방
$V_y$가 존재하여 $f$의 제한 $f^{-1}(V_y)\to V_y$가 유한
사상임을 (각각 닫힌 몰입 사상임을, 동형사상임을) 보이면 충분하다.
실제로 $V$를 $V_y$들의 합집합이라 하면 $f$의 제한
$f^{-1}(V)\to V$도 유한 사상이다 (각각 닫힌 몰입 사상이다,
동형사상이다) (II, 6.1.1 및 I, 4.2.4). $h$가 닫힌 사상이므로
$z$의 어떤 열린 근방 $U$가 존재하여 $h^{-1}(U)\subset V$이고,
명제가 증명된다.

\noindent\textit{(i)}
먼저 $f$가 고유 사상임을 주의하자 (II, 5.4.3). $f_z$가 유한이라고
가정하면, 모든 $y\in h^{-1}(z)$에 대해
$f^{-1}(V_y)\to V_y$가 유한이 되는 근방 $V_y$의 존재는
(\hyperref[III.4.4.11-fr]{4.4.11})에서 따른다.

\oldpage[III]{143}

$f_z$가 닫힌 몰입 사상인 경우를 다루기 위해서는 이미 $f$가 유한
사상이라고 가정할 수 있다. 따라서
$X=\operatorname{Spec}(\mathcal B)$이고, 여기서 $\mathcal B$는 연접
$\mathcal O_Y$-대수이며, $f$는 표준 준동형사상
$u:\mathcal O_Y\to\mathcal B$에 대응한다 (II, 1.2.7).
모든 $y\in h^{-1}(z)$에 대해
$u_y:\mathcal O_y\to\mathcal B_y$가 전사임을 보이면, 층
$\operatorname{Coker}u$가 연접이므로 (0\textsubscript{I}, 5.3.4
및 5.2.2) $y$의 어떤 근방 $V_y$에서 $u|V_y$가 전사임이 따를
것이다. 한편 유한 사상 $f_z$는 준동형사상
\[
v=u\otimes1:
\mathcal O_Y\otimes_{\mathcal O_Z}k(z)
\longrightarrow
\mathcal B\otimes_{\mathcal O_Z}k(z)
\]
에 대응하고, $f_z$가 닫힌 몰입 사상이라는 가정에서
\[
u_y\otimes1:
\mathcal O_y\otimes_{\mathcal O_z}k(z)
\longrightarrow
\mathcal B_y\otimes_{\mathcal O_z}k(z)
\]
가 전사임이 따른다. $\mathcal B_y$는 유한 생성
$\mathcal O_y$-가군이고 $\mathcal O_y$는 뇌터 국소환이므로,
(\hyperref[III.4.6.3-fr]{4.6.3})에서와 같이 나카야마 보조정리에서
결론이 따른다.

\noindent\textit{(ii)}
위와 같은 논증에 의해, 이제는 $u_y\otimes1$이 전단사임을 알고
$u_y$가 전단사임을 보이면 충분하다. 이는 다음 보조정리에서 따른다.
\end{proof}

\begin{lemma}[4.6.7.1]
\label{III.4.6.7.1-fr}
$A$, $B$를 두 뇌터 국소환, $\rho:A\to B$를 국소 준동형사상,
$u:N\to M$을 $B$-가군의 준동형사상이라 하자. $M$은 평탄
$A$-가군이고 $N$은 유한 생성 $B$-가군이며
\[
u\otimes1:N\otimes_Ak\longrightarrow M\otimes_Ak
\]
가 단사라고 가정하자. 여기서 $k$는 $A$의 잉여체이다. 그러면
$N$은 평탄 $A$-가군이고 $u$는 단사이다.
\end{lemma}

\begin{proof}
첫째 주장을 보이려면 유한 생성 $A$-가군의 임의의 쌍 $P$, $Q$와
임의의 단사 $A$-준동형사상 $v:P\to Q$에 대해
\[
1_N\otimes v:N\otimes_AP\longrightarrow N\otimes_AQ
\]
가 단사임을 보이면 된다. 그런데 가환 그림
\[
\xymatrix{
N\otimes_AP\ar[r]^{1_N\otimes v}\ar[d]_{u\otimes1_P}
  &N\otimes_AQ\ar[d]^{u\otimes1_Q}\\
M\otimes_AP\ar[r]_{1_M\otimes v}
  &M\otimes_AQ
}
\]
이 있고, 가정에 의해 $1_M\otimes v$가 단사이므로
$u\otimes1_P$가 단사임을 보이면 충분하다.

$\mathfrak m$을 $A$의 극대 아이디얼이라 하자.
$A$-가군 $N\otimes_AP$ 위의 $\mathfrak m$-아딕 여과는 이를
$B$-가군으로 보았을 때의 $\mathfrak mB$-아딕 여과이기도 하다.
$B$는 뇌터환이고 $\mathfrak mB$는 $B$의 야코브슨 근기에 포함되며,
$N$이 유한 생성 $B$-가군이고 $P$가 유한 생성 $A$-가군이므로
$N\otimes_AP$는 유한 생성 $B$-가군이다. 따라서 이 여과가 정하는
위상은 분리이다 (0\textsubscript{I}, 7.3.5). 그러므로
\[
\operatorname{gr}(u\otimes1_P):
\operatorname{gr}_\bullet(N\otimes_AP)
\longrightarrow
\operatorname{gr}_\bullet(M\otimes_AP)
\]
가 단사임을 보이면 충분하다. 여기서 등급 가군들은
$\mathfrak m$-아딕 여과에 관한 것이다
(Bourbaki, \emph{Alg. comm.}, chap. III, \S\,2, no~8,
cor. 1 du th. 1).

$M$은 평탄 $A$-가군이므로 준동형사상
\[
M\otimes_A(\mathfrak m^nP)
\longrightarrow
\mathfrak m^n(M\otimes_AP)
\]
들은 전단사이다. 따라서 표준 준동형사상
\[
\varphi_M:
\operatorname{gr}_0(M)\otimes_A\operatorname{gr}_\bullet(P)
\longrightarrow
\operatorname{gr}_\bullet(M\otimes_AP)
\]
도 전단사이다.

\oldpage[III]{144}

이제 가환 그림
\[
\xymatrix{
\operatorname{gr}_0(N)\otimes_A\operatorname{gr}_\bullet(P)
  \ar[r]^{\operatorname{gr}_0(u)\otimes1}\ar[d]_{\varphi_N}
  &\operatorname{gr}_0(M)\otimes_A\operatorname{gr}_\bullet(P)
  \ar[d]^{\varphi_M}\\
\operatorname{gr}_\bullet(N\otimes_AP)
  \ar[r]_{\operatorname{gr}(u\otimes1)}
  &\operatorname{gr}_\bullet(M\otimes_AP)
}
\]
이 있다. 여기서 $\varphi_M$은 전단사이고 $\varphi_N$은 전사이다.
또한 가정에 의해 $\operatorname{gr}_0(u)$가 단사이고
\[
\begin{split}
\operatorname{gr}_0(N)\otimes_A\operatorname{gr}_\bullet(P)
&=\operatorname{gr}_0(N)\otimes_k\operatorname{gr}_\bullet(P),\\
\operatorname{gr}_0(M)\otimes_A\operatorname{gr}_\bullet(P)
&=\operatorname{gr}_0(M)\otimes_k\operatorname{gr}_\bullet(P)
\end{split}
\]
이므로 $\operatorname{gr}_0(u)\otimes1$도 단사이다. 따라서
$\operatorname{gr}(u\otimes1)$은 단사이고, 이로써 첫째 주장이
증명된다. 앞의 논증에서 $P=A$로 놓으면 둘째 주장이 따른다.
\end{proof}

\begin{proposition}[4.6.8]
\label{III.4.6.8-fr}
$Z$를 국소 뇌터 준스킴, $X$, $Y$를 두 $Z$-준스킴이라 하고, 구조
사상 $g:X\to Z$, $h:Y\to Z$가 고유라고 하자. $Z'$을 $Z$의 닫힌
부분집합, $X'=g^{-1}(Z')$, $Y'=h^{-1}(Z')$을 그 역상들이라 하고,
\[
\widehat X=X_{/X'},\qquad
\widehat Y=Y_{/Y'},\qquad
\widehat Z=Z_{/Z'}
\]
을 이 닫힌 부분집합들을 따른 $X$, $Y$, $Z$의 형식적 완비화라
하자. $f:X\to Y$를 $Z$-사상,
$\widehat f:\widehat X\to\widehat Y$를 그 완비화로의 연장이라
하자. $\widehat f$가 동형사상이기 위한 (각각 닫힌 몰입 사상이기
위한) 필요충분조건은, $Z'$의 어떤 열린 근방 $U$가 존재하여 $f$의
제한 $g^{-1}(U)\to h^{-1}(U)$가 동형사상인 것이다
(각각 닫힌 몰입 사상인 것이다).
\end{proposition}

\begin{proof}
조건의 충분성은 즉시 따른다 (I, 10.14.7). 필요성을 보이기
위해서는 (\hyperref[III.4.6.7-fr]{4.6.7})에서와 같은 논증에 의해
모든 $y\in Y'$에 대해 $y$의 어떤 열린 근방 $V_y$가 존재하여
$f$의 제한 $f^{-1}(V_y)\to V_y$가 동형사상임을 (각각 닫힌 몰입
사상임을) 보이면 충분하다. 따라서 $Y=Z=\operatorname{Spec}(A)$이고
$A$가 뇌터환인 경우로 귀착된다.

가정에 의해 (I, 10.9.1 및 10.14.2) 모든 $y\in Y'$에 대해 올
$f^{-1}(y)$는 한 점 이하이다. 따라서 $f$가 고유이므로
(II, 5.4.3), $y$의 어떤 열린 근방 $U$가 존재하여 $f$의 제한
$f^{-1}(U)\to U$가 유한 사상이다
(\hyperref[III.4.4.11-fr]{4.4.11}). 그러므로 $f$가 유한
사상이라고 가정할 수 있고, 따라서 $X=\operatorname{Spec}(B)$이며
$B$는 유한 $A$-대수이다. $Y'=V(\mathfrak J)$이면
\[
\widehat Y=\operatorname{Spf}(\widehat A),\qquad
\widehat X=\operatorname{Spf}(\widehat B),
\]
여기서 $\widehat A$는 $A$의 $\mathfrak J$-준아딕 위상에 대한
분리 완비화이고, $\widehat B$는 $B$의 $\mathfrak JB$-준아딕
위상에 대한 분리 완비화, 또는 이와 같은 것으로서 $A$-가군 $B$의
$\mathfrak J$-준아딕 위상에 대한 분리 완비화이다. 또한
$\widehat f$는 표준 환 준동형사상 $\varphi:A\to B$의 연속적 연장
\[
\widehat\varphi:\widehat A\longrightarrow\widehat B
\]
에 대응하는 아핀 형식 스킴의 사상이다. 가정은
$\widehat\varphi$가 전단사라는 것이다 (각각 전사라는 것이다)
(I, 10.14.2). 그런데 $\widehat\varphi$는 $\varphi$를
$A$-가군의 준동형사상으로 보았을 때의 연속적 연장이기도 하다.
따라서 (I, 10.8.14)에 의해 $Y'$의 어떤 열린 근방 $U$가 존재하여
$\mathcal O_Y$-가군의 준동형사상
$\widetilde\varphi:\widetilde A\to\widetilde B$의 $U$에 대한
제한은 전단사이다 (각각 전사이다). 이로써 증명이 끝난다.
\end{proof}

\oldpage[III]{145}

\subsection{풍부성 판정법}

\begin{theorem}[4.7.1]
\label{III.4.7.1-fr}
$Y$를 국소 뇌터 준스킴, $f:X\to Y$를 고유 사상,
$\mathcal L$을 가역 $\mathcal O_X$-가군, $y$를 $Y$의 점이라 하자.
$X_y=X\otimes_Yk(y)=f^{-1}(y)$라 하고 $g$를 $X_y$에서 $X$로 가는
사영이라 하자. $\mathcal L_y=g^*(\mathcal L)
=\mathcal L\otimes_{\mathcal O_Y}k(y)$가 $X_y$ 위에서 풍부하면,
$Y$ 안의 $y$의 어떤 열린 근방 $U$가 존재하여
$\mathcal L|f^{-1}(U)$는 $f$의 $f^{-1}(U)$에 대한 제한에 관해
풍부하다.
\end{theorem}

\begin{proof}
\noindent\textit{I)}
다음과 같이 놓자.
\[
Y'=\operatorname{Spec}(\mathcal O_y),\qquad
X'=X\times_YY',\qquad
\mathcal L'=\mathcal L\otimes_{\mathcal O_Y}\mathcal O_y.
\]
먼저 $\mathcal L'$이 $f'=f_{(Y')}$에 관해 풍부함을 보이겠다.
가환 그림
\[
\xymatrix{
X\ar[d]_f & X'\ar[l]\ar[d]_{f'} & X_y\ar[l]\ar[d]_{f_y}\\
Y & Y'\ar[l] & \operatorname{Spec}(k(y))\ar[l]
}
\]
이 있다. $f'$은 고유이고 (II, 5.4.2, (iii)) $\mathcal O_y$는
뇌터환이므로, $Y=Y'=\operatorname{Spec}(\mathcal O_y)$이고
따라서 $X=X'$인 경우로 한정할 수 있다. 이때 $\mathcal L_y$가
$f_y$에 관해 풍부하다고 가정하고 $\mathcal L$이 $f$에 관해
풍부함을 보이면 된다 (II, 4.6.6).

판정법 (\hyperref[III.2.6.1-fr]{2.6.1, c)})을 적용하겠다. 실제로
모든 연접 $\mathcal O_X$-가군 $\mathcal F$에 대해 어떤 정수 $N$이
존재하여
\[
H^1(X,\mathcal F(n))=0\quad\text{$n\geq N$인 모든 $n$에 대해},
\qquad
\mathcal F(n)=\mathcal F\otimes\mathcal L^{\otimes n}
\]
임을 보이겠다. $y$는 $\mathcal O_y$의 극대 아이디얼
$\mathfrak m$에 대응하는 $Y$의 닫힌점이다. 따라서 $X_y$는 연접
아이디얼
\[
\mathcal J=f^*(\widetilde{\mathfrak m})\mathcal O_X
=\mathfrak m\mathcal O_X
\]
에 의해 정의되는 $X$의 닫힌 부분준스킴이고 (I, 4.4.5), $g$는
표준 단사이다. 이제 등급 $k(y)$-대수
\[
S=\operatorname{gr}(\mathcal O_y)
=\bigoplus_{j\geq0}\mathfrak m^j/\mathfrak m^{j+1}
\]
를 생각하자. $\mathcal O_y$가 뇌터환이므로 $S$는 유한형이다.
그러므로 $\mathcal O_X$-대수
$\mathcal S=f^*(\widetilde S)$는 준연접이고 유한형이다.
또한 이는 분명히 $\mathcal J$에 의해 소멸된다. 따라서
$\mathcal S_y=g^*(\mathcal S)$로 놓으면 $\mathcal S_y$는 준연접
유한형 $\mathcal O_{X_y}$-대수이고
$\mathcal S=g_*(\mathcal S_y)$이다.

한편
\[
\mathcal M_j=\mathfrak m^j\mathcal F/
\mathfrak m^{j+1}\mathcal F,\qquad
\mathcal M=\bigoplus_{j\geq0}\mathcal M_j
=\operatorname{gr}(\mathcal F)
\]
로 놓자. $\mathcal F$가 연접이므로 $\mathcal M$은 준연접 유한형
$\mathcal S$-가군이고 (0, 10.1.1), 역시 $\mathcal J$에 의해
소멸된다. 따라서
\[
\mathcal M'_j=g^*(\mathcal M_j),\qquad
\mathcal M'=g^*(\mathcal M)
=\bigoplus_{j\geq0}\mathcal M'_j
\]
로 놓으면 $\mathcal M'$은 준연접 유한형 등급
$\mathcal S_y$-가군이고 $\mathcal M=g_*(\mathcal M')$이다.
더 나아가
$\mathcal M'_j(n)=\mathcal M'_j\otimes\mathcal L_y^{\otimes n}$으로
놓으면 $\mathcal M'_j(n)=g^*(\mathcal M_j(n))$이다.

$f_y$는 고유이고 (II, 5.4.2, (iii)) $\mathcal L_y$는
풍부하므로 $f_y$는 사영 사상이다 (II, 5.5.4 및 4.6.11).
따라서 $\operatorname{Spec}(k(y))$, $f_y$, $\mathcal S_y$,
$\mathcal L_y$, $\mathcal M'$에 정리
(\hyperref[III.2.4.1-fr]{2.4.1, (ii)})를 적용할 수 있다.
어떤 정수 $N$이 존재하여 $n\geq N$이면 모든 $q>0$과 모든 $j$에
대해
\[
H^q(X_y,\mathcal M'_j(n))=0
\]
이다. 따라서 모든 $q>0$과 모든 $j$에 대해
$H^q(X,\mathcal M_j(n))=0$이기도 하다 (G, II, 4.9.1).
이제
\[
\mathcal F(n)_j
=\mathcal F(n)/\mathfrak m^{j+1}\mathcal F(n)
\]
로 놓자. 그러면 $j\geq1$일 때
$\mathcal F(n)_{j-1}=\mathcal F(n)_j/\mathcal M_j(n)$이고
$\mathcal F(n)_0=\mathcal M_0(n)$이다.
$H^1(X,\mathcal F(n)_0)=0$이고, 코호몰로지 완전열에 의해 모든
$j\geq1$에 대해
\[
H^1(X,\mathcal F(n)_j)=H^1(X,\mathcal F(n)_{j-1}).
\]
따라서 모든 $j\geq0$에 대해
$H^1(X,\mathcal F(n)_j)=0$이다. 이제
(\hyperref[III.4.2.1-fr]{4.2.1})에 의해
$H^1(X,\mathcal F(n))=0$이고, 이로써 첫째 주장이 증명된다.

\oldpage[III]{146}

\noindent\textit{II)}
증명의 처음 표기로 돌아가자. 언제나 $Y=\operatorname{Spec}(A)$가
아핀이라고 가정할 수 있다. $f'$은 유한형이고 $\mathcal L'$은
$f'$에 관해 풍부하므로, 어떤 정수 $m>0$이 존재하여
$\mathcal L'^{\otimes m}$은 $f'$에 관해 매우 풍부하다
(II, 4.6.11). 필요하면 $\mathcal L$을
$\mathcal L^{\otimes m}$으로 바꾸어 $\mathcal L'$이 $f'$에 관해
매우 풍부한 경우로 한정할 수 있다. 이때
$\mathcal L|f^{-1}(U)$가 $f$에 관해 매우 풍부함을 보이면 된다.

$f'$은 고유이므로, 적당한 정수 $r>0$과 $Y'$-닫힌 몰입 사상
\[
j:X'\longrightarrow P=\mathbf P_{Y'}^r
\]
이 존재하여 $\mathcal L'$은 $j^*(\mathcal O_P(1))$과 동형이다
(II, 5.5.4, (ii)). 이 몰입 사상은 전사
$\mathcal O_{X'}$-준동형사상
\[
u:\mathcal O_{X'}^{r+1}\longrightarrow\mathcal L'
\]
에 표준적으로 대응한다 (II, 4.2.3). 후자는 $X'$ 위의
$\mathcal L'$의 $r+1$개 절단 $s'_i$ $(0\leq i\leq r)$로서 이
$\mathcal O_{X'}$-가군을 생성하는 것들의 데이터와 대응한다
(0\textsubscript{I}, 5.1.1). 정의에 의해 이 절단들은 $Y'$ 위의
$f'_*(\mathcal L')$의 절단들이기도 하다.
\[
f'_*(\mathcal L')
=f_*(\mathcal L)\otimes_{\mathcal O_Y}\mathcal O_{Y'}
\]
이다 (0\textsubscript{I}, 5.4.10). $Y$는 아핀이고
$\mathcal O_y$는 $A$의 소아이디얼 $\mathfrak j_y$에서의
국소환이므로
\[
s'_i=s''_i/t_i,
\]
여기서 $s''_i$들은 $Y$ 위의 $f_*(\mathcal L)$의 절단들이고
$t_i$들은 $\mathfrak j_y$에 속하지 않는 $A$의 원소들이다.
따라서 $Y$ 안의 $y$의 어떤 아핀 열린 근방 $V$와
$f_*(\mathcal L)|V$의 절단 $s_i$들이 존재하여 $s'_i=s_i/1$이다.
여기서 공간 $Y'$은 $V$에 포함됨을 상기한다 (I, 2.4.2).

$s_i$들은 $f^{-1}(V)$ 위의 $\mathcal L$의 절단들이므로
준동형사상
\[
v:\bigl(\mathcal O_X|f^{-1}(V)\bigr)^{r+1}
\longrightarrow\mathcal L|f^{-1}(V)
\]
을 정의한다. 가정에 의해 이는 $f^{-1}(y)$의 모든 점에서 전사이다.
$\operatorname{Coker}(v)$는 연접이므로
(0\textsubscript{I}, 5.3.4), 그 받침은 닫혀 있다
(0\textsubscript{I}, 5.2.2). 따라서 $f^{-1}(y)$의 어떤 열린 근방
$W\subset f^{-1}(V)$가 존재하여 $v$의 $W$에 대한 제한은 전사이다.
$f$는 닫힌 사상이므로 $W=f^{-1}(U)$라고 가정할 수 있다. 여기서
$U$는 $y$의 열린 근방이다. 이제 (II, 4.2.3)에서 결론이 따른다.
\end{proof}

\subsection{형식적 준스킴의 유한 사상}

\begin{proposition}[4.8.1]
\label{III.4.8.1-fr}
$\mathfrak Y$를 국소 뇌터 형식적 준스킴, $\mathcal K$를
$\mathfrak Y$의 정의 아이디얼, $f:\mathfrak X\to\mathfrak Y$를
형식적 준스킴의 사상이라 하자. 다음 조건들은 서로 동치이다.
\begin{enumerate}
\item[a)]
$\mathfrak X$는 국소 뇌터이고 $f$는 아딕 사상이며 (I, 10.12.1),
$\mathcal J=f^*(\mathcal K)\mathcal O_{\mathfrak X}$로 놓을 때
$f$에서 유도되는 사상
\[
f_0:(\mathfrak X,\mathcal O_{\mathfrak X}/\mathcal J)
\longrightarrow
(\mathfrak Y,\mathcal O_{\mathfrak Y}/\mathcal K)
\]
은 유한이다.
\item[b)]
$\mathfrak X$는 국소 뇌터이고, $X_0\to Y_0$가 유한인 어떤 아딕
$(Y_n)$-귀납계 $(X_n)$의 귀납극한이다.
\item[c)]
$\mathfrak Y$의 모든 점은 어떤 뇌터 아핀 형식적 열린 근방 $V$를
가지며, $f^{-1}(V)$는 아핀 형식적 열린집합이고
\[
\Gamma(f^{-1}(V),\mathcal O_{\mathfrak X})
\]
는 $\Gamma(V,\mathcal O_{\mathfrak Y})$ 위의 유한 생성 가군이다.
\end{enumerate}
\end{proposition}

\begin{proof}
(I, 10.12.3)에 의해 a)가 b)를 함의함은 즉시 따른다. b)가 c)를
함의함을 보이기 위해
\[
\mathfrak Y=\operatorname{Spf}(B),\qquad
\mathcal K=\mathfrak R^\Delta
\]
라고 가정할 수 있다. 여기서 $B$는 뇌터 아딕 환이고
$\mathfrak R$은 $B$의 정의 아이디얼이다. 가정에 의해 $X_0$은
아핀 스킴이고 그 환 $A_0$은 유한 생성
$B/\mathfrak R$-가군이다 (II, 6.1.3). (I, 5.1.9)에 의해 각
$X_n$은 아핀 스킴이다. 그 환을 $A_n$이라 하면, 가정 b)에 의해
$m\leq n$일 때
\[
A_m\simeq A_n/\mathfrak R^{m+1}A_n.
\]
따라서 $\mathfrak X$는 $\operatorname{Spf}(A)$와 동형이고, 여기서
$A=\varprojlim A_n$이다. 이제 (0\textsubscript{I}, 7.2.9)에서
결론이 따른다. 끝으로 c)가 a)를 함의함을 보이기 위해 다시

\oldpage[III]{147}

$\mathfrak Y=\operatorname{Spf}(B)$,
$\mathfrak X=\operatorname{Spf}(A)$이고 $A$가 유한
$B$-대수인 경우로 한정할 수 있다. 이때
$A/\mathfrak RA$는 유한 $B/\mathfrak R$-대수이므로
(I, 10.10.9)에 의해 a)의 조건들이 성립한다.
\end{proof}

\begin{definition}[4.8.2]
\label{III.4.8.2-fr}
(\hyperref[III.4.8.1-fr]{4.8.1})의 서로 동치인 성질 a), b), c)가
성립할 때 사상 $f$를 \emph{유한}이라 하고, 또는
$\mathfrak X$를 \emph{유한 $\mathfrak Y$-형식적 준스킴} 혹은
\emph{$\mathfrak Y$ 위에서 유한인 형식적 준스킴}이라 한다.
\end{definition}

\begin{proposition}[4.8.3]
\label{III.4.8.3-fr}
\begin{enumerate}
\item[(i)]
국소 뇌터 형식적 준스킴의 닫힌 몰입 사상은 유한 사상이다.
\item[(ii)]
국소 뇌터 형식적 준스킴의 두 유한 사상의 합성은 유한 사상이다.
\item[(iii)]
$\mathfrak X$, $\mathfrak Y$, $\mathfrak S$를 세 국소 뇌터
형식적 준스킴, $f:\mathfrak X\to\mathfrak S$를 유한 사상,
$g:\mathfrak Y\to\mathfrak S$를 사상이라 하자. 그러면 사상
\[
\mathfrak X\times_{\mathfrak S}\mathfrak Y\longrightarrow\mathfrak Y
\]
은 유한이다.
\item[(iv)]
$\mathfrak S$를 국소 뇌터 형식적 준스킴,
$\mathfrak X'$, $\mathfrak Y'$을 두 국소 뇌터 형식적 준스킴이라
하되
$\mathfrak X'\times_{\mathfrak S}\mathfrak Y'$이 국소 뇌터라고
하자. $\mathfrak X$, $\mathfrak Y$가 국소 뇌터
$\mathfrak S$-형식적 준스킴이고
$f:\mathfrak X\to\mathfrak X'$,
$g:\mathfrak Y\to\mathfrak Y'$이 두 유한 $\mathfrak S$-사상이면
$f\times_{\mathfrak S}g$는 유한 사상이다.
\item[(v)]
$f:\mathfrak X\to\mathfrak Y$,
$g:\mathfrak Y\to\mathfrak Z$를 국소 뇌터 형식적 준스킴의 두
사상이라 하고, $g$가 유한형이고 분리라고 하자. $g\circ f$가
유한 사상이면 $f$도 유한 사상이다.
\end{enumerate}
\end{proposition}

\begin{proof}
(i)은 자명하다. 나머지 주장들은
(\hyperref[III.4.8.1-fr]{4.8.1})의 판정법 a)를 이용하면 보통
준스킴의 사상에 대한 대응하는 명제들 (II, 6.1.5)로 곧바로
귀착된다. 세부사항은 (I, 10.13.5)를 본떠 독자에게 맡긴다.
\end{proof}

\begin{corollary}[4.8.4]
\label{III.4.8.4-fr}
(I, 10.9.9)의 가정 아래에서 $f$가 유한 사상이면, 그 완비화들로의
연장 $\widehat f$도 유한 사상이다.
\end{corollary}

\begin{corollary}[4.8.5]
\label{III.4.8.5-fr}
$\mathfrak X$가 $\mathfrak Y$ 위에서 유한인 형식적 준스킴이고
$f:\mathfrak X\to\mathfrak Y$가 구조 사상이면, 모든 열린집합
$U\subset\mathfrak Y$에 대해 $f^{-1}(U)$는 $U$ 위에서 유한이다.
\end{corollary}

\begin{proposition}[4.8.6]
\label{III.4.8.6-fr}
$f:\mathfrak X\to\mathfrak Y$가 국소 뇌터 형식적 준스킴의 유한
사상이면 $f_*(\mathcal O_{\mathfrak X})$는 연접
$\mathcal O_{\mathfrak Y}$-대수이다.
\end{proposition}

\begin{proof}
$f$를 사상 $f_n:X_n\to Y_n$들의 귀납계 $(f_n)$의 귀납극한으로
볼 수 있다. $f_n$들이 유한 사상이고
$f_*(\mathcal O_{\mathfrak X})$가
$(f_n)_*(\mathcal O_{X_n})$들의 사영극한과 동형임을 보이자.
그러면 (I, 10.10.5)에 의해 주장이 성립한다.
$\mathfrak Y=\operatorname{Spf}(B)$,
$\mathfrak X=\operatorname{Spf}(A)$인 경우로 한정하고 다음을
관찰하면 충분하다. $\mathfrak R$이 $B$의 정의 아이디얼이고
$A$가 유한 생성 $B$-가군이면
$A/\mathfrak R^{n+1}A$는
$B/\mathfrak R^{n+1}B$ 위의 유한 생성 가군이고, $A$는
$A/\mathfrak R^{n+1}A$들의 사영극한이다.
\end{proof}

역으로 다음이 성립한다.

\begin{proposition}[4.8.7]
\label{III.4.8.7-fr}
$\mathfrak Y$를 국소 뇌터 형식적 준스킴,
$\mathcal A$를 연접 $\mathcal O_{\mathfrak Y}$-대수라 하자.
$\mathfrak Y$-동형사상을 제외하고 유일한, $\mathfrak Y$ 위에서
유한인 형식적 준스킴 $\mathfrak X$가 존재하여
\[
f_*(\mathcal O_{\mathfrak X})=\mathcal A
\]
이다. 여기서 $f:\mathfrak X\to\mathfrak Y$는 구조 사상이다.
\end{proposition}

\begin{proof}
$\mathcal K$를 $\mathfrak Y$의 정의 아이디얼이라 하고
\[
Y_n=(\mathfrak Y,\mathcal O_{\mathfrak Y}/\mathcal K^{n+1}),
\qquad
\mathcal A_n=\mathcal A/\mathcal K^{n+1}\mathcal A
\]
로 놓자. $\mathcal A_n$은 분명히 유한
$\mathcal O_{Y_n}$-대수이고, 따라서 유한 $Y_n$-준스킴

\oldpage[III]{148}

$X_n=\operatorname{Spec}(\mathcal A_n)$을 정의한다 (II, 6.1.3).
$m\leq n$일 때 표준 전사 준동형사상
$h_{mn}:\mathcal A_n\to\mathcal A_m$은 사상
$u_{mn}:X_m\to X_n$을 정의하며, 이 사상은 구조 사상을 $f_n$이라
할 때 그림
\[
\xymatrix{
X_n\ar[d]_{f_n} & X_m\ar[l]_{u_{mn}}\ar[d]^{f_m}\\
Y_n & Y_m\ar[l]
}
\]
을 가환으로 만들고 $X_m$을
$X_n\times_{Y_n}Y_m$과 동일시한다 (II, 1.4.6).
귀납계 $(X_n)$의 귀납극한인 형식적 준스킴 $\mathfrak X$는 국소
뇌터이고, 귀납계 $(f_n)$의 귀납극한인 구조 사상
$f:\mathfrak X\to\mathfrak Y$는 유한이다
(\hyperref[III.4.8.1-fr]{4.8.1} 및 II, 10.12.3.1).
또한 (\hyperref[III.4.8.6-fr]{4.8.6})의 증명에서
$f_*(\mathcal O_{\mathfrak X})$가 $\mathcal A_n$들의 사영극한임을
보았으므로 이는 $\mathcal A$와 같다 (I, 10.10.6).
유일성 주장은 다음의 더 일반적인 결과의 귀결이다.
\end{proof}

\begin{proposition}[4.8.8]
\label{III.4.8.8-fr}
$\mathfrak Y$를 국소 뇌터 형식적 준스킴,
$\mathfrak X$, $\mathfrak X'$을 $\mathfrak Y$ 위에서 유한인 두
$\mathfrak Y$-형식적 준스킴이라 하고,
$f:\mathfrak X\to\mathfrak Y$,
$f':\mathfrak X'\to\mathfrak Y$을 구조 사상들이라 하자. 다음 두
집합 사이에는 표준 전단사 대응이 존재한다.
\[
\operatorname{Hom}_{\mathfrak Y}(\mathfrak X,\mathfrak X')
\quad\hbox{와}\quad
\operatorname{Hom}_{\mathcal O_{\mathfrak Y}}
\bigl(f'_*(\mathcal O_{\mathfrak X'}),
f_*(\mathcal O_{\mathfrak X})\bigr)
\footnotemark.
\]
\footnotetext{마지막 표현은
$\mathcal O_{\mathfrak Y}$-대수 준동형사상
$f'_*(\mathcal O_{\mathfrak X'})
\to f_*(\mathcal O_{\mathfrak X})$들의 집합을 뜻한다.}
\end{proposition}

\begin{proof}
이 대응 $h\mapsto\mathcal A(h)$의 정의는 (II, 1.1.2)에서와 같다.
그것이 전단사임을 보이기 위해
$\mathfrak Y=\operatorname{Spf}(B)$가 뇌터 아핀 형식적 스킴인
경우로 곧바로 귀착된다. 이때
\[
\mathfrak X=\operatorname{Spf}(A),\qquad
\mathfrak X'=\operatorname{Spf}(A'),
\]
여기서 $A$, $A'$은 두 유한 $B$-대수이고
\[
f_*(\mathcal O_{\mathfrak X})=A^\Delta,\qquad
f'_*(\mathcal O_{\mathfrak X'})=A'^\Delta.
\]
한편 $\mathfrak Y$-사상
$\mathfrak X\to\mathfrak X'$과 대수 준동형사상인
$B$-준동형사상 $A'\to A$ 사이에는 전단사 대응이 있다. 이
준동형사상들은 반드시 연속이다 (I, 10.2.2). 다른 한편
$B$-가군 준동형사상 $A'\to A$와
$\mathcal O_{\mathfrak Y}$-가군 준동형사상
$A'^\Delta\to A^\Delta$ 사이에도 전단사 대응이 있다
(I, 10.10.2.3). 이 두 대응에서 결론이 따른다.
\end{proof}

\begin{corollary}[4.8.9]
\label{III.4.8.9-fr}
(\hyperref[III.4.8.8-fr]{4.8.8})에서 정의된 표준 전단사 대응에서
닫힌 몰입 사상 $\mathfrak X\to\mathfrak X'$은 전사
$\mathcal O_{\mathfrak Y}$-대수 준동형사상
\[
f'_*(\mathcal O_{\mathfrak X'})
\longrightarrow f_*(\mathcal O_{\mathfrak X})
\]
과 대응한다.
\end{corollary}

\begin{proof}
문제는 다시 $\mathfrak Y$ 위에서 국소적이므로 국소 뇌터 형식적
준스킴의 닫힌 몰입 사상의 정의로 귀착된다 (I, 10.14.2).
\end{proof}

\begin{corollary}[4.8.10]
\label{III.4.8.10-fr}
(\hyperref[III.4.8.1-fr]{4.8.1})의 표기와 가정 아래에서, 아딕 사상
$f$가 닫힌 몰입 사상이기 위한 필요충분조건은 $f_0$가 보통
준스킴의 닫힌 몰입 사상인 것이다.
\end{corollary}

\begin{proof}
이는 (\hyperref[III.4.8.9-fr]{4.8.9})와 연접
$\mathcal O_{\mathfrak Y}$-가군의 준동형사상에 대한 전사 조건
(I, 10.11.5)에서 곧바로 따른다.
\end{proof}

\begin{proposition}[4.8.11]
\label{III.4.8.11-fr}
국소 뇌터 형식적 준스킴의 사상
$f:\mathfrak X\to\mathfrak Y$가 유한이기 위한 필요충분조건은

\oldpage[III]{149}

$f$가 고유이고 모든 $y\in\mathfrak Y$에 대해 올
$f^{-1}(y)$가 유한인 것이다.
\end{proposition}

\begin{proof}
정의 (3.4.1 및 4.8.2)에 의해 문제는
(\hyperref[III.4.8.1-fr]{4.8.1})의 표기에서 $f_0$에 대한 같은
명제로 곧바로 귀착되며, 이는 바로
(\hyperref[III.4.4.2-fr]{4.4.2})이다.
\end{proof}

\section{연접 대수적 층의 존재 정리}
\label{section:III.5-fr}

\subsection{정리의 진술}
\label{subsection:III.5.1-fr}

\begin{env}[5.1.1]
\label{III.5.1.1-fr}
$A$를 뇌터 아딕 환, $\mathfrak J$를 $A$의 정의 아이디얼이라 하자.
그러면 $A$는 $\mathfrak J$-아딕 위상에 대해 분리이고 완비이다.
$Y=\operatorname{Spec}(A)$이면 아핀 형식적 스킴
$\operatorname{Spf}(A)$는 닫힌부분집합
$Y'=\mathrm V(\mathfrak J)$을 따른 $Y$의 완비화
$\widehat Y$와 동일시된다 (I, 10.10.1).
$X$를 유한형 $Y$-준스킴, $f:X\to Y$를 구조 사상이라 하자.
$X'=f^{-1}(Y')$을 따른 $X$의 완비화, 또는 같은 것으로서
$\widehat Y$-형식적 준스킴 $X\times_Y\widehat Y$를
$\widehat X$로 나타내겠다. $f$의 완비화들로의 연장을
$\widehat f:\widehat X\to\widehat Y$로 나타내겠다. 끝으로 모든
연접 $\mathcal O_X$-가군 $\mathcal F$에 대해 그 완비화
$\mathcal F_{/X'}$를 $\widehat{\mathcal F}$로 나타내겠다. 이는
연접 $\mathcal O_{\widehat X}$-가군이다.
\end{env}

\begin{proposition}[5.1.2]
\label{III.5.1.2-fr}
(\hyperref[III.5.1.1-fr]{5.1.1})의 가정과 표기 아래에서,
$\mathcal F$를 그 받침이 $Y$ 위에서 고유인 연접
$\mathcal O_X$-가군이라 하자 (II, 5.4.10). 그러면 표준
준동형사상들 (4.1.4)
\[
\rho_i:\mathrm H^i(X,\mathcal F)\longrightarrow
\mathrm H^i(\widehat X,\widehat{\mathcal F})
\]
은 동형사상이다.
\end{proposition}

\begin{proof}
$\mathrm H^i(X,\mathcal F)$는 유한 생성 $A$-가군이고 (3.2.4),
따라서 자신의 $\mathfrak J$-준아딕 위상에 대한 분리 완비화와
같다 (0\textsubscript{I}, 7.3.6). 그러므로 이 명제는
(\hyperref[III.4.1.10-fr]{4.1.10})의 특수한 경우일 뿐이다.
\end{proof}

연접 $\mathcal O_X$-가군의 모든 완전열
\[
0\longrightarrow\mathcal F'\longrightarrow\mathcal F
\longrightarrow\mathcal F''\longrightarrow0
\]
에 대해 표준 동형사상 $\rho_i$들이 연결 준동형사상들과 가환함을
상기하자 (0, 12.1.6 및 I, 10.8.9).

\begin{corollary}[5.1.3]
\label{III.5.1.3-fr}
$\mathcal F$, $\mathcal G$를 두 연접 $\mathcal O_X$-가군이라 하되,
그 받침들의 교집합이 $Y$ 위에서 고유라고 하자. 그러면 모든
준동형사상 $u:\mathcal F\to\mathcal G$에 그 완비화
$\widehat u:\widehat{\mathcal F}\to\widehat{\mathcal G}$를
대응시키는 표준 준동형사상
\begin{equation}
\operatorname{Hom}_{\mathcal O_X}(\mathcal F,\mathcal G)
\longrightarrow
\operatorname{Hom}_{\mathcal O_{\widehat X}}
(\widehat{\mathcal F},\widehat{\mathcal G})
\tag{5.1.3.1}
\label{III.5.1.3.1.eq-fr}
\end{equation}
은 동형사상이다. 더 나아가 $f$가 닫힌 사상이면
$\widehat u$가 단사이기 위한 (각각 전사이기 위한) 필요충분조건은
$u$가 단사인 것이다 (각각 전사인 것이다).
\end{corollary}

\begin{proof}
첫째 주장은 (\ref{III.5.1.3.1.eq-fr})의 왼쪽 항이 유한 생성
$A$-가군이고 따라서 자신의 분리 완비화와 같다는 사실에 의한
(\hyperref[III.4.5.3-fr]{4.5.3})의 특수한 경우이다.
둘째 주장을 보이기 위해 (I, 10.8.14)에 의해
$\widehat u$가 단사이기 위한 (각각 전사이기 위한) 필요충분조건은
$X'$의 어떤 근방에서 $u$가 단사인 것이다 (각각 전사인 것이다).
따라서 결론은 다음 보조정리에서 따른다.
\end{proof}

\oldpage[III]{150}

\begin{lemma}[5.1.3.1]
\label{III.5.1.3.1-fr}
(\hyperref[III.5.1.1-fr]{5.1.1})의 가정 아래에서 $f$가 닫힌
사상이라고 더 가정하면, $X$ 안의 $X'$의 모든 근방은 $X$와 같다.
\end{lemma}

\begin{proof}
먼저 $f(X)=Y$인 경우로 귀착될 수 있다. 실제로 가정에 의해
$f(X)$는 $Y$의 닫힌부분집합 $Z$이다. 더 나아가 $f$를
$f_{\mathrm{red}}$로 바꾸어 (I, 6.3.4) $X$와 $Y$가 축약이라고
가정할 수 있다. 그러면 $Y$를 바탕공간이 $Z$인 $Y$의 축약 닫힌
부분준스킴으로 바꿀 수 있다 (I, 5.2.2). $A$의 모든 아이디얼은
닫혀 있으므로 $A$의 모든 몫환은 아딕이고 뇌터이기 때문이다.

이제 $f(X')=Y'$이다. $V$가 $X$ 안의 $X'$의 열린 근방이면
$f(X-V)$는 가정에 의해 $Y$에서 닫혀 있고 $Y'$과 만나지 않는다.
그런데 $\mathfrak J$가 $A$의 야코브슨 근기에 포함되므로
(I, 1.1.15 및 0\textsubscript{I}, 7.1.10), 이는 $X-V$가
공집합이지 않는 한 불가능하다. 따라서 $V=X$이다.
\end{proof}

받침이 $Y$ 위에서 고유인 연접 $\mathcal O_X$-가군들로 한정하면,
(\hyperref[III.5.1.3-fr]{5.1.3})은 범주의 언어로 다음과 같이
말할 수 있다. 함자
\[
\mathcal F\longmapsto\widehat{\mathcal F}
\]
는 그러한 $\mathcal O_X$-가군들의 범주에서 연접
$\mathcal O_{\widehat X}$-가군들의 범주로 가는 충만하고 충실한
함자이고, 따라서 전자의 범주와 후자의 어떤 충만한 부분범주
사이의 동치를 준다 (0, 8.1.6). 존재 정리는 $X$가 $Y$ 위에서
고유일 때 이 부분범주가 실제로 모든 연접
$\mathcal O_{\widehat X}$-가군들의 범주임을 보일 것이다. 정확히는
다음과 같다.

\begin{theorem}[5.1.4]
\label{III.5.1.4-fr}
$A$를 뇌터 아딕 환, $Y=\operatorname{Spec}(A)$,
$\mathfrak J$를 $A$의 정의 아이디얼,
$Y'=\mathrm V(\mathfrak J)$이라 하자.
$f:X\to Y$를 분리이고 유한형인 사상,
$X'=f^{-1}(Y')$이라 하자. $Y'$과 $X'$을 따른 $Y$와 $X$의
완비화들을
\[
\widehat Y=Y_{/Y'}=\operatorname{Spf}(A),\qquad
\widehat X=X_{/X'}
\]
로 나타내고, $f$의 완비화들로의 연장을
$\widehat f:\widehat X\to\widehat Y$라 하자. 그러면 함자
\[
\mathcal F\longmapsto
\mathcal F_{/X'}=\widehat{\mathcal F}
\]
는 받침이 $\operatorname{Spec}(A)$ 위에서 고유인 연접
$\mathcal O_X$-가군들의 범주와, 받침이
$\operatorname{Spf}(A)$ 위에서 고유인 연접
$\mathcal O_{\widehat X}$-가군들의 범주 사이의 동치이다.
\end{theorem}

다시 말해, (\hyperref[III.5.1.3-fr]{5.1.3})을 고려하면 다음과
같다.

\begin{corollary}[5.1.5]
\label{III.5.1.5-fr}
$\mathcal O_{\widehat X}$-가군이 받침이
$\operatorname{Spec}(A)$ 위에서 고유인 어떤 연접
$\mathcal O_X$-가군의 완비화와 동형이기 위한 필요충분조건은,
그 가군이 연접이고 그 받침이 $\operatorname{Spf}(A)$ 위에서
고유인 것이다.
\end{corollary}

가장 중요한 경우는 다음이다.

\begin{corollary}[5.1.6]
\label{III.5.1.6-fr}
$X$가 $Y=\operatorname{Spec}(A)$ 위에서 고유라고 가정하자.
그러면 함자
$\mathcal F\mapsto\widehat{\mathcal F}$는 연접
$\mathcal O_X$-가군들의 범주와 연접
$\mathcal O_{\widehat X}$-가군들의 범주 사이의 동치이다.
\end{corollary}

\begin{env}[5.1.7]
\label{III.5.1.7-fr}
\emph{보충.}
형식적 준스킴 위의 연접층의 특성화 (I, 10.11.3)를 고려하면,
(\hyperref[III.5.1.1-fr]{5.1.1})의 조건 아래에서 받침이
$\operatorname{Spec}(A)$ 위에서 고유인 연접
$\mathcal O_X$-가군을 주는 것은 다음 데이터를 주는 것과 동치임을
알 수 있다. 다음과 같이 놓자.
\[
Y_n=\operatorname{Spec}(A/\mathfrak J^{n+1}),\qquad
X_n=X\times_Y Y_n.
\]
각 $\mathcal F_n$이 연접 $\mathcal O_{X_n}$-가군인 사영계
$(\mathcal F_n)$을 주되, $m\leq n$일 때
\[
\mathcal F_m
=\mathcal F_n\otimes_{\mathcal O_{Y_n}}\mathcal O_{Y_m}
\quad\text{(또는 같은 것으로서 }\mathcal F_m
=\mathcal F_n/\mathfrak J^{m+1}\mathcal F_n\text{)}
\]
이고 $\mathcal F_0$의 받침이 $Y_0$ 위에서 고유인 $X_0$의
부분집합이어야 한다. (I, 10.11.4)를 이용하면 연접
$\mathcal O_X$-가군의 준동형사상들도 마찬가지로 연접
$\mathcal O_{X_n}$-가군들의 사영계 사이의 준동형사상으로
해석된다.

\oldpage[III]{151}

알려진 모든 응용의 경우 $A$는 실제로 뇌터 아딕 국소환이다.
따라서 $Y_n$들은 아르틴 국소환의 스펙트럼이고, 이 절과 앞 절들의
결과들은 뇌터 아딕 국소환 위의 대수기하학을 상당한 정도로 아르틴
국소환 위의 대수기하학으로 환원한다.
\end{env}

\begin{corollary}[5.1.8]
\label{III.5.1.8-fr}
(\hyperref[III.5.1.4-fr]{5.1.4})의 조건 아래에서 대응
\[
Z\longmapsto\widehat Z=Z_{/(Z\cap X')}
\]
은 $Y$ 위에서 고유인 $X$의 닫힌 부분준스킴들의 집합에서
$\widehat Y$ 위에서 고유인 $\widehat X$의 닫힌 형식적
부분준스킴들의 집합으로 가는 전단사이다.
\end{corollary}

\begin{proof}
실제로 $\widehat X$의 닫힌 형식적 부분준스킴은
\[
(T,(\mathcal O_{\widehat X}/\mathcal A)|T)
\]
의 꼴이다. 여기서 $\mathcal A$는 $\mathcal O_{\widehat X}$의 연접
아이디얼이다 (I, 10.14.2). $T$가 $\widehat Y$ 위에서 고유이면
(\hyperref[III.5.1.4-fr]{5.1.4})에 의해
$\mathcal O_{\widehat X}/\mathcal A$는
$\widehat{\mathcal F}$ 꼴의 $\mathcal O_{\widehat X}$-가군과
동형이다. 여기서 $\mathcal F$는 받침이 $Y$ 위에서 고유인 연접
$\mathcal O_X$-가군이다. 또한
(\hyperref[III.5.1.3-fr]{5.1.3})에 의해 표준 준동형사상
\[
\mathcal O_{\widehat X}\longrightarrow
\mathcal O_{\widehat X}/\mathcal A
\]
은 $\widehat u$의 꼴이고, 여기서
$u:\mathcal O_X\to\mathcal F$는 전사
$\mathcal O_X$-가군 준동형사상이다. 따라서
$\mathcal F=\mathcal O_X/\mathcal N$의 꼴이고, 여기서
$\mathcal N$은 $\mathcal O_X$의 연접 아이디얼이다. 또한
$\mathcal A=\widehat{\mathcal N}$이다 (I, 10.8.8).
그러므로 결론은 (I, 10.14.7)에서 따른다.
\end{proof}

\subsection{존재 정리의 증명: 사영인 경우와 준사영인 경우}
\label{subsection:III.5.2-fr}

\begin{env}[5.2.1]
\label{III.5.2.1-fr}
(\hyperref[III.5.1.4-fr]{5.1.4})의 조건 아래에서, 연접
$\mathcal O_{\widehat X}$-가군이 받침이 $Y$ 위에서 고유인 어떤
연접 $\mathcal O_X$-가군 $\mathcal F$의 완비화
$\widehat{\mathcal F}$와 동형일 때 잠정적으로 그 가군을
\emph{대수화 가능}이라고 하겠다.
\end{env}

\begin{lemma}[5.2.2]
\label{III.5.2.2-fr}
$\mathcal F'$, $\mathcal G'$을 두 대수화 가능한
$\mathcal O_{\widehat X}$-가군이라 하자. 모든 준동형사상
$u:\mathcal F'\to\mathcal G'$에 대해
$\operatorname{Ker}(u)$, $\operatorname{Im}(u)$,
$\operatorname{Coker}(u)$는 대수화 가능하다.
\end{lemma}

\begin{proof}
실제로
\[
\mathcal F'=\widehat{\mathcal F},\qquad
\mathcal G'=\widehat{\mathcal G}
\]
이고 $\mathcal F$, $\mathcal G$가 받침이 $Y$ 위에서 고유인 연접
$\mathcal O_X$-가군들이면, (\hyperref[III.5.1.3-fr]{5.1.3})에
의해 어떤 준동형사상 $v:\mathcal F\to\mathcal G$에 대해
$u=\widehat v$이다. 함자
$\mathcal F\mapsto\widehat{\mathcal F}$의 완전성에 의해
$\operatorname{Ker}(\widehat v)$는
$\widehat{\operatorname{Ker}(v)}$와 동형이다. 또한
$\operatorname{Ker}(v)$의 받침은 $\mathcal F$의 받침에
포함되므로 $\operatorname{Ker}(u)$는 대수화 가능하다.
$\operatorname{Im}(u)$와 $\operatorname{Coker}(u)$에도 같은
논증을 적용한다.
\end{proof}

\begin{proposition}[5.2.3]
\label{III.5.2.3-fr}
$A$를 뇌터 아딕 환, $\mathfrak J$를 $A$의 정의 아이디얼,
$\mathfrak Y=\operatorname{Spf}(A)$,
$f:\mathfrak X\to\mathfrak Y$를 형식적 준스킴의 고유 사상이라
하자. 다음과 같이 놓자.
\[
Y_k=\operatorname{Spec}(A/\mathfrak J^{k+1}),\qquad
X_k=\mathfrak X\times_{\mathfrak Y}Y_k,
\]
그리고 모든 $\mathcal O_{\mathfrak X}$-가군 $\mathcal F$에 대해
\[
\mathcal F_k
=\mathcal F\otimes_{\mathcal O_{\mathfrak X}}\mathcal O_{X_k}
=\mathcal F/\mathfrak J^{k+1}\mathcal F.
\]
$\mathcal L$을 가역 $\mathcal O_{\mathfrak X}$-가군이라 하고
$\mathcal L_0=\mathcal L/\mathfrak J\mathcal L$이 풍부한
$\mathcal O_{X_0}$-가군이라고 가정하자. 모든
$\mathcal O_{\mathfrak X}$-가군 $\mathcal F$와 모든 정수 $n$에
대해
\[
\mathcal F(n)=\mathcal F\otimes\mathcal L^{\otimes n}
\]
으로 놓자. 그러면 모든 연접 $\mathcal O_{\mathfrak X}$-가군
$\mathcal F$에 대해 어떤 정수 $n_0$가 존재하여, 모든
$n\geq n_0$에 대해 다음 성질들이 성립한다.
\begin{enumerate}
\item[(i)]
모든 $k\geq0$에 대해 표준 준동형사상
\[
\mathrm H^0(\mathfrak X,\mathcal F(n))
\longrightarrow
\mathrm H^0(\mathfrak X,\mathcal F_k(n))
\]
은 전사이다.
\item[(ii)]
모든 $q>0$에 대해
$\mathrm H^q(\mathfrak X,\mathcal F(n))=0$이다.
\end{enumerate}
\end{proposition}

\begin{proof}
$\mathfrak X$와 $X_0$의 바탕공간들은 같음이 알려져 있다. 층
\[
\mathcal M_k
=\mathfrak J^k\mathcal F/\mathfrak J^{k+1}\mathcal F
\]
는 $\mathfrak J$에 의해 소멸되므로 연접
$\mathcal O_{X_0}$-가군으로 볼 수 있다
(0\textsubscript{I}, 5.3.10). 또한
\[
\mathcal M_k(n)
=\mathcal M_k\otimes_{\mathcal O_{X_0}}\mathcal L_0^{\otimes n}
\]
으로 놓으면 즉시
\[
\mathcal M_k(n)
=\mathfrak J^k\mathcal F(n)/
\mathfrak J^{k+1}\mathcal F(n)
\]
임을 알 수 있다. $\mathcal L_0$은
$f_0:X_0\to Y_0$에 관해 풍부하고

\oldpage[III]{152}

$f_0$는 고유이므로 $f_0$는 사영 사상이다 (II, 5.5.4).
$A$의 $\mathfrak J$-아딕 여과에 결부된 등급 $A$-대수
\[
S=\bigoplus_{k\geq0}
\mathfrak J^k/\mathfrak J^{k+1}
\]
를 생각하자. $A$가 뇌터환이므로 $S$는 유한형이다.
\[
\mathcal S'=f_0^*(\widetilde S),\qquad
\mathcal M=\bigoplus_{k\geq0}\mathcal M_k
\]
로 놓으면 $\mathcal S'$은 준연접 유한형
$\mathcal O_{X_0}$-대수이고, $\mathcal M$은 준연접 유한형 등급
$\mathcal S'$-가군이다. 실제로 $\mathcal F_0$가 연접이고
$\mathcal S'$-가군 $\mathcal M$을 생성한다.

따라서 정리 (\hyperref[III.2.4.1-fr]{2.4.1, (ii)})를 적용할 수
있다. 어떤 $n_0$가 존재하여 $n\geq n_0$이면 모든 $k$와 모든
$q>0$에 대해
\begin{equation}
\mathrm H^q(X_0,\mathcal M_k(n))=0.
\tag{5.2.3.1}
\label{III.5.2.3.1-fr}
\end{equation}
$\mathcal M_k(n)$을 이번에는 $\mathcal O_{\mathfrak X}$-가군으로
보아도 $q>0$, $n\geq n_0$일 때
\[
\mathrm H^q(\mathfrak X,\mathcal M_k(n))=0
\]
이다. 완전열
\[
0\longrightarrow
\frac{\mathfrak J^k\mathcal F(n)}
{\mathfrak J^{k+1}\mathcal F(n)}
\longrightarrow
\frac{\mathfrak J^h\mathcal F(n)}
{\mathfrak J^{k+1}\mathcal F(n)}
\longrightarrow
\frac{\mathfrak J^h\mathcal F(n)}
{\mathfrak J^k\mathcal F(n)}
\longrightarrow0
\]
에 코호몰로지 완전열을 적용하자. 먼저 $0\leq h<k$,
$n\geq n_0$, $q>0$일 때 $k-h$에 관한 귀납법으로
\begin{equation}
\mathrm H^q\left(\mathfrak X,
\frac{\mathfrak J^h\mathcal F(n)}
{\mathfrak J^k\mathcal F(n)}\right)=0
\tag{5.2.3.2}
\label{III.5.2.3.2-fr}
\end{equation}
을 얻는다. 특히 $h=0$이면
\begin{equation}
\mathrm H^q(\mathfrak X,\mathcal F_k(n))=0
\qquad
(n\geq n_0,\ k\geq0,\ q>0).
\tag{5.2.3.3}
\label{III.5.2.3.3-fr}
\end{equation}
$h=0$일 때 코호몰로지 완전열의 다른 부분은 완전열
\begin{equation}
\mathrm H^0(\mathfrak X,\mathcal F_{k+1}(n))
\longrightarrow
\mathrm H^0(\mathfrak X,\mathcal F_k(n))
\longrightarrow
\mathrm H^1\left(\mathfrak X,
\frac{\mathfrak J^k\mathcal F(n)}
{\mathfrak J^{k+1}\mathcal F(n)}\right)=0
\tag{5.2.3.4}
\label{III.5.2.3.4-fr}
\end{equation}
을 준다. 따라서 $h\leq k$일 때 표준 사상
\begin{equation}
\mathrm H^0(\mathfrak X,\mathcal F_k(n))
\longrightarrow
\mathrm H^0(\mathfrak X,\mathcal F_h(n))
\tag{5.2.3.5}
\label{III.5.2.3.5-fr}
\end{equation}
은 전사이다.

그러므로 모든 $q$에 대해 사영계
\[
\bigl(\mathrm H^q(\mathfrak X,\mathcal F_k(n))\bigr)_{k\geq0}
\]
는 $n\geq n_0$일 때 조건 (ML)을 만족한다. 한편
$\mathfrak X$의 모든 아핀 형식적 열린집합 $U$는 각 $X_k$에서도
아핀 열린집합이므로 (I, 10.5.2), 모든 $q>0$에 대해
$\mathrm H^q(U,\mathcal F_k(n))=0$이고 (I, 1.3.1),
$h\leq k$일 때
\[
\mathrm H^0(U,\mathcal F_k(n))
\longrightarrow\mathrm H^0(U,\mathcal F_h(n))
\]
은 전사이다 (I, 1.3.9). 따라서 (0, 13.3.1)의 적용 조건들이
충족되며, $n\geq n_0$일 때 다음을 얻는다.
\begin{enumerate}
\item[$1^\circ$]
모든 $q>0$에 대해
\[
\mathrm H^q(\mathfrak X,\mathcal F(n))
\longrightarrow
\varprojlim_k\mathrm H^q(\mathfrak X,\mathcal F_k(n))
\]
은 전단사이다. 따라서 (\ref{III.5.2.3.3-fr})에 의해
$\mathrm H^q(\mathfrak X,\mathcal F(n))=0$이다.
\item[$2^\circ$]
준동형사상
\[
\mathrm H^0(\mathfrak X,\mathcal F(n))
\longrightarrow
\varprojlim_k\mathrm H^0(\mathfrak X,\mathcal F_k(n))
\]
은 전단사이다. 또한 준동형사상
(\ref{III.5.2.3.5-fr})들이 전사이므로 각 준동형사상
\[
\varprojlim_k\mathrm H^0(\mathfrak X,\mathcal F_k(n))
\longrightarrow
\mathrm H^0(\mathfrak X,\mathcal F_h(n))
\]
도 전사이다.
\end{enumerate}
이로써 증명이 끝난다.
\end{proof}

\begin{corollary}[5.2.4]
\label{III.5.2.4-fr}
(\hyperref[III.5.2.3-fr]{5.2.3})의 가정 아래에서 모든 연접
$\mathcal O_{\mathfrak X}$-가군 $\mathcal F$에 대해 어떤 정수
$N$이 존재하여, $n\geq N$이면 $\mathcal F(n)$은

\oldpage[III]{153}

$\mathfrak X$ 위의 자신의 절단들로 생성된다. 다시 말해
$\mathcal F$는
\[
(\mathcal O_{\mathfrak X}(-n))^k
\]
꼴의 $\mathcal O_{\mathfrak X}$-가군의 몫과 동형이다.
\end{corollary}

\begin{proof}
$X_0$가 뇌터이므로, $\mathcal L_0$에 대한 가정과 (II, 4.5.5)에
의해 어떤 $n_0$가 존재하여 $n\geq n_0$이면
$\mathcal F_0(n)$은 $\mathfrak X$ 위의 자신의 절단들로 생성된다.
또한 $n_0$를 충분히 크게 잡아 $n\geq n_0$일 때 준동형사상
\[
\Gamma(\mathfrak X,\mathcal F(n))
\longrightarrow
\Gamma(\mathfrak X,\mathcal F_0(n))
\]
이 전사가 되게 할 수 있다
(\hyperref[III.5.2.3-fr]{5.2.3}). 따라서 유한 개의 절단
$s_i\in\Gamma(\mathfrak X,\mathcal F(n))$가 존재하여, 그
$\Gamma(\mathfrak X,\mathcal F_0(n))$ 안의 상들이
$\mathcal F_0(n)$을 생성한다 (0\textsubscript{I}, 5.2.3).
$\mathfrak J$는 $\mathfrak X$의 모든 점에서 국소환의 극대
아이디얼에 포함되므로, 이 국소환들에 나카야마 보조정리를 적용하면
$s_i$들이 $\mathcal F(n)$을 생성함을 얻는다
(0\textsubscript{I}, 5.1.1).
\end{proof}

\begin{env}[5.2.5]
\label{III.5.2.5-fr}
\emph{존재 정리의 증명: 사영인 경우.}
(\hyperref[III.5.1.4-fr]{5.1.4})의 표기를 사용하고 $f$가
사영이라고 가정하자. 그러면 풍부한 $\mathcal O_X$-가군
$\mathcal L$이 존재한다 (I, 5.5.4). 정의에 의해
\[
X_n=\widehat X\times_{\widehat Y}Y_n
\]
은 $X$의 닫힌 부분준스킴
$X\times_Y Y_n=f^{-1}(Y_n)$과 같다.
$\mathcal L'=\mathcal L\otimes_{\mathcal O_X}
\mathcal O_{\widehat X}$이 $\mathcal L$의 완비화이면
$\mathcal L'_0=\mathcal L/\mathfrak J\mathcal L$이며, 이를
$\mathcal O_{X_0}$-가군으로 본다. $\mathcal L'_0$은 풍부함이
알려져 있다 (II, 4.6.13, (i bis)).

따라서 $\mathcal L'$과 임의의 연접
$\mathcal O_{\widehat X}$-가군 $\mathcal F$에 따름정리
(\hyperref[III.5.2.4-fr]{5.2.4})를 적용할 수 있다. 적당한 정수
$n>0$, $k>0$에 대해 $\mathcal F$는
\[
\mathcal G=(\mathcal L'^{\otimes(-n)})^k
\]
의 몫과 동형이다. 그런데 $\mathcal G$는 분명히
$(\mathcal L^{\otimes(-n)})^k$의 완비화이므로 (I, 10.8.10)
대수화 가능하다. 이제 표준 전사 준동형사상
$u:\mathcal G\to\mathcal F$를 생각하고
\[
\mathcal H=\operatorname{Ker}(u)
\]
로 놓자. 이는 연접 $\mathcal O_{\widehat X}$-가군이다
(0\textsubscript{I}, 5.3.4). 같은 방식으로 어떤 대수화 가능한
$\mathcal O_{\widehat X}$-가군 $\mathcal K$와 준동형사상
$v:\mathcal K\to\mathcal G$가 존재하여
$\mathcal H=\operatorname{Im}(v)$임을 알 수 있다. 그러면
$\mathcal F=\operatorname{Coker}(v)$이고,
(\hyperref[III.5.2.2-fr]{5.2.2})에 의해 $\mathcal F$는 대수화
가능하다.
\end{env}

\begin{env}[5.2.6]
\label{III.5.2.6-fr}
\emph{존재 정리의 증명: 준사영인 경우.}
계속해서 (\hyperref[III.5.1.4-fr]{5.1.4})의 표기를 사용하고,
이제 $f$가 준사영이라고 가정하자. 어떤 사영 사상
$g:Z\to Y$가 존재하여 $X$는 $Z$의 어떤 열린집합 위에 유도된
$Y$-준스킴과 동일시된다 (II, 5.3.2).
$Z'=g^{-1}(Y')$로 놓으면 $X'=X\cap Z'$이다. 따라서 완비화
$\widehat X=X_{/X'}$는 완비화
$\widehat Z=Z_{/Z'}$의 열린집합 $X\cap Z'$ 위에 유도된 형식적
준스킴과 동일시된다 (I, 10.8.5).

$\mathcal F'$을 연접 $\mathcal O_{\widehat X}$-가군이라 하고 그
받침 $T'$이 $\widehat Y$ 위에서 고유라고 하자. 이는 정의에 의해,
바탕공간이 $T'\subset X'$이고 $f$의 제한 $T'\to Y'$이 고유인
$X'$의 닫힌 부분준스킴이 존재한다는 뜻이다. 따라서 $T'$은
$Y$ 위에서 고유이고, 그러므로 $Z'$에서 닫혀 있다
(II, 5.4.10).

이로부터 $\mathcal F'$은 다음 연접
$\mathcal O_{\widehat Z}$-가군 $\mathcal G'$이
$\widehat X$ 위에 유도하는 층임이 따른다. 즉
$\widehat Z$의 열린집합 $\widehat X$ 위에 정의된
$\mathcal F'$과 열린집합 $\widehat Z-T'$ 위의 영층을 붙여
$\mathcal G'$을 얻는다. 이 두 층은 교집합
$\widehat X-T'$에서 일치한다. $\mathcal G'$이 연접임은
명백하다. (\hyperref[III.5.2.5-fr]{5.2.5})에 의해 어떤 연접
$\mathcal O_Z$-가군 $\mathcal G$가 존재하여
\[
\mathcal G'=\widehat{\mathcal G}.
\]
$T$를 $\mathcal G$의 받침이라 하면
$T'=T\cap Z'$이다 (I, 10.8.12). $h$를 바탕공간이 $T$인 $Z$의
축약 닫힌 부분준스킴에 대한 $g$의 제한이라 하면
\[
T'=h^{-1}(Y')=T\cap g^{-1}(Y').
\]
따라서 $X\cap T$는 $T'$을 포함하는 $T$의 열린집합이다.

\oldpage[III]{154}

$h$는 고유이고 (II, 5.4.2), 따라서 닫힌 사상이므로
(\hyperref[III.5.1.3.1-fr]{5.1.3.1})에 의해
$X\cap T=T$이다. 다시 말해 $T\subset X$이다. 또한 $T$는
$Z$에서 닫혀 있으므로 $Y$ 위에서 고유이다.
$\mathcal F$를 $\mathcal G$가 $X$ 위에 유도하는 층이라 하자.
그 완비화 $\widehat{\mathcal F}$은
$\widehat{\mathcal G}$가 $\widehat X$ 위에 유도하는 층이므로
(I, 10.8.4) $\mathcal F'$과 같다. 이로써 증명이 끝난다.
\end{env}

\subsection{존재 정리의 증명: 일반적인 경우}
\label{subsection:III.5.3-fr}

\begin{lemma}[5.3.1]
\label{III.5.3.1-fr}
(\hyperref[III.5.1.4-fr]{5.1.4})의 조건 아래에서,
\[
0\to\mathcal H\to\mathcal F\to\mathcal G\to0
\]
이 연접 $\mathcal O_{\widehat X}$-가군들의 완전열이고
$\mathcal G$와 $\mathcal H$가 대수화 가능하면, $\mathcal F$도
대수화 가능하다.
\end{lemma}

\begin{proof}
실제로 $\mathcal G=\widehat{\mathcal B}$,
$\mathcal H=\widehat{\mathcal C}$라고 하자. 여기서 $\mathcal B$와
$\mathcal C$는 받침이 $Y$ 위에서 고유인 연접
$\mathcal O_X$-가군들이다. $\mathcal F$는 표준적으로 $A$-가군
\[
\operatorname{Ext}^1_{\mathcal O_{\widehat X}}
(\widehat X;\widehat{\mathcal B},\widehat{\mathcal C})
\]
의 한 원소를 정한다 ($0$, 12.3.2). 가정에 의해 이 $A$-가군은
표준적으로
\[
\operatorname{Ext}^1_{\mathcal O_X}(X;\mathcal B,\mathcal C)
\]
와 동형이다 (4.5.2). 따라서 연접 $\mathcal O_X$-가군들의 완전열
\[
0\to\mathcal C\to\mathcal A\to\mathcal B\to0
\]
이 존재하여, $\mathcal A$에 대응하는
$\operatorname{Ext}^1_{\mathcal O_X}(X;\mathcal B,\mathcal C)$의
원소의 표준상은 $\mathcal F$에 대응하는
$\operatorname{Ext}^1_{\mathcal O_{\widehat X}}
(\widehat X;\widehat{\mathcal B},\widehat{\mathcal C})$의 원소이다.
그런데 정의에 의해, 그리고 (I, 10.8.8, (iii))를 고려하면, 이는
$\mathcal F$가 $\widehat{\mathcal A}$와 동형이라는 뜻이다. 또한
$\operatorname{Supp}(\mathcal A)$는
$\operatorname{Supp}(\mathcal B)$와
$\operatorname{Supp}(\mathcal C)$의 합집합에 포함되므로 $Y$ 위에서
고유이다. 이로써 보조정리가 증명된다.
\end{proof}

\begin{corollary}[5.3.2]
\label{III.5.3.2-fr}
(\hyperref[III.5.1.1-fr]{5.1.1})의 조건 아래에서,
$u:\mathcal F\to\mathcal G$를 연접
$\mathcal O_{\widehat X}$-가군들의 준동형사상이라 하자.
$\mathcal G$, $\operatorname{Ker}(u)$ 및
$\operatorname{Coker}(u)$가 대수화 가능하면 $\mathcal F$도
대수화 가능하다.
\end{corollary}

\begin{proof}
준동형사상 $\mathcal G\to\operatorname{Coker}(u)$에 보조정리
(\hyperref[III.5.2.2-fr]{5.2.2})를 적용하면
$\operatorname{Im}(u)$가 대수화 가능함을 알 수 있다. 이제 완전열
\[
0\to\operatorname{Ker}(u)\to\mathcal F
\to\operatorname{Im}(u)\to0
\]
에 보조정리 (\hyperref[III.5.3.1-fr]{5.3.1})를 적용하면 충분하다.
\end{proof}

\begin{lemma}[5.3.3]
\label{III.5.3.3-fr}
(\hyperref[III.5.1.1-fr]{5.1.1})의 조건 아래에서,
$h:Z\to Y$를 유한형 사상, $\widehat Z$를
$Z'=h^{-1}(Y')$을 따른 $Z$의 완비화, $g:Z\to X$를 고유
$Y$-사상, $\widehat g:\widehat Z\to\widehat X$를 완비화들로
연장한 사상이라 하자. 모든 대수화 가능한
$\mathcal O_{\widehat Z}$-가군 $\mathcal F'$에 대해
$\widehat g_*(\mathcal F')$은 대수화 가능한
$\mathcal O_{\widehat X}$-가군이다.
\end{lemma}

\begin{proof}
실제로 $\mathcal F'=\widehat{\mathcal F}$인 연접
$\mathcal O_Z$-가군 $\mathcal F$가 있으면, 제1 비교 정리
(4.1.5)에 의해 $\widehat g_*(\widehat{\mathcal F})$는
$g_*(\mathcal F)$의 완비화와 동형이다.
\end{proof}

\begin{lemma}[5.3.4]
\label{III.5.3.4-fr}
$X$를 (통상적인) 뇌터 스킴, $X'$을 $X$의 닫힌 부분,
$f:Z\to X$를 고유 사상, $Z'=f^{-1}(X')$라 하고,
\[
\widehat X=X_{/X'},\qquad \widehat Z=Z_{/Z'}
\]
라 하자. 또한 $\widehat f:\widehat Z\to\widehat X$를 $f$의
완비화들로의 연장이라 하자. $\mathcal M$을 다음 성질을 갖는
$\mathcal O_X$의 연접 아이디얼이라 하자. 즉
\[
U=X-\operatorname{Supp}(\mathcal O_X/\mathcal M)
\]
이면 $f$의 제한 $f^{-1}(U)\to U$는 동형사상이다. 그러면 모든
연접 $\mathcal O_{\widehat X}$-가군 $\mathcal F$에 대해 어떤 정수
$n>0$가 존재하여, 표준 준동형사상
\[
\rho_{\mathcal F}:\mathcal F\longrightarrow
\widehat f_*\bigl(\widehat f^{,*}(\mathcal F)\bigr)
\]
($0_{\mathrm I}$, 4.4.3)의 핵과 여핵은
$\widehat{\mathcal M}^{,n}$에 의해 소멸된다.
\end{lemma}

\begin{proof}
$B$가 뇌터 환인 $X=\operatorname{Spec}(B)$의 경우로 한정할 수
있다. 이때 어떤 $B$의 아이디얼 $\mathfrak K$에 대해
$X'=\mathrm V(\mathfrak K)$이다. 먼저 $B$가 뇌터 아딕 환이고
$\mathfrak K$가 $B$의 정의 아이디얼인 경우로 환원할 수 있음을
보이자. 실제로 $B_1$을 $\mathfrak K$-선아딕 위상에 대한 $B$의
분리 완비환이라 하자. $\mathfrak K_1=\mathfrak K B_1$이면
$B_1$은 $\mathfrak K_1$을 정의 아이디얼로 갖는

\oldpage[III]{155}

뇌터 아딕 환이다. $X_1=\operatorname{Spec}(B_1)$로 놓고,
$h:X_1\to X$를 표준 준동형사상 $B\to B_1$에 대응하는 사상이라
하자. $X'_1=h^{-1}(X')$이면
$X'_1=\mathrm V(\mathfrak K_1)$이다. 마지막으로
\[
Z_1=Z\times_X X_1=Z_{(X_1)},\qquad
f_1=f_{(X_1)}:Z_1\to X_1
\]
로 놓자. $f_1$은 고유 사상이다 (II, 5.4.2). $\widehat X_1$을
$X'_1$을 따른 $X_1$의 완비화,
\[
\widehat Z_1=Z_1\times_{X_1}\widehat X_1
\]
을 $Z'_1=f_1^{-1}(X'_1)$을 따른 $Z_1$의 완비화,
$\widehat f_1$을 $f_1$의 완비화들로의 연장이라 하자.

$h$의 완비화들로의 연장
$\widehat h:\widehat X_1\to\widehat X$는 $B_1$의 항등사상에
대응하는 동형사상임이 즉시 따른다 (I, 10.9.1). 따라서 이에
대응하는 준동형사상 $\widehat Z_1\to\widehat Z$도
동형사상이고, 이 동형사상들은 $\widehat f_1$과
$\widehat f$를 동일시한다. 한편
$\mathcal M_1=h^*(\mathcal M)$은 $\mathcal O_{X_1}$의 연접
아이디얼이고
\[
\operatorname{Supp}(\mathcal O_{X_1}/\mathcal M_1)
=h^{-1}(\operatorname{Supp}(\mathcal O_X/\mathcal M))
\]
이다 (I, 9.1.13). 따라서
\[
U_1=X_1-\operatorname{Supp}(\mathcal O_{X_1}/\mathcal M_1)
\]
로 놓으면 $U_1=h^{-1}(U)$이고, 곧바로 $f_1$의 제한
$f_1^{-1}(U_1)\to U_1$이 동형사상임을 알 수 있다
(I, 3.2.7). 또한 완비화 $\widehat{\mathcal M}$과
$\widehat{\mathcal M}_1$은 $\widehat h$에 의해 동일시된다
(I, 10.9.5). 그러므로 $X_1$, $X'_1$, $f_1$, $\mathcal M_1$은
(\hyperref[III.5.3.4-fr]{5.3.4})의 모든 가정을 만족한다. 이제부터
$B$가 뇌터 아딕 환이고 $\mathfrak K$가 $B$의 정의 아이디얼이라고
가정할 수 있다.

이 경우 $\widehat X=\operatorname{Spf}(B)$이고,
$N$이 유한형 $B$-가군인 $\mathcal F=N^\Delta$로 쓸 수 있다.
따라서 $\mathcal G$를 연접 $\mathcal O_X$-가군
$\widetilde N$이라 하면
$\mathcal F=\widehat{\mathcal G}$이다 (I, 10.10.5). 그러므로
\[
\widehat f^{,*}(\mathcal F)=\widehat{f^*(\mathcal G)}
\]
이다 (I, 10.9.5). 또한 제1 비교 정리 (4.1.5)에 의해
\[
\widehat f_*\bigl(\widehat{f^*(\mathcal G)}\bigr)
\simeq\widehat{f_*\bigl(f^*(\mathcal G)\bigr)},
\]
이고, (\hyperref[III.5.1.3-fr]{5.1.3})에 의해 표준 준동형사상
$\rho_{\mathcal F}$는 바로 $\widehat{\rho}_{\mathcal G}$이다.
그런데
\[
\rho_{\mathcal G}:\mathcal G\longrightarrow
f_*\bigl(f^*(\mathcal G)\bigr)
\]
의 핵 $\mathcal P$와 여핵 $\mathcal R$은 연접
$\mathcal O_X$-가군들이고, 가정에 의해 이들의 $U$에 대한 제한은
명백히 영이다. 따라서 어떤 정수 $n>0$가 존재하여
\[
\mathcal M^n\mathcal P=\mathcal M^n\mathcal R=0
\]
이다 (I, 9.3.4). 이로부터
\[
\widehat{\mathcal M}^{,n}\widehat{\mathcal P}=
\widehat{\mathcal M}^{,n}\widehat{\mathcal R}=0
\]
이 따른다 (I, 10.8.8 및 10.8.10).
\end{proof}

\begin{env}[5.3.5]
\label{III.5.3.5-fr}
\emph{존재 정리의 증명의 끝.}
(\hyperref[III.5.1.4-fr]{5.1.4})의 가정 아래에서 뇌터 귀납법의
원리 ($0_{\mathrm I}$, 2.2.2)를 사용하자. 즉 바탕공간이 $X$와
다른 $X$의 모든 닫힌 부분준스킴 $T$에 대해 정리가 참이라고
가정한다. 여기서 $\widehat T$는 물론 $T\cap X'$을 따른 $T$의
완비화이다. $X$는 공집합이 아니라고 가정할 수 있다. $f$가
분리이고 유한형이므로 초우 보조정리 (II, 5.6.1)를 적용할 수
있다. 따라서 $Y$-스킴 $Z$와 $Y$-사상 $g:Z\to X$가 존재하여,
구조 사상 $h:Z\to Y$는 준사영이고 $g$는 사영이며 전사이다.
더 나아가 $X$의 공집합이 아닌 열린집합 $U$가 존재하여 제한
$g^{-1}(U)\to U$는 동형사상이다.

$\mathcal M$을 바탕공간이 $X-U$인 닫힌 부분준스킴을 정의하는
$\mathcal O_X$의 연접 아이디얼이라 하고 (I, 5.2.2),
$\mathcal F$를 받침 $E$가 $Y$ 위에서 고유인 연접
$\mathcal O_{\widehat X}$-가군이라 하자. $\widehat Z$를
$h^{-1}(Y')$을 따른 $Z$의 완비화,
$\widehat g:\widehat Z\to\widehat X$를 $g$의 완비화들로의
연장이라 하자. 그러면 $\widehat g^{,*}(\mathcal F)$는 받침이
$g^{-1}(E)$에 포함되는 연접 $\mathcal O_{\widehat Z}$-가군이다.
$g$는 사영이고 따라서 고유이므로, 이 받침은 $Y$ 위에서 고유이다
(II, 5.4.6). $h$가 준사영이므로
$\widehat g^{,*}(\mathcal F)$는
(\hyperref[III.5.2.6-fr]{5.2.6})에 의해 대수화 가능하다. 따라서

\oldpage[III]{156}

$g$가 고유임과 (\hyperref[III.5.3.3-fr]{5.3.3})에 의해
\[
\widehat g_*\bigl(\widehat g^{,*}(\mathcal F)\bigr)
\]
은 대수화 가능한 $\mathcal O_{\widehat X}$-가군이다. 이제
$\mathcal F$와 $g$에 (\hyperref[III.5.3.4-fr]{5.3.4})를 적용할
수 있다. 준동형사상
\[
\rho_{\mathcal F}:\mathcal F\to
\widehat g_*\bigl(\widehat g^{,*}(\mathcal F)\bigr)
\]
의 핵 $\mathcal P$와 여핵 $\mathcal R$은 어떤 거듭제곱
$\widehat{\mathcal M}^{,n}$에 의해 소멸된다. $T$를
$\mathcal M^n$으로 정의되는 $X$의 닫힌 부분준스킴이라 하자.
그 바탕공간은 $X-U$이다. 또한 $j:T\to X$를 표준 단사라 하자.
그러면 완비화들로의 연장 $\widehat j:\widehat T\to\widehat X$는
표준 단사이다 (I, 10.14.7). 따라서
\[
\mathcal P=\widehat j_*\bigl(\widehat j^{,*}(\mathcal P)\bigr),
\qquad
\mathcal R=\widehat j_*\bigl(\widehat j^{,*}(\mathcal R)\bigr).
\]
$U$는 공집합이 아니므로 귀납가정에 의해
$\widehat j^{,*}(\mathcal P)$와
$\widehat j^{,*}(\mathcal R)$은 대수화 가능한
$\mathcal O_{\widehat T}$-가군들이다.
(\hyperref[III.5.3.3-fr]{5.3.3})에 의해 $\mathcal P$와
$\mathcal R$도 대수화 가능하다. 이제
(\hyperref[III.5.3.2-fr]{5.3.2})를 적용하면 마침내
$\mathcal F$가 대수화 가능함이 증명된다.

\hfill C.Q.F.D.
\end{env}

\subsection{응용: 통상적 스킴의 사상과 형식적 스킴의 사상의 비교.
대수화 가능한 형식적 스킴}
\label{subsection:III.5.4-fr}

\begin{theorem}[5.4.1]
\label{III.5.4.1-fr}
$A$를 뇌터 아딕 환, $\mathfrak J$를 $A$의 정의 아이디얼이라 하고
\[
S=\operatorname{Spec}(A),\qquad S'=\mathrm V(\mathfrak J)
\]
로 놓자. $u:X\to S$를 고유 사상, $v:Y\to S$를 분리 유한형
사상이라 하자. $\widehat S$, $\widehat X$, $\widehat Y$를 각각
$S'$, $u^{-1}(S')$, $v^{-1}(S')$을 따른 $S$, $X$, $Y$의
완비화라 하자. 모든 $S$-사상 $f:X\to Y$에 대해
$\widehat f:\widehat X\to\widehat Y$가 $f$의 완비화들로의
연장이면, 사상 $f\mapsto\widehat f$는 전단사
\[
\operatorname{Hom}_S(X,Y)\xrightarrow{\ \sim\ }
\operatorname{Hom}_{\widehat S}(\widehat X,\widehat Y)
\]
이다.
\end{theorem}

\begin{proof}
먼저 $f\mapsto\widehat f$가 단사임을 보이자. 실제로 두
$S$-사상 $f,g:X\to Y$가 $\widehat f=\widehat g$를 만족한다고
하자. 그러면 $f$와 $g$가 일치하는 $X'=u^{-1}(S')$의 어떤 열린
근방 $V$가 존재한다 (I, 10.9.4). 그런데 $u$는 닫힌 사상이므로
(\hyperref[III.5.1.3.1-fr]{5.1.3.1})에 의해 $V=X$이고, 따라서
$f=g$이다.

이제 $f\mapsto\widehat f$가 전사임을 증명하자. 이를 위해
$h:\widehat X\to\widehat Y$를 $\widehat S$-사상이라 하자.
$Z=X\times_S Y$로 놓고 $p:Z\to X$, $q:Z\to Y$를 표준 사영이라
하자. $Z$는 $S$ 위에서 유한형이므로 (I, 6.3.4) 뇌터이다.
$\widehat Z$를 $Z'=p^{-1}(u^{-1}(S'))$을 따른 $Z$의 완비화라
하자. $\widehat Z$는 표준적으로
$\widehat X\times_{\widehat S}\widehat Y$와 동일시되고,
$\widehat Z\to\widehat X$, $\widehat Z\to\widehat Y$의 사영들은
$\widehat p$, $\widehat q$와 동일시된다 (I, 10.9.7).
$Y$가 $S$ 위에서 분리이므로 $\widehat Y$는 $\widehat S$ 위에서
분리이다 (I, 10.15.7). 따라서 그래프 사상
\[
\Gamma_h=(1_{\widehat X},h):\widehat X\to\widehat Z
\]
은 닫힌 몰입이다 (I, 10.15.4). 이 몰입에 대응하는
$\widehat Z$의 닫힌 형식적 부분준스킴을 $\mathfrak T$라 하고,
$j:\mathfrak T\to\widehat Z$를 표준 단사라 하자. 그러면
\[
\Gamma_h=j\circ w,
\]
여기서 $w:\widehat X\to\mathfrak T$는 동형사상이고
(I, 10.14.3), 그 역동형사상은 $\widehat p\circ j$이다. 또한
$\widehat X$가 $\widehat S$ 위에서 고유이므로 $\mathfrak T$도
분명히 $\widehat S$ 위에서 고유이다. 따라서
(\hyperref[III.5.1.8-fr]{5.1.8})에 의해 $Z$의 닫힌 부분준스킴
$T$가 존재하여
\[
\mathfrak T=\widehat T=T_{/(T\cap Z')},
\]
이고 $j=\widehat i$이다. 여기서 $i:T\to Z$는 표준 단사이다
(I, 10.14.7). 사상 $p\circ i:T\to X$는 동형사상이다. 실제로
가정에 의해 그 완비화
\[
\widehat{p\circ i}=\widehat p\circ\widehat i
\]
가 동형사상이므로, 위에서와 같이 $S$가 $S'$의 $S$ 안에서의
유일한 근방임을 주의하고

\oldpage[III]{157}

(4.6.8)을 적용하면 된다. $g:X\to T$를 $p\circ i$의
역동형사상이라 하고 $f=q\circ i\circ g$로 놓자. 이는
$X\to Y$인 사상이며, 정의에 의해 그 그래프는
$\Gamma_f=i\circ g$이다. $\widehat g$는
$\widehat{p\circ i}=w$의 역동형사상이므로
\[
\widehat{\Gamma_f}=\widehat i\circ\widehat g
=j\circ w=\Gamma_h.
\]
한편 $\widehat{\Gamma_f}=\Gamma_{\widehat f}$임을 알고 있으므로
(I, 10.9.8), 마침내 $h=\widehat f$이다. 이로써 증명이 끝난다.
\end{proof}

따라서 범주의 언어로 말하면, 모든 뇌터 아딕 환 $A$에 대해
함자 $X\mapsto\widehat X$는 $\operatorname{Spec}(A)$ 위의 고유
스킴들의 범주에서 $\operatorname{Spf}(A)$ 위의 고유 형식적
스킴들의 범주로 가는 충만하고 충실한 함자이다 ($0$, 8.1.6).
그러므로 이 함자는 전자의 범주와 후자의 어떤 부분범주 사이의
동치를 정한다. 이 부분범주의 대상들을 
\emph{대수화 가능한 형식적 스킴}이라 하겠다. 그러한 형식적 스킴
$\mathfrak X$에 대해, $\mathfrak X$가 $\widehat X$와 동형이 되게
하는 $\operatorname{Spec}(A)$ 위의 고유한 통상적 스킴 $X$가
유일한 동형사상까지 존재한다.

\begin{env}[5.4.2]
\label{III.5.4.2-fr}
\emph{보충.}
(\hyperref[III.5.4.1-fr]{5.4.1})의 표기를 사용하여
\[
S_n=\operatorname{Spec}(A/\mathfrak J^{n+1}),\qquad
X_n=X\times_S S_n,\qquad Y_n=Y\times_S S_n
\]
로 놓자. (\hyperref[III.5.4.1-fr]{5.4.1})과 (I, 10.12.3)에
의하면, $S$-사상 $f:X\to Y$를 주는 것은 $S_n$-사상
$f_n:X_n\to Y_n$들로 이루어진 아딕 $(S_n)$-귀납계를 주는 것과
동치이다 (I, 10.12.2).
\end{env}

\begin{remark}[5.4.3]
\label{III.5.4.3-fr}
존재 정리 (\hyperref[III.5.1.6-fr]{5.1.6})가 암시하는 듯한 것과
달리, $\operatorname{Spf}(A)$ 위에서 고유이지만 대수화 가능하지
않은 형식적 스킴들이 존재한다. 이는 복소 대수다양체에서 나오지
않는 콤팩트 해석공간이 존재하는 것과 마찬가지이다. 뒤에서 이러한
스킴들을 ``모듈라이 이론''에서 만나게 될 것이다. 이 이론은
(기초체가 $\mathbb C$일 때) 완비 대수다양체의 복소 구조의 무한소
변형을 바로 다루며, 그러한 변형이 대수적이지 않은 해석다양체를
낳을 수 있음이 알려져 있다.
\end{remark}

\begin{proposition}[5.4.4]
\label{III.5.4.4-fr}
$A$를 뇌터 아딕 환, $\mathfrak S=\operatorname{Spf}(A)$라 하고,
$g:\mathfrak X\to\mathfrak S$와
$h:\mathfrak Y\to\mathfrak S$를 형식적 스킴들의 두 고유 사상,
$f:\mathfrak X\to\mathfrak Y$를 $\mathfrak S$-사상이라 하자.
$f$가 유한이고 $\mathfrak Y$가 대수화 가능하면
$\mathfrak X$도 대수화 가능하다.
\end{proposition}

\begin{proof}
$g$와 $h$에 대한 가정만으로도 이미 $f$는 고유임을 주의하자
(3.4.1). 또한 $f$가 유한이기 위해서는 모든
$y\in\mathfrak Y$에 대해 올 $f^{-1}(y)$이 유한이면 충분하다
(\hyperref[III.4.8.11-fr]{4.8.11}). 가정에 의해
\[
\mathcal B=f_*(\mathcal O_{\mathfrak X})
\]
는 연접 $\mathcal O_{\mathfrak Y}$-대수이다
(\hyperref[III.4.8.6-fr]{4.8.6}). 이제
$\mathfrak Y=\widehat Y$이고 $h=\widehat w$이며
$w:Y\to\operatorname{Spec}(A)$가 통상적 스킴들의 고유 사상이라
하면, 존재 정리에 의해
\[
\mathcal B=\widehat{\mathcal C}
\]
를 만족하는 연접 $\mathcal O_Y$-대수 $\mathcal C$가 존재한다.
$X=\operatorname{Spec}(\mathcal C)$라 하고 $u:X\to Y$를 구조
사상이라 하자. 그러면 $\mathcal B$로부터 $\mathfrak X$를
정의하는 방식 (\hyperref[III.4.8.7-fr]{4.8.7})에서 곧바로
$\mathfrak X$가 표준적으로 $\widehat X$와 동형이고 $f$가
$\widehat u$와 동일시됨을 알 수 있다. 이를 확인하려면 $Y$가
아핀인 경우를 생각하면 충분하다.

(\hyperref[III.5.1.8-fr]{5.1.8})은
(\hyperref[III.5.4.4-fr]{5.4.4})의 특수한 경우임을 주의하자.
\end{proof}

\begin{theorem}[5.4.5]
\label{III.5.4.5-fr}
$A$를 뇌터 아딕 환, $\mathfrak J$를 $A$의 정의 아이디얼이라 하고,
\[
S=\operatorname{Spec}(A),\qquad
\mathfrak S=\widehat S=\operatorname{Spf}(A)
\]
로 놓자. $f:\mathfrak X\to\mathfrak S$를 형식적 스킴들의 고유
사상이라 하자. 다음과 같이 놓는다.
\[
S_k=\operatorname{Spec}(A/\mathfrak J^{k+1}),\qquad
X_k=\mathfrak X\times_{\mathfrak S}S_k,
\]
그리고 모든 $\mathcal O_{\mathfrak X}$-가군 $\mathcal F$에 대해
\[
\mathcal F_k
=\mathcal F\otimes_{\mathcal O_{\mathfrak X}}\mathcal O_{X_k}
=\mathcal F/\mathfrak J^{k+1}\mathcal F.
\]

\oldpage[III]{158}

$\mathcal L$을 가역 $\mathcal O_{\mathfrak X}$-가군이라 하고,
$\mathcal L_0=\mathcal L/\mathfrak J\mathcal L$이 풍부한
$\mathcal O_{X_0}$-가군이라고 가정하자. 그러면 $\mathfrak X$는
대수화 가능하다. 또한 $X$가 $\mathfrak X\simeq\widehat X$를
만족하는 고유 $S$-스킴이면, $\mathcal L\simeq\widehat{\mathcal M}$을
만족하는 풍부한 $\mathcal O_X$-가군 $\mathcal M$이 존재한다.
따라서 $X$는 $S$ 위에서 사영이다.
\end{theorem}

\begin{proof}
(\hyperref[III.5.2.3-fr]{5.2.3})을
$\mathcal F=\mathcal O_{\mathfrak X}$에 적용하자. 그러면 어떤 정수
$n_0$가 존재하여 모든 $n\geq n_0$에 대해 표준 준동형사상
\[
\Gamma(\mathfrak X,\mathcal L^{\otimes n})
\longrightarrow\Gamma(X_0,\mathcal L_0^{\otimes n})
\]
은 전사이다. $\mathcal L_0^{\otimes n}$이 $S_0$에 대해 매우
풍부하도록 $n\geq n_0$를 충분히 크게 택할 수 있다
(II, 4.5.10). 사상 $f_0:X_0\to S_0$가 고유이므로
$\Gamma(X_0,\mathcal L_0^{\otimes n})$은 유한형 $A$-가군이다
(3.2.1). 따라서
$\Gamma(\mathfrak X,\mathcal L^{\otimes n})$의 어떤 유한형
부분 $A$-가군 $E$가 존재하여, 그
$\Gamma(X_0,\mathcal L_0^{\otimes n})$ 안에서의 상은 이 가군
전체이다.

이제 모든 $k\geq0$에 대해 표준 단사
$E\to\Gamma(\mathfrak X,\mathcal L^{\otimes n})$에서 유도되는
준동형사상
\[
u_k:E/\mathfrak J^{k+1}E\longrightarrow
\Gamma(X_k,\mathcal L_k^{\otimes n})
\]
을 생각하자. $(f_k)_*(\mathcal L_k^{\otimes n})$은 준연접이고
\[
\Gamma(S_k,(f_k)_*(\mathcal L_k^{\otimes n}))
=\Gamma(X_k,\mathcal L_k^{\otimes n})
\]
이므로, $u_k$는 준동형사상
\[
\widetilde u_k:(E/\mathfrak J^{k+1}E)^\sim
\longrightarrow(f_k)_*(\mathcal L_k^{\otimes n})
\]
를 정하고, 따라서 준동형사상
\[
\widetilde u_k^\sharp:
f_k^*((E/\mathfrak J^{k+1}E)^\sim)
\longrightarrow\mathcal L_k^{\otimes n}
\]
도 정한다. 한편
\[
\mathcal G_k=f_k^*((E/\mathfrak J^{k+1}E)^\sim)
\]
로 놓으면 $\mathcal G_0=\mathcal G_k/\mathfrak J\mathcal G_k$이다
(I, 9.1.5). 따라서
\[
\widetilde u_0^\sharp:
\mathcal G_k/\mathfrak J\mathcal G_k
\longrightarrow
\mathcal L_k^{\otimes n}/\mathfrak J\mathcal L_k^{\otimes n}
\]
은 $\widetilde u_k^\sharp$에서 몫으로 넘어가서 얻어진다. 그런데
$E$의 정의에 의해 $\widetilde u_0^\sharp$는 바로 표준 준동형사상
\[
\sigma:f_0^*((f_0)_*(\mathcal L_0^{\otimes n}))
\longrightarrow\mathcal L_0^{\otimes n}
\]
이다. $\mathcal L_0^{\otimes n}$이 매우 풍부하다는 가정에 의해
$\widetilde u_0^\sharp$는 전사이다 (II, 4.4.3). 그러므로 나카야마
보조정리에 의해 모든 $\widetilde u_k^\sharp$도 전사이다.

따라서 각각의 $\widetilde u_k^\sharp$는 $S_k$-사상
\[
g_k:X_k\longrightarrow P_k
=\mathbf P(E/\mathfrak J^{k+1}E)
\]
를 정한다 (II, 4.2.2). (II, 4.1.3)에 의해 $h\leq k$이면
$P_h=P_k\times_{S_k}S_h$이다. 그러므로
$\widetilde u_k^\sharp$들 사이의 관계와 (II, 4.2.10)에 의해
$(g_k)$는 아딕 $(S_k)$-귀납계이다 (I, 10.12.2). 따라서 $g_k$들은
형식적 스킴들의 $\mathfrak S$-사상
\[
g:\mathfrak X\to\mathfrak P
\]
를 정한다. 여기서 $\mathfrak P$는 계 $(P_k)$의 귀납극한이며,
동일하게 $P=\mathbf P(E)$의 완비화 $\widehat P$이다.
$\mathcal L_0^{\otimes n}$이 매우 풍부하다는 가정에 의해 $g_0$는
닫힌 몰입이다 (II, 4.4.3). 따라서 $g$는 형식적 스킴들의 닫힌
몰입이고 (4.8.10), 그러므로 $\mathfrak X$는 대수화 가능하다
(5.1.8).

이제 존재 정리 (5.1.6)에 의해 $\mathcal L$은 어떤 가역
$\mathcal O_X$-가군 $\mathcal M$의 완비화
$\widehat{\mathcal M}$과 동형이다. 또한
$\mathcal L^{\otimes n}$은 $\mathcal M^{\otimes n}$의 완비화이다
(I, 10.8.10). 준동형사상 $\widetilde u_k^\sharp$들은 잘 정의된
준동형사상
\[
v:f^*(\widetilde E)\longrightarrow\mathcal M^{\otimes n}
\]
을 정한다 (5.1.7). $\widetilde u_k^\sharp$가 전사이므로
$\widehat v$도 전사이고 (I, 10.11.5), 따라서 $v$도 전사이다
(5.1.3). $v$가 정하는 사상 $r:X\to P$의 (II, 4.2.2)
완비화들로의 연장은 $g$이다. $g$가 닫힌 몰입이므로 (5.1.8)과
(5.4.1)에 의해 $r$도 닫힌 몰입이다. 따라서
$\mathcal M^{\otimes n}$은 매우 풍부하고 (II, 4.4.6),
$\mathcal M$은 풍부하다 (II, 4.5.10).
\end{proof}

\begin{remark}[5.4.6]
\label{III.5.4.6-fr}
$A$를 뇌터 아딕 환, $\mathfrak S=\operatorname{Spf}(A)$라 하고,
$f:\mathfrak X\to\mathfrak S$를 형식적 스킴들의 고유 사상이라
하자. 모든 아핀 형식적 열린집합 $U$에 대해
$\Gamma(U,\mathcal N)$이
$\Gamma(U,\mathcal O_{\mathfrak X})$의 멱영근기가 되게 하는
$\mathcal O_{\mathfrak X}$의 연접 아이디얼을 $\mathcal N$이라
하자. 이 아이디얼의 존재는 (I, 10.10.2)와 모든 환 준동형사상
$B\to C$가 $B$의 멱영근기를 $C$의 멱영근기 안으로 보낸다는
사실에서 쉽게 따른다. $\mathfrak X'$을 $\mathcal N$으로 정의되는
$\mathfrak X$의 닫힌 형식적 부분스킴이라 하자 (I, 10.14.2).
$\mathfrak X'$이 대수화 가능할 때 $\mathfrak X$ 자체도 대수화
가능한지 아는 것은 흥미로운 문제이다. 구조층
$\mathcal O_{\mathfrak X}$을 제곱이 영인 아이디얼로 확장한 것들,
즉 $\mathcal J$가 $\mathcal A$의 제곱이 영인 아이디얼이고
\[
\mathcal O_{\mathfrak X}\simeq\mathcal A/\mathcal J
\]
인 $\mathcal O_{\mathfrak X}$-대수 $\mathcal A$들을, 예를 들어

\oldpage[III]{159}

코호몰로지적 성격의 불변량으로 분류할 수 있다면 이 문제의 해답에
도달할 수 있을 것이다. 여기서 $\mathfrak X$는 통상적 준스킴일
수도 있고 형식적 준스킴일 수도 있다.
\end{remark}

\subsection{어떤 스킴들의 분해}
\label{subsection:III.5.5-fr}

\begin{proposition}[5.5.1]
\label{III.5.5.1-fr}
$A$를 뇌터 아딕 환, $\mathfrak J$를 $A$의 정의 아이디얼이라 하고
$Y=\operatorname{Spec}(A)$로 놓자. $f:X\to Y$를 분리 유한형
사상이라 하고
\[
Y_0=\operatorname{Spec}(A/\mathfrak J),\qquad
X_0=X\times_Y Y_0=f^{-1}(Y_0)
\]
로 놓자. $Z_0$를 $Y_0$ 위에서 고유인 $X_0$의 열린 부분이라
하자. 그러면 $X$ 안에 $Y$ 위에서 고유이고
$Z\cap X_0=Z_0$를 만족하는 열린닫힌 부분 $Z$가 존재한다.
\end{proposition}

\begin{proof}
가정에 의해 $T\cap X_0=Z_0$인 $X$의 열린집합 $T$가 존재한다.
$\widehat T$를 열린집합 $T$ 위에 $X$가 유도하는 스킴의 $Z_0$을
따른 완비화라 하자. $\mathcal O_{\widehat T}$의 받침은 $Y_0$
위에서 고유인 $Z_0$이므로, $\widehat T$는
$\widehat Y=\operatorname{Spf}(A)$ 위에서 고유이다 (3.4.1).
(5.1.8)에 의해 $Y$ 위에서 고유인 $T$의 닫힌 부분스킴 $Z$가
존재하여, 표준 단사 $i:Z\to T$의 완비화
\[
\widehat i:\widehat Z\to\widehat T
\]
는 동형사상이다. 여기서 $\widehat Z$는 $Z_0$을 따른 $Z$의
완비화이다. 따라서 (4.6.8)에 의해 $Z_0$의 $T$ 안에서의 어떤
열린 근방 $V$가 존재하여 $i$의 제한
$i^{-1}(V)\to V$는 동형사상이다. 그런데 $i^{-1}(V)$는 $Z$ 안에서
$Z_0$의 근방이므로 반드시 $Z$와 같다 (5.1.3.1). 따라서 $Z$는
$T$에서 열려 있고, 그러므로 $X$에서도 열려 있다. 이로써 증명이
끝난다.
\end{proof}

\begin{corollary}[5.5.2]
\label{III.5.5.2-fr}
$X_0$가 $Y_0$ 위에서 고유이면 $X$는 서로소인 두 열린 부분
$Z$와 $Z'$의 합집합이다. 여기서 $Z$는 $Y$ 위에서 고유이고
$X_0$를 포함한다. 또한 $Y$ 위에서 고유인 $X$의 모든 닫힌 부분
$P$는 $Z$에 포함된다.
\end{corollary}

\begin{proof}
마지막 주장은 다음 사실에서 따른다. $P\cap Z'$은 $P$에서 닫혀
있으므로 $Y$ 위에서 고유이다. 만일 $P\cap Z'$이 공집합이 아니면
$f(P\cap Z')$은 $Y$의 공집합이 아닌 닫힌집합이므로 $Y_0$와
만나야 한다 (5.1.3.1). 이는 $Z$의 정의에 모순된다.
\end{proof}

\begin{flushright}
\emph{(계속.)}
\end{flushright}

% 인쇄본 504쪽은 백지이다.
\clearpage
\thispagestyle{empty}
\null
\clearpage
\thispagestyle{plain}

\oldpage[III]{161}

\section*{참고문헌 (계속)}

\begin{flushleft}
[27] P. Gabriel, \emph{Des catégories abéliennes}, thèse, Paris (1961).

\medskip
[28] J. E. Roos, \emph{Sur les foncteurs dérivés de $\varprojlim$.
Applications}, C. R. Acad. Sci. Paris, t. CCLII, 1961, p. 3702--3704.

\medskip
[29] H. Cartan, \emph{Séminaire de l'École Normale Supérieure},
13\textsuperscript{e} année (1960--1961), exposé n\textsuperscript{o} 11.
\end{flushleft}

\clearpage
\oldpage[III]{162}
\section*{기호 색인}

\begin{flushleft}
$\mathbf{Ens}$ : 0, 8.1.1.

$h_X(Y)$, $h_X(u)$ ($X$, $Y$는 범주 $C$의 대상, $u$는 $C$의
사상) : 0, 8.1.1.

$h_w(Y)$ ($Y$는 $C$의 대상, $w$는 $C$의 사상) : 0, 8.1.2.

$Z_k(E_r^{pq})$, $B_k(E_r^{pq})$ ($k\geq r$) : 0, 11.1.1.

$Z_\infty(E_2^{pq})$, $B_\infty(E_2^{pq})$ : 0, 11.1.1.

$K^\bullet$, $K_\bullet$ (복합체) : 0, 11.2.1.

$K^{\bullet\bullet}$, $K_{\bullet\bullet}$ (이중복합체) : 0, 11.2.1.

$B^i(K^\bullet)$, $Z^i(K^\bullet)$, $H^i(K^\bullet)$,
$H^\bullet(K^\bullet)$ : 0, 11.2.1.

$B_i(K_\bullet)$, $Z_i(K_\bullet)$, $H_i(K_\bullet)$,
$H_\bullet(K_\bullet)$ : 0, 11.2.1.

$T(K^\bullet)$, $T(K_\bullet)$ ($T$는 함자) : 0, 11.2.1.

$E(K^\bullet)$ ($K^\bullet$은 여과 복합체) : 0, 11.2.2.

$K^{i,\bullet}$, $K^{\bullet,j}$, $Z_{\mathrm{II}}^p(K^{i,\bullet})$,
$B_{\mathrm{II}}^p(K^{i,\bullet})$, $H_{\mathrm{II}}^p(K^{i,\bullet})$,
$Z_{\mathrm I}^p(K^{\bullet,j})$, $B_{\mathrm I}^p(K^{\bullet,j})$,
$H_{\mathrm I}^p(K^{\bullet,j})$ : 0, 11.3.1.

$H_{\mathrm{II}}^p(K^{\bullet\bullet})$,
$Z_{\mathrm I}^q(H_{\mathrm{II}}^p(K^{\bullet\bullet}))$,
$B_{\mathrm I}^q(H_{\mathrm{II}}^p(K^{\bullet\bullet}))$,
$H_{\mathrm I}^q(H_{\mathrm{II}}^p(K^{\bullet\bullet}))$,
$H^n(K^{\bullet\bullet})$ : 0, 11.3.1.

$F_{\mathrm I}^p(K^{\bullet\bullet})$,
$F_{\mathrm{II}}^p(K^{\bullet\bullet})$, ${}'E(K^{\bullet\bullet})$,
${}''E(K^{\bullet\bullet})$ : 0, 11.3.2.

$K_{i,\bullet}$, $K_{\bullet,j}$, $H_p^{\mathrm{II}}(K_{i,\bullet})$,
$H_p^{\mathrm I}(K_{\bullet,j})$, $H_p^{\mathrm{II}}(K_{\bullet\bullet})$,
$H_p^{\mathrm I}(K_{\bullet\bullet})$,
$H_q^{\mathrm I}(H_p^{\mathrm{II}}(K_{\bullet\bullet}))$,
$H_q^{\mathrm{II}}(H_p^{\mathrm I}(K_{\bullet\bullet}))$,
$H_n(K_{\bullet\bullet})$ : 0, 11.3.5.

$\mathrm R^\bullet T(K^\bullet)$ ($T$는 공변 가법 함자) : 0, 11.4.3.

$\mathrm R^\bullet T(K^\bullet,K'^{\bullet})$ ($T$는 공변 가법
쌍함자) : 0, 11.4.6.

$\mathrm L_\bullet T(K_\bullet)$ ($T$는 공변 가법 함자) : 0, 11.6.2.

$\mathrm L_\bullet T(K_\bullet,K'_\bullet)$ ($T$는 공변 가법
쌍함자) : 0, 11.6.6.

$\mathrm L_\bullet T(K_{\bullet\bullet})$ ($T$는 공변 가법 함자) :
0, 11.7.2.

$|\sigma|$, $\Sigma(A)$, $\mathrm C_\bullet(A)$, $\mathrm L_\bullet(A)$
($A$는 유한집합, $\sigma$는 $A$의 단체) : 0, 11.8.1.

$\mathrm C^\bullet(A,B;\mathscr S)$, $\mathrm L^\bullet(A,B;\mathscr S)$
($A$, $B$는 유한집합, $\mathscr S$는 계수계) : 0, 11.8.4.

$\mathrm P^\bullet(A,B,\mathscr S)$ ($A$, $B$는 유한집합,
$\mathscr S$는 계수계) : 0, 11.8.6.

$\mathrm C^\bullet(A,B;\mathscr S^\bullet)$,
$\mathrm L^\bullet(A,B;\mathscr S^\bullet)$ ($A$, $B$는 유한집합,
$\mathscr S^\bullet$은 계수계들의 복합체) : 0, 11.8.8.

$\mathrm P^\bullet(A,B;\mathscr S^\bullet)$ ($A$, $B$는 유한집합,
$\mathscr S^\bullet$은 계수계들의 복합체) : 0, 11.8.10.

$\chi(M^\bullet)$ ($M^\bullet$은 유한 길이 가군들로 이루어진 유한
복합체) : 0, 11.10.1.

$\check C^\bullet(\mathfrak U,\mathcal F)$ ($\mathcal F$는 $X$ 위의
아벨군의 층, $\mathfrak U$는 $X$의 열린 덮개) : 0, 12.1.2.

$\mathrm R^pf_*(\mathcal F)$ ($f:X\to Y$는 환 달린 공간들의 사상,
$\mathcal F$는 $\mathcal O_X$-가군의 층) : 0, 12.2.1.

$\mathcal H^p(f,\mathcal K^\bullet)$,
$\mathcal H_f^p(\mathcal K^\bullet)$, ${}'E(f,\mathcal K^\bullet)$,
${}''E(f,\mathcal K^\bullet)$ ($f:X\to Y$는 환 달린 공간들의 사상,
$\mathcal K^\bullet$은 $\mathcal O_X$-가군들의 복합체) : 0, 12.4.1.

$\mathbf H^\bullet(X,\mathcal K^\bullet)$,
$\mathbf H^\bullet(V,\mathcal K^\bullet)$ ($X$는 환 달린 공간,
$\mathcal K^\bullet$은 $\mathcal O_X$-가군들의 복합체, $V$는
$X$의 열린집합) : 0, 12.4.2.

$\operatorname{gr}^\bullet(A)$ ($A$는 조건 (ML)을 만족하는 사영계) :
0, 13.4.2.

$E(A)$ ($A$는 여과 대상) : 0, 13.6.4.

$K_\bullet(\mathbf f)$, $i_{\mathbf f}$ ($\mathbf f$는 환의 원소들로
이루어진 유한족) : III, 1.1.1.

$K^\bullet(\mathbf f,M)$, $K_\bullet(\mathbf f,M)$ ($\mathbf f$는 환
$A$의 원소들로 이루어진 유한족, $M$은 $A$-가군) : III, 1.1.2.

$H^\bullet(\mathbf f,M)$, $H_\bullet(\mathbf f,M)$ ($\mathbf f$는 환
$A$의 원소들로 이루어진 유한족, $M$은 $A$-가군) : III, 1.1.3.

$C^\bullet((\mathbf f),M)$, $H^\bullet((\mathbf f),M)$ ($\mathbf f$는
환 $A$의 원소들로 이루어진 유한족, $M$은 $A$-가군) : III, 1.1.6.

$H^\bullet(U,\mathcal F^{(*)})$,
$H^\bullet(\mathfrak U,\mathcal F^{(*)})$ ($\mathcal F$는 준연접
$\mathcal O_X$-가군, $U$는 $X$의 열린집합, $\mathfrak U$는
$X$의 열린 덮개) : III, 2.1.1.
\end{flushleft}

\clearpage
\oldpage[III]{163}

\section*{용어 색인}

\begin{flushleft}
스펙트럼 열의 수렴 대상 : 0, 11.1.1.

대수화 가능한 ($\mathcal O_X$-가군) : III, 5.2.1.

대수화 가능한 (형식적 스킴) : III, 5.4.2.

해석적으로 정역인 (환) : III, 4.3.6.

준콤팩트 사상 : 0, 9.1.1.

분해의 증대 : 0, 11.4.1.

복합체의 오일러--푸앵카레 특성 : 0, 11.10.1.

연접 $\mathcal O_X$-가군의 오일러--푸앵카레 특성
($X$는 $\operatorname{Spec}(A)$ 위에서 사영, $A$는 아르틴 환) :
III, 2.5.1.

쌍교대 공사슬 : 0, 11.8.4.

이중복합체가 정하는 복합체 : 0, 11.3.1.

외대수의 복합체 : III, 1.1.1.

조건 (ML) : 0, 13.1.2.

구성가능인 (부분, 집합) : 0, 9.1.2.

구성가능인 (함수) : 0, 9.3.1.

컵곱 : 0, 12.1.2.

범주에서의 대수적 구조에 관한 디-준동형 : 0, 8.2.1.

아벨 범주에서 완전한 (부분집합) : III, 3.1.1.

본질적으로 상수인 (사영계) : 0, 13.4.2.

스테인 인자분해 : III, 4.3.3.

공이산, 공분리, 이산, 소진적, 유한, 분리인 여과 : 0, 11.1.3.

범주의 끝 (대상) : 0, 8.1.10.

형식적 준스킴의 유한 (사상) : III, 4.8.2.

$\mathfrak Y$ 위에서 유한인 (형식적 준스킴),
$\mathfrak Y$-유한인 (형식적 준스킴) : III, 4.8.2.

표준 공변 함자 $C\to\operatorname{Hom}(C^o,\mathbf{Ens})$ : 0, 8.1.2.

산술 종수 : III, 2.5.1.

기하적으로 연결인 (올) : III, 4.3.4.

범주에서의 대수적 구조에 관한 준동형 : 0, 8.2.1.

복합체 $K^\bullet$에 관한 함자의 초코호몰로지 : 0, 11.4.3.

두 복합체 $K^\bullet$, $K'^{\bullet}$에 관한 쌍함자의
초코호몰로지 : 0, 11.4.6.

복합체 $K_\bullet$에 관한 함자의 초호몰로지 : 0, 11.6.2.

두 복합체 $K_\bullet$, $K'_\bullet$에 관한 쌍함자의 초호몰로지 :
0, 11.6.6.

이중복합체에 관한 함자의 초호몰로지 : 0, 11.7.2.

국소 구성가능인 (집합) : 0, 9.1.11.

범주의 대상 위의 외부 (내부) 합성법칙 : 0, 8.2.1.

\oldpage[III]{164}

여과 $S$-$C$-가군, 등급 $\operatorname{gr}^\bullet(S)$-$C$-가군,
이중등급 $\operatorname{gr}^{\bullet\bullet}(S)$-$C$-가군 : 0, 13.6.6.

스펙트럼 열의 사상 : 0, 11.1.2.

올의 기하적 연결 성분 수 : III, 4.3.4.

$C$의 군 대상, $C$-군, $C$-환, $C$-가군 : 0, 8.2.3.

경계 대상, 공경계 대상, 공순환 대상, 순환 대상 : 0, 11.2.1.

보편상 대상 : 0, 13.1.1.

조건 (ML)을 만족하는 사영계의 연관 등급 대상 : 0, 13.4.2.

아래로 (위로) 유계인 등급 대상 : 0, 11.2.1.

형식적 준스킴의 ($\mathfrak Y$ 위에서) 고유인 부분 : III, 3.4.1.

충만한 (부분범주) : 0, 8.1.5.

충만하고 충실한 (함자) : 0, 8.1.5.

힐베르트 다항식 : III, 2.5.3.

형식적 준스킴의 고유 (사상) : III, 3.4.1.

표현 가능한 (함자) : 0, 8.1.8.

코호몰로지, 오른쪽, 왼쪽, 호몰로지, 단사, 자유, 평탄, 사영 분해 :
0, 11.4.1.

오른쪽 카르탕--아일렌베르크 분해, 단사 카르탕--아일렌베르크 분해 :
0, 11.4.2.

왼쪽 카르탕--아일렌베르크 분해, 사영 카르탕--아일렌베르크 분해 :
0, 11.6.1.

역콤팩트인 (집합) : 0, 9.1.1.

엄밀한 (사영계) : 0, 13.4.2.

아벨 범주에서의 스펙트럼 열 : 0, 11.1.1.

약수렴, 정칙, 공정칙, 쌍정칙 스펙트럼 열 : 0, 11.1.3.

퇴화하는 스펙트럼 열 : 0, 11.1.6.

여과 대상에 관한 함자의 스펙트럼 열 : 0, 13.6.4.

이중복합체의 스펙트럼 열 : 0, 11.3.2 및 11.3.5.

복합체에 관한 함자의 초코호몰로지 스펙트럼 열 : 0, 11.4.3.

복합체에 관한 함자의 초호몰로지 스펙트럼 열 : 0, 11.6.2.

이중복합체에 관한 함자의 초호몰로지 스펙트럼 열 : 0, 11.7.2.

계수계 : 0, 11.8.4.

단일분기인(unibranche) (환, 점) : III, 4.3.6.

보편적으로 열린 (사상) : III, 4.3.9.
\end{flushleft}

\clearpage
\oldpage[III]{165}

\section*{목차}

\begin{flushleft}
\textsc{제0장.} --- \textbf{예비사항 (계속)} \dotfill 5

\medskip
\S~8. 표현 가능 함자 \dotfill 5

8.1. 표현 가능 함자 \dotfill 5

8.2. 범주에서의 대수적 구조 \dotfill 9

\medskip
\S~9. 구성가능 집합 \dotfill 12

9.1. 구성가능 집합 \dotfill 12

9.2. 뇌터 공간에서의 구성가능 집합 \dotfill 14

9.3. 구성가능 함수 \dotfill 16

\medskip
\S~10. 평탄 가군에 관한 보충 \dotfill 17

10.1. 평탄 가군과 자유 가군 사이의 관계 \dotfill 17

10.2. 평탄성의 국소 판정법 \dotfill 18

10.3. 국소환의 평탄 확대의 존재 \dotfill 20

\medskip
\S~11. 호몰로지 대수학에 관한 보충 \dotfill 23

11.1. 스펙트럼 열에 관한 복습 \dotfill 23

11.2. 여과 복합체의 스펙트럼 열 \dotfill 27

11.3. 이중복합체의 스펙트럼 열 \dotfill 29

11.4. 복합체 $K^\bullet$에 관한 함자의 초코호몰로지
\dotfill 32

11.5. 초코호몰로지에서 귀납극한으로의 이행 \dotfill 35

11.6. 복합체 $K_\bullet$에 관한 함자의 초호몰로지
\dotfill 39

11.7. 이중복합체 $K_{\bullet\bullet}$에 관한 함자의 초호몰로지
\dotfill 41

11.8. 단체 복합체의 코호몰로지에 관한 보충 \dotfill 43

11.9. 유한 생성 복합체에 관한 보조정리 \dotfill 46

11.10. 유한 길이 가군 복합체의 오일러--푸앵카레 특성
\dotfill 48

\medskip
\S~12. 층의 코호몰로지에 관한 보충 \dotfill 49

12.1. 환 달린 공간 위의 가군층의 코호몰로지 \dotfill 49

12.2. 고차 직상 \dotfill 57

12.3. 층의 Ext 함자에 관한 보충 \dotfill 60

12.4. 직상 함자의 초코호몰로지 \dotfill 62

\oldpage[III]{166}

\S~13. 호몰로지 대수학에서의 사영극한 \dotfill 64

13.1. 미타그--레플러 조건 \dotfill 64

13.2. 아벨군에 대한 미타그--레플러 조건 \dotfill 65

13.3. 응용: 층들의 사영극한의 코호몰로지 \dotfill 68

13.4. 미타그--레플러 조건과 사영계의 연관 등급 대상 \dotfill 69

13.5. 여과 복합체의 스펙트럼 열들의 사영극한 \dotfill 71

13.6. 유한 여과를 갖춘 대상에 관한 함자의 스펙트럼 열 \dotfill 73

13.7. 여러 변수의 사영극한의 유도 함자 \dotfill 75

\medskip
\textsc{제III장.} --- \textbf{연접층의 코호몰로지적 연구} \dotfill 81

\medskip
\S~1. 아핀 스킴의 코호몰로지 \dotfill 82

1.1. 외대수 복합체에 관한 복습 \dotfill 82

1.2. 열린 덮개의 체흐 코호몰로지 \dotfill 85

1.3. 아핀 스킴의 코호몰로지 \dotfill 88

1.4. 임의의 준스킴의 코호몰로지에 대한 응용 \dotfill 89

\medskip
\S~2. 사영 사상의 코호몰로지적 연구 \dotfill 95

2.1. 몇몇 코호몰로지 군의 명시적 계산 \dotfill 95

2.2. 사영 사상의 기본정리 \dotfill 100

2.3. 등급 대수 및 가군의 층에 대한 응용 \dotfill 102

2.4. 기본정리의 일반화 \dotfill 107

2.5. 오일러--푸앵카레 특성과 힐베르트 다항식 \dotfill 109

2.6. 응용: 풍부성 판정법 \dotfill 111

\medskip
\S~3. 고유 사상의 유한성 정리 \dotfill 115

3.1. 데비사주 보조정리 \dotfill 115

3.2. 유한성 정리: 통상적 스킴의 경우 \dotfill 116

3.3. 유한성 정리의 일반화 (통상적 스킴) \dotfill 118

3.4. 유한성 정리: 형식적 스킴의 경우 \dotfill 119

\medskip
\S~4. 고유 사상의 기본 정리와 그 응용 \dotfill 122

4.1. 기본 정리 \dotfill 122

4.2. 특수한 경우들과 변형들 \dotfill 129

4.3. 자리스키의 연결성 정리 \dotfill 130

4.4. 자리스키의 ``주정리'' \dotfill 135

4.5. 준동형사상 가군의 완비화 \dotfill 138

4.6. 형식 사상과 보통 사상의 관계 \dotfill 139

4.7. 풍부성 판정법 \dotfill 145

4.8. 형식적 준스킴의 유한 사상 \dotfill 146

\oldpage[III]{167}

\S~5. 연접 대수적 층의 존재 정리 \dotfill 149

5.1. 정리의 진술 \dotfill 149

5.2. 존재 정리의 증명: 사영인 경우와 준사영인 경우
\dotfill 151

5.3. 존재 정리의 증명: 일반적인 경우 \dotfill 154

5.4. 응용: 통상적 스킴의 사상과 형식적 스킴의 사상의 비교.
대수화 가능한 형식적 스킴 \dotfill 156

5.5. 어떤 스킴들의 분해 \dotfill 159

\medskip
\textsc{참고문헌 (계속)} \dotfill 161

\textsc{기호 색인} \dotfill 162

\textsc{용어 색인} \dotfill 163

\begin{flushright}
\emph{1961년 6월 28일 접수.}
\end{flushright}
\end{flushleft}
% ===== END EXACT INLINE INPUT: c3-1.tex =====

\end{document}
